\documentclass[11pt,a4paper,oneside]{report}
\usepackage[margin=20mm,headheight=15pt]{geometry}
\usepackage[no-math]{fontspec}
\setmainfont[Script=Arabic]{Amiri}
\newfontfamily\latinfont{Latin Modern Roman}
\usepackage{amsmath,amssymb,amsfonts,amsthm,mathrsfs,mathtools,mathdots,nicefrac,stmaryrd}
\usepackage{xparse,etoolbox,xpunctuate,mfirstuc}
\usepackage{capt-of}
\usepackage{graphicx,tikz,xcolor,framed,comment,longtable,multirow,array,booktabs,verbatim,rotating}
\usetikzlibrary{arrows,automata,positioning,calc,decorations.pathmorphing,intersections}
\usepackage{bussproofs}
\input{../upstream/sty/bussproofs-extra.sty}
\usepackage[tableaux]{prooftrees}
\usepackage[authoryear]{natbib}
\usepackage[unicode,colorlinks,linkcolor=blue,urlcolor=blue,citecolor=blue]{hyperref}
\usepackage[capitalize,nameinlink]{cleveref}
\usepackage{fancyhdr}
\usepackage{bidi}
\makeatletter
\renewcommand\tagform@[1]{\maketag@@@{\LR{\latinfont(\ignorespaces#1\unskip\@@italiccorr)}}}
\makeatother

% Exact additional-symbol declarations from frozen sty/open-logic.sty:228--231.
\DeclareSymbolFont{symbolsC}{U}{ntxsyc}{m}{n}
\DeclareMathSymbol{\boxright}{\mathbin}{symbolsC}{128}
\DeclareMathSymbol{\diamondright}{\mathbin}{symbolsC}{132}
\DeclareMathSymbol{\fishhookright}{\mathbin}{symbolsC}{74}

\hypersetup{
  pdftitle={خلاص منطق — پښتو (پاکستان)},
  pdfauthor={د اوپن لاجک پروژې ونډه‌وال؛ په مصنوعي ځيرکتيا پښتو ژباړه},
  pdfsubject={د منطق درسي متن: رسمي ژبې، ثبوتونه، مدلونه، محاسبه او د سټونو نظريه}
}

\newcommand*{\DeclareDocumentMacro}[2]{\def#1{#2}}
\NewDocumentCommand{\settexttoken}{m s m m o o}{}
\makeatletter
\input{../upstream/sty/open-logic-formulas.sty}
\input{../upstream/sty/open-logic-selective.sty}
\input{../upstream/open-logic-config.sty}
\makeatother

\renewcommand\TAss{\RL{فرض}}

\RenewDocumentCommand \indcase { s t{!} m m +m }{%
  \DeclareDocumentMacro \indfrm {#3}%
  \DeclareDocumentMacro \indfrmp {#3}%
  \DeclareDocumentMacro \indcomplex {#4}%
  \IfBooleanTF{#1}
    {\LR{$#3$} اټومي دے۔ }
    {\LR{$#3 \ident #4$}: }%
  \IfBooleanTF{#2}{تمرين۔}{#5}}


% Localized reader names for the frozen stage-principle macros; identifiers unchanged.
\renewcommand\stageshier{\emph{پړاوونه بنسټ دي}}
\renewcommand\stagesord{\emph{پړاوونه مرتب دي}}
\renewcommand\stagesacc{\emph{پړاوونه تجمع کوي}}
\renewcommand\stagessucc{\emph{پړاوونه دوام کوي}}
\renewcommand\stagesinf{\emph{پړاوونه نامتناهيت ته رسېږي}}
% Frozen source sphere-drawing definitions; package metadata omitted for inline use.
% upstream/sty/ptolemaicastronomy.sty SHA-256 450f79a006fbf105046f366f31d4243d6b8a4903783d8ae05717c4106c9aad85

\tikzset{
  sphere/.style = {dotted},
  sphere intersection/.style = {fill=lightgray},
  sphere layer/.style = {fill=lightgray},
  proposition/.style={smooth,tension=1.7},
}

\pgfkeyssetvalue{/tikz/layerwidth}{.5}
\pgfkeyssetvalue{/tikz/innerfactor}{.4}

% \sphereplot{n} gives the plot codes for the nth sphere

\newcommand{\sphereplot}[1]{
  circle
    [radius=(#1)*\pgfkeysvalueof{/tikz/layerwidth}-
      \pgfkeysvalueof{/tikz/layerwidth}*\pgfkeysvalueof{/tikz/innerfactor}]
}


% \spheresystem[options]{n} draws a sphere system centered at
% the origin with n number of layers

\newcommand{\spheresystem}[2][]{
  \foreach \i in {1,...,#2}{
    \draw[sphere,#1] \sphereplot{\i} ;
  }
}

% \spherelayer[options]{n} shades the n-th layer

\newcommand{\spherelayer}[2][]{
  \begin{scope}[even odd rule]
    \fill[#1,sphere layer]
    \sphereplot{#2-1} \sphereplot{#2} ;
  \end{scope}
}

% \spherefill[options]{n} fills the nth sphere

\newcommand{\spherefill}[2][]{
    \fill[sphere intersection,#1]
    \sphereplot{#2} ;
}

% \sphereintersect[options]{path}{n} draws the path and shades the area
% of the path in the n th layer. Options only apply to the sphere layer

\newcommand{\sphereintersect}[3][]{
  \begin{scope}[even odd rule]
    \path[clip] #3;
    \spherefill[#1]{#2}
  \end{scope}
  \draw #3;
}

% \sphereproposition[options]{direction}{n}{width}{length} draws a
% proposition intersecting the n-th layer in direction away from the
% center of the sphere system, with endpoints length away from the
% vertex at an angle of width/2. Note that options applies to \draw,
% not to plot. To change the tension on the plot, e.g., say
% \tikzset{proposition/.style={smooth,tension=2}} -- in a group to
% make it apply only to a single proposition.

\newcommand{\propositionplot}[4]{
  plot [proposition]
  coordinates {+(#1+#3/2:#4)
    +(#1:#2*\pgfkeysvalueof{/tikz/layerwidth}-
    \pgfkeysvalueof{/tikz/layerwidth}*.9
    -\pgfkeysvalueof{/tikz/layerwidth}*\pgfkeysvalueof{/tikz/innerfactor})
    +(#1-#3/2:#4)}
}

\newcommand{\proposition}[5][]{
  \draw[proposition,#1] \propositionplot {#2}{#3}{#4}{#5} ;
  }

% \spherepropositionintersect does the same as \sphereproposition but
% also shades the are of intersection with the closest sphere

\newcommand{\propositionintersect}[5][]{
  \begin{scope}
  \path[clip] \propositionplot{#2}{#3}{#4}{#5};
  \spherefill[#1]{#3}
  \end{scope}
  \draw[proposition,#1] \propositionplot{#2}{#3}{#4}{#5};
}

% \spherepos[options]{direction}{n}{code} shifts the scope to a
% position in the center of the nth layer in direction angle from the
% center -- and then puts code there.

\newcommand{\spherepos}[4][]{
  \begin{scope}[#1,shift=(#2:#3*\pgfkeysvalueof{/tikz/layerwidth}-
      \pgfkeysvalueof{/tikz/layerwidth}/2-
        \pgfkeysvalueof{/tikz/layerwidth}*\pgfkeysvalueof{/tikz/innerfactor})]
    \path #4 ;
  \end{scope}
}

\newbox\gnBoxA
\newdimen\gnCornerHgt
\setbox\gnBoxA=\hbox{$\ulcorner$}
\global\gnCornerHgt=\ht\gnBoxA
\newdimen\gnArgHgt
\def\gn #1{%
  \setbox\gnBoxA=\hbox{$#1$}%
  \gnArgHgt=\ht\gnBoxA
  \ifnum\gnArgHgt<\gnCornerHgt \gnArgHgt=0pt
  \else \advance\gnArgHgt by -\gnCornerHgt
  \fi
  \raise\gnArgHgt\hbox{$\ulcorner$}\box\gnBoxA
  \raise\gnArgHgt\hbox{$\urcorner$}}
\renewcommand{\MP}{\LR{\latinfont MP}}
\renewcommand{\QR}{\LR{\latinfont QR}}
\renewcommand{\Hyp}{\LR{\latinfont Hyp}}

\newcommand{\olfileid}[3]{\gdef\pspart{#1}\gdef\pschapter{#2}\gdef\pssection{#3}}
\NewDocumentCommand{\olpart}{o m m}{%
  \gdef\pspart{#2}\gdef\pschapter{}\gdef\pssection{}%
  \IfNoValueTF{#1}{\part{#3}}{\part[#1]{#3}}%
  \label{#2:::part}}
\NewDocumentCommand{\olchapter}{o m m m}{%
  \gdef\pspart{#2}\gdef\pschapter{#3}\gdef\pssection{}%
  \IfNoValueTF{#1}{\chapter{#4}}{\chapter[#1]{#4}}%
  \label{#2:#3::chap}}
\NewDocumentCommand{\olsection}{o m}{%
  \IfNoValueTF{#1}{\section{#2}}{\section[#1]{#2}}%
  \label{\pspart:\pschapter:\pssection:sec}}
\newcommand{\ollabel}[1]{\label{\pspart:\pschapter:\pssection:#1}}
\makeatletter
\newcommand{\olpseganchor}[2]{%
  \protected@write\@auxout{}{%
    \string\newlabel{olpseg:#1:#2}{{}{\thepage}{}{}{}}}}
\makeatother
\NewDocumentCommand{\olref}{o o o m}{\begingroup
  \IfNoValueTF{#1}{\edef\pskey{\pspart:\pschapter:\pssection:#4}}{%
    \IfNoValueTF{#2}{\edef\pskey{\pspart:\pschapter:#1:#4}}{%
      \IfNoValueTF{#3}{\edef\pskey{\pspart:#1:#2:#4}}{\edef\pskey{#1:#2:#3:#4}}}}%
  \ifcsname r@\pskey\endcsname
    \LR{\ref{\pskey}}%
  \else
    \emph{ناحل شوې بهرنۍ حواله}%
  \fi
  \endgroup}
\let\Olref\olref
\NewDocumentCommand{\oliflabeldef}{m m m}{\ifcsname r@#1\endcsname #2\else #3\fi}
\providecommand{\OLEndChapterHook}{}
\providecommand{\OLEndPartHook}{}
\newcommand{\gitissue}[1]{%
  \href{https://github.com/OpenLogicProject/OpenLogic/issues/#1}{د سرچينې مسئله \##1}}

\newlength{\olphotowidth}
\setlength{\olphotowidth}{.46\textwidth}
\NewDocumentCommand{\olasset}{O{\olphotowidth} m}{%
  \centering\begin{LTR}\centering\resizebox{#1}{!}{\input{#2}}\end{LTR}}

% Each repeated portrait gets a unique figure label; photo-ID names its first occurrence.
\newcounter{psphoto}
\NewDocumentCommand{\olphoto}{O{\olphotowidth} m m}{%
  \par\medskip\begin{center}%
  \relax%
  \captionof{figure}{#3}%
  \stepcounter{psphoto}\label{photo-#2-\arabic{psphoto}}%
  \ifcsname psphotoseen@#2\endcsname\else
    \label{photo-#2}\expandafter\gdef\csname psphotoseen@#2\endcsname{1}%
  \fi
  \end{center}\medskip}

\renewenvironment{prooftree}
  {\begin{center}\begin{LTR}\bottomAlignProof}
  {\DisplayProof\end{LTR}\end{center}}
\newenvironment{oltableau}
  {\center\LTR\tableau{}}
  {\endtableau\endLTR\endcenter}

\newtheorem{thm}{قضيه}[chapter]
\newtheorem{lem}[thm]{لم}
\newtheorem{prop}[thm]{قضيه}
\newtheorem{cor}[thm]{نتيجه}
\theoremstyle{definition}
\newtheorem{defn}[thm]{تعريف}
\newtheorem*{axiom}{بديهي اصل}
\newtheorem{ex}[thm]{بېلګه}
\newtheorem{prob}[thm]{تمرين}
\newtheorem{rem}[thm]{يادونه}
\newtheorem{conv}[thm]{د ليکنې دود}
\crefname{part}{برخه}{برخې}
\crefname{chapter}{باب}{بابونه}
\crefname{section}{برخه}{برخې}
\crefname{thm}{قضيه}{قضيې}
\crefname{lem}{لم}{لمې}
\crefname{prop}{قضيه}{قضيې}
\crefname{cor}{نتيجه}{نتيجې}
\crefname{defn}{تعريف}{تعريفونه}
\crefname{ex}{بېلګه}{بېلګې}
\crefname{prob}{تمرين}{تمرينونه}
\crefname{rem}{يادونه}{يادونې}
\crefname{conv}{د ليکنې دود}{د ليکنې دودونه}
\newcommand{\crefpairconjunction}{ او }
\newcommand{\crefmiddleconjunction}{، }
\newcommand{\creflastconjunction}{ او }
\newcommand{\crefpairgroupconjunction}{ او }
\newcommand{\crefmiddlegroupconjunction}{، }
\newcommand{\creflastgroupconjunction}{ او }
\newcommand{\crefrangeconjunction}{ تر }
\newenvironment{defish}{\begin{framed}}{\end{framed}}
\newenvironment{editorial}{\begin{framed}\textbf{ادارتي يادونه۔}\enspace}{\end{framed}}
\newenvironment{history}{\begin{framed}\textbf{تاريخي يادونه۔}\enspace}{\end{framed}}
\newenvironment{explain}{}{}
\newenvironment{intro}{}{}
\newenvironment{reading}{\par\medskip\textbf{د نور لوست سرچينې۔}\enspace}{\par\medskip}
\newenvironment{digress}{\par\medskip}{\par\medskip}
\newenvironment{derivation}{%
  \par\smallskip\begin{LTR}\begin{tabular}{@{}rll@{}}}
  {\end{tabular}\end{LTR}\par\smallskip}

\renewcommand{\proofname}{ثبوت}
\renewcommand{\figurename}{شکل}
\renewcommand{\tablename}{جدول}
\renewcommand{\partname}{برخه}
\renewcommand{\chaptername}{باب}
\renewcommand{\contentsname}{فهرست}
\renewcommand{\bibname}{\RL{\normalfont سرچينې}}
\makeatletter
\renewcommand{\@pnumwidth}{2.8em}
\renewcommand{\@tocrmarg}{4.6em}
\makeatother
\setlength{\parindent}{0pt}
\setlength{\parskip}{5pt}
\setlength{\emergencystretch}{8em}
\hfuzz=5pt
\linespread{1.12}
\pagestyle{fancy}
\fancyhf{}
\fancyhead[RE,LO]{\small خلاص منطق — پښتو (پاکستان)}
\fancyhead[LE,RO]{\small د منطق درسي متن}
\fancyfoot[C]{\thepage}

\begin{document}
\tracinglostchars=3
\XeTeXgenerateactualtext=1
\setRTL
% XeLaTeX resets this mathcode at begin-document; restore OLP's active ! metavariables.
\mathcode`\!="8000
\begin{titlepage}
\thispagestyle{empty}
\begin{center}
\vspace*{22mm}
{\Huge خلاص منطق}\par\bigskip
{\Large پښتو — پاکستان}\par\bigskip
{\LARGE د منطق درسي متن}\par\bigskip
د اوپن لاجک پروژې د ونډه‌والو اثر\\
په مصنوعي ځيرکتيا پښتو ته ژباړل شوے
\end{center}
\bigskip
\small
دا درسي متن له سټونو، اړيکو او تابعو پيلېږي۔ قضيوي منطق،
د لومړۍ او دويمې درجې منطق، رسمي ثبوتونه، بشپړتيا، مدلونه،
محاسبه کېدنه او ناتمامي څېړي۔ د لامبډا حساب، څو ارزښته،
مودالي او شهودي منطق، ثبوت تيوري، استقرا، د سټونو نظريه او
د منطق تاريخ هم پکې دي۔ تعريفونه، قضيې، بېلګې او تمرينونه
د همدغو موضوعاتو په باب دي۔ د اصلي اثر ۷۲۲ سمون وړ سرچينيز
واحدونه له کتاب سره شته؛ د بديلو بابونو او اختياري توکو
اړوند متنونه په خپلو موضوعي ځايونو کښې لوستل کېدے شي۔

د اصلي متن ليکوالان د
\href{https://openlogicproject.org/people/}{اوپن لاجک پروژې ونډه‌وال}
دي۔ دا پښتو ژباړه د
\href{https://github.com/OpenLogicProject/OpenLogic/tree/9620cc73f9c8e0ad003c514a5d3748f29611c4c0}{اصلي سرچينې د همدې نسخې}
پر بنسټ ده۔
د پښتو متن په جوړولو کښې
\LR{\latinfont OpenAI Codex — GPT-5.6 Sol, Ultra}،
\LR{\latinfont OpenAI Codex — GPT-6 Sol, Ultra} او
\LR{\latinfont OpenAI Codex — GPT-6.1 Sol, Ultra}
کارول شوي دي۔ د ژبې بنسټ پاکستانۍ پښتو ده۔

د متن جواز
\href{https://creativecommons.org/licenses/by/4.0/}{\LR{\latinfont Creative Commons Attribution 4.0 International}}
دے۔ د انځورونو جلا انتسابونه او اجازې د کتاب په پای کښې دي۔
\href{https://github.com/KokunoYumeto/OpenLogic-translations}{د ژباړو فهرست}؛
\href{https://github.com/KokunoYumeto/openlogic-ps-Arab-PK/releases}{د پخوانيو نسخو ثبت}۔
\end{titlepage}
\tableofcontents
\clearpage


\olpseganchor{OLP-0001-B001}{start}
\chapter*{د خلاص منطق پروژې په اړه}
\addcontentsline{toc}{chapter}{د خلاص منطق پروژې په اړه}
\olpseganchor{OLP-0001-B001}{end}


\olpseganchor{OLP-0001-B002}{start}
\textit{د خلاص منطق متن} د صوري مابعدمنطق او صوري طريقو د خلاصې سرچينې يو ګډ درسي کتاب دے، چې له منځنۍ کچې پيل کوي (يعنې د صوري منطق له يوه لومړني کورس وروسته)۔ که څه هم د رياضي د غيرمتخصصو لوستونکو دپاره دے (په تېره د فلسفې او کمپيوټر ساينس د زده کوونکو دپاره)، خو دقيق دے۔
\olpseganchor{OLP-0001-B002}{end}

\olpseganchor{OLP-0001-B003}{start}
د ځينو هغو موضوعاتو بيان چې اوس پکښې شامل دي، ښائي لا بشپړ نۀ وي، او ډېرې برخې لا هم پراخې بياکتنې ته اړتيا لري۔ مونږ اراده لرو چې په راتلونکي کښې متن پراخ کړو او نور موضوعات پکښې شامل کړو۔ دا اراده هم لرو چې متن ته نورې اسانتياوې ورزياتې کړو، لکه اصطلاحنامه، د نورې لوستنې لست، تاريخي يادښتونه، انځورونه، غوره وضاحتونه، هغه برخې چې د پايلو تړاو له فلسفې، کمپيوټر ساينس او رياضي سره بيانوي، او نور تمرينونه او بېلګې۔ که کومه تېروتنه ومومئ، يا وړانديز لرئ، \href{https://github.com/OpenLogicProject/OpenLogic/wiki/Contributing}{مهرباني وکړئ د پروژې ډلې ته ئې خبر ورکړئ}۔
\olpseganchor{OLP-0001-B003}{end}

\olpseganchor{OLP-0001-B004}{start}
پروژه د خلاصې سرچينې په روحيه کار کوي۔ يوازې متن په ازاده توګه نۀ وړاندې کوو، بلکې د لاټېک سرچينه هم د کريېټيو کامنز د نوم اخيستنې د جواز لاندې ورکوو۔ دا جواز هر چا ته حق ورکوي چې زمونږ کار ښکته کړي، وکاروي، بدل ئې کړي، بيا ئې ترتيب کړي، بڼې ئې واړوي او بيا ئې ووېشي، په دې شرط چې مناسب نسبت ورکړي۔ د نورو معلوماتو دپاره، مهرباني وکړئ په \href{http://openlogicproject.org/}{openlogicproject.org} د خلاص منطق پروژې وېبپاڼه وګورئ۔
\olpseganchor{OLP-0001-B004}{end}

\olpseganchor{OLP-0002-B004}{start}
\clearpage
\olpseganchor{OLP-0002-B004}{end}

\olpseganchor{OLP-0002-B005}{start}
\begin{editorial}
دا دوتنه ټول هغه مطالب رااخلي چې د خلاص منطق په پروژه کښې شامل دي۔ که دا ډول ادارتي يادښتونه ښکاري، نو دا څرګندوي چې دوتنه د دې مطالبو د وړاندې کولو طريقې ته له پام کولو پرته جوړه شوې ده۔ که دا لولئ، نو غالباً له دې پي ډي اېف څخه درس ورکول يا مطالعه کول \emph{مناسب نۀ دي}۔
\olpseganchor{OLP-0002-B005}{end}

\olpseganchor{OLP-0002-B006}{start}
د خلاص منطق پروژه ډېرې داسې طريقې وړاندې کوي چې د هغو په وسيله د درس ورکولو يا ځان زده کړې دپاره مناسب متن جوړېدل ممکن وي۔ د بېلګې په توګه، په اصلي تنظيم کښې متن ټول منطقي عملګر بنسټيز ګڼي، او په ثبوتونو کښې د ټولو عملګرو ټول حالتونه څېړي۔ خو دا ډېر غوره دي چې ځينې حالتونه د تمرينونو په توګه پرېښودل شي۔ د خلاص منطق پروژه پخپله هم روان کار دے۔ د ګډ کار او ښه والي د هڅولو دپاره، مطالب آن هغه وخت هم شاملېږي چې يوازې د مسودې په بڼه وي، تمرينونه نۀ لري، او داسې نور۔ هغه پي ډي اېف چې د يو کورس دپاره جوړېږي، دا برخې به ترې وباسي۔
\olpseganchor{OLP-0002-B006}{end}

\olpseganchor{OLP-0002-B007}{start}
د درس او مطالعې دپاره د لا مناسبو پي ډي اېفونو د موندلو دپاره، \href{http://builds.openlogicproject.org/}{د خلاص منطق پروژې په وېبپاڼه شته نمونوي کورسونه} وګورئ۔ د خپل متن د جوړولو دپاره، له \href{https://github.com/OpenLogicProject/OpenLogic/tree/master/courses/sample}{نمونوي چلوونکې دوتنې} څخه پيل کولے شئ، يا د لا پرمختللو او مفصلو بېلګو دپاره له دې پروژې څخه د جوړو شوو درسي کتابونو سرچينې کتلے شئ۔
\end{editorial}
\olpseganchor{OLP-0002-B007}{end}

\olpseganchor{OLP-0003-B004}{start}
\olpart{sfr}{د سټونو ساده نظريه}
\olpseganchor{OLP-0003-B004}{end}

\olpseganchor{OLP-0003-B005}{start}
\begin{editorial}
  د دې برخې مطالب د سټونو د بنسټيزې ساده نظريې يوه پېژندګلو ده۔ د تېم بټن د سټونو د خلاصې نظريې په شاملېدو سره، دا برخه د عددونو د نظامونو جوړول او د نامتناهي والي بحث هم رانغاړي، چې د خلاص منطق پروژې د منطقي برخو دپاره اړين نۀ دي۔
\end{editorial}
\olpseganchor{OLP-0003-B005}{end}

\olfileid{sfr}{outline}{alternate-outline}
\olpseganchor{OLP-0720-B004}{start}
\paragraph{سټونه، اړيکې او تابعې}
\olpseganchor{OLP-0720-B004}{end}

\olpseganchor{OLP-0720-B005}{start}
\begin{editorial}
  د سټونو د ساده نظريې دا بڼه د لومړنيو څلورو بابونو مطالب په کراره وړاندې کوي۔ بشپړه برخه په \verb|sets-functions-relations-complete.tex| کښې راغلې ده۔ دلته د ساده سټ نظريې د بنسټونو پوره پېژندګلو کېږي۔ که زده‌کوونکي له دې بنسټونو سره اشنا نۀ وي، ښه ده چې کورس د دغو مطالبو له بيا لوست نه پيل کړي، لږ تر لږه د سټونو، تابعو او اړيکو له بحثه۔ دا د رياضيکي نښليک هغه بنسټيزه بلدتيا ورکوي چې د منطق د نورو برخو دپاره پکار ده۔
\end{editorial}
\olpseganchor{OLP-0720-B005}{end}

\olpseganchor{OLP-0004-B004}{start}
\olchapter{sfr}{set}{سټونه}
\olpseganchor{OLP-0004-B004}{end}

\olpseganchor{OLP-0005-B004}{start}
\olfileid{sfr}{set}{bas}
\olsection{د غړو له مخې برابري}
\olpseganchor{OLP-0005-B004}{end}

\olpseganchor{OLP-0005-B005}{start}
\emph{سټ} د څيزونو يوه ټولګه ده چې د يوه څيز په توګه ورته کتل کېږي۔ هغه څيزونه چې سټ ترې جوړ وي، د سټ \emph{عناصر} يا \emph{غړي} بللے شي۔ که \LR{$x$} د يو سټ~\LR{$A$} غړے وي، \LR{$x \in A$} ليکو؛ که نۀ وي، \LR{$x \notin A$} ليکو۔ هغه سټ چې هېڅ غړي نۀ لري، \emph{تش} سټ بللے شي او په ``\LR{$\emptyset$}'' ښودلے کېږي۔
\olpseganchor{OLP-0005-B005}{end}

\olpseganchor{OLP-0005-B006}{start}
\begin{explain}
دا مهمه نۀ ده چې سټ څنګه \emph{ټاکو}، د هغۀ غړي څنګه \emph{ترتيبوو}، يا په رښتيا د هغۀ غړي \emph{څو ځله} شمارو۔ يوازې دا مهم دي چې د هغۀ غړي کوم دي۔ دا خبره په لاندې اصل کښې راټولوو۔
\end{explain}
\olpseganchor{OLP-0005-B006}{end}

\olpseganchor{OLP-0005-B007}{start}
\begin{defn}[د غړو له مخې برابري]
  که \LR{$A$} او \LR{$B$} سټونه وي، نو \LR{$A = B$} هله او يوازې هله وي چې د~\LR{$A$} هر غړے د~\LR{$B$} هم غړے وي، او برعکس۔
\end{defn}
\olpseganchor{OLP-0005-B007}{end}

\olpseganchor{OLP-0005-B008}{start}
د غړو له مخې برابري د ځينو نښو د کارولو جواز راکوي۔ په عمومي توګه، چې کله ځينې څيزونه \LR{$a_{1}$}، \dots، \LR{$a_{n}$} ولرو، نو \LR{$\{a_{1}, \dots, a_{n}\}$} \emph{هم هغه} سټ دے چې غړي ئې \LR{$a_1, \ldots, a_n$} دي۔ په ``\emph{هم هغه}'' ټينګار کوو، ځکه چې د غړو له مخې برابري راته وائي چې داسې سټ يوازې \emph{يو} کېدے شي۔ په حقيقت کښې، د غړو له مخې برابري د لاندې خبرې جواز هم راکوي:
  \[
    \{a, a, b\} = \{a, b\} = \{b,a\}.
  \]
دا هم هغه خبره څرګندوي چې د سټونو په پام کښې نيولو پر وخت، د هغو د غړو ترتيب يا دا چې څو ځله ښودل شوي دي، مهم نۀ دي۔
\olpseganchor{OLP-0005-B008}{end}

\olpseganchor{OLP-0005-B009}{start}

\begin{ex}
چې کله د څيزونو يوه ډله ولرئ، هغه په يو سټ کښې راټولولے شئ۔ د بېلګې په توګه، د رېچرډ د وروڼو او خوېندو سټ يو داسې سټ دے چې يو کس لري، او هغه د \LR{$S=\{\textrm{Ruth}\}$} په بڼه ليکلے شو۔ له \LR{$4$} څخه د کوچنيو مثبتو صحيح عددونو سټ \LR{$\{1, 2, 3\}$} دے، خو د \LR{$\{3, 2, 1\}$} يا آن د \LR{$\{1, 2, 1, 2, 3\}$} په بڼه هم ليکل کېدے شي۔ د غړو له مخې د برابرۍ له مخې، دا ټول هم يو سټ دے۔ ځکه چې د \LR{$\{1, 2, 3\}$} هر غړے د \LR{$\{3, 2, 1\}$} (او د \LR{$\{1, 2, 1, 2, 3\}$}) هم غړے دے، او برعکس۔
\end{ex}

\olpseganchor{OLP-0005-B009}{end}

\olpseganchor{OLP-0005-B010}{start}
ډېر ځله به يو سټ د يو داسې خاصيت په وسيله ټاکو چې د هغۀ غړي ئې په ګډه لري۔ د دې دپاره به دا لنډه نښه کاروو: \LR{$\Setabs{x}{\phi(x)}$}، چې پکښې \LR{$\phi(x)$} د هغه خاصيت استازيتوب کوي چې~\LR{$x$} ئې بايد ولري، څو د سټ په غړو کښې وشمېرل شي۔
\olpseganchor{OLP-0005-B010}{end}

\olpseganchor{OLP-0005-B011}{start}

\begin{ex}
زمونږ په بېلګه کښې، \LR{$S$} په دې بڼه هم ټاکلے شو:
\[
S = \Setabs{x}{x \text{\RL{ د رېچرډ ورور يا خور دے}}}.
\]
\end{ex}

\olpseganchor{OLP-0005-B011}{end}

\olpseganchor{OLP-0005-B012}{start}

\begin{ex}
يو عدد هله او يوازې هله \emph{کامل} بللے شي چې د خپلو حقيقي قاسمونو له مجموعې سره برابر وي (يعنې هغه عددونه چې دا عدد پرې پوره وېشل کېږي، خو پخپله له دې عدد سره برابر نۀ وي)۔ د بېلګې په توګه، \LR{$6$} کامل دے ځکه چې حقيقي قاسمونه ئې \LR{$1$}، \LR{$2$} او~\LR{$3$} دي، او \LR{$6 = 1 + 2 + 3$}۔ په حقيقت کښې، \LR{$6$} له \LR{$10$} څخه يوازينے کوچنے مثبت صحيح عدد دے چې کامل دے۔ نو د غړو له مخې د برابرۍ په کارولو سره وئيلے شو:
\[
	\{6\} = \Setabs{x}{x\text{\RL{ کامل دے او }}0 \leq x \leq 10}
\]
ښي لاس ته نښه داسې لولو: ``د هغو \LR{$x$} ګانو سټ چې \LR{$x$} کامل وي او \LR{$0 \leq x \leq 10$} وي''۔ دلته برابري دا خبره تائيدوي چې د سټونو په پام کښې نيولو پر وخت، د هغو د ټاکلو طريقه مهمه نۀ ده۔ او په لا عمومي توګه، د غړو له مخې برابري ډاډ راکوي چې تل د هغو \LR{$x$} ګانو يوازې يو سټ وي چې \LR{$\phi(x)$} وي۔ نو د غړو له مخې برابري دا جواز راکوي چې \LR{$\Setabs{x}{\phi(x)}$} ته د هغو \LR{$x$} ګانو \emph{هم هغه} سټ ووايو چې~\LR{$\phi(x)$} وي۔
\end{ex}

\olpseganchor{OLP-0005-B012}{end}

\olpseganchor{OLP-0005-B013}{start}
د غړو له مخې برابري راته د دې ښودلو يوه طريقه راکوي چې سټونه برابر دي: د \LR{$A = B$} د ښودلو دپاره، وښايئ چې هر کله \LR{$x \in A$} وي نو \LR{$x \in B$} هم وي، او هر کله \LR{$y \in B$} وي نو \LR{$y \in A$} هم وي۔
\olpseganchor{OLP-0005-B013}{end}

\olpseganchor{OLP-0005-B014}{start}
\begin{prob}
ثابت کړئ چې له يو څخه زيات تش سټونه نۀ شي کېدے، يعنې وښايئ چې که \LR{$A$} او \LR{$B$} داسې سټونه وي چې هېڅ غړي نۀ لري، نو \LR{$A = B$}۔
\end{prob}
\olpseganchor{OLP-0005-B014}{end}

\olpseganchor{OLP-0006-B004}{start}
\olfileid{sfr}{set}{sub}
\olsection{فرعي سټونه او د فرعي سټونو سټونه}
\olpseganchor{OLP-0006-B004}{end}

\olpseganchor{OLP-0006-B005}{start}
\begin{explain}
مونږ به ډېر ځله سټونه سره پرتله کول غواړو۔ د پرتله کولو يو څرګند ډول داسې دے: \emph{هر هغه څيز چې په يو سټ کښې وي، په بل کښې هم وي}۔ دا حالت دومره مهم دے چې د دې دپاره نوې نښې معرفي کړو۔
\end{explain}
\olpseganchor{OLP-0006-B005}{end}

\olpseganchor{OLP-0006-B006}{start}
\begin{defn}[فرعي سټ]
که د يو سټ \LR{$A$} هر غړے د~\LR{$B$} هم غړے وي، نو وايو چې \LR{$A$} د~\LR{$B$} يو \emph{فرعي سټ} دے، او \LR{$A \subseteq B$} ليکو۔ که \LR{$A$} د~\LR{$B$} فرعي سټ نۀ وي، \LR{$A \not\subseteq B$} ليکو۔ که \LR{$A \subseteq B$} وي خو \LR{$A \neq B$} وي، \LR{$A \subsetneq B$} ليکو او وايو چې \LR{$A$} د \LR{$B$} يو \emph{خاص فرعي سټ} دے۔
\end{defn}
\olpseganchor{OLP-0006-B006}{end}

\olpseganchor{OLP-0006-B007}{start}
\begin{ex}
هر سټ د خپل ځان فرعي سټ دے، او \LR{$\emptyset$} د هر سټ فرعي سټ دے۔ د جفتو طبيعي عددونو سټ د طبيعي عددونو د سټ فرعي سټ دے۔ همدارنګه، \LR{$\{ a, b \} \subseteq \{ a, b, c \}$}۔ خو \LR{$\{ a, b, e \}$} د \LR{$\{ a, b, c \}$} فرعي سټ نۀ دے۔
\end{ex}
\olpseganchor{OLP-0006-B007}{end}

\olpseganchor{OLP-0006-B008}{start}
\begin{ex}
عدد \LR{$2$} د صحيح عددونو د سټ يو غړے دے، حال دا چې د جفتو عددونو سټ د صحيح عددونو د سټ فرعي سټ دے۔ خو يو سټ کېدے شي د بل سټ \emph{هم} غړے وي او هم فرعي سټ، د بېلګې په توګه، \LR{$\{0\} \in \{0, \{0\}\}$} او هم \LR{$\{0\} \subseteq \{0, \{0\}\}$}۔
\end{ex}
\olpseganchor{OLP-0006-B008}{end}

\olpseganchor{OLP-0006-B009}{start}
د غړو له مخې برابري د سټونو د برابرۍ معيار راکوي: \LR{$A = B$} هله او يوازې هله وي چې د~\LR{$A$} هر غړے د~\LR{$B$} هم غړے وي، او برعکس۔ د ``فرعي سټ'' تعريف \LR{$A \subseteq B$} بالکل د دې معيار د لومړۍ نيمې په توګه تعريفوي: د~\LR{$A$} هر غړے د~\LR{$B$} هم غړے وي۔ البته، که \LR{$A$} او \LR{$B$} سره بدل کړو، تعريف بيا هم عملي کېږي: يعنې \LR{$B \subseteq A$} هله او يوازې هله وي چې د~\LR{$B$} هر غړے د~\LR{$A$} هم غړے وي۔ او دا بيا بالکل د غړو له مخې د برابرۍ ``برعکس'' برخه ده۔ په نورو ټکو، د غړو له مخې له برابرۍ نه دا نتيجه راوځي چې سټونه هله او يوازې هله برابر وي چې د يو بل فرعي سټونه وي۔
\olpseganchor{OLP-0006-B009}{end}

\olpseganchor{OLP-0006-B010}{start}
\begin{prop}
\LR{$A = B$} هله او يوازې هله وي چې \LR{$A \subseteq B$} او \LR{$B \subseteq A$} دواړه وي۔
\end{prop}
\olpseganchor{OLP-0006-B010}{end}

\olpseganchor{OLP-0006-B011}{start}
اوس د ځينو نورو ګټورو نښو د معرفي کولو دپاره هم ښه موقع ده۔ د دې په تعريف کښې چې \LR{$A$} کله د~\LR{$B$} فرعي سټ وي، مونږ ووئيل چې ``د~\LR{$A$} هر غړے \dots دے''، او د ``\LR{$\dots$}'' ځاے مو په ``د \LR{$B$} غړے'' ډک کړ۔ خو دا د عبارت دومره عامه \emph{بڼه} ده چې د دې دپاره د صوري نښو معرفي کول به ګټور وي۔
\olpseganchor{OLP-0006-B011}{end}

\olpseganchor{OLP-0006-B012}{start}
\begin{defn}\ollabel{forallxina}
\LR{$(\forall x \in A)\phi$} د \LR{$\forall x(x \in A \lif \phi)$} لنډه بڼه ده۔ همدارنګه، \LR{$(\exists x \in A)\phi$} د \LR{$\exists x(x \in A \land \phi)$} لنډه بڼه ده۔
\end{defn}
\olpseganchor{OLP-0006-B012}{end}

\olpseganchor{OLP-0006-B013}{start}
د دې نښو په کارولو سره وئيلے شو چې \LR{$A \subseteq B$} هله او يوازې هله وي چې \LR{$(\forall x \in A)x \in B$} وي۔
\olpseganchor{OLP-0006-B013}{end}

\olpseganchor{OLP-0006-B014}{start}
اوس د يو ځانګړي ډول سټ څېړنې ته ځو: د يو ټاکلي سټ د ټولو فرعي سټونو سټ۔
\olpseganchor{OLP-0006-B014}{end}

\olpseganchor{OLP-0006-B015}{start}
\begin{defn}[د فرعي سټونو سټ]
هغه سټ چې د يو سټ~\LR{$A$} ټول فرعي سټونه پکښې وي، د~\LR{$A$} \emph{د فرعي سټونو سټ} بللے شي، او \LR{$\Pow{A}$} ليکل کېږي۔
  \[
    \Pow{A} = \Setabs{B}{B \subseteq A}
  \]
\end{defn}
\olpseganchor{OLP-0006-B015}{end}

\olpseganchor{OLP-0006-B016}{start}
\begin{ex}
د \LR{$\{ a, b, c \}$} ټول ممکن فرعي سټونه کوم دي؟ دا دي: \LR{$\emptyset$}، \LR{$\{a \}$}، \LR{$\{b\}$}، \LR{$\{c\}$}، \LR{$\{a, b\}$}، \LR{$\{a, c\}$}، \LR{$\{b, c\}$}، \LR{$\{a, b, c\}$}۔ د دې ټولو فرعي سټونو سټ \LR{$\Pow{\{a,b,c\}}$} دے:
\[
\Pow{\{ a, b, c \}} = \{\emptyset, \{a \}, \{b\}, \{c\}, \{a, b\},
\{b, c\}, \{a, c\}, \{a, b, c\}\}
\]
\end{ex}
\olpseganchor{OLP-0006-B016}{end}

\olpseganchor{OLP-0006-B017}{start}
\begin{prob}
د \LR{$\{a, b, c, d\}$} ټول فرعي سټونه وليکئ۔
\end{prob}
\olpseganchor{OLP-0006-B017}{end}

\olpseganchor{OLP-0006-B018}{start}
\begin{prob}
وښايئ چې که \LR{$A$} د \LR{$n$} په شمېر غړي ولري، نو \LR{$\Pow{A}$} د \LR{$2^n$} په شمېر غړي لري۔
\end{prob}
\olpseganchor{OLP-0006-B018}{end}

\olpseganchor{OLP-0007-B004}{start}
\olfileid{sfr}{set}{imp}
\olsection{ځينې مهم سټونه}
\olpseganchor{OLP-0007-B004}{end}

\olpseganchor{OLP-0007-B005}{start}
\begin{ex}
مونږ به زياتره له هغو سټونو سره کار لرو چې غړي ئې رياضيکي څيزونه وي۔ څلور داسې سټونه دومره مهم دي چې ځانګړي نومونه لري:
\begin{multline*}
    \Nat = \{0, 1, 2, 3, \ldots\} \\
    \shoveright{\text{\RL{د طبيعي عددونو سټ}}}\\
    \shoveleft{\Int = \{\ldots, -2, -1, 0, 1, 2, \ldots\}} \\
    \shoveright{\text{\RL{د صحيح عددونو سټ}}}\\
    \shoveleft{\Rat = \Setabs{\nicefrac{m}{n}}{m, n \in \Int\text{\RL{ او }}n \neq 0}}\\
    \shoveright{\text{\RL{د ناطقو عددونو سټ}}}\\
    \shoveleft{\Real = (-\infty, \infty)}\\
    \text{\RL{د حقيقي عددونو سټ (پيوستار)}}
\end{multline*}
دا ټول \emph{نامتناهي} سټونه دي، يعنې د هر يوه غړي په نامتناهي شمېر دي۔
\olpseganchor{OLP-0007-B005}{end}

\olpseganchor{OLP-0007-B006}{start}
چې کله د دې سټونو له يوه نه بل ته ځو، خپلې زېرمې ته \emph{نور} عددونه ورزياتوو۔ په حقيقت کښې، دا بايد څرګنده وي چې \LR{$\Nat \subseteq \Int \subseteq \Rat \subseteq \Real$}: ځکه چې هر طبيعي عدد صحيح عدد دے؛ هر صحيح عدد ناطق عدد دے؛ او هر ناطق عدد حقيقي عدد دے۔ همدارنګه، دا هم بايد څرګنده وي چې \LR{$\Nat \subsetneq \Int \subsetneq \Rat$}، ځکه چې \LR{$-1$} صحيح عدد دے خو طبيعي عدد نۀ دے، او \LR{$\nicefrac{1}{2}$} ناطق عدد دے خو صحيح عدد نۀ دے۔ دا لږه ناڅرګنده خبره ده چې \LR{$\Rat \subsetneq \Real$}، يعنې داسې حقيقي عددونه شته چې ناطق نۀ دي۔
\olpseganchor{OLP-0007-B006}{end}

\olpseganchor{OLP-0007-B007}{start}
کله ناکله به د مثبتو صحيح عددونو سټ \LR{$\PosInt = \{1, 2, 3, \dots\}$} او هغه سټ هم کاروو چې يوازې لومړي دوه طبيعي عددونه لري، يعنې \LR{$\Bin = \{0, 1\}$}۔
\end{ex}
\olpseganchor{OLP-0007-B007}{end}

\olpseganchor{OLP-0007-B008}{start}

\begin{ex}[توريز تسلسلونه]
يوه بله په زړه پورې بېلګه د يوې الفبا \LR{$A$} د \emph{متناهي توريزو تسلسلونو} سټ \LR{$A^{*}$} دے: د~\LR{$A$} د عناصرو هره متناهي لړۍ د \LR{$A$} يو توريز تسلسل دے۔ د هرې الفبا~\LR{$A$} دپاره، \emph{تش توريز تسلسل \LR{$\Lambda$}} هم د~\LR{$A$} په توريزو تسلسلونو کښې شاملوو۔ د بېلګې په توګه،
\begin{multline*}
\Bin^*
=\{\Lambda,0,1,00,01,10,11,\\
000,001,010,011,100,101,110,111,0000,\ldots\}.
\end{multline*}
که \LR{$x=x_{1}\ldots x_{n}\in A^{*}$} يو داسې توريز تسلسل وي چې د \LR{$A$} له \LR{$n$} ``تورو'' جوړ وي، نو وايو چې د دې تسلسل \emph{اوږدوالے}~\LR{$n$} دے او \LR{$\len{x}=n$} ليکو۔
\end{ex}

\olpseganchor{OLP-0007-B008}{end}

\olpseganchor{OLP-0007-B009}{start}
\begin{ex}[نامتناهي لړۍ]
د هر سټ \LR{$A$} دپاره، د~\LR{$A$} د غړو د نامتناهي لړيو سټ~\LR{$A^\omega$} هم په پام کښې نيولے شو۔ يوه نامتناهي لړۍ \LR{$a_1a_2a_3a_4\dots$} د څيزونو له داسې لست نه جوړه وي چې په يو لوري نامتناهي غځېږي، او هر څيز ئې د~\LR{$A$} يو غړے وي۔
\end{ex}
\olpseganchor{OLP-0007-B009}{end}

\olpseganchor{OLP-0008-B004}{start}
\olfileid{sfr}{set}{uni}
\olsection{اتحادونه او تقاطعې}
\olpseganchor{OLP-0008-B004}{end}

\olpseganchor{OLP-0008-B005}{start}
\begin{explain}
په \olref[sfr][set][bas]{sec} کښې مو د تجريد په وسيله د سټونو تعريفونه معرفي کړل، يعنې د \LR{$\Setabs{x}{\phi(x)}$} په بڼه تعريفونه۔ دلته يو خاصيت~\LR{$\phi$} په کار راوړو، او دا خاصيت د هغو سټونو يادونه کولے شي چې مخکې مو تعريف کړي دي۔ نو د بېلګې په توګه، که \LR{$A$} او~\LR{$B$} سټونه وي، د \LR{$\Setabs{x}{x \in A \lor x \in B}$} سټ له هغو ټولو څيزونو جوړ دے چې د \LR{$A$} يا~\LR{$B$} غړي وي، يعنې دا هغه سټ دے چې د \LR{$A$} او~\LR{$B$} غړي سره يوځاے کوي۔ دا لکه په \olref{fig:union} کښې انځورولے شو، چې رنګ شوې ساحه د دواړو سټونو \LR{$A$} او~\LR{$B$} غړي په ګډه ښيي۔
\olpseganchor{OLP-0008-B005}{end}

\olpseganchor{OLP-0008-B006}{start}
\begin{figure}
  \olasset{../upstream/assets/diagrams/union.tikz}
  \caption{د دوو سټونو اتحاد \LR{$A \cup B$} د~\LR{$A$} د غړو او د~\LR{$B$} د غړو ګډ سټ دے۔}
  \ollabel{fig:union}
\end{figure}
\olpseganchor{OLP-0008-B006}{end}

\olpseganchor{OLP-0008-B007}{start}
په سټونو دا عمل، يعنې يوځاے کول ئې، ډېر ګټور او عام دے؛ نو يو صوري نوم او نښه ورکوو۔
\end{explain}
\olpseganchor{OLP-0008-B007}{end}

\olpseganchor{OLP-0008-B008}{start}
\begin{defn}[اتحاد]
د دوو سټونو \LR{$A$} او \LR{$B$} \emph{اتحاد}، چې \LR{$A \cup B$} ليکل کېږي، د ټولو هغو څيزونو سټ دے چې د \LR{$A$}، \LR{$B$}، يا د دواړو غړي وي۔
\[
A \cup B = \Setabs{x}{x \in A \lor x \in B}
\]
\end{defn}
\olpseganchor{OLP-0008-B008}{end}

\olpseganchor{OLP-0008-B009}{start}
\begin{ex}
څرنګه چې د غړو تکرار مهم نۀ دے، د هغو دوو سټونو اتحاد چې ګډ غړے ولري، هغه غړے يوازې يو ځل لري؛ د بېلګې په توګه، \LR{$\{ a, b, c\} \cup \{ a, 0, 1\} = \{a, b, c, 0, 1\}$}۔
\olpseganchor{OLP-0008-B009}{end}

\olpseganchor{OLP-0008-B010}{start}
د يو سټ او د هغۀ د يو فرعي سټ اتحاد هم هغه لوے سټ دے: \LR{$\{a, b, c \} \cup \{a \} = \{a, b, c\}$}۔
\olpseganchor{OLP-0008-B010}{end}

\olpseganchor{OLP-0008-B011}{start}
له تش سټ سره د يو سټ اتحاد له هم هغه سټ سره برابر دے: \LR{$\{a, b, c \} \cup \emptyset = \{a, b, c \}$}۔
\end{ex}
\olpseganchor{OLP-0008-B011}{end}

\olpseganchor{OLP-0008-B012}{start}
\begin{prob}
ثابت کړئ چې که \LR{$A \subseteq B$} وي، نو \LR{$A \cup B = B$}۔
\end{prob}
\olpseganchor{OLP-0008-B012}{end}

\olpseganchor{OLP-0008-B013}{start}
\begin{explain}
د اتحاد يو ``دوه ګونے'' عمل هم په پام کښې نيولے شو۔ دا هغه عمل دے چې د هغو ټولو غړو سټ جوړوي چې د~\LR{$A$} غړي وي او د~\LR{$B$} هم غړي وي۔ دې عمل ته \emph{تقاطع} وئيل کېږي، او لکه په \olref{fig:intersection} کښې ښودل کېدے شي۔
\begin{figure}
  \olasset{../upstream/assets/diagrams/intersection.tikz}
  \caption{د دوو سټونو تقاطع \LR{$A \cap B$} د هغو د ګډو غړو سټ دے۔}
  \ollabel{fig:intersection}
\end{figure}
\end{explain}
\olpseganchor{OLP-0008-B013}{end}

\olpseganchor{OLP-0008-B014}{start}
\begin{defn}[تقاطع]
د دوو سټونو \LR{$A$} او \LR{$B$} \emph{تقاطع}، چې \LR{$A \cap B$} ليکل کېږي، د ټولو هغو څيزونو سټ دے چې د \LR{$A$} او~\LR{$B$} دواړو غړي وي۔
\[
A \cap B = \Setabs{x}{x \in A \land x \in B}
\]
دوه سټونه هله \emph{بېل} بللے شي چې تقاطع ئې تشه وي۔ مطلب دا چې هېڅ ګډ غړي نۀ لري۔
\end{defn}
\olpseganchor{OLP-0008-B014}{end}

\olpseganchor{OLP-0008-B015}{start}
\begin{ex}
که دوه سټونه هېڅ ګډ غړي ونۀ لري، تقاطع ئې تشه ده: \LR{$\{ a, b, c\} \cap \{ 0, 1\} = \emptyset$}۔
\olpseganchor{OLP-0008-B015}{end}

\olpseganchor{OLP-0008-B016}{start}
که دوه سټونه ګډ غړي ولري، تقاطع ئې د هغو ټولو سټ دے: \LR{$\{a, b, c \} \cap \{a, b, d \} = \{a, b\}$}۔
\olpseganchor{OLP-0008-B016}{end}

\olpseganchor{OLP-0008-B017}{start}
د يو سټ او د هغۀ د يو فرعي سټ تقاطع هم هغه کوچنے سټ دے: \LR{$\{a, b, c\} \cap \{a, b\} = \{a, b\}$}۔
\olpseganchor{OLP-0008-B017}{end}

\olpseganchor{OLP-0008-B018}{start}
له تش سټ سره د هر سټ تقاطع تشه ده: \LR{$\{a, b, c \} \cap \emptyset = \emptyset$}۔
\end{ex}
\olpseganchor{OLP-0008-B018}{end}

\olpseganchor{OLP-0008-B019}{start}
\begin{prob}
په دقيقه توګه ثابت کړئ چې که \LR{$A \subseteq B$} وي، نو \LR{$A \cap B = A$}۔
\end{prob}
\olpseganchor{OLP-0008-B019}{end}

\olpseganchor{OLP-0008-B020}{start}
\begin{explain}
له دوو څخه د زياتو سټونو اتحاد يا تقاطع هم جوړولے شو۔ په عمومي توګه د دې د بيانولو يوه ښکلې طريقه دا ده: فرض کړئ چې ټول هغه سټونه چې اتحاد (يا تقاطع) ئې جوړول غواړئ، په يو سټ کښې راټول کړئ۔ بيا د ټولو اصلي سټونو اتحاد د هغو ټولو څيزونو د سټ په توګه تعريفولے شو چې لږ تر لږه د دې سټ د يو غړي غړي وي، او تقاطع د هغو ټولو څيزونو د سټ په توګه چې د دې سټ د هر غړي غړي وي۔
\end{explain}
\olpseganchor{OLP-0008-B020}{end}

\olpseganchor{OLP-0008-B021}{start}
\begin{defn}
که \LR{$A$} د سټونو يو سټ وي، نو \LR{$\bigcup A$} د~\LR{$A$} د غړو د غړو سټ دے:
\begin{align*}
\bigcup A & = \Setabs{x}{x \text{\RL{ د دې سټ د يو غړي غړے دے: }} A},
\text{\RL{ يعنې،}}\\
& = \Setabs{x}{\text{\RL{داسې يو }} B \in A
  \text{\RL{ شته چې }} x \in B}
\end{align*}
\end{defn}
\olpseganchor{OLP-0008-B021}{end}

\olpseganchor{OLP-0008-B022}{start}
\begin{defn}
که \LR{$A$} د سټونو يو سټ وي، نو \LR{$\bigcap A$} د هغو څيزونو سټ دے چې د~\LR{$A$} ټول عناصر ئې په ګډه لري:
\begin{align*}
\bigcap A & = \Setabs{x}{x \text{\RL{ د دې سټ د هر غړي غړے دے: }} A},
\text{\RL{ يعنې،}}\\
 & = \Setabs{x}{\text{\RL{د هر }} B \in A, x \in B}
\end{align*}
\end{defn}
\olpseganchor{OLP-0008-B022}{end}

\olpseganchor{OLP-0008-B023}{start}
\begin{ex}
فرض کړئ چې \LR{$A = \{ \{ a, b \}, \{ a, d, e \}, \{ a, d \} \}$}۔ بيا \LR{$\bigcup A = \{ a, b, d, e \}$} او \LR{$\bigcap A = \{ a \}$}۔
\end{ex}
\begin{prob}
	وښايئ چې که \LR{$A$} يو سټ وي او \LR{$A \in B$} وي، نو \LR{$A \subseteq \bigcup B$}۔
\end{prob}
\olpseganchor{OLP-0008-B023}{end}

\olpseganchor{OLP-0008-B024}{start}
د سټونو د يوې لړۍ \LR{$A_1$}، \LR{$A_2$}، \dots دپاره هم دا کار کولے شو:
\begin{align*}
\bigcup_i A_i & = \Setabs{x}{x \text{\RL{ له دې سټونو څخه د يوه غړے دے: }} A_i}\\
\bigcap_i A_i & = \Setabs{x}{x \text{\RL{ د دې هر سټ غړے دے: }} A_i}.
\end{align*}
\olpseganchor{OLP-0008-B024}{end}

\olpseganchor{OLP-0008-B025}{start}
چې کله د سټونو يو \emph{اندکس} ولرو، يعنې داسې يو سټ \LR{$I$} چې د هر \LR{$i \in I$} دپاره \LR{$A_i$} په پام کښې نيسو، دا لنډې بڼې هم کارولے شو:
\begin{align*}
	\bigcup_{i \in I} A_i & = \bigcup \Setabs{A_i }{i \in I}\\
	\bigcap_{i \in I} A_i & = \bigcap\Setabs{A_i}{i \in I}
\end{align*}
\olpseganchor{OLP-0008-B025}{end}

\olpseganchor{OLP-0008-B026}{start}
په پای کښې، کېدے شي د~\LR{$A$} د هغو ټولو غړو سټ په پام کښې نيول وغواړو چې په~\LR{$B$} کښې نۀ وي۔ دا لکه په \olref{difference} کښې انځورولے شو۔
\olpseganchor{OLP-0008-B026}{end}

\olpseganchor{OLP-0008-B027}{start}
\begin{figure}
  \olasset{../upstream/assets/diagrams/difference.tikz}
  \caption{د دوو سټونو تفاضل \LR{$A \setminus B$} د~\LR{$A$} د هغو غړو سټ دے چې د~\LR{$B$} غړي نۀ وي۔}
  \ollabel{difference}
\end{figure}
\olpseganchor{OLP-0008-B027}{end}

\olpseganchor{OLP-0008-B028}{start}
\begin{defn}[تفاضل]
د سټونو \emph{تفاضل}~\LR{$A \setminus B$} د \LR{$A$} د هغو ټولو غړو سټ دے چې د~\LR{$B$} غړي نۀ وي، يعنې،
\[
A\setminus B = \Setabs{x}{x\in A \text{\RL{ او }} x \notin B}.
\]
\end{defn}
\olpseganchor{OLP-0008-B028}{end}

\olpseganchor{OLP-0008-B029}{start}
\begin{prob}
	ثابت کړئ چې که \LR{$A \subsetneq B$} وي، نو \LR{$B \setminus A \neq \emptyset$}۔
\end{prob}
\olpseganchor{OLP-0008-B029}{end}

\olpseganchor{OLP-0009-B004}{start}
\olfileid{sfr}{set}{pai}
\olsection{جوړې، څوګونې او کارتيزي حاصل ضربونه}
\olpseganchor{OLP-0009-B004}{end}

\olpseganchor{OLP-0009-B005}{start}
\begin{explain}
د غړو له مخې له برابرۍ نه دا نتيجه راوځي چې د سټونو عناصر ترتيب نۀ لري۔ نو که د ترتيب استازيتوب کول غواړو، \emph{مرتبې جوړې} \LR{$\tuple{x, y}$} کاروو۔ په يوه نامرتبه جوړه \LR{$\{x, y\}$} کښې ترتيب مهم نۀ دے: \LR{$\{x, y\} = \{y, x\}$}۔ په مرتبه جوړه کښې ترتيب مهم دے: که \LR{$x \neq y$} وي، نو \LR{$\tuple{x, y} \neq \tuple{y, x}$}۔
\olpseganchor{OLP-0009-B005}{end}

\olpseganchor{OLP-0009-B006}{start}
د سټونو په نظريه کښې مرتبې جوړې څنګه په پام کښې ونيسو؟ بنيادي خبره دا ده چې دا نظر وساتو: مرتبې جوړې هله او يوازې هله برابرې وي چې لومړے عنصر ئې هم يو وي او دويم عنصر ئې هم يو وي، يعنې:
\[
  \tuple{a, b}= \tuple{c, d}\text{\RL{ هله او يوازې هله چې دواړه }}a = c \text{\RL{ او }}b=d.
\]
د وينر--کوراتوسکي تعريف په کارولو سره د سټونو په نظريه کښې مرتبې جوړې تعريفولے شو۔
\end{explain}
\olpseganchor{OLP-0009-B006}{end}

\olpseganchor{OLP-0009-B007}{start}
\begin{defn}[مرتبه جوړه]\ollabel{wienerkuratowski}
	\LR{$\tuple{a, b} = \{\{a\}, \{a, b\}\}$}۔
\end{defn}
\olpseganchor{OLP-0009-B007}{end}

\olpseganchor{OLP-0009-B008}{start}
\begin{prob}
	د \olref[sfr][set][pai]{wienerkuratowski} په کارولو سره ثابت کړئ چې \LR{$\tuple{a, b}= \tuple{c, d}$} هله او يوازې هله وي چې \LR{$a = c$} او \LR{$b=d$} دواړه وي۔
\end{prob}
\olpseganchor{OLP-0009-B008}{end}

\olpseganchor{OLP-0009-B009}{start}
\begin{explain}
چې د مرتبې جوړې تعريف مو وټاکلو، د نورو سټونو د تعريف دپاره ئې کارولے شو۔ د بېلګې په توګه، کله کله له دوو څخه د زياتو څيزونو مرتبې لړۍ هم غواړو، لکه \emph{درې ګونې} \LR{$\tuple{x, y, z}$}، \emph{څلورګونې} \LR{$\tuple{x, y, z, u}$}، او داسې نورې۔ درې ګونې د ځانګړو مرتبو جوړو په توګه ګڼلے شو، چې لومړے عنصر ئې پخپله يوه مرتبه جوړه وي: \LR{$\tuple{x, y, z}$} هم \LR{$\tuple{\tuple{x, y},z}$} دے۔ د څلورګونو دپاره هم همداسې ده: \LR{$\tuple{x,y,z,u}$} هم \LR{$\tuple{\tuple{\tuple{x,y},z},u}$} دے، او داسې نور۔ په عمومي توګه، د \emph{مرتبو \LR{$n$}-ګونو} \LR{$\tuple{x_1, \dots, x_n}$} خبره کوو۔
\olpseganchor{OLP-0009-B009}{end}

\olpseganchor{OLP-0009-B010}{start}
د مرتبو جوړو، يا د نورو مرتبو \LR{$n$}-ګونو، ځينې سټونه به ګټور وي۔
\end{explain}
\olpseganchor{OLP-0009-B010}{end}

\olpseganchor{OLP-0009-B011}{start}
\begin{defn}[کارتيزي حاصل ضرب]
د ورکړل شوو سټونو \LR{$A$} او \LR{$B$} \emph{کارتيزي حاصل ضرب} \LR{$A \times B$} داسې تعريفېږي:
\[
  A \times B = \Setabs{\tuple{x, y}}{x \in A \text{\RL{ او }} y \in B}.
\]
\end{defn}
\olpseganchor{OLP-0009-B011}{end}

\olpseganchor{OLP-0009-B012}{start}
\begin{ex}
که \LR{$A = \{0, 1\}$} او \LR{$B = \{1, a, b\}$} وي، نو حاصل ضرب ئې دا دے:
\[
A \times B = \{ \tuple{0, 1}, \tuple{0, a}, \tuple{0, b},
    \tuple{1, 1}, \tuple{1, a}, \tuple{1, b} \}.
\]
\end{ex}
\olpseganchor{OLP-0009-B012}{end}

\olpseganchor{OLP-0009-B013}{start}
\begin{ex}
که \LR{$A$} يو سټ وي، له خپل ځان سره د \LR{$A$} حاصل ضرب، \LR{$A \times A$}، د~\LR{$A^2$} په بڼه هم ليکل کېږي۔ دا د هغو \emph{ټولو} جوړو \LR{$\tuple{x, y}$} سټ دے چې \LR{$x, y \in A$} وي۔ د ټولو درې ګونو \LR{$\tuple{x, y, z}$} سټ \LR{$A^3$} دے، او داسې نور۔ يو بازګشتي تعريف ورکولے شو:
\begin{align*}
  A^1 & = A\\
  A^{k+1} & = A^k \times A
\end{align*}
\end{ex}
\olpseganchor{OLP-0009-B013}{end}

\olpseganchor{OLP-0009-B014}{start}
\begin{prob}
د \LR{$\{1, 2, 3\}^3$} ټول غړي وليکئ۔
\end{prob}
\olpseganchor{OLP-0009-B014}{end}

\olpseganchor{OLP-0009-B015}{start}
\begin{prop}\ollabel{cardnmprod}
که \LR{$A$} د \LR{$n$} په شمېر غړي او \LR{$B$} د \LR{$m$} په شمېر غړي ولري، نو \LR{$A \times B$} د \LR{$n\cdot m$} په شمېر عناصر لري۔
\end{prop}
\olpseganchor{OLP-0009-B015}{end}

\olpseganchor{OLP-0009-B016}{start}
\begin{proof}
په~\LR{$A$} کښې د هر غړي~\LR{$x$} دپاره، د \LR{$\tuple{x, y} \in A \times B$} په بڼه د \LR{$m$} په شمېر غړي شته۔ فرض کړئ چې \LR{$B_x = \Setabs{\tuple{x, y}}{y \in B}$}۔ څرنګه چې هر کله \LR{$x_1 \neq x_2$} وي، \LR{$\tuple{x_1, y} \neq \tuple{x_2, y}$} وي، نو \LR{$B_{x_1} \cap B_{x_2} = \emptyset$}۔ خو که \LR{$A = \{x_1, \dots, x_n\}$} وي، نو \LR{$A \times B = B_{x_1} \cup \dots \cup B_{x_n}$}، او ځکه د \LR{$n\cdot m$} په شمېر غړي لري۔
\olpseganchor{OLP-0009-B016}{end}

\olpseganchor{OLP-0009-B017}{start}
د دې د انځورولو دپاره، د~\LR{$A \times B$} غړي په يوه جالۍ کښې ترتيب کړئ:
\[
\begin{array}{rcccc}
  B_{x_1} = & \{\tuple{x_1, y_1} & \tuple{x_1, y_2} & \dots & \tuple{x_1, y_m}\}\\
  B_{x_2} = & \{\tuple{x_2, y_1} & \tuple{x_2, y_2} & \dots & \tuple{x_2, y_m}\}\\
  \vdots & & \vdots\\
  B_{x_n} = & \{\tuple{x_n, y_1} & \tuple{x_n, y_2} & \dots & \tuple{x_n, y_m}\}
\end{array}
\]
څرنګه چې ټول \LR{$x_i$} يو له بله بېل دي، او ټول \LR{$y_j$} يو له بله بېل دي، په دې جالۍ کښې هېڅ دوه جوړې يو شان نۀ دي، او شمېر ئې \LR{$n\cdot m$} دے۔
\end{proof}
\olpseganchor{OLP-0009-B017}{end}

\olpseganchor{OLP-0009-B018}{start}
\begin{prob}
په~\LR{$k$} د استقرا په وسيله وښايئ چې د هر \LR{$k \ge 1$} دپاره، که \LR{$A$} د \LR{$n$} په شمېر غړي ولري، نو \LR{$A^k$} د \LR{$n^k$} په شمېر غړي لري۔
\end{prob}
\olpseganchor{OLP-0009-B018}{end}

\olpseganchor{OLP-0009-B019}{start}
\begin{ex}
که \LR{$A$} يو سټ وي، په~\LR{$A$} يوه \emph{کلمه} د~\LR{$A$} د غړو هره لړۍ ده۔ يوه لړۍ د~\LR{$A$} د غړو د يوې \LR{$n$}-ګونې په توګه ګڼل کېدے شي۔ د بېلګې په توګه، که \LR{$A = \{a, b, c\}$} وي، نو لړۍ ``\LR{$bac$}'' د درې ګونې~\LR{$\tuple{b, a, c}$} په توګه ګڼل کېدے شي۔ کلمې، يعنې د نښو لړۍ، په کمپيوټر ساينس کښې بنيادي اهميت لري۔ د منل شوي دود له مخې، د~\LR{$A$} غړي د~\LR{$1$} اوږدوالي لرونکې لړۍ ګڼو، او \LR{$\emptyset$} د~\LR{$0$} اوږدوالي لرونکې لړۍ ګڼو۔ بيا په~\LR{$A$} د \emph{ټولو} کلمو سټ دا دے:
\[
A^* = \{\emptyset\} \cup A \cup A^2 \cup A^3 \cup \dots
\]
\end{ex}
\olpseganchor{OLP-0009-B019}{end}

\olpseganchor{OLP-0721-B004}{start}
\olfileid{sfr}{set}{prf}
\olsection{د سټونو په اړه ثبوتونه}
\olpseganchor{OLP-0721-B004}{end}

\olpseganchor{OLP-0721-B005}{start}
\begin{editorial}
د ثبوتونو د طريقو عمومي باب په \verb|content/methods/proofs| کښې دے۔ دلته پر سټونو د دغو طريقو يوه ځانګړې بېلګه لولو۔
\end{editorial}
\olpseganchor{OLP-0721-B005}{end}

\olpseganchor{OLP-0721-B006}{start}
\begin{explain}
سټونه او تر دې ځايه معرفي شوے نښليک د سټونو د ټاکلو او د هغو تر منځ د اړيکو د څرګندولو لنډه لاره ده۔ ډېر وخت به د دغو اړيکو په باب دعوې هم ثابتوو۔ که له رياضيکي ثبوتونو سره اشنا نۀ يئ، يوه ساده بېلګه ګورو۔ د هر دوو سټونو \LR{$X$} او \LR{$Y$} دپاره ښيو چې تل \LR{$X \cap
(X \cup Y) = X$}۔ د سټونو داسې برابري څنګه ثابتوو؟ په ياد ولرئ چې سټونه يوازې له خپلو غړي ټاکل کېږي؛ دوه سټونه هغه او يوازې هغه وخت برابر دي چې هماغه غړي ولري۔ نو دلته ثابتوو چې (a) د \LR{$X \cap (X \cup
Y)$} هر غړے د \LR{$X$} هم غړے دے، او برعکس (b) د \LR{$X$} هر غړے د \LR{$X \cap (X \cup Y)$} هم غړے دے۔ يعنې دواړه ښيو، (a) \LR{$X \cap (X \cup Y) \subseteq X$} او (b) \LR{$X \subseteq X \cap (X \cup Y)$}۔
\olpseganchor{OLP-0721-B006}{end}

\olpseganchor{OLP-0721-B007}{start}
د «د \LR{$X \cap (X
\cup Y)$} هر غړے~\LR{$z$} د \LR{$X$} هم غړے دے» په څېر د يوې عامې دعوې د ثابتولو دپاره لومړے فرضوو چې \LR{$z \in X \cap (X \cup Y)$} يو اختياري غړے راکړے شوے او له همدې فرضه \LR{$z \in X$} ثابتوو۔ دا د «عمومي شرطي ثبوت» په نوم هم پېژندلے شئ۔ د فرض او پايلې تعريفونه هم کاروو، لکه دلته چې د \LR{$\cap$} او \LR{$\cup$} تعريفونه کاروو۔ نو د (a) دليل داسې دے۔
\olpseganchor{OLP-0721-B007}{end}

\olpseganchor{OLP-0721-B008}{start}
\begin{quote}
(a) لومړے غواړو وښيو چې \LR{$X \cap (X \cup Y) \subseteq X$}۔ د \LR{$\subseteq$} د تعريف له مخې، د هر \LR{$z$} دپاره، که \LR{$z \in X \cap (X \cup Y)$} نو \LR{$z \in X$}۔ فرض کړئ \LR{$z \in X \cap (X \cup
Y)$}۔ څرنګه چې \LR{$z$} د دوو سټونو د تقاطع غړے هغه او يوازې هغه وخت دے چې د دواړو سټونو غړے وي، نو \LR{$z \in X$} او \LR{$z \in X \cup Y$} اخلو۔ په ځانګړي ډول \LR{$z \in X$}، هماغه چې ثابتول مو غوښتل۔
\end{quote}
\olpseganchor{OLP-0721-B008}{end}

\olpseganchor{OLP-0721-B009}{start}
د ثبوت لومړۍ نيمه پوره شوه۔ په وروستي ګام کښې مو وکارول چې که د او تړنه، يعنې \LR{$z \in X$} او \LR{$z \in X
\cup Y$}، له يوه فرضه راځي، نو هره تړلې برخه ئې له هماغه فرضه راځي۔ دغه قاعده «د او ايستل» يا \LR{$\land$} ايستل بولي۔ اوس (b) ثابتوو۔
\olpseganchor{OLP-0721-B009}{end}

\olpseganchor{OLP-0721-B010}{start}
\begin{quote}
(b) اوس ثابتوو چې \LR{$X \subseteq X \cap (X \cup Y)$}۔ د \LR{$\subseteq$} د تعريف له مخې، د هر \LR{$z$} دپاره که \LR{$z \in X$}، نو \LR{$z \in X \cap
(X \cup Y)$} هم۔ فرض کړئ \LR{$z \in X$}۔ د \LR{$z \in X
\cap (X \cup Y)$} د ښودلو دپاره د \LR{$\cap$} له تعريفه دوه خبرې غواړو، (i) \LR{$z \in X$} او (ii) \LR{$z \in X \cup Y$}۔ (i) پخپله زمونږ فرض دے، نو نور ثبوت نۀ غواړي۔ د (ii) دپاره په ياد ولرئ چې \LR{$z$} د سټونو د اتحاد غړے هغه او يوازې هغه وخت دے چې لږ تر لږه د يوه سټ غړے وي۔ دلته دا شرط پوره دے، ځکه \LR{$z \in X$} او \LR{$X \cup Y$} د \LR{$X$} او \LR{$Y$} اتحاد دے۔ نو \LR{$z \in X \cup Y$}۔ (i) \LR{$z \in X$} او (ii) \LR{$z \in X \cup Y$} دواړه مو وښودل، نو د \LR{$\cap$} د تعريف له مخې \LR{$z \in X \cap (X \cup Y)$}۔
\end{quote}
\olpseganchor{OLP-0721-B010}{end}

\olpseganchor{OLP-0721-B011}{start}
دا بيان لږ اوږد دے، خو ښيي چې د سټونو او د هغو د اړيکو په اړه څنګه استدلال کوو۔ عموماً دومره تفصيل نۀ ورکوو؛ په ځانګړي ډول ښايي ټول تعريفونه بيا وانۀ يوو۔ په يوه پرمخ تللي متن کښې همدا ثبوت ډېر لنډ داسې ليکلے شو۔
\end{explain}
\olpseganchor{OLP-0721-B011}{end}

\olpseganchor{OLP-0721-B012}{start}
\begin{prop}[جذب]
د هر دوو سټونو \LR{$X$}، \LR{$Y$} دپاره،
\[
X \cap (X \cup Y) = X
\]
\end{prop}
\olpseganchor{OLP-0721-B012}{end}

\olpseganchor{OLP-0721-B013}{start}
\begin{proof}
(a) فرض کړئ \LR{$z \in X \cap (X \cup Y)$}۔ نو \LR{$z \in X$}، له دې امله \LR{$X \cap (X
  \cup Y) \subseteq X$}۔
\olpseganchor{OLP-0721-B013}{end}

\olpseganchor{OLP-0721-B014}{start}
(b) اوس فرض کړئ \LR{$z \in X$}۔ نو \LR{$z \in X \cup Y$} هم، او له دې امله \LR{$z \in X \cap (X \cup Y)$} هم۔ په نتيجه کښې \LR{$X \subseteq X \cap (X \cup
  Y)$}۔
\end{proof}
\olpseganchor{OLP-0721-B014}{end}

\olpseganchor{OLP-0721-B015}{start}
\begin{prob}
په تفصيل ثابت کړئ چې \LR{$X \cup (X \cap Y) = X$}۔ بيا يو لنډ ثبوت ورکړئ۔ (اشاره۔ د \LR{$X \cup (X \cap Y) \subseteq X$} لوري دپاره په جلا حالتونو ثبوت، يعنې \LR{$\lor$} ايستل، پکار دي۔)
\end{prob}
\olpseganchor{OLP-0721-B015}{end}

\olpseganchor{OLP-0010-B004}{start}
\olfileid{sfr}{set}{rus}
\olsection{د رسل پاراډوکس}
\olpseganchor{OLP-0010-B004}{end}

\olpseganchor{OLP-0010-B005}{start}
د غړو له مخې برابري د هغو \LR{$x$} ګانو د \emph{هم هغه} سټ دپاره د \LR{$\Setabs{x}{\phi(x)}$} نښې د کارولو جواز راکوي چې~\LR{$\phi(x)$} وي۔ خو د غړو له مخې برابري \emph{په اصل کښې} يوازې د دې خبرې جواز راکوي: \emph{که} داسې يو سټ وي چې غړي ئې ټول او يوازې هغه څيزونه وي چې \LR{$\phi$} وي، \emph{نو} داسې سټ يوازې يو وي۔ په بل عبارت: د يو~\LR{$\phi$} له ټاکلو وروسته، سټ \LR{$\Setabs{x}{\phi(x)}$} يوازينے دے، \emph{که موجود وي}۔
\olpseganchor{OLP-0010-B005}{end}

\olpseganchor{OLP-0010-B006}{start}
خو دا شرط مهم دے!{} بنيادي خبره دا ده چې له هر خاصيت څخه د \emph{سټ جوړول} ممکن نۀ دي۔ يعنې ځينې خاصيتونه سټونه \emph{نۀ} تعريفوي۔ که ټولو تعريفولے، نو له ښکاره تناقضونو سره به مخ شوو۔ د دې تر ټولو مشهوره بېلګه د رسل پاراډوکس دے۔
\olpseganchor{OLP-0010-B006}{end}

\olpseganchor{OLP-0010-B007}{start}
سټونه د نورو سټونو غړي کېدے شي؛ د بېلګې په توګه، د يو سټ~\LR{$A$} د فرعي سټونو سټ له سټونو جوړ وي۔ نو دا پوښتل يا څېړل معقول دي چې يو سټ د بل سټ غړے دے که نۀ۔ آيا يو سټ د خپل ځان غړے کېدے شي؟ د سټ په تصور کښې ظاهراً هېڅ داسې خبره نشته چې دا ناممکنه کړي۔ د بېلګې په توګه، که \emph{ټول} سټونه د څيزونو يوه ټولګه جوړوي، يو څوک ښائي فکر وکړي چې دا ټول په يو سټ کښې راټولېدے شي، يعنې د ټولو سټونو په سټ کښې۔ او دا چې پخپله يو سټ دے، نو د ټولو سټونو د سټ غړے به وي۔
\olpseganchor{OLP-0010-B007}{end}

\olpseganchor{OLP-0010-B008}{start}
د رسل پاراډوکس هغه وخت راولاړېږي چې دا خاصيت په پام کښې ونيسو چې يو سټ خپل ځان د غړے په توګه ونۀ لري، يعنې \emph{د خپل ځان غړے نۀ وي}۔ که فرض کړو چې د ټولو هغو سټونو يو سټ شته چې خپل ځان د غړے په توګه نۀ لري، نو څه به وشي؟ آيا
\[
R = \Setabs{x}{x \notin x}
\]
شته؟ څرګندېږي چې ثابتولے شو دا نشته۔
\olpseganchor{OLP-0010-B008}{end}

\olpseganchor{OLP-0010-B009}{start}
\begin{thm}[د رسل پاراډوکس]\ollabel{thm:russells-paradox}
	د \LR{$R = \Setabs{x}{x \notin x}$} په بڼه هېڅ سټ نشته۔
\end{thm}
\olpseganchor{OLP-0010-B009}{end}

\olpseganchor{OLP-0010-B010}{start}
\begin{proof}
که \LR{$R = \Setabs{x}{x \notin x}$} موجود وي، نو \LR{$R \in R$} هله او يوازې هله وي چې \LR{$R \notin R$} وي، او دا يو تناقض دے۔
\end{proof}
\olpseganchor{OLP-0010-B010}{end}

\olpseganchor{OLP-0010-B011}{start}

\begin{explain}
راځئ دا ثبوت لږ ورو وګورو۔ که \LR{$R$} موجود وي، دا پوښتنه معقوله ده چې \LR{$R \in R$} دے که نۀ۔ فرض کړئ چې په رښتيا \LR{$R \in R$}۔ اوس، \LR{$R$}~د هغو ټولو سټونو د سټ په توګه تعريف شوے ؤ چې د خپلو ځانونو غړي نۀ دي۔ نو که \LR{$R \in R$} وي، بيا \LR{$R$} پخپله د \LR{$R$} تعريفوونکے خاصيت نۀ لري۔ خو يوازې هغه سټونه په~\LR{$R$} کښې دي چې دا خاصيت لري؛ ځکه \LR{$R$} د~\LR{$R$} غړے نۀ شي کېدے، يعنې \LR{$R \notin R$}۔ خو \LR{$R$} په يو وخت د~\LR{$R$} غړے هم نۀ شي کېدے او ناغړے هم، نو له تناقض سره مخ يو۔
\olpseganchor{OLP-0010-B011}{end}

\olpseganchor{OLP-0010-B012}{start}
څرنګه چې د \LR{$R \in R$} فرض تناقض ته رسوي، نو \LR{$R \notin R$} لرو۔ خو دا هم تناقض ته رسوي!{} ځکه که \LR{$R \notin R$} وي، نو \LR{$R$} پخپله د \LR{$R$} تعريفوونکے خاصيت لري، او ځکه \LR{$R$} به د هغو ټولو نورو سټونو په شان د \LR{$R$} غړے وي چې د خپل ځان غړي نۀ دي۔ او بيا هم، په يو وخت د~\LR{$R$} ناغړے او غړے دواړه نۀ شي کېدے۔
\end{explain}

\olpseganchor{OLP-0010-B012}{end}

\olpseganchor{OLP-0010-B013}{start}
\begin{digress}
د سټونو داسې نظريه څنګه جوړه کړو چې د رسل له پاراډوکس څخه ځان وساتي، يعنې له دې \emph{متناقضې} دعوې څخه ځان وساتي چې \LR{$R = \Setabs{x}{x \notin x}$} موجود دے؟ د دې دپاره داسې بديهي اصلونه ټاکل پکار دي چې دا بيان کړي چې سټونه کله موجود وي (او کله نۀ وي)، او ډېر دقيق شرطونه ورکړي۔
\olpseganchor{OLP-0010-B013}{end}

\olpseganchor{OLP-0010-B014}{start}
د سټونو هغه نظريه چې په دې باب کښې په لنډه توګه بيان شوې، دا کار نۀ کوي۔ دا \emph{په رښتيا ساده نظريه} ده۔ دا يوازې درته وائي چې سټونه د غړو له مخې د برابرۍ پابند دي، او که ځينې سټونه ولرئ، د هغو اتحاد، تقاطع او داسې نور جوړولے شئ۔ د سټونو نظريه له دې نه په زيات دقت هم جوړېدل ممکن دي۔ 
\end{digress}
\olpseganchor{OLP-0010-B014}{end}

\olpseganchor{OLP-0011-B004}{start}
\olchapter{sfr}{rel}{اړيکې}
\olpseganchor{OLP-0011-B004}{end}

\olfileid{sfr}{rel}{alternate-outline}
\olpseganchor{OLP-0719-B004}{start}
\paragraph{اړيکې}
\olpseganchor{OLP-0719-B004}{end}

\olpseganchor{OLP-0012-B004}{start}
\olfileid{sfr}{rel}{set}
\olsection{اړيکې د سټونو په توګه}
\olpseganchor{OLP-0012-B004}{end}

\olpseganchor{OLP-0012-B005}{start}
\begin{explain}
په \olref[sfr][set][imp]{sec} کښې مو د ځينو مهمو سټونو يادونه کړې وه:
\LR{$\Nat$}، \LR{$\Int$}، \LR{$\Rat$}، \LR{$\Real$}۔ بې شکه به تاسو ته د دې سټونو
د ځينو د غړو ترمنځ په زړه پورې اړيکې يادې وي۔
د بېلګې په توګه، په دې هر سټ باندې يوه پوره معياري \emph{ترتيبي اړيکه}
شته۔ د \emph{هماغه شے کېدو} اړيکه هم شته چې هر څيز ئې له خپل ځان
سره لري او له بل هېڅ څيز سره ئې نۀ لري۔ مونږ به له ډېرو نورو په زړه
پورو اړيکو سره هم مخامخ شو، او ممکنې اړيکې تر دې هم ډېرې دي۔ خو تر
کتنې وړاندې به په دې خبره پيل وکړو چې اړيکې د سټ يو خاص ډول ګڼلے شو۔
\olpseganchor{OLP-0012-B005}{end}

\olpseganchor{OLP-0012-B006}{start}
د دې دپاره له \olref[sfr][set][pai]{sec} څخه دوه خبرې رايادې کړئ۔
لومړے، د \emph{مرتبې جوړې} مفهوم: که \LR{$a$} او \LR{$b$} راکړل شي، نو
\LR{$\tuple{a, b}$} جوړولے شو۔ دلته د غړو ترتيب \emph{خامخا} اهميت لري۔
نو که \LR{$a \neq b$} وي، \LR{$\tuple{a, b} \neq \tuple{b, a}$}۔
(دا له نامرتبو جوړو، يعنې له \LR{$2$}-غړيزو سټونو، سره پرتله کړئ؛ هلته
\LR{$\{a, b\}=\{b, a\}$}۔) دويم، د \emph{کارتيزي حاصل ضرب} مفهوم
راياد کړئ: که \LR{$A$} او \LR{$B$} سټونه وي، نو \LR{$A \times B$} جوړولے شو۔
دا د ټولو هغو جوړو \LR{$\tuple{x, y}$} سټ دے چې \LR{$x \in A$} او \LR{$y \in B$}
وي۔ په تېره بيا، \LR{$A^{2}= A \times A$} د \LR{$A$} د ټولو مرتبو جوړو سټ دے۔
\olpseganchor{OLP-0012-B006}{end}

\olpseganchor{OLP-0012-B007}{start}
اوس به په يو سټ باندې يوه خاصه اړيکه وګورو: د طبيعي عددونو په سټ
\LR{$\Nat$} باندې د \LR{$<$} اړيکه۔ د عددونو د ټولو هغو جوړو \LR{$\tuple{n, m}$}
سټ په پام کښې ونيسئ چې \LR{$n<m$} وي، يعنې:
\[
  R=\Setabs{\tuple{n, m}}{n, m \in \Nat \text{\RL{ او }} n<m}.
\]
د \LR{$n$} له \LR{$m$} څخه د کوچني کېدو او په \LR{$R$} کښې د جوړې \LR{$\tuple{n, m}$}
د غړي کېدو ترمنځ نژدې تړاو شته، يعنې:
\[
      n<m\text{\RL{ هله او يوازې هله چې }}\tuple{n, m} \in R.
\]
په رښتيا، د معلوماتو له هېڅ ضايع کېدو پرته، سټ \LR{$R$} په \LR{$\Nat$} باندې
د \LR{$<$} اړيکه \emph{ګڼلے} شو۔
\olpseganchor{OLP-0012-B007}{end}

\olpseganchor{OLP-0012-B008}{start}
په همدې طريقه، د عددونو ترمنځ د هرې اړيکې دپاره د \LR{$\Nat^{2}$} يو
فرعي سټ جوړولے شو۔ برعکس، که د عددونو د جوړو هر سټ
\LR{$S \subseteq \Nat^{2}$} راکړل شي، نو د عددونو ترمنځ ورسره سمه يوه
اړيکه شته: \LR{$n$} له \LR{$m$} سره دا اړيکه هله او يوازې هله لري چې
\LR{$\tuple{n, m} \in S$} وي۔ دا خبره لاندې تعريف ته جواز ورکوي:
\end{explain}
\olpseganchor{OLP-0012-B008}{end}

\olpseganchor{OLP-0012-B009}{start}
\begin{defn}[دوه ځاييزه اړيکه]
په سټ \LR{$A$} باندې يوه \emph{دوه ځاييزه اړيکه} د \LR{$A^{2}$} يو فرعي سټ
دے۔ که \LR{$R \subseteq A^{2}$} په \LR{$A$} باندې يوه دوه ځاييزه اړيکه وي او
\LR{$x, y \in A$} وي، نو کله کله د \LR{$\tuple{x, y} \in R$} پر ځاے
\LR{$Rxy$} (يا \LR{$xRy$}) ليکو۔
\end{defn}
\olpseganchor{OLP-0012-B009}{end}

\olpseganchor{OLP-0012-B010}{start}
\begin{ex}
  \ollabel{relations}
د طبيعي عددونو د جوړو سټ \LR{$\Nat^{2}$} په يوه دوه بعدي جدول کښې داسې
ليکلے شو:
\[
  \begin{array}{ccccc}
  \mathbf{\tuple{ 0,0 }} & \tuple{ 0,1 } &
    \tuple{ 0,2 } & \tuple{ 0,3 } & \ldots\\
  \tuple{ 1,0 } & \mathbf{\tuple{ 1,1 }} &
    \tuple{ 1,2 } & \tuple{ 1,3 } & \ldots\\
  \tuple{ 2,0 } & \tuple{ 2,1 } &
    \mathbf{\tuple{ 2,2 }} & \tuple{ 2,3 } & \ldots\\
  \tuple{ 3,0 } & \tuple{ 3,1 } & \tuple{ 3,2 } &
    \mathbf{\tuple{ 3,3 }} & \ldots\\
  \vdots & \vdots & \vdots & \vdots & \mathbf{\ddots}
  \end{array}
\]
دلته مو قطر په پنډ ليک ښودلے، ځکه د \LR{$\Nat^2$} هغه فرعي سټ چې پر
قطر پرتې جوړې لري، يعنې:
\[
  \{\tuple{0,0 }, \tuple{ 1,1 }, \tuple{ 2,2 }, \dots\},
  \]
په \LR{$\Nat$} باندې \emph{د عينيت اړيکه} ده۔ (د عينيت اړيکه ډېره کارېږي؛
راځئ د هر سټ \LR{$A$} دپاره
\LR{$\Id{A}=\Setabs{\tuple{ x,x }}{x \in A}$} تعريف کړو۔) د قطر له پاسه
د ټولو پرتو جوړو فرعي سټ، يعنې:
\[
  L = \{\tuple{ 0,1 },\tuple{ 0,2 },\ldots,\tuple{ 1,2 },
  \tuple{ 1,3 }, \dots, \tuple{ 2,3 }, \tuple{ 2,4 },\ldots\},
\]
د \emph{کوچني کېدو} اړيکه ده؛ يعنې \LR{$Lnm$} هله او يوازې هله چې
\LR{$n<m$} وي۔ د قطر لاندې د پرتو جوړو فرعي سټ، يعنې:
\[
  G=\{ \tuple{ 1,0 },\tuple{ 2,0 },\tuple{
    2,1 }, \tuple{ 3,0 },\tuple{ 3,1 },\tuple{ 3,2 }, \dots\},
\]
د \emph{لويېدو} اړيکه ده؛ يعنې \LR{$Gnm$} هله او يوازې هله چې \LR{$n>m$}
وي۔ د \LR{$L$} او \LR{$I$} اتحاد، چې \LR{$K=L\cup I$} ئې بللے شو، د
\emph{کوچني يا برابر کېدو} اړيکه ده: \LR{$Knm$} هله او يوازې هله چې
\LR{$n \le m$} وي۔ همدارنګه، \LR{$H=G \cup I$} د \emph{لوي يا برابر کېدو
اړيکه} ده۔ دا اړيکې، \LR{$L$}، \LR{$G$}، \LR{$K$} او \LR{$H$}، د اړيکو خاص ډولونه دي
چې \emph{ترتيبونه} بللے شي۔ \LR{$L$} او \LR{$G$} دا خاصيت لري چې هېڅ عدد له
خپل ځان سره \LR{$L$} يا \LR{$G$} اړيکه نۀ لري (يعنې، د هر \LR{$n$} دپاره، نۀ \LR{$Lnn$}
شته او نۀ \LR{$Gnn$})۔ د دې خاصيت لرونکې اړيکې \emph{نفي انعکاسي}
بللے شي، او که ترتيبونه هم وي، نو \emph{سخت ترتيبونه} بللے شي۔
\end{ex}
\olpseganchor{OLP-0012-B010}{end}

\olpseganchor{OLP-0012-B011}{start}
\begin{explain}
که څه هم ترتيبونه او عينيت مهمې او طبيعي اړيکې دي، بايد ټينګار وشي
چې زمونږ د تعريف له مخې د \LR{$A^{2}$} \emph{هر} فرعي سټ په \LR{$A$} باندې
يوه اړيکه ده، که هر څومره غيرطبيعي يا په تکلف جوړه شوې ښکاري۔
په تېره بيا، \LR{$\emptyset$} په هر سټ باندې يوه اړيکه ده (\emph{تشه اړيکه}،
چې د غړو هېڅ جوړه ئې نۀ لري)۔ پخپله \LR{$A^{2}$} هم په \LR{$A$} باندې يوه
اړيکه ده (چې هره جوړه ئې لري)، او \emph{ټولشموله اړيکه} بللے شي۔
خو د دې په شان يو څيز هم اړيکه شمېرلے شي:
\LR{$E=\Setabs{\tuple{n, m}}{n>5 \text{\RL{ يا }} m \times n \ge 34}$}۔
\end{explain}
\olpseganchor{OLP-0012-B011}{end}

\olpseganchor{OLP-0012-B012}{start}
\begin{prob}
  په سټ \LR{$\Pow{\{a, b, c\}}$} باندې د \LR{$\subseteq$} اړيکې
  غړي په لړ کښې وليکئ۔
\end{prob}
\olpseganchor{OLP-0012-B012}{end}

\olpseganchor{OLP-0013-B003}{start}
\olfileid{sfr}{rel}{ref}
\olsection{فلسفي غور}
\olpseganchor{OLP-0013-B003}{end}

\olpseganchor{OLP-0013-B004}{start}
په \olref[set]{sec} کښې مو اړيکې د خاصو سټونو په توګه تعريف کړې۔
لږ تم کېږو او يوه لنډه فلسفي پوښتنه کوو: داسې تعريف څۀ \emph{کوي}؟
دا ډېره د شک خبره ده چې مونږ دې ووايو چې د ماوراءالطبيعتي عينيت ځينې
حقيقتونه مو \emph{کشف} کړي دي؛ مثلاً، دا چې په \LR{$\Nat$} باندې ترتيبي
اړيکه \emph{په اصل کښې} هغه سټ
\LR{$R=  \Setabs{\tuple{n,m}}{n, m \in \Nat\text{\RL{ او }} n < m}$}
\emph{وختله} چې په \olref[set]{sec} کښې مو تعريف کړے و۔ د دې خبرې
درې وجهې دا دي۔
\olpseganchor{OLP-0013-B004}{end}

\olpseganchor{OLP-0013-B005}{start}
لومړۍ: په \olref[set][pai]{wienerkuratowski} کښې مو
\LR{$\tuple{a, b} = \{\{a\}, \{a, b\}\}$} تعريف کړے و۔ د دې پر ځاے
دا تعريف وګورئ:
\LR{$\lVert a, b\rVert = \{\{b\}, \{a, b\}\} = \tuple{b,a}$}۔
کله چې \LR{$a \neq b$} وي، نو \LR{$\tuple{a, b} \neq \lVert a,b\rVert$}۔
خو د \LR{$\tuple{a,b}$} پر ځاے مو همداسې \LR{$\lVert a,b\rVert$} هم د مرتبې
جوړې تعريف ګڼلے شو۔ دواړو تعريفونو به يو شان ښۀ کار کړے و۔ نو اوس
د طبيعي عددونو ترتيبي اړيکه ``کېدو'' ته دوه يو شان ښۀ نوماندان لرو:
\begin{align*}
		R &= \Setabs{\tuple{n,m}}{n, m \in \Nat \text{\RL{ او }}n < m}\\
		S &= \Setabs{\lVert n,m\rVert}{n, m \in \Nat \text{\RL{ او }}n < m}.
\end{align*}
ځکه چې د غړو له مخې د برابرۍ د اصل له مخې \LR{$R \neq S$} دي، ښکاره
ده چې \emph{دواړه} په \LR{$\Nat$} باندې له ترتيبي اړيکې سره عينيت نۀ شي
لرلے۔ خو دا به محض خپله خوښه، او له دې امله يو څه د شرم خبره وي، چې
ادعا وکړو چې په واقعيت کښې \LR{$R$}، د \LR{$S$} پر ځاے (يا برعکس)، ترتيبي
اړيکه \emph{ده}۔ (دا د سټونو ته د راکمونې د نظر خلاف د هغه دليل يوه
ډېره ساده بېلګه ده چې بناسراف په \mbox{\LR{\latinfont \citeyear{Benacerraf1965}}} کښې مشهور
کړے و۔ مونږ به څو ځله بيا دې ته راګرځو۔)
\olpseganchor{OLP-0013-B005}{end}

\olpseganchor{OLP-0013-B006}{start}
دويمه: که فکر کوو چې \emph{هره} اړيکه بايد له يو سټ سره عينيت
ولري، نو پخپله د سټ د غړيتوب اړيکه، \LR{$\in$}، هم بايد يو خاص سټ وي۔
په رښتيا، هغه به خامخا سټ \LR{$\Setabs{\tuple{x,y}}{x \in y}$} وي۔
خو آيا دا سټ شته؟ د رسل پاراډوکس ته په کتو، د داسې سټ د موجوديت
ادعا ساده خبره نۀ ده۔ په حقيقت کښې،
د سټونو
نظريه په يوه دقيقه طريقه، د بديهي اصولو د يوې نظريې په توګه، جوړول
ممکن دي، او هغه نظريه به په رښتيا د دې سټ موجوديت رد کړي۔
نو که ځينې اړيکې د سټونو په توګه هم ګڼلے شو، د سټ د غړيتوب اړيکه
به يوه خاصه پېښه وي۔
\olpseganchor{OLP-0013-B006}{end}

\olpseganchor{OLP-0013-B007}{start}
درېيمه: کله چې اړيکې له سټونو سره ``يوې'' ګڼو، ومو وئيل چې ځان ته
به د \LR{$\tuple{x,y} \in R$} پر ځاے د \LR{$Rxy$} ليکلو اجازه ورکوو۔ دا سمه
ده، په دې شرط چې د غړيتوب اړيکه، ``\LR{$\in$}''، د يو پريديکات
\emph{په توګه} وکارول شي۔ خو که فکر کوو چې ``\LR{$\in$}'' د سټ د يو
خاص ډول نښه ده، نو عبارت ``\LR{$\tuple{x,y} \in R$}'' يوازې له درېيو
مفردو اصطلاحاتو جوړ دے چې سټونه ښيي: ``\LR{$\tuple{x,y}$}''، ``\LR{$\in$}''
او ``\LR{$R$}''۔ د نومونو دا ډول لړ، لکه بې معنا توريز تسلسل ``پياله
قلمدان مېز''، د يو بیان د څرګندولو وس نۀ لري۔ نو بيا هم، که ځينې
اړيکې د سټونو په توګه ګڼلے شو، د سټ د غړيتوب اړيکه بايد يوه خاصه
پېښه وي۔ (په دې کښې د فرېګه د \emph{آس} د مفهوم د پاراډوکس يوه
ساده بڼه او هغه مشهور اعتراض سره يوځاے شوي چې وېټګنشټاين يو وخت
پر رسل کړے و۔)
\olpseganchor{OLP-0013-B007}{end}

\olpseganchor{OLP-0013-B008}{start}
نو دې خبرې مونږ کوم ځاے ته ورسولو؟ دا وئيل چې په عددونو باندې
اړيکې سټونه دي، هېڅ \emph{غلط} نۀ دي۔ يوازې بايد په هغې موخې پوه
شو چې دا خبره ورسره کوو۔ مونږ د ماوراءالطبيعتي عينيت يو حقيقت نۀ
بيانوو۔ يوازې دا يادونه کوو چې په ځينو سياقونو کښې (ځينې) اړيکې
خاص سټونه \emph{ګڼلے} شو (او ګڼو به ئې)۔
\olpseganchor{OLP-0013-B008}{end}

\olpseganchor{OLP-0014-B004}{start}
\olfileid{sfr}{rel}{prp}
\olsection{د اړيکو خاص خاصيتونه}
\olpseganchor{OLP-0014-B004}{end}

\olpseganchor{OLP-0014-B005}{start}
\begin{intro}
د اړيکو ځينې ډولونه دومره عام دي چې خاص نومونه ورکړل شوي دي۔
د بېلګې په توګه، \LR{$\le$} او \LR{$\subseteq$} دواړه په خپلو خپلو ساحو کښې
(مثلاً، د \LR{$\le$} دپاره په \LR{$\Nat$} او د \LR{$\subseteq$} دپاره په \LR{$\Pow{A}$}
کښې) په يو شان طريقو اړيکه جوړوي۔ د دې اړيکو د ورته والي او توپير
د دقيق معلومولو دپاره، د هغو خاصو خاصيتونو له مخې ئې وېشو چې
اړيکې ئې لرلے شي۔ څرګندېږي چې د دې ځينو خاصو خاصيتونو (ترکيبونه)
په تېره بيا مهم دي: ترتيبونه او د هم ارزښتۍ اړيکې۔
\end{intro}
\olpseganchor{OLP-0014-B005}{end}

\olpseganchor{OLP-0014-B006}{start}
\begin{defn}[انعکاس]
اړيکه \LR{$R \subseteq A^2$} هله او يوازې هله \emph{انعکاسي} ده چې
د هر \LR{$x \in A$} دپاره \LR{$Rxx$} وي۔
\end{defn}
\olpseganchor{OLP-0014-B006}{end}

\olpseganchor{OLP-0014-B007}{start}
\begin{defn}[انتقال]
اړيکه \LR{$R \subseteq A^2$} هله او يوازې هله \emph{انتقالي} ده چې
کله هم \LR{$Rxy$} او \LR{$Ryz$} وي، نو \LR{$Rxz$} هم وي۔
\end{defn}
\olpseganchor{OLP-0014-B007}{end}

\olpseganchor{OLP-0014-B008}{start}
\begin{defn}[تناظر]
اړيکه \LR{$R \subseteq A^2$} هله او يوازې هله \emph{متناظره} ده چې
کله هم \LR{$Rxy$} وي، نو \LR{$Ryx$} هم وي۔
\end{defn}
\olpseganchor{OLP-0014-B008}{end}

\olpseganchor{OLP-0014-B009}{start}
\begin{defn}[ضد تناظر]
اړيکه \LR{$R \subseteq A^2$} هله او يوازې هله \emph{ضد متناظره} ده چې
کله هم \LR{$Rxy$} او \LR{$Ryx$} دواړه وي، نو \LR{$x=y$} وي (په نورو ټکو: که
\LR{$x\neq y$} وي، نو يا \LR{$\lnot Rxy$} وي او يا \LR{$\lnot Ryx$})۔
\end{defn}
\olpseganchor{OLP-0014-B009}{end}

\olpseganchor{OLP-0014-B010}{start}
\begin{explain}
په متناظره اړيکه کښې \LR{$Rxy$} او \LR{$Ryx$} تل يوځاے وي، يا يو هم نۀ وي۔
په ضد متناظره اړيکه کښې \LR{$Rxy$} او \LR{$Ryx$} يوازې هله يوځاے کېدلے شي
چې \LR{$x = y$} وي۔ پام وکړئ چې دا خبره د \LR{$x = y$} په وخت د \LR{$Rxy$} او
\LR{$Ryx$} موجوديت \emph{نۀ لازموي}؛ يوازې دا چې مخه ئې نۀ نيول کېږي۔
نو ضد متناظره اړيکه انعکاسي کېدلے شي، خو هره ضد متناظره اړيکه
انعکاسي نۀ وي۔ دا هم په پام کښې وساتئ چې ضد متناظره کېدل او محض
متناظره نۀ کېدل بېل شرطونه دي۔ په رښتيا، يوه اړيکه په يوه وخت
متناظره او ضد متناظره دواړه کېدلے شي (مثلاً، د عينيت اړيکه داسې ده)۔
\end{explain}
\olpseganchor{OLP-0014-B010}{end}

\olpseganchor{OLP-0014-B011}{start}
\begin{defn}[تړلتيا]
اړيکه \LR{$R \subseteq A^2$} \emph{تړلې} ده که د ټولو \LR{$x,y\in A$}
دپاره، که \LR{$x \neq y$} وي، نو يا \LR{$Rxy$} وي او يا \LR{$Ryx$}۔
\end{defn}
\olpseganchor{OLP-0014-B011}{end}

\olpseganchor{OLP-0014-B012}{start}
\begin{prob}
د داسې اړيکو بېلګې ورکړئ چې (الف) انعکاسي او متناظرې وي، خو انتقالي
نۀ وي؛ (ب) انعکاسي او ضد متناظرې وي؛ (ج) ضد متناظرې او انتقالي وي،
خو انعکاسي نۀ وي؛ او (د) انعکاسي، متناظرې او انتقالي وي۔ پر عددونو
يا سټونو باندې اړيکې مۀ کاروئ۔
\end{prob}
\olpseganchor{OLP-0014-B012}{end}

\olpseganchor{OLP-0014-B013}{start}
\begin{defn}[د انعکاس نفي]
اړيکه \LR{$R \subseteq A^2$} \emph{نفي انعکاسي} بللے شي که د هر
\LR{$x \in A$} دپاره \LR{$Rxx$} نۀ وي۔
\end{defn}
\olpseganchor{OLP-0014-B013}{end}

\olpseganchor{OLP-0014-B014}{start}
\begin{defn}[نامتناظرتيا]
اړيکه \LR{$R \subseteq A^2$} \emph{نامتناظره} بللے شي که د هېڅ جوړې
\LR{$x,y\in A$} دپاره \LR{$Rxy$} او \LR{$Ryx$} دواړه ونۀ لرو۔
\end{defn}
\olpseganchor{OLP-0014-B014}{end}

\olpseganchor{OLP-0014-B015}{start}
پام وکړئ: که \LR{$A \neq \emptyset$} وي، نو په \LR{$A$} باندې هېڅ نفي انعکاسي
اړيکه انعکاسي نۀ ده، او په \LR{$A$} باندې هره نامتناظره اړيکه ضد متناظره
هم ده۔ خو داسې \LR{$R \subseteq A^2$} شته چې نۀ انعکاسي دي او نۀ نفي
انعکاسي، او داسې ضد متناظرې اړيکې شته چې نامتناظرې نۀ دي۔
\olpseganchor{OLP-0014-B015}{end}

\olpseganchor{OLP-0015-B004}{start}
\olfileid{sfr}{rel}{eqv}
\olpseganchor{OLP-0015-B004}{end}

\olpseganchor{OLP-0015-B005}{start}
\olsection{د هم ارزښتۍ اړيکې}
\olpseganchor{OLP-0015-B005}{end}

\olpseganchor{OLP-0015-B006}{start}
په يوه سټ باندې د عينيت اړيکه انعکاسي، متناظره او انتقالي ده۔
هغه اړيکې \LR{$R$} چې دا درې واړه خاصيتونه لري، ډېرې عامې دي۔
\olpseganchor{OLP-0015-B006}{end}

\olpseganchor{OLP-0015-B007}{start}
\begin{defn}[د هم ارزښتۍ اړيکه]
هغه اړيکه \LR{$R \subseteq A^2$} چې انعکاسي، متناظره او انتقالي وي،
\emph{د هم ارزښتۍ اړيکه} بللے شي۔ د \LR{$A$} غړي، \LR{$x$} او
\LR{$y$}، د \LR{$R$} له مخې \emph{هم ارزښته} بللے شي که \LR{$Rxy$} وي۔
\end{defn}
\olpseganchor{OLP-0015-B007}{end}

\olpseganchor{OLP-0015-B008}{start}
د هم ارزښتۍ اړيکې د \emph{هم ارزښتۍ د ټولګي} مفهوم راپيدا کوي۔
د هم ارزښتۍ يوه اړيکه ساحه په بېلو برخو ``ټوټه کوي''۔ په هره برخه
کښې ټول څيزونه يو له بل سره اړيکه لري؛ او د بېلو برخو څيزونه يو
له بل سره اړيکه نۀ لري۔ کله کله د همدې برخو په اړه \emph{نېغ په نېغه}
خبرې کول ګټور وي۔ د دې دپاره يو تعريف راوړو:
\olpseganchor{OLP-0015-B008}{end}

\olpseganchor{OLP-0015-B009}{start}
\begin{defn}\ollabel{def:equivalenceclass}
فرض کړئ \LR{$R \subseteq A^2$} د هم ارزښتۍ اړيکه ده۔ د هر \LR{$x \in A$}
دپاره، په \LR{$A$} کښې د \LR{$x$} \emph{د هم ارزښتۍ ټولګے} دا سټ دے:
\LR{$\equivrep{x}{R} = \Setabs{y \in A}{Rxy}$}۔ د \LR{$R$} له مخې د \LR{$A$}
\emph{خارج قسمت} دا دے:
\LR{$\equivclass{A}{R} = \Setabs{\equivrep{x}{R}}{x \in A}$}، يعنې
د هم ارزښتۍ د دې ټولګيو سټ۔
\end{defn}
\olpseganchor{OLP-0015-B009}{end}

\olpseganchor{OLP-0015-B010}{start}
راتلونکې نتيجه د هم ارزښتۍ د ټولګي د تعريف سموالى ښيي: ثابتوي چې
د هم ارزښتۍ ټولګي په رښتيا د \LR{$A$} وېشونکې برخې دي۔
\olpseganchor{OLP-0015-B010}{end}

\olpseganchor{OLP-0015-B011}{start}
\begin{prop}
که \LR{$R \subseteq A^2$} د هم ارزښتۍ اړيکه وي، نو \LR{$Rxy$} هله او يوازې
هله چې \LR{$\equivrep{x}{R} = \equivrep{y}{R}$} وي۔
\end{prop}
\olpseganchor{OLP-0015-B011}{end}

\olpseganchor{OLP-0015-B012}{start}
\begin{proof}
د کيڼې خوا نه ښۍ خوا ته د لوري دپاره، فرض کړئ \LR{$Rxy$} او
\LR{$z \in \equivrep{x}{R}$} وي۔ نو د تعريف له مخې \LR{$Rxz$}۔ ځکه چې \LR{$R$}
د هم ارزښتۍ اړيکه ده، \LR{$Ryz$}۔ (په تفصيل سره: له \LR{$Rxy$} او د \LR{$R$}
له تناظر څخه \LR{$Ryx$} لرو؛ او له \LR{$Rxz$} او د \LR{$R$} له انتقالي والي څخه
\LR{$Ryz$} لرو۔) نو \LR{$z \in \equivrep{y}{R}$}۔ په عمومي توګه،
\LR{$\equivrep{x}{R} \subseteq \equivrep{y}{R}$}۔ خو بالکل په همدې
ډول، \LR{$\equivrep{y}{R} \subseteq \equivrep{x}{R}$}۔ نو د غړو له
مخې د برابرۍ د اصل له مخې، \LR{$\equivrep{x}{R} = \equivrep{y}{R}$}۔
\olpseganchor{OLP-0015-B012}{end}

\olpseganchor{OLP-0015-B013}{start}
د ښۍ خوا نه کيڼې خوا ته د لوري دپاره، فرض کړئ
\LR{$\equivrep{x}{R} = \equivrep{y}{R}$}۔ ځکه چې \LR{$R$} انعکاسي ده،
\LR{$Ryy$}، نو \LR{$y \in \equivrep{y}{R}$}۔ له دې فرضيې څخه چې
\LR{$\equivrep{x}{R} = \equivrep{y}{R}$}، دا هم لرو چې
\LR{$y \in \equivrep{x}{R}$}۔ نو \LR{$Rxy$}۔
\end{proof}
\olpseganchor{OLP-0015-B013}{end}

\olpseganchor{OLP-0015-B014}{start}
\begin{ex}
د هم ارزښتۍ د اړيکو يوه ښۀ بېلګه د باقي له حسابه راځي۔ د هر \LR{$a$}،
\LR{$b$} او \LR{$n \in \PosInt$} دپاره، وايو چې \LR{$a \equiv_n b$} هله او
يوازې هله چې د \LR{$a$} په \LR{$n$} د وېش باقي د \LR{$b$} په \LR{$n$} د وېش له باقي
سره يو شان وي۔ (په يو څه زياته نښيزه ژبه: \LR{$a \equiv_n b$} هله او
يوازې هله چې د کوم \LR{$k \in \Int$} دپاره \LR{$a - b = kn$} وي۔)
اوس د هر \LR{$n$} دپاره \LR{$\equiv_n$} د هم ارزښتۍ اړيکه ده۔ او د
\LR{$\equiv_n$} له مخې جوړ شوي د هم ارزښتۍ ټولګي په شمېر پوره \LR{$n$}
بېل ټولګي دي؛ يعنې \LR{$\equivclass{\Nat}{\equiv_n}$} په شمېر \LR{$n$}
غړي لري۔ دا ټولګي دي: د هغو عددونو سټ چې بې له باقي پر
\LR{$n$} وېشل کېږي، يعنې \LR{$\equivrep{0}{\equiv_n}$}؛ د هغو عددونو سټ
چې پر \LR{$n$} د وېش باقي ئې \LR{$1$} وي، يعنې \LR{$\equivrep{1}{\equiv_n}$}؛
\ldots؛ او د هغو عددونو سټ چې پر \LR{$n$} د وېش باقي ئې \LR{$n-1$} وي،
يعنې \LR{$\equivrep{n-1}{\equiv_n}$}۔
\end{ex}
\olpseganchor{OLP-0015-B014}{end}

\olpseganchor{OLP-0015-B015}{start}
\begin{prob}
وښيئ چې د هر \LR{$n \in \PosInt$} دپاره \LR{$\equiv_n$} د هم ارزښتۍ
اړيکه ده، او \LR{$\equivclass{\Nat}{\equiv_n}$} په شمېر پوره \LR{$n$}
غړي لري۔
\end{prob}
\olpseganchor{OLP-0015-B015}{end}

\olpseganchor{OLP-0016-B004}{start}
\olfileid{sfr}{rel}{ord}
\olsection{ترتيبونه}
\olpseganchor{OLP-0016-B004}{end}

\olpseganchor{OLP-0016-B005}{start}
\begin{explain}
په ډېرو پرتلو کښې ځينې څيزونه، له يوه پلوه، تر نورو ``کوچني''،
ورسره ``برابر'' يا ترې ``لوي'' بولو۔ په دې کښې د \emph{ترتيب}
اړيکې شاملې دي۔ خو د ترتيبي اړيکو بېل ډولونه شته۔ مثلاً، ځينې
لازموي چې هر دوه څيزونه د پرتلې وړ وي، خو نورې داسې نۀ دي۔ ځينې
عينيت هم رانغاړي (لکه \LR{$\le$})، او ځينې ئې باسي (لکه \LR{$<$})۔ دلته د
ډولونو وېش لرل راسره مرسته کوي۔
\end{explain}
\olpseganchor{OLP-0016-B005}{end}

\olpseganchor{OLP-0016-B006}{start}
\begin{defn}[مقدماتي ترتيب]
هغه اړيکه چې انعکاسي او انتقالي دواړه وي، \emph{مقدماتي ترتيب}
بللے شي۔
\end{defn}
\olpseganchor{OLP-0016-B006}{end}

\olpseganchor{OLP-0016-B007}{start}
\begin{defn}[جزوي ترتيب]
هغه مقدماتي ترتيب چې ضد متناظر هم وي، \emph{جزوي ترتيب} بللے شي۔
\end{defn}
\olpseganchor{OLP-0016-B007}{end}

\olpseganchor{OLP-0016-B008}{start}
\begin{defn}[خطي ترتيب]\ollabel{def:linearorder}
هغه جزوي ترتيب چې تړلے هم وي، \emph{کلي ترتيب} يا \emph{خطي ترتيب}
بللے شي۔
\end{defn}
\olpseganchor{OLP-0016-B008}{end}

\olpseganchor{OLP-0016-B009}{start}
\begin{ex}
هر خطي ترتيب جزوي ترتيب هم دے، او هر جزوي ترتيب مقدماتي ترتيب
هم دے؛ خو برعکس خبرې سمې نۀ دي۔ په \LR{$A$} باندې ټولشموله اړيکه
مقدماتي ترتيب دے، ځکه انعکاسي او انتقالي ده۔ خو که \LR{$A$} له يوه
غړے څخه زيات ولري، نو ټولشموله اړيکه ضد متناظره نۀ ده،
او له دې امله جزوي ترتيب هم نۀ دے۔
\end{ex}
\olpseganchor{OLP-0016-B009}{end}

\olpseganchor{OLP-0016-B010}{start}
\begin{ex}
په \LR{$\Bin^*$} باندې د \emph{تر بل نۀ اوږدېدو} اړيکه
\LR{$\preccurlyeq$} وګورئ: \LR{$x \preccurlyeq y$} هله او يوازې هله چې
\LR{$\len{x} \le \len{y}$} وي۔ دا مقدماتي ترتيب دے (انعکاسي او انتقالي
دے)، او تړلے هم دے؛ خو جزوي ترتيب نۀ دے، ځکه ضد متناظر نۀ دے۔
مثلاً، \LR{$01 \preccurlyeq 10$} او \LR{$10 \preccurlyeq 01$}، خو
\LR{$01 \neq 10$}۔
\end{ex}
\olpseganchor{OLP-0016-B010}{end}

\olpseganchor{OLP-0016-B011}{start}
\begin{ex}
د سټونو په يوه سټ باندې د \LR{$\subseteq$} اړيکه يو مهم جزوي ترتيب دے۔
دا په عمومي توګه خطي ترتيب نۀ دے، ځکه که \LR{$a \neq b$} وي او
\LR{$\Pow{\{a, b\}} = \{\emptyset, \{a\}, \{b\}, \{a,b\}\}$}
وګورو، نو وينو چې \LR{$\{a\} \nsubseteq \{b\}$}،
\LR{$\{a\} \neq \{b\}$} او \LR{$\{b\} \nsubseteq \{a\}$}۔
\end{ex}
\olpseganchor{OLP-0016-B011}{end}

\olpseganchor{OLP-0016-B012}{start}
\begin{ex}
د \emph{بې له باقي وېشل کېدو} اړيکه داسې جزوي ترتيب راکوي چې خطي
ترتيب نۀ دے۔ د صحيح عددونو \LR{$n$} او \LR{$m$} دپاره، \LR{$n \mid m$} په دې
معنا ليکو چې \LR{$n$}، \LR{$m$} (بې له باقي) وېشي؛ يعنې هله او يوازې هله
چې داسې صحيح عدد \LR{$k$} شته چې \LR{$m = kn$} وي۔ په \LR{$\Nat$} باندې دا جزوي
ترتيب دے، خو خطي ترتيب نۀ دے: مثلاً، \LR{$2 \nmid 3$} او \LR{$3 \nmid 2$}
هم۔ که وېشمنتيا په \LR{$\Int$} باندې اړيکه وګڼو، يوازې مقدماتي ترتيب
دے، ځکه ضد متناظره نۀ ده: \LR{$1 \mid -1$} او \LR{$-1 \mid 1$}، خو
\LR{$1 \neq -1$}۔
\end{ex}
\olpseganchor{OLP-0016-B012}{end}

\olpseganchor{OLP-0016-B013}{start}
\begin{ex}
د لړيو په سټ \LR{$A^*$} باندې د \emph{غځونې} اړيکه دا ده:
\LR{$s \sqsubseteq s'$} هله او يوازې هله چې \LR{$s = \emptyseq$} (تشه لړۍ)
وي، \LR{$s = s'$} وي، يا \LR{$s = \tuple{s_1, \dots, s_n}$} او
\LR{$s' = \tuple{s_1, \dots, s_n, s_{n+1}, \dots, s_m}$} وي۔ که
\LR{$s \sqsubseteq s'$} وي، دا هم وايو چې \LR{$s$} د \LR{$s'$} \emph{پيلنۍ برخه}
ده۔ په \LR{$A^*$} باندې د غځونې اړيکه جزوي ترتيب دے، خو خطي ترتيب نۀ
دے؛ مثلاً، که \LR{$a \neq b$} وي، نو \LR{$ab \not\sqsubseteq ba$} او
\LR{$ba \not\sqsubseteq ab$}۔
\end{ex}
\olpseganchor{OLP-0016-B013}{end}

\olpseganchor{OLP-0016-B014}{start}
\begin{defn}[سخت ترتيب]
\emph{سخت ترتيب} هغه اړيکه ده چې نفي انعکاسي، نامتناظره او انتقالي وي۔
\end{defn}
\olpseganchor{OLP-0016-B014}{end}

\olpseganchor{OLP-0016-B015}{start}
\begin{defn}[سخت خطي ترتيب]\ollabel{def:strictlinearorder}
هغه سخت ترتيب چې تړلے هم وي، \emph{سخت کلي ترتيب} يا
\emph{سخت خطي ترتيب} بللے شي۔
\end{defn}
\olpseganchor{OLP-0016-B015}{end}

\olpseganchor{OLP-0016-B016}{start}
\begin{ex}
\LR{$\le$} هغه خطي ترتيب دے چې له سخت خطي ترتيب \LR{$<$} سره سمون خوري۔
\LR{$\subseteq$} هغه جزوي ترتيب دے چې له سخت ترتيب \LR{$\subsetneq$} سره
سمون خوري۔
\end{ex}
\olpseganchor{OLP-0016-B016}{end}

\olpseganchor{OLP-0016-B017}{start}
په \LR{$A$} باندې هر سخت ترتيب \LR{$R$} د قطر \LR{$\Id{A}$} په ورزياتولو، يعنې
د ټولو جوړو \LR{$\tuple{x, x}$} په ورزياتولو، جزوي ترتيب ګرځولے شو۔
(دې ته د \LR{$R$} \emph{انعکاسي بندښت} وايي۔) برعکس، له جزوي ترتيبه
په پيل کولو، د \LR{$\Id{A}$} په لرې کولو سخت ترتيب ترلاسه کولے شو۔
لاندې دوه نتيجې دا خبره دقيقه کوي۔
\olpseganchor{OLP-0016-B017}{end}

\olpseganchor{OLP-0016-B018}{start}
\begin{prop}\ollabel{prop:stricttopartial}
که \LR{$R$} په \LR{$A$} باندې سخت ترتيب وي، نو \LR{$R^+ = R \cup \Id{A}$}
جزوي ترتيب دے۔ سربېره پر دې، که \LR{$R$} سخت خطي ترتيب وي، نو \LR{$R^+$}
خطي ترتيب دے۔
\end{prop}
\olpseganchor{OLP-0016-B018}{end}

\olpseganchor{OLP-0016-B019}{start}
\begin{proof}
فرض کړئ \LR{$R$} سخت ترتيب دے؛ يعنې \LR{$R \subseteq A^2$} او \LR{$R$} نفي
انعکاسي، نامتناظر او انتقالي دے۔ \LR{$R^+ = R \cup \Id{A}$} کېږدئ۔
بايد وښيو چې \LR{$R^+$} انعکاسي، ضد متناظر او انتقالي دے۔
\olpseganchor{OLP-0016-B019}{end}

\olpseganchor{OLP-0016-B020}{start}
ښکاره ده چې \LR{$R^+$} انعکاسي دے، ځکه د هر \LR{$x \in A$} دپاره
\LR{$\tuple{x, x} \in \Id{A} \subseteq R^+$}۔
\olpseganchor{OLP-0016-B020}{end}

\olpseganchor{OLP-0016-B021}{start}
د دې د ښودلو دپاره چې \LR{$R^+$} ضد متناظر دے، د تناقض له لارې د ثبوت
دپاره فرض کړئ \LR{$R^+xy$} او \LR{$R^+yx$}، خو \LR{$x \neq y$}۔ ځکه چې
\LR{$\tuple{x,y} \in R \cup \Id{A}$}، خو
\LR{$\tuple{x, y} \notin \Id{A}$}، نو خامخا \LR{$\tuple{x, y} \in R$}
لرو، يعنې \LR{$Rxy$}۔ همدارنګه، \LR{$Ryx$}۔ خو دا له دې فرضيې سره تناقض لري
چې \LR{$R$} نامتناظر دے۔
\olpseganchor{OLP-0016-B021}{end}

\olpseganchor{OLP-0016-B022}{start}
د انتقالي والي د ثابتولو دپاره، فرض کړئ \LR{$R^+xy$} او \LR{$R^+yz$}۔
که \LR{$\tuple{x, y} \in R$} او \LR{$\tuple{y,z} \in R$} دواړه وي، نو
\LR{$\tuple{x, z} \in R$}، ځکه \LR{$R$} انتقالي دے۔ که داسې نۀ وي، نو يا
\LR{$\tuple{x, y} \in \Id{A}$}، يعنې \LR{$x = y$}، يا
\LR{$\tuple{y, z} \in \Id{A}$}، يعنې \LR{$y = z$}۔ په لومړي حالت کښې د
فرضيې له مخې \LR{$R^+yz$} لرو، او \LR{$x = y$}، نو \LR{$R^+xz$}۔ په دويم
حالت کښې هم همداسې ده۔ په هر حالت کښې \LR{$R^+xz$}، نو \LR{$R^+$} انتقالي
هم دے۔
\olpseganchor{OLP-0016-B022}{end}

\olpseganchor{OLP-0016-B023}{start}
د ``سربېره پر دې'' برخې په اړه، فرض کړئ \LR{$R$} تړلے هم دے۔ نو د ټولو
\LR{$x \neq y$} دپاره يا \LR{$Rxy$} يا \LR{$Ryx$}؛ يعنې يا
\LR{$\tuple{x, y} \in R$} يا \LR{$\tuple{y, x} \in R$}۔ ځکه چې
\LR{$R \subseteq R^+$}، دا خبره د \LR{$R^+$} دپاره هم سمه پاتې کېږي۔ نو
\LR{$R^+$} تړلے هم دے۔
\end{proof}
\olpseganchor{OLP-0016-B023}{end}

\olpseganchor{OLP-0016-B024}{start}
\begin{prop}\ollabel{prop:partialtostrict}
که \LR{$R$} په \LR{$A$} باندې جزوي ترتيب وي، نو
\LR{$R^- = R \setminus \Id{A}$} سخت ترتيب دے۔ سربېره پر دې، که \LR{$R$}
خطي ترتيب وي، نو \LR{$R^-$} سخت خطي ترتيب دے۔
\end{prop}
\olpseganchor{OLP-0016-B024}{end}

\olpseganchor{OLP-0016-B025}{start}
\begin{proof}
دا د تمرين په توګه پرېښودل شوے دے۔
\end{proof}
\olpseganchor{OLP-0016-B025}{end}

\olpseganchor{OLP-0016-B026}{start}
\begin{prob}
د \olref[sfr][rel][ord]{prop:partialtostrict} ثبوت ورکړئ۔
\end{prob}
\olpseganchor{OLP-0016-B026}{end}

\olpseganchor{OLP-0016-B027}{start}
لاندې ساده نتيجه ثابتوي چې سخت خطي ترتيبونه د غړو له مخې د
برابرۍ په شان يو خاصيت پوره کوي:
\olpseganchor{OLP-0016-B027}{end}

\olpseganchor{OLP-0016-B028}{start}
\begin{prop}\ollabel{prop:extensionality-strictlinearorders}
که \LR{$<$} په \LR{$A$} باندې سخت خطي ترتيب وي، نو:
\[
  (\forall a, b \in A)((\forall x \in A)(x < a \liff x < b) \lif a = b).
\]
\end{prop}
\olpseganchor{OLP-0016-B028}{end}

\olpseganchor{OLP-0016-B029}{start}
\begin{proof}
فرض کړئ \LR{$(\forall x \in A)(x < a \liff x < b)$}۔ که \LR{$a < b$} وي،
نو \LR{$a < a$}، چې د \LR{$<$} له نفي انعکاسي والي سره تناقض لري؛ نو
\LR{$a \nless b$}۔ بالکل په همدې ډول \LR{$b \nless a$}۔ نو \LR{$a = b$}، ځکه
\LR{$<$} تړلے دے۔
\end{proof}
\olpseganchor{OLP-0016-B029}{end}

\olpseganchor{OLP-0017-B004}{start}
\olfileid{sfr}{rel}{grp}
\olsection{ګرافونه}
\olpseganchor{OLP-0017-B004}{end}

\olpseganchor{OLP-0017-B005}{start}
\emph{ګراف} داسې شکل دے چې ټکي—چې ``غوټې'' يا ``څوکې'' (د
``څوکه'' جمع) بللے شي—پکښې د څنډو په ذريعه سره نښتي وي۔ ګرافونه
په منفصل رياضي او کمپيوټرپوهنه کښې هر ځاے کارېدونکې وسيلې دي۔
د اړيکو او جوړښتونو د ښودلو او انځورولو دپاره ډېر ګټور دي؛ له
محسوسو څيزونو، لکه د بېلو ډولونو له شبکو، نيولې تر مجردو جوړښتونو،
لکه د پرېکړو تر ممکنو پايلو پورې۔ په علمي ليکنو کښې د ګرافونو ډېر
بېل ډولونه شته۔ مثلاً، په دې سره توپير کوي چې څنډې لوری لري که نه،
نوم نښې لري که نه، آيا له يوې غوټې بېرته هماغې غوټې ته څنډې کېدلے
شي، د هماغو غوټو ترمنځ څو څنډې کېدلے شي، او داسې نور۔
\emph{لوری لرونکي ګرافونه} له اړيکو سره يو خاص تړاو لري۔
\olpseganchor{OLP-0017-B005}{end}

\olpseganchor{OLP-0017-B006}{start}
\begin{defn}[لوری لرونکے ګراف]
\emph{لوری لرونکے ګراف} \LR{$G = \tuple{V, E}$} د \emph{څوکو} يو سټ
\LR{$V$} او د \emph{څنډو} يو سټ \LR{$E \subseteq V^2$} دے۔
\end{defn}
\olpseganchor{OLP-0017-B006}{end}

\olpseganchor{OLP-0017-B007}{start}
\begin{explain}
زمونږ د تعريف له مخې، ګراف يوازې يو سټ او په همدغه سټ باندې يوه
اړيکه ده۔ البته، د ګرافونو په اړه د خبرو پر وخت، په شکل کښې د
ښودلو تمه ئې طبيعي ده: ګراف داسې رسمولے شو چې دوه څوکې \LR{$v_1$} او
\LR{$v_2$} په يو غشي هله او يوازې هله سره ونښلوو چې
\LR{$\tuple{v_1, v_2} \in E$} وي۔ د يوازې اړيکې او ګراف ترمنځ توپير
يوازې دا دے چې ګراف د څوکو سټ ټاکي؛ يعنې ګراف ګوښې څوکې هم
لرلے شي۔ خو مهمه خبره دا ده چې په سټ \LR{$X$} باندې هره اړيکه \LR{$R$} د
يو لوری لرونکي ګراف \LR{$\tuple{X, R}$} په توګه کتلے شو۔ برعکس، يو
لوری لرونکے ګراف \LR{$\tuple{V, E}$} د اړيکې \LR{$E \subseteq V^2$} په
توګه کتلے شو چې سټ \LR{$V$} پکښې په څرګنده ټاکل شوے وي۔
\end{explain}
\olpseganchor{OLP-0017-B007}{end}

\olpseganchor{OLP-0017-B008}{start}
\begin{ex}
هغه ګراف \LR{$\tuple{V, E}$} چې
\LR{$V = \{1, 2, 3, 4\}$} او
\LR{$E = \{\tuple{1,1}, \allowbreak \tuple{1, 2}, \allowbreak \tuple{1, 3},
\allowbreak \tuple{2, 3}\}$} وي، داسې ښکاري:
\begin{align*}
& \begin{tikzpicture}[->,node distance=2cm]
  \node[draw,circle] (A) {\LR{$1$}};
  \node[draw,circle] (B) [right of=A] {\LR{$2$}};
  \node[draw,circle] (C) [below of=B] {\LR{$3$}};
  \node[draw,circle] (D) [right of=B] {\LR{$4$}};
  \draw (A) to [loop above]  (A);
  \draw (A) to  (B);
  \draw (A) to  (C);
  \draw (B) to  (C);
  \end{tikzpicture}
\intertext{\RL{دا له هغه ګراف \LR{$\tuple{V', E}$} څخه بېل دے چې
  \LR{$V' = \{1, 2, 3\}$} وي، او داسې ښکاري:}}
& \begin{tikzpicture}[->,node distance=2cm]
  \node[draw,circle] (A) {\LR{$1$}};
  \node[draw,circle] (B) [right of=A] {\LR{$2$}};
  \node[draw,circle] (C) [below of=B] {\LR{$3$}};
  \draw (A) to [loop above]  (A);
  \draw (A) to  (B);
  \draw (A) to  (C);
  \draw (B) to  (C);
  \end{tikzpicture}
\end{align*}
\end{ex}
\olpseganchor{OLP-0017-B008}{end}

\olpseganchor{OLP-0017-B009}{start}
\begin{prob}
  په سټ \LR{$\{1, 2, 3, 4\}$} باندې د کوچني يا برابر کېدو اړيکه
  \LR{$\le$} د ګراف په توګه وګورئ او ورسره سم شکل رسم کړئ۔
\end{prob}
\olpseganchor{OLP-0017-B009}{end}

\olpseganchor{OLP-0018-B004}{start}
\olfileid{sfr}{rel}{tre}
\olsection{ونې}
\olpseganchor{OLP-0018-B004}{end}

\olpseganchor{OLP-0018-B005}{start}
د جزوي ترتيب يو خاص ډول چې د منطق په ټولو برخو کښې مهم رول لري،
\emph{ونه} ده۔ متناهي ونې د منطق په لومړنيو برخو کښې راځي:
مثلاً، فارمولونه په نحوي ونه کښې د خپلو برخو په بېلولو پوهېدلے
شو، او په ډېرو اشتقاق نظامونو کښې اشتقاقونه هم د
متناهي ونو بڼه خپلوي۔
%
نامتناهي ونې د بياني منطق او د لومړۍ درجې منطق د بشپړتيا د قضيو
په ثبوت کښې لا راڅرګندېږي، او په ټول رياضي منطق کښې کارول کېږي۔
\olpseganchor{OLP-0018-B005}{end}

\olpseganchor{OLP-0018-B006}{start}
د سټونو د نظريې د ونې مفهوم د ګرافونو په نظريه کښې د ونې له
مفهوم سره نژدې تړاو لري۔ دا د يوې (متناهي) ونې انځور دے:
\olpseganchor{OLP-0018-B006}{end}

\olpseganchor{OLP-0018-B007}{start}
\begin{center}
\begin{tikzpicture}[nodes={draw, circle}, -]
\node{\LR{$r$}} [grow'=up]
    child { node {\LR{$a$}}
        child { node {\LR{$c$}} }
        child { node {\LR{$d$}} }
        child { node {\LR{$e$}} }
    }
    child { node {\LR{$b$}} };
\end{tikzpicture}
\end{center}
\olpseganchor{OLP-0018-B007}{end}

\olpseganchor{OLP-0018-B008}{start}
تر ټولو لاندې غوټه \LR{$r$} ريښه ده۔ له \LR{$r$} پرته د هرې غوټې سمدستي
لاندې پوره يوه مور غوټه شته۔ که له غوټې \LR{$x$} څخه د څنډو په تعقيب
پورته تلو سره غوټې \LR{$y$} ته رسېدلے شو، نو د \LR{$x$} او \LR{$y$} ترمنځ دا
اړيکه داسې ګڼلے شو چې \LR{$x$} د \LR{$y$} يو \emph{نيکۀ} دے۔
\olpseganchor{OLP-0018-B008}{end}

\olpseganchor{OLP-0018-B009}{start}
په يوه ونه کښې د نيکۀ اړيکه سخت جزوي ترتيب دے۔ دا د سټونو د
نظريې تعريف ته لاره هواروي۔ د هغه د بيانولو دپاره دوه مفهومونو ته
اړتيا لرو۔ په سټ \LR{$A$} کښې، چې د \LR{$\le$} په ذريعه جزوي مرتب شوے وي،
\emph{تر ټولو کوچنے غړے} هغه غړے \LR{$x \in A$} دے چې د
هر \LR{$y \in A$} دپاره \LR{$x \le y$} ولرو۔ يو سټ د \LR{$\le$} په ذريعه
\emph{ښۀ مرتب} دے که هر ناتش فرعي سټ ئې تر ټولو کوچنے غړے ولري۔
\olpseganchor{OLP-0018-B009}{end}

\olpseganchor{OLP-0018-B010}{start}
\begin{defn}[ونه]
\emph{ونه} داسې جوړه \LR{$T = \tuple{A, \le}$} ده چې \LR{$A$} يو سټ وي،
او \LR{$\le$} په \LR{$A$} باندې داسې جزوي ترتيب وي چې يوازينے تر ټولو
کوچنے غړے \LR{$r \in A$} ولري (چې \emph{ريښه} بللے شي)، او د هر
\LR{$x \in A$} دپاره سټ \LR{$\Setabs{y}{y \le x}$} د \LR{$\le$} په ذريعه
ښۀ مرتب وي۔
\end{defn}
\olpseganchor{OLP-0018-B010}{end}

\olpseganchor{OLP-0018-B011}{start}
\begin{defn}[ځايناستي]
فرض کړئ \LR{$T = \tuple{A, \le}$} يوه ونه ده۔ که \LR{$x,y \in A$}،
\LR{$x < y$}، او داسې \LR{$z \in A$} نۀ وي چې \LR{$x < z < y$} وي، نو وايو
چې \LR{$y$} د \LR{$x$} \emph{ځايناستے} دے۔
\end{defn}
\olpseganchor{OLP-0018-B011}{end}

\olpseganchor{OLP-0018-B012}{start}
د \LR{$x \in A$} ځايناستي د هغې \emph{اولادونه} هم بللے شي۔ که \LR{$y$}
د \LR{$x$} ځايناستے وي، نو \LR{$x$} د \LR{$y$} \emph{مخکينے} يا \emph{مور}
بولو۔
\olpseganchor{OLP-0018-B012}{end}

\olpseganchor{OLP-0018-B013}{start}
\begin{prop}
که \LR{$\tuple{A,\le}$} يوه ونه وي، نو له ريښې پرته د هر \LR{$x \in A$}
له يوه څخه زيات مخکيني نۀ شي کېدلے۔
\end{prop}
\olpseganchor{OLP-0018-B013}{end}

\olpseganchor{OLP-0018-B014}{start}
\begin{proof}
  فرض کړئ \LR{$y_1 < x$}، \LR{$y_2 < x$} او \LR{$y_1 \neq y_2$}۔ نو
  \LR{$\{y_1, y_2\} \subseteq \Setabs{z}{z<x}$}۔ ځکه چې
  \LR{$\Setabs{z}{z<x}$} د \LR{$\le$} په ذريعه ښۀ مرتب دے، د دې فرعي سټ
  \LR{$\{y_1, y_2\}$} تر ټولو کوچنے غړے لري، چې ښکاره ده يا به \LR{$y_1$}
  وي يا \LR{$y_2$}۔ نو يا \LR{$y_1 \le y_2$} يا \LR{$y_2 \le y_1$}۔ فرض مو
  کړے و چې \LR{$y_1 \neq y_2$}، نو په حقيقت کښې يا \LR{$y_1 < y_2$} يا
  \LR{$y_2 < y_1$}۔ ځکه چې \LR{$y_1 < x$} او \LR{$y_2 < x$} مو هم فرض کړي وو،
  دا هم لرو چې يا \LR{$y_1 < y_2 < x$} يا \LR{$y_2 < y_1 < x$}۔ نو
  \LR{$y_1$} او \LR{$y_2$} دواړه د \LR{$x$} مخکيني نۀ شي کېدلے۔
\end{proof}
\olpseganchor{OLP-0018-B014}{end}

\olpseganchor{OLP-0018-B015}{start}
\begin{defn}
ونه \LR{$T = \tuple{A, \le}$} \emph{نامتناهي} بللے شي که \LR{$A$} نامتناهي
سټ وي، او که نه نو \emph{متناهي} ده۔ که په \LR{$T$} کښې هر \LR{$x \in A$}
يوازې متناهي شمېر ځايناستي ولري، نو وايو چې \LR{$T$}
\emph{متناهي څانګېدونکې} ده۔
\end{defn}
\olpseganchor{OLP-0018-B015}{end}

\olpseganchor{OLP-0018-B016}{start}
\begin{defn}[څانګې]
که ونه \LR{$T = \tuple{A, \le}$} راکړل شي، نو د \LR{$T$} يوه \emph{څانګه}
په \LR{$T$} کښې يوه اعظمي ځنځير ده؛ يعنې داسې سټ \LR{$B \subseteq A$} چې
د هر \LR{$x, y \in B$} دپاره يا \LR{$x \le y$} يا \LR{$y \le x$} وي، او د
هر \LR{$z \in X \setminus B$} دپاره داسې \LR{$u \in B$} شته چې نۀ
\LR{$z \le u$} وي او نۀ \LR{$u \le z$}۔
%
د \LR{$T$} د ټولو څانګو د سټ د ښودلو دپاره \LR{$[T]$} کاروو۔
\end{defn}
\olpseganchor{OLP-0018-B016}{end}

\olpseganchor{OLP-0018-B017}{start}
\begin{ex}
د متناهي څانګېدونکې ونې يوه مشهوره بېلګه د \LR{$0$} ګانو او \LR{$1$} ګانو
د متناهي لړيو \emph{نامتناهي دوه ګونې ونه} ده۔ کله کله په
\LR{$\{0,1\}^*$} يا \LR{$\Bin^*$} ښودل کېږي، او د غځونې د اړيکې
\LR{$\sqsubseteq$} په ذريعه مرتب شوې ده (مثلاً، \LR{$101 \sqsubseteq 101101$})۔
ځکه چې هر دوه ګونے توريز تسلسل په پاے کښې د \LR{$0$} يا \LR{$1$} په
ورزياتولو تل غځولے شو، دا ونه نامتناهي ډېر غړي لري: هر غړے \LR{$s$}
پوره دوه ځايناستي لري، \LR{$s0$} او \LR{$s1$}۔ ريښه ئې تشه لړۍ \LR{$\emptyseq$} ده۔
\end{ex}
\olpseganchor{OLP-0018-B017}{end}

\olpseganchor{OLP-0018-B018}{start}
\begin{ex}
په يو څه زياته عمومي توګه، د طبيعي عددونو د متناهي لړيو سټ
\LR{$\Nat^*$} د غځونې له اړيکې \LR{$\sqsubseteq$} سره هم يوه ونه ده۔
ښکاره ده چې دا متناهي څانګېدونکې نۀ ده: هر \LR{$s \in \Nat^*$}
نامتناهي ډېر ځايناستي \LR{$sn$} لري، د هر \LR{$n \in \Nat$} دپاره يو۔
هر \LR{$A \subseteq \Nat^*$} چې د \LR{$\sqsubseteq$} لاندې تړلے وي، د
\LR{$\Nat^*$} يوه \emph{فرعي ونه} ده۔ (يعنې \LR{$A$} داسې دے چې که
\LR{$s \in A$} او \LR{$s' \sqsubseteq s$} وي، نو \LR{$s' \in A$} هم وي۔)
ټولې متناهي ونې د \LR{$\Nat^*$} د متناهي فرعي ونو په توګه ښودل کېدلے شي۔
\end{ex}
\olpseganchor{OLP-0018-B018}{end}

\olpseganchor{OLP-0018-B019}{start}
\begin{prop}[د کونيګ مرستندويه قضيه]
که \LR{$T = \tuple{A,\le}$} متناهي څانګېدونکې نامتناهي ونه وي، نو
\LR{$T$} يوه نامتناهي څانګه لري۔
\end{prop}
\olpseganchor{OLP-0018-B019}{end}

\olpseganchor{OLP-0018-B020}{start}
د کونيګ د مرستندويې قضيې يوه خاصه بڼه چې د محاسبه کېدنې په نظريه
کښې ډېره کارېږي، \emph{د کونيګ کمزورې مرستندويه قضيه} بللے شي:
د \LR{$\{0,1\}^*$} هره نامتناهي فرعي ونه يوه نامتناهي څانګه لري۔
\olpseganchor{OLP-0018-B020}{end}

\olpseganchor{OLP-0019-B004}{start}
\olfileid{sfr}{rel}{ops}
\olsection{په اړيکو باندې عمليات}
\olpseganchor{OLP-0019-B004}{end}

\olpseganchor{OLP-0019-B005}{start}
اړيکې بدلول يا سره يوځاے کول ډېر وخت ګټور وي۔ په
\olref[sfr][rel][ord]{prop:stricttopartial} کښې مو د اړيکو
\emph{اتحاد} په پام کښې نيولے و؛ دا يوازې د دوو اړيکو اتحاد دے
چې د جوړو د سټونو په توګه ګڼل شوې وي۔ همدارنګه، په
\olref[sfr][rel][ord]{prop:partialtostrict} کښې مو د اړيکو نسبي
تفاضل په پام کښې نيولے و۔ لاندې ځينې نور عمليات دي چې په اړيکو
باندې ئې کولے شو۔
\olpseganchor{OLP-0019-B005}{end}

\olpseganchor{OLP-0019-B006}{start}
\begin{defn}\ollabel{relationoperations}
فرض کړئ \LR{$R$} او \LR{$S$} اړيکې دي، او \LR{$A$} هر يو سټ دے۔
\olpseganchor{OLP-0019-B006}{end}

\olpseganchor{OLP-0019-B007}{start}
د \LR{$R$} \emph{معکوس} دا دے:
\LR{$R^{-1} = \Setabs{\tuple{y, x}}{\tuple{x, y} \in R}$}۔
\olpseganchor{OLP-0019-B007}{end}

\olpseganchor{OLP-0019-B008}{start}
د \LR{$R$} او \LR{$S$} \emph{نسبي حاصل ضرب} دا دے:
\LR{$(R \mid S) = \{\tuple{x, z} : \exists y(Rxy \land Syz)\}$}۔
\olpseganchor{OLP-0019-B008}{end}

\olpseganchor{OLP-0019-B009}{start}
تر \LR{$A$} پورې د \LR{$R$} \emph{محدودونه} دا ده:
\LR{$\funrestrictionto{R}{A}= R \cap A^2$}۔
\olpseganchor{OLP-0019-B009}{end}

\olpseganchor{OLP-0019-B010}{start}
په \LR{$A$} باندې د \LR{$R$} \emph{تطبيق} دا دے:
\LR{$\funimage{R}{A} = \{y : (\exists x \in A)Rxy\}$}
\end{defn}
\olpseganchor{OLP-0019-B010}{end}

\olpseganchor{OLP-0019-B011}{start}
\begin{ex}
فرض کړئ \LR{$S \subseteq \Int^2$} په \LR{$\Int$} باندې د ځايناستي اړيکه
ده؛ يعنې \LR{$S = \Setabs{\tuple{x, y} \in \Int^2}{x + 1 = y}$}،
نو \LR{$Sxy$} هله او يوازې هله چې \LR{$x + 1 = y$} وي۔
\olpseganchor{OLP-0019-B011}{end}

\olpseganchor{OLP-0019-B012}{start}
\LR{$S^{-1}$} په \LR{$\Int$} باندې د مخکيني اړيکه ده، يعنې:
\LR{$\Setabs{\tuple{x,y}\in\Int^2}{x -1 =y}$}۔
\olpseganchor{OLP-0019-B012}{end}

\olpseganchor{OLP-0019-B013}{start}
\LR{$S\mid S$} دا دے:
\LR{$ \Setabs{\tuple{x,y}\in\Int^2}{x + 2 =y}$}
\olpseganchor{OLP-0019-B013}{end}

\olpseganchor{OLP-0019-B014}{start}
\LR{$\funrestrictionto{S}{\Nat}$} په \LR{$\Nat$} باندې د ځايناستي اړيکه ده۔
\olpseganchor{OLP-0019-B014}{end}

\olpseganchor{OLP-0019-B015}{start}
\LR{$\funimage{S}{\{1,2,3\}}$} دا دے: \LR{$\{2, 3, 4\}$}۔
\end{ex}
\olpseganchor{OLP-0019-B015}{end}

\olpseganchor{OLP-0019-B016}{start}
\begin{defn}[انتقالي بندښت]فرض کړئ \LR{$R \subseteq A^2$} يوه دوه ځاييزه اړيکه ده۔
\olpseganchor{OLP-0019-B016}{end}

\olpseganchor{OLP-0019-B017}{start}
د \LR{$R$} \emph{انتقالي بندښت} دا دے:
\LR{$R^+ = \bigcup_{0 < n \in \Nat} R^n$}؛ دلته په بازګشتي ډول
\LR{$R^1 = R$} او \LR{$R^{n+1} = R^n \mid R$} تعريفوو۔
\olpseganchor{OLP-0019-B017}{end}

\olpseganchor{OLP-0019-B018}{start}
د \LR{$R$} \emph{انعکاسي انتقالي بندښت} دا دے:
\LR{$R^* = R^+ \cup \Id{A}$}۔
\end{defn}
\olpseganchor{OLP-0019-B018}{end}

\olpseganchor{OLP-0019-B019}{start}
\begin{ex}
د ځايناستي اړيکه \LR{$S \subseteq \Int^2$} واخلئ۔ \LR{$S^2xy$} هله او
يوازې هله چې \LR{$x + 2 = y$} وي، \LR{$S^3xy$} هله او يوازې هله چې
\LR{$x + 3 = y$} وي، او داسې نور۔ نو \LR{$S^+xy$} هله او يوازې هله چې
د کوم \LR{$n \geq 1$} دپاره \LR{$x + n = y$} وي۔ په نورو ټکو، \LR{$S^+xy$}
هله او يوازې هله چې \LR{$x < y$} وي، او \LR{$S^*xy$} هله او يوازې هله چې
\LR{$x \le y$} وي۔
\end{ex}
\olpseganchor{OLP-0019-B019}{end}

\olpseganchor{OLP-0019-B020}{start}
\begin{prob}
وښيئ چې د \LR{$R$} انتقالي بندښت په رښتيا انتقالي دے۔
\end{prob}
\olpseganchor{OLP-0019-B020}{end}

\olpseganchor{OLP-0020-B004}{start}
\olchapter{sfr}{fun}{تابعې}
\olpseganchor{OLP-0020-B004}{end}











\olpseganchor{OLP-0020-B010}{start}
%
\olpseganchor{OLP-0020-B010}{end}

\olpseganchor{OLP-0021-B004}{start}
\olfileid{sfr}{fun}{bas}
\olsection{بنسټونه}
\olpseganchor{OLP-0021-B004}{end}

\olpseganchor{OLP-0021-B005}{start}
\begin{explain}
\emph{تابع} داسې نگاشت دے چې د يوه ورکړل شوي سټ هر غړے
د کوم ورکړل شوي (بل) سټ له يوه ټاکلي غړي سره تړي۔
د بېلګې په توګه، د \LR{$1$} ورزياتولو عمل يوه تابع ټاکي: هر عدد~\LR{$n$}
له يوه يوازيني عدد~\LR{$n+1$} سره تړل کېږي۔
\olpseganchor{OLP-0021-B005}{end}
  
\olpseganchor{OLP-0021-B006}{start}
په عمومي ډول، تابعې جوړې، درې ګونې او داسې نور د ننوتونو په توګه
اخيستلے شي او يو ډول وت بېرته ورکوي۔ له بنسټيز حسابه ډېرې تابعې
پېژنو۔ د بېلګې په توګه، جمع او ضرب تابعې دي۔ دوې دوه عددونه
اخلي او درېيم عدد بېرته ورکوي۔
\olpseganchor{OLP-0021-B006}{end}

\olpseganchor{OLP-0021-B007}{start}
په دې رياضيکي، تجريدي معنٰی، تابع يو \emph{تور بکس} دے:
يوازې دا خبره اهميت لري چې کوم وت له کوم ننوت سره تړل شوے دے،
نۀ دا چې وت په کومه طريقه حسابېږي۔
\end{explain}
\olpseganchor{OLP-0021-B007}{end}

\olpseganchor{OLP-0021-B008}{start}
\begin{defn}[تابع]
\emph{تابع} \LR{$f \colon A \to B$} د~\LR{$A$} د هر غړي
د~\LR{$B$} له يوه غړي سره تړون دے۔
\olpseganchor{OLP-0021-B008}{end}

\olpseganchor{OLP-0021-B009}{start}
\LR{$A$} د~\LR{$f$} \emph{د تعريف ساحه} او \LR{$B$} د~\LR{$f$} \emph{هدف سټ} بولو۔
د~\LR{$A$} غړي د~\LR{$f$} ننوتونه يا \emph{ارګومېنټونه} بلل کېږي،
او د~\LR{$B$} هغه غړے چې~\LR{$f$} ئې له ارګومېنټ~\LR{$x$} سره تړي،
د ارګومېنټ~\LR{$x$} دپاره \emph{د~\LR{$f$} قيمت} بلل کېږي او په~\LR{$f(x)$} ليکل کېږي۔
\olpseganchor{OLP-0021-B009}{end}

\olpseganchor{OLP-0021-B010}{start}
د~\LR{$f$} \emph{د قيمتونو سټ} \LR{$\ran{f}$} د هدف سټ هغه فرعي سټ دے چې
د~\LR{$f$} له هغو قيمتونو جوړ دے چې د کوم ارګومېنټ دپاره ترلاسه کېږي؛
\LR{$\ran{f} = \Setabs{f(x)}{x \in A}$}۔
\end{defn}
\olpseganchor{OLP-0021-B010}{end}

\olpseganchor{OLP-0021-B011}{start}
په \olref{fig:function} کښې انځوريز شکل د تابعو په پوهېدو کښې مرسته
کولے شي۔ کيڼې خواته بيضوي شکل د تابعې \emph{د تعريف ساحه} ښيي؛
ښي خوا ته بيضوي شکل د تابعې \emph{هدف سټ} ښيي؛ او غشے د تعريف
په ساحه کښې له يوه \emph{ارګومېنټ} څخه په هدف سټ کښې ورته
\emph{قيمت} ته ځي۔
\olpseganchor{OLP-0021-B011}{end}

\olpseganchor{OLP-0021-B012}{start}
\begin{figure}
  \olasset{../upstream/assets/diagrams/function.tikz}
  \caption{تابع د يوه سټ د هر غړي د بل سټ له
    يوه غړي سره تړون دے۔ غشے د تعريف په ساحه کښې
    له يوه ارګومېنټ څخه په هدف سټ کښې ورته قيمت ته ځي۔}
  \ollabel{fig:function}
\end{figure}
\olpseganchor{OLP-0021-B012}{end}

\olpseganchor{OLP-0021-B013}{start}
\begin{ex}
ضرب د طبيعي عددونو جوړې د ننوتونو په توګه اخلي او طبيعي عددونه
د وتونو په توګه ورکوي؛ نو له \LR{$\Nat \times \Nat$} (د تعريف له ساحې)
څخه \LR{$\Nat$} (هدف سټ) ته ځي۔ د قيمتونو سټ هم \LR{$\Nat$} دے، ځکه چې
هر \LR{$n \in \Nat$} د \LR{$n \times 1$} په توګه ترلاسه کېږي۔
\end{ex}
\olpseganchor{OLP-0021-B013}{end}

\olpseganchor{OLP-0021-B014}{start}
\begin{ex}
ضرب تابع ده، ځکه چې هر ننوت---د طبيعي عددونو هره جوړه---له يوه
يوازيني وت سره تړي: \LR{$\times \colon \Nat^2 \to \Nat$}۔ خو له يوه
طبيعي عدد~\LR{$n$} سره داسې حقيقي عدد~\LR{$x$} تړل چې \LR{$x^2 = n$} وي،
تابع نۀ جوړوي، ځکه چې هر مثبت صحيح عدد~\LR{$n$} دوه مربع جذرونه لري:
\LR{$\sqrt{n}$} او~\LR{$-\sqrt{n}$}۔ که يوازې غير منفي (اصلي) مربع جذر بېرته ورکړو،
تابع ترې جوړولے شو: \LR{$\sqrt{\phantom{X}} \colon \Nat \to \Real$}۔
\end{ex}
\olpseganchor{OLP-0021-B014}{end}

\olpseganchor{OLP-0021-B015}{start}
\begin{ex}
هغه اړيکه چې د يوه ټولګي هر زده کوونکے د هغه له وروستۍ درجې سره
تړي، تابع ده---يو زده کوونکے په هماغه ټولګي کښې دوه بېلې وروستۍ
درجې نۀ شي اخيستلے۔ هغه اړيکه چې د يوه ټولګي هر زده کوونکے
د هغه له والدينو سره تړي، تابع نۀ ده: زده کوونکي کېدے شي هېڅ،
يا دوه، يا تر دوو زيات والدين ولري۔
\end{ex}
\olpseganchor{OLP-0021-B015}{end}

\olpseganchor{OLP-0021-B016}{start}
\begin{explain}
تابعې داسې تعريفولے شو چې په يوه کره طريقه د هر ممکن ارګومېنټ
دپاره د تابعې قيمت وټاکو۔ د دې بېلې لارې دا دي چې يو فورمول
ورکړو، د قيمت د حسابولو طريقه بيان کړو، يا د هر ارګومېنټ قيمت
په فهرست کښې وليکو۔ تابعې چې په هره طريقه تعريفوو، بايد ډاډمن
شو چې د هر ارګومېنټ دپاره يو، او يوازې يو، قيمت ټاکو۔
\end{explain}
\olpseganchor{OLP-0021-B016}{end}


\olpseganchor{OLP-0021-B017}{start}
\begin{ex}
فرض کړئ \LR{$f \colon \Nat \to \Nat$} داسې تعريف شوې ده چې \LR{$f(x) = x+1$}
وي۔ دا تعريف~\LR{$f$} د داسې تابعې په توګه ټاکي چې طبيعي عددونه اخلي
او طبيعي عددونه ورکوي۔ تعريف وايي چې د ورکړل شوي طبيعي عدد~\LR{$x$}
دپاره به~\LR{$f$} د هغه ځايناستے~\LR{$x+1$} ورکړي۔ په دې حالت کښې هدف
سټ \LR{$\Nat$} د~\LR{$f$} د قيمتونو سټ نۀ دے، ځکه چې طبيعي عدد~\LR{$0$} د
هېڅ طبيعي عدد ځايناستے نۀ دے۔ د~\LR{$f$} د قيمتونو سټ د ټولو مثبتو
صحيح عددونو سټ، \LR{$\Int^{+}$}، دے۔
\end{ex}
\olpseganchor{OLP-0021-B017}{end}

\olpseganchor{OLP-0021-B018}{start}
\begin{ex}\ollabel{examplefunext}
فرض کړئ \LR{$g \colon \Nat \to \Nat$} داسې تعريف شوې ده چې \LR{$g(x) = x+2-1$}
وي۔ دا وايي چې~\LR{$g$} يوه تابع ده چې طبيعي عددونه اخلي او طبيعي
عددونه ورکوي۔ د ورکړل شوي طبيعي عدد دپاره به~\LR{$g$} د~\LR{$x$} د
ځايناستي د ځايناستي مخکينے عدد، يعنې \LR{$x+1$}، ورکړي۔
\end{ex}
\olpseganchor{OLP-0021-B018}{end}

\olpseganchor{OLP-0021-B019}{start}
\begin{explain}
اوس مو دوه تابعې، \LR{$f$} او \LR{$g$}، له بېلو \emph{تعريفونو} سره په پام
کښې ونيولې۔ خو دا \emph{هماغه يوه تابع} ده۔ ځکه چې د هر طبيعي
عدد~\LR{$n$} دپاره \LR{$f(n) = n+1 = n+2-1 = g(n)$} لرو۔ په بله وينا:
د \LR{$f$} او~\LR{$g$} تعريفونه مو په بېلو معادلو هماغه يو نگاشت ټاکي۔ نو
په ضمني ډول د تابعو د قيمتونو له مخې د برابرۍ په يوه اصل تکيه کوو:
\[
  \text{\RL{که }}\forall x\, f(x) = g(x)\text{\RL{، نو }}f = g
\]
په دې شرط چې د \LR{$f$} او~\LR{$g$} د تعريف ساحه او هدف سټ يو شان وي۔
\end{explain}
\olpseganchor{OLP-0021-B019}{end}

\olpseganchor{OLP-0021-B020}{start}
\begin{ex}
تابعې د حالتونو په وېشلو هم تعريفولے شو۔ د بېلګې په توګه،
\LR{$h \colon \Nat \to \Nat$} داسې تعريفولے شو:
\[
h(x) =
\begin{cases}
  \frac{x}{2} & \text{\RL{که \LR{$x$} جفت وي}} \\
  \frac{x+1}{2} & \text{\RL{که \LR{$x$} طاق وي۔}}
\end{cases}
\]
ځکه چې هر طبيعي عدد يا جفت دے يا طاق، د دې تابعې وت به تل
طبيعي عدد وي۔ يوازې دا په ياد وساتئ چې که تابع د حالتونو په
وېشلو تعريفوئ، هر ممکن ننوت بايد په دقيقه توګه په يوه حالت کښې
راشي۔ ځينې وختونه د دې ثبوت ته اړتيا وي چې حالتونه ټول امکانات
رانغاړي او يو له بل سره ګډ نۀ دي۔
\end{ex}
\olpseganchor{OLP-0021-B020}{end}

\olpseganchor{OLP-0022-B004}{start}
\olfileid{sfr}{fun}{kin}
\olsection{د تابعو ډولونه}
\olpseganchor{OLP-0022-B004}{end}

\olpseganchor{OLP-0022-B005}{start}
\begin{explain}
د تابعو د هغو ځينو ډولونو دپاره چې ډېر ورسره مخ کېږو، د يو ډول
ډلبندۍ پېژندل به ګټور وي۔
\olpseganchor{OLP-0022-B005}{end}

\olpseganchor{OLP-0022-B006}{start}
په پيل کښې داسې تابعې په پام کښې نيول غواړو چې د هدف سټ هر
غړے د تابعې يو قيمت وي۔ دې تابعو ته پر هدف سټ بشپړه ويل کېږي،
او د \olref{fig:surjective} په شان ئې ښودلے شو۔
\olpseganchor{OLP-0022-B006}{end}

\olpseganchor{OLP-0022-B007}{start}
\begin{figure}
  \olasset{../upstream/assets/diagrams/surjective.tikz}
  \caption{پر هدف سټ بشپړه تابع د هدف سټ هر غړے د قيمت په توګه لري۔}
  \ollabel{fig:surjective}
\end{figure}
\end{explain}
\olpseganchor{OLP-0022-B007}{end}

\olpseganchor{OLP-0022-B008}{start}
\begin{defn}[پر هدف سټ بشپړه تابع]
تابع \LR{$f \colon A \rightarrow B$} هله او يوازې هله \emph{پر هدف سټ بشپړه}
ده چې \LR{$B$} د~\LR{$f$} د قيمتونو سټ هم وي؛ يعنې د هر \LR{$y \in B$}
دپاره لږ تر لږه يو \LR{$x \in A$} وي چې~\LR{$f(x) = y$} وي، يا په نښو:
\[
  (\forall y \in B)(\exists x \in A)f(x) = y.
\]
داسې تابع له \LR{$A$} څخه \LR{$B$} ته پر هدف سټ بشپړه تابع بولو۔
\end{defn}
\olpseganchor{OLP-0022-B008}{end}

\olpseganchor{OLP-0022-B009}{start}
\begin{explain}
که ښودل غواړئ چې \LR{$f$} يوه پر هدف سټ بشپړه تابع ده، نو بايد وښيئ چې د~\LR{$f$}
د هدف سټ هر څيز د کوم ننوت \LR{$x$} دپاره د \LR{$f(x)$} قيمت دے۔
\olpseganchor{OLP-0022-B009}{end}

\olpseganchor{OLP-0022-B010}{start}
پام وکړئ چې هره تابع يوه پر هدف سټ بشپړه تابع \emph{رامنځته کوي}۔ ځکه،
د ورکړل شوې تابعې \LR{$f \colon A \to B$} دپاره، فرض کړئ
\LR{$f' \colon A \to \ran{f}$} په \LR{$f'(x) = f(x)$} تعريف شوې ده۔ ځکه
چې \LR{$\ran{f}$} د \LR{$\Setabs{f(x) \in B}{x \in A}$} په توګه \emph{تعريف}
شوے دے، دا تضمين شوے دے چې تابع \LR{$f'$} به پر هدف سټ بشپړه تابع وي۔
\end{explain}
\olpseganchor{OLP-0022-B010}{end}

\olpseganchor{OLP-0022-B011}{start}
\begin{explain}
اوس، هره تابع هر ممکن ننوت له يوه يوازيني وت سره تړي۔ خو داسې
تابعې هم شته چې بېل ننوتونه هېڅکله له يو شان وت سره نۀ تړي۔ دې
تابعو ته يو پر يو ويل کېږي، او د \olref{fig:injective} په شان
ئې ښودلے شو۔
\begin{figure}
  \olasset{../upstream/assets/diagrams/injective.tikz}
  \caption{يو پر يو تابع هېڅکله دوه بېل ارګومېنټونه له يوه قيمت سره نۀ تړي۔}
  \ollabel{fig:injective}
\end{figure}
\end{explain}
\olpseganchor{OLP-0022-B011}{end}

\olpseganchor{OLP-0022-B012}{start}
\begin{defn}[يو پر يو تابع]
تابع \LR{$f \colon A \rightarrow B$} هله او يوازې هله \emph{يو پر يو}
ده چې د هر \LR{$y \in B$} دپاره تر يوه زيات داسې \LR{$x \in A$} نۀ وي
چې~\LR{$f(x) = y$} وي۔ داسې تابع له \LR{$A$} څخه~\LR{$B$} ته يو پر يو تابع بولو۔
\end{defn}
\olpseganchor{OLP-0022-B012}{end}

\olpseganchor{OLP-0022-B013}{start}
\begin{explain}
که ښودل غواړئ چې \LR{$f$} يوه يو پر يو تابع ده، بايد وښيئ چې د~\LR{$f$}
د تعريف د ساحې د هر دوو غړو \LR{$x$} او \LR{$y$} دپاره،
که \LR{$f(x)=f(y)$} وي، نو \LR{$x=y$} وي۔
\end{explain}
\olpseganchor{OLP-0022-B013}{end}

\olpseganchor{OLP-0022-B014}{start}
\begin{ex}
ثابته تابع \LR{$f\colon \Nat \to \Nat$} چې په \LR{$f(x) = 1$} ورکړل شوې،
نۀ يو پر يو ده او نۀ پر هدف سټ بشپړه۔
\olpseganchor{OLP-0022-B014}{end}

\olpseganchor{OLP-0022-B015}{start}
د عينيت تابع \LR{$f\colon \Nat \to \Nat$} چې په \LR{$f(x) = x$} ورکړل شوې،
هم يو پر يو ده او هم پر هدف سټ بشپړه۔
\olpseganchor{OLP-0022-B015}{end}

\olpseganchor{OLP-0022-B016}{start}
د ځايناستي تابع \LR{$f \colon \Nat \to \Nat$} چې په \LR{$f(x) = x+1$} ورکړل
شوې، يو پر يو ده خو پر هدف سټ بشپړه نۀ ده۔
\olpseganchor{OLP-0022-B016}{end}
  
\olpseganchor{OLP-0022-B017}{start}
تابع \LR{$f \colon \Nat \to \Nat$} چې داسې تعريف شوې ده:
\[
  f(x) =
  \begin{cases}
    \frac{x}{2} & \text{\RL{که \LR{$x$} جفت وي}} \\
    \frac{x+1}{2} & \text{\RL{که \LR{$x$} طاق وي۔}}
  \end{cases}
\]
پر هدف سټ بشپړه ده، خو يو پر يو نۀ ده۔
\end{ex}
\olpseganchor{OLP-0022-B017}{end}

\olpseganchor{OLP-0022-B018}{start}
\begin{explain}
ډېر وخت داسې تابعې په پام کښې نيول غواړو چې هم يو پر يو وي
او هم پر هدف سټ بشپړه۔ دا تابعې دوه اړخيزه يو پر يو بولو۔ دوې په
\olref{fig:bijective} کښې ښودل شوې تابعې ته ورته دي۔ دوه اړخيزې يو پر يو تابعې
کله کله \emph{يو پر يو مطابقتونه} هم بلل کېږي، ځکه چې د هدف سټ
غړي د تعريف د ساحې له غړو سره په يوازيني ډول جوړوي۔
\begin{figure}
  \olasset{../upstream/assets/diagrams/bijective.tikz}
  \caption{دوه اړخيزه يو پر يو تابع د هدف سټ غړي د تعريف د ساحې له غړو سره په يوازيني ډول جوړوي۔}
  \ollabel{fig:bijective}
\end{figure}
\end{explain}
\olpseganchor{OLP-0022-B018}{end}

\olpseganchor{OLP-0022-B019}{start}
\begin{defn}[دوه اړخيزه يو پر يو تابع]
تابع \LR{$f \colon A \to B$} هله او يوازې هله \emph{دوه اړخيزه يو پر يو} ده چې
هم پر هدف سټ بشپړه وي او هم يو پر يو۔ داسې تابع له \LR{$A$} څخه~\LR{$B$}
ته (يا د \LR{$A$} او~\LR{$B$} ترمنځ) دوه اړخيزه يو پر يو تابع بولو۔
\end{defn}
\olpseganchor{OLP-0022-B019}{end}

\olpseganchor{OLP-0023-B004}{start}
\olfileid{sfr}{fun}{rel}
\olpseganchor{OLP-0023-B004}{end}

\olpseganchor{OLP-0023-B005}{start}
\olsection{تابعې د اړيکو په توګه}
\olpseganchor{OLP-0023-B005}{end}

\olpseganchor{OLP-0023-B006}{start}
\begin{explain}
هغه تابع چې د~\LR{$A$} غړي د~\LR{$B$} له غړو سره
تړي، په ښکاره د \LR{$A$} او~\LR{$B$} ترمنځ يوه اړيکه ټاکي؛ يعنې هغه اړيکه
چې د \LR{$x$} او \LR{$y$} ترمنځ هله او يوازې هله وي چې \LR{$f(x) = y$} وي۔
په حقيقت کښې، که د بېلګې په توګه د رياضياتو د بنسټيزو توکو
په کمولو کښې دلچسپي لرو، ان تابع~\LR{$f$} له دې اړيکې سره، يعنې
د جوړو له يوه سټ سره، \emph{يوه ګڼلے} شو۔ بيا دا پوښتنه راپورته
کېږي: کومې اړيکې په دې ډول تابعې ټاکي؟
\end{explain}
\olpseganchor{OLP-0023-B006}{end}

\olpseganchor{OLP-0023-B007}{start}
\begin{defn}[د تابعې ګراف]فرض کړئ \LR{$f\colon A \to B$} يوه تابع ده۔
د~\LR{$f$} \emph{ګراف} هغه اړيکه \LR{$R_f \subseteq A \times B$} ده چې داسې
تعريف شوې ده:
\[
R_f = \Setabs{\tuple{x,y}}{f(x) = y}.
\]
\end{defn}
\olpseganchor{OLP-0023-B007}{end}

\olpseganchor{OLP-0023-B008}{start}
\begin{explain}
د غړو له مخې د برابرۍ په اصل، د تابعې ګراف په يوازيني ډول ټاکل
کېږي۔ سربېره پر دې، د سټونو د غړو له مخې د برابرۍ اصل سملاسي
د تابعو د قيمتونو له مخې د برابرۍ ضمني اصل هم توجيه کوي: که د
\LR{$f$} او~\LR{$g$} د تعريف ساحه او هدف سټ يو شان وي او په ټولو قيمتونو
کښې سره موافقې وي، نو هماغه يوه تابع ده۔
\olpseganchor{OLP-0023-B008}{end}

\olpseganchor{OLP-0023-B009}{start}
همدارنګه، که يوه اړيکه «د تابعې خاصيت ولري»، نو د يوې تابعې ګراف ده۔
\end{explain}
\olpseganchor{OLP-0023-B009}{end}

\olpseganchor{OLP-0023-B010}{start}
\begin{prop}\ollabel{prop:graph-function}
فرض کړئ \LR{$R \subseteq A \times B$} داسې ده چې:
\begin{enumerate}
\item که \LR{$Rxy$} او \LR{$Rxz$} وي، نو \LR{$y = z$} وي؛ او
\item د هر \LR{$x \in A$} دپاره داسې \LR{$y \in B$} شته چې \LR{$\tuple{x,
y} \in R$} وي۔
\end{enumerate}
نو \LR{$R$} د هغې تابعې \LR{$f\colon A \to B$} ګراف ده چې په دې ډول
تعريف شوې ده: \LR{$f(x) = y$} هله او يوازې هله چې \LR{$Rxy$} وي۔
\end{prop}
\olpseganchor{OLP-0023-B010}{end}

\olpseganchor{OLP-0023-B011}{start}
\begin{proof}
فرض کړئ داسې \LR{$y$} شته چې \LR{$Rxy$} وي۔ که يو بل \LR{$z \neq y$} هم واے
چې \LR{$Rxz$} واے، پر~\LR{$R$} لګېدلے شرط به مات شوے و۔ نو که داسې \LR{$y$}
شته چې \LR{$Rxy$} وي، دا \LR{$y$} يوازينے دے؛ له دې کبله~\LR{$f$} په سمه توګه
تعريف شوې ده۔ په ښکاره، \LR{$R_f = R$}۔
\end{proof}
\olpseganchor{OLP-0023-B011}{end}

\olpseganchor{OLP-0023-B012}{start}
\begin{explain}
هره تابع \LR{$f\colon A \to B$} ګراف لري؛ يعنې د همدې د تعريف ساحې
او هدف سټ ترمنځ يوه اړيکه، چې د \LR{$A \times B$} فرعي سټ دے او په \LR{$f(x) = y$}
تعريف شوې ده۔ بل خوا، هره اړيکه~\LR{$R
\subseteq A \times B$} چې په \olref{prop:graph-function} کښې ورکړل
شوي خاصيتونه ولري، د يوې تابعې~\LR{$f \colon A \to B$} ګراف ده۔ د
تابعو او د هغو د ګرافونو د دې نږدې تړاو له کبله، تابع يوازې د
هغې د ګراف په توګه ګڼلے شو۔ په بله وينا، تابعې له ځينو اړيکو،
يعنې د مرتبو څوګونو له ځينو سټونو، سره يوې ګڼلے شو۔
خو پام وکړئ چې د دې «يو ګڼلو» مطلب
هماغه دے چې په \olref[sfr][rel][ref]{sec} کښې و: دا د تابعو د
مابعدالطبيعي ماهيت په اړه ادعا نۀ ده، بلکې دا مشاهده ده چې
تابعې د ځينو سټونو په توګه \emph{ګڼل} اسانتيا لري۔ د دې اسانتيا
يو لامل دا دے چې اوس پر تابعو هم هغو ته ورته عمليات په
پام کښې نيولے شو چې پر اړيکو مو وکړل (وګورئ
\olref[sfr][rel][ops]{sec})۔ په ځانګړي ډول:
\end{explain}
\olpseganchor{OLP-0023-B012}{end}

\olpseganchor{OLP-0023-B013}{start}
\begin{defn}\ollabel{defn:funimage}
فرض کړئ \LR{$f \colon A \to B$} يوه تابع ده او \LR{$C\subseteq A$}۔
\olpseganchor{OLP-0023-B013}{end}

\olpseganchor{OLP-0023-B014}{start}
تر~\LR{$C$} پورې د~\LR{$f$} \emph{محدودونه} هغه تابع
\LR{$\funrestrictionto{f}{C}\colon C \to B$} ده چې د ټولو \LR{$x \in C$}
دپاره په \LR{$(\funrestrictionto{f}{C})(x) = f(x)$} تعريف شوې ده۔ په بله
وينا، \LR{$\funrestrictionto{f}{C} = \Setabs{\tuple{x, y} \in R_f}{x \in
C}$}۔
\olpseganchor{OLP-0023-B014}{end}

\olpseganchor{OLP-0023-B015}{start}
په~\LR{$C$} باندې د~\LR{$f$} \emph{تطبيق} دا دے: \LR{$\funimage{f}{C} =
\Setabs{f(x)}{x \in C}$}۔ دې ته د~\LR{$f$} لاندې د~\LR{$C$} \emph{انځور} هم وايو۔
\end{defn}
\olpseganchor{OLP-0023-B015}{end}

\olpseganchor{OLP-0023-B016}{start}
\begin{explain}
له دې تعريفونو څرګندېږي چې د هرې تابعې~\LR{$f$} دپاره
\LR{$\ran{f} = \funimage{f}{\dom{f}}$}۔
په \olref[sfr][rel][ops]{sec} کښې د
تعريفونو او له اړيکو سره د تابعو د يو ګڼلو له مخې، ورته
مفهومونه مو تمه لرله۔ دلته د تابعې محدودونه يوازې د ننوت
مختصات محدودوي؛ د اړيکې مخکښ تعريف دواړه مختصات محدودوي۔
خو دوه نور عمليات---معکوسونه او نسبي
حاصل ضربونه---لږ زيات تفصيل غواړي۔ هغه به په \olref[inv]{sec}
او \olref[cmp]{sec} کښې وړاندې کړو۔
\end{explain}
\olpseganchor{OLP-0023-B016}{end}

\olpseganchor{OLP-0024-B004}{start}
\olfileid{sfr}{fun}{inv}
\olsection{د تابعو معکوسونه}
\olpseganchor{OLP-0024-B004}{end}

\olpseganchor{OLP-0024-B005}{start}
\begin{explain}
تابعې د نگاشتونو په توګه ګڼو۔ نو د تابعو په اړه يوه څرګنده
پوښتنه دا ده چې ايا نگاشت «سرچپه» کېدے شي۔ د بېلګې په توګه،
د ځايناستي تابع \LR{$f(x) = x + 1$} په دې معنٰی سرچپه کېدے شي چې
تابع \LR{$g(y) = y - 1$} د \LR{$f$} کړے کار «بېرته اړوي»۔
\olpseganchor{OLP-0024-B005}{end}

\olpseganchor{OLP-0024-B006}{start}
خو بايد احتياط وکړو۔ که څه هم د~\LR{$g$} تعريف يوه تابع \LR{$\Int \to \Int$}
ټاکي، يوه \emph{تابع} \LR{$\Nat \to \Nat$} نۀ ټاکي، ځکه چې
\LR{$g(0) \notin \Nat$}۔ نو په ساده حالتونو کښې هم دا خبره پوره څرګنده
نۀ ده چې تابع سرچپه کېدے شي که نۀ؛ دا د تعريف په ساحه او هدف
سټ پورې تړلې کېدے شي۔
\olpseganchor{OLP-0024-B006}{end}

\olpseganchor{OLP-0024-B007}{start}
د تابعې د معکوس مفهوم دا خبره لا کره کوي۔
\end{explain}
\olpseganchor{OLP-0024-B007}{end}

\olpseganchor{OLP-0024-B008}{start}
\begin{defn}
تابع \LR{$g \colon B \to A$} د تابعې \LR{$f \colon A \to B$} \emph{معکوس}
ده که د ټولو \LR{$x \in A$} او \LR{$y \in B$} دپاره \LR{$f(g(y)) = y$} او
\LR{$g(f(x)) = x$} وي۔
\end{defn}
\olpseganchor{OLP-0024-B008}{end}

\olpseganchor{OLP-0024-B009}{start}
که \LR{$f$} معکوس~\LR{$g$} ولري، ډېر وخت د~\LR{$g$} پر ځاے \LR{$f^{-1}$} ليکو۔
\olpseganchor{OLP-0024-B009}{end}

\olpseganchor{OLP-0024-B010}{start}
\begin{explain}
اوس به معلومه کړو چې تابعې کله معکوسونه لري۔ د \LR{$f\colon A \to B$}
د معکوس دپاره يو مناسب نوماند \LR{$g\colon B \to A$} دے چې «په دې
تعريف شوے» وي:
\[
g(y) = \text{\RL{«هغه يوازينے» \LR{$x$} چې \LR{$f(x) = y$} وي۔}}
\]
خو د «په دې تعريف شوے» (او «هغه يوازينے») شاوخوا احتياطي نښې
اشاره کوي چې دا تعريف نۀ دے۔ لږ تر لږه، تل او په بشپړ عموميت
سره کار نۀ ورکوي۔ ځکه د دې دپاره چې دا تعريف يوه تابع وټاکي،
بايد يو او يوازې يو~\LR{$x$} وي چې \LR{$f(x) = y$} وي---د~\LR{$g$} وت بايد په
يوازيني ډول ټاکل شوے وي۔ سربېره پر دې، د هر \LR{$y \in B$} دپاره
بايد ټاکل شوے وي۔ که داسې \LR{$x_1$} او \LR{$x_2 \in A$} وي چې
\LR{$x_1 \neq x_2$} خو \LR{$f(x_1) = f(x_2)$} وي، نو د
\LR{$y = f(x_1) = f(x_2)$} دپاره به \LR{$g(y)$} په يوازيني ډول نۀ وي
ټاکل شوے۔ او که هېڅ داسې~\LR{$x$} نۀ وي چې \LR{$f(x) = y$} وي، نو \LR{$g(y)$}
به بېخي نۀ وي ټاکل شوے۔ په بله وينا، د~\LR{$g$} د تعريف کېدو دپاره
بايد~\LR{$f$} هم يو پر يو وي او هم پر هدف سټ بشپړه۔
\olpseganchor{OLP-0024-B010}{end}

\olpseganchor{OLP-0024-B011}{start}
راځئ ګام په ګام لاړ شو۔ پوښتنه به په دوو برخو ووېشو: د ورکړل
شوې تابعې~\LR{$f\colon A \to B$} دپاره، کله داسې تابع \LR{$g\colon B \to A$}
شته چې \LR{$g(f(x)) = x$} وي؟ داسې~\LR{$g$} د \LR{$f$} کړے کار «بېرته اړوي»
او د~\LR{$f$} \emph{کيڼ معکوس} بلل کېږي۔ دويمه دا چې کله داسې تابع
\LR{$h\colon B \to A$} شته چې \LR{$f(h(y)) = y$} وي؟ داسې~\LR{$h$} د~\LR{$f$}
\emph{ښي معکوس} بلل کېږي---\LR{$f$} د~\LR{$h$} کړے کار «بېرته اړوي»۔
\end{explain}
\olpseganchor{OLP-0024-B011}{end}

\olpseganchor{OLP-0024-B012}{start}
\begin{prop}
که د تعريف ساحه ناتشه وي او \LR{$f\colon A \to B$} يوه يو پر يو تابع وي، نو د~\LR{$f$} يو
\emph{کيڼ معکوس}~\LR{$g\colon B \to A$} شته چې د ټولو \LR{$x \in A$}
دپاره \LR{$g(f(x)) = x$} وي۔
\end{prop}
\olpseganchor{OLP-0024-B012}{end}

\olpseganchor{OLP-0024-B013}{start}
\begin{proof}
فرض کړئ \LR{$f\colon A \to B$} يوه يو پر يو تابع ده۔ يو \LR{$y \in B$}
په پام کښې ونيسئ۔ که \LR{$y \in \ran{f}$} وي، نو داسې \LR{$x \in A$} شته
چې \LR{$f(x) = y$} وي۔ ځکه چې \LR{$f$} يو پر يو ده، يوازې يو داسې
\LR{$x \in A$} شته۔ بيا داسې تعريفولے شو: \LR{$g(y) = x$}؛ يعنې \LR{$g(y)$}
«هغه يوازينے» \LR{$x \in A$} دے چې \LR{$f(x) = y$} وي۔ که
\LR{$y \notin \ran{f}$} وي، له هر يوه~\LR{$a \in A$} سره ئې تړلے شو۔ نو
يو \LR{$a \in A$} ټاکلے او \LR{$g\colon B \to A$} داسې تعريفولے شو:
\[
g(y) = \begin{cases}
    x & \text{\RL{که \LR{$f(x) = y$} وي}}\\
    a & \text{\RL{که \LR{$y \notin \ran{f}$} وي۔}}
\end{cases}
\]
دا د ټولو \LR{$y \in B$} دپاره تعريف شوې ده، ځکه چې د هر داسې
\LR{$y \in \ran{f}$} دپاره په دقيقه توګه يو \LR{$x \in A$} شته چې
\LR{$f(x) = y$} وي۔ د تعريف له مخې، که \LR{$y = f(x)$} وي، نو \LR{$g(y) = x$}
وي؛ يعنې \LR{$g(f(x)) = x$}۔
\end{proof}
\olpseganchor{OLP-0024-B013}{end}

\olpseganchor{OLP-0024-B014}{start}
\begin{prob}
وښيئ چې که \LR{$f\colon A \to B$} کيڼ معکوس~\LR{$g$} ولري، نو~\LR{$f$}
يو پر يو ده۔
\end{prob}
\olpseganchor{OLP-0024-B014}{end}

\olpseganchor{OLP-0024-B015}{start}
\begin{prop}
    که \LR{$f\colon A \to B$} يوه پر هدف سټ بشپړه تابع وي، نو د~\LR{$f$}
    يو \emph{ښي معکوس}~\LR{$h\colon B \to A$} شته چې د ټولو~\LR{$y \in B$}
    دپاره \LR{$f(h(y)) = y$} وي۔
\end{prop}
\olpseganchor{OLP-0024-B015}{end}

\olpseganchor{OLP-0024-B016}{start}
\begin{proof}
فرض کړئ \LR{$f\colon A \to B$} يوه پر هدف سټ بشپړه تابع ده۔ يو \LR{$y \in B$}
په پام کښې ونيسئ۔ ځکه چې~\LR{$f$} پر هدف سټ بشپړه ده، داسې \LR{$x_y \in A$}
شته چې \LR{$f(x_y) = y$} وي۔ بيا داسې تعريفولے شو: \LR{$h(y) = x_y$}؛
يعنې د هر \LR{$y \in B$} دپاره يو داسې \LR{$x \in A$} ټاکو چې \LR{$f(x) = y$}
وي؛ ځکه چې~\LR{$f$} پر هدف سټ بشپړه ده، تل لږ تر لږه يو د ټاکلو دپاره
شته۔\footnote{ځکه چې \LR{$f$} پر هدف سټ بشپړه ده، د هر~\LR{$y \in B$}
دپاره سټ \LR{$\Setabs{x}{f(x) = y}$} تش نۀ دے۔ زموږ د~\LR{$h$} تعريف غواړي
چې له دې هر سټ څخه يو يوازينے \LR{$x$} وټاکو۔ دا چې دا کار تل ممکن
دے، په حقيقت کښې څرګنده خبره نۀ ده---د دې ټاکنو امکان په ساده
ډول د يوه بديهي اصل په توګه فرض کېږي۔ په بله وينا، دا قضيه
د ټاکنې په نامه بديهي اصل فرض کوي؛ دا مسئله به
په \olref[sth][choice][]{chap} کښې
بيا وڅېړو۔ خو په ډېرو ځانګړو حالتونو کښې، لکه
چې \LR{$A = \Nat$} وي يا متناهي وي، يا \LR{$f$} يوه دوه اړخيزه يو پر يو تابع وي،
د ټاکنې بديهي اصل ته اړتيا نۀ وي۔ (په هغه ځانګړي حالت کښې چې
\LR{$f$} دوه اړخيزه يو پر يو وي، د هر \LR{$y \in B$} دپاره سټ
\LR{$\Setabs{x}{f(x) = y}$} په دقيقه توګه يو غړے لري؛ نو د
انتخاب دپاره څو بديلونه نشته۔)}
د تعريف له مخې، که \LR{$x = h(y)$} وي، نو \LR{$f(x) = y$} وي؛ يعنې
د هر \LR{$y \in B$} دپاره \LR{$f(h(y)) = y$}۔
\end{proof}
\olpseganchor{OLP-0024-B016}{end}

\olpseganchor{OLP-0024-B017}{start}
\begin{prob}
وښيئ چې که \LR{$f\colon A \to B$} ښي معکوس~\LR{$h$} ولري، نو~\LR{$f$}
پر هدف سټ بشپړه ده۔
\end{prob}
\olpseganchor{OLP-0024-B017}{end}

\olpseganchor{OLP-0024-B018}{start}
\begin{explain}
  د مخکيني ثبوت د مفکورو په يوځاے کولو اوس دا ترلاسه کوو چې
  هره دوه اړخيزه يو پر يو تابع معکوس لري؛ يعنې يوه واحده تابع شته چې
  هم د~\LR{$f$} کيڼ معکوس دے او هم ښي معکوس۔
\end{explain}
\olpseganchor{OLP-0024-B018}{end}

\olpseganchor{OLP-0024-B019}{start}
\begin{prop}\ollabel{prop:bijection-inverse}
که \LR{$f\colon A \to B$} يوه دوه اړخيزه يو پر يو تابع وي، داسې تابع
\LR{$f^{-1}\colon B \to A$} شته چې د ټولو \LR{$x \in A$} دپاره
\LR{$f^{-1}(f(x)) = x$} او د ټولو \LR{$y \in B$} دپاره \LR{$f(f^{-1}(y)) = y$} وي۔
\end{prop}
\olpseganchor{OLP-0024-B019}{end}

\olpseganchor{OLP-0024-B020}{start}
\begin{proof}
تمرين۔
\end{proof}
\olpseganchor{OLP-0024-B020}{end}

\olpseganchor{OLP-0024-B021}{start}
\begin{prob}
\olref[sfr][fun][inv]{prop:bijection-inverse} ثابته کړئ۔ بايد
\LR{$f^{-1}$} تعريف کړئ، وښيئ چې تابع ده، او وښيئ چې د~\LR{$f$} معکوس
ده؛ يعنې د ټولو \LR{$x \in A$} او \LR{$y \in B$} دپاره \LR{$f^{-1}(f(x)) = x$}
او \LR{$f(f^{-1}(y)) = y$} وي۔
\end{prob}
\olpseganchor{OLP-0024-B021}{end}

\olpseganchor{OLP-0024-B022}{start}
\begin{explain}
د معکوسونو د ترلاسه کولو يوه لږه زياته عمومي لار هم شته۔ په
\olref[kin]{sec} کښې مو وليدل چې هره تابع \LR{$f$} د ټولو \LR{$x \in A$}
دپاره د \LR{$f'(x) = f(x)$} په ټاکلو سره يوه پر هدف سټ بشپړه تابع
\LR{$f' \colon A \to \ran{f}$} رامنځته کوي۔ په ښکاره، که~\LR{$f$}
يو پر يو وي، نو~\LR{$f'$} دوه اړخيزه يو پر يو ده، او له دې کبله د
\olref{prop:bijection-inverse} له مخې يوازينے معکوس لري۔ د نښو
په کارولو کښې په يو ډېر واړه تسامح سره، کله کله د \LR{$f'$} معکوس
يوازې «د~\LR{$f$} معکوس» بولو۔
\end{explain}
\olpseganchor{OLP-0024-B022}{end}

\olpseganchor{OLP-0024-B023}{start}
\begin{prop}\ollabel{prop:left-right}%
  وښيئ چې که \LR{$f\colon A \to B$} کيڼ معکوس~\LR{$g$} او ښي معکوس~\LR{$h$}
  ولري، نو \LR{$h = g$}۔
\end{prop}
\olpseganchor{OLP-0024-B023}{end}

\olpseganchor{OLP-0024-B024}{start}
\begin{proof}
  تمرين۔
\end{proof}
\olpseganchor{OLP-0024-B024}{end}

\olpseganchor{OLP-0024-B025}{start}
\begin{prob}
  \olref[sfr][fun][inv]{prop:left-right} ثابته کړئ۔
\end{prob}
\olpseganchor{OLP-0024-B025}{end}

\olpseganchor{OLP-0024-B026}{start}
\begin{prop}\ollabel{prop:inverse-unique}
هره تابع~\LR{$f$} تر يوه زيات معکوس نۀ لري۔
\end{prop}
\olpseganchor{OLP-0024-B026}{end}

\olpseganchor{OLP-0024-B027}{start}
\begin{proof}
  فرض کړئ \LR{$g$} او \LR{$h$} دواړه د~\LR{$f$} معکوسونه دي۔ نو په ځانګړي ډول
  \LR{$g$} د~\LR{$f$} کيڼ معکوس او \LR{$h$} ښي معکوس دے۔ د
  \olref{prop:left-right} له مخې، \LR{$g = h$}۔
\end{proof}
\olpseganchor{OLP-0024-B027}{end}

\olpseganchor{OLP-0025-B004}{start}
\olfileid{sfr}{fun}{cmp}
\olsection{د تابعو ترکيب}
\olpseganchor{OLP-0025-B004}{end}

\olpseganchor{OLP-0025-B005}{start}
\begin{explain}
په \olref[inv]{sec} کښې مو وليدل چې
د يوې دوه اړخيزې يو پر يو تابعې~\LR{$f$} معکوس~\LR{$f^{-1}$} پخپله يوه تابع ده۔ پر
تابعو يو بل عمل ترکيب دے: موږ دوه تابعې، \LR{$f$} او~\LR{$g$}، سره
ترکيبولے او نوې تابع تعريفولے شو؛ يعنې لومړے~\LR{$f$} او بيا~\LR{$g$}
تطبيق کړو۔ البته، دا يوازې هغه وخت ممکنه ده چې د قيمتونو سټونه
او د تعريف ساحې سره برابر راشي؛ يعنې د~\LR{$f$} د قيمتونو سټ بايد د~\LR{$g$}
د تعريف د ساحې فرعي سټ وي۔ پر تابعو
دا عمل، پر اړيکو د نسبي حاصل ضرب له هغه عمل سره ورته دے چې په
\olref[rel][ops]{sec} کښې راغلے دے۔
\olpseganchor{OLP-0025-B005}{end}

\olpseganchor{OLP-0025-B006}{start}
انځوريز شکل د ترکيب د مفکورې په تشريح کښې مرسته کولے شي۔ په
\olref{fig:composition} کښې دوه تابعې \LR{$f \colon A \to B$} او
\LR{$g \colon B \to C$} او د هغو ترکيب~\LR{$(\comp{f}{g})$} ښيو۔ تابع
\LR{$(\comp{f}{g}) \colon A \to C$} د~\LR{$A$} هر غړے د~\LR{$C$} له يوه
غړي سره تړي۔ دا چې د \LR{$A$} يو غړے
د~\LR{$C$} له کوم غړي سره تړل کېږي، داسې ئې ټاکو:
د ورکړل شوي ننوت \LR{$x \in A$} دپاره، لومړے تابع \LR{$f$} پر~\LR{$x$} تطبيق
کړئ، چې يو \LR{$f(x) = y \in B$} به ورکړي؛ بيا تابع \LR{$g$} پر~\LR{$y$} تطبيق
کړئ، چې يو \LR{$g(f(x)) = g(y) = z \in C$} به ورکړي۔
\begin{figure}
  \olasset[2\olphotowidth]{../upstream/assets/diagrams/composition.tikz}
  \caption{د دوو تابعو \LR{$f$} او~\LR{$g$} ترکيب \LR{$g \circ f$}۔}
  \ollabel{fig:composition}
\end{figure}
\end{explain}
\olpseganchor{OLP-0025-B006}{end}

\olpseganchor{OLP-0025-B007}{start}
\begin{defn}[ترکيب]
فرض کړئ \LR{$f\colon A \to B$} او \LR{$g\colon B \to C$} تابعې دي۔
له~\LR{$g$} سره د \LR{$f$} \emph{ترکيب} دا دے: \LR{$\comp{f}{g} \colon A \to C$}،
چې \LR{$(\comp{f}{g})(x) = g(f(x))$} وي۔
\end{defn}
\olpseganchor{OLP-0025-B007}{end}

\olpseganchor{OLP-0025-B008}{start}
\begin{ex}
تابعې \LR{$f(x) = x + 1$} او \LR{$g(x) = 2x$} په پام کښې ونيسئ۔ ځکه چې
\LR{$(\comp{f}{g})(x) = g(f(x))$}، نو د هر ننوت~\LR{$x$} دپاره بايد لومړے
د هغه ځايناستے واخلئ، بيا پايله په دوو ضرب کړئ۔ نو د هغو ترکيب
په \LR{$(\comp{f}{g})(x) = 2(x+1)$} ورکړل شوے دے۔
\end{ex}
\olpseganchor{OLP-0025-B008}{end}

\olpseganchor{OLP-0025-B009}{start}
\begin{prob}
وښيئ چې که \LR{$f \colon A \to B$} او \LR{$g \colon B \to C$} دواړه
يو پر يو وي، نو \LR{$\comp{f}{g}\colon A \to C$} هم يو پر يو ده۔
\end{prob}
\olpseganchor{OLP-0025-B009}{end}

\olpseganchor{OLP-0025-B010}{start}
\begin{prob}
وښيئ چې که \LR{$f \colon A \to B$} او \LR{$g \colon B \to C$} دواړه
پر هدف سټ بشپړه وي، نو \LR{$\comp{f}{g}\colon A \to C$} هم پر هدف سټ بشپړه ده۔
\end{prob}
\olpseganchor{OLP-0025-B010}{end}

\olpseganchor{OLP-0025-B011}{start}
\begin{prob}
فرض کړئ \LR{$f \colon A \to B$} او \LR{$g \colon B \to C$}۔ وښيئ چې د
\LR{$\comp{f}{g}$} ګراف \LR{$R_f \mid R_g$} دے۔
\end{prob}
\olpseganchor{OLP-0025-B011}{end}

\olpseganchor{OLP-0717-B004}{start}
\olfileid{sfr}{fun}{iso}
\olsection{آيزومورفيزم}
\olpseganchor{OLP-0717-B004}{end}

\olpseganchor{OLP-0717-B005}{start}
\begin{explain}
\emph{آيزومورفيزم} داسې دوه‌اړخيزه يو پر يو تابع ده چې د هغو سټونو جوړښت ساتي چې سره تړي۔ جوړښت د سټونو د غړي تر منځ له اړيکو ټاکل کېږي۔ دا دوه سټونه \LR{$X=\{1,2,3\}$} او \LR{$Y=\{4,5,6\}$} په پام کښې ونيسئ۔ دواړه سټونه د ځاے‌ناستي، له ... نه کم او له ... نه زيات په اړيکو جوړښت مومي۔ د دواړو سټونو تر منځ آيزومورفيزم هغه دوه اړخيزه يو پر يو تابع ده چې همدا اړيکې ساتي۔ نو دوه اړخيزه يو پر يو تابع \LR{$f \colon X
\to Y$} آيزومورفيزم ده که \LR{$i<j$} هغه او يوازې هغه وخت چې \LR{$f(i)<f(j)$}، \LR{$i>j$} هغه او يوازې هغه وخت چې \LR{$f(i)>f(j)$}، او \LR{$j$} د \LR{$i$} ځاے‌ناستے وي هغه او يوازې هغه وخت چې \LR{$f(j)$} د \LR{$f(i)$} ځاے‌ناستے وي۔
\end{explain}
\olpseganchor{OLP-0717-B005}{end}

\olpseganchor{OLP-0717-B006}{start}
\begin{defn}[آيزومورفيزم]
فرض کړئ \LR{$U$} جوړه \LR{$\langle X, R\rangle$} او \LR{$V$} جوړه \LR{$\langle
Y, S\rangle$} ده، چې \LR{$X$} او \LR{$Y$} سټونه دي او \LR{$R$} او \LR{$S$} په ترتيب پر \LR{$X$} او \LR{$Y$} اړيکې دي۔ له \LR{$X$} نه \LR{$Y$} ته دوه اړخيزه يو پر يو تابع \LR{$f$} له \LR{$U$} نه \LR{$V$} ته \emph{آيزومورفيزم} ده هغه او يوازې هغه وخت چې اړيکيز جوړښت وساتي؛ يعنې په \LR{$X$} کښې د هر \LR{$x_{1}$} او \LR{$x_{2}$} دپاره \LR{$\tuple{x_1,x_2} \in R$} هغه او يوازې هغه وخت چې \LR{$\tuple{f(x_1),f(x_2)} \in S$}۔
\end{defn}
\olpseganchor{OLP-0717-B006}{end}

\olpseganchor{OLP-0717-B007}{start}
\begin{ex}
دا دوه سټونه \LR{$X=\{1,2,3\}$} او \LR{$Y=\{4,5,6\}$} او د کمښت او زياتوالي اړيکې په پام کښې ونيسئ۔ تابع \LR{$f\colon X \to
Y$}، چې \LR{$f(x) = 7-x$}، د \LR{$\tuple{X,<}$} او \LR{$\tuple{Y,>}$} تر منځ آيزومورفيزم ده۔
\end{ex}
\olpseganchor{OLP-0717-B007}{end}

\olpseganchor{OLP-0026-B004}{start}
\olfileid{sfr}{fun}{par}
\olpseganchor{OLP-0026-B004}{end}

\olpseganchor{OLP-0026-B005}{start}
\olsection{جزوي تابعې}
\olpseganchor{OLP-0026-B005}{end}

\olpseganchor{OLP-0026-B006}{start}
\begin{explain}
کله کله ګټوره وي چې د تابعې په تعريف کښې نرمي راولو، څو د ټولو
ممکنو ننوتونو دپاره د تابعې د وت تعريف کېدل لازم نۀ وي۔ داسې
نگاشتونه \emph{جزوي تابعې} بلل کېږي۔
\end{explain}
\olpseganchor{OLP-0026-B006}{end}

\olpseganchor{OLP-0026-B007}{start}
\begin{defn}
\emph{جزوي تابع} \LR{$f \colon A \pto B$} داسې نگاشت دے چې د~\LR{$A$} له
هر غړي سره د~\LR{$B$} تر يوه زيات غړے نۀ
تړي۔ که \LR{$f$} له \LR{$x \in A$} سره د~\LR{$B$} يو غړے وتړي، وايو چې \LR{$f(x)$}
\emph{تعريف شوے} دے، او که داسې نۀ وي، \emph{تعريف نۀ شوے} دے۔
که \LR{$f(x)$} تعريف شوے وي، \LR{$f(x) \fdefined$} ليکو؛ کنه،
\LR{$f(x) \fundefined$} ليکو۔ د جزوي تابعې~\LR{$f$} \emph{د تعريف ساحه}
د~\LR{$A$} هغه فرعي سټ دے چې تابع پرې تعريف شوې ده؛ يعنې
\LR{$\dom{f} = \Setabs{x \in A}{f(x) \fdefined}$}۔
\end{defn}
\olpseganchor{OLP-0026-B007}{end}

\olpseganchor{OLP-0026-B008}{start}
\begin{ex}
هره تابع \LR{$f\colon A \to B$} يوه جزوي تابع هم ده۔ هغه جزوي تابعې
چې پر~\LR{$A$} هرځاے تعريف شوې وي---يعنې هغه څه چې تر دې دمه مو
يوازې تابع بلل---\emph{هرځاے تعريف شوې} تابعې هم بلل کېږي۔
\end{ex}
\olpseganchor{OLP-0026-B008}{end}

\olpseganchor{OLP-0026-B009}{start}
\begin{ex}
جزوي تابع \LR{$f \colon \Real \pto \Real$} چې په \LR{$f(x) = 1/x$} ورکړل
شوې، د \LR{$x = 0$} دپاره تعريف شوې نۀ ده، او په نورو ټولو ځايونو
کښې تعريف شوې ده۔
\end{ex}
\olpseganchor{OLP-0026-B009}{end}

\olpseganchor{OLP-0026-B010}{start}
\begin{prob}
د ورکړل شوې \LR{$f\colon A \pto B$} دپاره، جزوي تابع \LR{$g\colon B \pto
A$} داسې تعريف کړئ: د هر \LR{$y \in B$} دپاره، که يو يوازينے
\LR{$x \in A$} وي چې \LR{$f(x) = y$} وي، نو \LR{$g(y) = x$}؛ کنه،
\LR{$g(y) \fundefined$}۔ وښيئ چې که \LR{$f$} يو پر يو وي، نو د ټولو
\LR{$x \in \dom{f}$} دپاره \LR{$g(f(x)) = x$}، او د ټولو \LR{$y \in \ran{f}$}
دپاره \LR{$f(g(y)) = y$} وي۔
\end{prob}
\olpseganchor{OLP-0026-B010}{end}

\olpseganchor{OLP-0026-B011}{start}
\begin{defn}[د جزوي تابعې ګراف]
فرض کړئ \LR{$f\colon A \pto B$} يوه جزوي تابع ده۔ د~\LR{$f$} \emph{ګراف}
هغه اړيکه \LR{$R_f \subseteq A \times B$} ده چې داسې تعريف شوې ده:
\[
R_f = \Setabs{\tuple{x,y}}{f(x) = y}.
\]
\end{defn}
\olpseganchor{OLP-0026-B011}{end}

\olpseganchor{OLP-0026-B012}{start}
\begin{prop}
فرض کړئ \LR{$R \subseteq A \times B$} دا خاصيت لري چې هرکله \LR{$Rxy$}
او \LR{$Rxy'$} وي، نو \LR{$y = y'$} وي۔ بيا \LR{$R$} د هغې جزوي تابعې
\LR{$f\colon A \pto B$} ګراف ده چې داسې تعريف شوې ده: که داسې \LR{$y$}
وي چې \LR{$Rxy$} وي، نو \LR{$f(x) = y$}؛ کنه، \LR{$f(x) \fundefined$}۔ که \LR{$R$}
\emph{مسلسله} هم وي، يعنې د هر \LR{$x \in A$} دپاره داسې \LR{$y \in B$}
وي چې \LR{$Rxy$} وي، نو \LR{$f$} هرځاے تعريف شوې ده۔
\end{prop}
\olpseganchor{OLP-0026-B012}{end}

\olpseganchor{OLP-0026-B013}{start}
\begin{proof}
فرض کړئ داسې \LR{$y$} شته چې \LR{$Rxy$} وي۔ که يو بل \LR{$y' \neq y$} هم واے
چې \LR{$Rxy'$} واے، پر \LR{$R$} لګېدلے شرط به مات شوے و۔ نو که داسې \LR{$y$}
شته چې \LR{$Rxy$} وي، هغه \LR{$y$} يوازينے دے؛ له دې کبله \LR{$f$} په سمه
توګه تعريف شوې ده۔ په ښکاره، \LR{$R_f = R$}، او که~\LR{$R$} مسلسله وي،
نو \LR{$f$} هرځاے تعريف شوې ده۔
\end{proof}
\olpseganchor{OLP-0026-B013}{end}

\olpseganchor{OLP-0027-B004}{start}
\olchapter{sfr}{siz}{د سټونو اندازه}
\olpseganchor{OLP-0027-B004}{end}

\olpseganchor{OLP-0027-B005}{start}
\begin{editorial}
دا باب شمېرنې، د شمېر وړ والے او د شمېر نۀ وړ والے څېړي۔
ځينې برخې دوه بڼې لري: يوه زياته ابتدايي بڼه، چې شمېرنې د
لړونو يا له \LR{$\PosInt$} څخه د هغو تابعو په توګه اخلي چې پر هدف
سټ بشپړې وي؛ او يوه زياته تجريدي بڼه، چې شمېرنې له \LR{$\Nat$} سره
د دوه اړخيزو يو پر يو تابعو په توګه تعريفوي۔
\end{editorial}
\olpseganchor{OLP-0027-B005}{end}





















\olpseganchor{OLP-0027-B016}{start}
\begin{editorial}
لاندې \olref[sfr][siz][enm-alt]{sec}،
\olref[sfr][siz][nen-alt]{sec} او \olref[sfr][siz][red-alt]{sec}
د \olref[sfr][siz][enm]{sec}، \olref[sfr][siz][nen]{sec} او
\olref[sfr][siz][red]{sec} بديلې بڼې دي چې ټيم بټن د خپل
«د سټونو خلاصه نظريه» متن کښې د کارولو دپاره برابرې کړې دي۔
دوې لږې زياتې پرمختللې دي او د شمېر وړ والي يو بېل تعريف
کاروي چې د سټونو د نظريې په سياق کښې زيات مناسب دے؛ يعنې
له \LR{$\Nat$} يا د هغې له پيلنۍ برخې سره دوه اړخيزه يو پر يو تابع
لرل، د دې پر ځاے چې په لړ کښې د راوړلو وړ وي يا له \LR{$\PosInt$}
څخه د يوې پر هدف سټ بشپړې تابعې د قيمتونو سټ وي۔
\end{editorial}
\olpseganchor{OLP-0027-B016}{end}

\olfileid{sfr}{siz}{alternate-outline}
\olpseganchor{OLP-0722-B004}{start}
\paragraph{د سټونو اندازه}
\olpseganchor{OLP-0722-B004}{end}

\olpseganchor{OLP-0722-B005}{start}
\begin{editorial}
دا باب شمېرنې، د شمېر وړ والے او د شمېر نۀ وړ والے څېړي۔ ځينې برخې دوه بڼې لري، يوه زياته ابتدايي او بله د سټ نظريې د تفصيلي بحث دپاره زياته پرمختللې۔ په دې بڼه کښې ابتدايي برخې دي۔ پرمختللې بڼې په \verb|size-of-sets-complete| کښې شته۔
\end{editorial}
\olpseganchor{OLP-0722-B005}{end}

\olpseganchor{OLP-0028-B004}{start}
\olfileid{sfr}{siz}{int}
\olpseganchor{OLP-0028-B004}{end}

\olpseganchor{OLP-0028-B005}{start}
\olsection{پېژندنه}
\olpseganchor{OLP-0028-B005}{end}

\olpseganchor{OLP-0028-B006}{start}
کله چې ګېورګ کانتور په 1870مو کلونو کښې د سټونو نظريې ته وده ورکړه،
د هغه يو مقصد دا و چې د نامتناهي مجموعې مفکوره د منلو وړ کړي---
هغه څه چې د منځنيو پېړيو پوهانو به فعلي نامتناهي والے بللو۔ د
دې يوه مهمه برخه د بېلو سټونو د \emph{اندازې} په اړه د هغه
څېړنه وه۔ که \LR{$a$}، \LR{$b$} او \LR{$c$} ټول بېل وي، نو په وجداني ډول
سټ \LR{$\{a, b, c\}$} له \LR{$\{a, b\}$} څخه \emph{لوے} دے۔ خو د
نامتناهي سټونو په اړه څه وايو؟ ايا ټول د يو بل هومره لوے دي؟
معلومېږي چې داسې نۀ ده۔
\olpseganchor{OLP-0028-B006}{end}

\olpseganchor{OLP-0028-B007}{start}
دلته لومړۍ مهمه مفکوره د شمېرنې ده۔ هر متناهي سټ د هغه د ټولو
غړو په ليکلو په لړ کښې راوړلے شو۔ د ځينو
نامتناهي سټونو ټول غړي هم په لړ کښې راوړلے شو، که
اجازه ورکړو چې لړ پخپله نامتناهي وي۔ داسې سټونه د شمېر وړ
بلل کېږي۔ د کانتور حيرانوونکې پايله، چې د دې باب تر پايه به
ئې پوره وپېژنو، دا وه چې ځينې نامتناهي سټونه د شمېر وړ نۀ دي۔
\olpseganchor{OLP-0028-B007}{end}

\olpseganchor{OLP-0029-B004}{start}
\olfileid{sfr}{siz}{enm}
\olpseganchor{OLP-0029-B004}{end}

\olpseganchor{OLP-0029-B005}{start}
\olsection{شمېرنې او د شمېر وړ سټونه}
\olpseganchor{OLP-0029-B005}{end}

\olpseganchor{OLP-0029-B006}{start}
\begin{editorial}
  دا برخه د سټونو شمېرنې څېړي او له \LR{$\PosInt$} څخه د پر هدف سټ
  بشپړو تابعو په توګه ئې تعريفوي۔ د هغو لوستونکو دپاره چې
  لږ رياضيکي شاليد لري، خبرې ګام په ګام وړاندې کوي۔ يوه
  بديله، لنډه بڼه په \olref[enm-alt]{sec} کښې ورکړل شوې، چې
  شمېرنې په بېل ډول تعريفوي: له \LR{$\Nat$} (يا د هغې له پيلنۍ
  برخې) سره د دوه اړخيزو يو پر يو تابعو په توګه۔
\end{editorial}
\olpseganchor{OLP-0029-B006}{end}

\olpseganchor{OLP-0029-B007}{start}
\begin{explain}
مخکښې مو د سټونو بېلګې د هغو د غړو په
ليکلو ورکړې دي۔ راځئ په زياته عمومي توګه بحث وکړو چې د يوه
سټ غړي څنګه او کله په لړ کښې راوړلے شو، که هغه سټ
نامتناهي هم وي۔
\end{explain}
\olpseganchor{OLP-0029-B007}{end}

\olpseganchor{OLP-0029-B008}{start}
\begin{defn}[شمېرنه، په غير صوري ډول]
په غير صوري ډول، د سټ~\LR{$A$} \emph{شمېرنه} د~\LR{$A$} د
غړو يو (کېدے شي نامتناهي) لړ دے، داسې چې د \LR{$A$}
هر غړے د لړ په يوه متناهي مقام کښې راشي۔ که \LR{$A$} شمېرنه
ولري، نو \LR{$A$} ته \emph{د شمېر وړ} ويل کېږي۔
\end{defn}
\olpseganchor{OLP-0029-B008}{end}

\olpseganchor{OLP-0029-B009}{start}
\begin{explain}
د شمېرنو په اړه څو ټکي:
\begin{enumerate}
\item يوازې هغه لړونه شمېرنې ګڼو چې پيل ولري او له لومړني
  غړي پرته د هر غړي سمدستي
  مخکښې يو يوازينے غړے وي۔ په بله وينا، د لړ د لومړي
  غړي او هر بل غړي ترمنځ
  يوازې متناهي شمېر غړي وي۔ په ځانګړي ډول، دا معنٰی لري چې
  د شمېرنې هر غړے متناهي مقام لري: لومړے غړے
  په مقام~\LR{$1$} کښې وي، دويم په مقام~\LR{$2$} کښې، او داسې نور۔
\item د هماغه سټ~\LR{$A$} بېلې شمېرنې لرلے شو چې د
  غړو د راتلو په ترتيب کښې سره بېلې وي:
  \LR{$4$}، \LR{$1$}، \LR{$25$}، \LR{$16$}،~\LR{$9$} د لومړيو پنځو مربع عددونو (د سټ)
  هومره ښه شمېرنه ده لکه \LR{$1$}، \LR{$4$}، \LR{$9$}، \LR{$16$}،~\LR{$25$}۔
\item تکراري شمېرنې هم شمېرنې دي: \LR{$1$}، \LR{$1$}، \LR{$2$}، \LR{$2$}، \LR{$3$}،
  \LR{$3$}،~\dots{} د هماغه سټ شمېرنه ده چې \LR{$1$}، \LR{$2$}، \LR{$3$}،~\dots{}
  ئې شمېرنه کوي۔
\item کله چې يوه شمېرنه ټاکو، ترتيب او تکرار \emph{اهميت لري}:
  د مثبتو صحيح عددونو شمېرنه په \LR{$1$}، \LR{$2$}، \LR{$3$}، \LR{$1$}، \dots{}
  پيلولے شو، خو نمونه هغه وخت په اسانه ښکاري چې په معمولي
  ډول ئې په \LR{$1$}، \LR{$2$}، \LR{$3$}، \LR{$4$}،~\dots وشمېرو۔
\item شمېرنې بايد پيل ولري: \dots، \LR{$3$}، \LR{$2$}، \LR{$1$} د مثبتو صحيح
  عددونو شمېرنه نۀ ده، ځکه چې لومړے غړے نۀ لري۔ د دې
  د ليدلو دپاره چې دا له غير صوري تعريفه څنګه راځي، له ځانه
  وپوښتئ: «عدد 76 د لړ په کوم مقام کښې راځي؟»
\item لاندې لړ د مثبتو صحيح عددونو شمېرنه نۀ ده:
  \LR{$1$}، \LR{$3$}، \LR{$5$}، \dots، \LR{$2$}، \LR{$4$}، \LR{$6$}، \dots\@ ستونزه دا ده چې
  جفت عددونه د متناهي مقامونو پر ځاے په \LR{$\infty + 1$}،
  \LR{$\infty + 2$}، \LR{$\infty + 3$} مقامونو کښې راځي۔
\item تش سټ د شمېر وړ دے: تش لړ ئې شمېري!{}
\end{enumerate}
\end{explain}
\olpseganchor{OLP-0029-B009}{end}

\olpseganchor{OLP-0029-B010}{start}
\begin{prop}
  که \LR{$A$} شمېرنه ولري، بې له تکراره شمېرنه هم لري۔
\end{prop}
\olpseganchor{OLP-0029-B010}{end}

\olpseganchor{OLP-0029-B011}{start}
\begin{proof}
  فرض کړئ \LR{$A$} شمېرنه \LR{$x_1$}، \LR{$x_2$}، \dots{} لري، چې هر \LR{$x_i$}
  د~\LR{$A$} يو غړے دے۔ له شمېرنې څخه د تکراري
  غړو په لرې کولو تکرارونه لرې کولے شو۔ د
  بېلګې په توګه، شمېرنه داسې نوې شمېرنې ته اړولے شو چې \LR{$x_i$}
  پکښې هله وليکو چې د~\LR{$A$} داسې غړے وي چې په \LR{$x_1$}،
  \dots، \LR{$x_{i-1}$} کښې نۀ وي، يا \LR{$x_i$} له لړه لرې کړو که له
  مخکښې په \LR{$x_1$}، \dots،~\LR{$x_{i-1}$} کښې راغلے وي۔
\end{proof}
\olpseganchor{OLP-0029-B011}{end}

\olpseganchor{OLP-0029-B012}{start}
وروستے استدلال ښيي چې د شمېرنو او د شمېر وړ سټونو د ښه
پوهېدو او د هغو په اړه د خبرو د ثابتولو دپاره يو زيات کره
تعريف ته اړتيا لرو۔ لاندې تعريف دا اړتيا پوره کوي۔
\olpseganchor{OLP-0029-B012}{end}

\olpseganchor{OLP-0029-B013}{start}
\begin{defn}[شمېرنه، په صوري ډول]
د سټ \LR{$A \neq \emptyset$} \emph{شمېرنه} هره پر هدف سټ بشپړه
تابع \LR{$f \colon \PosInt \to A$} ده۔
\end{defn}
\olpseganchor{OLP-0029-B013}{end}

\olpseganchor{OLP-0029-B014}{start}
\begin{explain}
راځئ ځان ډاډمن کړو چې صوري تعريف او هغه غير صوري تعريف چې
يو ممکن نامتناهي لړ کاروي، همارزښته دي۔ لومړے، له \LR{$\PosInt$}
څخه سټ~\LR{$A$} ته هره پر هدف سټ بشپړه تابع د~\LR{$A$} شمېرنه کوي۔ داسې
تابع د پاسني غير صوري تعريف له مخې يوه شمېرنه ټاکي: لړ
\LR{$f(1)$}، \LR{$f(2)$}، \LR{$f(3)$}، \dots۔ ځکه چې \LR{$f$} پر هدف سټ بشپړه ده،
دا تضمين شوے چې د~\LR{$A$} هر غړے به د کوم~\LR{$n \in \PosInt$}
دپاره د~\LR{$f(n)$} قيمت وي۔ نو د \LR{$A$} هر غړے په لړ کښې
په کوم متناهي مقام راځي۔ ځکه چې کېدے شي تابع يو پر يو
نۀ وي، لړ تکرار لرلے شي؛ خو دا د منلو وړ دي، لکه چې پورته
ياد شول۔
\olpseganchor{OLP-0029-B014}{end}

\olpseganchor{OLP-0029-B015}{start}
بل خوا، که داسې لړ راکړل شوے وي چې د~\LR{$A$} ټول غړي
شمېري، يوه پر هدف سټ بشپړه تابع \LR{$f\colon \PosInt \to A$} داسې
تعريفولے شو چې \LR{$f(n)$} د لړ \LR{$n$}م غړے وټاکو، يا که
\LR{$n$}م غړے نۀ وي، نو د لړ وروستے غړے وټاکو۔
يوازې په هغه حالت کښې له دې څخه پر هدف سټ بشپړه تابع نۀ
جوړېږي چې \LR{$A$} تش وي او له دې کبله لړ هم تش وي۔ نو هر ناتش
لړ يوه پر هدف سټ بشپړه تابع \LR{$f\colon \PosInt \to A$} ټاکي۔
\end{explain}
\olpseganchor{OLP-0029-B015}{end}

\olpseganchor{OLP-0029-B016}{start}
\begin{defn}
  \ollabel{defn:enumerable}
  سټ~\LR{$A$} هله او يوازې هله د شمېر وړ دے چې تش وي يا شمېرنه ولري۔
\end{defn}
\olpseganchor{OLP-0029-B016}{end}

\olpseganchor{OLP-0029-B017}{start}
\begin{ex}
د مثبتو صحيح عددونو (\LR{$\PosInt$}) شمېرونکې تابع يوازې د عينيت
هغه تابع ده چې په \LR{$f(n) = n$} ورکړل شوې ده۔ د طبيعي عددونو
\LR{$\Nat$} شمېرونکې تابع \LR{$g(n) = n - 1$} ده۔
\end{ex}
\olpseganchor{OLP-0029-B017}{end}

\olpseganchor{OLP-0029-B018}{start}
\begin{ex}
تابعې \LR{$f\colon \PosInt \to \PosInt$} او \LR{$g \colon \PosInt \to
\PosInt$} چې داسې ورکړل شوې دي:
\begin{align*}
f(n) & = 2n \text{\RL{ او}}\\
g(n) & = 2n - 1
\end{align*}
په ترتيب سره د جفتو مثبتو صحيح عددونو او د طاقو مثبتو صحيح
عددونو شمېرنه کوي۔ خو له دوو هېڅ يوه د \LR{$\PosInt$} شمېرنه نۀ
ده، ځکه چې هېڅ يوه پر هدف سټ بشپړه نۀ ده۔
\end{ex}
\olpseganchor{OLP-0029-B018}{end}

\olpseganchor{OLP-0029-B019}{start}
\begin{prob}
د مثبتو مربع عددونو \LR{$1$}، \LR{$4$}، \LR{$9$}، \LR{$16$}، \dots يوه شمېرنه تعريف کړئ۔
\end{prob}
\olpseganchor{OLP-0029-B019}{end}

\olpseganchor{OLP-0029-B020}{start}
\begin{ex}
تابع \LR{$f(n) = (-1)^{n} \lceil \frac{(n-1)}{2}\rceil$} (چې
\LR{$\lceil x \rceil$} د \emph{پاسني صحيح عدد} تابع ښيي، يعنې \LR{$x$}
پورته تر ټولو نږدې صحيح عدد ته ګردوي) د صحيح عددونو~\LR{$\Int$}
د سټ شمېرنه کوي۔ پام وکړئ چې \LR{$f$} د مثبتو او منفي صحيح عددونو
ترمنځ په تګ راتګ سره د \LR{$\Int$} قيمتونه څنګه راوباسي:
\[
\begin{array}{c c c c c c c c}
f(1) & f(2) & f(3) & f(4) & f(5) & f(6) & f(7) & \dots \\ \\
- \lceil \tfrac{0}{2} \rceil & \lceil \tfrac{1}{2}\rceil & - \lceil \tfrac{2}{2} \rceil & \lceil \tfrac{3}{2} \rceil & - \lceil \tfrac{4}{2} \rceil  & \lceil \tfrac{5}{2}
\rceil & - \lceil \tfrac{6}{2} \rceil & \dots \\ \\
0 & 1 & -1 & 2 & -2 & 3 & -3 & \dots
\end{array}
\]
\LR{$f$} د حالتونو له مخې داسې تعريف شوې هم ګڼلے شئ:
\[
f(n) = \begin{cases}
  0 & \text{\RL{که \LR{$n = 1$} وي}}\\
  n/2 & \text{\RL{که \LR{$n$} جفت وي}}\\
  -(n-1)/2 & \text{\RL{که \LR{$n$} طاق او \LR{$>1$} وي}}
  \end{cases}
\]
\end{ex}
\olpseganchor{OLP-0029-B020}{end}

\olpseganchor{OLP-0029-B021}{start}
\begin{prob}
  وښيئ چې که \LR{$A$} او \LR{$B$} د شمېر وړ وي، نو \LR{$A \cup B$} هم
  داسې دے۔ د دې دپاره فرض کړئ چې پر هدف سټ بشپړه تابعې
  \LR{$f\colon \PosInt \to A$} او \LR{$g\colon \PosInt \to B$} شته؛
  يوه پر هدف سټ بشپړه تابع~\LR{$h\colon \PosInt \to A \cup B$}
  تعريف کړئ او ثابته کړئ چې پر هدف سټ بشپړه ده۔ هغه حالتونه
  هم په پام کښې ونيسئ چې \LR{$A$} يا~\LR{$B = \emptyset$} وي۔
\end{prob}
\olpseganchor{OLP-0029-B021}{end}
  
\olpseganchor{OLP-0029-B022}{start}
\begin{prob}
  وښيئ چې که \LR{$B \subseteq A$} او \LR{$A$} د شمېر وړ وي، نو~\LR{$B$}
  هم داسې دے۔ د دې دپاره فرض کړئ چې يوه پر هدف سټ بشپړه
  تابع \LR{$f\colon \PosInt \to A$} شته۔ يوه پر هدف سټ بشپړه
  تابع~\LR{$g\colon \PosInt \to B$} تعريف کړئ او ثابته کړئ چې
  پر هدف سټ بشپړه ده۔ که \LR{$B = \emptyset$} وي، څه پېښېږي؟
\end{prob}
\olpseganchor{OLP-0029-B022}{end}
    
\olpseganchor{OLP-0029-B023}{start}
\begin{prob}
  پر \LR{$n$} د استقرا په وسيله وښيئ چې که \LR{$A_1$}، \LR{$A_2$}، \dots،
  \LR{$A_n$} ټول د شمېر وړ وي، نو \LR{$A_1 \cup \dots \cup A_n$}
  هم داسې دے۔ دا حقيقت فرضولے شئ چې که دوه سټونه \LR{$A$} او~\LR{$B$}
  د شمېر وړ وي، نو~\LR{$A \cup B$} هم داسې دے۔
\end{prob}
\olpseganchor{OLP-0029-B023}{end}

\olpseganchor{OLP-0029-B024}{start}
که څه هم د يوه سټ د غړو د لړ کولو پر وخت
ښايي له \LR{$1$}م غړي څخه شمېر پيلول زيات طبيعي
ښکاري، رياضي پوهان د څيزونو د شمېرلو دپاره طبيعي عددونه~\LR{$\Nat$}
کارول خوښوي۔ دوے د لړ د \LR{$0$}م، \LR{$1$}م، \LR{$2$}م، او داسې نورو
غړو خبرې کوي۔ له دې سره سم، شمېرنه له
\LR{$\Nat$} څخه~\LR{$A$} ته د يوې پر هدف سټ بشپړې تابعې په توګه تعريفولے
شو۔ البته، دواړه تعريفونه همارزښته دي۔
\olpseganchor{OLP-0029-B024}{end}

\olpseganchor{OLP-0029-B025}{start}
\begin{prop}\ollabel{prop:enum-shift}
  يوه پر هدف سټ بشپړه تابع \LR{$f\colon \PosInt \to A$} هله او يوازې هله
  شته چې يوه پر هدف سټ بشپړه تابع \LR{$g\colon \Nat \to A$} شته وي۔
\end{prop}
\olpseganchor{OLP-0029-B025}{end}

\olpseganchor{OLP-0029-B026}{start}
\begin{proof}
  د ورکړل شوې پر هدف سټ بشپړې تابعې \LR{$f\colon \PosInt \to A$} دپاره،
  د ټولو \LR{$n \in \Nat$} دپاره \LR{$g(n) = f(n+1)$} تعريفولے شو۔ په
  اسانه ښکاري چې \LR{$g\colon \Nat \to A$} يوه پر هدف سټ بشپړه تابع
  ده۔ برعکس، د ورکړل شوې پر هدف سټ بشپړې تابعې \LR{$g\colon \Nat \to A$}
  دپاره، \LR{$f(n) = g(n-1)$} تعريف کړئ۔
\end{proof}
\olpseganchor{OLP-0029-B026}{end}

\olpseganchor{OLP-0029-B027}{start}
له دې لاندې پايله ترلاسه کېږي:
\olpseganchor{OLP-0029-B027}{end}

\olpseganchor{OLP-0029-B028}{start}
\begin{cor}\ollabel{cor:enum-nat}
سټ \LR{$A$} هله او يوازې هله د شمېر وړ دے چې تش وي يا يوه
پر هدف سټ بشپړه تابع \LR{$f\colon \Nat \to A$} شته وي۔
\end{cor}
\olpseganchor{OLP-0029-B028}{end}

\olpseganchor{OLP-0029-B029}{start}
پورته مو بحث وکړ چې د سټ~\LR{$A$} د غړو لړ په
بې تکراره لړ اړول کېدے شي۔ دا د شمېرنو دپاره هم رښتيا ده، خو
په کره ډول ئې بيانول او ثابتول لږ ګران دي۔ هره تابع
\LR{$f\colon \PosInt \to A$} بايد د ټولو \LR{$n \in \PosInt$} دپاره
تعريف شوې وي۔ که په~\LR{$A$} کښې يوازې متناهي شمېر غړي وي،
نو په ښکاره داسې تابع نۀ شو لرلے چې د~\LR{$\PosInt$} په نامتناهي
شمېر غړو تعريف شوې وي، د~\LR{$A$} ټول
غړي د قيمتونو په توګه واخلي، خو هېڅکله يو قيمت دوه
ځله وانخلي۔ په دې حالت کښې، يعنې کله چې بې تکراره لړ متناهي
وي، بايد د~\LR{$f$} دپاره د تعريف يوه بېله ساحه وټاکو، چې يوازې
متناهي شمېر غړي ولري۔ د تکرارونو نۀ لرل دا معنٰی
لري چې \LR{$f$} بايد يو پر يو وي۔ ځکه چې پر هدف سټ بشپړه هم ده،
نو د کوم متناهي سټ \LR{$\{1, \dots, n\}$} يا \LR{$\PosInt$} او~\LR{$A$}
ترمنځ يوه دوه اړخيزه يو پر يو تابع لټوو۔
\olpseganchor{OLP-0029-B029}{end}

\olpseganchor{OLP-0029-B030}{start}
\begin{prop}\ollabel{prop:enum-bij}
که \LR{$f\colon \PosInt \to A$} يوه پر هدف سټ بشپړه تابع وي (يعنې
د~\LR{$A$} شمېرنه وي)، نو يوه دوه اړخيزه يو پر يو تابع \LR{$g\colon Z \to A$}
شته چې \LR{$Z$} يا~\LR{$\PosInt$} وي يا د کوم~\LR{$n \in \PosInt$} دپاره
\LR{$\{1, \dots, n\}$} وي۔
\end{prop}
\olpseganchor{OLP-0029-B030}{end}

\olpseganchor{OLP-0029-B031}{start}
\begin{proof}
  تابع \LR{$g$} په بازګشتي ډول تعريفوو: فرض کړئ \LR{$g(1) = f(1)$}۔ که
  \LR{$g(i)$} له مخکښې تعريف شوے وي، نو \LR{$g(i+1)$} د \LR{$f(1)$}،
  \LR{$f(2)$}، \dots{} لومړے هغه قيمت وټاکئ چې تر اوسه په \LR{$g(1)$}،
  \dots، \LR{$g(i)$} کښې نۀ وي، که داسې قيمت شته وي۔ که \LR{$A$} يوازې
  \LR{$n$} غړي ولري، نو \LR{$g(1)$}، \dots، \LR{$g(n)$} ټول تعريف
  شوي دي، او په دې ډول مو يوه تابع \LR{$g\colon \{1, \dots, n\}
  \to A$} تعريف کړې ده۔ که \LR{$A$} نامتناهي شمېر غړي
  ولري، نو د هر \LR{$i$} دپاره بايد د \LR{$f(1)$}، \LR{$f(2)$}، \dots په
  شمېرنه کښې د~\LR{$A$} يو داسې غړے وي چې تر اوسه په
  \LR{$g(1)$}، \dots، \LR{$g(i)$} کښې نۀ وي۔ په دې حالت کښې مو يوه
  تابع \LR{$g\colon \PosInt \to A$} تعريف کړې ده۔
\olpseganchor{OLP-0029-B031}{end}
  
\olpseganchor{OLP-0029-B032}{start}
  تابع \LR{$g$} پر هدف سټ بشپړه ده، ځکه چې د~\LR{$A$} هر غړے په \LR{$f(1)$}،
  \LR{$f(2)$}، \dots{} کښې دے (ځکه چې \LR{$f$} پر هدف سټ بشپړه ده) او له
  دې کبله به بالاخره د کوم~\LR{$i$} دپاره د~\LR{$g(i)$} قيمت شي۔ دا
  يو پر يو هم ده، ځکه چې که داسې \LR{$j < i$} واے چې
  \LR{$g(j) = g(i)$} واے، نو \LR{$g(i)$} به له مخکښې په \LR{$g(1)$}،
  \dots، \LR{$g(i-1)$} کښې و، چې زموږ د~\LR{$g$} له تعريف سره په ټکر کښې دے۔
\end{proof}
\olpseganchor{OLP-0029-B032}{end}

\olpseganchor{OLP-0029-B033}{start}
\begin{cor}\ollabel{cor:enum-nat-bij}
سټ \LR{$A$} هله او يوازې هله د شمېر وړ دے چې تش وي يا يوه
دوه اړخيزه يو پر يو تابع \LR{$f\colon N \to A$} شته وي، چې يا \LR{$N = \Nat$}
وي يا د کوم \LR{$n \in \Nat$} دپاره \LR{$N = \{0, \dots, n\}$} وي۔
\end{cor}
\olpseganchor{OLP-0029-B033}{end}

\olpseganchor{OLP-0029-B034}{start}
\begin{proof}
\LR{$A$} هله او يوازې هله د شمېر وړ دے چې \LR{$A$} تش وي يا يوه
پر هدف سټ بشپړه \LR{$f\colon \PosInt \to A$} شته وي۔ د
\olref{prop:enum-bij} له مخې، وروستۍ خبره هله او يوازې هله
سمه ده چې يوه دوه اړخيزه يو پر يو تابع~\LR{$f\colon Z \to A$} شته وي،
چې \LR{$Z = \PosInt$} يا د کوم \LR{$n \in \PosInt$} دپاره
\LR{$Z = \{1, \dots, n\}$} وي۔ د \olref{prop:enum-shift} د ثبوت
د هماغه استدلال له مخې، دا بيا هله او يوازې هله سمه ده چې
يوه دوه اړخيزه يو پر يو تابع \LR{$g\colon N \to A$} شته وي، چې يا \LR{$N = \Nat$}
وي يا \LR{$N = \{0, \dots, n-1\}$} وي۔
\end{proof}
\olpseganchor{OLP-0029-B034}{end}

\olpseganchor{OLP-0029-B035}{start}
\begin{prob}
  د \olref[sfr][siz][enm]{defn:enumerable} له مخې، سټ \LR{$A$} هله
  او يوازې هله د شمېر وړ دے چې \LR{$A = \emptyset$} وي يا يوه
  پر هدف سټ بشپړه \LR{$f\colon \PosInt \to A$} شته وي۔ «د شمېر وړ
  سټ» په کره ډول داسې هم تعريفولے شو: سټ هله او يوازې هله
  د شمېر وړ دے چې يوه يو پر يو تابع \LR{$g\colon A \to \PosInt$}
  شته وي۔ وښيئ چې تعريفونه همارزښته دي؛ يعنې وښيئ چې يوه
  يو پر يو تابع \LR{$g\colon A \to \PosInt$} هله او يوازې هله
  شته چې يا \LR{$A = \emptyset$} وي يا يوه پر هدف سټ بشپړه
  \LR{$f\colon \PosInt \to A$} شته وي۔
\end{prob}
\olpseganchor{OLP-0029-B035}{end}

\olpseganchor{OLP-0030-B004}{start}
\olfileid{sfr}{siz}{zigzag}
\olpseganchor{OLP-0030-B004}{end}

\olpseganchor{OLP-0030-B005}{start}
\olsection{د کانتور د زيګزاګ طريقه}
\olpseganchor{OLP-0030-B005}{end}

\olpseganchor{OLP-0030-B006}{start}
\begin{explain}
مخکښې مو ځينې «اسانې» شمېرنې په پام کښې نيولې دي۔ اوس به
يو لږ ګران کار په پام کښې ونيسو۔ د طبيعي عددونو د جوړو سټ
په پام کښې ونيسئ، چې په
\olref[set][pai]{sec} کښې مو داسې تعريف کړے و:
\[
\Nat \times \Nat = \Setabs{\tuple{n,m}}{n,m \in \Nat}
\]
دا مرتبې جوړې په يوه \emph{جدولي ترتيب} کښې داسې برابرولے شو:
\[
\begin{array}{ c | c | c | c | c | c}
& \mathbf 0 & \mathbf 1 & \mathbf 2 & \mathbf 3 & \dots \\
\hline
\mathbf 0 & \tuple{0,0} & \tuple{0,1} & \tuple{0,2} & \tuple{0,3} & \dots \\
\hline
\mathbf 1 & \tuple{1,0} & \tuple{1,1} & \tuple{1,2} & \tuple{1,3} & \dots \\
\hline
\mathbf 2 & \tuple{2,0} & \tuple{2,1} & \tuple{2,2} & \tuple{2,3} & \dots \\
\hline
\mathbf 3 & \tuple{3,0} & \tuple{3,1} & \tuple{3,2} & \tuple{3,3} & \dots \\
\hline
\vdots & \vdots & \vdots & \vdots & \vdots & \ddots\\
\end{array}
\]
په ښکاره، په \LR{$\Nat \times \Nat$} کښې هره مرتبه جوړه به په
دې جدول کښې په دقيقه توګه يو ځل راشي۔ په ځانګړي ډول،
\LR{$\tuple{n,m}$} به په \LR{$n$}م قطار او \LR{$m$}م ستون کښې راشي۔ خو د
داسې جدول غړي په يوه «يو بُعدي» لړ کښې څنګه برابر کړو؟ په
لاندې جدول کښې نمونه د دې کار يوه لار ښيي، که څه هم البته
ډېرې نورې لارې هم شته:
\[
\begin{array}{ c | c | c | c | c | c | c}
& \mathbf 0 & \mathbf 1 & \mathbf 2 & \mathbf 3 & \mathbf 4 &\dots \\
\hline
\mathbf 0 & 0  & 1& 3 & 6& 10 &\ldots \\
\hline
\mathbf 1 &2 & 4& 7 & 11 & \dots &\ldots \\
\hline
\mathbf 2 & 5 & 8 & 12 & \ldots & \dots&\ldots \\
\hline
\mathbf 3 & 9 & 13 & \ldots & \ldots & \dots & \ldots \\
\hline
\mathbf 4 & 14 & \ldots & \ldots & \ldots & \dots & \ldots \\
\hline
\vdots & \vdots & \vdots & \vdots & \vdots&\ldots & \ddots\\
\end{array}
\]\noindent
دا نمونه \emph{د کانتور د زيګزاګ طريقه} بلل کېږي۔ دا د
\LR{$\Nat \times \Nat$} شمېرنه داسې کوي:
\[
\tuple{0,0}, \tuple{0,1}, \tuple{1,0}, \tuple{0,2}, \tuple{1,1},
\tuple{2,0}, \tuple{0,3}, \tuple{1,2}, \tuple{2,1}, \tuple{3,0}, \dots
\]
او له دې لاندې خبره ثابتېږي:
\end{explain}
\olpseganchor{OLP-0030-B006}{end}

\olpseganchor{OLP-0030-B007}{start}
\begin{prop}\ollabel{natsquaredenumerable}
\LR{$\Nat \times \Nat$} يو د شمېر وړ سټ دے۔
\end{prop}
\olpseganchor{OLP-0030-B007}{end}

\olpseganchor{OLP-0030-B008}{start}
\begin{proof}
فرض کړئ \LR{$f \colon \Nat \to \Nat\times\Nat$} هر \LR{$k \in \Nat$}
له هغې مرتبې څوګونې \LR{$\tuple{n,m} \in \Nat \times \Nat$} سره
تړي چې \LR{$k$} د کانتور د زيګزاګ په جدول کښې د \LR{$n$}م قطار او
\LR{$m$}م ستون قيمت وي۔
\end{proof}
\olpseganchor{OLP-0030-B008}{end}

\olpseganchor{OLP-0030-B009}{start}
\begin{explain}
دا تخنيک په ښه ډول عمومي کېدے هم شي۔ د بېلګې په توګه، د
طبيعي عددونو د مرتبو درې ګونو د سټ د شمېرنې دپاره ئې
کارولے شو؛ يعنې:
\[
\Nat \times \Nat \times \Nat = \Setabs{\tuple{n,m,k}}{n,m,k \in \Nat}
\]
\LR{$\Nat \times \Nat \times \Nat$} له \LR{$\Nat$} سره د \LR{$\Nat \times \Nat$}
د کارتيزي حاصل ضرب په توګه ګڼو؛ يعنې:
\[
\Nat^3 = (\Nat \times \Nat) \times \Nat =
\Setabs{\tuple{\tuple{n,m},k}}{n, m, k
  \in \Nat }
\]
نو د يوه محور د \LR{$\Nat$} په شمېرنې او د بل محور د \LR{$\Nat^2$}
په شمېرنې په نښه کولو سره، د يوه جدول په وسيله د \LR{$\Nat^3$}
شمېرنه کولے شو:
\[
\begin{array}{ c | c | c | c | c | c}
& \mathbf 0 & \mathbf 1 & \mathbf 2 & \mathbf 3 & \dots \\
\hline
\mathbf{\tuple{0,0}} & \tuple{0,0,0} & \tuple{0,0,1} & \tuple{0,0,2} & \tuple{0,0,3} & \dots \\
\hline
\mathbf{\tuple{0,1}} & \tuple{0,1,0} & \tuple{0,1,1} & \tuple{0,1,2} & \tuple{0,1,3} & \dots \\
\hline
\mathbf{\tuple{1,0}} & \tuple{1,0,0} & \tuple{1,0,1} & \tuple{1,0,2} & \tuple{1,0,3} & \dots \\
\hline
\mathbf{\tuple{0,2}} & \tuple{0,2,0} & \tuple{0,2,1} & \tuple{0,2,2} & \tuple{0,2,3} & \dots\\
\hline
\vdots & \vdots & \vdots & \vdots & \vdots & \ddots \\
\end{array}
\]
نو د کانتور د زيګزاګ طريقې ته ورته طريقې په کارولو سره د~\LR{$\Nat^3$}
شمېرنه هم ترلاسه کولے شو۔ او همداسې مخ په وړاندې تللے شو،
چې د هر طبيعي عدد \LR{$n$} دپاره د \LR{$\Nat^n$} شمېرنې ترلاسه کړو۔ نو لرو:
\end{explain}
\olpseganchor{OLP-0030-B009}{end}

\olpseganchor{OLP-0030-B010}{start}
\begin{prop}
د هر \LR{$n \in \Nat$} دپاره، \LR{$\Nat^n$} يو د شمېر وړ سټ دے۔
\end{prop}
\olpseganchor{OLP-0030-B010}{end}

\olpseganchor{OLP-0030-B011}{start}
\begin{prob}\label{sfr:siz:zigzag:prob:posint-n}
وښيئ چې د هر \LR{$n \in \Nat$} دپاره، \LR{$(\PosInt)^n$} يو د شمېر وړ سټ دے۔
\end{prob}
\olpseganchor{OLP-0030-B011}{end}

\olpseganchor{OLP-0030-B012}{start}
\begin{prob}\label{sfr:siz:zigzag:prob:posint-star}
  وښيئ چې \LR{$(\PosInt)^*$} يو د شمېر وړ سټ دے۔
  \cref{sfr:siz:zigzag:prob:posint-n} فرضولے شئ۔
\end{prob}
\olpseganchor{OLP-0030-B012}{end}

\olpseganchor{OLP-0031-B004}{start}
\olfileid{sfr}{siz}{pai}
\olsection{جوړوونکې تابعې او کوډونه}
\olpseganchor{OLP-0031-B004}{end}

\olpseganchor{OLP-0031-B005}{start}
\begin{explain}
د کانتور د زيګزاګ طريقه د \LR{$\Nat^n$} د شمېر وړ والے په انځوريز
ډول څرګندوي۔ خو راځئ پر هغه جدول تمرکز وکړو چې \LR{$\Nat^2$} ښيي۔
په جدول کښې د زيګزاګ کرښه تعقيبول او مقامونه شمېرل راښيي چې
\LR{$\tuple{1,2}$} له عدد~\LR{$7$} سره تړل شوې ده۔ خو ښه به وي که دا
په لا نېغه توګه حسابولے شو۔ يعنې ښه به وي چې د زيګزاګ شمېرنې
\emph{معکوس} تابع \LR{$g\colon \Nat^2 \to \Nat$} په لاس کښې ولرو،
داسې چې:
\[
g(\tuple{0,0}) = 0, \;
g(\tuple{0,1}) = 1, \;
g(\tuple{1,0}) = 2, \; \dots,  
g(\tuple{1,2}) = 7, \; \dots
\]
په دې سره په دقيق ډول حسابولے شو چې \LR{$\tuple{n, m}$} به زموږ
په شمېرنه کښې چېرته راشي۔
\olpseganchor{OLP-0031-B005}{end}

\olpseganchor{OLP-0031-B006}{start}
په حقيقت کښې، په دوو مشاهدو سره \LR{$g$} په نېغه توګه تعريفولے شو۔
لومړے: که د \LR{$n$}م قطار او \LR{$m$}م ستون قيمت~\LR{$v$} وي، نو د
\LR{$(n+1)$}م قطار او \LR{$(m-1)$}م ستون قيمت \LR{$v + 1$} دے۔ دويم: زموږ
د شمېرنې لومړے قطار له مثلثي عددونو جوړ دے چې له \LR{$0$}، \LR{$1$}،
\LR{$3$}، \LR{$6$} او داسې نورو پيلېږي۔ \LR{$k$}م مثلثي عدد د هغو طبيعي
عددونو جمع ده چې \LR{$< k$} دي، او په \LR{$k(k+1)/2$} حسابېږي۔ د دې دوو
مشاهدو په يوځاے کولو لاندې تابع په پام کښې ونيسئ:
\[
  g(n,m) = \frac{(n+m+1)(n+m)}{2} + n
\]
ډېر وخت د \LR{$g(\tuple{n, m})$} پر ځاے يوازې \LR{$g(n, m)$} ليکو، ځکه
چې سترګو ته اسانه دے۔ دا لومړے د \LR{$(n+m)^\text{\RL{th}}$} مثلثي عدد
ټاکل او بيا \LR{$n$} ورزياتول درته وايي۔ او دا جدول په هماغه طريقه
ډکوي چې غواړو ئې۔ نو په ځانګړي ډول، جوړه \LR{$\tuple{1, 2}$} له
\LR{$\frac{4 \times 3}{2} + 1 = 7$} سره تړل کېږي۔
\olpseganchor{OLP-0031-B006}{end}

\olpseganchor{OLP-0031-B007}{start}
دا تابع \LR{$g$} د جوړو د يوه سټ د شمېرنې \emph{معکوس} ده۔ داسې
تابعې \emph{جوړوونکې تابعې} بلل کېږي۔
\end{explain}
\olpseganchor{OLP-0031-B007}{end}

\olpseganchor{OLP-0031-B008}{start}
\begin{defn}[جوړوونکې تابع]
  تابع \LR{$f\colon A \times B \to \Nat$} يوه حسابي
  \emph{جوړوونکې تابع} ده که \LR{$f$} يو پر يو وي۔ دا هم وايو چې
  \LR{$f$} د \LR{$A \times B$} \emph{کوډول} کوي، او \LR{$f(x,y)$} د
  \LR{$\tuple{x,y}$} \emph{کوډ} دے۔
\end{defn}
\olpseganchor{OLP-0031-B008}{end}

\olpseganchor{OLP-0031-B009}{start}
\begin{explain}
جوړوونکې تابعې د بېلګې په توګه د طبيعي عددونو د جوړو د
کوډولو دپاره کارولے شو؛ يا په بله وينا، د غړو هره \emph{جوړه}
په يوه \emph{واحد} عدد ښودلے شو۔ د جوړوونکې تابعې په معکوس
عدد بېرته لوستلے شو؛ يعنې معلومولے شو چې کومه جوړه ښيي۔
\end{explain}
\olpseganchor{OLP-0031-B009}{end}

\olpseganchor{OLP-0031-B010}{start}
\begin{prob}
د ټولو غير منفي ناطقو عددونو د سټ يوه شمېرنه ورکړئ۔
\end{prob}
\olpseganchor{OLP-0031-B010}{end}

\olpseganchor{OLP-0031-B011}{start}
\begin{prob}
وښيئ چې \LR{$\Rat$} د شمېر وړ دے۔ په ياد ولرئ چې هر ناطق عدد
د کسر \LR{$z/m$} په توګه ليکل کېدے شي، چې \LR{$z \in \Int$} او
\LR{$m \in \Nat^+$} وي۔
\end{prob}
\olpseganchor{OLP-0031-B011}{end}

\olpseganchor{OLP-0031-B012}{start}
\begin{prob}
د \LR{$\Bin^*$} يوه شمېرنه تعريف کړئ۔
\end{prob}
\olpseganchor{OLP-0031-B012}{end}

\olpseganchor{OLP-0031-B013}{start}
\begin{prob}
د منطق له پېژندنيز کورسه په ياد ولرئ چې د رښتياوالي هر ممکن
جدول د رښتياوالي يوه تابع څرګندوي۔ په بله وينا، د کوم~\LR{$k$}
دپاره د رښتياوالي تابعې ټولې تابعې له \LR{$\Bin^k \to \Bin$} څخه
دي۔ ثابته کړئ چې د رښتياوالي د ټولو تابعو سټ د شمېر وړ دے۔
\end{prob}
\olpseganchor{OLP-0031-B013}{end}

\olpseganchor{OLP-0031-B014}{start}
\begin{prob}
وښيئ چې د يوه خپل سري نامتناهي د شمېر وړ سټ د ټولو
متناهي فرعي سټونو سټ د شمېر وړ دے۔
\end{prob}
\olpseganchor{OLP-0031-B014}{end}

\olpseganchor{OLP-0031-B015}{start}
\begin{prob}
د \LR{$\Nat$} يوه فرعي سټ ته هله \emph{د متناهي متمم لرونکے} ويل
کېږي چې متمم ئې په \LR{$\Nat$} کښې متناهي سټ وي؛ يعنې
\LR{$A \subseteq \Nat$} هله د متناهي متمم لرونکے دے چې
\LR{$\Nat\setminus A$} متناهي وي۔
فرض کړئ \LR{$I$} هغه سټ دے چې
غړي ئې په دقيقه توګه د \LR{$\Nat$} متناهي او د متناهي
متمم لرونکي فرعي سټونه دي۔ وښيئ چې \LR{$I$} د شمېر وړ دے۔
\end{prob}
\olpseganchor{OLP-0031-B015}{end}

\olpseganchor{OLP-0031-B016}{start}
\begin{prob}
وښيئ چې د د شمېر وړ سټونو د شمېر وړ اتحاد
د شمېر وړ دے۔ يعنې، هرکله چې \LR{$A_1$}، \LR{$A_2$}، \dots{} سټونه
وي او هر \LR{$A_i$} د شمېر وړ وي، نو د هغو ټولو اتحاد
\LR{$\bigcup_{i=1}^\infty A_i$} هم د شمېر وړ دے۔ [يادونه:
دا ګران دے!]
\end{prob}
\olpseganchor{OLP-0031-B016}{end}

\olpseganchor{OLP-0031-B017}{start}
\begin{prob}
فرض کړئ \LR{$f \colon A \times B \to \Nat$} يوه خپله سري جوړوونکې
تابع ده او د جوړو سټ ناتش دے۔ وښيئ چې د~\LR{$f$} معکوس پر خپلې
د تعريف ساحې دوه اړخيزه يو پر يو تابع ده چې د قيمتونو سټ ئې
\LR{$A \times B$} دے، او د هغې له ساحې څخه بهر د يوې ثابتې جوړې
په ورکولو د طبيعي عددونو پر ټولو غړو غځېدے شي۔
\end{prob}
\olpseganchor{OLP-0031-B017}{end}

\olpseganchor{OLP-0031-B018}{start}
\begin{prob}
داسې تابع وټاکئ چې \LR{$\Nat^3$} کوډ کړي۔
\end{prob}
\olpseganchor{OLP-0031-B018}{end}

\olpseganchor{OLP-0032-B004}{start}
\olfileid{sfr}{siz}{pai-alt}
\olpseganchor{OLP-0032-B004}{end}

\olpseganchor{OLP-0032-B005}{start}
\olsection{يوه بله جوړوونکې تابع}
\olpseganchor{OLP-0032-B005}{end}

\olpseganchor{OLP-0032-B006}{start}
\begin{explain}
د \LR{$\Nat^2$} داسې نورې شمېرنې هم شته چې د هغو د معکوس تابعو
معلومول اسانه کوي۔ يوه ئې دا ده۔ د شمېرنې پر جدولي ترتيب د
انځورولو پر ځاے، د مثبتو صحيح عددونو له داسې لړ څخه پيل وکړئ
چې په سر کښې له تشو ځايونو سره تړلے وي۔ وانګېړئ چې دا ځايونه
يو په بل پسې په لاندې ډول په جوړو \LR{$\tuple{n,m}$} ډکوو۔ لومړے
له هغو جوړو پيل وکړئ چې په لومړي مقام کښې~\LR{$0$} لري (يعنې جوړې
\LR{$\tuple{0,m}$})؛ لومړۍ جوړه (يعنې \LR{$\tuple{0,0}$}) په لومړي تش
مقام کښې کېږدئ، بيا يو تش مقام پرېږدئ، دويمه جوړه (يعنې
\LR{$\tuple{0,1}$}) په ورپسې تش مقام کښې کېږدئ، بيا يو مقام پرېږدئ،
او همداسې مخکښې لاړ شئ۔
زموږ د شمېرنې (نيمګړے) پيل اوس داسې دے:
\[\small
\begin{array}{@{}c c c c c c c c c c c@{}}
\mathbf 1 & \mathbf 2 & \mathbf 3 & \mathbf 4 & \mathbf 5 & \mathbf 6 & \mathbf 7 & \mathbf 8 & \mathbf 9 & \mathbf{10} & \dots \\ \\
\tuple{0,0} &  & \tuple{0,1} &  & \tuple{0,2} &  & \tuple{0,3} & & \tuple{0,4} &  & \dots \\
\end{array}
\]
همدا کار د \LR{$\tuple{1,m}$} جوړو دپاره په هغو مقامونو کښې تکرار
کړئ چې لا تش پاتې دي؛ بيا هم هر بل تش مقام پرېږدئ:
\[\small
\begin{array}{@{}c c c c c c c c c c c@{}}
\mathbf 1 & \mathbf 2 & \mathbf 3 & \mathbf 4 & \mathbf 5 & \mathbf 6 & \mathbf 7 & \mathbf 8 & \mathbf 9 & \mathbf{10} & \dots \\ \\
\tuple{0,0} & \tuple{1,0} & \tuple{0,1} &  & \tuple{0,2} & \tuple{1,1} & 
\tuple{0,3} & & \tuple{0,4} &  \tuple{1,2} & \dots \\
\end{array}
\]
په همدې طريقه \LR{$\tuple{2,m}$}، \LR{$\tuple{3,m}$} او داسې نورې جوړې
ور دننه کړئ۔
نو زموږ بشپړه شمېرنه
داسې پيلېږي:
\[\small
\begin{array}{@{}cc c c c c c c c c c@{}}
\mathbf 1 & \mathbf 2 & \mathbf 3 & \mathbf 4 & \mathbf 5 & \mathbf 6 & \mathbf 7 & \mathbf 8 & \mathbf 9 & \mathbf{10} & \dots \\ \\
\tuple{0,0} & \tuple{1,0} & \tuple{0,1} & \tuple{2,0}  & \tuple{0,2} & 
\tuple{1,1} & \tuple{0,3} & \tuple{3,0}  & \tuple{0,4} &  \tuple{1,2} & \dots \\
\end{array}
\]
که د پورته جدول حجرې د همدې شمېرنې له مخې وشمېرو، د زيګزاګ
پاکه کرښه به و نۀ مومو، بلکې دا ترتيب به ومومو:
\[
\begin{array}{ c | c | c | c | c | c | c | c }
& \mathbf 0 & \mathbf 1 & \mathbf 2 & \mathbf 3 & \mathbf 4 & \mathbf 5 & \dots \\
\hline
\mathbf 0 & 1 & 3 & 5 & 7 & 9 & 11 & \dots \\
\hline
\mathbf 1 & 2 & 6 & 10 & 14 & 18 & \dots & \dots \\
\hline
\mathbf 2 & 4 & 12 & 20 & 28 & \dots & \dots & \dots \\
\hline
\mathbf 3 & 8 & 24 & 40 & \dots & \dots & \dots & \dots \\
\hline
\mathbf 4 & 16 & 48 & \dots & \dots & \dots & \dots & \dots \\
\hline
\mathbf 5 & 32 & \dots & \dots & \dots & \dots & \dots & \dots \\
\hline
\vdots & \vdots & \vdots & \vdots & \vdots & \vdots & \vdots & \ddots\\
\end{array}
\]
وينو چې په~\LR{$0$} قطار کښې جوړې زموږ د شمېرنې په طاق شمېره
لرونکو مقامونو کښې دي؛ يعنې جوړه \LR{$\tuple{0,m}$} په \LR{$2m+1$} مقام
کښې ده۔ د دويم قطار جوړې، \LR{$\tuple{1,m}$}، په هغو مقامونو کښې دي
چې شمېر ئې د يوه طاق عدد دوه چنده دے، په ځانګړي ډول
\LR{$2 \cdot (2m+1)$}؛ د درېيم قطار جوړې، \LR{$\tuple{2,m}$}، په هغو
مقامونو کښې دي چې شمېر ئې د يوه طاق عدد څلور چنده،
\LR{$4 \cdot (2m+1)$}، دے؛ او همداسې نور۔ په هر قطار کښې د
\LR{$(2m+1)$} ضربي عوامل، \LR{$1$}، \LR{$2$}، \LR{$4$}، \LR{$8$}، \dots، په دقيقه توګه
د~\LR{$2$} طاقتونه دي: \LR{$1= 2^0$}، \LR{$2 = 2^1$}، \LR{$4 = 2^2$}،
\LR{$8 = 2^3$}، \dots\@ په حقيقت کښې، اړوند توان ښودونکے تل د پام
وړ جوړې لومړے غړے وي۔ نو د جوړې \LR{$\tuple{n,m}$} دپاره عامل
\LR{$2^n$} دے۔ په دې توګه عمومي فارمول لاس ته راځي:
\LR{$2^n \cdot (2m+1)$}۔ خو دا له جوړو څخه \emph{مثبتو} صحيح عددونو
ته نقشه ده؛ يعنې \LR{$\tuple{0,0}$} مقام~\LR{$1$} لري۔ که وغواړو له
مقام~\LR{$0$} څخه پيل وکړو، بايد له پايلې څخه~\LR{$1$} منفي کړو۔ نو دا
لاس ته راوړو:
\end{explain}
\olpseganchor{OLP-0032-B006}{end}

\olpseganchor{OLP-0032-B007}{start}
\begin{ex}
تابع \LR{$h\colon \Nat^2 \to \Nat$} چې په لاندې ډول ورکړل شوې ده،
\[
h(n,m) = 2^n (2m+1) - 1
\]
د طبيعي عددونو د جوړو د سټ~\LR{$\Nat^2$} يوه جوړوونکې تابع ده۔
\end{ex}
\olpseganchor{OLP-0032-B007}{end}

\olpseganchor{OLP-0032-B008}{start}
\begin{explain}
له دې سره سم، زموږ د \LR{$\Nat^2$} په دويمه شمېرنه کښې جوړه
\LR{$\tuple{0,0}$} کوډ \LR{$h(0,0) = 2^0(2\cdot 0+1) - 1 = 0$} لري؛
\LR{$\tuple{1,2}$} کوډ \LR{$2^{1} \cdot (2 \cdot 2 + 1) - 1 = 2
\cdot 5 - 1 = 9$} لري؛ او \LR{$\tuple{2,6}$} کوډ \LR{$2^{2} \cdot (2
\cdot 6 + 1) - 1 = 51$} لري۔
\end{explain}
\olpseganchor{OLP-0032-B008}{end}

\olpseganchor{OLP-0032-B009}{start}
ځينې وخت د طبيعي عددونو د جوړو~\LR{$\Nat^2$} کوډول دومره بس وي،
بې له دې چې کوډونه پر هدف سټ بشپړ وي۔ د داسې کوډولو معکوسې
تابعې يوازې جزوي تابعې وي۔
\olpseganchor{OLP-0032-B009}{end}

\olpseganchor{OLP-0032-B010}{start}
\begin{ex}
تابع \LR{$j\colon \Nat^2 \to \Nat^+$} چې په لاندې ډول ورکړل شوې ده،
\[
j(n,m) = 2^n3^m
\]
يوه يو پر يو تابع \LR{$\Nat^2 \to \Nat$} ده۔
\end{ex}
\olpseganchor{OLP-0032-B010}{end}

\olpseganchor{OLP-0033-B004}{start}
\olfileid{sfr}{siz}{nen}
\olpseganchor{OLP-0033-B004}{end}

\olpseganchor{OLP-0033-B005}{start}
\olsection{د شمېر نۀ وړ سټونه}
\olpseganchor{OLP-0033-B005}{end}

\olpseganchor{OLP-0033-B006}{start}
\begin{editorial}
  دا برخه په \olref[enm]{sec} کښې د ورکړل شوي تعريف په کارولو د
  \LR{$\Bin^\omega$} او \LR{$\Pow{\PosInt}$} د شمېر نۀ وړ والے ثابتوي۔ دا
  د \olref[enm-alt]{sec} له بڼې څخه لږه ابتدايي او لږه زياته
  تفصيلي جوړه شوې ده۔
\end{editorial}
\olpseganchor{OLP-0033-B006}{end}

\olpseganchor{OLP-0033-B007}{start}
ځينې سټونه، لکه د مثبتو صحيح عددونو سټ~\LR{$\PosInt$}، نامتناهي دي۔
تر اوسه مو د نامتناهي سټونو داسې بېلګې ليدلې چې ټولې
د شمېر وړ وې۔ خو داسې نامتناهي سټونه هم شته چې دا خاصيت نۀ
لري۔ داسې سټونه \emph{د شمېر نۀ وړ} بلل کېږي۔
\olpseganchor{OLP-0033-B007}{end}

\olpseganchor{OLP-0033-B008}{start}
تر هر څه لومړے ښايي همدا خبره لا حيرانوونکې وي چې د شمېر نۀ وړ
سټونه شته۔ د هر د شمېر وړ سټ~\LR{$A$} دپاره يوه پر هدف سټ بشپړه
تابع \LR{$f \colon \PosInt \to A$} شته۔ که يو سټ د شمېر نۀ وړ وي،
داسې تابع نۀ شته۔ يعنې هېڅ هغه تابع چې د~\LR{$\PosInt$} بې‌شمېره
غړي پر~\LR{$A$} نقشوي، ټول~\LR{$A$} نۀ شي پوره کولے۔ نو په~\LR{$A$}
کښې د بې‌شمېره مثبتو صحيح عددونو په پرتله «زيات» غړي شته۔
\olpseganchor{OLP-0033-B008}{end}

\olpseganchor{OLP-0033-B009}{start}
څنګه به ثابتوو چې يو سټ د شمېر نۀ وړ دے؟ بايد وښيئ چې داسې
پر هدف سټ بشپړه تابع وجود نۀ شي لرلے۔ په مساوي ډول، بايد وښيئ
چې د~\LR{$A$} غړي په يو اړخيز نامتناهي لړ کښې نۀ شي شمېرل کېدے۔ د
دې تر ټولو ښه لار دا ده چې وښيئ د~\LR{$A$} د غړي هر لړ بايد
لږ تر لږه يو غړے پرېږدي؛ يا دا چې هېڅ تابع
\LR{$f\colon \PosInt \to A$} پر هدف سټ بشپړه نۀ شي کېدے۔ دا د کانتور
په \emph{قطري طريقه} کولے شو۔ د~\LR{$A$} د غړي يو لړ، د بېلګې
په ډول \LR{$x_1$}، \LR{$x_2$}، \dots، چې راکړل شي، موږ د~\LR{$A$} يو بل غړے
جوړوو چې د خپل جوړښت له مخې په هېڅ ډول په هغه لړ کښې نۀ شي راتلے۔
\olpseganchor{OLP-0033-B009}{end}

\olpseganchor{OLP-0033-B010}{start}
زموږ لومړۍ بېلګه د ټولو هغو بې‌تشې نامتناهي تسلسلونو سټ
\LR{$\Bin^\omega$} دے چې له \LR{$0$}ګانو او \LR{$1$}ګانو جوړ دي۔
\olpseganchor{OLP-0033-B010}{end}

\olpseganchor{OLP-0033-B011}{start}
\begin{thm}
\ollabel{thm:nonenum-bin-omega}
\LR{$\Bin^\omega$}~د شمېر نۀ وړ دے۔
\end{thm}
\olpseganchor{OLP-0033-B011}{end}

\olpseganchor{OLP-0033-B012}{start}
\begin{proof}
د تناقض دپاره فرض کړئ چې \LR{$\Bin^\omega$} د شمېر وړ دے؛ يعنې
فرض کړئ چې د~\LR{$\Bin^\omega$} د ټولو غړي يو لړ \LR{$s_{1}$}،
\LR{$s_{2}$}، \LR{$s_{3}$}، \LR{$s_{4}$}، \dots{} شته۔ له دې هر يو \LR{$s_i$} پخپله
د \LR{$0$}ګانو او~\LR{$1$}ګانو يو نامتناهي تسلسل دے۔ د دې لړ د \LR{$i$}م تسلسل
\LR{$j$}م غړے به \LR{$s_i(j)$} وبولو۔ نو \LR{$i$}م تسلسل~\LR{$s_i$} دا دے:
\[
s_i(1), s_i(2), s_i(3), \dots
\]
\olpseganchor{OLP-0033-B012}{end}

\olpseganchor{OLP-0033-B013}{start}
دا لړ، او په کښې د هر تسلسل~\LR{$s_i$} غړي، په يوه جدولي ترتيب کښې
اېښودلے شو:
\[
\begin{array}{c|c|c|c|c|c}
& 1 & 2 & 3 & 4 & \dots \\\hline
1 & \mathbf{s_{1}(1)} & s_{1}(2) & s_{1}(3) & s_1(4) & \dots \\\hline
2 & s_{2}(1)& \mathbf{s_{2}(2)} & s_2(3) & s_2(4) & \dots \\\hline
3 & s_{3}(1)& s_{3}(2) & \mathbf{s_3(3)} & s_3(4) & \dots \\\hline
4 & s_{4}(1)& s_{4}(2) & s_4(3) & \mathbf{s_4(4)} & \dots \\\hline
\vdots & \vdots & \vdots & \vdots & \vdots & \mathbf{\ddots}
\end{array}
\]
د څنګ پر خوا نښې په لړ کښې د تسلسلونو شمېرې \LR{$s_1$}، \LR{$s_2$}،
\dots راکوي؛ د پاسنۍ خوا شمېرې د هر يوه تسلسل غړي په
نښه کوي۔ د بېلګې په ډول، \LR{$s_{1}(1)$} د هغه عدد نوم دے چې، که \LR{$0$}
وي يا~\LR{$1$}، په تسلسل~\LR{$s_{1}$} کښې لومړے غړے دے؛ او همداسې
نور۔
\olpseganchor{OLP-0033-B013}{end}

\olpseganchor{OLP-0033-B014}{start}
اوس د \LR{$0$}ګانو او \LR{$1$}ګانو يو داسې نامتناهي تسلسل
\LR{$\overline{s}$} جوړوو چې په هېڅ ډول په دې لړ کښې نۀ شي راتلے۔ د
\LR{$\overline{s}$} تعريف به په لړ \LR{$s_1$}، \LR{$s_2$}، \dots پورې تړلے وي۔
د \LR{$0$}ګانو او \LR{$1$}ګانو د نامتناهي تسلسلونو هر نامتناهي لړ يو داسې
نامتناهي تسلسل~\LR{$\overline{s}$} راپيدا کوي چې په ډاډه توګه په هغه
لړ کښې نۀ راځي۔
\olpseganchor{OLP-0033-B014}{end}

\olpseganchor{OLP-0033-B015}{start}
د \LR{$\overline{s}$} د تعريفولو دپاره د هغه ټول غړي ټاکو؛
يعنې د ټولو \LR{$n \in \PosInt$} دپاره \LR{$\overline{s}(n)$} ټاکو۔ دا کار
د پورته جدولي ترتيب پر قطر د لاندې تللو (ځکه ورته «قطري طريقه»
وايو)، بيا د هر \LR{$1$} په \LR{$0$} او د هر \LR{$0$} په~\LR{$1$} بدلولو کوو۔ په لا
مجرد ډول، \LR{$\overline{s}(n)$} د دې له مخې \LR{$0$} يا \LR{$1$} تعريفوو چې د
قطر \LR{$n$}م غړے، \LR{$s_n(n)$}، \LR{$1$} دے که \LR{$0$}:
\[
\overline{s}(n) =
\begin{cases}
1 & \text{\RL{که \LR{$s_{n}(n) = 0$} وي}}\\
0 & \text{\RL{که \LR{$s_{n}(n) = 1$} وي}}.
\end{cases}
\]
که فارمولونه مو د حالتونو تر تعريف زيات خوښېږي، دا هم تعريفولے
شئ چې \LR{$\overline{s}(n) = 1 - s_n(n)$}۔
\olpseganchor{OLP-0033-B015}{end}

\olpseganchor{OLP-0033-B016}{start}
په څرګند ډول \LR{$\overline{s}$} د \LR{$0$}ګانو او \LR{$1$}ګانو يو نامتناهي
تسلسل دے، ځکه دا يوازې د هغو \LR{$0$}ګانو او \LR{$1$}ګانو مخالف تسلسل دے
چې زموږ د جدولي ترتيب پر قطر راځي۔ نو \LR{$\overline{s}$} د
~\LR{$\Bin^\omega$} يو غړے دے۔ خو دا په لړ \LR{$s_1$}، \LR{$s_2$}،
\dots{} کښې نۀ شي راتلے۔ ولې؟
\olpseganchor{OLP-0033-B016}{end}

\olpseganchor{OLP-0033-B017}{start}
دا په لړ کښې لومړے تسلسل~\LR{$s_1$} نۀ شي کېدے، ځکه په لومړي
غړے کښې له \LR{$s_1$} څخه توپير لري۔ \LR{$s_1(1)$} چې هر څه وي،
موږ \LR{$\overline{s}(1)$} د هغه مخالف تعريف کړے دے۔ دا په لړ کښې
دويم تسلسل نۀ شي کېدے، ځکه \LR{$\overline{s}$} په دويم غړي کښې له
\LR{$s_2$} څخه توپير لري: که \LR{$s_2(2)$}، \LR{$0$} وي، \LR{$\overline{s}(2)$}،
\LR{$1$} دے، او برعکس۔ او همداسې نور۔
\olpseganchor{OLP-0033-B017}{end}

\olpseganchor{OLP-0033-B018}{start}
په لا دقيقه توګه: که \LR{$\overline{s}$} په لړ کښې وے، يو \LR{$k$} به
داسې وے چې \LR{$\overline{s} = s_{k}$}۔ دوه تسلسلونه هله او يوازې
هله يو شان وي چې په هر مقام کښې سره سمون ولري؛ يعنې د هر~\LR{$n$}
دپاره \LR{$\overline{s}(n) = s_{k}(n)$} وي۔ نو په ځانګړي ډول چې
\LR{$n = k$} واخلو، بايد \LR{$\overline{s}(k) = s_{k}(k)$} وي۔ \LR{$s_k(k)$}
يا \LR{$0$} دے يا~\LR{$1$}۔ که \LR{$0$} وي، \LR{$\overline{s}(k)$} بايد~\LR{$1$} وي؛
همداسې مو \LR{$\overline{s}$} تعريف کړے دے۔ خو که \LR{$s_k(k) = 1$} وي،
بيا هم د \LR{$\overline{s}$} د تعريف له مخې \LR{$\overline{s}(k) = 0$}۔
په هر حالت کښې \LR{$\overline{s}(k) \neq s_{k}(k)$}۔
\olpseganchor{OLP-0033-B018}{end}

\olpseganchor{OLP-0033-B019}{start}
موږ د~\LR{$\Bin^\omega$} د غړي د يوه لړ \LR{$s_1$}، \LR{$s_2$}،
\dots{} په فرضولو پيل وکړ۔ له دې لړ څخه مو تسلسل~\LR{$\overline{s}$}
جوړ کړ چې ثابته مو کړه په لړ کښې نۀ شي راتلے۔ خو که ټول \LR{$s_i$}
د \LR{$0$}ګانو او \LR{$1$}ګانو تسلسلونه وي، دا خامخا د \LR{$0$}ګانو او
\LR{$1$}ګانو يو تسلسل دے؛ يعنې \LR{$\overline{s} \in \Bin^\omega$}۔ دا په
ځانګړي ډول ښيي چې د~\LR{$\Bin^\omega$} د \emph{ټولو} غړي هېڅ
لړ نۀ شي کېدے، ځکه د هر داسې لړ دپاره يو داسې تسلسل
~\LR{$\overline{s}$} هم جوړولے شو چې په ډاډه توګه په لړ کښې نۀ راځي؛
نو دا فرض چې په~\LR{$\Bin^\omega$} کښې د ټولو تسلسلونو يو لړ شته،
تناقض ته رسېږي۔
\end{proof}
\olpseganchor{OLP-0033-B019}{end}

\olpseganchor{OLP-0033-B020}{start}
\begin{explain}
د ثبوت دې طريقې ته «قطرول» ويل کېږي، ځکه د~\LR{$\overline{s}$} د
تعريف دپاره د جدولي ترتيب قطر کاروي۔ په قطرولو کښې د جدول شتون
لازم نۀ دے: په ورته مفکوره ښودلے شو چې سټونه د شمېر وړ نۀ
دي، ان هغه وخت چې هېڅ جدول او رښتينے قطر نۀ وي۔
\end{explain}
\olpseganchor{OLP-0033-B020}{end}

\olpseganchor{OLP-0033-B021}{start}
\begin{thm}
\ollabel{thm:nonenum-pownat}
\LR{$\Pow{\PosInt}$} د شمېر وړ نۀ دے۔
\end{thm}
\olpseganchor{OLP-0033-B021}{end}

\olpseganchor{OLP-0033-B022}{start}
\begin{proof}
په هماغه طريقه مخکښې ځو: ښيو چې د~\LR{$\PosInt$} د فرعي سټونو د هر
لړ دپاره د \LR{$\PosInt$} يو داسې فرعي سټ شته چې په هغه لړ کښې نۀ
شي راتلے۔ فرض کړئ دا د~\LR{$\PosInt$} د فرعي سټونو يو ورکړل شوے لړ دے:
\[
Z_{1}, Z_{2}, Z_{3}, \dots
\]
اوس يو سټ \LR{$\overline{Z}$} داسې تعريفوو چې د هر \LR{$n \in \PosInt$}
دپاره، \LR{$n \in \overline{Z}$} هله او يوازې هله چې \LR{$n \notin Z_{n}$}:
\[
\overline{Z} = \Setabs{n \in \PosInt}{n \notin Z_n}
\]
\LR{$\overline{Z}$} په څرګند ډول د مثبتو صحيح عددونو يو سټ دے، ځکه
د فرض له مخې هر~\LR{$Z_n$} داسې دے؛ او له دې امله
\LR{$\overline{Z} \in \Pow{\PosInt}$}۔ خو \LR{$\overline{Z}$} په لړ کښې نۀ
شي راتلے۔ د دې د ښودلو دپاره به ثابته کړو چې د هر
\LR{$k \in \PosInt$} دپاره \LR{$\overline{Z} \neq Z_k$}۔
\olpseganchor{OLP-0033-B022}{end}

\olpseganchor{OLP-0033-B023}{start}
نو \LR{$k \in \PosInt$} خپله سرے ونيسئ۔ موږ \LR{$\overline{Z}$} داسې تعريف
کړے چې د هر \LR{$n \in \PosInt$} دپاره، \LR{$n \in \overline{Z}$} هله او
يوازې هله چې \LR{$n \notin Z_n$}۔ په ځانګړي ډول چې \LR{$n=k$} واخلو،
\LR{$k \in \overline{Z}$} هله او يوازې هله چې \LR{$k \notin Z_k$}۔ خو دا
ښيي چې \LR{$\overline{Z} \neq Z_k$}، ځکه \LR{$k$} د يوه سټ غړے
دے خو د بل نۀ دے، او له دې امله \LR{$\overline{Z}$} او \LR{$Z_k$} بېل
غړي لري۔ ځکه چې \LR{$k$} خپله سرے و، \LR{$\overline{Z}$} په لړ
\LR{$Z_1$}، \LR{$Z_2$}، \dots کښې نۀ دے۔
\end{proof}
\olpseganchor{OLP-0033-B023}{end}

\olpseganchor{OLP-0033-B024}{start}
\begin{explain}
تېر ثبوت د قطر نوم وانۀ خيست، خو که داسې ئې انځور کړئ، قطري
طريقه په کښې ليدلے شئ: وانګېړئ چې سټونه \LR{$Z_1$}، \LR{$Z_2$}، \dots
په يوه جدولي ترتيب کښې ليکل شوي، چې هر غړے~\LR{$j \in Z_i$}
په~\LR{$j$}م ستون کښې ليکل کېږي۔ فرض کړئ په هغه لړ کښې لومړي څلور
سټونه \LR{$\{1,2,3,\dots\}$}، \LR{$\{2, 4, 6, \dots\}$}، \LR{$\{1,2,5\}$} او
\LR{$\{3,4,5,\dots\}$} دي۔ نو جدولي ترتيب به داسې پيل شي:
\[
\begin{array}{r@{}rrrrrrr}
  Z_1 = \{ & \mathbf{1}, & 2, & 3, & 4, & 5, & 6, & \dots\}\\
  Z_2 = \{ &  & \mathbf{2}, &  & 4, &  & 6, & \dots\}\\
  Z_3 = \{ & 1, & 2, &  &  & 5\phantom{,} &  & \}\\
  Z_4 = \{ &  &  & 3, & \mathbf{4}, & 5, & 6, & \dots\}\\
  \vdots & & & & & \ddots
\end{array}
\]
بيا \LR{$\overline{Z}$} هغه سټ دے چې د قطر پر لاندې تلو جوړېږي: هغه
شمېرې پرېږدو چې پر قطر راځي، او هغه \LR{$j$}ګانې ور شاملوو چې د جدول
په \LR{$j$}م قطار/ستون کښې تش ځاے لري۔ په پورته حالت کښې به \LR{$1$} او
\LR{$2$} پرېږدو،~\LR{$3$} به ور شاملوو،~\LR{$4$} به پرېږدو، او همداسې نور۔
\end{explain}
\olpseganchor{OLP-0033-B024}{end}

\olpseganchor{OLP-0033-B025}{start}
\begin{prob}
په قطري استدلال وښيئ چې \LR{$\Pow{\Nat}$} د شمېر نۀ وړ دے۔
\end{prob}
\olpseganchor{OLP-0033-B025}{end}

\olpseganchor{OLP-0033-B026}{start}
\begin{prob}\label{sfr:siz:nen:prob:f-posint}
په څرګند قطري استدلال وښيئ چې د تابعو سټ
\LR{$f \colon \PosInt \to \PosInt$} د شمېر نۀ وړ دے۔ يعنې وښيئ
چې که \LR{$f_1$}، \LR{$f_2$}، \dots د تابعو يو لړ وي او هره
\LR{$f_i\colon \PosInt \to \PosInt$} وي، نو يوه
\LR{$\overline{f}\colon \PosInt \to \PosInt$} شته چې په دې لړ کښې نۀ ده۔
\end{prob}
\olpseganchor{OLP-0033-B026}{end}

\olpseganchor{OLP-0034-B004}{start}
\olfileid{sfr}{siz}{red}
\olpseganchor{OLP-0034-B004}{end}

\olpseganchor{OLP-0034-B005}{start}
\olsection{راکمونه}
\olpseganchor{OLP-0034-B005}{end}

\olpseganchor{OLP-0034-B006}{start}
\begin{editorial}
  دا برخه په راکمونه د شمېر نۀ وړ والے ثابتوي، او پايلې ئې د
  \olref[nen]{sec} له پايلو سره سمون خوري۔ يوه بله، لږه لنډه
  بڼه چې د \olref[nen-alt]{sec} له پايلو سره سمون خوري، په
  \olref[red-alt]{sec} کښې ورکړل شوې ده۔
\end{editorial}
\olpseganchor{OLP-0034-B006}{end}

\olpseganchor{OLP-0034-B007}{start}
موږ په قطري استدلال وښودل چې \LR{$\Pow{\PosInt}$} د شمېر نۀ وړ
دے۔ له دې مخکښې مو ثابته کړې وه چې \LR{$\Bin^\omega$}، يعنې د
\LR{$0$}ګانو او \LR{$1$}ګانو د ټولو نامتناهي تسلسلونو سټ،
د شمېر نۀ وړ دے۔ دا د \LR{$\Pow{\PosInt}$} د د شمېر نۀ وړ
کېدو د ثابتولو يوه بله لار ده: وښيئ چې \emph{که \LR{$\Pow{\PosInt}$}
د شمېر وړ وي، نو \LR{$\Bin^\omega$} هم د شمېر وړ دے}۔ ځکه
پوهېږو چې \LR{$\Bin^\omega$} د شمېر وړ نۀ دے، \LR{$\Pow{\PosInt}$}
هم داسې نۀ شي کېدے۔ دې ته د يوې مسئلې بلې ته \emph{راکمونه}
ويل کېږي؛ په دې حالت کښې د \LR{$\Bin^\omega$} د شمېرلو مسئله د
\LR{$\Pow{\PosInt}$} د شمېرلو مسئلې ته راکموو۔ د وروستۍ مسئلې يو حل،
يعنې د \LR{$\Pow{\PosInt}$} يوه شمېرنه، به د لومړۍ مسئلې حل، يعنې د
\LR{$\Bin^\omega$} يوه شمېرنه، راکړي۔
\olpseganchor{OLP-0034-B007}{end}

\olpseganchor{OLP-0034-B008}{start}
د يوه سټ~\LR{$B$} د شمېرلو مسئله څنګه د يوه سټ~\LR{$A$} د شمېرلو مسئلې
ته راکموو؟ د~\LR{$A$} د يوې شمېرنې د~\LR{$B$} په شمېرنه د اړولو يوه لار
ورکوو۔ د دې تر ټولو اسانه لار د يوې پر هدف سټ بشپړه تابع
\LR{$f\colon A \to B$} تعريفول دي۔ که \LR{$x_1$}، \LR{$x_2$}، \dots{} د~\LR{$A$}
شمېرنه وي، نو \LR{$f(x_1)$}، \LR{$f(x_2)$}، \dots{} به د~\LR{$B$} شمېرنه وي۔
زموږ په حالت کښې د يوې پر هدف سټ بشپړې تابع
\LR{$f\colon \Pow{\PosInt} \to \Bin^\omega$} په لټه کښې يو۔
\olpseganchor{OLP-0034-B008}{end}

\olpseganchor{OLP-0034-B009}{start}
\begin{prob}
وښيئ چې که يوه يو پر يو تابع \LR{$g\colon B \to A$} وي او
\LR{$B$}~د شمېر نۀ وړ وي، نو \LR{$A$} هم همداسې دے۔ دا په ښودلو
وکړئ چې څنګه د~\LR{$g$} په کارولو د~\LR{$A$} يوه شمېرنه د~\LR{$B$} په شمېرنه
اړولے شئ۔
\end{prob}
\olpseganchor{OLP-0034-B009}{end}

\olpseganchor{OLP-0034-B010}{start}
\begin{proof}[په راکمونه د {\olref[nen]{thm:nonenum-pownat}} ثبوت]
فرض کړئ چې \LR{$\Pow{\PosInt}$} د شمېر وړ وے، او له دې امله د
هغه يوه شمېرنه \LR{$Z_{1}$}، \LR{$Z_{2}$}، \LR{$Z_{3}$}، \dots وے۔
\olpseganchor{OLP-0034-B010}{end}

\olpseganchor{OLP-0034-B011}{start}
تابع \LR{$f \colon \Pow{\PosInt} \to \Bin^\omega$} داسې تعريف کړئ
چې \LR{$f(Z)$} هغه تسلسل \LR{$s$} وي چې \LR{$s(n) = 1$} هله او يوازې هله
چې \LR{$n \in Z$} وي، او په بل حالت کښې \LR{$s(n) = 0$} وي۔
دا په څرګند ډول تابع تعريفوي، ځکه
کله چې \LR{$Z \subseteq \PosInt$} وي، هر \LR{$n \in \PosInt$} يا د \LR{$Z$}
يو غړے وي يا نۀ وي۔ د بېلګې په ډول، د مثبتو جفتو
عددونو سټ \LR{$2\PosInt = \{2, 4, 6, \dots\}$} تسلسل
\LR{$010101\dots$} ته نقشېږي، تش سټ \LR{$0000\dots$} ته او پخپله سټ
\LR{$\PosInt$}، \LR{$1111\dots$} ته نقشېږي۔
\olpseganchor{OLP-0034-B011}{end}

\olpseganchor{OLP-0034-B012}{start}
دا پر هدف سټ بشپړه هم ده: د \LR{$0$}ګانو او \LR{$1$}ګانو هر تسلسل د مثبتو
صحيح عددونو له يوه سټ سره سمون خوري، يعنې له هغه سټ سره چې غړي
ئې هغه صحيح عددونه دي کوم چې د تسلسل د~\LR{$1$} لرونکو مقامونو سره
سمون خوري۔ په لا دقيقه توګه، فرض کړئ \LR{$s \in \Bin^\omega$}۔
\LR{$Z \subseteq \PosInt$} په لاندې ډول تعريف کړئ:
\[
Z = \Setabs{n \in \PosInt}{s(n) = 1}
\]
بيا \LR{$f(Z) = s$}، لکه چې د~\LR{$f$} تعريف ته په کتلو تاييدولے شو۔
\olpseganchor{OLP-0034-B012}{end}

\olpseganchor{OLP-0034-B013}{start}
اوس دا لړ په پام کښې ونيسئ:
\[
f(Z_1), f(Z_2), f(Z_3), \dots
\]
ځکه چې \LR{$f$} پر هدف سټ بشپړه ده، د \LR{$\Bin^\omega$} هر غړے بايد د
کوم ورودي دپاره د~\LR{$f$} د قيمت په توګه راشي، او نو بايد په لړ کښې
راشي۔ له دې امله دا لړ بايد ټول~\LR{$\Bin^\omega$} وشمېري۔
\olpseganchor{OLP-0034-B013}{end}

\olpseganchor{OLP-0034-B014}{start}
نو که \LR{$\Pow{\PosInt}$} د شمېر وړ وے، \LR{$\Bin^\omega$} به
د شمېر وړ وے۔ خو \LR{$\Bin^\omega$} د شمېر نۀ وړ دے
(\olref[nen]{thm:nonenum-bin-omega})۔ له دې امله
\LR{$\Pow{\PosInt}$} د شمېر نۀ وړ دے۔
\end{proof}
\olpseganchor{OLP-0034-B014}{end}

\olpseganchor{OLP-0034-B015}{start}
\begin{explain}
د راکمونې په لوري کښې ګډوډېدل اسانه دي۔ د بېلګې په ډول، يوه
پر هدف سټ بشپړه تابع \LR{$g \colon \Bin^\omega \to B$} دا \emph{نۀ}
ثابتوي چې \LR{$B$} د شمېر نۀ وړ دے۔ (تابع
\LR{$g \colon \Bin^\omega \to \Bin$} په پام کښې ونيسئ چې په
\LR{$g(s) = s(1)$} تعريف شوې؛ دا تابع د \LR{$0$}ګانو او \LR{$1$}ګانو يو تسلسل
د هغه لومړي غړے ته نقشوي۔ دا پر هدف سټ بشپړه ده، ځکه ځينې
تسلسلونه په \LR{$0$} او ځينې په \LR{$1$} پيلېږي۔ خو \LR{$\Bin$} متناهي دے۔)
دا هم په ياد ولرئ چې تابع~\LR{$f$} بايد پر هدف سټ بشپړه وي، ګنې
استدلال مخکښې نۀ ځي: بيا به دا ډاډ نۀ وي چې \LR{$f(x_1)$}، \LR{$f(x_2)$}،
\dots{} د~\LR{$B$} ټول غړي لري۔ د بېلګې په ډول،
\[
h(n) = \underbrace{000\dots0}_{\text{\RL{\LR{$n$} \LR{$0$}ګانې}}}111\dots
\]
تابع \LR{$h\colon \PosInt \to \Bin^\omega$} تعريفوي، خو \LR{$\PosInt$}
د شمېر وړ دے۔
له لومړنيو صفرونو وروسته ټول رقمونه يو دي، نو هر حاصل
يو نامتناهي توريز تسلسل دے۔
\end{explain}
\olpseganchor{OLP-0034-B015}{end}

\olpseganchor{OLP-0034-B016}{start}
\begin{prob}
په راکمونې استدلال وښيئ چې د مثبتو صحيح عددونو د جوړو د
\emph{سټونو} سټ د شمېر نۀ وړ دے۔
\end{prob}
\olpseganchor{OLP-0034-B016}{end}

\olpseganchor{OLP-0034-B017}{start}
\begin{prob}\label{sfr:siz:red:prob:nat-nat}
  په راکمونې استدلال وښيئ چې د ټولو تابعو \LR{$f\colon \Nat \to \Nat$}
  سټ~\LR{$X$} د شمېر نۀ وړ دے۔ (لارښوونه: له \LR{$X$} څخه
  ~\LR{$\Bin^\omega$} ته يوه پر هدف سټ بشپړه تابع ورکړئ۔)
\end{prob}
\olpseganchor{OLP-0034-B017}{end}

\olpseganchor{OLP-0034-B018}{start}
\begin{prob}
وښيئ چې \LR{$\Nat^\omega$}، يعنې د طبيعي عددونو د نامتناهي
تسلسلونو سټ، په راکمونې استدلال د شمېر نۀ وړ دے۔
\end{prob}
\olpseganchor{OLP-0034-B018}{end}

\olpseganchor{OLP-0034-B019}{start}
\begin{prob}
فرض کړئ \LR{$P$} له مثبتو صحيح عددونو څخه سټ~\LR{$\{0\}$} ته د تابعو سټ
دے، او \LR{$Q$} له مثبتو صحيح عددونو څخه سټ~\LR{$\{0\}$} ته د
\emph{جزوي} تابعو سټ دے۔ وښيئ چې \LR{$P$}~د شمېر وړ دے او \LR{$Q$}
نۀ دے۔ (لارښوونه: د \LR{$\Bin^\omega$} د شمېرلو مسئله د~\LR{$Q$} شمېرلو
ته راکمه کړئ۔)
\end{prob}
\olpseganchor{OLP-0034-B019}{end}

\olpseganchor{OLP-0034-B020}{start}
\begin{prob}
فرض کړئ \LR{$S$} له مثبتو صحيح عددونو څخه سټ~\LR{$\{0,1\}$} ته د ټولو
پر هدف سټ بشپړه تابعو سټ دے؛ يعنې \LR{$S$} له ټولو پر هدف سټ بشپړه
~\LR{$f\colon \PosInt \to \Bin$} څخه جوړ دے۔ وښيئ چې \LR{$S$}
د شمېر نۀ وړ دے۔
\end{prob}
\olpseganchor{OLP-0034-B020}{end}

\olpseganchor{OLP-0034-B021}{start}
\begin{prob}
وښيئ چې د ټولو حقيقي عددونو سټ~\LR{$\Real$} د شمېر نۀ وړ دے۔
\end{prob}
\olpseganchor{OLP-0034-B021}{end}

\olpseganchor{OLP-0035-B004}{start}
\olfileid{sfr}{siz}{equ}
\olpseganchor{OLP-0035-B004}{end}

\olpseganchor{OLP-0035-B005}{start}
\olsection{همشمېر والی}
\olpseganchor{OLP-0035-B005}{end}

\olpseganchor{OLP-0035-B006}{start}
موږ د سټونو د «اندازې» يوه شهودي مفکوره لرو چې د متناهي سټونو
دپاره سمه کار کوي۔ خو د نامتناهي سټونو څه؟ که وغواړو د \emph{هرې}
اندازې د دوو سټونو د اندازو د پرتله کولو يوه صوري لار جوړه کړو،
ښه به وي چې لومړے تعريف کړو سټونه کله يوه اندازه لري۔ دا
د فرېګه د ۱۸۸۴م کال د جرمني متن وينا ده:
\begin{quote}
  که د خوراک مېز ته يو خدمتګار د پلېټونو هومره چاړې اېښودل غواړي او
  غواړي چې په دې شمېر ډاډه شي، د هېڅ يوه شمېرلو ته اړتيا نۀ لري۔
  د هر پلېټ ښي اړخ ته دې يوه چاړه کېږدي، داسې چې د مېز هره چاړه
  د يو پلېټ ښي اړخ ته پرته وي۔ په دې ډول د هماغې مکاني اړيکې له
  مخې د هر پلېټ چاړه او د هرې چاړې پلېټ په يوازېنۍ توګه ټاکل کېږي۔
  \mbox{\LR{\latinfont \citep[\S70]{Frege1884}}}
\end{quote}
د دې عبارت بصيرت په يوه صوري تعريف کښې څرګندولے شو:
\olpseganchor{OLP-0035-B006}{end}

\olpseganchor{OLP-0035-B007}{start}
\begin{defn}\ollabel{comparisondef}
  \LR{$A$} له \LR{$B$} سره \emph{همشمېره} دے، چې \LR{$\cardeq{A}{B}$} ئې ليکو،
  هله او يوازې هله چې يوه دوه اړخيزه يو پر يو تابع \LR{$f \colon A \to B$} وي۔
\end{defn}
\olpseganchor{OLP-0035-B007}{end}

\olpseganchor{OLP-0035-B008}{start}
\begin{prop}\ollabel{equinumerosityisequi}
همشمېر والی يوه د معادلتوب اړيکه ده۔
\end{prop}
\olpseganchor{OLP-0035-B008}{end}

\olpseganchor{OLP-0035-B009}{start}
\begin{proof} 
بايد وښيو چې همشمېر والی انعکاسي، متناظر او انتقالي دے۔ فرض
کړئ \LR{$A, B$} او \LR{$C$} سټونه دي۔
\olpseganchor{OLP-0035-B009}{end}

\olpseganchor{OLP-0035-B010}{start}
\emph{انعکاسيتوب۔} عيني نقشه \LR{$\Id{A} \colon A \to A$}، چې د ټولو
\LR{$x \in A$} دپاره \LR{$\Id{A} (x) = x$} وي، يوه دوه اړخيزه يو پر يو تابع ده۔ نو
\LR{$\cardeq{A}{A}$}۔
\olpseganchor{OLP-0035-B010}{end}

\olpseganchor{OLP-0035-B011}{start}
\emph{تناظر۔} فرض کړئ \LR{$\cardeq{A}{B}$}؛ يعنې يوه دوه اړخيزه يو پر يو تابع
\LR{$f\colon A \to B$} شته۔ ځکه چې \LR{$f$} دوه اړخيزه يو پر يو ده، د هغې معکوس
\LR{$f^{-1}$} شته او هغه هم دوه اړخيزه يو پر يو دے۔ له دې امله
\LR{$f^{-1}\colon B \to A$} يوه دوه اړخيزه يو پر يو تابع ده، نو \LR{$\cardeq{B}{A}$}۔
\olpseganchor{OLP-0035-B011}{end}

\olpseganchor{OLP-0035-B012}{start}
\emph{انتقاليتوب۔} فرض کړئ چې \LR{$\cardeq{A}{B}$} او
\LR{$\cardeq{B}{C}$}؛ يعنې دوه اړخيزې يو پر يو تابعې \LR{$f\colon A \to B$} او
\LR{$g\colon B \to C$} شته۔ بيا ترکيب \LR{$\comp{f}{g}\colon A \to C$}
دوه اړخيزه يو پر يو دے، نو \LR{$\cardeq{A}{C}$}۔
\end{proof}
\olpseganchor{OLP-0035-B012}{end}

\olpseganchor{OLP-0035-B013}{start}
\begin{prop}
که \LR{$\cardeq{A}{B}$} وي، نو \LR{$A$} هله او يوازې هله د شمېر وړ
دے چې \LR{$B$} داسې وي۔
\end{prop}
\olpseganchor{OLP-0035-B013}{end}

\olpseganchor{OLP-0035-B014}{start}
\begin{editorial}
لاندې ثبوت که \olref[enm]{sec} شامل وي،
\olref[enm]{defn:enumerable} کاروي، او که نۀ وي
\olref[enm-alt]{defn:enumerable} کاروي۔
\end{editorial}
\olpseganchor{OLP-0035-B014}{end}

\olpseganchor{OLP-0035-B015}{start}
\begin{proof}
فرض کړئ \LR{$\cardeq{A}{B}$}، نو کومه دوه اړخيزه يو پر يو تابع \LR{$f \colon A
\to B$} شته، او فرض کړئ چې \LR{$A$} د شمېر وړ دے۔

  بيا يا \LR{$A = \emptyset$} دے، يا يوه پر هدف سټ بشپړه تابع
  \LR{$g\colon \PosInt \to A$} شته۔ که \LR{$A = \emptyset$} وي، نو
  \LR{$B = \emptyset$} هم دے (ګنې يو غړے~\LR{$y \in B$} به وے،
  خو هېڅ \LR{$x \in A$} به نۀ وے چې \LR{$f(x) = y$})۔
  له بلې خوا، که \LR{$g\colon \PosInt \to A$}
  پر هدف سټ بشپړه وي، نو \LR{$\comp{g}{f} \colon \PosInt \to B$}
  پر هدف سټ بشپړه ده۔ د دې د ليدلو دپاره \LR{$y \in B$} ونيسئ۔ ځکه
  چې \LR{$f$} پر هدف سټ بشپړه ده، يو \LR{$x \in A$} شته چې \LR{$f(x) = y$}۔ ځکه
  چې \LR{$g$} پر هدف سټ بشپړه ده، يو \LR{$n \in \PosInt$} شته چې \LR{$g(n) =
  x$}۔ له دې امله
\[
(\comp{g}{f})(n) = f(g(n)) = f(x) = y
\]
او نو \LR{$\comp{g}{f}$} پر هدف سټ بشپړه ده۔ موږ لرو چې
\LR{$\comp{g}{f}$} د~\LR{$B$} يوه شمېرنه ده، نو \LR{$B$}~د شمېر وړ دے۔
\olpseganchor{OLP-0035-B015}{end}

\olpseganchor{OLP-0035-B016}{start}
که \LR{$B$} د شمېر وړ وي، د~\LR{$f$} پر ځاے د دوه اړخيزه يو پر يو تابع
\LR{$f^{-1}\colon B \to A$} په کارولو د استدلال په تکرارولو مومو چې
\LR{$A$} د شمېر وړ دے۔
\end{proof}
\olpseganchor{OLP-0035-B016}{end}

\olpseganchor{OLP-0035-B017}{start}
\begin{prob}
وښيئ چې که \LR{$\cardeq{A}{C}$} او \LR{$\cardeq{B}{D}$} وي، او
\LR{$A \cap B = C \cap D = \emptyset$} وي، نو
\LR{$\cardeq{A \cup B}{C \cup D}$}۔
\end{prob}
\olpseganchor{OLP-0035-B017}{end}

\olpseganchor{OLP-0035-B018}{start}
\begin{prob}
وښيئ چې که \LR{$A$} نامتناهي او د شمېر وړ وي، نو
\LR{$\cardeq{A}{\Nat}$}۔
\end{prob}
\olpseganchor{OLP-0035-B018}{end}

\olpseganchor{OLP-0036-B004}{start}
\olfileid{sfr}{siz}{car}
\olpseganchor{OLP-0036-B004}{end}

\olpseganchor{OLP-0036-B005}{start}
\olsection{د بېلو اندازو سټونه او د کانتور قضيه}
\olpseganchor{OLP-0036-B005}{end}

\olpseganchor{OLP-0036-B006}{start}
\begin{explain}
موږ دا مفکوره په دقيقه وينا کښې وړاندې کړې چې دوه سټونه يوه
اندازه لري۔ دا مفکوره هم په دقيقه وينا کښې وړاندې کولے شو چې يو
سټ له بل څخه کوچنے دے۔ زموږ د «له ... څخه کوچنے (يا همشمېره)
دے» تعريف به د سټونو تر منځ د دوه اړخيزه يو پر يو تابع پر ځاے، له لومړي
سټ څخه دويم ته يوه يو پر يو تابع وغواړي۔ که داسې تابع وي، د
لومړي سټ اندازه د دويم سټ له اندازې څخه کوچنۍ يا ورسره مساوي ده۔
په شهودي ډول، له يوه سټ څخه بل ته يو پر يو تابع دا ډاډ راکوي
چې د تابعې د قيمتونو سټ لږ تر لږه د تعريف د ساحې هومره
غړي لري، ځکه د تعريف د ساحې هېڅ دوه غړي د
قيمتونو د سټ يوه غړے ته نۀ نقشېږي۔
\end{explain}
\olpseganchor{OLP-0036-B006}{end}

\olpseganchor{OLP-0036-B007}{start}
\begin{defn}
\LR{$A$} له~\LR{$B$} څخه \emph{لوے نۀ دے}، چې \LR{$\cardle{A}{B}$} ئې ليکو،
هله او يوازې هله چې يوه يو پر يو تابع \LR{$f \colon A \to B$} وي۔
\end{defn}
\olpseganchor{OLP-0036-B007}{end}

\olpseganchor{OLP-0036-B008}{start}
څرګنده ده چې دا يوه انعکاسي او انتقالي اړيکه ده، خو متناظره نۀ
ده (دا د تمرين په توګه پرېښوول شوې ده)۔ يوه داسې مفکوره هم
معرفي کولے شو چې وايي يو سټ له بل څخه (په ټينګه) کوچنے دے۔
\olpseganchor{OLP-0036-B008}{end}

\olpseganchor{OLP-0036-B009}{start}
\begin{defn}
\LR{$A$} له~\LR{$B$} څخه \emph{کوچنے دے}، چې \LR{$\cardless{A}{B}$} ئې ليکو،
هله او يوازې هله چې يوه يو پر يو تابع~\LR{$f\colon A \to B$} وي،
خو هېڅ دوه اړخيزه يو پر يو تابع~\LR{$g\colon A \to B$} نۀ وي؛ يعنې
\LR{$\cardle{A}{B}$} او \LR{$\cardneq{A}{B}$} وي۔
\end{defn}
\olpseganchor{OLP-0036-B009}{end}

\olpseganchor{OLP-0036-B010}{start}
څرګنده ده چې دا اړيکه ناانعکاسي او انتقالي ده۔ (دا د تمرين په
توګه پرېښوول شوې ده۔) په دې نښه ويلے شو چې يو سټ \LR{$A$} هله او
يوازې هله د شمېر وړ دے چې \LR{$\cardle{A}{\Nat}$} وي، او \LR{$A$}
هله او يوازې هله د شمېر نۀ وړ دے چې
\LR{$\cardless{\Nat}{A}$} وي۔
وروستے هم‌ارز بيان پر دې فرض ولاړ دے چې هر نامتناهي سټ
د طبيعي عددونو يوه يو پر يو کاپي لري۔ د ټاکنې اصل دا فرض
تضمينوي۔ په دې سره دا بيا ويلے شو:
%
\olref[sfr][siz][nen-alt]{thm:nonenum-pownat}
د دې مشاهدې په توګه چې
\LR{$\cardless{\Nat}{\Pow{\Nat}}$}۔ په
حقيقت کښې، \mbox{\LR{\latinfont \citet{Cantor1892}}} ثابته کړه چې دا وروستۍ خبره
\emph{بشپړه عمومي} ده:
\olpseganchor{OLP-0036-B010}{end}

\olpseganchor{OLP-0036-B011}{start}
\begin{thm}[کانتور]\ollabel{thm:cantor}
د هر سټ \LR{$A$} دپاره \LR{$\cardless{A}{\Pow{A}}$}۔
\end{thm}
\olpseganchor{OLP-0036-B011}{end}

\olpseganchor{OLP-0036-B012}{start}
\begin{proof}
نقشه \LR{$f(x) = \{x\}$} يوه يو پر يو تابع
\LR{$f \colon A \to \Pow{A}$} ده، ځکه که \LR{$x \neq y$} وي، نو د غړو له
مخې د برابرۍ په سبب \LR{$\{x\} \neq \{y\}$} هم وي، او نو
\LR{$f(x) \neq f(y)$}۔ له دې امله \LR{$\cardle{A}{\Pow{A}}$} لرو۔
\olpseganchor{OLP-0036-B012}{end}

\olpseganchor{OLP-0036-B013}{start}
\begin{editorial}
که \olref[nen]{sec} موجوده وي، تفصيلي ثبوت وړاندې کوو؛ ګنې يو
چټک ثبوت وړاندې کوو چې د \olref[nen-alt]{sec} سره سمون خوري۔
\end{editorial}
\olpseganchor{OLP-0036-B013}{end}

\olpseganchor{OLP-0036-B014}{start}
% 
  اوس به وښيو چې هېڅ پر هدف سټ بشپړه تابع~\LR{$g\colon A \to
  \Pow{A}$} نۀ شي کېدے، نو يوه دوه اړخيزه يو پر يو تابع خو لا پرېږده؛
  او له دې امله \LR{$\cardneq{A}{\Pow{A}}$}۔ ځکه فرض کړئ چې
  \LR{$g\colon A \to \Pow{A}$}۔ دا چې \LR{$g$} بشپړه ده، هر \LR{$x \in A$} د
  يوه فرعي سټ \LR{$g(x) \subseteq A$} سره نقشېږي۔ موږ ښودلے شو چې
  \LR{$g$} پر هدف سټ بشپړه نۀ شي کېدے۔ د دې دپاره يو فرعي سټ
  ~\LR{$\overline{A} \subseteq A$} تعريفوو چې د تعريف له مخې د~\LR{$g$}
  په قيمتونو سټ کښې نۀ شي راتلے۔ کېږدئ:
  \[
  \overline{A} = \Setabs{x \in A}{x \notin g(x)}.
  \]
  ځکه چې \LR{$g(x)$} د هر \LR{$x \in A$} دپاره تعريف شوے، \LR{$\overline{A}$}
  په څرګند ډول د~\LR{$A$} يو سم تعريف شوے فرعي سټ دے۔ خو دا د~\LR{$g$}
  د قيمتونو په سټ کښې نۀ شي راتلے۔ يو خپل سری \LR{$x \in A$} ونيسئ؛
  موږ به وښيو چې \LR{$\overline{A} \neq g(x)$}۔ که \LR{$x \in g(x)$} وي،
  نو \LR{$x \notin g(x)$} نۀ پوره کوي، او د~\LR{$\overline{A}$} د تعريف له
  مخې \LR{$x \notin \overline{A}$} لرو۔ که \LR{$x \in \overline{A}$} وي،
  د~\LR{$\overline{A}$} تعريفي خاصيت بايد پوره کړي؛ يعنې \LR{$x \in A$}
  او \LR{$x \notin g(x)$} وي۔ دا چې \LR{$x$} خپل سری و، دا ښيي چې د هر
  \LR{$x \in A$} دپاره، \LR{$x \in g(x)$} هله او يوازې هله چې
  \LR{$x \notin \overline{A}$}، او نو \LR{$g(x) \neq \overline{A}$}۔
  په بله وينا، \LR{$\overline{A}$} د~\LR{$g$} په قيمتونو
  سټ کښې نۀ شي راتلے، او دا له دې فرض سره تناقض لري چې~\LR{$g$} پر
  هدف سټ بشپړه ده۔{}
\end{proof}
\olpseganchor{OLP-0036-B014}{end}

\olpseganchor{OLP-0036-B015}{start}
\begin{explain}
د
  \olref{thm:cantor} ثبوت د
  \olref[nen]{thm:nonenum-pownat} له ثبوت سره پرتله کول ګټور دي۔
  هلته مو وښودل چې د~\LR{$\PosInt$} د فرعي سټونو د هر لړ \LR{$Z_1$}،
  \LR{$Z_2$}، \dots دپاره د عددونو يو سټ~\LR{$\overline{Z}$} جوړولے شو
  چې په ډاډه توګه په لړ کښې نۀ وي۔ په لړ کښې نۀ و، ځکه چې د هر
  \LR{$n \in \PosInt$} دپاره، \LR{$n \in Z_n$} هله او يوازې هله چې
  \LR{$n \notin \overline{Z}$}۔ په دې توګه تل کوم عدد د \LR{$Z_n$} او
  \LR{$\overline{Z}$} له يوه څخه غړے وي خو د بل نۀ وي۔ دلته
  هماغه مفکوره تعقيبوو، خو شاخصونه~\LR{$n$} اوس د~\LR{$\PosInt$} پر ځاے
  د~\LR{$A$} غړي دي۔ سټ \LR{$\overline{A}$} داسې تعريف شوے چې د
  هر \LR{$x \in A$} دپاره له~\LR{$g(x)$} څخه بېل وي، ځکه \LR{$x \in g(x)$}
  هله او يوازې هله چې \LR{$x \notin \overline{A}$}۔ بيا هم تل د~\LR{$A$}
  يو غړے شته چې له \LR{$g(x)$} او \LR{$\overline{A}$} څخه د يوه
  غړے وي خو د بل نۀ وي۔ او لکه څنګه چې له دې امله
  \LR{$\overline{Z}$} په لړ \LR{$Z_1$}، \LR{$Z_2$}، \dots کښې نۀ شي راتلے،
  \LR{$\overline{A}$} د~\LR{$g$} په قيمتونو سټ کښې نۀ شي راتلے۔
\olpseganchor{OLP-0036-B015}{end}

\olpseganchor{OLP-0036-B016}{start}
د
  \olref{thm:cantor} ثبوت د
  \olref[nen-alt]{thm:nonenum-pownat} له ثبوت سره پرتله کول ګټور
  دي۔ هلته مو وښودل چې د~\LR{$\Nat$} د فرعي سټونو د هر لړ \LR{$N_0$}،
  \LR{$N_1$}، \LR{$N_2$}، \dots دپاره د عددونو يو سټ~\LR{$D$} جوړولے شو چې په
  ډاډه توګه په لړ کښې نۀ وي۔ په لړ کښې نۀ و، ځکه چې د هر
  \LR{$n \in \Nat$} دپاره \LR{$n \in N_n$} هله او يوازې هله چې
  \LR{$n \notin D$}۔ دلته هماغه مفکوره تعقيبوو، خو شاخصونه~\LR{$n$} اوس
  د~\LR{$\Nat$} پر ځاے د~\LR{$A$} غړي دي۔ سټ \LR{$D$} داسې تعريف
  شوے چې د هر \LR{$x \in A$} دپاره له~\LR{$g(x)$} څخه بېل وي، ځکه
  \LR{$x \in g(x)$} هله او يوازې هله چې \LR{$x \notin D$}۔
\olpseganchor{OLP-0036-B016}{end}

\olpseganchor{OLP-0036-B017}{start}
دا ثبوت د رسل د پاراډوکس له ثبوت،
\olref[sfr][set][rus]{thm:russells-paradox}، سره هم د پرتله کولو
وړ دے۔ په حقيقت کښې، د کانتور قضيه د رسل د خپل پاراډوکس الهام و۔
\end{explain}
\olpseganchor{OLP-0036-B017}{end}

\olpseganchor{OLP-0036-B018}{start}
\begin{prob}
  وښيئ چې د هر سټ~\LR{$A$} دپاره هېڅ يو پر يو تابع
  \LR{$g\colon \Pow{A} \to A$} نۀ شي کېدے۔ لارښوونه: فرض کړئ
  \LR{$g\colon \Pow{A} \to A$} يو پر يو ده۔
  \LR{$D = \Setabs{g(B)}{B \subseteq A \text{\RL{ او }} g(B) \notin B}$}
  په پام کښې ونيسئ۔ کېږدئ \LR{$x = g(D)$}۔ د تناقض د راايستلو دپاره
  دا حقيقت وکاروئ چې \LR{$g$} يو پر يو ده۔
\end{prob}
\olpseganchor{OLP-0036-B018}{end}

\olpseganchor{OLP-0037-B004}{start}
\olfileid{sfr}{siz}{sb}
\olpseganchor{OLP-0037-B004}{end}

\olpseganchor{OLP-0037-B005}{start}
\olsection{د اندازې مفکوره او شروډر–برنشتاين}
\olpseganchor{OLP-0037-B005}{end}

\olpseganchor{OLP-0037-B006}{start}
\begin{explain}
دا يوه شهودي مفکوره ده: که \LR{$A$} له \LR{$B$} څخه لوے نۀ وي او \LR{$B$} له
\LR{$A$} څخه لوے نۀ وي، نو \LR{$A$} او \LR{$B$} همشمېره دي۔ رښتيا دا چې که
دا مفکوره \emph{ناسمه} وے، موږ به په ډېره ګرانه توجيه کړے شوے
چې زموږ تعريف شوې د همشمېر والي مفکوره د سټونو تر منځ د «اندازو»
له پرتله کولو سره څه اړيکه لري!{} خو له نېکه مرغه، شهودي مفکوره
سمه ده۔ دا د شروډر–برنشتاين په قضيې توجيه کېږي۔
\end{explain}
\olpseganchor{OLP-0037-B006}{end}

\olpseganchor{OLP-0037-B007}{start}
\begin{thm}[شروډر–برنشتاين]
	\ollabel{thm:schroder-bernstein}
	که \LR{$\cardle{A}{B}$} او \LR{$\cardle{B}{A}$} وي،
	نو \LR{$\cardeq{A}{B}$}۔
\end{thm}
\olpseganchor{OLP-0037-B007}{end}

\olpseganchor{OLP-0037-B008}{start}
\begin{explain}
په بله وينا، که له \LR{$A$} څخه~\LR{$B$} ته يوه يو پر يو تابع، او له
\LR{$B$} څخه~\LR{$A$} ته يوه يو پر يو تابع وي، نو له \LR{$A$} څخه~\LR{$B$} ته
يوه دوه اړخيزه يو پر يو تابع شته۔
\olpseganchor{OLP-0037-B008}{end}

\olpseganchor{OLP-0037-B009}{start}
خو د دې پايلې ثبوت په رښتيا خورا \emph{ګران} دے۔ په حقيقت کښې،
که څه هم کانتور دا پايله بيان کړه، نورو ثابته کړه۔\footnote{د
تاريخ په اړه د نورو معلوماتو دپاره، د بېلګې په ډول
\mbox{\LR{\latinfont \citet[pp.~165--6]{Potter2004}}} وګورئ۔}
%
اوس مهال بايد دا په باور ومنئ۔
\olpseganchor{OLP-0037-B009}{end}

\olpseganchor{OLP-0037-B010}{start}
له نېکه مرغه، شروډر–برنشتاين \emph{سمه} ده او دا فکر مو تاييدوي
چې زموږ تعريف شوې اړيکې، يعنې \LR{$\cardeq{A}{B}$} او
\LR{$\cardle{A}{B}$}، له «اندازې» سره څه اړيکه لري۔ پر دې سربېره، د
شروډر–برنشتاين قضيه ډېره \emph{ګټوره} ده۔ د دوو همشمېره سټونو
تر منځ د دوه اړخيزه يو پر يو تابع تصورول ګرانېدلای شي۔ د شروډر–برنشتاين
قضيه اجازه راکوي چې پرتله په حالتونو ووېشو، څو يوازې له لومړي
څخه دويم ته د يو پر يو تابع، او بيا برعکس، تصور وکړو۔
\end{explain}
\olpseganchor{OLP-0037-B010}{end}

\olpseganchor{OLP-0038-B004}{start}
\olfileid{sfr}{siz}{enm-alt}
\olpseganchor{OLP-0038-B004}{end}

\olpseganchor{OLP-0038-B005}{start}
\olsection{شمېرنې او د شمېر وړ سټونه}
\olpseganchor{OLP-0038-B005}{end}

\olpseganchor{OLP-0038-B006}{start}
\begin{editorial}
  دا برخه شمېرنې له \LR{$\Nat$} (پيلنيو ټوټو) سره د دوه اړخيزو يو
  پر يو تابعو په توګه تعريفوي، لکه په سټ‌تيورۍ کښې چې کېږي۔ نو
  د \olref[enm]{sec} له تعريفونو سره لږ ټکر لري، او هلته ټولې
  بېلګې تکراروي۔ دا د هغې برخې په پرتله لږه لنډه هم ده۔
\end{editorial}
\olpseganchor{OLP-0038-B006}{end}

\olpseganchor{OLP-0038-B007}{start}
يو متناهي سټ يوازې د هغه د غړي په شمېرلو ټاکلے شو۔ دا
هغه وخت کوو چې يو سټ داسې تعريفوو:
\[
  A = \{a_1, a_2, \ldots, a_n\}.
\]
که فرض کړو چې غړي \LR{$a_1$}، \dots، \LR{$a_n$} ټول سره بېل دي،
دا د \LR{$A$} او لومړنيو \LR{$n$} طبيعي عددونو \LR{$0$}، \dots، \LR{$n-1$} تر منځ
يوه دوه اړخيزه يو پر يو تابع راکوي۔ برعکس، ځکه هر متناهي سټ يوازې متناهي
شمېر غړي لري، هر متناهي سټ په داسې سمون کښې اېښودلے
شو۔ په بله وينا، که \LR{$A$} متناهي وي، د \LR{$A$} او
\LR{$\{0, \dots, n-1\}$} تر منځ يوه دوه اړخيزه يو پر يو تابع شته، چې \LR{$n$} د~\LR{$A$}
د غړي شمېر دے۔
\olpseganchor{OLP-0038-B007}{end}

\olpseganchor{OLP-0038-B008}{start}
که د نامتناهي سټونو ځينې ډولونه ومنو، نو ځينې نامتناهي سټونه به
هم شمېرل کېدے شي۔ دا په دې وينا دقيقولے شو چې يو نامتناهي سټ د
هغه او ټول~\LR{$\Nat$} تر منځ د دوه اړخيزه يو پر يو تابع په وسيله شمېرل کېږي۔
\olpseganchor{OLP-0038-B008}{end}

\olpseganchor{OLP-0038-B009}{start}
\begin{defn}[شمېرنه، سټ‌تيوريکي]
د يوه سټ \LR{$A$} \emph{شمېرنه} يوه دوه اړخيزه يو پر يو تابع ده چې د قيمتونو
سټ ئې \LR{$A$} او د تعريف ساحه ئې يا د طبيعي عددونو يو پيلنی سټ
\LR{$\{0, 1, \ldots, n\}$} {يا} د طبيعي عددونو ټول سټ~\LR{$\Nat$} وي۔
\end{defn}
\olpseganchor{OLP-0038-B009}{end}

\olpseganchor{OLP-0038-B010}{start}
\begin{explain}
د \emph{شمېرنې} د ويي د دې کارونې تر شا يو شهودي بنسټ شته۔ دا
ويل چې موږ يو سټ \LR{$A$} شمېرلے دے، په دې معنا دي چې يوه
دوه اړخيزه يو پر يو تابع \LR{$f$} شته چې پرې د سټ \LR{$A$} غړي شمېرلے شو۔ \LR{$0$}م غړے
\LR{$f(0)$} دے، لومړے \LR{$f(1)$}، \ldots، \LR{$n$}م \LR{$f(n)$}،
\ldots۔\footnote{هو، موږ له \LR{$0$} څخه شمېرو۔ البته له~\LR{$1$} څخه هم
پيلولے شو۔ په دې سره به لوے توپير رانۀ شي۔ يوازې به مو~\LR{$\Nat$}
په~\LR{$\PosInt$} بدلولے۔} د دې دليل د لاندې خبرې په ورزياتولو لا
روښانه کېدے شي:
\end{explain}
\olpseganchor{OLP-0038-B010}{end}

\olpseganchor{OLP-0038-B011}{start}
\begin{defn}
  \ollabel{defn:enumerable}
  يو سټ~\LR{$A$} هله او يوازې هله د شمېر وړ دے چې يا
  \LR{$A = \emptyset$} وي يا د~\LR{$A$} يوه شمېرنه وي۔ وايو چې \LR{$A$} هله
  او يوازې هله د شمېر نۀ وړ دے چې \LR{$A$} د شمېر وړ نۀ وي۔
\end{defn}
\olpseganchor{OLP-0038-B011}{end}

\olpseganchor{OLP-0038-B012}{start}
\begin{explain}
نو يو سټ هله او يوازې هله د شمېر وړ دے چې تش وي يا ئې
غړي په يوه شمېرنه شمېرلے شئ۔
\end{explain}
\olpseganchor{OLP-0038-B012}{end}

\olpseganchor{OLP-0038-B013}{start}
\begin{ex}
هغه تابع چې طبيعي عددونه شمېري يوازې عيني تابع
\LR{$\Id{\Nat} \colon \Nat \to \Nat$} ده چې په
\LR{$\Id{\Nat}(n) = n$} ورکړل شوې۔ هغه تابع چې \emph{مثبت} طبيعي
عددونه شمېري، \LR{$\Nat^+ = \Nat \setminus \{0\}$}، تابع
\LR{$g(n) = n + 1$} ده؛ يعنې راتلونکې تابع۔
\end{ex}
\olpseganchor{OLP-0038-B013}{end}

\olpseganchor{OLP-0038-B014}{start}
\begin{prob}
وښيئ چې يو سټ \LR{$A$} هله او يوازې هله د شمېر وړ دے چې يا
\LR{$A = \emptyset$} وي يا يوه پر هدف سټ بشپړه تابع \LR{$f\colon \Nat \to A$}
وي۔ وښيئ چې \LR{$A$} هله او يوازې هله د شمېر وړ دے چې يوه
يو پر يو تابع \LR{$g\colon A \to \Nat$} وي۔
\end{prob}
\olpseganchor{OLP-0038-B014}{end}

\olpseganchor{OLP-0038-B015}{start}
\begin{ex}
تابعې \LR{$f\colon \Nat \to \Nat$} او \LR{$g \colon \Nat \to \Nat$} چې
په ترتيب سره په لاندې ډول ورکړل شوې دي:
\begin{align*}
  f(n) & = 2n \text{\RL{ او}}\\
  g(n) & = 2n+1
\end{align*}
جفت طبيعي عددونه او طاق طبيعي عددونه شمېري۔ خو له دواړو يوه هم
پر هدف سټ بشپړه نۀ ده، نو يوه هم د \LR{$\Nat$} شمېرنه نۀ ده۔
\end{ex}
\olpseganchor{OLP-0038-B015}{end}

\olpseganchor{OLP-0038-B016}{start}
\begin{prob}
د مربع عددونو \LR{$1$}، \LR{$4$}، \LR{$9$}، \LR{$16$}، \dots يوه شمېرنه تعريف کړئ۔
\end{prob}
\olpseganchor{OLP-0038-B016}{end}

\olpseganchor{OLP-0038-B017}{start}
\begin{ex}
فرض کړئ \LR{$\lceil x \rceil$} د \emph{پاسني صحيح عدد} تابع ده، چې
\LR{$x$} تر ټولو نږدې صحيح عدد ته پورته ګردوي۔ بيا تابع
\LR{$f \colon \Nat \to \Int$} چې په لاندې ډول ورکړل شوې ده:
\[
  f(n) = (-1)^{n} \left\lceil\tfrac{n}{2}\right\rceil
\]
د صحيح عددونو سټ~\LR{$\Int$} په لاندې ډول شمېري:
\[
\begin{array}{c c c c c c c c}
f(0) & f(1) & f(2) & f(3) & f(4) & f(5) & f(6) & \dots \\ \\
\big\lceil \tfrac{0}{2} \big\rceil & -\big\lceil \tfrac{1}{2}\big\rceil &  \big\lceil \tfrac{2}{2} \big\rceil & -\big\lceil \tfrac{3}{2} \big\rceil & \big\lceil \tfrac{4}{2} \big\rceil  & -\big\lceil \tfrac{5}{2}\big\rceil & \big\lceil \tfrac{6}{2} \big\rceil & \dots \\ \\
0 & -1 & 1 & -2 & 2 & -3 & 3& \dots
\end{array}
\]
پام وکړئ چې \LR{$f$} څنګه د مثبتو او منفي صحيح عددونو تر منځ په
«ټوپونو» د~\LR{$\Int$} قيمتونه پيدا کوي۔ \LR{$f$} په حالتونو هم داسې
تعريفولے شئ:
\[
f(n) = \begin{cases}
  \frac{n}{2} & \text{\RL{که \LR{$n$} جفت وي}}\\
  -\frac{n+1}{2} & \text{\RL{که \LR{$n$} طاق وي}}
  \end{cases}
\]
\end{ex}
\olpseganchor{OLP-0038-B017}{end}

\olpseganchor{OLP-0038-B018}{start}
\begin{prob}
وښيئ چې که \LR{$A$} او \LR{$B$} د شمېر وړ وي، نو \LR{$A \cup B$} هم دے۔
\end{prob}
\olpseganchor{OLP-0038-B018}{end}

\olpseganchor{OLP-0038-B019}{start}
\begin{prob}
پر \LR{$n$} په استقرا وښيئ چې که \LR{$A_1$}، \LR{$A_2$}، \dots، \LR{$A_n$} ټول
د شمېر وړ وي، نو \LR{$A_1 \cup \dots \cup A_n$} هم دے۔
\end{prob}
\olpseganchor{OLP-0038-B019}{end}

\olpseganchor{OLP-0039-B004}{start}
\olfileid{sfr}{siz}{nen-alt}
\olpseganchor{OLP-0039-B004}{end}

\olpseganchor{OLP-0039-B005}{start}
\olsection{د شمېر نۀ وړ سټونه}
\olpseganchor{OLP-0039-B005}{end}

\olpseganchor{OLP-0039-B006}{start}
\begin{editorial}
  دا برخه په \olref[enm-alt]{sec} کښې د تعريفونو په کارولو، يعنې
  له \LR{$\PosInt$} څخه د پر هدف سټ بشپړې تابع پر ځاے له~\LR{$\Nat$} سره
  د دوه اړخيزې يو پر يو تابع په غوښتلو، د \LR{$\Bin^\omega$} او
  \LR{$\Pow{\Nat}$} د شمېر نۀ وړ والے ثابتوي۔
\end{editorial}
\olpseganchor{OLP-0039-B006}{end}

\olpseganchor{OLP-0039-B007}{start}
\begin{explain}
د طبيعي عددونو سټ \LR{$\Nat$} نامتناهي دے۔ دا په څرګند ډول
د شمېر وړ هم دے۔ خو د پام وړ حقيقت دا دے چې
\emph{د شمېر نۀ وړ} سټونه، يعنې هغه سټونه چې د شمېر وړ
نۀ دي، شته (\olref[sfr][siz][enm-alt]{defn:enumerable} وګورئ)۔
\olpseganchor{OLP-0039-B007}{end}

\olpseganchor{OLP-0039-B008}{start}
دا ښايي حيرانوونکې وي۔ بالاخره، دا ويل چې \LR{$A$} د شمېر نۀ وړ
دے، دا ويل دي چې \emph{هېڅ} دوه اړخيزه يو پر يو تابع
\LR{$f \colon \Nat \to A$} نۀ شته؛ يعنې هېڅ هغه تابع چې د~\LR{$\Nat$}
بې‌شمېره غړي پر~\LR{$A$} نقشوي، ټول~\LR{$A$} نۀ شي پوره کولے۔
نو که \LR{$A$} د شمېر نۀ وړ وي، د طبيعي عددونو په پرتله د~\LR{$A$}
«زيات» غړي شته۔
\olpseganchor{OLP-0039-B008}{end}

\olpseganchor{OLP-0039-B009}{start}
د دې د ثابتولو دپاره چې يو سټ د شمېر نۀ وړ دے، بايد وښيئ چې هېڅ
مناسبه دوه اړخيزه يو پر يو تابع نۀ شي موجودېدلای۔ د دې تر ټولو ښه لار دا ده
چې وښيئ د~\LR{$A$} د غړي د شمېرلو هره هڅه بايد لږ تر لږه يو
غړے پرېږدي؛ دا ښيي چې هېڅ تابع \LR{$f\colon \Nat \to A$}
پر هدف سټ بشپړه نۀ ده۔ د دې د ثابتولو يوه عمومي تګلاره د کانتور
\emph{قطري طريقه} ده۔ د~\LR{$A$} د غړي يو لړ، د بېلګې په
ډول \LR{$x_1$}، \LR{$x_2$}، \dots، چې راکړل شي، موږ د~\LR{$A$} يو بل
غړے جوړوو چې د خپل جوړښت له مخې په هېڅ ډول په هغه لړ
کښې نۀ شي راتلے۔
\olpseganchor{OLP-0039-B009}{end}

\olpseganchor{OLP-0039-B010}{start}
خو دا هر څه په بېلګه تر ټولو ښه پوهېدل کېږي۔ نو زموږ لومړۍ
بېلګه د \LR{$0$}ګانو او \LR{$1$}ګانو د ټولو نامتناهي توريزو تسلسلونو سټ
~\LR{$\Bin^\omega$} دے۔ (`\LR{$\Bin$}' د دوه‌ګوني دپاره دے، او يوازې ئې د
دوه غړيز سټ \LR{$\{0,1\}$} په توګه ګڼلے شو۔) خو اوس مهال دا لږ نرم بيان بايد هېڅ ګډوډي جوړه نۀ کړي۔
\end{explain}
\olpseganchor{OLP-0039-B010}{end}

\olpseganchor{OLP-0039-B011}{start}
\begin{thm}
\ollabel{thm:nonenum-bin-omega}
\LR{$\Bin^\omega$}~د شمېر نۀ وړ دے۔
\end{thm}
\olpseganchor{OLP-0039-B011}{end}

\olpseganchor{OLP-0039-B012}{start}
\begin{proof}
د \LR{$\Bin^\omega$} د يوه فرعي سټ هره شمېرنه په پام کښې ونيسئ۔ نو
کوم لړ \LR{$s_{0}$}، \LR{$s_{1}$}، \LR{$s_{2}$}، \dots{} لرو چې هر \LR{$s_n$} د
\LR{$0$}ګانو او~\LR{$1$}ګانو يو نامتناهي توريز تسلسل دے۔ فرض کړئ
\LR{$s_n(m)$} په دې لړ کښې د \LR{$n$}م توريز تسلسل \LR{$m$}م رقم دے۔
نو اوس خپل لړ د يوه جدولي
ترتيب په توګه ګڼلے شو، چې \LR{$s_n(m)$} په \LR{$n$}م قطار او \LR{$m$}م ستون
کښې اېښوول شوے دے:
\[
\begin{array}{c|c|c|c|c|c}
& 0 & 1 & 2 & 3 & \dots \\\hline
0 & \mathbf{s_{0}(0)} & s_{0}(1) & s_{0}(2) & s_0(3) & \dots \\\hline
1 & s_{1}(0)& \mathbf{s_{1}(1)} & s_1(2) & s_1(3) & \dots \\\hline
2 & s_{2}(0)& s_{2}(1) & \mathbf{s_2(2)} & s_2(3) & \dots \\\hline
3 & s_{3}(0)& s_{3}(1) & s_3(2) & \mathbf{s_3(3)} & \dots \\\hline
\vdots & \vdots & \vdots & \vdots & \vdots & \mathbf{\ddots}
\end{array}
\]
اوس د \LR{$0$}ګانو او \LR{$1$}ګانو يو داسې نامتناهي توريز تسلسل \LR{$d$}
جوړوو چې په دې لړ کښې نۀ وي۔ دا د هغه د هرې خانې په ټاکلو؛
يعنې د ټولو \LR{$n \in \Nat$} دپاره د \LR{$d(n)$} په ټاکلو، کوو۔ په شهودي
ډول، د پورته جدول پر قطر لاندې ځو (ځکه ورته «قطري طريقه» وايو)،
او بيا هر \LR{$1$} په \LR{$0$} او هر \LR{$0$} په~\LR{$1$} بدلوو۔
په لا مجرد ډول، \LR{$d(n)$} د دې له
مخې \LR{$0$} يا \LR{$1$} تعريفوو چې د قطر \LR{$n$}م غړے، \LR{$s_n(n)$}،
\LR{$1$} دے که \LR{$0$}؛ يعنې:
\[
d(n) =
\begin{cases}
1 & \text{\RL{که \LR{$s_{n}(n) = 0$} وي}}\\
0 & \text{\RL{که \LR{$s_{n}(n) = 1$} وي}}
\end{cases}
\]
په څرګند ډول \LR{$d \in \Bin^\omega$}، ځکه دا د \LR{$0$}ګانو او \LR{$1$}ګانو
يو نامتناهي توريز تسلسل دے۔ خو موږ \LR{$d$} داسې جوړ کړے چې د هر
\LR{$n \in \Nat$} دپاره \LR{$d(n) \neq s_n(n)$}۔ يعنې \LR{$d$} له \LR{$s_n$} څخه
په خپل \LR{$n$}م مقام کښې توپير لري۔ نو د هر \LR{$n\in \Nat$} دپاره
\LR{$d \neq s_n$}۔ ځکه \LR{$d$} په لړ \LR{$s_0$}، \LR{$s_1$}، \LR{$s_2$}،
\dots کښې نۀ شي راتلے۔
\olpseganchor{OLP-0039-B012}{end}

\olpseganchor{OLP-0039-B013}{start}
موږ وښودل چې د \LR{$\Bin^\omega$} د کوم فرعي سټ هره خپله سرې شمېرنه
به د \LR{$\Bin^\omega$} کوم غړے پرېږدي۔ نو د سټ
\LR{$\Bin^\omega$} هېڅ شمېرنه نۀ شته؛ يعنې \LR{$\Bin^\omega$}
د شمېر نۀ وړ دے۔
\end{proof}
\olpseganchor{OLP-0039-B013}{end}

\olpseganchor{OLP-0039-B014}{start}
\begin{explain}
د ثبوت دې طريقې ته «قطرول» ويل کېږي، ځکه د~\LR{$d$} د تعريف دپاره
د جدول قطر کاروي۔ خو قطرول د جدول شتون ته اړتيا نۀ لري۔ په
حقيقت کښې، په ورته مفکوره ښودلے شو چې ځينې سټونه
د شمېر نۀ وړ دي، ان چې هېڅ جدول او رښتينے قطر نۀ وي۔
لاندې پايله ښيي چې څنګه۔
\end{explain}
\olpseganchor{OLP-0039-B014}{end}

\olpseganchor{OLP-0039-B015}{start}
\begin{thm}
\ollabel{thm:nonenum-pownat}
\LR{$\Pow{\Nat}$} د شمېر وړ نۀ دے۔
\end{thm}
\olpseganchor{OLP-0039-B015}{end}

\olpseganchor{OLP-0039-B016}{start}
\begin{proof}
په هماغه طريقه مخکښې ځو: ښيو چې د~\LR{$\Nat$} د فرعي سټونو هر لړ
د \LR{$\Nat$} کوم فرعي سټ پرېږدي۔ نو فرض کړئ چې د \LR{$\Nat$} د فرعي
سټونو کوم لړ \LR{$N_0, N_1, N_2, \ldots$} لرو۔ سټ \LR{$D$} داسې تعريفوو:
\LR{$n \in D$} هله او يوازې هله چې \LR{$n \notin N_{n}$}:
\[
D = \Setabs{n \in \Nat}{n \notin N_n}
\]
په څرګند ډول \LR{$D\subseteq \Nat$}۔ خو \LR{$D$} په لړ کښې نۀ شي راتلے۔
بالاخره، د جوړښت له مخې \LR{$n \in D$} هله او يوازې هله چې
\LR{$n\notin N_n$}؛ نو د هر \LR{$n \in \Nat$} دپاره \LR{$D \neq N_n$}۔
\end{proof}
\olpseganchor{OLP-0039-B016}{end}

\olpseganchor{OLP-0039-B017}{start}
\begin{explain}
تېر ثبوت د قطر نوم وانۀ خيست۔ بيا هم که داسې ئې انځور کړئ، قطري
طريقه په کښې ليدلے شئ: وانګېړئ چې سټونه \LR{$N_0$}، \LR{$N_1$}، \dots په
جدولي ترتيب کښې ليکل شوي، چې \LR{$N_n$} په \LR{$n$}م قطار کښې داسې ليکو
چې \LR{$m$} په \LR{$m$}م ستون کښې هله او يوازې هله ليکو که \LR{$m \in N_n$}۔
د بېلګې په ډول، فرض کړئ په هغه لړ کښې لومړي څلور سټونه
\LR{$\{0,1,2,\dots\}$}، \LR{$\{1, 3, 5, \dots\}$}، \LR{$\{0,1,4\}$} او
\LR{$\{2,3,4,\dots\}$} دي؛ نو زموږ جدولي ترتيب به داسې پيل شي:
\[
\begin{array}{r@{}rrrrrrr}
  N_0 = \{ & \mathbf{0}, & 1, & 2, & & & & \dots\}\\
  N_1 = \{ &  & \mathbf{1}, &  & 3, &  & 5, & \dots\}\\
  N_2 = \{ & 0, & 1, &  &  & 4\phantom{,} &  & \}\\
  N_3 = \{ &  &  & 2, & \mathbf{3}, & 4, & & \dots\}\\
  &\vdots & & & & & & \ddots\phantom{\}}
  \end{array}
\]
بيا \LR{$D$} هغه سټ دے چې پر قطر د لاندې تلو له لارې تر لاسه کېږي:
\LR{$n \in D$} هله او يوازې هله اېږدو چې \LR{$n$} پر قطر \emph{نۀ} وي۔
نو په پورته حالت کښې \LR{$0$} او \LR{$1$} پرېږدو،~\LR{$2$} شاملوو،~\LR{$3$}
پرېږدو، او همداسې نور۔
\end{explain}
\olpseganchor{OLP-0039-B017}{end}

\olpseganchor{OLP-0039-B018}{start}
\begin{prob}
په څرګند قطري استدلال وښيئ چې د ټولو تابعو
\LR{$f \colon \Nat \to \Nat$} سټ د شمېر نۀ وړ دے۔ يعنې وښيئ
چې که \LR{$f_1$}، \LR{$f_2$}، \dots د تابعو يو لړ وي او هره
\LR{$f_i\colon \Nat \to \Nat$} وي، نو کومه \LR{$g \colon \Nat \to
\Nat$} شته چې په دې لړ کښې نۀ ده۔
\end{prob}
\olpseganchor{OLP-0039-B018}{end}

\olpseganchor{OLP-0040-B004}{start}
\olfileid{sfr}{siz}{red-alt}
\olpseganchor{OLP-0040-B004}{end}

\olpseganchor{OLP-0040-B005}{start}
\olsection{راکمونه}
\olpseganchor{OLP-0040-B005}{end}

\olpseganchor{OLP-0040-B006}{start}
\begin{editorial}
  دا برخه په راکمونه د شمېر نۀ وړ والے ثابتوي، او پايلې ئې د
  \olref[nen-alt]{sec} له پايلو سره سمون خوري۔ يوه بله، لږه
  تفصيلي بڼه چې د \olref[nen]{sec} له پايلو سره سمون خوري، په
  \olref[red]{sec} کښې ورکړل شوې ده۔
\end{editorial}
\olpseganchor{OLP-0040-B006}{end}

\olpseganchor{OLP-0040-B007}{start}
موږ په يوه قطري استدلال وښودل چې \LR{$\Bin^\omega$} د شمېر نۀ وړ
دے۔ په ورته قطري استدلال مو وښودل چې \LR{$\Pow{\Nat}$}
د شمېر نۀ وړ دے۔ خو د \LR{$\Pow{\Nat}$} د د شمېر نۀ وړ
کېدو يوه بله لار دا ده: وښيئ چې \emph{که \LR{$\Pow{\Nat}$}
د شمېر وړ وي، نو \LR{$\Bin^\omega$} هم د شمېر وړ دے}۔ ځکه
چې پوهېږو \LR{$\Bin^\omega$} د شمېر نۀ وړ دے، دا به وښيي چې
\LR{$\Pow{\Nat}$} هم همداسې دے۔
\olpseganchor{OLP-0040-B007}{end}

\olpseganchor{OLP-0040-B008}{start}
دې ته د يوې مسئلې بلې ته \emph{راکمونه} ويل کېږي۔ په دې حالت
کښې د \LR{$\Bin^\omega$} د شمېرلو مسئله د \LR{$\Pow{\Nat}$} د شمېرلو
مسئلې ته راکموو۔ د وروستۍ مسئلې يو حل، يعنې د \LR{$\Pow{\Nat}$} يوه
شمېرنه، به د لومړۍ مسئلې حل، يعنې د \LR{$\Bin^\omega$} يوه شمېرنه،
راکړي۔
\olpseganchor{OLP-0040-B008}{end}

\olpseganchor{OLP-0040-B009}{start}
د يوه سټ~\LR{$B$} د شمېرلو مسئله د يوه سټ~\LR{$A$} د شمېرلو مسئلې ته د
راکمولو دپاره، د~\LR{$A$} د يوې شمېرنې د~\LR{$B$} په شمېرنه د اړولو يوه
لار ورکوو۔ د دې تر ټولو اسانه لار د يوې پر هدف سټ بشپړه تابع تابع
\LR{$f\colon A \to B$} تعريفول دي۔ که \LR{$x_1$}، \LR{$x_2$}، \dots{} د~\LR{$A$}
شمېرنه وي، نو \LR{$f(x_1)$}، \LR{$f(x_2)$}، \dots{} به د~\LR{$B$} شمېرنه وي۔
زموږ په حالت کښې د يوې پر هدف سټ بشپړه تابع تابع
\LR{$f\colon \Pow{\Nat} \to \Bin^\omega$} په لټه کښې يو۔
\olpseganchor{OLP-0040-B009}{end}

\olpseganchor{OLP-0040-B010}{start}
\begin{prob}
وښيئ چې که يوه يو پر يو تابع \LR{$g\colon B \to A$} وي او
\LR{$B$}~د شمېر نۀ وړ وي، نو \LR{$A$} هم همداسې دے۔ دا په ښودلو
وکړئ چې څنګه د~\LR{$g$} په کارولو د~\LR{$A$} يوه شمېرنه د~\LR{$B$} په شمېرنه
اړولے شئ۔
\end{prob}
\olpseganchor{OLP-0040-B010}{end}

\olpseganchor{OLP-0040-B011}{start}
\begin{proof}[په راکمونه د {\olref[nen-alt]{thm:nonenum-pownat}} ثبوت]
فرض کړئ چې \LR{$\Pow{\Nat}$} د شمېر وړ وے، او له دې امله د هغه
يوه شمېرنه \LR{$N_{1}$}، \LR{$N_{2}$}، \LR{$N_{3}$}، \dots وے۔
\olpseganchor{OLP-0040-B011}{end}

\olpseganchor{OLP-0040-B012}{start}
تابع \LR{$f \colon \Pow{\Nat} \to \Bin^\omega$} داسې تعريف کړئ چې
\LR{$f(N)$} هغه تسلسل \LR{$s$} وي چې \LR{$s(n) = 1$} هله او يوازې هله چې
\LR{$n \in N$} وي، او په بل حالت کښې \LR{$s(n) = 0$} وي۔
\olpseganchor{OLP-0040-B012}{end}

\olpseganchor{OLP-0040-B013}{start}
دا په څرګند ډول تابع تعريفوي، ځکه کله چې \LR{$N \subseteq \Nat$} وي،
هر \LR{$n \in \Nat$} يا د~\LR{$N$} يو غړے وي يا نۀ وي۔ د بېلګې
په ډول، د طبيعي جفتو عددونو سټ
\LR{$2\Nat = \Setabs{2n}{n \in \Nat} = \{0,2, 4, 6, \dots\}$}
تسلسل \LR{$1010101\dots$} ته نقشېږي، \LR{$\emptyset$} سټ \LR{$0000\dots$} ته او پخپله
سټ \LR{$\Nat$}، \LR{$1111\dots$} ته نقشېږي۔
\olpseganchor{OLP-0040-B013}{end}

\olpseganchor{OLP-0040-B014}{start}
دا پر هدف سټ بشپړه هم ده: د \LR{$0$}ګانو او \LR{$1$}ګانو هر تسلسل د طبيعي
عددونو له يوه سټ سره سمون خوري، يعنې له هغه سټ سره چې غړي ئې
هغه طبيعي عددونه دي کوم چې د تسلسل د~\LR{$1$} لرونکو مقامونو سره
سمون خوري۔ په لا دقيقه توګه، که \LR{$s \in \Bin^\omega$} وي، نو
\LR{$N \subseteq \Nat$} داسې تعريف کړئ:
\[
N = \Setabs{n \in \Nat}{s(n) = 1}
\]
بيا \LR{$f(N) = s$}، لکه چې د~\LR{$f$} تعريف ته په کتلو تاييدولے شو۔
\olpseganchor{OLP-0040-B014}{end}

\olpseganchor{OLP-0040-B015}{start}
اوس دا لړ په پام کښې ونيسئ:
\[
f(N_1), f(N_2), f(N_3), \dots
\]
ځکه چې \LR{$f$} پر هدف سټ بشپړه ده، د \LR{$\Bin^\omega$} هر غړے بايد د کوم
ورودي دپاره د~\LR{$f$} د قيمت په توګه راشي، او نو بايد په لړ کښې راشي۔
له دې امله دا لړ بايد ټول~\LR{$\Bin^\omega$} وشمېري۔
\olpseganchor{OLP-0040-B015}{end}

\olpseganchor{OLP-0040-B016}{start}
نو که \LR{$\Pow{\Nat}$} د شمېر وړ وے، \LR{$\Bin^\omega$} به
د شمېر وړ وے۔ خو \LR{$\Bin^\omega$} د شمېر نۀ وړ دے
(\olref[nen-alt]{thm:nonenum-bin-omega})۔ له دې امله
\LR{$\Pow{\Nat}$} د شمېر نۀ وړ دے۔
\end{proof}
\olpseganchor{OLP-0040-B016}{end}

\olpseganchor{OLP-0040-B017}{start}
%\begin{explain}
%د راکمونې په لوري کښې ګډوډېدل اسانه دي۔ د بېلګې په ډول، يوه
%پر هدف سټ بشپړه تابع $g \colon \Bin^\omega \to X$ دا \emph{نۀ}
%ثابتوي چې $X$ د شمېر نۀ وړ دے۔ (تابع $g
%\colon \Bin^\omega \to \Bin$ په پام کښې ونيسئ چې په
%$g(s) = s(1)$ تعريف شوې؛ دا تابع د $0$ګانو او $1$ګانو يو تسلسل
%د هغه لومړي غړے ته نقشوي۔ دا پر هدف سټ بشپړه ده، ځکه ځينې
%تسلسلونه په $0$ او ځينې په $1$ پيلېږي۔ خو $\Bin$ متناهي دے۔)
%دا هم په ياد ولرئ چې تابع~$f$ بايد پر هدف سټ بشپړه وي، ګنې
%استدلال مخکښې نۀ ځي: بيا به دا ډاډ نۀ وي چې $f(x_1)$، $f(x_2)$،
%\dots{} د~$Y$ ټول غړي لري۔ د بېلګې په ډول، تابع
%$h\colon \Nat \to
%\Bin^\omega$ په لاندې ډول تعريف کړئ:
%\[
%h(n) = \underbrace{000\dots0}_{\text{$n$ $0$ګانې}}
%\]
%دا يوه تابع ده، خو $\Nat$ د شمېر وړ دے۔
%\end{explain}
\olpseganchor{OLP-0040-B017}{end}

\olpseganchor{OLP-0040-B018}{start}
\begin{prob}\label{sfr:siz:red-alt:prob:nat-nat}
  په راکمونې استدلال وښيئ چې د ټولو تابعو \LR{$f\colon \Nat \to \Nat$}
  سټ~\LR{$X$} د شمېر نۀ وړ دے۔ (لارښوونه: له \LR{$X$} څخه
  ~\LR{$\Bin^\omega$} ته يوه پر هدف سټ بشپړه تابع ورکړئ۔)
\end{prob}
\olpseganchor{OLP-0040-B018}{end}

\olpseganchor{OLP-0040-B019}{start}
\begin{prob}
په راکمونې استدلال وښيئ چې د طبيعي عددونو د جوړو د
\emph{سټونو} سټ، يعنې \LR{$\Pow{\Nat \times \Nat}$}،
د شمېر نۀ وړ دے۔
\end{prob}
\olpseganchor{OLP-0040-B019}{end}

\olpseganchor{OLP-0040-B020}{start}
\begin{prob}
په راکمونې استدلال وښيئ چې \LR{$\Nat^\omega$}، يعنې د طبيعي عددونو
د نامتناهي تسلسلونو سټ، د شمېر نۀ وړ دے۔
\end{prob}
\olpseganchor{OLP-0040-B020}{end}

\olpseganchor{OLP-0040-B021}{start}
%\begin{prob}
%فرض کړئ $P$ له $\Nat$ څخه سټ~$\{0\}$ ته د تابعو سټ دے، او $Q$
%له مثبتو صحيح عددونو څخه سټ~$\{0\}$ ته د \emph{جزوي} تابعو سټ
%دے۔ وښيئ چې $P$~د شمېر وړ دے او $Q$ نۀ دے۔ (لارښوونه: د
%$\Bin^\omega$ د شمېرلو مسئله د~$Q$ شمېرلو ته راکمه کړئ۔)
%\end{prob}
\olpseganchor{OLP-0040-B021}{end}

\olpseganchor{OLP-0040-B022}{start}
\begin{prob}
فرض کړئ \LR{$S$} له \LR{$\Nat$} څخه سټ~\LR{$\{0,1\}$} ته د ټولو
پر هدف سټ بشپړې تابعې سټ دے؛ يعنې \LR{$S$} له ټولو پر هدف سټ بشپړې تابعې
~\LR{$f \colon \Nat \to \Bin$} څخه جوړ دے۔ وښيئ چې \LR{$S$}
د شمېر نۀ وړ دے۔
\end{prob}
\olpseganchor{OLP-0040-B022}{end}

\olpseganchor{OLP-0040-B023}{start}
\begin{prob}
وښيئ چې د ټولو حقيقي عددونو سټ~\LR{$\Real$} د شمېر نۀ وړ دے۔
\end{prob}
\olpseganchor{OLP-0040-B023}{end}

\olpseganchor{OLP-0041-B004}{start}
\olchapter{sfr}{arith}{د شمېر نظامونو سټي جوړونه}
\olpseganchor{OLP-0041-B004}{end}

\olpseganchor{OLP-0041-B005}{start}
\begin{editorial}
په دې څپرکي کښې مواد د ساده سټ نظريې په دننه کښې د شمېر نظامونو
جوړښت وړاندې کوي۔ دا د ټم بټن له \emph{اوپن سېټ تيوري} متن څخه
اخېستل شوي دي۔
\end{editorial}
\olpseganchor{OLP-0041-B005}{end}

\olpseganchor{OLP-0042-B004}{start}
\olfileid{sfr}{arith}{int}
\olsection{له \texorpdfstring{\LR{$\Nat$}}{ℕ} څخه \texorpdfstring{\LR{$\Int$}}{ℤ} ته}
\olpseganchor{OLP-0042-B004}{end}

\olpseganchor{OLP-0042-B005}{start}
دا دوه بنسټيزې خبرې په پام کښې ونيسئ:
	\begin{enumerate}
		\item هر صحيح عدد د \LR{$n - m$} په بڼه ليکل کېدے شي، چې \LR{$n, m \in\Nat$} وي۔
		\item په عبارت \LR{$n - m$} کښې زېرمه شوي معلومات په هماغه ډول په مرتبه جوړه \LR{$\tuple{n, m}$} کښې زېرمه کېدے شي۔
	\end{enumerate}
موږ لا پوهېږو چې د طبيعي عددونو مرتبې جوړې د \LR{$\Nat^2$} غړي دي۔ او دا فرضوو چې په \LR{$\Nat$} پوهېږو۔ نو پر دغو دوو خبرو ولاړ يو ساده وړانديز دا دے: \emph{صحيح عددونه دې د طبيعي عددونو د مرتبو جوړو په توګه وګڼو}۔
\olpseganchor{OLP-0042-B005}{end}

\olpseganchor{OLP-0042-B006}{start}
په حقيقت کښې دا وړانديز ډېر ساده دے۔ ښکاره ده چې غواړو \LR{$0- 2 = 4 - 6$} وي۔ خو په څرګند ډول \LR{$\tuple{0, 2 }\neq \tuple{4, 6}$} دے۔ نو په ساده ډول نۀ شو ويلے چې \LR{$\Nat^2$} د صحيح عددونو سټ دے۔
\olpseganchor{OLP-0042-B006}{end}

\olpseganchor{OLP-0042-B007}{start}
د مخکښې ستونزې په عمومي کولو، موږ لاندې خاصيت غواړو:
	$$a - b = c - d \text{\RL{ هله او يوازې هله چې }}a + d = c + b$$
(دا بايد څرګنده وي چې صحيح عددونه په همدې ډول \emph{چلېدل} غواړي: يوازې دواړو خواوو ته \LR{$b$} او \LR{$d$} ورزيات کړئ۔) د دې چلند د يقيني کولو اسانه لار دا ده چې د مرتبو جوړو تر منځ د هم ارزښتۍ يوه اړيکه \LR{$\Intequiv$} داسې تعريف کړو:
$$\tuple{ a, b } \Intequiv \tuple{c, d}\text{\RL{ هله او يوازې هله چې }}a + d = c + b$$
اوس بايد وښيو چې دا د هم ارزښتۍ اړيکه ده۔
\begin{prop} \LR{$\Intequiv$} د هم ارزښتۍ اړيکه ده۔
\end{prop}
	\begin{proof}
	بايد وښيو چې \LR{$\Intequiv$} انعکاسي، تناظري او انتقالي ده۔
\olpseganchor{OLP-0042-B007}{end}
	
\olpseganchor{OLP-0042-B008}{start}
	\emph{انعکاسيت:} په څرګند ډول \LR{$\tuple{a, b} \Intequiv \tuple{a, b}$}، ځکه چې \LR{$a + b = b + a$}۔
\olpseganchor{OLP-0042-B008}{end}
	
\olpseganchor{OLP-0042-B009}{start}
	\emph{تناظر:} فرض کړئ \LR{$\tuple{a, b} \Intequiv \tuple{c, d}$}، نو \LR{$a + d = c + b$}۔ بيا \LR{$c + b = a + d$}، نو \LR{$\tuple{c, d} \Intequiv \tuple{a, b}$}۔
\olpseganchor{OLP-0042-B009}{end}
	
\olpseganchor{OLP-0042-B010}{start}
	\emph{انتقاليت:} فرض کړئ \LR{$\tuple{a, b} \Intequiv \tuple{c, d}\Intequiv \tuple{m, n}$}۔ نو \LR{$a + d = c + b$} او \LR{$c + n = m + d$}۔ له دې امله \LR{$a + d + c + n = c + b + m + d$}، او نو \LR{$a + n = m + b$}۔ په پايله کښې \LR{$\tuple{a, b } \Intequiv \tuple{m, n}$}۔		
\end{proof}
\olpseganchor{OLP-0042-B010}{end}

\olpseganchor{OLP-0042-B011}{start}
اوس په دې د هم ارزښتۍ اړيکه د هم ارزښتۍ ټولګي اخيستلے شو:
\olpseganchor{OLP-0042-B011}{end}

\olpseganchor{OLP-0042-B012}{start}
\begin{defn}
		صحيح عددونه د \LR{$\Intequiv$} له مخې د طبيعي عددونو د مرتبو جوړو د هم ارزښتۍ ټولګي دي؛ يعنې \LR{$\Int = \equivclass{\Nat^2}{\Intequiv}$}۔
\end{defn}
\olpseganchor{OLP-0042-B012}{end}

\olpseganchor{OLP-0042-B013}{start}
د دې ټاکلي تعريف په اړه ښايي بېلابېل \emph{فلسفي} غبرګونونه وي۔ خو د هغو غبرګونونو تر څېړلو مخکښې، ښه ده چې ځينې فني جزئيات پر مخ يوسو۔
\olpseganchor{OLP-0042-B013}{end}

\olpseganchor{OLP-0042-B014}{start}
چې صحيح عددونه مو تعريف کړل، پر هغو به بنسټيزې تابعې او اړيکې هم تعريفول غواړو۔ \LR{$\equivrep{m, n}{\Intequiv}$} دې د \LR{$\Intequiv$} له مخې هغه د هم ارزښتۍ ټولګے وي چې \LR{$\tuple{m, n}$} ئې يو غړے دے۔\footnote{يادونه: په \olref[sfr][rel][eqv]{def:equivalenceclass} کښې د معرفي شوې نښې په کارولو به مو د همدې څيز دپاره \LR{$\equivrep{\tuple{m,n}}{\Intequiv}$} ليکل۔ خو هغه لوستل لږ ګران دي۔} يعنې:
	$$\equivrep{m, n}{\Intequiv} = \Setabs{\tuple{a, b} \in \Nat^2}{\tuple{a, b}\Intequiv \tuple{m, n}}$$
نو اوس لاندې تعريفونه وړاندې کوو:
	\begin{align*}
		\equivrep{a, b}{\Intequiv} + \equivrep{c, d}{\Intequiv} &= \equivrep{a + c, b + d}{\Intequiv}\\
		\equivrep{a, b}{\Intequiv} \times \equivrep{c, d}{\Intequiv} &= \equivrep{a c + b  d, a  d + b c}{\Intequiv}\\
		\equivrep{a, b}{\Intequiv} \leq \equivrep{c, d}{\Intequiv} &\text{\RL{ هله او يوازې هله چې }}a + d \leq b + c
	\end{align*}	
(لکه چې دود دے، د اصولو د اسانه لوستلو دپاره `\LR{$ab$}' د `\LR{$(a \times b)$}' پر ځاے کاروم۔) اوس بايد ډاډمن شو چې دا تعريفونه لکه څنګه چې \emph{بايد}، هماغسې چلېږي۔ د دې مطلب په تفصيل بيانول او بيا ئې ارزول ډېر اوږد کار دے؛ جزئيات \olref[check]{sec} ته پرېږدو۔ خو لنډه خبره دا ده: هر څه سم کار کوي!
\olpseganchor{OLP-0042-B014}{end}

\olpseganchor{OLP-0042-B015}{start}
يو وروستے کار پاتې دے۔ موږ صحيح عددونه د طبيعي عددونو په
کارولو جوړ کړي دي۔ خو د دې معنا به دا وي چې طبيعي عددونه \emph{پخپله
صحيح عددونه نۀ دي}۔ د دې فلسفي ارزښت ته به په \olref[ref]{sec}
کښې ورستانۀ شو۔ خو په سوچه فني کچه به يوې داسې لارې ته اړتيا لرو
چې طبيعي عددونه د صحيح عددونو \emph{په توګه} وکارولے شو۔ مفکوره
ډېره اسانه ده: د هر \LR{$n \in \Nat$} دپاره ټاکو چې
\LR{$n_\Int = \equivrep{n, 0}{\Intequiv}$}۔ بايد تاييد کړو چې دا تعريف
سم چلند کوي؛ يعنې د هر \LR{$m, n \in \Nat$} دپاره
\begin{align*}
	(m + n)_\Int &= m_\Int + n_\Int\\
	(m \times n)_\Int &= m_\Int \times n_\Int\\
	m \leq n &\liff m_\Int \leq n_\Int
\end{align*}
خو دا ټول پوره ساده دي۔ د بېلګې په ډول، د دې د ښودلو دپاره چې
دويم خاصيت سم دے، د طبيعي عددونو له چلنده په ساده ډول کار اخلو او
لاندې استدلال کوو:
%\begin{proof}
%	د طبيعي عددونو له چلنده په کار اخيستلو، په څرګند ډول:
	\begin{align*}
%		(m + n)_\Int  = \equivrep{m + n, 0}{\Intequiv} = \equivrep{m + n, 0 + 0}{\Intequiv} = \equivrep{m, 0}{\Intequiv} + \equivrep{n, 0}{\Intequiv} = m_\Int + n_\Int\\
		(m \times n)_\Int  &= \equivrep{m \times n, 0}{\Intequiv} \\
		&= \equivrep{m \times n + 0 \times 0, m \times 0 + 0 \times n}{\Intequiv} \\
		&= \equivrep{m, 0}{\Intequiv} \times \equivrep{n, 0}{\Intequiv} \\
		&= m_\Int \times n_\Int
	\end{align*}
%\end{proof}
د نورو دوو شرطونو تاييد د تمرين په توګه پرېږدو۔
\begin{prob}
	وښيئ چې د هر \LR{$m, n \in \Nat$} دپاره \LR{$(m + n)_\Int = m_\Int + n_\Int$} او \LR{$m \leq n \liff m_\Int \leq n_\Int$}۔
\end{prob}
\olpseganchor{OLP-0042-B015}{end}

\olpseganchor{OLP-0043-B004}{start}
\olfileid{sfr}{arith}{rat}
\olsection{له \texorpdfstring{\LR{$\Int$}}{ℤ} څخه \texorpdfstring{\LR{$\Rat$}}{ℚ} ته}
\olpseganchor{OLP-0043-B004}{end}

\olpseganchor{OLP-0043-B005}{start}
همدا اوس مو وليدل چې په ساده سټ نظريه څنګه له طبيعي عددونو څخه
صحيح عددونه جوړېږي۔ اوس به په ډېرې ورته لارې له صحيح عددونو څخه
ناطق عددونه جوړول وګورو۔ زموږ لومړنۍ خبرې دا دي:
\begin{enumerate}
	\item هر ناطق عدد د \LR{$\nicefrac{i}{j}$} په بڼه ليکل کېدے شي،
	چې \LR{$i$} او \LR{$j$} دواړه صحيح عددونه وي خو \LR{$j$} صفر نۀ وي۔
	\item په عبارت \LR{$\nicefrac{i}{j}$} کښې زېرمه شوي معلومات په هماغه
	ډول په مرتبه جوړه \LR{$\tuple{ i, j}$} کښې زېرمه کېدے شي۔
\end{enumerate}
ښکاره لار به دا وي چې ناطق عددونه د \LR{$\Int \times
(\Int \setminus \{0_\Int\})$} څخه د اخيستل شوو مرتبو جوړو \emph{په
توګه} وګڼو۔ خو لکه مخکښې، دا به لږ ډېر ساده وي، ځکه موږ غواړو
\LR{$\nicefrac{3}{2} = \nicefrac{6}{4}$} وي، خو \LR{$\tuple{ 3, 2}\neq
\tuple{ 6, 4}$}۔ په عمومي ډول لاندې خاصيت غواړو:
\[
	\nicefrac{a}{b} = \nicefrac{c}{d} \text{\RL{ هله او يوازې هله چې }} a \times d = b \times c
\]
د دې د ترلاسه کولو دپاره پر \LR{$\Int \times (\Int \setminus
\{0_\Int\})$} د هم ارزښتۍ يوه اړيکه داسې تعريفوو:
\[
	\tuple{ a, b }\Ratequiv \tuple{ c, d} \text{\RL{ هله او يوازې هله چې }}a \times d = b \times c
\]
بايد وارزوو چې دا د هم ارزښتۍ اړيکه ده۔ دا د \LR{$\Intequiv$} حالت ته
ډېر ورته دے، او موږ ئې د تمرين په توګه پرېږدو۔
\begin{prob}
وښيئ چې \LR{$\Ratequiv$} د هم ارزښتۍ اړيکه ده۔
\end{prob}
خو دا موږ ته د لاندې خبرې اجازه راکوي:
\begin{defn}
ناطق عددونه د \LR{$\Ratequiv$} له مخې د صحيح عددونو د جوړو د هم ارزښتۍ
ټولګي دي (چې دويم غړے ئې صفر نۀ وي)۔ يعنې \LR{$\Rat =
\equivclass{(\Int \times (\Int\setminus \{0_\Int\}))}{\Ratequiv}$}۔
\end{defn}
\olpseganchor{OLP-0043-B005}{end}

\olpseganchor{OLP-0043-B006}{start}
لکه د صحيح عددونو په شان، ځينې بنسټيزې عمليې هم تعريفول غواړو۔
چې \LR{$\equivrep{i,j}{\Ratequiv}$} د \LR{$\Ratequiv$} له مخې هغه د هم
ارزښتۍ ټولګے وي چې \LR{$\tuple{i, j}$} ئې يو غړے دے، وايو:
\begin{align*}
	\equivrep{a, b}{\Ratequiv} + \equivrep{c, d}{\Ratequiv} &= \equivrep{ad + bc,  bd}{\Ratequiv}\\
	\equivrep{a, b}{\Ratequiv} \times \equivrep{c, d}{\Ratequiv} &= \equivrep{a  c, b d}{\Ratequiv}.
\intertext{\RL{پر دې ناطقو عددونو د \LR{$r \leq s$} د تعريفولو دپاره دا حقيقت کاروو چې
\LR{$r \le s$} هله او يوازې هله چې \LR{$s - r$} منفي نۀ وي؛ يعنې \LR{$s - r$} د
\LR{$\nicefrac{i}{j}$} په بڼه ليکل کېدے شي چې \LR{$i$} نامنفي او \LR{$j$}~مثبت وي:}}
	\equivrep{a, b}{\Ratequiv} \leq \equivrep{c, d}{\Ratequiv} &\text{\RL{ هله او يوازې هله چې }}
	\equivrep{c, d}{\Ratequiv} - \equivrep{a, b}{\Ratequiv} =
	\equivrep{i_\Int, j_\Int}{\Ratequiv}
\end{align*}
د ځينو \LR{$i \in \Nat$} او \LR{$0 \neq j \in \Nat$} دپاره۔
\olpseganchor{OLP-0043-B006}{end}

\olpseganchor{OLP-0043-B007}{start}
بيا بايد وارزوو چې دا تعريفونه لکه څنګه چې \emph{بايد}، هماغسې
چلېږي؛ او دا کار \olref[check]{sec} ته پرېږدو۔ خو رښتيا هم سم
چلېږي!{} په پاے کښې صحيح عددونه د ناطقو عددونو \emph{په توګه}
د کارولو يوه لار غواړو؛ نو د هر \LR{$i \in \Int$} دپاره ټاکو چې
\LR{$i_\Rat = \equivrep{i, 1_\Int}{\Ratequiv}$}۔ يو ځل بيا په
\olref[check]{sec} کښې ارزوو چې دا ټول سم چلند کوي۔
\olpseganchor{OLP-0043-B007}{end}

\olpseganchor{OLP-0043-B008}{start}
\begin{prob}
وښيئ چې د هر \LR{$i, j \in \Int$} دپاره \LR{$(i + j)_\Rat = i_\Rat+ j_\Rat$}
او \LR{$(i \times j)_\Rat = i_\Rat \times j_\Rat$} او
\LR{$i \leq j \liff i_\Rat \leq j_\Rat$}۔
\end{prob}
\olpseganchor{OLP-0043-B008}{end}

\olpseganchor{OLP-0044-B004}{start}
\olfileid{sfr}{arith}{real}
\olsection{د حقيقي عددونو کرښه}
\olpseganchor{OLP-0044-B004}{end}

\olpseganchor{OLP-0044-B005}{start}
بل ګام دا ښودل دي چې له ناطقو عددونو څخه حقيقي عددونه څنګه
جوړېږي۔ تر هغې مخکښې بايد پوه شو چې د حقيقي عددونو
\emph{ځانګړتيا} څه ده۔
\olpseganchor{OLP-0044-B005}{end}

\olpseganchor{OLP-0044-B006}{start}
حقيقي عددونه ناطقو عددونو ته ډېر ورته چلند کوي۔ (په فني ډول، دواړه
د \emph{مرتبو ميدانونو} بېلګې دي؛ د دې د تعريف دپاره
\olref[check]{orderedfield} وګورئ۔) اوس که مو د
\olref[sfr][siz][]{chap} تمرينونه حل کړي وي، نو پوهېږئ چې حقيقي
عددونه له ناطقو عددونو په کلکه زيات دي؛ يعنې
\LR{$\cardless{\Rat}{\Real}$}۔ دا خبره لومړے کانتور ثابته کړه۔ خو شاوخوا
دوه نيم زره کاله کېږي چې غير ناطق عددونه، يعنې هغه حقيقي عددونه چې
ناطق نۀ دي، پېژندل شوي دي۔ په رښتيا:
%\footnote{اوس: ``$\sqrt{2}$'' مثبت او منفي دواړو جذرونو ته يو شان اشاره کوي؛ خو په راتلونکو څو مخونو کښې به دا هېر کړم او يوازې مثبت جذر به په پام کښې نيسم۔}
\begin{thm}\ollabel{root2irrational}
\LR{$\sqrt{2}$} ناطق نۀ دے؛ يعنې \LR{$\sqrt{2} \notin \Rat$}
\end{thm}
\olpseganchor{OLP-0044-B006}{end}

\olpseganchor{OLP-0044-B007}{start}
\begin{proof}
د تناقض دپاره فرض کړئ چې \LR{$\sqrt{2}$} ناطق دے۔ نو د ځينو طبيعي
عددونو \LR{$m$} او \LR{$n$} دپاره \LR{$\sqrt{2} = \nicefrac{m}{n}$}۔ په حقيقت
کښې \LR{$m$} او \LR{$n$} داسې ټاکلے شو چې کسر نور نۀ شي ساده کېدے۔ په بيا
ترتيبولو \LR{$m^{2} = 2n^{2}$}۔ له دې ځايه ثبوت په دوو لارو بشپړولے شو:
\olpseganchor{OLP-0044-B007}{end}

\olpseganchor{OLP-0044-B008}{start}
\emph{لومړے، په هندسي ډول} (د ټېننباوم په پيروۍ)۔\footnote{دا
ثبوت \mbox{\LR{\latinfont \cite{Conway2006}}} روايت کړے دے۔} دا مربعګانې په پام کښې ونيسئ:
\begin{center}	
	\begin{tikzpicture}
		\draw[thick] (0,0) rectangle (3,3);
		\draw[thick, fill=red!50] (0,0) rectangle (2.3,2.3);
		\draw[thick, fill=yellow!50] (0.7,0.7) rectangle (3, 3);
		\draw[thick, fill=orange!50] (0.7,0.7) rectangle (2.3, 2.3);
		\draw[<->] (4, 0.7)--(4, 3);
		\node at (4.25, 1.85) (n) {\LR{$n$}};
		\draw[<->] (5, 0)--(5, 3);
		\node at (5.25, 1.5) (m) {\LR{$m$}};
	\end{tikzpicture}
\end{center}
ځکه چې \LR{$m^2 = 2n^2$}، هغه سيمه چې د ضلعې \LR{$n$} دواړه مربعګانې
يو پر بل راځي، له هغې سيمې سره مساوي مساحت لري چې له دغو دواړو
مربعګانو هېڅ يو نۀ ده پوښلې؛ يعنې د نارنجي مربع مساحت د دواړو
بې رنګه مربعګانو د مساحتونو له مجموعې سره مساوي دے۔ نو که د نارنجي
مربع ضلعه \LR{$p$} او د هر بې رنګه مربع ضلعه \LR{$q$} وي، نو \LR{$p^2 = 2q^2$}۔
خو اوس \LR{$\sqrt{2} = \nicefrac{p}{q}$}، چې \LR{$p < m$}، \LR{$q < n$} او
\LR{$p, q \in \Nat$} دي۔ دا له دې خبرې سره تناقض لري چې \LR{$m$} او \LR{$n$}
تر ټولو واړه ممکن عددونه ټاکل شوي وو۔
\olpseganchor{OLP-0044-B008}{end}
		
\olpseganchor{OLP-0044-B009}{start}
\emph{دويم، په صوري ډول۔} ځکه چې \LR{$m^{2} = 2n^{2}$}، نو \LR{$m$} جفت
دے۔ (دا ښودل اسانه دي چې که \LR{$x$} طاق وي، نو \LR{$x^2$} هم طاق وي۔)
نو د کوم \LR{$r \in \Nat$} دپاره \LR{$m = 2r$}۔ په بيا ترتيبولو
\LR{$2r^2 = n^2$}،
		%				\begin{align*}
		%					(2q)^{2} &= 2n^{2}\\
		%					4q^{2} &= 2n^{2}\\
		%					2q^{2} &= n^{2}
		%				\end{align*}
نو \LR{$n$} هم جفت دے۔ په دې توګه \LR{$m$} او \LR{$n$} دواړه جفت دي، او له دې
امله کسر \LR{$\nicefrac{m}{n}$} نور \emph{ساده کېدے} شي۔ تناقض!
\end{proof}
\olpseganchor{OLP-0044-B009}{end}

\olpseganchor{OLP-0044-B010}{start}
په ضمن کښې دا انځوريز ثبوت موږ ته اجازه راکوي چې د \olref[his][set][mythology]{sec} مواد بيا وګورو۔ ټېننباوم (1927--2006) بشپړ معاصر رياضي پوه و؛ خو ثبوت بې شکه ښکلے، پوره دقيق او پر هندسي پوهېدنه ولاړ دے!
\olpseganchor{OLP-0044-B010}{end}

\olpseganchor{OLP-0044-B011}{start}
په هر حال، حقيقي عددونه تر ناطقو عددونو «ډېر پراخ» دي۔ په يوه معنا په ناطقو عددونو کښې «تشيالونه» شته چې حقيقي عددونه ئې ډکوي۔ وايرشټراس وپېژندله چې دا د حقيقي عددونو يو خاصيت بيانوي، کوم چې هغوے له ناطقو عددونو بېلوي؛ يعنې د بشپړتيا خاصيت: \emph{د حقيقي عددونو هر ناتش سټ چې پاسنی حد لري، تر ټولو لږ پاسنی حد هم لري۔}
\olpseganchor{OLP-0044-B011}{end}

\olpseganchor{OLP-0044-B012}{start}
دا ليدل اسانه دي چې ناطق عددونه د بشپړتيا خاصيت نۀ لري۔ د بېلګې په ډول له \LR{$\sqrt{2}$} څخه د وړو ناطقو عددونو دا سټ په پام کښې ونيسئ:
\[
	\Setabs{p \in \Rat}{p^2 < 2 \text{\RL{ يا }}p < 0}
\]
دا په ناطقو عددونو کښې پاسنی حد لري؛ د بېلګې په ډول د دې ټول
غړي له \LR{$3$} څخه واړه دي۔ خو د دې تر ټولو لږ
پاسنی حد څه دے؟ غواړو ووايو `\LR{$\sqrt{2}$}'؛ خو همدا اوس مو وليدل
چې \LR{$\sqrt{2}$} \emph{ناطق نۀ} دے۔ او له \LR{$\sqrt{2}$} څخه لوے تر ټولو
لږ ناطق عدد نشته۔ نو دا سټ پاسنی حد لري خو تر ټولو لږ پاسنی حد نۀ
لري۔ له دې امله ناطق عددونه د بشپړتيا خاصيت نۀ لري۔
\olpseganchor{OLP-0044-B012}{end}

\olpseganchor{OLP-0044-B013}{start}
د دې پر خلاف، پيوستار «له معنوي پلوه بايد» د بشپړتيا خاصيت ولري۔ موږ يوازې دا نۀ غواړو چې \LR{$\sqrt{2}$} يو حقيقي عدد وي؛ غواړو د ناطقو عددونو د کرښې ټول «تشيالونه» ډک کړو۔ په رښتيا غواړو چې پخپله پيوستار هېڅ «تشيال» ونۀ لري۔ همدا هغه څه دي چې د بشپړتيا له لارې به ئې ترلاسه کړو۔
\olpseganchor{OLP-0044-B013}{end}

\olpseganchor{OLP-0045-B003}{start}
\olfileid{sfr}{arith}{cuts}
\olsection{له \texorpdfstring{\LR{$\Rat$}}{ℚ} څخه \texorpdfstring{\LR{$\Real$}}{ℝ} ته}
\olpseganchor{OLP-0045-B003}{end}

\olpseganchor{OLP-0045-B004}{start}
په اصل کښې د بشپړتيا خاصيت ښيي چې د حقيقي عددونو د کرښې هر ټکی
\LR{$\alpha$} دا کرښه په بشپړ ډول پر دوو برخو وېشي: هغه عددونه چې
\LR{$\alpha$} ئې تر ټولو لږ پاسنی حد دے، او هغه عددونه چې \LR{$\alpha$} ئې
تر ټولو لوے لاندينی حد دے۔ ډېډېکېنډ وړانديز وکړ چې له ناطقو عددونو
څخه د حقيقي عددونو د \emph{جوړولو} دپاره، حقيقي عددونه يوازې د هغو
\emph{پرېکونونو} په توګه وګڼو چې ناطق عددونه وېشي۔ يعنې \LR{$\sqrt{2}$}
له هغه \emph{پرېکون} سره يو ګڼو چې ناطق عددونه \LR{$< \sqrt{2}$} له
ناطق عددونو \LR{$> \sqrt{2}$} څخه بېلوي۔
\olpseganchor{OLP-0045-B004}{end}

\olpseganchor{OLP-0045-B005}{start}
راځئ دا خبره سمه منظمه کړو۔ که ناطق عددونه پر دوو برخو پرې کړو،
نو يوازې د هغې \emph{لاندېنۍ} برخې په کتلو خپل وېش په يوازيني ډول
پېژندلے شو۔ نو په دقيق ډول لاندې تعريف وړاندې کوو:
\olpseganchor{OLP-0045-B005}{end}

\olpseganchor{OLP-0045-B006}{start}
\begin{defn}[پرېکون]
يو \emph{پرېکون} \LR{$\alpha$} د ناطقو عددونو هره ناتشه خاصه پيلنۍ
برخه ده چې تر ټولو لوے غړے نۀ لري۔ يعنې \LR{$\alpha$} هله او يوازې هله
پرېکون دے چې:
\begin{enumerate}
	\item \emph{ناتش، خاص}: \LR{$\emptyset \neq \alpha \subsetneq \Rat$}
	\item \emph{پيلنۍ}: د ټولو \LR{$p,q \in \Rat$} دپاره: که \LR{$p < q \in \alpha$} وي، نو \LR{$p \in \alpha$}
	\item \emph{تر ټولو لوے غړے نۀ لري}: د هر \LR{$p \in \alpha$} دپاره داسې \LR{$q \in \alpha$} شته چې \LR{$p < q$} وي
\end{enumerate}
بيا \LR{$\Real$} د پرېکونونو سټ دے۔
\end{defn}
\olpseganchor{OLP-0045-B006}{end}

\olpseganchor{OLP-0045-B007}{start}
نو اوس ويلے شو چې \LR{$\sqrt{2} = \Setabs{p \in \Rat}{p^2 < 2\text{\RL{
يا }}p < 0}$}۔ البته بايد وارزوو چې دا په رښتيا يو \emph{پرېکون}
دے، خو دا کار \olref[check]{sec} ته پرېږدو۔
\olpseganchor{OLP-0045-B007}{end}

\olpseganchor{OLP-0045-B008}{start}
لکه مخکښې، چې ځينې څيزونه مو تعريف کړل، بيا بايد پر هغو بنسټيزې
تابعې او اړيکې تعريف کړو۔ په يوه اسانه تعريف پيل کوو:
\begin{align*}
	\alpha \leq \beta \text{\RL{ هله او يوازې هله چې }}\alpha \subseteq \beta
\end{align*}
د مرتبې دا تعريف موږ ته اجازه راکوي چې مرکزي پايله \emph{بيان}
کړو: د پرېکونونو سټ د بشپړتيا خاصيت لري۔ په بشپړ ډول د بيان بڼه
داسې ده۔ که \LR{$S$} د پرېکونونو يو ناتش سټ وي چې پاسنی حد لري، نو
\LR{$S$} تر ټولو لږ پاسنی حد لري۔ په لا تفصيل: يو پرېکون \LR{$\lambda$} شته
چې د \LR{$S$} پاسنی حد دے، يعنې \LR{$(\forall \alpha \in S)\alpha \subseteq
\lambda$}، او \LR{$\lambda$} تر ټولو لږ داسې پرېکون دے، يعنې
\LR{$(\forall \beta \in \Real)((\forall \alpha \in S)\alpha \subseteq \beta \lif \lambda \subseteq \beta)$}۔ اوس د پايلې ثبوت دا دے:
\olpseganchor{OLP-0045-B008}{end}

\olpseganchor{OLP-0045-B009}{start}
\begin{thm}\ollabel{realcompleteness}
د پرېکونونو سټ د بشپړتيا خاصيت لري۔
\end{thm}
\olpseganchor{OLP-0045-B009}{end}

\olpseganchor{OLP-0045-B010}{start}
\begin{proof}
فرض کړئ \LR{$S$} د پرېکونونو هر ناتش سټ دے چې پاسنی حد لري۔ وټاکئ
\LR{$\lambda = \bigcup S$}۔
%\Setabs{p \in \Rat}{(\exists \alpha \in S)p \in \alpha}$$
لومړے ادعا کوو چې \LR{$\lambda$} يو پرېکون دے:
\begin{enumerate}
\item ځکه چې \LR{$S$} ناتش دے، لږ تر لږه يو پرېکون په \LR{$S$} کښې دے، نو
\LR{$\emptyset \neq \lambda$}۔ ځکه چې \LR{$S$} د پرېکونونو سټ
دے، \LR{$\lambda \subseteq \Rat$}۔ ځکه چې \LR{$S$} پاسنی حد لري، کوم
\LR{$p \in \Rat$} د هر پرېکون \LR{$\alpha \in S$} څخه بهر دے۔ نو
\LR{$p\notin \lambda$}، او له دې امله \LR{$\lambda \subsetneq \Rat$}۔
\item فرض کړئ \LR{$p < q \in \lambda$}۔ نو داسې \LR{$\alpha \in S$} شته
چې \LR{$q \in \alpha$}۔ ځکه چې \LR{$\alpha$} يو پرېکون دے،
\LR{$p \in \alpha$}۔ نو \LR{$p \in \lambda$}۔
\item فرض کړئ \LR{$p \in \lambda$}۔ نو داسې \LR{$\alpha \in S$} شته چې
\LR{$p \in \alpha$}۔ ځکه چې \LR{$\alpha$} يو پرېکون دے، داسې \LR{$q \in
\alpha$} شته چې \LR{$p < q$}۔ نو \LR{$q \in \lambda$}۔
\end{enumerate}
په دې سره ادعا ثابته شوه۔ برسېره پر دې، په څرګند ډول
\LR{$(\forall \alpha \in S)\alpha \subseteq \bigcup S = \lambda$}؛
يعنې \LR{$\lambda$} د \LR{$S$} يو پاسنی حد دے۔ اوس فرض کړئ
\LR{$\beta \in \mathbb{R}$} هم پاسنی حد دے؛ يعنې
\LR{$(\forall \alpha \in S)\alpha \subseteq \beta$}۔ د هر
\LR{$p \in \Rat$} دپاره، که \LR{$p \in \lambda$} وي، نو داسې
\LR{$\alpha \in S$} شته چې \LR{$p \in \alpha$}، او په دې توګه
\LR{$p \in \beta$}۔ په عمومي کولو \LR{$\lambda \subseteq \beta$}۔ نو
\LR{$\lambda$} د \LR{$S$} \emph{تر ټولو لږ} پاسنی حد دے۔
\end{proof}
\olpseganchor{OLP-0045-B010}{end}

\olpseganchor{OLP-0045-B011}{start}
نو موږ داسې ډېر څيزونه لرو چې د بشپړتيا خاصيت پوره کوي۔ د دې د
ويلو يوه لار دا ده: زموږ په پرېکونونو کښې هېڅ «تشيال» نشته۔
(نو د ناطقو عددونو پر ځاے د حقيقي عددونو نور «پرېکونونه» اخيستل
به هېڅ په زړه پورې نوي څيزونه رانۀ کړي۔)
\olpseganchor{OLP-0045-B011}{end}

\olpseganchor{OLP-0045-B012}{start}
بيا بايد پر حقيقي عددونو ځينې عمليې تعريف کړو۔ ناطق عددونه په
حقيقي عددونو کښې په دې ټاکنه ورګډوو چې د هر \LR{$p \in \Rat$} دپاره
\LR{$p_\Real = \Setabs{q \in \Rat}{q < p}$}۔ بيا تعريفوو:
\begin{align*}
	\alpha + \beta &= \Setabs{p + q}{p \in \alpha \land q \in \beta}\\
%	\alpha - \beta &\defis \Setabs{p - q}{p \in \alpha \land q \in \Rat \setminus \beta}\\
	\alpha \times \beta &=
	\Setabs{p \times q}{0 \leq p \in \alpha \land 0 \leq q \in \beta} \cup 0_\Real & \text{\RL{که }}\alpha, \beta \geq 0_\Real
%	\alpha \div \beta &\defis
%	\Setabs{\nicefrac{p}{q}}{p \in \alpha \land q \in \Rat \setminus \beta}  & \text{که }\alpha \geq 0_\Real, \beta > 0_\Real
\end{align*}
د هر پرېکون لپاره د صفر د ضرب قاعده دا ده چې
\LR{$0_\Real \times \alpha = 0_\Real = \alpha \times 0_\Real$}۔
د نفي لپاره وټاکئ: %چې $0_\Real \times \alpha = 0_\Real = \alpha \times 0_\Real$، او ورزيات کړئ:
\begin{align*}
	-\alpha &= \Setabs{p - q}{p < 0 \land q \notin \alpha}
\end{align*}
او بيا وټاکئ:
\begin{align*}
	\alpha \times \beta &\defis
	\begin{cases}
		\mathord{-}\alpha \times \mathord{-}\beta &\text{\RL{که }}\alpha < 0_\Real\text{\RL{ او }}\beta < 0_\Real\\
		\mathord{-}(\mathord{-}\alpha \times \beta) &\text{\RL{که }}\alpha < 0_\Real \text{\RL{ او }}\beta > 0_\Real\\
		\mathord{-}(\alpha \times \mathord{-}\beta) &\text{\RL{که }}\alpha > 0_\Real \text{\RL{ او }}\beta < 0_\Real
	\end{cases}
\end{align*}
د تفريق تعريف \LR{$\alpha-\beta=\alpha+(-\beta)$} دے۔ که \LR{$\beta>0_\Real$} وي،
\LR{$\beta^{-1}$} د هغو ناطقو عددونو \LR{$p$} پرېکون دے چې يا \LR{$p\leq 0$} وي،
يا داسې مثبت ناطق عدد \LR{$q\notin\beta$} وي چې \LR{$p<1/q$} وي۔ که
\LR{$\beta<0_\Real$} وي، \LR{$\beta^{-1}=-((-\beta)^{-1})$} وټاکئ۔
د ناصفر مخرج لپاره \LR{$\alpha\div\beta=\alpha\times\beta^{-1}$} دے۔
بيا بايد وارزوو چې دا هر تعريف تل يو پرېکون ورکوي۔ په پاے کښې بايد په يوه اسانه (خو اوږده)
څرګندونه کښې وښيو چې په دې ډول تعريف شوي پرېکونونه کټ مټ لکه څنګه
چې بايد، هماغسې چلېږي۔ خو دا ټول \olref[check]{sec} ته پرېږدو۔
\olpseganchor{OLP-0045-B012}{end}

\olpseganchor{OLP-0046-B004}{start}
\olfileid{sfr}{arith}{ref}
\olsection{ځينې فلسفي غورونه}
\olpseganchor{OLP-0046-B004}{end}

\olpseganchor{OLP-0046-B005}{start}
فني خبرې تر دې ځايه بس دي۔ خو هغو څه ترلاسه کړل؟
\olpseganchor{OLP-0046-B005}{end}

\olpseganchor{OLP-0046-B006}{start}
ښه، دا خبره تقريباً له اختلافه پاکه ده چې هغو موږ ته څه {ښکلې}
سوچه رياضي راکړه۔ برسېره پر دې، ځينې ژورې تصوري لاسته راوړنې هم
وې۔ دا پېژندنه ژوره وه چې د بشپړتيا خاصيت د حقيقي او ناطقو عددونو
تر منځ بنسټيز توپير بيانوي۔ همدارنګه د ډېډېکېنډ د پرېکونونو په توګه
د حقيقي عددونو ښکاره جوړښت د تحليل موضوع ته کلک بنسټ ورکوي۔ پوهېږو
چې د \emph{بشپړ مرتب ميدان} تصور بې تناقضه دے، ځکه پرېکونونه کټ
مټ همداسې يو ميدان جوړوي۔
\olpseganchor{OLP-0046-B006}{end}

\olpseganchor{OLP-0046-B007}{start}
له دې ټولو سره، د دغو لاسته راوړنو په اړه بايد څو تحفظات څرګند کړو۔
\olpseganchor{OLP-0046-B007}{end}

\olpseganchor{OLP-0046-B008}{start}
لومړے، دا څرګنده نۀ ده چې د پرېکونونو په توګه د حقيقي عددونو ګڼل
د هغو د پېژندل شوو (ښايي نامتناهي) اعشاري پراختياوو په توګه تر ګڼلو
څومره \emph{زيات} دقيق دي۔ د حقيقي عددونو دا وروستے «جوړښت» د
کوشي تسلسلونو له لارې د حقيقي عددونو جوړښت ته څه ورته والے لري؛ خو
په حقيقت کښې رياضي پوهان د اوولسمې پېړۍ له پيله په بنسټيز ډول پرې
پوهېدل (\olref[cauchy]{sec} وګورئ)۔ په دقت کښې رښتينے زياتوالے له
دې پېژندنې راغے چې حقيقي عددونه د بشپړتيا خاصيت لري؛ يوازې دا وړتيا
چې حقيقي عددونه د ځانګړو سټونو په توګه جوړ کړو، ښايي پخپله دومره
په زړه پورې نۀ وي۔
\olpseganchor{OLP-0046-B008}{end}

\olpseganchor{OLP-0046-B009}{start}
دا لا هم لږ څرګنده ده چې د صحيح يا ناطقو عددونو (ډېرې اسانه) سټي
جوړونه په دغو ډګرونو کښې دقت زياتوي۔ دلته يوه ساده خبره د يادولو وړ
ده۔ صحيح عددونه مو د طبيعي عددونو د مرتبو جوړو د هم ارزښتۍ ټولګيو
په توګه \emph{جوړ} کړل، بيا مو ناطق عددونه د صحيح عددونو د مرتبو
جوړو د هم ارزښتۍ ټولګيو په توګه جوړ کړل، او بيا مو حقيقي عددونه د
ناطقو عددونو د سټونو په توګه جوړ کړل؛ خو سمدستي دا جوړښتونه
\emph{هېروو}۔ په ځانګړي ډول هېڅوک به کله هم ونۀ غواړي چې په يوه
رياضياتي ثبوت کښې دا جوړښتونه \emph{وکاروي} (البته د دې ثبوت پرته
چې جوړښتونه لکه څنګه چې ټاکل شوي وو، هماغسې چلېږي)۔ د يوه حقيقي
عدد په اړه نېغ په نېغه خبرې کول د طبيعي عددونو د سټونو د سټونو د
سټونو د سټونو د سټونو د سټونو د کوم سټ په اړه تر خبرو ډېر اسانه دي۔
%(ځکه حقيقي عددونه د ناطقو عددونو د سټونو په توګه جوړ شول؛ او دا د صحيح عددونو د مرتبو جوړو د هم ارزښتۍ ټولګيو، يعنې د صحيح عددونو د سټونو د سټونو د سټونو په توګه جوړ شول؛ او داسې نور۔)
%2/3 = [2,3]
%	يعنې {<2,3>, ... }
%	{ {{2},{2,3}}, ... }
%
%حقيقي عددونه
%	د ناطقو عددونو سټونه دي
%
%ناطق عددونه
%	د صحيح عددونو د سټونو د سټونو سټونه دي
%
%صحيح عددونه
%	د طبيعي عددونو د سټونو د سټونو سټونه دي
\olpseganchor{OLP-0046-B009}{end}
	
\olpseganchor{OLP-0046-B010}{start}
تر ټولو زيات شک په دې دے چې دا تعريفونه موږ ته وايي چې صحيح، ناطق
يا حقيقي عددونه په \emph{ماوراءالطبيعتي معنا څه دي}۔ يعنې په دې شک
دے چې حقيقي عددونه (مثلاً) ځينې ځانګړي سټونه (د سټونو د سټونو
\ldots) \emph{دي}۔ د دې نظر پر وړاندې ستر خنډ دا دے چې جوړښت په
ډېرو بېلابېلو لارو کېدے شو۔ د حقيقي عددونو په حالت کښې ځينې په
رښتيا په زړه پورې بېلابېل جوړښتونه شته (\olref[cauchy]{sec}
وګورئ)۔ خو د بېلابېلو جوړښتونو د ترلاسه کولو يوه ډېره ساده لار دا
ده: لکه په \olref[sfr][rel][ref]{sec} کښې، مرتبې جوړې مو لږ په
بله توګه تعريفولے شوې؛ که مو د مرتبې جوړې دا بديل تصور کارولے وے،
زموږ جوړښتونو به کټ مټ همدومره سم کار کړے وے، خو په پاے کښې به مو
بېل څيزونه ترلاسه کړي وو۔ له دې امله د صحيح، ناطقو او حقيقي عددونو
ډېر سيال سټي جوړښتونه شته۔ اوس دا ادعا چې صحيح عددونه (مثلاً)
\emph{دا} سټونه دي، نۀ \emph{هغه}، يوازې خپلسرې (او شرموونکې) به
وي۔ (لکه په \olref[sfr][rel][ref]{sec} کښې، دا د هغه استدلال يوه
بېلګه ده چې \mbox{\LR{\latinfont \citealt{Benacerraf1965}}} مشهور کړ۔)
\olpseganchor{OLP-0046-B010}{end}

\olpseganchor{OLP-0046-B011}{start}
يو بل ټکی هم د راپورته کولو وړ دے: زموږ په جوړښتونو کښې يو ډېر
\emph{عجيب} څيز شته۔ موږ په طبيعي عددونو پيل وکړ۔ بيا صحيح عددونه
جوړوو، او «د صحيح عددونو \LR{$0$}»، يعنې
\LR{$ \equivrep{0,0}{\Intequiv}$}، جوړوو۔ خو
\LR{$0 \neq \equivrep{0,0}{\Intequiv}$}۔ په حقيقت کښې زموږ د جوړښتونو
له مخې \emph{هېڅ} طبيعي عدد صحيح عدد نۀ دے۔ خو دا د شهودي پوهې
سخت خلاف ښکاري۔ همدارنګه په \olref[sfr][set][imp]{sec} کښې مو له
ډېر استدلال پرته ادعا وکړه چې \LR{$\Nat \subseteq \Rat$}۔ که جوړښتونه
موږ ته په دقيق ډول وايي چې عددونه \emph{څه} دي، نو دا ادعا په
څرګند ډول ناسمه وه۔
\olpseganchor{OLP-0046-B011}{end}

\olpseganchor{OLP-0046-B012}{start}
نو چې يو ګام شاته ودرېږو، کوم ځاے ته رسېږو؟ په ساده سټ نظريه کښې
په کار کولو او طبيعي عددونو ته په ورکړل شوو څيزونو کتلو، صحيح، ناطق
او حقيقي عددونه د ځينو سټونو په توګه \emph{کارولے} شو۔ په دې معنا
د دغو څيزونو نظريې په سټ نظريه کښې \emph{ورګډولے} شو۔ خو د دې
ورګډولو فلسفي ارزښت دومره ساده نۀ دے۔
\olpseganchor{OLP-0046-B012}{end}

\olpseganchor{OLP-0046-B013}{start}
البته دا هېڅ د بحث وروستۍ خبره نۀ ده!{} ټکی يوازې دا دے: د دې د
ښودلو دپاره چې د حقيقي عددونو سټي جوړونه ژور فلسفي ارزښت \emph{لري}،
يو څه نور \emph{فلسفي} استدلال ته اړتيا ده۔
\olpseganchor{OLP-0046-B013}{end}

\olpseganchor{OLP-0046-B014}{start}
%برسېره پر دې، د دې ټول څپرکي په اوږدو کښې مو طبيعي عددونه يوازې
%د راکړل شوو څيزونو په توګه ګڼلي دي۔ خو له هغو سره څه کولے شو؟
%دا د بل څپرکي موضوع ده۔
\olpseganchor{OLP-0046-B014}{end}

\olpseganchor{OLP-0047-B003}{start}
\olfileid{sfr}{arith}{check}
\olsection{مرتبې حلقې او ميدانونه}
\olpseganchor{OLP-0047-B003}{end}

\olpseganchor{OLP-0047-B004}{start}
په دې ټول څپرکي کښې مو ادعا وکړه چې ځينې تعريفونه «لکه څنګه چې
بايد» هماغسې چلېږي۔ په دې فني ضميمه کښې به روښانه کړو چې مطلب مو
څه دے، او (لنډه خاکه به ورکړو چې څنګه) وښيو تعريفونه «سم» چلېږي۔
لومړے ګورو چې عمليې د استازي په بدلېدو سره نۀ بدلېږي۔
که د صحيح عدد دوه استازي \LR{$\tuple{a,b}$} او \LR{$\tuple{a',b'}$} يو
ټولګے ولري، نو \LR{$a+b'=a'+b$} دے۔ د لومړۍ جوړې دواړو برخو ته
\LR{$b'$} او د دويمې جوړې دواړو برخو ته \LR{$b$} ورزيات کړئ۔ په دواړو
لارو يوه جوړه تر لاسه کېږي۔ که د يوه استازي دواړو برخو ته
\LR{$k$} ورزيات کړو، د جمع د پايلې دواړو برخو ته هم \LR{$k$} ورزياتېږي۔
که بل عامل \LR{$\tuple{c,d}$} وي، د ضرب د پايلې دواړو برخو ته
\LR{$k(c+d)$} ورزياتېږي۔ دا زياتوالے د پايلې ټولګے نۀ بدلوي۔
د جمع او ضرب د فورمولونو د تقارن له مخې همدا دليل د بل استازي
لپاره هم کار کوي۔ د ترتيب لپاره د دوو جوړو د اړوندو برخو پرتله
هغه وخت هماغسې پاتې کېږي چې په دواړو خواوو کښې يو طبيعي عدد ورزيات شي۔
په ناطقو عددونو کښې مخرجونه ناصفر دي۔ د دوو ناصفرو صحيح عددونو
ضرب هم ناصفر دے، ځکه هر يو د يوه مثبت طبيعي عدد يا د هغه د
جمعي معکوس استازيتوب کوي، او د دوو مثبتو طبيعي عددونو ضرب مثبت دے۔
د استازو د برابرۍ له شرطه \LR{$ab'=a'b$} اخلو۔ د ثابت بل استازي
\LR{$\tuple{c,d}$} لپاره د جمع د استازو چليپايي ضربونه
\LR{$(ad+bc)(b'd)$} او \LR{$(a'd+b'c)(bd)$} دي او يو له بل سره برابر دي۔
د ضرب لپاره \LR{$(ac)(b'd)=(a'c)(bd)$} دے۔ د جمع او ضرب تقارن
د بل استازي بدلېدل هم رانغاړي۔ جمعي معکوس هم خپلواک دے، ځکه
د صورت دواړه اړوند چليپايي ضربونه په منفي کولو بيا برابر وي۔
نو توپير او د هغه پر ټولګي ولاړ ترتيب هم د استازي په ټاکلو پورې
تړلي نۀ دي۔
\olpseganchor{OLP-0047-B004}{end}

\olpseganchor{OLP-0047-B005}{start}
په \olref[int]{sec} کښې مو پر \LR{$\Int$} جمع او ضرب تعريف کړل۔ غواړو
وښيو چې دا عمليې، لکه څنګه چې تعريف شوې دي، \LR{$\Int$} ته هغه جوړښت
ورکوي چې موږ ئې «غواړو»۔ په ځانګړي ډول دا جوړښت يوه تبادلي حلقه ده۔
\olpseganchor{OLP-0047-B005}{end}

\olpseganchor{OLP-0047-B006}{start}
\begin{defn}
	يوه \emph{تبادلي حلقه} يو سټ \LR{$S$} دے چې ځانګړي غړي \LR{$0$} او \LR{$1$} او عمليې \LR{$+$} او \LR{$\times$} لري، او دا اته فارمولونه پوره کوي:
	\begin{align*}
		\emph{اتحاديت}&&a + (b+ c) & = (a + b) + c \\
		&& (a \times b) \times c & = a \times (b\times c)\\
		\emph{تبادليت}&&a + b &= b+ a  \\
		&&  a \times b&= b\times a\\
		\emph{عيني غړي}&&a + 0 &= a \\
		&& a \times 1 &= a\\
		\emph{جمعي معکوس}&&(\exists b\in S)0&=a + b\\
		\emph{توزيعيت}&&a \times (b+ c ) &= (a \times b) + (a \times c)
	\end{align*}
	په ضمني ډول دا ټول په هغو کلي مقداري نښو تړل شوي چې ساحه ئې تر \LR{$S$} محدوده ده۔ او پام وکړئ چې دلته د \LR{$0$} او~\LR{$1$} غړي لازماً په همدې نومونو طبيعي عددونه نۀ دي۔
\end{defn}
\olpseganchor{OLP-0047-B006}{end}

\olpseganchor{OLP-0047-B007}{start}
نو د دې د ارزولو دپاره چې صحيح عددونه تبادلي حلقه جوړوي، يوازې
دا اته شرطونه ارزول غواړو۔ د شرطونو ثابتول {ګران} نۀ دي، خو لږ
اوږد کار دے۔ د بېلګې په ډول د جمع په حالت کښې \emph{اتحاديت}
داسې ثابتوو:
\olpseganchor{OLP-0047-B007}{end}

\olpseganchor{OLP-0047-B008}{start}
\begin{proof}
\LR{$i, j, k \in \Int$} ثابت ونيسئ۔ نو داسې \LR{$a_1, b_1, a_2, b_2, a_3, b_3 \in
\Nat$} شته چې \LR{$i = \equivrep{a_1, b_1}{}$} او \LR{$j =
\equivrep{a_2,b_2}{}$} او \LR{$k = \equivrep{a_3, b_3}{}$}۔ (د لوستلو
د اسانتيا دپاره د ``\LR{$\equivrep{x, y}{\Intequiv}$}'' پر ځاے
``\LR{$\equivrep{x, y}{}$}'' ليکو؛ په دې ټوله برخه کښې به همداسې کوو۔)
اوس:
\begin{align*}
	i + (j + k) &= \equivrep{a_1, b_1}{}+(\equivrep{a_2, b_2}{} + \equivrep{a_3, b_3}{}) \\
	&= \equivrep{a_1,  b_1}{} + \equivrep{a_2+a_3, b_2+b_3}{}\\
	&= \equivrep{a_1 + (a_2 + a_3), b_1 + (b_2 + b_3)}{}\\
	&= \equivrep{(a_1 + a_2) + a_3, (b_1 + b_2) + b_3}{}\\
	&= \equivrep{a_1 + a_2, b_1 + b_2}{} + \equivrep{a_3, b_3}{}\\
	&= (\equivrep{a_1, b_1}{} + \equivrep{a_2, b_2}{}) + \equivrep{a_3, b_3}{}\\
	&= (i+j) + k
\end{align*}
او د \LR{$\Nat$} پر جمع له چلنده ازاد کار اخلو۔
\end{proof}
\olpseganchor{OLP-0047-B008}{end}

\olpseganchor{OLP-0047-B009}{start}
په همدې ډول \emph{جمعي معکوس} داسې ثابتوو:
\olpseganchor{OLP-0047-B009}{end}

\olpseganchor{OLP-0047-B010}{start}
\begin{proof}
\LR{$i \in \Int$} ثابت ونيسئ، نو د ځينو \LR{$a,b \in \Nat$} دپاره
\LR{$i = \equivrep{a,b}{}$}۔ وټاکئ \LR{$j = \equivrep{b,a}{} \in \Int$}۔
د طبيعي عددونو له چلنده په کار اخيستلو
\LR{$(a+b) + 0 = 0 + (a+b)$}، نو د تعريف له مخې
\LR{$\tuple{a+b, b+a} \sim_\Int \tuple{0,0}$}، او له دې امله
\LR{$\equivrep{a+b, b+a}{} = \equivrep{0, 0}{} = 0_\Int$}۔ نو اوس
\LR{$i + j = \equivrep{a,b}{}+\equivrep{b,a}{}=\equivrep{a+b, b+a}{}= \equivrep{0,
0}{} = 0_\Int$}۔
\end{proof}
\olpseganchor{OLP-0047-B010}{end}

\olpseganchor{OLP-0047-B011}{start}
او د \emph{توزيعيت} ثبوت دا دے:
\olpseganchor{OLP-0047-B011}{end}

\olpseganchor{OLP-0047-B012}{start}
\begin{proof}
لکه پورته، \LR{$i = \equivrep{a_1, b_1}{}$}،
\LR{$j = \equivrep{a_2,b_2}{}$} او \LR{$k = \equivrep{a_3, b_3}{}$} ثابت
ونيسئ۔ اوس:
\begin{align*}
	i \times (j + k)
	&= \equivrep{a_1, b_1}{} \times (\equivrep{a_2,b_2}{} + \equivrep{a_3, b_3}{})\\
	&= \equivrep{a_1, b_1}{} \times \equivrep{a_2 + a_3,b_2+b_3}{}\\
	&= \equivrep{a_1  (a_2 + a_3) + b_1  (b_2+b_3), a_1  (b_2 + b_3) + b_1 (a_2 + a_3)}{}\\
	&= \equivrep{a_1 a_2 + a_1a_3 + b_1 b_2+b_1b_3, a_1 b_2 + a_1b_3 + a_2b_1 + a_3b_1}{}\\		
%		&= \equivrep{a_1 a_2 + b_1b_2 + a_1a_3 + b_1 b_2+b_1b_3, a_1 b_2 + a_1b_3 + a_2b_1 + a_3b_1}{}\\		
	&= \equivrep{a_1a_2 + b_1b_2, a_1b_2 + a_2b_1}{} + \equivrep{a_1a_3 + b_1b_3, a_1b_3 + a_3b_1}{}\\
	&= (\equivrep{a_1, b_1}{} \times \equivrep{a_2,b_2}{}) + (\equivrep{a_1, b_1}{} \times  \equivrep{a_3, b_3}{})\\
	&= (i \times j) + (i \times k)
\end{align*}
\end{proof}
\olpseganchor{OLP-0047-B012}{end}

\olpseganchor{OLP-0047-B013}{start}
د پاتې پنځو شرطونو ثبوت د تمرين په توګه پرېږدو۔ چې دا مو وکړل،
ښودلي مو دي چې \LR{$\Int$} يوه تبادلي حلقه جوړوي؛ يعنې جمع او ضرب
(لکه څنګه چې تعريف شوي) هماغسې چلېږي چې بايد۔
\olpseganchor{OLP-0047-B013}{end}

\olpseganchor{OLP-0047-B014}{start}
\begin{prob}
ثابت کړئ چې \LR{$\Int$} يوه تبادلي حلقه ده۔
\end{prob}
\olpseganchor{OLP-0047-B014}{end}

\olpseganchor{OLP-0047-B015}{start}
خو زموږ کار لا ختم نۀ دے۔ پر \LR{$\Int$} مو له جمع او ضرب سره يوه
ترتيبي اړيکه \LR{$\leq$} هم تعريف کړې، او بايد وارزوو چې دا لکه څنګه
چې بايد، هماغسې چلېږي۔ په لا تفصيل بايد وښيو چې \LR{$\Int$} يوه
\emph{مرتبه} حلقه جوړوي۔\footnote{له
	\olref[sfr][rel][ord]{def:linearorder} څخه په ياد ولرئ چې کلي
	ترتيب هغه اړيکه ده چې انعکاسي، انتقالي، ضدتناظري او وصل شوې وي۔
	له وصل‌والي او ضدتناظر څخه د کلک ترتيب \emph{درېګونې}
	پايله اخلو: د هر \LR{$a$} او \LR{$b$} دپاره له \LR{$a < b$}، \LR{$a=b$}
	او \LR{$a>b$} څخه کټ مټ يو حالت صدق کوي۔}
\olpseganchor{OLP-0047-B015}{end}

\olpseganchor{OLP-0047-B016}{start}
\begin{defn}
يوه \emph{مرتبه حلقه} هغه تبادلي حلقه ده چې يوه کلي ترتيبي اړيکه
\LR{$\leq$} هم لري، داسې چې:
\begin{align*}
	a \leq b &\lif a + c \leq b + c\\
	(a \leq b \land 0 \leq c) &\lif a \times c \leq b \times c
\end{align*}
\end{defn}
\olpseganchor{OLP-0047-B016}{end}

\olpseganchor{OLP-0047-B017}{start}
\begin{prob}
ثابت کړئ چې \LR{$\Int$} يوه مرتبه حلقه ده۔
\end{prob}
\olpseganchor{OLP-0047-B017}{end}

\olpseganchor{OLP-0047-B018}{start}
لکه مخکښې، دا ښودل چې جوړ شوے~\LR{$\Int$} يوه مرتبه حلقه ده، اوږد
خو معمولي کار دے۔ دا تاسو ته پرېږدو۔
\olpseganchor{OLP-0047-B018}{end}

\olpseganchor{OLP-0047-B019}{start}
په دې سره د صحيح عددونو کار بشپړېږي۔ خو اوس د ناطقو عددونو په اړه
ورته خبرې ښودل غواړو۔ په ځانګړي ډول بايد وښيو چې زموږ د \LR{$+$}،
\LR{$\times$} او \LR{$\leq$} د ورکړل شوو تعريفونو له مخې ناطق عددونه يو
مرتب \emph{ميدان} جوړوي:
\begin{defn}\ollabel{orderedfield}
يو \emph{مرتب ميدان} هغه مرتبه حلقه ده چې \LR{$0\ne 1$} پکښې
وي او دا شرط هم پوره کوي:
\begin{align*}
	\emph{ضربي معکوس}& & (\forall a \in S \setminus \{0\})(\exists b \in S) a\times b& = 1
\end{align*}
\end{defn}
\olpseganchor{OLP-0047-B019}{end}

\olpseganchor{OLP-0047-B020}{start}
چې مو وښودل \LR{$\Int$} يوه مرتبه حلقه جوړوي، بيا دا ښودل اسانه خو
اوږد دي چې \LR{$\Rat$} يو مرتب ميدان جوړوي۔
\olpseganchor{OLP-0047-B020}{end}

\olpseganchor{OLP-0047-B021}{start}
\begin{prob}
ثابت کړئ چې \LR{$\Rat$} يو مرتب ميدان دے۔
\end{prob}
\olpseganchor{OLP-0047-B021}{end}

\olpseganchor{OLP-0047-B022}{start}
د صحيح او ناطقو عددونو له سمبالولو وروسته يوازې حقيقي عددونه پاتې
دي۔ په ځانګړي ډول بايد وښيو چې \LR{$\Real$} يو \emph{بشپړ} مرتب ميدان
جوړوي، يعنې يو مرتب ميدان چې د بشپړتيا خاصيت لري۔
\olref[cuts]{realcompleteness} ثابته کړه چې \LR{$\Real$} د بشپړتيا خاصيت
لري۔ خو لا د دې ارزولو (ستړے) کار پاتې دے چې \LR{$\Real$} يو مرتب ميدان
دے۔
\olpseganchor{OLP-0047-B022}{end}

\olpseganchor{OLP-0047-B023}{start}
مخکښې له دې چې \emph{هغه} اوږد تمرين پيل کړو، بايد ځينې نورې
«سمدستي» خبرې وارزوو۔ د بېلګې په ډول بايد ډاډ ولرو چې د هر دوو
پرېکونونو \LR{$\alpha$} او \LR{$\beta$} دپاره، تعريف شوے \LR{$\alpha + \beta$}
په رښتيا يو \emph{پرېکون} دے۔ د دې حقيقت ثبوت دا دے:
\olpseganchor{OLP-0047-B023}{end}

\olpseganchor{OLP-0047-B024}{start}
\begin{proof}	
ځکه چې \LR{$\alpha$} او \LR{$\beta$} دواړه پرېکونونه دي،
\LR{$\alpha + \beta = \Setabs{p + q}{p \in \alpha \land q \in \beta}$}
د \LR{$\Rat$} يو ناتش خاص فرعي سټ دے۔ د دواړو پرېکونونو غړي جمع
کړئ څو يو غړے پيدا شي۔ د هر پرېکون له بېرون څخه يو ناطق عدد
واخلئ؛ د دغو دوو عددونو جمع د ټولو ممکنو جمعو پاسنے حد دے او
پخپله په پرېکون کښې نۀ ده۔ اوس فرض کړئ د ځينو
\LR{$p \in \alpha$} او \LR{$q \in \beta$} دپاره \LR{$x < p + q$}۔ بيا
\LR{$x - p < q$}، نو \LR{$x - p \in \beta$}، او
\LR{$x = p + (x - p) \in \alpha + \beta$}۔ نو \LR{$\alpha + \beta$} د
\LR{$\Rat$} يوه پيلنۍ برخه ده۔ په پاے کښې، د هر
\LR{$p + q \in \alpha + \beta$} دپاره، ځکه چې \LR{$\alpha$} او \LR{$\beta$}
دواړه پرېکونونه دي، داسې \LR{$p_1 \in \alpha$} او \LR{$q_1 \in \beta$} شته
چې \LR{$p < p_1$} او \LR{$q < q_1$}؛ نو
\LR{$p + q < p_1 + q_1 \in \alpha + \beta$}؛ نو
\LR{$\alpha + \beta$} تر ټولو لوے غړے نۀ لري۔
\end{proof}
\olpseganchor{OLP-0047-B024}{end}

\olpseganchor{OLP-0047-B025}{start}
ورته هڅې به درته اجازه وکړي چې وارزوئ \LR{$\alpha - \beta$}،
\LR{$\alpha \times \beta$} او \LR{$\alpha \div \beta$} پرېکونونه دي (په
وروستي حالت کښې هغه حالت نۀ شمېرو چې \LR{$\beta$} صفر پرېکون وي)۔ خو
د دې د تفصيل يو څو ګامونه دلته ورکوو۔ د يوه پرېکون د نفي هر
غړے د هغه له بېرون څخه د يوه ناطق عدد تر منفي کم دے۔ نو نفي
پيلنۍ برخه ده او د ناطقو عددونو د ګڼوالي له مخې تر ټولو لوے غړے
نۀ لري۔ د بېرون يو عدد نفي ناتشه کوي، او د اصلي پرېکون د يوه غړي
منفي ئې پاسنے حد دے چې په نفي کښې نۀ دے۔ ځکه تفريق د دوو پرېکونونو
جمع ده۔ په مثبت ضرب کښې د هر پرېکون له بېرون څخه مثبت ناطق عدد
واخلئ۔ د دواړو عددونو ضرب داسې پاسنے حد دے چې په ضرب کښې نۀ
دے۔ هر منفي ناطق عدد د صفر له پرېکونه راځي۔ هر نامنفي عدد چې
تر يو مثبت حاصل‌ضرب کم وي، د يوه مثبت عامل په وېشلو سره د بل
عامل تر ارزښت کمېږي، نو په هماغه پرېکون کښې وي۔ او هر مثبت
حاصل‌ضرب د يوه عامل په زياتولو سره زياتېږي۔ که دواړه پرېکونونه
مثبت وي، مثبت غړي لري، نو صفر هم وروستے غړے نۀ دے۔
د مثبت پرېکون په ضربي معکوس کښې منفي عددونه او صفر شته۔ هر مثبت
بهرنے شاهد يو مثبت پاسنے حد ورکوي چې تر هغه لاندې ناطق عددونه
د هماغه شاهد له مخې منل کېږي۔ د اصلي پرېکون د هر مثبت غړي
معکوس د دغو ټولو حدونو پاسنے حد دے، ځکه هر بهرنے شاهد تر دغه غړي
لوے دے۔ نو ضربي معکوس ناتش، خاص، پيلنے او بې وروستي غړي دے۔
د منفي پرېکون ضربي معکوس له نفي، او وېش له ضرب څخه ترلاسه کېږي۔
په تمرين کښې د ميدان د قوانينو د ثابتولو لپاره دا طريقه هم
ګټوره ده: د پرېکون له دننه او بېرون څخه ناطق عددونه په هر
غوښتل شوي مثبت واټن کښې موندل کېږي۔ له يوه دننني عدد څخه تر يوه
بهرني عدده د هماغه واټن نه په وړو برابر ناطق ګامونو ولاړ شئ
او لومړنے بهرنے ګام واخلئ۔ دا دنننے او بهرنے جوړه د جمعي او
ضربي معکوس د قوانينو د ثابتولو لپاره وکاروئ۔
\olpseganchor{OLP-0047-B025}{end}

\olpseganchor{OLP-0047-B026}{start}
\begin{prob}
ثابت کړئ چې \LR{$\Real$} يو مرتب ميدان دے۔
\end{prob}
\olpseganchor{OLP-0047-B026}{end}

\olpseganchor{OLP-0047-B027}{start}
خو يوه وړه نيمګړې خبره لا سمول غواړي۔ په \olref[cuts]{sec} کښې
مو وړانديز وکړ چې \LR{$\sqrt{2} =
\Setabs{p \in \Rat}{p < 0 \text{\RL{ يا }}p^2 < 2}$} اخيستلے شو۔ خو
بايد وښيو چې دا سټ يو \emph{پرېکون} دے۔ د دې حقيقت ثبوت دا دے:
\olpseganchor{OLP-0047-B027}{end}

\olpseganchor{OLP-0047-B028}{start}
\begin{proof}
په څرګند ډول دا د ناطقو عددونو يوه ناتشه خاصه پيلنۍ برخه ده؛ نو
دومره ښودل بس دي چې تر ټولو لوے غړے نۀ لري۔ په ځانګړي ډول دا ښودل
بس دي چې که \LR{$p$} يو مثبت ناطق عدد وي، \LR{$p^2 < 2$} وي او
\LR{$q = \frac{2p+2}{p+2}$} وي، نو \LR{$p < q$} او \LR{$q^2 < 2$} دواړه دي۔
د \LR{$p < q$} د ليدلو دپاره يوازې پام وکړئ:
\begin{align*}
	p^2 &< 2\\
	p^2 + 2p &< 2 + 2p\\
	p(p + 2) &< 2 + 2p\\
	p &< \tfrac{2+2p}{p+2} = q
\end{align*}
د \LR{$q^2 < 2$} د ليدلو دپاره يوازې پام وکړئ:
\begin{align*}
	p^2 &< 2\\
	2p^2 + 4p + 2 &< p^2 + 4p+ 4\\
	4p^2 + 8p + 4 &< 2(p^2 + 4p + 4)\\
	(2p+2)^2 & <2(p+2)^2\\
	\tfrac{(2p+2)^2}{(p+2)^2} &< 2\\
	q^2 &< 2
\end{align*}
\end{proof}
\olpseganchor{OLP-0047-B028}{end}

\olpseganchor{OLP-0048-B004}{start}
\olfileid{sfr}{arith}{cauchy}
\olpseganchor{OLP-0048-B004}{end}

\olpseganchor{OLP-0048-B005}{start}
\olsection{ضميمه: حقيقي عددونه د کوشي د لړيو په توګه}
\olpseganchor{OLP-0048-B005}{end}

\olpseganchor{OLP-0048-B006}{start}
په \olref[cuts]{sec} کښې مو حقيقي عددونه د ډېډېکېنډ د پرېکونونو په
توګه جوړ کړل۔ په دې برخه کښې يوه بله جوړونه تشريح کوو۔ دا د کوشي پر
هغه تعريف ولاړه ده چې اوس ئې د کوشي لړۍ بولو؛ خو له دې تعريف څخه د
حقيقي عددونو د \emph{جوړولو} کار د نولسمې پېړۍ نورو ليکوالانو، په
ځانګړې توګه وايرشټراس، هاينې، مېراے او کانټور، کړے دے۔ (د ښۀ تاريخ
دپاره \mbox{\LR{\latinfont \citeauthor{OConnorRobertson:RN}}} \mbox{\LR{\latinfont \citeyear{OConnorRobertson:RN}}}
وګورئ۔)
\olpseganchor{OLP-0048-B006}{end}

\olpseganchor{OLP-0048-B007}{start}
مخکښې له دې چې نولسمې پېړۍ ته ورسېږو، د سايمن سټېوين (1548--1620)
خبره د پام وړ ده۔ په لنډ ډول، سټېوين پوه شو چې هر حقيقي عدد د هغه د
اعشاري پراختيا له مخې ګڼلے شو۔ نو آن يو غير ناطق عدد لکه \LR{$\sqrt{2}$}
هم يوه ښۀ اعشاري پراختيا لري، چې پيل ئې داسې دے:
\[
	1.41421356237\ldots
\]
په سټ نظريه کښې د اعشاري پراختياوو نمونه جوړول ډېر اسان دي: يوازې
داسې تابعې \LR{$d \colon \Nat \to \Nat$} ئې وګڼئ چې \LR{$d(n)$} هغه \LR{$n$}م
اعشاري رقم وي چې موږ ئې څېړو۔ د هرې خانې ارزښت له صفر څخه تر
نهو پورې يو طبيعي عدد دے۔ بيا لږ بدلون ته اړتيا لرو، څو د حقيقي
عدد هغه برخه هم سمبال کړو چې له اعشاري ټکي مخکښې راځي (دلته يوازې
\LR{$1$})۔ يو بل بدلون، يعنې د هم ارزښتۍ اړيکه، هم پکار ده څو د بېلګې په
توګه يقيني کړو چې \LR{$0.999\ldots = 1$}۔ خو په سټ نظريه کښې د سټېوين په
لار د حقيقي عددونو بشپړه دقيقه جوړونه وړاندې کول ګران نۀ دي۔
\olpseganchor{OLP-0048-B007}{end}

\olpseganchor{OLP-0048-B008}{start}
سټېوين زموږ اصلي موضوع نۀ ده۔ (د سټېوين په اړه نور معلومات په
\mbox{\LR{\latinfont \citealt{KatzKatz2012}}} کښې وګورئ۔) خو دا ډېره ورته خبره وګورئ۔ د
\LR{$\sqrt{2}$} اعشاري پراختيا په نېغه د څېړلو پر ځاے، د \LR{$\sqrt{2}$} د
لا کره پراختياوو په اخيستلو د هغو ناطقو اټکلونو يوه \emph{لړۍ}
کتلے شو چې کره‌والے ئې پرله‌پسې زياتېږي:
\[
1, 1.4, 1.414, 1.4142, 1.41421,\ldots
\]
دا نظر چې حقيقي عددونه د «پرله‌پسې ښو اټکلونو» له لارې ګڼلے شو، د
يوې بلې پوهېدنې بنسټ راکوي (هغو پوهېدنو ته ورته چې \LR{$\Rat$} مو له
\LR{$\Int$}، يا \LR{$\Int$} مو له \LR{$\Nat$} څخه جوړول)۔ بنسټيزې خبرې دا دي:
\begin{enumerate}
	\item هر حقيقي عدد د يوې (کېدے شي نامتناهي) اعشاري پراختيا په
	بڼه ليکل کېدے شي۔
	\item په يوه (کېدے شي نامتناهي) اعشاري پراختيا کښې زېرمه شوي
	معلومات په هماغه ډول د ناطقو عددونو په يوه لړۍ کښې زېرمه کېدے شي۔
	\item د ناطقو عددونو لړۍ له \LR{$\Nat$} څخه \LR{$\Rat$} ته يوه تابع ګڼلے
	شو؛ يوازې \LR{$f(n)$} په لړۍ کښې \LR{$n$}م ناطق عدد وټاکئ۔
\end{enumerate}
البته له \LR{$\Nat$} څخه \LR{$\Rat$} ته هره \emph{تابع} حقيقي عدد نۀ راکوي۔
د بېلګې په توګه دا تابع وګورئ:
\[
	f(n) =\begin{cases}
			1 & \text{\RL{که }}n\text{\RL{ طاق وي}}\\
			0 &\text{\RL{که }}n\text{\RL{ جفت وي}}
		\end{cases}
\]
اصلي اندېښنه دا ده چې د \LR{$0,1,0,1,0,1,0,\ldots$} لړۍ هېڅ حقيقي
عدد ته «نژدې کېدونکې» نۀ ښکاري۔ نو د دې د يقيني کولو دپاره چې هغه
لړۍ څېړو چې کوم حقيقي عدد ته ورنژدې کېږي، بايد خپل پام هغو لړيو ته
محدود کړو چې کوم \emph{حد} لري۔
\olpseganchor{OLP-0048-B008}{end}

\olpseganchor{OLP-0048-B009}{start}
موږ د حد له نظر سره په \olref[his][set][limits]{sec} کښې مخکښې هم
مخامخ شوي يو۔ خو کټ مټ هماغه پخوانے تعريف نشو کارولے۔ هلته عبارت
«\LR{$(\forall \epsilon>0)$}» په پټه پر حقيقي عددونو کميت‌ايښودنه
درلوده؛ او موږ د حقيقي عددونو پر تابعو حدونه څېړل؛ نو د هغه تعريف
کارول به دا وي چې حقيقي عددونه لا له مخکښې ومنو، حال دا چې موږ کټ
مټ د هغو \emph{جوړول} غواړو۔ نېکمرغي ده چې د حد له يوه ډېر نژدې
نظر سره کار کولے شو۔
\begin{defn}\ollabel{def:CauchySequence}
  يوه تابع \LR{$f: \Nat \to \Rat$} هله او يوازې هله د \emph{کوشي لړۍ}
  ده چې د هر مثبت \LR{$\epsilon \in \Rat$} دپاره
  \LR{$(\exists \ell \in \Nat)(\forall m, n > \ell)|f(m) - f(n)| <
  \epsilon$} وي۔
\end{defn}
\olpseganchor{OLP-0048-B009}{end}

\olpseganchor{OLP-0048-B010}{start}
د حد عمومي نظر د پخوا په شان دے: که د کره‌والي يوه ټاکلې کچه
(چې په~\LR{$\epsilon$} اندازه کېږي) وغواړئ، د کتلو يوه «سيمه» شته
(هره ورودي چې له~\LR{$\ell$} څخه لويه وي)۔ دا هم په اسانۍ ليدل کېږي چې
زموږ لړۍ \LR{$1$}، \LR{$1.4$}، \LR{$1.414$}، \LR{$1.4142$}، \LR{$1.41421$}\ldots حد لري:
که غواړئ \LR{$\sqrt{2}$} د \LR{$\nicefrac{1}{10^{n}}$} له خطا دننه اټکل
کړئ، نو يوازې له \LR{$n$}م غړي وروسته هر غړے وګورئ۔
\olpseganchor{OLP-0048-B010}{end}

\olpseganchor{OLP-0048-B011}{start}
نو څرګند فکر به دا وي چې حقيقي عدد هماغه هره د کوشي لړۍ \emph{ده}۔
خو لکه د \LR{$\Int$} او \LR{$\Rat$} په جوړونو کښې، دا خبره به بېخي ساده وي:
د هر ورکړل شوي حقيقي عدد دپاره څو بېلې د کوشي لړۍ هماغه حقيقي عدد
ښيي۔ د دې د ليدلو يوه اسانه لار داسې ده۔ يوه د کوشي لړۍ~\LR{$f$} واخلئ
او \LR{$g$} کټ مټ د~\LR{$f$} په شان تابع تعريف کړئ، يوازې دومره توپير چې
\LR{$g(0)\neq f(0)$}۔ ځکه دواړه لړۍ له لومړي عدد وروسته په هر ځاے کښې
سره برابرې دي، نو په پاے کښې به ووايو چې د
\olref{def:CauchySequence} په معنا يو حد لري، او بايد د هماغه حقيقي
عدد «تعريفوونکې» وګڼل شي۔ نو په رښتيا بايد دا د کوشي لړۍ هماغه يو
حقيقي عدد وګڼو۔
\olpseganchor{OLP-0048-B011}{end}

\olpseganchor{OLP-0048-B012}{start}
په پايله کښې بيا بايد د کوشي پر لړيو د هم ارزښتۍ يوه اړيکه تعريف
کړو او حقيقي عددونه د هم ارزښتۍ له ټولګيو سره يو وګڼو۔
لومړے داسې تابع پکار ده چې په حد کښې \LR{$0$} ته ورنژدې شي۔
د هرې تابع \LR{$h : \Nat \to \Rat$} دپاره ووايئ چې \emph{\LR{$h$} \LR{$0$} ته
ورنژدې کېږي} هله او يوازې هله چې د هر مثبت \LR{$\epsilon \in \Rat$}
دپاره \LR{$(\exists \ell \in \Nat)(\forall n > \ell)|h(n)| <
\epsilon$} وي۔\footnote{دا د \LR{$\lim_{x
\mathord{\rightarrow}\infty}f(x) = 0$} له تعريف سره پرتله کړئ چې په
\olref[his][set][limits]{sec} کښې راغلے۔} برسېره پر دې، چې \LR{$f$} او
\LR{$g$} د \LR{$\Nat \to \Rat$} تابعې وي، \LR{$(f-g)(n) = f(n) - g(n)$}
وټاکئ۔ اوس تعريف کړئ:
\[
	f \Realequiv g \text{\RL{ هله او يوازې هله چې \LR{$(f-g)$} \LR{$0$} ته ورنژدې کېږي}}.
\]
\LR{$\Realequiv$} د هم ارزښتۍ اړيکه ده۔ د هرې لړۍ له ځان څخه
توپير صفر دے او د توپير لوری د مطلق ارزښت په حساب اغېز نۀ لري۔
د انتقالي خاصيت لپاره هر يو له دوو توپيرونو په پوره وروستۍ برخه
کښې د غوښتل شوي زغم تر نيمې کم واخلئ؛ مثلثي نابرابري د دواړو
ګډ توپير تر بشپړ زغم کموي۔
بيا، که وغواړو، حقيقي عددونه د \LR{$\Nat \to \Rat$} د ټولو کوشي
لړيو هغه د هم ارزښتۍ ټولګي تعريفولے شو چې د \LR{$\Realequiv$} له مخې
جوړېږي۔
\olpseganchor{OLP-0048-B012}{end}

\olpseganchor{OLP-0048-B013}{start}
\begin{prob}
د هر \LR{$n$} دپاره \LR{$f(n) = 0$} وټاکئ۔ او \LR{$g(n) = \frac{1}{(n+1)^2}$}
وټاکئ۔ وښيئ چې دواړه د کوشي لړۍ دي او په رښتيا د دواړو تابعو حد
\LR{$0$} دے، نو \LR{$f \Realequiv g$} هم دے۔
\end{prob}
\olpseganchor{OLP-0048-B013}{end}

\olpseganchor{OLP-0048-B014}{start}
له دې وروسته، د معمول په شان، \LR{$\equivrep{f}{\Realequiv}$} د هغه د
هم ارزښتۍ ټولګي دپاره ليکو چې \LR{$f$} ئې يو غړے دے۔ خو د لوست
د اسانولو دپاره به په راتلونکو برخو کښې لاندې‌ليک پرېږدو او يوازې
\LR{$\equivrep{f}{}$} ليکو۔ دا هم ټاکو چې د هر \LR{$q \in \Rat$} دپاره
\LR{$q_{\Real} = \equivrep{c_{q}}{}$} وي، چې \LR{$c_{q}$} ثابته تابع ده او
د هر \LR{$n \in \Nat$} دپاره \LR{$c_q(n) = q$} وي۔ بيا پر حقيقي عددونو
بنسټيزې اړيکې او عمليې داسې تعريفوو، د بېلګې په توګه:
\begin{align*}
	\equivrep{f}{} + \equivrep{g}{} &= 	\equivrep{(f + g)}{} \\
	\equivrep{f}{} \times 	\equivrep{g}{} &= \equivrep{(f \times g)}{}
\end{align*}
چې \LR{$(f + g)(n) = f(n) + g(n)$} او \LR{$(f \times g)(n) = f(n) \times
g(n)$}۔ البته بايد دا هم وارزوو چې که \LR{$f$} او \LR{$g$} د کوشي لړۍ وي، نو
هره يوه \LR{$(f + g)$}، \LR{$(f-g)$} او \LR{$(f\times g)$} هم د کوشي لړۍ ده۔
د جمع او تفريق لپاره د مثلث نابرابري وکاروئ۔ د ضرب لپاره پام
وکړئ چې هره د کوشي لړۍ په پوره وروستۍ برخه کښې محدوده وي؛ ځکه
هلته هر غړے له يوه ټاکلي وروستي غړي سره تر يوه واحد کم توپير لري۔
د دوو حاصل‌ضربونو توپير د لومړي عامل د بدلون او د دويم عامل د
بدلون په دوو حدو ووېشئ او پاتې عاملونه په يو ثابت حد محدود کړئ۔
همدا اټکل ښيي چې د استازو په داسې لړيو بدلول چې توپيرونه يې صفر
ته ورنژدې کېږي، د ضرب ټولګے نۀ بدلوي۔ د جمع او تفريق لپاره بيا
مثلثي نابرابري بس ده۔
\olpseganchor{OLP-0048-B014}{end}

\olpseganchor{OLP-0048-B015}{start}
په پاے کښې د مرتبې يو مفهوم تعريفوو۔ ووايئ \LR{$\equivrep{f}{}$}
\emph{مثبت} دے هله او يوازې هله چې هم \LR{$\equivrep{f}{}\neq 0_\Real$}
وي او هم \LR{$(\exists \ell \in \Nat)(\forall n > \ell)0 < f(n)$} وي۔
دا دوه شرطونه په ګډه له دې سره همارزښته دي چې په پوره
وروستۍ برخه کښې د لړۍ غړي له يوه ثابت مثبت ناطق عدد څخه لوے وي۔
د ناصفر ټولګي لپاره داسې مثبت زغم شته چې په هره وروستۍ برخه کښې
يو غړے لږ تر لږه همدومره مطلق ارزښت لري۔ د کوشي خاصيت له مخې په
پوره وروستۍ برخه کښې د غړو تر منځ واټن تر دې زغم نيم کم دے۔
نو د يوه داسې غړي نښه د ټولې وروستۍ برخې نښه ټاکي او مطلق
ارزښت يې له نيم زغم څخه لوے پاتې کېږي۔ دا شرط د استازي له بدلېدو
سره ساتل کېږي، ځکه د دوو استازو توپير په پاى کښې له دې مثبت
حده هر څومره کمېږي۔ بيا ووايئ
\LR{$\equivrep{f}{} < \equivrep{g}{}$} هله او يوازې هله چې
\LR{$\equivrep{(g - f)}{}$} مثبت وي۔ بايد وارزوو چې دا سم تعريف شوے
دے (يعنې د هم ارزښتۍ له ټولګي څخه د «نماينده» تابعې پر ټاکنه پورې
اړه نۀ لري)۔
%\begin{align*}
%	\equivrep{f}{} \leq \equivrep{g}{} \text{ iff }&\equivrep{f}{} = \equivrep{g}{} \text{ or }\\
%	&(\exists \ell \in \Nat)(\forall n \in \Nat)(\ell < n \lif f(n) \leq g(n))
%\end{align*}
خو له دې ارزونې وروسته په ډېرې اسانۍ ښودل کېدے شي چې دا تعريفونه
سم الجبري خاصيتونه راکوي؛ يعنې:
\begin{thm}\ollabel{thm:cauchyorderedfield}
د کوشي د لړيو د هم ارزښتۍ ټولګي يو مرتب ميدان جوړوي۔
\end{thm}
\olpseganchor{OLP-0048-B015}{end}


\olpseganchor{OLP-0048-B016}{start}
\begin{proof}
د تمرين لپاره اشاره: د حلقې قوانين د ناطقو عددونو د هرې خانې له
قانونونو راځي۔ د مثبتو ټولګيو جمع او ضرب بيا مثبت دي، ځکه د
هرې وروستۍ برخې لپاره ثابت مثبت لانديني حدونه لرو۔ هر ناصفر
ټولګے پخپله يا يې جمعي معکوس مثبت دے، د هماغه وروستي ثابت نښان
له مخې۔ د ضربي معکوس لپاره د ناصفر ټولګي د استازي هر وروستے
غړے سرچپه کړئ؛ مخکني غړي هر ناطق ارزښت اخيستلے شي۔ په وروستۍ
برخه کښې مطلق ارزښت له ثابت مثبت حد څخه لوے دے، نو د دوو سرچپه
غړو واټن د اصلي غړو تر واټن د هماغه حد د مربع په وېشلو زيات
نۀ وي۔ نو نوې لړۍ د کوشي ده او له استازي سره يې ضرب په وروستۍ
برخه کښې يو دے۔
\end{proof}
\olpseganchor{OLP-0048-B016}{end}

\olpseganchor{OLP-0048-B017}{start}
\begin{prob}
ثابته کړئ چې د کوشي د لړيو د هم ارزښتۍ ټولګي يو مرتب ميدان جوړوي۔
\end{prob}
\olpseganchor{OLP-0048-B017}{end}

\olpseganchor{OLP-0048-B018}{start}
دا ثابتول لږ ګران دي چې په دې ډول جوړ شوي حقيقي عددونه د بشپړتيا
خاصيت لري، نو ثبوت ئې وړاندې کوو۔
\olpseganchor{OLP-0048-B018}{end}

\olpseganchor{OLP-0048-B019}{start}
\begin{thm}
د کوشي د لړيو د هم ارزښتۍ د ټولګيو هر ناتش سټ چې پاسنی حد لري،
تر ټولو لږ پاسنی حد هم لري۔
\end{thm}
\olpseganchor{OLP-0048-B019}{end}

\olpseganchor{OLP-0048-B020}{start}
\begin{proof}[د ثبوت خاکه] فرض کړئ \LR{$S$} د کوشي د لړيو د هم ارزښتۍ
د ټولګيو هر داسې ناتش سټ دے چې پاسنی حد لري۔ د هرې کوشي لړۍ د
کافي وروستيو غړو واټن له يوه ټاکلي وروستي غړي څخه تر يوه واحد
کم دے۔ د هماغه غړي څخه دوه واحده ليرې دواړه ناطق سرونه واخلئ۔
د ثابتې مثبتې فاصلې له امله، د لړۍ ټولګے د دغو دوو ناطقو
عددونو د ثابتو ټولګيو تر منځ په کلکه پروت دے۔ د راکړل شوي
پاسني حد لپاره پورته سر واخلئ؛ نو داسې \LR{$p \in \Rat$} شته چې \LR{$p_{\Real}$} د
\LR{$S$} پاسنی حد دے۔ يو \LR{$r \in S$} واخلئ او د هغه لپاره لاندې سر
واخلئ؛ بيا داسې \LR{$q \in \Rat$} شته
چې \LR{$q_{\Real} < r$} وي۔ نو که د \LR{$S$} تر ټولو لږ پاسنی حد موجود وي،
د \LR{$q_\Real$} او \LR{$p_\Real$} تر منځ (دواړه سرونه پکې شامل) دے۔
\olpseganchor{OLP-0048-B020}{end}


\olpseganchor{OLP-0048-B021}{start}
موږ به تر ټولو لږ پاسني حد ته په يوه وخت له لاندې او پورته څخه
ورنژدې شو۔ په ځانګړې توګه دوه تابعې \LR{$f, g \colon \Nat \to \Rat$}
داسې تعريفوو چې موخه مو دا وي چې \LR{$f$} له پورته او \LR{$g$} له لاندې څخه
تر ټولو لږ پاسني حد ته ورنژدې شي۔ په دې تعريف پيل کوو:
\begin{align*}
	f(0) &= p \\
	g(0) &= q
\end{align*}
بيا، چې \LR{$a_n = \frac{f(n) + g(n)}{2}$} وي، وټاکئ:\footnote{دا يو بازګشتي تعريف دے؛ د دې ډول تعريفونو د شتون
دليل په \olref[sth][spine][recursion]{sec} کښې ثابتېږي۔}
\begin{align*}
	f(n+1) &=
	\begin{cases}
		a_n &\text{\RL{که }}(\forall h \in S)h \leq (a_n)_\Real\\
		f(n)&\text{\RL{که نۀ}}
	\end{cases}\\
	g(n+1) &=
	\begin{cases}
		a_n &\text{\RL{که }}(\exists h \in S)h > (a_n)_\Real\\
	 	g(n) &\text{\RL{که نۀ}}
	\end{cases}
\end{align*}
په هر ګام کښې له دوو متقابلو شرطونو څخه يوازې يو صدق کوي۔
نو پورته سر پاسنے حد پاتې کېږي او د لاندې سر څخه تل د سټ يو
غړے لوے وي۔ وروستي واټنونه يو په بل کښې دننه دي، او په هر ګام
کښې د واټن اوږدوالے نيمېږي۔ ځکه د هر ټاکلي مثبت ناطق زغم
لپاره يو واټن شته چې اوږدوالے يې تر هغه کم وي؛ د دواړو لړيو ټول
وروستي غړي په هماغه واټن کښې وي۔ نو \LR{$f$} او \LR{$g$} دواړه د کوشي لړۍ دي۔
تابع \LR{$(f-g)$} \LR{$0$} ته ورنژدې کېږي، ځکه د \LR{$f$} او \LR{$g$} تر منځ توپير
په هر ګام کښې نيم کېږي۔ نو \LR{$\equivrep{f}{} = \equivrep{g}{}$}۔
\olpseganchor{OLP-0048-B021}{end}

\olpseganchor{OLP-0048-B022}{start}
لومړے ښيو چې \LR{$\equivrep{f}{}$} د \LR{$S$} پاسنی حد دے، يعنې
\LR{$(\forall h \in S)h \leq \equivrep{f}{}$}۔
(په لار کښې به \olref{thm:cauchyorderedfield} کاروو۔) يو استازے
واخلئ چې \LR{$\equivrep{h}{} \in S$} وي، او د نقيض لپاره فرض کړئ چې
\LR{$\equivrep{f}{} < \equivrep{h}{}$}، نو
\LR{$0_\Real < \equivrep{(h-f)}{}$}۔ ځکه \LR{$f$} يوه نۀ زياتېدونکې د
کوشي لړۍ ده، د هغې د وروستيو غړو واټنونه هر څومره کمېدے
شي۔ د مثبت ښي ټولګي د ثابت لانديني حد تر نيمې لږ واټن
واخلئ۔ بيا داسې \LR{$n \in \Nat$} شته چې
\LR{$\equivrep{(c_{f(n)} - f)}{} < \equivrep{(h-f)}{}$}۔ نو:
\[
	(f(n))_\Real = \equivrep{c_{f(n)}}{} < \equivrep{f}{} + \equivrep{(h-f)}{} = \equivrep{h}{},
\]
خو دا له دې جوړښتي حقيقت سره په ټکر کښې دے چې
\LR{$\equivrep{h}{} \leq (f(n))_\Real$}۔
\olpseganchor{OLP-0048-B022}{end}

\olpseganchor{OLP-0048-B023}{start}
اوس ښيو چې \LR{$\equivrep{f}{} = \equivrep{g}{}$} د \LR{$S$} \emph{تر ټولو
لږ} پاسنی حد دے۔ نو \LR{$j$} دې هره د کوشي لړۍ وي او فرض کړئ
\LR{$\equivrep{j}{} < \equivrep{g}{}$}۔ لکه پورته (د دې حقيقت په کارولو
چې \LR{$g$} \emph{نۀ کمېدونکې} ده) د هغې د وروستيو غړو
واټن د مثبت توپير له ثابت لانديني حده کم وټاکئ۔ داسې \LR{$n \in \Nat$}
شته چې \LR{$\equivrep{j}{} < (g(n))_\Real$}۔ خو د جوړونې له مخې
داسې \LR{$h \in S$} شته چې \LR{$(g(n))_\Real < h$}، نو
\LR{$\equivrep{j}{} < h$} او له دې امله
\LR{$\equivrep{j}{}$} د \LR{$S$} پاسنی حد نۀ دے۔
\end{proof}
\olpseganchor{OLP-0048-B023}{end}

\olpseganchor{OLP-0049-B004}{start}
\olchapter{sfr}{infinite}{نامتناهي سټونه}
\olpseganchor{OLP-0049-B004}{end}

\olpseganchor{OLP-0049-B005}{start}
\begin{editorial}
د نامتناهي سټونو دا څپرکے د ټيم بټن له \emph{اوپن سېټ تيوري}
څخه اخيستل شوے دے۔
\end{editorial}
\olpseganchor{OLP-0049-B005}{end}

\olpseganchor{OLP-0050-B003}{start}
\olfileid{sfr}{infinite}{hilbert}
\olsection{د هيلبرټ هوټل}
\olpseganchor{OLP-0050-B003}{end}

\olpseganchor{OLP-0050-B004}{start}
د طبيعي عددونو سټ په څرګند ډول نامتناهي دے۔ نو که نۀ غواړو طبيعي
عددونه له مخکښې \emph{ومنـو}، زموږ لومړے ګام بايد دا وي چې نامتناهي
سټ په داسې اصطلاحاتو کښې وپېژنو چې خپله د طبيعي عددونو يادولو ته
اړتيا ونۀ لري۔ هيلبرټ دغه توپير د دوو هوټلونو په بېلګه څرګندوي۔
\begin{explain}
په يوه متناهي هوټل کښې که په هره کوټه کښې يو مېلمه وي،
د مېلمنو بيا وېش چې د هرې کوټې د يوه کس حد پوره کوي،
ټولې کوټې بېرته ډکوي۔ د مېلمنو شمېر د کوټو له شمېر سره
برابر دے، نو د نوي مېلمه دپاره خالي کوټه نۀ پاتې کېږي۔
\olpseganchor{OLP-0050-B004}{end}

\olpseganchor{OLP-0050-B005}{start}
اوس داسې کوټې وګڼئ چې شمېرې يې \LR{$1$}، \LR{$2$}، \LR{$3$}،
\LR{$4$}، \LR{$5$}، \dots دي او هره يوه خپل مېلمه لري۔
ټول مېلمانه يو ګام وړاندې لېږو۔ د هرې کوټې مېلمه د خپلې
کوټې ورپسې کوټې ته ځي۔ نوي ځايونه ټول بېل دي او لومړۍ
کوټه~\LR{$1$} نوي مېلمه ته پاتې کېږي۔
	\begin{center}
		\begin{tikzpicture}[scale = .75]
		\foreach \x in {1, 2, 3, 4, 5, 6, 7, 8, 9}
		{
			\node (\x a) at (\x, 1) {\small{\x}};
			\node (\x b) at (\x, 2) {\small{\x}};
		}
		\node (dotsa) at (10, 1) {\small{\ldots}};
		\node (dotsb) at (10, 2) {\small{\ldots}};
		\draw[->] (1b)--(2a);
		\draw[->] (2b)--(3a);
		\draw[->] (3b)--(4a);
		\draw[->] (4b)--(5a);
		\draw[->] (5b)--(6a);
		\draw[->] (6b)--(7a);
		\draw[->] (7b)--(8a);
		\draw[->] (8b)--(9a);
		\draw[->] (9b)--(dotsa);
		\draw (1,1) circle (.4);
		\end{tikzpicture}
	\end{center}
	(د هيلبرټ د ۱۹۲۴–۲۵ درسونو بېلګه، \mbox{\LR{\latinfont \citealt[730]{EwaldSieg2013}}}۔)
\end{explain}
مهم ټکی دا دے چې د هيلبرټ هوټل نامتناهي ډېرې کوټې لري؛ او د هغه
تشريح د دې خبرې تعريف ګڼلے شو چې دا څه معنا لري۔ په رښتيا، د
ډېډېکېنډ لار همدا وه (البته دلته په لوے زماني بې‌ځايه‌والي سره
وړاندې شوې؛ د ډېډېکېنډ تعريف د \mbox{\LR{\latinfont \citeyear{Dedekind1888}}} کال دے):
\olpseganchor{OLP-0050-B005}{end}

\olpseganchor{OLP-0050-B006}{start}
\begin{defn}\ollabel{defn:DedekindInfinite}
يو سټ \LR{$A$} هله او يوازې هله \emph{د ډېډېکېنډ له مخې نامتناهي} دے
چې له~\LR{$A$} څخه د~\LR{$A$} يو حقيقي فرعي سټ ته کومه يو پر يو تابع موجوده
وي۔ يعنې داسې \LR{$o \in A$} او يو پر يو تابع \LR{$f \colon A \to A$} شته
چې \LR{$o \notin \ran{f}$} وي۔
\end{defn}
\olpseganchor{OLP-0050-B006}{end}

\olpseganchor{OLP-0051-B003}{start}
\olfileid{sfr}{infinite}{dedekind}
\olsection{د ډېډېکېنډ الجبرونه}
\olpseganchor{OLP-0051-B003}{end}

\olpseganchor{OLP-0051-B004}{start}
موږ يوازې دا نۀ غواړو چې طبيعي عددونه نامتناهي وي؛ غواړو چې ځينې
(الجبري) خاصيتونه هم ولري: بايد د جمع، ضرب او نورو عملياتو له مخې
سم چلند وکړي۔
\olpseganchor{OLP-0051-B004}{end}

\olpseganchor{OLP-0051-B005}{start}
د ډېډېکېنډ نظر دا و چې د \emph{تالي تابعې} مفهوم بنسټيز وګڼي، او
بيا عددونه د هغو څيزونو په توګه وپېژني چې لاندې خاصيتونه لري:
\begin{enumerate}
	\item يو عدد، \LR{$0$}، شته چې د هېڅ عدد تالي نۀ دے
	\\يعنې \LR{$0 \notin \ran{s}$}
	\\يعنې \LR{$\forall x\ s(x) \neq 0$}
	\item بېل عددونه بېل تالي لري
	\\يعنې \LR{$s$} يوه يو پر يو تابع ده
	\\يعنې \LR{$\forall x \forall y (s(x) = s(y) \lif x = y)$}
	\item\ollabel{repeatedapplication} هر عدد پر \LR{$0$} د تالي تابعې له
	بيا بيا تطبيق څخه ترلاسه کېږي۔
\end{enumerate}
لومړي دوه شرطونه د لومړۍ مرتبې د منطق په کارولو په اسانۍ سمبالېږي
(پورته وګورئ)۔ خو يوازې د لومړۍ مرتبې په منطق سره
\olref{repeatedapplication} نشو سمبالولے۔ د ډېډېکېنډ ستره بريا دا وه
چې شرط \olref{repeatedapplication} ئې په سټي بڼه داسې بيا بيان کړ:
\begin{enumerate}
	\item[3\LR{$'$}.] طبيعي عددونه تر ټولو کوچنے سټ دے چې صفر پکښې غړے وي او
	\emph{د تالي تابعې له مخې بند} وي: يعنې که \LR{$s$} د سټ پر هر غړے تطبيق
	کړو، د سټ يو بل غړے ترلاسه کوو۔
\end{enumerate}
خو دا خبره بايد ورو ورو په تفصيل وليکو۔
\olpseganchor{OLP-0051-B005}{end}

\olpseganchor{OLP-0051-B006}{start}
\begin{defn}\ollabel{Closure}
	فرض کړئ \LR{$f\colon A\to A$} يوه تابعه ده۔ د \LR{$A$} فرعي سټ
	\LR{$X$} د \LR{$f$} له مخې \emph{بند} هله او يوازې هله دے چې
	\LR{$(\forall x \in X)f(x) \in X$} وي۔ اوس د هر \LR{$o\in A$} دپاره تعريف کړئ:
	$$\closureofunder{f}{o} = \bigcap\Setabs{X\subseteq A}{o \in X\text{\RL{ او }}X
\text{\RL{ د \LR{$f$} له مخې بند دے}}}$$
\end{defn}
\olpseganchor{OLP-0051-B006}{end}

\olpseganchor{OLP-0051-B007}{start}
نو \LR{$\closureofunder{f}{o}$} د ټولو هغو سټونو اشتراک دے چې د \LR{$f$} له
مخې بند دي او \LR{$o$} ئې يو غړے دے۔ په وجداني ډول، نو
\LR{$\closureofunder{f}{o}$} تر ټولو \emph{کوچنے} داسې د \LR{$f$} له مخې بند
سټ دے چې \LR{$o$} ئې يو غړے دے۔ راتلونکې پايله دا وجداني خبره
دقيقه کوي؛
\begin{lem}\ollabel{closureproperties}
	د هرې تابعې \LR{$f\colon A\to A$} او هر \LR{$o \in A$} دپاره:
	\begin{enumerate}
		\item\ollabel{closurehaselem} \LR{$o \in \closureofunder{f}{o}$}؛ او
		\item\ollabel{closureclosed} \LR{$\closureofunder{f}{o}$} د \LR{$f$} له مخې بند دے؛ او
		\item\ollabel{closuresmallest} که \LR{$X$} د \LR{$f$} له مخې بند وي او \LR{$o \in
		X$} وي، نو \LR{$\closureofunder{f}{o} \subseteq X$}
	\end{enumerate}
\end{lem}
\olpseganchor{OLP-0051-B007}{end}

\olpseganchor{OLP-0051-B008}{start}
\begin{proof}
پام وکړئ چې لږ تر لږه يو داسې سټ شته چې د \LR{$f$} له مخې بند دے او \LR{$o$}
ئې يو غړے دے، يعنې \LR{$\ran{f}\cup \{o\}$}۔ دا د
ساحې فرعي سټ دے، او د هر غړي ارزښت يې د تابعې په حاصل کښې
دے۔ د ټولو مناسبو فرعي سټونو کورنۍ د فرعي سټونو له سټه په
بېلونه اخلو۔ نو \LR{$\closureofunder{f}{o}$}، چې د دې \emph{ټولو}
سټونو اشتراک دے، شتون لري۔ اوس بايد \olref{closurehaselem}--\olref{closuresmallest} وارزوو۔
\olpseganchor{OLP-0051-B008}{end}

\olpseganchor{OLP-0051-B009}{start}
د \olref{closurehaselem} په اړه: \LR{$o \in \closureofunder{f}{o}$}،
ځکه دا د هغو سټونو اشتراک دے چې په ټولو کښې \LR{$o$} يو غړے دے۔
\olpseganchor{OLP-0051-B009}{end}

\olpseganchor{OLP-0051-B010}{start}
د \olref{closureclosed} په اړه: فرض کړئ \LR{$x \in
\closureofunder{f}{o}$}۔ نو که \LR{$o \in X$} او \LR{$X$} د \LR{$f$} له مخې بند
وي، \LR{$x \in X$}؛ او ځکه \LR{$X$} د \LR{$f$} له مخې بند دے، اوس \LR{$f(x) \in X$}۔
نو \LR{$f(x) \in \closureofunder{f}{o}$}۔
\olpseganchor{OLP-0051-B010}{end}

\olpseganchor{OLP-0051-B011}{start}
د \olref{closuresmallest} په اړه: په بشپړ عموميت کښې، که \LR{$X
\in C$} وي، نو \LR{$\bigcap C \subseteq X$}۔
\end{proof}
\olpseganchor{OLP-0051-B011}{end}

\olpseganchor{OLP-0051-B012}{start}
د دې په کارولو ويلے شو:
\olpseganchor{OLP-0051-B012}{end}

\olpseganchor{OLP-0051-B013}{start}
\begin{defn}
د \emph{ډېډېکېنډ الجبر} يو سټ \LR{$A$}، يوه تابع \LR{$f \colon A \to A$} او
يو \LR{$o \in A$} دي چې لاندې شرطونه پوره کوي:
	\begin{enumerate}
		\item \ollabel{ded:proper} \LR{$o \notin \ran{f}$}
		\item \ollabel{ded:injection} \LR{$f$} يوه يو پر يو تابع ده
		\item \ollabel{ded:closure} \LR{$A = \closureofunder{f}{o}$}
	\end{enumerate}
\end{defn}
\olpseganchor{OLP-0051-B013}{end}

\olpseganchor{OLP-0051-B014}{start}
ځکه \LR{$A = \closureofunder{f}{o}$} دے، زموږ پخوانۍ پايله وايي چې \LR{$A$}
تر ټولو کوچنے داسې د \LR{$f$} له مخې بند سټ دے چې \LR{$o$} ئې يو غړے
دے۔ په څرګند ډول هر ډېډېکېنډ الجبر د ډېډېکېنډ له مخې نامتناهي دے؛
يوازې د تعريف \olref{ded:proper} او \olref{ded:injection} مادې
وګورئ۔ خو تر دې په زړه پورې حقيقت دا دے چې هر د ډېډېکېنډ له مخې
نامتناهي سټ په ډېډېکېنډ الجبر بدلېدے شي۔
\olpseganchor{OLP-0051-B014}{end}

\olpseganchor{OLP-0051-B015}{start}
\begin{thm}\ollabel{thm:DedekindInfiniteAlgebra}
که کوم د ډېډېکېنډ له مخې نامتناهي سټ وي، نو يو ډېډېکېنډ الجبر شته۔
\end{thm}
\olpseganchor{OLP-0051-B015}{end}

\olpseganchor{OLP-0051-B016}{start}
\begin{proof}
فرض کړئ \LR{$D$} د ډېډېکېنډ له مخې نامتناهي دے۔ نو يوه يو پر يو تابع
\LR{$g \colon D \to D$} او يو غړے \LR{$o \in D \setminus \ran{g}$} شته۔
اوس \LR{$A = \closureofunder{g}{o}$} وټاکئ؛ د
\olref{closureproperties} له مخې \LR{$A$} شته او \LR{$o \in A$}۔
\LR{$f = \funrestrictionto{g}{A}$} وټاکئ۔ ښيو چې \LR{$A, f, o$} يو
ډېډېکېنډ الجبر جوړوي۔
\olpseganchor{OLP-0051-B016}{end}

\olpseganchor{OLP-0051-B017}{start}
د \olref{ded:proper} په اړه: \LR{$o \notin \ran{g}$} او \LR{$\ran{f}
\subseteq \ran{g}$}، نو \LR{$o\notin \ran{f}$}۔
\olpseganchor{OLP-0051-B017}{end}

\olpseganchor{OLP-0051-B018}{start}
د \olref{ded:injection} په اړه: \LR{$g$} پر \LR{$D$} يوه يو پر يو تابع ده؛
نو \LR{$f \subseteq g$} هم بايد يوه يو پر يو تابع وي۔
\olpseganchor{OLP-0051-B018}{end}

\olpseganchor{OLP-0051-B019}{start}
د \olref{ded:closure} په اړه: د \olref{closureproperties} له مخې
\LR{$A$} د \LR{$g$} له مخې بند دے؛ نو په لا قوي ډول \LR{$A$} د \LR{$f$} له مخې بند
دے۔ په دې توګه د \olref{closureproperties} له مخې
\LR{$\closureofunder{f}{o} \subseteq A$}۔ ځکه \LR{$\closureofunder{f}{o}$}
هم د \LR{$f$} له مخې بند دے او \LR{$f = \funrestrictionto{g}{A}$}، نو
\LR{$\closureofunder{f}{o}$} د \LR{$g$} له مخې بند دے۔ بيا د
\olref{closureproperties} له مخې
\LR{$A \subseteq \closureofunder{f}{o}$}۔
\end{proof}
\olpseganchor{OLP-0051-B019}{end}

\olpseganchor{OLP-0052-B004}{start}
\olfileid{sfr}{infinite}{induction}
\olsection[رياضي استقرا]{د ډېډېکېنډ الجبرونه او رياضي استقرا}
\olpseganchor{OLP-0052-B004}{end}

\olpseganchor{OLP-0052-B005}{start}
اوس مهمه خبره دا ده چې يو ډېډېکېنډ الجبر---بلکې \emph{هر}
ډېډېکېنډ الجبر---د طبيعي عددونو د بدل په توګه کار کولے شي۔ دا د
لاندې اسانې نتيجې له برکته ده:
\olpseganchor{OLP-0052-B005}{end}

\olpseganchor{OLP-0052-B006}{start}
\begin{thm}[رياضي استقرا]\ollabel{thm:dedinfiniteinduction}
فرض کړئ \LR{$N, s, o$} يو ډېډېکېنډ الجبر جوړوي۔ بيا د هر سټ \LR{$X$} دپاره:
	\begin{center}
		که \LR{$o \in X$} او \LR{$(\forall {n} \in N \cap X){s}({n}) \in X$} وي، {نو} \LR{$N \subseteq X$}۔
	\end{center}
\end{thm}
\olpseganchor{OLP-0052-B006}{end}

\olpseganchor{OLP-0052-B007}{start}
\begin{proof}
د ډېډېکېنډ الجبر د تعريف له مخې \LR{$N = \closureofunder{s}{o}$}۔ اوس
که هم \LR{${o} \in X$} او هم \LR{$(\forall {n} \in N)(n \in X \lif
{s}({n}) \in X)$} وي، نو \LR{$N = \closureofunder{s}{o} \subseteq X$}۔
\end{proof}
\olpseganchor{OLP-0052-B007}{end}

\olpseganchor{OLP-0052-B008}{start}
ځکه استقرا د طبيعي عددونو ځانګړنه ده، مقصد دا دے۔ هر د ډېډېکېنډ
له مخې نامتناهي سټ راکړل شي، ترې يو ډېډېکېنډ الجبر جوړولے شو او دا
الجبر د طبيعي عددونو د بدل په توګه کارولے شو۔
\olpseganchor{OLP-0052-B008}{end}

\olpseganchor{OLP-0052-B009}{start}
دا منو چې \olref{thm:dedinfiniteinduction} استقرا په \emph{سټي}
اصطلاحاتو کښې بيانوي۔ خو دا اصل په اسانۍ په هغو اصطلاحاتو کښې هم
بيانولے شو چې ښايي ډېرې اشنا وي:
\olpseganchor{OLP-0052-B009}{end}

\olpseganchor{OLP-0052-B010}{start}
\begin{cor}\ollabel{natinductionschema}
فرض کړئ \LR{$N, s, o$} يو ډېډېکېنډ الجبر جوړوي۔ بيا د هر فارمول
\LR{$\phi(x)$} دپاره، چې پاراميټرونه لرلے شي:
\begin{center}
	که \LR{$\phi(o)$} او \LR{$(\forall {n} \in N)(\phi(n)\lif
	\phi({s}({n})))$} وي، {نو} \LR{$(\forall n \in N)\phi(n)$}
\end{center}
\end{cor}
\olpseganchor{OLP-0052-B010}{end}

\olpseganchor{OLP-0052-B011}{start}
\begin{proof}
\LR{$X = \Setabs{n \in N}{\phi(n)}$} وټاکئ، او اوس
\olref{thm:dedinfiniteinduction} وکاروئ۔
\end{proof}
\olpseganchor{OLP-0052-B011}{end}

\olpseganchor{OLP-0052-B012}{start}
په دې پايله کښې مو وويل چې يو فارمول «پاراميټرونه لري»۔ په لنډه
معنا ئې دا ده چې د هر ډول څيزونو \LR{$c_1, \ldots, c_k$} دپاره له
\LR{$\phi(x, c_1, \ldots, c_k)$} سره کار کولے شو۔ په لا دقيق ډول، پايله
بې له دې چې «پاراميټرونه» ياد کړو، داسې بيانولے شو۔ د هر فارمول
\LR{$\phi(x, v_1, \ldots, v_k)$} دپاره، چې ټول ازاد متغيرونه ئې ښودل
شوي وي، لرو:
	\begin{align*}
			\forall v_1 \ldots \forall v_k((&\phi(o, v_1,\ldots, v_k) \land {}\\
			&	(\forall x \in N)(\phi(x,v_1, \ldots, v_k) \lif \phi(s(x), v_1,\ldots, v_k))) \lif {}\\
			&\hspace{3em} (\forall x \in N)\phi(x, v_1,\ldots, v_k))
	\end{align*}
په څرګند ډول، د «پاراميټرونو لرلو» خبره لوستل ډېر اسانه کولے شي۔
(په \olref[sth][][]{part} کښې به دا طريقه ډېره وکاروو۔)
\olpseganchor{OLP-0052-B012}{end}

\olpseganchor{OLP-0052-B013}{start}
بېرته د ډېډېکېنډ الجبرونو خبرې ته راځو: هر ډېډېکېنډ الجبر راکړل
شي، د جمع، ضرب او طاقت معمول حسابي تابعې هم تعريفولے شو۔ خو دا
اسان کار نۀ دے او د \emph{بازګشتي تعريف} طريقه کاروي۔ دا هغه طريقه
ده چې ډېر وروسته، په ډېر عمومي چوکاټ کښې، معرفي او توجيه کوو۔
(مينه‌وال ښايي له \olref[sth][ord-arithmetic][]{chap} وروسته دې
موضوع ته بيا راشي، يا د \mbox{\LR{\latinfont \citealt[pp.~95--8]{Potter2004}}} په څېر
کومه بله شننه ولولي۔) خو چې \LR{$N, s, o$} يو ډېډېکېنډ الجبر جوړوي، په
پاے کښې به لاندې ټاکنې کولے شو:
\begin{align*}
	{a} + {o} &= {a} & & & {a} \times {o} &= {o} & & & {a}^{o} &= s(o)\\
{a} + {s}({b}) &= {s}({a}+{b}) &&& {a} \times {s}({b}) &= ({a}\times {b}) + {a}  & & & {a}^{{s}({b})} &= {a}^{b} \times {a}
\end{align*}
او ښودلے شو چې دا تابعې لکه څنګه چې هيله کېږي، هماغسې چلند کوي۔
\olpseganchor{OLP-0052-B013}{end}

\olpseganchor{OLP-0053-B003}{start}
\olfileid{sfr}{infinite}{dedekindsproof}
\olpseganchor{OLP-0053-B003}{end}

\olpseganchor{OLP-0053-B004}{start}
\olsection[د ډېډېکېنډ «ثبوت»]{د نامتناهي سټ د شتون په اړه د
ډېډېکېنډ «ثبوت»}
\olpseganchor{OLP-0053-B004}{end}

\olpseganchor{OLP-0053-B005}{start}
په دې څپرکي کښې مو د طبيعي عددونو سټي شننه د ډېډېکېنډ د الجبرونو
له مخې وړاندې کړه۔ په \olref[arith][ref]{sec} کښې مو (له نورو
خبرو سره) د تحليل د سټي جوړونې پر فلسفي اهميت غور وکړ۔ اوس بايد د
هغه څه اهميت وارزوو چې دلته مو ترلاسه کړل۔
\olpseganchor{OLP-0053-B005}{end}

\olpseganchor{OLP-0053-B006}{start}
د \olref[sfr][arith][]{chap} په ټوله موده کښې مو طبيعي عددونه
راکړل شوي وګڼل او په څرګند ډول مو ترې صحيح، ناطق او حقيقي عددونه
جوړ کړل۔ په دې څپرکي کښې مو د طبيعي عددونو کوم څرګند جوړښت نۀ دے
وړاندې کړے۔ يوازې مو وښودل چې \emph{هر د ډېډېکېنډ له مخې نامتناهي
سټ راکړل شي}، داسې سټ تعريفولے شو چې کټ مټ لکه څنګه چې غواړو
\LR{$\Nat$} وچلېږي، هماغسې به وچلېږي۔
\olpseganchor{OLP-0053-B006}{end}

\olpseganchor{OLP-0053-B007}{start}
نو څرګنده ده چې دا ادعا نشو کولے چې کومې ماوراءالطبيعتي پوښتنې ته
مو ځواب ورکړے، لکه دا چې \emph{طبيعي عددونه کوم څيزونه دي}۔ خو دا
ښۀ کار دے۔ په \olref[sfr][arith][ref]{sec} کښې مو ټينګار کړے و چې
د ډېډېکېنډ د پرېکونونو د سټ په توګه د \LR{$\Real$} تعريف بايد يوه اسانه
ټاکنه وګڼو، نۀ يوه \emph{موندنه}۔ مهمه مشاهده دا ده چې د
ډېډېکېنډ پرېکونونه د حقيقي عددونو بنسټيز رياضي خاصيتونه لري۔ دلته
هم همدا خبره ده: مهمه مشاهده دا ده چې \emph{هر} ډېډېکېنډ الجبر د
طبيعي عددونو بنسټيز رياضي خاصيتونه لري۔ (په رښتيا، ډېډېکېنډ دا
ټکے په دې ثبوت لا پياوړے کړ چې ټول ډېډېکېنډ الجبرونه
\emph{آيزومورف} دي (\mbox{\LR{\latinfont \citeyear[Theorems
132--3]{Dedekind1888}}})۔ نو
دا هېڅ د حيرانۍ خبره نۀ ده چې اوسني ډېر «جوړښت‌پال» ډېډېکېنډ خپل
مخکښ بولي۔)
\olpseganchor{OLP-0053-B007}{end}

\olpseganchor{OLP-0053-B008}{start}
برسېره پر دې، موږ وښودل چې د طبيعي عددونو نظريه څنګه په ساده سټ
نظريه کښې ورګډېږي؛ دا ساده نظريه خپله لا هم يو څه غېر رسمي ده، خو
(لکه چې ښکاري) طبيعي عددونه له مخکښې راکړل شوي نۀ ګڼي۔ نو کېدے شي
د ډېډېکېنډ د خپلې سترې پروژې د پوره کولو پر لار يو۔ هغه دا پروژه
د ۱۸۸۸م کال د جرمني متن په سريزه کښې داسې تشريح کړې وه:
\begin{quote}
  په علم کښې چې کومه خبره ثابتېدے شي، بايد يوازې په ثبوت ومنل
  شي۔ د دې غوښتنې معقوليت څرګند دے، خو زما په اند د تر ټولو ساده علم
  د بنسټ په اېښودلو کښې هم لا نۀ ده پوره شوې، آن په وروستيو
  بيانونو کښې۔ مقصد مې د منطق هغه برخه ده چې د عددونو نظريه
  څېړي۔ کله چې حساب (الجبر، تحليل) يوازې د منطق يوه برخه بولم، دا هم
  وايم چې د عدد مفهوم د مکان او زمان له تصورونو او شهودي ادراکونو
  څخه بشپړ خپلواک ګڼم او د فکر د خالصو قوانينو نېغه زېږنده ئې بولم۔
  \mbox{\LR{\latinfont \citep[preface]{Dedekind1888}}}
\end{quote}
د ډېډېکېنډ زړه‌وره مفکوره دا ده۔ همدا اوس مو وښودل چې يوازې د
(ساده) سټ نظريې په کارولو طبيعي عددونه څنګه جوړېږي۔ په
\olref[sfr][arith][]{chap} کښې مو وليدل چې د طبيعي عددونو او يوې
اندازې سټ نظريې په راکړه کېدو سره حقيقي عددونه څنګه جوړېږي۔ نو
ښايي «حساب (الجبر، تحليل)» «يوازې د منطق يوه برخه» وګرځي (د
ډېډېکېنډ د «منطق» کلمې په پراخه معنا)۔
\olpseganchor{OLP-0053-B008}{end}

\olpseganchor{OLP-0053-B009}{start}
نظر همدا دے۔ خو يوه شېبه تم شئ۔ زموږ د ډېډېکېنډ الجبر جوړول (چې
د طبيعي عددونو بدل دے) د يوه د ډېډېکېنډ له مخې نامتناهي سټ د شتون
پر شرط ولاړ دي۔ (يوازې \olref[sfr][infinite][dedekind]{thm:DedekindInfiniteAlgebra}
ته بېرته وګورئ۔) تر څو چې د ډېډېکېنډ له مخې د نامتناهي سټ شتون د
«منطق» يا «د فکر د خالصو قوانينو» له لارې ثابت نشي، پروژه درېږي۔
\olpseganchor{OLP-0053-B009}{end}

\olpseganchor{OLP-0053-B010}{start}
نو، \emph{ايا} د ډېډېکېنډ له مخې د نامتناهي سټ شتون «د فکر د
خالصو قوانينو» له لارې ثابتېدے شي؟ د ډېډېکېنډ هڅه دا وه:
\begin{quote}
  زما د فکر نړۍ، يعنې د ټولو هغو څيزونو ټولګه \LR{$S$} چې زما د فکر
  موضوع کېدے شي، نامتناهي ده۔ ځکه که \LR{$s$} د \LR{$S$} يو غړے وي، دا فکر
  \LR{$s'$} چې \LR{$s$} زما د فکر موضوع کېدے شي، خپله د \LR{$S$} يو غړے دے۔
  که دا د \LR{$s$} تصوير \LR{$\phi(s)$} وګڼو، نو \dots~\LR{$S$} نامتناهي دے،
  همدا مو ثابتول غوښتل۔
  \mbox{\LR{\latinfont \citep[\S66]{Dedekind1888}}}
\end{quote}
داسې خبره د يوه داسې کتاب په منځ کښې ليدل رښتيا هم ډېره حيرانوونکې
ده چې زياته برخه ئې له بشپړو دقيقو رياضي ثبوتونو جوړه ده۔ دوه
تبصرې د پام وړ دي۔
\olpseganchor{OLP-0053-B010}{end}

\olpseganchor{OLP-0053-B011}{start}
لومړۍ: دا «ثبوت» هغه بڼه نږدې هېڅ نۀ لري چې اوس به ئې «رياضي»
وبولو۔ دا د نفسياتي څيزونو (فکرونو) خبره کوي، او هغه هم يوازې
\emph{ممکن} څيزونه۔
\olpseganchor{OLP-0053-B011}{end}

\olpseganchor{OLP-0053-B012}{start}
دويمه: لږ تر لږه لکه څنګه چې موږ د ډېډېکېنډ الجبرونه وړاندې کړي،
دا «ثبوت» يوه څرګنده فني نيمګړتيا لري۔ که د ډېډېکېنډ استدلال بريالے
وي، يوازې دا ثابتوي چې نامتناهي ډېر څيزونه (په ځانګړې توګه
نامتناهي ډېر فکرونه) شته۔ خو ډېډېکېنډ بايد دا دليل هم راکړي چې~\LR{$S$}
ولې يو يوازينے \emph{سټ} وګڼو چې نامتناهي ډېر غړي لري،
نۀ دا چې~\LR{$S$} \emph{څه څيزونه} (په جمع بڼه) وګڼو۔
\olpseganchor{OLP-0053-B012}{end}

\olpseganchor{OLP-0053-B013}{start}
دا حقيقت چې ډېډېکېنډ دلته تشه نۀ ليدله، ښايي وښيي چې د هغه د
«بشپړتيا» کلمه کټ مټ زموږ د «سټ» کلمې له کارونې سره برابره نۀ
ده۔\footnote{په رښتيا، د دې د باور دپاره نور دليلونه هم لرو؛
\mbox{\LR{\latinfont \citet[p.~23]{Potter2004}}} وګورئ۔} خو دا به ډېره د حيرانۍ خبره نۀ
وي۔ هغه پروژه چې په تېرو دوو څپرکيو کښې مو مخ ته يوړه---د طبيعي
عددونو «جوړول»، او له هغو څخه د صحيح، حقيقي او ناطقو عددونو
«جوړول»---ټوله په ساده ډول شوې ده۔ کله مو چې کوم سټ ته اړتيا لرله،
هماغه سټ مو منلے؛ خو د دې په اړه مو ډېر عمومي اصلونه نۀ دي ټاکلي
چې کټ مټ کوم سټونه کله شتون لري۔ خو پوهېږو چې \emph{ځينې} عمومي
اصلونه پکار دي، ځکه له هغو پرته به د رسل په پاراډوکس کښې ولوېږو۔
\olpseganchor{OLP-0053-B013}{end}

\olpseganchor{OLP-0053-B014}{start}
اوس د دې وخت دے چې له خپلې ساده‌ګۍ څخه ووځو۔
\olpseganchor{OLP-0053-B014}{end}

\olpseganchor{OLP-0054-B003}{start}
\olfileid{sfr}{infinite}{card-sb}
\olpseganchor{OLP-0054-B003}{end}

\olpseganchor{OLP-0054-B004}{start}
\olsection{ضميمه: د شروډر--برنشتاين ثبوت}
\olpseganchor{OLP-0054-B004}{end}

\olpseganchor{OLP-0054-B005}{start}
له ساده سټ نظريې څخه د وتلو مخکښې يو وروستے ساده (خو ډېر فني!)
ثبوت څېړو۔ دا د شروډر--برنشتاين ثبوت دے
(\olref[sfr][siz][sb]{thm:schroder-bernstein}): که
\LR{$\cardle{A}{B}$} او \LR{$\cardle{B}{A}$} وي، نو \LR{$\cardeq{A}{B}$}؛ يعنې که
يو پر يو تابعې \LR{$f \colon A \to B$} او \LR{$g \colon B \to A$} راکړل
شي، نو يوه دوه اړخيزه يو پر يو تابع \LR{$h \colon A \to B$} شته۔
\olpseganchor{OLP-0054-B005}{end}

\olpseganchor{OLP-0054-B006}{start}
په دې څپرکي کښې مو د ډېډېکېنډ د \emph{بندښتونو} مفهوم وکاراوه۔ په
رښتيا، ډېډېکېنډ په همدې مفهوم د شروډر--برنشتاين يو ښکلی ثبوت ورکړ،
چې دلته ئې وړاندې کوو۔ که لږ بېله خو په اصل کښې ورته شننه غواړئ،
دا ثبوت په نژدې ډول د \mbox{\LR{\latinfont \citet[pp.~157--8]{Potter2004}}} پيروي کوي۔
په انټرنېټ کښې لږه پلټنه به هم درته وښيي چې دا---لکه د
\LR{$\sqrt{2}$} غير ناطق‌والے---هغه قضيه ده چې \emph{ډېر} په زړه پورې
او بېل ثبوتونه لري۔
\olpseganchor{OLP-0054-B006}{end}

\olpseganchor{OLP-0054-B007}{start}
د \olref[sfr][infinite][dedekind]{Closure} په شان نښه وکاروئ او د
هر سټ \LR{$B$} او تابعې~\LR{$f$} دپاره وټاکئ:
\[
\Closureofunder{f}{B} = \bigcap \Setabs{X}{B \subseteq X
\text{\RL{ او \LR{$X$} د \LR{$f$} له مخې بند دے}}}
\]
په دې ډول تعريف شوے \LR{$\Closureofunder{f}{B}$} تر ټولو کوچنے داسې د
\LR{$f$} له مخې بند سټ دے چې~\LR{$B$} رانغاړي، يعنې:
\olpseganchor{OLP-0054-B007}{end}

\olpseganchor{OLP-0054-B008}{start}
\begin{lem}\ollabel{Closureprops}
	د هرې تابع \LR{$f$} او هر \LR{$B$} دپاره:
	\begin{enumerate}
		\item\ollabel{Closurehaselem} \LR{$B \subseteq \Closureofunder{f}{B}$}؛ او
		\item\ollabel{Closureclosed} \LR{$\Closureofunder{f}{B}$} د \LR{$f$} له مخې بند دے؛ او
		\item\ollabel{Closuresmallest} که \LR{$X$} د \LR{$f$} له مخې بند وي او \LR{$B
		\subseteq X$} وي، نو \LR{$\Closureofunder{f}{B} \subseteq X$}۔
	\end{enumerate}
\end{lem}
\olpseganchor{OLP-0054-B008}{end}

\olpseganchor{OLP-0054-B009}{start}
\begin{proof}
کټ مټ لکه په \olref[sfr][infinite][dedekind]{closureproperties} کښې۔
\end{proof}
\olpseganchor{OLP-0054-B009}{end}

\olpseganchor{OLP-0054-B010}{start}
شروډر--برنشتاين ته د رسېدو دپاره يو وروستے حقيقت پکار دے:
\olpseganchor{OLP-0054-B010}{end}

\olpseganchor{OLP-0054-B011}{start}
\begin{prop}\ollabel{sbhelper}
که \LR{$A \subseteq B \subseteq C$} او \LR{$A \approx C$} وي، نو \LR{$\cardeq{\cardeq{A}{B}}{C}$}۔
\end{prop}
\olpseganchor{OLP-0054-B011}{end}

\olpseganchor{OLP-0054-B012}{start}
\begin{proof}
يوه دوه اړخيزه يو پر يو تابع \LR{$f \colon C \to A$} راکړل شي، وټاکئ
\LR{$F = \Closureofunder{f}{C \setminus B}$} او له ډومېن \LR{$C$} سره تابع
\LR{$g$} داسې تعريف کړئ:
\[
	g(x) =
	\begin{cases}
			f(x) &\text{\RL{که \LR{$x \in F$} وي}}\\
			x & \text{\RL{که نۀ}}
	\end{cases}
\]
ښيو چې \LR{$g$} له \LR{$C \to B$} يوه دوه اړخيزه يو پر يو تابع ده؛ له دې به راووځي
چې \LR{$\comp{f^{-1}}{g} \colon A \to B$} يوه دوه اړخيزه يو پر يو تابع ده، او
ثبوت بشپړېږي۔
\olpseganchor{OLP-0054-B012}{end}

\olpseganchor{OLP-0054-B013}{start}
لومړے ادعا کوو چې که \LR{$x \in F$} خو \LR{$y\notin F$} وي، نو
\LR{$g(x) \neq g(y)$}۔ د نقيض دپاره فرض کړئ چې برابر دي؛ نو
\LR{$y = g(y) = g(x) = f(x)$}۔ ځکه \LR{$x \in F$} او د
\olref{Closureprops} له مخې \LR{$F$} د \LR{$f$} له مخې بند دے،
\LR{$y = f(x) \in F$} لرو، چې ټکر دے۔
\olpseganchor{OLP-0054-B013}{end}

\olpseganchor{OLP-0054-B014}{start}
اوس فرض کړئ \LR{$g(x) = g(y)$}۔ د پورته خبرې له مخې، \LR{$x \in F$} هله او
يوازې هله چې \LR{$y \in F$} وي۔ که \LR{$x, y \in F$} وي، نو
\LR{$f(x) = g (x) = g(y) = f(y)$}؛ او ځکه \LR{$f$} يوه دوه اړخيزه يو پر يو تابع ده،
\LR{$x = y$}۔ که \LR{$x, y \notin F$} وي، نو \LR{$x = g(x) = g(y) = y$}۔ نو
\LR{$g$} يوه يو پر يو تابع ده۔
\olpseganchor{OLP-0054-B014}{end}

\olpseganchor{OLP-0054-B015}{start}
دا پاتې ده چې وښيو \LR{$\ran{g} = B$}۔ يو \LR{$x \in B \subseteq C$}
وټاکئ۔ که \LR{$x \notin F$} وي، نو \LR{$g(x) = x$}۔ که \LR{$x \in F$} وي، نو د
کوم \LR{$y \in F$} دپاره \LR{$x = f(y)$}؛ ځکه که داسې نۀ وي، نو
\LR{$F \setminus \{x\}$} به د \LR{$f$} له مخې بند وي او \LR{$C\setminus B$} به
وغځوي، چې د \olref{Closureprops} له مخې ناشونې ده؛ اوس
\LR{$g(y) = f(y) = x$}۔
\end{proof}
\olpseganchor{OLP-0054-B015}{end}

\olpseganchor{OLP-0054-B016}{start}
په پاے کښې د اصلي پايلې ثبوت دا دے۔ ياد کړئ چې د تابعې \LR{$h$} او سټ
\LR{$D$} دپاره \LR{$\funimage{h}{D} = \Setabs{h(x)}{x \in D}$} تعريفوو۔
\olpseganchor{OLP-0054-B016}{end}

\olpseganchor{OLP-0054-B017}{start}
\begin{proof}[د شروډر--برنشتاين ثبوت] فرض کړئ \LR{$f \colon A \to B$}
او \LR{$g \colon B \to A$} يو پر يو تابعې دي۔ ځکه
\LR{$\funimage{f}{A} \subseteq B$}، نو
\LR{$\funimage{g}{\funimage{f}{A}} \subseteq g[B] \subseteq A$}۔
برسېره پر دې، \LR{$\comp{f}{g} \colon A \to
\funimage{g}{\funimage{f}{A}}$} يوه يو پر يو تابع ده، ځکه \LR{$g$} او
\LR{$f$} دواړه همداسې دي؛ او په رښتيا \LR{$\comp{f}{g}$} يوه
دوه اړخيزه يو پر يو تابع ده، يوازې په دې سبب چې مقابل ډومېن مو همداسې تعريف
کړے۔ نو \LR{$\cardeq{\funimage{g}{\funimage{f}{A}}}{A}$}، او په دې توګه
د \olref{sbhelper} له مخې يوه دوه اړخيزه يو پر يو تابع \LR{$h \colon A \to
\funimage{g}{B}$} شته۔ برسېره پر دې، \LR{$g^{-1}$} يوه دوه اړخيزه يو پر يو تابع
\LR{$\funimage{g}{B} \to B$} ده۔ نو \LR{$\comp{h}{g^{-1}} \colon A \to B$}
يوه دوه اړخيزه يو پر يو تابع ده۔
\end{proof}
\olpseganchor{OLP-0054-B017}{end}

\olpseganchor{OLP-0055-B004}{start}
\olpart{pl}{بياني منطق}
\olpseganchor{OLP-0055-B004}{end}

\olpseganchor{OLP-0055-B005}{start}
\begin{editorial}
  په دې برخه کښې د کلاسيکي بياني منطق مواد دي۔ لومړے څپرکے
  نسبتاً بنسټيز دے او يوازې تعريفونه او پايلې راغونډوي؛ ډېر ثبوتونه
  نۀ دي بشپړ شوي، بلکې د تمرينونو په توګه پرېښودل شوي دي۔ د ثبوت د
  نظامونو او د بشپړتيا د قضيې مواد د لومړۍ درجې منطق له برخې څخه
  راوړل شوي، چې د لومړۍ درجې منطق ټېګ پکې ناسم ټاکل شوے دے۔ په دې
  توګه د
  محمولونو، اصطلاحاتو او سورونو اړوند هر څه ترې وځي، او د
  جوړښتونه~\LR{$\Struct{M}$} پر ځاے د
  ارزښت ټاکنې~\LR{$\pAssign{v}$} خبره راځي۔
\olpseganchor{OLP-0055-B005}{end}

\olpseganchor{OLP-0055-B006}{start}
  پلان دا دے چې دا برخه لا په تفصيل پراخه شي، او نورې موضوعګانې او
  پايلې، لکه د رښتياوالي تابعوي بشپړتيا، هم ورزياتې شي۔
\end{editorial}
\olpseganchor{OLP-0055-B006}{end}



\olpseganchor{OLP-0055-B008}{start}
\tagfalse{FOL}
\olpseganchor{OLP-0055-B008}{end}



\olpseganchor{OLP-0055-B010}{start}
%
  

\olpseganchor{OLP-0055-B010}{end}

\olpseganchor{OLP-0055-B011}{start}
%
  

\olpseganchor{OLP-0055-B011}{end}

\olpseganchor{OLP-0055-B012}{start}
%
  

\olpseganchor{OLP-0055-B012}{end}

\olpseganchor{OLP-0055-B013}{start}
%
  

\olpseganchor{OLP-0055-B013}{end}



\olpseganchor{OLP-0055-B015}{start}
\tagtrue{FOL}
\olpseganchor{OLP-0055-B015}{end}

\olpseganchor{OLP-0056-B004}{start}
\olchapter{pl}{syn}{نحو او معنٰی پوهنه}
\olpseganchor{OLP-0056-B004}{end}

\olpseganchor{OLP-0056-B005}{start}
\begin{editorial}
  دا يوازې د تعريفونو ډېره لنډه لنډيزه ده۔ بايد پراخه شي، څو د
  فارمولونو پر جوړښت د استقرا د ثبوتونو اسانه پېژندنه او ډېرې نورې
  بېلګې پکې راشي۔
\end{editorial}
\olpseganchor{OLP-0056-B005}{end}

\olpseganchor{OLP-0057-B004}{start}
\olfileid{pl}{syn}{int}
\olpseganchor{OLP-0057-B004}{end}

\olpseganchor{OLP-0057-B005}{start}
\olsection{پېژندنه}
\olpseganchor{OLP-0057-B005}{end}

\olpseganchor{OLP-0057-B006}{start}
بياني منطق له هغو فارمولونه سره کار لري چې له
بياني متغيرونه څخه د \LR{$\lnot$}، \LR{$\land$}، \LR{$\lor$}، \LR{$\lif$} او
\LR{$\liff$} بياني نښلوونکو په کارولو جوړېږي۔ د تصور له مخې،
بياني متغير~\LR{$p$} د داسې يوې جملې يا بيان ځاے نيسي
چې رښتونے يا دروغژن وي۔ هر کله چې په فارمول کښې د
بياني متغير ``د رښتياوالي قيمت'' وټاکل شي، نو د
بياني نښلوونکو په کارولو ترې د جوړو شويو ټولو فارمولونه د
رښتياوالي قيمت هم ټاکل کېږي۔ وايو چې بياني منطق \emph{د رښتياوالي
  تابعوي} دے، ځکه چې معنٰی پوهنه ئې د رښتياوالي د قيمتونو د تابعو
له مخې ورکول کېږي۔ په تېره بيا، په بياني منطق کښې دا نۀ څېړو چې
رښتياوالے او دروغوالے نور څنګه ټاکل کېږي؛ مثلاً دا چې يو څه د يوازې
اتفاقي رښتياوالي پر ځاے بالضروره رښتيا دي، يا څوک پرې پوهېږي چې
رښتيا دي، يا اوس رښتيا دي نۀ دا چې پخوا رښتيا وو يا وروسته به رښتيا
وي۔ موږ يوازې د رښتيا (\LR{$\True$}) او دروغو (\LR{$\False$}) دوه قيمتونه په
پام کښې نيسو، نو دا امکان له بحثه وباسو چې يو بيان دې نۀ رښتونے وي
نۀ دروغژن، يا دې نيم رښتونے وي۔ موږ يوازې هغو نښلوونکو ته هم پام
کوو چې ترې د جوړ شوي فارمول د رښتياوالي قيمت د هغې د برخو د
رښتياوالي په قيمتونو پوره ټاکل کېږي، نۀ مثلاً د هغې په معنٰی۔ په
تېره بيا، دا لانجمنه ده چې په انګرېزۍ کښې د شرطيو رښتياوالي قيمت په
دې معنٰی تابعوي دے که نۀ۔ مادي شرطي~\LR{$\lif$} تابعوي دے؛ نور منطقونه
له هغو شرطيو سره کار کوي چې د رښتياوالي تابعوي نۀ وي۔
\olpseganchor{OLP-0057-B006}{end}

\olpseganchor{OLP-0057-B007}{start}
د رښتياوالي د تابعوي بياني منطق د نظريې او مابعدمنطق د ودې دپاره
بايد لومړے د هغه د عبارتونو نحو او معنٰی پوهنه تعريف کړو۔ موږ به له
بياني متغيرونه څخه د نښلوونکو په کارولو د فارمولونه
د جوړولو يوه لار بيان کړو۔ نور تعريفونه هم ممکن دي۔ نور نظامونه به
بېلې نښې غوره کړي، د بنسټيزو نښلوونکو بېل سټونه به وټاکي او قوسونه
به په بېله توګه وکاروي، يا به ئې هېڅ و نۀ کاروي، لکه په تش په نامه
پولنډۍ نښه‌اېښودنه کښې۔ خو په ټولو لارو کښې يو ګډ شے دا دے چې د
جوړېدو قاعدې د فارمولونه سټ \emph{په استقرايي ډول} تعريفوي۔ که
کار په سمه توګه وشي، هر عبارت د جوړېدو د قاعدو له مخې په اصل کښې
يوازې په يوه لار ترلاسه کېدے شي۔ هغه استقرايي تعريف چې
\emph{په يکتا ډول لوستل کېدونکي} عبارتونه ورکوي، موږ ته اجازه راکوي
چې په هماغه طريقه، يعنې په استقرايي تعريف، دغو عبارتونو ته معنٰی
ورکړو۔
\olpseganchor{OLP-0057-B007}{end}

\olpseganchor{OLP-0057-B008}{start}
عبارتونو ته معنٰی ورکول د معنٰی پوهنې ډګر دے۔ د بياني منطق په
معنٰی پوهنه کښې مرکزي مفهوم په ارزښت ټاکنه کښې صدق دے۔
ارزښت ټاکنه~\LR{$\pAssign{v}$} بياني متغيرونه ته د
رښتياوالي \LR{$\True$} او \LR{$\False$} قيمتونه ورکوي۔ هر ارزښت ټاکنه د هر
فارمول~\LR{$!A$} دپاره د رښتياوالي قيمت \LR{$\pValue{v}(!A)$} ټاکي۔
فارمول په ارزښت ټاکنه~\LR{$\pAssign{v}$} کښې هله او يوازې هله
صادق دے چې \LR{$\pValue{v}(!A) = \True$} وي؛ دا په \LR{$\pSat{v}{!A}$} ليکو۔
دا اړيکه د~\LR{$!A$} پر جوړښت د استقرا له مخې هم تعريفېدے شي: د منطقي
نښلوونکو د رښتياوالي تابعې کاروو، څو مثلاً د \LR{$!A \land !B$} صدق د
\LR{$!A$} او~\LR{$!B$} د صدق يا نۀ صدق له مخې تعريف کړو۔
\olpseganchor{OLP-0057-B008}{end}

\olpseganchor{OLP-0057-B009}{start}
د جملو د صدق د اړيکې \LR{$\pSat{v}{!A}$} پر بنسټ بيا د تاوتولوژۍ،
استلزام او د صدق وړتيا بنسټيز معنايي مفهومونه تعريفولے شو۔
فارمول هله تاوتولوژي، \LR{$\Entails !A$}، ده چې هر ارزښت ټاکنه
ئې صادق کړي، يعنې د هر~\LR{$\pAssign{v}$} دپاره
\LR{$\pValue{v}(!A) = \True$} وي۔ د فارمولونه يو سټ هله هغې ته
استلزام ورکوي، \LR{$\Gamma \Entails !A$}، چې هر هغه ارزښت ټاکنه چې د
~\LR{$\Gamma$} ټول فارمولونه صادق کړي، \LR{$!A$} هم صادق کړي۔ او د
فارمولونه يو سټ هله د صدق وړ دے چې کوم ارزښت ټاکنه ئې ټول
فارمولونه په يو وخت صادق کړي۔ ځکه چې فارمولونه په استقرايي
ډول تعريف شوي او صدق بيا د فارمولونه پر جوړښت د استقرا له مخې
تعريف شوے، موږ استقرا کارولے شو چې د خپلې معنٰی پوهنې خاصيتونه ثابت
او تعريف شوي معنايي مفهومونه يو له بل سره وتړو۔
\olpseganchor{OLP-0057-B009}{end}

\olpseganchor{OLP-0058-B004}{start}
\olfileid{pl}{syn}{fml}
\olpseganchor{OLP-0058-B004}{end}

\olpseganchor{OLP-0058-B005}{start}
\olsection{بياني فارمولونه}
\olpseganchor{OLP-0058-B005}{end}

\olpseganchor{OLP-0058-B006}{start}
د بياني منطق فارمولونه له \emph{بياني متغيرونه}، د دروغوالي له
  بياني ثابت~\LR{$\lfalse$}
    او د رښتياوالي له بياني ثابت~\LR{$\ltrue$} څخه د \emph{منطقي
  نښلوونکو} په کارولو جوړېږي۔
\olpseganchor{OLP-0058-B006}{end}

\olpseganchor{OLP-0058-B007}{start}
\begin{enumerate}
\item د بياني متغيرونه \LR{$\Obj p_0$}،
  \LR{$\Obj p_1$}، \dots يو شمېرېدونکے لامتناهي سټ~\LR{$\PVar$}
\item د دروغوالی بياني ثابت~\LR{$\lfalse$}۔
\item د رښتياوالی بياني ثابت~\LR{$\ltrue$}۔
\item منطقي نښلوونکي:
  \startycommalist
  \ycomma \LR{$\lnot$} (نفي)%
  \ycomma \LR{$\land$} (عطف)%
  \ycomma \LR{$\lor$} (فصل)%
  \ycomma \LR{$\lif$} (شرطي)%
  \ycomma \LR{$\liff$} (دوه اړخيز شرطي)%
\item د وقف نښې: (، ) او کامه۔
\end{enumerate}
\olpseganchor{OLP-0058-B007}{end}

\olpseganchor{OLP-0058-B008}{start}
د بياني منطق دا ژبه په \LR{$\Lang L_0$} ښيو۔
\olpseganchor{OLP-0058-B008}{end}

\olpseganchor{OLP-0058-B009}{start}
۔
\olpseganchor{OLP-0058-B009}{end}



\olpseganchor{OLP-0058-B011}{start}
% Alternate symbols
\olpseganchor{OLP-0058-B011}{end}

\olpseganchor{OLP-0058-B012}{start}
\begin{intro}
کېدے شي تاسو له هغو اصطلاحاتو او نښو سره بلد يئ چې زموږ له پورته
کارولو سره توپير لري۔ د منطق متنونه، او ښوونکي، عموماً د ``نفي''
دپاره \LR{$\sim$}، \LR{$\neg$} او~!، او د ``عطف'' دپاره \LR{$\wedge$}، \LR{$\cdot$} او
\LR{$\&$} کاروي۔ د ``شرطي'' يا ``استلزام'' عامې نښې \LR{$\rightarrow$}،
\LR{$\Rightarrow$} او \LR{$\supset$} دي۔
د ``دوه اړخيز شرطي''، ``دوه اړخيز استلزام''
  يا ``(مادي) هم ارزښت'' نښې \LR{$\leftrightarrow$}،
  \LR{$\Leftrightarrow$} او \LR{$\equiv$} دي۔
د \LR{$\lfalse$} نښې ته په بېلو ځايونو کښې
  ``دروغوالے''، ``فالسوم''، ``محال'' يا ``ښکته'' ويل کېږي۔
د \LR{$\ltrue$} نښې ته په بېلو ځايونو کښې
  ``رښتياوالے''، ``ويروم'' يا ``پورته'' ويل کېږي۔
\end{intro}
\olpseganchor{OLP-0058-B012}{end}

\olpseganchor{OLP-0058-B013}{start}
\begin{defn}[فارمول]
\ollabel{defn:formulas}
د بياني منطق د \emph{فارمولونه} سټ~\LR{$\Frm[L_0]$} په استقرايي ډول
داسې تعريفېږي:
\begin{enumerate}
\item \LR{$\lfalse$} يو اټومي فارمول دے۔
\olpseganchor{OLP-0058-B013}{end}

\olpseganchor{OLP-0058-B014}{start}
\item \LR{$\ltrue$} يو اټومي فارمول دے۔
\olpseganchor{OLP-0058-B014}{end}

\olpseganchor{OLP-0058-B015}{start}
\item هر بياني متغير~\LR{$\Obj p_i$} يو اټومي
  فارمول دے۔
\olpseganchor{OLP-0058-B015}{end}

\olpseganchor{OLP-0058-B016}{start}
\item که \LR{$!A$} يو فارمول وي، نو \LR{$\lnot !A$} هم
  يو فارمول دے۔
\olpseganchor{OLP-0058-B016}{end}

\olpseganchor{OLP-0058-B017}{start}
\item که \LR{$!A$} او \LR{$!B$} فارمولونه وي، نو \LR{$(!A \land
  !B)$} يو فارمول دے۔
\olpseganchor{OLP-0058-B017}{end}

\olpseganchor{OLP-0058-B018}{start}
\item که \LR{$!A$} او \LR{$!B$} فارمولونه وي، نو \LR{$(!A \lor !B)$}
  يو فارمول دے۔
\olpseganchor{OLP-0058-B018}{end}

\olpseganchor{OLP-0058-B019}{start}
\item که \LR{$!A$} او \LR{$!B$} فارمولونه وي، نو \LR{$(!A \lif !B)$}
  يو فارمول دے۔
\olpseganchor{OLP-0058-B019}{end}

\olpseganchor{OLP-0058-B020}{start}
\item که \LR{$!A$} او \LR{$!B$} فارمولونه وي، نو \LR{$(!A \liff !B)$}
  يو فارمول دے۔
\olpseganchor{OLP-0058-B020}{end}

\olpseganchor{OLP-0058-B021}{start}
\item بل هېڅ شے فارمول نۀ دے۔
\end{enumerate}
\end{defn}
\olpseganchor{OLP-0058-B021}{end}

\olpseganchor{OLP-0058-B022}{start}
\begin{explain}
د فارمولونه تعريف يو \emph{استقرايي تعريف} دے۔ په اصل کښې موږ
د فارمولونه سټ په نامتناهي ډېرو پړاوونو کښې جوړوو۔ په لومړني
پړاو کښې ټول اټومي فارمولونه فارمولونه ګڼو؛ دا د تعريف له لومړنيو
څو حالتونو سره سمون لري، يعنې د
\LR{$\ltrue$}، %
\LR{$\lfalse$}، %
\LR{$\Obj p_i$} حالتونه۔ په دې توګه ``اټومي فارمول'' د همدې بڼې هر
فارمول ته وايي۔
\olpseganchor{OLP-0058-B022}{end}

\olpseganchor{OLP-0058-B023}{start}
د تعريف نور حالتونه له هغو فارمولونه څخه د نويو فارمولونه د
جوړولو قاعدې راکوي چې له مخکښې جوړ شوي وي۔ په دويم پړاو کښې له دغو
قاعدو څخه په کار اخيستو له اټومي فارمولونه څخه فارمولونه
جوړولے شو۔ په درېيم پړاو کښې له اټومي فارمولونو او په دويم پړاو کښې
له ترلاسه شويو فارمولونو څخه نوي فارمولونه جوړوو، او همداسې مخکښې
ځو۔ فارمول هر هغه شے دے چې بالاخره په داسې کوم پړاو کښې جوړ
شي، او بل هېڅ شے نۀ دے۔
\end{explain}
\olpseganchor{OLP-0058-B023}{end}

\olpseganchor{OLP-0058-B024}{start}
کله چې د دوه ځايه نښلوونکي~\LR{$\ast$} په کارولو له \LR{$!B$} او \LR{$!C$} څخه
جوړ شوے فارمول \LR{$(!B \ast !C)$} ليکو، ډېر ځله د قوسونو تر ټولو بهرنی
جوړ پرېږدو او يوازې~\LR{$!B \ast !C$} ليکو۔
\olpseganchor{OLP-0058-B024}{end}

\olpseganchor{OLP-0058-B034}{start}
\begin{defn}[نحوي يوشانوالی]
د \LR{$\ident$} نښه د نښو د توريزو لړيو ترمنځ نحوي يوشانوالی څرګندوي؛
يعنې \LR{$!A \ident !B$} هله او يوازې هله چې \LR{$!A$} او \LR{$!B$} د نښو د يو
شان اوږدوالي توريزې لړۍ وي او په هر ځاے کښې يوه نښه ولري۔
\end{defn}
\olpseganchor{OLP-0058-B034}{end}

\olpseganchor{OLP-0058-B035}{start}
د \LR{$\ident$} نښې دواړو خواوو ته د يو ځاے کولو له لارې ترلاسه شوې
توريزې لړۍ هم راتلے شي۔ مثلاً \LR{$!A \ident (!B \lor !C)$} دا معنٰی لري:
د نښو توريزه لړۍ~\LR{$!A$} هماغه توريزه لړۍ ده چې په همدې ترتيب د پرانيستي
قوس، د \LR{$!B$} توريزې لړۍ، د \LR{$\lor$} نښې، د~\LR{$!C$} توريزې لړۍ او د تړلي
قوس په يو ځاے کولو ترلاسه کېږي۔ که داسې وي، نو پوهېږو چې د \LR{$!A$}
لومړۍ نښه پرانيستے قوس دے، \LR{$!A$} د \LR{$!B$} توريزه لړۍ د فرعي توريزې
لړۍ په توګه لري چې له دويمې نښې پيلېږي، ورپسې \LR{$\lor$} راځي، او داسې نور۔
\olpseganchor{OLP-0058-B035}{end}

\olpseganchor{OLP-0059-B004}{start}
\olfileid{pl}{syn}{pre}
\olpseganchor{OLP-0059-B004}{end}

\olpseganchor{OLP-0059-B005}{start}
\olsection{لومړنۍ پايلې}
\olpseganchor{OLP-0059-B005}{end}

\olpseganchor{OLP-0059-B006}{start}
\begin{thm}[\emph{د فارمولونه دپاره د استقرا اصل}]
  \ollabel{thm:induction}  
  که کوم خاصيت~\LR{$P$} د ټولو اټومي فارمولونه دپاره صدق کوي او داسې
  وي چې
  \begin{enumerate}
    \item هر کله چې د~\LR{$!A$} دپاره صدق کوي، د \LR{$\lnot !A$}
      دپاره هم صدق کوي؛
    \item هر کله چې د \LR{$!A$} او~\LR{$!B$} دپاره صدق کوي، د
      \LR{$(!A \land !B)$} دپاره هم صدق کوي؛
    \item هر کله چې د \LR{$!A$} او~\LR{$!B$} دپاره صدق کوي، د
      \LR{$(!A \lor !B)$} دپاره هم صدق کوي؛
    \item هر کله چې د \LR{$!A$} او~\LR{$!B$} دپاره صدق کوي، د
      \LR{$(!A \lif !B)$} دپاره هم صدق کوي؛
    \item هر کله چې د \LR{$!A$} او~\LR{$!B$} دپاره صدق کوي، د
      \LR{$(!A \liff !B)$} دپاره هم صدق کوي؛
  \end{enumerate}
  نو \LR{$P$} د ټولو فارمولونه دپاره صدق کوي۔
\end{thm}
\olpseganchor{OLP-0059-B006}{end}

\olpseganchor{OLP-0059-B007}{start}
\begin{proof}
  \LR{$S$} دې د هغو ټولو فارمولونه مجموعه وي چې خاصيت~\LR{$P$} لري۔ په
  څرګنده \LR{$S \subseteq \Frm[L_0]$}۔ \LR{$S$} د
  \olref[fml]{defn:formulas} ټول شرطونه پوره کوي: ټول اټومي
  فارمولونه لري او د منطقي عملګرونه له مخې بند دے۔ \LR{$\Frm[L_0]$} د
  دې ډول تر ټولو وړوکے ټولګے دے، نو \LR{$\Frm[L_0] \subseteq S$}۔ له دې
  امله \LR{$\Frm[L_0] = S$} او هر فارمول خاصيت~\LR{$P$} لري۔
\end{proof}
\olpseganchor{OLP-0059-B007}{end}

\olpseganchor{OLP-0059-B008}{start}
\begin{prop}\ollabel{prop:balanced}
   په~\LR{$\Frm[L_0]$} کښې هر فارمول \emph{متوازن} دے؛ يعنې د کيڼو
   قوسونو شمېر ئې د ښيو قوسونو له شمېر سره برابر دے۔
\end{prop}
\olpseganchor{OLP-0059-B008}{end}

\olpseganchor{OLP-0059-B009}{start}
\begin{prob}
د \olref[pl][syn][pre]{prop:balanced} ثبوت ورکړئ۔
\end{prob}
\olpseganchor{OLP-0059-B009}{end}

\olpseganchor{OLP-0059-B010}{start}
\begin{prop} \ollabel{prop:noinit}
د فارمول هېڅ خاصه پيلنۍ برخه فارمول نۀ ده۔
\end{prop}
\olpseganchor{OLP-0059-B010}{end}

\olpseganchor{OLP-0059-B011}{start}
\begin{prob}
  د \olref[pl][syn][pre]{prop:noinit} ثبوت ورکړئ۔
\end{prob}
\olpseganchor{OLP-0059-B011}{end}

\olpseganchor{OLP-0059-B012}{start}
\begin{prop}[يکتا لوستونتيا]
په \LR{$ \Frm[L_0]$} کښې هر فارمول~\LR{$!A$} په لاندې بڼو کښې د يوې په
توګه پوره يوه تجزيه لري:
\begin{enumerate}
\item \LR{$\lfalse$}۔
\olpseganchor{OLP-0059-B012}{end}

\olpseganchor{OLP-0059-B013}{start}
\item \LR{$\ltrue$}۔
\olpseganchor{OLP-0059-B013}{end}

\olpseganchor{OLP-0059-B014}{start}
\item د کوم \LR{$\Obj p_n \in  \PVar$} دپاره \LR{$\Obj p_n$}۔
\olpseganchor{OLP-0059-B014}{end}
  
\olpseganchor{OLP-0059-B015}{start}
\item د کوم فارمول~\LR{$!B$} دپاره \LR{$\lnot !B$}۔
\olpseganchor{OLP-0059-B015}{end}

\olpseganchor{OLP-0059-B016}{start}
\item د کومو فارمولونه \LR{$!B$} او~\LR{$!C$} دپاره
  \LR{$(!B \land !C)$}۔
\olpseganchor{OLP-0059-B016}{end}

\olpseganchor{OLP-0059-B017}{start}
\item د کومو فارمولونه \LR{$!B$} او~\LR{$!C$} دپاره
  \LR{$(!B \lor !C)$}۔
\olpseganchor{OLP-0059-B017}{end}

\olpseganchor{OLP-0059-B018}{start}
\item د کومو فارمولونه \LR{$!B$} او~\LR{$!C$} دپاره
  \LR{$(!B \lif !C)$}۔
\olpseganchor{OLP-0059-B018}{end}

\olpseganchor{OLP-0059-B019}{start}
\item د کومو فارمولونه \LR{$!B$} او~\LR{$!C$} دپاره
  \LR{$(!B \liff !C)$}۔
\end{enumerate}
سربېره پر دې، دا تجزيه \emph{يکتا} ده۔
\end{prop}
\olpseganchor{OLP-0059-B019}{end}

\olpseganchor{OLP-0059-B020}{start}
\begin{proof}
پر \LR{$!A$} استقرا کوو۔ مثلاً فرض کړئ چې \LR{$!A$} د \LR{$(!B \lif !C)$} او
\LR{$(!B' \lif !C')$} په توګه دوه بېلې لوستنې لري۔ نو \LR{$!B$} او \LR{$!B'$}
بايد يو شان وي، ګنې يو به د بل خاصه پيلنۍ برخه وي؛ نو که د \LR{$!A$}
دواړه لوستنې بېلې وي، علت ئې خامخا دا دے چې \LR{$!C$} او \LR{$!C'$} د نښو د
هماغې يوې لړۍ بېلې لوستنې دي، چې د استقرايي فرض له مخې ناشونې ده۔
\end{proof}
\olpseganchor{OLP-0059-B020}{end}


\olpseganchor{OLP-0059-B021}{start}
\begin{defn}[يو شان ځايناستي]
که \LR{$!A$} او \LR{$!B$} فارمولونه وي او \LR{$\Obj p_i$} يو بياني متغير وي، نو \LR{$\Subst{!A}{!B}{\Obj p_i}$} هغه پايله ښيي چې په
\LR{$!A$} کښې د \LR{$\Obj p_i$} هر راتګ د \LR{$!B$} په يوه راتګ بدل شي؛ همدارنګه،
د \LR{$\Obj p_1$}، \dots،~\LR{$\Obj p_n$} پر ځاے په يو وخت د فارمولونه
\LR{$!B_1$}، \dots،~\LR{$!B_n$} ځايناستي په
\LR{$\SSubst{!A}{\subst{!B_1}{\Obj p_1},\dots,\subst{!B_n}{\Obj p_n}}$}
ښودل کېږي۔
\end{defn}
\olpseganchor{OLP-0059-B021}{end}

\olpseganchor{OLP-0059-B022}{start}
\begin{prob} د لاندې پنځو فارمولونه په هر يوه کښې وټاکئ چې آيا
 فارمول د \( \Subst{!A}{!B}{\Obj p_i} \) د ځايناستي په توګه
 څرګندېدے شي، چې \( !A \) په دې کښې يو وي: (الف) \( \Obj p_0 \)؛
 (ب) \( ( \lnot \Obj p_0 \land \Obj p_1) \)؛ او (ج)
 \( ( ( \lnot \Obj p_0 \lif \Obj p_1 ) \land \Obj p_2 ) \)۔ په هر
 حالت کښې اړونده ځايناستي په ګوته کړئ۔
  \begin{enumerate}
    \item \( \Obj p_1 \) 
    \item \( ( \lnot \Obj p_0 \land \Obj p_0 ) \)
    \item \( ( ( \Obj p_0 \lor \Obj p_1 ) \land \Obj p_2 )  \)
    \item \( \lnot ( ( \Obj p_0 \lif \Obj p_1 ) \land \Obj p_2 ) \)
    \item \( (( \lnot ( \Obj p_0 \lif \Obj p_1 ) \lif ( \Obj p_0 \lor \Obj p_1 )) \land \lnot ( \Obj p_0 \land \Obj p_1 )) \)
  \end{enumerate}
\end{prob}
\olpseganchor{OLP-0059-B022}{end}

\olpseganchor{OLP-0059-B023}{start}
\begin{prob}
د \LR{$\Subst{!A}{!B}{p}$} يو له رياضي پلوه کره تعريف د استقرا له مخې
ورکړئ۔
\end{prob}
\olpseganchor{OLP-0059-B023}{end}

\olpseganchor{OLP-0060-B004}{start}
\olfileid{pl}{syn}{fseq}
\olpseganchor{OLP-0060-B004}{end}

\olpseganchor{OLP-0060-B005}{start}
\olsection{جوړښتي لړۍ}
\olpseganchor{OLP-0060-B005}{end}

\olpseganchor{OLP-0060-B006}{start}
په استقرايي تعريف سره د فارمولونه تعريفول، او د استقرا له لارې
د فارمولونه د خاصيتونو د ثابتولو بشپړوونکې طريقه، ښکلې او اغېزمنه
لار ده۔ خو کله کله دا هم ګټوره وي چې د فارمولونه جوړول له بېخه
پورته، ګام په ګام وڅېړو۔ دلته د \emph{جوړښتي لړۍ} په مفهوم همدا کار
کوو۔
\olpseganchor{OLP-0060-B006}{end}

\olpseganchor{OLP-0060-B007}{start}
\begin{defn}[د فارمولونو جوړښتي لړۍ]
\ollabel{defn:fseq-frm}
د ژبې~\LR{$\Lang L_0$} د نښو د توريزو لړيو يوه متناهي لړۍ
\LR{$\tuple{!A_0,\dotsc,!A_n}$} د \LR{$!A$} دپاره \emph{جوړښتي لړۍ} ده که
\LR{$!A \ident !A_n$} وي او د هر \LR{$i \leq n$} دپاره يا \LR{$!A_i$} اټومي
فارمول وي، يا داسې \LR{$j,k < i$} موجود وي چې له لاندې حالتونو څخه يو
صدق وکړي:
\begin{enumerate}
    \item \LR{$!A_i \ident \lnot !A_j$}۔%
    \item \LR{$!A_i \ident (!A_j \land !A_k)$}۔%
    \item \LR{$!A_i \ident (!A_j \lor !A_k)$}۔%
    \item \LR{$!A_i \ident (!A_j \lif !A_k)$}۔%
    \item \LR{$!A_i \ident (!A_j \liff !A_k)$}۔%
\end{enumerate}
\end{defn}
\olpseganchor{OLP-0060-B007}{end}

\olpseganchor{OLP-0060-B008}{start}
\begin{ex}
\[
\tuple{
    \Obj p_0,
    \Obj p_1,
    (\Obj p_1 \land \Obj p_0),
    \lnot (\Obj p_1 \land \Obj p_0)
}
\]
د \LR{$\lnot (\Obj p_1 \land \Obj p_0)$} يوه جوړښتي لړۍ ده، او دا هم ده:
\[
\tuple{
    \Obj p_0,
    \Obj p_1,
    \Obj p_0,
    (\Obj p_1 \land \Obj p_0),
    (\Obj p_0 \lif \Obj p_1),
    \lnot (\Obj p_1 \land \Obj p_0)
}.
\]
%
لکه چې له دويمې بېلګې ښکاري، په جوړښتي لړيو کښې ``بې‌ګټې توکي''
هم راتلے شي: داسې فارمولونه چې زياتي وي يا په جوړونه کښې ونډه نۀ
لري۔
\end{ex}
\olpseganchor{OLP-0060-B008}{end}

\olpseganchor{OLP-0060-B009}{start}
\begin{prop}
\ollabel{prop:formed}
په~\LR{$\Frm[L_0]$} کښې هر فارمول~\LR{$!A$} يوه جوړښتي لړۍ لري۔
\end{prop}
\olpseganchor{OLP-0060-B009}{end}

\olpseganchor{OLP-0060-B010}{start}
\begin{proof}
فرض کړئ \LR{$!A$} اټومي دے۔ نو لړۍ~\LR{$\tuple{!A}$} د~\LR{$!A$} جوړښتي لړۍ ده۔
%
اوس فرض کړئ چې \LR{$!B$} او~\LR{$!C$} په ترتيب سره جوړښتي لړۍ
\LR{$\tuple{!B_0,\dotsc,!B_n}$} او \LR{$\tuple{!C_0,\dotsc,!C_m}$} لري۔
%
\begin{enumerate}
  \item که \LR{$!A \ident \lnot !B$} وي، نو
      \LR{$\tuple{!B_0,\dotsc,!B_n,\lnot !B_n}$}
      د~\LR{$!A$} جوړښتي لړۍ ده۔
  \item که \LR{$!A \ident (!B \land !C)$} وي، نو
      \LR{$\tuple{!B_0,\dotsc,!B_n,!C_0,\dotsc,!C_m,(!B_n \land !C_m)}$}
      د~\LR{$!A$} جوړښتي لړۍ ده۔
  \item که \LR{$!A \ident (!B \lor !C)$} وي، نو
      \LR{$\tuple{!B_0,\dotsc,!B_n,!C_0,\dotsc,!C_m,(!B_n \lor !C_m)}$}
      د~\LR{$!A$} جوړښتي لړۍ ده۔
  \item که \LR{$!A \ident (!B \lif !C)$} وي، نو
      \LR{$\tuple{!B_0,\dotsc,!B_n,!C_0,\dotsc,!C_m,(!B_n \lif !C_m)}$}
      د~\LR{$!A$} جوړښتي لړۍ ده۔
  \item که \LR{$!A \ident (!B \liff !C)$} وي، نو
      \LR{$\tuple{!B_0,\dotsc,!B_n,!C_0,\dotsc,!C_m,(!B_n \liff !C_m)}$}
      د~\LR{$!A$} جوړښتي لړۍ ده۔
\end{enumerate}
د فارمولونه دپاره د استقرا د اصل له مخې، هر فارمول يوه
جوړښتي لړۍ لري۔
\end{proof}
\olpseganchor{OLP-0060-B010}{end}

\olpseganchor{OLP-0060-B011}{start}
د دې عکس هم ثابتولے شو۔ دا ځکه مهم دے چې ښيي د فارمولونو د تعريف
زموږ دواړه لارې همارزښته دي: دواړه يو شان پايلې ورکوي۔ دا هم ترې
معلومېږي چې د جوړښتي لړيو پر اوږدوالي عادي استقرا کارولے شو او د
فارمولونو په اړه قضيې ثابتولے شو۔
\olpseganchor{OLP-0060-B011}{end}

\olpseganchor{OLP-0060-B012}{start}
\begin{lem}
\ollabel{lem:fseq-init}
فرض کړئ \LR{$\tuple{!A_0,\dotsc,!A_n}$} د~\LR{$!A_n$} جوړښتي لړۍ ده او
\LR{$k \leq n$} دے۔ نو \LR{$\tuple{!A_0,\dotsc,!A_k}$} د~\LR{$!A_k$} جوړښتي لړۍ
ده۔
\end{lem}
\olpseganchor{OLP-0060-B012}{end}

\olpseganchor{OLP-0060-B013}{start}
\begin{proof}
تمرين۔
\end{proof}
\olpseganchor{OLP-0060-B013}{end}

\olpseganchor{OLP-0060-B014}{start}
\begin{thm}
\ollabel{thm:fseq-frm-equiv}
\LR{$\Frm[L_0]$} د ژبې~\LR{$\Lang L_0$} د نښو د ټولو هغو توريزو لړيو سټ دے
چې يوه جوړښتي لړۍ لري۔
\end{thm}
\olpseganchor{OLP-0060-B014}{end}

\olpseganchor{OLP-0060-B015}{start}
\begin{proof}
\LR{$F$} دې د ژبې~\LR{$\Lang L_0$} د نښو د هغو ټولو توريزو لړيو سټ وي چې
جوړښتي لړۍ لري۔ په \olref[pl][syn][fseq]{prop:formed} کښې مو وليدل
چې \LR{$\Frm[L_0] \subseteq F$}؛ اوس ئې عکس ثابتوو۔
\olpseganchor{OLP-0060-B015}{end}

\olpseganchor{OLP-0060-B016}{start}
فرض کړئ \LR{$!A$} يوه جوړښتي لړۍ \LR{$\tuple{!A_0,\dotsc,!A_n}$} لري۔ پر~\LR{$n$}
د قوي استقرا له لارې ثابتوو چې \LR{$!A \in \Frm[L_0]$}۔ زموږ استقرايي
فرض دا دے چې د نښو هره هغه توريزه لړۍ چې د \LR{$m < n$} اوږدوالي
جوړښتي لړۍ لري، په~\LR{$\Frm[L_0]$} کښې ده۔ دلته اوږدوالے د لړۍ د وروستي شاخص په حساب اخلو۔ د جوړښتي لړۍ د تعريف له مخې،
يا \LR{$!A_n$} اټومي دے، يا بايد داسې \LR{$j,k < n$} موجود وي چې له لاندې
حالتونو څخه يو صدق وکړي:
\begin{enumerate}
    \item \LR{$!A_n \ident \lnot !A_j$}۔%
    \item \LR{$!A_n \ident (!A_j \land !A_k)$}۔%
    \item \LR{$!A_n \ident (!A_j \lor !A_k)$}۔%
    \item \LR{$!A_n \ident (!A_j \lif !A_k)$}۔%
    \item \LR{$!A_n \ident (!A_j \liff !A_k)$}۔%
\end{enumerate}
اوس حالت په حالت استدلال کوو۔ که \LR{$!A_n$} اټومي وي، نو
\LR{$!A_n \in \Frm[L_0]$}۔ بل حالت کښې فرض کړئ چې \LR{$!A \ident
(!A_j \land !A_k)$}۔ د \olref[pl][syn][fseq]{lem:fseq-init} له مخې،
\LR{$\tuple{!A_0,\dotsc,!A_j}$} او \LR{$\tuple{!A_0,\dotsc,!A_k}$} په ترتيب
سره د \LR{$!A_j$} او~\LR{$!A_k$} جوړښتي لړۍ دي۔ ځکه چې دا د~\LR{$!A$} د جوړښتي
لړۍ خاصې پيلنۍ فرعي لړۍ دي، د دواړو اوږدوالے له~\LR{$n$} څخه لږ دے۔ نو
د استقرايي فرض له مخې \LR{$!A_j$} او~\LR{$!A_k$} په~\LR{$\Frm[L_0]$} کښې دي؛ او
بيا د فارمول د تعريف له مخې \LR{$(!A_j \land !A_k)$} هم پکې دے۔
د نفي په حالت کښې د يوه کوچني شاخص پيلنۍ لړۍ بس ده۔
استقرايي فرض ښيي چې د هغې وروستے غړے فارمول دے، او د نفي
جوړښتي قاعده ښيي چې نفي يې هم فارمول ده۔ د پاتې دوه‌ځاييزو
نښلوونکو دپاره هم دواړه پيلنۍ لړۍ اخلو او اړونده
جوړښتي قاعده کاروو۔
\end{proof}
\olpseganchor{OLP-0060-B016}{end}

\olpseganchor{OLP-0061-B004}{start}
\olfileid{pl}{syn}{val}
\olpseganchor{OLP-0061-B004}{end}

\olpseganchor{OLP-0061-B005}{start}
\olsection{ارزښت ټاکنې او صدق}
\olpseganchor{OLP-0061-B005}{end}

\olpseganchor{OLP-0061-B006}{start}
\begin{defn}[ارزښت ټاکنې]
\LR{$\{\True, \False\}$} دې د رښتياوالي د دوو قيمتونو، ``رښتيا'' او
``دروغ''، سټ وي۔ د \LR{$\Lang{L_0}$} دپاره \emph{ارزښت ټاکنه} يوه
تابع~\LR{$\pAssign{v}$} ده چې د ژبې بياني متغيرونه ته يا
\LR{$\True$} ورکوي يا \LR{$\False$}؛ يعنې
\LR{$\pAssign{v} \colon \PVar \to \{\True, \False \}$}۔
\end{defn}
\olpseganchor{OLP-0061-B006}{end}

\olpseganchor{OLP-0061-B007}{start}
\begin{defn}
  چې ارزښت ټاکنه~\LR{$\pAssign{v}$} راکړل شوے وي، د ارزونې تابع
  \LR{$\pValue{v} \colon \Frm[L_0] \to \{\True, \False \}$} په استقرايي
  ډول داسې تعريف کړئ:
  \begin{align*}
    \pValue{v}(\lfalse) & = \False; \\
    \pValue{v}(\ltrue) & = \True; \\
    \pValue{v}(\Obj p_n) & = \pAssign{v}(\Obj p_n); \\
    \pValue{v}(\lnot !A) & = \begin{cases}
        \True & \text{\RL{که }} \pValue{v}(!A) = \False;\\
        \False & \text{\RL{که داسې نۀ وي۔}}
      \end{cases} \\ 
    \pValue{v}(!A \land !B) & = \begin{cases}
        \True &
        \text{\RL{که \LR{$\pValue{v}(!A) = \True$} او \LR{$\pValue{v}(!B) = \True$} وي؛}}\\
        \False &
        \text{\RL{که \LR{$\pValue{v}(!A) = \False$} يا \LR{$\pValue{v}(!B) = \False$} وي}}.
      \end{cases}\\
    \pValue{v}(!A \lor !B) & = \begin{cases}
        \True &
        \text{\RL{که \LR{$\pValue{v}(!A) = \True$} يا \LR{$\pValue{v}(!B) = \True$} وي؛}}\\
        \False &
        \text{\RL{که \LR{$\pValue{v}(!A) = \False$} او \LR{$\pValue{v}(!B) = \False$} وي}}.
      \end{cases}\\
    \pValue{v}(!A \lif !B) & = \begin{cases}
        \True &
        \text{\RL{که \LR{$\pValue{v}(!A) = \False$} يا \LR{$\pValue{v}(!B) = \True$} وي؛}}\\
        \False &
        \text{\RL{که \LR{$\pValue{v}(!A) = \True$} او \LR{$\pValue{v}(!B) = \False$} وي}}.
      \end{cases}\\
    \pValue{v}(!A \liff !B) & = \begin{cases}
        \True &
        \text{\RL{که \LR{$\pValue{v}(!A) = \pValue{v}(!B)$} وي؛}}\\
        \False &
        \text{\RL{که \LR{$\pValue{v}(!A) \neq \pValue{v}(!B)$} وي}}.
      \end{cases}
\end{align*}
\end{defn}
\olpseganchor{OLP-0061-B007}{end}

\olpseganchor{OLP-0061-B008}{start}
\begin{explain}
دا مادې د رښتياوالي له لاندې جدولونو سره سمون لري:
\begin{center}

  \begin{tabular}{|c||c|} \hline
    \LR{$!A$} & \LR{$ \lnot !A$} \\
    \hline \hline
    \LR{$\True$} & \LR{$\False$} \\
    \LR{$\False$} & \LR{$\True$} \\
    \hline
  \end{tabular}


  \begin{tabular}{|cc||c|} \hline
    \LR{$!A$} & \LR{$!B$} & \LR{$!A \land !B$} \\
    \hline \hline
    \LR{$\True$} & \LR{$\True$} & \LR{$\True$} \\
    \LR{$\True$} & \LR{$\False$} & \LR{$\False$} \\
    \LR{$\False$} & \LR{$\True$} & \LR{$\False$} \\
    \LR{$\False$} & \LR{$\False$} & \LR{$\False$} \\
    \hline
  \end{tabular}


  \begin{tabular}{|cc||c|} \hline
    \LR{$!A$} & \LR{$!B$} & \LR{$!A \lor !B$} \\
    \hline \hline
    \LR{$\True$} & \LR{$\True$} & \LR{$\True$} \\
    \LR{$\True$} & \LR{$\False$} & \LR{$\True$} \\
    \LR{$\False$} & \LR{$\True$} & \LR{$\True$} \\
    \LR{$\False$} & \LR{$\False$} & \LR{$\False$} \\
    \hline
  \end{tabular}
%
\\[1em]

  \begin{tabular}{|cc||c|} \hline
    \LR{$!A$} & \LR{$!B$} & \LR{$!A \lif !B$} \\
    \hline \hline
    \LR{$\True$} & \LR{$\True$} & \LR{$\True$} \\
    \LR{$\True$} & \LR{$\False$} & \LR{$\False$} \\
    \LR{$\False$} & \LR{$\True$} & \LR{$\True$} \\
    \LR{$\False$} & \LR{$\False$} & \LR{$\True$} \\
    \hline
  \end{tabular}


  \begin{tabular}{|cc||c|} \hline
    \LR{$!A$} & \LR{$!B$} & \LR{$!A \liff !B$} \\
    \hline \hline
    \LR{$\True$} & \LR{$\True$} & \LR{$\True$} \\
    \LR{$\True$} & \LR{$\False$} & \LR{$\False$} \\
    \LR{$\False$} & \LR{$\True$} & \LR{$\False$} \\
    \LR{$\False$} & \LR{$\False$} & \LR{$\True$} \\
    \hline
  \end{tabular}

\end{center}
\end{explain}
\olpseganchor{OLP-0061-B008}{end}

\olpseganchor{OLP-0061-B009}{start}
\begin{prob}
فرض کړئ چې \LR{$\Lang{L_0}$} ته يو درې ځايه نښلوونکے \LR{$\diamondsuit$}
ورزياتوو چې ارزونه ئې داسې ورکول شوې ده:
\begin{gather*}
  \pValue{v}(\diamondsuit( !A , !B , !C ) ) = \begin{cases}
    \pValue{v}( !B ) &
    \text{\RL{که \LR{$\pValue{v}(!A) = \True$} وي؛}}\\
    \pValue{v}( !C ) &
    \text{\RL{که \LR{$\pValue{v}(!A) = \False $} وي۔}}
  \end{cases}
\end{gather*}
د دې نښلوونکي د رښتياوالي جدول وليکئ۔
\end{prob}
\olpseganchor{OLP-0061-B009}{end}

\olpseganchor{OLP-0061-B010}{start}
\begin{thm}[ځايي ټاکنه]
\ollabel{thm:LocalDetermination} فرض کړئ \LR{$\pAssign{v_1}$} او
\LR{$\pAssign{v_2}$} داسې ارزښت ټاکنې دي چې په \LR{$!A$} کښې په راتلونکو
بياني متغيرونه يو له بل سره موافق دي؛ يعنې هر کله چې
\LR{$\Obj p_n$} په فارمول~\LR{$!A$} کښې راشي،
\LR{$\pAssign{v_1}(\Obj p_n) = \pAssign{v_2}(\Obj p_n)$} وي۔ نو
\LR{$\pValue{v_1}$} او \LR{$\pValue{v_2}$} هم پر~\LR{$!A$} موافق دي؛ يعنې
\LR{$\pValue{v_1}(!A) = \pValue{v_2}(!A)$}۔
\end{thm}
\olpseganchor{OLP-0061-B010}{end}


\olpseganchor{OLP-0061-B011}{start}
\begin{proof}
پر \LR{$!A$} جوړښتي استقرا کوو۔ که دا يو قضايايي متغير وي،
دواړه ټاکنې يې د فرض له مخې په يو ډول ارزوي۔ هر منطقي ثابت
هم له ټاکنې پرته يو ټاکلے ارزښت لري۔ د نفي په حالت کښې دواړه
ټاکنې د نفي په دننه فارمول کښې پر ټولو راغلو متغيرونو موافقې دي،
ځکه دا متغيرونه په ټول فارمول کښې هم راځي۔ د استقرايي فرض له مخې
د دننه فارمول ارزښتونه يو شان دي، نو د نفي ارزښتونه يې هم يو شان
دي۔ په هر دوه‌ځاييز حالت کښې هم د دواړو فرعي فارمولونو متغيرونه
په ټول فارمول کښې راځي۔ د استقرايي فرض له مخې اړوند ارزښتونه
دواړه يو شان دي، او د نښلوونکي د ارزونې قاعده له دغو ارزښتونو
هماغه يوه پايله ورکوي۔
\end{proof}
\olpseganchor{OLP-0061-B011}{end}

\olpseganchor{OLP-0061-B012}{start}
\begin{defn}[صدق]
\ollabel{defn:satisfaction} د \emph{فارمول~\LR{$!A$} صدق په
  ارزښت ټاکنه~\LR{$\pAssign{v}$} کښې}، يعنې \LR{$\pSat{v}{!A}$}، په
  استقرايي ډول داسې تعريفولے شو۔ (د ``\LR{$\pSat{v}{!A}$} نۀ دے'' د
  څرګندولو دپاره \LR{$\pSat/{v}{!A}$} ليکو۔)
\begin{enumerate}
\item %
  \indcase{!A}{\lfalse}{\LR{$\pSat/{v}{\indfrm}$}۔}
\olpseganchor{OLP-0061-B012}{end}

\olpseganchor{OLP-0061-B013}{start}
\item %
  \indcase{!A}{\ltrue}{\LR{$\pSat{v}{\indfrm}$}۔}
\olpseganchor{OLP-0061-B013}{end}

\olpseganchor{OLP-0061-B014}{start}
\item \indcase{!A}{\Obj p_i}{\LR{$\pSat{v}{\indfrm}$}
  هله او يوازې هله چې \LR{$\pAssign{v}(\Obj p_i) = \True$} وي۔}
\olpseganchor{OLP-0061-B014}{end}

\olpseganchor{OLP-0061-B015}{start}
\item %
  \indcase{!A}{\lnot !B}{\LR{$\pSat{v}{\indfrm}$} هله او يوازې هله چې
    \LR{$\pSat/{v}{!B}$} وي۔}
\olpseganchor{OLP-0061-B015}{end}

\olpseganchor{OLP-0061-B016}{start}
\item %
  \indcase{!A}{(!B \land !C)}{\LR{$\pSat{v}{\indfrm}$} هله او يوازې هله
    چې \LR{$\pSat{v}{!B}$} او \LR{$\pSat{v}{!C}$} دواړه وي۔}
\olpseganchor{OLP-0061-B016}{end}

\olpseganchor{OLP-0061-B017}{start}
\item %
  \indcase{!A}{(!B \lor !C)}{\LR{$\pSat{v}{\indfrm}$} هله او يوازې هله
    چې \LR{$\pSat{v}{!B}$} يا \LR{$\pSat{v}{!C}$}، يا دواړه، وي۔}
\olpseganchor{OLP-0061-B017}{end}

\olpseganchor{OLP-0061-B018}{start}
\item %
  \indcase{!A}{(!B \lif !C)}{\LR{$\pSat{v}{\indfrm}$} هله او يوازې هله
    چې \LR{$\pSat/{v}{!B}$} يا \LR{$\pSat{v}{!C}$}، يا دواړه، وي۔}
\olpseganchor{OLP-0061-B018}{end}

\olpseganchor{OLP-0061-B019}{start}
\item %
  \indcase{!A}{(!B \liff !C)}{\LR{$\pSat{v}{\indfrm}$} هله او يوازې هله
    چې يا \LR{$\pSat{v}{!B}$} او \LR{$\pSat{v}{!C}$} دواړه وي، يا نۀ
    \LR{$\pSat{v}{!B}$} وي او نۀ \LR{$\pSat{v}{!C}$}۔}
\end{enumerate}
که \LR{$\Gamma$} د فارمولونه يو سټ وي، نو \LR{$\pSat{v}{\Gamma}$} هله او
يوازې هله چې د هر~\LR{$!A \in \Gamma$} دپاره \LR{$\pSat{v}{!A}$} وي۔
\end{defn}
\olpseganchor{OLP-0061-B019}{end}

\olpseganchor{OLP-0061-B020}{start}
\begin{prop}\ollabel{prop:sat-value}
  \LR{$\pSat{v}{!A}$} هله او يوازې هله چې \LR{$\pValue{v}(!A) = \True$} وي۔
\end{prop}
\olpseganchor{OLP-0061-B020}{end}

\olpseganchor{OLP-0061-B021}{start}
\begin{proof}
  پر~\LR{$!A$} استقرا کوو۔
\end{proof}
\olpseganchor{OLP-0061-B021}{end}

\olpseganchor{OLP-0061-B022}{start}
\begin{prob}
  د \olref[pl][syn][val]{prop:sat-value} ثبوت ورکړئ۔
\end{prob}
\olpseganchor{OLP-0061-B022}{end}

\olpseganchor{OLP-0062-B004}{start}
\olfileid{pl}{syn}{sem}
\olpseganchor{OLP-0062-B004}{end}

\olpseganchor{OLP-0062-B005}{start}
\olsection{معنايي مفهومونه}
\olpseganchor{OLP-0062-B005}{end}

\olpseganchor{OLP-0062-B006}{start}
لاندې معنايي مفهومونه تعريفوو:
\olpseganchor{OLP-0062-B006}{end}

\olpseganchor{OLP-0062-B007}{start}
\begin{defn} 
\begin{enumerate}
\item فارمول~\LR{$!A$} هله \emph{د صدق وړ} دے چې د کوم
  \LR{$\pAssign{v}$} دپاره \LR{$\pSat{v}{!A}$} وي؛ او هله \emph{د صدق نۀ وړ}
  دے چې د هېڅ \LR{$\pAssign{v}$} دپاره \LR{$\pSat{v}{!A}$} نۀ وي؛
\item فارمول~\LR{$!A$} هله \emph{تاوتولوژي} ده چې د ټولو
  ارزښت ټاکنې~\LR{$\pAssign{v}$} دپاره \LR{$\pSat{v}{!A}$} وي؛
\item فارمول~\LR{$!A$} هله \emph{اتفاقي} دے چې د صدق وړ وي خو
  تاوتولوژي نۀ وي؛
\item که \LR{$\Gamma$} د فارمولونه يو سټ وي، نو \LR{$\Gamma \Entails !A$}
  (``\LR{$\Gamma$}، \LR{$!A$} ته استلزام ورکوي'') هله او يوازې هله چې د هر
  هغه ارزښت ټاکنه~\LR{$\pAssign{v}$} دپاره \LR{$\pSat{v}{!A}$} وي چې
  \LR{$\pSat{v}{\Gamma}$} وي۔
\item که \LR{$\Gamma$} د فارمولونه يو سټ وي، نو \LR{$\Gamma$} هله
  \emph{د صدق وړ} دے چې داسې ارزښت ټاکنه~\LR{$\pAssign{v}$} موجود وي
  چې \LR{$\pSat{v}{\Gamma}$} وي؛ او که داسې نۀ وي، \LR{$\Gamma$}
  \emph{د صدق نۀ وړ} دے۔
\end{enumerate} 
\end{defn}
\olpseganchor{OLP-0062-B007}{end}

\olpseganchor{OLP-0062-B008}{start}
\begin{prob} 
 د لاندې څلورو فارمولونه په هر يوه کښې وټاکئ چې آيا هغه
 (الف)~د صدق وړ، (ب)~تاوتولوژي او (ج)~اتفاقي دے۔
\begin{enumerate}
  \item \( ( \Obj p_0 \lif ( \lnot \Obj p_1 \lif \lnot \Obj p_0 ) ) \).
  \item \( ( ( \Obj p_0 \land \lnot \Obj p_1 ) \lif ( \lnot \Obj p_0 \land \Obj p_2 )) \liff ( ( \Obj p_2 \lif \Obj p_0 ) \lif ( \Obj p_0 \lif \Obj p_1 )) \).
  \item \( ( \Obj p_0 \liff \Obj p_1 ) \lif ( \Obj p_2 \liff \lnot \Obj p_1 ) \).
  \item \( (( \Obj p_0 \liff ( \lnot \Obj p_1 \land \Obj p_2 )) \lor ( \Obj p_2 \lif ( \Obj p_0 \liff \Obj p_1 ))) \).
\end{enumerate}
\end{prob}
\olpseganchor{OLP-0062-B008}{end}

\olpseganchor{OLP-0062-B009}{start}
\begin{prop}
\ollabel{prop:semanticalfacts} 
\begin{enumerate} 
\item \LR{$!A$} هله او يوازې هله تاوتولوژي ده چې
  \LR{$\emptyset \Entails !A$} وي؛ 
\item که \LR{$\Gamma \Entails !A$} او \LR{$\Gamma \Entails !A \lif !B$} وي،
  نو \LR{$\Gamma \Entails !B$}؛
\item که \LR{$\Gamma$} د صدق وړ وي، نو د \LR{$\Gamma$} هر متناهي فرعي سټ هم
  د صدق وړ دے؛ 
\item \ollabel{def:monotonicity} يو اړخيز زياتوالے: که
  \LR{$\Gamma \subseteq \Delta$} او \LR{$\Gamma \Entails !A$} وي، نو
  \LR{$\Delta \Entails !A$} هم؛
\item \ollabel{def:Cut} تعدي: که \LR{$\Gamma \Entails !A$} او
  \LR{$\Delta \cup \{ !A\} \Entails !B$} وي، نو
  \LR{$\Gamma \cup \Delta \Entails !B$}۔
\end{enumerate}
\end{prop}
\olpseganchor{OLP-0062-B009}{end}

\olpseganchor{OLP-0062-B010}{start}
\begin{proof}
تمرين۔
\end{proof}
\olpseganchor{OLP-0062-B010}{end}

\olpseganchor{OLP-0062-B011}{start}
\begin{prob}
د \olref[pl][syn][sem]{prop:semanticalfacts} ثبوت ورکړئ۔
\end{prob}
\olpseganchor{OLP-0062-B011}{end}

\olpseganchor{OLP-0062-B012}{start}
\begin{prop}\ollabel{prop:entails-unsat}
  \LR{$\Gamma \Entails !A$} هله او يوازې هله چې
  \LR{$\Gamma \cup \{\lnot !A\}$} د صدق نۀ وړ وي۔
\end{prop}
\olpseganchor{OLP-0062-B012}{end}

\olpseganchor{OLP-0062-B013}{start}
\begin{proof}
تمرين۔
\end{proof}
\olpseganchor{OLP-0062-B013}{end}

\olpseganchor{OLP-0062-B014}{start}
\begin{prob}
د \olref[pl][syn][sem]{prop:entails-unsat} ثبوت ورکړئ۔
\end{prob}
\olpseganchor{OLP-0062-B014}{end}

\olpseganchor{OLP-0062-B015}{start}
\begin{thm}[د معنايي استنتاج قضيه]
  \ollabel{thm:sem-deduction} \LR{$\Gamma \Entails !A \lif !B$} هله او
  يوازې هله چې \LR{$\Gamma \cup \{!A\} \Entails !B$} وي۔
\end{thm}
\olpseganchor{OLP-0062-B015}{end}

\olpseganchor{OLP-0062-B016}{start}
\begin{proof}
تمرين۔
\end{proof}
\olpseganchor{OLP-0062-B016}{end}

\olpseganchor{OLP-0062-B017}{start}
\begin{prob}
د \olref[pl][syn][sem]{thm:sem-deduction} ثبوت ورکړئ۔
\end{prob}
\olpseganchor{OLP-0062-B017}{end}

\olpseganchor{OLP-0715-B004}{start}
\olfileid{pl}{prp}{snd}
\olpseganchor{OLP-0715-B004}{end}

\olpseganchor{OLP-0715-B005}{start}
\olsection{د قضيوي منطق سموالے}
\olpseganchor{OLP-0715-B005}{end}

\olpseganchor{OLP-0715-B006}{start}
\begin{thm}[سموالے]
\ollabel{thm:soundness}
که \LR{$\Gamma\Proves!A$} نو \LR{$\Gamma \Entails !A$}۔
\end{thm}
\olpseganchor{OLP-0715-B006}{end}

\olpseganchor{OLP-0715-B007}{start}
\begin{proof}
د اشتقاق پر جوړښت استقرا کوو۔ که \LR{$!A$} بديهي اصل وي، نو \LR{$\Entails
!A$} او له دې څخه \LR{$\Gamma \Entails!A$} هم راځي۔ همداسې که \LR{$!A \in \Gamma$}، نو \LR{$\Gamma \Entails!A$}۔ په استقرايي ګام کښې فرض کړئ \LR{$!A$} له \LR{$!C$} او \LR{$!C\lif!A$} څخه په \emph{مودوس پوننس} اخيستل شوے دے۔ د استقرا د فرضيې له مخې قضيه د \LR{$!C\lif !A$} او \LR{$!C$} دپاره سمه ده۔ د \olref[pl][syn][sem]{prop:semanticalfacts} دويم توکے مطلوبه پايله ورکوي۔
\end{proof}
\olpseganchor{OLP-0715-B007}{end}

\olpseganchor{OLP-0715-B008}{start}
\begin{cor}
که \LR{$\Gamma$} د صدق وړ وي، نو \LR{$\Gamma$} سازګار دے۔ په نتيجه کښې قضيوي منطق سازګار دے۔
\end{cor}
\olpseganchor{OLP-0715-B008}{end}

\olpseganchor{OLP-0714-B004}{start}
\olfileid{pl}{prp}{com}
\olpseganchor{OLP-0714-B004}{end}

\olpseganchor{OLP-0714-B005}{start}
\olsection{د قضيوي منطق بشپړتيا}
\olpseganchor{OLP-0714-B005}{end}

\olpseganchor{OLP-0714-B006}{start}
\begin{defn}
\ollabel{def:MaxCon}
د فارمولونه سټ \LR{$\Gamma$} هغه وخت \emph{اعظمي سازګار} دے چې پخپله سازګار وي او، که \LR{$\Delta$} داسې سازګار سټ وي چې \LR{$\Gamma
\subseteq \Delta$}، نو \LR{$\Gamma = \Delta$}۔
\end{defn}
\olpseganchor{OLP-0714-B006}{end}

\olpseganchor{OLP-0714-B007}{start}
\begin{lem}[د رښتينولۍ لمه]
  \ollabel{lem:truth}
فرض کړئ \LR{$\Gamma$} اعظمي سازګار دے۔ بيا:
\begin{enumerate}
\item \ollabel{lem:truth-1} \LR{$!A \in \Gamma$} هغه او يوازې هغه وخت چې \LR{$\lnot !A
  \notin \Gamma$}؛
\item \ollabel{lem:truth-2} \LR{$\Gamma \Proves !A$} هغه او يوازې هغه وخت چې \LR{$!A \in \Gamma$}؛
\item \ollabel{lem:truth-3} \LR{$!A \lif !B \in \Gamma$} هغه او يوازې هغه وخت چې \LR{$!A \notin \Gamma$} يا \LR{$!B \in \Gamma$}۔
\end{enumerate}
\end{lem}
\olpseganchor{OLP-0714-B007}{end}

\olpseganchor{OLP-0714-B008}{start}
\begin{proof}
\begin{enumerate}
\item فرض کړئ \LR{$!A \in \Gamma$}۔ که \LR{$\lnot !A \in \Gamma$} هم وي، \LR{$\Gamma$} ناسازګار دے، نو دواړه غړي يو ځاے کېدے نۀ شي۔ که نۀ \LR{$!A$} او نۀ \LR{$\lnot!A$} په \LR{$\Gamma$} کښې وي، د \olref[fol][axd][prv]{prop:provability-exhaustive} له مخې له \LR{$\Gamma\cup\{!A\}$} او \LR{$\Gamma \cup \{\lnot !A\}$} څخه لږ تر لږه يو سازګار دے۔ دا د \LR{$\Gamma$} له اعظمي والي سره ټکر لري۔ نو له دغو دوو فورمولونو څخه کټ مټ يو په سټ کښې دے او دواړه لوري ثابت شول۔
\item که \LR{$!A \in \Gamma$}، نو ښکاره ده چې \LR{$\Gamma \Proves !A$}۔ اوس فرض کړئ \LR{$\Gamma \Proves !A$}۔ که \LR{$!A \notin \Gamma$}، د \olref{lem:truth-1} له مخې \LR{$\lnot !A \in \Gamma$}۔ بيا \LR{$\Gamma \Proves
  !A$} او \LR{$\Gamma \Proves \lnot !A$}، چې د سازګارۍ خلاف دي۔
\item تمرين۔
\end{enumerate}
\end{proof}
\olpseganchor{OLP-0714-B008}{end}

\olpseganchor{OLP-0714-B009}{start}
د \LR{\olref{lem:truth}(\olref{lem:truth-2})} هغه لورے چې له ثابتېدنې غړيتوب اخلي، ښيي چې \LR{$\Gamma$} \emph{له اشتقاقي پلوه تړلے} دے؛ يعنې که \LR{$\Gamma \Proves !A$} نو \LR{$!A
\in \Gamma$}، د هر \LR{$!A$} دپاره۔
\olpseganchor{OLP-0714-B009}{end}

\olpseganchor{OLP-0714-B010}{start}
\begin{prob}
د \olref{lem:truth} ثبوت بشپړ کړئ۔
\end{prob}
\olpseganchor{OLP-0714-B010}{end}

\olpseganchor{OLP-0714-B011}{start}
\begin{thm}[بشپړتيا]
\ollabel{thm:completeness}
که \LR{$\Gamma$} سازګار وي، نو \LR{$\Gamma$} د صدق وړ دے۔
\end{thm}
\olpseganchor{OLP-0714-B011}{end}

\olpseganchor{OLP-0714-B012}{start}
\begin{proof}
فرض کړئ \LR{$!A_0$}، \LR{$!A_1$}،~\dots د ژبې د ټولو فارمولونه بشپړ نوملړ دے۔ د فارمولونه مخ په زياتېدو سټونه \LR{$\Gamma_0$}، \LR{$\Gamma_1$}،~\dots په بازګښتي ډول داسې ټاکو:
\begin{align*}
\Gamma_0 & = \Gamma\\
\Gamma_{n+1} &  =
\begin{cases}
  \Gamma_n \cup \{!A_n\} & \text{\RL{که \LR{$\Gamma_n \cup \{!A_n\}$} سازګار وي؛}}\\
  \Gamma_n \cup \{\lnot !A_n\} & \text{\RL{که داسې نۀ وي}}.
\end{cases}
\end{align*}
بيا داسې تعريفوو:
\[
\Gamma^* = \bigcup_{0\le n}\Gamma_n.
\]
لاندې خبرې په خپل ترتيب ثابتوو:
\begin{enumerate}
\item د هر \LR{$n$} دپاره سټ \LR{$\Gamma_n$} سازګار دے؛ د \LR{$n$} پر سر استقرا او \olref[fol][axd][prv]{prop:provability-exhaustive} وکاروئ۔
\item \LR{$\Gamma^*$} سازګار دے؛ که تناقض پکې اشتقاقېدے، د اشتقاق متناهي فرضونه به ټول د سلسلې په يوه پړاو کښې وو، چې د هغه پړاو له سازګارۍ سره ټکر لري۔
\item \LR{$\Gamma^*$} اعظمي دے؛ د نوملړ د هر غړي په خپل پړاو کښې يا هغه غړے يا ئې نفي شامله شوې ده۔ د يوځاے شوي سټ په هره ريښتينې پراختيا کښې به يو نوے فورمول له خپلې لا موجودې نفي سره يو ځاے شي۔
\end{enumerate}
اوس ارزښت‌ټاکنه \LR{$v$} داسې ټاکو چې \LR{$v(p_i) = \True$} هغه او يوازې هغه وخت وي چې \LR{$p_i \in \Gamma^*$}۔ پر \LR{$!A$} جوړښتي استقرا ښيي چې په \LR{$\Gamma^*$} کښې غړيتوب د \LR{$v$} له مخې له رښتينولۍ سره برابر دے۔ د نفي او شرطي تړنې حالتونه د رښتينولۍ د لمې د غړيتوب له شرطونو اخلو۔ د او او يا حالتونه په ترتيب له \olref[fol][com][ccs]{prop:ccs-and} او \olref[fol][com][ccs]{prop:ccs-or} او د ارزښت‌ټاکنې له تعريفه اخلو۔ د دروغ ثابت په سازګار سټ کښې نشته او د رښتيا ثابت د \olref[fol][axd][prp]{ax:ltrue} او اشتقاقي تړلتيا له مخې پکې شته۔ په دې بديهي نظام کښې دوه‌اړخيز شرط د دوو متقابلو شرطونو د او تړنه ده۔ نو د پېچلو تړلي فارمولونه دپاره هم \LR{$\overline{v}(!A) = \True$} هغه او يوازې هغه وخت چې \LR{$!A \in \Gamma^*$}۔ په ځانګړي ډول \LR{$v$}، \LR{$\Gamma^*$} پوره کوي؛ او چې \LR{$\Gamma \subseteq \Gamma^*$}، نو \LR{$\Gamma$} هم د صدق وړ دے۔
\end{proof}
\olpseganchor{OLP-0714-B012}{end}

\olpseganchor{OLP-0714-B013}{start}
\begin{cor}
که \LR{$\Gamma \Entails !A$} نو \LR{$\Gamma \Proves !A$}۔
\end{cor}
\olpseganchor{OLP-0714-B013}{end}

\olpseganchor{OLP-0714-B014}{start}
\begin{proof}
که \LR{$\Gamma\Proves/ !A$}، د \olref[fol][axd][prv]{prop:prov-incons} له مخې \LR{$\Gamma \cup \{\lnot !A\}$} سازګار دے۔ د بشپړتيا د قضيې له مخې \LR{$\Gamma \cup
\{\lnot !A\}$} د صدق وړ دے، نو \LR{$\Gamma \Entails/ !A$}۔
\end{proof}
\olpseganchor{OLP-0714-B014}{end}

\olpseganchor{OLP-0714-B015}{start}
\begin{prop}[د تراکم قضيه]
\ollabel{prop:compactness}
\LR{$\Gamma$} هغه او يوازې هغه وخت د صدق وړ دے چې د \LR{$\Gamma$} هر \emph{متناهي} فرعي سټ \LR{$\Gamma_0$} د صدق وړ وي۔
\end{prop}
\olpseganchor{OLP-0714-B015}{end}

\olpseganchor{OLP-0714-B016}{start}
\begin{proof}
\LR{$\Gamma$} هغه او يوازې هغه وخت د صدق نۀ وړ دے چې ناسازګار وي۔ دا بيا هغه او يوازې هغه وخت وي چې د \LR{$\Gamma$} کوم متناهي فرعي سټ \LR{$\Gamma_0$} ناسازګار وي، ځکه اشتقاق يوازې متناهي فرضونه کاروي۔ په پای کښې همدا هغه او يوازې هغه وخت وي چې د \LR{$\Gamma$} کوم متناهي فرعي سټ \LR{$\Gamma_0$} د صدق نۀ وړ وي۔
\end{proof}
\olpseganchor{OLP-0714-B016}{end}

\olpseganchor{OLP-0714-B017}{start}
\begin{cor}
\LR{$\Gamma \Entails !A$} هغه او يوازې هغه وخت چې د \LR{$\Gamma$} د کوم متناهي فرعي سټ \LR{$\Gamma_0$} دپاره \LR{$\Gamma_0 \Entails !A$} وي۔
\end{cor}
\olpseganchor{OLP-0714-B017}{end}

\olpseganchor{OLP-0063-B004}{start}
\olchapter{fol}{prf}{د اشتقاق نظامونه}
\olpseganchor{OLP-0063-B004}{end}

\olpseganchor{OLP-0063-B005}{start}
\begin{editorial}
دا څپرکے د اشتقاق د نظامونو عمومي مواد راغونډوي۔ کوم درسي
کتاب چې يو ځانګړے نظام کاروي، د هغه نظام د معرفي کوونکي څپرکي په
پيل کښې د پېژندنې برخه او د اړوندې ټولکتنې برخه ورګډولے شي۔
\end{editorial}
\olpseganchor{OLP-0063-B005}{end}

\olpseganchor{OLP-0064-B004}{start}
\olfileid{fol}{prf}{int}
\olpseganchor{OLP-0064-B004}{end}

\olpseganchor{OLP-0064-B005}{start}
\olsection{پېژندنه}
\olpseganchor{OLP-0064-B005}{end}

\olpseganchor{OLP-0064-B006}{start}
منطقونه عموماً معنٰی پوهنه او اشتقاق نظام دواړه لري۔ معنٰی
پوهنه د رښتياوالي، د صدق وړتيا، اعتبار او استلزام په څېر مفهومونو
سره کار لري۔ د اشتقاق د نظامونو موخه دا ده چې د استلزام او
اعتبار د ثابتولو يوه بشپړه نحوي طريقه ورکړي۔ دا نظامونه په دې معنٰی
بشپړ نحوي دي چې په داسې يوه نظام کښې اشتقاق يو متناهي نحوي
څيز وي، چې عموماً د تړلي فارمولونه يا فارمولونه يوه لړۍ يا بل
متناهي ترتيب وي۔ د اشتقاق ښۀ نظامونه دا خاصيت لري چې د
تړلي فارمولونه يا فارمولونه د هرې ورکړل شوې لړۍ يا ترتيب
``سموالے'' په ميخانيکي توګه کتل کېدے شي۔
\olpseganchor{OLP-0064-B006}{end}

\olpseganchor{OLP-0064-B007}{start}
د لومړۍ درجې منطق تر ټولو ساده، او له تاريخي پلوه لومړني،
اشتقاق نظامونه \emph{بديهي} وو۔ په داسې يوه نظام کښې د
فارمولونه يوه لړۍ هله اشتقاق ګڼل کېږي چې په هغې کښې هر
فارمول يا د ``بديهي اصولو'' په يوه ټاکلي سټ کښې وي، يا د ټاکلو
``استنتاجي قاعدو'' له يوې څخه په کار اخيستو په لړۍ کښې له مخکښو
فارمولونه څخه ترې راووځي؛ او په ميخانيکي توګه کتل کېدے شي چې آيا
فارمول بديهي اصل دے او آيا د استنتاج له کومې قاعدې سره سم له
نورو فارمولونه څخه راوتلے دے۔ د اشتقاق بديهي نظامونه
تشريح کول اسانه دي، او له مابعدمنطقي پلوه پرې کار کول هم اسانه دي،
خو په هغو کښې اشتقاقونه لوستل او پوهېدل ګران دي او جوړول ئې هم
ګران دي۔
\olpseganchor{OLP-0064-B007}{end}

\olpseganchor{OLP-0064-B008}{start}
د اشتقاق نور نظامونه د دې موخې دپاره جوړ شوي چې د
اشتقاقونه جوړول يا له بشپړېدو وروسته پر اشتقاقونه
پوهېدل اسانه کړي۔
بېلګې ئې طبيعي استنتاج، د رښتياوالي ونې چې د تابلو ثبوتونو په نامه
هم پېژندل کېږي، او د سېکوېنټ حساب دي۔ د اشتقاق ځينې نظامونه
په ځانګړې توګه د ميخانيکي کولو دپاره طرح شوي دي؛ مثلاً د حل طريقه
په سافټوير کښې په اسانۍ پلې کېږي، خو پر اشتقاقونه ئې پوهېدل
تقريباً ناشوني دي۔ د اشتقاق دا نور نظامونه زياتره
اشتقاقونه د لړيو پر ځاے د فارمولونه د ونو په توګه ښيي۔ په
دې سره په اسانۍ ليدل کېږي چې د اشتقاق کومې برخې پر کومو
نورو برخو تکيه لري۔
\olpseganchor{OLP-0064-B008}{end}

\olpseganchor{OLP-0064-B009}{start}
نو د يوه ورکړل شوي منطق، لکه د لومړۍ درجې منطق، د اشتقاق
بېل نظامونه به دا په بېلو ډولونو روښانه کړي چې تړلے فارمول کله
\emph{قضيه} ده او د نورو له ځينو څخه د تړلے فارمول د
د اشتقاق وړ کېدو معنٰی څه ده۔ دا که په بديهي اشتقاقونه،
طبيعي استنتاجونو، سېکوېنټي اشتقاقونه، د رښتياوالي په ونو يا د
حل په ابطالونو ترسره شي، غواړو چې دا اړيکې د اعتبار او استلزام له
معنايي مفهومونو سره سمون ولري۔ د ``\LR{$!A$}~قضيه ده'' دپاره
\LR{$\Proves !A$} او د ``\LR{$!A$} له~\LR{$\Gamma$} څخه د اشتقاق وړ دے'' دپاره
``\LR{$\Gamma \Proves !A$}'' ليکو۔ \LR{$\Proves$} چې هر ډول تعريف شي، غواړو
له \LR{$\Entails$} سره سمون ولري؛ يعنې:
\begin{enumerate}
\item \LR{$\Proves !A$} هله او يوازې هله چې \LR{$\Entails !A$}
\item \LR{$\Gamma \Proves !A$} هله او يوازې هله چې \LR{$\Gamma \Entails !A$}
\end{enumerate}
د پورته دواړو د ``يوازې هله'' لوري ته \emph{سموالے} ويل کېږي۔
اشتقاق نظام هله سم دے چې د اشتقاق وړتيا استلزام يا اعتبار
تضمين کړي۔ د اشتقاق هر د کار وړ نظام بايد سم وي؛ د
اشتقاق ناسم نظامونه هېڅ ګټه نۀ لري۔ په اصل کښې د
اشتقاق ټوله موخه دا ده چې د اعتبار يا استلزام نحوي تضمين
ورکړي۔ موږ به د وړاندې شويو اشتقاق نظامونو سموالے ثابت کړو۔
\olpseganchor{OLP-0064-B009}{end}

\olpseganchor{OLP-0064-B010}{start}
د دې عکس، يعنې ``هله'' لورے، هم مهم دے: دې ته \emph{بشپړتيا} ويل
کېږي۔ د اشتقاق بشپړ نظام دومره ځواکمن وي چې هر کله \LR{$!A$}
معتبر وي، وښيي چې \LR{$!A$} قضيه ده؛ او هر کله چې
\LR{$\Gamma \Entails !A$} وي، وښيي چې \LR{$\Gamma \Proves !A$}۔ بشپړتيا ثابتول
ګران دي، او ځينې منطقونه د اشتقاق هېڅ بشپړ نظام نۀ لري۔ د
لومړۍ درجې منطق ئې لري۔ کورت ګوډل لومړے کس و چې د لومړۍ درجې منطق د
اشتقاق د يوه نظام بشپړتيا ئې په خپله ۱۹۲۹ز کال څېړنيزه
رساله کښې ثابته کړه۔
\olpseganchor{OLP-0064-B010}{end}

\olpseganchor{OLP-0064-B011}{start}
له اشتقاق نظامونو سره يو بل تړلے مفهوم \emph{سازګاري} ده۔
د تړلي فارمولونه يو سټ هله ناسازګار بلل کېږي چې هر څه ترې
اشتقاق کېدے شي؛ او که داسې نۀ وي، سازګار دے۔ ناسازګاري د صدق
د نۀ وړتيا نحوي سياله ده: لکه د صدق نۀ وړ سټونو غوندې، د
تړلي فارمولونه ناسازګار سټونه ښې نظريې نۀ جوړوي او په بنسټيز ډول
نيمګړي وي۔ د تړلي فارمولونه سازګار سټونه کېدے شي رښتوني يا ګټور نۀ
وي، خو لږ تر لږه د منطقي ګټورتيا له دې تر ټولو ټيټ بريد څخه تېرېږي۔
د اشتقاق په بېلو نظامونو کښې د تړلي فارمولونه د سټونو د
سازګارۍ ځانګړے تعريف توپير کولے شي؛ خو لکه~\LR{$\Proves$}، غواړو
سازګاري د خپل معنايي سيال، د صدق وړتيا، سره سمه وي۔ غواړو تل داسې
وي چې \LR{$\Gamma$} هله او يوازې هله سازګار وي چې د صدق وړ وي۔ دلته د
``يوازې هله'' لورے له بشپړتيا سره برابر دے، ځکه سازګاري د صدق
وړتيا تضمينوي؛ او ``هله'' لورے له سموالي سره برابر دے، ځکه د صدق
وړتيا سازګاري تضمينوي۔ په حقيقت کښې د کلاسيکي لومړۍ درجې منطق
دپاره د سموالي او بشپړتيا دواړه بڼې همارزښته دي۔
\olpseganchor{OLP-0064-B011}{end}

\olpseganchor{OLP-0065-B004}{start}
\olfileid{fol}{prf}{seq}
\olpseganchor{OLP-0065-B004}{end}

\olpseganchor{OLP-0065-B005}{start}
\olsection{د سېکوېنټ حساب}
\olpseganchor{OLP-0065-B005}{end}

\olpseganchor{OLP-0065-B006}{start}
د اشتقاق ډېر نظامونه د تړلي فارمولونه له ترتيبونو سره کار
کوي، خو د سېکوېنټ حساب له \emph{سېکوېنټونو} سره کار کوي۔ سېکوېنټ
د دې بڼې يو عبارت دے:
\[
!A_1, \dots, !A_m \Sequent !B_1, \dots, !B_n,
\]
دواړه لړۍ متناهي دي او اوږدوالے يې په خپلواکه توګه ټاکل کېږي۔
سېکوېنټ د تړلي فارمولونه د لړيو يوه جوړه ده چې د سېکوېنټ په
نښه~\LR{$\Sequent$} بېله شوې ده۔ له دواړو لړيو هره يوه تشه کېدے شي۔ د
سېکوېنټ په حساب کښې اشتقاق د سېکوېنټونو يوه ونه ده۔ تر
ټولو پاسني سېکوېنټونه ځانګړې بڼه لري، او ``لومړني سېکوېنټونه'' يا
``بديهي اصول'' بلل کېږي؛ هر بل سېکوېنټ د استنتاج د کومې قاعدې له
لارې سم له پاسه سېکوېنټونو څخه راځي۔ د استنتاج قاعدې يا په
سېکوېنټونو کښې تړلي فارمولونه سمبالوي، يعنې کيڼې يا ښۍ خوا ته ئې
ورزياتوي، ترې لرې کوي يا ئې بيا ترتيبوي؛ او يا د قاعدې په پايله کښې
يو مرکب فارمول معرفي کوي۔ مثلاً د \LR{$\LeftR{\land}$} قاعده اجازه
راکوي چې له \LR{$!A, \Gamma \Sequent \Delta$} څخه \LR{$!A \land !B, \Gamma
\Sequent \Delta$} استنتاج کړو، او \LR{$\RightR{\lif}$} اجازه راکوي چې د
هر \LR{$\Gamma$}، \LR{$\Delta$}، \LR{$!A$} او~\LR{$!B$} دپاره له \LR{$!A, \Gamma \Sequent
\Delta, !B$} څخه \LR{$\Gamma \Sequent \Delta, !A \lif !B$} استنتاج کړو۔
(په تېره بيا، \LR{$\Gamma$} او~\LR{$\Delta$} تش کېدے شي۔)
\olpseganchor{OLP-0065-B006}{end}

\olpseganchor{OLP-0065-B007}{start}
د سېکوېنټ پر حساب ولاړه \LR{$\Proves$} اړيکه داسې تعريفېږي:
\LR{$\Gamma \Proves !A$} هله او يوازې هله چې داسې کومه لړۍ \LR{$\Gamma_0$}
موجوده وي چې په \LR{$\Gamma_0$} کښې هر \LR{$!A$} په~\LR{$\Gamma$} کښې وي، او داسې
اشتقاق موجود وي چې سېکوېنټ~\LR{$\Gamma_0 \Sequent !A$} ئې په
جرړه کښې وي۔ \LR{$!A$} هله د سېکوېنټ په حساب کښې قضيه ده چې سېکوېنټ
~\LR{$\Sequent !A$} کوم اشتقاق ولري۔ مثلاً دا اشتقاق
ښيي چې \LR{$\Proves (!A \land !B) \lif !A$}:
\begin{prooftree}
  \Axiom$!A \fCenter !A$
  \RightLabel{\LeftR{\land}}
  \UnaryInf$!A \land !B \fCenter !A$
  \RightLabel{\RightR{\lif}}
  \UnaryInf$\fCenter (!A \land !B) \lif !A$
\end{prooftree}
\olpseganchor{OLP-0065-B007}{end}

\olpseganchor{OLP-0065-B008}{start}
يو سټ \LR{$\Gamma$} د سېکوېنټ په حساب کښې هله ناسازګار دے چې د
\LR{$\Gamma_0 \Sequent$} کوم اشتقاق موجود وي، چې په کښې هر
\LR{$!A \in \Gamma_0$} په~\LR{$\Gamma$} کښې وي او د سېکوېنټ ښۍ خوا تشه وي۔
د \RightR{\Weakening} قاعدې په کارولو له ناسازګار سټ څخه هر
تړلے فارمول اشتقاق کېدے شي۔
\olpseganchor{OLP-0065-B008}{end}

\olpseganchor{OLP-0065-B009}{start}
د سېکوېنټ حساب په ۱۹۳۰مو کلونو کښې ګيرهارډ ګېنتڅن جوړ کړ۔ د منظمې
او متناظرې طرحې له امله دا د اشتقاقونه د نظريې د ودې دپاره
ډېر ګټور صوري نظام دے۔ د سېکوېنټ په حساب کښې د اشتقاقونه
موندل نسبتاً اسانه دي، خو دا اشتقاقونه ډېر ځله په سختۍ لوستل
کېږي او له ثبوتونو سره ئې تړاو کله کله په اسانۍ نۀ ښکاري۔ بيا هم،
دا د اشتقاق نظامونو ډېره ښکلې لار ثابته شوې، او ډېر منطقونه
د سېکوېنټ حساب نظامونه لري۔
\olpseganchor{OLP-0065-B009}{end}

\olpseganchor{OLP-0066-B004}{start}
\olfileid{fol}{prf}{ntd}
\olpseganchor{OLP-0066-B004}{end}

\olpseganchor{OLP-0066-B005}{start}
\olsection{طبيعي استنتاج}
\olpseganchor{OLP-0066-B005}{end}

\olpseganchor{OLP-0066-B006}{start}
طبيعي استنتاج د اشتقاق داسې نظام دے چې موخه ئې د رښتيني
استدلال، په تېره بيا د رياضي پوهانو د منظم استدلال، انځورول دي۔
رښتينے استدلال د يو شمېر ``طبيعي'' بڼو له مخې مخکښې ځي۔ مثلاً د
حالتونو ثبوت موږ ته اجازه راکوي چې د يوه فصلي فرض پر بنسټ پايله
ثابته کړو، په دې چې وښيو پايله د فصل له هرې برخې څخه راځي۔ نامستقيم
ثبوت اجازه راکوي چې پايله په دې سره ثابته کړو چې نفي ئې تناقض ته
رسېږي۔ شرطي ثبوت د ``که \dots نو \dots'' شرطي ادعا په دې سره ثابتوي
چې تالي له مقدم څخه راځي۔ طبيعي استنتاج د دغو طبيعي استنتاجونو د
ځينو صوري کول دي۔ هر منطقي نښلوونکے او هر سور دوه قاعدې لري، يوه د
معرفي کولو او بله د لرې کولو قاعده، او هره يوه ئې د داسې طبيعي
استنتاج له يوې بڼې سره سمه ده۔ مثلاً \LR{$\Intro{\lif}$} له شرطي ثبوت
سره او \LR{$\Elim{\lor}$} د حالتونو له ثبوت سره سمون لري۔ يوه په تېره
ساده قاعده \LR{$\Elim{\land}$} ده، چې اجازه راکوي له \LR{$!A \land !B$} څخه
~\LR{$!A$} يا \LR{$!B$} استنتاج کړو۔
\olpseganchor{OLP-0066-B006}{end}

\olpseganchor{OLP-0066-B007}{start}
يو خاصيت چې طبيعي استنتاج د اشتقاق له نورو نظامونو څخه
بېلوي، د فرضونو کارول دي۔ په طبيعي استنتاج کښې اشتقاق د
فارمولونه يوه ونه ده۔ يو فارمول د فارمولونه د ونې په جرړه
کښې وي، او د ونې ``پاڼې'' هغه فارمولونه وي چې پايله ترې راايستل
کېږي۔ په طبيعي استنتاج کښې ځينې پاڼه-فارمولونه د اشتقاق
دننه ونډه لري، خو تر هغې ``کارول شوي وي'' چې اشتقاق پايلې ته
رسېږي۔ دا په رښتيني استدلال کښې د هغو فرضيو د معرفي کولو له دود سره
سمون لري چې يوازې لږ وخت نافذې پاتې کېږي۔ مثلاً د حالتونو په ثبوت
کښې د فصل د هرې برخې رښتياوالے فرضوو؛ په شرطي ثبوت کښې د مقدم
رښتياوالے فرضوو؛ او په نامستقيم ثبوت کښې د پايلې د نفي رښتياوالے
فرضوو۔ د فرضي فرضونو د معرفي کولو او بيا د يوه منځني ګام د ثابتولو
دپاره د هغو ختمولو دا طريقه د طبيعي استنتاج څرګنده نښه ده۔ د طبيعي
استنتاج د يوه اشتقاق په پاڼو کښې فارمولونه فرضونه بلل کېږي،
او د استنتاج ځينې قاعدې کېدے شي هغه ``ايستل'' کړي۔ مثلاً که
مو له ځينو فرضونو څخه د~\LR{$!B$} کوم اشتقاق درلود چې~\LR{$!A$} هم
پکې شامل و، نو د \LR{$\Intro{\lif}$} قاعده اجازه راکوي چې~\LR{$!A \lif !B$}
استنتاج کړو او د~\LR{$!A$} بڼې هر فرض ختم کړو۔ (د دې ثبتولو
دپاره چې کوم فرضونه په کومو استنتاجونو کښې ختمېږي، استنتاج او هغه
فرضونه چې ختموي ئې په يوه شمېره نښه کوو۔) هغه فرضونه چې د
اشتقاق په پاے کښې نۀ ايستل شوې پاتې کېږي، په ګډه د پايلې
د رښتياوالي دپاره بس دي؛ نو اشتقاق ثابتوي چې د هغه
نۀ ايستل شوي فرضونه پايلې ته استلزام ورکوي۔
\olpseganchor{OLP-0066-B007}{end}

\olpseganchor{OLP-0066-B008}{start}
د طبيعي استنتاج پر بنسټ اړيکه \LR{$\Gamma \Proves !A$} هله صدق کوي چې
داسې اشتقاق موجود وي چې \LR{$!A$} د ونې وروستۍ تړلے فارمول وي،
او هره نۀ ايستل شوې پاڼه په~\LR{$\Gamma$} کښې وي۔ \LR{$!A$} په طبيعي
استنتاج کښې هله قضيه ده چې داسې اشتقاق موجود وي چې \LR{$!A$}
وروستۍ تړلے فارمول وي او ټول فرضونه ايستل شوي وي۔ مثلاً دا
اشتقاق ښيي چې \LR{$\Proves (!A \land !B) \lif !A$}:
\begin{prooftree}
  \AxiomC{\LR{$\Discharge{!A \land !B}{1}$}}
  \RightLabel{\Elim{\land}}
  \UnaryInfC{\LR{$!A$}}
  \DischargeRule{\Intro{\lif}}{1}
  \UnaryInfC{\LR{$(!A \land !B) \lif !A$}}
\end{prooftree}
نښه~\LR{$1$} څرګندوي چې فرض \LR{$!A \land !B$} د \Intro{\lif} په استنتاج کښې
ايستل کېږي۔
\olpseganchor{OLP-0066-B008}{end}
  
\olpseganchor{OLP-0066-B009}{start}
يو سټ~\LR{$\Gamma$} هله ناسازګار دے چې په طبيعي استنتاج کښې
\LR{$\Gamma \Proves \lfalse$} وي۔ د \FalseInt{} قاعده دا يقيني کوي چې له
ناسازګار سټ څخه هر تړلے فارمول اشتقاق کېدے شي۔
\olpseganchor{OLP-0066-B009}{end}

\olpseganchor{OLP-0066-B010}{start}
د طبيعي استنتاج نظامونه ګيرهارډ ګېنتڅن او ستانيسواف ياشکوفسکي په
۱۹۳۰مو کلونو کښې جوړ کړل، او وروسته ډاګ پراوېڅ او فريډرېک فېچ لا
پراخ کړل۔ ځکه چې استنتاجونه ئې د ثبوت طبيعي لارې انځوروي، فلسفيان
ئې خوښوي۔ د فېچ جوړې شوې بڼې زياتره د منطق په پېژندنيزو درسي
کتابونو کښې کارول کېږي۔ د منطق په فلسفه کښې کله کله د طبيعي استنتاج
قاعدې د منطقي عملګرو معنٰی ورکوي، چې دې ته ``ثبوت-نظري معنايي
پوهنه'' ويل کېږي۔
\olpseganchor{OLP-0066-B010}{end}

\olpseganchor{OLP-0067-B004}{start}
\olfileid{fol}{prf}{tab}
\olpseganchor{OLP-0067-B004}{end}

\olpseganchor{OLP-0067-B005}{start}
\olsection{تابلوګانې}
\olpseganchor{OLP-0067-B005}{end}

\olpseganchor{OLP-0067-B006}{start}
د اشتقاق ډېر نظامونه د تړلي فارمولونه له ترتيبونو سره کار
کوي، خو تابلوګانې له نښه لرونکي فارمولونه سره کار کوي۔
نښه لرونکے فارمول يوه جوړه ده چې د رښتياوالي له يوې نښې
(\LR{$\True$} يا \LR{$\False$}) او تړلے فارمول څخه جوړه ده:
\[
\sFmla{\True}{!A} \text{\RL{ يا }} \sFmla{\False}{!A}.
\]
تابلو له هغو نښه لرونکي فارمولونه څخه جوړ وي چې د لاندې خوا
ته څانګېدونکې ونه کښې ترتيب شوي وي۔ دا له يو شمېر \emph{فرضونو}
پيلېږي او په هغو نښه لرونکي فارمولونه دوام کوي چې د استنتاج د کومې
قاعدې په پلي کولو تر ځان له پاسه له يوه نښه لرونکي فارمولونه څخه
راوتلي وي۔ هره قاعده اجازه راکوي چې د يوې څانګې په پاے کښې يو يا
څو نښه لرونکي فارمولونه ورزيات کړو، يا دوه نښه لرونکي فارمولونه يو د
بل تر څنګ ورزيات کړو؛ په وروستي حالت کښې څانګه په دوو وېشل کېږي او
دوه ورزيات شوي نښه لرونکي فارمولونه د دواړو څانګو سرونه جوړوي۔
\olpseganchor{OLP-0067-B006}{end}

\olpseganchor{OLP-0067-B007}{start}
پر يوه مرکب نښه لرونکے فارمول د قاعدې له پلي کولو څخه داسې
نښه لرونکي فارمولونه ورزياتېږي چې سمدستي فرعي-فارمولونه وي۔ د هرې
قاعدې يوه جوړه راځي، د دواړو نښو هرې يوې دپاره يوه قاعده۔ مثلاً د
\LR{$\TRule{\True}{\land}$} قاعده پر \LR{$\sFmla{\True}{!A \land !B}$} پلې
کېږي او اجازه راکوي چې د هرې هغې څانګې په پاے کښې دواړه
نښه لرونکي فارمولونه \LR{$\sFmla{\True}{!A}$} او~\LR{$\sFmla{\True}{!B}$}
ورزيات کړو چې \LR{$\sFmla{\True}{!A \land !B}$} لري؛ او د
\LR{$\TRule{\False}{\land}$} قاعده اجازه راکوي چې يوه څانګه د
\LR{$\sFmla{\False}{!A}$} او \LR{$\sFmla{\False}{!B}$} په بېلو
غځونو زيات کړو۔ هره نوې څانګه د اصلي څانګې ټول مخکيني
فارمولونه لري او له دغو دوو څخه يو ټاکلے فارمول ورسره زياتېږي۔
تابلو هله تړلے دے چې د هغه هره څانګه د نښه لرونکي فارمولونه
سره برابره جوړه \LR{$\sFmla{\True}{!A}$} او \LR{$\sFmla{\False}{!A}$} ولري۔
\olpseganchor{OLP-0067-B007}{end}

\olpseganchor{OLP-0067-B008}{start}
پر تابلوګانې ولاړه \LR{$\Proves$} اړيکه داسې تعريفېږي:
\LR{$\Gamma \Proves !A$} هله او يوازې هله چې داسې کوم متناهي سټ
~\LR{$\Gamma_0 = \{!B_1, \dots, !B_n\} \subseteq \Gamma$} موجود وي چې د
لاندې فرضونو دپاره يو تړلے تابلو موجود وي:
\[
\{\sFmla{\False}{!A}, \sFmla{\True}{!B_1}, \dots, \sFmla{\True}{!B_n}\} 
\]
مثلاً دا تړلے تابلو ښيي چې \LR{$\Proves (!A \land !B) \lif !A$}:
\begin{oltableau}
  [\sFmla{\False}{(\formula{A} \land \formula{B}) \lif \formula{A}}, just = \TAss
    [\sFmla{\True}{\formula{A} \land \formula{B}}, just = {\TRule{\False}{\lif}[1]}
      [\sFmla{\False}{\formula{A}}, just = {\TRule{\False}{\lif}[1]}
        [\sFmla{\True}{\formula{A}}, just = {\TRule{\True}{\lif}[2]}
          [\sFmla{\True}{\formula{B}}, just = {\TRule{\True}{\lif}[2]}, close]
        ]
      ]
    ]
  ]
\end{oltableau}
\olpseganchor{OLP-0067-B008}{end}

\olpseganchor{OLP-0067-B009}{start}
يو سټ \LR{$\Gamma$} د تابلو په حساب کښې هله ناسازګار دے چې د
لاندې فرضونو دپاره يو تړلے تابلو موجود وي:
\[
\{\sFmla{\True}{!B_1}, \dots, \sFmla{\True}{!B_n}\} 
\]
د کوم \LR{$!B_i \in \Gamma$} دپاره۔
\olpseganchor{OLP-0067-B009}{end}

\olpseganchor{OLP-0067-B010}{start}
تابلوګانې په ۱۹۵۰مو کلونو کښې اېورټ بېت او ياکو هېنتيکا هر يوه
په خپلواکه توګه جوړ کړل، او رېمونډ سموليان ساده او عام کړل۔ دا کارول
ډېر اسانه دي، ځکه چې د تابلو جوړول ډېره منظمه طريقه ده۔ د
تابلوګانې د منظم طبيعت له امله په کمپيوټر کښې هم په اسانۍ پلي
کېږي۔ خو تابلو زياتره په سختۍ لوستل کېږي او له ثبوتونو سره ئې
تړاو کله کله په اسانۍ نۀ ښکاري۔ دا تګلاره هم ډېره عامه ده، او ډېر
بېل منطقونه د تابلو نظامونه لري۔ تابلوګانې له موږ سره د
هغو جوړښتونه په موندلو کښې هم مرسته کوي چې ورکړل شوي
تړلي فارمولونه، يا د هغو سټونه، صادقوي: که سټ د صدق وړ وي، تړلے
تابلو به و نۀ لري؛ يعنې هر تابلو به يوه پرانيستې څانګه
ولري۔ په دې شرط چې پر هغې څانګه هره ممکنه قاعده پلې شوې وي،
صادقوونکے جوړښت له يوې پرانيستې څانګې څخه ``لوستل کېدے''
شي۔ له سېکوېنټ حساب سره هم ډېر نږدې تړاو شته: په اصل کښې يو تړلے
تابلو د سېکوېنټ په حساب کښې يو لنډ شوے اشتقاق دے چې
سرچپه ليکل شوے وي۔
\olpseganchor{OLP-0067-B010}{end}

\olpseganchor{OLP-0068-B004}{start}
\olfileid{fol}{prf}{axd}
\olpseganchor{OLP-0068-B004}{end}

\olpseganchor{OLP-0068-B005}{start}
\olsection{بديهي اشتقاقونه}
\olpseganchor{OLP-0068-B005}{end}

\olpseganchor{OLP-0068-B006}{start}
بديهي اشتقاقونه د اشتقاق تر ټولو زاړۀ او ساده منطقي
نظامونه دي۔ د هغو اشتقاقونه يوازې د تړلي فارمولونه لړۍ دي۔ د
تړلي فارمولونه يوه لړۍ هله سم اشتقاق ګڼل کېږي چې په هغې کښې
هره تړلے فارمول~\LR{$!A$} له لاندې شرطونو څخه يو پوره کړي:
\begin{enumerate}
\item \LR{$!A$} بديهي اصل وي؛ يا
\item \LR{$!A$} د تړلي فارمولونه د ورکړل شوي سټ~\LR{$\Gamma$} يو غړے
  وي؛ يا
\item \LR{$!A$} د استنتاج په کومه قاعده توجيه شوے وي۔
\end{enumerate}
د بديهي اصل کېدو دپاره \LR{$!A$} بايد د يو شمېر ټاکلو تړلے فارمول
شېماوو له يوې سره بڼه ولري۔ د بديهي اصلونو د شېماوو ډېر داسې سټونه
شته چې د لومړۍ درجې منطق دپاره د اشتقاق د قناعت وړ، سم او
بشپړ نظام ورکوي۔ ځينې د هغو نښلوونکو له مخې ترتيب شوي وي چې واک پرې
چلوي؛ مثلاً دا شېماوې
\[
!A \lif (!B \lif !A) \qquad !B \lif (!B \lor !C) \qquad (!B \land !C) \lif !B
\]
هغه عام بديهي اصول دي چې پر \LR{$\lif$}، \LR{$\lor$} او~\LR{$\land$} واک چلوي۔ د
بديهي اصولو ځينې نظامونه غواړي د اصولو شمېر تر ټولو لږ وي۔ دا چې
کوم نښلوونکي بنسټيز وټاکل شي، ان داسې نظامونه هم موندل کېدے شي چې
يوازې له يوه بديهي اصل څخه جوړ وي۔
\olpseganchor{OLP-0068-B006}{end}

\olpseganchor{OLP-0068-B007}{start}
د استنتاج قاعده يوه شرطي وينا ده چې په اشتقاق کښې د
تړلے فارمول د توجيه کېدو دپاره بس شرط ورکوي۔ مودوس پوننس يوه ډېره
عامه داسې قاعده ده: وايي که \LR{$!A$} او \LR{$!A \lif !B$} له مخکښې توجيه
شوي وي، نو \LR{$!B$} توجيه شوے دے۔ دا معنٰی لري چې په اشتقاق
کښې هغه کرښه چې تړلے فارمول~\LR{$!B$} لري، په دې شرط توجيه شوې ده چې
\LR{$!A$} او \LR{$!A \lif !B$} دواړه، د کوم تړلے فارمول~\LR{$!A$} دپاره، په
اشتقاق کښې له~\LR{$!B$} څخه مخکښې راغلي وي۔
\olpseganchor{OLP-0068-B007}{end}

\olpseganchor{OLP-0068-B008}{start}
پر بديهي اشتقاقونه ولاړه \LR{$\Proves$} اړيکه داسې تعريفېږي:
\LR{$\Gamma \Proves !A$} هله او يوازې هله چې داسې اشتقاق موجود
وي چې تړلے فارمول~\LR{$!A$} ئې وروستے فارمول وي، او \LR{$\Gamma$} په هغه
اشتقاق کښې د هغو تړلي فارمولونه سټ ونيول شي چې د پورته (۲) له
مخې توجيه شوي وي۔ \LR{$!A$} هله قضيه ده چې \LR{$!A$} داسې اشتقاق
ولري چې~\LR{$\Gamma$} پکې تش وي؛ يعنې په اشتقاق کښې هره
تړلے فارمول يا د (۱) له مخې توجيه شوې وي يا د~(۳) له مخې۔ مثلاً دا
اشتقاق ښيي چې \LR{$\Proves !A  \lif (!B \lif (!B \lor !A))$}:
\begin{derivation}
  1. & \LR{$!B \lif (!B \lor !A)$} \\
  2. & \LR{$(!B \lif (!B \lor !A)) \lif (!A  \lif (!B \lif (!B \lor !A)))$}\\
  3. & \LR{$!A  \lif (!B \lif (!B \lor !A))$}
\end{derivation}
په ۱مه کرښه تړلے فارمول د بديهي اصل \LR{$!A \lif (!A \lor !B)$} بڼه
لري، چې د \LR{$!A$} او \LR{$!B$} ونډې پکې بدلې شوې دي۔ په ۲مه کرښه جمله د
بديهي اصل \LR{$!A \lif (!B \lif !A)$} بڼه لري۔ نو دواړه کرښې توجيه شوې
دي۔ ۳مه کرښه د مودوس پوننس له مخې توجيه شوې ده: که په لنډه ئې
\LR{$!D$} وليکو، نو ۲مه کرښه د \LR{$!C \lif !D$} بڼه لري، چې \LR{$!C$} هم
\LR{$!B \lif (!B \lor !A)$}، يعنې ۱مه کرښه، ده۔
\olpseganchor{OLP-0068-B008}{end}

\olpseganchor{OLP-0068-B009}{start}
يو سټ \LR{$\Gamma$} هله ناسازګار دے چې \LR{$\Gamma \Proves \lfalse$}۔ د
بديهي اصولو يو بشپړ نظام به دا هم ثابتوي چې د هر~\LR{$!A$} دپاره
\LR{$\lfalse \lif !A$}؛ نو که \LR{$\Gamma$} ناسازګار وي، د هر~\LR{$!A$} دپاره
\LR{$\Gamma \Proves !A$}۔
\olpseganchor{OLP-0068-B009}{end}

\olpseganchor{OLP-0068-B010}{start}
د منطق د بديهي اشتقاقونه نظامونه لومړے ګوټلوب فرېګه په خپل
۱۸۷۹ز اثر \emph{بېګريفشريفت} کښې وړاندې کړل، چې له همدې امله زياتره
د معاصر منطق لومړے اثر ګڼل کېږي۔ بيا الفرېډ نارت وايټ هېډ او برټرېنډ
رسل په \emph{پرېنسيپيا ماتيماتيکا} کښې، او ډيوېډ هيلبرټ او شاګردانو
ئې په ۱۹۲۰مو کلونو کښې بشپړ کړل۔ ځکه ورته زياتره ``د فرېګه نظامونه''
يا ``د هيلبرټ نظامونه'' ويل کېږي۔ دا ډېر انعطاف منونکي دي، ځکه چې د
يو منطق دپاره بديهي نظام زياتره په اسانۍ موندل کېږي۔ دا چې
اشتقاقونه ډېر ساده جوړښت او يوازې يوه يا دوه استنتاجي قاعدې
لري، د هغو \emph{په اړه} د څيزونو ثابتول هم نسبتاً اسانه دي۔ خو په
عمل کښې ئې کارول ډېر ګران دي؛ يعنې ثبوتونه موندل او ليکل ګران دي۔
\olpseganchor{OLP-0068-B010}{end}

\olpseganchor{OLP-0069-B004}{start}
\olchapter{fol}{seq}{د سېکوېنټ حساب}
\olpseganchor{OLP-0069-B004}{end}

\olpseganchor{OLP-0069-B005}{start}
\begin{editorial}
  دا څپرکے د کلاسیک د لومړۍ درجې منطق دپاره د ګېنتڅن معياري
  سېکوېنټ حساب \texttt{LK} وړاندې کوي۔ دې ته نور مثالونه او تمرينونه
  هم پکار دي۔ د ثبوت د نظام په توګه د سېکوېنټ له حساب سره د اړوندو
  موادو د شاملولو يا اېستلو دپاره د \texttt{prfLK} ټېګ وکاروئ۔
\end{editorial}
\olpseganchor{OLP-0069-B005}{end}





\olpseganchor{OLP-0069-B008}{start}
%
  

\olpseganchor{OLP-0069-B008}{end}







\olpseganchor{OLP-0069-B012}{start}
%
  

\olpseganchor{OLP-0069-B012}{end}







\olpseganchor{OLP-0069-B016}{start}
%
  

\olpseganchor{OLP-0069-B016}{end}



\olpseganchor{OLP-0069-B018}{start}
%
  
  

\olpseganchor{OLP-0069-B018}{end}

\olpseganchor{OLP-0070-B004}{start}
\olfileid{fol}{seq}{rul}
\olpseganchor{OLP-0070-B004}{end}

\olpseganchor{OLP-0070-B005}{start}
\olsection{قاعدې او اشتقاقونه}
\olpseganchor{OLP-0070-B005}{end}

\olpseganchor{OLP-0070-B006}{start}
د راتلونکو خبرو دپاره، پرېږدئ چې \LR{$\Gamma, \Delta, \Pi, \Lambda$} د
تړلي فارمولونه متناهي لړۍ وښيي۔
\olpseganchor{OLP-0070-B006}{end}

\olpseganchor{OLP-0070-B007}{start}
\begin{defn}[سېکوېنټ]
\emph{سېکوېنټ} د لاندې بڼې څرګندنه ده:
\[
\Gamma \Sequent \Delta
\]
دلته \LR{$\Gamma$} او \LR{$\Delta$} د ژبې \LR{$\Lang L$} د تړلي فارمولونه متناهي
(او کېدے شي تشې) لړۍ دي۔ \LR{$\Gamma$} ته \emph{مقدم} او \LR{$\Delta$} ته
\emph{تالي} ويل کېږي۔
\end{defn}
\olpseganchor{OLP-0070-B007}{end}

\olpseganchor{OLP-0070-B008}{start}
\begin{explain}
د يوه سېکوېنټ تر شا شهودي مفکوره دا ده: که د مقدم ټولې تړلي فارمولونه
صدق کوي، نو د تالي لږ تر لږه يوه تړلي فارمولونه صدق کوي۔ يعنې که
\LR{$\Gamma = \tuple{!A_1, \dots, !A_m}$} او
\LR{$\Delta = \tuple{!B_1, \dots, !B_n}$} وي، نو \LR{$\Gamma \Sequent \Delta$}
هله صدق کوي چې
\[
(!A_1 \land \cdots \land !A_m) \lif (!B_1 \lor \cdots \lor
!B_n)
\]
صدق کوي۔ دوه خاص حالتونه شته: يو هغه چې \LR{$\Gamma$} تش وي، او بل هغه
چې \LR{$\Delta$}~تش وي۔ کله چې \LR{$\Gamma$} تش وي، يعنې \LR{$m = 0$}، نو \LR{$\quad
\Sequent \Delta$} هله صدق کوي چې \LR{$!B_1 \lor \dots \lor !B_n$} صدق
کوي۔ کله چې \LR{$\Delta$} تش وي، يعنې \LR{$n = 0$}، نو \LR{$\Gamma \Sequent \quad$}
هله صدق کوي چې \LR{$\lnot(!A_1 \land \dots \land !A_m)$} صدق کوي۔
موږ يو سېکوېنټ هله معتبر بولو چې ورسره سمون لرونکې تړلے فارمول
معتبره وي۔
\end{explain}
\olpseganchor{OLP-0070-B008}{end}

\olpseganchor{OLP-0070-B009}{start}
که \LR{$\Gamma$} د تړلي فارمولونه يوه لړۍ وي، نو \LR{$\Gamma, !A$} د هغې
پايلې دپاره ليکو چې~\LR{$!A$} د~\LR{$\Gamma$} ښي پاے ته ورزيات شي (او \LR{$!A,
\Gamma$} د هغې پايلې دپاره چې~\LR{$!A$} د~\LR{$\Gamma$} کيڼ پاے ته ورزيات
شي)۔ که \LR{$\Delta$} هم د تړلي فارمولونه يوه لړۍ وي، نو \LR{$\Gamma,
\Delta$} د دواړو لړيو نښلول دي۔
\olpseganchor{OLP-0070-B009}{end}

\olpseganchor{OLP-0070-B010}{start}
\begin{defn}[لومړنے سېکوېنټ]
\emph{لومړنے سېکوېنټ} داسې سېکوېنټ دے
چې له لاندې بڼو څخه يوه لري:
  \begin{enumerate}
    \item \LR{$!A \Sequent !A$}
    \item \LR{$\quad \Sequent \ltrue$}
    \item \LR{$\lfalse \Sequent \quad$}
  \end{enumerate}، او \LR{$!A$} د ژبې هره تړلے فارمول
کېدے شي۔
\end{defn}
\olpseganchor{OLP-0070-B010}{end}

\olpseganchor{OLP-0070-B011}{start}
په سېکوېنټ حساب کښې اشتقاقونه د سېکوېنټونو داسې ټاکلې ونې
دي چې تر ټولو بره سېکوېنټونه ئې لومړني سېکوېنټونه وي؛ او که يو
سېکوېنټ د يوه يا دوو نورو سېکوېنټونو لاندې راشي، بايد د استنتاج د
يوې قاعدې له مخې ترې په سمه توګه راووځي۔ د~\LR{$\Log{LK}$} قاعدې پر دوو
اصلي ډولونو وېشل کېږي: \emph{منطقي} قاعدې او \emph{جوړښتي} قاعدې۔
منطقي قاعدې د لاندې سېکوېنټ د هغې تړلے فارمول د اصلي عملګر په نوم نومول کېږي چې \LR{$!A$} او/يا \LR{$!B$} لري۔ هره قاعده
دوه بڼې لري: يوه د داسې سېکوېنټ د استنتاج دپاره چې د
منطقي عملګر لرونکې تړلے فارمول ئې کيڼ لوري ته وي، او بله د هغه
سېکوېنټ دپاره چې دا تړلے فارمول ئې ښي لوري ته وي۔
\olpseganchor{OLP-0070-B011}{end}

\olpseganchor{OLP-0071-B004}{start}
\olfileid{fol}{seq}{prl}
\olpseganchor{OLP-0071-B004}{end}

\olpseganchor{OLP-0071-B005}{start}
\olsection{بياني قاعدې}
\olpseganchor{OLP-0071-B005}{end}

\olpseganchor{OLP-0071-B006}{start}
\subsection{د \texorpdfstring{\LR{$\lnot$}}{¬} دپاره قاعدې}
\olpseganchor{OLP-0071-B006}{end}

\olpseganchor{OLP-0071-B007}{start}
\begin{defish}
\Axiom$ \Gamma \fCenter \Delta, !A $
\RightLabel{\LeftR{\lnot}}
\UnaryInf$ \lnot !A, \Gamma \fCenter \Delta$
\DisplayProof
\hfill
\Axiom$!A, \Gamma \fCenter \Delta$
\RightLabel{\RightR{\lnot}}
\UnaryInf$ \Gamma \fCenter \Delta, \lnot !A $
\DisplayProof
\end{defish}
\olpseganchor{OLP-0071-B007}{end}

\olpseganchor{OLP-0071-B008}{start}
\subsection{د \texorpdfstring{\LR{$\land$}}{∧} دپاره قاعدې}
\olpseganchor{OLP-0071-B008}{end}

\olpseganchor{OLP-0071-B009}{start}
\begin{defish}\noindent
\begin{tabular}{l}
\Axiom$ !A, \Gamma \fCenter \Delta$
\RightLabel{\LeftR{\land}}
\UnaryInf$ !A \land !B, \Gamma \fCenter \Delta$
\DisplayProof
\\[3ex]
\Axiom$!B, \Gamma \fCenter \Delta$
\RightLabel{\LeftR{\land}}
\UnaryInf$!A \land !B, \Gamma \fCenter \Delta$
\DisplayProof
\end{tabular}
\hfill
\Axiom$\Gamma \fCenter \Delta, !A$
\Axiom$ \Gamma \fCenter \Delta, !B$
\RightLabel{\RightR{\land}}
\BinaryInf$ \Gamma \fCenter \Delta, !A \land !B $
\DisplayProof
\end{defish}
\olpseganchor{OLP-0071-B009}{end}

\olpseganchor{OLP-0071-B010}{start}
\subsection{د \texorpdfstring{\LR{$\lor$}}{∨} دپاره قاعدې}
\olpseganchor{OLP-0071-B010}{end}

\olpseganchor{OLP-0071-B011}{start}
\begin{defish}
\Axiom$!A, \Gamma \fCenter \Delta$
\Axiom$!B, \Gamma \fCenter \Delta$
\RightLabel{\LeftR{\lor}}
\BinaryInf$!A \lor !B, \Gamma \fCenter \Delta$
\DisplayProof
\hfill
\begin{tabular}{r}
\Axiom$\Gamma \fCenter \Delta, !A$
\RightLabel{\RightR{\lor}}
\UnaryInf$ \Gamma \fCenter \Delta, !A \lor !B$
\DisplayProof
\\[3ex]
\Axiom$ \Gamma \fCenter \Delta, !B$
\RightLabel{\RightR{\lor}}
\UnaryInf$ \Gamma \fCenter \Delta, !A \lor !B$
\DisplayProof
\end{tabular}
\end{defish}
\olpseganchor{OLP-0071-B011}{end}

\olpseganchor{OLP-0071-B012}{start}
\subsection{د \texorpdfstring{\LR{$\lif$}}{→} دپاره قاعدې}
\olpseganchor{OLP-0071-B012}{end}

\olpseganchor{OLP-0071-B013}{start}
\begin{defish}
\Axiom$ \Gamma \fCenter \Delta, !A$
\Axiom$ !B, \Pi \fCenter \Lambda$
\RightLabel{\LeftR{\lif}}
\BinaryInf$ !A \lif !B, \Gamma, \Pi \fCenter \Delta, \Lambda$
\DisplayProof
\hfill
\Axiom$ !A, \Gamma \fCenter \Delta, !B$
\RightLabel{\RightR{\lif}}
\UnaryInf$ \Gamma \fCenter \Delta, !A \lif !B $
\DisplayProof
\end{defish}
\olpseganchor{OLP-0071-B013}{end}

\olpseganchor{OLP-0072-B004}{start}
\olfileid{fol}{seq}{qrl}
\olpseganchor{OLP-0072-B004}{end}

\olpseganchor{OLP-0072-B005}{start}
\olsection{د سورونو قاعدې}
\olpseganchor{OLP-0072-B005}{end}

\olpseganchor{OLP-0072-B006}{start}
\subsection{د \texorpdfstring{\LR{$\lforall$}}{∀} دپاره قاعدې}
\olpseganchor{OLP-0072-B006}{end}

\olpseganchor{OLP-0072-B007}{start}
\begin{defish}
\Axiom$ !A(t), \Gamma \fCenter \Delta$
\RightLabel{\LeftR{\lforall}}
\UnaryInf$ \lforall[x][!A(x)],\Gamma \fCenter \Delta$
\DisplayProof
\hfill
\Axiom$ \Gamma \fCenter \Delta, !A(a) $
\RightLabel{\RightR{\lforall}}
\UnaryInf$ \Gamma \fCenter \Delta, \lforall[x][!A(x)]$
\DisplayProof
\end{defish}
\olpseganchor{OLP-0072-B007}{end}

\olpseganchor{OLP-0072-B008}{start}
په \LeftR{\lforall} کښې \LR{$t$} يوه تړلې اصطلاح ده (يعنې داسې اصطلاح
چې متغيرونه نۀ لري)۔ په \RightR{\lforall} کښې \LR{$a$}~ثابت سمبول دے
چې بايد د \RightR{\lforall} قاعدې په لاندې سېکوېنټ کښې هېڅ ځاے نۀ
وي راغلے۔ موږ \LR{$a$} د \RightR{\forall} استنتاج \emph{ځانګړے متغير}
بولو۔\footnote{موږ د “ځانګړي متغير” اصطلاح کاروو، که څه هم په
پورتنۍ قاعده کښې \LR{$a$} ثابت سمبول دے۔ دا نوم تاريخي لاملونه لري۔}
\olpseganchor{OLP-0072-B008}{end}

\olpseganchor{OLP-0072-B009}{start}
\subsection{د \texorpdfstring{\LR{$\lexists$}}{∃} دپاره قاعدې}
\olpseganchor{OLP-0072-B009}{end}

\olpseganchor{OLP-0072-B010}{start}
\begin{defish}
\Axiom$ !A(a), \Gamma \fCenter \Delta $
\RightLabel{\LeftR{\lexists}}
\UnaryInf$ \lexists[x][!A(x)], \Gamma \fCenter \Delta$
\DisplayProof
\hfill
\Axiom$ \Gamma \fCenter \Delta, !A(t) $
\RightLabel{\RightR{\lexists}}
\UnaryInf$ \Gamma \fCenter \Delta, \lexists[x][!A(x)]$
\DisplayProof
\end{defish}
\olpseganchor{OLP-0072-B010}{end}

\olpseganchor{OLP-0072-B011}{start}
بيا هم، \LR{$t$}~يوه تړلې اصطلاح ده، او \LR{$a$}~داسې ثابت سمبول دے چې
د \LeftR{\lexists} قاعدې په لاندې سېکوېنټ کښې نۀ راځي۔ موږ \LR{$a$} د
\LeftR{\lexists} استنتاج \emph{ځانګړے متغير} بولو۔
\olpseganchor{OLP-0072-B011}{end}

\olpseganchor{OLP-0072-B012}{start}
دا شرط چې ځانګړے متغير د \RightR{\lforall} يا
\LeftR{\lexists} استنتاج په لاندې سېکوېنټ کښې نۀ راځي،
\emph{د ځانګړي متغير شرط} بلل کېږي۔
\olpseganchor{OLP-0072-B012}{end}

\olpseganchor{OLP-0072-B013}{start}
\begin{explain}
هغه دود راياد کړئ چې کله \LR{$!A$} داسې فارمول وي چې
متغير~\LR{$x$} پکې ازاد وي، نو دا په~\LR{$!A(x)$} ليکلو ښيو۔ په همدې
چوکاټ کښې بيا \LR{$!A(t)$} د~\LR{$\Subst{!A}{t}{x}$} لنډون دے۔ نو د
\RightR{\lexists} قاعده داسې هم ليکلے شو:
\begin{prooftree}
  \Axiom$\Gamma \fCenter \Delta, \Subst{!A}{t}{x}$
  \RightLabel{\RightR{\lexists}}
  \UnaryInf$\Gamma \fCenter \Delta, \lexists[x][!A]$
\end{prooftree}
پام وکړئ چې \LR{$t$} کېدے شي له مخکښې په~\LR{$!A$} کښې راغلے وي؛ مثلاً
\LR{$!A$}~کېدے شي~\LR{$\Atom{\Obj P}{t,x}$} وي۔ له دې امله، له
~\LR{$\Gamma \Sequent \Delta, \Atom{\Obj P}{t,t}$} څخه
\LR{$\Gamma \Sequent \Delta, \lexists[x][\Atom{\Obj P}{t,x}]$} استنتاجول
د~\RightR{\lexists} سمه کارونه ده---تاسو د مقدمې د~\LR{$t$} يوه يا
زياتې پېښې په تړلي متغير~\LR{$x$} “بدلولے” شئ، او اړينه نۀ ده
چې ټولې پېښې ئې بدلې کړئ۔ خو د \RightR{\lforall} او
~\LeftR{\lexists} د ځانګړي متغير شرطونه غواړي چې
ثابت سمبول~\LR{$a$} په~\LR{$!A$} کښې نۀ راځي۔ نو د~\LR{$\RightR{\lforall}$} په
کارولو سره له \LR{$\Gamma \Sequent \Delta, \Atom{\Obj P}{a,a}$} څخه
\LR{$\Gamma \Sequent \Delta, \lforall[x][\Atom{\Obj P}{a,x}]$} په سمه
توګه نۀ شئ استنتاجولے۔
\end{explain}
\olpseganchor{OLP-0072-B013}{end}

\olpseganchor{OLP-0072-B014}{start}
\begin{explain}
په \RightR{\lexists} او \LeftR{\lforall} کښې پر اصطلاح~\LR{$t$} هېڅ
بنديز نشته۔ بل لوري ته، د \LeftR{\lexists} او
\RightR{\lforall} په قاعدو کښې د ځانګړي متغير شرط غواړي چې
ثابت سمبول~\LR{$a$} په پورتني سېکوېنټ کښې د~\LR{$!A(a)$} څخه بهر هېڅ ځاے
نۀ راځي۔ دا شرط د نظام د سموالي د يقيني کولو دپاره اړين دے؛ يعنې
نظام يوازې معتبر سېکوېنټونه اشتقاق کړي۔ له دې شرط پرته لاندې
استنتاجونه روا کېدل:
\begin{prooftree}
  \Axiom$!A(a) \fCenter !A(a)$
  \RightLabel{*\LeftR{\lexists}}
  \UnaryInf$\lexists[x][!A(x)] \fCenter !A(a)$
  \RightLabel{\RightR{\lforall}}
  \UnaryInf$\lexists[x][!A(x)] \fCenter \lforall[x][!A(x)]$
  \DisplayProof\bottomAlignProof
  \qquad
  \Axiom$!A(a) \fCenter !A(a)$
  \RightLabel{*\RightR{\lforall}}
  \UnaryInf$!A(a) \fCenter \lforall[x][!A(x)]$
  \RightLabel{\LeftR{\lexists}}
  \UnaryInf$\lexists[x][!A(x)] \fCenter \lforall[x][!A(x)]$
\end{prooftree}
خو \LR{$\lexists[x][!A(x)] \Sequent \lforall[x][!A(x)]$} معتبر نۀ دے۔
\end{explain}
\olpseganchor{OLP-0072-B014}{end}

\olpseganchor{OLP-0073-B004}{start}
\olfileid{fol}{seq}{srl}
\olpseganchor{OLP-0073-B004}{end}

\olpseganchor{OLP-0073-B005}{start}
\olsection{جوړښتي قاعدې}
\olpseganchor{OLP-0073-B005}{end}

\olpseganchor{OLP-0073-B006}{start}
موږ څو داسې قاعدو ته هم اړتيا لرو چې د يوه سېکوېنټ په کيڼ او ښي
لوري کښې د تړلي فارمولونه د بيا ترتيبولو اجازه راکړي۔ ځکه چې منطقي
قاعدې غواړي هغه تړلي فارمولونه چې قاعده پرې کار کوي، يا په بشپړ کيڼ
او يا په بشپړ ښي پاے کښې وي، نو د “تبادلې” داسې قاعدې ته اړتيا
لرو چې تړلي فارمولونه سم ځاے ته يوسي۔ کله کله دا هم مهمه وي چې دوه
يو شان تړلي فارمولونه سره يو کړو او هر لوري ته تړلے فارمول
ورزياته کړو۔
\olpseganchor{OLP-0073-B006}{end}

\olpseganchor{OLP-0073-B007}{start}
\subsection{کمزوري کول}
\olpseganchor{OLP-0073-B007}{end}

\olpseganchor{OLP-0073-B008}{start}
\begin{defish}
\Axiom$ \Gamma \fCenter \Delta $
\RightLabel{\LeftR{\Weakening}}
\UnaryInf$ !A, \Gamma \fCenter \Delta$
\DisplayProof
\hfill
\Axiom$ \Gamma \fCenter \Delta$
\RightLabel{\RightR{\Weakening}}
\UnaryInf$ \Gamma \fCenter \Delta, !A$
\DisplayProof
\end{defish}
\olpseganchor{OLP-0073-B008}{end}

\olpseganchor{OLP-0073-B009}{start}
\subsection{انقباض}
\olpseganchor{OLP-0073-B009}{end}

\olpseganchor{OLP-0073-B010}{start}
\begin{defish}
\Axiom$ !A, !A, \Gamma \fCenter \Delta $
\RightLabel{\LeftR{\Contraction}}
\UnaryInf$ !A, \Gamma \fCenter \Delta$
\DisplayProof
\hfill
\Axiom$ \Gamma \fCenter \Delta, !A, !A$
\RightLabel{\RightR{\Contraction}}
\UnaryInf$ \Gamma \fCenter \Delta, !A$
\DisplayProof
\end{defish}
\olpseganchor{OLP-0073-B010}{end}

\olpseganchor{OLP-0073-B011}{start}
\subsection{تبادله}
\olpseganchor{OLP-0073-B011}{end}

\olpseganchor{OLP-0073-B012}{start}
\begin{defish}
\Axiom$ \Gamma, !A, !B, \Pi \fCenter \Delta $
\RightLabel{\LeftR{\Exchange}}
\UnaryInf$ \Gamma, !B, !A, \Pi \fCenter \Delta$
\DisplayProof
\hfill
\Axiom$ \Gamma \fCenter \Delta, !A, !B, \Lambda$
\RightLabel{\RightR{\Exchange}}
\UnaryInf$ \Gamma \fCenter \Delta, !B, !A, \Lambda$
\DisplayProof
\end{defish}
\olpseganchor{OLP-0073-B012}{end}

\olpseganchor{OLP-0073-B013}{start}
د کمزوري کولو، انقباض او تبادلې د استنتاجونو يوه لړۍ به زياتره د
استنتاج په دوو کرښو ښوول کېږي۔
\olpseganchor{OLP-0073-B013}{end}

\olpseganchor{OLP-0073-B014}{start}
لاندې قاعده چې “پرېکون” بلل کېږي، په کره توګه اړينه نۀ ده، خو د
اشتقاقونه بيا کارول او يو ځاے کول ډېر اسانوي۔
\begin{defish}
\[
\Axiom$ \Gamma \fCenter \Delta, !A$
\Axiom$ !A, \Pi \fCenter \Lambda $
\RightLabel{\Cut}
\BinaryInf$ \Gamma, \Pi \fCenter \Delta, \Lambda$
\DisplayProof
\]
\end{defish}
\olpseganchor{OLP-0073-B014}{end}

\olpseganchor{OLP-0074-B004}{start}
\olfileid{fol}{seq}{der}
\olpseganchor{OLP-0074-B004}{end}

\olpseganchor{OLP-0074-B005}{start}
\olsection{اشتقاقونه}
\olpseganchor{OLP-0074-B005}{end}

\olpseganchor{OLP-0074-B006}{start}
\begin{explain}
موږ وويل چې لومړنے سېکوېنټ څنګه وي او د استنتاج قاعدې مو ورکړې۔
په سېکوېنټ حساب کښې اشتقاقونه له دغو څخه په استقرايي توګه
جوړېږي: هر اشتقاق يا په يوازې ځان يو لومړنے سېکوېنټ وي،
او يا له يوه يا دوو اشتقاقونه او ورپسې يوه استنتاج څخه جوړ وي۔
\end{explain}
\olpseganchor{OLP-0074-B006}{end}

\olpseganchor{OLP-0074-B007}{start}
\begin{defn}[\LR{$\Log{LK}$} اشتقاق]
د يوه سېکوېنټ~\LR{$S$} \emph{\LR{$\Log{LK}$}-اشتقاق} د سېکوېنټونو
داسې متناهي ونه ده چې لاندې شرطونه پوره کوي:
\begin{enumerate}
\item د ونې تر ټولو بره سېکوېنټونه لومړني سېکوېنټونه دي۔
\item د ونې تر ټولو لاندې سېکوېنټ~\LR{$S$} دے۔
\item له~\LR{$S$} پرته د ونې هر سېکوېنټ د استنتاج د يوې قاعدې د سمې
  کارونې مقدمه ده، او د هغې کارونې پايله په ونه کښې د هماغه
  سېکوېنټ مخامخ لاندې ولاړه وي۔
\end{enumerate}
بيا وايو چې \LR{$S$} د اشتقاق \emph{وروستنے سېکوېنټ} دے او
\LR{$S$} \emph{په \LR{$\Log{LK}$} کښې د اشتقاق وړ} دے (يا
\LR{$\Log{LK}$}-د اشتقاق وړ دے)۔
\end{defn}
\olpseganchor{OLP-0074-B007}{end}

\olpseganchor{OLP-0074-B008}{start}
\begin{ex}
هر لومړنے سېکوېنټ، مثلاً \LR{$!C \Sequent !C$}، اشتقاق دے۔
له دې څخه د يوې بلې قاعدې، مثلاً د \LR{$\LeftR{\Weakening}$} قاعدې، په
کارولو نوے اشتقاق ترلاسه کولے شو:
\begin{prooftree}
\Axiom$ \Gamma \fCenter \Delta $
\RightLabel{\LeftR{\Weakening}}
\UnaryInf$ !A, \Gamma \fCenter \Delta$
\end{prooftree}
خو قاعده عمومي مراد لري: په قاعده کښې \LR{$!A$} د هرې تړلے فارمول،
مثلاً د~\LR{$!D$}، په ځاے راوستلے شو۔ که مقدمه زموږ له لومړني سېکوېنټ
\LR{$!C \Sequent !C$} سره سمون خوري، نو \LR{$\Gamma$} او \LR{$\Delta$} دواړه
يوازې~\LR{$!C$} دي، او پايله به بيا \LR{$!D, !C \Sequent !C$} وي۔ نو لاندې
اشتقاق دے:
\begin{prooftree}
\Axiom$ !C \fCenter !C $
\RightLabel{\LeftR{\Weakening}}
\UnaryInf$ !D, !C \fCenter !C$
\end{prooftree}
اوس يوه بله قاعده، مثلاً \LR{$\LeftR{\Exchange}$}، کارولے شو چې په کيڼ
لوري کښې د دوو تړلي فارمولونه د ځاے بدلولو اجازه راکوي۔ نو لاندې هم
سم اشتقاق دے:
\begin{prooftree}
\Axiom$ !C \fCenter !C $
\RightLabel{\LeftR{\Weakening}}
\UnaryInf$ !D, !C \fCenter !C$
\RightLabel{\LeftR{\Exchange}}
\UnaryInf$ !C, !D \fCenter !C$
\end{prooftree}
د دې قاعدې په دې کارونه کښې، چې داسې ورکړل شوې وه:
\begin{prooftree}
\Axiom$ \Gamma, !A, !B, \Pi \fCenter \Delta $
\RightLabel{\LeftR{\Exchange}}
\UnaryInf$ \Gamma, !B, !A, \Pi \fCenter \Delta,$
\end{prooftree}
\LR{$\Gamma$} او \LR{$\Pi$} دواړه تش وو، \LR{$\Delta$} هم \LR{$!C$} دے، او د \LR{$!A$} او
\LR{$!B$} ونډې په همدې ترتيب \LR{$!D$} او~\LR{$!C$} لوبوي۔ نژدې په همدې طريقه
وينو چې
\begin{prooftree}
\Axiom$ !D \fCenter !D $
\RightLabel{\LeftR{\Weakening}}
\UnaryInf$ !C, !D \fCenter !D$
\end{prooftree}
هم اشتقاق دے۔ اوس دا دواړه اشتقاقونه اخلو او د
\LR{$\RightR{\land}$} په کارولو ئې يو ځاے کوو۔ هغه قاعده داسې وه:
\begin{prooftree}
\Axiom$\Gamma \fCenter \Delta, !A$
\Axiom$ \Gamma \fCenter \Delta, !B$
\RightLabel{\RightR{\land}}
\BinaryInf$ \Gamma \fCenter \Delta, !A \land !B $
\end{prooftree}
زموږ په حالت کښې مقدمې بايد د هغو اشتقاقونه له وروستيو
سېکوېنټونو سره سمون وخوري چې په مقدمو پاے ته رسېږي۔ يعنې \LR{$\Gamma$}
د \LR{$!C, !D$}، \LR{$\Delta$} تشه، \LR{$!A$} د \LR{$!C$} او \LR{$!B$} د \LR{$!D$} سره برابره
ده۔ نو که استنتاج سم وي، پايله \LR{$!C, !D \Sequent !C
\land !D$} ده۔
\begin{prooftree}
\Axiom$ !C \fCenter !C $
\RightLabel{\LeftR{\Weakening}}
\UnaryInf$ !D, !C \fCenter !C$
\RightLabel{\LeftR{\Exchange}}
\UnaryInf$ !C, !D \fCenter !C$
\Axiom$ !D \fCenter !D $
\RightLabel{\LeftR{\Weakening}}
\UnaryInf$ !C, !D \fCenter !D$
\RightLabel{\RightR{\land}}
\BinaryInf$ !C, !D \fCenter !C \land !D $
\end{prooftree}
البته مقدمې سرچپه کولے هم شو؛ بيا \LR{$!A$} به \LR{$!D$} او \LR{$!B$} به~\LR{$!C$} وي۔
\begin{prooftree}
\Axiom$ !D \fCenter !D $
\RightLabel{\LeftR{\Weakening}}
\UnaryInf$ !C, !D \fCenter !D$
\Axiom$ !C \fCenter !C $
\RightLabel{\LeftR{\Weakening}}
\UnaryInf$ !D, !C \fCenter !C$
\RightLabel{\LeftR{\Exchange}}
\UnaryInf$ !C, !D \fCenter !C$
\RightLabel{\RightR{\land}}
\BinaryInf$ !C, !D \fCenter !D \land !C $
\end{prooftree}
\end{ex}
\olpseganchor{OLP-0074-B008}{end}

\olpseganchor{OLP-0075-B004}{start}
\olfileid{fol}{seq}{pro}
\olpseganchor{OLP-0075-B004}{end}

\olpseganchor{OLP-0075-B005}{start}
\olsection{د اشتقاقونه بېلګې}
\olpseganchor{OLP-0075-B005}{end}

\olpseganchor{OLP-0075-B006}{start}
\begin{ex}
د سېکوېنټ \LR{$!A \land !B \Sequent !A$} يو
\LR{$\Log{LK}$}-اشتقاق ورکړئ۔
\olpseganchor{OLP-0075-B006}{end}

\olpseganchor{OLP-0075-B007}{start}
د اشتقاق په لاندې برخه کښې د غوښتل شوي وروستي سېکوېنټ په
ليکلو پيل کوو۔
\begin{prooftree}
\AxiomC{}
\UnaryInf$!A\land !B \fCenter !A$
\end{prooftree}
بيا بايد معلومه کړو چې د کوم ډول استنتاج لاندې سېکوېنټ دا بڼه
لرلے شي۔ دا يوه جوړښتي قاعده کېدے شي، خو غوره ده چې لومړے کومه
منطقي قاعده ولټوؤ۔ په لاندې سېکوېنټ کښې يوازينے منطقي نښلوونکے
\LR{$\land$} دے؛ نو د \LR{$\land$} قاعده لټوؤ، او ځکه چې د \LR{$\land$} نښه په
مقدم کښې راځي، د \LeftR{\land} قاعده ګورو۔
\begin{prooftree}
\AxiomC{}
\RightLabel{\LeftR{\land}}
\UnaryInf$!A\land !B \fCenter !A$
\end{prooftree}
د \LeftR{\land} استنتاج د پورتني سېکوېنټ دپاره دوه امکانونه شته:
يا \LR{$!A \Sequent !A$} او يا \LR{$!B \Sequent !A$}۔ څرګنده ده چې
\LR{$!A \Sequent !A$} يو لومړنے سېکوېنټ دے، او دا ښه خبره ده؛ خو
\LR{$!B \Sequent !A$} په عمومي توګه د اشتقاق وړ نۀ دے۔ پورتنے سېکوېنټ
ورډکوو:
\begin{prooftree}
\Axiom$!A \fCenter !A$
\RightLabel{\LeftR{\land}}
\UnaryInf$!A\land !B \fCenter !A$
\end{prooftree}
اوس د سېکوېنټ \LR{$!A \land !B \Sequent !A$} سم
\LR{$\Log{LK}$}-اشتقاق لرو۔
\end{ex}
\olpseganchor{OLP-0075-B007}{end}

\olpseganchor{OLP-0075-B008}{start}
\begin{ex}
د سېکوېنټ \LR{$\lnot !A \lor !B \Sequent !A \lif !B$} يو
\LR{$\Log{LK}$}-اشتقاق ورکړئ۔
\olpseganchor{OLP-0075-B008}{end}

\olpseganchor{OLP-0075-B009}{start}
د اشتقاق په لاندې برخه کښې د غوښتل شوي وروستي سېکوېنټ په
ليکلو پيل وکړئ۔
\begin{prooftree}
\AxiomC{}
\UnaryInf$\lnot !A \lor !B \fCenter !A \lif !B$
\end{prooftree}
د داسې منطقي قاعدې د موندلو دپاره چې دا وروستنے سېکوېنټ راکړي، د
وروستي سېکوېنټ منطقي نښلوونکي ګورو: \LR{$\lnot$}، \LR{$\lor$} او \LR{$\lif$}۔
اوس مهال يوازې \LR{$\lor$} او \LR{$\lif$} مهم دي، ځکه چې د وروستي سېکوېنټ د
تړلي فارمولونه اصلي عملګرونه دي؛ خو \LR{$\lnot$} د بل نښلوونکي د
چاپېريال دننه دے، نو وروسته به ئې وڅېړو۔ له دې امله د وروستي
استنتاج د منطقي قاعدې امکانونه \LeftR{\lor} او \RightR{\lif} دي۔
په رښتيا هره يوه غوره کولے شو، خو \RightR{\lif} غوره کوو، يوازې
د دې دپاره چې په دوو څانګو وېشل لږ وروسته کړو۔ د هغو
\RightR{\lif} استنتاجونو د بڼې له مخې چې دا لاندې سېکوېنټ ورکولے
شي، بڼه بايد داسې وي:
\begin{prooftree}
\AxiomC{}
\UnaryInf$ !A, \lnot !A \lor !B \fCenter !B $
\RightLabel{\RightR{\lif}} \UnaryInf$ \lnot !A \lor !B \fCenter !A \lif !B $
\end{prooftree}
\olpseganchor{OLP-0075-B009}{end}

\olpseganchor{OLP-0075-B010}{start}
که \LR{$\lnot !A \lor !B$} د مقدم بهرني پاے ته يوسو، د
\LeftR{\lor} قاعده کارولے شو۔ د شېما له مخې دا بايد په لاندې ډول
پر دوو پورتنيو سېکوېنټونو ووېشل شي۔
د کيڼې خوا د تبادلې قاعده د مقدم د دوو ګاونډيو فارمولونو ترتيب بدلوي۔
\begin{prooftree}
\AxiomC{}
\UnaryInf$\lnot !A, !A \fCenter !B$
\AxiomC{}
\UnaryInf$!B, !A \fCenter !B$
\RightLabel{\LeftR{\lor}}
\BinaryInf$ \lnot !A \lor !B, !A \fCenter !B $
\RightLabel{\LeftR{\Exchange}}
\UnaryInf$ !A, \lnot !A \lor !B \fCenter !B $
\RightLabel{\RightR{\lif}}
\UnaryInf$ \lnot !A \lor !B \fCenter !A \lif !B $
\end{prooftree}
راياد کړئ چې موږ هڅه کوو تر لومړنيو سېکوېنټونو پورې په پورته لوري
لار پيدا کړو؛ داسې ښکاري چې ډېر ورنژدې يو!{} ښۍ څانګه له لومړني
سېکوېنټ څخه يوازې يو کمزوري کول او يوه تبادله لرې ده، او بيا
بشپړېږي:
\begin{prooftree}
\AxiomC{}
\UnaryInf$\lnot !A, !A \fCenter !B$
\Axiom$!B \fCenter !B$
\RightLabel{\LeftR{\Weakening}}
\UnaryInf$!A, !B \fCenter !B$
\RightLabel{\LeftR{\Exchange}}
\UnaryInf$!B, !A \fCenter !B$
\RightLabel{\LeftR{\lor}}
\BinaryInf$\lnot !A \lor !B, !A \fCenter !B $
\RightLabel{\LeftR{\Exchange}}
\UnaryInf$ !A, \lnot !A \lor !B \fCenter !B $
\RightLabel{\RightR{\lif}}
\UnaryInf$ \lnot !A \lor !B \fCenter !A \lif !B $
\end{prooftree}
\olpseganchor{OLP-0075-B010}{end}

\olpseganchor{OLP-0075-B011}{start}
اوس کيڼه څانګه ګورو۔ په هره تړلے فارمول کښې يوازينے منطقي
نښلوونکے د مقدم د تړلي فارمولونه د \LR{$\lnot$} نښه ده، نو د
\LeftR{\lnot} قاعدې يوه کارونه ګورو۔
\begin{prooftree}
\AxiomC{}
\UnaryInf$ !A \fCenter !B, !A$
\RightLabel{\LeftR{\lnot}}
\UnaryInf$\lnot !A, !A \fCenter !B$
\Axiom$!B \fCenter !B$
\RightLabel{\LeftR{\Weakening}}
\UnaryInf$!A, !B \fCenter !B$
\RightLabel{\LeftR{\Exchange}}
\UnaryInf$!B, !A \fCenter !B$
\RightLabel{\LeftR{\lor}}
\BinaryInf$\lnot !A \lor !B, !A \fCenter !B $
\RightLabel{\LeftR{\Exchange}}
\UnaryInf$ !A, \lnot !A \lor !B \fCenter !B $
\RightLabel{\RightR{\lif}}
\UnaryInf$ \lnot !A \lor !B \fCenter !A \lif !B $
\end{prooftree}
لکه ښۍ څانګه چې مو بشپړه کړه، له دې کيڼې څانګې څخه هم يوازې يو
کمزوري کول او يوه تبادله لرې يو۔
\begin{prooftree}
\Axiom$!A \fCenter !A$
\RightLabel{\RightR{\Weakening}}
\UnaryInf$ !A \fCenter !A, !B$
\RightLabel{\RightR{\Exchange}}
\UnaryInf$ !A \fCenter !B, !A$
\RightLabel{\LeftR{\lnot}}
\UnaryInf$\lnot !A, !A \fCenter !B$
\Axiom$!B \fCenter !B$
\RightLabel{\LeftR{\Weakening}}
\UnaryInf$!A, !B \fCenter !B$
\RightLabel{\LeftR{\Exchange}}
\UnaryInf$!B, !A \fCenter !B$
\RightLabel{\LeftR{\lor}}
\BinaryInf$\lnot !A \lor !B, !A \fCenter !B $
\RightLabel{\LeftR{\Exchange}}
\UnaryInf$ !A, \lnot !A \lor !B \fCenter !B $
\RightLabel{\RightR{\lif}}
\UnaryInf$ \lnot !A \lor !B \fCenter !A \lif !B $
\end{prooftree}
\end{ex}
\olpseganchor{OLP-0075-B011}{end}

\olpseganchor{OLP-0075-B012}{start}
\begin{ex}
د سېکوېنټ \LR{$\lnot !A \lor \lnot !B
\Sequent \lnot (!A \land !B)$} يو \LR{$\Log{LK}$}-اشتقاق ورکړئ۔
\olpseganchor{OLP-0075-B012}{end}

\olpseganchor{OLP-0075-B013}{start}
د پورته طريقو په کارولو، غوښتل شوے وروستنے سېکوېنټ په لاندې برخه
کښې ليکو۔
\begin{prooftree}
\AxiomC{}
\UnaryInf$ \lnot !A \lor \lnot !B \fCenter \lnot (!A \land !B) $
\end{prooftree}
په وروستي سېکوېنټ کښې د تړلي فارمولونه شته اصلي نښلوونکي د \LR{$\lor$}
او \LR{$\lnot$} نښې دي۔ دلته هم \LeftR{\lor} او هم \RightR{\lnot}
کار ورکوي، خو په \RightR{\lnot} پيل کوو، ځکه چې د يوې شېبې دپاره
پر دوو څانګو وېشل ځنډوي:
\begin{prooftree}
\AxiomC{}
\UnaryInf$!A \land !B, \lnot !A \lor \lnot !B \fCenter $
\RightLabel{\RightR{\lnot}}
\UnaryInf$\lnot !A \lor \lnot !B \fCenter \lnot (!A \land !B)$
\end{prooftree}
اوس د \LeftR{\land} او \LeftR{\lor} تر منځ ټاکنه لرو۔ ګورو چې د
\LeftR{\land} قاعدې په کارولو څه کېږي: يا په سېکوېنټ \LR{$!A,
\lnot !A \lor \lnot !B \Sequent \quad$} او يا په سېکوېنټ \LR{$!B, \lnot !A
\lor \lnot !B \Sequent \quad$} پيل کولے شو۔ ځکه چې اشتقاق د
\LR{$!A$} او \LR{$!B$} له پلوه متناظر دے، لومړنے غوره کوو۔
\begin{prooftree}
\AxiomC{}
\UnaryInf$!A, \lnot !A \lor \lnot !B \fCenter $
\RightLabel{\LeftR{\land}}
\UnaryInf$!A \land !B, \lnot !A \lor \lnot !B \fCenter $
\RightLabel{\RightR{\lnot}}
\UnaryInf$\lnot !A \lor \lnot !B \fCenter \lnot (!A \land !B)$
\end{prooftree}
د اشتقاق د ډکولو په دوام وينو چې له يوې ستونزې سره مخ کېږو:
\begin{prooftree}
\Axiom$!A \fCenter !A$
\RightLabel{\LeftR{\lnot}}
\UnaryInf$ \lnot !A, !A \fCenter$
\AxiomC{}
\RightLabel{?}
\UnaryInf$!A \fCenter !B$
\RightLabel{\LeftR{\lnot}}
\UnaryInf$ \lnot !B, !A \fCenter$
\RightLabel{\LeftR{\lor}}
\BinaryInf$\lnot !A \lor \lnot !B, !A \fCenter $
\RightLabel{\LeftR{\Exchange}}
\UnaryInf$!A, \lnot !A \lor \lnot !B \fCenter $
\RightLabel{\LeftR{\land}}
\UnaryInf$!A \land !B, \lnot !A \lor \lnot !B \fCenter $
\RightLabel{\RightR{\lnot}}
\UnaryInf$\lnot !A \lor \lnot !B \fCenter \lnot (!A \land !B)$
\end{prooftree}
د ښۍ څانګې پورته برخه نوره نۀ شي راکمېدلے او د جوړښتي استنتاجونو
له لارې لومړني سېکوېنټ ته نۀ شي رسول کېدے؛ نو دا سمه لار نۀ ده۔
له دې امله پورته د \LeftR{\land} قاعدې کارول تېروتنه وه۔ بېرته
مخکيني حالت ته ځو او پر ځاے ئې \LeftR{\lor} قاعده پلې کوو، نو
ترلاسه کوو:
\begin{prooftree}
\AxiomC{}
\UnaryInf$\lnot !A, !A \land !B \fCenter $
\olpseganchor{OLP-0075-B013}{end}

\olpseganchor{OLP-0075-B014}{start}
\AxiomC{}
\UnaryInf$\lnot !B, !A \land !B \fCenter $
\olpseganchor{OLP-0075-B014}{end}

\olpseganchor{OLP-0075-B015}{start}
\RightLabel{\LeftR{\lor}}
\BinaryInf$\lnot !A \lor \lnot !B, !A \land !B \fCenter $
\RightLabel{\LeftR{\Exchange}}
\UnaryInf$!A \land !B, \lnot !A \lor \lnot !B \fCenter $
\RightLabel{\RightR{\lnot}}
\UnaryInf$\lnot !A \lor \lnot !B \fCenter \lnot (!A \land !B)$
\end{prooftree}
هره څانګه د مخکښې په شان بشپړه کړو، نو ترلاسه کوو:
\begin{prooftree}
\Axiom$ !A \fCenter!A$
\RightLabel{\LeftR{\land}}
\UnaryInf$!A \land !B \fCenter !A$
\RightLabel{\LeftR{\lnot}}
\UnaryInf$\lnot !A, !A \land !B \fCenter $
\olpseganchor{OLP-0075-B015}{end}

\olpseganchor{OLP-0075-B016}{start}
\Axiom$ !B \fCenter !B$
\RightLabel{\LeftR{\land}}
\UnaryInf$!A \land !B \fCenter !B$
\RightLabel{\LeftR{\lnot}}
\UnaryInf$\lnot !B, !A \land !B \fCenter $
\olpseganchor{OLP-0075-B016}{end}

\olpseganchor{OLP-0075-B017}{start}
\RightLabel{\LeftR{\lor}}
\BinaryInf$\lnot !A \lor \lnot !B, !A \land !B \fCenter $
\RightLabel{\LeftR{\Exchange}}
\UnaryInf$!A \land !B, \lnot !A \lor \lnot !B \fCenter $
\RightLabel{\RightR{\lnot}}
\UnaryInf$\lnot !A \lor \lnot !B \fCenter \lnot (!A \land !B)$
\end{prooftree}
(په دغو ګامونو کښې د \LR{$\land$} قاعدې د \LR{$\lnot$} له قاعدو څخه لاندې
هم پلې کېدے شوې او بيا به هم سم اشتقاق ترلاسه شوے و۔)
\end{ex}
\olpseganchor{OLP-0075-B017}{end}

\olpseganchor{OLP-0075-B018}{start}
\begin{ex}
تر اوسه مو د انقباض قاعده نۀ ده کارولې، خو کله کله اړينه وي۔ دا
ئې يوه بېلګه ده۔ فرض کړئ غواړو \LR{$\quad \Sequent !A \lor \lnot !A$}
ثابت کړو۔ د \LR{$\RightR{\lor}$} سرچپه پلي کول به له دغو دوو
اشتقاقونه څخه يو راکړي:
\begin{prooftree}
\AxiomC{}
\UnaryInf$ \fCenter !A$
\RightLabel{\RightR{\lor}}
\UnaryInf$ \fCenter !A \lor \lnot !A$
\DisplayProof\qquad\bottomAlignProof
\AxiomC{}
\UnaryInf$!A \fCenter $
\RightLabel{\RightR{\lnot}}
\UnaryInf$ \fCenter \lnot !A$
\RightLabel{\RightR{\lor}}
\UnaryInf$ \fCenter !A \lor \lnot !A$
\end{prooftree}
البته له دغو څخه هېڅ يو په لومړني سېکوېنټ نۀ پاے ته رسېږي۔ چل دا
دے چې پوه شو د انقباض قاعده د تړلے فارمول دوه نقلونه سره يو
کولے شي---او کله چې د ثبوت لټون کوو، يعنې له لاندې څخه پورته ځو،
د \LR{$!A \lor \lnot !A$} يوه کاپي په مقدمه کښې ساتلے شو، مثلاً:
\begin{prooftree}
\AxiomC{}
\UnaryInf$ \fCenter !A \lor \lnot !A, !A$
\RightLabel{\RightR{\lor}}
\UnaryInf$ \fCenter !A \lor \lnot !A, !A \lor \lnot !A$
\RightLabel{\RightR{\Contraction}}
\UnaryInf$ \fCenter !A \lor \lnot !A$
\end{prooftree}
اوس \LR{$\RightR{\lor}$} دويم ځل کارولے شو او~\LR{$\lnot !A$} هم ترلاسه کوو،
چې بشپړ اشتقاق ته رسېږي۔
\begin{prooftree}
\Axiom$!A \fCenter !A$
\RightLabel{\RightR{\lnot}}
\UnaryInf$\fCenter !A, \lnot !A$
\RightLabel{\RightR{\lor}}
\UnaryInf$\fCenter !A, !A \lor \lnot !A$
\RightLabel{\RightR{\Exchange}}
\UnaryInf$ \fCenter !A \lor \lnot !A, !A$
\RightLabel{\RightR{\lor}}
\UnaryInf$ \fCenter !A \lor \lnot !A, !A \lor \lnot !A$
\RightLabel{\RightR{\Contraction}}
\UnaryInf$ \fCenter !A \lor \lnot !A$
\end{prooftree}
\end{ex}
\olpseganchor{OLP-0075-B018}{end}

\olpseganchor{OLP-0075-B019}{start}
\begin{prob}
د لاندې سېکوېنټونو اشتقاقونه ورکړئ:
\begin{enumerate}
\item \LR{$!A \land (!B \land !C) \Sequent (!A \land !B) \land !C$}.
\item \LR{$!A \lor (!B \lor !C) \Sequent (!A \lor !B) \lor !C$}.
\item \LR{$!A \lif (!B \lif !C) \Sequent !B \lif (!A \lif !C)$}.
\item \LR{$!A \Sequent \lnot\lnot !A$}.
\end{enumerate}
\end{prob}
\olpseganchor{OLP-0075-B019}{end}

\olpseganchor{OLP-0075-B020}{start}
\begin{prob}
د لاندې سېکوېنټونو اشتقاقونه ورکړئ:
\begin{enumerate}
\item \LR{$(!A \lor !B) \lif !C \Sequent !A \lif !C$}.
\item \LR{$(!A \lif !C) \land (!B \lif !C) \Sequent (!A \lor !B) \lif !C$}.
\item \LR{$\Sequent \lnot(!A \land \lnot !A)$}.
\item \LR{$!B \lif !A \Sequent \lnot !A \lif \lnot !B$}.
\item \LR{$\Sequent (!A \lif \lnot !A) \lif \lnot !A$}.
\item \LR{$\Sequent \lnot(!A \lif !B) \lif \lnot !B$}.
\item \LR{$!A \lif !C \Sequent \lnot (!A \land \lnot !C)$}.
\item \LR{$!A \land \lnot !C \Sequent \lnot (!A \lif !C)$}.
\item \LR{$!A \lor !B, \lnot !B \Sequent !A$}.
\item \LR{$\lnot !A \lor \lnot !B \Sequent \lnot(!A \land !B)$}.
\item \LR{$\Sequent (\lnot !A \land \lnot !B) \lif\lnot(!A \lor !B)$}.
\item \LR{$\Sequent \lnot(!A \lor !B) \lif (\lnot !A \land \lnot !B)$}.
\end{enumerate}
\end{prob}
\olpseganchor{OLP-0075-B020}{end}

\olpseganchor{OLP-0075-B021}{start}
\begin{prob}
د لاندې سېکوېنټونو اشتقاقونه ورکړئ:
\begin{enumerate}
\item \LR{$\lnot(!A \lif !B) \Sequent !A$}.
\item \LR{$\lnot(!A \land !B) \Sequent \lnot !A \lor \lnot !B$}.
\item \LR{$!A \lif !B \Sequent \lnot !A \lor !B$}.
\item \LR{$\Sequent \lnot \lnot !A \lif !A$}.
\item \LR{$!A \lif !B, \lnot !A \lif !B \Sequent !B$}.
\item \LR{$(!A \land !B) \lif !C \Sequent (!A \lif !C) \lor (!B \lif !C)$}.
\item \LR{$(!A \lif !B) \lif !A \Sequent !A$}.
\item \LR{$\Sequent (!A \lif !B) \lor (!B \lif !C)$}.
\end{enumerate}
(دا ټول د~\LR{$\RightR{\Contraction}$} قاعدې ته اړتيا لري۔)
\end{prob}
\olpseganchor{OLP-0075-B021}{end}

\olpseganchor{OLP-0076-B004}{start}
\olfileid{fol}{seq}{prq}
\olpseganchor{OLP-0076-B004}{end}

\olpseganchor{OLP-0076-B005}{start}
\olsection{له سورونو سره اشتقاقونه}
\olpseganchor{OLP-0076-B005}{end}

\olpseganchor{OLP-0076-B006}{start}
\begin{ex}
د سېکوېنټ \LR{$\lexists[x][\lnot !A(x)]
\Sequent \lnot \lforall[x][!A(x)]$} يو
\LR{$\Log{LK}$}-اشتقاق ورکړئ۔
\olpseganchor{OLP-0076-B006}{end}

\olpseganchor{OLP-0076-B007}{start}
کله چې له سورونو سره کار کوو، بايد ډاډمن شو چې د ځانګړي متغير شرط
نۀ ماتوو، او کله کله د دې دپاره اړ يو چې د ځينو استنتاجونو د پلي
کولو ترتيب بدل کړو۔ په عمومي توګه دا ګټوره ده چې لومړے هغه قاعدې
وڅېړو چې د ځانګړي متغير د شرط تابع دي (په بشپړ شوي ثبوت کښې به
لاندې وي)۔ همدارنګه، غوره ده لږ وړاندې وګورو او اټکل وکړو چې
لومړنے سېکوېنټ به څنګه وي۔ زموږ په حالت کښې بايد د
\LR{$!A(a) \Sequent !A(a)$} په شان وي۔ يعنې کله چې د سورونو قاعدې
“سرچپه” کوو، بايد د \LR{$\lforall$} او \LR{$\lexists$} دواړو قاعدو دپاره
هماغه اصطلاح---چې \LR{$a$} به ئې وبولو---غوره کړو۔ که د هرې قاعدې
دپاره بېله اصطلاح غوره کړو، د \LR{$!A(a) \Sequent !A(b)$} په شان څه
ته رسېږو، چې البته د اشتقاق وړ نۀ دے۔
\olpseganchor{OLP-0076-B007}{end}

\olpseganchor{OLP-0076-B008}{start}
د معمول په شان په دې ليکلو پيل کوو:
\begin{prooftree}
\AxiomC{}
\UnaryInf$\lexists[x][\lnot !A(x)] \fCenter \lnot \lforall[x][!A(x)]$
\end{prooftree}
يا د \LeftR{\exists} قاعده پلې کولے شو او يا د
\RightR{\lnot} قاعده۔ ځکه چې \LeftR{\exists} د ځانګړي متغير د
شرط تابع ده، ښه ده چې دا ژر تر ژره وڅېړو؛ نو لومړے همدا پلې کوو۔
\begin{prooftree}
\AxiomC{}
\UnaryInf$ \lnot !A(a) \fCenter \lnot \lforall[x][!A(x)]$
\RightLabel{\LeftR{\lexists}}
\UnaryInf$ \lexists[x][\lnot !A(x)] \fCenter \lnot \lforall[x][!A(x)]$
\end{prooftree}
د \LeftR{\lnot} او \RightR{\lnot} قاعدو په سرچپه پلي کولو ترلاسه
کوو:
\begin{prooftree}
\AxiomC{}
\UnaryInf$\lforall[x][!A(x)] \fCenter !A(a)$
\RightLabel{\LeftR{\lnot}}
\UnaryInf$\lnot !A(a), \lforall[x][!A(x)] \fCenter $
\RightLabel{\LeftR{\Exchange}}
\UnaryInf$\lforall[x][!A(x)], \lnot !A(a) \fCenter $
\RightLabel{\RightR{\lnot}}
\UnaryInf$ \lnot !A(a) \fCenter \lnot \lforall[x] !A(x)$
\RightLabel{\LeftR{\lexists}}
\UnaryInf$ \lexists[x] \lnot !A(x) \fCenter \lnot \lforall[x] !A(x)$
\end{prooftree}
په دې ځاے کښې يوازينے امکان د \LeftR{\forall} قاعدې پلي کول دي۔
ځکه چې دا قاعده د ځانګړي متغير د بنديز تابع نۀ ده، نو کومه ستونزه
نشته۔ راياد کړئ چې غواړو يو لومړنے سېکوېنټ، يعنې د
\LR{$!A(a) \Sequent !A(a)$} په بڼه، ترلاسه کړو؛ نو د قاعدې د کارولو پر
مهال بايد \LR{$a$} د~\LR{$!A$} د ارګومېنټ په توګه غوره کړو۔
\begin{prooftree}
\Axiom$!A(a) \fCenter !A(a)$
\RightLabel{\LeftR{\lforall}}
\UnaryInf$\lforall[x][!A(x)] \fCenter !A(a)$
\RightLabel{\LeftR{\lnot}}
\UnaryInf$\lnot !A(a), \lforall[x][!A(x)] \fCenter $
\RightLabel{\LeftR{\Exchange}}
\UnaryInf$\lforall[x][!A(x)], \lnot !A(a) \fCenter $
\RightLabel{\RightR{\lnot}}
\UnaryInf$ \lnot !A(a) \fCenter \lnot \lforall[x][!A(x)]$
\RightLabel{\LeftR{\lexists}}
\UnaryInf$ \lexists[x][ \lnot !A(x)] \fCenter \lnot \lforall[x][!A(x)]$
\end{prooftree}
په تېره بيا له سورونو سره د کار پر مهال مهمه ده چې په دې ځاے کښې
بيا وګورو چې د ځانګړي متغير شرط خو نۀ دے مات شوے۔ ځکه چې زموږ له
کارول شوو قاعدو څخه يوازې \LeftR{\exists} د ځانګړي متغير د شرط
تابع وه، او ځانګړے متغير~\LR{$a$} د هغې په لاندې سېکوېنټ (وروستي
سېکوېنټ) کښې نۀ راځي، نو دا سم اشتقاق دے۔
\end{ex}
\olpseganchor{OLP-0076-B008}{end}

\olpseganchor{OLP-0076-B009}{start}
\begin{prob}
د لاندې سېکوېنټونو اشتقاقونه ورکړئ:
\begin{enumerate}
\item \LR{$\Sequent (\lforall[x][!A(x)] \land \lforall[y][!B(y)]) \lif
\lforall[z][(!A(z) \land !B(z))]$}.
\item \LR{$\Sequent (\lexists[x][!A(x)] \lor \lexists[y][!B(y)]) \lif
\lexists[z][(!A(z) \lor !B(z))]$}.
\item \LR{$\lforall[x][(!A(x) \lif !B)] \Sequent \lexists[y][!A(y)] \lif !B$}.
\item \LR{$\lforall[x][\lnot !A(x)] \Sequent \lnot\lexists[x][!A(x)]$}.
\item \LR{$\Sequent \lnot\lexists[x][!A(x)] \lif \lforall[x][\lnot !A(x)]$}.
\item \LR{$\Sequent \lnot\lexists[x][\lforall[y][((!A(x,y) \lif \lnot
!A(y,y)) \land (\lnot !A(y,y) \lif !A(x,y)))]]$}.
\end{enumerate}
\end{prob}
\olpseganchor{OLP-0076-B009}{end}

\olpseganchor{OLP-0076-B010}{start}
\begin{prob}
د لاندې سېکوېنټونو اشتقاقونه ورکړئ:
\begin{enumerate}
\item \LR{$\Sequent \lnot\lforall[x][!A(x)] \lif \lexists[x][\lnot!A(x)]$}.
\item \LR{$(\lforall[x][!A(x)] \lif !B) \Sequent \lexists[y][(!A(y) \lif !B)]$}.
\item \LR{$\Sequent \lexists[x][(!A(x) \lif \lforall[y][!A(y)])]$}.
\end{enumerate}
(دا ټول د~\LR{$\RightR{\Contraction}$} قاعدې ته اړتيا لري۔)
\end{prob}
\olpseganchor{OLP-0076-B010}{end}

\olpseganchor{OLP-0077-B004}{start}
\begin{editorial}
  دا برخه د سېکوېنټ حساب د اشتقاق وړتيا د اړيکې او سازګارۍ
  تعريفونه راغونډوي۔
\end{editorial}
\olpseganchor{OLP-0077-B004}{end}

\olpseganchor{OLP-0077-B005}{start}
\olfileid{fol}{seq}{ptn}
\olpseganchor{OLP-0077-B005}{end}

\olpseganchor{OLP-0077-B006}{start}
\olsection{د ثبوت-تيوريکي مفاهيم}
\olpseganchor{OLP-0077-B006}{end}

\olpseganchor{OLP-0077-B007}{start}
\begin{explain}
لکه څنګه چې مو يو شمېر مهم معنايي مفاهيم (اعتبار، معنايي استلزام
او د صدق وړتيا) تعريف کړي دي، اوس ورسره سم \emph{ثبوت-تيوريکي
مفاهيم} تعريفوو۔ دا مفاهيم په جوړښتونه کښې د تړلي فارمولونه
د صدق پر بنسټ نۀ، بلکې د ځينو سېکوېنټونو د اشتقاق وړتيا يا
د اشتقاق نۀ وړتيا پر بنسټ تعريفېږي۔ دا يوه مهمه موندنه وه چې دغه
مفاهيم سره سمون خوري۔ د دې سمون منځپانګه د \emph{سموالي} او
\emph{بشپړتيا قضيه} ده۔
\end{explain}
\olpseganchor{OLP-0077-B007}{end}

\olpseganchor{OLP-0077-B008}{start}
\begin{defn}[قضيې]
تړلے فارمول~\LR{$!A$} هله \emph{قضيه} ده چې په~\LR{$\Log{LK}$} کښې د
سېکوېنټ \LR{$\quad \Sequent !A$} اشتقاق موجود وي۔ که \LR{$!A$}
قضيه وي، \LR{$\Proves !A$} ليکو؛ او که نۀ وي، \LR{$\Proves/ !A$} ليکو۔
\end{defn}
\olpseganchor{OLP-0077-B008}{end}

\olpseganchor{OLP-0077-B009}{start}
\begin{defn}[د اشتقاق وړتيا]
تړلے فارمول~\LR{$!A$} د تړلي فارمولونه له سټ~\LR{$\Gamma$} څخه
\emph{د اشتقاق وړ} ده، يعنې \LR{$\Gamma \Proves !A$}، هله او يوازې
هله چې يو متناهي فرعي سټ~\LR{$\Gamma_0 \subseteq \Gamma$} او په
\LR{$\Gamma_0$} کښې د تړلي فارمولونه يو لړ~\LR{$\Gamma_0'$} موجود وي، داسې
چې \LR{$\Log{LK}$} د \LR{$\Gamma_0' \Sequent !A$} سېکوېنټ اشتقاق کړي۔
که \LR{$!A$} له \LR{$\Gamma$} څخه د اشتقاق وړ نۀ وي، \LR{$\Gamma \Proves/ !A$}
ليکو۔
\end{defn}
\olpseganchor{OLP-0077-B009}{end}

\olpseganchor{OLP-0077-B010}{start}
د انقباض، کمزوري کولو او تبادلې د قاعدو له امله، په~\LR{$\Gamma_0'$}
کښې د تړلي فارمولونه ترتيب او شمېر اهميت نۀ لري: که سېکوېنټ
\LR{$\Gamma_0' \Sequent !A$} د اشتقاق وړ وي، نو د هر داسې
\LR{$\Gamma_0''$} دپاره \LR{$\Gamma_0'' \Sequent !A$} هم د اشتقاق وړ دے
چې د~\LR{$\Gamma_0'$} هماغه تړلي فارمولونه لري۔ مثلاً که
\LR{$\Gamma_0 = \{!B, !C\}$} وي، نو هم
\LR{$\Gamma_0' = \tuple{!B, !B, !C}$} او هم
\LR{$\Gamma_0'' = \tuple{!C, !C, !B}$} داسې لړونه دي چې يوازې د
\LR{$\Gamma_0$} تړلي فارمولونه لري۔ که د يوه لړ لرونکے سېکوېنټ
د اشتقاق وړ وي، د بل لړ لرونکے هم دے، مثلاً:
\begin{prooftree}
  \AxiomC{}
  \Deduce$!B, !B, !C \fCenter !A$
  \RightLabel{\LeftR{\Contraction}}
  \UnaryInf$!B, !C \fCenter !A$
  \RightLabel{\LeftR{\Exchange}}
  \UnaryInf$!C, !B \fCenter !A$
  \RightLabel{\LeftR{\Weakening}}
  \UnaryInf$!C, !C, !B \fCenter !A$
\end{prooftree}
له دې وروسته به وايو چې که \LR{$\Gamma_0$} د تړلي فارمولونه يو متناهي
سټ وي، نو \LR{$\Gamma_0 \Sequent !A$} هر هغه سېکوېنټ دے چې مقدم ئې
په~\LR{$\Gamma_0$} کښې د تړلي فارمولونه يو لړ وي؛ او د اړتيا په وخت کښې
انقباضونه، تبادلې او کمزوري کول په ناويلې توګه ورګډوو۔
\olpseganchor{OLP-0077-B010}{end}

\olpseganchor{OLP-0077-B011}{start}
\begin{defn}[سازګاري]
د جملو سټ~\LR{$\Gamma$} هله او يوازې هله \emph{ناسازګار} دے چې يو
متناهي فرعي سټ~\LR{$\Gamma_0 \subseteq \Gamma$} موجود وي، داسې چې
\LR{$\Log{LK}$} د \LR{$\Gamma_0 \Sequent \quad$} سېکوېنټ اشتقاق کړي۔
که \LR{$\Gamma$} ناسازګار نۀ وي، يعنې که د هر متناهي
\LR{$\Gamma_0 \subseteq \Gamma$} دپاره \LR{$\Log{LK}$} د
\LR{$\Gamma_0 \Sequent \quad$} سېکوېنټ اشتقاق نۀ کړي، نو هغه
\emph{سازګار} بولو۔
\end{defn}
\olpseganchor{OLP-0077-B011}{end}

\olpseganchor{OLP-0077-B012}{start}
\begin{prop}[انعکاسيت]
\ollabel{prop:reflexivity}
که \LR{$!A \in \Gamma$} وي، نو \LR{$\Gamma \Proves !A$}۔
\end{prop}
\olpseganchor{OLP-0077-B012}{end}

\olpseganchor{OLP-0077-B013}{start}
\begin{proof}
لومړنے سېکوېنټ \LR{$!A \Sequent !A$} د اشتقاق وړ دے، او
\LR{$\{!A\} \subseteq \Gamma$}۔
\end{proof}
\olpseganchor{OLP-0077-B013}{end}
  
\olpseganchor{OLP-0077-B014}{start}
\begin{prop}[يکنواختي]
\ollabel{prop:monotonicity}
که \LR{$\Gamma \subseteq \Delta$} او \LR{$\Gamma \Proves !A$} وي، نو
\LR{$\Delta \Proves !A$}۔
\end{prop}
\olpseganchor{OLP-0077-B014}{end}

\olpseganchor{OLP-0077-B015}{start}
\begin{proof}
فرض کړئ \LR{$\Gamma \Proves !A$}، يعنې يو متناهي
\LR{$\Gamma_0 \subseteq \Gamma$} شته چې
\LR{$\Gamma_0 \Sequent !A$} د اشتقاق وړ دے۔ ځکه چې
\LR{$\Gamma \subseteq \Delta$}، نو \LR{$\Gamma_0$} د~\LR{$\Delta$} هم يو متناهي
فرعي سټ دے۔ له دې امله د \LR{$\Gamma_0 \Sequent !A$} اشتقاق
هم \LR{$\Delta \Proves !A$} ښيي۔
\end{proof}
\olpseganchor{OLP-0077-B015}{end}

\olpseganchor{OLP-0077-B016}{start}
\begin{prop}[انتقاليت]
\ollabel{prop:transitivity}
که \LR{$\Gamma \Proves !A$} او \LR{$\{!A\} \cup \Delta \Proves !B$} وي، نو
\LR{$\Gamma \cup \Delta \Proves !B$}۔
\end{prop}
\olpseganchor{OLP-0077-B016}{end}

\olpseganchor{OLP-0077-B017}{start}
\begin{proof}
که \LR{$\Gamma \Proves !A$} وي، يو متناهي
\LR{$\Gamma_0 \subseteq \Gamma$} او د
\LR{$\Gamma_0 \Sequent !A$} اشتقاق~\LR{$\pi_0$} شته۔ که
\LR{$\{!A\} \cup \Delta \Proves !B$} وي، نو د کوم متناهي فرعي سټ
\LR{$\Delta_0 \subseteq \Delta$} دپاره د
\LR{$!A, \Delta_0 \Sequent !B$} اشتقاق~\LR{$\pi_1$} شته۔ لاندې
اشتقاق وګورئ:
\begin{prooftree}
  \AxiomC{}
  \RightLabel{\LR{$\pi_0$}}
  \Deduce$\Gamma_0 \fCenter !A$
  \AxiomC{}
  \RightLabel{\LR{$\pi_1$}}
  \Deduce$!A, \Delta_0 \fCenter !B$
  \RightLabel{\Cut}
  \BinaryInf$\Gamma_0, \Delta_0 \fCenter !B$
\end{prooftree}
ځکه چې \LR{$\Gamma_0 \cup \Delta_0 \subseteq \Gamma \cup \Delta$}، نو
دا \LR{$\Gamma \cup \Delta \Proves !B$} ښيي۔
\end{proof}
\olpseganchor{OLP-0077-B017}{end}

\olpseganchor{OLP-0077-B018}{start}
پام وکړئ، په ځانګړې توګه د دې معنٰی دا ده چې که
\LR{$\Gamma \Proves !A$} او \LR{$!A \Proves !B$} وي، نو
\LR{$\Gamma \Proves !B$}۔ همدارنګه ترې دا هم راځي چې که
\LR{$!A_1, \dots, !A_n \Proves !B$} وي او د هر~\LR{$i$} دپاره
\LR{$\Gamma \Proves !A_i$} وي، نو \LR{$\Gamma \Proves !B$}۔
\olpseganchor{OLP-0077-B018}{end}

\olpseganchor{OLP-0077-B019}{start}
\begin{prop}
\ollabel{prop:incons}
\LR{$\Gamma$} هله او يوازې هله ناسازګار دے چې د هرې جملې~\LR{$!A$} دپاره
\LR{$\Gamma \Proves {!A}$} وي۔
\end{prop}
\olpseganchor{OLP-0077-B019}{end}

\olpseganchor{OLP-0077-B020}{start}
\begin{proof}
تمرين۔
\end{proof}
\olpseganchor{OLP-0077-B020}{end}

\olpseganchor{OLP-0077-B021}{start}
\begin{prob}
\olref[fol][seq][ptn]{prop:incons} ثابت کړئ۔
\end{prob}
\olpseganchor{OLP-0077-B021}{end}

\olpseganchor{OLP-0077-B022}{start}
\begin{prop}[متناهي شاهد]
  \ollabel{prop:proves-compact}
  \begin{enumerate}
  \item که \LR{$\Gamma \Proves !A$} وي، نو يو متناهي فرعي سټ
    \LR{$\Gamma_0 \subseteq \Gamma$} شته چې \LR{$\Gamma_0 \Proves !A$}۔
  \item که د~\LR{$\Gamma$} هر متناهي فرعي سټ سازګار وي، نو \LR{$\Gamma$}
    سازګار دے۔
  \end{enumerate}
\end{prop}
\olpseganchor{OLP-0077-B022}{end}

\olpseganchor{OLP-0077-B023}{start}
\begin{proof}
  \begin{enumerate}
    \item که \LR{$\Gamma \Proves !A$} وي، نو يو متناهي فرعي سټ
      \LR{$\Gamma_0 \subseteq \Gamma$} شته، داسې چې د
      \LR{$\Gamma_0 \Sequent !A$} سېکوېنټ اشتقاق لري۔ له دې امله
      \LR{$\Gamma_0 \Proves !A$}۔
    \item که \LR{$\Gamma$} ناسازګار وي، نو يو متناهي فرعي سټ
      \LR{$\Gamma_0 \subseteq \Gamma$} شته، داسې چې \LR{$\Log{LK}$} د
      \LR{$\Gamma_0 \Sequent \quad$} سېکوېنټ اشتقاق کړي۔ خو بيا
      \LR{$\Gamma_0$} د~\LR{$\Gamma$} يو ناسازګار متناهي فرعي سټ دے۔
  \end{enumerate}
\end{proof}
\olpseganchor{OLP-0077-B023}{end}

\olpseganchor{OLP-0078-B004}{start}
\olfileid{fol}{seq}{prv}
\olpseganchor{OLP-0078-B004}{end}

\olpseganchor{OLP-0078-B005}{start}
\olsection{د اشتقاق وړتيا او سازګاري}
\olpseganchor{OLP-0078-B005}{end}

\olpseganchor{OLP-0078-B006}{start}
اوس به د اشتقاق وړتيا اړيکې يو شمېر خاصيتونه ثابت کړو۔ دا
خاصيتونه په خپلواک ډول هم په زړه پورې دي، خو هر يو به د بشپړتيا د
قضيې په ثبوت کښې ونډه ولري۔
\olpseganchor{OLP-0078-B006}{end}

\olpseganchor{OLP-0078-B007}{start}
\begin{prop}\ollabel{prop:provability-contr}
  که \LR{$\Gamma \Proves !A$} وي او \LR{$\Gamma \cup \{!A\}$} ناسازګار وي،
  نو \LR{$\Gamma$} ناسازګار دے۔
\end{prop}
\olpseganchor{OLP-0078-B007}{end}

\olpseganchor{OLP-0078-B008}{start}
\begin{proof}
داسې متناهي \LR{$\Gamma_0$} او \LR{$\Gamma_1 \subseteq \Gamma$} شته چې
\LR{$\Log{LK}$} د \LR{$\Gamma_0 \Sequent !A$} او
\LR{$!A, \Gamma_1 \Sequent \quad$} سېکوېنټونه اشتقاق کړي۔
د \LR{$\Gamma_0 \Sequent !A$} د \LR{$\Log{LK}$}-اشتقاق نوم~\LR{$\pi_0$}
او د \LR{$\Gamma_1, !A \Sequent \quad$} د
\LR{$\Log{LK}$}-اشتقاق نوم~\LR{$\pi_1$} کېږدو۔ بيا لاندې
اشتقاق کولے شو:
\begin{prooftree}
\AxiomC{}
\RightLabel{\LR{$\pi_0$}}
\Deduce$ \Gamma_0 \fCenter !A $
\AxiomC{}
\RightLabel{\LR{$\pi_1$}}
\Deduce$!A, \Gamma_1 \fCenter $
\RightLabel{\Cut}
\BinaryInf$ \Gamma_0,\Gamma_1 \fCenter $
\end{prooftree}
\olpseganchor{OLP-0078-B008}{end}

\olpseganchor{OLP-0078-B009}{start}
ځکه چې \LR{$\Gamma_0 \subseteq \Gamma$} او
\LR{$\Gamma_1 \subseteq \Gamma$}، نو
\LR{$\Gamma_0 \cup \Gamma_1 \subseteq \Gamma$}؛ له دې امله \LR{$\Gamma$}
ناسازګار دے۔
\end{proof}
\olpseganchor{OLP-0078-B009}{end}

\olpseganchor{OLP-0078-B010}{start}
\begin{prop}
\ollabel{prop:prov-incons}
\LR{$\Gamma \Proves !A$} هله او يوازې هله چې
\LR{$\Gamma \cup \{\lnot !A\}$} ناسازګار وي۔
\end{prop}
\olpseganchor{OLP-0078-B010}{end}

\olpseganchor{OLP-0078-B011}{start}
\begin{proof}
لومړے فرض کړئ \LR{$\Gamma \Proves !A$}، يعنې د
\LR{$\Gamma \Sequent !A$} اشتقاق~\LR{$\pi_0$} شته۔ د
\LR{$\LeftR{\lnot}$} قاعدې په ورزياتولو د
\LR{$\lnot !A, \Gamma \Sequent \quad$} اشتقاق ترلاسه کوو؛
يعنې \LR{$\Gamma \cup \{\lnot !A\}$} ناسازګار دے۔
\olpseganchor{OLP-0078-B011}{end}

\olpseganchor{OLP-0078-B012}{start}
که \LR{$\Gamma \cup \{\lnot !A\}$} ناسازګار وي، نو د
\LR{$\lnot !A, \Gamma \Sequent \quad$} اشتقاق~\LR{$\pi_1$} شته۔
لاندې د \LR{$\Gamma \Sequent !A$} اشتقاق دے:
\begin{prooftree}
  \Axiom$!A \fCenter !A$
  \RightLabel{\RightR{\lnot}}
  \UnaryInf$\fCenter !A, \lnot !A$
  \AxiomC{}
  \RightLabel{\LR{$\pi_1$}}
  \Deduce$\lnot !A, \Gamma \fCenter$
  \RightLabel{\Cut}
  \BinaryInf$\Gamma \fCenter !A$
\end{prooftree}
\end{proof}
\olpseganchor{OLP-0078-B012}{end}

\olpseganchor{OLP-0078-B013}{start}
\begin{prob}
ثابت کړئ چې \LR{$\Gamma \Proves \lnot !A$} هله او يوازې هله چې
\LR{$\Gamma \cup \{!A\}$} ناسازګار وي۔
\end{prob}
\olpseganchor{OLP-0078-B013}{end}

\olpseganchor{OLP-0078-B014}{start}
\begin{prop}\ollabel{prop:explicit-inc}
  که \LR{$\Gamma \Proves !A$} او \LR{$\lnot !A \in \Gamma$} وي، نو \LR{$\Gamma$}
  ناسازګار دے۔
\end{prop}
\olpseganchor{OLP-0078-B014}{end}

\olpseganchor{OLP-0078-B015}{start}
\begin{proof}
  فرض کړئ \LR{$\Gamma \Proves !A$} او \LR{$\lnot !A \in \Gamma$}۔ بيا د کوم
  سېکوېنټ \LR{$\Gamma_0 \Sequent !A$} اشتقاق~\LR{$\pi$} شته۔
  سېکوېنټ \LR{$\lnot !A, \Gamma_0 \Sequent \quad$} هم د اشتقاق وړ دے:
  \begin{prooftree}
    \AxiomC{}
    \RightLabel{\LR{$\pi$}}
    \Deduce$\Gamma_0 \fCenter !A$
    \Axiom$!A \fCenter !A$
    \RightLabel{\LeftR{\lnot}}
    \UnaryInf$\lnot !A, !A \fCenter$
    \RightLabel{\LeftR{\Exchange}}
    \UnaryInf$!A, \lnot !A \fCenter$
    \RightLabel{\Cut}
    \BinaryInf$\Gamma_0, \lnot !A \fCenter$
  \end{prooftree}
  ځکه چې \LR{$\lnot !A \in \Gamma$} او \LR{$\Gamma_0 \subseteq \Gamma$}،
  نو دا ښيي چې \LR{$\Gamma$} ناسازګار دے۔
\end{proof}
\olpseganchor{OLP-0078-B015}{end}

\olpseganchor{OLP-0078-B016}{start}
\begin{prop}\ollabel{prop:provability-exhaustive}
  که \LR{$\Gamma \cup \{!A\}$} او \LR{$\Gamma \cup \{\lnot !A\}$} دواړه
  ناسازګار وي، نو \LR{$\Gamma$} ناسازګار دے۔
\end{prop}
\olpseganchor{OLP-0078-B016}{end}

\olpseganchor{OLP-0078-B017}{start}
\begin{proof}
داسې متناهي سټونه \LR{$\Gamma_0 \subseteq \Gamma$} او
\LR{$\Gamma_1 \subseteq \Gamma$} او داسې
\LR{$\Log{LK}$}-اشتقاقونه \LR{$\pi_0$} او \LR{$\pi_1$} شته چې په همدې ترتيب
د \LR{$!A, \Gamma_0 \Sequent \quad$} او
\LR{$\lnot !A, \Gamma_1 \Sequent \quad$} اشتقاقونه دي۔ بيا لاندې
اشتقاق کولے شو:
\begin{prooftree}
\AxiomC{}
\RightLabel{\LR{$\pi_0$}}
\Deduce$ !A, \Gamma_0 \fCenter $
\RightLabel{\RightR{\lnot}}
\UnaryInf$ \Gamma_0 \fCenter \lnot !A$
\AxiomC{}
\RightLabel{\LR{$\pi_1$}}
\Deduce$\lnot !A, \Gamma_1 \fCenter  $
\RightLabel{\Cut}
\BinaryInf$ \Gamma_0, \Gamma_1 \fCenter $
\end{prooftree}
ځکه چې \LR{$\Gamma_0 \subseteq \Gamma$} او \LR{$\Gamma_1 \subseteq \Gamma$}،
نو \LR{$\Gamma_0 \cup \Gamma_1 \subseteq \Gamma$}۔ له دې امله \LR{$\Gamma$}
ناسازګار دے۔
\end{proof}
\olpseganchor{OLP-0078-B017}{end}

\olpseganchor{OLP-0079-B005}{start}
\olfileid{fol}{seq}{ppr}
\olpseganchor{OLP-0079-B005}{end}

\olpseganchor{OLP-0079-B006}{start}
\olsection{د اشتقاق وړتيا او قضييز نښلوونکي}
\olpseganchor{OLP-0079-B006}{end}

\olpseganchor{OLP-0079-B007}{start}
\begin{explain}
  موږ ثابتوو چې د سېکوېنټ حساب د اشتقاق وړتيا اړيکه~\LR{$\Proves$}
  دومره پياوړې ده چې د قضييزو نښلوونکو په اړه ځينې بنسټيز حقيقتونه
  ثابت کړي، لکه \LR{$!A \land !B \Proves !A$} او
  \LR{$!A, !A \lif !B \Proves !B$} (مودوس پونېنس)۔ دغه حقيقتونه د
  بشپړتيا د قضيې د ثبوت دپاره اړين دي۔
\end{explain}
\olpseganchor{OLP-0079-B007}{end}

\olpseganchor{OLP-0079-B008}{start}
\begin{prop}\ollabel{prop:provability-land}
  \begin{enumerate}
  \item \ollabel{prop:provability-land-left} هم
    \LR{$!A \land !B \Proves !A$} او هم
    \LR{$!A \land !B \Proves !B$}۔
  \item \ollabel{prop:provability-land-right}
    \LR{$!A, !B \Proves !A \land !B$}۔
  \end{enumerate}
\end{prop}
\olpseganchor{OLP-0079-B008}{end}

\olpseganchor{OLP-0079-B009}{start}
\begin{proof}
  \begin{enumerate}
  \item دواړه سېکوېنټونه \LR{$!A \land !B \Sequent !A$} او
    \LR{$!A \land !B \Sequent !B$} د اشتقاق وړ دي:
    \begin{prooftree}
      \Axiom$!A \fCenter !A$
      \RightLabel{\LeftR{\land}}
      \UnaryInf$!A \land !B \fCenter !A$
      \DisplayProof\qquad\bottomAlignProof
      \Axiom$!B \fCenter !B$
      \RightLabel{\LeftR{\land}}
      \UnaryInf$!A \land !B \fCenter !B$
    \end{prooftree}
    \item دا د سېکوېنټ \LR{$!A, !B \Sequent !A \land !B$} يو
    اشتقاق دے۔
    د دواړو مقدمو په کيڼه خوا کښې هماغه يوه لړۍ ده۔ دوه ګونې کرښې جوړښتي استنتاجونه لنډوي۔
    \begin{prooftree}
      \Axiom$!A \fCenter !A$
      \doubleLine
      \UnaryInf$!A, !B \fCenter !A$
      \Axiom$!B \fCenter !B$
      \doubleLine
      \UnaryInf$!A, !B \fCenter !B$
      \RightLabel{\RightR{\land}}
      \BinaryInf$!A, !B \fCenter !A \land !B$
    \end{prooftree}
  \end{enumerate}
\end{proof}
\olpseganchor{OLP-0079-B009}{end}

\olpseganchor{OLP-0079-B010}{start}
\begin{prop}\ollabel{prop:provability-lor}
  \begin{enumerate}
  \item \LR{$!A \lor !B, \lnot !A, \lnot !B$} ناسازګار دے۔
  \item هم \LR{$!A \Proves !A \lor !B$} او هم
    \LR{$!B \Proves !A \lor !B$}۔
  \end{enumerate}
\end{prop}
\olpseganchor{OLP-0079-B010}{end}

\olpseganchor{OLP-0079-B011}{start}
\begin{proof}
  \begin{enumerate}
  \item موږ د سېکوېنټ
    \LR{$!A \lor !B, \lnot !A, \lnot !B \Sequent$} يو
    اشتقاق ورکوو:
    \begin{prooftree}
      \Axiom$!A \fCenter !A$
      \RightLabel{\LeftR{\lnot}}
      \UnaryInf$\lnot !A, !A \fCenter$
      \doubleLine
      \UnaryInf$!A, \lnot !A, \lnot !B \fCenter$
      \Axiom$!B \fCenter !B$
      \RightLabel{\LeftR{\lnot}}
      \UnaryInf$\lnot !B, !B \fCenter$
      \doubleLine
      \UnaryInf$!B, \lnot !A, \lnot !B \fCenter$
      \RightLabel{\LeftR{\lor}}
      \BinaryInf$ !A \lor !B, \lnot !A, \lnot !B \fCenter $
    \end{prooftree}
    (راياد کړئ چې د استنتاج دوه ګونې کرښې څو کمزوري کول، انقباضونه
    او تبادلې ښيي۔)
  \item دواړه سېکوېنټونه \LR{$!A \Sequent !A \lor !B$} او
    \LR{$!B \Sequent !A \lor !B$} اشتقاقونه لري:
    \begin{prooftree}
      \Axiom$!A \fCenter !A$
      \RightLabel{\RightR{\lor}}
      \UnaryInf$!A \fCenter !A \lor !B$
      \DisplayProof\qquad\bottomAlignProof
      \Axiom$!B \fCenter !B$
      \RightLabel{\RightR{\lor}}
      \UnaryInf$!B \fCenter !A \lor !B$
    \end{prooftree}
  \end{enumerate}
\end{proof}
\olpseganchor{OLP-0079-B011}{end}

\olpseganchor{OLP-0079-B012}{start}
\begin{prop}\ollabel{prop:provability-lif}
  \begin{enumerate}
  \item \ollabel{prop:provability-lif-left}
    \LR{$!A, !A \lif !B \Proves !B$}۔
  \item \ollabel{prop:provability-lif-right}
    هم \LR{$\lnot !A \Proves !A \lif !B$} او هم
    \LR{$!B \Proves !A \lif !B$}۔
  \end{enumerate}
\end{prop}
\olpseganchor{OLP-0079-B012}{end}

\olpseganchor{OLP-0079-B013}{start}
\begin{proof}
  \begin{enumerate}
    \item سېکوېنټ \LR{$!A \lif !B, !A \Sequent !B$} د اشتقاق وړ دے:
      \begin{prooftree}
        \Axiom$!A \fCenter !A$
        \Axiom$!B \fCenter !B$
        \RightLabel{\LeftR{\lif}}
        \BinaryInf$!A \lif !B, !A  \fCenter !B$
      \end{prooftree}
    \item دواړه سېکوېنټونه
      \LR{$\lnot !A \Sequent !A \lif !B$} او
      \LR{$!B \Sequent !A \lif !B$} د اشتقاق وړ دي:
      \begin{prooftree}
        \Axiom$!A \fCenter !A$
        \RightLabel{\LeftR{\lnot}}
        \UnaryInf$\lnot !A, !A \fCenter$
        \RightLabel{\LeftR{\Exchange}}
        \UnaryInf$!A, \lnot !A \fCenter$
        \RightLabel{\RightR{\Weakening}}
        \UnaryInf$!A, \lnot !A \fCenter !B$
        \RightLabel{\RightR{\lif}}
        \UnaryInf$\lnot !A \fCenter !A \lif !B$
        \DisplayProof\qquad\bottomAlignProof
        \Axiom$!B \fCenter !B$
        \RightLabel{\LeftR{\Weakening}}
        \UnaryInf$!A, !B \fCenter !B$
        \RightLabel{\RightR{\lif}}
        \UnaryInf$!B \fCenter !A \lif !B$
      \end{prooftree}
  \end{enumerate}
\end{proof}
\olpseganchor{OLP-0079-B013}{end}

\olpseganchor{OLP-0080-B004}{start}
\olfileid{fol}{seq}{qpr}
\olsection{د اشتقاق وړتيا او سورونه}
\olpseganchor{OLP-0080-B004}{end}

\olpseganchor{OLP-0080-B005}{start}
\begin{explain}
  د بشپړتيا قضيه دا هم غواړي چې د سېکوېنټ حساب قاعدې په دې برخه
  کښې د~\LR{$\Proves$} په اړه ثابت شوي حقيقتونه راکړي۔
\end{explain}
\olpseganchor{OLP-0080-B005}{end}

\olpseganchor{OLP-0080-B006}{start}
\begin{thm}
\ollabel{thm:strong-generalization} که \LR{$c$} داسې ثابت وي چې په
\LR{$\Gamma$} يا \LR{$!A(x)$} کښې نۀ راځي، او \LR{$\Gamma \Proves !A(c)$} وي، نو
\LR{$\Gamma \Proves \lforall[x][!A(x)]$}۔
\end{thm}
\olpseganchor{OLP-0080-B006}{end}

\olpseganchor{OLP-0080-B007}{start}
\begin{proof}
د کوم متناهي \LR{$\Gamma_0 \subseteq \Gamma$} دپاره
\LR{$\Gamma_0 \Sequent !A(c)$} د \LR{$\Log{LK}$}-اشتقاق نوم~\LR{$\pi_0$}
کېږدئ۔ د \LR{$\RightR{\lforall}$} استنتاج په ورزياتولو د
\LR{$\Gamma_0 \Sequent \lforall[x][!A(x)]$} اشتقاق ترلاسه کوو؛
ځکه \LR{$c$} په \LR{$\Gamma$} يا \LR{$!A(x)$} کښې نۀ راځي، نو د ځانګړي متغير شرط
پوره کېږي۔
\end{proof}
\olpseganchor{OLP-0080-B007}{end}

\olpseganchor{OLP-0080-B008}{start}
\begin{prop}
\ollabel{prop:provability-quantifiers}
\begin{enumerate}
\item \LR{$!A(t) \Proves \lexists[x][!A(x)]$}.
\item \LR{$\lforall[x][!A(x)] \Proves !A(t)$}.
\end{enumerate}
\end{prop}
\olpseganchor{OLP-0080-B008}{end}

\olpseganchor{OLP-0080-B009}{start}
\begin{proof}
\begin{enumerate}
\item سېکوېنټ
  \LR{$!A(t) \Sequent \lexists[x][!A(x)]$} د اشتقاق وړ دے:
  \begin{prooftree}
    \Axiom$!A(t) \fCenter !A(t)$
    \RightLabel{\RightR{\lexists}}
    \UnaryInf$!A(t) \fCenter \lexists[x][!A(x)]$
    \end{prooftree}
\item سېکوېنټ
  \LR{$\lforall[x][!A(x)] \Sequent !A(t)$} د اشتقاق وړ دے:
  \begin{prooftree}
    \Axiom$!A(t) \fCenter !A(t)$
    \RightLabel{\LeftR{\lforall}}
    \UnaryInf$\lforall[x][!A(x)] \fCenter !A(t)$
    \end{prooftree}
\end{enumerate}
\end{proof}
\olpseganchor{OLP-0080-B009}{end}

\olpseganchor{OLP-0081-B004}{start}
\olfileid{fol}{seq}{sou}
\olpseganchor{OLP-0081-B004}{end}

\olpseganchor{OLP-0081-B005}{start}
\olsection{سموالے}
\olpseganchor{OLP-0081-B005}{end}

\olpseganchor{OLP-0081-B006}{start}
\begin{explain}
د سېکوېنټ حساب په شان اشتقاق نظام هله \emph{سم} دے چې
هغه څيزونه اشتقاق نۀ کړي چې په حقيقت کښې نۀ وي۔ نو سموالے د
اشتقاق نظامونو د تضمين شوې خونديتوب يو ډول خاصيت دے۔
د دې له مخې چې کوم ثبوت-تيوريکي خاصيت تر بحث لاندې وي، مثلاً غواړو
پوه شو چې:
\begin{enumerate}
\item هر د اشتقاق وړ~\LR{$!A$} معتبر ده؛
\item که تړلے فارمول له ځينو نورو څخه د اشتقاق وړ وي، د هغو
  معنايي پايله هم ده؛
\item که د تړلي فارمولونه يو سټ ناسازګار وي، د صدق وړ نۀ دے۔
\end{enumerate}
دا د اشتقاق نظام مهم خاصيتونه دي۔ که له دغو څخه کوم يو
سم نۀ وي، اشتقاق نظام نيمګړے دے---له حده زيات څيزونه به
اشتقاق کوي۔ له دې امله د اشتقاق نظام سموالے ثابتول
تر ټولو زيات اهميت لري۔
\olpseganchor{OLP-0081-B006}{end}

\olpseganchor{OLP-0081-B007}{start}
ځکه دغه ټول ثبوت-تيوريکي خاصيتونه په سېکوېنټ حساب کښې د ځينو
سېکوېنټونو د اشتقاق وړتيا له لارې تعريف شوي دي، د پورته
(۱)--(۳) ثابتول غواړي چې د د اشتقاق وړ سېکوېنټونو د معنايي
خاصيتونو په اړه يو څه ثابت کړو۔ لومړے به تعريف کړو چې د سېکوېنټ
\emph{معتبر} کېدل څه معنٰی لري، او بيا به وښيو چې هر
د اشتقاق وړ سېکوېنټ معتبر دے۔ وروسته (۱)--(۳) د دې پايلې د
لازمو نتيجو په توګه راځي۔
\end{explain}
\olpseganchor{OLP-0081-B007}{end}

\olpseganchor{OLP-0081-B008}{start}
\begin{defn}
جوړښت~\LR{$\Struct
  M$} يو سېکوېنټ
\LR{$\Gamma \Sequent \Delta$} هله او يوازې هله \emph{پوره کوي} چې يا
د کوم \LR{$!A \in \Gamma$} دپاره
\LR{$\Sat/{M}{!A}$} وي، يا د کوم
\LR{$!A \in \Delta$} دپاره
\LR{$\Sat{M}{!A}$} وي۔
\olpseganchor{OLP-0081-B008}{end}

\olpseganchor{OLP-0081-B009}{start}
يو سېکوېنټ هله او يوازې هله \emph{معتبر} دے چې هر
جوړښت~\LR{$\Struct
  M$} هغه پوره کړي۔
\end{defn}
\olpseganchor{OLP-0081-B009}{end}

\olpseganchor{OLP-0081-B010}{start}
\begin{thm}[سموالے]
\ollabel{thm:sequent-soundness} که \LR{$\Log{LK}$} د
\LR{$\Theta \Sequent \Xi$} سېکوېنټ اشتقاق کړي، نو
\LR{$\Theta \Sequent \Xi$} معتبر دے۔
\end{thm}
\olpseganchor{OLP-0081-B010}{end}

\olpseganchor{OLP-0081-B011}{start}
\begin{proof}
فرض کړئ \LR{$\pi$} د \LR{$\Theta \Sequent \Xi$} اشتقاق دے۔ موږ
په~\LR{$\pi$} کښې د استنتاجونو پر شمېر~\LR{$n$} استقرا کوو۔
\olpseganchor{OLP-0081-B011}{end}

\olpseganchor{OLP-0081-B012}{start}
که د استنتاجونو شمېر~\LR{$0$} وي، نو \LR{$\pi$} يوازې له يوه لومړني
سېکوېنټ څخه جوړ دے۔ هر لومړنے سېکوېنټ
\LR{$!A \Sequent !A$} په څرګند ډول معتبر دے؛ ځکه د هر
\LR{$\Struct M$} دپاره يا
\LR{$\Sat/{M}{!A}$} وي يا
\LR{$\Sat{M}{!A}$}۔
\olpseganchor{OLP-0081-B012}{end}

\olpseganchor{OLP-0081-B013}{start}
که د استنتاجونو شمېر له~0 څخه زيات وي، د تر ټولو لاندې استنتاج
د ډول له مخې حالتونه بېلوو۔ د استقرا د فرضيې له مخې منلے شو چې د
دغه استنتاج مقدمې معتبرې دي؛ ځکه د هرې مقدمې په اشتقاق
کښې د استنتاجونو شمېر له~\LR{$n$} څخه کم دے۔
\olpseganchor{OLP-0081-B013}{end}

\olpseganchor{OLP-0081-B014}{start}
لومړے هغه ممکن استنتاجونه څېړو چې يوازې يوه مقدمه لري۔
\begin{enumerate}
\item وروستنے استنتاج کمزوري کول دي۔ بيا
  \LR{$\Theta \Sequent \Xi$} يا \LR{$!A, \Gamma \Sequent \Delta$} دے
  (که وروستنے استنتاج \LeftR{\Weakening} وي) يا
  \LR{$\Gamma \Sequent \Delta, !A$} دے
  (که \RightR{\Weakening} وي)، او اشتقاق له لاندې دواړو
  بڼو څخه په يوه پاے ته رسېږي:
  \begin{prooftree}
    \AxiomC{}
    \Deduce$\Gamma \fCenter \Delta$
    \RightLabel{\LeftR{\Weakening}}
    \UnaryInf$!A, \Gamma \fCenter \Delta$
    \DisplayProof\qquad\bottomAlignProof
    \AxiomC{}
    \Deduce$\Gamma \fCenter \Delta$
    \RightLabel{\RightR{\Weakening}}
    \UnaryInf$\Gamma \fCenter \Delta, !A$
  \end{prooftree}
  د استقرا د فرضيې له مخې \LR{$\Gamma \Sequent \Delta$} معتبر دے؛ يعنې
  د هر جوړښت~\LR{$\Struct{M}$} دپاره يا کوم
  \LR{$!C \in \Gamma$} شته چې
  \LR{$\Sat/{M}{!C}$} وي، يا کوم
  \LR{$!C \in \Delta$} شته چې
  \LR{$\Sat{M}{!C}$} وي۔
\olpseganchor{OLP-0081-B014}{end}
  
\olpseganchor{OLP-0081-B015}{start}
  که د کوم \LR{$!C \in \Gamma$} دپاره
  \LR{$\Sat/{M}{!C}$} وي، نو
  \LR{$!C \in \Theta$} هم دے؛ ځکه \LR{$\Theta = !A, \Gamma$}، نو د کوم
  \LR{$!C \in \Theta$} دپاره
  \LR{$\Sat/{M}{!C}$} دے۔ همدارنګه، که د کوم
  \LR{$!C \in \Delta$} دپاره
  \LR{$\Sat{M}{!C}$} وي، ځکه
  \LR{$!C \in \Xi$}، نو د کوم \LR{$!C \in \Xi$} دپاره
  \LR{$\Sat{M}{!C}$} دے۔ له دې امله
  \LR{$\Theta \Sequent \Xi$} معتبر دے۔
\item وروستنے استنتاج \LeftR{\lnot} دے: بيا د وروستي استنتاج
  مقدمه \LR{$\Gamma \Sequent \Delta, !A$} او پايله
  \LR{$\lnot !A, \Gamma \Sequent \Delta$} ده، يعنې اشتقاق داسې
  پاے ته رسېږي:
  \begin{prooftree}
    \AxiomC{}
    \Deduce$\Gamma \fCenter \Delta, !A$
    \RightLabel{\LeftR{\lnot}}
    \UnaryInf$\lnot !A, \Gamma \fCenter \Delta$
  \end{prooftree}
  او \LR{$\Theta = \lnot !A, \Gamma$}، په داسې حال کښې چې
  \LR{$\Xi = \Delta$}۔
\olpseganchor{OLP-0081-B015}{end}

\olpseganchor{OLP-0081-B016}{start}
  د استقرا فرضيه وايي چې \LR{$\Gamma \Sequent \Delta, !A$} معتبر دے؛
  يعنې د هر \LR{$\Struct{M}$} دپاره يا
  (الف) د کوم \LR{$!C \in \Gamma$} دپاره
  \LR{$\Sat/{M}{!C}$}، يا
  (ب) د کوم \LR{$!C \in \Delta$} دپاره
  \LR{$\Sat{M}{!C}$}، يا
  (ج) \LR{$\Sat{M}{!A}$} وي۔ غواړو وښيو چې
  \LR{$\Theta \Sequent \Xi$} هم معتبر دے۔ فرض کړئ
  \LR{$\Struct{M}$} جوړښت دے۔ که (الف) سم وي، نو داسې
  \LR{$!C \in \Gamma$} شته چې
  \LR{$\Sat/{M}{!C}$}، خو
  \LR{$!C \in \Theta$} هم دے۔ که (ب) سم وي، نو داسې
  \LR{$!C \in \Delta$} شته چې
  \LR{$\Sat{M}{!C}$}، خو
  \LR{$!C \in \Xi$} هم دے۔ په پاے کښې، که
  \LR{$\Sat{M}{!A}$} وي، نو
  \LR{$\Sat/{M}{\lnot !A}$} دے۔ ځکه
  \LR{$\lnot !A \in \Theta$}، نو داسې \LR{$!C \in \Theta$} شته چې
  \LR{$\Sat/{M}{!C}$} وي۔ له دې امله
  \LR{$\Theta \Sequent \Xi$} معتبر دے۔
\item وروستنے استنتاج \RightR{\lnot} دے: تمرين۔
\item وروستنے استنتاج \LeftR{\land} دے: دوه بڼې شته؛
  \LR{$!A \land !B$} په کيڼ لوري کښې د مقدمې له \LR{$!A$} يا \LR{$!B$} څخه
  استنتاج کېدے شي۔ په لومړي حالت کښې \LR{$\pi$} داسې پاے ته رسېږي:
  \begin{prooftree}
    \AxiomC{}
    \Deduce$!A, \Gamma \fCenter \Delta$
    \RightLabel{\LeftR{\land}}
    \UnaryInf$!A \land !B, \Gamma \fCenter \Delta$
  \end{prooftree}
  او \LR{$\Theta = !A \land !B, \Gamma$}، په داسې حال کښې چې
  \LR{$\Xi = \Delta$}۔ يو
  جوړښت~\LR{$\Struct
    M$} په پام کښې ونيسئ۔ ځکه د
  استقرا د فرضيې له مخې \LR{$!A, \Gamma \Sequent \Delta$} معتبر دے،
  يا (الف) \LR{$\Sat/{M}{!A}$}،
  يا (ب) د کوم \LR{$!C \in \Gamma$} دپاره
  \LR{$\Sat/{M}{!C}$}، يا
  (ج) د کوم \LR{$!C \in \Delta$} دپاره
  \LR{$\Sat{M}{!C}$} وي۔
  په (الف) حالت کښې
  \LR{$\Sat/{M}{!A \land !B}$}، نو داسې
  \LR{$!C \in \Theta$} (يعنې \LR{$!A \land !B$}) شته چې
  \LR{$\Sat/{M}{!C}$} وي۔ په (ب) حالت کښې داسې
  \LR{$!C \in \Gamma$} شته چې
  \LR{$\Sat/{M}{!C}$}، او
  \LR{$!C \in \Theta$} هم دے۔ په (ج) حالت کښې داسې
  \LR{$!C \in \Delta$} شته چې
  \LR{$\Sat{M}{!C}$}، او ځکه
  \LR{$\Xi = \Delta$}، \LR{$!C \in \Xi$} هم دے۔ نو په هر حالت کښې
  \LR{$\Struct M$} د
  \LR{$!A \land !B, \Gamma \Sequent \Delta$} سېکوېنټ پوره کوي۔
  ځکه \LR{$\Struct{M}$} اختياري و،
  \LR{$!A \land !B, \Gamma \Sequent \Delta$} معتبر دے۔ هغه حالت چې
  \LR{$!A \land !B$} له \LR{$!B$} څخه استنتاج کېږي په همدې ډول، د \LR{$!A$} پر
  ځاے د \LR{$!B$} په راوړلو، حل کېږي۔
\item وروستنے استنتاج \RightR{\lor} دے: دوه بڼې شته؛
  \LR{$!A \lor !B$} په ښي لوري کښې د مقدمې له \LR{$!A$} يا \LR{$!B$} څخه
  استنتاج کېدے شي۔ په لومړي حالت کښې \LR{$\pi$} داسې پاے ته رسېږي:
  \begin{prooftree}
    \AxiomC{}
    \Deduce$\Gamma \fCenter \Delta, !A$
    \RightLabel{\RightR{\lor}}
    \UnaryInf$\Gamma \fCenter \Delta, !A \lor !B$
  \end{prooftree}
  اوس \LR{$\Theta = \Gamma$} او \LR{$\Xi = \Delta, !A \lor !B$} دي۔ يو
  جوړښت~\LR{$\Struct{M}$} په پام کښې ونيسئ۔ ځکه
  \LR{$\Gamma \Sequent \Delta, !A$} معتبر دے، يا
  (الف) \LR{$\Sat{M}{!A}$}،
  يا (ب) د کوم \LR{$!C \in \Gamma$} دپاره
  \LR{$\Sat/{M}{!C}$}، يا
  (ج) د کوم \LR{$!C \in \Delta$} دپاره
  \LR{$\Sat{M}{!C}$} وي۔
  په (الف) حالت کښې
  \LR{$\Sat{M}{!A \lor !B}$}۔ په (ب) حالت کښې
  داسې \LR{$!C \in \Gamma$} شته چې
  \LR{$\Sat/{M}{!C}$}۔ په (ج) حالت کښې داسې
  \LR{$!C \in \Delta$} شته چې
  \LR{$\Sat{M}{!C}$}۔ نو په هر حالت کښې
  \LR{$\Struct{M}$} د
  \LR{$\Gamma \Sequent \Delta, !A \lor !B$}، يعنې
  \LR{$\Theta \Sequent \Xi$}، پوره کوي۔ ځکه
  \LR{$\Struct{M}$} اختياري و،
  \LR{$\Theta \Sequent \Xi$} معتبر دے۔ هغه حالت چې \LR{$!A \lor !B$} له
  \LR{$!B$} څخه استنتاج کېږي په همدې ډول، د \LR{$!A$} پر ځاے د \LR{$!B$} په
  راوړلو، حل کېږي۔
\item وروستنے استنتاج \RightR{\lif} دے: بيا \LR{$\pi$} داسې پاے ته رسېږي:
  \begin{prooftree}
    \AxiomC{}
    \Deduce$!A, \Gamma \fCenter \Delta, !B$
    \RightLabel{\RightR{\lif}}
    \UnaryInf$\Gamma \fCenter \Delta, !A \lif !B$
  \end{prooftree}
  بيا هم د استقرا فرضيه وايي چې مقدمه معتبره ده؛ غواړو وښيو چې
  پايله هم معتبره ده۔ يو اختياري
  \LR{$\Struct{M}$} ونيسئ۔ ځکه
  \LR{$!A, \Gamma \Sequent \Delta, !B$} معتبر دے، لږ تر لږه يو له دغو
  حالتونو څخه سم دے: (الف)
  \LR{$\Sat/{M}{!A}$}، (ب)
  \LR{$\Sat{M}{!B}$}، (ج) د کوم
  \LR{$~!C \in \Gamma$} دپاره
  \LR{$\Sat/{M}{!C}$}، يا (د) د کوم
  \LR{$!C \in \Delta$} دپاره
  \LR{$\Sat{M}{!C}$}۔
  په (الف) او (ب) حالتونو کښې
  \LR{$\Sat{M}{!A \lif !B}$}، نو داسې
  \LR{$!C \in \Delta, !A \lif !B$} شته چې
  \LR{$\Sat{M}{!C}$} وي۔ په (ج) حالت کښې د کوم
  \LR{$!C \in \Gamma$} دپاره
  \LR{$\Sat/{M}{!C}$} دے۔ په (د) حالت کښې د
  کوم \LR{$!C \in \Delta$} دپاره
  \LR{$\Sat{M}{!C}$} دے۔ په هر حالت کښې
  \LR{$\Struct{M}$} د
  \LR{$\Gamma \Sequent \Delta, !A \lif !B$} سېکوېنټ پوره کوي۔ ځکه
  \LR{$\Struct{M}$} اختياري و،
  \LR{$\Gamma \Sequent \Delta, !A \lif !B$} معتبر دے۔
  %
\item وروستنے استنتاج \LeftR{\lforall} دے: بيا داسې
  فارمول~\LR{$!A(x)$} او تړلې اصطلاح~\LR{$t$} شته چې \LR{$\pi$} داسې
  پاے ته رسېږي:
  \begin{prooftree}
    \AxiomC{}
    \Deduce$!A(t), \Gamma \fCenter \Delta$
    \RightLabel{\LeftR{\lforall}}
    \UnaryInf$\lforall[x][!A(x)], \Gamma \fCenter \Delta$
  \end{prooftree}
  غواړو وښيو چې پايله
  \LR{$\lforall[x][!A(x)], \Gamma \Sequent \Delta$} معتبره ده۔ يو
  جوړښت~\LR{$\Struct M$} په پام کښې ونيسئ۔ ځکه چې مقدمه
  \LR{$!A(t), \Gamma \Sequent \Delta$} معتبره ده، يا
  (الف) \LR{$\Sat/{M}{!A(t)}$}، (ب) د کوم \LR{$!C \in \Gamma$} دپاره
  \LR{$\Sat/{M}{!C}$}، يا (ج) د کوم \LR{$!C \in \Delta$} دپاره
  \LR{$\Sat{M}{!C}$} دے۔ په (الف) حالت کښې، د
  \olref[syn][sem]{prop:quant-terms} له مخې، که
  \LR{$\Sat{M}{\lforall[x][!A(x)]}$} وي، نو \LR{$\Sat{M}{!A(t)}$} دے۔ ځکه
  \LR{$\Sat/{M}{!A(t)}$}، نو
  \LR{$\Sat/{M}{\lforall[x][!A(x)]}$}۔ په (ب) او (ج) حالتونو کښې
  \LR{$\Struct{M}$} هم \LR{$\lforall[x][!A(x)], \Gamma \Sequent \Delta$}
  پوره کوي۔ ځکه \LR{$\Struct M$} اختياري و،
  \LR{$\lforall[x][!A(x)], \Gamma \Sequent \Delta$} معتبر دے۔
\item وروستنے استنتاج \RightR{\lexists} دے: تمرين۔
\item وروستنے استنتاج \RightR{\lforall} دے: بيا داسې
  فارمول~\LR{$!A(x)$} او ثابت سمبول~\LR{$a$} شته چې \LR{$\pi$} داسې
  پاے ته رسېږي:
  \begin{prooftree}
    \AxiomC{}
    \Deduce$\Gamma \fCenter \Delta, !A(a)$
    \RightLabel{\RightR{\lforall}}
    \UnaryInf$\Gamma \fCenter \Delta, \lforall[x][!A(x)]$
  \end{prooftree}
  دلته د ځانګړي متغير شرط پوره کېږي؛ يعنې \LR{$a$} په \LR{$!A(x)$}،
  \LR{$\Gamma$} يا \LR{$\Delta$} کښې نۀ راځي۔ د استقرا د فرضيې له مخې د
  وروستي استنتاج مقدمه معتبره ده۔ بايد وښيو چې پايله هم معتبره ده؛
  يعنې د هر جوړښت~\LR{$\Struct M$} دپاره يا
  (الف) \LR{$\Sat{M}{\lforall[x][!A(x)]}$}، يا
  (ب) د کوم \LR{$!C \in \Gamma$} دپاره \LR{$\Sat/{M}{!C}$}، يا
  (ج) د کوم \LR{$!C \in \Delta$} دپاره \LR{$\Sat{M}{!C}$} وي۔
\olpseganchor{OLP-0081-B016}{end}

\olpseganchor{OLP-0081-B017}{start}
  فرض کړئ \LR{$\Struct{M}$} يو اختياري جوړښت دے۔ که (ب) يا (ج)
  سم وي، کار بشپړ دے؛ نو فرض کړئ يو هم سم نۀ دے: د هر
  \LR{$!C \in \Gamma$} دپاره \LR{$\Sat{M}{!C}$}، او د هر
  \LR{$!C \in \Delta$} دپاره \LR{$\Sat/{M}{!C}$}۔ بايد وښيو چې (الف) سم دے،
  يعنې \LR{$\Sat{M}{\lforall[x][!A(x)]}$}۔ د
  \olref[syn][ass]{prop:sat-quant} له مخې کافي ده وښيو چې د ټولو
  متغير-ټاکنو~\LR{$s$} دپاره \LR{$\Sat{M}{!A(x)}[s]$}۔ نو يو اختياري
  متغير-ټاکنه~\LR{$s$} ونيسئ۔ داسې جوړښت~\LR{$\Struct{M'}$} په پام کښې
  ونيسئ چې د~\LR{$\Struct{M}$} په شان دے، پرته له دې چې
  \LR{$\Assign{a}{M'} = s(x)$}۔ د
  \olref[syn][ext]{cor:extensionality-sent} له مخې، د هر
  \LR{$!C \in \Gamma$} دپاره \LR{$\Sat{M'}{!C}$}؛ ځکه \LR{$a$} په~\LR{$\Gamma$} کښې
  نۀ راځي، او د هر \LR{$!C \in \Delta$} دپاره \LR{$\Sat/{M'}{!C}$}۔ خو
  مقدمه معتبره ده، نو \LR{$\Sat{M'}{!A(a)}$}۔ د
  \olref[syn][ass]{prop:sentence-sat-true} له مخې
  \LR{$\Sat{M'}{!A(a)}[s]$}، ځکه \LR{$!A(a)$} يوه جمله ده۔ اوس
  \LR{$\varAssign{s}{s}{x}$} او
  \LR{$s(x) = \Value{a}{M'}[s]$}، ځکه \LR{$\Struct{M'}$} مو همداسې تعريف
  کړے دے۔ نو \olref[syn][ext]{prop:ext-formulas} پلي کېږي او
  \LR{$\Sat{M'}{!A(x)}[s]$} ترلاسه کوو۔ ځکه \LR{$a$} په~\LR{$!A(x)$} کښې نۀ راځي،
  د \olref[syn][ext]{prop:extensionality} له مخې
  \LR{$\Sat{M}{!A(x)}[s]$}۔ ځکه \LR{$s$} اختياري و، دا مو ثابت کړل چې د ټولو
  متغير-ټاکنو دپاره \LR{$\Sat{M}{!A(x)}[s]$}۔
\item وروستنے استنتاج \LeftR{\lexists} دے: تمرين۔

\end{enumerate}
اوس هغه ممکن استنتاجونه څېړو چې دوه مقدمې لري۔
\begin{enumerate}
\item وروستنے استنتاج پرېکون دے: بيا \LR{$\pi$} داسې پاے ته رسېږي:
  \begin{prooftree}
    \AxiomC{}
    \Deduce$\Gamma \fCenter \Delta, !A$
    \AxiomC{}
    \Deduce$!A, \Pi \fCenter \Lambda$
    \RightLabel{\Cut}
    \BinaryInf$\Gamma, \Pi \fCenter \Delta, \Lambda$
  \end{prooftree}
  فرض کړئ \LR{$\Struct{M}$} جوړښت دے۔ د استقرا د فرضيې له مخې دواړه
  مقدمې معتبرې دي، نو
  \LR{$\Struct{M}$} دواړه پوره کوي۔
  دوه حالتونه بېلوو: (الف)
  \LR{$\Sat/{M}{!A}$} او (ب)
  \LR{$\Sat{M}{!A}$}۔ په (الف) حالت کښې، د دې
  دپاره چې \LR{$\Struct{M}$} کيڼه مقدمه
  پوره کړي، بايد \LR{$\Gamma \Sequent \Delta$} پوره کړي۔ بيا پايله هم
  پوره کوي۔ په (ب) حالت کښې، د دې دپاره چې
  \LR{$\Struct{M}$} ښۍ مقدمه پوره کړي،
  بايد \LR{$\Pi \Sequent \Lambda$} پوره کړي۔ بيا هم
  \LR{$\Struct{M}$} پايله پوره کوي۔
\item وروستنے استنتاج \RightR{\land} دے۔ بيا \LR{$\pi$} داسې پاے ته رسېږي:
  \begin{prooftree}
    \AxiomC{}
    \Deduce$\Gamma \fCenter \Delta, !A$
    \AxiomC{}
    \Deduce$\Gamma \fCenter \Delta, !B$
    \RightLabel{\RightR{\land}}
    \BinaryInf$\Gamma \fCenter \Delta, !A \land !B$
  \end{prooftree}
  يو جوړښت~\LR{$\Struct
    M$} په پام کښې ونيسئ۔ که
  \LR{$\Struct{M}$} د
  \LR{$\Gamma \Sequent \Delta$} سېکوېنټ پوره کوي، کار بشپړ دے۔ نو فرض
  کړئ نۀ ئې پوره کوي۔ ځکه د استقرا د فرضيې له مخې
  \LR{$\Gamma \fCenter \Delta, !A$} معتبر دے،
  \LR{$\Sat{M}{!A}$}۔ همدارنګه، ځکه
  \LR{$\Gamma \Sequent \Delta, !B$} معتبر دے،
  \LR{$\Sat{M}{!B}$}۔ نو
  \LR{$\Sat{M}{!A \land !B}$}۔
\item وروستنے استنتاج \LeftR{\lor} دے: تمرين۔
\item وروستنے استنتاج \LeftR{\lif} دے۔ بيا \LR{$\pi$} داسې پاے ته رسېږي:
  \begin{prooftree}
    \AxiomC{}
    \Deduce$\Gamma \fCenter \Delta, !A$
    \AxiomC{}
    \Deduce$!B, \Pi \fCenter \Lambda$
    \RightLabel{\LeftR{\lif}}
    \BinaryInf$!A \lif !B, \Gamma, \Pi \fCenter \Delta, \Lambda$
  \end{prooftree}
  بيا هم يو
  جوړښت~\LR{$\Struct{M}$} په پام کښې ونيسئ او فرض کړئ چې
  \LR{$\Struct{M}$} د
  \LR{$\Gamma, \Pi \Sequent \Delta, \Lambda$} سېکوېنټ نۀ پوره کوي۔
  بايد وښيو چې
  \LR{$\Sat/{M}{!A \lif !B}$}۔ که
  \LR{$\Struct{M}$} د
  \LR{$\Gamma, \Pi \Sequent \Delta, \Lambda$} سېکوېنټ نۀ پوره کوي،
  نۀ \LR{$\Gamma \Sequent \Delta$} او نۀ \LR{$\Pi \Sequent \Lambda$} پوره
  کوي۔ ځکه \LR{$\Gamma \Sequent \Delta, !A$} معتبر دے،
  \LR{$\Sat{M}{!A}$}؛ او ځکه
  \LR{$!B, \Pi \Sequent \Lambda$} معتبر دے،
  \LR{$\Sat/{M}{!B}$}۔ خو بيا
  \LR{$\Sat/{M}{!A \lif !B}$}، او همدا مو
  غوښتل۔
\end{enumerate}
\end{proof}
\olpseganchor{OLP-0081-B017}{end}

\olpseganchor{OLP-0081-B018}{start}
\begin{prob}
د \olref[fol][seq][sou]{thm:sequent-soundness} ثبوت بشپړ کړئ۔
\end{prob}
\olpseganchor{OLP-0081-B018}{end}



\olpseganchor{OLP-0081-B020}{start}
\begin{cor}
\ollabel{cor:weak-soundness}
که \LR{$\Proves !A$} وي، نو \LR{$!A$} معتبر ده۔
\end{cor}
\olpseganchor{OLP-0081-B020}{end}

\olpseganchor{OLP-0081-B021}{start}
\begin{cor}
\ollabel{cor:entailment-soundness}
که \LR{$\Gamma \Proves !A$} وي، نو \LR{$\Gamma \Entails !A$}۔
\end{cor}
\olpseganchor{OLP-0081-B021}{end}

\olpseganchor{OLP-0081-B022}{start}
\begin{proof}
که \LR{$\Gamma \Proves !A$} وي، نو د کوم متناهي فرعي سټ
\LR{$\Gamma_0 \subseteq \Gamma$} دپاره د
\LR{$\Gamma_0 \Sequent !A$} اشتقاق شته۔ د
\olref{thm:sequent-soundness} له مخې، هر
جوړښت~\LR{$\Struct{M}$} يا کوم \LR{$!B \in \Gamma_0$}
ناسموي يا \LR{$!A$} رښتيا کوي۔ نو که
\LR{$\Sat{M}{\Gamma}$} وي، بيا
\LR{$\Sat{M}{!A}$} هم دے۔
\end{proof}
\olpseganchor{OLP-0081-B022}{end}

\olpseganchor{OLP-0081-B023}{start}
\begin{cor}
\ollabel{cor:consistency-soundness}
که \LR{$\Gamma$} د صدق وړ وي، نو سازګار دے۔
\end{cor}
\olpseganchor{OLP-0081-B023}{end}

\olpseganchor{OLP-0081-B024}{start}
\begin{proof}
موږ مقابل عکس ثابتوو۔ فرض کړئ \LR{$\Gamma$} سازګار نۀ دے۔ بيا يو
متناهي \LR{$\Gamma_0 \subseteq \Gamma$} او د
\LR{$\Gamma_0 \Sequent \quad$} اشتقاق شته۔ د
\olref{thm:sequent-soundness} له مخې
\LR{$\Gamma_0 \Sequent \quad$} معتبر دے۔ په بله وينا، د هر
جوړښت~\LR{$\Struct{M}$} دپاره داسې
\LR{$!C \in \Gamma_0$} شته چې
\LR{$\Sat/{M}{!C}$}، او ځکه
\LR{$\Gamma_0 \subseteq \Gamma$}، دغه \LR{$!C$} په~\LR{$\Gamma$} کښې هم دے۔
نو هېڅ \LR{$\Struct{M}$} د~\LR{$\Gamma$} سټ نۀ
پوره کوي، او \LR{$\Gamma$} د صدق وړ نۀ دے۔
\end{proof}
\olpseganchor{OLP-0081-B024}{end}

\olpseganchor{OLP-0082-B004}{start}
\olfileid{fol}{seq}{ide}
\olpseganchor{OLP-0082-B004}{end}

\olpseganchor{OLP-0082-B005}{start}
\olsection{له عينيت سره اشتقاقونه}
\olpseganchor{OLP-0082-B005}{end}

\olpseganchor{OLP-0082-B006}{start}
له عينيت سره اشتقاقونه اضافي لومړنيو سېکوېنټونو او
استنتاجي قاعدو ته اړتيا لري۔
\olpseganchor{OLP-0082-B006}{end}

\olpseganchor{OLP-0082-B007}{start}
\begin{defn}[د~\LR{$\eq$} دپاره لومړني سېکوېنټونه]
که \LR{$t$} تړلې اصطلاح وي، نو \LR{${} \Sequent \eq[t][t]$} يو لومړنے
سېکوېنټ دے۔
\end{defn}
\olpseganchor{OLP-0082-B007}{end}

\olpseganchor{OLP-0082-B008}{start}
د \LR{$\eq$} قاعدې دا دي (\LR{$t_1$} او \LR{$t_2$} تړلې اصطلاحګانې دي):
\olpseganchor{OLP-0082-B008}{end}

\olpseganchor{OLP-0082-B009}{start}
\begin{defish}
\Axiom$ \eq[t_1][t_2], \Gamma \fCenter \Delta, !A(t_1) $
\RightLabel{\LR{$\eq$}}
\UnaryInf$\eq[t_1][t_2], \Gamma \fCenter \Delta, !A(t_2)$
\DisplayProof
\hfill
\Axiom$\eq[t_1][t_2], \Gamma \fCenter \Delta, !A(t_2) $
\RightLabel{\LR{$\eq$}}
\UnaryInf$\eq[t_1][t_2], \Gamma  \fCenter \Delta, !A(t_1)$
\DisplayProof
\end{defish}
\olpseganchor{OLP-0082-B009}{end}

\olpseganchor{OLP-0082-B010}{start}
\begin{ex}
که \LR{$s$} او \LR{$t$} تړلې اصطلاحګانې وي، نو
\LR{$\eq[s][t], !A(s) \Proves !A(t)$}:
\begin{prooftree}
\Axiom$ !A(s) \fCenter !A(s)$
\RightLabel{\LeftR{\Weakening}}
\UnaryInf$\eq[s][t], !A(s)  \fCenter !A(s)$
\RightLabel{\LR{$\eq$}}
\UnaryInf$\eq[s][t], !A(s)  \fCenter !A(t)$
\end{prooftree}
دا به د يو شان څيزونو د تعويض د اصل يا د لايبنېڅ د قانون په نامه
اشنا وي۔
\olpseganchor{OLP-0082-B010}{end}

\olpseganchor{OLP-0082-B011}{start}
\LR{$\Log{LK}$} ثابتوي چې \LR{$\eq$} متناظره او انتقالي ده:
\begin{prooftree}
\Axiom$ \fCenter \eq[t_1][t_1] $
\RightLabel{\LeftR{\Weakening}}
\UnaryInf$ \eq[t_1][t_2] \fCenter \eq[t_1][t_1] $
\RightLabel{\LR{$\eq$}}
\UnaryInf$ \eq[t_1][t_2] \fCenter \eq[t_2][t_1]$
\DisplayProof\qquad\bottomAlignProof
\Axiom$ \eq[t_1][t_2] \fCenter \eq[t_1][t_2] $
\RightLabel{\LeftR{\Weakening}}
\UnaryInf$\eq[t_2][t_3], \eq[t_1][t_2]  \fCenter \eq[t_1][t_2] $
\RightLabel{\LR{$\eq$}}
\UnaryInf$\eq[t_2][t_3], \eq[t_1][t_2]  \fCenter \eq[t_1][t_3]$
\RightLabel{\LeftR{\Exchange}}
\UnaryInf$\eq[t_1][t_2], \eq[t_2][t_3]  \fCenter \eq[t_1][t_3]$
\end{prooftree}
په کيڼ اشتقاق کښې فارمول~\LR{$\eq[x][t_1]$} زموږ
\LR{$!A(x)$} دے۔ په ښي اړخ کښې \LR{$!A(x)$} د~\LR{$\eq[t_1][x]$} په توګه اخلو۔
\end{ex}
\olpseganchor{OLP-0082-B011}{end}

\olpseganchor{OLP-0082-B012}{start}
\begin{prob}
د لاندې سېکوېنټونو اشتقاقونه ورکړئ:
\begin{enumerate}
\item \LR{$\Sequent \lforall[x][\lforall[y][((x = y \land !A(x)) \lif !A(y))]]$}
\item \LR{$\lexists[x][!A(x)] \land \lforall[y][\lforall[z][((!A(y) \land
    !A(z)) \lif y = z)]] \Sequent 
\lexists[x][(!A(x) \land \lforall[y][(!A(y) \lif y = x)])]$}
\end{enumerate}
\end{prob}
\olpseganchor{OLP-0082-B012}{end}

\olpseganchor{OLP-0083-B004}{start}
\olfileid{fol}{seq}{sid}
\olpseganchor{OLP-0083-B004}{end}

\olpseganchor{OLP-0083-B005}{start}
\olsection{له عينيت سره سموالے}
\olpseganchor{OLP-0083-B005}{end}

\olpseganchor{OLP-0083-B006}{start}
\begin{prop}
\LR{$\Log{LK}$} د عينيت له لومړنيو سېکوېنټونو او قاعدو سره سم دے۔
\end{prop}
\olpseganchor{OLP-0083-B006}{end}

\olpseganchor{OLP-0083-B007}{start}
\begin{proof}
د \LR{${} \Sequent \eq[t][t]$} په بڼه لومړني سېکوېنټونه معتبر دي؛
ځکه د هر جوړښت~\LR{$\Struct M$} دپاره
\LR{$\Sat{M}{\eq[t][t]}$}۔ (پام وکړئ، فرض کوو چې اصطلاح~\LR{$t$} تړلې ده،
يعنې هېڅ متغير نۀ لري؛ نو متغير-ټاکنې اهميت نۀ لري۔)
\olpseganchor{OLP-0083-B007}{end}

\olpseganchor{OLP-0083-B008}{start}
فرض کړئ په اشتقاق کښې وروستنے استنتاج \LR{$=$} دے۔ بيا مقدمه
\LR{$\eq[t_1][t_2], \Gamma \Sequent \Delta, !A(t_1)$} او پايله
\LR{$\eq[t_1][t_2], \Gamma \Sequent \Delta, !A(t_2)$} ده۔ يو
جوړښت~\LR{$\Struct M$} په پام کښې ونيسئ۔ بايد وښيو چې پايله
معتبره ده؛ يعنې که \LR{$\Sat{M}{\eq[t_1][t_2]}$} او
\LR{$\Sat{M}{\Gamma}$} وي، نو يا د کوم \LR{$!C \in \Delta$} دپاره
\LR{$\Sat{M}{!C}$}، يا \LR{$\Sat{M}{!A(t_2)}$} دے۔
\olpseganchor{OLP-0083-B008}{end}

\olpseganchor{OLP-0083-B009}{start}
د استقرا د فرضيې له مخې مقدمه معتبره ده۔ د دې معنٰی دا ده چې که
\LR{$\Sat{M}{\eq[t_1][t_2]}$} او \LR{$\Sat{M}{\Gamma}$} وي، نو يا
(الف) د کوم \LR{$!C \in \Delta$} دپاره \LR{$\Sat{M}{!C}$}، يا
(ب) \LR{$\Sat{M}{!A(t_1)}$} دے۔ په (الف) حالت کښې کار بشپړ دے۔ (ب)
حالت وګورئ۔ داسې متغير-ټاکنه~\LR{$s$} ونيسئ چې
\LR{$s(x) = \Value{t_1}{M}$}۔ د
\olref[syn][ass]{prop:sentence-sat-true} له مخې
\LR{$\Sat{M}{!A(t_1)}[s]$}۔ ځکه \LR{$\varAssign{s}{s}{x}$}، د
\olref[syn][ext]{prop:ext-formulas} له مخې
\LR{$\Sat{M}{!A(x)}[s]$}۔ ځکه \LR{$\Sat{M}{\eq[t_1][t_2]}$}، نو
\LR{$\Value{t_1}{M} = \Value{t_2}{M}$}، او له دې امله
\LR{$s(x) = \Value{t_2}{M}$}۔ د
\olref[syn][ext]{prop:ext-formulas} په بيا پلي کولو
\LR{$\Sat{M}{!A(t_2)}[s]$} هم لرو۔ د
\olref[syn][ass]{prop:sentence-sat-true} له مخې
\LR{$\Sat{M}{!A(t_2)}$}۔
\end{proof}
\olpseganchor{OLP-0083-B009}{end}

\olpseganchor{OLP-0084-B004}{start}
\olchapter{fol}{ntd}{فطري استنتاج}
\olpseganchor{OLP-0084-B004}{end}

\olpseganchor{OLP-0084-B005}{start}
\begin{editorial}
دا څپرکے د ګېنتڅن/پراويڅ په طريقه د فطري استنتاج يو نظام وړاندې
کوي۔
\olpseganchor{OLP-0084-B005}{end}
  
\olpseganchor{OLP-0084-B006}{start}
د ثبوت د يوه نظام په توګه له فطري استنتاج سره اړوند مواد د شاملولو
يا ايستلو دپاره د~\texttt{prfND} ټګ وکاروئ۔
\end{editorial}
\olpseganchor{OLP-0084-B006}{end}





\olpseganchor{OLP-0084-B009}{start}
%
  

\olpseganchor{OLP-0084-B009}{end}





\olpseganchor{OLP-0084-B012}{start}
%
  

\olpseganchor{OLP-0084-B012}{end}







\olpseganchor{OLP-0084-B016}{start}
%
  

\olpseganchor{OLP-0084-B016}{end}



\olpseganchor{OLP-0084-B018}{start}
%
  
  

\olpseganchor{OLP-0084-B018}{end}

\olpseganchor{OLP-0085-B004}{start}
\olfileid{fol}{ntd}{rul}
\olpseganchor{OLP-0085-B004}{end}

\olpseganchor{OLP-0085-B005}{start}
\olsection{قاعدې او اشتقاقونه}
\olpseganchor{OLP-0085-B005}{end}

\olpseganchor{OLP-0085-B006}{start}
\begin{explain}
د فطري استنتاج له نظامونو څخه موخه دا ده چې په رياضيکي ثبوت کښې له
کارېدونکي غيررسمي استدلال سره نژدې موازي وي (همدا وجه ده چې تر يوه
حده “فطري” دے)۔ د فطري استنتاج ثبوتونه له فرضيو پيلېږي۔ بيا د
استنتاج قاعدې ورباندې کارول کېږي۔ فرضيې د~\Intro{\lnot}،
\Intro{\lif}، \Elim{\lor} او \Elim{\lexists} د استنتاج د قاعدو په وسيله “ايستل” کېږي،
او د روښانتيا دپاره د ايستل شوې فرضيې نښان د استنتاج تر
څنګ ليکل کېږي۔
\end{explain}
\olpseganchor{OLP-0085-B006}{end}

\olpseganchor{OLP-0085-B007}{start}
\begin{defn}[فرضيه]
يوه \emph{فرضيه} هره هغه تړلے فارمول ده چې د هرې څانګې په تر ټولو
پورتني ځاے کښې وي۔
\end{defn}
\olpseganchor{OLP-0085-B007}{end}

\olpseganchor{OLP-0085-B008}{start}
په فطري استنتاج کښې اشتقاقونه د تړلي فارمولونه ځانګړې ونې دي؛
تر ټولو پورتنۍ تړلي فارمولونه فرضيې وي، او که يوه تړلے فارمول تر
يوې، دوو يا دريو نورو جملو لاندې وي، بايد په سمه توګه د استنتاج له
يوې قاعدې څخه پايله شي۔ د استنتاج په سر کښې تړلي فارمولونه د هغې
\emph{مقدمې} او لاندې تړلے فارمول د استنتاج \emph{پايله} بلل کېږي۔
قاعدې په جوړو کښې راځي: د هر منطقي عملګر دپاره يوه د داخلولو او
يوه د ايستلو قاعده۔ دوئ په پايله کښې منطقي عملګر داخلوي يا
منطقي عملګر د قاعدې له يوې مقدمې څخه باسي۔ ځينې قاعدې اجازه
ورکوي چې د يوه ټاکلي
ډول فرضيه \emph{ايستل} شي۔ د دې د ښودلو دپاره چې کومه
فرضيه د کوم استنتاج په وسيله ايستل کېږي، فرضيې او استنتاج
دواړو ته نښانونه ورکوو۔ دا د فرضيې په “\LR{$\Discharge{!A}{n}$}” ليکلو
ښودل کېږي۔
\olpseganchor{OLP-0085-B008}{end}

\olpseganchor{OLP-0085-B009}{start}
% Only include this sentence is on of \land, lor, lif or lnot is defined.
\olpseganchor{OLP-0085-B009}{end}

\olpseganchor{OLP-0085-B010}{start}
دا دود دے چې د ټولو منطقي عملګرونه~\LR{$\land$}، \LR{$\lor$}، \LR{$\lif$}،
    \LR{$\lnot$} او~\LR{$\lfalse$} دپاره قاعدې په پام کښې ونيول شي، که څه هم
    ځينې ئې تعريف شوي وي۔
\olpseganchor{OLP-0085-B010}{end}

\olpseganchor{OLP-0086-B004}{start}
\olfileid{fol}{ntd}{prl}
\olpseganchor{OLP-0086-B004}{end}

\olpseganchor{OLP-0086-B005}{start}
\olsection{قضيې قاعدې}
\olpseganchor{OLP-0086-B005}{end}


\olpseganchor{OLP-0086-B006}{start}
\subsection{د \texorpdfstring{\LR{$\land$}}{∧} دپاره قاعدې}
\olpseganchor{OLP-0086-B006}{end}

\olpseganchor{OLP-0086-B007}{start}
\begin{defish}
\AxiomC{\LR{$!A$}}
\AxiomC{\LR{$!B$}}
\RightLabel{\Intro{\land}}
\BinaryInfC{\LR{$!A \land !B$}}
\DisplayProof
\hfill
\begin{tabular}{r}
\AxiomC{\LR{$!A \land !B$}}
\RightLabel{\Elim{\land}}
\UnaryInfC{\LR{$!A$}}
\DisplayProof
\\[3ex]
\AxiomC{\LR{$!A \land !B$}}
\RightLabel{\Elim{\land}}
\UnaryInfC{\LR{$!B$}}
\DisplayProof
\end{tabular}
\end{defish}
\olpseganchor{OLP-0086-B007}{end}

\olpseganchor{OLP-0086-B008}{start}
\subsection{د \texorpdfstring{\LR{$\lor$}}{∨} دپاره قاعدې}
\olpseganchor{OLP-0086-B008}{end}

\olpseganchor{OLP-0086-B009}{start}
\begin{defish}
\begin{tabular}{r}
\AxiomC{\LR{$!A$}}
\RightLabel{\Intro{\lor}}
\UnaryInfC{\LR{$!A \lor !B$}}
\DisplayProof
\\[3ex]
\AxiomC{\LR{$!B$}}
\RightLabel{\Intro{\lor}}
\UnaryInfC{\LR{$!A \lor !B$}}
\DisplayProof
\end{tabular}
\hfill
\AxiomC{\LR{$!A \lor !B$}}
\AxiomC{\LR{$\Discharge{!A}{n}$}}
\DeduceC{\LR{$!C$}}
\AxiomC{\LR{$\Discharge{!B}{n}$}}
\DeduceC{\LR{$!C$}}
\DischargeRule{\Elim{\lor}}{n}
\TrinaryInfC{\LR{$!C$}}
\DisplayProof
\end{defish}
\olpseganchor{OLP-0086-B009}{end}

\olpseganchor{OLP-0086-B010}{start}
\subsection{د \texorpdfstring{\LR{$\lif$}}{→} دپاره قاعدې}
\olpseganchor{OLP-0086-B010}{end}

\olpseganchor{OLP-0086-B011}{start}
\begin{defish}
\AxiomC{\LR{$\Discharge{!A}{n}$}}
\DeduceC{\LR{$!B$}}
\DischargeRule{\Intro{\lif}}{n}
\UnaryInfC{\LR{$!A \lif !B$}}
\DisplayProof
\hfill
\AxiomC{\LR{$!A \lif !B$}}
\AxiomC{\LR{$!A$}}
\RightLabel{\Elim{\lif}}
\BinaryInfC{\LR{$!B$}}
\DisplayProof
\end{defish}
\olpseganchor{OLP-0086-B011}{end}

\olpseganchor{OLP-0086-B012}{start}
\subsection{د \texorpdfstring{\LR{$\lnot$}}{¬} دپاره قاعدې}
\olpseganchor{OLP-0086-B012}{end}

\olpseganchor{OLP-0086-B013}{start}
\begin{defish}
\AxiomC{\LR{$\Discharge{!A}{n}$}}
\noLine
\DeduceC{\LR{$\lfalse$}}
\DischargeRule{\Intro{\lnot}}{n}
\UnaryInfC{\LR{$\lnot !A$}}
\DisplayProof
\hfill
\AxiomC{\LR{$\lnot !A$}}
\AxiomC{\LR{$!A$}}
\RightLabel{\Elim{\lnot}}
\BinaryInfC{\LR{$\lfalse$}}
\DisplayProof
\end{defish}
\olpseganchor{OLP-0086-B013}{end}

\olpseganchor{OLP-0086-B014}{start}
\subsection{د \texorpdfstring{\LR{$\lfalse$}}{⊥} دپاره قاعدې}
\olpseganchor{OLP-0086-B014}{end}

\olpseganchor{OLP-0086-B015}{start}
\begin{defish}
\AxiomC{\LR{$\lfalse$}}
\RightLabel{\FalseInt}
\UnaryInfC{\LR{$!A$}}
\DisplayProof
\hfill
\AxiomC{\LR{$\Discharge{\lnot !A}{n}$}}
\DeduceC{\LR{$\lfalse$}}
\DischargeRule{\FalseCl}{n}
\UnaryInfC{\LR{$!A$}}
\DisplayProof
\end{defish}
\olpseganchor{OLP-0086-B015}{end}

\olpseganchor{OLP-0086-B016}{start}
پام وکړئ چې \LR{$\Intro{\lnot}$} او~\LR{$\FalseCl$} ډېر سره ورته دي: توپير
دا دے چې \LR{$\Intro{\lnot}$} يوه نفي شوې تړلے فارمول~\LR{$\lnot !A$} راوباسي،
خو \LR{$\FalseCl$} يوه مثبته تړلے فارمول~\LR{$!A$} راوباسي۔
\olpseganchor{OLP-0086-B016}{end}

\olpseganchor{OLP-0086-B017}{start}
کله چې يوه قاعده ښيي چې کومه فرضيه ايستل کېدے شي، موږ دا اجازه
بولو، شرط ئې نۀ بولو۔ مثلاً د~\LR{$\Intro{\lif}$} په قاعده کښې د
مقدمې~\LR{$!B$} په اشتقاق کښې د~\LR{$!A$} په بڼه هر شمېر فرضيې
ايستلے شو، چې صفر هم پکې شامل دے۔
\olpseganchor{OLP-0086-B017}{end}

\olpseganchor{OLP-0087-B004}{start}
\olfileid{fol}{ntd}{qrl}
\olpseganchor{OLP-0087-B004}{end}

\olpseganchor{OLP-0087-B005}{start}
\olsection{د سورونو قاعدې}
\olpseganchor{OLP-0087-B005}{end}

\olpseganchor{OLP-0087-B006}{start}
\subsection{د \texorpdfstring{\LR{$\lforall$}}{∀} دپاره قاعدې}
\olpseganchor{OLP-0087-B006}{end}

\olpseganchor{OLP-0087-B007}{start}
\begin{defish}
\AxiomC{\LR{$!A(a)$}}
\RightLabel{\Intro{\lforall}}
\UnaryInfC{\LR{$\lforall[x][\Atom{!A}{x}]$}}
\DisplayProof
\hfill
\AxiomC{\LR{$\lforall[x][\Atom{!A}{x}]$}}
\RightLabel{\Elim{\lforall}}
\UnaryInfC{\LR{$!A(t)$}}
\DisplayProof
\end{defish}
\olpseganchor{OLP-0087-B007}{end}

\olpseganchor{OLP-0087-B008}{start}
د~\LR{$\lforall$} په قاعدو کښې \LR{$t$} يوه تړلې اصطلاح ده (يعنې داسې اصطلاح
چې هېڅ متغير نۀ لري)، او \LR{$a$}~داسې ثابت سمبول دے چې په
پايله~\LR{$\lforall[x][!A(x)]$} يا په هرې هغې فرضيې کښې نۀ راځي چې د
مقدمې~\LR{$!A(a)$} په پای ته رسېدونکي اشتقاق کښې
نۀ ايستل شوې پاتې وي۔ موږ \LR{$a$} د~\Intro{\lforall} استنتاج
\emph{ځانګړے متغير} بولو۔\footnote{موږ د “ځانګړي متغير” اصطلاح
کاروو، که څه هم په پورتنۍ قاعده کښې \LR{$a$} يو ثابت دے۔ دا نوم تاريخي
لاملونه لري۔}
\olpseganchor{OLP-0087-B008}{end}

\olpseganchor{OLP-0087-B009}{start}
\subsection{د \texorpdfstring{\LR{$\lexists$}}{∃} دپاره قاعدې}
\olpseganchor{OLP-0087-B009}{end}

\olpseganchor{OLP-0087-B010}{start}
\begin{defish}
\AxiomC{\LR{$\Atom{!A}{t}$}}
\RightLabel{\Intro{\lexists}}
\UnaryInfC{\LR{$\lexists[x][\Atom{!A}{x}]$}}
\DisplayProof
\hfill
\AxiomC{\LR{$\lexists[x][\Atom{!A}{x}]$}}
\AxiomC{[\LR{$\Atom{!A}{a}$}]\LR{$^n$}}
\DeduceC{\LR{$!C$}}
\DischargeRule{\Elim{\lexists}}{n}
\BinaryInfC{\LR{$!C$}}
\DisplayProof
\end{defish}
\olpseganchor{OLP-0087-B010}{end}

\olpseganchor{OLP-0087-B011}{start}
بيا هم، \LR{$t$} يوه تړلې اصطلاح ده، او \LR{$a$} داسې ثابت سمبول دے چې
په مقدمه~\LR{$\lexists[x][!A(x)]$}، پايله~\LR{$!C$}، يا په هرې هغې فرضيې
کښې نۀ راځي چې د دواړو مقدمو په پای ته رسېدونکو اشتقاقونه کښې
نۀ ايستل شوې پاتې وي (د~\LR{$!A(a)$} فرضيو پرته)۔ موږ \LR{$a$} د
\Elim{\lexists} استنتاج \emph{ځانګړے متغير} بولو۔
\olpseganchor{OLP-0087-B011}{end}

\olpseganchor{OLP-0087-B012}{start}
هغه شرط چې په~\Intro{\lforall} يا~\Elim{\lexists} استنتاج کښې
ځانګړے متغير په پايله يا په هرې هغې فرضيې کښې نۀ راځي چې مقدمو ته
په رسېدونکو اشتقاقونه کښې نۀ ايستل شوې وي، او د وجود د ايستلو په قاعده کښې
په وجودي مقدمه کښې هم نۀ راځي، \emph{د ځانګړي
متغير شرط} بلل کېږي۔
\olpseganchor{OLP-0087-B012}{end}

\olpseganchor{OLP-0087-B013}{start}
\begin{explain}
هغه دود راياد کړئ چې کله \LR{$!A$} داسې فارمول وي چې
متغير~\LR{$x$} پکې ازاد وي، نو دا په~\LR{$!A(x)$} ليکلو ښيو۔ په همدې
چوکاټ کښې بيا \LR{$!A(t)$} د~\LR{$\Subst{!A}{t}{x}$} لنډون دے۔ نو د
\LR{$\Intro\lexists$} قاعده داسې هم ليکلے شو:
\begin{prooftree}
\AxiomC{\LR{$\Subst{!A}{t}{x}$}}
\RightLabel{\Intro{\lexists}}
\UnaryInfC{\LR{$\lexists[x][!A]$}}
\end{prooftree}
پام وکړئ چې \LR{$t$} کېدے شي له مخکښې په~\LR{$!A$} کښې راغلے وي؛ مثلاً
\LR{$!A$}~کېدے شي~\LR{$\Atom{\Obj P}{t,x}$} وي۔ له دې امله، له
~\LR{$\Atom{\Obj P}{t,t}$} څخه \LR{$\lexists[x][\Atom{\Obj P}{t,x}]$}
استنتاجول د~\LR{$\Intro\lexists$} سمه کارونه ده---تاسو د مقدمې د~\LR{$t$}
يوه يا زياتې پېښې په تړلي متغير~\LR{$x$} “بدلولے” شئ، او اړينه
نۀ ده چې ټولې پېښې ئې بدلې کړئ۔ خو د~\LR{$\Intro\lforall$} او
~\LR{$\Elim\lexists$} د ځانګړي متغير شرطونه غواړي چې ثابت سمبول~\LR{$a$}
په~\LR{$!A$} کښې نۀ راځي۔ نو د~\LR{$\Intro\lforall$} په کارولو سره له
\LR{$\Atom{\Obj P}{a,a}$} څخه \LR{$\lforall[x][\Atom{\Obj P}{a,x}]$} په سمه
توګه نۀ شئ استنتاجولے۔
\end{explain}
\olpseganchor{OLP-0087-B013}{end}

\olpseganchor{OLP-0087-B014}{start}
\begin{explain}
په~\Intro{\lexists} او~\Elim{\lforall} کښې له تړلتيا پرته بل
بنديز نشته، او اصطلاح~\LR{$t$} هره تړلې اصطلاح کېدے شي؛ نو د بل شرط
اندېښنه نۀ شته۔ بل لوري ته، د~\Elim{\lexists} او
\Intro{\lforall} په قاعدو کښې د ځانګړي متغير شرط غواړي چې
ثابت سمبول~\LR{$a$} په پايله يا کومې نۀ ايستل شوې فرضيې کښې هېڅ
ځاے نۀ راځي۔ دا شرط د نظام د سموالي د يقيني کولو دپاره اړين دے؛
يعنې نظام له هغو نۀ ايستل شوې فرضيو څخه يوازې هغه تړلي فارمولونه
اشتقاق کړي چې ترې معنايي پايله کېږي۔ له دې شرط پرته لاندې به
روا و:
\begin{prooftree}
  \AxiomC{\LR{$\lexists[x][!A(x)]$}}
  \AxiomC{\LR{$\Discharge{!A(a)}{1}$}}
  \RightLabel{*\Intro{\lforall}}
  \UnaryInfC{\LR{$\lforall[x][!A(x)]$}}
  \RightLabel{\Elim{\lexists}}
  \BinaryInfC{\LR{$\lforall[x][!A(x)]$}}
\end{prooftree}
خو \LR{$\lexists[x][!A(x)] \Entails/ \lforall[x][!A(x)]$}۔
\olpseganchor{OLP-0087-B014}{end}

\olpseganchor{OLP-0087-B015}{start}
ځکه چې د سورونو د ايستلو قاعدې د متغيرونه پر ځاے يوازې د تړلو
اصطلاحاتو ايښوولو اجازه ورکوي، نو هره هغه فارمول چې د
تړلي فارمولونه له يوه سټ څخه وايستل شي، پخپله تړلے فارمول ده۔
\end{explain}
\olpseganchor{OLP-0087-B015}{end}

\olpseganchor{OLP-0088-B004}{start}
\olfileid{fol}{ntd}{der}
\olpseganchor{OLP-0088-B004}{end}

\olpseganchor{OLP-0088-B005}{start}
\olsection{اشتقاقونه}
\olpseganchor{OLP-0088-B005}{end}

\olpseganchor{OLP-0088-B006}{start}
\begin{explain}
موږ وويل چې فرضيه څه ده او د استنتاج قاعدې مو ورکړې۔ په فطري
استنتاج کښې اشتقاقونه له دغو څخه په استقرايي ډول جوړېږي: هر
اشتقاق يا يوازې يوه فرضيه وي، يا له يو، دوو يا دريو
اشتقاقونه او ورپسې له يوه سم استنتاج څخه جوړ وي۔
\end{explain}
\olpseganchor{OLP-0088-B006}{end}

\olpseganchor{OLP-0088-B007}{start}
\begin{defn}[اشتقاق]
 له فرضيو~\LR{$\Gamma$} څخه د تړلے فارمول~\LR{$!A$} يو \emph{اشتقاق}
د تړلي فارمولونه يوه متناهي ونه ده چې لاندې شرطونه پوره کوي:
\begin{enumerate}
\item د ونې تر ټولو پورتنۍ تړلي فارمولونه يا په~\LR{$\Gamma$} کښې وي يا
  په ونه کښې د يوه استنتاج په وسيله ايستل شوې وي۔
\item د ونې تر ټولو لاندې تړلے فارمول~\LR{$!A$} ده۔
\item په ونه کښې هره تړلے فارمول، پرته له لاندېنۍ جملې~\LR{$!A$}
  څخه، د استنتاج د يوې قاعدې د سمې کارونې مقدمه وي چې پايله ئې په
  ونه کښې نېغ د هغې تړلے فارمول لاندې ولاړه وي۔
\end{enumerate}
بيا وايو چې \LR{$!A$} د اشتقاق \emph{پايله} او~\LR{$\Gamma$} ئې
نۀ ايستل شوې فرضيې دي۔
\olpseganchor{OLP-0088-B007}{end}

\olpseganchor{OLP-0088-B008}{start}
که له~\LR{$\Gamma$} څخه د~\LR{$!A$} کوم اشتقاق وي، وايو چې \LR{$!A$}
له~\LR{$\Gamma$} څخه \emph{د اشتقاق وړ} دے، يا په نښو کښې:
\LR{$\Gamma \Proves !A$}۔ که د~\LR{$!A$} داسې اشتقاق وي چې هره
فرضيه پکې ايستل شوې وي، نو~\LR{$\Proves !A$} ليکو۔
\end{defn}
\olpseganchor{OLP-0088-B008}{end}

\olpseganchor{OLP-0088-B009}{start}
\begin{ex}
هره فرضيه په يوازې ځان اشتقاق ده۔ نو مثلاً \LR{$!A$} په يوازې
ځان اشتقاق دے او \LR{$!B$} هم په يوازې ځان همداسې دے۔
له دغو څخه د، مثلاً، \LR{$\Intro{\land}$} قاعدې په کارولو يو نوے
اشتقاق اخستلے شو:
\begin{prooftree}
\AxiomC{\LR{$!A$}}
\AxiomC{\LR{$!B$}}
\RightLabel{\Intro{\land}}
\BinaryInfC{\LR{$!A \land !B$}}
\end{prooftree}
دا قاعدې عامې دي: پکې \LR{$!A$} او~\LR{$!B$} په هر ډول تړلي فارمولونه، مثلاً
په~\LR{$!C$} او~\LR{$!D$}، بدلولے شو۔ بيا به پايله \LR{$!C \land !D$} وي؛ نو
\begin{prooftree}
\AxiomC{\LR{$!C$}}
\AxiomC{\LR{$!D$}}
\RightLabel{\Intro{\land}}
\BinaryInfC{\LR{$!C \land !D$}}
\end{prooftree}
يو سم اشتقاق دے۔ البته فرضيې سره اړولے هم شو، څو \LR{$!D$} د~\LR{$!A$}
او \LR{$!C$} د~\LR{$!B$} ونډه واخلي۔ له دې امله
\begin{prooftree}
\AxiomC{\LR{$!D$}}
\AxiomC{\LR{$!C$}}
\RightLabel{\Intro{\land}}
\BinaryInfC{\LR{$!D \land !C$}}
\end{prooftree}
هم يو سم اشتقاق دے۔
\olpseganchor{OLP-0088-B009}{end}

\olpseganchor{OLP-0088-B010}{start}
اوس بله قاعده، مثلاً \LR{$\Intro{\lif}$}، کارولے شو۔ دا اجازه راکوي چې
يو شرطي قضيه پايله کړو او هره هغه فرضيه ايستل کړو چې له
دغه شرطي قضيې له مقدم سره يو شان وي۔ نو لاندې دواړه سم
اشتقاقونه دي:
\begin{prooftree}
\AxiomC{\LR{$\Discharge{!C}{1}$}}
\AxiomC{\LR{$!D$}}
\RightLabel{\Intro{\land}}
\BinaryInfC{\LR{$!C \land !D$}}
\DischargeRule{\Intro{\lif}}{1}
\UnaryInfC{\LR{$!C \lif (!C \land !D)$}}
\DisplayProof\bottomAlignProof
\AxiomC{\LR{$!C$}}
\AxiomC{\LR{$\Discharge{!D}{1}$}}
\RightLabel{\Intro{\land}}
\BinaryInfC{\LR{$!C \land !D$}}
\DischargeRule{\Intro{\lif}}{1}
\UnaryInfC{\LR{$!D \lif (!C \land !D)$}}
\end{prooftree}
دوئ په ترتيب سره ښيي چې \LR{$!D \Proves !C \lif (!C \land !D)$} او
\LR{$!C \Proves !D \lif (!C \land !D)$}۔
\olpseganchor{OLP-0088-B010}{end}

\olpseganchor{OLP-0088-B011}{start}
راياد کړئ چې د فرضيو ايستل اجازه ده، شرط نۀ دے: لازم نۀ ده چې
فرضیې وباسو۔ په ځانګړي ډول، يوه قاعده هله هم کارولے شو چې فرضيې په
اشتقاق کښې موجودې نۀ وي۔ مثلاً لاندې اشتقاق روا دے، که څه
هم د ايستل کېدو دپاره هېڅ فرضيه~\LR{$!A$} نۀ شته:
\begin{prooftree}
  \AxiomC{\LR{$!B$}}
  \DischargeRule{\Intro{\lif}}{1}
  \UnaryInfC{\LR{$!A \lif !B$}}
\end{prooftree}
\end{ex}
\olpseganchor{OLP-0088-B011}{end}

\olpseganchor{OLP-0089-B004}{start}
\olfileid{fol}{ntd}{pro}
\olpseganchor{OLP-0089-B004}{end}

\olpseganchor{OLP-0089-B005}{start}
\olsection{د اشتقاقونه بېلګې}
\olpseganchor{OLP-0089-B005}{end}

\olpseganchor{OLP-0089-B006}{start}
\begin{ex}
راځئ چې د تړلے فارمول~\LR{$(!A \land !B) \lif !A$} يو اشتقاق
ورکړو۔
\olpseganchor{OLP-0089-B006}{end}

\olpseganchor{OLP-0089-B007}{start}
په پيل کښې مطلوبه پايله د اشتقاق په لاندې برخه کښې ليکو۔
\begin{prooftree}
\AxiomC{}
\UnaryInfC{\LR{$(!A\land !B) \lif !A$}}
\end{prooftree}
\olpseganchor{OLP-0089-B007}{end}

\olpseganchor{OLP-0089-B008}{start}
اوس بايد معلومه کړو چې د کوم ډول استنتاج پايله د دې بڼې
تړلے فارمول کېدے شي۔ د پايلې اصلي عملګر~\LR{$\lif$} دے، نو هڅه
کوو چې د~\Intro{\lif} قاعدې په کارولو پايلې ته ورسېږو۔ غوره ده چې
اړوندې فرضيې وليکو او د پرمختګ پر مهال د استنتاج قاعدو ته نښانونه
ورکړو، څو د ثبوت په پای کښې په اسانۍ ووينو چې ټولې فرضيې
ايستل شوې دي که نۀ۔
\begin{prooftree}
\AxiomC{\LR{$\Discharge{!A \land !B}{1}$}}
\DeduceC{\LR{$!A$}}
\DischargeRule{\Intro{\lif}}{1} 
\UnaryInfC{\LR{$(!A\land !B) \lif !A$}}
\end{prooftree}
\olpseganchor{OLP-0089-B008}{end}

\olpseganchor{OLP-0089-B009}{start}
اوس بايد له فرضيې~\LR{$!A \land !B$} څخه تر~\LR{$!A$} پورې پړاوونه ډک کړو۔
ځکه چې يوازې له يوه رابط~\LR{$\land$} سره کار لرو، بايد د~\LR{$\land$} د
ايستلو قاعده وکاروو۔ له دې څخه لاندې ثبوت لاس ته راځي:
\begin{prooftree}
\AxiomC{\LR{$\Discharge{!A \land !B}{1}$}}
\RightLabel{\Elim{\land}}
\UnaryInfC{\LR{$!A$}}
\DischargeRule{\Intro{\lif}}{1} 
\UnaryInfC{\LR{$(!A\land !B) \lif !A$}}
\end{prooftree}
اوس د~\LR{$(!A \land !B) \lif !A$} يو سم اشتقاق لرو۔
\end{ex}
\olpseganchor{OLP-0089-B009}{end}

\olpseganchor{OLP-0089-B010}{start}
\begin{ex}
اوس راځئ چې د~\LR{$(\lnot !A \lor !B) \lif (!A \lif !B)$} يو
اشتقاق ورکړو۔
\olpseganchor{OLP-0089-B010}{end}

\olpseganchor{OLP-0089-B011}{start}
په پيل کښې مطلوبه پايله د اشتقاق په لاندې برخه کښې ليکو۔
\begin{prooftree}
\AxiomC{}
\UnaryInfC{\LR{$(\lnot !A \lor !B) \lif (!A \lif !B)$}}
\end{prooftree}
د هغې منطقي قاعدې د موندلو دپاره چې دا پايله راکولے شي، په پايله
کښې منطقي رابطونو~\LR{$\lnot$}، \LR{$\lor$} او~\LR{$\lif$} ته ګورو۔ اوس يوازې د
\LR{$\lif$} لومړۍ پېښه راته مهمه ده، ځکه همدا د وروستۍ پايلې د
تړلے فارمول~اصلي عملګر دے؛ \LR{$\lnot$}، \LR{$\lor$} او د~\LR{$\lif$}
دويمه پېښه د بل رابط په ساحه کښې دي، نو وروسته به ئې وڅېړو۔ له دې
امله د~\Intro{\lif} له قاعدې پيل کوو۔ سمه کارونه ئې بايد داسې وي:
\begin{prooftree}
\AxiomC{\LR{$\Discharge{\lnot !A \lor !B}{1}$}}
\DeduceC{\LR{$!A \lif !B$}}
\DischargeRule{\Intro{\lif}}{1}
\UnaryInfC{\LR{$(\lnot !A \lor !B) \lif (!A \lif !B)$}}
\end{prooftree}
له دې وروسته د دوام دوه لارې لرو۔ يا له لاندې څخه پورته کار روان
ساتلے او د~\Intro{\lif} بله کارونه لټولے شو، يا له پاسه لاندې کار
کولے او د~\Elim{\lor} قاعده کارولے شو۔ دويمه لار اخلو۔ فرضيه
\LR{$\lnot !A \lor !B$} د~\Elim{\lor} د تر ټولو کيڼې مقدمې په توګه
کاروو۔ د~\Elim{\lor} د سمې کارونې دپاره نورې دوه مقدمې بايد له
پايلې~\LR{$!A \lif !B$} سره يو شان وي، خو هره يوه په خپل وار له بلې
فرضیې، يعنې د~\LR{$\lnot !A \lor !B$} له دوو بېلندونو څخه له يوه،
رايستل کېدے شي۔ نو اشتقاق به داسې وي:
\begin{prooftree}
\AxiomC{\LR{$\Discharge{\lnot !A \lor !B}{1}$}}
\AxiomC{\LR{$\Discharge{\lnot !A}{2}$}}
\DeduceC{\LR{$!A \lif !B$}}
\AxiomC{\LR{$\Discharge{!B}{2}$}}
\DeduceC{\LR{$!A \lif !B$}}
\DischargeRule{\Elim{\lor}}{2}
\TrinaryInfC{\LR{$!A \lif !B$}}
\DischargeRule{\Intro{\lif}}{1} 
\UnaryInfC{\LR{$(\lnot !A \lor !B) \lif (!A \lif !B)$}}
\end{prooftree}
\olpseganchor{OLP-0089-B011}{end}

\olpseganchor{OLP-0089-B012}{start}
په ښي لوري دواړو څانګو کښې غواړو \LR{$!A \lif !B$}~اشتقاق کړو؛
غوره لار ئې د~\Intro{\lif} کارول دي۔
\begin{prooftree}
\AxiomC{\LR{$\Discharge{\lnot !A \lor !B}{1}$}}
\AxiomC{\LR{$\Discharge{\lnot !A}{2}, \Discharge{!A}{3}$}}
\DeduceC{\LR{$!B$}}
\DischargeRule{\Intro{\lif}}{3}
\UnaryInfC{\LR{$!A \lif !B$}}
\AxiomC{\LR{$\Discharge{!B}{2}, \Discharge{!A}{4}$}}
\DeduceC{\LR{$!B$}}
\DischargeRule{\Intro{\lif}}{4}
\UnaryInfC{\LR{$!A \lif !B$}}
\DischargeRule{\Elim{\lor}}{2}
\TrinaryInfC{\LR{$!A \lif !B$}}
\DischargeRule{\Intro{\lif}}{1} 
\UnaryInfC{\LR{$(\lnot !A \lor !B) \lif (!A \lif !B)$}}
\end{prooftree}
\olpseganchor{OLP-0089-B012}{end}

\olpseganchor{OLP-0089-B013}{start}
د اشتقاق په دوو تشو برخو کښې له~\LR{$\lnot !A$} او~\LR{$!A$} څخه،
او له~\LR{$!A$} او~\LR{$!B$} څخه، د~\LR{$!B$} اشتقاقونه غواړو۔ لومړۍ برخه
اخلو۔ \LR{$\lnot !A$} او~\LR{$!A$} د~\Elim{\lnot} دوه مقدمې دي:
\begin{prooftree}
\AxiomC{\LR{$\Discharge{\lnot !A}{2}$}}
\AxiomC{\LR{$\Discharge{!A}{3}$}}
\RightLabel{\Elim{\lnot}}
\BinaryInfC{\LR{$\lfalse$}}
\DeduceC{\LR{$!B$}}
\end{prooftree}
د~\FalseInt په کارولو \LR{$!B$} د پايلې په توګه اخستلے او څانګه بشپړولے
شو۔
\begin{prooftree}
\AxiomC{\LR{$\Discharge{\lnot !A \lor !B}{1}$}}
\AxiomC{\LR{$\Discharge{\lnot !A}{2}$}}
\AxiomC{\LR{$\Discharge{!A}{3}$}}
\RightLabel{\Elim{\lnot}}
\BinaryInfC{\LR{$\lfalse$}}
\RightLabel{\FalseInt}
\UnaryInfC{\LR{$!B$}}
\DischargeRule{\Intro{\lif}}{3}
\UnaryInfC{\LR{$!A \lif !B$}}
\AxiomC{\LR{$\Discharge{!B}{2}, \Discharge{!A}{4}$}}
\DeduceC{\LR{$!B$}}
\DischargeRule{\Intro{\lif}}{4}
\UnaryInfC{\LR{$!A \lif !B$}}
\DischargeRule{\Elim{\lor}}{2}
\TrinaryInfC{\LR{$!A \lif !B$}}
\DischargeRule{\Intro{\lif}}{1} 
\UnaryInfC{\LR{$(\lnot !A \lor !B) \lif (!A \lif !B)$}}
\end{prooftree}
\olpseganchor{OLP-0089-B013}{end}

\olpseganchor{OLP-0089-B014}{start}
اوس تر ټولو ښۍ څانګې ته ګورو۔ دلته بايد پوه شو چې د اشتقاق
تعريف \emph{د فرضيو ايستل روا بولي} خو \emph{لازمي ئې نۀ بولي}۔ په
بله وينا، که له فرضيو~\LR{$!A$} او~\LR{$!B$} څخه د يوې په کارولو \LR{$!B$} راوباسو
او بله و نۀ کاروو، دا روا ده۔ له~\LR{$!B$} څخه د~\LR{$!B$}~اشتقاق کول
څرګند دي: \LR{$!B$} په يوازې ځان داسې اشتقاق دے او هېڅ استنتاج
نۀ غواړي۔ نو فرضيه~\LR{$!A$} ساده ړنګولے شو۔
\begin{prooftree}
\AxiomC{\LR{$\Discharge{\lnot !A \lor !B}{1}$}}
\AxiomC{\LR{$\Discharge{\lnot !A}{2}$}}
\AxiomC{\LR{$\Discharge{!A}{3}$}}
\RightLabel{\Elim{\lnot}}
\BinaryInfC{\LR{$\lfalse$}}
\RightLabel{\FalseInt}
\UnaryInfC{\LR{$!B$}}
\DischargeRule{\Intro{\lif}}{3}
\UnaryInfC{\LR{$!A \lif !B$}}
\AxiomC{\LR{$\Discharge{!B}{2}$}}
\RightLabel{\Intro{\lif}}
\UnaryInfC{\LR{$!A \lif !B$}}
\DischargeRule{\Elim{\lor}}{2}
\TrinaryInfC{\LR{$!A \lif !B$}}
\DischargeRule{\Intro{\lif}}{1} 
\UnaryInfC{\LR{$(\lnot !A \lor !B) \lif (!A \lif !B)$}}
\end{prooftree}
پام وکړئ چې په بشپړ اشتقاق کښې د تر ټولو ښي اړخ
\Intro{\lif} استنتاج په حقيقت کښې هېڅ فرضيه نۀ باسي۔
\end{ex}
\olpseganchor{OLP-0089-B014}{end}

\olpseganchor{OLP-0089-B015}{start}
\begin{ex}
تر اوسه مو د~\FalseCl{} قاعدې ته اړتيا نۀ ده پېښه شوې۔ دا ځکه
ځانګړې ده چې داسې فرضيه ايستل راکوي چې د قاعدې د پايلې فرعي
فارمول نۀ وي۔ دا له~\FalseInt{} قاعدې سره نژدې تړلې ده۔ په
حقيقت کښې \FalseInt{} د~\FalseCl{} قاعدې ځانګړے حالت دے---د
“شهودي منطق” په نامه يو منطق شته چې يوازې \FalseInt{} پکې روا دے۔
کله چې بله هېڅ لار کار نۀ کوي، \FalseCl{} وروستۍ وسيله ده۔ مثلاً
فرض کړئ غواړو \LR{$!A \lor \lnot !A$}~اشتقاق کړو۔ زموږ عادي تګلاره
به دا وي چې د~\LR{$!A \lor \lnot !A$} د~\LR{$\Intro{\lor}$} په کارولو
اشتقاق کړو۔ خو دا
غواړي چې له هېڅ فرضیې پرته يا~\LR{$!A$} يا~\LR{$\lnot !A$}~اشتقاق کړو،
او دا نۀ کېږي۔ \FalseCl{} مو ژغوري!
\begin{prooftree}
  \AxiomC{\LR{$\Discharge{\lnot(!A \lor \lnot !A)}{1}$}}
  \DeduceC{\LR{$\lfalse$}}
  \DischargeRule{\FalseCl}{1}
  \UnaryInfC{\LR{$!A \lor \lnot !A$}}
\end{prooftree}
اوس له~\LR{$\lnot(!A \lor \lnot !A)$} څخه د~\LR{$\lfalse$} يو
اشتقاق لټوو۔ ځکه \LR{$\lfalse$} د~\LR{$\Elim{\lnot}$} پايله ده، نو
کېدے شي دا هڅه وکړو:
\begin{prooftree}
  \AxiomC{\LR{$\Discharge{\lnot(!A \lor \lnot !A)}{1}$}}
  \DeduceC{\LR{$\lnot !A$}}
  \AxiomC{\LR{$\Discharge{\lnot(!A \lor \lnot !A)}{1}$}}
  \DeduceC{\LR{$!A$}}
  \RightLabel{\Elim{\lnot}}
  \BinaryInfC{\LR{$\lfalse$}}
  \DischargeRule{\FalseCl}{1}
  \UnaryInfC{\LR{$!A \lor \lnot !A$}}
\end{prooftree}
د~\LR{$\lnot !A$} د اشتقاق موندلو تګلاره د~\LR{$\Intro{\lnot}$}
کارونه غواړي:
\begin{prooftree}
  \AxiomC{\LR{$\Discharge{\lnot(!A \lor \lnot !A)}{1}, \Discharge{!A}{2}$}}
  \DeduceC{\LR{$\lfalse$}}
  \DischargeRule{\Intro{\lnot}}{2}
  \UnaryInfC{\LR{$\lnot !A$}}
  \AxiomC{\LR{$\Discharge{\lnot(!A \lor \lnot !A)}{1}$}}
  \DeduceC{\LR{$!A$}}
  \RightLabel{\Elim{\lnot}}
  \BinaryInfC{\LR{$\lfalse$}}
  \DischargeRule{\FalseCl}{1}
  \UnaryInfC{\LR{$!A \lor \lnot !A$}}
\end{prooftree}
دلته \LR{$\lfalse$} په اسانۍ اخستلے شو: \LR{$\Elim{\lnot}$} پر
\LR{$\lnot(!A \lor \lnot !A)$} او~\LR{$!A \lor \lnot !A$} کاروو؛ دويمه
جمله زموږ له نوې فرضيې~\LR{$!A$} څخه د~\LR{$\Intro{\lor}$} په وسيله راځي:
\begin{prooftree}
  \AxiomC{\LR{$\Discharge{\lnot(!A \lor \lnot !A)}{1}$}}
  \AxiomC{\LR{$\Discharge{!A}{2}$}}
  \RightLabel{\Intro{\lor}}
  \UnaryInfC{\LR{$!A \lor \lnot !A$}}
  \RightLabel{\Elim{\lnot}}
  \BinaryInfC{\LR{$\lfalse$}}
  \DischargeRule{\Intro{\lnot}}{2}
  \UnaryInfC{\LR{$\lnot !A$}}
  \AxiomC{\LR{$\Discharge{\lnot(!A \lor \lnot !A)}{1}$}}
  \DeduceC{\LR{$!A$}}
  \RightLabel{\Elim{\lnot}}
  \BinaryInfC{\LR{$\lfalse$}}
  \DischargeRule{\FalseCl}{1}
  \UnaryInfC{\LR{$!A \lor \lnot !A$}}
\end{prooftree}
په ښي لوري کښې همدا تګلاره کاروو، خو \LR{$!A$} د~\FalseCl په وسيله
اخلو:
\begin{prooftree}
  \AxiomC{\LR{$\Discharge{\lnot(!A \lor \lnot !A)}{1}$}}
  \AxiomC{\LR{$\Discharge{!A}{2}$}}
  \RightLabel{\Intro{\lor}}
  \UnaryInfC{\LR{$!A \lor \lnot !A$}}
  \RightLabel{\Elim{\lnot}}
  \BinaryInfC{\LR{$\lfalse$}}
  \DischargeRule{\Intro{\lnot}}{2}
  \UnaryInfC{\LR{$\lnot !A$}}
  \AxiomC{\LR{$\Discharge{\lnot(!A \lor \lnot !A)}{1}$}}
  \AxiomC{\LR{$\Discharge{\lnot !A}{3}$}}
  \RightLabel{\Intro{\lor}}
  \UnaryInfC{\LR{$!A \lor \lnot !A$}}
  \RightLabel{\Elim{\lnot}}
  \BinaryInfC{\LR{$\lfalse$}}
  \DischargeRule{\FalseCl}{3}
  \UnaryInfC{\LR{$!A$}}
  \RightLabel{\Elim{\lnot}}
  \BinaryInfC{\LR{$\lfalse$}}
  \DischargeRule{\FalseCl}{1}
  \UnaryInfC{\LR{$!A \lor \lnot !A$}}
\end{prooftree}
\end{ex}
\olpseganchor{OLP-0089-B015}{end}

\olpseganchor{OLP-0089-B016}{start}
\begin{prob}
داسې اشتقاقونه ورکړئ چې لاندې وښيي:
\begin{enumerate}
\item \LR{$!A \land (!B \land !C) \Proves (!A \land !B) \land !C$}.
\item \LR{$!A \lor (!B \lor !C) \Proves (!A \lor !B) \lor !C$}.
\item \LR{$!A \lif (!B \lif !C) \Proves !B \lif (!A \lif !C)$}.
\item \LR{$!A \Proves \lnot\lnot !A$}.
\end{enumerate}
\end{prob}
\olpseganchor{OLP-0089-B016}{end}

\olpseganchor{OLP-0089-B017}{start}
\begin{prob}
داسې اشتقاقونه ورکړئ چې لاندې وښيي:
\begin{enumerate}
\item \LR{$(!A \lor !B) \lif !C \Proves !A \lif !C$}.
\item \LR{$(!A \lif !C) \land (!B \lif !C) \Proves (!A \lor !B) \lif !C$}.
\item \LR{$\Proves \lnot(!A \land \lnot !A)$}.
\item \LR{$!B \lif !A \Proves \lnot !A \lif \lnot !B$}.
\item \LR{$\Proves (!A \lif \lnot !A) \lif \lnot !A$}.
\item \LR{$\Proves \lnot(!A \lif !B) \lif \lnot !B$}.
\item \LR{$!A \lif !C \Proves \lnot (!A \land \lnot !C)$}.
\item \LR{$!A \land \lnot !C \Proves \lnot (!A \lif !C)$}.
\item \LR{$!A \lor !B, \lnot !B \Proves !A$}.
\item \LR{$\lnot !A \lor \lnot !B \Proves \lnot(!A \land !B)$}.
\item \LR{$\Proves (\lnot !A \land \lnot !B) \lif\lnot(!A \lor !B)$}.
\item \LR{$\Proves \lnot(!A \lor !B) \lif (\lnot !A \land \lnot !B)$}.
\end{enumerate}
\end{prob}
\olpseganchor{OLP-0089-B017}{end}

\olpseganchor{OLP-0089-B018}{start}
\begin{prob}
داسې اشتقاقونه ورکړئ چې لاندې وښيي:
\begin{enumerate}
\item \LR{$\lnot(!A \lif !B) \Proves !A$}.
\item \LR{$\lnot(!A \land !B) \Proves \lnot !A \lor \lnot !B$}.
\item \LR{$!A \lif !B \Proves \lnot !A \lor !B$}.
\item \LR{$\Proves \lnot \lnot !A \lif !A$}.
\item \LR{$!A \lif !B, \lnot !A \lif !B \Proves !B$}.
\item \LR{$(!A \land !B) \lif !C \Proves (!A \lif !C) \lor (!B \lif !C)$}.
\item \LR{$(!A \lif !B) \lif !A \Proves !A$}.
\item \LR{$\Proves (!A \lif !B) \lor (!B \lif !C)$}.
\end{enumerate}
(دا ټول د~\LR{$\FalseCl$} قاعده غواړي۔)
\end{prob}
\olpseganchor{OLP-0089-B018}{end}

\olpseganchor{OLP-0090-B004}{start}
\olfileid{fol}{ntd}{prq}
\olpseganchor{OLP-0090-B004}{end}

\olpseganchor{OLP-0090-B005}{start}
\olsection{له سورونو سره اشتقاقونه}
\olpseganchor{OLP-0090-B005}{end}

\olpseganchor{OLP-0090-B006}{start}
\begin{ex}
له سورونو سره د کار پر مهال بايد ډاډ واخلو چې د ځانګړي متغير شرط
نۀ ماتوو، او کله کله د دې دپاره د ځينو استنتاجونو د عملي کولو ترتيب
بدلول را بدلول غواړي۔ په عمومي توګه ګټه کوي چې لومړے د هغو قاعدو
د حل هڅه وکړو چې د ځانګړي متغير شرط ورباندې لګېږي (په بشپړ ثبوت
کښې به دوئ لاندې راځي)۔
\olpseganchor{OLP-0090-B006}{end}

\olpseganchor{OLP-0090-B007}{start}
راځئ وګورو چې د فارمول~\LR{$\lexists[x][\lnot !A(x)] \lif \lnot
\lforall[x][!A(x)]$} يو اشتقاق څنګه ورکوو۔ د معمول په شان
داسې پيل کوو:
\begin{prooftree}
\AxiomC{}
\UnaryInfC{\LR{$\lexists[x][\lnot !A(x)]\lif \lnot \lforall[x][!A(x)]$}}
\end{prooftree}
لومړے ليکو چې د~\Intro{\lif} د قاعدې په وسيله د وروستي پړاو د
توجيه کولو دپاره څه پکار دي۔
\begin{prooftree}
\AxiomC{\LR{$\Discharge{\lexists[x][\lnot !A(x)]}{1}$}}
\DeduceC{\LR{$\lnot \lforall[x][!A(x)]$}}
\DischargeRule{\Intro{\lif}}{1}
\UnaryInfC{\LR{$\lexists[x][\lnot !A(x)]\lif \lnot \lforall[x][!A(x)]$}}
\end{prooftree}
ځکه پر~\LR{$\lnot \lforall[x][!A(x)]$} د کارولو دپاره څرګنده قاعده نۀ
شته، اشتقاق داسې برابروو چې د~\Elim{\lexists} قاعده
وکارولے شو۔ دلته بايد د ځانګړي متغير شرط ته پام وکړو او داسې ثابت
وټاکو چې په~\LR{$\lexists[x][!A(x)]$} يا د هغې په هرې اړوندې فرضيې کښې
نۀ راځي۔ (خو ځکه چې هېڅ ثابت سمبولونه نۀ دي راغلي، هر انتخاب سم
دے۔)
\begin{prooftree}
\AxiomC{\LR{$\Discharge{\lexists[x][\lnot !A(x)]}{1}$}}
\AxiomC{\LR{$\Discharge{\lnot !A(a)}{2}$}}
\DeduceC{\LR{$\lnot \lforall[x][!A(x)]$}}
\DischargeRule{\Elim{\lexists}}{2}
\BinaryInfC{\LR{$\lnot \lforall[x][!A(x)]$}}
\DischargeRule{\Intro{\lif}}{1}
\UnaryInfC{\LR{$\lexists[x][\lnot !A(x)] \lif \lnot \lforall[x][!A(x)]$}}
\end{prooftree}
د~\LR{$\lnot \lforall[x][!A(x)]$} د راايستلو دپاره د~\Intro{\lnot}
قاعده کارول غواړو: دا غواړي چې يو تناقض راوباسو، او کېدے شي
\LR{$\lforall[x][!A(x)]$} د يوې اضافي فرضيې په توګه وکاروو۔ البته په دې
تناقض کښې هغه فرضيه~\LR{$\lnot !A(a)$} هم راتلے شي چې د
\Elim{\lexists} استنتاج به ئې وباسي۔ داسې ئې برابروو:
\begin{prooftree}
\AxiomC{\LR{$\Discharge{\lexists[x][\lnot !A(x)]}{1}$}}
\AxiomC{\LR{$\Discharge{\lnot !A(a)}{2}, \Discharge{\lforall[x][!A(x)]}{3}$}}
\DeduceC{\LR{$\lfalse$}}
\DischargeRule{\Intro{\lnot}}{3}
\UnaryInfC{\LR{$\lnot \lforall[x][!A(x)]$}}
\DischargeRule{\Elim{\lexists}}{2}
\BinaryInfC{\LR{$\lnot \lforall[x][!A(x)]$}}
\DischargeRule{\Intro{\lif}}{1}
\UnaryInfC{\LR{$\lexists[x][\lnot !A(x)]\lif \lnot \lforall[x][!A(x)]$}}
\end{prooftree}
داسې ښکاري چې تناقض ته نژدې يو۔ تر ټولو اسانه قاعده
\Elim{\lforall} ده چې د ځانګړي متغير هېڅ شرط نۀ لري۔ ځکه د عمومي
سور لرونکي~\LR{$x$} پر ځاے هره تړلې اصطلاح کارولے شو، ښه ده چې هماغه
~\LR{$a$} وکاروو څو تناقض ته ورسېږو۔
\begin{prooftree}
\AxiomC{\LR{$\Discharge{\lexists[x][\lnot !A(x)]}{1}$}}
\AxiomC{\LR{$\Discharge{\lnot !A(a)}{2}$}}
\AxiomC{\LR{$\Discharge{\lforall[x][!A(x)]}{3}$}}
\RightLabel{\Elim{\lforall}}
\UnaryInfC{\LR{$!A(a)$}}
\RightLabel{\Elim{\lnot}}
\BinaryInfC{\LR{$\lfalse$}}
\DischargeRule{\Intro{\lnot}}{3}
\UnaryInfC{\LR{$\lnot \lforall[x][!A(x)]$}}
\DischargeRule{\Elim{\lexists}}{2}
\BinaryInfC{\LR{$\lnot \lforall[x][!A(x)]$}}
\DischargeRule{\Intro{\lif}}{1}
\UnaryInfC{\LR{$\lexists[x][\lnot !A(x)]\lif \lnot \lforall[x][!A(x)]$}}
\end{prooftree}
\olpseganchor{OLP-0090-B007}{end}

\olpseganchor{OLP-0090-B008}{start}
په ځانګړي ډول له سورونو سره د کار پر مهال اوس دا بيا کتل مهم دي
چې د ځانګړي متغير شرط نۀ دے مات شوے۔ دلته دغه شرط يوازې پر
\Elim{\exists} لګېدۀ، او ځانګړے متغير~\LR{$a$} په هغو فرضيو کښې نۀ
راځي چې استنتاج ورپورې تړلے دے؛ نو دا سم اشتقاق دے۔
\end{ex}
\olpseganchor{OLP-0090-B008}{end}

\olpseganchor{OLP-0090-B009}{start}
\begin{ex}
کله کله له نورو فارمولونه څخه فارمول راايستلے شو۔ په دغو
حالتونو کښې کېدے شي نۀ ايستل شوې فرضيې ولرو۔ د خپلې وروستۍ موخې
تر څنګ د فرضيو حساب ساتل هم مهم دي۔
\olpseganchor{OLP-0090-B009}{end}

\olpseganchor{OLP-0090-B010}{start}
راځئ وګورو چې له فرضيو~\LR{$\lexists[x][(!A(x) \land !B(x))]$} او
\LR{$\lforall[x][(!B(x) \lif !C(x,b))]$} څخه د فارمول
\LR{$\lexists[x][!C(x,b)]$} يو اشتقاق څنګه ورکوو۔ د معمول په
شان پايله لاندې ليکو۔
\begin{prooftree}
\AxiomC{}
\UnaryInfC{\LR{$\lexists[x][!C(x,b)]$}}
\end{prooftree}
\olpseganchor{OLP-0090-B010}{end}

\olpseganchor{OLP-0090-B011}{start}
د کار دپاره دوه مقدمې لرو۔ د لومړۍ د کارولو، يعنې له
\LR{$\lexists[x][(!A(x) \land !B(x))]$} څخه د
\LR{$\lexists[x][!C(x, b)]$} د اشتقاق موندلو دپاره به د
\Elim{\lexists} قاعده کاروو۔ ځکه د ځانګړي متغير شرط لري، دا قاعده
لومړۍ کاروو۔ لاندې لاس ته راځي:
\begin{prooftree}
\AxiomC{\LR{$\lexists[x][(!A(x) \land !B(x))]$}}
\AxiomC{\LR{$\Discharge{!A(a) \land !B(a)}{1}$}}
\DeduceC{\LR{$\lexists[x][!C(x,b)]$}}
\DischargeRule{\Elim{\lexists}}{1}
\BinaryInfC{\LR{$\lexists[x][!C(x,b)]$}}
\end{prooftree}
هغه دوه فرضيې چې ورسره کار کوو، \LR{$!B$} شريکوي۔ ښايي اوس د
\Elim{\land} په کارولو د~\LR{$!B(a)$} بېلول ګټور وي۔
\begin{prooftree}
\AxiomC{\LR{$\lexists[x][(!A(x) \land !B(x)])$}}
\AxiomC{\LR{$\Discharge{!A(a) 
\land !B(a)}{1}$}}
\RightLabel{\Elim{\land}}
\UnaryInfC{\LR{$!B(a)$}}
\DeduceC{\LR{$\lexists[x][!C(x,b)]$}}
\DischargeRule{\Elim{\lexists}}{1}
\BinaryInfC{\LR{$\lexists[x][!C(x,b)]$}}
\end{prooftree}
\olpseganchor{OLP-0090-B011}{end}

\olpseganchor{OLP-0090-B012}{start}
د کار دپاره دويمه فرضيه~\LR{$\lforall[x][(!B(x) \lif !C(x,b))]$} ده۔
ځکه د ځانګړي متغير شرط نۀ لري، د~\Elim{\lforall} په کارولو د
ثابت سمبول~\LR{$a$} په وسيله \LR{$x$} بېلګه کولے او~\LR{$!B(a) \lif !C(a, b)$}
اخستلے شو۔ اوس \LR{$!B(a) \lif !C(a,b)$} او~\LR{$!B(a)$} دواړه لرو۔ بل ګام
بايد د~\Elim{\lif} نېغه کارونه وي۔
\begin{prooftree}
\AxiomC{\LR{$\lexists[x][(!A(x) \land !B(x))]$}}
\AxiomC{\LR{$\lforall[x][(!B(x) \lif !C(x,b))]$}}
\RightLabel{\Elim{\lforall}}
\UnaryInfC{\LR{$!B(a) \lif !C(a,b)$}}
\AxiomC{\LR{$\Discharge{!A(a) 
\land !B(a)}{1}$}}
\RightLabel{\Elim{\land}}
\UnaryInfC{\LR{$!B(a)$}}
\RightLabel{\Elim{\lif}}
\BinaryInfC{\LR{$!C(a,b)$}}
\DeduceC{\LR{$\lexists[x][!C(x,b)]$}}
\DischargeRule{\Elim{\lexists}}{1}
\insertBetweenHyps{\hspace{-5em}}
\BinaryInfC{\LR{$\lexists[x][!C(x,b)]$}}
\end{prooftree}
موخې ته ډېر نژدې يو!{} د~\Intro{\lexists} يوه کارونه موخې ته مو
رسوي۔
\begin{prooftree}
\AxiomC{\LR{$\lexists[x][(!A(x) \land !B(x))]$}}
\AxiomC{\LR{$\lforall[x][(!B(x) \lif !C(x,b))]$}}
\RightLabel{\Elim{\lforall}}
\UnaryInfC{\LR{$!B(a) \lif !C(a,b)$}}
\AxiomC{\LR{$\Discharge{!A(a) 
\land !B(a)}{1}$}}
\RightLabel{\Elim{\land}}
\UnaryInfC{\LR{$!B(a)$}}
\RightLabel{\Elim{\lif}}
\BinaryInfC{\LR{$!C(a,b)$}}
\RightLabel{\Intro{\lexists}}
\UnaryInfC{\LR{$\lexists[x][!C(x,b)]$}}
\DischargeRule{\Elim{\lexists}}{1}
\insertBetweenHyps{\hspace{-5em}}
\BinaryInfC{\LR{$\lexists[x][!C(x,b)]$}}
\end{prooftree}
ځکه په هر پړاو کښې مو ډاډ اخستے چې د ځانګړي متغير شرطونه نۀ دي
مات شوي، باور کولے شو چې دا سم اشتقاق دے۔
\end{ex}
\olpseganchor{OLP-0090-B012}{end}

\olpseganchor{OLP-0090-B013}{start}
\begin{ex}
له فرضيو~\LR{$\lforall[x][!A(x)] \lif \lexists[y][!B(y)]$} او
\LR{$\lnot\lexists[y][!B(y)]$} څخه د فارمول
\LR{$\lnot\lforall[x][!A(x)]$} يو اشتقاق ورکړئ۔ د معمول په شان
موخې فارمول لاندې ليکو۔
\begin{prooftree}
\AxiomC{}
\UnaryInfC{\LR{$\lnot\lforall[x][!A(x)]$}}
\end{prooftree}
د اشتقاق وروستۍ کرښه نفي ده، نو \Intro{\lnot} کارول غواړو۔
دا غواړي چې معلومه کړو تناقض څنګه اشتقاق کړو۔
\begin{prooftree}
\AxiomC{\LR{$\Discharge{\lforall[x][!A(x)]}{1}$}}
\DeduceC{\LR{$\lfalse$}}
\DischargeRule{\Intro{\lnot}}{1}
\UnaryInfC{\LR{$\lnot\lforall[x][!A(x)]$}}
\end{prooftree}
تر دې ځايه ښه يو۔ \Elim{\lforall} کارولے شو، خو څرګنده نۀ ده چې
موخې ته به مو ورسوي۔ پر ځاے ئې خپله يوه فرضيه کاروو۔
\LR{$\lforall[x][!A(x)] \lif \lexists[y][!B(y)]$} له
\LR{$\lforall[x][!A(x)]$} سره د~\Elim{\lif} قاعدې کارول راکوي۔
\begin{prooftree}
\AxiomC{\LR{$\lforall[x][!A(x)] \lif \lexists[y][!B(y)]$}}
\AxiomC{\LR{$\Discharge{\lforall[x][!A(x)]}{1}$}}
\RightLabel{\Elim{\lif}}
\BinaryInfC{\LR{$\lexists[y][!B(y)]$}}
\DeduceC{\LR{$\lfalse$}}
\DischargeRule{\Intro{\lnot}}{1}
\UnaryInfC{\LR{$\lnot\lforall[x][!A(x)]$}}
\end{prooftree}
اوس د کار دپاره يوه وروستۍ فرضيه لرو، او ښکاري چې د~\Elim{\lnot}
په کارولو تناقض ته مو رسولے شي۔
\begin{prooftree}
\AxiomC{\LR{$\lnot\lexists[y][!B(y)]$}}
\AxiomC{\LR{$\lforall[x][!A(x)] \lif \lexists[y][!B(y)]$}}
\AxiomC{\LR{$\Discharge{\lforall[x][!A(x)]}{1}$}}
\RightLabel{\Elim{\lif}}
\BinaryInfC{\LR{$\lexists[y][!B(y)]$}}
\RightLabel{\Elim{\lnot}}
\BinaryInfC{\LR{$\lfalse$}}
\DischargeRule{\Intro{\lnot}}{1}
\UnaryInfC{\LR{$\lnot\lforall[x][!A(x)]$}}
\end{prooftree}
\end{ex}
\olpseganchor{OLP-0090-B013}{end}

\olpseganchor{OLP-0090-B014}{start}
\begin{prob}
داسې اشتقاقونه ورکړئ چې لاندې وښيي:
\begin{enumerate}
\item \LR{$\Proves (\lforall[x][!A(x)] \land \lforall[y][!B(y)]) \lif
\lforall[z][(!A(z) \land !B(z))]$}.
\item \LR{$\Proves (\lexists[x][!A(x)] \lor \lexists[y][!B(y)]) \lif
\lexists[z][(!A(z) \lor !B(z))]$}.
\item \LR{$\lforall[x][(!A(x) \lif !B)] \Proves \lexists[y][!A(y)] \lif !B$}.
\item \LR{$\lforall[x][\lnot !A(x)] \Proves \lnot\lexists[x][!A(x)]$}.
\item \LR{$\Proves \lnot\lexists[x][!A(x)] \lif \lforall[x][\lnot !A(x)]$}.
\item \LR{$\Proves \lnot\lexists[x][\lforall[y][((!A(x,y) \lif \lnot
!A(y,y)) \land (\lnot !A(y,y) \lif !A(x,y)))]]$}.
\end{enumerate}
\end{prob}
\olpseganchor{OLP-0090-B014}{end}

\olpseganchor{OLP-0090-B015}{start}
\begin{prob}
داسې اشتقاقونه ورکړئ چې لاندې وښيي:
\begin{enumerate}
\item \LR{$\Proves \lnot\lforall[x][!A(x)] \lif \lexists[x][\lnot!A(x)]$}.
\item \LR{$(\lforall[x][!A(x)] \lif !B) \Proves \lexists[y][(!A(y) \lif !B)]$}.
\item \LR{$\Proves \lexists[x][(!A(x) \lif \lforall[y][!A(y)])]$}.
\end{enumerate}
(دا ټول د~\LR{$\FalseCl$} قاعده غواړي۔)
\end{prob}
\olpseganchor{OLP-0090-B015}{end}

\olpseganchor{OLP-0091-B004}{start}
\olfileid{fol}{ntd}{ptn}
\olpseganchor{OLP-0091-B004}{end}

\olpseganchor{OLP-0091-B005}{start}
\olsection{ثبوت-تيوريکي مفاهيم}
\olpseganchor{OLP-0091-B005}{end}

\olpseganchor{OLP-0091-B006}{start}
\begin{editorial}
دا برخه د فطري استنتاج د اثبات‌پذيرۍ اړيکې او سازګارۍ تعريفونه
راټولوي۔
\end{editorial}
\olpseganchor{OLP-0091-B006}{end}

\olpseganchor{OLP-0091-B007}{start}
\begin{explain}
لکه څنګه چې مو يو شمېر مهم معنايي مفاهيم (اعتبار، معنايي استلزام،
د صدق وړتيا) تعريف کړل، اوس ورسره سم \emph{ثبوت-تيوريکي مفاهيم}
تعريفوو۔ دا مفاهيم په جوړښتونه کښې د تړلي فارمولونه د پوره
کېدو له لارې نۀ، بلکې له نورو څخه د ځينو تړلي فارمولونه د
د اشتقاق وړتيا يا د اشتقاق نۀ وړتيا له لارې تعريفېږي۔ دا مهمه
موندنه وه چې دغه مفاهيم سره سمون خوري۔ د دې سمون منځپانګه د
\emph{سموالي} او \emph{بشپړتيا قضيې} دي۔
\end{explain}
\olpseganchor{OLP-0091-B007}{end}


\olpseganchor{OLP-0091-B008}{start}
\begin{defn}[قضيې]
يوه تړلے فارمول~\LR{$!A$} هله \emph{قضيه} ده چې په فطري استنتاج کښې
د~\LR{$!A$} داسې اشتقاق وي چې ټولې فرضيې پکې ايستل
شوې وي۔ که \LR{$!A$} قضيه وي، \LR{$\Proves !A$} ليکو، او که نۀ وي
\LR{$\Proves/ !A$} ليکو۔
\end{defn}
\olpseganchor{OLP-0091-B008}{end}

\olpseganchor{OLP-0091-B009}{start}
\begin{defn}[د اشتقاق وړتيا]
تړلے فارمول~\LR{$!A$} د تړلي فارمولونه له سټ~\LR{$\Gamma$} څخه،
\LR{$\Gamma \Proves !A$}، هله \emph{د اشتقاق وړ} ده چې داسې
اشتقاق وي چې پايله ئې~\LR{$!A$} وي او هره فرضيه پکې يا
ايستل شوې وي يا په~\LR{$\Gamma$} کښې وي۔ که \LR{$!A$} له~\LR{$\Gamma$}
څخه د اشتقاق وړ نۀ وي، \LR{$\Gamma \Proves/ !A$} ليکو۔
\end{defn}
\olpseganchor{OLP-0091-B009}{end}

\olpseganchor{OLP-0091-B010}{start}
\begin{defn}[سازګاري]
د تړلي فارمولونه يو سټ~\LR{$\Gamma$} هله او يوازې هله \emph{ناسازګار}
دے چې \LR{$\Gamma \Proves \lfalse$} وي۔ که \LR{$\Gamma$} ناسازګار نۀ وي،
يعنې \LR{$\Gamma \Proves/ \lfalse$} وي، نو \emph{سازګار} ئې بولو۔
\end{defn}
\olpseganchor{OLP-0091-B010}{end}

\olpseganchor{OLP-0091-B011}{start}
\begin{prop}[انعکاسيت]
\ollabel{prop:reflexivity}
که \LR{$!A \in \Gamma$} وي، نو \LR{$\Gamma \Proves !A$}۔
\end{prop}
\olpseganchor{OLP-0091-B011}{end}

\olpseganchor{OLP-0091-B012}{start}
\begin{proof}
فرضيه~\LR{$!A$} په يوازې ځان د~\LR{$!A$} داسې اشتقاق دے چې هره
نۀ ايستل شوې فرضيه ئې (يعنې \LR{$!A$}) په~\LR{$\Gamma$} کښې ده۔
\end{proof}
\olpseganchor{OLP-0091-B012}{end}
  

\olpseganchor{OLP-0091-B013}{start}
\begin{prop}[يکنواختي]
\ollabel{prop:monotonicity}
که \LR{$\Gamma \subseteq \Delta$} او \LR{$\Gamma \Proves !A$} وي، نو
\LR{$\Delta \Proves !A$}۔
\end{prop}
\olpseganchor{OLP-0091-B013}{end}

\olpseganchor{OLP-0091-B014}{start}
\begin{proof}
له~\LR{$\Gamma$} څخه د~\LR{$!A$} هر اشتقاق له~\LR{$\Delta$} څخه هم د~\LR{$!A$}
اشتقاق دے۔
\end{proof}
\olpseganchor{OLP-0091-B014}{end}

\olpseganchor{OLP-0091-B015}{start}
\begin{prop}[انتقاليت]
\ollabel{prop:transitivity}
که \LR{$\Gamma \Proves !A$} او \LR{$\{!A\} \cup \Delta \Proves !B$} وي، نو
\LR{$\Gamma \cup \Delta \Proves !B$}۔
\end{prop}
\olpseganchor{OLP-0091-B015}{end}

\olpseganchor{OLP-0091-B016}{start}
\begin{proof}
که \LR{$\Gamma \Proves !A$} وي، د~\LR{$!A$} داسې اشتقاق~\LR{$\delta_0$}
شته چې ټولې نۀ ايستل شوې فرضيې ئې په~\LR{$\Gamma$} کښې دي۔ که
\LR{$\{!A\} \cup \Delta \Proves !B$} وي، د~\LR{$!B$} داسې
اشتقاق~\LR{$\delta_1$} شته چې ټولې نۀ ايستل شوې فرضيې ئې
په~\LR{$\{!A\} \cup \Delta$} کښې دي۔ اوس لاندې وګورئ:
\begin{prooftree}
  \AxiomC{\LR{$\Delta, \Discharge{!A}{1}$}}
  \RightLabel{\LR{$\delta_1$}}
  \DeduceC{\LR{$!B$}}
  \DischargeRule{\Intro{\lif}}{1}
  \UnaryInfC{\LR{$!A \lif !B$}}
  \AxiomC{\LR{$\Gamma$}}
  \RightLabel{\LR{$\delta_0$}}
  \DeduceC{\LR{$!A$}}
  \RightLabel{\Elim{\lif}}
  \BinaryInfC{\LR{$!B$}}
\end{prooftree}
اوس ټولې نۀ ايستل شوې فرضيې په~\LR{$\Gamma \cup \Delta$} کښې دي؛
نو دا \LR{$\Gamma \cup \Delta \Proves !B$} ښيي۔
\end{proof}
\olpseganchor{OLP-0091-B016}{end}

\olpseganchor{OLP-0091-B017}{start}
کله چې \LR{$\Gamma = \{!A_1, !A_2, \ldots, !A_k\}$} يو متناهي سټ وي،
د~\LR{$\Gamma \Proves !B$} دپاره لنډه نښه
\LR{$!A_1,!A_2,\ldots,!A_k \Proves !B$} کارولے شو؛ په ځانګړي ډول
\LR{$!A \Proves !B$} دا معنٰی لري چې \LR{$\{!A\} \Proves !B$}۔
\olpseganchor{OLP-0091-B017}{end}

\olpseganchor{OLP-0091-B018}{start}
پام وکړئ که \LR{$\Gamma \Proves !A$} او \LR{$!A \Proves !B$} وي، نو
\LR{$\Gamma \Proves !B$}۔ همدارنګه، که
\LR{$!A_1, \dots, !A_n \Proves !B$} او د هر~\LR{$i$} دپاره
\LR{$\Gamma \Proves !A_i$} وي، نو \LR{$\Gamma \Proves !B$}۔
\olpseganchor{OLP-0091-B018}{end}

\olpseganchor{OLP-0091-B019}{start}
\begin{prop}
\ollabel{prop:incons}
لاندې شرطونه معادل دي۔
    \begin{enumerate}
      \item \( \Gamma \) ناسازګار دے۔
      \item د هرې تړلے فارمول~\( {!A} \) دپاره \( \Gamma \Proves {!A} \)۔
      \item د کومې تړلے فارمول~\( {!A} \) دپاره \( \Gamma \Proves {!A} \) او \( \Gamma \Proves \lnot {!A} \)۔
    \end{enumerate}
\end{prop}
\olpseganchor{OLP-0091-B019}{end}

\olpseganchor{OLP-0091-B020}{start}
\begin{proof}
تمرين۔
\end{proof}
\olpseganchor{OLP-0091-B020}{end}

\olpseganchor{OLP-0091-B021}{start}
\begin{prob}
\olref[fol][ntd][ptn]{prop:incons} ثابت کړئ۔
\end{prob}
\olpseganchor{OLP-0091-B021}{end}



\olpseganchor{OLP-0091-B023}{start}
\begin{prop}[تراکم]
  \ollabel{prop:proves-compact}
  \begin{enumerate}
  \item که \LR{$\Gamma \Proves !A$} وي، نو داسې متناهي فرعي سټ
    \LR{$\Gamma_0 \subseteq \Gamma$} شته چې \LR{$\Gamma_0 \Proves !A$} وي۔
  \item که د~\LR{$\Gamma$} هر متناهي فرعي سټ سازګار وي، نو \LR{$\Gamma$}
    سازګار دے۔
  \end{enumerate}
\end{prop}
\olpseganchor{OLP-0091-B023}{end}

\olpseganchor{OLP-0091-B024}{start}
\begin{proof}
  \begin{enumerate}
    \item که \LR{$\Gamma \Proves !A$} وي، نو له~\LR{$\Gamma$} څخه د~\LR{$!A$}
      کوم اشتقاق~\LR{$\delta$} شته۔ \LR{$\Gamma_0$} د~\LR{$\delta$} د
      نۀ ايستل شوې فرضيو سټ ونيسئ۔ ځکه هر اشتقاق متناهي
      دے، \LR{$\Gamma_0$} يوازې متناهي ډېرې تړلي فارمولونه لرلے شي۔ نو
      \LR{$\delta$} له متناهي~\LR{$\Gamma_0 \subseteq \Gamma$} څخه د~\LR{$!A$}
      اشتقاق دے۔
    \item دا د~(۱) د هغه ځانګړي حالت نقيض عکس دے چې
      \LR{$!A \ident \lfalse$} وي۔
  \end{enumerate}
\end{proof}
\olpseganchor{OLP-0091-B024}{end}

\olpseganchor{OLP-0092-B004}{start}
\olfileid{fol}{ntd}{prv}
\olpseganchor{OLP-0092-B004}{end}
      
\olpseganchor{OLP-0092-B005}{start}
\olsection{د اشتقاق وړتيا او سازګاري}
\olpseganchor{OLP-0092-B005}{end}

\olpseganchor{OLP-0092-B006}{start}
اوس به د اشتقاق وړتيا اړيکې يو شمېر خاصيتونه ثابت کړو۔ دوئ په
خپلواکه توګه په زړه پورې دي، خو هر يو به د بشپړتيا قضيې په ثبوت
کښې ونډه ولري۔
\olpseganchor{OLP-0092-B006}{end}

\olpseganchor{OLP-0092-B007}{start}
\begin{prop}\ollabel{prop:provability-contr}
  که \LR{$\Gamma \Proves !A$} وي او \LR{$\Gamma \cup \{!A\}$} ناسازګار وي،
  نو \LR{$\Gamma$} ناسازګار دے۔
\end{prop}
\olpseganchor{OLP-0092-B007}{end}

\olpseganchor{OLP-0092-B008}{start}
\begin{proof}
له~\LR{$\Gamma$} څخه د~\LR{$!A$} اشتقاق~\LR{$\delta_1$} او له
\LR{$\Gamma \cup \{!A\}$} څخه د~\LR{$\lfalse$} اشتقاق~\LR{$\delta_2$}
ونوموئ۔ بيا داسې اشتقاق کولے شو:
\begin{prooftree}
\AxiomC{\LR{$\Gamma, \Discharge{!A}{1}$}}
\RightLabel{\LR{$\delta_2$}}
\DeduceC{\LR{$\lfalse$}}
\DischargeRule{\Intro{\lnot}}{1}
\UnaryInfC{\LR{$\lnot !A$}}
\AxiomC{\LR{$\Gamma$}}
\RightLabel{\LR{$\delta_1$}}
\DeduceC{\LR{$!A$}}
\RightLabel{\Elim{\lnot}}
\BinaryInfC{\LR{$\lfalse$}}
\end{prooftree}
په نوي اشتقاق کښې فرضيه~\LR{$!A$}~ايستل شوې ده، نو دا
له~\LR{$\Gamma$} څخه اشتقاق دے۔
\end{proof}
\olpseganchor{OLP-0092-B008}{end}

\olpseganchor{OLP-0092-B009}{start}
\begin{prop}
\ollabel{prop:prov-incons}
\LR{$\Gamma \Proves !A$} هله او يوازې هله چې
\LR{$\Gamma \cup \{\lnot !A\}$} ناسازګار وي۔
\end{prop}
\olpseganchor{OLP-0092-B009}{end}

\olpseganchor{OLP-0092-B010}{start}
\begin{proof}
لومړے فرض کړئ \LR{$\Gamma \Proves !A$}؛ يعنې د~\LR{$!A$} داسې
اشتقاق~\LR{$\delta_0$} شته چې نۀ ايستل شوې فرضيې ئې
\LR{$\Gamma$} دي۔ له~\LR{$\Gamma \cup \{\lnot !A\}$} څخه د~\LR{$\lfalse$} يو
اشتقاق داسې اخلو:
\begin{prooftree}
  \AxiomC{\LR{$\lnot !A$}}
  \AxiomC{\LR{$\Gamma$}}
  \RightLabel{\LR{$\delta_0$}}
  \DeduceC{\LR{$!A$}}
  \RightLabel{\Elim{\lnot}}
  \BinaryInfC{\LR{$\lfalse$}}
\end{prooftree}
\olpseganchor{OLP-0092-B010}{end}

\olpseganchor{OLP-0092-B011}{start}
اوس فرض کړئ \LR{$\Gamma \cup \{\lnot !A\}$} ناسازګار دے، او
\LR{$\delta_1$} له~\LR{$\Gamma \cup \{\lnot !A\}$} کښې له نۀ ايستل شوې
فرضيو څخه د~\LR{$\lfalse$} اړوند اشتقاق دے۔ يوازې له~\LR{$\Gamma$}
څخه د~\LR{$!A$} يو اشتقاق د~\LR{$\FalseCl$} په کارولو اخلو:
\begin{prooftree}
  \AxiomC{\LR{$\Gamma, \Discharge{\lnot !A}{1}$}}
  \RightLabel{\LR{$\delta_1$}}
  \DeduceC{\LR{$\lfalse$}}
  \DischargeRule{\FalseCl}{1}
  \UnaryInfC{\LR{$!A$}}
\end{prooftree}
\end{proof}
\olpseganchor{OLP-0092-B011}{end}

\olpseganchor{OLP-0092-B012}{start}
\begin{prob}
ثابت کړئ چې \LR{$\Gamma \Proves \lnot !A$} هله او يوازې هله چې
\LR{$\Gamma \cup \{!A\}$} ناسازګار وي۔
\end{prob}
\olpseganchor{OLP-0092-B012}{end}

\olpseganchor{OLP-0092-B013}{start}
\begin{prop}\ollabel{prop:explicit-inc}
  که \LR{$\Gamma \Proves !A$} او \LR{$\lnot !A \in \Gamma$} وي، نو \LR{$\Gamma$}
  ناسازګار دے۔
\end{prop}
\olpseganchor{OLP-0092-B013}{end}

\olpseganchor{OLP-0092-B014}{start}
\begin{proof}
  فرض کړئ \LR{$\Gamma \Proves !A$} او \LR{$\lnot !A \in \Gamma$}۔ بيا
  له~\LR{$\Gamma$} څخه د~\LR{$!A$} کوم اشتقاق~\LR{$\delta$} شته۔ د
  \LR{$\Elim{\lnot}$} قاعدې دا ساده کارونه وګورئ:
  \begin{prooftree}
    \AxiomC{\LR{$\lnot !A$}}
    \AxiomC{\LR{$\Gamma$}}
    \RightLabel{\LR{$\delta$}}
    \DeduceC{\LR{$!A$}}
    \RightLabel{\Elim{\lnot}}
    \BinaryInfC{\LR{$\lfalse$}}
  \end{prooftree}
  ځکه \LR{$\lnot !A \in \Gamma$} او ټولې نۀ ايستل شوې فرضيې
  په~\LR{$\Gamma$} کښې دي، دا \LR{$\Gamma \Proves \lfalse$} ښيي۔
\end{proof}
\olpseganchor{OLP-0092-B014}{end}

\olpseganchor{OLP-0092-B015}{start}
\begin{prop}\ollabel{prop:provability-exhaustive}
  که \LR{$\Gamma \cup \{!A\}$} او \LR{$\Gamma \cup \{\lnot !A\}$} دواړه
  ناسازګار وي، نو \LR{$\Gamma$} ناسازګار دے۔
\end{prop}
\olpseganchor{OLP-0092-B015}{end}

\olpseganchor{OLP-0092-B016}{start}
\begin{proof}
په ترتيب سره له~\LR{$\Gamma \cup \{ !A \}$} څخه د~\LR{$\lfalse$} او له
\LR{$\Gamma \cup \{ \lnot !A \}$} څخه د~\LR{$\lfalse$}~اشتقاقونه
\LR{$\delta_1$} او~\LR{$\delta_2$} شته۔ بيا داسې اشتقاق کولے شو:
\begin{prooftree}
\AxiomC{\LR{$\Gamma, \Discharge{\lnot !A}{2}$}}
\RightLabel{\LR{$\delta_2$}}
\DeduceC{\LR{$\lfalse$}}
\DischargeRule{\Intro{\lnot}}{2}
\UnaryInfC{\LR{$\lnot \lnot !A$}}
\AxiomC{\LR{$\Gamma, \Discharge{!A}{1}$}}
\RightLabel{\LR{$\delta_1$}}
\DeduceC{\LR{$\lfalse$}}
\DischargeRule{\Intro{\lnot}}{1}
\UnaryInfC{\LR{$\lnot !A$}}
\RightLabel{\Elim{\lnot}}
\BinaryInfC{\LR{$\lfalse$}}
\end{prooftree}
ځکه فرضيې~\LR{$!A$} او~\LR{$\lnot !A$}~ايستل شوې دي، دا يوازې
له~\LR{$\Gamma$} څخه د~\LR{$\lfalse$} يو اشتقاق دے۔ نو \LR{$\Gamma$}
ناسازګار دے۔
\end{proof}
\olpseganchor{OLP-0092-B016}{end}

\olpseganchor{OLP-0093-B005}{start}
\olfileid{fol}{ntd}{ppr}
\olpseganchor{OLP-0093-B005}{end}

\olpseganchor{OLP-0093-B006}{start}
\olsection{د اشتقاق وړتيا او قضيې رابطونه}
\olpseganchor{OLP-0093-B006}{end}

\olpseganchor{OLP-0093-B007}{start}
\begin{explain}
  موږ ثابتوو چې د فطري استنتاج د اشتقاق وړتيا اړيکه~\LR{$\Proves$}
  دومره پياوړې ده چې د قضيې رابطونو په اړه ځينې بنسټيز حقيقتونه،
  لکه \LR{$!A \land !B \Proves !A$} او
  \LR{$!A, !A \lif !B \Proves !B$} (قياسي استنتاج)، ثابت کړي۔ دا
  حقيقتونه د بشپړتيا قضيې د ثبوت دپاره پکار دي۔
\end{explain}
\olpseganchor{OLP-0093-B007}{end}

\olpseganchor{OLP-0093-B008}{start}
\begin{prop}\ollabel{prop:provability-land}
  \begin{enumerate}
  \item \ollabel{prop:provability-land-left} دواړه
    \LR{$!A \land !B \Proves !A$} او \LR{$!A \land !B \Proves !B$}۔
  \item \ollabel{prop:provability-land-right}
    \LR{$!A, !B \Proves !A \land !B$}۔
  \end{enumerate}
\end{prop}
\olpseganchor{OLP-0093-B008}{end}

\olpseganchor{OLP-0093-B009}{start}
\begin{proof}
  \begin{enumerate}
  \item دواړه داسې اشتقاق کولے شو:
    \begin{prooftree}
      \AxiomC{\LR{$!A \land !B$}}
      \RightLabel{\Elim{\land}}
      \UnaryInfC{\LR{$!A$}}
      \DisplayProof\qquad\bottomAlignProof
      \AxiomC{\LR{$!A \land !B$}}
      \RightLabel{\Elim{\land}}
      \UnaryInfC{\LR{$!B$}}
    \end{prooftree}
  \item داسې ئې اشتقاق کولے شو:
    \begin{prooftree}
      \AxiomC{\LR{$!A$}}
      \AxiomC{\LR{$!B$}}
      \RightLabel{\Intro{\land}}
      \BinaryInfC{\LR{$!A \land !B$}}
    \end{prooftree}
  \end{enumerate}
\end{proof}
\olpseganchor{OLP-0093-B009}{end}


\olpseganchor{OLP-0093-B010}{start}
\begin{prop}\ollabel{prop:provability-lor}
  \begin{enumerate}
  \item \LR{$!A \lor !B, \lnot !A, \lnot !B$} ناسازګار دے۔
  \item دواړه \LR{$!A \Proves !A \lor !B$} او \LR{$!B \Proves !A \lor !B$}۔
  \end{enumerate}
\end{prop}
\olpseganchor{OLP-0093-B010}{end}

\olpseganchor{OLP-0093-B011}{start}
\begin{proof}
  \begin{enumerate}
  \item لاندې اشتقاق وګورئ:
    \begin{prooftree}
      \AxiomC{\LR{$!A \lor !B$}}
      \AxiomC{\LR{$\lnot !A$}}
      \AxiomC{\LR{$\Discharge{!A}{1}$}}
      \RightLabel{\Elim{\lnot}}
      \BinaryInfC{\LR{$\lfalse$}}
      \AxiomC{\LR{$\lnot !B$}}
      \AxiomC{\LR{$\Discharge{!B}{1}$}}
      \RightLabel{\Elim{\lnot}}
      \BinaryInfC{\LR{$\lfalse$}}
      \DischargeRule{\Elim{\lor}}{1}
      \TrinaryInfC{\LR{$\lfalse$}}
    \end{prooftree}
    دا له نۀ ايستل شوې فرضيو~\LR{$!A \lor !B$}، \LR{$\lnot !A$} او
    \LR{$\lnot !B$} څخه د~\LR{$\lfalse$} يو اشتقاق دے۔
  \item دواړه داسې اشتقاق کولے شو:
    \begin{prooftree}
      \AxiomC{\LR{$!A$}}
      \RightLabel{\Intro{\lor}}
      \UnaryInfC{\LR{$!A \lor !B$}}
      \DisplayProof\qquad\bottomAlignProof
      \AxiomC{\LR{$!B$}}
      \RightLabel{\Intro{\lor}}
      \UnaryInfC{\LR{$!A \lor !B$}}
    \end{prooftree}
  \end{enumerate}
\end{proof}
\olpseganchor{OLP-0093-B011}{end}

\olpseganchor{OLP-0093-B012}{start}
\begin{prop}\ollabel{prop:provability-lif}
  \begin{enumerate}
  \item \ollabel{prop:provability-lif-left}
    \LR{$!A, !A \lif !B \Proves !B$}۔
  \item \ollabel{prop:provability-lif-right}
    دواړه \LR{$\lnot !A \Proves !A \lif !B$} او
    \LR{$!B \Proves !A \lif !B$}۔
  \end{enumerate}
\end{prop}
\olpseganchor{OLP-0093-B012}{end}

\olpseganchor{OLP-0093-B013}{start}
\begin{proof}
  \begin{enumerate}
  \item داسې ئې اشتقاق کولے شو:
    \begin{prooftree}
      \AxiomC{\LR{$!A \lif !B$}}
      \AxiomC{\LR{$!A$}}
      \RightLabel{\Elim{\lif}}
      \BinaryInfC{\LR{$!B$}}
    \end{prooftree}
\olpseganchor{OLP-0093-B013}{end}
    
\olpseganchor{OLP-0093-B014}{start}
  \item دا د لاندې دوو اشتقاقونه په وسيله ښودل کېږي:
    \begin{prooftree}
      \AxiomC{\LR{$\lnot !A$}}
      \AxiomC{\LR{$\Discharge{!A}{1}$}}
      \RightLabel{\Elim{\lnot}}
      \BinaryInfC{\LR{$\lfalse$}}
      \RightLabel{\FalseInt}
      \UnaryInfC{\LR{$!B$}}
      \DischargeRule{\Intro{\lif}}{1}
      \UnaryInfC{\LR{$!A \lif !B$}}
      \DisplayProof\qquad\bottomAlignProof
      \AxiomC{\LR{$!B$}}
      \RightLabel{\Intro{\lif}}
      \UnaryInfC{\LR{$!A \lif !B$}}
    \end{prooftree}
    پام وکړئ چې \LR{$\Intro{\lif}$} د فرضيې~\LR{$!A$}~ايستل کولے
    شي، خو لازمه نۀ ده چې و ئې باسي۔
  \end{enumerate}
\end{proof}
\olpseganchor{OLP-0093-B014}{end}

\olpseganchor{OLP-0094-B005}{start}
\olfileid{fol}{ntd}{qpr}
\olpseganchor{OLP-0094-B005}{end}

\olpseganchor{OLP-0094-B006}{start}
\olsection{د اشتقاق وړتيا او سورونه}
\olpseganchor{OLP-0094-B006}{end}

\olpseganchor{OLP-0094-B007}{start}
\begin{explain}
  د بشپړتيا قضيه دا هم غواړي چې د فطري استنتاج قاعدې د~\LR{$\Proves$}
  په اړه هغه حقيقتونه ورکړي چې په دې برخه کښې ثابتېږي۔
\end{explain}
\olpseganchor{OLP-0094-B007}{end}

\olpseganchor{OLP-0094-B008}{start}
\begin{thm}
\ollabel{thm:strong-generalization} که \LR{$c$} داسې ثابت وي چې په
\LR{$\Gamma$} يا~\LR{$!A(x)$} کښې نۀ راځي، او \LR{$\Gamma \Proves !A(c)$} وي، نو
\LR{$\Gamma \Proves \lforall[x][!A(x)]$}۔
\end{thm}
\olpseganchor{OLP-0094-B008}{end}

\olpseganchor{OLP-0094-B009}{start}
\begin{proof}
\LR{$\delta$} له~\LR{$\Gamma$} څخه د~\LR{$!A(c)$} يو اشتقاق ونيسئ۔ د
\Intro{\lforall} استنتاج په زياتولو د~\LR{$\lforall[x][!A(x)]$} يو
اشتقاق اخلو۔ ځکه \LR{$c$} په~\LR{$\Gamma$} يا~\LR{$!A(x)$} کښې نۀ راځي،
د ځانګړي متغير شرط پوره کېږي۔
\end{proof}
\olpseganchor{OLP-0094-B009}{end}

\olpseganchor{OLP-0094-B010}{start}
\begin{prop}
\ollabel{prop:provability-quantifiers}
\begin{enumerate}
\item \LR{$!A(t) \Proves \lexists[x][!A(x)]$}.
\olpseganchor{OLP-0094-B010}{end}

\olpseganchor{OLP-0094-B011}{start}
\item \LR{$\lforall[x][!A(x)] \Proves !A(t)$}.
\end{enumerate}
\end{prop}
\olpseganchor{OLP-0094-B011}{end}

\olpseganchor{OLP-0094-B012}{start}
\begin{proof}
\begin{enumerate}
\item لاندې له~\LR{$!A(t)$} څخه د~\LR{$\lexists[x][!A(x)]$} يو
اشتقاق دے:
\begin{prooftree}
\AxiomC{\LR{$!A(t)$}}
\RightLabel{\Intro{\lexists}}
\UnaryInfC{\LR{$\lexists[x][!A(x)]$}}
\end{prooftree}
\olpseganchor{OLP-0094-B012}{end}

\olpseganchor{OLP-0094-B013}{start}
\item لاندې له~\LR{$\lforall[x][!A(x)]$} څخه د~\LR{$!A(t)$} يو
اشتقاق دے:
\begin{prooftree}
\AxiomC{\LR{$\lforall[x][!A(x)]$}}
\RightLabel{\Elim{\lforall}}
\UnaryInfC{\LR{$!A(t)$}}
\end{prooftree}
\end{enumerate}
\end{proof}
\olpseganchor{OLP-0094-B013}{end}

\olpseganchor{OLP-0095-B004}{start}
\olfileid{fol}{ntd}{sou}
\olpseganchor{OLP-0095-B004}{end}
      
\olpseganchor{OLP-0095-B005}{start}
\olsection{سموالے}
\olpseganchor{OLP-0095-B005}{end}

\olpseganchor{OLP-0095-B006}{start}
\begin{explain}
د فطري استنتاج په شان اشتقاق نظام هله \emph{سم} دے چې
هغه څيزونه اشتقاق نۀ کړي چې په حقيقت کښې ترې نۀ لازمېږي۔
نو سموالے د اشتقاق نظامونو د تضمين شوې خونديتوب يو ډول
خاصيت دے۔ د دې له مخې چې کوم ثبوت-تيوريکي خاصيت تر بحث لاندې وي،
مثلاً غواړو پوه شو چې:
\begin{enumerate}
\item هره د اشتقاق وړ~تړلے فارمول معتبره ده؛
\item که تړلے فارمول له ځينو نورو څخه د اشتقاق وړ وي، د هغو
  معنايي پايله هم ده؛
\item که د تړلي فارمولونه يو سټ ناسازګار وي، د صدق وړ نۀ دے۔
\end{enumerate}
دا د اشتقاق نظام مهم خاصيتونه دي۔ که له دغو څخه کوم يو
سم نۀ وي، اشتقاق نظام نيمګړے دے---له حده زيات څيزونه به
اشتقاق کوي۔ له دې امله د اشتقاق نظام سموالے ثابتول
تر ټولو زيات اهميت لري۔
\end{explain}
\olpseganchor{OLP-0095-B006}{end}

\olpseganchor{OLP-0095-B007}{start}
\begin{thm}[سموالے]
\ollabel{thm:soundness}
که \LR{$!A$} له نۀ ايستل شوې فرضيو~\LR{$\Gamma$} څخه د اشتقاق وړ وي،
نو \LR{$\Gamma \Entails !A$}۔
\end{thm}
\olpseganchor{OLP-0095-B007}{end}

\olpseganchor{OLP-0095-B008}{start}
\begin{proof}
فرض کړئ \LR{$\delta$} د~\LR{$!A$} يو اشتقاق دے۔ موږ په~\LR{$\delta$}
کښې د استنتاجونو پر شمېر استقرا کوو۔
\olpseganchor{OLP-0095-B008}{end}

\olpseganchor{OLP-0095-B009}{start}
د استقرا د بنسټ دپاره هغه حالت ثابتوو چې د استنتاجونو شمېر~\LR{$0$}
وي۔ په دې حالت کښې \LR{$\delta$} يوازې له يوې تړلے فارمول~\LR{$!A$} څخه
جوړ دے، يعنې يوه فرضيه ده۔ دغه فرضيه نۀ ايستل شوې ده؛ ځکه
فرضیې يوازې د استنتاجونو په وسيله ايستل کېدے شي، او دلته
هېڅ استنتاج نشته۔ نو هر
جوړښت~\LR{$\Struct{M}$}
چې د ثبوت ټولې نۀ ايستل شوې فرضيې پوره کوي، \LR{$!A$} هم پوره کوي۔
\olpseganchor{OLP-0095-B009}{end}

\olpseganchor{OLP-0095-B010}{start}
اوس د استقرا پړاو اخلو۔ فرض کړئ \LR{$\delta$}، \LR{$n$} استنتاجونه لري۔
د تر ټولو لاندې استنتاج مقدمه يا مقدمې د فرعي-اشتقاقونه په
وسيله اشتقاق شوې دي، او په هر يوه کښې تر~\LR{$n$} کم استنتاجونه
دي۔ د استقرا فرضيه منو: د تر ټولو لاندې استنتاج مقدمې د هغو
فرعي-اشتقاقونه له نۀ ايستل شوې فرضيو څخه لازمېږي چې په
دغو مقدمو پاے ته رسېږي۔ بايد وښيو چې پايله~\LR{$!A$} د ټول ثبوت له
نۀ ايستل شوې فرضيو څخه لازمېږي۔
\olpseganchor{OLP-0095-B010}{end}

\olpseganchor{OLP-0095-B011}{start}
د تر ټولو لاندې استنتاج د ډول له مخې حالتونه بېلوو۔ لومړے هغه
ممکن استنتاجونه څېړو چې يوازې يوه مقدمه لري۔
\olpseganchor{OLP-0095-B011}{end}

\olpseganchor{OLP-0095-B012}{start}
\begin{enumerate}
\item فرض کړئ وروستنے استنتاج \Intro{\lnot} دے: اشتقاق
  دا بڼه لري:
  \begin{prooftree}
    \AxiomC{\LR{$\Gamma, \Discharge{!A}{n}$}}
    \RightLabel{\LR{$\delta_1$}}
    \DeduceC{\LR{$\lfalse$}}
    \DischargeRule{\Intro{\lnot}}{n}
    \UnaryInfC{\LR{$\lnot !A$}}
  \end{prooftree}
  د استقرا د فرضيې له مخې \LR{$\lfalse$} د~\LR{$\delta_1$} له
  نۀ ايستل شوې فرضيو~\LR{$\Gamma \cup \{!A\}$} څخه لازمېږي۔ يو
  جوړښت~\LR{$\Struct{M}$}
  ونيسئ۔ بايد وښيو چې که
  \LR{$\Sat{M}{\Gamma}$} وي، نو
  \LR{$\Sat{M}{\lnot !A}$} دے۔ د خلف په توګه
  فرض کړئ \LR{$\Sat{M}{\Gamma}$}، خو
  \LR{$\Sat/{M}{\lnot !A}$}، يعنې
  \LR{$\Sat{M}{!A}$}۔ له دې به
  \LR{$\Sat{M}{\Gamma \cup \{!A\}}$} لازم شي۔
  دا زموږ د استقرا د فرضيې خلاف دے۔ نو
  \LR{$\Sat{M}{\lnot !A}$}۔
\olpseganchor{OLP-0095-B012}{end}

\olpseganchor{OLP-0095-B013}{start}
\item وروستنے استنتاج \Elim{\land} دے: دوه بڼې لري؛ له مقدمې
  \LR{$!A \land !B$} څخه يا~\LR{$!A$} يا~\LR{$!B$} راايستل کېدے شي۔ لومړے حالت
  ونيسئ۔ اشتقاق~\LR{$\delta$} داسې ښکاري:
  \begin{prooftree}
    \AxiomC{\LR{$\Gamma$}}
    \RightLabel{\LR{$\delta_1$}}
    \DeduceC{\LR{$!A \land !B$}}
    \RightLabel{\Elim{\land}}
    \UnaryInfC{\LR{$!A$}}
  \end{prooftree}
  د استقرا د فرضيې له مخې \LR{$!A \land !B$} د~\LR{$\delta_1$} له
  نۀ ايستل شوې فرضيو~\LR{$\Gamma$} څخه لازمېږي۔
  يو جوړښت~\LR{$\Struct{M}$} ونيسئ۔
  بايد وښيو چې که
  \LR{$\Sat{M}{\Gamma}$} وي، نو
  \LR{$\Sat{M}{!A}$} دے۔ فرض کړئ
  \LR{$\Sat{M}{\Gamma}$}۔ زموږ د استقرا د
  فرضيې (\LR{$\Gamma \Entails !A \land !B$}) له مخې پوهېږو چې
  \LR{$\Sat{M}{!A \land !B}$}۔ د تعريف له مخې
  \LR{$\Sat{M}{!A \land !B}$} هله او يوازې هله
  چې \LR{$\Sat{M}{!A}$} او
  \LR{$\Sat{M}{!B}$} دواړه وي۔ (هغه حالت چې له
  \LR{$!A \land !B$} څخه~\LR{$!B$} راايستل کېږي، په ورته ډول څېړل کېږي۔)
\olpseganchor{OLP-0095-B013}{end}
  
\olpseganchor{OLP-0095-B014}{start}
\item وروستنے استنتاج \Intro{\lor} دے: دوه بڼې لري؛
  \LR{$!A \lor !B$} يا له مقدمې~\LR{$!A$} يا له مقدمې~\LR{$!B$} څخه راايستل
  کېدے شي۔ لومړے حالت ونيسئ۔ اشتقاق دا بڼه لري:
  \begin{prooftree}
    \AxiomC{\LR{$\Gamma$}}
    \RightLabel{\LR{$\delta_1$}}
    \DeduceC{\LR{$!A$}}
    \RightLabel{\Intro{\lor}}
    \UnaryInfC{\LR{$!A \lor !B$}}
  \end{prooftree}
  د استقرا د فرضيې له مخې \LR{$!A$} د~\LR{$\delta_1$} له
  نۀ ايستل شوې فرضيو~\LR{$\Gamma$} څخه لازمېږي۔ يو
  جوړښت~\LR{$\Struct{M}$}
  ونيسئ۔ بايد وښيو چې که
  \LR{$\Sat{M}{\Gamma}$} وي، نو
  \LR{$\Sat{M}{!A \lor !B}$} دے۔ فرض کړئ
  \LR{$\Sat{M}{\Gamma}$}؛ بيا
  \LR{$\Sat{M}{!A}$}، ځکه
  \LR{$\Gamma \Entails !A$} (د استقرا فرضيه)۔ نو
  \LR{$\Sat{M}{!A \lor !B}$} هم بايد وي۔
  (هغه حالت چې \LR{$!A \lor !B$} له~\LR{$!B$} څخه راايستل کېږي، په ورته ډول
  څېړل کېږي۔)
\olpseganchor{OLP-0095-B014}{end}
  
\olpseganchor{OLP-0095-B015}{start}
\item وروستنے استنتاج \Intro{\lif} دے: \LR{$!A \lif !B$} له داسې فرعي
  ثبوت څخه راايستل کېږي چې فرضيه ئې~\LR{$!A$} او پايله ئې~\LR{$!B$} ده،
  يعنې:
  \begin{prooftree}
    \AxiomC{\LR{$\Gamma, \Discharge{!A}{n}$}}
    \RightLabel{\LR{$\delta_1$}}
    \DeduceC{\LR{$!B$}}
    \DischargeRule{\Intro{\lif}}{n}
    \UnaryInfC{\LR{$!A \lif !B$}}
  \end{prooftree}
  د استقرا د فرضيې له مخې \LR{$!B$} د~\LR{$\delta_1$} له
  نۀ ايستل شوې فرضيو څخه لازمېږي، يعنې
  \LR{$\Gamma \cup \{!A\} \Entails !B$}۔ يو
  جوړښت~\LR{$\Struct{M}$}
  ونيسئ۔ د~\LR{$\delta$}~نۀ ايستل شوې فرضيې يوازې~\LR{$\Gamma$} دي؛ ځکه
  \LR{$!A$} په وروستي استنتاج کښې ايستل شوې ده۔ نو بايد وښيو
  چې \LR{$\Gamma \Entails !A \lif !B$}۔ د خلف په توګه فرض کړئ چې د کوم
  جوړښت~\LR{$\Struct{M}$}
  دپاره \LR{$\Sat{M}{\Gamma}$}، خو
  \LR{$\Sat/{M}{!A \lif !B}$} وي۔ نو
  \LR{$\Sat{M}{!A}$} او
  \LR{$\Sat/{M}{!B}$}۔ خو د فرضيې له مخې
  \LR{$!B$} د~\LR{$\Gamma \cup \{!A\}$} معنايي پايله ده، يعنې
  \LR{$\Sat{M}{!B}$}، او دا تناقض دے۔ نو
  \LR{$\Gamma \Entails !A \lif !B$}۔
\olpseganchor{OLP-0095-B015}{end}

\olpseganchor{OLP-0095-B016}{start}
\item وروستنے استنتاج \FalseInt دے: دلته \LR{$\delta$} داسې پاے ته رسېږي:
  \begin{prooftree}
    \AxiomC{\LR{$\Gamma$}}
    \RightLabel{\LR{$\delta_1$}}
    \DeduceC{\LR{$\lfalse$}}
    \RightLabel{\FalseInt}
    \UnaryInfC{\LR{$!A$}}
  \end{prooftree}
  د استقرا د فرضيې له مخې \LR{$\Gamma \Entails \lfalse$}۔ بايد وښيو چې
  \LR{$\Gamma \Entails !A$}۔ فرض کړئ داسې نۀ وي؛ بيا د کوم
  \LR{$\Struct{M}$} دپاره
  \LR{$\Sat{M}{\Gamma}$} او
  \LR{$\Sat/{M}{!A}$} دي۔ خو
  کاذبي هېڅ جوړښت يا ارزونه نۀ پوره کوي؛ يعنې تل
  \LR{$\Sat/{M}{\lfalse}$} وي؛ نو له دې به
  \LR{$\Gamma \Entails/ \lfalse$} لازم شي، چې د استقرا د فرضيې خلاف دے۔
\olpseganchor{OLP-0095-B016}{end}

\olpseganchor{OLP-0095-B017}{start}
\item وروستنے استنتاج \FalseCl دے: تمرين۔
\olpseganchor{OLP-0095-B017}{end}

\olpseganchor{OLP-0095-B018}{start}

\item وروستنے استنتاج \Intro{\lforall} دے: بيا \LR{$\delta$} دا بڼه لري:
  \begin{prooftree}
    \AxiomC{\LR{$\Gamma$}}
    \RightLabel{\LR{$\delta_1$}}
    \DeduceC{\LR{$!A(a)$}}
    \RightLabel{\Intro{\lforall}}
    \UnaryInfC{\LR{$\lforall[x][!A(x)]$}}
  \end{prooftree}
  د استقرا د فرضيې له مخې مقدمه~\LR{$!A(a)$} د~نۀ ايستل شوې
  فرضيو~\LR{$\Gamma$} معنايي پايله ده۔ داسې جوړښت~\LR{$\Struct{M}$} ونيسئ
  چې \LR{$\Sat{M}{\Gamma}$} وي۔ بايد وښيو چې
  \LR{$\Sat{M}{\lforall[x][!A(x)]}$}۔ ځکه
  \LR{$\lforall[x][!A(x)]$} يوه تړلے فارمول ده، د دې معنٰی دا ده چې
  د هرې متغير-ګومارنې~\LR{$s$} دپاره بايد
  \LR{$\Sat{M}{!A(x)}[s]$} وښيو
  (\olref[syn][ass]{prop:sat-quant})۔ ځکه \LR{$\Gamma$} ټول له جملو جوړ
  دے، د هر~\LR{$!B \in \Gamma$} دپاره
  \LR{$\Sat{M}{!B}[s]$} د~\olref[syn][sat]{defn:satisfaction} له مخې
  دے۔ \LR{$\Struct{M'}$} داسې جوړښت ونيسئ چې د~\LR{$\Struct{M}$} په شان وي،
  خو \LR{$\Assign{a}{M'} = s(x)$} وي۔ ځکه \LR{$a$} په~\LR{$\Gamma$} کښې نۀ راځي،
  د~\olref[syn][ext]{cor:extensionality-sent} له مخې
  \LR{$\Sat{M'}{\Gamma}$}۔ ځکه \LR{$\Gamma \Entails !A(a)$}،
  \LR{$\Sat{M'}{!A(a)}$}۔ ځکه \LR{$!A(a)$} يوه تړلے فارمول ده، د
  \olref[syn][ass]{prop:sentence-sat-true} له مخې
  \LR{$\Sat{M'}{!A(a)}[s]$}۔ د
  \olref[syn][ext]{prop:ext-formulas} له مخې
  \LR{$\Sat{M'}{!A(x)}[s]$} هله او يوازې هله چې
  \LR{$\Sat{M'}{!A(a)}$} (ياد ولرئ \LR{$!A(a)$} هماغه
  \LR{$\Subst{!A(x)}{a}{x}$} دے)۔ نو \LR{$\Sat{M'}{!A(x)}[s]$}۔ ځکه \LR{$a$}
  په~\LR{$!A(x)$} کښې نۀ راځي، د
  \olref[syn][ext]{prop:extensionality} له مخې
  \LR{$\Sat{M}{!A(x)}[s]$}۔ خو \LR{$s$} خپل‌سری متغير-ګومارنه وه؛ نو
  \LR{$\Sat{M}{\lforall[x][!A(x)]}$}۔
\olpseganchor{OLP-0095-B018}{end}
  
\olpseganchor{OLP-0095-B019}{start}
\item وروستنے استنتاج \Intro{\lexists} دے: تمرين۔
\olpseganchor{OLP-0095-B019}{end}

\olpseganchor{OLP-0095-B020}{start}
\item وروستنے استنتاج \Elim{\forall} دے: تمرين۔

\end{enumerate}
\olpseganchor{OLP-0095-B020}{end}

\olpseganchor{OLP-0095-B021}{start}
اوس هغه ممکن استنتاجونه څېړو چې څو مقدمې لري:
\Elim{\lor}، \Intro{\land}، \Elim{\lnot}،
\Elim{\lif} او~\Elim{\lexists}۔
\begin{enumerate}
\olpseganchor{OLP-0095-B021}{end}

\olpseganchor{OLP-0095-B022}{start}
\item وروستنے استنتاج \Intro{\land} دے۔ \LR{$!A \land !B$} له
  مقدمو~\LR{$!A$} او~\LR{$!B$} څخه راايستل کېږي او \LR{$\delta$} دا بڼه لري:
  \begin{prooftree}
    \AxiomC{\LR{$\Gamma_1$}}
    \RightLabel{\LR{$\delta_1$}}
    \DeduceC{\LR{$!A$}}
    \AxiomC{\LR{$\Gamma_2$}}
    \RightLabel{\LR{$\delta_2$}}
    \DeduceC{\LR{$!B$}}
    \RightLabel{\Intro{\land}}
    \BinaryInfC{\LR{$!A \land !B$}}
  \end{prooftree}
  د استقرا د فرضيې له مخې \LR{$!A$} د~\LR{$\delta_1$} له
  نۀ ايستل شوې فرضيو~\LR{$\Gamma_1$} څخه، او~\LR{$!B$} د~\LR{$\delta_2$} له
  نۀ ايستل شوې فرضيو~\LR{$\Gamma_2$} څخه لازمېږي۔ د~\LR{$\delta$}
  نۀ ايستل شوې فرضيې \LR{$\Gamma_1 \cup \Gamma_2$} دي؛ نو بايد
  \LR{$\Gamma_1 \cup \Gamma_2 \Entails !A \land !B$} وښيو۔ داسې
  جوړښت~\LR{$\Struct{M}$}
  ونيسئ چې
  \LR{$\Sat{M}{\Gamma_1 \cup \Gamma_2}$} وي۔
  ځکه \LR{$\Sat{M}{\Gamma_1}$}، او
  \LR{$\Gamma_1 \Entails !A$}، نو
  \LR{$\Sat{M}{!A}$}۔ همدارنګه، ځکه
  \LR{$\Sat{M}{\Gamma_2}$} او
  \LR{$\Gamma_2 \Entails !B$}، نو
  \LR{$\Sat{M}{!B}$}۔ دواړه يو ځاے
  \LR{$\Sat{M}{!A \land !B}$} راکوي۔
\olpseganchor{OLP-0095-B022}{end}
  
\olpseganchor{OLP-0095-B023}{start}
\item وروستنے استنتاج \Elim{\lor} دے: تمرين۔
\olpseganchor{OLP-0095-B023}{end}

\olpseganchor{OLP-0095-B024}{start}
\item وروستنے استنتاج \Elim{\lif} دے۔ \LR{$!B$} له مقدمو
  \LR{$!A \lif !B$} او~\LR{$!A$} څخه راايستل کېږي۔ اشتقاق~\LR{$\delta$}
  داسې ښکاري:
  \begin{prooftree}
    \AxiomC{\LR{$\Gamma_1$}}
    \RightLabel{\LR{$\delta_1$}}
    \DeduceC{\LR{$!A \lif !B$}}
    \AxiomC{\LR{$\Gamma_2$}}
    \RightLabel{\LR{$\delta_2$}}
    \DeduceC{\LR{$!A$}}
    \RightLabel{\Elim{\lif}}
    \BinaryInfC{\LR{$!B$}}
  \end{prooftree}
  د استقرا د فرضيې له مخې \LR{$!A \lif !B$} د~\LR{$\delta_1$} له
  نۀ ايستل شوې فرضيو~\LR{$\Gamma_1$} څخه، او~\LR{$!A$} د~\LR{$\delta_2$} له
  نۀ ايستل شوې فرضيو~\LR{$\Gamma_2$} څخه لازمېږي۔ يو
  جوړښت~\LR{$\Struct{M}$}
  ونيسئ۔ بايد وښيو چې که
  \LR{$\Sat{M}{\Gamma_1 \cup \Gamma_2}$} وي، نو
  \LR{$\Sat{M}{!B}$} دے۔ فرض کړئ
  \LR{$\Sat{M}{\Gamma_1 \cup \Gamma_2}$}۔
  ځکه \LR{$\Gamma_1 \Entails !A \lif !B$}، نو
  \LR{$\Sat{M}{!A \lif !B}$}۔ ځکه
  \LR{$\Gamma_2 \Entails !A$}، نو
  \LR{$\Sat{M}{!A}$}۔ له دې څخه
  \LR{$\Sat{M}{!B}$} لازمېږي (ځکه که
  \LR{$\Sat/{M}{!B}$} وے، نو له
  \LR{$\Sat{M}{!A}$} سره به
  \LR{$\Sat/{M}{!A \lif !B}$} وے، چې د
  \LR{$\Sat{M}{!A \lif !B}$} خلاف دے)۔
\olpseganchor{OLP-0095-B024}{end}

\olpseganchor{OLP-0095-B025}{start}
\item وروستنے استنتاج \Elim{\lnot} دے: تمرين۔
\olpseganchor{OLP-0095-B025}{end}
  
\olpseganchor{OLP-0095-B026}{start}
\item وروستنے استنتاج \Elim{\lexists} دے: تمرين۔
\end{enumerate}
\end{proof}
\olpseganchor{OLP-0095-B026}{end}

\olpseganchor{OLP-0095-B027}{start}
\begin{prob}
د~\olref[fol][ntd][sou]{thm:soundness} ثبوت بشپړ کړئ۔
\end{prob}
\olpseganchor{OLP-0095-B027}{end}



\olpseganchor{OLP-0095-B029}{start}
\begin{cor}
\ollabel{cor:weak-soundness}
که \LR{$\Proves !A$} وي، نو \LR{$!A$} معتبره ده۔
\end{cor}
\olpseganchor{OLP-0095-B029}{end}

\olpseganchor{OLP-0095-B030}{start}
\begin{cor}
\ollabel{cor:consistency-soundness}
که \LR{$\Gamma$} د صدق وړ وي، نو سازګار دے۔
\end{cor}
\olpseganchor{OLP-0095-B030}{end}

\olpseganchor{OLP-0095-B031}{start}
\begin{proof}
نقيض عکس ثابتوو۔ فرض کړئ \LR{$\Gamma$} سازګار نۀ دے۔ بيا
\LR{$\Gamma \Proves \lfalse$}، يعنې له~\LR{$\Gamma$} څخه د
نۀ ايستل شوې فرضيو پر بنسټ د~\LR{$\lfalse$} يو اشتقاق شته۔
د~\olref{thm:soundness} له مخې هر
جوړښت~\LR{$\Struct{M}$}
چې \LR{$\Gamma$} پوره کوي، بايد \LR{$\lfalse$} هم پوره کړي۔ ځکه د هر
جوړښت~\LR{$\Struct{M}$}
دپاره \LR{$\Sat/{M}{\lfalse}$}، نو هېڅ
\LR{$\Struct{M}$}، \LR{$\Gamma$} نۀ شي پوره
کولے؛ يعنې \LR{$\Gamma$} د صدق وړ نۀ دے۔
\end{proof}
\olpseganchor{OLP-0095-B031}{end}

\olpseganchor{OLP-0096-B004}{start}
\olfileid{fol}{ntd}{ide}
\olpseganchor{OLP-0096-B004}{end}

\olpseganchor{OLP-0096-B005}{start}
\olsection{له عينيت سره اشتقاقونه}
\olpseganchor{OLP-0096-B005}{end}

\olpseganchor{OLP-0096-B006}{start}
له عينيت سره اشتقاقونه اضافي استنتاجي قاعدو ته اړتيا لري۔
\olpseganchor{OLP-0096-B006}{end}

\olpseganchor{OLP-0096-B007}{start}
\begin{defish}
\AxiomC{}
\RightLabel{\Intro{\eq}}
\UnaryInfC{\LR{$\eq[t][t]$}}
\DisplayProof
\hfill
\begin{tabular}{r}
\AxiomC{\LR{$\eq[t_1][t_2]$}}
\AxiomC{\LR{$!A(t_1)$}}
\RightLabel{\Elim{\eq}}
\BinaryInfC{\LR{$!A(t_2)$}}
\DisplayProof
\\[3ex]
\AxiomC{\LR{$\eq[t_1][t_2]$}}
\AxiomC{\LR{$!A(t_2)$}}
\RightLabel{\Elim{\eq}}
\BinaryInfC{\LR{$!A(t_1)$}}
\DisplayProof
\end{tabular}
\end{defish}
\olpseganchor{OLP-0096-B007}{end}

\olpseganchor{OLP-0096-B008}{start}
په پورته قاعدو کښې \LR{$t$}، \LR{$t_1$} او~\LR{$t_2$} تړلې اصطلاحګانې دي۔
د~\Intro{\eq} قاعده راکوي چې د~\LR{$\eq[t][t]$} په بڼه هره د يووالي
جمله نېغ په نېغه، له هېڅ فرضيې پرته، اشتقاق کړو۔
\olpseganchor{OLP-0096-B008}{end}

\olpseganchor{OLP-0096-B009}{start}
\begin{ex}
که \LR{$s$} او~\LR{$t$} تړلې اصطلاحګانې وي، نو
\LR{$!A(s), \eq[s][t] \Proves !A(t)$}:
\begin{prooftree}
\AxiomC{\LR{$\eq[s][t]$}}
\AxiomC{\LR{$!A(s)$}}
\RightLabel{\LR{$\Elim{\eq}$}}
\BinaryInfC{\LR{$!A(t)$}}
\end{prooftree}
دا به د «يو شان څيزونو د تعويض د اصل» يا د لايبنېڅ د قانون په نامه
اشنا وي۔
\end{ex}
\olpseganchor{OLP-0096-B009}{end}

\olpseganchor{OLP-0096-B010}{start}
\begin{prob}
ثابت کړئ چې \LR{$=$} متناظره او انتقالي دواړه ده؛ يعنې د
\LR{$\lforall[x][\lforall[y][(\eq[x][y] \lif
    \eq[y][x])]]$} او~\LR{$\lforall[x][\lforall[y][\lforall[z]((\eq[x][y]
    \land \eq[y][z]) \lif \eq[x][z])]]$}
اشتقاقونه ورکړئ۔
\end{prob}
\olpseganchor{OLP-0096-B010}{end}

\olpseganchor{OLP-0096-B011}{start}
\begin{ex}
موږ دا تړلے فارمول~اشتقاق کوو:
\begin{align*}
& \lforall[x][\lforall[y][((!A(x) \land !A(y)) \lif \eq[x][y])]]
\intertext{\RL{له دې تړلے فارمول څخه:}}
& \lexists[x][\lforall[y][(!A(y) \lif \eq[y][x])]]
\end{align*}
اشتقاق له لاندې څخه پورته جوړوو:
\begin{prooftree}
\AxiomC{\LR{$\lexists[x][\lforall[y][(!A(y) \lif \eq[y][x])]]
\quad \Discharge{!A(a) \land !A(b)}{1}$}}
\DeduceC{\LR{$\eq[a][b]$}}
\DischargeRule{\Intro{\lif}}{1}
\UnaryInfC{\LR{$((!A(a) \land !A(b)) \lif \eq[a][b])$}}
\RightLabel{\Intro{\lforall}}
\UnaryInfC{\LR{$\lforall[y][((!A(a) \land !A(y)) \lif \eq[a][y])]$}}
\RightLabel{\Intro{\lforall}}
\UnaryInfC{\LR{$\lforall[x][\lforall[y][((!A(x) \land !A(y)) \lif \eq[x][y])]]$}}
\end{prooftree}
اوس بايد اصلي فرضيه وکاروو: ځکه دا وجودي فارمول ده، د
منځنۍ پايلې~\LR{$\eq[a][b]$} د~اشتقاق کولو دپاره
\Elim{\lexists} کاروو۔
\begin{prooftree}
\AxiomC{\LR{$\lexists[x][\lforall[y][(!A(y) \lif \eq[y][x])]]$}}
\AxiomC{\LR{$\Discharge{\lforall[y][(!A(y) \lif \eq[y][c]])}{2}$}}
\noLine
\UnaryInfC{\LR{$\Discharge{!A(a) \land !A(b)}{1}$}}
\DeduceC{\LR{$\eq[a][b]$}}
\DischargeRule{\Elim{\lexists}}{2}
\BinaryInfC{\LR{$\eq[a][b]$}}
\DischargeRule{\Intro{\lif}}{1}
\UnaryInfC{\LR{$((!A(a) \land !A(b)) \lif \eq[a][b])$}}
\RightLabel{\Intro{\lforall}}
\UnaryInfC{\LR{$\lforall[y][((!A(a) \land !A(y)) \lif \eq[a][y])]$}}
\RightLabel{\Intro{\lforall}}
\UnaryInfC{\LR{$\lforall[x][\lforall[y][((!A(x) \land !A(y)) \lif \eq[x][y])]]$}}
\end{prooftree}
د پاس په ښي لوري کښې فرعي-اشتقاق د خپلو فرضيو په کارولو
بشپړېږي، څو \LR{$\eq[a][c]$} او~\LR{$\eq[b][c]$} وښيي۔ دا دوه جلا
اشتقاقونه غواړي۔ د~\LR{$\eq[a][c]$}~اشتقاق داسې دے:
\begin{prooftree}
\AxiomC{\LR{$\Discharge{\lforall[y][(!A(y) \lif \eq[y][c]])}{2}$}}
\RightLabel{\Elim{\lforall}}
\UnaryInfC{\LR{$!A(a) \lif \eq[a][c]$}}
\AxiomC{\LR{$\Discharge{!A(a) \land !A(b)}{1}$}}
\RightLabel{\Elim{\land}}
\UnaryInfC{\LR{$!A(a)$}}
\RightLabel{\Elim{\lif}}
\BinaryInfC{\LR{$\eq[a][c]$}}
\end{prooftree}
له~\LR{$\eq[a][c]$} او~\LR{$\eq[b][c]$} څخه
\LR{$\eq[a][b]$} د~\Elim{\eq} په وسيله اشتقاق کوو۔
\end{ex}
\olpseganchor{OLP-0096-B011}{end}

\olpseganchor{OLP-0096-B012}{start}
\begin{prob}
د لاندې فارمولونه~اشتقاقونه ورکړئ:
\begin{enumerate}
\item \LR{$\lforall[x][\lforall[y][((\eq[x][y] \land !A(x)) \lif !A(y))]]$}
\item \LR{$\lexists[x][!A(x)] \land \lforall[y][\lforall[z][((!A(y) \land
    !A(z)) \lif \eq[y][z])]] \lif \lexists[x][(!A(x) \land
  \lforall[y][(!A(y) \lif \eq[y][x])])]$}
\end{enumerate}
\end{prob}
\olpseganchor{OLP-0096-B012}{end}

\olpseganchor{OLP-0097-B004}{start}
\olfileid{fol}{ntd}{sid}
\olpseganchor{OLP-0097-B004}{end}

\olpseganchor{OLP-0097-B005}{start}
\olsection{له عينيت سره سموالے}
\olpseganchor{OLP-0097-B005}{end}

\olpseganchor{OLP-0097-B006}{start}
\begin{prop}
د~\LR{$\eq$} له قاعدو سره فطري استنتاج سم دے۔
\end{prop}
\olpseganchor{OLP-0097-B006}{end}

\olpseganchor{OLP-0097-B007}{start}
\begin{proof}
د~\LR{$\eq[t][t]$} په بڼه هره فارمول معتبره ده؛ ځکه د هر
جوړښت~\LR{$\Struct M$} دپاره \LR{$\Sat{M}{\eq[t][t]}$}۔ (پام وکړئ
چې اصطلاح~\LR{$t$} تړلې فرضوو، يعنې هېڅ متغير نۀ لري؛ نو
متغير-ګومارنې اغېز نۀ لري۔)
\olpseganchor{OLP-0097-B007}{end}

\olpseganchor{OLP-0097-B008}{start}
فرض کړئ په کوم اشتقاق کښې وروستنے استنتاج \Elim{\eq} دے،
يعنې اشتقاق لاندې بڼه لري:
\begin{prooftree}
  \AxiomC{\LR{$\Gamma_1$}}
  \RightLabel{\LR{$\delta_1$}}
  \DeduceC{\LR{$\eq[t_1][t_2]$}}
  \AxiomC{\LR{$\Gamma_2$}}
  \RightLabel{\LR{$\delta_2$}}
  \DeduceC{\LR{$!A(t_1)$}}
  \RightLabel{\Elim{\eq}}
  \BinaryInfC{\LR{$!A(t_2)$}}
\end{prooftree}
مقدمې~\LR{$\eq[t_1][t_2]$} او~\LR{$!A(t_1)$} په خپل وار له
نۀ ايستل شوې فرضيو~\LR{$\Gamma_1$} او~\LR{$\Gamma_2$} څخه
اشتقاق شوې دي۔ غواړو وښيو چې \LR{$!A(t_2)$} له
\LR{$\Gamma_1 \cup \Gamma_2$} څخه لازمېږي۔ داسې
جوړښت~\LR{$\Struct{M}$} ونيسئ چې
\LR{$\Sat{M}{\Gamma_1 \cup \Gamma_2}$} وي۔ د استقرا د فرضيې له مخې
\LR{$\Sat{M}{!A(t_1)}$} او~\LR{$\Sat{M}{\eq[t_1][t_2]}$}۔ نو
\LR{$\Value{t_1}{M} = \Value{t_2}{M}$}۔ \LR{$s$} هره متغير-ګومارنه ونيسئ،
او \LR{$m = \Value{t_1}{M} = \Value{t_2}{M}$}۔ د
\olref[fol][syn][ext]{prop:ext-formulas} له مخې
\LR{$\Sat{M}{!A(t_1)}[s]$} هله او يوازې هله چې
\LR{$\Sat{M}{!A(x)}[\Subst{s}{m}{x}]$}، او دا هله او يوازې هله چې
\LR{$\Sat{M}{!A(t_2)}[s]$}۔ ځکه \LR{$\Sat{M}{!A(t_1)}$}، نو
\LR{$\Sat{M}{!A(t_2)}$}۔
\end{proof}
\olpseganchor{OLP-0097-B008}{end}

\olpseganchor{OLP-0098-B004}{start}
\olchapter{fol}{tab}{تابلوګانې}
\olpseganchor{OLP-0098-B004}{end}

\olpseganchor{OLP-0098-B005}{start}
\begin{editorial}
دا څپرکے يو نښه لرونکے تحليلي تابلو نظام وړاندې کوي۔
\olpseganchor{OLP-0098-B005}{end}
  
\olpseganchor{OLP-0098-B006}{start}
د ثبوت د يوه نظام په توګه له تابلو سره اړوند مواد د شاملولو يا
ايستلو دپاره د~\texttt{prfTab} ټګ وکاروئ۔
\end{editorial}
\olpseganchor{OLP-0098-B006}{end}





\olpseganchor{OLP-0098-B009}{start}
%
  

\olpseganchor{OLP-0098-B009}{end}





\olpseganchor{OLP-0098-B012}{start}
%
  

\olpseganchor{OLP-0098-B012}{end}







\olpseganchor{OLP-0098-B016}{start}
%
  

\olpseganchor{OLP-0098-B016}{end}



\olpseganchor{OLP-0098-B018}{start}
%
  
  

\olpseganchor{OLP-0098-B018}{end}

\olpseganchor{OLP-0099-B004}{start}
\olfileid{fol}{tab}{rul}
\olpseganchor{OLP-0099-B004}{end}

\olpseganchor{OLP-0099-B005}{start}
\olsection{قاعدې او تابلوګانې}
\olpseganchor{OLP-0099-B005}{end}

\olpseganchor{OLP-0099-B006}{start}
تابلو په جوړښت کښې د~تړلے فارمول د رښتيا يا
دروغ کېدو د ممکنو لارو منظمه پلټنه ده۔ د تابلو بنسټيز توکي
نښه لرونکي فارمولونه دي: تړلي فارمولونه له يوه رښتياوالي ارزښت
«نښې» سره، چې يا~\LR{$\True$} وي يا~\LR{$\False$}۔ دا نښه لرونکي فارمولونه
په داسې ونه کښې ترتيبېږي چې لاندې خوا ته غځېږي۔
\olpseganchor{OLP-0099-B006}{end}

\olpseganchor{OLP-0099-B007}{start}
\begin{defn}
  \emph{نښه لرونکے فارمول} له يوه رښتياوالي ارزښت او~تړلے فارمول
  څخه جوړه ده، يعنې له دې دواړو څخه يوه:
  \[
  \sFmla{\True}{!A} \text{\RL{ يا }} \sFmla{\False}{!A}.
  \]
\end{defn}
\olpseganchor{OLP-0099-B007}{end}

\olpseganchor{OLP-0099-B008}{start}
په شهودي ډول \LR{$\sFmla{\True}{!A}$} داسې لوستلے شو چې «کېدے شي
\LR{$!A$} رښتيا وي»، او~\LR{$\sFmla{\False}{!A}$} داسې چې «کېدے شي
\LR{$!A$} دروغ وي» (په کوم جوړښت کښې)۔
\olpseganchor{OLP-0099-B008}{end}

\olpseganchor{OLP-0099-B009}{start}
په ونه کښې هر نښه لرونکے فارمول يا \emph{فرض} دے (چې د ونې په
تر ټولو پاسنۍ برخه کښې ليکل کېږي)، يا تر ځان له پاسه له کوم
نښه لرونکے فارمول څخه د يوې استنتاجي قاعدې په وسيله اخستل شوے
دے۔ د مخکښې فارمول د هر ممکن اصلي عملګر دپاره دوه
قاعدې شته: يوه د هغه حالت دپاره چې نښه~\LR{$\True$} وي، او بله د هغه
حالت دپاره چې نښه~\LR{$\False$} وي۔ ځينې قاعدې ونې ته څانګې ورکوي او
ځينې يوازې څانګې ته نښه لرونکي فارمولونه ورزياتوي۔ يوه قاعده
کېدے شي (او ډېر ځله بايد) سمدستي پر مخکښ نښه لرونکے فارمول نۀ،
بلکې له ريښې څخه د قاعدې د پلي کېدو تر ځاے پورې په څانګه کښې پر هر
نښه لرونکے فارمول پلې شي۔
\olpseganchor{OLP-0099-B009}{end}

\olpseganchor{OLP-0099-B010}{start}
يوه څانګه هله \emph{تړلې} ده چې
\LR{$\sFmla{\True}{!A}$} او~\LR{$\sFmla{\False}{!A}$} دواړه ولري۔
تړلے تابلو هغه دے چې هره څانګه ئې تړلې وي۔ د شهودي تعبير له
مخې هره څانګه يو ګډ امکان بيانوي، خو
\LR{$\sFmla{\True}{!A}$} او~\LR{$\sFmla{\False}{!A}$} يو ځاے ممکن نۀ دي۔
په بله وينا، که يوه څانګه تړلې وي، هغه امکان رد شوے دے چې څانګه
ئې بيانوي۔ په ځانګړي ډول، د دې معنٰی دا ده چې تړلے تابلو هغه
ټول امکانات ردوي چې د~\LR{$\sFmla{\True}{!A}$} په بڼه هر فرض په يو وخت
رښتيا او د~\LR{$\sFmla{\False}{!A}$} په بڼه هر فرض دروغ کړي۔
\olpseganchor{OLP-0099-B010}{end}

\olpseganchor{OLP-0099-B011}{start}
\emph{د~\LR{$!A$} دپاره} تړلے تابلو داسې تړلے تابلو دے چې
ريښه ئې~\LR{$\sFmla{\False}{!A}$} وي۔ که داسې تړلے تابلو موجود
وي، د~\LR{$!A$} د دروغ کېدو ټول امکانات رد شوي؛ يعنې \LR{$!A$} بايد په هر
جوړښت کښې رښتيا وي۔
\olpseganchor{OLP-0099-B011}{end}

\olpseganchor{OLP-0100-B004}{start}
\olfileid{fol}{tab}{prl}
\olpseganchor{OLP-0100-B004}{end}

\olpseganchor{OLP-0100-B005}{start}
\olsection{قضيوي قاعدې}
\olpseganchor{OLP-0100-B005}{end}

\olpseganchor{OLP-0100-B006}{start}
\subsection{د~\texorpdfstring{\LR{$\lnot$}}{¬} قاعدې}
\olpseganchor{OLP-0100-B006}{end}

\olpseganchor{OLP-0100-B007}{start}
\begin{defish}
\AxiomC{\sFmla{\True}{\lnot !A}}
\RightLabel{\TRule{\True}{\lnot}}
\UnaryInfC{\sFmla{\False}{!A}}
\DisplayProof
\hfill
\AxiomC{\sFmla{\False}{\lnot !A}}
\RightLabel{\TRule{\False}{\lnot}}
\UnaryInfC{\sFmla{\True}{!A}}
\DisplayProof
\end{defish}
\olpseganchor{OLP-0100-B007}{end}

\olpseganchor{OLP-0100-B008}{start}
\subsection{د~\texorpdfstring{\LR{$\land$}}{∧} قاعدې}
\olpseganchor{OLP-0100-B008}{end}

\olpseganchor{OLP-0100-B009}{start}
\begin{defish}\noindent
\AxiomC{\sFmla{\True}{!A \land !B}}
\RightLabel{\TRule{\True}{\land}}
\UnaryInfC{\sFmla{\True}{!A}}
\noLine
\UnaryInfC{\sFmla{\True}{!B}}
\DisplayProof
\hfill
\AxiomC{\sFmla{\False}{!A \land !B}}
\RightLabel{\TRule{\False}{\land}}
\UnaryInfC{\LR{$\sFmla{\False}{!A} \quad \mid \quad \sFmla{\False}{!B}$}}
\DisplayProof
\end{defish}
\olpseganchor{OLP-0100-B009}{end}

\olpseganchor{OLP-0100-B010}{start}
\subsection{د~\texorpdfstring{\LR{$\lor$}}{∨} قاعدې}
\olpseganchor{OLP-0100-B010}{end}

\olpseganchor{OLP-0100-B011}{start}
\begin{defish}
\AxiomC{\sFmla{\True}{!A \lor !B}}
\RightLabel{\TRule{\True}{\lor}}
\UnaryInfC{\LR{$\sFmla{\True}{!A} \quad \mid \quad \sFmla{\True}{!B}$}}
\DisplayProof
\hfill
\AxiomC{\sFmla{\False}{!A \lor !B}}
\RightLabel{\TRule{\False}{\lor}}
\UnaryInfC{\sFmla{\False}{!A}}
\noLine
\UnaryInfC{\sFmla{\False}{!B}}
\DisplayProof
\end{defish}
\olpseganchor{OLP-0100-B011}{end}

\olpseganchor{OLP-0100-B012}{start}
\subsection{د~\texorpdfstring{\LR{$\lif$}}{→} قاعدې}
\olpseganchor{OLP-0100-B012}{end}

\olpseganchor{OLP-0100-B013}{start}
\begin{defish}
\AxiomC{\sFmla{\True}{!A \lif !B}}
\RightLabel{\TRule{\True}{\lif}}
\UnaryInfC{\LR{$\sFmla{\False}{!A} \quad \mid \quad \sFmla{\True}{!B}$}}
\DisplayProof
\hfill
\AxiomC{\sFmla{\False}{!A \lif !B}}
\RightLabel{\TRule{\False}{\lif}}
\UnaryInfC{\sFmla{\True}{!A}}
\noLine
\UnaryInfC{\sFmla{\False}{!B}}
\DisplayProof
\end{defish}
\olpseganchor{OLP-0100-B013}{end}

\olpseganchor{OLP-0100-B014}{start}
\subsection{د پرېکون قاعده}
\olpseganchor{OLP-0100-B014}{end}

\olpseganchor{OLP-0100-B015}{start}
\begin{defish}
\AxiomC{}
\RightLabel{\Cut}
\UnaryInfC{\LR{$\sFmla{\True}{!A} \quad \mid \quad \sFmla{\False}{!A}$}}
\DisplayProof
\end{defish}
\olpseganchor{OLP-0100-B015}{end}

\olpseganchor{OLP-0100-B016}{start}
د~\Cut{} قاعده پر مخکښ نښه لرونکے فارمول نۀ «پلې کېږي»؛ بلکې
راکوي چې په تابلو کښې هره څانګه په دوو ووېشو: يوه څانګه
\LR{$\sFmla{\True}{!A}$} او بله~\LR{$\sFmla{\False}{!A}$} لري۔ دا اړينه نۀ
ده---د نښه لرونکي فارمولونه هر هغه سټ چې تړلے تابلو ولري،
له~\Cut پرته هم يو تړلے تابلو لري---خو راکوي چې تابلوګانې په
اسانه طريقه سره يو ځاے کړو۔
\olpseganchor{OLP-0100-B016}{end}

\olpseganchor{OLP-0101-B004}{start}
\olfileid{fol}{tab}{qrl}
\olpseganchor{OLP-0101-B004}{end}

\olpseganchor{OLP-0101-B005}{start}
\olsection{د سورونو قاعدې}
\olpseganchor{OLP-0101-B005}{end}

\olpseganchor{OLP-0101-B006}{start}
\subsection{د~\texorpdfstring{\LR{$\lforall$}}{∀} قاعدې}
\olpseganchor{OLP-0101-B006}{end}

\olpseganchor{OLP-0101-B007}{start}
\begin{defish}
\AxiomC{\sFmla{\True}{\lforall[x][!A(x)]}}
\RightLabel{\TRule{\True}{\forall}}
\UnaryInfC{\sFmla{\True}{!A(t)}}
\DisplayProof
\hfill
\AxiomC{\sFmla{\False}{\lforall[x][!A(x)]}}
\RightLabel{\TRule{\False}{\lforall}}
\UnaryInfC{\sFmla{\False}{!A(a)}}
\DisplayProof
\end{defish}
\olpseganchor{OLP-0101-B007}{end}

\olpseganchor{OLP-0101-B008}{start}
په~\TRule{\True}{\lforall} کښې \LR{$t$} تړلې اصطلاح ده (يعنې متغيرونه
نۀ لري)۔ په~\TRule{\False}{\lforall} کښې \LR{$a$} داسې ثابت سمبول
دے چې بايد د~\TRule{\False}{\lforall} تر قاعدې پاس په څانګه کښې
هېڅ ځاے نۀ وي راغلے۔ \LR{$a$} د~\TRule{\False}{\forall} استنتاج
\emph{ځانګړے متغير} بولو۔\footnote{د «ځانګړي متغير» اصطلاح کاروو،
سره له دې چې په پورته قاعده کښې \LR{$a$} يو ثابت سمبول دے۔ دا
تاريخي لاملونه لري۔}
\olpseganchor{OLP-0101-B008}{end}

\olpseganchor{OLP-0101-B009}{start}
\subsection{د~\texorpdfstring{\LR{$\lexists$}}{∃} قاعدې}
\olpseganchor{OLP-0101-B009}{end}

\olpseganchor{OLP-0101-B010}{start}
\begin{defish}
\AxiomC{\sFmla{\True}{\lexists[x][!A(x)]}}
\RightLabel{\TRule{\True}{\lexists}}
\UnaryInfC{\sFmla{\True}{!A(a)}}
\DisplayProof
\hfill
\AxiomC{\sFmla{\False}{\lexists[x][!A(x)]}}
\RightLabel{\TRule{\False}{\lexists}}
\UnaryInfC{\sFmla{\False}{!A(t)}}
\DisplayProof
\end{defish}
\olpseganchor{OLP-0101-B010}{end}

\olpseganchor{OLP-0101-B011}{start}
بيا هم \LR{$t$} تړلې اصطلاح ده، او~\LR{$a$} داسې ثابت سمبول دے چې د
\TRule{\True}{\lexists} تر قاعدې پاس په څانګه کښې نۀ راځي۔
\LR{$a$} د~\TRule{\True}{\lexists} استنتاج \emph{ځانګړے متغير} بولو۔
\olpseganchor{OLP-0101-B011}{end}

\olpseganchor{OLP-0101-B012}{start}
هغه شرط چې ځانګړے متغير د~\TRule{\False}{\lforall} يا
\TRule{\True}{\lexists} استنتاج تر پاسه په څانګه کښې نۀ راځي،
\emph{د ځانګړي متغير شرط} بلل کېږي۔
\olpseganchor{OLP-0101-B012}{end}

\olpseganchor{OLP-0101-B013}{start}
\begin{explain}
هغه قرارداد را ياد کړئ چې کله \LR{$!A$} داسې فارمول وي چې
متغير~\LR{$x$} پکې ازاد وي، دا د~\LR{$!A(x)$} په ليکلو ښيو۔ په همدې
سياق کښې \LR{$!A(t)$} د~\LR{$\Subst{!A}{t}{x}$} لنډه بڼه ده۔ نو د
\TRule{\False}{\lexists} قاعده داسې هم ليکلے شو:
\begin{prooftree}
  \AxiomC{\sFmla{\False}{\lexists[x][!A]}}
  \RightLabel{\TRule{\False}{\lexists}}
  \UnaryInfC{\sFmla{\False}{\Subst{!A}{t}{x}}}
\end{prooftree}
پام وکړئ چې \LR{$t$} له مخکښې په~\LR{$!A$} کښې راتلے شي؛ مثلاً \LR{$!A$} کېدے
شي \LR{$\Atom{\Obj P}{t,x}$} وي۔ نو له
\LR{$ \sFmla{\False}{\lexists[x][\Atom{\Obj P}{t,x}]}$} څخه
\LR{$\sFmla{\False}{\Atom{\Obj P}{t,t}}$} استنتاجول د
\TRule{\False}{\lexists} سمه کارونه ده۔ خو په
\TRule{\False}{\lforall} او~\TRule{\True}{\lexists} کښې د ځانګړي
متغير شرطونه غواړي چې ثابت سمبول~\LR{$a$} په~\LR{$!A$} کښې نۀ راځي۔ نو له
\LR{$\sFmla{\False}{\lforall[x][\Atom{\Obj P}{a,x}]}$}
څخه د~\LR{$\TRule{\False}{\lforall}$} په کارولو
\LR{$\sFmla{\False}{\Atom{\Obj P}{a,a}}$} په سمه توګه نۀ شئ استنتاجولے۔
\end{explain}
\olpseganchor{OLP-0101-B013}{end}

\olpseganchor{OLP-0101-B014}{start}
\begin{explain}
په~\TRule{\True}{\lforall} او~\TRule{\False}{\lexists} کښې د
تړلې اصطلاح~\LR{$t$} پر انتخاب نور بنديز نشته۔ بل لوري ته، د
\TRule{\True}{\lexists} او~\TRule{\False}{\lforall} په قاعدو کښې
د ځانګړي متغير شرط غواړي چې ثابت سمبول~\LR{$a$} د اړوند استنتاج تر
پاسه په څانګو کښې هېڅ ځاے نۀ راځي۔ دا شرط د نظام د سموالي د
يقيني کولو دپاره اړين دے۔ له دې شرط پرته لاندې به د
\LR{$\lexists[x][\formula{A}(x)] \lif \lforall[x][\formula{A}(x)]$}
دپاره يو تړلے تابلو و:
\begin{center}
\begin{tableau}{}
  [\sFmla{\False}{\lexists[x][\formula{A}(x)] \lif \lforall[x][\formula{A}(x)]}, just=\TAss
    [\sFmla{\True}{\lexists[x][\formula{A}(x)]},
      just={\TRule{\False}{\lif}[1]}
      [\sFmla{\False}{\lforall[x][\formula{A}(x)]},
        just={\TRule{\False}{\lif}[1]}
        [\sFmla{\True}{\formula{A}(a)},
          just={\TRule{\True}{\lexists}[2]}
          [\sFmla{\False}{\formula{A}(a)},
            just={\TRule{\False}{\lforall}[3]}, close]
        ]
      ]
    ]
  ]
\end{tableau}
\end{center}
خو \LR{$\lexists[x][\formula{A}(x)] \lif
\lforall[x][\formula{A}(x)]$} معتبر نۀ دے۔
\end{explain}
\olpseganchor{OLP-0101-B014}{end}

\olpseganchor{OLP-0102-B004}{start}
\olfileid{fol}{tab}{der}
\olpseganchor{OLP-0102-B004}{end}

\olpseganchor{OLP-0102-B005}{start}
\olsection{تابلوګانې}
\olpseganchor{OLP-0102-B005}{end}

\olpseganchor{OLP-0102-B006}{start}
\begin{explain}
و مو ويل چې فرض څه دے، او د استنتاج قاعدې مو ورکړې۔ له دغو څخه
تابلوګانې په استقرايي ډول جوړېږي: هر تابلو يا يوه يوازېنۍ
څانګه ده چې له يوه يا زياتو فرضونو جوړه وي، يا پر يوه څانګه د
استنتاج د کومې قاعدې په پلي کولو له بل تابلو څخه اخستل کېږي۔
\end{explain}
\olpseganchor{OLP-0102-B006}{end}

\olpseganchor{OLP-0102-B007}{start}
\begin{defn}[تابلو]
د فرضونو~\LR{$\sFmla{S_1}{!A_1}$}، \dots،
\LR{$\sFmla{S_n}{!A_n}$} دپاره (چې هر~\LR{$S_i$} يا~\LR{$\True$} وي
يا~\LR{$\False$}) تابلو د~نښه لرونکي فارمولونه يوه متناهي ونه
ده چې لاندې شرطونه پوره کوي:
\begin{enumerate}
\item د ونې تر ټولو پاسني \LR{$n$}~نښه لرونکي فارمولونه
  \LR{$\sFmla{S_i}{!A_i}$} دي، يو تر بل لاندې۔
\item په ونه کښې هر نښه لرونکے فارمول چې له فرضونو څخه نۀ وي، تر
  ځان پاس په څانګه کښې پر نښه لرونکے فارمول د يوې استنتاجي قاعدې
  له سمې کارونې څخه راوتلے وي۔
\end{enumerate}
د~تابلو يوه څانګه هله او يوازې هله \emph{تړلې} ده چې
\LR{$\sFmla{\True}{!A}$} او~\LR{$\sFmla{\False}{!A}$} دواړه ولري؛ که نۀ وي
\emph{پرانيستې} ده۔ هغه تابلو چې هره څانګه ئې تړلې وي، د
خپلو فرضونو د سټ \emph{تړلے تابلو} دے۔ که تابلو تړلے
نۀ وي، يعنې لږ تر لږه يوه پرانيستې څانګه ولري، نو \emph{پرانيستے}
دے۔
\end{defn}
\olpseganchor{OLP-0102-B007}{end}

\olpseganchor{OLP-0102-B008}{start}
\begin{ex}
  د فرضونو هر سټ په يوازې ځان تابلو دے، خو په عام ډول به
  تړلے نۀ وي۔ (په څرګند ډول يوازې هله تړلے دے چې فرضونه له مخکښې
  د~نښه لرونکي فارمولونه يوه جوړه
  \LR{$\sFmla{\True}{!A}$} او~\LR{$\sFmla{\False}{!A}$} ولري۔)
\olpseganchor{OLP-0102-B008}{end}

\olpseganchor{OLP-0102-B009}{start}
  له يوه تابلو (پرانيستي يا تړلي) څخه د کومې استنتاجي قاعدې
  په پلي کولو يو نوے، لوے تابلو اخستلے شو۔ قاعده په هغه کښې پر
  نښه لرونکے فارمول~\LR{$!A$} پلې کوو۔ دا قاعده به د هرې هغې څانګې په
  پاے کښې يو يا زيات نښه لرونکي فارمولونه ورزيات کړي چې د~\LR{$!A$} هغه
  پېښه لري چې قاعده مو پرې پلې کړې ده۔
\olpseganchor{OLP-0102-B009}{end}

\olpseganchor{OLP-0102-B010}{start}
  مثلاً فرض~\LR{$\sFmla{\True}{!A \land \lnot !A}$} ونيسئ۔ لاندې هغه
  (پرانيستے) تابلو دے چې يوازې له همدې فرض څخه جوړ دے:
  \begin{center}
    \begin{tableau}{}
      [\sFmla{\True}{\formula{A} \land \lnot \formula{A}}, just=\TAss]
    \end{tableau}{}
  \end{center}
  د~\LR{$\TRule{\True}{\land}$} قاعده پر فرض پلې کوو او ترې نوے
  تابلو اخلو۔ دا قاعده راکوي چې تابلو ته دوه نوې کرښې،
  \LR{$\sFmla{\True}{!A}$} او~\LR{$\sFmla{\True}{\lnot !A}$}، ورزياتې کړو:
  \begin{center}
    \begin{tableau}{}
      [\sFmla{\True}{\formula{A} \land \lnot \formula{A}}, just=\TAss
        [\sFmla{\True}{\formula{A}}, just={\TRule{\True}{\land}[1]},
          [\sFmla{\True}{\lnot \formula{A}}, just={\TRule{\True}{\land}[1]}
          ]
        ]
      ]
    \end{tableau}{}
  \end{center}
  کله چې تابلوګانې ليکو، پلې شوې قاعدې په ښي لوري کښې ثبتوو
  (مثلاً \LR{$\TRule{\True}{\land} 1$} دا معنٰی لري چې پر دې کرښه
  نښه لرونکے فارمول د~\LR{$\TRule{\True}{\land}$} قاعدې د هغې کارونې
  پايله ده چې د~\LR{$1$} پر کرښه نښه لرونکے فارمول ته شوې ده)۔ دا نوے
  تابلو اوس نور نښه لرونکي فارمولونه لري، خو يوازې پر يوه
  (\LR{$\sFmla{\True}{\lnot !A}$}) قاعده پلې کولے شو (په دې حالت کښې
  \LR{$\TRule{\True}{\lnot}$} قاعده)۔ له دې څخه لاندې تړلے
  تابلو راځي:
  \begin{center}
    \begin{tableau}{}
      [\sFmla{\True}{\formula{A} \land \lnot \formula{A}}, just=\TAss
        [\sFmla{\True}{\formula{A}}, just={\TRule{\True}{\land}[1]}
          [\sFmla{\True}{\lnot \formula{A}}, just={\TRule{\True}{\land}[1]}
            [\sFmla{\False}{\formula{A}}, just={\TRule{\True}{\lnot}[3]}, close]
          ]
        ]
      ]
    \end{tableau}{}
  \end{center}
\end{ex}
\olpseganchor{OLP-0102-B010}{end}

\olpseganchor{OLP-0103-B004}{start}
\olfileid{fol}{tab}{pro}
\olpseganchor{OLP-0103-B004}{end}

\olpseganchor{OLP-0103-B005}{start}
\olsection{د~تابلوګانې بېلګې}
\olpseganchor{OLP-0103-B005}{end}

\olpseganchor{OLP-0103-B006}{start}
\begin{ex}
راځئ د~تړلے فارمول~\LR{$(!A \land !B) \lif !A$} دپاره يو تړلے
تابلو ومومو۔
\olpseganchor{OLP-0103-B006}{end}

\olpseganchor{OLP-0103-B007}{start}
لومړے اړوند فرض د~تابلو په سر کښې ليکو۔
\begin{oltableau}
  [\sFmla{\False}{(\formula{A} \land \formula{B}) \lif \formula{A}},
    just = \TAss]
\end{oltableau}
\olpseganchor{OLP-0103-B007}{end}

\olpseganchor{OLP-0103-B008}{start}
يوازې يو فرض شته، نو يوازې يو نښه لرونکے فارمول دے چې قاعده پرې
پلې کولے شو۔ (د هر نښه لرونکے فارمول دپاره تل تر ټولو زياته يوه
قاعده پلې کېدے شي: دا د~تړلے فارمول د اړوندې نښې او
اصلي عملګر قاعده ده۔) په دې حالت کښې بايد
\LR{$\TRule{\False}{\lif}$} پلې کړو۔
\begin{oltableau}
  [\sFmla{\False}{(\formula{A} \land \formula{B}) \lif \formula{A}},
    checked, just = \TAss
    [\sFmla{\True}{\formula{A} \land \formula{B}},
      just={\TRule{\False}{\lif}[1]}
      [\sFmla{\False}{\formula{A}}, just={\TRule{\False}{\lif}[1]}]
    ]
  ]
\end{oltableau}
د دې حساب ساتلو دپاره چې پر کومو نښه لرونکي فارمولونه مو اړوندې
قاعدې پلې کړې دي، د جملې تر څنګ د سم نښان ليکو۔ خو د سم نښان
\emph{يوازې} هله وليکئ چې قاعده پر ټولو پرانيستو څانګو پلې شوې
وي۔ کله چې پر کوم نښه لرونکے فارمول په هره پرانيستې څانګه کښې
اړونده قاعده پلې شي، بيا ورته د بېرته ورګرځېدو او د قاعدې د بيا
پلي کولو اړتيا نۀ لرو۔ په دې حالت کښې يوازې يوه څانګه ده، نو قاعده
يوازې يو ځل پلې کول غواړي۔ (پام وکړئ چې د سم نښانونه يوازې د
تابلوګانو د جوړولو اسانتيا ده او په رسمي ډول د تابلوګانو د نحوې
برخه نۀ ده۔)
\olpseganchor{OLP-0103-B008}{end}

\olpseganchor{OLP-0103-B009}{start}
يو نوے نښه لرونکے فارمول شته چې قاعده پرې پلې کولے شو: پر
\LR{$2$}~کرښه \LR{$\sFmla{\True}{!A \land !B}$}۔ د
\LR{$\TRule{\True}{\land}$} قاعدې له پلي کولو څخه دا راځي:
\begin{oltableau}
  [\sFmla{\False}{(\formula{A} \land \formula{B}) \lif \formula{A}},
    checked, just = \TAss
    [\sFmla{\True}{\formula{A} \land \formula{B}},
      just={\TRule{\False}{\lif}[1]}, checked
      [\sFmla{\False}{\formula{A}}, just={\TRule{\False}{\lif}[1]}
        [\sFmla{\True}{\formula{A}}, just={\TRule{\True}{\land}[2]}
          [\sFmla{\True}{\formula{B}}, just={\TRule{\True}{\land}[2]}, close
          ]
        ]
      ]
    ]
  ]
\end{oltableau}
ځکه اوس څانګه \LR{$\sFmla{\True}{!A}$} (پر~\LR{$4$} کرښه) او
\LR{$\sFmla{\False}{!A}$} (پر~\LR{$3$} کرښه) دواړه لري، څانګه تړلې ده۔
ځکه همدا يوازېنۍ څانګه ده، تابلو تړلے دے۔ د
\LR{$(!A \land !B) \lif !A$} دپاره مو يو تړلے تابلو وموند۔
\end{ex}
\olpseganchor{OLP-0103-B009}{end}

\olpseganchor{OLP-0103-B010}{start}
\begin{ex}
اوس راځئ د~\LR{$(\lnot !A \lor !B) \lif (!A
\lif !B)$} دپاره يو تړلے تابلو ومومو۔
\olpseganchor{OLP-0103-B010}{end}

\olpseganchor{OLP-0103-B011}{start}
له اړوند فرض څخه پيل کوو:
\begin{oltableau}
  [\sFmla{\False}{(\lnot \formula{A} \lor \formula{B}) \lif
      (\formula{A} \lif \formula{B})}, just=\TAss]
\end{oltableau}
په دې تابلو کښې د يوازېني نښه لرونکے فارمول
اصلي عملګر~\LR{$\lif$} او نښه~\LR{$\False$} ده، نو د
\LR{$\TRule{\False}{\lif}$} قاعده پرې پلې کوو او دا اخلو:
\begin{oltableau}
  [\sFmla{\False}{(\lnot \formula{A} \lor \formula{B})
      \lif (\formula{A} \lif \formula{B})}, just=\TAss, checked
    [\sFmla{\True}{\lnot \formula{A} \lor \formula{B}},
      just={\TRule{\False}{\lif}[1]}
      [\sFmla{\False}{(\formula{A} \lif \formula{B})},
        just={\TRule{\False}{\lif}[1]}
      ]
    ]
  ]
\end{oltableau}
اوس دا ټاکنه لرو چې پر~\LR{$2$} کرښه~\LR{$\TRule{\True}{\lor}$} پلې کړو
که پر~\LR{$3$} کرښه \LR{$\TRule{\False}{\lif}$}۔ په حقيقت کښې د ترتيب
ټاکنه توپير نۀ کوي، تر هغې چې د هر نښه لرونکے فارمول اړونده قاعده
په هره څانګه کښې پلې شي۔ نو لومړۍ اخلو۔ د
\LR{$\TRule{\True}{\lor}$} قاعده تابلو ته څانګې ورکوي، او د
قاعدې دوه پايلې به هغه نوي نښه لرونکي فارمولونه وي چې دوو نويو
څانګو ته ورزياتېږي۔ له دې څخه دا راځي:
\begin{oltableau}
  [\sFmla{\False}{(\lnot \formula{A} \lor \formula{B}) \lif
      (\formula{A} \lif \formula{B})}, just=\TAss, checked
    [\sFmla{\True}{\lnot \formula{A} \lor \formula{B}},
      just={\TRule{\False}{\lif}[1]}, checked
      [\sFmla{\False}{(\formula{A} \lif \formula{B})},
        just={\TRule{\False}{\lif}[1]}
        [\sFmla{\True}{\lnot \formula{A}}, just={\TRule{\True}{\lor}[2]}]
        [\sFmla{\True}{\formula{B}}, just={\TRule{\True}{\lor}[2]}]
      ]
    ]
  ]
\end{oltableau}
لا مو پر~\LR{$3$} کرښه د~\LR{$\TRule{\False}{\lif}$} قاعده نۀ ده پلې
کړې؛ اوس ئې کوو۔ د وخت د سپما دپاره ئې پر دواړو څانګو پلې کوو۔
را ياد کړئ چې د کوم نښه لرونکے فارمول تر څنګ يوازې هله د سم نښان
ليکو چې اړونده قاعده مو په هره پرانيستې څانګه کښې پلې کړې وي۔ نو
غوره ده چې قاعده د هرې هغې څانګې په پاے کښې پلې کړو چې اړوند
نښه لرونکے فارمول لري۔ په دې توګه به په بېلابېلو څانګو کښې لاندې
هماغه نښه لرونکے فارمول ته بېرته نۀ ورګرځو۔
\begin{oltableau}
  [\sFmla{\False}{(\lnot \formula{A} \lor \formula{B}) \lif
      (\formula{A} \lif \formula{B})}, just=\TAss, checked
    [\sFmla{\True}{\lnot \formula{A} \lor \formula{B}},
      just={\TRule{\False}{\lif}[1]}, checked
      [\sFmla{\False}{(\formula{A} \lif \formula{B})},
        just={\TRule{\False}{\lif}[1]}, checked
        [\sFmla{\True}{\lnot \formula{A}}, just={\TRule{\True}{\lor}[2]}
          [\sFmla{\True}{\formula{A}}, just={\TRule{\False}{\lif}[3]}
            [\sFmla{\False}{\formula{B}}, just={\TRule{\False}{\lif}[3]}]
          ]
        ]
        [\sFmla{\True}{\formula{B}}, just={\TRule{\True}{\lor}[2]}
          [\sFmla{\True}{\formula{A}}, just={\TRule{\False}{\lif}[3]}
            [\sFmla{\False}{\formula{B}}, just={\TRule{\False}{\lif}[3]}, close]
          ]
        ]
      ]
    ]
  ]
\end{oltableau}
اوس ښۍ څانګه تړلې ده۔ په کيڼه څانګه کښې لا هم پر~\LR{$4$} کرښه د
\LR{$\TRule{\True}{\lnot}$} قاعده پلې کولے شو۔ له دې څخه
\LR{$\sFmla{\False}{!A}$} راځي او کيڼه څانګه تړل کېږي:
\begin{oltableau}
  [\sFmla{\False}{(\lnot \formula{A} \lor \formula{B}) \lif
      (\formula{A} \lif \formula{B})}, just=\TAss, checked
    [\sFmla{\True}{\lnot \formula{A} \lor \formula{B}},
      just={\TRule{\False}{\lif}[1]}, checked
      [\sFmla{\False}{(\formula{A} \lif \formula{B})},
        just={\TRule{\False}{\lif}[1]}, checked
        [\sFmla{\True}{\lnot \formula{A}}, just={\TRule{\True}{\lor}[2]}
          [\sFmla{\True}{\formula{A}}, just={\TRule{\False}{\lif}[3]}
            [\sFmla{\False}{\formula{B}}, just={\TRule{\False}{\lif}[3]}
              [\sFmla{\False}{\formula{A}},
                just={\TRule{\True}{\lnot}[4]}, close]
            ]
          ]
        ]
        [\sFmla{\True}{\formula{B}}, just={\TRule{\True}{\lor}[2]}
          [\sFmla{\True}{\formula{A}}, just={\TRule{\False}{\lif}[3]}
            [\sFmla{\False}{\formula{B}},
              just={\TRule{\False}{\lif}[3]}, close]
          ]
        ]
      ]
    ]
  ]
\end{oltableau}
\end{ex}
\olpseganchor{OLP-0103-B011}{end}

\olpseganchor{OLP-0103-B012}{start}
\begin{ex}
د فرضونو په توګه د هر شمېر نښه لرونکي فارمولونه دپاره
تابلوګانې ورکولے شو۔ ډېر ځله له يوې زياتې داسې قاعدې پلي کول هم
اړين وي چې څانګې ورکوي؛ او په عام ډول تابلو هر شمېر څانګې
لرلے شي۔ مثلاً د
\LR{$\{\sFmla{\True}{!A \lor (!B \land !C)}, \sFmla{\False}{(!A \lor !B) \land (!A
  \lor !C)}\}$} دپاره تابلو ونيسئ۔ لومړے پر لومړي فرض د
\LR{$\TRule{\True}{\lor}$} قاعده پلې کوو:
\begin{oltableau}
  [\sFmla{\True}{\formula{A} \lor (\formula{B} \land \formula{C})},
    just=\TAss, checked
    [\sFmla{\False}{(\formula{A} \lor \formula{B}) \land
        (\formula{A} \lor \formula{C})}, just=\TAss
      [\sFmla{\True}{\formula{A}}, just={\TRule{\True}{\lor}[1]}]
      [\sFmla{\True}{\formula{B} \land \formula{C}},
        just={\TRule{\True}{\lor}[1]}]
    ]
  ]
\end{oltableau}
اوس پر~\LR{$2$} کرښه د~\LR{$\TRule{\False}{\land}$} قاعده پلې کولے
شو۔ دا په دواړو څانګو کښې په يو وخت کوو، نو پر~\LR{$2$} کرښه د سم نښان
وهلے شو:
\begin{oltableau}
  [\sFmla{\True}{\formula{A} \lor (\formula{B} \land \formula{C})},
    just=\TAss, checked
    [\sFmla{\False}{(\formula{A} \lor \formula{B}) \land
        (\formula{A} \lor \formula{C})}, just=\TAss, checked
      [\sFmla{\True}{\formula{A}}, just={\TRule{\True}{\lor}[1]}
        [\sFmla{\False}{\formula{A} \lor \formula{B}},
          just={\TRule{\False}{\land}[2]}]
        [\sFmla{\False}{\formula{A} \lor \formula{C}},
          just={\TRule{\False}{\land}[2]}]
      ]
      [\sFmla{\True}{\formula{B} \land \formula{C}},
        just={\TRule{\True}{\lor}[1]}
        [\sFmla{\False}{\formula{A} \lor \formula{B}},
          just={\TRule{\False}{\land}[2]}]
        [\sFmla{\False}{\formula{A} \lor \formula{C}},
          just={\TRule{\False}{\land}[2]}]
      ]
    ]
  ]
\end{oltableau}
اوس \LR{$\TRule{\False}{\lor}$} پر ټولو هغو څانګو پلې کولے شو چې
\LR{$!A \lor !B$} لري:
\begin{oltableau}
  [\sFmla{\True}{\formula{A} \lor (\formula{B} \land \formula{C})},
    just=\TAss, checked
    [\sFmla{\False}{(\formula{A} \lor \formula{B}) \land
        (\formula{A} \lor \formula{C})}, just=\TAss, checked
      [\sFmla{\True}{\formula{A}}, just={\TRule{\True}{\lor}[1]}
        [\sFmla{\False}{\formula{A} \lor \formula{B}},
          just={\TRule{\False}{\land}[2]}, checked
          [\sFmla{\False}{\formula{A}}, just={\TRule{\False}{\lor}[4]}
            [\sFmla{\False}{\formula{B}},
              just={\TRule{\False}{\lor}[4]}, close]
          ]
        ]
        [\sFmla{\False}{\formula{A} \lor \formula{C}},
          just={\TRule{\False}{\land}[2]}]
      ]
      [\sFmla{\True}{\formula{B} \land \formula{C}},
        just={\TRule{\True}{\lor}[1]}
        [\sFmla{\False}{\formula{A} \lor \formula{B}},
          just={\TRule{\False}{\land}[2]}, checked
          [\sFmla{\False}{\formula{A}}, just={\TRule{\False}{\lor}[4]}
            [\sFmla{\False}{\formula{B}},
              just={\TRule{\False}{\lor}[4]}]
          ]
        ]
        [\sFmla{\False}{\formula{A} \lor \formula{C}},
          just={\TRule{\False}{\land}[2]}]
      ]
    ]
  ]
\end{oltableau}
اوس تر ټولو کيڼه څانګه تړلې ده۔ راځئ
\LR{$\TRule{\False}{\lor}$} پر~\LR{$!A \lor !C$} پلې کړو:
\begin{oltableau}
  [\sFmla{\True}{\formula{A} \lor (\formula{B} \land \formula{C})},
    just=\TAss, checked
    [\sFmla{\False}{(\formula{A} \lor \formula{B}) \land
        (\formula{A} \lor \formula{C})}, just=\TAss, checked
      [\sFmla{\True}{\formula{A}}, just={\TRule{\True}{\lor}[1]}
        [\sFmla{\False}{\formula{A} \lor \formula{B}},
          just={\TRule{\False}{\land}[2]}, checked
          [\sFmla{\False}{\formula{A}},
            just={\TRule{\False}{\lor}[4]}
            [\sFmla{\False}{\formula{B}},
              just={\TRule{\False}{\lor}[4]}, close]
          ]
        ]
        [\sFmla{\False}{\formula{A} \lor \formula{C}},
          just={\TRule{\False}{\land}[2]}, checked
          [\sFmla{\False}{\formula{A}},
            just={\TRule{\False}{\lor}[4]}, move by=2
            [\sFmla{\False}{\formula{C}},
              just={\TRule{\False}{\lor}[4]}, close]
          ]
        ]
      ]
      [\sFmla{\True}{\formula{B} \land \formula{C}},
        just={\TRule{\True}{\lor}[1]}
        [\sFmla{\False}{\formula{A} \lor \formula{B}},
          just={\TRule{\False}{\land}[2]}, checked
          [\sFmla{\False}{\formula{A}},
            just={\TRule{\False}{\lor}[4]}
            [\sFmla{\False}{\formula{B}},
              just={\TRule{\False}{\lor}[4]}
            ]
          ]
        ]
        [\sFmla{\False}{\formula{A} \lor \formula{C}},
          just={\TRule{\False}{\land}[2]}, checked
          [\sFmla{\False}{\formula{A}},
            just={\TRule{\False}{\lor}[4]}, move by=2
            [\sFmla{\False}{\formula{C}},
              just={\TRule{\False}{\lor}[4]}]
          ]
        ]
      ]
    ]
  ]
\end{oltableau}
پام وکړئ چې د روښانتيا دپاره مو د
\LR{$\TRule{\False}{\lor}$} د دويم ځل کارونې پايله لاندې يوړه۔ په دې
حالت کښې دې ته اړتيا نۀ وه؛ ځکه توجيهونه به يو شان وو۔
\olpseganchor{OLP-0103-B012}{end}

\olpseganchor{OLP-0103-B013}{start}
دوه څانګې پرانيستې پاتې دي، او پر~\LR{$3$} کرښه
\LR{$\sFmla{\True}{!B \land !C}$} لا بې د سم نښانه دے۔ پر هغې
\LR{$\TRule{\True}{\land}$} پلې کوو او تړلے تابلو اخلو:
\begin{oltableau}
  [\sFmla{\True}{\formula{A} \lor (\formula{B} \land \formula{C})},
    just=\TAss, checked
    [\sFmla{\False}{(\formula{A} \lor \formula{B}) \land
        (\formula{A} \lor \formula{C})}, just=\TAss, checked
      [\sFmla{\True}{\formula{A}}, just={\TRule{\True}{\lor}[1]}
        [\sFmla{\False}{\formula{A} \lor \formula{B}},
          just={\TRule{\False}{\land}[2]}, checked
          [\sFmla{\False}{\formula{A}},
            just={\TRule{\False}{\lor}[4]}
            [\sFmla{\False}{\formula{B}},
              just={\TRule{\False}{\lor}[4]}, close]
          ]
        ]
        [\sFmla{\False}{\formula{A} \lor \formula{C}},
          just={\TRule{\False}{\land}[2]}, checked
          [\sFmla{\False}{\formula{A}},
            just={\TRule{\False}{\lor}[4]}
            [\sFmla{\False}{\formula{C}},
              just={\TRule{\False}{\lor}[4]}, close]
          ]
        ]
      ]
      [\sFmla{\True}{\formula{B} \land \formula{C}},
        just={\TRule{\True}{\lor}[1]}, checked
        [\sFmla{\False}{\formula{A} \lor \formula{B}},
          just={\TRule{\False}{\land}[2]}, checked
          [\sFmla{\False}{\formula{A}},
            just={\TRule{\False}{\lor}[4]}
            [\sFmla{\False}{\formula{B}},
              just={\TRule{\False}{\lor}[4]}
              [\sFmla{\True}{\formula{B}},
                just={\TRule{\True}{\land}[3]}
                [\sFmla{\True}{\formula{C}},
                  just={\TRule{\True}{\land}[3]},close
                ]
              ]
            ]
          ]
        ]
        [\sFmla{\False}{\formula{A} \lor \formula{C}},
          just={\TRule{\False}{\land}[2]}, checked
          [\sFmla{\False}{\formula{A}},
            just={\TRule{\False}{\lor}[4]}
            [\sFmla{\False}{\formula{C}},
              just={\TRule{\False}{\lor}[4]}
              [\sFmla{\True}{\formula{B}},
                just={\TRule{\True}{\land}[3]}
                [\sFmla{\True}{\formula{C}},
                  just={\TRule{\True}{\land}[3]},close
                ]
              ]
            ]
          ]
        ]
      ]
    ]
  ]
\end{oltableau}
\olpseganchor{OLP-0103-B013}{end}

\olpseganchor{OLP-0103-B014}{start}
د پرتلې دپاره، لاندې د فرضونو د همدې سټ يو تړلے تابلو دے
چې قاعدې پکې په بل ترتيب پلې شوې دي:
\begin{oltableau}
  [\sFmla{\True}{\formula{A} \lor (\formula{B} \land \formula{C})},
    just=\TAss, checked
    [\sFmla{\False}{(\formula{A} \lor \formula{B}) \land
        (\formula{A} \lor \formula{C})}, just=\TAss, checked
      [\sFmla{\False}{\formula{A} \lor \formula{B}},
        just={\TRule{\False}{\land}[2]}, checked
        [\sFmla{\False}{\formula{A}},
          just={\TRule{\False}{\lor}[3]}
          [\sFmla{\False}{\formula{B}},
            just={\TRule{\False}{\lor}[3]}
            [\sFmla{\True}{\formula{A}},
              just={\TRule{\True}{\lor}[1]},close
            ]
            [\sFmla{\True}{\formula{B} \land \formula{C}},
              just={\TRule{\True}{\lor}[1]}, checked
              [\sFmla{\True}{\formula{B}},
                just={\TRule{\True}{\land}[6]}
                [\sFmla{\True}{\formula{C}},
                  just={\TRule{\True}{\land}[6]},close
                ]
              ]
            ]
          ]
        ]
      ]
      [\sFmla{\False}{\formula{A} \lor \formula{C}},
        just={\TRule{\False}{\land}[2]}, checked
        [\sFmla{\False}{\formula{A}},
          just={\TRule{\False}{\lor}[3]}
          [\sFmla{\False}{\formula{C}},
            just={\TRule{\False}{\lor}[3]}
            [\sFmla{\True}{\formula{A}},
              just={\TRule{\True}{\lor}[1]},close
            ]
            [\sFmla{\True}{\formula{B} \land \formula{C}},
              just={\TRule{\True}{\lor}[1]}, checked
              [\sFmla{\True}{\formula{B}},
                just={\TRule{\True}{\land}[6]}
                [\sFmla{\True}{\formula{C}},
                  just={\TRule{\True}{\land}[6]},close
                ]
              ]
            ]
          ]
        ]
      ]
    ]
  ]
\end{oltableau}
\end{ex}
\olpseganchor{OLP-0103-B014}{end}

\olpseganchor{OLP-0103-B015}{start}
\begin{prob}
د لاندې دپاره تړلي تابلوګانې ورکړئ:
\begin{enumerate}
\item \LR{$\sFmla{\True}{!A \land (!B \land !C)}, \sFmla{\False}{(!A \land !B) \land !C}$}.
\item \LR{$\sFmla{\True}{!A \lor (!B \lor !C)}, \sFmla{\False}{(!A \lor !B) \lor !C}$}.
\item \LR{$\sFmla{\True}{!A \lif (!B \lif !C)}, \sFmla{\False}{!B \lif (!A \lif !C)}$}.
\item \LR{$\sFmla{\True}{!A}, \sFmla{\False}{\lnot\lnot !A}$}.
\end{enumerate}
\end{prob}
\olpseganchor{OLP-0103-B015}{end}

\olpseganchor{OLP-0103-B016}{start}
\begin{prob}
د لاندې دپاره تړلي تابلوګانې ورکړئ:
\begin{enumerate}
\item \LR{$\sFmla{\True}{(!A \lor !B) \lif !C}, \sFmla{\False}{!A \lif !C}$}.
\item \LR{$\sFmla{\True}{(!A \lif !C) \land (!B \lif !C)}, \sFmla{\False}{(!A \lor !B) \lif !C}$}.
\item \LR{$\sFmla{\False}{\lnot(!A \land \lnot !A)}$}.
\item \LR{$\sFmla{\True}{!B \lif !A}, \sFmla{\False}{\lnot !A \lif \lnot !B}$}.
\item \LR{$\sFmla{\False}{(!A \lif \lnot !A) \lif \lnot !A}$}.
\item \LR{$\sFmla{\False}{\lnot(!A \lif !B) \lif \lnot !B}$}.
\item \LR{$\sFmla{\True}{!A \lif !C}, \sFmla{\False}{\lnot (!A \land \lnot !C)}$}.
\item \LR{$\sFmla{\True}{!A \land \lnot !C}, \sFmla{\False}{\lnot (!A \lif !C)}$}.
\item \LR{$\sFmla{\True}{!A \lor !B}, \sFmla{\True}{\lnot !B}, \sFmla{\False}{!A}$}.
\item \LR{$\sFmla{\True}{\lnot !A \lor \lnot !B}, \sFmla{\False}{\lnot(!A \land !B)}$}.
\item \LR{$\sFmla{\False}{(\lnot !A \land \lnot !B) \lif\lnot(!A \lor !B)}$}.
\item \LR{$\sFmla{\False}{\lnot(!A \lor !B) \lif (\lnot !A \land \lnot !B)}$}.
\end{enumerate}
\end{prob}
\olpseganchor{OLP-0103-B016}{end}

\olpseganchor{OLP-0103-B017}{start}
\begin{prob}
د لاندې دپاره تړلي تابلوګانې ورکړئ:
\begin{enumerate}
\item \LR{$\sFmla{\True}{\lnot(!A \lif !B)}, \sFmla{\False}{!A}$}.
\item \LR{$\sFmla{\True}{\lnot(!A \land !B)}, \sFmla{\False}{\lnot !A \lor \lnot !B}$}.
\item \LR{$\sFmla{\True}{!A \lif !B}, \sFmla{\False}{\lnot !A \lor !B}$}.
\item \LR{$\sFmla{\False}{\lnot \lnot !A \lif !A}$}.
\item \LR{$\sFmla{\True}{!A \lif !B}, \sFmla{\True}{\lnot !A \lif !B}, \sFmla{\False}{!B}$}.
\item \LR{$\sFmla{\True}{(!A \land !B) \lif !C}, \sFmla{\False}{(!A \lif !C) \lor (!B \lif !C)}$}.
\item \LR{$\sFmla{\True}{(!A \lif !B) \lif !A}, \sFmla{\False}{!A}$}.
\item \LR{$\sFmla{\False}{(!A \lif !B) \lor (!B \lif !C)}$}.
\end{enumerate}
\end{prob}
\olpseganchor{OLP-0103-B017}{end}

\olpseganchor{OLP-0104-B004}{start}
\olfileid{fol}{tab}{prq}
\olpseganchor{OLP-0104-B004}{end}

\olpseganchor{OLP-0104-B005}{start}
\olsection{له سورونو سره تابلوګانې}
\olpseganchor{OLP-0104-B005}{end}

\olpseganchor{OLP-0104-B006}{start}
\begin{ex}
له سورونو سره د کار پر مهال بايد ډاډ واخلو چې د ځانګړي متغير
شرط نۀ ماتوو، او کله کله د دې دپاره د ځينو استنتاجونو د پلي کولو
ترتيب بدلول را بدلول غواړي۔ په عام ډول ګټه کوي چې لومړے د هغو
قاعدو د حل هڅه وکړو چې د ځانګړي متغير شرط ورباندې لګېږي (په بشپړ
تابلو کښې به دوئ پاس راځي)۔
\olpseganchor{OLP-0104-B006}{end}

\olpseganchor{OLP-0104-B007}{start}
راځئ وګورو چې د~تړلے فارمول
\LR{$\lexists[x][\lnot !A(x)] \lif \lnot \lforall[x][!A(x)]$}
دپاره تابلو څنګه ورکوو۔ د معمول په شان د فرض له ثبتولو
پيل کوو:
\begin{oltableau}
[\sFmla{\False}{\lexists[x][\lnot \formula{A}(x)]\lif \lnot
    \lforall[x][\formula{A}(x)]}, just=\TAss]
\end{oltableau}
ځکه اصلي عملګر~\LR{$\lif$} دے، نو
\LR{$\TRule{\False}{\lif}$} پلې کوو:
\begin{oltableau}
[\sFmla{\False}{\lexists[x][\lnot \formula{A}(x)]\lif \lnot
    \lforall[x][\formula{A}(x)]}, just=\TAss, checked
  [\sFmla{\True}{\lexists[x][\lnot \formula{A}(x)]},
    just={\TRule{\False}{\lif}[1]}
    [\sFmla{\False}{\lnot\lforall[x][\formula{A}(x)]},
      just={\TRule{\False}{\lif}[1]}]
  ]
]
\end{oltableau}
راتلونکې د کار کرښه~\LR{$2$} ده۔ \LR{$\TRule{\True}{\lexists}$} کاروو۔
دا يو نوے ثابت سمبول غواړي؛ ځکه لا هېڅ ثابت سمبولونه نۀ دي
راغلي، هر يو ټاکلے شو، مثلاً~\LR{$a$}۔
\begin{oltableau}
[\sFmla{\False}{\lexists[x][\lnot \formula{A}(x)]\lif \lnot
    \lforall[x][\formula{A}(x)]}, just=\TAss, checked
  [\sFmla{\True}{\lexists[x][\lnot \formula{A}(x)]},
    just={\TRule{\False}{\lif}[1]}, checked
    [\sFmla{\False}{\lnot\lforall[x][\formula{A}(x)]},
      just={\TRule{\False}{\lif}[1]}
      [\sFmla{\True}{\lnot\formula{A}(a)}, just={\TRule{\True}{\lexists}[2]}]
    ]
  ]
]
\end{oltableau}
اوس پر~\LR{$3$} کرښه \LR{$\TRule{\False}{\lnot}$} پلې کوو:
\begin{oltableau}
[\sFmla{\False}{\lexists[x][\lnot \formula{A}(x)]\lif \lnot
    \lforall[x][\formula{A}(x)]}, just=\TAss, checked
  [\sFmla{\True}{\lexists[x][\lnot \formula{A}(x)]},
    just={\TRule{\False}{\lif}[1]}, checked
    [\sFmla{\False}{\lnot\lforall[x][\formula{A}(x)]},
      just={\TRule{\False}{\lif}[1]}, checked
      [\sFmla{\True}{\lnot\formula{A}(a)}, just={\TRule{\True}{\lexists}[2]}
        [\sFmla{\True}{\lforall[x][\formula{A}(x)]},
          just = {\TRule{\False}{\lnot}[3]}
        ]
      ]
    ]
  ]
]
\end{oltableau}
پر~\LR{$4$} کرښه د~\LR{$\TRule{\True}{\lnot}$} او ورپسې پر~\LR{$5$} کرښه د
\LR{$\TRule{\True}{\lforall}$} په پلي کولو تړلے تابلو اخلو۔
\begin{oltableau}
[\sFmla{\False}{\lexists[x][\lnot \formula{A}(x)]\lif \lnot
    \lforall[x][\formula{A}(x)]}, just=\TAss, checked
  [\sFmla{\True}{\lexists[x][\lnot \formula{A}(x)]},
    just={\TRule{\False}{\lif}[1]}, checked
    [\sFmla{\False}{\lnot\lforall[x][\formula{A}(x)]},
      just={\TRule{\False}{\lif}[1]}, checked
      [\sFmla{\True}{\lnot\formula{A}(a)}, just={\TRule{\True}{\lexists}[2]}
        [\sFmla{\True}{\lforall[x][\formula{A}(x)]},
          just = {\TRule{\False}{\lnot}[3]}
          [\sFmla{\False}{\formula{A}(a)}, just={\TRule{\True}{\lnot}[4]}
            [\sFmla{\True}{\formula{A}(a)},
              just={\TRule{\True}{\lforall}[5]}, close
            ]
          ]
        ]
      ]
    ]
  ]
]
\end{oltableau}
\end{ex}
\olpseganchor{OLP-0104-B007}{end}

\olpseganchor{OLP-0104-B008}{start}
\begin{ex}
راځئ وګورو چې د لاندې سټ دپاره تابلو څنګه ورکوو:
\[
\sFmla{\False}{\lexists[x][!C(x,b)]},
\sFmla{\True}{\lexists[x][(!A(x) \land !B(x))]},
\sFmla{\True}{\lforall[x][(!B(x) \lif !C(x,b))]}.
\]
د معمول په شان له فرضونو پيل کوو:
\begin{oltableau}
  [\sFmla{\False}{\lexists[x][\formula{C}(x,b)]}, just=\TAss
    [\sFmla{\True}{\lexists[x][(\formula{A}(x) \land \formula{B}(x))]},
      just=\TAss
      [\sFmla{\True}{\lforall[x][(\formula{B}(x) \lif \formula{C}(x,b))]},
        just=\TAss]
    ]
  ]
\end{oltableau}
تل بايد لومړے هغه قاعده پلې کړو چې د ځانګړي متغير شرط لري؛ په
دې حالت کښې هغه پر~\LR{$2$} کرښه \LR{$\TRule{\True}{\lexists}$} ده۔ ځکه
فرضونه ثابت سمبول~\LR{$b$} لري، بايد بل يو وکاروو؛ بيا~\LR{$a$} اخلو۔
\begin{oltableau}
  [\sFmla{\False}{\lexists[x][\formula{C}(x,b)]}, just=\TAss
    [\sFmla{\True}{\lexists[x][(\formula{A}(x) \land \formula{B}(x))]},
      just=\TAss, checked
      [\sFmla{\True}{\lforall[x][(\formula{B}(x) \lif \formula{C}(x,b))]},
        just=\TAss
        [\sFmla{\True}{\formula{A}(a) \land \formula{B}(a)},
          just={\TRule{\True}{\lexists}[2]}]
      ]
    ]
  ]
\end{oltableau}
که اوس پر~\LR{$1$} کرښه \LR{$\TRule{\False}{\lexists}$} يا پر~\LR{$3$} کرښه
\LR{$\TRule{\True}{\lforall}$} پلې کړو، بايد وټاکو چې د~\LR{$x$} پر ځاے
کومه اصطلاح~\LR{$t$} کښېږدو۔ ځکه دا قاعدې د ځانګړي متغير شرط نۀ لري،
هره خوښه اصطلاح ټاکلے شو۔ په ځينو حالتونو کښې ښايي قاعده څو ځله
له بېلابېلو~\LR{$t$}s سره پلې کول هم لازم شي۔ خو د عمومي قاعدې په توګه
ګټه کوي چې له هغو اصطلاحګانو څخه يوه واخلو چې له مخکښې په
تابلو کښې راغلې وي---دلته \LR{$a$} او~\LR{$b$}---او په دې حالت کښې
اټکل کولے شو چې له~\LR{$a$} سره د تړلې څانګې لاس ته راتلل ډېر ممکن دي۔
\begin{oltableau}
  [\sFmla{\False}{\lexists[x][\formula{C}(x,b)]}, just=\TAss
    [\sFmla{\True}{\lexists[x][(\formula{A}(x) \land \formula{B}(x))]},
      just=\TAss, checked
      [\sFmla{\True}{\lforall[x][(\formula{B}(x) \lif \formula{C}(x,b))]}, just=\TAss
        [\sFmla{\True}{\formula{A}(a) \land \formula{B}(a)}, just={\TRule{\True}{\lexists}[2]}
          [\sFmla{\False}{\formula{C}(a,b)}, just={\TRule{\False}{\lexists}[1]}
            [\sFmla{\True}{\formula{B}(a) \lif \formula{C}(a,b)},
              just={\TRule{\True}{\lforall}[3]}
            ]
          ]
        ]
      ]
    ]
  ]
\end{oltableau}
پر~\LR{$1$} او~\LR{$3$} کرښو نښه لرونکي فارمولونه ته د سم نښان نۀ ورکوو؛
ځکه ښايي بيا ئې وکاروو۔ اوس پر~\LR{$4$} کرښه
\LR{$\TRule{\True}{\land}$} پلې کړئ:
\begin{oltableau}
  [\sFmla{\False}{\lexists[x][\formula{C}(x,b)]}, just=\TAss
    [\sFmla{\True}{\lexists[x][(\formula{A}(x) \land \formula{B}(x))]},
      just=\TAss, checked
      [\sFmla{\True}{\lforall[x][(\formula{B}(x) \lif \formula{C}(x,b))]}, just=\TAss
        [\sFmla{\True}{\formula{A}(a) \land \formula{B}(a)},
          just={\TRule{\True}{\lexists}[2]}, checked
          [\sFmla{\False}{\formula{C}(a,b)}, just={\TRule{\False}{\lexists}[1]}
            [\sFmla{\True}{\formula{B}(a) \lif \formula{C}(a,b)},
              just={\TRule{\True}{\lforall}[3]}
              [\sFmla{\True}{\formula{A}(a)}, just={\TRule{\True}{\land}[4]}
                [\sFmla{\True}{\formula{B}(a)}, just={\TRule{\True}{\land}[4]}
                ]
              ]
            ]
          ]
        ]
      ]
    ]
  ]
\end{oltableau}
که اوس پر~\LR{$6$} کرښه \LR{$\TRule{\True}{\lif}$} پلې کړو،
تابلو تړل کېږي:
\begin{oltableau}
  [\sFmla{\False}{\lexists[x][\formula{C}(x,b)]}, just=\TAss
    [\sFmla{\True}{\lexists[x][(\formula{A}(x) \land \formula{B}(x))]},
      just=\TAss, checked
      [\sFmla{\True}{\lforall[x][(\formula{B}(x) \lif \formula{C}(x,b))]},
        just=\TAss
        [\sFmla{\True}{\formula{A}(a) \land \formula{B}(a)},
          just={\TRule{\True}{\lexists}[2]}, checked
          [\sFmla{\False}{\formula{C}(a,b)},
            just={\TRule{\False}{\lexists}[1]}
            [\sFmla{\True}{\formula{B}(a) \lif \formula{C}(a,b)},
              just={\TRule{\True}{\lforall}[3]}, checked
              [\sFmla{\True}{\formula{A}(a)},
                just={\TRule{\True}{\land}[4]}
                [\sFmla{\True}{\formula{B}(a)},
                  just={\TRule{\True}{\land}[4]}
                  [\sFmla{\False}{\formula{B}(a)},
                    just={\TRule{\True}{\lif}[6]},close
                  ]
                  [\sFmla{\True}{\formula{C}(a,b)},
                    just={\TRule{\True}{\lif}[6]},close
                  ]
                ]
              ]
            ]
          ]
        ]
      ]
    ]
  ]
\end{oltableau}
\olpseganchor{OLP-0104-B008}{end}


\olpseganchor{OLP-0104-B009}{start}
\end{ex}
\olpseganchor{OLP-0104-B009}{end}

\olpseganchor{OLP-0104-B010}{start}
\begin{ex}
د لاندې سټ دپاره تابلو جوړوو:
\[
\sFmla{\True}{\lforall[x][!A(x)]}, \sFmla{\True}{\lforall[x][!A(x)] 
  \lif \lexists[y][!B(y)]}, \sFmla{\True}{\lnot\lexists[y][!B(y)]}.
\]
د معمول په شان فرضونه ليکو:
\begin{oltableau}
  [\sFmla{\True}{\lforall[x][\formula{A}(x)]}, just=\TAss
    [\sFmla{\True}{\lforall[x][\formula{A}(x)] \lif
        \lexists[y][\formula{B}(y)]}, just=\TAss
      [\sFmla{\True}{\lnot\lexists[y][\formula{B}(y)]}, just=\TAss
      ]
    ]
  ]
\end{oltableau}
لومړے پر~\LR{$3$} کرښه د~\LR{$\TRule{\True}{\lnot}$} قاعده پلې کوو۔
د «تل لومړے د ځانګړي متغير شرط لرونکې قاعدې پلې کړئ» قاعدې يوه
پايله دا ده: «بې د ځانګړي متغير شرطه د سورونو قاعدې تر اړتيا پورې
وځنډوئ۔» همدارنګه هغه قاعدې هم وځنډوئ چې وېش رامنځته کوي۔
\begin{oltableau}
  [\sFmla{\True}{\lforall[x][\formula{A}(x)]}, just=\TAss
    [\sFmla{\True}{\lforall[x][\formula{A}(x)] \lif
        \lexists[y][\formula{B}(y)]}, just=\TAss
      [\sFmla{\True}{\lnot\lexists[y][\formula{B}(y)]}, just=\TAss, checked
        [\sFmla{\False}{\lexists[y][\formula{B}(y)]}, just={\TRule{\True}{\lnot}[3]}]
      ]
    ]
  ]
\end{oltableau}
نوې \LR{$4$}~کرښه \LR{$\TRule{\False}{\lexists}$} غواړي، چې د ځانګړي
متغير له شرط پرته د سور قاعده ده۔ نو دا ځنډوو او پر~\LR{$2$} کرښه
\LR{$\TRule{\True}{\lif}$} کاروو۔
\begin{oltableau}
  [\sFmla{\True}{\lforall[x][\formula{A}(x)]}, just=\TAss
    [\sFmla{\True}{\lforall[x][\formula{A}(x)] \lif
        \lexists[y][\formula{B}(y)]}, just=\TAss, checked
      [\sFmla{\True}{\lnot\lexists[y][\formula{B}(y)]}, just=\TAss, checked
        [\sFmla{\False}{\lexists[y][\formula{B}(y)]},
          just={\TRule{\True}{\lnot}[3]},
          [\sFmla{\False}{\lforall[x][\formula{A}(x)]}, just={\TRule{\True}{\lif}[2]}]
          [\sFmla{\True}{\lexists[y][\formula{B}(y)]}, just={\TRule{\True}{\lif}[2]}]
        ]
      ]
    ]
  ]
\end{oltableau}
دواړه نوي نښه لرونکي فارمولونه د ځانګړي متغير د شرط لرونکې
قاعدې غواړي، نو راتلونکے ګام بايد همدا وي:
\begin{oltableau}
  [\sFmla{\True}{\lforall[x][\formula{A}(x)]}, just=\TAss
    [\sFmla{\True}{\lforall[x][\formula{A}(x)] \lif
        \lexists[y][\formula{B}(y)]}, just=\TAss, checked
      [\sFmla{\True}{\lnot\lexists[y][\formula{B}(y)]}, just=\TAss, checked
        [\sFmla{\False}{\lexists[y][\formula{B}(y)]},
          just={\TRule{\True}{\lnot}[3]}
          [\sFmla{\False}{\lforall[x][\formula{A}(x)]}, just={\TRule{\True}{\lif}[2]},checked
            [\sFmla{\False}{\formula{A}(b)}, just={\TRule{\False}{\lforall}[5]}]
          ]
          [\sFmla{\True}{\lexists[y][\formula{B}(y)]}, just={\TRule{\True}{\lif}[2]},checked
            [\sFmla{\True}{\formula{B}(c)}, just={\TRule{\True}{\lexists}[5]}]
          ]
        ]
      ]
    ]
  ]
\end{oltableau}
د څانګو د تړلو دپاره بايد پر~\LR{$1$} او~\LR{$4$} کرښو
نښه لرونکي فارمولونه وکاروو۔ د هغو اړوندې قاعدې
(\TRule{\True}{\lforall} او~\TRule{\False}{\lexists}) د ځانګړي
متغير شرط نۀ لري، نو هرې مناسبې اصطلاحګانې ټاکلے شو۔ په دې حالت
کښې په خپل وار همدا \LR{$b$} او~\LR{$c$} دي۔
\begin{oltableau}
  [\sFmla{\True}{\lforall[x][\formula{A}(x)]}, just=\TAss
    [\sFmla{\True}{\lforall[x][\formula{A}(x)] \lif
        \lexists[y][\formula{B}(y)]}, just=\TAss, checked
      [\sFmla{\True}{\lnot\lexists[y][\formula{B}(y)]}, just=\TAss, checked
        [\sFmla{\False}{\lexists[y][\formula{B}(y)]},
          just={\TRule{\True}{\lnot}[3]}
          [\sFmla{\False}{\lforall[x][\formula{A}(x)]},
            just={\TRule{\True}{\lif}[2]},checked
            [\sFmla{\False}{\formula{A}(b)},
              just={\TRule{\False}{\lforall}[5]}
              [\sFmla{\True}{\formula{A}(b)},
                just={\TRule{\True}{\lforall}[1]},close
              ]
            ]
          ]
          [\sFmla{\True}{\lexists[y][\formula{B}(y)]},
            just={\TRule{\True}{\lif}[2]},checked
            [\sFmla{\True}{\formula{B}(c)},
              just={\TRule{\True}{\lexists}[5]}
              [\sFmla{\False}{\formula{B}(c)},
                just={\TRule{\False}{\lexists}[4]},close
              ]
            ]
          ]
        ]
      ]
    ]
  ]
\end{oltableau}
\end{ex}
\olpseganchor{OLP-0104-B010}{end}

\olpseganchor{OLP-0104-B011}{start}
\begin{prob}
د لاندې دپاره تړلي تابلوګانې ورکړئ:
\begin{enumerate}
\item \LR{$\sFmla{\False}{(\lforall[x][!A(x)] \land \lforall[y][!B(y)])
\lif \lforall[z][(!A(z) \land !B(z))]}$}.
\item \LR{$\sFmla{\False}{(\lexists[x][!A(x)] \lor \lexists[y][!B(y)])
\lif \lexists[z][(!A(z) \lor !B(z))]}$}.
\item \LR{$\sFmla{\True}{\lforall[x][(!A(x) \lif !B)]},
\sFmla{\False}{\lexists[y][!A(y)] \lif !B}$}.
\item \LR{$\sFmla{\True}{\lforall[x][\lnot !A(x)]},
\sFmla{\False}{\lnot\lexists[x][!A(x)]}$}.
\item \LR{$\sFmla{\False}{\lnot\lexists[x][!A(x)] \lif \lforall[x][\lnot
!A(x)]}$}.
\item \LR{$\sFmla{\False}{\lnot\lexists[x][\lforall[y][((!A(x,y) \lif
\lnot !A(y,y)) \land (\lnot !A(y,y) \lif !A(x,y)))]]}$}.
\end{enumerate}
\end{prob}
\olpseganchor{OLP-0104-B011}{end}

\olpseganchor{OLP-0104-B012}{start}
\begin{prob}
د لاندې دپاره تړلي تابلوګانې ورکړئ:
\begin{enumerate}
\item \LR{$\sFmla{\False}{\lnot\lforall[x][!A(x)] \lif \lexists[x][\lnot!A(x)]}$}.
\item \LR{$\sFmla{\True}{(\lforall[x][!A(x)] \lif !B)}, \sFmla{\False}{\lexists[y][(!A(y) \lif !B)]}$}.
\item \LR{$\sFmla{\False}{\lexists[x][(!A(x) \lif \lforall[y][!A(y)])]}$}.
\end{enumerate}
\end{prob}
\olpseganchor{OLP-0104-B012}{end}

\olpseganchor{OLP-0105-B004}{start}
\olfileid{fol}{tab}{ptn}
\olpseganchor{OLP-0105-B004}{end}
      
\olpseganchor{OLP-0105-B005}{start}
\olsection{ثبوت-تيوريکي مفاهيم}
\olpseganchor{OLP-0105-B005}{end}

\olpseganchor{OLP-0105-B006}{start}
\begin{editorial}
دا برخه د تابلوګانو د اثبات‌پذيرۍ اړيکې او سازګارۍ تعريفونه
راټولوي۔
\end{editorial}
\olpseganchor{OLP-0105-B006}{end}

\olpseganchor{OLP-0105-B007}{start}
\begin{explain}
لکه څنګه چې مو يو شمېر مهم معنايي مفاهيم (اعتبار، استلزام او د
صدق وړتيا) تعريف کړل، اوس ورته \emph{ثبوت-تيوريکي مفاهيم}
تعريفوو۔ دا مفاهيم په جوړښتونه کښې د تړلي فارمولونه د پوره
کېدو له مخې نۀ، بلکې د ځينو تړلو تابلوګانې د موجوديت له مخې
تعريفېږي۔ دا مهمه موندنه وه چې دغه مفاهيم سره سمون خوري۔ د دې
سمون منځپانګه د \emph{سموالي} او \emph{بشپړتيا قضيې} دي۔
\end{explain}
\olpseganchor{OLP-0105-B007}{end}


\olpseganchor{OLP-0105-B008}{start}
\begin{defn}[قضيې]
يوه تړلے فارمول~\LR{$!A$} هله \emph{قضيه} ده چې د
\LR{$\sFmla{\False}{!A}$} دپاره تړلے تابلو موجود وي۔ که \LR{$!A$}
قضيه وي، \LR{$\Proves !A$} ليکو، او که نۀ وي \LR{$\Proves/ !A$} ليکو۔
\end{defn}
\olpseganchor{OLP-0105-B008}{end}

\olpseganchor{OLP-0105-B009}{start}
\begin{defn}[د اشتقاق وړتيا]
تړلے فارمول~\LR{$!A$} د تړلي فارمولونه له سټ~\LR{$\Gamma$} څخه،
\LR{$\Gamma \Proves !A$}، هله \emph{د اشتقاق وړ} ده چې داسې متناهي
سټ~\LR{$\{!B_1, \dots, !B_n\} \subseteq \Gamma$} او د لاندې سټ دپاره
تړلے تابلو موجود وي:
\[
\{
  \sFmla{\False}{!A},
  \sFmla{\True}{!B_1}, \dots,
  \sFmla{\True}{!B_n}
  \}.
\]
که \LR{$!A$} له~\LR{$\Gamma$} څخه د اشتقاق وړ نۀ وي،
\LR{$\Gamma \Proves/ !A$} ليکو۔
\end{defn}
\olpseganchor{OLP-0105-B009}{end}

\olpseganchor{OLP-0105-B010}{start}
\begin{defn}[سازګاري]
د تړلي فارمولونه يو سټ~\LR{$\Gamma$} هله او يوازې هله \emph{ناسازګار}
دے چې داسې متناهي سټ~\LR{$\{!B_1, \dots, !B_n\} \subseteq \Gamma$}
او د لاندې سټ دپاره تړلے تابلو موجود وي:
\[
\{
  \sFmla{\True}{!B_1}, \dots,
  \sFmla{\True}{!B_n}  
  \}.
\]
که \LR{$\Gamma$} ناسازګار نۀ وي، \emph{سازګار} ئې بولو۔
\end{defn}
\olpseganchor{OLP-0105-B010}{end}

\olpseganchor{OLP-0105-B011}{start}
\begin{prop}[انعکاسيت]
\ollabel{prop:reflexivity}
که \LR{$!A \in \Gamma$} وي، نو \LR{$\Gamma \Proves !A$}۔
\end{prop}
\olpseganchor{OLP-0105-B011}{end}

\olpseganchor{OLP-0105-B012}{start}
\begin{proof}
که \LR{$!A \in \Gamma$} وي، نو \LR{$\{!A\}$} د~\LR{$\Gamma$} يو متناهي فرعي سټ
دے، او لاندې تابلو تړلے دے:
\begin{oltableau}
  [\sFmla{\False}{\formula{A}}, just = \TAss
    [\sFmla{\True}{\formula{A}}, just = \TAss,close]
  ]
\end{oltableau}
\end{proof}
\olpseganchor{OLP-0105-B012}{end}
  

\olpseganchor{OLP-0105-B013}{start}
\begin{prop}[يکنواختي]
\ollabel{prop:monotonicity}
که \LR{$\Gamma \subseteq \Delta$} او \LR{$\Gamma \Proves !A$} وي، نو
\LR{$\Delta \Proves !A$}۔
\end{prop}
\olpseganchor{OLP-0105-B013}{end}

\olpseganchor{OLP-0105-B014}{start}
\begin{proof}
د~\LR{$\Gamma$} هر متناهي فرعي سټ د~\LR{$\Delta$} هم متناهي فرعي سټ دے۔
\end{proof}
\olpseganchor{OLP-0105-B014}{end}

\olpseganchor{OLP-0105-B015}{start}
\begin{prop}[انتقاليت]
\ollabel{prop:transitivity}
که \LR{$\Gamma \Proves !A$} او
\LR{$\{!A\} \cup \Delta \Proves !B$} وي، نو
\LR{$\Gamma \cup \Delta \Proves !B$}۔
\end{prop}
\olpseganchor{OLP-0105-B015}{end}

\olpseganchor{OLP-0105-B016}{start}
\begin{proof}
که \LR{$\{!A\} \cup \Delta \Proves !B$} وي، نو داسې متناهي فرعي سټ
\LR{$\Delta_0 = \{!C_1, \dots, !C_n\} \subseteq \Delta$} شته چې
\begin{align*}
\{\sFmla{\False}{!B}, & \sFmla{\True}{!A}, \sFmla{\True}{!C_1},
\dots, \sFmla{\True}{!C_n}\}
\intertext{\RL{تړلے تابلو لري۔ که \LR{$\Gamma \Proves !A$} وي، نو
  داسې \LR{$!D_1$}, \dots, \LR{$!D_m \in \Gamma$} شته چې}}
\{\sFmla{\False}{!A}, & \sFmla{\True}{!D_1},
\dots, \sFmla{\True}{!D_m}\}
\end{align*}
تړلے تابلو لري۔
\olpseganchor{OLP-0105-B016}{end}

\olpseganchor{OLP-0105-B017}{start}
اوس د لاندې فرضونو تابلو ونيسئ:
\[
\sFmla{\False}{!B},
\sFmla{\True}{!C_1}, \dots, \sFmla{\True}{!C_n},
\sFmla{\True}{!D_1}, \dots, \sFmla{\True}{!D_m}.
\]
پر~\LR{$!A$} د~\Cut{} قاعده پلې کړئ۔ دا دوه څانګې جوړوي: يوه پکې
\LR{$\sFmla{\True}{!A}$}، او بله پکې \LR{$\sFmla{\False}{!A}$} لري۔ نو په
يوه څانګه کښې دا ټول موجود دي:
\[
\{\sFmla{\False}{!B}, \sFmla{\True}{!A},
\sFmla{\True}{!C_1}, \dots, \sFmla{\True}{!C_n}\}
\]
ځکه د دغو فرضونو دپاره تړلے تابلو شته، هغه له دې څانګې سره
نښلولے شو؛ هره هغه څانګه تړل کېږي چې له
\LR{$\sFmla{\True}{!A}$} څخه تېرېږي۔ په بله څانګه کښې دا ټول موجود دي:
\[
\{\sFmla{\False}{!A}, \sFmla{\True}{!D_1}, \dots,
\sFmla{\True}{!D_m}\}
\]
نو بل لورے هم بشپړولے او تړلے تابلو اخستلے شو۔ دا
\LR{$\Gamma \cup \Delta \Proves !B$} ښيي۔
\end{proof}
\olpseganchor{OLP-0105-B017}{end}

\olpseganchor{OLP-0105-B018}{start}
پام وکړئ د دې يوه ځانګړې پايله دا ده: که \LR{$\Gamma \Proves !A$} او
\LR{$!A \Proves !B$} وي، نو \LR{$\Gamma \Proves !B$}۔ له دې څخه دا هم
لازمېږي چې که \LR{$!A_1,\dots,!A_n \Proves !B$} او د هر~\LR{$i$} دپاره
\LR{$\Gamma \Proves !A_i$} وي، نو \LR{$\Gamma \Proves !B$}۔
\olpseganchor{OLP-0105-B018}{end}

\olpseganchor{OLP-0105-B019}{start}
\begin{prop}
\ollabel{prop:incons}
\LR{$\Gamma$} هله او يوازې هله ناسازګار دے چې د هرې
تړلے فارمول~\LR{$!A$} دپاره \LR{$\Gamma \Proves !A$} وي۔
\end{prop}
\olpseganchor{OLP-0105-B019}{end}

\olpseganchor{OLP-0105-B020}{start}
\begin{proof}
تمرين۔
\end{proof}
\olpseganchor{OLP-0105-B020}{end}

\olpseganchor{OLP-0105-B021}{start}
\begin{prob}
\olref[fol][tab][ptn]{prop:incons} ثابت کړئ۔
\end{prob}
\olpseganchor{OLP-0105-B021}{end}



\olpseganchor{OLP-0105-B023}{start}
\begin{prop}[تراکم]
  \ollabel{prop:proves-compact}
  \begin{enumerate}
  \item که \LR{$\Gamma \Proves !A$} وي، نو داسې متناهي فرعي سټ
    \LR{$\Gamma_0 \subseteq \Gamma$} شته چې \LR{$\Gamma_0 \Proves !A$} وي۔
  \item که د~\LR{$\Gamma$} هر متناهي فرعي سټ سازګار وي، نو \LR{$\Gamma$}
    سازګار دے۔
  \end{enumerate}
\end{prop}
\olpseganchor{OLP-0105-B023}{end}

\olpseganchor{OLP-0105-B024}{start}
\begin{proof}
  \begin{enumerate}
    \item که \LR{$\Gamma \Proves !A$} وي، نو داسې متناهي فرعي سټ
      \LR{$\Gamma_0 = \{!B_1, \dots, !B_n\}$} او د لاندې دپاره تړلے
      تابلو شته:
      \[
      \{\sFmla{\False}{!A}, \sFmla{\True}{!B_1}, \dots, \sFmla{\True}{!B_n}\}
      \]
      همدا تابلو، \LR{$\Gamma_0 \Proves !A$} هم ښيي۔
    \item که \LR{$\Gamma$} ناسازګار وي، نو د کوم متناهي فرعي سټ
      \LR{$\Gamma_0 = \{!B_1, \dots, !B_n\}$} دپاره د لاندې يو تړلے
      تابلو شته:
      \[
      \{\sFmla{\True}{!B_1}, \dots, \sFmla{\True}{!B_n}\}
      \]
      دا تړلے تابلو ښيي چې \LR{$\Gamma_0$} ناسازګار دے۔
  \end{enumerate}
\end{proof}
\olpseganchor{OLP-0105-B024}{end}

\olpseganchor{OLP-0106-B004}{start}
\olfileid{fol}{tab}{prv}
\olpseganchor{OLP-0106-B004}{end}

\olpseganchor{OLP-0106-B005}{start}
\olsection{د اشتقاق وړتيا او سازګاري}
\olpseganchor{OLP-0106-B005}{end}

\olpseganchor{OLP-0106-B006}{start}
اوس د اشتقاق وړتيا اړيکې يو شمېر خاصيتونه ثابتوو۔ دوئ په
خپلواک ډول هم په زړه پورې دي، خو هر يو به د بشپړتيا قضیې په ثبوت
کښې ونډه ولري۔
\olpseganchor{OLP-0106-B006}{end}

\olpseganchor{OLP-0106-B007}{start}
\begin{prop}\ollabel{prop:provability-contr}
  که \LR{$\Gamma \Proves !A$} او \LR{$\Gamma \cup \{!A\}$} ناسازګار وي،
  نو \LR{$\Gamma$} ناسازګار دے۔
\end{prop}
\olpseganchor{OLP-0106-B007}{end}

\olpseganchor{OLP-0106-B008}{start}
\begin{proof}
داسې متناهي
\LR{$\Gamma_0 = \{!B_1, \dots, !B_n\}$} او
\LR{$\Gamma_1 =\{!C_1, \dots, !C_m\} \subseteq \Gamma$} شته چې
  \begin{align*}
    \{\sFmla{\False}{!A}, &
    \sFmla{\True}{!B_1}, \dots, \sFmla{\True}{!B_n}\} \\
    \{\sFmla{\True}{!A}, &
    \sFmla{\True}{!C_1}, \dots, \sFmla{\True}{!C_m}\}
  \end{align*}
  تړلي تابلوګانې لري۔ پر~\LR{$!A$} د~\Cut{} قاعدې په کارولو دا په
  يوه تړلي تابلو کښې يو ځاے کولے شو چې د
  \LR{$\Gamma_0 \cup \Gamma_1$} ناسازګاري ښيي۔ ځکه
  \LR{$\Gamma_0 \subseteq \Gamma$} او~\LR{$\Gamma_1 \subseteq \Gamma$}،
  نو \LR{$\Gamma_0 \cup \Gamma_1 \subseteq \Gamma$}؛ له دې امله
  \LR{$\Gamma$} ناسازګار دے۔
\end{proof}
\olpseganchor{OLP-0106-B008}{end}

\olpseganchor{OLP-0106-B009}{start}
\begin{prop}
\ollabel{prop:prov-incons}
\LR{$\Gamma \Proves !A$} هله او يوازې هله چې
\LR{$\Gamma \cup \{\lnot !A\}$} ناسازګار وي۔
\end{prop}
\olpseganchor{OLP-0106-B009}{end}

\olpseganchor{OLP-0106-B010}{start}
\begin{proof}
لومړے فرض کړئ \LR{$\Gamma \Proves !A$}، يعنې د لاندې دپاره تړلے
تابلو شته:
\[
\{\sFmla{\False}{!A},
\sFmla{\True}{!B_1}, \dots, \sFmla{\True}{!B_n}\}
\]
د~\LR{$\TRule{\True}{\lnot}$} قاعدې په کارولو دا د لاندې دپاره په تړلي
تابلو اړولے شو:
\[
\{\sFmla{\True}{\lnot !A},
\sFmla{\True}{!B_1}, \dots, \sFmla{\True}{!B_n}\}.
\]
\olpseganchor{OLP-0106-B010}{end}

\olpseganchor{OLP-0106-B011}{start}
بل لوري ته، که د وروستي سټ دپاره تړلے تابلو وي، د لومړي سټ
په تړلي تابلو ئې اړولے شو: هغه هر فارمول ړنګوو چې پر لومړي
فرض~\LR{$\sFmla{\True}{\lnot !A}$} د~\TRule{\True}{\lnot} له کارونې
راوتلے وي، خپله دغه فرض هم ړنګوو، او فرض
\LR{$\sFmla{\False}{!A}$} ورزياتوو۔ ځکه که کومه څانګه مخکښې د
\TRule{\True}{\lnot} د هغې پايلې له امله تړلې وه چې پر
\LR{$\sFmla{\True}{\lnot !A}$} پلې شوې، يعنې
\LR{$\sFmla{\False}{!A}$}، نو په نوي تابلو کښې اړونده څانګه هم
تړلې ده۔ که په زوړ تابلو کښې کومه څانګه د فرض
\LR{$\sFmla{\True}{\lnot !A}$} او~\LR{$\sFmla{\False}{\lnot !A}$} دواړو له
امله تړلې وه، پر~\LR{$\sFmla{\False}{\lnot !A}$} د
\LR{$\TRule{\False}{\lnot}$} په پلي کولو او د
\LR{$\sFmla{\True}{!A}$} په اخستلو ئې په تړلې څانګه اړولے شو۔ دا څانګه
ځکه تړل کېږي چې \LR{$\sFmla{\False}{!A}$} مو د فرض په توګه ورزيات کړے
دے۔
\end{proof}
\olpseganchor{OLP-0106-B011}{end}

\olpseganchor{OLP-0106-B012}{start}
\begin{prob}
ثابت کړئ چې \LR{$\Gamma \Proves \lnot !A$} هله او يوازې هله چې
\LR{$\Gamma \cup \{!A\}$} ناسازګار وي۔
\end{prob}
\olpseganchor{OLP-0106-B012}{end}

\olpseganchor{OLP-0106-B013}{start}
\begin{prop}\ollabel{prop:explicit-inc}
  که \LR{$\Gamma \Proves !A$} او~\LR{$\lnot !A \in \Gamma$} وي، نو
  \LR{$\Gamma$} ناسازګار دے۔
\end{prop}
\olpseganchor{OLP-0106-B013}{end}

\olpseganchor{OLP-0106-B014}{start}
\begin{proof}
  فرض کړئ \LR{$\Gamma \Proves !A$} او~\LR{$\lnot !A \in \Gamma$}۔ بيا داسې
  \LR{$!B_1$}، \dots، \LR{$!B_n \in \Gamma$} شته چې
  \[
  \{\sFmla{\False}{!A}, \sFmla{\True}{!B_1}, \dots, \sFmla{\True}{!B_n}\}
  \]
  تړلے تابلو لري۔ فرض~\sFmla{\False}{!A} په
  \sFmla{\True}{\lnot !A} بدل کړئ، او د فرضونو وروسته هغه
  \sFmla{\False}{!A} پايله ورننباسئ چې پر
  \sFmla{\True}{\lnot !A} د~\TRule{\True}{\lnot} له کارونې
  راځي۔ په زوړ تابلو کښې هره تړلے فارمول چې د~\LR{$1$} کرښې په
  حواله توجيه شوې وه، اوس د~\LR{$n+2$} کرښې په حواله توجيه کېږي۔ له فرضونو وروسته د زړې
  تابلو نورې کرښې هم يوه کرښه وروسته ځي، او حوالې يې په همدې حساب
  بدلېږي۔ که
  زوړ تابلو تړلے و، نوے تابلو هم تړلے دے۔ دا د~\LR{$\Gamma$} ناسازګاري
  ښيي؛ ځکه ټول فرضونه په~\LR{$\Gamma$} کښې دي۔
\end{proof}
\olpseganchor{OLP-0106-B014}{end}

\olpseganchor{OLP-0106-B015}{start}
\begin{prop}\ollabel{prop:provability-exhaustive}
  که \LR{$\Gamma \cup \{!A\}$} او~\LR{$\Gamma \cup \{\lnot !A\}$} دواړه
  ناسازګار وي، نو \LR{$\Gamma$} ناسازګار دے۔
\end{prop}
\olpseganchor{OLP-0106-B015}{end}

\olpseganchor{OLP-0106-B016}{start}
\begin{proof}
  که داسې \LR{$!B_1$}، \dots، \LR{$!B_n \in \Gamma$} او
  \LR{$!C_1$}، \dots، \LR{$!C_m \in \Gamma$} وي چې
  \begin{align*}
    \{\sFmla{\True}{!A}, &
    \sFmla{\True}{!B_1}, \dots, \sFmla{\True}{!B_n}\} \text{\RL{ او}}\\
    \{\sFmla{\True}{\lnot !A}, &
    \sFmla{\True}{!C_1}, \dots, \sFmla{\True}{!C_m}\}
  \end{align*}
  دواړه تړلي تابلوګانې ولري، يو داسې ګډ تابلو جوړولے شو
  چې د~\LR{$\Gamma$} ناسازګاري ښيي۔ د فرضونو په توګه
  \LR{$\sFmla{\True}{!B_1}$}، \dots، \LR{$\sFmla{\True}{!B_n}$} له
  \LR{$\sFmla{\True}{!C_1}$}، \dots، \LR{$\sFmla{\True}{!C_m}$} سره يو
  ځاے کاروو، او بيا د~\Cut{} قاعده پلې کوو۔ دا دوه څانګې ورکوي:
  يوه په~\LR{$\sFmla{\True}{!A}$} او بله په~\LR{$\sFmla{\False}{!A}$} پيلېږي۔
\olpseganchor{OLP-0106-B016}{end}
  
\olpseganchor{OLP-0106-B017}{start}
  په کيڼ لوري کښې د لومړي تابلو د فرضونو لاندې برخه ورزياته
  کړئ۔ دلته د قاعدې هره کارونه لا هم سمه ده؛ ځکه د لومړي
  تابلو هر فرض، د~\LR{$\sFmla{\True}{!A}$} په شمول، موجود دے۔
  نو د~\LR{$\sFmla{\True}{!A}$} لاندې هره څانګه تړل کېږي۔
\olpseganchor{OLP-0106-B017}{end}
  
\olpseganchor{OLP-0106-B018}{start}
  په ښي لوري کښې د دويم تابلو د فرض لاندې برخه ورزياته کړئ،
  خو د~\LR{$\sFmla{\True}{\lnot !A}$} پر فرض د
  \LR{$\TRule{\True}{\lnot}$} د هرې کارونې پايلې ترې لرې کړئ۔ پر
  \LR{$\sFmla{\True}{\lnot !A}$} د~\LR{$\TRule{\True}{\lnot}$} پايله
  \LR{$\sFmla{\False}{!A}$} ده، چې بيا هم موجوده ده؛ ځکه د ګډ
  تابلو په ښي لوري کښې د~\Cut{} قاعدې پايله ده۔
\olpseganchor{OLP-0106-B018}{end}

\olpseganchor{OLP-0106-B019}{start}
  که په دويم تابلو کښې کومه څانګه د فرض
  \LR{$\sFmla{\True}{\lnot !A}$} (چې په ګډ تابلو کښې نور فرض نۀ
  دے) او~\LR{$\sFmla{\False}{\lnot !A}$} دواړو له امله تړلې وه، پر
  \LR{$\sFmla{\False}{\lnot !A}$} د~\LR{$\TRule{\False}{\lnot}$} په پلي کولو
  \LR{$\sFmla{\True}{!A}$} اخستلے شو۔ اوس په ګډ تابلو کښې اړونده
  څانګه هم تړل کېږي؛ ځکه د~\Cut{} قاعدې ښۍ پايله،
  \LR{$\sFmla{\False}{!A}$}، پکې ده۔ که په دويم تابلو کښې يوه
  څانګه د بل کوم لامل له امله تړلې وه، په ګډ تابلو کښې
  اړونده څانګه هم تړلې ده؛ ځکه په زوړ، دويم تابلو کښې پر
  څانګه له~\LR{$\sFmla{\True}{\lnot !A}$} پرته هر بل نښه لرونکي فارمولونه
  د ګډ تابلو پر اړوندې څانګې هم راځي۔
  \end{proof}
\olpseganchor{OLP-0106-B019}{end}

\olpseganchor{OLP-0107-B005}{start}
\olfileid{fol}{tab}{ppr}
\olpseganchor{OLP-0107-B005}{end}

\olpseganchor{OLP-0107-B006}{start}
\olsection{د اشتقاق وړتيا او قضيوي رابطونه}
\olpseganchor{OLP-0107-B006}{end}

\olpseganchor{OLP-0107-B007}{start}
\begin{explain}
  ثابتوو چې د تابلوګانو د اشتقاق وړتيا اړيکه~\LR{$\Proves$} دومره
  پياوړې ده چې د قضيوي رابطونو په اړه ځينې بنسټيز حقيقتونه ثابت
  کړي؛ لکه \LR{$!A \land !B \Proves !A$} او
  \LR{$!A, !A \lif !B \Proves !B$} (مودوس پوننس)۔ دا حقيقتونه د
  بشپړتيا قضیې د ثبوت دپاره پکار دي۔
\end{explain}
\olpseganchor{OLP-0107-B007}{end}

\olpseganchor{OLP-0107-B008}{start}
\begin{prop}\ollabel{prop:provability-land}
  \begin{enumerate}
  \item \ollabel{prop:provability-land-left} دواړه
    \LR{$!A \land !B \Proves !A$} او~\LR{$!A \land !B \Proves !B$}۔
  \item \ollabel{prop:provability-land-right}
    \LR{$!A, !B \Proves !A \land !B$}۔
  \end{enumerate}
\end{prop}
\olpseganchor{OLP-0107-B008}{end}

\olpseganchor{OLP-0107-B009}{start}
\begin{proof}
  \begin{enumerate}
  \item د~\LR{$\{\sFmla{\False}{!A}, \sFmla{\True}{!A \land !B}\}$} او
    \LR{$\{\sFmla{\False}{!B}, \sFmla{\True}{!A \land !B}\}$} دواړو
    دپاره تړلي تابلوګانې شته:
    \begin{oltableau}{}
      [\sFmla{\False}{\formula{A}}, just=\TAss
        [\sFmla{\True}{\formula{A} \land \formula{B}}, just=\TAss
          [\sFmla{\True}{\formula{A}},just={\TRule{\True}{\land}[2]}
            [\sFmla{\True}{\formula{B}},just={\TRule{\True}{\land}[2]}, close
            ]
          ]
        ]
      ]
    \end{oltableau}
    \begin{oltableau}{}
      [\sFmla{\False}{\formula{B}}, just=\TAss
        [\sFmla{\True}{\formula{A} \land \formula{B}}, just=\TAss
          [\sFmla{\True}{\formula{A}},just={\TRule{\True}{\land}[2]}
            [\sFmla{\True}{\formula{B}},just={\TRule{\True}{\land}[2]}, close
            ]
          ]
        ]
      ]
    \end{oltableau}
    \item د~\LR{$\{\sFmla{\True}{!A},
      \sFmla{\True}{!B}, \sFmla{\False}{!A \land !B}\}$} دپاره
      لاندې تړلے تابلو دے:
      \begin{oltableau}
        [\sFmla{\False}{\formula{A} \land \formula{B}}, just = \TAss
          [\sFmla{\True}{\formula{A}}, just = \TAss
            [\sFmla{\True}{\formula{B}}, just=\TAss
              [\sFmla{\False}{\formula{A}}, just = {\TRule{\False}{\land}[1]}, close]
              [\sFmla{\False}{\formula{B}}, just = {\TRule{\False}{\land}[1]}, close]
            ]
          ]
        ]
      \end{oltableau}
  \end{enumerate}
\end{proof}
\olpseganchor{OLP-0107-B009}{end}

\olpseganchor{OLP-0107-B010}{start}
\begin{prop}\ollabel{prop:provability-lor}
  \begin{enumerate}
  \item \LR{$\{!A \lor !B, \lnot !A, \lnot !B\}$} ناسازګار دے۔
  \item دواړه \LR{$!A \Proves !A \lor !B$} او
    \LR{$!B \Proves !A \lor !B$}۔
  \end{enumerate}
\end{prop}
\olpseganchor{OLP-0107-B010}{end}

\olpseganchor{OLP-0107-B011}{start}
\begin{proof}
  \begin{enumerate}
  \item د~\LR{$\{\sFmla{\True}{!A \lor !B},
    \sFmla{\True}{\lnot !A}, \sFmla{\True}{\lnot !B}\}$} تړلے
    تابلو ورکوو:
    \begin{oltableau}
      [\sFmla{\True}{\formula{A} \lor \formula{B}}, just = \TAss
        [\sFmla{\True}{\lnot \formula{A}}, just = \TAss
          [\sFmla{\True}{\lnot \formula{B}}, just = \TAss
            [\sFmla{\False}{\formula{A}}, just = {\TRule{\True}{\lnot}[2]}
              [\sFmla{\False}{\formula{B}}, just = {\TRule{\True}{\lnot}[3]}
                [\sFmla{\True}{\formula{A}}, just = {\TRule{\True}{\lor}[1]}, close]
                [\sFmla{\True}{\formula{B}}, just = {\TRule{\True}{\lor}[1]}, close]
              ]
            ]
          ]
        ]
      ]
    \end{oltableau}
  \item د~\LR{$\{\sFmla{\False}{!A \lor !B}, \sFmla{\True}{!A}\}$} او
    \LR{$\{\sFmla{\False}{!A \lor !B}, \sFmla{\True}{!B}\}$} دواړو
    دپاره تړلي تابلوګانې شته:
    \begin{oltableau}{}
      [\sFmla{\False}{\formula{A} \lor \formula{B}}, just=\TAss
        [\sFmla{\True}{\formula{A}}, just=\TAss
          [\sFmla{\False}{\formula{A}},just={\TRule{\False}{\lor}[1]}
            [\sFmla{\False}{\formula{B}},just={\TRule{\False}{\lor}[1]}, close
            ]
          ]
        ]
      ]
    \end{oltableau}
    \begin{oltableau}{}
      [\sFmla{\False}{\formula{A} \lor \formula{B}}, just=\TAss
        [\sFmla{\True}{\formula{B}}, just=\TAss
          [\sFmla{\False}{\formula{A}},just={\TRule{\False}{\lor}[1]}
            [\sFmla{\False}{\formula{B}},just={\TRule{\False}{\lor}[1]}, close
            ]
          ]
        ]
      ]
    \end{oltableau}
  \end{enumerate}
\end{proof}
\olpseganchor{OLP-0107-B011}{end}

\olpseganchor{OLP-0107-B012}{start}
\begin{prop}\ollabel{prop:provability-lif}
  \begin{enumerate}
  \item \ollabel{prop:provability-lif-left}
    \LR{$!A, !A \lif !B \Proves !B$}۔
  \item \ollabel{prop:provability-lif-right}
    دواړه \LR{$\lnot !A \Proves !A \lif !B$} او
    \LR{$!B \Proves !A \lif !B$}۔
  \end{enumerate}
\end{prop}
\olpseganchor{OLP-0107-B012}{end}

\olpseganchor{OLP-0107-B013}{start}
\begin{proof}
  \begin{enumerate}
    \item د~\LR{$\{\sFmla{\False}{!B}, \sFmla{\True}{!A \lif !B},
      \sFmla{\True}{!A}\}$} يو تړلے تابلو شته:
      \begin{oltableau}
        [\sFmla{\False}{\formula{B}}, just=\TAss
          [\sFmla{\True}{\formula{A} \lif \formula{B}}, just=\TAss
            [\sFmla{\True}{\formula{A}}, just=\TAss
              [\sFmla{\False}{\formula{A}}, just = {\TRule{\True}{\lif}[2]}, close]
              [\sFmla{\True}{\formula{B}}, just = {\TRule{\True}{\lif}[2]}, close]
            ]
          ]
        ]
      \end{oltableau}
    \item د~\LR{$\{\sFmla{\False}{!A \lif !B},
      \sFmla{\True}{\lnot !A}\}$} او
      \LR{$\{\sFmla{\False}{!A \lif !B}, \sFmla{\True}{!B}\}$} دواړو
      دپاره تړلي تابلوګانې شته:
      \begin{oltableau}
        [\sFmla{\False}{\formula{A} \lif \formula{B}}, just = \TAss
          [\sFmla{\True}{\lnot \formula{A}}, just = \TAss
            [\sFmla{\True}{\formula{A}}, just = {\TRule{\False}{\lif}[1]}
              [\sFmla{\False}{\formula{B}}, just = {\TRule{\False}{\lif}[1]}
                [\sFmla{\False}{\formula{A}}, just = {\TRule{\True}{\lnot}[2]}, close]
              ]
            ]
          ]
        ]
      \end{oltableau}
      \begin{oltableau}
        [\sFmla{\False}{\formula{A} \lif \formula{B}}, just = \TAss
          [\sFmla{\True}{\formula{B}}, just = \TAss
            [\sFmla{\True}{\formula{A}}, just = {\TRule{\False}{\lif}[1]}
              [\sFmla{\False}{\formula{B}}, just = {\TRule{\False}{\lif}[1]},close
              ]
            ]
          ]
        ]
      \end{oltableau}
  \end{enumerate}
\end{proof}
\olpseganchor{OLP-0107-B013}{end}

\olpseganchor{OLP-0108-B004}{start}
\olfileid{fol}{tab}{qpr}
\olsection{د اشتقاق وړتيا او سورونه}
\olpseganchor{OLP-0108-B004}{end}

\olpseganchor{OLP-0108-B005}{start}
\begin{explain}
  د بشپړتيا قضيه دا هم غواړي چې د تابلو قاعدې د~\LR{$\Proves$} په اړه
  هغه حقيقتونه ورکړي چې په دې برخه کښې ثابتېږي۔
\end{explain}
\olpseganchor{OLP-0108-B005}{end}

\olpseganchor{OLP-0108-B006}{start}
\begin{thm}
\ollabel{thm:strong-generalization} که \LR{$c$} داسې ثابت وي چې په
\LR{$\Gamma$} يا~\LR{$!A(x)$} کښې نۀ راځي، او \LR{$\Gamma \Proves !A(c)$} وي،
نو \LR{$\Gamma \Proves \lforall[x][!A(x)]$}۔
\end{thm}
\olpseganchor{OLP-0108-B006}{end}

\olpseganchor{OLP-0108-B007}{start}
\begin{proof}
فرض کړئ \LR{$\Gamma \Proves !A(c)$}، يعنې داسې
\LR{$!B_1$}، \dots، \LR{$!B_n \in \Gamma$} او د لاندې دپاره تړلے
تابلو شته:
\begin{align*}
\{ \sFmla{\False}{!A(c)}, &
\sFmla{\True}{!B_1}, \dots, \sFmla{\True}{!B_n} \}.
\intertext{\RL{بايد وښيو چې د لاندې دپاره هم تړلے تابلو شته:}}
\{ \sFmla{\False}{\lforall[x][!A(x)]}, &
\sFmla{\True}{!B_1}, \dots, \sFmla{\True}{!B_n} \}.
\end{align*}
تړلے تابلو واخلئ، لومړے فرض ئې په
\LR{$\sFmla{\False}{\lforall[x][!A(x)]}$} بدل کړئ، او د فرضونو وروسته
\LR{$\sFmla{\False}{!A(c)}$} ورننباسئ۔
\begin{center}
\begin{tableau}{not line numbering}
  [\sFmla{\False}{\formula{A}(c)}
    [\sFmla{\True}{\formula{B}_1}
      [\vdots
      [\sFmla{\True}{\formula{B}_n}, 
        [][][]]]]]
\end{tableau}\qquad
\begin{tableau}{not line numbering}
  [\sFmla{\False}{\lforall[x][\formula{A}(x)]}
    [\sFmla{\True}{\formula{B}_1}
      [\vdots
        [\sFmla{\True}{\formula{B}_n},
          [\sFmla{\False}{\formula{A}(c)}
        [][][]]]]]]
\end{tableau}
\end{center}
تابلو لا هم تړلے دے؛ ځکه ټولې هغه تړلي فارمولونه چې مخکښې د فرضونو
په توګه موجودې وې، اوس هم د~تابلو په سر کښې موجودې دي۔
ننويستل شوې کرښه د~\LR{$\TRule{\False}{\lforall}$} قاعدې د سمې کارونې
پايله ده؛ ځکه ثابت سمبول~\LR{$c$} په~\LR{$!B_1$}، \dots، \LR{$!B_n$} يا
\LR{$\lforall[x][!A(x)]$} کښې نۀ راځي، يعنې په نوي تابلو کښې د
ننويستل شوې کرښې پاس نۀ راځي۔
\end{proof}
\olpseganchor{OLP-0108-B007}{end}

\olpseganchor{OLP-0108-B008}{start}
\begin{prop}
\ollabel{prop:provability-quantifiers}
\begin{enumerate}
\item \LR{$!A(t) \Proves \lexists[x][!A(x)]$}.
\item \LR{$\lforall[x][!A(x)] \Proves !A(t)$}.
\end{enumerate}
\end{prop}
\olpseganchor{OLP-0108-B008}{end}

\olpseganchor{OLP-0108-B009}{start}
\begin{proof}
\begin{enumerate}
\item د
  \LR{$\sFmla{\False}{\lexists[x][!A(x)]}, \sFmla{\True}{!A(t)}$} دپاره
  تړلے تابلو دا دے:
  \begin{oltableau}
    [\sFmla{\False}{\lexists[x][\formula{A}(x)]}, just = \TAss
      [\sFmla{\True}{\formula{A}(t)}, just = \TAss
        [\sFmla{\False}{\formula{A}(t)}, just = {\TRule{\False}{\lexists}[1]}, close]]]
    \end{oltableau}
\item د
  \LR{$\sFmla{\False}{\formula{A}(t)}, \sFmla{\True}{\lforall[x][!A(x)]}, $} دپاره
  تړلے تابلو دا دے:
  \begin{oltableau}
    [\sFmla{\False}{\formula{A}(t)}, just = \TAss
      [\sFmla{\True}{\lforall[x][\formula{A}(x)]}, just = \TAss
        [\sFmla{\True}{\formula{A}(t)}, just = {\TRule{\True}{\lforall}[2]}, close]]]
    \end{oltableau}
\end{enumerate}
\end{proof}
\olpseganchor{OLP-0108-B009}{end}

\olpseganchor{OLP-0109-B004}{start}
\olfileid{fol}{tab}{sou}
\olpseganchor{OLP-0109-B004}{end}
      
\olpseganchor{OLP-0109-B005}{start}
\olsection{سموالے}
\olpseganchor{OLP-0109-B005}{end}

\olpseganchor{OLP-0109-B006}{start}
\begin{explain}
د تابلوګانو په شان اشتقاق نظام هله \emph{سم} دے چې
هغه څيزونه اشتقاق نۀ کړي چې په حقيقت کښې نۀ وي۔ نو سموالے د
اشتقاق نظامونو د تضمين شوې خونديتوب يو ډول خاصيت دے۔
د دې له مخې چې کوم ثبوت-تيوريکي خاصيت تر بحث لاندې وي، مثلاً غواړو
پوه شو چې:
\begin{enumerate}
\item هر د اشتقاق وړ~\LR{$!A$} معتبر دے؛
\item که تړلے فارمول له ځينو نورو څخه د اشتقاق وړ وي، د هغو
  معنايي پايله هم ده؛
\item که د تړلي فارمولونه يو سټ ناسازګار وي، د صدق وړ نۀ دے۔
\end{enumerate}
دا د اشتقاق نظام مهم خاصيتونه دي۔ که له دغو څخه کوم يو
سم نۀ وي، اشتقاق نظام نيمګړے دے---له حده زيات څيزونه به
اشتقاق کوي۔ له دې امله د اشتقاق نظام سموالے ثابتول
تر ټولو زيات اهميت لري۔
\olpseganchor{OLP-0109-B006}{end}

\olpseganchor{OLP-0109-B007}{start}
ځکه دغه ټول ثبوت-تيوريکي خاصيتونه د يو يا بل ډول د تړلو
تابلوګانې له لارې تعريف شوي دي، د پورته (۱)--(۳) ثابتول غواړي
چې د تړلو تابلوګانې د معنايي خاصيتونو په اړه يو څه ثابت کړو۔
لومړے به تعريف کړو چې د يو نښه لرونکے فارمول پوره کېدل په يو
جوړښت کښې څه معنٰی لري؛ بيا به وښيو چې که يو تابلو تړلے وي،
هېڅ جوړښت ئې ټول فرضونه نۀ پوره کوي۔ وروسته (۱)--(۳) د دې پايلې
د لازمو نتيجو په توګه راځي۔
\end{explain}
\olpseganchor{OLP-0109-B007}{end}

\olpseganchor{OLP-0109-B008}{start}
\begin{defn}
جوړښت~\LR{$\Struct{M}$}
يو نښه لرونکے فارمول~\LR{$\sFmla{\True}{!A}$} هله او يوازې هله
\emph{پوره کوي} چې \LR{$\Sat{M}{!A}$}، او
\LR{$\sFmla{\False}{!A}$} هله او يوازې هله پوره کوي چې
\LR{$\Sat/{M}{!A}$}۔
\LR{$\Struct{M}$} د نښه لرونکي فارمولونه يو
سټ~\LR{$\Gamma$} هله او يوازې هله پوره کوي چې هر
\LR{$\sFmla{S}{!A} \in \Gamma$} پوره کړي۔ \LR{$\Gamma$} هله \emph{د صدق وړ}
دے چې کوم جوړښت ئې پوره کړي،
او که داسې نۀ وي \emph{د صدق وړ نۀ} دے۔
\end{defn}
\olpseganchor{OLP-0109-B008}{end}

\olpseganchor{OLP-0109-B009}{start}
\begin{thm}[سموالے]
  \ollabel{thm:tableau-soundness}
  که \LR{$\Gamma$} يو تړلے تابلو ولري، نو \LR{$\Gamma$} د صدق وړ نۀ دے۔
\end{thm}
\olpseganchor{OLP-0109-B009}{end}

\olpseganchor{OLP-0109-B010}{start}
\begin{proof}
د يو تابلو څانګه هله او يوازې هله د صدق وړ بولو چې پر هغې
د نښه لرونکي فارمولونه سټ د صدق وړ وي؛ او تابلو هله د صدق وړ
بولو چې لږ تر لږه يوه د صدق وړ څانګه ولري۔
\olpseganchor{OLP-0109-B010}{end}

\olpseganchor{OLP-0109-B011}{start}
دا خبره ښيو: د استنتاج د کومې قاعدې په وسيله د يو د صدق وړ
تابلو پراخول تل يو د صدق وړ تابلو ورکوي۔ دا به قضيه
ثابته کړي: هر تړلے تابلو د~\LR{$\Gamma$} يوازې له فرضونو څخه جوړ
تابلو ته د استنتاج د قاعدو په پلي کولو تر لاسه کېږي۔ نو که
\LR{$\Gamma$} د صدق وړ واے، د هغې هر تابلو به د صدق وړ و۔ خو تړلے
تابلو په څرګند ډول د صدق وړ نۀ دے: هره څانګه دواړه
\LR{$\sFmla{\True}{!A}$} او \LR{$\sFmla{\False}{!A}$} لري، او هېڅ جوړښت
نۀ شي کولے چې \LR{$!A$} هم پوره او هم نۀ پوره کړي۔
\olpseganchor{OLP-0109-B011}{end}

\olpseganchor{OLP-0109-B012}{start}
فرض کړئ يو د صدق وړ تابلو لرو، يعنې داسې تابلو چې لږ
تر لږه يوه د صدق وړ څانګه لري۔ د استنتاج قاعده يا يوې څانګې ته
نښه لرونکي فارمولونه ورزياتوي يا يوه څانګه په دوو وېشي۔ که په
تابلو کښې کومه د صدق وړ څانګه وي چې اړوند قاعدوي تطبيق ئې نۀ
پراخوي، هغه په پراخ شوي تابلو کښې هم د صدق وړ پاتې کېږي؛ نو
پراخ شوے تابلو د صدق وړ دے۔ ځکه يوازې هغه حالت څېړل پکار دي چې
قاعده پر يوې د صدق وړ څانګې پلي کېږي۔
\olpseganchor{OLP-0109-B012}{end}

\olpseganchor{OLP-0109-B013}{start}
\LR{$\Gamma$} پر هغې څانګې د نښه لرونکي فارمولونه سټ وبولئ، او
\LR{$\sFmla{S}{!A} \in \Gamma$} هغه نښه لرونکے فارمول وبولئ چې قاعده
پرې پلي کېږي۔ که قاعده څانګه نۀ وېشي، بايد وښيو چې پراخه شوې څانګه،
يعنې \LR{$\Gamma$} له قاعدوي پايلو سره، لا هم د صدق وړ ده۔ که قاعده
څانګه وېشي، بايد وښيو چې له دوو ترلاسه شوو څانګو لږ تر لږه يوه د
صدق وړ ده۔
\olpseganchor{OLP-0109-B013}{end}
  
\olpseganchor{OLP-0109-B014}{start}
لومړے هغه ممکن استنتاجونه څېړو چې څانګه نۀ وېشي۔
\begin{enumerate}
\item څانګه پر \LR{$\sFmla{\True}{\lnot !B} \in \Gamma$} د
  \LR{$\TRule{\True}{\lnot}$} په پلي کولو پراخېږي۔ بيا پراخه شوې څانګه
  د \LR{$\Gamma \cup \{\sFmla{\False}{!B}\}$} نښه لرونکي فارمولونه لري۔
  فرض کړئ \LR{$\Sat{M}{\Gamma}$}۔ په ځانګړي ډول
  \LR{$\Sat{M}{\lnot !B}$}۔ نو
  \LR{$\Sat/{M}{!B}$}، يعنې
  \LR{$\Struct{M}$}،
  \LR{$\sFmla{\False}{!B}$} پوره کوي۔
\item څانګه پر \LR{$\sFmla{\False}{\lnot !B} \in \Gamma$} د
  \LR{$\TRule{\False}{\lnot}$} په پلي کولو پراخېږي: تمرين۔
\item څانګه پر \LR{$\sFmla{\True}{!B \land !C} \in \Gamma$} د
  \LR{$\TRule{\True}{\land}$} په پلي کولو پراخېږي؛ له دې څخه پر څانګه
  دوه نوي نښه لرونکي فارمولونه ترلاسه کېږي: \LR{$\sFmla{\True}{!B}$} او
  \LR{$\sFmla{\True}{!C}$}۔ فرض کړئ
  \LR{$\Sat{M}{\Gamma}$}، او په ځانګړي ډول
  \LR{$\Sat{M}{!B \land !C}$}۔ بيا
  \LR{$\Sat{M}{!B}$} او
  \LR{$\Sat{M}{!C}$}۔ يعنې
  \LR{$\Struct{M}$} دواړه
  \LR{$\sFmla{\True}{!B}$} او \LR{$\sFmla{\True}{!C}$} پوره کوي۔
\item څانګه پر \LR{$\sFmla{\False}{!B \lor !C} \in \Gamma$} د
  \LR{$\TRule{\False}{\lor}$} په پلي کولو پراخېږي: تمرين۔
\item څانګه پر \LR{$\sFmla{\False}{!B \lif !C} \in \Gamma$} د
  \LR{$\TRule{\False}{\lif}$} په پلي کولو پراخېږي۔ له دې څخه پر څانګه
  دوه نوي نښه لرونکي فارمولونه ترلاسه کېږي: \LR{$\sFmla{\True}{!B}$} او
  \LR{$\sFmla{\False}{!C}$}۔ فرض کړئ
  \LR{$\Sat{M}{\Gamma}$}، او په ځانګړي ډول
  \LR{$\Sat/{M}{!B \lif !C}$}۔ بيا
  \LR{$\Sat{M}{!B}$} او
  \LR{$\Sat/{M}{!C}$}۔ يعنې
  \LR{$\Struct{M}$} دواړه
  \LR{$\sFmla{\True}{!B}$} او \LR{$\sFmla{\False}{!C}$} پوره کوي۔
\olpseganchor{OLP-0109-B014}{end}

\olpseganchor{OLP-0109-B015}{start}
%
\item څانګه پر \LR{$\sFmla{\True}{\lforall[x][!A(x)]} \in \Gamma$} د
  \LR{$\TRule{\True}{\lforall}$} په پلي کولو پراخېږي۔ له دې څخه پر څانګه
  يو نوے نښه لرونکے فارمول~\LR{$\sFmla{\True}{!A(t)}$} ترلاسه کېږي۔
  فرض کړئ \LR{$\Sat{M}{\Gamma}$}، او په ځانګړي ډول
  \LR{$\Sat{M}{\lforall[x][!A(x)]}$}۔ د
  \olref[syn][sem]{prop:quant-terms} له مخې \LR{$\Sat{M}{!A(t)}$}۔ نو
  \LR{$\Struct{M}$}، \LR{$\sFmla{\True}{!A(t)}$} پوره کوي۔
\olpseganchor{OLP-0109-B015}{end}

\olpseganchor{OLP-0109-B016}{start}
\item څانګه پر \LR{$\sFmla{\False}{\lforall[x][!A(x)]} \in \Gamma$} د
  \LR{$\TRule{\False}{\lforall}$} په پلي کولو پراخېږي۔ له دې څخه پر
  څانګه يو نوے نښه لرونکے فارمول~\LR{$\sFmla{\False}{!A(a)}$} ترلاسه
  کېږي، چې \LR{$a$} داسې ثابت سمبول دے چې په~\LR{$\Gamma$} کښې نۀ راځي۔
  ځکه \LR{$\Gamma$} د صدق وړ دے، داسې \LR{$\Struct{M}$} شته چې
  \LR{$\Sat{M}{\Gamma}$}، او په ځانګړي ډول
  \LR{$\Sat/{M}{\lforall[x][!A(x)]}$}۔ بايد وښيو چې
  \LR{$\Gamma \cup \{\sFmla{\False}{!A(a)}\}$} د صدق وړ دے۔ د دې دپاره
  مناسب~\LR{$\Struct{M'}$} داسې تعريفوو۔
\olpseganchor{OLP-0109-B016}{end}

\olpseganchor{OLP-0109-B017}{start}
  د \olref[syn][ass]{prop:sat-quant} له مخې
  \LR{$\Sat/{M}{\lforall[x][!A(x)]}$} هله او يوازې هله چې د کوم~\LR{$s$}
  دپاره \LR{$\Sat/{M}{!A(x)}[s]$}۔ اوس \LR{$\Struct{M'}$} د~\LR{$\Struct{M}$} په
  شان وبولئ، خو \LR{$\Assign{a}{M'} = s(x)$}۔ د
  \olref[syn][ext]{cor:extensionality-sent} له مخې، د هر
  \LR{$\sFmla{\True}{!C} \in \Gamma$} دپاره \LR{$\Sat{M'}{!C}$}، او د هر
  \LR{$\sFmla{\False}{!C} \in \Gamma$} دپاره \LR{$\Sat/{M'}{!C}$}؛ ځکه
  \LR{$a$} په~\LR{$\Gamma$} کښې نۀ راځي۔
\olpseganchor{OLP-0109-B017}{end}

\olpseganchor{OLP-0109-B018}{start}
  د \olref[syn][ext]{prop:extensionality} له مخې
  \LR{$\Sat/{M'}{!A(x)}[s]$}۔ د \olref[syn][ext]{prop:ext-formulas} له مخې
  \LR{$\Sat/{M'}{!A(a)}[s]$}۔ ځکه \LR{$!A(a)$} يوه تړلے فارمول ده، د
  \olref[syn][ass]{prop:sentence-sat-true} له مخې
  \LR{$\Sat/{M'}{!A(a)}$}؛ يعنې \LR{$\Struct{M'}$}،
  \LR{$\sFmla{\False}{!A(a)}$} پوره کوي۔
\item څانګه پر \LR{$\sFmla{\True}{\lexists[x][!B(x)]} \in \Gamma$} د
  \LR{$\TRule{\True}{\lexists}$} په پلي کولو پراخېږي: تمرين۔
\item څانګه پر \LR{$\sFmla{\False}{\lexists[x][!B(x)]} \in \Gamma$} د
  \LR{$\TRule{\False}{\lexists}$} په پلي کولو پراخېږي: تمرين۔
\end{enumerate}
اوس هغه ممکن استنتاجونه څېړو چې څانګه وېشي۔
\begin{enumerate}
\item څانګه پر \LR{$\sFmla{\False}{!B \land !C} \in \Gamma$} د
  \LR{$\TRule{\False}{\land}$} په پلي کولو پراخېږي؛ له دې څخه دوه څانګې
  ترلاسه کېږي: کيڼه څانګه د~\LR{$\sFmla{\False}{!B}$} له لارې او ښۍ
  څانګه د~\LR{$\sFmla{\False}{!C}$} له لارې دوام کوي۔ فرض کړئ
  \LR{$\Sat{M}{\Gamma}$}، او په ځانګړي ډول
  \LR{$\Sat/{M}{!B \land !C}$}۔ بيا يا
  \LR{$\Sat/{M}{!B}$} يا
  \LR{$\Sat/{M}{!C}$}۔ په لومړي حالت کښې
  \LR{$\Struct{M}$}،
  \LR{$\sFmla{\False}{!B}$} پوره کوي؛ يعنې
  \LR{$\Struct{M}$} د کيڼې څانګې فارمولونه
  پوره کوي۔ په دويم حالت کښې
  \LR{$\Struct{M}$}،
  \LR{$\sFmla{\False}{!C}$} پوره کوي؛ يعنې
  \LR{$\Struct{M}$} د ښۍ څانګې فارمولونه
  پوره کوي۔
\item څانګه پر \LR{$\sFmla{\True}{!B \lor !C} \in \Gamma$} د
  \LR{$\TRule{\True}{\lor}$} په پلي کولو پراخېږي: تمرين۔
\item څانګه پر \LR{$\sFmla{\True}{!B \lif !C} \in \Gamma$} د
  \LR{$\TRule{\True}{\lif}$} په پلي کولو پراخېږي: تمرين۔
\item څانګه د~\Cut{} په وسيله پراخېږي۔ له دې څخه دوه څانګې ترلاسه
  کېږي: يوه \LR{$\sFmla{\True}{!B}$} او بله \LR{$\sFmla{\False}{!B}$} لري۔
  ځکه \LR{$\Sat{M}{\Gamma}$} او يا
  \LR{$\Sat{M}{!B}$} يا
  \LR{$\Sat/{M}{!B}$}،
  \LR{$\Struct{M}$} يا کيڼه يا ښۍ څانګه پوره
  کوي۔
\end{enumerate}
\end{proof}
\olpseganchor{OLP-0109-B018}{end}

\olpseganchor{OLP-0109-B019}{start}
\begin{prob}
د~\olref[fol][tab][sou]{thm:tableau-soundness} ثبوت بشپړ کړئ۔
\end{prob}
\olpseganchor{OLP-0109-B019}{end}



\olpseganchor{OLP-0109-B021}{start}
\begin{cor}
\ollabel{cor:weak-soundness}
که \LR{$\Proves !A$} وي، نو \LR{$!A$} معتبر دے۔
\end{cor}
\olpseganchor{OLP-0109-B021}{end}

\olpseganchor{OLP-0109-B022}{start}
\begin{cor}
\ollabel{cor:entailment-soundness}
که \LR{$\Gamma \Proves !A$} وي، نو \LR{$\Gamma \Entails !A$}۔
\end{cor}
\olpseganchor{OLP-0109-B022}{end}

\olpseganchor{OLP-0109-B023}{start}
\begin{proof}
  که \LR{$\Gamma \Proves !A$} وي، نو د ځينو
  \LR{$!B_1$}، \dots، \LR{$!B_n \in \Gamma$} دپاره
  \LR{$\{\sFmla{\False}{!A}, \sFmla{\True}{!B_1}, \dots,
  \sFmla{\True}{!B_n}\}$} يو تړلے تابلو لري۔ د
  \olref{thm:tableau-soundness} له مخې، هر
  جوړښت~\LR{$\Struct{M}$}
  يا کوم \LR{$!B_i$} ناسموي يا \LR{$!A$} رښتيا کوي۔ نو که
  \LR{$\Sat{M}{\Gamma}$} وي، بيا
  \LR{$\Sat{M}{!A}$} هم دے۔
\end{proof}
\olpseganchor{OLP-0109-B023}{end}

\olpseganchor{OLP-0109-B024}{start}
\begin{cor}
\ollabel{cor:consistency-soundness}
که \LR{$\Gamma$} د صدق وړ وي، نو سازګار دے۔
\end{cor}
\olpseganchor{OLP-0109-B024}{end}

\olpseganchor{OLP-0109-B025}{start}
\begin{proof}
مقابل عکس ثابتوو۔ فرض کړئ \LR{$\Gamma$} سازګار نۀ دے۔ بيا
\LR{$!B_1$}، \dots، \LR{$!B_n \in \Gamma$} شته او د
\LR{$\{\sFmla{\True}{!B_1}, \dots, \sFmla{\True}{!B_n}\}$} دپاره يو
تړلے تابلو شته۔ د \olref{thm:tableau-soundness} له مخې، داسې
\LR{$\Struct{M}$} نشته چې د ټولو
\LR{$i=1$}، \dots،~\LR{$n$} دپاره \LR{$\Sat{M}{!B_i}$} وي۔
خو بيا \LR{$\Gamma$} د صدق وړ نۀ دے۔
\end{proof}
\olpseganchor{OLP-0109-B025}{end}

\olpseganchor{OLP-0110-B004}{start}
\olfileid{fol}{tab}{ide}
\olpseganchor{OLP-0110-B004}{end}

\olpseganchor{OLP-0110-B005}{start}
\olsection{له عينيت سره تابلوګانې}
\olpseganchor{OLP-0110-B005}{end}

\olpseganchor{OLP-0110-B006}{start}
له عينيت سره تابلوګانې د استنتاج اضافي قاعدې غواړي۔
د~\LR{$\eq$} قاعدې دا دي (\LR{$t$}، \LR{$t_1$} او \LR{$t_2$} تړلي ترمونه دي):
\olpseganchor{OLP-0110-B006}{end}

\olpseganchor{OLP-0110-B007}{start}
\begin{defish}
\AxiomC{}
\RightLabel{\LR{$\eq$}}
\UnaryInfC{\sFmla{\True}{\eq[t][t]}}
\DisplayProof
\hfill
\AxiomC{\sFmla{\True}{\eq[t_1][t_2]}}
\noLine
\UnaryInfC{\sFmla{\True}{!A(t_1)}}
\RightLabel{\LR{$\TRule{\True}{\eq}$}}
\UnaryInfC{\sFmla{\True}{!A(t_2)}}
\DisplayProof
\hfill
\AxiomC{\sFmla{\True}{\eq[t_1][t_2]}}
\noLine
\UnaryInfC{\sFmla{\False}{!A(t_1)}}
\RightLabel{\LR{$\TRule{\False}{\eq}$}}
\UnaryInfC{\sFmla{\False}{!A(t_2)}}
\DisplayProof
\end{defish}
پام وکړئ چې له نورو ټولو قاعدو سره په توپير کښې،
\LR{$\TRule{\True}{\eq}$} او \LR{$\TRule{\False}{\eq}$} غواړي چې \emph{دوه}
نښه لرونکي فارمولونه له وړاندې پر څانګه وي؛ يعنې دواړه
\LR{$\sFmla{\True}{\eq[t_1][t_2]}$} او \LR{$\sFmla{S}{!A(t_1)}$}۔
\olpseganchor{OLP-0110-B007}{end}

\olpseganchor{OLP-0110-B008}{start}
\begin{ex}
که \LR{$s$} او \LR{$t$} تړلي ترمونه وي، نو
\LR{$\eq[s][t], !A(s) \Proves !A(t)$}:
\begin{oltableau}
  [\sFmla{\False}{\formula{A}(t)}, just = \TAss
    [\sFmla{\True}{\eq[s][t]}, just = \TAss
      [\sFmla{\True}{\formula{A}(s)}, just = \TAss
        [\sFmla{\True}{\formula{A}(t)}, just={\TRule{\True}{\eq}[2, 3]}, close]
      ]
    ]
  ]
\end{oltableau}
کېدے شي دا د يو شان څيزونو د تعويض د اصل يا د لايبنېڅ د قانون په
نوم درته اشنا وي۔
\olpseganchor{OLP-0110-B008}{end}

\olpseganchor{OLP-0110-B009}{start}
تابلوګانې ثابتوي چې \LR{$\eq$} متناظره ده، يعنې
\LR{$\eq[s_1][s_2] \Proves \eq[s_2][s_1]$}:
\begin{oltableau}
  [\sFmla{\False}{\eq[s_2][s_1]}, just = \TAss
    [\sFmla{\True}{\eq[s_1][s_2]}, just = \TAss
      [\sFmla{\True}{\eq[s_1][s_1]}, just = {\LR{$\eq$}}
        [\sFmla{\True}{\eq[s_2][s_1]}, just = {\TRule{\True}{\eq}[2, 3]}, close]
      ]
    ]
  ]
\end{oltableau}
دلته کرښه~\LR{$2$} د~\LR{$\TRule{\True}{\eq}$} لومړنے اړين
فارمول~\LR{$\sFmla{\True}{\eq[s_1][s_2]}$} دے۔ کرښه~\LR{$3$} دويم
اړين فارمول دے، چې د~\LR{$\sFmla{\True}{!A(s_1)}$} بڼه لري---\LR{$!A(x)$}
د~\LR{$\eq[x][s_1]$} په توګه وګڼئ؛ بيا \LR{$!A(s_1)$}،
\LR{$\eq[s_1][s_1]$} او \LR{$!A(s_2)$}، \LR{$\eq[s_2][s_1]$} کېږي۔
\olpseganchor{OLP-0110-B009}{end}

\olpseganchor{OLP-0110-B010}{start}
تابلوګان دا هم ثابتوي چې \LR{$\eq$} انتقالي ده، يعنې
\LR{$\eq[s_1][s_2], \eq[s_2][s_3] \Proves \eq[s_1][s_3]$}:
\begin{oltableau}
  [\sFmla{\False}{\eq[s_1][s_3]}, just = \TAss
    [\sFmla{\True}{\eq[s_1][s_2]}, just = \TAss
      [\sFmla{\True}{\eq[s_2][s_3]}, just = \TAss
        [\sFmla{\True}{\eq[s_1][s_3]}, just = {\TRule{\True}{\eq}[3, 2]}, close]
      ]
    ]
  ]
\end{oltableau}
په دې تابلو کښې د~\LR{$\TRule{\True}{\eq}$} لومړنے اړين
فارمول کرښه~\LR{$3$}، \LR{$\sFmla{\True}{\eq[s_2][s_3]}$} ده
(\LR{$s_2$} د~\LR{$t_1$} او \LR{$s_3$} د~\LR{$t_2$} ونډه لري)۔ دويم اړين فارمول چې
د~\LR{$\sFmla{\True}{!A(s_2)}$} بڼه لري، کرښه~\LR{$2$} ده۔ دلته \LR{$!A(x)$}
د~\LR{$\eq[s_1][x]$} په توګه وګڼئ؛ په دې سره \LR{$!A(s_2)$}،
\LR{$\eq[s_1][s_2]$} (يعنې کرښه~\LR{$2$}) او \LR{$!A(s_3)$} په پايله کښې
فارمول~\LR{$\eq[s_1][s_3]$} کېږي۔
\end{ex}
\olpseganchor{OLP-0110-B010}{end}

\olpseganchor{OLP-0110-B011}{start}
\begin{prob}
د لاندې دپاره تړلي تابلوګانې ورکړئ:
\begin{enumerate}
\item \LR{$\sFmla{\False}{\lforall[x][\lforall[y][((x = y \land !A(x))
      \lif !A(y))]]}$}
\item \LR{$\sFmla{\False}{\lexists[x][(!A(x) \land
    \lforall[y][(!A(y) \lif y = x)])]}$},\\
  \LR{$\sFmla{\True}{\lexists[x][!A(x)] \land
    \lforall[y][\lforall[z][((!A(y) \land !A(z)) \lif y = z)]]}$}
\end{enumerate}
\end{prob}
\olpseganchor{OLP-0110-B011}{end}

\olpseganchor{OLP-0111-B004}{start}
\olfileid{fol}{tab}{sid}
\olpseganchor{OLP-0111-B004}{end}

\olpseganchor{OLP-0111-B005}{start}
\olsection{له عينيت سره سموالے}
\olpseganchor{OLP-0111-B005}{end}

\olpseganchor{OLP-0111-B006}{start}
\begin{prop}
د عينيت له قاعدو سره تابلوګانې سم دي: هېڅ تړلے تابلو
د صدق وړ نۀ دے۔
\end{prop}
\olpseganchor{OLP-0111-B006}{end}

\olpseganchor{OLP-0111-B007}{start}
\begin{proof}
يوازې بايد د پخوا په شان وښيو چې که يو تابلو د صدق وړ
څانګه ولري، نو د~\LR{$\eq$} له قاعدو څخه د يوې په پلي کولو ترلاسه شوې
څانګه هم د صدق وړ ده۔ \LR{$\Gamma$} پر څانګه د نښه لرونکو فارمولونه
سټ وبولئ، او \LR{$\Struct{M}$} داسې جوړښت وبولئ چې \LR{$\Gamma$}
پوره کوي۔
\olpseganchor{OLP-0111-B007}{end}

\olpseganchor{OLP-0111-B008}{start}
فرض کړئ څانګه د~\LR{$\eq$} په کارولو، يعنې د نښه لرونکي
فارمول~\LR{$\sFmla{\True}{\eq[t][t]}$} په زياتولو پراخېږي۔ په
څرګند ډول \LR{$\Sat{M}{\eq[t][t]}$}؛ نو \LR{$\Struct{M}$}،
\LR{$\Gamma \cup \{\sFmla{\True}{\eq[t][t]}\}$} هم پوره کوي۔
\olpseganchor{OLP-0111-B008}{end}

\olpseganchor{OLP-0111-B009}{start}
که څانګه د~\LR{$\TRule{\True}{\eq}$} په کارولو پراخېږي، موږ يو نښه
لرونکے فارمول~\LR{$\sFmla{\True}{!A(t_2)}$} ورزياتوو، خو \LR{$\Gamma$}
دواړه~\LR{$\sFmla{\True}{\eq[t_1][t_2]}$} او
\LR{$\sFmla{\True}{!A(t_1)}$} لري۔
نو \LR{$\Sat{M}{\eq[t_1][t_2]}$} او \LR{$\Sat{M}{!A(t_1)}$} لرو۔ داسې
متغير-ټاکنه~\LR{$s$} واخلئ چې \LR{$s(x) = \Value{t_1}{M}$}۔ د
\olref[syn][ass]{prop:sentence-sat-true} له مخې
\LR{$\Sat{M}{!A(t_1)}[s]$}۔ ځکه \LR{$\varAssign{s}{s}{x}$}، د
\olref[syn][ext]{prop:ext-formulas} له مخې \LR{$\Sat{M}{!A(x)}[s]$}۔
ځکه \LR{$\Sat{M}{\eq[t_1][t_2]}$}، نو
\LR{$\Value{t_1}{M} = \Value{t_2}{M}$}، او ځکه
\LR{$s(x) = \Value{t_2}{M}$}۔ د
\olref[syn][ext]{prop:ext-formulas} په بيا پلي کولو
\LR{$\Sat{M}{!A(t_2)}[s]$} هم لرو۔ د
\olref[syn][ass]{prop:sentence-sat-true} له مخې
\LR{$\Sat{M}{!A(t_2)}$}۔ د~\LR{$\TRule{\False}{\eq}$} حالت هم په همدې ډول
حل کېږي۔
\end{proof}
\olpseganchor{OLP-0111-B009}{end}

\olpseganchor{OLP-0112-B004}{start}
\olchapter{fol}{axd}{بديهي اشتقاقونه}
\olpseganchor{OLP-0112-B004}{end}

\olpseganchor{OLP-0112-B005}{start}
\begin{editorial}
  لا تر اوسه دا هڅه نۀ ده شوې چې ډاډ ترلاسه شي چې د دې څپرکي مواد
  د هغو بېلابېلو ټاګونو درناوے کوي چې ښيي کوم نښلوونکي او کميت
  ټاکونکي بنسټيز يا تعريف شوي دي: له~\LR{$\liff$} پرته، چې تعريف شوے
  ګڼل کېږي، ټول بنسټيز ګڼل شوي دي۔ که د لومړۍ درجې منطق ټاګ رښتونے
  وي، له کميت ټاکونکو سره بڼه جوړوو؛ که نۀ وي، بې له هغو بڼه جوړوو۔
\end{editorial}
\olpseganchor{OLP-0112-B005}{end}





\olpseganchor{OLP-0112-B008}{start}
%
  

\olpseganchor{OLP-0112-B008}{end}



\olpseganchor{OLP-0112-B010}{start}
%
  

\olpseganchor{OLP-0112-B010}{end}





\olpseganchor{OLP-0112-B013}{start}
%
  

\olpseganchor{OLP-0112-B013}{end}





\olpseganchor{OLP-0112-B016}{start}
%
  

\olpseganchor{OLP-0112-B016}{end}

\olpseganchor{OLP-0113-B004}{start}
\olfileid{fol}{axd}{rul}
\olpseganchor{OLP-0113-B004}{end}

\olpseganchor{OLP-0113-B005}{start}
\olsection{قاعدې او اشتقاقونه}
\olpseganchor{OLP-0113-B005}{end}

\olpseganchor{OLP-0113-B006}{start}
\begin{explain}
  بديهي اشتقاقونه ښايي د منطق د اشتقاق تر ټولو ساده
  نظام وي۔ اشتقاق يوازې د فارمولونه يوه لړۍ ده۔ د
  اشتقاق ګڼل کېدو دپاره، په لړۍ کښې هر فارمول بايد
  يا د کوم بديهي اصل يوه بېلګه وي، يا د استنتاج د کومې قاعدې له مخې
  له يوه يا زياتو هغو فارمولونه څخه راشي چې په لړۍ کښې ترې وړاندې
  دي۔ اشتقاق خپل وروستے فارمول، اشتقاق کوي۔
\end{explain}
\olpseganchor{OLP-0113-B006}{end}

\olpseganchor{OLP-0113-B007}{start}
\begin{defn}[د اشتقاق وړتيا]
که \LR{$\Gamma$} د~\LR{$\Lang L$} د فارمولونه يو سټ وي، نو له~\LR{$\Gamma$}
څخه يو \emph{اشتقاق} د فارمولونه يوه متناهي لړۍ
\LR{$!A_1$}، \dots،~\LR{$!A_n$} ده، چې د هر \LR{$i \le n$} دپاره له لاندې څخه
يو شرط پوره کېږي:
\begin{enumerate}
\item \LR{$!A_i \in \Gamma$}؛ يا
\item \LR{$!A_i$} يو بديهي اصل دے؛ يا
\item \LR{$!A_i$} د استنتاج د کومې قاعدې له مخې له کوم \LR{$!A_j$} (او
  \LR{$!A_k$}) څخه راځي، چې \LR{$j < i$} (او \LR{$k < i$})۔
\end{enumerate}
\end{defn}
\olpseganchor{OLP-0113-B007}{end}

\olpseganchor{OLP-0113-B008}{start}
دا چې سم اشتقاق څه ګڼل کېږي، په دې پورې اړه لري چې موږ
کومې استنتاجي قاعدې منو (او البته کوم څيزونه بديهي اصول ګڼو)۔
استنتاجي قاعده يوه که-نو وينا ده چې راته وايي د ځينو شرطونو لاندې
په اشتقاق کښې يو ګام~\LR{$!A_i$} سم استنتاجي ګام دے۔
\olpseganchor{OLP-0113-B008}{end}

\olpseganchor{OLP-0113-B009}{start}
\begin{defn}[د استنتاج قاعده]
د \emph{استنتاج قاعده} يو بس شرط ورکوي چې له~\LR{$\Gamma$} څخه په
اشتقاق کښې څه شی سم استنتاجي ګام ګڼل کېږي۔
\end{defn}
\olpseganchor{OLP-0113-B009}{end}

\olpseganchor{OLP-0113-B010}{start}
د بېلګې په توګه، ځکه چې هره يو-غړيزه لړۍ~\LR{$!A$} چې
\LR{$!A \in \Gamma$} وي، په څرګند ډول اشتقاق ګڼل کېږي، لاندې
ډېره ساده استنتاجي قاعده کېدے شي:
\begin{quote}
  که \LR{$!A \in \Gamma$} وي، نو \LR{$!A$} له~\LR{$\Gamma$} څخه په هر
  اشتقاق کښې تل سم استنتاجي ګام دے۔
\end{quote}
همدارنګه که \LR{$!A$} له بديهي اصولو څخه يو وي، نو يوازې \LR{$!A$} هم
اشتقاق دے؛ ځکه دا هم يوه استنتاجي قاعده ده:
\begin{quote}
  که \LR{$!A$} بديهي اصل وي، نو \LR{$!A$} سم استنتاجي ګام دے۔
\end{quote}
خبره هغه وخت لا په زړه پورې شي چې استنتاجي قاعده هغو فارمولونه
ته رجوع کوي چې تر څېړل شوي ګام وړاندې راغلي وي۔ لاندې قاعده
\emph{مودوس پوننس} بلل کېږي:
\begin{quote}
  که \LR{$!B \lif !A$} او \LR{$!B$} په اشتقاق کښې وړاندې راغلي وي،
  نو~\LR{$!A$} سم استنتاجي ګام دے۔
\end{quote}
که همدا يوازينۍ استنتاجي قاعده وي، نو د اشتقاق پورته تعريف
مو داسې کېږي: \LR{$!A_1$}، \dots،~\LR{$!A_n$} هله او يوازې هله
اشتقاق دے چې د هر \LR{$i \le n$} دپاره له لاندې څخه يو شرط
پوره کېږي:
\begin{enumerate}
\item \LR{$!A_i \in \Gamma$}؛ يا
\item \LR{$!A_i$} يو بديهي اصل دے؛ يا
\item د کوم \LR{$j < i$} دپاره \LR{$!A_j$}، \LR{$!B \lif !A_i$} وي، او د کوم
  \LR{$k < i$} دپاره \LR{$!A_k$}،~\LR{$!B$} وي۔
\end{enumerate}
وروستے بند وايي چې \LR{$!A_i$} له~\LR{$!A_j$} (\LR{$!B \lif !A_i$}) او
\LR{$!A_k$} (\LR{$!B$}) څخه د مودوس پوننس په وسيله راځي۔ که له~\LR{$1$} څخه
تر~\LR{$n$} پورې تېرېدے شو، او هر ځل داسې فارمول~\LR{$!A_i$} ومومو
چې يا په~\LR{$\Gamma$} کښې وي، يا بديهي اصل وي، يا استنتاجي قاعده راته
وايي چې سم استنتاجي ګام دے، نو ټوله لړۍ سم اشتقاق ګڼل کېږي۔
\olpseganchor{OLP-0113-B010}{end}

\olpseganchor{OLP-0113-B011}{start}
\begin{defn}[د اشتقاق وړتيا]
يو فارمول~\LR{$!A$} هله له~\LR{$\Gamma$} څخه \emph{د اشتقاق وړ} دے،
چې \LR{$\Gamma \Proves !A$} ليکو، چې له~\LR{$\Gamma$} څخه داسې
اشتقاق وي چې په~\LR{$!A$} پاے ته رسېږي۔
\end{defn}
\olpseganchor{OLP-0113-B011}{end}

\olpseganchor{OLP-0113-B012}{start}
\begin{defn}[قضيې]
فارمول~\LR{$!A$} هله \emph{قضيه} ده چې له تش سټ څخه د~\LR{$!A$} يو
اشتقاق وي۔ که \LR{$!A$} قضيه وي \LR{$\Proves !A$}، او که نۀ وي
\LR{$\Proves/ !A$} ليکو۔
\end{defn}
\olpseganchor{OLP-0113-B012}{end}

\olpseganchor{OLP-0114-B004}{start}
\olfileid{fol}{axd}{prp}
\olpseganchor{OLP-0114-B004}{end}

\olpseganchor{OLP-0114-B005}{start}
\olsection{د قضيې نښلوونکو دپاره بديهي اصول او قاعدې}
\olpseganchor{OLP-0114-B005}{end}

\olpseganchor{OLP-0114-B006}{start}
\begin{defn}[بديهي اصول]
د قضيې نښلوونکو د بديهي اصولو سټ~\LR{$\PAx$} د لاندې ټولو بڼو
فارمولونه لري:
\begin{align}
  & (!A \land !B) \lif !A \ollabel{ax:land1}\\
  & (!A \land !B) \lif !B \ollabel{ax:land2}\\
  & !A \lif (!B \lif (!A \land !B)) \ollabel{ax:land3}\\
  & !A \lif (!A \lor !B) \ollabel{ax:lor1}\\
  & !A \lif (!B \lor !A) \ollabel{ax:lor2}\\
  & (!A \lif !C) \lif ((!B \lif !C) \lif ((!A \lor !B) \lif !C)) \ollabel{ax:lor3}\\
  & !A \lif (!B \lif !A) \ollabel{ax:lif1}\\
  & (!A \lif (!B \lif !C)) \lif ((!A \lif !B) \lif (!A \lif !C)) \ollabel{ax:lif2}\\
  & (!A \lif !B) \lif ((!A \lif \lnot !B) \lif \lnot !A) \ollabel{ax:lnot1}\\
  & \lnot !A \lif (!A \lif !B) \ollabel{ax:lnot2}\\
  & \ltrue \ollabel{ax:ltrue}\\
  & \lfalse \lif !A \ollabel{ax:lfalse1}\\
  & (!A \lif \lfalse) \lif \lnot !A \ollabel{ax:lfalse2}\\
  & \lnot\lnot !A \lif !A \ollabel{ax:dne}
\end{align}
\end{defn}
\olpseganchor{OLP-0114-B006}{end}

\olpseganchor{OLP-0114-B007}{start}
\begin{defn}[مودوس پوننس]
  که \LR{$!B$} او \LR{$!B \lif !A$} له وړاندې په اشتقاق کښې راغلي
  وي، نو \LR{$!A$} سم استنتاجي ګام دے۔
\end{defn}
\olpseganchor{OLP-0114-B007}{end}

\olpseganchor{OLP-0114-B008}{start}
د مودوس پوننس قاعده به په ``\MP'' لنډوو۔
\olpseganchor{OLP-0114-B008}{end}

\olpseganchor{OLP-0115-B004}{start}
\olfileid{fol}{axd}{qua}
\olpseganchor{OLP-0115-B004}{end}

\olpseganchor{OLP-0115-B005}{start}
\olsection{د کميت ټاکونکو دپاره بديهي اصول او قاعدې}
\olpseganchor{OLP-0115-B005}{end}

\olpseganchor{OLP-0115-B006}{start}
\begin{defn}[د کميت ټاکونکو بديهي اصول]
پر کميت ټاکونکو واکمن \emph{بديهي اصول} د لاندې ټولو بڼو بېلګې دي:
\begin{align}
\ollabel{ax:q1} & \lforall[x][!B] \lif !B(t), \\
\ollabel{ax:q2} & !B(t) \lif \lexists[x][!B].
\end{align}
د هر تړلي ترم~\LR{$t$} دپاره۔
\end{defn}
\olpseganchor{OLP-0115-B006}{end}

\olpseganchor{OLP-0115-B007}{start}
\begin{defn}[د کميت ټاکونکو قاعدې]
\item که \LR{$!B \lif !A(a)$} له وړاندې په اشتقاق کښې راغلے وي
  او \LR{$a$} په~\LR{$\Gamma$} يا~\LR{$!B$} کښې نۀ راځي، نو
  \LR{$!B \lif \lforall[x][!A(x)]$} سم استنتاجي ګام دے۔
\item که \LR{$!A(a) \lif !B$} له وړاندې په اشتقاق کښې راغلے وي
  او \LR{$a$} په~\LR{$\Gamma$} يا~\LR{$!B$} کښې نۀ راځي، نو
  \LR{$\lexists[x][!A(x)] \lif !B$} سم استنتاجي ګام دے۔
\end{defn}
\olpseganchor{OLP-0115-B007}{end}

\olpseganchor{OLP-0115-B008}{start}
له دغو دواړو هره يوه به په ``\QR'' لنډوو۔
\olpseganchor{OLP-0115-B008}{end}

\olpseganchor{OLP-0116-B004}{start}
\olfileid{fol}{axd}{pro}
\olpseganchor{OLP-0116-B004}{end}

\olpseganchor{OLP-0116-B005}{start}
\olsection{د اشتقاقونه بېلګې}
\olpseganchor{OLP-0116-B005}{end}


\olpseganchor{OLP-0116-B006}{start}
\begin{ex}
فرض کړئ غواړو \LR{$(\lnot !D \lor !E) \lif (!D \lif !E)$} ثابت کړو۔
دا په څرګند ډول زموږ د هېڅ بديهي اصل بېلګه نۀ ده، نو د هغې د
اشتقاق کولو دپاره د~\MP{} قاعده کارول پکار دي۔ زموږ يوازينۍ
قاعده مودوس پوننس ده، چې \LR{$!A$} او \LR{$!A \lif !B$} راکړل شي، د~\LR{$!B$}
توجيه راکوي۔ يوه تګلاره دا ده چې \olref[prp]{ax:lor3} داسې وکاروو
چې \LR{$!A$}، \LR{$\lnot !D$}؛ \LR{$!B$}، \LR{$!E$}؛ او \LR{$!C$}، \LR{$!D \lif !E$} وي؛ يعنې
دا بېلګه:
\[
(\lnot !D \lif (!D \lif !E)) \lif ((!E \lif (!D \lif !E)) \lif ((\lnot
!D \lor !E) \lif (!D \lif !E))).
\]
ولې؟ د مودوس پوننس دوه تطبيقونه وروستۍ برخه راکوي، او همدا مو
غوښتې ده۔ په اسانۍ وينو چې \LR{$\lnot !D \lif (!D \lif !E)$} د
\olref[prp]{ax:lnot2}، او \LR{$!E \lif (!D \lif !E)$} د
\olref[prp]{ax:lif1} بېلګه ده۔ نو زموږ اشتقاق دا دے:
\begin{derivation}
  1. &  \LR{$\lnot !D \lif (!D \lif !E)$} & \olref[prp]{ax:lnot2} \\
  2. & \LR{$(\lnot !D \lif (!D \lif !E)) \lif {}$}\\
  &\qquad \LR{$((!E \lif (!D \lif !E)) \lif ((\lnot !D \lor !E) \lif (!D \lif !E)))$} & \olref[prp]{ax:lor3}\\
  3. & \LR{$(!E \lif (!D \lif !E)) \lif ((\lnot !D \lor !E) \lif (!D \lif !E))$} & 1, 2, \MP\\
  4. & \LR{$!E \lif (!D \lif !E)$} & \olref[prp]{ax:lif1}\\
  5. & \LR{$(\lnot !D \lor !E) \lif (!D \lif !E)$} & 3, 4, \MP
\end{derivation}
\end{ex}
\olpseganchor{OLP-0116-B006}{end}

\olpseganchor{OLP-0116-B007}{start}
\begin{ex}\ollabel{ex:identity}
راځئ د~\LR{$!D \lif !D$} يو اشتقاق ولټوو۔ دا د بديهي اصل بېلګه
نۀ ده، نو په~\MP{} ئې اشتقاق کول پکار دي۔
\olref[prp]{ax:lif1} د~\LR{$!A \lif !B$} په بڼه يو بديهي اصل دے چې
\MP{} پرې پلي کېدے شي۔ د ګټې دپاره، البته، هغه~\LR{$!B$} چې \MP{} ئې په
دې حالت کښې د سم ګام په توګه توجيه کوي بايد~\LR{$!D \lif !D$} وي؛ ځکه
همدا مو غوښتل چې اشتقاق ئې کړو۔ په دې سره \LR{$!A$} هم بايد \LR{$!D$} وي؛
يعنې د \olref[prp]{ax:lif1} دې بېلګې ته کتلے شو:
\[
!D \lif (!D \lif !D)
\]
د~\MP{} د پلي کولو دپاره به د دويمې اړوندې مقدمې، يعنې~\LR{$!A$}،
توجيه هم پکار وي۔ خو زموږ په حالت کښې هغه~\LR{$!D$} دے، او \LR{$!D$} په
يوازې ځان اشتقاق کولے نۀ شو۔ نو بله تګلاره پکار ده۔
\olpseganchor{OLP-0116-B007}{end}

\olpseganchor{OLP-0116-B008}{start}
بل بديهي اصل چې يوازې~\LR{$\lif$} لري \olref[prp]{ax:lif2} دے، يعنې
\[
(!A \lif (!B \lif !C)) \lif ((!A \lif !B) \lif (!A \lif !C))
\]
د~\MP{} په دوه ځله پلي کولو وروستي دننني شرطي ته رسېدے شو۔ بيا دا
معنٰی لري چې د \olref[prp]{ax:lif2} داسې بېلګه غواړو چې
\LR{$!A \lif !C$} پکې \LR{$!D \lif !D$}، يعنې زموږ موخه فارمول وي۔
نو \LR{$!A$} او \LR{$!C$} دواړه~\LR{$!D$} دي۔ \LR{$!B$} څنګه وټاکو چې دواړه
\LR{$!A \lif (!B \lif !C)$} او \LR{$!A \lif !B$}، يعنې زموږ په حالت کښې
\LR{$!D \lif (!B \lif !D)$} او \LR{$!D \lif !B$}، هم د اشتقاق وړ وي؟
له دغو لومړنے، هر~\LR{$!B$} چې وټاکو، لا د \olref[prp]{ax:lif1} بېلګه
ده۔ او که \LR{$!B$}، \LR{$(!D \lif !D)$} وي، نو \LR{$!D \lif !B$} به هم د
\olref[prp]{ax:lif1} بله بېلګه وي۔ نو زموږ اشتقاق دا دے:
\begin{derivation}
  1. & \LR{$!D \lif ((!D \lif !D) \lif !D)$} & \olref[prp]{ax:lif1}\\
  2. & \LR{$(!D \lif ((!D \lif !D) \lif !D)) \lif {}$}\\
  & \qquad \LR{$((!D \lif (!D \lif !D)) \lif (!D \lif !D))$} & \olref[prp]{ax:lif2}\\
  3. & \LR{$(!D \lif (!D \lif !D)) \lif (!D \lif !D)$} & 1, 2, \MP\\
  4. & \LR{$!D \lif (!D \lif !D)$} & \olref[prp]{ax:lif1}\\
  5. & \LR{$!D \lif !D$} & 3, 4, \MP
\end{derivation}
\end{ex}
\olpseganchor{OLP-0116-B008}{end}

\olpseganchor{OLP-0116-B009}{start}
\begin{ex}\ollabel{ex:chain}
کله کله غواړو وښيو چې له ځينو نورو فارمولونه~\LR{$\Gamma$} څخه د کوم
فارمول يو اشتقاق شته۔ مثلاً راځئ وښيو چې له
\LR{$\Gamma = \{!A \lif !B, !B \lif !C\}$} څخه
\LR{$!A \lif !C$}، اشتقاق کولے شو۔
\begin{derivation}
  1. & \LR{$!A \lif !B$} & \Hyp\\
  2. & \LR{$!B \lif !C$} & \Hyp\\
  3. & \LR{$(!B \lif !C) \lif (!A \lif (!B \lif !C))$} & \olref[prp]{ax:lif1} \\
  4. & \LR{$!A \lif (!B \lif !C)$} & 2, 3, \MP\\
  5. & \LR{$(!A \lif (!B \lif !C)) \lif {}$}\\
  & \qquad \LR{$((!A \lif !B) \lif (!A \lif !C))$} & \olref[prp]{ax:lif2}\\
  6. & \LR{$((!A \lif !B) \lif (!A \lif !C))$} & 4, 5, \MP\\
  7. & \LR{$!A \lif !C$} & 1, 6, \MP
\end{derivation}
هغه کرښې چې ``\Hyp'' (د ``فرضیې'' دپاره) پرې ليکل شوي دي، ښيي چې
په هغې کرښه فارمول د~\LR{$\Gamma$} يو غړے دے۔
\end{ex}
\olpseganchor{OLP-0116-B009}{end}

\olpseganchor{OLP-0116-B010}{start}
\begin{prop}
  \ollabel{prop:chain} که \LR{$\Gamma \Proves !A \lif !B$} او \LR{$\Gamma
  \Proves !B \lif !C$} وي، نو \LR{$\Gamma \Proves !A \lif !C$}
\end{prop}
\olpseganchor{OLP-0116-B010}{end}

\olpseganchor{OLP-0116-B011}{start}
\begin{proof}
  فرض کړئ \LR{$\Gamma \Proves !A \lif !B$} او \LR{$\Gamma \Proves !B \lif
  !C$}۔ بيا له~\LR{$\Gamma$} څخه د~\LR{$!A \lif !B$} يو اشتقاق شته؛
  او له~\LR{$\Gamma$} څخه د~\LR{$!B \lif !C$} يو اشتقاق هم شته۔ د
  هغو په يو بل پسې نښلولو ئې په يو اشتقاق کښې يو ځاے کړئ۔
  اوس د مخکښې بېلګې د~اشتقاق ۳--۷ کرښې ورزياتې کړئ۔ دا له~\LR{$\Gamma$} څخه د
  \LR{$!A \lif !C$} يو اشتقاق دے---او همدا د نوي
  اشتقاق وروستۍ کرښه ده۔ پام وکړئ چې د ۴مې او ۷مې کرښې
  توجيه لا هم سمه پاتې کېږي که د ۱مې کرښې حواله د~\LR{$!A \lif !B$}
  د~اشتقاق د وروستۍ کرښې په حواله، او د ۲مې کرښې حواله د
  \LR{$!B \lif !C$} د~اشتقاق د وروستۍ کرښې په حواله بدله شي۔
\end{proof}
\olpseganchor{OLP-0116-B011}{end}

\olpseganchor{OLP-0116-B012}{start}
\begin{prob} 
له بديهي اصولو څخه د اشتقاقونه په وړاندې کولو وښيئ چې لاندې
خبرې سمې دي:
\begin{enumerate} 
\item \LR{$(!A \land !B) \lif (!B \land !A)$}
\item \LR{$((!A \land !B) \lif !C) \lif (!A \lif (!B \lif !C))$}
\item \LR{$\lnot(!A \lor !B) \lif \lnot !A$}
\end{enumerate} 
\end{prob}
\olpseganchor{OLP-0116-B012}{end}

\olpseganchor{OLP-0117-B004}{start}
\olfileid{fol}{axd}{prq}
\olpseganchor{OLP-0117-B004}{end}

\olpseganchor{OLP-0117-B005}{start}
\olsection{له کميت ټاکونکو سره اشتقاقونه}
\olpseganchor{OLP-0117-B005}{end}

\olpseganchor{OLP-0117-B006}{start}
\begin{ex}
راځئ د~\LR{$(\lforall[x][!A(x)] \land
\lforall[y][!B(y)]) \lif \lforall[x][(!A(x) \land !B(x))]$} يو
اشتقاق ورکړو۔
\olpseganchor{OLP-0117-B006}{end}

\olpseganchor{OLP-0117-B007}{start}
لومړے پام وکړئ چې
\begin{align*}
  (\lforall[x][!A(x)] \land \lforall[y][!B(y)]) & \lif \lforall[x][!A(x)]\\
  \intertext{\RL{د \olref[prp]{ax:land1} يوه بېلګه ده، او}}
  \lforall[x][!A(x)] & \lif !A(a) \\
  \intertext{\RL{د \olref[qua]{ax:q1} بېلګه ده۔ نو د \olref[pro]{prop:chain} له مخې پوهېږو چې}}
  (\lforall[x][!A(x)] \land \lforall[y][!B(y)]) & \lif !A(a) \\
  \intertext{\RL{د اشتقاق وړ دے۔ همدارنګه، ځکه}}
  (\lforall[x][!A(x)] \land \lforall[y][!B(y)]) & \lif \lforall[y][!B(y)] \qquad\text{\RL{او}}\\
    \lforall[y][!B(y)] & \lif !B(a)\\
    \intertext{\RL{په ترتيب د \olref[prp]{ax:land2} او \olref[qua]{ax:q1} بېلګې دي،}}
    (\lforall[x][!A(x)] \land \lforall[y][!B(y)]) & \lif !B(a)\\
    \intertext{\RL{د \olref[pro]{prop:chain} په وسيله د ثبوت وړ دے۔ د \olref[prp]{ax:land3} د يوې مناسبې بېلګې او د~\MP{} د دوو تطبيقونو په کارولو وينو چې}}
    (\lforall[x][!A(x)] \land \lforall[y][!B(y)]) & \lif (!A(a) \land !B(a))\\
    \intertext{\RL{د ثبوت وړ دے۔ اوس \QR{} پلي کولے شو او دا ترلاسه کوو:}}
(\lforall[x][!A(x)] \land \lforall[y][!B(y)]) & \lif \lforall[x][(!A(x) \land !B(x))].
\end{align*}
\end{ex}
\olpseganchor{OLP-0117-B007}{end}

\olpseganchor{OLP-0118-B004}{start}
\olfileid{fol}{axd}{ptn}
\olpseganchor{OLP-0118-B004}{end}

\olpseganchor{OLP-0118-B005}{start}
\olsection{ثبوت-تيوريکي مفاهيم}
\olpseganchor{OLP-0118-B005}{end}

\olpseganchor{OLP-0118-B006}{start}
\begin{explain}
لکه څنګه چې مو يو شمېر مهم معنايي مفاهيم
(اعتبار، معنايي لازميت او د صدق وړتيا)
تعريف کړي دي، اوس ورسره سم \emph{ثبوت-تيوريکي مفاهيم} تعريفوو۔
دا د جوړښتونه په دننه کښې د تړلي فارمولونه پوره کېدو ته په
رجوع نۀ، بلکې د ځينو فارمولونو د اشتقاق وړتيا يا
د اشتقاق نۀ وړتيا ته په رجوع تعريفېږي۔ دا يوه مهمه موندنه وه چې
دغه مفاهيم يو شان راځي۔ د هغو يو شان راتلل د \emph{سموالي} او
\emph{بشپړوالي د قضيو} منځپانګه ده۔
\end{explain}
\olpseganchor{OLP-0118-B006}{end}

\olpseganchor{OLP-0118-B007}{start}
\begin{defn}[د اشتقاق وړتيا]
فارمول~\LR{$!A$} هله له~\LR{$\Gamma$} څخه \emph{د اشتقاق وړ} دے،
چې \LR{$\Gamma \Proves !A$} ليکو، چې له~\LR{$\Gamma$} څخه داسې
اشتقاق وي چې په~\LR{$!A$} پاے ته رسېږي۔
\end{defn}
\olpseganchor{OLP-0118-B007}{end}

\olpseganchor{OLP-0118-B008}{start}
\begin{defn}[قضيې]
فارمول~\LR{$!A$} هله \emph{قضيه} ده چې له تش سټ څخه د~\LR{$!A$} يو
اشتقاق وي۔ که \LR{$!A$} قضيه وي \LR{$\Proves !A$}، او که نۀ وي
\LR{$\Proves/ !A$} ليکو۔
\end{defn}
\olpseganchor{OLP-0118-B008}{end}

\olpseganchor{OLP-0118-B009}{start}
\begin{defn}[سازګاري]
د فارمولونه يو سټ~\LR{$\Gamma$} هله او يوازې هله \emph{سازګار} دے
چې \LR{$\Gamma\Proves/ \lfalse$}؛ که داسې نۀ وي \emph{ناسازګار} دے۔
\end{defn}
\olpseganchor{OLP-0118-B009}{end}

\olpseganchor{OLP-0118-B010}{start}
\begin{prop}[انعکاس]
\ollabel{prop:reflexivity}
که \LR{$!A \in \Gamma$} وي، نو \LR{$\Gamma \Proves !A$}۔
\end{prop}
\olpseganchor{OLP-0118-B010}{end}

\olpseganchor{OLP-0118-B011}{start}
\begin{proof}
  يوازې فارمول~\LR{$!A$} له~\LR{$\Gamma$} څخه د~\LR{$!A$} يو
  اشتقاق دے۔
\end{proof}
\olpseganchor{OLP-0118-B011}{end}

\olpseganchor{OLP-0118-B012}{start}
\begin{prop}[زياتېدنه]
\ollabel{prop:monotonicity}
که \LR{$\Gamma \subseteq \Delta$} او \LR{$\Gamma \Proves !A$} وي، نو
\LR{$\Delta \Proves !A$}۔
\end{prop}
\olpseganchor{OLP-0118-B012}{end}

\olpseganchor{OLP-0118-B013}{start}
\begin{proof}
له~\LR{$\Gamma$} څخه د~\LR{$!A$} هر اشتقاق له~\LR{$\Delta$} څخه هم د~\LR{$!A$}
يو اشتقاق دے۔
\end{proof}
\olpseganchor{OLP-0118-B013}{end}

\olpseganchor{OLP-0118-B014}{start}
\begin{prop}[انتقال]
\ollabel{prop:transitivity}
که \LR{$\Gamma \Proves !A$} او \LR{$\{!A\} \cup \Delta \Proves !B$} وي،
نو \LR{$\Gamma \cup \Delta \Proves !B$}۔
\end{prop}
\olpseganchor{OLP-0118-B014}{end}

\olpseganchor{OLP-0118-B015}{start}
\begin{proof}
  فرض کړئ \LR{$\{!A\} \cup \Delta \Proves !B$}۔ بيا له
  \LR{$\{!A\} \cup \Delta$} څخه يو اشتقاق
  \LR{$!B_1$}، \dots، \LR{$!B_l = !B$} شته۔ په دې اشتقاق کښې ځينې
  ګامونه به د هغې قاعدې له مخې سم وي چې کومې مخکښې کرښې
  \LR{$!B_i = !A$} ته رجوع کوي۔ د فرض له مخې له~\LR{$\Gamma$} څخه د~\LR{$!A$}
  يو اشتقاق شته؛ يعنې داسې اشتقاق
  \LR{$!A_1$}، \dots، \LR{$!A_k = !A$} چې هر \LR{$!A_i$} يا بديهي اصل وي، يا
  د~\LR{$\Gamma$} يو غړے وي، يا د استنتاج د قاعدې له مخې سم
  وي۔ اوس دا لړۍ په پام کښې ونيسئ:
  \[
  !A_1, \dots, !A_k = !A, !B_1, \dots, !B_l = !B.
  \]
  دا له~\LR{$\Gamma \cup \Delta$} څخه د~\LR{$!B$} سم اشتقاق دے؛ ځکه
  هر \LR{$!B_i = !A$} اوس د هماغې قاعدې له مخې توجيه کېږي چې
  \LR{$!A_k = !A$} توجيه کوي۔
\end{proof}
\olpseganchor{OLP-0118-B015}{end}

\olpseganchor{OLP-0118-B016}{start}
پام وکړئ چې په ځانګړي ډول دا معنٰی لري: که
\LR{$\Gamma \Proves !A$} او \LR{$!A \Proves !B$} وي، نو
\LR{$\Gamma \Proves !B$}۔ له دې څخه دا هم راځي چې که
\LR{$!A_1, \dots, !A_n \Proves !B$} او د هر~\LR{$i$} دپاره
\LR{$\Gamma \Proves !A_i$} وي، نو \LR{$\Gamma \Proves !B$}۔
\olpseganchor{OLP-0118-B016}{end}

\olpseganchor{OLP-0118-B017}{start}
\begin{prop}
\ollabel{prop:incons}
\LR{$\Gamma$} هله او يوازې هله ناسازګار دے چې د هر~\LR{$!A$} دپاره
\LR{$\Gamma \Proves !A$} وي۔
\end{prop}
\olpseganchor{OLP-0118-B017}{end}

\olpseganchor{OLP-0118-B018}{start}
\begin{proof}
تمرين۔
\end{proof}
\olpseganchor{OLP-0118-B018}{end}

\olpseganchor{OLP-0118-B019}{start}
\begin{prob}
\olref[fol][axd][ptn]{prop:incons} ثابت کړئ۔
\end{prob}
\olpseganchor{OLP-0118-B019}{end}



\olpseganchor{OLP-0118-B021}{start}
\begin{prop}[تراکم]
\ollabel{prop:proves-compact}
  \begin{enumerate}
  \item که \LR{$\Gamma \Proves !A$} وي، نو داسې متناهي فرعي سټ
    \LR{$\Gamma_0 \subseteq \Gamma$} شته چې \LR{$\Gamma_0 \Proves !A$}۔
  \item که د~\LR{$\Gamma$} هر متناهي فرعي سټ سازګار وي، نو \LR{$\Gamma$}
    سازګار دے۔
  \end{enumerate}
\end{prop}
\olpseganchor{OLP-0118-B021}{end}

\olpseganchor{OLP-0118-B022}{start}
\begin{proof}
  \begin{enumerate}
    \item که \LR{$\Gamma \Proves !A$} وي، نو د فارمولونه داسې متناهي
      لړۍ \LR{$!A_1$}، \dots،~\LR{$!A_n$} شته چې \LR{$!A \ident !A_n$}؛ او هر
      \LR{$!A_i$} يا منطقي بديهي اصل وي، يا د~\LR{$\Gamma$} يو غړے
      وي، يا د مودوس پوننس له مخې له مخکښو فارمولونه څخه راځي۔
      \LR{$\Gamma_0$} د هغو \LR{$!A_i$} سټ ونيسئ چې په~\LR{$\Gamma$} کښې دي۔ بيا
      همدا اشتقاق له~\LR{$\Gamma_0$} څخه هم يو اشتقاق
      دے، نو \LR{$\Gamma_0 \Proves !A$}۔
    \item دا د ځانګړي حالت~\LR{$!A \ident \lfalse$} دپاره د~(۱) مقابل
      عکس دے۔
\end{enumerate}
\end{proof}
\olpseganchor{OLP-0118-B022}{end}

\olpseganchor{OLP-0119-B005}{start}
\olfileid{fol}{axd}{ded}
\olpseganchor{OLP-0119-B005}{end}

\olpseganchor{OLP-0119-B006}{start}
\olsection{د استنتاج قضيه}
\olpseganchor{OLP-0119-B006}{end}

\olpseganchor{OLP-0119-B007}{start}
لکه چې مو وليدل، په بديهي نظام کښې اشتقاقونه ورکول ستونزمن
دي، او اشتقاقونه موندل هم ګرانېږي۔ د فارمولونه اوږدې لړۍ
په رښتيا ليکلو پر ځاے، زياتره دا استدلال اسانه وي چې داسې
اشتقاقونه شته؛ د دې دپاره څو ساده پايلې کاروو۔ درې داسې پايلې
مو لا ثابتې کړې دي: \olref[ptn]{prop:reflexivity} وايي کله چې
\LR{$!A \in \Gamma$} پوهېږو، تل \LR{$\Gamma \Proves !A$} ويلے شو۔
\olref[ptn]{prop:monotonicity} وايي که \LR{$\Gamma \Proves !A$} وي، نو
\LR{$\Gamma \cup \{!B\} \Proves !A$} هم دے۔ او
\olref[ptn]{prop:transitivity} ښيي که \LR{$\Gamma \Proves !A$} او
\LR{$!A \Proves !B$} وي، نو \LR{$\Gamma \Proves !B$}۔ دا د مودوس پوننس يوه
بله ساده، ``ميتا'' بڼه ده:
\olpseganchor{OLP-0119-B007}{end}

\olpseganchor{OLP-0119-B008}{start}
\begin{prop}
  \ollabel{prop:mp} که \LR{$\Gamma \Proves !A$} او
  \LR{$\Gamma \Proves !A \lif !B$} وي، نو \LR{$\Gamma \Proves !B$}۔
\end{prop}
\olpseganchor{OLP-0119-B008}{end}

\olpseganchor{OLP-0119-B009}{start}
\begin{proof}
  موږ \LR{$\{!A, !A \lif !B\} \Proves !B$} لرو:
  \begin{derivation}
    1. & \LR{$!A$} & فرض\\
    2. & \LR{$!A \lif !B$} & فرض\\
    3. & \LR{$!B$} & 1, 2, مودوس پوننس
  \end{derivation}
  د \olref[ptn]{prop:transitivity} له مخې \LR{$\Gamma \Proves !B$}۔
\end{proof}
\olpseganchor{OLP-0119-B009}{end}

\olpseganchor{OLP-0119-B010}{start}
په دې برخه کښې تر ټولو مهمه پايله چې کاروو ئې د استنتاج قضيه ده:
\olpseganchor{OLP-0119-B010}{end}

\olpseganchor{OLP-0119-B011}{start}
\begin{thm}[د استنتاج قضيه]
\ollabel{thm:deduction-thm} \LR{$\Gamma \cup \{!A\} \Proves !B$} هله
او يوازې هله چې \LR{$\Gamma \Proves !A \lif !B$}۔
\end{thm}
\olpseganchor{OLP-0119-B011}{end}

\olpseganchor{OLP-0119-B012}{start}
\begin{proof}
د ``که'' لوری سمدستي دے۔ که \LR{$\Gamma \Proves !A \lif !B$} وي، د
\olref[ptn]{prop:monotonicity} له مخې
\LR{$\Gamma \cup \{!A\} \Proves !A \lif !B$} هم دے۔ همدارنګه د
\olref[ptn]{prop:reflexivity} له مخې
\LR{$\Gamma \cup \{!A\} \Proves !A$}۔ نو د \olref{prop:mp} له مخې
\LR{$\Gamma \cup \{!A\} \Proves !B$}۔
\olpseganchor{OLP-0119-B012}{end}

\olpseganchor{OLP-0119-B013}{start}
د ``يوازې که'' لوري دپاره له~\LR{$\Gamma \cup \{!A\}$} څخه د~\LR{$!B$} د
اشتقاق پر اوږدوالي استقرا کوو۔
\olpseganchor{OLP-0119-B013}{end}

\olpseganchor{OLP-0119-B014}{start}
د استقرا د بنسټ دپاره دعوه د اوږدوالي~\LR{$1$} د هر اشتقاق
دپاره ثابتوو۔ له~\LR{$\Gamma \cup \{!A\}$} څخه د~\LR{$!B$} د اوږدوالي~\LR{$1$}
اشتقاق يوازې له~\LR{$!B$} څخه جوړ وي؛ او که سم وي، يا
\LR{$!B \in \Gamma \cup \{!A\}$} وي يا \LR{$!B$} بديهي اصل وي۔
که \LR{$!B \in \Gamma$} وي يا بديهي اصل وي، نو \LR{$\Gamma \Proves !B$}۔
همدارنګه د \olref[prp]{ax:lif1} له مخې
\LR{$\Gamma \Proves !B \lif (!A \lif !B)$}، او \olref{prop:mp} موږ ته
\LR{$\Gamma \Proves !A \lif !B$} راکوي۔ که \LR{$!B \in \{ !A\}$} وي، نو
\LR{$\Gamma \Proves !A \lif !B$}؛ ځکه وروستۍ تړلے فارمول
\LR{$!A \lif !B$} د~\LR{$!A \lif !A$} هماغه جمله ده، او په
\olref[pro]{ex:identity} کښې مو هغه اشتقاق کړې ده۔
\olpseganchor{OLP-0119-B014}{end}

\olpseganchor{OLP-0119-B015}{start}
د استقرا په پړاو کښې فرض کړئ له~\LR{$\Gamma \cup \{!A\}$} څخه د~\LR{$!B$}
يو اشتقاق په داسې ګام~\LR{$!B$} پاے ته رسېږي چې په مودوس پوننس
توجيه شوے دے۔ (که په مودوس پوننس توجيه شوے نۀ وي، نو
\LR{$!B \in \Gamma$}، \LR{$!B \ident !A$}، يا \LR{$!B$} بديهي اصل دے، او د استقرا
د بنسټ هماغه استدلال پرې پلي کېږي۔) بيا د کوم فارمول~\LR{$!C$}
دپاره په اشتقاق کښې ځينې مخکښې ګامونه \LR{$!C \lif !B$} او
\LR{$!C$} دي؛ يعنې \LR{$\Gamma \cup \{!A\} \Proves !C \lif !B$} او
\LR{$\Gamma \cup \{!A\} \Proves !C$}، او اړوند اشتقاقونه لنډ دي،
نو د استقرا فرضيه پرې پلي کېږي۔ ځکه دواړه دا لرو:
\begin{align*}
  & \Gamma \Proves !A \lif (!C \lif !B); \\
  & \Gamma \Proves !A \lif !C.
\end{align*}
خو د \olref[prp]{ax:lif2} له مخې دا هم لرو:
\[
\Gamma \Proves (!A \lif (!C \lif !B)) \lif
((!A\lif !C)  \lif (!A \lif !B)),
\]
او د \olref{prop:mp} دوه تطبيقونه، لکه چې پکار ده،
\LR{$\Gamma \Proves !A \lif !B$} راکوي۔
\end{proof}
\olpseganchor{OLP-0119-B015}{end}

\olpseganchor{OLP-0119-B016}{start}
پام وکړئ چې \olref[prp]{ax:lif1} او \olref[prp]{ax:lif2} په عين
دې موخه ټاکل شوي وو چې د استنتاج قضيه سمه وي۔
\olpseganchor{OLP-0119-B016}{end}

\olpseganchor{OLP-0119-B017}{start}
لاندې د اشتقاق وړتيا په اړه ځينې ګټور حقيقتونه دي چې د تمرينونو
په توګه ئې پرېږدو۔
\olpseganchor{OLP-0119-B017}{end}

\olpseganchor{OLP-0119-B018}{start}
\begin{prop}
  \ollabel{prop:derivfacts}
\begin{enumerate}
\item \LR{$\Proves (!A \lif !B) \lif ((!B \lif !C)
  \lif (!A \lif !C))$}؛ \ollabel{derivfacts:a}
\item که \LR{$\Gamma \cup \{ \lnot !A\}
  \Proves \lnot !B$} وي، نو \LR{$\Gamma \cup \{ !B\} \Proves
  !A$} (مقابل عکس)؛ \ollabel{derivfacts:b}
\item \LR{$\{ !A, \lnot!A\} \Proves
    !B$} (له کاذبۍ هر څه، انفجار)؛ \ollabel{derivfacts:c}
\item \LR{$\{ \lnot\lnot!A\} \Proves
  !A$} (دوه ګونې نفي حذف)؛\ollabel{derivfacts:d}
\item که \LR{$\Gamma \Proves \lnot\lnot!A$} وي، نو \LR{$\Gamma \Proves
  !A$}؛\ollabel{derivfacts:e}
\end{enumerate}
\end{prop}
\olpseganchor{OLP-0119-B018}{end}

\olpseganchor{OLP-0119-B019}{start}
\begin{prob}
  \olref[fol][axd][ded]{prop:derivfacts} ثابت کړئ۔
\end{prob}
\olpseganchor{OLP-0119-B019}{end}

\olpseganchor{OLP-0120-B004}{start}
\olfileid{fol}{axd}{ddq}
\olpseganchor{OLP-0120-B004}{end}

\olpseganchor{OLP-0120-B005}{start}
\olsection{له کميت ټاکونکو سره د استنتاج قضيه}
\olpseganchor{OLP-0120-B005}{end}

\olpseganchor{OLP-0120-B006}{start}
\begin{thm}[د استنتاج قضيه]
\ollabel{thm:deduction-thm-q} که
\LR{$\Gamma \cup \{!A\} \Proves !B$} وي، نو
\LR{$\Gamma \Proves !A \lif !B$}۔
\end{thm}
\olpseganchor{OLP-0120-B006}{end}

\olpseganchor{OLP-0120-B007}{start}
\begin{proof}
يو ځل بيا له~\LR{$\Gamma \cup \{!A\}$} څخه د~\LR{$!B$} د اشتقاق
پر اوږدوالي استقرا کوو۔
\olpseganchor{OLP-0120-B007}{end}

\olpseganchor{OLP-0120-B008}{start}
د استقرا د بنسټ ثبوت د
\olref[ded]{thm:deduction-thm} په ثبوت کښې له بنسټ سره يو شان دے۔
\olpseganchor{OLP-0120-B008}{end}

\olpseganchor{OLP-0120-B009}{start}
د استقرا په پړاو کښې بيا فرض کړئ له~\LR{$\Gamma \cup \{!A\}$} څخه د
\LR{$!B$} اشتقاق په داسې ګام~\LR{$!B$} پاے ته رسېږي چې په استنتاجي
قاعده توجيه شوے دے۔ که استنتاجي قاعده مودوس پوننس وي، د
\olref[ded]{thm:deduction-thm} د ثبوت په شان مخکښې ځو۔ که قاعده
\QR{} وي، پوهېږو چې \LR{$!B \ident !C \lif \lforall[x][!D(x)]$} او د
\LR{$!C \lif !D(a)$} بڼې فارمول په اشتقاق کښې له مخکښې
راغلے دے، چې \LR{$a$} په~\LR{$!C$}، \LR{$!A$} يا~\LR{$\Gamma$} کښې نۀ راځي۔ نو دا لرو:
\begin{align*}
  \Gamma \cup \{!A\} & \Proves !C \lif !D(a),\\
  \intertext{\RL{او د استقرا فرضيه پلي کېږي؛ يعنې دا لرو:}}
    \Gamma & \Proves !A \lif (!C \lif !D(a)).\\
  \intertext{\RL{له دې څخه}}
  & \Proves (!A \lif (!C \lif !D(a))) \lif ((!A \land !C) \lif !D(a))\\
  \intertext{\RL{او مودوس پوننس سره ترلاسه کوو:}}
  \Gamma & \Proves (!A \land !C) \lif !D(a).\\
  \intertext{\RL{ځکه د ځانګړي متغير شرط لا هم پلي کېږي، دې اشتقاق ته په \QR{} توجيه شوے ګام ورزياتولے شو، او ترلاسه کوو:}}
    \Gamma & \Proves (!A \land !C) \lif \lforall[x][!D(x)].\\
    \intertext{\RL{دا هم لرو:}}
    & \Proves ((!A \land !C) \lif \lforall[x][!D(x)]) \lif (!A \lif (!C \lif \lforall[x][!D(x)])),\\
    \intertext{\RL{نو د مودوس پوننس له مخې،}}
    \Gamma & \Proves !A \lif (!C \lif \lforall[x][!D(x)]),
\end{align*}
يعنې \LR{$\Gamma \Proves !A \lif !B$}۔
\olpseganchor{OLP-0120-B009}{end}

\olpseganchor{OLP-0120-B010}{start}
هغه حالت د تمرين په توګه پرېږدو چې \LR{$!B$} د~\QR{} قاعدې په وسيله
توجيه کېږي، خو د~\LR{$\lexists[x][!D(x)] \lif !C$} بڼه لري۔
\end{proof}
\olpseganchor{OLP-0120-B010}{end}

\olpseganchor{OLP-0120-B011}{start}
\begin{prob}
  د~\olref[fol][axd][ddq]{thm:deduction-thm-q} ثبوت بشپړ کړئ۔
\end{prob}
\olpseganchor{OLP-0120-B011}{end}

\olpseganchor{OLP-0121-B004}{start}
\olfileid{fol}{axd}{prv}
\olpseganchor{OLP-0121-B004}{end}

\olpseganchor{OLP-0121-B005}{start}
\olsection{د اشتقاق وړتيا او سازګاري}
\olpseganchor{OLP-0121-B005}{end}

\olpseganchor{OLP-0121-B006}{start}
اوس به د اشتقاق وړتيا اړيکې يو شمېر خاصيتونه ثابت کړو۔ دا په
خپلواکه توګه په زړه پورې دي، خو هر يو به د بشپړوالي د قضیې په
ثبوت کښې ونډه ولري۔
\olpseganchor{OLP-0121-B006}{end}

\olpseganchor{OLP-0121-B007}{start}
\begin{prop}\ollabel{prop:provability-contr}
  که \LR{$\Gamma \Proves !A$} او \LR{$\Gamma \cup \{!A\}$} ناسازګار وي، نو
  \LR{$\Gamma$} ناسازګار دے۔
\end{prop}
\olpseganchor{OLP-0121-B007}{end}

\olpseganchor{OLP-0121-B008}{start}
\begin{proof}
  که \LR{$\Gamma \cup \{!A\}$} ناسازګار وي، نو
  \LR{$\Gamma \cup \{!A\} \Proves \lfalse$}۔ د
  \olref[ptn]{prop:reflexivity} له مخې د هر \LR{$!B \in \Gamma$} دپاره
  \LR{$\Gamma \Proves !B$}۔ د فرض له مخې \LR{$\Gamma \Proves !A$} هم دے،
  نو د هر \LR{$!B \in \Gamma \cup \{!A\}$} دپاره
  \LR{$\Gamma \Proves !B$}۔ د \olref[ptn]{prop:transitivity} له مخې
  \LR{$\Gamma \Proves \lfalse$}؛ يعنې \LR{$\Gamma$} ناسازګار دے۔
\end{proof}
\olpseganchor{OLP-0121-B008}{end}

\olpseganchor{OLP-0121-B009}{start}
\begin{prop}
\ollabel{prop:prov-incons}
\LR{$\Gamma \Proves !A$} هله او يوازې هله چې
\LR{$\Gamma \cup \{\lnot !A\}$} ناسازګار وي۔
\end{prop}
\olpseganchor{OLP-0121-B009}{end}

\olpseganchor{OLP-0121-B010}{start}
\begin{proof}
لومړے فرض کړئ \LR{$\Gamma \Proves !A$}۔ د
\olref[ptn]{prop:monotonicity} له مخې
\LR{$\Gamma \cup \{\lnot !A\} \Proves !A$}۔ د
\olref[ptn]{prop:reflexivity} له مخې
\LR{$\Gamma \cup \{\lnot !A\} \Proves \lnot !A$}۔ د
\olref[prp]{ax:lnot2} له مخې
\LR{$\Proves \lnot !A \lif (!A \lif \lfalse)$} هم لرو۔ نو د
\olref[ded]{prop:mp} په دوه ځله پلي کولو
\LR{$\Gamma \cup \{\lnot !A\} \Proves \lfalse$} لرو۔
\olpseganchor{OLP-0121-B010}{end}

\olpseganchor{OLP-0121-B011}{start}
اوس فرض کړئ \LR{$\Gamma \cup \{\lnot !A\}$} ناسازګار دے؛ يعنې
\LR{$\Gamma \cup \{\lnot !A\} \Proves \lfalse$}۔ د استنتاج د قضیې له
مخې \LR{$\Gamma \Proves \lnot !A \lif \lfalse$}۔ د
\olref[prp]{ax:lfalse2} له مخې
\LR{$\Gamma \Proves (\lnot !A \lif \lfalse) \lif \lnot\lnot !A$}، نو
د \olref[ded]{prop:mp} له مخې \LR{$\Gamma \Proves \lnot\lnot !A$}۔
ځکه \LR{$\Gamma \Proves \lnot\lnot !A \lif !A$}
(\olref[prp]{ax:dne})، د \olref[ded]{prop:mp} په بيا پلي کولو
\LR{$\Gamma \Proves !A$} لرو۔
\end{proof}
\olpseganchor{OLP-0121-B011}{end}

\olpseganchor{OLP-0121-B012}{start}
\begin{prob}
ثابته کړئ چې \LR{$\Gamma \Proves \lnot !A$} هله او يوازې هله چې
\LR{$\Gamma \cup \{!A\}$} ناسازګار وي۔
\end{prob}
\olpseganchor{OLP-0121-B012}{end}

\olpseganchor{OLP-0121-B013}{start}
\begin{prop}\ollabel{prop:explicit-inc}
  که \LR{$\Gamma \Proves !A$} او \LR{$\lnot !A \in \Gamma$} وي، نو
  \LR{$\Gamma$} ناسازګار دے۔
\end{prop}
\olpseganchor{OLP-0121-B013}{end}

\olpseganchor{OLP-0121-B014}{start}
\begin{proof}
  د \olref[prp]{ax:lnot2} له مخې
  \LR{$\Gamma \Proves \lnot !A \lif (!A \lif \lfalse)$}۔ د
  \olref[ded]{prop:mp} په دوه ځله پلي کولو
  \LR{$\Gamma \Proves \lfalse$}۔
\end{proof}
\olpseganchor{OLP-0121-B014}{end}


\olpseganchor{OLP-0121-B015}{start}
\begin{prop}\ollabel{prop:provability-exhaustive}
  که \LR{$\Gamma \cup \{!A\}$} او \LR{$\Gamma \cup \{\lnot !A\}$} دواړه
  ناسازګار وي، نو \LR{$\Gamma$} ناسازګار دے۔
\end{prop}
\olpseganchor{OLP-0121-B015}{end}

\olpseganchor{OLP-0121-B016}{start}
\begin{proof}
تمرين۔
\end{proof}
\olpseganchor{OLP-0121-B016}{end}

\olpseganchor{OLP-0121-B017}{start}
\begin{prob}
  \olref[fol][axd][prv]{prop:provability-exhaustive} ثابت کړئ۔
\end{prob}
\olpseganchor{OLP-0121-B017}{end}

\olpseganchor{OLP-0122-B004}{start}
\olfileid{fol}{axd}{ppr}
\olpseganchor{OLP-0122-B004}{end}

\olpseganchor{OLP-0122-B005}{start}
\olsection{د اشتقاق وړتيا او د قضیې نښلوونکي}
\olpseganchor{OLP-0122-B005}{end}

\olpseganchor{OLP-0122-B006}{start}
\begin{explain}
  ثابتوو چې د بديهي استنتاج د~\LR{$\Proves$} د اشتقاق وړتيا اړيکه
  دومره پياوړې ده چې د قضيې د نښلوونکو ځينې بنسټيز حقيقتونه ثابت
  کړي؛ لکه \LR{$!A \land !B \Proves !A$} او
  \LR{$!A, !A \lif !B \Proves !B$} (مودوس پوننس)۔ دغه حقيقتونه د
  بشپړوالي د قضیې د ثبوت دپاره پکار دي۔
\end{explain}
\olpseganchor{OLP-0122-B006}{end}

\olpseganchor{OLP-0122-B007}{start}
\begin{prop}\ollabel{prop:provability-land}
  \begin{enumerate}
  \item \ollabel{prop:provability-land-left} دواړه
    \LR{$!A \land !B \Proves !A$} او \LR{$!A \land !B \Proves !B$}
  \item \ollabel{prop:provability-land-right}
    \LR{$!A, !B \Proves !A \land !B$}۔
  \end{enumerate}
\end{prop}
\olpseganchor{OLP-0122-B007}{end}

\olpseganchor{OLP-0122-B008}{start}
\begin{proof}
  \begin{enumerate}
    \item د \olref[prp]{ax:land1} او \olref[prp]{ax:land2} څخه په
      مودوس پوننس۔
  \item د \olref[prp]{ax:land3} څخه د مودوس پوننس په دوه تطبيقونو۔
  \end{enumerate}
\end{proof}
\olpseganchor{OLP-0122-B008}{end}


\olpseganchor{OLP-0122-B009}{start}
\begin{prop}\ollabel{prop:provability-lor}
  \begin{enumerate}
  \item \LR{$!A \lor !B, \lnot !A, \lnot !B$} ناسازګار دے۔
  \item دواړه \LR{$!A \Proves !A \lor !B$} او
    \LR{$!B \Proves !A \lor !B$}۔
  \end{enumerate}
\end{prop}
\olpseganchor{OLP-0122-B009}{end}

\olpseganchor{OLP-0122-B010}{start}
\begin{proof}
  \begin{enumerate}
  \item له \olref[prp]{ax:lnot2} څخه
    \LR{$\Proves \lnot !A \lif (!A \lif \lfalse)$} او
    \LR{$\Proves \lnot !B \lif (!B \lif \lfalse)$} ترلاسه کوو۔ نو د
    استنتاج د قضیې له مخې \LR{$\{\lnot !A\} \Proves !A \lif \lfalse$}
    او \LR{$\{\lnot !B\} \Proves !B \lif \lfalse$} لرو۔ له
    \olref[prp]{ax:lor3} څخه
    \LR{$\{\lnot !A, \lnot !B\} \Proves (!A \lor !B) \lif \lfalse$}
    ترلاسه کوو۔ د استنتاج د قضیې له مخې
    \LR{$\{!A \lor !B, \lnot !A, \lnot !B\} \Proves \lfalse$}۔
  \item له \olref[prp]{ax:lor1} او \olref[prp]{ax:lor2} څخه په
    مودوس پوننس۔
  \end{enumerate}
\end{proof}
\olpseganchor{OLP-0122-B010}{end}

\olpseganchor{OLP-0122-B011}{start}
\begin{prop}\ollabel{prop:provability-lif}
  \begin{enumerate}
  \item \ollabel{prop:provability-lif-left}
    \LR{$!A, !A \lif !B \Proves !B$}۔
  \item \ollabel{prop:provability-lif-right}
    دواړه \LR{$\lnot !A \Proves !A \lif !B$} او
    \LR{$!B \Proves !A \lif !B$}۔
  \end{enumerate}
\end{prop}
\olpseganchor{OLP-0122-B011}{end}

\olpseganchor{OLP-0122-B012}{start}
\begin{proof}
  \begin{enumerate}
  \item دا اشتقاق کولے شو:
    \begin{derivation}
      1. & \LR{$!A$} & \Hyp\\
      2. & \LR{$!A \lif !B$} & \Hyp\\
      3. & \LR{$!B$} & 1, 2, \MP
    \end{derivation}
\olpseganchor{OLP-0122-B012}{end}
    
\olpseganchor{OLP-0122-B013}{start}
  \item په ترتيب د \olref[prp]{ax:lnot2} او
    \olref[prp]{ax:lif1}، او د استنتاج د قضیې له مخې۔
  \end{enumerate}
\end{proof}
\olpseganchor{OLP-0122-B013}{end}

\olpseganchor{OLP-0123-B005}{start}
\olfileid{fol}{axd}{qpr}
\olpseganchor{OLP-0123-B005}{end}

\olpseganchor{OLP-0123-B006}{start}
\olsection{د اشتقاق وړتيا او کميت ټاکونکي}
\olpseganchor{OLP-0123-B006}{end}

\olpseganchor{OLP-0123-B007}{start}
\begin{explain}
  د بشپړوالي قضيه دا هم غواړي چې بديهي استنتاجونه د~\LR{$\Proves$} په
  اړه په دې برخه کښې ثابت شوي حقيقتونه ورکړي۔
\end{explain}
\olpseganchor{OLP-0123-B007}{end}

\olpseganchor{OLP-0123-B008}{start}
\begin{thm}
\ollabel{thm:strong-generalization} که \LR{$c$} داسې ثابت سمبول وي چې
په~\LR{$\Gamma$} يا~\LR{$!A(x)$} کښې نۀ راځي، او \LR{$\Gamma \Proves !A(c)$} وي،
نو \LR{$\Gamma \Proves \lforall[x][!A(x)]$}۔
\end{thm}
\olpseganchor{OLP-0123-B008}{end}

\olpseganchor{OLP-0123-B009}{start}
\begin{proof}
د استنتاج د قضیې له مخې \LR{$\Gamma \Proves \ltrue \lif !A(c)$}۔ ځکه
\LR{$c$} په~\LR{$\Gamma$} يا~\LR{$\top$} کښې نۀ راځي، ترلاسه کوو چې
\LR{$\Gamma \Proves \ltrue \lif \lforall[x][!A(x)]$}۔ ځکه رښتيا يو
بديهي اصل دے، د ميتا-مودوس پوننس له مخې
\LR{$\Gamma \Proves \lforall[x][!A(x)]$}۔
\end{proof}
\olpseganchor{OLP-0123-B009}{end}

\olpseganchor{OLP-0123-B010}{start}
\begin{prop}
\ollabel{prop:provability-quantifiers}
\begin{enumerate}
\item \LR{$!A(t) \Proves \lexists[x][!A(x)]$}.
\olpseganchor{OLP-0123-B010}{end}

\olpseganchor{OLP-0123-B011}{start}
\item \LR{$\lforall[x][!A(x)] \Proves !A(t)$}.
\end{enumerate}
\end{prop}
\olpseganchor{OLP-0123-B011}{end}

\olpseganchor{OLP-0123-B012}{start}
\begin{proof}
\begin{enumerate}
\item د \olref[qua]{ax:q2} او د استنتاج د قضیې له مخې۔
\item د \olref[qua]{ax:q1} او د استنتاج د قضیې له مخې۔
\end{enumerate}
\end{proof}
\olpseganchor{OLP-0123-B012}{end}

\olpseganchor{OLP-0124-B004}{start}
\olfileid{fol}{axd}{sou}
\olpseganchor{OLP-0124-B004}{end}
      
\olpseganchor{OLP-0124-B005}{start}
\olsection{سموالے}
\olpseganchor{OLP-0124-B005}{end}

\olpseganchor{OLP-0124-B006}{start}
\begin{explain}
د بديهي استنتاج په شان اشتقاق نظام هله \emph{سم} دے چې
هغه څيزونه اشتقاق نۀ کړي چې په حقيقت کښې نۀ صدق کېږي۔ نو
سموالے د اشتقاق نظامونو د تضمين شوې خونديتوب يو ډول خاصيت
دے۔ د دې له مخې چې کوم ثبوت-تيوريکي خاصيت تر بحث لاندې وي، مثلاً
غواړو پوه شو چې:
\begin{enumerate}
\item هر د اشتقاق وړ~\LR{$!A$} معتبر دے؛
\item که \LR{$!A$} له ځينو نورو~\LR{$\Gamma$} څخه د اشتقاق وړ وي، د هغو
  معنايي پايله هم ده؛
\item که د فارمولونه يو سټ~\LR{$\Gamma$} ناسازګار وي، د صدق وړ نۀ
  دے۔
\end{enumerate}
دا د اشتقاق نظام مهم خاصيتونه دي۔ که له دغو څخه کوم يو
سم نۀ وي، اشتقاق نظام نيمګړے دے---له حده زيات څيزونه به
اشتقاق کوي۔ له دې امله د اشتقاق نظام سموالے ثابتول
تر ټولو زيات اهميت لري۔
\end{explain}
\olpseganchor{OLP-0124-B006}{end}

\olpseganchor{OLP-0124-B007}{start}
\begin{prop}
  که \LR{$!A$} يو بديهي اصل وي، نو
  د هر جوړښت~\LR{$\Struct{M}$} او متغير-ګومارنې~\LR{$s$} دپاره \LR{$\Sat{M}{!A}[s]$}۔
\end{prop}
\olpseganchor{OLP-0124-B007}{end}

\olpseganchor{OLP-0124-B008}{start}
\begin{proof}
بايد تصديق کړو چې ټول بديهي اصول معتبر دي۔ د بېلګې
  په توګه د~\olref[qua]{ax:q1} حالت دا دے: فرض کړئ \LR{$t$} په~\LR{$!A$}
  کښې د~\LR{$x$} دپاره ازاد وي، او
  \LR{$\Sat{M}{\lforall[x][!A]}[s]$} ومنئ۔ بيا د پوره کېدو د تعريف له
  مخې، د هرې \LR{$\varAssign{s'}{s}{x}$} دپاره \LR{$\Sat{M}{!A}[s']$} هم
  دے، او په ځانګړي ډول دا هله سم دے چې
  \LR{$s'(x) = \Value{t}{M}[s]$} وي۔ د
  \olref[syn][ext]{prop:ext-formulas} له مخې
  \LR{$\Sat{M}{\Subst{!A}{t}{x}}[s]$}۔ دا ښيي چې
  \LR{$\Sat{M}{(\lforall[x][!A] \lif \Subst{!A}{t}{x})}[s]$}۔
\end{proof}
\olpseganchor{OLP-0124-B008}{end}

\olpseganchor{OLP-0124-B009}{start}
\begin{thm}[سموالے]
\ollabel{thm:soundness}
که \LR{$\Gamma \Proves !A$} وي، نو \LR{$\Gamma \Entails !A$}۔
\end{thm}
\olpseganchor{OLP-0124-B009}{end}

\olpseganchor{OLP-0124-B010}{start}
\begin{proof}
له~\LR{$\Gamma$} څخه د~\LR{$!A$} د اشتقاق پر اوږدوالي استقرا کوو۔
که داسې هېڅ ګام نۀ وي چې په استنتاج توجيه شوے وي، نو په
اشتقاق کښې ټول فارمولونه يا د بديهي اصولو بېلګې دي يا
په~\LR{$\Gamma$} کښې دي۔ د مخکښې قضیې له مخې ټول بديهي اصول
معتبر دي؛ نو که \LR{$!A$} بديهي اصل وي،
\LR{$\Gamma \Entails !A$}۔ که \LR{$!A \in \Gamma$} وي، نو په څرګنده
\LR{$\Gamma \Entails !A$}۔
\olpseganchor{OLP-0124-B010}{end}

\olpseganchor{OLP-0124-B011}{start}
که د~\LR{$!A$} د اشتقاق وروستنے ګام په مودوس پوننس توجيه شوے
وي، نو په اشتقاق کښې فارمولونه~\LR{$!B$} او~\LR{$!B \lif !A$}
شته، او د استقرا فرضيه د اشتقاق پر هغو برخو پلي کېږي چې په
دغو فارمولونه پاے ته رسېږي (ځکه په هره برخه کښې لږ تر لږه يو
په استنتاج توجيه شوے ګام کم دے)۔ نو د استقرا د فرضيې له مخې
\LR{$\Gamma \Entails !B$} او \LR{$\Gamma \Entails !B \lif !A$}۔ بيا د
\olref[syn][sem]{thm:sem-deduction}
له مخې \LR{$\Gamma \Entails !A$}۔
\olpseganchor{OLP-0124-B011}{end}

\olpseganchor{OLP-0124-B012}{start}
اوس فرض کړئ وروستنے ګام په~\QR توجيه شوے دے۔ بيا هغه
ګام د~\LR{$!C \lif \lforall[x][!B(x)]$} بڼه لري، او ترې مخکښې ګام
\LR{$!C \lif !B(c)$} دے، چې \LR{$c$} په~\LR{$\Gamma$}، \LR{$!C$} يا
\LR{$\lforall[x][!B(x)]$} کښې نۀ راځي۔ د استقرا د فرضيې له مخې
\LR{$\Gamma \Entails !C \lif !B(c)$}۔ د
\olref[syn][sem]{thm:sem-deduction} له مخې
\LR{$\Gamma \cup \{!C\} \Entails !B(c)$}۔
\olpseganchor{OLP-0124-B012}{end}

\olpseganchor{OLP-0124-B013}{start}
داسې يو جوړښت~\LR{$\Struct{M}$} په پام کښې ونيسئ چې
\LR{$\Sat{M}{\Gamma \cup \{!C\}}$} وي۔ بايد وښيو چې
\LR{$\Sat{M}{\lforall[x][!B(x)]}$}۔ ځکه
\LR{$\lforall[x][!B(x)]$} يوه تړلے فارمول ده، د دې معنٰی دا ده چې د
هرې متغير-ګومارنې~\LR{$s$} دپاره بايد \LR{$\Sat{M}{!B(x)}[s]$} وښيو
(\olref[syn][ass]{prop:sat-quant})۔ ځکه \LR{$\Gamma \cup \{!C\}$} ټول
له جملو جوړ دے، د هر~\LR{$!D \in \Gamma$} دپاره
\LR{$\Sat{M}{!D}[s]$} د~\olref[syn][sat]{defn:satisfaction} له مخې دے۔
\LR{$\Struct{M'}$} داسې جوړښت ونيسئ چې د~\LR{$\Struct{M}$} په شان وي، خو
\LR{$\Assign{c}{M'} = s(x)$} وي۔ ځکه \LR{$c$} په~\LR{$\Gamma$} يا~\LR{$!C$} کښې نۀ
راځي، د~\olref[syn][ext]{cor:extensionality-sent} له مخې
\LR{$\Sat{M'}{\Gamma \cup \{!C\}}$}۔ ځکه
\LR{$\Gamma \cup \{!C\} \Entails !B(c)$}، نو \LR{$\Sat{M'}{!B(c)}$}۔
ځکه \LR{$!B(c)$} يوه تړلے فارمول ده، د
\olref[syn][ass]{prop:sentence-sat-true} له مخې
\LR{$\Sat{M'}{!B(c)}[s]$}۔ د~\olref[syn][ext]{prop:ext-formulas} له مخې
\LR{$\Sat{M'}{!B(x)}[s]$} هله او يوازې هله چې
\LR{$\Sat{M'}{!B(c)}[s]$} (ياد ولرئ چې \LR{$!B(c)$} هماغه
\LR{$\Subst{!B(x)}{c}{x}$} دے)۔ نو \LR{$\Sat{M'}{!B(x)}[s]$}۔ ځکه \LR{$c$} په
\LR{$!B(x)$} کښې نۀ راځي، د~\olref[syn][ext]{prop:extensionality} له
مخې \LR{$\Sat{M}{!B(x)}[s]$}۔ خو \LR{$s$} اختياري متغير-ګومارنه وه؛ نو
\LR{$\Sat{M}{\lforall[x][!B(x)]}$}۔ له دې امله
\LR{$\Gamma \cup \{!C\} \Entails \lforall[x][!B(x)]$}۔ د
\olref[syn][sem]{thm:sem-deduction} له مخې
\LR{$\Gamma \Entails !C \lif \lforall[x][!B(x)]$}۔
\olpseganchor{OLP-0124-B013}{end}

\olpseganchor{OLP-0124-B014}{start}
هغه حالت د تمرين په توګه پرېږدو چې \LR{$!A$} په~\QR{} توجيه کېږي، خو
د~\LR{$\lexists[x][!B(x)] \lif !C$} بڼه لري۔
\end{proof}
\olpseganchor{OLP-0124-B014}{end}

\olpseganchor{OLP-0124-B015}{start}
\begin{prob}
  د~\olref[fol][axd][sou]{thm:soundness} ثبوت بشپړ کړئ۔
\end{prob}
\olpseganchor{OLP-0124-B015}{end}

\olpseganchor{OLP-0124-B016}{start}
\begin{cor}
\ollabel{cor:weak-soundness}
که \LR{$\Proves !A$} وي، نو \LR{$!A$} معتبر دے۔
\end{cor}
\olpseganchor{OLP-0124-B016}{end}

\olpseganchor{OLP-0124-B017}{start}
\begin{cor}
\ollabel{cor:consistency-soundness}
که \LR{$\Gamma$} د صدق وړ وي، نو سازګار دے۔
\end{cor}
\olpseganchor{OLP-0124-B017}{end}

\olpseganchor{OLP-0124-B018}{start}
\begin{proof}
نقيض عکس ثابتوو۔ فرض کړئ \LR{$\Gamma$} سازګار نۀ دے۔ بيا
\LR{$\Gamma \Proves \lfalse$}؛ يعنې له~\LR{$\Gamma$} څخه د~\LR{$\lfalse$} يو
اشتقاق شته۔ د~\olref{thm:soundness} له مخې هر
جوړښت~\LR{$\Struct{M}$}
چې \LR{$\Gamma$} پوره کوي، بايد \LR{$\lfalse$} هم پوره کړي۔ ځکه د هر
جوړښت~\LR{$\Struct{M}$} دپاره \LR{$\Sat/{M}{\lfalse}$}، نو هېڅ
\LR{$\Struct{M}$}، \LR{$\Gamma$} نۀ شي پوره
کولے؛ يعنې \LR{$\Gamma$} د صدق وړ نۀ دے۔
\end{proof}
\olpseganchor{OLP-0124-B018}{end}

\olpseganchor{OLP-0125-B004}{start}
\olfileid{fol}{axd}{ide}
\olpseganchor{OLP-0125-B004}{end}

\olpseganchor{OLP-0125-B005}{start}
\olsection{له عينيت سره اشتقاقونه}
\olpseganchor{OLP-0125-B005}{end}

\olpseganchor{OLP-0125-B006}{start}
په اشتقاقونه کښې د~\LR{$\eq$} د شاملولو دپاره يوازې نوې بديهي
شېماوې ورزياتوو۔ د~اشتقاق او~\LR{$\Proves$} تعريف هماغسې پاتې
کېږي؛ يوازې نوي بديهي اصول هم منو۔
\olpseganchor{OLP-0125-B006}{end}

\olpseganchor{OLP-0125-B007}{start}
\begin{defn}[د عينيت بديهي اصول]
\begin{align}
\ollabel{ax:id1} & \eq[t][t], \\
\ollabel{ax:id2} & \eq[t_1][t_2] \lif (!B(t_1) \lif
  !B(t_2)),
\end{align}
د هر تړلي ترم~\LR{$t$}، \LR{$t_1$} او~\LR{$t_2$} دپاره۔
\end{defn}
\olpseganchor{OLP-0125-B007}{end}

\olpseganchor{OLP-0125-B008}{start}
\begin{prop}\ollabel{prop:sound}
  بديهي اصول~\olref{ax:id1} او~\olref{ax:id2} معتبر دي۔
\end{prop}
\olpseganchor{OLP-0125-B008}{end}

\olpseganchor{OLP-0125-B009}{start}
\begin{proof}
  تمرين۔
\end{proof}
\olpseganchor{OLP-0125-B009}{end}

\olpseganchor{OLP-0125-B010}{start}
\begin{prob}
\olref[fol][axd][ide]{prop:sound} ثابت کړئ۔
\end{prob}
\olpseganchor{OLP-0125-B010}{end}

\olpseganchor{OLP-0125-B011}{start}
\begin{prop}\ollabel{prop:iden1}
 د هر تړلي ترم~\LR{$t$} او سټ~\LR{$\Gamma$} دپاره
 \LR{$\Gamma \Proves \eq[t][t]$}۔
\end{prop}
\olpseganchor{OLP-0125-B011}{end}

\olpseganchor{OLP-0125-B012}{start}
\begin{prop}\ollabel{prop:iden2}
  که \LR{$\Gamma \Proves !A(t_1)$} او \LR{$\Gamma \Proves
  \eq[t_1][t_2]$} وي، نو \LR{$\Gamma \Proves !A(t_2)$}، په دې شرط چې
  دواړه ښودل شوي ترمونه تړلي وي۔
\end{prop}
\olpseganchor{OLP-0125-B012}{end}

\olpseganchor{OLP-0125-B013}{start}
\begin{proof}
دا فارمول
\[
(\eq[t_1][t_2] \lif (!A(t_1) \lif !A(t_2)))
\]
د~\olref{ax:id2} يوه بېلګه ده۔ پايله د~\MP په دوه پلي کولو سره
ترلاسه کېږي۔
\end{proof}
\olpseganchor{OLP-0125-B013}{end}

\paragraph{د منجمدې سرچينې بشپړ بديل متن}
\olpseganchor{OLP-0643-B005}{start}
\olfileid{fol}{axd}{prv-olp0643}
\olpseganchor{OLP-0643-B005}{end}

\olpseganchor{OLP-0643-B006}{start}
\olsection{د اشتقاق وړتيا ځانګړتياوې}
دا بحث د هيلبرټ د استنتاج هغه بڼه کاروي چې د اصولو کلي تړنې هم اصول ګڼي۔ د عمومي کولو په ثبوت کښې درې قاعدوي شېماوې په کار دي: د کلي فارمول ځانګړي کول، پر شرطيت د کميت وېش، او پر داسې متغير عمومي کول چې په مقدمې کښې ازاد نۀ وي۔ په دې بحث کښې Ax1، Ax2 او Ax3 په همدې ترتيب دغو شېماوو ته اشاره کوي۔
\olpseganchor{OLP-0643-B006}{end}

\olpseganchor{OLP-0643-B007}{start}
\begin{prop}[يکنواختي] که \LR{$\Gamma \subseteq \Delta$} او \LR{$\Gamma \Proves !A$} وي، نو \LR{$\Delta \Proves !A$} هم دے۔ \end{prop}
\olpseganchor{OLP-0643-B007}{end}

\olpseganchor{OLP-0643-B008}{start}
\begin{proof} هر متناهي \LR{$\Gamma_0 \subseteq \Gamma$} د~\LR{$\Delta$} متناهي فرعي سټ هم دے؛ نو د \LR{$\Gamma_0 \Proves !A$} اشتقاق، \LR{$\Delta \Proves !A$} هم ښيي۔ \end{proof}
\olpseganchor{OLP-0643-B008}{end}

\olpseganchor{OLP-0643-B009}{start}
\begin{prop}\ollabel{prop:provability} \LR{$\Gamma \Proves !A$} هغه او يوازې هغه وخت دے چې \LR{$\Gamma\cup \{\lnot!A \}$} نامتسق وي۔ \end{prop}
\olpseganchor{OLP-0643-B009}{end}

\olpseganchor{OLP-0643-B010}{start}
\begin{prop} \begin{enumerate}
\item \ollabel{prop:provability-contr} که \LR{$\Gamma \Proves !A$} او \LR{$\Gamma
\cup \{ !A\} \Proves \lfalse$} وي، نو \LR{$\Gamma$} نامتسق دے۔
\olpseganchor{OLP-0643-B010}{end}

\olpseganchor{OLP-0643-B011}{start}
\item \ollabel{prop:provability-lnot} که \LR{$\Gamma \cup \{!A\}
\Proves \lfalse$} وي، نو \LR{$\Gamma \Proves \lnot !A$}۔
\olpseganchor{OLP-0643-B011}{end}

\olpseganchor{OLP-0643-B012}{start}
\item \ollabel{prop:provability-exhaustive} که \LR{$\Gamma \cup \{!A\}
\Proves \lfalse$} او \LR{$\Gamma \cup \{\lnot !A\} \Proves
\lfalse$} وي، نو \LR{$\Gamma \Proves \lfalse$}۔
\olpseganchor{OLP-0643-B012}{end}

\olpseganchor{OLP-0643-B013}{start}
\item \ollabel{prop:provability-lor-left} که \LR{$\Gamma \cup \{!A\}
\Proves \lfalse$} او \LR{$\Gamma \cup \{!B\} \Proves
\lfalse$} وي، نو \LR{$\Gamma \cup \{!A \lor !B\} \Proves \lfalse$}۔
\olpseganchor{OLP-0643-B013}{end}

\olpseganchor{OLP-0643-B014}{start}
\item \ollabel{prop:provability-lor-right} که \LR{$\Gamma
\Proves !A$} يا \LR{$\Gamma \Proves !B$} وي، نو \LR{$\Gamma
\Proves !A \lor !B$}۔
\olpseganchor{OLP-0643-B014}{end}

\olpseganchor{OLP-0643-B015}{start}
\item \ollabel{prop:provability-land-left} که \LR{$\Gamma
\Proves !A \land !B$} وي، نو \LR{$\Gamma \Proves !A$} او
\LR{$\Gamma \Proves !B$}۔
\olpseganchor{OLP-0643-B015}{end}

\olpseganchor{OLP-0643-B016}{start}
\item \ollabel{prop:provability-land-right} که \LR{$\Gamma
\Proves !A$} او \LR{$\Gamma \Proves !B$} وي، نو \LR{$\Gamma
\Proves !A \land !B$}۔
\olpseganchor{OLP-0643-B016}{end}

\olpseganchor{OLP-0643-B017}{start}
\item \ollabel{prop:provability-mp} که \LR{$\Gamma \Proves
!A$} او \LR{$\Gamma \Proves !A \lif !B$} وي، نو \LR{$\Gamma
\Proves !B$}۔
\olpseganchor{OLP-0643-B017}{end}

\olpseganchor{OLP-0643-B018}{start}
\item \ollabel{prop:provability-lif} که \LR{$\Gamma \Proves
\lnot !A$} يا \LR{$\Gamma \Proves !B$} وي، نو \LR{$\Gamma \Proves
!A \lif !B$}۔ \end{enumerate} \end{prop}
\olpseganchor{OLP-0643-B018}{end}

\olpseganchor{OLP-0643-B019}{start}
\begin{proof}
\begin{enumerate}
\item فرض کړئ چې \LR{$\Gamma_0 \cup \{!A\} \Proves \lfalse$}۔
\olpseganchor{OLP-0643-B019}{end}

\olpseganchor{OLP-0643-B020}{start}
\begin{derivation} 
1. & \LR{$\Gamma_0 \cup \{!A\} \Proves \lfalse$} & فرض \\
2. & \LR{$\Gamma_1 \Proves !A$} & فرض\\ 
3. & \LR{$\Gamma_0 \cup \Gamma_1 \Proves \lfalse$} & قطع \\ 
\end{derivation}
\olpseganchor{OLP-0643-B020}{end}

\olpseganchor{OLP-0643-B021}{start}
ځکه چې \LR{$\Gamma_0 \subseteq \Gamma$} او \LR{$\Gamma_1 \subseteq \Gamma$} دي، \LR{$\Gamma_0 \cup \Gamma_1 \subseteq \Gamma$}؛ ځکه نو \LR{$\Gamma \Proves \lfalse$}۔
\olpseganchor{OLP-0643-B021}{end}

\olpseganchor{OLP-0643-B022}{start}
\item فرض کوو چې \LR{$\Gamma_0 \cup\{!A\} \Proves \lfalse$}۔
\olpseganchor{OLP-0643-B022}{end}

\olpseganchor{OLP-0643-B023}{start}
\begin{derivation} 
1. & \LR{$\Gamma_0 \cup\{!A\} \Proves \lfalse$} & فرض \\ 
2. & \LR{$\Gamma_0 \Proves !A \lif \lfalse$} & د استنتاج قضيه\\
3. & \LR{$\Gamma_0 \Proves (!A \lif \lfalse) \lif \lnot !A$} & Ax0\\ 
4. & \LR{$\Gamma_0 \Proves \lnot !A$} & مودوس پوننس 2, 3 \\ 
\end{derivation}
\olpseganchor{OLP-0643-B023}{end}

\olpseganchor{OLP-0643-B024}{start}
\item فرض کوو چې \LR{$\Gamma_0 \cup \{ !A \} \Proves \lfalse$} او \LR{$\Gamma_1 \cup \{\lnot !A \} \Proves \lfalse$} دي۔
\olpseganchor{OLP-0643-B024}{end}

\olpseganchor{OLP-0643-B025}{start}
\begin{derivation} 
1. & \LR{$\Gamma_0 \cup\{!A\} \Proves \lfalse$} & فرض \\
2. & \LR{$\Gamma_0 \Proves !A \lif \lfalse$} & د استنتاج قضيه \\ 
3. & \LR{$\Gamma_1 \cup \{\lnot !A \} \Proves \lfalse$} & فرض\\ 
4. & \LR{$\Gamma_0 \Proves (!A \lif \lfalse) \lif \lnot !A$} & Ax0\\ 
5. & \LR{$\Gamma_0 \Proves \lnot !A$} & مودوس پوننس 2, 4\\ 
6. & \LR{$\Gamma_0 \cup \Gamma_1 \Proves \lfalse$} & قطع 3, 5 \\
\end{derivation}
\olpseganchor{OLP-0643-B025}{end}

\olpseganchor{OLP-0643-B026}{start}
ځکه چې \LR{$\Gamma_0 \subseteq \Gamma$} او \LR{$\Gamma_1 \subseteq \Gamma$} دي، \LR{$\Gamma_0 \cup \Gamma_1 \subseteq \Gamma$}۔ ځکه نو \LR{$\Gamma \Proves \lfalse$}۔
\olpseganchor{OLP-0643-B026}{end}

\olpseganchor{OLP-0643-B027}{start}
% prop:provability-lor-left 
\item 
تمرين۔
\olpseganchor{OLP-0643-B027}{end}

\olpseganchor{OLP-0643-B028}{start}
% prop:provability-lor-right 
\item  
فرض کوو چې \LR{$\Gamma_0 \Proves !A$}۔
\olpseganchor{OLP-0643-B028}{end}

\olpseganchor{OLP-0643-B029}{start}
\begin{derivation} 
1. & \LR{$\Gamma_0 \Proves !A$} & فرض \\ 
2. & \LR{$\Gamma_0 \Proves !A \lif (!A \lor !B)$} & Ax0\\ 
3. & \LR{$\Gamma_0 \Proves !A \lor !B$} & مودوس پوننس 1, 2 \\
\end{derivation}
\olpseganchor{OLP-0643-B029}{end}

\olpseganchor{OLP-0643-B030}{start}
ځکه نو \LR{$\Gamma \Proves !A \lor !B$}۔ هغه ثبوت چې \LR{$\Gamma \Proves !B$} وي، دې ته ورته دے۔
\olpseganchor{OLP-0643-B030}{end}

\olpseganchor{OLP-0643-B031}{start}
% prop:provability-land-left 
\item  
فرض کړئ چې \LR{$\Gamma_0 \Proves !A \land !B$}
\olpseganchor{OLP-0643-B031}{end}

\olpseganchor{OLP-0643-B032}{start}
\begin{derivation} 
1. & \LR{$\Gamma_0 \Proves !A \land !B$} & فرض \\ 
2. & \LR{$\Gamma_0 \Proves (!A\land !B) \lif !A $} & Ax0\\ 
3. & \LR{$\Gamma_0 \Proves !A$} & مودوس پوننس 1, 2 \\
\end{derivation}
\olpseganchor{OLP-0643-B032}{end}

\olpseganchor{OLP-0643-B033}{start}
ځکه نو \LR{$\Gamma \Proves !A$}۔ دې ته ورته اشتقاق ښيي چې \LR{$\Gamma \Proves !B$}۔

\olpseganchor{OLP-0643-B033}{end}

\olpseganchor{OLP-0643-B034}{start}
% prop:provability-land-right
\item تمرين۔
\olpseganchor{OLP-0643-B034}{end}

\olpseganchor{OLP-0643-B035}{start}
% prop:provability-mp 
\item   فرض کوو چې \LR{$\Gamma_0 \Proves !A$} او \LR{$\Gamma_0 \Proves !A \lif !B $} دواړه دي۔
\olpseganchor{OLP-0643-B035}{end}

\olpseganchor{OLP-0643-B036}{start}
\begin{derivation} 
1. & \LR{$\Gamma_0 \Proves !A$} & فرض \\ 
2. & \LR{$\Gamma_0 \Proves !A \lif !B $} & فرض\\ 
3. & \LR{$\Gamma_0 \Proves !B$} & مودوس پوننس 1, 2\\
\end{derivation}
\olpseganchor{OLP-0643-B036}{end}

\olpseganchor{OLP-0643-B037}{start}
ځکه چې \LR{$\Gamma_0 \cup \Gamma_1 \subseteq \Gamma$}، له دې څخه \LR{$\Gamma \Proves !B$} راځي۔
\olpseganchor{OLP-0643-B037}{end}

\olpseganchor{OLP-0643-B038}{start}
% prop:provability-lif 
\item لومړے فرض کړئ چې \LR{$\Gamma \Proves \lnot !A$}۔
\olpseganchor{OLP-0643-B038}{end}

\olpseganchor{OLP-0643-B039}{start}
\begin{derivation} 
1. & \LR{$\Gamma_0 \Proves \lnot !A$} & فرض \\ 
2. & \LR{$\Gamma_0 \Proves \lnot !A \lif (!A \lif !B) $} & Ax0\\ 
3. & \LR{$\Gamma_0 \Proves !A \lif !B$} & مودوس پوننس 1,
2\\ \end{derivation}
\olpseganchor{OLP-0643-B039}{end}

\olpseganchor{OLP-0643-B040}{start}
د \LR{$\Gamma \Proves !B$} لپاره ثبوت دې ته ورته دے:
\olpseganchor{OLP-0643-B040}{end}

\olpseganchor{OLP-0643-B041}{start}
\begin{derivation} 
1. & \LR{$\Gamma_0 \Proves !B$} & فرض \\ 
2. & \LR{$\Gamma_0 \Proves !B \lif (!A \lif !B) $} & Ax0\\ 
3. & \LR{$\Gamma_0 \Proves !A \lif !B$} & مودوس پوننس 1, 2\\
\end{derivation}
\olpseganchor{OLP-0643-B041}{end}

\olpseganchor{OLP-0643-B042}{start}
\end{enumerate} 
\end{proof}
\olpseganchor{OLP-0643-B042}{end}



\olpseganchor{OLP-0643-B044}{start}
\begin{thm}[عمومي کول] \ollabel{thm:Generalization} که \LR{$\Gamma \Proves !A$} وي او \LR{$x$} د \LR{$\Gamma$} په هېڅ فارمول کښې ازاد نۀ وي، نو \LR{$\Gamma \Proves \lforall[x][!A]$}۔
\end{thm}
\olpseganchor{OLP-0643-B044}{end}

\olpseganchor{OLP-0643-B045}{start}
\begin{proof} پر \LR{$\mathsf{Thm}_\Gamma$} د استقرا له مخې: که \LR{$!A$} اصل وي، نو \LR{$\lforall[x][!A]$} هم اصل دے۔ که \LR{$!A\in \Gamma$} وي، نو \LR{$x$} په \LR{$!A$} کښې ازاد نۀ دے؛ ځکه نو \LR{$!A \lif \lforall[x][!A]$} يو اصل (\textbf{Ax3}) دے، او \LR{$\Gamma \Proves \lforall[x][!A]$} د مودوس پوننس له مخې۔ اوس فرض کړئ چې \LR{$!A$} د مودوس پوننس په مرسته راځي، ځکه \LR{$\Gamma \Proves !B$} او \LR{$\Gamma \Proves !B \lif !A$} دي۔ د استقرا د فرض له مخې \LR{$\Gamma \Proves \lforall[x][!B]$} او \LR{$\Gamma \Proves \lforall[x][( !B \lif !A)]$} دي۔ خو د \textbf{Ax2} له مخې دا هم دے:
\[ \Gamma \Proves
\lforall[x][( !B \lif !A)] \lif (\lforall[x][ !B] \lif \lforall[x][!A]), \]
او د مودوس پوننس دوه کارونې \LR{$\Gamma \Proves \lforall[x][!A]$} راکوي، لکه چې غوښتل شوي وو۔ \end{proof}
\olpseganchor{OLP-0643-B045}{end}

\olpseganchor{OLP-0643-B046}{start}
\begin{thm}\ollabel{thm:WeakGen} (\emph{پر ثابتونو کمزورے عمومي کول})
که \LR{$\Gamma \Proves !A$} وي او \LR{$c$} ثابت سمبول وي چې په \LR{$\Gamma$} کښې نۀ پېښېږي، نو داسې متغير \LR{$x$} شته چې په \LR{$!A$} کښې نۀ وي، داسې چې \LR{$\Gamma \Proves \lforall[x][\Subst{!A}{x}{c}]$}، او ثبوت \LR{$c$} نۀ کاروي۔
\end{thm}
\olpseganchor{OLP-0643-B046}{end}

\olpseganchor{OLP-0643-B047}{start}
\begin{proof} فرض کړئ \LR{$!A_1,\ldots,!A_n$} د \LR{$!A$} ثبوت له \LR{$\Gamma$} څخه وي، نو \LR{$!A=!A_n$}۔ داسې متغير \LR{$x$} واخلئ چې په \LR{$!A_1,\ldots,!A_n$} کښې نۀ وي، او نوې لړۍ په پام کښې ونيسئ: \LR{$\Subst{!A_1}{x}{c},\ldots,\Subst{!A_n}{x}{c}.$} داسې لړۍ د \LR{$\Subst{!A}{x}{c}$} ثبوت له \LR{$\Gamma$} څخه دے۔ په رښتيا، د هر \LR{$i=1,\ldots,n$} لپاره:
\begin{itemize}
\item که \LR{$!A_i$} اصل وي، نو \LR{$\Subst{!A_i}{x}{c}$} هم اصل دے؛
\item که \LR{$!A_i \in \Gamma$} وي، نو ځکه چې \LR{$c$} په \LR{$\Gamma$} کښې نۀ دے، \LR{$\Subst{!A_i}{x}{c} = !A_i \in \Gamma$} لرو؛
\item که \LR{$!A_i$} له \LR{$!A_j$} او \LR{$!A_j \lif !A_i$} څخه ترلاسه کېږي، نو \LR{$\Subst{!A_i}{x}{c}$} د مودوس پوننس په مرسته له \LR{$\Subst{!A_j}{x}{c}$} او \LR{$\Subst{(!A_j \lif !A_i)}{x}{c}
= \Subst{!A_j}{x}{c} \lif \Subst{!A_i}{x}{c}$} څخه راځي۔
\end{itemize}
ښکاره ده چې ثابت سمبول \LR{$c$} په نوې لړۍ کښې نور نۀ پېښېږي۔ اوس \LR{$\Gamma'$} دې هغه فارمولونه ولري چې له \LR{$\Gamma$} څخه د \LR{$\Subst{!A}{x}{c}$} په اشتقاق کښې راغلي دي؛ نو \LR{$x$} د \LR{$\Gamma'$} په هېڅ فارمول کښې ازاد نۀ دے او \LR{$\Gamma' \Proves \Subst{!A}{x}{c}$}۔ د عمومي کولو له مخې (\olref{thm:Generalization}) \LR{$\Gamma'
\Proves \lforall[x][\Subst{!A}{x}{c}]$}؛ له دې څخه د يکنواختۍ له مخې \LR{$\Gamma \Proves
\lforall[x][\Subst{!A}{x}{c}]$}، لکه چې غوښتل شوي وو۔
\end{proof}
\olpseganchor{OLP-0643-B047}{end}

\olpseganchor{OLP-0643-B048}{start}
\begin{explain}
هدف مو دا دے چې په \olref{thm:WeakGen} کښې دا شرط، چې \LR{$x$} په \LR{$!A$} کښې بالکل نۀ پېښېږي، په کمزورو شرطونو بدل کړو: \LR{$x$} په \LR{$!A$} کښې ازاد نۀ وي او د \LR{$c$} لپاره په \LR{$!A$} کښې ازاد وي۔ ښکاره ده چې دا د تړلي متغير په بدلولو ترلاسه کېدے شي؛ نو لومړے هماغه ثابتوو۔
\end{explain}
\olpseganchor{OLP-0643-B048}{end}

\olpseganchor{OLP-0643-B049}{start}
\begin{lem}
\ollabel{lem:free-for}
که \LR{$x$} د \LR{$c$} لپاره په \LR{$!A$} کښې ازاد وي او \LR{$y$} په \LR{$!A$} کښې ازاد نۀ وي، نو \LR{$x$} د \LR{$y$} لپاره په \LR{$\Subst{!A}{y}{c}$} کښې ازاد دے۔
\end{lem}
\olpseganchor{OLP-0643-B049}{end}

\olpseganchor{OLP-0643-B050}{start}
\begin{proof} که \LR{$x$} د \LR{$y$} لپاره په \LR{$\Subst{!A}{y}{c}$} کښې ازاد نۀ وي، نو د \LR{$y$} کومه ازاده پېښه په \LR{$\Subst{!A}{y}{c}$} کښې د کميت ټاکونکي \LR{$\lforall[x]$} د ساحې دننه راځي۔ خو د فرض له مخې \LR{$y$} په \LR{$!A$} کښې ازاد نۀ دے، نو ټولې داسې پېښې له تعويض \LR{$\Subst{}{y}{c}$} څخه راځي۔ نو د \LR{$c$} کومه پېښه د \LR{$\lforall[x]$} د ساحې دننه راځي او \LR{$x$} د \LR{$c$} لپاره په \LR{$!A$} کښې ازاد نۀ دے۔
\end{proof}
\olpseganchor{OLP-0643-B050}{end}

\olpseganchor{OLP-0643-B051}{start}
\begin{lem}[د تړلي متغير بدلول]
\ollabel{lem:ChangeBdVar}
که \LR{$x$} او \LR{$y$} په \LR{$!A$} کښې ازاد نۀ وي او دواړه د \LR{$c$} لپاره په \LR{$!A$} کښې ازاد وي، نو \LR{$\Proves \lforall[x][\Subst{!A}{x}{c}] \equiv
\lforall[y][\Subst{!A}{y}{c}]$}، او ثبوت \LR{$c$} نۀ کاروي۔
\end{lem}
\olpseganchor{OLP-0643-B051}{end}

\olpseganchor{OLP-0643-B052}{start}
\begin{proof}
له فرضونو څخه \LR{$y$} د \LR{$c$} لپاره په \LR{$!A$} کښې ازاد دے او \LR{$x$} په \LR{$!A$} کښې ازاد نۀ دے؛ نو د مخکينۍ \olref{lem:free-for} له مخې \LR{$y$} د \LR{$x$} لپاره په \LR{$\Subst{!A}{x}{c}$} کښې ازاد دے۔ له دې راځي چې
\LR{$\lforall[x][\Subst{!A}{x}{c}] \lif \Subst{!A}{x}{c} \Subst{}{y}{x}$}
يو اصل (\textbf{Ax1}) دے۔ ځکه چې \LR{$x$} په \LR{$!A$} کښې ازاد نۀ دے، \LR{$\Subst{!A}{x}{c} \Subst{}{y}{x} = \Subst{!A}{y}{c}$} لرو؛ نو
\[
\Proves \lforall[x][\Subst{!A}{x}{c}] \lif \Subst{!A}{y}{c},
\]
او د استنتاج د قضيې له مخې \LR{$\lforall[x][\Subst{!A}{x}{c}] \Proves
\Subst{!A}{y}{c}$}۔ ځکه چې \LR{$y$} په \LR{$!A$} کښې ازاد نۀ دے، په \LR{$\lforall[x][\Subst{!A}{x}{c}]$} کښې هم ازاد نۀ دے؛ نو د عمومي کولو له مخې
\LR{$\lforall[x][\Subst{!A}{x}{c}] \Proves \lforall[y][\Subst{!A}{y}{c}]$}،
او د استنتاج قضيه بيا \LR{$\Proves
\lforall[x][\Subst{!A}{x}{c}] \lif \lforall[y][\Subst{!A}{y}{c}]$} راکوي۔ د \LR{$\Proves \lforall[y][\Subst{!A}{y}{c}] \lif \lforall[x]
[\Subst{!A}{x}{c}]$} ثبوت عين متقارن دے، نو پايله د تاوتولوژۍ له مخې راځي۔
\end{proof}
\olpseganchor{OLP-0643-B052}{end}

\olpseganchor{OLP-0643-B053}{start}
\begin{thm}[پر ثابتونو قوي عمومي کول]
\ollabel{thm:StrongGen}
که \LR{$\Gamma \Proves !A$} وي، ثابت سمبول \LR{$c$} په \LR{$\Gamma$} کښې نۀ پېښېږي، او \LR{$x$} په \LR{$!A$} کښې ازاد نۀ وي خو د \LR{$c$} لپاره په \LR{$!A$} کښې ازاد وي، نو \LR{$\Gamma \Proves \lforall[x][\Subst{!A}{x}{c}]$}۔
\end{thm}
\olpseganchor{OLP-0643-B053}{end}

\olpseganchor{OLP-0643-B054}{start}
\begin{proof}
ځکه چې \LR{$\Gamma \Proves !A$} دے او \LR{$c$} په \LR{$\Gamma$} کښې نۀ دے، د کمزوري عمومي کولو له مخې داسې متغير \LR{$y$} شته چې په \LR{$!A$} کښې نۀ وي او \LR{$\Gamma \Proves \lforall[y][\Subst{!A}{y}{c}]$}۔ ځکه چې \LR{$y$} په \LR{$!A$} کښې بالکل نۀ دے، په \LR{$!A$} کښې ازاد نۀ دے او د \LR{$c$} لپاره په \LR{$!A$} کښې ازاد هم دے؛ که ورسره (د فرض له مخې) \LR{$x$} په \LR{$!A$} کښې ازاد نۀ دے او د \LR{$c$} لپاره په \LR{$!A$} کښې ازاد دے، نو د تړلي متغير د بدلولو شرطونه پوره دي۔ ځکه نو \LR{$\Proves \lforall[x][\Subst{!A}{x}{c}] \equiv
\lforall[y][\Subst{!A}{y}{c}]$}، له کومه چې \LR{$\Gamma \Proves
\lforall[x][\Subst{!A}{x}{c}]$} راځي۔
\end{proof}
\olpseganchor{OLP-0643-B054}{end}

\olpseganchor{OLP-0643-B055}{start}
\begin{thm}
\ollabel{thm:provability-quantifiers}
\begin{enumerate}
\item که \LR{$\Gamma \Proves !A(t)$} وي، نو \LR{$\Gamma \Proves
  \lexists[x][!A(x)]$}۔
\item که \LR{$\Gamma \Proves \lforall[x][!A(x)]$} وي، نو \LR{$\Gamma
  \Proves !A(t)$}۔
\end{enumerate}
دواړه حکمونه هغه وخت کاروو چې ترم د اړوند ازاد متغير پر ځاے په ازاد ډول تعويض کېدے شي۔ تړلے ترم دا شرط تل پوره کوي۔
\end{thm}
\olpseganchor{OLP-0643-B055}{end}

\olpseganchor{OLP-0643-B056}{start}
\begin{proof}
\begin{enumerate}
%%Need axioms for exists
\item تمرين۔
\item فرض کړئ چې \LR{$\Gamma_0 \Proves \lforall[x][!A(x)]$}۔
\olpseganchor{OLP-0643-B056}{end}

\olpseganchor{OLP-0643-B057}{start}
\begin{derivation} 
1. & \LR{$\Gamma_0 \Proves \lforall[x][!A(x)]$} & فرض \\
2. & \LR{$\Gamma_0 \Proves \lforall[x][!A(x)] \lif \Subst{!A}{t}{x}$} & Ax1 \\ 
3. & \LR{$\Gamma_0 \Proves \Subst{!A}{t}{x}$} & مودوس پوننس 1, 2 \\ 
\end{derivation}
\olpseganchor{OLP-0643-B057}{end}

\olpseganchor{OLP-0643-B058}{start}
\end{enumerate} 
\end{proof}
\olpseganchor{OLP-0643-B058}{end}

\olpseganchor{OLP-0126-B004}{start}
\olchapter{fol}{com}{د بشپړتيا قضيه}
\olpseganchor{OLP-0126-B004}{end}







\olpseganchor{OLP-0126-B008}{start}
%
  

\olpseganchor{OLP-0126-B008}{end}





\olpseganchor{OLP-0126-B011}{start}
%
  

\olpseganchor{OLP-0126-B011}{end}







\olpseganchor{OLP-0126-B015}{start}
%
  

\olpseganchor{OLP-0126-B015}{end}

\olpseganchor{OLP-0127-B004}{start}
\olfileid{fol}{com}{int}
\olpseganchor{OLP-0127-B004}{end}

\olpseganchor{OLP-0127-B005}{start}
\olsection{سريزه}
\olpseganchor{OLP-0127-B005}{end}

\olpseganchor{OLP-0127-B006}{start}
د بشپړتيا قضيه د منطق په اړه له تر ټولو بنسټيزو پايلو څخه يوه ده۔
دا په دوو بڼو وړاندې کېږي، او موږ به د دواړو همارزښت ثابت کړو۔ په
لومړۍ بڼه کښې دا زموږ د معنٰی پوهنې د پايلې او د اشتقاق
نظام ترمنځ د اړيکې په اړه يو بنسټيز حقيقت وايي: که
تړلے فارمول~\LR{$!A$} د ځينو تړلي فارمولونه~\LR{$\Gamma$} معنايي پايله
وي، نو داسې اشتقاق هم شته چې
\LR{$\Gamma \Proves !A$} ثابتوي۔ په دې توګه، د اشتقاق نظام
هغومره پياوړے دے څومره چې د هغو خبرو له ثابتولو پرته کېدے شي چې
په رښتيا نۀ لازمېږي۔
\olpseganchor{OLP-0127-B006}{end}

\olpseganchor{OLP-0127-B007}{start}
په دويمه بڼه کښې دا د مدل د شتون د پايلې په توګه ويل کېدے شي: د
تړلي فارمولونه هر سازګار سټ د صدق وړ دے۔ سازګاري يو ثبوت-تيوريکي
مفهوم دے: دا وايي چې زموږ د اشتقاق نظام د يو ټاکلي ډول
اشتقاقونه نۀ شي جوړولے۔ خو دا څوک تضمينوي چې يوازې ځکه چې
له~\LR{$\Gamma$} څخه د يو ټاکلي ډول هېڅ اشتقاقونه نشته، خامخا
داسې جوړښت~\LR{$\Struct{M}$}؟ د بشپړتيا د
قضيې تر لومړي ثبوت وړاندې---بلکې د هغو اشتقاق نظامونو تر
رامنځته کېدو هم وړاندې چې اوس ئې لرو---ستر جرمن رياضي‌پوه ډېوېډ
هيلبرټ دا نظر درلود چې د رياضي نظريو سازګاري د هغو څيزونو شتون
تضمينوي چې نظريې ئې په اړه دي۔ هغه ګوتلوب فرېګه ته د 1899 د دسمبر د 29مې نېټې په ليک کښې داسې وويل:
\begin{quote}
  که په خپله خوښه ټاکل شوي بديهي اصول، له خپلو ټولو پايلو سره،
  يو له بل سره په ټکر کښې نۀ وي، نو دغه اصول رښتوني دي او هغه
  څيزونه شته چې اصول يې تعريفوي۔ زما لپاره د رښتياوالي او شتون
  معيار همدا دے۔
\end{quote}
فرېګه سخت مخالفت وکړ۔ د بشپړتيا د قضيې دويمه بڼه ښيي چې هيلبرټ
لږ تر لږه په دې معنٰی سم و چې که بديهي اصول سازګار وي، نو داسې
\emph{کوم}
جوړښت شته چې ټول رښتوني کوي۔
\olpseganchor{OLP-0127-B007}{end}

\olpseganchor{OLP-0127-B008}{start}
دا يوازيني لاملونه نۀ دي چې د بشپړتيا قضيه---يا لا کره، د هغې
ثبوت---مهمه ده۔ دا يو شمېر مهمې پايلې لري چې ځينې به ئې جلا وڅېړو۔
د بېلګې په توګه، ځکه هر اشتقاق چې
\LR{$\Gamma \Proves !A$} ښيي متناهي دے، نو له~\LR{$\Gamma$} څخه يوازې
متناهي ډېر تړلي فارمولونه کارولے شي۔ له دې امله د بشپړتيا له قضيې
څخه لازمېږي چې که~\LR{$!A$} د~\LR{$\Gamma$} معنايي پايله وي، نو لا له وړاندې
د~\LR{$\Gamma$} د کوم متناهي فرعي سټ معنايي پايله ده۔ دې ته
\emph{د متناهي شاهد خاصيت} ويل کېږي۔ په همارز ډول، که د~\LR{$\Gamma$}
هر متناهي فرعي سټ سازګار وي، نو په خپله~\LR{$\Gamma$} هم سازګار دے۔
\olpseganchor{OLP-0127-B008}{end}

\olpseganchor{OLP-0127-B009}{start}
که څه هم د متناهي شاهد قضيه د بشپړتيا له قضيې څخه د
اشتقاقونه د منځنۍ لارې په تېرېدو لازمېږي، د بشپړتيا د قضيې
\emph{خپله ثبوت} هم نېغ په نېغه د هغې د ثابتولو دپاره کارول کېدے
شي۔ ځکه ثبوت د تړلي فارمولونه داسې سټ اخلي چې يو ټاکلے خاصيت---
سازګاري---لري، او له دې سټه داسې جوړښت جوړوي چې ټاکلي
خاصيتونه لري (په دې حالت کښې، سټ پوره کوي)۔ کابو همدا جوړونه د
متناهي شاهد خاصيت د نېغ ثبوت دپاره کارېږي، که د سازګارو سټونو پر
ځاے د تړلي فارمولونه له «په متناهي ډول د صدق وړ» سټونو څخه پيل
وکړو۔
دا جوړونه نورې پايلې هم ورکوي؛ د بېلګې په توګه، د
  تړلي فارمولونه هر د صدق وړ سټ يو متناهي يا شمېرېدونکے لامتناهي مدل
  لري۔ (دې پايلې ته د لوېنهايم--سکولم قضيه ويل کېږي۔) په عمومي ډول،
  د تړلي فارمولونه له سټونو څخه د جوړښتونه جوړول په منطق کښې
  ډېر ځله، او کله کله حتٰی په فلسفه کښې هم، کارول کېږي۔
\olpseganchor{OLP-0127-B009}{end}

\olpseganchor{OLP-0128-B004}{start}
\olfileid{fol}{com}{out}
\olpseganchor{OLP-0128-B004}{end}

\olpseganchor{OLP-0128-B005}{start}
\olsection{د ثبوت لنډيز}
\olpseganchor{OLP-0128-B005}{end}

\olpseganchor{OLP-0128-B006}{start}
د بشپړتيا د قضيې ثبوت يو څه پېچلے دے، او په لومړي لوست کښې ترې
لار ورکېدل اسانه دي۔ نو راځئ د ثبوت لنډيز وړاندې کړو۔ لومړے ګام د
نظر د زاويې داسې بدلون دے چې د ثبوت لاره راښيي۔ که بشپړتيا داسې
وګڼو چې «هر کله \LR{$\Gamma \Entails !A$} وي، نو
\LR{$\Gamma \Proves !A$} وي»، ښايي د کومې مفکورې موندل هم ګران وي:
ځکه د~\LR{$\Gamma \Proves !A$} د ښودلو دپاره بايد اشتقاق
ومومو، او داسې نۀ ښکاري چې فرضيه~\LR{$\Gamma \Entails !A$} په دې کښې
څه مرسته راسره کوي۔ د ځينو ثبوتي نظامونو دپاره اشتقاق
نېغ جوړېدے شي، خو موږ لږه بله لار نيسو۔ د نظر اړين بدلون دا دے:
بشپړتيا داسې هم ويل کېدے شي چې «که~\LR{$\Gamma$} سازګار وي، نو د صدق
وړ دے»۔ ښايي د~\LR{$\Gamma$} معلومات له دې فرضيې سره يو ځاے وکاروو چې
هغه سازګار دے، او داسې جوړښت
جوړ کړو چې په~\LR{$\Gamma$} کښې هر
تړلے فارمول پوره کړي۔ بالاخره موږ پوهېږو
چې د کوم ډول جوړښت په لټه کښې
يو: هغه چې هماغسې وي لکه~\LR{$\Gamma$} چې ئې بيانوي!{}
\olpseganchor{OLP-0128-B006}{end}

\olpseganchor{OLP-0128-B007}{start}
که~\LR{$\Gamma$} يوازې اتومي
  تړلي فارمولونه ولري، نو د هغې دپاره
مدل جوړول اسانه دي۔ فرض کړئ ټولې اتومي تړلي فارمولونه
  د~\LR{$\Atom{P}{a_1,\dots,a_n}$} بڼه لري، چې \LR{$a_i$} ګانې
  ثابت سمبولونه دي۔ يوازې دومره پکار دي چې
%
  د تعريف ساحه~\LR{$\Domain{M}$} او د~\LR{$P$} دپاره داسې ګومارنه پيدا کړو چې
  \LR{$\Sat{M}{\Atom{P}{a_1,\ldots,a_n}}$} وي۔ دا ډېره ګرانه نۀ ده:
  \LR{$\Domain{M} = \Nat$}، \LR{$\Assign{\Obj c_i}{M} = i$} وټاکئ، او د هر
  \LR{$\Atom{P}{a_1,\ldots,a_n} \in \Gamma$} دپاره جوړه
  \LR{$\tuple{k_1, \dots, k_n}$} په~\LR{$\Assign{P}{M}$} کښې واچوئ، چې
  \LR{$k_i$} د ثابت سمبول~\LR{$a_i$} شمېر دے (يعنې
  \LR{$a_i \ident \Obj c_{k_i}$})۔
\olpseganchor{OLP-0128-B007}{end}

\olpseganchor{OLP-0128-B008}{start}
اوس فرض کړئ~\LR{$\Gamma$} کوم داسې فارمول~\LR{$\lnot !B$} لري چې
\LR{$!B$} اتومي دے۔ ښايي اندېښنه وکړو چې د
\LR{$\Struct{M}$} جوړونه د~\LR{$\lnot !B$} د
رښتوني کولو له امکان سره ټکر کوي۔ خو دلته د~\LR{$\Gamma$} سازګاري
کار راځي: که \LR{$\lnot !B \in \Gamma$} وي، نو
\LR{$!B \notin \Gamma$}، کنه~\LR{$\Gamma$} به ناسازګار و۔ او که
\LR{$!B \notin \Gamma$} وي، نو زموږ د
\LR{$\Struct{M}$} د جوړونې له مخې
\LR{$\Sat/{M}{!B}$}؛ نو
\LR{$\Sat{M}{\lnot !B}$}۔ تر دې ځايه هر څه سم
دي۔
\olpseganchor{OLP-0128-B008}{end}

\olpseganchor{OLP-0128-B009}{start}
که~\LR{$\Gamma$} پېچلي، نااتومي فارمولونه ولري نو څه؟ فرض کړئ
\LR{$!A \land !B$} لري۔ د دې د رښتوني کولو دپاره بايد داسې پر مخ لاړ
شو لکه دواړه~\LR{$!A$} او~\LR{$!B$} چې په~\LR{$\Gamma$} کښې وي۔ او که
\LR{$!A \lor !B \in \Gamma$} وي، بايد لږ تر لږه يو ئې رښتونے کړو؛
يعنې داسې پر مخ لاړ شو لکه يو ئې چې په~\LR{$\Gamma$} کښې وي۔
\olpseganchor{OLP-0128-B009}{end}

\olpseganchor{OLP-0128-B010}{start}
دا لاندې مفکوره راښيي: په~\LR{$\Gamma$} کښې نور فارمولونه داسې
ورزياتوو چې (الف)~ترې جوړېدونکے سټ سازګار وساتو، (ب)~ډاډ ترلاسه
کړو چې د هرې ممکنې تړلے فارمول~\LR{$!A$} دپاره يا~\LR{$!A$} په
جوړ شوي سټ کښې وي يا~\LR{$\lnot !A$}، او (ج)~که هر کله
\LR{$!A \land !B$} په سټ کښې وي، دواړه~\LR{$!A$} او~\LR{$!B$} هم پکښې وي؛ که
\LR{$!A \lor !B$} په سټ کښې وي، لږ تر لږه يو له~\LR{$!A$} او~\LR{$!B$} څخه هم
پکښې وي؛ او داسې نور۔ موږ دې کار ته (که اړينه
وي، تر ابده) دوام
ورکوو۔ پيلنے سټ له ټولو ورزياتو شوو فارمولونه سره يو ځاے سټ~\LR{$\Gamma^*$} وبولئ۔
بيا به پورته جوړونه داسې
جوړښت~\LR{$\Struct{M}$}
راکړي چې په استقرا سره ثابتولے شو په~\LR{$\Gamma^*$} کښې ټولې جملې
پوره کوي، او ځکه چې \LR{$\Gamma \subseteq \Gamma^*$}، په~\LR{$\Gamma$}
کښې ټولې جملې هم پوره کوي۔ څرګندېږي چې د (الف) او~(ب) تضمين بس
دے۔ د جملو هغه سټ چې (ب) پوره کوي \emph{بشپړ} بلل کېږي۔ نو زموږ
کار به دا وي چې سازګار سټ~\LR{$\Gamma$} يو سازګار او بشپړ
سټ~\LR{$\Gamma^*$} ته وغځوو۔
\olpseganchor{OLP-0128-B010}{end}

\olpseganchor{OLP-0128-B011}{start}
%
په دې پلان کښې يوه ستونزه شته: که
\LR{$\lexists[x][!A(x)] \in \Gamma$} وي، غواړو په دې بهير کښې کوم
ثابت سمبول~\LR{$c$} وټاکو او~\LR{$!A(c)$} ورزيات کړو۔ خو څنګه پوهېږو چې
دا کار تل کېدے شي؟ ښايي زموږ په ژبه کښې يوازې څو ثابت سمبولونه
وي، او د هر يوه دپاره \LR{$\lnot !A(c) \in \Gamma$} ولرو۔ ورسره
\LR{$!A(c)$} هم نۀ شو زياتولے، ځکه سټ به ناسازګار کړي؛ بيا به نۀ پوهېدو
چې~\LR{$\Struct{M}$} بايد~\LR{$!A(c)$} رښتونے کړي که~\LR{$\lnot !A(c)$}۔ لا دا
هم کېدے شي چې~\LR{$\Gamma$} يوازې د داسې ژبې جملې ولري چې هېڅ ثابت
سمبول نۀ لري (د بېلګې په توګه، د سټونو د تيورۍ ژبه)۔
\olpseganchor{OLP-0128-B011}{end}

\olpseganchor{OLP-0128-B012}{start}
د دې ستونزې حل دا دے چې په پيل کښې بېخي نامتناهي ډېر ثابتونه، او
ورسره داسې جملې، زياتې کړو چې په سم ډول ئې له کميت ټاکونکو سره
وتړي۔ (البته بايد تصديق کړو چې دا کار ناسازګاري نۀ شي رامنځته
کولے۔)
\olpseganchor{OLP-0128-B012}{end}

\olpseganchor{OLP-0128-B013}{start}
زموږ لومړنۍ جوړونه هغه وخت ښه کار کوي چې په اتومي جملو کښې يوازې
ثابت سمبولونه ولرو۔ خو ژبه ښايي تابعې هم ولري۔ په هغه
حالت کښې ښايي پر~\LR{$\Nat$} د سمې تابعې موندل، چې دغو تابعې
ته ئې وګومارو او هر څه سم شي، ستونزمن وي۔ نو بله چاره دا ده: د
\LR{$\Obj c_i$} د تفسير دپاره د~\LR{$i$} د کارولو پر ځاے، د ثابت سمبولونه
خپله سټ د تعريف ساحه وټاکئ۔ بيا~\LR{$\Struct M$} هر ثابت سمبول خپل
ځان ته ګومارلے شي: \LR{$\Assign{\Obj c_i}{M} = \Obj c_i$}۔ خو ولې دا
کار تر پايه نۀ رسوو: د ژبې ټول \emph{ترمونه}~\LR{$\Domain{M}$} وټاکئ!{}
په دې صورت کښې تابعې ته د تابعو څرګنده ګومارنه شته (دغه
تابعې ترمونه د ارګومېنټونو په توګه اخلي او ترمونه د قيمتونو په توګه
لري): تابع~\LR{$\Obj f^n_i$} ته هغه تابع ګومارو چې \LR{$n$} ترمونه
\LR{$t_1$}، \dots،~\LR{$t_n$} د ننوت په توګه واخلي او ترم
\LR{$\Obj f^n_i(t_1, \dots, t_n)$} د قيمت په توګه ورکړي۔
\olpseganchor{OLP-0128-B013}{end}

\olpseganchor{OLP-0128-B014}{start}
د معما وروستۍ ټوټه دا ده چې له~\LR{$\eq$} سره څه وکړو۔
پريديکات~\LR{$\eq$} ثابت تفسير لري:
\LR{$\Sat{M}{\eq[t][t']}$} هله او يوازې هله چې
\LR{$\Value{t}{M} = \Value{t'}{M}$} وي۔ اوس که هر څه داسې وټاکو چې د
ترم~\LR{$t$} قيمت په خپله~\LR{$t$} وي، نو دا جوړښت به د
\LR{$\eq[t][t']$} په بڼه \emph{هېڅ} جمله رښتونې نۀ کړي، مګر دا چې
\LR{$t$} او~\LR{$t'$} يو او هماغه ترم وي۔ څرګنده ده چې دا ستونزه ده، ځکه د
تابعې لرونکې ژبې کابو هره په زړه پورې تيوري به داسې جملې
\LR{$\eq[t][t']$} د قضيو په توګه ولري چې~\LR{$t$} او~\LR{$t'$} يو ترم نۀ وي (د
بېلګې په توګه، د حساب په تيوريو کښې:
\LR{$\eq[(\Obj 0+ \Obj 0)][\Obj 0]$})۔ د دې ستونزې د هوارولو دپاره د
\LR{$\Struct M$} تعريف ساحه بدلوو: په~\LR{$\Domain{M}$} کښې د څيزونو په
توګه د ترمونو پر ځاے د ترمونو سټونه کاروو، او هر سټ ټول هغه ترمونه
لري چې د~\LR{$\Gamma$} جملې ئې برابر ګڼل غواړي۔ نو د بېلګې په توګه، که
\LR{$\Gamma$} د حساب تيوري وي، له دغو سټونو څخه يو به~\LR{$\Obj 0$}،
\LR{$(\Obj 0 + \Obj 0)$}، \LR{$(\Obj 0 \times \Obj 0)$}، او داسې نور ولري۔
همدا سټ به~\LR{$\Obj 0$} ته ګومارو، او څرګنده به شي چې همدا سټ په دې
کښې د ټولو ترمونو قيمت هم دے، لکه د~\LR{$(\Obj 0 + \Obj 0)$} قيمت۔
نو په دې بدله شوې جوړښت کښې به جمله
\LR{$\eq[(\Obj 0+ \Obj 0)][\Obj 0]$} رښتونې وي۔
\olpseganchor{OLP-0128-B014}{end}

\olpseganchor{OLP-0128-B015}{start}
نو دا کارونه به کوو۔ لومړے د بشپړ سازګارو سټونو خاصيتونه
څېړو؛ په ځانګړي ډول ثابتوو چې بشپړ سازګار سټ
\LR{$!A \land !B$} هله او يوازې هله لري چې دواړه~\LR{$!A$} او~\LR{$!B$} ولري،
او~\LR{$!A \lor !B$} هله او يوازې هله لري چې لږ تر لږه يو ئې ولري؛
او داسې نور (\olref[ccs]{prop:ccs})۔ بيا د جملو
  «مشبوع» سټونه تعريفوو او څېړو۔ مشبوع سټ داسې شرطي فارمولونه لري
  چې هره کميت ټاکونکې تړلے فارمول د هغې له بېلګو سره تړي
  (\olref[hen]{defn:henkin-exp})۔ ښيو چې هر سازګار سټ~\LR{$\Gamma$}
  تل مشبوع سټ~\LR{$\Gamma'$} ته غځېدے شي (\olref[hen]{lem:henkin})۔ که
  يو سټ سازګار، مشبوع او بشپړ وي، دا خاصيت هم لري چې
  \LR{$\lexists[x][!A(x)]$} هله او يوازې هله لري چې د کوم
    تړلي ترم~\LR{$t$} دپاره~\LR{$!A(t)$} ولري او
  \LR{$\lforall[x][!A(x)]$} هله او يوازې هله لري چې د
    ټولو تړلو ترمونو~\LR{$t$} دپاره~\LR{$!A(t)$} ولري
  (\olref[hen]{prop:saturated-instances})۔
بيا مشبوع سازګار سټ
\LR{$\Gamma'$} اخلو، او ښيو چې داسې
مشبوع، سازګار او بشپړ
سټ~\LR{$\Gamma^*$} ته غځېدے شي (\olref[lin]{lem:lindenbaum})۔ همدا
سټ~\LR{$\Gamma^*$} د خپل ترمي مدل~\LR{$\Struct M(\Gamma^*)$} د تعريف دپاره کاروو۔
د ترمي مدل تعريف ساحه د تړلو ترمونو سټ دے، او د هغه
  د پريديکاتونه تفسير د هغو اتومي
  تړلي فارمولونه له مخې ټاکل کېږي چې په~\LR{$\Gamma^*$} کښې دي
(\olref[mod]{defn:termmodel})۔ د مشبوع، بشپړو
سازګارو سټونو خاصيتونه کاروو، څو وښيو چې په رښتيا
\LR{$\Sat{M(\Gamma^*)}{!A}$} هله او
يوازې هله چې \LR{$!A \in \Gamma^*$}
(\olref[mod]{lem:truth})؛ نو په ځانګړي ډول
\LR{$\Sat{M(\Gamma^*)}{\Gamma}$}۔
په پاے کښې دا هم څېړو چې که~\LR{$\Gamma$}،~\LR{$\eq$} هم ولري،
  ترمي مدل څنګه تعريف کړو (\olref[ide]{defn:term-model-factor})،
  او ښيو چې هغه~\LR{$\Gamma^*$} پوره کوي (\olref[ide]{lem:truth})۔
\olpseganchor{OLP-0128-B015}{end}

\olpseganchor{OLP-0129-B005}{start}
\olfileid{fol}{com}{ccs}
\olpseganchor{OLP-0129-B005}{end}
      
\olpseganchor{OLP-0129-B006}{start}
\olsection{د تړلي فارمولونه بشپړ سازګار سټونه}
\olpseganchor{OLP-0129-B006}{end}

\olpseganchor{OLP-0129-B007}{start}
\begin{defn}[بشپړ سټ]
\ollabel{def:complete-set} د تړلي فارمولونه يو سټ~\LR{$\Gamma$} هله او
يوازې هله \emph{بشپړ} دے چې د هرې تړلے فارمول~\LR{$!A$}
دپاره يا~\LR{$!A \in \Gamma$} وي يا~\LR{$\lnot !A \in \Gamma$}۔
\end{defn}
\olpseganchor{OLP-0129-B007}{end}

\olpseganchor{OLP-0129-B008}{start}
\begin{explain}
د جملو بشپړ سټونه هېڅ پوښتنه بې ځوابه نۀ پرېږدي۔ د هرې
تړلے فارمول~\LR{$!A$} دپاره~\LR{$\Gamma$} «وايي» چې~\LR{$!A$} رښتونے دے که
کاذب۔ د بشپړ سټونو ارزښت د بشپړتيا د قضيې له ثبوت څخه هم
اوړي۔ د بېلګې په توګه، هغه تيوري چې بشپړ او د بديهي اصولو
له لارې څرګندېدونکې وي، تل د پرېکړې وړ وي۔
\end{explain}
\olpseganchor{OLP-0129-B008}{end}

\olpseganchor{OLP-0129-B009}{start}
\begin{explain}
بشپړ سازګار سټونه د بشپړتيا په ثبوت کښې مهم دي، ځکه
تضمينولے شو چې د تړلي فارمولونه هر سازګار سټ~\LR{$\Gamma$} په
بشپړ سازګار سټ~\LR{$\Gamma^*$} کښې شامل دے۔ بشپړ
سازګار سټ د هرې تړلے فارمول~\LR{$!A$} دپاره يا~\LR{$!A$} لري يا د هغې
نفي~\LR{$\lnot !A$}، خو دواړه نۀ لري۔ دا په ځانګړي ډول د
اتومي تړلي فارمولونه دپاره هم سم دي؛
نو له بشپړ سازګار سټ څخه په داسې ژبه کښې چې
ثابت سمبولونه په مناسبه توګه ورزيات شوي وي، داسې
جوړښت جوړولے شو چې پکښې د
پريديکاتونه تفسير د هغو
اتومي تړلي فارمولونه له
مخې تعريفېږي چې په~\LR{$\Gamma^*$} کښې دي۔ بيا ښودل کېدے شي چې دا
جوړښت په~\LR{$\Gamma^*$} کښې ټولې
تړلي فارمولونه (او ځکه په~\LR{$\Gamma$} کښې ټولې) رښتونې کوي۔ د دې
وروستي حقيقت ثبوت دا غواړي چې
\LR{$\lnot !A \in \Gamma^*$} هله او يوازې هله چې
\LR{$!A \notin \Gamma^*$}، او~\LR{$(!A \lor !B) \in \Gamma^*$} هله او
يوازې هله چې \LR{$!A \in \Gamma^*$} يا~\LR{$!B \in \Gamma^*$}؛ او داسې
نور۔
\end{explain}
\olpseganchor{OLP-0129-B009}{end}

\olpseganchor{OLP-0129-B010}{start}
په لاندې بحث کښې به ډېر ځله د~\LR{$\Proves$} د انعکاسيت، يکنواختۍ او
انتقاليت خاصيتونه په ناڅرګند ډول کاروو (وګورئ
%
  \cref{fol:seq:ptn:sec,fol:ntd:ptn:sec,fol:axd:ptn:sec,fol:tab:ptn:sec})۔
\olpseganchor{OLP-0129-B010}{end}

\olpseganchor{OLP-0129-B011}{start}
\begin{prop}
\ollabel{prop:ccs}
فرض کړئ~\LR{$\Gamma$} بشپړ او سازګار دے۔ بيا:
\begin{enumerate}
\item \ollabel{prop:ccs-prov-in} که \LR{$\Gamma \Proves !A$} وي، نو
  \LR{$!A \in \Gamma$}۔
\olpseganchor{OLP-0129-B011}{end}

\olpseganchor{OLP-0129-B012}{start}
\item \ollabel{prop:ccs-and} \LR{$!A \land !B \in \Gamma$}
  هله او يوازې هله چې دواړه \LR{$!A \in \Gamma$} او
  \LR{$!B \in \Gamma$} وي۔
\olpseganchor{OLP-0129-B012}{end}

\olpseganchor{OLP-0129-B013}{start}
\item \ollabel{prop:ccs-or} \LR{$!A \lor !B \in \Gamma$} هله
  او يوازې هله چې يا~\LR{$!A \in \Gamma$} وي يا~\LR{$!B \in \Gamma$}۔
\olpseganchor{OLP-0129-B013}{end}

\olpseganchor{OLP-0129-B014}{start}
\item \ollabel{prop:ccs-if} \LR{$!A \lif !B \in \Gamma$} هله
  او يوازې هله چې يا~\LR{$!A \notin \Gamma$} وي يا~\LR{$!B \in \Gamma$}۔
\end{enumerate}
\end{prop}
\olpseganchor{OLP-0129-B014}{end}

\olpseganchor{OLP-0129-B015}{start}
\begin{proof}
په لاندې ټولو برخو کښې فرض کوو چې~\LR{$\Gamma$} بشپړ او
سازګار دے۔
\begin{enumerate}
\item که \LR{$\Gamma \Proves !A$} وي، نو \LR{$!A \in \Gamma$}۔
\olpseganchor{OLP-0129-B015}{end}

\olpseganchor{OLP-0129-B016}{start}
فرض کړئ \LR{$\Gamma \Proves !A$}۔ د نقيض په توګه فرض کړئ چې
\LR{$!A \notin \Gamma$}۔ ځکه~\LR{$\Gamma$} بشپړ دے، نو
\LR{$\lnot !A \in \Gamma$}۔ د
%
  \cref{fol:seq:prv:prop:explicit-inc,fol:ntd:prv:prop:explicit-inc,fol:axd:prv:prop:explicit-inc,fol:tab:prv:prop:explicit-inc}
له مخې~\LR{$\Gamma$} ناسازګار دے۔ دا له دې فرضيې سره تناقض دے چې
\LR{$\Gamma$} سازګار دے۔ نو \LR{$!A \notin \Gamma$} نۀ شي کېدے، او ځکه
\LR{$!A \in \Gamma$}۔
\olpseganchor{OLP-0129-B016}{end}

\olpseganchor{OLP-0129-B017}{start}
\item %
%
\LR{$!A \land !B \in \Gamma$} هله او يوازې هله چې دواړه
\LR{$!A \in \Gamma$} او~\LR{$!B \in \Gamma$} وي:
\olpseganchor{OLP-0129-B017}{end}

\olpseganchor{OLP-0129-B018}{start}
د مخامخ لوري دپاره، فرض کړئ \LR{$!A \land !B \in \Gamma$}۔ بيا د
%
  \cref{fol:seq:ppr:prop:provability-land,fol:ntd:ppr:prop:provability-land,fol:axd:ppr:prop:provability-land,fol:tab:ppr:prop:provability-land}، د توکي~(۱) له
مخې \LR{$\Gamma \Proves !A$} او~\LR{$\Gamma \Proves !B$}۔ د
\olref{prop:ccs-prov-in} له مخې \LR{$!A \in \Gamma$} او
\LR{$!B \in \Gamma$}، لکه څنګه چې پکار وو۔
\olpseganchor{OLP-0129-B018}{end}

\olpseganchor{OLP-0129-B019}{start}
د عکس لوري دپاره، \LR{$!A \in \Gamma$} او~\LR{$!B \in \Gamma$} ونيسئ۔ د
%
  \cref{fol:seq:ppr:prop:provability-land,fol:ntd:ppr:prop:provability-land,fol:axd:ppr:prop:provability-land,fol:tab:ppr:prop:provability-land}، د توکي~(۲) له
مخې \LR{$\Gamma \Proves !A \land !B$}۔ د
\olref{prop:ccs-prov-in} له مخې \LR{$!A \land !B \in \Gamma$}۔
\olpseganchor{OLP-0129-B019}{end}

\olpseganchor{OLP-0129-B020}{start}
\item %
%
لومړے ښيو چې که \LR{$!A \lor !B \in \Gamma$} وي، نو يا
\LR{$!A \in \Gamma$} وي يا~\LR{$!B \in \Gamma$}۔ فرض کړئ
\LR{$!A \lor !B \in \Gamma$}، خو \LR{$!A \notin \Gamma$} او
\LR{$!B \notin \Gamma$}۔ ځکه~\LR{$\Gamma$} بشپړ دے، نو
\LR{$\lnot !A \in \Gamma$} او~\LR{$\lnot !B \in \Gamma$}۔ د
%
  \cref{fol:seq:ppr:prop:provability-lor,fol:ntd:ppr:prop:provability-lor,fol:axd:ppr:prop:provability-lor,fol:tab:ppr:prop:provability-lor}، د توکي~(۱) له
مخې~\LR{$\Gamma$} ناسازګار دے؛ دا تناقض دے۔ نو يا
\LR{$!A \in \Gamma$} وي يا~\LR{$!B \in \Gamma$}۔
\olpseganchor{OLP-0129-B020}{end}

\olpseganchor{OLP-0129-B021}{start}
د عکس لوري دپاره، فرض کړئ يا~\LR{$!A \in \Gamma$} وي يا
\LR{$!B \in \Gamma$}۔ د
%
  \cref{fol:seq:ppr:prop:provability-lor,fol:ntd:ppr:prop:provability-lor,fol:axd:ppr:prop:provability-lor,fol:tab:ppr:prop:provability-lor}، د توکي~(۲) له
مخې \LR{$\Gamma \Proves !A \lor !B$}۔ د
\olref{prop:ccs-prov-in} له مخې \LR{$!A \lor !B \in \Gamma$}، لکه
څنګه چې پکار وو۔
\olpseganchor{OLP-0129-B021}{end}

\olpseganchor{OLP-0129-B022}{start}
\item %
%
د مخامخ لوري دپاره، فرض کړئ \LR{$!A \lif !B \in \Gamma$}، او د نقيض
په توګه فرض کړئ چې \LR{$!A \in \Gamma$} او~\LR{$!B \notin \Gamma$}۔ د
دغو فرضونو له مخې \LR{$!A \lif !B \in \Gamma$} او
\LR{$!A \in \Gamma$}۔ د
%
  \cref{fol:seq:ppr:prop:provability-lif,fol:ntd:ppr:prop:provability-lif,fol:axd:ppr:prop:provability-lif,fol:tab:ppr:prop:provability-lif}، د توکي~(۱) له
مخې \LR{$\Gamma \Proves !B$}۔ خو بيا د
\olref{prop:ccs-prov-in} له مخې \LR{$!B \in \Gamma$}؛ دا له
\LR{$!B \notin \Gamma$} سره تناقض دے۔
\olpseganchor{OLP-0129-B022}{end}

\olpseganchor{OLP-0129-B023}{start}
د عکس لوري دپاره، لومړے هغه حالت وڅېړئ چې
\LR{$!A \notin \Gamma$} وي۔ ځکه~\LR{$\Gamma$} بشپړ دے، نو
\LR{$\lnot !A \in \Gamma$}۔ د
%
  \cref{fol:seq:ppr:prop:provability-lif,fol:ntd:ppr:prop:provability-lif,fol:axd:ppr:prop:provability-lif,fol:tab:ppr:prop:provability-lif}، د توکي~(۲) له
مخې \LR{$\Gamma \Proves !A \lif !B$}۔ بيا د
\olref{prop:ccs-prov-in} له مخې \LR{$!A \lif !B \in \Gamma$}، لکه
څنګه چې پکار وو۔
\olpseganchor{OLP-0129-B023}{end}

\olpseganchor{OLP-0129-B024}{start}
اوس هغه حالت وڅېړئ چې \LR{$!B \in \Gamma$} وي۔ بيا د
%
  \cref{fol:seq:ppr:prop:provability-lif,fol:ntd:ppr:prop:provability-lif,fol:tab:ppr:prop:provability-lif,fol:axd:ppr:prop:provability-lif}، يو ځل بيا د
توکي~(۲) له مخې \LR{$\Gamma \Proves !A \lif !B$}۔ د
\olref{prop:ccs-prov-in} له مخې \LR{$!A \lif !B \in \Gamma$}۔
\end{enumerate}
\end{proof}
\olpseganchor{OLP-0129-B024}{end}

\olpseganchor{OLP-0129-B025}{start}
\begin{prob}
د~\olref[fol][com][ccs]{prop:ccs} ثبوت بشپړ کړئ۔
\end{prob}
\olpseganchor{OLP-0129-B025}{end}

\olpseganchor{OLP-0130-B004}{start}
\olfileid{fol}{com}{hen}
\olpseganchor{OLP-0130-B004}{end}

\olpseganchor{OLP-0130-B005}{start}
\olsection{د هنکين غځونه}
\olpseganchor{OLP-0130-B005}{end}

\olpseganchor{OLP-0130-B006}{start}
\begin{explain}
د بشپړتيا د قضيې د ثابتولو د ستونزې يوه برخه دا ده چې هغه مدل چې
له بشپړ سازګار سټ~\LR{$\Gamma$} څخه ئې جوړوو بايد په~\LR{$\Gamma$} کښې ټول
کميت ټاکونکي لرونکي فارمولونه رښتوني کړي۔ د دې د تضمين دپاره د
ليون هنکين يو چل کاروو۔ په اصل کښې دا چل د ژبې په لامتناهي ډېرو
ثابت سمبولونه غځول او د هر داسې فارمول~\LR{$!A(x)$} دپاره چې يو
ازاد متغير لري، د لاندې بڼې يو فارمول ورزياتول دي:
\LR{$\lexists[x][!A(x)] \lif !A(c)$}،
چې~\LR{$c$} يو له نوو ثابت سمبولونه څخه دے۔ کله چې هغه جوړښت
جوړوو چې~\LR{$\Gamma$} پوره کوي، دا به تضمين کړي چې د نوو ثابت سمبولونو
په منځ کښې هره
رښتونې وجودي جمله يو شاهد ولري۔
\end{explain}
\olpseganchor{OLP-0130-B006}{end}

\olpseganchor{OLP-0130-B007}{start}
\begin{prop}
\ollabel{prop:lang-exp}
که~\LR{$\Gamma$} په~\LR{$\Lang L$} کښې سازګار وي او~\LR{$\Lang L'$} له~\LR{$\Lang L$}
څخه د نوو ثابت سمبولونه د~شمېرېدونکے لامتناهي سټ~\LR{$\Obj d_0$}،
\LR{$\Obj d_1$}، \dots په ورزياتولو ترلاسه شي، نو~\LR{$\Gamma$} په~\LR{$\Lang L'$}
کښې سازګار دے۔
\end{prop}
\olpseganchor{OLP-0130-B007}{end}

\olpseganchor{OLP-0130-B008}{start}
\begin{defn}[مشبوع سټ]
د ژبې~\LR{$\Lang {L}$} د فارمولونه يو سټ~\LR{$\Gamma$} هله او يوازې هله
\emph{مشبوع} دے چې د هر داسې فارمول~\LR{$!A(x) \in \Frm[L]$}
دپاره چې يو ازاد متغير~\LR{$x$} لري، داسې ثابت سمبول
\LR{$c \in \Lang{L}$} موجود وي چې
\LR{$\lexists[x][!A(x)] \lif !A(c) \in \Gamma$}۔
\end{defn}
\olpseganchor{OLP-0130-B008}{end}

\olpseganchor{OLP-0130-B009}{start}
لاندې تعريف به د ورپسې قضيې په ثبوت کښې وکارول شي۔
\olpseganchor{OLP-0130-B009}{end}

\olpseganchor{OLP-0130-B010}{start}
\begin{defn}
\ollabel{defn:henkin-exp}
فرض کړئ~\LR{$\Lang L'$} هماغسې وي لکه په~\olref{prop:lang-exp} کښې۔
د~\LR{$\Lang L'$} د ټولو هغو فارمولونه~\LR{$!A_i(x_i)$} يوه شمېرنه
\LR{$!A_0(x_0)$}، \LR{$!A_1(x_1)$}، \dots وټاکئ چې يو متغير~(\LR{$x_i$}) پکښې
ازاد راځي۔ تړلي فارمولونه~\LR{$!D_n$} پر~\LR{$n$} د استقرا له لارې تعريفوو۔
\olpseganchor{OLP-0130-B010}{end}

\olpseganchor{OLP-0130-B011}{start}
\LR{$c_0$} دې د هغو~\LR{$\Obj d_i$} په منځ کښې لومړنے ثابت سمبول وي چې
مونږ~\LR{$\Lang{L}$} ته ورزيات کړي دي او په~\LR{$!A_0(x_0)$} کښې نۀ راځي۔
په دې فرض چې~\LR{$!D_0$}، \dots،~\LR{$!D_{n-1}$} له مخکښې تعريف شوي دي،
\LR{$c_n$} دې د نوو ثابت سمبولونه~\LR{$\Obj d_i$} په منځ کښې لومړنے هغه وي
چې نۀ په~\LR{$!D_0$}، \dots،~\LR{$!D_{n-1}$} کښې راځي او نۀ په~\LR{$!A_n(x_n)$}
کښې۔
\olpseganchor{OLP-0130-B011}{end}

\olpseganchor{OLP-0130-B012}{start}
اوس~\LR{$!D_{n}$} دا فارمول وټاکئ: \LR{$\lexists[x_{n}][!A_{n}(x_{n})] \lif
  !A_{n}(c_{n})$}۔
\end{defn}
\olpseganchor{OLP-0130-B012}{end}

\olpseganchor{OLP-0130-B013}{start}
\begin{lem}
\ollabel{lem:henkin}
هر سازګار سټ~\LR{$\Gamma$} يوه مشبوع سازګار سټ~\LR{$\Gamma'$} ته غځېدے
شي۔
\end{lem}
\olpseganchor{OLP-0130-B013}{end}

\olpseganchor{OLP-0130-B014}{start}
\begin{proof}
د ژبې~\LR{$\Lang{L}$} د جملو يو سازګار سټ~\LR{$\Gamma$} راکړل
شوے دے؛ ژبه د نوو ثابت سمبولونه د~شمېرېدونکے لامتناهي سټ په
ورزياتولو وغځوئ، څو~\LR{$\Lang{L'}$} جوړ شي۔ د~\olref{prop:lang-exp}
له مخې~\LR{$\Gamma$} لا هم په بډايه ژبه کښې سازګار دے۔ همدارنګه،
\LR{$!D_i$} دې هماغسې وي لکه په~\olref{defn:henkin-exp} کښې۔ وټاکئ:
\begin{align*}
\Gamma_0 & = \Gamma \\
\Gamma_{n+1} & = \Gamma_n \cup \{!D_n \}
\end{align*}
يعنې~\LR{$\Gamma_{n+1} = \Gamma \cup \{ !D_0, \dots, !D_n \}$}، او
\LR{$\Gamma' = \bigcup_{n} \Gamma_n$} وټاکئ۔ \LR{$\Gamma'$} په څرګند ډول
مشبوع دے۔
\olpseganchor{OLP-0130-B014}{end}

\olpseganchor{OLP-0130-B015}{start}
که~\LR{$\Gamma'$} ناسازګار وے، نو د کوم~\LR{$n$} دپاره به~\LR{$\Gamma_n$}
ناسازګار وے (تمرين: څرګنده کړئ چې ولې)۔ نو د~\LR{$\Gamma'$} د سازګارۍ
د ښودلو دپاره دا بسنه کوي چې پر~\LR{$n$} د استقرا له لارې وښيو چې هر
سټ~\LR{$\Gamma_n$} سازګار دے۔
\olpseganchor{OLP-0130-B015}{end}

\olpseganchor{OLP-0130-B016}{start}
د استقرا بنسټ يوازې دا ادعا ده چې~\LR{$\Gamma_0 = \Gamma$} سازګار دے،
او دا د قضيې فرضيه ده۔ د استقرا د ګام دپاره، فرض کړئ چې
\LR{$\Gamma_{n}$} سازګار دے، خو~\LR{$\Gamma_{n+1} =
\Gamma_n \cup \{!D_n\}$} ناسازګار دے۔ در په ياد کړئ چې~\LR{$!D_n$} دا
دے:
\LR{$\lexists[x_{n}][!A_{n}(x_n)] \lif !A_{n}(c_{n})$}،
چې~\LR{$!A_n(x_n)$} د~\LR{$\Lang{L'}$} داسې فارمول دے چې يوازې متغير
\LR{$x_n$} ئې ازاد دے۔ د~\LR{$c_n$} د غوره کولو د طريقې له مخې (وګورئ
\olref{defn:henkin-exp})، \LR{$c_n$} نۀ په~\LR{$!A_n(x_n)$} کښې راځي او نۀ
په~\LR{$\Gamma_n$} کښې۔
\olpseganchor{OLP-0130-B016}{end}

\olpseganchor{OLP-0130-B017}{start}
که~\LR{$\Gamma_n \cup \{!D_n\}$} ناسازګار وي، نو
\LR{$\Gamma_n \Proves \lnot !D_n$}، او ځکه لاندې دواړه سم دي:

\[
\Gamma_n \Proves \lexists[x_n][!A_n(x_n)]
\qquad
\Gamma_n \Proves \lnot !A_n(c_n)
\]
ځکه چې~\LR{$c_n$} نۀ په~\LR{$\Gamma_n$} کښې راځي او نۀ په~\LR{$!A_n(x_n)$}
کښې، د قوي تعميم قضيه
\cref{fol:axd:qpr:thm:strong-generalization,fol:seq:qpr:thm:strong-generalization,fol:ntd:qpr:thm:strong-generalization,fol:tab:qpr:thm:strong-generalization}
پلې کېږي۔ له \LR{$\Gamma_n \Proves \lnot !A_n(c_n)$}
څخه دا ترلاسه کوو:
\LR{$\Gamma_n \Proves \lforall[x_n][\lnot !A_n(x_n)]$}۔
نو دواړه دا لرو:
\LR{$\Gamma_n \Proves \lexists[x_n][!A_n(x_n)]$} او
\LR{$\Gamma_n \Proves \lforall[x_n][\lnot !A_n(x_n)]$}؛
نو~\LR{$\Gamma_n$} پخپله ناسازګار دے۔
(پام وکړئ چې
\LR{$\lforall[x_n][\lnot !A_n(x_n)] \Proves
  \lnot\lexists[x_n][!A_n(x_n)]$}۔)
تناقض: فرض دا و چې~\LR{$\Gamma_n$} سازګار دے۔ نو
\LR{$\Gamma_n \cup \{ !D_n\}$} سازګار دے۔
\end{proof}
\olpseganchor{OLP-0130-B017}{end}

\olpseganchor{OLP-0130-B018}{start}
\begin{explain}
اوس به وښيو چې مشبوع، سازګار او \emph{بشپړ} سټونه دا
خاصيت لري چې يوه عمومي کميت ټاکونکې تړلے فارمول
  هله او يوازې هله لري چې د هغې ټولې بېلګې ولري
  او يوه وجودي کميت ټاکونکې تړلے فارمول هله او
  يوازې هله لري چې لږ تر لږه د هغې يوه بېلګه ولري۔ له دې څخه به
د دې ښودلو دپاره کار واخلو چې هغه جوړښت چې له مشبوع،
سازګار او بشپړ سټ څخه ئې جوړوو، د هغه ټولې کميت ټاکونکې جملې
رښتونې کوي۔
\end{explain}
\olpseganchor{OLP-0130-B018}{end}

\olpseganchor{OLP-0130-B019}{start}
\begin{prop}\ollabel{prop:saturated-instances}
فرض کړئ~\LR{$\Gamma$} بشپړ، سازګار او مشبوع دے۔
\begin{enumerate}
\item \LR{$\lexists[x][!A(x)] \in \Gamma$} هله او يوازې هله
  چې لږ تر لږه د يوه تړلي ترم~\LR{$t$} دپاره~\LR{$!A(t) \in \Gamma$} وي۔
\item \LR{$\lforall[x][!A(x)] \in \Gamma$} هله او يوازې هله
  چې د ټولو تړلو ترمونو~\LR{$t$} دپاره~\LR{$!A(t) \in \Gamma$} وي۔
\end{enumerate}
\end{prop}
\olpseganchor{OLP-0130-B019}{end}

\olpseganchor{OLP-0130-B020}{start}
\begin{proof}
  \begin{enumerate}
  \item %
    لومړے فرض کړئ چې
      \LR{$\lexists[x][!A(x)] \in \Gamma$}۔ ځکه~\LR{$\Gamma$} مشبوع دے، نو
      د کوم ثابت سمبول~\LR{$c$} دپاره
      \LR{$(\lexists[x][!A(x)] \lif !A(c)) \in \Gamma$}۔ د
      \cref{fol:axd:ppr:prop:provability-lif,fol:seq:ppr:prop:provability-lif,fol:ntd:ppr:prop:provability-lif,fol:tab:ppr:prop:provability-lif}، د توکي~(۱)،
      او \LR{\olref[ccs]{prop:ccs}(\olref[ccs]{prop:ccs-prov-in})} له مخې
      \LR{$!A(c) \in \Gamma$}۔
\olpseganchor{OLP-0130-B020}{end}

\olpseganchor{OLP-0130-B021}{start}
    د بل لوري دپاره مشبوع والی لازم نۀ دے: فرض کړئ
    \LR{$!A(t) \in \Gamma$}۔ بيا د
      \cref{fol:axd:qpr:prop:provability-quantifiers,fol:seq:qpr:prop:provability-quantifiers,fol:ntd:qpr:prop:provability-quantifiers,fol:tab:qpr:prop:provability-quantifiers}، د توکي~(۱) له مخې
    \LR{$\Gamma \Proves \lexists[x][!A(x)]$}۔ د
    \LR{\olref[ccs]{prop:ccs}(\olref[ccs]{prop:ccs-prov-in})} له مخې
    \LR{$\lexists[x][!A(x)] \in \Gamma$}۔
  \item %
    فرض کړئ چې د ټولو تړلو ترمونو~\LR{$t$}
      دپاره~\LR{$!A(t) \in \Gamma$}۔ د نقيض په توګه فرض کړئ چې
      \LR{$\lforall[x][!A(x)] \notin \Gamma$}۔ ځکه~\LR{$\Gamma$} بشپړ دے،
      نو~\LR{$\lnot\lforall[x][!A(x)] \in \Gamma$}۔ د مشبوع والي له مخې،
      د کوم ثابت سمبول~\LR{$c$} دپاره
      \LR{$(\lexists[x][\lnot !A(x)] \lif \lnot !A(c)) \in
        \Gamma$}۔ د فرض له مخې، ځکه چې~\LR{$c$} يو تړلے ترم دے، نو
      \LR{$!A(c) \in \Gamma$}۔ خو له دې څخه به~\LR{$\Gamma$} ناسازګار شي۔
      (تمرين: هغه اشتقاق ورکړئ چې ښيي
      \[
      \lnot \lforall[x][!A(x)],
      \lexists[x][\lnot !A(x)] \lif \lnot
        !A(c), !A(c)
      \]
      ناسازګار دے۔)
\olpseganchor{OLP-0130-B021}{end}

\olpseganchor{OLP-0130-B022}{start}
      د عکس لوري دپاره مشبوع والی لازم نۀ دے: فرض کړئ
      \LR{$\lforall[x][!A(x)] \in \Gamma$}۔ بيا د
      \cref{fol:axd:qpr:prop:provability-quantifiers,fol:seq:qpr:prop:provability-quantifiers,fol:ntd:qpr:prop:provability-quantifiers,fol:tab:qpr:prop:provability-quantifiers}،
      د توکي~(۲) له مخې~\LR{$\Gamma \Proves !A(t)$}۔ د
      \olref[ccs]{prop:ccs} له مخې~\LR{$!A(t) \in \Gamma$}۔
  \end{enumerate}
\end{proof}
\olpseganchor{OLP-0130-B022}{end}

\olpseganchor{OLP-0131-B004}{start}
\olfileid{fol}{com}{lin}
\olpseganchor{OLP-0131-B004}{end}

\olpseganchor{OLP-0131-B005}{start}
\olsection{د لينډنباوم لمه}
\olpseganchor{OLP-0131-B005}{end}

\olpseganchor{OLP-0131-B006}{start}
\begin{explain}
اوس يوه مرستندويه قضيه ثابتوو چې ښيي د تړلي فارمولونه هر سازګار
سټ د جملو په داسې سټ کښې شامل دے چې يوازې سازګار نۀ، بلکې
بشپړ هم وي۔ ثبوت په هر ځل د يوې تړلے فارمول په زياتولو
کار کوي، او په هر پړاو کښې تضمينوي چې سټ سازګار پاتې شي۔ دا ځکه
کوو چې د هر~\LR{$!A$} دپاره يا~\LR{$!A$} يا~\LR{$\lnot !A$} په کوم پړاو کښې
ورزيات شي۔ بيا د دې جوړونې د ټولو پړاوونو اتحاد يا~\LR{$!A$} لري يا
ئې نفي~\LR{$\lnot !A$}، او په دې توګه بشپړ دے۔ دا سازګار هم دے، ځکه
په هر پړاو کښې پام کوو چې ناسازګاري رامنځته نۀ کړو۔
\end{explain}
\olpseganchor{OLP-0131-B006}{end}

\olpseganchor{OLP-0131-B007}{start}
\begin{lem}[د لينډنباوم لمه]
\ollabel{lem:lindenbaum} په ژبه~\LR{$\Lang{L}$} کښې هر سازګار
سټ~\LR{$\Gamma$} داسې غځېدے شي چې يو بشپړ او سازګار
سټ~\LR{$\Gamma^*$} شي۔
\end{lem}
\olpseganchor{OLP-0131-B007}{end}

\olpseganchor{OLP-0131-B008}{start}
\begin{proof}
فرض کړئ~\LR{$\Gamma$} سازګار دے۔ د~\LR{$\Lang L$} د ټولو تړلي فارمولونه يوه
شمېرنه~\LR{$!A_0$}، \LR{$!A_1$}، \dots{} وي۔ \LR{$\Gamma_0 = \Gamma$} تعريف
کړئ، او
\[
\Gamma_{n+1} =
\begin{cases}
\Gamma_n \cup \{ !A_n \} & \textrm{که \LR{$\Gamma_n \cup \{!A_n\}$}
  سازګار وي؛} \\
\Gamma_n \cup \{ \lnot !A_n \} & \textrm{په بل صورت کښې۔}
\end{cases}
\]
\LR{$\Gamma^* = \bigcup_{n \geq 0} \Gamma_n$} وټاکئ۔
\olpseganchor{OLP-0131-B008}{end}

\olpseganchor{OLP-0131-B009}{start}
هر~\LR{$\Gamma_n$} سازګار دے: \LR{$\Gamma_0$} د تعريف او فرض له مخې سازګار
دے۔ که~\LR{$\Gamma_{n+1} = \Gamma_n \cup \{!A_n\}$} وي، لامل ئې دا
دے چې وروستنے سټ سازګار دے۔ که دا سازګار نۀ وي، نو
\LR{$\Gamma_{n+1} = \Gamma_n \cup \{\lnot !A_n\}$}۔ بايد تصديق کړو
چې~\LR{$\Gamma_n \cup \{\lnot !A_n\}$} سازګار دے۔ فرض کړئ چې سازګار
نۀ دے۔ بيا \emph{دواړه} \LR{$\Gamma_n \cup \{!A_n\}$} او
\LR{$\Gamma_n \cup \{\lnot !A_n\}$} ناسازګار دي۔ دا په دې معنا ده چې
\LR{$\Gamma_n$} به د
%
  \cref{fol:axd:prv:prop:provability-exhaustive,fol:seq:prv:prop:provability-exhaustive,fol:ntd:prv:prop:provability-exhaustive,fol:tab:prv:prop:provability-exhaustive}
له مخې ناسازګار وے، چې دا د استقرا له فرضيې سره تناقض دے۔
\olpseganchor{OLP-0131-B009}{end}

\olpseganchor{OLP-0131-B010}{start}
د هر~\LR{$n$} او هر~\LR{$i < n$} دپاره، \LR{$\Gamma_i \subseteq \Gamma_n$}۔
دا پر~\LR{$n$} د ساده استقرا له لارې ترلاسه کېږي۔ د~\LR{$n=0$} دپاره هېڅ
\LR{$i < 0$} نشته، نو ادعا په خپله سمه ده۔ د استقرا د ګام دپاره فرض
کړئ چې ادعا د~\LR{$n$} دپاره سمه ده۔ ښيو چې که~\LR{$i < n+1$} وي، نو
\LR{$\Gamma_i \subseteq \Gamma_{n+1}$}۔ د جوړونې له مخې يا
\LR{$\Gamma_{n+1} = \Gamma_n \cup \{!A_n\}$} ده يا
\LR{$= \Gamma_n \cup \{\lnot !A_n\}$}۔ نو
\LR{$\Gamma_n \subseteq \Gamma_{n+1}$}۔ که~\LR{$i < n+1$} وي، نو د
استقرا د فرضيې له مخې~\LR{$\Gamma_i \subseteq \Gamma_n$} (که~\LR{$i < n$}
وي)، يا دا ساده حقيقت لرو چې~\LR{$\Gamma_n \subseteq \Gamma_n$}
(که~\LR{$i = n$} وي)۔ د~\LR{$\subseteq$} د انتقاليت له مخې
\LR{$\Gamma_i \subseteq \Gamma_{n+1}$} ترلاسه کوو۔
\olpseganchor{OLP-0131-B010}{end}

\olpseganchor{OLP-0131-B011}{start}
له دې څخه څرګندېږي چې~\LR{$\Gamma^*$} سازګار دے۔ لامل دا دے: فرض کړئ
\LR{$\Gamma' \subseteq \Gamma^*$} متناهي دے۔
که دا فرعي سټ تش وي، د پيلني سازګار سټ فرعي سټ دے او ځکه سازګار دے۔
که نا تش وي، هر~\LR{$!B \in \Gamma'$} د
کوم~\LR{$i$} دپاره په~\LR{$\Gamma_i$} کښې هم دے۔ \LR{$n$} دې د دغو شاخصونو تر
ټولو لوے وي۔ ځکه که~\LR{$i \le n$} وي نو
\LR{$\Gamma_i \subseteq \Gamma_n$}، هر~\LR{$!B \in \Gamma'$} هم
\LR{$\in \Gamma_n$} دے؛ يعنې~\LR{$\Gamma' \subseteq \Gamma_n$}، او
\LR{$\Gamma_n$} سازګار دے۔ نو هر متناهي فرعي سټ
\LR{$\Gamma' \subseteq \Gamma^*$} سازګار دے۔ د
%
  \cref{fol:axd:ptn:prop:proves-compact,fol:seq:ptn:prop:proves-compact,fol:ntd:ptn:prop:proves-compact,fol:tab:ptn:prop:proves-compact} له مخې، \LR{$\Gamma^*$}
سازګار دے۔
\olpseganchor{OLP-0131-B011}{end}

\olpseganchor{OLP-0131-B012}{start}
د~\LR{$\Frm[L]$} هره تړلے فارمول په هغه لړ کښې راځي چې د~\LR{$\Gamma^*$}
د تعريف دپاره کارول شوے دے۔ که~\LR{$!A_n \notin \Gamma^*$} وي، لامل
ئې دا دے چې~\LR{$\Gamma_n \cup \{!A_n\}$} ناسازګار و۔ خو بيا
\LR{$\lnot !A_n \in \Gamma^*$}، نو~\LR{$\Gamma^*$} بشپړ دے۔
\end{proof}
\olpseganchor{OLP-0131-B012}{end}

\olpseganchor{OLP-0132-B004}{start}
\olfileid{fol}{com}{mod}
\olpseganchor{OLP-0132-B004}{end}

\olpseganchor{OLP-0132-B005}{start}
\olsection{د يو مدل جوړونه}
\olpseganchor{OLP-0132-B005}{end}

\olpseganchor{OLP-0132-B006}{start}
\begin{explain}
اوس مهال مو له~\LR{$\eq$} سره کار نشته؛ يعنې يوازې دا ښودل
  غواړو چې د~تړلي فارمولونه يو سازګار سټ~\LR{$\Gamma$} چې~\LR{$\eq$} نۀ لري،
  د صدق وړ دے۔ لومړے~\LR{$\Gamma$} يوه سازګار، بشپړ او مشبوع
  سټ~\LR{$\Gamma^*$} ته غځوو۔ په دې حالت کښې د مدل
  \LR{$\Struct{M(\Gamma^*)}$} تعريف ساده دے: د~\LR{$\Lang{L'}$} د تړلو
  ترمونو سټ د تعريف ساحه ګرځوو۔ هر ثابت سمبول پخپله ځان ته
  ورګومارو، او په عمومي ډول دا باوري کوو چې د هر تړلي ترم~\LR{$t$}
  دپاره~\LR{$\Value{t}{M(\Gamma^*)} = t$}۔ پريديکاتونه داسې تفسيروو
  چې يوه اتومي تړلے فارمول په~\LR{$\Struct{M(\Gamma^*)}$} کښې هله او
  يوازې هله رښتونې وي چې په~\LR{$\Gamma^*$} کښې وي۔ دا به څرګنده توګه
  په~\LR{$\Gamma^*$} کښې ټولې اتومي تړلي فارمولونه
  په~\LR{$\Struct{M(\Gamma^*)}$} کښې رښتونې کړي۔ نورې جملې هم رښتونې
  دي، په دې شرط چې هغه~\LR{$\Gamma^*$} چې ترې پيل کوو، سازګار، بشپړ او
  مشبوع وي۔
\end{explain}
\olpseganchor{OLP-0132-B006}{end}

\olpseganchor{OLP-0132-B007}{start}
%
\begin{defn}[ترمي مدل]
\ollabel{defn:termmodel}
فرض کړئ~\LR{$\Gamma^*$} په ژبه~\LR{$\Lang L$} کښې د~تړلي فارمولونه يو
بشپړ، سازګار او مشبوع سټ دے۔ د~\LR{$\Gamma^*$}
\emph{ترمي مدل}~\LR{$\Struct M(\Gamma^*)$} هغه جوړښت ده چې په
لاندې ډول تعريفېږي:
\begin{enumerate}
\item د د تعريف ساحه~\LR{$\Domain{M(\Gamma^*)}$} د~\LR{$\Lang L$} د ټولو
  تړلو ترمونو سټ دے۔
\item د ثابت سمبول~\LR{$c$} تفسير پخپله~\LR{$c$} دے:
  \LR{$\Assign{c}{M(\Gamma^*)} = c$}۔
\item تابع~\LR{$f$} ته هغه تابع ګومارل کېږي چې تړلي ترمونه
  \LR{$t_1$}، \dots،~\LR{$t_n$} د ارګومېنټونو په توګه واخلي او تړلے
  ترم~\LR{$f(t_1, \dots, t_n)$} د قيمت په توګه ولري:
\[
\Assign{f}{M(\Gamma^*)}(t_1, \dots, t_n) = f(t_1,\dots, t_n)
\]
\item که~\LR{$R$} يو \LR{$n$}-مقامي پريديکات وي، نو
  \[
  \tuple{t_1, \dots,
  t_n} \in \Assign{R}{M(\Gamma^*)} \text{\RL{ هله او يوازې هله چې }} \Atom{R}{t_1, \dots,
    t_n} \in \Gamma^*.
  \]
\end{enumerate}
\end{defn}

\olpseganchor{OLP-0132-B007}{end}

\olpseganchor{OLP-0132-B008}{start}
%
اوس به وګورو چې په رښتيا~\LR{$\Value{t}{M(\Gamma^*)} = t$} لرو۔
\olpseganchor{OLP-0132-B008}{end}

\olpseganchor{OLP-0132-B009}{start}
\begin{lem}
\ollabel{lem:val-in-termmodel} فرض کړئ~\LR{$\Struct M(\Gamma^*)$} د
\olref{defn:termmodel} ترمي مدل دے؛ بيا د هر تړلي ترم دپاره
\LR{$\Value{t}{M(\Gamma^*)} = t$}۔
\end{lem}
\begin{proof}
 ثبوت پر~\LR{$t$} د استقرا له لارې دے۔ د بنسټ حالت، چې~\LR{$t$}
 يو ثابت سمبول وي، نېغ په نېغه د ترمي مدل له تعريف څخه ترلاسه
 کېږي۔ د استقرا د ګام دپاره فرض کړئ~\LR{$t_1, \ldots, t_n$} داسې تړلي
 ترمونه دي چې~\LR{$\Value{t_i}{M(\Gamma^*)} = t_i$}، او~\LR{$f$} يوه
 \LR{$n$}-مقامي تابع ده۔ بيا
\begin{align*}
\Value{f(t_1,\ldots,t_n)}{M(\Gamma^*)} &= \Assign{f}{M(\Gamma^*)}(\Value{t_1}
 {M(\Gamma^*)},
\ldots, \Value{t_n}{M(\Gamma^*)}) \\
&= \Assign{f}{M(\Gamma^*)}(t_1, \dots, t_n) \\
&= f(t_1,\dots, t_n),
\end{align*}
او نو د استقرا له مخې دا د هر تړلي ترم~\LR{$t$} دپاره سم دي۔
\end{proof}

\olpseganchor{OLP-0132-B009}{end}

\olpseganchor{OLP-0132-B010}{start}
%
\begin{explain}
  جوړښت~\LR{$\Struct{M}$} ښائي يوه په وجودي ډول
    کميت ټاکونکې تړلے فارمول~\LR{$\lexists[x][!A(x)]$} رښتونې کړي،
    بې له دې چې کومه بېلګه~\LR{$!A(t)$} ئې رښتونې کړې وي۔
  جوړښت~\LR{$\Struct{M}$} ښائي د يوې په عمومي
    ډول کميت ټاکونکې تړلے فارمول~\LR{$\lforall[x][!A(x)]$} ټولې
    بېلګې~\LR{$!A(t)$} رښتونې کړي، خو پخپله~\LR{$\lforall[x][!A(x)]$}
    رښتونې نۀ کړي۔ لامل دا دے چې په عمومي ډول د~\LR{$\Domain{M}$}
  هر غړے د کوم تړلي ترم قيمت نۀ وي
  (\LR{$\Struct{M}$} ښائي پوښلے نۀ وي)۔ همدا لامل دے چې د صدق اړيکه د
  متغير-ګومارنو له لارې تعريفېږي۔ خو زموږ د ترمي مدل
  \LR{$\Struct{M(\Gamma^*)}$} دپاره دا اړينه نۀ ده، ځکه هغه پوښلے دے۔
  راتلونکې نتيجه همدا وايي۔
\end{explain}
\olpseganchor{OLP-0132-B010}{end}

\olpseganchor{OLP-0132-B011}{start}
\begin{prop}
\ollabel{prop:quant-termmodel}
فرض کړئ~\LR{$\Struct M(\Gamma^*)$} د~\olref{defn:termmodel} ترمي مدل دے۔
\begin{enumerate}
\item \LR{$\Sat{M(\Gamma^*)}{\lexists[x][!A(x)]}$} هله او يوازې هله چې
  لږ تر لږه د يوه تړلي ترم~\LR{$t$} دپاره~\LR{$\Sat{M(\Gamma^*)}{!A(t)}$} وي۔
\item \LR{$\Sat{M(\Gamma^*)}{\lforall[x][!A(x)]}$} هله او يوازې هله چې
  د ټولو تړلو ترمونو~\LR{$t$} دپاره~\LR{$\Sat{M(\Gamma^*)}{!A(t)}$} وي۔
\end{enumerate}
\end{prop}
\olpseganchor{OLP-0132-B011}{end}

\olpseganchor{OLP-0132-B012}{start}
\begin{proof}
\begin{enumerate}
\item %
  د~\olref[syn][ass]{prop:sat-quant} له مخې،
    \LR{$\Sat{M(\Gamma^*)}{\lexists[x][!A(x)]}$} هله او يوازې هله چې لږ
    تر لږه د يوې متغير-ګومارنې~\LR{$s$} دپاره
    \LR{$\Sat{M(\Gamma^*)}{!A(x)}[s]$} وي۔ ځکه
    \LR{$\Domain{M(\Gamma^*)}$} د~\LR{$\Lang{L}$} له تړلو ترمونو جوړ دے، دا
    هله او يوازې هله کېږي چې لږ تر لږه يو تړلے ترم~\LR{$t$} وي، داسې چې
    \LR{$s(x) = t$} او~\LR{$\Sat{M(\Gamma^*)}{!A(x)}[s]$}۔ د
    \olref[fol][syn][ext]{prop:ext-formulas} له مخې، چې~\LR{$s(x) = t$}
    وي، \LR{$\Sat{M(\Gamma^*)}{!A(x)}[s]$} هله او يوازې هله چې
    \LR{$\Sat{M(\Gamma^*)}{!A(t)}[s]$} وي۔ د
    \olref[fol][syn][ass]{prop:sentence-sat-true} له مخې،
    \LR{$\Sat{M(\Gamma^*)}{!A(t)}[s]$} هله او يوازې هله چې
    \LR{$\Sat{M(\Gamma^*)}{!A(t)}$} وي، ځکه~\LR{$!A(t)$} يوه جمله ده۔
\item %
  د~\olref[syn][ass]{prop:sat-quant} له مخې،
    \LR{$\Sat{M(\Gamma^*)}{\lforall[x][!A(x)]}$} هله او يوازې هله چې د
    هرې متغير-ګومارنې~\LR{$s$} دپاره~\LR{$\Sat{M(\Gamma^*)}{!A(x)}[s]$} وي۔
    راياد کړئ چې~\LR{$\Domain{M(\Gamma^*)}$} د~\LR{$\Lang{L}$} له تړلو
    ترمونو جوړ دے؛ نو د هر تړلي ترم~\LR{$t$} دپاره~\LR{$s(x) = t$} داسې يوه
    متغير-ګومارنه ده، او د هرې متغير-ګومارنې دپاره~\LR{$s(x)$} کوم تړلے
    ترم~\LR{$t$} دے۔ د~\olref[fol][syn][ext]{prop:ext-formulas} له مخې،
    چې~\LR{$s(x) = t$} وي، \LR{$\Sat{M(\Gamma^*)}{!A(x)}[s]$} هله او يوازې
    هله چې~\LR{$\Sat{M(\Gamma^*)}{!A(t)}[s]$} وي۔ د
    \olref[fol][syn][ass]{prop:sentence-sat-true} له مخې،
    \LR{$\Sat{M(\Gamma^*)}{!A(t)}[s]$} هله او يوازې هله چې
    \LR{$\Sat{M(\Gamma^*)}{!A(t)}$} وي، ځکه~\LR{$!A(t)$} يوه جمله ده۔
\end{enumerate}
\end{proof}

\olpseganchor{OLP-0132-B012}{end}



\olpseganchor{OLP-0132-B014}{start}
\begin{lem}[د صدق لمه]
\ollabel{lem:truth} فرض کړئ~\LR{$!A$}،~\LR{$\eq$} نۀ لري۔ بيا
\LR{$\Sat{M(\Gamma^*)}{!A}$} هله او يوازې هله چې \LR{$!A \in \Gamma^*$} وي۔
\end{lem}
\olpseganchor{OLP-0132-B014}{end}

\olpseganchor{OLP-0132-B015}{start}
\begin{proof}
دواړه لوري په يو وخت، او پر~\LR{$!A$} د استقرا له لارې ثابتوو۔
\begin{enumerate}
\item \indcase{!A}{\lfalse}
  {\LR{$\Sat/{M(\Gamma^*)}{\lfalse}$}
    د صدق د تعريف له مخې۔ له بلې خوا، ځکه~\LR{$\Gamma^*$} سازګار دے،
    نو~\LR{$\lfalse \notin \Gamma^*$}۔}
\olpseganchor{OLP-0132-B015}{end}

\olpseganchor{OLP-0132-B016}{start}
\item \indcase{!A}{\ltrue}
  {\LR{$\Sat{M(\Gamma^*)}{\ltrue}$}
    د صدق د تعريف له مخې۔ له بلې خوا، ځکه~\LR{$\Gamma^*$} سازګار او
    بشپړ دے، او~\LR{$\Gamma^* \Proves \ltrue$}، نو
    \LR{$\ltrue \in \Gamma^*$}۔}
\olpseganchor{OLP-0132-B016}{end}

\olpseganchor{OLP-0132-B017}{start}
\item \indcase{!A}{R(t_1, \dots, t_n)}
  {\LR{$\Sat{M(\Gamma^*)}{\Atom{R}{t_1, \dots, t_n}}$} هله او يوازې هله چې \LR{$\tuple{t_1,
      \dots, t_n} \in \Assign{R}{M(\Gamma^*)}$} وي (د صدق د تعريف
    له مخې)، او دا هله او يوازې هله چې
    \LR{$R(t_1, \dots, t_n) \in \Gamma^*$} وي (د~\LR{$\Struct
    M(\Gamma^*)$} د جوړونې له مخې)۔}
\olpseganchor{OLP-0132-B017}{end}

\olpseganchor{OLP-0132-B018}{start}
\item %
  \indcase{!A}{\lnot !B}
    {\LR{$\Sat{M(\Gamma^*)}{\indfrm}$} هله او يوازې هله چې
    \LR{$\Sat/{M(\Gamma^*)}{!B}$}
    وي (د صدق د تعريف له مخې)۔ د استقرا د فرضيې له مخې،
    \LR{$\Sat/{M(\Gamma^*)}{!B}$}
    هله او يوازې هله چې~\LR{$!B \notin \Gamma^*$} وي۔ ځکه~\LR{$\Gamma^*$}
    سازګار او بشپړ دے، \LR{$!B \notin \Gamma^*$} هله او يوازې
    هله چې~\LR{$\lnot !B \in \Gamma^*$} وي۔}
\olpseganchor{OLP-0132-B018}{end}

\olpseganchor{OLP-0132-B019}{start}
\item %
%
  \indcase{!A}{!B \land
    !C}{\LR{$\Sat{M(\Gamma^*)}{\indfrm}$}
    هله او يوازې هله چې دواړه
    \LR{$\Sat{M(\Gamma^*)}{!B}$} او
    \LR{$\Sat{M(\Gamma^*)}{!C}$} ولرو
    (د صدق د تعريف له مخې)، او دا هله او يوازې هله چې دواړه
    \LR{$!B \in \Gamma^*$} او~\LR{$!C \in \Gamma^*$} وي (د استقرا د فرضيې
    له مخې)۔ د~\LR{\olref[ccs]{prop:ccs}(\olref[ccs]{prop:ccs-and})} له
    مخې، دا هله او يوازې هله کېږي چې~\LR{$(!B \land !C) \in \Gamma^*$}۔}
\olpseganchor{OLP-0132-B019}{end}

\olpseganchor{OLP-0132-B020}{start}
\item %
%
  \indcase{!A}{!B \lor !C}
    {\LR{$\Sat{M(\Gamma^*)}{\indfrm}$}
    هله او يوازې هله چې
    \LR{$\Sat{M(\Gamma^*)}{!B}$} يا
    \LR{$\Sat{M(\Gamma^*)}{!C}$} وي
    (د صدق د تعريف له مخې)، او دا هله او يوازې هله چې
    \LR{$!B \in \Gamma^*$} يا~\LR{$!C \in \Gamma^*$} وي (د استقرا د فرضيې
    له مخې)۔ دا هله او يوازې هله کېږي چې
    \LR{$(!B \lor !C) \in \Gamma^*$} وي (د
    \LR{\olref[ccs]{prop:ccs}(\olref[ccs]{prop:ccs-or})} له مخې)۔}
\olpseganchor{OLP-0132-B020}{end}

\olpseganchor{OLP-0132-B021}{start}
\item %
%
  \indcase{!A}{!B \lif !C}
    {\LR{$\Sat{M(\Gamma^*)}{\indfrm}$}
    هله او يوازې هله چې~\LR{$\Sat/{M(\Gamma^*)}{!B}$} يا~\LR{$\Sat{M (\Gamma^*)}{!C}$} وي (د
    صدق د تعريف له مخې)، او دا هله او يوازې هله چې
    \LR{$!B \notin \Gamma^*$} يا~\LR{$!C \in \Gamma^*$} وي (د استقرا د فرضيې
    له مخې)۔ دا هله او يوازې هله کېږي چې
    \LR{$(!B \lif !C) \in \Gamma^*$} وي (د
    \LR{\olref[ccs]{prop:ccs}(\olref[ccs]{prop:ccs-if})} له مخې)۔}
\olpseganchor{OLP-0132-B021}{end}

\olpseganchor{OLP-0132-B022}{start}
%
\item %
  %
    \indcase{!A}{\lforall[x][!B(x)]}{\LR{$\Sat{M(\Gamma^*)}{\indfrm}$} هله او يوازې هله چې
      د ټولو تړلو ترمونو~\LR{$t$} دپاره~\LR{$\Sat{M(\Gamma^*)}{!B(t)}$} وي
      (\olref{prop:quant-termmodel})۔ د استقرا د فرضيې له مخې، دا
      هله او يوازې هله چې د ټولو تړلو ترمونو~\LR{$t$} دپاره
      \LR{$!B(t) \in \Gamma^*$} وي؛ د
      \olref[hen]{prop:saturated-instances} له مخې بيا دا هله او
      يوازې هله چې~\LR{$\lforall[x][!B(x)] \in \Gamma^*$} وي۔
      }
\olpseganchor{OLP-0132-B022}{end}

\olpseganchor{OLP-0132-B023}{start}
\item %
  %
    \indcase{!A}{\lexists[x][!B(x)]}{\LR{$\Sat{M(\Gamma^*)}{\indfrm}$} هله او يوازې هله چې
      لږ تر لږه د يوه تړلي ترم~\LR{$t$} دپاره
      \LR{$\Sat{M(\Gamma^*)}{!B(t)}$} وي
      (\olref{prop:quant-termmodel})۔ د استقرا د فرضيې له مخې، دا
      هله او يوازې هله چې لږ تر لږه د يوه تړلي ترم~\LR{$t$} دپاره
      \LR{$!B(t) \in \Gamma^*$} وي۔ د
      \olref[hen]{prop:saturated-instances} له مخې بيا دا هله او
      يوازې هله چې~\LR{$\lexists[x][!B(x)] \in \Gamma^*$} وي۔}

\end{enumerate}
\end{proof}
\olpseganchor{OLP-0132-B023}{end}

\olpseganchor{OLP-0133-B004}{start}
\olfileid{fol}{com}{ide}
\olsection{عينيت}
\olpseganchor{OLP-0133-B004}{end}

\olpseganchor{OLP-0133-B005}{start}
\begin{explain}
په مخکښې برخه کښې د ترمي مدل ورکړل شوې جوړونه د هغو سټونو
\LR{$\Gamma$} دپاره د اول مرتبې د منطق د بشپړوالي ثابتولو ته بس ده
چې~\LR{$\eq$} نۀ لري۔ ترمي مدل هر داسې~\LR{$!A \in \Gamma^*$} رښتونے کوي
چې~\LR{$\eq$} نۀ لري (او په دې توګه ټول~\LR{$!A \in \Gamma$} هم)۔ خو که
\LR{$\eq$} موجود وي، دا کار نۀ کوي۔ لامل دا دے چې بيا ښائي~\LR{$\Gamma^*$}
داسې يوه تړلے فارمول~\LR{$\eq[t][t']$} ولري، خو په ترمي مدل کښې د هر
ترم قيمت پخپله هماغه ترم دے۔ نو که~\LR{$t$} او~\LR{$t'$} بېل ترمونه وي، په
ترمي مدل کښې ئې قيمتونه، يعنې په ترتيب سره~\LR{$t$} او~\LR{$t'$}، هم بېل دي؛
ځکه~\LR{$\eq[t][t']$} ناسم دے۔ خو دا ستونزه د هغې جوړونې په کارولو
هوارولے شو چې «خارج قسمت جوړونه» بللے شي۔
\end{explain}
\olpseganchor{OLP-0133-B005}{end}

\olpseganchor{OLP-0133-B006}{start}
\begin{defn}
  فرض کړئ~\LR{$\Gamma^*$} په~\LR{$\Lang L$} کښې د جملو يو سازګار او
  بشپړ سټ دے۔ د~\LR{$\Lang L$} د تړلو ترمونو پر سټ اړيکه
  \LR{$\approx$} داسې تعريفوو:
  \[
  t \approx t' \text{\RL{\quad هله او يوازې هله چې \quad}} \eq[t][t'] \in \Gamma^*
  \]
\end{defn}
\olpseganchor{OLP-0133-B006}{end}

\olpseganchor{OLP-0133-B007}{start}
\begin{prop}
\ollabel{prop:approx-equiv}
اړيکه~\LR{$\approx$} لاندې خاصيتونه لري:
\begin{enumerate}
\item \LR{$\approx$} انعکاسي ده۔
\item \LR{$\approx$} متناظره ده۔
\item \LR{$\approx$} انتقالي ده۔
\item که~\LR{$t \approx t'$} وي،~\LR{$f$} يوه تابع وي، او
  \LR{$t_1$}، \dots،~\LR{$t_{i-1}$}، \LR{$t_{i+1}$}، \dots،~\LR{$t_n$} تړلي ترمونه وي، نو
\[
\Atom{f}{t_1,\dots, t_{i-1}, t, t_{i+1}, \dots, t_n} \approx
\Atom{f}{t_1,\dots, t_{i-1}, t', t_{i+1}, \dots, t_n}.
\]
\item که~\LR{$t \approx t'$} وي،~\LR{$R$} يو پريديکات وي، او
  \LR{$t_1$}، \dots،~\LR{$t_{i-1}$}، \LR{$t_{i+1}$}، \dots،~\LR{$t_n$} تړلي ترمونه وي، نو
\begin{multline*}
\Atom{R}{t_1,\dots, t_{i-1}, t, t_{i+1}, \dots, t_n} \in \Gamma^* \text{\RL{ هله او يوازې هله چې }} \\
\Atom{R}{t_1,\dots, t_{i-1}, t', t_{i+1}, \dots, t_n} \in \Gamma^*.
\end{multline*}
\end{enumerate}
\end{prop}
\olpseganchor{OLP-0133-B007}{end}

\olpseganchor{OLP-0133-B008}{start}
\begin{proof}
ځکه~\LR{$\Gamma^*$} سازګار او بشپړ دے، نو
\LR{$\eq[t][t'] \in \Gamma^*$} هله او يوازې هله چې
\LR{$\Gamma^* \Proves \eq[t][t']$} وي۔ نو د لاندې خبرو ښودل کافي دي:
\begin{enumerate}
\item د هر تړلي ترم~\LR{$t$} دپاره~\LR{$\Gamma^* \Proves \eq[t][t]$}۔
\item که~\LR{$\Gamma^* \Proves \eq[t][t']$} وي، نو
  \LR{$\Gamma^* \Proves \eq[t'][t]$}۔
\item که~\LR{$\Gamma^* \Proves \eq[t][t']$} او
  \LR{$\Gamma^* \Proves \eq[t'][t'']$} وي، نو
  \LR{$\Gamma^* \Proves \eq[t][t'']$}۔
\item که~\LR{$\Gamma^* \Proves \eq[t][t']$} وي، نو
\[
\Gamma^* \Proves
\eq[\Atom{f}{t_1,\dots,t_{i-1},t,t_{i+1},\dots,t_n}][\Atom{f}{t_1,\dots,t_{i-1},t',t_{i+1},\dots,t_n}]
\]
دا د هرې
\LR{$n$}-مقامي تابع~\LR{$f$} او تړلو ترمونو \LR{$t_1$}، \dots،
\LR{$t_{i-1}$}، \LR{$t_{i+1}$}، \dots،~\LR{$t_n$} دپاره سم دي۔
\item که~\LR{$\Gamma^* \Proves \eq[t][t']$} او
\LR{$\Gamma^* \Proves
\Atom{R}{t_1,\dots,t_{i-1},t,t_{i+1},\dots,t_n}$} وي، نو
\LR{$\Gamma^* \Proves \Atom{R}{t_1,\dots,t_{i-1},t',t_{i+1},\dots,t_n}$}
د هر \LR{$n$}-مقامي پريديکات~\LR{$R$} او تړلو ترمونو \LR{$t_1$}، \dots،
\LR{$t_{i-1}$}، \LR{$t_{i+1}$}، \dots،~\LR{$t_n$} دپاره۔
\end{enumerate}
\end{proof}
\olpseganchor{OLP-0133-B008}{end}

\olpseganchor{OLP-0133-B009}{start}
\begin{prob}
د~\olref[fol][com][ide]{prop:approx-equiv} ثبوت بشپړ کړئ۔
\end{prob}
\olpseganchor{OLP-0133-B009}{end}

\olpseganchor{OLP-0133-B010}{start}
\begin{defn}
فرض کړئ~\LR{$\Gamma^*$} په يوه ژبه~\LR{$\Lang L$} کښې يو سازګار او
بشپړ سټ دے،~\LR{$t$} يو تړلے ترم دے، او~\LR{$\approx$} د مخکښې
تعريف اړيکه ده۔ بيا:
\[
\equivrep{t}{\approx} = \Setabs{t'}{t'\in \Trm[L], t \approx t'}
\]
او~\LR{$\equivclass{\Trm[L]}{\approx} = \Setabs{\equivrep{t}{\approx}}{t \in \Trm[L]}$}۔
\end{defn}
\olpseganchor{OLP-0133-B010}{end}

\olpseganchor{OLP-0133-B011}{start}
\begin{defn}
\ollabel{defn:term-model-factor}
فرض کړئ~\LR{$\Struct M = \Struct M(\Gamma^*)$} د~\LR{$\Gamma^*$} هغه ترمي
مدل دے چې په~\olref[mod]{defn:termmodel} کښې راغلے۔ بيا
\LR{$\Struct{\equivclass{M}{\approx}}$} لاندې جوړښت ده:
\begin{enumerate}
\item \LR{$\Domain{\equivclass{M}{\approx}} = \equivclass{\Trm[L]}{\approx}$}۔
\item \LR{$\Assign{c}{\equivclass{M}{\approx}} = \equivrep{c}{\approx}$}۔
\item \LR{$\Assign{f}{\equivclass{M}{\approx}}(\equivrep{t_1}{\approx}, \dots,
  \equivrep{t_n}{\approx}) = \equivrep{\Atom{f}{t_1,\dots, t_n}}{\approx}$}۔
\item \LR{$\tuple{\equivrep{t_1}{\approx}, \dots, \equivrep{t_n}{\approx}} \in
  \Assign{R}{\equivclass{M}{\approx}}$} هله او يوازې هله چې
  \LR{$\Sat{M}{\Atom{R}{t_1,\dots, t_n}}$} وي؛ يعنې هله او يوازې هله چې
  \LR{$\Atom{R}{t_1,\dots, t_n} \in \Gamma^*$} وي۔
\end{enumerate}
\end{defn}
\olpseganchor{OLP-0133-B011}{end}

\olpseganchor{OLP-0133-B012}{start}
\begin{explain}
پام وکړئ چې~\LR{$\Assign{f}{\equivclass{M}{\approx}}$} او
\LR{$\Assign{R}{\equivclass{M}{\approx}}$} مو د
\LR{$\equivclass{\Trm[L]}{\approx}$} پر عناصرو داسې تعريف کړي چې هغوی
مو د~\LR{$\equivrep{t}{\approx}$} په بڼه ياد کړي؛ يعنې د
\emph{استازو}~\LR{$t \in \equivrep{t}{\approx}$} له لارې۔ بايد باوري
کړو چې دا تعريفونه د دې استازو پر انتخاب پورې تړلي نۀ دي؛ يعنې د
نورو داسې انتخابونو~\LR{$t'$} دپاره چې هماغه د هم ارزښتۍ ټولګي ټاکي
(\LR{$\equivrep{t}{\approx} = \equivrep{t'}{\approx}$})، تعريفونه هماغه
پايله ورکړي۔ د بېلګې په توګه، که~\LR{$R$} يو يو-مقامي پريديکات
وي، د تعريف وروستۍ فقره وايي چې
\LR{$\equivrep{t}{\approx} \in \Assign{R}{\equivclass{M}{\approx}}$} هله
او يوازې هله چې~\LR{$\Sat{M}{\Atom{R}{t}}$} وي۔ که د کوم بل ترم~\LR{$t'$} دپاره
چې~\LR{$t \approx t'$} وي،~\LR{$\Sat/{M}{\Atom{R}{t'}}$}، نو تعريف به وغواړي
چې~\LR{$\equivrep{t'}{\approx} \notin \Assign{R}{\equivclass{M}{\approx}}$}۔
که
\LR{$t \approx t'$} وي، نو~\LR{$\equivrep{t}{\approx} =
\equivrep{t'}{\approx}$}؛ خو نۀ شو کولے چې هم
\LR{$\equivrep{t}{\approx} \in \Assign{R}{\equivclass{M}{\approx}}$} او هم
\LR{$\equivrep{t}{\approx} \notin \Assign{R}{\equivclass{M}{\approx}}$}
ولرو۔ خو~\olref{prop:approx-equiv} تضمينوي چې دا نۀ شي پېښېدلے۔
\end{explain}
\olpseganchor{OLP-0133-B012}{end}

\olpseganchor{OLP-0133-B013}{start}
\begin{prop}
\LR{$\Struct{\equivclass{M}{\approx}}$} په ښه ډول تعريف شوے دے؛ يعنې که
\LR{$t_1$}، \dots،~\LR{$t_n$}، \LR{$t_1'$}، \dots،~\LR{$t_n'$} تړلي ترمونه وي، او د هرې
اړوندې جوړې دپاره~\LR{$t_i \approx t_i'$} وي، نو
\begin{enumerate}
\item \LR{$\equivrep{\Atom{f}{t_1,\dots, t_n}}{\approx} =
    \equivrep{\Atom{f}{t_1',\dots, t_n'}}{\approx}$}؛ يعنې
  \[
  \Atom{f}{t_1,\dots, t_n} \approx \Atom{f}{t_1',\dots, t_n'}
  \]
  او
\item \LR{$\Sat{M}{\Atom{R}{t_1,\dots, t_n}}$} هله او يوازې هله چې
  \LR{$\Sat{M}{\Atom{R}{t_1',\dots, t_n'}}$} وي؛ يعنې
  \[
    \Atom{R}{t_1,\dots, t_n} \in \Gamma^* \text{\RL{ هله او يوازې هله چې }}
    \Atom{R}{t_1',\dots, t_n'} \in \Gamma^*.
  \]
\end{enumerate}
\end{prop}
\olpseganchor{OLP-0133-B013}{end}

\olpseganchor{OLP-0133-B014}{start}
\begin{proof}
له~\olref{prop:approx-equiv} څخه پر~\LR{$n$} د استقرا له لارې ترلاسه کېږي۔
\end{proof}
\olpseganchor{OLP-0133-B014}{end}

\olpseganchor{OLP-0133-B015}{start}
د ترمي مدل د حالت په څېر، د صدق لمې تر ثبوت مخکښې لاندې لمې ته
اړتيا لرو۔
\olpseganchor{OLP-0133-B015}{end}

\olpseganchor{OLP-0133-B016}{start}
\begin{lem}
\ollabel{lem:val-in-termmodel-factored} فرض کړئ~\LR{$\Struct M = \Struct M
 (\Gamma^*)$}؛ بيا~\LR{$\Value{t}{\equivclass{M}{\approx}} = \equivrep{t}
 {\approx}$}۔
\end{lem}
\begin{proof}
ثبوت د~\olref[mod]{lem:val-in-termmodel} له ثبوت سره ورته دے۔
\end{proof}
\olpseganchor{OLP-0133-B016}{end}

\olpseganchor{OLP-0133-B017}{start}
\begin{prob}
د~\olref[fol][com][ide]{lem:val-in-termmodel-factored} ثبوت بشپړ کړئ۔
\end{prob}
\olpseganchor{OLP-0133-B017}{end}

\olpseganchor{OLP-0133-B018}{start}
\begin{lem}
\ollabel{lem:truth}
د ټولو جملو~\LR{$!A$} دپاره~\LR{$\Sat{\equivclass{M}{\approx}}{!A}$} هله او
يوازې هله چې~\LR{$!A \in \Gamma^*$} وي۔
\end{lem}
\olpseganchor{OLP-0133-B018}{end}

\olpseganchor{OLP-0133-B019}{start}
\begin{proof}
پر~\LR{$!A$} د استقرا له لارې، کټ مټ د~\olref[mod]{lem:truth} د ثبوت
په څېر۔ يوازې هغه حالت زيات پام غواړي چې~\LR{$!A \ident \eq[t][t']$} وي۔
\begin{align*}
\Sat{\equivclass{M}{\approx}}{\eq[t][t']} & \text{\RL{ هله او يوازې هله چې }} \equivrep{t}{\approx} = \equivrep{t'}{\approx}
\text{\RL{ وي (د \LR{$\Struct{\equivclass{M}{\approx}}$} د تعريف له مخې)}}\\
& \text{\RL{ هله او يوازې هله چې }} t \approx t' \text{\RL{ وي (د \LR{$\equivrep{t}{\approx}$} د تعريف له مخې)}}\\
& \text{\RL{ هله او يوازې هله چې }} \eq[t][t'] \in \Gamma^* \text{\RL{ وي (د \LR{$\approx$} د تعريف له مخې)۔}}
\end{align*}
\end{proof}
\olpseganchor{OLP-0133-B019}{end}

\olpseganchor{OLP-0133-B020}{start}
\begin{digress}
پام وکړئ: سره له دې چې~\LR{$\Struct{M(\Gamma^*)}$} تل د شمېر وړ او
نامتناهي دے،~\LR{$\Struct{\equivclass{M}{\approx}}$} ښائي متناهي وي، ځکه
کېدلے شي چې د~\LR{$\equivrep{t}{\approx}$} ټولګي يوازې متناهي شمېر وي۔
دا تمه کېدونکې ده، ځکه ښائي~\LR{$\Gamma$} داسې تړلي فارمولونه ولري چې هره
هغه جوړښت چې دوى پکښې رښتونې وي، بايد متناهي وي۔ د بېلګې
په توګه،~\LR{$\lforall[x][\lforall[y][x = y]]$} يوه سازګاره تړلے فارمول
ده، خو يوازې په هغو جوړښتونه کښې رښتونې کېږي چې د تعريف ساحه
ئې کټ مټ يو غړے لري۔
\end{digress}
\olpseganchor{OLP-0133-B020}{end}

\olpseganchor{OLP-0134-B004}{start}
\olfileid{fol}{com}{cth}
\olpseganchor{OLP-0134-B004}{end}

\olpseganchor{OLP-0134-B005}{start}
\olsection{د بشپړتيا قضيه}
\olpseganchor{OLP-0134-B005}{end}

\olpseganchor{OLP-0134-B006}{start}
\begin{explain}
راځئ خپلې پايلې سره يوځاے کړو: د بشپړتيا قضيې ته رسېږو۔
\end{explain}
\olpseganchor{OLP-0134-B006}{end}

\olpseganchor{OLP-0134-B007}{start}
\begin{thm}[د بشپړتيا قضيه]
\ollabel{thm:completeness}
فرض کړئ~\LR{$\Gamma$} د~تړلي فارمولونه يو سټ دے۔ که~\LR{$\Gamma$} سازګار وي،
نو د صدق وړ دے۔
\end{thm}
\olpseganchor{OLP-0134-B007}{end}

\olpseganchor{OLP-0134-B008}{start}
\begin{proof}
فرض کړئ~\LR{$\Gamma$} سازګار دے۔ د
  \olref[hen]{lem:henkin} له مخې، يو مشبوع سازګار سټ
  \LR{$\Gamma' \supseteq \Gamma$} شته۔ د~\olref[lin]{lem:lindenbaum} له
مخې، داسې~\LR{$\Gamma^* \supseteq \Gamma'$} شته چې
سازګار او بشپړ دے۔

  ځکه~\LR{$\Gamma' \subseteq \Gamma^*$}، نو د هرې
  فارمول~\LR{$!A(x)$} دپاره~\LR{$\Gamma^*$} په لاندې بڼه يوه تړلے فارمول
  لري:
  \LR{$\lexists[x][!A(x)] \lif !A(c)$}
  او په دې توګه~\LR{$\Gamma^*$} مشبوع دے۔ که~\LR{$\Gamma$}،~\LR{$\eq$} نۀ لري،
  نو د~\olref[mod]{lem:truth} له مخې،
\LR{$\Sat{M(\Gamma^*)}{!A}$} هله او
يوازې هله چې~\LR{$!A \in \Gamma^*$} وي۔ له دې څخه په تېره بيا دا ترلاسه
کېږي چې د ټولو~\LR{$!A \in \Gamma$} دپاره
\LR{$\Sat{M(\Gamma^*)}{!A}$}؛ نو
\LR{$\Gamma$} د صدق وړ دے۔
  که~\LR{$\Gamma$}،~\LR{$\eq$} ولري، نو د~\olref[ide]{lem:truth} له مخې د ټولو
  تړلي فارمولونه~\LR{$!A$} دپاره
  \LR{$\Sat{\equivclass{M}{\approx}}{!A}$} هله او يوازې هله چې
  \LR{$!A \in \Gamma^*$} وي۔ په تېره بيا، د ټولو~\LR{$!A \in \Gamma$} دپاره
  \LR{$\Sat{\equivclass{M}{\approx}}{!A}$}؛ نو~\LR{$\Gamma$} د صدق وړ دے۔
\end{proof}
\olpseganchor{OLP-0134-B008}{end}

\olpseganchor{OLP-0134-B009}{start}
\begin{cor}[د بشپړتيا قضيه، دويمه بڼه]
\ollabel{cor:completeness}
د ټولو~\LR{$\Gamma$} او تړلي فارمولونه~\LR{$!A$} دپاره: که
\LR{$\Gamma \Entails !A$} وي، نو~\LR{$\Gamma \Proves !A$}۔
\end{cor}
\olpseganchor{OLP-0134-B009}{end}

\olpseganchor{OLP-0134-B010}{start}
\begin{proof}
پام وکړئ چې په~\olref{cor:completeness} او
\olref{thm:completeness} کښې~\LR{$\Gamma$} ګانې په کلي ډول کميت ټاکل شوې
دي۔ د دې دپاره چې ځان ګډوډ نۀ کړو،~\olref{thm:completeness} د يوه
بېل متغير په کارولو بيا بيانوو: د~تړلي فارمولونه د هر سټ~\LR{$\Delta$}
دپاره، که~\LR{$\Delta$} سازګار وي، نو د صدق وړ دے۔ د نقيض عکس له مخې،
که~\LR{$\Delta$} د صدق وړ نۀ وي، نو~\LR{$\Delta$} ناسازګار دے۔ دا به د لازمې
نتيجې په ثبوت کښې وکاروو۔
\olpseganchor{OLP-0134-B010}{end}

\olpseganchor{OLP-0134-B011}{start}
فرض کړئ~\LR{$\Gamma \Entails !A$}۔ بيا د
\olref[syn][sem]{prop:entails-unsat} له مخې
\LR{$\Gamma \cup \{\lnot !A\}$} د صدق وړ نۀ دے۔ که
\LR{$\Gamma \cup \{\lnot !A\}$} خپل~\LR{$\Delta$} ونيسو، د
\olref{thm:completeness} مخکښې بڼه راکوي چې
\LR{$\Gamma \cup \{\lnot !A\}$} ناسازګار دے۔ د
%
  \cref{fol:axd:prv:prop:prov-incons,fol:seq:prv:prop:prov-incons,fol:ntd:prv:prop:prov-incons,fol:tab:prv:prop:prov-incons},
له مخې~\LR{$\Gamma \Proves !A$}۔
\end{proof}
\olpseganchor{OLP-0134-B011}{end}

\olpseganchor{OLP-0134-B012}{start}
\begin{prob}
\olref[fol][com][cth]{cor:completeness} وکاروئ چې
\olref[fol][com][cth]{thm:completeness} ثابت کړئ، او په دې توګه
وښيئ چې د بشپړتيا قضيې دواړه بيانونه هم ارزښته دي۔
\end{prob}
\olpseganchor{OLP-0134-B012}{end}



\olpseganchor{OLP-0134-B014}{start}
\begin{prob}
د دې دپاره چې د~اشتقاق نظام بشپړ وي، قاعدې ئې بايد دومره
پياوړې وي چې هر د صدق وړ نۀ سټ ناسازګار ثابت کړي۔ د بشپړتيا د ثبوت
دپاره د~اشتقاق کومې قاعدې اړينې وې؟ آيا له دې قاعدو څخه
کومه يوه د ثبوت په هېڅ ځاے کښې نۀ ده کارول شوې؟ د دې پوښتنو د ځواب
دپاره داسې يو لړ يا شکل جوړ کړئ چې وښيي د~اشتقاق کومې
قاعدې په هغو کومو پايلو کښې کارول شوې چې د
\olref[fol][com][cth]{thm:completeness} ثبوت ته رسېږي۔ په دې ثبوتونو
کښې د قاعدو هر ضمني استعمال هم خامخا ياد کړئ۔
\end{prob}
\olpseganchor{OLP-0134-B014}{end}

\olpseganchor{OLP-0135-B004}{start}
\olfileid{fol}{com}{com}
\olpseganchor{OLP-0135-B004}{end}

\olpseganchor{OLP-0135-B005}{start}
\olsection{د متناهي شاهد قضيه}
\olpseganchor{OLP-0135-B005}{end}

\olpseganchor{OLP-0135-B006}{start}
د بشپړتيا قضيې يوه مهمه پايله د متناهي شاهد قضيه ده۔ د متناهي
شاهد قضيه وايي چې که د~تړلي فارمولونه د يوه سټ هر
\emph{متناهي} فرعي سټ د صدق وړ وي، نو ټول سټ د صدق وړ دے، حتٰی
که سټ پخپله نامتناهي وي۔ دا هېڅ څرګنده خبره نۀ ده۔ لږ تر لږه په
لومړي نظر داسې هېڅ څۀ نۀ ښکاري چې د~تړلي فارمولونه د داسې نامتناهي
سټونو امکان رد کړي چې متناقض وي، خو تناقض ئې، په اصطلاح، يوازې له
نامتناهي شمېر څخه راولاړېږي۔ د متناهي شاهد قضيه وايي چې داسې حالت
ردېدلے شي: د~تړلي فارمولونه داسې د صدق وړ نۀ نامتناهي سټونه نشته
چې هر متناهي فرعي سټ ئې د صدق وړ وي۔ د بشپړتيا قضيې په څېر، د دې
يوه بڼه هم له استلزام سره تړاو لري: که د~تړلي فارمولونه يو نامتناهي
سټ کوم څۀ لازموي، نو لا له وړاندې ئې يو متناهي فرعي سټ لازموي۔
\olpseganchor{OLP-0135-B006}{end}

\olpseganchor{OLP-0135-B007}{start}
\begin{defn}
  د~فارمولونه يو سټ~\LR{$\Gamma$} هله او يوازې هله
  \emph{په متناهي ډول د صدق وړ} دے چې هر متناهي
  \LR{$\Gamma_0 \subseteq \Gamma$} د صدق وړ وي۔
\end{defn}
\olpseganchor{OLP-0135-B007}{end}

\olpseganchor{OLP-0135-B008}{start}
\begin{thm}[د متناهي شاهد قضيه]
\ollabel{thm:compactness}
د هرې د جملو ټولګې~\LR{$\Gamma$} او هرې جملې~\LR{$!A$} دپاره لاندې خبرې سمې دي:
\begin{enumerate}
  \item \LR{$\Gamma \Entails !A$} هله او يوازې هله چې يو متناهي
    \LR{$\Gamma_0 \subseteq \Gamma$} وي، داسې چې
    \LR{$\Gamma_0 \Entails !A$}۔
  \item \LR{$\Gamma$} هله او يوازې هله د صدق وړ دے چې په متناهي ډول
    د صدق وړ وي۔
\end{enumerate}
\end{thm}
\olpseganchor{OLP-0135-B008}{end}

\olpseganchor{OLP-0135-B009}{start}
\begin{proof}
موږ (2) ثابتوو۔ که~\LR{$\Gamma$} د صدق وړ وي، نو داسې
جوړښت~\LR{$\Struct{M}$}
شته چې د ټولو~\LR{$!A \in \Gamma$} دپاره
\LR{$\Sat{M}{!A}$}۔ څرګنده ده چې همدا
\LR{$\Struct{M}$} د~\LR{$\Gamma$} هر متناهي فرعي
سټ هم پوره کوي؛ نو~\LR{$\Gamma$} په متناهي ډول د صدق وړ دے۔
\olpseganchor{OLP-0135-B009}{end}

\olpseganchor{OLP-0135-B010}{start}
اوس فرض کړئ~\LR{$\Gamma$} په متناهي ډول د صدق وړ دے۔ بيا هر متناهي
فرعي سټ~\LR{$\Gamma_0 \subseteq \Gamma$} د صدق وړ دے۔ د سموالي له مخې
(%
  \cref{fol:axd:sou:cor:consistency-soundness,fol:seq:sou:cor:consistency-soundness,fol:ntd:sou:cor:consistency-soundness,fol:tab:sou:cor:consistency-soundness})،
هر متناهي فرعي سټ سازګار دے۔ بيا د لاندې پايلې له مخې~\LR{$\Gamma$}
پخپله هم بايد سازګار وي:
%
  \cref{fol:axd:ptn:prop:proves-compact,fol:seq:ptn:prop:proves-compact,fol:ntd:ptn:prop:proves-compact,fol:tab:ptn:prop:proves-compact}.
د بشپړتيا له مخې (\olref[cth]{thm:completeness})، ځکه~\LR{$\Gamma$}
سازګار دے، نو د صدق وړ دے۔
\end{proof}
\olpseganchor{OLP-0135-B010}{end}

\olpseganchor{OLP-0135-B011}{start}
\begin{prob}
د~\olref[fol][com][com]{thm:compactness} برخه (1) ثابته کړئ۔
\end{prob}
\olpseganchor{OLP-0135-B011}{end}



\olpseganchor{OLP-0135-B013}{start}
%
\begin{ex}
د يوې تيورۍ~\LR{$\Gamma$} په هر مدل~\LR{$\Struct{M}$} کښې، هر ترم~\LR{$t$} خامخا
د~\LR{$\Domain{M}$} يو غړے ټاکي۔ آيا دا هم تضمينولے شو چې د
\LR{$\Domain{M}$} هر غړے کوم ترم ټاکلے وي؟ په بله وينا، آيا داسې
تيورۍ~\LR{$\Gamma$} شته چې ټول مدلونه ئې پوښلي وي؟ د متناهي شاهد قضيه
ښيي چې که~\LR{$\Gamma$} نامتناهي مدلونه ولري، داسې نۀ ده۔ دا داسې
ليدلے شو: فرض کړئ~\LR{$\Struct{M}$} د~\LR{$\Gamma$} يو نامتناهي مدل دے، او
\LR{$c$} داسې يو ثابت سمبول دے چې د~\LR{$\Gamma$} په ژبه کښې نشته۔ فرض
کړئ~\LR{$\Delta$} د ټولو هغو جملو~\LR{$\eq/[c][t]$} سټ دے چې~\LR{$t$} د~\LR{$\Gamma$}
په ژبه~\LR{$\Lang{L}$} کښې يو تړلے ترم وي؛ يعنې
  \[
  \Delta = \Setabs{\eq/[c][t]}{t \in \Trm[L]}.
  \]
د~\LR{$\Gamma \cup \Delta$} يو متناهي فرعي سټ د
\LR{$\Gamma' \cup \Delta'$} په بڼه ليکل کېدے شي، چې
\LR{$\Gamma' \subseteq \Gamma$} او~\LR{$\Delta' \subseteq \Delta$} وي۔ ځکه
\LR{$\Delta'$} متناهي دے، يوازې متناهي ډېر ترمونه لرلے شي۔ داسې
\LR{$a \in \Domain{M}$} واخلئ چې د~\LR{$\Domain{M}$} يو غړے وي او له
هغوى څخه ئې هېڅ يو نۀ ټاکي؛ او~\LR{$\Struct{M'}$} داسې جوړښت
واخلئ چې کټ مټ د~\LR{$\Struct{M}$} په څېر وي، خو
\LR{$\Assign{c}{M'} = a$} هم ولري۔ ځکه په~\LR{$\Delta'$} کښې د راغلو ټولو
\LR{$t$} دپاره~\LR{$a \neq \Value{t}{M}$}، نو~\LR{$\Sat{M'}{\Delta'}$}۔ ځکه
\LR{$\Sat{M}{\Gamma}$}،~\LR{$\Gamma' \subseteq \Gamma$}، او~\LR{$c$} په~\LR{$\Gamma$}
کښې نۀ راځي، نو~\LR{$\Sat{M'}{\Gamma'}$} هم۔ دواړه سره يوځاے راکوي چې
د~\LR{$\Gamma \cup \Delta$} د هر متناهي فرعي سټ
\LR{$\Gamma' \cup \Delta'$} دپاره~\LR{$\Sat{M'}{\Gamma' \cup \Delta'}$}۔
نو د~\LR{$\Gamma \cup \Delta$} هر متناهي فرعي سټ د صدق وړ دے۔ د
متناهي شاهد قضيې له مخې،~\LR{$\Gamma \cup \Delta$} پخپله د صدق وړ دے۔
نو داسې مدلونه شته چې~\LR{$\Sat{M}{\Gamma \cup \Delta}$}۔ هر داسې
\LR{$\Struct{M}$} د~\LR{$\Gamma$} مدل دے، خو پوښلے نۀ دے، ځکه د~\LR{$\Lang{L}$}
د ټولو ترمونو~\LR{$t$} دپاره
\LR{$\Value{c}{M} \neq \Value{t}{M}$}۔
\end{ex}
\olpseganchor{OLP-0135-B013}{end}

\olpseganchor{OLP-0135-B014}{start}
\begin{ex}
داسې يوه ژبه~\LR{$\Lang{L}$} په پام کښې ونيسئ چې
پريديکات~\LR{$<$}، ثابت سمبولونه~\LR{$\Obj{0}$} او~\LR{$\Obj{1}$}، او
تابعې~\LR{$+$}،~\LR{$\times$} او~\LR{$-$} ولري۔ فرض کړئ~\LR{$\Gamma$} په دې
ژبه کښې د ټولو هغو تړلي فارمولونه سټ دے چې په
جوړښت~\LR{$\Struct{Q}$} کښې رښتونې دي، چې د تعريف ساحه ئې
\LR{$\Rat$} او تفسيرونه ئې څرګند دي۔~\LR{$\Gamma$} د~\LR{$\Lang{L}$} د ټولو هغو
تړلي فارمولونه سټ دے چې د ناطقو عددونو په اړه رښتونې دي۔ څرګنده ده
چې په~\LR{$\Rat$} (او حتٰی په~\LR{$\Real$}) کښې داسې عددونه~\LR{$r$} نشته چې
له~\LR{$0$} څخه لوی، خو د ټولو~\LR{$k \in \PosInt$} دپاره له~\LR{$1/k$} څخه
کوچني وي۔ که داسې عدد موجود واے، يو \emph{بې حده کوچنے عدد} به
و: له صفره بېل، خو بې حده کوچنے۔ د متناهي شاهد قضيه د دې ښودلو
دپاره کارېدلے شي چې د~\LR{$\Gamma$} داسې مدلونه شته چې بې حده کوچني
عددونه پکښې موجود دي۔ زموږ په ژبه کښې د تقسيم تابع نشته
(پر صفر تقسيم ناتعريف دے، او تابعې بايد په بشپړو تابعو
تفسير شي)۔ خو بيا هم څرګندولے شو چې~\LR{$r < 1/k$}، ځکه دا هله او يوازې
هله سم دي چې~\LR{$r \cdot k < 1$}۔ اوس~\LR{$c$} يو نوے ثابت سمبول ونيسئ او
\LR{$\Delta$} داسې وټاکئ:
\[
\{\Obj{0} < c\}
\cup \Setabs{ c \times \num k < \Obj{1} }{k \in \PosInt}
\]
(چې~\LR{$\num{k} = (\Obj{1} + (\Obj{1} + \dots + (\Obj{1} +
\Obj{1})\dots))$} کښې \LR{$k$} ځله~\LR{$\Obj{1}$} راغلے)۔ د~\LR{$\Delta$} د هر
متناهي فرعي سټ~\LR{$\Delta_0$} دپاره داسې~\LR{$K$} شته چې په~\LR{$\Delta_0$}
کښې د ټولو تړلي فارمولونه~\LR{$c \times \num{k} < \Obj{1}$} دپاره
\LR{$k < K$} وي۔ که~\LR{$\Struct{Q}$} داسې~\LR{$\Struct{Q'}$} ته پراخه کړو چې
\LR{$\Assign{c}{Q'} = 1/K$} ولري، نو د هر متناهي
\LR{$\Gamma_0 \subseteq \Gamma$} دپاره
\LR{$\Sat{Q'}{\Gamma_0 \cup \Delta_0}$} لرو؛ نو
\LR{$\Gamma \cup \Delta$} په متناهي ډول د صدق وړ دے (تمرين: دا په
تفصيل ثابت کړئ)۔ د متناهي شاهد قضيې له مخې،
\LR{$\Gamma \cup \Delta$} د صدق وړ دے۔ د~\LR{$\Gamma \cup \Delta$} هر
مدل~\LR{$\Struct{S}$} يو بې حده کوچنے عدد لري؛ يعنې~\LR{$\Assign{c}{S}$}۔
\end{ex}

\olpseganchor{OLP-0135-B014}{end}

\olpseganchor{OLP-0135-B015}{start}
\begin{prob}
د حساب په معياري مدل~\LR{$\Struct{N}$} کښې داسې هېڅ
غړے~\LR{$k \in \Domain{N}$} نشته چې هره صيغه~\LR{$\num{n} < x$}
پوره کړي (چې~\LR{$\num{n}$} هغه~\LR{$\Obj{0}^{\prime\dots\prime}$} دے چې
\LR{$n$} دانې~\LR{$\prime$} لري)۔ د متناهي شاهد قضيه وکاروئ او وښيئ چې د
حساب په ژبه کښې هغه~تړلي فارمولونه چې د حساب په معياري
مدل~\LR{$\Struct{N}$} کښې رښتونې دي، په داسې جوړښت~\LR{$\Struct{N'}$}
کښې هم رښتونې دي چې يو داسې غړے لري چې په رښتيا هره
صيغه~\LR{$\num{n} < x$} پوره کوي۔
\end{prob}
\olpseganchor{OLP-0135-B015}{end}

\olpseganchor{OLP-0135-B016}{start}

\begin{ex}
پوهېږو چې له~عينيت سره د اول مرتبې منطق دا څرګندولے شي چې
د~د تعريف ساحه کچه بايد لږ تر لږه ټاکلې اندازه ولري: جمله
\LR{$!A_{\ge n}$} (چې وايي «لږ تر لږه~\LR{$n$} بېل څيزونه شته») يوازې په
هغو جوړښتونو کښې رښتونې ده چې~\LR{$\Domain{M}$} لږ تر لږه~\LR{$n$} څيزونه
ولري۔ نو که داسې وټاکو:
  \[
  \Delta = \Setabs{!A_{\ge n}}{n \ge 1}
  \]
بيا د~\LR{$\Delta$} هر مدل بايد نامتناهي وي۔ په دې توګه يوې تيورۍ ته
د~\LR{$\Delta$} په ورزياتولو تضمينولے شو چې يوازې نامتناهي مدلونه
ولري: د~\LR{$\Gamma \cup \Delta$} مدلونه کټ مټ د~\LR{$\Gamma$} نامتناهي
مدلونه دي۔
\olpseganchor{OLP-0135-B016}{end}

\olpseganchor{OLP-0135-B017}{start}
نو د اول مرتبې منطق نامتناهي‌والے څرګندولے شي۔ خو د متناهي شاهد
قضيه ښيي چې متناهي‌والے نۀ شي څرګندولے۔ ځکه فرض کړئ د جملو کوم
سټ~\LR{$\Lambda$} په ټولو او يوازې متناهي جوړښتونه کښې رښتونے
وي۔ بيا~\LR{$\Delta \cup \Lambda$} په متناهي ډول د صدق وړ دے۔ ولې؟
فرض کړئ
\LR{$\Delta' \cup \Lambda' \subseteq \Delta \cup \Lambda$} متناهي دے،
چې~\LR{$\Delta' \subseteq \Delta$} او~\LR{$\Lambda' \subseteq \Lambda$} وي۔
که لومړے فرعي سټ خالي وي، د اندېکس ارزښت يو ټاکو؛ او که خالي نۀ وي،
\LR{$n$} دې تر ټولو لوی هغه عدد وي چې~\LR{$!A_{\ge n} \in \Delta'$}۔
ځکه~\LR{$\Lambda$} په ټولو متناهي جوړښتونو کښې
رښتونے دے، داسې مدل~\LR{$\Struct{M}$} لري چې متناهي ډېر، خو
\LR{$\ge n$}،~غړي ولري۔ خو بيا
\LR{$\Sat{M}{\Delta' \cup \Lambda'}$}۔ د متناهي شاهد قضيې له مخې،
\LR{$\Delta \cup \Lambda$} يو نامتناهي مدل لري، او دا له هغې فرضيې سره
تناقض لري چې~\LR{$\Lambda$} يوازې په متناهي جوړښتونه کښې رښتونے
دے۔
\end{ex}

\olpseganchor{OLP-0135-B017}{end}

\olpseganchor{OLP-0136-B004}{start}
\olfileid{fol}{com}{cpd}
\olpseganchor{OLP-0136-B004}{end}

\olpseganchor{OLP-0136-B005}{start}
\olsection{د متناهي شاهد قضيې نېغ ثبوت}
\olpseganchor{OLP-0136-B005}{end}

\olpseganchor{OLP-0136-B006}{start}
د بشپړتيا قضيې له کارولو پرته، د بشپړتيا قضيې د ثبوت د هماغو
نظرونو په کارولو، د متناهي شاهد قضيه نېغ په نېغه ثابتولے شو۔ د
بشپړتيا قضيې په ثبوت کښې مو د~تړلي فارمولونه له يوه سازګار
سټ~\LR{$\Gamma$} څخه پيل وکړ، هغه مو د~تړلي فارمولونه يوه
سازګار، مشبوع او بشپړ سټ~\LR{$\Gamma^*$} ته
پراخ کړ، او بيا مو وښودله چې له~\LR{$\Gamma^*$} څخه په جوړ شوي
ترمي
  مدل~\LR{$\Struct{M(\Gamma^*)}$}
کښې د~\LR{$\Gamma$} ټولې تړلي فارمولونه رښتونې دي؛ نو~\LR{$\Gamma$} د صدق وړ دے۔
\olpseganchor{OLP-0136-B006}{end}

\olpseganchor{OLP-0136-B007}{start}
همدا طريقه د دې ښودلو دپاره کارولے شو چې د جملو يو په متناهي ډول
د صدق وړ سټ، د صدق وړ دے۔ يوازې د هغو پايلو اړوندې بڼې ثابتول
پکار دي چې د صدق لمې ته رسېږي، او پکښې «سازګار» په «په متناهي
ډول د صدق وړ» بدلوو۔
\olpseganchor{OLP-0136-B007}{end}

\olpseganchor{OLP-0136-B008}{start}
\begin{prop}
\ollabel{prop:fsat-ccs}
فرض کړئ~\LR{$\Gamma$}،~بشپړ او په متناهي ډول د صدق وړ دے۔ بيا:
\begin{enumerate}
\item \LR{$(!A \land !B) \in \Gamma$}
  هله او يوازې هله چې دواړه \LR{$!A \in \Gamma$} او \LR{$!B \in \Gamma$} وي۔
\olpseganchor{OLP-0136-B008}{end}

\olpseganchor{OLP-0136-B009}{start}
\item \LR{$(!A \lor !B) \in \Gamma$} هله او يوازې هله چې
  يا~\LR{$!A \in \Gamma$} يا~\LR{$!B \in \Gamma$} وي۔
\olpseganchor{OLP-0136-B009}{end}

\olpseganchor{OLP-0136-B010}{start}
\item \LR{$(!A \lif !B) \in \Gamma$} هله او يوازې هله چې
  يا~\LR{$!A \notin \Gamma$} يا~\LR{$!B \in \Gamma$} وي۔
\end{enumerate}
\end{prop}
\olpseganchor{OLP-0136-B010}{end}

\olpseganchor{OLP-0136-B011}{start}
\begin{prob}
\olref[fol][com][cpd]{prop:fsat-ccs} ثابت کړئ۔~\LR{$\Proves$} مۀ کاروئ۔
\end{prob}
\olpseganchor{OLP-0136-B011}{end}



\olpseganchor{OLP-0136-B013}{start}
%
\begin{lem}
\ollabel{lem:fsat-henkin} هر په متناهي ډول د صدق وړ سټ~\LR{$\Gamma$}
يو مشبوع، په متناهي ډول د صدق وړ سټ~\LR{$\Gamma'$} ته غځېدلے شي۔
\end{lem}

\olpseganchor{OLP-0136-B013}{end}

\olpseganchor{OLP-0136-B014}{start}
\begin{prob}
\olref[fol][com][cpd]{lem:fsat-henkin} ثابت کړئ۔ (اشاره: مهم ګام
دا ښودل دي چې که~\LR{$\Gamma_n$} په متناهي ډول د صدق وړ وي، نو
\LR{$\Gamma_n \cup \{!D_n\}$} هم همداسې دے؛ بې له دې چې
اشتقاقونه يا سازګارۍ ته مراجعه وشي۔)
\end{prob}
\olpseganchor{OLP-0136-B014}{end}

\olpseganchor{OLP-0136-B015}{start}

\begin{prop}\ollabel{prop:fsat-instances}
فرض کړئ~\LR{$\Gamma$} بشپړ، په متناهي ډول د صدق وړ، او مشبوع دے۔
\begin{enumerate}
\item \LR{$\lexists[x][!A(x)] \in \Gamma$} هله او يوازې هله چې
لږ تر لږه د يوه تړلي ترم~\LR{$t$} دپاره~\LR{$!A(t) \in \Gamma$} وي۔
\item \LR{$\lforall[x][!A(x)] \in \Gamma$} هله او يوازې هله چې
د ټولو تړلو ترمونو~\LR{$t$} دپاره~\LR{$!A(t) \in \Gamma$} وي۔
\end{enumerate}
\end{prop}

\olpseganchor{OLP-0136-B015}{end}

\olpseganchor{OLP-0136-B016}{start}
\begin{prob}
\olref[fol][com][cpd]{prop:fsat-instances} ثابت کړئ۔
\end{prob}
\olpseganchor{OLP-0136-B016}{end}

\olpseganchor{OLP-0136-B017}{start}
\begin{lem}
\ollabel{lem:fsat-lindenbaum} هر په متناهي ډول د صدق وړ سټ~\LR{$\Gamma$}
يو بشپړ او په متناهي ډول د صدق وړ سټ~\LR{$\Gamma^*$} ته
غځېدلے شي۔
\end{lem}
\olpseganchor{OLP-0136-B017}{end}

\olpseganchor{OLP-0136-B018}{start}
\begin{prob}
\olref[fol][com][cpd]{lem:fsat-lindenbaum} ثابت کړئ۔ (اشاره: مهم
ګام دا ښودل دي چې که~\LR{$\Gamma_n$} په متناهي ډول د صدق وړ وي، نو يا
\LR{$\Gamma_n \cup \{!A_n\}$} يا~\LR{$\Gamma_n \cup \{\lnot !A_n\}$} په
متناهي ډول د صدق وړ دے۔)
\end{prob}
\olpseganchor{OLP-0136-B018}{end}



\olpseganchor{OLP-0136-B020}{start}
\begin{thm}[د متناهي شاهد خاصيت]
\ollabel{thm:compactness-direct} \LR{$\Gamma$} هله او يوازې هله د صدق
وړ دے چې په متناهي ډول د صدق وړ وي۔
\end{thm}
\olpseganchor{OLP-0136-B020}{end}

\olpseganchor{OLP-0136-B021}{start}
\begin{proof}
که~\LR{$\Gamma$} د صدق وړ وي، نو داسې
جوړښت~\LR{$\Struct{M}$}
شته چې د ټولو~\LR{$!A \in \Gamma$} دپاره
\LR{$\Sat{M}{!A}$} وي۔ څرګنده ده چې همدا
\LR{$\Struct{M}$} د~\LR{$\Gamma$} هر متناهي
فرعي سټ هم پوره کوي؛ نو~\LR{$\Gamma$} په متناهي ډول د صدق وړ دے۔
\olpseganchor{OLP-0136-B021}{end}

\olpseganchor{OLP-0136-B022}{start}
اوس فرض کړئ~\LR{$\Gamma$} په متناهي ډول د صدق وړ دے۔ د
  \olref{lem:fsat-henkin} له مخې، داسې په متناهي ډول د صدق وړ
  مشبوع سټ~\LR{$\Gamma' \supseteq \Gamma$} شته۔ د
\olref{lem:fsat-lindenbaum} له مخې،
\LR{$\Gamma'$} يو بشپړ او په متناهي ډول د
صدق وړ سټ~\LR{$\Gamma^*$} ته غځېدلے شي، او~\LR{$\Gamma^*$} لا
هم مشبوع دے۔ که عينيت موجود نۀ وي، د
ترمي مدل~\LR{$\Struct{M(\Gamma^*)}$}؛ او که عينيت موجود وي، د مخکښې
برخې خارج قسمت شوے ترمي مدل د
\olref[mod]{defn:termmodel} د جوړونې پر بنسټ جوړ کړئ۔
پام وکړئ چې~\olref[mod]{prop:quant-termmodel} د~\LR{$\Gamma^*$} پر
سازګارۍ (يا، په دې حساب، پر بشپړ والي يا مشبوع والي) تکيه
نۀ کوله؛ يوازې پر دې ئې تکيه کوله چې~\LR{$\Struct{M(\Gamma^*)}$} پوښلے
دے۔ د عينيت
په حالت کښې د هم ارزښتۍ او سازګارتيا شرطونه هم نېغ له متناهي صدق
وړتيا څخه ترلاسه کېږي: د هر شرط ناکامي به يو د صدق وړ نۀ متناهي
فرعي سټ ورکړي۔ د صدق لمې ثبوت
(\olref[mod]{lem:truth}) هغه وخت هم سم پر مخ ځي چې
\olref[ccs]{prop:ccs} او
\olref[hen]{prop:saturated-instances} ته مراجعې په ترتيب سره
\olref{prop:fsat-ccs} او~\olref{prop:fsat-instances} ته په مراجعو
بدلې کړو۔
د عينيت په حالت کښې د برابرۍ استقرايي حالت د خارج قسمت
له تعريف څخه نېغ ترلاسه کېږي۔
\end{proof}
\olpseganchor{OLP-0136-B022}{end}

\olpseganchor{OLP-0136-B023}{start}
\begin{prob}
د صدق لمې~(\olref[fol][com][mod]{lem:truth}) بشپړ ثبوت په هغه بڼه
وليکئ چې د~\olref[fol][com][cpd]{thm:compactness-direct} د ثبوت
دپاره اړينه ده۔
\end{prob}
\olpseganchor{OLP-0136-B023}{end}

\olpseganchor{OLP-0137-B004}{start}
\olfileid{fol}{com}{dls}
\olsection{د لوېنهايم--سکولم قضيه}
\olpseganchor{OLP-0137-B004}{end}

\olpseganchor{OLP-0137-B005}{start}
د لوېنهايم--سکولم قضيه وايي چې که يوه تيوري نامتناهي مدل ولري،
نو داسې يو مدل هم لري چې تر ډېره شمېرېدونکے لامتناهي وي۔ د دې خبرې يوه
سملاسي پايله دا ده چې د اول مرتبې منطق دا نۀ شي څرګندولے چې د
جوړښت کچه د شمېر نۀ وړ ده: هره تړلے فارمول يا د
تړلي فارمولونه هر سټ چې په ټولو د شمېر نۀ وړ جوړښتونه
کښې رښتونے وي، په کوم د شمېر وړ جوړښت کښې هم رښتونے دے۔
\olpseganchor{OLP-0137-B005}{end}

\olpseganchor{OLP-0137-B006}{start}
\begin{thm}
\ollabel{thm:downward-ls} که~\LR{$\Gamma$} سازګار وي، نو يو
د شمېر وړ مدل لري؛ يعنې په داسې جوړښت کښې د صدق وړ
دے چې د تعريف ساحه ئې يا متناهي يا شمېرېدونکے لامتناهي وي۔
\end{thm}
\olpseganchor{OLP-0137-B006}{end}

\olpseganchor{OLP-0137-B007}{start}
\begin{proof}
که~\LR{$\Gamma$} سازګار وي، د بشپړتيا قضيې ثبوت چې کوم
جوړښت~\LR{$\Struct M$} ورکوي، د هغه د تعريف ساحه~\LR{$\Domain{M}$}
د ژبې~\LR{$\Lang L$} د ترمونو له سټ څخه لويه نۀ ده۔ نو~\LR{$\Struct M$} تر
ډېره شمېرېدونکے لامتناهي دے۔
\end{proof}
\olpseganchor{OLP-0137-B007}{end}

\olpseganchor{OLP-0137-B008}{start}
\begin{thm}
\ollabel{noidentity-ls} که~\LR{$\Gamma$} د عينيت پرته د اول مرتبې د
منطق په ژبه کښې د~تړلي فارمولونه يو سازګار سټ وي، نو يو
شمېرېدونکے لامتناهي مدل لري؛ يعنې په داسې جوړښت کښې د صدق وړ
دے چې د تعريف ساحه ئې نامتناهي او د شمېر وړ وي۔
\end{thm}
\olpseganchor{OLP-0137-B008}{end}

\olpseganchor{OLP-0137-B009}{start}
\begin{proof}
که~\LR{$\Gamma$} سازګار وي او داسې هېڅ جمله نۀ لري چې عينيت پکښې
راغلے وي، نو د بشپړتيا قضيې ثبوت چې کوم
جوړښت~\LR{$\Struct M$} ورکوي، د هغه د تعريف ساحه~\LR{$\Domain{M}$}
کټ مټ د ژبې~\LR{$\Lang L'$} د ترمونو سټ دے۔ نو~\LR{$\Struct{M}$}
شمېرېدونکے لامتناهي دے، ځکه~\LR{$\Trm[L']$} هم داسې دے۔
\end{proof}
\olpseganchor{OLP-0137-B009}{end}

\olpseganchor{OLP-0137-B010}{start}
\begin{ex}[د سکولم پاراډوکس]
د زرملو--فرېنکل سټ تيوري~\LR{$\Log{ZFC}$} يو ډېر پياوړے چوکاټ دے چې
کابو ټولې رياضيکي خبرې پکښې څرګندېدلے شي، د سټونو د کچو په اړه
حقيقتونه هم۔ نو د بېلګې په توګه،~\LR{$\Log{ZFC}$} ثابتولے شي چې د
حقيقي عددونو سټ~\LR{$\Real$}،~د شمېر نۀ وړ دے؛ د کانټور قضيه
ثابتولے شي چې د هر سټ تواني سټ پخپله له هغه سټ څخه لوی دے، او داسې
نور۔ که~\LR{$\Log{ZFC}$} سازګاره وي، ټول مدلونه ئې نامتناهي دي؛ او
سربېره پر دې، ټول داسې غړي لري چې تيوري ئې په اړه وايي
چې~د شمېر نۀ وړ دي، لکه هغه عنصر چې د~\LR{$\Log{ZFC}$} دا قضيه
رښتونې کوي چې د طبيعي عددونو تواني سټ موجود دے۔ د
لوېنهايم--سکولم قضيې له مخې،~\LR{$\Log{ZFC}$}،~د شمېر وړ مدلونه هم
لري: داسې مدلونه چې «د شمېر نۀ وړ» سټونه لري، خو پخپله
د شمېر وړ دي۔
\end{ex}
\olpseganchor{OLP-0137-B010}{end}

\paragraph{د منجمدې سرچينې بشپړ بديل متن}
\olpseganchor{OLP-0644-B005}{start}
\olfileid{fol}{com}{mcs-olp0644}
\olsection{د تړلي فارمولونه اعظمي سازګار سټونه}
\olpseganchor{OLP-0644-B005}{end}

\olpseganchor{OLP-0644-B006}{start}
\begin{defn}[اعظمي سازګار سټ]
\ollabel{mcs}
د تړلي فارمولونه سټ~\LR{$\Gamma$} \emph{اعظمي سازګار} دے هله او يوازې هله چې
\begin{enumerate}
\item \LR{$\Gamma$} سازګار وي، او
\item که \LR{$\Gamma \subsetneq \Gamma'$} وي، نو \LR{$\Gamma'$} ناسازګار وي۔
\end{enumerate}
\end{defn}
\olpseganchor{OLP-0644-B006}{end}

\olpseganchor{OLP-0644-B007}{start}
\begin{explain}
د پاسني تعريف يو همارز بديل دا دے: د جملو سټ \LR{$\Gamma$} \emph{اعظمي سازګار} دے هله او يوازې هله چې
\begin{enumerate}
\item \LR{$\Gamma$} سازګار وي، او
\item که \LR{$\Gamma \cup \{ !A \}$} سازګار وي، نو \LR{$!A \in \Gamma$}۔
\end{enumerate}
په بل عبارت، اعظمي سازګار سټ~\LR{$\Gamma$} ته هغه تړلي فارمولونه چې له وړاندې په~\LR{$\Gamma$} کښې نۀ وي، له دې پرته نۀ شو زياتولے چې پايله کښې ترې جوړ شوے لوے سټ ناسازګار شي۔
\end{explain}
\olpseganchor{OLP-0644-B007}{end}

\olpseganchor{OLP-0644-B008}{start}
\begin{explain}
اعظمي سازګار سټونه د بشپړتيا په ثبوت کښې ځکه مهم دي چې تضمينولے شو د تړلي فارمولونه هر سازګار سټ~\LR{$\Gamma$} په يو اعظمي سازګار سټ~\LR{$\Gamma^*$} کښې پروت وي، او اعظمي سازګار سټ د هرې تړلے فارمول~\LR{$!A$} دپاره يا \LR{$!A$} يا د هغې نفي \LR{$\lnot !A$} لري۔ دا په ځانګړي ډول د اتومي تړلي فارمولونه دپاره هم رښتيا دي؛ نو په داسې ژبه کښې چې په مناسبو ثابت سمبولونه غځول شوې وي، له اعظمي سازګار سټ څخه جوړښت جوړولے شو، چې د پريديکاتونه تفسير ئې د دې له مخې تعريفېږي چې کومې اتومي تړلي فارمولونه په~\LR{$\Gamma^*$} کښې دي۔ بيا ښودلے شو چې دا جوړښت د~\LR{$\Gamma^*$} ټولې تړلي فارمولونه (او ځکه د \LR{$\Gamma$} هم) رښتيا کوي۔ د دې وروستي حقيقت ثبوت غواړي چې \LR{$\lnot !A \in \Gamma^*$} هله او يوازې هله چې \LR{$!A \notin \Gamma^*$}، \LR{$(!A \lor !B) \in \Gamma^*$} هله او يوازې هله چې \LR{$!A \in \Gamma^*$} يا \LR{$!B \in \Gamma^*$}، او داسې نور۔
\end{explain}
\olpseganchor{OLP-0644-B008}{end}

\olpseganchor{OLP-0644-B009}{start}
\begin{prop}
\ollabel{prop:mcs}
فرض کړئ چې \LR{$\Gamma$} اعظمي سازګار دے۔ نو:
\begin{enumerate}
\item \ollabel{prop:mcs-prov-in} که \LR{$\Gamma \Proves !A$} وي، نو \LR{$!A \in \Gamma$}۔
\olpseganchor{OLP-0644-B009}{end}

\olpseganchor{OLP-0644-B010}{start}
\item \ollabel{prop:mcs-either-or} د هر \LR{$!A$} دپاره يا \LR{$!A \in \Gamma$} يا \LR{$\lnot !A \in \Gamma$}۔
\olpseganchor{OLP-0644-B010}{end}

\olpseganchor{OLP-0644-B011}{start}
\item \ollabel{prop:mcs-and} \LR{$(!A \land !B) \in \Gamma$} هله او يوازې هله چې \LR{$!A \in \Gamma$} او \LR{$!B \in \Gamma$} دواړه وي۔
\olpseganchor{OLP-0644-B011}{end}

\olpseganchor{OLP-0644-B012}{start}
\item \ollabel{prop:mcs-or} \LR{$(!A \lor !B) \in \Gamma$} هله او يوازې هله چې يا \LR{$!A \in \Gamma$} يا \LR{$!B \in \Gamma$}۔
\olpseganchor{OLP-0644-B012}{end}

\olpseganchor{OLP-0644-B013}{start}
\item \ollabel{prop:mcs-if} \LR{$(!A \lif !B) \in \Gamma$} هله او يوازې هله چې يا \LR{$!A \notin \Gamma$} يا \LR{$!B \in \Gamma$}۔
\end{enumerate}
\end{prop}
\olpseganchor{OLP-0644-B013}{end}

\olpseganchor{OLP-0644-B014}{start}
\begin{proof}
په ټولو لاندې برخو کښې فرض کوو چې \LR{$\Gamma$} اعظمي سازګار دے۔
\begin{enumerate}
\item که \LR{$\Gamma \Proves !A$} وي، نو \LR{$!A \in \Gamma$}۔
\olpseganchor{OLP-0644-B014}{end}

\olpseganchor{OLP-0644-B015}{start}
فرض کړئ چې \LR{$\Gamma \Proves !A$}۔ برعکس فرض کړئ چې \LR{$!A \notin \Gamma$}: نو ځکه چې \LR{$\Gamma$} اعظمي سازګار دے، \LR{$\Gamma \cup \{!A\}$} ناسازګار دے، ځکه نو \LR{$\Gamma \cup \{!A\} \Proves \lfalse$}۔ د
\cref{fol:seq:prv:prop:provability-contr,fol:ntd:prv:prop:provability-contr}
له مخې \LR{$\Gamma$} ناسازګار دے۔ دا د دې فرض خلاف دے چې \LR{$\Gamma$} سازګار دے۔ ځکه نو داسې نۀ شي کېدے چې \LR{$!A \notin \Gamma$}؛ نو \LR{$!A \in \Gamma$}۔
\olpseganchor{OLP-0644-B015}{end}

\olpseganchor{OLP-0644-B016}{start}
\item د هر \LR{$!A$} دپاره يا \LR{$!A \in \Gamma$} يا \LR{$\lnot !A \in \Gamma$}۔
\olpseganchor{OLP-0644-B016}{end}

\olpseganchor{OLP-0644-B017}{start}
برعکس فرض کړئ چې د کوم~\LR{$!A$} دپاره \LR{$!A \notin \Gamma$} او \LR{$\lnot !A \notin \Gamma$} دواړه وي۔ ځکه چې \LR{$\Gamma$} اعظمي سازګار دے، \LR{$\Gamma \cup \{!A\}$} او \LR{$\Gamma \cup \{\lnot !A\}$} دواړه ناسازګار دي؛ نو \LR{$\Gamma \cup \{!A\} \Proves \lfalse$} او \LR{$\Gamma \cup \{\lnot !A\} \Proves \lfalse$}۔ د
\cref{fol:seq:prv:prop:provability-exhaustive,fol:ntd:prv:prop:provability-exhaustive}
له مخې \LR{$\Gamma$} ناسازګار دے، چې تناقض دے۔ نو داسې تړلے فارمول~\LR{$!A$} نۀ شي کېدے، او د هر \LR{$!A$} دپاره \LR{$!A \in \Gamma$} يا \LR{$\lnot !A \in \Gamma$}۔
\olpseganchor{OLP-0644-B017}{end}

\olpseganchor{OLP-0644-B018}{start}
\item %
%
\LR{$(!A \land !B) \in \Gamma$} هله او يوازې هله چې \LR{$!A \in \Gamma$} او \LR{$!B \in \Gamma$} دواړه وي:
\olpseganchor{OLP-0644-B018}{end}

\olpseganchor{OLP-0644-B019}{start}
د مخامخ جهت دپاره فرض کړئ چې \LR{$(!A \land !B) \in \Gamma$}۔ نو \LR{$\Gamma \Proves !A \land !B$}۔ د
\cref{fol:seq:ppr:prop:provability-land,fol:ntd:ppr:prop:provability-land} (لومړۍ فقره)
له مخې \LR{$\Gamma \Proves !A$} او \LR{$\Gamma \Proves !B$}۔ د
\olref{prop:mcs-prov-in} له مخې \LR{$!A \in \Gamma$} او \LR{$!B \in \Gamma$}، لکه چې غوښتل شوي وو۔
\olpseganchor{OLP-0644-B019}{end}

\olpseganchor{OLP-0644-B020}{start}
د معکوس جهت دپاره \LR{$!A \in \Gamma$} او \LR{$!B \in \Gamma$} ونيسئ۔ نو \LR{$\Gamma \Proves !A$} او \LR{$\Gamma \Proves !B$}۔ د
\cref{fol:seq:ppr:prop:provability-land,fol:ntd:ppr:prop:provability-land} (دويمه فقره)
له مخې \LR{$\Gamma \Proves !A \land !B$}۔ د \olref{prop:mcs-prov-in} له مخې \LR{$(!A \land !B) \in \Gamma$}۔
\olpseganchor{OLP-0644-B020}{end}

\olpseganchor{OLP-0644-B021}{start}
\item %
%
\LR{$(!A \lor !B) \in \Gamma$} هله او يوازې هله چې يا \LR{$!A \in \Gamma$} يا \LR{$!B \in \Gamma$}۔
\olpseganchor{OLP-0644-B021}{end}

\olpseganchor{OLP-0644-B022}{start}
د مخامخ جهت د مقابل نقيض دپاره فرض کړئ چې \LR{$!A \notin \Gamma$} او \LR{$!B \notin \Gamma$}۔ غواړو وښيو چې \LR{$(!A \lor !B) \notin \Gamma$}۔ ځکه چې \LR{$\Gamma$} اعظمي سازګار دے، \LR{$\Gamma \cup \{!A\} \Proves \lfalse$} او \LR{$\Gamma \cup \{!B\} \Proves \lfalse$}۔ د
\cref{fol:seq:ppr:prop:provability-lor,fol:ntd:ppr:prop:provability-lor} (لومړۍ فقره)
له مخې \LR{$\Gamma \cup \{(!A \lor !B)\}$} ناسازګار دے۔ نو \LR{$(!A \lor !B) \notin \Gamma$}، لکه چې غوښتل شوي وو۔
\olpseganchor{OLP-0644-B022}{end}

\olpseganchor{OLP-0644-B023}{start}
د معکوس جهت دپاره فرض کړئ چې \LR{$!A \in \Gamma$} يا \LR{$!B \in \Gamma$}۔ نو \LR{$\Gamma \Proves !A$} يا \LR{$\Gamma \Proves !B$}۔ د
\cref{fol:seq:ppr:prop:provability-lor,fol:ntd:ppr:prop:provability-lor} (دويمه فقره)
له مخې \LR{$\Gamma \Proves !A \lor !B$}۔ د \olref{prop:mcs-prov-in} له مخې \LR{$(!A \lor !B) \in \Gamma$}، لکه چې غوښتل شوي وو۔
\olpseganchor{OLP-0644-B023}{end}

\olpseganchor{OLP-0644-B024}{start}
\item %
%
\LR{$(!A \lif !B) \in \Gamma$} هله او يوازې هله چې يا \LR{$!A \notin \Gamma$} يا \LR{$!B \in \Gamma$}:
\olpseganchor{OLP-0644-B024}{end}

\olpseganchor{OLP-0644-B025}{start}
د مخامخ جهت دپاره \LR{$(!A \lif !B) \in \Gamma$} ونيسئ، او برعکس فرض کړئ چې \LR{$!A \in \Gamma$} او \LR{$!B \notin \Gamma$}۔ په دغو فرضونو \LR{$\Gamma \Proves !A \lif !B$} او \LR{$\Gamma \Proves !A$}۔ د
\cref{fol:seq:ppr:prop:provability-lif,fol:ntd:ppr:prop:provability-lif} (لومړۍ فقره)
له مخې \LR{$\Gamma \Proves !B$}۔ خو بيا د \olref{prop:mcs-prov-in} له مخې \LR{$!B \in \Gamma$}، چې د \LR{$!B \notin \Gamma$} فرض سره په تناقض کښې دے۔
\olpseganchor{OLP-0644-B025}{end}

\olpseganchor{OLP-0644-B026}{start}
د معکوس جهت دپاره، اول هغه حالت وګورئ چې \LR{$!A \notin \Gamma$}۔ د \olref{prop:mcs-either-or} له مخې \LR{$\lnot !A \in \Gamma$}، نو \LR{$\Gamma \Proves \lnot !A$}۔ د
\cref{fol:seq:ppr:prop:provability-lif,fol:ntd:ppr:prop:provability-lif} (دويمه فقره)
له مخې \LR{$\Gamma \Proves !A \lif !B$}۔ بيا د \olref{prop:mcs-prov-in} له مخې راځي چې \LR{$(!A \lif !B) \in \Gamma$}، لکه چې غوښتل شوي وو۔
\olpseganchor{OLP-0644-B026}{end}

\olpseganchor{OLP-0644-B027}{start}
اوس هغه حالت وګورئ چې \LR{$!B \in \Gamma$}۔ نو \LR{$\Gamma \Proves !B$}، او د
\cref{fol:seq:ppr:prop:provability-lif,fol:ntd:ppr:prop:provability-lif} (دويمه فقره)
له مخې \LR{$\Gamma \Proves !A \lif !B$}۔ د \olref{prop:mcs-prov-in} له مخې \LR{$(!A \lif !B) \in \Gamma$}۔
\end{enumerate}
\end{proof}
\olpseganchor{OLP-0644-B027}{end}

\olpseganchor{OLP-0644-B028}{start}
\begin{prob}
د \olref[fol][com][mcs-olp0644]{prop:mcs} ثبوت بشپړ کړئ۔
\end{prob}
\olpseganchor{OLP-0644-B028}{end}

\olpseganchor{OLP-0138-B004}{start}
\olpart{fol}{د لومړۍ درجې منطق}
\olpseganchor{OLP-0138-B004}{end}

\olpseganchor{OLP-0138-B005}{start}
\begin{editorial}
  په دې برخه کښې د لومړۍ درجې منطق مابعدمنطق تر بشپړتيا پورې
  څېړل کېږي۔ اوس مهال دا برخه د بياني منطق پر جلا بيان تکيه نۀ
  کوي؛ هر څۀ پکښې ثابتېږي۔ که د \texttt{FOL} ټېګ ناسم وي، د سرچينې
  فايلونه به د کميت ټاکونکو اړوند مواد وباسي (او
  ``جوړښت'' به په ``ارزښت ټاکنه''، \LR{$\Struct{M}$} به په
  \LR{$\pAssign{v}$}، او داسې نورو بدل کړي)۔ په حقيقت کښې، د بياني منطق
  د برخې زياتره مواد همدا د لومړۍ درجې منطق مواد دي چې د \texttt{FOL}
  ټېګ پکښې بند کړے شوے دے۔
\olpseganchor{OLP-0138-B005}{end}

\olpseganchor{OLP-0138-B006}{start}
  که د بياني منطق برخه هم شامله وي، نو ډېر تکرار ترې جوړېږي۔ خو
  پلان دا دے چې دا برخه د بياني منطق مواد په پام کښې نيولے شي (او
  مخکښې څېړل شوي مواد وباسي، او همدارنګه ثبوتونه د بياني برخې
  اړوندو ځايونو ته په حوالو لنډ کړي)۔
\end{editorial}
\olpseganchor{OLP-0138-B006}{end}











\olpseganchor{OLP-0138-B012}{start}
%
  

\olpseganchor{OLP-0138-B012}{end}

\olpseganchor{OLP-0138-B013}{start}
%
  

\olpseganchor{OLP-0138-B013}{end}

\olpseganchor{OLP-0138-B014}{start}
%
  

\olpseganchor{OLP-0138-B014}{end}

\olpseganchor{OLP-0138-B015}{start}
%
  

\olpseganchor{OLP-0138-B015}{end}

\olpseganchor{OLP-0139-B004}{start}
\olchapter{fol}{int}{د لومړۍ درجې منطق پېژندنه}
\olpseganchor{OLP-0139-B004}{end}

\olpseganchor{OLP-0140-B004}{start}
\olfileid{fol}{int}{fol}
\olpseganchor{OLP-0140-B004}{end}

\olpseganchor{OLP-0140-B005}{start}
\olsection{د لومړۍ درجې منطق}
\olpseganchor{OLP-0140-B005}{end}

\olpseganchor{OLP-0140-B006}{start}
تاسو غالباً د صوري منطق په لومړۍ پېژندنه کښې د لومړۍ درجې له منطق
سره بلد يئ۔\footnote{په حقيقت کښې موږ لږ و ډېر دا فرضوو چې تاسو
ورسره بلد يئ!{} که نۀ يئ، نو يو لا ابتدايي درسي کتاب، لکه
\emph{\texttt{forall x}} \mbox{\LR{\latinfont \citep{Magnus2021}}}، بيا کتلے شئ۔} ښايي تاسو دا د
``کميت ټاکونکو منطق'' يا ``پريديکاتي منطق'' په نامه پېژنئ۔ د لومړۍ
درجې منطق تر هر څه مخکښې يوه صوري ژبه ده۔ يعنې يو ټاکلے لغتځونډ
لري، او عبارتونه ئې د همدې لغتځونډ تارونه دي۔ خو هر تار جايز نۀ
دے۔ د جايزو عبارتونو بېلابېل ډولونه شته: ترمونه، فارمولونه او
تړلي فارمولونه۔ موږ په اصل کښې د لومړۍ درجې منطق له تړلي فارمولونه
سره لېوال يو: هغوئ د انګرېزي ژبې د جملو صوري سيال راکوي، او د هغوئ
په باب هغه پوښتنې کولے شو چې منطقپوه عموماً ورسره لېوال وي۔ مثلاً:
\begin{itemize}
    \item ايا \LR{$!B$} له~\LR{$!A$} څخه په منطقي ډول لازمېږي؟
    \item ايا \LR{$!A$} په منطقي ډول رښتونې، په منطقي ډول دروغجنه، که
اتفاقي ده؟
    \item ايا \LR{$!A$} او \LR{$!B$} هم‌ارزښته دي؟
\end{itemize}
\olpseganchor{OLP-0140-B006}{end}

\olpseganchor{OLP-0140-B007}{start}
دا پوښتنې په بنسټيز ډول د لومړۍ درجې منطق د تړلي فارمولونه د
``معنٰی'' په باب دي۔ مثلاً يو فيلسوف به دا پوښتنه، چې ايا \LR{$!B$} له
~\LR{$!A$} څخه په منطقي ډول لازمېږي، داسې وڅېړي: ايا کوم داسې حالت شته
چې \LR{$!A$} پکښې رښتونې خو~\LR{$!B$} دروغجنه وي (\LR{$!B$} له~\LR{$!A$} څخه نۀ
لازمېږي)، که هر هغه حالت چې \LR{$!A$} رښتونې کوي~\LR{$!B$} هم رښتونې کوي
(\LR{$!B$} له~\LR{$!A$} څخه لازمېږي)؟ خو لا راته نۀ دي ويل شوي چې ``حالت''
څۀ دے---دا د \emph{معنٰی پوهنې} کار دے۔ د لومړۍ درجې منطق معنٰی
پوهنه د فيلسوف د ``حالت'' د وجداني مفکورې، او هم---چې دا مهمه
ده---په يوه حالت کښې د تړلے فارمول~\LR{$!A$} د \emph{رښتونې کېدو}
رياضيکي دقيق مدل راکوي۔ هغه رياضيکي دقيق مدل چې جوړوو ئې
جوړښت بولو۔ هغه اړيکه چې ``پکښې رښتونې'' په دقيق ډول
بيانوي، د \emph{صدق} اړيکه بلل کېږي۔ نو موږ به د
تړلي فارمولونه~\LR{$!A$} او جوړښتونه~\LR{$\Struct{M}$} دپاره دا تعريفوو
چې ``\LR{$!A$} په~\LR{$\Struct{M}$} کښې صدق کوي'' (په سمبولونو کښې:
\LR{$\Sat{M}{!A}$})۔ له دې وروسته نور معنايي اصطلاحات، لکه ``ترې
لازمېږي'' يا ``په منطقي ډول رښتونې ده''، هم په دقيق ډول تعريفولے
شو۔ دا تعريفونه به ممکنه کړي چې په رياضيکي دقت سره، مثلاً، فيصله
وکړو چې ايا
\LR{$\lforall[x][(!A(x) \lif !B(x))], \lexists[x][!A(x)] \Entails \lexists[x][!B(x)]$}
ده که نۀ۔ ځواب به طبعاً ``هو'' وي۔
د لومړۍ مقدمې عمومي کميت ټاکونکے يوازې پر شرطي فارمول حاکم دے۔
که تاسو مخکښې د
انګرېزي جملو د لومړۍ درجې منطق په سمبولونو اړولو زده کړه کړې وي، دا
به مثلاً د ``ټولې مېږي حشرات دي، مېږي شته، نو حشرات شته'' د
سمبولونو بڼه وپېژنئ۔ دا په څرګند ډول معتبر استدلال دے؛ نو زموږ د
صوري ژبې د ``ترې لازمېږي'' رياضيکي مدل هم بايد همدا ځواب ورکړي۔
\olpseganchor{OLP-0140-B007}{end}

\olpseganchor{OLP-0140-B008}{start}
بله موضوع چې غالباً د صوري منطق له لومړۍ پېژندنې څخه مو ياده ده،
\emph{اشتقاقونه} دي۔ که مو د صوري منطق لومړے کورس لوستلے وي،
ښوونکي به درباندې د داسې اشتقاقونه د موندلو تمرين کړے وي؛
ښايي حتٰی داسې اشتقاق مو موندلے وي چې پورته استلزام
ښيي۔ د اشتقاقونه د ورکولو ډېرې بېلابېلې لارې شته: ښايي تاسو
``فطري استنتاج'' يا ``د رښتياوالي ونې'' کارولې وي، خو نورې ډېرې
لارې هم شته۔ د اشتقاق د نظامونو موخه داسې وسيلې برابرول دي
چې د منطقپوهانو پورته پوښتنې پرې ځواب شي: مثلاً د فطري استنتاج داسې
اشتقاق چې \LR{$\lforall[x][(!A(x) \lif !B(x))]$} او
\LR{$\lexists[x][!A(x)]$} ئې مقدمې وي او \LR{$\lexists[x][!B(x)]$} ئې پايله
(وروستۍ کرښه) وي، دا \emph{تصديقوي} چې \LR{$\lexists[x][!B(x)]$} په
منطقي ډول له \LR{$\lforall[x][(!A(x) \lif !B(x))]$} او
\LR{$\lexists[x][!A(x)]$} څخه لازمېږي۔
\olpseganchor{OLP-0140-B008}{end}

\olpseganchor{OLP-0140-B009}{start}
خو دا ولې؟ په ظاهره د اشتقاق نظامونه له معنٰی پوهنې سره
هېڅ تړاو نۀ لري: د يوه صوري اشتقاق ورکول يوازې دا غواړي چې
سمبولونه د ځينو قاعدو له مخې سره منظم شي؛ هلته د ``حالتونو'' يا
``پکښې رښتونې'' هېڅ يادونه نشته۔ د اشتقاق د نظامونو او
معنٰی پوهنې تر منځ تړاو بايد په مابعدمنطقي څېړنه ثابت شي۔ پکار دا
ده چې مثلاً په رياضيکي ډول ثابته کړو: له مقدمو
\LR{$\lforall[x][(!A(x) \lif !B(x))]$} او \LR{$\lexists[x][!A(x)]$} څخه د
\LR{$\lexists[x][!B(x)]$} صوري اشتقاق هله او يوازې هله ممکن دے
چې \LR{$\lforall[x][(!A(x) \lif !B(x))]$} او
\LR{$\lexists[x][!A(x)]$} په ګډه \LR{$\lexists[x][!B(x)]$} لازم کړي۔ خو تر دې مخکښې چې دا کار وشي، د نحو او معنٰی پوهنې د تعريفونو
د سمولو دپاره ډېر زيارکښ کار کول پکار دي۔
\olpseganchor{OLP-0140-B009}{end}

\olpseganchor{OLP-0141-B004}{start}
\olfileid{fol}{int}{syn}
\olpseganchor{OLP-0141-B004}{end}

\olpseganchor{OLP-0141-B005}{start}
\olsection{نحو}
\olpseganchor{OLP-0141-B005}{end}

\olpseganchor{OLP-0141-B006}{start}
لومړے بايد دا په دقيق ډول وټاکو چې د سمبولونو کوم تارونه د لومړۍ
درجې منطق تړلي فارمولونه ګڼل کېږي۔ دا کار به وروسته وکړو؛ اوس به
يوازې د بېلګو له مخې مخکښې ځو۔ د لومړۍ درجې منطق بنسټيزې جوړښتي
خښتې---يعنې لغتځونډ---په دوو برخو وېشل کېږي۔ لومړۍ برخه هغه
سمبولونه دي چې په وسيله ئې ځانګړي څيزونه ټاکو يا د هغوئ په باب
ځانګړې خبرې کوو۔ څيزونه په ثابت سمبولونه ټاکو، او بيا د ټاکلو شوو
څيزونو په باب په پريديکاتونه څه وايو۔ مثلاً ښايي~\LR{$\Obj a$} د
ثابت سمبول په توګه د يوه څيز د ټاکلو دپاره وکاروو، او بيا د
تړلے فارمول~\LR{$\Atom{\Obj P}{\Obj a}$} په وسيله د هغه په باب څه
ووايو۔ که د ``\LR{$\Obj a$}'' او ``\LR{$\Obj P$}'' معناوې په ذهن کښې ولرئ،
\LR{$\Atom{\Obj P}{\Obj a}$} د انګرېزي ژبې د يوې جملې په توګه لوستلے
شئ (او غالباً مو د صوري منطق د لومړۍ زده کړې پر وخت همداسې کړي
دي)۔ کله چې د لومړۍ درجې منطق داسې ساده تړلي فارمولونه ولرو، د
لغتځونډ په دويمه برخه، يعنې منطقي سمبولونو (نښلوونکو او کميت
ټاکونکو)، ترې نورې پېچلې جملې جوړولے شو۔ نو مثلاً داسې عبارتونه
جوړولے شو لکه \LR{$(\Atom{\Obj P}{\Obj a} \land \Atom{\Obj Q}{\Obj b})$}
يا~\LR{$\lexists[\Obj x][\Atom{\Obj P}{\Obj x}]$}۔
\olpseganchor{OLP-0141-B006}{end}

\olpseganchor{OLP-0141-B007}{start}
د معنٰی پوهنې د دقيقو تعريفونو او د هغو د اشتقاق د نظامونو
د قاعدو د ورکولو دپاره چې سخته مابعدمنطقي څېړنه ئې غواړي، تر هر څه
مخکښې بايد په دقيق ډول تعريف کړو چې د لومړۍ درجې منطق
تړلے فارمول څه ګڼل کېږي۔ بنسټيزه مفکوره دومره اسانه ده: ځينې ساده
تړلي فارمولونه يوازې له پريديکاتونه او ثابت سمبولونه څخه جوړولے
شو، لکه~\LR{$\Atom{\Obj P}{\Obj a}$}۔ بيا له دغو څخه د نښلوونکو او کميت
ټاکونکو په وسيله پېچلې جملې جوړوو۔ خو هغه قواعد په دقيق ډول کوم دي
چې له مخې ئې د پېچلو تړلي فارمولونه جوړول راته جايز دي؟ دا بايد
وټاکل شي، کنه ``د لومړۍ درجې منطق تړلے فارمول'' مو په کافي دقت نۀ
دے تعريف کړے۔ څو مسئلې شته۔ لومړۍ دا چې سم تارونه د تړلي فارمولونه
په توګه وشمېرو۔ دويمه دا چې کار په داسې ډول وکړو چې د
\emph{ټولو} تړلي فارمولونه په باب رياضيکي ثبوتونه ورکولے شو۔ په پاے
کښې بايد پر تړلي فارمولونه د ځينو لومړنيو عملياتو دقيق تعريفونه هم
ورکړو، لکه ``په~\LR{$!A$} کښې هر~\LR{$\Obj x$} په~\LR{$\Obj b$} بدل کړه''۔ ستونزه
دا ده چې د لومړۍ درجې منطق کميت ټاکونکي او متغيرونه دا کار
بشپړ څرګند نۀ پرېږدي۔ مثلاً ايا
\LR{$\lexists[\Obj x][\Atom{\Obj P}{\Obj a}]$} بايد تړلے فارمول وګڼل
شي؟ او \LR{$\lexists[\Obj x][\lexists[\Obj x][\Atom{\Obj P}{\Obj x}]]$}
څنګه؟ ``په \LR{$(\Atom{\Obj P}{\Obj x} \land \lexists[\Obj x][\Atom{\Obj P}{\Obj x}])$}
کښې \LR{$\Obj x$} په~\LR{$\Obj b$} بدل کړه'' پايله بايد څه وي؟
\olpseganchor{OLP-0141-B007}{end}

\olpseganchor{OLP-0142-B004}{start}
\olfileid{fol}{int}{fml}
\olpseganchor{OLP-0142-B004}{end}

\olpseganchor{OLP-0142-B005}{start}
\section{فارمولونه}
\olpseganchor{OLP-0142-B005}{end}

\olpseganchor{OLP-0142-B006}{start}
د لومړۍ درجې منطق د تړلي فارمولونه د سخت تعريف او د متغيرونه
له کارولو څخه راولاړو شوو مسئلو د هوارولو دپاره به لاندې لاره
کاروو۔ لومړے د عبارتونو يو \emph{بل} سټ تعريفوو: فارمولونه۔ له
دې وروسته پکښې د متغيرونه ونډه څېړلے شو---او د ځينو نورو
مفکورو، يعنې ``ازادو'' او ``تړلو'' متغيرونه، پر بنسټ تعريفولے
شو چې تړلے فارمول څه ده (يعنې داسې فارمول چې ازاد
متغيرونه نۀ لري)۔ دا کار يوازې د ``تړلے فارمول'' د تعريف د
اسانه کولو دپاره نۀ کوو؛ بلکې د صدق د معنايي مفهوم د تعريف په
طريقه کښې به هم ډېر اړين وي۔
\olpseganchor{OLP-0142-B006}{end}

\olpseganchor{OLP-0142-B007}{start}
راځئ ``فارمول'' د لومړۍ درجې د يوې ساده ژبې دپاره تعريف کړو:
داسې ژبه چې يوازې يو پريديکات~\LR{$\Obj P$}، يو
ثابت سمبول~\LR{$\Obj a$}، او يوازې منطقي سمبولونه \LR{$\lnot$}، \LR{$\land$}
او~\LR{$\lexists$} لري۔ زموږ بشپړ تعريفونه به ډېر عام وي: نامتناهي ډېر
پريديکاتونه او ثابت سمبولونه به جايز کړو۔ په حقيقت کښې هغه
تابعې به هم څېړو چې له ثابت سمبولونه او متغيرونه سره
يوځاے ``ترمونه'' جوړوي۔ اوس دپاره~\LR{$\Obj a$} او متغيرونه زموږ يوازيني
ترمونه دي۔ البته نامتناهي ډېرو متغيرونه ته اړتيا لرو۔ په رسمي
ډول به \LR{$\Obj v_0$}، \LR{$\Obj v_1$}، \dots، د متغيرونو په توګه کاروو۔
\olpseganchor{OLP-0142-B007}{end}

\olpseganchor{OLP-0142-B008}{start}
\begin{defn}
د \emph{فارمولونه} سټ~\LR{$\Frm$} داسې تعريفېږي:
\begin{enumerate}
\item\ollabel{fmls-atom} \LR{$\Atom{\Obj P}{\Obj a}$} او \LR{$\Atom{\Obj
  P}{\Obj v_i}$} فارمولونه دي (\LR{$i \in \Nat$})۔
\olpseganchor{OLP-0142-B008}{end}

\olpseganchor{OLP-0142-B009}{start}
\item \ollabel{fmls-not}که \LR{$!A$} فارمول وي، نو \LR{$\lnot !A$}
  فارمول ده۔
\olpseganchor{OLP-0142-B009}{end}

\olpseganchor{OLP-0142-B010}{start}
\item که \LR{$!A$} او \LR{$!B$} فارمولونه وي، نو \LR{$(!A \land
  !B)$} فارمول ده۔
\olpseganchor{OLP-0142-B010}{end}

\olpseganchor{OLP-0142-B011}{start}
\item \ollabel{fmls-ex}که \LR{$!A$} فارمول او \LR{$x$}
  متغير وي، نو \LR{$\lexists[x][!A]$} فارمول ده۔
\olpseganchor{OLP-0142-B011}{end}

\olpseganchor{OLP-0142-B012}{start}
\item \ollabel{fmls-limit}بل هېڅ څيز فارمول نۀ دے۔
\end{enumerate}
\end{defn}
\olpseganchor{OLP-0142-B012}{end}

\olpseganchor{OLP-0142-B013}{start}
\olref{fmls-atom} راته وايي چې د هر \LR{$i \in \Nat$} دپاره
\LR{$\Atom{\Obj P}{\Obj a}$} او \LR{$\Atom{\Obj P}{\Obj v_i}$}
فارمولونه دي۔ دا تش په نامه \emph{اټومي} فارمولونه دي۔ هغوئ
د پيل دپاره څه راکوي۔ نور شرطونه له هغو څخه چې لا مو
جوړ کړي، د نويو فارمولونه د جوړولو لارې راکوي۔ نو مثلاً د
\olref{fmls-not} له مخې \LR{$\lnot \Atom{\Obj P}{\Obj v_2}$}
فارمول ده، ځکه \LR{$\Atom{\Obj P}{\Obj v_2}$} د
\olref{fmls-atom} له مخې لا فارمول ده۔ بيا د
\olref{fmls-ex} له مخې \LR{$\lexists[\Obj v_2][\lnot \Atom{\Obj P}{\Obj v_2}]$}
بله فارمول ده، او همداسې مخکښې۔ \olref{fmls-limit} راته وايي
چې \emph{يوازې} هغه تارونه فارمولونه ګڼل کېږي چې په همدې طريقه
ئې جوړولے شو۔ په ځانګړي ډول
\LR{$\lexists[\Obj v_0][\Atom{\Obj P}{\Obj a}]$} او
\LR{$\lexists[\Obj v_0][\lexists[\Obj v_0][\Atom{\Obj P}{\Obj a}]]$}
په رښتيا فارمولونه ګڼل کېږي، خو \LR{$(\lnot \Atom{\Obj P}{\Obj a})$}
نۀ ګڼل کېږي، ځکه دباندې قوسونه ئې زيات دي۔
\olpseganchor{OLP-0142-B013}{end}

\olpseganchor{OLP-0142-B014}{start}
د فارمولونه د تعريف دا طريقه \emph{استقرايي تعريف} بلل کېږي،
او اجازه راکوي چې د استقرايي ثبوت په هغه بڼه د فارمولونه په باب
څه ثابت کړو چې \emph{جوړښتي استقرا} بلل کېږي۔ دا دواړه په عمومي
ډول په \olref[mth][ind][idf]{sec} او \olref[mth][ind][sti]{sec}
کښې څېړل شوي دي؛ وروسته ثبوتونو ته له ننوتلو مخکښې ئې بيا وګورئ۔
بنسټيزه مفکوره دا ده: که غواړئ ثابته کړئ چې يو څۀ د ټولو
فارمولونه دپاره رښتيا دي، لومړے ښيئ چې د اټومي فارمولونه
دپاره رښتيا دي؛ بيا ښيئ چې \emph{که} د هر فارمول~\LR{$!A$}
(او~\LR{$!B$}) دپاره رښتيا وي، د \LR{$\lnot !A$}، \LR{$(!A \land !B)$} او
\LR{$\lexists[x][!A]$} دپاره \emph{هم} رښتيا دي۔ مثلاً په دې سره ثابتېږي
چې دا د \LR{$\lexists[\Obj v_2][\lnot \Atom{\Obj P}{\Obj v_2}]$} دپاره
رښتيا دي: له لومړۍ برخې پوهېږئ چې د اټومي
فارمول~\LR{$\Atom{\Obj P}{\Obj v_2}$} دپاره رښتيا دي۔ بيا له دويمې
برخې ترلاسه کوئ چې د \LR{$\lnot \Atom{\Obj P}{\Obj v_2}$} دپاره رښتيا
دي، او ورپسې بيا د \LR{$\lexists[\Obj v_2][\lnot \Atom{\Obj P}{\Obj v_2}]$}
په خپله دپاره رښتيا دي۔ ځکه چې ټول فارمولونه په استقرايي ډول له
اټومي فارمولونه څخه جوړېږي، دا طريقه د هر فارمول دپاره کار کوي۔
\olpseganchor{OLP-0142-B014}{end}

\olpseganchor{OLP-0143-B004}{start}
\olfileid{fol}{int}{sat}
\olpseganchor{OLP-0143-B004}{end}

\olpseganchor{OLP-0143-B005}{start}
\olsection{صدق}
\olpseganchor{OLP-0143-B005}{end}

\olpseganchor{OLP-0143-B006}{start}
کله چې پوه شو فارمولونه څه دي، سمدستي د لومړۍ درجې منطق معنٰی
پوهنې ته وړاندې تللے شو: دلته بنسټيز تعريف د جوړښت دے۔
زموږ د ساده ژبې دپاره جوړښت~\LR{$\Struct M$} يوازې درې برخې
لري: يو ناتش سټ~\LR{$\Domain{M}$} چې \emph{د تعريف ساحه} بلل کېږي، هغه
څيز چې~\LR{$\Obj a$} ئې په~\LR{$\Struct M$} کښې ټاکي، او هغه څيزونه چې
\LR{$\Obj P$} ورباندې په~\LR{$\Struct M$} کښې رښتونے دے۔ د~\LR{$\Obj a$} ټاکلے
څيز په~\LR{$\Assign{\Obj a}{M}$}، او د هغو څيزونو سټ چې~\LR{$\Obj P$}
ورباندې رښتونے دے په~\LR{$\Assign{\Obj P}{M}$} ښودل کېږي۔
جوړښت~\LR{$\Struct{M}$} همدا درې څيزونه لري: \LR{$\Domain{M}$}،
\LR{$\Assign{\Obj a}{M} \in \Domain{M}$} او
\LR{$\Assign{\Obj P}{M} \subseteq \Domain{M}$}۔ عمومي حالت به لا پېچلے
وي، ځکه ډېر پريديکاتونه او ثابت سمبولونه به وي،
پريديکاتونه له يوه څخه زيات ځايونه لرلے شي، او تابعې به
هم وي۔ لکه د دې څپرکي دوه‌ځاييزه اړيکه، څوځاييز پريديکات
د څيزونو پر مرتب لړ صدق کوي؛ ثابت سمبولونه هر يو يو څيز ټاکي۔
\olpseganchor{OLP-0143-B006}{end}

\olpseganchor{OLP-0143-B007}{start}
دا د هغو فارمولونه د صدق د تعريف دپاره بس ده چې متغيرونه
نۀ لري۔ مفکوره دا ده چې داسې استقرايي تعريف ورکړو چې د
فارمولونه زموږ د تعريف طريقه منعکس کړي۔ ټاکو چې يو اټومي فارمول
کله په~\LR{$\Struct{M}$} کښې صدق کوي؛ بيا مثلاً دا ټاکو چې \LR{$\lnot !A$}
کله په~\LR{$\Struct{M}$} کښې صدق کوي، د دې پر بنسټ چې~\LR{$!A$} په~\LR{$\Struct{M}$} کښې صدق
کوي که نۀ۔ مثلاً داسې تعريفولے شو:
\begin{enumerate}
  \item \LR{$\Atom{\Obj P}{\Obj a}$} په~\LR{$\Struct{M}$} کښې هله صدق کوي
  چې \LR{$\Assign{\Obj a}{M} \in \Assign{\Obj P}{M}$} وي۔
  \item \LR{$\lnot !A$} په~\LR{$\Struct{M}$} کښې هله صدق کوي چې \LR{$!A$} په
  ~ \LR{$\Struct{M}$} کښې صدق نۀ کوي۔
  \item \LR{$(!A \land !B)$} په~\LR{$\Struct{M}$} کښې هله صدق کوي چې \LR{$!A$}
  په~\LR{$\Struct{M}$} کښې صدق وکړي، او \LR{$!B$} هم په~\LR{$\Struct{M}$} کښې
  صدق وکړي۔
\end{enumerate}
فرض کړئ \LR{$\Domain{M} = \{0, 1, 2\}$}، \LR{$\Assign{\Obj a}{M} = 1$}، او
\LR{$\Assign{\Obj P}{M} = \{1, 2\}$} دي۔ دا تعريف راته وايي چې
\LR{$\Atom{\Obj P}{\Obj a}$} په~\LR{$\Struct{M}$} کښې صدق کوي (ځکه
\LR{$\Assign{\Obj a}{M} =  1 \in \{1,2\} = \Assign{\Obj P}{M}$})۔
ورپسې راته وايي چې \LR{$\lnot \Atom{\Obj P}{\Obj a}$} په~\LR{$\Struct{M}$}
کښې صدق نۀ کوي؛ نو په خپل وار \LR{$\lnot\lnot \Atom{\Obj P}{\Obj a}$}
صدق کوي او \LR{$(\lnot \Atom{\Obj P}{\Obj a} \land \Atom{\Obj P}{\Obj a})$}
صدق نۀ کوي، او همداسې مخکښې۔
\olpseganchor{OLP-0143-B007}{end}

\olpseganchor{OLP-0143-B008}{start}
ستونزه هغه وخت پيدا کېږي چې د کميت ټاکونکو دپاره تعريف ورکول
غواړو: زړه مو غواړي داسې ووايو چې
``\LR{$\lexists[\Obj v_0][\Atom{\Obj P}{\Obj v_0}]$} هله صدق کوي چې
\LR{$\Atom{\Obj P}{\Obj v_0}$} صدق وکړي''۔ خو جوړښت~\LR{$\Struct{M}$}
راته نۀ وايي چې له متغيرونه سره څه وکړو۔ په اصل کښې غواړو
ووايو چې \LR{$\Atom{\Obj P}{\Obj v_0}$} \emph{د~\LR{$\Obj v_0$} د کوم قيمت
دپاره} صدق کوي۔ د دې د دقيق کولو دپاره داسې لارې ته اړتيا لرو چې
د~\LR{$\Domain{M}$} غړي يوازې~\LR{$\Obj a$} ته نۀ، بلکې~\LR{$\Obj v_0$}
ته هم وګومارو۔ د همدې دپاره د متغير \emph{ګومارنې} ورپېژنو۔
د متغير ګومارنه يوازې داسې تابع~\LR{$s$} ده چې متغيرونه
د~\LR{$\Domain{M}$} غړي ته وړي (زموږ په بېلګه کښې له~\LR{$0$}، \LR{$1$}
يا~\LR{$2$} څخه يوه ته)۔ ځکه مخکښې نۀ
پوهېږو چې په فارمول کښې به کوم متغيرونه راشي، دا نۀ شو
محدودولے چې~\LR{$s$} کومو متغيرونه ته قيمتونه وټاکي۔ ساده حل دا
دے چې وغواړو~\LR{$s$} \emph{ټولو} متغيرونه~\LR{$\Obj v_0$}،
\LR{$\Obj v_1$}، \dots\@ ته قيمتونه وګوماري۔ موږ به يوازې د اړينو هغو
قيمتونه کاروو۔
\olpseganchor{OLP-0143-B008}{end}

\olpseganchor{OLP-0143-B009}{start}
د دې پر ځاے چې د فارمولونه صدق يوازې د جوړښت په نسبت
تعريف کړو، دا به د جوړښت~\LR{$\Struct{M}$} \emph{او} د
متغير ګومارنې~\LR{$s$} دواړو په نسبت تعريف کړو، او په لنډ ډول به
\LR{$\Sat{M}{!A}[s]$} ليکو۔ اوس به زموږ تعريف د هغو اټومي فارمولونه
دپاره يو اضافي شرط هم ولري چې متغيرونه لري:
\begin{enumerate}
  \item \LR{$\Sat{M}{\Atom{\Obj P}{\Obj a}}[s]$} هله او يوازې هله چې
  \LR{$\Assign{\Obj a}{M} \in \Assign{\Obj P}{M}$} وي۔
  \item \LR{$\Sat{M}{\Atom{\Obj P}{\Obj v_i}}[s]$} هله او يوازې هله چې
  \LR{$s(\Obj v_i) \in \Assign{\Obj P}{M}$} وي۔
  \item \LR{$\Sat{M}{\lnot !A}[s]$} هله او يوازې هله چې
  \LR{$\Sat{M}{!A}[s]$} نۀ وي۔
  \item \LR{$\Sat{M}{(!A \land !B)}[s]$} هله او يوازې هله چې
  \LR{$\Sat{M}{!A}[s]$} او \LR{$\Sat{M}{!B}[s]$} دواړه وي۔
\end{enumerate}
ښه، په دې سره يوه ستونزه هواره شوه: اوس ويلے شو چې~\LR{$\Struct{M}$} د
\LR{$s(\Obj v_0)$} قيمت دپاره \LR{$\Atom{\Obj P}{\Obj v_0}$} کله پوره کوي۔
د \LR{$\lexists[\Obj v_0][\Atom{\Obj P}{\Obj v_0}]$} د تعريف د سمولو
دپاره يو بل کار پکار دے: غواړو
\LR{$\Sat{M}{\lexists[\Obj v_0][\Atom{\Obj P}{\Obj v_0}]}[s]$} هله او
يوازې هله وي چې د~\LR{$\Obj v_0$} د قيمت د ګومارلو د \emph{کومې} لارې
\LR{$s'$} دپاره \LR{$\Sat{M}{\Atom{\Obj P}{\Obj v_0}}[s']$} وي۔ خو~\LR{$\Obj v_0$}
ته ګومارلے قيمت لازماً هغه قيمت نۀ دے چې~\LR{$s(\Obj v_0)$} ئې ټاکي۔
د دې دپاره يوه نښه ورپېژنو: که \LR{$m \in \Domain{M}$} وي، نو
\LR{$\Subst{s}{m}{\Obj v_0}$} هغه ګومارنه ده چې (له~\LR{$\Obj v_0$} پرته د
ټولو متغيرونه دپاره) کټ مټ د~\LR{$s$} په شان وي، خو~\LR{$\Obj v_0$} ته
\LR{$m$} ګوماري۔ اوس زموږ تعريف داسې کېدے شي:
\begin{enumerate}\setcounter{enumi}{4}
  \item \LR{$\Sat{M}{\lexists[\Obj v_i][!A]}[s]$} هله او يوازې هله چې
  د کوم~\LR{$m \in \Domain{M}$} دپاره
  \LR{$\Sat{M}{!A}[\Subst{s}{m}{\Obj v_i}]$} وي۔
\end{enumerate}
ايا دا کار کوي؟ فرض کړئ د ټولو~\LR{$i \in \Nat$} دپاره
\LR{$s(\Obj v_i) = 0$} ټاکو۔
\LR{$\Sat{M}{\lexists[\Obj v_0][\Atom{\Obj P}{\Obj v_0}]}[s]$} هله او
يوازې هله چې داسې~\LR{$m \in \Domain{M}$} وي چې
\LR{$\Sat{M}{\Atom{\Obj P}{\Obj v_0}}[\Subst{s}{m}{\Obj v_0}]$} وي۔ او
داسې غړے شته: \LR{$m = 1$} يا \LR{$m = 2$} ټاکلے شو۔ پام وکړئ دا حتٰی هغه
وخت رښتيا دي چې~\LR{$s$} پخپله~\LR{$\Obj v_0$} ته ګومارلے قيمت
\LR{$s(\Obj v_0)$}---دلته~\LR{$0$}---کار نۀ کوي۔
\LR{$\Sat{M}{\Atom{\Obj P}{\Obj v_0}}[\Subst{s}{1}{\Obj v_0}]$} لرو، خو
\LR{$\Sat{M}{\Atom{\Obj P}{\Obj v_0}}[s]$} نۀ لرو۔
\olpseganchor{OLP-0143-B009}{end}

\olpseganchor{OLP-0143-B010}{start}
که دا مغشوش او دروند ښکاري، همداسې دے۔ خو دا زياته پېچلتيا د ټولو
فارمولونه د صدق د دقيق، استقرايي تعريف دپاره اړينه ده، او د
معنايي مفاهيمو د دقيق تعريف دپاره ورته څه پکار دي۔ نورې لارې هم
شته، خو ټولې په همدې کچه ښکلې (يا ناشاعرانه) دي۔
\olpseganchor{OLP-0143-B010}{end}

\olpseganchor{OLP-0144-B004}{start}
\olfileid{fol}{int}{snt}
\olpseganchor{OLP-0144-B004}{end}

\olpseganchor{OLP-0144-B005}{start}
\olsection{تړلي فارمولونه}
\olpseganchor{OLP-0144-B005}{end}

\olpseganchor{OLP-0144-B006}{start}
ښه، اوس د جوړښتونه او فارمولونه دپاره د صدق (``پکښې
رښتونې'') د تعريف يوه طرحه لرو۔ خو دا يو اضافي څيز---د
متغير ګومارنه---غواړي، حال دا چې زموږ غوښتنه د
تړلي فارمولونه تعريف و۔ له ګومارنو څنګه ځان خلاصوو، او تړلي فارمولونه
څه دي؟
\olpseganchor{OLP-0144-B006}{end}

\olpseganchor{OLP-0144-B007}{start}
غالباً د صوري منطق له لومړۍ پېژندنې څخه مو د متغيرونه او کميت
ټاکونکو د اړيکې بحث ياد وي۔ له کميت ټاکونکي څخه تل يو متغير
وروسته راځي، او بيا د تړلے فارمول په هغه برخه کښې چې کميت ټاکونکے
ورباندې پلي کېږي (د هغه ``ساحه'')، متغير د هماغه کميت ټاکونکي
له خوا ``تړلے'' ګڼو۔ په فارمولونه کښې دا شرط نۀ و چې هر
متغير دې ورته برابر کميت ټاکونکے ولري؛ بې له برابر کميت
ټاکونکي متغيرونه ``ازاد'' يا ``ناتړلي'' دي۔ موږ به ټول هغه
فارمولونه تړلي فارمولونه بولو چې ازاد متغيرونه نۀ لري۔
\olpseganchor{OLP-0144-B007}{end}

\olpseganchor{OLP-0144-B008}{start}
په فارمول~\LR{$!A$} کښې د متغير د يوې پېښې د تړل کېدو، د
هغې د تړونکي کميت ټاکونکي، او د هغې د ازاد پاتې کېدو وجداني مفکوره
بيا هم ګرانه نۀ ده۔ ښايي د دې د ازمېښت کومه طريقه مو زده کړې وي،
چې ښايي د قوسونو شمېرل هم پکښې وي۔ خو موږ بايد پر دقيق تعريف ټينګار
وکړو---او ځکه چې فارمولونه مو په استقرا تعريف کړي، په
فارمول~\LR{$!A$} کښې د متغير~\LR{$x$} ازادې او تړلې پېښې هم په
استقرا تعريفولے شو۔ مثلاً زموږ د ساده شوې ژبې دپاره داسې کېدے شي:
\begin{enumerate}
  \item که \LR{$!A$} اټومي وي، د~\LR{$x$} ټولې پېښې پکښې ازادې دي (يعنې په
  ~ \LR{$\Atom{\Obj P}{x}$} کښې د~\LR{$x$} پېښه ازاده ده)۔
  \item که \LR{$!A$} د~\LR{$\lnot !B$} بڼه ولري، نو په~\LR{$\lnot !B$} کښې د~\LR{$x$}
  يوه پېښه هله او يوازې هله ازاده ده چې په~\LR{$!B$} کښې د~\LR{$x$} ورته
  پېښه ازاده وي (يعنې په~\LR{$!B$} کښې د متغيرونو ازادې پېښې کټ مټ په
  ~ \LR{$\lnot !B$} کښې ورته پېښې دي)۔
  \item که \LR{$!A$} د~\LR{$(!B \land !C)$} بڼه ولري، نو په~\LR{$(!B \land !C)$}
  کښې د~\LR{$x$} يوه پېښه هله او يوازې هله ازاده ده چې په~\LR{$!B$} يا~\LR{$!C$}
  کښې د~\LR{$x$} ورته پېښه ازاده وي۔
  \item که \LR{$!A$} د~\LR{$\lexists[x][!B]$} بڼه ولري، نو په~\LR{$!A$} کښې د~\LR{$x$}
  هېڅ پېښه ازاده نۀ ده؛ او که د~\LR{$\lexists[y][!B]$} بڼه ولري، چې~\LR{$y$}
  له~\LR{$x$} څخه بېل متغير وي، نو په~\LR{$\lexists[y][!B]$} کښې د~\LR{$x$}
  يوه پېښه هله او يوازې هله ازاده ده چې په~\LR{$!B$} کښې د~\LR{$x$} ورته
  پېښه ازاده وي۔
\end{enumerate}
\olpseganchor{OLP-0144-B008}{end}

\olpseganchor{OLP-0144-B009}{start}
کله چې د متغيرونو د ازادو او تړلو پېښو دقيق تعريف ولرو، په ساده
ډول ويلے شو: تړلے فارمول هر هغه فارمول ده چې د
متغيرونه ازادې پېښې نۀ لري۔
\olpseganchor{OLP-0144-B009}{end}

\olpseganchor{OLP-0145-B004}{start}
\olfileid{fol}{int}{sem}
\olpseganchor{OLP-0145-B004}{end}

\olpseganchor{OLP-0145-B005}{start}
\olsection{معنايي مفاهيم}
\olpseganchor{OLP-0145-B005}{end}

\olpseganchor{OLP-0145-B006}{start}
پورته مو يادونه وکړه چې کله ګورو ايا \LR{$\Sat{M}{!A}[s]$} صدق کوي، د
اسانتيا دپاره~\LR{$s$} ته پرېږدو چې ټولو متغيرونه ته قيمتونه
وګوماري، خو يوازې هغه قيمتونه کارېږي چې په~\LR{$!A$} کښې
متغيرونه ته ئې ګوماري۔ په حقيقت کښې، يوازې په~\LR{$!A$} کښې د
\emph{ازادو} متغيرونو قيمتونه اهميت لري۔ البته ځکه چې موږ احتياط
کوو، دا حقيقت به ثابتوو۔ ځکه چې تړلي فارمولونه ازاد متغيرونه نۀ
لري، نو په جوړښت کښې د هغوئ د صدق دپاره~\LR{$s$} هېڅ اهميت نۀ
لري۔ نو کله چې~\LR{$!A$} تړلے فارمول وي، \LR{$\Sat{M}{!A}$} داسې تعريفولے
شو چې معنٰی ئې ``د ټولو~\LR{$s$} دپاره \LR{$\Sat{M}{!A}[s]$}'' ده؛ او تصادفاً
دا هله او يوازې هله رښتيا دي چې لږ تر لږه د يوه~\LR{$s$} دپاره
\LR{$\Sat{M}{!A}[s]$} وي۔ د فارمولونه د صدق د کار وړ تعريف دپاره د
متغير ګومارنې ورپېژندل پکار دي، خو د تړلي فارمولونه صدق د
متغير له ګومارنو څخه خپلواک دے۔
\olpseganchor{OLP-0145-B006}{end}

\olpseganchor{OLP-0145-B007}{start}
کله چې د ``\LR{$\Sat{M}{!A}$}'' تعريف ولرو، د لومړۍ درجې منطق د
تړلي فارمولونه په باب پوهېږو چې ``حالت'' او ``پکښې رښتونې'' څه معنٰی
لري۔ د تړلي فارمولونه دپاره د~\LR{$\Sat{M}{!A}$} د تعريف پر بنسټ بيا د
اعتبار، استلزام او صدق وړتيا بنسټيز معنايي مفاهيم تعريفولے شو۔ يوه
جمله، \LR{$\Entails !A$}، هله معتبره ده چې هره جوړښت ئې پوره
کړي۔ د تړلي فارمولونه د يوه سټ له خوا، \LR{$\Gamma \Entails !A$}،
هله لازمېږي چې هره هغه جوړښت چې په~\LR{$\Gamma$} کښې ټولې
تړلي فارمولونه پوره کوي،~\LR{$!A$} هم پوره کړي۔ او د تړلي فارمولونه يو سټ
هله د صدق وړ دے چې کومه جوړښت پکښې ټولې تړلي فارمولونه په
يو وخت پوره کړي۔
\olpseganchor{OLP-0145-B007}{end}

\olpseganchor{OLP-0145-B008}{start}
ځکه چې فارمولونه په استقرا تعريف شوي، او صدق بيا د فارمولونه
پر جوړښت په استقرا تعريفېږي، د خپلې معنٰی پوهنې د خاصيتونو د ثبوت
او د تعريف شوو معنايي مفاهيمو د تړلو دپاره استقرا کارولے شو۔ ځينې
دغه خاصيتونه به راټول او ثابت کړو؛ يوه وجه دا ده چې هر يو پخپله په
زړه پورې دے، خو اصلي وجه دا ده چې وروسته د معنٰی پوهنې او د
اشتقاق د نظامونو د اړيکې په څېړنه کښې ډېر راپکارېږي۔ د دې
کار دپاره به ځينې داسې نحوي مفاهيم او عمليات هم په دقيق ډول (يعنې
په استقرا) تعريفول پکار وي چې لا مو نۀ دي ياد کړي۔
\olpseganchor{OLP-0145-B008}{end}

\olpseganchor{OLP-0146-B004}{start}
\olfileid{fol}{int}{sub}
\olpseganchor{OLP-0146-B004}{end}

\olpseganchor{OLP-0146-B005}{start}
\olsection{ځايناستي}
\olpseganchor{OLP-0146-B005}{end}

\olpseganchor{OLP-0146-B006}{start}
يوه بېلګه څېړو څو څرګنده شي چې خبرې څنګه سره تړلې دي، او د نحو او
معنٰی پوهنې وده زموږ د وروستيو پرمختللو څېړنو بنسټ څنګه جوړوي۔ زموږ
د اشتقاق نظامونه بايد اجازه راکړي چې
\LR{$\Atom{\Obj P}{\Obj a}$} له
\LR{$\lforall[\Obj v_0][\Atom{\Obj P}{\Obj v_0}]$} څخه اشتقاق کړو۔
په مقدمه کښې د عمومي کميت ټاکونکي تر واک لاندې بشپړ اټومي فارمول دے۔
ښايي وغواړو دا
حتٰی د استنتاج د قاعدې په توګه بيان کړو۔ خو د دې دپاره بايد په تر
ټولو عامو اصطلاحاتو کښې ئې ويلے شو: يوازې د~\LR{$\Obj P$}، \LR{$\Obj a$} او
\LR{$\Obj v_0$} دپاره نۀ، بلکې د هر فارمول~\LR{$!A$}، ترم~\LR{$t$} او
متغير~\LR{$x$} دپاره۔ (ياد مو وي چې ثابت سمبولونه ترمونه دي، خو
له ثابت سمبولونه او تابعې څخه جوړ شوي لا پېچلي ترمونه به هم
څېړو۔) نو غواړو داسې څه ووايو: ``هر کله چې مو
\LR{$\lforall[x][!A(x)]$} اشتقاق کړے وي، د~\LR{$!A(t)$} په استنتاج کښې په
حقه يئ---يعنې هغه پايله چې~\LR{$\lforall[x]$} لرې او~\LR{$x$} په~\LR{$t$} بدل
کړي''۔ خو ``\LR{$x$} په~\LR{$t$} بدلول'' په دقيق ډول څه معنٰی لري؟ د~\LR{$!A(x)$}
او~\LR{$!A(t)$} تر منځ اړيکه څه ده؟ ايا دا تل کار کوي؟
\olpseganchor{OLP-0146-B006}{end}

\olpseganchor{OLP-0146-B007}{start}
د دې د دقيق کولو دپاره د \emph{ځايناستي} عمل تعريفوو۔ ځايناستي په
اصل کښې نازکه ده، ځکه په~\LR{$!A$} کښې ټول~\LR{$x$} ګانې په ساده ډول په~\LR{$t$}
نۀ شو بدلولے، او هر~\LR{$t$} د هر~\LR{$x$} پر ځاے نۀ شي راتلے۔ دا به بيا په
استقرايي تعريفونو حل کوو۔ خو چې دا کار وشي، د استنتاج د قاعدې داسې
ټاکل چې ``له~\LR{$\lforall[x][!A(x)]$} څخه~\LR{$!A(t)$} استنتاج کړه'' په دقيق
تعريف بدلېږي۔ سربېره پر دې به وښودلے شو چې دا په دې معنٰی د
استنتاج ښه قاعده ده چې \LR{$\lforall[x][!A(x)]$}،~\LR{$!A(t)$} لازموي۔ خو د
دې د ثبوت دپاره بيا داسې څه ثابتول پکار دي چې په لومړي نظر مو دې
پوښتنې ته هڅوي: ``دا ولې کوو؟'' له~\LR{$\lforall[x][!A(x)]$} څخه د
~\LR{$!A(t)$} لازمېدل پر دې حقيقت ولاړ دي چې \LR{$\Sat{M}{!A(t)}$} صدق کوي
که نۀ، دا يوازې د ترم~\LR{$t$} په قيمت پورې تړلې ده؛ يعنې که~\LR{$m$} د
~\LR{$\Domain{M}$} هر هغه غړے وټاکو چې~\LR{$t$} ئې ټاکي، نو
\LR{$\Sat{M}{!A(t)}[s]$} هله او يوازې هله چې
\LR{$\Sat{M}{!A(x)}[\Subst{s}{m}{x}]$} وي۔ دا حتٰی هغه وخت هم صدق کوي
چې~\LR{$t$} متغيرونه ولري، خو بايد د پايلې په دقيق بيان کښې احتياط
وکړو۔
\olpseganchor{OLP-0146-B007}{end}

\olpseganchor{OLP-0147-B004}{start}
\olfileid{fol}{int}{mod}
\olpseganchor{OLP-0147-B004}{end}

\olpseganchor{OLP-0147-B005}{start}
\olsection{مدلونه او تيورۍ}
\olpseganchor{OLP-0147-B005}{end}

\olpseganchor{OLP-0147-B006}{start}
کله چې د لومړۍ درجې منطق نحو او معنٰی پوهنه تعريف کړو، د
جوړښتونه او معنايي مفاهيمو د خاصيتونو څېړل پيلولے شو۔ د
اشتقاق نظامونه هم تعريف او وڅېړلے شو۔ د تړلي فارمولونه د
يوه سټ په باب پوښتلے شو: کومې جوړښتونه د هغه سټ ټولې
تړلي فارمولونه رښتونې کوي؟ د تړلي فارمولونه يو سټ~\LR{$\Gamma$} چې راکړل
شي، هغه جوړښت~\LR{$\Struct{M}$} چې هغوئ پوره کوي د~\LR{$\Gamma$}
\emph{مدل} بلل کېږي۔ ښايي له~\LR{$\Gamma$} څخه پيل او د هغې مدلونه
ولټوو---هغوئ څنګه دي؟ بايد څومره لوے يا واړه وي؟ خو ښايي له يوې
جوړښت يا د جوړښتونه له يوې ټولګې څخه هم پيل او
پوښتنه وکړو: کومې تړلي فارمولونه پکښې رښتونې دي؟ ايا داسې
تړلي فارمولونه شته چې دا جوړښتونه په دې معنٰی
\emph{ځانګړې} کړي چې هغوئ، او يوازې هغوئ، پکښې رښتونې وي؟ دا ډول
پوښتنې د \emph{مدل تيورۍ} ډګر دے۔ همدغه پوښتنې د \emph{بديهي
طريقې} بنسټ هم جوړوي: د جوړښتونه د يوې ټولګې بيانول د
تړلي فارمولونه په يوه سټ، يعنې د يوې تيورۍ په بديهي اصولو۔ دا د دې
مشاهدې له امله ممکنېږي چې کټ مټ هغه تړلي فارمولونه د بديهي اصولو په
ټولو مدلونو کښې رښتونې دي چې د لومړۍ درجې منطق کښې له هغو بديهي
اصولو څخه لازمېږي۔
\olpseganchor{OLP-0147-B006}{end}

\olpseganchor{OLP-0147-B007}{start}
د يوې ډېرې ساده بېلګې په توګه مخکښ‌مرتبې وګورئ۔ مخکښ‌مرتبه پر کوم
سټ~\LR{$A$} داسې اړيکه~\LR{$R$} ده چې انعکاسي او انتقالي دواړه وي۔ سټ~\LR{$A$}
چې دوه‌ځاييزه اړيکه~\LR{$R \subseteq A \times A$} ورباندې وي، کټ مټ
هغه څه دي چې د يوه دوه‌ځاييز اړوند سمبول~\LR{$\Obj P$} لرونکې د لومړۍ
درجې ژبې د جوړښت دپاره ئې غواړو: \LR{$\Domain{M} = A$} او
\LR{$\Assign{\Obj P}{M} = R$} ټاکو۔ ځکه~\LR{$R$} مخکښ‌مرتبه ده، انعکاسي او
انتقالي ده، او د لومړۍ درجې منطق د تړلي فارمولونه داسې سټ~\LR{$\Gamma$}
موندلے شو چې دا خبره وايي:
\begin{align*}
  & \lforall[\Obj v_0][\Atom{\Obj P}{\Obj v_0,\Obj v_0}]\\
  & \lforall[\Obj v_0][\lforall[\Obj v_1][\lforall[\Obj v_2][((\Atom{\Obj P}{\Obj v_0,\Obj v_1} \land \Atom{\Obj P}{\Obj v_1,\Obj v_2}) \lif \Atom{\Obj P}{\Obj v_0,\Obj v_2})]]]
\end{align*}
دا تړلي فارمولونه يوازې د ``د هر~\LR{$x$} دپاره~\LR{$Rxx$}'' (\LR{$R$} انعکاسي
ده) او ``هر کله چې~\LR{$Rxy$} او~\LR{$Ryz$} وي، نو~\LR{$Rxz$} هم وي'' (\LR{$R$}
انتقالي ده) سمبولي بڼې دي۔ وينو چې جوړښت~\LR{$\Struct{M}$} د
دغو دوو تړلي فارمولونه~\LR{$\Gamma$} هله او يوازې هله مدل دے چې~\LR{$R$}
(يعنې~\LR{$\Assign{\Obj P}{M}$}) پر~\LR{$A$} (يعنې~\LR{$\Domain{M}$}) مخکښ‌مرتبه
وي۔ په بل عبارت، د~\LR{$\Gamma$} مدلونه کټ مټ مخکښ‌مرتبې دي۔ د ټولو
مخکښ‌مرتبو هر هغه خاصيت چې د~\LR{$\Obj P$} يوازيني پريديکات لرونکې
د لومړۍ درجې ژبه کښې بيانېدے شي (لکه پورته انعکاسيت او انتقاليت)،
په~\LR{$\Gamma$} کښې له دوو تړلي فارمولونه څخه لازمېږي، او برعکس۔ نو د
~\LR{$\Gamma$} د مدلونو په باب چې څه ثابتولے شو، د ټولو مخکښ‌مرتبو په
باب مو ثابت کړي دي۔
\olpseganchor{OLP-0147-B007}{end}

\olpseganchor{OLP-0147-B008}{start}
د هرې ځانګړې تيورۍ او د مدلونو د ټولګي (لکه~\LR{$\Gamma$} او ټولې
مخکښ‌مرتبې) دپاره به په اړونده د لومړۍ درجې ژبه کښې د بيانېدونکو او
نۀ بيانېدونکو څيزونو په باب په زړه پورې پوښتنې وي۔ د
جوړښتونه ځينې داسې خاصيتونه شته چې د ټولو ژبو او د مدلونو د
ټولګيو دپاره په زړه پورې دي، يعنې د د تعريف ساحه د کچې اړوند خاصيتونه۔
مثلاً د هر~\LR{$n \in \PosInt$} دپاره تل بيانولے شو چې د تعريف ساحه کټ مټ
\LR{$n$}~غړي لري۔ د نامتناهي ډېرو تړلي فارمولونه د يوه سټ په
کارولو دا هم بيانولے شو چې د تعريف ساحه نامتناهي ده۔ خو دا نۀ شو
بيانولے چې د تعريف ساحه متناهي يا د شمېر نۀ وړ ده۔ د لومړۍ
درجې ژبو د محدوديتونو دا پايلې د متناهي شاهد او
د لوېنهايم--سکولم له قضيو څخه راځي۔
\olpseganchor{OLP-0147-B008}{end}

\olpseganchor{OLP-0148-B004}{start}
\olfileid{fol}{int}{scp}
\olpseganchor{OLP-0148-B004}{end}

\olpseganchor{OLP-0148-B005}{start}
\olsection{سموالے او بشپړتيا}
\olpseganchor{OLP-0148-B005}{end}

\olpseganchor{OLP-0148-B006}{start}
د لومړۍ درجې منطق دپاره به د اشتقاق نظامونه هم ورپېژنو۔
منطقپوهانو ډېر د اشتقاق نظامونه جوړ کړي دي، خو هغوئ ټول د
تړلي فارمولونه تر منځ هماغه د اشتقاق وړتيا اړيکه تعريفوي۔ وايو
چې~\LR{$\Gamma$}،~\LR{$!A$} اشتقاق کوي، \LR{$\Gamma \Proves !A$}، که د يوه
ټاکلي، په دقيق ډول تعريف شوي ډول اشتقاق موجود وي۔
اشتقاقونه تل د سمبولونو متناهي ترتيبونه دي---ښايي د
تړلي فارمولونه لړ وي، يا کوم لا پېچلے جوړښت۔ د اشتقاق د
نظامونو موخه داسې وسيله برابرول دي چې وټاکي ايا کوم سټ~\LR{$\Gamma$}
تړلے فارمول لازموي که نۀ۔ د دې موخې د پوره کولو دپاره بايد
\LR{$\Gamma \Entails !A$} هله او يوازې هله وي چې \LR{$\Gamma \Proves !A$}
وي۔
\olpseganchor{OLP-0148-B006}{end}

\olpseganchor{OLP-0148-B007}{start}
که \LR{$\Gamma \Proves !A$} وي خو \LR{$\Gamma \Entails !A$} نۀ وي، زموږ د
اشتقاق نظام به ډېر پياوړے وي او له حده زيات څه به ثابتوي۔
هغه خاصيت چې له~\LR{$\Gamma \Proves !A$} څخه \LR{$\Gamma \Entails !A$}
لازمېږي، \emph{سموالے} بلل کېږي، او د اشتقاق د هر ښه نظام
لږ تر لږه غوښتنه ده۔ له بلې خوا، که \LR{$\Gamma \Entails !A$} وي خو
\LR{$\Gamma \Proves !A$} نۀ وي، نو زموږ د اشتقاق نظام ډېر کمزورے
دے او کافي څه نۀ ثابتوي۔ هغه خاصيت چې له~\LR{$\Gamma \Entails !A$} څخه
\LR{$\Gamma \Proves !A$} لازمېږي، \emph{بشپړتيا} بلل کېږي۔ سموالے
عموماً نسبتاً اسانه ثابتېږي (د اشتقاقونه پر جوړښت د استقرا له
لارې، ځکه هغوئ په استقرا تعريف شوي دي)۔ بشپړتيا ثابتول ګران دي۔
\olpseganchor{OLP-0148-B007}{end}

\olpseganchor{OLP-0148-B008}{start}
سموالے او بشپړتيا څو مهمې پايلې لري۔ که د تړلي فارمولونه يو
سټ~\LR{$\Gamma$} کوم تناقض (لکه~\LR{$!A \land \lnot !A$}) اشتقاق کړي،
\emph{ناسازګار} بلل کېږي۔ ناسازګار~\LR{$\Gamma$} هېڅ مدل نۀ شي لرلے؛ د
صدق وړ نۀ دے۔ له بشپړتيا څخه ئې عکس لازمېږي: هر هغه~\LR{$\Gamma$} چې
ناسازګار نۀ وي---يا، لکه موږ چې به وايو، \emph{سازګار} وي---يو مدل
لري۔ په حقيقت کښې دا له بشپړتيا سره هم‌ارزښته ده، او موږ به د
بشپړتيا همدا بڼه په اصل کښې ثابتوو۔ دا ژوره او ښايي حيرانوونکې
پايله ده: يوازې دا چې له~\LR{$\Gamma$} څخه~\LR{$!A \land \lnot !A$} نۀ شئ
ثابتولے، تضمينوي چې داسې جوړښت شته چې کټ مټ د~\LR{$\Gamma$} د
بيان په شان ده۔ نو بشپړتيا دې پوښتنې ته ځواب ورکوي: د
تړلي فارمولونه کوم سټونه مدلونه لري؟ ځواب: ټول او يوازې سازګار
سټونه۔
\olpseganchor{OLP-0148-B008}{end}

\olpseganchor{OLP-0148-B009}{start}
د سموالي او بشپړتيا قضيې دوه مهمې پايلې لري: د متناهي شاهد قضيه او
د لوېنهايم--سکولم قضيه۔ دا د مدلونو په تيورۍ کښې مهمې پايلې دي
او د ډېرو په زړه پورې پايلو د ثبوت دپاره کارېدے شي۔ دوه مو لا يادې
کړې دي: د لومړۍ درجې منطق دا نۀ شي بيانولے چې د جوړښت
د تعريف ساحه متناهي يا د شمېر نۀ وړ ده۔
\olpseganchor{OLP-0148-B009}{end}

\olpseganchor{OLP-0148-B010}{start}
په تاريخي ډول د دې ټولو کارونو سم معلومول ډېر وخت ونيو---د لومړۍ
درجې منطق نحو او معنٰی پوهنه څنګه تعريف شي، د اشتقاق ښه
نظامونه څنګه تعريف شي، د هغوئ سموالے او بشپړتيا څنګه ثابته شي، او
د لومړۍ درجې ژبو د بيانېدونکو او نۀ بيانېدونکو څيزونو توپير څنګه
روښانه شي۔ اوس پوهېږو چې دا هر څه څنګه وکړو، خو په ټولو جزئياتو
تېرېدل لا هم مغشوشوونکي او ستومانوونکي کېدے شي۔ بيا هم مهم دي، ځکه
دلته د لومړۍ درجې منطق د صوري ژبې دپاره جوړې شوې طريقې په منطق،
کمپيوټر ساينس او ژبپوهنه کښې هر ځاے پلي کېږي۔ نو د جزئياتو زيار په
اوږده موده کښې ګټه کوي۔
\olpseganchor{OLP-0148-B010}{end}

\olpseganchor{OLP-0149-B004}{start}
\olchapter{fol}{syn}{د لومړۍ درجې منطق نحوه}
\olpseganchor{OLP-0149-B004}{end}

\olfileid{fol}{syn}{combined-driver-olp0646}
\paragraph{د ګډې سرچينې ترتيب}
دا ګډ سرليک د لاندې نحو او معنٰی پوهنې دواړو پرلهپسې بابونو سرچينه‌يي ترتيب ښيي؛ اصلي پنځلس واردات په سمون وړ ګډ فايل کښې خوندي دي، او اړوند برخې دلته بې له تکراري بابونو راځي۔
\olpseganchor{OLP-0646-B004}{start}
\paragraph{د سرچينې ګډ باب: نحو او معنٰی پوهنه}
\olpseganchor{OLP-0646-B004}{end}

\paragraph{د منجمدې سرچينې بشپړ بديل متن}
\olpseganchor{OLP-0645-B004}{start}
\olfileid{fol}{syn}{int-olp0645}
\olpseganchor{OLP-0645-B004}{end}

\olpseganchor{OLP-0645-B005}{start}
\olsection{سريزه}
\olpseganchor{OLP-0645-B005}{end}

\olpseganchor{OLP-0645-B006}{start}
د لومړۍ درجې منطق د تيورۍ او فرا تيورۍ د جوړولو دپاره بايد اول د هغې د عبارتونو نحو او معنٰی پوهنه تعريف کړو۔ د لومړۍ درجې منطق عبارتونه ترمونه او فارمولونه دي۔ ترمونه له متغيرونه، ثابت سمبولونه او تابعې جوړېږي۔ بيا فارمولونه له پريديکاتونه او ترمونو سره يوځاے جوړېږي (همدا تر ټولو کوچنۍ، «اتومي» فارمولونه جوړوي)، او له اتومي فارمولونه څخه د منطقي نښلوونکو او کمیت ټاکونکو په مرسته لا پېچلي فارمولونه جوړولے شو۔ د جوړېدو د قواعدو د ليکلو ډېرې بېلابېلې طريقې شته؛ موږ يوازې يوه ممکنه طريقه ورکوو۔ نور نظامونه به نورې نښې غوره کوي، د لومړنيو نښلوونکو نور سټونه ټاکي، قوسونه به په بله طريقه کاروي (يا به يې بيخي نۀ کاروي، لکه په تش په نامه پولنډۍ نښه ليکنه کښې)۔ خو د ټولو طريقو ګډه خبره دا ده چې د جوړېدو قواعد د ترمونو او فارمولونه سټ \emph{په استقرايي ډول} تعريفوي۔ که کار په سمه توګه وشي، هر عبارت د جوړېدو د قواعدو له مخې اساساً يوازې په يوه طريقه منځ ته راتلے شي۔ هغه استقرايي تعريف چې داسې عبارتونه جوړوي چې \emph{يکتا لوستونتيا} لري، دا معنا لري چې دغو عبارتونو ته په هماغه طريقه، يعنې په استقرايي تعريف، مانا ورکولے شو۔
\olpseganchor{OLP-0645-B006}{end}

\olpseganchor{OLP-0645-B007}{start}
عبارتونو ته مانا ورکول د معنٰی پوهنې کار دے۔ په معنٰی پوهنه کښې مرکزي مفهوم په جوړښت کښې صدق دے۔ جوړښت د ژبې بنسټيزو اجزاوو ته مانا ورکوي: د تعريف ساحه د څيزونو يو غيرتشه سټ دے۔ کمیت ټاکونکي داسې تفسيرېږي چې پر همدې ساحه ګرځي؛ ثابت سمبولونه ته د ساحې غړي، تابعې ته له د تعريف ساحه څخه همدغې ساحې ته تابعې، او پريديکاتونه ته پر د تعريف ساحه اړيکې ګومارل کېږي۔ د تعريف ساحه او بنسټيزو لغتونو ته شوې ګومارنې په ګډه جوړښت جوړوي۔ متغيرونه په فارمولونه کښې راتلے شي، او د معنٰی پوهنې د ورکولو دپاره بايد هغو ته هم د د تعريف ساحه غړي ګومارل شي؛ همدا د متغير ګومارنه ده۔ په پاې کښې د صدق اړيکه دغه څيزونه يوځاے کوي۔ فارمول په جوړښت~\LR{$\Struct{M}$} کښې د متغير ګومارنې~\LR{$s$} په نسبت صدق کولے شي، چې \LR{$\Sat{M}{!A}[s]$} ليکل کېږي۔ دا اړيکه هم د~\LR{$!A$} پر جوړښت په استقرا تعريفېږي؛ د منطقي نښلوونکو د رښتياوالي جدولونه کاروي، څو مثلاً د \LR{$!A \land !B$} صدق د \LR{$!A$} او ~\LR{$!B$} د صدق (يا نۀ صدق) له مخې تعريف کړي۔ بيا معلومېږي چې که فارمول~\LR{$!A$}، تړلے فارمول وي، يعنې ازاد متغيرونه ونۀ لري، د متغير ګومارنه بې اغېزې ده؛ نو د تړلي فارمولونه په اړه يوازې دا ويلے شو چې په جوړښتونه کښې صدق کوي (يا نۀ کوي)۔
\olpseganchor{OLP-0645-B007}{end}

\olpseganchor{OLP-0645-B008}{start}
د جملو دپاره د صدق اړيکې \LR{$\Sat{M}{!A}$} پر بنسټ بيا د اعتبار، استلزام او د صدق وړتيا بنسټيز معنايي مفهومونه تعريفولے شو۔ جمله هله معتبره ده، \LR{$\Entails !A$}، چې په هر جوړښت کښې صدق وکړي۔ د تړلي فارمولونه يو سټ ئې هله لازموي، \LR{$\Gamma \Entails !A$}، چې هر هغه جوړښت چې د~\LR{$\Gamma$} ټولو تړلي فارمولونه ته صدق ورکوي، د~\LR{$!A$} صدق هم ورکړي۔ او د جملو يو سټ هله د صدق وړ دے چې کوم جوړښت د هغه ټولو تړلي فارمولونه ته په يوه وخت کښې صدق ورکړي۔ ځکه چې فارمولونه په استقرايي ډول تعريفېږي، او صدق بيا د فارمولونه پر جوړښت په استقرا تعريفېږي، استقرا د خپلې معنٰی پوهنې د ځانګړتياوو د ثابتولو او د تعريف شويو معنايي مفهومونو د تړلو دپاره کارولے شو۔
\olpseganchor{OLP-0645-B008}{end}

\olpseganchor{OLP-0150-B004}{start}
\olfileid{fol}{syn}{itx}
\olpseganchor{OLP-0150-B004}{end}

\olpseganchor{OLP-0150-B005}{start}
\olsection{پېژندنه}
\olpseganchor{OLP-0150-B005}{end}

\olpseganchor{OLP-0150-B006}{start}
د لومړۍ درجې منطق د تيورۍ او ميتاتيورۍ د پراختيا دپاره بايد تر هر څه
مخکښې د دې منطق د عبارتونو نحوه او معنٰی پوهنه تعريف کړو۔ د لومړۍ درجې
منطق عبارتونه ترمونه او فارمولونه دي۔ ترمونه له متغيرونه،
ثابت سمبولونه او تابعې څخه جوړېږي۔ فارمولونه بيا له
پريديکاتونه او ترمونو څخه جوړېږي (په دې ډول تر ټولو واړۀ، يعنې
``اټومي'' فارمولونه جوړېږي)، او بيا له اټومي فارمولونه څخه د
منطقي تړونکو او کميت ټاکونکو په وسيله نور پېچلي فارمولونه جوړولے شو۔ د
جوړېدو د قاعدو د ټاکلو ډېرې بېلابېلې لارې شته؛ موږ يوازې يوه ممکنه لاره
وړاندې کوو۔ نور نظامونه بېل سمبولونه غوره کوي، د تړونکو بېلې ټولګې د
لومړنيو په توګه ټاکي، او قوسونه په بله طريقه کاروي (يا، لکه د تش په نامه
پولنډۍ ليکنې په حالت کښې، ئې بيخي نۀ کاروي)۔ خو ټولې لارې په دې کښې
شريکې دي چې د جوړېدو قواعد د ترمونو او فارمولونه ټولګه
\emph{په استقرايي ډول} تعريفوي۔ که دا کار سم وشي، نو هر عبارت د جوړېدو
د قواعدو له مخې په بنسټيز ډول يوازې په يوه طريقه تر لاسه کېږي۔ هغه
استقرايي تعريف چې \emph{يوازېنۍ لوستنې وړ} عبارتونه رامنځته کوي، موږ ته
دا امکان راکوي چې همدې طريقې---يعنې استقرايي تعريف---سره دغو عبارتونو ته
معناوې هم ورکړو۔
\olpseganchor{OLP-0150-B006}{end}

\olpseganchor{OLP-0151-B004}{start}
\olfileid{fol}{syn}{fol}
\olpseganchor{OLP-0151-B004}{end}

\olpseganchor{OLP-0151-B005}{start}
\olsection{د لومړۍ درجې ژبې}
\olpseganchor{OLP-0151-B005}{end}


\olpseganchor{OLP-0151-B006}{start}
د لومړۍ درجې د منطق عبارتونه له يوه بنسټيز لغت څخه جوړېږي چې
\emph{متغيرونه}، \emph{ثابت سمبولونه}، \emph{پريديکاتونه}
او کله ناکله \emph{تابعې} لري۔ له دغو سمبولونو سره، د منطقي
نښلوونکو، سورو او د وقف د نښو لکه قوسونو او کامو په کارولو،
\emph{ترمونه} او \emph{فارمولونه} جوړېږي۔
\olpseganchor{OLP-0151-B006}{end}

\olpseganchor{OLP-0151-B007}{start}
\begin{explain}
په غير رسمي ډول، پريديکاتونه د ځانګړنو او اړيکو نومونه دي،
ثابت سمبولونه د انفرادي څيزونو نومونه دي، او تابعې د نقشو
نومونه دي۔ دا ټول، له عينيت~\LR{$\eq$} پرته، \emph{غير منطقي
سمبولونه} دي او په ګډه يوه ژبه جوړوي۔ هره د لومړۍ درجې ژبه~\LR{$\Lang L$}
د خپلو غير منطقي سمبولونو له مخې ټاکل کېږي۔ په تر ټولو عام حالت کښې،
\LR{$\Lang L$} د هر ډول سمبولونه په بې نهايته شمېر لري۔
\end{explain}
\olpseganchor{OLP-0151-B007}{end}

\olpseganchor{OLP-0151-B008}{start}
په عام حالت کښې موږ په لومړۍ درجې منطق کښې لاندې سمبولونه کاروو:
\olpseganchor{OLP-0151-B008}{end}

\olpseganchor{OLP-0151-B009}{start}
\begin{enumerate}
\item منطقي سمبولونه
\begin{enumerate}
\item منطقي نښلوونکي:
  \startycommalist
  \ycomma \LR{$\lnot$} (نفي)%
  \ycomma \LR{$\land$} (عطف)%
  \ycomma \LR{$\lor$} (فصل)%
  \ycomma \LR{$\lif$} (شرطي)%
  \ycomma \LR{$\liff$} (دوه اړخيز شرطي)%
  \ycomma \LR{$\lforall$} (عمومي سور)%
  \ycomma \LR{$\lexists$} (وجودي سور).
\item د دروغوالی~\LR{$\lfalse$} لپاره بياني ثابت۔
\item د رښتياوالی~\LR{$\ltrue$} لپاره بياني ثابت۔
\item دوه ځايه عينيت~\LR{$\eq$}۔
\item د متغيرونه يو شمېرېدونکے لامتناهي سټ: \LR{$\Obj v_0$}، \LR{$\Obj v_1$}، \LR{$\Obj
  v_2$}، \dots
\end{enumerate}
\item غير منطقي سمبولونه، چې د لومړۍ درجې منطق \emph{معياري
  ژبه} جوړوي
\begin{enumerate}
\item د هر \LR{$n>0$} لپاره د \LR{$n$}-ځايه پريديکاتونه يو شمېرېدونکے لامتناهي سټ: \LR{$\Obj
  A^n_0$}، \LR{$\Obj A^n_1$}، \LR{$\Obj A^n_2$}، \dots
\item د ثابت سمبولونه يو شمېرېدونکے لامتناهي سټ: \LR{$\Obj c_0$}، \LR{$\Obj c_1$}، \LR{$\Obj
  c_2$}، \dots۔
\item د هر \LR{$n>0$} لپاره د \LR{$n$}-ځايه تابعې يو شمېرېدونکے لامتناهي سټ:
  \LR{$\Obj f^n_0$}، \LR{$\Obj f^n_1$}، \LR{$\Obj f^n_2$}، \dots
\end{enumerate}
\item د وقف نښې: پرانيستے قوس (، تړلے قوس )، او کامه۔
\end{enumerate}
\olpseganchor{OLP-0151-B009}{end}

\olpseganchor{OLP-0151-B010}{start}
زموږ ډېری تعريفونه او پايلې به د لومړۍ درجې منطق د بشپړې معياري ژبې
لپاره بيانېږي۔ خو د کارونې له ډول سره سم ژبه يوازې څو پريديکاتونه،
ثابت سمبولونه او تابعې ته هم محدودولے شو۔
\olpseganchor{OLP-0151-B010}{end}

\olpseganchor{OLP-0151-B011}{start}
\begin{ex}
د حساب ژبه~\LR{$\Lang L_A$} يو دوه ځايه پريديکات~\LR{$<$}، يو
ثابت سمبول~\LR{$\Obj 0$}، يو يو ځايه تابع~\LR{$\prime$}، او دوه دوه
ځايه تابعې~\LR{$+$} او~\LR{$\times$} لري۔
\end{ex}
\olpseganchor{OLP-0151-B011}{end}

\olpseganchor{OLP-0151-B012}{start}
\begin{ex}
د سټ تيورۍ ژبه~\LR{$\Lang L_Z$} يوازې يو دوه ځايه پريديکات~\LR{$\in$}
لري۔
\end{ex}
\olpseganchor{OLP-0151-B012}{end}

\olpseganchor{OLP-0151-B013}{start}
\begin{ex}
د ترتيبونو ژبه~\LR{$\Lang L_\le$} يوازې يو دوه ځايه پريديکات~\LR{$\le$}
لري۔
\end{ex}
\olpseganchor{OLP-0151-B013}{end}

\olpseganchor{OLP-0151-B014}{start}
دا بيا هم يوازې دودونه دي: په رسمي ډول دا سمبولونه مستعار نومونه دي؛
د بېلګې په توګه، \LR{$<$}، \LR{$\in$} او \LR{$\le$} د \LR{$\Obj A^2_0$}، \LR{$\Obj 0$} د
\LR{$\Obj c_0$}، \LR{$\prime$} د \LR{$\Obj f^1_0$}، \LR{$+$} د \LR{$\Obj f^2_0$}، او \LR{$\times$} د
\LR{$\Obj f^2_1$} مستعار نومونه دي۔
\olpseganchor{OLP-0151-B014}{end}

\olpseganchor{OLP-0151-B015}{start}
.
\olpseganchor{OLP-0151-B015}{end}



\olpseganchor{OLP-0151-B017}{start}
% Alternate symbols
\olpseganchor{OLP-0151-B017}{end}

\olpseganchor{OLP-0151-B018}{start}
\begin{intro}
کېدے شي تاسو له هغو اصطلاحاتو او سمبولونو سره بلد يئ چې زموږ له پورته
کارولو سره توپير لري۔ د منطق متنونه، او ښوونکي، عموماً د ``نفي''
دپاره \LR{$\sim$}، \LR{$\neg$} او~!، او د ``عطف'' دپاره \LR{$\wedge$}، \LR{$\cdot$} او
\LR{$\&$} کاروي۔ د ``شرطي'' يا ``استلزام'' عام سمبولونه \LR{$\rightarrow$}،
\LR{$\Rightarrow$} او \LR{$\supset$} دي۔
د ``دوه اړخيز شرطي''، ``دوه اړخيز استلزام'' يا
  ``(مادي) هم ارزښت'' سمبولونه \LR{$\leftrightarrow$}،
  \LR{$\Leftrightarrow$} او \LR{$\equiv$} دي۔
د \LR{$\lfalse$} سمبول ته په بېلو ځايونو کښې
  ``دروغوالے''، ``فالسوم''، ``محال'' يا ``ښکته'' ويل کېږي۔
د \LR{$\ltrue$} سمبول ته په بېلو ځايونو کښې
  ``رښتياوالے''، ``ويروم'' يا ``پورته'' ويل کېږي۔
\olpseganchor{OLP-0151-B018}{end}

\olpseganchor{OLP-0151-B019}{start}
دود دا دے چې د لاتيني الفبا د پيل واړه توري (لکه \LR{$a$}، \LR{$b$}، \LR{$c$}) د
ثابت سمبولونه لپاره (چې کله ناکله نومونه بلل کېږي)، او د الفبا د پای
واړه توري (لکه \LR{$x$}، \LR{$y$}، \LR{$z$}) د متغيرونه لپاره وکارول شي۔ سورونه
له متغيرونه سره يوځای راځي، لکه \LR{$x$}؛ د ليکدود په بديلونو کښې د
عمومي سور لپاره \LR{$\forall x$}، \LR{$(\forall x)$}، \LR{$(x)$}، \LR{$\Pi x$}،
\LR{$\bigwedge_x$}، او د وجودي سور لپاره \LR{$\exists x$}، \LR{$(\exists
x)$}، \LR{$(Ex)$}، \LR{$\Sigma x$}، \LR{$\bigvee_x$} شامل دي۔
\end{intro}
\olpseganchor{OLP-0151-B019}{end}

\olpseganchor{OLP-0151-B020}{start}
\begin{explain}
موږ ټول بياني عملګرونه او دواړه سورونه د ژبې بنسټيز سمبولونه ګڼلے شو۔
يا د بنسټيزو سمبولونو يوه کوچنۍ ټولګه غوره کولے او نور منطقي عملګرونه
تعريف شوي ګڼلے شو۔ د بوليني عملګرونو په هغو ټولګو کښې چې ``د رښتياوالي
د تابع له مخې بشپړې'' دي، \LR{$\{ \lnot, \lor \}$}، \LR{$\{ \lnot, \land \}$}
او \LR{$\{ \lnot, \lif\}$} شامل دي---په بياني ځواک کښې د بشپړې لومړۍ
درجې ژبې دپاره له دې هره ټولګه له هر يوه سور سره يوځای کېدے شي۔
\olpseganchor{OLP-0151-B020}{end}

\olpseganchor{OLP-0151-B021}{start}
کېدے شي تاسو له دوو نورو منطقي عملګرونه سره هم اشنا يئ: د شېفر
کرښه~\LR{$|$} (چې د هنري شېفر په نوم نومول شوې ده)، او د پيرس
غشی~\LR{$\downarrow$}، چې د کواین خنجر هم بلل کېږي۔ کله چې دې عملګرونو ته په
ترتيب سره د ``نېنډ'' او ``نور'' معمولې لوستنې ورکړل شي، هر يو ئې په
يوازې ځان د رښتياوالي د تابع له مخې بشپړ وي۔
\end{explain}
\olpseganchor{OLP-0151-B021}{end}

\olpseganchor{OLP-0152-B004}{start}
\olfileid{fol}{syn}{frm}
\olpseganchor{OLP-0152-B004}{end}

\olpseganchor{OLP-0152-B005}{start}
\olsection{ترمونه او فارمولونه}
\olpseganchor{OLP-0152-B005}{end}

\olpseganchor{OLP-0152-B006}{start}
يوه د لومړۍ درجې ژبه~\LR{$\Lang L$} چې راکړل شي، د~\LR{$\Lang L$} له بنسټيز
لغت څخه جوړ شوي عبارتونه تعريفولے شو۔ په دغو کښې په ځانګړي ډول
\emph{ترمونه} او \emph{فارمولونه} شامل دي۔
\olpseganchor{OLP-0152-B006}{end}

\olpseganchor{OLP-0152-B007}{start}
\begin{defn}[ترمونه]
\ollabel{defn:terms}
د~\LR{$\Lang L$} د \emph{ترمونو} سټ~\LR{$\Trm[L]$} په استقرايي ډول داسې
تعريفېږي:
\begin{enumerate}
\item هر متغير يو ترم دے۔
\item د~\LR{$\Lang L$} هر ثابت سمبول يو ترم دے۔
\item که \LR{$f$} يو \LR{$n$}-ځايه تابع وي او \LR{$t_1$}، \dots، \LR{$t_n$}
  ترمونه وي، نو \LR{$\Atom{f}{t_1, \ldots, t_n}$} يو ترم دے۔
\item 
  بل هېڅ څيز ترم نۀ دے۔
\end{enumerate}
هغه ترم چې هېڅ متغيرونه نۀ لري، \emph{تړلے ترم} بلل کېږي۔
\end{defn}
\olpseganchor{OLP-0152-B007}{end}

\olpseganchor{OLP-0152-B008}{start}
\begin{explain}
ثابت سمبولونه زموږ د ژبې په مشخصاتو او د ترمونو په تعريف کښې د
سمبولونو د يوې جلا ډلې په توګه راځي، خو د صفر ځايه تابعې په
توګه هم ور شاملېدے شول۔ بيا د ترمونو د تعريف دويم شرط ته اړتيا نۀ
وه۔ يوازې دا بايد ومنو چې که \LR{$n = 0$} وي، نو
\LR{$\Atom{f}{t_1, \ldots, t_n}$} پخپله يوازې \LR{$f$} معنا لري۔
\end{explain}
\olpseganchor{OLP-0152-B008}{end}

\olpseganchor{OLP-0152-B009}{start}
\begin{defn}[فارمولونه]
\ollabel{defn:formulas}
د ژبې~\LR{$\Lang L$} د \emph{فارمولونه} سټ~\LR{$\Frm[L]$} په استقرايي
ډول داسې تعريفېږي:
\begin{enumerate}
\item \LR{$\lfalse$} يو اټومي فارمول دے۔
\olpseganchor{OLP-0152-B009}{end}

\olpseganchor{OLP-0152-B010}{start}
\item \LR{$\ltrue$} يو اټومي فارمول دے۔
\olpseganchor{OLP-0152-B010}{end}

\olpseganchor{OLP-0152-B011}{start}
\item که \LR{$R$} د~\LR{$\Lang L$} يو \LR{$n$}-ځايه پريديکات وي او \LR{$t_1$}، \dots،
  \LR{$t_n$} د~\LR{$\Lang L$} ترمونه وي، نو \LR{$\Atom{R}{t_1,\ldots, t_n}$} يو
  اټومي فارمول دے۔
\olpseganchor{OLP-0152-B011}{end}

\olpseganchor{OLP-0152-B012}{start}
\item که \LR{$t_1$} او \LR{$t_2$} د~\LR{$\Lang L$} ترمونه وي، نو
  \LR{$\Atom{\eq}{t_1, t_2}$} يو اټومي فارمول دے۔
\olpseganchor{OLP-0152-B012}{end}

\olpseganchor{OLP-0152-B013}{start}
\item که \LR{$!A$} فارمول وي، نو \LR{$\lnot !A$} هم
  فارمول دے۔
\olpseganchor{OLP-0152-B013}{end}

\olpseganchor{OLP-0152-B014}{start}
\item که \LR{$!A$} او \LR{$!B$} فارمولونه وي، نو \LR{$(!A \land
  !B)$} فارمول دے۔
\olpseganchor{OLP-0152-B014}{end}

\olpseganchor{OLP-0152-B015}{start}
\item که \LR{$!A$} او \LR{$!B$} فارمولونه وي، نو \LR{$(!A \lor !B)$}
  فارمول دے۔
\olpseganchor{OLP-0152-B015}{end}

\olpseganchor{OLP-0152-B016}{start}
\item که \LR{$!A$} او \LR{$!B$} فارمولونه وي، نو \LR{$(!A \lif !B)$}
  فارمول دے۔
\olpseganchor{OLP-0152-B016}{end}

\olpseganchor{OLP-0152-B017}{start}
\item که \LR{$!A$} او \LR{$!B$} فارمولونه وي، نو \LR{$(!A \liff !B)$}
  فارمول دے۔
\olpseganchor{OLP-0152-B017}{end}

\olpseganchor{OLP-0152-B018}{start}
\item که \LR{$!A$} فارمول او \LR{$x$} متغير وي، نو
  \LR{$\lforall[x][!A]$} فارمول دے۔
\olpseganchor{OLP-0152-B018}{end}

\olpseganchor{OLP-0152-B019}{start}
\item که \LR{$!A$} فارمول او \LR{$x$} متغير وي، نو
  \LR{$\lexists[x][!A]$} فارمول دے۔
\olpseganchor{OLP-0152-B019}{end}

\olpseganchor{OLP-0152-B020}{start}
\item بل هېڅ څيز فارمول نۀ دے۔
\end{enumerate}
\end{defn}
\olpseganchor{OLP-0152-B020}{end}

\olpseganchor{OLP-0152-B021}{start}
\begin{explain}
د ترمونو د سټ او د فارمولونه د سټ تعريفونه \emph{استقرايي
تعريفونه} دي۔ په اصل کښې موږ د فارمولونه سټ په نامتناهي ډېرو
پړاوونو کښې جوړوو۔ په لومړني پړاو کښې ټول اټومي فارمولونه فارمولونه
ګڼو؛ دا د تعريف له لومړنيو څو حالتونو سره سمون لري، يعنې د
\LR{$\ltrue$}، %
\LR{$\lfalse$}، %
\LR{$\Atom{R}{t_1,\dots,t_n}$} او \LR{$\Atom{\eq}{t_1,t_2}$} حالتونه۔ نو
``اټومي فارمول'' د همدې بڼو هر فارمول ته وايي۔
\olpseganchor{OLP-0152-B021}{end}

\olpseganchor{OLP-0152-B022}{start}
د تعريف نور حالتونه له هغو فارمولونه څخه چې لا جوړ شوي وي، د نويو
فارمولونه د جوړولو قاعدې راکوي۔ په دويم پړاو کښې د دغو قاعدو په
کارولو له اټومي فارمولونه څخه فارمولونه جوړولے شو۔ په درېيم
پړاو کښې له اټومي فارمولونو او په دويم پړاو کښې له ترلاسه شوو فارمولونو
څخه نوي فارمولونه جوړوو، او همداسې مخکښې۔ فارمول هر هغه څيز
دے چې بالاخره په داسې يوه پړاو کښې جوړ شي، او بل هېڅ څيز نۀ دے۔
\end{explain}
\olpseganchor{OLP-0152-B022}{end}

\olpseganchor{OLP-0152-B023}{start}
د دود له مخې، \LR{$\eq$} د خپلو مدخلو تر منځ ليکو او قوسونه پرېږدو:
\LR{$\eq[t_1][t_2]$} د \LR{$\Atom{\eq}{t_1,t_2}$} لنډون دے۔ همدارنګه
\LR{$\lnot \Atom{\eq}{t_1,t_2}$} د \LR{$\eq/[t_1][t_2]$} په بڼه لنډوو۔ کله
چې د دوه ځايه نښلوونکي~\LR{$\ast$} په کارولو له \LR{$!B$} او \LR{$!C$} څخه جوړ
شوے فارمول \LR{$(!B \ast !C)$} ليکو، ډېر ځله تر ټولو دباندې جوړه قوسونه
پرېږدو او يوازې~\LR{$!B \ast !C$} ليکو۔
\olpseganchor{OLP-0152-B023}{end}

\olpseganchor{OLP-0152-B024}{start}
\begin{intro}
د منطق ځينې متنونه د دې دپاره چې
\LR{$\lexists[x][!A]$} %
او %
\LR{$\lforall[x][!A]$} %
فارمولونه
وګڼل شي، دا شرط لګوي چې متغير~\LR{$x$} بايد په~\LR{$!A$} کښې راشي۔
که دا شرط ونۀ لګوئ، کومه ستونزه نۀ پيدا کېږي او کار اسانه کېږي۔
\end{intro}
\olpseganchor{OLP-0152-B024}{end}

\olpseganchor{OLP-0152-B036}{start}
کله چې د يوې ځانګړې کارونې په ژبه کښې کار کوو، ډېر ځله دوه ځايه
پريديکاتونه او تابعې د اړوندو ترمونو تر منځ ليکو؛ مثلاً د
حساب په ژبه کښې \LR{$t_1 < t_2$} او \LR{$(t_1 + t_2)$}، او د سټ تيورۍ په ژبه
کښې \LR{$t_1 \in t_2$} ليکو۔ د حساب په ژبه کښې د ناستي تابع د دود له مخې
آن د خپل مدخل \emph{وروسته} ليکل کېږي:~\LR{$t'$}۔ خو په رسمي ډول دا په
ترتيب سره يوازې د \LR{$\Atom{\Obj A^2_0}{t_1, t_2}$}،
\LR{$\Obj f^2_0(t_1, t_2)$}، \LR{$\Atom{\Obj A^2_0}{t_1, t_2}$} او
\LR{$\Obj f^1_0(t)$} دوديز لنډونونه دي۔
\olpseganchor{OLP-0152-B036}{end}

\olpseganchor{OLP-0152-B037}{start}
\begin{defn}[نحوي يوشانوالے]
سمبول \LR{$\ident$} د سمبولونو د تارونو تر منځ نحوي يوشانوالے څرګندوي؛
يعنې \LR{$!A \ident !B$} هله او يوازې هله چې \LR{$!A$} او \LR{$!B$} د يو شان
اوږدوالي لرونکي د سمبولونو تارونه وي او په هر ځاے کښې يو شان سمبول
ولري۔
\end{defn}
\olpseganchor{OLP-0152-B037}{end}

\olpseganchor{OLP-0152-B038}{start}
د \LR{$\ident$} سمبول دواړو خواوو ته هغه تارونه هم ليکل کېدے شي چې په نښلولو
ترلاسه شوي وي؛ مثلاً \LR{$!A \ident (!B \lor !C)$} دا معنا لري: د
سمبولونو تار~\LR{$!A$} له هغه تار سره يو شان دے چې په همدې ترتيب د پرانيستي
قوس، تار \LR{$!B$}، سمبول \LR{$\lor$}، تار~\LR{$!C$} او تړلي قوس له نښلولو ترلاسه
کېږي۔ که دا حالت وي، نو پوهېږو چې د~\LR{$!A$} لومړے سمبول پرانيستے قوس
دے، \LR{$!A$} د دويم سمبول له ځايه \LR{$!B$} د يوه فرعي تار په توګه لري، له
هغه فرعي تار وروسته~\LR{$\lor$} راځي، او همداسې نور۔
\olpseganchor{OLP-0152-B038}{end}

\olpseganchor{OLP-0152-B039}{start}
څرنګه چې ترمونه او فارمولونه د استقرايي تعريفونو له لارې له بنسټيزو
اجزاوو څخه جوړېږي، د هغو په باب د څه ثابتولو دپاره لاندې استقرايي
اصول کارولے شو۔
\olpseganchor{OLP-0152-B039}{end}

\olpseganchor{OLP-0152-B040}{start}
\begin{lem}[\emph{پر ترمونو د استقرا اصل}]
    \ollabel{lem:trmind}
    فرض کړئ \LR{$\Lang L$} د لومړۍ درجې يوه ژبه ده۔
    که يو خاصيت~\LR{$P$} داسې وي چې
    %
    \begin{enumerate}
        \item د هر متغير~\LR{$v$} دپاره صدق کوي؛
        %
        \item د~\LR{$\Lang L$} د هر ثابت سمبول~\LR{$a$} دپاره صدق کوي؛ او
        %
        \item کله چې د \LR{$t_1$}،~\dots، \LR{$t_n$} دپاره صدق کوي او \LR{$f$}
        د~\LR{$\Lang L$} يو \LR{$n$}-ځايه تابع وي، د
        \LR{$f(t_1,\dotsc,t_n)$} دپاره هم صدق کوي،
    \end{enumerate}
    (په دې فرض چې \LR{$t_1$}،~\dots، \LR{$t_n$} د~\LR{$\Lang{L}$} ترمونه دي)،
    نو \LR{$P$} د~\LR{$\Trm[L]$} د هر ترم دپاره صدق کوي۔
\end{lem}
\olpseganchor{OLP-0152-B040}{end}

\olpseganchor{OLP-0152-B041}{start}
\begin{prob}
    \olref[fol][syn][frm]{lem:trmind} ثابته کړئ۔
\end{prob}
\olpseganchor{OLP-0152-B041}{end}

\olpseganchor{OLP-0152-B042}{start}
\begin{lem}[\emph{پر فارمولونه د استقرا اصل}]
    \ollabel{thm:frmind}
    فرض کړئ \LR{$\Lang L$} د لومړۍ درجې يوه ژبه ده۔
    که يو خاصيت~\LR{$P$} د ټولو اټومي فارمولونه دپاره صدق کوي او داسې
    وي چې
    %
    \begin{enumerate}
        \item کله چې د~\LR{$!A$} دپاره صدق کوي، د
            \LR{$\lnot !A$} دپاره هم صدق کوي؛
        \item کله چې د \LR{$!A$} او~\LR{$!B$} دپاره صدق کوي، د
            \LR{$(!A \land !B)$} دپاره هم صدق کوي؛
        \item کله چې د \LR{$!A$} او~\LR{$!B$} دپاره صدق کوي، د
            \LR{$(!A \lor !B)$} دپاره هم صدق کوي؛
        \item کله چې د \LR{$!A$} او~\LR{$!B$} دپاره صدق کوي، د
            \LR{$(!A \lif !B)$} دپاره هم صدق کوي؛
        \item کله چې د \LR{$!A$} او~\LR{$!B$} دپاره صدق کوي، د
            \LR{$(!A \liff !B)$} دپاره هم صدق کوي؛
        \item کله چې د~\LR{$!A$} دپاره صدق کوي، د
            \LR{$\lexists[x][!A]$} دپاره هم صدق کوي؛
        \item کله چې د~\LR{$!A$} دپاره صدق کوي، د
            \LR{$\lforall[x][!A]$} دپاره هم صدق کوي؛
  \end{enumerate}
  (په دې فرض چې \LR{$!A$} او \LR{$!B$} د~\LR{$\Lang{L}$} فارمولونه دي)،
  نو \LR{$P$} د~\LR{$\Frm[L]$} د ټولو فارمولونو دپاره صدق کوي۔
\end{lem}
\olpseganchor{OLP-0152-B042}{end}

\olpseganchor{OLP-0153-B004}{start}
\olfileid{fol}{syn}{unq}
\olpseganchor{OLP-0153-B004}{end}

\olpseganchor{OLP-0153-B005}{start}
\olsection{يوازينے لوست}
\olpseganchor{OLP-0153-B005}{end}

\olpseganchor{OLP-0153-B006}{start}
\begin{explain}
موږ چې فارمولونه په کومه طريقه تعريف کړي دي، هغه دا تضمينوي چې هر
فارمول يو \emph{يوازينے لوست} لري؛ يعنې د فارمولونه د جوړولو
زموږ د قاعدو له مخې ئې د جوړولو په اصل کښې يوازې يوه طريقه شته، او د
``تفسيرولو'' ئې هم يوازې يوه طريقه شته۔ که داسې نۀ واے، مبهم
فارمولونه به مو لرل؛ يعنې داسې فارمولونه چې له يوه څخه زيات
لوستونه يا تفسيرونه لري---او څرګنده ده چې موږ ترې ډډه کول غواړو۔ خو
تر دې مهمه دا ده چې له دې خاصيت پرته به زموږ ډېری راتلونکي تعريفونه
او ثبوتونه کار ورنۀ کړي۔
\olpseganchor{OLP-0153-B006}{end}

\olpseganchor{OLP-0153-B007}{start}
د دې خبرې د روښانولو ښايي غوره لار دا وي چې وګورو، که د فارمولونه
د جوړولو داسې ناسمې قاعدې مو ورکړې واے چې يوازينے لوست نۀ تضمينوي،
څه به پېښ شوي وو۔ مثلاً کېدے شي د نښلوونکو د جوړولو په قاعدو کښې مو
قوسونه هېر کړي او دا مو روا ګڼلي واے:
\begin{quote}
که \LR{$!A$} او \LR{$!B$} فارمولونه وي، نو \LR{$!A \lif !B$} هم فارمول دے۔
\end{quote}
له يوه اټومي فارمول \LR{$!D$} څخه په پيلولو به دې قاعدې د \LR{$!D \lif !D$}
جوړول را ته روا کړي وو۔ له دې او \LR{$!D$} څخه به مو \LR{$!D \lif !D \lif
!D$} ترلاسه کړے و۔ خو دا په دوو طريقو جوړېدے شي:
\begin{enumerate}
\item \LR{$!D$} د \LR{$!A$} په توګه، او \LR{$!D \lif !D$} د \LR{$!B$} په توګه اخلو۔
\item \LR{$!A$} د \LR{$!D \lif !D$} په توګه اخلو، او \LR{$!B$} پخپله~\LR{$!D$} دے۔
\end{enumerate}
په همدې ډول د فارمول~\LR{$!D \lif !D \lif !D$} د ``لوستلو'' دوه
طريقې شته۔ دا د \LR{$!B \lif !C$} بڼه لري چې \LR{$!B$} پکښې \LR{$!D$} او \LR{$!C$}
پکښې \LR{$!D \lif !D$} دے، خو \emph{همدارنګه} د \LR{$!B \lif !C$} بڼه لري
چې \LR{$!B$} پکښې \LR{$!D \lif !D$} او \LR{$!C$} پکښې~\LR{$!D$} دے۔
\olpseganchor{OLP-0153-B007}{end}

\olpseganchor{OLP-0153-B008}{start}
که داسې وشي، زموږ تعريفونه به تل کار نۀ کوي۔ مثلاً کله چې د يوه فارمول
اصلي عملګر تعريفوو، وايو: د \LR{$!B \lif !C$} بڼه لرونکي فارمول
اصلي عملګر د~\LR{$\lif$} همدا ښودل شوے مورد دے۔ خو که فارمول \LR{$!D
\lif !D \lif !D$} له \LR{$!B \lif !C$} سره په پورته يادو شوو دوو بېلو
طريقو برابرولے شو، نو په يوه حالت کښې د \LR{$\lif$} لومړے مورد او په بل
حالت کښې ئې دويم مورد د اصلي عملګر په توګه ترلاسه کوو۔ خو موږ
غواړو چې اصلي عملګر د فارمول يوه \emph{تابع} وي؛ يعنې هر
فارمول بايد کټ مټ د اصلي عملګر يو مورد ولري۔
\end{explain}
\olpseganchor{OLP-0153-B008}{end}

\olpseganchor{OLP-0153-B009}{start}
\begin{lem}
په فارمول~\LR{$!A$} کښې د کيڼو او ښيو قوسونو شمېر سره برابر دے۔
\end{lem}
\olpseganchor{OLP-0153-B009}{end}

\olpseganchor{OLP-0153-B010}{start}
\begin{proof}
دا د \LR{$!A$} د جوړېدو پر طريقه په استقرا ثابتوو۔ دې ته دوه څيزونه پکار
دي: (الف) لومړے بايد ثابته کړو چې ټول اټومي فارمولونه ياد خاصيت لري
(د استقرا بنسټ)۔ (ب) بيا بايد ثابته کړو چې کله له راکړل شوو فارمولونو
څخه نوي فارمولونه جوړوو، که زاړه فارمولونه ياد خاصيت ولري، نو نوي ئې
هم لري۔
\olpseganchor{OLP-0153-B010}{end}

\olpseganchor{OLP-0153-B011}{start}
\LR{$l(!A)$} دې په~\LR{$!A$} کښې د کيڼو قوسونو شمېر وي او \LR{$r(!A)$} دې د ښيو
قوسونو شمېر وي؛ \LR{$l(t)$} او \LR{$r(t)$} دې په همدې ډول په يوه ترم~\LR{$t$} کښې
د کيڼو او ښيو قوسونو شمېرونه وي۔
\olpseganchor{OLP-0153-B011}{end}

\olpseganchor{OLP-0153-B012}{start}
\begin{prob}
    ثابته کړئ چې د هر ترم~\LR{$t$} دپاره \LR{$l(t) = r(t)$}۔
\end{prob}
\olpseganchor{OLP-0153-B012}{end}

\olpseganchor{OLP-0153-B013}{start}
\begin{enumerate}
\item \indcase{!A}{\lfalse}{\LR{$\indfrm$} \LR{$0$} کيڼ او \LR{$0$}
    ښي قوسونه لري۔}
\olpseganchor{OLP-0153-B013}{end}

\olpseganchor{OLP-0153-B014}{start}
\item \indcase{!A}{\ltrue}{\LR{$\indfrm$} \LR{$0$} کيڼ او \LR{$0$}
    ښي قوسونه لري۔}
\olpseganchor{OLP-0153-B014}{end}

\olpseganchor{OLP-0153-B015}{start}
\item \indcase{!A}{\Atom{R}{t_1,\dots,t_n}}{\LR{$l(\indfrm) = 1 + l(t_1) +
  \dots + l(t_n) = 1 + r(t_1) + \dots + r(t_n) = r(\indfrm)$}۔ دلته
  هغه حقيقت کاروو چې د تمرين په توګه پرېښودل شوے دے: د هر ترم~\LR{$t$}
  دپاره \LR{$l(t) = r(t)$}۔}
\olpseganchor{OLP-0153-B015}{end}

\olpseganchor{OLP-0153-B016}{start}
\item \indcase{!A}{\eq[t_1][t_2]}{\LR{$l(\indfrm) = l(t_1) + l(t_2) =
  r(t_1) + r(t_2) = r(\indfrm)$}۔}
\olpseganchor{OLP-0153-B016}{end}

\olpseganchor{OLP-0153-B017}{start}
\item \indcase{!A}{\lnot !B}{د استقرايي فرض له مخې
    \LR{$l(!B) = r(!B)$}۔ نو \LR{$l(\indfrm) = l(!B) = r(!B) =
    r(\indfrm)$}۔}
\olpseganchor{OLP-0153-B017}{end}

\olpseganchor{OLP-0153-B018}{start}
\item \indcase{!A}{(!B \ast !C)}{د استقرايي فرض له مخې \LR{$l(!B) =
  r(!B)$} او \LR{$l(!C) = r(!C)$}۔ نو \LR{$l(\indfrm) = 1 + l(!B) + l(!C) =
  1 + r(!B) + r(!C) = r(\indfrm)$}۔}
\olpseganchor{OLP-0153-B018}{end}

\olpseganchor{OLP-0153-B019}{start}
\item \indcase{!A}{\lforall[x][!B]}{د استقرايي فرض له مخې
    \LR{$l(!B) = r(!B)$}۔ نو \LR{$l(\indfrm) = l(!B) = r(!B) =
    r(\indfrm)$}۔}
\olpseganchor{OLP-0153-B019}{end}

\olpseganchor{OLP-0153-B020}{start}
% Print case for \lexists if prvEx and not prvAll

\olpseganchor{OLP-0153-B020}{end}

\olpseganchor{OLP-0153-B021}{start}
% Just say similarly if prvEx and prvAll
\item \indcase{!A}{\lexists[x][!B]}{په همدې ډول۔}
\end{enumerate}
\end{proof}
\olpseganchor{OLP-0153-B021}{end}

\olpseganchor{OLP-0153-B022}{start}
\begin{defn}[حقيقي مخکښ]
د سمبولونو تار \LR{$!B$} د سمبولونو د تار~\LR{$!A$} \emph{حقيقي مخکښ} دے، که
له \LR{$!B$} سره د سمبولونو د يوه ناتش تار نښلول~\LR{$!A$} ترلاسه کړي۔
\end{defn}
\olpseganchor{OLP-0153-B022}{end}

\olpseganchor{OLP-0153-B023}{start}
\begin{lem}\ollabel{lem:no-prefix}
که \LR{$!A$} فارمول وي او \LR{$!B$} د \LR{$!A$} حقيقي مخکښ وي، نو \LR{$!B$}
فارمول نۀ دے۔
\end{lem}
\olpseganchor{OLP-0153-B023}{end}

\olpseganchor{OLP-0153-B024}{start}
\begin{proof}
تمرين۔
\end{proof}
\olpseganchor{OLP-0153-B024}{end}

\olpseganchor{OLP-0153-B025}{start}
\begin{prob}
\olref[fol][syn][unq]{lem:no-prefix} ثابته کړئ۔
\end{prob}
\olpseganchor{OLP-0153-B025}{end}

\olpseganchor{OLP-0153-B026}{start}
\begin{prop}
\ollabel{prop:unique-atomic}
که \LR{$!A$} يو اټومي فارمول وي، نو له لاندې شرطونو څخه يو، او يوازې
يو، پوره کوي۔
\begin{enumerate}
\item \LR{$!A \ident \lfalse$}۔
\item \LR{$!A \ident \ltrue$}۔
\item \LR{$!A \ident \Atom{R}{t_1,\dots,t_n}$}، چې \LR{$R$} يو \LR{$n$}-ځايه
  پريديکات، \LR{$t_1$}، \dots، \LR{$t_n$} ترمونه، او هر يو \LR{$R$}،
  \LR{$t_1$}، \dots، \LR{$t_n$} په يوازيني ډول ټاکل شوي دي۔
\item \LR{$!A \ident \eq[t_1][t_2]$}، چې \LR{$t_1$} او \LR{$t_2$} په يوازيني ډول
  ټاکل شوي ترمونه دي۔
\end{enumerate}
\end{prop}
\olpseganchor{OLP-0153-B026}{end}

\olpseganchor{OLP-0153-B027}{start}
\begin{proof}
تمرين۔
\end{proof}
\olpseganchor{OLP-0153-B027}{end}

\olpseganchor{OLP-0153-B028}{start}
\begin{prob}
\olref[fol][syn][unq]{prop:unique-atomic} ثابته کړئ۔ (لارښوونه: د
ترمونو دپاره د \olref[fol][syn][unq]{lem:no-prefix} يوه بڼه بيان او
ثابته کړئ۔)
\end{prob}
\olpseganchor{OLP-0153-B028}{end}

\olpseganchor{OLP-0153-B029}{start}
\begin{prop}[يوازينے لوست]
هر فارمول له لاندې شرطونو څخه يو، او يوازې يو، پوره کوي۔
\begin{enumerate}
\item \LR{$!A$} اټومي دے۔
\olpseganchor{OLP-0153-B029}{end}

\olpseganchor{OLP-0153-B030}{start}
\item \LR{$!A$} د \LR{$\lnot !B$} بڼه لري۔
\olpseganchor{OLP-0153-B030}{end}

\olpseganchor{OLP-0153-B031}{start}
\item \LR{$!A$} د \LR{$(!B \land !C)$} بڼه لري۔
\olpseganchor{OLP-0153-B031}{end}

\olpseganchor{OLP-0153-B032}{start}
\item \LR{$!A$} د \LR{$(!B \lor !C)$} بڼه لري۔
\olpseganchor{OLP-0153-B032}{end}

\olpseganchor{OLP-0153-B033}{start}
\item \LR{$!A$} د \LR{$(!B \lif !C)$} بڼه لري۔
\olpseganchor{OLP-0153-B033}{end}

\olpseganchor{OLP-0153-B034}{start}
\item \LR{$!A$} د \LR{$(!B \liff !C)$} بڼه لري۔
\olpseganchor{OLP-0153-B034}{end}

\olpseganchor{OLP-0153-B035}{start}
\item \LR{$!A$} د \LR{$\lforall[x][!B]$} بڼه لري۔
\olpseganchor{OLP-0153-B035}{end}

\olpseganchor{OLP-0153-B036}{start}
\item \LR{$!A$} د \LR{$\lexists[x][!B]$} بڼه لري۔
\end{enumerate}
سربېره پر دې، په هر حالت کښې \LR{$!B$}، يا \LR{$!B$} او \LR{$!C$}، په يوازيني ډول
ټاکل شوي دي۔ يعنې مثلاً د \LR{$!B$}، \LR{$!C$} او \LR{$!B'$}، \LR{$!C'$} داسې بېلې جوړې
نشته چې \LR{$!A$} په يو وخت د دواړو بڼه ولري:
\LR{$(!B \lif !C)$} او \LR{$(!B' \lif !C')$}۔
\end{prop}
\olpseganchor{OLP-0153-B036}{end}

\olpseganchor{OLP-0153-B037}{start}
\begin{proof}
د جوړولو قاعدې غواړي چې که فارمول اټومي نۀ وي، بايد په
پرانيستي قوس~(، \LR{$\lnot$}، يا په يوه سور پيل شي۔ بل
خوا، هر فارمول چې له لاندې سمبولونو څخه په يوه پيلېږي، بايد
اټومي وي: پريديکات، تابع، ثابت سمبول
، \LR{$\lfalse$}، \LR{$\ltrue$}۔
\olpseganchor{OLP-0153-B037}{end}

\olpseganchor{OLP-0153-B038}{start}
نو په اصل کښې يوازې دا ښودل پکار دي چې که \LR{$!A$} د \LR{$(!B \ast !C)$}
او هم د \LR{$(!B' \mathbin{\ast'} !C')$} بڼه لري، نو \LR{$!B \ident !B'$}،
\LR{$!C \ident !C'$}، او \LR{$\ast = {\ast'}$}۔
\olpseganchor{OLP-0153-B038}{end}

\olpseganchor{OLP-0153-B039}{start}
نو فرض کړئ چې هم \LR{$!A \ident (!B \ast !C)$} او هم \LR{$!A \ident (!B'
\mathbin{\ast'} !C')$}۔ اوس يا \LR{$!B \ident !B'$} دے، يا نۀ دے۔ که وي،
نو څرګنده ده چې \LR{$\ast = {\ast'}$} او \LR{$!C \ident !C'$}، ځکه بيا دا د
\LR{$!A$} هغه فرعي تارونه دي چې له يوه ځايه پيلېږي او يو شان اوږدوالے لري۔
بل حالت \LR{$!B \not\ident !B'$} دے۔ څرنګه چې \LR{$!B$} او \LR{$!B'$} دواړه د
\LR{$!A$} هغه فرعي تارونه دي چې له يوه ځايه پيلېږي، يو بايد د بل حقيقي
مخکښ وي۔ خو دا د \olref{lem:no-prefix} له مخې ناشونې ده۔
\end{proof}
\olpseganchor{OLP-0153-B039}{end}

\olpseganchor{OLP-0154-B004}{start}
\olfileid{fol}{syn}{mai}
\olpseganchor{OLP-0154-B004}{end}

\olpseganchor{OLP-0154-B005}{start}
\olsection{د فارمول اصلي عملګر}
\olpseganchor{OLP-0154-B005}{end}

\olpseganchor{OLP-0154-B006}{start}
\begin{explain}
ډېر ځله د هغه وروستي عملګر په باب خبرې کول ګټور وي چې د
فارمول~\LR{$!A$} په جوړولو کښې کارول شوے وي۔ دا عملګر د~\LR{$!A$}
\emph{اصلي عملګر} بلل کېږي۔ په شهودي ډول دا د~\LR{$!A$} ``تر ټولو دباندې''
عملګر دے۔ مثلاً د \LR{$\lnot !A$} اصلي عملګر \LR{$\lnot$} دے، د \LR{$(!A \lor
!B)$} اصلي عملګر \LR{$\lor$} دے، او همداسې نور۔
\end{explain}
\olpseganchor{OLP-0154-B006}{end}


\olpseganchor{OLP-0154-B007}{start}
\begin{defn}[اصلي عملګر]
\ollabel{def:main-op}
د فارمول~\LR{$!A$} \emph{اصلي عملګر} په لاندې ډول تعريفېږي:
\begin{enumerate}
\item \indcase*{!A}{!A}{\LR{$\indfrm$} هېڅ اصلي عملګر نۀ لري۔}
\olpseganchor{OLP-0154-B007}{end}

\olpseganchor{OLP-0154-B008}{start}
\item \indcase{!A}{\lnot !B}{د \LR{$\indfrm$} اصلي عملګر
  سمبول~\LR{$\lnot$} دے۔}
\olpseganchor{OLP-0154-B008}{end}

\olpseganchor{OLP-0154-B009}{start}
\item \indcase{!A}{(!B \land !C)}{د \LR{$\indfrm$}
  اصلي عملګر سمبول~\LR{$\land$} دے۔}
\olpseganchor{OLP-0154-B009}{end}

\olpseganchor{OLP-0154-B010}{start}
\item \indcase{!A}{(!B \lor !C)}{د \LR{$\indfrm$}
  اصلي عملګر سمبول~\LR{$\lor$} دے۔}
\olpseganchor{OLP-0154-B010}{end}

\olpseganchor{OLP-0154-B011}{start}
\item \indcase{!A}{(!B \lif !C)}{د \LR{$\indfrm$}
  اصلي عملګر سمبول~\LR{$\lif$} دے۔}
\olpseganchor{OLP-0154-B011}{end}

\olpseganchor{OLP-0154-B012}{start}
\item \indcase{!A}{(!B \liff !C)}{د \LR{$\indfrm$}
  اصلي عملګر سمبول~\LR{$\liff$} دے۔}
\olpseganchor{OLP-0154-B012}{end}

\olpseganchor{OLP-0154-B013}{start}
\item \indcase{!A}{\lforall[x][!B]}{د \LR{$\indfrm$}
  اصلي عملګر سمبول~\LR{$\lforall$} دے۔}
\olpseganchor{OLP-0154-B013}{end}

\olpseganchor{OLP-0154-B014}{start}
\item \indcase{!A}{\lexists[x][!B]}{د \LR{$\indfrm$}
  اصلي عملګر سمبول~\LR{$\lexists$} دے۔}
\end{enumerate}
\end{defn}
\olpseganchor{OLP-0154-B014}{end}

\olpseganchor{OLP-0154-B015}{start}
په هر حالت کښې زموږ موخه په فارمول کښې د اصلي عملګر همدا ښودل
شوے ځانګړے \emph{مورد} دے۔ مثلاً فارمول \LR{$((!D \lif !E) \lif (!E
\lif !D))$} د \LR{$(!B \lif !C)$} بڼه لري، چې \LR{$!B$} پکښې \LR{$(!D \lif !E)$}
او \LR{$!C$} پکښې \LR{$(!E \lif !D)$} دے؛ نو د \LR{$\lif$} دويم مورد ئې
اصلي عملګر دے۔
\olpseganchor{OLP-0154-B015}{end}

\olpseganchor{OLP-0154-B016}{start}
\begin{explain}
دا د هغې تابع يو \emph{بازګشتي} تعريف دے چې غير اټومي
فارمولونه د هغو د اصلي عملګر موردونو ته لېږي۔ څرنګه چې
فارمولونه په استقرايي ډول تعريف شوي دي، هر فارمول~\LR{$!A$} په
\olref{def:main-op} کښې له ورکړل شوو حالتونو څخه يو پوره کوي۔ دا
تضمينوي چې د هر غير اټومي فارمول~\LR{$!A$} دپاره اصلي عملګر موجود دے۔ څرنګه چې هر فارمول له دغو شرطونو څخه
يوازې يو پوره کوي، او په هر حالت کښې هغه کوچني فارمولونه په
يوازيني ډول ټاکل شوي دي چې \LR{$!A$} ترې جوړ شوے، د~\LR{$!A$} د اصلي عملګر مورد هم يوازينے دے؛ نو په رښتيا مو يوه تابع تعريف کړې ده۔
\end{explain}
\olpseganchor{OLP-0154-B016}{end}

\olpseganchor{OLP-0154-B017}{start}
موږ فارمولونه په
\olref{tab:main-op} کښې د هغو د اصلي عملګر د سمبول له مخې په
ورکړل شوو نومونو يادوو۔

\olpseganchor{OLP-0154-B017}{end}

\olpseganchor{OLP-0154-B018}{start}
\begin{table}[!h]
\centering
\begin{tabular}{c | c | c}
اصلي عملګر & د فارمول ډول & بېلګه\\
\hline
هېڅ & اټومي (فارمول) &
\LR{$\lfalse$},
\LR{$\ltrue$},
\LR{$\Atom{R}{t_1, \dots, t_n}$},
\LR{$\eq[t_1][t_2]$}\\
\LR{$\lnot$} & نفي & \LR{$\lnot !A$} \\
\LR{$\land$} & عطف & \LR{$(!A \land !B)$} \\
\LR{$\lor$} & فصل & \LR{$(!A \lor !B)$} \\
\LR{$\lif$} & شرطي & \LR{$(!A \lif !B)$} \\
\LR{$\liff$} & دوه اړخيز شرطي & \LR{$(!A \liff !B)$} \\
\LR{$\lforall[][]$} & عمومي (فارمول)& \LR{$\lforall[x][!A]$} \\
\LR{$\lexists[][]$} & وجودي (فارمول)& \LR{$\lexists[x][!A]$}
\end{tabular}
\caption{اصلي عملګر او د فارمولونه نومونه}
\ollabel{tab:main-op}
\end{table}
\olpseganchor{OLP-0154-B018}{end}

\olpseganchor{OLP-0155-B004}{start}
\olfileid{fol}{syn}{sbf}
\olpseganchor{OLP-0155-B004}{end}

\olpseganchor{OLP-0155-B005}{start}
\olsection{فرعي فارمولونه}
\olpseganchor{OLP-0155-B005}{end}

\olpseganchor{OLP-0155-B006}{start}
\begin{explain}
ډېر ځله د هغو فارمولونه په باب خبرې کول ګټور وي چې يو راکړل شوے
فارمول ``جوړوي''۔ دې ته د هغه \emph{فرعي فارمولونه} وايو۔ هر
فارمول~\LR{$!A$} د خپل ځان فرعي فارمول هم ګڼل کېږي؛ د \LR{$!A$} هر هغه
فرعي فارمول چې پخپله \LR{$!A$} نۀ وي، \emph{حقيقي فرعي فارمول} دے۔
\end{explain}
\olpseganchor{OLP-0155-B006}{end}

\olpseganchor{OLP-0155-B007}{start}
\begin{defn}[سمدستي فرعي فارمول]
که \LR{$!A$} فارمول وي، د \LR{$!A$} \emph{سمدستي فرعي فارمولونه} په
استقرايي ډول داسې تعريفېږي:
\begin{enumerate}
\item اټومي فارمولونه سمدستي فرعي فارمولونه نۀ لري۔
\olpseganchor{OLP-0155-B007}{end}

\olpseganchor{OLP-0155-B008}{start}
\item \indcase{!A}{\lnot !B}{د \LR{$\indfrm$} يوازينے سمدستي
    فرعي فارمول همدا~\LR{$!B$} دے۔}
\olpseganchor{OLP-0155-B008}{end}

\olpseganchor{OLP-0155-B009}{start}
\item \indcase{!A}{(!B \ast !C)}{د \LR{$\indfrm$} سمدستي
  فرعي فارمولونه همدا \LR{$!B$} او \LR{$!C$} دي (\LR{$\ast$} له دوه ځايه
  نښلوونکو څخه هر يو کېدے شي)۔}
\olpseganchor{OLP-0155-B009}{end}

\olpseganchor{OLP-0155-B010}{start}
\item \indcase{!A}{\lforall[x][!B]}{د \LR{$\indfrm$} يوازينے
    سمدستي فرعي فارمول همدا~\LR{$!B$} دے۔}
\olpseganchor{OLP-0155-B010}{end}

\olpseganchor{OLP-0155-B011}{start}
\item \indcase{!A}{\lexists[x][!B]}{د \LR{$\indfrm$} يوازينے
    سمدستي فرعي فارمول همدا~\LR{$!B$} دے۔}
\end{enumerate}
\end{defn}
\olpseganchor{OLP-0155-B011}{end}

\olpseganchor{OLP-0155-B012}{start}
\begin{defn}[حقيقي فرعي فارمول]
که \LR{$!A$} فارمول وي، د \LR{$!A$} \emph{حقيقي فرعي فارمولونه} په
بازګشتي ډول داسې تعريفېږي:
\begin{enumerate}
\item اټومي فارمولونه حقيقي فرعي فارمولونه نۀ لري۔
\olpseganchor{OLP-0155-B012}{end}

\olpseganchor{OLP-0155-B013}{start}
\item \indcase{!A}{\lnot !B}{د \LR{$\indfrm$} حقيقي
    فرعي فارمولونه همدا~\LR{$!B$} او د~\LR{$!B$} ټول حقيقي فرعي فارمولونه دي۔}
\olpseganchor{OLP-0155-B013}{end}

\olpseganchor{OLP-0155-B014}{start}
\item \indcase{!A}{(!B \ast !C)}{د \LR{$\indfrm$} حقيقي
  فرعي فارمولونه همدا \LR{$!B$}، \LR{$!C$}، د \LR{$!B$} ټول حقيقي فرعي فارمولونه
  او د~\LR{$!C$} ټول حقيقي فرعي فارمولونه دي۔}
\olpseganchor{OLP-0155-B014}{end}

\olpseganchor{OLP-0155-B015}{start}
\item \indcase{!A}{\lforall[x][!B]}{د \LR{$\indfrm$} حقيقي
    فرعي فارمولونه همدا~\LR{$!B$} او د~\LR{$!B$} ټول حقيقي
    فرعي فارمولونه دي۔}
\olpseganchor{OLP-0155-B015}{end}

\olpseganchor{OLP-0155-B016}{start}
\item \indcase{!A}{\lexists[x][!B]}{د \LR{$\indfrm$} حقيقي
    فرعي فارمولونه همدا~\LR{$!B$} او د~\LR{$!B$} ټول حقيقي
    فرعي فارمولونه دي۔}
\end{enumerate}
\end{defn}
\olpseganchor{OLP-0155-B016}{end}

\olpseganchor{OLP-0155-B017}{start}
\begin{defn}[فرعي فارمول]
د \LR{$!A$} فرعي فارمولونه پخپله \LR{$!A$} او د هغه ټول حقيقي
فرعي فارمولونه دي۔
\end{defn}
\olpseganchor{OLP-0155-B017}{end}

\olpseganchor{OLP-0155-B018}{start}
\begin{explain}
پام وکړئ چې سمدستي فرعي فارمولونه او حقيقي فرعي فارمولونه مو په
لږ توپير تعريف کړي دي۔ په لومړي حالت کښې مو د~\LR{$!A$} د هرې ممکنې بڼې
دپاره د هغه سمدستي فرعي فارمولونه نېغ په نېغه تعريف کړي دي۔ دا د
حالتونو له مخې يو څرګند تعريف دے، او حالتونه ئې د فارمولونه د سټ
له استقرايي تعريف سره سمون لري۔ په دويم حالت کښې مو هم د ټولو
فارمولونه د سټ د تعريف طريقه تعقيب کړې ده، خو په هر حالت کښې مو د
کوچنيو فارمولونه \LR{$!B$} او \LR{$!C$} د حقيقي فرعي فارمولونه تر څنګ،
پخپله دا فارمولونه هم ور شامل کړي دي۔ په دې سره تعريف \emph{بازګشتي}
کېږي۔ په عمومي ډول، پر يوه استقرايي تعريف شوي سټ (زموږ په حالت کښې
فارمولونه) د تابع تعريف هله بازګشتي وي چې د تابع د تعريف حالتونه
پخپله هماغه تابع وکاروي۔ خو د ښه تعريف دپاره بايد ډاډه شو چې د تابع
قيمتونه تل يوازې د هغو مدخلو دپاره کاروو چې تر تعريفېدونکې مدخل
``مخکښې'' راځي---زموږ په حالت کښې، د \LR{$(!B \ast !C)$} د ``حقيقي
فرعي فارمول'' د تعريف پر وخت يوازې د ``مخکښو'' فارمولونه \LR{$!B$}
او \LR{$!C$} حقيقي فرعي فارمولونه کاروو۔
\end{explain}
\olpseganchor{OLP-0155-B018}{end}

\olpseganchor{OLP-0155-B019}{start}
\begin{prop}
\ollabel{prop:subfrm-trans}
فرض کړئ \LR{$!B$} د \LR{$!A$} يو فرعي فارمول او \LR{$!C$} د \LR{$!B$} يو فرعي فارمول
دے۔ بيا \LR{$!C$} د \LR{$!A$} يو فرعي فارمول دے۔ په بله وينا، د فرعي فارمول
اړيکه انتقالي ده۔
\end{prop}
\olpseganchor{OLP-0155-B019}{end}

\olpseganchor{OLP-0155-B020}{start}
\begin{prob}
\olref[fol][syn][sbf]{prop:subfrm-trans} ثابته کړئ۔
\end{prob}
\olpseganchor{OLP-0155-B020}{end}

\olpseganchor{OLP-0155-B021}{start}
\begin{prop}
\ollabel{prop:count-subfrms}
فرض کړئ \LR{$!A$} داسې فارمول دے چې \LR{$n$} نښلوونکي او سورونه لري۔ بيا \LR{$!A$}
تر ډېره \LR{$2n+1$} فرعي فارمولونه لري۔
\end{prop}
\olpseganchor{OLP-0155-B021}{end}

\olpseganchor{OLP-0155-B022}{start}
\begin{prob}
\olref[fol][syn][sbf]{prop:count-subfrms} ثابته کړئ۔
\end{prob}
\olpseganchor{OLP-0155-B022}{end}

\olpseganchor{OLP-0156-B004}{start}
\olfileid{fol}{syn}{fseq}
\olpseganchor{OLP-0156-B004}{end}

\olpseganchor{OLP-0156-B005}{start}
\olsection{جوړښتي لړۍ}
\olpseganchor{OLP-0156-B005}{end}

\olpseganchor{OLP-0156-B006}{start}
په استقرايي تعريف سره د فارمولونه تعريفول، او د استقرا له لارې د
فارمولونه د خاصيتونو د ثابتولو بشپړوونکې طريقه، ښکلې او اغېزمنه لار
ده۔ خو کله کله دا هم ګټوره وي چې د فارمولونه جوړول له بېخه پورته،
ګام په ګام وڅېړو۔ دلته د \emph{جوړښتي لړۍ} په مفهوم همدا کار کوو۔
%
د دې د ښودلو دپاره چې ترمونه او فارمولونه څنګه له خپلو استقرايي
تعريفونو پرته په دې طريقه معرفي کېدے شي، لومړے د يوې ژبې~\LR{$\Lang L$}
له سمبولونو څخه د جوړ شوي خپل سري تار مفهوم معرفي کوو۔
\olpseganchor{OLP-0156-B006}{end}

\olpseganchor{OLP-0156-B007}{start}
\begin{defn}[تارونه]
\ollabel{defn:string}
فرض کړئ \LR{$\Lang L$} د لومړۍ درجې يوه ژبه ده۔ \emph{د~\LR{$\Lang L$} تار}
د~\LR{$\Lang L$} د سمبولونو يوه متناهي لړۍ ده۔ چې ژبه~\LR{$\Lang L$} له
چاپېريال څخه په څرګند ډول ټاکلې وي، د~\LR{$\Lang L$} تار ته به ډېر ځله
يوازې \emph{تار} وايو۔
\end{defn}
\olpseganchor{OLP-0156-B007}{end}

\olpseganchor{OLP-0156-B008}{start}
\begin{ex}
د لومړۍ درجې د هرې ژبې \LR{$\Lang L$} دپاره ټول
\LR{$\Lang L$}-فارمولونه د~\LR{$\Lang L$} تارونه دي، خو عکس ئې سم نۀ دے۔
مثلاً \[)(\Obj v_0\lif\lexists\] د~\LR{$\Lang L$} يو تار دے، خو د~\LR{$\Lang L$}
فارمول نۀ دے۔
\end{ex}
\olpseganchor{OLP-0156-B008}{end}

\olpseganchor{OLP-0156-B009}{start}
\begin{defn}[د ترمونو جوړښتي لړۍ]
\ollabel{defn:fseq-trm}
د~\LR{$\Lang L$} د تارونو متناهي لړۍ \LR{$\tuple{t_0,\dotsc,t_n}$} د يوه
ترم \LR{$t$} \emph{جوړښتي لړۍ} ده، که \LR{$t \ident t_n$} وي او د هر \LR{$i \leq
n$} دپاره يا \LR{$t_i$} متغير يا ثابت سمبول وي، او يا~\LR{$\Lang L$}
يو \LR{$k$}-ځايه تابع~\LR{$f$} ولري او داسې \LR{$m_1,\dotsc,m_k < i$}
موجود وي چې \LR{$t_i \ident f(t_{m_1},\dotsc,t_{m_k})$}۔ چې توپير ئې اړين
وي، د ترمونو جوړښتي لړۍ به \emph{د ترم جوړښتي لړۍ} بولو۔
\end{defn}
\olpseganchor{OLP-0156-B009}{end}

\olpseganchor{OLP-0156-B010}{start}
\begin{ex}
لړۍ
\[
    \tuple{\Obj c_0, \Obj v_0, \Atom{\Obj f^2_0}{\Obj c_0, \Obj v_0}, \Atom{\Obj f^1_0}{\Atom{\Obj f^2_0}{\Obj c_0, \Obj v_0}}}
\]
د ترم \LR{$\Atom{\Obj f^1_0}{\Atom{\Obj f^2_0}{\Obj c_0, \Obj v_0}}$}
يوه جوړښتي لړۍ ده، او دا هم ده:
\[
    \tuple{\Obj v_0, \Obj c_0, \Atom{\Obj f^2_0}{\Obj c_0, \Obj v_0}, \Atom{\Obj f^1_0}{\Atom{\Obj f^2_0}{\Obj c_0, \Obj v_0}}}.
\]
\end{ex}
\olpseganchor{OLP-0156-B010}{end}

\olpseganchor{OLP-0156-B011}{start}
\begin{defn}[د فارمولونو جوړښتي لړۍ]
\ollabel{defn:fseq-frm}
د~\LR{$\Lang L$} د تارونو متناهي لړۍ \LR{$\tuple{!A_0,\dotsc,!A_n}$} د~\LR{$!A$}
\emph{جوړښتي لړۍ} ده، که \LR{$!A \ident !A_n$} وي او د هر \LR{$i \leq n$}
دپاره يا \LR{$!A_i$} يو اټومي فارمول وي، يا داسې \LR{$j,k < i$} او
متغير~\LR{$x$} موجود وي چې له لاندې حالتونو څخه يو صدق وکړي:
\begin{enumerate}
    \item \LR{$!A_i \ident \lnot !A_j$}۔%
    \item \LR{$!A_i \ident (!A_j \land !A_k)$}۔%
    \item \LR{$!A_i \ident (!A_j \lor !A_k)$}۔%
    \item \LR{$!A_i \ident (!A_j \lif !A_k)$}۔%
    \item \LR{$!A_i \ident (!A_j \liff !A_k)$}۔%
    \item \LR{$!A_i \ident \lforall[x][!A_j]$}۔%
    \item \LR{$!A_i \ident \lexists[x][!A_j]$}۔%
\end{enumerate}
چې توپير ئې اړين وي، د فارمولونو جوړښتي لړۍ به \emph{د فارمول
جوړښتي لړۍ} بولو۔
\end{defn}
\olpseganchor{OLP-0156-B011}{end}

\olpseganchor{OLP-0156-B012}{start}
\begin{ex}
\[
\tuple{
    \Atom{\Obj A^1_0}{\Obj v_0},
    \Atom{\Obj A^1_1}{\Obj c_1},
    (\Atom{\Obj A^1_1}{\Obj c_1} \land \Atom{\Obj A^1_0}{\Obj v_0}),
    \lexists[\Obj v_0][(\Atom{\Obj A^1_1}{\Obj c_1} \land \Atom{\Obj A^1_0}{\Obj v_0})]
}
\]
د \LR{$\lexists[\Obj v_0][(\Atom{\Obj A^1_1}{\Obj c_1} \land \Atom{\Obj
A^1_0}{\Obj v_0})]$} يوه جوړښتي لړۍ ده، او دا هم ده:
\begin{multline*}
\tuple{
    \Atom{\Obj A^1_0}{\Obj v_0},
    \Atom{\Obj A^1_1}{\Obj c_1},
    (\Atom{\Obj A^1_1}{\Obj c_1} \land \Atom{\Obj A^1_0}{\Obj v_0}),
    \Atom{\Obj A^1_1}{\Obj c_1},\\
    \lforall[\Obj v_1][\Atom{\Obj A^1_0}{\Obj v_0}],
    \lexists[\Obj v_0][(\Atom{\Obj A^1_1}{\Obj c_1} \land \Atom{\Obj A^1_0}{\Obj v_0})]
}.
\end{multline*}
%
لکه چې له دويمې بېلګې ښکاري، په جوړښتي لړيو کښې ``بې‌ګټې توکي''
هم راتلے شي: داسې فارمولونه چې زياتي وي يا په جوړونه کښې ونډه نۀ
لري۔
\end{ex}
\olpseganchor{OLP-0156-B012}{end}

\olpseganchor{OLP-0156-B013}{start}
\begin{prop}\ollabel{prop:formed}
په~\LR{$\Frm[L]$} کښې هر فارمول~\LR{$!A$} يوه جوړښتي لړۍ لري۔
\end{prop}
\olpseganchor{OLP-0156-B013}{end}

\olpseganchor{OLP-0156-B014}{start}
\begin{proof}
فرض کړئ \LR{$!A$} اټومي دے۔ نو لړۍ~\LR{$\tuple{!A}$} د~\LR{$!A$} جوړښتي لړۍ ده۔
%
اوس فرض کړئ چې \LR{$!B$} او~\LR{$!C$} په ترتيب سره جوړښتي لړۍ
\LR{$\tuple{!B_0,\dotsc,!B_n}$} او \LR{$\tuple{!C_0,\dotsc,!C_m}$} لري۔
%
\begin{enumerate}
  \item که \LR{$!A \ident \lnot !B$} وي، نو
      \LR{$\tuple{!B_0,\dotsc,!B_n,\lnot !B_n}$}
      د~\LR{$!A$} جوړښتي لړۍ ده۔
  \item که \LR{$!A \ident (!B \land !C)$} وي، نو
      \LR{$\tuple{!B_0,\dotsc,!B_n,!C_0,\dotsc,!C_m,(!B_n \land !C_m)}$}
      د~\LR{$!A$} جوړښتي لړۍ ده۔
  \item که \LR{$!A \ident (!B \lor !C)$} وي، نو
      \LR{$\tuple{!B_0,\dotsc,!B_n,!C_0,\dotsc,!C_m,(!B_n \lor !C_m)}$}
      د~\LR{$!A$} جوړښتي لړۍ ده۔
  \item که \LR{$!A \ident (!B \lif !C)$} وي، نو
      \LR{$\tuple{!B_0,\dotsc,!B_n,!C_0,\dotsc,!C_m,(!B_n \lif !C_m)}$}
      د~\LR{$!A$} جوړښتي لړۍ ده۔
  \item که \LR{$!A \ident (!B \liff !C)$} وي، نو
      \LR{$\tuple{!B_0,\dotsc,!B_n,!C_0,\dotsc,!C_m,(!B_n \liff !C_m)}$}
      د~\LR{$!A$} جوړښتي لړۍ ده۔
  \item که \LR{$!A \ident \lforall[x][!B]$} وي، نو
      \LR{$\tuple{!B_0,\dotsc,!B_n,\lforall[x][!B_n]}$}
      د~\LR{$!A$} جوړښتي لړۍ ده۔
  \item که \LR{$!A \ident \lexists[x][!B]$} وي، نو
      \LR{$\tuple{!B_0,\dotsc,!B_n,\lexists[x][!B_n]}$}
      د~\LR{$!A$} جوړښتي لړۍ ده۔
\end{enumerate}
پر فارمولونه د استقرا د اصل له مخې، هر فارمول يوه جوړښتي لړۍ
لري۔
\end{proof}
\olpseganchor{OLP-0156-B014}{end}

\olpseganchor{OLP-0156-B015}{start}
د دې عکس هم ثابتولے شو۔ دا ځکه مهم دے چې ښيي د فارمولونو د تعريف
زموږ دواړه لارې هم ارزښته دي: دواړه يو شان پايلې ورکوي۔ دا هم ترې
معلومېږي چې د جوړښتي لړيو پر اوږدوالي عادي استقرا کارولے او د
فارمولونو په باب قضيې ثابتولے شو۔
\olpseganchor{OLP-0156-B015}{end}

\olpseganchor{OLP-0156-B016}{start}
\begin{lem}
\ollabel{lem:fseq-init}
فرض کړئ \LR{$\tuple{!A_0,\dotsc,!A_n}$} د~\LR{$!A_n$} جوړښتي لړۍ ده او \LR{$k
\leq n$} دے۔ نو \LR{$\tuple{!A_0,\dotsc,!A_k}$} د~\LR{$!A_k$} جوړښتي لړۍ ده۔
\end{lem}
\olpseganchor{OLP-0156-B016}{end}

\olpseganchor{OLP-0156-B017}{start}
\begin{proof}
تمرين۔
\end{proof}
\olpseganchor{OLP-0156-B017}{end}

\olpseganchor{OLP-0156-B018}{start}
\begin{prob}
\olref[fol][syn][fseq]{lem:fseq-init} ثابته کړئ۔
\end{prob}
\olpseganchor{OLP-0156-B018}{end}

\olpseganchor{OLP-0156-B019}{start}
\begin{thm}
\ollabel{thm:fseq-frm-equiv}
\LR{$\Frm[L]$} د~\LR{$\Lang L$} د ټولو هغو تارونو~\LR{$!A$} سټ دے چې د~\LR{$!A$} دپاره
د فارمول جوړښتي لړۍ موجوده وي۔
\end{thm}
\olpseganchor{OLP-0156-B019}{end}

\olpseganchor{OLP-0156-B020}{start}
\begin{proof}
\LR{$F$} دې د ژبې~\LR{$\Lang L$} د سمبولونو د ټولو هغو تارونو سټ وي چې جوړښتي
لړۍ لري۔ په
\olref[fol][syn][fseq]{prop:formed} کښې مو وليدل چې \LR{$\Frm[L] \subseteq
F$}؛ نو اوس ئې عکس ثابتوو۔
\olpseganchor{OLP-0156-B020}{end}

\olpseganchor{OLP-0156-B021}{start}
فرض کړئ \LR{$!A$} يوه جوړښتي لړۍ \LR{$\tuple{!A_0,\dotsc,!A_n}$} لري۔ پر~\LR{$n$}
د قوي استقرا له لارې ثابتوو چې \LR{$!A \in \Frm[L]$}۔ زموږ استقرايي فرض
دا دے چې د سمبولونو هر هغه تار چې د وروستي شاخص \LR{$m < n$} لرونکې
جوړښتي لړۍ لري، په~\LR{$\Frm[L]$} کښې دے۔ د جوړښتي لړۍ د تعريف له مخې، يا
\LR{$!A \ident !A_n$} وي او وروستے تار اټومي وي، يا بايد داسې \LR{$j,k < n$}
موجود وي چې له لاندې حالتونو څخه يو صدق وکړي:
\begin{enumerate}
    \item \LR{$!A \ident \lnot !A_j$}۔%
    \item \LR{$!A \ident (!A_j \land !A_k)$}۔%
    \item \LR{$!A \ident (!A_j \lor !A_k)$}۔%
    \item \LR{$!A \ident (!A_j \lif !A_k)$}۔%
    \item \LR{$!A \ident (!A_j \liff !A_k)$}۔%
    \item \LR{$!A \ident \lforall[x][!A_j]$}۔%
    \item \LR{$!A \ident \lexists[x][!A_j]$}۔%
\end{enumerate}
اوس حالت په حالت استدلال کوو۔ که \LR{$!A$} اټومي وي، نو
\LR{$!A_n \in \Frm[L]$}۔ بل حالت کښې فرض کړئ چې
\LR{$!A \ident (!A_j \land !A_k)$}۔ د
\olref[fol][syn][fseq]{lem:fseq-init} له مخې،
\LR{$\tuple{!A_0,\dotsc,!A_j}$} او \LR{$\tuple{!A_0,\dotsc,!A_k}$} په ترتيب
سره د \LR{$!A_j$} او~\LR{$!A_k$} جوړښتي لړۍ دي۔ څرنګه چې دا د~\LR{$!A$} د جوړښتي
لړۍ حقيقي پيلنۍ فرعي لړۍ دي، د دواړو وروستے شاخص له~\LR{$n$} څخه کم دے۔
نو د استقرايي فرض له مخې \LR{$!A_j$} او~\LR{$!A_k$} په~\LR{$\Frm[L]$} کښې دي؛ او
بيا د فارمول د تعريف له مخې \LR{$(!A_j \land !A_k)$} هم پکې دے۔ نور
حالتونه په ورته استدلال ثابتېږي۔
\end{proof}
\olpseganchor{OLP-0156-B021}{end}

\olpseganchor{OLP-0156-B022}{start}
د ترمونو جوړښتي لړۍ د فارمولونه د جوړښتي لړيو په شان خاصيتونه لري۔
\olpseganchor{OLP-0156-B022}{end}

\olpseganchor{OLP-0156-B023}{start}
\begin{prop}
\ollabel{prop:fseq-trm-equiv}
\LR{$\Trm[L]$} د~\LR{$\Lang L$} د ټولو هغو تارونو \LR{$t$} سټ دے چې د~\LR{$t$} دپاره
د ترم جوړښتي لړۍ موجوده وي۔
\end{prop}
\olpseganchor{OLP-0156-B023}{end}

\olpseganchor{OLP-0156-B024}{start}
\begin{proof}
تمرين۔
\end{proof}
\olpseganchor{OLP-0156-B024}{end}

\olpseganchor{OLP-0156-B025}{start}
\begin{prob}
\olref[fol][syn][fseq]{prop:fseq-trm-equiv} ثابته کړئ۔
لارښوونه: د
\olref[fol][syn][fseq]{thm:fseq-frm-equiv} په ثبوت کښې کارول شوې
طريقې ته ورته تګلاره وکاروئ۔
\end{prob}
\olpseganchor{OLP-0156-B025}{end}

\olpseganchor{OLP-0156-B026}{start}
په جوړښتي لړيو کښې دوه ډوله ``بې‌ګټې توکي'' راتلے شي: تکراري غړي، او
هغه غړي چې د فارمول يا ترم په جوړولو کښې ونډه نۀ لري۔ دواړه د تر ټولو
لنډو جوړښتي لړيو په کتلو لرې کولے شو۔
\olpseganchor{OLP-0156-B026}{end}

\olpseganchor{OLP-0156-B027}{start}
\begin{defn}[تر ټولو لنډې جوړښتي لړۍ]
\ollabel{defn:minimal-fseq}
د فارمول~\LR{$!A$} جوړښتي لړۍ \LR{$\tuple{!A_0, \ldots, !A_n}$} د~\LR{$!A$}
\emph{تر ټولو لنډه جوړښتي لړۍ} ده، که د~\LR{$!A$} هره بله جوړښتي
لړۍ~\LR{$s$} داسې وي چې د~\LR{$s$} اوږدوالے له~\LR{$n+1$} سره برابر يا ترې زيات وي۔
\olpseganchor{OLP-0156-B027}{end}

\olpseganchor{OLP-0156-B028}{start}
په همدې ډول، د ترم~\LR{$t$} جوړښتي لړۍ \LR{$\tuple{t_0, \ldots, t_n}$} د~\LR{$t$}
\emph{تر ټولو لنډه جوړښتي لړۍ} ده، که د~\LR{$t$} هره بله جوړښتي
لړۍ~\LR{$s$} داسې وي چې د~\LR{$s$} اوږدوالے له~\LR{$n+1$} سره برابر يا ترې زيات وي۔
\end{defn}
\olpseganchor{OLP-0156-B028}{end}

\olpseganchor{OLP-0156-B029}{start}
پام وکړئ چې يو فارمول يا ترم له يوې څخه زياتې تر ټولو لنډې جوړښتي لړۍ
لرلے شي، خو په ټولو کښې به کټ مټ يو شان تارونه وي۔
\olpseganchor{OLP-0156-B029}{end}

\olpseganchor{OLP-0156-B030}{start}
\begin{prop}
\ollabel{prop:subformula-equivs}
لاندې شرطونه هم ارزښته دي:
\begin{enumerate}
    \item \LR{$!B$} د~\LR{$!A$} يو فرعي-فارمول دے۔
    \item \LR{$!B$} د~\LR{$!A$} په هره جوړښتي لړۍ کښې راځي۔
    \item \LR{$!B$} د~\LR{$!A$} په کومه تر ټولو لنډه جوړښتي لړۍ کښې راځي۔
\end{enumerate}
\end{prop}
\olpseganchor{OLP-0156-B030}{end}

\olpseganchor{OLP-0156-B031}{start}
\begin{proof}
تمرين۔
\end{proof}
\olpseganchor{OLP-0156-B031}{end}

\olpseganchor{OLP-0156-B032}{start}
\begin{prob}
\olref[fol][syn][fseq]{prop:subformula-equivs} ثابته کړئ۔
\end{prob}
\olpseganchor{OLP-0156-B032}{end}

\olpseganchor{OLP-0156-B033}{start}
\begin{history}
جوړښتي لړۍ ريمونډ سمولين په خپل درسي کتاب
\emph{First-Order Logic} \mbox{\LR{\latinfont \citep{Smullyan1968}}} کښې معرفي کړې۔
د جوړښتي لړيو نور خاصيتونه \mbox{\LR{\latinfont \citet{Zuckerman1973}}} ثابت کړل۔
\end{history}
\olpseganchor{OLP-0156-B033}{end}

\olpseganchor{OLP-0157-B004}{start}
\olfileid{fol}{syn}{fvs}
\olpseganchor{OLP-0157-B004}{end}

\olpseganchor{OLP-0157-B005}{start}
\olsection{ازاد متغيرونه او تړلي فارمولونه}
\olpseganchor{OLP-0157-B005}{end}

\olpseganchor{OLP-0157-B006}{start}
\begin{defn}[د متغير ازاد موردونه]
\ollabel{defn:free-occ}
په فارمول کښې د متغير \emph{ازاد} موردونه په استقرايي
ډول داسې تعريفېږي:
\begin{enumerate}
\item \indcase*{!A}{\LR{$!A$} اټومي دے}{په \LR{$\indfrm$} کښې د
  متغير ټول موردونه ازاد دي۔}
\olpseganchor{OLP-0157-B006}{end}

\olpseganchor{OLP-0157-B007}{start}
\item \indcase{!A}{\lnot !B}{د \LR{$\indfrm$} د متغير
  ازاد موردونه کټ مټ د \LR{$!B$} ازاد موردونه دي۔}
\olpseganchor{OLP-0157-B007}{end}

\olpseganchor{OLP-0157-B008}{start}
\item \indcase{!A}{(!B \ast !C)}{د \LR{$\indfrm$} د متغير ازاد
  موردونه په \LR{$!B$} کښې ازاد موردونه او ورسره په~\LR{$!C$} کښې ازاد موردونه
  دي۔}
\olpseganchor{OLP-0157-B008}{end}

\olpseganchor{OLP-0157-B009}{start}
\item \indcase{!A}{\lforall[x][!B]}{په \LR{$\indfrm$} کښې د
  متغير ازاد موردونه په~\LR{$!B$} کښې ټول ازاد موردونه دي، د~\LR{$x$}
  له موردونو پرته۔}
\olpseganchor{OLP-0157-B009}{end}

\olpseganchor{OLP-0157-B010}{start}
\item \indcase{!A}{\lexists[x][!B]}{په \LR{$\indfrm$} کښې د
  متغير ازاد موردونه په~\LR{$!B$} کښې ټول ازاد موردونه دي، د~\LR{$x$}
  له موردونو پرته۔}
\end{enumerate}
\end{defn}
\olpseganchor{OLP-0157-B010}{end}

\olpseganchor{OLP-0157-B011}{start}
\begin{defn}[تړلي متغيرونه]
په يوه فارمول~\LR{$!A$} کښې د متغير يو مورد هله \emph{تړلے} دے
چې ازاد نۀ وي۔
\end{defn}
\olpseganchor{OLP-0157-B011}{end}

\olpseganchor{OLP-0157-B012}{start}
\begin{prob}
په \olref[fol][syn][fvs]{defn:free-occ} پسې د تړلو متغيرونو د موردونو
يو استقرايي تعريف ورکړئ۔
\end{prob}
\olpseganchor{OLP-0157-B012}{end}

\olpseganchor{OLP-0157-B013}{start}
\begin{defn}[ساحه]
که \LR{$\lforall[x][!B]$} په يوه فارمول~\LR{$!A$} کښې د يوه
  فرعي فارمول مورد وي، نو په~\LR{$!A$} کښې د~\LR{$!B$} اړوند مورد د
  \LR{$\lforall[x]$} د اړوند مورد \emph{ساحه} بلل کېږي۔
  د \LR{$\lexists[x]$} دپاره هم همدا شان۔
\olpseganchor{OLP-0157-B013}{end}

\olpseganchor{OLP-0157-B014}{start}
که \LR{$!B$} په~\LR{$!A$} کښې د سور د يوه مورد
\LR{$\lforall[x]$} يا
    \LR{$\lexists[x]$} ساحه وي، نو په~\LR{$!B$} کښې د
\LR{$x$} ازاد موردونه په \LR{$\lforall[x][!B]$} او
\LR{$\lexists[x][!B]$} کښې تړلي دي۔ وايو چې دا
موردونه د ياد شوي سور د مورد \emph{له خوا تړل شوي} دي۔
\end{defn}
\olpseganchor{OLP-0157-B014}{end}

\olpseganchor{OLP-0157-B015}{start}
\begin{ex}
لاندې فارمول په پام کښې ونيسئ:
\[
\lexists[\Obj v_0][\underbrace{\Atom{\Obj A^2_0}{\Obj v_0,\Obj v_1}}_{!B}]
\]
\LR{$!B$} د \LR{$\lexists[\Obj v_0]$} ساحه ده۔ سور په \LR{$!B$} کښې د \LR{$\Obj v_0$}
مورد تړي، خو د \LR{$\Obj v_1$} مورد نۀ تړي۔ نو په دې حالت کښې \LR{$\Obj v_1$}
ازاد متغير دے۔
\olpseganchor{OLP-0157-B015}{end}


\olpseganchor{OLP-0157-B016}{start}
اوس وينو چې په يوه زيات پېچلي فارمول~\LR{$!A$} کښې دا څنګه کار کوي:
\[
\lforall[\Obj v_0][\underbrace{(\Atom{\Obj A^1_0}{\Obj v_0} \lif
    \Atom{\Obj A^2_0}{\Obj v_0, \Obj v_1})}_{!B}] \lif \lexists[\Obj
  v_1][\underbrace{(\Atom{\Obj A^2_1}{\Obj v_0, \Obj v_1} \lor \lforall[\Obj v_0][\overbrace{\lnot \Atom{\Obj A^1_1}{\Obj v_0}}^{!D}])}_{!C}]
\]
\LR{$!B$} د لومړي \LR{$\lforall[\Obj v_0]$} ساحه، \LR{$!C$} د
\LR{$\lexists[\Obj v_1]$} ساحه، او \LR{$!D$} د دويم \LR{$\lforall[\Obj v_0]$} ساحه
ده۔ لومړے \LR{$\lforall[\Obj v_0]$} په~\LR{$!B$} کښې د \LR{$\Obj v_0$} موردونه
تړي، \LR{$\lexists[\Obj v_1]$} په \LR{$!C$} کښې د \LR{$\Obj v_1$} مورد تړي، او دويم
\LR{$\lforall[\Obj v_0]$} په~\LR{$!D$} کښې د \LR{$\Obj v_0$} مورد تړي۔ د \LR{$\Obj v_1$}
لومړے مورد او د \LR{$\Obj v_0$} څلورم مورد په~\LR{$!A$} کښې ازاد دي۔ د
\LR{$\Obj v_0$} وروستے مورد په \LR{$!D$} کښې ازاد، خو په \LR{$!C$} او~\LR{$!A$} کښې
تړلے دے۔
\end{ex}
\olpseganchor{OLP-0157-B016}{end}

\olpseganchor{OLP-0157-B017}{start}
\begin{defn}[جمله]
فارمول~\LR{$!A$} هله او يوازې هله يوه 
\emph{تړلے فارمول} ده چې د متغيرونه هېڅ ازاد مورد ونۀ لري۔
\end{defn}
\olpseganchor{OLP-0157-B017}{end}

\olpseganchor{OLP-0157-B018}{start}
% add examples!
\olpseganchor{OLP-0157-B018}{end}

\olpseganchor{OLP-0158-B004}{start}
\olfileid{fol}{syn}{sub}
\olpseganchor{OLP-0158-B004}{end}

\olpseganchor{OLP-0158-B005}{start}
\olsection{ځايناستي}
\olpseganchor{OLP-0158-B005}{end}

\olpseganchor{OLP-0158-B006}{start}
\begin{defn}[په ترم کښې ځايناستي]
\LR{$\Subst{s}{t}{x}$}، يعنې په \LR{$s$} کښې د~\LR{$x$} د هر مورد پر ځاے د \LR{$t$}
د \emph{ځايناستي} پايله، په بازګشتي ډول داسې تعريفوو:
\begin{enumerate}
\item \indcase{s}{c}{\LR{$\Subst{\indfrm}{t}{x}$} پخپله \LR{$s$} دے۔}
\olpseganchor{OLP-0158-B006}{end}

\olpseganchor{OLP-0158-B007}{start}
\item \indcase{s}{y}{که \LR{$y$} يو متغير او \LR{$y \not\ident x$} وي، نو
  \LR{$\Subst{\indfrm}{t}{x}$} هم پخپله~\LR{$s$} دے۔}
\olpseganchor{OLP-0158-B007}{end}

\olpseganchor{OLP-0158-B008}{start}
\item \indcase{s}{x}{\LR{$\Subst{\indfrm}{t}{x}$} همدا~\LR{$t$} دے۔}
\olpseganchor{OLP-0158-B008}{end}

\olpseganchor{OLP-0158-B009}{start}
\item\indcase{s}{\Atom{f}{t_1, \dots, t_n}}{\LR{$\Subst{\indfrmp}{t}{x}$}
  همدا \LR{$\Atom{f}{\Subst{t_1}{t}{x}, \dots, \Subst{t_n}{t}{x}}$} دے۔}
\end{enumerate}
\end{defn}
\olpseganchor{OLP-0158-B009}{end}

\olpseganchor{OLP-0158-B010}{start}
\begin{defn}
يو ترم~\LR{$t$} په \LR{$!A$} کښې د \LR{$x$} دپاره هله \emph{ازاد} دے چې
په \LR{$!A$} کښې د~\LR{$x$} هېڅ ازاد مورد د داسې سور په ساحه کښې نۀ راځي چې
په~\LR{$t$} کښې کوم متغير تړي۔
\end{defn}
\olpseganchor{OLP-0158-B010}{end}

\olpseganchor{OLP-0158-B011}{start}
\begin{ex} ~
\begin{enumerate}
\item \LR{$\Obj v_8$} په \LR{$\lexists[\Obj
  v_3]\Atom{\Obj A^2_4}{\Obj v_3,\Obj v_1}$} کښې د \LR{$\Obj v_1$} دپاره
  ازاد دے۔
\olpseganchor{OLP-0158-B011}{end}

\olpseganchor{OLP-0158-B012}{start}
\item \LR{$\Obj f^2_1(\Obj v_1, \Obj v_2)$} په
  \LR{$\lforall[\Obj v_2]\Atom{\Obj A^2_4}{\Obj v_0,\Obj v_2}$} کښې د
  \LR{$\Obj v_0$} دپاره \emph{ازاد نۀ} دے۔
\end{enumerate}
\end{ex}
\olpseganchor{OLP-0158-B012}{end}

\olpseganchor{OLP-0158-B013}{start}
\begin{defn}[په فارمول کښې ځايناستي]
که \LR{$!A$} فارمول، \LR{$x$}~متغير، او \LR{$t$}~داسې ترم وي چې
په~\LR{$!A$} کښې د~\LR{$x$} دپاره ازاد وي، نو
\LR{$\Subst{!A}{t}{x}$} په~\LR{$!A$} کښې د~\LR{$x$} د ټولو ازادو موردونو پر ځاے
د~\LR{$t$} د ځايناستي پايله ده۔
\begin{enumerate}
\item \indcase{!A}{\lfalse}{\LR{$\Subst{\indfrm}{t}{x}$}
    همدا \LR{$\lfalse$} دے۔}
\olpseganchor{OLP-0158-B013}{end}

\olpseganchor{OLP-0158-B014}{start}
\item \indcase{!A}{\ltrue}{\LR{$\Subst{\indfrm}{t}{x}$}
    همدا \LR{$\ltrue$} دے۔}
\olpseganchor{OLP-0158-B014}{end}

\olpseganchor{OLP-0158-B015}{start}
\item \indcase{!A}{\Atom{P}{t_1,\dots,
    t_n}}{\LR{$\Subst{\indfrm}{t}{x}$} همدا
    \LR{$\Atom{P}{\Subst{t_1}{t}{x},
    \dots, \Subst{t_n}{t}{x}}$} دے۔}
\olpseganchor{OLP-0158-B015}{end}

\olpseganchor{OLP-0158-B016}{start}
\item \indcase{!A}{\eq[t_1][t_2]}{\LR{$\Subst{\indfrmp}{t}{x}$} همدا
  \LR{$\Subst{t_1}{t}{x} = \Subst{t_2}{t}{x}$} دے۔}
\olpseganchor{OLP-0158-B016}{end}

\olpseganchor{OLP-0158-B017}{start}
\item \indcase{!A}{\lnot !B}{\LR{$\Subst{\indfrmp}{t}{x}$}
    همدا \LR{$\lnot \Subst{!B}{t}{x}$} دے۔}
\olpseganchor{OLP-0158-B017}{end}

\olpseganchor{OLP-0158-B018}{start}
\item \indcase{!A}{(!B \land
    !C)}{\LR{$\Subst{\indfrmp}{t}{x}$} همدا \LR{$(\Subst{!B}{t}{x} \land
    \Subst{!C}{t}{x})$} دے۔}
\olpseganchor{OLP-0158-B018}{end}

\olpseganchor{OLP-0158-B019}{start}
\item \indcase{!A}{(!B \lor
    !C)}{\LR{$\Subst{\indfrmp}{t}{x}$} همدا \LR{$(\Subst{!B}{t}{x} \lor
    \Subst{!C}{t}{x})$} دے۔}
\olpseganchor{OLP-0158-B019}{end}

\olpseganchor{OLP-0158-B020}{start}
\item \indcase{!A}{(!B \lif
    !C)}{\LR{$\Subst{\indfrmp}{t}{x}$} همدا \LR{$(\Subst{!B}{t}{x} \lif
    \Subst{!C}{t}{x})$} دے۔}
\olpseganchor{OLP-0158-B020}{end}

\olpseganchor{OLP-0158-B021}{start}
\item \indcase{!A}{(!B \liff
    !C)}{\LR{$\Subst{\indfrmp}{t}{x}$} همدا \LR{$(\Subst{!B}{t}{x} \liff
    \Subst{!C}{t}{x})$} دے۔}
\olpseganchor{OLP-0158-B021}{end}

\olpseganchor{OLP-0158-B022}{start}
\item 
\indcase{!A}{\lforall[y][!B]}{که \LR{$y$} له \LR{$x$} څخه بېل متغير وي، نو
    \LR{$\Subst{\indfrmp}{t}{x}$} همدا
    \LR{$\lforall[y][\Subst{!B}{t}{x}]$} دے؛ که نۀ وي، نو
    \LR{$\Subst{\indfrmp}{t}{x}$} پخپله \LR{$\indfrm$} دے۔}
\olpseganchor{OLP-0158-B022}{end}

\olpseganchor{OLP-0158-B023}{start}
\item 
\indcase{!A}{\lexists[y][!B]}{که \LR{$y$} له \LR{$x$} څخه بېل متغير وي، نو
    \LR{$\Subst{\indfrmp}{t}{x}$} همدا
    \LR{$\lexists[y][\Subst{!B}{t}{x}]$} دے؛ که نۀ وي، نو
    \LR{$\Subst{\indfrmp}{t}{x}$} پخپله \LR{$\indfrm$} دے۔}
\end{enumerate}
\end{defn}
\olpseganchor{OLP-0158-B023}{end}

\olpseganchor{OLP-0158-B024}{start}
\begin{explain}
پام وکړئ چې ځايناستي کېدے شي بې‌اغېزې وي: که \LR{$x$} په \LR{$!A$} کښې هېڅ
نۀ وي راغلے، نو \LR{$\Subst{!A}{t}{x}$} پخپله~\LR{$!A$} دے۔
\olpseganchor{OLP-0158-B024}{end}

\olpseganchor{OLP-0158-B025}{start}
دا محدوديت چې \LR{$t$} بايد په~\LR{$!A$} کښې د~\LR{$x$} دپاره ازاد وي، د
لاندې حالتونو د مخنيوي دپاره اړين دے۔ که \LR{$!A \ident
\lexists[y][x < y]$} او \LR{$t \ident y$} وي، نو \LR{$\Subst{!A}{t}{x}$} به
\LR{$\lexists[y][y < y]$} شي۔ په دې حالت کښې ازاد متغير \LR{$y$} د ځايناستي
پر مهال د سور \LR{$\lexists[y]$} له خوا «راګېرېږي»، او دا ناغوښتې پايله
ده۔ مثلاً غواړو چې هر کله \LR{$\lforall[x][!B]$} صدق کوي،
\LR{$\Subst{!B}{t}{x}$} هم صدق وکړي۔ خو
\LR{$\lforall[x][\lexists[y][x < y]]$} په پام کښې ونيسئ (دلته \LR{$!B$} همدا
\LR{$\lexists[y][x < y]$} دے)۔ دا د بېلګې په توګه د طبيعي عددونو په باب
رښتونې جمله ده: د هر عدد~\LR{$x$} دپاره تر هغه لوے يو عدد~\LR{$y$} شته۔ که
مو~\LR{$y$} د~\LR{$x$} د ممکنې ځايناستۍ په توګه روا ګڼلے واے، پايله به مو
\LR{$\Subst{!B}{y}{x} \ident \lexists[y][y < y]$} وه، چې دروغ ده۔ د دې
مخنيوے په دې کوو چې په~\LR{$t$} کښې هېڅ ازاد متغير بايد په~\LR{$!A$} کښې د
کوم سور له خوا تړلے نۀ شي۔
\end{explain}
\olpseganchor{OLP-0158-B025}{end}

\olpseganchor{OLP-0158-B026}{start}
د پېچلې نښې د مخنيوي دپاره ډېر ځله لاندې دود کاروو: که \LR{$!A$} داسې
فارمول وي چې متغير~\LR{$x$} پکښې ازاد راتلے شي، د دې د
ښودلو دپاره~\LR{$!A(x)$} هم ليکو۔ چې موخه مو کوم \LR{$!A$} او~\LR{$x$} دي څرګنده
وي، او \LR{$t$} يو ترم وي (چې فرض کېږي په \LR{$!A(x)$} کښې د \LR{$x$} دپاره ازاد
دے)، نو \LR{$!A(t)$} د \LR{$\Subst{!A}{t}{x}$} لنډون ليکو۔ مثلاً ويلے شو:
«\LR{$!A(t)$} د~\LR{$\lforall[x][!A(x)]$} يو مورد بولو۔» له دې څخه موخه دا
ده چې که \LR{$!A$} هر فارمول، \LR{$x$}~متغير، او \LR{$t$}~داسې ترم وي
چې په~\LR{$!A$} کښې د~\LR{$x$} دپاره ازاد وي، نو \LR{$\Subst{!A}{t}{x}$} د
\LR{$\lforall[x][!A]$} يو مورد دے۔
\olpseganchor{OLP-0158-B026}{end}

\olpseganchor{OLP-0159-B004}{start}
\olchapter{fol}{sem}{د لومړۍ درجې منطق معنٰی پوهنه}
\olpseganchor{OLP-0159-B004}{end}

\olpseganchor{OLP-0160-B004}{start}
\olfileid{fol}{syn}{its}
\olpseganchor{OLP-0160-B004}{end}

\olpseganchor{OLP-0160-B005}{start}
\olsection{پېژندنه}
\olpseganchor{OLP-0160-B005}{end}

\olpseganchor{OLP-0160-B006}{start}
عبارتونو ته معنا ورکول د معنٰی پوهنې کار دے۔ د معنٰی پوهنې مرکزي
مفهوم په جوړښت کښې د صدق مفهوم دے۔ جوړښت د ژبې
بنسټيزو اجزاوو ته معنا ورکوي: د تعريف ساحه د څيزونو يو ناتش سټ دے۔
سورونه داسې تفسيرېږي چې پر همدې ساحه ګرځي، ثابت سمبولونه د ساحې
له غړو سره ګومارل کېږي، تابعې له د تعريف ساحه څخه
همدې ساحې ته له تابعو سره ګومارل کېږي، او پريديکاتونه پر
د تعريف ساحه له اړيکو سره ګومارل کېږي۔ د تعريف ساحه او د بنسټيز لغت
ګومارنې په ګډه جوړښت جوړوي۔ متغيرونه په فارمولونه
کښې راتلے شي، او د معنٰی پوهنې د ورکولو دپاره بايد هغو ته هم د
د تعريف ساحه غړي وګومارو---دا د متغير ګومارنه ده۔ په پای کښې
د صدق اړيکه دا ټول سره يوځاے کوي۔ فارمول په
جوړښت~\LR{$\Struct{M}$} کښې د متغير ګومارنې~\LR{$s$} په نسبت
صدق کولے شي، چې \LR{$\Sat{M}{!A}[s]$} ليکل کېږي۔ دا اړيکه هم د~\LR{$!A$} پر
جوړښت په استقرا تعريفېږي: د منطقي نښلوونکو د رښتياوالي جدولونه
کاروو، څو مثلاً د \LR{$(!A \land !B)$} صدق د \LR{$!A$} او~\LR{$!B$} د صدق يا نه
صدق له مخې تعريف کړو۔ بيا څرګندېږي چې که فارمول~\LR{$!A$}
تړلے فارمول وي، يعنې ازاد متغيرونه ونۀ لري، د متغير ګومارنه
اړونده نۀ ده؛ نو ويلے شو چې تړلي فارمولونه په ساده ډول په
جوړښتونه کښې صدق کوي يا نۀ کوي۔
\olpseganchor{OLP-0160-B006}{end}

\olpseganchor{OLP-0160-B007}{start}
د تړلي فارمولونه دپاره د صدق د اړيکې \LR{$\Sat{M}{!A}$} پر بنسټ بيا د
اعتبار، معنايي استلزام، او د صدق وړتيا بنسټيز معنايي مفهومونه
تعريفولے شو۔ تړلے فارمول هله معتبره ده، \LR{$\Entails !A$}، چې هر
جوړښت پکښې صدق وکړي۔ د تړلي فارمولونه يو سټ هله هغه لازموي،
\LR{$\Gamma \Entails !A$}، چې هر هغه جوړښت چې د~\LR{$\Gamma$} په ټولو
تړلي فارمولونه کښې صدق کوي، په~\LR{$!A$} کښې هم صدق وکړي۔ او د
تړلي فارمولونه يو سټ هله د صدق وړ دے چې کوم جوړښت په هغه کښې
په يو وخت په ټولو تړلي فارمولونه کښې صدق وکړي۔ څرنګه چې فارمولونه
په استقرايي ډول تعريف شوي، او صدق بيا د فارمولونه پر جوړښت په
استقرا تعريفېږي، استقرا کارولے شو څو د خپلې معنٰی پوهنې خاصيتونه ثابت
او تعريف شوي معنايي مفهومونه سره وتړو۔
\olpseganchor{OLP-0160-B007}{end}

\olpseganchor{OLP-0161-B004}{start}
\olfileid{fol}{syn}{str}
\olpseganchor{OLP-0161-B004}{end}

\olpseganchor{OLP-0161-B005}{start}
\olsection{د لومړۍ درجې ژبو جوړښتونه}
\olpseganchor{OLP-0161-B005}{end}

\olpseganchor{OLP-0161-B006}{start}
\begin{explain}
د لومړۍ درجې ژبې پخپله \emph{بې‌تفسيره} دي: له ثابت سمبولونه،
تابعې او پريديکاتونه سره کومه ځانګړې معنا نۀ وي تړلې۔
معنا د  \emph{جوړښت} په ټاکلو ورکول کېږي۔
هغه \emph{د تعريف ساحه} ټاکي؛ يعنې هغه څيزونه چې ثابت سمبولونه ئې
ټاکي، تابعې پرې عمل کوي، او سورونه پرې ګرځي۔ سربېره پر دې،
دا هم ټاکي چې کوم ثابت سمبولونه کوم څيزونه ټاکي، تابع
څيزونه څنګه نورو څيزونو ته لېږي، او پريديکاتونه پر کومو څيزونو
صدق کوي۔ جوړښتونه په منطق کښې د \emph{معنايي} مفهومونو بنسټ
دي، لکه لازميت، اعتبار، او د صدق وړتيا۔ په ليکنو کښې ورته په بېلابېل
ډول «جوړښتونه»، «تفسيرونه» يا «مدلونه» ويل کېږي۔
\end{explain}
\olpseganchor{OLP-0161-B006}{end}

\olpseganchor{OLP-0161-B007}{start}
\begin{defn}[جوړښتونه]
د لومړۍ درجې منطق د يوې ژبې \LR{$\Lang{L}$} دپاره 
\emph{جوړښت}~\LR{$\Struct M$} له لاندې اجزاوو جوړ دے:
\begin{enumerate}
\item \emph{د تعريف ساحه:} يو ناتش سټ، \LR{$\Domain M$}
\item \emph{د ثابت سمبولونه تفسير:} د \LR{$\Lang{L}$} د هر
  ثابت سمبول~\LR{$c$} دپاره يو غړے \LR{$\Assign{c}{M} \in \Domain M$}
\item \emph{د پريديکاتونه تفسير:} د \LR{$\Lang{L}$} د هر \LR{$n$}-ځايه
  پريديکات~\LR{$R$} دپاره (له \LR{$\eq$} پرته) يوه \LR{$n$}-ځايه اړيکه
  \LR{$\Assign{R}{M} \subseteq \Domain{M}^n$}
\item \emph{د تابعې تفسير:} د \LR{$\Lang{L}$} د هر \LR{$n$}-ځايه
  تابع~\LR{$f$} دپاره يوه \LR{$n$}-ځايه تابع \LR{$\Assign{f}{M}
  \colon \Domain{M}^n \to \Domain{M}$}
\end{enumerate}
\end{defn}
\olpseganchor{OLP-0161-B007}{end}

\olpseganchor{OLP-0161-B008}{start}
\begin{ex}
د حساب د ژبې دپاره جوړښت~\LR{$\Struct M$} له يوه سټ؛ د
\LR{$\Domain M$} له يوه غړي \LR{$\Assign{\Obj 0}{M}$}، چې د
ثابت سمبول~\LR{$\Obj 0$} تفسير دے؛ له يوې يو ځايه تابع
\LR{$\Assign{\Obj \prime}{M} \colon \Domain{M} \to \Domain M$}؛ له دوو
دوه ځايه تابعو \LR{$\Assign{\Obj +}{M}$} او \LR{$\Assign{\Obj
  \times}{M}$}، چې دواړه \LR{$\Domain M^2 \to \Domain M$} دي؛ او له يوې
دوه ځايه اړيکې \LR{$\Assign{\Obj <}{M} \subseteq \Domain{M}^2$} جوړ دے۔
\olpseganchor{OLP-0161-B008}{end}

\olpseganchor{OLP-0161-B009}{start}
د داسې جوړښت يوه څرګنده بېلګه دا ده:
\begin{enumerate}
\item \LR{$\Domain N = \Nat$}
\item \LR{$\Assign{\Obj 0}{N} = 0$}
\item د هر \LR{$n \in \Nat$} دپاره \LR{$\Assign{\Obj \prime}{N}(n) = n + 1$}
\item د هر \LR{$n, m \in \Nat$} دپاره \LR{$\Assign{\Obj +}{N}(n, m) = n + m$}
\item د هر \LR{$n, m \in \Nat$} دپاره \LR{$\Assign{\Obj \times}{N}(n, m) = n\cdot m$}
\item \LR{$\Assign{\Obj <}{N} = \Setabs{\tuple{n, m}}{n \in \Nat, m \in
  \Nat, n < m}$}
\end{enumerate}
د \LR{$\Lang L_A$} دپاره په دې ډول تعريف شوے جوړښت~\LR{$\Struct N$} د
\emph{حساب معياري مدل} بلل کېږي، ځکه د~\LR{$\Lang L_A$} غير منطقي ثابتونه
کټ مټ هغسې تفسيروي چې تمه ئې کوئ۔
\olpseganchor{OLP-0161-B009}{end}

\olpseganchor{OLP-0161-B010}{start}
خو د~\LR{$\Lang L_A$} دپاره ډېر نور ممکن جوړښتونه هم شته۔ مثلاً د
ساحې په توګه د~\LR{$\Nat$} پر ځاے د صحيح عددونو سټ~\LR{$\Int$} اخيستلے، او د
\LR{$\Obj 0$}، \LR{$\Obj \prime$}، \LR{$\Obj +$}، \LR{$\Obj \times$} او \LR{$\Obj <$} تفسيرونه
ورسره سم تعريفولے شو۔ خو د~\LR{$\Lang L_A$} داسې جوړښتونه هم تعريفولے شو
چې له عددونو سره هېڅ لرې اړيکه هم نۀ لري۔
\end{ex}
\olpseganchor{OLP-0161-B010}{end}

\olpseganchor{OLP-0161-B011}{start}
\begin{ex}
د سټ تيورۍ د ژبې~\LR{$\Lang L_Z$} يو جوړښت~\LR{$\Struct M$} يوازې يوه ساحه
او يوه دوه ځايه اړيکه غواړي۔ نو په فني لحاظ مثلاً د خلکو سټ او اړيکه
«\LR{$x$} له \LR{$y$} څخه مشر دے» هم د \LR{$\Lang L_Z$} دپاره جوړښت
کېدے شي، او هم~\LR{$\Nat$} له \LR{$n \ge m$} سره د \LR{$n, m \in \Nat$} دپاره۔
\olpseganchor{OLP-0161-B011}{end}

\olpseganchor{OLP-0161-B012}{start}
د \LR{$\Lang L_Z$} يو په ځانګړي ډول په زړه پورې جوړښت، چې د ساحې
غړي ئې په رښتيا سټونه او د \LR{$\Obj \in$} تفسير ئې په رښتيا
اړيکه «\LR{$x$} د~\LR{$y$} غړے دے» وي، د \emph{موروثي متناهي
سټونو} جوړښت~\LR{$\Struct{{HF}}$} دے:
\begin{enumerate}
\item \LR{$\Domain{{{HF}}} = \emptyset \cup \Pow{\emptyset} \cup
  \Pow{\Pow{\emptyset}} \cup \Pow{\Pow{\Pow{\emptyset}}} \cup \dots$}؛
\item \LR{$\Assign{\Obj \in}{{{HF}}} = \Setabs{\tuple{x, y}}{x, y \in
  \Domain{{{HF}}}, x \in y}$}۔
\end{enumerate}
\end{ex}
\olpseganchor{OLP-0161-B012}{end}

\olpseganchor{OLP-0161-B013}{start}
\begin{digress}
دا چې د څه شي د جوړښت ګڼلو دپاره کوم شرطونه ږدو، زموږ پر
منطق اغېز کوي۔ مثلاً د تشو ساحو نۀ منل، د سورو لرونکو جملو د صدق د
معمول حساب له مخې، دا تضمينوي چې
\LR{$\lexists[x][(!A(x) \lor \lnot !A(x))]$} معتبر دے---يعنې يو منطقي صدق
دے۔ او دا شرط چې ټول ثابت سمبولونه بايد د ساحې کوم څيز ته اشاره وکړي،
تضمينوي چې وجودي تعميم د استنتاج يوه سمه بڼه ده:
\LR{$!A(a)$}، نو \LR{$\lexists[x][!A(x)]$}۔ که مو نومونو ته د ساحې بهر يا هېڅ
څيز ته اشاره کول روا ګڼلي واے، د \emph{ازاد منطق} پر لور به تللي وو؛
په هغه کښې وجودي تعميم يوه اضافي مقدمې ته اړتيا لري:
\LR{$!A(a)$} او \LR{$\lexists[x][\eq[x][a]]$}، نو \LR{$\lexists[x][!A(x)]$}۔
\end{digress}
\olpseganchor{OLP-0161-B013}{end}

\olpseganchor{OLP-0162-B004}{start}
\olfileid{fol}{syn}{cov}
\olpseganchor{OLP-0162-B004}{end}

\olpseganchor{OLP-0162-B005}{start}
\olsection{پوښل شوي جوړښتونه د لومړۍ درجې ژبو دپاره}
\olpseganchor{OLP-0162-B005}{end}

\olpseganchor{OLP-0162-B006}{start}
\begin{explain}
په ياد ولرئ چې يو ترم هله \emph{تړلے} دے چې هېڅ متغيرونه ونۀ لري۔
\end{explain}
\olpseganchor{OLP-0162-B006}{end}

\olpseganchor{OLP-0162-B007}{start}
\begin{defn}[د تړلو ترمونو قيمت]
که \LR{$t$} د ژبې~\LR{$\Lang L$} يو تړلے ترم او \LR{$\Struct M$} د~\LR{$\Lang L$} يو
جوړښت وي، \emph{قيمت}~\LR{$\Value{t}{M}$} ئې په لاندې ډول
تعريفېږي:
\begin{enumerate}
\item که \LR{$t$} يوازې ثابت سمبول~\LR{$c$} وي، نو
  \LR{$\Value{c}{M} = \Assign{c}{M}$}۔
\item که \LR{$t$} د \LR{$\Atom{f}{t_1, \ldots, t_n}$} بڼه ولري، نو
  \[
  \Value{t}{M} = \Assign{f}{M}(\Value{t_1}{M}, \ldots,
  \Value{t_n}{M}).
  \]
\end{enumerate}
\end{defn}
\olpseganchor{OLP-0162-B007}{end}

\olpseganchor{OLP-0162-B008}{start}
\begin{defn}[پوښل شوے جوړښت]
يو جوړښت هله \emph{پوښل شوے} دے چې د ساحې هر غړے د کوم تړلي
ترم قيمت وي۔
\end{defn}
\olpseganchor{OLP-0162-B008}{end}

\olpseganchor{OLP-0162-B009}{start}
\begin{ex}
فرض کړئ~\LR{$\Lang L$} هغه ژبه ده چې ثابت سمبولونه \LR{$\Obj{zero}$}،
\LR{$\Obj{one}$}، \LR{$\Obj{two}$}، \dots، دوه ځايه پريديکات~\LR{$<$}، او
دوه ځايه تابعې \LR{$+$} او \LR{$\times$} لري۔ د~\LR{$\Lang L$}
جوړښت~\LR{$\Struct M$} دې هغه وي چې ساحه ئې
\LR{$\Domain M = \{0, 1, 2, \ldots \}$} او ګومارنې ئې
\LR{$\Assign{\Obj{zero}}{M} = 0$}، \LR{$\Assign{\Obj{one}}{M} = 1$}،
\LR{$\Assign{\Obj{two}}{M} = 2$}، او همداسې نورې دي۔ د دوه ځايه اړيکې د
سمبول \LR{$<$} دپاره، سټ \LR{$\Assign{<}{M}$} د ټولو هغو جوړو
\LR{$\tuple{c_1, c_2} \in \Domain{M}^2$} سټ دے چې \LR{$c_1$} له~\LR{$c_2$} څخه
کوچنے وي: مثلاً \LR{$\tuple{1, 3} \in \Assign{<}{M}$}، خو
\LR{$\tuple{2, 2} \notin \Assign{<}{M}$}۔ د دوه ځايه تابع \LR{$+$}
دپاره \LR{$\Assign{+}{M}$} په معمول ډول تعريف کړئ---مثلاً
\LR{$\Assign{+}{M}(2,3)$} عدد~\LR{$5$} ته لېږي؛ د دوه ځايه
تابع~\LR{$\times$} دپاره هم همداسې وکړئ۔ نو د \LR{$\Obj{four}$}
قيمت پخپله~\LR{$4$} دے، او د \LR{$\times(\Obj{two},
+(\Obj{three},\Obj{zero}))$} قيمت (يا په منځ‌ليکي نښه کښې،
\LR{$\Obj{two} \times (\Obj{three} + \Obj{zero})$}) دا دے:
\begin{multline*}
\Value{\times(\Obj{two}, +(\Obj{three},\Obj{zero}))}{M} \\
\begin{aligned}
& =\Assign{\times}{M}(\Value{\Obj{two}}{M}, \Value{+(\Obj{three}, \Obj{zero})}{M})\\
& = \Assign{\times}{M}(\Value{\Obj{two}}{M}, \Assign{+}{M}(\Value{\Obj{three}}{M},
\Value{\Obj{zero}}{M})) \\
& = \Assign{\times}{M}(\Assign{\Obj{two}}{M}, \Assign{+}{M}(\Assign{\Obj{three}}{M},
\Assign{\Obj{zero}}{M})) \\
& = \Assign{\times}{M}(2, \Assign{+}{M}(3, 0)) \\
& = \Assign{\times}{M}(2, 3) \\
& = 6
\end{aligned}
\end{multline*}
\end{ex}
\olpseganchor{OLP-0162-B009}{end}

\olpseganchor{OLP-0162-B010}{start}
\begin{prob}
ايا \LR{$\Struct N$}، يعنې د حساب معياري مدل، پوښل شوے دے؟ څرګنده ئې کړئ۔
\end{prob}
\olpseganchor{OLP-0162-B010}{end}

\olpseganchor{OLP-0163-B004}{start}
\olfileid{fol}{syn}{sat}
\olpseganchor{OLP-0163-B004}{end}

\olpseganchor{OLP-0163-B005}{start}
\olsection{په  جوړښت کښې د
   فارمول صدق}
\olpseganchor{OLP-0163-B005}{end}

\olpseganchor{OLP-0163-B006}{start}
\begin{explain}
هغه بنسټيز مفهومونه چې يوې خوا ته عبارتونه لکه ترمونه او
فارمولونه، او بلې خوا ته جوړښتونه سره تړي، د يوه ترم
\emph{قيمت} او د فارمول \emph{صدق} دي۔ په غير رسمي ډول،
د يوه ترم قيمت د جوړښت يو غړے دے---که ترم
يوازې يو ثابت وي، قيمت ئې هغه څيز دے چې جوړښت له هغه
ثابت سره ګومارلے وي؛ او که د تابعې په کارولو جوړ شوے وي، د
هغه قيمت د ثابتونو له قيمتونه او هغو تابعو څخه حسابېږي چې
جوړښت په ترم کښې تابعو ته ګومارلې وي۔ فارمول په يوه
جوړښت کښې هله \emph{صدق کوي} چې اړوندو ته ورکړل شوے تفسير
فارمول د جوړښت په ساحه کښې رښتونے کړي۔ د صدق دا مفهوم په
استقرا ټاکل کېږي: د جوړښت مشخصات نېغ په نېغه وايي چې اټومي
فارمولونه کله صدق کوي، او دا چې يو پېچلے فارمول کله صدق کوي،
د هغه له اصلي نښلوونکي يا سور او د هغه د سمدستي فرعي فارمولونه له
صدق يا نه صدق سره سم تعريفوو۔
\olpseganchor{OLP-0163-B006}{end}

\olpseganchor{OLP-0163-B007}{start}
دلته د سورو حالت لږ نازک دے، ځکه د سور لرونکي فارمول سمدستي
فرعي فارمول يو ازاد متغير لري، او جوړښتونه د
متغيرونه قيمتونه نۀ ټاکي۔ د دې ستونزې د حل دپاره د
\emph{متغير ګومارنې} هم معرفي کوو، او صدق يوازې د
جوړښت په نسبت نۀ، بلکې د جوړښت او متغير
ګومارنې په نسبت تعريفوو۔
\end{explain}
\olpseganchor{OLP-0163-B007}{end}

\olpseganchor{OLP-0163-B008}{start}
\begin{defn}[د متغير ګومارنه]
د جوړښت~\LR{$\Struct{M}$} دپاره \emph{د متغير ګومارنه}~\LR{$s$} هغه
تابع ده چې هر متغير د~\LR{$\Domain M$} کوم غړے ته لېږي؛
يعنې \LR{$s\colon \Var \to \Domain M$}۔
\end{defn}
\olpseganchor{OLP-0163-B008}{end}

\olpseganchor{OLP-0163-B009}{start}
\begin{explain}
جوړښت هر ثابت سمبول ته قيمت، او د متغير ګومارنه
هر متغير ته يو قيمت ورکوي۔ خو غواړو چې له دغو څخه جوړ شوي ترمونه هم
د د تعريف ساحه غړي ونوموي۔ د دې دپاره د ترمونو قيمت په
استقرا تعريفوو۔ د ثابت سمبولونه او متغيرونو قيمت هماغه دے چې
جوړښت يا د متغير ګومارنه ئې ټاکي؛ د زيات پېچلو ترمونو قيمت
د هغو تابعو په بازګشتي کارولو حسابېږي چې جوړښت له
تابعې سره ګومارلې وي۔
\end{explain}
\olpseganchor{OLP-0163-B009}{end}

\olpseganchor{OLP-0163-B010}{start}
\begin{defn}[د ترمونو قيمت]
که \LR{$t$} د ژبې~\LR{$\Lang L$} يو ترم، \LR{$\Struct M$} د~\LR{$\Lang L$} يو
جوړښت، او \LR{$s$} د~\LR{$\Struct M$} دپاره متغير ګومارنه وي،
\emph{قيمت}~\LR{$\Value{t}{M}[s]$} ئې په لاندې ډول تعريفېږي:
\begin{enumerate}
\item \indcase{t}{c}{\LR{$\Value{\indfrm}{M}[s] = \Assign{\indcomplex}{M}$}۔}
\item \indcase{t}{x}{\LR{$\Value{\indfrm}{M}[s] = s(\indcomplex)$}۔}
\item \indcase{t}{\Atom{f}{t_1, \ldots, t_n}}{
\[
\Value{\indfrm}{M}[s] = \Assign{f}{M}(\Value{t_1}{M}[s], \ldots,
\Value{t_n}{M}[s]).
\]}
\end{enumerate}
\end{defn}
\olpseganchor{OLP-0163-B010}{end}

\olpseganchor{OLP-0163-B011}{start}
\begin{defn}[د~\LR{$x$} بديله ګومارنه]
که \LR{$s$} د جوړښت~\LR{$\Struct M$} دپاره متغير ګومارنه وي،
د~\LR{$\Struct M$} هره هغه متغير ګومارنه~\LR{$s'$} چې له~\LR{$s$} سره تر
ډېره يوازې د~\LR{$x$} په ګومارل شوي څيز کښې توپير لري، د~\LR{$s$}
\emph{د~\LR{$x$} بديله ګومارنه} بلل کېږي۔ که \LR{$s'$} د~\LR{$s$} د~\LR{$x$} بديله
ګومارنه وي، \LR{$\varAssign{s'}{s}{x}$} ليکو۔
\end{defn}
\olpseganchor{OLP-0163-B011}{end}

\olpseganchor{OLP-0163-B012}{start}
\begin{explain}
پام وکړئ چې د يوې ګومارنې~\LR{$s$} د~\LR{$x$} بديله ګومارنه اړ نۀ ده چې
په رښتيا~\LR{$x$} ته بېل څيز وګوماري۔ په حقيقت کښې هره ګومارنه د خپل ځان
د~\LR{$x$} بديله ګومارنه هم ګڼل کېږي۔
\end{explain}
\olpseganchor{OLP-0163-B012}{end}

\olpseganchor{OLP-0163-B013}{start}
\begin{defn}
  که \LR{$s$} د جوړښت~\LR{$\Struct M$} دپاره متغير ګومارنه
  او \LR{$m \in \Domain{M}$} وي، نو ګومارنه~\LR{$\Subst{s}{m}{x}$} هغه د
  متغير ګومارنه ده چې داسې تعريفېږي:
  \[\Subst{s}{m}{x}(y) = \begin{cases}
    m & \text{\RL{که }} y \ident x\\
    s(y) & \text{\RL{که نۀ وي}}.
  \end{cases}\]
\end{defn}
\olpseganchor{OLP-0163-B013}{end}

\olpseganchor{OLP-0163-B014}{start}
په بله وينا، \LR{$\Subst{s}{m}{x}$} د~\LR{$s$} د~\LR{$x$} هغه ځانګړې بديله ګومارنه
ده چې د ساحې غړے~\LR{$m$} له~\LR{$x$} سره ګوماري، او له~\LR{$x$} پرته
متغيرونه ته هماغه څيزونه ګوماري چې \LR{$s$} ئې ګوماري۔
\olpseganchor{OLP-0163-B014}{end}

\olpseganchor{OLP-0163-B015}{start}
\begin{defn}[صدق]
\ollabel{defn:satisfaction}
په جوړښت~\LR{$\Struct M$} کښې د فارمول~\LR{$!A$} صدق د
متغير ګومارنې~\LR{$s$} په نسبت، چې په سمبولونو کښې
\LR{$\Sat{M}{!A}[s]$} دے، په بازګشتي ډول داسې تعريفېږي۔ (موږ
\LR{$\Sat/{M}{!A}[s]$} د «\LR{$\Sat{M}{!A}[s]$} نۀ دے» په معنا ليکو۔)
\begin{enumerate}
\item %
  \indcase{!A}{\lfalse}{\LR{$\Sat/{M}{\indfrm}[s]$}۔}
\olpseganchor{OLP-0163-B015}{end}

\olpseganchor{OLP-0163-B016}{start}
\item %
  \indcase{!A}{\ltrue}{\LR{$\Sat{M}{\indfrm}[s]$}۔}
\olpseganchor{OLP-0163-B016}{end}

\olpseganchor{OLP-0163-B017}{start}
\item \indcase{!A}{\Atom{R}{t_1, \dots, t_n}}{\LR{$\Sat{M}{\indfrm}[s]$}
  هله او يوازې هله چې \LR{$\langle \Value{t_1}{M}[s], \dots,
  \Value{t_n}{M}[s] \rangle \in \Assign{R}{M}$}۔}
\olpseganchor{OLP-0163-B017}{end}

\olpseganchor{OLP-0163-B018}{start}
\item \indcase{!A}{\eq[t_1][t_2]}{\LR{$\Sat{M}{\indfrm}[s]$} هله او
  يوازې هله چې \LR{$\Value{t_1}{M}[s] = \Value{t_2}{M}[s]$}۔}
\olpseganchor{OLP-0163-B018}{end}

\olpseganchor{OLP-0163-B019}{start}
\item %
  \indcase{!A}{\lnot !B}{\LR{$\Sat{M}{\indfrm}[s]$} هله او يوازې هله چې
    \LR{$\Sat/{M}{!B}[s]$}۔}
\olpseganchor{OLP-0163-B019}{end}

\olpseganchor{OLP-0163-B020}{start}
\item %
  \indcase{!A}{(!B \land !C)}{\LR{$\Sat{M}{\indfrm}[s]$} هله او يوازې
    هله چې \LR{$\Sat{M}{!B}[s]$} او \LR{$\Sat{M}{!C}[s]$}۔}
\olpseganchor{OLP-0163-B020}{end}

\olpseganchor{OLP-0163-B021}{start}
\item %
  \indcase{!A}{(!B \lor !C)}{\LR{$\Sat{M}{\indfrm}[s]$} هله او يوازې
    هله چې \LR{$\Sat{M}{!B}[s]$} يا \LR{$\Sat{M}{!C}[s]$} (يا دواړه)۔}
\olpseganchor{OLP-0163-B021}{end}

\olpseganchor{OLP-0163-B022}{start}
\item %
  \indcase{!A}{(!B \lif !C)}{\LR{$\Sat{M}{\indfrm}[s]$} هله او يوازې
    هله چې \LR{$\Sat/{M}{!B}[s]$} يا \LR{$\Sat{M}{!C}[s]$} (يا دواړه)۔}
\olpseganchor{OLP-0163-B022}{end}

\olpseganchor{OLP-0163-B023}{start}
\item %
  \indcase{!A}{(!B \liff !C)}{\LR{$\Sat{M}{\indfrm}[s]$} هله او يوازې
    هله چې يا دواړه \LR{$\Sat{M}{!B}[s]$} او \LR{$\Sat{M}{!C}[s]$} وي، يا
    نۀ \LR{$\Sat{M}{!B}[s]$} وي او نۀ \LR{$\Sat{M}{!C}[s]$}۔}
\olpseganchor{OLP-0163-B023}{end}

\olpseganchor{OLP-0163-B024}{start}
\item %
  \indcase{!A}{\lforall[x][!B]}{\LR{$\Sat{M}{\indfrm}[s]$} هله او يوازې
    هله چې د هر غړے~\LR{$m \in \Domain M$} دپاره
    \LR{$\Sat{M}{!B}[\Subst{s}{m}{x}]$}۔}
\olpseganchor{OLP-0163-B024}{end}

\olpseganchor{OLP-0163-B025}{start}
\item %
  \indcase{!A}{\lexists[x][!B]}{\LR{$\Sat{M}{\indfrm}[s]$} هله او يوازې
  هله چې لږ تر لږه د يوه غړے~\LR{$m \in \Domain M$} دپاره
  \LR{$\Sat{M}{!B}[\Subst{s}{m}{x}]$}۔}
\end{enumerate}
\end{defn}
\olpseganchor{OLP-0163-B025}{end}

\olpseganchor{OLP-0163-B026}{start}
\begin{explain}
د متغير ګومارنې په وروستي
دوو شرطونو کښې مهمې
دي۔ د \LR{$\lforall[x][!B(x)]$} صدق داسې نۀ شو تعريفولے:
«د ټولو \LR{$m \in \Domain{M}$} دپاره \LR{$\Sat{M}{!B(m)}$}۔»
د \LR{$\lexists[x][!B(x)]$} صدق داسې نۀ شو تعريفولے: «لږ تر لږه د يوه
\LR{$m \in \Domain{M}$} دپاره \LR{$\Sat{M}{!B(m)}$}۔» لامل دا دے چې که
\LR{$m \in \Domain M$} وي، هغه د ژبې سمبول نۀ دے؛ نو \LR{$!B(m)$}
فارمول نۀ دے (يعنې \LR{$\Subst{!B}{m}{x}$} ناتعريف دے)۔ دا هم نۀ
شو فرضولے چې داسې ثابت سمبولونه يا ترمونه لرو چې د~\LR{$\Struct{M}$} هر
غړے ونوموي، ځکه د جوړښتونه په تعريف کښې هېڅ شرط دا
نۀ غواړي۔ په معياري ژبه کښې د ثابت سمبولونه سټ شمېرېدونکے لامتناهي دے؛
نو که \LR{$\Domain{M}$} د شمېر وړ نۀ وي، د هر څيز د نومولو دپاره
کافي ثابت سمبولونه هم نشته۔
\olpseganchor{OLP-0163-B026}{end}

\olpseganchor{OLP-0163-B027}{start}
دا ستونزه د متغير ګومارنو په معرفي کولو حلوو؛ په هغو سره
متغيرونه نېغ په نېغه د ساحې له غړي سره تړلے شو۔ بيا د دې
پر ځاے چې مثلاً ووايو
\LR{$\lexists[x][!B(x)]$} په
\LR{$\Struct M$} کښې هله او يوازې هله صدق کوي چې
لږ تر لږه د يوه \LR{$m \in \Domain{M}$} دپاره
\LR{$\Sat{M}{!B(m)}$} وي، وايو چې دا په~\LR{$\Struct M$} کښې د~\LR{$s$}
\emph{په نسبت} هله او يوازې هله صدق کوي چې \LR{$!B(x)$} د
\LR{$\Subst{s}{m}{x}$} په نسبت لږ تر لږه د يوه
\LR{$m \in \Domain M$} دپاره صدق وکړي۔
\end{explain}
\olpseganchor{OLP-0163-B027}{end}

\olpseganchor{OLP-0163-B028}{start}
\begin{ex}
فرض کړئ \LR{$\Lang{L} = \{a, b, f, R\}$}، چې \LR{$a$} او \LR{$b$}
ثابت سمبولونه، \LR{$f$}~دوه ځايه تابع، او \LR{$R$}~دوه ځايه
پريديکات دي۔ لاندې تعريف شوے جوړښت~\LR{$\Struct{M}$} په پام
کښې ونيسئ:
\begin{enumerate}
\item \LR{$\Domain M = \{1, 2, 3, 4\}$}
\item \LR{$\Assign{a}{M} = 1$}
\item \LR{$\Assign{b}{M} = 2$}
\item که \LR{$x+y \le 3$} وي، \LR{$\Assign{f}{M}(x, y) = x+y$}، او که نۀ وي
  \LR{$= 3$}۔
\item \LR{$\Assign{R}{M} = \{\tuple{1, 1}, \tuple{1, 2}, \tuple{2, 3}, \tuple{2, 4}\}$}
\end{enumerate}
تابع \LR{$s(x) = 1$}، چې هر متغير ته \LR{$1 \in \Domain{M}$} ګوماري،
د~\LR{$\Struct{M}$} دپاره د متغير ګومارنه ده۔
\olpseganchor{OLP-0163-B028}{end}

\olpseganchor{OLP-0163-B029}{start}
بيا
\begin{align*}
\Value{f(a,b)}{M}[s] & = \Assign{f}{M}(\Value{a}{M}[s], \Value{b}{M}[s]).
\intertext{\RL{څرنګه چې \LR{$a$} او \LR{$b$} ثابت سمبولونه دي،
  \LR{$\Value{a}{M}[s] = \Assign{a}{M} = 1$} او
  \LR{$\Value{b}{M}[s] = \Assign{b}{M} = 2$}۔ نو}}
\Value{f(a,b)}{M}[s] & = \Assign{f}{M}(1, 2) = 1+2 = 3.
\intertext{\RL{د \LR{$f(f(a,b),a)$} د قيمت د حسابولو دپاره بايد دا وګورو:}}
\Value{f(f(a,b),a)}{M}[s] & = \Assign{f}{M}(\Value{f(a, b)}{M}[s],
\Value{a}{M}[s]) = \Assign{f}{M}(3, 1) = 3,
\intertext{\RL{ځکه \LR{$3+1 > 3$}۔ څرنګه چې \LR{$s(x) = 1$} او
  \LR{$\Value{x}{M}[s] = s(x)$}، دا هم لرو:}}
\Value{f(f(a,b),x)}{M}[s] & = \Assign{f}{M}(\Value{f(a, b)}{M}[s],
\Value{x}{M}[s]) = \Assign{f}{M}(3, 1) = 3,
\end{align*}
\olpseganchor{OLP-0163-B029}{end}

\olpseganchor{OLP-0163-B030}{start}
يو اټومي فارمول~\LR{$R(t_1, t_2)$} هله صدق کوي چې د مدخلو د قيمتونو
جوړه، يعنې \LR{$\tuple{\Value{t_1}{M}[s],
  \Value{t_2}{M}[s]}$}، د~\LR{$\Assign{R}{M}$} غړے وي۔ نو مثلاً
\LR{$\Sat{M}{R(b,f(a,b))}[s]$} لرو، ځکه
\LR{$\tuple{\Value{b}{M}, \Value{f(a,b)}{M}} = \tuple{2, 3} \in
\Assign{R}{M}$}؛ خو \LR{$\Sat/{M}{R(x, f(a,b))}[s]$}، ځکه
\LR{$\tuple{1, 3} \notin \Assign{R}{M}$}۔
\olpseganchor{OLP-0163-B030}{end}

\olpseganchor{OLP-0163-B031}{start}
د دې د معلومولو دپاره چې يو غير اټومي فارمول~\LR{$!A$} صدق کوي که نۀ، د
استقرايي تعريف هغه شرط کاروئ چې پر اصلي نښلوونکي عملي کېږي۔ مثلاً د
\LR{$R(a, a) \lif (R(b, x) \lor R(x, b))$} اصلي نښلوونکے~\LR{$\lif$} دے، او
\begin{align*}
  & \Sat{M}{R(a, a) \lif (R(b, x) \lor R(x, b))}[s] \text{\RL{ هله او يوازې هله چې }}\\
  & \qquad
  \Sat/{M}{R(a,a)}[s] \text{\RL{ يا }} \Sat{M}{R(b, x) \lor R(x, b)}[s]
  \intertext{\RL{څرنګه چې \LR{$\Sat{M}{R(a,a)}[s]$} (ځکه
    \LR{$\tuple{1,1} \in \Assign{R}{M}$})، ځواب لا نۀ شو ټاکلے او لومړے
    بايد معلومه کړو چې \LR{$\Sat{M}{R(b, x) \lor R(x, b)}[s]$} دے که نۀ:}}
  & \Sat{M}{R(b, x) \lor R(x, b)}[s] \text{\RL{ هله او يوازې هله چې }}\\
  & \qquad \Sat{M}{R(b,x)}[s] \text{\RL{ يا }} \Sat{M}{R(x, b)}[s]
  \intertext{\RL{او دا حالت سم دے، ځکه \LR{$\Sat{M}{R(x, b)}[s]$}
    (ځکه \LR{$\tuple{1,2} \in \Assign{R}{M}$})۔}}
\end{align*}
\olpseganchor{OLP-0163-B031}{end}

\olpseganchor{OLP-0163-B032}{start}
په ياد ولرئ چې د~\LR{$s$} د~\LR{$x$} بديله ګومارنه هغه د متغير ګومارنه ده چې
له~\LR{$s$} سره تر ډېره يوازې د~\LR{$x$} په ګومارنه کښې توپير لري۔ د
\LR{$\Domain{M}$} د هر غړے دپاره د~\LR{$s$} يوه د~\LR{$x$} بديله ګومارنه
شته:
\begin{align*}
  s_1 & = \Subst{s}{1}{x}, &
  s_2 & = \Subst{s}{2}{x},\\
  s_3 & = \Subst{s}{3}{x}, &
  s_4 & = \Subst{s}{4}{x}.
\end{align*}
نو مثلاً \LR{$s_2(x) = 2$}، او له~\LR{$x$} پرته د هر متغير~\LR{$y$} دپاره
\LR{$s_2(y) = s(y) = 1$}۔ د جوړښت~\LR{$\Struct{M}$} دپاره دا د~\LR{$s$} د~\LR{$x$}
ټولې بديلې ګومارنې دي، ځکه \LR{$\Domain{M} = \{1, 2, 3, 4\}$}۔ په ځانګړي
ډول پام وکړئ چې \LR{$s_1 = s$} (هره~\LR{$s$} تل د خپل ځان د~\LR{$x$} بديله ګومارنه
هم وي)۔
\olpseganchor{OLP-0163-B032}{end}

\olpseganchor{OLP-0163-B033}{start}
د دې د معلومولو دپاره چې د وجودي سور لرونکے
  فارمول~\LR{$\lexists[x][!A(x)]$} صدق کوي که نۀ، بايد معلومه کړو چې
  لږ تر لږه د يوه \LR{$m \in \Domain M$} دپاره
  \LR{$\Sat{M}{!A(x)}[\Subst{s}{m}{x}]$} دے که نۀ۔ نو
  \[
  \Sat{M}{\lexists[x][(R(b,x) \lor R(x,b))]}[s],
  \]
  ځکه \LR{$\Sat{M}{R(b,x) \lor R(x, b)}[\Subst{s}{1}{x}]$}
  (\LR{$\Subst{s}{3}{x}$} هم شرط پوره کوي)۔ خو
  \[
  \Sat/{M}{\lexists[x][(R(b,x) \land R(x,b))]}[s]
  \]
  ځکه هر \LR{$m \in \Domain{M}$} چې واخلو،
  \LR{$\Sat/{M}{R(b,x) \land R(x,b)}[\Subst{s}{m}{x}]$}۔
\olpseganchor{OLP-0163-B033}{end}

\olpseganchor{OLP-0163-B034}{start}
د دې د معلومولو دپاره چې د عمومي سور لرونکے
  فارمول~\LR{$\lforall[x][!A(x)]$} صدق کوي که نۀ، بايد معلومه کړو چې
  د ټولو \LR{$m \in \Domain M$} دپاره
  \LR{$\Sat{M}{!A(x)}[\Subst{s}{m}{x}]$} دے که نۀ۔ نو
  \[
  \Sat{M}{\lforall[x][(R(x,a) \lif R(a,x))]}[s],
  \]
  ځکه د ټولو \LR{$m \in \Domain M$} دپاره
  \LR{$\Sat{M}{R(x,a) \lif R(a,x)}[\Subst{s}{m}{x}]$}۔ د \LR{$m = 1$} دپاره
  \LR{$\Sat{M}{R(a,x)}[\Subst{s}{1}{x}]$} لرو، نو تالي رښتونے دے؛ او د
  \LR{$m = 2$}، \LR{$3$} او~\LR{$4$} دپاره
  \LR{$\Sat/{M}{R(x,a)}[\Subst{s}{m}{x}]$}، نو مقدم دروغ دے۔ خو
  \[
  \Sat/{M}{\lforall[x][(R(a,x) \lif R(x,a))]}[s]
  \]
  ځکه \LR{$\Sat/{M}{R(a,x) \lif R(x,a)}[\Subst{s}{2}{x}]$}
  (ځکه \LR{$\Sat{M}{R(a, x)}[\Subst{s}{2}{x}]$} او
  \LR{$\Sat/{M}{R(x, a)}[\Subst{s}{2}{x}]$})۔
\olpseganchor{OLP-0163-B034}{end}





\olpseganchor{OLP-0163-B037}{start}
د يو زيات پېچلي حالت دپاره دا فارمول په پام کښې ونيسئ:
\[
\lforall[x][(R(a,x) \lif \lexists[y][R(x,y)])].
\]
څرنګه چې \LR{$\Sat/{M}{R(a,x)}[\Subst{s}{3}{x}]$} او
\LR{$\Sat/{M}{R(a,x)}[\Subst{s}{4}{x}]$}، هغه په زړه پورې حالتونه چې د
شرطي د تالي په باب فکر پکښې پکار دے يوازې \LR{$m = 1$} او \LR{$m = 2$} دي۔
ايا \LR{$\Sat{M}{\lexists[y][R(x,y)]}[\Subst{s}{1}{x}]$}؟ دا هله صدق کوي
چې لږ تر لږه يو \LR{$n \in \Domain M$} داسې وي چې
\LR{$\Sat{M}{R(x,y)}[\Subst{\Subst{s}{1}{x}}{n}{y}]$}۔ په حقيقت کښې که
\LR{$n = 1$} واخلو، نو
\LR{$\Subst{\Subst{s}{1}{x}}{n}{y} = \Subst{s}{1}{y} = s$}۔ څرنګه چې
\LR{$s(x) = 1$}، \LR{$s(y) = 1$}، او \LR{$\tuple{1,1} \in \Assign{R}{M}$}، ځواب هو
دے۔
\olpseganchor{OLP-0163-B037}{end}

\olpseganchor{OLP-0163-B038}{start}
د دې د معلومولو دپاره چې
\LR{$\Sat{M}{\lexists[y][R(x,y)]}[\Subst{s}{2}{x}]$}، بايد د
متغير ګومارنې \LR{$\Subst{\Subst{s}{2}{x}}{n}{y}$} وګورو۔ دلته
د \LR{$n = 1$} دپاره دا ګومارنه~\LR{$s_2 = \Subst{s}{2}{x}$} ده؛ او
\LR{$R(x,y)$} پکښې صدق نۀ کوي (\LR{$s_2(x) = 2$}، \LR{$s_2(y) = 1$}، او
\LR{$\tuple{2,1}\notin \Assign{R}{M}$})۔ خو
\LR{$\Subst{\Subst{s}{2}{x}}{3}{y} = \Subst{s_2}{3}{y}$} په پام کښې
ونيسئ۔ \LR{$\Sat{M}{R(x,y)}[\Subst{s_2}{3}{y}]$}، ځکه
\LR{$\tuple{2,3} \in \Assign{R}{M}$}؛ نو
\LR{$\Sat{M}{\lexists[y][R(x,y)]}[s_2]$}۔
\olpseganchor{OLP-0163-B038}{end}

\olpseganchor{OLP-0163-B039}{start}
نو د ټولو \LR{$m \in \Domain M$} دپاره يا
\LR{$\Sat/{M}{R(a,x)}[\Subst{s}{m}{x}]$} (که \LR{$m = 3$} يا \LR{$4$} وي)، يا
\LR{$\Sat{M}{\lexists[y][R(x,y)]}[\Subst{s}{m}{x}]$} (که \LR{$m = 1$} يا \LR{$2$}
وي)؛ نو
\[
\Sat{M}{\lforall[x][(R(a,x) \lif \lexists[y][R(x,y)])]}[s].
\]
بل خوا،
\[
\Sat/{M}{\lexists[x][(R(a,x) \land \lforall[y][R(x,y)])]}[s].
\]
\LR{$\Sat{M}{R(a,x)}[\Subst{s}{m}{x}]$} يوازې د \LR{$m = 1$} او \LR{$m = 2$}
دپاره لرو۔ خو د~\LR{$m$} د دواړو دغو قيمتونو دپاره داسې
\LR{$n \in \Domain M$} شته، يعنې \LR{$n = 4$}، چې
\LR{$\Sat/{M}{R(x,y)}[\Subst{\Subst{s}{m}{x}}{n}{y}]$}؛ نو د \LR{$m = 1$} او
\LR{$m = 2$} دپاره
\LR{$\Sat/{M}{\lforall[y][R(x,y)]}[\Subst{s}{m}{x}]$}۔ په ټوله کښې داسې
\LR{$m \in \Domain M$} نشته چې
\LR{$\Sat{M}{R(a,x) \land \lforall[y][R(x,y)]}[\Subst{s}{m}{x}]$}۔
\end{ex}
\olpseganchor{OLP-0163-B039}{end}

\olpseganchor{OLP-0163-B046}{start}
\begin{prob}
فرض کړئ \LR{$\Lang L = \{c, f, A\}$}، چې يو ثابت سمبول، يوه يو ځايه
تابع او يو دوه ځايه پريديکات لري، او
جوړښت~\LR{$\Struct{M}$} داسې ورکړل شوے دے:
\begin{enumerate}
\item \LR{$\Domain M = \{1, 2, 3\}$}
\item \LR{$\Assign{c}{M} = 3$}
\item \LR{$\Assign{f}{M}(1) = 2, \Assign{f}{M}(2) = 3, \Assign{f}{M}(3) = 2$}
\item \LR{$\Assign{A}{M} = \{\tuple{1, 2}, \tuple{2, 3}, \tuple{3, 3}\}$}
\end{enumerate}
(الف) د ټولو متغيرونه~\LR{$v$} دپاره \LR{$s(v) = 1$} وټاکئ۔ معلومه کړئ
چې آيا
\[
\Sat{M}{\lexists[x][(A(f(z), c) \lif \lforall[y][(A(y, x) \lor A(f(y),
      x))])]}[s]
\]
صدق کوي که نۀ، او لامل ئې څرګند کړئ۔
\olpseganchor{OLP-0163-B046}{end}

\olpseganchor{OLP-0163-B047}{start}
(ب) يو بېل جوړښت او متغير ګومارنه ورکړئ چې فارمول پکښې
صدق ونۀ کړي۔
\end{prob}
\olpseganchor{OLP-0163-B047}{end}

\olpseganchor{OLP-0164-B004}{start}
\olfileid{fol}{syn}{ass}
\olpseganchor{OLP-0164-B004}{end}

\olpseganchor{OLP-0164-B005}{start}
\olsection{د متغيرونو ګومارنې}
\olpseganchor{OLP-0164-B005}{end}

\olpseganchor{OLP-0164-B006}{start}
\begin{explain}
د متغير ګومارنه~\LR{$s$} \emph{هر} متغير ته يو ارزښت ورکوي---او
متغيرونه بې‌شمېره دي۔ البته دا اړينه نۀ ده۔ موږ له متغير
ګومارنو څخه يوازې د اسانتيا دپاره غواړو چې ټولو متغيرونه ته
ارزښتونه وګوماري۔ د يوه ترم~\LR{$t$} ارزښت، او دا چې فارمول~\LR{$!A$} د
ګومارنې~\LR{$s$} په نسبت په جوړښت کښې صدق کوي که نۀ، يوازې په هغو
ارزښتونو پورې تړلي دي چې~\LR{$s$} ئې د~\LR{$t$} متغيرونه او د~\LR{$!A$} ازادو
متغيرونه ته ګوماري۔ د لاندې دوو قضيو منځپانګه همدا ده۔ د «پورې
تړلي دي» مانا دقيقه کولو دپاره ښيو چې هرې دوې هغه متغير ګومارنې چې
د~\LR{$t$} پر ټولو متغيرونو سره موافقې وي، هغه ته يو شان ارزښت ورکوي؛ او
که دوې متغير ګومارنې د~\LR{$!A$} پر ټولو ازادو متغيرونو سره موافقې وي، نو
\LR{$!A$} د يوې په نسبت هله صدق کوي چې د بلې په نسبت هم صدق وکړي۔
\end{explain}
\olpseganchor{OLP-0164-B006}{end}

\olpseganchor{OLP-0164-B007}{start}
\begin{prop}\ollabel{prop:valindep}
که په ترم~\LR{$t$} کښې متغيرونه د \LR{$x_1$}، \dots،~\LR{$x_n$} په منځ کښې وي،
او \LR{$s_1(x_i) = s_2(x_i)$} د \LR{$i = 1$}، \dots،~\LR{$n$} دپاره وي، نو
\LR{$\Value{t}{M}[s_1] = \Value{t}{M}[s_2]$}۔
\end{prop}
\olpseganchor{OLP-0164-B007}{end}

\olpseganchor{OLP-0164-B008}{start}
\begin{proof}
د~\LR{$t$} پر پېچلتيا د استقرا له لارې۔ په بنسټيز حالت کښې~\LR{$t$} يا
ثابت سمبول وي يا له \LR{$x_1$}، \dots،~\LR{$x_n$} متغيرونو څخه يو وي۔ که
\LR{$t = c$} وي، نو
\LR{$\Value{t}{M}[s_1] = \Assign{c}{M} = \Value{t}{M}[s_2]$}۔ که
\LR{$t = x_i$} وي، د قضيې د فرض له مخې \LR{$s_1(x_i) = s_2(x_i)$}؛ نو
\LR{$\Value{t}{M}[s_1] = s_1(x_i) = s_2(x_i) = \Value{t}{M}[s_2]$}۔
\olpseganchor{OLP-0164-B008}{end}

\olpseganchor{OLP-0164-B009}{start}
د استقرايي ګام دپاره فرض کړئ چې
\LR{$t = \Atom{f}{t_1, \dots, t_k}$}، او ادعا د
\LR{$t_1$}، \dots، \LR{$t_k$} دپاره صدق کوي۔ نو
\begin{align*}
  \Value{t}{M}[s_1] & = \Value{\Atom{f}{t_1, \dots, t_k}}{M}[s_1] \\
  & = \Assign{f}{M}(\Value{t_1}{M}[s_1], \dots, \Value{t_k}{M}[s_1]).
\intertext{\RL{د \LR{$j = 1$}، \dots،~\LR{$k$} دپاره، د~\LR{$t_j$} متغيرونه د
    \LR{$x_1$}، \dots،~\LR{$x_n$} په منځ کښې دي۔ د استقرايي فرض له مخې
    \LR{$\Value{t_j}{M}[s_1] = \Value{t_j}{M}[s_2]$}۔ نو}}
\Value{t}{M}[s_1] & = \Value{\Atom{f}{t_1, \dots, t_k}}{M}[s_1] \\
  & = \Assign{f}{M}(\Value{t_1}{M}[s_1], \dots, \Value{t_k}{M}[s_1]) \\
  & = \Assign{f}{M}(\Value{t_1}{M}[s_2], \dots, \Value{t_k}{M}[s_2]) \\
  & = \Value{\Atom{f}{t_1, \dots, t_k}}{M}[s_2] = \Value{t}{M}[s_2].
\end{align*}
\end{proof}
\olpseganchor{OLP-0164-B009}{end}

\olpseganchor{OLP-0164-B010}{start}
\begin{prop}\ollabel{prop:satindep}
که په~\LR{$!A$} کښې ازاد متغيرونه د \LR{$x_1$}، \dots،~\LR{$x_n$} په منځ کښې
وي، او \LR{$s_1(x_i) = s_2(x_i)$} د \LR{$i = 1$}، \dots،~\LR{$n$} دپاره وي، نو
\LR{$\Sat{M}{!A}[s_1]$} هله او يوازې هله چې \LR{$\Sat{M}{!A}[s_2]$}۔
\end{prop}
\olpseganchor{OLP-0164-B010}{end}

\olpseganchor{OLP-0164-B011}{start}
\begin{proof}
د~\LR{$!A$} پر پېچلتيا استقرا کاروو۔ په بنسټيز حالت کښې، چې~\LR{$!A$} اټومي
وي، \LR{$!A$} له لاندې څخه يو کېدے شي:
\LR{$\ltrue$}،
\LR{$\lfalse$}،
د يوه \LR{$k$}-ځايه اړوند~\LR{$R$} او ترمونو \LR{$t_1$}، \dots،~\LR{$t_k$} دپاره
\LR{$\Atom{R}{t_1, \dots, t_k}$}، يا د ترمونو~\LR{$t_1$} او~\LR{$t_2$} دپاره
\LR{$\eq[t_1][t_2]$}۔ په وروستيو دوو حالتونو کښې يوازې د
دوه اړخيز شرطي مخامخ لور ثابتوو، ځکه د شاتني لور ثبوت ئې هم‌شکله
دے۔
\olpseganchor{OLP-0164-B011}{end}

\olpseganchor{OLP-0164-B012}{start}
\begin{enumerate}
\item %
  \indcase{!A}{\ltrue}{دواړه \LR{$\Sat{M}{\indfrm}[s_1]$} او
    \LR{$\Sat{M}{\indfrm}[s_2]$} صدق کوي۔}
\olpseganchor{OLP-0164-B012}{end}

\olpseganchor{OLP-0164-B013}{start}
\item %
  \indcase{!A}{\lfalse}{دواړه \LR{$\Sat/{M}{!A}[s_1]$} او
    \LR{$\Sat/{M}{!A}[s_2]$}۔}
\olpseganchor{OLP-0164-B013}{end}

\olpseganchor{OLP-0164-B014}{start}
\item
  \indcase{!A}{\Atom{R}{t_1, \ldots, t_k}}{فرض کړئ
    \LR{$\Sat{M}{\indfrm}[s_1]$}۔ نو
    \[
    \langle \Value{t_1}{M}[s_1], \ldots, \Value{t_k}{M}[s_1] \rangle
    \in \Assign{R}{M}.
    \]
    د \LR{$i = 1$}، \dots،~\LR{$k$} دپاره، د \olref{prop:valindep} له مخې
    \LR{$\Value{t_i}{M}[s_1] = \Value{t_i}{M}[s_2]$}۔ نو دا هم لرو چې
    \LR{$\langle \Value{t_1}{M}[s_2], \ldots, \Value{t_k}{M}[s_2] \rangle
    \in \Assign{R}{M}$}؛ ځکه \LR{$\Sat{M}{\indfrm}[s_2]$}۔}
\item
  \indcase{!A}{\eq[t_1][t_2]}{فرض کړئ \LR{$\Sat{M}{\indfrm}[s_1]$}۔
    بيا \LR{$\Value{t_1}{M}[s_1] = \Value{t_2}{M}[s_1]$}۔ نو
    \begin{align*}
      \Value{t_1}{M}[s_2] & = \Value{t_1}{M}[s_1]
      & \text{\RL{(د \olref{prop:valindep} له مخې)}} \\
      & = \Value{t_2}{M}[s_1]
      & \text{\RL{(ځکه چې \LR{$\Sat{M}{\eq[t_1][t_2]}[s_1]$})}}\\
      &= \Value{t_2}{M}[s_2]
      & \text{\RL{(د \olref{prop:valindep} له مخې)،}}
    \end{align*}
    نو \LR{$\Sat{M}{\eq[t_1][t_2]}[s_2]$}۔}
\end{enumerate}
\olpseganchor{OLP-0164-B014}{end}

\olpseganchor{OLP-0164-B015}{start}
اوس فرض کړئ چې د~\LR{$!A$} څخه د لږ پېچليو ټولو فارمولونه~\LR{$!B$} دپاره
\LR{$\Sat{M}{!B}[s_1]$} هله او يوازې هله چې \LR{$\Sat{M}{!B}[s_2]$}۔ استقرايي
ګام د~\LR{$!A$} د اصلي چلوونکي له مخې په حالتونو وېشل کېږي۔ په هر
حالت کښې يوازې د دوه اړخيز شرطي مخامخ لور ثابتوو؛ د شاتني لور
ثبوت ئې هم‌شکله دے۔ د سورو له حالتونو پرته، په ټولو حالتونو کښې
استقرايي فرض د~\LR{$!A$} پر فرعي فارمولونه~\LR{$!B$} کاروو۔ د~\LR{$!B$} ازاد
متغيرونه د~\LR{$!A$} د ازادو متغيرونو په منځ کښې دي۔ نو که~\LR{$s_1$} او~\LR{$s_2$}
د~\LR{$!A$} پر ازادو متغيرونو سره موافقې وي، د~\LR{$!B$} پر هغو هم موافقې دي،
او استقرايي فرض پر~\LR{$!B$} کارول کېدے شي۔
\olpseganchor{OLP-0164-B015}{end}

\olpseganchor{OLP-0164-B016}{start}
\begin{enumerate}
\item %
  %
    \indcase{!A}{\lnot !B}{که \LR{$\Sat{M}{\indfrm}[s_1]$} وي، نو
      \LR{$\Sat/{M}{!B}[s_1]$}؛ د استقرايي فرض له مخې
      \LR{$\Sat/{M}{!B}[s_2]$}، ځکه \LR{$\Sat{M}{\indfrm}[s_2]$}۔}
\olpseganchor{OLP-0164-B016}{end}

\olpseganchor{OLP-0164-B017}{start}
\item %
  %
    \indcase{!A}{!B \land !C}{که \LR{$\Sat{M}{\indfrm}[s_1]$} وي، نو
      \LR{$\Sat{M}{!B}[s_1]$} او \LR{$\Sat{M}{!C}[s_1]$}؛ د استقرايي فرض له
      مخې \LR{$\Sat{M}{!B}[s_2]$} او \LR{$\Sat{M}{!C}[s_2]$}۔ ځکه
      \LR{$\Sat{M}{\indfrm}[s_2]$}۔}
\olpseganchor{OLP-0164-B017}{end}

\olpseganchor{OLP-0164-B018}{start}
\item %
  %
    \indcase{!A}{!B \lor !C}{که \LR{$\Sat{M}{\indfrm}[s_1]$} وي، نو
      \LR{$\Sat{M}{!B}[s_1]$} يا \LR{$\Sat{M}{!C}[s_1]$}۔ د استقرايي فرض له
      مخې \LR{$\Sat{M}{!B}[s_2]$} يا \LR{$\Sat{M}{!C}[s_2]$}؛ نو
      \LR{$\Sat{M}{\indfrm}[s_2]$}۔}
\olpseganchor{OLP-0164-B018}{end}

\olpseganchor{OLP-0164-B019}{start}
\item %
  %
    \indcase{!A}{!B \lif !C}{که \LR{$\Sat{M}{\indfrm}[s_1]$} وي، نو
    \LR{$\Sat/{M}{!B}[s_1]$} يا \LR{$\Sat{M}{!C}[s_1]$}۔ د استقرايي فرض له مخې
    \LR{$\Sat/{M}{!B}[s_2]$} يا \LR{$\Sat{M}{!C}[s_2]$}؛ نو
    \LR{$\Sat{M}{\indfrm}[s_2]$}۔}
\olpseganchor{OLP-0164-B019}{end}

\olpseganchor{OLP-0164-B020}{start}
\item %
  %
    \indcase{!A}{!B \liff !C}{که \LR{$\Sat{M}{\indfrm}[s_1]$} وي، نو يا
      \LR{$\Sat{M}{!B}[s_1]$} او \LR{$\Sat{M}{!C}[s_1]$} دواړه، يا
      \LR{$\Sat/{M}{!B}[s_1]$} او \LR{$\Sat/{M}{!C}[s_1]$} دواړه صدق نۀ کوي۔
      د استقرايي فرض له مخې، يا \LR{$\Sat{M}{!B}[s_2]$} او
      \LR{$\Sat{M}{!C}[s_2]$} دواړه، يا \LR{$\Sat/{M}{!B}[s_2]$} او
      \LR{$\Sat/{M}{!C}[s_2]$} دواړه صدق نۀ کوي۔ په هر حالت کښې
      \LR{$\Sat{M}{\indfrm}[s_2]$}۔}
\olpseganchor{OLP-0164-B020}{end}

\olpseganchor{OLP-0164-B021}{start}
\item %
  %
    \indcase{!A}{\lexists[x][!B]}{که \LR{$\Sat{M}{\indfrm}[s_1]$} وي، داسې
      \LR{$m \in \Domain{M}$} شته چې \LR{$\Sat{M}{!B}[\Subst{s_1}{m}{x}]$}۔
      \LR{$s_1' = \Subst{s_1}{m}{x}$} او \LR{$s_2' =
      \Subst{s_2}{m}{x}$} ونيسئ۔ د~\LR{$!B$} ازاد متغيرونه د \LR{$x_1$}،
      \dots، \LR{$x_n$} او~\LR{$x$} په منځ کښې دي۔ \LR{$s_1'(x_i) = s_2'(x_i)$}،
      ځکه \LR{$s_1'$} او~\LR{$s_2'$} په ترتيب سره د~\LR{$s_1$} او~\LR{$s_2$}
      \LR{$x$}-ډولونه دي، او د فرض له مخې \LR{$s_1(x_i) = s_2(x_i)$}۔
      د~\LR{$s_1'$} او~\LR{$s_2'$} د تعريف له مخې
      \LR{$s_1'(x) = s_2'(x) = m$}۔ نو استقرايي فرض پر~\LR{$!B$} او
      \LR{$s_1'$}، \LR{$s_2'$} کارول کېږي، او \LR{$\Sat{M}{!B}[s_2']$} لرو۔ څرنګه
      چې \LR{$s_2' = \Subst{s_2}{m}{x}$}، داسې \LR{$m \in \Domain{M}$} شته چې
      \LR{$\Sat{M}{!B}[\Subst{s_2}{m}{x}]$}؛ ځکه
      \LR{$\Sat{M}{\indfrm}[s_2]$}۔}
\olpseganchor{OLP-0164-B021}{end}

\olpseganchor{OLP-0164-B022}{start}
\item %
  %
    \indcase{!A}{\lforall[x][!B]}{که \LR{$\Sat{M}{\indfrm}[s_1]$} وي، نو
      د هر \LR{$m \in \Domain{M}$} دپاره
      \LR{$\Sat{M}{!B}[\Subst{s_1}{m}{x}]$}۔ غواړو وښيو چې د هر
      \LR{$m \in \Domain{M}$} دپاره
      \LR{$\Sat{M}{!B}[\Subst{s_2}{m}{x}]$} هم صدق کوي۔ يو خپل‌سري
      \LR{$m \in \Domain{M}$} ونيسئ، او
      \LR{$s_1' = \Subst{s_1}{m}{x}$} او \LR{$s_2' =
      \Subst{s_2}{m}{x}$} په پام کښې ونيسئ۔
      \LR{$\Sat{M}{!B}[s_1']$} لرو۔ د~\LR{$!B$} ازاد متغيرونه د \LR{$x_1$}،
      \dots، \LR{$x_n$} او~\LR{$x$} په منځ کښې دي۔
      \LR{$s_1'(x_i) = s_2'(x_i)$}، ځکه \LR{$s_1'$} او~\LR{$s_2'$} په ترتيب سره
      د~\LR{$s_1$} او~\LR{$s_2$} \LR{$x$}-ډولونه دي، او د فرض له مخې
      \LR{$s_1(x_i) = s_2(x_i)$}۔ د~\LR{$s_1'$} او~\LR{$s_2'$} د تعريف له مخې
      \LR{$s_1'(x) = s_2'(x) = m$}۔ نو استقرايي فرض پر~\LR{$!B$} او
      \LR{$s_1'$}، \LR{$s_2'$} کارول کېږي، او \LR{$\Sat{M}{!B}[s_2']$} لرو۔ دا د هر
      \LR{$m \in \Domain{M}$} دپاره صدق کوي، يعنې د ټولو
      \LR{$m \in \Domain{M}$} دپاره
      \LR{$\Sat{M}{!B}[\Subst{s_2}{m}{x}]$}؛ نو
      \LR{$\Sat{M}{\indfrm}[s_2]$}۔}
\end{enumerate}
د استقرا له مخې پايله اخلو چې هر کله د~\LR{$!A$} ازاد متغيرونه د
\LR{$x_1$}، \dots، \LR{$x_n$} په منځ کښې وي او
\LR{$s_1(x_i)=s_2(x_i)$} د \LR{$i = 1$}، \dots،~\LR{$n$} دپاره وي، نو
\LR{$\Sat{M}{!A}[s_1]$} هله او يوازې هله چې \LR{$\Sat{M}{!A}[s_2]$}۔
\end{proof}
\olpseganchor{OLP-0164-B022}{end}



\olpseganchor{OLP-0164-B024}{start}
\begin{explain}
تړلي فارمولونه ازاد متغيرونه نۀ لري؛ نو هرې دوې متغير ګومارنې د هرې
جملې ټولو (صفر) ازادو متغيرونو ته يو شان څيزونه ګوماري۔ له دې کبله
اوس ثابته شوې قضيه وايي چې تړلے فارمول د متغير ګومارنې په نسبت په
جوړښت کښې صدق کوي که نۀ، له هغې ګومارنې څخه په بشپړ ډول خپلواک دي۔
دا حقيقت ثبتوو۔ همدا د لاندې تعريف بنسټ دے، چې د متغير ګومارنې له
يادولو پرته په جوړښت کښې د تړلے فارمول صدق تعريفوي۔
\end{explain}
\olpseganchor{OLP-0164-B024}{end}

\olpseganchor{OLP-0164-B025}{start}
\begin{cor}
\ollabel{cor:sat-sentence}
که~\LR{$!A$} تړلے فارمول او~\LR{$s$} يوه متغير ګومارنه وي، نو
\LR{$\Sat{M}{!A}[s]$} هله او يوازې هله چې د هرې متغير ګومارنې~\LR{$s'$} دپاره
\LR{$\Sat{M}{!A}[s']$}۔
\end{cor}
\olpseganchor{OLP-0164-B025}{end}

\olpseganchor{OLP-0164-B026}{start}
\begin{proof}
  \LR{$s'$} هره متغير ګومارنه ونيسئ۔ څرنګه چې~\LR{$!A$} تړلے فارمول ده، ازاد
  متغيرونه نۀ لري؛ ځکه هره متغير ګومارنه~\LR{$s'$} په څرګند ډول د~\LR{$!A$}
  ټولو ازادو متغيرونو ته هماغه څيزونه ګوماري چې~\LR{$s$} ئې ګوماري۔ نو د
  \olref{prop:satindep} شرط پوره کېږي، او
  \LR{$\Sat{M}{!A}[s]$} هله او يوازې هله چې \LR{$\Sat{M}{!A}[s']$}۔
\end{proof}
\olpseganchor{OLP-0164-B026}{end}

\olpseganchor{OLP-0164-B027}{start}
\begin{defn}
\ollabel{defn:satisfaction}
که~\LR{$!A$} تړلے فارمول وي، وايو چې جوړښت~\LR{$\Struct M$}
په~\LR{$!A$} کښې \emph{صدق کوي}، \LR{$\Sat{M}{!A}$}، هله او يوازې هله چې د
ټولو متغير ګومارنو~\LR{$s$} دپاره \LR{$\Sat{M}{!A}[s]$}۔
\end{defn}
\olpseganchor{OLP-0164-B027}{end}

\olpseganchor{OLP-0164-B028}{start}
که \LR{$\Sat{M}{!A}$} وي، په ساده ډول دا هم وايو چې \emph{\LR{$!A$}
په~\LR{$\Struct{M}$} کښې رښتيا ده۔} د صدق مفهوم په طبيعي ډول له انفرادي
تړلي فارمولونه څخه د تړلي فارمولونه سټونو ته غځېږي۔
\olpseganchor{OLP-0164-B028}{end}

\olpseganchor{OLP-0164-B029}{start}
\begin{defn}
\ollabel{defn:sat}
که~\LR{$\Gamma$} د تړلي فارمولونه يو سټ وي، وايو چې
جوړښت~\LR{$\Struct M$} په~\LR{$\Gamma$} کښې \emph{صدق کوي}،
\LR{$\Sat{M}{\Gamma}$}، هله او يوازې هله چې د ټولو
\LR{$!A \in \Gamma$} دپاره \LR{$\Sat{M}{!A}$}۔
\end{defn}
\olpseganchor{OLP-0164-B029}{end}

\olpseganchor{OLP-0164-B030}{start}
\begin{prop}\ollabel{prop:sentence-sat-true}
  فرض کړئ \LR{$\Struct{M}$} جوړښت، \LR{$!A$} تړلے فارمول، او~\LR{$s$}
  يوه متغير ګومارنه ده۔ \LR{$\Sat{M}{!A}$} هله او يوازې هله چې
  \LR{$\Sat{M}{!A}[s]$}۔
\end{prop}
\olpseganchor{OLP-0164-B030}{end}

\olpseganchor{OLP-0164-B031}{start}
\begin{proof}
تمرین۔
\end{proof}
\olpseganchor{OLP-0164-B031}{end}

\olpseganchor{OLP-0164-B032}{start}
\begin{prob}
\olref[fol][syn][ass]{prop:sentence-sat-true} ثابته کړئ۔
\end{prob}
\olpseganchor{OLP-0164-B032}{end}

\olpseganchor{OLP-0164-B033}{start}
\begin{prop}\ollabel{prop:sat-quant}
  فرض کړئ \LR{$!A(x)$} يوازې~\LR{$x$} ازاد لري، او~\LR{$\Struct M$}
  جوړښت دے۔ نو:
  \begin{enumerate}
    \item \LR{$\Sat{M}{\lexists[x][!A(x)]}$} هله او يوازې هله چې
      لږ تر لږه د يوې متغير ګومارنې~\LR{$s$} دپاره
      \LR{$\Sat{M}{!A(x)}[s]$}۔
    \item \LR{$\Sat{M}{\lforall[x][!A(x)]}$} هله او يوازې هله چې
      د ټولو متغير ګومارنو~\LR{$s$} دپاره \LR{$\Sat{M}{!A(x)}[s]$}۔
\end{enumerate}
\end{prop}
\olpseganchor{OLP-0164-B033}{end}

\olpseganchor{OLP-0164-B034}{start}
\begin{proof}
تمرین۔
\end{proof}
\olpseganchor{OLP-0164-B034}{end}

\olpseganchor{OLP-0164-B035}{start}
\begin{prob}
\olref[fol][syn][ass]{prop:sat-quant} ثابته کړئ۔
\end{prob}
\olpseganchor{OLP-0164-B035}{end}

\olpseganchor{OLP-0164-B036}{start}
\begin{prob}
\DeclareRobustCommand{\VDash}{\mathrel{||}\joinrel\Relbar}
فرض کړئ~\LR{$\Lang L$} له تابعې پرته ژبه ده۔ د
جوړښت~\LR{$\Struct M$}، ثابت سمبول~\LR{$c$}، او
\LR{$a \in \Domain M$} په لرلو سره، \LR{$\Struct M[a/c]$} هغه جوړښت
تعريف کړئ چې کټ مټ د~\LR{$\Struct M$} په شان دے، خو
\LR{$\Assign{c}{M[a/c]} = a$}۔ د تړلي فارمولونه~\LR{$!A$} دپاره
\LR{$\Struct M \VDash !A$} داسې تعريف کړئ:
\begin{enumerate}
\item %
  \indcase{!A}{\lfalse}{\LR{$\Struct{M} \VDash \indfrm$} نۀ۔}
\olpseganchor{OLP-0164-B036}{end}

\olpseganchor{OLP-0164-B037}{start}
\item %
  \indcase{!A}{\ltrue}{\LR{$\Struct{M} \VDash \indfrm$}۔}
\olpseganchor{OLP-0164-B037}{end}

\olpseganchor{OLP-0164-B038}{start}
\item \indcase{!A}{\Atom{R}{d_1, \dots, d_n}}{\LR{$\Struct M \VDash \indfrm$}
  هله او يوازې هله چې
  \LR{$\langle \Assign{d_1}{M}, \dots, \Assign{d_n}{M} \rangle \in
  \Assign{R}{M}$}۔}
\olpseganchor{OLP-0164-B038}{end}

\olpseganchor{OLP-0164-B039}{start}
\item \indcase{!A}{\eq[d_1][d_2]}{\LR{$\Struct M \VDash \indfrm$} هله او
  يوازې هله چې \LR{$\Assign{d_1}{M} = \Assign{d_2}{M}$}۔}
\olpseganchor{OLP-0164-B039}{end}

\olpseganchor{OLP-0164-B040}{start}
\item %
  \indcase{!A}{\lnot !B}{\LR{$\Struct{M} \VDash \indfrm$} هله او يوازې
    هله چې \LR{$\Struct{M} \VDash {!B}$} نۀ وي۔}
\olpseganchor{OLP-0164-B040}{end}

\olpseganchor{OLP-0164-B041}{start}
\item %
  \indcase{!A}{(!B \land !C)}{\LR{$\Struct{M} \VDash
    \indfrm$} هله او يوازې هله چې \LR{$\Struct{M} \VDash!B$} او
    \LR{$\Struct{M} \VDash !C$}۔}
\olpseganchor{OLP-0164-B041}{end}

\olpseganchor{OLP-0164-B042}{start}
\item %
  \indcase{!A}{(!B \lor !C)}{\LR{$\Struct{M} \VDash \indfrm$} هله او
    يوازې هله چې \LR{$\Struct{M} \VDash !B$} يا
    \LR{$\Struct{M} \VDash !C$} (يا دواړه)۔}
\olpseganchor{OLP-0164-B042}{end}

\olpseganchor{OLP-0164-B043}{start}
\item %
  \indcase{!A}{(!B \lif !C)}{\LR{$\Struct{M} \VDash \indfrm$} هله او
    يوازې هله چې \LR{$\Struct {M} \VDash !B$} نۀ وي يا
    \LR{$\Struct M \VDash !C$} وي (يا دواړه)۔}
\olpseganchor{OLP-0164-B043}{end}

\olpseganchor{OLP-0164-B044}{start}
\item %
  \indcase{!A}{(!B \liff !C)}{\LR{$\Struct{M} \VDash {\indfrm}$} هله او
    يوازې هله چې يا \LR{$\Struct{M} \VDash {!B}$} او
    \LR{$\Struct{M} \VDash {!C}$} دواړه وي، يا نۀ
    \LR{$\Struct{M} \VDash {!B}$} وي او نۀ \LR{$\Struct{M} \VDash {!C}$}۔}
\olpseganchor{OLP-0164-B044}{end}

\olpseganchor{OLP-0164-B045}{start}
\item %
  \indcase{!A}{\lforall[x][!B]}{\LR{$\Struct{M} \VDash {\indfrm}$} هله او
    يوازې هله چې د ټولو \LR{$a \in \Domain{M}$} دپاره
    \LR{$\Struct{M[a/c]} \VDash \Subst{!B}{c}{x}$}، په دې شرط چې~\LR{$c$}
    په~\LR{$!B$} کښې نۀ وي راغلے۔}
\olpseganchor{OLP-0164-B045}{end}

\olpseganchor{OLP-0164-B046}{start}
\item %
  \indcase{!A}{\lexists[x][!B]}{\LR{$\Struct{M} \VDash
  {\indfrm}$} هله او يوازې هله چې داسې \LR{$a \in \Domain M$} وي چې
  \LR{$\Struct{M[a/c]} \VDash \Subst{!B}{c}{x}$}، په دې شرط چې~\LR{$c$}
  په~\LR{$!B$} کښې نۀ وي راغلے۔}
\end{enumerate}
فرض کړئ \LR{$x_1$}، \dots، \LR{$x_n$} په~\LR{$!A$} کښې ټول ازاد متغيرونه،
\LR{$c_1$}، \dots، \LR{$c_n$} هغه ثابت سمبولونه دي چې په~\LR{$!A$} کښې نۀ دي،
\LR{$a_1$}، \dots، \LR{$a_n \in \Domain M$}، او \LR{$s(x_i) = a_i$}۔
\olpseganchor{OLP-0164-B046}{end}

\olpseganchor{OLP-0164-B047}{start}
وښيئ چې \LR{$\Sat{M}{!A}[s]$} هله او يوازې هله چې
\LR{$\Struct M[a_1/c_1,\dots,a_n/c_n]
\VDash \Subst{\Subst{!A}{c_1}{x_1}\dots}{c_n}{x_n}$}۔
\olpseganchor{OLP-0164-B047}{end}

\olpseganchor{OLP-0164-B048}{start}
(دا مسئله ښيي چې د لومړۍ درجې منطق دپاره داسې معنٰی پوهنه ورکول
ممکن دي چې له متغير ګومارنو پرته کار وکړي۔)
\end{prob}
\olpseganchor{OLP-0164-B048}{end}

\olpseganchor{OLP-0164-B049}{start}
\begin{prob}
فرض کړئ~\LR{$f$} يو داسې تابعي سمبول دے چې په~\LR{$!A(x,y)$} کښې نۀ وي۔ وښيئ
چې داسې جوړښت~\LR{$\Struct{M}$} شته چې
\LR{$\Sat{M}{\lforall[x][\lexists[y][!A(x,y)]]}$} هله او يوازې هله چې داسې
\LR{$\Struct M'$} شته چې \LR{$\Sat{M'}{\lforall[x][!A(x,f(x))]}$}۔
\olpseganchor{OLP-0164-B049}{end}

\olpseganchor{OLP-0164-B050}{start}
(دا مسئله د سکولم د قضيې په نامه د پېژندل شوې قضيې يو ځانګړے حالت
دے؛ \LR{$\lforall[x][!A(x,f(x))]$} ته د
\LR{$\lforall[x][\lexists[y][!A(x,y)]]$} \emph{د سکولم نورماله بڼه} ويل
کېږي۔)
\end{prob}
\olpseganchor{OLP-0164-B050}{end}

\olpseganchor{OLP-0165-B004}{start}
\olfileid{fol}{syn}{ext}
\olpseganchor{OLP-0165-B004}{end}

\olpseganchor{OLP-0165-B005}{start}
\olsection{امتداديتوب}
\olpseganchor{OLP-0165-B005}{end}

\olpseganchor{OLP-0165-B006}{start}
\begin{explain}
امتداديتوب، چې کله کله اړوندتيا هم بلل کېږي، په غير رسمي ډول داسې
څرګندولے شو: يوازېني عوامل چې د فارمول~\LR{$!A$} پر صدق د
جوړښت~\LR{$\Struct M$} او د متغير ګومارنې~\LR{$s$} په نسبت
اغېز لري، د د تعريف ساحه کچه او هغه ګومارنې دي چې~\LR{$\Struct M$} او~\LR{$s$} ئې
د ژبې هغو اجزاوو ته کوي چې په رښتيا په~\LR{$!A$} کښې راځي۔
\olpseganchor{OLP-0165-B006}{end}

\olpseganchor{OLP-0165-B007}{start}
د امتداديتوب يوه سمدستي پايله دا ده: که دوه
جوړښتونه~\LR{$\Struct M$} او~\LR{$\Struct M'$} د ژبې پر ټولو هغو اجزاوو
سره موافق وي چې په جمله~\LR{$!A$} کښې راځي، او يوه ساحه ولري، نو
\LR{$\Struct M$} او~\LR{$\Struct M'$} بايد پر دې هم موافق وي چې پخپله~\LR{$!A$}
رښتيا ده که نۀ۔
\end{explain}
\olpseganchor{OLP-0165-B007}{end}

\olpseganchor{OLP-0165-B008}{start}
\begin{prop}[امتداديتوب]
  \ollabel{prop:extensionality}
  فرض کړئ~\LR{$!A$} فارمول، او~\LR{$\Struct M_1$} او~\LR{$\Struct M_2$} داسې
  جوړښتونه دي چې \LR{$\Domain{M_1} = \Domain{M_2}$}، او~\LR{$s$} پر
  \LR{$\Domain{M_1} = \Domain{M_2}$} يوه متغير ګومارنه ده۔ که په~\LR{$!A$} کښې
  د راتلونکو هر ثابت سمبول~\LR{$c$}، اړوند سمبول~\LR{$R$}، او
  تابع~\LR{$f$} دپاره
  \LR{$\Assign{c}{M_1} = \Assign{c}{M_2}$}،
  \LR{$\Assign{R}{M_1}=\Assign{R}{M_2}$}، او
  \LR{$\Assign{f}{M_1} = \Assign{f}{M_2}$} وي، نو
  \LR{$\Sat{M_1}{!A}[s]$} هله او يوازې هله چې \LR{$\Sat{M_2}{!A}[s]$}۔
\end{prop}
\olpseganchor{OLP-0165-B008}{end}

\olpseganchor{OLP-0165-B009}{start}
\begin{proof}
  لومړے (پر~\LR{$t$} د استقرا له لارې) ثابته کړئ چې د هر ترم دپاره
  \LR{$\Value{t}{M_1}[s] = \Value{t}{M_2}[s]$}۔ بيا قضيه پر~\LR{$!A$} د استقرا
  له لارې ثابته کړئ، او د استقرا په بنسټ کښې (چې~\LR{$!A$} اټومي وي) له
  همدې ثابتې شوې ادعا څخه کار واخلئ۔
\end{proof}
\olpseganchor{OLP-0165-B009}{end}

\olpseganchor{OLP-0165-B010}{start}
\begin{prob}
د \olref[fol][syn][ext]{prop:extensionality} ثبوت په تفصيل بشپړ کړئ۔
\end{prob}
\olpseganchor{OLP-0165-B010}{end}

\olpseganchor{OLP-0165-B011}{start}
\begin{cor}[د تړلي فارمولونه دپاره امتداديتوب]
  \ollabel{cor:extensionality-sent}
  فرض کړئ~\LR{$!A$} تړلے فارمول، او~\LR{$\Struct{M_1}$}، \LR{$\Struct{M_2}$} لکه
  په \olref{prop:extensionality} کښې وي۔ نو \LR{$\Sat{M_1}{!A}$} هله او
  يوازې هله چې \LR{$\Sat{M_2}{!A}$}۔
\end{cor}
\olpseganchor{OLP-0165-B011}{end}

\olpseganchor{OLP-0165-B012}{start}
\begin{proof}
له \olref{prop:extensionality} او \olref[ass]{cor:sat-sentence} څخه
لاس ته راځي۔
\end{proof}
\olpseganchor{OLP-0165-B012}{end}

\olpseganchor{OLP-0165-B013}{start}
پر دې سربېره، د يوه ترم ارزښت، او دا چې جوړښت په
فارمول کښې صدق کوي که نۀ، يوازې د هغه د فرعي ترمونو په
ارزښتونو پورې تړلي دي۔
\olpseganchor{OLP-0165-B013}{end}

\olpseganchor{OLP-0165-B014}{start}
\begin{prop}\ollabel{prop:ext-terms}
فرض کړئ~\LR{$\Struct M$} جوړښت، \LR{$t$} او~\LR{$t'$} ترمونه، او~\LR{$s$} يوه
متغير ګومارنه ده۔ نو \LR{$\Value{\Subst{t}{t'}{x}}{M}[s] =
\Value{t}{M}[\Subst{s}{\Value{t'}{M}[s]}{x}]$}۔
\end{prop}
\olpseganchor{OLP-0165-B014}{end}

\olpseganchor{OLP-0165-B015}{start}
\begin{proof}
پر~\LR{$t$} د استقرا له لارې۔
\begin{enumerate}
\item که~\LR{$t$} يو ثابت وي، يعنې \LR{$t\ident c$}، نو
  \LR{$\Subst{t}{t'}{x} = c$}، او \LR{$\Value{c}{M}[s] = \Assign{c}{M} =
  \Value{c}{M}[\Subst{s}{\Value{t'}{M}[s]}{x}]$}۔
\olpseganchor{OLP-0165-B015}{end}

\olpseganchor{OLP-0165-B016}{start}
\item که~\LR{$t$} له~\LR{$x$} پرته بل متغير وي، يعنې \LR{$t \ident y$}، نو
  \LR{$\Subst{t}{t'}{x} = y$}، او \LR{$\Value{y}{M}[s] =
  \Value{y}{M}[\Subst{s}{\Value{t'}{M}[s]}{x}]$}، ځکه
  \LR{$\varAssign{s}{\Subst{s}{\Value{t'}{M}[s]}{x}}{x}$}۔
\olpseganchor{OLP-0165-B016}{end}

\olpseganchor{OLP-0165-B017}{start}
\item که \LR{$t \ident x$} وي، نو \LR{$\Subst{t}{t'}{x} = t'$}۔ خو د
  \LR{$\Subst{s}{\Value{t'}{M}[s]}{x}$} د تعريف له مخې
  \LR{$\Value{x}{M}[\Subst{s}{\Value{t'}{M}[s]}{x}] = \Value{t'}{M}[s]$}۔
\olpseganchor{OLP-0165-B017}{end}

\olpseganchor{OLP-0165-B018}{start}
\item که \LR{$t \ident \Atom{f}{t_1,\dots,t_n}$} وي، نو:
\begin{multline*}
  \Value{\Subst{t}{t'}{x}}{M}[s]  \\
  \begin{aligned}[b]
& = \Value{\Atom{f}{\Subst{t_1}{t'}{x}, \dots, \Subst{t_n}{t'}{x}}}{M}[s]\\
& \qquad    \text{\RL{د \LR{$\Subst{t}{t'}{x}$} د تعريف له مخې}}\\
& = \Assign{f}{M}(\Value{\Subst{t_1}{t'}{x}}{M}[s], \dots,
    \Value{\Subst{t_n}{t'}{x}}{M}[s])\\
    & \qquad  \text{\RL{د \LR{$\Value{\Atom{f}{\dots}}{M}[s]$} د تعريف له مخې}}\\
& = \Assign{f}{M}(\Value{t_1}{M}[\Subst{s}{\Value{t'}{M}[s]}{x}], \dots,
   \Value{t_n}{M}[\Subst{s}{\Value{t'}{M}[s]}{x}])\\
& \qquad    \text{\RL{د استقرايي فرض له مخې}}\\
& = \Value{t}{M}[\Subst{s}{\Value{t'}{M}[s]}{x}]
    \text{\RL{د \LR{$\Value{\Atom{f}{\dots}}{M}[\Subst{s}{\Value{t'}{M}[s]}{x}]$} د تعريف له مخې}}
  \end{aligned}
\end{multline*}
\end{enumerate}
\end{proof}
\olpseganchor{OLP-0165-B018}{end}

\olpseganchor{OLP-0165-B019}{start}
\begin{prop}\ollabel{prop:ext-formulas}
فرض کړئ~\LR{$\Struct M$} جوړښت، \LR{$!A$} فارمول، \LR{$t'$} يو ترم،
او~\LR{$s$} يوه متغير ګومارنه ده۔ نو
\LR{$\Sat{M}{\Subst{!A}{t'}{x}}[s]$} هله او يوازې هله چې
\LR{$\Sat{M}{!A}[\Subst{s}{\Value{t'}{M}[s]}{x}]$}۔
\end{prop}
\olpseganchor{OLP-0165-B019}{end}

\olpseganchor{OLP-0165-B020}{start}
\begin{proof}
تمرین۔
\end{proof}
\olpseganchor{OLP-0165-B020}{end}

\olpseganchor{OLP-0165-B021}{start}
\begin{prob}
\olref[fol][syn][ext]{prop:ext-formulas} ثابته کړئ۔
\end{prob}
\olpseganchor{OLP-0165-B021}{end}

\olpseganchor{OLP-0165-B022}{start}
\begin{explain}
  د \cref{fol:syn:ext:prop:ext-terms,fol:syn:ext:prop:ext-formulas}
  مطلب دا دے۔ فرض کړئ يو ترم~\LR{$t$} يا فارمول~\LR{$!A$}، او بل ترم~\LR{$t'$}
  لرو، او غواړو د~\LR{$\Subst{t}{t'}{x}$} ارزښت، يا دا معلوم کړو چې
  \LR{$\Subst{!A}{t'}{x}$} د متغير ګومارنې~\LR{$s$} په نسبت په
  جوړښت~\LR{$\Struct M$} کښې صدق کوي که نۀ۔ يوې لارې کښې لومړے
  تعويض ترسره کولے، او بيا د~\LR{$\Struct{M}$} او~\LR{$s$} په نسبت ارزښت يا
  صدق څېړلے شو۔ بلې لارې کښې لومړے په~\LR{$\Struct{M}$} کښې د~\LR{$s$} په نسبت
  د~\LR{$t'$} ارزښت~\LR{$m = \Value{t'}{M}[s]$} معلومولے، بيا د
  متغير ګومارنه په~\LR{$\Subst{s}{m}{x}$} بدلولے، او ورپسې په
  \LR{$\Struct{M}$} کښې د~\LR{$\Subst{s}{m}{x}$} په نسبت د~\LR{$t$} ارزښت، يا
  \LR{$\Sat{M}{!A}[\Subst{s}{m}{x}]$} څېړلے شو۔
  \Cref{fol:syn:ext:prop:ext-terms,fol:syn:ext:prop:ext-formulas}
  تضمينوي چې په دواړو لارو کښې ځواب يو شان وي۔
\end{explain}
\olpseganchor{OLP-0165-B022}{end}

\olpseganchor{OLP-0166-B004}{start}
\olfileid{fol}{syn}{sem}
\olsection{معنايي مفهومونه}
\olpseganchor{OLP-0166-B004}{end}

\olpseganchor{OLP-0166-B005}{start}
\begin{explain}
د لومړۍ درجې ژبو د جوړښتونه د تعريف په لرلو سره، د جملو ځينې
بنسټيز معنايي خاصيتونه او د هغو تر منځ اړيکې تعريفولے شو۔ تر ټولو
ساده ئې د جملې د \emph{اعتبار} مفهوم دے۔ يوه جمله هله معتبره ده چې
په هر جوړښت کښې صدق وکړي۔ معتبرې جملې هغه دي چې د هغو د
نامنطقي سمبولونو د تفسير له څرنګوالي پرته صدق کوي۔ ځکه معتبرې جملې
\emph{منطقي حقيقتونه} هم بلل کېږي---هغوی په هر جوړښت کښې
رښتيا دي، يعنې صدق کوي؛ نو رښتينولي ئې يوازې په هغو کښې پر راغلو
منطقي سمبولونو او د هغو پر نحوي جوړښت پورې تړلې ده، پر
نامنطقي سمبولونو يا د هغو پر تفسير نۀ۔
\end{explain}
\olpseganchor{OLP-0166-B005}{end}

\olpseganchor{OLP-0166-B006}{start}
\begin{defn}[اعتبار]
جمله~\LR{$!A$} هله \emph{معتبره} ده، \LR{$\Entails !A$}، چې د هر
جوړښت~\LR{$\Struct M$} دپاره \LR{$\Sat{M}{!A}$}۔
\end{defn}
\olpseganchor{OLP-0166-B006}{end}

\olpseganchor{OLP-0166-B007}{start}
\begin{defn}[استلزام]
د جملو سټ~\LR{$\Gamma$} هله جمله~\LR{$!A$} \emph{لازموي}،
\LR{$\Gamma \Entails !A$}، چې د هر هغه جوړښت~\LR{$\Struct M$} دپاره
چې \LR{$\Sat{M}{\Gamma}$} وي، \LR{$\Sat{M}{!A}$} هم وي۔
\end{defn}
\olpseganchor{OLP-0166-B007}{end}


\olpseganchor{OLP-0166-B008}{start}
\begin{defn}[د صدق وړتيا]
د جملو سټ~\LR{$\Gamma$} هله \emph{د صدق وړ} دے چې د کوم
جوړښت~\LR{$\Struct M$} دپاره \LR{$\Sat{M}{\Gamma}$} وي۔ که~\LR{$\Gamma$} د
صدق وړ نۀ وي، \emph{د صدق نۀ وړ} بلل کېږي۔
\end{defn}
\olpseganchor{OLP-0166-B008}{end}

\olpseganchor{OLP-0166-B009}{start}
\begin{prop}
جمله~\LR{$!A$} هله او يوازې هله معتبره ده چې د جملو د هر سټ~\LR{$\Gamma$}
دپاره \LR{$\Gamma \Entails !A$}۔
\end{prop}
\olpseganchor{OLP-0166-B009}{end}

\olpseganchor{OLP-0166-B010}{start}
\begin{proof}
د مخامخ لور دپاره، فرض کړئ~\LR{$!A$} معتبره ده، او~\LR{$\Gamma$} د جملو يو
سټ دے۔ \LR{$\Struct M$} داسې جوړښت ونيسئ چې
\LR{$\Sat{M}{\Gamma}$}۔ څرنګه چې~\LR{$!A$} معتبره ده، \LR{$\Sat{M}{!A}$}؛ نو
\LR{$\Gamma \Entails !A$}۔
\olpseganchor{OLP-0166-B010}{end}

\olpseganchor{OLP-0166-B011}{start}
د شاتني لور د عکس نقيض دپاره، فرض کړئ~\LR{$!A$} نامعتبره ده؛ نو داسې
جوړښت~\LR{$\Struct M$} شته چې \LR{$\Sat/{M}{!A}$}۔ کله چې
\LR{$\Gamma = \{ \ltrue \}$} وي، د~\LR{$\ltrue$} د اعتبار له کبله
\LR{$\Sat{M}{\Gamma}$}۔ ځکه داسې جوړښت~\LR{$\Struct M$} شته چې
\LR{$\Sat{M}{\Gamma}$} خو \LR{$\Sat/{M}{!A}$}؛ نو~\LR{$\Gamma$}، \LR{$!A$} نۀ لازموي۔
\end{proof}
\olpseganchor{OLP-0166-B011}{end}

\olpseganchor{OLP-0166-B012}{start}
\begin{prop}
\ollabel{prop:entails-unsat}
\LR{$\Gamma \Entails !A$} هله او يوازې هله چې
\LR{$\Gamma \cup \{\lnot !A\}$} د صدق نۀ وړ وي۔
\end{prop}
\olpseganchor{OLP-0166-B012}{end}

\olpseganchor{OLP-0166-B013}{start}
\begin{proof}
د مخامخ لور دپاره، فرض کړئ \LR{$\Gamma \Entails !A$}، او د عکس فرض وکړئ
چې داسې جوړښت~\LR{$\Struct M$} شته چې
\LR{$\Sat{M}{\Gamma \cup \{ \lnot !A \}}$}۔ څرنګه چې
\LR{$\Sat{M}{\Gamma}$} او \LR{$\Gamma \Entails !A$}، نو \LR{$\Sat{M}{!A}$}۔
له بلې خوا، له \LR{$\Sat{M}{\Gamma\cup \{ \lnot !A \}}$} څخه
\LR{$\Sat{M}{\lnot !A}$}؛ نو هم \LR{$\Sat{M}{!A}$} او هم \LR{$\Sat/{M}{!A}$} لرو،
چې تناقض دے۔ ځکه داسې جوړښت~\LR{$\Struct M$} نۀ شي موجودېدے؛ نو
\LR{$\Gamma \cup \{ \lnot !A \}$} د صدق نۀ وړ دے۔
\olpseganchor{OLP-0166-B013}{end}

\olpseganchor{OLP-0166-B014}{start}
د شاتني لور دپاره، فرض کړئ \LR{$\Gamma \cup \{ \lnot !A \}$} د صدق نۀ وړ
دے۔ نو د هر جوړښت~\LR{$\Struct M$} دپاره يا
\LR{$\Sat/{M}{\Gamma}$}، يا \LR{$\Sat{M}{!A}$}۔ ځکه د هر هغه
جوړښت~\LR{$\Struct M$} دپاره چې \LR{$\Sat{M}{\Gamma}$} وي،
\LR{$\Sat{M}{!A}$}؛ نو \LR{$\Gamma \Entails !A$}۔
\end{proof}
\olpseganchor{OLP-0166-B014}{end}

\olpseganchor{OLP-0166-B015}{start}
\begin{prob}
\begin{enumerate}
\item وښيئ چې \LR{$\Gamma \Entails \bot$} هله او يوازې هله چې~\LR{$\Gamma$} د
صدق نۀ وړ وي۔
\item وښيئ چې \LR{$\Gamma \cup \{!A\} \Entails \bot$} هله او يوازې هله چې
\LR{$\Gamma \Entails \lnot !A$}۔
\item فرض کړئ~\LR{$c$} په~\LR{$!A$} يا~\LR{$\Gamma$} کښې نۀ راځي۔ وښيئ چې
  \LR{$\Gamma \Entails \lforall[x][!A]$} هله او يوازې هله چې
  \LR{$\Gamma \Entails \Subst{!A}{c}{x}$}۔
\end{enumerate}
\end{prob}
\olpseganchor{OLP-0166-B015}{end}

\olpseganchor{OLP-0166-B016}{start}
\begin{prop}
که \LR{$\Gamma \subseteq \Gamma'$} او \LR{$\Gamma \Entails !A$} وي، نو
\LR{$\Gamma' \Entails !A$}۔
\end{prop}
\olpseganchor{OLP-0166-B016}{end}

\olpseganchor{OLP-0166-B017}{start}
\begin{proof}
فرض کړئ \LR{$\Gamma \subseteq \Gamma'$} او \LR{$\Gamma \Entails !A$}۔
\LR{$\Struct M$} داسې جوړښت ونيسئ چې \LR{$\Sat{M}{\Gamma'}$}؛ نو
\LR{$\Sat{M}{\Gamma}$}، او له \LR{$\Gamma \Entails !A$} څخه
\LR{$\Sat{M}{!A}$} اخلو۔ ځکه هر کله چې \LR{$\Sat{M}{\Gamma'}$} وي،
\LR{$\Sat{M}{!A}$}؛ نو \LR{$\Gamma' \Entails !A$}۔
\end{proof}
\olpseganchor{OLP-0166-B017}{end}


\olpseganchor{OLP-0166-B018}{start}
\begin{thm}[د معنايي استنتاج قضيه]
\ollabel{thm:sem-deduction}
\LR{$\Gamma \cup \{!A\} \Entails !B$} هله او يوازې هله چې
\LR{$\Gamma \Entails !A \lif !B$}۔
\end{thm}
\olpseganchor{OLP-0166-B018}{end}

\olpseganchor{OLP-0166-B019}{start}
\begin{proof}
د مخامخ لور دپاره، فرض کړئ
\LR{$\Gamma \cup \{ !A \} \Entails !B$}، او~\LR{$\Struct M$} داسې
جوړښت ونيسئ چې \LR{$\Sat{M}{\Gamma}$}۔ که \LR{$\Sat{M}{!A}$} وي، نو
\LR{$\Sat{M}{\Gamma \cup \{ !A \} }$}؛ او څرنګه چې
\LR{$\Gamma \cup \{ !A \}$}، \LR{$!B$} لازموي، \LR{$\Sat{M}{!B}$} اخلو۔ ځکه
\LR{$\Sat{M}{!A \lif !B}$}، نو \LR{$\Gamma \Entails !A \lif !B$}۔
\olpseganchor{OLP-0166-B019}{end}

\olpseganchor{OLP-0166-B020}{start}
د شاتني لور دپاره، فرض کړئ \LR{$\Gamma \Entails !A \lif !B$}، او
\LR{$\Struct M$} داسې جوړښت ونيسئ چې
\LR{$\Sat{M}{\Gamma \cup \{ !A \}}$}۔ بيا \LR{$\Sat{M}{\Gamma}$}؛ نو
\LR{$\Sat{M}{!A \lif !B}$}، او له \LR{$\Sat{M}{!A}$} څخه \LR{$\Sat{M}{!B}$}۔
ځکه هر کله چې \LR{$\Sat{M}{\Gamma \cup \{ !A \} }$} وي،
\LR{$\Sat{M}{!B}$}؛ نو \LR{$\Gamma \cup \{ !A \} \Entails !B$}۔
\end{proof}
\olpseganchor{OLP-0166-B020}{end}

\olpseganchor{OLP-0166-B021}{start}
\begin{prop}\ollabel{prop:quant-terms}
  فرض کړئ~\LR{$\Struct{M}$} جوړښت، \LR{$!A(x)$} داسې فارمول دے
  چې يو ازاد متغير~\LR{$x$} لري، او~\LR{$t$} يو تړلے ترم دے۔ نو:
  \begin{enumerate}
    \item \LR{$!A(t) \Entails \lexists[x][!A(x)]$}
    \item \LR{$\lforall[x][!A(x)] \Entails !A(t)$}
  \end{enumerate}
\end{prop}
\olpseganchor{OLP-0166-B021}{end}

\olpseganchor{OLP-0166-B022}{start}
\begin{proof}
  \begin{enumerate}
    \item %
      فرض کړئ \LR{$\Sat{M}{!A(t)}$}۔ \LR{$s$} داسې
        متغير ګومارنه ونيسئ چې \LR{$s(x) = \Value{t}{M}$}۔ څرنګه چې
        \LR{$!A(t)$} تړلے فارمول ده، \LR{$\Sat{M}{!A(t)}[s]$}۔ د
        \olref[ext]{prop:ext-formulas} له مخې
        \LR{$\Sat{M}{!A(x)}[s]$}۔ د \olref[ass]{prop:sat-quant} له مخې
        \LR{$\Sat{M}{\lexists[x][!A(x)]}$}۔
    \item %
      فرض کړئ
        \LR{$\Sat{M}{\lforall[x][!A(x)]}$}۔ \LR{$s$} داسې متغير ګومارنه ونيسئ
        چې \LR{$s(x) = \Value{t}{M}$}۔ د
        \olref[ass]{prop:sat-quant} له مخې
        \LR{$\Sat{M}{!A(x)}[s]$}۔ د \olref[ext]{prop:ext-formulas} له مخې
        \LR{$\Sat{M}{!A(t)}[s]$}۔ څرنګه چې \LR{$!A(t)$} تړلے فارمول ده، د
        \olref[ass]{prop:sentence-sat-true} له مخې
        \LR{$\Sat{M}{!A(t)}$}۔
  \end{enumerate}
\end{proof}
\olpseganchor{OLP-0166-B022}{end}

\olpseganchor{OLP-0167-B004}{start}
\olchapter{fol}{mat}{تيورۍ او د هغو مدلونه}
\olpseganchor{OLP-0167-B004}{end}

\olpseganchor{OLP-0168-B004}{start}
\olfileid{fol}{mat}{int}
\olsection{پېژندنه}
\olpseganchor{OLP-0168-B004}{end}

\olpseganchor{OLP-0168-B005}{start}
\begin{explain}
د بديهي طريقې وده د ساينس په تاريخ کښې يوه ستره لاسته راوړنه ده،
او د رياضۍ په تاريخ کښې ځانګړے اهميت لري۔ د يوه ډګر بديهي وده ډېرې
پوښتنې روښانوي: دا ډګر د څه په باب دے؟ تر ټولو بنسټيز مفاهيم ئې کوم
دي؟ هغوئ له يو بل سره څنګه تړلي دي؟ ايا د ډګر ټول مفاهيم د دغو
بنسټيزو مفاهيمو له مخې تعريفېدے شي؟ دا مفاهيم د کومو قوانينو پيروي
کوي، او بايد ئې وکړي؟
\olpseganchor{OLP-0168-B005}{end}

\olpseganchor{OLP-0168-B006}{start}
بديهي طريقه او منطق د يو بل دپاره جوړ شوي دي۔ صوري منطق د بديهي
تيوريو د جوړولو، د تيورۍ له بديهي اصولو څخه په کره ټاکلې طريقه د
قضيو د ثابتولو، او د بديهي اصولو پوره کوونکو ټولو نظامونو د خاصيتونو
د منظمې څېړنې وسيلې برابروي۔
\end{explain}
\olpseganchor{OLP-0168-B006}{end}

\olpseganchor{OLP-0168-B007}{start}
\begin{defn}
د تړلي فارمولونه يو سټ~\LR{$\Gamma$} هله او يوازې هله \emph{بند} دے چې
هر کله \LR{$\Gamma \Entails !A$} وي، نو \LR{$!A \in \Gamma$} هم وي۔ د
تړلي فارمولونه د يوه سټ~\LR{$\Gamma$} \emph{بندښت} دا دے:
\LR{$\Setabs{!A}{\Gamma \Entails !A}$}۔
\olpseganchor{OLP-0168-B007}{end}

\olpseganchor{OLP-0168-B008}{start}
وايو چې~\LR{$\Gamma$} د جملو د يوه سټ~\LR{$\Delta$} له خوا \emph{په بديهي
اصولو ټاکل شوے} دے، که~\LR{$\Gamma$} د~\LR{$\Delta$} بندښت وي۔
\end{defn}
\olpseganchor{OLP-0168-B008}{end}

\olpseganchor{OLP-0168-B009}{start}
\begin{explain}
يوه بديهي تيوري د هغو جملو سټ ګڼلے شو چې د هغې د بديهي اصولو
سټ~\LR{$\Delta$} په بديهي اصولو ټاکلے وي۔ په بل عبارت، کله چې د لومړۍ درجې ژبه د هغه
ساينس د بنسټيزو مفاهيمو دپاره غيرمنطقي سمبولونه ولري چې غواړو په
بديهي ډول ئې وڅېړو، او ورسره د تړلي فارمولونه داسې سټ ولرو چې د
ساينس بنسټيز قوانين بيانوي، نو تيوري په دې ژبه کښې د ټولو هغو
تړلي فارمولونه په توګه انځورولے شو چې له بديهي اصولو څخه لازمېږي۔
دا لړۍ له هغو ساده بېلګو څخه، چې يوازې يو بنسټيز مفهوم او ساده
بديهي اصول لري، لکه د جزوي ترتيبونو تيوري، تر نيوټني ميخانيک غوندې
پېچلو تيوريو پورې غځېږي۔
\olpseganchor{OLP-0168-B009}{end}

\olpseganchor{OLP-0168-B010}{start}
لاندې مهم منطقي حقيقتونه دي چې دغه صوري چلند د بديهي طريقې دپاره
دومره ارزښتمن کوي۔ فرض کړئ~\LR{$\Gamma$} د يوې تيورۍ بديهي نظام، يعنې د
جملو يو سټ دے۔
\begin{enumerate}
\item په کره توګه ويلے شو چې يو بديهي نظام د جوړښتونه مطلوبه
  ټولګه کله رانغاړي۔ يعنې که د جوړښتونه يوه ټاکلې ټولګه مو په
  پام کښې وي، بديهي نظام~\LR{$\Gamma$} هغه هله او يوازې هله په برياليتوب
  رانغاړي چې مطلوب جوړښتونه کټ مټ هغه~\LR{$\Struct M$} وي چې
  \LR{$\Sat{M}{\Gamma}$} وي۔
\item کېدے شي په دې کار کښې ځکه پاتې راشو چې داسې~\LR{$\Struct M$} وي
  چې \LR{$\Sat{M}{\Gamma}$}، خو~\LR{$\Struct M$} زموږ له مطلوبو
  جوړښتونه څخه نۀ وي۔ دا ښايي دې ته مو اړ کړي چې داسې بديهي
  اصول ورزيات کړو چې په~\LR{$\Struct M$} کښې رښتوني نۀ وي۔
\item که لږ تر لږه په دې کښې بريالي يو چې~\LR{$\Gamma$} په ټولو مطلوبو
  جوړښتونه کښې رښتونے وي، نو هر کله چې
  \LR{$\Gamma \Entails !A$} وي، جمله~\LR{$!A$} هم په ټولو مطلوبو
  جوړښتونه کښې رښتونې ده۔ نو د منطقي وسيلو (لکه د
  اشتقاق طريقو) په کارولو ښودلے شو چې جملې په ټولو مطلوبو
  جوړښتونه کښې رښتونې دي: بس دا وښيو چې له بديهي اصولو څخه
  لازمېږي۔
\item کله کله مو کوم مطلوب جوړښتونه په پام کښې نۀ وي، بلکې
  له خپلو بديهي اصولو پيل کوو: ځينې بنسټيز مفاهيم اخلو چې غواړو
  ټاکلي قوانين پوره کړي، او دغه قوانين په يوه بديهي نظام کښې ثبتوو۔
  سمدستي دا تصديق کول غواړو چې بديهي اصول يو بل نۀ نقضوي؛ که ئې
  وکړي، هېڅ مفاهيم دغه قوانين نۀ شي پوره کولے او تيوري مو ناسازګاره
  جوړه کړې ده۔ د~\LR{$\Gamma$} د يوه مدل په موندلو تصديق کولے شو چې داسې
  نۀ دي شوي۔ او که تيوري مو مدلونه ولري، په منطقي طريقو ئې څېړلے او
  مدلونه هم جوړولے شو۔
\item د بديهي اصولو استقلال هم يوه مهمه پوښتنه ده۔ کېدے شي يو بديهي
  اصل په رښتيا د نورو پايله وي او له همدې امله زائد وي۔ ثابتولے شو
  چې په~\LR{$\Gamma$} کښې بديهي اصل~\LR{$!A$} زائد دے، که
  \LR{$\Gamma \setminus \{!A\} \Entails !A$} ثابت کړو۔ دا هم ثابتولے شو
  چې بديهي اصل زائد نۀ دے، که وښيو چې
  \LR{$(\Gamma \setminus \{!A\}) \cup \{\lnot !A\}$} د صدق وړ دے۔ د
  بېلګې په توګه، د هندسې د نورو بديهي اصولو څخه د موازي کرښو د اصل
  استقلال همداسې ښودل شوے و۔
\item بله مهمه پوښتنه په يوه تيورۍ کښې د مفاهيمو د تعريفېدنې په باب ده: د
  ژبې ټاکنه دا ټاکي چې د تيورۍ مدلونه له څه جوړ وي۔ خو د تيورۍ هر اړخ
  لازم نۀ دے چې په مدلونو کښې جلا استازيتوب ولري۔ مثلاً هر ترتيب
  \LR{$\le$} ورسره سم سخت ترتيب~\LR{$<$} ټاکي---چې يو راکړل شي، بل ترې تعريفولے
  شو۔ نو د داسې ترتيب لرونکې تيورۍ په مدل کښې د اړوند سخت ترتيب
  \emph{جلا} شتون لازم نۀ دے۔ په عمومي ډول، يوه اړيکه د نورو له مخې
  کله تعريفېدے شي؟ يوه اړيکه د نورو له مخې کله نۀ شي تعريفېدے (او له
  همدې امله بايد د ژبې په بنسټيزو سمبولونو کښې ورزياته شي)؟
\end{enumerate}
\end{explain}
\olpseganchor{OLP-0168-B010}{end}

\olpseganchor{OLP-0169-B004}{start}
\olfileid{fol}{mat}{exs}
\olpseganchor{OLP-0169-B004}{end}

\olpseganchor{OLP-0169-B005}{start}
\olsection{د جوړښتونه د خاصيتونو بيانول}
\olpseganchor{OLP-0169-B005}{end}

\olpseganchor{OLP-0169-B006}{start}
\begin{explain}
د تابعو او اړيکو شرطونه، يا په لا عمومي ډول دا خبره بيانول چې په
يوه جوړښت کښې تابعې او اړيکې دغه شرطونه پوره کوي، ډېر ځله ګټور او
مهم وي۔ د بېلګې په توګه، غواړو د يوې ژبې هغه جوړښتونه چې
پريديکاتونه پکښې زموږ مطلوبه ``معنٰی'' لري، له هغو جوړښتونو څخه
بېل کړو چې داسې نۀ دي۔ البته بشپړه ازادي لرو چې خپل ``مطلوب''
جوړښتونه وټاکو؛ مثلاً ويلے شو چې د پريديکات~\LR{$\le$} تفسير بايد يو
ترتيب وي، يا يوازې د~\LR{$\Lang L$} له هغو تفسيرونو سره دلچسپي لرو چې
ساحه ئې له سټونو جوړه وي او~\LR{$\Obj \in$} د ``د غړے په توګه
شتون'' په اړيکه تفسير شي۔ خو ايا دا د ژبې په تړلي فارمولونه کولے شو؟
په بل عبارت، پر جوړښت~\LR{$\Struct M$} کوم شرطونه د~\LR{$\Struct M$}
د ژبې په يوه تړلے فارمول (يا ښايي د تړلي فارمولونه په يوه سټ)
بيانولے شو؟ ځينې شرطونه نۀ شو بيانولے۔ مثلاً د~\LR{$\Lang L_A$} داسې
هېڅ جمله نشته چې يوازې په هغه جوړښت~\LR{$\Struct M$} کښې رښتونې
وي چې~\LR{$\Domain M = \Nat$} وي۔ ``ساحه يوازې طبيعي عددونه لري'' نۀ شو
بيانولے۔ خو ښايي د جوړښتونه ځينې ``جوړښتي خاصيتونه'' بيانولے
شو۔ د جوړښتونه کوم خاصيتونه په تړلي فارمولونه بيانولے شو؟ يا
په بله وينا، د جوړښتونه کومې ټولګې د هغو جوړښتونو په توګه
بيانولے شو چې يوه تړلے فارمول (يا د تړلي فارمولونه يو سټ) رښتونې
کوي؟
\end{explain}
\olpseganchor{OLP-0169-B006}{end}

\olpseganchor{OLP-0169-B007}{start}
\begin{defn}[د يوه سټ مدل]
فرض کړئ~\LR{$\Gamma$} په ژبه~\LR{$\Lang L$} کښې د تړلي فارمولونه يو سټ دے۔
وايو چې جوړښت~\LR{$\Struct M$} د~\LR{$\Gamma$} \emph{مدل دے}، که د
هر~\LR{$!A \in \Gamma$} دپاره \LR{$\Sat{M}{!A}$} وي۔
\end{defn}
\olpseganchor{OLP-0169-B007}{end}

\olpseganchor{OLP-0169-B008}{start}
\begin{ex}
جمله~\LR{$\lforall[x][x \le x]$} په~\LR{$\Struct M$} کښې هله او يوازې هله
رښتونې ده چې~\LR{$\Assign{\le}{M}$} انعکاسي اړيکه وي۔ جمله
\LR{$\lforall[x][\lforall[y][((x \le y \land y \le x) \lif x = y)]]$}
په~\LR{$\Struct M$} کښې هله او يوازې هله رښتونې ده چې
\LR{$\Assign{\le}{M}$} ضد متناظره وي۔ جمله
\LR{$\lforall[x][\lforall[y][\lforall[z][((x \le y \land y \le z)
      \lif x \le z)]]]$} په~\LR{$\Struct M$} کښې هله او يوازې هله رښتونې
ده چې~\LR{$\Assign{\le}{M}$} انتقالي وي۔ نو د لاندې سټ مدلونه
\begin{align*}
\{\quad &\lforall[x][x \le x], \\
   & \lforall[x][\lforall[y][((x \le y \land y \le
    x) \lif x = y)]], \\
   &\lforall[x][\lforall[y][\lforall[z][((x \le y
      \land y \le z) \lif x \le z)]]] \quad \}
\end{align*}
کټ مټ هغه جوړښتونه دي چې~\LR{$\Assign{\le}{M}$} پکښې انعکاسي، ضد
متناظره او انتقالي، يعنې جزوي ترتيب وي۔ له همدې امله، دا د
\emph{جزوي ترتيبونو د لومړۍ درجې تيورۍ} د بديهي اصولو په توګه
اخستلے شو۔
\end{ex}
\olpseganchor{OLP-0169-B008}{end}

\olpseganchor{OLP-0170-B004}{start}
\olfileid{fol}{mat}{the}
\olpseganchor{OLP-0170-B004}{end}

\olpseganchor{OLP-0170-B005}{start}
\olsection{د لومړۍ درجې تيوريو بېلګې}
\olpseganchor{OLP-0170-B005}{end}

\olpseganchor{OLP-0170-B006}{start}
\begin{ex}
په ژبه~\LR{$\Lang L_<$} کښې د سختو خطي ترتيبونو تيوري د لاندې سټ له خوا
په بديهي اصولو ټاکل کېږي:
\begin{align*}
\{\quad & \lforall[x][\lnot x < x], \\
& \lforall[x][\lforall[y][((x < y \lor y <
    x) \lor x = y)]], \\
& \lforall[x][\lforall[y][\lforall[z][((x < y
      \land y < z) \lif x < z)]]] \quad \}
\end{align*}
دا مطلوب جوړښتونه په بشپړ ډول رانغاړي: هر سخت خطي ترتيب د دې
بديهي نظام مدل دے؛ او برعکس، که~\LR{$R$} پر يوه سټ~\LR{$X$} خطي ترتيب وي، نو
هغه جوړښت~\LR{$\Struct M$} چې \LR{$\Domain M = X$} او
\LR{$\Assign{<}{M} = R$} لري، د دې تيورۍ مدل دے۔
\end{ex}
\olpseganchor{OLP-0170-B006}{end}

\olpseganchor{OLP-0170-B007}{start}
\begin{ex}
د ګروپونو تيوري په هغه ژبه کښې چې~\LR{$\Obj 1$} (ثابت سمبول) او~\LR{$\cdot$}
(دوه‌ځاييزه تابع) لري، د لاندې فورمولونو له خوا په بديهي
اصولو ټاکل کېږي:
\begin{align*}
& \lforall[x][\eq[(x \cdot \Obj 1)][x]]\\
& \lforall[x][\lforall[y][\lforall[z][\eq[(x \cdot (y \cdot z))][((x
          \cdot y) \cdot z)]]]]\\
& \lforall[x][\lexists[y][\eq[(x \cdot y)][\Obj 1]]]
\end{align*}
\end{ex}
\olpseganchor{OLP-0170-B007}{end}

\olpseganchor{OLP-0170-B008}{start}
\begin{ex}
د پيانو د حساب تيوري د حساب په ژبه~\LR{$\Lang L_A$} کښې د لاندې جملو له
خوا په بديهي اصولو ټاکل کېږي۔
\begin{align*}
& \lforall[x][\lforall[y][(\eq[x'][y'] \lif \eq[x][y])]]\\
& \lforall[x][\eq/[\Obj 0][x']]\\
& \lforall[x][\eq[(x + \Obj 0)][x]]\\
& \lforall[x][\lforall[y][\eq[(x + y')][(x + y)']]]\\
& \lforall[x][\eq[(x \times \Obj 0)][\Obj 0]]\\
& \lforall[x][\lforall[y][\eq[(x \times y')][((x \times y) + x)]]]\\
& \lforall[x][\lforall[y][(x < y \liff \lexists[z][\eq[(z' + x)][y])]]]\\
\intertext{\RL{او ورسره د لاندې بڼې ټولې جملې}}
& (!A(\Obj 0) \land \lforall[x][(!A(x) \lif !A(x'))]) \lif \lforall[x][!A(x)]
\end{align*}
ځکه چې د وروستۍ بڼې نامتناهي ډېرې جملې شته، دا بديهي نظام
نامتناهي دے۔ وروستۍ بڼه \emph{د استقرا شېما} بلل کېږي۔ (په اصل کښې
د استقرا شېما تر هغه لږه پېچلې ده چې دلته مو ښودلې ده۔)
\olpseganchor{OLP-0170-B008}{end}

\olpseganchor{OLP-0170-B009}{start}
وروستے بديهي اصل د~\LR{$<$} يو \emph{صريح تعريف} دے۔
\end{ex}
\olpseganchor{OLP-0170-B009}{end}

\olpseganchor{OLP-0170-B010}{start}
\begin{ex}
د خالصو سټونو تيوري د رياضۍ په بنسټونو (او فلسفه) کښې مهمه ونډه
لري۔ يو سټ هله خالص دے چې ټول غړي ئې هم خالص سټونه وي۔ نو
تش سټ خالص ګڼل کېږي، خو هغه سټ چې کوم داسې څيز ئې غړے وي
چې خپله سټ نۀ وي، خالص نۀ دے۔ په دې توګه خالص سټونه يوازې له تش سټ
څخه، بې له ``غيرسټي بنسټيزو عناصرو''، جوړ شوي سټونه دي؛ يعنې داسې
څيزونه نۀ لري چې خپله سټونه نۀ وي۔
\olpseganchor{OLP-0170-B010}{end}

\olpseganchor{OLP-0170-B011}{start}
لاندې فورمولونه د خالصو سټونو د يوې تيورۍ بديهي نظام ګڼل کېدے شي:
\begin{align*}
& \lexists[x][\lnot \lexists[y][y \in x]]\\
& \lforall[x][\lforall[y][(\lforall[z](z \in x \liff z \in y) \lif
\eq[x][y])]]\\
& \lforall[x][\lforall[y][\lexists[z][\lforall[u][(u \in z \liff
(\eq[u][x] \lor \eq[u][y]))]]]]\\
& \lforall[x][\lexists[y][\lforall[z][(z \in y \liff \lexists[u][(z \in
        u \land u \in x)])]]]\\
\intertext{\RL{او ورسره د لاندې بڼې ټولې جملې}} &
\lexists[x][\lforall[y][(y \in x \liff !A(y))]]
\end{align*}
لومړے بديهي اصل وايي چې داسې يو سټ شته چې هېڅ غړي نۀ لري
(يعنې~\LR{$\emptyset$} شته)؛ دويم وايي چې سټونه امتدادي دي؛ درېيم وايي چې د هر
دوو سټونو~\LR{$X$} او~\LR{$Y$} دپاره سټ~\LR{$\{X, Y\}$} شته؛ څلورم وايي چې د هر
سټ~\LR{$X$} دپاره سټ~\LR{$\cup X$} شته، چېرته چې~\LR{$\cup X$} د~\LR{$X$} د ټولو
عناصرو اتحاد دے۔
\olpseganchor{OLP-0170-B011}{end}

\olpseganchor{OLP-0170-B012}{start}
وروستۍ يادې شوې تړلي فارمولونه په ګډه \emph{د ساده ادراک شېما} بلل
کېږي۔ په اصل کښې وايي چې د هر~\LR{$!A(x)$} دپاره سټ
\LR{$\Setabs{x}{!A(x)}$} شته---نو په لومړي نظر يو رښتونے، ګټور او ښايي
حتا ضروري بديهي اصل ښکاري۔ ``ساده'' ځکه بلل کېږي چې په پايله کښې دا
تيوري د صدق وړ نۀ پرېږدي: که~\LR{$!A(y)$} د~\LR{$\lnot y \in y$} په توګه
واخلئ، لاندې تړلے فارمول لاس ته راځي:
\[
\lexists[x][\lforall[y][(y \in x \liff \lnot y \in y)]]
\]
او دا تړلے فارمول په هېڅ جوړښت کښې نۀ صدق کوي۔
\end{ex}
\olpseganchor{OLP-0170-B012}{end}

\olpseganchor{OLP-0170-B013}{start}
\begin{ex}
د \emph{مېرېولوژۍ} يا جزپوهنې په ډګر کښې د \emph{برخه‌والي} اړيکه
بنسټيزه اړيکه ده۔ لکه د سټونو د تيوريو په شان، د برخه‌والي داسې
تيورۍ هم شته چې د دې اړيکې بېلابېل (او کله کله يو بل سره ټکر لرونکي)
تصورات په بديهي اصولو ټاکي۔
\olpseganchor{OLP-0170-B013}{end}

\olpseganchor{OLP-0170-B014}{start}
د مېرېولوژۍ ژبه يو دوه‌ځاييز اړوند سمبول~\LR{$\Obj P$} لري، او
\LR{$\Atom{\Obj P}{x, y}$} دا ``معنٰی'' لري چې~\LR{$x$} د~\LR{$y$} يوه برخه ده۔
کله چې دا تفسير په پام کښې ولرو، د دې ژبې جوړښت د
\emph{برخه‌والي جوړښت} بلل کېږي۔ البته د يوه دوه‌ځاييز اړوند هر
جوړښت په رښتيا د دې نوم وړ نۀ دے۔ د ``برخه‌والي'' د رانغاړلو دپاره
بايد~\LR{$\Assign{\Obj P}{M}$} ځينې شرطونه پوره کړي، او دا شرطونه د
برخه‌والي د تيورۍ د بديهي اصولو په توګه ټاکلے شو۔ مثلاً برخه‌والے
پر څيزونو جزوي ترتيب دے: هر څيز د خپل ځان (که څه هم \emph{غيرحقيقي})
برخه ده؛ دوه بېل څيزونه د يو بل برخې نۀ شي کېدے؛ او د يوه څيز د
برخې برخه خپله هم د هغه څيز برخه ده۔ پام وکړئ چې په دې معنٰی ``د
يوه څيز برخه ده'' له ``فرعي سټ دے'' سره ورته ده، خو له ``عنصر دے''
سره نۀ ده ورته، ځکه عنصرول نۀ انعکاسي دي او نۀ انتقالي۔
\begin{align*}
& \lforall[x][\Atom{\Obj P}{x,x}] \\
& \lforall[x][\lforall[y][((\Part{x}{y} \land \Part{y}{x})
      \lif \eq[x][y])]] \\
& \lforall[x][\lforall[y][\lforall[z][((\Part{x}{y} \land
        \Part{y}{z}) \lif \Part{x}{z})]]]\\
\intertext{\RL{سربېره پر دې، هر دوه څيزونه مېرېولوژيکي مجموعه لري (داسې
  څيز چې دغه دواړه څيزونه ئې برخې وي او په دې لحاظ تر ټولو کوچنے وي)۔}}  &
\lforall[x][\lforall[y][\lexists[z][\lforall[u][(\Part{z}{u} \liff
        (\Part{x}{u} \land \Part{y}{u}))]]]]
\end{align*}
دا يوازې د برخه‌والي ځينې هغه بنسټيز اصول دي چې مابعدالطبيعتپوهان ئې
څېړي۔ خو نور اصول ډېر ژر داسې پېچلي شي چې تر ځينو تعريف شويو اړيکو
وړاندې ئې جوړول يا ليکل ستونزمن وي۔ مثلاً په مېرېولوژۍ کښې زياتره
دلچسپي لرونکي مابعدالطبيعتپوهان لاندې اصل هم معتبر ګڼي: هر کله چې
څيز~\LR{$x$} حقيقي برخه~\LR{$y$} ولري، نو داسې برخه~\LR{$z$} هم لري چې له~\LR{$y$} سره
هېڅ ګډه برخه ونۀ لري، او د~\LR{$y$} او~\LR{$z$} امتزاج~\LR{$x$} وي۔
\end{ex}
\olpseganchor{OLP-0170-B014}{end}

\olpseganchor{OLP-0171-B004}{start}
\olfileid{fol}{mat}{exr}
\olpseganchor{OLP-0171-B004}{end}

\olpseganchor{OLP-0171-B005}{start}
\olsection{په  جوړښت کښې د
  اړيکو بيانول}
\olpseganchor{OLP-0171-B005}{end}

\olpseganchor{OLP-0171-B006}{start}
\begin{explain}
د فارمولونه يوه اصلي ګټه دا ده چې په جوړښت~\LR{$\Struct M$}
کښې د~\LR{$\Struct M$} د ژبې~\LR{$\Lang L$} د بنسټيزو سمبولونو له مخې
خاصيتونه او اړيکې بيانوي۔ مطلب مو دا دے: د~\LR{$\Struct M$} د تعريف ساحه
د څيزونو يو سټ دے۔ ثابت سمبولونه، تابعې او پريديکاتونه
په~\LR{$\Struct M$} کښې په ترتيب سره د~\LR{$\Domain M$} په ځينو څيزونو،
پر~\LR{$\Domain M$} تابعو او پر~\LR{$\Domain M$} اړيکو تفسيرېږي۔ مثلاً که
\LR{$\Obj A^2_0$} په~\LR{$\Lang L$} کښې وي، نو~\LR{$\Struct M$} ورته
اړيکه~\LR{$R = \Assign{{\Obj A^2_0}}{M}$} ګوماري۔ بيا فارمول
\LR{$\Atom{\Obj A^2_0}{\Obj v_1, \Obj v_2}$} همدغه اړيکه په لاندې معنٰی
\emph{بيانوي}: که يوه متغير-ګومارنه~\LR{$s$}، \LR{$\Obj v_1$} له
\LR{$a \in \Domain{M}$} او~\LR{$\Obj v_2$} له~\LR{$b \in \Domain M$} سره وتړي، نو
\[
Rab \text{\RL{\quad هله او يوازې هله چې\quad}} \Sat{M}{\Atom{\Obj A^2_0}{\Obj v_1, \Obj v_2}}[s].
\]
پام وکړئ چې دلته بايد متغير-ګومارنې شاملې کړو: يوازې دا نۀ شو ويلے
چې ``\LR{$Rab$} هله او يوازې هله چې
\LR{$\Sat{M}{\Atom{\Obj A^2_0}{a, b}}$}''، ځکه~\LR{$a$} او~\LR{$b$} زموږ د ژبې
سمبولونه نۀ، بلکې د~\LR{$\Domain{M}$} غړي دي۔
\olpseganchor{OLP-0171-B006}{end}

\olpseganchor{OLP-0171-B007}{start}
موږ يوازې اټومي فارمولونه نۀ لرو، بلکې په منطقي نښلوونکو او
سورونو ئې يوځاے کولے شو؛ نو لا پېچلي فارمولونه نورې هغه اړيکې
تعريفولے شي چې نېغ په نېغه په~\LR{$\Struct M$} کښې نۀ دي شاملې۔ موږ دې
ته ګورو چې دا څنګه کېږي، او په ځانګړي ډول په جوړښت کښې کومې
اړيکې تعريفولے شو۔
\end{explain}
\olpseganchor{OLP-0171-B007}{end}

\olpseganchor{OLP-0171-B008}{start}
\begin{defn}
فرض کړئ~\LR{$!A(\Obj v_1,\dots, \Obj v_n)$} د~\LR{$\Lang L$} داسې
فارمول دے چې يوازې~\LR{$\Obj v_1$}،\dots،~\LR{$\Obj v_n$} پکښې ازاد
راځي، او~\LR{$\Struct M$} د~\LR{$\Lang L$} يو جوړښت دے۔
\LR{$!A(\Obj v_1,\dots, \Obj v_n)$} هله او يوازې هله
\emph{اړيکه}~\LR{$R \subseteq \Domain M^n$} \emph{بيانوي} چې
\[
Ra_1\dots a_n \text{\RL{\quad هله او يوازې هله چې\quad}} \Sat{M}{\Atom{!A}{\Obj
    v_1,\dots, \Obj v_n}}[s]
\]
د هرې داسې متغير-ګومارنې~\LR{$s$} دپاره چې
\LR{$s(\Obj v_i) = a_i$} (\LR{$i = 1, \dots, n$}) وي۔
\end{defn}
\olpseganchor{OLP-0171-B008}{end}

\olpseganchor{OLP-0171-B009}{start}
\begin{ex}
د حساب په معياري مدل~\LR{$\Struct N$} کښې فارمول
\LR{$\Obj v_1 < \Obj v_2 \lor \eq[\Obj v_1][\Obj v_2]$} پر~\LR{$\Nat$} د
\LR{$\le$} اړيکه بيانوي۔ فارمول~\LR{$\eq[\Obj v_2][\Obj v_1']$} د ناستي
اړيکه بيانوي؛ يعنې هغه~\LR{$R \subseteq \Nat^2$} چې~\LR{$Rnm$} هله وي چې~\LR{$m$}
د~\LR{$n$} ناستے وي۔ فارمول~\LR{$\eq[\Obj v_1][\Obj v_2']$} د مخکښ اړيکه
بيانوي۔ دواړه فارمولونه
\LR{$\lexists[\Obj v_3][(\eq/[\Obj v_3][\Obj 0] \land
  \eq[\Obj v_2][(\Obj v_1 + \Obj v_3)])]$} او~\LR{$\lexists[\Obj
  v_3][\eq[(\Obj v_1 + {\Obj v_3}')][\Obj v_2]]$} د~\LR{$\Obj <$} اړيکه
بيانوي۔ د دې معنٰی دا ده چې اړوند سمبول~\LR{$<$} په رښتيا د حساب په ژبه
کښې زائد دے؛ تعريفېدے شي۔
\end{ex}
\olpseganchor{OLP-0171-B009}{end}

\olpseganchor{OLP-0171-B010}{start}
\begin{explain}
دا مفکوره يوازې په ځانګړو جوړښتونه کښې په زړه پورې نۀ ده،
بلکې هر کله مهمه ده چې يوه ژبه د يوه يا څو مطلوبو مدلونو د بيانولو
دپاره کاروو، يعنې تيورۍ څېړو۔ دغه تيورۍ زياتره يوازې څو
پريديکاتونه د بنسټيزو سمبولونو په توګه لري، خو په هغه ساحه کښې
چې بيانوي ئې، ډېرې نورې اړيکې مهمې ونډې لري۔ که دغه نورې اړيکې د ژبې
د بنسټيزو پريديکاتونه د تفسيروونکو اړيکو له مخې په منظمه توګه
بيانېدے شي، وايو چې هغوئ په ژبه کښې \emph{تعريفولے} شو۔
\end{explain}
\olpseganchor{OLP-0171-B010}{end}

\olpseganchor{OLP-0171-B011}{start}
\begin{prob}
په~\LR{$\Lang L_A$} کښې داسې فارمولونه ومومئ چې لاندې اړيکې تعريف کړي:
\begin{enumerate}
\item \LR{$n$} د~\LR{$i$} او~\LR{$j$} تر منځ دے؛
\item \LR{$n$}، \LR{$m$} په پوره ډول وېشي (يعنې~\LR{$m$} د~\LR{$n$} مضرب دے)؛
\item \LR{$n$} اول عدد دے (يعنې له~\LR{$1$} او~\LR{$n$} پرته بل هېڅ عدد~\LR{$n$} په پوره
  ډول نۀ وېشي)۔
\end{enumerate}
\end{prob}
\olpseganchor{OLP-0171-B011}{end}

\olpseganchor{OLP-0171-B012}{start}
\begin{prob}
فرض کړئ فارمول~\LR{$!A(\Obj v_1, \Obj v_2)$} په
جوړښت~\LR{$\Struct M$} کښې اړيکه~\LR{$R \subseteq \Domain M^2$}
بيانوي۔ داسې فارمولونه ومومئ چې لاندې اړيکې بيان کړي:
\begin{enumerate}
\item د~\LR{$R$} معکوسه اړيکه~\LR{$R^{-1}$}؛
\item نسبي حاصل‌ضرب~\LR{$R \mid R$}؛
\end{enumerate}
ايا د~\LR{$R$} د انتقالي بندښت~\LR{$R^+$} د بيانولو کومه طريقه موندلے شئ؟
\end{prob}
\olpseganchor{OLP-0171-B012}{end}

\olpseganchor{OLP-0171-B013}{start}
\begin{prob}
فرض کړئ~\LR{$\Lang{L}$} هغه ژبه ده چې يوازې دوه‌ځاييز اړوند سمبول~\LR{$<$}
لري (بل هېڅ ثابت سمبولونه، تابعې يا پريديکاتونه نۀ لري---
البته له~\LR{$\eq$} پرته)۔ فرض کړئ~\LR{$\Struct{N}$} داسې جوړښت دے چې
\LR{$\Domain{N} = \Nat$}، او
\LR{$\Assign{<}{N} = \Setabs{\tuple{n,m}}{n < m}$}۔ لاندې ثابت کړئ:
\begin{enumerate}
\item \LR{$\{ 0 \}$} په~\LR{$\Struct{N}$} کښې تعريفېدے شي؛
\item \LR{$\{ 1 \}$} په~\LR{$\Struct{N}$} کښې تعريفېدے شي؛
\item \LR{$\{ 2 \}$} په~\LR{$\Struct{N}$} کښې تعريفېدے شي؛
\item د هر~\LR{$n \in \Nat$} دپاره سټ~\LR{$\{ n \}$} په~\LR{$\Struct{N}$} کښې
  تعريفېدے شي؛
\item د~\LR{$\Domain{N}$} هر متناهي فرعي سټ په~\LR{$\Struct{N}$} کښې
  تعريفېدے شي؛
\item د~\LR{$\Domain{N}$} هر د متناهي متمم لرونکے فرعي سټ په~\LR{$\Struct{N}$} کښې
  تعريفېدے شي (چېرته چې~\LR{$X \subseteq \Nat$} هله او يوازې هله
  د متناهي متمم لرونکے دے چې~\LR{$\Nat \setminus X$} متناهي وي)۔
\end{enumerate}
\end{prob}
\olpseganchor{OLP-0171-B013}{end}

\olpseganchor{OLP-0172-B004}{start}
\olfileid{fol}{mat}{set}
\olpseganchor{OLP-0172-B004}{end}

\olpseganchor{OLP-0172-B005}{start}
\olsection{د سټونو تيوري}
\olpseganchor{OLP-0172-B005}{end}

\olpseganchor{OLP-0172-B006}{start}
د رياضۍ کابو ټولې برخې د سټونو په تيورۍ کښې جوړېدے شي۔ په دې
تيورۍ کښې د رياضۍ جوړول څو کارونه غواړي۔ لومړے، د~\LR{$\in$} اړيکې دپاره
د بديهي اصولو يو سټ غواړي۔ بېلابېل بديهي نظامونه جوړ شوي دي چې کله
کله د~\LR{$\in$} په باب يو بل سره ټکر لرونکي خاصيتونه لري۔ د~\LR{$\Log{ZFC}$}
په نوم بديهي نظام، يعنې د انتخاب له بديهي اصل سره د
زرملو--فرېنکل سټ تيوري، ځانګړے مقام لري: دا تر بل هر بديهي نظام ډېر
کارول کېږي او څېړل کېږي، ځکه بديهي اصول ئې د
کابو ټولو هغو خبرو د ثبوت دپاره بس کېږي چې رياضيپوهان ئې ثابتول
غواړي۔ خو د دې له څرګندولو وړاندې بايد لومړے روښانه کړو چې د
رياضيپوهانو مطلوبې خبرې اصلًا څنګه \emph{بيانولے} شو۔ د پيل دپاره،
ژبه هېڅ ثابت سمبولونه يا تابعې نۀ لري؛ نو په لومړي نظر دا
روښانه نۀ ده چې د ځانګړو سټونو (لکه~\LR{$\emptyset$} يا~\LR{$\Nat$})، پر
سټونو د عملونو (لکه~\LR{$X \cup Y$} او~\LR{$\Pow{X}$})، يا لا د هغو جوړښتونو
په باب څنګه خبرې وکړو چې له سټونو پرته نور څيزونه هم رانغاړي، لکه
اړيکې او تابعې۔
\olpseganchor{OLP-0172-B006}{end}

\olpseganchor{OLP-0172-B007}{start}
په پيل کښې، ``عنصر دے'' يوازېنۍ مهمه اړيکه نۀ ده؛ ``فرعي سټ دے''
کابو همدومره مهمه ښکاري۔ خو ``فرعي سټ دے'' د ``عنصر دے'' له مخې
\emph{تعريفولے} شو۔ د دې دپاره بايد د سټ تيورۍ په ژبه کښې داسې
فارمول~\LR{$!A(x, y)$} ومومو چې د سټونو يوه جوړه~\LR{$\tuple{X, Y}$}
هله او يوازې هله پوره کړي چې~\LR{$X \subseteq Y$} وي۔ خو~\LR{$X$} هله او
يوازې هله د~\LR{$Y$} فرعي سټ دے چې د~\LR{$X$} ټول غړي د~\LR{$Y$} هم
غړي وي۔ نو~\LR{$\subseteq$} په لاندې فارمول تعريفولے شو:
\[
\lforall[z][(z \in x \lif z \in y)]
\]
اوس چې هر کله په يوه فارمول کښې اړيکه~\LR{$\subseteq$} کارول غواړو، د
هغې پر ځاے همدا فارمول کارولے شو (په دې شرط چې~\LR{$x$} او~\LR{$y$} په مناسب
ډول بدل او تړلے متغير~\LR{$z$} د اړتيا پر وخت بيا ونومول شي)۔ مثلاً د
سټونو امتداديتوب وايي چې که هر دوه سټونه~\LR{$x$} او~\LR{$y$} په دوه اړخيزه
توګه په يو بل کښې شامل وي، نو~\LR{$x$} او~\LR{$y$} بايد هماغه يو سټ وي۔ دا په
\LR{$\lforall[x][\lforall[y][((x \subseteq y \land y
    \subseteq x) \lif x = y)]]$} بيانېدے شي؛ يا که~\LR{$\subseteq$} په
پورته تعريف بدل کړو، په لاندې ډول:
\[
\lforall[x][\lforall[y][((\lforall[z][(z \in x \lif z \in y)] \land
    \lforall[z][(z \in y \lif z \in x)]) \lif x = y)]].
\]
دا په رښتيا د~\LR{$\Log{ZFC}$} يو بديهي اصل، يعنې ``د امتداديتوب بديهي
اصل'' دے۔
\olpseganchor{OLP-0172-B007}{end}

\olpseganchor{OLP-0172-B008}{start}
د~\LR{$\emptyset$} دپاره هېڅ ثابت سمبول نشته، خو ``\LR{$x$} تش دے'' په
\LR{$\lnot \lexists[y][y \in x]$} بيانولے شو۔ بيا ``\LR{$\emptyset$} شته''
په تړلے فارمول~\LR{$\lexists[x][\lnot \lexists[y][y \in x]]$} اوړي۔
دا د~\LR{$\Log{ZFC}$} يو بل بديهي اصل دے۔ (پام وکړئ چې د امتداديتوب
بديهي اصل دا لازموي چې يوازې يو تش سټ وي۔) هر کله چې د سټ تيورۍ په
ژبه کښې د~\LR{$\emptyset$} په باب خبرې کول غواړو، داسې به ليکو: ``داسې
يو سټ شته چې تش دے او~\dots''۔ د بېلګې په توګه، دا حقيقت چې
\LR{$\emptyset$} د هر سټ فرعي سټ دے، په لاندې ډول بيانولے شو:
\[
\lexists[x][(\lnot\lexists[y][y \in x] \land \lforall[z][x \subseteq
    z])]
\]
چېرته چې، البته، \LR{$x \subseteq z$} به بيا په خپل تعريف بدلېږي۔
\olpseganchor{OLP-0172-B008}{end}

\olpseganchor{OLP-0172-B009}{start}
د سټونو د عملونو، لکه~\LR{$X \cup Y$} او~\LR{$\Pow{X}$}، په باب د خبرو دپاره
ورته چل کارول پکار دي۔ د سټ تيورۍ په ژبه کښې د تابعې سمبولونه نشته،
خو تابعوي اړيکې~\LR{$X \cup Y = Z$} او~\LR{$\Pow{X} = Y$} په لاندې ډول
بيانولے شو:
\begin{align*}
& \lforall[u][((u \in x \lor u \in y) \liff u \in z)]\\
& \lforall[u][(u \subseteq x \liff u \in y )]
\end{align*}
ځکه د~\LR{$X \cup Y$} غړي کټ مټ هغه سټونه دي چې يا د~\LR{$X$}
غړي وي يا د~\LR{$Y$} (يا د دواړو) غړي وي، او د~\LR{$\Pow{X}$}
غړي کټ مټ
د~\LR{$X$} فرعي سټونه دي۔ خو دا اجازه نۀ راکوي چې~\LR{$x \cup y$} يا~\LR{$\Pow{x}$}
د ترمونو په توګه وکاروو: يوازې ټول هغه فارمولونه کارولے شو چې
اړيکې~\LR{$X \cup Y = Z$} او~\LR{$\Pow{X} = Y$} تعريفوي۔ په اصل کښې لا دا هم
نۀ پوهېږو چې دغه اړيکې کله هم پوره کېږي، يعنې نۀ پوهېږو چې اتحادونه
او د فرعي سټونو سټونه تل شته۔ مثلاً تړلے فارمول
\LR{$\lforall[x][\lexists[y][\Pow{x} = y]]$} د~\LR{$\Log{ZFC}$} يو بل بديهي
اصل دے (د فرعي سټونو د سټ بديهي اصل)۔
\olpseganchor{OLP-0172-B009}{end}

\olpseganchor{OLP-0172-B010}{start}
اوس د مرتبو جوړو يا تابعو په باب خبرې څنګه کوو؟ دلته بايد څرګنده کړو
چې مرتبې جوړې او تابعې څنګه د سټونو ځانګړي ډولونه ګڼلے شو۔ د مرتبې
جوړې~\LR{$\tuple{x, y}$} د تعريف يوه طريقه سټ~\LR{$\{\{x\}, \{x, y\}\}$} ده۔
خو لکه مخکښې، داسې تابع نۀ شو معرفي کولے چې دغه سټ ونوموي؛
يوازې اړيکه~\LR{$\tuple{x, y} = z$}، يعنې
\LR{$\{\{x\}, \{x, y\}\} = z$}، تعريفولے شو:
\[
\lforall[u][(u \in z \liff (\lforall[v][(v \in u \liff v = x)] \lor
  \lforall[v][(v \in u \liff (v = x \lor v = y))]))]
\]
دا وايي چې د~\LR{$z$} غړي~\LR{$u$} کټ مټ هغه سټونه دي چې يا يوازې
\LR{$x$} ئې غړے وي، يا يوازې~\LR{$x$} او~\LR{$y$} ئې غړي وي (په
بل عبارت، هغه سټونه چې يا له~\LR{$\{x\}$} سره يو شان وي يا
له~\LR{$\{x, y\}$} سره)۔ چې دا ولرو، نورې خبرې هم کولے شو؛ مثلاً
\LR{$X \times Y = Z$}:
\[
\lforall[z][(z \in Z \liff \lexists[x][\lexists[y][(x \in
        X \land y \in Y \land \tuple{x, y} = z)]])]
\]
\olpseganchor{OLP-0172-B010}{end}

\olpseganchor{OLP-0172-B011}{start}
تابعه~\LR{$f \colon X \to Y$} د اړيکې~\LR{$f(x) = y$} په توګه، يعنې د جوړو
د سټ~\LR{$\Setabs{\tuple{x,y}}{f(x) = y}$} په توګه، ګڼلے شو۔ بيا ويلے شو
چې سټ~\LR{$f$} له~\LR{$X$} څخه~\LR{$Y$} ته تابعه ده، که (الف) يوه اړيکه
\LR{$\subseteq X \times Y$} وي؛ (ب) ټوله وي، يعنې د هر~\LR{$x \in X$} دپاره
داسې~\LR{$y \in Y$} وي چې~\LR{$\tuple{x, y} \in f$}؛ او (ج) تابعوي وي، يعنې
هر کله چې~\LR{$\tuple{x, y}, \tuple{x, y'} \in f$} وي، نو~\LR{$y = y'$} وي
(ځکه د تابعې قيمتونه بايد يوازيني وي)۔ نو ``\LR{$f$} له~\LR{$X$} څخه~\LR{$Y$} ته
تابعه ده'' داسې ليکل کېدے شي:
\begin{align*}
 \lforall[u][(u \in f \lif {} & \lexists[x][\lexists[y][(x \in X \land y \in
      Y \land \tuple{x, y} = u)]])] \land {}\\
 \lforall[x][(x \in X \lif {} &
      (\lexists[y][(y \in Y \land \mathrm{maps}(f, x, y))] \land {}\\
& (\lforall[y][\lforall[y'][((\mathrm{maps}(f, x, y) \land
    \mathrm{maps}(f, x, y')) \lif y = y')]])))]
\end{align*}
چېرته چې~\LR{$\mathrm{maps}(f, x, y)$} د
\LR{$\lexists[v][(v \in f \land \tuple{x, y} = v)]$} لنډون دے (دا
فارمول ``\LR{$f(x) = y$}'' بيانوي)۔
لومړے کلي شرط د تابعې ټولې جوړې په ټاکلې ساحه او د ارزښتونو په سټ کښې
ځايوي۔ دويم کلي شرط وايي چې د ساحې هر غړے کټ مټ يو ارزښت لري۔
\olpseganchor{OLP-0172-B011}{end}

\olpseganchor{OLP-0172-B012}{start}
اوس دا هم ستونزمنه نۀ ده چې مثلاً بيان کړو چې~\LR{$f\colon X \to Y$}
يو پر يو ده:
\begin{multline*}
 f \colon X \to Y \land \lforall[x][\lforall[x'][((x \in X \land x' \in
     X \land {} \\
 \lexists[y][(\mathrm{maps}(f, x, y) \land \mathrm{maps}(f,
        x', y))]) \lif x = x')]]
\end{multline*}
تابعه~\LR{$f\colon X \to Y$} هله او يوازې هله يو پر يو ده چې هر
کله~\LR{$f$}، \LR{$x, x' \in X$} يوه~\LR{$y$} ته ولېږي، نو~\LR{$x = x'$} وي۔ که دا
فارمول په~\LR{$\mathrm{inj}(f, X, Y)$} لنډ کړو، لا له مخکښې په داسې مقام
کښې يو چې د سټ تيورۍ په ژبه د کانتور د قضيې غوندې نۀ معمولي خبره
بيان کړو: له~\LR{$\Pow{X}$} څخه~\LR{$X$} ته هېڅ يو پر يو تابعه نشته:
\[
\lforall[X][\lforall[Y][(\Pow{X} = Y \lif
    \lnot\lexists[f][\mathrm{inj}(f, Y, X)])]]
\]
\olpseganchor{OLP-0172-B012}{end}

\olpseganchor{OLP-0172-B013}{start}
کېدے شي څوک فکر وکړي چې سټ تيوري يو بل داسې بديهي اصل غواړي چې د هر
تعريفوونکي خاصيت دپاره د يوه سټ شتون تضمين کړي۔ که~\LR{$!A(x)$} د سټ
تيورۍ داسې فارمول وي چې متغير~\LR{$x$} پکښې ازاد وي، لاندې تړلے فارمول
په پام کښې نيولے شو:
\[
\lexists[y][\lforall[x][(x \in y \liff !A(x))]].
\]
دا تړلے فارمول وايي چې داسې سټ~\LR{$y$} شته چې غړي ئې ټول او
يوازې هغه~\LR{$x$} دي چې~\LR{$!A(x)$} پوره کوي۔ دا شېما ``د ادراک اصل'' بلل
کېږي۔ ډېره ګټوره ښکاري؛ خو له بده مرغه ناسازګاره ده۔
\LR{$!A(x) \ident \lnot x \in x$} واخلئ؛ بيا د ادراک اصل وايي:
\[
\lexists[y][\lforall[x][(x \in y \liff x \notin x)]],
\]
يعنې د ټولو هغو سټونو د سټ شتون وايي چې د خپل ځان غړي نۀ
وي۔ داسې سټ نۀ شي موجودېدے---دا د راسل تناقض دے۔ په اصل کښې
\LR{$\Log{ZFC}$} د دې اصل يوه محدوده---او سازګاره---بڼه لري، يعنې د
بېلتون اصل:
\[
\lforall[z][\lexists[y][\lforall[x][(x \in y \liff (x \in z \land
      !A(x))]]].
\]
\olpseganchor{OLP-0172-B013}{end}

\olpseganchor{OLP-0172-B014}{start}
\begin{prob}
داسې اشتقاق ورکړئ چې لاندې وښيي، او په دې توګه ثابته کړئ
چې د ادراک اصل ناسازګار دے:
\[
\lexists[y][\lforall[x][(x \in y \liff x \notin x)]] \Proves \lfalse.
\]
ښايي لومړے د لاندې ثبوت ګټور وي:
\LR{$(A \lif \lnot A) \land (\lnot A \lif A) \Proves \lfalse$}۔
\end{prob}
\olpseganchor{OLP-0172-B014}{end}

\olpseganchor{OLP-0173-B004}{start}
\olfileid{fol}{mat}{siz}
\olpseganchor{OLP-0173-B004}{end}

\olpseganchor{OLP-0173-B005}{start}
\olsection{د جوړښتونه د کچې بيانول}
\olpseganchor{OLP-0173-B005}{end}

\olpseganchor{OLP-0173-B006}{start}
\begin{explain}
د جوړښتونو ځينې خاصيتونه د ژبې د غيرمنطقي سمبولونو له کارولو پرته
هم بيانولے شو۔ مثلاً داسې تړلي فارمولونه شته چې په جوړښت
کښې هله او يوازې هله رښتونې وي چې د تعريف ساحه، يعنې د جوړښت
ساحه، لږ تر لږه، زيات نه زيات، يا کټ مټ يو ټاکلے شمېر~\LR{$n$}
غړي ولري۔
\end{explain}
\olpseganchor{OLP-0173-B006}{end}

\olpseganchor{OLP-0173-B007}{start}
\begin{prop}
تړلے فارمول
\begin{multline*}
  !A_{\ge n} \ident \lexists[x_1][\lexists[x_2][\dots\lexists[x_n][{}\\
  \begin{aligned}
  (\eq/[x_1][x_2] \land {}
  \eq/[x_1][x_3] \land \eq/[x_1][x_4] \land \dots \land \eq/[x_1][x_n] \land {}\\
\eq/[x_2][x_3] \land \eq/[x_2][x_4] \land \dots \land {} \eq/[x_2][x_n] \land {} \\
\vdots\\
\eq/[x_{n-1}][x_n])
\end{aligned}]]]
\end{multline*}
په جوړښت~\LR{$\Struct M$} کښې هله او يوازې هله رښتونې ده چې
\LR{$\Domain M$} لږ تر لږه~\LR{$n$} غړي ولري۔ په پايله کښې،
\LR{$\Sat{M}{\lnot !A_{\ge n+1}}$} هله او يوازې هله چې~\LR{$\Domain M$}
زيات نه زيات~\LR{$n$} غړي ولري۔
وجودي سوري د داسې غړو د شتون خبره کوي چې هر دوه يې سره بېل وي۔
که يوازې يو غړے وغواړو، د غړو د بېلوالي لپاره هېڅ شرط نشته۔
\olpseganchor{OLP-0173-B007}{end}

\olpseganchor{OLP-0173-B008}{start}
\end{prop}
\olpseganchor{OLP-0173-B008}{end}

\olpseganchor{OLP-0173-B009}{start}
\begin{prop}
تړلے فارمول
\begin{multline*}
!A_{= n} \ident \lexists[x_1][\lexists[x_2][\dots\lexists[x_n][{} \\
\begin{aligned}
  (\eq/[x_1][x_2] \land {}
  \eq/[x_1][x_3] \land \eq/[x_1][x_4] \land \dots \land \eq/[x_1][x_n] \land {}\\
\eq/[x_2][x_3] \land \eq/[x_2][x_4] \land \dots \land {} \eq/[x_2][x_n] \land {} \\
\vdots\\
\eq/[x_{n-1}][x_n] \land {} \\
\lforall[y][(\eq[y][x_1] \lor \dots \lor \eq[y][x_n])])
\end{aligned}]]]
\end{multline*}
په جوړښت~\LR{$\Struct M$} کښې هله او يوازې هله رښتونې ده چې
\LR{$\Domain M$} کټ مټ~\LR{$n$} غړي ولري۔
وجودي سوري د بېلو غړو شتون غواړي، او وروستے کلي سور وايي چې
د ساحې هر غړے له همدغو غړو څخه يو دے۔
دلته ټول شرط د وجودي سورو په ساحه کښې او د کلي سور قوسونه جوړې دي۔
\end{prop}
\olpseganchor{OLP-0173-B009}{end}

\olpseganchor{OLP-0173-B010}{start}
\begin{prop}
يو جوړښت هله او يوازې هله نامتناهي دے چې د لاندې سټ مدل وي:
\[
\{!A_{\ge 1}, !A_{\ge 2}, !A_{\ge 3}, \dots \}.
\]
\end{prop}
\olpseganchor{OLP-0173-B010}{end}

\olpseganchor{OLP-0173-B011}{start}
هېڅ يوه خالصه منطقي جمله نشته چې په~\LR{$\Struct M$} کښې هله او يوازې
هله رښتونې وي چې~\LR{$\Domain M$} نامتناهي وي۔ خو له غيرمنطقي
پريديکاتونه سره داسې تړلي فارمولونه ورکولے شو چې يوازې نامتناهي
مدلونه لري (که څه هم هر نامتناهي جوړښت ئې مدل نۀ دے)۔ د
متناهي جوړښت خاصيت او د د شمېر نۀ وړ جوړښت خاصيت حتا د
تړلي فارمولونه په نامتناهي سټ هم نۀ شي بيانېدے۔ دا حقيقتونه د
متناهي شاهد او لوېنهايم--سکولم له قضيو څخه راځي۔
\olpseganchor{OLP-0173-B011}{end}

\olpseganchor{OLP-0174-B004}{start}
\olchapter{fol}{byd}{د لومړۍ درجې له منطق څخه وړاندې}
\olpseganchor{OLP-0174-B004}{end}

\olpseganchor{OLP-0174-B005}{start}
\begin{editorial}
  دا څپرکى، چې د جېرمي اېوېګېډ د منطق له يادښتونو څخه برابر شوے، د
  شته نورو منطقي نظامونو يوازې تر ټولو لنډه کتنه وړاندې کوي۔ موخه ئې دا
  ده چې په هغه کورس کښې چې دا موضوعات نۀ رانغاړي، د نورو مطالعو موضوعات
  وروښيي۔ هيله ده چې دلته ياده شوې هره موضوع به بالاخره په پرانيستې
  منطق پروژه کښې د يوې مستقلې برخې په کچه بشپړ بحث ومومي۔
\end{editorial}
\olpseganchor{OLP-0174-B005}{end}

\olpseganchor{OLP-0175-B004}{start}
\olfileid{fol}{byd}{int}
\olpseganchor{OLP-0175-B004}{end}

\olpseganchor{OLP-0175-B005}{start}
\olsection{ټوليزه کتنه}
\olpseganchor{OLP-0175-B005}{end}

\olpseganchor{OLP-0175-B006}{start}
د لومړۍ درجې منطق د پام وړ يوازينے منطقي نظام نۀ دے: د دې منطق ډېر
غځېدلي او بدل ډولونه شته۔ يو منطق په عام ډول د يوې ژبې صوري مشخصات،
اکثره، خو تل نۀ، د استنتاج يو نظام، او اکثره، خو تل نۀ، يوه ټاکل شوې
معنٰی پوهنه لري۔ خو د دې ويي تخنيکي کارونه يوه څرګنده پوښتنه راپورته
کوي: هغه منطقونه چې د لومړۍ درجې منطق نۀ دي، په وجداني يا فلسفي معنا
د ``منطق'' له ويي سره څه اړيکه لري؟ لاندې ټول بيان شوي نظامونه د يوه
يا بل ډول استدلال د مدل کولو دپاره جوړ شوي؛ آيا ويلے شو چې څه شي دوى
منطقي کوي؟
\olpseganchor{OLP-0175-B006}{end}

\olpseganchor{OLP-0175-B007}{start}
اسانه ځوابونه نۀ تر سترګو کېږي۔ د ``منطق'' ويے په بېلابېلو لارو او
بېلابېلو چاپېريالونو کښې کارول کېږي، او دا مفهوم، لکه د ``رښتيا''
مفهوم، له ډېرو فلسفي دريځونو څېړل شوے دے۔ د بېلګې په توګه، څوک ښايي
د منطقي استدلال موخه دا وګڼي چې معلومه کړي کومې ويناوې په لازمي ډول
رښتونې، له تجربې وړاندې رښتونې، د غيرمنطقي اصطلاحاتو له تفسير څخه
خپلواکې رښتونې، د خپلې بڼې له مخې رښتونې، يا د ژبني تړون له مخې
رښتونې دي؛ او له دغو مفکورو څخه هره يوه ډېر وضاحت غواړي۔ که پام يوازې
په رياضياتو کښې کارېدونکي منطق ته هم محدود کړو، د هغه د ساحې په هکله لږ
توافق شته۔ د بېلګې په توګه، رسل او وايټ‌هېډ په
\textit{پرېنسيپيا مېتېماتيکا} کښې هڅه وکړه چې رياضيات د منطق پر بنسټ،
د فرېګې له پيل کړې {\em منطقپالې} دود سره سم، جوړ کړي۔ د دوى منطقي
نظام د لوړو ډولونو د منطق يوه بڼه وه چې لاندې بيان شوي نظام ته ورته
وه۔ په پاى کښې اړ شول داسې بديهي اصول ورزيات کړي چې د ډېرو معيارونو
له مخې سوچه منطقي نۀ ښکاري (په ځانګړي ډول د لامتناهي والي بديهي اصل
او د راکمېدنې بديهي اصل)؛ خو بيا هم څوک دا دريځ لرلے شي چې د لوړې
درجې استدلال ځينې بڼې بايد منطقي ومنل شي۔ د دې پر خلاف، کواين، چې
هستپوهنه ئې ``قضيې'' د بحث د روا څيزونو په توګه نۀ مني، استدلال کوي
چې د دويمې او لوړې درجې منطق په اصل کښې د سټ تيوري ده چې د پسه پوست
ئې اغوستے؛ په نورو ټکو، هغه نظامونه چې پر پريديکاتونو کميت ټاکنه لري
سوچه منطقي نۀ دي۔
\olpseganchor{OLP-0175-B007}{end}

\olpseganchor{OLP-0175-B008}{start}
اوس دپاره غوره ده چې دا ډول فلسفي مسئلې بل وخت ته پرېږدو او لاندې
نظامونه يوازې د بېلابېلو ډولونو استدلال، که منطقي وي که بل څه، صوري
ايډيالې بڼې وګڼو۔
\olpseganchor{OLP-0175-B008}{end}

\olpseganchor{OLP-0176-B004}{start}
\olfileid{fol}{byd}{msl}
\olpseganchor{OLP-0176-B004}{end}

\olpseganchor{OLP-0176-B005}{start}
\olsection{د ډېرو ډولونو منطق}
\olpseganchor{OLP-0176-B005}{end}

\olpseganchor{OLP-0176-B006}{start}
د لومړۍ درجې په منطق کښې متغيرونه او کميت ټاکونکي پر يوه
د تعريف ساحه ګرځي۔ خو ډېر ځله د څو (يو له بله بېلو) د تعريف ساحې درلودل
ګټور وي: د بېلګې په توګه، ښايي د عددونو د تعريف ساحه، د هندسي څيزونو
د تعريف ساحه، له عددونو څخه عددونو ته د تابعو د تعريف ساحه، د ابېلي
ګروپونو د تعريف ساحه او داسې نور وغواړئ۔
\olpseganchor{OLP-0176-B006}{end}

\olpseganchor{OLP-0176-B007}{start}
د ډېرو ډولونو منطق دا ډول چوکاټ برابروي۔ پيل د ``ډولونو'' له يوه لړ
څخه کېږي---د يوه څيز ``ډول'' هغه ``د تعريف ساحه'' ښيي چې څيز بايد پکې
اوسېږي۔ بيا د هر ډول دپاره متغيرونه او کميت ټاکونکي، او (په عام
ډول) د هر ډول دپاره يو عينيت لرو۔ تابعې او اړيکې هم د هغو
څيزونو د ډولونو له مخې ``ټايپ'' کېږي چې د ارګومېنټونو په توګه ئې
اخيستلے شي۔ له دې پرته د لومړۍ درجې د منطق عادي قاعدې ساتل کېږي، او
د کميت ټاکونکو د قاعدو بڼې د هر ډول دپاره تکرارېږي۔
\olpseganchor{OLP-0176-B007}{end}

\olpseganchor{OLP-0176-B008}{start}
د بېلګې په توګه، د نړيوالو اړيکو د مطالعې دپاره ښايي داسې ژبه وټاکو
چې دوه ډوله څيزونه، فرانسوي اتباع او جرمن اتباع، ولري۔ ښايي د لومړي
ډول د څيزونو دپاره يوه يوځاييزه اړيکه، ``شراب څښي''، د دويم ډول د
څيزونو دپاره بله يوځاييزه اړيکه، ``ورسټ خوري''، او يوه دوه ځاييزه
اړيکه، ``څو ملتيزه واده شوې جوړه جوړوي''، ولرو۔ وروستۍ اړيکه دوه
ارګومېنټونه اخلي: لومړے د لومړي ډول او دويم د دويم ډول وي۔ که \LR{$a$}،
\LR{$b$}، \LR{$c$} متغيرونه پر فرانسوي اتباعو او \LR{$x$}، \LR{$y$}، \LR{$z$} متغيرونه پر
جرمن اتباعو وګرځوو، نو
\[
\lforall[a][\lforall[x][(\Atom{\Obj{MarriedTo}}{a,x} \lif
(\Atom{\Obj{DrinksWine}}{a} \lor \lnot \Atom{\Obj{EatsWurst}}{x}))]]
\]
دا وايي چې که کوم فرانسوے له يوه جرمن سره وادۀ شوے وي، نو يا هغه
فرانسوے شراب څښي يا هغه جرمن ورسټ نۀ خوري۔
\olpseganchor{OLP-0176-B008}{end}

\olpseganchor{OLP-0176-B009}{start}
د ډېرو ډولونو منطق په طبيعي ډول د لومړۍ درجې په منطق کښې ځايېدلے شي:
د ډېرو ډولونو د د تعريف ساحې ټول څيزونه په يوه د لومړۍ درجې د تعريف ساحه
کښې يوځاے کوو، د ډولونو د پېژندلو دپاره يوځاييز پريديکاتونه کاروو
او کميت ټاکونکي نسبي کوو۔ د بېلګې په توګه، د پاسنۍ بېلګې اړونده د
لومړۍ درجې ژبه به د نورو بيان شويو اړيکو تر څنګ ``\LR{$\Obj{German}$}''
او ``\LR{$\Obj{French}$}'' يوځاييز پريديکاتونه ولري، او د ډول شرطونه
به ترې لرې وي۔ يو ډول‌لرونکے کميت ټاکونکے \LR{$\lforall[x][!A]$}، چې \LR{$x$}
د جرمن ډول متغير وي، داسې ژباړل کېږي:
\[
\lforall[x][(\Atom{\Obj{German}}{x} \lif !A)].
\]
بايد داسې بديهي اصول هم ورزيات کړو چې ډولونه بېل تضمين کړي---لکه
\LR{$\lforall[x][\lnot (\Atom{\Obj{German}}{x} \land
  \Atom{\Obj{French}}{x})]$}---او هم داسې بديهي اصول چې تضمين کړي
``شراب څښي'' يوازې د \LR{$\Atom{\Obj{French}}{x}$} پريديکات پوره کوونکو
څيزونو دپاره صدق کوي، او داسې نور۔ له دې تړونونو او بديهي اصولو سره
دا ښودل ګران نۀ دي چې د ډېرو ډولونو تړلي فارمولونه د لومړۍ درجې
تړلي فارمولونه ته، او د ډېرو ډولونو اشتقاقونه د لومړۍ درجې
اشتقاقونه ته ژباړل کېږي۔ همدارنګه، د ډېرو ډولونو جوړښتونه
د لومړۍ درجې اړوندو جوړښتونه ته او برعکس ``ژباړل'' کېږي؛ نو د
ډېرو ډولونو د منطق دپاره هم د بشپړتيا قضيه لرو۔
\olpseganchor{OLP-0176-B009}{end}

\olpseganchor{OLP-0177-B004}{start}
\olfileid{fol}{byd}{sol}
\olpseganchor{OLP-0177-B004}{end}

\olpseganchor{OLP-0177-B005}{start}
\olsection{د دويمې درجې منطق}
\olpseganchor{OLP-0177-B005}{end}

\olpseganchor{OLP-0177-B006}{start}
د دويمې درجې منطق ژبه يوازې د افرادو پر د تعريف ساحه کميت ټاکنې ته
اجازه نۀ ورکوي، بلکې د هغه د تعريف ساحه پر اړيکو هم کميت ټاکل پرې کېدے
شي۔ د لومړۍ درجې يوې ژبې~\LR{$\Lang{L}$} ته د هر~\LR{$k$} دپاره داسې
متغيرونه~\LR{$R$} ورزياتېږي چې پر \LR{$k$}-ځاييزو اړيکو ګرځي، او پر دغو
متغيرونو کميت ټاکنه هم روا کېږي۔ که \LR{$R$} د يوې \LR{$k$}-ځاييزې اړيکې
متغير وي او \LR{$t_1$}، \dots،~\LR{$t_k$} عادي (د لومړۍ درجې) ترمونه
وي، نو \LR{$\Atom{R}{t_1,\dots,t_k}$} يو اټومي فارمول دے۔ له دې پرته
د فارمولونه سټ هماغسې تعريفېږي لکه د لومړۍ درجې په منطق کښې، خو د
دويمې درجې کميت ټاکنې اضافي شرطونه لري۔ پام وکړئ چې عينيت
يوازې د لومړۍ درجې د ترمونو دپاره لرو: که \LR{$R$} او \LR{$S$} د يوه شان
ځايښت~\LR{$k$} اړيکيز متغيرونه وي، \LR{$\eq[R][S]$} د لاندې عبارت د لنډون
په توګه تعريفولے شو:
\[
\lforall[x_1 \dots][\lforall[x_k][(\Atom{R}{x_1, \dots, x_k} \liff
  \Atom{S}{x_1, \dots, x_k})]].
\]
\olpseganchor{OLP-0177-B006}{end}

\olpseganchor{OLP-0177-B007}{start}
د دويمې درجې منطق قاعدې يوازې د کميت ټاکونکو قاعدې د دويمې درجې
نويو متغيرونو ته غځوي۔ خو دلته بايد لږ پام وشي چې دا متغيرونه د
ژبې~\LR{$\Lang{L}$} له پريديکاتونه، او په عام ډول د~\LR{$\Lang{L}$} له
فارمولونه سره څنګه چلند کوي۔ اړيکيز متغيرونه د اړيکو د ډول
ترمونه دي، او د افرادو ترمونه د خپل ډول غړي ښيي۔ نو د لاندې بڼې
استنتاجونه لرو:
\[
!A(R) \Proves \lexists[R][!A(R)]
\]
خو که \LR{$\Lang{L}$} د حساب ژبه وي او د ثابتې اړيکې نښه~\LR{$<$} ولري، تمه
کېږي چې لاندې استنتاج هم معتبر وي:
\[
x < y \Proves \lexists[R][\Atom{R}{x,y}]
\]
يا د يوه ورکړل شوي فارمول~\LR{$!A$} دپاره:
\[
!A(x_1, \dots, x_k) \Proves \lexists[R][\Atom{R}{x_1,\dots,x_k}]
\]
په لا عام ډول، ښايي د لاندې بڼې استنتاجونه روا کول وغواړو:
\[
\Subst{!A}{\lambd[\vec x][!B(\vec x)]}{R} \Proves
\lexists[R][!A]
\]
دلته \LR{$\Subst{!A}{\lambd[\vec x][!B(\vec x)]}{R}$} هغه پايله ښيي چې
په~\LR{$!A$} کښې د \LR{$\Atom{R}{t_1,\dots,t_k}$} بڼې هر اټومي فارمول په
\LR{$!B(t_1, \dots, t_k)$} بدل شي۔ دا وروستۍ قاعده د {\em ادراک شېما} له
شتون سره برابره ده، يعنې د لاندې بڼې يو بديهي اصل:
\[
\lexists[R][\lforall[x_1, \dots, x_k][(!A(x_1, \dots, x_k) \liff
\Atom{R}{x_1, \dots, x_k})]],
\]
د دويمې درجې په ژبه کښې د هر فارمول~\LR{$!A$} دپاره يوه داسې بېلګه
چې \LR{$R$} پکې ازاد متغير نۀ وي۔ (تمرين: وښيئ چې که \LR{$R$} ته اجازه ورکړل
شي چې په~\LR{$!A$} کښې ازاد راشي، دا شېما ناسازګاره ده۔ د اړيکې د
اټومي فارمول نفي په پام کښې ونيسئ۔)
\olpseganchor{OLP-0177-B007}{end}

\olpseganchor{OLP-0177-B008}{start}
کله چې منطقپوهان د ``دويمې درجې منطق بديهي اصول'' يادوي، په عام ډول
د دويمې درجې د کميت ټاکونکو پر قاعدو او د ادراک پر شېما د لومړۍ درجې
منطق لږ تر لږه غځونه مرادوي۔ خو د دې بديهي اصولو او قاعدو د کمزورو
فرعي نظامونو مطالعه هم ډېر ځله په زړه پورې وي۔ د بېلګې په توګه، پام
وکړئ چې د ادراک بديهي شېما په خپل بشپړ عموميت کښې
\emph{امپريديکاتيف} ده: دا د داسې اړيکې
\LR{$\Atom{R}{x_1, \dots, x_k}$} د شتون دعوه روا کوي چې په داسې
فارمول ``تعريف'' شوې وي چې د دويمې درجې کميت ټاکونکي لري؛ او دا
کميت ټاکونکي د ټولو داسې اړيکو پر سټ ګرځي---هغه سټ چې خپله \LR{$R$} هم پکې
ده!{} د
شلمې پېړۍ د پيل شاوخوا د راسل تناقض ته يو عام غبرګون دا و چې دا ډول
تعريفونه ملامت وګڼل شي او د رياضياتو د بنسټونو په جوړولو کښې ترې ډډه
وشي۔ که په فارمول~\LR{$!A$} کښې د دويمې درجې کميت ټاکونکو کارونه منع
شي، د ادراک يوه {\em پريديکاتيف} بڼه ترلاسه کېږي چې څه ناڅه کمزورې
ده۔
\olpseganchor{OLP-0177-B008}{end}

\olpseganchor{OLP-0177-B009}{start}
د معنايي نظر له مخې، د دويمې درجې يو جوړښت داسې ګڼلے شو چې
د ژبې دپاره د لومړۍ درجې يو جوړښت او ورسره پر د تعريف ساحه د
هغو اړيکو يو سټ لري چې د دويمې درجې کميت ټاکونکي پرې ګرځي (په لا
دقيق ډول، د هر~\LR{$k$} دپاره د ځايښت~\LR{$k$} اړيکو يو سټ شته)۔ البته که
ادراک په اشتقاق نظام کښې شامل وي، دا اضافي شرط لرو چې د
``دويمې درجې په برخه'' کښې د ادراک د بديهي اصولو د پوره کولو دپاره
کافي اړيکې وي---ګنې د اشتقاق نظام سم نۀ دے!{} د کافي اړيکو د
تضمين يوه اسانه لار دا ده چې د دويمې درجې برخه د لومړۍ درجې پر برخه
له \emph{ټولو} اړيکو جوړه وګڼو۔ داسې جوړښت ته \emph{بشپړ}
ويل کېږي، او په يوه معنا په رښتيا د ژبې ``ټاکل شوے جوړښت''
دے۔ که خپل پام بشپړو جوړښتونه ته محدود کړو، هغه څه ترلاسه کوو
چې د دويمې درجې \emph{بشپړه} معنٰی پوهنه بلل کېږي۔ په دې حالت کښې د
يوه جوړښت مشخصول يوازې د لومړۍ درجې د برخې مشخصولو ته راکمېږي،
ځکه د دويمې درجې د برخې منځپانګه په ضمني ډول ترې معلومېږي۔
\olpseganchor{OLP-0177-B009}{end}

\olpseganchor{OLP-0177-B010}{start}
په لنډيز کښې، د دويمې درجې منطق په خبرو کښې څه ابهام شته۔ د
اشتقاق نظام له مخې ښايي له لاندې دواړو څخه يو مراد وي:
\begin{enumerate}
\item د دويمې درجې يو ``لږ تر لږه'' اشتقاق نظام، له ځينو
  ادراکي بديهي اصولو سره۔
\item د دويمې درجې ``معياري'' اشتقاق نظام، له بشپړ ادراک سره۔
\end{enumerate}
د معنٰی پوهنې له مخې هم ښايي له لاندې دواړو څخه يوه د پام وړ وي:
\begin{enumerate}
\item ``کمزورې'' معنٰی پوهنه، چې جوړښت پکې د لومړۍ درجې يوه
  برخه او د ادراک د بديهي اصولو د پوره کولو دپاره کافي لويه د دويمې
  درجې برخه لري۔
\item د دويمې درجې ``معياري'' معنٰی پوهنه، چې يوازې بشپړ
  جوړښتونه پکې څېړل کېږي۔
\end{enumerate}
کله چې منطقپوهان خپل مراد اشتقاق نظام يا معنٰی پوهنه نۀ
مشخصوي، اکثره د هر لړ دويم توکے مرادوي۔ لکه چې وبه وينو، د دې معنٰی
پوهنې ګټه دا ده چې د ډېرو طبيعي رياضيکي جوړښتونو قطعي توصيفونه راکوي؛
په ورته وخت کښې اشتقاق نظام ډېر پياوړے او د دې معنٰی پوهنې
دپاره سم دے۔ نيمګړتيا ئې دا ده چې اشتقاق نظام د دې معنٰی
پوهنې دپاره \emph{بشپړ نۀ} دے؛ په حقيقت کښې د دويمې درجې د بشپړې
معنٰی پوهنې دپاره \emph{هېڅ} په اغېزمن ډول ورکړل شوے اشتقاق
نظام بشپړ نۀ دے۔ له بلې خوا به ووينو چې اشتقاق نظام د کمزورې
معنٰی پوهنې دپاره \emph{بشپړ} دے؛ نو که يوه جمله د اثبات وړ نۀ وي،
يو \emph{داسې} جوړښت شته، چې لازم نۀ دے بشپړ وي، او جمله پکې
دروغ ده۔
\olpseganchor{OLP-0177-B010}{end}

\olpseganchor{OLP-0177-B011}{start}
د دويمې درجې منطق ژبه ډېره بډايه ده۔ يوځاييزې اړيکې د د تعريف ساحه له
فرعي سټونو سره يو شان ګڼلے شو، نو په ځانګړي ډول پر دغو سټونو کميت
ټاکلے شو؛ د بېلګې په توګه، د طبيعي عددونو دپاره استقرا په يوه بديهي
اصل بيانولے شو:
\[
\lforall[R][((\Atom{R}{\Obj{0}} \land \lforall[x][(\Atom{R}{x} \lif
    \Atom{R}{x'})]) \lif \lforall[x][\Atom{R}{x}])].
\]
که د حساب ژبه د \LR{$\Obj 0, \Obj \prime, +, \times$} او~\LR{$<$} نښې ولري، د
هغوى د چلند د بيانولو دپاره لاندې بديهي اصول ورزياتولے شو:
\begin{enumerate}
\item \LR{$\lforall[x][\lnot x' = \Obj 0]$}
\item \LR{$\lforall[x][\lforall[y][(x' = y' \lif x = y)]]$}
\item \LR{$\lforall[x][(x + \Obj 0) = x]$}
\item \LR{$\lforall[x][\lforall[y][(x + y') = (x + y)']]$}
\item \LR{$\lforall[x][(x \times \Obj 0) = \Obj 0]$}
\item \LR{$\lforall[x][\lforall[y][(x \times y') = ((x \times y) + x)]]$}
\item \LR{$\lforall[x][\lforall[y][(x < y \liff \lexists[z][y = (x + z')])]]$}
\end{enumerate}
دا ښودل ګران نۀ دي چې دغه بديهي اصول د پاسني استقرا له بديهي اصل سره،
په دې شرط چې د دويمې درجې بشپړه معنٰی پوهنه کاروو، د حساب د معياري
جوړښت~\LR{$\Struct{N}$}، د حساب د معياري مدل، قطعي توصيف وړاندې
کوي۔ که کوم جوړښت~\LR{$\Struct{M}$}
دا بديهي اصول رښتوني کوي، له~\LR{$\Nat$} څخه د~\LR{$\Struct{M}$} د تعريف ساحه ته
تابع~\LR{$f$} پر~\LR{$\Nat$} د عادي بازګښت په وسيله داسې تعريف کړئ چې
\LR{$f(0) = \Assign{\Obj 0}{M}$} او \LR{$f(x+1) = \Assign{\prime}{M}(f(x))$} وي۔
پر~\LR{$\Nat$} د عادي استقرا او په~\LR{$\Struct M$} کښې د بديهي اصولو (1)
او~(2) له صدق څخه وينو چې \LR{$f$} يو پر يو ده۔ د دې د ښودلو دپاره
چې \LR{$f$} پر هدف سټ بشپړه ده، \LR{$P$} د~\LR{$\Domain{M}$} د هغو غړو سټ ونيسئ چې
د~\LR{$f$} د قيمتونو په سټ کښې دي۔ ځکه چې \LR{$\Struct M$} بشپړ دے، \LR{$P$} د
دويمې درجې په د تعريف ساحه کښې دے۔ د~\LR{$f$} له جوړونې پوهېږو چې
\LR{$\Assign{\Obj 0}{M}$} په~\LR{$P$} کښې دے او \LR{$P$} د~\LR{$\Assign{\prime}{M}$}
لاندې تړلے دے۔ دا چې د استقرا بديهي اصل په~\LR{$\Struct M$} کښې (په ځانګړي
ډول د~\LR{$P$} دپاره) صدق کوي، تضمينوي چې \LR{$P$} د~\LR{$\Struct M$} د لومړۍ درجې
له ټول د تعريف ساحه سره برابر دے۔ نو \LR{$f$} دوه اړخيزه يو پر يو تابع ده۔ دا
ښودل هم زيات ګران نۀ دي چې \LR{$f$} هومومورفيزم دے؛ پر~\LR{$\Nat$} عادي استقرا
بيا بيا کاروو۔ پر صفر او جانشين د تابعې چلند د جوړونې له مخې
معلوم دے۔ د جمع او ضرب بازګښتي اصول د دواړو عملياتو د ساتنې
دپاره، او وروستے اصل د ترتيب د ساتنې او انعکاس دپاره کارېږي۔
\olpseganchor{OLP-0177-B011}{end}

\olpseganchor{OLP-0177-B012}{start}
د سټ تيورۍ په اصطلاحاتو کښې تابع يوازې د اړيکې يو ځانګړے ډول دے؛ د
بېلګې په توګه، يوځاييزه تابع~\LR{$f$} له داسې دوه ځاييزې اړيکې~\LR{$R$} سره يو
شان ګڼل کېدے شي چې \LR{$\lforall[x][\lexists![y][R(x,y)]]$} پوره کوي۔ نو
پر تابعو هم کميت ټاکلے شو۔ بيا د بشپړې معنٰی پوهنې په کارولو د
لامتناهي جوړښتونه ټولګے د هغو جوړښتونه~\LR{$\Struct M$} د
ټولګي په توګه تعريفولے شو چې د~\LR{$\Struct M$} له د تعريف ساحه څخه د هغه
يو مناسب فرعي سټ ته يو پر يو تابع لري:
\[
\lexists[f][(\lforall[x][\lforall[y][(\eq[f(x)][f(y)] \lif \eq[x][y])]]
  \land \lexists[y][\lforall[x][\eq/[f(x)][y]]])].
\]
دا جمله په کره ډول د ډېدېکېند په معنا نامتناهيت ښيي۔ که د ټاکنې
اصل وکاروو، هر نامتناهي سټ په دې معنا نامتناهي دے۔ په همدې فرض
د دې جملې نفي د متناهي جوړښتونه ټولګے تعريفوي۔
\olpseganchor{OLP-0177-B012}{end}

\olpseganchor{OLP-0177-B013}{start}
سربېره پر دې، د خطي ترتيب تعريف ته د لاندې شرط په زياتولو د ښو
ترتيبونو ټولګے تعريفولے شو:
\[
\lforall[P][(\lexists[x][\Atom{P}{x}] \lif \lexists[x][(\Atom{P}{x}
    \land \lforall[y][(y < x \lif \lnot \Atom{P}{y})])])].
\]
دا وايي چې هر ناتش سټ تر ټولو کوچنے غړے لري، که ``سټ'' له ``يوځاييزې
اړيکې'' سره يو شان وګڼو۔ د بلې بېلګې دپاره، د بې‌لوري ګراف تړلتيا داسې بيانولے شو چې څوکې په دوو ناتشو
برخو نۀ وېشل کېږي چې تر منځ ئې هېڅ څنډه ونۀ وي:
\[
\lnot \lexists[A][(\lexists[x][A(x)] \land \lexists[y][\lnot A(y)]
  \land \lforall[w][\lforall[z][((\Atom{A}{w} \land \lnot \Atom{A}{z})
      \lif \lnot \Atom{R}{w,z})]])].
\]
د يوې بلې بېلګې په توګه، په تمرين کښې هڅه کولے شئ د هغو متناهي
جوړښتونه ټولګے تعريف کړئ چې د تعريف ساحه ئې جفته اندازه لري۔ تر
دې هم په زړه پورې دا چې حقيقي عددونه د ناطقو عددونو لرونکي بشپړ مرتب
ميدان په توګه په قطعي ډول توصيفولے شو۔
\olpseganchor{OLP-0177-B013}{end}

\olpseganchor{OLP-0177-B014}{start}
په لنډ ډول، د دويمې درجې منطق د لومړۍ درجې له منطق څخه ډېر زيات
بيانوونکے دے۔ دا ښه خبر و؛ اوس بد خبر ته راځو۔ مخکې مو وويل چې د
دويمې درجې د بشپړې معنٰی پوهنې دپاره هېڅ بشپړ اغېزمن اشتقاق
نظام نشته۔ که ښه ئې وګڼو که بد، د لومړۍ درجې منطق ډېر خاصيتونه دلته
نشته، چې د متناهي شاهد خاصيت او د لوېنهايم--سکولم قضيې هم پکې دي۔
\olpseganchor{OLP-0177-B014}{end}

\olpseganchor{OLP-0177-B015}{start}
له بلې خوا، که څوک د دويمې درجې بشپړه معنٰی پوهنه پرېږدي او کمزورې
بڼه ئې ومني، نو د دويمې درجې لږ تر لږه اشتقاق نظام د دې معنٰی
پوهنې دپاره بشپړ دے۔ په نورو ټکو، که \LR{$\Proves$} داسې ولولو چې ``په لږ
تر لږه نظام کښې ثابتوي'' او \LR{$\Entails$} داسې چې ``په کمزورې معنٰی پوهنه
کښې معنايي استلزام لري''، ښودلے شو چې هر کله \LR{$\Gamma \Entails !A$} وي
نو \LR{$\Gamma \Proves !A$} وي۔ که څوک په اشتقاق نظام کښې ځانګړي
ادراکي بديهي اصول شاملول غواړي، بايد معنٰی پوهنه د دويمې درجې هغو
جوړښتونو ته محدوده کړي چې دا اصول پوره کوي: د بېلګې په توګه، که
\LR{$\Delta$} د ادراکي بديهي اصولو يو سټ وي (کېدے شي ټول وي)، نو که
\LR{$\Gamma \cup \Delta \Entails !A$} وي، \LR{$\Gamma \cup \Delta \Proves !A$}
هم وي۔ په ځانګړي ډول، که \LR{$!A$} د هغو ادراکي بديهي اصولو په کارولو د
اثبات وړ نۀ وي چې موږ ئې څېړو، د~\LR{$\lnot !A$} داسې يو مدل شته چې بيا هم
دغه ادراکي بديهي اصول پوره کوي۔
\olpseganchor{OLP-0177-B015}{end}

\olpseganchor{OLP-0177-B016}{start}
دا چې د بشپړتيا قضيه د کمزورې معنٰی پوهنې دپاره ولې صدق کوي، تر ټولو
اسانه ئې دا ده چې د دويمې درجې منطق په لاندې ډول د ډېرو ډولونو منطق
وګڼو۔ يو ډول د عادي ``لومړۍ درجې'' د تعريف ساحه په توګه تفسيرېږي، او
بيا د هر~\LR{$k$} دپاره د ``ځايښت~\LR{$k$} اړيکو'' د تعريف ساحه لرو۔ ژبې ته
جوړې دننه اړيکيزې نښې ``\LR{$\Atom{\Obj{true}_k}{R,x_1,\dots,x_k}$}''
ورکوو؛ مراد ئې دا دے چې \LR{$R$} د~\LR{$x_1$}، \dots،~\LR{$x_k$} دپاره صدق کوي،
چرته چې \LR{$R$} د ``\LR{$k$}-ځاييزې اړيکې'' د ډول متغير او \LR{$x_1$}،
\dots،~\LR{$x_k$} د لومړۍ درجې د ډول څيزونه دي۔ که د اړيکو ډول
د غړو په توګه اړيکې ولري، دوه غړي چې پر عين ټوپلونو صدق کوي يو
غړے دے۔ په څوډوليزه ژبه کښې دا شرط د هر ځايښت د پراختيايي
برابرۍ له اصل سره څرګندولے شو۔
\olpseganchor{OLP-0177-B016}{end}

\olpseganchor{OLP-0177-B017}{start}
د دې يو شان ګڼلو له مخې، د دويمې درجې کمزورې معنٰی پوهنه په اصل کښې د
ډېرو ډولونو د منطق عادي معنٰی پوهنه ده؛ او مخکې مو ليدلي چې د ډېرو
ډولونو منطق د لومړۍ درجې په منطق کښې ځايېدلے شي۔ نو د دواړو خواوو
ژباړو ته په پام، د دويمې درجې منطق کمزورے تصور په رښتيا د لومړۍ درجې
منطق يوه پټه بڼه ده چې د تعريف ساحه ئې هم ``څيزونه'' او هم ``اړيکې''
لري او مناسب بديهي اصول پرې واکمن دي۔
\olpseganchor{OLP-0177-B017}{end}

\olpseganchor{OLP-0178-B004}{start}
\olfileid{fol}{byd}{hol}
\olpseganchor{OLP-0178-B004}{end}

\olpseganchor{OLP-0178-B005}{start}
\olsection{د لوړې درجې منطق}
\olpseganchor{OLP-0178-B005}{end}

\olpseganchor{OLP-0178-B006}{start}
له لومړۍ درجې منطق څخه دويمې درجې منطق ته په تېرېدو مو وکړے شول چې
په صوري ژبه کښې د لومړۍ درجې د د تعريف ساحه د څيزونو د سټونو په هکله
خبرې وکړو۔ ولې همدلته ودرېږو؟ د بېلګې په توګه، د درېيمې درجې منطق
بايد موږ ته د څيزونو د سټونو د سټونو، يا ښايي آن د هغو سټونو، په
هکله د خبرو وس راکړي چې هم څيزونه او هم د څيزونو سټونه لري۔ د څلورمې
درجې منطق بيا د دې ډول څيزونو د سټونو په هکله خبرې راکوي۔ لکه چې
اټکل مو کړے وي، دا مفکوره په خپله خوښه هر څومره تکرارولے شو۔
\olpseganchor{OLP-0178-B006}{end}

\olpseganchor{OLP-0178-B007}{start}
په عملي کار کښې د لوړې درجې منطق ډېر ځله د اړيکو پر ځاے د تابعو په
اصطلاحاتو کښې د فارمول په بڼه وړاندې کېږي۔ (له طبيعي يو شان ګڼلو
سره دا توپير بنسټيز نۀ دے۔) چې ځينې بنسټيز ``ډولونه'' \LR{$A$}، \LR{$B$}،
\LR{$C$}،~\dots راکړل شي (چې اوس به ئې ``ټايپونه'' بولو)، د لاندې شرط په
ورکولو نوي ټايپونه جوړولے شو:
\begin{quote}
که \LR{$\sigma$} او \LR{$\tau$} متناهي ټايپونه وي، \LR{$\sigma \to \tau$} هم
متناهي ټايپ دے۔
\end{quote}
ټايپونه داسې نحوي ``نښکې'' وګڼئ چې هغه څيزونه ډلبندي کوي چې په خپل
د تعريف ساحه کښې ئې غواړو؛ \LR{$\sigma \to \tau$} هغه څيزونه بيانوي چې د
ټايپ~\LR{$\sigma$} څيزونه د ټايپ~\LR{$\tau$} څيزونو ته وړونکې تابعې دي۔ د
بېلګې په توګه، ښايي د ``رښتوني'' او ``دروغ'' د رښتياوالي قيمتونو يو
ټايپ~\LR{$\Omega$} او د طبيعي عددونو يو ټايپ~\LR{$\Nat$} وغواړو۔ په دې حالت
کښې د ټايپ~\LR{$\Nat \to \Omega$} څيزونه د يوځاييزو اړيکو يا د~\LR{$\Nat$}
فرعي سټونو په توګه ګڼلے شئ؛ د ټايپ~\LR{$\Nat \to \Nat$} څيزونه له طبيعي
عددونو څخه طبيعي عددونو ته تابعې دي؛ او د ټايپ
\LR{$(\Nat \to \Nat) \to \Nat$} څيزونه ``فانکشنلونه'' دي، يعنې د لوړې
درجې داسې تابعې چې تابعې اخلي او عددونه ورکوي۔
\olpseganchor{OLP-0178-B007}{end}

\olpseganchor{OLP-0178-B008}{start}
لکه د دويمې درجې منطق، د لوړې درجې منطق هم د ډېرو ډولونو يو داسې
منطق ګڼلے شو چې د څيز د هر مطلوب ټايپ دپاره يو ډول لري۔ خو اکثره دا
لا روښانه وي چې د لوړو ټايپونو نحو له بنسټه تعريف کړو۔ د بېلګې په
توګه، د متناهي ټايپونو يو سټ په استقرايي ډول داسې تعريفولے شو:
\begin{enumerate}
\item \LR{$\Nat$} يو متناهي ټايپ دے۔
\item که \LR{$\sigma$} او \LR{$\tau$} متناهي ټايپونه وي، \LR{$\sigma
  \to \tau$} هم متناهي ټايپ دے۔
\item که \LR{$\sigma$} او \LR{$\tau$} متناهي ټايپونه وي، \LR{$\sigma \times
  \tau$} هم متناهي ټايپ دے۔
\end{enumerate}
په وجداني ډول، \LR{$\Nat$} د طبيعي عددونو ټايپ، \LR{$\sigma \to \tau$} له
\LR{$\sigma$} څخه~\LR{$\tau$} ته د تابعو ټايپ، او \LR{$\sigma \times \tau$} د هغو
جوړو ټايپ ښيي چې يو څيز ئې له~\LR{$\sigma$} او بل له~\LR{$\tau$} څخه وي۔ بيا د
ترمونو يو سټ په استقرايي ډول داسې تعريفولے شو:
\begin{enumerate}
\item د هر ټايپ~\LR{$\sigma$} دپاره د \LR{$x$}، \LR{$y$}، \LR{$z$}، \dots په څېر د
  ټايپ~\LR{$\sigma$} متغيرونو يوه زېرمه شته۔
\item \LR{$\Obj 0$} د ټايپ~\LR{$\Nat$} يو ترم دے۔
\item \LR{$\Obj S$} (ناستي) د ټايپ~\LR{$\Nat \to \Nat$} يو ترم دے۔
\item که \LR{$s$} د ټايپ~\LR{$\sigma$} ترم او \LR{$t$} د ټايپ
  \LR{$\Nat \to (\sigma \to \sigma)$} ترم وي، \LR{$\Obj{R}_{st}$} د ټايپ
  \LR{$\Nat \to \sigma$} يو ترم دے۔
\item که \LR{$s$} د ټايپ~\LR{$\tau \to \sigma$} ترم او \LR{$t$} د ټايپ~\LR{$\tau$}
  ترم وي، \LR{$s(t)$} د ټايپ~\LR{$\sigma$} يو ترم دے۔
\item که \LR{$s$} د ټايپ~\LR{$\sigma$} ترم او \LR{$x$} د ټايپ~\LR{$\tau$} متغير وي،
  \LR{$\lambd[x][s]$} د ټايپ~\LR{$\tau \to \sigma$} يو ترم دے۔
\item که \LR{$s$} د ټايپ~\LR{$\sigma$} ترم او \LR{$t$} د ټايپ~\LR{$\tau$} ترم وي،
  \LR{$\tuple{s, t}$} د ټايپ~\LR{$\sigma \times \tau$} يو ترم دے۔
\item که \LR{$s$} د ټايپ~\LR{$\sigma \times \tau$} يو ترم وي، \LR{$p_1(s)$} د
  ټايپ~\LR{$\sigma$} يو ترم او \LR{$p_2(s)$} د ټايپ~\LR{$\tau$} يو ترم دے۔
\end{enumerate}
په وجداني ډول، \LR{$\Obj{R}_{st}$} هغه تابع ښيي چې په بازګشتي ډول داسې
تعريف شوې وي:
\begin{align*}
\Obj{R}_{st}(0) & = s \\
\Obj{R}_{st}(x+1) & = t(x, R_{st}(x)),
\end{align*}
\LR{$\tuple{s, t}$} هغه جوړه ښيي چې لومړے جز ئې~\LR{$s$} او دويم جز ئې~\LR{$t$}
دے، او \LR{$p_1(s)$} او~\LR{$p_2(s)$} د~\LR{$s$} لومړے او دويم غړي
(``پروجيکشنونه'') ښيي۔ په پاى کښې، \LR{$\lambd[x][s]$} هغه تابع~\LR{$f$} ښيي
چې داسې تعريف شوې وي:
\[
f(x) = s
\]
د ټايپ~\LR{$\tau$} د هر~\LR{$x$} دپاره؛ نو توکے~(6) د ادراک يوه بڼه راکوي او
د ترمونو په وسيله د تابعو تعريفول شوني کوي۔ فارمولونه د يوه شان
ټايپ د ترمونو تر منځ له عينيت ويناوو~\LR{$\eq[s][t]$}، عادي بياني
تړونکو او د لوړو ټايپونو له کميت ټاکنې څخه جوړېږي۔ بيا د نظام بديهي
اصول هغه بنسټيزې معادلې ګڼلے شو چې پر پاس تعريف شويو ترمونو واکمنې
دي، او ورسره د کميت ټاکونکو او عينيت لرونکي منطق عادي قاعدې۔
\olpseganchor{OLP-0178-B008}{end}

\olpseganchor{OLP-0178-B009}{start}
که د متناهي ټايپونو نظام ته د رښتياوالي د قيمتونو ټايپ~\LR{$\Omega$}
ورزيات شي، بايد د هغه پر کارونه واکمن بديهي اصول هم شامل شي۔ په حقيقت
کښې، که لږ ځير واوسو، پېچلي فارمولونه په بشپړ ډول لرې کولے او د
ټايپ~\LR{$\Omega$} په ترمونو ئې بدلولے شو!{} بيا د ثبوت نظام هم ورسره سم
بدلېدے شي۔ پايله په اصل کښې د \emph{ټايپونو ساده تيوري} ده چې الونزو
چرچ په ۱۹۳۰مو کلونو کښې وړاندې کړې وه۔
\olpseganchor{OLP-0178-B009}{end}

\olpseganchor{OLP-0178-B010}{start}
لکه د دويمې درجې منطق، د لوړو ټايپونو د معنٰی پوهنې هم بېلابېلې بڼې
شته چې ښايي کارول ئې وغواړو۔ په بشپړه بڼه کښې د ټايپ
\LR{$\sigma \to \tau$} متغيرونه د ټايپ~\LR{$\sigma$} له څيزونو څخه د
ټايپ~\LR{$\tau$} څيزونو ته د \emph{ټولو} تابعو پر سټ ګرځي۔ لکه چې تمه
کېږي، دا معنٰی پوهنه دومره پياوړې ده چې يو بشپړ او اغېزمن
اشتقاق نظام نۀ مني۔ خو کمزورې معنٰی پوهنه هم څېړلے شو چې
جوړښت پکې د هر ټايپ~\LR{$\tau$} دپاره د غړو سټونه~\LR{$T_\tau$} او
ورسره د تطبيق، پروجيکشن او داسې نورو مناسب عملونه لري۔ که جزئيات سم
بشپړ شي، د پورته بيان شويو اشتقاق نظامونو د ډولونو دپاره د
بشپړتيا قضيې ترلاسه کولے شو۔
\olpseganchor{OLP-0178-B010}{end}

\olpseganchor{OLP-0178-B011}{start}
د لوړو ټايپونو منطق ځکه زړه راښکونکے دے چې يو داسې چوکاټ برابروي چې
د رياضياتو لويه برخه پکې په طبيعي ډول ځايولے شو: له~\LR{$\Nat$} څخه په
پيل حقيقي عددونه، مسلسل تابعې او داسې نور تعريفولے شو۔ دا منطق د
شهودي منطق په چاپېريال کښې په ځانګړي ډول زړه راښکونکے هم دے، ځکه
ټايپونه روښانه ``تعميري'' تفسيرونه لري۔ په حقيقت کښې د لوړو ټايپونو
د معنٰی پوهنې تعميري بڼې (چې د کلاسیک منطق پر ځاے پر شهودي منطق
ولاړې وي) جوړولے شو؛ دا بڼې دغه تعميري تفسيرونه ډېر ښه روښانوي او په
ډېرو لارو د خپلو کلاسیکو سيالانو په پرتله په زړه پورې دي۔
\olpseganchor{OLP-0178-B011}{end}

\olpseganchor{OLP-0179-B004}{start}
\olfileid{fol}{byd}{il}
\olpseganchor{OLP-0179-B004}{end}

\olpseganchor{OLP-0179-B005}{start}
\olsection{شهودي منطق}
\olpseganchor{OLP-0179-B005}{end}

\olpseganchor{OLP-0179-B006}{start}
د دويمې او لوړې درجې د منطق پر خلاف، د لومړۍ درجې شهودي منطق د
کلاسیکې بڼې يو محدود شوے نظام دے چې د لا ``تعميري'' ډول استدلال د مدل
کولو دپاره جوړ شوے دے۔ لاندې بېلګې ښايي ځينې بنسټيزې انګېزې روښانه کړي۔
\olpseganchor{OLP-0179-B006}{end}

\olpseganchor{OLP-0179-B007}{start}
فرض کړئ يوه ورځ يو څوک درته راشي او اعلان وکړي چې طبيعي عدد~\LR{$x$} ئې
ټاکلے دے او دا خاصيت لري: که \LR{$x$} اولي وي، د ريمان فرضيه رښتونې ده؛
او که \LR{$x$} مرکب وي، د ريمان فرضيه دروغ ده۔ دا خو لوے زېرے دے!{} د
ريمان د فرضيې رښتياوالے يا دروغوالے د رياضياتو له لويو ناحل شويو
پوښتنو څخه دے، او داسې ښکاري چې دې کس مسئله محاسبې ته راکمه کړې،
يعنې يوازې بايد معلومه شي چې يو ځانګړے عدد اولي دے که نۀ۔
\olpseganchor{OLP-0179-B007}{end}

\olpseganchor{OLP-0179-B008}{start}
د~\LR{$x$} جادويي قيمت څو دے؟ هغه ئې داسې بيانوي: \LR{$x$} هغه طبيعي عدد دے
چې که د ريمان فرضيه رښتونې وي له~\LR{$7$} سره برابر دے، او که نۀ،
له~\LR{$9$} سره۔
\olpseganchor{OLP-0179-B008}{end}

\olpseganchor{OLP-0179-B009}{start}
تاسو په غوسه خپلې پيسې بېرته غواړئ۔ له کلاسیک نظره پاسنے بيان په
رښتيا د~\LR{$x$} يو يوازينے قيمت ټاکي؛ خو څه چې تاسو ئې په اصل کښې غواړئ
د~\LR{$x$} داسې قيمت دے چې \emph{په څرګند ډول} درکړل شوے وي۔
\olpseganchor{OLP-0179-B009}{end}

\olpseganchor{OLP-0179-B010}{start}
د يوې بلې، ښايي لږې مصنوعي، بېلګې دپاره لاندې پوښتنه وګورئ۔ پوهېږو
چې يو غير ناطق عدد ناطق توان ته پورته کېدے او ناطقه پايله ورکولے شي؛
د بېلګې په توګه، \LR{$\sqrt{2}^2 = 2$}۔ دومره څرګنده نۀ ده چې يو غير ناطق
عدد \emph{غير ناطق} توان ته پورته شي او بيا هم ناطقه پايله ورکړي که
نۀ۔ لاندې قضيه مثبت ځواب ورکوي:
\olpseganchor{OLP-0179-B010}{end}

\olpseganchor{OLP-0179-B011}{start}
\begin{thm}
داسې غير ناطق عددونه \LR{$a$} او \LR{$b$} شته چې \LR{$a^b$} ناطق وي۔
\end{thm}
\olpseganchor{OLP-0179-B011}{end}

\olpseganchor{OLP-0179-B012}{start}
\begin{proof}
\LR{$\sqrt{2}^{\sqrt{2}}$} وګورئ۔ که دا ناطق وي، کار بشپړ دے: کولے شو
\LR{$a = b = \sqrt{2}$} ونيسو۔ که نۀ وي، دا غير ناطق دے۔ بيا لرو:
\[
(\sqrt{2}^{\sqrt{2}})^{\sqrt{2}} = \sqrt{2}^{\sqrt{2} \cdot
  \sqrt{2}} = \sqrt{2}^2 = 2,
\]
چې خامخا ناطق دے۔ نو په دې حالت کښې \LR{$a$} له
\LR{$\sqrt{2}^{\sqrt{2}}$} سره او \LR{$b$} له~\LR{$\sqrt 2$} سره برابر ونيسئ۔
\end{proof}
\olpseganchor{OLP-0179-B012}{end}

\olpseganchor{OLP-0179-B013}{start}
آيا دا يو معتبر ثبوت دے؟ زياتره رياضيپوهان ئې معتبر ګڼي۔ خو دلته بيا
هم څه ناڅه نۀ قانع کوونکې خبره شته: د حقيقي عددونو د يوې جوړې شتون
مو په يوه خاصيت سره ثابت کړے، بې له دې چې ويلے شو دا \emph{کومه}
جوړه ده۔ همدا پايله داسې هم ثابتېدے شي چې جوړه~\LR{$a$}، \LR{$b$} پخپله په
ثبوت کښې \emph{ورکړل شوې} وي: \LR{$a = \sqrt{3}$} او \LR{$b = \log_3 4$}
ونيسئ۔ بيا
\[
a^b = \sqrt{3}^{\log_3 4} = 3^{1/2 \cdot \log_3 4} = (3^{\log_3
  4})^{1/2} = 4^{1/2}= 2,
\]
ځکه چې \LR{$3^{\log_3 x} = x$}۔
\olpseganchor{OLP-0179-B013}{end}

\olpseganchor{OLP-0179-B014}{start}
شهودي منطق د داسې استدلال د مدل کولو دپاره جوړ شوے چې د لومړي ثبوت
په څېر ګامونه نۀ مني۔ د~\LR{$!A(x)$} پوره کوونکي~\LR{$x$} د شتون ثابتول په دې
معنا دي چې يو ځانګړے~\LR{$x$} او د هغه له خوا د~\LR{$!A$} د پوره کولو ثبوت
ورکړئ، لکه په دويم ثبوت کښې۔ دا ثابتول چې \LR{$!A$} يا \LR{$!B$} صدق کوي،
غواړي چې له دواړو څخه يو په رښتيا ثابت کړے شئ۔
\olpseganchor{OLP-0179-B014}{end}

\olpseganchor{OLP-0179-B015}{start}
په صوري ډول، د لومړۍ درجې شهودي منطق هغه وخت ترلاسه کېږي چې د لومړۍ
درجې د منطق اشتقاق نظام په يوه ټاکلې طريقه محدود کړئ۔ همداسې
د دويمې يا لوړې درجې منطق شهودي بڼې هم شته۔ له رياضيکي نظره دا يوازې
صوري استنتاجي نظامونه دي، خو لکه مخکې چې وويل شول، موخه ئې د رياضيکي
استدلال يو ډول مدل کول دي۔ څوک دا هغه استدلال ګڼلے شي چې د رياضياتو
يو ځانګړے فلسفي دريځ ئې روا بولي (لکه د بروور شهوديت)؛ څوک ئې د
رياضيکي استدلال يو لا ``ملموس'' او قانع کوونکے ډول ګڼلے شي (د بيشپ د
تعميريت په دود)؛ او دا بحث هم کېدے شي چې صوري بيان غيرصوري انګېزه په
رښتيا رانغاړي که نۀ۔ خو هر فلسفي دريځ چې ولرو، شهودي منطق د يوه صوري
وړاندې شوي منطق په توګه مطالعه کولے شو؛ او د هر دليل له مخې، ډېر
رياضيکي منطقپوهان دا مطالعه په زړه پورې ګڼي۔
\olpseganchor{OLP-0179-B015}{end}

\olpseganchor{OLP-0179-B016}{start}
د شهودي تړونکو يو غيرصوري تعميري تفسير شته چې په عام ډول د بي‌اېچ‌کې
تفسير بلل کېږي (د بروور، هايتنګ او کولموګوروف په نومونو نومول شوے)۔
دا داسې دے: د~\LR{$!A \land !B$} ثبوت د~\LR{$!A$} له يوه ثبوت او د~\LR{$!B$} له
يوه ثبوت څخه جوړه وي؛ د~\LR{$!A \lor !B$} ثبوت يا د~\LR{$!A$} ثبوت وي يا
د~\LR{$!B$}، او څرګند معلومات لرو چې کوم يو دے؛ د~\LR{$!A \lif !B$} ثبوت يوه
کړنلاره وي چې د~\LR{$!A$} ثبوت د~\LR{$!B$} په ثبوت بدلوي؛ د
\LR{$\lforall[x][!A(x)]$} ثبوت يوه کړنلاره وي چې د~\LR{$x$} د هر قيمت دپاره د
\LR{$!A(x)$} ثبوت بېرته ورکوي؛ او د~\LR{$\lexists[x][!A(x)]$} ثبوت د~\LR{$x$} له
يوه قيمت او دې ثبوت څخه جوړ وي چې هغه قيمت~\LR{$!A$} پوره کوي۔ دا تفسير
په هغو محاسبوي اصطلاحاتو هم بيانېدے شي چې د ``کري--هاورډ
آيزومورفيزم'' يا ``د ټايپونو په توګه د فارمولونه بېلګه'' بلل
کېږي: فارمول د معلوماتو يو ټاکلے ټايپ مشخصوونکے وګڼئ، او ثبوتونه
د همدې معلوماتي ټايپونو محاسبوي څيزونه وګڼئ چې راښيي اړوند
فارمول رښتونے دے۔
\olpseganchor{OLP-0179-B016}{end}

\olpseganchor{OLP-0179-B017}{start}
شهودي منطق ډېر ځله کلاسیک منطق ``د خارج وسط له اصل پرته'' ګڼل کېږي۔
لاندې قضيه دا خبره لا دقيقه کوي۔
\begin{thm}
په شهودي ډول لاندې بديهي شېمې يو له بل سره برابرې دي:
\begin{enumerate}
\item \LR{$(\lnot !A \lif \lfalse) \lif !A$}.
\item \LR{$!A \lor \lnot !A$}
\item \LR{$\lnot \lnot !A \lif !A$}
\end{enumerate}
\end{thm}
له دوو نورو څخه د هرې شېمې بېلګې ترلاسه کول په شهودي منطق کښې يو ښه
تمرين دے۔
\olpseganchor{OLP-0179-B017}{end}

\olpseganchor{OLP-0179-B018}{start}
د شهودي بياني منطق لومړني استنتاجي نظامونه، چې د بروور د شهوديت د
صوري بڼو په توګه وړاندې شوي وو، کولموګوروف، ګليوېنکو او هايتنګ په
خپلواکه توګه جوړ کړل۔ د لومړۍ درجې شهودي منطق (او د شهودي رياضياتو
د ځينو برخو) لومړۍ صوري بڼه هايتنګ وړاندې کړه۔ که څه هم يو شمېر په
کلاسیک ډول معتبرې شېمې په شهودي ډول معتبرې نۀ دي، ډېرې نورې معتبرې
دي۔
\olpseganchor{OLP-0179-B018}{end}

\olpseganchor{OLP-0179-B019}{start}
\emph{دوه‌ګونې نفي ژباړه} د کلاسیک او شهودي منطق تر منځ يوه مهمه
اړيکه بيانوي۔ دا په استقرايي ډول داسې تعريفېږي (\LR{$!A^N$} د کلاسیک
فارمول~\LR{$!A$} ``شهودي'' ژباړه وګڼئ):
\begin{align*}
!A^N & \ident \lnot\lnot !A \quad \text{\RL{د اټومي فارمولونه \LR{$!A$} دپاره}} \\
(!A \land !B)^N & \ident (!A^N \land !B^N) \\
(!A \lor !B)^N & \ident  \lnot\lnot (!A^N \lor !B^N) \\
(!A \lif !B)^N & \ident (!A^N \lif !B^N) \\
(\lforall[x][!A])^N & \ident \lforall[x][!A^N] \\
(\lexists[x][!A])^N & \ident \lnot\lnot\lexists[x][!A^N]
\end{align*}
کولموګوروف او ګليوېنکو د بياني منطق دپاره د دې ژباړې بڼې لرلې؛ د
پريديکاتي منطق دپاره ګوډل او ګېنڅن په خپلواکه توګه وړاندې کړه۔ لرو:
\olpseganchor{OLP-0179-B019}{end}

\olpseganchor{OLP-0179-B020}{start}
\begin{thm}
\begin{enumerate}
\item \LR{$!A \liff !A^N$} په کلاسیک ډول د اثبات وړ دے۔
\item که \LR{$!A$} په کلاسیک ډول د اثبات وړ وي، \LR{$!A^N$} په شهودي ډول د
  اثبات وړ دے۔
\end{enumerate}
\end{thm}
\olpseganchor{OLP-0179-B020}{end}

\olpseganchor{OLP-0179-B021}{start}
اوس لاندې خبرې اترې تصورولے شو۔ کلاسیک رياضيپوه: ``ما \LR{$!A$} ثابت
کړے!'' شهودي رياضيپوه: ``ستا ثبوت معتبر نۀ دے۔ تا په اصل کښې
\LR{$!A^N$} ثابت کړے دے۔'' کلاسیک رياضيپوه: ``ماته سمه ده!'' د کلاسیک
رياضيپوه له نظره شهودي رياضيپوه يوازې په وړو توپيرونو نښتے، ځکه دواړه
سره برابر دي۔ خو شهودي رياضيپوه ټينګار کوي چې توپير شته۔
\olpseganchor{OLP-0179-B021}{end}

\olpseganchor{OLP-0179-B022}{start}
پام وکړئ چې پاسنۍ ژباړه يوازې له سوچه منطق سره کار لري؛ دا پوښتنه نۀ
څېړي چې د کلاسیکو او شهودي رياضياتو دپاره مناسب \emph{غيرمنطقي}
بديهي اصول کوم دي يا د دواړو تر منځ څه اړيکه ده۔ خو د پاسنۍ قضيې
لاندې لږ غځونه څه ګټور معلومات راکوي:
\olpseganchor{OLP-0179-B022}{end}

\olpseganchor{OLP-0179-B023}{start}
\begin{thm}
که \LR{$\Gamma$} په کلاسیک ډول \LR{$!A$} ثابتوي، \LR{$\Gamma^N$} په شهودي ډول
\LR{$!A^N$} ثابتوي۔
\end{thm}
\olpseganchor{OLP-0179-B023}{end}

\olpseganchor{OLP-0179-B024}{start}
په نورو ټکو، که \LR{$!A$} له ځينو فرضيو څخه په کلاسیک ډول د اثبات وړ وي،
\LR{$!A^N$} د هغو له دوه‌ګونې نفي ژباړو څخه د اثبات وړ دے۔
\olpseganchor{OLP-0179-B024}{end}

\olpseganchor{OLP-0179-B025}{start}
د دې د ښودلو دپاره چې يوه جمله يا بياني فارمول په شهودي ډول
معتبر دے، يوازې يو ثبوت ورکول پکار دي۔ خو دا به څنګه ښيئ چې معتبر نۀ
دے؟ د دې دپاره داسې معنٰی پوهنې ته اړ يو چې سمه، او غوره دا چې بشپړه
وي۔ د کرېپکي معنٰی پوهنه دا اړتيا ښه پوره کوي۔
\olpseganchor{OLP-0179-B025}{end}

\olpseganchor{OLP-0179-B026}{start}
هماغه کار کولے شو چې د کلاسیک منطق دپاره مو وکړ: معنٰی پوهنه تعريف
کړو او سموالے او بشپړتيا ثابته کړو۔ خو لاندې توپير ته پام ارزښت لري۔
د کلاسیک منطق معنٰی پوهنه په يوه معنا ``څرګنده'' وه او د تړونکو په
معنا کښې ضمني پرته وه۔ که څه هم د کرېپکي د معنٰی پوهنې دپاره څه
وجداني انګېزه ورکولے شو، دا د هماغه ډول حتمي والي احساس نۀ راکوي۔
سربېره پر دې، د يوه کلاسیک جوړښت مفهوم يو طبيعي رياضيکي
مفهوم دے؛ نو يا د جوړښت مفهوم د کلاسیک د لومړۍ درجې منطق د
مطالعې وسيله ګڼلے شو، يا کلاسیک د لومړۍ درجې منطق د رياضيکي
جوړښتونه د مطالعې وسيله۔ د دې پر خلاف، د کرېپکي جوړښتونه
يوازې منطقي جوړونې ګڼل کېدے شي؛ داسې نۀ ښکاري چې خپلواک رياضيکي ارزښت
ولري۔
\olpseganchor{OLP-0179-B026}{end}

\olpseganchor{OLP-0179-B027}{start}
د بياني ژبې دپاره د کرېپکي جوړښت~\LR{$\mModel M = \tuple{W, R, V}$}
د يوه سټ~\LR{$W$}، پر~\LR{$W$} د يوه جزوي ترتيب~\LR{$R$} چې تر ټولو کوچنے
غړے ولري، او د بياني متغيرونو له يوې ``يکنواختې'' ګومارنې
څخه د~\LR{$W$} غړي ته جوړ دے۔ وجداني تصور دا دے چې د~\LR{$W$}
غړي ``نړۍګانې'' يا ``د پوهې حالتونه'' ښيي؛ غړے
\LR{$v \geq u$} د~\LR{$u$} يو ``ممکن راتلونکے حالت'' ښيي؛ او~\LR{$u$} ته ګومارل شوي
بياني متغيرونه هغه قضيې دي چې په حالت~\LR{$u$} کښې رښتونې پېژندل شوې دي۔
بيا د تحميل اړيکه~\LR{$\mSat{M}{!A}[w]$} دا اړيکه د ژبې ټولو
فارمولونه ته غځوي؛ \LR{$\mSat{M}{!A}[w]$} داسې ولولئ چې ``\LR{$!A$} په
حالت~\LR{$w$} کښې رښتونے دے''۔ اړيکه په استقرايي ډول داسې تعريفېږي:
\begin{enumerate}
\item \LR{$\mSat{M}{\Obj p_i}[w]$} هله او يوازې هله چې \LR{$\Obj p_i$} يو له
  هغو بياني متغيرونو څخه وي چې~\LR{$w$} ته ګومارل شوي۔
\item \LR{$\mSat/{M}{\lfalse}[w]$}۔
\item \LR{$\mSat{M}{(!A \land !B)}[w]$} هله او يوازې هله چې
  \LR{$\mSat{M}{!A}[w]$} او \LR{$\mSat{M}{!B}[w]$} وي۔
\item \LR{$\mSat{M}{(!A \lor !B)}[w]$} هله او يوازې هله چې
  \LR{$\mSat{M}{!A}[w]$} يا \LR{$\mSat{M}{!B}[w]$} وي۔
\item \LR{$\mSat{M}{(!A \lif !B)}[w]$} هله او يوازې هله چې هر کله
  \LR{$w' \geq w$} او \LR{$\mSat{M}{!A}[w']$} وي، نو \LR{$\mSat{M}{!B}[w']$} هم وي۔
\end{enumerate}
دا ښودل يو ښه تمرين دے چې
\LR{$\lnot (p \land q) \lif (\lnot p \lor \lnot q)$} په شهودي ډول معتبر
نۀ دے؛ د کرېپکي داسې جوړښت جوړ کړئ چې مقابله بېلګه ورکړي۔
\olpseganchor{OLP-0179-B027}{end}

\olpseganchor{OLP-0180-B004}{start}
\olfileid{fol}{byd}{mod}
\olpseganchor{OLP-0180-B004}{end}

\olpseganchor{OLP-0180-B005}{start}
\olsection{مودالي منطقونه}
\olpseganchor{OLP-0180-B005}{end}

\olpseganchor{OLP-0180-B006}{start}
د يوې شرطي جملې دا بېلګه وګورئ:
\begin{quote}
  که جېرمي په هغې کوټه کښې يوازې وي، نو هغه مست او بربنډ دے او پر
  څوکيو ګډېږي۔
\end{quote}
دا داسې شرطي وينا ده چې ښايي په مادي لحاظ رښتونې خو بيا هم
غولوونکې وي، ځکه عبارت داسې ښيي چې د مقدم او پايلې تر منځ تر دې
پياوړې اړيکه شته چې يوازې يا مقدم دروغ وي يا تالي رښتونے۔ يعنې، عبارت
دا ښيي چې دعوه يوازې په دې ځانګړې نړۍ کښې رښتونې نۀ ده (چرته چې
ښايي ځکه په تش ډول رښتونې وي چې جېرمي په کوټه کښې يوازې نۀ دے)، بلکې
که مقدم رښتونے \emph{وے} نو پايله به هم رښتونې \emph{وه}۔ په نورو
ټکو، دا وينا داسې اخيستل کېدے شي چې دعوه يوازې په دې نړۍ کښې نۀ، بلکې
په هره ``ممکنه'' نړۍ کښې رښتونې ده؛ يا دا چې د دې ځانګړې نړۍ د محض
رښتياوالي پر خلاف، \emph{لازماً} رښتونې ده۔
\olpseganchor{OLP-0180-B006}{end}

\olpseganchor{OLP-0180-B007}{start}
مودالي منطق د دې ډول لزوم د معنا کولو دپاره جوړ شوے۔ مودالي بياني
منطق له عادي بياني منطق څخه د بکس عملګر په زياتولو ترلاسه کېږي؛ يعنې
که \LR{$!A$} فارمول وي، نو \LR{$\Box !A$} هم دے۔ په وجداني ډول \LR{$\Box !A$}
وايي چې \LR{$!A$} \emph{لازماً} رښتونے، يا په هره ممکنه نړۍ کښې رښتونے
دے۔ \LR{$\Diamond !A$} په عام ډول د \LR{$\lnot \Box \lnot !A$} لنډون ګڼل کېږي
او داسې لوستل کېږي چې \LR{$!A$} \emph{ممکن} رښتونے دے۔ البته موداليت په
پريديکاتي منطق کښې هم ورزياتېدلے شي۔
\olpseganchor{OLP-0180-B007}{end}

\olpseganchor{OLP-0180-B008}{start}
کرېپکي جوړښتونه مودالي منطق ته معنٰی پوهنه ورکولے شي؛ په حقيقت
کښې کرېپکي دا معنٰی پوهنه لومړے د مودالي منطق په پام کښې نيولو سره
جوړه کړې وه۔ دلته پر جزوي ترتيبونو د محدودېدو پر ځاے په عام ډول د
``ممکنو نړۍګانو'' يو سټ~\LR{$P$} او د نړۍګانو تر منځ دوه ځاييزه
``د لاسرسي'' اړيکه \LR{$\Atom{R}{x,y}$} لرو۔ په وجداني ډول
\LR{$\Atom{R}{p,q}$} وايي چې نړۍ~\LR{$q$} له~\LR{$p$} سره سازګاره ده؛ يعنې که موږ
په نړۍ~\LR{$p$} کښې ``يو''، بايد دا امکان ومنو چې نړۍ به د~\LR{$q$} په څېر وه۔
\olpseganchor{OLP-0180-B008}{end}

\olpseganchor{OLP-0180-B009}{start}
مودالي منطق کله د ``مفهومي'' منطق په نامه يادېږي، د ``مصداقي'' منطق
پر خلاف۔ د کلاسیک منطق په څېر د يوه مصداقي منطق ټاکل شوې معنٰی پوهنه
يوازې يوې، ``واقعي'' نړۍ ته اشاره کوي؛ خو د يوه ``مفهومي'' منطق معنٰی
پوهنه پر يوې پراخې هستپوهنې تکيه کوي۔ د لزوم د جوړښت کولو تر
څنګ، موداليت د نورو ژبنيو جوړښتونو د جوړښت کولو دپاره هم
کارېدلے شي، او \LR{$\Box$} او \LR{$\Diamond$} د اړوند کارونې له مخې بيا تفسيرېږي۔
د بېلګې په توګه:
\begin{enumerate}
\item د اثبات‌وړتيا په منطق کښې \LR{$\Box !A$} داسې لوستل کېږي چې ``\LR{$!A$}
  د اثبات وړ دے'' او \LR{$\Diamond !A$} داسې چې ``\LR{$!A$} سازګار دے''۔
\item په معرفتي منطق کښې \LR{$\Box !A$} داسې لوستل کېدے شي چې ``زه په
  \LR{$!A$} پوهېږم'' يا ``زه پر \LR{$!A$} باور لرم''۔
\item په زماني منطق کښې \LR{$\Box !A$} داسې لوستل کېدے شي چې ``\LR{$!A$} تل
  رښتونے وي'' او \LR{$\Diamond !A$} داسې چې ``\LR{$!A$} کله کله رښتونے وي''۔
\end{enumerate}
\olpseganchor{OLP-0180-B009}{end}

\olpseganchor{OLP-0180-B010}{start}
غواړو منطق په هغو قاعدو او بديهي اصولو پياوړے کړو چې له موداليت سره
کار لري۔ د بېلګې په توګه، نظام~\LR{$\Log{S4}$} د بياني منطق عادي بديهي
اصول او قاعدې، او ورسره لاندې بديهي اصول لري:
\begin{align*}
& \Box (!A \lif !B) \lif (\Box !A \lif \Box !B)\\
& \Box !A \lif !A \\
& \Box !A \lif \Box \Box !A
\intertext{\RL{او دا قاعده هم: ``له \LR{$!A$} څخه \LR{$\Box !A$} پايله واخلئ''۔
  \LR{$\Log{S5}$} لاندې بديهي اصل ورزياتوي:}}
& \Diamond !A \lif \Box \Diamond !A
\end{align*}
د دې بديهي اصولو بدل ډولونه ښايي د بېلابېلو کارونو دپاره مناسب وي؛ د
بېلګې په توګه، S5 په عام ډول د منطقي لزوم د مفهوم ځانګړوونکے ګڼل
کېږي۔ ښه خبره دا ده چې اکثره داسې معنٰی پوهنه موندلے شو چې د کرېپکي
په جوړښتونه کښې د لاسرسي اړيکه په طبيعي لارو محدودوي او د هغې
په نسبت د اشتقاق نظام سم او بشپړ وي۔ د بېلګې په توګه،
\LR{$\Log{S4}$} د کرېپکي د هغو جوړښتونه له ټولګي سره سمون لري چې د
لاسرسي اړيکه ئې انعکاسي او انتقالي وي۔ \LR{$\Log{S5}$} د کرېپکي د هغو
جوړښتونه له ټولګي سره سمون لري چې د لاسرسي اړيکه ئې {\em عمومي}
وي، يعنې هره نړۍ له هرې بلې نړۍ څخه د لاسرسي وړ وي؛ نو \LR{$\Box !A$} هله
او يوازې هله صدق کوي چې \LR{$!A$} په هره نړۍ کښې صدق وکړي۔
\olpseganchor{OLP-0180-B010}{end}

\olpseganchor{OLP-0181-B004}{start}
\olfileid{fol}{byd}{oth}
\olpseganchor{OLP-0181-B004}{end}

\olpseganchor{OLP-0181-B005}{start}
\olsection{نور منطقونه}
\olpseganchor{OLP-0181-B005}{end}

\olpseganchor{OLP-0181-B006}{start}
لکه چې ښايي تر اوسه مو معلومه کړې وي، د نوي منطق جوړول ګران نۀ دي۔
تاسو هم خپله نحو جوړولے، د استنتاج يو نظام پنځولے، او ورته معنٰی پوهنه
ترتيبولے شئ۔ که غواړئ د اشتقاق نظام مو د معنٰی پوهنې په نسبت
بشپړ وي، ښايي لږ ځيرکتيا ته اړ شئ؛ او د نړۍ د قانع کولو دپاره هم څه
هڅه پکار ده چې منطق مو په رښتيا په زړه پورې دے۔ خو په بدل کښې ئې د
خپل منطق د رياضيکي او محاسبوي خاصيتونو په سپړلو ډېر ساعتونه پاکه او
خوندوره بوختيا لرلے شئ۔
\olpseganchor{OLP-0181-B006}{end}

\olpseganchor{OLP-0181-B007}{start}
وروستيو لسيزو د صوري منطقونو رښتينې چاودنه ليدلې ده۔ فزي منطق د
مبهمو خاصيتونو په هکله د استدلال د مدل کولو دپاره جوړ شوے۔ احتمالي
منطق د ناڅرګندتيا په هکله استدلال مدل کوي۔ ډيفالټ منطقونه او
نايکنواخته منطقونه د ماتېدونکو استدلالي بڼو د مدل کولو دپاره جوړ شوي؛
يعنې داسې ``معقول'' استنتاجونه چې د نويو معلوماتو په راتلو وروسته
ردېدلے شي۔ معرفتي منطقونه د پوهې په هکله استدلال، سببي منطقونه د سبب
او مسبب د اړيکو په هکله استدلال، او آن ``تکليفي'' منطقونه د اخلاقي
مکلفيتونو په هکله استدلال مدل کوي۔ دا نظامونه چې د فلسفي، رياضيکي يا
محاسبوي انګېزې له مخې معرفي شوي وي، ښايي د رياضيکي منطق، فلسفي منطق،
مصنوعي ځيرکتيا، ادراکي پوهنې يا بلې څانګې تر سرليک لاندې مطالعه شي۔
\olpseganchor{OLP-0181-B007}{end}

\olpseganchor{OLP-0181-B008}{start}
لړ همداسې اوږدېږي او امکانونه بې‌پايه ښکاري۔ ښايي هېڅکله د لايبنېڅ
دې خوب ته و نۀ رسېږو چې ټول انساني عقل محاسبې ته راکم کړو---خو دا مو
له هڅې څخه نۀ شي منع کولے۔
\olpseganchor{OLP-0181-B008}{end}

\olpseganchor{OLP-0182-B004}{start}
\olpart{mod}{مدل تيوري}
\olpseganchor{OLP-0182-B004}{end}

\olpseganchor{OLP-0182-B005}{start}
\begin{editorial}
  د مدل تيورۍ مواد نيمګړي او ازمېښتي دي۔ دا مهال دا يوازې د مدل
  تيورۍ په هکله د الډو انتونېلي له يادښتونو څخه اخيستل شوې بڼه ده، له
  هغو موضوعاتو پرته چې د لومړۍ درجې منطق په برخه کښې راغلې دي (تيورۍ،
  بشپړتيا او د متناهي شاهد خاصيت)۔ د درسي کتاب دپاره د ګټورتيا په موخه
  دا برخه لا ډېرې سريزې، انګېزې، سپړنې او تمرينونو ته اړتيا لري۔ اېنډي
  ارانا پلان لري چې په ځانګړي ډول پر دې برخه کار وکړي (\gitissue{65})۔
\end{editorial}
\olpseganchor{OLP-0182-B005}{end}

\olpseganchor{OLP-0183-B004}{start}
\olchapter{mod}{bas}{د مدل تيورۍ بنسټونه}
\olpseganchor{OLP-0183-B004}{end}















\olpseganchor{OLP-0183-B012}{start}
%
\olpseganchor{OLP-0183-B012}{end}

\olpseganchor{OLP-0184-B004}{start}
\olfileid{mod}{bas}{red}
\section{راکمونې او غځونې}
\olpseganchor{OLP-0184-B004}{end}

\olpseganchor{OLP-0184-B005}{start}
ډېر ځله د هغو ژبو، او د دغو ژبو د جوړښتونه پرتله ګټوره يا
اړينه وي چې ځينې نښې سره ګډې لري۔ تر ټولو عام حالت دا دے چې د يوې
ژبه~\LR{$\Lang{L}$} ټولې نښې د بلې ژبه~\LR{$\Lang{L'}$}
برخه هم وي، يعنې \LR{$\Lang{L} \subseteq \Lang{L'}$}۔ بيا يو
\LR{$\Lang{L}$}-جوړښت~\LR{$\Struct{M}$} تل د اضافي نښو په تفسيرولو
\LR{$\Lang{L'}$}-جوړښت ته غځېدلے شي، په داسې حال کښې چې د ګډو نښو
تفسيرونه هماغسې پاتې شي۔ له بلې خوا، له يوه
\LR{$\Lang{L'}$}-جوړښت~\LR{$\Struct{M'}$} څخه د هغو نښو د تفسيرونو په
``هېرولو'' يو \LR{$\Lang{L}$}-جوړښت ترلاسه کولے شو چې په~\LR{$\Lang{L}$} کښې
نۀ راځي۔
\olpseganchor{OLP-0184-B005}{end}

\olpseganchor{OLP-0184-B006}{start}
\begin{defn}
\ollabel{defn:reduct}
فرض کړئ \LR{$\Lang L \subseteq \Lang L'$}، \LR{$\Struct M$} يو
\LR{$\Lang L$}-جوړښت او \LR{$\Struct M'$} يو \LR{$\Lang L'$}-جوړښت
دے۔ \LR{$\Struct M$} د~\LR{$\Struct M'$} تر~\LR{$\Lang L$} پورې \emph{راکمونه}، او
\LR{$\Struct M'$} د~\LR{$\Struct M$} تر~\LR{$\Lang L'$} پورې \emph{غځونه} هله او
يوازې هله ده چې
\begin{enumerate}
\item \LR{$\Domain{M} = \Domain{M'}$}
\item د هر داسې ثابت سمبول~\LR{$c \in \Lang L$} دپاره،
  \LR{$\Assign{c}{M} = \Assign{c}{M'}$} وي۔
\item د هر داسې تابع~\LR{$f \in \Lang L$} دپاره،
  \LR{$\Assign{f}{M} = \Assign{f}{M'}$} وي۔
\item د هر داسې پريديکات~\LR{$P \in \Lang L$} دپاره،
  \LR{$\Assign{P}{M} = \Assign{P}{M'}$} وي۔
\end{enumerate}
\end{defn}
\olpseganchor{OLP-0184-B006}{end}

\olpseganchor{OLP-0184-B007}{start}
\begin{prop}
\ollabel{prop:reduct}
که يو \LR{$\Lang{L}$}-جوړښت~\LR{$\Struct{M}$} د يوه
\LR{$\Lang{L'}$}-جوړښت~\LR{$\Struct{M'}$} راکمونه وي، نو د~\LR{$\Lang{L}$}
د ټولو تړلي فارمولونه~\LR{$!A$} دپاره
\[
\Sat{M}{!A} \text{\RL{ هله او يوازې هله چې }} \Sat{M'}{!A}.
\]
\end{prop}
\olpseganchor{OLP-0184-B007}{end}

\olpseganchor{OLP-0184-B008}{start}
\begin{proof}
  تمرين۔
\end{proof}
\olpseganchor{OLP-0184-B008}{end}

\olpseganchor{OLP-0184-B009}{start}
\begin{prob}
\olref[mod][bas][red]{prop:reduct} ثابته کړئ۔
\end{prob}
\olpseganchor{OLP-0184-B009}{end}

\olpseganchor{OLP-0184-B010}{start}
\begin{defn}
که يو \LR{$\Lang{L}$}-جوړښت \LR{$\Struct{M}$} ولرو، او \LR{$\Lang{L'} =
\Lang{L} \cup \{P\}$} د~\LR{$\Lang{L}$} هغه غځونه وي چې يوازينے
\LR{$n$}-ځاييز پريديکات~\LR{$P$} ورزياتوي، او \LR{$R \subseteq \Domain{M}^n$}
يوه \LR{$n$}-ځاييزه اړيکه وي، نو د~\LR{$\Struct{M}$} هغه
غځونه~\LR{$\Struct{M'}$} چې \LR{$\Assign{P}{M'} = R$} وي، په
\LR{$\Expan{M}{R}$} ليکو۔
\end{defn}
\olpseganchor{OLP-0184-B010}{end}

\olpseganchor{OLP-0185-B004}{start}
\olfileid{mod}{bas}{sub}
\olsection{فرعي جوړښتونه}
\olpseganchor{OLP-0185-B004}{end}

\olpseganchor{OLP-0185-B005}{start}
د يوه جوړښت~\LR{$\Struct{M}$} د تعريف ساحه ښايي د بل
جوړښت~\LR{$\Struct{M'}$} فرعي سټ وي۔ خو څرګنده ده چې \LR{$\Struct{M}$} بايد
يوازې هغه وخت د~\LR{$\Struct{M'}$} يوه ``برخه'' وګڼو چې نۀ يوازې
\LR{$\Domain{M} \subseteq \Domain{M'}$} وي، بلکې \LR{$\Struct{M}$} او
\LR{$\Struct{M'}$} لږ تر لږه پر ګډې برخې~\LR{$\Domain{M}$} د ژبې د نښو په
تفسيرولو کښې هم ``سره سلا'' وي۔
\olpseganchor{OLP-0185-B005}{end}

\olpseganchor{OLP-0185-B006}{start}
\begin{defn}
\ollabel{defn:substructure}
د يوې ژبې~\LR{$\Lang L$} جوړښتونه \LR{$\Struct M$} او \LR{$\Struct M'$} راکړل شي۔
وايو \LR{$\Struct M$} د~\LR{$\Struct M'$} \emph{فرعي جوړښت}، او
\LR{$\Struct M'$} د~\LR{$\Struct M$} \emph{غځونه} ده، چې
\LR{$\Struct M \substruct \Struct M'$} ئې ليکو، هله او يوازې هله چې
\begin{enumerate}
\item \LR{$\Domain{M} \subseteq \Domain{M'}$} وي؛
\item د هر داسې ثابت \LR{$c \in \Lang L$} دپاره، \LR{$\Assign{c}{M} =
    \Assign{c}{M'}$} وي؛
\item د هر داسې \LR{$n$}-ځاييز تابع~\LR{$f \in \Lang L$} دپاره، د
  ټولو \LR{$a_1$}، \dots،~\LR{$a_n \in \Domain{M}$} دپاره
  \LR{$\Assign{f}{M}(a_1, \dots, a_n) = \Assign{f}{M'}(a_1, \dots, a_n)$}
  وي؛
\item د هر داسې \LR{$n$}-ځاييز پريديکات~\LR{$R \in \Lang L$} دپاره، د
  ټولو \LR{$a_1$}، \dots،~\LR{$a_n \in \Domain{M}$} دپاره
  \LR{$\langle a_1, \dots, a_n\rangle \in \Assign{R}{M}$} هله او يوازې
  هله چې \LR{$\langle a_1, \dots, a_n\rangle \in \Assign{R}{M'}$} وي۔
\end{enumerate}
\end{defn}
\olpseganchor{OLP-0185-B006}{end}

\olpseganchor{OLP-0185-B007}{start}
\begin{rem}
\ollabel{rem:substructure}
که ژبه ثابت يا تابعې ونۀ لري، نو هر
\LR{$N \subseteq \Domain{M}$} د~\LR{$\Struct M$} يو فرعي
جوړښت~\LR{$\Struct{N}$} ټاکي چې د تعريف ساحه ئې
\LR{$\Domain{N} = N$} دے؛ د دې دپاره \LR{$\Assign{R}{N} =
\Assign{R}{M} \cap N^n$} ږدو۔
\end{rem}
\olpseganchor{OLP-0185-B007}{end}

\olpseganchor{OLP-0185-B008}{start}
% prove something about this? Examples?
\olpseganchor{OLP-0185-B008}{end}

\olpseganchor{OLP-0186-B004}{start}
\olfileid{mod}{bas}{ove}
\olsection{له بريده اوښتنه}
\olpseganchor{OLP-0186-B004}{end}

\olpseganchor{OLP-0186-B005}{start}
\begin{thm}
\ollabel{overspill} که د جملو يو سټ~\LR{$\Gamma$} د هرې کچې متناهي
مدلونه ولري، نو يو نامتناهي مدل هم لري۔
\end{thm}
\olpseganchor{OLP-0186-B005}{end}

\olpseganchor{OLP-0186-B006}{start}
\begin{proof}
د~\LR{$\Gamma$} ژبه د شمېر وړ نويو ثابتونو \LR{$c_0$}، \LR{$c_1$}، \dots په
ورزياتولو وغځوئ او سټ \LR{$\Gamma \cup \{c_i \neq c_j : i \neq j\}$}
وګورئ۔ دا ويل چې \LR{$\Gamma$} د هرې کچې متناهي مدلونه لري، معنا ئې دا ده
چې د هر \LR{$m >0$} دپاره داسې \LR{$n\ge m$} شته چې \LR{$\Gamma$} د شمېر~\LR{$n$} يو مدل
لري۔ له دې پايله کېږي چې \LR{$\Gamma \cup \{c_i \neq c_j : i \neq j\}$}
په متناهي ډول د صدق وړ دے۔ د متناهي شاهد قضيې له مخې، \LR{$\Gamma \cup
\{c_i \neq c_j : i \neq j\}$} يو مدل~\LR{$\Struct M$} لري چې د تعريف ساحه
ئې بايد نامتناهي وي، ځکه ټولې نابرابرۍ \LR{$c_i \neq c_j$} پوره کوي۔
\end{proof}
\olpseganchor{OLP-0186-B006}{end}

\olpseganchor{OLP-0186-B007}{start}
\begin{prop}
\ollabel{inf-not-fo}
د لومړۍ درجې د هېڅ ژبې داسې جمله~\LR{$!A$} نشته چې په يوه
جوړښت~\LR{$\Struct M$} کښې هله او يوازې هله رښتونې وي چې د دغه
جوړښت د تعريف ساحه~\LR{$\Domain{M}$} نامتناهي وي۔
\end{prop}
\olpseganchor{OLP-0186-B007}{end}

\olpseganchor{OLP-0186-B008}{start}
\begin{proof}
که داسې~\LR{$!A$} موجوده وے، نفي~\LR{$\lnot !A$} به ئې په ټولو او يوازې په
متناهي جوړښتونه کښې رښتونې وه؛ نو د هرې کچې متناهي مدلونه به
ئې لرل خو نامتناهي مدل به ئې نۀ درلود، چې دا له~\olref{overspill}
سره په ټکر کښې ده۔
\end{proof}
\olpseganchor{OLP-0186-B008}{end}

\olpseganchor{OLP-0187-B004}{start}
\olfileid{mod}{bas}{iso}
\olsection{آيزومورف جوړښتونه}
\olpseganchor{OLP-0187-B004}{end}

\olpseganchor{OLP-0187-B005}{start}
د لومړۍ درجې جوړښتونه په دوو لارو يو شان کېدے شي۔ يوه لار دا
ده چې يو شان تړلي فارمولونه رښتونې کړي۔ داسې جوړښتونه
\emph{ابتدايي هم‌ارز} بولو۔ خو جوړښتونه ډېر سره توپير لرلے او بيا هم
يو شان تړلي فارمولونه رښتونې کولے شي---د بېلګې په توګه، يو ئې
د شمېر وړ کېدے شي او بل نۀ۔ لامل دا دے چې د يوه
جوړښت ډېر خاصيتونه د لومړۍ درجې په ژبو کښې نۀ شي بيانېدلے؛
يا ځکه چې ژبه پوره بډايه نۀ وي، يا د لومړۍ درجې منطق د بنسټيزو
محدوديتونو، لکه د لوېنهايم--سکولم د قضيې، له امله۔ نو بله او لا سخته
معنا چې جوړښتونه پکې يو شان کېدے شي دا ده چې له بنسټه هماغه وي:
يوازې په هغو څيزونو کښې توپير ولري چې ترې جوړ دي، خو په جوړښتي
خاصيتونو کښې نۀ۔ د دې خبرې د دقيقولو يوه لار د \emph{آيزومورفيزم}
مفهوم دے۔
\olpseganchor{OLP-0187-B005}{end}

\olpseganchor{OLP-0187-B006}{start}
\begin{defn}
\ollabel{defn:elem-equiv}
د يوې ژبه~\LR{$\Lang{L}$} دوه جوړښتونه \LR{$\Struct{M}$} او
\LR{$\Struct M'$} راکړل شي۔ وايو \LR{$\Struct{M}$} له~\LR{$\Struct M'$} سره
\emph{ابتدايي هم‌ارز} دے، چې \LR{$\Struct{M} \equiv \Struct M'$} ئې
ليکو، هله او يوازې هله چې د~\LR{$\Lang{L}$} د هرې تړلے فارمول~\LR{$!A$}
دپاره \LR{$\Sat{M}{!A}$} هله او يوازې هله چې \LR{$\Sat{M'}{!A}$} وي۔
\end{defn}
\olpseganchor{OLP-0187-B006}{end}

\olpseganchor{OLP-0187-B007}{start}
\begin{defn}
\ollabel{defn:isomorphism} د يوې ژبه~\LR{$\Lang L$} دوه
جوړښتونه \LR{$\Struct{M}$} او \LR{$\Struct M'$} راکړل شي۔ وايو
\LR{$\Struct{M}$} له~\LR{$\Struct M'$} سره \emph{آيزومورف} دے، چې
\LR{$\Struct{M} \simeq \Struct M'$} ئې ليکو، هله او يوازې هله چې داسې
تابع \LR{$h \colon \Domain{M} \to \Domain{M'}$} شته چې:
\begin{enumerate}
\item \LR{$h$} يو پر يو ده: که \LR{$h(x) = h(y)$} وي، نو \LR{$x = y$}؛
\item \LR{$h$} پر هدف سټ بشپړه ده: د هر \LR{$y \in \Domain{M'}$} دپاره داسې
  \LR{$x \in \Domain{M}$} شته چې \LR{$h(x) = y$}؛
\item \ollabel{defn:iso-const}د هر ثابت سمبول~\LR{$c$} دپاره:
  \LR{$h(\Assign{c}{M}) = \Assign{c}{M'}$}؛
\item \ollabel{defn:iso-pred}د هر \LR{$n$}-ځاييز پريديکات~\LR{$P$} دپاره:
  \[
  \tuple{a_1, \dots, a_n}\in \Assign{P}{M} \quad\text{\RL{هله او يوازې هله چې}}\quad
  \tuple{h(a_1), \dots, h(a_n)} \in \Assign{P}{M'};
  \]
\item \ollabel{defn:iso-func}د هر \LR{$n$}-ځاييز تابع~\LR{$f$} دپاره:
  \[
  h(\Assign{f}{M}(a_1, \dots, a_n)) =
  \Assign{f}{M'}(h(a_1), \dots, h(a_n)).
  \]
\end{enumerate}
\end{defn}
\olpseganchor{OLP-0187-B007}{end}

\olpseganchor{OLP-0187-B008}{start}
\begin{thm}
\ollabel{thm:isom}
که \LR{$\Struct{M} \iso \Struct M'$} وي، نو \LR{$\Struct{M} \elemequiv
\Struct{M'}$}۔
\end{thm}
\olpseganchor{OLP-0187-B008}{end}

\olpseganchor{OLP-0187-B009}{start}
\begin{proof}
فرض کړئ \LR{$h$} له~\LR{$\Struct{M}$} څخه پر~\LR{$\Struct M'$} يو آيزومورفيزم دے۔
د هرې ګومارنې~\LR{$s$} دپاره \LR{$h \circ s$} د~\LR{$h$} او~\LR{$s$} ترکيب، يعنې
په~\LR{$\Struct{M'}$} کښې هغه ګومارنه ده چې
\LR{$(h \circ s)(x) = h(s(x))$} وي۔ پر~\LR{$t$} او~\LR{$!A$} په استقرا لاندې
پياوړې دعوې ثابتولے شو:
\begin{enumerate}
  \item[a.] \LR{$h(\Value{t}{M}[s]) = \Value{t}{M'}[h\circ s]$}۔
  \item[b.] \LR{$\Sat{M}{!A}[s]$} هله او يوازې هله چې
    \LR{$\Sat{M'}{!A}[h \circ s]$}۔
\end{enumerate}
لومړۍ دعوه د~\LR{$t$} پر پېچلتيا په استقرا ثابتېږي۔
\begin{enumerate}
\item که \LR{$t \ident c$} وي، نو \LR{$\Value{c}{M}[s] = \Assign{c}{M}$} او
  \LR{$\Value{c}{M'}[h \circ s] = \Assign{c}{M'}$}۔ نو
  \LR{$h(\Value{t}{M}[s]) = h(\Assign{c}{M}) = \Assign{c}{M'}$} (د
  \olref{defn:isomorphism} له~\olref{defn:iso-const} څخه)
  \LR{$= \Value{t}{M'}[h \circ s]$}۔
\item که \LR{$t \ident x$} وي، نو \LR{$\Value{x}{M}[s] = s(x)$} او
  \LR{$\Value{x}{M'}[h \circ s] = h(s(x))$}۔ نو
  \LR{$h(\Value{x}{M}[s]) = h(s(x)) = \Value{x}{M'}[h \circ s]$}۔
\item که \LR{$t \ident f(t_1, \dots, t_n)$} وي، نو
  \begin{align*}
    \Value{t}{M}[s] & = \Assign{f}{M}(\Value{t_1}{M}[s], \dots, \Value{t_n}{M}[s]) \quad\text{\RL{او}}\\
  \Value{t}{M'}[h \circ s] & = \Assign{f}{M'}(\Value{t_1}{M'}[h \circ
    s], \dots, \Value{t_n}{M'}[h \circ s]).
  \end{align*}
  د استقرا
  فرض دا دے چې د هر~\LR{$i$} دپاره \LR{$h(\Value{t_i}{M}[s]) =
  \Value{t_i}{M'}[h\circ s]$}۔ نو
  \begin{align}
    h(\Value{t}{M}[s])
    & = h(\Assign{f}{M}(\Value{t_1}{M}[s], \dots, \Value{t_n}{M}[s])) \notag\\
    & = \Assign{f}{M'}(h(\Value{t_1}{M}[s]), \dots,
    h(\Value{t_n}{M}[s])) \ollabel{iso-1}\\
    & = \Assign{f}{M'}(\Value{t_1}{M'}[h \circ s], \dots,
    \Value{t_n}{M'}[h \circ s]) \ollabel{iso-2}\\
    & = \Value{t}{M'}[h\circ s] \notag
  \end{align}
  دلته \olref{iso-1} د~\olref{defn:isomorphism} له
  \olref{defn:iso-func} او \olref{iso-2} د استقرا له فرض څخه پايله
  کېږي۔
\end{enumerate}
د (b) برخې ثبوت د تمرين دپاره پرېښودل شوے دے۔
\olpseganchor{OLP-0187-B009}{end}

\olpseganchor{OLP-0187-B010}{start}
که \LR{$!A$} جمله وي، ګومارنې~\LR{$s$} او~\LR{$h \circ s$} بې اغېزې دي، او
\LR{$\Sat{M}{!A}$} هله او يوازې هله چې \LR{$\Sat{M'}{!A}$} لرو۔
\end{proof}
\olpseganchor{OLP-0187-B010}{end}

\olpseganchor{OLP-0187-B011}{start}
\begin{prob}
د~\olref[mod][bas][iso]{thm:isom} د (b) برخې ثبوت په تفصيل بشپړ
کړئ۔ پام وکړئ چې د~\olref[mod][bas][iso]{defn:isomorphism} د
آيزومورفيزم ځانګړوونکي پنځه خاصيتونه هر يو چېرته کارول کېږي۔
\end{prob}
\olpseganchor{OLP-0187-B011}{end}

\olpseganchor{OLP-0187-B012}{start}
\begin{defn}
د يوه جوړښت~\LR{$\Struct{M}$} \emph{اتومورفيزم} له~\LR{$\Struct{M}$} څخه پر
خپل ځان يو آيزومورفيزم دے۔
\end{defn}
\olpseganchor{OLP-0187-B012}{end}

\olpseganchor{OLP-0187-B013}{start}
\begin{prob}
وښيئ چې د هر جوړښت~\LR{$\Struct{M}$} دپاره، که \LR{$X$} د~\LR{$\Struct{M}$} يو
تعريفېدونکی فرعي سټ او \LR{$h$} د~\LR{$\Struct{M}$} يو اتومورفيزم وي، نو
\LR{$X = \Setabs{h(x)}{x \in X}$} (يعنې \LR{$X$} د~\LR{$h$} لاندې ثابت پاتې کېږي)۔
\end{prob}
\olpseganchor{OLP-0187-B013}{end}

\olpseganchor{OLP-0188-B004}{start}
\olfileid{mod}{bas}{thm}
\olpseganchor{OLP-0188-B004}{end}

\olpseganchor{OLP-0188-B005}{start}
\section{د يوه جوړښت تيوري}
\olpseganchor{OLP-0188-B005}{end}

\olpseganchor{OLP-0188-B006}{start}
هر جوړښت~\LR{$\Struct{M}$} ځينې تړلي فارمولونه رښتونې او ځينې
دروغې کوي۔ د هغو ټولو تړلي فارمولونه سټ چې دا جوړښت ئې رښتونې کوي، د
هغه \emph{تيوري} بلل کېږي۔ دا سټ په رښتيا يوه تيوري ده، ځکه د هغې هر
معنايي استلزام بايد د هغې په ټولو مدلونو، او له دې ډلې
په~\LR{$\Struct{M}$} کښې، رښتونے وي۔
\olpseganchor{OLP-0188-B006}{end}

\olpseganchor{OLP-0188-B007}{start}
\begin{defn}
  يو جوړښت~\LR{$\Struct M$} راکړل شي۔ د~\LR{$\Struct{M}$}
  \emph{تيوري} د هغو تړلي فارمولونه سټ~\LR{$\Theory{M}$} دے چې
  په~\LR{$\Struct{M}$} کښې رښتونې وي، يعنې
  \LR{$\Theory{M} = \Setabs{!A}{\Sat{M}{!A}}$}۔
\end{defn}
\olpseganchor{OLP-0188-B007}{end}

\olpseganchor{OLP-0188-B008}{start}
د ``تيورۍ'' اصطلاح په غيرصوري ډول د هغو تړلي فارمولونه سټونو دپاره
هم کاروو چې يو ټاکلے مطلوب تفسير لري، که د استنتاج لاندې تړلي وي او
که نۀ۔
\olpseganchor{OLP-0188-B008}{end}

\olpseganchor{OLP-0188-B009}{start}
\begin{prop}
د هر~\LR{$\Struct{M}$} دپاره \LR{$\Theory{M}$} بشپړه ده۔
\end{prop}
\olpseganchor{OLP-0188-B009}{end}

\olpseganchor{OLP-0188-B010}{start}
\begin{proof}
د هرې تړلے فارمول~\LR{$!A$} دپاره يا \LR{$\Sat{M}{!A}$} وي يا
\LR{$\Sat{M}{\lnot !A}$}؛ نو يا \LR{$!A \in \Theory{M}$} وي يا
\LR{$\lnot !A \in \Theory{M}$}۔
\end{proof}
\olpseganchor{OLP-0188-B010}{end}

\olpseganchor{OLP-0188-B011}{start}
\begin{prop}\ollabel{prop:equiv}
  که د هر~\LR{$!A \in \Theory{M}$} دپاره \LR{$\Struct{N} \models !A$} وي، نو
  \LR{$\Struct{M} \elemequiv \Struct{N}$}۔
\end{prop}
\olpseganchor{OLP-0188-B011}{end}

\olpseganchor{OLP-0188-B012}{start}
\begin{proof}
دا چې د ټولو \LR{$!A \in \Theory{M}$} دپاره \LR{$\Sat{N}{!A}$} دے،
\LR{$\Theory{M} \subseteq \Theory{N}$}۔ که \LR{$\Sat{N}{!A}$} وي، نو
\LR{$\Sat/{N}{\lnot !A}$}؛ ځکه نو \LR{$\lnot !A \notin \Theory{M}$}۔ دا چې
\LR{$\Theory{M}$} بشپړه ده، \LR{$!A \in \Theory{M}$}۔ نو
\LR{$\Theory{N} \subseteq \Theory{M}$}، او له دې
\LR{$\Struct{M} \elemequiv \Struct{N}$} لرو۔
\end{proof}
\olpseganchor{OLP-0188-B012}{end}

\olpseganchor{OLP-0188-B013}{start}
\begin{rem}\ollabel{remark:R}
  \LR{$\Struct{R} = \langle\Real, <\rangle$} وګورئ: دا هغه جوړښت
  دے چې د تعريف ساحه ئې د حقيقي عددونو سټ~\LR{$\Real$} او ژبه ئې
  يوازې يو دوه ځاييز پريديکات لري چې پر حقيقي عددونو د~\LR{$<$}
  اړيکې په توګه تفسير شوے۔ څرګنده ده چې \LR{$\Struct{R}$}
  د شمېر نۀ وړ دے؛ خو دا چې \LR{$\Theory{R}$} څرګنده سازګاره ده، د
  لوېنهايم--سکولم د قضيې له مخې د شمېر وړ مدل لري، مثلاً
  \LR{$\Struct{S}$}، او د~\olref{prop:equiv} له مخې \LR{$\Struct{R}
  \equiv \Struct{S}$}۔ سربېره پر دې، دا چې \LR{$\Struct{R}$} او
  \LR{$\Struct{S}$} آيزومورف نۀ دي، ښيي چې د~\olref[iso]{thm:isom} معکوس
  په عام ډول نۀ صدق کوي۔
\end{rem}
\olpseganchor{OLP-0188-B013}{end}

\olpseganchor{OLP-0189-B004}{start}
\olfileid{mod}{bas}{pis}
\section{جزوي آيزومورفيزمونه}
\olpseganchor{OLP-0189-B004}{end}

\olpseganchor{OLP-0189-B005}{start}
\begin{defn}
  دوه جوړښتونه \LR{$\Struct{M}$} او \LR{$\Struct{N}$} راکړل شي۔
  له~\LR{$\Struct{M}$} څخه~\LR{$\Struct{N}$} ته \emph{جزوي آيزومورفيزم} يوه
  متناهي جزوي تابع~\LR{$p$} ده چې ارګومېنټونه ئې په~\LR{$\Domain M$} کښې او
  قيمتونه ئې په~\LR{$\Domain N$} کښې وي، او د~\olref[iso]{defn:isomorphism}
  د آيزومورفيزم شرطونه پر خپله د تعريف ساحه پوره کوي:
  \begin{enumerate}
  \item \LR{$p$} يو پر يو ده؛
  \item د هر ثابت سمبول~\LR{$c$} دپاره: که \LR{$p(\Assign{c}{M})$} تعريف
    شوے وي، نو \LR{$p(\Assign{c}{M}) = \Assign{c}{N}$}؛
  \item د هر \LR{$n$}-ځاييز پريديکات~\LR{$P$} دپاره: که \LR{$a_1$}،
    \dots،~\LR{$a_n$} د~\LR{$p$} د تعريف په ساحه کښې وي، نو
    \LR{$\langle a_1, \dots, a_n\rangle \in \Assign P M$} هله او يوازې
    هله چې \LR{$\langle p(a_1), \dots, p(a_n) \rangle \in \Assign P N$}؛
  \item د هر \LR{$n$}-ځاييز تابع~\LR{$f$} دپاره: که \LR{$a_1$}،
    \dots،~\LR{$a_n$} د~\LR{$p$} د تعريف په ساحه کښې وي، نو
    \LR{$p(\Assign f M (a_1, \dots,a_n)) =
    \Assign f N (p(a_1), \dots, p(a_n))$}۔
  \end{enumerate}
  د~\LR{$p$} متناهي والی په دې معنا دے چې \LR{$\dom{p}$} متناهي ده۔
\end{defn}
\olpseganchor{OLP-0189-B005}{end}

\olpseganchor{OLP-0189-B006}{start}
پام وکړئ چې تشه تابع~\LR{$\emptyset$} تل د هر دوو جوړښتونه تر منځ
يو جزوي آيزومورفيزم وي۔
\olpseganchor{OLP-0189-B006}{end}

\olpseganchor{OLP-0189-B007}{start}
\begin{defn}\ollabel{defn:partialisom}
  دوه جوړښتونه \LR{$\Struct{M}$} او \LR{$\Struct{N}$}
  \emph{په جزوي ډول آيزومورف} دي، چې \LR{$\Struct{M} \iso[p]
  \Struct{N}$} ئې ليکو، هله او يوازې هله چې د~\LR{$\Struct{M}$} او
  \LR{$\Struct{N}$} تر منځ د جزوي آيزومورفيزمونو يو ناتش سټ~\LR{$I$} وي چې د
  \emph{مخ او شا} خاصيت پوره کوي:
  \begin{enumerate}
  \item (\emph{مخ}) د هر \LR{$p \in I$} او \LR{$a \in \Domain M$} دپاره داسې
    \LR{$q \in I$} شته چې \LR{$p \subseteq q$} او~\LR{$a$} د~\LR{$q$} د تعريف په ساحه
    کښې وي؛
  \item (\emph{شا}) د هر \LR{$p \in I$} او \LR{$b \in \Domain N$} دپاره داسې
    \LR{$q \in I$} شته چې \LR{$p \subseteq q$} او~\LR{$b$} د~\LR{$q$} د قيمتونو په ساحه
    کښې وي۔
  \end{enumerate}
\end{defn}
\olpseganchor{OLP-0189-B007}{end}

\olpseganchor{OLP-0189-B008}{start}
\begin{thm}\ollabel{thm:p-isom1}
  که \LR{$\Struct{M} \iso[p] \Struct{N}$} وي او \LR{$\Struct{M}$} او
  \LR{$\Struct{N}$} د شمېر وړ وي، نو \LR{$\Struct{M} \iso \Struct{N}$}۔
\end{thm}
\olpseganchor{OLP-0189-B008}{end}

\olpseganchor{OLP-0189-B009}{start}
\begin{proof}
  دا چې \LR{$\Struct{M}$} او \LR{$\Struct{N}$} د شمېر وړ دي، فرض کړئ
  \LR{$\Domain{M} = \{a_0, a_1, \ldots \}$} او
  \LR{$\Domain{N} = \{b_0, b_1, \ldots \}$}۔ له يوه خپل سري
  \LR{$p_0 \in I$} څخه په پيل، د جزوي آيزومورفيزمونو زياتېدونکې لړۍ
  \LR{$p_0 \subseteq p_1 \subseteq p_2 \subseteq \cdots$} داسې تعريفوو:
  \begin{enumerate}
  \item که \LR{$n+1$} طاق وي، يعنې \LR{$n = 2r$}، نو د مخ خاصيت په کارولو
    داسې \LR{$p_{n+1} \in I$} ومومئ چې \LR{$p_n \subseteq p_{n+1}$} او~\LR{$a_r$}
    د~\LR{$p_{n+1}$} د تعريف په ساحه کښې وي؛
  \item که \LR{$n+1$} جفت وي، يعنې \LR{$n+1 =2r$}، نو د شا خاصيت په کارولو
    داسې \LR{$p_{n+1} \in I$} ومومئ چې \LR{$p_n \subseteq p_{n+1}$} او~\LR{$b_r$}
    د~\LR{$p_{n+1}$} د قيمتونو په ساحه کښې وي۔
  \end{enumerate}
اوس که
\[
p = \bigcup_{n\ge 0} p_n,
\]
کېږدو، نو \LR{$p$} د~\LR{$\Struct{M}$} او~\LR{$\Struct{N}$} تر منځ يو آيزومورفيزم
دے۔
\end{proof}
\olpseganchor{OLP-0189-B009}{end}

\olpseganchor{OLP-0189-B010}{start}
\begin{prob}
  په تفصيل وښيئ چې په~\olref[mod][bas][pis]{thm:p-isom1} کښې تعريف
  شوے~\LR{$p$} په رښتيا يو آيزومورفيزم دے۔
\end{prob}
\olpseganchor{OLP-0189-B010}{end}

\olpseganchor{OLP-0189-B011}{start}
\begin{thm}\ollabel{thm:p-isom2}
  فرض کړئ \LR{$\Struct{M}$} او \LR{$\Struct{N}$} د يوې سوچه اړيکيزې
  ژبه جوړښتونه دي (داسې ژبه چې يوازې
  پريديکاتونه لري او تابعې يا ثابتونه نۀ لري)۔ بيا که
  \LR{$\Struct{M} \iso[p] \Struct{N}$} وي، نو \LR{$\Struct{M} \elemequiv
  \Struct{N}$} هم وي۔
\end{thm}
\olpseganchor{OLP-0189-B011}{end}

\olpseganchor{OLP-0189-B012}{start}
\begin{proof}
  پر فارمولونه په استقرا ښودل کېږي چې که \LR{$a_1$}، \dots،~\LR{$a_n$} او
  \LR{$b_1$}، \dots،~\LR{$b_n$} داسې وي چې يو جزوي آيزومورفيزم~\LR{$p$} هر~\LR{$a_i$}
  پر~\LR{$b_i$} وړي، او \LR{$s_1(x_i) =a_i$} او \LR{$s_2(x_i) =b_i$} وي (د
  \LR{$i =1$}، \dots،~\LR{$n$} دپاره)، نو \LR{$\Sat{M}{!A}[s_1]$} هله او يوازې هله
  چې \LR{$\Sat{N}{!A}[s_2]$} وي۔ د~\LR{$n=0$} حالت
  \LR{$\Struct{M} \elemequiv \Struct{N}$} راکوي۔
\end{proof}
\olpseganchor{OLP-0189-B012}{end}

\olpseganchor{OLP-0189-B013}{start}
\begin{rem}
که تابعې موجودې وي، مخکينۍ پايله بيا هم رښتونې ده، خو د هغه
آيزومورفيزم څېړل پکار دي چې~\LR{$p$} ئې د~\LR{$a_1$}، \dots،~\LR{$a_n$} له خوا د
\LR{$\Struct{M}$} پر توليد شوي فرعي جوړښت او د~\LR{$b_1$}،
\dots،~\LR{$b_n$} له خوا د~\LR{$\Struct{N}$} پر توليد شوي فرعي جوړښت
رامنځته کوي۔
\end{rem}
\olpseganchor{OLP-0189-B013}{end}

\olpseganchor{OLP-0189-B014}{start}
مخکينۍ پايله په پړاوونو هم ``ماتېدے'' شي: په يوه فارمول کښې د
يو بل دننه کميت ټاکونکو د شمېر او د اړوندو جزوي آيزومورفيزمونو د
غځېدو د ځلونو تر منځ اړيکه ټينګوو۔
\olpseganchor{OLP-0189-B014}{end}

\olpseganchor{OLP-0189-B015}{start}
\begin{defn}
  د هر فارمول~\LR{$!A$} دپاره، د~\LR{$!A$} \emph{د کميت ټاکونکو رتبه}، چې
  \LR{$\QuantRank{!A} \in \Nat$} ئې ليکو، په بازګشتي ډول په~\LR{$!A$} کښې د
  يو بل دننه کميت ټاکونکو تر ټولو لوړه شمېره تعريفېږي۔ دوه
  جوړښتونه \LR{$\Struct{M}$} او \LR{$\Struct{N}$} \emph{\LR{$n$}-هم‌ارز} دي،
  چې \LR{$\Struct{M} \elemequiv[n] \Struct{N}$} ئې ليکو، که د کميت
  ټاکونکو د مرتبې تر~\LR{$n$} پورې پر ټولو جملو سره سلا وي۔
\end{defn}
\olpseganchor{OLP-0189-B015}{end}

\olpseganchor{OLP-0189-B016}{start}
\begin{prop}\ollabel{prop:qr-finite}
  فرض کړئ \LR{$\Lang{L}$} يوه متناهي سوچه اړيکيزه ژبه ده، يعنې
  داسې ژبه چې متناهي شمېر پريديکاتونه او ثابت سمبولونه
  لري او تابعې نۀ لري۔ بيا د هر \LR{$n \in \Nat$} دپاره، له
  منطقي هم‌ارزۍ پرته، د~\LR{$\Lang{L}$} په ژبه کښې يوازې متناهي ډېرې داسې
  د لومړۍ درجې تړلي فارمولونه شته چې د کميت ټاکونکو رتبه ئې تر~\LR{$n$}
  زياته نۀ وي۔
\end{prop}
\olpseganchor{OLP-0189-B016}{end}

\olpseganchor{OLP-0189-B017}{start}
\begin{proof}
  پر~\LR{$n$} په استقرا۔
\end{proof}
\olpseganchor{OLP-0189-B017}{end}

\olpseganchor{OLP-0189-B018}{start}
\begin{defn}
  يو جوړښت~\LR{$\Struct{M}$} راکړل شي۔ \LR{$\Domain M^{<\omega}$} د
  \LR{$\Domain{M}$} پر غړو د ټولو متناهي لړيو سټ ونيسئ۔ په~\LR{$\mathbf{a}, \mathbf{b}, \mathbf{c}, \ldots$} د غړو متناهي لړۍ ښيو۔ که
  \LR{$\mathbf{a} \in \Domain{M}^{<\omega}$} او \LR{$a \in \Domain{M}$} وي، نو
  \LR{$\mathbf{a}a$} له~\LR{$a$} سره د~\LR{$\mathbf{a}$} \emph{نښلون} ښيي۔
\end{defn}
\olpseganchor{OLP-0189-B018}{end}

\olpseganchor{OLP-0189-B019}{start}
\begin{defn}
  جوړښتونه \LR{$\Struct{M}$} او \LR{$\Struct{N}$} راکړل شي۔ د يو شان
  اوږدوالي لړيو تر منځ اړيکې
  \LR{$I_n \subseteq \Domain M^{<\omega} \times \Domain N^{<\omega}$}
  پر~\LR{$n$} په بازګښت داسې تعريفوو: شاخص د پاتې غځونو شمېر دے؛ هره جوړه کېدے شي له مخکې هر متناهي اوږدوالے ولري۔
   \begin{enumerate}
   \item \LR{$I_0(\mathbf{a},\mathbf{b})$} هله او يوازې هله چې
     \LR{$\mathbf{a}$} او \LR{$\mathbf{b}$} په~\LR{$\Struct{M}$} او~\LR{$\Struct{N}$}
     کښې يو شان اټومي فارمولونه پوره کړي؛ يعنې که د دواړو لړيو
     اوږدوالی~\LR{$k$} وي، \LR{$s_1(x_i) = a_i$} او \LR{$s_2(x_i) = b_i$} وي، او
     \LR{$!A$} داسې اټومي فارمول وي چې ټول متغيرونه ئې په
     \LR{$x_1$}، \dots،~\LR{$x_k$} کښې وي، نو \LR{$\Sat{M}{!A}[s_1]$} هله او يوازې
     هله چې \LR{$\Sat{N}{!A}[s_2]$} وي۔
   \item \LR{$I_{n+1} (\mathbf{a},\mathbf{b})$} هله او يوازې هله چې د هر
     \LR{$a\in \Domain M$} دپاره داسې \LR{$b\in \Domain N$} وي چې
     \LR{$I_n(\mathbf{a}a,\mathbf{b}b)$}، او برعکس۔
   \end{enumerate}
\end{defn}
\olpseganchor{OLP-0189-B019}{end}


\olpseganchor{OLP-0189-B020}{start}
\begin{defn}
  \LR{$\Struct{M} \approx_n \Struct{N}$} ليکو که
  \LR{$I_n(\emptyseq,\emptyseq)$} د~\LR{$\Struct{M}$} او~\LR{$\Struct{N}$} دپاره
  صدق وکړي (دلته \LR{$\emptyseq$} تشه لړۍ ده)۔
\end{defn}
\olpseganchor{OLP-0189-B020}{end}

\olpseganchor{OLP-0189-B021}{start}
\begin{thm}\ollabel{thm:b-n-f}
  فرض کړئ \LR{$\Lang{L}$} يوه سوچه اړيکيزه ژبه ده۔ بيا
  \LR{$I_n(\mathbf{a},\mathbf{b})$} دا لازموي چې د هر داسې~\LR{$!A$} دپاره چې
  \LR{$\QuantRank{!A} \le n$} وي، \LR{$\Sat{M}{!A}[\mathbf{a}]$} هله او يوازې
  هله چې \LR{$\Sat{N}{!A}[\mathbf{b}]$} وي (دلته هم، \LR{$\mathbf{a}$} هغه
  وخت~\LR{$!A$} پوره کوي چې هره داسې~\LR{$s$} چې \LR{$s(x_i) = a_i$} وي، \LR{$!A$} پوره
  کړي)۔ سربېره پر دې، که \LR{$\Lang{L}$} متناهي وي، معکوس هم صدق کوي۔
\end{thm}
\olpseganchor{OLP-0189-B021}{end}

\olpseganchor{OLP-0189-B022}{start}
\begin{proof}
  دا چې \LR{$I_n(\mathbf{a},\mathbf{b})$} لازموي چې \LR{$\mathbf{a}$} او
  \LR{$\mathbf{b}$} د کميت ټاکونکو تر~\LR{$n$} پورې يو شان فارمولونه پوره
  کوي، پر~\LR{$!A$} په اسانه استقرا ثابتېږي۔ د معکوس دپاره پر~\LR{$n$} استقرا
  کوو او \olref{prop:qr-finite} کاروو، چې تضمينوي د هر~\LR{$n$} دپاره د
  هغې مرتبې تر ډېره متناهي ډېر يو له بله نۀ هم‌ارز فارمولونه شته۔
\olpseganchor{OLP-0189-B022}{end}

\olpseganchor{OLP-0189-B023}{start}
  د~\LR{$n=0$} دپاره، دا فرض چې \LR{$\mathbf{a}$} او \LR{$\mathbf{b}$} يو شان بې
  کميت ټاکونکو فارمولونه پوره کوي، راکوي چې يو شان اټومي فارمولونه
  پوره کوي؛ نو \LR{$I_0(\mathbf{a},\mathbf{b})$}۔
\olpseganchor{OLP-0189-B023}{end}

\olpseganchor{OLP-0189-B024}{start}
  د~\LR{$n+1$} په حالت کښې فرض کړئ \LR{$\mathbf{a}$} او \LR{$\mathbf{b}$} د کميت
  ټاکونکو تر~\LR{$n+1$} پورې يو شان فارمولونه پوره کوي۔ د
  \LR{$I_{n+1}(\mathbf{a},\mathbf{b})$} د ښودلو دپاره کافي ده چې د هر
  \LR{$a \in \Domain M$} دپاره داسې \LR{$b \in \Domain N$} وښيو چې
  \LR{$I_n(\mathbf{a}a,\mathbf{b}b)$}؛ او د استقرا له فرضه بيا کافي ده
  وښيو چې د هر \LR{$a \in \Domain M$} دپاره داسې \LR{$b \in \Domain N$} شته
  چې \LR{$\mathbf{a}a$} او \LR{$\mathbf{b}b$} د کميت ټاکونکو تر~\LR{$n$} پورې يو
  شان فارمولونه پوره کوي۔
\olpseganchor{OLP-0189-B024}{end}

\olpseganchor{OLP-0189-B025}{start}
  يو \LR{$a \in \Domain M$} راکړل شي۔ \LR{$!T^a_n$} د رتبې تر~\LR{$n$} پورې د هغو
  فارمولونه~\LR{$!B(x,\mathbf{y})$} سټ ونيسئ چې \LR{$\mathbf{a}a$} ئې
  په~\LR{$\Struct{M}$} کښې پوره کوي؛ \LR{$!T^a_n$} له منطقي هم‌ارزۍ پرته
  متناهي دے، نو د غړو د تړون په اخيستو ئې يو د لومړۍ درجې
  فارمول ګڼلے شو۔ له دې پايله کېږي چې \LR{$\mathbf{a}$}
  \LR{$\lexists[x][!T^a_n(x,\mathbf{y})]$} پوره کوي، چې د کميت ټاکونکو
  رتبه ئې تر~\LR{$n+1$} زياته نۀ ده۔ د فرض له مخې \LR{$\mathbf{b}$} هماغه
  فارمول په~\LR{$\Struct{N}$} کښې پوره کوي؛ نو داسې
  \LR{$b \in \Domain N$} شته چې \LR{$\mathbf{b}b$}، \LR{$!T^a_n$} پوره کوي۔ په
  ځانګړي ډول، \LR{$\mathbf{b}b$} د کميت ټاکونکو تر~\LR{$n$} پورې هماغه
  فارمولونه پوره کوي چې \LR{$\mathbf{a}a$} ئې پوره کوي۔ په ورته ډول
  ښودل کېږي چې د هر \LR{$b \in \Domain N$} دپاره داسې \LR{$a\in \Domain M$}
  شته چې \LR{$\mathbf{a}a$} او \LR{$\mathbf{b}b$} د کميت ټاکونکو تر~\LR{$n$} پورې
  يو شان فارمولونه پوره کوي؛ په دې سره ثبوت بشپړېږي۔
\end{proof}
\olpseganchor{OLP-0189-B025}{end}

\olpseganchor{OLP-0189-B026}{start}
\begin{cor}\ollabel{cor:b-n-f}
  که \LR{$\Struct{M}$} او \LR{$\Struct{N}$} په يوه متناهي ژبه کښې
  سوچه اړيکيز جوړښتونه وي، نو \LR{$\Struct{M} \approx_n\Struct{N}$}
  هله او يوازې هله چې \LR{$\Struct{M} \elemequiv[n] \Struct{N}$} وي۔ په
  ځانګړي ډول، \LR{$\Struct{M} \elemequiv \Struct{N}$} هله او يوازې هله
  چې د هر~\LR{$n$} دپاره \LR{$\Struct{M} \approx_n \Struct{N}$} وي۔
\end{cor}
\olpseganchor{OLP-0189-B026}{end}

\olpseganchor{OLP-0190-B004}{start}
\olfileid{mod}{bas}{dlo}
\section{بې له سرحدي ټکو ګڼ خطي ترتيبونه}
\olpseganchor{OLP-0190-B004}{end}

\olpseganchor{OLP-0190-B005}{start}
\begin{defn}
  \emph{بې له سرحدي ټکو ګڼ خطي ترتيب} د هغې ژبه يو جوړښت
  \LR{$\Struct{M}$} دے چې يوازينے دوه ځاييز پريديکات~\LR{$<$} لري او لاندې
  جملې پوره کوي:
  \begin{enumerate}
  \item \LR{$\lforall[x][\lnot x < x]$}؛
  \item \LR{$\lforall[x][\lforall[y][\lforall[z][(x < y \lif (y < z \lif x
    <z ))]]]$}؛
  \item \LR{$\lforall[x][\lforall[y][(x< y \lor \eq[x][y] \lor y < x)]]$}؛
  \item \LR{$\lforall[x][\lexists[y][x < y]]$}؛
  \item \LR{$\lforall[x][\lexists[y][y < x]]$}؛
  \item \LR{$\lforall[x][\lforall[y][(x < y \lif \lexists[z][(x < z \land
        z < y))]]]$}۔
 \end{enumerate}
\end{defn}
\olpseganchor{OLP-0190-B005}{end}

\olpseganchor{OLP-0190-B006}{start}
\begin{thm}\ollabel{thm:cantorQ}
  هر دوه د شمېر وړ، بې له سرحدي ټکو ګڼ خطي ترتيبونه آيزومورف دي۔
\end{thm}
\olpseganchor{OLP-0190-B006}{end}

\olpseganchor{OLP-0190-B007}{start}
\begin{proof}
  فرض کړئ \LR{$\Struct{M_1}$} او \LR{$\Struct{M_2}$} د شمېر وړ، بې له
  سرحدي ټکو ګڼ خطي ترتيبونه دي، \LR{${<_1} = \Assign{<}{M_1}$} او
  \LR{${<_2} = \Assign{<}{M_2}$}، او \LR{$\PIso{I}$} د دوى تر منځ د ټولو جزوي
  آيزومورفيزمونو سټ دے۔ \LR{$\PIso{I}$} تش نۀ دے، ځکه لږ تر لږه
  \LR{$\emptyset \in \PIso{I}$}۔ ښيو چې \LR{$\PIso{I}$} د مخ او شا خاصيت پوره
  کوي۔ بيا \LR{$\Struct{M_1} \iso[p] \Struct{M_2}$}، او قضيه
  له~\olref[pis]{thm:p-isom1} څخه پايله کېږي۔
\olpseganchor{OLP-0190-B007}{end}

\olpseganchor{OLP-0190-B008}{start}
  د دې د ښودلو دپاره چې \LR{$\PIso{I}$} د مخ خاصيت پوره کوي، فرض کړئ
  \LR{$p \in \PIso{I}$} او د~\LR{$i = 1$}، \dots،~\LR{$n$} دپاره
  \LR{$p(a_i) = b_i$}، او د عموميت له زيان پرته فرض کړئ
  \LR{$a_1 <_1 a_2 <_1 \cdots <_1 a_n$}۔ يو \LR{$a \in \Domain{M_1}$} چې
  راکړل شي، \LR{$b \in \Domain{M_2}$} داسې ومومئ:
  \begin{enumerate}
  \item که \LR{$a <_1 a_1$} وي، \LR{$b \in \Domain{M_2}$} داسې ونيسئ چې
    \LR{$b <_2 b_1$}؛
  \item که \LR{$a_n <_1 a$} وي، \LR{$b \in \Domain{M_2}$} داسې ونيسئ چې
    \LR{$b_n <_2 b$}؛
 \item که د کوم~\LR{$i$} دپاره \LR{$a_i <_1 a <_1 a_{i+1}$} وي، نو
   \LR{$b \in \Domain{M_2}$} داسې ونيسئ چې \LR{$b_i <_2 b <_2 b_{i+1}$}۔
  \end{enumerate}
  د مطلوب خاصيت لرونکے~\LR{$b$} تل موندل کېدے شي، ځکه \LR{$\Struct{M_2}$} بې
  له سرحدي ټکو ګڼ خطي ترتيب دے۔ \LR{$q = p \cup \{ \langle a, b \rangle
  \}$} داسې تعريف کړئ چې \LR{$q \in \PIso{I}$} د~\LR{$p$} مطلوبه غځونه وي۔ په
  دې سره د مخ خاصيت ثابت شو۔ د شا خاصيت هم ورته دے۔ نو
  \LR{$\Struct{M_1} \iso[p] \Struct{M_2}$}؛ او د
  \olref[pis]{thm:p-isom1} له مخې \LR{$\Struct{M_1} \iso \Struct{M_2}$}۔
\end{proof}
\olpseganchor{OLP-0190-B008}{end}

\olpseganchor{OLP-0190-B009}{start}
\begin{prob}
  د~\olref[mod][bas][dlo]{thm:cantorQ} ثبوت په دې ښودلو بشپړ کړئ چې
  \LR{$\PIso{I}$} د شا خاصيت پوره کوي۔
\end{prob}
\olpseganchor{OLP-0190-B009}{end}

\olpseganchor{OLP-0190-B010}{start}
\begin{rem}
  \LR{$\Struct{S}$}، بې له سرحدي ټکو هر د شمېر وړ ګڼ خطي ترتيب،
  ونيسئ۔ بيا (د~\olref{thm:cantorQ} له مخې) \LR{$\Struct{S} \iso
  \Struct{Q}$}، چې \LR{$\Struct{Q} = (\Rat, <)$} هغه د شمېر وړ ګڼ
  خطي ترتيب دے چې د تعريف ساحه ئې د ناطقو عددونو سټ~\LR{$\Rat$} دے۔ اوس د
  \olref[thm]{remark:R} جوړښت~\LR{$\Struct{R} = (\Real, <)$} بيا وګورئ۔ مو
  وليدل چې داسې د شمېر وړ جوړښت~\LR{$\Struct{S}$} شته چې
  \LR{$\Struct{R} \elemequiv \Struct{S}$}۔ خو \LR{$\Struct{S}$}، بې له سرحدي
  ټکو يو د شمېر وړ ګڼ خطي ترتيب دے؛ نو له جوړښت~\LR{$\Struct{Q}$}
  سره آيزومورف (او ځکه ابتدايي هم‌ارز) جوړښت دے۔ د ابتدايي هم‌ارزۍ له
  انتقاليت څخه \LR{$\Struct{R} \elemequiv \Struct{Q}$} پايله کېږي۔ (دا مو
  په مستقيم ډول هم ښودلے شول: د هماغه مخ او شا دليل په وسيله مو
  \LR{$\Struct{R} \iso[p] \Struct{Q}$} ټينګولے شو۔)
\end{rem}
\olpseganchor{OLP-0190-B010}{end}

\paragraph{د منجمدې سرچينې بشپړه اړونده برخه}
\olpseganchor{OLP-0653-B004}{start}
\olfileid{mod}{bas}{nsa-olp0653}
\section{د حساب نامعياري مدلونه}
\olpseganchor{OLP-0653-B004}{end}

\olpseganchor{OLP-0653-B005}{start}
\begin{defn}
د حساب ژبه دې \LR{$\Lang{L_N}$} وي، چې يو ثابت سمبول \LR{$\Obj{0}$}، دوه‌ځايي پريديکات \LR{$<$}، يوځايي تابع \LR{$\prime$}، او دوه‌ځايي تابعې \LR{$+$} او \LR{$\times$} لري۔
\begin{enumerate}
\item د حساب \emph{معياري مدل} د \LR{$\Lang{L_N}$} هغه جوړښت \LR{$\Struct{N}$} دے چې ساحه ئې \LR{$\Nat = \{ 0, 1, 2, \dots\}$} ده؛ \LR{$\Obj{0}$} د \LR{$0$} په توګه، \LR{$<$} پر~\LR{$\Nat$} د کوچنيوالي د اړيکې په توګه، او \LR{$\prime$}، \LR{$+$} او \LR{$\times$} په ترتيب سره پر \LR{$\Nat$} د جانشين، جمع او ضرب په توګه تعبيروي۔
\item \emph{رښتينے حساب} تيوري \LR{$\Theory{N}$} ده۔
\end{enumerate}
\end{defn}
\olpseganchor{OLP-0653-B005}{end}

\olpseganchor{OLP-0653-B006}{start}
په \LR{$\Lang{L_N}$} کښې د کار پر مهال د \LR{$\Obj{0}^{\prime\dots\prime}$} په بڼه هر ترم، چې د جانشين تابع ئې پر \LR{$\Obj{0}$} باندې \LR{$n$} ځله لګېدلې وي، په~\LR{$\num{n}$} لنډوو۔
\olpseganchor{OLP-0653-B006}{end}

\olpseganchor{OLP-0653-B007}{start}
\begin{defn}
د \LR{$\Lang{L_N}$} يو جوړښت~\LR{$\Struct{M}$} هله او يوازې هله \emph{معياري} دے چې \LR{$\Struct{N} \iso \Struct{M}$}۔
\end{defn}
\olpseganchor{OLP-0653-B007}{end}

\olpseganchor{OLP-0653-B008}{start}
\begin{thm}\ollabel{thm:non-std}
د رښتيني حساب نامعياري د شمېر وړ مدلونه شته۔
\end{thm}
\olpseganchor{OLP-0653-B008}{end}

\olpseganchor{OLP-0653-B009}{start}
\begin{proof}
د نوي ثابت سمبول~\LR{$c$} په زياتولو \LR{$\Lang{L_N}$} غځوو او دا تيوري ګورو:
\[
\Theory{N} \cup \Setabs{\num{n} < c}{n \in \Nat}.
\]
تيوري په متناهي ډول د صدق وړ ده؛ نو د کمپکتوالي له مخې يو مدل \LR{$\Struct{M}$} لري، چې د ښکته لوېنهايم--سکولم قضيې له مخې ئې د شمېر وړ اخيستلے شو۔ که \LR{$\Domain{M}$} د \LR{$\Struct{M}$} ساحه وي، نو د \LR{$\Lang{L}$} د غيرمنطقي نښو تعبيرونه دې په \LR{$\Struct{M}$} کښې داسې وي: \LR{$\mathbf{z} = \Assign{\Obj{0}}{M} \in \Domain{M}$}، \LR{${\prec} =
\Assign{<}{M} \subseteq M^2$}، \LR{$* = \Assign{\prime}{M} \colon
\Domain{M} \to \Domain{M}$}، او \LR{$\oplus = \Assign{+}{M}, \otimes =
\Assign{\times}{M} \colon \Domain{M}^2 \to \Domain{M}$}۔ د هر \LR{$x
\in \Domain{M}$} دپاره \LR{$x^*$} د \LR{$\Domain{M}$} هغه عنصر ښيي چې له \LR{$x$} څخه د~\LR{$*$} په لګولو ترلاسه کېږي۔
\olpseganchor{OLP-0653-B009}{end}

\olpseganchor{OLP-0653-B010}{start}
اوس، که \LR{$h$} د \LR{$\Struct{N}$} او د \LR{$\Struct{M}$} د اصلي حسابي راکمونې تر منځ ايزومورفيزم وے، نو داسې \LR{$n \in \Nat$} به وے چې \LR{$h(n) = \Assign{c}{M}$}۔ نو په \LR{$\Struct{N}$} کښې داسې هره ارزښت ټاکنه \LR{$s$} واخلو چې \LR{$s(x) = n$}۔ بيا \LR{$\Sat{N}{\eq[\num{n}][x]}[s]$}؛ د \olref[iso]{thm:isom} د ثبوت له مخې \LR{$\Sat{M}{\eq[\num{n}][x]}[h \circ s]$} هم، نو \LR{$\Assign{c}{M} =
\mathbf{z}^{*\cdots *}$} (چې \LR{$*$} پکښې \LR{$n$} ځله تکرار شوے دے)۔ خو دا ناشوني دي، ځکه د فرض له مخې \LR{$\Sat{M}{\num{n} < c}$} او \LR{$\prec$} غيرانعکاسي ده۔ نو \LR{$\Struct{M}$} نامعياري دے۔
\end{proof}
\olpseganchor{OLP-0653-B010}{end}

\olpseganchor{OLP-0653-B011}{start}
\begin{prob}
پر سټ \LR{$X$} يوه اړيکه \LR{$R$} هله او يوازې هله \emph{ښه‌بنسټه} ده چې په~\LR{$R$} کښې نامتناهي ښکته تلونکې لړۍ نۀ وي؛ يعنې په \LR{$X$} کښې داسې \LR{$x_0$}، \LR{$x_1$}، \LR{$x_2$}،~\dots نۀ وي چې \LR{$\dots x_2Rx_1Rx_0$}۔ د دې فرض لاندې چې د زېرملو--فرېنکل سټ تيوري~\LR{$ZF$} سازګاره ده، وښيئ چې د \LR{$ZF$} نۀ ښه‌بنسټه مدلونه شته؛ يعنې داسې مدلونه \LR{$\mathfrak{M}$} چې \LR{$\dots x_2 \in x_1 \in x_0$}۔
\end{prob}
\olpseganchor{OLP-0653-B011}{end}

\olpseganchor{OLP-0653-B012}{start}
په \olref{thm:non-std} کښې نامعياري مدل \LR{$\Struct{M}$} له معياري مدل سره ابتدايي هم‌ارز دے؛ ځکه ئې څو خاصيتونه راايستلے شو۔ د دې برخې پاتې متن همدې کار ته ځانګړے دے، چې د \LR{$\Struct{M}$} د خاصيتونو له لارې د \LR{$\Theory{N}$} د د شمېر وړ نامعياري مدلونو دقيقه ځانګړونه ترلاسه کړو۔
\olpseganchor{OLP-0653-B012}{end}

\olpseganchor{OLP-0653-B013}{start}
\begin{enumerate}
\item د \LR{$\Domain{M}$} هېڅ غړے له خپل ځان څخه د \LR{$\prec$} له مخې کوچنے نۀ دے: جمله \LR{$\lforall[x][\lnot x < x]$} په \LR{$\Struct{N}$} کښې رښتينې ده، نو په \LR{$\Struct{M}$} کښې هم۔
\item په ورته استدلال ترلاسه کوو چې \LR{$\prec$} د \LR{$\Domain{M}$} يو \emph{خطي ترتيب} دے؛ يعنې پر \LR{$\Domain{M}$} يوه کلي، غيرانعکاسي او متعدي اړيکه ده۔
\item عنصر \LR{$\mathbf{z}$} د \LR{$\Domain{M}$} د \LR{$\prec$} له مخې تر ټولو کوچنے عنصر دے۔
\item د \LR{$\Domain{M}$} هر غړے د \LR{$\prec$} له مخې له خپل \LR{$*$}-جانشين څخه کوچنے دے، او \LR{$x^*$} د \LR{$\Domain{M}$} هغه تر ټولو کوچنے غړے دے چې د \LR{$\prec$} له مخې له \LR{$x$} څخه لوے دے۔
\item \LR{$\Struct{M}$} د \LR{$\prec$} يوه ابتدائي ټوټه لري چې له \LR{$\Nat$} سره ايزومورف ده: \LR{$\mathbf{z}, \mathbf{z}^*, \mathbf{z}^{**}, \dots$}۔ دې ته د \LR{$\Domain{M}$} \emph{معياري برخه} وايو۔ د \LR{$\Domain{M}$} هر بل غړے \emph{نامعياري} دے۔ د \LR{$\Domain{M}$} نامعياري غړي خامخا شته؛ کنه په \olref{thm:non-std} کښې د ثبوت تابع \LR{$h$} به ايزومورفيزم وے۔ \LR{$n, m, \dots$} د هغو متغيرونه په توګه کاروو چې د \LR{$\Struct{M}$} پر همدې معياري برخې ګرځي۔
\item هر نامعياري عنصر له هر معياري عنصر څخه لوے دے؛ ځکه د هر \LR{$n \in \Nat$} دپاره
  \[
  \Sat{N}{\lforall[z][(\lnot(\eq[z][\Obj{0}] \lor
  \dots \lor \eq[z][\num{n}]) \lif \num{n} < z)]},
  \]
  نو که \LR{$z \in \Domain{M}$} له ټولو معياري عناصرو بېل وي، بايد له ټولو څخه \emph{لوے} وي۔
\item له \LR{$\mathbf{z}$} پرته د \LR{$\Domain{M}$} هر غړے د \LR{$\Domain{M}$} د يوه يکتا عنصر \LR{$*$}-جانشين دے؛ دغه مخکښېنی په \LR{$^*x$} ښيو۔ که \LR{$x = y^*$} وي، نو که د \LR{$x$} او \LR{$y$} يو معياري وي، دواړه معياري دي (او که يو نامعياري وي، دواړه نامعياري دي)۔
\item پر \LR{$\Domain{M}$} د هم‌ارزۍ اړيکه \LR{$\approx$} داسې تعريفوو: \LR{$x \approx y$} هله او يوازې هله چې د کوم \emph{معياري} \LR{$n$} دپاره \LR{$x \oplus n = y$} يا \LR{$y \oplus n =x$}۔ په نورو ټکو، \LR{$x \approx y$} هله او يوازې هله چې د \LR{$x$} او \LR{$y$} تر منځ واټن متناهي وي۔ که \LR{$n$} او \LR{$m$} معياري وي، نو \LR{$n \approx m$}۔ د \LR{$x$} \emph{بلاک} د هم‌ارزۍ ټولګے \LR{$[x] = \Setabs{y}{x
  \approx y}$} تعريفوو۔
\item فرض کړئ چې \LR{$x \prec y$} او \LR{$x \not\approx y$}۔ څرنګه چې \LR{$\Sat{N}{\lforall[x][\lforall[y][(x < y \lif (x' < y \lor x' =
      y))]]}$}، نو \LR{$x^* \prec y$} يا \LR{$x^* = y$}۔ دويم حالت ناشونے دے، ځکه ترې \LR{$x \approx y$} راځي؛ نو \LR{$x^* \prec y$}۔ همداسې، که \LR{$x \prec y$} او \LR{$x \not\approx y$}، نو \LR{$x \prec {^*y}$}۔ د دغه مخکيني غړي سره مساوات به دواړه په يوه بلاک کښې واچوي۔ په دواړو لوريو دغه ګام هر متناهي شمېر تکرارولے شو۔ نو که \LR{$x \prec y$} او \LR{$x \not\approx y$}، نو هر \LR{$w \approx x$} د \LR{$\prec$} له مخې له هر \LR{$v \approx y$} څخه کوچنے دے۔ هر \emph{نامعياري} بلاک \LR{$[x]$} داسې دوه‌خوا نامتناهي لړۍ جوړوي:
  \[
  \cdots \prec  {^{**}x} \prec {^*}x \prec x \prec x^* \prec x^{**}
  \prec \cdots
  \]
  دې ته د \LR{$Z$} لړۍ وايو، ځکه د صحيح عددونو ترتيبي ډول لري۔ معياري بلاک له صفر څخه پيلېږي او د طبيعي عددونو په شان ترتيب لري۔
\item د \LR{$\prec$} ترتيب بلاکونو ته پورته کېدے شي: که \LR{$x \prec y$} وي او دواړه په بېلو بلاکونو کښې وي، نو د \LR{$x$} بلاک د \LR{$y$} له بلاک څخه کوچنے دے۔ بلاک هله \emph{نامعياري} دے چې نامعياري عنصر ولري۔ معياري بلاک تر ټولو کوچنے بلاک دے۔
\item تر ټولو کوچنے نامعياري بلاک نشته: که \LR{$y$} نامعياري وي، نو داسې \LR{$x \prec y$} شته چې \LR{$x$} هم نامعياري وي او \LR{$x \not\approx y$}۔ ثبوت: په معياري مدل \LR{$\Struct{N}$} کښې هر عدد پر دوو وېشل کېږي، او ښايي يو پاتې شي؛ يعنې
  \LR{$\Sat{N}{\lforall[y][\lexists[x][(y = x + x \lor y = x + x +
        \Obj{0}')]]}$}۔ د ابتدايي هم‌ارزۍ له مخې د هر \LR{$y \in \Domain{M}$} دپاره داسې \LR{$x \in \Domain{M}$} شته چې \LR{$x \oplus x = y$} يا \LR{$x \oplus x \oplus \mathbf{z}^*= y$}۔ که \LR{$x$} معياري وے، \LR{$y$} به هم معياري وے؛ نو \LR{$x$} نامعياري دے۔ سربېره پر دې، \LR{$x$} او \LR{$y$} په بېلو بلاکونو کښې دي؛ يعنې \LR{$x \not\approx y$}۔ د دې د ليدلو دپاره فرض کړئ چې دواړه په يوه بلاک کښې وي؛ يعنې د کوم معياري \LR{$n$} دپاره \LR{$x \oplus n = y$}۔ که \LR{$y = x \oplus x$} وي، نو \LR{$x \oplus n = x
  \oplus x$}، چې د جمع د حذف قانون له مخې ترې \LR{$x = n$} راځي (دغه قانون په \LR{$\Struct{N}$} کښې صدق کوي، نو په \LR{$\Struct{M}$} کښې هم)؛ بيا خو \LR{$x$} معياري شو۔ که \LR{$y = x \oplus x \oplus \mathbf{z}^*$} وي، استدلال همداسې دے۔
\item تر ټولو لوے بلاک هم نشته۔ هر نامعياري غړے مثبت دے؛ دوه‌چنده ئې ترې لوے دے۔ که دواړه په يوه بلاک کښې وای، د جمع د حذف له مخې خپله غړے به له يوې معياري عددنښې سره برابر و۔
\item د بلاکونو ترتيب ګڼ دے: که \LR{$[x]$} له \LR{$[y]$} څخه کوچنے وي (چې \LR{$x \not\approx y$})، نو د دواړو تر منځ يو له دواړو بېل بلاک \LR{$[z]$} شته۔ فرض کړئ چې \LR{$x \prec y$}۔ لکه مخکې، \LR{$x \oplus y$} پر دوو وېشل کېږي (ښايي پاتې شونے ولري)؛ نو داسې \LR{$u \in \Domain{M}$} شته چې \LR{$x \oplus y = u \oplus u$} يا \LR{$x
  \oplus y = u \oplus u \oplus \mathbf{z}^*$}۔ عنصر \LR{$u$} د \LR{$x$} او \LR{$y$} اوسط دے، نو د دواړو تر منځ دے۔ د بېلو بلاکونو له امله واټن له يوه زيات دے۔ که اوسط له کيڼ غړي څخه لوے نۀ وي، دوه‌برابره او پاتې شونے به ښي غړے د کيڼ غړي تر تالي زيات نۀ پرېږدي۔ که اوسط له ښي غړي څخه کوچنے نۀ وي، د دواړو مجموعه به د اوسط له دوه‌برابره کمه وي۔ فرض کړئ چې \LR{$x \oplus y = u \oplus u$} (بل حالت ورته دے): که \LR{$u \approx x$} وي، نو د کوم معياري \LR{$n$} دپاره
  \[
  x \oplus y = x \oplus n \oplus x \oplus n,
  \]
  نو \LR{$y = x \oplus n \oplus n$}، او \LR{$x \approx y$} به ولرو، چې له فرض سره خلاف دے۔ پايله دا ده چې \LR{$u \not\approx x$}۔ که \LR{$u \approx y$} وای، د اوسط په معادله کښې د ښي غړي معياري واټن د کيڼ غړي له واټن سره برابرېده، يا به ترې يو کم و۔ په دواړو حالتونو کښې د کيڼ او ښي غړو ټول واټن معياري کېده۔ له دې امله اوسط په درېيم بلاک کښې دے۔
\end{enumerate}
نو په شمېرېدونکي نامعياري مدل کښې نامعياري بلاکونه د ناطقو عددونو په شان مرتب دي: يو د شمېر وړ خطي ترتيب جوړوي چې سرحدي ټکي نۀ لري، او ګڼوالے ئې هم ثابت شو۔ ترې راځي چې د رښتيني حساب د هر دوو د شمېر وړ نامعياري مدلونو \LR{$\Struct{M}_1$} او \LR{$\Struct{M_2}$} راکمونې هغې ژبې ته چې يوازې \LR{$<$} او \LR{$=$} لري، ايزومورف دي۔ په رښتيا، ايزومورفيزم \LR{$h$} داسې ټاکلے شو: د \LR{$\Struct{M_1}$} او \LR{$\Struct{M_2}$} معياري برخې له معياري مدل \LR{$\Struct{N}$} سره ايزومورف دي، نو له يو بل سره هم۔ د نامعياري برخې بلاکونه پخپله د ناطقو عددونو په شان مرتب دي، نو د \olref[dlo]{thm:cantorQ} له مخې ايزومورف دي؛ د بلاکونو ايزومورفيزم د بلاکونو \emph{دننه} ايزومورفيزم ته غځولے شو: په هرې جوړې کښې په خپل سر ټاکل شوي عناصر يو له بل سره برابر کړو، او بيا په \LR{$\Struct{M_1}$} کښې د \LR{$x$} د جانشين انځور په \LR{$\Struct{M_2}$} کښې د \LR{$x$} د انځور جانشين وټاکو۔ پام وکړئ چې ترې دا \emph{نۀ} راځي چې \LR{$\mathfrak{M}_1$} او \LR{$\mathfrak{M}_2$} د حساب په بشپړه ژبه کښې ايزومورف دي (ايزومورفيزم تل له يوې ټاکلې نښلنډې سره تړلے دے)، ځکه پر \LR{$\Domain{M_1}$} او \LR{$\Domain{M_2}$} د جمع او ضرب د تعريفولو نۀ ايزومورف لارې شته۔ (دا د واوت له مشهورې قضيې څخه هم راځي چې د يوې بشپړې تيورۍ د شمېرېدونکو مدلونو شمېر دوه نۀ شي کېدے۔)
\olpseganchor{OLP-0653-B013}{end}

\olpseganchor{OLP-0653-B014}{start}
\begin{prob}
وښيئ چې د حساب په نامعياري مدل کښې تر ټولو لوے بلاک نۀ شي کېدے۔
\end{prob}
\olpseganchor{OLP-0653-B014}{end}

\olpseganchor{OLP-0653-B015}{start}
\begin{prob}
د لومړۍ درجې ژبه \LR{$\Lang{L}$} دې (له \LR{$\eq$} پرته) يوازې پريديکات \LR{$<$} ولري، او \LR{$\Struct{N} = (\Nat,
<)$} دې وي۔ د \LR{$\Struct{N}$} ټول متناهي يا هم‌متناهي فرعي سټونه د تعريف وړ دي۔ وښيئ چې د \LR{$\Struct{N}$} \emph{يوازې} همدا فرعي سټونه د تعريف وړ دي۔
\olpseganchor{OLP-0653-B015}{end}

\olpseganchor{OLP-0653-B016}{start}
(اشاره: لومړے، \LR{$\Obj{prc}(x,y)$} دې د \LR{$\Lang{L}$} هغه فارمول لنډ کړي چې وايي «\LR{$x$} د \LR{$y$} بې‌واسطې مخکښېنی دے»:
\[
x<y \land \lnot \lexists[z][(x<z \land z < y)].
\]
اوس د \LR{$\Struct{N}$} د هر تعريفېدونکي فرعي سټ سره په \LR{$\Lang{L}$} کښې يو فارمول \LR{$!A(x)$} تړلے دے۔ د هر داسې \LR{$!A$} دپاره دا جمله \LR{$!D$} وګورئ:
\[
\lexists[x][\lforall[y][\lforall[z][((x<y \land x<z \land \Obj{prc}(y,z)
  \land !A(y)) \lif !A(z))]]].
\]
وښيئ چې \LR{$\Sat{N}{!D}$} هله او يوازې هله چې د \LR{$\Struct{N}$} هغه فرعي سټ چې \LR{$!A$} ئې تعريفوي، متناهي يا هم‌متناهي وي۔
\olpseganchor{OLP-0653-B016}{end}

\olpseganchor{OLP-0653-B017}{start}
اوس \LR{$\Struct{M}$} دې له \LR{$\Struct{N}$} سره ابتدايي هم‌ارز نامعياري مدل وي۔ که \LR{$a \in \Domain{M}$} نامعياري وي، نو \LR{$b, c \in \Domain{M}$} دې له \LR{$a$} څخه لوے وي، او \LR{$b$} دې د~\LR{$c$} بې‌واسطې مخکښېنی وي۔ بيا د \LR{$\Domain{M}$} د ترتيبي جوړښت يو اتومورفيزم \LR{$h$} شته چې \LR{$h(b)=c$} (ولې؟)۔ ځکه که \LR{$b$} د \LR{$!A$} د صدق وړ کړي، نو \LR{$c$} هم همداسې کوي (ولې؟)۔ ترې راځي چې \LR{$!D$} په \LR{$\Struct{M}$} کښې رښتينې ده، نو په \LR{$\Struct{N}$} کښې هم۔ خو دا وايي چې د \LR{$\Struct{N}$} هغه فرعي سټ چې \LR{$!A$} ئې تعريفوي، متناهي يا هم‌متناهي دے۔
\end{prob}
\olpseganchor{OLP-0653-B017}{end}

\olpseganchor{OLP-0191-B004}{start}
\olchapter{mod}{mar}{د حساب مدلونه}
\olpseganchor{OLP-0191-B004}{end}

\olpseganchor{OLP-0192-B004}{start}
\olfileid{mod}{mar}{int}
\section{سريزه}
\olpseganchor{OLP-0192-B004}{end}

\olpseganchor{OLP-0192-B005}{start}
د حساب \emph{معياري مدل} هغه جوړښت~\LR{$\Struct{N}$} دے چې
\LR{$\Domain{N} = \Nat$} وي او \LR{$\Obj{0}$}، \LR{$\prime$}، \LR{$+$}، \LR{$\times$} او \LR{$<$}
هغسې تفسير شوي وي لکه تمه چې ترې کېږي۔ يعنې \LR{$\Obj{0}$}، \LR{$0$} دے،
\LR{$\prime$} د جانشين تابع ده، \LR{$+$} د جمع په توګه او \LR{$\times$} د ضرب په
توګه د~\LR{$\Nat$} د شمېرو تفسير لري۔ په ځانګړي ډول:
\begin{align*}
  \Assign{\Obj{0}}{N} & = 0\\
  \Assign{\prime}{N}(n) & = n + 1\\
  \Assign{+}{N}(n, m) & = n + m\\
  \Assign{\times}{N}(n, m) & = nm\\
  \tuple{n,m}\in\Assign{<}{N} & \liff n<m
\end{align*}
البته د~\LR{$\Lang{L_A}$} داسې جوړښتونه هم شته چې تعريف ساحې ئې له~\LR{$\Nat$}
څخه بېلې دي۔ د بېلګې په توګه، \LR{$\Struct{M}$} داسې ونيسئ چې
\LR{$\Domain{M} = \{a\}^*$} وي (د يوازېنۍ نښې~\LR{$a$} متناهي لړۍ؛ يعنې
\LR{$\emptyseq$}، \LR{$a$}، \LR{$aa$}، \LR{$aaa$}، \dots)، او تفسيرونه ئې داسې وي:
\begin{align*}
  \Assign{\Obj{0}}{M} & = \emptyseq\\
  \Assign{\prime}{M}(s) & = s \concat a\\
  \Assign{+}{M}(a^n, a^m) & = a^{n + m}\\
  \Assign{\times}{M}(a^n, a^m) & = a^{nm}\\
  \tuple{a^n,a^m}\in\Assign{<}{M} & \liff n<m
\end{align*}
دا دوه جوړښتونه په دې معنا «په اصل کښې يو شان» دي چې يوازينے
توپير ئې د غړي او د تعريف ساحې په څيزونو کښې دے، نه دا چې
د تفسير د تابعو له مخې د د تعريف ساحې غړي له يو بل سره څنګه تړاو لري۔
موږ وايو چې دا دوه جوړښتونه \emph{آيزومورف} دي۔
\olpseganchor{OLP-0192-B005}{end}

\olpseganchor{OLP-0192-B006}{start}
د متناهي شاهد له قضيه څخه په اسانۍ پايله کېږي چې هره تيوري چې په
\LR{$\Struct{N}$} کښې رښتونې وي، داسې مدلونه هم لري چې له~\LR{$\Struct{N}$} سره
آيزومورف نۀ وي۔ دا ډول جوړښتونه \emph{نامعياري} بلل کېږي۔ په معياري
مدل کښې (يعنې په~\LR{$\Struct{N}$} او له هغې سره په ټولو آيزومورفو
جوړښتونه کښې) ټول غړي د معياري عددنښو په ارزښتونو
\LR{$\num{n}$} کښې راټولېږي، يعنې:
\[
\Domain{N} = \Setabs{\Value{\num{n}}{N}}{n \in \Nat}
\]
خو په نامعياري مدلونو کښې داسې نۀ ده: که \LR{$\Struct{M}$} نامعياري وي، نو
لږ تر لږه يو \LR{$x \in \Domain{M}$} شته چې د هر~\LR{$n$} دپاره
\LR{$x \neq \Value{\num{n}}{M}$} وي۔
\olpseganchor{OLP-0192-B006}{end}

\olpseganchor{OLP-0192-B007}{start}
دا نامعياري عناصر په زړه پورې دي: «نامتناهي طبيعي شمېرې» دي۔
خو د هغو شتون د ناتکميلۍ پديدې هم په يوه معنا روښانوي۔ د بېلګې په
توګه د پيانو د حساب د سازګارۍ جمله، \LR{$\OCon[\Th{PA}]$}، يعنې
\LR{$\lnot \lexists[x][\OPrf[\Th{PA}](x, \gn{\lfalse})]$}، راواخلئ۔
د ناتکميلۍ د دويمې قضيې له مخې، \LR{$\Th{PA}$}، \LR{$\OCon[\Th{PA}]$} نۀ ثابتوي؛ د معياري مدل د سموالي له مخې ئې نفي هم نۀ ثابتوي، يعنې
\LR{$\lnot \OCon[\Th{PA}]$} نۀ ثابتوي۔ نو هره يوه جمله په سازګار ډول له~\LR{$\Th{PA}$}
سره يوځاے کېدے شي۔ څرنګه چې \LR{$\Th{PA}$} سازګاره ده،
\LR{$\Sat{N}{\OCon[\Th{PA}]}$} او په پايله کښې
\LR{$\Sat/{N}{\lnot \OCon[\Th{PA}]}$} لرو۔ نو \LR{$\Struct{N}$} د
\LR{$\Th{PA} \cup \{\lnot \OCon[\Th{PA}]\}$} مدل \emph{نۀ} دے، او ټول
مدلونه ئې بايد نامعياري وي۔ د \LR{$\Th{PA} \cup
\{\lnot \OCon[\Th{PA}]\}$} مدلونه بايد داسې يو غړے ولري چې د
\LR{$\lexists[x][\OPrf[\Th{PA}](x,\gn{\lfalse})]$} د رښتينولۍ شاهد وي، يعنې
له~\LR{$\Th{PA}$} څخه د تناقض د اشتقاق په توګه د دغه مدل له مخې شمېرل کېږي۔ داسې
غړے معياري کېدے نۀ شي، ځکه د هر~\LR{$n$} دپاره
\LR{$\Th{PA} \Proves \lnot \OPrf[\Th{PA}](\num{n}, \gn{\lfalse})$}۔
\olpseganchor{OLP-0192-B007}{end}

\olpseganchor{OLP-0193-B004}{start}
\olfileid{mod}{mar}{stm}
\olsection{د حساب معياري مدلونه}
\olpseganchor{OLP-0193-B004}{end}

\olpseganchor{OLP-0193-B005}{start}
د حساب ژبه~\LR{$\Lang{L_A}$} په څرګند ډول د شمېرو، په ځانګړي ډول د
طبيعي شمېرو، په هکله ده۔ نو «معياري مدل»~\LR{$\Struct{N}$} ځانګړے دے:
همدا هغه مدل دے چې موږ ئې په هکله خبرې کول غواړو۔ خو په منطق کښې
زموږ پام ډېر ځله يوازې جوړښتي ځانګړتياوو ته وي، او هر دوه آيزومورف
جوړښتونه دغه ځانګړتياوې ګډې لري۔ نو لږ پراخ فکر کولے شو او
هر هغه جوړښت چې له~\LR{$\Struct{N}$} سره آيزومورف وي «معياري»
وبولو۔
\olpseganchor{OLP-0193-B005}{end}

\olpseganchor{OLP-0193-B006}{start}
\begin{defn}
د~\LR{$\Lang{L_A}$} يو جوړښت \emph{معياري} دے، که له~\LR{$\Struct{N}$}
سره آيزومورف وي۔
\end{defn}
\olpseganchor{OLP-0193-B006}{end}

\olpseganchor{OLP-0193-B007}{start}
\begin{prop}
\ollabel{prop:standard-domain} که جوړښت~\LR{$\Struct{M}$} معياري
وي، نو تعريف ساحه ئې د معياري عددنښو د ارزښتونو سټ دے، يعنې:
\[
\Domain{M} = \Setabs{\Value{\num{n}}{M}}{n \in \Nat}
\]
\end{prop}
\olpseganchor{OLP-0193-B007}{end}

\olpseganchor{OLP-0193-B008}{start}
\begin{proof}
څرګنده ده چې هر \LR{$\Value{\num{n}}{M} \in \Domain{M}$} دے۔ يوازې دا
ښودل پاتې دي چې هر \LR{$x \in \Domain{M}$} د کوم~\LR{$n$} دپاره
\LR{$\Value{\num{n}}{M}$} دے۔ څرنګه چې \LR{$\Struct{M}$} معياري دے، له
\LR{$\Struct{N}$} سره آيزومورف دے۔ فرض کړئ \LR{$g\colon \Nat \to \Domain{M}$}
آيزومورفيزم وي۔ بيا \LR{$g(n) = g(\Value{\num{n}}{N}) =
\Value{\num{n}}{M}$}۔ خو د هر \LR{$x \in \Domain{M}$} دپاره داسې~\LR{$n \in
\Nat$} شته چې \LR{$g(n) = x$} وي، ځکه \LR{$g$} پر هدف سټ بشپړه ده۔
\end{proof}
\olpseganchor{OLP-0193-B008}{end}

\olpseganchor{OLP-0193-B009}{start}
\begin{explain}
که د~\LR{$\Lang{L_A}$} يو جوړښت~\LR{$\Struct{M}$} معياري وي، د~د تعريف ساحه
ټول غړي ئې د معياري عددنښو \LR{$\num{0}$}، \LR{$\num{1}$}، \LR{$\num{2}$}،
\dots په نوم يادېدے شي، يعنې د \LR{$\Obj{0}$}، \LR{$\Obj{0}'$}،
\LR{$\Obj{0}''$} او داسې نورو ترمونو په وسيله۔ البته دا معنا نۀ لري چې د
\LR{$\Domain{M}$} غړي پخپله شمېرې دي؛ يوازې دومره ده چې د هغو
په نښه کولے شو، لکه څنګه چې په~\LR{$\Domain{N}$} کښې شمېرې نښه کوو۔
\end{explain}
\olpseganchor{OLP-0193-B009}{end}

\olpseganchor{OLP-0193-B010}{start}
\begin{prob}
وښيئ چې د~\olref[mod][mar][stm]{prop:standard-domain} معکوس ناسم دے؛
يعنې داسې يو جوړښت~\LR{$\Struct{M}$} ورکړئ چې
\LR{$\Domain{M} = \Setabs{\Value{\num{n}}{M}}{n \in \Nat}$} وي، خو له
\LR{$\Struct{N}$} سره آيزومورف نۀ وي۔
\end{prob}
\olpseganchor{OLP-0193-B010}{end}

\olpseganchor{OLP-0193-B011}{start}
\begin{prop}
\ollabel{prop:thq-standard}
که \LR{$\Sat{M}{\Th{Q}}$} او \LR{$\Domain{M} =
\Setabs{\Value{\num{n}}{M}}{n \in \Nat}$} وي، نو \LR{$\Struct{M}$} معياري دے۔
\end{prop}
\olpseganchor{OLP-0193-B011}{end}

\olpseganchor{OLP-0193-B012}{start}
\begin{proof}
بايد وښيو چې \LR{$\Struct{M}$} له~\LR{$\Struct{N}$} سره آيزومورف دے۔ تابع
\LR{$g\colon \Nat \to \Domain{M}$} داسې ونيسئ: \LR{$g(n) =
\Value{\num{n}}{M}$}۔ د فرض له مخې \LR{$g$} پر هدف سټ بشپړه ده۔ دا
يو پر يو هم ده: هر کله چې \LR{$n \neq m$} وي، \LR{$\Th{Q} \Proves
\eq/[\num{n}][\num{m}]$}۔ نو له \LR{$\Sat{M}{\Th{Q}}$} څخه
\LR{$\Sat{M}{\eq/[\num{n}][\num{m}]}$}، هر کله چې \LR{$n \neq m$}۔ په پايله
کښې \LR{$n \neq m$} وي، نو \LR{$\Value{\num{n}}{M} \neq
\Value{\num{m}}{M}$}، يعنې \LR{$g(n) \neq g(m)$}۔
\olpseganchor{OLP-0193-B012}{end}

\olpseganchor{OLP-0193-B013}{start}
اوس بايد دا هم وڅېړو چې \LR{$g$} آيزومورفيزم دے۔
\begin{enumerate}
\item \LR{$g(\Assign{\Obj{0}}{N}) = g(0)$}، ځکه \LR{$\Assign{\Obj{0}}{N} = 0$}۔
  د~\LR{$g$} له تعريفه \LR{$g(0) = \Value{\num{0}}{M}$}۔ خو \LR{$\num{0}$} هماغه
  \LR{$\Obj{0}$} دے، او د داسې ترم ارزښت چې ثابت سمبول وي، هغه څه دے
  چې جوړښت ئې دغه ثابت سمبول ته ټاکي، يعنې \LR{$\Value{\Obj{0}}{M} =
  \Assign{\Obj{0}}{M}$}۔ نو \LR{$g(\Assign{\Obj{0}}{N}) =
  \Assign{\Obj{0}}{M}$}، لکه اړين چې و۔
\item \LR{$g(\Assign{\prime}{N}(n)) = g(n+1)$}، ځکه په~\LR{$\Struct{N}$} کښې
  \LR{$\prime$} پر~\LR{$\Nat$} د جانشين تابع ده۔ بيا د~\LR{$g$} له تعريفه
  \LR{$g(n+1) = \Value{\num{n+1}}{M}$}۔ خو \LR{$\num{n+1}$} له~\LR{$\num{n}'$}
  سره يو ترم دے، نو \LR{$\Value{\num{n+1}}{M} =
  \Value{\num{n}'}{M}$}۔ د ارزښت د تابعې له تعريفه دا
  \LR{$= \Assign{\prime}{M}(\Value{\num{n}}{M})$} دے۔ څرنګه چې
  \LR{$\Value{\num{n}}{M} = g(n)$}، نو
  \LR{$g(\Assign{\prime}{N}(n)) = \Assign{\prime}{M}(g(n))$}۔
\item \LR{$g(\Assign{+}{N}(n,m)) = g(n+m)$}، ځکه په~\LR{$\Struct{N}$} کښې \LR{$+$}
  د~\LR{$\Nat$} د جمع تابع ده۔ بيا د~\LR{$g$} له تعريفه
\LR{$g(n+m) = \Value{\num{n+m}}{M}$}۔
  خو \LR{$\Th{Q} \Proves \num{n+m} = (\num{n} + \num{m})$}، نو
  \LR{$\Value{\num{n+m}}{M} = \Value{\num{n}+\num{m}}{M}$}۔ د ارزښت د
  تابعې له تعريفه دا
  \LR{$= \Assign{+}{M}(\Value{\num{n}}{M},\Value{\num{m}}{M})$} دے۔
  څرنګه چې \LR{$\Value{\num{n}}{M} = g(n)$} او
  \LR{$\Value{\num{m}}{M} = g(m)$}، نو
  \LR{$g(\Assign{+}{N}(n, m)) = \Assign{+}{M}(g(n), g(m))$}۔
\item \LR{$g(\Assign{\times}{N}(n, m)) = \Assign{\times}{M}(g(n), g(m))$}:
  تمرين۔
\item \LR{$\tuple{n,m} \in \Assign{<}{N}$} هله او يوازې هله چې \LR{$n < m$}۔
  که \LR{$n < m$} وي، نو \LR{$\Th{Q} \Proves \num{n} < \num{m}$} او
  \LR{$\Sat{M}{\num{n} < \num{m}}$} هم لرو۔ نو
  \LR{$\tuple{\Value{\num{n}}{M}, \Value{\num{m}}{M}} \in
  \Assign{<}{M}$}، يعنې \LR{$\tuple{g(n), g(m)} \in \Assign{<}{M}$}۔
  که \LR{$n \not< m$} وي، نو \LR{$\Th{Q} \Proves \lnot \num{n} < \num{m}$} او
  په پايله کښې \LR{$\Sat/{M}{\num{n} < \num{m}}$}۔ نو، لکه مخکې،
  \LR{$\tuple{g(n), g(m)} \notin \Assign{<}{M}$}۔ په ګډه:
  \LR{$\tuple{n,m} \in \Assign{<}{N}$} هله او يوازې هله چې
  \LR{$\tuple{g(n), g(m)} \in \Assign{<}{M}$}۔
\end{enumerate}
\end{proof}
\olpseganchor{OLP-0193-B013}{end}

\olpseganchor{OLP-0193-B014}{start}
\begin{explain}
تابع~\LR{$g$} د~\LR{$\Lang{L_A}$} د هر بل جوړښت~\LR{$\Struct{M}$} د تعريف
ساحې ته له~\LR{$\Nat$} څخه د نقشې تر ټولو څرګنده لار ده، ځکه هر داسې
\LR{$\Struct{M}$} هغه غړي لري چې \LR{$\num{0}$}، \LR{$\num{1}$}،
\LR{$\num{2}$} او داسې نورو په نومونو يادېږي۔ نو د حيرانتيا خبره نۀ ده چې
که \LR{$\Struct{M}$} د~\LR{$\num{n}$} په باب لږ تر لږه ځينې بنسټيزې ويناوې
هماغسې رښتونې کړي لکه~\LR{$\Struct{N}$} او \LR{$g$} هم دوه اړخيزه يو په يو
وي، نو \LR{$g$} آيزومورفيزم ګرځي۔ په حقيقت کښې که \LR{$\Domain{M}$} له هغو
غړي پرته چې \LR{$\num{n}$} ئې نوموي بل څه نۀ لري، نو همدا
يوازينے آيزومورفيزم دے۔
\end{explain}
\olpseganchor{OLP-0193-B014}{end}

\olpseganchor{OLP-0193-B015}{start}
\begin{prop}
\ollabel{prop:thq-unique-iso} که \LR{$\Struct{M}$} معياري وي، نو د
\olref{prop:thq-standard} له ثبوت څخه \LR{$g$} له~\LR{$\Struct{N}$} څخه
\LR{$\Struct{M}$} ته يوازينے آيزومورفيزم دے۔
\end{prop}
\olpseganchor{OLP-0193-B015}{end}

\olpseganchor{OLP-0193-B016}{start}
\begin{proof}
فرض کړئ \LR{$h\colon \Nat \to \Domain{M}$} له~\LR{$\Struct{N}$} څخه
\LR{$\Struct{M}$} ته آيزومورفيزم وي۔ د~\LR{$n$} پر بنسټ استقرا ښيو چې \LR{$g=h$}۔
که \LR{$n=0$} وي، نو د~\LR{$g$} له تعريفه \LR{$g(0)=\Assign{\Obj{0}}{M}$}۔ خو
څرنګه چې \LR{$h$} آيزومورفيزم دے، \LR{$h(0)=h(\Assign{\Obj{0}}{N}) =
\Assign{\Obj{0}}{M}$}، نو \LR{$g(0)=h(0)$}۔
\olpseganchor{OLP-0193-B016}{end}

\olpseganchor{OLP-0193-B017}{start}
اوس د \LR{$n+1$} حالت راواخلئ:
\begin{align*}
  g(n+1) & = \Value{\num{n+1}}{M} \text{\RL{ د \LR{$g$} له تعريفه}}\\
  & = \Value{\num{n}'}{M} \text{\RL{ ځکه \LR{$\num{n+1}\ident \num{n}'$}}}\\
  & = \Assign{\prime}{M}(\Value{\num{n}}{M})
    \text{\RL{ د \LR{$\Value{t'}{M}$} له تعريفه}}\\
  & = \Assign{\prime}{M}(g(n)) \text{\RL{ د \LR{$g$} له تعريفه}}\\
  & = \Assign{\prime}{M}(h(n)) \text{\RL{ د استقرايي فرض له مخې}}\\
  & = h(\Assign{\prime}{N}(n)) \text{\RL{ ځکه \LR{$h$} آيزومورفيزم دے}}\\
  & = h(n+1)
\end{align*}
\end{proof}
\olpseganchor{OLP-0193-B017}{end}

\olpseganchor{OLP-0193-B018}{start}
\begin{explain}
د هر شمېرېدونکے لامتناهي سټ~\LR{$M$} دپاره د~\LR{$\Nat$} او~\LR{$M$} تر منځ يوه
دوه اړخيزه يو پر يو تابع شته، نو هر داسې~\LR{$M$} په بالقوه ډول د معياري
مدل~\LR{$\Struct{M}$} د تعريف ساحه کېدے شي۔ په حقيقت کښې، کله چې يو څيز
\LR{$z \in M$} او مناسبه تابع~\LR{$s$} د~\LR{$\Assign{\Obj{0}}{M}$} او
\LR{$\Assign{\prime}{M}$} په توګه وټاکئ، د \LR{$+$}، \LR{$\times$} او \LR{$<$} تفسيرونه
لا له مخکې ټاکل شوي وي۔ يوازې هغه تابعې~\LR{$s\colon M \to M \setminus
\{z\}$} په معياري مدل کښې د~\LR{$\Assign{\prime}{M}$} په توګه مناسبې دي چې
يو پر يو او پر هدف سټ بشپړه دواړه وي۔ د~\LR{$s$} د قيمتونو ساحه
\LR{$z$} لرلے نۀ شي، ځکه که داسې واے نو
\LR{$\lforall[x][\eq/[\Obj 0][x']]$} به دروغ واے۔ دا تړلے فارمول
په~\LR{$\Struct{N}$} کښې رښتونې ده، نو \LR{$\Struct{M}$} ئې هم بايد رښتونې کړي۔ تابع~\LR{$s$} بايد
يو پر يو وي، ځکه په~\LR{$\Struct{N}$} کښې د جانشين
تابع~\LR{$\Assign{\prime}{N}$} همداسې ده، او د هغې
\LR{$\Assign{\prime}{N}$} يو پر يو توب په~\LR{$\Struct{N}$} کښې د
تړلے فارمول له خوا بيانېږي۔ دا بايد
پر هدف سټ بشپړه هم وي، ځکه که نه نو داسې~\LR{$x \in M \setminus \{z\}$}
به وي چې د~\LR{$s$} د قيمتونو په ساحه کښې نۀ وي، يعنې دا تړلے فارمول
\LR{$\lforall[x][(\eq[x][\Obj 0] \lor
\lexists[y][\eq[y'][x]])]$} به په~\LR{$\Struct{M}$} کښې دروغ وي، حال دا چې
په~\LR{$\Struct{N}$} کښې رښتونې ده۔ د صفر له بيا بيا تالي اخستلو څخه بايد د ډومېن هر غړے ترلاسه شي۔
يوازې د تابعې يو پر يو او پر شتون کافي نۀ دي، ځکه له صفر څخه جلا د تالي لړۍ هم پاتې کېدے شي۔ کله چې د صفر له تالي څخه ټول غړي ترلاسه شي، همدغه لړۍ د طبيعي شمېرو او ډومېن تر منځ آيزومورفيزم ټاکي، او جمع، ضرب او ترتيب د هغې له مخې تعبيروو۔
\end{explain}
\olpseganchor{OLP-0193-B018}{end}

\olpseganchor{OLP-0194-B004}{start}
\olfileid{mod}{mar}{nst}
\section{نامعياري مدلونه}
\olpseganchor{OLP-0194-B004}{end}

\olpseganchor{OLP-0194-B005}{start}
\begin{explain}
د~\LR{$\Lang{L_A}$} يو جوړښت هغه وخت معياري بولو چې له
\LR{$\Struct{N}$} سره آيزومورف وي۔ که يو جوړښت له~\LR{$\Struct{N}$}
سره آيزومورف نۀ وي، نامعياري بلل کېږي۔
\end{explain}
\olpseganchor{OLP-0194-B005}{end}

\olpseganchor{OLP-0194-B006}{start}
\begin{defn}
د~\LR{$\Lang{L_A}$} يو جوړښت~\LR{$\Struct{M}$} \emph{نامعياري} دے،
که له~\LR{$\Struct{N}$} سره آيزومورف نۀ وي۔ هغه غړي
\LR{$x \in \Domain{M}$} چې د کوم \LR{$n \in \Nat$} دپاره
\LR{$\Value{\num{n}}{M}$} وي، د~\LR{$\Struct{M}$} \emph{معياري شمېرې} بلل
کېږي؛ او نورې \emph{نامعياري شمېرې} بلل کېږي۔
\end{defn}
\olpseganchor{OLP-0194-B006}{end}

\olpseganchor{OLP-0194-B007}{start}
\begin{explain}
د~\olref[stm]{prop:standard-domain} له مخې د~\LR{$\Lang{L_A}$} هر
معياري جوړښت يوازې معياري غړي لري۔ نو يو نامعياري
جوړښت بايد لږ تر لږه يو نامعياري عنصر ولري۔ په حقيقت کښې،
د يوه نامعياري غړے شتون دا تضمينوي چې جوړښت
نامعياري دے۔
\end{explain}
\olpseganchor{OLP-0194-B007}{end}

\olpseganchor{OLP-0194-B008}{start}
\begin{prop}
که د~\LR{$\Lang{L_A}$} يو جوړښت~\LR{$\Struct{M}$} يوه نامعياري شمېره
ولري، نو \LR{$\Struct{M}$} نامعياري دے۔
\end{prop}
\olpseganchor{OLP-0194-B008}{end}

\olpseganchor{OLP-0194-B009}{start}
\begin{proof}
برعکس فرض کړئ، يعنې \LR{$\Struct{M}$} معياري وي خو نامعياري شمېره~\LR{$x$}
ولري۔ \LR{$g\colon \Nat \to \Domain{M}$} يو آيزومورفيزم ونيسئ۔ په
اسانه (پر~\LR{$n$} په استقرا) ليدل کېږي چې
\LR{$g(\Value{\num{n}}{N}) = \Value{\num{n}}{M}$}۔ په بله وينا، \LR{$g$} د
\LR{$\Struct{N}$} معياري شمېرې د~\LR{$\Struct{M}$} معياري شمېرو ته وړي۔ که
\LR{$\Struct{M}$} يوه نامعياري شمېره ولري، \LR{$g$} پر هدف سټ بشپړه کېدے نۀ
شي، چې له فرض سره ټکر دے۔
\end{proof}
\olpseganchor{OLP-0194-B009}{end}

\olpseganchor{OLP-0194-B010}{start}
\begin{prob}
په ياد راوړئ چې \LR{$\Th{Q}$} لاندې بديهي اصول لري:
\begin{align*}
& \lforall[x][\lforall[y][(\eq[x'][y'] \lif \eq[x][y])]] \tag{\LR{$!Q_1$}}\\
& \lforall[x][\eq/[\Obj 0][x']] \tag{\LR{$!Q_2$}}\\
& \lforall[x][(\eq[x][\Obj 0] \lor \lexists[y][\eq[x][y']])] \tag{\LR{$!Q_3$}}
\end{align*}
داسې جوړښتونه~\LR{$\Struct{M_1}$}، \LR{$\Struct{M_2}$}، \LR{$\Struct{M_3}$}
ورکړئ چې:
\begin{enumerate}
\item \LR{$\Sat{M_1}{!Q_1}$}، \LR{$\Sat{M_1}{!Q_2}$}، \LR{$\Sat/{M_1}{!Q_3}$}؛
\item \LR{$\Sat{M_2}{!Q_1}$}، \LR{$\Sat/{M_2}{!Q_2}$}، \LR{$\Sat{M_2}{!Q_3}$}؛ او
\item \LR{$\Sat/{M_3}{!Q_1}$}، \LR{$\Sat{M_3}{!Q_2}$}، \LR{$\Sat{M_3}{!Q_3}$}۔
\end{enumerate}
څرګنده ده چې د هر يوه دپاره يوازې~\LR{$\Assign{\Obj{0}}{M_i}$} او
\LR{$\Assign{\prime}{M_i}$} ټاکل پکار دي۔
\end{prob}
\olpseganchor{OLP-0194-B010}{end}

\olpseganchor{OLP-0194-B011}{start}
\begin{explain}
د~\LR{$\Lang{L_A}$} دپاره نامعياري جوړښتونه په اسانه ټاکل کېدے
شي۔ د بېلګې په توګه، داسې جوړښت ونيسئ چې د تعريف ساحه ئې~\LR{$\Int$} وي او
ټولې غيرمنطقي نښې په معمولي ډول تفسير کړي۔ څرنګه چې منفي شمېرې د
هېڅ~\LR{$n$} دپاره د~\LR{$\num{n}$} ارزښتونه نۀ دي، دا جوړښت نامعياري دے۔
البته دا به د حساب \emph{مدل} په دې معنا نۀ وي چې هماغه جملې رښتونې
کړي چې~\LR{$\Struct{N}$} ئې رښتونې کوي۔ مثلاً
\LR{$\lforall[x][\eq/[x'][\Obj{0}]]$} دروغ ده۔ خو د متناهي شاهد له قضيې
څخه په اسانۍ ثابتولے شو چې د حساب نامعياري مدلونه شته۔
\end{explain}
\olpseganchor{OLP-0194-B011}{end}

\olpseganchor{OLP-0194-B012}{start}
\begin{prop}
فرض کړئ \LR{$\Th{TA} = \Setabs{!A}{\Sat{N}{!A}}$} د~\LR{$\Struct{N}$} تيوري
ده۔ \LR{$\Th{TA}$} يو د شمېر وړ نامعياري مدل لري۔
\end{prop}
\olpseganchor{OLP-0194-B012}{end}

\olpseganchor{OLP-0194-B013}{start}
\begin{proof}
ژبه~\LR{$\Lang{L_A}$} په يوه نوي ثابت سمبول~\LR{$c$} وغځوئ او د
تړلي فارمولونه لاندې سټ ونيسئ:
\[
\Gamma = \Th{TA} \cup \{\eq/[c][\num{0}], \eq/[c][\num{1}],
\eq/[c][\num{2}], \dots\}
\]
د~\LR{$\Gamma$} هر مدل~\LR{$\Struct{M^c}$} به داسې يو غړے
\LR{$x = \Assign{c}{M}$} ولري چې نامعياري وي، ځکه د هر~\LR{$n \in \Nat$}
دپاره \LR{$x \neq \Value{\num{n}}{M}$}۔ دا هم څرګنده ده چې
\LR{$\Sat{M^c}{\Th{TA}}$}، ځکه \LR{$\Th{TA} \subseteq \Gamma$}۔ که له
\LR{$\Struct{M^c}$} څخه يوازې د~\LR{$c$} تفسير هېر کړو او هغه د~\LR{$\Lang{L_A}$}
يو جوړښت~\LR{$\Struct{M}$} وګرځوو، تعريف ساحه ئې لا هم نامعياري
\LR{$x$} لري او \LR{$\Sat{M}{\Th{TA}}$} هم وي۔ وروستۍ خبره ځکه تضمينېږي چې
\LR{$c$} په~\LR{$\Th{TA}$} کښې نۀ راځي۔ نو کافي ده وښيو چې \LR{$\Gamma$} مدل لري۔
\olpseganchor{OLP-0194-B013}{end}

\olpseganchor{OLP-0194-B014}{start}
د متناهي شاهد قضيه کاروو۔ که د~\LR{$\Gamma$} هر متناهي فرعي سټ د صدق وړ
وي، نو \LR{$\Gamma$} هم د صدق وړ دے۔ يو متناهي فرعي سټ
\LR{$\Gamma_0 \subseteq \Gamma$} ونيسئ۔ په~\LR{$\Gamma_0$} کښې د~\LR{$\Th{TA}$}
ځينې تړلي فارمولونه او ښايي د~\LR{$\eq/[c][\num{n}]$} د بڼې متناهي ډېرې جملې وي۔
که د دې بڼې کومه جمله وي، فرض کړئ~\LR{$k$} تر ټولو لوے عدد دے چې
\LR{$\eq/[c][\num{k}] \in \Gamma_0$} وي؛ او که هېڅ داسې جمله نه وي، د
k ارزښت صفر ونيسئ۔ \LR{$\Struct{N_k}$} د~\LR{$\Struct{N}$} داسې غځونه تعريف کړئ چې
\LR{$\Assign{c}{N_k} = k+1$} وي۔ بيا \LR{$\Sat{N_k}{\Gamma_0}$}: که
\LR{$!A \in \Th{TA}$} وي، \LR{$\Sat{N_k}{!A}$}، ځکه \LR{$\Struct{N_k}$} له~\LR{$c$}
پرته په هر څه کښې د~\LR{$\Struct{N}$} په شان دے او \LR{$c$} په~\LR{$!A$} کښې نۀ
راځي۔ او د دغې بڼې هرې راغلې جملې دپاره
\LR{$\Sat{N_k}{\eq/[c][\num{n}]}$}، ځکه \LR{$n \le k$} او
\LR{$\Value{c}{N_k} = k+1$}۔ نو د~\LR{$\Gamma$} هر متناهي فرعي سټ د صدق وړ
دے۔ د متناهي شاهد له قضيې \LR{$\Gamma$} يو مدل لري؛ او څرنګه چې ژبه
شمېرېدونکې او دغه مدل نامتناهي دے، د ښکته لوېنهايم--سکولم له قضيې
يو شمېرېدونکے داسې مدل اخلو۔
\end{proof}
\olpseganchor{OLP-0194-B014}{end}

\olpseganchor{OLP-0195-B004}{start}
\olfileid{mod}{mar}{mdq}
\section{د~\texorpdfstring{\LR{$\Th{Q}$}}{Q} مدلونه}
\olpseganchor{OLP-0195-B004}{end}

\olpseganchor{OLP-0195-B005}{start}
\begin{explain}
موږ پوهېږو چې داسې نامعياري جوړښتونه شته چې هماغه
تړلي فارمولونه رښتونې کوي چې~\LR{$\Struct{N}$} ئې رښتونې کوي؛ يعنې د
\LR{$\Th{TA}$} مدلونه دي۔ څرنګه چې \LR{$\Sat{N}{\Th{Q}}$}، د~\LR{$\Th{TA}$} هر
مدل د~\LR{$\Th{Q}$} هم مدل دے۔ \LR{$\Th{Q}$} له~\LR{$\Th{TA}$} څخه ډېره کمزورې
ده؛ مثلاً
\LR{$\Th{Q} \Proves/ \lforall[x][\lforall[y][\eq[(x + y)][(y+x)]]]$}۔
کمزورې تيورۍ په اسانه پوره کېږي: ډېر مدلونه لري۔ مثلاً \LR{$\Th{Q}$}
داسې مدلونه لري چې
\LR{$\lforall[x][\lforall[y][\eq[(x + y)][(y+x)]]]$} دروغوي، خو هغه د
\LR{$\Th{TA}$} يا حتا د~\LR{$\Th{PA}$} مدلونه نۀ شي کېدے۔ د~\LR{$\Th{Q}$} مدلونه
نسبتاً ساده هم دي: په څرګند ډول ئې ټاکلے شو۔
\end{explain}
\olpseganchor{OLP-0195-B005}{end}

\olpseganchor{OLP-0195-B006}{start}
\begin{ex}
\ollabel{ex:model-K-of-Q}
هغه جوړښت~\LR{$\Struct{K}$} ونيسئ چې تعريف ساحه ئې
\LR{$\Domain{K} = \Nat \cup \{a\}$} او تفسيرونه ئې دا وي:
\begin{align*}
  \Assign{\Obj{0}}{K} & = 0\\
  \Assign{\prime}{K}(x) & =
  \begin{cases}
    x+1 & \text{\RL{که \LR{$x\in \Nat$} وي}}\\
    a & \text{\RL{که \LR{$x = a$} وي}}
  \end{cases}\\
  \Assign{+}{K}(x, y) & =
  \begin{cases}
    x+y & \text{\RL{که \LR{$x$}، \LR{$y \in\Nat$} وي}}\\
    a & \text{\RL{په نورو حالاتو کښې}}
  \end{cases}\\
  \Assign{\times}{K}(x, y) & =
  \begin{cases}
    xy & \text{\RL{که \LR{$x$}، \LR{$y \in\Nat$} وي}}\\
    0 & \text{\RL{که \LR{$x = 0$} يا \LR{$y = 0$} وي}}\\
    a & \text{\RL{په نورو حالاتو کښې}}\\
  \end{cases}\\
  \Assign{<}{K} & =
  \Setabs{\tuple{x,y}}{x, y \in \Nat \text{\RL{ او }} x<y} \cup
  \Setabs{\tuple{x,a}}{x \in \Domain{K}}
\end{align*}
د~\LR{$\Sat{K}{\Th{Q}}$} ښودلو دپاره بايد ثابته کړو چې د~\LR{$\Th{Q}$} ټول
بديهي اصول په~\LR{$\Struct{K}$} کښې رښتوني دي۔ د اسانتيا دپاره
\LR{$x^\nssucc$} د~\LR{$\Assign{\prime}{K}(x)$} په ځاے ليکو (په~\LR{$\Struct{K}$}
کښې د~\LR{$x$} «جانشين»)، \LR{$x \nsplus y$} د~\LR{$\Assign{+}{K}(x,y)$} په ځاے
(په~\LR{$\Struct{K}$} کښې د~\LR{$x$} او~\LR{$y$} «مجموعه»)، \LR{$x \nstimes y$} د
\LR{$\Assign{\times}{K}(x,y)$} په ځاے (په~\LR{$\Struct{K}$} کښې د~\LR{$x$}
او~\LR{$y$} «حاصل ضرب»)، او
\LR{$x \nsless y$} د~\LR{$\tuple{x,y}\in\Assign{<}{K}$} په ځاے ليکو۔ په دې
لنډونو سره د~\LR{$\Struct{K}$} عمليات په لا څرګند ډول داسې ښودلے شو:
\[
\begin{array}{c|c}
  x & x^\nssucc \\
  \hline
  n & n+1 \\
  a & a
\end{array}
\qquad
\begin{array}{c|ccc}
  x \nsplus y & 0 & m & a \\
  \hline
  0 & 0 & m & a \\
  n & n & n+m & a \\
  a & a & a & a \\
\end{array}
\qquad
\begin{array}{c|ccc}
  x \nstimes y & 0 & m & a \\
  \hline
  0 & 0 & 0 & 0 \\
  n & 0 & nm & a \\
  a & 0 & a & a \\
\end{array}
\]
\LR{$n\nsless m$} هله او يوازې هله چې \LR{$n<m$}، د~\LR{$n$}، \LR{$m\in\Nat$}
دپاره؛ او د هر~\LR{$x\in\Domain{K}$} دپاره \LR{$x\nsless a$}۔
\olpseganchor{OLP-0195-B006}{end}

\olpseganchor{OLP-0195-B007}{start}
\LR{$\Sat{K}{\lforall[x][\lforall[y][(\eq[x'][y'] \lif \eq[x][y])]]}$}،
ځکه \LR{$\nssucc$} يو پر يو ده۔
\LR{$\Sat{K}{\lforall[x][\eq/[\Obj 0][x']]}$}، ځکه په~\LR{$\Struct{K}$} کښې
\LR{$0$} د~\LR{$\nssucc$} جانشين نۀ دے۔
\LR{$\Sat{K}{\lforall[x][(\eq[x][\Obj 0] \lor
\lexists[y][\eq[x][y']])]}$}، ځکه د هر~\LR{$n>0$} دپاره
\LR{$n=(n-1)^\nssucc$} او \LR{$a=a^\nssucc$}۔
\olpseganchor{OLP-0195-B007}{end}

\olpseganchor{OLP-0195-B008}{start}
\LR{$\Sat{K}{\lforall[x][\eq[(x + \Obj 0)][x]]}$}، ځکه
\LR{$n\nsplus0=n+0=n$} او د~\LR{$\nsplus$} له تعريفه \LR{$a\nsplus0=a$}۔
\LR{$\Sat{K}{\lforall[x][\lforall[y][\eq[(x + y')][(x + y)']]]}$} لږه
پېچلې ده۔ که \LR{$n$} او~\LR{$m$} دواړه معياري وي:
\begin{align*}
(n \nsplus m^\nssucc) & = (n+(m+1)) = (n+m)+1 = (n \nsplus m)^\nssucc
\intertext{\RL{ځکه \LR{$\nsplus$} او \LR{$^\nssucc$} پر معياري شمېرو له~\LR{$+$} او
  \LR{$\prime$} سره سلا دي۔ اوس فرض کړئ \LR{$x \in \Domain{K}$}۔ بيا}}
(x \nsplus a^\nssucc) & = (x \nsplus a) = a = a^\nssucc = (x \nsplus a)^\nssucc
\intertext{\RL{پاتې حالت دا دے چې \LR{$y \in \Domain{K}$} وي خو \LR{$x=a$}۔
  دلته هم بايد وڅېړو چې \LR{$y=n$} معياري دے يا \LR{$y=a$}:}}
(a \nsplus n^\nssucc) & = (a \nsplus (n+1)) = a = a^\nssucc = (a \nsplus n)^\nssucc\\
(a \nsplus a^\nssucc) & = (a \nsplus a) = a = a^\nssucc = (a \nsplus a)^\nssucc
\end{align*}
البته دا له اړتيا زيات تفصيل دے۔ مثلاً څرنګه چې د هر~\LR{$z$} دپاره
\LR{$a\nsplus z=a$}، نو سمدستي \LR{$a\nsplus a^\nssucc=a$} پايله کولے شو۔
پاتې بديهي اصول هم په همدې ډول کتل کېدے شي۔
\olpseganchor{OLP-0195-B008}{end}

\olpseganchor{OLP-0195-B009}{start}
نو \LR{$\Struct{K}$} د~\LR{$\Th{Q}$} يو مدل دے۔ «جمع»~\LR{$\nsplus$} ئې تبادلي هم
ده۔ خو داسې نورې جملې شته چې په~\LR{$\Struct{N}$} کښې رښتونې او په
\LR{$\Struct{K}$} کښې دروغ دي، او برعکس۔ مثلاً \LR{$a\nsless a$}، نو
\LR{$\Sat{K}{\lexists[x][x < x]}$} او
\LR{$\Sat/{K}{\lforall[x][\lnot x<x]}$}۔ له دې څرګندېږي چې
\LR{$\Th{Q} \Proves/ \lforall[x][\lnot x < x]$}۔
\end{ex}
\olpseganchor{OLP-0195-B009}{end}

\olpseganchor{OLP-0195-B010}{start}
\begin{prob}
ثابته کړئ چې د~\olref[mod][mar][mdq]{ex:model-K-of-Q} جوړښت
\LR{$\Struct{K}$} د~\LR{$\Th{Q}$} پاتې بديهي اصول پوره کوي:
\begin{align*}
  & \lforall[x][\eq[(x \times \Obj 0)][\Obj 0]] \tag{\LR{$!Q_6$}}\\
  & \lforall[x][\lforall[y][\eq[(x \times y')][((x \times y) + x)]]] \tag{\LR{$!Q_7$}}\\
  & \lforall[x][\lforall[y][(x < y \liff \lexists[z][\eq[(z' + x)][y]])]] \tag{\LR{$!Q_8$}}
\end{align*}
داسې تړلے فارمول ومومئ چې يوازې~\LR{$\prime$} ولري، په~\LR{$\Struct{N}$}
کښې رښتونې او په~\LR{$\Struct{K}$} کښې دروغ وي۔
\end{prob}
\olpseganchor{OLP-0195-B010}{end}

\olpseganchor{OLP-0195-B011}{start}
\begin{ex}
\ollabel{ex:model-L-of-Q} هغه جوړښت~\LR{$\Struct{L}$} ونيسئ چې
تعريف ساحه ئې \LR{$\Domain{L} = \Nat \cup \{a,b\}$} او تفسيرونه
\LR{$\Assign{\prime}{L} = \nssucc$}، \LR{$\Assign{+}{L} = \nsplus$} د لاندې
جدولونو له مخې وي:
\[
\begin{array}{c|c}
  x & x^\nssucc \\
  \hline
  n & n+1 \\
  a & a\\
  b & b
\end{array}
\qquad
\begin{array}{c|ccc}
  x \nsplus y & m & a & b\\
  \hline
  n & n+m & b & a\\
  a & a & b & a\\
  b & b & b & a
\end{array}
\]
څرنګه چې \LR{$\nssucc$} يو پر يو ده، \LR{$0$} ئې د قيمتونو په ساحه کښې
نۀ دے، او له~\LR{$0$} پرته هر \LR{$x\in\Domain{L}$} پکې دے، نو
\LR{$!Q_1$}--\LR{$!Q_3$} په~\LR{$\Struct{L}$} کښې رښتوني دي۔ د هر~\LR{$x$} دپاره
\LR{$x\nsplus0=x$}، نو \LR{$!Q_4$} هم رښتونے دے۔ د~\LR{$!Q_5$} دپاره
\LR{$x\nsplus y^\nssucc$} او \LR{$(x\nsplus y)^\nssucc$} پرتله کړئ۔ که \LR{$x$}
او~\LR{$y$} دواړه معياري وي، برابر دي، ځکه بيا \LR{$\nssucc$} او~\LR{$\nsplus$}
له~\LR{$\prime$} او~\LR{$+$} سره سلا دي۔ که \LR{$x$} نامعياري او~\LR{$y$} معياري وي،
\LR{$x\nsplus y^\nssucc=x=x^\nssucc=(x\nsplus y)^\nssucc$}۔ که \LR{$x$} او
\LR{$y$} دواړه نامعياري وي، څلور حالتونه لرو:
\begin{align*}
&  a \nsplus a^\nssucc  = b = b^\nssucc = (a \nsplus a)^\nssucc\\
&  b \nsplus b^\nssucc  = a = a^\nssucc = (b \nsplus b)^\nssucc\\
&  b \nsplus a^\nssucc  = b = b^\nssucc = (b \nsplus a)^\nssucc\\
&  a \nsplus b^\nssucc  = a = a^\nssucc = (a \nsplus b)^\nssucc\\
\intertext{\RL{که \LR{$x$} معياري خو \LR{$y$} نامعياري وي، لرو:}}
&  n \nsplus a^\nssucc  = n \nsplus a = b = b^\nssucc = (n \nsplus a)^\nssucc\\
&  n \nsplus b^\nssucc  = n \nsplus b = a = a^\nssucc = (n \nsplus b)^\nssucc
\end{align*}
نو \LR{$\Sat{L}{!Q_5}$}۔ خو \LR{$a\nsplus0\neq0\nsplus a$}، نو
\LR{$\Sat/{L}{\lforall[x][\lforall[y][\eq[(x+y)][(y+x)]]]}$}۔
\end{ex}
\olpseganchor{OLP-0195-B011}{end}

\olpseganchor{OLP-0195-B012}{start}
\begin{prob}
د~\olref[mod][mar][mdq]{ex:model-L-of-Q} جوړښت~\LR{$\Struct{L}$} په
\LR{$\nstimes$} او~\LR{$\nsless$} سره وغځوئ، څو د~\LR{$\times$} او~\LR{$<$} تفسيرونه
وي۔ وښيئ چې جوړښت مو د~\LR{$\Th{Q}$} پاتې بديهي اصول پوره کوي:
\begin{align*}
& \lforall[x][\eq[(x \times \Obj 0)][\Obj 0]] \tag{\LR{$!Q_6$}}\\ &
  \lforall[x][\lforall[y][\eq[(x \times y')][((x \times y) + x)]]]
  \tag{\LR{$!Q_7$}}\\ & \lforall[x][\lforall[y][(x < y \liff \lexists[z][\eq[(z'+x)][y]])]] \tag{\LR{$!Q_8$}}
\end{align*}
\end{prob}
\olpseganchor{OLP-0195-B012}{end}

\olpseganchor{OLP-0195-B013}{start}
\begin{prob}
د~\olref[mod][mar][mdq]{ex:model-L-of-Q} په~\LR{$\Struct{L}$} کښې
\LR{$a^\nssucc=a$} او \LR{$b^\nssucc=b$} دي۔ ايا د~\LR{$\Th{Q}$} داسې مدل شته چې
\LR{$a^\nssucc=b$} او \LR{$b^\nssucc=a$} وي؟
\end{prob}
\olpseganchor{OLP-0195-B013}{end}

\olpseganchor{OLP-0195-B014}{start}
\begin{explain}
موږ د~\LR{$\Th{Q}$} داسې مدلونه په څرګند ډول جوړ کړل چې نامعياري
غړي ئې د معياري عناصرو «هاخوا» پراتۀ دي۔ په حقيقت کښې،
بديهي اصول همدومره غواړي۔ يو نامعياري غړے~\LR{$x$} د~\LR{${} \nsless 0$} شرط
پوره کولے نۀ شي، ځکه \LR{$\Th{Q} \Proves
\lforall[x][\lnot x<0]$} (وګورئ
\olref[inc][req][min]{lem:less-zero})۔ همدارنګه د هر~\LR{$n$} دپاره
\LR{$\Th{Q} \Proves \lforall[x][(x < \num{n}' \lif
(\eq[x][\num{0}] \lor \eq[x][\num{1}] \lor \dots \lor
\eq[x][\num{n}]))]$} (\olref[inc][req][min]{lem:less-nsucc})، نو د
هېڅ~\LR{$n>0$} دپاره \LR{$a\nsless n$} نۀ شو لرلے۔
\end{explain}
\olpseganchor{OLP-0195-B014}{end}

\olpseganchor{OLP-0196-B004}{start}
\olfileid{mod}{mar}{mpa}
\section{د~\texorpdfstring{\LR{$\Th{PA}$}}{PA} مدلونه}
\olpseganchor{OLP-0196-B004}{end}

\olpseganchor{OLP-0196-B005}{start}
\begin{explain}
د~\LR{$\Th{TA}$} هر نامعياري مدل د~\LR{$\Th{PA}$} هم نامعياري مدل دے۔ موږ
پوهېږو چې د~\LR{$\Th{TA}$}، او په پايله کښې د~\LR{$\Th{PA}$}، نامعياري
مدلونه شته۔ دا هم پوهېږو چې داسې نامعياري مدلونه نامعياري «شمېرې»
لري، يعنې د تعريف ساحې داسې غړي چې له ټولو معياري «شمېرو»
هاخوا دي۔ خو دا څنګه ترتيب شوي دي؟ څومره دي؟ ومو ليدل چې د کمزورې
تيورۍ~\LR{$\Th{Q}$} مدلونه حتا يوازې يوه نامعياري شمېره لرلے شي۔ خو دغه
ساده جوړښتونه د~\LR{$\Th{PA}$} يا~\LR{$\Th{TA}$} مدلونه نۀ دي۔
\olpseganchor{OLP-0196-B005}{end}

\olpseganchor{OLP-0196-B006}{start}
د~\LR{$\Th{PA}$} يا~\LR{$\Th{TA}$} د مدلونو د جوړښت د پوهېدو کلي دا ده چې
وګورو په دغو تيوريو کښې کوم حقيقتونه د اشتقاق وړ دي۔ مثلاً
\LR{$\Th{PA}$} لا له مخکې
\LR{$\lforall[x][\eq/[x][x']]$} او
\LR{$\lforall[x][\lforall[y][\eq[(x+y)][(y+x)]]]$} ثابتوي، نو هغه ساده
جوړښتونه چې دا تړلي فارمولونه دروغوي د~\LR{$\Th{PA}$} د مدلونو په توګه
ردوي۔
\olpseganchor{OLP-0196-B006}{end}

\olpseganchor{OLP-0196-B007}{start}
فرض کړئ~\LR{$\Struct{M}$} د~\LR{$\Th{PA}$} مدل دے۔ که
\LR{$\Th{PA} \Proves !A$} وي، نو \LR{$\Sat{M}{!A}$}۔ بيا د
\LR{$\Assign{\Obj{0}}{M}$} دپاره~\LR{$\nszero$}، د~\LR{$\Assign{\prime}{M}$} دپاره
\LR{$\nssucc$}، د~\LR{$\Assign{+}{M}$} دپاره~\LR{$\nsplus$}، د
\LR{$\Assign{\times}{M}$} دپاره~\LR{$\nstimes$} او د~\LR{$\Assign{<}{M}$} دپاره
\LR{$\nsless$} کاروو۔ هره تړلے فارمول~\LR{$!A$} بيا د~\LR{$\nszero$}، \LR{$\nssucc$}،
\LR{$\nsplus$}، \LR{$\nstimes$} او~\LR{$\nsless$} په هکله کوم شرط بيانوي، او که
\LR{$\Sat{M}{!A}$} وي نو هغه شرط بايد پوره شي۔ مثلاً که
\LR{$\Sat{M}{!Q_1}$}، يعنې
\LR{$\Sat{M}{\lforall[x][\lforall[y][(\eq[x'][y'] \lif \eq[x][y])]]}$}،
نو \LR{$\nssucc$} بايد يو پر يو وي۔
\end{explain}
\olpseganchor{OLP-0196-B007}{end}

\olpseganchor{OLP-0196-B008}{start}
\begin{prop}
په~\LR{$\Struct{M}$} کښې \LR{$\nsless$} يو خطي سخت ترتيب دے، يعنې:
\begin{enumerate}
\item د هېڅ~\LR{$x \in \Domain{M}$} دپاره \LR{$x \nsless x$} نۀ وي۔
\item که \LR{$x \nsless y$} او \LR{$y \nsless z$} وي، نو \LR{$x \nsless z$}۔
\item د هر \LR{$x\neq y$} دپاره يا \LR{$x\nsless y$}، يا \LR{$y\nsless x$}۔
\end{enumerate}
\end{prop}
\olpseganchor{OLP-0196-B008}{end}

\olpseganchor{OLP-0196-B009}{start}
\begin{proof}
\LR{$\Th{PA}$} دا ثابتوي:
\begin{enumerate}
\item \LR{$\lforall[x][\lnot x < x]$}
\item \LR{$\lforall[x][\lforall[y][\lforall[z][((x < y \land y < z) \lif x < z)]]]$}
\item \LR{$\lforall[x][\lforall[y][((x < y \lor y < x) \lor \eq[x][y]))]]$}
\end{enumerate}
\end{proof}
\olpseganchor{OLP-0196-B009}{end}


\olpseganchor{OLP-0196-B010}{start}
\begin{prop}
\ollabel{prop:M-discrete} \LR{$\nszero$} د~\LR{$\Domain{M}$} تر ټولو کوچنے
غړے د~\LR{$\nsless$} په ترتيب کښې دے۔ د هر~\LR{$x$} دپاره
\LR{$x\nsless x^\nssucc$}، او \LR{$x^\nssucc$} تر ټولو کوچنے
\LR{$\nsless$}-لوي غړے دے۔ له صفر څخه پرته د هر~\LR{$x$} دپاره داسې يوازينے~\LR{$y$} شته چې
\LR{$y^\nssucc=x$} وي۔ (موږ~\LR{$y$} په~\LR{$\Struct{M}$} کښې
د~\LR{$x$} «مخکښېنی» بولو او په~\LR{$^\nssucc x$} ئې ښيو۔)
\end{prop}
\olpseganchor{OLP-0196-B010}{end}

\olpseganchor{OLP-0196-B011}{start}
\begin{proof}
تمرين۔
\end{proof}
\olpseganchor{OLP-0196-B011}{end}

\olpseganchor{OLP-0196-B012}{start}
\begin{prob}
په~\LR{$\Lang{L_A}$} کښې داسې تړلي فارمولونه ومومئ چې په~\LR{$\Th{PA}$} کښې
د اشتقاق وړ وي (او ځکه په~\LR{$\Struct{N}$} کښې رښتونې وي) او په
\olref[mod][mar][mpa]{prop:M-discrete} کښې د~\LR{$\nszero$}، \LR{$\nssucc$}
او~\LR{$\nsless$} ځانګړتياوې تضمين کړي۔
\end{prob}
\olpseganchor{OLP-0196-B012}{end}

\olpseganchor{OLP-0196-B013}{start}
\begin{prop}
د~\LR{$\Struct{M}$} ټول معياري غړي د~\LR{$\nsless$} له مخې له ټولو
نامعياري غړي څخه کوچني دي۔
\end{prop}
\olpseganchor{OLP-0196-B013}{end}

\olpseganchor{OLP-0196-B014}{start}
\begin{proof}
\LR{$n$} د معياري غړے~\LR{$\Value{\num{n}}{M}$} د
\LR{$\Struct{M}$} غړي د لنډون په توګه کاروو۔ لا~\LR{$\Th{Q}$} د هر~\LR{$n\in\Nat$} دپاره
\LR{$\lforall[x][(x < \num{n}' \lif (\eq[x][\num{0}] \lor
  \eq[x][\num{1}] \lor \dots \lor \eq[x][\num{n}]))]$} ثابتوي۔
د~\LR{$\nsless \nszero$} شرط پوره کوونکي غړي نشته۔ نو که~\LR{$n$} معياري او~\LR{$x$}
نامعياري وي، \LR{$x\nsless n$} نۀ شو لرلے۔ د تعريف له مخې نامعياري
عنصر د هېڅ~\LR{$n\in\Nat$} دپاره \LR{$\Value{\num{n}}{M}$} نۀ وي، نو
\LR{$x\neq n$} هم دے۔ څرنګه چې \LR{$\nsless$} خطي ترتيب دے، بايد
\LR{$n\nsless x$} وي۔
\end{proof}
\olpseganchor{OLP-0196-B014}{end}

\olpseganchor{OLP-0196-B015}{start}
\begin{prop}
د~\LR{$\Domain{M}$} هر نامعياري غړے~\LR{$x$} د لاندې فرعي سټ يو غړے
دے:
\[
\dots ^{\nssucc\nssucc\nssucc}x \nsless ^{\nssucc\nssucc}x \nsless
^{\nssucc}x \nsless x \nsless x^{\nssucc} \nsless x^{\nssucc\nssucc}
\nsless x^{\nssucc\nssucc\nssucc} \nsless \dots
\]
موږ دا فرعي سټ د~\LR{$x$} \emph{بلاک} بولو او په~\LR{$[x]$} ئې ليکو۔ نه تر
ټولو کوچنے او نه تر ټولو لوے غړے لري۔ دا د هغو
\LR{$y\in\Domain{M}$} د سټ په توګه ځانګړې کېدے شي چې د کوم معياري~\LR{$n$}
دپاره \LR{$x\nsplus n=y$} يا \LR{$y\nsplus n=x$} وي۔
\end{prop}
\olpseganchor{OLP-0196-B015}{end}

\olpseganchor{OLP-0196-B016}{start}
\begin{proof}
داسې سټ~\LR{$[x]$} تل شته، ځکه د~\LR{$\Domain{M}$} هر غړے~\LR{$y$}
يوازينے جانشين~\LR{$y^\nssucc$} لري او که دغه غړے نامعياري وي،
يوازينے مخکښېنی~\LR{$^\nssucc y$} لري۔ د پرله‌پسې غړي \LR{$y$} او~\LR{$y^\nssucc$}
دپاره \LR{$y\nsless y^\nssucc$}، او \LR{$y^\nssucc$} د~\LR{$\nsless$} له مخې
د~\LR{$\Domain{M}$} تر ټولو کوچنے هغه غړے دے چې \LR{$y$} ترې
\LR{$\nsless$} وي۔ څرنګه چې تل
\LR{$^\nssucc y\nsless y$} او \LR{$y\nsless y^\nssucc$}، نو \LR{$[x]$} نه تر ټولو
کوچنے او نه تر ټولو لوے غړے لري۔ که \LR{$y\in[x]$} وي، نو
\LR{$x\in[y]$}، ځکه بيا يا \LR{$y^{\nssucc\dots\nssucc}=x$} يا
\LR{$x^{\nssucc\dots\nssucc}=y$}۔ که
\LR{$y^{\nssucc\dots\nssucc}=x$} وي (له~\LR{$n$} ځله~\LR{$\nssucc$} سره)، نو
\LR{$y\nsplus n=x$} او برعکس، ځکه (که~\LR{$n$} د~\LR{$\prime$} نښو شمېر وي)
\LR{$\Th{PA} \Proves \lforall[x][\eq[x^{\prime\dots\prime}][(x +
    \num{n})]]$}۔
\end{proof}
\olpseganchor{OLP-0196-B016}{end}

\olpseganchor{OLP-0196-B017}{start}
\begin{prop}
که \LR{$[x]\neq[y]$} او \LR{$x\nsless y$} وي، نو د هر \LR{$u\in[x]$} او
\LR{$v\in[y]$} دپاره \LR{$u\nsless v$}۔
\end{prop}
\olpseganchor{OLP-0196-B017}{end}

\olpseganchor{OLP-0196-B018}{start}
\begin{proof}
پام وکړئ چې \LR{$\Th{PA} \Proves
\lforall[x][\lforall[y][(x < y \lif (x' < y \lor x' = y))]]$}۔
نو که \LR{$u\nsless v$} وي او \LR{$[u]\neq[v]$}، د هر~\LR{$n$} دپاره
\LR{$u\nsplus n^\nssucc\nsless v$} هم لرو۔
\olpseganchor{OLP-0196-B018}{end}

\olpseganchor{OLP-0196-B019}{start}
هر \LR{$u\in[x]$} د~\LR{$\nsless y$} شرط پوره کوي: د فرض له مخې \LR{$x\nsless y$}۔ که
\LR{$u\nsless x$} وي، د انتقال له مخې \LR{$u\nsless y$}۔ او که
\LR{$x\nsless u$} خو \LR{$u\in[x]$} وي، نو د کوم~\LR{$n$} دپاره
\LR{$u=x\nsplus n^\nssucc$}؛ نو د هغه حقيقت له مخې چې اوس مو ثابت کړ
\LR{$u\nsless y$}۔
\olpseganchor{OLP-0196-B019}{end}

\olpseganchor{OLP-0196-B020}{start}
اوس فرض کړئ \LR{$v\in[y]$} د~\LR{$\nsless y$} شرط پوره کوي، يعنې د کوم معياري~\LR{$m$}
دپاره \LR{$v\nsplus m^\nssucc=y$}۔ بيا \LR{$v\nsless x$} کېدے نۀ شي، ګنې
\LR{$y=v\nsplus m^\nssucc\nsless x$} به و۔ همدارنګه \LR{$x\neq v$}، ګنې
\LR{$x\nsplus m^\nssucc=v\nsplus m^\nssucc=y$} او \LR{$[x]=[y]$} به وو۔
نو \LR{$x\nsless v$}۔ بيا د هر~\LR{$n$} دپاره
\LR{$x\nsplus n^\nssucc\nsless v$} هم۔ له دې امله که \LR{$x\nsless u$} او
\LR{$u\in[x]$} وي، نو \LR{$u\nsless v$}۔ که \LR{$u\nsless x$} وي، د انتقال له
مخې \LR{$u\nsless v$}۔
\olpseganchor{OLP-0196-B020}{end}

\olpseganchor{OLP-0196-B021}{start}
په پای کښې که \LR{$y\nsless v$} وي، نو \LR{$u\nsless v$}، ځکه لکه چې ومو
ښودل \LR{$u\nsless y$} او \LR{$y\nsless v$}۔
\end{proof}
\olpseganchor{OLP-0196-B021}{end}

\olpseganchor{OLP-0196-B022}{start}
\begin{cor}
که \LR{$[x]\neq[y]$} وي، نو \LR{$[x]\cap[y]=\emptyset$}۔
\end{cor}
\olpseganchor{OLP-0196-B022}{end}

\olpseganchor{OLP-0196-B023}{start}
\begin{proof}
فرض کړئ \LR{$z\in[x]$} او \LR{$x\nsless y$}۔ بيا د هر \LR{$u\in[y]$} دپاره
\LR{$z\nsless u$}۔ که \LR{$z\in[y]$} هم واے، \LR{$z\nsless z$} به مو لرلے۔
د~\LR{$y\nsless x$} حالت هم ورته دے۔
\end{proof}
\olpseganchor{OLP-0196-B023}{end}

\olpseganchor{OLP-0196-B024}{start}
\begin{explain}
دا معنا لري چې بلاکونه پخپله داسې ترتيبېدے شي چې د~\LR{$\nsless$} درناوی
وکړي: \LR{$[x]\nsless[y]$} هله او يوازې هله چې \LR{$x\nsless y$}؛ يا په
هم‌ارزه توګه، چې د هر \LR{$u\in[x]$} او \LR{$v\in[y]$} دپاره
\LR{$u\nsless v$} وي۔ څرګنده ده چې معياري بلاک~\LR{$[0]$} تر ټولو کوچنے بلاک
دے۔ له هېڅ نامعياري بلاک سره تقاطع نۀ لري، او هېڅ دوه نامعياري
بلاکونه هم يو بل نۀ پرې کوي۔ په ځانګړي ډول، د پرله‌پسې جانشينونو يا
مخکښينو په اخيستو بل بلاک ته نۀ شئ رسېدے۔
\end{explain}
\olpseganchor{OLP-0196-B024}{end}

\olpseganchor{OLP-0196-B025}{start}
\begin{prop}
که \LR{$x$} او~\LR{$y$} نامعياري وي، نو \LR{$x\nsless x\nsplus y$} او
\LR{$x\nsplus y\notin[x]$}۔
\end{prop}
\olpseganchor{OLP-0196-B025}{end}

\olpseganchor{OLP-0196-B026}{start}
\begin{proof}
که \LR{$y$} نامعياري وي، نو \LR{$y\neq\nszero$}۔
\LR{$\Th{PA} \Proves \lforall[x][\lforall[y][(y \neq \Obj{0} \lif x < (x+y))]]$}۔
اوس فرض کړئ \LR{$x\nsplus y\in[x]$}۔ څرنګه چې
\LR{$x\nsless x\nsplus y$}، د کوم معياري عدد دپاره
\LR{$x\nsplus n^\nssucc=x\nsplus y$}۔ خو د جمع د اختصار قانون
\LR{$\Th{PA} \Proves \lforall[x][\lforall[y][\lforall[z][(\eq[(x+y)][(x+z)] \lif y = z)]]]$}
دے۔ له دې به د کوم معياري~\LR{$n$} دپاره \LR{$y=n^\nssucc$} پايله شي؛ خو
\LR{$y$} نامعياري فرض شوې ده۔
\end{proof}
\olpseganchor{OLP-0196-B026}{end}

\olpseganchor{OLP-0196-B027}{start}
\begin{prop}
تر ټولو کوچنے نامعياري بلاک نشته۔
\end{prop}
\olpseganchor{OLP-0196-B027}{end}

\olpseganchor{OLP-0196-B028}{start}
\begin{proof}
\LR{$\Th{PA} \Proves \lforall[x][\lexists[y][(\eq[(y+y)][x] \lor
      \eq[(y+y)'][x])]]$}؛ يعنې هر~\LR{$x$} پر~\LR{$2$} وېشل کېږي (ښايي پاتې
شونې~\LR{$1$} ولري)۔ که \LR{$x$} نامعياري وي، \LR{$y$} هم نامعياري دے۔ د مخکينۍ
قضيې له مخې \LR{$y\nsless y\nsplus y$} او \LR{$y\nsplus y\notin[y]$}۔
بيا \LR{$y\nsless(y\nsplus y)^\nssucc$} او
\LR{$(y\nsplus y)^\nssucc\notin[y]$} هم۔ خو
\LR{$x=y\nsplus y$} يا \LR{$x=(y\nsplus y)^\nssucc$}، نو \LR{$y\nsless x$} او
\LR{$y\notin[x]$}۔
\end{proof}
\olpseganchor{OLP-0196-B028}{end}

\olpseganchor{OLP-0196-B029}{start}
\begin{prop}
تر ټولو لوے بلاک نشته۔
\end{prop}
\olpseganchor{OLP-0196-B029}{end}

\olpseganchor{OLP-0196-B030}{start}
\begin{proof}
تمرين۔
\end{proof}
\olpseganchor{OLP-0196-B030}{end}

\olpseganchor{OLP-0196-B031}{start}
\begin{prob}
وښيئ چې د~\LR{$\Th{PA}$} په يوه نامعياري مدل کښې تر ټولو لوے بلاک
نشته۔
\end{prob}
\olpseganchor{OLP-0196-B031}{end}

\olpseganchor{OLP-0196-B032}{start}
\begin{prop}
\ollabel{prop:blocks-dense}
د بلاکونو ترتيب ګڼ دے۔ يعنې که \LR{$x\nsless y$} او \LR{$[x]\neq[y]$} وي،
نو داسې بلاک~\LR{$[z]$} شته چې له دواړو بېل او د هغو تر منځ وي۔
\end{prop}
\olpseganchor{OLP-0196-B032}{end}

\olpseganchor{OLP-0196-B033}{start}
\begin{proof}
فرض کړئ \LR{$x\nsless y$} او د دواړو بلاکونه بېل دي۔ لکه مخکې، \LR{$x\nsplus y$} پر دوه وېشل کېږي
(ښايي پاتې شونې ولري): داسې \LR{$z\in\Domain{M}$} شته چې يا
\LR{$x \nsplus y = z \nsplus z$} يا
\LR{$x \nsplus y = (z \nsplus z)^\nssucc$}۔ عنصر~\LR{$z$} د~\LR{$x$} او~\LR{$y$}
«اوسط» دے۔ پر دوه د وېش شتون په استقرا راځي: د جفت عدد تالي د هماغه خارج قسمت سره پاتې شونے يو لري، او د تاق عدد تالي د راتلونکي خارج قسمت دوه‌برابره ده۔ د دواړو بلاکونو د بېلوالي له کبله د دواړو غړو تر منځ له يو زيات واټن شته۔ که اوسط د کيڼ غړي سره برابر يا ترې کوچنے وای، ښي غړے به د کيڼ له تالي څخه لوے نۀ و۔ او که اوسط د ښي غړي سره برابر يا ترې لوے وای، د دواړو غړو مجموعه به د اوسط له دوه‌برابر څخه کوچنۍ وه۔ نو \LR{$x\nsless z$} او \LR{$z\nsless y$}۔ ځکه چې د~\LR{$x$} او~\LR{$y$} بلاکونه بېل دي، \LR{$y$} د~\LR{$x$} له هر معياري شمېر تالي څخه لوے دے۔ که~\LR{$z$} د~\LR{$x$} په بلاک کښې وای، د کوم معياري~\LR{$k$} دپاره به \LR{$z \mathrel{\leq_{\!M}} x \nsplus \Value{\num{k}}{M}$} و۔ د دوه‌برابره کولو او د يو د زياتولو وروسته به~\LR{$y$} له~\LR{$x$} څخه تر~\LR{$2k+1$} زيات لرې نۀ و، چې تناقض دے۔ همداسې که~\LR{$z$} د~\LR{$y$} په بلاک کښې وای، د کوم معياري~\LR{$k$} دپاره به \LR{$y \mathrel{\leq_{\!M}} z \nsplus \Value{\num{k}}{M}$} و۔ د اوسط معادله به بيا وښيي چې \LR{$z$} له~\LR{$x$} څخه تر~\LR{$k$} زيات لرې نۀ دے، نو~\LR{$y$} به له~\LR{$x$} څخه تر~\LR{$2k$} زيات لرې نۀ وي۔ دواړه ناممکن دي۔
\end{proof}
\olpseganchor{OLP-0196-B033}{end}

\olpseganchor{OLP-0196-B034}{start}
\begin{prob}
د~\olref[mod][mar][mpa]{prop:blocks-dense} مفصل ثبوت وليکئ۔
د~\LR{$z$} د شتون د تضمين دپاره \LR{$\Th{PA}$} بايد کومه تړلے فارمول
اشتقاق کړي؟ ولې \LR{$x\nsless z$} او \LR{$z\nsless y$}، او ولې
\LR{$[x]\neq[z]$} او \LR{$[z]\neq[y]$} دي؟
\end{prob}
\olpseganchor{OLP-0196-B034}{end}

\olpseganchor{OLP-0196-B035}{start}
\begin{explain}
له دې امله، په هر شمېرېدونکي نامعياري مدل کښې نامعياري بلاکونه د
ناطقو په شان ترتيب لري: يو شمېرېدونکے لامتناهي ګڼ خطي ترتيب بې له
سرحدي ټکو جوړوي۔ ښودل کېدے شي چې هر دوه داسې شمېرېدونکے لامتناهي
ترتيبونه آيزومورف دي۔ نو د رښتيني حساب د هر دوه د شمېر وړ
نامعياري مدلونو \LR{$\Struct{M}_1$} او~\LR{$\Struct{M_2}$} هغه راکمونې چې
يوازې~\LR{$<$} او~\LR{$=$} لرونکې ژبې ته کېږي، آيزومورفې دي۔ په رښتيا،
آيزومورفيزم~\LR{$h$} داسې تعريفېدے شي: د~\LR{$\Struct{M_1}$} او
\LR{$\Struct{M_2}$} معياري برخې د معياري مدل~\LR{$\Struct{N}$} سره، او ځکه له
يو بل سره، آيزومورفې دي۔ د نامعياري برخې بلاکونه پخپله د ناطقو په
شان ترتيب شوي او ځکه آيزومورف دي؛ د بلاکونو يو آيزومورفيزم د هر
بلاک دننه د خپل‌سرو عناصرو په سره برابرولو غځېدے شي، او بيا د
\LR{$\Struct{M_1}$} کښې د~\LR{$x$} د جانشين انځور د~\LR{$\Struct{M_2}$} کښې د
\LR{$x$} د انځور جانشين وټاکل شي۔ له دې دا \emph{نۀ} راځي چې
\LR{$\mathfrak{M}_1$} او~\LR{$\mathfrak{M}_2$} د حساب په بشپړه ژبه کښې
آيزومورف دي (آيزومورفيزم تل د يوې ژبه په نسبت وي)، ځکه پر
\LR{$\Domain{M_1}$} او~\LR{$\Domain{M_2}$} د جمع او ضرب د تعريفولو نامتجانسې
لارې شته۔ (دا د واټ د يوې مشهورې قضيې څخه هم راځي چې د يوې بشپړې
تيورۍ د شمېرېدونکو مدلونو شمېر 2 کېدے نۀ شي۔)
\end{explain}
\olpseganchor{OLP-0196-B035}{end}

\olpseganchor{OLP-0197-B004}{start}
\olfileid{mod}{mar}{cmp}
\section{د حساب محاسبه کېدونکي مدلونه}
\olpseganchor{OLP-0197-B004}{end}

\olpseganchor{OLP-0197-B005}{start}
\begin{explain}
معياري مدل~\LR{$\Struct{N}$} دوه ښې ځانګړتياوې لري۔ تعريف ساحه ئې طبيعي
شمېرې~\LR{$\Nat$} دي، يعنې عناصر ئې هماغه ډول څيزونه دي چې د حساب
په ژبه ئې په هکله خبرې کول غواړو، او معياري عددنښه~\LR{$\num{n}$} په
رښتيا~\LR{$n$} په نښه کوي۔ بله ښه ځانګړتيا دا ده چې د~\LR{$\Lang{L_A}$} د
غيرمنطقي نښو ټول تفسيرونه \emph{محاسبه کېدونکي} دي۔ د جانشين، جمع
او ضرب تابعې چې \LR{$\Assign{\prime}{N}$}، \LR{$\Assign{+}{N}$} او
\LR{$\Assign{\times}{N}$} دي، د شمېرو محاسبه کېدونکې تابعې دي (مثلاً د
ټيورينګ ماشينونو له خوا محاسبه کېدونکې، يا په ابتدايي بازګښت
تعريفېدونکې)۔ او په~\LR{$\Struct{N}$} کښې د کوچنيوالي اړيکه، يعنې
\LR{$\Assign{<}{N}$}، پرېکړه کېدونکې ده۔
\olpseganchor{OLP-0197-B005}{end}

\olpseganchor{OLP-0197-B006}{start}
د~\LR{$\Th{Q}$} او~\LR{$\Th{PA}$} په شان د حسابي تيوريو نامعياري مدلونه بايد
نامعياري عناصر ولري۔ نو تعريف ساحې ئې عموماً له~\LR{$\Nat$} سره نور
غړي هم لري۔ خو هر نامتناهي شمېرېدونکے جوړښت پر هر
شمېرېدونکے لامتناهي سټ، د~\LR{$\Nat$} په ګډون، جوړېدے شي۔ نو داسې نامعياري
مدلونه هم شته چې تعريف ساحه ئې~\LR{$\Nat$} وي۔ البته په داسې
مدل~\LR{$\Struct{M}$} کښې لږ تر لږه ځينې شمېرې خپل معمولي نقش نۀ شي
لوبولے، ځکه داسې~\LR{$k$} بايد وي چې د ټولو~\LR{$n \in \Nat$} دپاره
له~\LR{$\Value{\num{n}}{M}$} څخه بېل وي۔
\end{explain}
\olpseganchor{OLP-0197-B006}{end}

\olpseganchor{OLP-0197-B007}{start}
\begin{defn}
د~\LR{$\Lang{L_A}$} يو جوړښت~\LR{$\Struct{M}$} \emph{محاسبه
کېدونکے} دے هله او يوازې هله چې \LR{$\Domain{M}=\Nat$}؛
\LR{$\Assign{\prime}{M}$}، \LR{$\Assign{+}{M}$} او~\LR{$\Assign{\times}{M}$}
محاسبه کېدونکې تابعې؛ او \LR{$\Assign{<}{M}$} پرېکړه کېدونکې اړيکه وي۔
\end{defn}
\olpseganchor{OLP-0197-B007}{end}

\olpseganchor{OLP-0197-B008}{start}
\begin{ex}\ollabel{ex:comp-model-q}
د~\olref[mdq]{ex:model-K-of-Q} جوړښت~\LR{$\Struct{K}$} په ياد راوړئ۔
تعريف ساحه ئې \LR{$\Domain{K}=\Nat\cup\{a\}$} او تفسيرونه ئې دا وو:
\begin{align*}
  \Assign{\Obj{0}}{K} & = 0\\
  \Assign{\prime}{K}(x) & =
  \begin{cases}
    x+1 & \text{\RL{که \LR{$x\in \Nat$} وي}}\\
    a & \text{\RL{که \LR{$x = a$} وي}}
  \end{cases}\\
  \Assign{+}{K}(x, y) & =
  \begin{cases}
    x+y & \text{\RL{که \LR{$x$}، \LR{$y \in\Nat$} وي}}\\
    a & \text{\RL{په نورو حالاتو کښې}}
  \end{cases}\\
  \Assign{\times}{K}(x, y) & =
  \begin{cases}
    xy & \text{\RL{که \LR{$x$}، \LR{$y \in\Nat$} وي}}\\
    0 & \text{\RL{که \LR{$x=0$} يا \LR{$y=0$} وي}}\\
    a & \text{\RL{په نورو حالاتو کښې}}\\
  \end{cases}\\
  \Assign{<}{K} & =
  \Setabs{\tuple{x,y}}{x, y \in \Nat \text{\RL{ او }} x<y} \cup
  \Setabs{\tuple{x,a}}{x \in \Domain{K}}
\end{align*}
خو \LR{$\Domain{K}$} شمېرېدونکے لامتناهي او له~\LR{$\Nat$} سره همشمېره دے۔
مثلاً \LR{$g\colon\Nat\to\Domain{K}$} چې \LR{$g(0)=a$} او د~\LR{$n>0$} دپاره
\LR{$g(n)=n-1$} وي، يوه دوه اړخيزه يو پر يو تابع ده۔ دا د~\LR{$\Th{Q}$} د نوي
مدل~\LR{$\Struct{K'}$} او~\LR{$\Struct{K}$} تر منځ پر آيزومورفيزم بدلولے شو۔
په~\LR{$\Struct{K'}$} کښې بايد د~\LR{$\Lang{L_A}$} نښو ته بېلې تابعې او
اړيکې وټاکو، ځکه د~\LR{$\Nat$} بېل غړي د معياري او نامعياري
شمېرو نقشونه لوبوي۔
\olpseganchor{OLP-0197-B008}{end}

\olpseganchor{OLP-0197-B009}{start}
په ځانګړي ډول، اوس~\LR{$0$} د~\LR{$a$} نقش لوبوي، نه د تر ټولو کوچنۍ معياري
شمېرې۔ تر ټولو کوچنۍ معياري شمېره اوس~\LR{$1$} ده۔ نو
\LR{$\Assign{\Obj{0}}{K'}=1$} ټاکو۔ د جانشين تابع هم اوس بېله ده: يوه
معياري شمېره، يعنې \LR{$n>0$}، ورکړل شي، بيا هم~\LR{$n+1$} راګرځوي۔ خو~\LR{$0$}
اوس د~\LR{$a$} نقش لوبوي، چې خپل جانشين دے۔ نو
\LR{$\Assign{\prime}{K'}(0)=0$}۔ د جمع او ضرب دپاره هم لرو:
\begin{align*}
\Assign{+}{K'}(x, y) & =
  \begin{cases}
    x+y-1 & \text{\RL{که \LR{$x$}، \LR{$y >0$} وي}}\\
    0 & \text{\RL{په نورو حالاتو کښې}}
  \end{cases}\\
  \Assign{\times}{K'}(x, y) & =
  \begin{cases}
    1 & \text{\RL{که \LR{$x = 1$} يا \LR{$y = 1$} وي}}\\
    xy - x - y + 2 & \text{\RL{که \LR{$x$}، \LR{$y > 1$} وي}}\\
    0 & \text{\RL{په نورو حالاتو کښې}}\\
  \end{cases}
\end{align*}
او \LR{$\tuple{x,y}\in\Assign{<}{K'}$} هله او يوازې هله چې \LR{$x<y$}،
\LR{$x>0$} او \LR{$y>0$} وي، يا \LR{$y=0$} وي۔
\olpseganchor{OLP-0197-B009}{end}

\olpseganchor{OLP-0197-B010}{start}
دا ټولې تابعې د طبيعي شمېرو محاسبه کېدونکې تابعې دي او
\LR{$\Assign{<}{K'}$} پر~\LR{$\Nat$} يوه پرېکړه کېدونکې اړيکه ده؛ خو دغه
تابعې پر~\LR{$\Nat$} د جانشين، جمع او ضرب له معمولي تابعو سره يو شان نۀ
دي، او \LR{$\Assign{<}{K'}$} هم پر~\LR{$\Nat$} له معمولي~\LR{$<$} سره يو شان نۀ ده۔
\end{ex}
\olpseganchor{OLP-0197-B010}{end}

\olpseganchor{OLP-0197-B011}{start}
\begin{prob}
داسې جوړښت~\LR{$\Struct{L'}$} ورکړئ چې \LR{$\Domain{L'}=\Nat$} وي او
د~\olref[mod][mar][mdq]{ex:model-L-of-Q} له~\LR{$\Struct{L}$} سره
آيزومورف وي۔
\end{prob}
\olpseganchor{OLP-0197-B011}{end}

\olpseganchor{OLP-0197-B012}{start}
\begin{explain}
\Olref{ex:comp-model-q} ښيي چې \LR{$\Th{Q}$} پر~\LR{$\Nat$} د تعريف ساحې
محاسبه کېدونکي نامعياري مدلونه لري۔ خو لاندې پايله ښيي چې د
\LR{$\Th{PA}$} د مدلونو دپاره دا خبره رښتونې نۀ ده (او ځکه د~\LR{$\Th{TA}$}
د مدلونو دپاره هم نۀ ده)۔
\end{explain}
\olpseganchor{OLP-0197-B012}{end}

\olpseganchor{OLP-0197-B013}{start}
\begin{thm}[د ټېننباوم قضيه]
\LR{$\Struct{N}$} تر آيزومورفيزم پورې د~\LR{$\Th{PA}$} يوازينے محاسبه
کېدونکے مدل دے۔
\end{thm}
\olpseganchor{OLP-0197-B013}{end}

\olpseganchor{OLP-0198-B004}{start}
\olchapter{mod}{int}{د منځګړيتوب قضيه}
\olpseganchor{OLP-0198-B004}{end}

\olpseganchor{OLP-0199-B004}{start}
\olfileid{mod}{int}{int}
\olpseganchor{OLP-0199-B004}{end}

\olpseganchor{OLP-0199-B005}{start}
\olsection{سريزه}
\olpseganchor{OLP-0199-B005}{end}

\olpseganchor{OLP-0199-B006}{start}
د منځګړيتوب قضيه لاندې پايله ده: فرض کړئ
\LR{$\Entails !A \lif !B$}۔ بيا داسې تړلے فارمول~\LR{$!C$} شته چې
\LR{$\Entails !A \lif !C$} او \LR{$\Entails !C \lif !B$} وي۔ سربېره پر دې،
په~\LR{$!C$} کښې هر ثابت سمبول، تابع او پريديکات (له
\LR{$\eq$} پرته) په~\LR{$!A$} او~\LR{$!B$} دواړو کښې راځي۔ تړلے فارمول~\LR{$!C$} د
\LR{$!A$} او~\LR{$!B$} يو \emph{منځګړی} بلل کېږي۔
\olpseganchor{OLP-0199-B006}{end}

\olpseganchor{OLP-0199-B007}{start}
د منځګړيتوب قضيه په خپل ذات کښې په زړه پورې ده، خو اصلي ارزښت ئې
په دې کښې دے چې د يوې تيورۍ دننه د تعريف‌وړتيا په هکله پايلې، او
هغه شرطونه پرې ثابتولے شو چې لاندې ئې د دوو سازګارو تيوريو يوځاے
کول يوه سازګاره تيوري ورکوي۔ لومړۍ پايله د بېت د تعريف‌وړتيا قضيه،
او دويمه د رابنسن د ګډې سازګارۍ قضيه بلل کېږي۔
\olpseganchor{OLP-0199-B007}{end}

\olpseganchor{OLP-0200-B004}{start}
\olfileid{mod}{int}{sep}
\olpseganchor{OLP-0200-B004}{end}

\olpseganchor{OLP-0200-B005}{start}
\olsection{د تړلي فارمولونه بېلول}
\olpseganchor{OLP-0200-B005}{end}

\olpseganchor{OLP-0200-B006}{start}
د منځګړيتوب د قضيې ثبوت ته تر تلو مخکې لږ بنسټيز کار پکار دے۔ د
\LR{$!A$} او~\LR{$!B$} منځګړی داسې تړلے فارمول~\LR{$!C$} ده چې
\LR{$!A \Entails !C$} او \LR{$!C \Entails !B$} وي۔ د نقيض عکس له مخې،
وروستۍ خبره هله او يوازې هله رښتونې ده چې
\LR{$\lnot !B \Entails \lnot !C$} وي۔ هغه تړلے فارمول~\LR{$!C$} چې دا
خاصيت ولري، ويل کېږي چې \LR{$!A$} او~\LR{$\lnot !B$} \emph{بېلوي}۔ نو د
\LR{$!A$} او~\LR{$!B$} د يو منځګړي موندل د داسې تړلے فارمول له موندلو سره
برابر دي چې \LR{$!A$} او~\LR{$\lnot !B$} بېل کړي۔ د معمول په څېر به د يوې
عمومي شوې بڼې څېړل ګټور وي: داسې جمله چې د تړلي فارمولونه دوه
\emph{سټونه} بېل کړي۔
\olpseganchor{OLP-0200-B006}{end}

\olpseganchor{OLP-0200-B007}{start}
\begin{defn}
يوه جمله~\LR{$!C$} د جملو سټونه~\LR{$\Gamma$} او~\LR{$\Delta$} هله او يوازې هله
\emph{بېلوي} چې \LR{$\Gamma \Entails !C$} او
\LR{$\Delta \Entails \lnot !C$} وي۔ که داسې هېڅ تړلے فارمول نۀ وي،
نو~\LR{$\Gamma$} او~\LR{$\Delta$} \emph{نه بېلېدونکي} دي۔
\end{defn}
\olpseganchor{OLP-0200-B007}{end}

\olpseganchor{OLP-0200-B008}{start}
د~\LR{$\Gamma$}، \LR{$\Delta$} او~\LR{$!C$} د مدلونو د ټولګيو تر منځ د شاملېدو
اړيکې لاندې ښودل شوې دي:
\begin{figure}[h]
  \centering
  \begin{tikzpicture}[node distance=2cm, auto, thick]
    \draw [rounded corners] (0,0) -- (6,0) -- (6,3) -- (0,3) --  cycle;
    \draw (1.5,1.5) circle (0.9cm);
    \draw (4.5,1.5) circle (0.9cm);
    \path node at (1.5,1.5) {\Large \LR{$\Gamma$}};
    \path node at (4.5,1.5) {\Large \LR{$\Delta$}};
    \path node at (0.4,2.6) {\Large \LR{$\formula{C}$}};
    \path node at (3.4,0.4) {\Large \LR{$\lnot \formula{C}$}};
    \draw (2.5,0) .. controls (2.5,1.5) and (3.5,1.5) .. (3.5,3);
  \end{tikzpicture}
  \caption{\LR{$!C$}، \LR{$\Gamma$} او~\LR{$\Delta$} بېلوي}
  \ollabel{fig:sep}
\end{figure}
\olpseganchor{OLP-0200-B008}{end}


\olpseganchor{OLP-0200-B009}{start}
\begin{lem}\ollabel{lem:sep1}
فرض کړئ~\LR{$\Lang{L}_0$} هغه ژبه ده چې هر هغه ثابت سمبول،
تابع او پريديکات (له~\LR{$\doteq$} پرته) لري چې په
\LR{$\Gamma$} او~\LR{$\Delta$} \emph{دواړو} کښې راځي؛ او~\LR{$\Lang{L}'_0$} د
هر~\LR{$n \ge 0$} دپاره د نامتناهي ډېرو نويو ثابت سمبولونه~\LR{$\Obj c_n$}
په ورزياتولو ترې ترلاسه شوې وي۔ بيا که~\LR{$\Gamma$} او~\LR{$\Delta$} په
\LR{$\Lang{L}_0$} کښې نه بېلېدونکي وي، په~\LR{$\Lang{L}'_0$} کښې هم نه
بېلېدونکي دي۔
\end{lem}
\olpseganchor{OLP-0200-B009}{end}

\olpseganchor{OLP-0200-B010}{start}
\begin{proof}
په غيرمستقيم ډول پر مخ ځو: د تناقض په موخه فرض کړئ چې~\LR{$\Gamma$} او
\LR{$\Delta$} په~\LR{$\Lang{L}'_0$} کښې بېلېږي۔ د بېلوونکې جملې يو نوی ثابت سمبول~\LR{$c$} واخلئ او په يوه تازه ازاده متغيره نښه ئې بدل کړئ۔ ترلاسه شوې بڼه~\LR{$!C$} وبولئ، نو بېلوونکې جمله~\LR{$\Subst{!C}{c}{x}$} ده او \LR{$\Gamma \Entails \Subst{!C}{c}{x}$} او \LR{$\Delta \Entails \lnot \Subst{!C}{c}{x}$} لرو۔ تر څو چې بل نوی ثابت سمبول پاته وي، همدا ګام بيا کاروو۔ د متناهي شاهد له مخې د~\LR{$\Gamma$} يو
متناهي فرعي سټ~\LR{$\Gamma_0$} او د~\LR{$\Delta$} يو متناهي فرعي
سټ~\LR{$\Delta_0$} شته، داسې چې
\LR{$\Gamma_0 \Entails \Subst{!C}{c}{x}$} او
\LR{$\Delta_0 \Entails \lnot \Subst{!C}{c}{x}$} وي۔ پرېږدئ~\LR{$!G$} په
\LR{$\Gamma_0$} کښې د ټولو فارمولونه عطف، او~\LR{$!H$} په~\LR{$\Delta_0$}
کښې د ټولو فارمولونه عطف وي۔ بيا
\begin{align*}
  !G & \Entails \Subst{!C}{c}{x}, & !H  \Entails \lnot \Subst{!C}{c}{x}.
\end{align*}
ځکه چې~\LR{$c$} په~\LR{$!G$} او~\LR{$!H$} کښې نۀ راځي، په هر مدل کښې هر غړے د~\LR{$c$} تفسير کېدے شي۔ له لومړي استلزام څخه~\LR{$!G \Entails \lforall[x][!C]$} ترلاسه کوو۔ له دويم څخه~\LR{$!H \Entails \lforall[x][\lnot !C]$} ترلاسه کوو، او ځکه چې د مدل ساحه تشه نۀ ده، \LR{$!H \Entails \lnot \lforall[x][!C]$} هم لرو۔ د يکنواختۍ له مخې
\begin{align*}
  \Gamma &\Entails \lforall[x][!C], &
  \Delta & \Entails \lnot \lforall[x][!C],
\end{align*}
نو~\LR{$\lforall[x][!C]$} په~\LR{$\Lang{L}_0$} کښې~\LR{$\Gamma$} او~\LR{$\Delta$}
بېلوي۔
\end{proof}
\olpseganchor{OLP-0200-B010}{end}

\olpseganchor{OLP-0200-B011}{start}
\begin{lem}\ollabel{lem:sep2}
فرض کړئ~\LR{$\Gamma \cup \{ \lexists[x][!S] \}$} او~\LR{$\Delta$} نه
بېلېدونکي دي، او~\LR{$c$} يو نوی ثابت سمبول دے چې په~\LR{$\Gamma$}،
\LR{$\Delta$} يا~\LR{$!S$} کښې نۀ راځي۔ بيا
\LR{$\Gamma \cup \{ \lexists[x][!S], \Subst{!S}{c}{x} \}$} او
\LR{$\Delta$} هم نه بېلېدونکي دي۔
\end{lem}
\olpseganchor{OLP-0200-B011}{end}

\olpseganchor{OLP-0200-B012}{start}
\begin{proof}
د تناقض دپاره فرض کړئ چې~\LR{$!C$}،
\LR{$\Gamma \cup \{\lexists[x][!S], \Subst{!S}{c}{x}\}$} او
\LR{$\Delta$} بېلوي، خو په عين وخت کښې
\LR{$\Gamma \cup \{\lexists[x][!S] \}$} او~\LR{$\Delta$} نه بېلېدونکي دي۔
دوه حالتونه بېلوو:
\begin{enumerate}
\item \LR{$c$} په~\LR{$!C$} کښې نۀ راځي: په دې حالت کښې
  \LR{$\Gamma \cup \{\lexists[x][!S], \lnot!C \}$} د صدق وړ دے (که
  داسې نۀ وے، \LR{$!C$} به~\LR{$\Gamma \cup \{\lexists[x][!S] \}$} او
  \LR{$\Delta$} بېل کړي وو)۔ په داسې مدل کښې د وجودي جملې يو شاهد د نوي ثابت تفسير
  وګرځوئ؛ نو د شاهد د جملې تر ورزياتولو وروسته هم~\LR{$!C$}
  \LR{$\Gamma \cup \{ \lexists[x][!S], \Subst{!S}{c}{x} \}$} او
  \LR{$\Delta$} نۀ بېلوي۔
\item \LR{$c$} په بېلوونکې جمله کښې راځي۔ هغه د~\LR{$\Subst{!C}{c}{x}$} په بڼه ليکو،
  چې~\LR{$!C$} د~\LR{$c$} پر ځاے له يو نوي ازاد متغير څخه جوړ شوے دے۔ بيا لرو چې
  \[
  \Gamma \cup \{ \lexists[x][!S], \Subst{!S}{c}{x}\} \Entails \Subst{!C}{c}{x},
  \]
  نو د استنتاج د قضيې او تعميم له مخې
  \LR{$\Gamma, \lexists[x][!S] \Entails \lforall[x][(!S \lif !C)]$}،
  او په پای کښې
  \LR{$\Gamma \cup \{ \lexists[x][!S] \} \Entails \lexists[x][!C]$}۔
  له بلې خوا، \LR{$\Delta \Entails \lnot \Subst{!C}{c}{x}$} او ځکه د
  تعميم له مخې~\LR{$\Delta \Entails \lnot \lexists[x][!C]$}۔ نو
  \LR{$\Gamma \cup \{\lexists[x][!S] \}$} او~\LR{$\Delta$} د بېلېدو وړ دي؛
  دا تناقض دے۔
\end{enumerate}
\end{proof}
\olpseganchor{OLP-0200-B012}{end}

\olpseganchor{OLP-0201-B004}{start}
\olfileid{mod}{int}{prf}
\olpseganchor{OLP-0201-B004}{end}

\olpseganchor{OLP-0201-B005}{start}
\olsection{د کرېګ د منځګړيتوب قضيه}
\olpseganchor{OLP-0201-B005}{end}

\olpseganchor{OLP-0201-B006}{start}
\begin{thm}[د کرېګ د منځګړيتوب قضيه]
\ollabel{thm:interpol}
که~\LR{$\Entails !A \lif !B$} وي، نو داسې تړلے فارمول~\LR{$!C$} شته چې
\LR{$\Entails !A \lif !C$} او \LR{$\Entails !C \lif !B$} وي، او په~\LR{$!C$}
کښې هر ثابت سمبول، تابع او پريديکات (له~\LR{$\eq$} پرته)
په~\LR{$!A$} او~\LR{$!B$} دواړو کښې راځي۔ تړلے فارمول~\LR{$!C$} د~\LR{$!A$} او~\LR{$!B$}
يو \emph{منځګړی} بلل کېږي۔
\end{thm}
\olpseganchor{OLP-0201-B006}{end}

\olpseganchor{OLP-0201-B007}{start}
\begin{proof}
فرض کړئ~\LR{$\Lang{L}_1$} د~\LR{$!A$} ژبه او~\LR{$\Lang{L}_2$} د~\LR{$!B$} ژبه ده۔
\LR{$\Lang{L}_0 = \Lang{L}_1 \cap \Lang{L}_2$} پرېږدئ۔ د هر
\LR{$i \in \{0, 1, 2 \}$} دپاره~\LR{$\Lang{L}'_i$} له~\LR{$\Lang{L}_i$} څخه د
نامتناهي ډېرو نويو ثابت سمبولونه~\LR{$\Obj c_0, \Obj c_1, \Obj c_2,
\dots$} په ورزياتولو ترلاسه شوې ژبه وي۔
\olpseganchor{OLP-0201-B007}{end}

\olpseganchor{OLP-0201-B008}{start}
که~\LR{$!A$} د صدق وړ نۀ وي، \LR{$\lexists[x][\eq/[x][x]]$} يو منځګړی دے۔
که~\LR{$\lnot !B$} د صدق وړ نۀ وي (او له همدې امله~\LR{$!B$} معتبره وي)،
\LR{$\lexists[x][\eq[x][x]]$} يو منځګړی دے۔ نو دا هم فرض کولے شو چې
\LR{$!A$} او~\LR{$\lnot !B$} دواړه د صدق وړ دي۔
\olpseganchor{OLP-0201-B008}{end}

\olpseganchor{OLP-0201-B009}{start}
د منځګړيتوب د قضيې د نقيض عکس د ثابتولو دپاره فرض کړئ چې د~\LR{$!A$}
او~\LR{$!B$} هېڅ منځګړی نشته۔ په بله وينا، فرض کړئ چې~\LR{$\{!A\}$} او
\LR{$\{\lnot !B\}$} په~\LR{$\Lang{L}_0$} کښې نه بېلېدونکي دي۔
\olpseganchor{OLP-0201-B009}{end}

\olpseganchor{OLP-0201-B010}{start}
موخه مو دا ده چې جوړه~\LR{$(\{ !A \}, \{\lnot!B\})$} يوې اعظمي نه
بېلېدونکې جوړې~\LR{$(\Gamma^*, \Delta^*)$} ته وغځوو۔ پرېږدئ~\LR{$!A_0$}،
\LR{$!A_1$}، \LR{$!A_2$}، \dots، د~\LR{$\Lang{L}'_1$} تړلي فارمولونه، او~\LR{$!B_0$}،
\LR{$!B_1$}، \LR{$!B_2$}، \dots، د~\LR{$\Lang{L}'_2$} تړلي فارمولونه په لړ کښې
راوړي۔
د~\LR{$n \ge 0$} دپاره د تړلي فارمولونه د سټونو دوه زياتېدونکې لړۍ
\LR{$(\Gamma_n, \Delta_n)$} داسې تعريفوو۔
\LR{$\Gamma_0 = \{ !A\}$} او~\LR{$\Delta_0 = \{\lnot !B \}$} کښېږدئ۔ که
\LR{$(\Gamma_n, \Delta_n)$} لا تعريف شوي وي، \LR{$\Gamma_{n+1}$} او
\LR{$\Delta_{n+1}$} له مخکنيو سټونو پيلوئ او لاندې جملې ورزياتوئ۔ د شاهد د ټاکلو پر وخت په پراخه ګډه ژبه کښې تل نوى ثابت موندل کېږي، ځکه هر پړاو يوازې متناهي جملې لري۔
\begin{enumerate}
\item که~\LR{$\Gamma_n \cup \{!A_n \}$} او~\LR{$\Delta_n$} په
  \LR{$\Lang{L}'_0$} کښې نه بېلېدونکي وي، \LR{$!A_n$} په~\LR{$\Gamma_{n+1}$} کښې
  واچوئ۔ سربېره پر دې، که~\LR{$!A_n$} يو وجودي فارمول
  \LR{$\lexists[x][!S]$} وي، داسې نوی ثابت سمبول~\LR{$c$} وټاکئ چې په
  \LR{$\Gamma_n$}، \LR{$\Delta_n$}، \LR{$!A_n$} يا~\LR{$!B_n$} کښې نۀ راځي، او
  \LR{$\Subst{!S}{c}{x}$} په~\LR{$\Gamma_{n+1}$} کښې واچوئ۔
\item که~\LR{$\Gamma_{n+1}$} او~\LR{$\Delta_n \cup \{!B_n \}$} په
  \LR{$\Lang{L}'_0$} کښې نه بېلېدونکي وي، \LR{$!B_n$} په~\LR{$\Delta_{n+1}$} کښې
  واچوئ۔ سربېره پر دې، که~\LR{$!B_n$} يو وجودي فارمول
  \LR{$\lexists[x][!S]$} وي، داسې نوی ثابت سمبول~\LR{$c$} وټاکئ چې په
  \LR{$\Gamma_{n+1}$}، \LR{$\Delta_n$}، \LR{$!A_n$} يا~\LR{$!B_n$} کښې نۀ راځي، او
  \LR{$\Subst{!S}{c}{x}$} په~\LR{$\Delta_{n+1}$} کښې واچوئ۔
\end{enumerate}
په پای کښې داسې تعريف کړئ:
\begin{align*}
  \Gamma^* & = \bigcup_{n\ge 0} \Gamma_n, &
  \Delta^* & = \bigcup_{n\ge 0} \Delta_n.
\end{align*}
اوس پر~\LR{$n$} په هممهاله استقرا ثابتولے شو چې:
\begin{enumerate}
\item\ollabel{part-a} \LR{$\Gamma_n$} او~\LR{$\Delta_n$} په~\LR{$\Lang{L}'_0$}
  کښې نه بېلېدونکي دي؛
\item\ollabel{part-b} \LR{$\Gamma_{n+1}$} او~\LR{$\Delta_n$} په
  \LR{$\Lang{L}'_0$} کښې نه بېلېدونکي دي۔
\end{enumerate}
د~\olref{part-a} بنسټ له~\olref[sep]{lem:sep1} څخه ترلاسه کېږي۔ د
\olref{part-b} دپاره بايد درې حالتونه بېل کړو:
\begin{enumerate}
\item که~\LR{$\Gamma_0 \cup \{!A_0 \}$} او~\LR{$\Delta_0$} بېلېدونکي وي،
  نو~\LR{$\Gamma_1 = \Gamma_0$} او~\olref{part-b} يوازې هماغه
  \olref{part-a} ده؛
\item که~\LR{$\Gamma_1 = \Gamma_0 \cup\{ !A_0\}$} وي، نو~\LR{$\Gamma_1$}
  او~\LR{$\Delta_0$} د جوړونې له مخې نه بېلېدونکي دي؛
\item هغه حالت پاتې کېږي چې~\LR{$!A_0$} وجودي وي، نو
  \LR{$\Gamma_1 = \Gamma_0 \cup \{ \lexists[x][!S],
  \Subst{!S}{c}{x} \}$}۔ د جوړونې له مخې
  \LR{$\Gamma_0 \cup \{ \lexists[x][!S]\}$} او~\LR{$\Delta_0$} نه
  بېلېدونکي دي، نو د~\olref[sep]{lem:sep2} له مخې
  \LR{$\Gamma_0 \cup \{ \lexists[x][!S], \Subst{!S}{c}{x} \}$} او
  \LR{$\Delta_0$} هم نه بېلېدونکي دي۔
\end{enumerate}
په دې سره د پورته~\olref{part-a} او~\olref{part-b} د استقرا بنسټ
بشپړېږي۔ همدا درې حالتونه په هر پړاو کښې کاروو: يا جمله نۀ ورزياتېږي، يا يوازې جمله ورزياتېږي، يا جمله له خپل وجودي شاهد سره ورزياتېږي۔ د~\olref{part-a} د بل پړاو لپاره له~\olref{part-b} څخه پيلوئ او
د دويم سټ د ورزياتونې درې حالتونه کاروئ۔ په وجودي حالت کښې د بېلېدو اړيکه په دواړو لوريو يو شان ده: د بېلوونکې جملې نفي د اړخونو ځاے بدلوي۔ له همدې امله \olref[sep]{lem:sep2} د دويم سټ شاهد ته هم پلي کېږي۔
د~\olref{part-b} د بل پړاو لپاره له همدې ثابت شوي~\olref{part-a} څخه پيلوئ او د لومړي سټ د ورزياتونې درې حالتونه کاروئ۔ په وجودي حالت کښې بيا \olref[sep]{lem:sep2} پلي کېږي۔
په دې سره پر \olref{part-a} او~\olref{part-b} استقرا پاے ته رسېږي۔
\olpseganchor{OLP-0201-B010}{end}

\olpseganchor{OLP-0201-B011}{start}
له دې څخه پايله کېږي چې~\LR{$\Gamma^*$} او~\LR{$\Delta^*$} نه بېلېدونکي دي؛
که داسې نۀ وے، د متناهي شاهد له مخې به داسې~\LR{$n \ge 0$} وای چې
\LR{$\Gamma_n$} او~\LR{$\Delta_n$} بېلېدونکي وي، د~\olref{part-a} خلاف۔ په
ځانګړي ډول، \LR{$\Gamma^*$} او~\LR{$\Delta^*$} سازګار دي: ځکه که لومړنے يا
دويم سټ ناسازګار وے، په ترتيب سره به
\LR{$\lexists[x][\eq/[x][x]]$} يا~\LR{$\lforall[x][\eq[x][x]]$} بېل کړي وو۔
\olpseganchor{OLP-0201-B011}{end}

\olpseganchor{OLP-0201-B012}{start}
اوس ښيو چې~\LR{$\Gamma^*$} په~\LR{$\Lang{L}'_1$} کښې او په همدې ډول
\LR{$\Delta^*$} په~\LR{$\Lang{L}'_2$} کښې اعظمي سازګار دے۔ د لومړي سټ دپاره
فرض کړئ چې د کوم~\LR{$n \ge 0$} دپاره \LR{$!A_n \notin \Gamma^*$} او
\LR{$\lnot !A_n \notin \Gamma^*$} وي۔ که~\LR{$!A_n \notin \Gamma^*$} وي،
نو~\LR{$\Gamma_n \cup \{!A_n \}$} له~\LR{$\Delta_n$} څخه بېلېدونکے دے، او
ځکه داسې~\LR{$!C \in \Lang{L}'_0$} شته چې دواړه لاندې خبرې سمې وي:
\begin{align*}
  \Gamma^* & \Entails !A_n \lif !C, &
  \Delta^* & \Entails \lnot !C.
\end{align*}
همدا ډول، که~\LR{$\lnot !A_n \notin \Gamma^*$} وي، داسې
\LR{$!C' \in \Lang{L}'_0$} شته چې دواړه لاندې خبرې سمې وي:
\begin{align*}
  \Gamma^* & \Entails \lnot !A_n \lif !C', &
  \Delta^* & \Entails \lnot !C'.
\end{align*}
د بياني منطق له مخې~\LR{$\Gamma^* \Entails !C \lor !C'$} او
\LR{$\Delta^* \Entails \lnot (!C \lor !C')$}، نو~\LR{$!C \lor !C'$}،
\LR{$\Gamma^*$} او~\LR{$\Delta^*$} بېلوي۔ ورته استدلال ثابتوي چې~\LR{$\Delta^*$}
هم اعظمي دے۔
\olpseganchor{OLP-0201-B012}{end}

\olpseganchor{OLP-0201-B013}{start}
په پای کښې ښيو چې~\LR{$\Gamma^* \cap \Delta^*$} په~\LR{$\Lang{L}'_0$} کښې
اعظمي سازګار دے۔ دا په څرګنده سازګار دے، ځکه د دوو سازګارو سټونو
تقاطع ده۔ د اعظميت د ښودلو دپاره~\LR{$!S \in \Lang{L}'_0$} ونيسئ۔ اوس
\LR{$\Gamma^*$} په~\LR{$\Lang{L}'_1 \supseteq \Lang{L}'_0$} کښې اعظمي دے،
او په همدې ډول~\LR{$\Delta^*$} په
\LR{$\Lang{L}'_2 \supseteq \Lang{L}'_0$} کښې اعظمي دے۔ نو يا
\LR{$!S \in \Gamma^*$} يا~\LR{$\lnot !S \in \Gamma^*$}، او يا
\LR{$!S \in \Delta^*$} يا~\LR{$\lnot !S \in \Delta^*$} وي۔ که
\LR{$!S \in \Gamma^*$} او~\LR{$\lnot !S \in \Delta^*$} وي، نو~\LR{$!S$} به
\LR{$\Gamma^*$} او~\LR{$\Delta^*$} بېل کړي؛ او که~\LR{$\lnot !S \in \Gamma^*$}
او~\LR{$!S \in \Delta^*$} وي، نو~\LR{$\lnot !S$} به~\LR{$\Gamma^*$} او
\LR{$\Delta^*$} بېل کړي۔ له دې امله، يا
\LR{$!S \in \Gamma^* \cap \Delta^*$} يا
\LR{$\lnot !S \in \Gamma^* \cap \Delta^*$} وي، او
\LR{$\Gamma^* \cap \Delta^*$} اعظمي دے۔
\olpseganchor{OLP-0201-B013}{end}

\olpseganchor{OLP-0201-B014}{start}
څرنګه چې~\LR{$\Gamma^*$} اعظمي سازګار دے، داسې يو مدل~\LR{$\Struct{M}'_1$}
لري چې د تعريف ساحه~\LR{$\Domain{M'_1}$} ئې په هغو او يوازې هغو عناصرو
مشتمله ده چې ثابت سمبولونه تفسيروي، يعنې~\LR{$\Assign{c}{M'_1}$}؛ کټ مټ
لکه د بشپړتيا د قضيې په ثبوت کښې
(\olref[fol][com][cth]{thm:completeness})۔ د يو تړلي ترم د برابرۍ وجودي جمله په اعظمي سازګار سټ کښې وي، ځکه جمله معتبره ده۔ د شمېرنې په خپل وار کښې دا جمله ورزياتېږي: که يې ورزياتونه بېلوونکې وای، هماغه بېلوونکې جمله به وروستي سټونه بېل کړي وو۔ ټاکل شوے شاهد يو ګډ ثابت دے، نو هر تړلے ترم د يوه ګډ ثابت سره برابرېږي۔ په همدې ډول، \LR{$\Delta^*$}
داسې يو مدل~\LR{$\Struct{M}'_2$} لري چې د تعريف ساحه~\LR{$\Domain{M'_2}$} ئې د
ثابت سمبولونه په تفسيرونو~\LR{$\Assign{c}{M'_2}$} جوړه ده۔
\olpseganchor{OLP-0201-B014}{end}

\olpseganchor{OLP-0201-B015}{start}
\LR{$\Struct{M_1}$} له~\LR{$\Struct{M'_1}$} څخه داسې ترلاسه کړئ چې په
\LR{$\Lang{L}'_1 \setminus \Lang{L}'_0$} کښې د ثابت سمبولونه،
تابعې او پريديکاتونه تفسيرونه ترې وغورځوئ؛ او
\LR{$\Struct{M_2}$} هم په همدې ډول ترلاسه کړئ۔ بيا
\LR{$h \colon M_1 \to M_2$} چې د
\LR{$h(\Assign{c}{M'_1}) = \Assign{c}{M'_2}$} له مخې تعريف شوے، په
\LR{$\Lang{L}'_0$} کښې يو آيزومورفيزم دے، ځکه لکه چې ومو ښودل
\LR{$\Gamma^* \cap \Delta^*$} په~\LR{$\Lang{L}'_0$} کښې اعظمي سازګار دے۔
دا ځکه پايله کېږي چې هره~\LR{$\Lang{L}'_0$}-تړلے فارمول يا په
\LR{$\Gamma^*$} او~\LR{$\Delta^*$} دواړو کښې ده، يا په هېڅ يوه کښې نۀ ده:
نو~\LR{$\Assign{c}{M'_1} \in \Assign{P}{M'_1}$} هله او يوازې هله چې
\LR{$\Atom{P}{c} \in \Gamma^*$}، هله او يوازې هله چې
\LR{$\Atom{P}{c} \in \Delta^*$}، او هله او يوازې هله چې
\LR{$\Assign{c}{M'_2} \in \Assign{P}{M'_2}$} وي۔ د دوو ګډو ثابتونو برابري په دواړو مدلونو کښې يو شان ده، نو نقشه ښه تعريف شوې او يو پر يو ده۔ هر غړے د يو ګډ ثابت تفسير دے، نو نقشه پر سر ده۔ د هر تابع ارزښت هم د يو ګډ ثابت په وسيله نومولے شو؛ د دغه ارزښت د برابرۍ جمله په دواړو مدلونو کښې يو شان صدق لري۔ د څو ځاييزو اړيکو لپاره هماغه اټومي دليل پر ټوله جوړه پلي کوو۔
\olpseganchor{OLP-0201-B015}{end}

\olpseganchor{OLP-0201-B016}{start}
اوس به د ژبه~\LR{$\Lang{L}'_1 \cup \Lang{L}'_2$} يو
مدل~\LR{$\Struct{M}$} داسې تعريف کړو:
\begin{enumerate}
\item د تعريف ساحه~\LR{$\Domain{M}$} هماغه~\LR{$\Domain{M_2}$} دے، يعنې د ټولو
  \LR{$\Assign{c}{M'_2}$} عناصرو سټ؛
\item که يو پريديکات~\LR{$P$} په
  \LR{$\Lang{L}'_2 \setminus \Lang{L}'_1$} کښې وي، نو
  \LR{$\Assign{P}{M} = \Assign{P}{M'_2}$}؛
\item که يو پريديکات~\LR{$P$} په
  \LR{$\Lang{L}'_1\setminus \Lang{L}'_2$} کښې وي، نو
  \LR{$\Assign{P}{M} = h(\Assign{P}{M'_1})$}؛ يعنې
  \LR{$\tuple{\Assign{c_1}{M'_2}, \dots, \Assign{c_n}{M'_2}} \in
  \Assign{P}{M}$} هله او يوازې هله چې
  \LR{$\tuple{\Assign{c_1}{M'_1}, \dots,
  \Assign{c_n}{M'_1}} \in \Assign{P}{M'_1}$} وي۔
\item که يو پريديکات~\LR{$P$} په~\LR{$\Lang{L}'_0$} کښې وي، نو
  \LR{$\Assign{P}{M} = \Assign{P}{M'_2} = h(\Assign{P}{M'_1})$}۔
\item د~\LR{$\Lang{L}'_1 \cup \Lang{L}'_2$} تابعې، د
  ثابت سمبولونه په ګډون، په ورته ډول اداره کېږي۔ د دويم مدل د تابعو تفسير اخلو؛ د هغو تابعو لپاره چې يوازې په لومړۍ ژبه کښې دي، آرګومېنټونه د معکوسې نقشې په وسيله لومړي مدل ته وړو او د تابع حاصل بېرته د~\LR{$h$} په وسيله راوړو۔ په ګډو تابعو دواړه طريقې د آيزومورفيزم له امله يو تفسير ورکوي۔
\end{enumerate}
\olpseganchor{OLP-0201-B016}{end}

\olpseganchor{OLP-0201-B017}{start}
په پای کښې پر فارمولونه په استقرا ښودل کېږي چې~\LR{$\Struct{M}$} د
\LR{$\Lang{L}'_1$} پر ټولو فارمولونه له~\LR{$\Struct{M'_1}$} سره، او د
\LR{$\Lang{L}'_2$} پر ټولو فارمولونه له~\LR{$\Struct{M'_2}$} سره موافق دے۔
په ځانګړي ډول، \LR{$\Struct{M} \Entails \Gamma^* \cup \Delta^*$}؛ نو
\LR{$\Struct{M} \Entails !A$} او~\LR{$\Struct{M} \Entails \lnot!B$}، او
\LR{$\not\Entails !A \lif !B$}۔ په دې سره د کرېګ د منځګړيتوب د قضيې
ثبوت پاے ته رسېږي۔
\end{proof}
\olpseganchor{OLP-0201-B017}{end}

\olpseganchor{OLP-0202-B004}{start}
\olfileid{mod}{int}{def}
\olpseganchor{OLP-0202-B004}{end}

\olpseganchor{OLP-0202-B005}{start}
\olsection{د تعريف‌وړتيا قضيه}
\olpseganchor{OLP-0202-B005}{end}

\olpseganchor{OLP-0202-B006}{start}
د منځګړيتوب د قضيې يو مهم استعمال د بېت د تعريف‌وړتيا قضيه ده۔ د
يوې~\LR{$n$}-ځاييزې اړيکې~\LR{$R$} د تعريفولو دپاره داسې فارمول~\LR{$!C$}
ورکولے شو چې~\LR{$n$} ازاد متغيرونه ولري او~\LR{$R$} پکې نۀ راځي۔ دا به
د~\LR{$!C$} له مخې د~\LR{$R$} يو \emph{څرګند} تعريف وي۔ بيا دا هم ويلے شو چې
په داسې ژبه کښې يوه تيوري~\LR{$\Sigma(P)$} چې~\LR{$n$}-ځاييز
پريديکات~\LR{$P$} لري، \LR{$P$} په څرګنده توګه تعريفوي که يو صوري
څرګند تعريف ولري (يا لږ تر لږه ترې لازم راشي)، يعنې
\[
\Sigma(P) \Entails \lforall[x_1][\dots
  \lforall[x_n][(\Atom{P}{x_1,\dots, x_n} \liff !C(x_1, \dots,
    x_n))]].
\]
خو څرګند تعريف د يوې اړيکې د تعريفولو---يعنې په بشپړ ډول د هغې د
ټاکلو---يوازې يوه لار ده۔ يوه تيوري داسې هم کېدے شي چې په هر مدل
کښې د~\LR{$P$} تفسير د پاتې ژبه په تفسير سره ټاکل شوے وي۔ د
تعريف‌وړتيا قضيه وايي چې هر کله يوه تيوري د~\LR{$P$} تفسير په دې ډول
ټاکي---يعنې~\LR{$P$} \emph{په ضمني توګه تعريفوي}---نو هغه په څرګنده
توګه هم تعريفوي۔
\olpseganchor{OLP-0202-B006}{end}

\olpseganchor{OLP-0202-B007}{start}
\begin{defn}
فرض کړئ~\LR{$\Lang{L}$} داسې ژبه ده چې پريديکات~\LR{$P$} نۀ
لري۔ د~\LR{$\Lang{L} \cup \{P\}$} د تړلي فارمولونه يو سټ~\LR{$\Sigma(P)$}،
\LR{$P$} هله او يوازې هله \emph{په څرګنده توګه تعريفوي} چې د
\LR{$\Lang{L}$} داسې فارمول~\LR{$!C(x_1, \dots, x_n)$} وي چې
\[
\Sigma(P) \Entails \lforall[x_1][\dots
  \lforall[x_n][(\Atom{P}{x_1,\dots, x_n} \liff !C(x_1, \dots,
    x_n))]].
\]
\end{defn}
\olpseganchor{OLP-0202-B007}{end}

\olpseganchor{OLP-0202-B008}{start}
\begin{defn}
فرض کړئ~\LR{$\Lang{L}$} داسې ژبه ده چې پريديکاتونه~\LR{$P$} او
\LR{$P'$} نۀ لري۔ د~\LR{$\Lang{L} \cup \{P\}$} د تړلي فارمولونه يو
سټ~\LR{$\Sigma(P)$}، \LR{$P$} هله او يوازې هله \emph{په ضمني توګه تعريفوي}
چې
\[
\Sigma(P) \cup \Sigma(P') \Entails \lforall[x_1][\dots
  \lforall[x_n][(\Atom{P}{x_1,\dots, x_n} \liff \Atom{P'}{x_1,\dots,
      x_n})]],
\]
چې~\LR{$\Sigma(P')$} په~\LR{$\Sigma(P)$} کښې د~\LR{$P$} پر ځاے په يو شان ډول د
\LR{$P'$} له اېښودلو څخه ترلاسه کېږي۔
\end{defn}
\olpseganchor{OLP-0202-B008}{end}

\olpseganchor{OLP-0202-B009}{start}
په بله وينا، د هر مدل~\LR{$\Struct{M}$} او
\LR{$R, R' \subseteq \Domain{M}^n$} دپاره، که دواړه
\LR{$\Expan{M}{R} \Entails \Sigma(P)$} او
\LR{$\Expan{M}{R'} \Entails \Sigma(P')$} وي، نو~\LR{$R=R'$}؛ دلته
\LR{$\Expan{M}{R}$} هغه جوړښت~\LR{$\Struct{M'}$} دے چې~\LR{$\Lang{L}$}
\LR{$\Lang{L} \cup \{P\}$} ته غځوي او~\LR{$\Assign{P}{M'} = R$} کوي، او د
\LR{$\Expan{M}{R'}$} دپاره هم همداسې ده۔
\olpseganchor{OLP-0202-B009}{end}

\olpseganchor{OLP-0202-B010}{start}
\begin{thm}[د بېت د تعريف‌وړتيا قضيه]
د~\LR{$\Lang{L} \cup\{P\}$} د جملو يو سټ~\LR{$\Sigma(P)$}، \LR{$P$}
هله او يوازې هله په ضمني توګه تعريفوي چې~\LR{$\Sigma(P)$}، \LR{$P$} په
څرګنده توګه تعريف کړي۔
\end{thm}
\olpseganchor{OLP-0202-B010}{end}

\olpseganchor{OLP-0202-B011}{start}
\begin{proof}
که~\LR{$\Sigma(P)$}، \LR{$P$} په څرګنده توګه تعريف کړي، نو دواړه
\begin{align*}
  \Sigma(P) & \Entails & \lforall[x_1][\dots \lforall[x_n]
    [(\Atom{P}{x_1,\dots, x_n} \liff !C(x_1,\dots,x_n))]]\\
  \Sigma(P') & \Entails & \lforall[x_1][\dots \lforall[x_n]
    [(\Atom{P'}{x_1,\dots, x_n} \liff !C(x_1,\dots,x_n))]]
\end{align*}
دي، او پايله ترې ترلاسه کېږي۔ د بل لوري دپاره فرض کړئ چې
\LR{$\Sigma(P)$}، \LR{$P$} په ضمني توګه تعريفوي۔ لومړے نوي او له يو بله بېل ثابت سمبولونه~\LR{$c_1$}،
\dots،~\LR{$c_n$} په~\LR{$\Lang{L}$} کښې ورزياتوو۔ بيا
\[
\Sigma(P) \cup \Sigma(P') \Entails
\Atom{P}{c_1, \dots, c_n} \to  \Atom{P'}{c_1, \dots, c_n}.
\]
د متناهي شاهد له مخې داسې متناهي سټونه
\LR{$\Delta_0 \subseteq \Sigma(P)$} او~\LR{$\Delta_1 \subseteq \Sigma(P')$}
شته چې
\[
\Delta_0 \cup \Delta_1 \Entails
\Atom{P}{c_1, \dots, c_n} \to \Atom{P'}{c_1, \dots, c_n}.
\]
پرېږدئ~\LR{$!D(P)$} د ټولو هغو تړلي فارمولونه~\LR{$!A(P)$} عطف وي چې يا
\LR{$!A(P) \in \Delta_0$} يا~\LR{$!A(P') \in \Delta_1$}؛ او~\LR{$!D(P')$} د
ټولو هغو تړلي فارمولونه~\LR{$!A(P')$} عطف وي چې يا
\LR{$!A(P) \in \Delta_0$} يا~\LR{$!A(P') \in \Delta_1$}۔ بيا
\LR{$!D(P) \land !D(P') \Entails \Atom{P}{c_1, \dots, c_n} \to
\Atom{P'}{c_1,\dots,c_n}$}۔
دا داسې بيا ترتيبولے شو چې هر پريديکات د~\LR{$\Entails$} په يوه
اړخ کښې راشي:
\[
!D(P) \land \Atom{P}{c_1, \dots, c_n} \Entails
!D(P') \to \Atom{P'}{c_1, \dots, c_n}.
\]
د کرېګ د منځګړيتوب له قضيې داسې تړلے فارمول
\LR{$!C(c_1,\dots, c_n)$} شته چې~\LR{$P$} يا~\LR{$P'$} نۀ لري او:
\begin{align*}
  !D(P) \land \Atom{P}{c_1, \dots, c_n} & \Entails !C(c_1,\dots, c_n); \\
  !C(c_1,\dots, c_n) & \Entails !D(P') \to \Atom{P'}{c_1, \dots, c_n}.
\end{align*}
له دې دوو استلزامونو له لومړي څخه لرو چې
\LR{$!D(P) \Entails \Atom{P}{c_1,\dots, c_n} \lif
!C(c_1,\dots, c_n)$}۔ او له دويم څخه، ځکه چې د~\LR{$\Lang{L} \cup
\{P\}$} يو مدل~\LR{$\Expan{M}{R} \Entails !A(P)$} هله او يوازې هله چې
ورسره اړوند د~\LR{$\Lang{L} \cup \{P'\}$} مدل
\LR{$\Expan{M}{R} \models !A(P')$} وي، لرو چې
\LR{$!C(c_1,\dots, c_n) \Entails !D(P) \lif
\Atom{P}{c_1,\dots, c_n}$}؛ له دې څخه:
\[
!D(P) \Entails !C(c_1,\dots,c_n) \to \Atom{P}{c_1,\dots, c_n}.
\]
دواړه سره يوځاے کول راکوي چې
\LR{$!D(P) \Entails (\Atom{P}{c_1,\dots, c_n}
\liff !C(c_1, \dots, c_n))$}۔ \LR{$\Sigma(P)$} د عطف هره جمله رانغاړي۔ ځکه چې نوي ثابتونه په~\LR{$\Sigma(P)$} کښې نۀ راځي، په هر مدل کښې د هرې مرتبې ډلې غړي د دغو ثابتونو تفسيرونه ګرځولے شو۔ له همدې امله د تعميم له مخې
\[
\Sigma(P) \Entails
\lforall[x_1][\dots\lforall[x_n][(\Atom{P}{x_1,\dots, x_n} \liff
    !C(x_1,\dots, x_n))]].
\]
\end{proof}
\olpseganchor{OLP-0202-B011}{end}

\olpseganchor{OLP-0203-B004}{start}
\olchapter{mod}{lin}{د ليندستروم قضيه}
\olpseganchor{OLP-0203-B004}{end}

\olpseganchor{OLP-0204-B004}{start}
\olfileid{mod}{lin}{int}
\section{پېژندګلو}
\olpseganchor{OLP-0204-B004}{end}

\olpseganchor{OLP-0204-B005}{start}
په دې څپرکي کښې موخه دا ده چې د لومړۍ درجې منطق په اړه د ليندستروم
ځانګړنه ثابته کړو: د هغو منطقونو په ډله کښې چې ځينې نور شرطونه هم
پوره کوي، دا هغه اعظمي منطق دے چې د متناهي شاهد او د لوېنهايم--سکولم
ښکته قضیې پکښې صدق کوي
(\olref[fol][com][com]{thm:compactness} او
\olref[fol][com][dls]{thm:downward-ls})۔ لومړے بايد د منطقونو هغه عمومي
ټولګی لا په عمومي ډول ځانګړے کړو چې قضيه پرې عملي کېږي۔ بحث به مو
تر \emph{اړيکيزو} ژبو محدود وي، يعنې هغو ژبو ته چې يوازې
پريديکاتونه او فردي ثابتونه لري، خو هېڅ تابعې نۀ لري۔
\olpseganchor{OLP-0204-B005}{end}

\olpseganchor{OLP-0205-B004}{start}
\olfileid{mod}{lin}{alg}
\olsection{انتزاعي منطقونه}
\olpseganchor{OLP-0205-B004}{end}

\olpseganchor{OLP-0205-B005}{start}
\begin{defn}
يو \emph{انتزاعي منطق} جوړه~\LR{$\tuple{L, \models_L}$} ده، چې~\LR{$L$} داسې
تابع ده چې هرې ژبه~\LR{$\Lang{L}$} ته د تړلي فارمولونه يو
سټ~\LR{$L(\Lang{L})$} ورکوي، او~\LR{$\models_L$} د~\LR{$\Lang{L}$} د
جوړښتونه او د~\LR{$L(\Lang{L})$} د غړي تر منځ اړيکه ده۔
په ځانګړي ډول،~\LR{$\tuple{F, \models}$} عادي د لومړۍ درجې منطق دے؛ يعنې
\LR{$F$} هغه تابع ده چې ژبه~\LR{$\Lang{L}$} ته د~\LR{$\Lang{L}$} له
ثابتونو څخه د جوړو شوو د لومړۍ درجې تړلي فارمولونه سټ ورکوي، او
\LR{$\models$} د لومړۍ درجې منطق د پوره کولو اړيکه ده۔
\end{defn}
\olpseganchor{OLP-0205-B005}{end}

\olpseganchor{OLP-0205-B006}{start}
پام وکړئ چې د يوې ورکړې شوې ژبه د جوړښت هماغه مفهوم
کاروو چې د لومړۍ درجې منطق دپاره ئې کاروو؛ خو دا له مخکښې نۀ فرضوو
چې تړلي فارمولونه د~\LR{$\Lang{L}$} له بنسټيزو نښو څخه په معمولي ډول جوړې
شوې دي، او نه دا چې~\LR{$\models_L$} اړيکه د لومړۍ درجې منطق په شان په
بازګشتي ډول تعريف شوې ده۔ نو دا بشپړ عمومي تعريف، د بېلګې په توګه،
هغه حالت هم رانغاړي چې د~\LR{$\tuple{L,\models_L}$} تړلي فارمولونه نامتناهي
اوږده تړونونه يا بېلتونونه ولري، له~\LR{$\lexists$} او~\LR{$\lforall$} پرته نور
کميت ټاکونکي ولري (لکه «نامتناهي ډېر داسې~\LR{$x$} شته چې \dots»)، يا
ښايي د کميت ټاکونکو نامتناهي اوږده مخ‌لړۍ ولري۔ د دې د ټينګار دپاره
چې په~\LR{$L(\Lang{L})$} کښې «تړلي فارمولونه» اړينه نۀ ده د لومړۍ درجې
منطق معمولي تړلي فارمولونه وي، په دې څپرکي کښې پرې~\LR{$!E$}،~\LR{$!F$}،~\dots
متغيرونه چلوو، او~\LR{$!A$}،~\LR{$!B$}،~\dots د لومړۍ درجې معمولي
فارمولونه ته ساتو۔
\olpseganchor{OLP-0205-B006}{end}

\olpseganchor{OLP-0205-B007}{start}
\begin{defn}
\LR{$\Mod(L){!E}$} دې ټولګی~\LR{$\Setabs{\Struct{M}}{\Struct{M}
  \models_L !E}$} وښيي۔ که ژبه په څرګنده ښودل پکار وي،
\LR{$\Mod[L](L){!E}$} ليکو۔ د~\LR{$\Lang{L}$} دوه جوړښتونه~\LR{$\Struct{M}$}
او~\LR{$\Struct{N}$} په~\LR{$\tuple{L, \models_L}$} کښې \emph{ابتدايي هم‌ارز}
دي، چې~\LR{$\Struct{M} \elemequiv[L] \Struct{N}$} ئې ليکو، که د
\LR{$L(\Lang{L})$} هماغه تړلي فارمولونه په دواړو کښې رښتونې وي۔
\end{defn}
\olpseganchor{OLP-0205-B007}{end}

\olpseganchor{OLP-0205-B008}{start}
\begin{defn}
د~ژبه~\LR{$\Lang{L}$} يو انتزاعي منطق~\LR{$\tuple{L,\models_L}$}
\emph{نورمال} دے که لاندې خاصيتونه پوره کړي:
\begin{enumerate}
\item (\emph{\LR{$L$}-زياتېدنه}) د~ژبې~\LR{$\Lang{L}$} او
  \LR{$\Lang{L'}$} دپاره، که~\LR{$\Lang{L} \subseteq \Lang{L'}$} وي، نو
  \LR{$L(\Lang{L}) \subseteq L(\Lang{L'})$}۔
\item (\emph{د غځونې خاصيت}) د هر~\LR{$!E \in L(\Lang{L})$} دپاره د
  \LR{$\Lang{L}$} يو \emph{متناهي} فرعي ژبه~\LR{$\Lang{L'}$} شته چې~\LR{$!E \in L(\Lang{L'})$} او
  اړيکه~\LR{$\Struct{M} \models_L !E$} يوازې د~\LR{$\Lang{L'}$} پر
  راکمونې پورې اړه لري؛ يعنې که~\LR{$\Struct{M}$} او~\LR{$\Struct{N}$} پر
  \LR{$\Lang{L'}$} يوه راکمونه ولري، نو~\LR{$\Struct{M} \models_L !E$} هله او
  يوازې هله چې~\LR{$\Struct{N} \models_L !E$} وي۔
\item (\emph{د آيزومورفيزم خاصيت}) که~\LR{$\Struct{M} \models_L !E$} او
  \LR{$\Struct{M} \simeq \Struct{N}$} وي، نو~\LR{$\Struct{N} \models_L !E$}
  هم وي۔
\item (\emph{د نوم‌اړونې خاصيت}) اړيکه~\LR{$\models_L$} د نوم‌اړونې له
  مخې ساتل کېږي: که ژبه~\LR{$\Lang{L}'$} له~\LR{$\Lang{L}$} څخه داسې
  لاس ته راشي چې هره نښه~\LR{$P$} د هماغې ځاييزې کچې په نښه~\LR{$P'$} او هر
  ثابت~\LR{$c$} په بېل ثابت~\LR{$c'$} بدل شي، نو د هر
  جوړښت~\LR{$\Struct{M}$} او هرې تړلے فارمول~\LR{$!E$} دپاره،
  \LR{$\Struct{M} \models_L !E$} هله او يوازې هله چې
  \LR{$\Struct{M}' \models_L !E'$} وي؛ دلته~\LR{$\Struct{M}'$} د~\LR{$\Struct{M}$}
  اړوند~\LR{$\Lang{L}'$}-جوړښت او~\LR{$!E' \in L(\Lang{L}')$} د~\LR{$!E$}
  نوم‌اړول شوې جمله ده۔
\item (\emph{بولي خاصيت}) انتزاعي منطق~\LR{$\tuple{L,
  \models_L}$} د بولي رابطونو له مخې بند دے، په دې معنا چې د هر
  \LR{$!E \in L(\Lang{L})$} دپاره داسې~\LR{$!F \in L(\Lang{L})$} شته چې
  \LR{$\Struct{M} \models_L !F$} هله او يوازې هله چې
  \LR{$\Struct{M} \not\models_L !E$} وي؛ او د هر~\LR{$!E$} او~\LR{$!F$} دپاره داسې
  \LR{$!G$} شته چې~\LR{$\Mod(L){!G} = \Mod(L){!E} \cap
  \Mod(L){!F}$} وي۔ د اټومي فارمولونه او نورو رابطونو دپاره هم
  همداسې ده۔
\item (\emph{د کميت ټاکونکي خاصيت}) په~\LR{$\Lang{L}$} کښې د هر ثابت~\LR{$c$}
  او هر~\LR{$!E \in L(\Lang{L})$} دپاره داسې~\LR{$!F \in L(\Lang{L'})$} شته
  چې
  \[
  \Mod[L'](L){!F} = \Setabs{\Struct{M}}{\Expan{M}{a} \in
  \Mod[L](L){!E} \text{\RL{ چې }} a \in \Domain{M}},
  \]
  چې دلته~\LR{$\Lang{L'} = \Lang{L} \setminus \{c\}$} او
  \LR{$\Expan{M}{a}$} د~\LR{$\Struct{M}$} هغه~\LR{$\Lang{L}$}-غځونه ده چې~\LR{$a$}
  ثابت~\LR{$c$} ته ورکوي۔
\item (\emph{د نسبي کولو خاصيت}) يوه تړلے فارمول
  \LR{$!E \in L(\Lang{L})$}، يو نوے~\LR{$(n+1)$}-ځاييز محمول~\LR{$R$}، او
  داسې نښې~\LR{$c_1$}، \dots،~\LR{$c_n$} راکړل شي چې په~\LR{$\Lang{L}$} کښې نۀ وي؛
  بيا يوه تړلے فارمول~\LR{$!F \in L(\Lang{L} \cup
  \{R,c_1,\ldots,c_n\})$} شته چې د~\LR{$!E$} \emph{پر
  \LR{$\Atom{R}{x, c_1, \dots c_n}$} نسبي کول} بلل کېږي، او د هر
  جوړښت~\LR{$\Struct{M}$} دپاره:
  \[
  \Expan{M}{X, b_1, \ldots, b_n} \models_L !F \text{\RL{ هله او
    يوازې هله چې }} \Struct{N} \models_L !E,
  \]
  چې~\LR{$\Struct{N}$} د~\LR{$\Struct{M}$} هغه فرعي جوړښت دے چې د تعريف ساحه ئې
  \LR{$\Domain{N} = \Setabs{a\in \Domain{M}}{\tuple{a,b_1,\ldots,b_n}\in X}$}
  دے (\olref[bas][sub]{rem:substructure} وګورئ)، او
  \LR{$\Expan{M}{X, b_1, \ldots, b_n}$} د~\LR{$\Struct{M}$} هغه غځونه ده چې
  \LR{$R$}،~\LR{$c_1$}، \dots،~\LR{$c_n$} په ترتيب سره په~\LR{$X$}،~\LR{$b_1,$}
  \dots،~\LR{$b_n$} تفسيروي (چې~\LR{$X \subseteq \Domain{M}^{n+1}$})۔
\end{enumerate}
\end{defn}
\olpseganchor{OLP-0205-B008}{end}

\olpseganchor{OLP-0205-B009}{start}
\begin{defn}
دوه انتزاعي منطقونه~\LR{$\tuple{L_1, \models_{L_1}}$} او
\LR{$\tuple{L_2, \models_{L_2}}$} راکړل شي۔ وايو چې دويم \emph{لږ تر
لږه د لومړي هومره څرګندوونکی} دے، چې~\LR{$\tuple{L_1, \models_{L_1}}
\leq \tuple{L_2, \models_{L_2}}$} ئې ليکو، که د هرې
ژبه~\LR{$\Lang{L}$} او هرې تړلے فارمول
\LR{$!E \in L_1(\Lang{L})$} دپاره داسې تړلے فارمول
\LR{$!F \in L_2(\Lang{L})$} شته چې~\LR{$\Mod[L](L_1){!E} =
\Mod[L](L_2){!F}$} وي۔ منطقونه~\LR{$\tuple{L_1, \models_{L_1}}$} او
\LR{$\tuple{L_2, \models_{L_2}}$} \emph{هم‌ارز} دي که
\LR{$\tuple{L_1, \models_{L_1}} \leq \tuple{L_2, \models_{L_2}}$} او
\LR{$\tuple{L_2, \models_{L_2}} \leq \tuple{L_1, \models_{L_1}}$} وي۔
\end{defn}
\olpseganchor{OLP-0205-B009}{end}

\olpseganchor{OLP-0205-B010}{start}
\begin{rem}
  د لومړۍ درجې منطق، يعنې انتزاعي منطق~\LR{$\tuple{F, \models}$}، نورمال
  دے۔ په حقيقت کښې، پورته خاصيتونه تر ډېره د لومړۍ درجې منطق دپاره
  سملاسي دي۔ يوازې دومره يادوو چې د غځونې خاصيت له مصداقيت څخه
  راځي، او پر~\LR{$\Atom{R}{x, c_1, \dots, c_n}$} د تړلے فارمول~\LR{$!E$}
  نسبي کول د هرې فرعي فارمول~\LR{$\lforall[x][!F]$} په
  \LR{$\lforall[x][(\Atom{R}{x, c_1, \dots, c_n} \lif \ !F)]$} بدلولو
  سره تر لاسه کېږي۔ سربېره پر دې، که~\LR{$\tuple{L, \models_L}$} نورمال
  وي، نو~\LR{$\tuple{F, \models} \leq \tuple{L, \models_{L}}$}، لکه چې
  دا خبره د لومړۍ درجې پر فارمولونه په استقرا ښودل کېدے شي۔ نو له
  عموميت څخه بې له کوم زيان څخه فرض کولے شو چې د لومړۍ درجې هره
  تړلے فارمول د هر نورمال منطق غړې ده۔
\end{rem}
\olpseganchor{OLP-0205-B010}{end}

\olpseganchor{OLP-0206-B004}{start}
\olfileid{mod}{lin}{lsp}
\olpseganchor{OLP-0206-B004}{end}

\olpseganchor{OLP-0206-B005}{start}
\olsection{د متناهي شاهد او لوېنهايم--سکولم خاصيتونه}
\olpseganchor{OLP-0206-B005}{end}

\olpseganchor{OLP-0206-B006}{start}
اوس د متناهي شاهد او لوېنهايم--سکولم خاصيتونه د انتزاعي منطقونو حالت
ته په څرګند ډول غځوو۔
\olpseganchor{OLP-0206-B006}{end}

\olpseganchor{OLP-0206-B007}{start}
\begin{defn}
يو انتزاعي منطق~\LR{$\tuple{L, \models_L}$} \emph{د متناهي شاهد خاصيت}
لري که د~\LR{$L(\Lang{L})$}-تړلي فارمولونه هر سټ~\LR{$\Gamma$} هغه وخت د صدق
وړ وي چې هر متناهي~\LR{$\Gamma_0 \subseteq \Gamma$} ئې د صدق وړ وي۔
\end{defn}
\olpseganchor{OLP-0206-B007}{end}

\olpseganchor{OLP-0206-B008}{start}
\begin{defn}
\LR{$\tuple{L, \models_L}$} \emph{د لوېنهايم--سکولم ښکته خاصيت} لري که
هر د صدق وړ~\LR{$\Gamma$} يو د شمېر وړ مدل ولري۔
\end{defn}
\olpseganchor{OLP-0206-B008}{end}


\olpseganchor{OLP-0206-B009}{start}
د~\olref[bas][pis]{defn:partialisom} د جزوي آيزومورفيزم مفهوم سوچه
«الجبري» دے؛ يعنې د ژبې د تړلي فارمولونه له يادونې پرته او يوازې د
جوړښتونه د~ژبه~\LR{$\Lang{L}$} د ثابتونو له مخې ورکړل شوے
دے۔ نو دا مفهوم پر انتزاعي منطقونو هم عملي کېږي۔ د لومړۍ درجې منطق
په حالت کښې له~\olref[bas][pis]{thm:p-isom2} څخه پوهېږو چې که دوه
جوړښتونه په جزوي ډول آيزومورف وي، نو ابتدايي هم‌ارز دي۔ هغه
ثبوت انتزاعي منطقونو ته نۀ غځېږي، ځکه پر فارمولونه استقرا د خپل
سري~\LR{$!E \in L(\Lang{L})$} دپاره ښايي موجوده نۀ وي؛ خو که د
لوېنهايم--سکولم خاصيت صدق وکړي، قضيه بيا هم رښتونې ده۔
\olpseganchor{OLP-0206-B009}{end}

\olpseganchor{OLP-0206-B010}{start}
\begin{thm}
\ollabel{thm:abstract-p-isom}
فرض کړئ~\LR{$\tuple{L, \models_L}$} يو نورمال منطق دے چې د
لوېنهايم--سکولم خاصيت لري۔ بيا هر دوه په جزوي ډول آيزومورف
جوړښتونه په~\LR{$\tuple{L, \models_L}$} کښې ابتدايي هم‌ارز دي۔
\end{thm}
\olpseganchor{OLP-0206-B010}{end}

\olpseganchor{OLP-0206-B011}{start}
\begin{proof}
فرض کړئ~\LR{$\Struct{M} \simeq_p \Struct{N}$}، خو د کوم~\LR{$!E$} دپاره
\LR{$\Struct{M} \models_L !E$} او~\LR{$\Struct{N} \not\models_L !E$} هم وي۔ د
آيزومورفيزم د خاصيت له مخې فرض کولے شو چې~\LR{$\Domain{M}$} او
\LR{$\Domain{N}$} سره بېل دي؛ او د غځونې د خاصيت له مخې فرض کولے شو چې
\LR{$!E \in L(\Lang{L})$} د يوې متناهي ژبه~\LR{$\Lang{L}$} دپاره
وي۔ \LR{$\mathcal{I}$} دې د~\LR{$\Struct{M}$} او~\LR{$\Struct{N}$} تر منځ د جزوي
آيزومورفيزمونو يو سټ وي؛ له عموميت څخه بې له زيان څخه دا هم فرض کړئ
چې که~\LR{$p \in \mathcal{I}$} او~\LR{$q \subseteq p$} وي، نو
\LR{$q \in \mathcal{I}$} هم وي۔
\olpseganchor{OLP-0206-B011}{end}

\olpseganchor{OLP-0206-B012}{start}
\LR{$\Domain{M}^{<\omega}$} د~\LR{$\Domain{M}$} د غړي د متناهي لړيو
سټ دے۔ \LR{$S$} دې پر~\LR{$\Domain{M}^{<\omega}$} هغه درې‌ځاييزه اړيکه وي چې
نښلون ښيي؛ يعنې که~\LR{$\mathbf{a}, \mathbf{b},
\mathbf{c} \in \Domain{M}^{<\omega}$} وي، نو
\LR{$S(\mathbf{a}, \mathbf{b}, \mathbf{c})$} هله او يوازې هله صدق کوي
چې~\LR{$\mathbf{c}$} د~\LR{$\mathbf{a}$} او~\LR{$\mathbf{b}$} نښلون وي۔ همدارنګه،
\LR{$T$} دې هغه درې‌ځاييزه اړيکه وي چې د~\LR{$b \in M$} او
\LR{$\mathbf{a}, \mathbf{c} \in \Domain{M}^{<\omega}$} دپاره
\LR{$T(\mathbf{a}, b, \mathbf{c})$} هله او يوازې هله صدق کوي چې
\LR{$\mathbf{a} = a_1, \dots a_n$} او
\LR{$\mathbf{c} = a_1, \dots a_n, b$} وي۔ دوه نوي درې‌ځاييز
پريديکاتونه~\LR{$P$} او~\LR{$Q$} واخلئ، او جوړښت~\LR{$\Struct{M}^*$}
داسې جوړ کړئ چې جهان ئې~\LR{$\Domain{M} \cup \Domain{M}^{<\omega}$} وي،
\LR{$\Struct{M}$} د فرعي جوړښت په توګه ولري، او~\LR{$P$} او~\LR{$Q$} په نښلوونکو
اړيکو~\LR{$S$} او~\LR{$T$} تفسير کړي؛ نو~\LR{$\Struct{M}^*$} د
ژبه~\LR{$\Lang{L} \cup \{ P, Q\}$} جوړښت دے۔ د اړتيا پر وخت د
جهان د دوو برخو بېل نقلونه اخلو، څو دا اتحاد سره بېل کوډ شوے وي۔
\olpseganchor{OLP-0206-B012}{end}

\olpseganchor{OLP-0206-B013}{start}
\LR{$\Domain{N}^{<\omega}$}،~\LR{$S'$}،~\LR{$T'$}،~\LR{$P'$}،~\LR{$Q'$} او
\LR{$\Struct{N}^*$} په ورته ډول تعريف کړئ۔ ځکه چې د فرض له مخې
\LR{$\Struct{M} \simeq_p \Struct{N}$}، د~\LR{$\Domain{M}^{<\omega}$} او
\LR{$\Domain{N}^{<\omega}$} تر منځ داسې اړيکه~\LR{$I$} شته چې
\LR{$I(\mathbf{a}, \mathbf{b})$} هله او يوازې هله صدق کوي چې
\LR{$\mathbf{a}$} او~\LR{$\mathbf{b}$} د~\LR{$\mathcal{I}$} يوه جزوي نقشه وښيي؛ دغه سټ د
\olref[bas][pis]{defn:partialisom} د مخ او شا شرط پوره کړي۔ اوس
\LR{$\Struct{K}$} هغه جوړښت ونيسئ چې د تعريف ساحه ئې د
\LR{$\Struct{M}^*$} او~\LR{$\Struct{N}^*$} د د تعريف ساحې سره بېل کوډ شوے اتحاد
وي،~\LR{$\Struct{M}^*$} او~\LR{$\Struct{N}^*$} د فرعي جوړښتونه په توګه
ولري، او ژبه ئې يو اضافي دوه‌ځاييز پريديکات~\LR{$R$} ولري چې په
اړيکه~\LR{$I$} تفسيرېږي، او داسې پريديکاتونه هم ولري چې د دواړو برخو
د تعريف ساحې~\LR{$\Domain{M}^*$} او~\LR{$\Domain{N}^*$} ټاکي۔
\olpseganchor{OLP-0206-B013}{end}

\olpseganchor{OLP-0206-B014}{start}
\begin{figure}[h]
  \centering
  \begin{tikzpicture}[node distance=2cm, auto, thick, >=stealth']
    \draw [rounded corners] (0,0) -- (8,0) -- (8,4) -- (0,4) --  cycle;
    \draw (2,2) circle (0.5cm);
    \draw (2,2) circle (1.25cm);
    \draw (6,2) circle (0.5cm);
    \draw (6,2) circle (1.25cm);
    \path node at (0.75,3.5) {\large \LR{$\Struct{K}$}};
    \path node at (2,2) {\large \LR{$\Struct{M}$}};
    \path node at (6,2) {\large \LR{$\Struct{N}$}};
    \path node at (3.5,1) {\large \LR{$\Struct{M}^*$}};
    \path node at (7.5,1) {\large \LR{$\Struct{N}^*$}};
    \node (Idom) at (2.8,2) {};
    \node (Irng) at (5.2,2) {};
    \draw[<->, bend left] (Idom) to node {\large \LR{$I$}} (Irng) ;
  \end{tikzpicture}
  \caption{جوړښت~\LR{$\Struct{K}$} او د هغه دننه جزوي
    آيزومورفيزم۔}
\end{figure}
\olpseganchor{OLP-0206-B014}{end}

\olpseganchor{OLP-0206-B015}{start}
مهمه کتنه دا ده چې د~جوړښت~\LR{$\Struct{K}$} په ژبه کښې يوه د
انتزاعي منطق تړلے فارمول~\LR{$!D_1$} شته چې په~\LR{$\Struct{K}$} کښې رښتونې
ده او وايي~\LR{$\Struct{M} \models_L !E$} او
\LR{$\Struct{N} \not\models_L !E$}؛ د دې دپاره د نسبي کولو خاصيت پکار
دے۔ همدارنګه، يوه د لومړۍ درجې تړلے فارمول~\LR{$!D_2$} په~\LR{$\Struct{K}$}
کښې رښتونې ده او وايي چې~\LR{$\Struct{M} \simeq_p \Struct{N}$} د جزوي
آيزومورفيزم~\LR{$I$} له لارې۔ په دې جمله کښې د تشو لړيو جوړه شاملوو، له هرې جوړې د اټومي معلوماتو د ساتنې غوښتنه کوو، او د دواړو خواوو د هر غړي د ورزياتولو دپاره په هماغه اړيکه کښې يوه غځېدلې جوړه غواړو۔ د متناهي ژبې د هر محمول دپاره د لړيو د برابر ځايونو غړي پرتله کوو؛ د نښلون ترڅنګ د ځاے لوستلو محمول هم په پراخه ژبه کښې کاروو۔ د لوېنهايم--سکولم د خاصيت له مخې،~\LR{$!D_1$}
او~\LR{$!D_2$} په ګډه په يوه د شمېر وړ مدل~\LR{$\Struct{K}_0$} کښې
رښتونې دي؛ په دې مدل کښې په جزوي ډول آيزومورف فرعي جوړښتونه
\LR{$\Struct{M}_0$} او~\LR{$\Struct{N}_0$} شته، داسې چې
\LR{$\Struct{M}_0 \models_L !E$} او~\LR{$\Struct{N}_0 \not\models_L !E$}۔ خو
د شمېر وړ او په جزوي ډول آيزومورف جوړښتونه د
\olref[bas][pis]{thm:p-isom1} له مخې په حقيقت کښې آيزومورف دي؛ دا د
نورمالو انتزاعي منطقونو د آيزومورفيزم له خاصيت سره تناقض دے۔
\end{proof}
\olpseganchor{OLP-0206-B015}{end}

\olpseganchor{OLP-0207-B004}{start}
\olfileid{mod}{lin}{prf}
\olpseganchor{OLP-0207-B004}{end}

\olpseganchor{OLP-0207-B005}{start}
\olsection{د ليندستروم قضيه}
\olpseganchor{OLP-0207-B005}{end}

\olpseganchor{OLP-0207-B006}{start}
\begin{lem}
\ollabel{lem:lindstrom}
فرض کړئ~\LR{$!E \in L(\Lang{L})$} او~\LR{$\Lang{L}$} متناهي ده؛ دا هم فرض
کړئ چې داسې~\LR{$n \in \Nat$} شته چې د هر دوو جوړښتونه
\LR{$\Struct{M}$} او~\LR{$\Struct{N}$} دپاره، که~\LR{$\Struct{M} \equiv_n
\Struct{N}$} او~\LR{$\Struct{M} \models_L !E$} وي، نو
\LR{$\Struct{N} \models_L !E$} هم وي۔ بيا~\LR{$!E$} د لومړۍ درجې له يوې
تړلے فارمول سره هم‌ارزه ده؛ يعنې د لومړۍ درجې داسې~\LR{$!D$} شته چې
\LR{$\Mod(L){!E} = \Mod(L){!D}$} وي۔
\end{lem}
\olpseganchor{OLP-0207-B006}{end}

\olpseganchor{OLP-0207-B007}{start}
\begin{proof}
\LR{$n$} داسې ونيسئ چې هر دوه~\LR{$n$}-هم‌ارز جوړښتونه~\LR{$\Struct{M}$} او
\LR{$\Struct{N}$}،~\LR{$!E$} ته د ورکړي ارزښت په اړه سره سلا وي۔
\olref[bas][pis]{prop:qr-finite} راياد کړئ: په يوه متناهي
ژبه کښې، له منطقي هم‌ارزۍ پرته، د لومړۍ درجې يوازې متناهي
ډېرې داسې تړلي فارمولونه شته چې د کميت ټاکونکو رتبه ئې تر~\LR{$n$} زياته
نۀ وي۔ اوس د هر ټاکلي جوړښت~\LR{$\Struct{M}$} دپاره
\LR{$!D_{\Struct{M}}$} د هغو د لومړۍ درجې تړلي فارمولونه~\LR{$!E$} د هر منطقي
هم‌ارزۍ له ټولګي څخه د يوه استازي تړون ونيسئ چې په~\LR{$\Struct{M}$} کښې
رښتونې وي او~\LR{$\QuantRank{!E} \le n$} ولري؛ دا تړون متناهي دے، او
\LR{$\Struct{N} \models !D_{\Struct{M}}$} هله او يوازې هله چې
\LR{$\Struct{N} \equiv_n \Struct{M}$} وي۔ بيا
\LR{$!D = \textstyle\bigvee
\Setabs{!D_{\Struct{M}}}{\Struct{M} \models_L !E}$} کېږدئ؛ له هر
منطقي هم‌ارزۍ ټولګي څخه د يوه استازي په ساتلو دا بېلتون هم متناهي
دے۔
\olpseganchor{OLP-0207-B007}{end}

\olpseganchor{OLP-0207-B008}{start}
پايله~\LR{$\Mod(L){!E} = \Mod(L){!D}$} تر لاسه کېږي۔ په حقيقت کښې، که
\LR{$\Struct{N} \models_L !D$} وي، نو د کوم
\LR{$\Struct{M} \models_L !E$} دپاره~\LR{$\Struct{N} \models
!D_{\Struct{M}}$}؛ نو د لمې د فرض له مخې
\LR{$\Struct{N} \models_L !E$} هم۔ برعکس، که
\LR{$\Struct{N} \models_L !E$} وي، نو~\LR{$!D_\Struct{N}$} په~\LR{$!D$} کښې يوه
بېلېدونکې برخه ده، او دا چې~\LR{$\Struct{N} \models !D_\Struct{N}$}، نو
\LR{$\Struct{N} \models_L !D$} هم۔
\end{proof}
\olpseganchor{OLP-0207-B008}{end}

\olpseganchor{OLP-0207-B009}{start}
\begin{thm}[د ليندستروم قضيه]
  \ollabel{thm:lindstrom} فرض کړئ~\LR{$\tuple{L, \models_L}$} يو نورمال
  انتزاعي منطق دے چې د متناهي شاهد او لوېنهايم--سکولم خاصيتونه لري۔
  بيا~\LR{$\tuple{L, \models_L} \le \tuple{F, \models}$}؛ نو
  \LR{$\tuple{L, \models_L}$} د لومړۍ درجې له منطق سره هم‌ارز دے۔
\end{thm}
\olpseganchor{OLP-0207-B009}{end}

\olpseganchor{OLP-0207-B010}{start}
\begin{proof}
د~\olref{lem:lindstrom} له مخې، دومره ښودل کافي دي چې د هر
\LR{$!E \in L(\Lang{L})$} دپاره، چې~\LR{$\Lang{L}$} متناهي وي، داسې
\LR{$n \in \Nat$} شته چې د هر دوو جوړښتونه~\LR{$\Struct{M}$} او
\LR{$\Struct{N}$} دپاره: که~\LR{$\Struct{M} \equiv_n \Struct{N}$} وي، نو
\LR{$\Struct{M}$} او~\LR{$\Struct{N}$} پر~\LR{$!E$} سره سلا وي۔ بيا~\LR{$!E$} د لومړۍ
درجې له يوې تړلے فارمول سره هم‌ارزه ده، او له دې څخه
\LR{$\tuple{L, \models_L} \le \tuple{F, \models}$} پايله کېږي۔ ځکه چې په
يوه متناهي او سوچه اړيکيزه ژبه کښې کار کوو، د
\olref[bas][pis]{thm:b-n-f} له مخې دا وينا چې
\LR{$\Struct{M} \equiv_n \Struct{N}$} ده، په اړوندې الجبري وينا
\LR{$I_n(\emptyset,\emptyset)$} بدلولے شو۔
\olpseganchor{OLP-0207-B010}{end}

\olpseganchor{OLP-0207-B011}{start}
يو~\LR{$!E$} راکړل شي، او د تناقض دپاره فرض کړئ چې د هر~\LR{$n$} دپاره داسې
جوړښتونه~\LR{$\Struct{M}_n$} او~\LR{$\Struct{N}_n$} شته چې
\LR{$I_n(\emptyset, \emptyset)$}، خو، د بېلګې په توګه،
\LR{$\Struct{M}_n \models_L !E$} او~\LR{$\Struct{N}_n \not\models_L !E$} وي۔
د هرې خوا په متناهي ژبه کښې د ثابتونو تر منځ د برابرۍ او د اټومي جملو يوازې متناهي بڼې شته۔ د دواړو خواوو يوه ګډه نامتناهي فرعي لړۍ واخلو چې دغه بڼې پکې ثابتې وي۔ د بيا شمېرلو پر وخت د هر نوي شاخص دپاره داسې پخوانے شاخص اخلو چې تر هغه کوچنے نۀ وي؛ د لوړې رتبې د مخ او شا شرط ټيټه رتبه هم ورکوي۔ په هره خوا کښې ثابتونو ته هماغه ټاکلي غړي ورکوو، او د نورو غړو د هرې برخې دپاره بېل نقلونه اخلو۔ ځکه چې اټومي بڼې يو شان دي، د اړيکو تفسيرونه پر ګډو ثابتونو سره سمون لري۔ \LR{$\Struct{M}$} د ټولو~\LR{$\Struct{M}_n$}ګانو اتحاد
ونيسئ؛ يعنې هغه يوازينی کوچنی جوړښت چې هر~\LR{$\Struct{M}_n$} د
فرعي جوړښت په توګه لري۔ لکه د~\olref[lsp]{thm:abstract-p-isom} په
ثبوت کښې،~\LR{$\Struct{M}^*$} د~\LR{$\Struct{M}$} هغه غځونه ونيسئ چې د تعريف ساحه
ئې~\LR{$\Domain{M} \cup \Domain{M}^{<\omega}$} وي، او په پراخه شوې
ژبه کښې د نښلون محمولونه~\LR{$P$} او~\LR{$Q$} ولري۔
\olpseganchor{OLP-0207-B011}{end}

\olpseganchor{OLP-0207-B012}{start}
په ورته ډول~\LR{$\Struct{N}_n$}،~\LR{$\Struct{N}$} او~\LR{$\Struct{N}^*$} تعريف
کړئ۔ د ښيې خوا د محمولونو او ثابتونو دپاره بېل نومونه کاروو، څو د دواړو خواوو نومونه په ګډ جوړښت کښې ټکر ونۀ کړي۔ اوس~\LR{$\Struct{K}$} هغه جوړښت ونيسئ چې د تعريف ساحه ئې د
\LR{$\Struct{M}^*$} او~\LR{$\Struct{N}^*$} د د تعريف ساحې تر څنګ طبيعي
عددونه~\LR{$\Nat$} د خپل طبيعي ترتيب~\LR{$\le$} سره رانغاړي۔ ژبه ئې داسې
اضافي محمولونه لري چې د اړوندو د تعريف ساحې~\LR{$\Domain{M}$}،~\LR{$\Domain{N}$}،
\LR{$\Domain{M}^{<\omega}$} او~\LR{$\Domain{N}^{<\omega}$} ښيي، او داسې
محمولونه هم لري چې د~\LR{$\Struct{M}_n$} او~\LR{$\Struct{N}_n$} د تعريف
ساحې په لاندې معنا کوډوي:
\begin{align*}
  \Domain{M_n} & = \Setabs{a \in \Domain{M}}{R(a, n)}; &
  \Domain{N_n} & = \Setabs{a \in \Domain{N}}{S(a,n)}; \\
  \Domain{M}^{<\omega}_n & = \Setabs{a \in \Domain{M}^{<\omega}}{R(a,n)}; &
  \Domain{N}^{<\omega}_n & = \Setabs{a \in \Domain{N}^{<\omega}}{S(a,n)}.
\end{align*}
جوړښت~\LR{$\Struct{K}$} يوه څلورځاييزه اړيکه~\LR{$J$} هم لري، داسې چې
\LR{$J(n,r,\mathbf{a},\mathbf{b})$} هله او يوازې هله صدق کوي چې لړۍ د~\LR{$n$}مې جوړې په ساحو کښې وي او د هماغې جوړې لپاره \LR{$I_r(\mathbf{a},\mathbf{b})$} صدق وکړي۔
\olpseganchor{OLP-0207-B012}{end}

\olpseganchor{OLP-0207-B013}{start}
اوس د~\LR{$\Lang{L}$} په هغه ژبه کښې چې په~\LR{$R$}،~\LR{$S$}،~\LR{$J$} او
نورو نښو پراخه شوې ده، داسې تړلے فارمول~\LR{$!D$} شته چې وايي:~\LR{$\le$}
يو منفصل خطي ترتيب دے چې لومړے خو وروستی غړے نۀ لري؛ د ترتيب
د هر~\LR{$n$} دپاره~\LR{$\Struct{M}_n \models_L !E$} او
\LR{$\Struct{N}_n \not\models_L !E$}؛ او
\LR{$J(n,r,\mathbf{a},\mathbf{b})$} د~\LR{$n$}مې جوړې د~\LR{$r$} پاتې ګامونو اړيکه ده۔ په هر پړاو کښې د تشو لړيو جوړه په هماغه پړاو هومره پاتې ګامونو سره شامله ده۔ د تشې لړۍ اوږدوالے صفر دے؛ د يوه غړي په ورزياتولو د اوږدوالي سملاسي تالي اخلو، پخواني ځايونه خپل غړي ساتي او په نوي ځاے کښې هماغه ورزيات شوے غړے راځي۔ که يوه جوړه په دې اړيکه کښې وي، لړۍ يو شان اوږدوالے لري، غړي ئې د هماغه پړاو په ساحو کښې دي، او د اټومي فورمولونو ارزښت ساتي۔ د هرې مثبتې پاتې رتبې دپاره د دې رتبې مخکينے شاخص شته۔ د کيڼې خوا هر غړے له ښيې خوا د داسې غړي په ملګرتيا ورزياتېدے شي چې غځېدلې لړۍ د هماغه پړاو په ټيټه رتبه کښې بيا په اړيکه کښې وي؛ او د ښيې خوا هر غړے هم همداسې۔ دا د لومړۍ درجې شرطونه دي: د لړيو غځول په نښلوونکو محمولونو بيانوو، او د هر محمول اټومي ساتنه د هغو غړو پر ټولو جوړو ازمويو چې د لړيو په يو شان ځايونو کښې راغلي وي۔ د ځايونو د لوستلو اړيکه هم د لړيو په کوډ کښې شاملوو، څو د ثابتونو، برابرۍ او د متناهي ژبې د هر محمول د ساتنې شرط په يوه متناهي تړون کښې وليکل شي۔ د مخ او شا تړلي فارمولونه د پيل او غځولو شرطونه ټاکي او د اټومي ساتنې تړلي فارمولونه د جوړو جزوي آيزومورفيزم ټاکي۔ د انتزاعي جملې اړوندې برخې د نسبي کولو، د ښيې خوا د نوم‌اړونې، او د کميت ټاکونکي له خاصيتونو څخه تر لاسه کېږي۔ د ساحو تش نه‌والی او د خپل هر ثابت شموليت د لومړۍ درجې په برخې کښې ټاکو۔
\olpseganchor{OLP-0207-B013}{end}

\olpseganchor{OLP-0207-B014}{start}
ژبه په نويو ثابتونو~\LR{$c$}،~\LR{$c_0$}،~\LR{$c_1$}،~\dots پراخه کړئ، او له
\LR{$!D$} سره د لومړۍ درجې هغه جملې يوځاے کړئ چې~\LR{$c_0$} د ترتيب لومړے
غړے ټاکي، هر~\LR{$c_{k+1}$} د~\LR{$c_k$} سملاسي جانشين ټاکي، او د هر معياري
\LR{$k \in \Nat$} دپاره~\LR{$c_k < c$} وايي۔ د دې تيورۍ هر متناهي فرعي سټ
په~\LR{$\Struct{K}$} کښې داسې پوره کېدے شي چې~\LR{$c$} د هغو متناهي ډېرو
يادو شوو جانشينانو څخه پورته يو طبيعي عدد وټاکي۔ د متناهي شاهد د
خاصيت له مخې ټوله تيوري يو مدل~\LR{$\Struct{K}^*$} لري۔ نو
\LR{$n^* = \Assign{c}{K^*}$} د هر معياري~\LR{$k$} دپاره د~\LR{$c_k$} له تفسير څخه
پورته دے، او د کوډ شوي منفصل ترتيب يو نامعياري غړے دے۔ په
ځانګړي ډول،~\LR{$\Struct{K}^*$} داسې فرعي جوړښتونه
\LR{$\Struct{M_{n^*}}$} او~\LR{$\Struct{N_{n^*}}$} لري چې
\LR{$\Struct{M_{n^*}} \models_L !E$} او
\LR{$\Struct{N_{n^*}} \not\models_L !E$}۔ اوس د~\LR{$\Domain{M_{n^*}}$} او
\LR{$\Domain{N_{n^*}}$} د~\LR{$k$}-توکيزو لړيو د جوړو يو سټ~\LR{$\mathcal{I}$}
داسې تعريفولے شو:~\LR{$\tuple{\mathbf{a}, \mathbf{b}} \in \mathcal{I}$}
هله او يوازې هله چې~\LR{$J(n^*,n^*-k,\mathbf{a},\mathbf{b})$} وي؛ دلته~\LR{$k$}
د~\LR{$\mathbf{a}$} او~\LR{$\mathbf{b}$} اوږدوالی دے۔ ځکه~\LR{$n^*$} نامعياري
دے، د هر معياري~\LR{$k$} دپاره~\LR{$n^* - k >0$}؛ دلته منفي کول د سملاسي مخکيني غړي د ټاکلو تکرار دے۔ تشه جوړه د ټاکلي پړاو د پيل له شرطه راځي۔ له هرې جوړې څخه په دواړو لوريو غځونه په هماغه پړاو او له يو کم پاتې شاخص سره راځي، او د اټومي ساتنې له امله هره جوړه جزوي آيزومورفيزم ورکوي۔ نو سټ~\LR{$\mathcal{I}$} د دې
خبرې شاهد دے چې~\LR{$\Struct{M_{n^*}} \simeq_p
\Struct{N_{n^*}}$}۔ خو د~\olref[lsp]{thm:abstract-p-isom} له مخې
\LR{$\Struct{M_{n^*}}$} له~\LR{$\Struct{N_{n^*}}$} سره~\LR{$L$}-هم‌ارز دے، او دا
تناقض دے۔
\end{proof}
\olpseganchor{OLP-0207-B014}{end}

\olpseganchor{OLP-0208-B004}{start}
\olpart{cmp}{محاسبه کېدنه}
\olpseganchor{OLP-0208-B004}{end}

\olpseganchor{OLP-0208-B005}{start}
\begin{editorial}
  دا برخه د محاسبه کېدنې د تيورۍ په اړه د جرېمي اېوېګېډ پر يادښتونو
  ولاړه ده۔ اوس‌مهال يوازې د بازګشتي تابعو څپرکی تمرينونه لري، او
  ټوله برخه لا د هڅونې، بېلګو، جزيياتو او تمرينونو په زياتولو سره
  پراخېدو ته اړتيا لري۔
\end{editorial}
\olpseganchor{OLP-0208-B005}{end}

\olpseganchor{OLP-0209-B004}{start}
\olchapter{cmp}{rec}{بازګشتي تابعې}
\olpseganchor{OLP-0209-B004}{end}

\olpseganchor{OLP-0209-B005}{start}
\begin{editorial}
  دا د بازګشتي تابعو په اړه د جرېمي اېوېګېډ يادښتونه دي چې ريچرډ
  زاک بيا کتلي او پراخ کړي دي۔ دا څپرکی ځينې تمرينونه لري، او په
  خپلواکه توګه هم شاملېدے شي څو د نحو د حسابي کولو د بحث بنسټ برابر
  کړي۔
\end{editorial}
\olpseganchor{OLP-0209-B005}{end}

\olpseganchor{OLP-0210-B004}{start}
\olfileid{cmp}{rec}{int}
\olsection{پېژندګلو}
\olpseganchor{OLP-0210-B004}{end}

\olpseganchor{OLP-0210-B005}{start}
د محاسبه کېدنې د يوې رياضي تيورۍ د جوړولو دپاره، تر هر څه وړاندې د
محاسبه کېدنې يو \emph{مدل} جوړول پکار دي۔ اوس موږ محاسبه کېدنه د هغه
ډول کار په توګه ګڼو چې کمپيوټرونه ئې کوي، او کمپيوټرونه له نښو سره
کار کوي۔ خو د محاسبه کېدنې د تيوريو د پرمختګ په پيل کښې د محاسبې
نمونه‌يي بېلګه \emph{عددي} محاسبه وه۔ رياضيپوهانو تل د عددونو د
تيورۍ له تابعو سره مينه لرلې، يعنې له داسې تابعو \LR{$f\colon \Nat^n \to
\Nat$} سره چې محاسبه کېدے شي۔ نو دا د حيرانتيا خبره نۀ ده چې د
محاسبه کېدنې د تيورۍ په پيل کښې همدا تابعې څېړل کېدې۔ د محاسبه
کېدونکو عددي تابعو تر ټولو پېژندل شوې بېلګې، لکه جمع، ضرب او توان
(د طبيعي عددونو)، يو په زړه پورې خاصيت لري: دا په
\emph{بازګشتي} ډول تعريفېدے شي۔ ځکه نو دا ډېره طبيعي ده چې د
بازګشتي تعريفونو پر بنسټ د \emph{محاسبه کېدونکې تابع} يو عمومي تعريف
جوړ کړو۔ د عددونو د تيورۍ د تابعو د بازګشتي تعريف له ډېرو ممکنو لارو
څخه دلته يوه په ځانګړي ډول ساده نمونه مرکزي کېږي: تش په نامه
\emph{بنسټيز بازګښت}۔
\olpseganchor{OLP-0210-B005}{end}

\olpseganchor{OLP-0210-B006}{start}
له محاسبه کېدونکو تابعو سره سره ښايي محاسبه کېدونکي سټونه او اړيکې
هم راته د پام وړ وي۔ يو سټ هغه وخت محاسبه کېدونکی دے چې دا محاسبه
کولے شو چې يو ورکړل شوے عدد د سټ غړے دے که نۀ؛ او يوه
اړيکه هله او يوازې هله محاسبه کېدونکې ده چې دا محاسبه کولے شو چې
لړۍ \LR{$\tuple{n_1, \dots, n_k}$} د اړيکې غړے ده که نۀ۔ د يو
سټ يا اړيکې د \emph{ځانګړونکې تابع} په پام کښې نيولو سره، د محاسبه
کېدونکو سټونو او اړيکو بحث د محاسبه کېدونکو تابعو تر بحث لاندې
راوستل کېدے شي۔ په دې توګه بنسټيزې بازګشتي اړيکې هم تعريفولے شو؛ د
بېلګې په توګه، دا اړيکه چې «\LR{$n$}،~\LR{$m$} بې له پاتې شوني وېشي» يوه
بنسټيزه بازګشتي اړيکه ده۔
\olpseganchor{OLP-0210-B006}{end}

\olpseganchor{OLP-0210-B007}{start}
خو بنسټيزې بازګشتي تابعې---يعنې هغه چې له بنسټيزو تابعو څخه د
ترکيب او بنسټيز بازګښت په کارولو تعريفېدے شي---يوازينۍ محاسبه
کېدونکې د عددونو د تيورۍ تابعې نۀ دي۔ د بنسټيز بازګښت ډېرې عمومونې
څېړل شوې دي، خو د تابعو د محاسبې تر ټولو پياوړې او په پراخه کچه منل
شوې زياته لاره بې‌پوله پلټنه ده۔ دا د \emph{جزوي بازګشتي تابعو}
تعريف ته، او له دې سره تړلي د \emph{عمومي بازګشتي تابعو} تعريف ته
لار هواروي۔ عمومي بازګشتي تابعې محاسبه کېدونکې او هرځاے تعريف شوې
دي، او دا تعريف په کره توګه هغه جزوي بازګشتي تابعې ځانګړې کوي چې
هرځاے تعريف شوې وي۔ بازګشتي تابعې د محاسبې د هر بل مدل (ټيورينګ
ماشينونو، لمبډا حساب، او داسې نورو) تقليد کولے شي؛ ځکه نو د محاسبې
له ډېرو منل شوو مدلونو څخه يو مدل وړاندې کوي۔
\olpseganchor{OLP-0210-B007}{end}

\olpseganchor{OLP-0211-B004}{start}
\olfileid{cmp}{rec}{pre}
\olsection{بنسټيز بازګښت}
\olpseganchor{OLP-0211-B004}{end}

\olpseganchor{OLP-0211-B005}{start}
د طبيعي عددونو يو خاصيت دا دے چې هر طبيعي عدد پر~\LR{$0$} د راتلونکي
عمل~\LR{$+1$} په متناهي څو ځله کارولو سره تر لاسه کېدے شي---هر طبيعي عدد
يا~\LR{$0$} دے، يا د \dots{} د~\LR{$0$} راتلونکی دے۔ د دې حقيقت په کارولو سره
د يوې تابع~\LR{$h\colon\Nat \to \Nat$} د ټاکلو يوه لاره داسې ده:
(الف)~وټاکئ چې د آرګومېنټ~\LR{$0$} دپاره د~\LR{$h$} ارزښت څۀ دے، او
(ب)~دا هم وټاکئ چې د~\LR{$h(x)$} ارزښت له ورکړې وروسته د~\LR{$h(x+1)$} ارزښت
څنګه محاسبه شي۔ ځکه (الف) په نېغه راته وايي چې~\LR{$h(0)$} څۀ دے، نو
\LR{$h$} د~\LR{$0$} دپاره تعريف شوه۔ اوس د~\LR{$x=0$} دپاره د (ب) لارښوونه په
کارولو سره، له~\LR{$h(0)$} څخه \LR{$h(1) = h(0+1)$} محاسبه کولے شو۔ هماغه
لارښوونه د~\LR{$x=1$} دپاره په کارولو سره، له~\LR{$h(1)$} څخه
\LR{$h(2) = h(1+1)$} محاسبه کوو، او همداسې مخکښې ځو۔ د هر طبيعي
عدد~\LR{$x$} دپاره به بالاخره هغه پړاو ته ورسېږو چې~\LR{$h(x+1)$} له~\LR{$h(x)$}
څخه تعريفوو؛ نو~\LR{$h(x)$} د ټولو~\LR{$x \in \Nat$} دپاره تعريف شوې ده۔
\olpseganchor{OLP-0211-B005}{end}

\olpseganchor{OLP-0211-B006}{start}
د بېلګې په توګه، فرض کړئ~\LR{$h\colon \Nat \to \Nat$} په لاندې دوو
معادلو ټاکو:
\begin{align*}
h(0) & =  1\\
h(x+1) & =  2 \cdot h(x)
\end{align*}
که د ضرب طريقه له مخکښې را معلومه وي، نو دا معادلې د پورته (الف)
او~(ب) دپاره اړين معلومات راکوي۔ د دويمې معادلې پرله‌پسې په کارولو
سره لاس ته راوړو چې
\begin{align*}
  h(1) & = 2\cdot h(0) = 2,\\
  h(2) & = 2\cdot h(1) = 2\cdot 2,\\
  h(3) & = 2 \cdot h(2) = 2\cdot 2 \cdot 2,\\
  & \vdots
\end{align*}
وينو چې هغه تابع~\LR{$h$} چې ټاکلې مو ده، \LR{$h(x) = 2^x$} ده۔
\olpseganchor{OLP-0211-B006}{end}

\olpseganchor{OLP-0211-B007}{start}
د طبيعي عددونو ځانګړے خاصيت ضمانت کوي چې يوازې يوه تابع~\LR{$h$} شته
چې دا دواړه معيارونه پوره کوي۔ د دې غوندې معادلو جوړې ته د
تابع~\LR{$h$} \emph{په بنسټيز بازګښت تعريف} وايي۔ دې ته ځکه
«بازګشتي» وايي چې~\LR{$h$} په بازګشتي ډول تعريفوو، يعنې تعريف، او په
ځانګړي ډول دويمه معادله، پخپله~\LR{$h$}
په ښي اړخ کښې لري۔ دې ته ځکه «بنسټيز» وايي چې د~\LR{$h(x+1)$} په تعريف
کښې يوازې د~\LR{$h(x)$} ارزښت، يعنې سملاسي مخکينی ارزښت، کاروو۔ دا
پر~\LR{$\Nat$} د يوې تابع د بازګشتي تعريف تر ټولو ساده لاره ده۔
\olpseganchor{OLP-0211-B007}{end}

\olpseganchor{OLP-0211-B008}{start}
د جمع او ضرب غوندې لا بنسټيزې تابعې هم په بنسټيز بازګښت تعريفولے
شو۔ خو په دې حالتونو کښې اړوندې تابعې~\LR{$2$}-ځاييزې دي۔ د آرګومېنټونو
يو ځاے ثابت نيسو او بل د بازګښت دپاره کاروو۔ د بېلګې په توګه، د
\LR{$\Add(x, y)$} د تعريف دپاره~\LR{$x$} ثابت نيولے شو، ارزښت لومړے د~\LR{$y=0$}
دپاره تعريفولے شو، او بيا د~\LR{$y+1$} دپاره ئې د~\LR{$y$} له مخې تعريفولے
شو۔ ځکه چې~\LR{$x$} ثابت دے، د تعريفي معادلو په کيڼ او ښي اړخ دواړو
کښې به راځي۔
\begin{align*}
\Add(x,0) & =  x\\
\Add(x,y+1) & =  \Add(x,y)+1
\end{align*}
دا معادلې د ټولو~\LR{$x$} \emph{او}~\LR{$y$} دپاره د~\LR{$\Add$} ارزښت ټاکي۔ د
بېلګې په توګه، د~\LR{$\Add(2,3)$} د موندلو دپاره تعريفي معادلې د~\LR{$x = 2$}
دپاره کاروو: لومړۍ معادله کاروو چې~\LR{$\Add(2,0) = 2$} ومومو، بيا
دويمه معادله پرله‌پسې کاروو چې~\LR{$\Add(2,1) = 2 + 1 = 3$}،
\LR{$\Add(2, 2) = 3 + 1 = 4$}، او~\LR{$\Add(2, 3) = 4 + 1 = 5$} ومومو۔
\olpseganchor{OLP-0211-B008}{end}

\olpseganchor{OLP-0211-B009}{start}
د~\LR{$\Add$} په تعريف کښې مو د دويمې معادلې په ښي اړخ کښې~\LR{$+$} وکاراوه،
خو يوازې د~\LR{$1$} د زياتولو دپاره۔ په بله وينا، د راتلونکي
تابع~\LR{$\Succ(z) = z+1$} مو وکاروله او پر مخکيني ارزښت~\LR{$\Add(x,y)$} مو
عملي کړه، څو~\LR{$\Add(x,y+1)$} تعريف کړو۔ نو بازګشتي تعريف داسې ګڼلے شو
چې د يوې تابع له مخې ورکړل شوے دے او موږ ئې پر مخکيني ارزښت عملي
کوو۔ خو که تابع يوازې پر مخکيني ارزښت نۀ، بلکې پر~\LR{$x$} او~\LR{$y$} هم
تابع وګرځوو، زيان نۀ لري---او کله ناکله اړينه وي۔ دا تعريف وګورئ:
\begin{align*}
  \Mult(x,0) & =  0 \\
  \Mult(x,y+1) & =  \Add(\Mult(x,y),x)
\end{align*}
دا د تابع~\LR{$\Mult$} يو بنسټيز بازګشتي تعريف دے چې تابع~\LR{$\Add$} هم پر
مخکيني ارزښت~\LR{$\Mult(x,y)$} او هم پر لومړي آرګومېنټ~\LR{$x$} عملي کوي۔
دا هم تابع~\LR{$\Mult(x,y)$} د ټولو آرګومېنټونو~\LR{$x$} او~\LR{$y$} دپاره
تعريفوي۔ د بېلګې په توګه، \LR{$\Mult(2,3)$} د~\LR{$\Mult(2,0)$}،
\LR{$\Mult(2,1)$}، \LR{$\Mult(2,2)$} او~\LR{$\Mult(2,3)$} په پرله‌پسې محاسبې
سره ټاکل کېږي:
\begin{align*}
  \Mult(2,0) & = 0\\
  \Mult(2,1) & = \Mult(2,0+1) =
  \Add(\Mult(2,0), 2) = \Add(0, 2) = 2\\
  \Mult(2,2) & = \Mult(2,1+1) =
  \Add(\Mult(2,1), 2) = \Add(2, 2) = 4\\
  \Mult(2,3) & = \Mult(2,2+1) =
  \Add(\Mult(2,2), 2) = \Add(4, 2) = 6
\end{align*}
\olpseganchor{OLP-0211-B009}{end}

\olpseganchor{OLP-0211-B010}{start}
نو عمومي نمونه داسې ده: د يوې تابع~\LR{$h(x_0, \dots, x_{k-1}, y)$}
د بنسټيز بازګشتي تعريف دپاره دوه معادلې ورکوو۔ لومړۍ د~\LR{$h(x_0,
\dots, x_{k-1}, 0)$} ارزښت~\LR{$h$} ته له اشارې پرته تعريفوي۔ دويمه د
\LR{$h(x_0, \dots, x_{k-1}, y+1)$} ارزښت د~\LR{$h(x_0, \dots,
x_{k-1}, y)$}، نورو آرګومېنټونو~\LR{$x_0$}، \dots،~\LR{$x_{k-1}$}، او~\LR{$y$}
له مخې تعريفوي۔ په دې دويمه معادله کښې يوازې د~\LR{$h$} سملاسي مخکينی
ارزښت کارېدے شي۔ که د دې دوو معادلو د ښي اړخونو ورکړي عملونه هم
تابعې~\LR{$f$} او~\LR{$g$} وګڼو، نو په بنسټيز بازګښت د نوې تابع~\LR{$h$} د تعريف
عمومي نمونه داسې ده:
\begin{align*}
  h(x_0, \dots, x_{k-1}, 0) & = f(x_0, \dots, x_{k-1})\\
  h(x_0, \dots, x_{k-1}, y+1) & =
  g(x_0, \dots, x_{k-1}, y, h(x_0, \dots, x_{k-1}, y))
\end{align*}
د~\LR{$\Add$} په حالت کښې، \LR{$k=1$} او~\LR{$f(x_0) = x_0$} (عينيت تابع) دي،
او~\LR{$g(x_0, y, z) = z + 1$} (هغه~\LR{$3$}-ځاييزه تابع چې د خپل درېيم
آرګومېنټ راتلونکی ورکوي) ده:
\begin{align*}
  \Add(x_0, 0) & = f(x_0) = x_0\\
  \Add(x_0, y+1) & = g(x_0, y, \Add(x_0, y)) =
  \Succ(\Add(x_0, y))
\end{align*}
د~\LR{$\Mult$} په حالت کښې، \LR{$f(x_0) = 0$} (هغه ثابت تابع چې تل~\LR{$0$}
ورکوي) او~\LR{$g(x_0, y, z) = \Add(z,x_0)$} (هغه~\LR{$3$}-ځاييزه تابع چې د
خپل وروستي او لومړي آرګومېنټ مجموعه ورکوي) دي:
\begin{align*}
  \Mult(x_0,0) & =  f(x_0) = 0 \\
  \Mult(x_0,y+1) & =  g(x_0, y, \Mult(x_0,y)) =
  \Add(\Mult(x_0,y), x_0)
\end{align*}
\olpseganchor{OLP-0211-B010}{end}

\olpseganchor{OLP-0212-B004}{start}
\olfileid{cmp}{rec}{com}
\olsection{ترکيب}
\olpseganchor{OLP-0212-B004}{end}

\olpseganchor{OLP-0212-B005}{start}
که~\LR{$f$} او~\LR{$g$} د طبيعي عددونو دوه يوځاييزې تابعې وي، نو ترکيب ئې
کولے شو: \LR{$h(x) = f(g(x))$}۔ په دې حالت کښې نوې تابع~\LR{$h(x)$} له
تابعو~\LR{$f$} او~\LR{$g$} څخه په \emph{ترکيب} تعريف شوې ده۔ غواړو دا د يو
څخه د زياتو آرګومېنټونو تابعو ته عمومي کړو۔
\olpseganchor{OLP-0212-B005}{end}

\olpseganchor{OLP-0212-B006}{start}
د دې يوه لاره داسې ده: فرض کړئ~\LR{$f$} يوه~\LR{$k$}-ځاييزه تابع ده، او
\LR{$g_0$}، \dots،~\LR{$g_{k-1}$} هغه~\LR{$k$} تابعې دي چې ټولې~\LR{$n$}-ځاييزې دي۔
بيا يوه نوې~\LR{$n$}-ځاييزه تابع~\LR{$h$} داسې تعريفولے شو:
\[
h(x_0, \dots, x_{n-1}) =
f(g_0(x_0, \dots, x_{n-1}), \dots, g_{k-1}(x_0, \dots, x_{n-1}))
\]
که~\LR{$f$} او ټولې~\LR{$g_i$} محاسبه کېدونکې وي، نو~\LR{$h$} هم محاسبه کېدونکې
ده: د~\LR{$h(x_0, \dots, x_{n-1})$} د محاسبې دپاره لومړے د هر
\LR{$i = 0$}، \dots،~\LR{$k-1$} دپاره ارزښتونه~\LR{$y_i = g_i(x_0, \dots,
x_{n-1})$} محاسبه کړئ۔ بيا دا ارزښتونه~\LR{$f$} ته ورکړئ، څو
\LR{$h(x_0, \dots, x_{n-1}) = f(y_0, \dots, y_{k-1})$} محاسبه کړئ۔
\olpseganchor{OLP-0212-B006}{end}

\olpseganchor{OLP-0212-B007}{start}
ښايي دا د دې خبرې له حده زيات محدود بيان ښکاره شي چې له ځينو شته
تابعو څخه د نوې تابع د محاسبې پر مهال څۀ پېښېږي۔ يوه خبره دا ده چې
کله ناکله د يوې تابع ټول آرګومېنټونه نۀ کاروو؛ لکه هغه وخت چې د
\LR{$\Add$} په بنسټيز بازګشتي تعريف کښې مو د کارولو دپاره
\LR{$g(x, y, z) = \Succ(z)$} تعريف کړه۔ فرض کړئ د لاندې تابعو د کارولو
اجازه لرو:
\[
\Proj{n}{i}(x_0, \dots, x_{n-1}) = x_i
\]
تابعو~\LR{$\Proj{n}{i}$} ته \emph{پروجکشن} تابعې وايي:
\LR{$\Proj{n}{i}$} يوه~\LR{$n$}-ځاييزه تابع ده۔ بيا~\LR{$g$} داسې تعريفېدے شي:
\[
g(x, y, z) = \Succ(\Proj{3}{2}(x, y, z)).
\]
دلته د~\LR{$f$} ونډه~\LR{$1$}-ځاييزه تابع~\LR{$\Succ$} ترسره کوي، نو~\LR{$k=1$} دے۔ او
يوه~\LR{$3$}-ځاييزه تابع~\LR{$\Proj{3}{2}$} لرو چې د~\LR{$g_0$} ونډه ترسره کوي۔
پايله يوه~\LR{$3$}-ځاييزه تابع ده چې د درېيم آرګومېنټ راتلونکی ورکوي۔
\olpseganchor{OLP-0212-B007}{end}

\olpseganchor{OLP-0212-B008}{start}
پروجکشن تابعې دا اجازه هم راکوي چې د آرګومېنټونو په بيا ترتيبولو يا
سره يو کولو نوې تابعې تعريف کړو۔ د بېلګې په توګه، تابع
\LR{$h(x) = \Add(x, x)$} داسې تعريفېدے شي:
\[
h(x_0) = \Add(\Proj{1}{0}(x_0),\Proj{1}{0}(x_0)).
\]
دلته~\LR{$k=2$}، \LR{$n=1$}، د~\LR{$f(y_0,y_1)$} ونډه~\LR{$\Add$} ترسره کوي، او د
\LR{$g_0(x_0)$} او~\LR{$g_1(x_0)$} دواړو ونډې~\LR{$\Proj{1}{0}(x_0)$}، يعنې
يوځاييزه پروجکشن تابع (يا عينيت تابع)، ترسره کوي۔
\olpseganchor{OLP-0212-B008}{end}

\olpseganchor{OLP-0212-B009}{start}
که~\LR{$f(y_0, y_1)$} هغه تابع وي چې له مخکښې ئې لرو، نو تابع
\LR{$h(x_0, x_1) = f(x_1, x_0)$} داسې تعريفولے شو:
\[
h(x_0, x_1) = f(\Proj{2}{1}(x_0, x_1),\Proj{2}{0}(x_0, x_1)).
\]
دلته~\LR{$k=2$}، \LR{$n = 2$}، او د~\LR{$g_0$} او~\LR{$g_1$} ونډې په ترتيب سره
\LR{$\Proj{2}{1}$} او~\LR{$\Proj{2}{0}$} ترسره کوي۔
\olpseganchor{OLP-0212-B009}{end}

\olpseganchor{OLP-0212-B010}{start}
ښايي دا اندېښنه هم ولرئ چې د ټولو~\LR{$g_0$}، \dots،~\LR{$g_{k-1}$} ځاييزه
کچه بايد يو شان~\LR{$n$} وي۔ (په ياد ولرئ چې د يوې تابع
\emph{ځاييزه کچه} د آرګومېنټونو شمېر دے؛ د~\LR{$n$}-ځاييزې تابع ځاييزه
کچه~\LR{$n$} ده۔) خو د پروجکشن تابعو زياتول اړينه انعطاف‌پذيري برابروي۔
د بېلګې په توګه، فرض کړئ~\LR{$f$} او~\LR{$g$}، \LR{$3$}-ځاييزې تابعې دي او~\LR{$h$}
هغه~\LR{$2$}-ځاييزه تابع ده چې داسې تعريف شوې ده:
\[
h(x,y) = f(x,g(x,x,y),y).
\]
د~\LR{$h$} تعريف له پروجکشن تابعو سره داسې بيا ليکل کېدے شي:
\[
h(x,y) = f(\Proj{2}{0}(x,y), g(\Proj{2}{0}(x,y), \Proj{2}{0}(x,y),
\Proj{2}{1}(x,y)), \Proj{2}{1}(x,y)).
\]
بيا~\LR{$h$} له~\LR{$\Proj{2}{0}$}، \LR{$l$} او~\LR{$\Proj{2}{1}$} سره د~\LR{$f$} ترکيب
دے، چې
\[
l(x,y) = g(\Proj{2}{0}(x,y),\Proj{2}{0}(x,y),\Proj{2}{1}(x,y)),
\]
يعنې~\LR{$l$} له~\LR{$\Proj{2}{0}$}، \LR{$\Proj{2}{0}$} او~\LR{$\Proj{2}{1}$} سره د
\LR{$g$} ترکيب دے۔
\olpseganchor{OLP-0212-B010}{end}

\olpseganchor{OLP-0213-B004}{start}
\olfileid{cmp}{rec}{prf}
\olsection{بنسټيزې بازګشتي تابعې}
\olpseganchor{OLP-0213-B004}{end}

\olpseganchor{OLP-0213-B005}{start}
راځئ يو ځل بيا وليکو چې له شته تابعو څخه د بنسټيز بازګښت او ترکيب
په کارولو نوې تابعې څنګه تعريفولے شو۔
\olpseganchor{OLP-0213-B005}{end}

\olpseganchor{OLP-0213-B006}{start}
\begin{defn}
  \ollabel{defn:primitive-recursion} فرض کړئ~\LR{$f$} يوه~\LR{$k$}-ځاييزه
  تابع (\LR{$k\ge 1$}) او~\LR{$g$} يوه~\LR{$(k+2)$}-ځاييزه تابع ده۔ هغه تابع چې
  \emph{له~\LR{$f$} او~\LR{$g$} څخه په بنسټيز بازګښت تعريفېږي}،
  \LR{$(k+1)$}-ځاييزه تابع~\LR{$h$} ده چې لاندې معادلې ئې تعريفوي:
  \begin{align*}
    h(x_0,\dots,x_{k-1},0) & =  f(x_0,\dots,x_{k-1}) \\
    h(x_0,\dots,x_{k-1},y+1) & = g(x_0,\dots,x_{k-1}, y, h(x_0,\dots,x_{k-1}, y))
  \end{align*}
\end{defn}
\olpseganchor{OLP-0213-B006}{end}

\olpseganchor{OLP-0213-B007}{start}
\begin{defn}
  \ollabel{defn:composition} فرض کړئ~\LR{$f$} يوه~\LR{$k$}-ځاييزه تابع ده، او
  \LR{$g_0$}، \dots،~\LR{$g_{k-1}$} هغه~\LR{$k$} تابعې دي چې ټولې~\LR{$n$}-ځاييزې
  دي۔ هغه تابع چې \emph{له~\LR{$f$} او~\LR{$g_0$}، \dots،~\LR{$g_{k-1}$} څخه په
    ترکيب تعريفېږي}، \LR{$n$}-ځاييزه تابع~\LR{$h$} ده چې داسې تعريفېږي:
  \[
  h(x_0, \dots, x_{n-1}) =
  f(g_0(x_0, \dots, x_{n-1}), \dots, g_{k-1}(x_0, \dots, x_{n-1})).
  \]
\end{defn}
\olpseganchor{OLP-0213-B007}{end}

\olpseganchor{OLP-0213-B008}{start}
له~\LR{$\Succ$} او پروجکشن تابعو
\[
\Proj{n}{i}(x_0,\dots,x_{n-1}) = x_i,
\]
سره سره، د هر طبيعي عدد~\LR{$n$} او~\LR{$i < n$} دپاره تابع~\LR{$\Zero(x) = 0$}
هم په بنسټيزو بازګشتي تابعو کښې شاملوو۔
\olpseganchor{OLP-0213-B008}{end}

\olpseganchor{OLP-0213-B009}{start}
\begin{defn}
  د بنسټيزو بازګشتي تابعو سټ له~\LR{$\Nat^n$} څخه~\LR{$\Nat$} ته د تابعو هغه
  سټ دے چې په لاندې مادو استقرايي ډول تعريفېږي:
  \begin{enumerate}
  \item \LR{$\Zero$} بنسټيزه بازګشتي ده۔
  \item \LR{$\Succ$} بنسټيزه بازګشتي ده۔
  \item هره پروجکشن تابع~\LR{$\Proj{n}{i}$} بنسټيزه بازګشتي ده۔
  \item که~\LR{$f$} يوه~\LR{$k$}-ځاييزه بنسټيزه بازګشتي تابع وي او~\LR{$g_0$}،
    \dots،~\LR{$g_{k-1}$}، \LR{$n$}-ځاييزې بنسټيزې بازګشتي تابعې وي، نو له
    \LR{$g_0$}، \dots،~\LR{$g_{k-1}$} سره د~\LR{$f$} ترکيب بنسټيزه بازګشتي دے۔
  \item که~\LR{$f$} يوه~\LR{$k$}-ځاييزه بنسټيزه بازګشتي تابع او~\LR{$g$} يوه
    \LR{$k+2$}-ځاييزه بنسټيزه بازګشتي تابع وي، نو له~\LR{$f$} او~\LR{$g$} څخه په
    بنسټيز بازګښت تعريف شوې تابع بنسټيزه بازګشتي ده۔
\end{enumerate}
\end{defn}
\olpseganchor{OLP-0213-B009}{end}

\olpseganchor{OLP-0213-B010}{start}
\begin{explain}
په لنډه وينا، د بنسټيزو بازګشتي تابعو سټ هغه تر ټولو وړوکے سټ دے
چې~\LR{$\Zero$}، \LR{$\Succ$} او پروجکشن تابعې~\LR{$\Proj{n}{j}$} لري، او د ترکيب
او بنسټيز بازګښت له مخې بند دے۔
\olpseganchor{OLP-0213-B010}{end}

\olpseganchor{OLP-0213-B011}{start}
د بنسټيزو بازګشتي تابعو سټ د «پړاوونو» له مخې هم بيانولے شو۔
\LR{$S_0$} دې د پيلوونکو تابعو، يعنې~\LR{$\Zero$}، \LR{$\Succ$} او پروجکشنونو، سټ
وښيي۔ دا د~\LR{$0$} پړاو بنسټيزې بازګشتي تابعې دي۔ کله چې يو
پړاو~\LR{$S_i$} تعريف شو، \LR{$S_{i+1}$} دې د پخپله~\LR{$S_i$} د تابعو او هغو
ټولو تابعو سټ وي چې د همدې پړاو پر تابعو د ترکيب يا بنسټيز بازګښت
يوه بېلګه په کارولو لاس ته راځي۔ بيا
\[
S = \bigcup_{i \in \Nat} S_i
\]
د ټولو بنسټيزو بازګشتي تابعو سټ دے۔
د هر پړاو تابعې په ورپسې پړاو کښې هم وي۔ له دې کبله د ترکيب ټولې دننۍ تابعې په يوه ګډ پړاو کښې اخستلے شو۔
\end{explain}
\olpseganchor{OLP-0213-B011}{end}

\olpseganchor{OLP-0213-B012}{start}
راځئ تصديق کړو چې~\LR{$\Add$} بنسټيزه بازګشتي تابع ده۔
\olpseganchor{OLP-0213-B012}{end}

\olpseganchor{OLP-0213-B013}{start}
\begin{prop}
  د جمع تابع~\LR{$\Add(x,y) = x+y$} بنسټيزه بازګشتي ده۔
\end{prop}
\olpseganchor{OLP-0213-B013}{end}

\olpseganchor{OLP-0213-B014}{start}
\begin{proof}
د~\LR{$\Add$} يو بنسټيز بازګشتي تعريف مو لا له دوو تابعو~\LR{$f$} او~\LR{$g$} څخه
لرو چې د~\olref{defn:primitive-recursion} له شکل سره سمون خوري:
\begin{align*}
  \Add(x_0, 0) & = f(x_0) = x_0\\
  \Add(x_0, y+1) & = g(x_0, y, \Add(x_0, y)) =
  \Succ(\Add(x_0, y))
\end{align*}
نو~\LR{$\Add$} په دې شرط بنسټيزه بازګشتي ده چې~\LR{$f$} او~\LR{$g$} هم داسې وي۔
\LR{$f(x_0) = x_0 = \Proj{1}{0}(x_0)$} ده، او پروجکشن تابعې بنسټيزې
بازګشتي شمېرل کېږي، نو~\LR{$f$} بنسټيزه بازګشتي ده۔ تابع~\LR{$g$} هغه
درې‌ځاييزه تابع~\LR{$g(x_0, y, z)$} ده چې داسې تعريفېږي:
\[
g(x_0, y, z) = \Succ(z).
\]
دا لا نۀ ښيي چې~\LR{$g$} بنسټيزه بازګشتي ده، ځکه~\LR{$g$} او~\LR{$\Succ$} په کره
توګه يوه تابع نۀ دي: \LR{$\Succ$} يوځاييزه ده، خو~\LR{$g$} بايد درې‌ځاييزه
وي۔ مګر~\LR{$g$} په ترکيب په «رسمي» ډول داسې تعريفولے شو:
\[
g(x_0, y, z) = \Succ(\Proj{3}{2}(x_0, y, z))
\]
څرنګه چې~\LR{$\Succ$} او~\LR{$\Proj{3}{2}$} بنسټيزې بازګشتي تابعې شمېرل
کېږي، نو~\LR{$g$} هم داسې ده، ځکه چې له بنسټيزو بازګشتي تابعو څخه په
ترکيب تعريفېدے شي۔
\end{proof}
\olpseganchor{OLP-0213-B014}{end}

\olpseganchor{OLP-0213-B015}{start}
\begin{prop}
  \ollabel{prop:mult-pr}
  د ضرب تابع~\LR{$\Mult(x,y) = x \cdot y$} بنسټيزه بازګشتي ده۔
\end{prop}
\olpseganchor{OLP-0213-B015}{end}

\olpseganchor{OLP-0213-B016}{start}
\begin{proof}
  تمرين۔
\end{proof}
\olpseganchor{OLP-0213-B016}{end}

\olpseganchor{OLP-0213-B017}{start}
\begin{prob}
  \olref[cmp][rec][prf]{prop:mult-pr} په دې ښودلو ثابته کړئ چې د
  \LR{$\Mult$} بنسټيز بازګشتي تعريف د
  \olref[cmp][rec][prf]{defn:primitive-recursion} غوښتل شوي شکل ته
  راوستل کېدے شي، او اړوندې تابعې~\LR{$f$} او~\LR{$g$} بنسټيزې بازګشتي دي۔
\end{prob}
\olpseganchor{OLP-0213-B017}{end}

\olpseganchor{OLP-0213-B018}{start}
\begin{ex}
دا مو د بنسټيز بازګشتي تعريف تر ټولو لومړۍ بېلګه ده:
\begin{align*}
h(0) & =  1 \\
h(y+1) & =  2 \cdot h(y).
\end{align*}
دا تابع د~\olref{defn:primitive-recursion} غوښتل شوي شکل ته نۀ شي
راتلے، ځکه~\LR{$k=0$} دے۔ په تعريف کښې ثابتونه~\LR{$1$} او~\LR{$2$} هم راځي۔ د
لومړۍ ستونزې د هوارولو دپاره يو بې‌کاره آرګومېنټ ورزياتوو او
تابع~\LR{$h'$} تعريفوو:
\begin{align*}
h'(x_0, 0) & =  f(x_0) = 1 \\
h'(x_0, y+1) & =  g(x_0, y, h'(x_0, y)) = 2 \cdot h'(x_0, y).
\end{align*}
تابع~\LR{$f(x_0) = 1$} له~\LR{$\Succ$} او~\LR{$\Zero$} څخه په ترکيب تعريفېدے شي:
\LR{$f(x_0) = \Succ(\Zero(x_0))$}۔ تابع~\LR{$g$} له~\LR{$g'(z) = 2 \cdot z$} او
پروجکشنونو څخه په ترکيب تعريفېدے شي:
\begin{align*}
g(x_0, y, z) & = g'(\Proj{3}{2}(x_0, y, z))
\intertext{\RL{او~\LR{$g'$} بيا پخپله په ترکيب داسې تعريفېدے شي:}}
g'(z) & = \Mult(g''(z), \Proj{1}{0}(z))
\intertext{\RL{او}}
g''(z) & = \Succ(f(z)),
\end{align*}
چې~\LR{$f$} هماغه پورته تابع ده: \LR{$f(z) = \Succ(\Zero(z))$}۔ اوس
چې~\LR{$h'$} لرو، بيا ترکيب کارولے شو او~\LR{$h(y) =
h'(\Proj{1}{0}(y),\Proj{1}{0}(y))$} ټاکلے شو۔ دا ښيي چې~\LR{$h$} له
بنسټيزو تابعو څخه د ترکيبونو او بنسټيزو بازګښتونو د يوې لړۍ په
کارولو تعريفېدے شي، نو~\LR{$h$} بنسټيزه بازګشتي ده۔
\end{ex}
\olpseganchor{OLP-0213-B018}{end}

\olpseganchor{OLP-0214-B004}{start}
\olfileid{cmp}{rec}{not}
\olsection{د بنسټيز بازګښت نښې}
\olpseganchor{OLP-0214-B004}{end}

\olpseganchor{OLP-0214-B005}{start}
د بنسټيزو بازګشتي تابعو د کره استقرايي بيان يوه ګټه دا ده چې دا
تابعې په منظم ډول بيانولے شو۔ د بېلګې په توګه، هرې داسې تابع ته
لاندې يوه «نښه» ټاکلے شو۔ د صفر، راتلونکي او پروجکشنونو دپاره په
ترتيب سره نښې~\LR{$\Zero$}، \LR{$\Succ$} او~\LR{$\Proj{n}{i}$} وکاروئ۔ اوس فرض
کړئ~\LR{$h$} له يوې~\LR{$k$}-ځاييزې تابع~\LR{$f$} او~\LR{$n$}-ځاييزو تابعو~\LR{$g_0$}،
\dots،~\LR{$g_{k-1}$} څخه په ترکيب تعريف شوې ده، او وروستيو تابعو ته
مو نښې~\LR{$F$}، \LR{$G_0$}، \dots،~\LR{$G_{k-1}$} ټاکلې دي۔ بيا د يوې نوې
نښې~\LR{$\fn{Comp}_{k,n}$} په کارولو سره، تابع~\LR{$h$} په
\LR{$\fn{Comp}_{k,n}[F,G_0,\dots,G_{k-1}]$} ښودلے شو۔
\olpseganchor{OLP-0214-B005}{end}

\olpseganchor{OLP-0214-B006}{start}
په بنسټيز بازګښت د تعريف شوو تابعو دپاره ورته نښې کارولے شو۔ فرض
کړئ~\LR{$(k+1)$}-ځاييزه تابع~\LR{$h$} له~\LR{$k$}-ځاييزې تابع~\LR{$f$} او
\LR{$(k+2)$}-ځاييزې تابع~\LR{$g$} څخه په بنسټيز بازګښت تعريف شوې ده، او~\LR{$f$}
او~\LR{$g$} ته ټاکل شوې نښې په ترتيب سره~\LR{$F$} او~\LR{$G$} دي۔ بيا~\LR{$h$} ته
ټاکل شوې نښه~\LR{$\fn{Rec}_k[F,G]$} ده۔
\olpseganchor{OLP-0214-B006}{end}

\olpseganchor{OLP-0214-B007}{start}
په ياد ولرئ چې د جمع تابع په بنسټيز بازګښت داسې تعريف شوې ده:
\begin{align*}
  \Add(x_0, 0) & = \Proj{1}{0}(x_0) = x_0\\
  \Add(x_0, y+1) & = \Succ(\Proj{3}{2}(x_0, y, \Add(x_0, y))) = \Add(x_0, y) +1
\end{align*}
دلته د~\LR{$f$} ونډه~\LR{$\Proj{1}{0}$} ترسره کوي، او د~\LR{$g$} ونډه
\LR{$\Succ(\Proj{3}{2}(x_0, y, z))$} ترسره کوي۔ وروستۍ تابع ته
\LR{$\fn{Comp}_{1,3}[\Succ,\Proj{3}{2}]$} نښه ټاکل شوې، ځکه دا له
يوې~\LR{$1$}-ځاييزې تابع~\LR{$\Succ$} او~\LR{$3$}-ځاييزې تابع~\LR{$\Proj{3}{2}$} څخه په
ترکيب د يوې تابع د تعريف پايله ده۔ په دې ترتيب، د جمع تابع په لاندې
نښه ښودلے شو:
\[
\fn{Rec}_1[\Proj{1}{0},\fn{Comp}_{1,3}[\Succ,\Proj{3}{2}]].
\]
د دې نښو لرل کله ناکله ګټور وي؛ د بېلګې په توګه، د بنسټيزو
بازګشتي تابعو د شمېرنې پر مهال۔
\olpseganchor{OLP-0214-B007}{end}

\olpseganchor{OLP-0214-B008}{start}
\begin{prob}
  د~\LR{$\Mult$} بشپړه بنسټيزه بازګشتي نښه ورکړئ۔
\end{prob}
\olpseganchor{OLP-0214-B008}{end}

\olpseganchor{OLP-0215-B004}{start}
\olfileid{cmp}{rec}{cmp}
\olsection{بنسټيزې بازګشتي تابعې محاسبه کېدونکې دي}
\olpseganchor{OLP-0215-B004}{end}

\olpseganchor{OLP-0215-B005}{start}
فرض کړئ يوه تابع~\LR{$h$} په بنسټيز بازګښت تعريف شوې ده:
\begin{eqnarray*}
h(\vec x, 0)     & = & f(\vec x) \\
h(\vec x, y + 1) & = & g(\vec x, y, h(\vec x, y))
\end{eqnarray*}
او فرض کړئ تابعې~\LR{$f$} او~\LR{$g$} محاسبه کېدونکې دي۔ (موږ~\LR{$\vec x$} د
\LR{$x_0$}، \dots،~\LR{$x_{k-1}$} د لنډون دپاره کاروو۔) بيا
\LR{$h(\vec x, 0)$} په څرګند ډول محاسبه کېدے شي، ځکه دا يوازې
\LR{$f(\vec x)$} ده چې محاسبه کېدونکې مو فرض کړې ده۔ بيا
\LR{$h(\vec x, 1)$} هم محاسبه کېدے شي، ځکه~\LR{$1 = 0 + 1$} دے او په دې ډول
\LR{$h(\vec x, 1)$} يوازې دا ده:
\begin{align*}
 h(\vec x, 1) & = g(\vec x, 0, h(\vec x, 0)) =  g(\vec x, 0, f(\vec x)).
\intertext{\RL{په همدې ډول مخکښې تللے شو او دا محاسبه کولے شو:}}
h(\vec x, 2) & = g(\vec x, 1, h(\vec x, 1)) = g(\vec x, 1, g(\vec x, 0, f(\vec x)))\\
h(\vec x, 3) & = g(\vec x, 2, h(\vec x, 2)) = g(\vec x, 2, g(\vec x, 1, g(\vec x, 0, f(\vec x))))\\
h(\vec x, 4) & = g(\vec x, 3, h(\vec x, 3)) = g(\vec x, 3, g(\vec x, 2, g(\vec x, 1, g(\vec x, 0, f(\vec x)))))\\
& \vdots
\end{align*}
نو په عمومي ډول د~\LR{$h(\vec x, y)$} د محاسبې دپاره په پرله‌پسې توګه
\LR{$h(\vec x, 0)$}، \LR{$h(\vec x, 1)$}، \dots محاسبه کوو، تر هغه چې
\LR{$h(\vec x, y)$} ته ورسېږو۔
\olpseganchor{OLP-0215-B005}{end}

\olpseganchor{OLP-0215-B006}{start}
نو که تابعې~\LR{$f$} او~\LR{$g$} محاسبه کېدونکې وي، يو بنسټيز بازګشتي تعريف
يوه نوې محاسبه کېدونکې تابع ورکوي۔ همدارنګه، که تابعې~\LR{$f$} او~\LR{$g_i$}
محاسبه کېدونکې وي، د تابعو ترکيب هم محاسبه کېدونکې تابع ورکوي۔
\olpseganchor{OLP-0215-B006}{end}

\olpseganchor{OLP-0215-B007}{start}
څرنګه چې بنسټيزې تابعې~\LR{$\Zero$}، \LR{$\Succ$} او~\LR{$\Proj{n}{i}$} محاسبه
کېدونکې دي، او ترکيب او بنسټيز بازګښت له محاسبه کېدونکو تابعو څخه
محاسبه کېدونکې تابعې ورکوي، نو هره بنسټيزه بازګشتي تابع محاسبه
کېدونکې ده۔
\olpseganchor{OLP-0215-B007}{end}

\olpseganchor{OLP-0216-B004}{start}
\olfileid{cmp}{rec}{exa}
\olsection{د بنسټيزو بازګشتي تابعو بېلګې}
\olpseganchor{OLP-0216-B004}{end}

\olpseganchor{OLP-0216-B005}{start}
موږ لا د بنسټيزو بازګشتي تابعو ځينې بېلګې لرو: د جمع او ضرب تابعې~\LR{$\Add$} او~\LR{$\Mult$}۔ د عينيت تابع~\LR{$\fn{id}(x) = x$} بنسټيزه بازګشتي ده، ځکه دا يوازې~\LR{$\Proj{1}{0}$} ده۔ ثابتې تابعې \LR{$\fn{const}_n(x) = n$} هم بنسټيزې بازګشتي دي، ځکه هغوئ له~\LR{$\Zero$} او~\LR{$\Succ$} څخه په پرله‌پسې ترکيب تعريفېدے شي۔ دا هغه وخت ګټور دي چې په بنسټيزو بازګشتي تعريفونو کښې ثابتونه کارول غواړو۔ مثلاً، که تابع~\LR{$f(x) = 2 \cdot x$} تعريفول غواړو، نو له~\LR{$\fn{const}_n(x)$} او ضرب څخه ئې په ترکيب داسې ترلاسه کولے شو:
\LR{$f(x) = \Mult(\fn{const}_2(x), \Proj{1}{0}(x))$}۔ له دې وروسته به دا چل کاروو۔
\olpseganchor{OLP-0216-B005}{end}

\olpseganchor{OLP-0216-B006}{start}
\begin{prop}
  د توان تابع~\LR{$\fn{exp}(x, y) = x^y$} بنسټيزه بازګشتي ده۔
\end{prop}
\olpseganchor{OLP-0216-B006}{end}

\olpseganchor{OLP-0216-B007}{start}
\begin{proof}
  \LR{$\fn{exp}$} په بنسټيز بازګښت داسې تعريفولے شو:
  \begin{align*}
    \fn{exp}(x, 0) & = 1\\
    \fn{exp}(x, y+1) & = \Mult(x, \fn{exp}(x,y)).
    \intertext{\RL{په کره توګه دا لا له بنسټيزو بازګشتي تابعو څخه يو بازګشتي تعريف نۀ دے۔ خو رسمي بڼه ئې داسې ده:}}
    \fn{exp}(x, 0) & = f(x)\\
    \fn{exp}(x, y+1) & = g(x, y, \fn{exp}(x,y)).
    \intertext{\RL{چې پکښې}}
    f(x) & = \Succ(\Zero(x)) = 1\\
    g(x, y, z) & = \Mult(\Proj{3}{0}(x, y, z), \Proj{3}{2}(x, y, z)) = x \cdot z
  \end{align*}
  او په دې ډول~\LR{$f$} او~\LR{$g$} له بنسټيزو بازګشتي تابعو څخه په ترکيب تعريف شوې دي۔
\end{proof}
\olpseganchor{OLP-0216-B007}{end}

\olpseganchor{OLP-0216-B008}{start}
\begin{prop}
  مخکينۍ تابع~\LR{$\fn{pred}(y)$} چې داسې تعريف شوې ده،
  \[
  \fn{pred}(y) = \begin{cases}
    0 & \text{\RL{که \LR{$y=0$} وي}}\\
    y-1 & \text{\RL{که نۀ}}
  \end{cases}
  \]
  بنسټيزه بازګشتي ده۔
\end{prop}
\olpseganchor{OLP-0216-B008}{end}

\olpseganchor{OLP-0216-B009}{start}
\begin{proof}
 پام وکړئ چې
 \begin{align*}
   \fn{pred}(0) & = 0 \text{\RL{ او}}\\
   \fn{pred}(y+1) & = y.
 \end{align*}
 دا نژدې يو بنسټيز بازګشتي تعريف دے۔ په کره توګه د بنسټيز بازګښت د تعريف له نمونې سره ځکه نۀ سمېږي چې هغه نمونه لږ تر لږه يو زيات آرګومېنټ~\LR{$x$} غواړي۔ دا په دې لحاظ هم نابلده ده چې د \LR{$\fn{pred}(y+1)$} په تعريف کښې په رښتيا~\LR{$\fn{pred}(y)$} نۀ کاروي۔ خو لومړے~\LR{$\fn{pred}'(x, y)$} داسې تعريفولے شو:
 \begin{align*}
   \fn{pred}'(x, 0) & = \Zero(x) = 0,\\
   \fn{pred}'(x, y+1) & = \Proj{3}{1}(x, y, \fn{pred'}(x, y)) = y.
 \end{align*}
او بيا~\LR{$\fn{pred}$} ترې په ترکيب تعريفولے شو، مثلاً داسې:
\LR{$\fn{pred}(x) = \fn{pred}'(\Zero(x), \Proj{1}{0}(x))$}۔
\end{proof}
\olpseganchor{OLP-0216-B009}{end}

\olpseganchor{OLP-0216-B010}{start}
\begin{prop}
  د فاکتوريل تابع~\LR{$\fn{fac}(x) = \fact{x} = 1 \cdot 2 \cdot 3 \cdot \dots \cdot x$} بنسټيزه بازګشتي ده۔
\end{prop}
\olpseganchor{OLP-0216-B010}{end}

\olpseganchor{OLP-0216-B011}{start}
\begin{proof}
  څرګند بنسټيز بازګشتي تعريف دا دے:
  \begin{align*}
    \fn{fac}(0) &= 1\\
    \fn{fac}(y+1) & = \fn{fac}(y) \cdot (y+1).
    \intertext{\RL{په رسمي ډول لومړے دوه ځاييزه تابع~\LR{$h$} داسې تعريفوو:}}
    h(x, 0) & = \fn{const}_1(x)\\
    h(x, y+1) & = g(x, y, h(x, y))
    \intertext{\RL{چې پکښې \LR{$g(x, y, z) = \Mult(\Proj{3}{2}(x, y, z), \Succ(\Proj{3}{1}(x, y, z)))$} ده؛ بيا ږدو چې}}
    \fn{fac}(y) & = h(\Proj{1}{0}(y), \Proj{1}{0}(y)) = h(y,y).
  \end{align*}
  له دې وروسته به لږ کم رسمي واوسو او هر ځل به د ترکيب او بنسټيز بازګښت بشپړ رسمي تعريفونه نۀ ورکوو۔
\end{proof}
\olpseganchor{OLP-0216-B011}{end}

\olpseganchor{OLP-0216-B012}{start}
\begin{prop}
  پرې شوې تفريق~\LR{$x \tsub y$}، چې داسې تعريف شوې ده،
  \[
  x \tsub y = \begin{cases}
    0 & \text{\RL{که \LR{$x < y$} وي}}\\
    x-y & \text{\RL{که نۀ}}
  \end{cases}
  \]
  بنسټيزه بازګشتي ده۔
\end{prop}
\olpseganchor{OLP-0216-B012}{end}

\olpseganchor{OLP-0216-B013}{start}
\begin{proof}
  لرو چې:
  \begin{align*}
    x \tsub 0 & = x\\
    x \tsub (y+1) & = \fn{pred}(x \tsub y)
  \end{align*}
\end{proof}
\olpseganchor{OLP-0216-B013}{end}

\olpseganchor{OLP-0216-B014}{start}
\begin{prop}
 د~\LR{$x$} او~\LR{$y$} تر منځ واټن~\LR{$\left|x-y\right|$} بنسټيز بازګشتي دے۔
\end{prop}
\olpseganchor{OLP-0216-B014}{end}

\olpseganchor{OLP-0216-B015}{start}
\begin{proof}
  لرو چې~\LR{$\left| x-y \right| = (x \tsub y) + (y \tsub x)$}؛ نو واټن له~\LR{$+$} او~\LR{$\tsub$} څخه په ترکيب تعريفېدے شي، او دواړه بنسټيزې بازګشتي دي۔
\end{proof}
\olpseganchor{OLP-0216-B015}{end}

\olpseganchor{OLP-0216-B016}{start}
\begin{prop}
  د~\LR{$x$} او~\LR{$y$} تر ټولو لوے ارزښت~\LR{$\fn{max}(x,y)$} بنسټيز بازګشتي دے۔
\end{prop}
\olpseganchor{OLP-0216-B016}{end}

\olpseganchor{OLP-0216-B017}{start}
\begin{proof}
  \LR{$\fn{max}(x,y)$} له~\LR{$+$} او~\LR{$\tsub$} څخه په ترکيب داسې تعريفولے شو:
  \[
  \fn{max}(x,y) \defis x + (y \tsub x).
  \]
  که~\LR{$x$} تر ټولو لوے وي، يعنې~\LR{$x \ge y$}، نو~\LR{$y \tsub x = 0$}؛ ځکه \LR{$x + (y \tsub x) = x + 0 = x$}۔ که~\LR{$y$} تر ټولو لوے وي، نو \LR{$y \tsub x = y - x$}، او ځکه~\LR{$x + (y \tsub x) = x + (y - x) = y$}۔
\end{proof}
\olpseganchor{OLP-0216-B017}{end}

\olpseganchor{OLP-0216-B018}{start}
\begin{prop}
  \ollabel{prop:min-pr}
  د~\LR{$x$} او~\LR{$y$} تر ټولو وړوکے ارزښت~\LR{$\fn{min}(x,y)$} بنسټيز بازګشتي دے۔
\end{prop}
\olpseganchor{OLP-0216-B018}{end}

\olpseganchor{OLP-0216-B019}{start}
\begin{proof}
  تمرين۔
\end{proof}
\olpseganchor{OLP-0216-B019}{end}

\olpseganchor{OLP-0216-B020}{start}
\begin{prob}
  \olref[cmp][rec][exa]{prop:min-pr} ثابت کړئ۔
\end{prob}
\olpseganchor{OLP-0216-B020}{end}

\olpseganchor{OLP-0216-B021}{start}
\begin{prob}
وښيئ چې
\[
f(x, y) =
2^{(2^{\iddots^{2^{x}}})}\raisebox{1ex}{\bigg\rbrace}
\raisebox{1ex}{\text {\LR{$y$} ځله \LR{$2$}}}
\]
بنسټيزه بازګشتي ده۔
\end{prob}
\olpseganchor{OLP-0216-B021}{end}

\olpseganchor{OLP-0216-B022}{start}
\begin{prob}
وښيئ چې صحيح وېش~\LR{$d(x, y) = \lfloor x/y \rfloor$} (يعنې هغه وېش چې د اعشاري ټکي وروسته هر څه پکښې غورځوئ) بنسټيز بازګشتي دے۔ کله چې~\LR{$y = 0$} وي، شرط ږدو چې~\LR{$d(x, y) = 0$} وي۔ د بنسټيز بازګښت او ترکيب په کارولو د~\LR{$d$} يو څرګند تعريف ورکړئ۔
\end{prob}
\olpseganchor{OLP-0216-B022}{end}

\olpseganchor{OLP-0216-B023}{start}
\begin{prop}
د بنسټيزو بازګشتي تابعو سټ د لاندې دوو عملونو له مخې بند دے:
\begin{enumerate}
\item متناهي جمعې: که~\LR{$f(\vec x, z)$} بنسټيزه بازګشتي وي، نو لاندې تابع هم داسې ده:
\[
g(\vec x, y) \defis \sum_{z = 0}^y f(\vec x, z).
\]
\item متناهي حاصل‌ضربونه: که~\LR{$f(\vec x, z)$} بنسټيزه بازګشتي وي، نو لاندې تابع هم داسې ده:
\[
h(\vec x, y) \defis \prod_{z = 0}^y f(\vec x, z).
\]
\end{enumerate}
\end{prop}
\olpseganchor{OLP-0216-B023}{end}

\olpseganchor{OLP-0216-B024}{start}
\begin{proof}
د بېلګې په توګه، متناهي جمعې په لاندې معادلو بازګشتي تعريفېږي:
\begin{align*}
  g(\vec x, 0) & = f(\vec x, 0)\\
  g(\vec x, y+1) & = g(\vec x, y) + f(\vec x, y+1).
\end{align*}
\end{proof}
\olpseganchor{OLP-0216-B024}{end}

\olpseganchor{OLP-0217-B004}{start}
\olfileid{cmp}{rec}{prr}
\olsection{بنسټيزې بازګشتي اړيکې}
\olpseganchor{OLP-0217-B004}{end}

\olpseganchor{OLP-0217-B005}{start}
\begin{defn}
اړيکه~\LR{$R(\vec x)$} هله بنسټيزه بازګشتي بلل کېږي چې ځانګړونکې تابع ئې
\[
\Char{R}(\vec x) = \left\{
  \begin{array}{ll}
  1 & \mbox{که \LR{$R(\vec x)$} وي} \\
  0 & \mbox{که نۀ}
  \end{array}
\right.
\]
بنسټيزه بازګشتي وي۔
\end{defn}
\olpseganchor{OLP-0217-B005}{end}

\olpseganchor{OLP-0217-B006}{start}
په بله وينا، کله چې د بنسټيزې بازګشتي اړيکې~\LR{$R(\vec x)$} خبره کوو، مطلب مو د~\LR{$\Char{R}(\vec x) = 1$} په بڼه اړيکه وي، چې~\LR{$\Char{R}$} يوه بنسټيزه بازګشتي تابع ده او پر هر ننوت 1 يا 0 راګرځوي۔ د بېلګې په توګه، اړيکه~\LR{$\fn{IsZero}(x)$} هله او يوازې هله صدق کوي چې \LR{$x = 0$} وي، او له هغې سره سمون لرونکې تابع~\LR{$\Char{\fn{IsZero}}$} په بنسټيز بازګښت داسې تعريفېږي:
\begin{align*}
\Char{\fn{IsZero}}(0) & = 1,\\
\Char{\fn{IsZero}}(x+1) & = 0.
\end{align*}
\olpseganchor{OLP-0217-B006}{end}

\olpseganchor{OLP-0217-B007}{start}
بايد څرګنده وي چې اړيکې له نورو بنسټيزو بازګشتي تابعو سره ترکيبولے شو۔ نو لاندې اړيکې هم بنسټيزې بازګشتي دي:
\begin{enumerate}
\item د برابرۍ اړيکه~\LR{$x = y$}، چې په~\LR{$\fn{IsZero}(\left|x-y\right|)$} تعريفېږي؛
\item د کوچني-يا-برابر اړيکه~\LR{$x \leq y$}، چې په~\LR{$\fn{IsZero}(x \tsub y)$} تعريفېږي۔
\end{enumerate}
\olpseganchor{OLP-0217-B007}{end}

\olpseganchor{OLP-0217-B008}{start}
\begin{prop}
  د بنسټيزو بازګشتي اړيکو سټ د بولي عملونو له مخې بند دے؛ يعنې که \LR{$P(\vec x)$} او~\LR{$Q(\vec x)$} بنسټيزې بازګشتي وي، نو لاندې اړيکې هم داسې دي:
  \begin{enumerate}
  \item \LR{$\lnot P(\vec x)$}
  \item \LR{$P(\vec x) \land Q(\vec x)$}
  \item \LR{$P(\vec x) \lor Q(\vec x)$}
  \item \LR{$P(\vec x) \lif Q(\vec x)$}
  \end{enumerate}
\end{prop}
\olpseganchor{OLP-0217-B008}{end}

\olpseganchor{OLP-0217-B009}{start}
\begin{proof}
  فرض کړئ~\LR{$P(\vec x)$} او~\LR{$Q(\vec x)$} بنسټيزې بازګشتي دي، يعنې ځانګړونکې تابعې~\LR{$\Char{P}$} او~\LR{$\Char{Q}$} ئې بنسټيزې بازګشتي دي۔ بايد وښيو چې د~\LR{$\lnot P(\vec x)$} او نورو ځانګړونکې تابعې هم بنسټيزې بازګشتي دي۔
  \[
  \Char{\lnot P}(\vec x) = \begin{cases}
    0 & \text{\RL{که \LR{$\Char{P}(\vec x) = 1$} وي}}\\
    1 & \text{\RL{که نۀ}}
  \end{cases}
  \]
  \LR{$\Char{\lnot P}(\vec x)$} د~\LR{$1 \tsub \Char{P}(\vec x)$} په توګه تعريفولے شو۔
  \[
  \Char{P \land Q}(\vec x) = \begin{cases}
    1 & \text{\RL{که \LR{$\Char{P}(\vec x) = \Char{Q}(\vec x) = 1$} وي}}\\
    0 & \text{\RL{که نۀ}}
  \end{cases}
  \]
  \LR{$\Char{P \land Q}(\vec x)$} د~\LR{$\Char{P}(\vec x) \cdot \Char{Q}(\vec x)$}، يا د~\LR{$\fn{min}(\Char{P}(\vec x), \Char{Q}(\vec x))$} په توګه تعريفولے شو۔ په همدې ډول،
  \begin{align*}
    \Char{P \lor Q}(\vec x) & = \fn{max}(\Char{P}(\vec x), \Char{Q}(\vec x)) \text{\RL{ او}}\\
    \Char{P \lif Q}(\vec x) & = \fn{max}(1 \tsub \Char{P}(\vec x), \Char{Q}(\vec x)).
  \end{align*}
\end{proof}
\olpseganchor{OLP-0217-B009}{end}

\olpseganchor{OLP-0217-B010}{start}
\begin{prop}
  د بنسټيزو بازګشتي اړيکو سټ د محدودې کميت ټاکنې له مخې بند دے؛ يعنې که~\LR{$R(\vec x, z)$} بنسټيزه بازګشتي اړيکه وي، نو لاندې اړيکې هم داسې دي:
  \begin{align*}
    & \bforall{z < y}{R(\vec x, z)} \text{\RL{ او}}\\
    & \bexists{z < y}{R(\vec x, z)}.
  \end{align*}
  \LR{$\bforall{z < y}{R(\vec x, z)}$} د~\LR{$\vec x$} او~\LR{$y$} په باب هله او يوازې هله صدق کوي چې~\LR{$R(\vec x, z)$} له~\LR{$y$} څخه د هر کوچني~\LR{$z$} دپاره صدق وکړي؛ د~\LR{$\bexists{z < y}{R(\vec x, z)}$} دپاره هم ورته تشريح ده۔
\end{prop}
\olpseganchor{OLP-0217-B010}{end}

\olpseganchor{OLP-0217-B011}{start}
\begin{proof}
  د دود له مخې~\LR{$\bforall{z < 0}{R(\vec x, z)}$} رښتونے نيسو (په ساده دليل چې له~\LR{$0$} څخه کوچنے~\LR{$z$} نشته)، او \LR{$\bexists{z < 0}{R(\vec x, z)}$} دروغ نيسو۔ محدود کلي کميت ټاکونکی کټ مټ د متناهي حاصل‌ضرب يا تکراري کمينې په شان کار کوي؛ يعنې که \LR{$P(\vec x, y) \defiff \bforall{z < y}{R(\vec x, z)}$} وي، نو \LR{$\Char{P}(\vec x, y)$} داسې تعريفېدے شي:
  \begin{align*}
    \Char{P}(\vec x, 0) & = 1\\
    \Char{P}(\vec x, y+1) & = \fn{min}(\Char{P}(\vec x, y), \Char{R}(\vec x, y)).
  \end{align*}
  محدوده وجودي کميت ټاکنه هم په ورته ډول د~\LR{$\fn{max}$} په کارولو تعريفېدے شي۔ يا ئې له محدودې کلي کميت ټاکنې څخه د دې برابرۍ په کارولو تعريفولے شو:
  \LR{$\bexists{z < y}{R(\vec x, z)} \liff \lnot \bforall{z < y}{\lnot R(\vec x, z)}$}۔ پام وکړئ چې مثلاً د~\LR{$\bexists{x \leq y}{\dots x\dots}$} په بڼه محدود کميت ټاکونکی له~\LR{$\bexists{x < y+1}{\dots x \dots}$} سره برابر دے۔
\end{proof}
\olpseganchor{OLP-0217-B011}{end}

\olpseganchor{OLP-0217-B012}{start}
\begin{prob}
  وښيئ چې درې ځاييزه اړيکه~\LR{$x \equiv y \mod n$} (مودولو~\LR{$n$} تطابق) بنسټيزه بازګشتي ده۔
\end{prob}
\olpseganchor{OLP-0217-B012}{end}

\olpseganchor{OLP-0217-B013}{start}
يوه بله ګټوره بنسټيزه بازګشتي تابع شرطي تابع~\LR{$\fn{cond}(x,y,z)$} ده، چې داسې تعريفېږي:
\begin{align*}
  \fn{cond}(x,y,z) & = \begin{cases}
  y & \text{\RL{که \LR{$x = 0$} وي}} \\
  z & \text{\RL{که نۀ}}.
\end{cases}
\intertext{\RL{دا په بازګشتي ډول داسې تعريفېږي:}}
\fn{cond}(0,y,z) & = y,\\
\fn{cond}(x+1,y,z) & = z.
\end{align*}
له دې څخه د بنسټيزو بازګشتي اړيکو پر بنسټ په حالتونو د بنسټيزو بازګشتي تابعو تعريفونه توجيه کولے شو:
\olpseganchor{OLP-0217-B013}{end}

\olpseganchor{OLP-0217-B014}{start}
\begin{prop}
که~\LR{$g_0(\vec x)$}، \dots،~\LR{$g_m(\vec x)$} بنسټيزې بازګشتي تابعې وي، او~\LR{$R_0(\vec x)$}، \dots،~\LR{$R_{m-1}(\vec x)$} بنسټيزې بازګشتي اړيکې وي، نو لاندې تعريف شوې تابع~\LR{$f$} هم بنسټيزه بازګشتي ده:
\[
f(\vec x) = \begin{cases}
    g_0(\vec x) & \text{\RL{که \LR{$R_0(\vec{x})$} وي}} \\
    g_1(\vec x) & \text{\RL{که \LR{$R_1(\vec{x})$} وي او \LR{$R_0(\vec{x})$} نۀ وي}} \\
    \vdots & \\
    g_{m-1}(\vec x) & \text{\RL{که \LR{$R_{m-1}(\vec{x})$} وي او مخکيني هېڅ يو نۀ وي}} \\
    g_m(\vec x) & \mbox{که نۀ}
\end{cases}
\]
\end{prop}
\olpseganchor{OLP-0217-B014}{end}

\olpseganchor{OLP-0217-B015}{start}
\begin{proof}
  کله چې~\LR{$m = 1$} وي، دا يوازې لاندې تابع ده:
  \[
  f(\vec x) = \fn{cond}(\Char{\lnot R_0}(\vec x),g_0(\vec x),g_1(\vec x)).
  \]
  کله چې~\LR{$m$} له~\LR{$1$} څخه لوے وي، يوازې د همدې بڼې تعريفونه سره ترکيبولے شو۔
\end{proof}
\olpseganchor{OLP-0217-B015}{end}

\olpseganchor{OLP-0218-B004}{start}
\olfileid{cmp}{rec}{bmi}
\olsection{محدوده کمينه‌موندنه}
\olpseganchor{OLP-0218-B004}{end}

\olpseganchor{OLP-0218-B005}{start}
\begin{explain}
ډېر وخت دا ګټوره وي چې يوه تابع د هغه تر ټولو کوچني عدد په توګه تعريف کړو چې کوم خاصيت يا اړيکه~\LR{$P$} پوره کوي۔ که~\LR{$P$} پرېکړه کېدونکې وي، دا تابع يوازې په دې ډول محاسبه کولے شو چې ټول ممکن عددونه~\LR{$0$}، \LR{$1$}، \LR{$2$}، \dots يو په بل پسې وازمېيو، تر هغه چې تر ټولو کوچنے عدد ومومو چې~\LR{$P$} پوره کوي۔ دا ډول بې‌پولې پلټنه مو د بنسټيزو بازګشتي تابعو له چوکاټه باسي۔ خو که يوازې په هغه تر ټولو کوچني عدد کښې دلچسپي ولرو چې
\emph{له کوم خپلواکه ورکړل شوي حد څخه کوچنے وي}، نو بنسټيز بازګشتي پاتې کېږو۔ په بله او لږ عمومي وينا، فرض کړئ يوه بنسټيزه بازګشتي اړيکه~\LR{$R(x,z)$} لرو۔ هغه تابع په پام کښې ونيسئ چې \LR{$x$} او~\LR{$y$} هغه تر ټولو کوچني~\LR{$z < y$} ته لېږي چې~\LR{$R(x, z)$} پوره کوي۔ دا تابع هم د~\LR{$R(x, 0)$}، \LR{$R(x, 1)$}، \dots، \LR{$R(x, y-1)$} په ازمېيلو محاسبه کېدے شي۔ خو دا ولې بنسټيزه بازګشتي ده؟
\end{explain}
\olpseganchor{OLP-0218-B005}{end}

\olpseganchor{OLP-0218-B006}{start}
\begin{prop}
که~\LR{$R(\vec x, z)$} بنسټيزه بازګشتي وي، نو تابع~\LR{$m_R(\vec{x}, y)$} هم داسې ده؛ دا تابع، که داسې عدد وي، له~\LR{$y$} څخه تر ټولو کوچنے~\LR{$z$} راګرځوي چې~\LR{$R(\vec x, z)$} پوره کړي، او که داسې عدد نۀ وي نو~\LR{$y$} راګرځوي۔ تابع~\LR{$m_R$} به داسې ليکو:
\[
\bmin{z < y}{R(\vec{x}, z)},
\]
\end{prop}
\olpseganchor{OLP-0218-B006}{end}

\olpseganchor{OLP-0218-B007}{start}
\begin{proof}
داسې هېڅ~\LR{$z < 0$} نشته چې~\LR{$R(\vec{x}, z)$} پوره کړي، ځکه په اصل کښې هېڅ~\LR{$z < 0$} نشته۔ نو~\LR{$m_R(\vec x, 0) = 0$}۔
\olpseganchor{OLP-0218-B007}{end}

\olpseganchor{OLP-0218-B008}{start}
کله چې حد د~\LR{$y + 1$} په بڼه وي، درې حالتونه لرو:
\begin{enumerate}
\item داسې~\LR{$z < y$} شته چې~\LR{$R(\vec{x}, z)$} پوره کوي؛ په دې حالت کښې \LR{$m_R(\vec{x}, y+1) = m_R(\vec{x}, y)$}۔
\item داسې~\LR{$z<y$} نشته، خو~\LR{$R(\vec{x}, y)$} پوره کېږي؛ بيا \LR{$m_R(\vec{x}, y+1) = y$}۔
\item داسې~\LR{$z < y+1$} نشته چې~\LR{$R(\vec{x}, z)$} پوره کړي؛ بيا \LR{$m_R(\vec{x}, y+1) = y+1$}۔
\end{enumerate}
په دغو درې واړو حالتونو کښې د پاراميټرونو لومړنی وکټور ثابت ساتو۔
نو د~\LR{$m_R(\vec x, 0)$} پېلنی ارزښت او د راتلونکي ارزښت قاعده په يوه بنسټيز بازګښت کښې سره يوځاے کوو۔ د برابرۍ ازموينه او د اړيکې ځانګړونکې تابع دواړه بنسټيزې بازګشتي دي، نو د حالتونو له مخې دا ټاکنه هم بنسټيزه بازګشتي ده:
\begin{align*}
m_R(\vec{x}, 0) & = 0\\
m_R(\vec{x}, y+1) & =
\begin{cases}
m_R(\vec{x}, y) & \text{\RL{که \LR{$m_R(\vec x, y) \neq y$} وي}}\\
y & \text{\RL{که \LR{$m_R(\vec x, y) = y$} او \LR{$R(\vec{x}, y)$} وي}}\\
y+1 & \text{\RL{که نۀ۔}}
\end{cases}
\end{align*}
پام وکړئ چې داسې~\LR{$z<y$} شته چې~\LR{$R(\vec x, z)$} پوره کړي، هله او يوازې هله چې~\LR{$m_R(\vec x, y) \neq y$} وي۔
\end{proof}
\olpseganchor{OLP-0218-B008}{end}

\olpseganchor{OLP-0218-B009}{start}
\begin{prob}
فرض کړئ~\LR{$R(\vec x, z)$} بنسټيزه بازګشتي ده۔ تابع~\LR{$m'_R(\vec{x}, y)$} چې، که داسې عدد وي، له~\LR{$y$} څخه تر ټولو کوچنے~\LR{$z$} راګرځوي چې \LR{$R(\vec{x}, z)$} پوره کړي، او که داسې عدد نۀ وي نو~\LR{$0$} راګرځوي، له~\LR{$\Char{R}$} څخه په بنسټيز بازګښت تعريف کړئ۔
\end{prob}
\olpseganchor{OLP-0218-B009}{end}

\olpseganchor{OLP-0219-B004}{start}
\olfileid{cmp}{rec}{pri}
\olsection{اوليه عددونه}
\olpseganchor{OLP-0219-B004}{end}

\olpseganchor{OLP-0219-B005}{start}
محدوده کميت ټاکنه او محدوده کمينه‌موندنه د دې دپاره ډېر وسايل راکوي چې وښيو طبيعي تابعې او اړيکې بنسټيزې بازګشتي دي۔ د بېلګې په توګه اړيکه «\LR{$x$}، \LR{$y$} وېشي» په پام کښې ونيسئ، چې په~\LR{$x \mid y$} ليکل کېږي۔ اړيکه~\LR{$x \mid y$} هله صدق کوي چې~\LR{$y$} پر~\LR{$x$} بې له پاتې شوني ووېشل شي، يعنې~\LR{$y$} د~\LR{$x$} صحيح مضرب وي۔ که دا اړيکه صدق ونۀ کړي، يعنې کله چې د~\LR{$y$} پر~\LR{$x$} وېش د پاتې شوني په ژبه بيانېدے شي او پاتې شونی~\LR{$> 0$} وي، نو~\LR{$x \nmid y$} ليکو۔ په بله وينا، \LR{$x \mid y$} هله او يوازې هله چې د کوم~\LR{$z$} دپاره~\LR{$x \cdot z = y$} وي۔ کله چې~\LR{$x>0$} وي، له~\LR{$x \cdot z = y$} څخه~\LR{$z \leq y$} راځي۔ که~\LR{$x=0$} وي، معادله يوازې د~\LR{$y=0$} لپاره شاهد لري، او په دې حالت کښې~\LR{$z=0$} ټاکلے شو۔ نو~\LR{$x \mid y$} هله او يوازې هله چې د کوم~\LR{$z \le y$} دپاره~\LR{$x \cdot z = y$} وي۔ له~\LR{$=$} او ضرب څخه په محدودې وجودي کميت ټاکنې اړيکه~\LR{$x \mid y$} داسې تعريفولے شو:
\[
x \mid y \defiff \bexists{z \leq y}{(x \cdot z) = y}.
\]
نو مو وښودله چې~\LR{$x \mid y$} بنسټيزه بازګشتي ده۔
\olpseganchor{OLP-0219-B005}{end}

\olpseganchor{OLP-0219-B006}{start}
طبيعي عدد~\LR{$x$} هله
\emph{اوليه} دے چې نۀ~\LR{$0$} وي، نۀ~\LR{$1$}، او يوازې پر~\LR{$1$} او پر خپل ځان وېشل کېږي۔ په بله وينا، اوليه عددونه داسې دي چې کله هم~\LR{$y \mid x$} وي، نو يا~\LR{$y = 1$} وي يا~\LR{$y=x$}۔ د دې د ازمېيلو دپاره چې~\LR{$x$} اوليه دے که نۀ، يوازې د ټولو~\LR{$y \le x$} دپاره \LR{$y \mid x$} ازمېيل پکار دي، ځکه د~\LR{$x$} د مثبت‌والي په فرض که~\LR{$y > x$} وي، نو په خپله \LR{$y \nmid x$} وي۔ نو اړيکه~\LR{$\fn{Prime}(x)$}، چې هله او يوازې هله صدق کوي چې~\LR{$x$} اوليه وي، داسې تعريفېدے شي:
\[
\fn{Prime}(x) \defiff x \geq 2 \land \bforall{y \leq x}{(y \mid x \lif y = 1 \lor y = x)}
\]
او ځکه بنسټيزه بازګشتي ده۔
\olpseganchor{OLP-0219-B006}{end}

\olpseganchor{OLP-0219-B007}{start}
اوليه عددونه~\LR{$2$}، \LR{$3$}، \LR{$5$}، \LR{$7$}، \LR{$11$}، او داسې نور دي۔ تابع~\LR{$p(x)$} په پام کښې ونيسئ چې په دې لړۍ کښې~\LR{$x$}م اوليه عدد راګرځوي؛ يعنې \LR{$p(0) = 2$}، \LR{$p(1) = 3$}، \LR{$p(2) = 5$}، او داسې نور۔ (د اسانتيا دپاره به ډېر وخت~\LR{$p(x)$} د~\LR{$p_x$} په توګه ليکو: \LR{$p_0=2$}، \LR{$p_1=3$}، او داسې نور۔)
\olpseganchor{OLP-0219-B007}{end}

\olpseganchor{OLP-0219-B008}{start}
که يوه تابع~\LR{$\fn{nextPrime}(x)$} مو لرله چې له~\LR{$x$} څخه لومړنے لوے اوليه عدد راګرځوي، نو~\LR{$p$} په بنسټيز بازګښت په اسانۍ تعريفېدے شي:
\begin{align*}
  p(0) & = 2\\
  p(x+1) & = \fn{nextPrime}(p(x))
\end{align*}
څرنګه چې~\LR{$\fn{nextPrime}(x)$} هغه تر ټولو کوچنے~\LR{$y$} دے چې~\LR{$y > x$} وي او~\LR{$y$} اوليه وي، دا په بې‌پولې پلټنه په اسانۍ محاسبه کېدے شي۔ خو د اقليدس د يوې پايلې له امله ئې په محدودې کمينه‌موندنې هم تعريفولے شو: تل داسې يو اوليه عدد شته چې له~\LR{$x$} څخه لوے او له \LR{$\fact{x}+1$} څخه کوچنے يا برابر وي۔
\[
  \fn{nextPrime}(x) =
  \bmin{y \leq \fact{x}+1}{(y > x \land \fn{Prime}(y))}.
\]
دا ښيي چې~\LR{$\fn{nextPrime}(x)$}، او ځکه~\LR{$p(x)$}، يوازې محاسبه کېدونکې نۀ بلکې بنسټيزه بازګشتي ده۔
\olpseganchor{OLP-0219-B008}{end}

\olpseganchor{OLP-0219-B009}{start}
(که لېواله يئ، دا د اقليدس د قضيې يو لنډ ثبوت دے۔ که ورکړل شوے عدد له دوو کوچنے وي، 2 نېغ په نېغه مطلوب اوليه عدد دے۔ په پاتې حالت کښې فرض کړئ~\LR{$p_n$} تر ټولو لوے اوليه عدد دے چې~\LR{$\le x$} وي۔ د ټولو هغو اوليه عددونو حاصل‌ضرب \LR{$p = p_0\cdot p_1 \cdot \dots \cdot p_n$} په پام کښې ونيسئ چې \LR{$\le x$} دي۔ يا~\LR{$p+1$} اوليه دے، يا د~\LR{$x$} او~\LR{$p+1$} تر منځ بل اوليه عدد شته۔ ولې؟ فرض کړئ~\LR{$p+1$} اوليه نۀ دے۔ بيا کوم اوليه عدد \LR{$q \mid p+1$} شته چې~\LR{$q < p+1$} وي۔ اوليه عددونه~\LR{$\le x$} يو هم~\LR{$p+1$} نۀ وېشي۔ (د~\LR{$p$} د تعريف له مخې هر اوليه \LR{$p_i \le x$}، \LR{$p$} وېشي، يعنې پاتې شونی~\LR{$0$} ورکوي۔ نو هر \LR{$p_i \le x$} د~\LR{$p+1$} په وېش کښې پاتې شونی~\LR{$1$} ورکوي، او ځکه \LR{$p_i \nmid p+1$}۔) نو~\LR{$q$} يو اوليه عدد دے چې~\LR{$>x$} او~\LR{$<p+1$} دے۔ او \LR{$p \le \fact{x}$}، نو داسې اوليه عدد شته چې~\LR{$>x$} او \LR{$\le \fact{x}+1$} وي۔)
\olpseganchor{OLP-0219-B009}{end}

\olpseganchor{OLP-0219-B010}{start}
\begin{prob}
صحيح وېش~\LR{$d(x, y)$} د محدودې کمينه‌موندنې په کارولو تعريف کړئ۔
\end{prob}
\olpseganchor{OLP-0219-B010}{end}

\olpseganchor{OLP-0220-B004}{start}
\olfileid{cmp}{rec}{seq}
\olsection{لړۍ}
\olpseganchor{OLP-0220-B004}{end}

\olpseganchor{OLP-0220-B005}{start}
د بنسټيزو بازګشتي تابعو سټ د پام وړ ټينګ دے۔ خو کله چې د \emph{لړيو} د سمبالولو يوه مناسبه طريقه جوړه کړو، لا ډېر څه کولے شو۔ د طبيعي عددونو متناهي لړۍ به له طبيعي عددونو سره په لاندې ډول وپېژنو: لړۍ~\LR{$\langle a_0, a_1, a_2, \dots, a_k \rangle$} له دې عدد سره سمون لري:
\[
p_0^{a_0+1} \cdot p_1^{a_1+1} \cdot p_2^{a_2+1} \cdot \dots \cdot p_k^{a_k+1}.
\]
توانونو ته يو ځکه زياتوو چې مثلاً لړۍ~\LR{$\langle 2, 7, 3\rangle$} او \LR{$\langle 2, 7, 3, 0, 0 \rangle$} بېل عددي کوډونه ولري۔ \LR{$0$} او~\LR{$1$} دواړه د تشې لړۍ کوډونه نيولے شو؛ د کره‌والي دپاره دې~\LR{$\emptyseq$} د~\LR{$0$} ښودنه وکړي۔
\olpseganchor{OLP-0220-B005}{end}

\olpseganchor{OLP-0220-B006}{start}
د لړيو دا کوډول د حساب د تش په نامه بنسټيزې قضيې له امله کار کوي: هر طبيعي عدد~\LR{$n \ge 2$} په پوره يوه طريقه په لاندې بڼه ليکل کېدے شي:
\[
n = p_0^{a_0} \cdot p_1^{a_1} \cdot \dots \cdot p_k^{a_k}
\]
چې~\LR{$a_k \ge 1$} وي۔ دا ضمانت کوي چې نقشه~\LR{$\tuple{}(a_0, \dots, a_k) = \tuple{a_0, \dots, a_k}$} يو په يو ده: بېلې لړۍ بېلو عددونو ته لېږل کېږي، او له هر عدد سره تر ټولو زياته يوه لړۍ سمون لري۔
\olpseganchor{OLP-0220-B006}{end}

\olpseganchor{OLP-0220-B007}{start}
اوس به وښيو چې د يوې لړۍ د اوږدوالي موندل، د هغې د~\LR{$i$}م غړي موندل، د لړۍ په پاے کښې د يوه غړي زياتول، او د دوو لړيو نښلول ټولې بنسټيزې بازګشتي چارې دي۔
\olpseganchor{OLP-0220-B007}{end}

\olpseganchor{OLP-0220-B008}{start}
\begin{prop}
  تابع~\LR{$\len{s}$}، چې د لړۍ~\LR{$s$} اوږدوالی راګرځوي، بنسټيزه بازګشتي ده۔
\end{prop}
\olpseganchor{OLP-0220-B008}{end}

\olpseganchor{OLP-0220-B009}{start}
\begin{proof}
  اړيکه~\LR{$R(i, s)$} داسې تعريف کړئ:
  \[
  R(i, s) \text{\RL{ هله او يوازې هله چې }} p_i \mid s \land p_{i+1} \nmid s.
  \]
  \LR{$R$} په څرګند ډول بنسټيزه بازګشتي ده۔ کله چې~\LR{$s$} د يوې ناتشې لړۍ کوډ وي، يعنې
  \[
  s = p_0^{a_0+1} \cdot \dots \cdot p_{k}^{a_{k}+1},
  \]
  \LR{$R(i,s)$} هله صدق کوي چې~\LR{$p_i$} تر ټولو لوے اوليه عدد وي چې \LR{$p_i \mid s$}، يعنې~\LR{$i = k$} وي۔ نو د~\LR{$s$} اوږدوالی هله او يوازې هله~\LR{$i+1$} دے چې~\LR{$p_i$} تر ټولو لوے اوليه عدد وي چې~\LR{$s$} وېشي؛ ځکه دا تعريف ورکولے شو:
  \[
  \len{s} =
  \begin{cases}
    0 & \text{\RL{که \LR{$s = 0$} يا \LR{$s = 1$} وي}} \\
    1 + \bmin{i < s}{R(i, s)} & \text{\RL{که نۀ}}
  \end{cases}
  \]
  دلته محدوده کمينه‌موندنه کارولے شو، ځکه کله چې~\LR{$s$} د يوې لړۍ کوډ وي، يوازې يو~\LR{$i$} اړيکه~\LR{$R(i,s)$} پوره کوي، او که~\LR{$i$} موجود وي نو پخپله له~\LR{$s$} څخه کوچنے دے۔
\end{proof}
\olpseganchor{OLP-0220-B009}{end}

\olpseganchor{OLP-0220-B010}{start}
\begin{prop}
  تابع~\LR{$\fn{append}(s,a)$}، چې د لړۍ~\LR{$s$} په پاے کښې د~\LR{$a$} د ورزياتولو پايله راګرځوي، بنسټيزه بازګشتي ده۔
\end{prop}
\olpseganchor{OLP-0220-B010}{end}

\olpseganchor{OLP-0220-B011}{start}
\begin{proof}
  \LR{$\fn{append}$} داسې تعريفېدے شي:
  \[
  \fn{append}(s,a) =
  \begin{cases}
    2^{a+1} & \text{\RL{که \LR{$s = 0$} يا \LR{$s = 1$} وي}} \\
    s \cdot p_{\len{s}}^{a+1} & \text{\RL{که نۀ۔}}
  \end{cases}
  \]
\end{proof}
\olpseganchor{OLP-0220-B011}{end}

\olpseganchor{OLP-0220-B012}{start}
\begin{prop}
  تابع~\LR{$\fn{element}(s,i)$}، چې د~\LR{$s$} \LR{$i$}م غړے راګرځوي (لومړني غړي ته~\LR{$0$}م وايي)، او که~\LR{$i$} د~\LR{$s$} له اوږدوالي څخه لوے يا ورسره برابر وي نو~\LR{$0$} راګرځوي، بنسټيزه بازګشتي ده۔
\end{prop}
\olpseganchor{OLP-0220-B012}{end}

\olpseganchor{OLP-0220-B013}{start}
\begin{proof}
پام وکړئ چې~\LR{$a$} هله او يوازې هله د~\LR{$s$} \LR{$i$}م غړے دے چې \LR{$p_i^{a+1}$} د~\LR{$p_i$} تر ټولو لوے توان وي چې~\LR{$s$} وېشي؛ يعنې \LR{$p_i^{a+1} \mid s$} وي خو~\LR{$p_i^{a+2} \nmid s$} وي۔ نو:
  \[
  \fn{element}(s,i) =
  \begin{cases}
    0 & \mbox{که \LR{$i \geq \len{s}$} وي} \\
    \bmin{a < s}{(p_i^{a+2} \nmid s)} & \text{\RL{که نۀ۔}}
  \end{cases}
  \]
\end{proof}
\olpseganchor{OLP-0220-B013}{end}

\olpseganchor{OLP-0220-B014}{start}
د پورته تعريف شوو تابعو د رسمي نومونو پر ځاے لا لنډې نښې کاروو۔ د~\LR{$\fn{element}(s,i)$} پر ځاے به~\LR{$(s)_i$} کاروو، او \LR{$\tuple{s_0, \dots, s_k}$} به د لاندې عبارت لنډون وي:
\[
\fn{append}(\fn{append}(\dots \fn{append}(\emptyseq,s_0) \dots),s_k).
\]
پام وکړئ که~\LR{$s$} د~\LR{$k$} اوږدوالی ولري، نو د~\LR{$s$} غړي \LR{$(s)_0$}، \dots،~\LR{$(s)_{k-1}$} دي۔
\olpseganchor{OLP-0220-B014}{end}

\olpseganchor{OLP-0220-B015}{start}
\begin{prop}
تابع~\LR{$\fn{concat}(s,t)$}، چې دوه لړۍ سره نښلوي، بنسټيزه بازګشتي ده۔
\end{prop}
\olpseganchor{OLP-0220-B015}{end}

\olpseganchor{OLP-0220-B016}{start}
\begin{proof}
  داسې تابع~\LR{$\fn{concat}$} غواړو چې دا خاصيت ولري:
  \[
    \fn{concat}(\tuple{a_0, \dots, a_k}, \tuple{b_0, \dots, b_l}) = \tuple{a_0, \dots, a_k, b_0, \dots, b_l}.
  \]
  مرستندويه تابع~\LR{$\fn{hconcat}(s,t,n)$} کاروو چې د~\LR{$t$} لومړۍ~\LR{$n$} نښې له~\LR{$s$} سره نښلوي۔ دا تابع په بنسټيز بازګښت داسې تعريفېدے شي:
  \begin{align*}
    \fn{hconcat}(s,t,0) & = s\\
    \fn{hconcat}(s,t,n+1) & = \fn{append}(\fn{hconcat}(s,t,n),(t)_n)
    \intertext{\RL{بيا~\LR{$\fn{concat}$} داسې تعريفولے شو:}}
    \fn{concat}(s,t) & = \fn{hconcat}(s,t,\len{t}).
  \end{align*}
\end{proof}
\olpseganchor{OLP-0220-B016}{end}

\olpseganchor{OLP-0220-B017}{start}
د~\LR{$\fn{concat}(s,t)$} پر ځاے به~\LR{$s \concat t$} ليکو۔
\olpseganchor{OLP-0220-B017}{end}

\olpseganchor{OLP-0220-B018}{start}
دا به راته ګټوره وي چې د يوې لړۍ عددي کوډ د هغې د اوږدوالي او تر ټولو لوے غړي له مخې محدود کړے شو۔ فرض کړئ~\LR{$s$} د~\LR{$k$} اوږدوالي يوه لړۍ ده او هر غړے ئې له کوم عدد~\LR{$x$} څخه کوچنے يا ورسره برابر دے۔ که لړۍ ناتشه وي، نو~\LR{$s$} تر ټولو زيات~\LR{$k$} اوليه عوامل لري، هر يو ئې تر ټولو زيات~\LR{$p_{k-1}$} دے، او هر يو د~\LR{$s$} په اوليه تجزيه کښې تر ټولو زيات د~\LR{$x+1$} توان ته پورته شوے دے۔ نو داسې تعريف کوو:
\[
\fn{sequenceBound}(x,k) =
\fn{cond}(k,1,p_{\fn{pred}(k)}^{k \cdot (x+1)})
\]
بيا د پورته بيان شوې لړۍ~\LR{$s$} عددي کوډ تر ټولو زيات \LR{$\fn{sequenceBound}(x,k)$} دے۔
د تشې لړۍ په حالت کښې حد يو دے؛ دواړه د صفر او يو کوډونه تر دغه حد لاندې يا ورسره برابر دي۔
\olpseganchor{OLP-0220-B018}{end}

\olpseganchor{OLP-0220-B019}{start}
پر لړيو د داسې حد په لرلو نوې تابعې په محدودې پلټنې تعريفولے شو۔ مثلاً~\LR{$\fn{concat}$} په محدودې پلټنې تعريفولے شو۔ يوازې د هغه څيز (د نښلول شوې لړۍ د عدد) يو بنسټيز بازګشتي \emph{ځانګړتيايي شرط} ليکل او د پلټنې حد ټاکل پکار دي۔ لاندې تعريف کار کوي:
\begin{align*}
  \fn{concat}(s,t) = {} & \bmin{v \leq \fn{sequenceBound}(s+t,\len{s} + \len{t})}{} \\
  & \quad(\len{v} = \len{s} + \len{t} \land {}\\
  & \qquad \bforall{i < \len{s}}{((v)_i = (s)_i)} \land {} \\
  & \qquad \bforall{j < \len{t}}{((v)_{\len{s}+j} = (t)_j)})
\end{align*}
د اوږدوالي له شرط سره د دواړو لړيو د غړو شرطونه په جلا توګه لګېږي۔ د تشې لومړۍ لړۍ په حالت کښې د دويمې لړۍ شرط لا هم ټول غړي ټاکي۔
\olpseganchor{OLP-0220-B019}{end}

\olpseganchor{OLP-0220-B020}{start}
\begin{prob}
وښيئ چې داسې بنسټيزه بازګشتي تابع~\LR{$\fn{sconcat}(s)$} شته چې دا خاصيت ولري:
\[
\fn{sconcat}(\tuple{s_0, \dots, s_k}) = s_0 \concat \dots \concat s_k.
\]
\end{prob}
\olpseganchor{OLP-0220-B020}{end}

\olpseganchor{OLP-0220-B021}{start}
\begin{prob}
وښيئ چې داسې بنسټيزه بازګشتي تابع~\LR{$\fn{tail}(s)$} شته چې دا خاصيت ولري:
\begin{align*}
  \fn{tail}(\emptyseq) & = 0 \text{\RL{ او}}\\
  \fn{tail}(\tuple{s_0, \dots, s_{k}}) & = \tuple{s_1, \dots, s_{k}}.
\end{align*}
\end{prob}
\olpseganchor{OLP-0220-B021}{end}

\olpseganchor{OLP-0220-B022}{start}
\begin{prop}
  \ollabel{prop:subseq}
  تابع~\LR{$\fn{subseq}(s, i, n)$} چې د~\LR{$s$} له~\LR{$i$}م غړي څخه پيلېدونکې د \LR{$n$} اوږدوالي فرعي لړۍ راګرځوي، بنسټيزه بازګشتي ده۔
\end{prop}
\olpseganchor{OLP-0220-B022}{end}

\olpseganchor{OLP-0220-B023}{start}
\begin{proof}
  تمرين۔
\end{proof}
\olpseganchor{OLP-0220-B023}{end}

\olpseganchor{OLP-0220-B024}{start}
\begin{prob}
  \olref[cmp][rec][seq]{prop:subseq} ثابت کړئ۔
\end{prob}
\olpseganchor{OLP-0220-B024}{end}

\olpseganchor{OLP-0221-B004}{start}
\olfileid{cmp}{rec}{tre}
\olsection{ونې}
\olpseganchor{OLP-0221-B004}{end}

\olpseganchor{OLP-0221-B005}{start}
ځينې وخت د ونو ښودل د طبيعي عددونو په توګه ګټور وي، کټ مټ لکه لړۍ چې په عددونو ښيو او د هغو خاصيتونه او عملونه د کوډونو پر بنسټيزو بازګشتي اړيکو او تابعو ښيو۔ د دې دپاره لړۍ او د هغو کوډونه کاروو۔ يوه ونه يا يوازې يوه غوټه (چې ښايي لېبل ولري) وي، يا يوه غوټه (چې ښايي لېبل ولري) وي چې له يو شمېر فرعي ونو سره تړلې وي۔ دې غوټې ته د ونې \emph{ريښه}، او هغو فرعي ونو ته چې ورسره تړلې وي د هغې \emph{سمدستي فرعي ونې} وايي۔
\olpseganchor{OLP-0221-B005}{end}

\olpseganchor{OLP-0221-B006}{start}
ونې په بازګشتي ډول د~\LR{$\tuple{k, d_1, \dots, d_k}$} په لړۍ کوډوو، چې \LR{$k$} د سمدستي فرعي ونو شمېر او~\LR{$d_1$}، \dots،~\LR{$d_k$} د سمدستي فرعي ونو کوډونه دي۔ که غوټې لېبلونه ولري، له سمدستي فرعي ونو وروسته ئې شاملولے شو۔ نو هغه ونه چې يوازې يوه د~\LR{$l$} لېبل لرونکې غوټه ولري، په~\LR{$\tuple{0,l}$} کوډېږي؛ او هغه ونه چې يوه د~\LR{$l_1$} لېبل لرونکې ريښه ئې له دوو يوازېنيو غوټو سره تړلې وي، چې لېبلونه ئې~\LR{$l_2$} او~\LR{$l_3$} دي، په~\LR{$\tuple{2, \tuple{0, l_2}, \tuple{0, l_3}, l_1}$} کوډېږي۔
\olpseganchor{OLP-0221-B006}{end}

\olpseganchor{OLP-0221-B007}{start}
\begin{prop}
  \ollabel{prop:subtreeseq}
  تابع~\LR{$\fn{SubtreeSeq}(t)$}، چې د داسې لړۍ کوډ راګرځوي چې غړي ئې د \LR{$t$} کوډ لرونکې ونې د ټولو فرعي ونو کوډونه دي، بنسټيزه بازګشتي ده۔
\end{prop}
\olpseganchor{OLP-0221-B007}{end}

\olpseganchor{OLP-0221-B008}{start}
\begin{proof}
  لومړے پام وکړئ چې~\LR{$\fn{ISubtrees}(t) = \fn{subseq}(t, 1, (t)_0)$} بنسټيزه بازګشتي ده او د ونې~\LR{$t$} د سمدستي فرعي ونو کوډونه راګرځوي۔ اوس مرستندويه تابع~\LR{$\fn{hSubtreeSeq}(t,n)$} داسې تعريفوو چې د ټولو هغو فرعي ونو لړۍ محاسبه کړي چې ريښې ئې د~\LR{$t$} له ريښې څخه تر ټولو زيات~\LR{$n$} څنډې لرې وي۔ د~\LR{$t$} د هغو فرعي ونو لړۍ چې له ريښې څخه تر ټولو زيات~\LR{$0$} څنډې لرې دي---يعنې پخپله د~\LR{$t$} په ريښه پيلېږي---يوازې له~\LR{$t$} څخه جوړه لړۍ ده۔ د~\LR{$n+1$} تر کچې د ټولو فرعي ونو د لړۍ د ترلاسه کولو دپاره، د~\LR{$n$} تر کچې شته فرعي ونې د هغو ټولو سمدستي فرعي ونو له لړۍ سره نښلوو چې د~\LR{$n$} تر کچې له فرعي ونو راځي۔
  د تر يوې کچې پورې لړيو په نښلون کښې د يوې فرعي ونې کوډ تکرارېدے شي۔
  د ټولو د لست د ترلاسه کولو دپاره پام وکړئ: که~\LR{$f(x)$} يوه بنسټيزه بازګشتي تابع وي چې د لړيو کوډونه راګرځوي، نو د غړو د مثبت شمېر دپاره تابع~\LR{$g_f(s,k)=f((s)_0) \concat \dots \concat f((s)_{k-1})$} هم بنسټيزه بازګشتي ده؛ د صفر غړو نښلون تشه لړۍ نيسو:
    \begin{align*}
      g(s, 0) & = \emptyseq\\
      g(s, k+1) & = g(s, k) \concat f((s)_k)
    \end{align*}
    مثلاً، که~\LR{$s$} د ونو يوه لړۍ وي، نو \LR{$h(s) = g_{\fn{ISubtrees}}(s, \len{s})$} د~\LR{$s$} د غړو د ټولو سمدستي فرعي ونو لړۍ راکوي۔
    پاراميټر د اخستل کېدونکو لومړنيو غړو شمېر دے۔ د صفر غړو نښلون تشه لړۍ ده، او هر ګام ورپسې يو غړے ورزياتوي۔
    له دې څخه~\LR{$\fn{hSubtreeSeq}$} داسې تعريفولے شو:
    \begin{align*}
      \fn{hSubtreeSeq}(t, 0) & = \tuple{t} \\
      \fn{hSubtreeSeq}(t, n+1) & = \fn{hSubtreeSeq}(t, n) \concat h(\fn{hSubtreeSeq}(t, n)).
    \end{align*}
    د~\LR{$t$} په کوډ شوې ونه کښې د فرعي ونو تر ټولو لوړه کچه، يعنې د ريښې او يوې پاڼه غوټې تر منځ تر ټولو لوے واټن، د کوډ~\LR{$t$} له مخې محدوده ده۔ نو د~\LR{$t$} په کوډ شوې ونه کښې د ټولو فرعي ونو د کوډونو لړۍ~\LR{$\fn{hSubtreeSeq}(t, t)$} ورکوي۔
\end{proof}
\olpseganchor{OLP-0221-B008}{end}

\olpseganchor{OLP-0221-B009}{start}
\begin{prob}
  د~\olref[cmp][rec][tre]{prop:subtreeseq} په ثبوت کښې د \LR{$\fn{hSubtreeSeq}$} تعريف په عمومي ډول تکرارونه راولي۔ يو داسې بل تعريف ورکړئ چې ضمانت وکړي د يوې فرعي ونې کوډ په پايله لست کښې يوازې يو ځل راځي۔
\end{prob}
\olpseganchor{OLP-0221-B009}{end}

\olpseganchor{OLP-0222-B004}{start}
\olfileid{cmp}{rec}{ore}
\olsection{نور بازګښتونه}
\olpseganchor{OLP-0222-B004}{end}

\olpseganchor{OLP-0222-B005}{start}
د جوړو او لړيو په کارولو د بنسټيز بازګښت لا نابلده (او ګټورې) بڼې توجيه کولے شو۔ مثلاً ډېر وخت دوه تابعې يوځاے تعريفول ګټور وي، لکه په لاندې تعريف کښې:
\begin{align*}
h_0(\vec x, 0) & = f_0(\vec x) \\
h_1(\vec x, 0) & = f_1(\vec x) \\
h_0(\vec x, y+1) & = g_0(\vec x, y, h_0(\vec x, y), h_1(\vec x, y)) \\
h_1(\vec x, y+1) & = g_1(\vec x, y, h_0(\vec x, y), h_1(\vec x, y))
\end{align*}
دا د \emph{هممهاله بازګښت} يوه بېلګه ده۔ د تابعو د تعريف بله ګټوره لار دا ده چې د~\LR{$h(\vec x, y+1)$} ارزښت د ټولو ارزښتونو \LR{$h(\vec x, 0)$}، \dots،~\LR{$h(\vec x, y)$} له مخې ورکړو، لکه په لاندې تعريف کښې:
\begin{align*}
h(\vec x, 0) & = f(\vec x) \\
h(\vec x, y+1) & = g(\vec x, y, \tuple{h(\vec x, 0), \dots, h(\vec x, y)}).
\end{align*}
لاندې چوکاټ دا مفکوره لا لنډه بيانوي:
\[
h(\vec x, y) = g(\vec x, y, \tuple{h(\vec x, 0), \dots, h(\vec x, y-1)})
\]
په دې پوهېدو سره چې کله~\LR{$y$}، \LR{$0$} وي، د~\LR{$g$} وروستے آرګومېنټ يوازې تشه لړۍ وي۔ په دواړو بڼو کښې مفکوره دا ده چې د «راتلونکي ګام» د محاسبې پر مهال تابع~\LR{$h$} تر دې ځايه د محاسبه شوو ارزښتونو ټوله لړۍ کارولے شي۔ دې ته \emph{د ارزښتونو د بهير بازګښت} وايي۔ د يوې ځانګړې بېلګې دپاره لاندې ډول تعريف توجيه کولے شي:
\begin{align*}
h(\vec x, y) & = \begin{cases}
  g(\vec x, y, h(\vec x, k(\vec x, y))) & \text{\RL{که \LR{$k(\vec x, y) < y$} وي}} \\
  f(\vec x) & \text{\RL{که نۀ}}
\end{cases}
\end{align*}
په بله وينا، په~\LR{$y$} د~\LR{$h$} ارزښت د~\LR{$h$} د هر مخکيني ارزښت له مخې محاسبه کېدے شي، او هغه مخکينی ارزښت~\LR{$k$} ورکوي۔
\olpseganchor{OLP-0222-B005}{end}

\olpseganchor{OLP-0222-B006}{start}
\begin{prob}
  د پاتې شوني تابع~\LR{$r(x,y)$} د ارزښتونو د بهير په بازګښت تعريف کړئ۔ (که~\LR{$x$} او~\LR{$y$} طبيعي عددونه او~\LR{$y > 0$} وي، نو~\LR{$r(x,y)$} له~\LR{$y$} څخه هغه کوچنے عدد دے چې د کوم~\LR{$z$} دپاره \LR{$x = z\times y + r(x,y)$} وي۔ د کره‌والي دپاره وايو که~\LR{$y=0$} وي، نو~\LR{$r(x,0) = 0$}۔)
\end{prob}
\olpseganchor{OLP-0222-B006}{end}

\olpseganchor{OLP-0222-B007}{start}
بايد په دې فکر وکړئ چې دا تابعې په عادي بنسټيز بازګښت څنګه ترلاسه کېږي۔ د بنسټيز بازګښت يوه وروستۍ بڼه په دې لحاظ لا انعطاف منونکې ده چې د لارې په اوږدو کښې د \emph{پاراميټرونو} (اړخيزو ارزښتونو) بدلولو اجازه ورکوي:
\begin{align*}
h(\vec x, 0) & = f(\vec x) \\
h(\vec x, y+1) & = g(\vec x, y, h(k(\vec x), y))
\end{align*}
دا هم په عادي بنسټيز بازګښت تقليدېدے شي۔ (دا کول ستونزمن دي۔ د يوې لارښوونې دپاره، محاسبه په لاس له پاے څخه بېرته پرانيزئ۔)
\olpseganchor{OLP-0222-B007}{end}

\olpseganchor{OLP-0223-B004}{start}
\olfileid{cmp}{rec}{npr}
\olsection{نابنسټيزې بازګشتي تابعې}
\olpseganchor{OLP-0223-B004}{end}

\olpseganchor{OLP-0223-B005}{start}
بنسټيزې بازګشتي تابعې ټولې هغه تابعې نۀ رانغاړي چې په شهودي ډول محاسبه کېدونکې دي۔ په شهودي ډول بايد څرګنده وي چې د ټولو يوځاييزو بنسټيزو بازګشتي تابعو يو لست~\LR{$f_0$}، \LR{$f_1$}، \LR{$f_2$}،~\dots جوړولے شو، داسې چې د~\LR{$y$} پر ننوت د~\LR{$f_x$} ارزښت په اغېزمنه توګه محاسبه کړے شو؛ په بله وينا، لاندې تعريف شوې تابع~\LR{$g(x,y)$} محاسبه کېدونکې ده:
\[
g(x,y) = f_x(y)
\]
خو بيا لاندې تابع هم محاسبه کېدونکې ده:
\begin{eqnarray*}
h(x) & = & g(x,x) + 1 \\
& = & f_x(x) +1.
\end{eqnarray*}
د هرې بنسټيزې بازګشتي تابع~\LR{$f_i$} دپاره د~\LR{$h$} او~\LR{$f_i$} ارزښتونه په~\LR{$i$} کښې سره بېل دي۔ نو~\LR{$h$} محاسبه کېدونکې خو بنسټيزه بازګشتي نۀ ده؛ د~\LR{$g$} په باب هم همدا ويلے شو۔ دا د کانتور د قطري‌کولو د استدلال يوه «اغېزمنه» بڼه ده۔
\olpseganchor{OLP-0223-B005}{end}

\olpseganchor{OLP-0223-B006}{start}
د هغو محاسبه کېدونکو تابعو لا څرګندې بېلګې هم ورکولے شو چې بنسټيزې بازګشتي نۀ دي۔ مثلاً نښه~\LR{$g^n(x)$} دې~\LR{$g(g(\dots g(x)))$} وښيي، چې ټولې~\LR{$n$} دانې~\LR{$g$}ګانې لري؛ او د تابعو لړۍ~\LR{$g_0,g_1,\dots$} داسې تعريف کړئ:
\begin{eqnarray*}
g_0(x) & = & x+1 \\
g_{n + 1}(x) & = & g_n^x(x)
\end{eqnarray*}
تاسو ارزولے شئ چې هره تابع~\LR{$g_n$} بنسټيزه بازګشتي ده۔ هره ورپسې تابع له مخکينۍ څخه ډېره چټکه لوېږي؛ \LR{$g_1(x)$} له~\LR{$2x$} سره برابره ده، \LR{$g_2(x)$} له~\LR{$2^x \cdot x$} سره برابره ده، او~\LR{$g_3(x)$} نژدې د~\LR{$x$} دانو~\LR{$2$}ګانو د تواني برج په شان لوېږي۔ د اکرمن--پېټر تابع په اصل کښې همدا تابع~\LR{$G(x) = g_x(x)$} ده، او ښودل کېدے شي چې دا له هرې بنسټيزې بازګشتي تابع څخه چټکه لوېږي۔
\olpseganchor{OLP-0223-B006}{end}

\olpseganchor{OLP-0223-B007}{start}
راځئ د بنسټيزو بازګشتي تابعو د شمېرنې مسئلې ته بېرته راشو۔ په ياد ولرئ چې هرې بنسټيزې بازګشتي تابع ته مو سمبوليکه نښه ټاکلې ده؛ نو د نښو شمېرل بس دي۔ هرې نښې~\LR{$F$} ته يو طبيعي عدد~\LR{$\#(F)$} په بازګشتي ډول داسې ټاکلے شو:
\begin{eqnarray*}
\#(\Zero) & = & \langle 0 \rangle \\
\#(\Succ) & = & \langle 1 \rangle \\
\#(\Proj{n}{i}) & = & \langle 2, n, i \rangle \\
\#(\fn{Comp}_{k,l}[H,G_0,\dots,G_{k-1}]) & = & \langle 3,k,l,\#(H),\#(G_0),\dots,\#(G_{k-1}) \rangle \\
\#(\fn{Rec}_l[G,H]) & = & \langle 4, l, \#(G), \#(H) \rangle
\end{eqnarray*}
دلته دا حقيقت کاروو چې د عددونو هره لړۍ د لړيو د مخکيني کوډ له لارې د يوه طبيعي عدد په توګه ګڼل کېدے شي۔ پايله دا ده چې هرې نښې ته يو طبيعي عددي کوډ ټاکل شوے دے۔ البته ځينې لړۍ (او ځکه ځينې عددونه) له هېڅ نښې سره سمون نۀ لري؛ خو که~\LR{$i$} د يوې يوځاييزې بنسټيزې بازګشتي تابع د نښې کوډ وي، \LR{$f_i$} هماغه تابع نيسو، او که~\LR{$i$} داسې نښه نۀ کوډوي نو د ثابتې صفر تابع ئې نيسو۔ وروستۍ پايله دا ده چې د يوځاييزو بنسټيزو بازګشتي تابعو د شمېرلو يوه څرګنده طريقه لرو۔
\olpseganchor{OLP-0223-B007}{end}

\olpseganchor{OLP-0223-B008}{start}
(په حقيقت کښې ځينې تابعې، لکه ثابته~\LR{$0$} تابع، په لست کښې تر يو ځل زياتې راځي۔ دا يوازې زموږ د کوډولو مصنوعي پايله نۀ ده، بلکې د دې حقيقت پايله هم ده چې ثابته صفر تابع له يوې څخه زياتې نښې لري۔ وروسته به ووينو چې دا تکرارونه په محاسبه کېدونکې توګه نۀ شي مخنيوے کېدے؛ مثلاً هېڅ داسې محاسبه کېدونکې تابع نشته چې پرېکړه وکړي يوه ورکړل شوې نښه ثابته صفر تابع ښيي که نۀ۔)
\olpseganchor{OLP-0223-B008}{end}

\olpseganchor{OLP-0223-B009}{start}
اوس تابع~\LR{$g(x,y)$} د~\LR{$f_x(y)$} په توګه اخلو، چې~\LR{$f_x$} هماغې شمېرنې ته اشاره کوي چې اوس مو بيان کړه۔ څنګه پوهېږو چې~\LR{$g(x,y)$} محاسبه کېدونکې ده؟ په شهودي ډول خبره څرګنده ده: د~\LR{$g(x,y)$} د محاسبې دپاره لومړے~\LR{$x$} «پرانيزئ» او وګورئ چې ايا دا د يوې يوځاييزې تابع نښه ده که نۀ۔ که وي، پر ننوت~\LR{$y$} د هغې تابع ارزښت محاسبه کړئ۔ که نۀ وي، د شمېرنې د تعريف له مخې صفر ورکړئ۔
\olpseganchor{OLP-0223-B009}{end}

\olpseganchor{OLP-0223-B010}{start}

\begin{digress}
کېدے شي لا مو منلې وي چې (په لږ کار سره!) داسې يو پروګرام، مثلاً په جاوا يا~C++ کښې، ليکلے شو چې دا کار وکړي؛ بيا د چرچ--ټيورينګ اصل ته مراجعه کولے شو، چې وايي هر هغه څه چې په شهودي ډول محاسبه کېدونکې وي، په ټيورينګ ماشين محاسبه کېدے شي۔
\olpseganchor{OLP-0223-B010}{end}

\olpseganchor{OLP-0223-B011}{start}
البته د~\LR{$g(x,y)$} د محاسبه کېدو د ښودلو لا نېغه لار دا ده چې هغه ټيورينګ ماشين په څرګند ډول بيان کړو چې دا تابع محاسبه کوي۔ په ځانګړي ډول، دا به د چرچ--ټيورينګ اصل او شهود ته د مراجعې مخه ونيسي۔ ډېر ژر به دومره وسايل جوړ کړي وي چې د~\LR{$g(x,y)$} محاسبه کېدل د محاسبې د داسې مدل په کارولو وښيو چې پر ټيورينګ ماشين \emph{تقليدېدے} شي: يعنې بازګشتي تابعې۔
\end{digress}

\olpseganchor{OLP-0223-B011}{end}

\olpseganchor{OLP-0224-B004}{start}
\olfileid{cmp}{rec}{par}
\olsection{جزوي بازګشتي تابعې}
\olpseganchor{OLP-0224-B004}{end}


\olpseganchor{OLP-0224-B005}{start}
د بازګشتي تابعو د تعريف د انګېزې دپاره پام وکړئ چې زموږ هغه ثبوت، چې
محاسبه کېدونکې خو نابنسټيزې بازګشتي تابعې شته، په حقيقت کښې تر دې ډېر
څه ثابتوي۔ استدلال ساده و: موږ يوازې دا حقيقت وکاراوه چې تابعې
\LR{$f_0,f_1,\dots$} په داسې ډول په لړ کښې راوستل کېدے شي چې د~\LR{$x$} او~\LR{$y$}
د تابع په توګه \LR{$f_x(y)$} محاسبه کېدونکې وي۔ نو دا استدلال د تابعو پر
\emph{هر هغه ټولګي عملي کېږي چې په دې ډول په لړ کښې راوستل کېدے شي}۔
له دې سره له يوې ستونزې سره مخ کېږو: غواړو محاسبه کېدونکې تابعې په
څرګند ډول بيان کړو؛ خو د محاسبه کېدونکو تابعو د يوې ټولګې هېڅ څرګند
بيان بشپړ کېدے نۀ شي!
\olpseganchor{OLP-0224-B005}{end}

\olpseganchor{OLP-0224-B006}{start}
د وتلو لار دا ده چې \emph{جزوي} تابعو ته هم ځای ورکړو۔ وبه وينو چې
جزوي محاسبه کېدونکې تابعې \emph{په رښتيا} په لړ کښې راوستل کېدے شي۔
په حقيقت کښې نږدې له وړاندې پوهېږو چې داسې ده، ځکه ټيورينګ
  ماشينونه په منظم ډول په لړ کښې راوستل کېدے شي۔ وروسته به خپل قطري
استدلال ته راوګرځو او وبه څېړو چې ولې د جزوي تابعو په ګډون هغه استدلال
مخ ته نۀ ځي۔
\olpseganchor{OLP-0224-B006}{end}

\olpseganchor{OLP-0224-B007}{start}
اوس پوښتنه دا ده: بنسټيزو بازګشتي تابعو ته څه ورزيات کړو چې ټولې جزوي
بازګشتي تابعې تر لاسه کړو؟ دوه کارونه پکار دي:
\begin{enumerate}
\item د بنسټيزو بازګشتي تابعو خپل تعريف داسې بدل کړو چې جزوي تابعو ته
  هم ځای ورکړي۔
\item تعريف ته يو څه \emph{ورزيات} کړو، څو ځينې نوې جزوي تابعې هم پکې
  شاملې شي۔
\end{enumerate}
\olpseganchor{OLP-0224-B007}{end}

\olpseganchor{OLP-0224-B008}{start}
لومړے کار اسان دے۔ د پخوا په شان له صفر، جانشين او پروجکشنونو څخه پيل
کوو، او ټولګے د ترکيب او بنسټيز بازګښت لاندې بندوو۔ يوازينے توپير دا
دے چې د ترکيب او بنسټيز بازګښت تعريفونه بايد د دې امکان دپاره بدل کړو
چې په تعريف کښې ځينې ترمونه تعريف شوي نۀ وي۔ که \LR{$f$} او~\LR{$g$} جزوي تابعې
وي، \LR{$f(x) \fdefined$} به په دې مانا ليکو چې \LR{$f$} پر~\LR{$x$} تعريف شوې ده،
يعنې \LR{$x$} د~\LR{$f$} په تعريف ساحه کښې دے؛ او \LR{$f(x) \fundefined$} به د دې
مخالفې مانا، يعنې دا چې \LR{$f$} پر~\LR{$x$} تعريف شوې نۀ ده، څرګندوي۔
\LR{$f(x) \simeq g(x)$} به په دې مانا کاروو چې يا \LR{$f(x)$} او~\LR{$g(x)$} دواړه
تعريف شوي نۀ دي، يا دواړه تعريف شوي او سره برابر دي۔ دا نښې به د لا
پېچلو ترمونو دپاره هم کاروو۔ دا دود به منو چې که \LR{$h$} او \LR{$g_0$}،
\dots،~\LR{$g_k$} ټولې جزوي تابعې وي، نو
\[
h(g_0(\vec x),\dots,g_k(\vec x))
\]
هله او يوازې هله تعريف شوے دے چې هر~\LR{$g_i$} پر~\LR{$\vec x$} تعريف شوے وي،
او \LR{$h$} پر \LR{$g_0(\vec x)$}، \dots،~\LR{$g_k(\vec x)$} تعريف شوې وي۔ په دې
پوهه د جزوي تابعو د ترکيب او بنسټيز بازګښت تعريفونه هماغه پورته
تعريفونه دي، خو ``\LR{$=$}'' بايد پر ``\LR{$\simeq$}'' بدل کړو۔
\olpseganchor{OLP-0224-B008}{end}

\olpseganchor{OLP-0224-B009}{start}
بنسټيزو بازګشتي تابعو ته چې څه ورزياتوو څو جزوي تابعې تر لاسه کړو،
هغه د \emph{بې‌حده پلټنې عامل} دے۔ که \LR{$f(x,\vec z)$} پر طبيعي عددونو
هره جزوي تابع وي، \LR{$\mu x \; f(x,\vec z)$} داسې تعريف کړئ:
\begin{quote}
  تر ټولو وړوکے \LR{$x$} چې \LR{$f(0,\vec z), f(1,\vec z), \dots, f(x,\vec
  z)$} ټول تعريف شوي وي، او \LR{$f(x,\vec z) = 0$} وي، که داسې~\LR{$x$} موجود وي؛
\end{quote}
او که داسې نۀ وي، \LR{$\mu x \; f(x,\vec z)$} تعريف شوے نۀ دے۔ دا
\LR{$\mu x \; f(x,\vec z)$} په يکتا ډول تعريفوي۔
\olpseganchor{OLP-0224-B009}{end}

\olpseganchor{OLP-0224-B010}{start}
\begin{explain}
پام وکړئ چې زموږ تعريف ټيورينګ ماشينونو، يا الګوريتمونو
يا د محاسبې کوم ځانګړي مدل ته هېڅ اشاره نۀ کوي۔ خو د ترکيب او بنسټيز
بازګښت په شان، د بې‌حده پلټنې تر شا هم عملياتي، محاسبوي شهود شته۔ د
يوې جزوي تابع د محاسبه کېدو په بحث کښې هغه آرګومېنټونه، چې تابع پرې
تعريف شوې نۀ وي، له هغو ننوتونو سره سمون خوري چې محاسبه پرې نۀ درېږي۔
د~\LR{$\mu x \; f(x,\vec z)$} د محاسبې کړنلار داسې ده:
\LR{$f(0,\vec z), f(1,\vec z), f(2,\vec z)$} او ورپسې قيمتونه محاسبه کړئ، تر څو د
صفر قيمت راوګرځي۔ خو که له منځنيو محاسبو څخه کومه يوه ونۀ درېږي، د
\LR{$\mu x \; f(x,\vec z)$} محاسبه هم نۀ درېږي۔
\end{explain}
\olpseganchor{OLP-0224-B010}{end}

\olpseganchor{OLP-0224-B011}{start}
که \LR{$R(x,\vec z)$} هره اړيکه وي، \LR{$\mu x \; R(x,\vec z)$} داسې تعريفېږي:
\LR{$\mu x \; (1 \tsub \Char{R}(x,\vec z))$}۔ په بله وينا، \LR{$\mu x \;
R(x,\vec z)$} د~\LR{$x$} تر ټولو وړوکے قيمت راګرځوي چې \LR{$R(x,\vec z)$} پوره
کوي۔ نو که \LR{$f(x,\vec z)$} هرځاے تعريف شوې تابع وي، \LR{$\mu x \;
f(x,\vec z)$} له \LR{$\mu x \; (f(x,\vec z) = 0)$} سره يو شان دے۔ خو پام
وکړئ چې زموږ لومړنی تعريف لا عمومي دے، ځکه دا امکان هم پرېږدي چې
\LR{$f(x,\vec z)$} په هر ځای کښې تعريف شوې نۀ وي (خو د دې پر خلاف د يوې
اړيکې ځانګړونکې تابع تل هرځاے تعريف شوې وي)۔
\olpseganchor{OLP-0224-B011}{end}

\olpseganchor{OLP-0224-B012}{start}
\begin{defn}
د \emph{جزوي بازګشتي تابعو} سټ له طبيعي عددونو څخه طبيعي عددونو ته د
بېلابېلو ځاييزو کچو د جزوي تابعو تر ټولو کوچنے سټ دے چې صفر، جانشين او
پروجکشنونه لري، او د ترکيب، بنسټيز بازګښت او بې‌حده پلټنې لاندې بند دے۔
\end{defn}
\olpseganchor{OLP-0224-B012}{end}

\olpseganchor{OLP-0224-B013}{start}
البته ځينې جزوي بازګشتي تابعې به په اتفاقي ډول هرځاے تعريف شوې وي،
يعنې د هر آرګومېنټ دپاره تعريف شوې وي۔
\olpseganchor{OLP-0224-B013}{end}

\olpseganchor{OLP-0224-B014}{start}
\begin{defn}
\ollabel{defn:recursive-fn}
د \emph{بازګشتي تابعو} سټ د هغو جزوي بازګشتي تابعو سټ دے چې هرځاے
تعريف شوې وي۔
\end{defn}
\olpseganchor{OLP-0224-B014}{end}

\olpseganchor{OLP-0224-B015}{start}
پر دې د ټينګار دپاره چې يوه بازګشتي تابع په هر ځای کښې تعريف شوې ده،
کله کله ورته ``هرځاے تعريف شوې بازګشتي'' تابع هم ويل کېږي۔
\olpseganchor{OLP-0224-B015}{end}

\olpseganchor{OLP-0225-B004}{start}
\olfileid{cmp}{rec}{nft}
\olsection{د نورمال شکل قضيه}
\olpseganchor{OLP-0225-B004}{end}

\olpseganchor{OLP-0225-B005}{start}
\begin{thm}[د کليني د نورمال شکل قضيه]
\ollabel{thm:kleene-nf}
يوه بنسټيزه بازګشتي اړيکه \LR{$T(e, x, s)$} او يوه بنسټيزه بازګشتي تابع
\LR{$U(s)$} شته چې دا خاصيت لري: که \LR{$f$} هره جزوي بازګشتي تابع وي، نو د
کوم~\LR{$e$} دپاره
\[
f(x) \simeq U(\umin{s}{T(e, x, s)})
\]
د هر~\LR{$x$} دپاره پوره کېږي۔
\end{thm}
\olpseganchor{OLP-0225-B005}{end}

\olpseganchor{OLP-0225-B006}{start}
\begin{explain}
د نورمال شکل د قضيې ثبوت اوږد دے، خو بنسټيزه مفکوره ئې ساده ده۔ هره
جزوي بازګشتي تابع يو \emph{شاخص}~\LR{$e$} لري؛ په شهودي ډول دا هغه عدد دے
چې د تابع پروګرام يا تعريف کوډوي۔ که \LR{$f(x) \fdefined$} وي، محاسبه په
منظم ډول ثبتېدے او په کوم عدد~\LR{$s$} کوډېدے شي، او دا حقيقت، چې \LR{$s$} پر
ننوت~\LR{$x$} د~\LR{$f$} محاسبه کوډوي، يوازې د~\LR{$x$} او تعريف~\LR{$e$} په کارولو سره
په بنسټيز بازګشتي ډول کتل کېدے شي۔ نو اړيکه~\LR{$T$}، يعنې ``د شاخص~\LR{$e$}
تابع د ننوت~\LR{$x$} دپاره يوه محاسبه لري، او \LR{$s$} دا محاسبه کوډوي''،
بنسټيزه بازګشتي ده۔ د محاسبې د بشپړ ريکارډ~\LR{$s$} په ورکړې سره د~\LR{$s$}
``پايله'' د~\LR{$f(x)$} قيمت دے، او هغه هم له~\LR{$s$} څخه په بنسټيز بازګشتي
ډول تر لاسه کېدے شي۔
\olpseganchor{OLP-0225-B006}{end}

\olpseganchor{OLP-0225-B007}{start}
د نورمال شکل قضيه ښيي چې د هرې جزوي بازګشتي تابع د تعريف دپاره يوازې
يوه بې‌حده پلټنه پکار ده۔ په بنسټيز ډول د ټولو عددونو له منځه پلټنه
کوو تر څو داسې عدد ومومو چې د شاخص~\LR{$e$} د تابع محاسبه پر ننوت~\LR{$x$}
کوډوي۔ عددونه~\LR{$e$} د جزوي بازګشتي تابعو د ``نومونو'' په توګه کارولے
شو، او د قضيې په معادله تعريف شوې تابع~\LR{$f$} ته \LR{$\cfind{e}$} ليکلے شو۔
پام وکړئ چې يوه جزوي بازګشتي تابع تر يوه ډېر شاخص لرلے شي---په حقيقت
کښې هره جزوي بازګشتي تابع بې‌شمېره ډېر شاخصونه لري۔
\end{explain}
\olpseganchor{OLP-0225-B007}{end}

\olpseganchor{OLP-0226-B004}{start}
\olfileid{cmp}{rec}{hlt}
\olsection{د درېدنې مسئله}
\olpseganchor{OLP-0226-B004}{end}

\olpseganchor{OLP-0226-B005}{start}
په عمومي ډول د \emph{درېدنې مسئله} دا ده چې د يوې محاسبه کېدونکې
تابع د مشخصې~\LR{$e$} (لکه پروګرام) او يو عدد~\LR{$n$} په ورکړې سره پرېکړه وشي
چې د دې تابع محاسبه پر ننوت~\LR{$n$} درېږي، يعنې کومه پايله ورکوي، که نه۔
اېلن ټيورينګ په نوميالي ډول ثابته کړه چې دا مسئله پخپله د يوې محاسبه
کېدونکې تابع له لارې نۀ شي حل کېدے؛ يعنې تابع
\[
h(e, n) =
\begin{cases}
1 & \text{\RL{که محاسبه \LR{$e$} پر ننوت \LR{$n$} ودرېږي}}\\
0 & \text{\RL{که داسې نۀ وي،}}
\end{cases}
\]
محاسبه کېدونکې نۀ ده۔
\olpseganchor{OLP-0226-B005}{end}

\olpseganchor{OLP-0226-B006}{start}
د جزوي بازګشتي تابعو په چوکاټ کښې د يو پروګرام د مشخصې رول د کليني
د نورمال شکل په قضيه کښې ورکړل شوے شاخص~\LR{$e$} لوبولے شي۔ که \LR{$f$} جزوي
بازګشتي تابع وي، هر هغه~\LR{$e$} چې د نورمال شکل د قضيې معادله پوره کوي،
د~\LR{$f$} يو شاخص دے۔ د عدد~\LR{$e$} په ورکړې سره د نورمال شکل قضيه وايي چې
\[
\cfind{e}(x) \simeq U(\mu s \; T(e, x, s))
\]
جزوي بازګشتي ده، او د هرې جزوي بازګشتي \LR{$f\colon \Nat \to \Nat$}
دپاره يو \LR{$e \in \Nat$} شته چې د ټولو~\LR{$x \in \Nat$} دپاره
\LR{$\cfind{e}(x) \simeq f(x)$} وي۔ په حقيقت کښې د هرې داسې~\LR{$f$} دپاره
يوازې يو نه، بلکې بې‌شمېره ډېر داسې~\LR{$e$} شته۔ \emph{د درېدنې تابع}~\LR{$h$}
داسې تعريفېږي:
\[
h(e, x) =
\begin{cases}
1 & \text{\RL{که \LR{$\cfind{e}(x) \fdefined$} وي}}\\
0 & \text{\RL{که داسې نۀ وي۔}}
\end{cases}
\]
نو \LR{$h(e,x)=0$} هله او يوازې هله چې \LR{$\cfind{e}(x) \fundefined$} وي؛ د
نورمال شکل په ټاکلې شمېرنه کښې هر طبيعي عدد~\LR{$e$} د کومې جزوي بازګشتي
تابع شاخص دے، که هغه تابع په هر ننوت هم تعريف شوې نۀ وي۔
\olpseganchor{OLP-0226-B006}{end}

\olpseganchor{OLP-0226-B007}{start}
\begin{thm}
\ollabel{thm:halting-problem}
د درېدنې تابع~\LR{$h$} جزوي بازګشتي نۀ ده۔
\end{thm}
\olpseganchor{OLP-0226-B007}{end}

\olpseganchor{OLP-0226-B008}{start}
\begin{proof}
که \LR{$h$} جزوي بازګشتي وای، دا تابع مو تعريفولے شوه:
\[
d(y) =
\begin{cases}
1 & \text{\RL{که \LR{$h(y, y) = 0$} وي}}\\
\umin{x}{x \neq x} & \text{\RL{که داسې نۀ وي۔}}
\end{cases}
\]
څرنګه چې هېڅ عدد~\LR{$x$} شرط \LR{$x \neq x$} نۀ پوره کوي،
\LR{$\umin{x}{x \neq x}$} نشته؛ نو \LR{$d(y) \fundefined$} هله او يوازې هله
چې \LR{$h(y,y) \neq 0$} وي۔ له دې تعريف څخه دا پايلې راځي:
\begin{enumerate}
\item د شاخص~\LR{$y$} دپاره، \LR{$d(y) \fdefined$} هله او يوازې هله چې
  \LR{$\cfind{y}(y) \fundefined$} وي۔
\item \LR{$d(y) \fundefined$} هله او يوازې هله چې \LR{$\cfind{y}(y)
  \fdefined$} وي۔
\end{enumerate}
که \LR{$h$} جزوي بازګشتي وای، \LR{$d$} به هم جزوي بازګشتي وه۔ نو د کليني د
نورمال شکل د قضيې له مخې دا يو شاخص~\LR{$e_d$} لري۔ د~\LR{$h(e_d,e_d)$} قيمت
وګورئ۔ دوه ممکن حالتونه شته، \LR{$0$} او~\LR{$1$}۔
\begin{enumerate}
\item که \LR{$h(e_d, e_d) = 1$} وي، نو \LR{$\cfind{e_d}(e_d) \fdefined$} دے۔
  خو \LR{$\cfind{e_d} \simeq d$}، او \LR{$d(e_d)$} هله تعريف شوے دے چې
  \LR{$h(e_d,e_d)=0$} وي۔ نو \LR{$h(e_d,e_d) \neq 1$}۔
\item که \LR{$h(e_d, e_d) = 0$} وي، نو د شاخص~\LR{$e_d$} تابع
  \LR{$\cfind{e_d}(e_d) \fundefined$} ده۔ خو بيا هم
  \LR{$\cfind{e_d} \simeq d$}، حال دا چې \LR{$d(e_d)$} تعريف شوے دے؛ دغه
  مساوات به بيا \LR{$\cfind{e_d}(e_d) \fdefined$} غوښتل۔
\end{enumerate}
د شاخص~\LR{$e_d$} په دواړو ممکنو حالتونو کښې تناقض راځي۔ که \LR{$h$} جزوي بازګشتي وای، \LR{$d$} به هم
جزوي بازګشتي وه او د هغې شاخص~\LR{$e_d$} به منل شوی وای؛ خو هېڅ ممکن قيمت
د دغه شاخص له تعريف سره نۀ جوړېږي۔ نو بايد دې پايلې ته ورسېږو چې \LR{$h$}
جزوي بازګشتي کېدے نۀ شي۔
\end{proof}
\olpseganchor{OLP-0226-B008}{end}

\olpseganchor{OLP-0227-B004}{start}
\olfileid{cmp}{rec}{gen}
\olsection{عمومي بازګشتي تابعې}
\olpseganchor{OLP-0227-B004}{end}

\olpseganchor{OLP-0227-B005}{start}
د هرځاے تعريف شويو تابعو د يو سټ د تر لاسه کولو بله لار هم شته۔
هرځاے تعريف شوې تابع \LR{$f(x,\vec z)$} هغه وخت \emph{منظمه} وبولئ چې د
طبيعي عددونو د هرې لړۍ~\LR{$\vec z$} دپاره يو~\LR{$x$} موجود وي چې
\LR{$f(x,\vec z)=0$} وي۔ په بله وينا، منظمې تابعې کټ مټ هغه تابعې دي چې
بې‌حده پلټنه پرې عملي کېدے شي او په پايله کښې يوه هرځاے تعريف شوې
تابع تر لاسه کېږي۔ بې‌حده پلټنه په محتاط ډول يوازې منظمو تابعو ته
محدودولے شو:
\olpseganchor{OLP-0227-B005}{end}

\olpseganchor{OLP-0227-B006}{start}
\begin{defn}
\ollabel{defn:general-recursive}
د \emph{عمومي بازګشتي تابعو} سټ له طبيعي عددونو څخه طبيعي عددونو ته د
بېلابېلو ځاييزو کچو د تابعو تر ټولو کوچنے سټ دے چې صفر، جانشين او
پروجکشنونه لري، او د ترکيب، بنسټيز بازګښت او پر \emph{منظمو} تابعو د
عملي شوې بې‌حده پلټنې لاندې بند دے۔
\end{defn}
\olpseganchor{OLP-0227-B006}{end}

\olpseganchor{OLP-0227-B007}{start}
په څرګند ډول هره عمومي بازګشتي تابع هرځاے تعريف شوې ده۔ د
\olref{defn:general-recursive} او \olref[par]{defn:recursive-fn} تر
منځ توپير دا دے چې په وروستي تعريف کښې د جوړونې په منځ کښې د جزوي
بازګشتي تابعو کارول هم روا دي؛ يوازينے شرط دا دے چې په پاے کښې تر لاسه
شوې تابع هرځاے تعريف شوې وي۔ نو د ``عمومي'' ټکی، چې يو تاريخي پاتې
شونی دے، سم نوم نۀ دے؛ په ظاهره \olref{defn:general-recursive} تر
\olref[par]{defn:recursive-fn} \emph{لږ} عمومي دے۔ خو له نېکه مرغه دا
توپير ظاهري دے؛ که څه هم تعريفونه بېل دي، د عمومي بازګشتي تابعو سټ او
د بازګشتي تابعو سټ يو او هماغه سټ دے۔
\olpseganchor{OLP-0227-B007}{end}

\olpseganchor{OLP-0228-B004}{start}
\olchapter{cmp}{thy}{د محاسبه کېدنې تيوري}
\olpseganchor{OLP-0228-B004}{end}

\olpseganchor{OLP-0228-B005}{start}
\begin{editorial}
  د دې څپرکي مواد بايد وکتل او پراخ کړل شي۔ په ځانګړي ډول تر اوسه
  هېڅ تمرين نۀ لري۔
\end{editorial}
\olpseganchor{OLP-0228-B005}{end}

\olpseganchor{OLP-0229-B004}{start}
\olfileid{cmp}{thy}{int}
\olsection{پېژندنه}
\olpseganchor{OLP-0229-B004}{end}

\olpseganchor{OLP-0229-B005}{start}
د منطق هغه څانګه چې \emph{د محاسبه کېدنې تيوري} بلل کېږي، د تابعو او
سټونو د محاسبه کېدنې يا نسبي محاسبه کېدنې اړوندې مسئلې څېړي۔ دا د
کليني د اغېز نښه ده چې دې مضمون ته پخوا \emph{د بازګښت تيوري} ويل
کېدل، او نن دواړه نومونه په عام ډول کارېږي۔
\olpseganchor{OLP-0229-B005}{end}

\olpseganchor{OLP-0229-B006}{start}
تابع~\LR{$f\colon \Nat \pto \Nat$} هغه وخت \emph{جزوي محاسبه کېدونکې}
وبولئ چې د محاسبې په کوم مدل کښې محاسبه کېدے شي۔ که \LR{$f$} هرځاے تعريف
شوې وي، په ساده ډول به وايو چې \LR{$f$} \emph{محاسبه کېدونکې} ده۔ هغه
اړيکه~\LR{$R$} هم محاسبه کېدونکې بلل کېږي چې ځانګړونکې تابع~\LR{$\Char{R}$} ئې
محاسبه کېدونکې وي۔ که \LR{$f$} او~\LR{$g$} جزوي تابعې وي، \LR{$f(x) \fdefined$} به
په دې مانا ليکو چې \LR{$f$} پر~\LR{$x$} تعريف شوې ده، يعنې \LR{$x$} د~\LR{$f$} په تعريف
ساحه کښې دے؛ او \LR{$f(x) \fundefined$} به د دې مخالفه مانا، يعنې دا چې
\LR{$f$} پر~\LR{$x$} تعريف شوې نۀ ده، څرګندوي۔ \LR{$f(x) \simeq g(x)$} به په دې
مانا کاروو چې يا \LR{$f(x)$} او~\LR{$g(x)$} دواړه تعريف شوي نۀ دي، يا دواړه
تعريف شوي او سره برابر دي۔
\olpseganchor{OLP-0229-B006}{end}

\olpseganchor{OLP-0229-B007}{start}
د محاسبه کېدنې تيوري د محاسبې کوم ځانګړي مدل ته له مراجعې پرته هم
څېړل کېدے شي۔ د دې دپاره ښودل کېږي چې يوه نړيواله جزوي محاسبه
کېدونکې تابع~\LR{$\fn{Un}(k,x)$} شته۔ په دې سره جزوي محاسبه کېدونکې تابعې
په لړ کښې راوستل کېدے شي۔ موږ به \LR{$\cfind{k}$} د~\LR{$k$}-مې يوځاييزې جزوي
محاسبه کېدونکې تابع دپاره وکاروو، چې د
\LR{$\cfind{k}(x) \simeq \fn{Un}(k,x)$} له لارې تعريفېږي۔ (کليني د دې
موخې دپاره \LR{$\{ k \}$} کارولی و، خو په وروستيو کښې دا نښه هومره نۀ ده
کارول شوې۔) لږ لا په عمومي ډول، د هرې ځاييزې کچې جزوي محاسبه کېدونکې
تابعې په يو شان ډول په لړ کښې راوستل کېدے شي، او \LR{$\cfind{k}[n]$} به د
\LR{$k$}-مې \LR{$n$}-ځاييزې جزوي محاسبه کېدونکې تابع دپاره کاروو۔
\olpseganchor{OLP-0229-B007}{end}

\olpseganchor{OLP-0229-B008}{start}
که \LR{$f(\vec x,y)$} هرځاے تعريف شوې يا جزوي تابع وي، نو
\LR{$\umin{y}{f(\vec x,y)}$} د~\LR{$\vec x$} هغه تابع ده چې تر ټولو وړوکے~\LR{$y$}
راګرځوي چې \LR{$f(\vec x,y)=0$} وي، په دې شرط چې \LR{$f(\vec x,0)$}، \dots،
\LR{$f(\vec x,y-1)$} ټول تعريف شوي وي؛ که داسې~\LR{$y$} نۀ وي،
\LR{$\umin{y}{f(\vec x,y)}$} تعريف شوے نۀ دے۔ که \LR{$R(\vec x,y)$} يوه اړيکه
وي، \LR{$\umin{y}{R(\vec x,y)}$} تر ټولو وړوکے~\LR{$y$} تعريفېږي چې
\LR{$R(\vec x,y)$} رښتيا وي؛ په بله وينا، تر ټولو وړوکے~\LR{$y$} چې
\LR{$1 \tsub \Char{R}(\vec x,y)=0$} وي۔
\olpseganchor{OLP-0229-B008}{end}

\olpseganchor{OLP-0229-B009}{start}
د دې د ښودلو دپاره چې يوه تابع محاسبه کېدونکې ده، دوه لارې شته:
\begin{enumerate}
\item په کلک ډول: يو ټيورينګ ماشين يا جزوي بازګشتي تابع په څرګند ډول
  بيان کړئ، او وښيئ چې همدا موخه شوې تابع محاسبه کوي؛
\item په غير رسمي ډول: هغه الګوريتم بيان کړئ چې تابع محاسبه کوي، او د
  چرچ اصل ته مراجعه وکړئ۔
\end{enumerate}
د دواړو تر منځ تېره پوله نشته؛ د يو الګوريتم مفصل بيان بايد دومره
معلومات ورکړي چې په اصل کښې د سم ټيورينګ ماشين يا د جزوي بازګشتي
تعريفونو د لړۍ طرحه نسبتاً څرګنده وي۔ بشپړ کلک تعريفونه ښايي معلوماتي
نۀ وي، او هڅه به کوو چې د دغو دوو لارو تر منځ مناسب منځنی دريځ ومومو؛
په لنډ ډول، هڅه به کوو داسې شهودي خو کلک ثبوتونه ومومو چې ترې څرګنده
وي کره تعريفونه تر لاسه کېدے شي۔
\olpseganchor{OLP-0229-B009}{end}

\olpseganchor{OLP-0230-B004}{start}
\olfileid{cmp}{thy}{cod}
\olsection{د محاسبو کوډول}
\olpseganchor{OLP-0230-B004}{end}

\olpseganchor{OLP-0230-B005}{start}
د محاسبې په هر مدل کښې لاندې کارونه کېدے شي:
\begin{enumerate}
\item د محاسبه کېدونکو تابعو \emph{تعريفونه} په منظم ډول بيانول۔ د
  بېلګې په توګه، د ټيورينګ ماشين مشخصې، بازګشتي تعريفونه يا د پروګرام
  ليکلو په يوه ژبه کښې پروګرامونه د دې تعريفونو ورکوونکي ګڼلے شئ۔
\item د يو ورکړل شوي ننوت دپاره، د کوم تعريف له مخې د يوې تابع د
  محاسبې بشپړ ريکارډ بيانول۔ د بېلګې په توګه، د ټيورينګ ماشين محاسبه
  د محاسبې په هر ګام کښې د بڼو د لړۍ له لارې بيانېدے شي (د ماشين
  حالت او د ټېپ منځپانګه)۔
\item دا کتل چې د يوې محاسبې يو اټکلي ريکارډ په رښتيا د دې ريکارډ دے
  که نه چې د ورکړل شوي تعريف محاسبه کېدونکې تابع به پر ورکړل شوي ننوت
  څنګه محاسبه شي (هغه ننوت چې تابع پرې تعريف شوې وي، يعنې محاسبه
  درېږي)۔
\item د يوې محاسبې د بشپړ ريکارډ له داسې بيان څخه د ورکړل شوي ننوت
  دپاره د تابع قيمت راايستل۔ د بېلګې په توګه، د يوې درېدونکې ټيورينګ
  ماشين محاسبې په وروستي ګام کښې د ټېپ منځپانګه قيمت دے۔
\end{enumerate}
\olpseganchor{OLP-0230-B005}{end}

\olpseganchor{OLP-0230-B006}{start}
د کوډولو په کارولو سره د محاسبه کېدونکې تابع هر بيان ته داسې عددي
\emph{شاخص} ټاکل کېدے شي چې لارښوونې په محاسبه کېدونکي ډول له شاخصه
بېرته تر لاسه شي۔ په ورته ډول د يوې محاسبې بشپړ ريکارډ هم په يوه عدد
کوډېدے شي۔ په پايله کښې پيدا شوې حسابي اړيکه، ``\LR{$s$} د شاخص~\LR{$e$} د تابع
د محاسبې ريکارډ پر ننوت~\LR{$x$} کوډوي''، او تابع، ``د کوډ~\LR{$s$} لرونکې د
محاسبې لړۍ وت''، بيا محاسبه کېدونکې دي؛ په حقيقت کښې دواړه بنسټيز
بازګشتي دي۔
\olpseganchor{OLP-0230-B006}{end}

\olpseganchor{OLP-0230-B007}{start}
دا بنسټيز حقيقت ډېر ځواکمن دے، او پرېږدي چې د محاسبه کېدنې په اړه ډېرې
حيرانوونکې او مهمې پايلې د محاسبې له غوره شوي مدل څخه په خپلواک ډول
ثابتې کړو۔
\olpseganchor{OLP-0230-B007}{end}

\olpseganchor{OLP-0231-B004}{start}
\olfileid{cmp}{thy}{nfm}
\olsection{د نورمال شکل قضيه}
\olpseganchor{OLP-0231-B004}{end}

\olpseganchor{OLP-0231-B005}{start}
فرض کړئ چې د محاسبه کېدونکو تابعو تعريفونه بيانولے شو، او کتلے شو
چې پر کوم ننوت د هغې تابع د محاسبې د بشپړ ريکارډ يو اټکلي بيان سم دے
که نه۔ نو له دې سره سمه ښکاري چې د محاسبې له مدل څخه په خپلواک ډول د
هرې محاسبه کېدونکې تابع~\LR{$f$} قيمت پر هر ننوت~\LR{$x$} داسې ټاکلے شو:
\begin{enumerate}
  \item د محاسبې د ريکارډونو د ټولو ممکنه بيانونو له منځه پلټنه وکړئ۔
  \item وازمويئ چې ورکړل شوے ريکارډ د~\LR{$f(x)$} د محاسبې ريکارډ دے که نه۔
  \item که سم ريکارډ مو موندلے وي، د~\LR{$f(x)$} قيمت ترې راوباسئ۔
\end{enumerate}
دا چې خبره په رښتيا همداسې ده، د کليني د نورمال شکل د قضيې منځپانګه
ده۔
\olpseganchor{OLP-0231-B005}{end}

\olpseganchor{OLP-0231-B006}{start}
\begin{thm}[د کليني د نورمال شکل قضيه]
\ollabel{thm:normal-form}
يوه بنسټيزه بازګشتي اړيکه~\LR{$T(e,x,s)$} او يوه بنسټيزه بازګشتي تابع
~\LR{$U(s)$} شته چې دا خاصيت لري: که \LR{$f$} هره جزوي محاسبه کېدونکې تابع وي،
نو د کوم~\LR{$e$} دپاره
\[
f(x) \simeq U(\umin{s}{T(e, x, s)})
\]
د هر~\LR{$x$} دپاره پوره کېږي۔
\end{thm}
\olpseganchor{OLP-0231-B006}{end}

\olpseganchor{OLP-0231-B007}{start}
\begin{proof}[د ثبوت خاکه]
د محاسبې د هر مدل دپاره د محاسبه کېدونکې تابع~\LR{$f$} يو بيان په کلک ډول
تعريفېدے او داسې بيان په يوه طبيعي عدد~\LR{$e$} کوډېدے شي۔ د ``محاسبې
لړۍ'' يو مفهوم هم په کلک ډول تعريفېدے شي چې د شاخص~\LR{$e$} تابع پر
ننوت~\LR{$x$} د محاسبه کولو بهير ثبتوي۔ دا ډول د محاسبې لړۍ هم په ورته ډول
په يوه عدد~\LR{$s$} کوډېدے شي۔ دا هر څه په داسې ډول کېدے شي چې
\begin{enumerate}
  \item اړيکه \LR{$T(e,x,s)$}، چې هله او يوازې هله پوره کېږي چې عدد~\LR{$s$} د
  شاخص~\LR{$e$} د تابع د محاسبې لړۍ پر ننوت~\LR{$x$} کوډوي، او
  \item تابع \LR{$U(s)$}، چې په~\LR{$s$} کوډ شوې د محاسبې لړۍ د هغې محاسبې
  وروستۍ پايلې ته وړي،
\end{enumerate}
دواړه محاسبه کېدونکي وي۔ په حقيقت کښې اړيکه~\LR{$T$} او تابع~\LR{$U$} دواړه
بنسټيز بازګشتي دي۔
\end{proof}
\olpseganchor{OLP-0231-B007}{end}

\olpseganchor{OLP-0231-B008}{start}
\begin{explain}
د نورمال شکل د قضيې د کلک ثبوت دپاره بايد د محاسبې يو مدل وټاکو، د
محاسبه کېدونکو تابعو د بيانونو او د محاسبې د لړيو کوډول په تفصيل عملي
کړو، او تصديق کړو چې اړيکه~\LR{$T$} او تابع~\LR{$U$} بنسټيز بازګشتي دي۔ د ډېرو
کارونو دپاره دومره بسنه کوي چې \LR{$T$} او~\LR{$U$} محاسبه کېدونکي او \LR{$U$}
هرځاے تعريف شوې وي۔
\olpseganchor{OLP-0231-B008}{end}

\olpseganchor{OLP-0231-B009}{start}
ښايي د نورمال شکل د قضيې ثبوت په لنډ شعار کښې داسې ښه ياد پاتې شي:
\LR{$\umin{s}{T(e,x,s)}$} پر ننوت~\LR{$x$} د شاخص~\LR{$e$} د تابع د محاسبې يوه لړۍ
لټوي، او که داسې لړۍ وموندل شي، \LR{$U$} ئې وت راګرځوي۔
\end{explain}
\olpseganchor{OLP-0231-B009}{end}

\olpseganchor{OLP-0231-B010}{start}
که د محاسبې مدل جزوي بازګشتي تابعې وي (چې کليني په اصل کښې همدا
کارولې)، قضيه ښيي چې د هرې تابع په تعريف کښې د بې‌حده پلټنې يوازې يوه
کارونه، يعنې يوازې يو \LR{$\umin{y}{f(\vec x,y)}$} عامل، پکار دے۔ په دې
مانا قضيه ښيي چې هره جزوي بازګشتي تابع يو نورمال شکل لري۔
\olpseganchor{OLP-0231-B010}{end}

\olpseganchor{OLP-0231-B011}{start}
\LR{$T$} او~\LR{$U$} د شمېرنې \LR{$\cfind{0}$}، \LR{$\cfind{1}$}، \LR{$\cfind{2}$}، \dots
د تعريف دپاره کارېدے شي۔ له دې وروسته به فرض کوو چې د~\LR{$T$} او~\LR{$U$} يو
مناسب انتخاب مو ثابت کړے دے، او معادله
\[
\cfind{e}(x) \simeq U(\umin{s}{T(e,x,s)})
\]
به د~\LR{$\cfind{e}$} \emph{تعريف} ګڼو۔
\olpseganchor{OLP-0231-B011}{end}

\olpseganchor{OLP-0231-B012}{start}
دا يو بل ګټور حقيقت دے:
\olpseganchor{OLP-0231-B012}{end}

\olpseganchor{OLP-0231-B013}{start}
\begin{thm}
هره جزوي محاسبه کېدونکې تابع بې‌شمېره ډېر شاخصونه لري۔
\end{thm}
\olpseganchor{OLP-0231-B013}{end}

\olpseganchor{OLP-0231-B014}{start}
دا خبره بيا هم په شهودي ډول څرګنده ده۔ د يوې محاسبه کېدونکې تابع له
هر ورکړل شوي بيان څخه يو بل بيان جوړولے شو چې هماغه تابع (د ننوت او
وت هماغه جوړه) محاسبه کوي، خو مثلاً لومړے داسې څه کوي چې پر محاسبه
هېڅ اغېز نۀ لري (لکه دا کتل چې \LR{$0=0$} دے، يا تر~\LR{$5$} پورې شمېرل، او
داسې نور)۔ د بدل شوي بيان شاخص به تل له اصلي شاخصه بېل وي۔ دواړه د
هماغې تابع شاخصونه دي، خو تابع په لږ بېل ډول محاسبه کوي۔
\olpseganchor{OLP-0231-B014}{end}

\olpseganchor{OLP-0232-B004}{start}
\olfileid{cmp}{thy}{smn}
\olsection{د \texorpdfstring{\LR{$s$}}{s}-\texorpdfstring{\LR{$m$}}{m}-\texorpdfstring{\LR{$n$}}{n} قضيه}
\olpseganchor{OLP-0232-B004}{end}

\olpseganchor{OLP-0232-B005}{start}
\begin{explain}
راتلونکې قضيه د داسې دليل له مخې د ``\LR{$s$}-\LR{$m$}-\LR{$n$} قضيې'' په نوم
پېژندل کېږي چې لږ وروسته به څرګند شي۔ ستونزمنه برخه دا پوهېدل دي چې
قضيه په کره ډول څه وايي؛ کله چې ئې بيان درک کړئ، قضيه به بيخي څرګنده
وبرېښي۔
\end{explain}
\olpseganchor{OLP-0232-B005}{end}

\olpseganchor{OLP-0232-B006}{start}
\begin{thm}
\ollabel{thm:s-m-n}
  د طبيعي عددونو د هرې جوړې \LR{$n$} او~\LR{$m$} دپاره يوه بنسټيزه بازګشتي
  تابع~\LR{$s^m_n$} شته، داسې چې د هرې لړۍ
  \LR{$e$}، \LR{$a_0$}، \dots، \LR{$a_{m-1}$}، \LR{$y_0$}، \dots، \LR{$y_{n-1}$} دپاره لرو:
  \[
  \cfind{s^m_n(e, a_0, \dots, a_{m-1})}[n](y_0, \dots, y_{n-1}) \simeq
  \cfind{e}[m+n](a_0, \dots, a_{m-1}, y_0, \dots, y_{n-1}).
\]
\end{thm}
\olpseganchor{OLP-0232-B006}{end}

\olpseganchor{OLP-0232-B007}{start}
\begin{explain}
ګټوره ده چې \LR{$s^m_n$} پر \emph{پروګرامونو} د عمل کوونکې تابع په توګه
وګڼو۔ يعنې \LR{$s^m_n$} د يوې \LR{$(m+n)$}-ځاييزې تابع پروګرام~\LR{$e$} او ثابت
ننوتونه \LR{$a_0$}، \dots، \LR{$a_{m-1}$} اخلي؛ بيا د پاتې دليلونو د
\LR{$n$}-ځاييزې تابع پروګرام \LR{$s^m_n(e, a_0, \dots, a_{m-1})$} راګرځوي۔
که~\LR{$e$} د يوه ټيورينګ ماشين تشريح وګڼئ، نو
\LR{$s^m_n(e, a_0, \dots, a_{m-1})$} د داسې ټيورينګ ماشين کوډ دے چې د
\LR{$y_0$}، \dots،~\LR{$y_{n-1}$} په ننوت کښې \LR{$a_0$}، \dots،~\LR{$a_{m-1}$} د
ننوت د لړۍ په سر کښې ور زياتوي، او~\LR{$e$} چلوي۔ نو هره~\LR{$s^m_n$} يوازې
هغه بنسټيزه بازګشتي تابع ده چې د اړوند ټيورينګ ماشين کوډ پيدا کوي۔
\end{explain}
\olpseganchor{OLP-0232-B007}{end}

\olpseganchor{OLP-0233-B004}{start}
\olfileid{cmp}{thy}{uni}
\olsection{نړيواله جزوي محاسبه کېدونکې تابع}
\olpseganchor{OLP-0233-B004}{end}

\olpseganchor{OLP-0233-B005}{start}
\begin{thm}
\ollabel{thm:univ-comp}
يوه نړيواله جزوي محاسبه کېدونکې تابع~\LR{$\fn{Un}(e,x)$} شته۔ په بله
وينا، داسې تابع \LR{$\fn{Un}(e,x)$} شته چې:
\begin{enumerate}
\item \LR{$\fn{Un}(e,x)$} جزوي محاسبه کېدونکې ده۔
\item که \LR{$f(x)$} هره جزوي محاسبه کېدونکې تابع وي، نو يو طبيعي عدد
\LR{$e$} شته، داسې چې د هر~\LR{$x$} دپاره \LR{$f(x) \simeq \fn{Un}(e,x)$} وي۔
\end{enumerate}
\end{thm}
\olpseganchor{OLP-0233-B005}{end}

\olpseganchor{OLP-0233-B006}{start}
\begin{proof}
\LR{$\fn{Un}(e,x) \simeq U(\umin{s}{T(e,x,s)})$} وټاکئ، چې~\LR{$U$} او~\LR{$T$}
د کليني د نورمال شکل په قضيه کښې د
(\olref[nfm]{thm:normal-form}) په شان دي۔
\end{proof}
\olpseganchor{OLP-0233-B006}{end}

\olpseganchor{OLP-0233-B007}{start}
\begin{explain}
دا يوازې د دې خبرې کره بڼه ده چې موږ د جزوي محاسبه کېدونکو تابعو
اغېزمنه شمېرنه لرو۔ نظر دا دے چې که د هغې تابع دپاره~\LR{$f_e$} وليکو چې
\LR{$f_e(x) = \fn{Un}(e,x)$} ئې تعريفوي، نو لړۍ \LR{$f_0$}، \LR{$f_1$}، \LR{$f_2$}،
\dots ټولې جزوي محاسبه کېدونکې تابعې رانغاړي، او \LR{$f_e(x)$} د~\LR{$e$}
او~\LR{$x$} په نسبت ``په يو شان'' محاسبه کېدے شي۔ د ساده‌والي دپاره داسې
دوه‌ځاييزه تابع کاروو چې د يوځاييزو تابعو دپاره نړيواله ده؛ خو د
عددونو د لړيو په کوډولو سره دا خبره په اسانۍ ډېرو دليلونو ته عمومي
کولے شو۔ د بېلګې په توګه، پام وکړئ چې که \LR{$f(x,y,z)$} يوه \LR{$3$}-ځاييزه
جزوي بازګشتي تابع وي، نو تابع
\LR{$g(x) \simeq f((x)_0, (x)_1, (x)_2)$} يوه يوځاييزه جزوي بازګشتي
تابع ده۔
د نوې تابع تعريف په دې پورې تړلے دے چې درې‌ځاييزه تابع پر راايستل شويو ننوتونو تعريف شوې ده که نه۔
\end{explain}
\olpseganchor{OLP-0233-B007}{end}

\olpseganchor{OLP-0234-B004}{start}
\olfileid{cmp}{thy}{nou}
\olsection{هېڅ نړيواله محاسبه کېدونکې تابع نشته}
\olpseganchor{OLP-0234-B004}{end}

\olpseganchor{OLP-0234-B005}{start}
که څه هم داسې جزوي محاسبه کېدونکې تابع شته چې د جزوي محاسبه کېدونکو
تابعو دپاره نړيواله ده، داسې هرځاے تعريف شوې محاسبه کېدونکې تابع
نشته چې د ټولو هرځاے تعريف شويو محاسبه کېدونکو تابعو دپاره نړيواله
وي۔
\olpseganchor{OLP-0234-B005}{end}

\olpseganchor{OLP-0234-B006}{start}
\begin{thm}
\ollabel{thm:no-univ}
د ټولو هرځاے تعريف شويو محاسبه کېدونکو تابعو لپاره هېڅ هرځاے تعريف
شوې محاسبه کېدونکې نړيواله تابع نشته۔ په بله وينا، هره هرځاے تعريف شوې تابع
\LR{$\fn{Un}'(k, x)$} چې دا ځانګړنه ولري چې که \LR{$f(x)$} هرځاے تعريف شوې
محاسبه کېدونکې تابع وي، نو يو طبيعي عدد~\LR{$k$} شته، داسې چې د هر~\LR{$x$}
دپاره \LR{$f(x) = \fn{Un}'(k,x)$} وي، محاسبه کېدونکې نۀ ده۔
\end{thm}
\olpseganchor{OLP-0234-B006}{end}

\olpseganchor{OLP-0234-B007}{start}
\begin{proof}
ثبوت يو ساده قطري استدلال دے: که \LR{$\fn{Un}'(k,x)$} هرځاے تعريف شوې او
محاسبه کېدونکې وای، نو
\[
d(x) = \fn{Un}'(x, x) + 1
\]
به هم هرځاے تعريف شوې او محاسبه کېدونکې وه۔ خو د تعريف له مخې
\LR{$d(k)$} له~\LR{$\fn{Un}'(k,k)$} سره برابره نۀ ده۔ نو د هر~\LR{$k$} دپاره د
\LR{$d(x)$} او~\LR{$\fn{Un}'(k, x)$} قيمتونه لږ تر لږه پر يوه~\LR{$x$}، يعنې
\LR{$x = k$}، سره بېل دي۔
\end{proof}
\olpseganchor{OLP-0234-B007}{end}

\olpseganchor{OLP-0234-B008}{start}
\begin{explain}
پورته \olref[uni]{thm:univ-comp} ښيي چې له دې قطري استدلاله وتلے شو،
خو يوازې په دې بيه چې نړيواله تابع جزوي ومنو۔ يعنې~\LR{$\fn{Un}$} د هرځاے
تعريف شويو محاسبه کېدونکو تابعو دپاره نړيواله ده، خو پخپله هرځاے
تعريف شوې نۀ ده۔ په جزوي حالت کښې پر قطر محاسبه ناتعريف پاتې کېدے شي، نو د يو زياتولو له لارې
تل بېل ارزښت نۀ شو تر لاسه کولے۔
\end{explain}
\olpseganchor{OLP-0234-B008}{end}

\olpseganchor{OLP-0234-B009}{start}
\begin{prob}
  د دې د پوهېدو دپاره چې د \olref{thm:no-univ} د ثبوت قطري استدلال
  ولې په جزوي حالت کښې کار نۀ کوي، تابع
  \LR{$f(x) \simeq \fn{Un}(x,x)+1$} وګورئ۔ ايا دا جزوي محاسبه کېدونکې
  ده؟ که وي، يو شاخص~\LR{$e$} لري، يعنې
  \LR{$f(x) \simeq \fn{Un}(e,x)$}۔ د~\LR{$f(e)$} په اړه څه ويلے شئ؟
\end{prob}
\olpseganchor{OLP-0234-B009}{end}

\olpseganchor{OLP-0235-B004}{start}
\olfileid{cmp}{thy}{hlt}
\olsection{د درېدنې مسئله}
\olpseganchor{OLP-0235-B004}{end}

\olpseganchor{OLP-0235-B005}{start}
د جوړښت له مخې نړيواله جزوي محاسبه کېدونکې تابع~\LR{$\fn{Un}(e,x)$} هله
او يوازې هله تعريف شوې ده چې په~\LR{$e$} کوډ شوې د تابع محاسبه د ننوت~\LR{$x$}
دپاره يو قيمت پيدا کړي۔ طبيعي پوښتنه دا ده چې ايا موږ پرېکړه کولے شو
چې داسې ده که نه۔ په حقيقت کښې دا پرېکړه نۀ شو کولے۔ د محاسبې د
ټيورينګ ماشين په مدل کښې دا په دې مانا ده چې ايا ورکړل شوے ټيورينګ
ماشين پر ورکړل شوي ننوت درېږي که نه، په محاسبوي ډول د پرېکړې وړ نۀ
ده۔ له همدې کبله لاندې قضيه د ``درېدنې د مسئلې د ناپرېکړتيا'' په نوم
پېژندل کېږي۔ لاندې به دوه ثبوتونه ورکړو۔ لومړے زموږ د مخکنۍ بحث لړۍ
پسې غځوي، او دويم ډېر نېغ دے۔
\olpseganchor{OLP-0235-B005}{end}

\olpseganchor{OLP-0235-B006}{start}
\begin{thm}
\ollabel{thm:halting-problem}
فرض کړئ
\[
h(e, x)  =
\begin{cases}
1 & \text{\RL{که \LR{$\fn{Un}(e, x)$} تعريف شوې وي}} \\
0 & \text{\RL{که داسې نۀ وي۔}}
\end{cases}
\]
نو \LR{$h$} محاسبه کېدونکې نۀ ده۔
\end{thm}
\olpseganchor{OLP-0235-B006}{end}

\olpseganchor{OLP-0235-B007}{start}
\begin{proof}
فرض کړئ \LR{$h$} محاسبه کېدونکې ده۔ ښيو چې په دې سره يوه نړيواله محاسبه
کېدونکې تابع تعريفولے شو۔ داسې تعريف کړئ:
\[
\fn{Un'}(e,x) =
\begin{cases}
\fn{Un}(e,x) & \text{\RL{که \LR{$h(e,x) = 1$} وي}} \\
0 & \text{\RL{که داسې نۀ وي۔}}
\end{cases}
\]
خو اوس \LR{$\fn{Un'}(e, x)$} هرځاے تعريف شوې تابع ده، او که~\LR{$h$} محاسبه
کېدونکې وي نو دا هم محاسبه کېدونکې ده۔ د بېلګې په توګه، د بنسټيز
بازګښت په کارولو سره~\LR{$g$} داسې تعريفولے شو:
\begin{align*}
g(0, e, x) & \simeq 0 \\
g(y+1, e, x) & \simeq \fn{Un}(e,x);
\end{align*}
بيا
\[
\fn{Un'}(e,x) \simeq g(h(e,x),e,x).
\]
څرنګه چې \LR{$\fn{Un'}(e,x)$} په هر هغه ځاے کښې له~\LR{$\fn{Un}(e,x)$} سره
برابره ده چې وروستۍ تابع تعريف شوې وي، نو~\LR{$\fn{Un'}$} د هغو جزوي
محاسبه کېدونکو تابعو دپاره نړيواله ده چې هرځاے تعريف شوې وي۔ خو دا له
\olref[nou]{thm:no-univ} سره تناقض دے۔
\end{proof}
\olpseganchor{OLP-0235-B007}{end}

\olpseganchor{OLP-0235-B008}{start}
\begin{proof}
فرض کړئ \LR{$h(e,x)$} محاسبه کېدونکې ده۔ تابع~\LR{$g$} داسې تعريف کړئ:
\[
g(x) =
\begin{cases}
  0                & \text{\RL{که \LR{$h(x,x) = 0$} وي}} \\
  \fundefined & \text{\RL{که داسې نۀ وي۔}}
\end{cases}
\]
تابع~\LR{$g$} جزوي محاسبه کېدونکې ده۔ د بېلګې په توګه، دا د
\LR{$\umin{y}{h(x,x) = 0}$} په توګه تعريفولے شو۔ نو د کوم~\LR{$e$} دپاره د هر
~\LR{$x$} په نسبت \LR{$g(x) \simeq \fn{Un}(e,x)$} ده۔ ايا~\LR{$g$} پر~\LR{$e$} تعريف
شوې ده؟ که تعريف شوې وي، نو د~\LR{$g$} د تعريف له مخې \LR{$h(e,e) = 0$} دے
(کله چې~\LR{$g$} تعريف شوې وي، يوازې د~\LR{$0$} قيمت اخلي)۔ د~\LR{$h$} د تعريف له
مخې دا په دې مانا ده چې \LR{$\fn{Un}(e, e)$} تعريف شوې نۀ ده۔ زموږ د دې
فرض له مخې چې د هر~\LR{$x$} دپاره \LR{$g(x) \simeq \fn{Un}(e, x)$} ده،
\LR{$g(e)$} هم تعريف شوې نۀ ده، چې تناقض دے۔ بل پلو، که \LR{$g(e)$} تعريف شوې
نۀ وي، نو \LR{$h(e,e) \neq 0$}، او ځکه \LR{$h(e,e) = 1$} دے۔ له دې څخه راځي
چې \LR{$\fn{Un}(e, e)$} تعريف شوې ده۔ خو څرنګه چې
\LR{$g(x) \simeq \fn{Un}(e, x)$} ده، نو \LR{$g(e)$} به هم تعريف شوې وه۔ بيا
تناقض راغے۔
\end{proof}
\olpseganchor{OLP-0235-B008}{end}

\olpseganchor{OLP-0235-B009}{start}

\begin{explain}
دا استدلال د ټيورينګ ماشينونو په ژبه هم بيانولے شو۔ فرض کړئ داسې
ټيورينګ ماشين~\LR{$H$} شته چې د يو ټيورينګ ماشين~\LR{$E$} تشريح او يو
ننوت~\LR{$x$} اخلي، او پرېکړه کوي چې~\LR{$E$} پر~\LR{$x$} درېږي که نه۔ بيا به يو
بل ټيورينګ ماشين~\LR{$G$} جوړولے شو چې يو ننوت~\LR{$x$} اخلي، \LR{$H$} چلوي څو
پرېکړه وکړي چې د شاخص~\LR{$x$} ماشين~\LR{$M_x$} پر ننوت~\LR{$x$} درېږي که نه، او
بيا د هغې مخالف کار کوي۔ په بله وينا، که~\LR{$H$} ووايي چې~\LR{$M_x$} پر~\LR{$x$}
درېږي، \LR{$G$} په ناپايه کړۍ کښې ننوځي؛ او که~\LR{$H$} ووايي چې~\LR{$M_x$} پر~\LR{$x$}
نۀ درېږي، نو~\LR{$G$} بس درېږي۔ ايا~\LR{$G$} د خپل شاخص په ننوت درېږي؟ پورته
استدلال ښيي چې دا هله او يوازې هله درېږي چې نۀ درېږي—يو تناقض۔ نو
زموږ دا فرض چې داسې ټيورينګ ماشين~\LR{$H$} شته، بايد ناسم وي۔
\end{explain}

\olpseganchor{OLP-0235-B009}{end}

\olpseganchor{OLP-0236-B004}{start}
\olfileid{cmp}{thy}{rus}
\olsection{د رسل له تناقض سره پرتله}
\olpseganchor{OLP-0236-B004}{end}

\olpseganchor{OLP-0236-B005}{start}
ګټوره ده چې د دې برخې استدلالونه د رسل له تناقض سره پرتله او بېل
کړو:
\begin{enumerate}
\item د رسل تناقض: فرض کړئ \LR{$S = \Setabs{x}{x \notin x}$}۔ بيا \LR{$S
  \in S$} هله او يوازې هله چې \LR{$S \notin S$} وي، چې تناقض دے۔
\olpseganchor{OLP-0236-B005}{end}

\olpseganchor{OLP-0236-B006}{start}
  \emph{پايله:} داسې سټ~\LR{$S$} نشته۔ د ``ټولو سټونو د سټ'' شتون منل د
  سټ تيورۍ له نورو بديهي اصولو سره ناسازګار دي۔
\olpseganchor{OLP-0236-B006}{end}

\olpseganchor{OLP-0236-B007}{start}
\item د رسل د تناقض يوه بدله بڼه: \LR{$F$} له ټولو تابعو له سټ څخه
  \LR{$\{ 0, 1 \}$} ته هغه ``تابع'' وګڼئ چې داسې تعريفېږي:
  \[
  F(f) =
  \begin{cases}
    1 & \text{\RL{که \LR{$f$} د~\LR{$f$} د تعريف ساحې غړے وي، او \LR{$f(f) = 0$} وي}} \\
    0 & \text{\RL{که داسې نۀ وي}}
  \end{cases}
  \]
  ورته استدلال ښيي چې \LR{$F(F) = 0$} هله او يوازې هله چې \LR{$F(F) = 1$}
  وي، چې تناقض دے۔
\olpseganchor{OLP-0236-B007}{end}

\olpseganchor{OLP-0236-B008}{start}
  \emph{پايله:} \LR{$F$} تابع نۀ ده۔ د ``ټولو تابعو سټ'' دومره لوے دے
  چې د يوې تابع تعريف ساحه نۀ شي کېدے۔
\olpseganchor{OLP-0236-B008}{end}

\olpseganchor{OLP-0236-B009}{start}
\item قطري استدلال: \LR{$f_0$}، \LR{$f_1$}، \dots د جزوي محاسبه کېدونکو تابعو
  شمېرنه وي، او \LR{$G \colon \Nat \to \{ 0, 1 \}$} داسې تعريف کړئ:
  \[
  G(x) =
  \begin{cases}
    1 & \text{\RL{که \LR{$f_x(x)\downarrow = 0$} وي}} \\
    0 & \text{\RL{که داسې نۀ وي}}
  \end{cases}
  \]
  که \LR{$G$} محاسبه کېدونکې وي، نو د کوم~\LR{$k$} دپاره تابع~\LR{$f_k$} ده۔ خو
  بيا \LR{$G(k) = 1$} هله او يوازې هله چې \LR{$G(k) = 0$} وي، چې تناقض دے۔
\olpseganchor{OLP-0236-B009}{end}

\olpseganchor{OLP-0236-B010}{start}
  \emph{پايله:} \LR{$G$} محاسبه کېدونکې نۀ ده۔ پام وکړئ چې د سټ تيورۍ د
  بديهي اصولو له مخې~\LR{$G$} بيا هم تابع ده؛ دلته تناقض نشته، يوازې يوه
  څرګندونه ده۔
\end{enumerate}
\olpseganchor{OLP-0236-B010}{end}

\olpseganchor{OLP-0236-B011}{start}
د جزوي تابعو، محاسبه کېدونکو تابعو، جزوي محاسبه کېدونکو تابعو او
داسې نورو يادونه ګډوډوونکې کېدے شي۔ له~\LR{$\Nat$} څخه~\LR{$\Nat$} ته د ټولو
جزوي تابعو سټ د څيزونو يوه لويه ټولګه ده۔ ځينې ئې هرځاے تعريف شوې،
ځينې ئې محاسبه کېدونکې، ځينې دواړه هرځاے تعريف شوې او محاسبه
کېدونکې، او ځينې يوه هم نۀ دي۔ په ياد ولرئ چې کله بې له قيده
``تابع'' وايو، موخه مو هرځاے تعريف شوې تابع وي۔ نو دا ډولونه لرو:
\begin{enumerate}
\item محاسبه کېدونکې تابعې
\item هغه جزوي محاسبه کېدونکې تابعې چې هرځاے تعريف شوې نۀ دي
\item هغه تابعې چې محاسبه کېدونکې نۀ دي
\item هغه جزوي تابعې چې نۀ هرځاے تعريف شوې او نۀ محاسبه کېدونکې دي
\end{enumerate}
د دې د بېلولو دپاره ښايي داسې لوے مربع رسمول ګټور وي چې له~\LR{$\Nat$}
څخه~\LR{$\Nat$} ته ټولې جزوي تابعې وښيي؛ بيا پکښې دوه يو پر بل ورغلي
سيمې په نښه کړئ، چې يوه د هرځاے تعريف شويو تابعو او بله د جزوي
محاسبه کېدونکو تابعو وي۔ ښه تمرين دے چې وګورئ د نقشې په هره جوړه شوې
سيمه کښې يو څيز بيانولے شئ که نه۔
\olpseganchor{OLP-0236-B011}{end}

\olpseganchor{OLP-0237-B004}{start}
\olfileid{cmp}{thy}{cps}
\olsection{محاسبه کېدونکي سټونه}
\olpseganchor{OLP-0237-B004}{end}

\olpseganchor{OLP-0237-B005}{start}
د محاسبه کېدنې مفهوم له محاسبه کېدونکو تابعو څخه محاسبه کېدونکو
سټونو ته غځولے شو:
\olpseganchor{OLP-0237-B005}{end}

\olpseganchor{OLP-0237-B006}{start}
\begin{defn}
  فرض کړئ \LR{$S$} د طبيعي عددونو يو سټ دے۔ \LR{$S$} هله او يوازې هله
  \emph{محاسبه کېدونکے} دے چې ځانګړونکې تابع~\LR{$\Char{S}$} ئې محاسبه
  کېدونکې وي۔ په بله وينا، \LR{$S$} هله او يوازې هله محاسبه کېدونکے دے
  چې تابع
\[
\Char{S}(x) =
\begin{cases}
1 & \text{\RL{که \LR{$x \in S$} وي}} \\
0 & \text{\RL{که داسې نۀ وي}}
\end{cases}
\]
محاسبه کېدونکې وي۔ همدارنګه، اړيکه
\LR{$R(x_0, \dots, x_{k-1})$} هله او يوازې هله محاسبه کېدونکې ده چې
ځانګړونکې تابع ئې محاسبه کېدونکې وي۔
\olpseganchor{OLP-0237-B006}{end}

\olpseganchor{OLP-0237-B007}{start}
محاسبه کېدونکي سټونه او اړيکې \emph{د پرېکړې وړ} هم بلل کېږي۔
\end{defn}
\olpseganchor{OLP-0237-B007}{end}

\olpseganchor{OLP-0237-B008}{start}
\begin{explain}
پام وکړئ چې اوس د محاسبه کېدنې څو مفهومونه لرو: د جزوي تابعو، تابعو
او سټونو دپاره۔ دا سره ګډ نۀ کړئ!{} هغه ټيورينګ ماشين چې
  جزوي تابع محاسبه کوي، د تابع په هغو ننوتونو کښې ئې وت راګرځوي چې
  تابع پرې تعريف شوې وي؛ هغه ټيورينګ ماشين چې يو سټ محاسبه کوي، تر
  دې پرېکړې وروسته چې ننوت په سټ کښې دے که نه، يا~\LR{$1$} يا~\LR{$0$}
  راګرځوي۔
\end{explain}
\olpseganchor{OLP-0237-B008}{end}

\olpseganchor{OLP-0238-B004}{start}
\olfileid{cmp}{thy}{ces}
\olsection{په محاسبوي ډول د شمېر وړ سټونه}
\olpseganchor{OLP-0238-B004}{end}

\olpseganchor{OLP-0238-B005}{start}
\begin{defn}
يو سټ هغه وخت \emph{په محاسبوي ډول د شمېر وړ} دے چې تش وي يا د يوې
محاسبه کېدونکې تابع د قيمتونو سټ وي۔
\end{defn}
\olpseganchor{OLP-0238-B005}{end}

\olpseganchor{OLP-0238-B006}{start}
\begin{history}
په محاسبوي ډول د شمېر وړ سټونه د دې پر ځاے \emph{په بازګشتي ډول د
  شمېر وړ} سټونه هم بلل کېږي۔ دا اصلي اصطلاح ده، او نن دواړه نومونه،
او ورسره لنډونونه ``c.e.'' او ``r.e.''، په عام ډول کارېږي۔
\end{history}
\olpseganchor{OLP-0238-B006}{end}

\olpseganchor{OLP-0238-B007}{start}
\begin{explain}
بايد په دې فکر وکړئ چې تعريف څه مانا لري او اصطلاح ولې مناسبه ده۔ نظر
دا دے چې که~\LR{$S$} د محاسبه کېدونکې تابع~\LR{$f$} د قيمتونو سټ وي، نو
\[
S = \{ f(0), f(1), f(2), \dots \},
\]
او ځکه~\LR{$f$} د~\LR{$S$} د غړو د ``شمېرلو'' په توګه ګڼل کېدے شي۔ پام وکړئ
چې د تعريف له مخې~\LR{$f$} لازمه نۀ ده زياتېدونکې تابع وي، يعنې شمېرنه
لازمه نۀ ده په زياتېدونکي ترتيب وي۔ په حقيقت کښې~\LR{$f$} لازمه نۀ ده
يو په يو هم وي، يعنې د~\LR{$S$} په شمېرنه \LR{$f(0)$}، \LR{$f(1)$}، \LR{$f(2)$}،
\dots{} کښې تکرار منل شوے دے۔ د بېلګې په توګه، ثابته تابع
\LR{$f(x) = 0$} سټ~\LR{$\{ 0 \}$} شمېري۔
\end{explain}
\olpseganchor{OLP-0238-B007}{end}

\olpseganchor{OLP-0238-B008}{start}
هر محاسبه کېدونکے سټ په محاسبوي ډول د شمېر وړ دے۔ د دې د ليدلو
دپاره فرض کړئ~\LR{$S$} محاسبه کېدونکے دے۔ که~\LR{$S$} تش وي، نو د تعريف له
مخې په محاسبوي ډول د شمېر وړ دے۔ که نه، \LR{$a$} د~\LR{$S$} هر يو غړے وي۔
تابع~\LR{$f$} داسې تعريف کړئ:
\[
f(x) =
\begin{cases}
x & \text{\RL{که \LR{$\Char{S}(x) = 1$} وي}} \\
a & \text{\RL{که داسې نۀ وي۔}}
\end{cases}
\]
بيا~\LR{$f$} محاسبه کېدونکې تابع ده، او~\LR{$S$} د~\LR{$f$} د قيمتونو سټ دے۔
\olpseganchor{OLP-0238-B008}{end}

\olpseganchor{OLP-0239-B004}{start}
\olfileid{cmp}{thy}{eqc}
\olpseganchor{OLP-0239-B004}{end}

\olpseganchor{OLP-0239-B005}{start}
\olsection[د c.e. سټونو تعريفونه]{په محاسبوي ډول د شمېر وړ سټونو
  معادل تعريفونه}
\olpseganchor{OLP-0239-B005}{end}


\olpseganchor{OLP-0239-B006}{start}
لاندې قضيه د دې خبرې څو مهم معادل بيانونه راکوي چې يو سټ په محاسبوي
ډول د شمېر وړ وي۔
\olpseganchor{OLP-0239-B006}{end}

\olpseganchor{OLP-0239-B007}{start}
\begin{thm}
\ollabel{thm:ce-equiv}
فرض کړئ \LR{$S$} د طبيعي عددونو يو سټ دے۔ بيا لاندې خبرې سره معادلې دي:
\begin{enumerate}
\item\ollabel{case:ce} \LR{$S$} په محاسبوي ډول د شمېر وړ دے۔
\item\ollabel{case:ran-pc} \LR{$S$} د يوې \emph{جزوي} محاسبه کېدونکې تابع
د قيمتونو سټ دے۔
\item\ollabel{case:ran-prim} \LR{$S$} تش دے يا د يوې بنسټيزې بازګشتي تابع
د قيمتونو سټ دے۔
\item\ollabel{case:ce-domain} \LR{$S$} د يوې جزوي محاسبه کېدونکې تابع
\emph{د تعريف ساحه} ده۔
\end{enumerate}
\end{thm}
\olpseganchor{OLP-0239-B007}{end}

\olpseganchor{OLP-0239-B008}{start}
\begin{explain}
لومړۍ درې فقرې وايي چې هر ناتش په محاسبوي ډول د شمېر وړ سټ په معادل
ډول د يوې محاسبه کېدونکې تابع، جزوي محاسبه کېدونکې تابع، يا بنسټيزې
بازګشتي تابع په وسيله شمېرل کېدے شي۔ څلورمه فقره وايي چې که~\LR{$S$} په
محاسبوي ډول د شمېر وړ وي، نو د کوم شاخص~\LR{$e$} دپاره
\[
S = \Setabs{x}{\cfind{e}(x) \fdefined}.
\]
په بله وينا، \LR{$S$} د هغو ننوتونو سټ دے چې د~\LR{$\cfind{e}$} محاسبه پرې
درېږي۔ له همدې کبله په محاسبوي ډول د شمېر وړ سټونه کله ناکله
\emph{نيمه د پرېکړې وړ} بلل کېږي: که يو عدد په سټ کښې وي، بالاخره
``هو'' تر لاسه کوئ؛ خو که پکښې نۀ وي، هېڅکله ``نه'' نۀ تر لاسه
کوئ!{}
\end{explain}
\olpseganchor{OLP-0239-B008}{end}

\olpseganchor{OLP-0239-B009}{start}
\begin{proof}
څرنګه چې هره بنسټيزه بازګشتي تابع محاسبه کېدونکې ده او هره محاسبه
کېدونکې تابع جزوي محاسبه کېدونکې ده، \olref{case:ran-prim} پر
\olref{case:ce} دلالت کوي او \olref{case:ce} پر
\olref{case:ran-pc} دلالت کوي۔ (پام وکړئ چې که~\LR{$S$} تش وي، نو~\LR{$S$} د
هغې جزوي محاسبه کېدونکې تابع د قيمتونو سټ دے چې هېڅ ځاے تعريف شوې
نۀ ده۔) که وښيو چې \olref{case:ran-pc} پر
\olref{case:ran-prim} دلالت کوي، لومړۍ درې فقرې مو سره معادلې ښودلې
دي۔
\olpseganchor{OLP-0239-B009}{end}

\olpseganchor{OLP-0239-B010}{start}
نو فرض کړئ \LR{$S$} د جزوي محاسبه کېدونکې تابع~\LR{$\cfind{e}$} د قيمتونو سټ
دے۔ که~\LR{$S$} تش وي، کار بشپړ دے۔ که نه، \LR{$a$} د~\LR{$S$} هر يو غړے وي۔ د
کليني د نورمال شکل د قضيې له مخې ليکلے شو:
\[
\cfind{e}(x) = U(\umin{s}{T(e, x, s)}).
\]
په ځانګړي ډول، \LR{$\cfind{e}(x) \fdefined$} او \LR{$= y$} هله او يوازې هله
چې داسې~\LR{$s$} وي چې \LR{$T(e, x, s)$} او \LR{$U(s) = y$} وي۔ تابع~\LR{$f(z)$} داسې
تعريف کړئ:
\[
f(z) = \begin{cases}
  U((z)_1) & \text{\RL{که \LR{$T(e, (z)_0, (z)_1)$} وي}} \\
  a        & \text{\RL{که داسې نۀ وي۔}}
\end{cases}
\]
بيا~\LR{$f$} بنسټيزه بازګشتي ده، ځکه~\LR{$T$} او~\LR{$U$} داسې دي۔
د ټيورينګ ماشينونو په ژبه، که~\LR{$z$} داسې جوړه
  \LR{$\tuple{(z)_0, (z)_1}$} کوډوي چې \LR{$(z)_1$} د ماشين~\LR{$M_e$} پر
  ننوت~\LR{$(z)_0$} يوه درېدلې محاسبه وي، نو~\LR{$f$} د محاسبې وت راګرځوي؛
  که نه، \LR{$a$} راګرځوي۔
\olpseganchor{OLP-0239-B010}{end}

\olpseganchor{OLP-0239-B011}{start}
بايد وښيو چې~\LR{$S$} د~\LR{$f$} د قيمتونو سټ دے؛ يعنې د هر طبيعي عدد~\LR{$y$}
دپاره، \LR{$y \in S$} هله او يوازې هله چې د~\LR{$f$} په قيمتونو کښې وي۔ په
مخکښ لوري کښې فرض کړئ \LR{$y \in S$}۔ بيا~\LR{$y$} د~\LR{$\cfind{e}$} په قيمتونو
کښې دے، نو د ځينو~\LR{$x$} او~\LR{$s$} دپاره \LR{$T(e,x,s)$} پوره کېږي او
\LR{$U(s) = y$} ده۔ خو بيا \LR{$y = f(\tuple{x,s})$}۔ بل پلو، فرض کړئ~\LR{$y$} د
~\LR{$f$} په قيمتونو کښې دے۔ بيا يا \LR{$y = a$}، يا د کوم~\LR{$z$} دپاره
\LR{$T(e,(z)_0,(z)_1)$} او \LR{$U((z)_1) = y$} دي۔ څرنګه چې په وروستي حالت
کښې \LR{$\cfind{e}((z)_0) \fdefined = y$} ده، په دواړو حالونو کښې~\LR{$y$}
په~\LR{$S$} کښې دے۔
\olpseganchor{OLP-0239-B011}{end}

\olpseganchor{OLP-0239-B012}{start}
(نښه \LR{$\cfind{e}(x) \fdefined = y$} دا مانا لري چې ``\LR{$\cfind{e}(x)$}
تعريف شوې او له~\LR{$y$} سره برابره ده۔'' موږ په همدې اسانتيا
\LR{$\cfind{e}(x) = y$} هم کارولے شو؛ خو زيات غشے کله ناکله دا رايادوي
چې له يوې جزوي تابع سره کار لرو۔)
\olpseganchor{OLP-0239-B012}{end}

\olpseganchor{OLP-0239-B013}{start}
د \olref{thm:ce-equiv} د ثبوت د بشپړولو دپاره دومره بسنه کوي چې
وښيو \olref{case:ce} او~\olref{case:ce-domain} سره معادل دي۔ لومړے
ښيو چې \olref{case:ce} پر~\olref{case:ce-domain} دلالت کوي۔ که ورکړل
شوے سټ تش وي، د هغې جزوي محاسبه کېدونکې تابع د تعريف ساحه ده چې هېڅ
ځاے تعريف شوې نۀ ده۔ اوس د ناتش حالت دپاره فرض کړئ~\LR{$S$} د يوې محاسبه
کېدونکې تابع~\LR{$f$} د قيمتونو سټ دے، يعنې
\[
S = \Setabs{y}{\text{\RL{د کوم \LR{$x$} دپاره، }} f(x) = y}.
\]
وټاکئ:
\[
g(y) = \umin{x}{(f(x) = y)}.
\]
بيا~\LR{$g$} جزوي محاسبه کېدونکې تابع ده، او \LR{$g(y)$} هله او يوازې هله
تعريف شوې ده چې د کوم~\LR{$x$} دپاره \LR{$f(x) = y$} وي۔ په بله وينا، د~\LR{$g$}
د تعريف ساحه د~\LR{$f$} د قيمتونو سټ دے۔
د ټيورينګ ماشينونو په ژبه: که ټيورينګ ماشين~\LR{$F$} د~\LR{$S$}
غړي شمېري، نو ټيورينګ ماشين~\LR{$G$} داسې وټاکئ چې په~\LR{$S$} کښې د ورکړل
شوي غړي د موندلو دپاره د~\LR{$F$} په وتونو کښې پلټنه وکړي؛ که ئې ومومي
ودرېږي، او که ئې نۀ مومي د تل دپاره پلټنې ته دوام ورکړي۔
\olpseganchor{OLP-0239-B013}{end}

\olpseganchor{OLP-0239-B014}{start}
په پای کښې، د دې د ښودلو دپاره چې \olref{case:ce-domain} پر
\olref{case:ce} دلالت کوي، فرض کړئ~\LR{$S$} د جزوي محاسبه کېدونکې
تابع~\LR{$\cfind{e}$} د تعريف ساحه ده، يعنې
\[
S = \Setabs{x}{\cfind{e}(x) \fdefined}.
\]
که~\LR{$S$} تش وي، کار بشپړ دے؛ که نه، \LR{$a$} د~\LR{$S$} هر يو غړے وي۔ تابع~\LR{$f$}
داسې تعريف کړئ:
\[
f(z) = \begin{cases}
(z)_0 & \text{\RL{که \LR{$T(e,(z)_0,(z)_1)$} وي}} \\
a & \text{\RL{که داسې نۀ وي۔}}
\end{cases}
\]
بيا، لکه پورته، عدد~\LR{$x$} هله او يوازې هله د~\LR{$f$} په قيمتونو کښې دے
چې \LR{$\cfind{e}(x) \fdefined$} وي، يعنې هله او يوازې هله چې \LR{$x \in S$}
وي۔ د ټيورينګ ماشينونو په ژبه: که ماشين~\LR{$M_e$} د~\LR{$S$}
نيمه پرېکړه کوي، د~\LR{$S$} غړي د ټيورينګ ماشينونو د ټولو ممکنه محاسبو
په کتلو او د درېدلو محاسبو د اړوندو ننوتونو په راګرځولو سره وشمېرئ۔
\end{proof}
\olpseganchor{OLP-0239-B014}{end}

\olpseganchor{OLP-0239-B015}{start}
د \olref{thm:ce-equiv} فقره~\olref{case:ce-domain} موږ ته په
محاسبوي ډول د شمېر وړ سټونو د شمېرلو يوه اسانه لار راکوي: د هر~\LR{$e$}
دپاره~\LR{$W_e$} د~\LR{$\cfind{e}$} د تعريف ساحه وټاکئ، يعنې
\[
W_e = \Setabs{x}{\cfind{e}(x) \fdefined}.
\]
بيا که~\LR{$A$} هر په محاسبوي ډول د شمېر وړ سټ وي، د کوم~\LR{$e$} دپاره
\LR{$A = W_e$} وي۔
\olpseganchor{OLP-0239-B015}{end}

\olpseganchor{OLP-0239-B016}{start}
لاندې قضيه په محاسبوي ډول د شمېر وړ سټونو يوه بله ځانګړنه راکوي۔
\olpseganchor{OLP-0239-B016}{end}

\olpseganchor{OLP-0239-B017}{start}
\begin{thm}
\ollabel{thm:exists-char}
سټ~\LR{$S$} هله او يوازې هله په محاسبوي ډول د شمېر وړ دے چې داسې محاسبه
کېدونکې اړيکه~\LR{$R(x,y)$} وي چې
\[
S = \Setabs{ x }{ \lexists[y][R(x,y)] }.
\]
\end{thm}
\olpseganchor{OLP-0239-B017}{end}

\olpseganchor{OLP-0239-B018}{start}
\begin{proof}
په مخکښ لوري کښې فرض کړئ~\LR{$S$} په محاسبوي ډول د شمېر وړ دے۔ نو د کوم
~\LR{$e$} دپاره \LR{$S = W_e$}۔ د~\LR{$e$} د همدې قيمت دپاره~\LR{$S$} داسې ليکلے شو:
\[
S = \Setabs{ x }{ \lexists[y][T(e, x, y)] }.
\]
په بېرته لوري کښې فرض کړئ \LR{$S = \Setabs{ x }{
  \lexists[y][R(x, y)] }$}۔ تابع~\LR{$f$} داسې تعريف کړئ:
\[
f(x) \simeq \umin{y}{R(x, y)}.
\]
بيا~\LR{$f$} جزوي محاسبه کېدونکې ده، او~\LR{$S$} د~\LR{$f$} د تعريف ساحه ده۔
\end{proof}
\olpseganchor{OLP-0239-B018}{end}

\olpseganchor{OLP-0240-B004}{start}
\olfileid{cmp}{thy}{ncp}
\olsection{محاسبه نۀ کېدونکي سټونه شته}
\olpseganchor{OLP-0240-B004}{end}


\olpseganchor{OLP-0240-B005}{start}
پورته مو وليدل چې هر محاسبه کېدونکے سټ په محاسبوي ډول د شمېر وړ
دے۔ ايا د دې معکوس هم رښتيا دے؟ لاندې قضيه ښيي چې په عمومي ډول داسې
نۀ ده۔
\olpseganchor{OLP-0240-B005}{end}

\olpseganchor{OLP-0240-B006}{start}
\begin{thm}\ollabel{thm:K-0}
فرض کړئ \LR{$K_0$} سټ \LR{$\Setabs{\tuple{e, x}}{\cfind{e}(x) \fdefined}$}
دے۔ بيا \LR{$K_0$} په محاسبوي ډول د شمېر وړ دے، خو محاسبه کېدونکے نۀ دے۔
\end{thm}
\olpseganchor{OLP-0240-B006}{end}

\olpseganchor{OLP-0240-B007}{start}
\begin{proof}
د دې د ليدلو دپاره چې \LR{$K_0$} په محاسبوي ډول د شمېر وړ دے، پام وکړئ
چې دا د هغې تابع~\LR{$f$} د تعريف ساحه ده چې داسې تعريفېږي:
\[
f(z) = \umin{y}{(\len{z} = 2 \land T((z)_0, (z)_1, y))}.
\]
ځکه که \LR{$\cfind{e}(x)$} تعريف شوې وي، \LR{$f(\tuple{e, x})$} د درېدلې
محاسبې يوه لړۍ مومي؛ که \LR{$\cfind{e}(x)$} تعريف شوې نۀ وي، نو
\LR{$f(\tuple{e, x})$} هم تعريف شوې نۀ ده؛ او که~\LR{$z$} ان يوه جوړه هم نۀ
کوډوي، نو \LR{$f(z)$} بيا هم تعريف شوې نۀ ده۔
\olpseganchor{OLP-0240-B007}{end}

\olpseganchor{OLP-0240-B008}{start}
دا حقيقت چې \LR{$K_0$} محاسبه کېدونکے نۀ دے، هماغه د درېدنې د مسئلې
ناپرېکړتيا ده، \olref[hlt]{thm:halting-problem}۔
\end{proof}
\olpseganchor{OLP-0240-B008}{end}

\olpseganchor{OLP-0240-B009}{start}
سټ~\LR{$K_0$} د هغو جوړو~\LR{$\tuple{e,x}$} سټ دے چې
\LR{$\cfind{e}(x) \fdefined$} وي، يعنې \LR{$\tuple{e,x} \in K_0$} هله او
يوازې هله چې~\LR{$\cfind{e}$} پر ننوت~\LR{$x$} تعريف شوې وي (ودرېږي)، نو دې
ته ``د درېدنې سټ'' هم ويل کېږي۔ سټ
\LR{$K = \Setabs{e}{\cfind{e}(e) \fdefined}$} د ``پر خپل ځان د درېدنې
سټ'' دے۔ دا ډېر ځله د يوې معياري ناپرېکړه کېدونکې بېلګې په توګه
کارول کېږي۔
\olpseganchor{OLP-0240-B009}{end}

\olpseganchor{OLP-0240-B010}{start}
\begin{thm}\ollabel{thm:K}
پر خپل ځان د درېدنې سټ
\LR{$K = \Setabs{e}{\cfind{e}(e) \fdefined}$}، c.e. دے خو د پرېکړې
وړ نۀ دے۔
\end{thm}
\olpseganchor{OLP-0240-B010}{end}

\olpseganchor{OLP-0240-B011}{start}
\begin{proof}
  فرض کړئ \LR{$K$} د پرېکړې وړ دے، يعنې ځانګړونکې تابع~\LR{$\Char{K}$} ئې
  محاسبه کېدونکې ده۔ وټاکئ:
  \[d(e) = \begin{cases}
  1 & \text{\RL{که \LR{$\Char{K}(e) = 0$} وي}}\\
  \fundefined & \text{\RL{که داسې نۀ وي۔}}
  \end{cases}
  \]
  فرض کړئ \LR{$k$} د~\LR{$d$} شاخص دے، يعنې \LR{$d \simeq \cfind{k}$}۔ بيا
  \LR{$d(k) \simeq \cfind{k}(k)$} ده۔ دا له هغه حقيقت سره تناقض لري چې
  \LR{$d(k) \fdefined$} هله او يوازې هله چې
  \LR{$\cfind{k}(k) \fundefined$} وي، او دا د~\LR{$d$} له تعريفه راځي۔
\olpseganchor{OLP-0240-B011}{end}

\olpseganchor{OLP-0240-B012}{start}
  \LR{$K$} د تابع \LR{$f(x) = \umin{y}{T(x,x,y)}$} د تعريف ساحه ده او ځکه
  c.e. دے۔
\end{proof}
\olpseganchor{OLP-0240-B012}{end}

\olpseganchor{OLP-0241-B004}{start}
\olfileid{cmp}{thy}{clo}
\olpseganchor{OLP-0241-B004}{end}

\olpseganchor{OLP-0241-B005}{start}
\olsection[د c.e. سټونو اتحاد او اشتراک]{په محاسبوي ډول د شمېر وړ
  سټونه د اتحاد او اشتراک لاندې تړلي دي}
\olpseganchor{OLP-0241-B005}{end}

\olpseganchor{OLP-0241-B006}{start}
لاندې قضيه د محاسبوي ډول د شمېر وړ سټونو د ټولګي ځينې تړلتيايي
خاصيتونه راکوي۔
\olpseganchor{OLP-0241-B006}{end}

\olpseganchor{OLP-0241-B007}{start}
\begin{thm}
فرض کړئ \LR{$A$} او~\LR{$B$} په محاسبوي ډول د شمېر وړ دي۔ بيا \LR{$A \cap B$} او
\LR{$A \cup B$} هم په محاسبوي ډول د شمېر وړ دي۔
\end{thm}
\olpseganchor{OLP-0241-B007}{end}

\olpseganchor{OLP-0241-B008}{start}
\begin{proof}
\olref[eqc]{thm:ce-equiv} اجازه راکوي چې په محاسبوي ډول د شمېر وړ
سټونو بېلابېلې ځانګړنې وکاروو۔ د روښانتيا دپاره به څو بېل ثبوتونه
وړاندې کړو۔
\olpseganchor{OLP-0241-B008}{end}

\olpseganchor{OLP-0241-B009}{start}
د لومړي ثبوت دپاره فرض کړئ~\LR{$A$} د محاسبه کېدونکې تابع~\LR{$f$} په وسيله
شمېرل کېږي، او~\LR{$B$} د محاسبه کېدونکې تابع~\LR{$g$} په وسيله شمېرل کېږي۔
وټاکئ:
\begin{align*}
h(x) & = \umin{y}{(f(y) = x \lor g(y) = x)} \text{\RL{ او}}\\
j(x) & = \umin{y}{(f((y)_0) = x \land g((y)_1) = x)}.
\end{align*}
بيا \LR{$A \cup B$} د~\LR{$h$} د تعريف ساحه، او \LR{$A \cap B$} د~\LR{$j$} د تعريف
ساحه ده۔
\olpseganchor{OLP-0241-B009}{end}

\olpseganchor{OLP-0241-B010}{start}
\begin{explain}
په محاسبوي ژبه کښې خبره داسې ده: که داسې کړنلارې راکړل شوې وي چې
\LR{$A$} او~\LR{$B$} شمېري، نو دا چې يو غړے~\LR{$x$} په \LR{$A \cup B$} کښې دے که نه،
په دواړو شمېرنو کښې د~\LR{$x$} په لټولو نيمه پرېکړه کولے شو؛ او دا چې
\LR{$x$} په \LR{$A \cap B$} کښې دے که نه، په دواړو شمېرنو کښې په يو وخت د~\LR{$x$}
په لټولو نيمه پرېکړه کولے شو۔
\end{explain}
\olpseganchor{OLP-0241-B010}{end}

\olpseganchor{OLP-0241-B011}{start}
د دويم ثبوت دپاره بيا فرض کړئ چې~\LR{$A$} د~\LR{$f$} او~\LR{$B$} د~\LR{$g$} په وسيله
شمېرل کېږي۔ وټاکئ:
\[
k(x) = \begin{cases}
f(x/2) & \text{\RL{که \LR{$x$} جفت وي}} \\
g((x-1)/2) & \text{\RL{که \LR{$x$} طاق وي۔}}
\end{cases}
\]
بيا~\LR{$k$}، \LR{$A \cup B$} شمېري؛ نظر دا دے چې~\LR{$k$} په نوبت د~\LR{$f$} او~\LR{$g$}
شمېرنې کاروي۔ د~\LR{$A \cap B$} شمېرل لږ ستونزمن دي۔ که \LR{$A \cap B$} تش
وي، په څرګند ډول په محاسبوي ډول د شمېر وړ دے۔ که نه، \LR{$c$} د
\LR{$A \cap B$} هر يو غړے وي، او~\LR{$l$} داسې تعريف کړئ:
\[
l(x) = \begin{cases}
f((x)_0) & \text{\RL{که \LR{$f((x)_0) = g((x)_1)$} وي}} \\
c & \text{\RL{که داسې نۀ وي۔}}
\end{cases}
\]
په محاسبوي ژبه کښې~\LR{$l$} د~\LR{$f$} او~\LR{$g$} د شمېرنو د غړو پر جوړو تېرېږي
او هره برابره جوړه چې ومومي، وت ئې راګرځوي؛ که جوړه برابره نۀ وي،
ثابت غړے~\LR{$c$} راګرځوي۔
\olpseganchor{OLP-0241-B011}{end}

\olpseganchor{OLP-0241-B012}{start}
د وروستي ثبوت دپاره فرض کړئ~\LR{$A$} د جزوي تابع~\LR{$m(x)$} \emph{د تعريف
ساحه} او~\LR{$B$} د جزوي تابع~\LR{$n(x)$} د تعريف ساحه ده۔ بيا \LR{$A \cap B$} د
جزوي تابع \LR{$m(x) + n(x)$} د تعريف ساحه ده۔
\olpseganchor{OLP-0241-B012}{end}

\olpseganchor{OLP-0241-B013}{start}
\begin{explain}
په محاسبوي ژبه کښې، که~\LR{$A$} د هغو قيمتونو سټ وي چې~\LR{$m$} پرې درېږي او
\LR{$B$} د هغو قيمتونو سټ وي چې~\LR{$n$} پرې درېږي، نو \LR{$A \cap B$} د هغو
قيمتونو سټ دے چې دواړه کړنلارې پرې درېږي۔
\end{explain}
\olpseganchor{OLP-0241-B013}{end}

\olpseganchor{OLP-0241-B014}{start}
د درېدنې د قيمتونو د سټ په توګه د~\LR{$A \cup B$} څرګندول لږ ستونزمن دي،
ځکه~\LR{$m$} او~\LR{$n$} بايد په موازي ډول تقليد شي۔ \LR{$d$} د~\LR{$m$} يو شاخص او~\LR{$e$}
د~\LR{$n$} يو شاخص وي؛ په بله وينا، \LR{$m = \cfind{d}$} او
\LR{$n = \cfind{e}$}۔ بيا \LR{$A \cup B$} د دې تابع د تعريف ساحه ده:
\[
p(x) = \umin{y}{(T(d,x,y) \lor T(e,x,y))}.
\]
\begin{explain}
په محاسبوي ژبه کښې، پر ننوت~\LR{$x$} تابع~\LR{$p$} يا د~\LR{$m$} يا د~\LR{$n$} درېدلې
محاسبه لټوي، او که يوه ئې ومومي درېږي۔
\end{explain}
\end{proof}
\olpseganchor{OLP-0241-B014}{end}

\olpseganchor{OLP-0242-B004}{start}
\olfileid{cmp}{thy}{cmp}
\olsection{په محاسبوي ډول د شمېر وړ سټونه د متمم لاندې تړلي نۀ دي}
\olpseganchor{OLP-0242-B004}{end}

\olpseganchor{OLP-0242-B005}{start}
فرض کړئ \LR{$A$} په محاسبوي ډول د شمېر وړ دے۔ ايا د~\LR{$A$} متمم،
\LR{$\Complement{A} = \Nat \setminus A$}، همېشه په محاسبوي ډول د شمېر وړ
هم دے؟ لاندې قضيه او نتيجه ښيي چې ځواب ``نه'' دے۔
\olpseganchor{OLP-0242-B005}{end}

\olpseganchor{OLP-0242-B006}{start}
\begin{thm}
\ollabel{thm:ce-comp}
فرض کړئ \LR{$A$} د طبيعي عددونو هر يو سټ دے۔ بيا \LR{$A$} هله او يوازې هله
محاسبه کېدونکے دے چې \LR{$A$} او~\LR{$\Complement{A}$} دواړه په محاسبوي ډول
د شمېر وړ وي۔
\end{thm}
\olpseganchor{OLP-0242-B006}{end}

\olpseganchor{OLP-0242-B007}{start}
\begin{proof}
مخکښ لورے اسان دے: که~\LR{$A$} محاسبه کېدونکے وي، نو~\LR{$\Complement{A}$}
هم محاسبه کېدونکے دے (\LR{$\Char{A} = 1 \tsub
\Char{\Complement{A}}$})، او ځکه دواړه په محاسبوي ډول د شمېر وړ دي۔
\olpseganchor{OLP-0242-B007}{end}

\olpseganchor{OLP-0242-B008}{start}
په بل لوري کښې فرض کړئ~\LR{$A$} او~\LR{$\Complement{A}$} دواړه په محاسبوي ډول
د شمېر وړ دي۔ \LR{$A$} د~\LR{$\cfind{d}$} د تعريف ساحه، او
\LR{$\Complement{A}$} د~\LR{$\cfind{e}$} د تعريف ساحه وي۔ تابع~\LR{$h$} داسې تعريف
کړئ:
\[
h(x) = \umin{s}{(T(d,x,s) \lor T(e,x,s))}.
\]
په بله وينا، پر ننوت~\LR{$x$} تابع~\LR{$h$} يا د~\LR{$\cfind{d}$} يا
د~\LR{$\cfind{e}$} يوه درېدلې محاسبه لټوي۔ اوس که \LR{$x \in A$} وي، په
لومړي حالت کښې به بريالۍ شي، او که \LR{$x \in \Complement{A}$} وي، په
دويم حالت کښې به بريالۍ شي۔ نو~\LR{$h$} هرځاے تعريف شوې محاسبه کېدونکې
تابع ده۔ خو اوس د هر~\LR{$x$} دپاره \LR{$x \in A$} هله او يوازې هله چې
\LR{$T(d, x, h(x))$} وي، يعنې که~\LR{$\cfind{d}$} پر همدې ننوت تعريف شوې وي۔
څرنګه چې \LR{$T(d, x, h(x))$} محاسبه کېدونکې اړيکه ده، نو~\LR{$A$}
محاسبه کېدونکے دے۔
\end{proof}
\olpseganchor{OLP-0242-B008}{end}

\olpseganchor{OLP-0242-B009}{start}
\begin{explain}
په غير رسمي محاسبوي ژبه کښې خبره لا اسانه پوهېدل کېږي: د~\LR{$A$} د
پرېکړې دپاره پر ننوت~\LR{$x$} د~\LR{$\cfind{d}$} او~\LR{$\cfind{e}$} محاسبې په
نوبت، يوه يوه ګام، پر مخ بوځئ۔ يوه به خامخا ودرېږي؛ که~\LR{$\cfind{d}$} ودرېږي، نو~\LR{$x$}
په~\LR{$A$} کښې دے، او که نه، \LR{$x$} په~\LR{$\Complement{A}$} کښې دے۔
\end{explain}
\olpseganchor{OLP-0242-B009}{end}

\olpseganchor{OLP-0242-B010}{start}
\begin{cor}
\ollabel{cor:comp-k}
\LR{$\Complement{K_0}$} په محاسبوي ډول د شمېر وړ نۀ دے۔
\end{cor}
\olpseganchor{OLP-0242-B010}{end}

\olpseganchor{OLP-0242-B011}{start}
\begin{proof}
موږ پوهېږو چې~\LR{$K_0$} په محاسبوي ډول د شمېر وړ دے، خو محاسبه کېدونکے
نۀ دے۔ که~\LR{$\Complement{K_0}$} په محاسبوي ډول د شمېر وړ وای، نو د
\olref{thm:ce-comp} له مخې به~\LR{$K_0$} محاسبه کېدونکے و، چې له
\olref[ncp]{thm:K-0} سره تناقض دے۔
\end{proof}
\olpseganchor{OLP-0242-B011}{end}

\olpseganchor{OLP-0243-B004}{start}
\olfileid{cmp}{thy}{red}
\olsection{راکمېدنه}
\olpseganchor{OLP-0243-B004}{end}

\olpseganchor{OLP-0243-B005}{start}
\begin{explain}
اوس پوهېږو چې لږ تر لږه يو سټ، \LR{$K_0$}، شته چې په محاسبوي ډول د شمېر
وړ دے خو محاسبه کېدونکے نۀ دے۔ بايد څرګنده وي چې نور هم شته۔ د
راکمېدنې طريقه يوه پياوړې لار راکوي چې د نورو سټونو دا خاصيتونه، بې
له دې چې هر ځل بنسټيزو اصولو ته ستانۀ شو، ثابت کړو۔
\olpseganchor{OLP-0243-B005}{end}

\olpseganchor{OLP-0243-B006}{start}
په عمومي ډول، له سټ~\LR{$A$} څخه سټ~\LR{$B$} ته يوه ``راکمونه'' داسې طريقه ده
چې د دې پوښتنې ځوابونه چې غړي په~\LR{$B$} کښې دي که نه، د دې
پوښتنې په ځوابونو بدلوي چې غړي په~\LR{$A$} کښې دي که نه۔ موږ به
د ``ډېر-پر-يو راکمېدنې'' په نوم مفهوم ته پام وکړو، خو د بېلابېلو
خاصيتونو لرونکي د راکمېدنې ډېر نور مفهومونه هم شته۔ د راکمېدنې
مفهومونه د محاسبوي پېچلتيا په څېړنه کښې هم مرکزي دي، هلته د اغېزمنتيا
مسئلې هم بايد په پام کښې ونيول شي۔ د بېلګې په توګه، يو سټ هغه وخت
``اېن‌پي-بشپړ'' بلل کېږي چې په اېن‌پي کښې وي او د اېن‌پي هره مسئله
ورته راکمول کېدے شي؛ دلته د لاندې مفهوم ته ورته راکمونه کارېږي، خو
دا زيات شرط لري چې راکمونه په څوحدي وخت کښې محاسبه شي۔
\olpseganchor{OLP-0243-B006}{end}

\olpseganchor{OLP-0243-B007}{start}
موږ د راکمونې مفهوم لا له مخکښې په ضمني ډول کارولے دے۔ سټ~\LR{$K$} داسې
تعريف کړئ:
\[
K = \Setabs{x}{\cfind{x}(x) \fdefined},
\]
يعنې \LR{$K = \Setabs{x}{x \in W_x}$}۔ زموږ هغه ثبوت چې د درېدنې مسئله
نۀ حل کېږي (\olref[hlt]{thm:halting-problem})، تر ټولو نېغ دا ښيي چې
\LR{$K$} محاسبه کېدونکے نۀ دے۔ راياد کړئ چې~\LR{$K_0$} دا سټ دے:
\[
K_0 = \Setabs{\tuple{e, x}}{\cfind{e}(x) \fdefined },
\]
يعنې \LR{$K_0 = \Setabs{\tuple{e,x}}{x \in W_e}$}۔ د~\LR{$K$} د نۀ محاسبه
کېدنې هر ثبوت د~\LR{$K_0$} نۀ محاسبه کېدنې ته غځول اسان دي: که~\LR{$K_0$}
محاسبه کېدونکے وای، نو دا مو پرېکړه کولے شوه چې غړے~\LR{$x$} په
\LR{$K$} کښې دے که نه، يوازې په دې پوښتنې سره چې جوړه~\LR{$\tuple{x, x}$} په
\LR{$K_0$} کښې ده که نه۔ هغه تابع~\LR{$f$} چې~\LR{$x$} په~\LR{$\tuple{x, x}$} بدلوي،
له~\LR{$K$} څخه~\LR{$K_0$} ته د \emph{راکمونې} يوه بېلګه ده۔
\end{explain}
\olpseganchor{OLP-0243-B007}{end}

\olpseganchor{OLP-0243-B008}{start}
\begin{defn}
فرض کړئ \LR{$A$} او~\LR{$B$} د طبيعي عددونو سټونه دي۔ محاسبه کېدونکې تابع
~\LR{$f\colon \Nat \to \Nat$} له~\LR{$A$} څخه~\LR{$B$} ته \emph{ډېر-پر-يو راکمونه}
ده، هله او يوازې هله چې د هر طبيعي عدد~\LR{$x$} دپاره
\[
x \in A \quad \text{\RL{هله او يوازې هله چې}} \quad f(x) \in B.
\]
که داسې راکمونه~\LR{$f$} شته، وايو~\LR{$A$}، \LR{$B$} ته \emph{ډېر-پر-يو راکمېدونکے}
دے، چې \LR{$A \leq_m B$} ليکل کېږي۔ که~\LR{$A$}، \LR{$B$} ته ډېر-پر-يو راکمېدونکے
وي او برعکس هم، نو~\LR{$A$} او~\LR{$B$} \emph{ډېر-پر-يو معادل} بلل کېږي، چې
\LR{$A \equiv_m B$} ليکل کېږي۔
\end{defn}
\olpseganchor{OLP-0243-B008}{end}

\olpseganchor{OLP-0243-B009}{start}
که د پورته تعريف تابع~\LR{$f$} اتفاقاً يو په يو وي، نو~\LR{$A$}، \LR{$B$} ته
\emph{يو-پر-يو راکمېدونکے} بلل کېږي۔ لاندې ډېرې راکمونې دا پياوړے
شرط پوره کوي، خو موږ به دا حقيقت نۀ کاروو۔
\olpseganchor{OLP-0243-B009}{end}

\olpseganchor{OLP-0243-B010}{start}
\begin{digress}
دا رښتيا ده، خو هېڅ څرګنده نۀ ده، چې يو-پر-يو راکمېدنه په رښتيا د
ډېر-پر-يو راکمېدنې په پرتله پياوړے شرط دے۔ په بله وينا، داسې نامتناهي
سټونه~\LR{$A$} او~\LR{$B$} شته چې~\LR{$A$}، \LR{$B$} ته ډېر-پر-يو راکمېدونکے وي خو
~\LR{$B$} ته يو-پر-يو راکمېدونکے نۀ وي۔
\end{digress}
\olpseganchor{OLP-0243-B010}{end}

\olpseganchor{OLP-0244-B004}{start}
\olfileid{cmp}{thy}{ppr}
\olsection{د راکمېدنې خاصيتونه}
\olpseganchor{OLP-0244-B004}{end}

\olpseganchor{OLP-0244-B005}{start}
که~\LR{$A$}، \LR{$B$} ته راکمېږي، \LR{$A \leq_m B$} ليکو۔ دا نښه وړانديز کوي چې
که \LR{$A \leq_m B$} وي، نو~\LR{$A$} له~\LR{$B$} څخه ``زيات ستونزمن نۀ دے''، او
\LR{$B$} د~\LR{$A$} ``هومره يا ترې زيات ستونزمن دے''؛ او~\LR{$\le_m$} د عددونو د
معمولي~\LR{$\le$} په شان سټونه (يا د پرېکړې مسئلې) د پېچلتيا له مخې
مرتبوي۔ لاندې دوه قضيې د دې شهود ملاتړ کوي۔ لومړۍ وايي چې~\LR{$\le_m$}
انتقالي ده۔
\olpseganchor{OLP-0244-B005}{end}

\olpseganchor{OLP-0244-B006}{start}
\begin{prop}
\ollabel{prop:trans-red}
که \LR{$A \leq_m B$} او \LR{$B \leq_m C$} وي، نو \LR{$A \leq_m C$}۔
\end{prop}
\olpseganchor{OLP-0244-B006}{end}

\olpseganchor{OLP-0244-B007}{start}
\begin{proof}
له~\LR{$A$} څخه~\LR{$B$} ته د راکمونې او له~\LR{$B$} څخه~\LR{$C$} ته د راکمونې ترکيب
له~\LR{$A$} څخه~\LR{$C$} ته راکمونه ورکوي۔
\end{proof}
\olpseganchor{OLP-0244-B007}{end}

\olpseganchor{OLP-0244-B008}{start}
\begin{prob}
\olref{prop:trans-red} په دې ښودلو ثابت کړئ چې که~\LR{$f$} او~\LR{$g$} په
ترتيب سره له~\LR{$A$} څخه~\LR{$B$} او له~\LR{$B$} څخه~\LR{$C$} ته ډېر-پر-يو راکمونې
وي، نو~\LR{$\comp{f}{g}$} له~\LR{$A$} څخه~\LR{$C$} ته ډېر-پر-يو راکمونه ده۔
\end{prob}
\olpseganchor{OLP-0244-B008}{end}

\olpseganchor{OLP-0244-B009}{start}
\begin{prop}
\ollabel{prop:reduce}
فرض کړئ \LR{$A$} او~\LR{$B$} هر سټونه دي، او \LR{$A \leq_m B$}۔
\begin{enumerate}
\item که~\LR{$B$} په محاسبوي ډول د شمېر وړ وي، نو~\LR{$A$} هم داسې دے۔
\item که~\LR{$B$} محاسبه کېدونکے وي، نو~\LR{$A$} هم داسې دے۔
\end{enumerate}
\end{prop}
\olpseganchor{OLP-0244-B009}{end}

\olpseganchor{OLP-0244-B010}{start}
\begin{proof}
\LR{$f$} له~\LR{$A$} څخه~\LR{$B$} ته ډېر-پر-يو راکمونه وي۔ د لومړۍ دعوې دپاره بس
دا وازمويئ چې که~\LR{$B$} د جزوي تابع~\LR{$g$} د تعريف ساحه وي، نو~\LR{$A$} د
~\LR{$\comp{f}{g}$} د تعريف ساحه ده:
\begin{align*}
x \in A & \text{\RL{ هله او يوازې هله چې }} f(x) \in B \\
& \text{\RL{ هله او يوازې هله چې }}  g(f(x)) \fdefined.
\end{align*}
\olpseganchor{OLP-0244-B010}{end}

\olpseganchor{OLP-0244-B011}{start}
د دويمې دعوې دپاره راياد کړئ چې که~\LR{$B$} محاسبه کېدونکے وي، نو~\LR{$B$}
او~\LR{$\Complement{B}$} په محاسبوي ډول د شمېر وړ دي
(\olref[cmp]{thm:ce-comp})۔ دا ازمويل ستونزمن نۀ دي چې~\LR{$f$} له
\LR{$\Complement{A}$} څخه~\LR{$\Complement{B}$} ته هم ډېر-پر-يو راکمونه ده؛
نو د دې ثبوت د لومړۍ برخې له مخې~\LR{$A$} او~\LR{$\Complement{A}$} دواړه
په محاسبوي ډول د شمېر وړ دي۔ ځکه~\LR{$A$} هم د
\olref[cmp]{thm:ce-comp} له مخې محاسبه کېدونکے دے۔ (په بديل ډول،
کتلے شئ چې \LR{$\Char{A} = \comp{f}{\Char{B}}$}؛ نو که~\LR{$\Char{B}$} محاسبه
کېدونکې وي، \LR{$\Char{A}$} هم محاسبه کېدونکې ده۔)
\end{proof}
\olpseganchor{OLP-0244-B011}{end}

\olpseganchor{OLP-0244-B012}{start}
\begin{prob}
فرض کړئ~\LR{$f$} له~\LR{$A$} څخه~\LR{$B$} ته ډېر-پر-يو راکمونه ده۔ وښيئ چې~\LR{$f$} له
\LR{$\Complement{A}$} څخه~\LR{$\Complement{B}$} ته هم ډېر-پر-يو راکمونه ده۔
\end{prob}
\olpseganchor{OLP-0244-B012}{end}

\olpseganchor{OLP-0244-B013}{start}
\begin{prob}
وښيئ چې که \LR{$f\colon \Nat \to \Nat$} ډېر-پر-يو راکمونه وي، نو
\LR{$\Char{A} = \comp{f}{\Char{B}}$}۔
\end{prob}
\olpseganchor{OLP-0244-B013}{end}

\olpseganchor{OLP-0244-B014}{start}
\begin{digress}
د \emph{ټيورينګ راکمېدنې} په نوم د راکمېدنې يو ډېر عمومي مفهوم په
نورو سياقونو کښې ګټور دے، په ځانګړي ډول د ناپرېکړتيا د پايلو په
ثبوت کښې۔ پام وکړئ چې د \olref[cmp]{cor:comp-k} له مخې د~\LR{$K_0$} متمم
~\LR{$K_0$} ته نۀ راکمېږي، ځکه په محاسبوي ډول د شمېر وړ نۀ دے۔ خو په
شهودي ډول، که د~\LR{$K_0$} په اړه د پوښتنو ځوابونه مو پېژندل، د هغې د
متمم په اړه د پوښتنو ځوابونه به مو هم پېژندل۔ سټ~\LR{$A$} هغه وخت~\LR{$B$} ته
د ټيورينګ له مخې راکمېدونکے بلل کېږي چې په~\LR{$A$} کښې د پوښتنو ځوابونه
د داسې محاسبه کېدونکې کړنلارې په وسيله وټاکل شي چې د~\LR{$B$} په اړه
پوښتنې کولے شي۔ دا له ډېر-پر-يو راکمېدنې زيات ازاد دے، ځکه هلته
(1)~يوازې د~\LR{$B$} په اړه يوه پوښتنه کولے شئ، او (2)~د ``هو'' ځواب بايد
د~\LR{$A$} د پوښتنې په ``هو'' ځواب، او همداسې ``نه'' په ``نه''، بدل شي۔
بيا هم که~\LR{$A$}، \LR{$B$} ته د ټيورينګ له مخې راکمېدونکے وي او~\LR{$B$} محاسبه
کېدونکے وي، نو~\LR{$A$} هم محاسبه کېدونکے دے (که څه هم، لکه وليدل مو،
ورته خبره د محاسبوي شمېر وړتيا دپاره نۀ پوره کېږي)۔
\olpseganchor{OLP-0244-B014}{end}

\olpseganchor{OLP-0244-B015}{start}
بايد د راکمېدنې د هغو بېلابېلو مفهومونو په اړه فکر وکړئ چې بحث مو
پرې وکړ، او توپيرونه ئې درک کړئ۔ خو په دې څپرکي کښې به يوازې له
ډېر-پر-يو راکمېدنې سره کار کوو۔ په ضمني ډول، په تېر پراګراف کښې د
راکمېدنې دواړه ډولونه په محاسبوي پېچلتيا کښې ورته بڼې لري، خو دا
زيات شرط ورسره وي چې ټيورينګ ماشينونه په څوحدي وخت کښې وچلېږي: د
ډېر-پر-يو راکمېدنې پېچلتيايي بڼه \emph{کارپ راکمېدنه}، او د ټيورينګ
راکمېدنې پېچلتيايي بڼه \emph{کوک راکمېدنه} بلل کېږي۔
\end{digress}
\olpseganchor{OLP-0244-B015}{end}

\olpseganchor{OLP-0245-B004}{start}
\olfileid{cmp}{thy}{cce}
\olsection{بشپړ په محاسبوي ډول د شمېر وړ سټونه}
\olpseganchor{OLP-0245-B004}{end}

\olpseganchor{OLP-0245-B005}{start}
\begin{defn}
سټ~\LR{$A$} (د ډېر-پر-يو راکمېدنې لاندې) يو
\emph{بشپړ په محاسبوي ډول د شمېر وړ سټ} دے، که
\begin{enumerate}
\item \LR{$A$} په محاسبوي ډول د شمېر وړ وي، او
\item د هر بل په محاسبوي ډول د شمېر وړ سټ~\LR{$B$} دپاره \LR{$B \leq_m A$} وي۔
\end{enumerate}
\end{defn}
\olpseganchor{OLP-0245-B005}{end}

\olpseganchor{OLP-0245-B006}{start}
په بله وينا، بشپړ په محاسبوي ډول د شمېر وړ سټونه د ټولو ممکنه په
محاسبوي ډول د شمېر وړ سټونو ``تر ټولو ستونزمن'' دي۔ دوي اجازه راکوي
چې د \emph{هر} په محاسبوي ډول د شمېر وړ سټ په اړه پوښتنو ته ځواب
ووايو۔
\olpseganchor{OLP-0245-B006}{end}

\olpseganchor{OLP-0245-B007}{start}
\begin{thm}
\LR{$K$}، \LR{$K_0$} او~\LR{$K_1$} ټول بشپړ په محاسبوي ډول د شمېر وړ سټونه دي۔
\end{thm}
\olpseganchor{OLP-0245-B007}{end}

\olpseganchor{OLP-0245-B008}{start}
\begin{proof}
د دې د ليدلو دپاره چې~\LR{$K_0$} بشپړ دے، \LR{$B$} هر په محاسبوي ډول د شمېر
وړ سټ وي۔ نو د کوم شاخص~\LR{$e$} دپاره
\[
B = W_e = \Setabs{x}{\cfind{e}(x) \fdefined}.
\]
\LR{$f$} هغه تابع وي چې \LR{$f(x) = \tuple{e, x}$}۔ بيا د هر طبيعي عدد~\LR{$x$}
دپاره \LR{$x \in B$} هله او يوازې هله چې \LR{$f(x) \in K_0$} وي۔ په بله
وينا، \LR{$f$}، \LR{$B$}، \LR{$K_0$} ته راکموي۔
\olpseganchor{OLP-0245-B008}{end}

\olpseganchor{OLP-0245-B009}{start}
د دې د ليدلو دپاره چې~\LR{$K_1$} بشپړ دے، پام وکړئ چې د
\olref[k1]{prop:k1} په ثبوت کښې مو~\LR{$K_0$} ورته راکم کړے و۔ نو د
\olref[ppr]{prop:trans-red} له مخې هر په محاسبوي ډول د شمېر وړ سټ
~\LR{$K_1$} ته هم راکمول کېدے شي۔
\olpseganchor{OLP-0245-B009}{end}

\olpseganchor{OLP-0245-B010}{start}
\LR{$K_0$} د درېدنې د جوړو سټ دے او~\LR{$K$} ته راکمول کېدے شي۔
د هرې داسې جوړې لپاره، د پاراميټرونو له قضيې څخه د داسې
پروګرام شاخص په محاسبه موندلے شو چې خپل ننوت پرېږدي او د
لومړي مختص پروګرام پر دويم مختص چلوي۔ نو نوی پروګرام پر خپل
شاخص هله او يوازې هله درېږي چې د لومړي مختص پروګرام پر دويم
مختص ودرېږي۔ په دې ډول هرې جوړې ته همدا نوی شاخص ورمنسوبوو؛
له انتقاليت څخه د قطري سټ بشپړتيا راځي۔
\end{proof}
\olpseganchor{OLP-0245-B010}{end}

\olpseganchor{OLP-0245-B011}{start}
\begin{prob}
له~\LR{$K$} څخه~\LR{$K_0$} ته يوه راکمونه ورکړئ۔
\end{prob}
\olpseganchor{OLP-0245-B011}{end}

\olpseganchor{OLP-0245-B012}{start}
\begin{digress}
نو څرګنده شوه چې تر اوسه مو د په محاسبوي ډول د شمېر وړ سټونو هره
کتلې بېلګه يا محاسبه کېدونکې ده، يا بشپړه۔ دا بايد عجيبه وبرېښي!{}
ايا داسې په محاسبوي ډول د شمېر وړ سټونه شته چې نۀ محاسبه کېدونکي او
نۀ بشپړ وي؟ ځواب هو دے، خو دا تر هغه وخته نۀ وه ثابته شوې چې د
۱۹۵۰مو کلونو په منځ کښې فريډبرګ او موچنيک په خپلواکه توګه ثابته کړه۔
\end{digress}
\olpseganchor{OLP-0245-B012}{end}

\olpseganchor{OLP-0246-B004}{start}
\olfileid{cmp}{thy}{k1}
\olsection{د راکمېدنې يوه بېلګه}
\olpseganchor{OLP-0246-B004}{end}

\olpseganchor{OLP-0246-B005}{start}
راځئ د \olref[ppr]{prop:reduce} يوه کارونه وګورو۔
\olpseganchor{OLP-0246-B005}{end}

\olpseganchor{OLP-0246-B006}{start}
\begin{prop}
\ollabel{prop:k1}
فرض کړئ
\[
K_1 = \Setabs{e}{\cfind{e}(0) \fdefined}.
\]
بيا \LR{$K_1$} په محاسبوي ډول د شمېر وړ دے، خو محاسبه کېدونکے نۀ دے۔
\end{prop}
\olpseganchor{OLP-0246-B006}{end}

\olpseganchor{OLP-0246-B007}{start}
\begin{proof}
څرنګه چې \LR{$K_1 = \Setabs{e}{\lexists[s][T(e,0,s)]}$}، نو د
\olref[eqc]{thm:exists-char} له مخې~\LR{$K_1$} په محاسبوي ډول د شمېر وړ
دے۔
\olpseganchor{OLP-0246-B007}{end}

\olpseganchor{OLP-0246-B008}{start}
د دې د ښودلو دپاره چې~\LR{$K_1$} محاسبه کېدونکے نۀ دے، راځئ وښيو چې
\LR{$K_0$} ورته راکمېدونکے دے۔
\olpseganchor{OLP-0246-B008}{end}

\olpseganchor{OLP-0246-B009}{start}
\begin{explain}
دا لږ ستونزمن دے، ځکه د~\LR{$K_1$} په کارولو سره يوازې د هغو محاسبو په
اړه پوښتنې کولے شو چې له ځانګړي ننوت، \LR{$0$}، څخه پيلېږي۔ فرض کړئ يو
هوښيار ملګرے لرئ چې د دې ډول پوښتنو ځوابونه درکولے شي (داسې ملګري د
``اوراکلونو'' په نوم پېژندل کېږي)۔ بيا فرض کړئ څوک درته راشي او
وپوښتي چې~\LR{$\tuple{e, x}$} په~\LR{$K_0$} کښې دے که نه، يعنې ماشين~\LR{$e$} پر
ننوت~\LR{$x$} درېږي که نه۔ يو کار دا کولے شئ چې بل ماشين~\LR{$e_x$} جوړ کړئ؛
دا ماشين د \emph{هر} ننوت دپاره هغه ننوت له پامه غورځوي او پر ځاے ئې
~\LR{$e$} پر ننوت~\LR{$x$} چلوي۔ نو څرګنده ده چې د ماشين~\LR{$e$} پر ننوت~\LR{$x$}
درېدل له دې پوښتنې سره معادل دي چې ماشين~\LR{$e_x$} پر ننوت~\LR{$0$} (يا هر
بل ننوت) درېږي که نه۔ بيا له خپل ملګري وپوښتئ چې دا نوی ماشين~\LR{$e_x$}
پر ننوت~\LR{$0$} درېږي که نه؛ بدلې شوې پوښتنې ته د ملګري ځواب د اصلي
پوښتنې ځواب درکوي۔ دا له~\LR{$K_0$} څخه~\LR{$K_1$} ته موخه شوې راکمونه راکوي۔
\end{explain}
\olpseganchor{OLP-0246-B009}{end}

\olpseganchor{OLP-0246-B010}{start}
د نړيوالې جزوي محاسبه کېدونکې تابع په کارولو سره، \LR{$f$} هغه
درې‌ځاييزه تابع وي چې داسې تعريفېږي:
\[
f(x,y,z) \simeq \cfind{x}(y).
\]
پام وکړئ چې~\LR{$f$} خپل درېيم ننوت بيخي له پامه غورځوي۔ داسې شاخص~\LR{$e$}
وټاکئ چې \LR{$f = \cfind{e}[3]$} وي؛ نو لرو:
\[
\cfind{e}[3](x,y,z) \simeq \cfind{x}(y).
\]
د \LR{$s$}-\LR{$m$}-\LR{$n$} قضيې له مخې داسې تابع~\LR{$s(e,x,y)$} شته چې د هر~\LR{$z$}
دپاره
\begin{align*}
\cfind{s(e,x,y)}(z) & \simeq \cfind{e}[3](x,y,z) \\
& \simeq \cfind{x}(y).
\end{align*}
\olpseganchor{OLP-0246-B010}{end}

\olpseganchor{OLP-0246-B011}{start}
\begin{explain}
د پورته غير رسمي استدلال په ژبه، \LR{$s(e,x,y)$} د هغه ماشين شاخص دے چې
د هر ننوت~\LR{$z$} دپاره هغه ننوت له پامه غورځوي او~\LR{$\cfind{x}(y)$}
محاسبه کوي۔
\end{explain}
\olpseganchor{OLP-0246-B011}{end}

\olpseganchor{OLP-0246-B012}{start}
په ځانګړي ډول، لرو:
\[
\cfind{s(e,x,y)}(0) \fdefined \quad \text{\RL{هله او يوازې هله چې}} \quad
\cfind{x}(y) \fdefined.
\]
په بله وينا، \LR{$\tuple{x, y} \in K_0$} هله او يوازې هله چې
\LR{$s(e,x,y) \in K_1$} وي۔ نو هغه تابع~\LR{$g$} چې داسې تعريفېږي:
\[
g(w) = s(e,(w)_0,(w)_1)
\]
له~\LR{$K_0$} څخه~\LR{$K_1$} ته راکمونه ده۔
\end{proof}
\olpseganchor{OLP-0246-B012}{end}

\olpseganchor{OLP-0247-B004}{start}
\olfileid{cmp}{thy}{tot}
\olsection{د هرځاے تعريف کېدو مسئله ناپرېکړه کېدونکې ده}
\olpseganchor{OLP-0247-B004}{end}

\olpseganchor{OLP-0247-B005}{start}
راځئ د~\LR{$s$}-\LR{$m$}-\LR{$n$} قضيې د کارولو يوه بله بېلګه وګورو، څو وښيو چې
يو څيز محاسبه کېدونکے نۀ دے۔ \LR{$\fn{Tot}$} د هرځاے تعريف شويو محاسبه
کېدونکو تابعو د شاخصونو سټ وي، يعنې
\[
\fn{Tot} = \Setabs{x}{\text{\RL{د هر \LR{$y$} دپاره، \LR{$\cfind{x}(y)\fdefined$}}}}.
\]
\olpseganchor{OLP-0247-B005}{end}

\olpseganchor{OLP-0247-B006}{start}
\begin{prop}
\ollabel{prop:total}
\LR{$\fn{Tot}$} محاسبه کېدونکے نۀ دے۔
\end{prop}
\olpseganchor{OLP-0247-B006}{end}

\olpseganchor{OLP-0247-B007}{start}
\begin{proof}
د دې د ليدلو دپاره چې~\LR{$\fn{Tot}$} محاسبه کېدونکے نۀ دے، دا ښودل بس
دي چې~\LR{$K$} ورته راکمېدونکے دے۔ \LR{$h(x,y)$} داسې تعريف کړئ:
\[
h(x,y) \simeq
\begin{cases}
0 & \text{\RL{که \LR{$x \in K$} وي}} \\
\fundefined & \text{\RL{که داسې نۀ وي۔}}
\end{cases}
\]
پام وکړئ چې~\LR{$h(x,y)$} بيخي پر~\LR{$y$} تکيه نۀ لري۔ دا ليدل بايد ستونزمن
نۀ وي چې~\LR{$h$} جزوي محاسبه کېدونکې ده: پر ننوتونو~\LR{$x, y$} دا~\LR{$h$} په
دې ډول محاسبه کوو چې لومړے تابع~\LR{$\cfind{x}$} پر ننوت~\LR{$x$} تقليدوو؛
که دا محاسبه ودرېږي، \LR{$h(x,y)$}، \LR{$0$} راګرځوي او درېږي۔ نو \LR{$h(x,y)$} يوازې
\LR{$\Zero(\umin{s}{T(x,x,s)})$} ده، چېرته چې~\LR{$\Zero$} ثابت صفر تابع ده۔
\olpseganchor{OLP-0247-B007}{end}

\olpseganchor{OLP-0247-B008}{start}
د~\LR{$s$}-\LR{$m$}-\LR{$n$} قضيې په کارولو داسې بنسټيزه بازګشتي تابع~\LR{$k(x)$} شته
چې د هر~\LR{$x$} او~\LR{$y$} دپاره
\[
\cfind{k(x)}(y) =
\begin{cases}
0 & \text{\RL{که \LR{$x \in K$} وي}} \\
\fundefined & \text{\RL{که داسې نۀ وي۔}}
\end{cases}
\]
نو~\LR{$\cfind{k(x)}$} هله هرځاے تعريف شوې ده چې~\LR{$x \in K$} وي، او که نه
نو تعريف شوې نۀ ده۔ په دې توګه~\LR{$k$} له~\LR{$K$} څخه~\LR{$\fn{Tot}$} ته راکمونه
ده۔
\end{proof}
\olpseganchor{OLP-0247-B008}{end}

\olpseganchor{OLP-0247-B009}{start}
\begin{digress}
په حقيقت کښې~\LR{$\fn{Tot}$} آن په محاسبوي ډول د شمېر وړ هم نۀ دے؛
پېچلتيا ئې د ``حسابي پوړيز نظام'' په لا لوړه برخه کښې پرته ده۔ خو
دلته به د دې پياوړتيا په اړه اندېښنه نۀ کوو۔
\end{digress}
\olpseganchor{OLP-0247-B009}{end}

\olpseganchor{OLP-0248-B004}{start}
\olfileid{cmp}{thy}{rce}
\olsection{د رايس قضيه}
\olpseganchor{OLP-0248-B004}{end}

\olpseganchor{OLP-0248-B005}{start}
که لږ فکر وکړئ، وبه وينئ چې د~\LR{$\fn{Tot}$} ځانګړي جزئيات د
\olref[tot]{prop:total} په ثبوت کښې هېڅ رول نۀ لري۔ موږ~\LR{$h(x,y)$}
داسې جوړ کړے و چې کټ مټ هغه وخت د ثابتې تابع~\LR{$j(y) = 0$} په شان عمل
وکړي چې~\LR{$x$} په~\LR{$K$} کښې وي؛ خو په همدغو حالاتو کښې مو د هرې بلې
جزوي محاسبه کېدونکې تابع په شان هم عمل پرې کولے شو۔ دا مشاهده موږ ته
اجازه راکوي يوه لا عمومي قضيه بيان کړو چې، په لنډ ډول، وايي د محاسبه
کېدونکو تابعو هېڅ غيرابتذالي خاصيت د پرېکړې وړ نۀ دے۔
\olpseganchor{OLP-0248-B005}{end}

\olpseganchor{OLP-0248-B006}{start}
په ياد ولرئ چې~\LR{$\cfind{0}$}، \LR{$\cfind{1}$}، \LR{$\cfind{2}$}،~\dots زموږ د
جزوي محاسبه کېدونکو تابعو معياري شمېرنه ده۔
\olpseganchor{OLP-0248-B006}{end}

\olpseganchor{OLP-0248-B007}{start}
\begin{thm}[د رايس قضيه]
  \LR{$C$} د جزوي محاسبه کېدونکو تابعو هر يو سټ وي، او \LR{$A =
  \Setabs{n}{\cfind{n} \in C}$} وي۔ که~\LR{$A$} محاسبه کېدونکے وي، نو يا
  خو~\LR{$C$} تش دے، يا~\LR{$C$} د ټولو جزوي محاسبه کېدونکو تابعو سټ دے۔
\end{thm}
\olpseganchor{OLP-0248-B007}{end}

\olpseganchor{OLP-0248-B008}{start}
\emph{شاخصي سټ} داسې سټ~\LR{$A$} دے چې که~\LR{$n$} او~\LR{$m$} هغه شاخصونه وي چې
يوه تابع ``محاسبه کوي''، نو يا~\LR{$n$} او~\LR{$m$} دواړه په~\LR{$A$} کښې دي، يا
يو هم نۀ دے۔ دا ليدل ستونزمن نۀ دي چې په قضيه کښې سټ~\LR{$A$} دا خاصيت
لري۔ برعکس، که~\LR{$A$} شاخصي سټ وي او~\LR{$C$} د هغو تابعو سټ وي چې دا
شاخصونه ئې محاسبه کوي، نو \LR{$A = \Setabs{n}{\cfind{n} \in C}$}۔
\olpseganchor{OLP-0248-B008}{end}

\olpseganchor{OLP-0248-B009}{start}
\begin{explain}
په دې اصطلاح کښې د رايس قضيه له دې وينا سره معادله ده چې هېڅ
غيرابتذالي شاخصي سټ د پرېکړې وړ نۀ دے۔ د قضيې د مطلب د پوهېدو دپاره
ګټوره ده چې د \emph{پروګرامونو} (مثلاً ستاسو په خوښه پروګرامي ژبه
کښې) او د هغو د محاسبه کوونکو تابعو تر منځ توپير روښانه کړو۔ د
پروګرامونو (شاخصونو) په اړه، چې نحوي څيزونه دي، يقيناً داسې پوښتنې
شته چې محاسبه کېدونکې دي: ايا دا پروګرام تر~150 زياتې نښې لري؟ ايا
تر~22 زياتې کرښې لري؟ ايا د ``وايل'' جمله لري؟ ايا رشته ``هېلو
ورلډ'' کله هم د ``پرنټ'' جملې د دليل په توګه راځي؟ د رايس قضيه وايي
چې د پروګرام د \emph{چلند} په اړه هېڅ غيرابتذالي پوښتنه محاسبه
کېدونکې نۀ ده۔ په دې کښې دا پوښتنې هم راځي: ايا پروګرام پر ننوت~\LR{$0$}
درېږي؟ ايا کله هم درېږي؟ ايا کله هم جفت عدد راګرځوي؟
\end{explain}
\olpseganchor{OLP-0248-B009}{end}

\olpseganchor{OLP-0248-B010}{start}
\begin{proof}[د رايس د قضيې ثبوت]
فرض کړئ~\LR{$C$} نۀ تش دے او نۀ د ټولو جزوي محاسبه کېدونکو تابعو سټ، او
~\LR{$A$} په~\LR{$C$} کښې د تابعو د شاخصونو سټ وي۔ وبه ښيو چې که~\LR{$A$} محاسبه
کېدونکے وای، د درېدنې مسئله مو حل کولے شوه؛ نو~\LR{$A$} محاسبه کېدونکے
نۀ دے۔
\olpseganchor{OLP-0248-B010}{end}

\olpseganchor{OLP-0248-B011}{start}
بې له دې چې عموميت له لاسه ورکړو، فرض کولے شو چې هغه تابع~\LR{$f$} چې
هېچرته تعريف شوې نۀ ده، په~\LR{$C$} کښې نۀ ده (که نه، په لاندې استدلال
کښې~\LR{$C$} او د هغې متمم سره بدل کړئ)۔ \LR{$g$} په~\LR{$C$} کښې هره تابع وي۔ نظر
دا دے چې که مو~\LR{$A$} پرېکړه کولے شو، د هغو شاخصونو چې~\LR{$f$} محاسبه کوي
او هغو چې~\LR{$g$} محاسبه کوي، توپير مو کولے شو؛ بيا مو همدغه وړتيا د
درېدنې د مسئلې د حل دپاره کارولے شوه۔
\olpseganchor{OLP-0248-B011}{end}

\olpseganchor{OLP-0248-B012}{start}
طريقه ئې دا ده۔ د نړيوالو جزوي محاسبه کېدونکو تابعو په کارولو دا
تابع تعريفولے شو:
\[
h(x,y) \simeq
\begin{cases}
\text{\RL{تعريف شوې نۀ ده}} & \text{\RL{که \LR{$\cfind{x}(x) \fundefined$} وي}} \\
g(y) & \text{\RL{که داسې نۀ وي۔}}
\end{cases}
\]
د~\LR{$h$} د محاسبې دپاره لومړے د~\LR{$\cfind{x}(x)$} د محاسبې هڅه کوو؛ که
دا محاسبه ودرېږي، بيا~\LR{$g(y)$} محاسبه کوو؛ او که \emph{هغه} محاسبه
ودرېږي، وت ئې راګرځوو۔ په لا صوري ډول ليکلے شو:
\[
h(x,y) \simeq \Proj{2}{0}(g(y),\fn{Un}(x,x)).
\]
چېرته چې~\LR{$\Proj{2}{0}(z_0, z_1) = z_0$} هغه~\LR{$2$}-ځاييزه پروجکشن تابع
ده چې~\LR{$0$}-م دليل راګرځوي، او محاسبه کېدونکې ده۔
\olpseganchor{OLP-0248-B012}{end}

\olpseganchor{OLP-0248-B013}{start}
بيا~\LR{$h$} د جزوي محاسبه کېدونکو تابعو ترکيب دے، او ښے اړخ کټ مټ هغه
وخت تعريف شوے او له~\LR{$g(y)$} سره برابر دے چې~\LR{$\fn{Un}(x,x)$} او~\LR{$g(y)$}
دواړه تعريف شوي وي۔
\olpseganchor{OLP-0248-B013}{end}

\olpseganchor{OLP-0248-B014}{start}
پام وکړئ چې د هر ثابت~\LR{$x$} دپاره، که~\LR{$\cfind{x}(x)$} تعريف شوې نۀ
وي، نو~\LR{$h(x,y)$} د هر~\LR{$y$} دپاره تعريف شوې نۀ ده؛ او که~\LR{$\cfind{x}(x)$}
تعريف شوې وي، نو \LR{$h(x,y) \simeq g(y)$}۔ نو د~\LR{$x$} د هر ثابت قيمت
دپاره~\LR{$h(x,y)$} يا کټ مټ د~\LR{$f$}، يا کټ مټ د~\LR{$g$} په شان عمل کوي، او
دا پرېکړه چې~\LR{$\cfind{x}(x)$} تعريف شوې ده که نه، د دې پرېکړې سره
برابره ده چې له دې دوو څخه کوم حالت دے۔ خو دا بيا د دې پرېکړې سره
برابره ده چې \LR{$h_x(y) \simeq h(x,y)$} په~\LR{$C$} کښې ده که نه، او که~\LR{$A$}
محاسبه کېدونکے وای، همدا کار مو کولے شو۔
\olpseganchor{OLP-0248-B014}{end}

\olpseganchor{OLP-0248-B015}{start}
په لا صوري ډول، څرنګه چې~\LR{$h$} جزوي محاسبه کېدونکې ده، د کوم شاخص~\LR{$e$}
دپاره له تابع~\LR{$\cfind{e}$} سره برابره ده۔ د~\LR{$s$}-\LR{$m$}-\LR{$n$} قضيې له مخې
داسې بنسټيزه بازګشتي تابع~\LR{$s$} شته چې د هر~\LR{$x$} دپاره
\LR{$\cfind{s(e,x)}(y) \simeq h_x(y)$} وي۔\LR{$x$} دپاره لرو چې که
\LR{$\cfind{x}(x) \fdefined$} وي، نو~\LR{$\cfind{s(e,x)}$} هماغه تابع ده چې
~\LR{$g$} ده، او ځکه~\LR{$s(e,x)$} په~\LR{$A$} کښې دے۔ له بلې خوا، که
\LR{$\cfind{x}(x) \uparrow$} وي، نو~\LR{$\cfind{s(e,x)}$} هماغه تابع ده چې
~\LR{$f$} ده، او ځکه~\LR{$s(e,x)$} په~\LR{$A$} کښې نۀ دے۔ په بله وينا، د هر~\LR{$x$}
دپاره \LR{$x \in K$} هله او يوازې هله چې \LR{$s(e,x) \in A$} وي۔ که~\LR{$A$}
محاسبه کېدونکے وای، ~\LR{$K$} هم محاسبه کېدونکے و، چې تناقض دے۔ نو~\LR{$A$}
محاسبه کېدونکے نۀ دے۔
\end{proof}
\olpseganchor{OLP-0248-B015}{end}

\olpseganchor{OLP-0248-B016}{start}
د رايس قضيه ډېره پياوړې ده۔ لاندې سمدستي نتيجه ئې د کارونې څو
بېلګې ښيي۔
\begin{cor}
لاندې سټونه د پرېکړې وړ نۀ دي۔
\begin{enumerate}
\item \LR{$\Setabs{x}{\text{\RL{\LR{$17$} د \LR{$\cfind{x}$} په قيمتونو کښې دے}}}$}
\item \LR{$\Setabs{x}{\text{\RL{\LR{$\cfind{x}$} ثابته ده}}}$}
\item \LR{$\Setabs{x}{\text{\RL{\LR{$\cfind{x}$} هرځاے تعريف شوې ده}}}$}
\item \LR{$\Setabs{x}{\text{\RL{هر کله چې \LR{$y < y'$}، \LR{$\cfind{x}(y) \fdefined$}، او
    که \LR{$\cfind{x}(y') \fdefined$}، نو \LR{$\cfind{x}(y) < \cfind{x}(y')$}}}}$}
\end{enumerate}
\end{cor}
\olpseganchor{OLP-0248-B016}{end}

\olpseganchor{OLP-0248-B017}{start}
\begin{proof}
  دا ټول غيرابتذالي شاخصي سټونه دي۔
\end{proof}
\olpseganchor{OLP-0248-B017}{end}

\olpseganchor{OLP-0249-B004}{start}
\olfileid{cmp}{thy}{fix}
\olsection{د ثابت ټکي قضيه}
\olpseganchor{OLP-0249-B004}{end}

\olpseganchor{OLP-0249-B005}{start}
راځئ يو ځل بيا د درېدنې مسئله وګورو۔ د لنډمهالې نښې په توګه
~\LR{$\gn{\cfind{x}(y)}$} د~\LR{$\tuple{x, y}$} دپاره وليکو؛ دا د
\LR{$\cfind{x}(y)$} د قيمت د يو ``نوم'' استازيتوب وګڼئ۔ په دې نښه د
درېدنې د مسئلې د ناپرېکړتيا يو ثبوت بيا بيانولے شو۔
\olpseganchor{OLP-0249-B005}{end}

\olpseganchor{OLP-0249-B006}{start}
پوښتنه: ايا داسې محاسبه کېدونکې تابع~\LR{$h$} شته چې لاندې خاصيت ولري؟ د
هر~\LR{$x$} او~\LR{$y$} دپاره
\[
h(\gn{\cfind{x}(y)}) =
\begin{cases}
1 & \text{\RL{که \LR{$\cfind{x}(y) \fdefined$} وي}} \\
0 & \text{\RL{که داسې نۀ وي۔}}
\end{cases}
\]
\olpseganchor{OLP-0249-B006}{end}

\olpseganchor{OLP-0249-B007}{start}
ځواب: نه؛ که نه، دا جزوي تابع
\[
g(x) \simeq
\begin{cases}
0 & \text{\RL{که \LR{$h(\gn{\cfind{x}(x)}) = 0$} وي}} \\
\text{\RL{تعريف شوې نۀ ده}} & \text{\RL{که داسې نۀ وي}}
\end{cases}
\]
به محاسبه کېدونکې وه، او ځکه به ئې کوم شاخص~\LR{$e$} درلود۔ خو بيا لرو:
\[
\cfind{e}(e) \simeq
\begin{cases}
0 & \text{\RL{که \LR{$h(\gn{\cfind{e}(e)}) = 0$} وي}} \\
\text{\RL{تعريف شوې نۀ ده}} & \text{\RL{که داسې نۀ وي،}}
\end{cases}
\]
چې په دې صورت کښې~\LR{$\cfind{e}(e)$} هله او يوازې هله تعريف شوې ده چې
تعريف شوې نۀ وي، او دا تناقض دے۔
\olpseganchor{OLP-0249-B007}{end}

\olpseganchor{OLP-0249-B008}{start}
اوس هغه معادله وګورئ چې~\LR{$\cfind{e}$} لري۔ دلته په يوه معنا د
ځان-مراجعې يوه بېلګه شته: موږ د~\LR{$\cfind{e}(e)$} قيمت داسې برابر کړے
چې په يوه ځانګړې طريقه پر~\LR{$\gn{\cfind{e}(e)}$} تکيه وکړي۔ د ثابت ټکي
قضيه وايي چې دا کار په عمومي ډول \emph{کولے شو}---يوازې د تناقضونو د
ثابتولو دپاره نۀ۔
\olpseganchor{OLP-0249-B008}{end}

\olpseganchor{OLP-0249-B009}{start}
\olref{lem:fixed-equiv} د ثابت ټکي د قضيې د بيانولو دوه معادل شکلونه
راکوي۔ د منطقي نظره، د دې بيانونو معادلتوب له دې راځي چې دواړه رښتيا
دي؛ خو زموږ اصلي موخه دا ده چې هر يو په نېغه توګه له بل څخه راځي، نو
د همدې قضيې د بديلو بيانونو په توګه اخيستل کېدے شي۔
\olpseganchor{OLP-0249-B009}{end}

\olpseganchor{OLP-0249-B010}{start}
\begin{lem}
\ollabel{lem:fixed-equiv}
لاندې بيانونه معادل دي:
\begin{enumerate}
\item د هرې جزوي محاسبه کېدونکې تابع~\LR{$g(x,y)$} دپاره داسې شاخص~\LR{$e$}
  شته چې د هر~\LR{$y$} دپاره
\[
\cfind{e}(y) \simeq g(e,y).
\]
\item د هرې محاسبه کېدونکې تابع~\LR{$f(x)$} دپاره داسې شاخص~\LR{$e$} شته چې
  د هر~\LR{$y$} دپاره
\[
\cfind{e}(y) \simeq \cfind{f(e)}(y).
\]
\end{enumerate}
\end{lem}
\olpseganchor{OLP-0249-B010}{end}

\olpseganchor{OLP-0249-B011}{start}
\begin{proof}
\LR{$(1) \Rightarrow (2)$}: ورکړې~\LR{$f$} ته~\LR{$g$} داسې تعريف کړئ چې
\LR{$g(x,y) \simeq \fn{Un}(f(x),y)$} وي۔ (1) وکاروئ څو داسې شاخص~\LR{$e$}
تر لاسه کړئ چې د هر~\LR{$y$} دپاره
\begin{align*}
\cfind{e}(y) & \simeq \fn{Un}(f(e),y) \\
& \simeq \cfind{f(e)}(y).
\end{align*}
\olpseganchor{OLP-0249-B011}{end}

\olpseganchor{OLP-0249-B012}{start}
\LR{$(2) \Rightarrow (1)$}: ورکړې~\LR{$g$} ته د~\LR{$s$}-\LR{$m$}-\LR{$n$} قضيه وکاروئ څو
داسې~\LR{$f$} تر لاسه کړئ چې د هر~\LR{$x$} او~\LR{$y$} دپاره
\LR{$\cfind{f(x)}(y) \simeq g(x,y)$} وي۔ (2) وکاروئ څو داسې شاخص~\LR{$e$} تر
لاسه کړئ چې
\begin{align*}
\cfind{e}(y) & \simeq \cfind{f(e)}(y) \\
& \simeq g(e,y).
\end{align*}
په دې سره ثبوت بشپړېږي۔
په دې لړيو کښې جزوي برابري په دې مانا ده چې د هر ننوت لپاره
يا دواړه اړخونه تعريف دي او يو ارزښت لري، يا دواړه ناتعريف دي۔
\end{proof}
\olpseganchor{OLP-0249-B012}{end}

\olpseganchor{OLP-0249-B013}{start}
\begin{explain}
مخکښې له دې چې وښيو بيان~(1) رښتيا دے (او ځکه~(2) هم)، فکر وکړئ چې
دا څومره عجيبه خبره ده۔ \LR{$e$} يو کمپيوټري پروګرام وګڼئ؛ بيان~(1) وايي
چې د هرې ورکړې جزوي محاسبه کېدونکې~\LR{$g(x,y)$} دپاره داسې کمپيوټري
پروګرام~\LR{$e$} موندلے شئ چې \LR{$g_e(y) \simeq g(e,y)$} محاسبه کوي۔ په بله
وينا، داسې کمپيوټري پروګرام موندلے شئ چې يوه پر خپل پروګرام مراجعه
کوونکې تابع محاسبه کوي۔
\end{explain}
\olpseganchor{OLP-0249-B013}{end}

\olpseganchor{OLP-0249-B014}{start}
\begin{thm}
په~\olref{lem:fixed-equiv} کښې دواړه بيانونه رښتيا دي۔ په ځانګړي
ډول، د هرې جزوي محاسبه کېدونکې تابع~\LR{$g(x,y)$} دپاره داسې شاخص~\LR{$e$}
شته چې د هر~\LR{$y$} دپاره
\[
\cfind{e}(y) \simeq g(e,y).
\]
\end{thm}
\olpseganchor{OLP-0249-B014}{end}

\olpseganchor{OLP-0249-B015}{start}
\begin{proof}
اړين توکي لا د پورته درېدنې مسئلې په بحث کښې نغښتي دي۔
\LR{$\fn{diag}(x)$} داسې محاسبه کېدونکې تابع وي چې د هر~\LR{$x$} دپاره د تابع
\LR{$f_x(y) \simeq \cfind{x}(x,y)$} يو شاخص راګرځوي، يعنې
\[
\cfind{\fn{diag}(x)}(y) \simeq \cfind{x}(x,y).
\]
\LR{$\fn{diag}$} داسې تابع وګڼئ چې د يوې دوه‌ځاييزې تابع پروګرام د يوې
يوځاييزې تابع په پروګرام بدلوي، او دا د اصلي پروګرام په لومړي دليل
کښې د هغه خپل شاخص په ثابتولو تر لاسه کېږي۔ تابع~\LR{$\fn{diag}$} په صوري
ډول داسې تعريفېدے شي: لومړے~\LR{$s$} داسې تعريف کړئ:
\[
s(x,y) \simeq \fn{Un}^2(x,x,y),
\]
چېرته چې~\LR{$\fn{Un}^2$} يوه درې‌ځاييزه تابع ده چې د جزوي محاسبه
کېدونکو دوه‌ځاييزو تابعو دپاره نړيواله ده۔ بيا د~\LR{$s$}-\LR{$m$}-\LR{$n$} قضيې
له مخې داسې بنسټيزه بازګشتي تابع~\LR{$\fn{diag}$} موندلے شو چې دا پوره
کړي:
\[
\cfind{\fn{diag}(x)}(y) \simeq s(x,y).
\]
\olpseganchor{OLP-0249-B015}{end}

\olpseganchor{OLP-0249-B016}{start}
اوس تابع~\LR{$l$} داسې تعريف کړئ:
\[
l(x,y) \simeq g(\fn{diag}(x),y).
\]
او~\LR{$\gn{l}$} د~\LR{$l$} يو شاخص وي۔ په پای کښې
\LR{$e = \fn{diag}(\gn{l})$} وټاکئ۔ بيا د هر~\LR{$y$} دپاره لرو:
\begin{align*}
\cfind{e}(y) & \simeq \cfind{\fn{diag}(\gn{l})}(y) \\
& \simeq \cfind{\gn{l}}(\gn{l}, y) \\
& \simeq l(\gn{l}, y) \\
& \simeq g(\fn{diag}(\gn{l}),y) \\
& \simeq g(e, y),
\end{align*}
لکه چې غوښتل شوي وو۔
\end{proof}
\olpseganchor{OLP-0249-B016}{end}

\olpseganchor{OLP-0249-B017}{start}
\begin{explain}
دلته څه روان دي؟ فرض کړئ تاسو ته د داسې کمپيوټري پروګرام د ليکلو
دنده درکړل شوې چې خپل متن چاپ کړي۔ بيا فرض کړئ چې په داسې پروګرامي
ژبه کار کوئ چې د رشتو يو بډای او عجيب کتابتون لري۔ په ځانګړي ډول،
فرض کړئ ستاسو پروګرامي ژبه داسې تابع~\LR{$\fn{diag}$} لري: تابع
\LR{$\fn{diag}$} يوه ننوت رشته~\LR{$s$} اخلي، په~\LR{$s$} کښې د `x' نښې هره بېلګه
مومي، او هغه د اصلي
رشتې په اقتباس شوې بڼه بدلوي۔ د بېلګې په توګه، که دا رشته
\begin{quote}
\begin{verbatim}
hello x world
\end{verbatim}
\end{quote}
د ننوت په توګه ورکړل شي، تابع دا
\begin{quote}
\begin{verbatim}
hello 'hello x world' world
\end{verbatim}
\end{quote}
د وت په توګه راګرځوي۔ په دې صورت کښې موخه شوے پروګرام ليکل اسان
دي؛ ازمويلے شئ چې
\begin{quote}
\begin{verbatim}
print(diag('print(diag(x))'))
\end{verbatim}
\end{quote}
کار سر ته رسوي۔ د~C++ او جاوا په شان ډېرو عامو پروګرامي ژبو کښې هم
همدغه نظر (په يو څه پېچلې پليينې سره) کار کوي۔
\olpseganchor{OLP-0249-B017}{end}

\olpseganchor{OLP-0249-B018}{start}
موږ د ثابت ټکي د قضيې له ثبوت څخه يوازې څو ګامه لرې يو۔ فرض کړئ د
چاپ تابعې يوه بڼه~\LR{$\fn{print}(x,y)$} يوه رشته~\LR{$x$} او بل عددي دليل~\LR{$y$}
اخلي، او رشته~\LR{$x$}، \LR{$y$} ځله چاپ کوي۔ بيا دا ``پروګرام''
\begin{quote}
\begin{verbatim}
getinput(y);
print(diag('getinput(y); print(diag(x), y)'), y)
\end{verbatim}
\end{quote}
پر ننوت~\LR{$y$} خپل ځان~\LR{$y$} ځله چاپ کوي۔ که د
\LR{$\fn{getinput}$}---\LR{$\fn{print}$}---\LR{$\fn{diag}$} چوکاټ په هرې تابع
\LR{$g(x,y)$} بدل کړو، دا تر لاسه کېږي:
\begin{quote}
\begin{verbatim}
g(diag('g(diag(x), y)'), y)
\end{verbatim}
\end{quote}
دا داسې پروګرام دے چې پر ننوت~\LR{$y$}، تابع~\LR{$g$} پر خپل پروګرام او~\LR{$y$}
چلوي۔ که ``اقتباس کول'' د ``يو شاخص کارولو'' په توګه وګڼو، پورته
ثبوت تر لاسه کوو۔
\olpseganchor{OLP-0249-B018}{end}

\olpseganchor{OLP-0249-B019}{start}
اوس دپاره که غواړئ ثبوت صوري چل، يا توره جادوګري، وګڼئ پروا نۀ کوي۔
خو بايد د پورته ورکړل شوي استدلال جزئيات بيا جوړ کړے شئ۔ کله چې د
نابشپړتيا قضيې (او اړونده ``د ثابت ټکي قضيه'') ثابتوو، د دې د کار
کولو د علت د پوهېدو په نورو لارو به بحث وکړو۔
\end{explain}
\olpseganchor{OLP-0249-B019}{end}

\olpseganchor{OLP-0249-B020}{start}

\begin{digress}
همدغه نظر د ``ثابت ټکي'' يو ترکيبګر هم راکولے شي۔ فرض کړئ يو لامبډا
ترم~\LR{$g$} لرئ، او داسې بل ترم~\LR{$k$} غواړئ چې~\LR{$k$} له~\LR{$gk$} سره
\LR{$\beta$}-معادل وي۔ زموږ د نښيزو دودونو په کارولو ترمونه داسې تعريف
کړئ:
\[
\fn{diag}(x) = xx
\]
او
\[
l(x) = g(\fn{diag}(x))
\]
يعنې~\LR{$l$} دا ترم دے: \LR{$\lambd[x][g(xx)]$}۔ \LR{$k$} ترم~\LR{$ll$} وي۔ بيا لرو:
\begin{align*}
k & = (\lambd[x][g(xx)])(\lambd[x][g(xx)]) \\
& \red  g((\lambd[x][g(xx)])(\lambd[x][g(xx)])) \\
& = gk.
\end{align*}
که
\[
Y = \lambd[g][((\lambd[x][g(xx)])(\lambd[x][g(xx)]))]
\]
واخلو، نو~\LR{$Yg$} او~\LR{$g(Yg)$} يو ګډ ترم ته راکمېږي؛ ځکه
\LR{$Yg \equiv_\beta g(Yg)$}۔ دې ته ``د کري ترکيبګر'' ويل کېږي۔ که پر
ځاے ئې
\[
Y = (\lambd[xg][g(xxg)])(\lambd[xg][g(xxg)])
\]
واخلو، نو په حقيقت کښې~\LR{$Yg$}، \LR{$g(Yg)$} ته راکمېږي، چې لا پياوړے بيان
دے۔ د~\LR{$Y$} دا وروستۍ بڼه ``د ټيورينګ ترکيبګر'' بلل کېږي۔
\end{digress}

\olpseganchor{OLP-0249-B020}{end}

\olpseganchor{OLP-0250-B004}{start}
\olfileid{cmp}{thy}{apf}
\olsection{د ثابت ټکي د قضيې کارول}
\olpseganchor{OLP-0250-B004}{end}

\olpseganchor{OLP-0250-B005}{start}
د ثابت ټکي قضيه په اصل کښې اجازه راکوي چې جزوي محاسبه کېدونکې تابعې
د هغو د شاخصونو په وسيله تعريف کړو۔ د بېلګې په توګه، داسې شاخص~\LR{$e$}
موندلے شو چې د هر~\LR{$y$} دپاره
\[
\cfind{e}(y) = e + y.
\]
د بلې بېلګې په توګه، د ثابت ټکي د قضيې ثبوت کارولے شو څو په جاوا يا
C++ کښې داسې پروګرام جوړ کړو چې خپل متن چاپ کړي۔
\olpseganchor{OLP-0250-B005}{end}

\olpseganchor{OLP-0250-B006}{start}
راياد کړئ چې که د هر~\LR{$e$} دپاره~\LR{$W_e$} د~\LR{$\cfind{e}$} د تعريف ساحه
وټاکو، نو لړۍ~\LR{$W_0$}، \LR{$W_1$}، \LR{$W_2$}،~\dots په محاسبوي ډول د شمېر وړ
سټونه شمېري۔ ځينې دا سټونه محاسبه کېدونکي دي۔ پوښتلے شو چې ايا داسې
الګوريتم شته چې قيمت~\LR{$x$} د ننوت په توګه واخلي، او که~\LR{$W_x$} محاسبه
کېدونکے وي، د هغه د ځانګړونکې تابع يو شاخص را وګرځوي؟ ځواب ``نه''
دے، داسې الګوريتم نشته:
\olpseganchor{OLP-0250-B006}{end}

\olpseganchor{OLP-0250-B007}{start}
\begin{thm}
داسې جزوي محاسبه کېدونکې تابع~\LR{$f$} نشته چې دا خاصيت ولري: هر کله چې
~\LR{$W_e$} محاسبه کېدونکے وي، نو~\LR{$f(e)$} تعريف شوے وي او~\LR{$\cfind{f(e)}$}
ئې ځانګړونکې تابع وي۔
\end{thm}
\olpseganchor{OLP-0250-B007}{end}

\olpseganchor{OLP-0250-B008}{start}
\begin{proof}
\LR{$f$} هره جزوي محاسبه کېدونکې تابع وي؛ داسې~\LR{$e$} به جوړ کړو چې~\LR{$W_e$}
محاسبه کېدونکے وي، خو~\LR{$\cfind{f(e)}$} ئې ځانګړونکې تابع نۀ وي۔ د ثابت
ټکي د قضيې په کارولو داسې شاخص~\LR{$e$} موندلے شو چې
\[
\cfind{e}(y) \simeq
\begin{cases}
0 & \text{\RL{که \LR{$y=0$} او \LR{$\cfind{f(e)}(0) \fdefined = 0$} وي}} \\
\text{\RL{تعريف شوې نۀ ده}} & \text{\RL{که داسې نۀ وي۔}}
\end{cases}
\]
يعنې~\LR{$e$} د ثابت ټکي د قضيې په دې تابع له پلي کولو تر لاسه کېږي:
\[
g(x,y) \simeq
\begin{cases}
0 & \text{\RL{که \LR{$y=0$} او \LR{$\cfind{f(x)}(0) \fdefined = 0$} وي}} \\
\text{\RL{تعريف شوې نۀ ده}} & \text{\RL{که داسې نۀ وي۔}}
\end{cases}
\]
په غير رسمي ډول داسې ليدلے شو چې~\LR{$g$} جزوي محاسبه کېدونکې ده: پر
ننوتونو~\LR{$x$} او~\LR{$y$} الګوريتم لومړے ګوري چې~\LR{$y$} له~\LR{$0$} سره برابر دے
که نه۔ که برابر وي، الګوريتم~\LR{$f(x)$} محاسبه کوي، او بيا نړيوال ماشين
کاروي څو~\LR{$\cfind{f(x)}(0)$} محاسبه کړي۔ که دا وروستۍ محاسبه ودرېږي او
~\LR{$0$} را وګرځوي، الګوريتم~\LR{$0$} راګرځوي؛ که نه، الګوريتم نۀ درېږي۔
که د شاخص محاسبه پاے ته نۀ رسېږي، دغه الګوريتم هم نۀ درېږي۔
پر ناصفر ننوت هم نۀ درېږي۔
\olpseganchor{OLP-0250-B008}{end}

\olpseganchor{OLP-0250-B009}{start}
خو اوس پام وکړئ چې که~\LR{$\cfind{f(e)}(0)$} تعريف شوې او له~\LR{$0$} سره
برابره وي، نو~\LR{$\cfind{e}(y)$} کټ مټ هغه وخت تعريف شوې ده چې~\LR{$y$} له
~\LR{$0$} سره برابر وي، نو \LR{$W_e = \{ 0 \}$}۔ که~\LR{$\cfind{f(e)}(0)$} تعريف
شوې نۀ وي، يا تعريف شوې خو له~\LR{$0$} سره برابره نۀ وي، نو
\LR{$W_e = \emptyset$}۔ په دواړو حالاتو کښې~\LR{$\cfind{f(e)}$} د~\LR{$W_e$}
ځانګړونکې تابع نۀ ده: پر ننوت~\LR{$0$} يا ناسم ځواب ورکوي، يا ځواب نۀ ورکوي۔
د تعريف ساحه په هر صورت محاسبه کېدونکې ده۔ که د نوماندې تابع
ارزښت پر ټاکلي شاخص ناتعريف وي، نو د قضيې شرط همدلته ماتېږي؛
که ارزښت تعريف وي، پورته استدلال ښيي چې راوتلے شاخص د اړينې
ځانګړونکې تابع شاخص نۀ دے۔
\end{proof}
\olpseganchor{OLP-0250-B009}{end}

\olpseganchor{OLP-0251-B004}{start}
\olfileid{cmp}{thy}{slf}
\olsection{د ځان-مراجعې په کارولو د تابعو تعريف}
\olpseganchor{OLP-0251-B004}{end}

\olpseganchor{OLP-0251-B005}{start}
په عمومي ډول دا ګټوره ده چې تابعې د هغو په خپلو اصطلاحاتو کښې تعريف
کړے شو۔ د بېلګې په توګه، ورکړو محاسبه کېدونکو تابعو~\LR{$k$}، \LR{$l$} او~\LR{$m$}
ته د ثابت ټکي لمه وايي چې داسې جزوي محاسبه کېدونکې تابع~\LR{$f$} شته چې
د هر~\LR{$y$} دپاره دا معادله پوره کوي:
\[
f(y) \simeq
\begin{cases}
k(y) & \text{\RL{که \LR{$l(y) = 0$} وي}} \\
f(m(y)) & \text{\RL{که داسې نۀ وي۔}}
\end{cases}
\]
په لا ځانګړي ډول، ~\LR{$f$} بيا داسې تر لاسه کېږي:
\[
g(x,y) \simeq
\begin{cases}
k(y) & \text{\RL{که \LR{$l(y) = 0$} وي}} \\
\cfind{x}(m(y)) & \text{\RL{که داسې نۀ وي}}
\end{cases}
\]
او بيا د ثابت ټکي لمه کارول کېږي څو داسې شاخص~\LR{$e$} وموندل شي چې
\LR{$\cfind{e}(y) \simeq g(e,y)$} وي۔
\olpseganchor{OLP-0251-B005}{end}

\olpseganchor{OLP-0251-B006}{start}
د يوې محسوسې بېلګې دپاره د ``تر ټولو لوی ګډ مقسوم عليه'' تابع
\LR{$\fn{gcd}(u,v)$} داسې تعريفېدے شي:
\[
\fn{gcd}(u,v) \simeq
\begin{cases}
v & \text{\RL{که \LR{$u = 0$} وي}} \\
\fn{gcd}(\fn{mod}(v, u), u) & \text{\RL{که داسې نۀ وي}}
\end{cases}
\]
چېرته چې~\LR{$\fn{mod}(v, u)$} پر~\LR{$u$} د~\LR{$v$} د وېش پاتې شونې څرګندوي۔ د
ثابت ټکي لمې ته مراجعه ښيي چې~\LR{$\fn{gcd}$} جزوي محاسبه کېدونکې ده۔
(په حقيقت کښې دا په پورته بڼه راوستل کېدے شي، چې~\LR{$y$} د جوړې
\LR{$\tuple{u, v}$} کوډ وګڼل شي۔) بيا پر~\LR{$u$} استقرا ښيي چې په حقيقت کښې
~\LR{$\fn{gcd}$} هرځاے تعريف شوې ده۔
\olpseganchor{OLP-0251-B006}{end}

\olpseganchor{OLP-0251-B007}{start}
البته داسې ځان-مراجع تعريفونه جوړولے شو چې له پورته بېلګو ډېر پېچلي
وي۔ زياتره پروګرامي ژبې په يوه يا بله طريقه د تابعو تعريف په خپلو
اصطلاحاتو کښې مني۔ پام وکړئ چې دا د هغې وړتيا په پرتله لږ حيرانوونکے
دے چې يوه تابع د هغه الګوريتم د \emph{شاخص} په وسيله تعريف کړو چې
هماغه تابع محاسبه کوي؛ او د ثابت ټکي قضيه په بشپړ عموميت کښې همدا
اجازه راکوي۔
\olpseganchor{OLP-0251-B007}{end}

\olpseganchor{OLP-0252-B004}{start}
\olpart{tur}{ټيورينګ ماشينونه}
\olpseganchor{OLP-0252-B004}{end}

\olpseganchor{OLP-0253-B004}{start}
\olchapter{tur}{mac}{د ټيورينګ ماشين محاسبې}
\olpseganchor{OLP-0253-B004}{end}

\olpseganchor{OLP-0254-B004}{start}
\olfileid{tur}{mac}{int}
\olsection{پېژندنه}
\olpseganchor{OLP-0254-B004}{end}

\olpseganchor{OLP-0254-B005}{start}
دا څه معنا لري چې، مثلاً، له~\LR{$\Nat$} څخه~\LR{$\Nat$} ته يوه تابع
\emph{محاسبه کېدونکې} وي؟ د لومړنيو او تر ټولو نامتو ځوابونو څخه يو
دا دے چې تابع هغه وخت محاسبه کېدونکې ده چې د ټيورينګ ماشين په وسيله
محاسبه کېدے شي۔ دا مفهوم الن ټيورينګ په~1936 کښې وړاندې کړ۔ د ټيورينګ
ماشينونه د \emph{محاسبې د يو مدل} بېلګه ده---دوي د ``محاسبوي کړنلارې''
نظر د تعريفولو يوه رياضياتي کره لار ده۔ د دې کره معنا تر بحث لاندې
ده، خو په پراخه کچه منل شوې چې ټيورينګ ماشينونه د محاسبوي کړنلارو د
ټاکلو يوه لار ده۔ که څه هم د ``ټيورينګ ماشين'' اصطلاح د خوځنده پرزو
لرونکي فزيکي ماشين انځور ذهن ته راولي، په کره معنا ټيورينګ ماشين يو
سوچه رياضياتي جوړښت دے، او په همدې توګه د محاسبوي کړنلارې نظر مثالي
کوي۔ مثلاً، موږ د ټيورينګ ماشين د وخت يا حافظې پر اړتياوو هېڅ بنديز
نۀ ږدو: ټيورينګ ماشينونه يو څيز ان هغه وخت محاسبه کولے شي چې محاسبه
له هغو ټولو ذرو څخه زياته زېرمه يا زيات ګامونه وغواړي چې په کاينات
کښې شته۔
\olpseganchor{OLP-0254-B005}{end}

\olpseganchor{OLP-0254-B006}{start}
\begin{explain}
ښايي غوره وي چې ټيورينګ ماشين د يو ځانګړي ډول خيالي ميکانيزم د
پروګرام په توګه وګڼو۔ دا ميکانيزم له يوې \emph{پټې} او يو
\emph{لوست-ليک سر} څخه جوړ دے۔ زموږ د ټيورينګ ماشينونو په بڼه کښې
پټه په يوه لوري (ښي لوري ته) نامتناهي ده، او په \emph{مربعونو} وېشل
شوې، چې هر يو ئې د يوې متناهي \emph{الفبا} يوه نښه لرلے شي۔ دا ډول
الفباوې هر شمېر بېلابېلې نښې لرلے شي، خو موږ به تر ډېره په درېيو بسنه
کوو: \LR{$\TMendtape$}، \LR{$\TMblank$} او~\LR{$\TMstroke$}۔ کله چې ميکانيزم پيل
شي، پټه تشه وي (يعنې هر مربع د~\LR{$\TMblank$} نښه لري)، پرته له تر ټولو
کيڼ مربع څخه چې~\LR{$\TMendtape$} لري، او يو متناهي شمېر مربعونو څخه چې
\emph{ننوت} لري۔ په هر وخت کښې ميکانيزم له متناهي شمېر
\emph{حالتونو} څخه په يوه کښې وي۔ په پيل کښې سر د ننوت تر ټولو کيڼ
مربع، يعنې د پټې د پای نښې ښي لوري ته مربع، لولي او په يو ټاکلي
\emph{لومړني حالت} کښې وي۔ که ننوت تش وي، سر د پای نښې ښي لوري ته هماغه تش مربع لولي۔ د ميکانيزم د چلېدو په هر ګام کښې د اوس
لوستل کېدونکي مربع منځپانګه، د ميکانيزم حالت او د ټيورينګ ماشين
پروګرام په ګډه ټاکي چې ورپسې څه کېږي۔ د ټيورينګ ماشين پروګرام د يوې
جزوي تابع په توګه ورکول کېږي چې حالت~\LR{$q$} او نښه~\LR{$\sigma$} د ننوت په
توګه اخلي، او درې‌يز~\LR{$\tuple{q', \sigma', D}$} راګرځوي۔ هر کله چې
ميکانيزم په حالت~\LR{$q$} کښې وي او نښه~\LR{$\sigma$} لولي، د اوسني مربع نښه
په~\LR{$\sigma'$} بدلوي، سر د دې له مخې کيڼ، ښي لوري ته ځي يا پر ځاے پاتې
کېږي چې~\LR{$D$} په ترتيب سره~\LR{$\TMleft$}، \LR{$\TMright$} يا~\LR{$\TMstay$} وي، او
ميکانيزم حالت~\LR{$q'$} ته ځي۔
\olpseganchor{OLP-0254-B006}{end}

\olpseganchor{OLP-0254-B007}{start}
د بېلګې په توګه، په~\olref{fig:tm} کښې حالت وګورئ۔
\begin{figure}
  \olasset{../upstream/assets/diagrams/turing-machine.tikz}
  \caption{يو ټيورينګ ماشين چې خپل پروګرام عملي کوي۔}
  \ollabel{fig:tm}
\end{figure}
د ټيورينګ ماشين د پټې ليدونکې برخه په تر ټولو کيڼ مربع کښې د پټې د
پای نښه~\LR{$\TMendtape$} لري، ورپسې درې~\LR{$1$} ګانې، يو~\LR{$0$}، او بيا څلور
نورې~\LR{$1$} ګانې دي۔ سر له کيڼ لوري درېيم مربع لولي، چې~\LR{$1$} لري، او په
حالت~\LR{$q_1$} کښې دے---وايو ``ماشين په حالت~\LR{$q_1$} کښې~\LR{$1$} لولي۔'' که
د ټيورينګ ماشين پروګرام د ننوت~\LR{$\tuple{q_1, 1}$} دپاره درې‌يز
\LR{$\tuple{q_2, 0, \TMstay}$} را وګرځوي، ماشين به اوس په درېيم مربع کښې
~\LR{$1$} په~\LR{$0$} بدل کړي، لوست-ليک سر به پر خپل ځاے پرېږدي، او حالت~\LR{$q_2$}
ته به واوړي۔ که بيا پروګرام د ننوت~\LR{$\tuple{q_2, 0}$} دپاره
\LR{$\tuple{q_3, 0, \TMright}$} را وګرځوي، ماشين به~\LR{$0$} په بل~\LR{$0$} وليکي
(يعنې په عمل کښې به د لوست-ليک سر لاندې د پټې منځپانګه هماغسې
پرېږدي)، يو مربع ښي لوري ته به لاړ شي، او حالت~\LR{$q_3$} ته به ننوځي۔ او
همداسې نور۔
\olpseganchor{OLP-0254-B007}{end}

\olpseganchor{OLP-0254-B008}{start}
وايو ماشين هغه وخت \emph{درېږي} چې له داسې حالت~\LR{$q_n$} او نښې
~\LR{$\sigma$} سره مخ شي چې د~\LR{$\tuple{q_n, \sigma}$} دپاره هېڅ لارښوونه
نۀ وي، يعنې د ننوت~\LR{$\tuple{q_n,\sigma}$} دپاره د حالت بدلون تابع
تعريف شوې نۀ وي۔ په بله وينا، ماشين د عملي کولو دپاره هېڅ لارښوونه
نۀ لري، او په هماغه ځاے کار بندوي۔ کله ناکله درېدل د ځانګړي درېدني
حالت~\LR{$h$} په وسيله څرګندېږي۔ دا به وروسته په زياتو جزئياتو وښودل شي۔
\end{explain}
\olpseganchor{OLP-0254-B008}{end}

\olpseganchor{OLP-0254-B009}{start}
\begin{digress}
د ټيورينګ د مقالې، ``د محاسبه کېدونکو عددونو په اړه،'' ښکلا دا ده چې
هغه نۀ يوازې يو صوري تعريف وړاندې کوي، بلکې دا استدلال هم کوي چې
تعريف د محاسبه کېدنې شهودي مفهوم رانغاړي۔ له تعريفه بايد څرګنده وي
چې هره تابع چې د ټيورينګ ماشين په وسيله محاسبه کېږي، په شهودي معنا
محاسبه کېدونکې ده۔ ټيورينګ درې ډوله استدلالونه وړاندې کوي چې معکوسه
خبره هم رښتيا ده، يعنې هره تابع چې موږ ئې په طبيعي ډول محاسبه کېدونکې
ګڼو، د داسې ماشين په وسيله محاسبه کېدونکې ده۔ د ټيورينګ په خپلو
ټکو، دا دي:
\begin{enumerate}
\item شهود ته نېغه مراجعه۔
\item د دوو تعريفونو د معادلتوب ثبوت (په هغه صورت کښې چې نوی تعريف
  زيات شهودي جاذبيت ولري)۔
\item د محاسبه کېدونکو عددونو د لويو ټولګيو بېلګې ورکول۔
\end{enumerate}
زموږ موخه دا ده چې د محاسبه کېدنې مفهوم ``په اصل کښې'' تعريف کړو،
يعنې د وخت او ځاے عملي محدوديتونه په پام کښې وانۀ خلو۔ البته، کله چې
د محاسبه کېدنې تر ټولو پراخ تعريف ولرو، بيا د محدودو سرچينو لرونکې
محاسبه څېړلے شو؛ دا د ``محاسبوي پېچلتيا'' په نوم مضمون زړه جوړوي۔
\end{digress}
\olpseganchor{OLP-0254-B009}{end}

\olpseganchor{OLP-0254-B010}{start}
\begin{history}
الن ټيورينګ په~1936 کښې ټيورينګ ماشينونه اختراع کړل۔ که څه هم په هغه
وخت کښې د هغه پاملرنه د لومړۍ درجې منطق پرېکړه کېدنې ته وه، دا مقاله
د کمپيوټر د جوړښت د بنسټونو يوه پرېکنده مقاله بلل شوې ده۔ په مقاله
کښې ټيورينګ پر محاسبه کېدونکو حقيقي عددونو، يعنې هغو حقيقي عددونو چې
اعشاري غځونې ئې محاسبه کېدونکې دي، تمرکز کوي؛ خو يادونه کوي چې د هغه
مفهومونه پر طبيعي عددونو محاسبه کېدونکو تابعو او نورو ته اړول ستونزمن
نۀ دي۔ پام وکړئ چې دا په~1941 کښې د لومړي کارکوونکي عمومي-موخې
کمپيوټر له جوړېدو پوره پنځه کاله مخکښې وو (الماني کونراد زوزه دا د
خپلو والدينو د کور په ناسته کوټه کښې جوړ کړ)، په بلېچلي پارک کښې د
ټيورينګ او د هغه د ملګرو د کوډ ماتوونکي کولوسس له جوړېدو اووه کاله
مخکښې (1943)، د امريکايي اېنياک څخه نهه کاله مخکښې (1945)، په
مانچسټر کښې د لومړي برتانوي عمومي-موخې کمپيوټر---د مانچسټر د کوچنۍ
کچې ازمېښتي ماشين---له جوړېدو دولس کاله مخکښې (1948)، او د
امريکايانو له خوا د بيناک له لومړۍ ازموينې ديارلس کاله مخکښې (1949)۔
د مانچسټر کوچنۍ کچې ازمېښتی ماشين د لومړي زېرمه-پروګرام کمپيوټر وياړ
لري---پخواني ماشينونه د هرې نوې دندې دپاره په لاس بيا سيمبندي کېدل۔
\end{history}
\olpseganchor{OLP-0254-B010}{end}

\olpseganchor{OLP-0255-B004}{start}
\olfileid{tur}{mac}{rep}
\olsection{د ټيورينګ ماشينونو څرګندونه}
\olpseganchor{OLP-0255-B004}{end}

\olpseganchor{OLP-0255-B005}{start}
\begin{explain}
ټيورينګ ماشينونه په ليديز ډول د \emph{حالت ډياګرامونو} په وسيله
ښودل کېدے شي۔ دا ډياګرامونه د حالت له هغو خانو جوړ دي چې غشي ئې سره
نښلوي۔ لکه چې تمه کېږي، د حالت هره خانه د ماشين د يو حالت استازيتوب
کوي۔ هر غشی د هغې لارښوونې استازيتوب کوي چې له هماغه حالت څخه عملي
کېدے شي، او د لارښوونې جزئيات د اړوند غشي دپاسه يا لاندې ليکل شوي
وي۔ لاندې ماشين وګورئ، چې يوازې دوه دننني حالتونه، \LR{$q_0$} او~\LR{$q_1$}،
او يوه لارښوونه لري:
\[
\begin{tikzpicture}[->,>=stealth',shorten >=1pt,auto,node distance=2.8cm,
                    semithick]
  \tikzstyle{every state}=[fill=none,draw=black,text=black]
\olpseganchor{OLP-0255-B005}{end}

\olpseganchor{OLP-0255-B006}{start}
  \node[initial,state]   (A)              {\LR{$q_0$}};
  \node[state]   (B) [right of=A] {\LR{$q_1$}};
\olpseganchor{OLP-0255-B006}{end}

\olpseganchor{OLP-0255-B007}{start}
  \path (A) edge  node {\TMtrans{\TMblank}{\TMstroke}{\TMright}} (B);
\end{tikzpicture}
\]
په ياد ولرئ چې ټيورينګ ماشين يو لوست-ليک سر او يوه پټه لري چې ننوت
پرې ليکل شوے وي۔ لارښوونه داسې لوستل کېږي: \emph{که په حالت~\LR{$q_0$}
کښې~\LR{$\TMblank$} لولئ، يو~\LR{$\TMstroke$} وليکئ، ښي لوري ته لاړ شئ، او
حالت~\LR{$q_1$} ته واوړئ}۔ دا له دې سره معادل ده چې د حالت بدلون تابع
\LR{$\tuple{q_0, \TMblank}$} په~\LR{$\tuple{q_1, \TMstroke,
\TMright}$} واړوي۔
\end{explain}
\olpseganchor{OLP-0255-B007}{end}

\olpseganchor{OLP-0255-B008}{start}
\begin{ex}
\emph{جفت ماشين}: لاندې ټيورينګ ماشين هله او يوازې هله درېږي چې پر
پټه د~\LR{$\TMstroke$} ګانو شمېر جفت وي (په دې فرض چې د پټې ټول
\LR{$\TMstroke$} ګانې له لومړي~\LR{$\TMblank$} څخه مخکښې راځي)۔
\[
\begin{tikzpicture}[->,>=stealth',shorten >=1pt,auto,node distance=2.8cm,
                    semithick]
  \tikzstyle{every state}=[fill=none,draw=black,text=black]
\olpseganchor{OLP-0255-B008}{end}

\olpseganchor{OLP-0255-B009}{start}
  \node[initial,state]         (A)              {\LR{$q_0$}};
  \node[state]         (B) [right of=A] {\LR{$q_1$}};
\olpseganchor{OLP-0255-B009}{end}

\olpseganchor{OLP-0255-B010}{start}
  \path (A) edge [bend left] node {\TMtrans{\TMstroke}{\TMstroke}{\TMright}} (B)
        (B) edge [loop above] node {\TMtrans{\TMblank}{\TMblank}{\TMright}} (B)
            edge [bend left] node {\TMtrans{\TMstroke}{\TMstroke}{\TMright}} (A);
\end{tikzpicture}
\]
\olpseganchor{OLP-0255-B010}{end}

\olpseganchor{OLP-0255-B011}{start}
د حالت ډياګرام له لاندې حالت بدلون تابع سره سمون لري:
\begin{align*}
\delta(q_0, \TMstroke) & = \tuple{q_1, \TMstroke, \TMright},\\
\delta(q_1, \TMstroke) & = \tuple{q_0, \TMstroke, \TMright},\\
\delta(q_1, \TMblank)  & = \tuple{q_1, \TMblank, \TMright}
\end{align*}
\end{ex}
\olpseganchor{OLP-0255-B011}{end}

\olpseganchor{OLP-0255-B012}{start}
\begin{explain}
پورته ماشين يوازې هغه وخت درېږي چې ننوت د کرښو جفت شمېر وي۔ که نه،
ماشين (په نظري ډول) تر نامعلومې مودې کار ته دوام ورکوي۔ د هر ماشين
او ننوت دپاره د ماشين \emph{تشکيلونه} پرله‌پسې څارلے شو، څو وت ئې
وټاکو۔ د تشکيلونو صوري تعريف به وروسته ورکړو۔ اوس د تشکيلونو په اړه
په شهودي ډول د ډياګرامونو د يوې لړۍ په توګه فکر کولے شو، چې د کار په
بهير کښې په هره شېبه د ماشين حالت ښيي۔ تشکيلونه د پټې منځپانګه، د
ماشين حالت او د لوست-ليک سر ځاے ښيي۔
\olpseganchor{OLP-0255-B012}{end}

\olpseganchor{OLP-0255-B013}{start}
راځئ د جفت ماشين تشکيلونه هغه وخت پرله‌پسې وڅارو چې د څلورو
\LR{$\TMstroke$} ګانو پر ننوت پيل شي۔ په دې صورت کښې تمه لرو چې ماشين
ودرېږي۔ بيا به ماشين د درېيو~\LR{$\TMstroke$} ګانو پر ننوت وچلوو، چې
هلته به ماشين تل روان وي۔
\olpseganchor{OLP-0255-B013}{end}

\olpseganchor{OLP-0255-B014}{start}
ماشين په حالت~\LR{$q_0$} کښې، د تر ټولو کيڼ~\LR{$\TMstroke$} په لوستلو پيل
کوي۔ د ماشين لومړنی حالت داسې ښودلے شو:
\[
\TMendtape \TMstroke_0 \TMstroke \TMstroke \TMstroke \TMblank \ldots
\]
پورته تشکيل روښانه دے۔ لکه چې ليدل کېږي، ماشين په صفرم حالت کښې د
تر ټولو کيڼ~\LR{$\TMstroke$} په لوستلو پيل کوي۔ دا پر لومړي~\LR{$\TMstroke$}
د حالت د نوم په فرعي نښه ښودل کېږي۔
په دې ځاے کښې د عملي کېدو وړ لارښوونه
\LR{$\delta(q_0, \TMstroke) = \tuple{q_1, \TMstroke, \TMright}$} ده، نو
ماشين پر پټه ښي لوري ته ځي او حالت~\LR{$q_1$} ته اوړي۔
\[
\TMendtape \TMstroke \TMstroke_1 \TMstroke \TMstroke \TMblank \ldots
\]
څرنګه چې ماشين اوس په حالت~\LR{$q_1$} کښې يو~\LR{$\TMstroke$} لولي، بايد
لارښوونه~\LR{$\delta(q_1, \TMstroke) = \tuple{q_0,
  \TMstroke, \TMright}$} ``تعقيب'' کړو۔ په نتيجه کښې دا تشکيل لاس ته
راځي:
\[
\TMendtape \TMstroke \TMstroke \TMstroke_0 \TMstroke \TMblank \ldots
\]
د ماشين د دوام سره قواعد يو ځل بيا په هماغه ترتيب عملي کېږي، او
لاندې دوه تشکيلونه لاس ته راځي:
\[
\TMendtape \TMstroke \TMstroke \TMstroke \TMstroke_1 \TMblank \ldots
\]
\[
\TMendtape \TMstroke \TMstroke \TMstroke \TMstroke \TMblank_0 \ldots
\]
ماشين اوس په حالت~\LR{$q_0$} کښې يو~\LR{$\TMblank$} لولي۔ د حالت بدلون د
ډياګرام له مخې په اسانه وينو چې د عملي کولو هېڅ لارښوونه نشته، نو
ماشين درېدلے دے۔ معنا دا چې ننوت منل شوے دے۔
\olpseganchor{OLP-0255-B014}{end}

\olpseganchor{OLP-0255-B015}{start}
اوس فرض کړئ ماشين د درېيو~\LR{$\TMstroke$} ګانو پر ننوت پيلوو۔ لومړني
څو تشکيلونه ورته دي، ځکه هماغه لارښوونې عملي کېږي او يوازې د پټې په
ننوت کښې لږ توپير شته:
\[
\TMendtape \TMstroke_0 \TMstroke \TMstroke \TMblank \ldots
\]
\[
\TMendtape \TMstroke \TMstroke_1 \TMstroke \TMblank \ldots
\]
\[
\TMendtape \TMstroke \TMstroke \TMstroke_0 \TMblank \ldots
\]
\[
\TMendtape \TMstroke \TMstroke \TMstroke \TMblank_1 \ldots
\]
ماشين اوس له ټولو~\LR{$\TMstroke$} ګانو څخه تېر شوے، او په حالت~\LR{$q_1$}
کښې يو~\LR{$\TMblank$} لولي۔ لکه چې په ډياګرام کښې ښودل شوې، د
\LR{$\delta(q_1, \TMblank) =\tuple{q_1,
\TMblank, \TMright}$} په بڼه يوه لارښوونه شته۔ څرنګه چې پټه ښي لوري
ته تر نامعلومې مودې پورې له~\LR{$\TMblank$} څخه ډکه ده، ماشين به دا
لارښوونه \emph{تل} عملي کوي، په حالت~\LR{$q_1$} کښې به پاتې کېږي او لا
پسې به ښي لوري ته ځي۔ ماشين به هېڅکله ونۀ درېږي، او ننوت نۀ مني۔
\end{explain}
\olpseganchor{OLP-0255-B015}{end}

\olpseganchor{OLP-0255-B016}{start}
\begin{explain}
دا خبره مهمه ده چې ټول ماشينونه نۀ درېږي۔ که درېدل په دې معنا وي چې
د ماشين د عملي کولو لارښوونې خلاصې شي، نو داسې ماشين جوړولے شو چې
هېڅکله نۀ درېږي: بس دا باوري کړو چې په هر حالت کښې د هرې نښې دپاره
يو وتونکے غشی شته۔ په حالت~\LR{$q_0$} کښې د~\LR{$\TMblank$} د لوستلو يوه
لارښوونه په ور زياتولو جفت ماشين داسې بدلولے شو چې تل روان وي۔
\end{explain}
\olpseganchor{OLP-0255-B016}{end}

\olpseganchor{OLP-0255-B017}{start}
\begin{ex}
\[
\begin{tikzpicture}[->,>=stealth',shorten >=1pt,auto,node distance=2.8cm,
                    semithick]
  \tikzstyle{every state}=[fill=none,draw=black,text=black]
\olpseganchor{OLP-0255-B017}{end}

\olpseganchor{OLP-0255-B018}{start}
  \node[initial,state] (A)              {\LR{$q_0$}};
  \node[state]         (B) [right of=A] {\LR{$q_1$}};
\olpseganchor{OLP-0255-B018}{end}

\olpseganchor{OLP-0255-B019}{start}
  \path
  (A) edge [bend left] node {\TMtrans{\TMstroke}{\TMstroke}{\TMright}} (B)
      edge [loop above] node {\TMtrans{\TMblank}{\TMblank}{\TMright}} (A)
  (B) edge [loop above] node {\TMtrans{\TMblank}{\TMblank}{\TMright}} (B)
      edge [bend left] node {\TMtrans{\TMstroke}{\TMstroke}{\TMright}} (A);
\end{tikzpicture}
\]
\end{ex}
\olpseganchor{OLP-0255-B019}{end}

\olpseganchor{OLP-0255-B020}{start}
\begin{explain}
د ماشين جدولونه د ټيورينګ ماشينونو د څرګندولو بله لار ده۔ د ماشين
جدولونه د پټې الفبا پر~\LR{$x$}-محور، او د ماشين د حالتونو سټ پر~\LR{$y$}-محور
ښيي۔ د جدول دننه، د هر حالت او نښې په ګډ ځاے کښې د لارښوونې پاتې
برخه---نوې نښه، نوی حالت او د خوځښت لوری---ليکل کېږي۔ د ماشين جدولونه
دا اسانه کوي چې ماشين په کوم حالت او د کومې نښې دپاره درېږي۔ په جدول
کښې هر تش ځاے د ماشين د درېدو يو ممکن ټکے دے۔ د حالت د ډياګرامونو او
د لارښوونو د سټونو پر خلاف، چې د ماشين د درېدو ځايونه پکښې تل سملاسي
څرګند نۀ وي، د ماشين په جدول کښې د تشو ځايونو په موندلو هر د درېدو
ټکے ژر پېژندل کېږي۔
\end{explain}
\olpseganchor{OLP-0255-B020}{end}

\olpseganchor{OLP-0255-B021}{start}
\begin{ex}
د جفت ماشين جدول دا دے:
\[
\centering
\begin{tabular}{lllll}
\cline{2-4}
\multicolumn{1}{l|}{}      & \multicolumn{1}{c|}{\LR{$\TMblank$}}
& \multicolumn{1}{c|}{\LR{$\TMstroke$}}     & \multicolumn{1}{c|}{\LR{$\TMendtape$}}   &  \\ \cline{1-4}
\multicolumn{1}{|l|}{\LR{$q_0$}} & \multicolumn{1}{l|}{}
& \multicolumn{1}{l|}{\TMtrans{\TMstroke}{q_1}{\TMright}} &
\multicolumn{1}{l|}{\phantom{\TMtrans{\TMstroke}{q_1}{\TMright}}} &  \\ \cline{1-4}
\multicolumn{1}{|l|}{\LR{$q_1$}} & \multicolumn{1}{l|}{\TMtrans{\TMblank}{q_1}{\TMright}}
& \multicolumn{1}{l|}{\TMtrans{\TMstroke}{q_0}{\TMright}} &
\multicolumn{1}{l|}{\phantom{\TMtrans{\TMstroke}{q_1}{\TMright}}} &  \\ \cline{1-4}
\end{tabular}
\]
لکه چې وينو، ماشين هغه وخت درېږي چې په حالت~\LR{$q_0$} کښې يو~\LR{$\TMblank$} لولي۔
\end{ex}
\olpseganchor{OLP-0255-B021}{end}

\olpseganchor{OLP-0255-B022}{start}
\begin{explain}
تر اوسه مو يوازې هغه ماشينونه کتلي چې ننوت لولي او مني ئې۔ خو
ټيورينګ ماشينونه د لوستلو او ليکلو دواړو وړتيا لري۔ د داسې ماشين يوه
بېلګه (که څه هم ډېرې، ډېرې بېلګې شته) \emph{دوه‌چنده کوونکے} دے۔
دوه‌چنده کوونکے چې پر پټه د~\LR{$n$} شمېر~\LR{$\TMstroke$} ګانو له يوې ډلې سره
پيل شي، د~\LR{$2n$} شمېر~\LR{$\TMstroke$} ګانو يوه ډله د وت په توګه ورکوي۔
\end{explain}
\olpseganchor{OLP-0255-B022}{end}

\olpseganchor{OLP-0255-B023}{start}
\begin{ex}
  \ollabel{ex:doubler} مخکښې له دې چې دوه‌چنده کوونکے ماشين جوړ کړو،
  د مسئلې د حل دپاره يوه \emph{تګلاره} جوړول مهم دي۔ څرنګه چې ماشين
  (لکه چې موږ تعريف کړے) دا نۀ شي ياد ساتلے چې څو~\LR{$\TMstroke$} ګانې
  ئې لوستې دي، بايد د پټې د ټولو~\LR{$\TMstroke$} ګانو د شمېر ساتلو يوه
  لار پيدا کړو۔ يوه داسې لار دا ده چې وت له ننوت څخه په يو~\LR{$\TMblank$}
  جلا کړو۔ بيا ماشين له ننوت څخه لومړے~\LR{$\TMstroke$} پاکولے، د ننوت پر
  پاتې برخې تېرېدلے، يو~\LR{$\TMblank$} پرېښودلے، او دوه نوي~\LR{$\TMstroke$}
  ګانې ليکلے شي۔ بيا به ماشين بېرته لاړ شي، په ننوت کښې به دويم
  \LR{$\TMstroke$} پيدا کړي، او هغه به هم دوه‌چنده کړي۔ د ننوت د هر يو
  \LR{$\TMstroke$} په بدل کښې به د وت دوه~\LR{$\TMstroke$} ګانې وليکي۔ د
  ماشين د تګ په ترڅ کښې د ننوت په پاکولو باوري کولے شو چې هېڅ
  \LR{$\TMstroke$} پاتې نۀ شي او نۀ دوه ځله دوه‌چنده شي۔ کله چې ټول ننوت
  پاک شي، پر پټه به~\LR{$2n$} شمېر~\LR{$\TMstroke$} ګانې پاتې وي۔ د لاس ته
  راغلي ټيورينګ ماشين د حالت ډياګرام په~\olref{fig:doubler} کښې
  ښودل شوے دے۔
  \begin{figure}
\[
\begin{tikzpicture}[->,>=stealth',shorten >=1pt,auto,node distance=2.8cm,
                    semithick]
  \tikzstyle{every state}=[fill=none,draw=black,text=black]
\olpseganchor{OLP-0255-B023}{end}

\olpseganchor{OLP-0255-B024}{start}
  \node[initial,state] (1)              {\LR{$q_0$}};
  \node[state]         (2) [right of=1] {\LR{$q_1$}};
  \node[state]         (3) [right of=2] {\LR{$q_2$}};
  \node[state]         (4) [below of=3] {\LR{$q_3$}};
  \node[state]         (5) [left of=4]  {\LR{$q_4$}};
  \node[state]         (6) [left of=5]  {\LR{$q_5$}};
\olpseganchor{OLP-0255-B024}{end}

\olpseganchor{OLP-0255-B025}{start}
  \path (1) edge node {\TMtrans{\TMstroke}{\TMblank}{\TMright}} (2)
    (2) edge [loop above] node {\TMtrans{\TMstroke}{\TMstroke}{\TMright}} (2)
      edge node {\TMtrans{\TMblank}{\TMblank}{\TMright}} (3)
    (3) edge [loop above] node {\TMtrans{\TMstroke}{\TMstroke}{\TMright}} (3)
        edge node {\TMtrans{\TMblank}{\TMstroke}{\TMright}} (4)
    (4) edge [loop below] node {\TMtrans{\TMblank}{\TMstroke}{\TMleft}} (4)
        edge node {\TMtrans{\TMstroke}{\TMstroke}{\TMleft}} (5)
    (5) edge [loop below]  node {\TMtrans{\TMstroke}{\TMstroke}{\TMleft}} (5)
        edge              node {\TMtrans{\TMblank}{\TMblank}{\TMleft}} (6)
    (6) edge [loop below] node {\TMtrans{\TMstroke}{\TMstroke}{\TMleft}} (6)
        edge              node {\TMtrans{\TMblank}{\TMblank}{\TMright}} (1);
\end{tikzpicture}
\]
\caption{يو دوه‌چنده کوونکے ماشين}
\ollabel{fig:doubler}
\end{figure}
\end{ex}
\olpseganchor{OLP-0255-B025}{end}

\olpseganchor{OLP-0255-B026}{start}
\begin{prob}
يو خپل‌سری ننوت وټاکئ او په~\olref[tur][mac][rep]{ex:doubler} کښې
د دوه‌چنده کوونکي ماشين تشکيلونه پرله‌پسې وڅارئ۔
\end{prob}
\olpseganchor{OLP-0255-B026}{end}

\olpseganchor{OLP-0255-B027}{start}
\begin{prob}
د~\LR{$\{\TMendtape,\TMblank, A, B\}$} الفبا لرونکے ټيورينګ ماشين جوړ
کړئ چې د~\LR{$A$} ګانو او~\LR{$B$} ګانو هره هغه رشته مني، يعنې پرې درېږي، چې
د~\LR{$A$} ګانو شمېر ئې د~\LR{$B$} ګانو له شمېر سره برابر وي \emph{او} ټول
\LR{$A$} ګانې له ټولو~\LR{$B$} ګانو مخکښې راشي؛ او هره هغه رشته ردوي، يعنې
پرې نۀ درېږي، چې د~\LR{$A$} ګانو شمېر د~\LR{$B$} ګانو له شمېر سره برابر نۀ وي
يا ټولې~\LR{$A$} ګانې له ټولو~\LR{$B$} ګانو مخکښې رانۀ شي۔ (مثلاً، ماشين بايد
\LR{$AABB$} او~\LR{$AAABBB$} ومني، خو~\LR{$AAB$} او~\LR{$AABBAABB$} دواړه رد کړي۔)
\end{prob}
\olpseganchor{OLP-0255-B027}{end}

\olpseganchor{OLP-0255-B028}{start}
\begin{prob}
د~\LR{$\{\TMendtape,\TMblank, A, B\}$} الفبا لرونکے ټيورينګ ماشين جوړ
کړئ چې د~\LR{$A$} ګانو او~\LR{$B$} ګانو هره رشته~\LR{$\alpha$} د ننوت په توګه
واخلي او نقل ئې کړي، څو د~\LR{$\alpha\alpha$} په بڼه وت پيدا کړي۔ (مثلاً،
د~\LR{$ABBA$} ننوت بايد د~\LR{$ABBAABBA$} وت ورکړي۔)
\end{prob}
\olpseganchor{OLP-0255-B028}{end}

\olpseganchor{OLP-0255-B029}{start}
\begin{prob}
\emph{الفبايي؟}: د~\LR{$\{\TMendtape,\TMblank, A, B\}$} الفبا لرونکے
ټيورينګ ماشين جوړ کړئ چې د~\LR{$A$} ګانو او~\LR{$B$} ګانو يوه متناهي لړۍ د
ننوت په توګه ورکړل شي، نو وګوري چې ټولې~\LR{$A$} ګانې د ټولو~\LR{$B$} ګانو
کيڼ لوري ته راځي که نه۔ ماشين بايد د ننوت رشته پر پټه پرېږدي، او که
رشته ``الفبايي'' وي ودرېږي، او که نۀ وي تل په کړۍ کښې وچلېږي۔
\end{prob}
\olpseganchor{OLP-0255-B029}{end}

\olpseganchor{OLP-0255-B030}{start}
\begin{prob}
\emph{الفبايي کوونکے}: د~\LR{$\{\TMendtape,\TMblank, A, B\}$} الفبا
لرونکے ټيورينګ ماشين جوړ کړئ چې د~\LR{$A$} ګانو او~\LR{$B$} ګانو يوه متناهي
لړۍ د ننوت په توګه واخلي او داسې ئې بيا ترتيب کړي چې ټولې~\LR{$A$} ګانې
د ټولو~\LR{$B$} ګانو کيڼ لوري ته راشي۔ (مثلاً، لړۍ~\LR{$BABAA$} بايد په
\LR{$AAABB$}، او لړۍ~\LR{$ABBABB$} بايد په~\LR{$AABBBB$} بدله شي۔)
\end{prob}
\olpseganchor{OLP-0255-B030}{end}

\olpseganchor{OLP-0256-B004}{start}
\olfileid{tur}{mac}{tur}
\olsection{ټيورينګ ماشينونه}
\olpseganchor{OLP-0256-B004}{end}

\olpseganchor{OLP-0256-B005}{start}
\begin{explain}
د ټيورينګ ماشين د جوړښت صوري تعريف مجرد ښکاري، خو په اصل کښې ساده
دے: دا يوازې هغه ټول معلومات په يوه رياضياتي جوړښت کښې رانغاړي چې د
ټيورينګ ماشين د کار د ټاکلو دپاره پکار دي۔ په دې کښې دا شامل دي: (1)
ماشين په کومو حالتونو کښې اوسېدلے شي، (2) پر پټه کومې نښې روا دي، (3)
ماشين بايد په کوم حالت کښې پيل وکړي، او (4) د ماشين د لارښوونو سټ څه
دے۔
\end{explain}
\olpseganchor{OLP-0256-B005}{end}

\olpseganchor{OLP-0256-B006}{start}
\begin{defn}[ټيورينګ ماشين]
يو \emph{ټيورينګ ماشين} \LR{$M$} هغه څلور‌يز \LR{$\langle Q, \Sigma, q_0,
\delta\rangle$} دے چې له لاندې توکو جوړ دے:
\begin{enumerate}
\item د \emph{حالتونو} يو متناهي سټ~\LR{$Q$}،
\item يوه متناهي \emph{الفبا} \LR{$\Sigma$} چې \LR{$\TMendtape$} او
  \LR{$\TMblank$} پکښې شامل دي،
\item يو \emph{لومړنی حالت}~\LR{$q_0 \in Q$}،
\item د \emph{لارښوونو} يو متناهي سټ~\LR{$\delta\colon Q \times \Sigma
  \pto Q \times \Sigma \times \{\TMleft, \TMright, \TMstay\}$}۔
\end{enumerate}
جزوي تابع~\LR{$\delta$} ته د~\LR{$M$} \emph{د حالت بدلون تابع} هم ويل کېږي۔
\end{defn}
\olpseganchor{OLP-0256-B006}{end}

\olpseganchor{OLP-0256-B007}{start}
\begin{explain}
موږ فرض کوو چې پټه يوازې په يوه لوري کښې نامتناهي ده۔ ځکه نو ګټوره
ده چې ځانګړې نښه~\LR{$\TMendtape$} د پټې د کيڼ پای د نښې په توګه وټاکو۔
په دې سره د ټيورينګ ماشين پروګرامونه په اسانۍ پوهېدے شي چې کله له پټې
څخه د ``لوېدو په خطر'' کښې دي۔ موږ دا فرض کولے شو چې دا نښه هېڅکله
نۀ بدلېږي، يعنې که \LR{$\delta(q,\TMendtape)$} تعريف شوې وي، نو
\LR{$\delta(q,\TMendtape) = \tuple{q', \TMendtape, x}$} وي۔ ځينې درسي
کتابونه دا کار کوي؛ موږ ئې نۀ کوو۔ د خپل ټيورينګ ماشين د جوړولو پر
وخت يوازې دومره پام وکړئ چې~\LR{$\TMendtape$} هېڅکله بدله نۀ کړي۔ سربېره
پر دې، داسې حالتونه هم شته چې د دې نښې د بدلولو اجازه څه ګټور
انعطاف برابروي۔
\end{explain}
\olpseganchor{OLP-0256-B007}{end}

\olpseganchor{OLP-0256-B008}{start}
\begin{ex}
\emph{جفت ماشين:} جفت ماشين په صوري ډول هغه څلور‌يز
\LR{$\tuple{Q, \Sigma, q_0, \delta}$} دے چې
\begin{align*}
Q & = \{ q_0, q_1 \} \\
\Sigma & = \{ \TMendtape, \TMblank, \TMstroke \}, \\
\delta(q_0, \TMstroke) & = \tuple{q_1, \TMstroke, \TMright},\\
\delta(q_1, \TMstroke) & = \tuple{q_0, \TMstroke, \TMright},\\
\delta(q_1, \TMblank)  & = \tuple{q_1, \TMblank, \TMright}.
\end{align*}
\end{ex}
\olpseganchor{OLP-0256-B008}{end}

\olpseganchor{OLP-0257-B004}{start}
\olfileid{cmp}{tur}{con}
\olsection{تشکيلونه او محاسبې}
\olpseganchor{OLP-0257-B004}{end}

\olpseganchor{OLP-0257-B005}{start}
\begin{explain}
هغه مخکښېنی بحث را ياد کړئ چې د جفت ماشين تشکيلونه مو پکښې يو په بل
پسې څارل۔ د پټې، لوست-ليک سر او ټيورينګ ماشين پروګرام خيالي ميکانيزم
په حقيقت کښې يوازې د ټيورينګ ماشين د محاسبې د ليدلو يوه شهودي لار ده۔
په صوري ډول، پر يوه ټاکلي ننوت د ټيورينګ ماشين محاسبه د
\emph{تشکيلونو} د يوې لړۍ په توګه تعريفولے شو---او تشکيل په خپل وار د
نښو يوه لړۍ (چې د محاسبې په يوه ټاکلي پړاو کښې د پټې له منځپانګې سره
سمون لري)، د لوست-ليک سر د ځاے ښودونکے يو عدد، او يو حالت دے۔ د دې
توکو په وسيله ټاکلے شو چې ټيورينګ ماشين~\LR{$M$} پر يوه ټاکلي ننوت څه
محاسبه کوي۔
\end{explain}
\olpseganchor{OLP-0257-B005}{end}

\olpseganchor{OLP-0257-B006}{start}
\begin{defn}[تشکيل]
د ټيورينګ ماشين \LR{$M = \tuple{Q, \Sigma, q_0,
\delta}$} يو \emph{تشکيل} هغه درې‌يز \LR{$\tuple{C, m, q}$} دے چې
\begin{enumerate}
\item \LR{$C \in \Sigma^*$} د~\LR{$\Sigma$} د نښو يوه متناهي لړۍ ده،
\item \LR{$m \in \Nat$} يو داسې عدد دے چې \LR{$m < \len{C}$}، او
\item \LR{$q \in Q$}۔
\end{enumerate}
په شهودي ډول، لړۍ~\LR{$C$} د پټې منځپانګه ده (له تر ټولو کيڼ مربع څخه تر
وروستي ناتش يا مخکښې ليدل شوي مربع پورې د ټولو مربعونو نښې)، \LR{$m$} د
هغه مربع شمېره ده چې لوست-ليک سر ئې لولي (تر ټولو کيڼ مربع په~\LR{$0$}
شمېرلو سره)، او~\LR{$q$} د ماشين اوسنی حالت دے۔
\end{defn}
\olpseganchor{OLP-0257-B006}{end}

\olpseganchor{OLP-0257-B007}{start}
\begin{explain}
د ټيورينګ ماشين ممکن ننوت د نښو يوه لړۍ ده، او عموماً داسې لړۍ وي چې
يو عدد په يوه بڼه کوډوي۔ د ټيورينګ ماشين لومړنی تشکيل هغه تشکيل دے
چې ماشين پکښې پر همدغه ننوت کار پيلوي: پټه د پټې د پای نښه لري او
سملاسي ورپسې ښي لوري ته پر مربعونو ننوت ليکل شوے وي، لوست-ليک سر د
ننوت تر ټولو کيڼ مربع (يعنې د کيڼ پای د نښې ښي لوري ته مربع) لولي،
او ميکانيزم په ټاکلي لومړني حالت~\LR{$q_0$} کښې وي۔
\end{explain}
\olpseganchor{OLP-0257-B007}{end}

\olpseganchor{OLP-0257-B008}{start}
\begin{defn}[لومړنی تشکيل]
د ناتش ننوت \LR{$I \in \Sigma^*$} دپاره د~\LR{$M$} \emph{لومړنی تشکيل} دا دے:
\[
\tuple{\TMendtape \frown I, 1, q_0}.
\]
که ننوت تشه لړۍ وي، لومړنی تشکيل دا دے:
\[
\tuple{\TMendtape \frown \TMblank, 1, q_0}.
\]
\end{defn}
\olpseganchor{OLP-0257-B008}{end}

\olpseganchor{OLP-0257-B009}{start}
\begin{explain}
نښه~\LR{$\frown$} د \emph{يوځاے کولو} دپاره ده---د ننوت رشته د کيڼ پای د
نښې سملاسي ښي لوري ته پيلېږي۔
\end{explain}
\olpseganchor{OLP-0257-B009}{end}

\olpseganchor{OLP-0257-B010}{start}
\begin{defn}
وايو چې تشکيل \LR{$\tuple{C, m, q}$} \emph{په يوه ګام کښې تشکيل
\LR{$\tuple{C', m', q'}$} ورکوي} (د~\LR{$M$} له مخې)، هله او يوازې هله چې
\begin{enumerate}
\item د~\LR{$C$} \LR{$m$}-مه نښه~\LR{$\sigma$} وي،
\item د~\LR{$M$} د لارښوونو سټ وټاکي چې \LR{$\delta(q, \sigma) =
  \tuple{q', \sigma', D}$}،
\item د~\LR{$C'$} \LR{$m$}-مه نښه~\LR{$\sigma'$} وي، او
\item
\begin{enumerate}
\item \LR{$D = L$} وي، او که \LR{$m>0$} وي نو \LR{$m' = m - 1$}، او که \LR{$m=0$} وي نو \LR{$m'=0$}، يا
\item \LR{$D = R$} او \LR{$m' = m + 1$} وي، يا
\item \LR{$D = N$} او \LR{$m' = m$} وي،
\end{enumerate}
\item که \LR{$m' = \len{C}$} وي، نو \LR{$\len{C'} = \len{C} + 1$} او د~\LR{$C'$}
  \LR{$m'$}-مه نښه~\LR{$\TMblank$} ده۔ که نۀ وي، \LR{$\len{C'}=\len{C}$}۔
\item د هر داسې~\LR{$i$} دپاره چې \LR{$i < \len{C}$} او \LR{$i \neq m$} وي، \LR{$C'(i) = C(i)$}،
\end{enumerate}
\end{defn}
\olpseganchor{OLP-0257-B010}{end}

\olpseganchor{OLP-0257-B011}{start}
\begin{defn}\ollabel{defn:run-output}
پر ننوت~\LR{$I$} د~\LR{$M$} يوه \emph{چلونه} د~\LR{$M$} د تشکيلونو داسې لړۍ~\LR{$C_i$}
ده چې~\LR{$C_0$} د ننوت~\LR{$I$} دپاره د~\LR{$M$} لومړنی تشکيل وي، او هر~\LR{$C_i$} په
يوه ګام کښې~\LR{$C_{i+1}$} ورکوي۔
\olpseganchor{OLP-0257-B011}{end}

\olpseganchor{OLP-0257-B012}{start}
وايو چې~\LR{$M$} \emph{پر ننوت~\LR{$I$} له~\LR{$k$} ګامونو وروسته درېږي} که \LR{$C_k =
\tuple{C, m, q}$} وي، د~\LR{$C$} \LR{$m$}-مه نښه~\LR{$\sigma$} وي، او \LR{$\delta(q,
\sigma)$} تعريف شوې نۀ وي۔ په دې صورت کښې د ننوت~\LR{$I$} دپاره د~\LR{$M$}
\emph{وت}~\LR{$O$} هله ټاکل کېږي چې~\LR{$O$} د نښو داسې رشته وي چې
يا تشه وي او يا په~\LR{$\TMblank$} نۀ ختمېږي، او د کوم~\LR{$j \in \Nat$} دپاره \LR{$C =
\TMendtape \concat O \concat \TMblank^j$} وي۔ که داسې رشته او طبيعي
عدد شتون ونۀ لري، دې درېدو ته وت نۀ ټاکل کېږي۔ که وت موجود وي، يکتا دے: د پټې له لومړۍ نښې وروسته يوازې د پای تشې نښې لرې کوو۔ (\LR{$\TMblank^j$} د~\LR{$j$}
شمېر~\LR{$\TMblank$} نښو يوه لړۍ ده۔)
\end{defn}
\olpseganchor{OLP-0257-B012}{end}

\olpseganchor{OLP-0257-B013}{start}
\begin{explain}
د دې تعريف له مخې د~\LR{$M$} هر ټاکل شوی وت~\LR{$O$} تل له~\LR{$\TMblank$} پرته په
بلې نښې ختمېږي؛ يا که په وخت~\LR{$k$} کښې ټوله پټه له~\LR{$\TMblank$} څخه ډکه
وي (له تر ټولو کيڼې~\LR{$\TMendtape$} پرته)، نو~\LR{$O$} تشه رشته ده۔ هغه
درېدلې چلونه چې د پای نښې له اړينې بڼې سره نۀ وي، هېڅ وت نۀ لري۔
\end{explain}
\olpseganchor{OLP-0257-B013}{end}

\olpseganchor{OLP-0258-B004}{start}
\olfileid{tur}{mac}{una}
\olsection{د عددونو يوګونې څرګندونه}
\olpseganchor{OLP-0258-B004}{end}

\olpseganchor{OLP-0258-B005}{start}
\begin{explain}
ټيورينګ ماشينونه پر خپله پټه د ليکل شوو نښو پر لړيو کار کوي۔ د
ټيورينګ ماشين د الفبا له مخې، د نښو دا لړۍ بېلابېل ننوتونه او وتونه
څرګندولے شي۔ البته، هغه ټيورينګ ماشينونه ځانګړې پاملرنه غواړي چې
\emph{حسابي} تابعې، يعنې د طبيعي عددونو تابعې، محاسبه کوي۔ د مثبتو
صحيح عددونو د څرګندولو يوه ساده لار دا ده چې د يوې نښې~\LR{$\TMstroke$}
د لړيو په توګه کوډ شي۔ که \LR{$n \in \Nat$} وي، نو \LR{$\TMstroke^n$} دې هغه
وخت تشه لړۍ وي چې \LR{$n = 0$} وي؛ او که نۀ وي، هغه لړۍ دې وي چې کره
\LR{$n$} شمېر~\LR{$\TMstroke$} نښې لري۔
\end{explain}
\olpseganchor{OLP-0258-B005}{end}

\olpseganchor{OLP-0258-B006}{start}
\begin{defn}[محاسبه]
ټيورينګ ماشين~\LR{$M$} هله او يوازې هله تابع \LR{$f\colon \Nat^k \to \Nat$}
\emph{محاسبه کوي} چې~\LR{$M$} پر دې ننوت ودرېږي:
\[
\TMstroke^{n_1} \TMblank \TMstroke^{n_2} \TMblank \dots \TMblank \TMstroke^{n_k}
\]
او وت ئې \LR{$\TMstroke^{f(n_1, \dots, n_k)}$} وي۔
\end{defn}
\olpseganchor{OLP-0258-B006}{end}

\olpseganchor{OLP-0258-B007}{start}
\begin{prob}
  دا تعريف کړئ چې ټيورينګ ماشين~\LR{$M$} کله تابع
  \LR{$f\colon \Nat^k \to \Nat^m$} محاسبه کوي۔
\end{prob}
\olpseganchor{OLP-0258-B007}{end}

\olpseganchor{OLP-0258-B008}{start}
\begin{ex}\ollabel{ex:adder}
\emph{جمع:}
راځئ داسې ماشين جوړ کړو چې تابع \LR{$f(n,m) = n + m$} محاسبه کړي۔ د دې
دپاره داسې ماشين پکار دے چې پر پټه د~\LR{$n$} او~\LR{$m$} اوږدوالي لرونکو د
\LR{$\TMstroke$} دوو غونډونو سره پيل وکړي، او د~\LR{$n+m$} شمېر~\LR{$\TMstroke$}
ګانو په يوه غونډ ودرېږي۔ د~\LR{$\TMstroke$} د ننوت دواړه غونډونه په يوه
\LR{$\TMblank$} سره بېلېږي؛ نو يوه لار دا ده چې په هغه مربع کښې يوه کرښه
وليکو چې~\LR{$\TMblank$} لري، او وروستۍ~\LR{$\TMstroke$} پاکه کړو۔
\begin{figure}\[
\begin{tikzpicture}[->,>=stealth',shorten >=1pt,auto,node distance=2.8cm,
                    semithick]
  \tikzstyle{every state}=[fill=none,draw=black,text=black]
\olpseganchor{OLP-0258-B008}{end}

\olpseganchor{OLP-0258-B009}{start}
  \node[initial,state]         (A)              {\LR{$q_0$}};
  \node[state]         (B) [right of=A] {\LR{$q_1$}};
  \node[state]         (C) [right of=B] {\LR{$q_2$}};
\olpseganchor{OLP-0258-B009}{end}

\olpseganchor{OLP-0258-B010}{start}
  \path (A) edge node {\TMtrans{\TMblank}{\TMstroke}{\TMstay}} (B)
           edge [loop above] node {\TMtrans{\TMstroke}{\TMstroke}{\TMright}} (B)
        (B) edge [loop above] node {\TMtrans{\TMstroke}{\TMstroke}{\TMright}} (B)
            edge node {\TMtrans{\TMblank}{\TMblank}{\TMleft}} (C)
        (C) edge [loop above] node {\TMtrans{\TMstroke}{\TMblank}{\TMstay}} (C);
\end{tikzpicture}
\]\caption{هغه ماشين چې \LR{$f(x,y) = x+y$} محاسبه کوي}
\ollabel{fig:adder}
\end{figure}
\end{ex}
\olpseganchor{OLP-0258-B010}{end}

\olpseganchor{OLP-0258-B011}{start}
\begin{prob}
په~\olref[tur][mac][una]{ex:adder} کښې د ماشين تشکيلونه د
ننوت~\LR{$\tuple{3,2}$} دپاره يو په بل پسې وڅارئ۔ که ماشين \LR{$0+0$} محاسبه
کړي، څه پېښېږي؟
\end{prob}
\olpseganchor{OLP-0258-B011}{end}

\olpseganchor{OLP-0258-B012}{start}
\begin{explain}
په~\olref[rep]{ex:doubler} کښې مو د داسې ټيورينګ ماشين بېلګه ورکړه
چې د~\LR{$\TMstroke$} نښو يوه لړۍ د ننوت په توګه اخلي او پر پټه د هغې د
دوه‌چنده شمېر~\LR{$\TMstroke$} نښو په لړۍ درېږي---يعنې دوه‌چنده کوونکے
ماشين۔ خو څرنګه چې په وت کښې د~\LR{$\TMstroke$} نښو د دوه‌چنده غونډ کيڼ
لوري ته~\LR{$\TMblank$} نښې شته، نو دا په حقيقت کښې تابع \LR{$f(x) = 2x$} نۀ
محاسبه کوي، که څه هم ښايي تاسو داسې ګڼلې وي۔ موږ به د دې د سمولو
دوه لارې وښيو۔
\end{explain}
\olpseganchor{OLP-0258-B012}{end}

\olpseganchor{OLP-0258-B013}{start}
\begin{ex}
په~\olref{fig:doubler-disc} کښې ماشين تابع \LR{$f(x) = 2x$} محاسبه کوي۔
دا ماشين د~\olref[rep]{ex:doubler} د ماشين په شان د ننوت د هرې
\LR{$\TMstroke$} په بدل کښې ننوت نۀ پاکوي او په لرې ښي لوري کښې دوه
\LR{$\TMstroke$} نښې نۀ ليکي، بلکې د ننوت د هرې~\LR{$\TMstroke$} دپاره ښي
لوري ته يوه~\LR{$\TMstroke$} ور زياتوي۔ دا بايد په ياد وساتي چې ننوت چېرته
ختمېږي؛ نو د ننوت او ور زياتو شوو کرښو تر منځ يوه~\LR{$\TMblank$} پرېږدي
او په پای کښې ئې په~\LR{$\TMstroke$} ډکوي۔ موږ بايد دا هم ``په ياد'' ولرو
چې په ننوت کښې چېرته يو؛ نو د ننوت په غونډ کښې يوه~\LR{$\TMstroke$} په
لنډمهاله توګه په~\LR{$\TMblank$} بدلوو۔
  \begin{figure}
\[
\begin{tikzpicture}[->,>=stealth',shorten >=1pt,auto,node distance=2.8cm,
                    semithick]
  \tikzstyle{every state}=[fill=none,draw=black,text=black]
\olpseganchor{OLP-0258-B013}{end}

\olpseganchor{OLP-0258-B014}{start}
  \node[initial,state] (0)              {\LR{$q_0$}};
  \node[state]         (1) [above of=0] {\LR{$q_1$}};
  \node[state]         (2) [above of=1] {\LR{$q_2$}};
  \node[state]         (3) [right of=2] {\LR{$q_3$}};
  \node[state]         (4) [below of=3] {\LR{$q_4$}};
  \node[state]         (5) [below of=4] {\LR{$q_5$}};
  \node[state]         (6) [above right of=4] {\LR{$q_6$}};
  \node[state]         (7) [below of=6] {\LR{$q_7$}};
  \node[state]         (8) [below of=7] {\LR{$q_8$}};
\olpseganchor{OLP-0258-B014}{end}

\olpseganchor{OLP-0258-B015}{start}
  \path (0) edge node {\TMtrans{\TMstroke}{\TMblank}{\TMright}} (1)
%    (0) edge node {\TMtrans{\TMblank}{\TMblank}{\TMstay}} (h)
    (1) edge [loop left] node {\TMtrans{\TMstroke}{\TMstroke}{\TMright}} (1)
      edge node {\TMtrans{\TMblank}{\TMblank}{\TMright}} (2)
    (2) edge [loop above] node {\TMtrans{\TMstroke}{\TMstroke}{\TMright}} (2)
        edge node {\TMtrans{\TMblank}{\TMstroke}{\TMleft}} (3)
    (3) edge [loop above] node {\TMtrans{\TMstroke}{\TMstroke}{\TMleft}} (3)
        edge node[left] {\TMtrans{\TMblank}{\TMblank}{\TMleft}} (4)
    (4) edge        node[left] {\TMtrans{\TMstroke}{\TMstroke}{\TMleft}} (5)
    (5) edge [loop below] node {\TMtrans{\TMstroke}{\TMstroke}{\TMleft}} (5)
        edge              node {\TMtrans{\TMblank}{\TMstroke}{\TMright}} (0)
    (4) edge      node[sloped] {\TMtrans{\TMblank}{\TMstroke}{\TMright}} (6)
    (6) edge              node {\TMtrans{\TMblank}{\TMstroke}{\TMright}} (7)
    (7) edge [loop right] node {\TMtrans{\TMstroke}{\TMstroke}{\TMright}} (7)
        edge              node {\TMtrans{\TMblank}{\TMblank}{\TMleft}} (8)
    (8) edge [loop right] node {\TMtrans{\TMstroke}{\TMblank}{\TMstay}} (8);
    \end{tikzpicture}
\]
\caption{هغه ماشين چې \LR{$f(x) = 2x$} محاسبه کوي}
\ollabel{fig:doubler-disc}
\end{figure}
\end{ex}
\olpseganchor{OLP-0258-B015}{end}

\olpseganchor{OLP-0258-B016}{start}
\begin{ex}\ollabel{ex:mover}
د تابع \LR{$f(x) = 2x$} د محاسبې دويمه شونتيا دا ده چې اصلي دوه‌چنده
کوونکے ماشين وساتو، خو په پای کښې داسې حالتونه او لارښوونې ور زياتې
کړو چې د کرښو دوه‌چنده غونډ د پټې تر ټولو کيڼ لوري ته يوسي۔ په
\olref{fig:mover} کښې ماشين همدا وروستۍ برخه ترسره کوي: که دا پر داسې
پټه پيل شي چې لومړے د~\LR{$\TMblank$} نښو يو غونډ او بيا د~\LR{$\TMstroke$}
نښو يو غونډ لري (او سر د~\LR{$\TMblank$} نښو د غونډ په هر ځاے کښې وي)،
نو~\LR{$\TMstroke$} نښې يوه يوه پاکوي او د پټې په پيل کښې ئې ليکي۔ د دې
د معلومولو دپاره چې کار کله بشپړ شو، ماشين لومړے د~\LR{$\TMstroke$} نښو
د غونډ پای په~\LR{$\TMendtape$} نښه نښانوي، او دا نښه په پای کښې پاکېږي۔
موږ د حالتونو شمېرل له~\LR{$q_6$} څخه پيل کړي، څو دوي دوه‌چنده کوونکي
ماشين ته ور زياتېدے شي۔ تاسو يوازې يوې اضافي لارښوونې ته اړتيا لرئ:
\LR{$\delta(q_0, \TMblank) = \tuple{q_6,\TMblank,\TMstay}$}؛ يعنې له~\LR{$q_0$}
څخه~\LR{$q_6$} ته يو داسې غشی چې نښه ئې~\LR{$\TMtrans{\TMblank}{\TMblank}{\TMstay}$}
وي۔ (يوه نرۍ ستونزه شته: لاس ته راغلے ماشين د ننوت~\LR{$x=0$} دپاره کار
نۀ کوي۔ موږ دا د تمرين په توګه پرېږدو۔)
\begin{figure}
  \[
  \begin{tikzpicture}[->,>=stealth',shorten >=1pt,auto,node distance=2.8cm,
                      semithick]
    \tikzstyle{every state}=[fill=none,draw=black,text=black]
    \node[initial,state] (6)              {\LR{$q_6$}};
    \node[state]         (7) [right of=6] {\LR{$q_7$}};
    \node[state]         (8) [right of=7] {\LR{$q_8$}};
    \node[state]         (9) [below of=8] {\LR{$q_9$}};
    \node[state]         (10) [left of=9] {\LR{$q_{10}$}};
    \node[state]         (11) [left of=10]  {\LR{$q_{11}$}};
    \node[state]         (12) [below of=11]  {\LR{$q_{12}$}};
    \node[state]         (13) [right of=12] {\LR{$q_{13}$}};
    \node[state]         (14) [right of=13] {\LR{$q_{14}$}};
    \path
    (6)  edge node {\TMtrans{\TMblank}{\TMblank}{\TMright}} (7)
    (7)  edge [loop above] node {\TMtrans{\TMblank}{\TMblank}{\TMright}} (7)
         edge node {\TMtrans{\TMstroke}{\TMstroke}{\TMright}} (8)
    (8)  edge [loop above] node {\TMtrans{\TMstroke}{\TMstroke}{\TMright}} (8)
         edge node {\TMtrans{\TMblank}{\TMendtape}{\TMleft}} (9)
    (9)  edge [loop right] node {\TMtrans{\TMstroke}{\TMstroke}{\TMleft}} (9)
         edge node {\TMtrans{\TMblank}{\TMblank}{\TMright}} (10)
    (10) edge node[above] {\TMtrans{\TMstroke}{\TMblank}{\TMleft}} (11)
    (11) edge [loop above] node {\TMtrans{\TMblank}{\TMblank}{\TMleft}} (11)
         edge node[left] {\begin{tabular}{@{}l@{}}
          \TMtrans{\TMendtape}{\TMendtape}{\TMright}\\
          \TMtrans{\TMstroke}{\TMstroke}{\TMright}
         \end{tabular}} (12)
    (12) edge node[] {\TMtrans{\TMblank}{\TMstroke}{\TMright}} (13)
    (13) edge [loop below] node {\TMtrans{\TMblank}{\TMblank}{\TMright}} (13)
         edge node[sloped] {\TMtrans{\TMstroke}{\TMblank}{\TMleft}} (11)
    (13) edge node {\TMtrans{\TMendtape}{\TMblank}{\TMstay}} (14);
  \end{tikzpicture}
  \]
  \caption{د~\LR{$\TMstroke$} نښو د يو غونډ کيڼ لوري ته وړل}
  \ollabel{fig:mover}
\end{figure}
\end{ex}
\olpseganchor{OLP-0258-B016}{end}

\olpseganchor{OLP-0258-B017}{start}
\begin{prob}
په~\olref[tur][mac][una]{ex:mover} کښې مو داسې ماشين بيان کړ چې د
\olref[tur][mac][rep]{fig:doubler} د دوه‌چنده کوونکي ماشين او د
\olref[tur][mac][una]{fig:mover} د لېږدوونکي ماشين له يوځاے کولو څخه
جوړ دے۔ که دا يوځاے شوے ماشين پر ننوت~\LR{$x=0$}، يعنې پر تشه پټه پيل
کړئ، څه پېښېږي؟ ماشين به څنګه سم کړئ چې په دې صورت کښې د~\LR{$2x=0$} له
وت سره ودرېږي؟ (تاسو بايد دا د يوه حالت او يوه حالت بدلون په ور
زياتولو وکړے شئ۔)
\end{prob}
\olpseganchor{OLP-0258-B017}{end}

\olpseganchor{OLP-0258-B018}{start}
\begin{prob}
\emph{تفريق:} داسې ټيورينګ ماشين طرح کړئ چې د~\LR{$n$} او~\LR{$m$} اوږدوالي
لرونکې د کرښو دوه ناتشې رشتې د ننوت په توګه ورکړل شي، چې \LR{$n > m$}،
نو تابع \LR{$f(n,m) = n - m$} محاسبه کړي۔
\end{prob}
\olpseganchor{OLP-0258-B018}{end}

\olpseganchor{OLP-0258-B019}{start}
\begin{prob}
\emph{مساوات:} داسې ټيورينګ ماشين طرح کړئ چې لاندې تابع محاسبه کړي:
\[
\fn{equality}(n,m) =
\begin{cases}
  \text{\RL{1}} & \text{\RL{که~\LR{$n = m$} وي}} \\
  \text{\RL{0}} & \text{\RL{که~\LR{$n \neq m$} وي}}
\end{cases}
\]
چې پکښې~\LR{$n$} او~\LR{$m \in \PosInt$} دي۔
\end{prob}
\olpseganchor{OLP-0258-B019}{end}

\olpseganchor{OLP-0258-B020}{start}
\begin{prob}
د تابع \LR{$\min(x,y)$} د محاسبې دپاره يو ټيورينګ ماشين طرح کړئ، چې
\LR{$x$} او~\LR{$y$} داسې مثبت صحيح عددونه وي چې پر پټه د~\LR{$\TMstroke$} نښو په
رشتو څرګند شوي او په يوه~\LR{$\TMblank$} بېل شوي وي۔ تاسو د ماشين په
الفبا کښې اضافي نښې کارولے شئ۔
\olpseganchor{OLP-0258-B020}{end}

\olpseganchor{OLP-0258-B021}{start}
تابع~\LR{$\min$} له خپلو ارزښتمنو څخه تر ټولو کوچنی ارزښت ټاکي؛ نو
\LR{$\min(3,5)=3$}، \LR{$\min(20,16)=16$}، او \LR{$\min(4,4)=4$}، او داسې نور۔
\end{prob}
\olpseganchor{OLP-0258-B021}{end}

\olpseganchor{OLP-0258-B022}{start}
\begin{defn}
  ټيورينګ ماشين~\LR{$M$} هله او يوازې هله جزوي تابع \LR{$f\colon \Nat^k
  \pto \Nat$} محاسبه کوي چې
  \begin{enumerate}
    \item که \LR{$f(n_1, \dots, n_k) = m$} وي، نو~\LR{$M$} پر ننوت
    \LR{$\TMstroke^{n_1}\concat\TMblank\concat
    \dots \concat\TMblank\concat\TMstroke^{n_k}$} د وت~\LR{$\TMstroke^{m}$} سره درېږي۔
    \item که \LR{$f(n_1, \dots, n_k)$} تعريف شوې نۀ وي، نو~\LR{$M$} بېخي نۀ
    درېږي، يا داسې درېږي چې معتبر وت ورته نۀ ټاکل کېږي، يا له داسې
    وت سره درېږي چې د~\LR{$\TMstroke$} نښو يوازينی غونډ نۀ وي۔
  \end{enumerate}
\end{defn}
\olpseganchor{OLP-0258-B022}{end}

\olpseganchor{OLP-0259-B004}{start}
\olfileid{tur}{mac}{hal}
\olsection{درېدني حالتونه}
\olpseganchor{OLP-0259-B004}{end}

\olpseganchor{OLP-0259-B005}{start}
\begin{explain}
که څه هم موږ خپل ماشينونه داسې تعريف کړي چې يوازې د عملي کېدونکې
لارښوونې د نشتوالي پر وخت درېږي، د ټيورينګ ماشينونو په عامو
څرګندونو کښې يو ځانګړی \emph{درېدنی حالت}~\LR{$h$} وي، داسې چې
\LR{$h \in Q$}۔
\olpseganchor{OLP-0259-B005}{end}

\olpseganchor{OLP-0259-B006}{start}
د درېدني حالت تر شا نظر ساده دے: کله چې ماشين خپل کار بشپړ کړي (د
ننوت منلو ته چمتو وي، يا د وت ليکل ئې بشپړ کړي وي)، داسې حالت~\LR{$h$}
ته ننوځي چې هلته درېږي۔ ځينې ماشينونه دوه درېدني حالتونه لري: يو
ننوت مني او بل ننوت ردوي۔
\end{explain}
\olpseganchor{OLP-0259-B006}{end}

\olpseganchor{OLP-0259-B007}{start}
\begin{ex}\emph{درېدني حالتونه}۔
د دې مفهوم د روښانولو دپاره راځئ د جفت ماشين له يو بدلون څخه پيل
وکړو۔ د دې پر ځاے چې ماشين د جفت ننوت په صورت کښې په حالت~\LR{$q_0$}
کښې ودرېږي، يوه لارښوونه ور زياتولے شو چې ماشين درېدني حالت ته
واستوي۔
\[
\begin{tikzpicture}[->,>=stealth',shorten >=1pt,auto,node distance=2.8cm,
                    semithick]
  \tikzstyle{every state}=[fill=none,draw=black,text=black]
\olpseganchor{OLP-0259-B007}{end}

\olpseganchor{OLP-0259-B008}{start}
  \node[initial, state]         (A)                     {\LR{$q_0$}};
  \node[state]         (B) [right of=A] {\LR{$q_1$}};
  \node[state]         (C) [below of=A] {\LR{$h$}};
\olpseganchor{OLP-0259-B008}{end}

\olpseganchor{OLP-0259-B009}{start}
  \path (A) edge [bend left] node {\TMtrans{\TMstroke}{\TMstroke}{R}} (B)
  	    edge node {\TMtrans{\TMblank}{\TMblank}{N}} (C)
        (B) edge [loop above] node {\TMtrans{\TMblank}{\TMblank}{R}} (B)
            edge [bend left] node {\TMtrans{\TMstroke}{\TMstroke}{R}} (A);
\end{tikzpicture}
\]
\olpseganchor{OLP-0259-B009}{end}

\olpseganchor{OLP-0259-B010}{start}
راځئ بېلګه نوره هم پراخه کړو۔ کله چې ماشين ومومي چې ننوت طاق دے،
هېڅکله نۀ درېږي۔ موږ د کړۍ وهونکې لارښوونې پر ځاے داسې لارښوونه
ايښودلے شو چې رد حالت~\LR{$r$} ته ځي، او په دې توګه ماشين ته يو
\emph{رد} حالت ور زياتولے شو۔
\[
\begin{tikzpicture}[->,>=stealth',shorten >=1pt,auto,node distance=2.8cm,
                    semithick]
  \tikzstyle{every state}=[fill=none,draw=black,text=black]
\olpseganchor{OLP-0259-B010}{end}

\olpseganchor{OLP-0259-B011}{start}
  \node[initial,state]         (A)                     {\LR{$q_0$}};
  \node[state]         (B) [right of=A] {\LR{$q_1$}};
  \node[state]         (C) [below of=A] {\LR{$h$}};
  \node[state]         (D) [below of=B] {\LR{$r$}};
\olpseganchor{OLP-0259-B011}{end}

\olpseganchor{OLP-0259-B012}{start}
  \path (A) edge [bend left] node {\TMtrans{\TMstroke}{\TMstroke}{R}} (B)
  	    edge node {\TMtrans{\TMblank}{\TMblank}{N}} (C)
        (B) edge node {\TMtrans{\TMblank}{\TMblank}{N}} (D)
            edge [bend left] node {\TMtrans{\TMstroke}{\TMstroke}{R}} (A);
\end{tikzpicture}
\]
\end{ex}
\olpseganchor{OLP-0259-B012}{end}

\olpseganchor{OLP-0259-B013}{start}
\begin{explain}
د ځانګړي درېدني حالت زياتول د دې په څېر حالتونو کښې ګټور وي، ځکه په
څرګنده ښيي چې ماشين ټاکلي ننوتونه کله مني يا ردوي۔ خو پام کول مهم دي
چې د ځانګړي درېدني حالت په زياتولو هېڅ محاسبوي ځواک نۀ زياتېږي۔ همداراز
د درېدو لږ صوري مفهوم خپلې ګټې لري۔ د درېدو هغه تعريف چې تر اوسه مو
په دې څپرکي کښې کارولے، د \emph{درېدو مسئلې} ثبوت شهودي او څرګندول ئې
اسانه کوي۔ له همدې امله موږ خپل اصلي تعريف ته دوام ورکوو۔
\end{explain}
\olpseganchor{OLP-0259-B013}{end}

\olpseganchor{OLP-0260-B004}{start}
\olfileid{tur}{mac}{dis}
\olsection{منضبط ماشينونه}
\olpseganchor{OLP-0260-B004}{end}

\olpseganchor{OLP-0260-B005}{start}
\begin{explain}
په برخه~\olref[hal]{sec} کښې مو هغه ټيورينګ ماشينونه وڅېړل چې يو
ټاکلی درېدنی حالت~\LR{$h$} لري---داسې ماشينونه، که بېخي درېږي، هرومرو په
حالت~\LR{$h$} کښې درېږي۔ په دې معنا د يوه درېدني حالت لرونکي ماشينونه د
هغو ټيورينګ ماشينونو په پرتله ډېر ``منضبط'' دي چې موږ ئې په عمومي
ډول منو۔ نور بنديزونه هم د ټيورينګ ماشينونو پر چلند لګولے شو۔ مثلاً،
موږ ټيورينګ ماشينونه له دې هم نۀ دي منع کړي چې کله د پټې د پای نښه
پر مربع~\LR{$0$} پاکه کړي، يا له مربع~\LR{$0$} څخه کيڼ لوري ته د تلو هڅه وکړي۔
(زموږ تعريف وايي چې په دې صورت کښې سر بس پر مربع~\LR{$0$} پاتې کېږي؛ په
نورو تعريفونو کښې ماشين درېږي۔) همداراز کله ناکله دا فرض هم ګټور وي
چې ټيورينګ ماشين، که بېخي درېږي، پر مربع~\LR{$1$} ودرېږي۔
\end{explain}
\olpseganchor{OLP-0260-B005}{end}

\olpseganchor{OLP-0260-B006}{start}
\begin{defn}\ollabel{defn:disciplined}
ټيورينګ ماشين~\LR{$M$} هله او يوازې هله \emph{منضبط} دے چې
\begin{enumerate}
    \item يو ټاکلی يوازينی درېدنی حالت~\LR{$h$} ولري،
    \item که بېخي درېږي، د مربع~\LR{$1$} د لوستلو پر وخت ودرېږي،
    \item پر مربع~\LR{$0$} د~\LR{$\TMendtape$} نښه هېڅکله پاکه نۀ کړي، او
    \item له مربع~\LR{$0$} څخه کيڼ لوري ته د تلو هڅه هېڅکله ونۀ کړي۔
\end{enumerate}
\end{defn}
\olpseganchor{OLP-0260-B006}{end}

\olpseganchor{OLP-0260-B007}{start}
\begin{explain}
موږ مخکښې بحث کړے چې هر ټيورينګ ماشين په داسې ماشين بدلېدے شي چې
هماغه چلند او يو ټاکلی درېدنی حالت ولري۔ دا کار يوازې د نوي حالت~\LR{$h$}
په زياتولو، او د هر داسې جوړې~\LR{$\tuple{q,\sigma}$} دپاره د لارښوونې
\LR{$\delta(q, \sigma) = \tuple{h, \sigma, N}$} په زياتولو کېږي چې اصلي
\LR{$\delta$} پرې تعريف شوې نۀ وي۔ که څه هم ثبوت ئې ستړے کوونکے دے، خو
دا رښتيا ده چې هر ټيورينګ ماشين~\LR{$M$} په داسې منضبط ټيورينګ ماشين~\LR{$M'$}
بدلېدے شي چې پر هماغو ننوتونو درېږي او هماغه وت پيدا کوي۔ مثلاً، که
ټيورينګ ماشين ودرېږي او پر مربع~\LR{$1$} نۀ وي، داسې لارښوونې ور زياتولے
شو چې سر تر هغه کيڼ لوري ته بوځي څو د پټې د پای نښه پيدا کړي، بيا يو
مربع ښي لوري ته لاړ شي، او بيا ودرېږي۔ د نورو شرطونو د پوره کولو پر
لار به تاسو پخپله فکر وکړئ۔
\end{explain}
\olpseganchor{OLP-0260-B007}{end}

\olpseganchor{OLP-0260-B008}{start}
\begin{ex}
    په~\olref{fig:adder-disc} کښې د جمع ماشين چې په
    \olref[una]{ex:adder} کښې و، په منضبط ماشين بدلوو۔
    \begin{figure}\[
\begin{tikzpicture}[->,>=stealth',shorten >=1pt,auto,node distance=2.8cm,
                    semithick]
  \tikzstyle{every state}=[fill=none,draw=black,text=black]
\olpseganchor{OLP-0260-B008}{end}

\olpseganchor{OLP-0260-B009}{start}
  \node[initial,state]         (A)              {\LR{$q_0$}};
  \node[state]         (B) [right of=A] {\LR{$q_1$}};
  \node[state]         (C) [below right of=B] {\LR{$q_2$}};
  \node[state]         (D) [below left of=C] {\LR{$q_3$}};
  \node[state]         (H) [left of=D] {\LR{$h$}};
\olpseganchor{OLP-0260-B009}{end}

\olpseganchor{OLP-0260-B010}{start}
  \path (A) edge node {\TMtrans{\TMblank}{\TMstroke}{\TMstay}} (B)
           edge [loop below] node {\TMtrans{\TMstroke}{\TMstroke}{\TMright}} (B)
        (B) edge [loop below] node {\TMtrans{\TMstroke}{\TMstroke}{\TMright}} (B)
            edge node {\TMtrans{\TMblank}{\TMblank}{\TMleft}} (C)
        (C) edge node {\TMtrans{\TMstroke}{\TMblank}{\TMleft}} (D)
        (D) edge [loop above] node {\TMtrans{\TMstroke}{\TMstroke}{\TMleft}} (D)
        edge node {\TMtrans{\TMendtape}{\TMendtape}{\TMright}} (H);
\end{tikzpicture}
\]\caption{يو منضبط جمع کوونکے ماشين}
\ollabel{fig:adder-disc}
\end{figure}
\end{ex}
\olpseganchor{OLP-0260-B010}{end}

\olpseganchor{OLP-0260-B011}{start}
\begin{prop}\ollabel{prop:disciplined} د هر ټيورينګ ماشين~\LR{$M$} دپاره
داسې منضبط ټيورينګ ماشين~\LR{$M'$} شته چې که~\LR{$M$} د وت~\LR{$O$} سره درېږي،
هغه هم د وت~\LR{$O$} سره درېږي؛ او که~\LR{$M$} نۀ درېږي، هغه هم نۀ درېږي۔ په
ځانګړې توګه، هره تابع \LR{$f\colon\Nat^n \to \Nat$} چې د ټيورينګ ماشين
په وسيله محاسبه کېدونکې وي، د منضبط ټيورينګ ماشين په وسيله هم محاسبه
کېدونکې ده۔
\end{prop}
\olpseganchor{OLP-0260-B011}{end}

\olpseganchor{OLP-0260-B012}{start}
\begin{prob}\label{tur:mac:dis:prob:disc-succ}
داسې منضبط ماشين ورکړئ چې \LR{$f(x) = x+1$} محاسبه کړي۔
\end{prob}
\olpseganchor{OLP-0260-B012}{end}


\olpseganchor{OLP-0260-B013}{start}
\begin{prob}\label{tur:mac:dis:prob:copier}
داسې منضبط ماشين پيدا کړئ چې پر ننوت \LR{$\TMstroke^n$} پيل شي او وت
\LR{$\TMstroke^n \concat \TMblank \concat \TMstroke^n$} پيدا کړي۔
\end{prob}
\olpseganchor{OLP-0260-B013}{end}

\olpseganchor{OLP-0261-B004}{start}
\olfileid{tur}{mac}{cmb}
\olsection{د ټيورينګ ماشينونو يوځاے کول}
\olpseganchor{OLP-0261-B004}{end}

\olpseganchor{OLP-0261-B005}{start}
\begin{explain}
د ټيورينګ ماشينونو هغه بېلګې چې تر اوسه مو ليدلې، نسبتاً ساده وې۔
خو په حقيقت کښې هره مسئله چې په هره اوسنۍ پروګرام‌ليکنې ژبه حل کېدے
شي، د ټيورينګ ماشينونو په وسيله هم حل کېدے شي۔ د لا پېچلو ټيورينګ
ماشينونو د جوړولو دپاره بايد ځان ډاډه کړو چې ماشينونه يوځاے کولے شو؛
په دې توګه به يوه کړنلاره په ساده برخو ووېشو او د لا پېچلو مسئلو د
حل ماشينونه به جوړ کړو۔ که د پېچلې مسئلې په رغوونکو برخو د وېشلو
طبيعي لار پيدا کړو، مسئله په څو پړاوونو کښې حل کولے شو: څو ساده
ټيورينګ ماشينونه جوړوو او بيا ئې په يوه داسې ماشين کښې يوځاے کوو چې
ټوله مسئله حل کړي۔ دا خبره په بله برخه کښې د درېدو مسئلې د څېړلو پر
وخت ځانګړې مهمه ده۔
\olpseganchor{OLP-0261-B005}{end}

\olpseganchor{OLP-0261-B006}{start}
ټيورينګ ماشينونه \LR{$M = \tuple{Q, \Sigma, q_0, \delta}$} او
~\LR{$M' = \tuple{Q', \Sigma', q_0', \delta'}$} څنګه يوځاے کړو؟ اوس د
~\LR{$M$} له درېدو وروسته د پټې تشکيل د ماشين~\LR{$M'$} د يوې چلونې د ننوت
تشکيل په توګه کاروو۔ د داسې يوه ټيورينګ ماشين \LR{$M \frown M'$} د لاس
ته راوړلو دپاره لاندې کارونه وکړئ:
\begin{enumerate}
    \item د~\LR{$M'$} ټول حالتونه~\LR{$Q'$} بيا وشمېرئ (يا بيا ونوموئ)، څو
    \LR{$M$} او~\LR{$M'$} هېڅ ګډ حالت ونۀ لري (\LR{$Q \cap Q' = \emptyset$})۔
    \item د~\LR{$M \frown M'$} حالتونه \LR{$Q \cup Q'$} دي۔
    \item د پټې الفبا \LR{$\Sigma \cup \Sigma'$} ده۔
    \item لومړنی حالت~\LR{$q_0$} دے۔
    \item د حالت بدلون تابع هغه تابع~\LR{$\delta''$} ده چې داسې ورکول کېږي:
    \[\delta''(q,\sigma) =
    \begin{cases}
      \delta(q,\sigma) & \text{\RL{که \LR{$q \in Q$} وي او دا ارزښت تعريف شوے وي}}\\
      \delta'(q,\sigma) & \text{\RL{که \LR{$q \in Q'$} وي}}\\
      \tuple{q_0', \sigma, \TMstay} & \text{\RL{که \LR{$q \in Q$} وي او
      \LR{$\delta(q,\sigma)$} تعريف شوې نۀ وي}}
    \end{cases}\]
\end{enumerate}
لاس ته راغلے ماشين چې په حالت \LR{$q \in Q$} کښې وي، د~\LR{$M$} لارښوونې
کاروي، او چې په حالت~\LR{$q \in Q'$} کښې وي، د~\LR{$M'$} لارښوونې کاروي۔ کله
چې دا په داسې حالت \LR{$q \in Q$} کښې وي او داسې نښه~\LR{$\sigma$} لولي چې
~\LR{$M$} ورته د حالت بدلون نۀ لري (يعنې~\LR{$M$} به درېدلے و)، د~\LR{$M'$} لومړني
حالت ته ننوځي (او د پټې منځپانګه او د سر ځاے هماغسې پرېږدي)۔ که د دويم ماشين په حالت کښې د لوستل شوې نښې دپاره لارښوونه نۀ وي، ګډ ماشين هماغلته درېږي۔ هره جزوي تابع د ګډې الفبا پر هغو نويو نښو هم ناتعريف شوې ګڼو چې په خپلې الفبا کښې ئې نۀ وي۔
\olpseganchor{OLP-0261-B006}{end}

\olpseganchor{OLP-0261-B007}{start}
پام وکړئ، تر څو چې ماشين~\LR{$M$} منضبط نۀ وي، موږ نۀ پوهېږو چې د~\LR{$M$}
د درېدو پر وخت د پټې سر چېرته وي؛ نو د~\LR{$M$} په درېدني تشکيل کښې سر
اړين نۀ دے چې مربع~\LR{$1$} ولولي۔ د ماشينونو د يوځاے کولو پر وخت بايد
دا خبره په پام کښې وساتو۔
\end{explain}
\olpseganchor{OLP-0261-B007}{end}

\olpseganchor{OLP-0261-B008}{start}
\begin{ex}
\emph{د ماشينونو يوځاے کول:} موږ به داسې ماشين طرح کړو چې د~\LR{$n$} او
~\LR{$m$} اوږدوالي لرونکو د~\LR{$\TMstroke$} نښو له دوو غونډونو جوړ ننوت باندې
پيل شي او پر پټه د~\LR{$2(m+n)$} شمېر~\LR{$\TMstroke$} نښو په يوه غونډ
ودرېږي۔ د دې ماشين د جوړولو دپاره دوه اشنا ماشينونه يوځاے کولے شو:
د جمع ماشين او دوه‌چنده کوونکے۔ لومړے د جمع ماشين د حالت ډياګرام
کاږو۔
\[
\begin{tikzpicture}[->,>=stealth',shorten >=1pt,auto,node distance=2.8cm,
                    semithick]
  \tikzstyle{every state}=[fill=none,draw=black,text=black]
\olpseganchor{OLP-0261-B008}{end}

\olpseganchor{OLP-0261-B009}{start}
  \node[initial,state] (A)              {\LR{$q_0$}};
  \node[state]         (B) [right of=A] {\LR{$q_1$}};
  \node[state]         (C) [right of=B] {\LR{$q_2$}};
\olpseganchor{OLP-0261-B009}{end}

\olpseganchor{OLP-0261-B010}{start}
  \path (A) edge node {\TMtrans{\TMblank}{\TMstroke}{\TMstay}} (B)
            edge [loop above] node {\TMtrans{\TMstroke}{\TMstroke}{\TMright}} (B)
        (B) edge [loop above] node {\TMtrans{\TMstroke}{\TMstroke}{\TMright}} (B)
            edge node {\TMtrans{\TMblank}{\TMblank}{\TMleft}} (C)
        (C) edge [loop above] node {\TMtrans{\TMstroke}{\TMblank}{\TMstay}} (C);
\end{tikzpicture}
\]
موږ نۀ غواړو چې ماشين په حالت~\LR{$q_2$} کښې ودرېږي، بلکې د وت د
دوه‌چنده کولو دپاره بايد کار ته دوام ورکړي۔ را ياد کړئ چې دوه‌چنده
کوونکے ماشين د ننوت لومړۍ کرښه پاکوي او په بېل وت کښې دوه کرښې ليکي۔
راځئ يوه لارښوونه ور زياته کړو، څو ډاډمن شو چې د پټې سر د جمع ماشين
د وت لومړۍ کرښه لولي۔
\[
\begin{tikzpicture}[->,>=stealth',shorten >=1pt,auto,node distance=2.8cm,
                    semithick]
  \tikzstyle{every state}=[fill=none,draw=black,text=black]
\olpseganchor{OLP-0261-B010}{end}

\olpseganchor{OLP-0261-B011}{start}
  \node[initial,state] (A)              {\LR{$q_0$}};
  \node[state]         (B) [right of=A] {\LR{$q_1$}};
  \node[state]         (C) [right of=B] {\LR{$q_2$}};
  \node[state]         (D) [below left of=C] {\LR{$q_3$}};
  \node[state]         (E) [below left of=D] {\LR{$q_4$}};
\olpseganchor{OLP-0261-B011}{end}

\olpseganchor{OLP-0261-B012}{start}
  \path (A) edge node {\TMtrans{\TMblank}{\TMstroke}{\TMstay}} (B)
            edge [loop above] node {\TMtrans{\TMstroke}{\TMstroke}{\TMright}} (B)
        (B) edge [loop above] node {\TMtrans{\TMstroke}{\TMstroke}{\TMright}} (B)
            edge node {\TMtrans{\TMblank}{\TMblank}{\TMleft}} (C)
        (C) edge node {\TMtrans{\TMstroke}{\TMblank}{\TMleft}} (D)
        (D) edge [loop left] node {\TMtrans{\TMstroke}{\TMstroke}{\TMleft}} (D)
            edge node[left, xshift=-2mm] {\TMtrans{\TMendtape}{\TMendtape}{\TMright}} (E);
\end{tikzpicture}
\]
اوس د ننوت دوه‌چنده کول اسانه دي---يوازې دوه‌چنده کوونکے ماشين له
حالت~\LR{$q_4$} سره نښلوو۔ د دې دپاره د دوه‌چنده کوونکي حالتونه داسې بيا
نومول پکار دي چې د~\LR{$q_0$} پر ځاے له~\LR{$q_4$} څخه پيل شي؛ په دې توګه دوه
لومړني حالتونه نۀ جوړېږي۔ وروستنی ډياګرام بايد د
\olref{fig:combined} په شان وي۔
\begin{figure}
\[
\begin{tikzpicture}[->,>=stealth',shorten >=1pt,auto,node distance=2.8cm,
                    semithick]
  \tikzstyle{every state}=[fill=none,draw=black,text=black]
  \node[initial,state] (A)              {\LR{$q_0$}};
  \node[state]         (B) [right of=A] {\LR{$q_1$}};
  \node[state]         (C) [right of=B] {\LR{$q_2$}};
  \node[state]         (D) [below left of=C] {\LR{$q_3$}};
  \node[state]         (E) [below left of=D] {\LR{$q_4$}};
  \node[state]         (2) [right of=E] {\LR{$q_5$}};
  \node[state]         (3) [right of=2] {\LR{$q_6$}};
  \node[state]         (4) [below of=3] {\LR{$q_7$}};
  \node[state]         (5) [left of=4]  {\LR{$q_8$}};
  \node[state]         (6) [left of=5]  {\LR{$q_9$}};
\olpseganchor{OLP-0261-B012}{end}

\olpseganchor{OLP-0261-B013}{start}
  \path (A) edge node {\TMtrans{\TMblank}{\TMstroke}{\TMstay}} (B)
            edge [loop above] node {\TMtrans{\TMstroke}{\TMstroke}{\TMright}} (B)
        (B) edge [loop above] node {\TMtrans{\TMstroke}{\TMstroke}{\TMright}} (B)
            edge node {\TMtrans{\TMblank}{\TMblank}{\TMleft}} (C)
        (C) edge node {\TMtrans{\TMstroke}{\TMblank}{\TMleft}} (D)
        (D) edge [loop left] node {\TMtrans{\TMstroke}{\TMstroke}{\TMleft}} (D)
            edge node[left, xshift=-2mm] {\TMtrans{\TMendtape}{\TMendtape}{\TMright}} (E)
    (E) edge node {\TMtrans{\TMstroke}{\TMblank}{\TMright}} (2)
    (2) edge [loop above] node {\TMtrans{\TMstroke}{\TMstroke}{\TMright}} (2)
      edge node {\TMtrans{\TMblank}{\TMblank}{\TMright}} (3)
    (3) edge [loop above] node {\TMtrans{\TMstroke}{\TMstroke}{\TMright}} (3)
        edge node {\TMtrans{\TMblank}{\TMstroke}{\TMright}} (4)
    (4) edge [loop below] node {\TMtrans{\TMblank}{\TMstroke}{\TMleft}} (4)
        edge node {\TMtrans{\TMstroke}{\TMstroke}{\TMleft}} (5)
    (5) edge [loop below]  node {\TMtrans{\TMstroke}{\TMstroke}{\TMleft}} (5)
        edge              node {\TMtrans{\TMblank}{\TMblank}{\TMleft}} (6)
    (6) edge [loop below] node {\TMtrans{\TMstroke}{\TMstroke}{\TMleft}} (6)
        edge              node {\TMtrans{\TMblank}{\TMblank}{\TMright}} (E);
\end{tikzpicture}
\]
\caption{د جمع کوونکي او دوه‌چنده کوونکي ماشينونو يوځاے کول}
\ollabel{fig:combined}
\end{figure}
\end{ex}
\olpseganchor{OLP-0261-B013}{end}

\olpseganchor{OLP-0261-B014}{start}
\begin{prop}
    که~\LR{$M$} او~\LR{$M'$} منضبط وي او په ترتيب سره تابعې \LR{$f\colon
    \Nat^k \to \Nat$} او \LR{$f'\colon \Nat \to \Nat$} محاسبه کوي، نو
    \LR{$M \frown M'$} منضبط دے او~\LR{$\comp{f}{f'}$} محاسبه کوي۔
\end{prop}
\olpseganchor{OLP-0261-B014}{end}

\olpseganchor{OLP-0261-B015}{start}
\begin{proof}
    څرنګه چې~\LR{$M$} منضبط دے، کله چې د وت
    ~\LR{$f(n_1,\dots,n_k) = m$} سره درېږي، سر مربع~\LR{$1$} لولي۔ که اوس د
    ~\LR{$M'$} لومړني حالت ته ننوځو، ماشين به د وت~\LR{$f'(m)$} سره ودرېږي او
    بيا به هم مربع~\LR{$1$} لولي۔ د~\olref[dis]{defn:disciplined} نور
    شرطونه هم پوره کېږي۔
\end{proof}
\olpseganchor{OLP-0261-B015}{end}

\olpseganchor{OLP-0261-B016}{start}
\begin{prob}
    د~\cref{tur:mac:dis:prob:disc-succ} ماشين~\LR{$M$} واخلئ،
    \LR{$M \frown M$} جوړ کړئ، او په دې توګه داسې منضبط ټيورينګ ماشين
    ورکړئ چې \LR{$f(x) = x+2$} محاسبه کوي۔
\end{prob}
\olpseganchor{OLP-0261-B016}{end}

\olpseganchor{OLP-0262-B004}{start}
\olfileid{tur}{mac}{var}
\olsection{د ټيورينګ ماشينونو بڼې}
\olpseganchor{OLP-0262-B004}{end}

\olpseganchor{OLP-0262-B005}{start}
په حقيقت کښې د ټيورينګ ماشينونو د تعريفولو ډېرې شونې لارې شته، او
زموږ تعريف يوازې يوه لار ده۔ زموږ تعريف له ځينو اړخونو څخه تر نورو
ازاد دے۔ موږ هره متناهي الفبا منو؛ يو لا محدود تعريف ښايي يوازې د
پټې دوې نښې، \LR{$\TMstroke$} او~\LR{$\TMblank$}، ومني۔ موږ ماشين ته اجازه
ورکوو چې يوه نښه پر پټه وليکي او په هماغه وخت کښې وخوځېږي؛ نور
تعريفونه يا ليکل او يا خوځېدل مني۔ موږ د سر له خوځولو پرته د نښې د
ليکلو شونتيا منو؛ نور تعريفونه د~\LR{$\TMstay$} ``لارښوونه'' نۀ لري۔ زموږ
تعريف له نورو اړخونو څخه زيات محدود دے۔ موږ فرض کړې چې پټه يوازې په
يوه لوري کښې نامتناهي ده؛ نور تعريفونه پټه کيڼ او ښي دواړو لوريو ته
نامتناهي پرېږدي۔ ان هر شمېر بېلې پټې، يا د مربعونو نامتناهي جال هم
منل کېدے شي۔ موږ د ټيورينګ ماشين د لارښوونو سټ د حالت بدلون تابع په
وسيله څرګندوو؛ نور تعريفونه د حالت بدلون اړيکه کاروي چې ماشين ته په
هر ټاکلي حالت کښې تر يوې زياتې ممکنې لارښوونې ورکوي۔
\olpseganchor{OLP-0262-B005}{end}

\olpseganchor{OLP-0262-B006}{start}
د تعريف دا وروستنۍ نرمونه ځانګړې په زړه پورې ده۔ زموږ په تعريف کښې
چې ماشين په حالت~\LR{$q$} کښې نښه~\LR{$\sigma$} لولي، \LR{$\delta(q, \sigma)$}
ټاکي چې نوې نښه، حالت او د پټې د سر ځاے څه دے۔ خو که د لارښوونو سټ
د اوسنيو حالت-نښه جوړو \LR{$\tuple{q,
  \sigma}$} او نويو حالت-نښه-لوري درې‌يزونو \LR{$\tuple{q', \sigma',
  D}$} تر منځ يوه اړيکه وګڼو، نو د ټيورينګ ماشين عمل ښايي په يوازيني
ډول ټاکلے نۀ وي---د لارښوونو اړيکه هم \LR{$\tuple{q,
  \sigma, q', \sigma', D}$} او هم \LR{$\tuple{q, \sigma, q'', \sigma'', D'}$}
لرلے شي۔ په دې صورت کښې موږ يو \emph{ناتعييني} ټيورينګ ماشين لرو۔
دا ماشينونه د محاسبوي پېچلتيا په تيورۍ کښې مهم رول لري۔
\olpseganchor{OLP-0262-B006}{end}

\olpseganchor{OLP-0262-B007}{start}
د ټيورينګ ماشين د درېدو دپاره هم بېلابېل دودونه شته: موږ وايو هغه
وخت درېږي چې د حالت بدلون تابع تعريف شوې نۀ وي؛ نور تعريفونه غواړي
چې ماشين په يوه ځانګړي ټاکلي درېدني حالت کښې وي۔ په~\olref[hal]{sec}
کښې مو څرګنده کړې چې د ټاکلي درېدني حالت غوښتل داسې بنديز نۀ دے چې
د ټيورينګ ماشينونو محاسبه کېدونکي څيزونه بدل کړي۔ څرنګه چې زموږ د
ټيورينګ ماشينونو پټې يوازې په يوه لوري کښې نامتناهي دي، داسې حالتونه
شته چې ټيورينګ ماشين يوه لارښوونه په سمه توګه نۀ شي عملي کولے: کله
چې تر ټولو کيڼ مربع لولي او بايد کيڼ لوري ته لاړ شي۔ زموږ د تعريف له
مخې ماشين د ``لوېدو'' پر ځاے هماغلته پاتې کېږي، خو موږ داسې تعريفولے
شو چې په دې حالت کښې ودرېږي۔ دا تعريف هم معادل دے: د هغه ټيورينګ
ماشين چلند چې له مربع~\LR{$0$} څخه کيڼ لوري ته د تلو پر هڅه درېږي، د هر
حالت بدلون \LR{$\delta(q,\TMendtape) =
\tuple{q',\sigma,\TMleft}$} په ړنګولو تقليدولے شو---بيا ماشين پر
~\LR{$\TMendtape$} د کيڼ لوري ته د تلو د هڅې پر ځاے درېږي۔\footnote{دا
\emph{بالکل} کار نۀ کوي، ځکه هېڅ څيز موږ د مربع~\LR{$0$} پرته پر نورو
مربعونو د~\LR{$\TMendtape$} له ليکلو او لوستلو څخه نۀ منع کوي
(\olref[una]{ex:mover} وګورئ)۔ د دې د مخنيوي دپاره د دې موخې لپاره
يوه دويمه نښه~\LR{$\TMendtape'$} ور زياتولے شو۔}
\olpseganchor{OLP-0262-B007}{end}

\olpseganchor{OLP-0262-B008}{start}
د عددونو د څرګندولو بېلابېلې لارې هم شته (او په دې توګه د ټيورينګ
ماشين محاسبه شوې ننوت-وت تابع هم بدلېږي): موږ يوګونې څرګندونه کاروو،
خو دوه‌ګونې څرګندونه هم کارېدے شي۔ د دې دپاره پر~\LR{$\TMblank$} او
~\LR{$\TMendtape$} سربېره دوو نښو ته اړتيا ده۔
\olpseganchor{OLP-0262-B008}{end}

\olpseganchor{OLP-0262-B009}{start}
اوس يوه په زړه پورې خبره دا ده: له دې بڼو څخه يوه هم دا نۀ بدلوي چې
کومې تابعې د ټيورينګ په معنا محاسبه کېدونکې دي۔ \emph{که يوه تابع د
يو تعريف له مخې د ټيورينګ په معنا محاسبه کېدونکې وي، نو د ټولو
تعريفونو له مخې د ټيورينګ په معنا محاسبه کېدونکې ده۔}
\olpseganchor{OLP-0262-B009}{end}

\olpseganchor{OLP-0262-B010}{start}
موږ د دې د پخلي جزئياتو ته نۀ ځو۔ يوازې يوه بېلګه وګورئ: کيڼ او ښي
دواړو لوريو ته د نامتناهي پټې، يا د څو پټو، په منلو هېڅ اضافي
محاسبوي ځواک نۀ لاس ته راوړو۔ په لنډه توګه دليل دا دے چې د يوې
يو-لوري نامتناهي پټې لرونکے ټيورينګ ماشين د څو پټو يا دوه-لوري
نامتناهي پټو تقليد کولے شي۔ مثلاً، د اضافي حالتونو او لارښوونو په
کارولو د څو پټو يا دوه-لوري نامتناهي پټې لرونکي ماشين پروګرام په
داسې پروګرام ``ژباړلے'' شو چې يوه يو-لوري نامتناهي پټه لري۔ ژباړل
شوے ماشين جفت مربعونه د پټې~\LR{$1$} مربعونو دپاره (يا د دوه-لوري
نامتناهي پټې د ``مثبتو'' مربعونو دپاره)، او طاق مربعونه د پټې~\LR{$2$}
مربعونو دپاره (يا د ``منفي'' مربعونو دپاره) کارولے شي۔
\olpseganchor{OLP-0262-B010}{end}

\olpseganchor{OLP-0263-B004}{start}
\olfileid{tur}{mac}{ctt}
\olsection{د چرچ--ټيورينګ اصل}
\olpseganchor{OLP-0263-B004}{end}

\olpseganchor{OLP-0263-B005}{start}
ټيورينګ ماشينونه بايد د اغېزمنې کړنلارې د مفهوم کره بديل وي۔ ټيورينګ
فکر کاوه چې هر څوک چې هم د اغېزمنې کړنلارې او هم د ټيورينګ ماشين په
مفهوم پوه شي، دا شهود به ولري چې هر هغه کار چې د اغېزمنې کړنلارې په
وسيله کېدے شي، د ټيورينګ ماشين په وسيله هم کېدے شي۔ دا ادعا له دې
خبرې ملاتړ اخلي چې د اغېزمنې کړنلارې د مفهوم ټول نور وړانديز شوي کره
بديلونه په امتدادي ډول د ټيورينګ ماشين له مفهوم سره معادل راووځي---
يعنې د تابعو کټ مټ هماغه سټ محاسبه کولے شي۔ دې ادعا ته
\emph{د چرچ--ټيورينګ اصل} ويل کېږي۔
\olpseganchor{OLP-0263-B005}{end}

\olpseganchor{OLP-0263-B006}{start}
\begin{defn}[د چرچ--ټيورينګ اصل]
\emph{د چرچ--ټيورينګ اصل} وايي چې هر هغه څيز چې د اغېزمنې کړنلارې
په وسيله محاسبه کېدے شي، د ټيورينګ په معنا محاسبه کېدونکے دے۔
\end{defn}
\olpseganchor{OLP-0263-B006}{end}

\olpseganchor{OLP-0263-B007}{start}
د چرچ--ټيورينګ اصل په دوو لارو کارول کېږي۔ لومړے ډول کارول ئې د
سستۍ يوه پلمه ده۔ فرض کړئ د يوه څيز د محاسبې د اغېزمنې کړنلارې
توضيح لرو، مثلاً په ``کاذب کوډ'' کښې۔ بيا د چرچ--ټيورينګ اصل کارولے
شو، څو دا ادعا توجيه کړو چې هماغه تابع يو ټيورينګ ماشين محاسبه کوي،
ان که موږ هغه ماشين په عمل کښې نۀ وي جوړ کړے۔
\olpseganchor{OLP-0263-B007}{end}

\olpseganchor{OLP-0263-B008}{start}
د چرچ--ټيورينګ اصل بل استعمال له فلسفي پلوه لا په زړه پورې دے۔
ثابتېدے شي چې داسې تابعې شته چې ټيورينګ ماشينونه ئې نۀ شي محاسبه
کولے۔ له دې څخه د چرچ--ټيورينګ اصل په کارولو پايله اخيستلے شو چې
هغه تابع په هېڅ ډول کړنلارې سره په اغېزمنه توګه نۀ شي محاسبه کېدے۔
ځکه که داسې کړنلاره موجوده وے، د چرچ--ټيورينګ اصل له مخې به ورته يو
ټيورينګ ماشين هم موجود وے۔ نو که ثابت کړے شو چې هېڅ ټيورينګ ماشين
هغه نۀ شي محاسبه کولے، اغېزمنه کړنلاره هم ورته نشته۔ په ځانګړې توګه،
د چرچ--ټيورينګ اصل د دې ادعا دپاره کارول کېږي چې د درېدو په نوم
مسئله نۀ يوازې د ټيورينګ ماشينونو په وسيله نۀ شي حل کېدے، بلکې په
هېڅ ډول په اغېزمنه توګه نۀ شي حل کېدے۔
\olpseganchor{OLP-0263-B008}{end}

\olpseganchor{OLP-0264-B004}{start}
\olchapter{tur}{und}{ناپرېکړتيا}
\olpseganchor{OLP-0264-B004}{end}

\olpseganchor{OLP-0265-B004}{start}
\olfileid{tur}{und}{int}
\olsection{پېژندنه}
\olpseganchor{OLP-0265-B004}{end}

\olpseganchor{OLP-0265-B005}{start}
ښايي دا څرګنده وبرېښي چې هره تابع، ان هره حسابي تابع، محاسبه کېدونکې
نۀ شي کېدے۔ تابعې دومره ډېرې دي او چلند ئې دومره پېچلے دے۔ مثلاً هغه
تابعې چې د راډيو‌اکټيفو ذرو له تجزيې، يا له بل ګډوډ يا تصادفي چلند
څخه تعريفېږي۔ فرض کړئ له يوه ټاکلي وخت څخه يوه-ثانيه‌يي واټنونه شمېرو،
او تابع~\LR{$f(n)$} د کاينات د هغو ذرو د شمېر په توګه تعريفوو چې له هغې
لومړنۍ شېبې وروسته په~\LR{$n$}-م يوه-ثانيه‌يي واټن کښې تجزيه کېږي۔ دا د
داسې تابعې غوندې ښکاري چې هېڅکله ئې د محاسبې هيله نۀ شو کولے۔
\olpseganchor{OLP-0265-B005}{end}

\olpseganchor{OLP-0265-B006}{start}
خو دا چې د داسې تابعو د محاسبې لار تصور نۀ شو کولے يوه خبره ده، او
دا چې په رښتيا ثابته کړو چې دوي نۀ محاسبه کېږي، بله خبره ده۔ په اصل
کښې ان هغه تابعې چې بېخي ناهيلې کوونکې پېچلې ښکاري، په مجردې معنا
محاسبه کېدونکې کېدے شي۔ مثلاً فرض کړئ کاينات د وخت له پلوه متناهي دے
---لکه ځينې کيهان‌پوهنيزې تيورۍ چې وړاندوينه کوي، يوه ورځ به په ډېر
لرې راتلونکي کښې کاينات په يوه ټکي کښې راټول شي۔ بيا له هغې لومړنۍ
شېبې څخه يوازې د متناهي (خو بېخي لوے) شمېر ثانيو دپاره~\LR{$f(n)$} تعريف
شوې وي۔ او هره تابع چې يوازې د متناهي شمېر ننوتونو دپاره تعريف شوې
وي محاسبه کېدونکې ده: وتونه په يوه لوے جدول کښې لست کولے، يا ئې د
ټيورينګ ماشين د حالت بدلون په يوه ډېر لوے ډياګرام کښې کوډولے شو۔
\olpseganchor{OLP-0265-B006}{end}

\olpseganchor{OLP-0265-B007}{start}
موږ ډېر ځله د هغو تابعو ځانګړو حالتونو ته پام کوو چې ارزښتونه ئې د
هو/نه پوښتنو ځوابونه ورکوي۔ مثلاً پوښتنه ``ايا~\LR{$n$} لومړنی عدد دے؟''
له دې تابع سره تړلې ده:
\[
\fn{isprime}(n) = \begin{cases}
  1 & \text{\RL{که \LR{$n$} لومړنی وي}}\\
  0 & \text{\RL{که نۀ وي۔}}
  \end{cases}
\]
وايو يوه هو/نه پوښتنه په \emph{اغېزمنه توګه پرېکړه کېدے} شي، که ور
سره تړلې د~\LR{$1/0$} ارزښت لرونکې تابع په اغېزمنه توګه محاسبه کېدونکې وي۔
\olpseganchor{OLP-0265-B007}{end}

\olpseganchor{OLP-0265-B008}{start}
د دې د رياضياتي ثبوت دپاره چې داسې تابعې شته چې په اغېزمنه توګه نۀ
شي محاسبه کېدے، يا داسې مسئلې شته چې په اغېزمنه توګه نۀ شي پرېکړه
کېدے، اړينه ده چې د محاسبې يو ټاکلی مدل وټاکو، او وښيو چې داسې تابعې
شته چې دا ئې نۀ شي محاسبه کولے يا داسې مسئلې شته چې دا ئې نۀ شي
پرېکړه کولے۔ مثلاً ښودلے شو چې هره تابع د ټيورينګ ماشينونو په وسيله
نۀ شي محاسبه کېدے او هره مسئله د ټيورينګ ماشينونو په وسيله نۀ شي
پرېکړه کېدے۔ بيا د چرچ--ټيورينګ اصل کارولے شو، څو پايله واخلو چې نۀ
يوازې ټيورينګ ماشينونه د هرې تابع د محاسبې دپاره پوره ځواکمن نۀ دي،
بلکې هېڅ اغېزمنه کړنلاره دا کار نۀ شي کولے۔
\olpseganchor{OLP-0265-B008}{end}

\olpseganchor{OLP-0265-B009}{start}
د داسې منفي پايلو د ثبوت کلي دا حقيقت دے چې موږ پخپله ټيورينګ
ماشينونو ته عددونه ټاکلے شو۔ تر ټولو اسانه لار ئې دا ده چې دوي وشمېرو:
ښايي د ټيورينګ ماشينونو او د هغو د پروګرامونو د ليکلو يوه ټاکلې لار
وټاکو، او بيا ئې په منظم ډول لست کړو۔ کله چې ووينو دا کار کېدے شي،
نو د ټيورينګ-نامحاسبه کېدونکو تابعو شتون د شمېر کچې له ساده دلايلو
څخه راځي: له~\LR{$\Nat$} څخه~\LR{$\Nat$} ته د تابعو سټ (په اصل کښې ان يوازې
له~\LR{$\Nat$} څخه~\LR{$\{0, 1\}$} ته) د شمېر نۀ وړ دے، خو څرنګه چې ټول
ټيورينګ ماشينونه شمېرلے شو، د ټيورينګ-محاسبه کېدونکو تابعو سټ يوازې
شمېرېدونکے لامتناهي دے۔
\olpseganchor{OLP-0265-B009}{end}

\olpseganchor{OLP-0265-B010}{start}
موږ داسې \emph{ټاکلې} تابعې او مسئلې هم تعريفولے شو چې په ترتيب سره
ئې نۀ محاسبه کېدل او ناپرېکړتيا ثابتولے شو۔ يوه داسې مسئله د
\emph{درېدو مسئلې} په نوم ده۔ د ټيورينګ ماشين لارښوونې په لست کولو
هغه په متناهي ډول بيانېدے شي۔ د ټيورينګ ماشين دا ډول بيان، يعنې د
ټيورينګ ماشين پروګرام، البته بل ټيورينګ ماشين ته د ننوت په توګه
کارېدے شي۔ نو هغه ټيورينګ ماشينونه څېړلے شو چې د نورو ټيورينګ
ماشينونو په اړه پوښتنې پرېکړه کوي۔ يوه په زړه پورې پوښتنه دا ده:
``ايا ورکړل شوے ټيورينګ ماشين به پر ننوت~\LR{$n$} له پيل وروسته بالاخره
ودرېږي؟'' ښه به وے که داسې ټيورينګ ماشين موجود وے چې دا پوښتنه ئې
پرېکړه کولے: دا د کيفيت-کنټرول يو ټيورينګ ماشين وګڼئ چې ډاډ ورکوي
ټيورينګ ماشينونه په نامتناهي کړيو او داسې نورو کښې نۀ بندېږي۔ په زړه
پورې حقيقت، چې ټيورينګ ثابت کړ، دا دے چې داسې ټيورينګ ماشين نشي
موجودېدے۔ هېڅ يو ټيورينګ ماشين نشته چې د يوه ټيورينګ ماشين~\LR{$M$} له
بيان او يوه عدد~\LR{$n$} څخه جوړ ننوت باندې پيل شي او تل له~\LR{$1$} يا~\LR{$0$}
وت سره د دې له مخې ودرېږي چې~\LR{$M$} به پر ننوت~\LR{$n$} درېدلے و که نه۔
\olpseganchor{OLP-0265-B010}{end}

\olpseganchor{OLP-0265-B011}{start}
کله چې د ټاکلو ناپرېکړه کېدونکو مسئلو بېلګې ولرو، د نورو مسئلو د
ناپرېکړتيا د ښودلو دپاره ئې هم کارولے شو۔ مثلاً يوه مشهوره
ناپرېکړه کېدونکې مسئله دا پوښتنه ده: ``ايا د لومړۍ درجې
فارمول~\LR{$!A$} معتبر دے؟'' داسې ټيورينګ ماشين نشته چې د لومړۍ درجې
فارمول~\LR{$!A$} د ننوت په توګه ورکړل شي او هرومرو له~\LR{$1$} يا~\LR{$0$} وت
سره د دې له مخې ودرېږي چې~\LR{$!A$} معتبر دے که نه۔ په تاريخي ډول د دې
مسئلې د اغېزمن حل د کړنلارې د موندلو پوښتنې ته يوازې ``د'' پرېکړې
مسئله ويل کېده؛ نو وايو چې د پرېکړې مسئله ناحل کېدونکې ده۔ ټيورينګ
او چرچ دا پايله نږدې په يوه وخت کښې په خپلواکه توګه ثابته کړه، ځکه
نو دې ته د چرچ--ټيورينګ قضيه هم ويل کېږي۔
\olpseganchor{OLP-0265-B011}{end}

\olpseganchor{OLP-0266-B004}{start}
\olfileid{tur}{und}{enu}
\olsection{د ټيورينګ ماشينونو شمېرل}
\olpseganchor{OLP-0266-B004}{end}

\olpseganchor{OLP-0266-B005}{start}
\begin{explain}
موږ ښودلے شو چې د ټولو ټيورينګ ماشينونو سټ د شمېر وړ دے۔ دا
له دې حقيقت څخه راځي چې هر ټيورينګ ماشين په متناهي ډول بيانېدے شي۔
د حالتونو سټ او د پټې لغت‌زېرمه متناهي سټونه دي۔ د حالت بدلون تابع
له~\LR{$Q \times \Sigma$} څخه~\LR{$Q \times \Sigma \times \{\TMleft,
\TMright, \TMstay\}$} ته يوه قسمي تابع ده، نو دا هم د هغو متناهي
شمېر آرګوماني جوړو د ارزښتونو په لست کولو ټاکل کېدے شي چې تابع پرې
تعريف شوې ده۔
\olpseganchor{OLP-0266-B005}{end}

\olpseganchor{OLP-0266-B006}{start}
تر دې ځايه دا خبره سمه ده، خو يو نازک توپير شته۔ د ټيورينګ ماشينونو
تعريف پر دې هېڅ بنديز نۀ لګاوه چې د حالتونو د سټ او د پټې د الفبا
غړي څه شيان کېدے شي۔ نو مثلاً د هر حقيقي عدد دپاره له تخنيکي
پلوه يو ټيورينګ ماشين شته چې هماغه عدد د حالت په توګه کاروي۔ خو د
ټيورينګ ماشين \emph{چلند} له دې څخه خپلواک دے چې کوم څيزونه د حالتونو
او لغت‌زېرمه په توګه کارول کېږي۔ په~\olref{fig:variants} کښې دا دوه
ټيورينګ ماشينونه وګورئ۔
\begin{figure}
\begin{center}
\begin{tikzpicture}[->,>=stealth',shorten >=1pt,auto,node distance=2.8cm,
                    semithick]
  \tikzstyle{every state}=[fill=none,draw=black,text=black]
\olpseganchor{OLP-0266-B006}{end}

\olpseganchor{OLP-0266-B007}{start}
  \node[initial,state]         (A)              {\LR{$q_0$}};
  \node[state]         (B) [right of=A] {\LR{$q_1$}};
\olpseganchor{OLP-0266-B007}{end}

\olpseganchor{OLP-0266-B008}{start}
  \path (A) edge [bend left] node {\TMtrans{\TMstroke}{\TMstroke}{\TMright}} (B)
        (B) edge [loop above] node {\TMtrans{\TMblank}{\TMblank}{\TMright}} (B)
            edge [bend left] node {\TMtrans{\TMstroke}{\TMstroke}{\TMright}} (A);
\end{tikzpicture}\\
\begin{tikzpicture}[->,>=stealth',shorten >=1pt,auto,node distance=2.8cm,
  semithick]
\tikzstyle{every state}=[fill=none,draw=black,text=black]
\olpseganchor{OLP-0266-B008}{end}

\olpseganchor{OLP-0266-B009}{start}
\node[initial,state]         (A)              {\LR{$s$}};
\node[state]         (B) [right of=A] {\LR{$h$}};
\olpseganchor{OLP-0266-B009}{end}

\olpseganchor{OLP-0266-B010}{start}
\path (A) edge [bend left] node {\TMtrans{A}{A}{\TMright}} (B)
(B) edge [loop above] node {\TMtrans{\TMblank}{\TMblank}{\TMright}} (B)
edge [bend left] node {\TMtrans{A}{A}{\TMright}} (A);
\end{tikzpicture}
\end{center}
\caption{د \emph{جفت} ماشين بڼې}
\ollabel{fig:variants}
\end{figure}
دا دوه ډياګرامونه له دوو ماشينونو سره سمون لري: يو~\LR{$M$} چې د پټې
الفبا ئې~\LR{$\Sigma = \{\TMendtape,\TMblank,\TMstroke\}$} او د حالتونو
سټ ئې~\LR{$\{q_0,q_1\}$} دے، او بل~\LR{$M'$} چې الفبا ئې~\LR{$\Sigma' =
\{\TMendtape,\TMblank,A\}$} او حالتونه ئې~\LR{$\{s,h\}$} دي۔ خو نورې
لارښوونې ئې يو شان دي: \LR{$M$} به د~\LR{$n$} دانو~\LR{$\TMstroke$} پر لړۍ هله او
يوازې هله ودرېږي چې~\LR{$n$} جفت وي، او~\LR{$M'$} به د~\LR{$n$} دانو~\LR{$A$} پر لړۍ
هله او يوازې هله ودرېږي چې~\LR{$n$} جفت وي۔ موږ يوازې~\LR{$\TMstroke$} په~\LR{$A$}،
\LR{$q_0$} په~\LR{$s$}، او~\LR{$q_1$} په~\LR{$h$} بدل نومولي دي۔ دا بېلګه عموميت مومي:
موږ ټيورينګ ماشينونه تر هغه يو شان ګڼلے شو چې يو د نښو او حالتونو د
داسې نوم‌بدلون له لارې له بل څخه تر لاسه کېږي۔ په اصل کښې، موږ د
ټيورينګ ماشين نښې او حالتونه په ساده توګه مثبت صحيح اعداد ګڼلے شو:
د~\LR{$\sigma_0$} پر ځاے~\LR{$1$}، د~\LR{$\sigma_1$} پر ځاے~\LR{$2$}، او داسې نور
وګڼئ؛ \LR{$\TMendtape$}~\LR{$1$} دے، \LR{$\TMblank$}~\LR{$2$} دے، او داسې نور۔ په دې
ډول د \emph{جفت} ماشين په~\olref{fig:standard-even} کښې انځور شوے
ماشين ګرځي۔
\begin{figure}
\[\begin{tikzpicture}[->,>=stealth',shorten >=1pt,auto,node distance=2.8cm,
  semithick]
\tikzstyle{every state}=[fill=none,draw=black,text=black]
\olpseganchor{OLP-0266-B010}{end}

\olpseganchor{OLP-0266-B011}{start}
\node[initial,state]         (A)              {\LR{$1$}};
\node[state]         (B) [right of=A] {\LR{$2$}};
\olpseganchor{OLP-0266-B011}{end}

\olpseganchor{OLP-0266-B012}{start}
\path (A) edge [bend left] node {\TMtrans{3}{3}{\TMright}} (B)
(B) edge [loop above] node {\TMtrans{2}{2}{\TMright}} (B)
edge [bend left] node {\TMtrans{3}{3}{\TMright}} (A);
\end{tikzpicture}
\]
\caption{يو معياري \emph{جفت} ماشين}
\ollabel{fig:standard-even}
\end{figure}
هغه ټيورينګ ماشين چې حالتونه او نښې ئې مثبت صحيح اعداد وي ښايي
\emph{معياري} ماشين وبولو، او له دې وروسته يوازې معياري ماشينونه
وڅېړو۔\footnote{د ``معياري ماشين'' اصطلاح معياري نۀ ده۔}
\olpseganchor{OLP-0266-B012}{end}

\olpseganchor{OLP-0266-B013}{start}
موږ غوښتل وښيو چې د ټيورينګ ماشينونو سټ د شمېر وړ دے، او پورتنيو
خبرو ته په پام کافي ده چې وښيو د معياري ټيورينګ ماشينونو سټ
د شمېر وړ دے۔ فرض کړئ يو معياري ټيورينګ ماشين~\LR{$M = \tuple{Q,
\Sigma, q_0, \delta}$} راکړل شوے دے۔ هغه د مثبتو صحيح اعدادو په يوه
متناهي مزي څنګه بيانولے شو؟ لومړے به د حالتونو شمېر، پخپله حالتونه،
د نښو شمېر، پخپله نښې، او پيلنی حالت لست کړو۔ (په ياد ولرئ، دا ټول
مثبت صحيح اعداد دي، ځکه~\LR{$M$} معياري ماشين دے۔) د~\LR{$\delta$} په اړه څه؟
د ممکنو آرګومانونو سټ، يعنې د~\LR{$\tuple{q,\sigma}$} جوړې، متناهي دے،
ځکه~\LR{$Q$} او~\LR{$\Sigma$} متناهي دي۔ نو په~\LR{$\delta$} کښې معلومات يوازې د
ټولو هغو \LR{$5$}-تاييونو~\LR{$\tuple{q, \sigma, q', \sigma', d}$} متناهي
لست دے چې~\LR{$\delta(q,\sigma) = \tuple{q', \sigma', D}$} وي، او~\LR{$d$}
هغه عدد دے چې لوری~\LR{$D$} کوډوي (مثلاً د~\LR{$\TMleft$} دپاره~\LR{$1$}،
د~\LR{$\TMright$} دپاره~\LR{$2$}، او د~\LR{$\TMstay$} دپاره~\LR{$3$})۔
\olpseganchor{OLP-0266-B013}{end}

\olpseganchor{OLP-0266-B014}{start}
په دې ډول هر معياري ټيورينګ ماشين د مثبتو صحيح اعدادو په يوه متناهي
لست، يعنې د~\LR{$s_M \in (\PosInt)^*$} په لړۍ، بيانېدے شي۔ مثلاً معياري
\emph{جفت} ماشين په دې لړۍ کوډ کېږي:
\[
2, \underbrace{1, 2}_Q, 3, \overbrace{1, 2, 3}^\Sigma, 1, \underbrace{1, 3, 2, 3, 2}_{\delta(1,3) = \tuple{2,3,R}},
\overbrace{2, 2, 2, 2, 2}^{\delta(2,2) = \tuple{2,2,R}},
\underbrace{2, 3, 1, 3, 2}_{\delta(2,3) = \tuple{1,3,R}}.
\]
\end{explain}
\olpseganchor{OLP-0266-B014}{end}

\olpseganchor{OLP-0266-B015}{start}
\begin{thm}
له~\LR{$\Nat$} څخه~\LR{$\Nat$} ته داسې تابعې شته چې د ټيورينګ په وسيله
محاسبه کېدونکې نۀ دي۔
\end{thm}
\olpseganchor{OLP-0266-B015}{end}

\olpseganchor{OLP-0266-B016}{start}
\begin{proof}
موږ پوهېږو چې د مثبتو صحيح اعدادو د متناهي لړيو سټ~\LR{$(\PosInt)^*$}
د شمېر وړ دے (\cref{sfr:siz:zigzag:prob:posint-star})۔ له دې څخه
راځي چې د معياري ټيورينګ ماشينونو د بيانونو سټ، د~\LR{$(\PosInt)^*$} د
يوه فرعي سټ په توګه، پخپله هم شمېرل کېدونکے دے۔ له~\LR{$\Nat$} څخه~\LR{$\Nat$}
ته هره د ټيورينګ په وسيله محاسبه کېدونکې تابع د کوم (په اصل کښې د
ډېرو) ټيورينګ ماشين له خوا محاسبه کېږي۔ د هغې د حالتونو او نښو نومونه
په مثبتو صحيح اعدادو بدلولو سره (په ځانګړي ډول~\LR{$\TMendtape$} په~\LR{$1$}،
\LR{$\TMblank$} په~\LR{$2$}، او~\LR{$\TMstroke$} په~\LR{$3$}) وينو چې هره د ټيورينګ په
وسيله محاسبه کېدونکې تابع د يوه معياري ټيورينګ ماشين له خوا محاسبه
کېږي۔ معنا دا چې له~\LR{$\Nat$} څخه~\LR{$\Nat$} ته د ټولو د ټيورينګ په وسيله
محاسبه کېدونکو تابعو سټ هم شمېرل کېدونکے دے۔
\olpseganchor{OLP-0266-B016}{end}

\olpseganchor{OLP-0266-B017}{start}
له بلې خوا، له~\LR{$\Nat$} څخه~\LR{$\Nat$} ته د ټولو تابعو سټ د شمېر وړ
نۀ دے (\cref{sfr:siz:red:prob:nat-nat})۔ که ټولې تابعې د کوم ټيورينګ
ماشين په وسيله محاسبه کېدے، نو د هغو ټيورينګ ماشينونو د ټولو بيانونو
په لست کولو مو د ټولو تابعو سټ شمېرلے شو چې دا تابعې محاسبه کوي۔ نو
ځينې تابعې شته چې د ټيورينګ په وسيله محاسبه کېدونکې نۀ دي۔
\end{proof}
\olpseganchor{OLP-0266-B017}{end}

\olpseganchor{OLP-0266-B018}{start}
\begin{prob}
  ايا د ټيورينګ ماشينونو د بيانولو داسې لار درته ذهن ته راځي چې د
  حالتونو او د الفبا د نښو څرګند لست کولو ته اړتيا ونۀ لري؟ تاسو د
  ``معياري'' ماشين خپل مفهوم تعريفولے شئ، خو د دې په اړه څه ووايئ چې
  ولې هر ټيورينګ ماشين ستاسو په نوې معنا د يوه ``معياري'' ماشين له
  خوا تمثيلېدے شي۔
  څرګنده کړئ چې ستاسو په بيان کښې هره لارښوونه څنګه ټاکل کېږي او
  ماشين له بيا نومولو وروسته هم پر هماغه ننوت هماغه کړنه کوي۔
\end{prob}
\olpseganchor{OLP-0266-B018}{end}

\olpseganchor{OLP-0267-B004}{start}
\olfileid{tur}{und}{uni}
\olsection{نړيوال ټيورينګ ماشينونه}
\olpseganchor{OLP-0267-B004}{end}

\olpseganchor{OLP-0267-B005}{start}
په~\olref[enu]{sec} کښې مو وڅېړل چې هر ټيورينګ ماشين د صحيح اعدادو
په يوه متناهي لړۍ بيانېدے شي۔ دا لړۍ د ټيورينګ ماشين حالتونه، الفبا،
پيلنی حالت، او لارښوونې کوډوي۔ موږ دا هم په ګوته کړه چې د دې ټولو
بيانونو سټ د شمېر وړ دے۔ څرنګه چې د داسې بيانونو سټ
شمېرېدونکے لامتناهي دے، معنا ئې دا ده چې له~\LR{$\Nat$} څخه دې بيانونو ته يوه
پر هدف سټ بشپړه تابع شته۔ داسې پر هدف سټ بشپړه تابع مثلاً د کانتور د
زېګ‌زېګ ميتود په کارولو تر لاسه کېدے شي۔ دا موږ ته د ټيورينګ
ماشينونو د ټولو (بيانونو) د شمېرلو لار راکوي۔ که يوه داسې شمېره
ثابته کړو، نو اوس د~\LR{$1$}-م، \LR{$2$}-م، \dots، او~\LR{$e$}-م ټيورينګ ماشين په
اړه خبرې معنا لري۔ دې عددونو ته \emph{شاخصونه} ويل کېږي۔
\olpseganchor{OLP-0267-B005}{end}

\olpseganchor{OLP-0267-B006}{start}
\begin{defn}
که~\LR{$M$} (زموږ په ثابته شوې شمېره کښې) \LR{$e$}-م ټيورينګ ماشين وي، وايو
چې~\LR{$e$} د~\LR{$M$} يو \emph{شاخص} دے۔ د~\LR{$e$}-م ټيورينګ ماشين دپاره~\LR{$M_e$}
ليکو۔
\end{defn}
\olpseganchor{OLP-0267-B006}{end}

\olpseganchor{OLP-0267-B007}{start}
يو ماشين ښايي له يوه څخه ډېر شاخصونه ولري؛ مثلاً د~\LR{$M$} دوه بيانونه
ښايي د هغو د لارښوونو د لست کولو په ترتيب کښې توپير ولري، او دا
بېلابېل بيانونه به بېلابېل شاخصونه ولري۔
\olpseganchor{OLP-0267-B007}{end}

\olpseganchor{OLP-0267-B008}{start}
مهمه دا ده چې د ټيورينګ ماشين د بيانونو شمېره داسې ورکول کېدے شي
چې موږ د~\LR{$M$} بيان د هغه له شاخص څخه، او د يوه ماشين~\LR{$M$} يو شاخص د
هغه له بيان څخه په اغېزمنه توګه محاسبه کړے شو۔ بيا د چرچ--ټيورينګ د
اصل له مخې داسې ټيورينګ ماشين موندل کېدے شي چې د شاخص~\LR{$e$} لرونکي
ټيورينګ ماشين بيان بېرته تر لاسه کړي او اړوند بيان پر خپله پټه د وت
په توګه وليکي۔ بيان به د~\LR{$\TMstroke$} د بلاکونو يوه لړۍ وي (چې د~\LR{$M_e$}
په بيانوونکې لړۍ کښې مثبت صحيح اعداد تمثيلوي)۔
\olpseganchor{OLP-0267-B008}{end}

\olpseganchor{OLP-0267-B009}{start}
دې ته په پام اوس دا پوښتنه طبيعي ده: د ټيورينګ ماشين د شاخصونو کومې
تابعې پخپله د ټيورينګ ماشينونو په وسيله محاسبه کېدونکې دي؟ د ټيورينګ
ماشين د شاخصونو کوم خاصيتونه د ټيورينګ ماشينونو په وسيله پرېکړه کېدے
شي؟ يوه بېلګه: هغه تابع چې شاخص~\LR{$e$} د هغه شاخص~\LR{$e$} لرونکي ټيورينګ ماشين
د حالتونو شمېر ته وړي، د يوه ټيورينګ ماشين په وسيله محاسبه کېدونکې
ده۔ داسې ټيورينګ ماشين به دا کار کوي: پر هغې پټې له پيل وروسته چې د
\LR{$e$} دانو~\LR{$\TMstroke$} يو بلاک لري، لومړے به~\LR{$e$} په خپل بيان ناکوډ
کړي۔ بيان اوس پر پټه د~\LR{$\TMstroke$} د بلاکونو په يوه لړۍ تمثيلېږي۔
څرنګه چې په دې لړۍ کښې لومړے غړے د حالتونو شمېر دے، اوس
يوازې دا پاتې دي چې د~\LR{$\TMstroke$} له لومړي بلاک پرته نور هر څه پاک
کړي او بيا ودرېږي۔
\olpseganchor{OLP-0267-B009}{end}

\olpseganchor{OLP-0267-B010}{start}
يوه د پام وړ پايله دا ده:
\olpseganchor{OLP-0267-B010}{end}

\olpseganchor{OLP-0267-B011}{start}
\begin{thm}\ollabel{thm:universal-tm} داسې يو \emph{نړيوال ټيورينګ
  ماشين}~\LR{$U$} شته چې کله پر ننوت~\LR{$\tuple{e,n}$} پيل شي،
  \begin{enumerate}
    \item هله او يوازې هله درېږي چې~\LR{$M_e$} پر ننوت~\LR{$n$} ودرېږي، او
    \item که~\LR{$M_e$} له وت~\LR{$m$} سره ودرېږي، نو~\LR{$U$} هم همداسې کوي۔
  \end{enumerate}
  نو~\LR{$U$} هغه تابع~\LR{$f\colon \Nat \times \Nat \pto \Nat$} محاسبه کوي
  چې~\LR{$f(e,n) = m$} وي که~\LR{$M_e$} پر ننوت~\LR{$n$} له پيل وروسته له وت~\LR{$m$}
  سره ودرېږي، او که نۀ وي نو ناتعريفه ده۔
\end{thm}
\olpseganchor{OLP-0267-B011}{end}

\olpseganchor{OLP-0267-B012}{start}
\begin{proof}
  په عمل کښې د~\LR{$U$} جوړول بنسټيز ډول ناشوني دي، ځکه دا يو بېخي پېچلے
  ماشين دے۔ خو موږ د هغه د کار عمومي طرحه بيانولے، او بيا د
  چرچ--ټيورينګ اصل رابللے شو۔ کله چې~\LR{$U$} پيل شي، پټه ئې د~\LR{$e$} دانو
  \LR{$\TMstroke$} يو بلاک لري چې ورپسې د~\LR{$n$} دانو~\LR{$\TMstroke$} يو بلاک
  دے۔ لومړے شاخص~\LR{$e$} د ننوت~\LR{$n$} ښي لور ته ``ناکوډوي''۔ له دې څخه د
  اعدادو يو لست جوړېږي (يعنې د~\LR{$\TMstroke$} هغه بلاکونه چې په
  \LR{$\TMblank$} سره بېلېږي) چې د~\LR{$M_e$} ماشين لارښوونې بيانوي۔ بيا~\LR{$U$}
  د~\LR{$M_e$} د پيلني حالت عدد او عدد~\LR{$1$} د~\LR{$M_e$} د بيان ښي لور ته پر
  پټه ليکي۔ (دا هم په يوګوني ډول، د~\LR{$\TMstroke$} د بلاکونو په توګه،
  تمثيلېږي۔) ورپسې ننوت (د~\LR{$n$} دانو~\LR{$\TMstroke$} بلاک) ښي لور ته
  کاپي کوي---خو هر~\LR{$\TMstroke$} د درې دانو~\LR{$\TMstroke$} په يوه بلاک
  بدلوي (په ياد ولرئ، د~\LR{$\TMstroke$} نښې عدد~\LR{$3$} دے، \LR{$1$} د~\LR{$\TMendtape$}
  عدد او~\LR{$2$} د~\LR{$\TMblank$} عدد دے)۔ د دې بلاکونو د لړۍ په کيڼ سر کښې
  (چې د~\LR{$U$} پر پټه د~\LR{$\TMblank$} نښې سره بېلېږي) يو~\LR{$\TMstroke$}، يعنې
  د~\LR{$\TMendtape$} کوډ، ليکي۔
\olpseganchor{OLP-0267-B012}{end}

\olpseganchor{OLP-0267-B013}{start}
  اوس د~\LR{$U$} پر پټه دا شيان دي: شاخص~\LR{$e$}، عدد~\LR{$n$}، د پيلني حالت کوډ
  عدد (``اوسنی حالت'')، د سر د پيلني موقعيت عدد~\LR{$1$} (``د سر اوسنی
  موقعيت'')، او د ``پټې'' پيلنۍ منځپانګه (د~\LR{$\TMstroke$} د بلاکونو
  يوه لړۍ چې د~\LR{$M_e$} د نښو کوډ عددونه---``نښې''---تمثيلوي او په
  \LR{$\TMblank$} سره بېلېږي)۔
\olpseganchor{OLP-0267-B013}{end}

\olpseganchor{OLP-0267-B014}{start}
  اوس هغه څه په دې ډول مشابه چلوي چې~\LR{$M_e$} به پر ننوت~\LR{$n$} له پيل
  وروسته کړي وے:
  \begin{enumerate}
    \item د ``سر د اوسني موقعيت'' عدد~\LR{$k$} پيدا کول (په پيل کښې دا~\LR{$1$}
    دے)،
    \item د ``پټې'' \LR{$k$}-م بلاک ته تلل، څو وګوري چې هلته کومه ``نښه''
    ده،
    \item\ollabel{find-inst}%
    له اوسني ``حالت'' او ``نښې'' سره سمه لارښوونه پيدا کول،
    \item د ``پټې'' \LR{$k$}-م بلاک ته بېرته تلل او هلته ``نښه'' د هغې
    نښې په کوډ عدد بدلول چې~\LR{$M_e$} به ليکلې وے،
    \item سر هغه ځاے ته وړل چې اوسنی ``حالت'' پکې ثبتېږي، او هلته
    عدد د نوي حالت په عدد بدلول،
    \item هغه ځاے ته تلل چې د ``پټې موقعيت'' پکې ثبتېږي، او يو
    \LR{$\TMstroke$} پاکول يا يو~\LR{$\TMstroke$} زياتول (په ترتيب سره که
    لارښوونه د کيڼ يا ښي لور د حرکت خبره کوي)۔
    \item تکرارول۔\footnote{موږ دلته پر ځينو نازکو ستونزو سرسري
    تېرېږو۔ مثلاً کله چې~\LR{$U$} هغه شمېرند زياتوي چې د ``سر اوسنی
    موقعيت'' پرې ساتي، ښايي څه زيات ځاے ته اړ شي---په هغه حالت کښې
    به ټوله ``پټه'' ښي لور ته وخوځوي۔}
  \end{enumerate}
  که~\LR{$M_e$} پر ننوت~\LR{$n$} له پيل وروسته هېڅکله ونۀ درېږي، نو~\LR{$U$} هم
  هېڅکله نۀ درېږي، ځکه وت ئې ناتعريفه دے۔
\olpseganchor{OLP-0267-B014}{end}

\olpseganchor{OLP-0267-B015}{start}
  که په ګام~\olref{find-inst} کښې څرګنده شي چې د~\LR{$M_e$} بيان د اوسني
  ``حالت''/``نښې'' جوړې دپاره هېڅ لارښوونه نۀ لري، نو~\LR{$M_e$} به
  درېدلے وے۔ که دا پېښه شي، \LR{$U$} د خپلې پټې هغه برخه پاکوي چې د
  ``پټې'' کيڼ لور ته ده۔ د درې دانو~\LR{$\TMstroke$} د هر بلاک دپاره
  (چې د~\LR{$M_e$} پر پټه يو~\LR{$\TMstroke$} تمثيلوي) د خپلې پټې پر کيڼ سر
  يو~\LR{$\TMstroke$} ليکي، او ``پټه'' پرله‌پسې پاکوي۔ کله چې دا کار بشپړ
  شي، د~\LR{$U$} پټه د~\LR{$m$} اوږدوالي د~\LR{$\TMstroke$} يوازې يو بلاک لري۔
\olpseganchor{OLP-0267-B015}{end}

\olpseganchor{OLP-0267-B016}{start}
  که~\LR{$U$} پر ``پټه'' د درې دانو~\LR{$\TMstroke$} له بلاک پرته له بل څه
  سره مخ شي، سملاسي درېږي۔ څرنګه چې په دې حالت کښې د~\LR{$U$} پټه د
  \LR{$\TMstroke$} يوازې يو بلاک نۀ لري، وت ئې طبيعي عدد نۀ دے؛ يعنې په
  دې حالت کښې~\LR{$f(e,n)$} ناتعريفه ده۔
\end{proof}
\olpseganchor{OLP-0267-B016}{end}

\olpseganchor{OLP-0268-B004}{start}
\olfileid{tur}{und}{hal}
\olsection{د درېدنې مسئله}
\olpseganchor{OLP-0268-B004}{end}

\olpseganchor{OLP-0268-B005}{start}
\begin{explain}
فرض کړئ موږ د ټيورينګ ماشين د بيانونو کومه شمېره ثابته کړې ده۔ نو
هر ټيورينګ ماشين يو \emph{شاخص} اخلي: د ټيورينګ ماشين د بيانونو په
شمېره~\LR{$M_1$}، \LR{$M_2$}، \LR{$M_3$}، \dots{} کښې د هغه ځاے۔
\olpseganchor{OLP-0268-B005}{end}

\olpseganchor{OLP-0268-B006}{start}
موږ پوهېږو چې د ټيورينګ په وسيله نۀ محاسبه کېدونکې تابعې بايد موجودې
وي: د ټيورينګ ماشين د بيانونو سټ---او له دې سره د ټيورينګ ماشينونو
سټ---د شمېر وړ دے، خو له~\LR{$\Nat$} څخه~\LR{$\Nat$} ته د ټولو تابعو سټ
داسې نۀ دے۔ خو موږ د نامحاسبه کېدونکو تابعو ټاکلې بېلګې هم موندلے
شو۔ يوه داسې تابع د درېدنې تابع ده۔
\end{explain}
\olpseganchor{OLP-0268-B006}{end}

\olpseganchor{OLP-0268-B007}{start}
\begin{defn}[د درېدنې تابع]
 د \emph{درېدنې تابع}~\LR{$h$} داسې تعريفېږي:
\[
h(e,n) =
\begin{cases}
  \text{\RL{0}} & \text{\RL{که ماشين~\LR{$M_e$} د ننوت \LR{$n$} دپاره ونۀ درېږي}} \\
  \text{\RL{1}} & \text{\RL{که ماشين~\LR{$M_e$} د ننوت \LR{$n$} دپاره ودرېږي}}
\end{cases}
\]
\end{defn}
\olpseganchor{OLP-0268-B007}{end}

\olpseganchor{OLP-0268-B008}{start}
\begin{defn}[د درېدنې مسئله]
د \emph{درېدنې مسئله} د دې ټاکلو مسئله ده چې (د هر~\LR{$e$} او~\LR{$n$} دپاره)
ايا ټيورينګ ماشين~\LR{$M_e$} د~\LR{$n$} دانو کرښو پر ننوت درېږي که نه۔
\end{defn}
\olpseganchor{OLP-0268-B008}{end}

\olpseganchor{OLP-0268-B009}{start}
\begin{explain}
موږ د دې په ښودلو ثابتوو چې~\LR{$h$} د ټيورينګ په وسيله محاسبه کېدونکې
نۀ ده چې ور پورې تړلې تابع~\LR{$s$} د ټيورينګ په وسيله محاسبه کېدونکې
نۀ ده۔ دا ثبوت پر دې حقيقت تکيه کوي چې هر هغه څه چې يو ټيورينګ ماشين
ئې محاسبه کولے شي د يوه منضبط ټيورينګ ماشين په وسيله هم محاسبه کېدے
شي (\olref[mac][dis]{sec})، او پر دې حقيقت چې دوه ټيورينګ ماشينونه
سره نښلول کېدے شي، څو يو واحد ماشين جوړ کړي (\olref[mac][cmb]{sec})۔
\end{explain}
\olpseganchor{OLP-0268-B009}{end}

\olpseganchor{OLP-0268-B010}{start}
\begin{defn} تابع~\LR{$s$} داسې تعريفېږي:
\[
s(e) =
\begin{cases}
  \text{\RL{0}} & \text{\RL{که ماشين~\LR{$M_e$} د ننوت \LR{$e$} دپاره ونۀ درېږي}} \\
  \text{\RL{1}} & \text{\RL{که ماشين~\LR{$M_e$} د ننوت \LR{$e$} دپاره ودرېږي}}
\end{cases}
\]
\end{defn}
\olpseganchor{OLP-0268-B010}{end}

\olpseganchor{OLP-0268-B011}{start}
\begin{lem}
تابع~\LR{$s$} د ټيورينګ په وسيله محاسبه کېدونکې نۀ ده۔
\end{lem}
\olpseganchor{OLP-0268-B011}{end}

\olpseganchor{OLP-0268-B012}{start}
\begin{proof}
د تناقض دپاره فرضوو چې تابع~\LR{$s$} د ټيورينګ په وسيله محاسبه کېدونکې
ده۔ بيا به داسې ټيورينګ ماشين~\LR{$S$} موجود وے چې~\LR{$s$} محاسبه کوي۔ د
عموميت له زيان پرته فرضولے شو چې کله~\LR{$S$} درېږي، د لومړۍ مربع د
لوستلو پر مهال درېږي (يعنې منضبط دے)۔ دا ماشين له بل ماشين~\LR{$J$} سره
``نښلېدے'' شي چې که پر ننوت~\LR{$0$} پيل شي درېږي (يعنې که په پيلني
حالت کښې د پټې د پاي-نښې ښي لور ته مربع لوستلو پر مهال~\LR{$\TMblank$}
ولولي)، او که نۀ وي ښي لور ته روانېږي او هېڅکله نۀ درېږي۔
\LR{$S \concat J$}، يعنې هغه ماشين چې د~\LR{$S$} له~\LR{$J$} سره په نښلولو جوړېږي،
يو ټيورينګ ماشين دے، نو د کوم~\LR{$e$} دپاره~\LR{$M_e$} دے (يعنې په شمېره کښې
يو ځاے راځي)۔ \LR{$M_e$} د~\LR{$e$} دانو~\LR{$\TMstroke$} پر ننوت پيل کړئ۔ دوه
امکانونه دي: يا~\LR{$M_e$} درېږي يا نۀ درېږي۔
\begin{enumerate}
\item فرض کړئ~\LR{$M_e$} د~\LR{$e$} دانو~\LR{$\TMstroke$} پر ننوت درېږي۔ بيا~\LR{$s(e)
  = 1$}۔ نو کله چې~\LR{$S$} پر~\LR{$e$} پيل شي، د وت په توګه پر پټه له يوه
  \LR{$\TMstroke$} سره درېږي۔ بيا~\LR{$J$} پر داسې پټه پيلېږي چې يو~\LR{$\TMstroke$}
  لري۔ په هغه حالت کښې~\LR{$J$} نۀ درېږي۔ خو~\LR{$M_e$} هماغه ماشين~\LR{$S
  \concat J$} دے، نو بايد کټ مټ هغه څه وکړي چې لومړے~\LR{$S$} او ورپسې~\LR{$J$}
  ئې کوي (يعنې په دې حالت کښې ښي لور ته روان شي او هېڅکله ونۀ درېږي)۔
  نو~\LR{$M_e$} د~\LR{$e$} دانو~\LR{$\TMstroke$} پر ننوت نۀ شي درېدې۔
\olpseganchor{OLP-0268-B012}{end}

\olpseganchor{OLP-0268-B013}{start}
\item اوس فرض کړئ~\LR{$M_e$} د~\LR{$e$} دانو~\LR{$\TMstroke$} پر ننوت نۀ درېږي۔
  بيا~\LR{$s(e) = 0$}، او کله چې~\LR{$S$} پر ننوت~\LR{$e$} پيل شي، له تشې پټې سره
  درېږي۔ کله چې~\LR{$J$} پر تشه پټه پيل شي، سملاسي درېږي۔ بيا هم~\LR{$M_e$}
  هغه څه کوي چې لومړے~\LR{$S$} او ورپسې~\LR{$J$} ئې کوي، نو~\LR{$M_e$} بايد د~\LR{$e$}
  دانو~\LR{$\TMstroke$} پر ننوت ودرېږي۔
\end{enumerate}
په هر حالت کښې له خپل فرض سره تناقض ته رسېږو۔ دا ښيي چې داسې ټيورينګ
ماشين~\LR{$S$} نشي موجودېدے: \LR{$s$} د ټيورينګ په وسيله محاسبه کېدونکې نۀ
ده۔
\end{proof}
\olpseganchor{OLP-0268-B013}{end}

\olpseganchor{OLP-0268-B014}{start}
\begin{thm}[د درېدنې د مسئلې ناحلي]
\ollabel{thm:halting-problem} د درېدنې مسئله ناحل ده، يعنې تابع~\LR{$h$}
د ټيورينګ په وسيله محاسبه کېدونکې نۀ ده۔
\end{thm}
\olpseganchor{OLP-0268-B014}{end}

\olpseganchor{OLP-0268-B015}{start}
\begin{proof}
فرض کړئ~\LR{$h$} د ټيورينګ په وسيله محاسبه کېدونکې وے، مثلاً د ټيورينګ
ماشين~\LR{$H$} په وسيله۔ موږ له~\LR{$H$} څخه داسې ټيورينګ ماشين جوړولے شو چې
\LR{$s$} محاسبه کړي: لومړے د ننوت يوه کاپي جوړه کړئ (چې په~\LR{$\TMblank$}
نښه بېلېږي)۔ بيا بېرته پيل ته لاړ شئ او~\LR{$H$} وچلوئ۔ موږ په څرګند ډول
داسې ماشين جوړولے شو چې لومړے کار کوي
(\cref{tur:mac:dis:prob:copier} وګورئ)، او که~\LR{$H$} موجود وے، له داسې
کاپي‌کوونکي ماشين سره مو هغه ``نښلولے'' او داسې نوے ماشين مو ترې
جوړولے شو چې ټاکي ايا~\LR{$M_e$} پر ننوت~\LR{$e$} درېږي؛ يعنې~\LR{$s$} محاسبه کوي۔
خو موږ لا ښودلې ده چې داسې هېڅ ماشين نشي موجودېدے۔ نو~\LR{$h$} هم د
ټيورينګ په وسيله محاسبه کېدونکې نۀ ده۔
\end{proof}
\olpseganchor{OLP-0268-B015}{end}

\olpseganchor{OLP-0268-B016}{start}
\begin{prob}
د درې-درېدنې مسئله د داسې پرېکړيزې کړنلارې د ورکولو مسئله ده چې
وټاکي ايا په خپلسر غوره شوے ټيورينګ ماشين پر داسې پټه چې له دې پرته
تشه وي، د
درې دانو~\LR{$\TMstroke$} پر ننوت درېږي که نه۔ ثابته کړئ چې د درې-درېدنې
مسئله ناحل ده۔
\end{prob}
\olpseganchor{OLP-0268-B016}{end}

\olpseganchor{OLP-0268-B017}{start}
\begin{prob}
وښيئ چې که د ټيورينګ ماشين او ننوت د هغو جوړو~\LR{$M_e$} او~\LR{$n$} دپاره
چې~\LR{$e \neq n$} وي د درېدنې مسئله حل کېدونکې وي، نو د هغو حالتونو
دپاره هم حل کېدونکې ده چې~\LR{$e = n$} وي۔
\end{prob}
\olpseganchor{OLP-0268-B017}{end}

\olpseganchor{OLP-0268-B018}{start}
\begin{prob}
موږ ثابته کړه چې که ننوت داسې عدد~\LR{$e$} وي چې د ټولو ټيورينګ ماشينونو
د يوې شمېرې له لارې ټيورينګ ماشين~\LR{$M_e$} پېژني، نو د درېدنې مسئله
ناحل ده۔ که پر ځاے ئې د~\olref[tur][und][enu]{sec} د ټيورينګ
ماشينونو بيانونه نېغ په نېغه د ننوت په توګه ومنو، نو څه؟ ايا داسې
ټيورينګ ماشين موجودېدے شي چې د درېدنې مسئله پرېکړه کړي خو د شاخصونو
پر ځاے د ټيورينګ ماشينونو بيانونه د ننوت په توګه واخلي؟ څرګنده کړئ
چې ولې يا ولې نه۔
\end{prob}
\olpseganchor{OLP-0268-B018}{end}

\olpseganchor{OLP-0268-B019}{start}
\begin{prob} وښيئ چې \emph{قسمي} تابع~\LR{$s'$} چې داسې تعريفېږي:
  \[
  s'(e) =
  \begin{cases}
    \text{\RL{1}} & \text{\RL{که ماشين~\LR{$M_e$} د ننوت \LR{$e$} دپاره ودرېږي}}\\
    \text{\RL{ناتعريفه}} & \text{\RL{که ماشين~\LR{$M_e$} د ننوت \LR{$e$} دپاره ونۀ درېږي}}
  \end{cases}
  \]
  د ټيورينګ په وسيله محاسبه کېدونکې \emph{ده}۔
\end{prob}
\olpseganchor{OLP-0268-B019}{end}

\olpseganchor{OLP-0269-B004}{start}
\olfileid{tur}{und}{dec}
\olsection{د پرېکړې مسئله}
\olpseganchor{OLP-0269-B004}{end}

\olpseganchor{OLP-0269-B005}{start}
وايو د لومړۍ درجې منطق هله او يوازې هله \emph{پرېکړه کېدونکے} دے
چې د دې ټاکلو اغېزمن ميتود موجود وي چې يو ورکړل شوے تړلے فارمول
معتبر دے که نه۔ لکه چې څرګندېږي، داسې ميتود نشته: د لومړۍ درجې د
غونډلو د اعتبار د پرېکړې مسئله ناحل ده۔
\olpseganchor{OLP-0269-B005}{end}

\olpseganchor{OLP-0269-B006}{start}
د دې مهمې منفي پايلې د ټينګولو دپاره ثابتوو چې د پرېکړې مسئله د يوه
ټيورينګ ماشين په وسيله نۀ شي حل کېدے۔ يعنې ښيو چې داسې ټيورينګ ماشين
نشته چې هر کله پر داسې پټه پيل شي چې د لومړۍ درجې يو تړلے فارمول
لري، بالاخره ودرېږي او د دې له مخې~\LR{$1$} يا~\LR{$0$} وباسي چې تړلے فارمول
معتبر دے که نه۔ د چرچ--ټيورينګ د اصل له مخې هره محاسبه کېدونکې تابع
د ټيورينګ په وسيله محاسبه کېدونکې ده۔ نو که دا ``د اعتبار تابع''
بيخي په اغېزمنه توګه محاسبه کېدونکې وے، د ټيورينګ په وسيله به محاسبه
کېدونکې وے۔ که د ټيورينګ په وسيله محاسبه کېدونکې نۀ وي، نو په
اغېزمنه توګه هم نۀ شي محاسبه کېدے۔
\olpseganchor{OLP-0269-B006}{end}

\olpseganchor{OLP-0269-B007}{start}
د پرېکړې مسئلې د ناحلي د ثبوت تګلاره مو دا ده چې د درېدنې مسئله هغې
ته راکمه کړو۔ معنا ئې دا ده: موږ ثابته کړې چې هغه تابع~\LR{$h(e,w)$} چې
که په~\LR{$e$} بيان شوے ټيورينګ ماشين پر ننوت~\LR{$w$} ودرېږي، له وت~\LR{$1$} سره
درېږي، او که نۀ وي~\LR{$0$} وباسي، د ټيورينګ په وسيله محاسبه کېدونکې نۀ
ده۔ موږ به وښيو چې که داسې ټيورينګ ماشين وے چې د لومړۍ درجې د غونډلو
اعتبار پرېکړه کوي، نو داسې ټيورينګ ماشين هم شته چې~\LR{$h$} محاسبه کوي۔
څرنګه چې~\LR{$h$} د ټيورينګ ماشين په وسيله نۀ شي محاسبه کېدے، داسې
ټيورينګ ماشين هم نشته چې اعتبار پرېکړه کړي۔
\olpseganchor{OLP-0269-B007}{end}

\olpseganchor{OLP-0269-B008}{start}
په دې تګلاره کښې لومړے ګام دا ښودل دي چې د هر ننوت~\LR{$w$} او يوه
ټيورينګ ماشين~\LR{$M$} دپاره موږ په اغېزمنه توګه هغه تړلے فارمول
\LR{$!T(M, w)$} بيانولے شو چې د~\LR{$M$} د لارښوونو سټ او ننوت~\LR{$w$} تمثيلوي،
او هغه تړلے فارمول~\LR{$!E(M, w)$} چې ``\LR{$M$} بالاخره درېږي'' څرګندوي،
داسې چې:
\begin{quote}
  \LR{$\Entails !T(M, w) \lif !E(M,w)$} هله او يوازې هله چې~\LR{$M$} د
  ننوت~\LR{$w$} دپاره ودرېږي۔
\end{quote}
زموږ د ثبوت لويه برخه به د دې غونډلو~\LR{$!T(M, w)$} او~\LR{$!E(M, w)$} په
بيانولو او د دې خبرې په تاييدولو کښې وي چې~\LR{$!T(M, w) \lif !E(M, w)$}
هله او يوازې هله معتبر دے چې~\LR{$M$} پر ننوت~\LR{$w$} ودرېږي۔
\olpseganchor{OLP-0269-B008}{end}

\olpseganchor{OLP-0270-B004}{start}
\olfileid{tur}{und}{rep}
\olsection{د ټيورينګ ماشينونو تمثيلول}
\olpseganchor{OLP-0270-B004}{end}

\olpseganchor{OLP-0270-B005}{start}
\begin{explain}
د لومړۍ درجې منطق په يوه تړلے فارمول د ټيورينګ ماشينونو او د هغو
د چلند د تمثيلولو دپاره بايد يوه مناسبه ژبه تعريف کړو۔ ژبه دوه برخې
لري: د ماشين د تشکيلونو د بيانولو دپاره پريديکاتونه، او د اجرا د
ګامونو (``شېبو'') او پر پټه د موقعيتونو د شمېرلو دپاره تعبيرونه۔
\olpseganchor{OLP-0270-B005}{end}

\olpseganchor{OLP-0270-B006}{start}
موږ دوه ډوله پريديکاتونه ورپېژنو چې دواړه دوه-ځايه دي: د هر
حالت~\LR{$q$} دپاره يو پريديکات~\LR{$\Obj Q_q$}، او د پټې د هرې
نښې~\LR{$\sigma$} دپاره يو پريديکات~\LR{$\Obj S_\sigma$}۔ لومړي موږ ته
اجازه راکوي چې د~\LR{$M$} حالت او د هغه د پټې د سر موقعيت بيان کړو، او
دويم موږ ته اجازه راکوي چې د پټې منځپانګه بيان کړو۔
\olpseganchor{OLP-0270-B006}{end}

\olpseganchor{OLP-0270-B007}{start}
د پټې د سر د موقعيتونو او د اجرا شوو ګامونو د شمېر د څرګندولو دپاره
د اعدادو د څرګندولو يوې لارې ته اړ يو۔ دا د يوه ثابت سمبول~\LR{$\Obj
0$}، او د يوې \LR{$1$}-ځايه تابع~\LR{$\prime$}، يعنې د تالي تابع، په وسيله کېږي۔
د دود له مخې دا د خپل آرګومان \emph{وروسته} ليکل کېږي (او قوسونه
ترې پرېږدو)۔
\olpseganchor{OLP-0270-B007}{end}

\olpseganchor{OLP-0270-B008}{start}
د هر عدد~\LR{$n$} دپاره يو معياري ترم~\LR{$\num{n}$} شته، يعنې د~\LR{$n$}
\emph{عددنښه}، چې په~\LR{$\Lang L_M$} کښې ئې تمثيلوي۔ \LR{$\num{0}$} هم
\LR{$\Obj 0$} دے، \LR{$\num{1}$} هم~\LR{$\Obj 0'$} دے، \LR{$\num{2}$} هم~\LR{$\Obj 0''$}
دے، او داسې نور۔ په لا صوري ډول:
\begin{align*}
\num{0} & = \Obj 0 \\
\num{n+1} &= \num{n}'
\end{align*}
ترم~\LR{$\num{0}$}، يعنې~\LR{$\Obj 0$}، هم پر پټه تر ټولو کيڼ موقعيت او هم د
اجرا له لومړي ګام څخه مخکې وخت (پيلنی تشکيل) نوموي۔ ترم~\LR{$\num{1}$}،
يعنې~\LR{$\Obj 0'$}، د تر ټولو کيڼې مربع ښي لور ته مربع او د اجرا له لومړي
ګام څخه وروسته وخت نوموي، او داسې نور۔
\olpseganchor{OLP-0270-B008}{end}

\olpseganchor{OLP-0270-B009}{start}
موږ يو پريديکات~\LR{$<$} هم ورپېژنو، څو هم د پټې د موقعيتونو ترتيب
(چې هلته معنا ئې ``کيڼ لور ته'' ده) او هم د اجرا د ګامونو ترتيب
(چې هلته معنا ئې ``مخکې'' ده) څرګند کړو۔
\olpseganchor{OLP-0270-B009}{end}

\olpseganchor{OLP-0270-B010}{start}
کله چې ژبه مو تياره شوه، د~\LR{$!T(M, w)$} ``اکسيومونه'' لست کوو؛ يعنې
هغه تړلي فارمولونه چې په ګډه د~\LR{$M$} چلند بيانوي کله چې پر ننوت~\LR{$w$}
وچلول شي۔ داسې تړلي فارمولونه به وي چې پر~\LR{$\Obj 0$}، \LR{$\prime$}، او~\LR{$<$}
شرطونه ږدي، داسې تړلي فارمولونه چې د ننوت تشکيل بيانوي، او داسې
تړلي فارمولونه چې د~\LR{$M$} تشکيل د يوې ځانګړې لارښوونې له اجرا وروسته
بيانوي۔
\end{explain}
\olpseganchor{OLP-0270-B010}{end}

\olpseganchor{OLP-0270-B011}{start}
\begin{defn}
  \ollabel{defn:tm-descr}
يو ټيورينګ ماشين~\LR{$M = \tuple{Q, \Sigma, q_0, \delta}$} چې راکړل شي،
ژبه~\LR{$\Lang L_M$} له دې شيانو جوړه ده:
\begin{enumerate}
\item د هر حالت~\LR{$q \in Q$} دپاره يو دوه-ځايه پريديکات
  \LR{$\Obj Q_q(x, y)$}۔ په شهودي ډول~\LR{$\Obj Q_q(\num{m}, \num{n})$} دا
  څرګندوي: ``له~\LR{$n$} ګامونو وروسته، \LR{$M$} په حالت~\LR{$q$} کښې دے او~\LR{$m$}-مه
  مربع لولي۔''
\item د هرې نښې~\LR{$\sigma\in \Sigma$} دپاره يو دوه-ځايه پريديکات
  \LR{$\Obj S_\sigma(x, y)$}۔ په شهودي ډول~\LR{$\Obj S_\sigma(\num{m},
  \num{n})$} دا څرګندوي: ``له~\LR{$n$} ګامونو وروسته، \LR{$m$}-مه مربع نښه
  \LR{$\sigma$} لري۔''
\item يو ثابت سمبول~\LR{$\Obj 0$}
\item يوه يو-ځايه تابع~\LR{$\prime$}
\item يو دوه-ځايه پريديکات~\LR{$<$}
\end{enumerate}
\end{defn}
\olpseganchor{OLP-0270-B011}{end}

\olpseganchor{OLP-0270-B012}{start}
هغه تړلي فارمولونه چې پر ننوت~\LR{$w = \sigma_{i_1}\dots\sigma_{i_k}$}
د ټيورينګ ماشين~\LR{$M$} عمل بيانوي دا دي:
\begin{enumerate}
\item هغه اکسيومونه چې اعداد او~\LR{$<$} بيانوي:
\begin{enumerate}
\item يو تړلے فارمول چې وايي هر عدد له خپل تالي څخه کوچنی دے:
\[
\lforall[x][x < x']
\]
\item يو تړلے فارمول چې د~\LR{$<$} تعدي تضمينوي:
\[
\lforall[x][\lforall[y][\lforall[z][
      ((x < y \land y < z) \lif x < z)]]]
\]
\end{enumerate}
\item هغه اکسيومونه چې د ننوت تشکيل بيانوي:
\begin{enumerate}
\item له~\LR{$0$} ګامونو وروسته---مخکې له دې چې ماشين پيل شي---\LR{$M$} په
  پيلني حالت~\LR{$q_0$} کښې دے او مربع~\LR{$1$} لولي:
\[
\Obj Q_{q_0}(\num{1}, \num{0})
\]
\item لومړۍ \LR{$k+1$} مربعې نښې~\LR{$\TMendtape$}، \LR{$\sigma_{i_1}$}، \dots،
  \LR{$\sigma_{i_k}$} لري:
\[
\Obj S_\TMendtape(\num{0}, \num{0}) \land
\Obj S_{\sigma_{i_1}}(\num{1}, \num{0}) \land
\dots \land
\Obj S_{\sigma_{i_k}}(\num{k}, \num{0})
\]
\item له دې پرته پټه تشه ده:
\[
\lforall[x][(\num{k} < x \lif \Obj S_\TMblank(x, \num{0}))]
\]
\end{enumerate}
\item هغه اکسيومونه چې له يوه تشکيل څخه بل ته د حالت بدلون بيانوي:
\olpseganchor{OLP-0270-B012}{end}

\olpseganchor{OLP-0270-B013}{start}
د لاندې برخې دپاره، \LR{$!A(x, y)$} د ټولو هغو تړلي فارمولونه عطف وي چې
دا بڼه لري:
\[
\lforall[z][
  (((z < x \lor x < z) \land \Obj S_\sigma(z, y))
  \lif \Obj S_\sigma(z, y'))]
\]
چې~\LR{$\sigma \in \Sigma$} وي۔ موږ~\LR{$!A(\num{m},\num{n})$} د دې څرګندولو
دپاره کاروو: ``له مربع~\LR{$m$} پرته، له~\LR{$n+1$} ګامونو وروسته پټه هماغسې
ده لکه له~\LR{$n$} ګامونو وروسته۔''
\begin{enumerate}
\item \ollabel{rep-right} د هرې لارښوونې~\LR{$\delta(q_i, \sigma) =
  \tuple{q_j, \sigma', \TMright}$} دپاره دا تړلے فارمول:
\begin{align*}
& \lforall[x][\lforall[y][(
   (\Obj Q_{q_i}(x, y) \land \Obj S_{\sigma}(x, y)) \lif {}]] \\
&\qquad   (\Obj Q_{q_j}(x', y') \land \Obj S_{\sigma'}(x, y') \land
!A(x, y)))
\end{align*}
دا وايي چې که ماشين له~\LR{$y$} ګامونو وروسته په حالت~\LR{$q_i$} کښې وي او
هغه مربع~\LR{$x$} لولي چې نښه~\LR{$\sigma$} لري، نو له~\LR{$y+1$} ګامونو وروسته
مربع~\LR{$x+1$} لولي، په حالت~\LR{$q_j$} کښې وي، مربع~\LR{$x$} اوس~\LR{$\sigma'$} لري،
او له~\LR{$x$} پرته هره مربع هماغه نښه لري چې له~\LR{$y$} ګامونو وروسته ئې
درلوده۔
\olpseganchor{OLP-0270-B013}{end}

\olpseganchor{OLP-0270-B014}{start}
\item \ollabel{rep-left} د هرې لارښوونې~\LR{$\delta(q_i, \sigma) =
  \tuple{q_j, \sigma', \TMleft}$} دپاره دا تړلے فارمول:
\begin{align*}
& \lforall[x][\lforall[y][
    ((\Obj Q_{q_i}(x', y) \land \Obj S_{\sigma}(x', y)) \lif {}]]\\
& \qquad   (\Obj Q_{q_j}(x, y') \land \Obj S_{\sigma'}(x', y') \land
!A(x', y))) \land {}\\
& \lforall[y][((\Obj Q_{q_i}(\num{0}, y) \land \Obj S_{\sigma}(\num{0},
    y)) \lif {}]\\
& \qquad (\Obj Q_{q_j}(\num{0}, y') \land \Obj S_{\sigma'}(\num{0},
  y') \land !A(\num{0}, y)))
\end{align*}
يوه شېبه فکر وکړئ چې دا څنګه کار کوي: اوس له ``که مربع~\LR{$x$} لولي
\dots'' څخه نۀ، بلکې له ``که مربع~\LR{$x+1$} لولي~\dots'' څخه پيل
کوو۔ کيڼ لور ته حرکت معنا دا چې په بل ګام کښې ماشين مربع~\LR{$x$} لولي۔
خو هغه مربع~\LR{$x+1$} ده چې څه پرې ليکل کېږي۔ دا کار ځکه داسې کوو چې
منفي کول يا د مخکيني تابع نۀ لرو۔
د نابدلو نښو شرط له همدې ليکل شوې مربع پرته پر ټولو نورو مربعو لګېږي۔
\olpseganchor{OLP-0270-B014}{end}

\olpseganchor{OLP-0270-B015}{start}
په ياد ولرئ چې د~\LR{$x+1$} په بڼه اعداد~\LR{$1$}، \LR{$2$}، \dots، دي؛ يعنې دا
هغه حالت نۀ رانغاړي چې ماشين مربع~\LR{$0$} لولي او بايد کيڼ لور ته وخوځېږي
(چې البته نۀ شي کولے---يوازې پر ځاے پاتې کېږي)۔ هغه ځانګړے حالت د
دويم عطف له خوا رانغاړل کېږي: وايي چې که ماشين له~\LR{$y$} ګامونو وروسته
په حالت~\LR{$q_i$} کښې مربع~\LR{$0$} لولي او مربع~\LR{$0$} نښه~\LR{$\sigma$} لري، نو
له~\LR{$y+1$} ګامونو وروسته لا هم مربع~\LR{$0$} لولي، اوس په حالت~\LR{$q_j$} کښې
دے، پر مربع~\LR{$0$} نښه~\LR{$\sigma'$} ده، او له مربع~\LR{$0$} پرته مربعې هماغه
نښې لري چې له~\LR{$y$} ګامونو وروسته ئې لرلې۔
\item \ollabel{rep-stay} د هرې لارښوونې~\LR{$\delta(q_i, \sigma) =
  \tuple{q_j, \sigma', \TMstay}$} دپاره دا تړلے فارمول:
\begin{align*}
& \lforall[x][\lforall[y][(
   (\Obj Q_{q_i}(x, y) \land \Obj S_{\sigma}(x, y)) \lif {}]] \\
&\qquad   (\Obj Q_{q_j}(x, y') \land \Obj S_{\sigma'}(x, y') \land
!A(x, y)))
\end{align*}
\end{enumerate}
\end{enumerate}
\LR{$!T(M, w)$} د ټيورينګ ماشين~\LR{$M$} او ننوت~\LR{$w$} دپاره د پورتنيو ټولو
تړلي فارمولونه عطف وي۔
\olpseganchor{OLP-0270-B015}{end}

\olpseganchor{OLP-0270-B016}{start}
د دې څرګندولو دپاره چې~\LR{$M$} بالاخره درېږي، بايد داسې تړلے فارمول
ومومو چې وايي: ``له يو شمېر ګامونو وروسته به د حالت بدلون تابع
ناتعريفه وي۔'' \LR{$X$} د ټولو هغو جوړو~\LR{$\tuple{q, \sigma}$} سټ وي چې
\LR{$\delta(q, \sigma)$} ناتعريفه وي۔ بيا~\LR{$!E(M, w)$} دا تړلے فارمول وي:
\[
\lexists[x][\lexists[y][(\bigvee_{\tuple{q, \sigma} \in
      X}(\Obj Q_q(x, y) \land \Obj S_\sigma(x, y)))]]
\]
\olpseganchor{OLP-0270-B016}{end}

\olpseganchor{OLP-0270-B017}{start}
که داسې ټيورينګ ماشين وکاروو چې ټاکل شوے درېدنی حالت~\LR{$h$} ولري، خبره
لا اسانه ده: بيا دا تړلے فارمول~\LR{$!E(M, w)$}
\[
\lexists[x][\lexists[y][\Obj Q_h(x, y)]]
\]
دا څرګندوي چې ماشين بالاخره درېږي۔
\olpseganchor{OLP-0270-B017}{end}

\olpseganchor{OLP-0270-B018}{start}
\begin{prop}
\ollabel{prop:mlessk}
که~\LR{$m < k$} وي، نو~\LR{$!T(M, w) \Entails \num{m} < \num{k}$}۔
\end{prop}
\olpseganchor{OLP-0270-B018}{end}

\olpseganchor{OLP-0270-B019}{start}
\begin{proof}
تمرین۔
\end{proof}
\olpseganchor{OLP-0270-B019}{end}

\olpseganchor{OLP-0270-B020}{start}
\begin{prob}
\olref[tur][und][rep]{prop:mlessk} ثابت کړئ۔
(اشاره: پر~\LR{$k-m$} استقرا وکاروئ)۔
\end{prob}
\olpseganchor{OLP-0270-B020}{end}

\olpseganchor{OLP-0271-B004}{start}
\olfileid{tur}{und}{ver}
\olsection{د تمثيل تاييدول}
\olpseganchor{OLP-0271-B004}{end}

\olpseganchor{OLP-0271-B005}{start}
\begin{explain}
د دې د تاييدولو دپاره چې زموږ تمثيل کار کوي، بايد دوه خبرې ثابتې
کړو۔ لومړے بايد وښيو چې که~\LR{$M$} پر ننوت~\LR{$w$} ودرېږي، نو~\LR{$!T(M, w)
\lif !E(M, w)$} معتبر دے۔ بيا بايد معکوسه خبره وښيو، يعنې که~\LR{$!T(M,
w) \lif !E(M, w)$} معتبر وي، نو~\LR{$M$} په رښتيا هم پر ننوت~\LR{$w$} د چلولو
پر مهال بالاخره درېږي۔
\olpseganchor{OLP-0271-B005}{end}

\olpseganchor{OLP-0271-B006}{start}
د دې دواړو د ثبوت تګلارې بېخي سره توپير لري۔ د لومړۍ پايلې دپاره
بايد وښيو چې د لومړۍ درجې منطق يو تړلے فارمول (يعنې~\LR{$!T(M, w)
\lif !E(M, w)$}) معتبر دے۔ د دې تر ټولو اسانه لار د يوه
اشتقاق ورکول دي۔ خو زموږ ثبوت بايد د ټولو~\LR{$M$} او~\LR{$w$} دپاره
کار وکړي، نو په اصل کښې داسې يو واحد تړلے فارمول نشته چې موږ ئې
دپاره اشتقاق ورکړو، بلکې نامتناهي ډېر دي۔ نو تر ټولو ښه کار
چې کولے ئې شو دا دے چې په استقرا ثابته کړو: \LR{$M$} او~\LR{$w$} هر ډول وي، او
پر ننوت~\LR{$w$} د~\LR{$M$} درېدل هر څومره ګامونه اخلي، د~\LR{$!T(M, w) \lif
!E(M, w)$} يو اشتقاق به موجود وي۔
\olpseganchor{OLP-0271-B006}{end}

\olpseganchor{OLP-0271-B007}{start}
په طبيعي ډول به زموږ استقرا د هغو ګامونو پر شمېر مخته ځي چې~\LR{$M$} ئې
درېدني تشکيل ته تر رسېدو مخکې اخلي۔ په خپل استقرايي ثبوت کښې به ټينګه
کړو چې پر ننوت~\LR{$w$} د~\LR{$M$} د چلولو د هر ګام~\LR{$n$} دپاره~\LR{$!T(M, w)
\Entails !C(M, w, n)$}، چې~\LR{$!C(M, w, n)$} له~\LR{$n$} ګامونو وروسته پر~\LR{$w$}
د چلېدلي~\LR{$M$} تشکيل په سمه توګه بيانوي۔ اوس که~\LR{$M$} پر ننوت~\LR{$w$}، مثلاً
له~\LR{$n$} ګامونو وروسته، ودرېږي، \LR{$!C(M, w, n)$} به يو درېدنی تشکيل بيان
کړي۔ دا به هم وښيو چې هر کله~\LR{$!C(M, w, n)$} يو درېدنی تشکيل بيانوي،
\LR{$!C(M, w, n) \Entails !E(M, w)$}۔ نو که~\LR{$M$} پر ننوت~\LR{$w$} ودرېږي، د
کوم~\LR{$n$} دپاره به~\LR{$M$} له~\LR{$n$} ګامونو وروسته په درېدني تشکيل کښې وي۔
له همدې امله~\LR{$!T(M, w) \Entails !C(M, w, n)$}، چې~\LR{$!C(M, w, n)$} يو
درېدنی تشکيل بيانوي؛ او څرنګه چې په هغه حالت کښې~\LR{$!C(M, w, n)
\Entails !E(M, w)$}، نو~\LR{$!T(M, w) \Entails !E(M, w)$} تر لاسه کوو؛
يعنې~\LR{$\Entails !T(M, w) \lif !E(M, w)$}۔
\olpseganchor{OLP-0271-B007}{end}

\olpseganchor{OLP-0271-B008}{start}
د معکوسې خبرې تګلاره بېخي توپير لري۔ دلته فرضوو چې~\LR{$\Entails !T(M,
w) \lif !E(M, w)$}، او بايد ثابته کړو چې~\LR{$M$} پر ننوت~\LR{$w$} درېږي۔ له
فرض څخه مو تر لاسه کېږي چې~\LR{$!T(M, w) \Entails !E(M, w)$}، يعنې
\LR{$!E(M, w)$} په هر هغه جوړښت کښې رښتيا دے چې~\LR{$!T(M, w)$} پکې
رښتيا دے۔ نو داسې يو جوړښت~\LR{$\Struct{M}$} بيان کوو چې~\LR{$!T(M,
w)$} پکې رښتيا دے: ډومېن به ئې~\LR{$\Nat$} وي، او د ټولو~\LR{$\Obj Q_q$} او
\LR{$\Obj S_\sigma$} تعبير به پر ننوت~\LR{$w$} د~\LR{$M$} د چلولو پر مهال د تشکيلونو
له خوا ورکول کېږي۔ نو مثلاً~\LR{$\Sat{M}{\Obj Q_q(\num{m}, \num{n})}$}
هله او يوازې هله چې~\LR{$M$} پر ننوت~\LR{$w$} د~\LR{$n$} ګامونو دپاره وچلول شي،
په حالت~\LR{$q$} کښې وي او مربع~\LR{$m$} لولي۔ اوس څرنګه چې د فرض له مخې
\LR{$!T(M, w) \Entails !E(M, w)$}، او د جوړښت له مخې~\LR{$\Sat{M}{!T(M, w)}$}،
نو~\LR{$\Sat{M}{!E(M, w)}$}۔ خو~\LR{$\Sat{M}{!E(M, w)}$} هله او يوازې هله چې
کوم~\LR{$n \in \Domain{M} = \Nat$} موجود وي داسې چې پر ننوت~\LR{$w$} چلېدلے
\LR{$M$} له~\LR{$n$} ګامونو وروسته په درېدني تشکيل کښې وي۔
\end{explain}
\olpseganchor{OLP-0271-B008}{end}

\olpseganchor{OLP-0271-B009}{start}
\begin{defn}
\LR{$!C(M, w, n)$} دا تړلے فارمول وي:
\[
\Obj Q_q(\num{m}, \num{n}) \land \Obj S_{\sigma_0}(\num{0}, \num{n})
\land \dots \land \Obj S_{\sigma_k}(\num{k}, \num{n}) \land
\lforall[x][(\num{k} < x \lif \Obj S_\TMblank(x, \num{n}))]
\]
چې~\LR{$q$} په وخت~\LR{$n$} کښې د~\LR{$M$} حالت دے، \LR{$M$} په وخت~\LR{$n$} کښې مربع~\LR{$m$}
لولي، مربع~\LR{$i$} په وخت~\LR{$n$} کښې نښه~\LR{$\sigma_i$} لري چې~\LR{$0 \le i \le
k$}، او~\LR{$k$} د دې دوو حدونو لويه شمېره ده: په پيل کښې د پټې تر ټولو ښۍ ناتشه
مربع، او له پيل څخه تر~\LR{$n$} ګامونو پورې د سر تر ټولو ښۍ ليدلې مربع۔ د سر د پيل مربع هم شمېرل کېږي۔
\end{defn}
\olpseganchor{OLP-0271-B009}{end}

\olpseganchor{OLP-0271-B010}{start}
\begin{lem}\ollabel{lem:halt-config-implies-halt}
که پر ننوت~\LR{$w$} چلېدلے~\LR{$M$} له~\LR{$n$} ګامونو وروسته په يوه درېدني تشکيل
کښې وي، نو~\LR{$!C(M, w, n) \Entails !E(M, w)$}۔
\end{lem}
\olpseganchor{OLP-0271-B010}{end}

\olpseganchor{OLP-0271-B011}{start}
\begin{proof}
فرض کړئ~\LR{$M$} د ننوت~\LR{$w$} دپاره له~\LR{$n$} ګامونو وروسته درېږي۔ داسې يو
حالت~\LR{$q$}، مربع~\LR{$m$}، او نښه~\LR{$\sigma$} شته چې:
\begin{enumerate}
\item له~\LR{$n$} ګامونو وروسته~\LR{$M$} په حالت~\LR{$q$} کښې دے، هغه مربع~\LR{$m$}
  لولي چې~\LR{$\sigma$} پرې ښکاري۔
\item د حالت بدلون تابع~\LR{$\delta(q, \sigma)$} ناتعريفه ده۔
\end{enumerate}
\LR{$!C(M, w, n)$} د دې تشکيل بيان دے او دا فقرې~\LR{$\Obj Q_{q}(\num{m},
\num{n})$} او~\LR{$\Obj S_{\sigma}(\num{m}, \num{n})$} به پکې شاملې وي۔
دا فقرې په ګډه~\LR{$!E(M, w)$} راايجابوي:
\[
\lexists[x][\lexists[y][(\bigvee_{\tuple{q, \sigma} \in
      X}(\Obj Q_q(x, y) \land \Obj S_\sigma(x, y)))]]
\]
ځکه~\LR{$\Obj Q_{q}(\num{m}, \num{n}) \land \Obj S_{\sigma}(\num{m},
\num{n}) \Entails \bigvee_{\tuple{q, \sigma} \in X} (\Obj Q_q(\num{m},
\num{n}) \land \Obj S_{\sigma}(\num{m}, \num{n}))$}، ځکه
\LR{$\tuple{q,\sigma} \in X$}۔
\end{proof}
\olpseganchor{OLP-0271-B011}{end}

\olpseganchor{OLP-0271-B012}{start}
\begin{explain}
نو که~\LR{$M$} د ننوت~\LR{$w$} دپاره ودرېږي، داسې~\LR{$n$} شته چې~\LR{$!C(M, w, n)
\Entails !E(M,w)$}۔ اوس به وښيو چې د هر وخت~\LR{$n$} دپاره~\LR{$!T(M, w)
\Entails !C(M, w, n)$}۔
\end{explain}
\olpseganchor{OLP-0271-B012}{end}

\olpseganchor{OLP-0271-B013}{start}
\begin{lem}
\ollabel{lem:config}
د هر~\LR{$n$} دپاره، که~\LR{$M$} د~\LR{$n$}-م تشکيل تر رسېدو مخکې نۀ وي درېدلے، نو
\LR{$!T(M, w) \Entails !C(M, w, n)$}۔
\end{lem}
\olpseganchor{OLP-0271-B013}{end}

\olpseganchor{OLP-0271-B014}{start}
\begin{proof}
د استقرا بنسټ: که~\LR{$n = 0$} وي، د~\LR{$!C(M, w, 0)$} عاطفې د~\LR{$!T(M, w)$} له پيليزو اکسيومونو څخه راځي۔ د تش ننوت په حالت کښې سر په لومړۍ مربع پيل کوي۔ د پټې د تشوالي له اکسيوم او د هر عدد له خپل تالي څخه د کوچنيوالي له اکسيوم څخه راځي چې هم هغه مربع او هم تر هغې ښۍ مربعې تشې دي۔
\olpseganchor{OLP-0271-B014}{end}

\olpseganchor{OLP-0271-B015}{start}
استقرايي فرض: که~\LR{$M$} له~\LR{$n$}-م ګام څخه مخکې نۀ وي درېدلے، نو
\LR{$!T(M,w) \Entails !C(M, w, n)$}۔ بايد وښيو چې (تر هغه چې~\LR{$!C(M, w,
n)$} يو درېدنی تشکيل نۀ بيانوي) \LR{$!T(M, w) \Entails !C(M, w, n+1)$}۔
\olpseganchor{OLP-0271-B015}{end}

\olpseganchor{OLP-0271-B016}{start}
فرض کړئ~\LR{$n \ge 0$} او پر~\LR{$w$} پيل شوي~\LR{$M$} له~\LR{$n$} ګامونو وروسته په
حالت~\LR{$q$} کښې مربع~\LR{$m$} لولي۔ څرنګه چې~\LR{$M$} له~\LR{$n$} ګامونو وروسته نۀ
درېږي، د~\LR{$M$} په پروګرام کښې بايد له لاندې درېو بڼو څخه يوه لارښوونه
وي:
\olpseganchor{OLP-0271-B016}{end}

\olpseganchor{OLP-0271-B017}{start}
\begin{enumerate}
\item \ollabel{right} \LR{$\delta(q, \sigma) = \tuple{q', \sigma', \TMright}$}
\olpseganchor{OLP-0271-B017}{end}

\olpseganchor{OLP-0271-B018}{start}
\item \ollabel{left} \LR{$\delta(q, \sigma) = \tuple{q', \sigma', \TMleft}$}
\olpseganchor{OLP-0271-B018}{end}

\olpseganchor{OLP-0271-B019}{start}
\item \ollabel{stay} \LR{$\delta(q, \sigma) = \tuple{q', \sigma', \TMstay}$}
\end{enumerate}
\olpseganchor{OLP-0271-B019}{end}

\olpseganchor{OLP-0271-B020}{start}
موږ به دا درې واړه حالتونه په وار سره وڅېړو۔
\olpseganchor{OLP-0271-B020}{end}

\olpseganchor{OLP-0271-B021}{start}
\begin{enumerate}
\item فرض کړئ د~\olref{right} په بڼه يوه لارښوونه شته۔ د
  \LR{\olref[rep]{defn:tm-descr}(\olref[rep]{rep-right})} له مخې معنا ئې دا
  ده چې
\begin{align*}
& \lforall[x][\lforall[y][((\Obj Q_{q}(x,
  y) \land \Obj S_\sigma(x, y)) \lif {}]]\\
& \qquad (\Obj Q_{q'}(x',
  y') \land \Obj S_{\sigma'}(x, y') \land !A(x, y)))
\intertext{\RL{د~\LR{$!T(M,w)$} يوه عاطفه ده۔ دا لاندې تړلے فارمول راايجابوي
  (کلي منځته راوړنه، \LR{$\num{m}$} د~\LR{$x$} دپاره او~\LR{$\num{n}$} د~\LR{$y$}
  دپاره):}}
& (\Obj Q_{q}(\num{m}, \num{n}) \land \Obj S_{\sigma}(\num{m},
\num{n})) \lif {}\\
& \qquad (\Obj Q_{q'}(\num{m}', \num{n}') \land
\Obj S_{\sigma'}(\num{m}, \num{n}') \land !A(\num{m}, \num{n})).
\intertext{\RL{د استقرايي فرض له مخې~\LR{$!T(M, w) \Entails !C(M, w, n)$}؛
  يعنې:}}
& \Obj Q_q(\num{m}, \num{n}) \land \Obj S_{\sigma_0}(\num{0}, \num{n})
\land \dots \land \Obj S_{\sigma_k}(\num{k}, \num{n}) \land \\
& \qquad\lforall[x][(\num{k} < x \lif \Obj S_\TMblank(x, \num{n}))]\\
\intertext{\RL{څرنګه چې له~\LR{$n$} ګامونو وروسته د پټې مربع~\LR{$m$} نښه~\LR{$\sigma$}
  لري، اړونده عاطفه~\LR{$\Obj S_\sigma(\num{m}, \num{n})$} ده، نو دا
  راايجابوي:}}
& \Obj Q_{q}(\num{m}, \num{n}) \land \Obj S_{\sigma}(\num{m},
   \num{n})
\intertext{\RL{اوس تر لاسه کوو:}}
& \Obj Q_{q'}(\num{m}', \num{n}') \land \Obj S_{\sigma'}(\num{m},
  \num{n}') \land {}\\
& \qquad \bigwedge_{\substack{0 \le i \le k \\ i \ne m}}\Obj S_{\sigma_i}(\num{i}, \num{n}') \land {}\\
& \qquad \lforall[x][(\num{k} < x \lif \Obj S_\TMblank(x, \num{n}'))]
\end{align*}
په دې ډول: لومړۍ کرښه نېغ په نېغه د مخکيني شرطي د تالي څخه د مودوس
پونېنس په وسيله راځي۔ د منځنۍ کرښې هره عاطفه په~\LR{$!C(M, w, n)$} کښې له اړوندې عاطفې او~\LR{$!A(\num{m}, \num{n})$} څخه راځي۔ د~\LR{$i < m$} او~\LR{$m < i$} دواړو حالتونو کښې د تالي او تعدي اکسيومونه د اړوندو عددنښو تر منځ ترتيب راکوي۔
\olpseganchor{OLP-0271-B021}{end}

\olpseganchor{OLP-0271-B022}{start}
که~\LR{$m < k$} وي، \LR{$!T(M,w) \Entails \num{m} < \num {k}$}
(\olref[rep]{prop:mlessk}) او د~\LR{$<$} د تعدي له مخې
\LR{$\lforall[x][(\num{k} < x \lif \num{m} < x)]$} لرو۔ که~\LR{$m = k$} وي،
نو~\LR{$\lforall[x][(\num{k} < x \lif \num{m} < x)]$} يوازې د منطق له
مخې راځي۔ بيا وروستۍ کرښه په~\LR{$!C(M, w, n)$} کښې له اړوندې عاطفې،
\LR{$\lforall[x][(\num{k} < x \lif \num{m} < x)]$}، او~\LR{$!A(\num{m},
\num{n})$} څخه راځي۔ که~\LR{$m<k$} وي، دا لا همدا~\LR{$!C(M, w, n+1)$} دے۔
\olpseganchor{OLP-0271-B022}{end}

\olpseganchor{OLP-0271-B023}{start}
اوس فرض کړئ~\LR{$m=k$}۔ په هغه حالت کښې د پټې سر له~\LR{$n+1$} ګامونو وروسته
مربع~\LR{$k+1$} هم ليدلې ده، چې اوس تر ټولو ښي ليدل شوې مربع ده۔ نو
\LR{$!C(M, w, n+1)$} يوه نوې عاطفه~\LR{$\Obj S_\TMblank(\num{k}',\num{n}')$}
لري، او وروستۍ عاطفه ئې~\LR{$\lforall[x][(\num{k}' < x \lif \Obj
S_\TMblank(x, \num{n}'))]$} ده۔ بايد تاييد کړو چې دا دوه تړلي فارمولونه
هم ايجابېږي۔
\olpseganchor{OLP-0271-B023}{end}

\olpseganchor{OLP-0271-B024}{start}
موږ لا~\LR{$\lforall[x][(\num{k} < x \lif \Obj S_\TMblank(x,
  \num{n}'))]$} لرو۔ په ځانګړي ډول، له دې څخه~\LR{$\num{k} < \num{k}'
\lif \Obj S_\TMblank(\num{k}', \num{n}')$} تر لاسه کوو۔ له اکسيوم
\LR{$\lforall[x][x < x']$} څخه~\LR{$\num{k} < \num{k}'$} تر لاسه کوو۔ د مودوس
پونېنس له مخې~\LR{$\Obj S_\TMblank(\num{k}',\num{n}')$} راځي۔
\olpseganchor{OLP-0271-B024}{end}

\olpseganchor{OLP-0271-B025}{start}
همدارنګه، څرنګه چې~\LR{$!T(M,w) \Entails \num{k} < \num{k}'$}، د~\LR{$<$}
د تعدي اکسيوم موږ ته~\LR{$\lforall[x][(\num{k}' < x \lif \Obj
  S_\TMblank(x, \num{n}'))]$} راکوي۔ (د دې تاييد د تمرین په توګه
پرېږدو۔)
\olpseganchor{OLP-0271-B025}{end}

\olpseganchor{OLP-0271-B026}{start}
\item فرض کړئ د~\olref{left} په بڼه يوه لارښوونه شته۔ بيا د
  \LR{\olref[rep]{defn:tm-descr}(\olref[rep]{rep-left})} له مخې،
\begin{align*}
& \lforall[x][\lforall[y][((\Obj Q_{q}(x', y) \land \Obj
    S_{\sigma}(x', y)) \lif {}]]\\
& \qquad (\Obj Q_{q'}(x, y') \land \Obj
  S_{\sigma'}(x', y') \land !A(x', y))) \land {}\\
& \lforall[y][((\Obj Q_{q}(\num{0}, y) \land \Obj S_{\sigma}(\num{0},
    y)) \lif {}]\\
& \qquad (\Obj Q_{q'}(\num{0}, y') \land \Obj S_{\sigma'}(\num{0}, y')
  \land !A(\num{0}, y)))
\intertext{\RL{د~\LR{$!T(M,w)$} يوه عاطفه ده۔ که~\LR{$m>0$} وي، نو~\LR{$l = m - 1$}
  پرېږدئ (يعنې~\LR{$m = l+1$})۔ د پورتني تړلے فارمول لومړۍ عاطفه دا
  راايجابوي:}}
& (\Obj Q_{q}(\num{l}', \num{n}) \land \Obj S_{\sigma}(\num{l}', \num{n}))
\lif {} \\
& \qquad (\Obj Q_{q'}(\num{l}, \num{n}') \land \Obj S_{\sigma'}(\num{l}',
\num{n}') \land !A(\num{l}', \num{n}))
\intertext{\RL{که نۀ وي، \LR{$l = m = 0$} پرېږدئ او لاندې هغه تړلے فارمول
  وګورئ چې دويمه عاطفه ئې راايجابوي:}}
& ((\Obj Q_{q}(\num{0}, \num{n}) \land \Obj S_{\sigma}(\num{0}, \num{n})) \lif {}\\
& \qquad   (\Obj Q_{q'}(\num{0}, \num{n}') \land \Obj S_{\sigma'}(\num{0}, \num{n}') \land
!A(\num{0}, \num{n})))
\intertext{\RL{په دواړو حالتونو کښې اړوند شرطي فارمول د اوسنۍ تشکيل له بيان سره دا راايجابوي:}}
&  \Obj Q_{q'}(\num{l}, \num{n}') \land \Obj S_{\sigma'}(\num{m},
  \num{n}') \land {}\\
& \qquad \bigwedge_{\substack{0 \le i \le k \\ i \ne m}}\Obj S_{\sigma_i}(\num{i}, \num{n}') \land {}\\
& \qquad  \lforall[x][(\num{k} < x
  \lif \Obj S_\TMblank(x, \num{n}'))]
\end{align*}
لکه مخکې۔ (په ياد ولرئ چې په لومړي حالت کښې~\LR{$\num{l}' \ident
\num{l+1} \ident \num{m}$}، او په دويم حالت کښې~\LR{$\num{l} \ident
\num{0}$}۔) په هر حالت کښې د ليکل شوې مربع نوې نښه او د پاتې مربعو پخوانۍ نښې يوځاے د~\LR{$!C(M, w, n+1)$} عاطفې جوړوي۔
\olpseganchor{OLP-0271-B026}{end}

\olpseganchor{OLP-0271-B027}{start}
\item د~\olref{stay} په حالت کښې سر په هماغه مربع پاتې کېږي او د پټې ښے حد هماغه وي۔ د هماغې مربع نوې نښه د لارښوونې د شرطي فارمول له تالي څخه اخلو۔ د نورو مربعو نښې، د پخواني حد ښۍ تشې مربعې هم، د نابدلې پټې له شرط څخه اخلو۔ جزئيات په تمرين کښې بشپړ کړئ۔
\end{enumerate}
نو د هر هغه~\LR{$n$} دپاره چې ماشين پکښې تشکيل ته رسېدلے وي، \LR{$!T(M, w) \Entails !C(M, w, n)$}۔ که دغه تشکيل درېدنے وي، استقرا همدلته پای ته رسېږي۔
\end{proof}
\olpseganchor{OLP-0271-B027}{end}

\olpseganchor{OLP-0271-B028}{start}
\begin{prob}
د~\olref[tur][und][ver]{lem:config} د ثبوت حالت
\olref[tur][und][ver]{stay} بشپړ کړئ۔
\end{prob}
\olpseganchor{OLP-0271-B028}{end}

\olpseganchor{OLP-0271-B029}{start}
\begin{prob}
له~\LR{$\Obj S_{\sigma_i}(\num{i}, \num{n})$} او~\LR{$!A(m, n)$} څخه د
\LR{$\Obj S_{\sigma_i}(\num{i}, \num{n}')$} يو اشتقاق ورکړئ
(په دې فرض چې~\LR{$i \neq m$}؛ يعنې يا~\LR{$i < m$} يا~\LR{$m < i$})۔
\end{prob}
\olpseganchor{OLP-0271-B029}{end}

\olpseganchor{OLP-0271-B030}{start}
\begin{prob}
له~\LR{$\lforall[x][(\num{k} < x \lif \Obj S_\TMblank(x, \num{n}'))]$}،
\LR{$\lforall[x][x < x']$}، او
\LR{$\lforall[x][\lforall[y][\lforall[z][ ((x < y \land y < z) \lif x <
      z)]]]$} څخه د~\LR{$\lforall[x][(\num{k}' < x \lif \Obj
  S_\TMblank(x, \num{n}'))]$} يو اشتقاق ورکړئ۔)
\end{prob}
\olpseganchor{OLP-0271-B030}{end}


\olpseganchor{OLP-0271-B031}{start}
\begin{lem}
\ollabel{lem:valid-if-halt}
که~\LR{$M$} پر ننوت~\LR{$w$} ودرېږي، نو~\LR{$!T(M, w) \lif
!E(M, w)$} معتبر دے۔
\end{lem}
\olpseganchor{OLP-0271-B031}{end}

\olpseganchor{OLP-0271-B032}{start}
\begin{proof}
د~\olref{lem:config} له مخې، په هر وخت~\LR{$n$} کښې چې~\LR{$M$} ترې مخکې نۀ وي درېدلے، د هغه د تشکيل بيان~\LR{$!C(M, w, n)$} د~\LR{$!T(M, w)$} له خوا ايجابېږي۔ فرض کړئ~\LR{$M$} له~\LR{$k$} ګامونو وروسته درېږي۔ په هغه وخت کښې به
د کوم~\LR{$m \in \Nat$} دپاره مربع~\LR{$m$} لولي۔ بيا~\LR{$!C(M, w, k)$} د~\LR{$M$}
يو درېدنی تشکيل بيانوي؛ يعنې~\LR{$\Obj Q_q(\num{m}, \num{k})$} او~\LR{$\Obj
S_\sigma(\num{m}, \num{k})$} دواړه د عاطفو په توګه لري، چې
\LR{$\delta(q,\sigma)$} ناتعريفه ده۔ نو د
\olref{lem:halt-config-implies-halt} له مخې~\LR{$!C(M, w, k) \Entails
!E(M, w)$}۔ خو څرنګه چې~\LR{$!T(M, w) \Entails !C(M, w, k)$}، نو~\LR{$!T(M,
w) \Entails !E(M, w)$}، او ځکه~\LR{$!T(M, w) \lif !E(M, w)$} معتبر دے۔
\end{proof}
\olpseganchor{OLP-0271-B032}{end}


\olpseganchor{OLP-0271-B033}{start}
\begin{explain}
د خپلې ادعا د تاييد بشپړولو دپاره بايد معکوس لوری هم ټينګ کړو: که
\LR{$!T(M, w) \lif !E(M, w)$} معتبر وي، نو~\LR{$M$} په رښتيا هم پر ننوت~\LR{$w$}
له پيل وروسته درېږي۔
\end{explain}
\olpseganchor{OLP-0271-B033}{end}

\olpseganchor{OLP-0271-B034}{start}
\begin{lem}
\ollabel{lem:halt-if-valid}
که~\LR{$\Entails !T(M, w) \lif !E(M, w)$}، نو~\LR{$M$} پر ننوت~\LR{$w$} درېږي۔
\end{lem}
\olpseganchor{OLP-0271-B034}{end}

\olpseganchor{OLP-0271-B035}{start}
\begin{proof}
د~\LR{$\Lang L_M$} هغه جوړښت~\LR{$\Struct M$} وګورئ چې ډومېن ئې
\LR{$\Nat$} دے، \LR{$\Obj 0$} په~\LR{$0$}، \LR{$\prime$} په تالي تابع، او~\LR{$<$} د کوچنيوالي
په اړيکه تعبيروي، او محمولونه~\LR{$\Obj Q_q$} او~\LR{$\Obj S_\sigma$} په دې
ډول تعبيروي:
\begin{align*}
  \Assign{\Obj Q_q}{M} & =
\Setabs{\tuple{m, n}}{\begin{array}{ll}\text{\RL{پر~\LR{$w$} له پيل وروسته، له~\LR{$n$} ګامونو وروسته،}}\\ \text{\RL{\LR{$M$} په حالت \LR{$q$}
  کښې دے او مربع~\LR{$m$} لولي}}
\end{array}} \\
\Assign{\Obj S_\sigma}{M} & = \Setabs{\tuple{m, n}}{\begin{array}{ll}
\text{\RL{پر~\LR{$w$} له پيل وروسته، له \LR{$n$} ګامونو وروسته،}}\\ \text{\RL{مربع~\LR{$m$} د~\LR{$M$}
  نښه~\LR{$\sigma$} لري}}
\end{array}}
\end{align*}
له درېدني تشکيل وروسته د هېڅ حالت محمول صدق نۀ کوي، او د پټې نښې د درېدنې په حالت کښې ساتو۔ په بله وينا، موږ جوړښت~\LR{$\Struct{M}$} داسې جوړوو چې هغه څه ګام
په ګام بيان کړي چې پر ننوت~\LR{$w$} پيل شوے~\LR{$M$} په رښتيا کوي۔ په څرګند
ډول~\LR{$\Sat{M}{!T(M, w)}$}۔ که~\LR{$\Entails !T(M, w) \lif !E(M, w)$}، نو
\LR{$\Sat{M}{!E(M, w)}$} هم؛ يعنې:
\[
\Sat{M}{\lexists[x][\lexists[y][(\bigvee_{\tuple{q, \sigma} \in
      X}(\Obj Q_q(x, y) \land \Obj S_\sigma(x, y)))]]}.
\]
څرنګه چې~\LR{$\Domain{M} = \Nat$}، بايد داسې~\LR{$m$}، \LR{$n \in \Nat$} موجود
وي چې د کوم~\LR{$q$} او~\LR{$\sigma$} دپاره~\LR{$\Sat{M}{\Obj Q_q(\num{m},
\num{n}) \land \Obj S_\sigma(\num{m}, \num{n})}$}، داسې چې
\LR{$\delta(q, \sigma)$} ناتعريفه وي۔ د~\LR{$\Struct M$} د تعريف له مخې، معنا
ئې دا ده چې پر ننوت~\LR{$w$} پيل شوے~\LR{$M$} له~\LR{$n$} ګامونو وروسته په
حالت~\LR{$q$} کښې دے او نښه~\LR{$\sigma$} لولي، او د حالت بدلون تابع ناتعريفه
ده؛ يعنې~\LR{$M$} درېدلے دے۔
\end{proof}
\olpseganchor{OLP-0271-B035}{end}

\olpseganchor{OLP-0272-B004}{start}
\olfileid{tur}{und}{uns}
\olsection{د پرېکړې مسئله ناحل ده}
\olpseganchor{OLP-0272-B004}{end}

\olpseganchor{OLP-0272-B005}{start}
\begin{thm}
\ollabel{thm:decision-prob}
د پرېکړې مسئله ناحل ده: داسې ټيورينګ ماشين~\LR{$D$} نشته چې پر هغې پټې
له پيلېدو وروسته چې د لومړۍ درجې منطق يو تړلے فارمول~\LR{$!B$} د ننوت
په توګه لري، \LR{$D$}~بالاخره ودرېږي او هله او يوازې هله~\LR{$1$} وباسي چې~\LR{$!B$}
معتبر وي، او که نۀ وي~\LR{$0$} وباسي۔
\end{thm}
\olpseganchor{OLP-0272-B005}{end}

\olpseganchor{OLP-0272-B006}{start}
\begin{proof}
فرض کړئ د پرېکړې مسئله حل کېدونکې وے، يعنې فرض کړئ چې يو ټيورينګ
ماشين~\LR{$D$} شته۔ بيا به مو د درېدنې مسئله په لاندې ډول حل کړې وے۔ يو
ټيورينګ ماشين~\LR{$E$} جوړوو چې د ټيورينګ ماشين~\LR{$M_e$} عدد~\LR{$e$} او
ننوت~\LR{$w$} ور کړل شي، اړوند تړلے فارمول~\LR{$!T(M_e, w) \lif !E(M_e,
w)$} محاسبه کړي او په داسې حال کښې ودرېږي چې پر پټه تر ټولو کيڼه
مربع لولي۔ بيا به ماشين~\LR{$E \concat D$} ته چې ننوتونه~\LR{$e$} او~\LR{$w$} ور
کړل شي، لومړے~\LR{$!T(M_e, w) \lif !E(M_e, w)$} محاسبه کړي او ورپسې به
د پرېکړې مسئلې ماشين~\LR{$D$} پر همدې ننوت وچلوي۔ \LR{$D$} هله او يوازې هله
له وت~\LR{$1$} سره درېږي چې~\LR{$!T(M_e, w) \lif !E(M_e, w)$} معتبر وي، او
که نۀ وي~\LR{$0$} وباسي۔ د~\olref[ver]{lem:halt-if-valid} او
\olref[ver]{lem:valid-if-halt} له مخې، \LR{$!T(M_e, w) \lif !E(M_e, w)$}
هله او يوازې هله معتبر دے چې~\LR{$M_e$} پر ننوت~\LR{$w$} ودرېږي۔ نو~\LR{$E
\concat D$} چې ننوتونه~\LR{$e$} او~\LR{$w$} ور کړل شي، هله او يوازې هله له
وت~\LR{$1$} سره درېږي چې~\LR{$M_e$} پر ننوت~\LR{$w$} ودرېږي، او که نۀ وي له
وت~\LR{$0$} سره درېږي۔ په بله وينا، \LR{$E \concat D$} به د درېدنې مسئله حل
کړې وے۔ خو د~\olref[hal]{thm:halting-problem} له مخې پوهېږو چې داسې
ټيورينګ ماشين شتون نۀ شي لرلے۔
\end{proof}
\olpseganchor{OLP-0272-B006}{end}

\olpseganchor{OLP-0272-B007}{start}
\begin{cor}\ollabel{cor:undecidable-sat}%
دا پرېکړه کېدونکې نۀ ده چې آيا د لومړۍ درجې منطق کوم خپل‌سری
تړلے فارمول د صدق وړ دے۔
\end{cor}
\olpseganchor{OLP-0272-B007}{end}

\olpseganchor{OLP-0272-B008}{start}
\begin{proof}
  فرض کړئ چې د صدق وړتيا يو ټيورينګ ماشين~\LR{$S$} پرېکړه کولے شوه۔ بيا
  به مو د پرېکړې مسئله په لاندې ډول حل کړې وے: يو
  تړلے فارمول~\LR{$!B$} د ننوت په توګه واخلئ او~\LR{$!B$} يوه مربع ښي لوري
  ته يوسئ۔ بېرته مربع~\LR{$1$} ته راشئ او نښه~\LR{$\lnot$} وليکئ۔
\olpseganchor{OLP-0272-B008}{end}

\olpseganchor{OLP-0272-B009}{start}
  اوس ټيورينګ ماشين~\LR{$S$} وچلوئ۔ دا بالاخره درېږي او پر پټه يا~\LR{$1$}
  وباسي (که~\LR{$\lnot !B$} د صدق وړ وي) يا~\LR{$0$} (که~\LR{$\lnot !B$} د صدق
  وړ نۀ وي)۔ که پر مربع~\LR{$1$} يوه~\LR{$\TMstroke$} وي، پاکه ئې کړئ؛ که
  مربع~\LR{$1$} تشه وي، يوه~\LR{$\TMstroke$} پرې وليکئ، او بيا ودرېږئ۔
\olpseganchor{OLP-0272-B009}{end}

\olpseganchor{OLP-0272-B010}{start}
  دا ټيورينګ ماشين تل درېږي، او وت ئې هله او يوازې هله~\LR{$1$} وي چې
  \LR{$\lnot !B$} د صدق وړ نۀ وي، او که نۀ وي~\LR{$0$} وي۔ څرنګه چې~\LR{$!B$} هله
  او يوازې هله معتبر دے چې~\LR{$\lnot !B$} د صدق وړ نۀ وي، ماشين هله او
  يوازې هله~\LR{$1$} وباسي چې~\LR{$!B$} معتبر وي، او که نۀ وي~\LR{$0$} وباسي؛
  يعنې د پرېکړې مسئله به حل کړي۔
\end{proof}
\olpseganchor{OLP-0272-B010}{end}

\olpseganchor{OLP-0272-B011}{start}
\begin{explain}
نو داسې ټيورينګ ماشين نشته چې دې پوښتنې ته تل سم ``هو'' يا ``نه''
ځواب ورکړي: ``آيا~\LR{$!B$} د لومړۍ درجې منطق يو معتبر تړلے فارمول
دے؟'' خو داسې ټيورينګ ماشين \emph{شته} چې تل سم ``هو'' ځواب
ورکوي---مګر که ځواب ``نه'' وي، هېڅ نۀ درېږي۔ دا د لومړۍ درجې منطق
د سموالي او بشپړتيا له قضيې، او له دې حقيقت څخه راځي چې
اشتقاقونه په اغېزمنه توګه شمېرل کېدے شي۔
\end{explain}
\olpseganchor{OLP-0272-B011}{end}

\olpseganchor{OLP-0272-B012}{start}
\begin{thm}
  \ollabel{thm:valid-ce}%
  د لومړۍ درجې د تړلي فارمولونه اعتبار نيمه د پرېکړې وړ دے: داسې
  ټيورينګ ماشين~\LR{$E$} شته چې پر هغې پټې له پيلېدو وروسته چې د لومړۍ
  درجې منطق يو تړلے فارمول~\LR{$!B$} د ننوت په توګه لري، \LR{$E$}~هله او يوازې
  هله بالاخره درېږي او~\LR{$1$} وباسي چې~\LR{$!B$} معتبر وي، او که نۀ وي نۀ
  درېږي۔
\end{thm}
\olpseganchor{OLP-0272-B012}{end}

\olpseganchor{OLP-0272-B013}{start}
\begin{proof}
  د لومړۍ درجې منطق ټول ممکن اشتقاقونه د يوه اغېزمن الګوريتم
  په وسيله يو په بل پسې توليدېدے شي۔ ماشين~\LR{$E$} همدا کار کوي، او کله
  چې داسې اشتقاق ومومي چې~\LR{$\Proves !B$} ښيي، له وت~\LR{$1$} سره
  درېږي۔ د سموالي د قضيې له مخې، که~\LR{$E$} له وت~\LR{$1$} سره درېږي، لامل
  ئې دا دے چې~\LR{$\Entails !B$}۔ د بشپړتيا د قضيې له مخې، که~\LR{$\Entails
  !B$} وي، داسې اشتقاق شته چې~\LR{$\Proves !B$} ښيي۔ څرنګه
  چې~\LR{$E$} ټول ممکن اشتقاقونه په منظم ډول توليدوي، بالاخره به
  داسې يو ومومي چې~\LR{$\Proves !B$} ښيي؛ نو بالاخره به له وت~\LR{$1$} سره
  ودرېږي۔
\end{proof}
\olpseganchor{OLP-0272-B013}{end}

\olpseganchor{OLP-0273-B004}{start}
\olfileid{tur}{und}{tra}
\olsection{د تراختن‌بروت قضيه}
\olpseganchor{OLP-0273-B004}{end}

\olpseganchor{OLP-0273-B005}{start}
\begin{explain}
  په~\olref[rep]{sec} کښې مو د يوه ټيورينګ ماشين~\LR{$M$} او ننوت
  رشته~\LR{$w$} دپاره تړلي فارمولونه~\LR{$!T(M,w)$} او~\LR{$!E(M,w)$} تعريف کړل۔
  بيا مو په~\olref[ver]{lem:valid-if-halt} او
  \olref[ver]{lem:halt-if-valid} کښې وښودل چې~\LR{$!T(M,w) \lif
  !E(M,w)$} هله او يوازې هله معتبر دے چې~\LR{$M$} پر ننوت~\LR{$w$} له پيلېدو
  وروسته بالاخره ودرېږي۔ څرنګه چې د درېدنې مسئله پرېکړه کېدونکې نۀ
  ده، له دې څخه پايله اخلو چې د لومړۍ درجې منطق د تړلي فارمولونه
  اعتبار او د صدق وړتيا پرېکړه کېدونکي نۀ دي
  (\cref{tur:und:uns:thm:decision-prob,tur:und:uns:cor:undecidable-sat})۔
\olpseganchor{OLP-0273-B005}{end}

\olpseganchor{OLP-0273-B006}{start}
  خو د جملو اعتبار او د صدق وړتيا د خپلو‌سرو
  جوړښتونه دپاره تعريف شوي، که هغه متناهي وي او که نامتناهي۔
  ښايي تاسو ګومان وکړئ چې دا پرېکړه اسانه ده چې آيا يو
  تړلے فارمول په يوه متناهي جوړښت کښې د صدق وړ دے (يا په
  ټولو متناهي جوړښتونه کښې معتبر دے)۔ د پرېکړې مسئلې د ناحلي
  ثبوت داسې برابرولے شو چې وښيي خبره داسې نۀ ده۔
\olpseganchor{OLP-0273-B006}{end}

\olpseganchor{OLP-0273-B007}{start}
  لومړے، که د~\olref[ver]{lem:halt-if-valid} ثبوت ته بېرته ورشئ،
  وبه وينئ چې هلته مو د~\LR{$!T(M,w)$} داسې مدل~\LR{$\Struct M$} جوړ کړ چې
  بالکل دا بيانوي چې ماشين~\LR{$M$} پر ننوت~\LR{$w$} له پيلېدو وروسته څه
  کوي۔ د هغه مدل ډومېن~\LR{$\Nat$} و، يعنې نامتناهي و۔ خو که~\LR{$M$} په
  رښتيا پر ننوت~\LR{$w$} ودرېږي، په همدې ډول يو متناهي مدل~\LR{$\Struct M'$}
  جوړولے شو۔ فرض کړئ~\LR{$M$} پر ننوت~\LR{$w$} له پيلېدو وروسته له~\LR{$k$}
  ګامونو وروسته درېږي۔ د ډومېن~\LR{$\Domain{M'}$} دپاره سټ~\LR{$\{0,
  \dots, n\}$} واخلئ۔ \LR{$n$} دې د~\LR{$k+1$} او د~\LR{$w$} د اوږدوالي له دواړو څخه لوے وي، او داسې وټاکئ:
  \[
    \Assign{\prime}{M'}(x) =
    \begin{cases}
      x + 1 &\text{\RL{که \LR{$x < n$} وي}}\\
      n &\text{\RL{که نۀ وي،}}
    \end{cases}
  \]
  او~\LR{$\tuple{x,y} \in \Assign{<}{M'}$} هله او يوازې هله چې~\LR{$x < y$}
  يا~\LR{$x = y = n$} وي۔ له دې پرته~\LR{$\Struct{M'}$} بالکل د~\LR{$\Struct{M}$}
  په شان تعريفېږي۔ د~\LR{$\Struct{M'}$} د تعريف له مخې، لکه د
  \olref[ver]{lem:halt-if-valid} په ثبوت کښې،
  \LR{$\Sat{M'}{!T(M,w)}$}۔ او څرنګه چې مو فرض کړې ده چې~\LR{$M$} پر
  ننوت~\LR{$w$} درېږي، \LR{$\Sat{M'}{!E(M,w)}$}۔ نو~\LR{$\Struct{M'}$} د~\LR{$!T(M,w)
  \land !E(M,w)$} يو متناهي مدل دے (پام وکړئ چې~\LR{$\lif$} مو
  په~\LR{$\land$} بدل کړے دے)۔
\olpseganchor{OLP-0273-B007}{end}

\olpseganchor{OLP-0273-B008}{start}
  موږ د ثبوت نيمې لارې ته رسېدلي يو: ښودلې مو ده چې که~\LR{$M$} پر
  ننوت~\LR{$w$} ودرېږي، نو~\LR{$!T(M,w) \land !E(M,w)$} يو متناهي مدل لري۔
  له بده مرغه د دې نقيض عکس سم نۀ دے؛ يعنې داسې ټيورينګ ماشينونه
  شته چې پر کوم ننوت~\LR{$w$} نۀ درېږي، خو~\LR{$!T(M,w) \land !E(M,w)$} بيا
  هم يو متناهي مدل لري۔ د بېلګې په توګه هغه ماشين~\LR{$M$} په پام کښې
  ونيسئ چې يوازې يو حالت~\LR{$q_0$} او دا لارښوونه لري:
  \LR{$\delta(q_0,\TMblank) = \tuple{q_0,\TMblank,\TMstay}$}۔ دا ماشين
  پر تش ننوت~\LR{$w = \emptyseq$} له پيلېدو وروسته هېڅکله نۀ درېږي: په
  يوه نامتناهي کړۍ کښې دے، خو پټه نۀ بدلوي او سر نۀ خوځوي۔ ټول
  تشکيلونه يو شان دي (هماغه حالت، هماغه د سر موقعيت، هماغه د پټې
  منځپانګه)۔ داسې متناهي جوړښت~\LR{$\Struct{M''}$} تعريفولے شو
  چې~\LR{$!T(M,\emptyseq) \land !E(M,\emptyseq)$} پوره کړي (تمرين)۔ خو
  \LR{$!T(M,w)$} په مناسبه توګه بدلولے شو، څو داسې جوړښتونه ترې
  بهر کړو۔
\olpseganchor{OLP-0273-B008}{end}

\olpseganchor{OLP-0273-B009}{start}
  \begin{prob}
    فرض کړئ~\LR{$M$} داسې ټيورينګ ماشين دے چې يوازې حالت~\LR{$q_0$} او يوازې
    لارښوونه~\LR{$\delta(q_0,\TMblank) =
    \tuple{q_0,\TMblank,\TMstay}$} لري۔ داسې وټاکئ:
    \LR{$\Domain{M''} = \{0, 1, 2\}$}، \LR{$\Assign{\prime}{M''}(0) =
    \Assign{\prime}{M''}(1) = 1$} او~\LR{$\Assign{\prime}{M''}(2) = 2$}،
    او~\LR{$\Assign{<}{M''} = \{\tuple{0,1}, \tuple{1,1},
    \tuple{2,2}\}$}۔ \LR{$\Assign{\Obj Q_{q_0}}{M''}$}، \LR{$\Assign{\Obj
    S_{\TMblank}}{M''}$}، او~\LR{$\Assign{\Obj S_{\TMendtape}}{M''}$}
    داسې تعريف کړئ چې~\LR{$!T(M,\emptyseq)$} او~\LR{$!E(M,\emptyseq)$} رښتوني
    شي، او څرګنده کړئ چې ولې رښتوني دي۔ اشاره: پام وکړئ چې
    \LR{$\delta(q_0, \TMendtape)$} ناتعريفه ده۔ ډاډ ترلاسه کړئ چې
    \begin{align*}
    & \Obj Q_{q_0}(\num{1}, \num{n}) \land \Obj S_{\TMendtape}(\num{0}, \num{n})
      \land
      \lforall[x][(\num{0} < x \lif \Obj S_\TMblank(x, \num{n}))] \text{\RL{\quad د هر \LR{$n \in \Nat$} دپاره}}\\
    & \lexists[y][(\Obj Q_{q_0}(\num{0}, y) \land \Obj S_{\TMendtape}(\num{0}, y))]
    \end{align*}
    دواړه په~\LR{$\Struct{M''}$} کښې رښتوني وي۔
  \end{prob}
\end{explain}
\olpseganchor{OLP-0273-B009}{end}

\olpseganchor{OLP-0273-B010}{start}
هغه تړلي فارمولونه په پام کښې ونيسئ چې د ټيورينګ ماشين~\LR{$M$} عمل پر
ننوت~\LR{$w = \sigma_{i_1}\dots\sigma_{i_k}$} بيانوي:
\begin{enumerate}
\item هغه اکسيومونه چې اعداد او~\LR{$<$} بيانوي (بالکل لکه په
\olref[rep]{sec} کښې د~\LR{$!T(M,w)$} په تعريف کښې)۔
\item هغه اکسيومونه چې د ننوت تشکيل بيانوي: بالکل لکه د~\LR{$!T(M,w)$}
په تعريف کښې۔
\item هغه اکسيومونه چې له يوه تشکيل څخه بل ته د حالت بدلون بيانوي:
\olpseganchor{OLP-0273-B010}{end}

\olpseganchor{OLP-0273-B011}{start}
د لاندې برخې دپاره، \LR{$!A(x, y)$} دې د پخوا په شان وي، او داسې دې وي:
\[
  !B(y) \ident \lforall[x][(x < y \lif \eq/[x][y])].
\]
\begin{enumerate}
\item \ollabel{rep-right} د هرې لارښوونې~\LR{$\delta(q_i, \sigma) =
  \tuple{q_j, \sigma', \TMright}$} دپاره دا تړلے فارمول:
\begin{align*}
& \lforall[x][\lforall[y][(
   (\Obj Q_{q_i}(x, y) \land \Obj S_{\sigma}(x, y)) \lif {}]] \\
&\qquad   (\Obj Q_{q_j}(x', y') \land \Obj S_{\sigma'}(x, y') \land
!A(x, y) \land !B(y')))
\end{align*}
\item \ollabel{rep-left} د هرې لارښوونې~\LR{$\delta(q_i, \sigma) =
  \tuple{q_j, \sigma', \TMleft}$} دپاره دا تړلے فارمول:
\begin{align*}
& \lforall[x][\lforall[y][
    ((\Obj Q_{q_i}(x', y) \land \Obj S_{\sigma}(x', y)) \lif {}]]\\
& \qquad   (\Obj Q_{q_j}(x, y') \land \Obj S_{\sigma'}(x', y') \land
!A(x', y) \land !B(y'))) \land {}\\
& \lforall[y][((\Obj Q_{q_i}(\num{0}, y) \land \Obj S_{\sigma}(\num{0},
    y)) \lif {}]\\
& \qquad (\Obj Q_{q_j}(\num{0}, y') \land \Obj S_{\sigma'}(\num{0},
  y') \land !A(\num{0}, y) \land !B(y')))
\end{align*}
\item \ollabel{rep-stay} د هرې لارښوونې~\LR{$\delta(q_i, \sigma) =
  \tuple{q_j, \sigma', \TMstay}$} دپاره دا تړلے فارمول:
\begin{align*}
& \lforall[x][\lforall[y][(
   (\Obj Q_{q_i}(x, y) \land \Obj S_{\sigma}(x, y)) \lif {}]] \\
&\qquad (\Obj Q_{q_j}(x, y') \land \Obj S_{\sigma'}(x, y') \land
!A(x, y) \land !B(y')))
\end{align*}
\end{enumerate}
لکه چې وينئ، هغه تړلي فارمولونه چې د~\LR{$M$} حالت-بدلونونه بيانوي د
\LR{$!T(M,w)$} له اړوندو تړلے فارمول سره يو شان دي، يوازې په پای کښې
\LR{$!B(y')$} ور زياتوو۔ \LR{$!B(y')$} دا تضمينوي چې د ``بل'' تشکيل عدد~\LR{$y'$}
له ټولو مخکښېنيو عددونو~\LR{$\Obj 0$}، \LR{$\Obj 0'$}، \dots، څخه بېل وي۔
% \item Sentences that express consistency conditions:
% \begin{enumerate}
%   \item\ollabel{rep-con-symb}%
%   تړلے فارمول that says that no square can have more than one
%   symbol on it at any given time:
%   \[\lforall[x][\lforall[y][(\Obj S_{\sigma}(x, y) \lif \lnot \Obj
%   S_{\sigma'}(x, y))]],\]
%   for every pair $\sigma \neq \sigma'$ of tape symbols of~$T$.
%   \item\ollabel{rep-con-state}%
%   تړلے فارمول that says that $M$ cannot be in more than one state
%   at any given time:
%   \[\lforall[x][\lforall[y][(\Obj Q_{q}(x, y) \lif
%   \lforall[z][\lnot \Obj
%   Q_{q'}(z, y))]],\]
%   for every pair $q \neq q'$ of states of~$T$.
%   \item\ollabel{rep-con-tape}%
%   تړلے فارمول that says that $M$ cannot be on more than one tape
%   square at any given time:
%   \[\lforall[x][\lforall[y][(\Obj Q_{q}(x, y) \lif
%   \lforall[z][\eq/[z][y]\lif \lnot \Obj
%   Q_{q'}(z, y))]],\]
%   for every pair $q$, $q'$ of states of~$T$.
% \end{enumerate}
\end{enumerate}
\LR{$!T'(M, w)$} د ټيورينګ ماشين~\LR{$M$} او ننوت~\LR{$w$} دپاره د پورتنيو ټولو
تړلي فارمولونه عطف وي۔
\olpseganchor{OLP-0273-B011}{end}

\olpseganchor{OLP-0273-B012}{start}
\begin{lem}\ollabel{lem:halts-sat}
  که~\LR{$M$} پر ننوت~\LR{$w$} له پيلېدو وروسته ودرېږي، نو~\LR{$!T'(M,w) \land
  !E(M,w)$} يو متناهي مدل لري۔
\end{lem}
\olpseganchor{OLP-0273-B012}{end}

\olpseganchor{OLP-0273-B013}{start}
\begin{proof}
  \LR{$\Struct{M'}$} د~\olref[ver]{lem:halt-if-valid} په ثبوت کښې د
  جوړښت په شان واخلئ، خو دا توپيرونه ولري:
  \begin{align*}
    \Domain{M'} & = \{0, \dots, n\},\\
    \Assign{\prime}{M'}(x) & =
    \begin{cases}
      x + 1 &\text{\RL{که \LR{$x < n$} وي}}\\
      n &\text{\RL{که نۀ وي،}}
    \end{cases} \\
    \tuple{x,y} \in \Assign{<}{M'} &\text{\RL{هله او يوازې هله چې \LR{$x < y$} يا \LR{$x = y = n$} وي،}}
  \end{align*}
  چې~\LR{$n = \max(k+2,\len{w}+1)$} او~\LR{$k$} تر ټولو کوچنی داسې عدد دے چې~\LR{$M$}
  پر ننوت~\LR{$w$} له پيلېدو وروسته له~\LR{$k$} ګامونو وروسته درېدلے وي۔ د پټې سر په لومړۍ مربع پيل کوي او په هر ګام کښې تر يوې مربع زيات ښي لوري ته نۀ ځي۔ نو د اجرا ټول وختونه او د سر ټول ځايونه د ډومېن له وروستۍ شمېرې څخه کوچني دي۔ د اجرا پر مهال د حالت او پټې محمولونه له واقعي تشکيلونو اخلو، او له درېدنې وروسته د حالت ټول محمولونه دروغ ګڼو۔ د هرې واقع شوې لارښوونې لپاره تالي، د پټې د نښو ساتل او د نوې شېبې د عدد بېلوالے پر خپل ځاے پوره کېږي۔ په وروستي تشکيل کښې د درېدنې شرط پوره دے۔ نو \LR{$\Sat{M'}{!T'(M,w) \land !E(M,w)}$}۔
\end{proof}
\olpseganchor{OLP-0273-B013}{end}

\olpseganchor{OLP-0273-B014}{start}
\begin{prob}
  د~\olref[tur][und][tra]{lem:halts-sat} ثبوت په دې ښودلو بشپړ کړئ
  چې~\LR{$\Sat{M'}{!T'(M,w) \land !E(M,w)}$}۔
\end{prob}
\olpseganchor{OLP-0273-B014}{end}

\olpseganchor{OLP-0273-B015}{start}
\begin{lem}\ollabel{lem:sat-halts}
  که~\LR{$!T'(M,w) \land !E(M,w)$} يو متناهي مدل ولري، نو~\LR{$M$} پر
  ننوت~\LR{$w$} له پيلېدو وروسته درېږي۔
\end{lem}
\olpseganchor{OLP-0273-B015}{end}

\olpseganchor{OLP-0273-B016}{start}
\begin{proof}
  نقيض عکس ښيو۔ فرض کړئ~\LR{$M$} پر ننوت~\LR{$w$} له پيلېدو وروسته نۀ
  درېږي۔ که~\LR{$!T'(M,w) \land !E(M,w)$} هېڅ مدل ونۀ لري، خبره بشپړه
  ده۔ نو فرض کړئ~\LR{$\Struct{M}$} د~\LR{$!T'(M,w) \land !E(M, w)$} يو مدل
  دے۔ بايد وښيو چې دا متناهي کېدے نۀ شي۔
\olpseganchor{OLP-0273-B016}{end}

\olpseganchor{OLP-0273-B017}{start}
  د~\olref[ver]{lem:config} استقرا د صفرم تشکيل بيان له پيليزو اکسيومونو اخلي۔ که~\LR{$M$} پر ننوت~\LR{$w$} د~\LR{$n$}-م تشکيل تر رسېدو مخکې نۀ وي درېدلے او~\LR{$n>0$} وي، د دغه تشکيل نه مخکې وروستۍ لارښوونه د تشکيل له نوي بيان سره د وخت د عدد بېلوالے هم ورکوي۔ څرنګه چې~\LR{$M$} پر ننوت~\LR{$w$} نۀ درېږي، د هر مثبت~\LR{$n \in \Nat$} دپاره
  \LR{$!T'(M,w) \Entails !C(M, w, n) \land !B(\num{n})$}۔ پام وکړئ چې
  د~\olref[rep]{prop:mlessk} له مخې، د ټولو~\LR{$k < n$} دپاره~\LR{$!T'(M,w)
  \Entails \num{k} < \num{n}$}۔ همدارنګه~\LR{$!B(\num{n}) \Entails
  \num{k} < \num{n} \lif \eq/[\num{k}][\num{n}]$}۔ نو د ټولو~\LR{$k <
  n$} دپاره~\LR{$\Sat{M}{\eq/[\num{k}][\num{n}]}$}؛ يعنې نامتناهي ډېر
  ترمونه~\LR{$\num{k}$} بايد په~\LR{$\Struct{M}$} کښې ټول بېلابېل ارزښتونه
  ولري۔ خو دا غوښتنه کوي چې~\LR{$\Domain{M}$} نامتناهي وي؛ نو
  \LR{$\Struct{M}$} د~\LR{$!T'(M,w) \land !E(M, w)$} متناهي مدل نۀ شي
  کېدے۔
\end{proof}
\olpseganchor{OLP-0273-B017}{end}

\olpseganchor{OLP-0273-B018}{start}
\begin{prob}
  د~\olref[tur][und][tra]{lem:sat-halts} ثبوت په دې ښودلو بشپړ کړئ
  چې که~\LR{$n>0$} وي او~\LR{$M$} پر ننوت~\LR{$w$} د~\LR{$n$}-م تشکيل تر رسېدو مخکې نۀ وي درېدلے، نو~\LR{$!T'(M,w) \Entails !B(\num{n})$}۔
\end{prob}
\olpseganchor{OLP-0273-B018}{end}

\olpseganchor{OLP-0273-B019}{start}
\begin{thm}[د تراختن‌بروت قضيه]
  \ollabel{thm:trakhtenbrodt}
  دا پرېکړه کېدونکې نۀ ده چې آيا د لومړۍ درجې منطق کوم خپل‌سری
  تړلے فارمول يو متناهي مدل لري (يعنې په يوه متناهي جوړښت کښې د
  صدق وړ دے)۔
\end{thm}
\olpseganchor{OLP-0273-B019}{end}

\olpseganchor{OLP-0273-B020}{start}
\begin{proof}
  فرض کړئ داسې ټيورينګ ماشين~\LR{$F$} وے چې په يوه متناهي جوړښت کښې د
  صدق وړتيا پرېکړه کوي۔ بيا به مو د هر ټيورينګ ماشين~\LR{$M$} او
  ننوت~\LR{$w$} دپاره جمله~\LR{$!T'(M,w) \land !E(M,w)$} محاسبه کړې وے، او
  د~\LR{$F$} په وسيله به مو پرېکړه کړې وے چې آيا دا يو متناهي مدل لري۔
  د~\cref{tur:und:tra:lem:halts-sat,tur:und:tra:lem:sat-halts} له
  مخې، دا هله او يوازې هله متناهي مدل لري چې~\LR{$M$} پر ننوت~\LR{$w$} له
  پيلېدو وروسته ودرېږي۔ نو د~\LR{$F$} په وسيله به مو د درېدنې مسئله حل
  کړې وے، حال دا چې پوهېږو دا ناحل ده۔
\end{proof}
\olpseganchor{OLP-0273-B020}{end}

\olpseganchor{OLP-0273-B021}{start}
\begin{cor}\ollabel{cor:fproof-incomp}%
  داسې هېڅ اشتقاق نظام چې ثبوتونه ئې په اغېزمنه توګه شمېرل کېدے شي، د \emph{متناهي}
  اعتبار دپاره سم او بشپړ نۀ شي کېدے۔ يعنې داسې اشتقاق نظام چې
  \LR{$\Proves !B$} هله او يوازې هله ولري چې د هرې متناهي
  جوړښت~\LR{$\Struct{M}$} دپاره~\LR{$\Sat{M}{!B}$} وي۔
\end{cor}
\olpseganchor{OLP-0273-B021}{end}

\olpseganchor{OLP-0273-B022}{start}
\begin{proof}
  د هرې جملې د متناهي مدلونو لټون او د هغې د نفي د ثبوت لټون يو په بل پسې يو يو ګام پر مخ وړو۔ د لومړي لټون دپاره د جملې د متناهي ژبې ټول متناهي جوړښتونه د ډومېن د اندازې په ترتيب شمېرو، او په هر جوړښت کښې د جملې صدق فيصله کوو۔ که جمله متناهي مدل ولري، لومړے لټون بريالے کېږي۔ که ئې ونۀ لري، د فرض شوي بشپړ نظام له مخې د هغې نفي يو ثبوت لري او دويم لټون به بريالے شي۔ موندل شوے مدل په خپله د شتون شاهد دے، او د نظام سموالى د دويم لټون منفي ځواب تضمينوي۔ په دې توګه به مو د متناهي مدل شتون پرېکړے شوے ؤ، چې د تراختن‌بروت له قضيې سره په ټکر کښې دے۔
\end{proof}
\olpseganchor{OLP-0273-B022}{end}

\olpseganchor{OLP-0273-B023}{start}
\begin{prob}
  \olref[tur][und][tra]{cor:fproof-incomp} ثابت کړئ۔ پام وکړئ چې~\LR{$!B$}
  په هره متناهي جوړښت کښې هله او يوازې هله پوره کېږي چې
  \LR{$\lnot !B$} په يوه متناهي جوړښت کښې د صدق وړ نۀ وي۔ څرګنده کړئ چې
  ولې په يوه متناهي جوړښت کښې د صدق وړتيا د
  \olref[tur][und][uns]{thm:valid-ce} په معنا نيمه د پرېکړې وړ ده۔
  له دې څخه دا استدلال وکړئ چې که د متناهي اعتبار دپاره داسې
  اشتقاق نظام وے چې ثبوتونه ئې په اغېزمنه توګه شمېرل کېدے شي، نو په يوه متناهي جوړښت کښې د صدق وړتيا
  به پرېکړه کېدونکې وه۔
\end{prob}
\olpseganchor{OLP-0273-B023}{end}

\olpseganchor{OLP-0274-B004}{start}
\olpart{inc}{ناتکميلي}
\olpseganchor{OLP-0274-B004}{end}

\olpseganchor{OLP-0274-B005}{start}
\begin{editorial}
  د دې برخې مواد د ناتکميلۍ قضيې رانغاړي۔ دا د لومړۍ درجې منطق د
  برخو (په ځانګړي ډول د ثبوت نظام) او د محاسبه‌پوهې په برخه کښې د
  بازګشتي تابعو پر موادو تکيه کوي۔ دا د جېرېمي اېوېګاډ پر يادښتونو
  ولاړ دي چې ريچرډ زاک سمونونه پکې کړي دي۔
\end{editorial}
\olpseganchor{OLP-0274-B005}{end}

\olpseganchor{OLP-0275-B004}{start}
\olchapter{inc}{int}{ناتکميلۍ ته پېژندنه}
\olpseganchor{OLP-0275-B004}{end}

\olpseganchor{OLP-0276-B004}{start}
\olfileid{inc}{int}{bgr}
\olpseganchor{OLP-0276-B004}{end}

\olpseganchor{OLP-0276-B005}{start}
\olsection{تاريخي شاليد}
\olpseganchor{OLP-0276-B005}{end}

\olpseganchor{OLP-0276-B006}{start}
په دې برخه کښې به پر هغو تاريخي پرمختګونو لنډ بحث وکړو چې د
ناتکميلۍ د قضيو د چاپېريال په پوهېدو کښې مرسته کوي۔ په ځانګړي ډول به
د رياضياتي منطق د تاريخ يوه ډېره لنډه ټوليزه کتنه وړاندې کړو، او بيا
به د رياضياتو د بنسټونو د تاريخ په باب څو خبرې وکړو۔
\olpseganchor{OLP-0276-B006}{end}

\olpseganchor{OLP-0276-B007}{start}
\begin{digress}
  د «رياضياتي منطق» غونډله دوه‌معناييزه ده۔ د «رياضياتي» ويے د موضوع
  د بيان په توګه تعبيرېدے شي، لکه «د رياضياتو منطق»، چې د رياضياتي
  استدلال اصول ښيي؛ يا د لارو د بيان په توګه، لکه «د منطق رياضيات»،
  چې د استدلال د اصولو رياضياتي مطالعه ښيي۔ لاندې روايت رياضياتي منطق
  په دواړو معناوو، او ډېر ځله په يوه وخت کښې، رانغاړي۔
\end{digress}
\olpseganchor{OLP-0276-B007}{end}

\olpseganchor{OLP-0276-B008}{start}
د منطق مطالعه په اصل کښې له ارسطو څخه پيل شوه، چې نږدې په
384--322 \textsc{ق‌م} کښې ژوندے و۔ د هغه \emph{مقولات}،
\emph{لومړۍ تحليلات} او \emph{دويمې تحليلات} د علمي استدلال د اصولو
منظمې مطالعې رانغاړي، چې د قياس بشپړه او منظمه مطالعه هم پکې شامله ده۔
\olpseganchor{OLP-0276-B008}{end}

\olpseganchor{OLP-0276-B009}{start}
د ارسطو منطق د منځنيو پېړيو په اوږدو کښې پر مدرسي فلسفې واکمن و؛ ان
تر اتلسمې پېړۍ پورې کانټ ټينګار کاوه چې د ارسطو منطق بشپړ دے او
سمون ته اړتيا نۀ لري۔ خو د قياس تيوري ډېره محدوده ده او د رياضياتي
استدلال له تر ټولو سطحي اړخونو پرته نور څه نۀ شي مدل کولے۔ يوه پېړۍ
تر دې وړاندې، د نيوټن هم‌مهالي لايبنيڅ د منطقي استدلال دپاره د يوه
بشپړ «حساب» تصور وکړ او د داسې حساب د جوړولو پر لور ئې ځينې لومړني
ګامونه واخيستل؛ په اصل کښې ئې د قضييز منطق يوه بڼه بيان کړه۔
\olpseganchor{OLP-0276-B009}{end}

\olpseganchor{OLP-0276-B010}{start}
نولسمه پېړۍ د منطق دپاره يو بنسټيز بدلون و۔ په 1854 کښې جورج بول
\emph{د فکر قانونونه} وليکل، چې د قضييز منطق بشپړه الجبري مطالعه
پکې وه او له اوسنيو وړاندېونو څخه ډېره لرې نۀ ده۔ په 1879 کښې ګوټلوب
فرېګې خپله \emph{بېګريفشريفټ} (مفهومي ليکدود) خپره کړه، چې قضييز
منطق د کميت‌ټاکونکو او اړيکو په ورزياتولو غځوي او په دې توګه د لومړۍ
درجې منطق هم رانغاړي۔ په حقيقت کښې د فرېګې منطقي نظامونو د لوړې درجې
منطق او نور څه هم رانغاړل۔ فرېګې په خپلو \emph{د حساب بنسټيزو
قانونونو} کښې هڅه وکړه وښيي چې ټول حساب د هغه په بېګريفشريفټ کښې له
سوچه منطقي فرضيو څخه استنتاجېدے شي۔ له بده مرغه دا فرضيې ناسازګارې
وختې، لکه رسل چې په 1902 کښې وښودل۔ خو که ناسازګار بديهي اصل څنګ ته
کړو، فرېګې نږدې په يوازې ځان اوسنی منطق اختراع کړ، چې يوه حيرانوونکې
لاسته راوړنه ده۔ له بول وروسته الجبري فکر لرونکو پوهانو، لکه پيرس او
شروډر، کميتي منطق په خپلواکه توګه هم جوړ کړ۔
\olpseganchor{OLP-0276-B010}{end}

\olpseganchor{OLP-0276-B011}{start}
اوس راځئ د رياضياتو د بنسټونو پرمختګونو ته مخه کړو۔ البته څرنګه چې
منطق په رياضياتو کښې مهم رول لري، له پاس بيان شويو پرمختګونو سره ئې
ډېر متقابل اغېز شته۔ د بېلګې په توګه، فرېګې خپل منطق په دې څرګنده
موخه جوړ کړ چې وښيي ټول رياضيات يوازې پر هغه منطقي چوکاټ ولاړېدے شي؛
په ځانګړي ډول غوښتل ئې وښيي چې رياضيات، د کانټ د نظر پر خلاف، له
تجربې وړاندې \emph{تحليلي} حقايق دي، نه له تجربې وړاندې
\emph{ترکيبي} حقايق۔
\olpseganchor{OLP-0276-B011}{end}

\olpseganchor{OLP-0276-B012}{start}
ډېر خلک د رياضياتو اصلي زېږون له يونانيانو سره تړي۔ د اقليدس
\emph{اصول}، چې نږدې 300 \textsc{ق‌م} کښې ليکل شوي، لا له وړاندې د
يوناني رياضياتو يوه پخه بېلګه ده او پر کره‌والي او دقت ټينګار کوي۔
د اقليدس د \emph{اصولو} تعريفونه او ثبوتونه د نن ورځې د لېسې د هندسې
په درسي کتابونو کښې نږدې هماغسې پاتې دي (تر هغه ځايه چې لا په لېسو
کښې هندسه تدريسېږي)۔ د رياضياتي استدلال دا مدل نه يوازې په رياضياتو
بلکې د فلسفې په څانګو کښې هم د کره استدلال بېلګه ګڼل شوې ده۔
(سپينوزا ان اخلاقي او ديني استدلالونه د اقليدسي سبک له مخې وړاندې
کړل، چې ليدل ئې عجيبه دي!)
\olpseganchor{OLP-0276-B012}{end}

\olpseganchor{OLP-0276-B013}{start}
نيوټن او لايبنيڅ په اوولسمه پېړۍ کښې کالکولس اختراع کړ۔ (د لومړيتوب
يوه سخته شخړه پېړۍ پېړۍ روانه وه، خو اوس زياتره پوهان مني چې دواړه
پرمختګونه تر ډېره خپلواک وو۔) په کالکولس کښې، د بېلګې په توګه، د
نامتناهي وړو مقدارونو د نامتناهي جمعو په باب استدلال راځي؛ دغو
ځانګړتياوو د بيشپ برکلي نيوکې راوپارولې، چې استدلال ئې کاوه پر خدای
باور د هغه د وخت تر رياضياتو لږ معقول نۀ دے۔ په اتلسمه پېړۍ کښې د
کالکولس لارې په پراخه توګه وکارول شوې؛ مثلاً ليونارډ اويلر د
نامتناهي جمعو محاسبې په حيرانوونکو پايلو وکارولې۔
\olpseganchor{OLP-0276-B013}{end}

\olpseganchor{OLP-0276-B014}{start}
په نولسمه پېړۍ کښې رياضيپوهانو هڅه وکړه د کالکولس د لا ټينګ بنسټ په
جوړولو د برکلي نيوکو ته ځواب ووايي۔ د کوشي، وايرشتراس، بولتسانو او
نورو هڅو موږ د حد، پيوستون، تفاضل او تکامل اوسنيو تعريفونو ته
ورسولو، چې د «اېپسيلونونو او ډېلټاګانو» په اصطلاحاتو کښې دي؛ يعنې
نامتناهي وړو مقدارونو ته هېڅ مراجعه نۀ کوي۔ د پېړۍ په وروستيو کښې
رياضيپوهانو هڅه وکړه لا وړاندې لاړ شي او د حقيقي عددونو په ګډون د
کالکولس ټول اړخونه د طبيعي عددونو په اصطلاحاتو کښې بيان کړي۔
(کرونېکر: «خدای صحيح عددونه پيدا کړل، نور هر څه د انسان کار دے۔»)
په 1872 کښې ډېدېکينډ «پيوستون او غيرناطق عددونه» وليکل، چې وئې ښودل
څنګه حقيقي عددونه د ناطقو عددونو د سټونو په توګه «جوړېدے» شي (او
ناطق عددونه، لکه چې پوهېږئ، د طبيعي عددونو د جوړو په توګه ليدل کېدے
شي)؛ په 1888 کښې ئې «Was sind und was sollen die Zahlen» (نږدې معنا
ئې «طبيعي عددونه څه دي، او څه بايد وي؟») وليکل، چې موخه ئې په سوچه
«منطقي» اصطلاحاتو کښې د طبيعي عددونو بيان و۔ په 1887 کښې کرونېکر
«\"Uber den Zahlbegriff» («د عدد د مفهوم په باب») وليکل او د صحيح
عددونو په اصطلاحاتو کښې ئې د ټولو رياضياتي څيزونو د تمثيل خبره وکړه؛
په 1889 کښې جوزېپې پيانو د طبيعي عددونو صوري، نښيز بديهي اصول ورکړل۔
\olpseganchor{OLP-0276-B014}{end}

\olpseganchor{OLP-0276-B015}{start}
د نولسمې پېړۍ پاى د نامتناهي په چلند کښې يو نوی زړورتوب هم راوړ۔
تر دې وړاندې له نامتناهي څيزونو او جوړښتونو (لکه د طبيعي عددونو له
سټ) سره ډېر په احتياط چلند کېده؛ «نامتناهي ډېر» د «څومره چې وغواړئ
هومره» په معنا، او «په حد کښې ورنږدې کېږي» د «څومره چې وغواړئ هومره
نږدې کېږي» په معنا ګڼل کېده۔ خو ګيورګ کانتور وښودل چې نامتناهي په
خپله ظاهري معنا اخيستل کېدے شي۔ د کانتور، ډېدېکينډ او نورو کار د
رياضياتو هغه عمومي سټ‌تيوريکي پوهاوی رامنځته کړ چې اوس په پراخه توګه
منل شوے دے۔
\olpseganchor{OLP-0276-B015}{end}

\olpseganchor{OLP-0276-B016}{start}
دا مو په منطق او بنسټونو کښې د شلمې پېړۍ پرمختګونو ته رسوي۔ په 1902
کښې رسل د فرېګې په منطقي نظام کښې پاراډوکس وموند۔ په 1904 کښې
څېر مېلو د کانتور د ښه‌ترتيب اصل د تش په نوم «د انتخاب بديهي اصل» په
کارولو ثابت کړ؛ د دې بديهي اصل د رواوالي په باب ډېر بحث وشو۔ د 1910
او 1913 تر منځ د رسل او وايټ‌هېډ د \emph{پرېنسيپيا مېتېماتيکا} درې
ټوکونه خپاره شول، چې د رياضياتو پر منطقي بنسټونو د درولو د فرېګې
پروګرام ئې وغځاوه۔ له بده مرغه رسل او وايټ‌هېډ اړ شول دوه داسې اصول
ومني چې د سوچه منطقي اصولو په توګه ئې توجيه ستونزمنه ښکارېده: د
نامتناهي والي بديهي اصل او د «راکمېدنې» بديهي اصل۔ په 1900مو کلونو
کښې پوانکاره په رياضياتو کښې د «غيرپريديکاتي تعريفونو» پر کارولو
نيوکه وکړه، او په 1910مو کلونو کښې بروور د ټولو رياضياتو د بيا بنسټيز
کولو دپاره د داسې «شهودي» بنسټ وړانديز پيل کړ چې د خارج وسط اصل
(\LR{$!A \lor \lnot !A$}) پکې نۀ کارېده۔
\olpseganchor{OLP-0276-B016}{end}

\olpseganchor{OLP-0276-B017}{start}
په رښتيا هم عجيب وختونه وو!{} منطق ته د ټولو رياضياتو د راکمولو
پروګرام اوس «منطقپالنه» بلل کېږي او د پاسنيو ستونزو له امله عموماً
ناکام ګڼل کېږي۔ د شهودي ذهني جوړونو په اصطلاحاتو کښې د رياضياتو د
جوړولو پروګرام «شهوديت» بلل کېږي او پر ورځنيو رياضياتو ډېر سخت
محدوديتونه ايښوونکی ګڼل کېږي۔ د پېړۍ د بدلېدو شاوخوا ډېوېډ هېلبرټ،
چې د هر وخت له تر ټولو اغېزمنو رياضيپوهانو څخه و، د کانتور او
ډېدېکينډ د نويو مجردو لارو ټينګ ملاتړی و: «هېڅوک به موږ له هغه جنت
څخه ونه باسي چې کانتور راته جوړ کړے دے۔» په همدې وخت کښې د دې نويو
لارو (چې عجيبه دا ده اوس «کلاسیکې» بلل کېږي) بنسټيزو نيوکو ته هم
متوجه و۔ هغه داسې لار وړانديز کړه چې دواړه موخې تر لاسه کړي:
\begin{enumerate}
\item کلاسیکې لارې په صوري بديهي اصولو او قاعدو تمثيل کړئ؛ رياضياتي
  پوښتنې په يوه بديهي نظام کښې د فارمولونه په توګه تمثيل کړئ۔
\item په خوندي، «متناهي‌پالو» لارو ثابته کړئ چې دا صوري استنتاجي
  نظامونه سازګار دي۔
\end{enumerate}
\olpseganchor{OLP-0276-B017}{end}

\olpseganchor{OLP-0276-B018}{start}
د هېلبرټ کار د لومړۍ موخې د پوره کولو پر لور ډېر وړاندې لاړ۔ په 1899
کښې ئې دا کار په خپل نامتو کتاب \emph{د هندسې بنسټونه} کښې د هندسې
دپاره کړے و۔ په ورپسې کلونو کښې هغه او يو شمېر زده‌کوونکو او
همکارانو ئې د رياضياتو پر نورو څانګو کار وکړ، څو هغه څه وکړي چې
هېلبرټ د هندسې دپاره کړي وو۔ هېلبرټ پخپله د حساب او تحليل دپاره
بديهي نظامونه ورکړل۔ څېر مېلو د سټ تيوري بديهي کول ورکړل، چې
فرېنکل، سکولېم، فون نويمان او نورو وغځول۔ د 1920مو کلونو تر نيمايي
دوه داسې لارې وې چې د «ټولو» رياضياتو د بديهي کولو نوم ئې غوښت:
د رسل او وايټ‌هېډ \emph{پرېنسيپيا مېتېماتيکا} او هغه څه چې وروسته
د څېر مېلو--فرېنکل سټ تيوري په نوم وپېژندل شول۔
\olpseganchor{OLP-0276-B018}{end}

\olpseganchor{OLP-0276-B019}{start}
په 1921 کښې هېلبرټ يوه څېړنيزه پروژه پيل کړه څو د دې نظامونو د
سازګارۍ د ثبوت موخه پوره کړي۔ د هغه څو زده‌کوونکو، په ځانګړي ډول
برنايس، اکرمان، او وروسته ګېنڅن، په دې پروژه کښې ورسره مرسته وکړه۔
د دې موخې بنسټيز نظر دا و چې په رياضياتو کښې د ناسازګارۍ د
اشتقاق د امکان پوښتنه د نښو د ممکنه لړيو په باب يو
ترکيبي مسئله وګرځوي؛ يعنې د هغو جملو ممکنه لړۍ چې د حساب، تحليل يا
سټ تيورۍ د يوه بديهي نظام له بديهي اصولو څخه د، مثلاً،
\LR{$!A \land \lnot !A$} د سم اشتقاق معيار پوره کوي۔ د نښو د
داسې لړۍ د ناممکن والي ثبوت، څرنګه چې پخپله رياضياتي ثبوت دے، په دې
بديهي نظامونو کښې صوري کېدے شي۔ په بله وينا، داسې يوه جمله~\LR{$\OCon$}
به وي چې وايي، مثلاً، حساب سازګار دے۔ سربېره پر دې، دا جمله بايد په
اړوندو نظامونو کښې د اثبات وړ وي، په ځانګړي ډول که ثبوت ئې يوازې ډېرې
محدودې، «متناهي‌پالې» لارې غواړي۔
\olpseganchor{OLP-0276-B019}{end}

\olpseganchor{OLP-0276-B020}{start}
دويمه موخه، چې جوړ شوي بديهي نظامونه به هره رياضياتي پوښتنه حل کړي،
په دوو لارو کره کېدے شي۔ په يوه لار داسې بيانېدے شي: د رياضياتو د
بديهي نظام د ژبې د هرې جملې~\LR{$!A$} دپاره يا~\LR{$!A$} يا~\LR{$\lnot !A$} له
بديهي اصولو څخه د اثبات وړ وي۔ که دا رښتيا وے، داسې جمله به نۀ وه چې
د بديهي اصولو پر بنسټ نه ثابته او نه ردېدے شي، او داسې پوښتنه به نۀ
وه چې بديهي اصول ئې حل نۀ کړي۔ دې خاصيت لرونکی بديهي نظام
\emph{بشپړ} بلل کېږي۔ البته د هرې ټاکلې جملې دپاره لا هم دا معلومول
ستونزمن کار کېدے شي چې له دوو بديلونو څخه کوم يو سم دے۔ خو په اصل
کښې بايد د معلومولو يوه لار وي۔ په حقيقت کښې د هېلبرټ تر څېړنې لاندې
د بديهي اصولو او اشتقاق نظامونو دپاره بشپړتيا دا لازموله چې
داسې يوه لار شته، که څه هم هېلبرټ په دې نۀ پوهېده۔ د پوښتنې دويم
تعبير دا لا پياوړې غوښتنه ده: داسې ميخانيکي، محاسبوي لار موجوده وي
چې د يوې ورکړل شوې جملې~\LR{$!A$} دپاره معلومه کړي له بديهي اصولو څخه
استنتاجېدونکې ده که نۀ۔
\olpseganchor{OLP-0276-B020}{end}

\olpseganchor{OLP-0276-B021}{start}
په 1931 کښې ګوډل د «ناتکميلۍ دوه قضيې» ثابتې کړې، چې وئې ښودل دا
پروګرام بريالی کېدے نۀ شي۔ د رياضياتو دپاره هېڅ داسې بديهي نظام نشته
چې بشپړ وي؛ په ځانګړي ډول هغه جمله چې د بديهي اصولو سازګاري بيانوي
نه ثابتېدے او نه ردېدے شي۔
\olpseganchor{OLP-0276-B021}{end}

\olpseganchor{OLP-0276-B022}{start}
دې د هېلبرټ اصلي پروګرام ته وژونکی ګوزار ورکړ۔ خو لکه په رياضياتو
کښې چې ډېر ځله کېږي، د څېړنې په زړه پورې نوې لارې ئې هم پرانيستې۔
که د ټولو رياضياتو يو واحد، هر اړخيز صوري نظام نشته، نو معنا لري چې
لا محدود نظامونه جوړ او وڅېړل شي چې څه پکې ثابتېدے شي۔ همدارنګه معنا
لري چې د دې نظامونو د سازګارۍ د ثابتولو دپاره لږ محدودې ثبوتي لارې
جوړې، او د هغو د سازګارۍ د ثبوت د سختوالي د اندازه کولو لارې وموندل
شي۔ څرنګه چې ګوډل وښودل (نږدې) هر صوري نظام داسې پوښتنې لري چې نۀ
ئې شي حل کولے، نو معنا لري چې د هغو «په زړه پورې» پوښتنو لټون وشي
چې يو ورکړل شوے صوري نظام ئې نۀ شي حل کولے، او معلومه شي چې د حل
دپاره ئې صوري نظام بايد څومره پياوړے وي۔ منطقپوهان تر نن ورځې دا
پوښتنې د رياضياتو په يوه نوې څانګه، د ثبوتونو په تيورۍ، کښې څېړي۔
\olpseganchor{OLP-0276-B022}{end}

\olpseganchor{OLP-0277-B004}{start}
\olfileid{inc}{int}{def}
\olpseganchor{OLP-0277-B004}{end}

\olpseganchor{OLP-0277-B005}{start}
\olsection{تعريفونه}
\olpseganchor{OLP-0277-B005}{end}

\olpseganchor{OLP-0277-B006}{start}
د رياضياتو د صوري کولو، او د داسې صوري کولو د سازګارۍ او بشپړتيا د
ښودلو د هېلبرټ پروژې د عملي کولو دپاره لومړنی کار دا و چې يوه ژبه،
منطقي چوکاټ او د بديهي اصولو نظام وټاکل شي۔ د خپلې موخې دپاره به فرض
کړو چې رياضيات د لومړۍ درجې په يوه ژبه کښې صوري کېدے شي؛ يعنې د
ثابت سمبولونه، تابعې او پريديکاتونه داسې يو سټ شته چې د
لومړۍ درجې منطق له تړونکو او کميت‌ټاکونکو سره يوځاے موږ ته د رياضياتو
د ادعاوو د بيان وس راکوي۔ زياتره خلک مني چې داسې ژبه شته: د سټ تيورۍ
ژبه، چې~\LR{$\in$} پکې يوازينۍ غيرمنطقي نښه ده۔ دا چې دومره ساده ژبه دومره
بيانوونکې ده، البته په لومړي نظر ډېره نامعقوله ادعا ښکاري، او د دې د
ثابتولو دپاره ډېر کار وشو چې نږدې ټول رياضيات په همدې ډېرې محدودې
اصطلاح‌زېرمه کښې بيانېدے شي۔ د ساده‌والي دپاره به اوس خپل بحث حساب
ته محدود کړو، يعنې د رياضياتو هغې برخې ته چې يوازې له طبيعي
عددونو~\LR{$\Nat$} سره کار لري۔ د حساب د حقايقو د بيان طبيعي ژبه~\LR{$\Lang
L_A$} ده۔ \LR{$\Lang L_A$} يو دوه‌ځاييز پريديکات~\LR{$<$}، يو
ثابت سمبول~\LR{$\Obj 0$}، يوه يوځاييزه تابع~\LR{$\prime$}، او دوه
دوه‌ځاييزې تابعې~\LR{$+$} او~\LR{$\times$} لري۔
\olpseganchor{OLP-0277-B006}{end}

\olpseganchor{OLP-0277-B007}{start}
\begin{defn}
د تړلي فارمولونه سټ~\LR{$\Gamma$} يوه \emph{تيوري} ده که د لازمې پايلې له
مخې تړلې وي، يعنې که
\LR{$\Gamma = \Setabs{!A}{\Gamma \Entails !A}$} وي۔
\end{defn}
\olpseganchor{OLP-0277-B007}{end}

\olpseganchor{OLP-0277-B008}{start}
د تيوريو د مشخصولو دوه اسانه لارې شته۔ يوه دا ده چې تيوري په يوه
جوړښت کښې د رښتينو تړلي فارمولونه سټ وټاکو۔ د بېلګې په توګه،
د~\LR{$\Lang L_A$} هغه جوړښت وګورئ چې د تعريف ساحه ئې~\LR{$\Nat$} وي او
ټولې غيرمنطقي نښې پکې هغسې تفسير شوې وي لکه تمه چې ترې کېږي۔
\olpseganchor{OLP-0277-B008}{end}

\olpseganchor{OLP-0277-B009}{start}
\begin{defn}\ollabel{def:standard-model}
د حساب \emph{معياري مدل} هغه جوړښت~\LR{$\Struct N$} دے چې داسې
تعريفېږي:
\begin{enumerate}
\item \LR{$\Domain N = \Nat$}
\item \LR{$\Assign{\Obj 0}{N} = 0$}
\item \LR{$\Assign{\Obj \prime}{N}(n) = n + 1$} د ټولو \LR{$n \in \Nat$} دپاره
\item \LR{$\Assign{\Obj +}{N}(n, m) = n + m$} د ټولو \LR{$n, m \in \Nat$} دپاره
\item \LR{$\Assign{\Obj \times}{N}(n, m) = n\cdot m$} د ټولو \LR{$n, m \in \Nat$} دپاره
\item \LR{$\Assign{\Obj <}{N} = \Setabs{\tuple{n, m}}{n \in \Nat, m \in
  \Nat, n < m}$}
\end{enumerate}
\end{defn}
\olpseganchor{OLP-0277-B009}{end}

\olpseganchor{OLP-0277-B010}{start}
د~`\LR{$\times$}' او~`\LR{$\cdot$}' تر منځ توپير ته پام وکړئ: \LR{$\times$} د حساب
د ژبې يوه نښه ده۔ البته دا مو داسې ټاکلې چې ضرب راياد کړي، خو
\LR{$\times$} د ضرب عمل نۀ دے، بلکې يوه دوه‌ځاييزه تابع‌نښه ده (په رسمي
ډول~\LR{$\Obj f^2_1$})۔ د دې پر خلاف، \LR{$\cdot$} پخپله د ضرب عادي تابع
\emph{ده}۔ کله چې د~\LR{$n \cdot m$} په شان څه وينئ، موخه مو د~\LR{$n$} او~\LR{$m$}
عددونو حاصل‌ضرب وي؛ کله چې د~\LR{$x \times y$} په شان څه وينئ، د حساب د
ژبې د يوې اصطلاح خبره کوو۔ په معياري مدل کښې تابع‌نښه~\LR{$\times$} پر
طبيعي عددونو د تابع~\LR{$\cdot$} په توګه تفسيرېږي۔ د جمع دپاره~\LR{$+$} هم د
حساب د ژبې د تابع‌نښې او هم پر طبيعي عددونو د جمع د تابع په توګه
کاروو۔ دلته بايد له چاپېرياله معلومه کړئ چې کومه معنا موخه ده۔
\olpseganchor{OLP-0277-B010}{end}

\olpseganchor{OLP-0277-B011}{start}
\begin{defn}
د \emph{رښتيني حساب} تيوري د هغو تړلي فارمولونه سټ دے چې د حساب په
معياري مدل کښې پوره کېږي، يعنې
\[
\Th{TA} = \Setabs{!A}{\Sat{N}{!A}}.
\]
\end{defn}
\olpseganchor{OLP-0277-B011}{end}

\olpseganchor{OLP-0277-B012}{start}
\LR{$\Th{TA}$} يوه تيوري ده، ځکه هر کله چې~\LR{$\Th{TA} \Entails !A$} وي،
\LR{$!A$} په هر هغه جوړښت کښې پوره کېږي چې~\LR{$\Th{TA}$} پوره کوي۔
څرنګه چې~\LR{$\Sat{N}{\Th{TA}}$}، نو~\LR{$\Sat{N}{!A}$} لرو، او له دې امله
\LR{$!A \in \Th{TA}$}۔
\olpseganchor{OLP-0277-B012}{end}

\olpseganchor{OLP-0277-B013}{start}
د تيورۍ~\LR{$\Gamma$} د مشخصولو بله لار دا ده چې هغه د هغو تړلي فارمولونه سټ وټاکو
چې د جملو له يوه سټ~\LR{$\Gamma_0$} څخه لازمېږي۔ په دې حالت کښې~\LR{$\Gamma$}
د لازمې پايلې له مخې د~\LR{$\Gamma_0$} «تړون» دے۔ د تيورۍ دا ډول مشخصول
يوازې هغه وخت په زړه پورې دي چې~\LR{$\Gamma_0$} په څرګند ډول مشخص شوے وي؛
مثلاً د~\LR{$\Gamma_0$} غړي لست شوي وي۔ لږ تر لږه~\LR{$\Gamma_0$}
بايد د پرېکړې وړ وي؛ يعنې داسې محاسبه کېدونکې ازموينه بايد وي چې يوه
تړلے فارمول د~\LR{$\Gamma_0$} غړې ده که نۀ۔ د~\LR{$\Gamma_0$} تړلي فارمولونه
د~\LR{$\Gamma$} دپاره \emph{بديهي اصول} بولو، او وايو چې~\LR{$\Gamma$} د
\LR{$\Gamma_0$} په وسيله \emph{بديهي شوې} ده۔
\olpseganchor{OLP-0277-B013}{end}

\olpseganchor{OLP-0277-B014}{start}
\begin{defn}
تيوري~\LR{$\Gamma$} هله او يوازې هله د~\LR{$\Gamma_0$} په وسيله بديهي شوې ده
چې
\[
\Gamma = \Setabs{!A}{\Gamma_0 \Entails !A}
\]
\end{defn}
\olpseganchor{OLP-0277-B014}{end}

\olpseganchor{OLP-0277-B015}{start}
\begin{defn}
هغه تيوري~\LR{$\Th{Q}$} چې په لاندې جملو بديهي شوې ده د «رابنسن
\LR{$\Th{Q}$}» په نوم پېژندل کېږي او د حساب يوه ډېره ساده تيوري ده۔
\begin{align*}
& \lforall[x][\lforall[y][(\eq[x'][y'] \lif \eq[x][y])]] \tag{\LR{$!Q_1$}}\\
& \lforall[x][\eq/[\Obj 0][x']] \tag{\LR{$!Q_2$}}\\
& \lforall[x][(\eq[x][\Obj 0] \lor \lexists[y][\eq[x][y']])] \tag{\LR{$!Q_3$}}\\
& \lforall[x][\eq[(x + \Obj 0)][x]] \tag{\LR{$!Q_4$}}\\
& \lforall[x][\lforall[y][\eq[(x + y')][(x + y)']]] \tag{\LR{$!Q_5$}}\\
& \lforall[x][\eq[(x \times \Obj 0)][\Obj 0]] \tag{\LR{$!Q_6$}}\\
& \lforall[x][\lforall[y][\eq[(x \times y')][((x \times y) + x)]]] \tag{\LR{$!Q_7$}}\\
& \lforall[x][\lforall[y][(x < y \liff \lexists[z][\eq[(z' + x)][y]])]] \tag{\LR{$!Q_8$}}
\end{align*}
د تړلي فارمولونه سټ~\LR{$\{!Q_1, \dots, !Q_8\}$} د~\LR{$\Th{Q}$} بديهي اصول
دي، نو~\LR{$\Th{Q}$} له هغو څخه لازمېدونکي ټول تړلي فارمولونه رانغاړي:
\[
\Th{Q} = \Setabs{!A}{\{!Q_1, \dots, !Q_8\} \Entails !A}.
\]
\end{defn}
\olpseganchor{OLP-0277-B015}{end}

\olpseganchor{OLP-0277-B016}{start}
\begin{defn}
فرض کړئ~\LR{$!A(x)$} په~\LR{$\Lang L_A$} کښې داسې فارمول ده چې آزاد
متغيرونه~\LR{$x$} او~\LR{$y_1$}، \dots،~\LR{$y_n$} لري۔ بيا د لاندې بڼې هره
تړلے فارمول
\[
\lforall[y_1][\dots\lforall[y_n][((!A(\Obj 0) \land \lforall[x][(!A(x)
\lif !A(x'))]) \lif \lforall[x][!A(x)])]]
\]
د \emph{استقرا شېما} يوه بېلګه ده۔
\olpseganchor{OLP-0277-B016}{end}

\olpseganchor{OLP-0277-B017}{start}
\emph{د پيانو حساب}~\LR{$\Th{PA}$} هغه تيوري ده چې د~\LR{$\Th{Q}$} په بديهي
اصولو او د استقرا شېما په ټولو بېلګو بديهي شوې ده۔
\end{defn}
\olpseganchor{OLP-0277-B017}{end}

\olpseganchor{OLP-0277-B018}{start}
\begin{explain}
د استقرا شېما هره بېلګه په~\LR{$\Struct{N}$} کښې رښتونې ده۔ دا خبره تر
ټولو اسانه هغه وخت ليدل کېږي چې فارمول~\LR{$!A$} يوازې يو آزاد
متغير~\LR{$x$} ولري۔ بيا~\LR{$!A(x)$} په~\LR{$\Struct{N}$} کښې د~\LR{$\Nat$} يو
فرعي سټ~\LR{$X_{!A}$} تعريفوي۔ \LR{$X_{!A}$} د ټولو هغو~\LR{$n \in \Nat$} سټ دے
چې~\LR{$s(x) = n$} وي او~\LR{$\Sat{N}{!A(x)}[s]$}۔ د استقرا شېما اړونده بېلګه
دا ده:
\[
((!A(\Obj 0) \land \lforall[x][(!A(x) \lif !A(x'))]) \lif
  \lforall[x][!A(x)]).
\]
که مقدمه ئې په~\LR{$\Struct{N}$} کښې رښتونې وي، نو~\LR{$0 \in X_{!A}$}، او
هر کله چې~\LR{$n \in X_{!A}$} وي، \LR{$n+1$} هم پکې وي۔ له~\LR{$0 \in X_{!A}$}
څخه~\LR{$1 \in X_{!A}$} ترلاسه کوو۔ له~\LR{$1 \in X_{!A}$} څخه~\LR{$2 \in X_{!A}$}
ترلاسه کوو، او همداسې وړاندې۔ نو د هر~\LR{$n \in \Nat$} دپاره
\LR{$n \in X_{!A}$}۔ خو دا په دې معنا ده چې~\LR{$\lforall[x][!A(x)]$} په
\LR{$\Struct{N}$} کښې پوره کېږي۔
\end{explain}
\olpseganchor{OLP-0277-B018}{end}

\olpseganchor{OLP-0277-B019}{start}
\LR{$\Th{Q}$} او~\LR{$\Th{PA}$} دواړه بديهي شوې تيورۍ دي۔ لويه پوښتنه دا ده
چې څومره پياوړې دي؟ مثلاً آيا~\LR{$\Th{PA}$} د~\LR{$\Nat$} په باب ټول هغه
حقايق ثابتولے شي چې په~\LR{$\Lang L_A$} کښې بيانېدے شي؟ په کره توګه، آيا
د~\LR{$\Th{PA}$} بديهي اصول ټولې هغه پوښتنې حل کوي چې په~\LR{$\Lang L_A$} کښې
جوړېدے شي؟
\olpseganchor{OLP-0277-B019}{end}

\olpseganchor{OLP-0277-B020}{start}
د دې پوښتنې بله بڼه دا ده: آيا~\LR{$\Th{PA} = \Th{TA}$}؟ \LR{$\Th{TA}$} په
څرګند ډول د~\LR{$\Nat$} په باب ټول حقايق ثابتوي (يعنې ټول رانغاړي)، او
ټولې هغه پوښتنې حل کوي چې په~\LR{$\Lang L_A$} کښې جوړېدے شي؛ ځکه که~\LR{$!A$}
په~\LR{$\Lang L_A$} کښې تړلے فارمول وي، نو يا~\LR{$\Sat{N}{!A}$} يا
\LR{$\Sat{N}{\lnot !A}$}، او له دې امله يا~\LR{$\Th{TA} \Entails !A$} يا
\LR{$\Th{TA} \Entails \lnot !A$}۔ داسې تيوري \emph{بشپړ} بولو۔
\olpseganchor{OLP-0277-B020}{end}

\olpseganchor{OLP-0277-B021}{start}
\begin{defn}
تيوري~\LR{$\Gamma$} هله او يوازې هله \emph{بشپړ} ده چې د هغې د
ژبې د هرې تړلے فارمول~\LR{$!A$} دپاره يا~\LR{$\Gamma \Entails !A$} يا
\LR{$\Gamma \Entails \lnot !A$} وي۔
\end{defn}
\olpseganchor{OLP-0277-B021}{end}

\olpseganchor{OLP-0277-B022}{start}
\begin{explain}
د بشپړتيا د قضيې له مخې~\LR{$\Gamma \Entails !A$} هله او يوازې هله چې
\LR{$\Gamma \Proves !A$}، نو~\LR{$\Gamma$} هله او يوازې هله بشپړه ده چې د
هغې د ژبې د هرې تړلے فارمول~\LR{$!A$} دپاره يا~\LR{$\Gamma \Proves !A$} يا
\LR{$\Gamma \Proves \lnot !A$} وي۔
\end{explain}
\olpseganchor{OLP-0277-B022}{end}

\olpseganchor{OLP-0277-B023}{start}
بله پوښتنه چې راپورته کېږي دا ده: آيا داسې محاسبوي کړنلاره شته چې
وآزمويو يوه تړلے فارمول په~\LR{$\Th{TA}$}، په~\LR{$\Th{PA}$}، يا ان يوازې
په~\LR{$\Th{Q}$} کښې ده که نۀ؟ دا پوښتنه د دې په تعريفولو لا کره کولے شو
چې يو سټ (مثلاً د تړلي فارمولونه سټ) کله د پرېکړې وړ وي۔
\olpseganchor{OLP-0277-B023}{end}

\olpseganchor{OLP-0277-B024}{start}
\begin{defn}
سټ~\LR{$X$} هله او يوازې هله \emph{د پرېکړې وړ} دے چې داسې محاسبوي
کړنلاره وي چې پر ننوت~\LR{$x$}، که~\LR{$x \in X$} وي~\LR{$1$} او که نۀ وي~\LR{$0$}
بېرته ورکوي۔
\end{defn}
\olpseganchor{OLP-0277-B024}{end}

\olpseganchor{OLP-0277-B025}{start}
نو زموږ پوښتنه داسې شوه: آيا~\LR{$\Th{TA}$} (\LR{$\Th{PA}$}، \LR{$\Th{Q}$})
د پرېکړې وړ ده؟
\olpseganchor{OLP-0277-B025}{end}

\olpseganchor{OLP-0277-B026}{start}
د دې ټولو پوښتنو ځواب به «نه» وي۔ له دې تيوريو څخه يوه هم د پرېکړې
وړ نۀ ده۔ خو دا پديده يوازې په همدې ځانګړو تيوريو پورې محدوده نۀ ده۔
په حقيقت کښې \emph{هره} هغه تيوري چې ځينې شرطونه پوره کوي له همدغو
پايلو سره مخ کېږي۔ له دغو شرطونو څخه يو، چې~\LR{$\Th{Q}$} او~\LR{$\Th{PA}$}
ئې پوره کوي، دا دے چې د بديهي اصولو د د پرېکړې وړ سټ په وسيله
بديهي شوې وي۔
\olpseganchor{OLP-0277-B026}{end}

\olpseganchor{OLP-0277-B027}{start}
\begin{defn}
يوه تيوري \emph{په محاسبوي ډول بديهي کېدونکې} ده که د بديهي اصولو د
د پرېکړې وړ سټ په وسيله بديهي شوې وي۔
\end{defn}
\olpseganchor{OLP-0277-B027}{end}

\olpseganchor{OLP-0277-B028}{start}
\begin{ex}
هره تيوري چې د تړلي فارمولونه د يوه متناهي سټ په وسيله بديهي شوې وي
په محاسبوي ډول بديهي کېدونکې ده، ځکه هر متناهي سټ د پرېکړې وړ دے۔ نو مثلاً
\LR{$\Th{Q}$}، په محاسبوي ډول بديهي کېدونکې ده۔
\olpseganchor{OLP-0277-B028}{end}

\olpseganchor{OLP-0277-B029}{start}
د~\LR{$\Th{PA}$} په شان په شېمايي ډول بديهي شوې تيورۍ هم
په محاسبوي ډول بديهي کېدونکې دي۔ د دې د ازمويلو دپاره چې~\LR{$!B$} د~\LR{$\Th{PA}$} په
بديهي اصولو کښې شامله ده که نۀ، يعنې د تابع~\LR{$\Char{X}$} د محاسبې
دپاره، چې که~\LR{$!B$} د~\LR{$\Th{PA}$} يو بديهي اصل وي
\LR{$\Char{X}(!B) = 1$} او که نۀ وي~\LR{$=0$}، داسې کولے شو: لومړے وګورئ
\LR{$!B$} د~\LR{$\Th{Q}$} له بديهي اصولو څخه ده که نۀ۔ که وي، ځواب «هو» دے
او~\LR{$\Char{X}(!B) = 1$}۔ که نۀ وي، وازمويئ چې د استقرا شېما بېلګه ده
که نۀ۔ دا په منظمه توګه کېدے شي؛ په دې حالت کښې ښايي تر ټولو اسانه
دا وي چې داسې ئې ووينو: د استقرا شېما هره بېلګه په يو شمېر کلي
کميت‌ټاکونکو پيلېږي، او بيا داسې فرعي-فارمول راځي چې يو شرطي
دے۔ د هغه شرطي پايله~\LR{$\lforall[x][!A(x, y_1, \dots, y_n)]$} ده، چې
\LR{$x$} او~\LR{$y_1$}، \dots،~\LR{$y_n$} د~\LR{$!A$} ټول آزاد متغيرونه دي او د~\LR{$!B$}
په پيل کښې کميت‌ټاکونکي متغيرونه~\LR{$y_1$}، \dots،~\LR{$y_n$} تړي۔ کله چې دا
\LR{$!A$} راوباسو او وګورو چې آزاد متغيرونه ئې د لومړنيو کلي
کميت‌ټاکونکو او~\LR{$\lforall[x]$} له تړلو متغيرونو سره برابر دي، بيا
ګورو چې د شرطي مقدمه له لاندې سره برابره ده که نۀ:
\[
!A(\Obj 0, y_1, \dots, y_n) \land \lforall[x][(!A(x, y_1, \dots, y_n)
\lif !A(x', y_1, \dots, y_n))]
\]
بيا هم، که برابره وي، \LR{$!B$} د استقرا شېما بېلګه ده؛ او که برابره نۀ
وي، \LR{$!B$} د استقرا شېما بېلګه نۀ ده۔
\end{ex}
\olpseganchor{OLP-0277-B029}{end}

\olpseganchor{OLP-0277-B030}{start}
د دې پوښتنې، او د دې لا عامې پوښتنې د ځوابولو دپاره چې کومې تيورۍ
بشپړې يا د پرېکړې وړ دي، لاندې تعريف هم ګټور دے۔ په ياد ولرئ چې سټ~\LR{$X$}
هله او يوازې هله د شمېر وړ دے چې تش وي يا داسې پر هدف سټ بشپړه
تابع~\LR{$f \colon \Nat \to X$} موجوده وي۔ داسې تابع د~\LR{$X$} شمېره بلل کېږي۔
\olpseganchor{OLP-0277-B030}{end}

\olpseganchor{OLP-0277-B031}{start}
\begin{defn}
سټ~\LR{$X$} هله او يوازې هله \emph{په محاسبوي ډول د شمېر وړ} (په لنډ ډول
c.e.) بلل کېږي چې تش وي يا محاسبه کېدونکې شمېره ولري۔
\end{defn}
\olpseganchor{OLP-0277-B031}{end}

\olpseganchor{OLP-0277-B032}{start}
د د محاسبوي بديهي کېدو وړتيا تر څنګ، پر هغو تيوريو بل شرط چې د ناتکميلۍ
قضيې پرې پلي کېږي دا دے چې د محاسبه کېدونکو تابعو او
د پرېکړې وړ اړيکو په باب بنسټيز حقايق ثابتولو ته کافي پياوړې وي۔
له «بنسټيزو حقايقو» څخه موخه هغه تړلي فارمولونه دي چې د هرې تابعې د
ارګومېنټونو دپاره ئې ارزښتونه بيانوي۔ او له «کافي پياوړې» څخه موخه
دا ده چې اړوندې تيورۍ دا جملې په خپلو قضيو کښې شمېري۔ د بېلګې په
توګه يوه نمونه‌يي محاسبه کېدونکې تابع، يعنې جمع، وګورئ۔ د~\LR{$2$} او~\LR{$3$}
ارګومېنټونو دپاره د~\LR{$+$} ارزښت~\LR{$5$} دے، يعنې~\LR{$2+3 = 5$}۔ د حساب په
ژبه کښې يوه جمله چې وايي د~\LR{$2$} او~\LR{$3$} ارګومېنټونو دپاره د~\LR{$+$} ارزښت
\LR{$5$} دے، دا ده: \LR{$(\num{2} + \num{3}) = \num{5}$}۔ او مثلاً~\LR{$\Th{Q}$}
دا جمله ثابتوي۔ په عام ډول غواړو د هرې محاسبه کېدونکې تابع
\LR{$f(x_1, x_2)$} دپاره په~\LR{$\Lang{L_A}$} کښې داسې فارمول
\LR{$!A_f(x_1, x_2, y)$} وي چې هر کله~\LR{$f(n_1, n_2) = m$}،
\LR{$\Th{Q} \Proves !A_f(\num{n_1}, \num{n_2}, \num{m})$} وي۔ په دې ډول
\LR{$\Th{Q}$} ثابتوي چې د~\LR{$n_1$}، \LR{$n_2$} ارګومېنټونو دپاره د~\LR{$f$} ارزښت~\LR{$m$}
دے۔ په حقيقت کښې لږ څه زيات غواړو، يعنې دا هم ثابته کړي چې د~\LR{$n_1$}،
\LR{$n_2$} ارګومېنټونو دپاره بل هېڅ عدد د~\LR{$f$} ارزښت نۀ دے۔ د
د پرېکړې وړ اړيکو دپاره هم همداسې ده۔ لاندې دوه تعريفونه دا خبره
کره کوي۔
\olpseganchor{OLP-0277-B032}{end}

\olpseganchor{OLP-0277-B033}{start}
\begin{defn}
فارمول~\LR{$!A(x_1, \dots, x_k, y)$} په~\LR{$\Gamma$} کښې تابع
\LR{$f\colon \Nat^k \to \Nat$} هله او يوازې هله \emph{تمثيل}
کوي چې هر کله~\LR{$f(n_1, \dots, n_k) = m$} وي، نو
\begin{enumerate}
\item \LR{$\Gamma \Proves !A(\num{n_1}, \dots, \num{n_k}, \num{m})$}، او
\item \LR{$\Gamma \Proves \lforall[y](!A(\num{n_1}, \dots, \num{n_k},
y) \lif y = \num{m})$}۔
\end{enumerate}
\end{defn}
\olpseganchor{OLP-0277-B033}{end}

\olpseganchor{OLP-0277-B034}{start}
\begin{defn}
فارمول~\LR{$!A(x_1, \dots, x_k)$} اړيکه~\LR{$R \subseteq \Nat^k$} هله او
يوازې هله \emph{تمثيل} کوي چې
\begin{enumerate}
\item هر کله~\LR{$R(n_1, \dots, n_k)$} وي،
\LR{$\Gamma \Proves !A(\num{n_1}, \dots, \num{n_k})$}، او
\item هر کله~\LR{$R(n_1, \dots, n_k)$} نۀ وي،
\LR{$\Gamma \Proves \lnot !A(\num{n_1}, \dots, \num{n_k})$}۔
\end{enumerate}
\end{defn}
\olpseganchor{OLP-0277-B034}{end}

\olpseganchor{OLP-0277-B035}{start}
يوه تيوري د ناتکميلۍ د قضيو د پلي کېدو دپاره «کافي پياوړې» ده که
ټولې محاسبه کېدونکې تابعې او ټولې د پرېکړې وړ اړيکې تمثيل
کوي۔ \LR{$\Th{Q}$} او د هغې غځونې دا شرط پوره کوي، خو د دې ثابتول به لږ
وخت واخلي؛ دا د هغو خبرو په باب يو غيرابتذالي حقيقت دے چې~\LR{$\Th{Q}$}
ئې ثابتولے شي، او ښودل ئې ځکه ستونزمن دي چې~\LR{$\Th{Q}$} يوازې څو بديهي
اصول لري چې بايد دا ټول حقايق ترې ثابت کړو۔ خو~\LR{$\Th{Q}$} ډېره کمزورې
تيوري ده۔ نو که څه هم دا ثابتول ستونزمن دي چې~\LR{$\Th{Q}$} ټولې محاسبه
کېدونکې تابعې تمثيلوي، زياتره په زړه پورې تيورۍ تر~\LR{$\Th{Q}$} پياوړې
دي؛ يعنې تر~\LR{$\Th{Q}$} زيات څه ثابتوي۔ او هر څه چې~\LR{$\Th{Q}$} ثابتوي،
هره پياوړې تيوري ئې هم ثابتوي؛ څرنګه چې~\LR{$\Th{Q}$} ټولې محاسبه کېدونکې
تابعې تمثيلوي، هره پياوړې تيوري ئې هم تمثيلوي۔ په دې معنا ډېرې په
زړه پورې تيورۍ د ناتکميلۍ د قضيو دا شرط پوره کوي۔ نو زموږ سخت کار به
ګټه وکړي، ځکه ښيي چې د ناتکميلۍ قضيې پر ډېرو تيوريو پلي کېږي۔ په
يقيني ډول هره هغه تيوري چې موخه ئې د «ټولو رياضياتو» صوري کول وي،
بايد هر څه ثابت کړي چې~\LR{$\Th{Q}$} ئې ثابتوي، ځکه لږ تر لږه بايد د
لومړنيو محاسبو پايلې رانغاړلے شي۔ نو هره هغه تيوري چې د «ټولو
رياضياتو» د تيورۍ نومانده وي، د ناتکميلۍ قضيې به پرې پلي کېږي۔
\olpseganchor{OLP-0277-B035}{end}

\olpseganchor{OLP-0278-B004}{start}
\olfileid{inc}{int}{ovr}
\olpseganchor{OLP-0278-B004}{end}

\olpseganchor{OLP-0278-B005}{start}
\olsection{د ناتکميلۍ د پايلو ټوليزه کتنه}
\olpseganchor{OLP-0278-B005}{end}

\olpseganchor{OLP-0278-B006}{start}
هېلبرټ تمه لرله چې رياضيات په داسې په محاسبوي ډول بديهي کېدونکې تيورۍ کښې
صوري کېدے شي چې بشپړ او د پرېکړې وړ ثابتېدے شي۔ سربېره پر
دې، موخه ئې دا وه چې د دې تيورۍ سازګاري په ډېرو کمزورو،
«متناهي‌پالو» لارو ثابته کړي، څو د شهوديت د ننګونو پر وړاندې له
کلاسیکو رياضياتو دفاع وکړي۔ د ګوډل د ناتکميلۍ قضيو وښودل چې دا موخې
نۀ شي تر لاسه کېدے۔
\olpseganchor{OLP-0278-B006}{end}

\olpseganchor{OLP-0278-B007}{start}
د ګوډل د ناتکميلۍ لومړۍ قضيې وښودل چې د رسل او وايټ‌هېډ د
\emph{پرېنسيپيا مېتېماتيکا} يوه بڼه بشپړ نۀ ده۔ خو ثبوت په
اصل کښې ډېر عام و او پر ډېرو بېلابېلو تيوريو پلي کېږي۔ معنا ئې دا ده
چې ستونزه يوازې دا نۀ وه چې \emph{پرېنسيپيا مېتېماتيکا} رياضيات په
بشپړ ډول رانۀ غښتل، بلکې \emph{هېڅ} منل کېدونکې تيوري ئې نۀ شي
رانغاړلے۔ څه وخت ولګېد چې د تيوريو هغه ځانګړتياوې بېلې شي چې د
ناتکميلۍ د قضيو د پلي کېدو دپاره بس دي، او د ګوډل ثبوت داسې عام شي
چې يوازې په همدې ځانګړتياوو تکيه وکړي۔ خو اوس کولے شو د حساب د
ژبې~\LR{$\Lang L_A$} د تيوريو دپاره د ناتکميلۍ د لومړۍ قضيې يوه ډېره عامه
بڼه بيان کړو۔
\olpseganchor{OLP-0278-B007}{end}

\olpseganchor{OLP-0278-B008}{start}
\begin{thm}
که~\LR{$\Gamma$} په~\LR{$\Lang L_A$} کښې داسې سازګاره او په محاسبوي ډول بديهي کېدونکې
تيوري وي چې ټولې محاسبه کېدونکې تابعې او ټولې د پرېکړې وړ اړيکې
تمثيل کوي، نو~\LR{$\Gamma$} بشپړ نۀ ده۔
\end{thm}
\olpseganchor{OLP-0278-B008}{end}

\olpseganchor{OLP-0278-B009}{start}
دا ويل چې~\LR{$\Gamma$} بشپړ نۀ ده، په دې معنا دي چې لږ تر لږه
د يوې تړلے فارمول~\LR{$!A$} دپاره~\LR{$\Gamma \Proves/ !A$} او
\LR{$\Gamma \Proves/ \lnot !A$} وي۔ داسې تړلے فارمول (له~\LR{$\Gamma$}
څخه) \emph{خپلواکه} بلل کېږي۔
په حقيقت کښې په نسبتاً چټکۍ ثابتولے شو
چې خپلواکې جملې بايد موجودې وي۔ خو د قضيې د ګوډلي ثبوت ځواک په دې کښې
دے چې د داسې خپلواکې تړلے فارمول يوه \emph{ټاکلې بېلګه} وړاندې کوي۔
دا په زړه پورې جوړونه يوه تړلے فارمول~\LR{$!G_\Gamma$} پيدا کوي، چې د
\LR{$\Gamma$} \emph{ګوډل جمله} بلل کېږي او ځکه ناثابتېدونکې ده چې په
\LR{$\Gamma$} کښې~\LR{$!G_\Gamma$} له دې ادعا سره برابره ده چې~\LR{$!G_\Gamma$} په
\LR{$\Gamma$} کښې ناثابتېدونکې ده۔ جوړونه دا کار \emph{تعميري} کوي؛ يعنې
که د~\LR{$\Gamma$} بديهي کول او د اشتقاق نظام بيان راکړل شي، ثبوت
په رښتيا د~\LR{$!G_\Gamma$} د ليکلو يوه لار راکوي۔
\olpseganchor{OLP-0278-B009}{end}

\olpseganchor{OLP-0278-B010}{start}
د ګوډل په ثبوت کښې جوړونه غواړي چې د~\LR{$\Lang L_A$} د اصطلاحاتو او
فارمولونه ځانګړتياوې او عملونه پخپله په~\LR{$\Lang L_A$} کښې د بيانولو
لار پيدا کړو۔ په دې کښې داسې ځانګړتياوې شاملې دي لکه «\LR{$!A$}
تړلے فارمول ده»، «\LR{$\delta$} د~\LR{$!A$} اشتقاق ده»، او داسې
عملونه لکه~\LR{$\Subst{!A}{t}{x}$}۔ دا لار بايد (الف) دا ځانګړتياوې او
اړيکې د نښو او د هغو د لړيو د «کوډولو» په وسيله د طبيعي عددونو په
توګه بيان کړي (ځکه اصطلاحات، فارمولونه، اشتقاقونه او نور د
نښو لړۍ دي، او~\LR{$\Lang L_A$} د طبيعي عددونو په باب خبرې کوي)۔ بايد
(ب) دا کار داسې وکړي چې~\LR{$\Gamma$} اړوند حقايق ثابت کړي؛ نو بايد وښيو
چې دا ځانګړتياوې د طبيعي عددونو په د پرېکړې وړ خاصيتونو کوډېږي او
عملونه د طبيعي عددونو له محاسبه کېدونکو تابعو سره سمون لري۔ دې ته
«د نحو حسابي کول» ويل کېږي۔
\olpseganchor{OLP-0278-B010}{end}

\olpseganchor{OLP-0278-B011}{start}
خو د نحو د حسابي کولو تر څېړلو وړاندې به دا شرط وګورو چې~\LR{$\Gamma$}
«کافي پياوړې» وي؛ يعنې ټولې محاسبه کېدونکې تابعې او ټولې
د پرېکړې وړ اړيکې تمثيل کړي۔ دا غواړي چې «محاسبه کېدونکي»
ته يو کره تعريف ورکړو۔ دا په بېلابېلو لارو کېدے شي؛ مثلاً د ټيورينګ
ماشينونو د مدل په وسيله، يا د هغو تابعو په توګه چې په کومې عمومي
موخې پروګرام‌ليک ژبه کښې پروګرامونه ئې محاسبه کوي۔ خو څرنګه چې زموږ
موخه د~\LR{$\Lang L_A$} د ژبې په يوه تيورۍ کښې د دې تابعو او اړيکو
تمثيل کول دي، غوره ده چې يوازې د عددي تابعو د محاسبه کېدنې يو ساده تعريف
وټاکو۔ دا د \emph{بازګشتي تابعې} مفهوم دے۔ نو لومړے به د بازګشتي
تابعو بحث وکړو۔ بيا به وښيو چې~\LR{$\Th{Q}$} لا له وړاندې ټولې بازګشتي
تابعې او اړيکې تمثيل کوي۔ دا به موږ ته اجازه راکړي چې د
ناتکميلۍ قضيه پر~\LR{$\Th{Q}$} او~\LR{$\Th{PA}$} غوندې ټاکلو تيوريو پلې کړو،
ځکه ثابته کړې به مو وي چې دا د «کافي پياوړو» تيوريو بېلګې دي۔
\olpseganchor{OLP-0278-B011}{end}

\olpseganchor{OLP-0278-B012}{start}
د نحو د حسابي کولو وروستۍ پايله داسې فارمول~\LR{$\OProv[\Gamma](x)$}
ده چې د فارمولونه د عددونو په توګه د کوډولو په وسيله د~\LR{$\Gamma$}
له بديهي اصولو څخه د اثبات وړتيا بيانوي۔ په ځانګړي ډول، که~\LR{$!A$} په
عدد~\LR{$n$} کوډ شوې وي او~\LR{$\Gamma \Proves !A$}، نو
\LR{$\Gamma \Proves \Prov[\Gamma](\num{n})$}۔ د~\LR{$\Gamma$} دپاره دا
«اثبات‌وړتيا پريديکات» موږ ته دا وس هم راکوي چې په يوه معنا د~\LR{$\Gamma$}
سازګاري د~\LR{$\Lang L_A$} د يوې تړلے فارمول په توګه بيان کړو: د
\LR{$\Gamma$} دپاره د «سازګارۍ جمله» په نوم هغه تړلے فارمول
\LR{$\lnot \Prov[\Gamma](\num{n})$} ونيسئ، چې~\LR{$n$} پکې د يوه تناقض،
مثلاً~\LR{$\lfalse$}، کوډ وي۔ د ناتکميلۍ
دويمه قضيه وايي چې سازګارې په محاسبوي ډول بديهي کېدونکې تيورۍ خپلې د سازګارۍ
جملې هم نۀ ثابتوي۔ د دې قضيې د پلي کېدو شرطونه تر دې لږ سخت دي چې
تيوري دې يوازې ټولې محاسبه کېدونکې تابعې او د پرېکړې وړ اړيکې
تمثيل کړي، خو وبه ښيو چې~\LR{$\Th{PA}$} دا شرطونه پوره کوي۔
\olpseganchor{OLP-0278-B012}{end}

\olpseganchor{OLP-0279-B004}{start}
\olfileid{inc}{int}{dec}
\olpseganchor{OLP-0279-B004}{end}

\olpseganchor{OLP-0279-B005}{start}
\olsection{ناپرېکړتيا او ناتکميلي}
\olpseganchor{OLP-0279-B005}{end}

\olpseganchor{OLP-0279-B006}{start}
د ناتکميلۍ د قضيو د ګوډلي ثبوت د نحو حسابي کول غواړي۔ خو ان له دې
پرته، يوازې په دې فرض چې يوه تيوري ټولې د پرېکړې وړ اړيکې
تمثيل کوي، ځينې ښې پايلې تر لاسه کولے شو۔ ثبوت يو قطري
استدلال دے چې د درېدنې د مسئلې د ناپرېکړتيا ثبوت ته ورته دے۔
\olpseganchor{OLP-0279-B006}{end}

\olpseganchor{OLP-0279-B007}{start}
\begin{thm}
که~\LR{$\Gamma$} يوه داسې سازګاره تيوري وي چې هره د پرېکړې وړ اړيکه
تمثيل کوي، نو~\LR{$\Gamma$} د پرېکړې وړ نۀ ده۔
\end{thm}
\olpseganchor{OLP-0279-B007}{end}

\olpseganchor{OLP-0279-B008}{start}
\begin{proof}
فرض کړئ~\LR{$\Gamma$} د پرېکړې وړ وي۔ ښيو چې که~\LR{$\Gamma$} هره
د پرېکړې وړ اړيکه تمثيل کړي، نو بايد ناسازګاره وي۔
\olpseganchor{OLP-0279-B008}{end}

\olpseganchor{OLP-0279-B009}{start}
د پرېکړې وړ خاصيتونه (يوځاييزې اړيکې) په هغو فارمولونه تمثيلېږي
چې يو آزاد متغير لري۔ فرض کړئ~\LR{$!A_0(x)$}، \LR{$!A_1(x)$}، \dots، د ټولو
داسې فارمولونه يوه محاسبه کېدونکې شمېره وي۔ اوس لاندې سټ
\LR{$D \subseteq \Nat$} وګورئ:
\[
D = \Setabs{n}{\Gamma \Proves \lnot !A_n(\num{n})}
\]
سټ~\LR{$D$} د پرېکړې وړ دے، ځکه دا ازمويلے شو چې~\LR{$n \in D$} دے که نۀ:
لومړے~\LR{$!A_n(x)$}، او له هغې څخه~\LR{$\lnot !A_n(\num{n})$} محاسبه کوو۔ په
څرګند ډول په~\LR{$!A_n(x)$} کښې د~\LR{$x$} د هر آزاد پېښېدنې پر ځاے د
اصطلاح~\LR{$\num{n}$} ايښوول
او د~\LR{$!A_n(\num{n})$} مخې ته~\LR{$\lnot$} ليکل يو ميخانيکي کار دے۔ د فرض
له مخې~\LR{$\Gamma$} د پرېکړې وړ ده، نو دا ازمويلے شو چې
\LR{$\lnot !A_n(\num{n}) \in \Gamma$} ده که نۀ۔ که وي، \LR{$n \in D$}؛ او که
نۀ وي، \LR{$n \notin D$}۔ نو~\LR{$D$} هم د پرېکړې وړ دے۔
\olpseganchor{OLP-0279-B009}{end}

\olpseganchor{OLP-0279-B010}{start}
څرنګه چې~\LR{$\Gamma$} ټول د پرېکړې وړ خاصيتونه تمثيل کوي، نو~\LR{$D$} هم
تمثيل کوي۔ او هغه فارمولونه چې~\LR{$D$} په~\LR{$\Gamma$} کښې
تمثيل کوي ټول د~\LR{$!A_0(x)$}، \LR{$!A_1(x)$}، \dots، په لړۍ کښې
دي۔ نو داسې عدد~\LR{$d$} ونيسئ چې~\LR{$!A_d(x)$} په~\LR{$\Gamma$} کښې~\LR{$D$}
تمثيل کوي۔ که~\LR{$d \notin D$} وي، نو څرنګه چې~\LR{$!A_d(x)$}
تمثيلوي~\LR{$D$}، \LR{$\Gamma \Proves \lnot !A_d(\num{d})$}۔ خو دا په دې معنا ده
چې~\LR{$d$} د~\LR{$D$} ټاکونکی شرط پوره کوي، او له دې امله~\LR{$d \in D$}۔ دا له
\LR{$d \notin D$} سره تناقض لري۔ نو د غيرمستقيم ثبوت له مخې~\LR{$d \in D$}۔
\olpseganchor{OLP-0279-B010}{end}

\olpseganchor{OLP-0279-B011}{start}
څرنګه چې~\LR{$d \in D$}، د~\LR{$D$} د تعريف له مخې
\LR{$\Gamma \Proves \lnot !A_d(\num{d})$}۔ له بلې خوا، څرنګه چې~\LR{$!A_d(x)$}
په~\LR{$\Gamma$} کښې~\LR{$D$} تمثيل کوي،
\LR{$\Gamma \Proves !A_d(\num{d})$}۔ نو~\LR{$\Gamma$} ناسازګاره ده۔
\end{proof}
\olpseganchor{OLP-0279-B011}{end}

\olpseganchor{OLP-0279-B012}{start}
\begin{explain}
پاسنۍ قضيه ښيي چې هېڅ سازګاره تيوري چې ټولې د پرېکړې وړ اړيکې
تمثيل کوي، د پرېکړې وړ کېدے نۀ شي۔ موږ به وښيو چې~\LR{$\Th{Q}$}
ټولې د پرېکړې وړ اړيکې تمثيل کوي؛ معنا ئې دا ده چې ټولې
هغه تيورۍ چې~\LR{$\Th{Q}$} رانغاړي، لکه~\LR{$\Th{PA}$} او~\LR{$\Th{TA}$}، هم دا
کار کوي او له دې امله هغوى هم د پرېکړې وړ نۀ دي۔ (څرنګه چې دا ټولې
تيورۍ په معياري مدل کښې رښتونې دي، ټولې سازګارې دي۔)
\olpseganchor{OLP-0279-B012}{end}

\olpseganchor{OLP-0279-B013}{start}
دا پايله د ناتکميلۍ د لومړۍ قضيې د يوې کمزورې بڼې د ترلاسه کولو
دپاره هم کارولے شو۔ هره هغه تيوري چې په محاسبوي ډول بديهي کېدونکې او
بشپړ وي، د پرېکړې وړ ده۔ نو سازګارې تيورۍ چې
په محاسبوي ډول بديهي کېدونکې وي او ټول د پرېکړې وړ خاصيتونه تمثيل
کوي، بشپړ کېدے نۀ شي۔
\end{explain}
\olpseganchor{OLP-0279-B013}{end}

\olpseganchor{OLP-0279-B014}{start}
\begin{thm}
که~\LR{$\Gamma$} په محاسبوي ډول بديهي کېدونکې او بشپړ وي، نو د پرېکړې وړ
ده۔
\end{thm}
\olpseganchor{OLP-0279-B014}{end}

\olpseganchor{OLP-0279-B015}{start}
\begin{proof}
هره ناسازګاره تيوري د پرېکړې وړ ده، ځکه ناسازګارې تيورۍ ټول
تړلي فارمولونه رانغاړي؛ نو د «آيا~\LR{$!A \in \Gamma$}؟» پوښتنې ځواب تل
«هو» وي، يعنې پرېکړه ئې کېدے شي۔
\olpseganchor{OLP-0279-B015}{end}

\olpseganchor{OLP-0279-B016}{start}
نو فرض کړئ~\LR{$\Gamma$} سازګاره، او ورسره په محاسبوي ډول بديهي کېدونکې او
بشپړ وي۔ څرنګه چې~\LR{$\Gamma$} په محاسبوي ډول بديهي کېدونکې ده، نو
په محاسبوي ډول د شمېر وړ ده۔ ځکه له~\LR{$\Gamma$} د بديهي اصولو څخه ټول
سم اشتقاقونه په يوه محاسبه کېدونکې تابع شمېرلے شو۔ له يوه سم
اشتقاق څخه هغه تړلے فارمول محاسبه کولے شو چې دا ئې
اشتقاق؛ نو په ګډه داسې محاسبه کېدونکې تابع شته چې د~\LR{$\Gamma$} ټولې
قضيې شمېري۔ څرنګه چې~\LR{$\Gamma$} سازګاره او بشپړ ده،
تړلے فارمول هله او يوازې هله د~\LR{$\Gamma$} قضيه ده چې~\LR{$\lnot !A$}
قضيه نۀ وي۔ نو داسې پرېکړه کولے شو چې~\LR{$!A \in \Gamma$} ده که نۀ۔ د
\LR{$\Gamma$} ټولې قضيې وشمېرئ۔ کله چې~\LR{$!A$} په لست کښې راشي، پوهېږو چې
\LR{$\Gamma \Proves !A$}۔ کله چې~\LR{$\lnot !A$} په لست کښې راشي، پوهېږو چې
\LR{$\Gamma \Proves/ !A$}۔ څرنګه چې~\LR{$\Gamma$} بشپړ ده، له دې دوو
حالتونو يو بالاخره رامنځته کېږي؛ نو کړنلاره بالاخره ځواب ورکوي۔
\end{proof}
\olpseganchor{OLP-0279-B016}{end}

\olpseganchor{OLP-0279-B017}{start}
\begin{cor}
\ollabel{cor:incompleteness}
که~\LR{$\Gamma$} سازګاره او په محاسبوي ډول بديهي کېدونکې وي او هر د پرېکړې وړ
خاصيت تمثيل کړي، نو بشپړ نۀ ده۔
\end{cor}
\olpseganchor{OLP-0279-B017}{end}

\olpseganchor{OLP-0279-B018}{start}
\begin{proof}
که~\LR{$\Gamma$} بشپړ وے، د تېرې قضيې له مخې به د پرېکړې وړ وه
(ځکه په محاسبوي ډول بديهي کېدونکې او سازګاره ده)۔ خو څرنګه چې~\LR{$\Gamma$} هر د پرېکړې وړ
خاصيت تمثيل کوي، د لومړۍ قضيې له مخې د پرېکړې وړ نۀ ده۔
\end{proof}
\olpseganchor{OLP-0279-B018}{end}

\olpseganchor{OLP-0279-B019}{start}
\begin{prob}
وښيئ چې~\LR{$\Th{TA} = \Setabs{!A}{\Sat{N}{!A}}$}، په محاسبوي ډول بديهي کېدونکې نۀ
ده۔ دا فرضولے شئ چې~\LR{$\Th{TA}$} ټول د پرېکړې وړ خاصيتونه تمثيلوي۔
\end{prob}
\olpseganchor{OLP-0279-B019}{end}

\olpseganchor{OLP-0279-B020}{start}
کله چې ثابته کړو، مثلاً، \LR{$\Th{Q}$} ټول د پرېکړې وړ خاصيتونه
تمثيل کوي، ملازمه راته وايي چې~\LR{$\Th{Q}$} بايد ناتکميله وي۔
خو ثبوت ئې د خپلواکې تړلے فارمول بېلګه نۀ راکوي؛ يوازې ښيي چې داسې
تړلے فارمول بايد موجوده وي۔ د بېلګې د ترلاسه کولو دپاره بايد نحو
حسابي کړو او د ګوډل د اصلي ثبوت نظر تعقيب کړو۔ او البته لا هم بايد
لومړۍ ادعا ثابته کړو، يعنې دا چې~\LR{$\Th{Q}$} په رښتيا ټول د پرېکړې وړ
خاصيتونه تمثيل کوي۔
\olpseganchor{OLP-0279-B020}{end}

\olpseganchor{OLP-0279-B021}{start}
بايد پام وشي چې هره \emph{په زړه پورې} تيوري ناتکميله يا ناپرېکړه
کېدونکې نۀ ده۔ ډېرې داسې تيورۍ شته چې د په زړه پورې رياضياتي حقايقو
بيانولو ته کافي پياوړې دي، خو د ګوډل د پايلې شرطونه نۀ پوره کوي۔ د
بېلګې په توګه،
\LR{$\Th{Pres} = \Setabs{!A \in \Lang{L_{A^+}}}{\Sat{N}{!A}}$}، يعنې د
ضرب~\LR{$\times$} پرته د حساب د ژبې د هغو تړلي فارمولونه سټ چې په معياري
مدل کښې رښتونې دي، هم بشپړه او هم د پرېکړې وړ ده۔ دې تيورۍ ته د
پرېسبورګر حساب ويل کېږي، او د طبيعي عددونو په باب ټول هغه حقايق
ثابتوي چې يوازې په~\LR{$\Obj{0}$}، \LR{$\prime$} او~\LR{$+$} جوړېدے شي۔
\olpseganchor{OLP-0279-B021}{end}

\olpseganchor{OLP-0280-B004}{start}
\olchapter{inc}{art}{د نحو حسابي کول}
\olpseganchor{OLP-0280-B004}{end}

\olpseganchor{OLP-0280-B005}{start}
\begin{editorial}
  يادونه وکړئ چې د نښه‌لرونکو تابلوګانو دپاره حسابي کول لا نۀ دي
  چمتو شوي۔
\end{editorial}
\olpseganchor{OLP-0280-B005}{end}











\olpseganchor{OLP-0280-B011}{start}
%
  

\olpseganchor{OLP-0280-B011}{end}

\olpseganchor{OLP-0280-B012}{start}
%
  

\olpseganchor{OLP-0280-B012}{end}

\olpseganchor{OLP-0281-B004}{start}
\olfileid{inc}{art}{int}
\olsection{پېژندګلو}
\olpseganchor{OLP-0281-B004}{end}

\olpseganchor{OLP-0281-B005}{start}
د محاسبه‌پوهې او منطق د نښلولو دپاره داسې لار پکار ده چې د منطق د
څيزونو (سمبولونو، ترمونو، فارمولونه او اشتقاقونه)، پر هغو د
عملياتو، او د هغو د خاصيتونو او اړيکو په اړه داسې خبرې پرې وکړو چې
محاسبوي چلند ته برابرې وي۔ دا کار نېغ په نېغه کولے شو: پر سمبولونو،
د سمبولونو پر لړيو او له هغو څخه پر جوړو شويو نورو څيزونو محاسبه
کېدونکې تابعې او اړيکې څېړو۔ ځکه چې د منطقي نحو ټول څيزونه متناهي دي
او د سمبولونو له د شمېر وړ سټ څخه جوړ شوي، نو د محاسبې په ځينو
مدلونو کښې دا شوني دي۔ خو د محاسبې نور مدلونه---لکه بازګشتي تابعې---
يوازې عددونو او د هغو اړيکو او تابعو ته محدود دي۔ برسېره پر دې، په
پاې کښې غواړو چې نحو د ځينو تيوريو دننه هم وڅېړو، په ځانګړي ډول په
هغو تيوريو کښې چې د حساب په ژبه کښې بيان شوې وي۔ په دې حالتونو کښې
د نحو \emph{حسابي کول} اړين دي؛ يعنې نحوي څيزونه، پر هغو عمليات او د
هغو اړيکې په ترتيب د عددونو، حسابي تابعو او حسابي اړيکو په توګه
تمثيل کړو۔ دا مفکوره چې لايبنېڅ ته رسېږي، نحوي څيزونو ته د عددونو
ټاکل دي۔
\olpseganchor{OLP-0281-B005}{end}

\olpseganchor{OLP-0281-B006}{start}
سمبولونو ته د هغو د «کوډونو» په توګه د عددونو ټاکل نسبتاً ساده دي۔
ځينې سمبولونه لږه ستونزه پيدا کوي، ځکه د بېلګې په توګه بې‌شمېره
متغيرونه، او ان د هرې ځاييزې کچې~\LR{$n$} بې‌شمېره تابعې شته۔ خو
البته سمبولونو ته په منظم ډول داسې عددونه ټاکل کېدے شي چې، د بېلګې
په توګه، \LR{$\Obj v_2$} او \LR{$\Obj v_3$} بېل کوډونه ولري۔ د سمبولونو لړۍ
(لکه ترمونه او فارمولونه) تر دې لويه ستونزه ده۔ خو که د عددونو له
لړيو سره سوچه حسابي چلند کولے شو (د بېلګې په توګه د لړيو د اولي
عددونو-د-قوتونو کوډول)، نو د جلا سمبولونو کوډول د سمبولونو د لړيو تر
کوډولو غځولے شو، او بيا ئې د فارمولونه تر لړيو يا نورو ترتيبونو،
لکه اشتقاقونه، نور هم غځولے شو۔ دا غځېدلے کوډول «ګوډل شمېرنه»
بلل کېږي۔ هر ترم، فارمول او اشتقاق ته يوه ګوډل شمېره
ټاکل کېږي۔
\olpseganchor{OLP-0281-B006}{end}

\olpseganchor{OLP-0281-B007}{start}
کله چې د سمبولونو لړۍ د هغو د کوډونو د لړۍ په توګه کوډوو، او د لړيو
د کوډولو داسې نظام غوره کوو چې په محاسبه کېدونکو تابعو پرې چلند کېدے
شي، نو پر ګوډل شمېرو هم په محاسبه کېدونکو تابعو چلند کولے شو۔ په عمل
کښې به ټولې اړوندې تابعې بنسټيزې بازګشتي وي۔ د بېلګې په توګه، د يوې
لړۍ اوږدوالے محاسبه کول او د لړۍ له کوډ څخه د هغې \LR{$i$}-م غړے محاسبه
کول دواړه بنسټيز بازګشتي دي۔ که د لړۍ کوډوونکے عدد، د بېلګې په توګه،
د فارمول~\LR{$!A$} ګوډل شمېره وي، سمدستي وينو چې د فارمول
اوږدوالے او په فارمول کښې د \LR{$i$}-م سمبول (کوډ) هم د~\LR{$!A$} له
ګوډل شمېرې څخه محاسبه کېدے شي۔ د دې ثابتول لږ ستونزمن دي چې، د بېلګې
په توګه، د سم جوړ شوي ترم يا د سم اشتقاق د ګوډل شمېرې په
توګه د يو عدد خاصيت بنسټيز بازګشتي دے۔ بيا هم دا شوني دي، ځکه د پام
وړ لړۍ (ترمونه، فارمولونه او اشتقاقونه) په استقرايي ډول
تعريف شوې دي۔
\olpseganchor{OLP-0281-B007}{end}

\olpseganchor{OLP-0281-B008}{start}
د بېلګې په توګه د تعويض عمليه وګورئ۔ که \LR{$!A$} يو فارمول، \LR{$x$} يو
متغير او \LR{$t$} يو ترم وي، نو \LR{$\Subst{!A}{t}{x}$} په~\LR{$!A$} کښې د~\LR{$x$} د
هرې ازادې پېښې په~\LR{$t$} بدلولو پايله ده۔ اوس فرض کړئ چې \LR{$!A$}، \LR{$x$} او
\LR{$t$} ته مو په ترتيب \LR{$k$}، \LR{$l$} او \LR{$m$} ګوډل شمېرې ټاکلې دي۔ هماغه
طرحه \LR{$\Subst{!A}{t}{x}$} ته، مثلاً، ګوډل شمېره~\LR{$n$} ټاکي۔ دا تطبيق---
چې \LR{$k$}، \LR{$l$} او \LR{$m$} پر~\LR{$n$} تطبيقوي---د تعويض د عمليې حسابي سيال دے۔
کله چې د تعويض عمليه \LR{$!A$}، \LR{$x$} او \LR{$t$} پر
\LR{$\Subst{!A}{t}{x}$} تطبيقوي، د حسابي شوي تعويض تابع د ګوډل شمېرو
\LR{$k$}، \LR{$l$} او \LR{$m$} پر ګوډل شمېرې~\LR{$n$} تطبيقوي۔ وبه وينو چې دا تابع
بنسټيزه بازګشتي ده۔
\olpseganchor{OLP-0281-B008}{end}

\olpseganchor{OLP-0281-B009}{start}
د نحو حسابي کول يوازې انتزاعي ارزښت نۀ لري، که څه هم په لومړي ځل دا
يوه نااشنا ژوره مفکوره وه چې د حساب د ژبې په شان ژبې، چې د ژبو «په
اړه د خبرو» کومه اله ورسره نۀ وي، بيا هم د عباراتو پېچلي خاصيتونه
صوري کولے شي۔ له دې وروسته يوازې يو وړوکے ګام پاتې کېږي چې وپوښتو په
دې ژبه کښې يوه تيوري، لکه د پېانو حساب، د خپلې ژبې په اړه څه
\emph{ثابتولے} شي (د بېلګې په توګه دا چې تړلي فارمولونه د ثبوت وړ يا
رښتيا دي که نۀ)۔ دا موږ د ګوډل (د ناثبوت‌وړتيا په اړه) او تارسکي (د
رښتيا د ناتعريفېدنې په اړه) نومياليو محدودوونکو قضيو ته رسوي۔ خو د
نحو د حسابي کولو چل د محاسبه‌پوهې د ځينو مهمو پايلو د ثبوت دپاره هم
مهم دے، د بېلګې په توګه د تيوريو د محاسبوي ځواک يا د محاسبې د بېلو
مدلونو ترمنځ د اړيکې په اړه۔ د نحو حسابي کول د نورو څيزونو او
خاصيتونو د حسابي کولو د يو مدل په توګه کار کوي۔ د بېلګې په توګه، په
ورته ډول د حالتونو او محاسبو (مثلاً د ټيورينګ ماشينونو د محاسبو)
حسابي کول هم شوني دي۔ په دې سره د محاسبې په يو مدل (لکه ټيورينګ
ماشينونو) کښې محاسبې په بل مدل (لکه بازګشتي تابعو) کښې شبيه کولے شو۔
\olpseganchor{OLP-0281-B009}{end}

\olpseganchor{OLP-0282-B004}{start}
\olfileid{inc}{art}{cod}
\olsection{د سمبولونو کوډول}
\olpseganchor{OLP-0282-B004}{end}

\olpseganchor{OLP-0282-B005}{start}
د لومړۍ درجې منطق بنسټيزه ژبه~\LR{$\Lang L$} دا سمبولونه کاروي:
\[
\lfalse \quad \lnot \quad \lor \quad \land \quad \lif \quad \lforall
\quad \lexists \quad \eq \quad ( \quad ) \quad ,
\]
له دې سره د متغيرونو او ثابت سمبولونه د شمېر وړ سټونه، او د هرې
ځاييزې کچې د تابعې او پريديکاتونه، د شمېر وړ سټونه هم
کاروي۔ موږ
دې سمبولونو هر يوه ته داسې \emph{کوډونه} ټاکلے شو چې هر سمبول ته د
هغه د کوډ په توګه يو ځانګړے عدد ورسېږي، او دوو بېلو سمبولونو ته يو
شان عدد ور نۀ کړل شي۔ پوهېږو چې دا شوني دي، ځکه د ټولو سمبولونو سټ
د شمېر وړ دے، نو د هغه او د طبيعي عددونو د سټ ترمنځ
دوه اړخيزه يو پر يو تابع شته۔ خو دا هم باوري کول غواړو چې سمبول (او د هغه په
اړه ځينې معلومات، لکه د تابع ځاييزه کچه) له کوډ څخه په
محاسبه کېدونکې لار بېرته ترلاسه کړے شو۔ البته د دې کار ډېرې ممکنې
لارې شته۔ دلته يوه داسې لار ده چې بنسټيزې بازګشتي تابعې کاروي۔
(ياد کړئ چې \LR{$\tuple{n_0, \dots, n_k}$} هغه عدد دے چې د عددونو لړۍ
\LR{$n_0$}، \dots، \LR{$n_k$} کوډوي۔)
\olpseganchor{OLP-0282-B005}{end}

\olpseganchor{OLP-0282-B006}{start}
\begin{defn}
که \LR{$s$} د~\LR{$\Lang L$} يو سمبول وي، د هغه \emph{سمبول-کوډ}~\LR{$\scode s$}
داسې تعريف کړئ:
\begin{enumerate}
\item که \LR{$s$} په منطقي سمبولونو کښې وي، \LR{$\scode s$} د لاندې جدول له
  مخې ورکول کېږي:
\[
\begin{array}{cccccccccc}
  \lfalse & \lnot & \lor & \land & \lif & \lforall \\
\tuple{0, 0} & \tuple{0, 1} & \tuple{0, 2} & \tuple{0, 3} &
\tuple{0, 4} & \tuple{0, 5} \\
\lexists & \eq & ( & ) & ,\\
\tuple{0, 6} & \tuple{0, 7} &
\tuple{0, 8} & \tuple{0, 9} & \tuple{0, 10}
\end{array}
\]
\item که \LR{$s$}، \LR{$i$}-م متغير~\LR{$\Obj v_i$} وي، نو
  \LR{$\scode s = \tuple{1, i}$}۔
\item که \LR{$s$}، \LR{$i$}-م ثابت سمبول~\LR{$\Obj c_i$} وي، نو
  \LR{$\scode s = \tuple{2, i}$}۔
\item که \LR{$s$}، \LR{$i$}-م \LR{$n$}-ځايه تابع~\LR{$\Obj f_i^n$} وي، نو
  \LR{$\scode s = \tuple{3, n, i}$}۔
\item که \LR{$s$}، \LR{$i$}-م \LR{$n$}-ځايه پريديکات~\LR{$\Obj P_i^n$} وي، نو
  \LR{$\scode s = \tuple{4, n, i}$}۔
\end{enumerate}
\end{defn}
\olpseganchor{OLP-0282-B006}{end}

\olpseganchor{OLP-0282-B007}{start}
\begin{prop}
لاندې اړيکې بنسټيزې بازګشتي دي:
\begin{enumerate}
\item \LR{$\fn{Fn}(x, n)$} هله صدق کوي چې \LR{$x$} د کوم~\LR{$i$} دپاره د
  \LR{$\Obj f^n_i$} کوډ وي؛ يعنې \LR{$x$} د يوې \LR{$n$}-ځايه تابع کوډ وي۔
\item \LR{$\fn{Pred}(x, n)$} هله صدق کوي چې \LR{$x$} د کوم~\LR{$i$} دپاره د
  \LR{$\Obj P^n_i$} کوډ وي، يا \LR{$x$} د~\LR{$\eq$} کوډ او \LR{$n = 2$} وي؛ يعنې \LR{$x$}
  د يوه \LR{$n$}-ځايه پريديکات کوډ وي۔
\end{enumerate}
\end{prop}
\olpseganchor{OLP-0282-B007}{end}

\olpseganchor{OLP-0282-B008}{start}
\begin{defn}
که \LR{$s_0, \dots, s_{n-1}$} د سمبولونو يوه لړۍ وي، د هغې
\emph{ګوډل شمېره} \LR{$\tuple{\scode{s_0}, \dots, \scode{s_{n-1}}}$} ده۔
\end{defn}
\olpseganchor{OLP-0282-B008}{end}

\olpseganchor{OLP-0282-B009}{start}
\begin{explain}
پام وکړئ چې \emph{کوډونه} او \emph{ګوډل شمېرې} بېل څيزونه دي۔ د
بېلګې په توګه متغير~\LR{$\Obj v_5$} کوډ~\LR{$\scode{\Obj v_5} = \tuple{1, 5}
= 2^2\cdot 3^6$} لري۔ خو که متغير~\LR{$\Obj v_5$} د يوه ترم په توګه
وګڼو، دا د سمبولونو يوه لړۍ (د اوږدوالي~\LR{$1$}) هم ده۔ د
\emph{ترم}~\LR{$\Obj v_5$} \emph{ګوډل شمېره}~\LR{$\Gn{\Obj v_5}$} دا ده:
\LR{$\tuple{\scode{\Obj v_5}} = 2^{\scode{\Obj v_5} + 1} =
2^{2^2\cdot 3^6 + 1}$}۔
\end{explain}
\olpseganchor{OLP-0282-B009}{end}

\olpseganchor{OLP-0282-B010}{start}
\begin{ex}
ياد کړئ چې که \LR{$k_0$}، \dots، \LR{$k_{n-1}$} د عددونو يوه لړۍ وي، نو د
اولي عددونو-د-قوتونو په کوډول کښې د لړۍ کوډ
\LR{$\tuple{k_0, \dots, k_{n-1}}$} دا دے:
\[
2^{k_0+1}\cdot3^{k_1+1}\cdot \dots \cdot p_{n-1}^{k_{n-1}+1},
\]
چې \LR{$p_i$}، \LR{$i$}-م اولي عدد دے (له \LR{$p_0 = 2$} څخه پيلېږي)۔ نو د بېلګې
په توګه فارمول~\LR{$\eq[\Obj v_0][\Obj 0]$}، يا په څرګنده توګه
\LR{${\eq}(\Obj v_0,\Obj c_0)$}، دا ګوډل شمېره لري:
\[
\tuple{\scode{\eq},\scode{(},\scode{\Obj v_0},\scode{,}, \scode{\Obj
    c_0},\scode{)}}.
\]
دلته \LR{$\scode{\eq}$} دا دے: \LR{$\tuple{0,7} = 2^{0+1}\cdot
3^{7+1}$}؛ \LR{$\scode{\Obj v_0}$} دا دے: \LR{$\tuple{1,0} =
2^{1+1}\cdot3^{0+1}$}؛ او همداسې نور۔ نو \LR{$\Gn{=(\Obj v_0,\Obj c_0)}$}
دا دے:
\begin{multline*}
2^{\scode{=} + 1}\cdot 3^{\scode{(}+1}\cdot 5^{\scode{\Obj v_0}+1}
\cdot 7^{\scode{,} + 1} \cdot 11^{\scode{\Obj c_0}+1} \cdot
13^{\scode{)}+1} = \\
2^{2^1\cdot 3^8 + 1}\cdot 3^{2^1\cdot 3^9+1}\cdot 5^{2^2\cdot 3^1+1}
\cdot 7^{2^1\cdot 3^{11} + 1} \cdot 11^{2^3\cdot3^1+1} \cdot
13^{2^1\cdot3^{10}+1} = \\
2^{13\,123}\cdot 3^{39\,367}\cdot 5^{13}\cdot 7^{354\,295}\cdot11^{25}\cdot13^{118\,099}.
\end{multline*}
\end{ex}
\olpseganchor{OLP-0282-B010}{end}

\olpseganchor{OLP-0283-B004}{start}
\olfileid{inc}{art}{trm}
\olsection{د ترمونو کوډول}
\olpseganchor{OLP-0283-B004}{end}

\olpseganchor{OLP-0283-B005}{start}
\begin{explain}
يو ترم يوازې د سمبولونو يو ځانګړے ډول لړۍ ده: دا د ترمونو د جوړښتي
قاعدو له مخې له ثابتونو او متغيرونو څخه په استقرايي ډول جوړېږي۔ ځکه
چې د سمبولونو لړۍ د عددونو په توګه کوډېدلے شي---د سمبولونو د کوډولو
له يوې طرحې او د عددونو د لړيو د کوډولو له يوې لارې سره---نو ترمونو
ته د ګوډل شمېرو ټاکل ستونزمن نۀ دي۔ اصلي ستونزه دا ښودل دي چې يو عدد
چې د سم جوړ شوي ترم ګوډل شمېره وي، دا خاصيت ئې محاسبه کېدونکے، بلکې
بنسټيز بازګشتي دے۔
\end{explain}
\olpseganchor{OLP-0283-B005}{end}

\olpseganchor{OLP-0283-B006}{start}
متغيرونه او ثابت سمبولونه تر ټولو ساده ترمونه دي، او دا ازمويل
چې \LR{$x$} د داسې يوه ترم ګوډل شمېره ده، اسانه دي: \LR{$\fn{Var}(x)$} هله
صدق کوي چې \LR{$x$} د کوم~\LR{$i$} دپاره \LR{$\Gn{\Obj v_i}$} وي۔ په بله وينا،
\LR{$x$} د اوږدوالي~\LR{$1$} يوه لړۍ ده او يوازينے غړے ئې \LR{$(x)_0$} د کوم
متغير~\LR{$\Obj v_i$} کوډ دے؛ يعنې \LR{$x$} د کوم~\LR{$i$} دپاره
\LR{$\tuple{\tuple{1, i}}$} دے۔ په ورته ډول \LR{$\fn{Const}(x)$} هله صدق کوي
چې \LR{$x$} د کوم~\LR{$i$} دپاره \LR{$\Gn{\Obj c_i}$} وي۔ دا دواړه اړيکې بنسټيزې
بازګشتي دي، ځکه که داسې \LR{$i$} موجود وي، بايد \LR{$<x$} وي:
\begin{align*}
  \fn{Var}(x) & \defiff \bexists{i<x}{x = \tuple{\tuple{1, i}}}\\
  \fn{Const}(x) & \defiff \bexists{i<x}{x = \tuple{\tuple{2, i}}}
\end{align*}
\olpseganchor{OLP-0283-B006}{end}

\olpseganchor{OLP-0283-B007}{start}
\begin{prop}
\ollabel{prop:term-primrec}
هغه اړيکې \LR{$\fn{Term}(x)$} او \LR{$\fn{ClTerm}(x)$} چې په ترتيب هله صدق کوي
چې \LR{$x$} د يوه ترم يا تړلي ترم ګوډل شمېره وي، بنسټيزې بازګشتي دي۔
\end{prop}
\olpseganchor{OLP-0283-B007}{end}

\olpseganchor{OLP-0283-B008}{start}
\begin{proof}
د سمبولونو يوه لړۍ~\LR{$s$} هله ترم دے چې د ترمونو داسې لړۍ~\LR{$s_0$}،
\dots، \LR{$s_{k-1}=s$} موجوده وي چې ثبتوي ترم~\LR{$s$} د ترمونو د جوړښتي
قاعدو له مخې له ثابت سمبولونه او متغيرونه څخه څنګه جوړ شوے دے۔
د دې د بيانولو دپاره چې داسې يوه اټکلي جوړښتي لړۍ جوړښتي قاعدې
تعقيبوي، بايد د هر \LR{$i<k$} دپاره له لاندې څخه يو حالت سم وي:
\begin{enumerate}
\item \LR{$s_i$} يو متغير~\LR{$\Obj v_j$} دے؛ يا
\item \LR{$s_i$} يو ثابت سمبول~\LR{$\Obj c_j$} دے؛ يا
\item \LR{$s_i$} له هغو \LR{$n$} ترمونو \LR{$t_1$}، \dots، \LR{$t_n$} څخه، چې له
  ځاے~\LR{$i$} وړاندې راغلي وي، د يوه \LR{$n$}-ځايه
  تابع~\LR{$\Obj f^n_j$} په کارولو جوړ شوے دے۔
\end{enumerate}
د دې د ښودلو دپاره چې د ګوډل شمېرو اړونده اړيکه بنسټيزه بازګشتي
ده، دا شرط بايد په بنسټيز بازګشتي ډول بيان کړو؛ يعنې بنسټيزې
بازګشتي تابعې، اړيکې او محدود کمیت ايښودنه وکاروو۔
\olpseganchor{OLP-0283-B008}{end}

\olpseganchor{OLP-0283-B009}{start}
فرض کړئ \LR{$y$} هغه عدد دے چې لړۍ~\LR{$s_0$}، \dots، \LR{$s_{k-1}$} کوډوي؛ يعنې
\LR{$y = \tuple{\Gn{s_0}, \dots, \Gn{s_{k-1}}}$}۔ دا د ګوډل شمېرې~\LR{$x$}
لرونکي ترم دپاره هله جوړښتي لړۍ کوډوي چې د هر \LR{$i<k$} دپاره:
\begin{enumerate}
\item \LR{$\fn{Var}((y)_i)$}؛ يا
\item \LR{$\fn{Const}((y)_i)$}؛ يا
\item داسې \LR{$n$} او عدد~\LR{$z=\tuple{z_1, \dots, z_n}$} شته چې هر \LR{$z_l$}
  د کوم \LR{$i'<i$} دپاره له \LR{$(y)_{i'}$} سره برابر وي، او
\[
(y)_i = \Gn{\Obj f^n_j(} \concat \fn{flatten}(z) \concat \Gn{)},
\]
\end{enumerate}
او سربېره پر دې \LR{$(y)_{k-1}=x$} وي۔ (تابع \LR{$\fn{flatten}(z)$} لړۍ
\LR{$\tuple{\Gn{t_1}, \dots, \Gn{t_n}}$} په \LR{$\Gn{t_1, \dots, t_n}$}
بدلوي او بنسټيزه بازګشتي ده۔)
\olpseganchor{OLP-0283-B009}{end}

\olpseganchor{OLP-0283-B010}{start}
په (3) کښې شاخصونه \LR{$j$} او \LR{$n$}، د ترمونو \LR{$t_l$} ګوډل شمېرې~\LR{$z_l$}،
او د لړۍ~\LR{$\tuple{z_1, \dots, z_n}$} کوډ~\LR{$z$} ټول له~\LR{$y$} څخه واړه دي۔
پورته \LR{$k$} په \LR{$\len{y}$} بدلولے شو۔ نو «\LR{$y$} د ګوډل شمېرې~\LR{$x$} لرونکي
ترم د جوړښتي لړۍ کوډ دے» په داسې ډول بيانولے شو چې وښيي دا اړيکه
بنسټيزه بازګشتي ده۔
\olpseganchor{OLP-0283-B010}{end}

\olpseganchor{OLP-0283-B011}{start}
اوس يوازې دا باور ترلاسه کول پکار دي چې پر~\LR{$y$} يو بنسټيز بازګشتي حد
شته۔ خو که \LR{$x$} د يوه ترم ګوډل شمېره وي، نو هغه بايد داسې جوړښتي لړۍ
ولري چې تر ټولو زيات \LR{$\len{x}$} ترمونه ولري (ځکه د~\LR{$s$} په جوړښتي
لړۍ کښې هر ترم بايد د~\LR{$s$} په کوم ځاے کښې پيل شي، او دوه فرعي ترمونه
په يوه ځاے نۀ شي پيلېدے)۔ البته د~\LR{$s$} د هر فرعي ترم ګوډل شمېره
\LR{$\le x$} ده۔ نو تل داسې جوړښتي لړۍ شته چې کوډ ئې
\LR{$\le p_{k-1}^{k(x+1)}$} وي، چې \LR{$k=\len{x}$} دے۔
\olpseganchor{OLP-0283-B011}{end}

\olpseganchor{OLP-0283-B012}{start}
د \LR{$\fn{ClTerm}$} دپاره يوازې د متغيرونه فقره پرېږدئ۔
\end{proof}
\olpseganchor{OLP-0283-B012}{end}

\olpseganchor{OLP-0283-B013}{start}
\begin{prob}
وښيئ چې تابع \LR{$\fn{flatten}(z)$}، چې لړۍ
\LR{$\tuple{\Gn{t_1}, \dots, \Gn{t_n}}$} په \LR{$\Gn{t_1, \dots, t_n}$}
بدلوي، بنسټيزه بازګشتي ده۔
\end{prob}
\olpseganchor{OLP-0283-B013}{end}

\olpseganchor{OLP-0283-B014}{start}
\begin{prop}\ollabel{prop:num-primrec}
  تابع \LR{$\fn{num}(n) = \Gn{\num{n}}$} بنسټيزه بازګشتي ده۔
\end{prop}
\olpseganchor{OLP-0283-B014}{end}

\olpseganchor{OLP-0283-B015}{start}
\begin{proof}
  \LR{$\fn{num}(n)$} په بنسټيز بازګښت داسې تعريفوو:
  \begin{align*}
    \fn{num}(0) & = \Gn{\Obj 0}\\
    \fn{num}(n+1) & = \Gn{\prime(} \concat \fn{num}(n) \concat \Gn{)}.
  \end{align*}
\end{proof}
\olpseganchor{OLP-0283-B015}{end}

\olpseganchor{OLP-0284-B004}{start}
\olfileid{inc}{art}{frm}
\olsection{د فارمولونه کوډول}
\olpseganchor{OLP-0284-B004}{end}

\olpseganchor{OLP-0284-B005}{start}
کله چې اړيکه \LR{$\fn{Term}(x)$} مو په بنسټيز بازګشتي ډول تعريف کړه، د
هغې په کارولو د فارمولونه اړونده اړيکه \LR{$\fn{Frm}(x)$} هم په
بنسټيز بازګشتي ډول تعريفولے شو۔
\olpseganchor{OLP-0284-B005}{end}

\olpseganchor{OLP-0284-B006}{start}
\begin{prop}
هغه اړيکه \LR{$\fn{Atom}(x)$} چې هله صدق کوي چې \LR{$x$} د يوه اټومي
فارمول ګوډل شمېره وي، بنسټيزه بازګشتي ده۔
\end{prop}
\olpseganchor{OLP-0284-B006}{end}

\olpseganchor{OLP-0284-B007}{start}
\begin{proof}
عدد \LR{$x$} هله د يوه اټومي فارمول ګوډل شمېره ده چې له لاندې څخه
يو حالت صدق وکړي:
\begin{enumerate}
\item داسې \LR{$n$}، \LR{$j<x$} او \LR{$z<x$} شته چې د هر \LR{$i<n$} دپاره
  \LR{$\fn{Term}((z)_i)$} وي او \LR{$x = $}
\[
\Gn{\Obj P^n_j(} \concat \fn{flatten}(z) \concat \Gn{)}.
\]
\item داسې \LR{$z_1,z_2<x$} شته چې \LR{$\fn{Term}(z_1)$}،
  \LR{$\fn{Term}(z_2)$} وي، او \LR{$x = $}
\[
\Gn{{\eq}(} \concat z_1 \concat \Gn{,} \concat z_2 \concat{\Gn{)}}.
\]
  \item \LR{$x = \Gn{\lfalse}$}۔
  \item \LR{$x = \Gn{\ltrue}$}۔
\end{enumerate}
\end{proof}
\olpseganchor{OLP-0284-B007}{end}

\olpseganchor{OLP-0284-B008}{start}
\begin{prop}
\ollabel{prop:frm-primrec}
هغه اړيکه \LR{$\fn{Frm}(x)$} چې هله صدق کوي چې \LR{$x$} د فارمول ګوډل
شمېره وي، بنسټيزه بازګشتي ده۔
\end{prop}
\olpseganchor{OLP-0284-B008}{end}

\olpseganchor{OLP-0284-B009}{start}
\begin{proof}
د سمبولونو لړۍ~\LR{$s$} هله فارمول ده چې د فارمول داسې
جوړښتي لړۍ~\LR{$s_0$}، \dots، \LR{$s_{k-1}=s$} موجوده وي چې ثبتوي~\LR{$s$} د
جوړښتي قاعدو له مخې له اټومي فارمولونه څخه څنګه جوړ شوے دے۔ د هر
\LR{$s_i$} کوډ (او په حقيقت کښې د لړۍ \LR{$\tuple{s_0, \dots, s_{k-1}}$} کوډ
هم) د~\LR{$s$} له کوډ~\LR{$x$} څخه وړوکے دے۔
\end{proof}
\olpseganchor{OLP-0284-B009}{end}

\olpseganchor{OLP-0284-B010}{start}
\begin{prob}
د \olref[inc][art][frm]{prop:frm-primrec} تفصيلي ثبوت د
\olref[inc][art][trm]{prop:term-primrec} د لومړي ثبوت په څېر ورکړئ۔
\end{prob}
\olpseganchor{OLP-0284-B010}{end}

\olpseganchor{OLP-0284-B011}{start}
\begin{prop}
\ollabel{prop:freeocc-primrec}
اړيکه \LR{$\fn{FreeOcc}(x,z,i)$}، چې هله صدق کوي چې د ګوډل شمېرې~\LR{$x$}
لرونکي فارمول \LR{$i$}-م سمبول د ګوډل شمېرې~\LR{$z$} لرونکي متغير يوه ازاده
پېښه وي، بنسټيزه بازګشتي ده۔
\end{prop}
\olpseganchor{OLP-0284-B011}{end}

\olpseganchor{OLP-0284-B012}{start}
\begin{proof}
تمرین۔
\end{proof}
\olpseganchor{OLP-0284-B012}{end}

\olpseganchor{OLP-0284-B013}{start}
\begin{prob}
\olref[inc][art][frm]{prop:freeocc-primrec} ثابت کړئ۔ دا حقيقت
کارولے شئ چې د فارمول هر هغه فرعي تار چې پخپله فارمول
وي، د هغه فرعي-فارمول دے۔
\end{prob}
\olpseganchor{OLP-0284-B013}{end}

\olpseganchor{OLP-0284-B014}{start}
\begin{prop}
هغه خاصيت \LR{$\fn{Sent}(x)$} چې هله صدق کوي چې \LR{$x$} د تړلے فارمول
ګوډل شمېره وي، بنسټيز بازګشتي دے۔
\end{prop}
\olpseganchor{OLP-0284-B014}{end}

\olpseganchor{OLP-0284-B015}{start}
\begin{proof}
تړلے فارمول داسې فارمول دے چې د متغيرونه ازادې پېښې
نۀ لري۔ نو \LR{$\fn{Sent}(x)$} هله صدق کوي چې
\begin{multline*}
\bforall{i<\len{x}}{\bforall{z<x}{}}\\
(\bexists{j<z}{z=\Gn{\Obj v_j}} \lif \lnot\fn{FreeOcc}(x,z,i)).
\end{multline*}
\end{proof}
\olpseganchor{OLP-0284-B015}{end}

\olpseganchor{OLP-0285-B004}{start}
\olfileid{inc}{art}{sub}
\olsection{تعويض}
\olpseganchor{OLP-0285-B004}{end}

\olpseganchor{OLP-0285-B005}{start}
ياد کړئ چې تعويض هغه عمليه ده چې په فارمول~\LR{$!A$} کښې د
متغير~\LR{$u$} ټولې ازادې پېښې په ترم~\LR{$t$} بدلوي، او
\LR{$\Subst{!A}{t}{u}$} ليکل کېږي۔ کله چې دا عمليه د متغيرونه،
فارمولونه او ترمونو پر ګوډل شمېرو ترسره شي، بنسټيزه بازګشتي ده۔
\olpseganchor{OLP-0285-B005}{end}

\olpseganchor{OLP-0285-B006}{start}
\begin{prop}\ollabel{prop:subst-primrec}
داسې بنسټيزه بازګشتي تابع \LR{$\fn{Subst}(x,y,z)$} شته چې دا خاصيت لري:
\[
\fn{Subst}(\Gn{!A}, \Gn{t}, \Gn{u}) = \Gn{\Subst{!A}{t}{u}}.
\]
\end{prop}
\olpseganchor{OLP-0285-B006}{end}

\olpseganchor{OLP-0285-B007}{start}
\begin{proof}
بيا تابع \LR{$\fn{hSubst}$} په بنسټيز بازګښت داسې تعريفولے شو:
\begin{multline*}
\begin{aligned}
\fn{hSubst}(x, y, z, 0) & = \emptyseq \\
\fn{hSubst}(x, y, z, i+1) & =
\end{aligned}\\
\begin{cases}
\fn{hSubst}(x, y, z, i) \concat y & \text{\RL{که \LR{$\fn{FreeOcc}(x, z, i)$} وي}} \\
\fn{append}(\fn{hSubst}(x, y, z, i), (x)_{i}) & \text{\RL{که نۀ وي۔}}
\end{cases}
\end{multline*}
اوس \LR{$\fn{Subst}(x,y,z)$} د \LR{$\fn{hSubst}(x,y,z,\len{x})$} په توګه
تعريفېدے شي۔
\end{proof}
\olpseganchor{OLP-0285-B007}{end}

\olpseganchor{OLP-0285-B008}{start}
\begin{prop}
\ollabel{prop:free-for}
اړيکه \LR{$\fn{FreeFor}(x,y,z)$}، چې هله صدق کوي چې د ګوډل شمېرې~\LR{$y$}
لرونکے ترم د ګوډل شمېرې~\LR{$z$} لرونکي متغير دپاره د ګوډل شمېرې~\LR{$x$}
لرونکي فارمول کښې ازاد وي، بنسټيزه بازګشتي ده۔
\end{prop}
\olpseganchor{OLP-0285-B008}{end}

\olpseganchor{OLP-0285-B009}{start}
\begin{proof} تمرین۔ \end{proof}
\olpseganchor{OLP-0285-B009}{end}

\olpseganchor{OLP-0285-B010}{start}
\begin{prob}
\olref[inc][art][sub]{prop:free-for} ثابت کړئ۔
\end{prob}
\olpseganchor{OLP-0285-B010}{end}

\olpseganchor{OLP-0286-B004}{start}
\olfileid{inc}{art}{plk}
\olsection{په \texorpdfstring{\LR{$\Log{LK}$}}{LK} کښې اشتقاقونه}
\olpseganchor{OLP-0286-B004}{end}

\olpseganchor{OLP-0286-B005}{start}
\begin{explain}
د اشتقاقونه د حسابي کولو دپاره بايد اشتقاقونه د عددونو
په توګه تمثيل کړو۔ ځکه چې اشتقاقونه د سېکوېنټونو داسې ونې دي
چې هر استنتاج يوه نښه هم لري، نو بازګشتي تمثيل تر ټولو څرګنده لار
ده: يو اشتقاق د يوې څوګونې په توګه تمثيلوو چې اجزا ئې
پاې-سېکوېنټ، نښه، او د وروستي استنتاج مقدماتو ته د رسېدونکو
فرعي-اشتقاقونه تمثيلونه دي۔
\end{explain}
\olpseganchor{OLP-0286-B005}{end}

\olpseganchor{OLP-0286-B006}{start}
\begin{defn}
که \LR{$\Gamma$} د تړلي فارمولونه يوه متناهي لړۍ وي،
\LR{$\Gamma=\tuple{!A_1, \dots, !A_n}$}، نو
\LR{$\Gn{\Gamma}=\tuple{\Gn{!A_1}, \dots, \Gn{!A_n}}$}۔
\olpseganchor{OLP-0286-B006}{end}

\olpseganchor{OLP-0286-B007}{start}
که \LR{$\Gamma \Sequent \Delta$} يو سېکوېنټ وي، نو د
\LR{$\Gamma \Sequent \Delta$} ګوډل شمېره دا ده:
\[
\Gn{\Gamma \Sequent \Delta} = \tuple{\Gn{\Gamma}, \Gn{\Delta}}
\]
\olpseganchor{OLP-0286-B007}{end}

\olpseganchor{OLP-0286-B008}{start}
که \LR{$\pi$} په \LR{$\Log{LK}$} کښې اشتقاق وي، \LR{$\Gn{\pi}$} داسې
تعريفېږي:
\begin{enumerate}
\item که \LR{$\pi$} يوازې له ابتدايي سېکوېنټ \LR{$\Gamma \Sequent \Delta$}
  څخه جوړ وي، نو \LR{$\Gn{\pi}$} دا دے:
  \[
    \tuple{0, \Gn{\Gamma \Sequent \Delta}}.
  \]
\item که \LR{$\pi$} په داسې استنتاج پاې ته رسېږي چې يو يا دوه مقدمات
  لري، پايله ئې \LR{$\Gamma \Sequent \Delta$} وي، او \LR{$\pi_1$} او \LR{$\pi_2$}
  هغه سمدستي فرعي ثبوتونه وي چې د وروستي استنتاج پر مقدماتو پاې ته
  رسېږي، نو \LR{$\Gn{\pi}$} په ترتيب له دې دوو څخه يو دے:
  \begin{align*}
    & \tuple{1, \Gn{\pi_1}, \Gn{\Gamma \Sequent \Delta}, k} \text{\RL{ يا}}\\
    & \tuple{2, \Gn{\pi_1}, \Gn{\pi_2}, \Gn{\Gamma \Sequent \Delta}, k},
  \end{align*}
  چې \LR{$k$} د لاندې جدول له مخې د وروستي استنتاج د کارول شوې قاعدې
  له مخې ټاکل کېږي:
\olpseganchor{OLP-0286-B008}{end}

\olpseganchor{OLP-0286-B009}{start}
  \begin{tabular}{lccccccc}
    \text{\RL{قاعده:}} & \LeftR{\Weakening} & \RightR{\Weakening} &
    \LeftR{\Contraction} & \RightR{\Contraction} &
    \LeftR{\Exchange} & \RightR{\Exchange} \\
    \LR{$k$}: & 1 & 2 & 3 & 4 & 5 & 6 \\[2ex]
    \text{\RL{قاعده:}} & \LeftR{\lnot} & \RightR{\lnot} &
    \LeftR{\land} & \RightR{\land} &
    \LeftR{\lor} & \RightR{\lor} \\
    \LR{$k$}: & 7 & 8 & 9 & 10 & 11 & 12 \\[2ex]
    \text{\RL{قاعده:}} & \LeftR{\lif} & \RightR{\lif} &
    \LeftR{\lforall} & \RightR{\lforall} &
    \LeftR{\lexists} & \RightR{\lexists} \\
    \LR{$k$}: & 13 & 14 & 15 & 16 & 17 & 18 \\[2ex]
    \text{\RL{قاعده:}} & \Cut & = \\
    \LR{$k$}: & 19 & 20
  \end{tabular}
\end{enumerate}
\end{defn}
\olpseganchor{OLP-0286-B009}{end}

\olpseganchor{OLP-0286-B010}{start}
\begin{ex}
  دا ډېر ساده اشتقاق وګورئ:
  \begin{prooftree}
    \Axiom$!A \fCenter !A$
    \RightLabel{\LeftR{\land}}
    \UnaryInf$!A \land !B \fCenter !A$
    \RightLabel{\RightR{\lif}}
    \UnaryInf$\fCenter (!A \land !B) \lif !A$
  \end{prooftree}
  د ابتدايي سېکوېنټ ګوډل شمېره به
  \LR{$p_0=\tuple{0, \Gn{!A \Sequent !A}}$} وي۔ د هغه اشتقاق
  ګوډل شمېره چې د~\LR{$\LeftR{\land}$} پر پايله ختمېږي،
  \LR{$p_1=\tuple{1,p_0,\Gn{!A \land !B \Sequent !A},9}$} ده
  (\LR{$1$} ځکه چې \LR{$\LeftR{\land}$} يو مقدمه لري؛ بيا د پايلې
  \LR{$!A \land !B \Sequent !A$} ګوډل شمېره؛ او \LR{$9$} هغه عدد دے چې
  \LR{$\LeftR{\land}$} کوډوي)۔ بيا د ټول اشتقاق ګوډل شمېره
  \LR{$\tuple{1,p_1,\Gn{\Sequent (!A \land !B) \lif !A},14}$} ده؛ يعنې
  \[
  \tuple{1, \tuple{1, \tuple{0, \Gn{!A \Sequent !A}}, \Gn{!A \land !B \Sequent !A}, 9},
    \Gn{\Sequent (!A \land !B) \lif !A}, 14}.
  \]
\end{ex}
\olpseganchor{OLP-0286-B010}{end}

\olpseganchor{OLP-0286-B011}{start}
\begin{explain}
چې د اشتقاقونه پر يو تمثيل هوکړې ته ورسېدو، دا هم بايد وښيو چې
داسې اشتقاقونه په بنسټيز بازګشتي ډول اداره کولے شو، او د هغو
اړين خاصيتونه او اړيکې هم په همدې ډول بيانولے شو۔ ځينې عمليات ساده
دي: د بېلګې په توګه، د اشتقاق له ګوډل شمېرې~\LR{$p$} څخه
\LR{$\fn{EndSequent}(p)=(p)_{(p)_0+1}$} د هغه د پاې-سېکوېنټ ګوډل شمېره
راکوي، او \LR{$\fn{LastRule}(p)=(p)_{(p)_0+2}$} د هغه د وروستۍ قاعدې
کوډ راکوي۔ خاصيت \LR{$\fn{Sequent}(s)$} چې داسې تعريفېږي:
\[
  \len{s} = 2 \land \bforall{i<\len{(s)_0} + \len{(s)_1}}{\fn{Sent}(((s)_0 \concat (s)_1)_i)}
\]
هله د~\LR{$s$} په اړه صدق کوي چې \LR{$s$} د تړلي فارمولونه لرونکي يوه
سېکوېنټ ګوډل شمېره وي۔ ځينې نور عمليات ډېر ستونزمن دي۔ لږ تر لږه د
هغو د ترسره کولو لار به راونغاړو۔ موخه دا ده چې وښيو اړيکه «\LR{$\pi$}
له~\LR{$\Gamma$} څخه د~\LR{$!A$} يو اشتقاق دے» د~\LR{$\pi$} او~\LR{$!A$} د
ګوډل شمېرو بنسټيزه بازګشتي اړيکه ده۔
\end{explain}
\olpseganchor{OLP-0286-B011}{end}

\olpseganchor{OLP-0286-B012}{start}
\begin{prop}\ollabel{prop:followsby}
  خاصيت \LR{$\fn{Correct}(p)$}، چې هله صدق کوي چې د ګوډل شمېرې~\LR{$p$}
  لرونکي اشتقاق~\LR{$\pi$} وروستے استنتاج سم وي، بنسټيز بازګشتي
  دے۔
\end{prop}
\olpseganchor{OLP-0286-B012}{end}

\olpseganchor{OLP-0286-B013}{start}
\begin{proof}
  \LR{$\Gamma \Sequent \Delta$} هله ابتدايي سېکوېنټ دے چې يا داسې
  تړلے فارمول~\LR{$!A$} وي چې \LR{$\Gamma \Sequent \Delta$} همدا
  \LR{$!A \Sequent !A$} وي، يا داسې ترم~\LR{$t$} وي چې
  \LR{$\Gamma \Sequent \Delta$} همدا \LR{$\emptyset \Sequent \eq[t][t]$} وي۔
  د ګوډل شمېرو په ژبه، \LR{$\fn{InitSeq}(s)$} هله صدق کوي چې
  \begin{align*}
  \bexists{x < s}{} (\fn{Sent}(x) & \land
  s = \tuple{\tuple{x},\tuple{x}})
  \lor {}\\
  \bexists{t<s}{} (\fn{Term}(t) & \land
  s = \tuple{0, \tuple{\Gn{{\eq}(} \concat t \concat \Gn{,} \concat t \concat \Gn{)}}}).
  \end{align*}
\olpseganchor{OLP-0286-B013}{end}

\olpseganchor{OLP-0286-B014}{start}
  د استنتاج د هرې قاعدې~\LR{$R$} دپاره دا هم بايد وښيو چې اړيکه
  \LR{$\fn{FollowsBy}_R(p)$} بنسټيزه بازګشتي ده؛
  \LR{$\fn{FollowsBy}_R(p)$} هله صدق کوي چې \LR{$p$}
  د اشتقاق~\LR{$\pi$} ګوډل شمېره وي، او د~\LR{$\pi$} پاې-سېکوېنټ د
  \LR{$R$} په سم پلي کولو د~\LR{$\pi$} د سمدستي فرعي-اشتقاقونه له
  پاې-سېکوېنټونو څخه راشي۔
\olpseganchor{OLP-0286-B014}{end}

\olpseganchor{OLP-0286-B015}{start}
  يو ساده حالت د \RightR{\land} قاعده ده۔ که \LR{$\pi$} په يوه سم
  \LR{$\RightR{\land}$} استنتاج ختم شي، داسې ښکاري:
  \begin{prooftree}
    \AxiomC{}
    \RightLabel{\LR{$\pi_1$}}
    \Deduce$\Gamma \fCenter \Delta, !A$
\olpseganchor{OLP-0286-B015}{end}

\olpseganchor{OLP-0286-B016}{start}
    \AxiomC{}
    \RightLabel{\LR{$\pi_2$}}
    \Deduce$\Gamma \fCenter \Delta, !B$
\olpseganchor{OLP-0286-B016}{end}

\olpseganchor{OLP-0286-B017}{start}
    \RightLabel{\RightR\land}
    \BinaryInf$\Gamma \fCenter \Delta, !A \land !B$
  \end{prooftree}
  نو په اشتقاق~\LR{$\pi$} کښې وروستے استنتاج هله د
  \LR{$\RightR{\land}$} سم پلي کول دي چې د تړلي فارمولونه داسې لړۍ
  \LR{$\Gamma$} او \LR{$\Delta$}، او داسې دوه تړلي فارمولونه \LR{$!A$} او~\LR{$!B$} شته
  چې د~\LR{$\pi_1$} پاې-سېکوېنټ \LR{$\Gamma \Sequent \Delta,!A$}، د~\LR{$\pi_2$}
  پاې-سېکوېنټ \LR{$\Gamma \Sequent \Delta,!B$}، او د~\LR{$\pi$} پاې-سېکوېنټ
  \LR{$\Gamma \Sequent \Delta,!A\land !B$} وي۔ دا يوازې د ګوډل شمېرو
  ژبې ته ژباړل پکار دي۔ که \LR{$s=\Gn{\Gamma \Sequent \Delta}$} وي، نو
  \LR{$(s)_0=\Gn{\Gamma}$} او \LR{$(s)_1=\Gn{\Delta}$} دي۔ نو
  \LR{$\fn{FollowsBy}_{\RightR{\land}}(p)$} هله صدق کوي چې
\begin{align*}
& \bexists{g < p}{\bexists{d < p}{\bexists{a < p}{\bexists{b < p}{\quad}}}} \\
& \qquad \fn{EndSequent}(p) =
  \tuple{g, d \concat
    \tuple{\Gn{(} \concat a \concat \Gn{\land} \concat b \concat \Gn{)}}} \land {} \\
& \qquad \fn{EndSequent}((p)_1) = \tuple{g, d \concat \tuple{a}} \land {}\\
& \qquad \fn{EndSequent}((p)_2) = \tuple{g, d \concat \tuple{b}} \land {}\\
& \qquad (p)_0 = 2 \land \fn{LastRule}(p) = 10.
\end{align*}
جلا کرښې په ترتيب دا بيانوي: «داسې لړۍ~(\LR{$\Gamma$}) شته چې ګوډل
شمېره ئې~\LR{$g$} ده، داسې لړۍ~(\LR{$\Delta$}) شته چې ګوډل شمېره ئې~\LR{$d$} ده،
داسې فارمول~(\LR{$!A$}) شته چې ګوډل شمېره ئې~\LR{$a$} ده، او داسې
فارمول~(\LR{$!B$}) شته چې ګوډل شمېره ئې~\LR{$b$} ده»؛ داسې چې «د~\LR{$\pi$}
پاې-سېکوېنټ \LR{$\Gamma \Sequent \Delta,!A\land !B$} دے»، «د~\LR{$\pi_1$}
پاې-سېکوېنټ \LR{$\Gamma \Sequent \Delta,!A$} دے»، «د~\LR{$\pi_2$}
پاې-سېکوېنټ \LR{$\Gamma \Sequent \Delta,!B$} دے»، او «\LR{$\pi$} دوه
سمدستي فرعي ثبوتونه لري او وروستۍ استنتاجي قاعده
\LR{$\RightR\land$} (په شمېره~\LR{$10$}) ده»۔
\olpseganchor{OLP-0286-B017}{end}

\olpseganchor{OLP-0286-B018}{start}
په~\LR{$\pi$} کښې وروستے استنتاج هله د \LR{$\RightR\lexists$} سم پلي کول
دي چې داسې لړۍ \LR{$\Gamma$} او \LR{$\Delta$}، داسې فارمول~\LR{$!A$}، متغير
\LR{$x$} او ترم~\LR{$t$} شته چې د~\LR{$\pi$} پاې-سېکوېنټ
\LR{$\Gamma \Sequent \Delta,\lexists[x][!A]$} او د~\LR{$\pi_1$}
پاې-سېکوېنټ \LR{$\Gamma \Sequent \Delta,\Subst{!A}{t}{x}$} وي۔ نو د
ګوډل شمېرو په ژبه \LR{$\fn{FollowsBy}_{\RightR\lexists}(p)$} هله صدق کوي
چې
\begin{align*}
  & \bexists{g<p}{\bexists{d<p}{\bexists{a<p}{\bexists{x<p}{\bexists{t<p}{\quad}}}}}\\
  & \qquad \fn{EndSequent}(p) =
  \tuple{
    g,
    d \concat \tuple{\Gn{\lexists} \concat x \concat a}
  } \land {}\\
  & \qquad \fn{EndSequent}((p)_1) =
  \tuple{
    g,
    d \concat \tuple{\fn{Subst}(a, t, x)}
  } \land {}\\
  & \qquad (p)_0 = 1 \land \fn{LastRule}(p) = 18.
\end{align*}
بيا \LR{$\fn{Correct}(p)$} داسې تعريفوو:
\begin{multline*}
  \fn{Sequent}(\fn{EndSequent}(p)) \land {}\\
  [(\fn{LastRule}(p) = 1 \land
  \fn{FollowsBy}_{\LeftR\Weakening}(p)) \lor \dots \lor {}\\
  (\fn{LastRule}(p) = 20 \land \fn{FollowsBy}_{\eq}(p)) \lor {}\\
  (p)_0 = 0 \land \fn{InitialSeq}(\fn{EndSequent}(p))]
\end{multline*}
لومړۍ کرښه باوري کوي چې د~\LR{$p$} پاې-سېکوېنټ په رښتيا د
تړلي فارمولونه يو سېکوېنټ دے۔ وروستۍ کرښه هغه حالت رانغاړي چې \LR{$p$}
يوازې يو ابتدايي سېکوېنټ وي۔
\end{proof}
\olpseganchor{OLP-0286-B018}{end}

\olpseganchor{OLP-0286-B019}{start}
\begin{prob}
  لاندې خاصيتونه د \olref[inc][art][plk]{prop:followsby} په څېر
  تعريف کړئ:
  \begin{enumerate}
  \item \LR{$\fn{FollowsBy}_{\Cut}(p)$}؛
  \item \LR{$\fn{FollowsBy}_{\LeftR\lif}(p)$}؛
  \item \LR{$\fn{FollowsBy}_{\eq}(p)$}؛
  \item \LR{$\fn{FollowsBy}_{\RightR{\lforall}}(p)$}۔
  \end{enumerate}
  د وروستي خاصيت دپاره دا هم بايد وښيئ چې په بنسټيز بازګشتي ډول
  ازمويلے شئ چې د ګوډل شمېرې~\LR{$p$} لرونکي اشتقاق وروستے
  استنتاج د ځانګړي متغير شرط پوره کوي؛ يعنې د~\LR{$\RightR{\lforall}$}
  ځانګړے متغير~\LR{$a$} په پاې-سېکوېنټ کښې نۀ راځي۔
\end{prob}
\olpseganchor{OLP-0286-B019}{end}

\olpseganchor{OLP-0286-B020}{start}
\begin{prop}
  \ollabel{prop:deriv}
  اړيکه \LR{$\fn{Deriv}(p)$}، چې هله صدق کوي چې \LR{$p$} د يوه سم
  اشتقاق~\LR{$\pi$} ګوډل شمېره وي، بنسټيزه بازګشتي ده۔
\end{prop}
\olpseganchor{OLP-0286-B020}{end}

\olpseganchor{OLP-0286-B021}{start}
\begin{proof}
  اشتقاق~\LR{$\pi$} هله سم دے چې هر استنتاج ئې د يوې قاعدې سم
  پلي کول وي؛ يعنې د هغه هر فرعي-اشتقاقونه په سم استنتاج پاې ته
  ورسېږي۔ نو \LR{$\fn{Deriv}(p)$} هله صدق کوي چې
  \[
  \bforall{i<\len{\fn{SubtreeSeq}(p)}}{\fn{Correct}((\fn{SubtreeSeq}(p))_i)}.
  \]
\end{proof}
\olpseganchor{OLP-0286-B021}{end}

\olpseganchor{OLP-0286-B022}{start}
\begin{prop}
فرض کړئ \LR{$\Gamma$} د تړلي فارمولونه يو بنسټيز بازګشتي سټ دے۔ بيا
اړيکه \LR{$\Prf[\Gamma](x,y)$}، چې وايي «\LR{$x$} د
\LR{$\Gamma_0 \Sequent !A$} د اشتقاق~\LR{$\pi$} کوډ دے، چې د کوم
متناهي \LR{$\Gamma_0\subseteq\Gamma$} دپاره وي، او \LR{$y$} د~\LR{$!A$} ګوډل
شمېره ده»، بنسټيزه بازګشتي ده۔
\end{prop}
\olpseganchor{OLP-0286-B022}{end}

\olpseganchor{OLP-0286-B023}{start}
\begin{proof}
فرض کړئ «\LR{$y\in\Gamma$}» د بنسټيز بازګشتي
پريديکات~\LR{$R_\Gamma(y)$} په واسطه ورکول کېږي۔ بايد وښيو چې
\LR{$\Prf[\Gamma](x,y)$} بنسټيزه بازګشتي ده؛ دا اړيکه هله صدق کوي چې
\LR{$y$} د يوې جملې~\LR{$!A$} ګوډل شمېره وي او \LR{$x$} د داسې
\LR{$\Log{LK}$}-اشتقاق کوډ وي چې پاې-سېکوېنټ ئې
\LR{$\Gamma_0\Sequent !A$} وي۔
\olpseganchor{OLP-0286-B023}{end}

\olpseganchor{OLP-0286-B024}{start}
د تېرې قضيې له مخې خاصيت \LR{$\fn{Deriv}(x)$}، چې هله صدق کوي چې \LR{$x$} په
\LR{$\Log{LK}$} کښې د يوه سم اشتقاق~\LR{$\pi$} کوډ وي، بنسټيز بازګشتي
دے۔ که \LR{$x$} داسې کوډ وي، \LR{$\fn{EndSequent}(x)$} د~\LR{$\pi$} د
پاې-سېکوېنټ کوډ دے، نو \LR{$(\fn{EndSequent}(x))_0$} د پاې-سېکوېنټ د
کيڼ لوري کوډ او \LR{$(\fn{EndSequent}(x))_1$} د ښي لوري کوډ دے۔ نو «د
~\LR{$\pi$} د پاې-سېکوېنټ ښے لور~\LR{$!A$} دے» داسې بيانولے شو:
\LR{$\len{(\fn{EndSequent}(x))_1}=1\land
((\fn{EndSequent}(x))_1)_0=y$}۔
د~\LR{$\pi$} د پاې-سېکوېنټ کيڼ لور البته پخپله متناهي دے؛ يوازې دا بيانول
پکار دي چې هره جمله ئې په~\LR{$\Gamma$} کښې ده۔ نو
\LR{$\Prf[\Gamma](x,y)$} داسې تعريفولے شو:
\begin{align*}
\Prf[\Gamma](x, y) \defiff {}&
\fn{Deriv}(x) \land {} \\
& \bforall{i <
  \len{(\fn{EndSequent}(x))_0}}{R_\Gamma(((\fn{EndSequent}(x))_0)_i)} \land {}\\
& \len{(\fn{EndSequent}(x))_1} = 1 \land ((\fn{EndSequent}(x))_1)_0 = y.
\end{align*}
\end{proof}
\olpseganchor{OLP-0286-B024}{end}

\olpseganchor{OLP-0287-B004}{start}
\olfileid{inc}{art}{pnd}
\olsection{په طبيعي استنتاج کښې اشتقاقونه}
\olpseganchor{OLP-0287-B004}{end}

\olpseganchor{OLP-0287-B005}{start}
\begin{explain}
د اشتقاقونه د حسابي کولو دپاره بايد اشتقاقونه د عددونو
په توګه تمثيل کړو۔ ځکه چې اشتقاقونه د فارمولونه داسې ونې
دي چې هر استنتاج يوه يا دوه نښې لري، نو بازګشتي تمثيل تر ټولو څرګنده
لار ده: يو اشتقاق د يوې څوګونې په توګه تمثيلوو چې اجزا ئې د
وروستي استنتاج مقدماتو ته د رسېدونکو سمدستي فرعي-اشتقاقونه
شمېر، د دغو فرعي-اشتقاقونه تمثيلونه، پاې-فارمول، د وروستي
استنتاج د ختمولو نښه، او يو داسې عدد دي چې د وروستي استنتاج ډول ښيي۔
\end{explain}
\olpseganchor{OLP-0287-B005}{end}

\olpseganchor{OLP-0287-B006}{start}
\begin{defn}
  که \LR{$\delta$} په طبيعي استنتاج کښې اشتقاق وي، نو
  \LR{$\Gn{\delta}$} په استقرايي ډول داسې تعريفېږي:
  \begin{enumerate}
  \item
    که \LR{$\delta$} يوازې له فرض~\LR{$!A$} څخه جوړ وي، نو \LR{$\Gn{\delta}$}
    دا دے: \LR{$\tuple{0,\Gn{!A},n}$}۔ عدد~\LR{$n$} که دا يو
    نۀ ايستل شوې فرض وي، \LR{$0$} دے؛ او که نۀ، عددي نښه ده۔
  \item
    که \LR{$\delta$} په داسې استنتاج پاې ته رسېږي چې صفر، يو، دوه يا درې
    مقدمات لري، نو \LR{$\Gn{\delta}$} په ترتيب له دې څخه يو دے:
    \begin{align*}
      & \tuple{0, \Gn{!A}, n, k}, \\
      & \tuple{1, \Gn{\delta_1}, \Gn{!A}, n, k}, \\
      & \tuple{2, \Gn{\delta_1}, \Gn{\delta_2}, \Gn{!A}, n, k}, \text{\RL{ يا}}\\
      & \tuple{3, \Gn{\delta_1}, \Gn{\delta_2}, \Gn{\delta_3}, \Gn{!A}, n,
        k},
    \end{align*}
    دلته \LR{$\delta_1$}، \LR{$\delta_2$} او \LR{$\delta_3$} هغه
    فرعي-اشتقاقونه دي چې په~\LR{$\delta$} کښې د وروستي استنتاج پر
    مقدمه يا مقدماتو پاې ته رسېږي؛ \LR{$!A$} په~\LR{$\delta$} کښې د وروستي
    استنتاج پايله ده؛
    \LR{$n$} د وروستي استنتاج د ختمولو نښه ده (که استنتاج هېڅ فرض نۀ
    ختموي، \LR{$0$} ده)؛ او \LR{$k$} د لاندې جدول له مخې د وروستي استنتاج د
    کارول شوې قاعدې له مخې ټاکل کېږي۔
\olpseganchor{OLP-0287-B006}{end}

\olpseganchor{OLP-0287-B007}{start}
    \begin{tabular}{lccccccc}
      \text{\RL{قاعده:}} & \Intro{\land} & \Elim{\land} & \Intro{\lor} & \Elim{\lor} \\
      \LR{$k$}: & 1 & 2 & 3 & 4  \\[1.5ex]
      \text{\RL{قاعده:}} &  \Intro{\lif} & \Elim{\lif} & \Intro{\lnot} & \Elim{\lnot} \\
      \LR{$k$}: & 5 & 6 & 7 & 8 \\[1.5ex]
      \text{\RL{قاعده:}}  &  \FalseInt & \FalseCl & \Intro{\lforall} & \Elim{\lforall} \\
      \LR{$k$}: & 9 & 10 & 11 & 12 \\[1.5ex]
      \text{\RL{قاعده:}} & \Intro{\lexists} & \Elim{\lexists} & \Intro{\eq} & \Elim{\eq} \\
      \LR{$k$}: & 13 & 14 & 15 & 16
    \end{tabular}
  \end{enumerate}
\end{defn}
\olpseganchor{OLP-0287-B007}{end}

\olpseganchor{OLP-0287-B008}{start}
\begin{ex}
  دا ډېر ساده اشتقاق وګورئ:
  \begin{prooftree}
    \AxiomC{\LR{$\Discharge{!A \land !B}{1}$}}
    \RightLabel{\Elim{\land}}
    \UnaryInfC{\LR{$!A$}}
    \DischargeRule{\Intro{\lif}}{1}
    \UnaryInfC{\LR{$(!A \land !B) \lif !A$}}
  \end{prooftree}
  د فرض ګوډل شمېره به
  \LR{$d_0=\tuple{0,\Gn{!A \land !B},1}$} وي۔ د هغه اشتقاق ګوډل
  شمېره چې د~\LR{$\Elim{\land}$} پر پايله ختمېږي،
  \LR{$d_1=\tuple{1,d_0,\Gn{!A},0,2}$} ده (\LR{$1$} ځکه چې
  \LR{$\Elim{\land}$} يوه مقدمه لري؛ بيا د پايلې~\LR{$!A$} ګوډل شمېره؛ \LR{$0$}
  ځکه چې هېڅ فرض نۀ ختمېږي؛ او \LR{$2$} هغه عدد دے چې \LR{$\Elim{\land}$}
  کوډوي)۔ بيا د ټول اشتقاق ګوډل شمېره
  \LR{$\tuple{1,d_1,\Gn{((!A \land !B) \lif !A)},1,5}$} ده؛ يعنې
  \[
  \tuple{1, \tuple{1, \tuple{0, \Gn{(!A \land !B)}, 1}, \Gn{!A}, 0, 2},
    \Gn{((!A \land !B) \lif !A)}, 1, 5}.
  \]
\end{ex}
\olpseganchor{OLP-0287-B008}{end}

\olpseganchor{OLP-0287-B009}{start}
\begin{explain}
چې د اشتقاقونه پر يو تمثيل هوکړې ته ورسېدو، دا هم بايد وښيو
چې د داسې اشتقاقونه ګوډل شمېرې په بنسټيز بازګشتي ډول اداره
کولے شو، او د هغو اړين خاصيتونه او اړيکې هم بيانولے شو۔ ځينې عمليات
ساده دي: د بېلګې په توګه، د اشتقاق له ګوډل شمېرې~\LR{$d$} څخه
\LR{$\fn{EndFmla}(d)=(d)_{(d)_0+1}$} د هغه د پاې-فارمول ګوډل
شمېره راکوي، \LR{$\fn{DischargeLabel}(d)=(d)_{(d)_0+2}$} د ختمولو نښه
راکوي، او \LR{$\fn{LastRule}(d)=(d)_{(d)_0+3}$} هغه عدد راکوي چې د وروستي
استنتاج ډول ښيي۔ ځينې نور عمليات ډېر ستونزمن دي۔ لږ تر لږه د هغو د
ترسره کولو لار به راونغاړو۔ موخه دا ده چې وښيو اړيکه «\LR{$\delta$}
له~\LR{$\Gamma$} څخه د~\LR{$!A$} يو اشتقاق دے» د~\LR{$\delta$} او~\LR{$!A$}
د ګوډل شمېرو بنسټيزه بازګشتي اړيکه ده۔
\end{explain}
\olpseganchor{OLP-0287-B009}{end}

\olpseganchor{OLP-0287-B010}{start}
\begin{prop}
  لاندې اړيکې بنسټيزې بازګشتي دي:
  \begin{enumerate}
  \item \LR{$!A$} په~\LR{$\delta$} کښې د نښې~\LR{$n$} لرونکي فرض په توګه راځي۔
  \item په~\LR{$\delta$} کښې د نښې~\LR{$n$} لرونکي ټول فرضونه د~\LR{$!A$} په بڼه
    دي (يعنې په~\LR{$\delta$} کښې د نښې~\LR{$n$} په کارولو فرض~\LR{$!A$}
    ايستل کولے شو)۔
  \end{enumerate}
\end{prop}
\olpseganchor{OLP-0287-B010}{end}

\olpseganchor{OLP-0287-B011}{start}
\begin{proof}
  بايد وښيو چې د فارمولونه د ګوډل شمېرو او د اشتقاقونه د
  ګوډل شمېرو ترمنځ اړوندې اړيکې بنسټيزې بازګشتي دي۔
  \begin{enumerate}
  \item غواړو وښيو چې \LR{$\fn{Assum}(x,d,n)$} بنسټيزه بازګشتي ده؛ دا
    هله صدق کوي چې \LR{$x$} د ګوډل شمېرې~\LR{$d$} لرونکي اشتقاق د
    نښې~\LR{$n$} لرونکي فرض ګوډل شمېره وي۔ دا هله سم دي چې د ګوډل شمېرې
    \LR{$\tuple{0,x,n}$} لرونکے اشتقاق د~\LR{$d$} يو
    فرعي-اشتقاق وي۔ پام وکړئ، زموږ د اشتقاقونه د کوډولو
    لار د هغو ونو د کوډولو ځانګړے حالت دے چې په
    \olref[cmp][rec][tre]{sec} کښې معرفي شوې؛ نو بنسټيزه بازګشتي
    تابع \LR{$\fn{SubtreeSeq}(d)$} د~\LR{$d$} د ټولو فرعي-اشتقاقونه د
    ګوډل شمېرو لړۍ (چې اوږدوالے ئې تر ټولو زيات~\LR{$d$} دے) راکوي۔ نو
    داسې تعريفولے شو:
    \[
    \fn{Assum}(x, d, n) \defiff \bexists{i<d}{(\fn{SubtreeSeq}(d))_i =
      \tuple{0, x, n}}.
    \]
  \item غواړو وښيو چې \LR{$\fn{Discharge}(x,d,n)$} بنسټيزه بازګشتي ده؛
    دا هله صدق کوي چې د ګوډل شمېرې~\LR{$d$} لرونکي اشتقاق ټول
    نښه~\LR{$n$} لرونکي فرضونه د ګوډل شمېرې~\LR{$x$} لرونکي فارمول وي۔
    خو دا اړيکه هله صدق کوي چې
    \LR{$\bforall{y<d}{(\fn{Assum}(y,d,n)\lif y=x)}$}۔
  \end{enumerate}
\end{proof}
\olpseganchor{OLP-0287-B011}{end}

\olpseganchor{OLP-0287-B012}{start}
\begin{prop}
  \ollabel{prop:followsby}
  خاصيت \LR{$\fn{Correct}(d)$}، چې هله صدق کوي چې د ګوډل شمېرې~\LR{$d$}
  لرونکي اشتقاق~\LR{$\delta$} وروستے استنتاج سم وي، بنسټيز
  بازګشتي دے۔
\end{prop}
\olpseganchor{OLP-0287-B012}{end}

\olpseganchor{OLP-0287-B013}{start}
\begin{proof}
  دلته د استنتاج د هرې قاعدې~\LR{$R$} دپاره بايد وښيو چې اړيکه
  \LR{$\fn{FollowsBy}_R(d)$} بنسټيزه بازګشتي ده؛
  \LR{$\fn{FollowsBy}_R(d)$} هله صدق کوي چې \LR{$d$} د
  اشتقاق~\LR{$\delta$} ګوډل شمېره وي او د~\LR{$\delta$}
  پاې-فارمول د~\LR{$R$} په سم پلي کولو د~\LR{$\delta$} د سمدستي
  فرعي-اشتقاقونه له پاې-فارمولونو څخه راشي۔
\olpseganchor{OLP-0287-B013}{end}

\olpseganchor{OLP-0287-B014}{start}
  يو ساده حالت د \Intro{\land} قاعده ده۔ که \LR{$\delta$} په يوه سم
  \LR{$\Intro{\land}$} استنتاج ختم شي، داسې ښکاري:
  \begin{prooftree}
    \AxiomC{}
    \RightLabel{\LR{$\delta_1$}}
    \DeduceC{\LR{$!A$}}
\olpseganchor{OLP-0287-B014}{end}

\olpseganchor{OLP-0287-B015}{start}
    \AxiomC{}
    \RightLabel{\LR{$\delta_2$}}
    \DeduceC{\LR{$!B$}}
\olpseganchor{OLP-0287-B015}{end}

\olpseganchor{OLP-0287-B016}{start}
    \RightLabel{\Intro\land}
    \BinaryInfC{\LR{$!A \land !B$}}
  \end{prooftree}
  بيا د~\LR{$\delta$} ګوډل شمېره~\LR{$d$} دا ده:
  \LR{$\tuple{2,d_1,d_2,\Gn{(!A \land !B)},0,k}$}، چې
  \LR{$\fn{EndFmla}(d_1)=\Gn{!A}$}، \LR{$\fn{EndFmla}(d_2)=\Gn{!B}$}،
  \LR{$n=0$} او \LR{$k=1$} دي۔ نو \LR{$\fn{FollowsBy}_{\Intro\land}(d)$} داسې
  تعريفولے شو:
  \begin{multline*}
    (d)_0 = 2 \land \fn{DischargeLabel}(d) = 0 \land \fn{LastRule}(d) = 1 \land {}\\
    \fn{EndFmla}(d) = {}\\
    \Gn{(} \concat \fn{EndFmla}((d)_1) \concat \Gn{\land}
    \concat \fn{EndFmla}((d)_2) \concat \Gn{)}.
  \end{multline*}
\olpseganchor{OLP-0287-B016}{end}

\olpseganchor{OLP-0287-B017}{start}
  بله ساده بېلګه د \LR{$\Intro\eq$} قاعده ده۔ دا هېڅ مقدمه نۀ لري، نو د
  فرضونو په شان \LR{$(d)_0=0$} دے۔ دا د ختمولو نښه هم نۀ لري؛ يعنې \LR{$n=0$}
  دے۔ خو \LR{$!A$} بايد د کوم تړلي ترم~\LR{$t$} دپاره د~\LR{$\eq[t][t]$} په بڼه
  وي۔ يو بنسټيز بازګشتي تعريف دا دے:
  \begin{multline*}
    (d)_0 = 0 \land \fn{DischargeLabel}(d) = 0 \land {}\\
    \bexists{t<d}{(\fn{ClTerm}(t) \land
      \fn{EndFmla}(d) = {}\\
      \Gn{{\eq}(} \concat t \concat \Gn{,} \concat t \concat \Gn{)})}.
  \end{multline*}
\olpseganchor{OLP-0287-B017}{end}

\olpseganchor{OLP-0287-B018}{start}
  د يوې لا پېچلې بېلګې دپاره \LR{$\fn{FollowsBy}_{\Intro{\lif}}(d)$} هله
  صدق کوي چې د~\LR{$\delta$} پاې-فارمول د~\LR{$(!A\lif !B)$} په بڼه وي،
  د~\LR{$\delta_1$} پاې-فارمول هم~\LR{$!B$} وي، او په~\LR{$\delta$} کښې هر
  نښه~\LR{$n$} لرونکے فرض د~\LR{$!A$} په بڼه وي۔ دا په بنسټيز بازګشتي ډول
  داسې بيانولے شو:
  \begin{multline*}
    (d)_0 = 1 \land {}\\
    \bexists{a<d}{(\fn{Discharge}(a, (d)_1, \fn{DischargeLabel}(d)) \land {}}\\
      \fn{EndFmla}(d) = (\Gn{(} \concat a \concat \Gn{\lif}
      \concat \fn{EndFmla}((d)_1) \concat \Gn{)}))
  \end{multline*}
  (دلته \LR{$a$} د~\LR{$!A$} ګوډل شمېره وګڼئ)۔
\olpseganchor{OLP-0287-B018}{end}

\olpseganchor{OLP-0287-B019}{start}
  د بلې بېلګې دپاره \Intro{\lexists} وګورئ۔ دلته په~\LR{$\delta$} کښې
  وروستے استنتاج هله سم دے چې داسې فارمول~\LR{$!A$}، تړلے ترم~\LR{$t$}
  او متغير~\LR{$x$} شته چې \LR{$\Subst{!A}{t}{x}$} د
  اشتقاق~\LR{$\delta_1$} پاې-فارمول وي او
  \LR{$\lexists[x][!A]$} د وروستي استنتاج پايله وي۔ نو
  \LR{$\fn{FollowsBy}_{\Intro{\lexists}}(d)$} هله صدق کوي چې
  \begin{multline*}
    (d)_0 = 1 \land \fn{DischargeLabel}(d) = 0 \land {} \\
    \bexists{a < d}{\bexists{x<d}{\bexists{t<d}{
          (\fn{ClTerm}(t) \land \fn{Var}(x)  \land {}}}}\\
    \fn{Subst}(a,t,x) = \fn{EndFmla}((d)_1) \land
    \fn{EndFmla}(d) = (\Gn{\lexists} \concat x \concat a)).
  \end{multline*}
\olpseganchor{OLP-0287-B019}{end}

\olpseganchor{OLP-0287-B020}{start}
  بيا \LR{$\fn{Correct}(d)$} داسې تعريفوو:
  \begin{multline*}
    \fn{Sent}(\fn{EndFmla}(d)) \land {}\\
    [(\fn{LastRule}(d) = 1 \land
    \fn{FollowsBy}_{\Intro\land}(d)) \lor \dots \lor {}\\
    (\fn{LastRule}(d) = 16 \land \fn{FollowsBy}_{\Elim\eq}(d)) \lor {}\\
    \bexists{n<d}{\bexists{x<d}{(d = \tuple{0, x, n})}}].
  \end{multline*}
  لومړۍ کرښه باوري کوي چې د~\LR{$d$} پاې-فارمول يوه جمله ده۔ وروستۍ
  کرښه هغه حالت رانغاړي چې \LR{$d$} يوازې يو فرض وي۔
\end{proof}
\olpseganchor{OLP-0287-B020}{end}

\olpseganchor{OLP-0287-B021}{start}
\begin{prob}
  لاندې خاصيتونه د \olref[inc][art][pnd]{prop:followsby} په څېر
  تعريف کړئ:
  \begin{enumerate}
  \item \LR{$\fn{FollowsBy}_{\Elim{\lif}}(d)$}؛
  \item \LR{$\fn{FollowsBy}_{\Elim{\eq}}(d)$}؛
  \item \LR{$\fn{FollowsBy}_{\Elim{\lor}}(d)$}؛
  \item \LR{$\fn{FollowsBy}_{\Intro{\lforall}}(d)$}۔
  \end{enumerate}
  د وروستي خاصيت دپاره دا هم بايد وښيئ چې په بنسټيز بازګشتي ډول
  ازمويلے شئ چې د ګوډل شمېرې~\LR{$d$} لرونکي اشتقاق وروستے
  استنتاج د ځانګړي متغير شرط پوره کوي؛ يعنې د
  \LR{$\Intro{\lforall}$} استنتاج ځانګړے متغير~\LR{$a$} نۀ د~\LR{$d$} په
  پاې-فارمول کښې راځي او نۀ د~\LR{$d$} په يوه پرانيستي فرض کښې۔ د
  دې دپاره د \olref[inc][art][pnd]{prop:openassum} بنسټيز بازګشتي
  پريديکات~\LR{$\fn{OpenAssum}$} کارولے شئ۔
\end{prob}
\olpseganchor{OLP-0287-B021}{end}

\olpseganchor{OLP-0287-B022}{start}
\begin{prop}
  \ollabel{prop:deriv}
  اړيکه \LR{$\fn{Deriv}(d)$}، چې هله صدق کوي چې \LR{$d$} د يوه سم
  اشتقاق~\LR{$\delta$} ګوډل شمېره وي، بنسټيزه بازګشتي ده۔
\end{prop}
\olpseganchor{OLP-0287-B022}{end}

\olpseganchor{OLP-0287-B023}{start}
\begin{proof}
  اشتقاق~\LR{$\delta$} هله سم دے چې هر استنتاج ئې د يوې قاعدې
  سم پلي کول وي؛ يعنې د هغه هر فرعي-اشتقاقونه په سم استنتاج پاې
  ته ورسېږي۔ نو \LR{$\fn{Deriv}(d)$} هله صدق کوي چې
  \[
  \bforall{i<\len{\fn{SubtreeSeq}(d)}}{\fn{Correct}((\fn{SubtreeSeq}(d))_i)}
  \]
\end{proof}
\olpseganchor{OLP-0287-B023}{end}

\olpseganchor{OLP-0287-B024}{start}
\begin{prop}
  \ollabel{prop:openassum} اړيکه \LR{$\fn{OpenAssum}(z,d)$}، چې هله صدق
  کوي چې \LR{$z$} د ګوډل شمېرې~\LR{$d$} لرونکي اشتقاق~\LR{$\delta$} د يوه
  نۀ ايستل شوې فرض~\LR{$!A$} ګوډل شمېره وي، بنسټيزه بازګشتي ده۔
\end{prop}
\olpseganchor{OLP-0287-B024}{end}

\olpseganchor{OLP-0287-B025}{start}
\begin{proof}
  د يوه فرض پېښه هله ايستل ده چې د نښې~\LR{$n$} سره د~\LR{$\delta$}
  په داسې فرعي-اشتقاق کښې راشي چې د ختمولو په نښه~\LR{$n$} لرونکې
  قاعده ختمېږي۔ نو \LR{$!A$} هله د~\LR{$\delta$} يو نۀ ايستل شوې فرض دے
  چې لږ تر لږه يوه پېښه ئې په~\LR{$\delta$} کښې ايستل نۀ وي۔
  بايد پام وکړو: \LR{$\delta$} کېدے شي د~\LR{$!A$} دواړه ايستل او
  نۀ ايستل شوې پېښې ولري۔
\olpseganchor{OLP-0287-B025}{end}

\olpseganchor{OLP-0287-B026}{start}
  داسې لړۍ \LR{$\delta_0$}، \dots، \LR{$\delta_k$} وګورئ چې
  \LR{$\delta_0=\delta$}، \LR{$\delta_k$} فرض \LR{$\Discharge{!A}{n}$} (د کوم~\LR{$n$}
  دپاره) وي، او \LR{$\delta_{i+1}$} د~\LR{$\delta_i$} سمدستي
  فرعي-اشتقاق وي۔ که داسې لړۍ موجوده وي چې په هغې کښې هېڅ
  \LR{$\delta_i$} د ختمولو نښه~\LR{$n$} لرونکي استنتاج باندې نۀ ختمېږي، نو
  \LR{$!A$} د~\LR{$\delta$} يو نۀ ايستل شوې فرض دے۔
\olpseganchor{OLP-0287-B026}{end}

\olpseganchor{OLP-0287-B027}{start}
  بنسټيزه بازګشتي تابع \LR{$\fn{SubtreeSeq}(d)$} د~\LR{$\delta$} د ټولو
  فرعي-اشتقاقونه د ګوډل شمېرو لړۍ راکوي۔ د ريښې او پاڼې تر منځ د فرعي-اشتقاقونه د لارې هر کوډ د دغې لړۍ غړے دے۔ دلته يوازې د هر غړي شتون ازمويو؛ ترتيب او تکرار د غړيتوب په دغه ازموينه کښې شرط نۀ دي۔
  \LR{$\fn{Subseq}(s,s')$} هله صدق کوي چې
  \LR{$\bforall{i<\len{s}}{\bexists{j<\len{s'}}{(s)_i=(s')_j}}$}۔ د
  سمدستي فرعي-اشتقاق کېدل هم داسې دي:
  \LR{$\fn{Subderiv}(d,d')$} هله صدق کوي چې
  \LR{$\bexists{j<(d')_0}{d=(d')_{j+1}}$}۔
  د لارې اوږدوالے د کوډ د پرله‌پسې کمېدو له کبله تر ټولو زيات د ريښې کوډ جمع يو دے، او د هرې غوټې کوډ تر ريښې زيات نۀ دے۔ د لړۍ د کوډ حد له دغو دوو حدونو راځي۔ صفر د هغو فرضونو نښه ده چې له پيله پرانيستي دي؛ يوازې مثبته نښه د ختمولو په شرط کښې ګورو۔ نو \LR{$\fn{OpenAssum}(z,d)$} داسې تعريفولے شو:
  \begin{multline*}
    \bexists{s<\fn{sequenceBound}(d,d+1)+1}{(\len{s}>0 \land (s)_0=d \land {}} \\
    \bexists{n<d}{((s)_{\len{s}\tsub 1}=\tuple{0,z,n} \land {}} \\
    \bforall{i<(\len{s}\tsub 1)}{(\fn{Subderiv}((s)_{i+1},(s)_i) \land {}} \\
    (n=0 \lor \fn{DischargeLabel}((s)_i)\neq n)))).
  \end{multline*}
\end{proof}
\olpseganchor{OLP-0287-B027}{end}

\olpseganchor{OLP-0287-B028}{start}
\begin{prop}
  \ollabel{prop:prf-prim-rec}
  فرض کړئ \LR{$\Gamma$} د تړلي فارمولونه يو بنسټيز بازګشتي سټ دے۔ بيا
  اړيکه \LR{$\Prf[\Gamma](x,y)$}، چې وايي «\LR{$x$} له~\LR{$\Gamma$} کښې
  نۀ ايستل شوې فرضونو څخه د~\LR{$!A$} د اشتقاق~\LR{$\delta$} کوډ
  دے او \LR{$y$} د~\LR{$!A$} ګوډل شمېره ده»، بنسټيزه بازګشتي ده۔
\end{prop}
\olpseganchor{OLP-0287-B028}{end}

\olpseganchor{OLP-0287-B029}{start}
\begin{proof}
  فرض کړئ «\LR{$y\in\Gamma$}» د بنسټيز بازګشتي
  پريديکات~\LR{$R_\Gamma(y)$} په واسطه ورکول کېږي۔ بايد وښيو چې
  \LR{$\Prf[\Gamma](x,y)$} بنسټيزه بازګشتي ده؛ دا اړيکه هله صدق کوي چې
  \LR{$y$} د يوې جملې~\LR{$!A$} ګوډل شمېره وي، او \LR{$x$} د طبيعي استنتاج د داسې
  اشتقاق کوډ وي چې پاې-فارمول ئې~\LR{$!A$} او ټول
  نۀ ايستل شوي فرضونه ئې په~\LR{$\Gamma$} کښې وي۔
\olpseganchor{OLP-0287-B029}{end}

\olpseganchor{OLP-0287-B030}{start}
  د \olref{prop:deriv} له مخې خاصيت \LR{$\fn{Deriv}(x)$}، چې هله صدق کوي
  چې \LR{$x$} په طبيعي استنتاج کښې د يوه سم اشتقاق~\LR{$\delta$} ګوډل
  شمېره وي، بنسټيز بازګشتي دے۔ نو \LR{$\Prf[\Gamma](x,y)$} داسې تعريفولے
  شو:
  \begin{align*}
    \Prf[\Gamma](x, y) \defiff {}
    & \fn{Sent}(y) \land \fn{Deriv}(x) \land \fn{EndFmla}(x) = y \land {} \\
    & \bforall{z < x}{(\fn{OpenAssum}(z, x) \lif R_\Gamma(z))}.
  \end{align*}
\end{proof}
\olpseganchor{OLP-0287-B030}{end}

\olpseganchor{OLP-0288-B004}{start}
\olfileid{inc}{art}{pax}
\olsection{بديهي اشتقاقونه}
\olpseganchor{OLP-0288-B004}{end}

\olpseganchor{OLP-0288-B005}{start}
\begin{explain}
د بديهي اشتقاقونه د حسابي کولو دپاره بايد اشتقاقونه د
عددونو په توګه تمثيل کړو۔ ځکه چې اشتقاقونه يوازې د
فارمولونه لړۍ دي، نو څرګنده لار دا ده چې هر اشتقاق د هغه
د فارمولونه د کوډونو د لړۍ په کوډ سره کوډ کړو۔
\end{explain}
\olpseganchor{OLP-0288-B005}{end}

\olpseganchor{OLP-0288-B006}{start}
\begin{defn}
که \LR{$\delta$} داسې بديهي اشتقاق وي چې له فارمولونه
\LR{$!A_1$}، \dots،~\LR{$!A_n$} څخه جوړ وي، نو \LR{$\Gn{\delta}$} دا دے:
\[
\tuple{\Gn{!A_1}, \dots, \Gn{!A_n}}.
\]
\end{defn}
\olpseganchor{OLP-0288-B006}{end}

\olpseganchor{OLP-0288-B007}{start}
\begin{ex}
  دا ډېر ساده اشتقاق وګورئ:
  \begin{derivation}
  1. & \LR{$!B \lif (!B \lor !A)$} \\
  2. & \LR{$(!B \lif (!B \lor !A)) \lif (!A  \lif (!B \lif (!B \lor !A)))$}\\
  3. & \LR{$!A  \lif (!B \lif (!B \lor !A))$}
  \end{derivation}
  د دې اشتقاق ګوډل شمېره به دا وي:
  \begin{align*}
  \openTuple\,
    & \Gn{!B \lif (!B \lor !A)}, \\
    &\Gn{(!B \lif (!B \lor !A)) \lif (!A  \lif (!B \lif (!B \lor !A)))},\\
    & \Gn{!A  \lif (!B \lif (!B \lor !A))} \,\closeTuple.
  \end{align*}
\end{ex}
\olpseganchor{OLP-0288-B007}{end}

\olpseganchor{OLP-0288-B008}{start}
\begin{explain}
چې د اشتقاقونه پر يو تمثيل هوکړې ته ورسېدو، دا هم بايد وښيو
چې داسې اشتقاقونه په بنسټيز بازګشتي ډول اداره کولے شو، او د
هغو اړين خاصيتونه او اړيکې هم په همدې ډول بيانولے شو۔ ځينې عمليات
ساده دي: د بېلګې په توګه، د اشتقاق له ګوډل شمېرې~\LR{$d$} څخه
\LR{$(d)_{\len{d}-1}$} د هغه د پاې-فارمول ګوډل شمېره راکوي۔ ځينې
نور عمليات ډېر ستونزمن دي۔ لږ تر لږه د هغو د ترسره کولو لار به
راونغاړو۔ موخه دا ده چې وښيو اړيکه «\LR{$\delta$} له~\LR{$\Gamma$} څخه د~\LR{$!A$}
يو اشتقاق دے» د~\LR{$\delta$} او~\LR{$!A$} په ګوډل شمېرو کښې
بنسټيزه بازګشتي ده۔
\end{explain}
\olpseganchor{OLP-0288-B008}{end}

\olpseganchor{OLP-0288-B009}{start}
\begin{prop}
\ollabel{prop:followsby}
لاندې اړيکې بنسټيزې بازګشتي دي:
\begin{enumerate}
\item \LR{$!A$} يو بديهي اصل دے۔
\item په~\LR{$\delta$} کښې \LR{$i$}-مه کرښه په مودوس پوننس توجيه شوې ده۔
\item په~\LR{$\delta$} کښې \LR{$i$}-مه کرښه په \QR{} توجيه شوې ده۔
\item \LR{$\delta$} يو سم اشتقاق دے۔
\end{enumerate}
\end{prop}
\olpseganchor{OLP-0288-B009}{end}

\olpseganchor{OLP-0288-B010}{start}
\begin{proof}
بايد وښيو چې د فارمولونه د ګوډل شمېرو او د اشتقاقونه د
ګوډل شمېرو ترمنځ اړوندې اړيکې بنسټيزې بازګشتي دي۔
\begin{enumerate}
\item موږ ته د بديهي شېماوو يو ټاکلے لست راکړل شوے، او \LR{$!A$} هله يو
  بديهي اصل دے چې د دې شېماوو څخه د يوې په ورکړې بڼه وي۔ ځکه چې د
  شېماوو لست متناهي دے، بس ده چې د هرې بديهي شېما دپاره په بنسټيز
  بازګشتي ډول وازميوو چې \LR{$!A$} د هغې په بڼه دے که نۀ۔ د بېلګې په
  توګه دا بديهي شېما وګورئ:
  \[
  !B \lif (!C \lif !B).
  \]
  \LR{$!A$} هله د دې بديهي شېما يوه بېلګه ده چې داسې فارمولونه \LR{$!B$}
  او \LR{$!C$} شته چې \LR{$!A$} د «\LR{$($}»، بيا \LR{$!B$}، بيا «\LR{$\lif$}»، بيا «\LR{$($}»،
  بيا \LR{$!C$}، بيا «\LR{$\lif$}»، بيا \LR{$!B$}، او بيا «\LR{$))$}» په نښلولو ترلاسه
  شي۔ د~\LR{$!A$} د ګوډل شمېرې~\LR{$n$} اړوند خاصيت ازمويلے شو، ځکه د لړيو
  نښلول بنسټيز بازګشتي دي، او د~\LR{$!B$} او~\LR{$!C$} ګوډل شمېرې بايد د~\LR{$!A$}
  له ګوډل شمېرې څخه وړې وي؛ ځکه کله چې اړيکه صدق کوي، \LR{$!B$} او \LR{$!C$}
  دواړه د~\LR{$!A$} فرعي-فارمولونه دي۔ نو داسې تعريفولے شو:
  \begin{multline*}
  \fn{IsAx}_{!B \lif (!C \lif !B)}(n) \defiff \bexists{b< n}{
    \bexists{c < n}{(\fn{Sent}(b) \land \fn{Sent}(c) \land {}}}\\ n =
  \Gn{(} \concat b \concat \Gn{\lif} \concat \Gn{(} \concat c \concat
  \Gn{\lif} \concat b \concat \Gn{))}).
  \end{multline*}
  که د هرې بديهي شېما دپاره داسې تعريف ولرو، د هغو د فصل خاصيت
  \LR{$\fn{IsAx}(n)$} تعريفوي: «\LR{$n$} د يوه بديهي اصل ګوډل شمېره ده»۔
\item په~\LR{$\delta$} کښې \LR{$i$}-مه کرښه هله په مودوس پوننس توجيه شوې ده
  چې داسې کرښې \LR{$j$} او \LR{$k<i$} شته چې د~\LR{$j$} پر کرښه تړلے فارمول کوم
  فارمول~\LR{$!A$}، د~\LR{$k$} پر کرښه جمله \LR{$!A\lif !B$}، او د~\LR{$i$} پر کرښه
  جمله~\LR{$!B$} وي۔
  \begin{multline*}
    \fn{MP}(d, i) \defiff \bexists{j < i}{\bexists{k < i}{}}\\
    (d)_k = \Gn{(} \concat (d)_j
      \concat \Gn{\lif} \concat (d)_i \concat \Gn{)}
  \end{multline*}
  ځکه محدود کمیت ايښودنه، نښلول او~\LR{$=$} بنسټيز بازګشتي دي، نو دا
  يوه بنسټيزه بازګشتي اړيکه تعريفوي۔
\item په~\LR{$\delta$} کښې يوه کرښه هله په \QR{} توجيه شوې ده چې د
  \LR{$!B\lif\lforall[x][!A(x)]$} په بڼه وي، يوه مخکښې کرښه د کوم
  ثابت سمبول~\LR{$c$} دپاره \LR{$!B\lif !A(c)$} وي، او \LR{$c$} په~\LR{$!B$} کښې نۀ
  راځي۔ دا هله کېږي چې:
  \begin{enumerate}
  \item داسې تړلے فارمول~\LR{$!B$} وي؛ او
  \item داسې فارمول~\LR{$!A(x)$} وي چې يوازې يو ازاد متغير~\LR{$x$}
    ولري؛ داسې چې
  \item کرښه~\LR{$i$}، \LR{$!B\lif\lforall[x][!A(x)]$} ولري؛
  \item کومه کرښه \LR{$j<i$} د کوم ثابت~\LR{$c$} دپاره
    \LR{$!B\lif\Subst{!A}{c}{x}$} ولري؛ او
  \item \LR{$c$} په~\LR{$!B$} کښې نۀ راځي۔
  \end{enumerate}
  دا ټول په بنسټيز بازګشتي ډول ازمويلے شو، ځکه د~\LR{$!B$}، \LR{$!A(x)$} او
  \LR{$x$} ګوډل شمېرې د~\LR{$i$} پر کرښه د فارمول له ګوډل شمېرې څخه وړې دي،
  او د~\LR{$c$} ګوډل شمېره د~\LR{$j$} پر کرښه د فارمول له ګوډل شمېرې څخه
  وړه ده:
  \begin{multline*}
    \fn{QR}_1(d, i) \defiff \bexists{j<i}{}\bexists{b < (d)_i}{\bexists{x <
        (d)_i}{\bexists{a < (d)_i}{\bexists{c < (d)_j}{(}}}}
    \\ \fn{Var}(x) \land \fn{Const}(c) \land {} \\
    (d)_i = \Gn{(} \concat
    b \concat \Gn{\lif} \concat \Gn{\lforall} \concat x \concat a
    \concat \Gn{)} \land {}\\
    (d)_j = \Gn{(} \concat b \concat
    \Gn{\lif} \concat \fn{Subst}(a,c,x) \concat \Gn{)} \land {}\\ \fn{Sent}(b)
    \land \fn{Sent}(\fn{Subst}(a,c,x)) \land {} \bforall{k <
      \len{b}}{(b)_k \neq (c)_0})
  \end{multline*}
  دلته فرض کوو چې \LR{$c$} او \LR{$x$} په ترتيب د ثابت او متغير ګوډل شمېرې
  دي چې د ترمونو په توګه ورته کتل شوي (يعنې د هغو سمبول-کوډونه نۀ
  دي)۔ دا چې \LR{$x$} د~\LR{$!A(x)$} يوازينے ازاد متغير دے، په دې ازميوو چې
  \LR{$\Subst{!A(x)}{c}{x}$} يوه تړلے فارمول ده؛ او دا چې \LR{$c$} په~\LR{$!B$}
  کښې نۀ راځي، په دې شرط باوروو چې د~\LR{$!B$} هر سمبول له~\LR{$c$} څخه بېل
  وي۔
\olpseganchor{OLP-0288-B010}{end}

\olpseganchor{OLP-0288-B011}{start}
  د \QR{} بله بڼه د تمرین دپاره پرېږدو۔
\item \LR{$d$} هله د يوه سم اشتقاق ګوډل شمېره ده چې هره کرښه ئې
  بديهي اصل وي، يا په مودوس پوننس يا~\QR{} توجيه شوې وي۔ اشتقاق لږ تر لږه يوه کرښه لري۔ نو:
  \[
  \fn{Deriv}(d) \defiff \len{d}>0 \land \bforall{i < \len{d}}{(\fn{IsAx}((d)_i) \lor
    \fn{MP}(d,i) \lor \fn{QR}(d, i))}
  \]
\end{enumerate}
\end{proof}
\olpseganchor{OLP-0288-B011}{end}

\olpseganchor{OLP-0288-B012}{start}
\begin{prob}
لاندې اړيکې د \olref[inc][art][pax]{prop:followsby} په څېر تعريف
کړئ:
\begin{enumerate}
\item \LR{$\fn{IsAx}_{!A \lif (!B \lif (!A \land !B))}(n)$}؛
\item \LR{$\fn{IsAx}_{\lforall[x][!A(x)] \lif !A(t)}(n)$}؛
\item \LR{$\fn{QR}_{2}(d, i)$} (د \QR{} د بلې بڼې دپاره)۔
\end{enumerate}
\end{prob}
\olpseganchor{OLP-0288-B012}{end}

\olpseganchor{OLP-0288-B013}{start}
\begin{prop}
فرض کړئ \LR{$\Gamma$} د تړلي فارمولونه يو بنسټيز بازګشتي سټ دے۔ بيا
اړيکه \LR{$\Prf[\Gamma](x,y)$}، چې وايي «\LR{$x$} له~\LR{$\Gamma$} څخه د~\LR{$!A$} د
اشتقاق~\LR{$\delta$} کوډ دے او \LR{$y$} د~\LR{$!A$} ګوډل شمېره ده»،
بنسټيزه بازګشتي ده۔
\end{prop}
\olpseganchor{OLP-0288-B013}{end}

\olpseganchor{OLP-0288-B014}{start}
\begin{proof}
فرض کړئ «\LR{$y\in\Gamma$}» د بنسټيز بازګشتي
پريديکات~\LR{$R_\Gamma(y)$} په واسطه ورکول کېږي۔ بايد وښيو چې اړيکه
\LR{$\Prf[\Gamma](x,y)$} بنسټيزه بازګشتي ده؛
\LR{$\Prf[\Gamma](x,y)$} هله صدق کوي چې
\LR{$y$} د تړلے فارمول~\LR{$!A$} ګوډل شمېره وي او \LR{$x$} له~\LR{$\Gamma$} څخه
د~\LR{$!A$} د اشتقاق کوډ وي۔
\olpseganchor{OLP-0288-B014}{end}

\olpseganchor{OLP-0288-B015}{start}
د تېرې قضيې له مخې خاصيت \LR{$\fn{Deriv}(x)$}، چې هله صدق کوي چې \LR{$x$} د
يوه سم اشتقاق~\LR{$\delta$} کوډ وي، بنسټيز بازګشتي دے۔ خو هغه
تعريف سټ~\LR{$\Gamma$} د اشتقاق د کرښو د توجيه د يوې اضافي لارې
په توګه په پام کښې نۀ و نيولے۔ زموږ بنسټيز بازګشتي ازموينې چې يوه
کرښه په \QR{} توجيه شوې که نۀ، دا شرط هم پرېښے و چې ثابت~\LR{$c$} ته په
\LR{$\Gamma$} کښې د راتلو اجازه نشته۔ تعريف داسې سمولے شو چې \LR{$\Gamma$}
نېغ په نېغه په پام کښې ونيسي، خو د \LR{$\fn{Deriv}$} او د استنتاج د قضيې
کارول اسانه دي۔ \LR{$\Gamma\Proves !A$} هله صدق کوي چې د تړلي فارمولونه
داسې کوم متناهي لست \LR{$!B_1$}، \dots، \LR{$!B_n\in\Gamma$} وي چې
\LR{$\{!B_1,\dots,!B_n\}\Proves !A$}۔ او د استنتاج د قضيې له مخې دا هله
سم دي چې
\LR{$\Proves(!B_1\lif(!B_2\lif\cdots(!B_n\lif !A)\cdots))$} وي۔ دا چې
د ګوډل شمېرې~\LR{$z$} لرونکې تړلے فارمول په دې بڼه ده که نۀ، په
بنسټيز بازګشتي ډول ازمويلے شو۔ نو \LR{$x$} له~\LR{$\Gamma$} څخه د ګوډل
شمېرې~\LR{$y$} لرونکې تړلے فارمول د اشتقاق د ګوډل شمېرې په توګه
نۀ اخلو، بلکې \LR{$x$} له~\LR{$\emptyset$} څخه د پورته بڼې د يوه ځالي شرطي د
اشتقاق ګوډل شمېره اخلو۔
\olpseganchor{OLP-0288-B015}{end}

\olpseganchor{OLP-0288-B016}{start}
لومړے، که د تړلي فارمولونه يوه لړۍ ولرو، د ټولو دغو جملو په مقدماتو
او د يوې ورکړې تړلے فارمول په تالي سره شرطي په بنسټيز بازګشتي ډول
جوړولے شو:
\begin{align*}
  \fn{hCond}(s, y, 0) & = y \\
  \fn{hCond}(s, y, n+1) & = \Gn{(} \concat (s)_{n} \concat \Gn{\lif}
  \concat \fn{hCond}(s, y, n) \concat \Gn{)}\\
  \fn{Cond}(s, y) & = \fn{hCond}(s, y, \len{s})\\
  \intertext{\RL{نو \LR{$\Prf[\Gamma](x,y)$} داسې تعريفولے شو:}}
  \Prf[\Gamma](x, y) & \defiff \fn{Sent}(y) \land \bexists{s < \fn{sequenceBound}(x,x)+1}{(} \\
  &\qquad (x)_{\len{x}-1} = \fn{Cond}(s,y) \land {} \\
  &\qquad\bforall{i<\len{s}}{(s)_i \in \Gamma} \land {} \\
  &\qquad\fn{Deriv}(x)).
\end{align*}
د~\LR{$s$} حد له دې څخه راځي چې هر \LR{$(s)_i$} د اشتقاق د وروستۍ
کرښې د يوه فرعي-فارمول ګوډل شمېره ده؛ يعنې له
\LR{$(x)_{\len{x}-1}$} څخه وړه ده۔ د مقدماتو~\LR{$!B \in \Gamma$} شمېر،
يعنې د~\LR{$s$} اوږدوالے، د~\LR{$x$} د وروستۍ کرښې له اوږدوالي څخه کم دے۔
\end{proof}
\olpseganchor{OLP-0288-B016}{end}

\olpseganchor{OLP-0289-B004}{start}
\olchapter{inc}{req}{په~\texorpdfstring{\LR{$\Th{Q}$}}{Q} کښې تمثيلېدنه}
\olpseganchor{OLP-0289-B004}{end}

\olpseganchor{OLP-0290-B004}{start}
\olfileid{inc}{req}{int}
\olsection{پېژندګلو}
\olpseganchor{OLP-0290-B004}{end}

\olpseganchor{OLP-0290-B005}{start}
د ناتکميلۍ قضيې پر هغو تيوريو پلي کېږي چې د محاسبه کېدونکو تابعو
په اړه بنسټيز حقايق پکې بيان او ثابتېدے شي۔ موږ به يوه ډېره لږه
داسې تيوري بيان کړو چې «\LR{$\Th{Q}$}» (يا کله کله د رافايل رابنسن په
نوم «د رابنسن~\LR{$Q$}») بلل کېږي۔ وبه وايو چې يوه تابعه په~\LR{$\Th{Q}$}
کښې
\emph{تمثيلېدونکې} وي څه معنا لري، او بيا به لاندې خبره ثابته کړو:
\begin{quote}
  يوه تابعه په~\LR{$\Th{Q}$} کښې هله او يوازې هله تمثيلېدونکې ده چې
  محاسبه کېدونکې وي۔
\end{quote}
دا له يوې خوا د محاسبې يو بل مدل راکوي۔ خو له بلې خوا به ئې د دې د
ښودلو دپاره هم وکاروو چې سټ~\LR{$\Setabs{!A}{\Th{Q} \Proves !A}$} د
پرېکړې وړ نۀ دے؛ د درېدنې مسئله به ورته راکمه کړو۔ تر هغه وخته چې
بحث بشپړوو، تر دې به ډېرې پياوړې پايلې مو ثابتې کړې وي۔
\olpseganchor{OLP-0290-B005}{end}

\olpseganchor{OLP-0290-B006}{start}
د~\LR{$\Th{Q}$} ژبه د حساب ژبه ده؛ \LR{$\Th{Q}$} له لاندې بديهي اصولو څخه
جوړه ده (چې د عينيت لرونکي لومړۍ درجې منطق له نورو بديهي
اصولو او قاعدو سره يوځای کارول کېږي):
\begin{align*}
& \lforall[x][\lforall[y][(\eq[x'][y'] \lif \eq[x][y])]] \tag{\LR{$!Q_1$}}\\
& \lforall[x][\eq/[\Obj 0][x']] \tag{\LR{$!Q_2$}}\\
& \lforall[x][(\eq[x][\Obj 0] \lor \lexists[y][\eq[x][y']])] \tag{\LR{$!Q_3$}}\\
& \lforall[x][\eq[(x + \Obj 0)][x]] \tag{\LR{$!Q_4$}}\\
& \lforall[x][\lforall[y][\eq[(x + y')][(x + y)']]] \tag{\LR{$!Q_5$}}\\
& \lforall[x][\eq[(x \times \Obj 0)][\Obj 0]] \tag{\LR{$!Q_6$}}\\
& \lforall[x][\lforall[y][\eq[(x \times y')][((x \times y) + x)]]] \tag{\LR{$!Q_7$}}\\
& \lforall[x][\lforall[y][(x < y \liff \lexists[z][\eq[(z' + x)][y]])]] \tag{\LR{$!Q_8$}}
\end{align*}
د هر طبيعي عدد~\LR{$n$} دپاره عددنښه~\LR{$\num{n}$} داسې تعريف کړئ چې ترم
\LR{$\Obj{0}^{\prime\prime\ldots\prime}$} وي او د پريم نښو ټول شمېر ئې~\LR{$n$} وي۔ نو
\LR{$\num{0}$} يوازې ثابت سمبول~\LR{$\Obj{0}$} دے، \LR{$\num{1}$}،
\LR{$\Obj{0}'$} دے، \LR{$\num{2}$}، \LR{$\Obj{0}''$} دے، او همداسې وړاندې۔
\olpseganchor{OLP-0290-B006}{end}

\olpseganchor{OLP-0290-B007}{start}
د حساب د يوې تيورۍ په توګه~\LR{$\Th{Q}$} \emph{بېخي} کمزورې ده؛ د
بېلګې په توګه ان د~\LR{$\lforall[x][\eq/[x][x']]$} يا
\LR{$\lforall[x][\lforall[y][(x + y) = (y + x)]]$} په شان ډېر ساده
حقايق هم پکې نۀ شئ ثابتولے۔ خو وبه وينو چې د~\LR{$\Th{Q}$} د پام وړتيا
تر ډېره \emph{له همدې} کمزورۍ څخه راځي۔ په حقيقت کښې دومره پياوړې
ده چې د ناتکميلۍ قضيه پرې پلې شي، خو له دې زياته نۀ ده۔ د~\LR{$\Th{Q}$}
د پام وړتيا بل لامل دا دے چې د بديهي اصولو سټ ئې \emph{متناهي} دے۔
\olpseganchor{OLP-0290-B007}{end}

\olpseganchor{OLP-0290-B008}{start}
له~\LR{$\Th{Q}$} څخه يوه پياوړې تيوري (چې د \emph{پيانو حساب}
\LR{$\Th{PA}$} بلل کېږي) په~\LR{$\Th{Q}$} کښې د استقرا شېما په زياتولو
ترلاسه کېږي:
\[
(!A(\Obj 0) \land \lforall[x][(!A(x) \lif !A(x'))]) \lif \lforall[x][!A(x)]
\]
دلته~\LR{$!A(x)$} هر فارمول دے۔ که~\LR{$!A(x)$} له~\LR{$x$} پرته نور ازاد
متغيرونه ولري، په سر کښې کلي کمیت ايښوونکي ورزياتوو چې ټول وتړي
(څو د استقرا د شېما اړونده بېلګه تړلے فارمول وي)۔ د بېلګې په
توګه، که~\LR{$!A(x, y)$} کښې متغير~\LR{$y$} هم ازاد وي، اړونده بېلګه
دا ده:
\[
\lforall[y][((!A(\Obj 0) \land \lforall[x][(!A(x) \lif !A(x'))]) \lif
  \lforall[x][!A(x)])]
\]
د استقرا د شېما د بېلګو په کارولو سره د~\LR{$\Th{PA}$} له بديهي اصولو
څخه د~\LR{$\Th{Q}$} په پرتله ډېر څه ثابتېدے شي۔ په حقيقت کښې د طبيعي
عددونو په اړه د داسې «طبيعي» بيانونو موندل ډېر کار غواړي چې د پيانو
په حساب کښې نۀ ثابتېږي!{}
\olpseganchor{OLP-0290-B008}{end}

\olpseganchor{OLP-0290-B009}{start}
\begin{defn}
\ollabel{defn:representable-fn}
  له طبيعي عددونو څخه طبيعي عددونو ته تابعه
  \LR{$f(x_0,\ldots,x_k)$} هله \emph{په~\LR{$\Th{Q}$} کښې تمثيلېدونکې}
  بلل کېږي چې داسې فارمول~\LR{$!A_f(x_0,\dots,x_k,y)$} وي چې هر
  کله~\LR{$f(n_0,\dots,n_k) = m$} وي، \LR{$\Th{Q}$} لاندې دواړه ثابت کړي:
\begin{enumerate}
\item\ollabel{defn:rep:a} \LR{$!A_f(\num{n_0}, \dots, \num{n_k}, \num{m})$}
\item\ollabel{defn:rep:b} \LR{$\lforall[y][(!A_f(\num{n_0}, \dots,
\num{n_k}, y) \lif \num{m} = y)]$}۔
\end{enumerate}
\end{defn}
\olpseganchor{OLP-0290-B009}{end}

\olpseganchor{OLP-0290-B010}{start}
د تعريف د بيان نورې برابرې لارې هم شته؛ د بېلګې په توګه په برابر
ډول دا غوښتلے شو چې~\LR{$\Th{Q}$}،
\LR{$\lforall[y][(!A_f(\num{n_0}, \dots,
\num{n_k}, y) \liff \eq[y][\num{m}])]$} ثابت کړي۔
\olpseganchor{OLP-0290-B010}{end}

\olpseganchor{OLP-0290-B011}{start}
\begin{thm}
\ollabel{thm:representable-iff-comp}
يوه تابعه په~\LR{$\Th{Q}$} کښې هله او يوازې هله تمثيلېدونکې ده چې
محاسبه کېدونکې وي۔
\end{thm}
\olpseganchor{OLP-0290-B011}{end}

\olpseganchor{OLP-0290-B012}{start}
د قضيې ثبوت دوه لوري لري۔ له کيڼ څخه ښي لور ته ثبوت، کله چې د نحو
حسابي کول جوړ شي، نسبتاً ساده دے۔ بل لور زيات کار غواړي۔ بنسټيزه
مفکوره دا ده: «عمومي بازګشتي» د «محاسبه کېدونکي» د کره کولو يوه لار
ټاکو، او ښيو چې هره عمومي بازګشتي تابعه په~\LR{$\Th{Q}$} کښې
تمثيلېدونکې ده۔ په ياد ولرئ، تابعه هله عمومي بازګشتي ده چې له
\LR{$\Zero$}، تالي تابع~\LR{$\Succ$}، او پروجکشن تابعو~\LR{$\Proj{n}{i}$} څخه د
ترکيب، بنسټيز بازګښت، او منظم اقل موندلو په کارولو سره تعريفېدے شي۔
نو د دې د ښودلو يوه لار چې هره عمومي بازګشتي تابعه په~\LR{$\Th{Q}$} کښې
تمثيلېدونکې ده، دا ده چې وښيو بنسټيزې تابعې تمثيلېدونکې دي، او هر
کله چې ځينې تابعې تمثيلېدونکې وي، د هغو له ترکيب، بنسټيز بازګښت، او
منظم اقل موندلو څخه تعريف شوې تابعې هم تمثيلېدونکې وي۔ په بله وينا،
ښودلے شو چې بنسټيزې تابعې تمثيلېدونکې دي، او تمثيلېدونکې تابعې د
ترکيب، بنسټيز بازګښت، او منظم اقل موندلو «لاندې تړلې» دي۔ دا خبره
باور راکوي چې هره عمومي بازګشتي تابعه تمثيلېدونکې ده۔
\olpseganchor{OLP-0290-B012}{end}

\olpseganchor{OLP-0290-B013}{start}
خو دا ګام چې وښيو تمثيلېدونکې تابعې د بنسټيز بازګښت لاندې تړلې دي،
ستونزمن راوځي۔ له دې ګام څخه د ډډې دپاره لومړے ښيو چې په حقيقت کښې
له بنسټيز بازګښت پرته کار کولے شو۔ يعنې ښيو چې هره عمومي بازګشتي
تابعه يوازې د بنسټيزو تابعو، ترکيب، او منظم اقل موندلو په کارولو
تعريفېدے شي۔ د دې دپاره ښيو چې بنسټيز بازګښت په رښتيا د يوه ځانګړي
منظم اقل موندلو په وسيله ترسره کېدے شي۔ خو د دې د سم کار دپاره بايد
يو شمېر نورې بنسټيزې تابعې هم ورزياتې کړو: جمع، ضرب، او د عينيت د
اړيکې ځانګړونکې تابعه~\LR{$\Char{=}$}۔ بيا قضيه په دې ښودلو ثابتولے شو
چې دا \emph{ټولې} بنسټيزې تابعې په~\LR{$\Th{Q}$} کښې تمثيلېدونکې دي، او
تمثيلېدونکې تابعې د ترکيب او منظم اقل موندلو لاندې تړلې دي۔
\olpseganchor{OLP-0290-B013}{end}

\olpseganchor{OLP-0291-B004}{start}
\olfileid{inc}{req}{rpc}
\olsection{په~\texorpdfstring{\LR{$\Th{Q}$}}{Q} کښې تمثيلېدونکې تابعې محاسبه کېدونکې دي}
\olpseganchor{OLP-0291-B004}{end}

\olpseganchor{OLP-0291-B005}{start}
موږ به ثابته کړو چې هره هغه تابعه چې په~\LR{$\Th{Q}$} کښې تمثيلېدونکې
وي، محاسبه کېدونکې ده۔ لومړے بايد په~\LR{$\Th{Q}$} کښې د تمثيلېدونکو
تابعو په اړه يوه لمه ثابته کړو۔
\olpseganchor{OLP-0291-B005}{end}

\olpseganchor{OLP-0291-B006}{start}
\begin{lem}\ollabel{lem:rep-q}
  که~\LR{$f(x_0,
  \dots, x_k)$} په~\LR{$\Th{Q}$} کښې تمثيلېدونکې وي، داسې
  فارمول~\LR{$!A_f(x_0, \dots, x_k, y)$} شته چې
  \[
  \Th{Q} \Proves !A_f(\num{n_0}, \dots, \num{n_k}, \num{m})
  \quad\text{\RL{هله او يوازې هله چې}}\quad m = f(n_0, \dots, n_k).
  \]
\end{lem}
\olpseganchor{OLP-0291-B006}{end}

\olpseganchor{OLP-0291-B007}{start}
\begin{proof}
  «که» برخه
\LR{\olref[int]{defn:representable-fn}(\olref[int]{defn:rep:a})} ده۔ «يوازې
که» برخه داسې ليدل کېږي: فرض کړئ
\LR{$\Th{Q} \Proves !A_f(\num{n_0},
\dots, \num{n_k}, \num{m})$}، خو~\LR{$m \neq f(n_0, \dots, n_k)$}۔
\LR{$l = f(n_0, \dots, n_k)$} ونيسئ۔ د
\LR{\olref[int]{defn:representable-fn}(\olref[int]{defn:rep:a})} له مخې،
\LR{$\Th{Q} \Proves !A_f(\num{n_0}, \dots, \num{n_k}, \num{l})$}۔ د
\LR{\olref[int]{defn:representable-fn}(\olref[int]{defn:rep:b})} له مخې،
\LR{$\lforall[y][(!A_f(\num{n_0}, \dots, \num{n_k}, y) \lif \num{l} =
y)]$}۔ دا فارمول په تيورۍ کښې ثابت دے۔ د منطق او دې فرض په کارولو سره چې~\LR{$\Th{Q} \Proves
!A_f(\num{n_0}, \dots, \num{n_k}, \num{m})$}، ترلاسه کوو چې
\LR{$\Th{Q} \Proves \eq[\num{l}][\num{m}]$}۔ له بلې خوا، د
\olref[bre]{lem:q-proves-neq} له مخې، \LR{$\Th{Q} \Proves
\eq/[\num{l}][\num{m}]$}۔ نو~\LR{$\Th{Q}$} ناسازګاره ده۔ خو دا ناشونې
ده، ځکه~\LR{$\Th{Q}$} په معياري مدل کښې پوره کېږي (وګورئ
\olref[int][def]{def:standard-model})، \LR{$\Sat{N}{\Th{Q}}$}، او د صحت
د قضيې له مخې هره د پوره کېدو وړ تيوري تل سازګاره وي
(\cref{fol:axd:sou:cor:consistency-soundness,fol:seq:sou:cor:consistency-soundness,fol:ntd:sou:cor:consistency-soundness,fol:tab:sou:cor:consistency-soundness})۔
\end{proof}
\olpseganchor{OLP-0291-B007}{end}

\olpseganchor{OLP-0291-B008}{start}
\begin{lem}
هره هغه تابعه چې په~\LR{$\Th{Q}$} کښې تمثيلېدونکې وي، محاسبه کېدونکې
ده۔
\end{lem}
\olpseganchor{OLP-0291-B008}{end}

\olpseganchor{OLP-0291-B009}{start}
\begin{proof}
لومړے د دې خبرې شهودي مفکوره وړاندې کړو۔ د~\LR{$f$} د محاسبې دپاره
لاندې کار کوو۔ د حساب د ژبې ټول ممکن اشتقاقونه~\LR{$\delta$} په
لړ کښې راوړو۔ دا کار په ميخانيکي ډول کېدے شي۔ د هر يوه دپاره ګورو
چې آيا د~\LR{$!A_f(\num{n_0}, \dots, \num{n_k}, \num{m})$} په بڼه د
فارمول کوم اشتقاق دے که نۀ (هماغه فارمول چې د
\olref{lem:rep-q} له مخې په~\LR{$\Th{Q}$} کښې~\LR{$f$} تمثيلوي)۔ که وي، د
\olref{lem:rep-q} له مخې~\LR{$m = f(n_0, \dots, n_k)$}، نو د~\LR{$f$} ارزښت
مو وموند۔ پلټنه پاې ته رسېږي، ځکه
\LR{$\Th{Q} \Proves !A_f(\num{n_0}, \dots, \num{n_k},
\num{f(n_0, \dots, n_k)})$}؛ نو بالاخره د سم ډول کوم~\LR{$\delta$} مومو۔
\olpseganchor{OLP-0291-B009}{end}

\olpseganchor{OLP-0291-B010}{start}
دا لا بشپړ کره نۀ دي، ځکه زموږ کړنلاره يوازې پر عددونو د کار پر ځاے
پر اشتقاقونه او فارمولونه کار کوي، او لا مو په کره توګه نۀ
دي ويلي چې ولې «د ټولو ممکنو اشتقاقونه لړل» په ميخانيکي ډول
شوني دي۔ خو لکه چې ليدلي مو دي، ترمونه، فارمولونه او
اشتقاقونه په ګوډل شمېرو کوډېدے شي۔ د محاسبې يو کره مدل، يعنې
عمومي بازګشتي تابعې، مو هم معرفي کړے دے۔ او ليدلي مو دي چې اړيکه
\LR{$\Prf[\Th{Q}](d,y)$}، چې هله صدق کوي چې~\LR{$d$} د~\LR{$\Th{Q}$} له بديهي
اصولو څخه د ګوډل شمېرې~\LR{$y$} لرونکي فارمول د کوم
اشتقاق ګوډل شمېره وي، (بنسټيزه) بازګشتي ده۔ نورې بنسټيزې
بازګشتي تابعې چې پکار مو دي~\LR{$\fn{num}$}
(\olref[art][trm]{prop:num-primrec}) او~\LR{$\fn{Subst}$}
(\olref[art][sub]{prop:subst-primrec}) دي۔ له دغو څخه~\LR{$f$} د اقل
موندلو په وسيله تعريفېدے شي؛ نو~\LR{$f$} بازګشتي ده۔
\olpseganchor{OLP-0291-B010}{end}

\olpseganchor{OLP-0291-B011}{start}
لومړے دا تعريف کړئ:
\begin{multline*}
  A(n_0, \dots, n_k, m) = \\
  \fn{Subst}(\fn{Subst}(\dots\fn{Subst}(\Gn{!A_f}, \fn{num}(n_0), \Gn{x_0}),\\ \dots),
  \fn{num}(n_k),  \Gn{x_k}), \fn{num}(m), \Gn{y})
\end{multline*}
دا پېچلې ښکاري، خو يوازې همدا تابعه ده:
\LR{$A(n_0, \dots, n_k,
m) = \Gn{!A_f(\num{n_0}, \dots, \num{n_k}, \num{m})}$}۔
\olpseganchor{OLP-0291-B011}{end}

\olpseganchor{OLP-0291-B012}{start}
اوس اړيکه~\LR{$R(n_0, \dots, n_k, s)$} وګورئ چې هله صدق کوي چې
\LR{$(s)_0$} د~\LR{$\Th{Q}$} څخه د
\LR{$!A_f(\num{n_0}, \dots, \num{n_k}, \num{(s)_1})$} د کوم
اشتقاق ګوډل شمېره وي:
\[
R(n_0, \dots, n_k, s) \quad\text{\RL{هله او يوازې هله چې}}\quad \Prf[\Th{Q}]((s)_0, A(n_0,
\dots, n_k, (s)_1))
\]
که داسې~\LR{$s$} ومومو چې~\LR{$R(n_0, \dots, n_k, s)$} پرې صدق کوي، د عددونو
داسې جوړه---\LR{$(s)_0$} او~\LR{$(s)_1$}---مو موندلې چې~\LR{$(s)_0$} د
\LR{$!A_f(\num{n_0}, \dots, \num{n_k}, \num{(s)_1})$} د کوم اشتقاق
ګوډل شمېره ده۔ نو د~\LR{$s$} لټول په غيرصوري ثبوت کښې د~\LR{$d$} او~\LR{$m$} د
جوړې د لټولو په شان دي۔ يوه محاسبه کېدونکې تابعه چې داسې~\LR{$s$}
«لټوي»، په منظم اقل موندلو تعريفېدے شي۔ په پام کښې ونيسئ چې~\LR{$R$}
منظمه ده: د هر~\LR{$n_0$}، \dots،~\LR{$n_k$} دپاره د
\LR{$\Th{Q} \Proves !A_f(\num{n_0}, \dots, \num{n_k},
\num{f(n_0, \dots, n_k)})$} کوم اشتقاق~\LR{$\delta$} شته؛ نو
\LR{$R(n_0, \dots, n_k, s)$} د
\LR{$s = \tuple{\Gn{\delta}, f(n_0, \dots, n_k)}$} دپاره صدق کوي۔ نو~\LR{$f$}
داسې ليکلے شو:
\[
f(n_0,\dots,n_{k}) = (\umin{s}{R(n_0, \dots, n_k, s)})_1.
\]
\end{proof}
\olpseganchor{OLP-0291-B012}{end}

\paragraph{د منجمدې سرچينې بشپړ بديل متن}
\olpseganchor{OLP-0648-B004}{start}
\olfileid{inc}{req}{cee-olp0648}
\olsection{د \texorpdfstring{\LR{$C$}}{C} تابعې}
\olpseganchor{OLP-0648-B004}{end}

\olpseganchor{OLP-0648-B005}{start}
\LR{$C$} دې د تابعو تر ټولو کوچنے سټ وي چې دا تابعې رانغاړي:
\begin{enumerate}
\item \LR{$0$}،
\item تالي،
\item جمع،
\item ضرب،
\item پروجکشنونه، او
\item د عينيت مشخصه تابعه، \LR{$\Char{=}$}؛
\end{enumerate}
او د دغو عمليو لاندې تړلے وي:
\begin{enumerate}
\item ترکيب، او
\item بې‌حده پلټنه، چې پر منظمو تابعو تطبيقېږي۔
\end{enumerate}
ياد ولرئ چې دا وروستنے قيد يوازې دا معنا لري چې د \LR{$\mu$} عمليه هله کارولے شئ چې ترې ترلاسه شوې تابعه تامه وي۔ دا د \emph{عمومي بازګشتي} تابعو له تعريف سره پرتله کړئ: دلته مو جمع، ضرب او \LR{$\Char{=}$} ورزيات کړي دي، خو بنسټيز بازګښت مو لرې کړے دے۔
\olpseganchor{OLP-0648-B005}{end}

\olpseganchor{OLP-0648-B006}{start}
څرګنده ده چې په \LR{$C$} کښې هره تابعه بازګشتي ده، ځکه جمع، ضرب او \LR{$\Char{=}$} بازګشتي دي۔ وبه ښيو چې برعکس هم صحيح دے؛ يعنې د \LR{$C$} د نورو اجزاوو په مرسته بنسټيز بازګښت ترسره کولے شو۔
\olpseganchor{OLP-0648-B006}{end}

\olpseganchor{OLP-0292-B004}{start}
\olfileid{inc}{req}{bet}
\olsection{د بېټا تابعې لمه}
\olpseganchor{OLP-0292-B004}{end}


\olpseganchor{OLP-0292-B005}{start}
د دې د ښودلو دپاره چې که جمع، ضرب، او~\LR{$\Char{=}$} راسره وي نو
بنسټيز بازګښت ترسره کولے شو، بايد داسې تابعې جوړې کړو چې لړۍ سمبال
کړي۔ (که توان هم راسره وے، کار به مو اسان و۔) کله چې بنسټيز بازګښت
راسره و، د «\LR{$n$}م اول عدد» په شان څيزونه مو تعريفولے شول او نسبتاً
ساده کوډونه مو غوره کولے شول۔ خو دلته بنسټيز بازګښت راسره نشته---په
حقيقت کښې غواړو وښيو چې بنسټيز بازګښت د اقل موندلو په کارولو ترسره
کولے شو---نو بايد څه زياته ځيرکتيا وکاروو۔
\olpseganchor{OLP-0292-B005}{end}

\olpseganchor{OLP-0292-B006}{start}
\begin{lem}
\ollabel{lem:beta}
داسې تابعه~\LR{$\beta(d,i)$} شته چې د هرې لړۍ \LR{$a_0$}، \dots،~\LR{$a_n$}
دپاره داسې عدد~\LR{$d$} شته چې د هر~\LR{$i \le n$} دپاره
\LR{$\beta(d,i) = a_i$} وي۔ سربېره پر دې،~\LR{$\beta$} له بنسټيزو تابعو څخه
يوازې د ترکيب او منظم اقل موندلو په کارولو تعريفېدے شي۔
\end{lem}
\olpseganchor{OLP-0292-B006}{end}

\olpseganchor{OLP-0292-B007}{start}
\LR{$d$} داسې وګڼئ چې لړۍ~\LR{$\tuple{a_0, \dots, a_n}$} کوډوي، او
\LR{$\beta(d,i)$} ئې~\LR{$i$}م غړے بېرته ورکوي۔ (پام وکړئ چې دا «کوډول» هغه
د اول عددونو د توانونو کوډول \emph{نۀ} کاروي چې لا له وړاندې ورسره
اشنا يو!) لمه ډېره لږه ادعا کوي؛ دا نۀ وايي چې لړۍ يو ځاے کولے شو،
غړي ورزياتولے شو، يا ان دا چې د ترکيب او منظم اقل موندلو له لارې
تعريفېدونکو تابعو په کارولو له~\LR{$a_0$}، \dots،~\LR{$a_n$} څخه~\LR{$d$}
\emph{محاسبه} کولے شو۔ يوازې دومره وايي چې د «کوډ پرانيستلو» داسې
تابعه شته چې هره لړۍ «کوډ شوې» وي۔
\olpseganchor{OLP-0292-B007}{end}

\olpseganchor{OLP-0292-B008}{start}
د~\LR{$\beta$} نښه د ګوډل ده۔ بيا وايو، د لمې د ثبوت ګرانه برخه دا ده
چې د ظاهراً محدودو وسيلو، يعنې يوازې د ترکيب او اقل موندلو، په
کارولو مناسبه~\LR{$\beta$} تعريف کړو---خو اجازه لرو چې جمع، ضرب،
او~\LR{$\Char{=}$} وکاروو۔ د دې لمې د ثبوت بېلابېلې لارې شته، خو يوه له
تر ټولو پاکو لارو څخه لا هم د ګوډل اصلي لاره ده، چې د عددونو د تيورۍ
هغه حقيقت ئې کاراوه چې د سونزي قضيه بلل کېږي (په دوديز ډول، «د چين
د پاتې شونو قضيه»)۔
\olpseganchor{OLP-0292-B008}{end}

\olpseganchor{OLP-0292-B009}{start}
\begin{defn}
دوه طبيعي عددونه~\LR{$a$} او~\LR{$b$} هله \emph{خپلمنځي اول} دي چې ستر
ګډ مقسوم‌عليه ئې~\LR{$1$} وي؛ په بله وينا، له دې پرته بل ګډ مقسوم‌عليه
نۀ لري۔
\end{defn}
\olpseganchor{OLP-0292-B009}{end}

\olpseganchor{OLP-0292-B010}{start}
\begin{defn}
طبيعي عددونه~\LR{$a$} او~\LR{$b$} \emph{د~\LR{$c$} په مودولو سره مطابق} دي،
\LR{$a \equiv b \mod c$}، هله چې~\LR{$c \mid (a-b)$}؛ يعنې کله چې پر~\LR{$c$}
ووېشل شي،~\LR{$a$} او~\LR{$b$} يو شان پاتې شونې لري۔
\end{defn}
\olpseganchor{OLP-0292-B010}{end}

\olpseganchor{OLP-0292-B011}{start}
دا د سونزي قضيه ده:
\begin{thm}
فرض کړئ \LR{$x_0$}، \dots،~\LR{$x_n$} (جوړه په جوړه) خپلمنځي اول دي۔
\LR{$y_0$}، \dots،~\LR{$y_n$} هر عددونه وي۔ بيا داسې عدد~\LR{$z$} شته چې
\begin{align*}
z & \equiv y_0 \mod x_0 \\
z & \equiv y_1 \mod x_1 \\
& \vdots  \\
z & \equiv y_n \mod x_n.
\end{align*}
\end{thm}
\olpseganchor{OLP-0292-B011}{end}

\olpseganchor{OLP-0292-B012}{start}
د سونزي قضيه به داسې کاروو: که \LR{$x_0$}، \dots،~\LR{$x_n$} په ترتيب سره
له~\LR{$y_0$}، \dots،~\LR{$y_n$} څخه لوی وي،~\LR{$z$} د لړۍ
\LR{$\tuple{y_0, \dots,y_n}$} کوډ ګڼلے شو۔ د~\LR{$y_i$} د بېرته موندلو
دپاره يوازې~\LR{$z$} پر~\LR{$x_i$} وېشو او پاتې شونې ئې اخلو۔ د دې کوډ کارولو
دپاره بايد د~\LR{$x_0$}، \dots،~\LR{$x_n$} مناسب ارزښتونه ومومو۔
\olpseganchor{OLP-0292-B012}{end}

\olpseganchor{OLP-0292-B013}{start}
په دې برخه کښې به څو کتنې راسره مرسته وکړي۔ د ورکړل شوو
\LR{$y_0$}، \dots،~\LR{$y_n$} دپاره دا ونيسئ:
\begin{align*}
j &= \max(n, y_0 + 1, \dots, y_n + 1), \\
m &= \lcm(1,\dots,j),
\end{align*}
او دا ونيسئ:
\begin{align*}
x_0 & = 1 + m \\
x_1 & = 1 + 2 \cdot m \\
x_2 & = 1 + 3 \cdot m \\
& \vdots  \\
x_n & = 1 + (n+1) \cdot m
\end{align*}
بيا دوه خبرې سمې دي:
\begin{enumerate}
\item\ollabel{rel-prime} \LR{$x_0,\dots,x_n$} خپلمنځي اول دي۔
\item\ollabel{less} د هر~\LR{$i$} دپاره، \LR{$y_i < x_i$}۔
\end{enumerate}
د دې ليدلو دپاره چې~\olref{rel-prime} سمه ده، پام وکړئ چې که~\LR{$p$}
اول عدد وي او \LR{$p \mid x_i$} او~\LR{$p \mid x_k$} وي، نو
\LR{$p \mid 1 + (i+1) m$} او~\LR{$p \mid 1 + (k+1) m$} دي۔ خو بيا~\LR{$p$} د
دواړو توپير وېشي:
\[
(1 + (i+1)m) - (1+ (k+1)m) = (i-k) m.
\]
څرنګه چې~\LR{$p$}، \LR{$1 + (i+1)m$} وېشي، په عين وخت کښې~\LR{$m$} نشي وېشلے
(که ئې وېشلے، لومړۍ وېشنې به~\LR{$1$} پاتې شونې پرېښې وه)۔ نو ځکه چې
\LR{$p$}، \LR{$(i-k)m$} وېشي،~\LR{$p$} بايد~\LR{$i-k$} ووېشي۔ خو~\LR{$\left|i-k\right|$} تر
ډېره~\LR{$n$} دے، او موږ~\LR{$j \geq n$} غوره کړے دے، نو له دې څخه بيا
\LR{$p \mid m$} لازمېږي، چې تناقض دے۔ نو هېڅ اول عدد~\LR{$x_i$} او~\LR{$x_k$}
دواړه نۀ وېشي۔ فقره~\olref{less} اسانه ده:
\LR{$y_i < j \leq m < x_i$} لرو۔
\olpseganchor{OLP-0292-B013}{end}

\olpseganchor{OLP-0292-B014}{start}
اوس د~\LR{$\beta$} تابعې لمه ثابته کړو۔ په ياد ولرئ چې~\LR{$0$}، تالي،
جمع، ضرب،~\LR{$\Char{=}$}، پروجکشنونه، او هره هغه تابعه کارولے شو چې
له دوى څخه د ترکيب او پر منظمو تابعو د اقل موندلو په کارولو تعريف
شوې وي۔ که د کومې اړيکې ځانګړونکې تابعه همداسې تعريفېدونکې وي، هغه
اړيکه هم کارولے شو۔ د پخوا په شان ښودلے شو چې دا اړيکې د بولي
ترکيبونو او محدود کمیت ايښودنې په وړاندې تړلې دي؛ د بېلګې په توګه:
\begin{align*}
\fn{not}(x) & \defis \Char{=}(x,0)\\
\bmin{x \leq z}{R(x,y)} & \defis \umin{x}{(R(x,y) \lor x = z)}\\
\bexists{x \leq z}{R(x,y)} & \defiff R(\bmin{x \leq z}{R(x,y)}, y)
\end{align*}
بيا ښودلے شو چې دا ټول هم له بنسټيز بازګښت پرته تعريفېدونکي دي:
\begin{enumerate}
\item جوړوونکې تابعه، \LR{$J(x,y) = \frac{1}{2}[(x+y)(x+y+1)] + x$}؛
% maybe explain more what is going on here, a bit confusing.
\item د پروجکشن تابعې
\begin{align*}
K(z) & = \bmin{x \leq z}{\bexists{y \leq z}{z = J(x,y)}},\\
L(z) & = \bmin{y \leq z}{\bexists{x \leq z}{z = J(x,y)}};
\end{align*}
\item د کوچني والي اړيکه~\LR{$x < y$}؛
\item د وېش اړيکه~\LR{$x \mid y$}؛
% \item $x \tsub y$
% \item $\fn{Prime}(x)$
% \item Assuming $p$ is prime, the relation ``$x$ is a power of $p$'':
% \[
% \bforall{y \leq x}{(y \mid x \lif y = 1 \lor y = x)}.
% \]
\item تابعه~\LR{$\fn{rem}(x,y)$}، چې د~\LR{$y$} پر~\LR{$x$} د وېش پاتې شونې
  بېرته ورکوي۔
\end{enumerate}
اوس تعريف کړئ:
\begin{align*}
\beta^*(d_0,d_1,i) & = \fn{rem}(1+(i+1) d_1,d_0) \text{\RL{ او}}\\
\beta(d,i) & = \beta^*(K(d),L(d),i).
\end{align*}
دا هماغه تابعه ده چې غواړو ئې۔ پورته غوندې ورکړل شوو
\LR{$a_0,\dots,a_n$} دپاره دا ونيسئ:
\[
j = \max(n,a_0+1,\dots,a_n+1),
\]
او~\LR{$d_1 = \lcm(1,\dots,j)$} ونيسئ۔ له پورته
\olref{rel-prime} څخه پوهېږو چې \LR{$1+d_1$}، \LR{$1+2 d_1$}، \dots،
\LR{$1+(n+1) d_1$} خپلمنځي اول دي، او له~\olref{less} څخه چې دا ټول
له~\LR{$a_0,\dots,a_n$} څخه لوی دي۔ د سونزي د قضيې له مخې داسې ارزښت
\LR{$d_0$} شته چې د هر~\LR{$i$} دپاره
\[
d_0 \equiv a_i \mod (1+(i+1)d_1)
\]
وي، او نو (ځکه~\LR{$d_1$} له~\LR{$a_i$} څخه لوی دے)،
\[
a_i = \fn{rem}(1+(i+1)d_1,d_0).
\]
\LR{$d = J(d_0,d_1)$} ونيسئ۔ بيا د هر~\LR{$i \le n$} دپاره لرو:
\begin{align*}
\beta(d,i) & =  \beta^*(d_0,d_1,i) \\
& =  \fn{rem}(1+(i+1) d_1,d_0) \\
& =  a_i
\end{align*}
او دا هماغه څه دي چې پکار مو وو۔ په دې سره د~\LR{$\beta$} تابعې د لمې
ثبوت بشپړېږي۔
\olpseganchor{OLP-0292-B014}{end}

\olpseganchor{OLP-0292-B015}{start}
\begin{prob}
  وښيئ چې اړيکې~\LR{$x < y$}، \LR{$x \mid y$}، او تابعه
  \LR{$\fn{rem}(x,y)$} له بنسټيز بازګښت پرته تعريفېدے شي۔ \LR{$0$}، تالي،
  جمع، ضرب،~\LR{$\Char{=}$}، پروجکشنونه، او محدوده کمينه‌موندنه او کمیت
  ايښودنه کارولے شئ۔
\end{prob}
\olpseganchor{OLP-0292-B015}{end}

\olpseganchor{OLP-0293-B004}{start}
\olfileid{inc}{req}{pri}
\olsection{د بنسټيز بازګښت شبيه کول}
\olpseganchor{OLP-0293-B004}{end}

\olpseganchor{OLP-0293-B005}{start}
اوس ښودلے شو چې د بنسټيز بازګښت په وسيله تعريف د بېټا تابعې په
کارولو سره د منظم اقل موندلو له لارې «شبيه» کېدے شي۔ فرض کړئ
\LR{$f(\vec x)$} او~\LR{$g(\vec x, y, z)$} لرو۔ بيا هغه تابعه~\LR{$h(\vec x,y)$}
چې له~\LR{$f$} او~\LR{$g$} څخه په بنسټيز بازګښت تعريف شوې وي، دا ده:
\begin{align*}
h(\vec x, 0) & =  f(\vec x) \\
h(\vec x, y+1) & =  g(\vec x, y, h(\vec x, y)).
\end{align*}
بايد وښيو چې~\LR{$h$} يوازې د ترکيب او منظم اقل موندلو په کارولو، د
بنسټيزو تابعو او هغو تابعو په مرسته چې له بنسټيزو تابعو څخه د ترکيب
او منظم اقل موندلو له لارې تعريف شوې وي (لکه~\LR{$\beta$})، له~\LR{$f$} او
\LR{$g$} څخه تعريفېدے شي۔
\olpseganchor{OLP-0293-B005}{end}

\olpseganchor{OLP-0293-B006}{start}
\begin{lem}
\ollabel{lem:prim-rec}
که~\LR{$h$} له~\LR{$f$} او~\LR{$g$} څخه په بنسټيز بازګښت تعريفېدے شي، نو له~\LR{$f$}،
\LR{$g$}، تابعو~\LR{$\Zero$}، \LR{$\Succ$}، \LR{$\Proj{n}{i}$}، \LR{$\Add$}، \LR{$\Mult$}،
\LR{$\Char{=}$} څخه د ترکيب او منظم اقل موندلو په کارولو تعريفېدے شي۔
\end{lem}
\olpseganchor{OLP-0293-B006}{end}

\olpseganchor{OLP-0293-B007}{start}
\begin{proof}
لومړے يوه مرستندويه تابعه~\LR{$\hat h(\vec x, y)$} تعريف کړئ چې تر ټولو
وړوکے عدد~\LR{$d$} بېرته ورکوي چې~\LR{$d$} داسې لړۍ کوډوي چې لاندې شرطونه پوره
کوي:
\begin{enumerate}
\item \LR{$(d)_0 = f(\vec x)$}، او
\item د هر~\LR{$i < y$} دپاره، \LR{$(d)_{i+1} = g(\vec x, i, (d)_i)$}،
\end{enumerate}
چېرته چې اوس~\LR{$(d)_i$} د~\LR{$\beta(d,i)$} لنډون دے۔ په بله وينا،
\LR{$\hat h$} د دې لړۍ تر ټولو وړوکے کوډ بېرته ورکوي:
\LR{$\tuple{h(\vec x, 0), h(\vec x, 1), \dots, h(\vec x, y)}$}۔
\LR{$\hat h$} داسې ليکلے شو:
\[
\hat h(\vec x, y) = \umin{d}{(\beta(d,0) = f(\vec x) \land \bforall{i <
  y}{\beta(d,i+1) = g(\vec x, i,\beta(d,i))})}.
\]
پام وکړئ: دلته هېڅ بنسټيز بازګښت پکار نۀ دے، يوازې اقل موندل دي۔
محدود کميت ټاکل د ترکيب او منظم اقل موندلو په کارولو تعريفېږي، لکه د بېټا تابعې په برخه کښې۔ د هرې ورودي جوړې د بازګښت د ارزښتونو لړۍ متناهي ده۔ د بېټا تابعې لمه \olref[bet]{lem:beta} د دغې لړۍ يو کوډ تضمينوي، نو د اقل موندلو شرط تل يو شاهد لري۔
\olpseganchor{OLP-0293-B007}{end}

\olpseganchor{OLP-0293-B008}{start}
خو اوس لرو:
\[
h(\vec x, y) = \beta(\hat h(\vec x, y), y),
\]
نو~\LR{$h$} له بنسټيزو تابعو څخه يوازې د ترکيب او منظم اقل موندلو په
کارولو تعريفېدے شي۔
\end{proof}
\olpseganchor{OLP-0293-B008}{end}

\olpseganchor{OLP-0294-B004}{start}
\olfileid{inc}{req}{bre}
\olsection{بنسټيزې تابعې په~\texorpdfstring{\LR{$\Th{Q}$}}{Q} کښې تمثيلېدونکې دي}
\olpseganchor{OLP-0294-B004}{end}

\olpseganchor{OLP-0294-B005}{start}
لومړے بايد وښيو چې ټولې بنسټيزې تابعې په~\LR{$\Th{Q}$} کښې
تمثيلېدونکې دي۔ په پاې کښې بايد وښيو چې هرې~\LR{$k$}-ځايه بنسټيزې
تابعې~\LR{$f(x_0,\dots,x_{k-1})$} ته داسې فارمول
\LR{$!A_f(x_0,\dots,x_{k-1},y)$} څنګه وټاکو چې هغه تمثيل کړي۔
\olpseganchor{OLP-0294-B005}{end}

\olpseganchor{OLP-0294-B006}{start}
صفر، تالي، جمع، ضرب، د مساواتو ځانګړونکې تابعه، او پروجکشنونه
تمثيلولے شو۔ په هره قضيه کښې مناسب تمثيلوونکی فارمول په بشپړه توګه
ساده دے؛ د بېلګې په توګه، صفر د~\LR{$y = \Obj 0$} فارمول په وسيله
تمثيلېږي، تالي د~\LR{$x_0' = y$} فارمول په وسيله تمثيلېږي، او جمع
د~\LR{$(x_0 + x_1) = y$} فارمول په وسيله تمثيلېږي۔ کار په دې کښې دے چې وښيو
\LR{$\Th{Q}$} اړوند تړلي فارمولونه ثابتولے شي؛ د بېلګې په توګه، دا ويل
چې جمع د پورته فارمول په وسيله تمثيل شوې ده، دا غواړي چې وښيو
د طبيعي عددونو د هرې جوړې~\LR{$m$} او~\LR{$n$} دپاره~\LR{$\Th{Q}$} دا ثابتوي:
\begin{align*}
& \eq[\num n + \num m][\num {n+m}] \text{\RL{ او}}\\
& \lforall[y][(\eq[(\num n + \num m)][y] \lif \eq[y][\num{n+m}])].
\end{align*}
\olpseganchor{OLP-0294-B006}{end}

\olpseganchor{OLP-0294-B007}{start}
\begin{prop}
\ollabel{prop:rep-zero}
صفر تابعه~\LR{$\Zero(x) = 0$} په~\LR{$\Th{Q}$} کښې د
\LR{$!A_{\Zero}(x,y) \ident \eq[y][\Obj 0]$} په وسيله تمثيلېږي۔
\end{prop}
\olpseganchor{OLP-0294-B007}{end}

\olpseganchor{OLP-0294-B008}{start}
\begin{prop}
\ollabel{prop:rep-succ}
تالي تابعه~\LR{$\Succ(x) = x+1$} په~\LR{$\Th{Q}$} کښې د
\LR{$!A_{\Succ}(x,y) \ident \eq[y][x']$} په وسيله تمثيلېږي۔
\end{prop}
\olpseganchor{OLP-0294-B008}{end}

\olpseganchor{OLP-0294-B009}{start}
\begin{prop}
\ollabel{prop:rep-proj}
پروجکشن تابعه~\LR{$\Proj{n}{i}(x_0, \dots, x_{n-1}) = x_i$} په
\LR{$\Th{Q}$} کښې د دې په وسيله تمثيلېږي:
\[!A_{\Proj{n}{i}}(x_0, \dots, x_{n-1}, y) \ident \eq[y][x_i].\]
\end{prop}
\olpseganchor{OLP-0294-B009}{end}

\olpseganchor{OLP-0294-B010}{start}
\begin{prob}
ثابته کړئ چې~\LR{$\eq[y][\Obj 0]$}، \LR{$\eq[y][x']$}، او~\LR{$\eq[y][x_i]$} په
ترتيب سره~\LR{$\Zero$}، \LR{$\Succ$}، او~\LR{$\Proj{n}{i}$} تمثيلوي۔
\end{prob}
\olpseganchor{OLP-0294-B010}{end}

\olpseganchor{OLP-0294-B011}{start}
\begin{prop}
\ollabel{prop:rep-id}
د~\LR{$=$} ځانګړونکې تابعه،
\[
\Char{=}(x_0, x_1) =
\begin{cases}
  1 & \text{\RL{که }} x_0 =x_1\\
  0 & \text{\RL{بل ډول}}
\end{cases}
\]
په~\LR{$\Th{Q}$} کښې د دې په وسيله تمثيلېږي:
\[
  !A_{\Char{=}}(x_0, x_1, y) \ident (\eq[x_0][x_1] \land \eq[y][\num{1}]) \lor (\eq/[x_0][x_1] \land
\eq[y][\num{0}]).
\]
\end{prop}
\olpseganchor{OLP-0294-B011}{end}

\olpseganchor{OLP-0294-B012}{start}
ثبوت لاندې لمې ته اړتيا لري۔
\olpseganchor{OLP-0294-B012}{end}

\olpseganchor{OLP-0294-B013}{start}
\begin{lem}
\ollabel{lem:q-proves-neq} که طبيعي عددونه~\LR{$n$} او~\LR{$m$} داسې وي چې
\LR{$n \neq m$}، نو \LR{$\Th{Q} \Proves \eq/[\num n][\num m]$}۔
\end{lem}
\olpseganchor{OLP-0294-B013}{end}

\olpseganchor{OLP-0294-B014}{start}
\begin{proof}
پر~\LR{$n$} استقرا وکاروئ څو وښيئ چې د هر~\LR{$m$} دپاره، که~\LR{$n \neq m$}
وي، نو \LR{$Q \Proves \eq/[\num n][\num m]$}۔
\olpseganchor{OLP-0294-B014}{end}

\olpseganchor{OLP-0294-B015}{start}
په بنسټيزه قضيه کښې~\LR{$n = 0$} دے۔ که~\LR{$m$} له~\LR{$0$} سره مساوي نۀ وي، نو
د کوم طبيعي عدد~\LR{$k$} دپاره~\LR{$m = k + 1$}۔ يو داسې بديهي اصل لرو چې
\LR{$\lforall[x][\eq/[0][x']]$} وايي۔ د کمیت ايښوونکي د يوه بديهي اصل له
مخې، د~\LR{$x$} پر ځاے د~\LR{$\num k$} په ايښودلو، دا پايله اخلو:
\LR{$\eq/[0][\num k']$}۔ خو~\LR{$\num k'$} هماغه~\LR{$\num m$} دے۔
\olpseganchor{OLP-0294-B015}{end}

\olpseganchor{OLP-0294-B016}{start}
د استقرا په پړاو کښې فرضولے شو چې ادعا د~\LR{$n$} دپاره سمه ده، او
\LR{$n+1$} ګورو۔ \LR{$m$} هر طبيعي عدد ونيسئ۔ دوه امکانونه شته: يا~\LR{$m = 0$}
دے يا د کوم~\LR{$k$} دپاره~\LR{$m = k+1$}۔ لومړۍ قضيه لکه پورته حل کېږي۔ په
دويمه قضيه کښې فرض کړئ~\LR{$n+1 \neq k+1$}۔ بيا~\LR{$n \neq k$}۔ د~\LR{$n$}
دپاره د استقرايي فرضيې له مخې
\LR{$\Th{Q} \Proves \eq/[\num n][\num k]$}۔ داسې بديهي اصل لرو چې
\LR{$\lforall[x][\lforall[y][\eq[x'][y'] \lif \eq[x][y]]]$} وايي۔ د
کمیت ايښوونکي د يوه بديهي اصل په کارولو، لرو:
\LR{$\eq[\num n'][\num k'] \lif \eq[\num n][\num k]$}۔ د قضيو د منطق په
کارولو، په~\LR{$\Th{Q}$} کښې دا پايله اخلو:
\LR{$\eq/[\num n][\num k] \lif \eq/[\num n'][\num k']$}۔ د وضع مقدم په
کارولو، \LR{$\eq/[\num n'][\num k']$} اخلو، چې هماغه غوښتنه ده، ځکه
\LR{$\num k'$}، \LR{$\num m$} دے۔
\end{proof}
\olpseganchor{OLP-0294-B016}{end}

\olpseganchor{OLP-0294-B017}{start}
\begin{explain}
پام وکړئ چې لمه ډېره خبره نۀ کوي: په اصل کښې وايي چې~\LR{$\Th{Q}$}
ثابتولے شي چې بېلابېلې عددنښې بېلابېل څيزونه ښيي۔ د بېلګې په
توګه،~\LR{$\Th{Q}$}، \LR{$0'' \neq 0'''$} ثابتوي۔ خو دا ښودل چې خبره په عام
ډول سمه ده څه پاملرنه غواړي۔ دا هم په پام کښې ونيسئ چې که څه هم
استقرا کاروو، خو دا د~\LR{$\Th{Q}$} \emph{بهر} استقرا ده۔
\end{explain}
\olpseganchor{OLP-0294-B017}{end}

\olpseganchor{OLP-0294-B018}{start}
\begin{proof}[د \olref{prop:rep-id} ثبوت]
که~\LR{$n = m$} وي، نو~\LR{$\num{n}$} او~\LR{$\num{m}$} يو شان ترمونه دي، او
\LR{$\Char{=}(n, m) = 1$}۔ خو \LR{$\Th{Q} \Proves (\eq[\num{n}][\num{m}] \land
\eq[\num{1}][\num{1}])$}، نو~\LR{$!A_=(\num{n}, \num{m},
\num{1})$} ثابتوي۔ که~\LR{$n \neq m$} وي، نو~\LR{$\Char=(n, m) = 0$}۔ د
\olref{lem:q-proves-neq} له مخې،
\LR{$\Th{Q} \Proves \eq/[\num{n}][\num{m}]$}، او نو دا هم ثابتوي:
\LR{$(\eq/[\num{n}][\num{m}] \land \Obj 0 = \Obj 0)$}۔ له دې امله
\LR{$\Th{Q} \Proves !A_=(\num{n}, \num{m}, \num{0})$}۔
\olpseganchor{OLP-0294-B018}{end}

\olpseganchor{OLP-0294-B019}{start}
د دويمې برخې دپاره هم دوه قضيې لرو۔ که~\LR{$n = m$} وي، بايد وښيو:
\LR{$\Th{Q} \Proves \lforall[y][(!A_=(\num{n}, \num{m}, y)
  \lif \eq[y][\num{1}])]$}۔ په غيرصوري ډول استدلال کوو، فرض کړئ
\LR{$!A_=(\num{n}, \num{m}, y)$}، يعنې
\[
(\eq[\num{n}][\num{n}] \land \eq[y][\num{1}]) \lor
(\eq/[\num{n}][\num{n}] \land \eq[y][\num{0}])
\]
کيڼ منفصل د منطق له مخې~\LR{$\eq[y][\num{1}]$} لازموي؛ ښے منفصل له
\LR{$\eq[\num{n}][\num{n}]$} سره تناقض لري، چې په منطق ثابتېدونکے دے۔
\olpseganchor{OLP-0294-B019}{end}

\olpseganchor{OLP-0294-B020}{start}
بل پلو، فرض کړئ~\LR{$n \neq m$}۔ بيا~\LR{$!A_=(\num{n},
\num{m}, y)$} دا دے:
\[
(\eq[\num{n}][\num{m}] \land \eq[y][\num{1}]) \lor
(\eq/[\num{n}][\num{m}] \land \eq[y][\num{0}])
\]
دلته کيڼ منفصل له~\LR{$\eq/[\num{n}][\num{m}]$} سره تناقض لري، چې د
\olref{lem:q-proves-neq} له مخې په~\LR{$\Th{Q}$} کښې ثابتېدونکے دے؛ ښے
منفصل~\LR{$\eq[y][\num{0}]$} لازموي۔
\end{proof}
\olpseganchor{OLP-0294-B020}{end}

\olpseganchor{OLP-0294-B021}{start}
\begin{prop}
\ollabel{prop:rep-add}
د جمع تابعه~\LR{$\Add(x_0, x_1) = x_0+x_1$} په~\LR{$\Th{Q}$} کښې د دې په
وسيله تمثيلېږي:
\[
  !A_{\Add}(x_0, x_1, y) \ident \eq[y][(x_0 + x_1)].
\]
\end{prop}
\olpseganchor{OLP-0294-B021}{end}

\olpseganchor{OLP-0294-B022}{start}
\begin{lem}
\ollabel{lem:q-proves-add}
\LR{$\Th{Q} \Proves \eq[(\num{n} + \num{m})][\num{n+m}]$}
\end{lem}
\olpseganchor{OLP-0294-B022}{end}

\olpseganchor{OLP-0294-B023}{start}
\begin{proof}
دا پر~\LR{$m$} په استقرا ثابتوو۔ که~\LR{$m = 0$} وي، ادعا دا ده:
\LR{$\Th{Q} \Proves \eq[(\num{n} + \Obj 0)][\num{n}]$}۔ دا له بديهي
اصل~\LR{$!Q_4$} څخه راځي۔ اوس د~\LR{$m$} دپاره ادعا فرض کړئ؛ راځئ د~\LR{$m+1$}
دپاره ادعا ثابته کړو، يعنې وښيو:
\LR{$\Th{Q} \Proves \eq[(\num{n} +
  \num{m+1})][\num{n+m+1}]$}۔ پام وکړئ چې~\LR{$\num{m+1}$} هماغه
\LR{$\num{m}'$} دے، او~\LR{$\num{n+m+1}$} هماغه~\LR{$\num{n+m}'$} دے۔ د بديهي
اصل~\LR{$!Q_5$} له مخې، \LR{$\Th{Q} \Proves
\eq[(\num{n} + \num{m}')][(\num{n}+\num{m})']$}۔ د استقرايي فرضيې
له مخې، \LR{$\Th{Q} \Proves
\eq[(\num{n} + \num{m})][\num{n+m}]$}۔ نو
\LR{$\Th{Q} \Proves \eq[(\num{n} + \num{m}')][\num{n+m}']$}۔
\end{proof}
\olpseganchor{OLP-0294-B023}{end}

\olpseganchor{OLP-0294-B024}{start}
\begin{proof}[د \olref{prop:rep-add} ثبوت]
هغه فارمول~\LR{$!A_\Add(x_0, x_1, y)$} چې~\LR{$\Add$} تمثيلوي،
\LR{$\eq[y][(x_0 + x_1)]$} ده۔ لومړے ښيو چې که~\LR{$\Add(n, m) = k$} وي،
نو \LR{$\Th{Q} \Proves !A_\Add(\num{n}, \num{m}, \num{k})$}؛ يعنې
\LR{$\Th{Q} \Proves \eq[\num{k}][(\num{n} + \num{m})]$}۔ خو ځکه چې
\LR{$k = n + m$}، نو~\LR{$\num{k}$} هماغه~\LR{$\num{n+m}$} دے، او په
\olref{lem:q-proves-add} کښې مو ښودلې ده چې
\LR{$\Th{Q} \Proves \eq[(\num{n} + \num{m})][\num{n+m}]$}۔
\olpseganchor{OLP-0294-B024}{end}

\olpseganchor{OLP-0294-B025}{start}
دا هم بايد وښيو چې که~\LR{$\Add(n, m) = k$} وي، نو
\[
\Th{Q} \Proves \lforall[y][(!A_\Add(\num{n}, \num{m}, y) \lif
  \eq[y][\num{k}])].
\]
فرض کړئ~\LR{$\eq[(\num{n} + \num{m})][y]$} لرو۔ ځکه چې
\[
\Th{Q} \Proves \eq[(\num{n}+\num{m})][\num{n+m}],
\]
کيڼ اړخ په~\LR{$\num{n+m}$} بدلولے شو او د خپل سري~\LR{$y$} دپاره
\LR{$\eq[\num{n+m}][y]$} تر لاسه کوو۔
\end{proof}
\olpseganchor{OLP-0294-B025}{end}

\olpseganchor{OLP-0294-B026}{start}
\begin{prop}
\ollabel{prop:rep-mult}
د ضرب تابعه~\LR{$\Mult(x_0, x_1) = x_0 \cdot x_1$} په~\LR{$\Th{Q}$} کښې د
دې په وسيله تمثيلېږي:
\[
  !A_{\Mult}(x_0, x_1, y) \ident y = (x_0 \times x_1).
\]
\end{prop}
\olpseganchor{OLP-0294-B026}{end}

\olpseganchor{OLP-0294-B027}{start}
\begin{proof} تمرین۔ \end{proof}
\olpseganchor{OLP-0294-B027}{end}

\olpseganchor{OLP-0294-B028}{start}
\begin{lem}
\ollabel{lem:q-proves-mult}
\LR{$\Th{Q} \Proves \eq[(\num{n} \times \num{m})][\num{n \cdot m}]$}
\end{lem}
\olpseganchor{OLP-0294-B028}{end}

\olpseganchor{OLP-0294-B029}{start}
\begin{proof} تمرین۔ \end{proof}
\olpseganchor{OLP-0294-B029}{end}

\olpseganchor{OLP-0294-B030}{start}
\begin{prob}
\olref[inc][req][bre]{lem:q-proves-mult} ثابته کړئ۔
\end{prob}
\olpseganchor{OLP-0294-B030}{end}

\olpseganchor{OLP-0294-B031}{start}
\begin{prob}
له \olref[inc][req][bre]{lem:q-proves-mult} څخه د
\olref[inc][req][bre]{prop:rep-mult} په ثبوت کښې کار واخلئ۔
\end{prob}
\olpseganchor{OLP-0294-B031}{end}

\olpseganchor{OLP-0294-B032}{start}
\begin{explain}
  په ياد ولرئ چې د حساب د ژبې د تابعې نښې دپاره~\LR{$\times$}، او پر
  عددونو د عادي ضرب د عمل دپاره~\LR{$\cdot$} کاروو۔ نو~\LR{$\cdot$} د عددونو
  د تعبيرونو ترمنځ راتلے شي (لکه په~\LR{$m \cdot n$} کښې)، خو~\LR{$\times$}
  يوازې د حساب د ژبې د ترمونو ترمنځ راځي (لکه په
  \LR{$(\num{m} \times \num{n})$} کښې)۔ تر دې هم مغشوشوونکې دا ده چې
  \LR{$+$} هم د تابع او هم د جمع د عمل دپاره کارېږي۔ کله چې د
  ترمونو ترمنځ راځي---لکه په~\LR{$(\num{n} + \num{m})$} کښې---دا د حساب
  د ژبې~\LR{$2$}-ځايه تابع ده، او کله چې د عددونو ترمنځ راځي---
  لکه په~\LR{$n+m$} کښې---دا د جمع عمل دے۔ په دې کښې د~\LR{$\num{n+m}$}
  قضيه هم شامله ده: دا له عدد~\LR{$n+m$} سره سمون لرونکې معياري عددنښه
  ده۔
\end{explain}
\olpseganchor{OLP-0294-B032}{end}

\olpseganchor{OLP-0295-B004}{start}
\olfileid{inc}{req}{cmp}
\olsection{ترکيب په~\texorpdfstring{\LR{$\Th{Q}$}}{Q} کښې تمثيلېدونکے دے}
\olpseganchor{OLP-0295-B004}{end}

\olpseganchor{OLP-0295-B005}{start}
فرض کړئ~\LR{$h$} داسې تعريف شوې ده:
\[
h(x_0,\dots,x_{l-1}) = f(g_0(x_0,\dots,x_{l-1}), \dots,
g_{k-1}(x_0,\dots,x_{l-1})).
\]
چېرته چې موږ لا وړاندې داسې فارمولونه \LR{$!A_f, !A_{g_0}, \dots,
!A_{g_{k-1}}$} موندلي دي چې په ترتيب سره تابعې~\LR{$f$}، او~\LR{$g_0$}،
\dots،~\LR{$g_{k-1}$} تمثيلوي۔ بايد داسې فارمول~\LR{$!A_h$} ومومو چې
\LR{$h$} تمثيل کړي۔
\olpseganchor{OLP-0295-B005}{end}

\olpseganchor{OLP-0295-B006}{start}
له يوې ساده قضيې پيل وکړو، چې ټولې تابعې~\LR{$1$}-ځايه وي؛ يعنې
\LR{$h(x) = f(g(x))$} وګورئ۔ که~\LR{$!A_f(y, z)$}، \LR{$f$} تمثيلوي، او
\LR{$!A_g(x, y)$}، \LR{$g$} تمثيلوي، داسې فارمول~\LR{$!A_h(x, z)$} پکار
ده چې~\LR{$h$} تمثيل کړي۔ پام وکړئ چې~\LR{$h(x) = z$} هله او يوازې هله چې
داسې~\LR{$y$} شته چې \LR{$z = f(y)$} او~\LR{$y = g(x)$} دواړه وي۔ (که
\LR{$h(x) = z$} وي، نو~\LR{$g(x)$} داسې~\LR{$y$} دے؛ که داسې~\LR{$y$} موجود وي، نو
څرنګه چې \LR{$y = g(x)$} او~\LR{$z = f(y)$} دي، \LR{$z = f(g(x))$}۔) له دې څخه
ښکاري چې \LR{$\lexists[y][(!A_g(x, y) \land !A_f(y,
  z))]$} د~\LR{$!A_h(x, z)$} دپاره ښه نوماند دے۔ يوازې بايد دا تصديق کړو
چې~\LR{$\Th{Q}$} اړوند فارمولونه ثابتوي۔
\olpseganchor{OLP-0295-B006}{end}

\olpseganchor{OLP-0295-B007}{start}
\begin{prop}
\ollabel{prop:rep1}
که~\LR{$h(n) = m$} وي، نو \LR{$\Th{Q} \Proves !A_h(\num{n}, \num{m})$}۔
\end{prop}
\olpseganchor{OLP-0295-B007}{end}

\olpseganchor{OLP-0295-B008}{start}
\begin{proof}
فرض کړئ~\LR{$h(n) = m$}، يعنې~\LR{$f(g(n)) = m$}۔ \LR{$k = g(n)$} ونيسئ۔ بيا
\begin{align*}
  \Th{Q} & \Proves !A_g(\num{n}, \num{k})
  \intertext{\RL{ځکه~\LR{$!A_g$}، \LR{$g$} تمثيلوي، او}}
  \Th{Q} & \Proves !A_f(\num{k}, \num{m})
  \intertext{\RL{ځکه~\LR{$!A_f$}، \LR{$f$} تمثيلوي۔ نو}}
  \Th{Q} & \Proves !A_g(\num{n}, \num{k}) \land !A_f(\num{k}, \num{m})
  \intertext{\RL{او په پايله کښې دا هم}}
  \Th{Q} & \Proves \lexists[y][(!A_g(\num{n}, y) \land !A_f(y, \num{m}))],
\end{align*}
يعنې \LR{$\Th{Q} \Proves !A_h(\num{n}, \num{m})$}۔
\end{proof}
\olpseganchor{OLP-0295-B008}{end}

\olpseganchor{OLP-0295-B009}{start}
\begin{prop}
\ollabel{prop:rep2}
که~\LR{$h(n) = m$} وي، نو \LR{$\Th{Q} \Proves \lforall[z][(!A_h(\num{n}, z) \lif
  z = \num{m})]$}۔
\end{prop}
\olpseganchor{OLP-0295-B009}{end}

\olpseganchor{OLP-0295-B010}{start}
\begin{proof}
فرض کړئ~\LR{$h(n) = m$}، يعنې~\LR{$f(g(n)) = m$}۔ \LR{$k = g(n)$} ونيسئ۔ بيا
\begin{align*}
  \Th{Q} & \Proves \lforall[y][(!A_g(\num{n}, y) \lif \eq[y][\num{k}])]
  \intertext{\RL{ځکه~\LR{$!A_g$}، \LR{$g$} تمثيلوي، او}}
  \Th{Q} & \Proves \lforall[z][(!A_f(\num{k}, z) \lif \eq[z][\num{m}])]
  \intertext{\RL{ځکه~\LR{$!A_f$}، \LR{$f$} تمثيلوي۔ يوازې د لږ منطق په کارولو
    ښودلے شو چې دا هم}}
  \Th{Q} & \Proves \lforall[z][(\lexists[y][(!A_g(\num{n}, y) \land
      !A_f(y, z))] \lif \eq[z][\num{m}])].
\end{align*}
يعنې \LR{$\Th{Q} \Proves \lforall[y][(!A_h(\num n, y) \lif \eq[y][\num m])]$}۔
\end{proof}
\olpseganchor{OLP-0295-B010}{end}

\olpseganchor{OLP-0295-B011}{start}
همدا مفکوره په هغې زياتې پېچلې قضيې کښې هم کار کوي چې~\LR{$f$} او
\LR{$g_i$} له~\LR{$1$} څخه زيات ځایونه لري۔
\olpseganchor{OLP-0295-B011}{end}

\olpseganchor{OLP-0295-B012}{start}
\begin{prop}
\ollabel{prop:rep-composition}
که~\LR{$!A_f(y_0, \dots, y_{k-1}, z)$} په~\LR{$\Th{Q}$} کښې
\LR{$f(y_0, \dots, y_{k-1})$} تمثيلوي، او
\LR{$!A_{g_i}(x_0, \dots, x_{l-1}, y)$} په~\LR{$\Th{Q}$} کښې
\LR{$g_i(x_0, \dots, x_{l-1})$} تمثيلوي، نو
\begin{multline*}
  \lexists[y_0\dots][\lexists[y_{k-1}][(!A_{g_0}(x_0,\dots,x_{l-1},y_0) \land
      \dots \land {}]]\\
  !A_{g_{k-1}}(x_0,\dots,x_{l-1},y_{k-1}) \land !A_f(y_0,\dots,y_{k-1},z))
\end{multline*}
دا تمثيلوي:
\[
h(x_0, \dots, x_{l-1}) = f(g_0(x_0, \dots, x_{l-1}), \dots, g_{k-1}(x_0,
\dots, x_{l-1})).
\]
\end{prop}
\olpseganchor{OLP-0295-B012}{end}

\olpseganchor{OLP-0295-B013}{start}
\begin{proof}
تمرین۔
\end{proof}
\olpseganchor{OLP-0295-B013}{end}

\olpseganchor{OLP-0295-B014}{start}
\begin{prob}
د \olref[inc][req][cmp]{prop:rep1} او
\olref[inc][req][cmp]{prop:rep2} ثبوتونه د لارښود په توګه وکاروئ او
د \olref[inc][req][cmp]{prop:rep-composition} ثبوت په تفصيل ترسره
کړئ۔
\end{prob}
\olpseganchor{OLP-0295-B014}{end}

\olpseganchor{OLP-0296-B004}{start}
\olfileid{inc}{req}{min}
\olsection{منظم اقل موندل په~\texorpdfstring{\LR{$\Th{Q}$}}{Q} کښې تمثيلېدونکے دے}
\olpseganchor{OLP-0296-B004}{end}

\olpseganchor{OLP-0296-B005}{start}
بې‌حده پلټنه وګورو۔ فرض کړئ~\LR{$g(x, z)$} منظمه او په~\LR{$\Th{Q}$} کښې
تمثيلېدونکې ده، د بېلګې په توګه د~فارمول~\LR{$!A_g(x, z, y)$} په
وسيله۔ فرض کړئ~\LR{$f$} د~\LR{$f(z) = \umin{x}{[g(x, z) = 0]}$} په ډول تعريف
شوې ده۔ غواړو داسې فارمول~\LR{$!A_f(z, y)$} ومومو چې~\LR{$f$} تمثيل
کړي۔ د~\LR{$f(z)$} ارزښت هغه عدد~\LR{$x$} دے چې (الف)~\LR{$g(x, z) = 0$} پوره
کوي او (ب) تر ټولو وړوکے داسې عدد وي، يعنې د هر~\LR{$w < x$} دپاره
\LR{$g(w, z) \neq 0$} وي۔ نو دا طبيعي انتخاب دے:
\[
!A_f(z,y) \ident !A_g(y, z, \Obj 0) \land \lforall[w][(w < y \lif \lnot
  !A_g(w, z, \Obj 0))].
\]
البته، په عامه قضيه کښې بايد~\LR{$z$} په~\LR{$z_0$}، \dots،~\LR{$z_k$} بدل کړو۔
\olpseganchor{OLP-0296-B005}{end}

\olpseganchor{OLP-0296-B006}{start}
ثبوت به يو ځل بيا د هغو څيزونو په اړه څو لمې وکاروي چې~\LR{$\Th{Q}$}
ئې د ثابتولو دپاره پوره ځواک لري۔
\olpseganchor{OLP-0296-B006}{end}

\olpseganchor{OLP-0296-B007}{start}
\begin{lem}
\ollabel{lem:succ} د هر~ثابت سمبول~\LR{$a$} او هر طبيعي عدد~\LR{$n$} دپاره،
\[
\Th{Q} \Proves \eq[(a' + \num n)][(a + \num n)'].
\]
\end{lem}
\olpseganchor{OLP-0296-B007}{end}

\olpseganchor{OLP-0296-B008}{start}
\begin{proof}
ثبوت، د معمول په شان، پر~\LR{$n$} په استقرا دے۔ په بنسټيزه قضيه کښې
\LR{$n = 0$}؛ بايد وښيو چې~\LR{$\Th{Q}$}،
\LR{$\eq[(a' + \Obj 0)][(a + \Obj 0)']$} ثابتوي۔ خو لرو:
\begin{align}
  \Th{Q} & \Proves \eq[(a' + \Obj 0)][a'] \quad \text{\RL{د بديهي اصل~\LR{$Q_4$} له مخې}}
  \ollabel{step1}\\
  \Th{Q} & \Proves  \eq[(a + \Obj 0)][a] \quad \text{\RL{د بديهي اصل~\LR{$Q_4$} له مخې}}
  \ollabel{step2} \\
  \Th{Q} & \Proves \eq[(a + \Obj 0)'][a'] \quad \text{\RL{د \olref{step2} له مخې}}
  \ollabel{step3} \\
  \Th{Q} & \Proves \eq[(a' + \Obj 0)][(a + \Obj 0)'] \quad
  \text{\RL{د \olref{step1} او \olref{step3} له مخې}}\notag
\end{align}
د استقرا په پړاو کښې فرضولے شو چې ښودلې مو ده:
\LR{$\Th{Q} \Proves \eq[(a' + \num n)][(a + \num n)']$}۔ ځکه چې
\LR{$\num{n+1}$}، \LR{$\num{n}'$} دے، بايد وښيو چې~\LR{$\Th{Q}$}،
\LR{$\eq[(a' + \num{n}')][(a + \num{n}')']$} ثابتوي۔ لرو:
\begin{align}
  \Th{Q} & \Proves \eq[(a' + \num n')][(a' + \num n)'] \quad
  \text{\RL{د بديهي اصل~\LR{$!Q_5$} له مخې}} \ollabel{step5}\\
  \Th{Q} & \Proves \eq[(a' + \num n)'][(a + \num n')'] \quad
  \text{\RL{د استقرايي فرضيې او بديهي اصل~\LR{$!Q_5$} له مخې}} \ollabel{step6}\\
  \Th{Q} & \Proves \eq[(a' + \num n')][(a + \num n')'] \quad
  \text{\RL{د \olref{step5} او \olref{step6} له مخې۔}} \notag
\end{align}
\end{proof}
\olpseganchor{OLP-0296-B008}{end}

\olpseganchor{OLP-0296-B009}{start}
بيا يادول په کار دي چې دا له دې خبرې کمزورې ده چې~\LR{$\Th{Q}$}،
\LR{$\lforall[x][\lforall[y][(x' + y) = (x + y)']]$} ثابتوي۔ که څه هم
دا تړلے فارمول په~\LR{$\Struct{N}$} کښې رښتينې ده،~\LR{$\Th{Q}$} ئې نۀ
ثابتوي۔
\olpseganchor{OLP-0296-B009}{end}

\olpseganchor{OLP-0296-B010}{start}
\begin{lem}
\ollabel{lem:less-zero}
\LR{$\Th{Q} \Proves \lforall[x][\lnot x < \Obj 0]$}.
\end{lem}
\olpseganchor{OLP-0296-B010}{end}

\olpseganchor{OLP-0296-B011}{start}
\begin{proof}
ثبوت په غيرصوري ډول ورکوو (يعنې يوازې دومره نښې ورکوو چې صوري
اشتقاق څنګه جوړېږي)۔
\olpseganchor{OLP-0296-B011}{end}

\olpseganchor{OLP-0296-B012}{start}
د خپل سري~\LR{$a$} دپاره بايد~\LR{$\lnot a < \Obj 0$} ثابت کړو۔ د~\LR{$<$} د
تعريف له مخې، بايد په~\LR{$\Th{Q}$} کښې
\LR{$\lnot \lexists[y][\eq[(y'+a)][\Obj 0]]$} ثابت کړو۔
\LR{$\lexists[y][\eq[(y'+a)][\Obj 0]]$} فرضوو او تناقض ثابتوو۔ فرض کړئ
\LR{$\eq[(b'+a)][\Obj 0]$}۔ د~\LR{$!Q_3$} په کارولو لرو:
\LR{$\eq[a][\Obj 0] \lor \lexists[y][\eq[a][y']]$}۔ قضيې بېلوو۔
\olpseganchor{OLP-0296-B012}{end}

\olpseganchor{OLP-0296-B013}{start}
لومړۍ قضيه:~\LR{$\eq[a][\Obj 0]$} صدق کوي۔ له
\LR{$\eq[(b'+a)][\Obj 0]$} څخه~\LR{$\eq[(b' + \Obj 0)][\Obj 0]$} لرو۔ د
\LR{$\Th{Q}$} د بديهي اصل~\LR{$!Q_4$} له مخې،
\LR{$\eq[(b' + \Obj 0)][b']$} لرو، او له دې امله
\LR{$\eq[b'][\Obj 0]$}۔ خو د بديهي اصل~\LR{$!Q_2$} له مخې
\LR{$\eq/[b'][\Obj 0]$} هم لرو، چې تناقض دے۔
\olpseganchor{OLP-0296-B013}{end}

\olpseganchor{OLP-0296-B014}{start}
دويمه قضيه: د کوم~\LR{$c$} دپاره~\LR{$\eq[a][c']$}۔ خو بيا
\LR{$\eq[(b' + c')][\Obj 0]$} لرو۔ د بديهي اصل~\LR{$!Q_5$} له مخې،
\LR{$\eq[(b' + c)'][\Obj 0]$} لرو، چې يو ځل بيا له بديهي اصل~\LR{$Q_2$} سره
تناقض لري۔
\end{proof}
\olpseganchor{OLP-0296-B014}{end}

\olpseganchor{OLP-0296-B015}{start}
\begin{lem}
\ollabel{lem:less-nsucc}
د هر طبيعي عدد~\LR{$n$} دپاره،
  \[
  \Th{Q} \Proves
  \lforall[x][(x < \num {n+1} \lif (\eq[x][\Obj 0] \lor \dots \lor
    \eq[x][\num n]))].
  \]
\end{lem}
\olpseganchor{OLP-0296-B015}{end}

\olpseganchor{OLP-0296-B016}{start}
\begin{proof}
پر~\LR{$n$} استقرا کاروو۔ لومړے بنسټيزه قضيه وګورو، چې~\LR{$n = 0$} وي۔ په
دې قضيه کښې بايد د خپل سري~\LR{$a$} دپاره
\LR{$a < \num 1 \lif \eq[a][\Obj 0]$} وښيو۔ فرض کړئ~\LR{$a < \num 1$}۔ بيا
د~\LR{$<$} د تعريفوونکي بديهي اصل له مخې،
\LR{$\lexists[y][\eq[(y'+a)][\Obj 0']]$} لرو (ځکه~\LR{$\num 1 \ident
\Obj 0'$})۔
\olpseganchor{OLP-0296-B016}{end}

\olpseganchor{OLP-0296-B017}{start}
فرض کړئ~\LR{$b$} دا ځانګړنه لري، يعنې
\LR{$\eq[(b'+a)][\Obj 0']$} لرو۔ بايد~\LR{$\eq[a][\Obj 0]$} وښيو۔ د بديهي
اصل~\LR{$!Q_3$} له مخې، يا~\LR{$\eq[a][\Obj 0]$} لرو يا داسې~\LR{$c$} شته چې
\LR{$\eq[a][c']$}۔ په لومړۍ قضيه کښې د ښودلو څه نشته۔ نو
\LR{$\eq[a][c']$} فرض کړئ۔ بيا~\LR{$\eq[(b' + c')][\Obj 0']$} لرو۔ د
\LR{$\Th{Q}$} د بديهي اصل~\LR{$!Q_5$} له مخې،
\LR{$\eq[(b'+c)'][\Obj 0']$} لرو۔ د بديهي اصل~\LR{$!Q_1$} له مخې،
\LR{$\eq[(b' + c)][\Obj 0]$} لرو۔ خو د بديهي اصل~\LR{$!Q_8$} له مخې د دې معنا
دا ده چې~\LR{$c < \Obj 0$}، چې له~\olref{lem:less-zero} سره تناقض لري۔
\olpseganchor{OLP-0296-B017}{end}

\olpseganchor{OLP-0296-B018}{start}
اوس د استقرا پړاو ته راځو۔ د~\LR{$n$} قضيه فرضوو او د~\LR{$n+1$} قضيه
ثابتوو۔ نو فرض کړئ~\LR{$a < \num {n+2}$}۔ يو ځل بيا د~\LR{$!Q_3$} په کارولو
دوه قضيې بېلولے شو:~\LR{$\eq[a][\Obj 0]$}، او د کوم~\LR{$b$} دپاره
\LR{$\eq[a][b']$}۔ په لومړۍ قضيه کښې،
\LR{$\eq[a][\Obj 0] \lor \dots \lor \eq[a][\num{n+1}]$} په ښکاره لازمېږي۔
په دويمه قضيه کښې~\LR{$b' < \num {n+2}$} لرو، يعنې
\LR{$b' < \num{n+1}'$}۔ د بديهي اصل~\LR{$!Q_8$} له مخې، د کوم~\LR{$c$} دپاره
\LR{$\eq[(c'+b')][\num{n+1}']$}۔ د بديهي اصل~\LR{$!Q_5$} له مخې،
\LR{$\eq[(c'+b)'][\num{n+1}']$}۔ د بديهي اصل~\LR{$!Q_1$} له مخې،
\LR{$\eq[(c'+b)][\num{n+1}]$}، او نو د بديهي اصل~\LR{$!Q_8$} له مخې
\LR{$b < \num{n+1}$}۔ د استقرايي فرضيې له مخې،
\LR{$\eq[b][\Obj 0] \lor \dots \lor \eq[b][\num{n}]$}۔ له دې څخه د منطق
په وسيله~\LR{$\eq[b'][\Obj 0'] \lor \dots \lor \eq[b'][\num{n}']$} اخلو،
او ځکه چې~\LR{$\eq[a][b']$}، نو
\LR{$\eq[a][\num{1}] \lor \dots \lor \eq[a][\num{n+1}]$}۔
\end{proof}
\olpseganchor{OLP-0296-B018}{end}

\olpseganchor{OLP-0296-B019}{start}
\begin{lem}
  \ollabel{lem:trichotomy} د هر طبيعي عدد~\LR{$m$} دپاره،
  \[
  \Th{Q} \Proves
  \lforall[y][((y < \num{m} \lor \num{m} < y) \lor \eq[y][\num{m}])].
  \]
\end{lem}
\olpseganchor{OLP-0296-B019}{end}

\olpseganchor{OLP-0296-B020}{start}
\begin{proof}
پر~\LR{$m$} په استقرا۔ لومړے د~\LR{$m=0$} قضيه وګورئ۔ د~\LR{$!Q_3$} له مخې
\LR{$\Th{Q} \Proves
\lforall[y][(\eq[y][\Obj 0] \lor \lexists[z][\eq[y][z']])]$}۔
\LR{$a$} خپل سرے ونيسئ۔ بيا يا~\LR{$\eq[a][\Obj 0]$} وي يا د کوم~\LR{$b$} دپاره
\LR{$\eq[a][b']$}۔ په لومړۍ قضيه کښې
\LR{$(a < \Obj 0 \lor \Obj 0 < a) \lor \eq[a][\Obj 0]$} هم لرو۔ خو که
\LR{$\eq[a][b']$} وي، نو د~\LR{$\eq$} د منطق له مخې
\LR{$\eq[(b' + \Obj 0)][(a + \Obj 0)]$}۔ د~\LR{$!Q_4$} له مخې
\LR{$\eq[(a + \Obj 0)][a]$}، نو~\LR{$\eq[(b' + \Obj 0)][a]$} لرو، او له دې
امله~\LR{$\lexists[z][\eq[(z' + \Obj 0)][a]]$}۔ په~\LR{$!Q_8$} کښې د~\LR{$<$} د
تعريف له مخې،~\LR{$\Obj 0 < a$}۔ که~\LR{$\Obj 0 < a$} وي، نو
\LR{$(\Obj 0 < a \lor a < \Obj 0) \lor \eq[a][\Obj 0]$} هم لرو۔
\olpseganchor{OLP-0296-B020}{end}

\olpseganchor{OLP-0296-B021}{start}
اوس فرض کړئ چې لرو:
\begin{align*}
  \Th{Q} & \Proves \lforall[y][((y < \num{m} \lor \num{m} < y) \lor
    \eq[y][\num{m}])]
  \intertext{\RL{او غواړو وښيو:}}
  \Th{Q} & \Proves \lforall[y][((y < \num{m+1} \lor \num{m+1} < y) \lor
    \eq[y][\num{m+1}])]
\end{align*}
\LR{$a$} خپل سرے ونيسئ۔ د~\LR{$!Q_3$} له مخې، يا~\LR{$\eq[a][\Obj 0]$} وي يا د
کوم~\LR{$b$} دپاره~\LR{$\eq[a][b']$}۔ په لومړۍ قضيه کښې، د~\LR{$!Q_4$} له مخې
\LR{$\eq[\num{m}' + a][\num{m+1}]$} لرو، او نو د~\LR{$!Q_8$} له مخې
\LR{$a < \num{m+1}$}۔
\olpseganchor{OLP-0296-B021}{end}

\olpseganchor{OLP-0296-B022}{start}
اوس دويمه قضيه،~\LR{$\eq[a][b']$}، وګورئ۔ د استقرايي فرضيې له مخې،
\LR{$(b < \num{m} \lor \num{m} < b) \lor
    \eq[b][\num{m}]$}۔
\olpseganchor{OLP-0296-B022}{end}

\olpseganchor{OLP-0296-B023}{start}
لومړے منفصل~\LR{$b < \num{m}$} د~\LR{$!Q_8$} له مخې له
\LR{$\lexists[z][\eq[(z' + b)][\num{m}]]$} سره معادل دے۔ فرض کړئ~\LR{$c$}
دا ځانګړنه لري۔ که~\LR{$\eq[(c' + b)][\num{m}]$} وي، نو
\LR{$\eq[(c' + b)'][\num{m}']$} هم وي۔ د~\LR{$!Q_5$} له مخې،
\LR{$\eq[(c' + b)'][(c' + b')]$}۔ له دې امله
\LR{$\eq[(c' + b')][\num{m}']$}۔ پر~\LR{$c'$} د وجودي تعميم په کولو او په دې
پام سره چې~\LR{$\num{m}' \ident \num{m+1}$}،
\LR{$\lexists[u][\eq[(u' + b')][\num{m+1}]]$} اخلو۔ له دې امله، که
\LR{$b < \num{m}$} وي، نو~\LR{$b' < \num{m+1}$}، او نو
\LR{$a < \num{m+1}$}۔
\olpseganchor{OLP-0296-B023}{end}

\olpseganchor{OLP-0296-B024}{start}
اوس فرض کړئ~\LR{$\num{m} < b$}، يعنې
\LR{$\lexists[z][\eq[(z' + \num{m})][b]]$}۔ فرض کړئ~\LR{$c$} داسې~\LR{$z$} دے،
يعنې~\LR{$\eq[(c' + \num{m})][b]$}۔ د منطق له مخې،
\LR{$\eq[(c' + \num{m})'][b']$}۔ د~\LR{$!Q_5$} له مخې،
\LR{$\eq[(c' + \num{m}')][b']$}۔ ځکه چې~\LR{$\eq[a][b']$} او
\LR{$\num{m}' \ident \num{m+1}$}، نو~\LR{$\eq[(c' + \num{m+1})][a]$}۔ د
\LR{$!Q_8$} له مخې،~\LR{$\num{m+1} < a$}۔
\olpseganchor{OLP-0296-B024}{end}

\olpseganchor{OLP-0296-B025}{start}
په پاې کښې،~\LR{$\eq[b][\num{m}]$} فرض کړئ۔ بيا د منطق له مخې
\LR{$\eq[b'][\num{m}']$}، او نو~\LR{$\eq[a][\num{m+1}]$}۔
\olpseganchor{OLP-0296-B025}{end}

\olpseganchor{OLP-0296-B026}{start}
نو د~\LR{$m$} او~\LR{$b$} د قضيې له هر منفصل څخه د~\LR{$m+1$} او~\LR{$a$} اړوند منفصل
تر لاسه کولے شو۔
\end{proof}
\olpseganchor{OLP-0296-B026}{end}

\olpseganchor{OLP-0296-B027}{start}
\begin{prop}
\ollabel{prop:rep-minimization}
که~\LR{$!A_g(x, z, y)$} په~\LR{$\Th{Q}$} کښې~\LR{$g(x, z)$} تمثيلوي، نو
\[
!A_f(z,y) \ident !A_g(y, z, \Obj 0) \land \lforall[w][(w < y \lif \lnot
  !A_g(w, z, \Obj 0))]
\]
\LR{$f(z) = \umin{x}{[g(x, z) = 0]}$} تمثيلوي۔
\end{prop}
\olpseganchor{OLP-0296-B027}{end}

\olpseganchor{OLP-0296-B028}{start}
\begin{proof}
لومړے ښيو چې که~\LR{$f(n) = m$} وي، نو
\LR{$\Th{Q} \Proves !A_f(\num n, \num m)$}؛ يعنې
\begin{align}
  \Th{Q} & \Proves !A_g(\num{m}, \num{n}, \Obj 0) \land \lforall[w][(w <
    \num{m} \lif \lnot !A_g(w, \num{n}, \Obj 0))]. \notag
  \intertext{\RL{ځکه چې~\LR{$!A_g(x, z, y)$}، \LR{$g(x, z)$} تمثيلوي او
    \LR{$g(m, n) = 0$} که~\LR{$f(n) = m$} وي، لرو:}}
\Th{Q} & \Proves !A_g(\num{m}, \num{n}, \Obj 0). \notag
\intertext{\RL{که~\LR{$f(n) = m$} وي، نو د هر~\LR{$k < m$} دپاره~\LR{$g(k, n) \neq 0$}۔ نو}}
\Th{Q} & \Proves \lnot !A_g(\num{k}, \num{n}, \Obj 0). \notag
\intertext{\RL{تر لاسه کوو:}}
\Th{Q} & \Proves \lforall[w][(w < \num{m} \lif \lnot
  !A_g(w, \num{n}, \Obj 0))]. \ollabel{rep-less}
\end{align}
که~\LR{$m = 0$} وي، دا د~\olref{lem:less-zero} له مخې، او که نۀ وي د
\olref{lem:less-nsucc} له مخې راځي۔
\olpseganchor{OLP-0296-B028}{end}

\olpseganchor{OLP-0296-B029}{start}
اوس وښيو چې که~\LR{$f(n) = m$} وي، نو
\LR{$\Th{Q} \Proves
\lforall[y][(!A_f(\num{n}, y) \lif \eq[y][\num{m}])]$}۔ يو ځل بيا
استدلال په غيرصوري ډول لنډوو او صوري کول ئې لوستونکي ته پرېږدو۔
\olpseganchor{OLP-0296-B029}{end}

\olpseganchor{OLP-0296-B030}{start}
فرض کړئ~\LR{$!A_f(\num{n}, b)$}۔ له دې څخه (الف)
\LR{$!A_g(b, \num{n}, \Obj 0)$} او (ب)
\LR{$\lforall[w][(w < b \lif \lnot !A_g(w, \num{n}, \Obj
  0))]$} اخلو۔ د~\olref{lem:trichotomy} له مخې،
\LR{$(b < \num{m} \lor \num{m} < b) \lor \eq[b][\num{m}]$}۔ وبه ښيو چې
\LR{$b < \num{m}$} او~\LR{$\num{m} < b$} دواړه تناقض ته رسېږي۔
\olpseganchor{OLP-0296-B030}{end}

\olpseganchor{OLP-0296-B031}{start}
که~\LR{$\num{m} < b$} وي، نو له~(ب) څخه
\LR{$\lnot !A_g(\num{m}, \num{n}, \Obj 0)$}۔ خو~\LR{$m = f(n)$}، نو
\LR{$g(m, n) = 0$}؛ او ځکه چې~\LR{$!A_g$}، \LR{$g$} تمثيلوي،
\LR{$\Th{Q} \Proves !A_g(\num{m}, \num{n}, \Obj 0)$}۔ نو تناقض لرو۔
\olpseganchor{OLP-0296-B031}{end}

\olpseganchor{OLP-0296-B032}{start}
اوس~\LR{$b < \num{m}$} فرض کړئ۔ بيا ځکه چې د~\olref{rep-less} له مخې
\LR{$\Th{Q} \Proves \lforall[w][(w <
  \num{m} \lif \lnot !A_g(w, \num{n}, \Obj 0))]$}، نو
\LR{$\lnot !A_g(b, \num{n}, \Obj 0)$} اخلو۔ دا يو ځل بيا له~(الف) سره
تناقض لري۔
\end{proof}
\olpseganchor{OLP-0296-B032}{end}

\olpseganchor{OLP-0297-B004}{start}
\olfileid{inc}{req}{crq}
\olsection{محاسبه کېدونکې تابعې په~\texorpdfstring{\LR{$\Th{Q}$}}{Q} کښې تمثيلېدونکې دي}
\olpseganchor{OLP-0297-B004}{end}

\olpseganchor{OLP-0297-B005}{start}
\begin{thm}
هره محاسبه کېدونکې تابعه په~\LR{$\Th{Q}$} کښې تمثيلېدونکې ده۔
\end{thm}
\olpseganchor{OLP-0297-B005}{end}

\olpseganchor{OLP-0297-B006}{start}
\begin{proof}
د کره والي دپاره، او د چرچ--ټيورينګ د دعوې په کارولو سره، ووايو چې
يوه تابعه هله او يوازې هله محاسبه کېدونکې ده چې عمومي بازګشتي وي۔
عمومي بازګشتي تابعې هغه دي چې له صفر تابع~\LR{$\Zero$}، تالي
تابع~\LR{$\Succ$}، او پروجکشن تابع~\LR{$\Proj{n}{i}$} څخه د ترکيب، بنسټيز
بازګښت، او منظم اقل موندلو په کارولو تعريفېدے شي۔ د
\olref[pri]{lem:prim-rec} له مخې، هره تابعه~\LR{$h$} چې له~\LR{$f$} او~\LR{$g$}
څخه تعريفېدے شي، له~\LR{$f$}، \LR{$g$}، او~\LR{$\Zero$}، \LR{$\Succ$}،
\LR{$\Proj{n}{i}$}، \LR{$\Add$}، \LR{$\Mult$}، \LR{$\Char{=}$} څخه هم د ترکيب او منظم
اقل موندلو په کارولو تعريفېدے شي۔ له دې امله، يوه تابعه هله او يوازې
هله عمومي بازګشتي ده چې له~\LR{$\Zero$}، \LR{$\Succ$}، \LR{$\Proj{n}{i}$}،
\LR{$\Add$}، \LR{$\Mult$}، \LR{$\Char{=}$} څخه د ترکيب او منظم اقل موندلو په
کارولو تعريفېدے شي۔
\olpseganchor{OLP-0297-B006}{end}

\olpseganchor{OLP-0297-B007}{start}
سربېره پر دې، ښودلې مو ده چې د پام وړ بنسټيزې تابعې په~\LR{$\Th{Q}$}
کښې تمثيلېدونکې دي
(\cref{inc:req:bre:prop:rep-zero,inc:req:bre:prop:rep-succ,inc:req:bre:prop:rep-proj,inc:req:bre:prop:rep-id,inc:req:bre:prop:rep-add,inc:req:bre:prop:rep-mult})،
او هره هغه تابعه چې له تمثيلېدونکو تابعو څخه په ترکيب يا منظم اقل
موندلو تعريف شوې وي
(\olref[cmp]{prop:rep-composition}،
\olref[min]{prop:rep-minimization}) هم تمثيلېدونکې ده۔ نو هره عمومي
بازګشتي تابعه په~\LR{$\Th{Q}$} کښې تمثيلېدونکې ده۔
\end{proof}
\olpseganchor{OLP-0297-B007}{end}

\olpseganchor{OLP-0297-B008}{start}
\begin{explain}
موږ ښودلې ده چې د محاسبه کېدونکو تابعو سټ د هغو تابعو د سټ په توګه
ځانګړے کېدے شي چې په~\LR{$\Th{Q}$} کښې تمثيلېدونکې دي۔ په حقيقت کښې
ثبوت تر دې عام دے۔ د تمثيلېدنې له تعريف څخه په اسانۍ ليدل کېږي چې
هره هغه تيوري چې~\LR{$\Th{Q}$} غځوي (يا د عددنښو او د تمثيل د دواړو شرطونو ساتونکے تفسير د~\LR{$\Th{Q}$} لري)،
محاسبه کېدونکې تابعې تمثيلولے شي۔ برعکس، په سازګاره تيورۍ کښې چې د بېلو عددنښو نابرابري ثابتوي او اشتقاق نظام ئې د اشتقاق اغېزمن لټون پرېږدي، هره تمثيلېدونکې تابعه محاسبه کېدونکې ده۔ د ورکړلو ارګومېنټونو دپاره د ثبوتونو او د پايلې د ممکنو عددونو جوړې په اغېزمنه توګه وشمېرئ۔ د هرې جوړې دپاره وګورئ چې هماغه ثبوت پر ورکړلو عددنښو د تمثيلوونکي فارمول او د وړانديز شوې پايلې پر عددنښه پاے ته رسېږي که نه۔ د تمثيل لومړے شرط د سمې پايلې د ثبوت شتون تضمينوي۔ که د بلې پايلې دپاره هم ثبوت وموندل شي، د تمثيل دويم شرط به د دواړو عددنښو برابري ثابته کړي۔ د بېلو عددنښو د نابرابرۍ اثبات او سازګاري دا امکان ردوي۔ په ناسازګار کلاسيک نظام کښې له تناقض څخه د هرې تابعې د تمثيل دواړه شرطونه ثابتېږي۔ نو د بېلګې په توګه، د
محاسبه کېدونکو تابعو سټ د هغو تابعو د سټ په توګه ځانګړے کېدے شي چې
د پيانو په حساب، يا ان د زرمیلو--فرېنکل په سټ تيورۍ کښې تمثيلېدونکې
دي۔ لکه ګوډل چې يادونه کړې، دا يو څه هېښوونکې ده۔ وبه وينو چې د
ثبوت‌وړتيا په برخه کښې پوښتنې د ټاکلې تيورۍ پر وړاندې ډېرې حساسې
دي؛ په ټوليز ډول هر څومره چې بديهي اصول پياوړي وي، هغومره ډېر څه
ثابتولے شئ۔ خو د بديهي تيوريو په يوه پراخه ډله کښې تمثيلېدونکې
تابعې کټ مټ محاسبه کېدونکې تابعې دي؛ تر هغه چې تيورۍ
اغېزمنې بديهي کېدونکې او سازګارې وي، پياوړې تيورۍ زياتې تابعې نۀ تمثيلوي۔
\end{explain}
\olpseganchor{OLP-0297-B008}{end}

\paragraph{د منجمدې سرچينې بشپړ بديل متن}
\olpseganchor{OLP-0647-B004}{start}
\olfileid{inc}{req}{cre-olp0647}
\olsection{د \texorpdfstring{\LR{$C$}}{C} تابعې په \texorpdfstring{\LR{$\Th{Q}$}}{Q} کښې تمثيلېدونکې دي}
\olpseganchor{OLP-0647-B004}{end}

\olpseganchor{OLP-0647-B005}{start}
بايد وښيو چې په \LR{$C$} کښې هره تابعه په \LR{$\Th{Q}$} کښې تمثيلېدونکې ده۔ په پاے کښې دا ښودل پکار دي چې د \LR{$C$} هرې \LR{$k$}-ځايه تابعې \LR{$f(x_0,\dots,x_{k-1})$} ته داسې فارمول \LR{$!A_f(x_0,\dots,x_{k-1},y)$} څنګه ګومارلے شو چې هغه تمثيل کړي۔
\olpseganchor{OLP-0647-B005}{end}

\olpseganchor{OLP-0647-B006}{start}
\begin{lem}
که طبيعي عددونه \LR{$n$} او \LR{$m$} راکړل شوي وي او \LR{$n \neq m$} وي، نو \LR{$Q \Proves \num n \neq \num m$}۔
\end{lem}
\olpseganchor{OLP-0647-B006}{end}

\olpseganchor{OLP-0647-B007}{start}
\begin{proof}
پر \LR{$n$} استقرا وکاروئ، څو وښيئ چې د هر \LR{$m$} دپاره، که \LR{$n \neq m$} وي، نو \LR{$Q \Proves \eq/[\num n][\num m]$}۔
\olpseganchor{OLP-0647-B007}{end}

\olpseganchor{OLP-0647-B008}{start}
په بنسټيز حالت کښې \LR{$n = 0$} دے۔ که \LR{$m$} له \LR{$0$} سره برابر نۀ وي، نو د کوم طبيعي عدد \LR{$k$} دپاره \LR{$m = k + 1$} دے۔ يو اصل لرو چې وايي \LR{$\lforall[x][\eq/[0][x']]$}۔ د کمیت ټاکونکي د يو اصل په مرسته، د \LR{$x$} پر ځاے \LR{$\num k$} په اېښودلو سره، \LR{$\eq/[0][\num k']$} ترلاسه کولے شو۔ خو \LR{$\num k'$} هماغه \LR{$\num m$} دے۔
\olpseganchor{OLP-0647-B008}{end}

\olpseganchor{OLP-0647-B009}{start}
په استقرايي ګام کښې فرض کولے شو چې دعوه د \LR{$n$} دپاره صحيح ده، او بيا \LR{$n+1$} په پام کښې نيسو۔ \LR{$m$} دې هر طبيعي عدد وي۔ دوه امکانونه شته: يا \LR{$m = 0$} دے، يا د کوم \LR{$k$} دپاره \LR{$m = k+1$} دے۔ لومړے حالت د پورته طريقې په شان حل کېږي۔ په دويم حالت کښې فرض کړئ چې \LR{$n+1 \neq k+1$} دے۔ نو \LR{$n \neq k$} دے۔ د \LR{$n$} دپاره د استقرايي فرض له مخې \LR{$\Th{Q} \Proves \eq/[\num n][\num k]$} لرو۔ يو اصل وايي \LR{$\lforall[x][\lforall[y][\eq[x'][y'] \lif \eq[x][y]]]$}۔ د کمیت ټاکونکي د يو اصل په کارولو \LR{$\eq[\num n'][\num k'] \lif \eq[\num n][\num k]$} لرو۔ د قضيو د منطق په مرسته په \LR{$\Th{Q}$} کښې \LR{$\eq/[\num n][\num k] \lif \eq/[\num n'][\num k']$} ترلاسه کولے شو۔ د مودوس پوننس په کارولو \LR{$\eq/[\num n'][\num k']$} ترلاسه کېږي، چې همدا مو مطلوبه پايله ده، ځکه \LR{$\num k'$} هماغه \LR{$\num m$} دے۔
\end{proof}
\olpseganchor{OLP-0647-B009}{end}

\olpseganchor{OLP-0647-B010}{start}
پام وکړئ چې دا مرستندويه قضيه ډېره لويه دعوه نۀ کوي: په اصل کښې وايي چې \LR{$\Th{Q}$} ثابتولے شي چې بېلې عددنښې بېل څيزونه ښيي۔ مثلاً \LR{$\Th{Q}$} ثابتوي چې \LR{$0'' \neq 0'''$}۔ خو د دې عمومي ثبوت يو څه احتياط غواړي۔ دا هم په پام کښې وساتئ چې که څه هم استقرا کاروو، دا استقرا د \LR{$\Th{Q}$} \emph{څخه بهر} ده۔
\olpseganchor{OLP-0647-B010}{end}

\olpseganchor{OLP-0647-B011}{start}
صفر، تالي، جمع، ضرب، د عينيت مشخصه تابعه او پروجکشنونه تمثيلولے شو۔ په هر حالت کښې مناسب تمثيلوونکے فارمول ډېر نېغ په نېغه دے؛ مثلاً صفر په دې فارمول تمثيلېږي:
\[
y = 0,
\]
تالي په دې فارمول تمثيلېږي:
\[
x_0' = y,
\]
او جمع په دې فارمول تمثيلېږي:
\[
x_0 + x_1 = y.
\]
کار دا دے چې وښيو \LR{$\Th{Q}$} اړوند بيانونه ثابتولے شي؛ مثلاً دا ويل چې جمع په پورته فارمول تمثيلېږي، دا ښودل غواړي چې د طبيعي عددونو د هرې جوړې \LR{$m$} او \LR{$n$} دپاره \LR{$\Th{Q}$} دا ثابتوي:
\[
\eq[\num n + \num m][\num {n+m}]
\]
او
\[
\lforall[y][(\eq[(\num n + \num m)][y] \lif \eq[y][\num {n+m}])].
\]
\olpseganchor{OLP-0647-B011}{end}

\olpseganchor{OLP-0647-B012}{start}
د ترکيب په اړه څۀ ويلے شو؟ فرض کړئ چې \LR{$h$} داسې تعريف شوے دے:
\[
h(x_0,\dots,x_{l-1}) = f(g_0(x_0,\dots,x_{l-1}), \dots,
g_{k-1}(x_0,\dots,x_{l-1})).
\]
او هغه فارمولونه \LR{$!A_f, !A_{g_0}, \dots, !A_{g_{k-1}}$} مو له مخکې موندلي دي چې په ترتيب سره تابعې \LR{$f,g_0,\dots,g_{k-1}$} تمثيلوي۔ بيا د \LR{$h$} د تمثيل دپاره \LR{$!A_h$} تعريفولے شو، په دې توګه چې \LR{$!A_h(x_0,\dots,x_{l-1},y)$} دا فارمول وټاکو:
\begin{multline*}
  \lexists[z_0\dots][\lexists[z_{k-1}][(!A_{g_0}(x_0,\dots,x_{l-1},z_0) \land
  \dots \land !A_{g_{k-1}}(x_0,\dots,x_{l-1},z_{k-1}) \land]]\\
  !A_f(z_0,\dots,z_{k-1},y)).
\end{multline*}
\olpseganchor{OLP-0647-B012}{end}

\olpseganchor{OLP-0647-B013}{start}
په پاے کښې بې‌حده پلټنه په پام کښې نيسو۔ فرض کړئ چې \LR{$g(x,\vec z)$} منظمه ده او په \LR{$\Th{Q}$} کښې، مثلاً په فارمول \LR{$!A_g(x,\vec z,y)$}، تمثيلېدونکې ده۔ \LR{$f$} دې په \LR{$f(\vec z) = \mu x \; g(x,\vec z)$} تعريف شي۔ غواړو داسې فارمول \LR{$!A_f(\vec z,y)$} ومومو چې \LR{$f$} تمثيل کړي۔ يو طبيعي انتخاب دا دے:
\[
!A_f(\vec z,y) \ident !A_g(y,\vec z,0) \land \lforall[w][(w < y
\lif \lnot !A_g(w,\vec z,0))].
\]
\olpseganchor{OLP-0647-B013}{end}

\olpseganchor{OLP-0647-B014}{start}
د ځينو نورو مرستندويه قضيو په کارولو ښودل کېدے شي چې دا انتخاب کار کوي؛ مثلاً:
\olpseganchor{OLP-0647-B014}{end}

\olpseganchor{OLP-0647-B015}{start}
\begin{lem}
د هر متغير \LR{$x$} او هر طبيعي عدد \LR{$n$} دپاره \LR{$\Th{Q}$} دا ثابتوي: \LR{$\eq[(x'
  + \num n)][(x + \num n)']$}۔
\end{lem}
\olpseganchor{OLP-0647-B015}{end}

\olpseganchor{OLP-0647-B016}{start}
بيا هم دا يادونه پکار ده چې دا تر هغې دعوې کمزورې خبره ده چې \LR{$\Th{Q}$}، \LR{$\lforall[x][\lforall[y][(x' + y) = (x +y)']]$} ثابتوي؛ د اثبات دغه دعوه ناسمه ده، نۀ په معياري طبيعي عددونو کښې خپله حسابي مساوات۔
\olpseganchor{OLP-0647-B016}{end}

\olpseganchor{OLP-0647-B017}{start}
\begin{proof}
ثبوت د معمول په شان پر \LR{$n$} په استقرا دے۔ په بنسټيز حالت کښې \LR{$n = 0$} دے؛ بايد وښيو چې \LR{$\Th{Q}$}، \LR{$\eq[(x'+0)][(x+0)']$} ثابتوي۔ خو دا لرو:
\begin{eqnarray*}
& & \eq[(x' + 0)][x'] \quad \text{\RL{له اصل ۴ څخه}} \\
& & \eq[(x + 0)][x] \quad \text{\RL{له اصل ۴ څخه}} \\
& & \eq[(x+0)'][x'] \quad \text{\RL{له کرښې ۲ څخه}} \\
& & \eq[(x' + 0)][(x + 0)'] \quad \text{\RL{له کرښو ۱ او ۳ څخه}}
\end{eqnarray*}
په استقرايي ګام کښې فرض کولے شو چې \LR{$\eq[(x' + \num n)][(x + \num n)']$} مو په \LR{$\Th{Q}$} کښې استنتاج کړے دے۔ ځکه چې \LR{$\num{n+1}$} هماغه \LR{$\num n'$} دے، بايد وښيو چې \LR{$\Th{Q}$}، \LR{$\eq[(x' + \num n')][(x + \num n')']$} ثابتوي۔ د مساواتو دا لړۍ په \LR{$\Th{Q}$} کښې استنتاج کېدے شي:
\begin{eqnarray*}
(x' + \num n') & \eq & (x' + \num n)' \quad \text{\RL{اصل ۵}} \\
& \eq & (x + \num n')' \quad \text{\RL{له استقرايي فرض او اصل ۵ څخه}}
\end{eqnarray*}
\end{proof}
\olpseganchor{OLP-0647-B017}{end}

\olpseganchor{OLP-0647-B018}{start}
\begin{lem}
\begin{enumerate}
\item \LR{$\Th{Q}$}، \LR{$\lnot (x < \num 0)$} ثابتوي۔
\item د هر طبيعي عدد \LR{$n$} دپاره \LR{$\Th{Q}$} دا ثابتوي:
\[
x < \num {n+1} \lif (\eq[x][\num 0] \lor \dots \lor \eq[x][\num n]).
\]
\end{enumerate}
\end{lem}
\olpseganchor{OLP-0647-B018}{end}

\olpseganchor{OLP-0647-B019}{start}
\begin{proof}
لومړۍ برخه او د دويمې برخې يوه برخه په غير رسمي ډول کوو، يعنې يوازې دا اشارې ورکوو چې رسمي اشتقاق څنګه جوړ شي۔
\olpseganchor{OLP-0647-B019}{end}

\olpseganchor{OLP-0647-B020}{start}
د لومړۍ برخې دپاره، د \LR{$<$} د تعريف له مخې، بايد په \LR{$\Th{Q}$} کښې \LR{$\lnot \lexists[y][\eq[(y'+x)][0]]$} ثابت کړو؛ دا د لومړۍ درجې منطق د اصولو او قواعدو په کارولو له \LR{$\lforall[y][\eq/[(y' + x)][0]]$} سره همارزښته ده۔ مفکوره دا ده: فرض کړئ چې \LR{$\eq[(y'+x)][0]$} دے۔ که \LR{$x$}، 0 وي، \LR{$\eq[(y' + 0)][0]$} لرو۔ خو د \LR{$\Th{Q}$} د اصل ۴ له مخې \LR{$\eq[(y' + 0)][y']$} لرو، او د اصل ۲ له مخې \LR{$\eq/[y'][0]$} لرو، چې تناقض دے۔ نو \LR{$\lforall[y][\eq/[(y' +x)][0]]$}۔ که \LR{$x$}، \LR{$0$} نۀ وي، د اصل ۳ له مخې داسې \LR{$z$} شته چې \LR{$\eq[x][z']$} دے۔ خو بيا \LR{$\eq[(y' + z')][0]$} لرو۔ د اصل~۵ له مخې \LR{$\eq[(y' + z)'][0]$} لرو، چې بيا هم د اصل~۲ سره په تناقض کښې دے۔
\olpseganchor{OLP-0647-B020}{end}

\olpseganchor{OLP-0647-B021}{start}
د دويمې برخې دپاره پر \LR{$n$} استقرا وکاروئ۔ بنسټيز حالت په پام کښې نيسو، چې \LR{$n =0$} دے۔ په دې حالت کښې بايد \LR{$x < \num 1 \lif \eq[x][\num 0]$} وښيو۔ فرض کړئ چې \LR{$x < \num 1$} دے۔ نو د \LR{$<$} د تعريفي اصل له مخې \LR{$\lexists[y][\eq[(y'+x)][0']]$} لرو۔ فرض کړئ چې \LR{$y$} دغه خاصيت لري؛ نو \LR{$\eq[y'+x][0']$} لرو۔
\olpseganchor{OLP-0647-B021}{end}

\olpseganchor{OLP-0647-B022}{start}
بايد \LR{$\eq[x][0]$} وښيو۔ د اصل ۳ له مخې، که \LR{$x$}، 0 نۀ وي، د کوم \LR{$z$} دپاره له \LR{$z'$} سره برابر دے۔ بيا \LR{$\eq[(y' + z')][0']$} لرو۔ د \LR{$\Th{Q}$} د اصل~۵ له مخې \LR{$\eq[(y'+z)'][0']$} لرو۔ د اصل ۱ له مخې \LR{$\eq[(y' + z)][0]$} لرو۔ خو دا د تعريف له مخې \LR{$z < 0$} معنا لري، چې له لومړۍ برخې سره په تناقض کښې دے۔ په استقرايي ګام کښې هم يا غړے صفر دے او يا د بل غړي تالي دے۔ د تالي په حالت کښې، د جمع د تالي اصل او د تالي يوپر‌يووالے د کوچنيوالي فرض د مخکيني غړي او مخکينۍ عددنښې تر منځ کوچنيوالي ته راکموي۔ استقرايي فرض وايي چې مخکينے غړے له کومې وړې عددنښې سره برابر دے، نو خپله غړے د هغې عددنښې له تالي سره برابر دے۔
\end{proof}
\olpseganchor{OLP-0647-B022}{end}

\olpseganchor{OLP-0647-B023}{start}
\begin{explain}
ښودلې مو ده چې د محاسبه کېدونکو تابعو سټ د هغو تابعو د سټ په توګه ځانګړے کېدے شي چې په \LR{$\Th{Q}$} کښې تمثيلېدونکې دي۔ په حقيقت کښې ثبوت تر دې عام دے۔ د تمثيلېدنې له تعريف څخه په اسانۍ ليدل کېږي چې هره تيوري چې \LR{$\Th{Q}$} غځوي (يا د عددنښو او د تمثيل د دواړو شرطونو ساتونکے تفسير د~\LR{$\Th{Q}$} لري)، محاسبه کېدونکې تابعې تمثيلولے شي؛ برعکس، هره تابعه چې په سازګاره تيورۍ کښې تمثيلېږي محاسبه کېدونکې ده، که د اشتقاق نظام د ثبوتونو اغېزمنه شمېرنه ورکړي او تيوري د بېلو عددنښو نابرابري ثابتوي۔ د ورکړلو ارګومېنټونو دپاره د ثبوتونو او د پايلې د ممکنو عددونو جوړې په اغېزمنه توګه وشمېرئ۔ د هرې جوړې دپاره وګورئ چې هماغه ثبوت پر ورکړلو عددنښو د تمثيلوونکي فارمول او د وړانديز شوې پايلې پر عددنښه پاے ته رسېږي که نه۔ د تمثيل لومړے شرط د سمې پايلې د ثبوت شتون تضمينوي۔ که د بلې پايلې دپاره هم ثبوت وموندل شي، د تمثيل دويم شرط به د دواړو عددنښو برابري ثابته کړي۔ د بېلو عددنښو د نابرابرۍ اثبات او سازګاري دا امکان ردوي۔ په ناسازګار کلاسيک نظام کښې له تناقض څخه د هرې تابعې د تمثيل دواړه شرطونه ثابتېږي۔ نو مثلاً د محاسبه کېدونکو تابعو سټ د هغو تابعو د سټ په توګه ځانګړے کېدے شي چې د پيانو په حساب، يا ان د زرمیلو--فرېنکل په سټ تيورۍ کښې تمثيلېږي۔ لکه ګوډل چې يادونه کړې، دا يو څه هېښوونکې ده۔ وبه وينو چې د ثبوت‌وړتيا پوښتنې د ټاکلې تيورۍ پر وړاندې ډېرې حساسې دي؛ په ټوليز ډول، هر څومره چې اصول پياوړي وي، هغومره ډېر څه ثابتولے شئ۔ خو د اصولي تيوريو په يوه پراخه سازګاره ډله کښې تمثيلېدونکې تابعې کټ مټ محاسبه کېدونکې تابعې دي۔ 
\end{explain}
\olpseganchor{OLP-0647-B023}{end}

\olpseganchor{OLP-0298-B004}{start}
\olfileid{inc}{req}{rel}
\olpseganchor{OLP-0298-B004}{end}

\olpseganchor{OLP-0298-B005}{start}
\olsection{د اړيکو تمثيلول}
\olpseganchor{OLP-0298-B005}{end}

\olpseganchor{OLP-0298-B006}{start}
دا بيان کړو چې يوه \emph{اړيکه} تمثيلېدونکې وي څه معنا لري۔
\olpseganchor{OLP-0298-B006}{end}

\olpseganchor{OLP-0298-B007}{start}
\begin{defn}
\ollabel{defn:representing-relations} پر طبيعي عددونو اړيکه
\LR{$R(x_0,\dots,x_k)$} هله \emph{په~\LR{$\Th{Q}$} کښې تمثيلېدونکې} ده چې
داسې فارمول~\LR{$!A_R(x_0,\dots,x_k)$} وي چې هر کله
\LR{$R(n_0,\dots,n_k)$} رښتيا وي، \LR{$\Th{Q}$}،
\LR{$!A_R(\num{n_0},\dots,\num{n_k})$} ثابت کړي، او هر کله
\LR{$R(n_0,\dots,n_k)$} ناسم وي، \LR{$\Th{Q}$}،
\LR{$\lnot !A_R(\num{n_0}, \dots, \num{n_k})$} ثابت کړي۔
\end{defn}
\olpseganchor{OLP-0298-B007}{end}

\olpseganchor{OLP-0298-B008}{start}
\begin{thm}
\ollabel{thm:representing-rels} يوه اړيکه په~\LR{$\Th{Q}$} کښې هله او
يوازې هله تمثيلېدونکې ده چې محاسبه کېدونکې وي۔
\end{thm}
\olpseganchor{OLP-0298-B008}{end}

\olpseganchor{OLP-0298-B009}{start}
\begin{proof}
د مخکښ لوري دپاره فرض کړئ اړيکه~\LR{$R(x_0,\dots,x_k)$} په
فارمول~\LR{$!A_R(x_0,\dots,x_k)$} تمثيلېږي۔ د~\LR{$R$} د محاسبې الګوريتم دا
دے: پر ننوتونو~\LR{$n_0$}، \dots،~\LR{$n_k$} په هممهاله ډول د
\LR{$!A_R(\num{n_0}, \dots, \num{n_k})$} د ثبوت او د
\LR{$\lnot !A_R(\num{n_0}, \dots, \num{n_k})$} د ثبوت پلټنه وکړئ۔ زموږ
د فرض له مخې پلټنه خامخا يو له دغو دواړو مومي؛ که لومړے وي، «هو»
راپور کړئ، او که نۀ، «نه» راپور کړئ۔
\olpseganchor{OLP-0298-B009}{end}

\olpseganchor{OLP-0298-B010}{start}
په بل لوري کښې فرض کړئ~\LR{$R(x_0, \dots, x_k)$} محاسبه کېدونکې ده۔ د
تعريف له مخې دا په دې معنا ده چې تابعه~\LR{$\Char{R}(x_0, \dots, x_k)$}
محاسبه کېدونکې ده۔ د~\olref[int]{thm:representable-iff-comp} له
مخې، \LR{$\Char{R}$} په کوم فارمول، مثلاً
\LR{$!A_{\Char{R}}(x_0, \dots, x_k, y)$}، تمثيلېږي۔
\LR{$!A_R(x_0, \dots, x_k)$} هغه فارمول ونيسئ چې
\LR{$!A_{\Char{R}}(x_0, \dots, x_k, \num{1})$} وي۔ بيا د هر~\LR{$n_0$}،
\dots،~\LR{$n_k$} دپاره، که~\LR{$R(n_0, \dots, n_k)$} رښتيا وي، نو
\LR{$\Char{R}(n_0, \dots, n_k) = 1$}؛ په دې حالت کښې~\LR{$\Th{Q}$}،
\LR{$!A_{\Char{R}}(\num{n_0}, \dots, \num{n_k}, \num{1})$} ثابتوي، او
نو~\LR{$\Th{Q}$}، \LR{$!A_R(\num{n_0}, \dots, \num{n_k})$} ثابتوي۔ له بلې
خوا، که~\LR{$R(n_0, \dots, n_k)$} ناسم وي، نو
\LR{$\Char{R}(n_0, \dots, n_k) = 0$}۔ دا په دې معنا ده چې~\LR{$\Th{Q}$}
لاندې ثابتوي:
\[
\lforall[y][(!A_{\Char{R}}(\num{n_0}, \dots, \num{n_k}, y) \lif y =
  \num{0})].
\]
څرنګه چې~\LR{$\Th{Q}$}، \LR{$\eq/[\num{0}][\num{1}]$} ثابتوي، \LR{$\Th{Q}$}،
\LR{$\lnot !A_{\Char{R}}(\num{n_0}, \dots, \num{n_k}, \num{1})$}
ثابتوي، او نو~\LR{$\lnot !A_R(\num{n_0}, \dots, \num{n_k})$} هم
ثابتوي۔
\end{proof}
\olpseganchor{OLP-0298-B010}{end}

\olpseganchor{OLP-0298-B011}{start}
\begin{prob}
وښيئ چې که~\LR{$R$} په~\LR{$\Th{Q}$} کښې تمثيلېدونکې وي، نو~\LR{$\Char{R}$} هم
همداسې ده۔
\end{prob}
\olpseganchor{OLP-0298-B011}{end}

\olpseganchor{OLP-0299-B004}{start}
\olfileid{inc}{req}{und}
\olsection{ناپرېکړتيا}
\olpseganchor{OLP-0299-B004}{end}

\olpseganchor{OLP-0299-B005}{start}
تيوري~\LR{$\Th{T}$} هغه وخت \emph{ناپرېکړه‌وړ} بولو چې داسې محاسبوي
کړنلاره نۀ وي چې له متناهي شمېر ګامونو وروسته د~\LR{$\Th{T}$} د ژبې د
هرې جملې~\LR{$!A$} دپاره په باوري ډول دې پوښتنې ته سم ځواب ورکړي:
«آيا~\LR{$\Th{T}$}، \LR{$!A$} ثابتوي؟» نو~\LR{$\Th{Q}$} به هله او يوازې هله د
پرېکړې وړ وه چې داسې محاسبوي کړنلاره وای چې د حساب د ژبې د ورکړل شوې
جملې~\LR{$!A$} دپاره پرېکړه وکړي چې~\LR{$\Th{Q} \Proves !A$} ده که نۀ۔ دا
پوښتنه داسې لا کره کولے شو: آيا اړيکه~\LR{$\Prov[\Th{Q}](y)$}، چې پر~\LR{$y$}
هله صدق کوي چې~\LR{$y$} په~\LR{$\Th{Q}$} کښې د ثابتېدونکې جملې ګوډل شمېره
وي، بازګشتي ده؟ ځواب «نه» دے۔
\olpseganchor{OLP-0299-B005}{end}

\olpseganchor{OLP-0299-B006}{start}
\begin{thm}
\LR{$\Th{Q}$} ناپرېکړه‌وړ ده؛ يعنې اړيکه
\[
\Prov[\Th{Q}](y) \defiff \fn{Sent}(y) \land
\lexists[x][\Prf[\Th{Q}](x, y)]
\]
بازګشتي نۀ ده۔
\end{thm}
\olpseganchor{OLP-0299-B006}{end}

\olpseganchor{OLP-0299-B007}{start}
\begin{proof}
فرض کړئ بازګشتي وای۔ بيا د درېدنې مسئله داسې حلولے شو: د ورکړو
\LR{$e$} او~\LR{$n$} دپاره پوهېږو چې~\LR{$\cfind{e}(n) \fdefined$} هله او يوازې
هله چې داسې~\LR{$s$} وي چې~\LR{$T(e, n, s)$}، چې~\LR{$T$} د
\olref[cmp][rec][nft]{thm:kleene-nf} څخه د کلېني پريديکات دے۔ ځکه
\LR{$T$} بنسټيزه بازګشتي ده، په~\LR{$\Th{Q}$} کښې د کوم فارمول~\LR{$!B_T$} په
وسيله تمثيلېږي؛ يعنې~\LR{$\Th{Q} \Proves
!B_T(\num{e}, \num{n}, \num{s})$} هله او يوازې هله چې~\LR{$T(e,n,s)$}۔
که~\LR{$\Th{Q} \Proves !B_T(\num{e}, \num{n}, \num{s})$}، نو
\LR{$ \Th{Q} \Proves \lexists[y][!B_T(\num{e}, \num{n}, y)]$} هم۔ که
داسې~\LR{$s$} نۀ وي، نو د هر~\LR{$s$} دپاره
\LR{$\Th{Q} \Proves \lnot !B_T(\num{e}, \num{n}, \num{s})$}۔ خو~\LR{$\Th{Q}$}،
\LR{$\omega$}-سازګاره ده؛ يعنې که د هر~\LR{$n \in \Nat$} دپاره
\LR{$\Th{Q} \Proves \lnot !A(\num{n})$}، نو
\LR{$\Th{Q} \Proves/ \lexists[y][!A(y)]$}۔ دا ځکه پوهېږو چې د~\LR{$\Th{Q}$}
بديهي اصول په معياري مدل~\LR{$\Struct{N}$} کښې رښتوني دي۔ نو
\LR{$\Th{Q} \Proves/ \lexists[y][!B_T(\num{e}, \num{n}, y)]$}۔ په بله
وينا، \LR{$\Th{Q} \Proves \lexists[y][!B_T(\num{e}, \num{n}, y)]$}
هله او يوازې هله چې داسې~\LR{$s$} وي چې~\LR{$T(e,n,s)$}؛ يعنې هله او يوازې
هله چې~\LR{$\cfind{e}(n) \fdefined$}۔ له~\LR{$e$} او~\LR{$n$} څخه
\LR{$\Gn{\lexists[y][!B_T(\num{e}, \num{n}, y)]}$} محاسبه کولے شو؛
\LR{$g(e,n)$} هغه بنسټيزه بازګشتي تابعه ونيسئ چې دا کار کوي۔ نو
\[
h(e, n) =
\begin{cases}
1 & \text{\RL{که~\LR{$\Prov[\Th{Q}](g(e, n))$} وي}}\\
0 & \text{\RL{که نۀ وي}}.
\end{cases}
\]
دا به ښودلې وای چې که~\LR{$\Prov[\Th{Q}]$} بازګشتي وي، \LR{$h$} هم بازګشتي
ده۔ خو د~\olref[cmp][rec][hlt]{thm:halting-problem} له مخې~\LR{$h$}
بازګشتي نۀ ده، نو~\LR{$\Prov[\Th{Q}]$} هم بازګشتي نۀ شي کېدے۔
\end{proof}
\olpseganchor{OLP-0299-B007}{end}

\olpseganchor{OLP-0299-B008}{start}
\begin{cor}
د لومړۍ درجې منطق ناپرېکړه‌وړ دے۔
\end{cor}
\olpseganchor{OLP-0299-B008}{end}

\olpseganchor{OLP-0299-B009}{start}
\begin{proof}
که د لومړۍ درجې منطق د پرېکړې وړ وای، په~\LR{$\Th{Q}$} کښې ثبوت‌وړتيا
هم د پرېکړې وړ وه، ځکه~\LR{$\Th{Q} \Proves !A$} هله او يوازې هله چې
\LR{$\Proves !T \lif !A$}، چېرته چې~\LR{$!T$} د~\LR{$\Th{Q}$} د بديهي اصولو
عطف دے۔
\end{proof}
\olpseganchor{OLP-0299-B009}{end}

\olpseganchor{OLP-0300-B004}{start}
\olfileid{inc}{inp}{s1c}
\olsection{د~\texorpdfstring{\texorpdfstring{\LR{$\Sigma_1$}}{Σ₁}}{سيګما-۱} بشپړتيا}
\olpseganchor{OLP-0300-B004}{end}

\olpseganchor{OLP-0300-B005}{start}
د~\LR{$\Th{Q}$} او د هغې د سازګارو، بديهي کېدونکو غځونو د ناتکميلۍ
سره سره، ليدلي مو دي چې~\LR{$\Th{Q}$} د عددنښو په اړه ډېر بنسټيز
حقيقتونه ثابتوي۔ په حقيقت کښې دا خبره په پام وړ ډول غځېدے شي۔ د دې
د پوهېدو دپاره چې په~\LR{$\Th{Q}$} کښې کوم څيزونه ثابتېدے شي، د
\LR{$\Delta_0$}، \LR{$\Sigma_1$}، او~\LR{$\Pi_1$} فارمولونه مفهومونه معرفي
کوو۔ په نا کره ډول،~\LR{$\Sigma_1$} فارمول د
\LR{$\lexists[x][!B(x)]$} په بڼه وي، چې~\LR{$!B$} يوازې د قضيو د نښلوونکو او
محدودو کمیت ايښوونکو په کارولو جوړ شوے وي۔ وبه ښيو چې که~\LR{$!A$} داسې
\LR{$\Sigma_1$} تړلے فارمول وي چې په~\LR{$\Struct{N}$} کښې رښتينې وي، نو
\LR{$\Th{Q} \Proves !A$} (\olref{thm:sigma1-completeness})۔
\olpseganchor{OLP-0300-B005}{end}

\olpseganchor{OLP-0300-B006}{start}
\begin{defn}
\ollabel{defn:bd-quant}
يوه \emph{محدوده وجودي فارمول} د
\LR{$\lexists[x][(x < t \land !A(x))]$} په بڼه وي، چې~\LR{$t$} يو ترم دے او تړل کېدونکے متغير پکښې نۀ راځي، او
موږ ئې په دوديز ډول~\LR{$\bexists{x < t}{!A(x)}$} ليکو۔
%
يوه \emph{محدوده کلي فارمول} د
\LR{$\lforall[x][(x < t \lif !A(x))]$} په بڼه وي، چې~\LR{$t$} يو ترم دے او تړل کېدونکے متغير پکښې نۀ راځي، او
موږ ئې په دوديز ډول~\LR{$\bforall{x < t}{!A(x)}$} ليکو۔
\end{defn}
\olpseganchor{OLP-0300-B006}{end}

\olpseganchor{OLP-0300-B007}{start}
\begin{defn}
\ollabel{defn:delta0-sigma1-pi1-frm}
يوه فارمول~\LR{$!B$} هله~\LR{$\Delta_0$} ده چې له اټومي
فارمولونه څخه يوازې د قضيو د نښلوونکو او محدودې کمیت ايښودنې په
کارولو جوړه شوې وي۔
%
يوه فارمول~\LR{$!A$} هله~\LR{$\Sigma_1$} ده چې
\LR{$!A \ident \lexists[x][!B(x)]$} وي او~\LR{$!B$}، \LR{$\Delta_0$} وي۔
%
يوه فارمول~\LR{$!A$} هله~\LR{$\Pi_1$} ده چې
\LR{$!A \ident \lforall[x][!B(x)]$} وي او~\LR{$!B$}، \LR{$\Delta_0$} وي۔
\end{defn}
\olpseganchor{OLP-0300-B007}{end}

\olpseganchor{OLP-0300-B008}{start}
\begin{lem}
\ollabel{lem:q-proves-clterm-id} فرض کړئ~\LR{$t$} تړلے ترم دے چې
\LR{$\Value{t}{N} = n$}۔ بيا \LR{$\Th{Q} \Proves \eq[t][\num n]$}۔
\end{lem}
\olpseganchor{OLP-0300-B008}{end}

\olpseganchor{OLP-0300-B009}{start}
\begin{proof}
دا د~\LR{$t$} پر پېچلتيا په استقرا ثابتوو۔ په بنسټيزه قضيه کښې
\LR{$\Value{\Obj 0}{N} = 0$}، او
\LR{$\Th{Q} \Proves \eq[\Obj 0][\num 0]$}، ځکه
\LR{$\num 0 \ident \Obj 0$}۔
%
د استقرايي قضيې دپاره، فرض کړئ~\LR{$t_1$} او~\LR{$t_2$} داسې ترمونه دي چې
\LR{$\Value{t_1}{N} = n_1$}، \LR{$\Value{t_2}{N} = n_2$}،
\LR{$\Th{Q} \Proves \eq[t_1][\num n_1]$}، او
\LR{$\Th{Q} \Proves \eq[t_2][\num n_2]$}۔
\olpseganchor{OLP-0300-B009}{end}

\olpseganchor{OLP-0300-B010}{start}
بيا~\LR{$\Value{(t_1')}{N} = n_1 + 1$}، او د عينيت د لومړي-رتبه قواعدو
له مخې، چې پر استقرايي فرضيه او
فارمول~\LR{$\eq[\num{n_1}'][\num{n_1}']$} پلي شوي، لرو چې
\LR{$\Th{Q} \Proves \eq[t_1'][{\num n_1}']$}؛ نو د عددنښو د تعريف له
مخې \LR{$\Th{Q} \Proves \eq[t_1'][\num{n_1 + 1}]$}۔
\olpseganchor{OLP-0300-B010}{end}

\olpseganchor{OLP-0300-B011}{start}
د جمعو دپاره لرو:
\[
      \Value{(t_1 + t_2)}{N}
    = \Value{t_1}{N} + \Value{t_2}{N}
    = n_1 + n_2.
\]
د استقرايي فرضيې او د عينيت د قواعدو له مخې،
\LR{$\Th{Q} \Proves \eq[t_1 + t_2][\num{n_1} + t_2]$}، او بيا د عينيت
د قواعدو په دويم پلي کولو
\LR{$\Th{Q} \Proves \eq[t_1 + t_2][\num{n_1} + \num{n_2}]$}۔ د
\olref[inc][req][bre]{lem:q-proves-add} له مخې،
\LR{$\Th{Q} \Proves \eq[\num{n_1} + \num{n_2}][\num{n_1 + n_2}]$}،
نو \LR{$\Th{Q} \Proves \eq[t_1 + t_2][\num{n_1 + n_2}]$}۔
\olpseganchor{OLP-0300-B011}{end}

\olpseganchor{OLP-0300-B012}{start}
همداسې استدلال د~\LR{$\times$} دپاره هم کار کوي، چې
\olref[inc][req][bre]{lem:q-proves-mult} پکې کاروو۔
%
ځکه چې دا د حساب ټول تړلي ترمونه رانغاړي، نو د هر داسې تړلي
ترم~\LR{$t$} دپاره چې~\LR{$\Value{t}{N} = n$} وي،
\LR{$\Th{Q} \Proves \eq[t][\num n]$}۔
\end{proof}
\olpseganchor{OLP-0300-B012}{end}

\olpseganchor{OLP-0300-B013}{start}
\begin{prob}
په \olref{lem:q-proves-clterm-id} کښې د~\LR{$\times$} اړوند برخه په
تفصيل ثابته کړئ۔
\end{prob}
\olpseganchor{OLP-0300-B013}{end}

\olpseganchor{OLP-0300-B014}{start}
\begin{lem}
\ollabel{lem:atomic-completeness}
فرض کړئ~\LR{$t_1$} او~\LR{$t_2$} تړلي ترمونه دي۔ بيا
\begin{enumerate}
\item که~\LR{$\Value{t_1}{N} = \Value{t_2}{N}$} وي،
    نو \LR{$\Th{Q} \Proves \eq[t_1][t_2]$}۔
\item که~\LR{$\Value{t_1}{N} \neq \Value{t_2}{N}$} وي،
    نو \LR{$\Th{Q} \Proves \eq/[t_1][t_2]$}۔
\item که~\LR{$\Value{t_1}{N} < \Value{t_2}{N}$} وي،
    نو \LR{$\Th{Q} \Proves t_1 < t_2$}۔
\item که~\LR{$\Value{t_2}{N} \leq \Value{t_1}{N}$} وي،
    نو \LR{$\Th{Q} \Proves \lnot(t_1 < t_2)$}۔
\end{enumerate}
\end{lem}
\olpseganchor{OLP-0300-B014}{end}

\olpseganchor{OLP-0300-B015}{start}
\begin{proof}
د ورکړل شوو ترمونو~\LR{$t_1$} او~\LR{$t_2$} دپاره
\LR{$n = \Value{t_1}{N}$} او~\LR{$m = \Value{t_2}{N}$} ټاکو۔
\olpseganchor{OLP-0300-B015}{end}

\olpseganchor{OLP-0300-B016}{start}
فرض کړئ~\LR{$!A \ident t_1 = t_2$}۔ د
\olref{lem:q-proves-clterm-id} له مخې،
\LR{$\Th{Q} \Proves \eq[t_1][\num n]$} او
\LR{$\Th{Q} \Proves \eq[t_2][\num m]$}۔ که~\LR{$n = m$} وي، نو
\LR{$\Th{Q} \Proves \eq[\num n][\num m]$}، او له دې امله د عينيت د
تعدي له مخې~\LR{$\Th{Q} \Proves \eq[t_1][t_2]$}۔ که~\LR{$n \neq m$} وي، نو
\LR{$\Th{Q} \Proves \eq/[\num n][\num m]$}، او يو ځل بيا د عينيت د
تعدي له مخې~\LR{$\Th{Q} \Proves \eq/[t_1][t_2]$}۔
\olpseganchor{OLP-0300-B016}{end}

\olpseganchor{OLP-0300-B017}{start}
اوس~\LR{$!A \ident t_1 < t_2$} ونيسئ۔ په دواړو قضيو کښې پر بديهي
اصل~\LR{$!Q_8$} تکيه کوو، چې وايي د ټولو~\LR{$x,y$} دپاره
\LR{$x < y \liff \lexists[z][\eq[z' + x][y]]$}۔
\olpseganchor{OLP-0300-B017}{end}

\olpseganchor{OLP-0300-B018}{start}
فرض کړئ~\LR{$\Sat{N}{t_1 < t_2}$}۔ بيا داسې~\LR{$k \in \Nat$} شته چې
\LR{$n + k + 1 = m$}۔ د~\olref{lem:q-proves-clterm-id} له مخې،
\LR{$\Th{Q} \Proves \eq[t_1][\num n]$} او
\LR{$\Th{Q} \Proves \eq[t_2][\num m]$}، او د دې لمې د لومړۍ برخې له
مخې، \LR{$\Th{Q} \Proves \eq[{\num k}' + \num n][\num m]$}۔ د عينيت د
تعدي له مخې، \LR{$\Th{Q} \Proves \eq[{\num k}' + t_1][t_2]$}، نو
\LR{$\Th{Q} \Proves \lexists[z][\eq[z' + t_1][t_2]]$}۔ د~\LR{$!Q_8$} د
ښي-څخه-کيڼ لوري له مخې،~\LR{$\Th{Q} \Proves t_1 < t_2$}۔
\olpseganchor{OLP-0300-B018}{end}

\olpseganchor{OLP-0300-B019}{start}
د دې پر ځاے فرض کړئ~\LR{$\Sat/{N}{t_1 < t_2}$}، يعنې~\LR{$m \leq n$}۔
%
په~\LR{$\Th{Q}$} کښې کار کوو او~\LR{$t_1 < t_2$} فرضوو۔ د~\LR{$!Q_8$} د
کيڼ-څخه-ښي لوري له مخې، داسې~\LR{$z$} شته چې
\LR{$\eq[z' + t_1][t_2]$}۔ ځکه چې
\LR{$\Th{Q} \Proves \eq[t_1][\num n]$} او
\LR{$\Th{Q} \Proves \eq[t_2][\num m]$}، نو
\LR{$\eq[z' + \num n][\num m]$}۔
%
پر~\LR{$m$} د~\LR{$!Q_5$} په کارولو د يوې بهرنۍ استقرا له مخې،
\LR{$\eq[z' + \num{n - m}][\Obj 0]$}۔
که~\LR{$m = n$} وي، نو~\LR{$\eq[z'][\Obj 0]$}، چې د~\LR{$!Q_2$} له لارې تناقض
ورکوي۔ که~\LR{$m < n$} وي، نو د~\LR{$!Q_5$} په بيا کارولو
\LR{$\eq[(z' + \num{n - m - 1})'][\Obj 0]$}، چې د~\LR{$!Q_2$} له لارې
تناقض ورکوي۔ نو~\LR{$\Th{Q} \Proves \lnot(t_1 < t_2)$}۔ په دغه بهرنۍ استقرا کښې د ښي عددنښې هر تالي د جمع د تالي له اصل سره پر کيڼ اړخ هم ښکاري؛ د تالي يوپر‌يووالے دواړه تالي نښې لرې کوي۔
\end{proof}
\olpseganchor{OLP-0300-B019}{end}

\olpseganchor{OLP-0300-B020}{start}
\begin{lem}
\ollabel{lem:bounded-quant-equiv}
فرض کړئ~\LR{$!A$} کومه فارمول ده،~\LR{$t$} تړلے ترم دے، او
\LR{$k=\Value{t}{N}$}۔ بيا
\begin{enumerate}
\item \LR{$\Th{Q} \Proves \bforall{x<t}{!A(x)}$} هله او يوازې هله چې
\LR{$\Th{Q} \Proves !A(\num 0) \land \dots \land !A(\num{k-1})$}۔
\item \LR{$\Th{Q} \Proves \bexists{x<t}{!A(x)}$} هله او يوازې هله چې
\LR{$\Th{Q} \Proves !A(\num 0) \lor \dots \lor !A(\num{k-1})$}۔
\end{enumerate}
\end{lem}
\olpseganchor{OLP-0300-B020}{end}

\olpseganchor{OLP-0300-B021}{start}
\begin{proof}
    د محدود کلي کمیت ايښوونکي قضيه ثابتوو۔
    که~\LR{$\Value{t}{N} = 0$} وي، نو د معادلۍ کيڼ اړخ په~\LR{$\Th{Q}$}
    کښې ثابتېدونکے دے، ځکه د
    \olref[inc][req][min]{lem:less-zero} له مخې هېڅ~\LR{$x<\num 0$}
    نشته۔ همداسې، تش عطف په ساده ډول~\LR{$\ltrue$} ګڼلے شو، چې په
    \LR{$\Th{Q}$} کښې هم ثابتېدونکے دے۔
    %
    نو فرضوو چې د کوم طبيعي عدد~\LR{$k$} دپاره
    \LR{$\Value{t}{N} = k+1$}۔ د
    \olref{lem:q-proves-clterm-id} له مخې فرضولے شو چې د
    \LR{$\bforall{x<\num{k+1}}{!A(x)}$} په بڼه پر فارمول کار کوو۔
\olpseganchor{OLP-0300-B021}{end}

\olpseganchor{OLP-0300-B022}{start}
    فرض کړئ~\LR{$\Th{Q} \Proves \bforall{x<\num{k+1}}{!A(x)}$}، او
    \LR{$n \leq k$} ونيسئ۔ ځکه چې د~\olref{lem:atomic-completeness} له
    مخې~\LR{$\Th{Q} \Proves \num n < \num{k+1}$}، نو د منطق له مخې
    \LR{$\Th{Q} \Proves !A(\num n)$}۔ د هر~\LR{$n \leq k$} دپاره د دې حقيقت
    په~\LR{$k+1$} ځله پلي کولو، لکه څنګه چې غوښتل شوي وو، اخلو:
    \LR{$\Th{Q} \Proves !A(\num 0) \land \dots \land !A(\num k)$}۔
\olpseganchor{OLP-0300-B022}{end}

\olpseganchor{OLP-0300-B023}{start}
    د بل لوري دپاره فرض کړئ چې
    \LR{$\Th{Q} \Proves !A(\num 0) \land \dots \land !A(\num k)$}۔ په
    \LR{$\Th{Q}$} کښې کار کوو؛ فرض کړئ~\LR{$x < \num{k+1}$}۔ د
    \olref[inc][req][min]{lem:less-nsucc} له مخې لرو:
    \LR{$x = \num 0 \lor \dots \lor x = \num k$}، نو د منطق له مخې
    \LR{$!A(x)$}؛ او له دې څخه کلي ادعا
    \LR{$\bforall{x<\num{k+1}}{!A(x)}$} راځي۔
\olpseganchor{OLP-0300-B023}{end}

\olpseganchor{OLP-0300-B024}{start}
    د محدودې وجودي کمیت ايښودنې لرونکو فارمولونه معادلي هم له همدې دوو حقايقو راځي۔ د حد لاندې هره عددنښه د کوچنيوالي شرط پوره کوي، او هر غړے چې له حد څخه کوچنے وي له دغو عددنښو څخه له يوې سره برابر دے۔ په لومړي لوري کښې د وجودي شاهد له دغې برابرۍ څخه د عددنښو منفصل اخلو، او په دويم لوري کښې له هر منفصل غړي څخه وجودي جمله اخلو۔ د صفر په حد کښې وجودي جمله ناشونې ده او تش منفصل هم دروغ دے۔
\end{proof}
\olpseganchor{OLP-0300-B024}{end}

\olpseganchor{OLP-0300-B025}{start}
\begin{prob}
په \olref{lem:bounded-quant-equiv} کښې د وجودي قضيې تفصيلي ثبوت
ورکړئ۔
\end{prob}
\olpseganchor{OLP-0300-B025}{end}

\olpseganchor{OLP-0300-B026}{start}
\begin{lem}
\ollabel{lem:delta0-completeness}
که~\LR{$!A$} داسې~\LR{$\Delta_0$} تړلے فارمول وي چې په~\LR{$\Struct{N}$} کښې
رښتينې وي، نو~\LR{$\Th{Q} \Proves !A$}۔ که دروغ وي، نو تيوري ئې نفي ثابتوي۔
\end{lem}
\olpseganchor{OLP-0300-B026}{end}

\olpseganchor{OLP-0300-B027}{start}
\begin{proof}
دواړه دعوې د~فارمول پر جوړښت په يوځايي استقرا ثابتوو۔ د محدودې کمیت ټاکنې په حالت کښې د ماتريس هره تړلې عددنښه‌يي بېلګه د استقرا تر فرض لاندې راځي۔
%
بنسټيزه قضيه~\olref{lem:atomic-completeness} ورکوي، نو استقرايي
پړاو ته ځو۔ د اسانۍ دپاره د نفي قضيه د هغې~فارمول د جوړښت له
مخې پر فرعي قضيو وېشو چې نفي ورباندې پلي شوې ده۔
\olpseganchor{OLP-0300-B027}{end}

\olpseganchor{OLP-0300-B028}{start}
\begin{enumerate}
\item فرض کړئ~\LR{$(!A \land !B)$} په~\LR{$\Struct{N}$} کښې رښتينې ده، نو
\LR{$!A$} او~\LR{$!B$} په~\LR{$\Struct{N}$} کښې رښتينې دي۔ د استقرايي فرضيې له
مخې، \LR{$\Th{Q} \Proves !A$} او~\LR{$\Th{Q} \Proves !B$}، نو د منطق له مخې
\LR{$\Th{Q} \Proves (!A \land !B)$}۔
%
\item فرض کړئ~\LR{$\lnot (!A \land !B)$} په~\LR{$\Struct{N}$} کښې رښتينې ده،
نو يا~\LR{$\lnot !A$} يا~\LR{$\lnot !B$} په~\LR{$\Struct{N}$} کښې رښتينې ده۔ د
عموميت له زيان پرته، لومړۍ فرض کړئ۔ د استقرايي فرضيې له مخې
\LR{$\Th{Q} \Proves \lnot !A$}، او له دې امله د منطق له مخې
\LR{$\Th{Q} \Proves \lnot (!A \land !B)$}۔
%
\item فرض کړئ~\LR{$(!A \lor !B)$} په~\LR{$\Struct{N}$} کښې رښتينې ده، نو يا
\LR{$!A$} په~\LR{$\Struct{N}$} کښې يا~\LR{$!B$} په~\LR{$\Struct{N}$} کښې رښتينې ده۔ د عموميت له زيان
پرته، لومړۍ فرض کړئ۔ د استقرايي فرضيې له مخې
\LR{$\Th{Q} \Proves !A$}، او له دې امله د منطق له مخې
\LR{$\Th{Q} \Proves (!A \lor !B)$}۔
%
\item فرض کړئ~\LR{$\lnot(!A \lor !B)$} په~\LR{$\Struct{N}$} کښې رښتينې ده،
نو~\LR{$\lnot !A$} او~\LR{$\lnot !B$} په~\LR{$\Struct{N}$} کښې رښتينې دي۔ د
استقرايي فرضيې له مخې، \LR{$\Th{Q} \Proves \lnot !A$} او
\LR{$\Th{Q} \Proves \lnot !B$}۔ په پايله کښې، د منطق له مخې
\LR{$\Th{Q} \Proves \lnot(!A \lor !B)$}۔
%
\item فرض کړئ~\LR{$\bforall{x<t}{!A(x)}$} په~\LR{$\Struct{N}$} کښې رښتينې
ده، چې~\LR{$t$} تړلے ترم او~\LR{$k=\Value{t}{N}$} دے۔ د استقرايي فرضيې او
منطق له مخې، که~\LR{$!A(\num n)$} د هر~\LR{$n < \Value{t}{N}$} دپاره په
\LR{$\Struct{N}$} کښې رښتينې وي، نو
\LR{$\Th{Q} \Proves !A(\num 0) \land \dots \land !A(\num{k-1})$}۔ د
\olref{lem:bounded-quant-equiv} له مخې،
\LR{$\Th{Q} \Proves \bforall{x<t}{!A(x)}$}۔
%
\item د محدود وجودي کمیت ايښوونکي قضيه، چې د
\LR{$\bexists{x < t}{!A(x)}$} په بڼه تړلے فارمول لرو، د محدود کلي
کمیت ايښوونکي قضيې ته ورته ده۔ که جمله رښتينې وي، د تړلي حد له ارزښت څخه کوچنے طبيعي شاهد شته۔ د ماتريس د هماغې عددنښې بېلګه د استقرايي فرضيې له مخې ثابتېږي۔ د اټومي لمې له مخې د شاهد د کوچنيوالي شرط هم ثابتېږي؛ بيا وجودي تعميم مطلوبه جمله ورکوي۔
%
\item فرض کړئ~\LR{$\lnot \bforall{x<t}{!A(x)}$} په~\LR{$\Struct{N}$} کښې
رښتينې ده، چې~\LR{$t$} تړلے ترم دے۔ دا تړلے فارمول له
تړلے فارمول~\LR{$\bexists{x<t}{\lnot !A(x)}$} سره معادله ده، او دا معادله په
\LR{$\Th{Q}$} کښې استنتاجېدے شي؛ نو د محدودو وجودي کمیت ايښوونکو
استدلال پلي کولے شو۔ په دې حالت کښې د شاهد دپاره د استقرا له دويمې دعوې کار اخلو، ځکه د ماتريس دغه بېلګه دروغ ده۔
%
\item همداسې، فرض کړئ~\LR{$\lnot \bexists{x<t}{!A(x)}$} په~\LR{$\Struct{N}$}
کښې رښتينې ده، چې~\LR{$t$} تړلے ترم دے۔ دا تړلے فارمول په~\LR{$\Th{Q}$} کښې
له~\LR{$\bforall{x<t}{\lnot!A(x)}$} سره معادله ده، نو د محدودو کلي کمیت
ايښوونکو استدلال پلي کولے شو۔ د حد لاندې د ماتريس ټولې عددنښه‌يي بېلګې دروغ دي، نو د استقرا له دويمې دعوې د هرې بېلګې نفي او د محدود کلي کمیت ايښوونکي له معادلۍ کلي پايله اخلو۔
%
\item په پاې کښې، فرض کړئ~\LR{$\lnot !A$} په~\LR{$\Struct{N}$} کښې رښتينې
ده۔ د شرطي او دوه‌اړخيز نښلوونکي دپاره د استقرا له فرضه د دواړو غړو اثبات يا د نفي اثبات اخلو، د هر يوه د رښتينوالي له مخې۔ د دواړو غړو دغه ثبوتونه د قضيو د منطق له لارې د ټولې جملې اثبات ورکوي که رښتينې وي، او که دروغ وي د نفي اثبات ورکوي۔ رښتيا ثابتېږي او د دروغو نفي هم ثابتېږي۔ اوس د نفي دوه قضيې پاتې دي: يا~\LR{$!A$} اټومي وي، او يا د کومې
\LR{$\Delta_0$} تړلے فارمول~\LR{$!B$} دپاره
\LR{$\lnot !A \ident \lnot\lnot !B$} وي۔ که~\LR{$!A$} اټومي وي، نو د
\olref{lem:atomic-completeness} له مخې
\LR{$\Th{Q} \Proves \lnot !A$}۔ که
\LR{$\lnot !A \ident \lnot\lnot !B$} وي، نو د منطق له مخې دا په
\LR{$\Th{Q}$} کښې په ثابتېدونکي ډول له~\LR{$!B$} سره معادله ده، چې په
\LR{$\Struct{N}$} کښې رښتينې ده، ځکه~\LR{$\lnot !A$} په~\LR{$\Struct{N}$} کښې
رښتينې ده۔ نو د استقرايي فرضيې له مخې
\LR{$\Th{Q} \Proves \lnot !A$}۔
\end{enumerate}
\end{proof}
\olpseganchor{OLP-0300-B028}{end}

\olpseganchor{OLP-0300-B029}{start}
\begin{prob}
په \olref{lem:delta0-completeness} کښې د وجودي قضيې تفصيلي ثبوت
ورکړئ۔
\end{prob}
\olpseganchor{OLP-0300-B029}{end}

\olpseganchor{OLP-0300-B030}{start}
\begin{thm}
\ollabel{thm:sigma1-completeness}
که~\LR{$!A$} داسې~\LR{$\Sigma_1$} تړلے فارمول وي چې په~\LR{$\Struct{N}$} کښې
رښتينې وي، نو~\LR{$\Th{Q} \Proves !A$}۔
\end{thm}
\olpseganchor{OLP-0300-B030}{end}

\olpseganchor{OLP-0300-B031}{start}
\begin{proof}
که~\LR{$\lexists[x][!A(x)]$} داسې~\LR{$\Sigma_1$} تړلے فارمول وي چې په
\LR{$\Struct{N}$} کښې رښتينې وي، نو داسې طبيعي عدد~\LR{$n$} او د متحولونو
د ارزښت ټاکنه~\LR{$s$} شته چې~\LR{$s(x) = n$} او~\LR{$\Sat{N}{!A(x)}[s]$}۔ د
پوره کېدو د اړيکې په اړه له معياري حقيقتونو څخه
\LR{$\Sat{N}{!A(\num n)}$} لازمېږي۔
خو~\LR{$!A(\num n)$} يوه
\LR{$\Delta_0$} فارمول ده، نو د~\olref{lem:delta0-completeness} له
مخې~\LR{$\Th{Q} \Proves !A(\num n)$}؛ او له دې امله د منطق له مخې
\LR{$\Th{Q} \Proves \lexists[x][!A(x)]$} هم لرو۔
\end{proof}
\olpseganchor{OLP-0300-B031}{end}

\olpseganchor{OLP-0301-B004}{start}
\begin{editorial}
  دا څپرکے د محاسبه کېدنې د تيورۍ د څپرکي پر موادو تکيه لري، خو که
  هغه څپرکے نۀ وي لوستل شوے، دا پرېښودل کېدے شي۔ اوس مهال دا د
  جېرېمي اېويګېډ د يادښتونو يوه بنسټيزه اړونه ده، بياکتنه ئې نۀ ده
  شوې، او تمرينونه پکې نشته۔
\end{editorial}
\olpseganchor{OLP-0301-B004}{end}

\olpseganchor{OLP-0301-B005}{start}
\olchapter{inc}{tcp}{تيورۍ او محاسبه کېدنه}
\olpseganchor{OLP-0301-B005}{end}

\olpseganchor{OLP-0302-B004}{start}
\olfileid{inc}{tcp}{int}
\olpseganchor{OLP-0302-B004}{end}

\olpseganchor{OLP-0302-B005}{start}
\olsection{پېژندګلو}
\olpseganchor{OLP-0302-B005}{end}

\olpseganchor{OLP-0302-B006}{start}
\begin{editorial}
دا برخه بايد له سره وليکل شي۔
\end{editorial}
\olpseganchor{OLP-0302-B006}{end}

\olpseganchor{OLP-0302-B007}{start}
موږ لاندې شيان لرو:
\begin{enumerate}
\item دا تعريف چې يوه تابعه په~\LR{$\Th{Q}$} کښې تمثيلېدونکې وي څه معنا
  لري (\olref[req][int]{defn:representable-fn})؛
\item دا تعريف چې يوه اړيکه په~\LR{$\Th{Q}$} کښې تمثيلېدونکې وي څه معنا
  لري (\olref[req][rel]{defn:representing-relations})؛
\item هغه قضيه چې وايي د~\LR{$\Th{Q}$} تمثيلېدونکې تابعې کټ مټ هماغه
  محاسبه کېدونکې تابعې دي
  (\olref[req][int]{thm:representable-iff-comp})؛
\item او هغه قضيه چې وايي د~\LR{$\Th{Q}$} تمثيلېدونکې اړيکې کټ مټ هماغه
  محاسبه کېدونکې اړيکې دي
  \olref[req][rel]{thm:representing-rels}۔
\end{enumerate}
\olpseganchor{OLP-0302-B007}{end}

\olpseganchor{OLP-0302-B008}{start}
يوه \emph{تيوري} د جملو داسې سټ دے چې د اشتقاق له مخې تړلے وي؛
يعنې دا خاصيت لري چې هر کله~\LR{$T$}، \LR{$!A$} ثابتوي، نو~\LR{$!A$} په~\LR{$T$} کښې
وي۔ غوره به دا وي چې تيوري د جملو د يوې ټولګې په توګه وګڼو چې د
هغو جملو ټولې لازمې پايلې هم ورسره وي۔ له دې وروسته به~\LR{$\Th{Q}$} د
هغې \emph{تيورۍ} دپاره کاروو چې په~\olref[req][int]{sec} کښې له
اتو بديهي اصولو څخه د اشتقاق وړ جملو له سټ څخه جوړه ده۔ په ياد ولرئ
چې د~\LR{$\Th{Q}$} فارمولونه د عددونو په توګه کوډولے شو؛ که~\LR{$!A$} داسې
فارمول وي، پرېږدئ~\LR{$\Gn{!A}$} هغه عدد وښيي چې~\LR{$!A$} کوډوي۔ د همدې
کوډونې له مخې اوس پوښتلے شو چې د فارمولونو بېلابېل سټونه محاسبه
کېدونکي دي که نۀ۔
\olpseganchor{OLP-0302-B008}{end}

\olpseganchor{OLP-0303-B004}{start}
\olfileid{inc}{tcp}{qce}
\olpseganchor{OLP-0303-B004}{end}

\olpseganchor{OLP-0303-B005}{start}
\olsection{\texorpdfstring{\LR{$\Th{Q}$}}{Q}، c.e.-بشپړه ده}
\olpseganchor{OLP-0303-B005}{end}


\olpseganchor{OLP-0303-B006}{start}
\begin{thm}
\LR{$\Th{Q}$}، c.e. ده خو د پرېکړې وړ نۀ ده۔ په حقيقت کښې دا يو
بشپړ c.e. سټ دے۔
\end{thm}
\olpseganchor{OLP-0303-B006}{end}

\olpseganchor{OLP-0303-B007}{start}
\begin{proof}
دا ليدل ستونزمن نۀ دي چې~\LR{$\Th{Q}$}، c.e. ده، ځکه دا د هغو
جملو~\LR{$y$} د (کوډونو) سټ دے چې د ثبوت يو کوډ~\LR{$x$} ولري:
\[
\Th{Q} = \Setabs{y}{\lexists[x][\Prf[\Th{Q}](x,y)]}.
\]
خو پوهېږو چې~\LR{$\Prf[\Th{Q}](x,y)$} محاسبه کېدونکې ده (په حقيقت کښې
بنسټيزه بازګشتي ده)، او هر هغه سټ چې په پورتنۍ بڼه وليکل شي، په
محاسبوي ډول د شمېر وړ وي۔
\olpseganchor{OLP-0303-B007}{end}

\olpseganchor{OLP-0303-B008}{start}
دا ويل چې دا يو بشپړ په محاسبوي ډول د شمېر وړ سټ دے، له دې خبرې
سره برابر دي چې~\LR{$K \leq_m \Th{Q}$}، چېرته چې
\LR{$K = \Setabs{x}{!A_x(x) \downarrow}$}۔ نو راځئ وښيو چې~\LR{$K$}،
\LR{$\Th{Q}$} ته راکمېدونکے دے۔ ځکه چې د کلېني پريديکات~\LR{$T(e,x,s)$}
بنسټيزه بازګشتي دے، په~\LR{$\Th{Q}$} کښې، فرض کړئ، د~\LR{$!A_T$} په وسيله
تمثيلېږي۔ بيا د هر~\LR{$x$} دپاره لرو:
\begin{align*}
x \in K & \lif \lexists[s][T(x,x,s)] \\
& \lif \lexists[s][(\Th{Q} \Proves !A_T(\num x,\num x, \num s))]
  \\
& \lif \Th{Q} \Proves \lexists[s][!A_T(\num x, \num x, s)].
\end{align*}
برعکس، که~\LR{$\Th{Q} \Proves \lexists[s][!A_T(\num x, \num x, s)]$}،
نو په حقيقت کښې د کوم طبيعي عدد~\LR{$n$} دپاره بايد فارمول
\LR{$!A_T(\num x, \num x, \num n)$} رښتيا وي۔ اوس که~\LR{$T(x,x,n)$} ناسم
وای، نو ځکه چې~\LR{$!A_T$}، \LR{$T$} تمثيلوي، \LR{$\Th{Q}$} به
\LR{$\lnot !A_T(\num x, \num x, \num n)$} ثابت کړے وای۔ خو بيا به
\LR{$\Th{Q}$} يو ناسم فارمول ثابت کړے وای، چې تناقض دے۔ نو~\LR{$T(x,x,n)$}
بايد رښتيا وي، او له دې څخه~\LR{$!A_x(x) \downarrow$} لازمېږي۔
\olpseganchor{OLP-0303-B008}{end}

\olpseganchor{OLP-0303-B009}{start}
په لنډ ډول، د هر~\LR{$x$} دپاره، \LR{$x$} هله او يوازې هله په~\LR{$K$} کښې دے
چې~\LR{$\Th{Q}$}، \LR{$\lexists[s][!A_T(\num x,\num x,s)]$} ثابت کړي۔ نو هغه
تابعه~\LR{$f$} چې~\LR{$x$} د جملې~\LR{$\lexists[s][!A_T(\num x,
  \num x, s)]$} (يو کوډ) ته وړي، له~\LR{$K$} څخه~\LR{$\Th{Q}$} ته راکمونه ده۔ د فورمول جوړول، د عددنښې اېښودل او د جملې کوډ اخيستل ټول محاسبه کېدونکي ګامونه دي۔
\end{proof}
\olpseganchor{OLP-0303-B009}{end}

\olpseganchor{OLP-0304-B004}{start}
\olfileid{inc}{tcp}{oqn}
\olpseganchor{OLP-0304-B004}{end}

\olpseganchor{OLP-0304-B005}{start}
\olsection{د~\texorpdfstring{\LR{$\Th{Q}$}}{Q}، \texorpdfstring{\LR{$\omega$}}{ω}-سازګارې غځونې ناپرېکړه‌وړې دي}
\olpseganchor{OLP-0304-B005}{end}

\olpseganchor{OLP-0304-B006}{start}
\begin{explain}
دا ثبوت چې~\LR{$\Th{Q}$} په محاسبوي ډول د شمېر وړ بشپړه ده، پر دې حقيقت
تکيه درلوده چې په~\LR{$\Th{Q}$} کښې هره ثابتېدونکې جمله د طبيعي عددونو
په اړه «رښتيا» ده۔ راتلونکے تعريف او قضيه دا پايله پياوړې کوي، او
کټ مټ د «رښتياوالي» هغه اړخونه ټاکي چې په پورتني ثبوت کښې ورته
اړتيا وه۔ خو پر دې قضيه ډېر مه درېږئ، ځکه ژر به ئې لا نوره هم
پياوړې کړو۔ موږ دا زياتره د تاريخي موخې دپاره راوړو: د ګوډل اصلي
مقالې د~\LR{$\omega$}-سازګارۍ مفهوم کارولے و، خو لږ وروسته ئې پايله د
\LR{$\omega$}-سازګارۍ پر ځاے د عادي سازګارۍ په کارولو پياوړې شوه۔
\end{explain}
\olpseganchor{OLP-0304-B006}{end}

\olpseganchor{OLP-0304-B007}{start}
\begin{defn}
\ollabel{thm:oconsis-q}
يوه تيوري~\LR{$\Th{T}$} هله~\LR{$\omega$}-سازګاره ده چې لاندې شرط پوره کړي:
که~\LR{$\lexists[x][!A(x)]$} هره جمله وي او~\LR{$\Th{T}$}،
\LR{$\lnot !A(\num 0)$}، \LR{$\lnot !A(\num 1)$}، \LR{$\lnot !A(\num 2)$}،
\dots ثابتوي، نو~\LR{$\Th{T}$}، \LR{$\lexists[x][!A(x)]$} نۀ ثابتوي۔
\end{defn}
\olpseganchor{OLP-0304-B007}{end}

\olpseganchor{OLP-0304-B008}{start}
\begin{thm}
  پرېږدئ~\LR{$\Th{T}$} هره هغه~\LR{$\omega$}-سازګاره تيوري وي چې~\LR{$\Th{Q}$}
  پکې شامله وي۔ بيا~\LR{$\Th{T}$} د پرېکړې وړ نۀ ده۔
\end{thm}
\olpseganchor{OLP-0304-B008}{end}

\olpseganchor{OLP-0304-B009}{start}
\begin{proof}
که~\LR{$\Th{Q}$} په~\LR{$\Th{T}$} کښې شامله وي، نو~\LR{$\Th{T}$} محاسبه کېدونکې
تابعې او اړيکې تمثيلوي۔ يوازې پخوانی ثبوت لږ بدلول پکار دي۔ لکه
پورته، که~\LR{$x \in K$} وي، نو~\LR{$\Th{T}$}،
\LR{$\lexists[s][!A_T(\num x,\num x, s)]$} ثابتوي۔ برعکس، فرض کړئ
\LR{$\Th{T}$}، \LR{$\lexists[s][!A_T(\num x, \num x, s)]$} ثابتوي۔ بيا بايد
\LR{$x$} په~\LR{$K$} کښې وي: که نۀ وي، د ماشين~\LR{$x$} داسې محاسبه نشته چې پر
ننوت~\LR{$x$} ودرېږي؛ ځکه~\LR{$!A_T$} د کلېني اړيکه~\LR{$T$} تمثيلوي، \LR{$\Th{T}$}،
\LR{$\lnot !A_T(\num x, \num x, \num 0)$}،
\LR{$\lnot !A_T(\num x, \num x, \num 1)$}، \dots ثابتوي، او په دې سره
\LR{$\Th{T}$}، \LR{$\omega$}-ناسازګاره کېږي۔
\end{proof}
\olpseganchor{OLP-0304-B009}{end}

\olpseganchor{OLP-0305-B004}{start}
\olfileid{inc}{tcp}{cqn}
\olpseganchor{OLP-0305-B004}{end}

\olpseganchor{OLP-0305-B005}{start}
\olsection{د~\texorpdfstring{\LR{$\Th{Q}$}}{Q} سازګارې غځونې ناپرېکړه‌وړې دي}
\olpseganchor{OLP-0305-B005}{end}

\olpseganchor{OLP-0305-B006}{start}
\begin{explain}
په ياد ولرئ چې يوه تيوري هله~\emph{سازګاره} ده چې د هېڅ
فارمول~\LR{$!A$} دپاره هم~\LR{$!A$} او هم~\LR{$\lnot !A$} دواړه ثابت نۀ کړي۔ ځکه
له تناقض څخه هر څه لازمېږي، يوه ناسازګاره تيوري ابتذالي وي: هره
جمله پکې ثابتېدونکې وي۔ څرګنده ده چې که يوه تيوري~\LR{$\omega$}-سازګاره
وي، نو سازګاره هم ده۔ خو يوازې سازګاره اوسېدل کمزورے شرط دے (يعنې
داسې تيورۍ شته چې سازګارې دي خو~\LR{$\omega$}-سازګارې نۀ دي)۔ په
\olref[oqn]{thm:oconsis-q} کښې فرضيه ساده سازګارۍ ته کمزورې کولے
شو، او په دې توګه يوه پياوړې قضيه ترلاسه کوو۔
\end{explain}
\olpseganchor{OLP-0305-B006}{end}

\olpseganchor{OLP-0305-B007}{start}
\begin{lem}
هېڅ «نړيواله محاسبه کېدونکې اړيکه» نشته۔ يعنې داسې دوه‌ځاييزه
محاسبه کېدونکې اړيکه~\LR{$R(x,y)$} نشته چې لاندې خاصيت ولري: هر کله
\LR{$S(y)$} يوه يوځاييزه محاسبه کېدونکې اړيکه وي، داسې~\LR{$k$} وي چې د هر
\LR{$y$} دپاره~\LR{$S(y)$} هله او يوازې هله رښتيا وي چې~\LR{$R(k,y)$} رښتيا وي۔
\end{lem}
\olpseganchor{OLP-0305-B007}{end}

\olpseganchor{OLP-0305-B008}{start}
\begin{proof}
فرض کړئ~\LR{$R(x,y)$} يوه نړيواله محاسبه کېدونکې اړيکه وي۔ پرېږدئ
\LR{$S(y)$}، اړيکه~\LR{$\lnot R(y,y)$} وي۔ ځکه~\LR{$S(y)$} محاسبه کېدونکې ده، د
کوم~\LR{$k$} دپاره~\LR{$S(y)$} له~\LR{$R(k,y)$} سره هم‌ارزه ده۔ خو بيا~\LR{$S(k)$} هم
له~\LR{$R(k,k)$} او هم له~\LR{$\lnot R(k,k)$} سره هم‌ارزه ده، چې تناقض دے۔
\end{proof}
\olpseganchor{OLP-0305-B008}{end}


\olpseganchor{OLP-0305-B009}{start}
\begin{thm}
پرېږدئ~\LR{$\Th{T}$} هره هغه سازګاره تيوري وي چې~\LR{$\Th{Q}$} پکې شامله
وي۔ بيا~\LR{$\Th{T}$} د پرېکړې وړ نۀ ده۔
\end{thm}
\olpseganchor{OLP-0305-B009}{end}


\olpseganchor{OLP-0305-B010}{start}
\begin{proof}
فرض کړئ~\LR{$\Th{T}$} د~\LR{$\Th{Q}$} يوه سازګاره او د پرېکړې وړ غځونه وي۔
موږ به د~\LR{$\Th{T}$} په کارولو د يوې نړيوالې محاسبه کېدونکې اړيکې په
تعريفولو سره تناقض ترلاسه کړو۔
\olpseganchor{OLP-0305-B010}{end}

\olpseganchor{OLP-0305-B011}{start}
پرېږدئ~\LR{$R(x,y)$} هله او يوازې هله صدق وکړي چې
\begin{quote}
\LR{$x$} د~\LR{$!D(u)$} فارمول کوډوي، او~\LR{$\Th{T}$}، \LR{$!D(\num y)$} ثابتوي۔
\end{quote}
ځکه چې فرض کوو~\LR{$\Th{T}$} د پرېکړې وړ ده، \LR{$R$} محاسبه کېدونکې ده۔
راځئ وښيو چې~\LR{$R$} نړيواله ده۔ که~\LR{$S(y)$} هره محاسبه کېدونکې اړيکه
وي، نو په~\LR{$\Th{Q}$} (او له دې کبله په~\LR{$\Th{T}$}) کښې د کوم
فارمول~\LR{$!D_S(u)$} په وسيله تمثيلېږي۔ بيا د هر~\LR{$n$} دپاره لرو:
\begin{eqnarray*}
S(n) & \lif & T \vdash !D_S(\num n) \\
& \lif & R(\Gn{!D_S(u)}, n)
\end{eqnarray*}
او
\begin{eqnarray*}
\lnot S(n) & \lif & T \vdash \lnot !D_S(\num n) \\
& \lif & T \not\vdash !D_S(\num n) \quad \text{\RL{(ځکه~\LR{$\Th{T}$} سازګاره ده)}} \\
& \lif & \lnot R(\Gn{!D_S(u)}, n).
\end{eqnarray*}
يعنې د هر~\LR{$y$} دپاره~\LR{$S(y)$} هله او يوازې هله رښتيا وي چې
\LR{$R(\Gn{!D_S(u)}, y)$} رښتيا وي۔ نو~\LR{$R$} نړيواله ده، او هغه تناقض مو
ترلاسه کړ چې غوښت۔
\end{proof}
\olpseganchor{OLP-0305-B011}{end}

\olpseganchor{OLP-0305-B012}{start}
پرېږدئ «رښتينی حساب» تيوري~\LR{$\Setabs{!A}{ \Sat{\Nat}{!A}}$} وي؛
يعنې د حساب د ژبې د هغو جملو سټ چې په معياري تفسير کښې رښتيا دي۔
\olpseganchor{OLP-0305-B012}{end}

\olpseganchor{OLP-0305-B013}{start}
\begin{cor}
رښتينی حساب د پرېکړې وړ نۀ دے۔
\end{cor}
\olpseganchor{OLP-0305-B013}{end}

\olpseganchor{OLP-0305-B014}{start}
%In Section~\olref{undefinability:truth:section} we will state a stronger
%result, due to Tarski.
\olpseganchor{OLP-0305-B014}{end}

\olpseganchor{OLP-0306-B004}{start}
\olfileid{inc}{tcp}{cax}
\olpseganchor{OLP-0306-B004}{end}

\olpseganchor{OLP-0306-B005}{start}
\olsection{په محاسبوي ډول بديهي کېدونکې تيورۍ}
\olpseganchor{OLP-0306-B005}{end}

\olpseganchor{OLP-0306-B006}{start}
يوه تيوري~\LR{$\Th{T}$} هله~\emph{په محاسبوي ډول بديهي کېدونکې} بلل کېږي چې د بديهي
اصولو يو محاسبه کېدونکے سټ~\LR{$A$} ولري۔ (دا ويل چې~\LR{$A$} د~\LR{$\Th{T}$}
د بديهي اصولو سټ دے، معنا ئې دا ده چې
\LR{$T = \Setabs{!B}{A \Proves !B}$}۔)
د طبيعي عددونو هره «معقوله» بديهي کول به دا خاصيت لري۔ په ځانګړې
توګه، هره هغه تيوري چې د بديهي اصولو متناهي سټ ولري،
په محاسبوي ډول بديهي کېدونکې وي۔
\olpseganchor{OLP-0306-B006}{end}

\olpseganchor{OLP-0306-B007}{start}
\begin{lem}
فرض کړئ~\LR{$\Th{T}$}، په محاسبوي ډول بديهي کېدونکې ده۔ بيا~\LR{$\Th{T}$}،
په محاسبوي ډول د شمېر وړ ده۔
\end{lem}
\olpseganchor{OLP-0306-B007}{end}

\olpseganchor{OLP-0306-B008}{start}
\begin{proof}
فرض کړئ~\LR{$A$} د~\LR{$\Th{T}$} د بديهي اصولو يو محاسبه کېدونکے سټ دے۔ د
دې دپاره چې د~\LR{$!A \in T$} لپاره يو مثبت شاهد ومومو، له بديهي اصولو څخه
د~\LR{$!A$} د اشتقاق پلټنه کوو۔ د ممکنو ثبوتونو کوډونه په ترتيب ګورو؛ په هر کوډ کښې کتلے شو چې آيا ثبوت سم دے او کارول شوي بديهي اصول په ټاکلي سټ کښې دي۔ که جمله ثابتېدونکې وي، پلټنه ئې يو ثبوت مومي۔ که ثبوت ونۀ لري، پلټنه روانه پاتې کېږي۔
\olpseganchor{OLP-0306-B008}{end}

\olpseganchor{OLP-0306-B009}{start}
لږ په بله وينا، \LR{$!A$} هله او يوازې هله په~\LR{$\Th{T}$} کښې ده چې
په~\LR{$A$} کښې د بديهي اصولو يو متناهي لست~\LR{$!B_1$}، \dots، \LR{$!B_k$} او
په لومړۍ درجې منطق کښې د~\LR{$(!B_1 \land \dots \land !B_k) \lif !A$}
يو اشتقاق وي۔ خو موږ لا پوهېږو چې هر هغه سټ چې تعريف ئې د
«داسې \dots شته چې \dots» په بڼه وي، c.e. وي، په دې شرط چې
دويم «\dots» محاسبه کېدونکے وي۔
\end{proof}
\olpseganchor{OLP-0306-B009}{end}

\olpseganchor{OLP-0307-B004}{start}
\olfileid{inc}{tcp}{cdc}
\olpseganchor{OLP-0307-B004}{end}

\olpseganchor{OLP-0307-B005}{start}
\olsection{په محاسبوي ډول بديهي کېدونکې بشپړې تيورۍ د پرېکړې وړ دي}
\olpseganchor{OLP-0307-B005}{end}

\olpseganchor{OLP-0307-B006}{start}
يوه تيوري هله~\emph{بشپړه} بلل کېږي چې د هرې جملې~\LR{$!A$} دپاره يا
\LR{$!A$} او يا~\LR{$\lnot !A$} ثابتېدونکې وي۔
\olpseganchor{OLP-0307-B006}{end}

\olpseganchor{OLP-0307-B007}{start}
\begin{lem}
فرض کړئ يوه تيوري~\LR{$\Th{T}$} بشپړه او په محاسبوي ډول بديهي کېدونکې ده۔ بيا
\LR{$\Th{T}$} د پرېکړې وړ ده۔
\end{lem}
\olpseganchor{OLP-0307-B007}{end}

\olpseganchor{OLP-0307-B008}{start}
\begin{proof}
فرض کړئ~\LR{$\Th{T}$} بشپړه ده او~\LR{$A$} د بديهي اصولو يو محاسبه کېدونکے
سټ دے۔ که~\LR{$\Th{T}$} ناسازګاره وي، نو څرګنده ده چې محاسبه کېدونکې
ده۔ (الګوريتم: «تل هو ووايه۔») نو فرض کولے شو چې~\LR{$\Th{T}$} سازګاره
هم ده۔
\olpseganchor{OLP-0307-B008}{end}

\olpseganchor{OLP-0307-B009}{start}
د دې پرېکړې دپاره چې يوه جمله~\LR{$!A$} په~\LR{$\Th{T}$} کښې ده که نۀ،
په يو وخت له~\LR{$\Th{T}$} څخه هم د~\LR{$!A$} د اشتقاق او هم د
\LR{$\lnot !A$} د اشتقاق پلټنه وکړئ۔ ځکه~\LR{$\Th{T}$} بشپړه ده،
خامخا به يو يا بل ومومئ؛ او ځکه~\LR{$\Th{T}$} سازګاره ده، که د
\LR{$\lnot !A$} يو اشتقاق ومومئ، نو د~\LR{$!A$} هېڅ اشتقاق
نشته۔
\olpseganchor{OLP-0307-B009}{end}

\olpseganchor{OLP-0307-B010}{start}
په بله وينا، موږ لا پوهېږو چې~\LR{$\Th{T}$}، c.e. ده؛ نو د هغې
قضيې له مخې چې مخکښې مو ثابته کړې، بس دا ښودل پکار دي چې د~\LR{$\Th{T}$}
متمم هم~c.e. دے۔ خو فارمول~\LR{$!A$} هله او يوازې هله په
\LR{$\Th{\bar T}$} کښې دے چې~\LR{$\lnot !A$} په~\LR{$\Th{T}$} کښې وي؛ نو
\LR{$\Th{\bar T} \leq_m \Th{T}$}۔
\end{proof}
\olpseganchor{OLP-0307-B010}{end}

\olpseganchor{OLP-0308-B004}{start}
\olfileid{inc}{tcp}{inc}
\olpseganchor{OLP-0308-B004}{end}

\olpseganchor{OLP-0308-B005}{start}
\olsection{\texorpdfstring{\LR{$\Th{Q}$}}{Q} هېڅ بشپړه، سازګاره او په محاسبوي ډول بديهي کېدونکې
  غځونه نۀ لري}
\olpseganchor{OLP-0308-B005}{end}

\olpseganchor{OLP-0308-B006}{start}
\begin{thm}
\ollabel{thm:first-incompleteness}
د~\LR{$\Th{Q}$} هېڅ بشپړه، سازګاره او په محاسبوي ډول بديهي کېدونکې غځونه نشته۔
\end{thm}
\olpseganchor{OLP-0308-B006}{end}

\olpseganchor{OLP-0308-B007}{start}
\begin{proof}
موږ لا پوهېږو چې د~\LR{$\Th{Q}$} هېڅ سازګاره او د پرېکړې وړ غځونه
نشته۔ خو که~\LR{$\Th{T}$} بشپړه او بديهي‌کړې وي، نو د پرېکړې وړ
ده۔
\end{proof}
\olpseganchor{OLP-0308-B007}{end}

\olpseganchor{OLP-0308-B008}{start}
\begin{explain}
دا قضيه د ګوډل د ناتکميلۍ د لومړۍ قضيې له اصلي ۱۹۳۱ز بيان څخه ډېره
لرې نۀ ده۔ له نوې اصطلاح پرته، مهم توپيرونه دا دي: ګوډل د
«سازګار» پر ځاے «\LR{$\omega$}-سازګار» لري؛ او هغه نۀ شو کولے په بشپړ
عموميت «په محاسبوي ډول بديهي کېدونکې» ووايي، ځکه هغه مهال د محاسبه کېدنې صوري
مفهوم لا جوړ شوے نۀ و۔ (د محاسبه کېدنې صوري مدلونه په راتلونکې
لسيزه کښې رامنځته شول، چې ګوډل هم پکې برخه لرله، او تر ډېره د دې
دپاره چې د هغو تيوريو ډولونه ځانګړي کړي چې د ګوډل ښکارندې پرې پلې
کېږي۔)
\olpseganchor{OLP-0308-B008}{end}

\olpseganchor{OLP-0308-B009}{start}
قضيه وايي چې هر څه يوځای نۀ شئ لرلے؛ يعنې بشپړتيا، سازګاري او
د محاسبوي بديهي کېدو وړتيا۔ خو که له دې هر يو پرېږدئ، نور دوه لرلے شئ:
\LR{$\Th{Q}$} سازګاره او د محاسبوي بديهي اصولو لرونکې ده، خو بشپړه نۀ
ده؛ ناسازګاره تيوري بشپړه ده او محاسبوي بديهي اصول لري (مثلاً د
\LR{$\{ \eq/[0][0] \}$} په وسيله)، خو سازګاره نۀ ده؛ او د حساب د رښتينو
جملو سټ بشپړ او سازګار دے، خو د محاسبوي بديهي اصولو لرونکے نۀ دے۔
\end{explain}
\olpseganchor{OLP-0308-B009}{end}

\olpseganchor{OLP-0309-B004}{start}
\olfileid{inc}{tcp}{ins}
\olpseganchor{OLP-0309-B004}{end}

\olpseganchor{OLP-0309-B005}{start}
\olsection{په~\texorpdfstring{\LR{$\Th{Q}$}}{Q} کښې ثابتېدونکې او ردېدونکې جملې په محاسبوي
  ډول نه بېلېدونکې دي}
\olpseganchor{OLP-0309-B005}{end}


\olpseganchor{OLP-0309-B006}{start}
پرېږدئ~\LR{$\Th{\bar Q}$} د هغو جملو سټ وي چې \emph{نفيانې} ئې په
\LR{$\Th{Q}$} کښې ثابتېدونکې دي؛ يعنې
\LR{$\Th{\bar Q} = \Setabs{!A}{\Th{Q} \Proves \lnot !A}$}۔ په ياد ولرئ
چې بېل سټونه~\LR{$A$} او~\LR{$B$} هله په محاسبوي ډول نه بېلېدونکي بلل کېږي
چې داسې محاسبه کېدونکے سټ~\LR{$C$} نۀ وي چې~\LR{$A \subseteq C$} او
\LR{$B \subseteq \Complement{C}$}۔
\olpseganchor{OLP-0309-B006}{end}

\olpseganchor{OLP-0309-B007}{start}
\begin{lem}
\LR{$\Th{Q}$} او~\LR{$\Th{\bar Q}$} په محاسبوي ډول نه بېلېدونکي دي۔
\end{lem}
\olpseganchor{OLP-0309-B007}{end}

\olpseganchor{OLP-0309-B008}{start}
\begin{proof}
فرض کړئ~\LR{$C$} داسې محاسبه کېدونکے سټ دے چې~\LR{$\Th{Q} \subseteq C$} او
\LR{$\Th{\bar Q} \subseteq \Complement{C}$}۔ پرېږدئ~\LR{$R(x,y)$} لاندې
اړيکه وي:
\begin{quote}
\LR{$x$} د~\LR{$!D(u)$} فارمول کوډوي او~\LR{$!D(\num y)$} په~\LR{$C$} کښې دے۔
\end{quote}
موږ به وښيو چې~\LR{$R(x,y)$} يوه نړيواله محاسبه کېدونکې اړيکه ده، چې
تناقض ترې ترلاسه کېږي۔
\olpseganchor{OLP-0309-B008}{end}

\olpseganchor{OLP-0309-B009}{start}
فرض کړئ~\LR{$S(y)$} محاسبه کېدونکې ده، او په~\LR{$\Th{Q}$} کښې د~\LR{$!D_S(u)$}
په وسيله تمثيلېږي۔ بيا
\begin{eqnarray*}
S(n) & \lif & \Th{Q} \Proves !D_S(\num n) \\
& \lif & !D_S(\num n) \in C
\end{eqnarray*}
او
\begin{eqnarray*}
\lnot S(n) & \lif & \Th{Q} \Proves \lnot !D_S(\num n) \\
& \lif & !D_S(\num n) \in \Th{\bar Q} \\
& \lif & !D_S(\num n) \not\in C
\end{eqnarray*}
نو~\LR{$S(y)$} له~\LR{$R(\Gn{!D_S(u)},y)$} سره هم‌ارزه ده۔
\end{proof}
\olpseganchor{OLP-0309-B009}{end}

\olpseganchor{OLP-0310-B004}{start}
\olfileid{inc}{tcp}{con}
\olpseganchor{OLP-0310-B004}{end}

\olpseganchor{OLP-0310-B005}{start}
\olsection{له~\texorpdfstring{\LR{$\Th{Q}$}}{Q} سره سازګارې تيورۍ ناپرېکړه‌وړې دي}
\olpseganchor{OLP-0310-B005}{end}

\olpseganchor{OLP-0310-B006}{start}
لاندې قضيه وايي چې يوازې~\LR{$\Th{Q}$} ناپرېکړه‌وړې نۀ ده، بلکې په
حقيقت کښې هره هغه تيوري ناپرېکړه‌وړې ده چې له~\LR{$\Th{Q}$} سره اختلاف
ونۀ لري۔
\olpseganchor{OLP-0310-B006}{end}

\olpseganchor{OLP-0310-B007}{start}
\begin{thm}
پرېږدئ~\LR{$\Th{T}$} د حساب په ژبه کښې هره هغه تيوري وي چې له~\LR{$\Th{Q}$}
سره سازګاره وي (يعنې~\LR{$\Th{T} \cup \Th{Q}$} سازګاره وي)۔ بيا
\LR{$\Th{T}$} ناپرېکړه‌وړې ده۔
\end{thm}
\olpseganchor{OLP-0310-B007}{end}

\olpseganchor{OLP-0310-B008}{start}
\begin{proof}
په ياد ولرئ چې~\LR{$\Th{Q}$} د بديهي اصولو متناهي سټ~\LR{$!Q_1$}، \dots،
\LR{$!Q_8$} لري۔ ان دا ټول په يوه بديهي اصل بدلولے شو:
\LR{$!E = !Q_1 \land \dots \land !Q_8$}۔
\olpseganchor{OLP-0310-B008}{end}

\olpseganchor{OLP-0310-B009}{start}
فرض کړئ~\LR{$\Th{T}$} يوه د پرېکړې وړ تيوري ده چې له~\LR{$\Th{Q}$} سره
سازګاره ده۔ پرېږدئ
\[
C = \Setabs{!A}{\Th{T} \Proves !E \lif !A}.
\]
موږ ښيو چې~\LR{$C$} به د~\LR{$\Th{Q}$} او~\LR{$\Th{\bar Q}$} يوه محاسبه کېدونکې
بېلوونکې وي، چې تناقض دے۔ لومړے، که~\LR{$!A$} په~\LR{$\Th{Q}$} کښې وي، نو
\LR{$!A$} د~\LR{$\Th{Q}$} له بديهي اصولو څخه ثابتېدونکې ده؛ د استنتاج د
قضيې له مخې په لومړۍ درجې منطق کښې د~\LR{$!E \lif !A$} يو
اشتقاق شته۔ نو~\LR{$!A$} په~\LR{$C$} کښې ده۔
\olpseganchor{OLP-0310-B009}{end}

\olpseganchor{OLP-0310-B010}{start}
له بلې خوا، که~\LR{$!A$} په~\LR{$\Th{\bar Q}$} کښې وي، نو په لومړۍ درجې
منطق کښې د~\LR{$!E \lif \lnot !A$} يو ثبوت شته۔ که~\LR{$\Th{T}$}،
\LR{$!E \lif !A$} هم ثابت کړي، نو~\LR{$\Th{T}$}، \LR{$\lnot !E$} ثابتوي، او په
دې حالت کښې~\LR{$\Th{T} \cup \Th{Q}$} ناسازګاره ده۔ خو موږ فرض کوو
چې~\LR{$\Th{T} \cup \Th{Q}$} سازګاره ده، نو~\LR{$\Th{T}$}، \LR{$!E \lif !A$}
نۀ ثابتوي، او له دې امله~\LR{$!A$} په~\LR{$C$} کښې نۀ ده۔
\olpseganchor{OLP-0310-B010}{end}

\olpseganchor{OLP-0310-B011}{start}
موږ وښودله چې که~\LR{$!A$} په~\LR{$\Th{Q}$} کښې وي، نو په~\LR{$C$} کښې ده، او
که~\LR{$!A$} په~\LR{$\Th{\bar Q}$} کښې وي، نو په~\LR{$\Complement{C}$} کښې ده۔
نو~\LR{$C$} يوه محاسبه کېدونکې بېلوونکې ده، او همدا هغه تناقض دے چې
غوښت مو۔
\end{proof}
\olpseganchor{OLP-0310-B011}{end}

\olpseganchor{OLP-0310-B012}{start}
دا قضيه ډېره پياوړې ده۔ د بېلګې په توګه، له دې څخه لاندې پايله
لازمېږي:
\olpseganchor{OLP-0310-B012}{end}

\olpseganchor{OLP-0310-B013}{start}
\begin{cor}
  د حساب د ژبې دپاره لومړۍ درجې منطق (يعنې سټ
  \LR{$\Setabs{!A}{\text{\RL{\LR{$!A$} په لومړۍ درجې منطق کښې ثابتېدونکې ده}}}$})
  ناپرېکړه‌وړ دے۔
\end{cor}
\olpseganchor{OLP-0310-B013}{end}

\olpseganchor{OLP-0310-B014}{start}
\begin{proof}
لومړۍ درجې منطق د~\LR{$\emptyset$} د پايلو سټ دے، چې له~\LR{$\Th{Q}$} سره
سازګار دے۔
\end{proof}
\olpseganchor{OLP-0310-B014}{end}

\olpseganchor{OLP-0311-B004}{start}
\olfileid{inc}{tcp}{itp}
\olpseganchor{OLP-0311-B004}{end}

\olpseganchor{OLP-0311-B005}{start}
\olsection{هغه تيورۍ ناپرېکړه‌وړې دي چې~\texorpdfstring{\LR{$\Th{Q}$}}{Q} پکې تفسيرېدونکې وي}
\olpseganchor{OLP-0311-B005}{end}

\olpseganchor{OLP-0311-B006}{start}
دا پايلې لا نورې هم پياوړې کولے شو۔ په غيرصوري ډول، په بله
ژبه~\LR{$\Lang{L_2}$} کښې د يوې ژبې~\LR{$\Lang{L_1}$} تفسيرول دا غواړي چې د
\LR{$\Lang{L_1}$} کائنات، اړيکيزې نښې او تابعي نښې په~\LR{$\Lang{L_2}$} کښې
د فارمولونه په وسيله تعريف شي۔ که څه هم دلته ورته وخت نۀ ورکوو،
دا تعريف په کره بڼه جوړېدے شي۔
\olpseganchor{OLP-0311-B006}{end}

\olpseganchor{OLP-0311-B007}{start}
\begin{thm}
  فرض کړئ~\LR{$\Th{T}$} په داسې ژبه کښې يوه تيوري ده چې د حساب ژبه پکې
  داسې تفسيرېدے شي چې~\LR{$\Th{T}$} د~\LR{$\Th{Q}$} له تفسير سره سازګاره وي۔
  بيا~\LR{$\Th{T}$} ناپرېکړه‌وړې ده۔ که~\LR{$\Th{T}$} د~\LR{$\Th{Q}$} د بديهي
  اصولو تفسير ثابتوي، نو د~\LR{$\Th{T}$} هېڅ سازګاره غځونه د پرېکړې وړ
  نۀ ده۔
\end{thm}
\olpseganchor{OLP-0311-B007}{end}

\olpseganchor{OLP-0311-B008}{start}
فرض کړئ چې د~\LR{$\Th{T}$} پرېکړه کولے شو۔ د تفسير شويو بديهي اصولو متناهي سټ~\LR{$\Th{T}$} ته ورزياتوو۔ د فرض له مخې غځونه سازګاره ده۔ د دې معلومولو لپاره چې يوه جمله له غځونې لازمېږي که نه، په اصلي تيورۍ کښې هغه شرطي جمله ازمويو چې مقدم ئې د زياتو اصولو عطف او تالي ئې هماغه جمله ده۔ نو غځونه هم د پرېکړې وړ ده۔ د حساب هرې جملې ته د هغې تفسير شوې بڼه په محاسبوي ډول جوړوو، او د هغو حسابي جملو سټ اخلو چې تفسير ئې په غځونه کښې ثابتېږي۔ دا سټ~\LR{$\Th{Q}$} رانغاړي، ځکه تفسير د استنتاج ګامونه ساتي۔ دا سټ د~\LR{$\Th{\bar Q}$} هېڅ جمله نۀ رانغاړي، ځکه د هغې د نفي تفسير په غځونه کښې ثابتېږي، او غځونه سازګاره ده۔ نو محاسبه کېدونکے بېلوونکے مو ترلاسه کړ، چې تناقض دے۔ که تفسير شوي اصول له وړاندې په اصلي تيورۍ کښې ثابتېږي، د هرې سازګارې غځونې لپاره همدغه دليل کاروو۔
\LR{$\Th{ZFC}$}، يعنې د انتخاب له بديهي اصل سره د زر مېلو--فرېنکل سټ
تيوري، داسې بديهي بنسټ ګڼلے شو چې د عادي رياضياتو ډېره برخه پکې
ترسره کېدے شي۔ په~\LR{$\Th{ZFC}$} کښې طبيعي عددونه تعريفولے شو، او د
دې تفسير له مخې د~\LR{$\Th{Q}$} بديهي اصول رښتيا دي۔ نو لرو:
\olpseganchor{OLP-0311-B008}{end}

\olpseganchor{OLP-0311-B009}{start}
\begin{cor}
د~\LR{$\Th{ZFC}$} هېڅ سازګاره او د پرېکړې وړ غځونه نشته۔
\end{cor}
\olpseganchor{OLP-0311-B009}{end}

\olpseganchor{OLP-0311-B010}{start}
\begin{cor}
د~\LR{$\Th{ZFC}$} هېڅ بشپړه، سازګاره او په محاسبوي ډول
په محاسبوي ډول بديهي کېدونکې غځونه نشته۔
\end{cor}
\olpseganchor{OLP-0311-B010}{end}

\olpseganchor{OLP-0311-B011}{start}
د~\LR{$\Th{ZFC}$} ژبه يوازې يوه دوه‌ځاييزه اړيکه~\LR{$\in$} لري۔ (په حقيقت
کښې ان عينيت ته هم اړتيا نشته۔) نو لرو:
\olpseganchor{OLP-0311-B011}{end}

\olpseganchor{OLP-0311-B012}{start}
\begin{cor}
د هرې هغې ژبې دپاره لومړۍ درجې منطق ناپرېکړه‌وړ دے چې يوه
دوه‌ځاييزه اړيکيزه نښه ولري۔
\end{cor}
\olpseganchor{OLP-0311-B012}{end}

\olpseganchor{OLP-0311-B013}{start}
دا پايله هرې هغې ژبې ته هم غځېږي چې دوه يوځاييزې تابعي نښې ولري،
ځکه د هغو په وسيله يوه دوه‌ځاييزه اړيکيزه نښه شبيه کېدے شي۔ همدا
اوس يادې شوې پايلې کره پولې دي: په حقيقت کښې د هرې هغې ژبې دپاره
لومړۍ درجې منطق د پرېکړې وړ دے چې يوازې \emph{يوځاييزې} اړيکيزې
نښې او تر ډېره يوه \emph{يوځاييزه} تابعي نښه ولري۔
\olpseganchor{OLP-0311-B013}{end}

\olpseganchor{OLP-0311-B014}{start}
يو بل په زړه پورې حقيقت۔ پوهېږو چې په معياري مدل کښې د ژبې
\LR{$\Obj 0$}، \LR{$'$}، \LR{$+$}، \LR{$\times$}، \LR{$<$} د رښتينو جملو سټ ناپرېکړه‌وړ
دے۔ په حقيقت کښې~\LR{$<$} د نورو نښو په وسيله تعريفولے شو، او بيا~\LR{$+$}
د~\LR{$\times$} او~\LR{$'$} په وسيله تعريفولے شو۔ نو د ژبې~\LR{$\Obj 0$}، \LR{$'$}،
\LR{$\times$} د رښتينو جملو سټ ناپرېکړه‌وړ دے۔ له بلې خوا، پرېسبورګر
ښودلې ده چې د ژبې~\LR{$\Obj 0$}، \LR{$'$}، \LR{$+$} د هغو جملو سټ چې د حساب په
معياري تفسير کښې رښتيا دي، د پرېکړې وړ دے۔
خو دا کړنلاره له محاسبوي پلوه عملي نۀ ده۔
\olpseganchor{OLP-0311-B014}{end}

\olpseganchor{OLP-0312-B004}{start}
\olchapter{inc}{inp}{ناتکميلي او اثبات‌وړتيا}
\olpseganchor{OLP-0312-B004}{end}

\olpseganchor{OLP-0313-B004}{start}
\olfileid{inc}{inp}{int}
\olpseganchor{OLP-0313-B004}{end}

\olpseganchor{OLP-0313-B005}{start}
\olsection{پېژندنه}
\olpseganchor{OLP-0313-B005}{end}

\olpseganchor{OLP-0313-B006}{start}
هېلبرټ په دې اند و چې د يو رياضيکي جوړښت، لکه د طبيعي عددونو، دپاره
د بديهي اصولو نظام تر هغه وخته بسيا نۀ دے چې د هماغه جوړښت ټولې
رښتينې جملې ترې استنتاج نۀ شي۔ د صوري استنتاج له نظامونو سره د هغه
وروستۍ لېوالتيا چې ورسره يوځاے کړو، نو داسې ښکاري چې هغه غوښتل د
طبيعي عددونو په باب زموږ صوري استدلالي نظامونه يوازې سازګار نۀ، بلکې
\emph{بشپړ} هم وي؛ يعنې د نظام د ژبې هره جمله يا پخپله د اشتقاق وړ
وي يا ئې نفي همداسې وي۔ د ګوډل د ناتکميلۍ لومړۍ قضيه ښيي چې
داسې د بديهي اصولو نظام نشته: د حساب دپاره هېڅ بشپړ، سازګار او
په محاسبوي ډول بديهي کېدونکې صوري نظام نشته۔ په حقيقت کښې، هېڅ «په کافي اندازه
پياوړې»، سازګاره او په محاسبوي ډول بديهي کېدونکې رياضيکي تيوري بشپړه نۀ ده۔
\olpseganchor{OLP-0313-B006}{end}

\olpseganchor{OLP-0313-B007}{start}
د هېلبرټ يو لا مهم هدف، چې د نوې («کلاسيکې») رياضي د توجيه دپاره
د هغه د پروګرام منځنے ټکی و، د کلاسيک استدلال استازيتوب کوونکو صوري
نظامونو د سازګارۍ متناهي‌پاله ثبوتونه موندل وو۔ نو د هېلبرټ د
پروګرام له پلوه د ګوډل د ناتکميلۍ دويمه قضيه تر لومړۍ ډېر دروند ګوزار
و۔ دويمه قضيه هم د لومړۍ په شان په لږ ناڅرګندو ټکو بيانېدای شي:
نږدې معنا ئې دا ده چې د حساب هېڅ په کافي اندازه پياوړې تيوري خپله
سازګاري نۀ شي ثابتولے۔ د «په کافي اندازه پياوړې» مطلب به له~\LR{$\Th{Q}$}
څخه لږ څه زيات وغواړي۔
\olpseganchor{OLP-0313-B007}{end}

\olpseganchor{OLP-0313-B008}{start}
د ناتکميلۍ د قضيې د ګوډل د اصلي ثبوت مفکوره د اپيمنيدس په تناقض
کښې موندل کېږي۔ اپيمنيدس، چې د کريت اوسېدونکے و، وويل چې د کريت
ټول اوسېدونکي دروغجن دي؛ د دې تناقض لا نېغه بڼه «دا جمله ناسمه ده»
ويل دي۔ په اصل کښې ګوډل د رښتياوالي پر ځاے د اشتقاق وړتيا وکاروله
او تړلے فارمول ئې صوري کړه چې په غير مستقيم ډول وايي: زه پخپله
د اشتقاق وړ نۀ يم۔ که دغه تړلے فارمول د اشتقاق وړ واے، تيوري به
ناسازګاره وه۔ ګوډل دا هم وښودله چې د هغه د بديهي اصولو له نظامه د
دغې تړلے فارمول نفي هم د اشتقاق وړ نۀ ده۔ (د دې دويمې برخې دپاره
ګوډل اړ و چې فرض کړي تيوري~\LR{$\Th{T}$} «\LR{$\omega$}-سازګاره» ده۔
\LR{$\omega$}-سازګاري له سازګارۍ سره تړاو لري، خو تر هغې پياوړے خاصيت
دے۔\footnote{يعنې هره \LR{$\omega$}-سازګاره تيوري سازګاره ده، خو
برعکس لازمه نۀ ده۔} له ګوډل څو کاله وروسته
روسر وښودله چې يوازې د~\LR{$\Th{T}$} عادي سازګاري فرض کول بس دي۔)
\olpseganchor{OLP-0313-B008}{end}

\olpseganchor{OLP-0313-B009}{start}
لومړۍ ستونزه دا پوهېدل دي چې داسې تړلے فارمول څنګه جوړېږي چې
ځان ته مراجعه وکړي۔ د~\LR{$\Th{Q}$} په ژبه کښې د هر فارمول~\LR{$!A$}
دپاره دې~\LR{$\gn{!A}$} هغه عدديه نښه وي چې له~\LR{$\Gn{!A}$} سره سمون خوري۔
د دې معنا ته ځير شئ: \LR{$!A$}~د~\LR{$\Th{Q}$} په ژبه کښې فارمول دے؛
\LR{$\Gn{!A}$}~يو طبيعي عدد دے؛ او \LR{$\gn{!A}$}~د~\LR{$\Th{Q}$} په ژبه کښې يو
\emph{ترم} دے۔ نو د~\LR{$\Th{Q}$} د ژبې هر فارمول~\LR{$!A$} يو
\emph{نوم}، يعنې~\LR{$\gn{!A}$}، لري چې د~\LR{$\Th{Q}$} په ژبه کښې ترم
دے۔ دا موږ ته
هغه مفهومي چوکاټ راکوي چې په هغه کښې د~\LR{$\Th{Q}$} د ژبې
فارمولونه د نورو فارمولونه په باب څه «ويلے» شي۔ لاندې لمه
د ثابت ټکي د لمې په نوم يادېږي۔
\olpseganchor{OLP-0313-B009}{end}

\olpseganchor{OLP-0313-B010}{start}
\begin{lem}
فرض کړئ \LR{$\Th{T}$} هره هغه تيوري ده چې~\LR{$\Th{Q}$} غځوي، او \LR{$!B(x)$}
هر هغه فارمول دے چې يوازې متغير~\LR{$x$} پکې آزاد وي۔ نو داسې
تړلے فارمول~\LR{$!A$} شته چې
\LR{$\Th{T} \Proves!A \liff !B(\gn{!A})$}۔
\end{lem}
\olpseganchor{OLP-0313-B010}{end}

\olpseganchor{OLP-0313-B011}{start}
لمه وايي چې د هر خاصيت~\LR{$!B(x)$} دپاره داسې تړلے فارمول~\LR{$!A$}
شته چې وايي «\LR{$!B(x)$} زما په باب رښتيا دے»، او~\LR{$\Th{T}$} هم دغه
برابري «پېژني»۔
\olpseganchor{OLP-0313-B011}{end}

\olpseganchor{OLP-0313-B012}{start}
داسې تړلے فارمول څنګه جوړولے شو؟ د اپيمنيدس د تناقض لاندې بڼه
وګورئ چې \href{https://www.scientificamerican.com/article/paradox/}{کواين په ۱۹۶۲ کښې} وړاندې کړې ده۔
\begin{explain}
يو عبارت واخلو چې د خپل نقل له يوې ټاکلې نحوي چارې سره د
دروغوالي حکم تړي۔ د عبارت خبره دا ده چې «که دا عبارت له خپل نقل وروسته وليکل شي، پايله يې يوه ناسمه جمله ده»۔
اوس لومړے عبارت په نقل‌نښو کښې وليکئ او سمدستي يې وروسته هماغه
عبارت بيا، بې له نقل‌نښو، وليکئ۔ په دې توګه يوه بشپړه جمله جوړېږي۔
\end{explain}
دا جمله په نېغه ځان ته مراجعه نۀ کوي۔ يوازې د نقل‌قول د نښو ترمنځ
د نحوي شيانو په باب خبره کوي، نو له دې جملو سره په يوه کچه ده:
\begin{enumerate}
\item «رابرټ» يو ښۀ نوم دے۔
\item «زه وځغاستم۔» يوه لنډه جمله ده۔
\item «درې کلمې لري» درې کلمې لري۔
\end{enumerate}
د حکم موضوع هماغه عبارت دے چې په نقل‌نښو کښې ليکل شوے دے۔
له نقل وروسته د عبارت بيا ليکل هماغه بشپړه جمله جوړوي؛ نو جمله
د همدې جوړې شوې جملې د دروغوالي حکم کوي۔ همدا د تناقض لامل دے۔
\olpseganchor{OLP-0313-B012}{end}

\olpseganchor{OLP-0314-B004}{start}
\olfileid{inc}{inp}{fix}
\olpseganchor{OLP-0314-B004}{end}

\olpseganchor{OLP-0314-B005}{start}
\olsection{د ثابت ټکي لمه}
\olpseganchor{OLP-0314-B005}{end}

\olpseganchor{OLP-0314-B006}{start}
\begin{explain}
د ثابت ټکي لمه وايي چې د هر فارمول~\LR{$!B(x)$} دپاره داسې
تړلے فارمول~\LR{$!A$} شته چې
\LR{$\Th{T} \Proves !A \liff !B(\gn{!A})$}، په دې شرط چې~\LR{$\Th{T}$}
تيوري~\LR{$\Th{Q}$} غځوي۔ د دروغجنې جملې په بېلګه کښې غواړو چې
\LR{$!A$} په~\LR{$\Th{T}$} کښې په ثابتېدونکي ډول له «\LR{$\gn{!A}$} ناسمه ده»
سره معادله وي؛ يعنې له هغې جملې سره چې وايي~\LR{$\Gn{!A}$} د يوې
ناسمې تړلے فارمول ګوډل شمېره ده۔ د ثبوت د مفکورې د پوهېدو دپاره
دا ګټوره ده چې هغه د~\LR{$!A$} په باب د کواین له دې غير صوري څرګندونې
سره پرتله کړو:
«که يې خپله نقل شوې بڼه ترې مخکې کېښودل شي، دروغ وايي» که يې
خپله نقل شوې بڼه ترې مخکې کېښودل شي، دروغ وايي۔ له يوه عبارت څخه
د هغه خپله نقل شوې بڼه مخکې ايښودل او بيا ترې جمله جوړول د هغه
عبارت \emph{قطرول} بللے شو، او پايله ئې د هغه قطري بڼه ده۔ نو د
«که يې خپله نقل شوې بڼه ترې مخکې کېښودل شي، دروغ وايي» قطري بڼه
همدا «که يې خپله نقل شوې بڼه ترې مخکې کېښودل شي، دروغ وايي» که يې
خپله نقل شوې بڼه ترې مخکې کېښودل شي، دروغ وايي، ده۔ اوس پام وکړئ
چې د کواین دروغجنه جمله د «دروغ وايي» قطري بڼه نۀ ده، بلکې د
«که يې خپله نقل شوې بڼه ترې مخکې کېښودل شي، دروغ وايي» قطري بڼه
ده۔ نو هغه خاصيت چې د دروغجنې جملې د جوړولو دپاره قطرول کېږي،
پخپله هم قطرول رانغاړي!{}
\olpseganchor{OLP-0314-B006}{end}

\olpseganchor{OLP-0314-B007}{start}
د حساب په ژبه کښې د يو آزاد متغير لرونکي فارمول نقل شوې بڼه
د هغه د ګوډل شمېرې په حسابولو او بيا د دغې شمېرې د معياري عدديې
په آزاد متغير کښې په ځاے کولو جوړوو۔ د~\LR{$!E(x)$} قطري بڼه
\LR{$!E(\num{n})$} ده، چېرته چې \LR{$n = \Gn{!E(x)}$} دے۔ (له دې وروسته
\LR{$\num{\Gn{!E(x)}}$} په~\LR{$\gn{!E(x)}$} لنډوو۔) نو که \LR{$!B(x)$} د
«ناسمه ده» خاصيت وي، «که خپله نقل شوې بڼه ترې مخکې راشي، دروغ وايي»
به دا وي: «کله چې د خپلې قطري بڼې پر ګوډل شمېرې پلے شي، دروغ وايي»۔
که د هغې تابع~\LR{$\fn{diag}(n)$} دپاره چې د ګوډل شمېرې~\LR{$n$} لرونکي
فارمول د قطري بڼې ګوډل شمېره حسابوي، يوه نښه~\LR{$\Obj{diag}$}
ولرو، \LR{$!E(x)$} به د~\LR{$!B(\Obj{diag}(x))$} په بڼه وليکو۔ د کواین
د دروغجنې جملې بڼه به بيا د همدې عبارت قطري بڼه وي، يعنې
\LR{$!E(\gn{!E(x)})$} يا
\LR{$!B(\Obj{diag}(\gn{!B(\Obj{diag}(x))}))$}۔ البته اوس \LR{$!B(x)$}
هر بل خاصيت هم کېدے شي، او همدا جوړښت به بيا هم کار وکړي۔ د
ناتکميلۍ د قضيې دپاره \LR{$!B(x)$} داسې اخلو: «\LR{$x$} په~\LR{$\Th{T}$} کښې
د اشتقاق وړ نۀ دے»۔ نو \LR{$!E(x)$} به ووايي: «کله چې د خپلې قطري
بڼې پر ګوډل شمېرې پلے شي، داسې تړلے فارمول جوړوي چې په~\LR{$\Th{T}$}
کښې د اشتقاق وړ نۀ ده»۔
\olpseganchor{OLP-0314-B007}{end}

\olpseganchor{OLP-0314-B008}{start}
د دې خبرې په~\LR{$\Th{T}$} کښې د صوري کولو دپاره بايد د~\LR{$\fn{diag}$}
د صوري کولو لار ومومو۔ تابع~\LR{$\fn{diag}(n)$} محاسبه کېدونکې ده، بلکې
ابتدايي بازګشتي ده: که \LR{$n$} د يوه فورمول~\LR{$!E(x)$} ګوډل شمېره وي،
\LR{$\fn{diag}(n)$} د~\LR{$!E(\gn{!E(x)})$} ګوډل شمېره ورکوي۔ (ياد ولرئ
چې~\LR{$\gn{!E(x)}$} د~\LR{$!E(x)$} د ګوډل شمېرې معياري عدديه، يعنې
\LR{$\num{\Gn{!E(x)}}$}، ده۔) که \LR{$\Obj{diag}$} په~\LR{$\Th{T}$} کښې د
تابع~\LR{$\fn{diag}$} تمثيلوونکې تابع‌نښه واے، موږ به \LR{$!A$} د
\LR{$!B(\Obj{diag}(\gn{!B(\Obj{diag}(x))}))$} فورمول اخيستے واے۔
پام وکړئ چې
\begin{align*}
\fn{diag}(\Gn{!B(\Obj{diag}(x))}) & =
\Gn{!B(\Obj{diag}(\gn{!B(\Obj{diag}(x))}))} \\
& = \Gn{!A}.
\end{align*}
که فرض کړو چې~\LR{$\Th{T}$}
\[
\Obj{diag}(\gn{!B(\Obj{diag}(x))}) = \gn{!A},
\]
اشتقاق کولے شي، نو دا هم اشتقاق کولے شي چې
\LR{$!B(\Obj{diag}(\gn{!B(\Obj{diag}(x))}))
\liff !B(\gn{!A})$}۔ خو د دې معادلۍ کيڼ اړخ د تعريف له مخې~\LR{$!A$}
دے۔
\olpseganchor{OLP-0314-B008}{end}

\olpseganchor{OLP-0314-B009}{start}
البته \LR{$\Obj{diag}$} په عمومي توګه د~\LR{$\Th{T}$} تابع‌نښه نۀ ده، او
يقيناً د~\LR{$\Th{Q}$} تابع‌نښه نۀ ده۔ خو ځکه چې~\LR{$\fn{diag}$} محاسبه
کېدونکې ده، په~\LR{$\Th{Q}$} کښې د کوم فورمول
\LR{$!D_{\fn{diag}}(x,y)$} په وسيله \emph{تمثيلېږي}۔ نو د
\LR{$!B(\Obj{diag}(x))$} پر ځاے
\LR{$\lexists[y][(!D_{\fn{diag}}(x,y) \land !B(y))]$} ليکلے شو۔
له دې پرته، پورته لنډ بيان شوے ثبوت بشپړېږي، او په حقيقت کښې
لا په~\LR{$\Th{Q}$} کښې بشپړېږي۔
\end{explain}
\olpseganchor{OLP-0314-B009}{end}

\olpseganchor{OLP-0314-B010}{start}
\begin{lem}
\ollabel{lem:fixed-point} فرض کړئ \LR{$!B(x)$} هر هغه فورمول دے چې
يو آزاد متغير~\LR{$x$} لري۔ نو داسې جمله~\LR{$!A$} شته چې
\LR{$\Th{Q} \Proves !A \liff !B(\gn{!A})$}۔
\end{lem}
\olpseganchor{OLP-0314-B010}{end}

\olpseganchor{OLP-0314-B011}{start}
\begin{proof}
د ورکړي~\LR{$!B(x)$} دپاره دې \LR{$!E(x)$} فورمول
\LR{$\lexists[y][(!D_{\fn{diag}}(x,y) \land !B(y))]$} وي، او
\LR{$!A$} دې د هغه قطري بڼه، يعنې \LR{$!E(\gn{!E(x)})$} فورمول وي۔
\olpseganchor{OLP-0314-B011}{end}

\olpseganchor{OLP-0314-B012}{start}
ځکه چې \LR{$!D_{\fn{diag}}$} د~\LR{$\fn{diag}$} تمثيل کوي، او
\LR{$\fn{diag}(\Gn{!E(x)}) = \Gn{!A}$}، نو~\LR{$\Th{Q}$} دا
اشتقاق کولے شي:
\begin{align}
  & !D_{\fn{diag}}(\gn{!E(x)}, \gn{!A}) \ollabel{repdiag1} \\
  & \lforall[y][(!D_{\fn{diag}}(\gn{!E(x)},y) \lif
  \eq[y][\gn{!A}])]. \ollabel{repdiag2}
\end{align}
اوس ښيو چې \LR{$\Th{Q} \Proves !A \liff !B(\gn{!A})$}۔ موږ يوازې د
منطق او په~\LR{$\Th{Q}$} کښې د د اشتقاق وړ حقايقو په کارولو، په غير
صوري ډول استدلال کوو۔
\olpseganchor{OLP-0314-B012}{end}

\olpseganchor{OLP-0314-B013}{start}
لومړے فرض کړئ~\LR{$!A$}، يعنې \LR{$!E(\gn{!E(x)})$}۔ د~\LR{$!E(x)$} تعريف
ته په ستنېدو وينو چې \LR{$!E(\gn{!E(x)})$} همدا دے:
\[
\lexists[y][(!D_{\fn{diag}}(\gn{!E(x)},y) \land !B(y))].
\]
داسې يو~\LR{$y$} وګورئ۔ ځکه چې
\LR{$!D_{\fn{diag}}(\gn{!E(x)},y)$} دے، د~\olref{repdiag2} له مخې
\LR{$y = \gn{!A}$}۔ نو له~\LR{$!B(y)$} څخه~\LR{$!B(\gn{!A})$} اخلو۔
\olpseganchor{OLP-0314-B013}{end}

\olpseganchor{OLP-0314-B014}{start}
اوس فرض کړئ \LR{$!B(\gn{!A})$}۔ د~\olref{repdiag1} له مخې لرو:
\begin{align*}
& !D_{\fn{diag}}(\gn{!E(x)}, \gn{!A}) \land !B(\gn{!A}).
\intertext{\RL{نو ورڅخه راځي چې}}
& \lexists[y][(!D_{\fn{diag}}(\gn{!E(x)},y) \land !B(y))].
\end{align*}
خو دا هماغه \LR{$!E(\gn{!E(x)})$}، يعنې~\LR{$!A$}، دے۔
\end{proof}
\olpseganchor{OLP-0314-B014}{end}

\olpseganchor{OLP-0314-B015}{start}
\begin{digress}
دا ثبوت د محاسبه کېدنې په تيورۍ کښې د ثابت ټکي د لمې له ثبوت سره
پرتله کړئ۔ توپير دا دے چې دلته غواړو يو \emph{بيان} د خپل ځان په
وسيله تعريف کړو، خو هلته مو يوه \emph{تابعه} د خپل ځان په وسيله
تعريفوله؛ له دې توپيره پرته مفکوره په رښتيا هماغه ده۔
\end{digress}
\olpseganchor{OLP-0314-B015}{end}

\olpseganchor{OLP-0314-B016}{start}
\begin{prob}
  فارمول~\LR{$!A(x)$} د \emph{رښتياوالي تعريف} دے، که د هرې
  تړلي فارمولونه~\LR{$!B$} دپاره \LR{$\Th{Q} \Proves !B \liff !A(\gn{!B})$} وي۔
  د ثابت ټکي د لمې په کارولو وښيئ چې هېڅ فارمول د رښتياوالي
  تعريف نۀ دے۔
\end{prob}
\olpseganchor{OLP-0314-B016}{end}

\olpseganchor{OLP-0315-B004}{start}
\olfileid{inc}{inp}{1in}
\olpseganchor{OLP-0315-B004}{end}

\olpseganchor{OLP-0315-B005}{start}
\olsection{د ناتکميلۍ لومړۍ قضيه}
\olpseganchor{OLP-0315-B005}{end}

\olpseganchor{OLP-0315-B006}{start}
اوس د ګوډل د ناتکميلۍ د لومړۍ قضيې اصلي ثبوت بيانولے شو۔ فرض کړئ
\LR{$\Th{T}$} په هغې ژبه کښې هره په محاسبوي ډول بديهي‌کړې تيوري ده چې
د حساب ژبه غځوي، او~\LR{$\Th{T}$} د~\LR{$\Th{Q}$} بديهي اصول رانغاړي۔ دا په
ځانګړي ډول دا معنا لري چې~\LR{$\Th{T}$} محاسبه کېدونکې تابعې او اړيکې
تمثيلوي۔
\olpseganchor{OLP-0315-B006}{end}

\olpseganchor{OLP-0315-B007}{start}
موږ استدلال کړے دے چې که فورمولونه او ثبوتونه په مناسب ډول د عددونو
په توګه کوډ شي، نو اړيکه~\LR{$\Prf[\Th{T}](x,y)$} محاسبه کېدونکې ده۔
\LR{$\Prf[\Th{T}](x,y)$} هله او يوازې هله برقرارېږي چې~\LR{$x$} په~\LR{$\Th{T}$}
کښې د هغه فارمول د اشتقاق ګوډل شمېره وي چې ګوډل شمېره
ئې~\LR{$y$} ده۔ په حقيقت کښې ګوډل د خپلې ټاکلې تيورۍ دپاره وښودله چې
دا اړيکه ابتدايي بازګشتي ده؛ د دې دپاره ئې په خپله مقاله کښې د
۴۵ تابعو او اړيکو لړ وکاراوه۔ پينځه څلوېښتمه اړيکه، \LR{$x B y$}، د
ګوډل د~\LR{$\Th{T}$} د هماغې ټاکنې دپاره همدا~\LR{$\Prf[\Th{T}](x,y)$} ده۔
ياد ولرئ چې ګوډل په خپله مقاله کښې «بازګشتي» وايي، خو موږ به اوس
ورته «ابتدايي بازګشتي» وايو۔
\olpseganchor{OLP-0315-B007}{end}

\olpseganchor{OLP-0315-B008}{start}
څرنګه چې~\LR{$\Prf[\Th{T}](x,y)$} محاسبه کېدونکې ده، په~\LR{$\Th{T}$} کښې
تمثيلېدے شي۔ د هغې تمثيلوونکي فورمول دپاره~\LR{$\OPrf[\Th{T}](x,y)$}
کاروو۔ \LR{$\OProv[\Th{T}](y)$} دې فورمول
\LR{$\lexists[x][\OPrf[\Th{T}](x,y)]$} وي۔ دا د ګوډل په لړ کښې
شپږڅلوېښتمه اړيکه، \LR{$\fn{Bew}(y)$}، بيانوي۔ لکه څنګه چې ګوډل يادونه
کوي، دا يوازينۍ اړيکه ده چې «بازګشتي نۀ شو بللے»۔ ښايي د هغه مطلب
دا و: له تعريفه ئې محاسبه کېدنه نۀ څرګندېږي؛ او وروستيو پرمختګونو
په حقيقت کښې وښودله چې محاسبه کېدونکې هم نۀ ده۔
\olpseganchor{OLP-0315-B008}{end}

\olpseganchor{OLP-0315-B009}{start}
فرض کړئ \LR{$\Th{T}$} داسې په محاسبوي ډول بديهي کېدونکې تيوري ده چې~\LR{$\Th{Q}$}
رانغاړي۔ نو~\LR{$\Prf[\Th{T}](x, y)$} د پرېکړې وړ ده، او ځکه په~\LR{$\Th{Q}$}
کښې د يوه فارمول~\LR{$\OPrf[\Th{T}](x, y)$} په وسيله تمثيلېږي۔
\LR{$\OProv[\Th{T}](y)$} دې هغه فورمول وي چې پورته مو بيان کړ۔ د ثابت
ټکي د لمې له مخې داسې فورمول~\LR{$!G_\Th{T}$} شته چې~\LR{$\Th{Q}$}، او ځکه
\LR{$\Th{T}$} هم، دا اشتقاق کوي:
\begin{equation}
\ollabel{eqn:qpf}
!G_\Th{T} \liff \lnot \OProv[\Th{T}](\gn{!G_\Th{T}}).
\end{equation}
پام وکړئ چې~\LR{$!G_\Th{T}$} په اصل کښې وايي: «\LR{$!G_\Th{T}$} په~\LR{$\Th{T}$}
کښې د اشتقاق وړ نۀ دے»۔
\olpseganchor{OLP-0315-B009}{end}

\olpseganchor{OLP-0315-B010}{start}
\begin{lem}\ollabel{lem:cons-G-unprov}
که~\LR{$\Th{T}$} سازګاره او په محاسبوي ډول بديهي کېدونکې تيوري وي چې~\LR{$\Th{Q}$}
غځوي، نو \LR{$\Th{T} \Proves/ !G_\Th{T}$}۔
\end{lem}
\olpseganchor{OLP-0315-B010}{end}

\olpseganchor{OLP-0315-B011}{start}
\begin{proof}
فرض کړئ \LR{$\Th{T}$}، \LR{$!G_\Th{T}$} اشتقاق کوي۔ نو د هغه
اشتقاق \emph{شته} دے، او په دې توګه د کوم عدد~\LR{$m$} دپاره
اړيکه~\LR{$\Prf[\Th{T}](m,
\Gn{!G_\Th{T}})$} برقرارېږي۔ خو بيا~\LR{$\Th{Q}$} جمله
\LR{$\OPrf[\Th{T}](\num m, \gn{!G_\Th{T}})$} اشتقاق کوي۔ نو~\LR{$\Th{Q}$}
\LR{$\lexists[x][\OPrf[\Th{T}](x,\gn{!G_\Th{T}})]$} هم اشتقاق کوي،
چې د تعريف له مخې~\LR{$\OProv[\Th{T}](\gn{!G_\Th{T}})$} دے۔ د
\olref{eqn:qpf} له مخې~\LR{$\Th{Q}$}، \LR{$\lnot !G_\Th{T}$} اشتقاق کوي؛
او ځکه چې~\LR{$\Th{T}$}، \LR{$\Th{Q}$} غځوي، \LR{$\Th{T}$} ئې هم ثابتوي۔
موږ وښودله چې که~\LR{$\Th{T}$}، \LR{$!G_\Th{T}$} اشتقاق کړي، نو د هغه
نفي~\LR{$\lnot !G_\Th{T}$} هم اشتقاق کوي؛ په دې صورت کښې به تيوري
ناسازګاره وي۔
\end{proof}
\olpseganchor{OLP-0315-B011}{end}

\olpseganchor{OLP-0315-B012}{start}
\begin{defn}
\ollabel{thm:oconsis-q}
تيوري~\LR{$\Th{T}$} \emph{\LR{$\omega$}-سازګاره} ده که لاندې شرط پوره کړي:
که~\LR{$\lexists[x][!A(x)]$} هره جمله وي او~\LR{$\Th{T}$} د
\LR{$\lnot !A(\num 0)$}، \LR{$\lnot !A(\num 1)$}، \LR{$\lnot !A(\num 2)$}،
\dots هره يوه اشتقاق کړي، نو~\LR{$\Th{T}$}،
\LR{$\lexists[x][!A(x)]$} نۀ ثابتوي۔
\end{defn}
\olpseganchor{OLP-0315-B012}{end}

\olpseganchor{OLP-0315-B013}{start}
پام وکړئ چې هره \LR{$\omega$}-سازګاره تيوري سازګاره هم ده۔ دا نېغ له
دې خبرې راځي چې که~\LR{$\Th{T}$} ناسازګاره وي، نو د هر~\LR{$!A$} دپاره
\LR{$\Th{T} \Proves !A$}۔ په ځانګړي ډول که~\LR{$\Th{T}$} ناسازګاره وي، د
هر~\LR{$n$} دپاره \LR{$\lnot !A(\num n)$} اشتقاق کوي او ورسره
\LR{$\lexists[x][!A(x)]$} هم اشتقاق کوي۔ نو ناسازګاره~\LR{$\Th{T}$}
\LR{$\omega$}-ناسازګاره ده۔ د متقابل عکس له مخې، که~\LR{$\Th{T}$}
\LR{$\omega$}-سازګاره وي، نو بايد سازګاره وي۔
\olpseganchor{OLP-0315-B013}{end}

\olpseganchor{OLP-0315-B014}{start}
\begin{lem}\ollabel{lem:omega-cons-G-unref}
که~\LR{$\Th{T}$}، \LR{$\omega$}-سازګاره او په محاسبوي ډول بديهي کېدونکې تيوري وي چې
\LR{$\Th{Q}$} غځوي، نو \LR{$\Th{T} \Proves/ \lnot !G_\Th{T}$}۔
\end{lem}
\olpseganchor{OLP-0315-B014}{end}

\olpseganchor{OLP-0315-B015}{start}
\begin{proof}
ښيو چې که~\LR{$\Th{T}$}، \LR{$\lnot !G_\Th{T}$} اشتقاق کړي، نو
\LR{$\omega$}-ناسازګاره ده۔ فرض کړئ~\LR{$\Th{T}$}، \LR{$\lnot !G_\Th{T}$}
اشتقاق کوي۔ که~\LR{$\Th{T}$} ناسازګاره وي، نو \LR{$\omega$}-ناسازګاره
ده، او ثبوت بشپړ دے۔ که نۀ، نو~\LR{$\Th{T}$} سازګاره ده او د
\olref{lem:cons-G-unprov} له مخې~\LR{$!G_\Th{T}$} نۀ اشتقاق کوي۔
څرنګه چې په~\LR{$\Th{T}$} کښې د~\LR{$!G_\Th{T}$} هېڅ اشتقاق نشته،
نو~\LR{$\Th{Q}$} دا اشتقاق کوي:
\[
\lnot \OPrf[\Th{T}](\num 0, \gn{!G_\Th{T}}), \lnot \OPrf[\Th{T}](\num 1,
\gn{!G_\Th{T}}), \lnot \OPrf[\Th{T}](\num 2, \gn{!G_\Th{T}}), \dots
\]
او~\LR{$\Th{T}$} هم همداسې کوي۔ له بلې خوا د~\olref{eqn:qpf} له مخې
\LR{$\lnot !G_\Th{T}$} له
\LR{$\lexists[x][\OPrf[\Th{T}](x,\gn{!G_\Th{T}})]$} سره معادل دے۔
نو~\LR{$\Th{T}$}، \LR{$\omega$}-ناسازګاره ده۔
\end{proof}
\olpseganchor{OLP-0315-B015}{end}

\olpseganchor{OLP-0315-B016}{start}
\begin{prob}
  هره \LR{$\omega$}-سازګاره تيوري سازګاره ده۔ وښيئ چې برعکس نۀ ده؛
  يعنې داسې تيورۍ شته چې سازګارې خو \LR{$\omega$}-ناسازګارې دي۔ دا په
  ښودلو ثابت کړئ چې \LR{$\Th{Q} \cup
  \{\lnot !G_\Th{Q}\}$} سازګاره خو \LR{$\omega$}-ناسازګاره ده۔
\end{prob}
\olpseganchor{OLP-0315-B016}{end}

\olpseganchor{OLP-0315-B017}{start}
\begin{thm}
\ollabel{thm:first-incompleteness} فرض کړئ \LR{$\Th{T}$} هره
\LR{$\omega$}-سازګاره، په محاسبوي ډول بديهي کېدونکې تيوري ده چې~\LR{$\Th{Q}$} غځوي۔ نو
\LR{$\Th{T}$} بشپړه نۀ ده۔
\end{thm}
\olpseganchor{OLP-0315-B017}{end}

\olpseganchor{OLP-0315-B018}{start}
\begin{proof}
  که~\LR{$\Th{T}$}، \LR{$\omega$}-سازګاره وي، نو سازګاره ده؛ ځکه د
  \olref{lem:cons-G-unprov} له مخې \LR{$\Th{T}
  \Proves/ !G_\Th{T}$}۔ د~\olref{lem:omega-cons-G-unref} له مخې
  \LR{$\Th{T} \Proves/ \lnot !G_\Th{T}$}۔ دا معنا لري چې~\LR{$\Th{T}$}
  ناتکميله ده، ځکه نۀ~\LR{$!G_\Th{T}$} او نۀ~\LR{$\lnot !G_\Th{T}$}
  اشتقاق کوي۔
\end{proof}
\olpseganchor{OLP-0315-B018}{end}

\olpseganchor{OLP-0316-B004}{start}
\olfileid{inc}{inp}{ros}
\olpseganchor{OLP-0316-B004}{end}

\olpseganchor{OLP-0316-B005}{start}
\olsection{د روسر قضيه}
\olpseganchor{OLP-0316-B005}{end}

\olpseganchor{OLP-0316-B006}{start}
ايا د ګوډل ثبوت داسې بدلولے شو چې د «\LR{$\omega$}-سازګارې» پر ځاے
يوازې «سازګاره» تيوري فرض کړو او پياوړې پايله واخلو؟ ځواب «هو» دے؛
د روسر موندلې چاره دا کار کوي۔ د روسر چاره دا ده چې د
\LR{$\OProv[\Th{T}](y)$} پر ځاے د اشتقاق وړتيا يو «بدل شوے»
پريديکات~\LR{$\ORProv_T(y)$} وکاروو۔
\olpseganchor{OLP-0316-B006}{end}

\olpseganchor{OLP-0316-B007}{start}
\begin{thm}
\ollabel{thm:rosser}
فرض کړئ \LR{$\Th{T}$} هره سازګاره او په محاسبوي ډول بديهي کېدونکې تيوري ده چې
\LR{$\Th{Q}$} غځوي۔ نو~\LR{$\Th{T}$} بشپړه نۀ ده۔
\end{thm}
\olpseganchor{OLP-0316-B007}{end}

\olpseganchor{OLP-0316-B008}{start}
\begin{proof}
ياد ولرئ چې \LR{$\OProv[\Th{T}](y)$} د
\LR{$\lexists[x][\OPrf[\Th{T}](x,
  y)]$} په توګه تعريف شوے دے۔ دلته~\LR{$\OPrf[\Th{T}](x, y)$} هغه
د پرېکړې وړ اړيکه تمثيلوي چې هله او يوازې هله برقرارېږي چې~\LR{$x$}
د ګوډل شمېرې~\LR{$y$} لرونکې تړلے فارمول د اشتقاق ګوډل
شمېره وي۔ هغه اړيکه هم د پرېکړې وړ ده چې د~\LR{$x$} او~\LR{$y$} تر منځ
هله برقرارېږي چې~\LR{$x$} د ګوډل شمېرې~\LR{$y$} لرونکې جملې د
\emph{ردونې} ګوډل شمېره وي۔ \LR{$\fn{not}(x)$} دې هغه ابتدايي
بازګشتي تابعه وي چې داسې کار کوي: که~\LR{$x$} د يوه فورمول~\LR{$!A$} کوډ
وي، نو~\LR{$\fn{not}(x)$} د~\LR{$\lnot
!A$} کوډ وي۔ په دې توګه \LR{$\Refut[\Th{T}](x, y)$} هله او يوازې هله
برقرارېږي چې \LR{$\Prf[\Th{T}](x, \fn{not}(y))$} برقرار وي۔
\LR{$\ORefut[\Th{T}](x, y)$} دې دا اړيکه تمثيل کړي۔ نو که
\LR{$\Th{T} \Proves \lnot !A$} وي او~\LR{$\delta$} ئې اړوند اشتقاق
وي، \LR{$\Th{Q} \Proves
\ORefut[\Th{T}](\gn{\delta}, \gn{!A})$}۔ موږ
\LR{$\ORProv[\Th{T}](y)$} داسې تعريفوو:
\[
\lexists[x][(\OPrf[\Th{T}](x,y) \land \lforall[z][(z < x \lif \lnot
  \ORefut[\Th{T}](z,y))])].
\]
په تقريبي ډول~\LR{$\ORProv[\Th{T}](y)$} وايي: «په~\LR{$\Th{T}$} کښې د~\LR{$y$}
ثبوت شته، او د~\LR{$y$} تر دې لنډه ردونه نشته»۔ که~\LR{$\Th{T}$} سازګاره
فرض کړو، نو~\LR{$\ORProv[\Th{T}](y)$} د هماغو عددونو دپاره رښتيا دے
چې~\LR{$\OProv[\Th{T}](y)$} ئې دپاره رښتيا دے؛ خو په~\LR{$\Th{T}$} کښې د
\emph{اثبات‌وړتيا} له پلوه (او اوس پوهېږو چې د رښتياوالي او
اثبات‌وړتيا تر منځ توپير شته!) دواړه بېل خاصيتونه لري۔ که
\LR{$\Th{T}$} \emph{نا}سازګاره وي، نو دواړه د هماغو عددونو دپاره
\emph{نۀ} برقرارېږي!{} (\LR{$\ORProv[\Th{T}](y)$} ته ډېر ځله داسې
لوستل کېږي: «\LR{$y$} د روسر له مخې اثبات‌وړ دے»۔ خو لکه چې مو همدا
اوس وليدل، د روسر اثبات‌وړتيا د اثبات‌وړتيا کوم ځانګړے ډول نۀ
دے: په ناسازګارو تيوريو کښې داسې تړلي فارمولونه شته چې اثبات‌وړ
دي، خو د روسر له مخې اثبات‌وړ نۀ دي۔ دا نوم لوستونکے ګډوډولے شي؛
د ګډوډۍ د مخنيوي دپاره ورته يو بېل مصنوعي نوم، لکه «\LR{$y$} شمووړ دے»، هم
ويلے شئ۔)
\olpseganchor{OLP-0316-B008}{end}

\olpseganchor{OLP-0316-B009}{start}
د ثابت ټکي د لمې له مخې داسې فورمول~\LR{$!R_\Th{T}$} شته چې
\begin{equation}
  \Th{Q} \Proves !R_\Th{T} \liff \lnot \ORProv[\Th{T}](\gn{!R_\Th{T}}).
  \ollabel{RT}
\end{equation}
د~\olref[1in]{thm:first-incompleteness} د ثبوت برعکس، دلته دعوه
دا ده چې که~\LR{$\Th{T}$} سازګاره وي، نو~\LR{$\Th{T}$}، \LR{$!R_\Th{T}$}
نۀ اشتقاق کوي، او~\LR{$\Th{T}$}، \LR{$\lnot !R_\Th{T}$} هم نۀ اشتقاق کوي۔
(په نورو ټکو، د \LR{$\omega$}-سازګارۍ فرض ته اړتيا نشته۔)
\olpseganchor{OLP-0316-B009}{end}

\olpseganchor{OLP-0316-B010}{start}
لومړے ښيو چې \LR{$\Th{T} \Proves/ !R_{\Th{T}}$}۔ که ئې ثابت کړے
واے، نو له~\LR{$T$} څخه به د~\LR{$!R_{\Th{T}}$} يو اشتقاق و؛ د
هغې ګوډل شمېره دې~\LR{$n$} وي۔ بيا
\LR{$\Th{Q} \Proves \OPrf[\Th{T}](\num{n}, \gn{!R_{\Th{T}}})$}، ځکه چې
\LR{$\OPrf[\Th{T}]$} په~\LR{$\Th{Q}$} کښې \LR{$\Prf[\Th{T}]$} تمثيلوي۔ همدا
راز، د هر~\LR{$k < n$} دپاره~\LR{$k$} د~\LR{$\lnot !R_{\Th{T}}$} د يوه
اشتقاق ګوډل شمېره نۀ ده، ځکه~\LR{$\Th{T}$} سازګاره ده۔ نو د
هر~\LR{$k < n$} دپاره
\LR{$\Th{Q} \Proves \lnot
\ORefut[\Th{T}](\num{k}, \gn{!R_{\Th{T}}})$}۔ د
\olref[req][min]{lem:less-nsucc} له مخې
\LR{$\Th{Q} \Proves \lforall[z][(z < \num{n} \lif \lnot \ORefut[\Th{T}](z,
  \gn{!R_{\Th{T}}}))]$}۔ نو
\[
\Th{Q} \Proves \lexists[x][(\OPrf[\Th{T}](x,\gn{!R_{\Th{T}}}) \land \lforall[z][(z
    < x \lif \lnot \ORefut[\Th{T}](z,\gn{!R_{\Th{T}}}))])],
\]
خو دا هماغه~\LR{$\ORProv[\Th{T}](\gn{!R_{\Th{T}}})$} دے۔ د~\olref{RT}
له مخې \LR{$\Th{Q}
\Proves \lnot !R_{\Th{T}}$}۔ ځکه چې~\LR{$\Th{T}$}، \LR{$\Th{Q}$} غځوي،
نو \LR{$\Th{T}
\Proves \lnot !R_{\Th{T}}$} هم دے۔ موږ فرض کړې وه چې
\LR{$\Th{T} \Proves !R_{\Th{T}}$}؛ نو~\LR{$\Th{T}$} به ناسازګاره شي، چې د
قضيې له فرض سره ټکر لري۔
\olpseganchor{OLP-0316-B010}{end}

\olpseganchor{OLP-0316-B011}{start}
اوس ښيو چې \LR{$\Th{T} \Proves/ \lnot !R_{\Th{T}}$}۔ بيا فرض کړئ چې
دا نفي ثابتوي او~\LR{$n$} د~\LR{$\lnot !R_{\Th{T}}$} د يوه
اشتقاق ګوډل شمېره ده۔ نو
\LR{$\Refut[\Th{T}](n, \Gn{!R_{\Th{T}}})$} برقرارېږي۔ ځکه چې
\LR{$\ORefut[\Th{T}]$} په~\LR{$\Th{Q}$} کښې \LR{$\Refut[\Th{T}]$} تمثيلوي،
نو \LR{$\Th{Q} \Proves
\ORefut[\Th{T}](\num{n}, \gn{!R_{\Th{T}}})$}۔ بيا به وښيو چې
\LR{$\Th{T}$} ناسازګاره کېږي، ځکه~\LR{$!R_{\Th{T}}$} به هم اشتقاق کوي۔
څرنګه چې
\begin{align*}
\Th{Q} & \Proves !R_{\Th{T}} \liff \lnot \ORProv[\Th{T}](\gn{!R_{\Th{T}}}),
\intertext{\RL{او ځکه چې \LR{$\Th{T}$}، \LR{$\Th{Q}$} غځوي، بس دا وښيو چې}}
\Th{Q} & \Proves \lnot \ORProv[\Th{T}](\gn{!R_{\Th{T}}}).
\end{align*}
تړلے فارمول~\LR{$\lnot
\ORProv[\Th{T}](\gn{!R_{\Th{T}}})$}، يعنې
\begin{align*}
  \lnot & \lexists[x][(\OPrf[\Th{T}](x,\gn{!R_{\Th{T}}}) \land
    \lforall[z][(z < x \lif \lnot \ORefut[\Th{T}](z,\gn{!R_{\Th{T}}}))])],
  \intertext{\RL{په منطقي ډول له دې سره معادله ده:}}
  & \lforall[x][(\OPrf[\Th{T}](x,\gn{!R_{\Th{T}}}) \lif
    \lexists[z][(z < x \land \ORefut[\Th{T}](z,\gn{!R_{\Th{T}}}))])].
\end{align*}
موږ د منطق او په~\LR{$\Th{Q}$} کښې د اشتقاق کېدونکو حقايقو په
کارولو غير صوري استدلال کوو۔ فرض کړئ~\LR{$x$} خپل‌سرے دے او
\LR{$\OPrf[\Th{T}](x,
\gn{!R_{\Th{T}}})$} برقرارېږي۔ موږ لا پوهېږو چې
\LR{$\Th{T} \Proves/ !R_{\Th{T}}$}، نو د هر~\LR{$k$} دپاره
\LR{$\Th{Q} \Proves \lnot \OPrf[\Th{T}](\num{k}, \gn{!R_{\Th{T}}})$}۔
له دې د هر~\LR{$k$} دپاره \LR{$\eq/[x][\num{k}]$} راځي۔ په ځانګړي ډول (الف)
\LR{$\eq/[x][\num{n}]$} لرو۔ همدا راز
\LR{$\lnot(\eq[x][\num{0}]
\lor \eq[x][\num{1}] \lor \dots \lor \eq[x][\num{n-1}])$} لرو،
نو د~\olref[req][min]{lem:less-nsucc} له مخې (ب)
\LR{$\lnot(x < \num{n})$}۔ د~\olref[req][min]{lem:trichotomy} له مخې
\LR{$\num{n} < x$}۔ ځکه چې
\LR{$\Th{Q} \Proves
\ORefut[\Th{T}](\num{n}, \gn{!R_{\Th{T}}})$}، لرو چې
\LR{$\num{n} < x \land
\ORefut[\Th{T}](\num{n}, \gn{!R_{\Th{T}}})$}؛ له دې څخه
\LR{$\lexists[z][(z < x
\land \ORefut[\Th{T}](z,\gn{!R_{\Th{T}}}))]$} راځي۔ ځکه چې~\LR{$x$}
خپل‌سرے و، اړينه پايله دا ده:
\[
\lforall[x][(\OPrf[\Th{T}](x,\gn{!R_{\Th{T}}}) \lif
  \lexists[z][(z < x \land \ORefut[\Th{T}](z,\gn{!R_{\Th{T}}}))])].
\]
\end{proof}
\olpseganchor{OLP-0316-B011}{end}

\olpseganchor{OLP-0316-B012}{start}
\begin{prob}
د طبيعي عددونو دوه سټونه~\LR{$A$} او~\LR{$B$} هغه وخت
\emph{په محاسبوي ډول نه بېلېدونکي} بلل کېږي چې هېڅ د پرېکړې وړ
سټ~\LR{$X$} داسې نۀ وي چې \LR{$A \subseteq X$} او
\LR{$B \subseteq \Complement{X}$} وي (\LR{$\Complement{X}$} د~\LR{$X$}
متمم، يعنې~\LR{$\Nat \setminus X$}، دے)۔ فرض کړئ~\LR{$\Th{T}$} د~\LR{$\Th{Q}$}
سازګاره او په محاسبوي ډول بديهي کېدونکې غځونه ده۔ \LR{$A$} دې په~\LR{$\Th{T}$} کښې د
اثبات‌وړو تړلي فارمولونه د ګوډل شمېرو سټ وي، او~\LR{$B$} دې په~\LR{$\Th{T}$}
کښې د ردېدونکو جملو د ګوډل شمېرو سټ وي۔ ثابت کړئ چې~\LR{$A$} او~\LR{$B$}
په محاسبوي ډول نه بېلېدونکي دي۔
\end{prob}
\olpseganchor{OLP-0316-B012}{end}

\olpseganchor{OLP-0317-B004}{start}
\olfileid{inc}{inp}{gop}
\olpseganchor{OLP-0317-B004}{end}

\olpseganchor{OLP-0317-B005}{start}
\olsection{د ګوډل له اصلي مقالې سره پرتله}
\olpseganchor{OLP-0317-B005}{end}

\olpseganchor{OLP-0317-B006}{start}
د ګوډل د ۱۹۳۱ز کال مقالې ته لږ وخت ورکول ګټور دي۔ د هغې پېژندنه
د هغو مفکورو لنډ بيان کوي چې موږ همدا اوس وڅېړلې۔ که مقاله يوازې
په تېره هم وګورئ، بيا هم د هر پړاو کار څرګندېږي: لومړے ګوډل صوري
نظام~\LR{$P$} بيانوي (نحو، بديهي اصول او د ثبوت قاعدې)؛ بيا ابتدايي
بازګشتي تابعې او اړيکې تعريفوي؛ ورپسې ښيي چې~\LR{$x B y$} ابتدايي
بازګشتي ده او استدلال کوي چې ابتدايي بازګشتي تابعې او اړيکې
په~\LR{$\Th{P}$} کښې تمثيلېږي۔ بيا د ناتکميلۍ قضيه لکه پورته ثابتوي۔
په درېيمه برخه کښې ښيي چې ناثابتېدونکې دعوه د حساب د ژبې يوه جمله
اخيستے شو۔ د~\LR{$\beta$}-لمې اصل همدلته دے؛ موږ هم دا لمه د لړيو
د سمبالولو دپاره وکاروله، کله چې مو ښودله چې بازګشتي تابعې
په~\LR{$\Th{Q}$} کښې تمثيلېږي۔ ګوډل دومره وړاندې نۀ ځي چې د بسيا
بديهي اصولو يو لږترين سټ بېل کړي، خو اوس پوهېږو چې~\LR{$\Th{Q}$} د
دې کار دپاره بسيا ده۔ په پاې کښې، په څلورمه برخه کښې د ناتکميلۍ
د دويمې قضيې د ثبوت لنډه نقشه ورکوي۔
\olpseganchor{OLP-0317-B006}{end}

\olpseganchor{OLP-0318-B004}{start}
\olfileid{inc}{inp}{prc}
\olpseganchor{OLP-0318-B004}{end}

\olpseganchor{OLP-0318-B005}{start}
\olsection{د~\texorpdfstring{\LR{$\Th{PA}$}}{PA} د اشتقاق وړتيا شرطونه}
\olpseganchor{OLP-0318-B005}{end}

\olpseganchor{OLP-0318-B006}{start}
د پېانو حساب، يا~\LR{$\Th{PA}$}، هغه تيوري ده چې~\LR{$\Th{Q}$} د ټولو
فارمولونه دپاره د استقرا د بديهي اصولو په زياتولو غځوي۔ په نورو
ټکو، \LR{$\Th{Q}$} ته د دې بڼې بديهي اصول ورزياتېږي:
\[
(!A(0) \land \lforall[x][(!A(x) \lif !A(x'))]) \lif \lforall[x][!A(x)]
\]
د هر فارمول~\LR{$!A$} دپاره۔ پام وکړئ چې دا په حقيقت کښې يوه
\emph{شېما} ده، يعنې نامتناهي ډېر بديهي اصول دي (او دا هم څرګندېږي
چې~\LR{$\Th{PA}$} په متناهي ډول {\em نۀ} شي بديهي کېدے)۔ خو ځکه چې
په مؤثره طريقه معلومولے شو چې د نښو يوه لړۍ د استقرا د کوم بديهي
اصل بېلګه ده که نۀ، د~\LR{$\Th{PA}$} د بديهي اصولو سټ محاسبه کېدونکے
دے۔ \LR{$\Th{PA}$} له~\LR{$\Th{Q}$} څخه ډېره پياوړې تيوري ده۔ مثلاً د
عادي استقرا په کارولو په آسانۍ ثابتولے شو چې جمع او ضرب تبادلي
دي۔ په حقيقت کښې د عددونو د تيورۍ او ترکيبپوهنې زياتره
متناهي‌پاله استدلالونه په~\LR{$\Th{PA}$} کښې ترسره کېدے شي۔
\olpseganchor{OLP-0318-B006}{end}

\olpseganchor{OLP-0318-B007}{start}
څرنګه چې~\LR{$\Th{PA}$} په محاسبوي ډول بديهي‌کړې ده، د
د اشتقاق وړتيا پريديکات~\LR{$\Prf[\Th{PA}](x,y)$} محاسبه کېدونکے
دے او ځکه په~\LR{$\Th{Q}$} کښې، او په دې توګه په~\LR{$\Th{PA}$} کښې هم،
تمثيلېږي۔ د پخوا په شان د دې اړيکې د تمثيلوونکي فورمول دپاره
\LR{$\OPrf[\Th{PA}](x,y)$} کاروو۔ \LR{$\OProv[\Th{PA}](y)$} دې فورمول
\LR{$\lexists[x][\OPrf[\Th{PA}](x,y)]$} وي؛ دا په مفهومي ډول وايي:
«\LR{$y$} د~\LR{$\Th{PA}$} له بديهي اصولو څخه د اشتقاق وړ دے»۔
د عددونو تر منځ اړيکه~\LR{$\Prf[\Th{PA}](x,y)$} په څيز ژبه کښې د~\LR{$\OPrf[\Th{PA}](x,y)$} په وسيله تمثيلېږي۔ له~\LR{$\Th{Q}$} د بديهي اصولو څخه لږ څه زيات ته ځکه اړ يو
چې خپله تيوري بايد دومره پياوړې وي چې د دې د اشتقاق وړتيا
پريديکات په باب څو بنسټيز حقايق اشتقاق کړي۔ اړين حقايق دا دي:
\begin{enumerate}
\item[P1.] که \LR{$\Th{PA} \Proves !A$} وي، نو \LR{$\Th{PA} \Proves
  \OProv[\Th{PA}](\gn{!A})$}۔
\item[P2.] د ټولو فارمولونه \LR{$!A$} او \LR{$!B$} دپاره:
  \[
  \Th{PA} \Proves \OProv[\Th{PA}](\gn{!A \lif !B}) \lif
  (\OProv[\Th{PA}](\gn{!A}) \lif \OProv[\Th{PA}](\gn{!B})).
  \]
\item[P3.] د هر فارمول~\LR{$!A$} دپاره:
  \[
  \Th{PA} \Proves \OProv[\Th{PA}](\gn{!A})
  \lif \OProv[\Th{PA}](\gn{\OProv[\Th{PA}](\gn{!A})}).
  \]
\end{enumerate}
دا درې خاصيتونه يوازې په دې توګه کتلے شو چې فارمول
\LR{$\OProv[\Th{PA}](y)$} په دقت بيان کړو او د~\LR{$\Th{PA}$} بديهي اصول
د اړوندو صوري اشتقاقونه د بيانولو دپاره وکاروو۔ شرطونه (۱)
او~(۲) آسانه دي؛ اصلي کار شرط~(۳) غواړي۔ (فکر وکړئ چې دا به
څه ډول کار وي \dots) د تفصيلاتو بشپړول ستړي کوونکي او زړه راکښونکي
نۀ وي، نو اوس له تاسو غواړو چې دا ومنئ چې~\LR{$\Th{PA}$} همدا
درې پورته خاصيتونه لري۔ د~\LR{$\OProv[\Th{PA}](y)$} يوه مناسبه ټاکنه
به دا هم پوره کوي:
\begin{enumerate}
\item[P4.] که \LR{$\Th{PA} \Proves \OProv[\Th{PA}](\gn{!A})$} وي، نو
  \LR{$\Th{PA} \Proves !A$}۔
\end{enumerate}
خو موږ به دې حقيقت ته اړ نۀ شو۔
\olpseganchor{OLP-0318-B007}{end}

\olpseganchor{OLP-0318-B008}{start}
\begin{digress}
په ضمن کښې ګوډل هم دلته زموږ په شان د تفصيلاتو په ليکلو کښې
تنبلي وکړه۔ د خپلې ۱۹۳۱ز مقالې په پاې کښې ئې د ناتکميلۍ د دويمې
قضيې د ثبوت لنډه نقشه ورکړه او ژمنه ئې وکړه چې تفصيلات به په بله
مقاله کښې راوړي۔ بيا ورته وخت ورنۀ غے؛ ځکه هر څوک چې په استدلال
پوهېدل، باور ئې درلود چې تفصيلات بشپړېدے شي (او هغه اړ نۀ و چې
پخپله ئې وليکي۔)
\end{digress}
\olpseganchor{OLP-0318-B008}{end}

\olpseganchor{OLP-0319-B004}{start}
\olfileid{inc}{inp}{2in}
\olpseganchor{OLP-0319-B004}{end}

\olpseganchor{OLP-0319-B005}{start}
\olsection{د ناتکميلۍ دويمه قضيه}
\olpseganchor{OLP-0319-B005}{end}

\olpseganchor{OLP-0319-B006}{start}
دا دعوه څنګه بيان کړو چې~\LR{$\Th{PA}$} خپله سازګاري نۀ شي ثابتولے؟
دا ويل چې~\LR{$\Th{PA}$} ناسازګاره ده، د دې ويلو برابر دي چې
\LR{$\Th{PA} \Proves \eq[0][1]$}۔ نو د سازګارۍ جمله~\LR{$\OCon[\Th{PA}]$}
د تړلے فارمول~\LR{$\lnot
\OProv[\Th{PA}](\gn{\eq[0][1]})$} په توګه اخيستے شو، او بيا
لاندې قضيه اړينه پايله راکوي:
\olpseganchor{OLP-0319-B006}{end}

\olpseganchor{OLP-0319-B007}{start}
\begin{thm}
\ollabel{thm:second-incompleteness}
که~\LR{$\Th{PA}$} سازګاره وي، نو~\LR{$\Th{PA}$}،
\LR{$\OCon[\Th{PA}]$} نۀ اشتقاق کوي۔
\end{thm}
\olpseganchor{OLP-0319-B007}{end}

\olpseganchor{OLP-0319-B008}{start}
دا خبره مهمه ده چې قضيه د~\LR{$\OCon[\Th{PA}]$} په ټاکلي تمثيل پورې
تړلې ده، يعنې د~\LR{$\OProv[\Th{PA}](y)$} په ټاکلي تمثيل پورې۔ موږ
يوازې دومره کاروو چې د~\LR{$\OProv[\Th{PA}](y)$} تمثيل د
د اشتقاق وړتيا درې شرطونه پوره کوي؛ نو قضيه هرې هغې تيورۍ ته
غځېږي چې د اشتقاق وړتيا پريديکات ئې دغه خاصيتونه ولري۔
\olpseganchor{OLP-0319-B008}{end}

\olpseganchor{OLP-0319-B009}{start}
د ګوډل د استدلال لنډه نقشه لوستل ګټور دي، ځکه قضيه پکې لکه د ښې
ټوکې اغېزمنه وروستۍ جمله راځي۔ استدلال داسې دے: \LR{$!G_\Th{PA}$}
دې د ګوډل هغه جمله وي چې د~\olref[1in]{thm:first-incompleteness}
په ثبوت کښې مو جوړه کړه۔ موږ وښودله: «که~\LR{$\Th{PA}$} سازګاره وي،
نو~\LR{$\Th{PA}$}، \LR{$!G_\Th{PA}$} نۀ اشتقاق کوي»۔ که دا خبره
پخپله \emph{په}~\LR{$\Th{PA}$} کښې صوري کړو، لاندې ثبوت لرو:
\[
\OCon[\Th{PA}] \lif \lnot \OProv[\Th{PA}](\gn{!G_\Th{PA}}).
\]
اوس فرض کړئ~\LR{$\Th{PA}$}، \LR{$\OCon[\Th{PA}]$} اشتقاق کوي۔ بيا
\LR{$\lnot
\OProv[\Th{PA}](\gn{!G_\Th{PA}})$} هم اشتقاق کوي۔
\LR{$\Prov[\Th{PA}]$} د عددونو بهرنۍ اړيکه او \LR{$\OProv[\Th{PA}]$} د ثبوتېدنې داخلي فورمول دے۔ ځکه چې~\LR{$!G_\Th{PA}$} د ګوډل جمله ده، دا له
\LR{$!G_\Th{PA}$} سره معادله ده۔ نو~\LR{$\Th{PA}$}، \LR{$!G_\Th{PA}$}
اشتقاق کوي۔
\olpseganchor{OLP-0319-B009}{end}

\olpseganchor{OLP-0319-B010}{start}
خو موږ پوهېږو چې که~\LR{$\Th{PA}$} سازګاره وي، \LR{$!G_\Th{PA}$} نۀ
اشتقاق کوي!{} نو که~\LR{$\Th{PA}$} سازګاره وي،
\LR{$\OCon[\Th{PA}]$} هم نۀ اشتقاق کولے شي۔
\olpseganchor{OLP-0319-B010}{end}

\olpseganchor{OLP-0319-B011}{start}
د استدلال د لا کره کولو دپاره~\LR{$!G_\Th{PA}$} د~\LR{$\Th{PA}$} د ګوډل
جمله اخلو او د اشتقاق وړتيا شرطونه (P1)--(P3) کاروو، څو
وښيو چې~\LR{$\Th{PA}$}، \LR{$\OCon[\Th{PA}] \lif
!G_\Th{PA}$} اشتقاق کوي۔ له دې به څرګنده شي چې~\LR{$\Th{PA}$}،
\LR{$\OCon[\Th{PA}]$} نۀ اشتقاق کوي۔ دا په~\LR{$\Th{PA}$} کښې د
ثبوت لنډه نقشه ده۔ (د آسانتيا دپاره د~\LR{$\Th{PA}$} لاندې ليکونه
پرېږدو۔)
\begin{align}
& !G \liff \lnot \OProv(\gn{!G}) \ollabel{G2-1}\\
& \qquad\text{\RL{\LR{$!G$} د ګوډل جمله ده}}\notag \\
& !G \lif \lnot \OProv(\gn{!G}) \ollabel{G2-2}\\
  & \qquad\text{\RL{له \olref{G2-1} څخه}} \notag\\
& !G \lif
  (\OProv(\gn{!G}) \lif \lfalse) \ollabel{G2-3}\\
  & \qquad\text{\RL{له \olref{G2-2} څخه د منطق له مخې}}\notag\\
& \OProv(\gn{
    !G \lif
    (\OProv(\gn{!G}) \lif \lfalse)
  }) \ollabel{G2-4}\\
  & \qquad\text{\RL{له \olref{G2-3} څخه د P1 شرط له مخې}} \notag\\
& \OProv(\gn{!G}) \lif
  \OProv(\gn{
    (\OProv(\gn{!G}) \lif \lfalse)
    }) \ollabel{G2-5}\\
  & \qquad\text{\RL{له \olref{G2-4} څخه د P2 شرط له مخې}} \notag\\
& \OProv(\gn{!G}) \lif (\OProv(\gn{\OProv(\gn{!G})}) \lif \OProv(\gn{\lfalse})) \ollabel{G2-6}\\
  & \qquad\text{\RL{له \olref{G2-5} څخه د P2 شرط او منطق له مخې}} \notag\\
& \OProv(\gn{!G}) \lif
  \OProv(\gn{\OProv(\gn{!G})}) \ollabel{G2-7}\\
   & \qquad\text{\RL{د P3 له مخې}} \notag\\
& \OProv(\gn{!G}) \lif \OProv(\gn{\lfalse}) \ollabel{G2-8}\\
  & \qquad \text{\RL{له \olref{G2-6} او \olref{G2-7} څخه د منطق له مخې}}\notag\\
& \OCon \lif \lnot \OProv(\gn{!G}) \ollabel{G2-9}\\
  & \qquad\text{\RL{د \olref{G2-8} متقابل عکس او \LR{$\OCon \ident \lnot \OProv(\gn{\lfalse})$}}}\notag \\
& \OCon \lif !G \notag\\
  & \qquad\text{\RL{له \olref{G2-1} او \olref{G2-9} څخه د منطق له مخې}}\notag
\end{align}
په پورته ثبوت کښې د منطق کارول يوازې د قضيې د منطق بنسټيز
حقايق دي: مثلاً~\olref{G2-3} له
\LR{$\Proves \lnot!A \liff (!A\lif
\lfalse)$} څخه کار اخلي، او~\olref{G2-8} له
\LR{$!A \lif (!B \lif !C), !A \lif !B
\Proves !A \lif !C$} څخه۔ په~\olref{G2-5} او~\olref{G2-6}
کښې د P2 شرط کارول د هغه پر بېلګو
\LR{$\OProv(\gn{!A \lif !B}) \lif
(\OProv(\gn{!A}) \lif \OProv(\gn{!B}))$} ولاړ دي۔ په لومړۍ
بېلګه کښې \LR{$!A \ident
!G$} او \LR{$!B \ident \OProv(\gn{!G}) \lif \lfalse$} دي؛ په
دويمه کښې \LR{$!A
\ident \OProv(\gn{!G})$} او \LR{$!B \ident \lfalse$} دي۔
\olpseganchor{OLP-0319-B011}{end}

\olpseganchor{OLP-0319-B012}{start}
د ناتکميلۍ د دويمې قضيې لا انتزاعي بڼه دا ده:
\olpseganchor{OLP-0319-B012}{end}

\olpseganchor{OLP-0319-B013}{start}
\begin{thm}
\ollabel{thm:second-incompleteness-gen} فرض کړئ~\LR{$\Th{T}$} هره
سازګاره، بديهي‌کړې تيوري ده چې~\LR{$\Th{Q}$} غځوي، او
\LR{$\OProv[\Th{T}](y)$} هر هغه فورمول دے چې د~\LR{$\Th{T}$} دپاره د
د اشتقاق وړتيا شرطونه P1--P3 پوره کوي۔ نو~\LR{$\Th{T}$}،
\LR{$\OCon[T]$} نۀ اشتقاق کوي۔
\end{thm}
\olpseganchor{OLP-0319-B013}{end}

\olpseganchor{OLP-0319-B014}{start}
\begin{prob}
وښيئ چې~\LR{$\Th{PA}$}، \LR{$!G_{\Th{PA}} \lif \OCon[\Th{PA}]$}
اشتقاق کوي۔
\end{prob}
\olpseganchor{OLP-0319-B014}{end}

\olpseganchor{OLP-0319-B015}{start}
\begin{digress}
د دې کيسې درس دا دے چې د رياضي هېڅ «معقوله» سازګاره تيوري
خپله د سازګارۍ جمله نۀ شي اشتقاق کولے۔ فرض کړئ~\LR{$\Th{T}$}
د رياضي داسې تيوري ده چې~\LR{$\Th{Q}$} او د هېلبرټ «متناهي‌پاله»
استدلال (هر څه چې مطلب ترې وي) رانغاړي۔ نو ټوله~\LR{$\Th{T}$} د
\LR{$\Th{T}$} د سازګارۍ جمله نۀ شي اشتقاق کولے؛ ځکه ئې متناهي‌پاله
برخه هم په لا پياوړي دليل سره د~\LR{$\Th{T}$} د سازګارۍ جمله نۀ شي
اشتقاق کولے۔ په دې معنا د «ټولې رياضي» دپاره د سازګارۍ
متناهي‌پاله ثبوت نشته۔
\olpseganchor{OLP-0319-B015}{end}

\olpseganchor{OLP-0319-B016}{start}
د «متناهي‌پاله» لفظ په تفسير کښې لږ ګنجايش شته۔ ګوډل په خپلې
۱۹۳۱ز مقالې کښې دا شونتيا مني چې هغه څه چې موږ ئې «متناهي‌پاله»
بولو، ښايي د رياضي له هغو ډولونو بهر وي چې هېلبرټ غوښتل صوري
کړي۔ خو ګوډل په دې کښې د نرمۍ څخه کار اخيست؛ نن دا ګرانه ده چې
داسې څه ومومو چې په معقول ډول متناهي‌پاله وبلل شي، خو مثلاً
په~\LR{$\Th{ZFC}$}، يعنې د ټاکنې له بديهي اصل سره د زرملو--فرانکل
د سټونو په تيورۍ کښې، صوري نۀ شي۔
\end{digress}
\olpseganchor{OLP-0319-B016}{end}

\olpseganchor{OLP-0320-B004}{start}
\olfileid{inc}{inp}{lob}
\olpseganchor{OLP-0320-B004}{end}

\olpseganchor{OLP-0320-B005}{start}
\olsection{د لوب قضيه}
\olpseganchor{OLP-0320-B005}{end}

\olpseganchor{OLP-0320-B006}{start}
د يوې تيورۍ~\LR{$\Th{T}$} د ګوډل جمله د~\LR{$\lnot
\OProv[\Th{T}](y)$} يو ثابت ټکی دے؛ يعنې داسې تړلے فارمول~\LR{$!G$}
چې
\[
\Th{T} \Proves \lnot \OProv[\Th{T}](\gn{!G}) \liff !G.
\]
دا جمله د اشتقاق وړ نۀ ده۔ ځکه که~\LR{$\Th{T} \Proves !G$} وي، نو
(الف) د اشتقاق وړتيا د شرط~(۱) له مخې
\LR{$\Th{T} \Proves \OProv[\Th{T}](\gn{!G})$}؛ او (ب)
\LR{$\Th{T}
\Proves !G$} له
\LR{$\Th{T} \Proves \lnot \OProv[\Th{T}](\gn{!G})
\liff !G$} سره يوځاے دا راکوي چې
\LR{$\Th{T} \Proves \lnot \OProv[\Th{T}](\gn{!G})$}۔ په دې صورت کښې
به~\LR{$\Th{T}$} ناسازګاره وي۔ اوس طبيعي پوښتنه دا ده چې د
\LR{$\OProv[\Th{T}](y)$} د ثابت ټکي حالت څه دے؛ يعنې د داسې
تړلے فارمول~\LR{$!H$} حالت چې
\[
\Th{T} \Proves \OProv[\Th{T}](\gn{!H}) \liff !H.
\]
که دا جمله د اشتقاق وړ واے، د شرط~(۱) له مخې به
\LR{$\Th{T} \Proves \OProv[\Th{T}](\gn{!H})$} و؛ خو همدا پايله د
پورته معادلۍ د مودس پوننس په کارولو هم ترلاسه کېږي۔ نو لږ تر لږه
د ګوډل د جملې د پخواني استدلال له لارې د~\LR{$\Th{T}$} ناسازګاري نۀ
راځي۔ دا، البته، نۀ ښيي چې~\LR{$\Th{T}$} په رښتيا~\LR{$!H$}
اشتقاق \emph{کوي}۔
\olpseganchor{OLP-0320-B006}{end}

\olpseganchor{OLP-0320-B007}{start}
که پوښتنه لږ عمومي کړو، پرمختګ کولے شو۔ د ثابت ټکي د معادلۍ
له کيڼ څخه ښي ته لوری،
\LR{$\OProv[\Th{T}](\gn{!H}) \lif !H$}، د يوې عمومي شېما بېلګه ده
چې د \emph{انعکاس اصل} بلل کېږي:
\LR{$\OProv[\Th{T}](\gn{!A}) \lif !A$}۔ دا نوم ځکه ورکول کېږي چې
په يوه معنا ښيي~\LR{$\Th{T}$} د هغو څه په باب «انعکاس» کولے شي چې
اشتقاق کوي؛ په بنسټيز ډول د هر~\LR{$!A$} دپاره وايي: «که~\LR{$\Th{T}$}،
\LR{$!A$} اشتقاق کولے شي، نو~\LR{$!A$} رښتيا دے»۔ دا خبره، البته،
يوازې د سمې تيورۍ دپاره رښتيا ده؛ نو داسې ښکاري چې تيورۍ به
په عمومي توګه د دې شېما ټولې بېلګې نۀ اشتقاق کوي۔ نو يوه
کافي پياوړې تيوري چې د اشتقاق وړتيا شرطونه پوره کوي، کومې
بېلګې اشتقاق کولے شي؟ هغه ټولې بېلګې چې~\LR{$!A$} پخپله
د اشتقاق وړ وي، يقيناً۔ لاندې پايله ښيي چې بس همدومره دي۔
\olpseganchor{OLP-0320-B007}{end}

\olpseganchor{OLP-0320-B008}{start}
\begin{thm}\ollabel{thm:lob}
فرض کړئ~\LR{$\Th{T}$} داسې په محاسبوي ډول بديهي کېدونکې تيوري ده چې~\LR{$\Th{Q}$}
غځوي، او~\LR{$\OProv[\Th{T}](y)$} هغه فورمول دے چې د
\olref[2in]{sec} شرطونه P1--P3 پوره کوي۔ که~\LR{$\Th{T}$}،
\LR{$\OProv[\Th{T}](\gn{!A}) \lif !A$} اشتقاق کړي، نو په حقيقت
کښې~\LR{$\Th{T}$}، \LR{$!A$} هم اشتقاق کوي۔
\end{thm}
\olpseganchor{OLP-0320-B008}{end}

\olpseganchor{OLP-0320-B009}{start}
په بل ډول: که \LR{$\Th{T} \Proves/ !A$} وي، نو
\LR{$\Th{T} \Proves/
\OProv[\Th{T}](\gn{!A}) \lif !A$}۔ دا پايله د لوب د قضيې په نوم
يادېږي۔
\olpseganchor{OLP-0320-B009}{end}

\olpseganchor{OLP-0320-B010}{start}
\begin{explain}
د لوب د قضيې د ثبوت لارښوده مفکوره د دې يو ځيرک «ثبوت» دے چې
سانتا کلاز شته۔ (که دغه پايله مو نۀ خوښېږي، هره بله خوښه پايله
پر ځاے کولے شئ۔) استدلال داسې دے:
\begin{enumerate}
\item \LR{$X$} دې دا جمله وي: «که~\LR{$X$} رښتيا وي، نو سانتا کلاز شته»۔
\item فرض کړئ~\LR{$X$} رښتيا ده۔
\item نو د هغې ويل شوې خبره برقرار ده؛ يعنې که~\LR{$X$} رښتيا وي،
  سانتا کلاز شته۔
\item ځکه چې فرض مو کړے~\LR{$X$} رښتيا ده، له (۲) او~(۳) څخه د
  مودس پوننس په کارولو پايله اخلو چې سانتا کلاز شته۔
\item موږ (۴)، «سانتا کلاز شته»، له (۲)، «\LR{$X$} رښتيا ده»، څخه
  وايستله۔ د شرطي ثبوت له مخې مو وښودله: «که~\LR{$X$} رښتيا وي، نو
  سانتا کلاز شته»۔
\item خو دا هماغه جمله~\LR{$X$} ده۔ نو ومو ښودله چې~\LR{$X$} رښتيا ده۔
\item بيا د پورته (۲)--(۴) استدلال له مخې سانتا کلاز شته۔
\end{enumerate}
که په دې مفکوره کښې «رښتيا ده» په «د اشتقاق وړ ده» او «سانتا
کلاز شته» په~\LR{$!A$} واړوو او هغه صوري کړو، د لوب د قضيې ثبوت
ترې جوړېږي۔ چاره دا ده چې د ثابت ټکي لمه په
فارمول~\LR{$\OProv[\Th{T}](y) \lif !A$} پلې کړو۔ د دې فورمول
ثابت ټکی په پورته لنډ استدلال کښې له~\LR{$X$} جملې سره سمون لري۔
\end{explain}
\olpseganchor{OLP-0320-B010}{end}

\olpseganchor{OLP-0320-B011}{start}
\begin{proof}[د \olref{thm:lob} ثبوت]
فرض کړئ~\LR{$!A$} داسې تړلے فارمول ده چې~\LR{$\Th{T}$}،
\LR{$\OProv[\Th{T}](\gn{!A}) \lif !A$} اشتقاق کوي۔ \LR{$!B(y)$}
دې فارمول~\LR{$\OProv[\Th{T}](y)
\lif !A$} وي، او د ثابت ټکي د لمې په کارولو داسې
تړلے فارمول~\LR{$!D$} ومومئ چې~\LR{$\Th{T}$}،
\LR{$!D \liff !B(\gn{!D})$} اشتقاق کوي۔ بيا لاندې هره يوه
په~\LR{$\Th{T}$} کښې د اشتقاق وړ ده:
\begin{align}
  & !D \liff (\OProv[\Th{T}](\gn{!D}) \lif !A) \ollabel{L-1}\\
  & \qquad \text{\RL{\LR{$!D$} د~\LR{$!B(y)$} ثابت ټکی دے}}\notag \\
  & !D \lif (\OProv[\Th{T}](\gn{!D}) \lif !A) \ollabel{L-2}\\
  & \qquad\text{\RL{له \olref{L-1} څخه}}\notag\\
  & \OProv[\Th{T}](\gn{!D \lif (\OProv[\Th{T}](\gn{!D}) \lif !A)}) \ollabel{L-3}\\
  & \qquad \text{\RL{له \olref{L-2} څخه د P1 شرط له مخې}}\notag \\
  & \OProv[\Th{T}](\gn{!D}) \lif \OProv[\Th{T}](\gn{\OProv[\Th{T}](\gn{!D}) \lif !A})
  \ollabel{L-4}\\
  &\qquad \text{\RL{له \olref{L-3} څخه د P2 شرط له مخې}}\notag \\
  & \OProv[\Th{T}](\gn{!D}) \lif (\OProv[\Th{T}](\gn{\OProv[\Th{T}](\gn{!D})}) \lif \OProv[\Th{T}](\gn{!A})) \ollabel{L-5}\\
  &\qquad \text{\RL{له \olref{L-4} څخه د P2 شرط له مخې بيا}} \notag\\
& \OProv[\Th{T}](\gn{!D}) \lif \OProv[\Th{T}](\gn{\OProv[\Th{T}](\gn{!D})}) \ollabel{L-6}\\
  & \qquad\text{\RL{د اشتقاق وړتيا د P3 شرط له مخې}} \notag\\
  & \OProv[\Th{T}](\gn{!D}) \lif \OProv[\Th{T}](\gn{!A}) \ollabel{L-7} \\
  &\qquad\text{\RL{له \olref{L-5} او \olref{L-6} څخه}}\notag\\
  & \OProv[\Th{T}](\gn{!A}) \lif !A \ollabel{L-8}\\
  &\qquad\text{\RL{د قضيې د فرض له مخې}} \notag\\
  & \OProv[\Th{T}](\gn{!D}) \lif !A \ollabel{L-9}\\
  &\qquad\text{\RL{له \olref{L-7} او \olref{L-8} څخه}}\notag\\
  & (\OProv[\Th{T}](\gn{!D}) \lif !A) \lif !D \ollabel{L-10}\\
  & \qquad \text{\RL{له \olref{L-1} څخه}}\notag \\
  & !D \ollabel{L-11}\\
  & \qquad\text{\RL{له \olref{L-9} او \olref{L-10} څخه}}\notag \\
  & \OProv[\Th{T}](\gn{!D}) \ollabel{L-12}\\
  & \qquad\text{\RL{له \olref{L-11} څخه د P1 شرط له مخې}}\notag \\
  & !A \qquad\qquad\text{\RL{له \olref{L-9} او \olref{L-12} څخه}}\notag
\end{align}
\end{proof}
\olpseganchor{OLP-0320-B011}{end}

\olpseganchor{OLP-0320-B012}{start}
د لوب قضيه چې ولرو، د ناتکميلۍ د دويمې قضيې لنډ ثبوت ترې راځي
(د هغو تيوريو دپاره چې د P1--P3 شرطونه پوره کوونکے
د اشتقاق وړتيا پريديکات لري): که
\LR{$\Th{T} \Proves
\OProv[\Th{T}](\gn{\lfalse}) \lif \lfalse$} وي، نو
\LR{$\Th{T} \Proves \lfalse$}۔ که~\LR{$\Th{T}$} سازګاره وي،
\LR{$\Th{T} \Proves/ \lfalse$}۔ نو
\LR{$\Th{T}
\Proves/ \OProv[\Th{T}](\gn{\lfalse}) \lif \lfalse$}، يعنې
\LR{$\Th{T} \Proves/
\OCon[\Th{T}]$}۔ له دې څخه دا هم ښودلے شو چې~\LR{$!H$}، د
\LR{$\OProv[\Th{T}](x)$} ثابت ټکی، د اشتقاق وړ دے۔ ځکه چې
\begin{align*}
  \Th{T} & \Proves \OProv[\Th{T}](\gn{!H}) \liff !H\\
  \intertext{\RL{په ځانګړي ډول}}
    \Th{T} & \Proves \OProv[\Th{T}](\gn{!H}) \lif !H
\end{align*}
او ځکه د لوب د قضيې له مخې \LR{$\Th{T} \Proves !H$}۔
\olpseganchor{OLP-0320-B012}{end}

\olpseganchor{OLP-0320-B013}{start}
% Going in the other direction, for homework I
% may ask you to work through a short proof of L\"ob's theorem, using
% the second incompleteness theorem instead of the fixed-point lemma.
\olpseganchor{OLP-0320-B013}{end}

\olpseganchor{OLP-0320-B014}{start}
\begin{prob}
فرض کړئ~\LR{$\Th{T}$} په محاسبوي ډول بديهي‌کړې تيوري ده، او
\LR{$\OProv[\Th{T}]$} د~\LR{$\Th{T}$} دپاره د اشتقاق وړتيا پريديکات
دے۔ لاندې څلور بيانونه وګورئ:
\begin{enumerate}
\item که \LR{$T \Proves !A$} وي، نو \LR{$T \Proves \OProv[\Th{T}](\gn{!A})$}۔
\item \LR{$T \Proves !A \lif \OProv[\Th{T}](\gn{!A})$}۔
\item که \LR{$T \Proves \OProv[\Th{T}](\gn{!A})$} وي، نو \LR{$T \Proves !A$}۔
\item \LR{$T \Proves \OProv[\Th{T}](\gn{!A}) \lif !A$}۔
\end{enumerate}
له کومو شرطونو سره له دې بيانونو هر يو رښتيا دے؟
\end{prob}
\olpseganchor{OLP-0320-B014}{end}

\olpseganchor{OLP-0321-B004}{start}
\olfileid{inc}{inp}{tar}
\olpseganchor{OLP-0321-B004}{end}

\olpseganchor{OLP-0321-B005}{start}
\olsection{د رښتياوالي نه تعريفېدنه}
\olpseganchor{OLP-0321-B005}{end}

\olpseganchor{OLP-0321-B006}{start}
د \emph{تعريفېدنې} مفهوم د حساب د ژبې پر صوري معناپوهنې ولاړ دے۔
موږ د حساب په ژبه کښې د فورمولونو او جملو يو سټ بيان کړے دے۔
«مراد تفسير» دا دے چې دغه جملې د طبيعي عددونو په باب دعوې وبولو؛
داسې دعوه رښتيا يا ناسمه کېدے شي۔ \LR{$\Struct{N}$} دې هغه
جوړښت وي چې تعريف‌سيمه ئې~\LR{$\Nat$} او د حساب د ژبې
نښو تفسير ئې معياري دے۔ بيا~\LR{$\Sat{N}{!A}$} معنا لري: «\LR{$!A$}
په معياري تفسير کښې رښتيا ده»۔
\olpseganchor{OLP-0321-B006}{end}

\olpseganchor{OLP-0321-B007}{start}
\begin{defn}
د طبيعي عددونو اړيکه~\LR{$R(x_1,\dots,x_k)$} په~\LR{$\Struct{N}$} کښې
هله او يوازې هله \emph{تعريفېدونکې} ده چې د حساب د ژبې داسې
فورمول~\LR{$!A(x_1,\dots,x_k)$} وي چې د ټولو~\LR{$n_1,\dots,n_k$} دپاره
\LR{$R(n_1,\dots,n_k)$} هله او يوازې هله برقرار وي چې
\LR{$\Sat{N}{!A(\num n_1,\dots,\num
  n_k)}$}۔
\end{defn}
\olpseganchor{OLP-0321-B007}{end}

\olpseganchor{OLP-0321-B008}{start}
په بل ډول، يوه اړيکه په~\LR{$\Struct{N}$} کښې هله او يوازې هله
تعريفېدونکې ده چې په تيورۍ~\LR{$\Th{TA}$} کښې تمثيلېدونکې وي؛ دلته
\LR{$\Th{TA} =
\Setabs{!A}{\Sat{N}{!A}}$} د حساب د رښتينو جملو سټ دے۔ (که دا
خبره سملاسي روښانه نۀ وي، تعريفونو ته بېرته ورشئ او ځان پرې
قانع کړئ۔)
\olpseganchor{OLP-0321-B008}{end}

\olpseganchor{OLP-0321-B009}{start}
\begin{lem}
هره محاسبه کېدونکې اړيکه په~\LR{$\Struct{N}$} کښې تعريفېدونکې ده۔
\end{lem}
\olpseganchor{OLP-0321-B009}{end}

\olpseganchor{OLP-0321-B010}{start}
\begin{proof}
په آسانۍ کتل کېدے شي چې په~\LR{$\Th{Q}$} کښې د يوې اړيکې
تمثيلوونکے فورمول په~\LR{$\Struct{N}$} کښې هماغه اړيکه تعريفوي۔
\end{proof}
\olpseganchor{OLP-0321-B010}{end}

\olpseganchor{OLP-0321-B011}{start}
اوس پوښتنه دا ده چې آيا برعکس هم رښتيا دے؟ يعنې آيا په~\LR{$\Struct{N}$}
کښې هره تعريفېدونکې اړيکه محاسبه کېدونکې ده؟ ځواب نۀ دے۔ مثلاً:
\olpseganchor{OLP-0321-B011}{end}

\olpseganchor{OLP-0321-B012}{start}
\begin{lem}
د درېدنې اړيکه په~\LR{$\Struct{N}$} کښې تعريفېدونکې ده۔
\end{lem}
\olpseganchor{OLP-0321-B012}{end}

\olpseganchor{OLP-0321-B013}{start}
\begin{proof}
ياد ولرئ چې د کلين د عادي بڼې قضيه وايي: هره جزوي محاسبه
کېدونکې تابعه~\LR{$f$} داسې شاخص~\LR{$e$} لري چې د ټولو~\LR{$x \in \Nat$}
دپاره \LR{$f(x) =
U(\umin{s}{T(e,x,s)})$} وي؛ دلته~\LR{$U$} او~\LR{$T$} ابتدايي بازګشتي
او ځکه کلي دي۔ نو~\LR{$f(x)$} هله او يوازې هله تعريف شوې وي (يعنې
محاسبه ودرېږي) چې داسې~\LR{$s$} وي چې~\LR{$T(e,x,s)$} برقرار وي۔
\olpseganchor{OLP-0321-B013}{end}

\olpseganchor{OLP-0321-B014}{start}
اوس~\LR{$H$} د درېدنې اړيکه واخلو، يعنې
\[
H = \Setabs{\tuple{e,x}}{\lexists[s][T(e, x, s)]}.
\]
\LR{$!D_T$} دې په~\LR{$\Struct{N}$} کښې~\LR{$T$} تعريف کړي۔ نو
\[
H = \Setabs{\tuple{e,x}}{\Sat{N}{\lexists[s][!D_T(\num e, \num x, s)]}},
\]
او په دې توګه~\LR{$\lexists[s][!D_T(z, x, s)]$} په~\LR{$\Struct{N}$}
کښې~\LR{$H$} تعريفوي۔
\end{proof}
\olpseganchor{OLP-0321-B014}{end}

\olpseganchor{OLP-0321-B015}{start}
\begin{prob}
وښيئ چې \LR{$Q(n) \defiff n \in \Setabs{\Gn{!A}}{\Th{Q} \Proves !A}$}
  په حساب کښې تعريفېدونکې ده۔
\end{prob}
\olpseganchor{OLP-0321-B015}{end}

\olpseganchor{OLP-0321-B016}{start}
خو پخپله~\LR{$\Th{TA}$} څنګه ده؟ آيا په حساب کښې تعريفېدونکې ده؟
يعنې آيا سټ~\LR{$\Setabs{\Gn{!A}}{\Sat{N}{!A}}$} په حساب کښې
تعريفېدونکے دے؟ د تارسکي قضيه منفي ځواب ورکوي۔
\olpseganchor{OLP-0321-B016}{end}

\olpseganchor{OLP-0321-B017}{start}
\begin{thm}
\ollabel{thm:tarski}
د حساب د رښتينو تړلي فارمولونه سټ په حساب کښې نۀ شي تعريفېدے۔
\end{thm}
\olpseganchor{OLP-0321-B017}{end}

\olpseganchor{OLP-0321-B018}{start}
\begin{proof}
فرض کړئ~\LR{$!D(x)$} دا سټ تعريفوي؛ يعنې~\LR{$\Sat{N}{!A}$} هله او
يوازې هله چې~\LR{$\Sat{N}{!D(\gn{!A})}$}۔ د ثابت ټکي د لمې له مخې
داسې فورمول~\LR{$!A$} شته چې
\LR{$\Th{Q} \Proves !A \liff \lnot !D(\gn{!A})$}، او ځکه
\LR{$\Sat{N}{!A \liff \lnot !D(\gn{!A})}$}۔ خو بيا~\LR{$\Sat{N}{!A}$}
هله او يوازې هله چې~\LR{$\Sat{N}{\lnot !D(\gn{!A})}$}؛ دا له دې
فرض سره ټکر لري چې~\LR{$!D(y)$} د حساب د رښتينو جملو سټ تعريفوي۔
\end{proof}
\olpseganchor{OLP-0321-B018}{end}

\olpseganchor{OLP-0321-B019}{start}
تارسکي دا څېړنه د رښتياوالي پر يوه لا عمومي فلسفي مفهوم پلې کړه۔
د هرې ژبې~\LR{$L$} دپاره ئې استدلال وکړ چې د~\LR{$L$} د رښتياوالي مناسب
مفهوم بايد د هرې جملې~\LR{$X$} دپاره دا شرط پوره کړي:
\begin{quote}
«\LR{$X$}» هله او يوازې هله رښتيا ده چې \LR{$X$}۔
\end{quote}
د \href{https://doi.org/10.2307/2102968}{تارسکي د ۱۹۴۴ مقالې په څلورمه برخه} کښې دا شرط د واورې په بېلګه څرګندېږي۔
\begin{explain}
د واورې د سپينوالي خبره په پام کښې ونيسئ۔ که واوره سپينه وي،
د دې خبرې په نقل پسې د رښتياوالي حکم هم رښتيا دے؛ که واوره
سپينه نۀ وي، د رښتياوالي حکم ناسم دے۔ د جملې نقل د جملې نوم
دے، او د واورې په باب خبره خپله جمله ده۔
\end{explain}
خو د هرې هغې ژبې دپاره چې د قطري تابع د تمثيل توان ولري، او د
هر ژبني پريديکات~\LR{$T(x)$} دپاره، داسې جمله~\LR{$X$} جوړولے شو چې وايي:
«\LR{$X$} هله او يوازې هله چې \LR{$T(\text{\RL{«\LR{$X$}»}})$} نۀ وي»۔ که موږ نۀ
غواړو چې د رښتياوالي پريديکات يوه جمله په يو وخت کښې هم رښتيا او
هم ناسمه وبولي، نو د تارسکي پايله دا وه چې د يوې ژبې د ټولو جملو
دپاره د رښتياوالي پريديکات د ژبې له پولو څخه له وتلو پرته نۀ شو
ټاکلے۔ په نورو ټکو، د يوې ژبې د رښتياوالي پريديکات په هماغه ژبه
کښې نۀ شي تعريفېدے۔
\olpseganchor{OLP-0321-B019}{end}

\olpseganchor{OLP-0322-B004}{start}
\olpart{sol}{د دويمې درجې منطق}
\olpseganchor{OLP-0322-B004}{end}

\olpseganchor{OLP-0322-B005}{start}
\begin{editorial}
دا د دويمې درجې منطق د برخې لومړنۍ بڼه ده۔
\end{editorial}
\olpseganchor{OLP-0322-B005}{end}

\olpseganchor{OLP-0323-B004}{start}
\olchapter{sol}{syn}{نحوه او معنٰی پوهنه}
\olpseganchor{OLP-0323-B004}{end}

\olpseganchor{OLP-0323-B005}{start}
\begin{editorial}
تر اوسه د دويمې درجې منطق بنسټيزه نحوه او معنٰی پوهنه څېړل شوې ده۔
د يو فصل په توګه دا لا ډېره لنډه ده۔ د دويمې درجې متغيرونو پر ځاے
اېښودنه هم بايد وڅېړل شي، څو د دويمې درجې منطق د اشتقاق
نظامونو په باب خبرې وشي؛ په دې کښې ځينې نازک اړخونه شته۔
\end{editorial}
\olpseganchor{OLP-0323-B005}{end}

\olpseganchor{OLP-0324-B004}{start}
\olfileid{sol}{syn}{int}
\olpseganchor{OLP-0324-B004}{end}

\olpseganchor{OLP-0324-B005}{start}
\olsection{پېژندنه}
\olpseganchor{OLP-0324-B005}{end}

\olpseganchor{OLP-0324-B006}{start}
په لومړۍ درجې منطق کښې د يوې ټاکلې ژبې غيرمنطقي سمبولونه، يعنې
ثابت سمبولونه، تابعې او پريديکاتونه، له منطقي سمبولونو
سره يوځاے کوو څو د لومړۍ درجې جوړښتونه په باب خبرې بيانې کړو۔
دا کار د صدق د مفهوم له لارې کېږي: جوړښت~\LR{$\Struct{M}$}، د
متغيرونو ګومارنه~\LR{$s$} او فارمول~\LR{$!A$} سره تړي۔ \LR{$\Sat{M}{!A}[s]$}
هله برقرار وي چې هغه څۀ رښتيا وي چې \LR{$!A$} يې وايي، کله چې د هغې
ثابت سمبولونه، تابعې او پريديکاتونه د~\LR{$\Struct{M}$} له
مخې، او ازاد متغيرونه يې د~\LR{$s$} له مخې تفسير شي۔ د عينيت~\LR{$\eq$}
تفسير د \LR{$\Sat{M}{!A}[s]$} په تعريف کښې دننه دے، او د~\LR{$\lforall$} او
\LR{$\lexists$} تفسير هم همداسې دے۔ لومړے تل د جوړښت د
د تعريف ساحه~\LR{$\Domain{M}$} پر سر د عينيت اړيکه ده؛ کميت ټاکونکي تل د
ټولې د تعريف ساحه پر غړو ګرځي۔ خو مهمه دا ده چې کميت ټاکنه يوازې د
د تعريف ساحه پر غړو روا ده؛ له کميت ټاکونکي وروسته يوازې د څيز
متغيرونه راتلے شي۔
\olpseganchor{OLP-0324-B006}{end}

\olpseganchor{OLP-0324-B007}{start}
په دويمې درجې منطق کښې ژبه او د صدق تعريف دواړه غځېږي: ازاد او تړلي
د تابع او محمول متغيرونه، او پر هغو کميت ټاکنه، هم پکې راځي۔ دا
متغيرونه له تابعې او پريديکاتونه سره هماغه اړيکه لري چې
د څيز متغيرونه يې له ثابت سمبولونه سره لري۔ د دويمې درجې منطق د
ترمونو او فارمولونه په جوړولو کښې هماغسې ونډه لري، او کميت ټاکنه
پرې په ورته ډول کېږي۔ په \emph{معياري} معنٰی پوهنه کښې د دويمې درجې
کميت ټاکونکي د اړوند ډول پر ټولو ممکنو څيزونو ګرځي: د تابع متغيرونو
دپاره له \LR{$\Domain{M}$} څخه~\LR{$\Domain{M}$} ته پر \LR{$n$}-ځايه تابعو، او د
محمول متغيرونو دپاره پر \LR{$n$}-ځايه اړيکو۔ د بېلګې په توګه،
\LR{$\lforall[\Obj{v_0}][(\Obj{P^1_0}(\Obj{v_0}) \lor \lnot
  \Obj{P^1_0}(\Obj{v_0}))]$} د لومړۍ او دويمې درجې دواړو منطقونو
فارمولا ده؛ خو په دويمې درجې منطق کښې دا هم کتلے شو:
\LR{$\lforall[\Obj{V^1_0}][\lforall[\Obj{v_0}][(\Obj{V^1_0}(\Obj{v_0}) \lor
    \lnot \Obj{V^1_0}(\Obj{v_0}))]]$} او
\LR{$\lexists[\Obj{V^1_0}][\lforall[\Obj{v_0}][(\Obj{V^1_0}(\Obj{v_0}) \lor
    \lnot \Obj{V^1_0}(\Obj{v_0}))]]$}۔ ځکه چې په دغو کښې ازاد
متغيرونه نشته، دا د دويمې درجې منطق تړلي فارمولونه دي۔ دلته
\LR{$\Obj{V^1_0}$} يو دويمې درجې \LR{$1$}-ځايه محمول متغير دے۔ د
\LR{$\Obj{V^1_0}$} روا تفسيرونه هماغه دي چې د \LR{$\Obj{P^1_0}$} په شان
يو \LR{$1$}-ځايه پريديکات ته يې ټاکلے شو، يعنې د~\LR{$\Domain{M}$}
فرعي سټونه۔ نو پر هغو کميت ټاکنه دا وايي چې
\LR{$\lforall[\Obj{v_0}][(\Obj{V^1_0}(\Obj v_0) \lor \lnot
  \Obj{V^1_0}(v_0))]$} د~\LR{$\Domain{M}$} د يو فرعي سټ د
\LR{$\Obj{V^1_0}$} د ارزښت په توګه د ټاکلو په ټولو لارو کښې برقرار دے،
يا لږ تر لږه په يوه لاره کښې برقرار دے۔ ځکه چې هر سټ يو ورکړے څيز
يا لري يا يې نۀ لري، دواړه جملې په هر جوړښت کښې رښتيا دي۔
\olpseganchor{OLP-0324-B007}{end}

\olpseganchor{OLP-0716-B004}{start}
\olfileid{sol}{syn}{lan}
\olpseganchor{OLP-0716-B004}{end}

\olpseganchor{OLP-0716-B005}{start}
\olsection{د دويمې درجې منطق ژبه}
\olpseganchor{OLP-0716-B005}{end}

\olpseganchor{OLP-0716-B006}{start}
لکه په لومړۍ درجه منطق کښې، د دويمې درجې منطق عبارتونه له داسې بنسټيزې زېرمې جوړېږي چې \emph{متغيرونه}، \emph{ثابت سمبولونه}، \emph{پريديکاتونه} او کله کله \emph{تابعې} لري۔ له دې سمبولونو، منطقي تړنونو، کميت‌ښودونکو او د قوسونو او کامو په څېر د ليک نښو څخه \emph{ترمونه} او \emph{فارمولونه} جوړېږي۔ فرق دا دے چې د څيزونو پر متغيرونو سربېره په دويمه درجه منطق کښې د اړيکو او تابعو متغيرونه هم شته او کميت‌ښودنه پرې کېدے شي۔ نو د لومړۍ درجې منطقي سمبولونه شته، او ورسره دا هم شته:
\olpseganchor{OLP-0716-B006}{end}

\olpseganchor{OLP-0716-B007}{start}
\begin{enumerate}
\item د اړيکې د ځايونو د هر شمېر~\LR{$n$} دپاره د اړيکو د دويمې درجې متغيرونه يو شمېرېدونکے لامتناهي سټ: \LR{$\Obj V_0^n$}، \LR{$\Obj V_1^n$}، \LR{$\Obj V_2^n$}، \dots
\item د تابعو د دويمې درجې متغيرونه يو شمېرېدونکے لامتناهي سټ: \LR{$\Obj u_0^n$}، \LR{$\Obj u_1^n$}، \LR{$\Obj u_2^n$}، \dots
\end{enumerate}
\olpseganchor{OLP-0716-B007}{end}

\olpseganchor{OLP-0716-B008}{start}
په لومړۍ درجه منطق کښې عموماً د عينيت~\LR{$\eq$} نښه هم شامله وي۔ د ژبې~\LR{$\Lang{L}$} نامنطقي سمبولونه د مشخصو څيزونو، اړيکو او تابعو په باب د خبرو دپاره مهم دي۔ داسې تړلي فارمولونه هم شته چې نامنطقي سمبول نۀ کاروي، خو يوازې له~\LR{$\eq$} سره د ژبې څرګندوونکے وس محدود وي۔ په دويمه درجه منطق کښې د اړيکو او تابعو د متغيرونو زېرمه بې‌پايه ده۔ نو د لومړۍ درجې په عبارت کښې د هر نامنطقي ثابت، محمول او تابعې سمبول پر ځاے هم‌ډوله ازاد متغير کارولے شو، چې د اړيکې يا تابعې د ځايونو شمېر ئې هماغه وي۔ متغيرونه داسې ټاکو چې کميت‌ښودونکے ئې ونۀ تړي او د متغيرونو ټاکنه هماغه څيز، اړيکه يا تابع ورکړي چې سمبول يې تعبيراوه۔ په دې ډول هماغه خبره د ځانګړو نامنطقي سمبولونو له زېرمې پرته څرګندېږي۔
\olpseganchor{OLP-0716-B008}{end}

\olpseganchor{OLP-0325-B004}{start}
\olfileid{sol}{syn}{frm}
\olpseganchor{OLP-0325-B004}{end}

\olpseganchor{OLP-0325-B005}{start}
\olsection{ترمونه او فارمولونه}
\olpseganchor{OLP-0325-B005}{end}

\olpseganchor{OLP-0325-B006}{start}
د لومړۍ درجې منطق په څېر، د دويمې درجې منطق عبارتونه له داسې بنسټيز
لغت څخه جوړېږي چې \emph{متغيرونه}، \emph{ثابت سمبولونه}،
\emph{پريديکاتونه} او کله ناکله \emph{تابعې} لري۔ له
دغو څخه د منطقي نښلوونکو، کميت ټاکونکو او د لېک نښو، لکه قوسونو او
کاماګانو، په مرسته \emph{ترمونه} او \emph{فارمولونه} جوړېږي۔
توپير دا دے چې د دويمې درجې منطق د څيزونو د متغيرونو تر څنګ د اړيکو
او تابعو متغيرونه هم لري او پر هغو هم کميت ټاکنه روا ګڼي۔ نو د
دويمې درجې منطق منطقي سمبولونه د لومړۍ درجې منطق سمبولونه دي، له
لاندې اضافو سره:
\olpseganchor{OLP-0325-B006}{end}

\olpseganchor{OLP-0325-B007}{start}
\begin{enumerate}
\item د هرې ځايتيا~\LR{$n$} د دويمې درجې د اړيکو متغيرونه يو
  شمېرېدونکے لامتناهي سټ: \LR{$\Obj V_0^n$}، \LR{$\Obj V_1^n$}، \LR{$\Obj V_2^n$}، \dots
\item د دويمې درجې د تابعو متغيرونه يو شمېرېدونکے لامتناهي سټ:
  \LR{$\Obj u_0^n$}، \LR{$\Obj u_1^n$}، \LR{$\Obj u_2^n$}، \dots
\end{enumerate}
\olpseganchor{OLP-0325-B007}{end}

\olpseganchor{OLP-0325-B008}{start}
لکه څرنګه چې \LR{$x$}، \LR{$y$}، \LR{$z$} د لومړۍ درجې د \LR{$\Obj v_i$} متغيرونو دپاره
فرا-متغيرونه کاروو، همداسې به \LR{$X$}، \LR{$Y$}، \LR{$Z$} او نور د \LR{$\Obj V_i^n$}
دپاره، او \LR{$u$}، \LR{$v$} او نور د~\LR{$\Obj u_i^n$} دپاره فرا-متغيرونه کاروو۔
\olpseganchor{OLP-0325-B008}{end}

\olpseganchor{OLP-0325-B009}{start}
\begin{explain}
د دويمې درجې ژبې غيرمنطقي سمبولونه د لومړۍ درجې ژبې په څېر ټاکل
کېږي: د هغې ثابت سمبولونه، تابعې او پريديکاتونه
لست کېږي۔
\olpseganchor{OLP-0325-B009}{end}

\olpseganchor{OLP-0325-B010}{start}
په لومړۍ درجې منطق کښې عينيت~\LR{$\eq$} عموماً شامله وي۔ هلته
د ژبې~\LR{$\Lang{L}$} غيرمنطقي سمبولونه اړين دي څو په زړه پورې څۀ بيان
کړے شو۔ بېشکه داسې تړلي فارمولونه شته چې هېڅ غيرمنطقي سمبول نۀ
کاروي، خو يوازې له~\LR{$\eq$} سره په زړه پورې خبره کول ګران دي۔ په
دويمې درجې منطق کښې، ځکه چې د اړيکو او تابعو د متغيرونو نامحدود
زېرمتون لرو، له ځانګړيو غيرمنطقي سمبولونو پرته هم هر هغه څۀ ويلے شو
چې په لومړۍ درجې ژبه کښې يې ويلے شو۔
\end{explain}
\olpseganchor{OLP-0325-B010}{end}

\olpseganchor{OLP-0325-B011}{start}
\begin{defn}[د دويمې درجې ترمونه]
د~\LR{$\Lang L$} د \emph{دويمې درجې ترمونو} سټ، \LR{$\TrmSOL[L]$}، د
\olref[fol][syn][frm]{defn:terms} په تعريف کښې د لاندې شرط په زياتولو
تعريفېږي:
\begin{enumerate}
\item که \LR{$u$} يو \LR{$n$}-ځايه تابع متغير وي او \LR{$t_1$}، \dots، \LR{$t_n$}
  ترمونه وي، نو \LR{$\Atom{u}{t_1, \ldots, t_n}$} يو ترم دے۔
\end{enumerate}
\end{defn}
\olpseganchor{OLP-0325-B011}{end}

\olpseganchor{OLP-0325-B012}{start}
\begin{explain}
نو د دويمې درجې ترم د لومړۍ درجې ترم ته ورته ښکاري؛ يوازې دا توپير
لري چې د لومړۍ درجې په ترم کښې د تابع~\LR{$\Obj{f^n_i}$} پر
ځاے د دويمې درجې په ترم کښې د تابع متغير~\LR{$\Obj{u^n_i}$} راتلے شي۔
\end{explain}
\olpseganchor{OLP-0325-B012}{end}

\olpseganchor{OLP-0325-B013}{start}
\begin{defn}[د دويمې درجې فارمول]
د ژبې~\LR{$\Lang L$} د \emph{دويمې درجې فارمولونه} سټ~\LR{$\FrmSOL[L]$}
د \olref[fol][syn][frm]{defn:formulas} په تعريف کښې د لاندې شرطونو
په زياتولو تعريفېږي:
\begin{enumerate}
\item که \LR{$X$} يو \LR{$n$}-ځايه محمول متغير وي او \LR{$t_1$}، \dots، \LR{$t_n$}
  د~\LR{$\Lang L$} د دويمې درجې ترمونه وي، نو
  \LR{$\Atom{X}{t_1,\ldots, t_n}$} يو اټومي فارمول دے۔
\olpseganchor{OLP-0325-B013}{end}

\olpseganchor{OLP-0325-B014}{start}
\item که \LR{$!A$} فارمول او \LR{$u$} يو تابع متغير وي، نو
  \LR{$\lforall[u][!A]$} فارمول دے۔
\olpseganchor{OLP-0325-B014}{end}

\olpseganchor{OLP-0325-B015}{start}
\item که \LR{$!A$} فارمول او \LR{$X$} يو محمول متغير وي، نو
  \LR{$\lforall[X][!A]$} فارمول دے۔
\olpseganchor{OLP-0325-B015}{end}

\olpseganchor{OLP-0325-B016}{start}
\item که \LR{$!A$} فارمول او \LR{$u$} يو تابع متغير وي، نو
  \LR{$\lexists[u][!A]$} فارمول دے۔
\olpseganchor{OLP-0325-B016}{end}

\olpseganchor{OLP-0325-B017}{start}
\item که \LR{$!A$} فارمول او \LR{$X$} يو محمول متغير وي، نو
  \LR{$\lexists[X][!A]$} فارمول دے۔
\end{enumerate}
\end{defn}
\olpseganchor{OLP-0325-B017}{end}

\olpseganchor{OLP-0326-B004}{start}
\olfileid{sol}{syn}{sat}
\olpseganchor{OLP-0326-B004}{end}

\olpseganchor{OLP-0326-B005}{start}
\olsection{صدق}
\olpseganchor{OLP-0326-B005}{end}

\olpseganchor{OLP-0326-B006}{start}
\begin{explain}
د دويمې درجې فارمولونه دپاره د صدق اړيکې~\LR{$\Sat{M}{!A}[s]$} د
تعريف دپاره بايد تعريفونه داسې وغځوو چې د دويمې درجې متغيرونه
هم پکې راشي۔ د جوړښت مفهوم په دويمې درجې منطق کښې هماغه
د لومړۍ درجې منطق مفهوم دے۔ توپير يوازې د متغيرونو په
ګومارنو~\LR{$s$} کښې دے: اوس بايد نۀ يوازې د لومړۍ درجې متغيرونه،
بلکې د دويمې درجې متغيرونه ته هم ارزښتونه ورکړي۔
\end{explain}
\olpseganchor{OLP-0326-B006}{end}

\olpseganchor{OLP-0326-B007}{start}
\begin{defn}[د متغير ګومارنه]
د جوړښت~\LR{$\Struct{M}$} دپاره \emph{د متغير ګومارنه}~\LR{$s$}
هغه تابع ده چې هر لاندې څيز لېږي:
\begin{enumerate}
\item د څيز متغير~\LR{$\Obj{v_i}$} د~\LR{$\Domain M$} يوه غړي ته،
  يعنې \LR{$s(\Obj{v_i}) \in \Domain{M}$}؛
\item \LR{$n$}-ځايه اړيکيز متغير~\LR{$\Obj{V_i^n}$} پر~\LR{$\Domain{M}$} يوه
  \LR{$n$}-ځايه اړيکې ته، يعنې \LR{$s(\Obj{V_i^n}) \subseteq \Domain{M}^n$}؛
\item \LR{$n$}-ځايه تابع متغير~\LR{$\Obj{u_i^n}$} له \LR{$\Domain{M}$} څخه
  \LR{$\Domain{M}$} ته يوې \LR{$n$}-ځايه تابع ته، يعنې
  \LR{$s(\Obj{u_i^n})\colon \Domain{M}^n \to \Domain{M}$}؛
\end{enumerate}
\end{defn}
\olpseganchor{OLP-0326-B007}{end}

\olpseganchor{OLP-0326-B008}{start}
\begin{explain}
جوړښت هر ثابت سمبول او تابع ته قيمت
ورکوي، او د دويمې درجې د متغير ګومارنه د هر څيز او تابع متغير دپاره
څيزونه او تابعې ټاکي۔ دواړه په ګډه هر ترم ته يو ارزښت ټاکلو ته لار
هواروي۔
\end{explain}
\olpseganchor{OLP-0326-B008}{end}

\olpseganchor{OLP-0326-B009}{start}
\begin{defn}[د ترم قيمت]
که \LR{$t$} د ژبې~\LR{$\Lang L$} يو ترم، \LR{$\Struct M$} د~\LR{$\Lang L$} يو
جوړښت، او \LR{$s$} د~\LR{$\Struct M$} دپاره متغير ګومارنه
وي، \emph{قيمت}~\LR{$\Value{t}{M}[s]$} د لومړۍ درجې ترمونو په څېر
تعريفېږي، له لاندې اضافي شرط سره:
\begin{quote}
\indcase{t}{\Atom{u}{t_1, \ldots, t_n}}{
\[
\Value{\indfrm}{M}[s] = s(u)(\Value{t_1}{M}[s], \ldots,
\Value{t_n}{M}[s]).
\]}
\end{quote}
\end{defn}
\olpseganchor{OLP-0326-B009}{end}

\olpseganchor{OLP-0326-B010}{start}
\begin{defn}[\LR{$x$}-بڼه]
که \LR{$s$} د جوړښت~\LR{$\Struct M$} دپاره متغير ګومارنه
وي، نو د~\LR{$\Struct M$} هره هغه متغير ګومارنه \LR{$s'$} چې له~\LR{$s$}
څخه زيات نه زيات د~\LR{$x$} په ټاکلي ارزښت کښې توپير لري، د~\LR{$s$}
\emph{\LR{$x$}-بڼه} بلل کېږي۔ که \LR{$s'$} د~\LR{$s$} يوه \LR{$x$}-بڼه وي،
\LR{$\varAssign{s'}{s}{x}$} ليکو۔ (د دويمې درجې متغيرونو \LR{$X$} او \LR{$u$}
دپاره هم همداسې ده۔)
\end{defn}
\olpseganchor{OLP-0326-B010}{end}

\olpseganchor{OLP-0326-B011}{start}
\begin{defn}
  که \LR{$s$} د جوړښت~\LR{$\Struct M$} دپاره متغير ګومارنه
  وي او \LR{$m \in \Domain{M}$}، نو ګومارنه~\LR{$\Subst{s}{m}{x}$} هغه
  متغير ګومارنه ده چې داسې تعريفېږي:
  \[\Subst{s}{m}{x}(y) = \begin{cases}
    m & \text{\RL{که }} y \ident x\\
    s(y) & \text{\RL{کنه}},
  \end{cases}\]
  که \LR{$X$} يو \LR{$n$}-ځايه اړيکيز متغير او \LR{$M \subseteq
  \Domain{M}^n$} وي، نو \LR{$\Subst{s}{M}{X}$} هغه متغير ګومارنه ده
  چې داسې تعريفېږي:
  \[\Subst{s}{M}{X}(y) = \begin{cases}
    M & \text{\RL{که }} y \ident X\\
    s(y) & \text{\RL{کنه}}. 
  \end{cases}\]
  که \LR{$u$} يو \LR{$n$}-ځايه تابع متغير او \LR{$f\colon \Domain{M}^n
  \to \Domain{M}$} وي، نو \LR{$\Subst{s}{f}{u}$} هغه متغير ګومارنه
  ده چې داسې تعريفېږي:
  \[\Subst{s}{f}{u}(y) = \begin{cases}
    f & \text{\RL{که }} y \ident u\\
    s(y) & \text{\RL{کنه}}.
  \end{cases}\]
  په هر حالت کښې \LR{$y$} د لومړۍ يا دويمې درجې هر متغير کېدے شي۔
\end{defn}
\olpseganchor{OLP-0326-B011}{end}

\olpseganchor{OLP-0326-B012}{start}
\begin{defn}[صدق]
د دويمې درجې فارمولونه~\LR{$!A$} دپاره د صدق تعريف د
\olref[fol][syn][sat]{defn:satisfaction} په څېر دے، له لاندې اضافو
سره:
\begin{enumerate}
\item \indcase{!A}{\Atom{X^n}{t_1, \dots, t_n}}{\LR{$\Sat{M}{\indfrm}[s]$}
  هله او يوازې هله چې \LR{$\langle \Value{t_1}{M}[s], \dots, \Value{t_n}{M}[s] \rangle \in
  s(X^n)$}۔}
\item %
  \indcase{!A}{\lforall[X][!B]}{\LR{$\Sat{M}{\indfrm}[s]$} هله او يوازې
  هله چې د هر \LR{$M \subseteq \Domain{M}^n$} دپاره
  \LR{$\Sat{M}{!B}[\Subst{s}{M}{X}]$}۔}
\olpseganchor{OLP-0326-B012}{end}

\olpseganchor{OLP-0326-B013}{start}
\item %
  \indcase{!A}{\lexists[X][!B]}{\LR{$\Sat{M}{\indfrm}[s]$} هله او يوازې
  هله چې لږ تر لږه د يوه \LR{$M \subseteq \Domain{M}^n$} دپاره
  \LR{$\Sat{M}{!B}[\Subst{s}{M}{X}]$} وي۔}
\olpseganchor{OLP-0326-B013}{end}

\olpseganchor{OLP-0326-B014}{start}
\item %
  \indcase{!A}{\lforall[u][!B]}{\LR{$\Sat{M}{\indfrm}[s]$} هله او يوازې
  هله چې د هرې \LR{$f\colon \Domain{M}^n \to \Domain{M}$} دپاره
  \LR{$\Sat{M}{!B}[\Subst{s}{f}{u}]$} وي۔}
\olpseganchor{OLP-0326-B014}{end}

\olpseganchor{OLP-0326-B015}{start}
\item %
  \indcase{!A}{\lexists[u][!B]}{\LR{$\Sat{M}{\indfrm}[s]$} هله او يوازې
  هله چې لږ تر لږه د يوې \LR{$f\colon \Domain{M}^n \to \Domain{M}$} دپاره
  \LR{$\Sat{M}{!B}[\Subst{s}{f}{u}]$} وي۔}
\end{enumerate}
\end{defn}
\olpseganchor{OLP-0326-B015}{end}

\olpseganchor{OLP-0326-B016}{start}
\begin{ex}
  فارمول \LR{$\lforall[z][(\Atom{X}{z} \liff \lnot
    \Atom{Y}{z})]$} وګورئ۔ په دې کښې د دويمې درجې کميت ټاکونکي
  نشته، خو د دويمې درجې متغيرونه \LR{$X$} او~\LR{$Y$} پکښې دي (دلته
  يو ځايه ګڼل کېږي)۔ اړونده د لومړۍ درجې تړلے فارمول
  \LR{$\lforall[z][(\Atom{P}{z} \liff \lnot \Atom{R}{z})]$} وايي چې د~\LR{$P$}
  تر تفسير لاندې راتلونکے څيز د~\LR{$R$} تر تفسير لاندې نۀ راځي، او
  برعکس۔ په جوړښت کښې د پريديکات~\LR{$P$} تفسير د
  ~\LR{$\Assign{P}{M}$} له لارې ورکول کېږي۔ خو د دويمې درجې د \LR{$X$} او~\LR{$Y$}
  په شان متغيرونه تفسير پخپله جوړښت نۀ، بلکې د
  متغير ګومارنه ورکوي۔ ځکه چې د دويمې درجې دا فارمول
  تړلے فارمول نۀ ده (ازاد متغيرونه \LR{$X$} او~\LR{$Y$} لري)، يوازې
  د جوړښت~\LR{$\Struct{M}$} او د متغير ګومارنې~\LR{$s$} په
  نسبت صدق کولے شي۔
\olpseganchor{OLP-0326-B016}{end}

\olpseganchor{OLP-0326-B017}{start}
  \LR{$\Sat{M}{\lforall[z][(\Atom{X}{z} \liff \lnot \Atom{Y}{z})]}[s]$}
  هله وي چې د~\LR{$s(X)$} غړي د~\LR{$s(Y)$} غړي نۀ وي،
  او برعکس؛ يعنې هله او يوازې هله چې
  \LR{$s(Y) = \Domain{M} \setminus s(X)$}۔ مثلاً \LR{$\Domain{M} = \{1, 2, 3\}$}
  ونيسئ۔ ځکه چې دلته هېڅ پريديکاتونه، تابعې يا
  ثابت سمبولونه نشته، يوازې د~\LR{$\Struct{M}$} ساحه اړونده ده۔ اوس
  د \LR{$s_1(X) = \{1, 2\}$} او \LR{$s_1(Y) = \{3\}$} دپاره
  \LR{$\Sat{M}{\lforall[z][(\Atom{X}{z}
      \liff \lnot \Atom{Y}{z})]}[s_1]$} لرو۔
\olpseganchor{OLP-0326-B017}{end}

\olpseganchor{OLP-0326-B018}{start}
  برعکس، که \LR{$s_2(X) = \{1, 2\}$} او \LR{$s_2(Y) = \{2, 3\}$} وي،
  \LR{$\Sat/{M}{\lforall[z][(\Atom{X}{z} \liff \lnot
      \Atom{Y}{z})]}[s_2]$}۔ ځکه
  \LR{$\Sat{M}{\Atom{X}{z}}[\Subst{s_2}{2}{z}]$} (ځکه \LR{$2 \in \Subst{s_2}{2}{z}(X)$})،
  خو \LR{$\Sat/{M}{\lnot \Atom{Y}{z}}[\Subst{s_2}{2}{z}]$} (ځکه \LR{$2 \in
  \Subst{s_2}{2}{z}(Y)$} هم دے)۔
\end{ex}
\olpseganchor{OLP-0326-B018}{end}

\olpseganchor{OLP-0326-B019}{start}
\begin{ex}
  \LR{$\Sat{M}{\lexists[Y][(\lexists[y][\Atom{Y}{y}] \land
  \lforall[z][(\Atom{X}{z} \liff \lnot \Atom{Y}{z})])]}[s]$} هله
  وي چې داسې \LR{$N \subseteq \Domain{M}$} وي چې
  \LR{$\Sat{M}{(\lexists[y][\Atom{Y}{y}] \land \lforall[z][(\Atom{X}{z}
  \liff \lnot \Atom{Y}{z})])}[\Subst{s}{N}{Y}]$}۔ او دا د هغه
  \LR{$N \neq \emptyset$} دپاره وي (څو
  \LR{$\Sat{M}{\lexists[y][\Atom{Y}{y}]}[\Subst{s}{N}{Y}]$} وي) چې، د
  تېرې بېلګې په څېر، \LR{$N = \Domain{M} \setminus s(X)$} وي۔ په بل
  عبارت، \LR{$\Sat{M}{\lexists[Y][(\lexists[y][\Atom{Y}{y}] \land
  \lforall[z][(\Atom{X}{z} \liff \lnot \Atom{Y}{z})])]}[s]$}
  هله او يوازې هله چې \LR{$\Domain{M} \setminus s(X)$} ناتش وي، يعنې
  \LR{$s(X) \neq \Domain{M}$}۔ نو فارمول مثلاً هله صدق کوي چې
  \LR{$\Domain{M} = \{1, 2, 3\}$} او \LR{$s(X) = \{1, 2\}$} وي، خو هله
  نۀ چې \LR{$s(X) = \{1, 2, 3\} = \Domain{M}$} وي۔
\olpseganchor{OLP-0326-B019}{end}

\olpseganchor{OLP-0326-B020}{start}
  ځکه چې کله \LR{$s(X) = \Domain{M}$} وي دا فارمول صدق نۀ کوي،
  نو تړلے فارمول
  \[
  \lforall[X][\lexists[Y][(\lexists[y][\Atom{Y}{y}] \land
      \lforall[z][(\Atom{X}{z} \liff \lnot \Atom{Y}{z})])]]
  \]
  هېڅکله صدق نۀ کوي: د هر جوړښت~\LR{$\Struct{M}$} دپاره
  ګومارنه~\LR{$s(X) = \Domain{M}$} دا تړلے فارمول دروغ کوي۔ بل لور ته
  جمله
  \[
  \lexists[X][\lexists[Y][(\lexists[y][\Atom{Y}{y}] \land
      \lforall[z][(\Atom{X}{z} \liff \lnot \Atom{Y}{z})])]]
  \]
  د هرې ګومارنې~\LR{$s$} په نسبت صدق کوي، ځکه تل داسې
  \LR{$M \subseteq \Domain{M}$} موندلے شو چې \LR{$M \neq \Domain{M}$} وي
  (مثلاً \LR{$M = \emptyset$})۔
\end{ex}
\olpseganchor{OLP-0326-B020}{end}

\olpseganchor{OLP-0326-B021}{start}
\begin{ex}
  د دويمې درجې تړلے فارمول~\LR{$\lforall[X][\lforall[y][X(y)]]$} وايي
  چې هره \LR{$1$}-ځايه اړيکه، يعنې هر خاصيت، د هر څيز دپاره برقرار دے۔
  دا په څرګند ډول هېڅکله رښتيا نۀ ده، ځکه په هر~\LR{$\Struct{M}$} کښې،
  د داسې متغير ګومارنې~\LR{$s$} دپاره چې \LR{$s(X) = \emptyset$} او
  \LR{$s(y) = a \in \Domain{M}$} وي، \LR{$\Sat/{M}{X(y)}[s]$} لرو۔ دا معنا
  لري چې \LR{$!A \lif \lforall[X][\lforall[y][X(y)]]$} په دويمې درجې
  منطق کښې له \LR{$\lnot !A$} سره معادل دے، يعنې
  \LR{$\Sat{M}{!A \lif
  \lforall[X][\lforall[y][X(y)]]}$} هله او يوازې هله چې
  \LR{$\Sat{M}{\lnot !A}$}۔ په بل عبارت، په دويمې درجې منطق کښې
  \LR{$\lnot$} د \LR{$\lforall$} او~\LR{$\lif$} په مرسته تعريفولے شو۔
\end{ex}
\olpseganchor{OLP-0326-B021}{end}

\olpseganchor{OLP-0326-B022}{start}
\begin{prob}
  وښيئ چې په دويمې درجې منطق کښې \LR{$\lforall$} او~\LR{$\lif$} نور
  نښلوونکي تعريفولے شي:
  \begin{enumerate}
    \item يو يوځاييز اړيکيز متغير \LR{$X$} وټاکئ چې په \LR{$!A$} او \LR{$!B$} کښې ازاد نۀ وي۔ ثابته کړئ چې په دويمې درجې منطق کښې \LR{$!A \land !B$} له
    \LR{$\lforall[X][(!A \lif (!B \lif \lforall[x][X(x)])\lif
    \lforall[x][X(x)])]$} سره معادل دے۔
    \item داسې د دويمې درجې فارمولا ومومئ چې يوازې \LR{$\lforall$} او
    \LR{$\lif$} کاروي او له \LR{$!A \lor !B$} سره معادله وي۔
  \end{enumerate}
\end{prob}
\olpseganchor{OLP-0326-B022}{end}

\olpseganchor{OLP-0327-B004}{start}
\olfileid{sol}{syn}{sem}
\olsection{معنايي مفهومونه}
\olpseganchor{OLP-0327-B004}{end}

\olpseganchor{OLP-0327-B005}{start}
\begin{explain}
د \emph{اعتبار}، \emph{استلزام} او \emph{د صدق وړتيا} مرکزي منطقي
مفهومونه د دويمې درجې منطق دپاره د لومړۍ درجې منطق په شان تعريفېږي؛
يوازې بنسټيزه د صدق اړيکه اوس د دويمې درجې د فارمولونه اړيکه ده۔
د دويمې درجې تړلے فارمول هغه فارمول ده چې ټول متغيرونه
يې، د محمول او تابع متغيرونو په ګډون، تړلي وي۔
\end{explain}
\olpseganchor{OLP-0327-B005}{end}

\olpseganchor{OLP-0327-B006}{start}
\begin{defn}[اعتبار]
جمله~\LR{$!A$} هله \emph{معتبره} ده، \LR{$\Entails !A$}، چې د هر
جوړښت~\LR{$\Struct M$} دپاره \LR{$\Sat{M}{!A}$}۔
\end{defn}
\olpseganchor{OLP-0327-B006}{end}

\olpseganchor{OLP-0327-B007}{start}
\begin{defn}[استلزام]
د جملو سټ~\LR{$\Gamma$} هله جمله~\LR{$!A$} \emph{لازموي}، \LR{$\Gamma
\Entails !A$}، چې د هر هغه جوړښت~\LR{$\Struct M$} دپاره چې
\LR{$\Sat{M}{\Gamma}$} وي، \LR{$\Sat{M}{!A}$} هم وي۔
\end{defn}
\olpseganchor{OLP-0327-B007}{end}

\olpseganchor{OLP-0327-B008}{start}
\begin{defn}[د صدق وړتيا]
د جملو سټ~\LR{$\Gamma$} هله \emph{د صدق وړ} دے چې د کوم
جوړښت~\LR{$\Struct M$} دپاره \LR{$\Sat{M}{\Gamma}$} وي۔ که~\LR{$\Gamma$}
د صدق وړ نۀ وي، \emph{د صدق نۀ وړ} بلل کېږي۔
\end{defn}
\olpseganchor{OLP-0327-B008}{end}

\olpseganchor{OLP-0328-B004}{start}
\olfileid{sol}{syn}{exp}
\olsection{د بيان ځواک}
\olpseganchor{OLP-0328-B004}{end}

\olpseganchor{OLP-0328-B005}{start}
\begin{explain}
پر دويمې درجې متغيرونو کميت ټاکنه د لومړۍ درجې منطق په پرتله د ژبې د
بيان ځواک ډېر زياتوي۔ د دويمې درجې وجودي کميت ټاکنه پرېږدي چې ووايو
د ټاکلو خاصيتونو لرونکې تابعې يا اړيکې شته۔ په لومړۍ درجې منطق کښې
د دې يوازينۍ لار د همدې موخې دپاره د يو غيرمنطقي سمبول، يعنې
تابع يا پريديکات، ټاکل دي۔ د دويمې درجې عمومي کميت
ټاکنه پرېږدي چې ووايو د د تعريف ساحه ټول فرعي سټونه، پر هغې ټولې
اړيکې، يا له د تعريف ساحه څخه هغې ته ټولې تابعې يو خاصيت لري۔ په
لومړۍ درجې منطق کښې يوازې دا ويلے شو چې د ژبې يوه غيرمنطقي سمبول ته
ګومارل شوے فرعي سټ، اړيکه يا تابع هغه خاصيت لري۔ د دې خاصيت په
ټاکلو کښې هم د دويمې درجې کميت ټاکنه کارولے شو، که وايو چې ځينې
فرعي سټونه، اړيکې او تابعې يې لري، او که وايو چې ټول يې لري۔ دا
موږ ته اجازه راکوي چې داسې اړيکې تعريف کړو چې په لومړۍ درجې منطق
کښې نۀ تعريفېږي، او د ساحې داسې خاصيتونه بيان کړو چې هلته نۀ شي
بيانېدلے۔
\end{explain}
\olpseganchor{OLP-0328-B005}{end}

\olpseganchor{OLP-0328-B006}{start}
\begin{defn}
که \LR{$\Struct{M}$} د ژبې~\LR{$\Lang{L}$} يو جوړښت وي، اړيکه
\LR{$R \subseteq \Domain{M}^2$} هله په~\LR{$\Lang{L}$} کښې \emph{تعريفېدونکې}
ده چې داسې فارمول~\LR{$!A_R(x, y)$} وي چې يوازې \LR{$x$} او~\LR{$y$} يې ازاد
متغيرونه وي، او \LR{$R(a, b)$}، يعنې \LR{$\tuple{a,b} \in R$}، هله او يوازې
هله برقرار وي چې د \LR{$s(x) = a$} او \LR{$s(y) = b$} دپاره
\LR{$\Sat{M}{!A_R(x, y)}[s]$} وي۔
\end{defn}
\olpseganchor{OLP-0328-B006}{end}

\olpseganchor{OLP-0328-B007}{start}
\begin{ex}
په لومړۍ درجې منطق کښې د عينيت اړيکه~\LR{$\Id{\Domain{M}}$}، يعنې
\LR{$\Setabs{\tuple{a,a}}{a \in
  \Domain{M}}$}، د \LR{$\eq[x][y]$} په فارمولا تعريفولے شو۔ په دويمې
درجې منطق کښې دا اړيکه \emph{له~\LR{$\eq$} پرته} هم تعريفولے شو۔ ځکه
چې که \LR{$a$} او \LR{$b$} د~\LR{$\Domain{M}$} يو شان غړے وي، د
~\LR{$\Domain{M}$} د هماغو فرعي سټونو غړي دي (ځکه سټونه د خپلو
غړي له مخې ټاکل کېږي)۔ برعکس، که \LR{$a$} او \LR{$b$} بېل وي، په
هماغو فرعي سټونو کښې غړي نۀ دي: مثلاً که \LR{$a \neq b$} وي، \LR{$a \in
\{a\}$} خو \LR{$b \notin \{a\}$}۔ نو ``د~\LR{$\Domain{M}$} په هماغو
فرعي سټونو کښې د غړي کېدل'' داسې اړيکه ده چې پر \LR{$a$} او
\LR{$b$} هله او يوازې هله برقرار وي چې \LR{$a = b$}۔ دا اړيکه په دويمې درجې
منطق کښې بيانېدے شي، ځکه د~\LR{$\Domain{M}$} پر ټولو فرعي سټونو کميت
ټاکلے شو۔ نو لاندې فارمول \LR{$\Id{\Domain{M}}$} تعريفوي:
\[
\lforall[X][(X(x) \liff X(y))]
\]
\end{ex}
\olpseganchor{OLP-0328-B007}{end}

\olpseganchor{OLP-0328-B008}{start}
\begin{prob}
وښيئ چې \LR{$\lforall[X][(X(x) \lif X(y))]$} (پام: \LR{$\lif$}،
نۀ~\LR{$\liff$}!) \LR{$\Id{\Domain{M}}$} تعريفوي۔
\end{prob}
\olpseganchor{OLP-0328-B008}{end}

\olpseganchor{OLP-0328-B009}{start}
\begin{ex}
که \LR{$R$} دوه ځايه پريديکات وي، \LR{$\Assign{R}{M}$} پر~\LR{$\Domain{M}$}
دوه ځايه اړيکه ده۔ ښايي لږ مغشوشوونکې وي، خو \LR{$R$} به هم د
پريديکات~\LR{$R$} دپاره او هم د اړيکې~\LR{$\Assign{R}{M}$} پخپله دپاره
کاروو۔ د~\LR{$R$}
\emph{متعدي تړون}~\LR{$R^*$} هغه اړيکه ده چې د \LR{$a$} او \LR{$b$} تر منځ هله
برقراره وي چې د ځينو \LR{$c_1$}، \dots، \LR{$c_k$} دپاره \LR{$R(a,c_1)$}،
\LR{$R(c_1, c_2)$}، \dots، \LR{$R(c_k,b)$} برقرار وي۔ په دې کښې د \LR{$k = 0$}
حالت هم شامل دے: که \LR{$R(a,b)$} وي، \LR{$R^*(a,b)$} هم وي۔ نو
\LR{$R \subseteq R^*$}۔ په حقيقت کښې \LR{$R^*$} تر ټولو کوچنۍ اړيکه ده چې
~\LR{$R$} رانغاړي او متعدي ده۔ په دويمې درجې منطق کښې ويلے شو چې \LR{$X$}
هغه متعدي اړيکه ده چې~\LR{$R$} رانغاړي:
\begin{multline*}
  !B_R(X) \ident \lforall[x][\lforall[y][(R(x,y) \lif X(x, y))]] \land {}\\
\lforall[x][\lforall[y][\lforall[z][((X(x,y) \land X(y,z)) \lif X(x,
      z))]]].
\end{multline*}
لومړے عطف وايي چې \LR{$R \subseteq X$}، او دويم وايي چې \LR{$X$} متعدي ده۔
\olpseganchor{OLP-0328-B009}{end}

\olpseganchor{OLP-0328-B010}{start}
دا ويل چې \LR{$X$} د دې ډول تر ټولو کوچنۍ اړيکه ده، معنا يې دا ده چې
دا پخپله په هرې داسې اړيکې کښې شامله ده چې \LR{$R$} رانغاړي او متعدي ده۔
نو د~\LR{$R$} متعدي تړون په لاندې فارمول تعريفولے شو:
\[
R^*(X) \ident !B_R(X) \land \lforall[Y][(!B_R(Y) \lif
  \lforall[x][\lforall[y][(X(x, y) \lif Y(x,y))]])].
\]
\LR{$\Sat{M}{R^*(X)}[s]$} هله او يوازې هله وي چې \LR{$s(X) = R^*$}۔ د~\LR{$R$}
متعدي تړون په لومړۍ درجې منطق کښې نۀ شي بيانېدلے۔
\end{ex}
\olpseganchor{OLP-0328-B010}{end}

\olpseganchor{OLP-0329-B004}{start}
\olfileid{sol}{syn}{siz}
\olpseganchor{OLP-0329-B004}{end}

\olpseganchor{OLP-0329-B005}{start}
\olsection{د نامتناهي او د شمېر وړ
  د تعريف ساحې بيانول}
\olpseganchor{OLP-0329-B005}{end}

\olpseganchor{OLP-0329-B006}{start}
سټ~\LR{$M$} هله او يوازې هله (د ډېدېکېند په معنا) نامتناهي دے چې داسې
يو پر يو تابع~\LR{$f\colon M \to M$} وي چې پر هدف سټ بشپړه نۀ وي،
يعنې \LR{$\ran{f} \neq M$}۔ په لومړۍ درجې منطق کښې يو ځايه
تابع~\LR{$f$} کتلے شو او ويلے شو چې په
جوړښت~\LR{$\Struct{M}$} کښې هغې ته ګومارل شوې
تابع~\LR{$\Assign{f}{M}$} يو پر يو ده او \LR{$\ran{f} \neq
\Domain{M}$}:
\[
\lforall[x][\lforall[y][(\eq[f(x)][f(y)] \to \eq[x][y])]] \land
\lexists[y][\lforall[x][\eq/[y][f(x)]]].
\]
که \LR{$\Struct{M}$} په دې تړلے فارمول کښې صدق وکړي،
\LR{$\Assign{f}{M}: \Domain{M} \to \Domain{M}$} يو پر يو ده؛ نو
\LR{$\Domain{M}$} بايد نامتناهي وي۔ که \LR{$\Domain{M}$} نامتناهي وي، او
داسې تابع وي، \LR{$\Assign{f}{M}$} هغه تابع ټاکلے شو، او \LR{$\Struct{M}$} به
په تړلے فارمول کښې صدق وکړي۔ خو د دې دپاره اړينه ده چې زموږ ژبه
غيرمنطقي سمبول~\LR{$f$} ولري چې دې موخې ته يې کاروو۔ په دويمې درجې
منطق کښې په ساده ډول ويلے شو چې داسې تابع \emph{شته}۔ بيا~\LR{$f$} ته
اړتيا نۀ وي، او د سوچه دويمې درجې منطق دا تړلے فارمول ترلاسه کوو:
\[
\fn{Inf} \ident \lexists[u][(\lforall[x][\lforall[y][(\eq[u(x)][u(y)]
      \lif \eq[x][y])]] \land \lexists[y][\lforall[x][\eq/[y][u(x)]]])].
\]
\LR{$\Sat{M}{\fn{Inf}}$} هله او يوازې هله چې \LR{$\Domain{M}$} نامتناهي
وي۔ بيا \LR{$\fn{Fin} \ident \lnot \fn{Inf}$} تعريفولے شو؛
\LR{$\Sat{M}{\fn{Fin}}$} هله او يوازې هله چې \LR{$\Domain{M}$} متناهي وي۔
د سوچه لومړۍ درجې منطق يوه تړلے فارمول دا نۀ شي بيانولے چې
د تعريف ساحه نامتناهي ده، که څۀ هم د هغو يو نامتناهي سټ يې بيانولے
شي۔ د سوچه لومړۍ درجې منطق د تړلي فارمولونه هېڅ سټ نشته چې هله او
يوازې هله په جوړښت کښې صدق وکړي چې د هغې ساحه متناهي وي۔
\olpseganchor{OLP-0329-B006}{end}

\olpseganchor{OLP-0329-B007}{start}
\begin{prop}
\LR{$\Sat{M}{\fn{Inf}}$} هله او يوازې هله چې \LR{$\Domain{M}$} نامتناهي وي۔
\end{prop}
\olpseganchor{OLP-0329-B007}{end}

\olpseganchor{OLP-0329-B008}{start}
\begin{proof}
\LR{$\Sat{M}{\fn{Inf}}$} هله او يوازې هله چې
  \LR{$\Sat{M}{\lforall[x][\lforall[y][(\eq[u(x)][u(y)] \lif \eq[x][y])]]
    \land \lexists[y][\lforall[x][\eq/[y][u(x)]]]}[s]$} د کوم~\LR{$s$}
  دپاره وي۔ که داسې وي، \LR{$s(u)$} يوه يو پر يو تابع ده، او داسې
  \LR{$y \in \Domain{M}$} شته چې د~\LR{$s(u)$} په مدى کښې نۀ راځي۔ برعکس،
  که داسې يو پر يو \LR{$f\colon \Domain{M} \to \Domain{M}$} وي چې
  \LR{$\ran{f} \neq \Domain{M}$}، نو \LR{$s(u) = f$} د دې ډول متغير ګومارنه
  ده۔
\end{proof}
\olpseganchor{OLP-0329-B008}{end}

\olpseganchor{OLP-0329-B009}{start}
غيرتشه سټ \LR{$M$} هله د شمېر وړ دے چې د غړي يې داسې شمېرنه وي:
\[
m_0, m_1, m_2, \dots
\]
چې تکرار نۀ لري، خو متناهي کېدے شي۔ دا ډول شمېرنه هله او يوازې
هله شته چې يو غړے \LR{$z \in M$} او يوه تابع \LR{$f\colon M \to M$}
وي، داسې چې \LR{$z$}، \LR{$f(z)$}، \LR{$f(f(z))$}، \dots، د~\LR{$M$} ټول غړي
وي۔ که شمېرنه وي، \LR{$z = m_0$} او \LR{$f(m_k) = m_{k+1}$} (يا، که \LR{$m_k$}
د شمېرنې وروستے غړے وي، \LR{$f(m_k) = m_k$}) اړين غړے او تابع
ورکوي۔ بل لور ته، که داسې \LR{$z$} او \LR{$f$} وي، نو \LR{$z$}، \LR{$f(z)$}،
\LR{$f(f(z))$}، \dots، د تکرارونو په لرې کولو د~\LR{$M$} يوه شمېرنه ورکوي، او \LR{$M$} د شمېر وړ دے۔
دلته هره د تعريف ساحه ناتشه ده، نو لومړنے غړے اخيستلے شو۔ د \LR{$z$} او~\LR{$f$} شتون په دويمې درجې منطق کښې بيانولے شو، څو داسې
تړلے فارمول جوړه کړو چې هله او يوازې هله په جوړښت کښې
رښتيا وي چې جوړښت د شمېر وړ وي:
\[
\fn{Count} \ident
\lexists[z][\lexists[u][\lforall[X][((X(z) \land
      \lforall[x][(X(x) \lif X(u(x)))]) \lif \lforall[x][X(x)])]]]
\]
\olpseganchor{OLP-0329-B009}{end}

\olpseganchor{OLP-0329-B010}{start}
\begin{prop}
\LR{$\Sat{M}{\fn{Count}}$} هله او يوازې هله چې \LR{$\Domain{M}$}
د شمېر وړ وي۔
\end{prop}
\olpseganchor{OLP-0329-B010}{end}

\olpseganchor{OLP-0329-B011}{start}
\begin{proof}
فرض کړئ \LR{$\Domain{M}$} د شمېر وړ ده، او \LR{$m_0$}، \LR{$m_1$}، \dots،
يې يوه شمېرنه وي۔ د تکرار په لرې کولو ډاډمنېدے شو چې هېڅ \LR{$m_k$} دوه
ځلې نۀ راځي۔ له وروستي غړي پرته \LR{$f(m_k) = m_{k+1}$} تعريف کړئ او \LR{$s(z) = m_0$} او
\LR{$s(u) = f$} وټاکئ۔ که شمېرنه متناهي وي، نو د وروستي غړي د تابعې ارزښت هماغه وروستے غړے وټاکئ۔ ښيو چې
\[
\Sat{M}{\lforall[X][((X(z) \land \lforall[x][(X(x) \lif X(u(x)))])
    \lif \lforall[x][X(x)])]}[s]
\]
فرض کړئ \LR{$M \subseteq \Domain{M}$} يو په خپل سر ټاکل شوے فرعي سټ دے۔ بيا فرض
کړئ چې
\LR{$\Sat{M}{(X(z) \land \lforall[x][(X(x) \lif
X(u(x)))])}[\Subst{s}{M}{X}]$}۔ نو \LR{$\Subst{s}{M}{X}(z) \in M$} او
هر کله چې \LR{$x \in M$} وي، \LR{$(\Subst{s}{M}{X}(u))(x) \in M$} هم وي۔ په
بل عبارت، ځکه چې \LR{$\varAssign{\Subst{s}{M}{X}}{s}{X}$}، \LR{$m_0 \in M$}
او که \LR{$x \in M$} وي نو \LR{$f(x) \in M$}؛ نو \LR{$m_0 \in M$}،
\LR{$m_1 = f(m_0) \in M$}، \LR{$m_2 = f(f(m_0)) \in M$}، او داسې نور۔ نو
\LR{$M = \Domain{M}$}، او په نتيجه کښې
\LR{$\Sat{M}{\lforall[x][X(x)]}[\Subst{s}{M}{X}]$}۔ ځکه چې
\LR{$M \subseteq \Domain{M}$} په خپل سر ټاکل شوے و، ثبوت بشپړ شو:
\LR{$\Sat{M}{\fn{Count}}$}۔
\olpseganchor{OLP-0329-B011}{end}

\olpseganchor{OLP-0329-B012}{start}
اوس فرض کړئ \LR{$\Sat{M}{\fn{Count}}$}، يعنې
\[
\Sat{M}{\lforall[X][((X(z) \land \lforall[x][(X(x) \lif X(u(x)))])
    \lif \lforall[x][X(x)])]}[s]
\]
د کوم~\LR{$s$} دپاره۔ \LR{$m = s(z)$} او \LR{$f = s(u)$} وټاکئ او
\LR{$M = \{m, f(m), f(f(m)), \dots\}$} وګورئ۔ داسې تعريف شوے \LR{$M$} په
څرګند ډول د شمېر وړ دے۔ نو
\[
\Sat{M}{(X(z) \land \lforall[x][(X(x) \lif X(u(x)))])
    \lif \lforall[x][X(x)]}[\Subst{s}{M}{X}]
\]
د فرض له مخې۔ همدارنګه،
\LR{$\Sat{M}{X(z)}[\Subst{s}{M}{X}]$}، ځکه چې \LR{$M \ni m =
\Subst{s}{M}{X}(z)$}، او
\LR{$\Sat{M}{\lforall[x][(X(x) \lif
X(u(x)))]}[\Subst{s}{M}{X}]$} هم، ځکه چې که \LR{$x \in M$} وي نو
\LR{$f(x) \in M$} هم وي۔ نو ځکه چې مقدم او شرطي دواړه صدق کوي، تالي
هم بايد صدق وکړي:
\LR{$\Sat{M}{\lforall[x][X(x)]}[\Subst{s}{M}{X}]$}۔ خو دا معنا لري چې
\LR{$M = \Domain{M}$}؛ نو \LR{$\Domain{M}$} د شمېر وړ ده، ځکه \LR{$M$}
د تعريف له مخې همداسې دے۔
\end{proof}
\olpseganchor{OLP-0329-B012}{end}

\olpseganchor{OLP-0329-B013}{start}
\begin{prob}
تړلے فارمول \LR{$\fn{Inf} \land \fn{Count}$} په هغو او يوازې هغو
ساحو کښې رښتيا ده چې شمېرېدونکے لامتناهي وي۔ د~\LR{$\fn{Count}$} تعريف داسې
واړوئ چې يوه بله تړلے فارمول شي او په مستقيم ډول وښيي چې ساحه
شمېرېدونکے لامتناهي ده، او ثابته يې کړئ۔
\end{prob}
\olpseganchor{OLP-0329-B013}{end}

\olpseganchor{OLP-0330-B004}{start}
\olchapter{sol}{met}{د دويمې درجې منطق فرا تيوري}
\olpseganchor{OLP-0330-B004}{end}

\olpseganchor{OLP-0331-B004}{start}
\olfileid{sol}{met}{int}
\olpseganchor{OLP-0331-B004}{end}

\olpseganchor{OLP-0331-B005}{start}
\olsection{سريزه}
\olpseganchor{OLP-0331-B005}{end}

\olpseganchor{OLP-0331-B006}{start}
د لومړۍ درجې منطق څو ښې ځانګړنې لري۔ پوهېږو چې د پرېکړې وړ نۀ دے،
خو لږ تر لږه د بديهي اصولو په بڼه وړاندې کېدے شي۔ يعنې د لومړۍ درجې
منطق دپاره داسې اغېزمن ثبوتي نظامونه شته چې سم او بشپړ دي: د اشتقاق وړتيا
اړيکه~\LR{$\Proves$} داسې ټاکي چې د تړلي فارمولونه د هر سټ~\LR{$\Gamma$} او هرې
تړلے فارمول~\LR{$!A$}
دپاره \LR{$\Gamma \Entails !A$} هله او يوازې هله چې \LR{$\Gamma \Proves !A$}۔
له دې په ځانګړې توګه دا راځي چې د لومړۍ درجې منطق معتبرې جملې
په محاسبوي ډول د شمېر وړ دي۔ يوه محاسبه کېدونکې تابع~\LR{$f\colon \Nat \to
\Sent[L]$} شته چې د~\LR{$f$} ارزښتونه د~\LR{$\Lang{L}$} ټولې او يوازې معتبرې
تړلي فارمولونه دي۔ ځکه اشتقاقونه شمېرلے شو او هغه اشتقاقونه چې يوه
واحده تړلے فارمول اشتقاق کوي، هماغې تړلے فارمول ته لېږو۔ د دويمې درجې منطق
د لومړۍ درجې له منطق څخه د بيان زيات ځواک لري؛ نو د هغه د معتبرو
جملو راټولول عموماً لا پېچلي دي۔ په حقيقت کښې به وښيو چې د دويمې
درجې منطق يوازې ناپرېکړه‌وړ نۀ دے، بلکې معتبرې جملې ئې حتا
په محاسبوي ډول د شمېر وړ هم نۀ دي۔ له دې امله د دويمې درجې منطق
د بشپړې معنٰی پوهنې دپاره اغېزمن، سم او بشپړ ثبوتي نظام نۀ شي کېدے (که څه هم سم خو نابشپړ
ثبوتي نظامونه شته او په څېړنه کښې مهم دي)۔
\olpseganchor{OLP-0331-B006}{end}

\olpseganchor{OLP-0331-B007}{start}
د لومړۍ درجې منطق دوه نورې ځانګړنې هم لري: کمپکت دے (که د تړلي فارمولونه
د سټ~\LR{$\Gamma$} هر متناهي فرعي سټ د صدق وړ وي، نو \LR{$\Gamma$} پخپله هم
د صدق وړ دے)؛ او د لوېنهايم--سکولم قضيه پکې صدق کوي (که \LR{$\Gamma$}
نامتناهي مدل ولري، نو شمېرېدونکے لامتناهي مدل هم لري)۔ دا دواړه پايلې
د دويمې درجې منطق په بشپړه معنٰی پوهنه کښې په عمومي ډول نۀ صدق کوي۔ لامل بيا هم دا دے چې د دويمې درجې
منطق د د تعريف ساحې د اندازې په باب داسې واقعيتونه بيانولے شي چې د
لومړۍ درجې منطق ئې نۀ شي بيانولے۔
\olpseganchor{OLP-0331-B007}{end}

\olpseganchor{OLP-0332-B004}{start}
\olfileid{sol}{met}{spa}
\olpseganchor{OLP-0332-B004}{end}

\olpseganchor{OLP-0332-B005}{start}
\olsection{د دويمې درجې حساب}
\olpseganchor{OLP-0332-B005}{end}

\olpseganchor{OLP-0332-B006}{start}
په ياد ولرئ چې د پېانو د حساب تيوري~\LR{$\Th{PA}$} د~\LR{$\Th{Q}$} اته
بديهي اصول رانغاړي:
\begin{align*}
& \lforall[x][\eq/[x'][\Obj 0]]\\
& \lforall[x][\lforall[y][(\eq[x'][y'] \lif \eq[x][y])]]\\
& \lforall[x][(\eq[x][\Obj 0] \lor \lexists[y][\eq[x][y']])]\\ 
& \lforall[x][\eq[(x + \Obj 0)][x]]\\
& \lforall[x][\lforall[y][\eq[(x + y')][(x + y)']]]\\
& \lforall[x][\eq[(x \times \Obj 0)][\Obj 0]]\\
& \lforall[x][\lforall[y][\eq[(x \times y')][((x \times y) + x)]]]\\
& \lforall[x][\lforall[y][(x < y \liff \lexists[z][\eq[(z' + x)][y]])]]\\
\intertext{\RL{او ورسره د لاندې بڼې ټولې جملې:}}
& (!A(\Obj 0) \land \lforall[x][(!A(x) \lif !A(x'))]) \lif \lforall[x][!A(x)].
\end{align*}
وروستۍ بڼه يو «شېما» دے: داسې قالب چې د حساب د ژبې نامتناهي ډېرې
تړلي فارمولونه جوړوي، د هر فارمول~\LR{$!A(x)$} دپاره يوه۔ دې قالب ته
د استقرا (د لومړۍ درجې) \emph{بديهي شېما} وايو۔ د پېانو په
\emph{دويمې درجې} حساب~\LR{$\Th{PA^2}$} کښې استقرا په يوه جمله بيانېدے
شي۔ \LR{$\Th{PA^2}$} د پورته لومړنيو اتو بديهي اصولو او د (دويمې درجې)
\emph{استقرا له بديهي اصل} څخه جوړه ده:
\[
\lforall[X][((X(\Obj 0) \land \lforall[x][(X(x) \lif X(x'))]) \lif \lforall[x][X(x)])].
\]
دا وايي: که د د تعريف ساحه يو فرعي سټ~\LR{$X$}، \LR{$\Assign{\Obj{0}}{M}$} ولري
او له هر~\LR{$x \in \Domain{M}$} سره \LR{$\Assign{\prime}{M}(x)$} هم ولري
(يعنې «د جانشين تر عمل لاندې تړلے» وي)، نو د د تعريف ساحه ټول څيزونه
لري (يعنې \LR{$X = \Domain{M}$})۔
\olpseganchor{OLP-0332-B006}{end}

\olpseganchor{OLP-0332-B007}{start}
د استقرا بديهي اصل تضمينوي چې هر جوړښت چې دا اصل پوره کوي،
په~\LR{$\Domain{M}$} کښې يوازې هغه غړي لري چې بديهي اصول ئې
شتون غواړي؛ يعنې د~\LR{$\num{n}$} ارزښتونه د \LR{$n \in \Nat$} دپاره۔
د~\LR{$\Th{PA^2}$} هېڅ مدل نامعياري شمېرې نۀ لري۔ دلته د دويمې درجې کميت په بشپړه معنٰی پوهنه کښې د ساحې پر ټولو فرعي سټونو چلېږي۔
\olpseganchor{OLP-0332-B007}{end}

\olpseganchor{OLP-0332-B008}{start}
\begin{thm}
\ollabel{thm:sol-pa-standard}
که \LR{$\Sat{M}{\Th{PA^2}}$}، نو \LR{$\Domain{M} =
\Setabs{\Value{\num{n}}{M}}{n \in \Nat}$}.
\end{thm}
\olpseganchor{OLP-0332-B008}{end}

\olpseganchor{OLP-0332-B009}{start}
\begin{proof}
دې \LR{$N = \Setabs{\Value{\num{n}}{M}}{n \in \Nat}$} وي، او فرض کړئ
\LR{$\Sat{M}{\Th{PA^2}}$}۔ څرګنده ده چې د هر \LR{$n \in \Nat$} دپاره
\LR{$\Value{\num{n}}{M} \in \Domain{M}$}؛ نو \LR{$N \subseteq \Domain{M}$}۔
\olpseganchor{OLP-0332-B009}{end}

\olpseganchor{OLP-0332-B010}{start}
اوس د بل لوري شمول وښيو۔ داسې ارزښت ټاکنه~\LR{$s$} وګورئ چې
\LR{$s(X) = N$}۔ د فرض له مخې:
\begin{align*}
\Sat{M}{\lforall[X][((X(\Obj 0) \land \lforall[x][(X(x)
      \lif X(x'))]) \lif \lforall[x][X(x)])]}, & \text{\RL{ نو}}\\
\Sat{M}{(X(\Obj 0) \land \lforall[x][(X(x) \lif X(x'))]) \lif
  \lforall[x][X(x)]}[s]. & 
\end{align*}
د دې شرطي جملې مقدم وګورئ۔ \LR{$\Value{\Obj{0}}{M} \in N$}، نو
\LR{$\Sat{M}{X(\Obj{0})}[s]$}۔ دويم عطف~\LR{$\lforall[x][(X(x) \lif X(x'))]$}
هم صدق کوي۔ ځکه فرض کړئ \LR{$x \in N$}۔ د~\LR{$N$} له تعريفه، د کوم~\LR{$n$}
دپاره \LR{$x = \Value{\num{n}}{M}$}۔ له دې څخه
\LR{$\Assign{\prime}{M}(x) = \Value{\num{n+1}}{M} \in N$}۔ نو
\LR{$\Assign{\prime}{M}(x) \in N$}۔
\olpseganchor{OLP-0332-B010}{end}

\olpseganchor{OLP-0332-B011}{start}
نو \LR{$\Sat{M}{X(\Obj 0) \land \lforall[x][(X(x) \lif X(x'))]}[s]$} لرو۔
له دې امله \LR{$\Sat{M}{\lforall[x][X(x)]}[s]$}۔ دا وايي چې د هر
\LR{$x \in \Domain{M}$} دپاره \LR{$x \in s(X) = N$}؛ نو \LR{$\Domain{M}
\subseteq N$}۔
\end{proof}
\olpseganchor{OLP-0332-B011}{end}

\olpseganchor{OLP-0332-B012}{start}
\begin{cor}
\ollabel{cor:sol-pa-categorical}
د~\LR{$\Th{PA^2}$} هر دوه مدلونه آيزومورف دي۔
\end{cor}
\olpseganchor{OLP-0332-B012}{end}

\olpseganchor{OLP-0332-B013}{start}
\begin{proof}
د \olref{thm:sol-pa-standard} له مخې، د~\LR{$\Th{PA^2}$} د هر مدل تعريف ساحه
د~\LR{$\Value{\num{n}}{M}$} له ارزښتونو پرته بل څه نۀ لري۔ هر داسې مدل
د~\LR{$\Th{Q}$} مدل هم دے۔ د \olref[mod][mar][stm]{prop:thq-standard} له
مخې، هر داسې مدل معياري دے؛ يعنې له~\LR{$\Struct{N}$} سره آيزومورف دے۔
\end{proof}
\olpseganchor{OLP-0332-B013}{end}

\olpseganchor{OLP-0332-B014}{start}
پورته مو \LR{$\Th{PA^2}$} هغه تيوري وبلله چې لومړني اته حسابي بديهي
اصول او د دويمې درجې د استقرا بديهي اصل لري۔ په حقيقت کښې د دويمې
درجې منطق د بيان د ځواک له برکته، د دويمې درجې پېانو حساب دپاره
يوازې \emph{لومړني دوه} حسابي بديهي اصول او استقرا بس دي۔
\olpseganchor{OLP-0332-B014}{end}

\olpseganchor{OLP-0332-B015}{start}
\begin{prop}\ollabel{prop:sol-pa-definable}
دې \LR{$\Th{PA^{2\dagger}}$} هغه د دويمې درجې تيوري وي چې لومړني دوه
حسابي بديهي اصول (د جانشين بديهي اصول) او د دويمې درجې د استقرا
بديهي اصل لري۔ نو \LR{$\le$}، \LR{$+$} او \LR{$\times$} په~\LR{$\Th{PA^{2\dagger}}$}
کښې تعريفېدے شي۔
\end{prop}
\olpseganchor{OLP-0332-B015}{end}

\olpseganchor{OLP-0332-B016}{start}
\begin{proof}
د دې ښودلو دپاره چې \LR{$\le$} تعريفېدے شي، داسې فورمول~\LR{$!A_\le(x, y)$}
موندل غواړو چې \LR{$\Sat{N}{!A_\le(\num n, \num m)}$} هله او يوازې هله
چې \LR{$n \le m$}۔ دا فورمول وګورئ:
\[
!B(x, Y) \ident Y(x) \land \lforall[y][(Y(y) \lif Y(y'))]
\]
که د~\LR{$Y \subseteq \Nat$} سټ \LR{$!B(\num n, Y)$} د صدق وړ کړي، نو \LR{$\Setabs{m}{n \le m} \subseteq Y$}۔ ځکه چې د جانشين شرط د پيل له غړي څخه هر ورپسې غړے رانغاړي۔ له بلې خوا، له هماغې شمېرې پيلېدونکے سټ پخپله د پيل غړے لري او د جانشين لاندې تړلے دے۔ نو هر هغه عدد چې له پيل څخه کوچنے وي، په دغه سټ کښې نشته؛ له دې امله
\LR{$!A_\le(x, y) \ident \lforall[Y][(!B(x, Y) \lif Y(y))]$} اخيستلے شو۔
\olpseganchor{OLP-0332-B016}{end}

\olpseganchor{OLP-0332-B017}{start}
د دې ليدلو دپاره چې جمع تعريفېدے شي، پام وکړئ چې \LR{$k+l=m$} هله او
يوازې هله چې داسې تابع~\LR{$u$} شته چې \LR{$u(0)=k$}، د هر \LR{$n$} دپاره
\LR{$u(n')=u(n)'$}، او \LR{$m = u(l)$}۔ دا معادلت د لاندې فورمول په وسيله
په~\LR{$\Th{PA^{2\dagger}}$} کښې د جمع د تعريف دپاره کارولے شو:
\[
!A_+(x,y,z) \ident \lexists[u][(u(\Obj 0)=x \land \lforall[w][u(w')=u
 (w)'] \land u(y)=z)]
\] 
څرګنده بايد وي چې \LR{$\Sat{N}{!A_+(\num k, \num l, \num m)}$} هله او
يوازې هله چې \LR{$k+l=m$}۔
\end{proof}
\olpseganchor{OLP-0332-B017}{end}

\olpseganchor{OLP-0332-B018}{start}
\begin{prob}
د \olref[sol][met][spa]{prop:sol-pa-definable} ثبوت بشپړ کړئ۔
\end{prob}
\olpseganchor{OLP-0332-B018}{end}

\olpseganchor{OLP-0333-B004}{start}
\olfileid{sol}{met}{nax}
\olpseganchor{OLP-0333-B004}{end}

\olpseganchor{OLP-0333-B005}{start}
\olsection{د دويمې درجې منطق د بديهي اصولو په بڼه بشپړ نۀ شي وړاندې کېدے}
\olpseganchor{OLP-0333-B005}{end}


\olpseganchor{OLP-0333-B006}{start}
\begin{thm}
\ollabel{thm:sol-undecidable}
د دويمې درجې منطق ناپرېکړه‌وړ دے۔
\end{thm}
\olpseganchor{OLP-0333-B006}{end}

\olpseganchor{OLP-0333-B007}{start}
\begin{proof}
د لومړۍ درجې تړلے فارمول په لومړۍ درجې منطق کښې هله او يوازې هله
معتبره ده چې په دويمې درجې منطق کښې هم معتبره وي؛ او د لومړۍ درجې
منطق ناپرېکړه‌وړ دے۔
\end{proof}
\olpseganchor{OLP-0333-B007}{end}

\olpseganchor{OLP-0333-B008}{start}
\begin{thm}
\ollabel{cor:sol-not-axiomatizable}
د دويمې درجې منطق د بشپړې معنٰی پوهنې دپاره هېڅ اغېزمن، سم او بشپړ اشتقاق نظام نشته۔
\end{thm}
\olpseganchor{OLP-0333-B008}{end}

\olpseganchor{OLP-0333-B009}{start}
\begin{proof}
فرض کړئ \LR{$!A$} د حساب په ژبه کښې تړلے فارمول ده۔ \LR{$\Sat{N}{!A}$} هله
او يوازې هله چې \LR{$\Th{PA^2} \Entails !A$}۔ دې \LR{$!P$} د~\LR{$\Th{PA^2}$} د نهو
بديهي اصولو عطف وي۔ \LR{$\Th{PA^2} \Entails !A$} هله او يوازې هله چې
\LR{$\Entails !P \lif !A$}، يعنې د هر مدل په اړه \LR{$\Sat{M}{!P \lif !A}$}۔
اوس هغه تړلے فارمول
\LR{$\lforall[z][\lforall[u][\lforall[u'][\lforall[u''][\lforall[L][(!P' \lif !A')]]]]]$}
وګورئ چې \LR{$\Obj{0}$} پکې په \LR{$z$}، \LR{$\prime$} په يوځايه تابع متغير~\LR{$u$}،
\LR{$+$} او \LR{$\times$} په ترتيب په دوه‌ځايه تابع متغيرونو~\LR{$u'$} او~\LR{$u''$}،
او \LR{$<$} په دوه‌ځايه اړيکيز متغير~\LR{$L$} بدل شوي او بيا پرې کلي کميت
ايښودل شوے دے۔ دا د خالص دويمې درجې منطق معتبره جمله هله او يوازې
هله ده چې اصلي جمله معتبره وي؛ يعنې هله او يوازې هله چې
\LR{$\Th{PA^2} \Entails !A$}، او دا هله او يوازې هله چې \LR{$\Sat{N}{!A}$}۔
نو که د دويمې درجې منطق د بشپړې معنٰی پوهنې دپاره اغېزمن، سم او بشپړ ثبوتي نظام واے، د هغه په
وسيله به مو د ثبوتونو د کوډونو په لړ کښې د هرې حسابي جملې د تفسير شوي اعتبار لپاره ثبوت لټولے۔ ټولې پلټنې په وار وار يو ګام وړاندې وړو، او هر کله چې ثبوت ومومو، اړونده حسابي جمله خپروو۔ په دې توګه به مو په~\LR{$\Struct{N}$} کښې د رښتينو تړلي فارمولونه يوه محاسبه
کېدونکې شمېرنه~\LR{$f\colon \Nat \to \Sent[L_A]$} تعريف کړې واے۔ دلته د حسابي تمثيل دپاره د جملو له مخکې ټاکل شوي عددي کوډونه کارېږي۔ دا تابع
به په~\LR{$\Th{Q}$} کښې د لومړۍ درجې د يوه فورمول~\LR{$!B_f(x, y)$} په وسيله
تمثيلېدله۔ بيا به فارمول~\LR{$\lexists[x][!B_f(x, y)]$} په~\LR{$\Struct{N}$}
کښې د لومړۍ درجې د رښتينو تړلي فارمولونه سټ تعريف کړے واے، چې د
تارسکي له قضيې سره په ټکر کښې دے۔
\end{proof}
\olpseganchor{OLP-0333-B009}{end}

\olpseganchor{OLP-0334-B004}{start}
\olfileid{sol}{met}{com}
\olpseganchor{OLP-0334-B004}{end}

\olpseganchor{OLP-0334-B005}{start}
\olsection{د دويمې درجې منطق د متناهي شاهد خاصيت نۀ لري}
\olpseganchor{OLP-0334-B005}{end}

\olpseganchor{OLP-0334-B006}{start}
\begin{explain}
د جملو سټ~\LR{$\Gamma$} هغه وخت \emph{په متناهي ډول د صدق وړ} بلئ چې هر
متناهي فرعي سټ ئې د صدق وړ وي۔ د لومړۍ درجې منطق دا خاصيت لري چې
که د تړلي فارمولونه سټ~\LR{$\Gamma$} په متناهي ډول د صدق وړ وي، نو پخپله
هم د صدق وړ دے۔ دا \emph{د متناهي شاهد خاصيت (کمپکتوالی)} بلل کېږي۔ يوه هم‌ارزښته بڼه
ئې د معنايي استلزام په هکله ده: که \LR{$\Gamma \Entails !A$}، نو د کوم
متناهي فرعي سټ~\LR{$\Gamma_0 \subseteq \Gamma$} دپاره لا له وړاندې
\LR{$\Gamma_0 \Entails !A$}۔ په دې بڼه دا د بشپړتيا د قضيې سمدستي پايله
ده: که \LR{$\Gamma \Entails !A$}، نو د بشپړتيا له مخې \LR{$\Gamma \Proves !A$}۔
خو يو اشتقاق د~\LR{$\Gamma$} له متناهي شمېر تړلي فارمولونه څخه
ګټه اخيستلے شي۔
\olpseganchor{OLP-0334-B006}{end}

\olpseganchor{OLP-0334-B007}{start}
د متناهي شاهد خاصيت د دويمې درجې منطق په بشپړه معنٰی پوهنه کښې نۀ صدق کوي۔ د دويمې درجې د
تړلي فارمولونه داسې سټونه شته چې په متناهي ډول د صدق وړ دي، خو
پخپله د صدق وړ نۀ دي؛ او داسې سټونه هم شته چې له هغو څخه~\LR{$!A$}
منطقي لازم راځي، خو له هېڅ متناهي فرعي سټ څخه ئې~\LR{$!A$} نۀ راځي۔
\end{explain}
\olpseganchor{OLP-0334-B007}{end}


\olpseganchor{OLP-0334-B008}{start}
\begin{thm}
\ollabel{thm:sol-undecidable}
د دويمې درجې منطق په بشپړه معنٰی پوهنه کښې د متناهي شاهد خاصيت نۀ لري۔
\end{thm}
\olpseganchor{OLP-0334-B008}{end}

\olpseganchor{OLP-0334-B009}{start}
\begin{proof}
په ياد ولرئ چې
\[
\fn{Inf} \ident \lexists[u][(\lforall[x][\lforall[y][(\eq[u(x)][u(y)]
      \lif \eq[x][y])]] \land \lexists[y][\lforall[x][\eq/[y][u(x)]]])]
\]
په جوړښت کښې هله او يوازې هله صدق کوي چې تعريف ساحه ئې
نامتناهي وي۔ دې \LR{$!A^{\ge n}$} هغه جمله وي چې وايي تعريف ساحه لږ تر
لږه \LR{$n$} غړي لري، لکه:
\[
!A^{\ge n} \ident \lexists[x_1][\dots\lexists[x_n][(\eq/[x_1][x_2]
    \land \eq/[x_1][x_3] \land \dots \land \eq/[x_{n-1}][x_n])]].
\]
په دې فورمول کښې د درې ټکو نښه د شاهدانو د ټولو بېلو جوړو نابرابري رانغاړي۔ د يو شاهد لپاره دغه عطف خالي او رښتيا دے۔
د تړلي فارمولونه لاندې سټ وګورئ:
\[
\Gamma = \{\lnot \fn{Inf}, !A^{\ge 1}, !A^{\ge 2}, !A^{\ge 3}, \dots\}.
\]
دا په متناهي ډول د صدق وړ دے، ځکه چې د هر متناهي فرعي سټ~\LR{$\Gamma_0
\subseteq \Gamma$} دپاره داسې \LR{$k \ge 1$} غوره کولے شو چې له~\LR{$k$} څخه د لوړو
اندازو هېڅ جمله په دغه فرعي سټ کښې نۀ وي۔
که \LR{$\Domain{M}$}
\LR{$k$} غړي ولري، نو \LR{$\Sat{M}{\Gamma_0}$}۔ خو \LR{$\Gamma$} د صدق وړ
نۀ دے: که \LR{$\Sat{M}{\lnot \fn{Inf}}$}، نو \LR{$\Domain{M}$} بايد متناهي وي،
مثلاً د اندازې~\LR{$k$}۔ بيا \LR{$\Sat/{M}{!A^{\ge k+1}}$}۔
\end{proof}
\olpseganchor{OLP-0334-B009}{end}

\olpseganchor{OLP-0334-B010}{start}
\begin{prob}
داسې سټ~\LR{$\Gamma$} او تړلے فارمول~\LR{$!A$} مثال ورکړئ چې
\LR{$\Gamma \Entails !A$} وي، خو د هر متناهي فرعي سټ~\LR{$\Gamma_0 \subseteq
\Gamma$} دپاره \LR{$\Gamma_0 \Entails/ !A$} وي۔
\end{prob}
\olpseganchor{OLP-0334-B010}{end}

\olpseganchor{OLP-0335-B004}{start}
\olfileid{sol}{met}{lst}
\olpseganchor{OLP-0335-B004}{end}

\olpseganchor{OLP-0335-B005}{start}
\olsection{د لوېنهايم--سکولم قضيه د دويمې درجې منطق دپاره نۀ صدق کوي}
\olpseganchor{OLP-0335-B005}{end}

\olpseganchor{OLP-0335-B006}{start}
\begin{explain}
د (ښکته) لوېنهايم--سکولم قضيه وايي چې د تړلي فارمولونه هر هغه سټ
چې په شمېر وړ لومړۍ درجې ژبه کښې نامتناهي مدل لري، د شمېر وړ مدل هم لري۔ دا هم د بشپړتيا
د قضيې پايله ده: د بشپړتيا ثبوت د تړلي فارمولونه د هر سازګار سټ
دپاره يو مدل جوړوي او هغه مدل د شمېر وړ وي۔ د پورته
لوېنهايم--سکولم قضيه هم شته، چې ضمانت کوي: که د تړلي فارمولونه يو سټ
شمېرېدونکے لامتناهي مدل ولري، نو د شمېر نۀ وړ مدل هم لري۔ دا دواړه
قضيې د دويمې درجې منطق په بشپړه معنٰی پوهنه کښې نۀ صدق کوي۔
\end{explain}
\olpseganchor{OLP-0335-B006}{end}


\olpseganchor{OLP-0335-B007}{start}
\begin{thm}
\ollabel{thm:sol-no-ls} د لوېنهايم--سکولم قضيه د دويمې درجې منطق
په بشپړه معنٰی پوهنه کښې نۀ صدق کوي: داسې تړلي فارمولونه شته چې نامتناهي مدلونه لري،
خو د شمېر وړ مدل نۀ لري۔
\end{thm}
\olpseganchor{OLP-0335-B007}{end}

\olpseganchor{OLP-0335-B008}{start}
\begin{proof}
په ياد ولرئ چې
\[
\fn{Count} \ident \lexists[z][\lexists[u][\lforall[X][((X(z) \land
      \lforall[x][(X(x) \lif X(u(x)))]) \lif \lforall[x][X(x)])]]]
\]
په جوړښت~\LR{$\Struct{M}$} کښې هله او يوازې هله رښتيا ده چې
\LR{$\Domain{M}$} د شمېر وړ وي؛ نو \LR{$\lnot\fn{Count}$} په~\LR{$\Struct{M}$}
کښې هله او يوازې هله رښتيا ده چې \LR{$\Domain{M}$} د شمېر نۀ وړ وي۔
داسې جوړښتونه شته: هر د شمېر نۀ وړ سټ د د تعريف ساحه په توګه
واخلئ، لکه \LR{$\Pow{\Nat}$} يا~\LR{$\Real$}۔ نو \LR{$\lnot \fn{Count}$} نامتناهي
مدلونه لري، خو د شمېر وړ مدلونه نۀ لري۔
\end{proof}
\olpseganchor{OLP-0335-B008}{end}

\olpseganchor{OLP-0335-B009}{start}
\begin{thm}
داسې تړلي فارمولونه شته چې شمېرېدونکے لامتناهي مدلونه لري، خو
د شمېر نۀ وړ مدلونه نۀ لري۔
\end{thm}
\olpseganchor{OLP-0335-B009}{end}

\olpseganchor{OLP-0335-B010}{start}
\begin{proof}
\LR{$\fn{Count} \land \fn{Inf}$} په~\LR{$\Nat$} کښې رښتيا ده، خو په هېڅ داسې
  جوړښت~\LR{$\Struct{M}$} کښې نۀ ده رښتيا چې \LR{$\Domain{M}$} ئې
  د شمېر نۀ وړ وي۔
\end{proof}
\olpseganchor{OLP-0335-B010}{end}

\olpseganchor{OLP-0336-B004}{start}
\olchapter{sol}{set}{دويمې درجې منطق او د سټونو تيوري}
\olpseganchor{OLP-0336-B004}{end}

\olpseganchor{OLP-0336-B005}{start}
\begin{editorial}
دا څپرکے په دويمې درجې منطق کښې د فرعي سټونو د سټونو او د پيوستار د
کوډ کولو په هکله دے۔ پايلې بيان شوې دي، خو ثبوتونه ئې لا نۀ دي
بشپړ شوي۔ تر اوسه مسئلې هم نشته، او کېدے شي پخپلو تعريفونو او
پايلو کښې هم ستونزې وي۔ په احتياط ترې کار واخلئ او هر ناسم يا
ناڅرګند ټکے راپور کړئ۔
\end{editorial}
\olpseganchor{OLP-0336-B005}{end}

\olpseganchor{OLP-0337-B004}{start}
\olfileid{sol}{set}{int}
\olpseganchor{OLP-0337-B004}{end}

\olpseganchor{OLP-0337-B005}{start}
\olsection{سريزه}
\olpseganchor{OLP-0337-B005}{end}

\olpseganchor{OLP-0337-B006}{start}
څرنګه چې دويمې درجې منطق د تعريف ساحې پر فرعي سټونو او هم پر
تابعو کميت اچولے شي، تمه کېږي چې د سټونو د تيورۍ لږ تر لږه يوه
برخه په دې منطق کښې بيانېدے شي۔ له «بيانېدو» څخه مو مطلب دا دے
چې د سټونو د تيورۍ خاصيتونه او ويناوې په دويمې درجې منطق کښې،
بې له دې چې د سټونو دپاره کوم ځانګړے غيرمنطقي لغت وکاروو
(لکه د سټونو د تيورۍ د غړيتوب پريديکات)، څرګندولے شو۔
د بېلګې په توګه، د سټونو اتحاد، اشتراک او د فرعي سټ اړيکه
تعريفولے شو؛ د سټونو اندازې هم سره پرتله کولے شو او د
کانتور د قضيې غوندې پايلې بيانولے شو۔
\olpseganchor{OLP-0337-B006}{end}

\olpseganchor{OLP-0338-B004}{start}
\olfileid{sol}{set}{cmp}
\olpseganchor{OLP-0338-B004}{end}

\olpseganchor{OLP-0338-B005}{start}
\olsection{د سټونو پرتله کول}
\olpseganchor{OLP-0338-B005}{end}

\olpseganchor{OLP-0338-B006}{start}
\begin{prop}
فارمول \LR{$\lforall[x][(X(x) \lif Y(x))]$} د فرعي سټ اړيکه
تعريفوي؛ يعنې \LR{$\Sat{M}{\lforall[x][(X(x) \lif Y(x))]}[s]$} هله او
يوازې هله چې \LR{$s(X) \subseteq s(Y)$}۔
\end{prop}
\olpseganchor{OLP-0338-B006}{end}

\olpseganchor{OLP-0338-B007}{start}
\begin{prop}
فارمول \LR{$\lforall[x][(X(x) \liff Y(x))]$} د سټونو د يووالي
اړيکه تعريفوي؛ يعنې \LR{$\Sat{M}{\lforall[x][(X(x) \liff Y(x))]}[s]$}
هله او يوازې هله چې \LR{$s(X) = s(Y)$}۔
\end{prop}
\olpseganchor{OLP-0338-B007}{end}

\olpseganchor{OLP-0338-B008}{start}
\begin{prop}
فارمول \LR{$\lexists[x][X(x)]$} د ناتشوالي خاصيت تعريفوي؛ يعنې
\LR{$\Sat{M}{\lexists[x][X(x)]}[s]$} هله او يوازې هله چې \LR{$s(X) \neq
\emptyset$}۔
\end{prop}
\olpseganchor{OLP-0338-B008}{end}

\olpseganchor{OLP-0338-B009}{start}
سټ~\LR{$X$} له سټ~\LR{$Y$} څخه لوے نۀ دے، \LR{$\cardle{X}{Y}$}، هله او يوازې
هله چې يو پر يو تابع \LR{$f\colon X \to Y$} وي۔ دا چې د تابع
يو پر يو والے او په~\LR{$X$} کښې د ورودي غړو د ارزښتونو په~\LR{$Y$}
کښې شتون څرګندولے شو، د تعريف ساحې د فرعي سټونو د اندازې د
«لوے نۀ دے» اړيکه هم تعريفولے شو۔
\olpseganchor{OLP-0338-B009}{end}

\olpseganchor{OLP-0338-B010}{start}
\begin{prop}
لاندې فورمول
\[
\lexists[u][(\lforall[x][(X(x) \lif Y(u(x)))] \land \lforall[x][\lforall[y][((X(x)\land X(y)\land \eq[u(x)][u(y)]) \lif
      \eq[x][y])]])]
\]
د «لوے نۀ دے» اړيکه تعريفوي۔ د تابع محدودول لومړي سټ ته يو پر يو داخلونه ورکوي۔ له بلې خوا، د لومړي سټ يوه داخلونه د ټولې تعريف ساحې تابع ته په دې ډول غځوو چې له سټ څخه بهر هر غړے يوه ټاکلې پايلې ته ولېږو۔ د تعريف ساحه ناتشه ده۔
\end{prop}
\olpseganchor{OLP-0338-B010}{end}

\olpseganchor{OLP-0338-B011}{start}
دوه سټونه يو شان اندازه لري، يا «همشمېره» دي، \LR{$\cardeq{X}{Y}$}،
هله او يوازې هله چې دوه اړخيزه يو پر يو تابع~\LR{$f\colon X \to Y$} وي۔
\olpseganchor{OLP-0338-B011}{end}

\olpseganchor{OLP-0338-B012}{start}
\begin{prop}
لاندې فورمول
\begin{multline*}
  \lexists[u][(\lforall[x][(X(x) \lif Y(u(x)))] \land {}\\
    \lforall[x][\lforall[y][((X(x) \land X(y) \land \eq[u(x)][u(y)]) \lif \eq[x][y])]]
  \land {} \\\lforall[y][(Y(y) \lif \lexists[x][(X(x)
      \land \eq[y][u(x)])])])]
\end{multline*}
د همشمېره‌والي اړيکه تعريفوي۔
\end{prop}
\olpseganchor{OLP-0338-B012}{end}

\olpseganchor{OLP-0338-B013}{start}
دا فارمولونه به په ترتيب سره په \LR{$X \subseteq Y$}، \LR{$X = Y$}،
\LR{$X \neq \emptyset$}، \LR{$\cardle{X}{Y}$} او \LR{$\cardeq{X}{Y}$} لنډوو۔
(دا لږ څه ګډوډوونکې کېدے شي، ځکه د سټونو~\LR{$X$} او~\LR{$Y$} په
غيررسمي خبرو کښې هم همدا نښې کاروو؛ خو دلته دا نښې د دويمې درجې
منطق د هغو فارمولونه لنډې بڼې دي چې يوځايي اړيکيز متغيرونه
\LR{$X$} او~\LR{$Y$} پکښې راځي۔)
\olpseganchor{OLP-0338-B013}{end}

\olpseganchor{OLP-0338-B014}{start}
\begin{prop}
تړلے فارمول \LR{$\lforall[X][\lforall[Y][((\cardle{X}{Y} \land
    \cardle{Y}{X}) \lif \cardeq{X}{Y})]]$} معتبره ده۔
\end{prop}
\olpseganchor{OLP-0338-B014}{end}

\olpseganchor{OLP-0338-B015}{start}
\begin{proof}
دا تړلے فارمول په جوړښت~\LR{$\Struct{M}$} کښې هله صدق کوي
چې د هر دوو فرعي سټونو \LR{$X \subseteq \Domain{M}$} او
\LR{$Y \subseteq \Domain{M}$} دپاره، که \LR{$\cardle{X}{Y}$} او
\LR{$\cardle{Y}{X}$} وي، نو \LR{$\cardeq{X}{Y}$} هم وي۔ خو دا د
\emph{هر ډول} سټونو~\LR{$X$} او~\LR{$Y$} دپاره رښتيا ده؛ همدا د
شرودرـبرنشتاين قضيه ده۔
\end{proof}
\olpseganchor{OLP-0338-B015}{end}

\olpseganchor{OLP-0339-B004}{start}
\olfileid{sol}{set}{crd}
\olpseganchor{OLP-0339-B004}{end}

\olpseganchor{OLP-0339-B005}{start}
\olsection{د سټونو شمېرنيزې اندازې}
\olpseganchor{OLP-0339-B005}{end}

\olpseganchor{OLP-0339-B006}{start}
\begin{explain}
لکه څنګه چې د تعريف ساحې متناهي يا نامتناهي، د شمېر وړ يا
د شمېر نۀ وړ والے څرګندولے شو، همداسې د~\LR{$\Domain{M}$} د يوه
فرعي سټ متناهي يا نامتناهي، د شمېر وړ يا د شمېر نۀ وړ
والے هم تعريفولے شو۔
\end{explain}
\olpseganchor{OLP-0339-B006}{end}

\olpseganchor{OLP-0339-B007}{start}
\begin{prop}
فورمول \LR{$\fn{Inf}(X) \ident$}
\begin{multline*}
\lexists[u][(\lforall[x][\lforall[y][((X(x)\land X(y)\land \eq[u(x)][u(y)]) \lif
      \eq[x][y])]] \land {} \\\lforall[x][(X(x) \lif X(u(x)))] \land {} \\
\lexists[y][(X(y) \land \lforall[x][(X(x)
      \lif \eq/[y][u(x)])])])]
\end{multline*}
د متغيرونو د ارزښت‌ټاکنې~\LR{$s$} له مخې هله او يوازې هله صدق کوي
چې \LR{$s(X)$} نامتناهي وي، د ټاکنې د اصل تر معمول فرض لاندې۔ په متناهي سټ کښې هره يو پر يو ځانته تابع پر سر ده۔ په نامتناهي سټ کښې يوه د شمېر وړ نامتناهي لړۍ غوره کوو، د هغې هر غړے خپل تالي ته لېږو، او نور غړي پر خپل ځاے پرېږدو۔ لومړنے غړے په انځور کښې نۀ راځي۔
\end{prop}
\olpseganchor{OLP-0339-B007}{end}

\olpseganchor{OLP-0339-B008}{start}
\begin{prop}
فورمول \LR{$\fn{Count}(X) \ident $}
\begin{multline*}
\lnot \lexists[x][X(x)] \lor {} \\
\lexists[z][\lexists[u][(X(z) \land
    \lforall[x][(X(x) \lif X(u(x)))] \land {} \\ \lforall[Y][((Y(z) \land
      \lforall[x][(Y(x) \lif Y(u(x)))]) \lif X \subseteq Y)])]]
\end{multline*}
د متغيرونو د ارزښت‌ټاکنې~\LR{$s$} له مخې هله او يوازې هله صدق کوي
چې \LR{$s(X)$} د شمېر وړ وي۔ په ناتش حالت کښې د ټاکلي لومړي غړي مدار د تابعې پر وړاندې تړلے دے، نو ټول سټ رانغاړي؛ د ځانته تابعې له شرطه مدار پخپله په ټاکلي سټ کښې دے۔ برعکس، متناهي سټ په يوه لړۍ کښې راوړئ او وروستے غړے پر خپل ځاے پرېږدئ۔ په نامتناهي شمېر وړ سټ کښې هر غړے په لړۍ کښې خپل تالي ته ولېږئ۔
\end{prop}
\olpseganchor{OLP-0339-B008}{end}

\olpseganchor{OLP-0339-B009}{start}
د کانتور له قضيې پوهېږو چې د شمېر نۀ وړ سټونه شته، او د
نامتناهي اندازو بېلابېلې ډېرې مرتبې هم شته۔ د سټونو تيوري
د سټونو د اندازو بشپړ حساب جوړوي او سټونو ته نامتناهي
شمېرنيز عددونه ټاکي۔ طبيعي عددونه د متناهي سټونو اندازې
اندازه کوي۔ د شمېرېدونکے لامتناهي سټونو اندازه لومړۍ نامتناهي
شمېرنيزه اندازه ده چې~\LR{$\aleph_0$} («الف‌صفر») بلل کېږي۔
ورپسې نامتناهي اندازه~\LR{$\aleph_1$} ده۔ دا هغه تر ټولو کوچنۍ
اندازه ده چې يو سټ ئې لرلے شي خو د شمېر وړ نۀ وي (يعنې
اندازه ئې~\LR{$\aleph_0$} نۀ وي)۔ «\LR{$X$} د~\LR{$\aleph_0$} اندازه لري»
داسې تعريفولے شو: \LR{$\fn{Aleph}_0(X) \liff \fn{Inf}(X) \land \fn{Count}(X)$}۔
\LR{$X$} د~\LR{$\aleph_1$} اندازه هله لري چې پخپله د شمېر وړ نۀ وي،
خو هر هغه فرعي سټ ئې چې له ټاکلي سټ څخه په شمېرنيزه اندازه
کوچنے وي، متناهي يا د~\LR{$\aleph_0$} اندازه ولري۔ نو دا په
دې توګه د~\LR{$\aleph_0$} تر کچې د کوچنيو اندازو شرط دے۔ دا په
لاندې فورمول څرګندېږي:
\LR{$\fn{Aleph_1}(X) \ident \lnot \fn{Count}(X) \land
\lforall[Y][((Y \subseteq X \land \lnot \cardeq{X}{Y}) \lif
  (\lnot\fn{Inf}(Y) \lor \fn{Aleph}_0(Y)))]$}۔ د~\LR{$\aleph_2$}
اندازه هم په ورته ډول تعريفېږي۔
\olpseganchor{OLP-0339-B009}{end}

\olpseganchor{OLP-0339-B010}{start}
يوه اندازه ځانګړې پاملرنه غواړي: د پيوستار شمېرنيزه اندازه۔
دا د~\LR{$\Pow{\Nat}$} اندازه ده، يا په معادل ډول د~\LR{$\Real$} اندازه۔
دا چې يو سټ د پيوستار اندازه لري، په دويمې درجې منطق کښې هم
څرګندېږي، خو لږ زيات کار غواړي۔
\olpseganchor{OLP-0339-B010}{end}

\olpseganchor{OLP-0340-B004}{start}
\olfileid{sol}{set}{pow}
\olpseganchor{OLP-0340-B004}{end}

\olpseganchor{OLP-0340-B005}{start}
\olsection{د پيوستار اندازه}
\olpseganchor{OLP-0340-B005}{end}

\olpseganchor{OLP-0340-B006}{start}
\begin{explain}
په دويمې درجې منطق کښې د د تعريف ساحه پر فرعي سټونو کميت اچولے شو،
خو د د تعريف ساحه د فرعي سټونو پر سټونو ئې نۀ شو اچولے۔ دا کار
مستقيم ترسره کولو ته \emph{د درېيمې درجې} منطق پکار دے۔ د بېلګې
په توګه، که د کانتور دا قضيه بيانول غواړو چې د يوه سټ د ټولو
فرعي سټونو له سټ څخه هماغه سټ ته هېڅ يو پر يو تابع نشته،
کېدے شي داسې ئې ووايو: «د هر سټ~\LR{$X$} او هر سټ~\LR{$P$} دپاره، که
\LR{$P$} د~\LR{$X$} د ټولو فرعي سټونو سټ وي، نو \LR{$\cardle{P}{X}$} نۀ دے»۔
خو دا چې~\LR{$P$} د~\LR{$X$} د ټولو فرعي سټونو سټ دے، بايد دا هم
رسمي کړو چې د~\LR{$P$} غړي د~\LR{$X$} ټول او يوازې فرعي
سټونه دي، مثلاً \LR{$\lforall[Y][(P(Y) \liff Y \subseteq X)]$}۔
ستونزه په~\LR{$P(Y)$} کښې ده: دا د دويمې درجې منطق فارمول
نۀ ده، ځکه د~\LR{$P$} غوندې يوځايي اړيکيز متغير ته يوازې
ترمونه د ارګومېنټونو په توګه ورکولے شو۔
\olpseganchor{OLP-0340-B006}{end}

\olpseganchor{OLP-0340-B007}{start}
خو که د تعريف ساحه کافي لويه وي، د سټونو پر سټونو کميت
اچول \emph{تمثيلولے} شو۔ دا له دې خبرې ګټه اخلي چې دوه‌ځايي
اړيکه~\LR{$R$} د د تعريف ساحه غړي د همدې د تعريف ساحه
له غړي سره تړي۔ د داسې~\LR{$R$} دپاره ټول هغه غړي
چې له يوه~\LR{$x$} سره په~\LR{$R$} تړلي وي، يوځاے کولے شو:
\LR{$\Setabs{y \in
  \Domain{M}}{R(x,y)}$} هغه سټ دے چې~\LR{$x$} ئې «کوډوي»۔
برعکس، که \LR{$Z \subseteq \Pow{\Domain{M}}$} د~\LR{$\Domain{M}$}
د فرعي سټونو يوه ټولګه وي او په~\LR{$\Domain{M}$} کښې لږ تر لږه
د~\LR{$Z$} د سټونو هومره غړي وي، نو اړيکه
\LR{$R \subseteq \Domain{M}^2$} هم شته چې هر \LR{$Y \in Z$}
د کوم~\LR{$x$} په وسيله د~\LR{$R$} له مخې کوډ شي۔
\end{explain}
\olpseganchor{OLP-0340-B007}{end}

\olpseganchor{OLP-0340-B008}{start}
\begin{defn}
که \LR{$R \subseteq \Domain{M}^2$} وي، نو \emph{\LR{$x$} د~\LR{$R$} له مخې}
سټ \LR{$\Setabs{y
  \in \Domain{M}}{R(x, y)}$} \emph{کوډوي}۔
\end{defn}
\olpseganchor{OLP-0340-B008}{end}

\olpseganchor{OLP-0340-B009}{start}
که غړے~\LR{$x \in \Domain{M}$} د~\LR{$R$} له مخې سټ
\LR{$Z \subseteq \Domain{M}$} کوډوي، نو سټ \LR{$Y \subseteq \Domain{M}$}
د سټونو يو سټ کوډوي: هغه سټونه چې د~\LR{$Y$} غړي ئې
کوډوي۔ نو سټ~\LR{$Y$} کولے شي د~\LR{$R$} له مخې~\LR{$\Pow{X}$} کوډ کړي۔
دا هله کېږي چې د هر \LR{$Z \subseteq X$} دپاره کوم \LR{$x \in Y$}
د~\LR{$R$} له مخې~\LR{$Z$} کوډ کړي، او هر \LR{$x \in Y$} د~\LR{$R$} له مخې
کوم فرعي سټ~\LR{$Z \subseteq X$} کوډ کړي۔
\olpseganchor{OLP-0340-B009}{end}

\olpseganchor{OLP-0340-B010}{start}
\begin{prop}
فارمول
  \[
  \fn{Codes}(x, R, Z) \ident \lforall[y][(Z(y) \liff
    R(x, y))]
  \]
دا څرګندوي چې \LR{$s(x)$} د~\LR{$s(R)$} له مخې~\LR{$s(Z)$} کوډوي۔
فارمول
\begin{multline*}
  \fn{Pow}(Y, R, X) \ident \\
  \lforall[Z][(Z \subseteq X \lif \lexists[x][(Y(x) \land
    \fn{Codes}(x, R, Z))])] \land {} \\ \lforall[x][(Y(x) \lif
  \lforall[Z][(\fn{Codes}(x, R, Z) \lif Z \subseteq X)])]
\end{multline*}
دا څرګندوي چې \LR{$s(Y)$} د~\LR{$s(R)$} له مخې د~\LR{$s(X)$} د فرعي
سټونو سټ کوډوي؛ يعنې د~\LR{$s(Y)$} غړي د~\LR{$s(R)$} له مخې
د~\LR{$s(X)$} ټول او يوازې فرعي سټونه کوډوي۔
\end{prop}
\olpseganchor{OLP-0340-B010}{end}

\olpseganchor{OLP-0340-B011}{start}
\begin{explain}
په دې طريقه د فرعي سټونو د خپلو کوډونو پر سر د کميت
اچولو له لارې د فرعي سټونو د سټ په هکله ويناوې څرګندولے
شو۔ د بېلګې په توګه، د کانتور قضيه اوس داسې ويلے شو
چې د کوډونو له داسې سټ څخه، چې د~\LR{$X$} د
ټولو فرعي سټونو سټ کوډوي، پخپله~\LR{$X$} ته هېڅ
يو پر يو تابع نشته۔
\end{explain}
\olpseganchor{OLP-0340-B011}{end}

\olpseganchor{OLP-0340-B012}{start}
\begin{prop}
لاندې جمله
\begin{multline*}
  \lforall[X][\lforall[Y][\lforall[R][(\fn{Pow}(Y, R, X) \lif \\
      \lnot \lexists[u][(\lforall[x][\lforall[y][((Y(x)\land Y(y)\land \eq[u(x)][u(y)]) \lif
            \eq[x][y])]] \land {}\\
        \lforall[x][(Y(x) \lif X(u(x)))])])]]]
\end{multline*}
معتبره ده۔ فرض کړئ داسې يوه داخلونه شته۔ د هغې پر انځور هر غړي ته هغه فرعي سټ ورکړئ چې د هغه غړي يوازينے مخکينے کوډ ئې ټاکي، او د انځور له باندې تش سټ ورکړئ۔ له ټاکلي سټ څخه د هغه د ټولو فرعي سټونو پر سټ يوه پر سر تابع ترلاسه کېږي۔ خو هغه قطري فرعي سټ چې کټ مټ هغه غړي لري چې د خپل ټاکل شوي فرعي سټ غړي نۀ دي، د دې تابعې په انځور کښې نۀ شي راتللے۔
\end{prop}
\olpseganchor{OLP-0340-B012}{end}

\olpseganchor{OLP-0340-B013}{start}
\begin{explain}
د شمېرېدونکے لامتناهي سټ د ټولو فرعي سټونو سټ
د شمېر نۀ وړ دے، نو اندازه ئې د هر شمېرېدونکے لامتناهي
سټ تر اندازې لوړه ده (چې~\LR{$\aleph_0$} ده)۔ د~\LR{$\Pow{\Nat}$}
اندازې ته «د پيوستار اندازه» ويل کېږي، ځکه دا د حقيقي
عددونو پر کرښه د نقطو د سټ~\LR{$\Real$} له اندازې سره برابره
ده۔ که د تعريف ساحه دومره لويه وي چې د کوم
شمېرېدونکے لامتناهي سټ د ټولو فرعي سټونو سټ کوډ کړے شي،
نو د پيوستار په اندازه سټ داسې بيانولے شو چې هغه له
هر داسې سټ~\LR{$Y$} سره همشمېره وي چې د~\LR{$\aleph_0$}
په اندازه سټ~\LR{$X$} د ټولو فرعي سټونو سټ کوډوي او بېلابېل
غړي ئې بېلابېل فرعي سټونه کوډوي۔
(که د تعريف ساحه دومره لويه نۀ وي، يعنې د~\LR{$\Real$}
سره همشمېره فرعي سټ ونۀ لري، نو د~\LR{$\Pow{X}$} د کوډ
کولو اړيکه هم نۀ شي لرلے۔)
\end{explain}
\olpseganchor{OLP-0340-B013}{end}

\olpseganchor{OLP-0340-B014}{start}
\begin{prop}
که \LR{$\cardle{\Real}{\Domain{M}}$} وي، نو فارمول
\begin{multline*}
\fn{Cont}(Y) \ident
\lexists[X][\lexists[R][((\fn{Aleph}_0(X) \land
      \fn{Pow}(Y, R, X)) \land \\
      \lforall[x][\lforall[y][((Y(x) \land {}
      Y(y) \land \lforall[z][R(x,z) \liff R(y,z)]) \lif \eq[x][y])]])]]
\end{multline*}
دا څرګندوي چې \LR{$\cardeq{s(Y)}{\Real}$}۔
\end{prop}
\olpseganchor{OLP-0340-B014}{end}

\olpseganchor{OLP-0340-B015}{start}
\begin{proof}
\LR{$\fn{Pow}(Y, R, X)$} دا څرګندوي چې \LR{$s(Y)$} د~\LR{$s(R)$} له مخې
د~\LR{$s(X)$} د ټولو فرعي سټونو سټ کوډوي، او
\LR{$\fn{Aleph}_0(X)$} وايي چې ټاکلے سټ د شمېر وړ نامتناهي
سټ دے۔ نو د~\LR{$s(Y)$} اندازه لږ تر لږه د پيوستار هومره ده؛
کېدے شي لويه هم وي که د~\LR{$s(Y)$} ډېر غړي د~\LR{$X$}
يو فرعي سټ کوډ کړي۔ وروستے عطف دا حالت ردوي: هغه غواړي
چې د~\LR{$s(Y)$} د غړو او د~\LR{$s(X)$} د فرعي سټونو ترمنځ د~\LR{$s(R)$}
له مخې اړيکه يو پر يو وي۔ برعکس، که د کوډونو سټ د پيوستار په اندازه وي، د تعريف په ساحه کښې يو شمېر وړ نامتناهي سټ غوره کړو۔ د کوډونو سټ د هغه له ټولو فرعي سټونو سره يو پر يو او پر سر برابر کړو او اړونده کوډوونکې اړيکه له همدې برابرۍ تعريف کړو۔
\end{proof}
\olpseganchor{OLP-0340-B015}{end}

\olpseganchor{OLP-0340-B016}{start}
\begin{prop}
\LR{$\cardeq{\Domain{M}}{\Real}$} هله او يوازې هله چې
\begin{multline*}
  \Sat{M}{\lexists[Y][(\fn{Cont}(Y) \land {} \\
    \lexists[u][(\lforall[x][Y(u(x))] \land {}
      \lforall[x][\lforall[y][(\eq[u(x)][u(y)] \lif \eq[x][y])]] \land {} \\
      \lforall[y][(Y(y) \lif \lexists[x][\eq[y][u(x)]])])])]}.
\end{multline*}
\end{prop}
\olpseganchor{OLP-0340-B016}{end}

\olpseganchor{OLP-0340-B017}{start}
\begin{explain}
د پيوستار فرضيه وايي چې د پيوستار اندازه لومړۍ
د شمېر نۀ وړ شمېرنيزه اندازه ده؛ يعنې
\LR{$\Pow{\Nat}$} د~\LR{$\aleph_1$} اندازه لري۔
\end{explain}
\olpseganchor{OLP-0340-B017}{end}

\olpseganchor{OLP-0340-B018}{start}
\begin{prop}
د پيوستار فرضيه هله او يوازې هله رښتيا ده چې
\[\fn{CH} \ident
\lforall[X][(\fn{Aleph}_1(X) \liff \fn{Cont}(X))]\]
معتبره وي۔
\end{prop}
\olpseganchor{OLP-0340-B018}{end}

\olpseganchor{OLP-0340-B019}{start}
دا نۀ ده رښتيا چې که د پيوستار فرضيه غلطه وي، نو
\LR{$\lnot \fn{CH}$} به معتبره وي، يا برعکس۔ په
د شمېر وړ تعريف ساحه کښې نۀ د~\LR{$\aleph_1$} په
اندازه فرعي سټ شته او نۀ د پيوستار په اندازه؛ نو
\LR{$\fn{CH}$} په هره د شمېر وړ تعريف ساحه کښې
رښتيا وي۔ خو داسې بله جمله ورکولے شو چې هله او يوازې
هله معتبره وي چې د پيوستار فرضيه غلطه وي:
\olpseganchor{OLP-0340-B019}{end}

\olpseganchor{OLP-0340-B020}{start}
\begin{prop}
د پيوستار فرضيه هله او يوازې هله غلطه ده چې
\[\fn{NCH} \ident
\lforall[X][(\fn{Cont}(X) \lif \lexists[Y][(Y \subseteq X \land \lnot
    \fn{Count}(Y) \land \lnot \cardeq{X}{Y})])]\]
معتبره وي۔
\end{prop}
\olpseganchor{OLP-0340-B020}{end}

\olpseganchor{OLP-0341-B004}{start}
\olpart{lam}{د لامبډا حساب}
\olpseganchor{OLP-0341-B004}{end}

\olpseganchor{OLP-0341-B005}{start}
\begin{editorial}
دا برخه د لامبډا حساب څېړي۔ د پېژندنې څپرکے د \LR{\latinfont Jeremy Avigad} پر
يادښتونو ولاړ دے؛ د دې څپرکي يوه برخه اوس تکراري ده او په
ورپسې څپرکو کښې هم راغلې ده۔ د نحو، د چرچ--روسر خاصيت او د لامبډا له مخې
د تعريفېدنې څپرکي \LR{\latinfont Zesen Qian} د خپل Mitacs د اوړي د کارآموزۍ پر مهال
چمتو کړي وو۔ د هغو څپرکو بياکتنه او سمونه لا پاتې ده۔
\end{editorial}
\olpseganchor{OLP-0341-B005}{end}

\olpseganchor{OLP-0342-B004}{start}
\olchapter{lam}{int}{پېژندنه}
\olpseganchor{OLP-0342-B004}{end}

\olpseganchor{OLP-0342-B005}{start}
\begin{editorial}
دا څپرکے د لامبډا حساب په هکله د \LR{\latinfont Jeremy Avigad} له اصلي لنډو
يادښتونو جوړ دے۔ برخې ئې يوځاے کول غواړي، او د لامبډا له مخې د
تعريفېدنې مواد ئې بايد د همدې موضوع له بېل، ډېر تفصيلي څپرکي سره
يوځاے شي۔
\end{editorial}
\olpseganchor{OLP-0342-B005}{end}

\olpseganchor{OLP-0343-B004}{start}
\olfileid{lam}{int}{ovr}
\olsection{لنډه کتنه}
\olpseganchor{OLP-0343-B004}{end}

\olpseganchor{OLP-0343-B005}{start}
الونزو چرچ په لومړيو ۱۹۳۰مو کلونو کښې لامبډا حساب په اصل کښې د
تعميري منطق د بنسټ په توګه جوړ کړے و، \emph{نۀ} د محاسبه کېدونکو
تابعو د نمونې په توګه۔ خو ډېر ژر وښودل شو چې دا د محاسبه کېدنې له
نورو تعريفونو، لکه د ټيورينګ له مخې محاسبه کېدونکو تابعو او جزوي
بازګشتي تابعو، سره معادل دے۔ دا چې دغه معادلتوب په لومړي سر کښې يو
څه حيرانوونکے و، د محاسبه کېدنې دا پېژندنه لا په زړه پورې کوي۔
\olpseganchor{OLP-0343-B005}{end}

\olpseganchor{OLP-0343-B006}{start}
د لامبډا نښه د تابعې د نوم له جلا تعريفولو پرته، د هغې په تعريفوونکي
نښيز عبارت سره نېغې تابعې ته د اشارې يوه آسانه لاره ده۔ د دې پر ځاے
چې ووايو ``\LR{$f$} هغه تابع وګڼئ چې \LR{$f(x) = x + 3$} ئې تعريفوي''، ويلے
شو ``\LR{$f$} هغه تابع وګڼئ چې \LR{$\lambd[x][(x + 3)]$} ئې تعريفوي''۔ په بله
وينا، \LR{$\lambd[x][(x+3)]$} يوازې د هغې تابعې يو \emph{نوم} دے چې په
خپل آرګومېنټ درې ورزياتوي۔ په دې عبارت کښې \LR{$x$} يو تش‌ځايے متغير دے:
هماغه تابع په \LR{$\lambd[y][(y + 3)]$} هم ښودل کېدے شي۔ دا نښه د نورو
پاراميټرونو په شتون کښې هم کار کوي۔ مثلاً، که \LR{$g(x, y)$} يوه دوه‌ځاييزه
تابع او \LR{$k$} يو طبيعي عدد وي، نو \LR{$\lambd[x][g(x,k)]$} هغه تابع ده چې
هر~\LR{$x$}،~\LR{$g(x, k)$} ته وړي۔
\olpseganchor{OLP-0343-B006}{end}

\olpseganchor{OLP-0343-B007}{start}
له نښيز عبارت څخه د تابعې د تعريف دغه طريقه \emph{لامبډا تجريد}
بلل کېږي۔ د لامبډا تجريد بل اړخ \emph{تطبيق} دے: که يوه تابع~\LR{$f$}
ولرو (مثلاً پر طبيعي عددونو تعريف شوې)، پر هر قيمت، لکه ۲، ئې تطبيقولے
شو۔ په دوديزه نښه کښې د پايلې دپاره البته~\LR{$f(2)$} ليکو۔
\olpseganchor{OLP-0343-B007}{end}

\olpseganchor{OLP-0343-B008}{start}
که لامبډا تجريد او تطبيق سره يوځاے کړو، څۀ پېښېږي؟ بيا حاصل عبارت
په دې ساده کېدے شي چې د تجريد شوي متغير په ځاے د تابعې ورکړل شوے
آرګومېنټ کېږدو۔ د بېلګې په توګه،
\[
(\lambd[x][(x + 3)])(2)
\]
تر~\LR{$2 + 3$} پورې ساده کېدے شي۔
\olpseganchor{OLP-0343-B008}{end}

\olpseganchor{OLP-0343-B009}{start}
تر دې ځايه مو دوديزو مفاهيمو ته يوازې نوې نښې ورکړې دي۔ خو لامبډا
حساب د سټونو د تيورۍ له ليدلوري څخه ډېر بنسټيز بېلتون لري۔ په دې
چوکاټ کښې:
\begin{enumerate}
\item هر څۀ يوه تابع ښيي۔
\item تابعې د لامبډا تجريد په وسيله تعريفېدے شي۔
\item هر څۀ پر هر بل څيز تطبيقېدے شي۔
\end{enumerate}
مثلاً، که \LR{$F$} د لامبډا حساب يو ترم وي، \LR{$F(F)$} تل د معنٰی لرونکے
ګڼل کېږي۔ دې پراخ چوکاټ ته \emph{بې‌ټايپه} لامبډا حساب وايي؛ دلته
``بې‌ټايپه'' معنا دا ده چې پر څۀ د څۀ شي د تطبيقولو محدوديت نشته۔
\olpseganchor{OLP-0343-B009}{end}

\olpseganchor{OLP-0343-B010}{start}
\begin{digress}
يو \emph{ټايپ‌لرونکے} لامبډا حساب هم شته چې د بې‌ټايپې بڼې يو مهم
بدلون دے۔ که څۀ هم ټايپ‌لرونکے لامبډا حساب په ډېرو اړخونو کښې
بې‌ټايپې بڼې ته ورته دے، له کلاسيکي سټي چوکاټ سره ئې سمون ډېر آسان
دے او ځينې خاصيتونه ئې ډېر توپير لري۔
\olpseganchor{OLP-0343-B010}{end}

\olpseganchor{OLP-0343-B011}{start}
د لامبډا حساب څېړنه په نظري کمپيوټر ساينس او د پروګرامي ژبو په
طرحه کښې د مرکزي اهميت لرونکې ثابته شوې ده۔ LISP، چې جان مک‌کارتي په ۱۹۵۰مو
کلونو کښې جوړه کړې وه، د داسې مفکورو تر اغېز لاندې د يوې لومړنۍ
پروګرامي ژبې بېلګه ده۔
\end{digress}
\olpseganchor{OLP-0343-B011}{end}

\olpseganchor{OLP-0344-B004}{start}
\olfileid{lam}{int}{syn}
\olsection{د لامبډا حساب نحو}
\olpseganchor{OLP-0344-B004}{end}

\olpseganchor{OLP-0344-B005}{start}
لومړے د متغيرونو يوه لړۍ \LR{$x$}، \LR{$y$}، \LR{$z$}،~\dots او ځينې ثابتې نښې
\LR{$a$}، \LR{$b$}، \LR{$c$}،~\dots اخلو۔ د ترمونو سټ په استقرايي ډول داسې
تعريفېږي:
\begin{enumerate}
\item هر متغير يو ترم دے۔
\item هر ثابت يو ترم دے۔
\item که \LR{$M$} او \LR{$N$} ترمونه وي، \LR{$(MN)$} هم ترم دے۔
\item که \LR{$M$} ترم او \LR{$x$} متغير وي، \LR{$(\lambd[x][M])$} هم ترم دے۔
\end{enumerate}
د \LR{$(MN)$} په بڼه ترمونه \emph{تطبيقونه} او د
\LR{$(\lambd[x][M])$} په بڼه ترمونه \emph{تجريدونه} بلل کېږي۔
\olpseganchor{OLP-0344-B005}{end}

\olpseganchor{OLP-0344-B006}{start}
هغه نظام چې هېڅ ثابت نۀ لري، \emph{خالص} لامبډا حساب بلل کېږي۔ موږ به
زياتره په خالص \LR{$\lambd$}-حساب کښې کار کوو؛ نو ټول واړه لاتيني توري
به متغيرونه ښيي۔ د لويو لاتيني تورو (\LR{$M$}، \LR{$N$} او داسې نورو) په وسيله
د \LR{$\lambd$}-حساب ترمونه ښيو۔
\olpseganchor{OLP-0344-B006}{end}

\olpseganchor{OLP-0344-B007}{start}
د نښو څو دودونه به خپل کړو:
\olpseganchor{OLP-0344-B007}{end}

\olpseganchor{OLP-0344-B008}{start}
\begin{conv}
\begin{enumerate}
\item چې قوسونه پرېښودل شي، تطبيق له کيڼ لوري ښي لوري ته ترسره کېږي۔
  مثلاً، که \LR{$M$}، \LR{$N$}، \LR{$P$} او \LR{$Q$} ترمونه وي، نو \LR{$MNPQ$} د
  \LR{$(((MN)P)Q)$} لنډه بڼه ده۔
\item همداسې، چې قوسونه پرېښودل شي، لامبډا تجريد ته تر وسه پراخه
  ساحه ورکوو۔ مثلاً، \LR{$\lambd[x][MNP]$} د
  \LR{$(\lambd[x][((MN)P)])$} په معنا لولو۔
\item يوه لامبډا د څو متغيرونو د تجريد دپاره کارېدے شي۔ مثلاً،
  \LR{$\lambd[xyz][M]$} د
  \LR{$\lambd[x][\lambd[y][\lambd[z][M]]]$} لنډه بڼه ده۔
\end{enumerate}
\end{conv}
\olpseganchor{OLP-0344-B008}{end}

\olpseganchor{OLP-0344-B009}{start}
د بېلګې په توګه،
\[
\lambd[xy][xxyx \lambd[z][xz]]
\]
د دې لنډه بڼه ده:
\[
\lambd[x][\lambd[y][((((xx)y)x)(\lambd[z][(xz)]))]].
\]
دا دودونه په ياد کړئ۔ په لومړي سر کښې به درته ډېر ستړي کوونکي وي،
خو ورسره به عادت شئ؛ څه موده وروسته به تر بې‌شمېره قوسونو کم ستړي
کوونکي درته ښکاري۔
\olpseganchor{OLP-0344-B009}{end}

\olpseganchor{OLP-0344-B010}{start}
دوه ترمونه چې يوازې د تړلو متغيرونو په نومونو کښې توپير ولري،
\LR{$\alpha$}-معادل بلل کېږي؛ مثلاً \LR{$\lambd[x][x]$} او
\LR{$\lambd[y][y]$}۔ آسانه به وي چې دوئ ``هماغه'' ترم وګڼو؛ په بله
وينا، چې کله \LR{$M$} او \LR{$N$} يو شان بولو، د تړلو متغيرونو تر نوم بدلولو
پورې يو شانوالے هم مراد لرو۔ د متغير هغه راتګونه چې د داسې يوې \LR{$\lambd$} په ساحه
کښې وي چې هماغه متغير تړي، ``تړلي''، او نور راتګونه ``ازاد''
بلل کېږي۔ په مخکينۍ بېلګه کښې
ازاد متغير نشته؛ خو په
\[
(\lambd[z][yz])x
\]
کښې \LR{$y$} او \LR{$x$} ازاد دي او \LR{$z$} تړلے دے۔
\olpseganchor{OLP-0344-B010}{end}

\olpseganchor{OLP-0345-B004}{start}
\olfileid{lam}{int}{red}
\olsection{د لامبډا ترمونو راکمول}
\olpseganchor{OLP-0345-B004}{end}

\olpseganchor{OLP-0345-B005}{start}
له لامبډا ترمونو سره څه کولے شو؟ ساده کولے ئې شو۔ که \LR{$M$} او \LR{$N$} هر ډول
لامبډا ترمونه او \LR{$x$} هر متغير وي، نو له دې وروسته چې د~\LR{$M$} هغه تړلي
متغيرونه بيا ونوموو چې تر تعويض وروسته د~\LR{$N$} له ازادو متغيرونو سره
ټکرېږي، په~\LR{$M$} کښې د~\LR{$x$} پر ځاے د~\LR{$N$} د کېښودلو پايله په
\LR{$\Subst{M}{N}{x}$} ښيو۔ مثلاً،
\[
\Subst{(\lambd[w][xxw])}{yyz}{x} = \lambd[w][(yyz)(yyz)w].
\]
\olpseganchor{OLP-0345-B005}{end}

\olpseganchor{OLP-0345-B006}{start}
\begin{digress}
د تعويض نورې نښې \LR{$[N/x]M$}، \LR{$[x/N]M$} او \LR{$M[x/N]$} هم دي۔ پام
ورته وکړئ!
\end{digress}
\olpseganchor{OLP-0345-B006}{end}

\olpseganchor{OLP-0345-B007}{start}
په وجداني ډول، \LR{$(\lambd[x][M])N$} او \LR{$\Subst{M}{N}{x}$} يوه معنا
لري؛ لومړے ترم په دويم بدلول \emph{\LR{$\beta$}-انقباض} بلل کېږي۔
\LR{$(\lambd[x][M])N$} ته \emph{د راکمېدو وړ عبارت} (redex) او
\LR{$\Subst{M}{N}{x}$} ته ئې \emph{د انقباض پايله} (contractum) وايي۔
په عمومي ډول، که د~\LR{$P$} د کوم فرعي ترم په \LR{$\beta$}-انقباض سره~\LR{$P'$}
ته تلل شوني وي، وايو چې \LR{$P$} په \emph{يو ګام کښې تر~\LR{$P'$} پورې
\LR{$\beta$}-راکمېږي}، او \LR{$P \redone P'$} ليکو۔ که له~\LR{$P$} څخه د يو-ګامي
راکمولو د څو ګامونو په وسيله (ښايي هېڅ ګام نۀ وي)~\LR{$P'$} ترلاسه کولے
شو، نو وايو \LR{$P$} تر~\LR{$P'$} پورې \emph{\LR{$\beta$}-راکمېږي}؛ په نښو کښې
\LR{$P \red P'$}۔ هغه ترم چې نور \LR{$\beta$}-راکمول ئې ناشوني وي،
\emph{\LR{$\beta$}-ناراکمېدونکے} يا \emph{\LR{$\beta$}-عادي} بلل کېږي۔ چې
سياق روښانه وي، د ``\LR{$\beta$}-راکمېږي'' پر ځاے يوازې ``راکمېږي'' او
داسې نور وايو۔
\olpseganchor{OLP-0345-B007}{end}

\olpseganchor{OLP-0345-B008}{start}
څو بېلګې وګورو۔
\begin{enumerate}
\item لرو:
\begin{align*}
(\lambd[x][xxy]) \lambd[z][z] & \redone (\lambd[z][z])(\lambd[z][z]) y \\
& \redone (\lambd[z][z]) y \\
& \redone y.
\end{align*}
\item د يوه ترم ``ساده کول'' هغه لا پېچلے هم کولے شي:
\begin{align*}
(\lambd[x][xxy])(\lambd[x][xxy]) & \redone (\lambd[x][xxy])(\lambd[x][xxy])y \\
& \redone (\lambd[x][xxy])(\lambd[x][xxy])yy \\
& \redone \dots
\end{align*}
\item ترم بې له بدلون څخه هم پرېښودے شي:
\[
(\lambd[x][xx])(\lambd[x][xx]) \redone (\lambd[x][xx])(\lambd[x][xx]).
\]
\item ځينې ترمونه په څو لارو راکمېدے شي؛ مثلاً،
\[
(\lambd[x][(\lambd[y][yx]) z)] v \redone (\lambd[y][yv]) z
\]
د تر ټولو بهرني تطبيق په انقباض؛ او
\[
(\lambd[x][(\lambd[y][yx]) z)] v \redone (\lambd[x][zx]) v
\]
د تر ټولو دننني تطبيق په انقباض۔ خو پام وکړئ چې په دې حالت کښې
دواړه ترمونه وروسته هماغه يو ترم~\LR{$zv$} ته راکمېږي۔
\end{enumerate}
\olpseganchor{OLP-0345-B008}{end}

\olpseganchor{OLP-0345-B009}{start}
په وروستۍ بېلګه کښې دغه ګډه وروستۍ پايله تصادف نۀ دے، بلکې د
لامبډا حساب يو ژور او مهم خاصيت ښيي چې د ``چرچ--روسر خاصيت'' په
نامه يادېږي۔
\olpseganchor{OLP-0345-B009}{end}

\olpseganchor{OLP-0346-B004}{start}
\olfileid{lam}{int}{cr}
\olsection{د چرچ--روسر خاصيت}
\olpseganchor{OLP-0346-B004}{end}

\olpseganchor{OLP-0346-B005}{start}
\begin{thm}
\ollabel{thm:church-rosser}
فرض کړئ \LR{$M$}، \LR{$N_1$} او \LR{$N_2$} داسې ترمونه دي چې \LR{$M \red N_1$} او
\LR{$M \red N_2$}۔ بيا داسې ترم \LR{$P$} شته چې \LR{$N_1 \red P$} او \LR{$N_2 \red P$}۔
\end{thm}
\olpseganchor{OLP-0346-B005}{end}

\olpseganchor{OLP-0346-B006}{start}
\begin{cor}
که \LR{$M$} عادي بڼې ته راکمېدے شي، دغه عادي بڼه يوازينۍ ده۔
\end{cor}
\olpseganchor{OLP-0346-B006}{end}

\olpseganchor{OLP-0346-B007}{start}
\begin{proof}
که \LR{$M \red N_1$} او \LR{$M \red N_2$} وي، د مخکينۍ قضيې له مخې داسې
ترم~\LR{$P$} شته چې \LR{$N_1$} او \LR{$N_2$} دواړه~\LR{$P$} ته راکمېږي۔ که \LR{$N_1$} او
\LR{$N_2$} دواړه په عادي بڼه کښې وي، دا يوازې هغه وخت شوني دي چې
\LR{$N_1 \ident P \ident N_2$} وي۔
\end{proof}
\olpseganchor{OLP-0346-B007}{end}

\olpseganchor{OLP-0346-B008}{start}
په پاے کښې، دوه ترمونه \LR{$M$} او~\LR{$N$} \emph{\LR{$\beta$}-معادل}، يا يوازې
\emph{معادل}، هغه وخت بولو چې دواړه يوه ګډ ترم ته راکمېږي؛ په بله
وينا، داسې \LR{$P$} شته چې \LR{$M \red P$} او \LR{$N \red P$}۔ دا په
\LR{$M \equal[\beta] N$} ليکو۔ د \olref{thm:church-rosser} په کارولو
څېړلے شئ چې \LR{$\equal[\beta]$} د معادلتوب اړيکه ده، او دا اضافي
خاصيت هم لري چې د هر \LR{$M$} او~\LR{$N$} دپاره، که \LR{$M \red N$} يا
\LR{$N \red M$} وي، نو \LR{$M \equal[\beta] N$}۔ (په حقيقت کښې، ښودل کېدے
شي چې \LR{$\equal[\beta]$} د دې خاصيت لرونکې \emph{تر ټولو کوچنۍ}
د معادلتوب اړيکه ده۔)
\olpseganchor{OLP-0346-B008}{end}

\olpseganchor{OLP-0347-B004}{start}
\olfileid{lam}{rep}{cur}
\olsection{کَري کول}
\olpseganchor{OLP-0347-B004}{end}

\olpseganchor{OLP-0347-B005}{start}
يو \LR{$\lambd$}-تجريد \LR{$\lambd[x][M]$} د يوه آرګومېنټ تابع ښيي۔ دا هغه
وخت څرګند محدوديت دے چې څو آرګومېنټه تابع تعريفول غواړو۔ يوه لار دا
ده چې \LR{$\lambd$}-حساب داسې پراخ کړو چې جوړې، درې‌ګونې او داسې نور پکې
جوړول شوني شي؛ په دې صورت کښې، مثلاً، درې‌ځاييزه تابع
\LR{$\lambd[x][M]$} به د خپل آرګومېنټ په توګه درې‌ګونې ته انتظار کوي۔
خو په \emph{کَري کولو} دا کار ډېر آسان دے۔
\olpseganchor{OLP-0347-B005}{end}

\olpseganchor{OLP-0347-B006}{start}
يوه بېلګه وګورو۔ د يوې شېبې دپاره فرضوو چې په \LR{$\lambd$}-حساب کښې
\LR{$+$} عمل لرو۔ د جمع تابع \LR{$2$}-ځاييزه ده، يعنې دوه آرګومېنټونه اخلي۔
خو يو \LR{$\lambd$}-تجريد يوازې يوځاييزې تابعې راکوي: نحو د
\LR{$\lambd[(x,y)][(x+y)]$} په څېر عبارتونه نۀ پرېږدي۔ خو هغه يوځاييزه
تابع~\LR{$f_x(y)$} څېړلے شو چې \LR{$\lambd[y][(x+y)]$} ئې راکوي او په خپل
يو آرګومېنټ~\LR{$y$} کښې~\LR{$x$} ورزياتوي۔ په حقيقت کښې دا يوه تابع نۀ ده،
بلکې د بېلابېلو ``\LR{$x$} ورزيات کړه'' تابعو يوه کورنۍ ده، د هر عدد~\LR{$x$}
دپاره يوه۔ اوس بله يوځاييزه تابع~\LR{$g$} د \LR{$\lambd[x][f_x]$} په توګه
تعريفولے شو۔ چې پر آرګومېنټ~\LR{$x$} تطبيق شي، \LR{$g(x)$} تابع~\LR{$f_x$}
بېرته ورکوي؛ نو د هغې قيمتونه نورې تابعې دي۔ اوس که \LR{$g$} پر \LR{$x$} او
بيا پايله پر~\LR{$y$} تطبيق کړو، دا تر لاسه کوو:
\LR{$(g(x))y = f_x(y) = x+y$}۔ په دې ډول يوځاييزه تابع~\LR{$g$} د جمع د
دوه‌ځاييزې تابعې هماغه کار کولے شي۔ ``کَري کول'' د دوه‌ځاييزو تابعو
په يوځاييزو تابعو د بدلولو همدغه چل ته وايي، چې قيمتونه ئې پخپله
يوځاييزې تابعې وي۔
\olpseganchor{OLP-0347-B006}{end}

\olpseganchor{OLP-0347-B007}{start}
اوس د \LR{$\lambd$}-حساب د نحو دننه يوه سمه بېلګه وګورو۔ تابع
\LR{$f(x,y) = x$} څنګه ښيو؟ که داسې تابع تعريفول غواړو چې دوه
آرګومېنټونه اخلي او لومړے بېرته ورکوي،
\LR{$\lambd[x][\lambd[y][x]]$} ليکلے شو۔ دا په لفظي توګه هغه تابع ده
چې آرګومېنټ~\LR{$x$} اخلي او تابع~\LR{$\lambd[y][x]$} بېرته ورکوي۔ تابع
\LR{$\lambd[y][x]$} بل آرګومېنټ~\LR{$y$} اخلي، خو له پامه ئې غورځوي او تل~\LR{$x$}
بېرته ورکوي۔ وګورو چې که \LR{$\lambd[x][\lambd[y][x]]$} پر دوو
آرګومېنټونو تطبيق کړو، څۀ کېږي:
د رسم لپاره تړلي پاراميټرونه د الفا-بدلون په وسيله د ورکړو آرګومېنټونو له ازادو متغيرونو تازه ټاکو۔
\begin{align*}
  (\lambd[x][\lambd[y][x]])MN
  \bredone & (\lambd[y][M])N \\
  \bredone & M
\end{align*}
\olpseganchor{OLP-0347-B007}{end}

\olpseganchor{OLP-0347-B008}{start}
په عمومي ډول، که په يوه ترم~\LR{$N$} تعريف شوې تابع د \LR{$x_1$}، \dots،~\LR{$x_n$}
پاراميټرونه ولري، داسې ئې ليکلے شو:
\LR{$\lambd[x_1][\lambd[x_2][\ldots\lambd[x_n][N]]]$}۔ که پرې \LR{$n$}
آرګومېنټونه تطبيق کړو، نو دا تر لاسه کوو:
په عمومي رسم کښې هم تړلي پاراميټرونه تازه ټاکو۔ په دې توګه وروستے تعويض د مخکيني آرګومېنټ ازاد متغير نۀ بدلوي۔
\begin{multline*}
  (\lambd[x_1][\lambd[x_2][\ldots\lambd[x_n][N]]]) M_1 \dots M_n \bredone\\
  \begin{aligned}
  \bredone {} & (\Subst{(\lambd[x_2][\ldots\lambd[x_n][N]])}{M_1}{x_1}) M_2
  \dots M_n\\
   \eqs {} & (\lambd[x_2][\ldots\lambd[x_n][\Subst{N}{M_1}{x_1}]]) M_2
            \dots M_n \\
  \vdots & \\
  \bredone {} &\Subst{\Subst{N}{M_1}{x_1}\ldots}{M_n}{x_n}
  \end{aligned}
\end{multline*}
وروستۍ کرښه په لفظي توګه دا معنا لري چې د تابعې د تعريف په بدنه کښې
د~\LR{$x_i$} پر ځاے~\LR{$M_i$} کېږدو؛ دا کټ مټ هغه څه دي چې پر تابعې د څو
آرګومېنټونو د تطبيق پر مهال ئې غواړو۔
\olpseganchor{OLP-0347-B008}{end}

\olpseganchor{OLP-0348-B004}{start}
\olfileid{lam}{int}{rep}
\olsection{د لامبډا په وسيله د تعريف وړ حسابي تابعې}
\olpseganchor{OLP-0348-B004}{end}

\olpseganchor{OLP-0348-B005}{start}
لامبډا حساب څنګه د محاسبې نمونه کېدے شي؟ په لومړي سر کښې د دې
پوښتنې معنا هم روښانه نۀ ده۔ د طبيعي عددونو پر سر د محاسبه کېدنې د
څېړنې دپاره بايد د هغو عددونو مناسب تمثيل ومومو۔ لاندې يو داسې
تمثيل دے چې په حيرانوونکې توګه ښه کار کوي۔
\olpseganchor{OLP-0348-B005}{end}

\olpseganchor{OLP-0348-B006}{start}
\begin{defn}
د هر طبيعي عدد~\LR{$n$} دپاره \emph{د چرچ عددنښه} \LR{$\num{n}$} هغه لامبډا
ترم \LR{$\lambd[x][\lambd[y][(x(x(x(\dots x(y)))))]]$} تعريف کړئ چې
په بدنه کښې \LR{$n$} ځله \LR{$x$} راغلے وي۔
\end{defn}
\olpseganchor{OLP-0348-B006}{end}

\olpseganchor{OLP-0348-B007}{start}
ترمونه \LR{$\num{n}$} ``تکراروونکي'' دي: پر ننوت~\LR{$f$}، \LR{$\num{n}$} هغه
تابع ورکوي چې \LR{$y$}،~\LR{$f^n(y)$} ته وړي۔ پام وکړئ چې هره عددنښه په عادي
بڼه کښې ده۔ اوس ويلے شو چې يو لامبډا ترم پر طبيعي عددونو يوه تابع
``محاسبه کوي'' يعنې څۀ۔
\olpseganchor{OLP-0348-B007}{end}

\olpseganchor{OLP-0348-B008}{start}
\begin{defn}
فرض کړئ \LR{$f(x_0, \dots, x_{k-1})$} له~\LR{$\Nat$} څخه~\LR{$\Nat$} ته يوه
\LR{$k$}-ځاييزه جزوي تابع ده۔ وايو چې يو \LR{$\lambd$}-ترم~\LR{$F$} هغه وخت، او
يوازې هغه وخت،~\LR{$f$} \emph{لامبډا-تعريفوي} چې د طبيعي عددونو د
هرې لړۍ \LR{$n_0$}، \dots،~\LR{$n_{k-1}$} دپاره،
\[
F\, \num{n_0}\, \num{n_1} \dots \num{n_{k-1}} \red \num{f(n_0, n_1, \dots,
  n_{k-1})}
\]
وي، که \LR{$f(n_0, \dots, n_{k-1})$} تعريف شوې وي؛ او که نۀ وي، نو
\LR{$F\, \num{n_0}\, \num{n_1} \dots \num{n_{k-1}}$} هېڅ عادي بڼه
نۀ لري۔
\end{defn}
\olpseganchor{OLP-0348-B008}{end}

\olpseganchor{OLP-0348-B009}{start}
\begin{thm}
\ollabel{thm:lambda-def}
يوه تابع~\LR{$f$} هغه وخت، او يوازې هغه وخت، جزوي محاسبه کېدونکې تابع
ده چې يو لامبډا ترم ئې لامبډا-تعريف کړي۔
\end{thm}
\olpseganchor{OLP-0348-B009}{end}

\olpseganchor{OLP-0348-B010}{start}
\begin{explain}
دا قضيه يو څه حيرانوونکې ده۔ د محاسبې د نمونې په توګه لامبډا حساب
يو ډېر ساده حساب دے؛ يوازې لامبډا تجريد او تطبيق ئې عملونه دي!{} خو
له دغو لږو وسيلو څخه هره محاسبوي کړنلاره پلې کېدے شي۔
\end{explain}
\olpseganchor{OLP-0348-B010}{end}

\olpseganchor{OLP-0349-B004}{start}
\olfileid{lam}{int}{cmp}
\olsection{د لامبډا په وسيله د تعريف وړ تابعې محاسبه کېدونکې دي}
\olpseganchor{OLP-0349-B004}{end}

\olpseganchor{OLP-0349-B005}{start}
\begin{thm}
\ollabel{thm:lambda-computable}
که يوه جزوي تابع~\LR{$f$} يو لامبډا ترم لامبډا-تعريف کړې وي، نو
هغه محاسبه کېدونکې ده۔
\end{thm}
\olpseganchor{OLP-0349-B005}{end}

\olpseganchor{OLP-0349-B006}{start}
\begin{proof}
فرض کړئ تابع~\LR{$f$} يو لامبډا ترم~\LR{$X$} لامبډا-تعريف کړې ده۔ د~\LR{$f$}
د محاسبې يوه غير صوري کړنلاره بيان کړو۔ پر ننوت \LR{$m_0$}،
\dots،~\LR{$m_{n-1}$}، ترم \LR{$X \num{m_0} \ldots \num{m_{n-1}}$} وليکئ۔ يوه ونه جوړه کړئ: لومړے د اصلي ترم ټول يو-ګامي
راکمول وليکئ؛ تر هغو لاندې د هر يوه ټول يو-ګامي راکمول (يعنې د اصلي
ترم دوه-ګامي راکمول) وليکئ؛ او همداسې دوام ورکړئ۔ که هر کله يوې
عددنښې ته ورسېدئ، له هغې سره سم طبيعي عدد د ځواب په توګه ورکړئ۔
لټون د ونې په پرله‌پسې کچو کوو، څو کومه محدودې ګامونو لرونکې
لار پاتې نۀ شي۔ که هېڅ عددنښې ته نۀ رسېږو، لټون همداسې دوام کوي
او په دغه ننوت تابع ناتعريفه ده؛ ناتعريفه‌والے د لټون د پاے
يو جلا ځواب نۀ دے۔
\olpseganchor{OLP-0349-B006}{end}

\olpseganchor{OLP-0349-B007}{start}
د چرچ اصل ته په رجوع سره ويلے شو چې دا تابع محاسبه کېدونکې ده۔ د
قضيې د ثبوت ښه لاره به دا وي چې دغه لټون‌کړنلاره په بازګشتي ډول
بيان کړو۔ مثلاً، د بنسټيزو بازګشتي تابعو او اړيکو يوه لړۍ تعريفولے
شو، لکه ``\LR{$\fn{IsASubterm}$}''، ``\LR{$\fn{Substitute}$}''،
``\LR{$\fn{ReducesToInOneStep}$}''، ``\LR{$\fn{ReductionSequence}$}''،
``\LR{$\fn{Numeral}$}'' او داسې نور۔ بيا د~\LR{$f(m_0, \dots, m_{n-1})$}
د محاسبې جزوي بازګشتي کړنلاره دا ده چې د يو-ګامي راکمول داسې لړۍ
ولټوو چې له \LR{$X \num{m_0} \dots \num{m_{n-1}}$} څخه پيل او پر يوه
عددنښه پاے ته رسېږي، او له هغې عددنښې سره سم عدد بېرته ورکړو۔
جزئيات اوږده او ستړي کوونکي دي، خو له دې پرته عادي دي۔
\end{proof}
\olpseganchor{OLP-0349-B007}{end}

\olpseganchor{OLP-0350-B004}{start}
\olfileid{lam}{int}{lrp}
\olsection{محاسبه کېدونکې تابعې د لامبډا په وسيله د تعريف وړ دي}
\olpseganchor{OLP-0350-B004}{end}

\olpseganchor{OLP-0350-B005}{start}
\begin{thm}
\ollabel{thm:computable-lambda}
هره محاسبه کېدونکې جزوي تابع د لامبډا په وسيله د تعريف وړ ده۔
\end{thm}
\olpseganchor{OLP-0350-B005}{end}

\olpseganchor{OLP-0350-B006}{start}
\begin{proof}
غواړو وښيو چې جزوي محاسبه کېدونکې تابع يو لامبډا ترم لامبډا-تعريف کوي۔ د \olref[cmp][rec][nft]{thm:kleene-nf} له مخې کافي ده چې د يوې ټولې بنسټيزې بازګشتي ازموينې لومړنے بريالے شاهد ولټوو او د شاهد له کوډ څخه ځواب واخلو۔ د څو ورودي عددونو په حالت کښې لومړے د هغوی کوډ جوړوو؛ د لړۍ کوډول بنسټيز بازګښت دے۔ د لومړنيو تابعو، د ټوليزو تابعو د ترکيب او د ټوليزو تابعو د بنسټيز بازګښت لاندې جوړونې ښيي چې لومړنۍ تابعې د لامبډا په وسيله د تعريف وړ دي او په دغو دوو عملونو هره ټوله د لامبډا په وسيله د تعريف وړ تابع بيا د لامبډا په وسيله د تعريف وړ پاتې کېږي۔ د ازموينې ځانګړونکې تابعه داسې اخلو چې د بريا پر وخت صفر او کنه يو ورکوي۔ د شاهد څخه د ځواب اخستونکے ترم~\LR{$U$} دې وي، او ازموينه دې~\LR{$P$} لامبډا-تعريف کړي۔ د ورودي عددونو لړۍ په~\LR{$\vec x$} ښيو۔
د چرچ د عددنښو لپاره دا ترمونه وکاروئ:
\[
\mathsf{T}=\lambda a.\lambda b.a,\quad\mathsf{F}=\lambda a.\lambda b.b,\quad Z=\lambda n.n(\lambda u.\mathsf{F})\mathsf{T}.
\]
پر صفر~\LR{$Z$} لومړے انتخاب او پر هر مثبت عدد دويم انتخاب راګرځوي۔ \LR{$S$} د تالي ترم دے او~\LR{$Y$} د ثابت ټکي د تورينګ ترم دے۔ د ورکړو ورودي عددونو لپاره دا لټون جوړوو:
\[
L_{\vec x}=Y(\lambda h.\lambda s.(Z(P(\vec x,s)))(U(s))(h(Ss))).
\]
د تعريفوونکي ترم بدنه~\LR{$L_{\vec x}\num{0}$} ده؛ د ورودي لړۍ متغيرونه په بهرني تجريد کښې تړو۔ د هرې شېبې په پيل کښې د ثابت ټکي د لګېدلو تر ټولو کيڼ رېډکس راکموو۔ بيا ازموينه د صفر يا يو عددنښې ته رسوو او اړوند انتخاب اخلو۔ د بريا په صورت کښې~\LR{$U$} د شاهد ځواب ورکوي۔ د ناکامۍ په صورت کښې د تالي عددنښه اخلو او د هماغه لټون بل پړاو ته ځو۔ دا د لومړني بريالي شاهد دپاره د غوښتل شوې پايلې عددنښه ورکوي، نو په دې حالت کښې تابع همدا ترم لامبډا-تعريف کوي۔
که شاهد هېڅکله پيدا نۀ شي، همدا راکمونه نامتناهي ده۔ د هغې په هر پړاو کښې لږ تر لږه يو تر ټولو کيڼ رېډکس راکمېږي۔ د \href{https://pages.jh.edu/rrynasi1/FoundationsOFMath/Literature/LambdaCalculus/Barendregt1984TheLambdaCalculus.ItsSyntax+Semantics.RevisedEdition.pdf}{بارېنډرېخټ، ۱۹۸۴، قضيه ۸٫۴٫۹ او پايله ۸٫۴٫۱۰} له مخې، داسې نامتناهي لړۍ لرونکے ترم عادي بڼه نۀ لري۔ نو په ناتعريفه ننوت کښې هم د لامبډا-تعريف شرط پوره کېږي۔
\end{proof}
\olpseganchor{OLP-0350-B006}{end}

\olpseganchor{OLP-0350-B007}{start}
د ثبوت پاتې برخه د لا آسانه لوستلو دپاره به د نښو يو لا دوديز
طرز وکاروو۔ مثلاً، د \LR{$Mxyz$} پر ځاے \LR{$M(x, y, z)$} ليکو۔ که څۀ هم دا
نښه څو آرګومېنټونه راښيي، په ياد ولرئ چې د بې‌ټايپه لامبډا حساب
ترمونه ځانته ټاکلې ځاييزه کچه نۀ لري۔ نو د هماغه ترم~\LR{$M$} دپاره د
\LR{$M(x,y)$} او \LR{$M(x,y,z,w)$} ليکل هم معنا لري۔ خو دا نښه ښيي چې موږ
\LR{$M$} د درې متغيرونو د تابعې په توګه ګورو، او د تعريفونو تر شا موخه
روښانوي۔ همداسې به د ``\LR{$M$} په \LR{$M =
\lambd[x][\lambd[y][\lambd[z][\dots]]]$} تعريف کړئ'' پر ځاے
وايو ``\LR{$M$} په \LR{$M(x,y,z) = \dots$} تعريف کړئ''۔
\olpseganchor{OLP-0350-B007}{end}

\olpseganchor{OLP-0351-B004}{start}
\olfileid{lam}{int}{bas}
\olsection{لومړنۍ بنسټيزې بازګشتي تابعې د لامبډا په وسيله د تعريف وړ دي}
\olpseganchor{OLP-0351-B004}{end}

\olpseganchor{OLP-0351-B005}{start}
\begin{lem}
تابعې \LR{$\Zero$}، \LR{$\Succ$} او \LR{$\Proj{n}{i}$} د لامبډا په وسيله د تعريف وړ دي۔
\end{lem}
\olpseganchor{OLP-0351-B005}{end}

\olpseganchor{OLP-0351-B006}{start}
\begin{proof}
\LR{$\Zero$} يوازې \LR{$\lambd[u][\lambd[x][\lambd[y][y]]]$} دے۔
بهرنی تړون د تابعې ورودي آرګومېنټ اخلي۔ بدنه تل د صفر عددنښه ورکوي۔
\olpseganchor{OLP-0351-B006}{end}

\olpseganchor{OLP-0351-B007}{start}
د تالي تابع~\LR{$\Succ$} تعريف دا دے:
\LR{$\fn{Succ}(u) = \lambd[x][\lambd[y][x(uxy)]]$}۔ ځير شئ چې دا
ولې کار کوي: هره عددنښه~\LR{$\num{n}$} د تکراروونکي په توګه وګڼئ۔ د
هرې تابعې~\LR{$f$} دپاره \LR{$\fn{Succ}(\num{n},f)$} داسې تابع ده چې پر
ننوت~\LR{$y$}، لومړے له~\LR{$y$} څخه په پيل سره~\LR{$f$}،~\LR{$n$} ځله تطبيقوي او
بيا ئې يو ځل نور هم تطبيقوي۔
\olpseganchor{OLP-0351-B007}{end}

\olpseganchor{OLP-0351-B008}{start}
د پروجيکشنونو په هکله نور څۀ نشته: \LR{$\fn{Proj}^n_i(x_0, \dots,
x_{n-1}) = x_i$}۔ په بله وينا، زموږ د نښيزو دودونو له مخې
\LR{$\fn{Proj}^n_i$} لامبډا ترم
\LR{$\lambd[x_0][\dots \lambd[x_{n-1}][x_i]]$} دے۔
\end{proof}
\olpseganchor{OLP-0351-B008}{end}

\olpseganchor{OLP-0352-B004}{start}
\olfileid{lam}{int}{com}
\olsection{د لامبډا په وسيله د تعريف وړ تابعې د ترکيب تر عمل لاندې تړلې دي}
\olpseganchor{OLP-0352-B004}{end}

\olpseganchor{OLP-0352-B005}{start}
\begin{lem}
د لامبډا په وسيله د تعريف وړ تابعې د ترکيب تر عمل لاندې تړلې دي۔
\end{lem}
\olpseganchor{OLP-0352-B005}{end}

\olpseganchor{OLP-0352-B006}{start}
\begin{proof}
فرض کړئ \LR{$f$} د \LR{$h$}، \LR{$g_0,$}
\dots،~\LR{$g_{k-1}$} له ترکيب څخه تعريف شوې ده۔ که \LR{$h$}، \LR{$g_0$}،
\dots،~\LR{$g_{k-1}$} په ترتيب سره د \LR{$H$}، \LR{$G_0$}،
\dots،~\LR{$G_{k-1}$} په وسيله لامبډا-تعريف وي، بايد داسې
ترم~\LR{$F$} ومومو چې~\LR{$f$} لامبډا-تعريفوي۔ خو~\LR{$F$} په ساده ډول داسې
تعريفولے شو:
\[
F(x_0, \dots, x_{l-1}) = H(G_0(x_0, \dots, x_{l-1}),
\dots, G_{k-1}(x_0, \dots, x_{l-1})).
\]
په بله وينا، د لامبډا حساب ژبه د ترکيب د تمثيل دپاره مناسبه ده۔
دا محاسبه هغه وخت کار کوي چې ټول دنني ارزښتونه تعريف شوي وي۔ د بهرنۍ تابعې ترم کولے شي يو آرګومېنټ له پامه وغورځوي، نو دغه مستقيم ترکيب په ناتعريفه دنني ارزښت کښې هم کله ناکله عددنښه ورکوي۔ د جزوي تابعو عمومي تړلتيا د \olref[cmp]{thm:lambda-computable} او \olref[lrp]{thm:computable-lambda} له مخې راځي۔ د ترکيب د محاسبې کړنلاره لومړے ټول دنني ارزښتونه محاسبه کوي، بيا بهرنۍ تابع پرې چلوي۔ که کوم لازم ارزښت ناتعريف وي، همدا کړنلاره نۀ درېږي۔ نو ترکيب جزوي محاسبه کېدونکے او لامبډا-تعريفېدونکے دے۔
\end{proof}
\olpseganchor{OLP-0352-B006}{end}

\olpseganchor{OLP-0353-B004}{start}
\olfileid{lam}{int}{pr}
\olsection{د لامبډا په وسيله د تعريف وړ تابعې د بنسټيز بازګښت تر عمل لاندې تړلې دي}
\olpseganchor{OLP-0353-B004}{end}

\olpseganchor{OLP-0353-B005}{start}
د بنسټيز بازګښت په برخه کښې بالاخره يو څۀ کار ته اړتيا لرو۔ دا کار
به په پړاوونو کښې کوو۔ د پخوا په شان، فرض کوو چې ترمونه \LR{$G$} او \LR{$H$}
په ترتيب سره تابعې \LR{$g$} او~\LR{$h$} لامبډا-تعريف کوي؛ غواړو داسې
ترم \LR{$F$} ومومو چې لاندې تعريف شوې تابع~\LR{$f$} لامبډا-تعريفوي:
\begin{align*}
f(0, \vec z) & = g(\vec z) \\
f(x+1, \vec z) & = h(x, f(x,\vec z), \vec z).
\end{align*}
نو په عمومي توګه، چې لامبډا ترمونه \LR{$G'$} او~\LR{$H'$} راکړل شوي وي، دومره
بس ده چې داسې ترم~\LR{$F$} ومومو چې
\begin{align*}
F(\num{0}, \vec z) & \equiv G'(\vec z) \\
F(\overline{n+1}, \vec z) & \equiv H'(\num{n}, F(\num{n}, \vec z), \vec z)
\end{align*}
د هر طبيعي عدد~\LR{$n$} دپاره پوره کړي؛ دا چې \LR{$G'$} او~\LR{$H'$} په ترتيب سره
\LR{$g$} او~\LR{$h$} لامبډا-تعريف کوي، معنا ئې دا ده چې کله د \LR{$\vec z$}
په ځاے عددنښې \LR{$\num{\vec m}$} کښېږدو، نو تر هغه چې ټول لازم بازګښتي ارزښتونه تعريف شوي وي، \LR{$F(\num{n+1},\num{\vec m})$} سم ځواب ته عادي کېږي۔
\olpseganchor{OLP-0353-B005}{end}

\olpseganchor{OLP-0353-B006}{start}
خو د دې دپاره دومره بس ده چې داسې ترم~\LR{$F$} ومومو چې
\begin{align*}
F(\num 0) & \equiv G \\
F(\overline {n+1}) & \equiv H(\num{n},F(\num{n}))
\intertext{\RL{د هر طبيعي عدد \LR{$n$} دپاره پوره کړي، چېرته چې}}
G & = \lambd[\vec z][G'(\vec z)] \text{\RL{او}}\\
H(u,v) & = \lambd[\vec z][H'(u,v(\vec z),\vec z)].
\end{align*}
\LR{$\vec z$} سره د جلا چلند اړتيا نشته---هغوي په لامبډا نښه کښې جذبېږي۔
\olpseganchor{OLP-0353-B006}{end}

\olpseganchor{OLP-0353-B007}{start}
د ترم \LR{$F$} تر تعريفولو وړاندې د مرتبو جوړو د سمبالولو يوه وسيله پکار
ده۔ لاندې لمه دا وسيله راکوي۔
\olpseganchor{OLP-0353-B007}{end}

\olpseganchor{OLP-0353-B008}{start}
\begin{lem}
داسې لامبډا ترم \LR{$D$} شته چې د هر دوو لامبډا ترمونو \LR{$M$} او~\LR{$N$}
دپاره \LR{$D(M,N)(\num{0}) \red M$} او \LR{$D(M,N)(\num{1}) \red N$} وي۔
\end{lem}
\olpseganchor{OLP-0353-B008}{end}

\olpseganchor{OLP-0353-B009}{start}
\begin{proof}
لومړے لامبډا ترم \LR{$K$} داسې تعريف کړئ:
\[
K(y) = \lambd[x][y].
\]
په بله وينا، \LR{$K$} ترم \LR{$\lambd[y][\lambd[x][y]]$} دے۔ په بل نظر، د
هر~\LR{$M$} دپاره \LR{$K(M)$} ثابته تابع ده چې پر هر ننوت~\LR{$M$} ورکوي۔
\olpseganchor{OLP-0353-B009}{end}

\olpseganchor{OLP-0353-B010}{start}
اوس \LR{$D(x,y,z)$} په \LR{$D(x,y,z) = z (K(y))x$} تعريف کړئ۔ بيا لرو
\begin{align*}
D(M,N,\num 0) & \red \num 0 (K(N)) M \red M \text{\RL{او}}\\
D(M,N,\num 1) & \red \num 1 (K(N)) M \red K(N) M \red N,
\end{align*}
لکه څنګه چې غوښتل شوي وو۔
\end{proof}
\olpseganchor{OLP-0353-B010}{end}

\olpseganchor{OLP-0353-B011}{start}
مطلب دا دے چې \LR{$D(M,N)$} د مرتبې جوړې \LR{$\tuple{M, N}$} استازيتوب کوي؛
که فرض کړو چې \LR{$P$} هم د داسې جوړې استازے وي، نو \LR{$P(\num 0)$} او
\LR{$P(\num 1)$} په ترتيب سره د چپ او ښي پروجيکشنونو، \LR{$(P)_0$} او
\LR{$(P)_1$}، استازيتوب کوي۔ له دې وروسته همدا وروستۍ نښې کاروو۔
\olpseganchor{OLP-0353-B011}{end}

\olpseganchor{OLP-0353-B012}{start}
\begin{lem}
د لامبډا په وسيله د تعريف وړ تابعې د بنسټيز بازګښت تر عمل لاندې تړلې دي۔
\end{lem}
\olpseganchor{OLP-0353-B012}{end}

\olpseganchor{OLP-0353-B013}{start}
\begin{proof}
بايد وښيو چې د هر دوو راکړل شويو ترمونو \LR{$G$} او \LR{$H$} دپاره داسې ترم
\LR{$F$} موندل کېدے شي چې
\begin{align*}
F(\num{0}) & \equiv G \\
F(\num{n+1}) & \equiv H(\num{n}, F(\num{n}))
\end{align*}
د هر طبيعي عدد~\LR{$n$} دپاره پوره کړي۔ مفکوره په لنډ ډول دا ده چې
د \emph{جوړو} لاندې لړۍ محاسبه کړو:
\[
\tuple{\num 0, F(\num 0)}, \tuple{\num 1, F(\num 1)}, \dots,
\]
او عددنښې د تکراروونکو په توګه وکاروو۔ پام وکړئ چې لومړۍ جوړه
يوازې \LR{$\tuple{\num 0, G}$} ده۔ له جوړې \LR{$\tuple{\num{n},
F(\num{n})}$} څخه راتلونکې جوړه \LR{$\tuple{\num{n+1}, F(\num{n+1})}$}
بايد له \LR{$\tuple{\num{n+1}, H(\num{n}, F(\num{n}))}$} سره برابره وي۔
داسې لامبډا ترم~\LR{$T$} جوړوو چې دا يو-ګامي بدلون ترسره کړي۔
\olpseganchor{OLP-0353-B013}{end}

\olpseganchor{OLP-0353-B014}{start}
تفصيل دا دے۔ \LR{$T(u)$} داسې تعريف کړئ:
\[
T(u) = \tuple{S((u)_0), H((u)_0,(u)_1)}.
\]
دلته د S په نښه د جانشين تابع هغه لامبډا
ترم کارېږي چې په مخکني بنسټيزو تابعو برخه کښې تعريف شوے دے۔
اوس په آسانۍ سره کتلے شو چې د هر عدد~\LR{$n$} دپاره
\[
T(\tuple{\num{n}, M}) \red \tuple{\num{n+1}, H(\num{n}, M)}.
\]
لکه پاس چې وويل شول، د \LR{$G$} او \LR{$H$} په لرلو سره \LR{$F(u)$} داسې
تعريف کړئ:
\[
F(u) = (u(T,\tuple{\num 0, G}))_1.
\]
په بله وينا، د ننوت~\LR{$\num{n}$} دپاره \LR{$F$}، \LR{$T$} پر
\LR{$\tuple{\num 0, G}$} باندې \LR{$n$} ځله تکراروي او بيا دويم غړے
ورکوي۔ په پيل کښې لرو
\begin{enumerate}
\item \LR{$\num{0} (T, \tuple{\num 0, G}) \equiv \tuple{\num{0}, G}$}
\item \LR{$F(\num{0}) \equiv G$}
\end{enumerate}
پر~\LR{$n$} د استقرا په وسيله د هر طبيعي عدد دپاره لاندې ښودلے شو:
\begin{enumerate}
\item \LR{$\num{n+1}(T, \tuple{\num{0}, G}) \equiv \tuple{\num{n+1}, F(\num{n+1})}$}
\item \LR{$F(\num{n+1}) \equiv H(\num{n}, F(\num{n}))$}
\end{enumerate}
د دويمې فقرې دپاره لرو
\begin{align*}
F(\num{n+1}) & \red (\num{n+1}(T, \tuple{\num 0, G}))_1 \\
& \equiv (T(\num{n} (T, \tuple{\num 0, G})))_1 \\
& \equiv (T(\tuple{\num{n}, F(\num{n})}))_1 \\
& \equiv (\tuple{\num{n+1}, H(\num{n}, F(\num{n}))})_1 \\
& \equiv H(\num{n}, F(\num{n})).
\end{align*}
په درېيمه کرښه کښې مو د جوړې دپاره د استقرا فرضيه کارولې ده۔
د لومړۍ فقرې دپاره لرو
\begin{align*}
\num{n+1} (T, \tuple{\num 0, G}) &
\equiv T(\num{n} (T, \tuple{\num 0, G})) \\
& \equiv T( \tuple{\num{n}, F(\num{n})}) \\
& \equiv \tuple{\num{n+1}, H(\num{n}, F(\num{n}))} \\
& \equiv \tuple{\num{n+1}, F(\num{n+1})}.
\end{align*}
په وروستۍ کرښه کښې مو دويمه فقره کارولې ده۔ نو ومو ښودله چې
\LR{$F(\num 0) \equiv G$} او د هر \LR{$n$} دپاره \LR{$F(\num {n+1}) \equiv H(\num
n, F(\num{n}))$}، او همدا زموږ غوښتنه وه۔
د جوړو تکرار هغه ارزښت سم ورکوي چې د بازګښت د هر لازم ګام ارزښت پکښې تعريف شوے وي۔ د ناتعريفه ګام ترم کله ناکله له پامه غورځېږي، ځکه د لامبډا تطبيق خپل آرګومېنټ حتمي نۀ راکموي۔ د جزوي تابعو د تړلتيا لپاره \olref[cmp]{thm:lambda-computable} او \olref[lrp]{thm:computable-lambda} وکاروئ۔ د بازګښت محاسبه د پيليز ارزښت په موندلو پيل کوي او بيا د ګام تابع په ترتيب چلوو۔ که يو لازم ګام نۀ درېږي، ټوله کړنلاره هم نۀ درېږي۔ نو د جزوي محاسبه کېدنې او لامبډا-تعريف له برابرۍ څخه عمومي تړلتيا راځي۔
\end{proof}
\olpseganchor{OLP-0353-B014}{end}

\olpseganchor{OLP-0354-B004}{start}
\olfileid{lam}{int}{fix}
\olsection{د ثابت ټکي ترکيب کوونکي}
\olpseganchor{OLP-0354-B004}{end}

\olpseganchor{OLP-0354-B005}{start}
فرض کړئ يو لامبډا ترم \LR{$g$} لرئ او داسې بل ترم \LR{$k$} غواړئ چې \LR{$k$} له
\LR{$gk$} سره \LR{$\beta$}-هم‌ارزه وي۔ زموږ د نښيز دود له مخې ترمونه
\[
\fn{diag}(x) = xx
\]
او
\[
l(x) = g(\fn{diag}(x))
\]
تعريف کړئ؛ په بله وينا، \LR{$l$} ترم \LR{$\lambd[x][g(xx)]$} دے۔ فرض کړئ
\LR{$k$} ترم \LR{$ll$} وي۔ نو لرو
\begin{align*}
k & = (\lambd[x][g(xx)])(\lambd[x][g(xx)]) \\
& \red  g((\lambd[x][g(xx)])(\lambd[x][g(xx)])) \\
& = gk.
\end{align*}
که
\[
Y = \lambd[g][((\lambd[x][g(xx)])(\lambd[x][g(xx)]))]
\]
واخلو، نو \LR{$Yg$} او \LR{$g(Yg)$} يو ګډ ترم ته تقليلېږي؛ ځکه نو \LR{$Yg
\equiv_\beta g(Yg)$}۔ دا ``د کري ترکيب کوونکے'' بلل کېږي۔ که د دې
پر ځاے
\[
Y = (\lambd[xg][g(xxg)])(\lambd[xg][g(xxg)])
\]
واخلو، نو په رښتيا \LR{$Yg$}، \LR{$g(Yg)$} ته تقليلېږي، چې تر دې پياوړې
خبره ده۔ د \LR{$Y$} دا وروستۍ بڼه ``د تورينګ ترکيب کوونکے'' بلل کېږي۔
\olpseganchor{OLP-0354-B005}{end}

\olpseganchor{OLP-0355-B004}{start}
\olfileid{lam}{int}{min}
\olsection{د لامبډا په وسيله د تعريف وړ تابعې د کمينه‌موندنې تر عمل لاندې تړلې دي}
\olpseganchor{OLP-0355-B004}{end}

\olpseganchor{OLP-0355-B005}{start}
\begin{lem}
فرض کړئ \LR{$f(x,y)$} د لامبډا په وسيله د تعريف وړ ده۔ \LR{$g$} داسې تعريف کړئ:
\[
g(x) \simeq \umin{y}{f(x,y)}.
\]
نو \LR{$g$} هم د لامبډا په وسيله د تعريف وړ ده۔
\end{lem}
\olpseganchor{OLP-0355-B005}{end}

\olpseganchor{OLP-0355-B006}{start}
\begin{proof}
مفکوره په لنډ ډول داسې ده۔ د \LR{$x$} په لرلو سره به د ثابت ټکي لامبډا
ترم \LR{$Y$} په وسيله داسې تابع \LR{$h_x(n)$} تعريف کړو چې له~\LR{$n$} څخه د~\LR{$y$}
لټون پيلوي؛ بيا \LR{$g(x)$} يوازې \LR{$h_x(0)$} دے۔ تابع~\LR{$h_x$} د ثابت ټکي
د يوې معادلې د حل په توګه ښودلے شو:
\[
h_x(n) \simeq
\begin{cases}
n & \text{\RL{که \LR{$f(x,n) = 0$} وي}} \\
h_x(n+1) & \text{\RL{کنه.}}
\end{cases}
\]
\olpseganchor{OLP-0355-B006}{end}

\olpseganchor{OLP-0355-B007}{start}
تفصيل ئې دا دے۔ څرنګه چې \LR{$f$} لامبډا-تعريفېدونکې ده، يو ترم
\LR{$F$} ئې لامبډا-تعريف کوي۔
په ياد ولرئ چې يو لامبډا ترم~\LR{$D$} هم لرو، داسې چې
\LR{$D(M, N, \bar{0}) \red M$} او \LR{$D(M, N, \bar{1})
\red N$}۔ د اوس دپاره \LR{$x$} ثابت ونيسئ۔ د \LR{$h_x$} د لامبډا-تعريف
کولو دپاره داسې ترم~\LR{$H$} غواړو چې پر~\LR{$x$} تړلے وي او دا معادله پوره کړي:
\[
H(\num{n}) \equiv D(\num{n}, H(S(\num{n})), F(x, \num{n})).
\]
دا کار د ثابت ټکي ترم~\LR{$Y$} په وسيله کولے شو۔ لومړے \LR{$U$} داسې ترم
ونيسئ:
\[
\lambd[h][\lambd[z][D(z,(h(Sz)),F(x,z))]],
\]
او بيا \LR{$H$} ترم~\LR{$YU$} ونيسئ۔ پام وکړئ چې په~\LR{$H$} کښې يوازنے آزاد
متغير~\LR{$x$} دے۔ راځئ وښيو چې \LR{$H$} پاسنۍ معادله پوره کوي۔
\olpseganchor{OLP-0355-B007}{end}

\olpseganchor{OLP-0355-B008}{start}
د \LR{$Y$} له تعريفه لرو
\[
H = YU \equiv U(YU) = U(H).
\]
په ځانګړې توګه، د هر طبيعي عدد~\LR{$n$} دپاره لرو
\begin{align*}
H(\num{n}) & \equiv U(H, \num{n}) \\
& \red D(\num{n}, H(S(\num{n})), F(x, \num{n})),
\end{align*}
لکه څنګه چې غوښتل شوي وو۔ پام وکړئ چې که په وروستۍ کرښه کښې د~\LR{$x$}
په ځاے عددنښه \LR{$\num{m}$} کښېږدئ، نو که \LR{$F(\num{m}, \num{n})$}
\LR{$\num{0}$} ته تقليل شي، دا عبارت \LR{$\num{n}$} ته تقليلېږي؛ او که
\LR{$F(\num{m}, \num{n})$} بلې هرې عددنښې ته تقليل شي، نو عبارت
\LR{$H(S(\num{n}))$} ته تقليلېږي۔
\olpseganchor{OLP-0355-B008}{end}

\olpseganchor{OLP-0355-B009}{start}
په تعريف شوي حالت کښې \LR{$G$} ترم \LR{$\lambd[x][H(\num 0)]$} ونيسئ۔ د هر~\LR{$m$} دپاره که \LR{$g(m)$} تعريف شوې وي، له صفره تر شاهد پورې ټولې ازموينې تعريف شوې دي۔ نو \LR{$G(\num m)$}، \LR{$\overline {g(m)}$} ته تقليلېږي۔ په دې حالت کښې \LR{$G$}،~\LR{$g$} لامبډا-تعريفوي، يعنې غوښتل شوې عددنښه ورکوي۔
د ناتعريفه ننوت لپاره د دې ټاکونکي مستقيمه بڼه په عمومي توګه بس نۀ ده۔ يو بې‌عادي‌بڼې ترم کولے شي په تطبيق کښې هغه برخه وغورځوي چې عادي بڼه نۀ لري۔ د عمومي لمې لپاره د \olref[cmp]{thm:lambda-computable} له مخې د ازموينې جزوي محاسبه واخلو۔ پر صفر يې پيلوو، او تر هغه پر بل عدد ازموينه کوو چې اوسنی ارزښت تعريف شوے او غيرصفر وي۔ د لومړي صفر پر موندلو هماغه شاخص ورکوو۔ که کوم اړين ارزښت ناتعريف وي، يا تل غيرصفر ارزښتونه راشي، لټون نۀ درېږي۔ دا د جزوي کمينه‌موندنې تعريف دے۔ د \olref[lrp]{thm:computable-lambda} له مخې د دغې جزوي محاسبې لپاره يو لامبډا ترم شته چې پر ناتعريفه ننوت عادي بڼه نۀ لري۔
\end{proof}
\olpseganchor{OLP-0355-B009}{end}

\olpseganchor{OLP-0356-B004}{start}
\olchapter{lam}{syn}{نحو}
\olpseganchor{OLP-0356-B004}{end}

\olpseganchor{OLP-0357-B004}{start}
\olfileid{lam}{syn}{trm}
\olsection{ترمونه}
\olpseganchor{OLP-0357-B004}{end}

\olpseganchor{OLP-0357-B005}{start}
د لامبډا حساب ترمونه د متغيرونو له نامتناهي زېرمي \LR{$\Obj{v_0}$}،
\LR{$\Obj{v_1}$}، \dots، له~``\LR{$\lambd$}'' نښې او قوسونو څخه په استقرايي
ډول جوړېږي۔ د متغيرونو د ښودلو دپاره \LR{$x$}، \LR{$y$}، \LR{$z$}، \dots{} او
د ترمونو د ښودلو دپاره \LR{$M$}، \LR{$N$}، \LR{$P$}، \dots{} کاروو۔
\olpseganchor{OLP-0357-B005}{end}

\olpseganchor{OLP-0357-B006}{start}
\begin{defn}[ترمونه] \ollabel{defn:term}
د لامبډا حساب د \emph{ترمونو} سټ په استقرايي ډول داسې تعريفېږي:
\begin{enumerate}
  \item \ollabel{defn:term-var} که \LR{$x$} متغير وي، نو \LR{$x$} ترم دے۔
  \item \ollabel{defn:term-abs} که \LR{$x$} متغير او \LR{$M$} ترم وي، نو
    \LR{$(\lambd[x][M])$} ترم دے۔
  \item \ollabel{defn:term-app} که \LR{$M$} او \LR{$N$} دواړه ترمونه وي، نو
    \LR{$(MN)$} ترم دے۔
\end{enumerate}
\end{defn}
\olpseganchor{OLP-0357-B006}{end}

\olpseganchor{OLP-0357-B007}{start}
که ترم \LR{$(\lambd[x][M])$} د \olref{defn:term-abs} له مخې جوړ شوے
وي، نو وايو چې دا د \emph{تجريد} پايله ده، او په \LR{$\lambd[x]$} کښې
\LR{$x$} ته \emph{پارامېټر} وايو۔ که ترم \LR{$(MN)$} د
\olref{defn:term-app} له مخې جوړ شوے وي، نو دا د \emph{تطبيق}
پايله ده۔
\olpseganchor{OLP-0357-B007}{end}

\olpseganchor{OLP-0357-B008}{start}
پاس تعريف شوي ترمونه ټولې قوسې په بشپړه توګه لري۔ دا کار يو څه
دروند کېدے شي، لکه ترم
\LR{$(\lambd[x][((\lambd[x][x])(\lambd[x][(xx)]))])$} چې ښيي۔ د
قوسونو د پرېښودلو دپاره به نښيز دودونه وټاکو۔ خو رسمي تعريف دا
آسانوي چې د \olref{defn:term} له مخې د هر ترم د جوړېدو طريقه
وپېژنو۔ د بېلګې په توګه، د ترم
\LR{$(\lambd[x][((\lambd[x][x])(\lambd[x][(xx)]))])$} د جوړېدو وروستے
ګام بايد تجريد وي چې پارامېټر ئې~\LR{$x$} دے۔ دا ترم له
\LR{$((\lambd[x][x])(\lambd[x][(xx)]))$} څخه د تجريد په وسيله جوړېږي؛
او دغه ترم د دوو ترمونو تطبيق دے۔ له دغو دوو څخه هر ترم بيا د
تجريد پايله ده، او همداسې نور۔
\olpseganchor{OLP-0357-B008}{end}

\olpseganchor{OLP-0357-B009}{start}
\begin{prob}
د \LR{$(\lambd[g][((\lambd[x][(g (x x))])
  (\lambd[x][(g (x x))]))])$} د جوړېدو طريقه بيان کړئ۔
\end{prob}
\olpseganchor{OLP-0357-B009}{end}

\olpseganchor{OLP-0358-B004}{start}
\olfileid{lam}{syn}{unq}
\olsection{يکتا لوستونتيا}
\olpseganchor{OLP-0358-B004}{end}

\olpseganchor{OLP-0358-B005}{start}
کېدے شي وپوښتو چې ايا د هر ترم د جوړېدو يوازې يوه لاره شته؛ هو،
همداسې ده۔ د هر لامبډا ترم د جوړولو او تفسير يوازې يوه لاره شته۔
په لاندې بحث کښې \emph{جوړونه} هغه کړنلاره ده چې د
\olref[trm]{defn:term} د جوړولو قاعدې يو يا څو ځله کاروي او ترم
ترې جوړوي۔
\olpseganchor{OLP-0358-B005}{end}

\olpseganchor{OLP-0358-B006}{start}
\begin{lem}\ollabel{lem:term-start}
ترم يا په متغير پيلېږي يا په قوس۔
\end{lem}
\olpseganchor{OLP-0358-B006}{end}

\olpseganchor{OLP-0358-B007}{start}
\begin{proof}
يو څيز هله ترم ګڼل کېږي چې د \olref[trm]{defn:term} له مخې جوړ
شوے وي۔ که د \olref[trm]{defn:term-var} پايله وي، نو بايد متغير
وي۔ که د \olref[trm]{defn:term-abs} يا
\olref[trm]{defn:term-app} پايله وي، نو په قوس پيلېږي۔
\end{proof}
\olpseganchor{OLP-0358-B007}{end}

\olpseganchor{OLP-0358-B008}{start}
\begin{lem}\ollabel{lem:app-start}
د تطبيق پايله يا په دوو قوسونو پيلېږي يا په يوه قوس او بيا متغير۔
\end{lem}
\olpseganchor{OLP-0358-B008}{end}

\olpseganchor{OLP-0358-B009}{start}
\begin{proof}
که \LR{$M$} د تطبيق پايله وي، نو بڼه ئې \LR{$(PQ)$} ده؛ ځکه په قوس پيلېږي۔
څرنګه چې \LR{$P$} ترم دے، د \olref{lem:term-start} له مخې يا په قوس
پيلېږي يا په متغير۔
\end{proof}
\olpseganchor{OLP-0358-B009}{end}

\olpseganchor{OLP-0358-B010}{start}
\begin{lem}\ollabel{lem:initial}
د ترم هېڅ خاصه پيلنۍ برخه پخپله ترم نۀ ده۔
\end{lem}
\olpseganchor{OLP-0358-B010}{end}

\olpseganchor{OLP-0358-B011}{start}
\begin{prob}
\olref[lam][syn][unq]{lem:initial} د ترمونو پر اوږدوالي په استقرا
ثابت کړئ۔
\end{prob}
\olpseganchor{OLP-0358-B011}{end}

\olpseganchor{OLP-0358-B012}{start}
\begin{prop}[يکتا لوستونتيا] \ollabel{prop:unq}
د هر ترم جوړونه يکتا ده۔ په بله وينا، که ترم \LR{$M$} د يوې جوړونې په
وسيله جوړ شي، نو د همدې ترم د جوړولو بله جوړونه نشته۔
\end{prop}
\olpseganchor{OLP-0358-B012}{end}

\olpseganchor{OLP-0358-B013}{start}
\begin{proof}
  دا د ترمونو پر جوړونه په استقرا ثابتوو۔
  \begin{enumerate}
    \item \LR{$M$} د \LR{$x$} په بڼه دے، چې \LR{$x$} يو متغير دے۔ د تجريد او
      تطبيق پايلې تل په قوسونو پيلېږي، نو د \LR{$M$} د جوړولو دپاره نۀ شي
      کارېدے؛ ځکه د \LR{$M$} جوړونه بايد د
      \LR{\olref[trm]{defn:term}(\olref[trm]{defn:term-var})} يوازې يو ګام
      وي۔
    \item \LR{$M$} د \LR{$(\lambd[x][N])$} په بڼه دے، چې \LR{$x$} متغير او \LR{$N$}
      ترم دے۔ د \LR{\olref[trm]{defn:term}(\olref[trm]{defn:term-var})}
      له مخې نۀ شي جوړېدے، ځکه يوازې يو متغير نۀ دے۔ د
      \olref{lem:app-start} له مخې د تطبيق پايله هم نۀ ده۔ نو \LR{$M$}
      يوازې پر~\LR{$N$} د تجريد پايله کېدے شي۔ د استقرا د فرضيې له مخې
      د \LR{$N$} جوړونه پخپله يکتا ده۔
    \item \LR{$M$} د \LR{$(PQ)$} په بڼه دے، چې \LR{$P$} او \LR{$Q$} ترمونه دي۔ څرنګه
      چې په قوس پيلېږي، د
      \LR{\olref[trm]{defn:term}(\olref[trm]{defn:term-var})} له مخې هم
      نۀ شي جوړېدے۔ د \olref{lem:term-start} له مخې \LR{$P$} په
      \LR{$\lambd$} نۀ شي پيلېدے، نو \LR{$(PQ)$} د تجريد پايله نۀ شي کېدے۔
      اوس فرض کړئ چې \LR{$M$} د تطبيق له لارې د جوړېدو بله طريقه هم لري؛
      يعنې د \LR{$(P'Q')$} په بڼه هم دے۔ بيا \LR{$P$} د \LR{$P'$} خاصه پيلنۍ برخه
      ده، يا برعکس؛ خو دا د \olref{lem:initial} له مخې ناشوني دي۔
      نو \LR{$P$} او \LR{$Q$} په يکتا توګه ټاکل شوي دي، او د استقرا د فرضيې
      له مخې د \LR{$P$} او \LR{$Q$} جوړونې يکتا دي۔
  \end{enumerate}
\end{proof}
\olpseganchor{OLP-0358-B013}{end}

\olpseganchor{OLP-0358-B014}{start}
د پاسنۍ قضيې د لوستلو آسانه بڼه دا ده:
\begin{prop}
  ترم \LR{$M$} يوازې له لاندې بڼو څخه يوه لرلے شي:
  \begin{enumerate}
    \item \LR{$x$}، چې \LR{$x$} د \LR{$M$} له مخې په يکتا توګه ټاکل شوے متغير دے۔
    \item \LR{$(\lambd[x][N])$}، چې \LR{$x$} متغير او \LR{$N$} بل ترم دے، او
      دواړه د \LR{$M$} له مخې په يکتا توګه ټاکل شوي دي۔
    \item \LR{$(PQ)$}، چې \LR{$P$} او \LR{$Q$} د \LR{$M$} له مخې په يکتا توګه ټاکل شوي
      دوه ترمونه دي۔
  \end{enumerate}
\end{prop}
\olpseganchor{OLP-0358-B014}{end}

\olpseganchor{OLP-0359-B004}{start}
\olfileid{lam}{syn}{abb}
\olsection{د نحو لنډې نښې}
\olpseganchor{OLP-0359-B004}{end}

\olpseganchor{OLP-0359-B005}{start}
د \olref[trm]{defn:term} له مخې تعريف شوي ترمونه کله کله په ليکلو
کښې درانۀ وي؛ نو د نحو د لا لنډې بڼې ټاکل ګټور دي۔ البته بايد پام
وکړو چې په لنډه نښه کښې ترمونه هم په يکتا توګه ولوستل شي۔ يوه ډېره
کارېدونکې بڼه، چې \emph{لنډ شوي ترمونه} بلل کېږي، داسې ده:
\olpseganchor{OLP-0359-B005}{end}

\olpseganchor{OLP-0359-B006}{start}
\begin{enumerate}
\item چې قوسونه پرېښودل شي، تطبيق له کيڼ نه ښي ته کېږي۔ مثلاً، که
  \LR{$M$}، \LR{$N$}، \LR{$P$} او~\LR{$Q$} ترمونه وي، نو \LR{$MNPQ$} د \LR{$(((MN)P)Q)$}
  لنډه نښه ده۔
\item همداسې، چې قوسونه پرېښودل شي، د لامبډا تجريد د امکان تر
  ټولو پراخه ساحه اخلي۔ مثلاً، \LR{$\lambd[x][MNP]$} د
  \LR{$(\lambd[x][MNP])$} په توګه لوستل کېږي۔
\item يوه لامبډا د څو متغيرونو د تجريد دپاره کارېدے شي۔ مثلاً،
  \LR{$\lambd[xyz][M]$} د \LR{$\lambd[x][\lambd[y][\lambd[z][M]]]$}
  لنډه نښه ده۔
\end{enumerate}
\olpseganchor{OLP-0359-B006}{end}

\olpseganchor{OLP-0359-B007}{start}
د بېلګې په توګه،
\[
\lambd[xy][xxyx \lambd[z][xz]]
\]
د دې لنډه نښه ده:
\[
(\lambd[x][(\lambd[y][((((xx)y)x)(\lambd[z][(xz)]))])]).
\]
\olpseganchor{OLP-0359-B007}{end}

\olpseganchor{OLP-0359-B008}{start}
\begin{prob}
لنډ شوے ترم \LR{$\lambd[g][(\lambd[x][g (x x)])
  \lambd[x][g (x x)]]$} په بشپړه بڼه وليکئ۔
\end{prob}
\olpseganchor{OLP-0359-B008}{end}

\olpseganchor{OLP-0360-B004}{start}
\olfileid{lam}{syn}{fv}
\olsection{ازاد متغيرونه}
\olpseganchor{OLP-0360-B004}{end}

\olpseganchor{OLP-0360-B005}{start}
لامبډا حساب د تابعو په اړه دے، او تابعې د لامبډا تجريد له لارې
رامنځته کېږي۔ په شهودي ډول، \LR{$\lambd[x][M]$} هغه تابع ده چې قيمتونه
ئې~\LR{$M$} ټاکي، کله چې د تابعې آرګومېنټ~\LR{$x$} ته ورکړل شي۔ خو په~\LR{$M$}
کښې د~\LR{$x$} هر مورد اړوند نۀ دے: که په~\LR{$M$} کښې بل تجريد
\LR{$\lambd[x][N]$} وي، نو په~\LR{$N$} کښې د \LR{$x$} موردونه د
\LR{$\lambd[x][N]$} دپاره اړوند دي، خو د \LR{$\lambd[x][M]$} دپاره نۀ دي۔
نو په \LR{$\lambd[x][M]$} کښې دننه يو لامبډا تجريد \LR{$\lambd[x]$} د~\LR{$M$}
هغه د \LR{$x$} موردونه \emph{تړي} چې لا د بل لامبډا تجريد له خوا تړل
شوي نۀ وي---يعنې په~\LR{$M$} کښې د~\LR{$x$} \emph{ازاد} موردونه۔
\olpseganchor{OLP-0360-B005}{end}

\olpseganchor{OLP-0360-B006}{start}
\begin{defn}[ساحه]
که \LR{$\lambd[x][M]$} په يوه ترم~\LR{$N$} کښې واقع شي، نو د~\LR{$M$} اړوند
مورد د~\LR{$\lambd[x]$} د اړوند مورد \emph{ساحه} ده۔
\end{defn}
\olpseganchor{OLP-0360-B006}{end}

\olpseganchor{OLP-0360-B007}{start}
\begin{defn}[ازاد او تړلے مورد]
  په يوه ترم~\LR{$M$} کښې د متغير \LR{$x$} مورد هله \emph{ازاد} دے چې د
  \LR{$\lambd[x]$} په ساحه کښې نۀ وي، او که وي نو \emph{تړلے} دے۔ په
  \LR{$\lambd[x][M]$} کښې د متغير~\LR{$x$} يو مورد هله او يوازې هله د لومړي
  \LR{$\lambd[x]$} له خوا تړلے دے چې په~\LR{$M$} کښې د \LR{$x$} همدا مورد ازاد وي۔
\end{defn}
\olpseganchor{OLP-0360-B007}{end}

\olpseganchor{OLP-0360-B008}{start}
\begin{ex}
  په \LR{$\lambd[x][x y]$} کښې \LR{$x$} او \LR{$y$} دواړه د \LR{$\lambd[x]$} په
  ساحه کښې دي، نو \LR{$x$} د \LR{$\lambd[x]$} له خوا تړلے دے۔ څرنګه چې \LR{$y$}
  د هېڅ \LR{$\lambd[y]$} په ساحه کښې نۀ دے، ازاد دے۔ په
  \LR{$\lambd[x][x x]$} کښې د \LR{$x$} دواړه موردونه د~\LR{$\lambd[x]$} له خوا
  تړلي دي، ځکه دواړه په \LR{$xx$} کښې ازاد دي۔ په \LR{$((\lambd[x][xx])x)$}
  کښې د~\LR{$x$} وروستے مورد ازاد دے، ځکه د \LR{$\lambd[x]$} په ساحه کښې نۀ
  دے۔ په \LR{$\lambd[x][(\lambd[x][x])x]$} کښې د لومړي \LR{$\lambd[x]$}
  ساحه \LR{$(\lambd[x][x])x$} ده او د دويم \LR{$\lambd[x]$} ساحه د~\LR{$x$} له
  پای څخه دويم مورد دے۔ په \LR{$(\lambd[x][x])x$} کښې د \LR{$x$} وروستے
  مورد ازاد او له پای څخه دويم مورد تړلے دے۔ نو په
  \LR{$\lambd[x][(\lambd[x][x])x]$} کښې د \LR{$x$} له پای څخه دويم مورد د
  دويم \LR{$\lambd[x]$} له خوا او وروستے مورد د لومړي \LR{$\lambd[x]$}
  له خوا تړلے دے۔
\end{ex}
\olpseganchor{OLP-0360-B008}{end}

\olpseganchor{OLP-0360-B009}{start}
په يوه ترم \LR{$P$} کښې د متغيرونو ټول موردونه کتلے شو او د ازادو
متغيرونو سټ ترې موندلے شو۔ دا سټ په \LR{$\FV{P}$} ښيو؛ طبيعي تعريف
ئې داسې دے:
\olpseganchor{OLP-0360-B009}{end}

\olpseganchor{OLP-0360-B010}{start}
\begin{defn}[د ترم ازاد متغيرونه] \ollabel{def:fv}
  د ترم د \emph{ازادو متغيرونو} سټ په استقرايي ډول داسې تعريفېږي:
  \begin{enumerate}
    \item \LR{$\FV{x} = \{x\}$} \ollabel{def:fv1}
    \item \LR{$\FV{\lambd[x][N]} = \FV{N} \setminus \{x\}$}    \ollabel{def:fv2}
    \item \LR{$\FV{PQ} = \FV{P} \cup \FV{Q}$} \ollabel{def:fv3}
  \end{enumerate}
\end{defn}
\olpseganchor{OLP-0360-B010}{end}

\olpseganchor{OLP-0360-B011}{start}
\begin{prob}
  \begin{enumerate}
  \item په دې ترم کښې د \LR{$\lambd[g]$} او دواړو \LR{$\lambd[x]$} ساحې
    وټاکئ: \LR{$\lambd[g][(\lambd[x][g (x x)]) \lambd[x][g (x x)]]$}۔
  \item په \LR{$\lambd[g][(\lambd[x][g (x x)]) \lambd[x][g (x x)]]$}
    کښې ايا د متغيرونو ټول موردونه تړلي دي؟ هر يو د کوم تجريد له
    خوا تړلے دے؟
    \item \LR{$\FV{\lambd[x][(\lambd[y][(\lambd[z][x y]) z]) y]}$}
    ومومئ۔
  \end{enumerate}
\end{prob}
\olpseganchor{OLP-0360-B011}{end}

\olpseganchor{OLP-0360-B012}{start}
\begin{explain}
ازاد متغير له باندنۍ نړۍ (\emph{چاپېريال}) سره د يوې حوالې په شان
دے، او هغه ترم چې ازاد متغيرونه لري د نيمګړي ټاکل شوي ترم په توګه
کتلے شو، ځکه چلند ئې د چاپېريال په ټاکلو پورې تړلے دے۔ مثلاً، په
ترم \LR{$\lambd[x][f x]$} کښې، چې آرګومېنټ~\LR{$x$} اخلي او پر هغه آرګومېنټ
د~\LR{$f$} د تطبيق پايله ورکوي، متغير~\LR{$f$} ازاد دے۔ د دې ترم قيمت په خپل چاپېريال پورې
تړلے دے، په تېره بيا په هغه چاپېريال کښې د \LR{$f$} په قيمت پورې۔
\olpseganchor{OLP-0360-B012}{end}

\olpseganchor{OLP-0360-B013}{start}
که پر دې ترم تجريد وکړو، \LR{$\lambd[f][\lambd[x][f
    x]]$} ترلاسه کوو۔ دا ترم نور د چاپېريال پر متغير \LR{$f$} تړلے نۀ
دے، ځکه اوس داسې تابع ښيي چې دوه آرګومېنټونه اخلي او د لومړي پر
دويم د تطبيق پايله ورکوي۔ په چاپېريال کښې د \LR{$f$} بدلون د دې ترم پر
چلند هېڅ اغېز نۀ لري؛ ترم يوازې هغه څۀ کاروي چې د آرګومېنټ په توګه
ورکول کېږي، نۀ په چاپېريال کښې د \LR{$f$} قيمت۔
\end{explain}
\olpseganchor{OLP-0360-B013}{end}

\olpseganchor{OLP-0360-B014}{start}
\begin{defn}[تړلے ترم، ترکيب کوونکے]
  هغه ترم چې ازاد متغيرونه ونۀ لري، \emph{تړلے ترم} يا
  \emph{ترکيب کوونکے} بلل کېږي۔
\end{defn}
\olpseganchor{OLP-0360-B014}{end}

\olpseganchor{OLP-0360-B015}{start}
\begin{lem}\ollabel{lem:fv}
  \begin{enumerate}
    \item \ollabel{lem:fv-abs} که \LR{$y \neq x$} وي، نو \LR{$y \in
      \FV{\lambd[x][N]}$} هله او يوازې هله چې \LR{$y \in \FV{N}$} وي۔
    \item \ollabel{lem:fv-app} \LR{$y \in \FV{PQ}$} هله او يوازې هله
      چې \LR{$y \in \FV{P}$} يا \LR{$y \in \FV{Q}$} وي۔
    \end{enumerate}
\end{lem}
\olpseganchor{OLP-0360-B015}{end}

\olpseganchor{OLP-0360-B016}{start}
\begin{proof}
  تمرين۔
\end{proof}
\olpseganchor{OLP-0360-B016}{end}

\olpseganchor{OLP-0360-B017}{start}
\begin{prob}
  \olref[lam][syn][fv]{lem:fv} ثابت کړئ۔
\end{prob}
\olpseganchor{OLP-0360-B017}{end}

\olpseganchor{OLP-0361-B004}{start}
\olfileid{lam}{syn}{sub}
\olsection{تعويض}
\olpseganchor{OLP-0361-B004}{end}

\olpseganchor{OLP-0361-B005}{start}
\begin{explain}
ازاد متغيرونه د چاپېريال متغيرونو ته حوالې دي؛ نو دا معنا لري چې د
ازاد متغير پر ځاے کوم ټاکلے قيمت وکاروو۔ مثلاً، کېدے شي په
\LR{$\lambd[x][f x]$} کښې د \LR{$f$} پر ځاے يو ټاکلے ترم، لکه د هماني تابع
\LR{$\lambd[y][y]$}، کېږدو۔ په دې سره \LR{$\lambd[x][(\lambd[y][y]) x]$}
ترلاسه کېږي۔ د لامبډا ترمونو په وسيله د ازادو متغيرونو د بدلولو دې
بهير ته \emph{تعويض} وايي۔
\end{explain}
\olpseganchor{OLP-0361-B005}{end}

\olpseganchor{OLP-0361-B006}{start}
\begin{defn}[تعويض] \ollabel{defn:substitution}
  په ترم~\LR{$M$} کښې د متغير~\LR{$x$} پر ځاے د ترم~\LR{$N$} \emph{تعويض}،
  \LR{$\Subst{M}{N}{x}$}، په استقرايي ډول داسې تعريفېږي:
  \begin{enumerate}
    \item \LR{$\Subst{x}{N}{x} = N$}. \ollabel{defn:substitution-1}
    \item \LR{$\Subst{y}{N}{x}  = y$} که \LR{$x \neq y$} وي۔ \ollabel{defn:substitution-2}
    \item \LR{$\Subst{PQ}{N}{x}  = (\Subst{P}{N}{x}) (\Subst{Q}{N}{x})$}.
      \ollabel{def:substitution-3}
    \item \LR{$\Subst{(\lambd[y][P])}{N}{x} = \lambd[y][\Subst{P}{N}{x}]$}،
      که \LR{$x \neq y$} او \LR{$y \notin \FV{N}$} وي؛ که نه، تعويض ناټاکلے دے۔
      \ollabel{defn:substitution-4}
  \end{enumerate}
\end{defn}
\olpseganchor{OLP-0361-B006}{end}

\olpseganchor{OLP-0361-B007}{start}
\begin{explain}
د \olref{defn:substitution} د \olref{defn:substitution-4} بند له مخې، \LR{$x
\neq y$} ځکه شرط دے چې په~\LR{$M$} کښې د متغير~\LR{$x$} \emph{تړلي} موردونه
په~\LR{$N$} بدل نۀ کړو۔ مثلاً، که د
\LR{$\Subst{\lambd[x][x]}{y}{x}$} پر معنا غور وکړو، د تړلي متغير بدلول
به \LR{$\lambd[x][y]$} ورکړي، حال دا چې د ازادو موردونو د بدلولو
مقصودي پايله هماغه \LR{$\lambd[x][x]$} ده۔ د پورته رسمي تعريف له مخې
خو په دې حالت کښې تعويض ناټاکلے دے۔
\olpseganchor{OLP-0361-B007}{end}

\olpseganchor{OLP-0361-B008}{start}
کله چې په \LR{$\lambd[y][P]$} کښې د~\LR{$x$} پر ځاے~\LR{$N$} تعويضوو، دا هم شرط
دے چې \LR{$y \notin \FV{N}$} وي۔ مثلاً، په \LR{$\lambd[y][x]$} کښې د \LR{$x$}
پر ځاے \LR{$y$}، يعنې \LR{$\Subst{\lambd[y][x]}{y}{x}$}، نۀ شو تعويضولے،
ځکه پايله به ئې \LR{$\lambd[y][y]$} شي، چې داسې تابع ده چې آرګومېنټ
اخلي او هماغه سمدستي بېرته ورکوي۔ خو \LR{$\lambd[y][x]$} داسې تابع ده
چې تل ترم~\LR{$x$} (يا هغه څۀ چې \LR{$x$} ورته اشاره کوي) ورکوي۔ نو هغه
پايله چې په رښتيا ئې غواړو داسې تابع ده چې آرګومېنټ واخلي، له
پامه ئې وغورځوي او د چاپېريال متغير~\LR{$y$} ورکړي۔ د دې د سم ترسره
کولو دپاره بايد لومړے تړلے متغير~\LR{$y$} «بيا ونوموو»۔
\end{explain}
\olpseganchor{OLP-0361-B008}{end}

\olpseganchor{OLP-0361-B009}{start}
\begin{prob}
  د لاندې تعويضونو پايله څۀ ده؟
  \begin{enumerate}
  \item \LR{$\Subst{\lambd[y][x(\lambd[w][vwx])]}{(uv)}{x}$}
  \item \LR{$\Subst{\lambd[y][x(\lambd[x][x])]}{(\lambd[y][xy])}{x}$}
  \item \LR{$\Subst{y(\lambd[v][xv])}{(\lambd[y][vy])}{x}$}
  \end{enumerate}
\end{prob}
\olpseganchor{OLP-0361-B009}{end}

\olpseganchor{OLP-0361-B010}{start}
\begin{thm} \ollabel{thm:notinfv}
  که \LR{$x \notin \FV{M}$} وي، نو، په دې شرط چې کيڼ لورے ئې ټاکلے
  وي، \LR{$\FV{\Subst{M}{N}{x}} = \FV{M}$} کېږي۔
\end{thm}
\olpseganchor{OLP-0361-B010}{end}

\olpseganchor{OLP-0361-B011}{start}
\begin{proof}
  د~\LR{$M$} د جوړښت پر بنسټ استقرا کوو۔
  \begin{enumerate}
  \item \LR{$M$} يو متغير دے: تمرين۔
  \item \LR{$M$} د \LR{$(PQ)$} په بڼه دے: تمرين۔
  \item \LR{$M$} د \LR{$\lambd[y][P]$} په بڼه دے؛ څرنګه چې
    \LR{$\Subst{\lambd[y][P]}{N}{x}$} ټاکلے دے، نو دا بايد
    \LR{$\lambd[y][\Subst{P}{N}{x}]$} وي۔ له دې کبله
    \LR{$\Subst{P}{N}{x}$} هم ټاکلے دے؛ سربېره پر دې، \LR{$x \neq y$} او
    \LR{$x \notin \FV{P}$}۔ نو:
    \begin{multline*}
      \FV{\Subst{\lambd[y][P]}{N}{x}} \\
    \begin{aligned}
      & = \FV{\lambd[y][\Subst{P}{N}{x}]} &&
      \text{\RL{د \olref{defn:substitution-4} له مخې}}\\
      & = \FV{\Subst{P}{N}{x}} \setminus \{y\} &&
      \text{\RL{د \LR{\olref[fv]{def:fv}(\olref[fv]{def:fv2})} له مخې}}\\
      & = \FV{P} \setminus \{y\} && \text{\RL{د استقرا د فرضيې له مخې}} \\
      & = \FV{\lambd[y][P]} && \text{\RL{د \LR{\olref[fv]{def:fv}(\olref[fv]{def:fv2})} له مخې}}
    \end{aligned}
    \end{multline*}
  \end{enumerate}
\end{proof}
\olpseganchor{OLP-0361-B011}{end}

\olpseganchor{OLP-0361-B012}{start}
\begin{prob}
  د \olref[lam][syn][sub]{thm:notinfv} ثبوت بشپړ کړئ۔
\end{prob}
\olpseganchor{OLP-0361-B012}{end}

\olpseganchor{OLP-0361-B013}{start}
\begin{thm} \ollabel{thm:infv}
  که \LR{$x \in \FV{M}$} وي، نو، په دې شرط چې کيڼ لورے ئې ټاکلے وي،
  \LR{$\FV{\Subst{M}{N}{x}} = (\FV{M} \setminus
  \{x\}) \cup \FV{N}$} کېږي۔
\end{thm}
\olpseganchor{OLP-0361-B013}{end}

\olpseganchor{OLP-0361-B014}{start}
\begin{proof}
  د~\LR{$M$} د جوړښت پر بنسټ استقرا کوو۔
  \begin{enumerate}
  \item \LR{$M$} يو متغير دے: تمرين۔
  \item \LR{$M$} د \LR{$PQ$} په بڼه دے: څرنګه چې
    \LR{$\Subst{(PQ)}{N}{x}$} ټاکلے دے، نو دا بايد
    \LR{$(\Subst{P}{N}{x})(\Subst{Q}{N}{x})$} وي او دواړه دننني
    تعويضونه هم ټاکلي وي۔ همدارنګه، چې \LR{$x \in \FV{PQ}$} دے، نو
    \LR{$x \in \FV{P}$} يا \LR{$x \in \FV{Q}$} يا دواړه۔ پاتې ثبوت تمرين دے۔
  \item \LR{$M$} د \LR{$\lambd[y][P]$} په بڼه دے۔ څرنګه چې
    \LR{$\Subst{\lambd[y][P]}{N}{x}$} ټاکلے دے، نو دا بايد
    \LR{$\lambd[y][\Subst{P}{N}{x}]$} وي؛ \LR{$\Subst{P}{N}{x}$} هم
    ټاکلے دے، \LR{$x \neq y$} او \LR{$y \notin \FV{N}$} دي۔ له دې سره،
    ځکه چې \LR{$x \in \FV{\lambd[y][P]}$} دے، \LR{$x \in \FV{P}$} هم دے۔
    اوس:
    \begin{multline*}
      \FV{\Subst{(\lambd[y][P])}{N}{x}} \\
      \begin{aligned}
      & = \FV{\lambd[y][\Subst{P}{N}{x}]} \\
      & = \FV{\Subst{P}{N}{x}} \setminus \{y\} \\
      & = ((\FV{P} \setminus \{x\}) \cup \FV{N}) \setminus \{y\}
        && \text{\RL{د استقرا د فرضيې له مخې}} \\
      & = (\FV{P} \setminus \{x, y\}) \cup \FV{N}
        && y \notin \FV{N} \\
      & = (\FV{\lambd[y][P]} \setminus \{x\}) \cup \FV{N}
      \end{aligned}
    \end{multline*}
  \end{enumerate}
\end{proof}
\olpseganchor{OLP-0361-B014}{end}

\olpseganchor{OLP-0361-B015}{start}
\begin{prob}
  د \olref[lam][syn][sub]{thm:infv} ثبوت بشپړ کړئ۔
\end{prob}
\olpseganchor{OLP-0361-B015}{end}

\olpseganchor{OLP-0361-B016}{start}
\begin{thm}\ollabel{thm:clr}
  که ښی لورے ټاکلے وي او \LR{$x \notin \FV{N}$} وي، نو
  \LR{$x \notin \FV{\Subst{M}{N}{x}}$}۔
\end{thm}
\olpseganchor{OLP-0361-B016}{end}

\olpseganchor{OLP-0361-B017}{start}
\begin{proof}
  تمرين۔
\end{proof}
\olpseganchor{OLP-0361-B017}{end}

\olpseganchor{OLP-0361-B018}{start}
\begin{prob}
  \olref[lam][syn][sub]{thm:clr} ثابت کړئ۔
\end{prob}
\olpseganchor{OLP-0361-B018}{end}


\olpseganchor{OLP-0361-B019}{start}
\begin{thm}\ollabel{thm:inv}
  که \LR{$\Subst{M}{y}{x}$} ټاکلے وي او \LR{$y \notin \FV{M}$} وي، نو
  \LR{$\Subst{\Subst{M}{y}{x}}{x}{y} = M$}۔
\end{thm}
\olpseganchor{OLP-0361-B019}{end}

\olpseganchor{OLP-0361-B020}{start}
\begin{proof}
  د~\LR{$M$} د جوړښت پر بنسټ استقرا کوو۔
  \begin{enumerate}
  \item \LR{$M$} يو متغير \LR{$z$} دے: تمرين۔
  \item \LR{$M$} د \LR{$(PQ)$} په بڼه دے۔ نو:
    \begin{align*}
      \Subst{\Subst{(PQ)}{y}{x}}{x}{y}
      &=\Subst{((\Subst{P}{y}{x})(\Subst{Q}{y}{x}))}{x}{y} \\
      &= (\Subst{\Subst{P}{y}{x}}{x}{y})(\Subst{\Subst{Q}{y}{x}}{x}{y}) \\
      &= (PQ) \text{\RL{ د استقرا د فرضيې له مخې}}
    \end{align*}
  \item \LR{$M$} د \LR{$\lambd[z][N]$} په بڼه دے۔ څرنګه چې
    \LR{$\Subst{\lambd[z][N]}{y}{x}$} ټاکلے دے، نو \LR{$z \neq x$} او
    \LR{$z \neq y$} دي او دنننے تعويض هم ټاکلے دے۔ د قضيې له بل فرض
    او د تړونکي له نابرابرۍ څخه راځي چې y د N ازاد متغير نۀ دے۔
    په معکوس تعويض کښې هم د تړونکي شرطونه د همدې دوو نابرابريو
    له مخې پوره دي۔ نو:
    \begin{align*}
      \Subst{\Subst{(\lambd[z][N])}{y}{x}}{x}{y}\\
      & = \Subst{(\lambd[z][\Subst{N}{y}{x}])}{x}{y} \\
      & = \lambd[z][\Subst{\Subst{N}{y}{x}}{x}{y}] \\
      &= \lambd[z][N] \text{\RL{ د استقرا د فرضيې له مخې}}
    \end{align*}
  \end{enumerate}
\end{proof}
\olpseganchor{OLP-0361-B020}{end}

\olpseganchor{OLP-0361-B021}{start}
\begin{prob}
  د \olref[lam][syn][sub]{thm:inv} ثبوت بشپړ کړئ۔
\end{prob}
\olpseganchor{OLP-0361-B021}{end}

\paragraph{د منجمدې سرچينې بشپړه اړونده برخه}
\olpseganchor{OLP-0651-B004}{start}
\olfileid{lam}{int}{con-olp0651}
\olsection{بدلون او راکمول}
\olpseganchor{OLP-0651-B004}{end}

\olpseganchor{OLP-0651-B005}{start}
په لامبډا حساب کښې د بدلون يا راکمول درې ډولونه شته۔ د بدلون او راکمول توپير دا دے چې «اړول» په يو لوري کېږي که په دواړو لوريو؛ دا به وروسته ووينو۔
\olpseganchor{OLP-0651-B005}{end}

\olpseganchor{OLP-0651-B006}{start}
\LR{$\alpha$}-بدلون د متغيرونو د نومونو بدلول دي۔ دا طبيعي ده چې \LR{$\lambda x. x$} او \LR{$\lambda y.y$} د يوې تابعې، يعنې د عينيت تابعې، تمثيل وګڼو۔ په راتلونکې برخه کښې د تړلي متغير د نوم بدلولو کره شرطونه ټاکو۔
\olpseganchor{OLP-0651-B006}{end}

\olpseganchor{OLP-0362-B004}{start}
\olfileid{lam}{syn}{alp}
\olsection{\texorpdfstring{\LR{$\alpha$}}{α}-بدلون}
\olpseganchor{OLP-0362-B004}{end}

\olpseganchor{OLP-0362-B005}{start}
د \LR{$\lambd[x][x]$} او \LR{$\lambd[y][y]$} تر منځ څۀ اړيکه ده؟ دواړه د
هماني تابع استازيتوب کوي۔ البته د جوړښت له پلوه بېل ترمونه دي۔
توپير ئې يوازې د تړلي متغير په نوم کښې دے؛ يو د بل د تړلي متغير
د «بيا نومولو» پايله ده۔ دې ته \emph{\LR{$\alpha$}-بدلون} وايي۔
\olpseganchor{OLP-0362-B005}{end}

\olpseganchor{OLP-0362-B006}{start}
\begin{defn}[د تړلي متغير د نوم بدلون، \LR{$\aconvone$}]
  که ترم~\LR{$M$} د \LR{$\lambd[x][N]$} يو مورد ولري، \LR{$x\neq y$} وي،
  \LR{$y \notin \FV{N}$} وي او \LR{$\Subst{N}{y}{x}$} ټاکلے وي، نو د دې مورد
  په
  \begin{equation*}
    \lambd[y][\Subst{N}{y}{x}]
  \end{equation*}
  بدلول، چې پايله ئې~\LR{$M'$} ده، \emph{د تړلي متغير د نوم بدلون}
  بلل کېږي او \LR{$M \redone[\alpha] M'$} ليکل کېږي۔
\end{defn}
\olpseganchor{OLP-0362-B006}{end}

\olpseganchor{OLP-0362-B007}{start}
\begin{defn}[د اړيکې سازګارتيا]
  پر ترمونو يوه اړيکه~\LR{$R$} \emph{سازګاره} بلل کېږي که دا شرطونه
  پوره کړي:
  \begin{enumerate}
  \item که \LR{$R N N'$} وي، نو \LR{$R \lambd[x][N] \lambd[x][N']$} وي۔
  \item که \LR{$R P P'$} وي، نو \LR{$R (PQ) (P'Q)$} وي۔
  \item که \LR{$R Q Q'$} وي، نو \LR{$R (PQ) (PQ')$} وي۔
  \end{enumerate}
\end{defn}
\olpseganchor{OLP-0362-B007}{end}

\olpseganchor{OLP-0362-B008}{start}
نو تعريف بيا داسې ويلے شو:
\begin{defn}[د تړلي متغير د نوم بدلون، \LR{$\aconvone$}]
  \emph{د تړلي متغير د نوم بدلون} (\LR{$\redone[\alpha]$}) پر ترمونو
  تر ټولو کوچنۍ سازګاره اړيکه ده چې دا شرط پوره کوي:
  \begin{align*}
    &\lambd[x][N] \redone[\alpha] \lambd[y][\Subst{N}{y}{x}] && \text{\RL{که
      \LR{$x \neq y$}، \LR{$y \notin \FV{N}$} وي}} \\
    & &&\text{\RL{او \LR{$\Subst{N}{y}{x}$} ټاکلے وي}}
  \end{align*}
\end{defn}
\olpseganchor{OLP-0362-B008}{end}

\olpseganchor{OLP-0362-B009}{start}
دلته «تر ټولو کوچنۍ» معنا دا ده چې اړيکه يوازې هغه جوړې لري
چې د سازګارتيا او زيات شوي شرط له مخې لازمي دي؛ نورې جوړې نۀ
لري۔ له همدې امله اړيکه داسې هم تعريفېږي:
\olpseganchor{OLP-0362-B009}{end}

\olpseganchor{OLP-0362-B010}{start}
\begin{defn}[د تړلي متغير د نوم بدلون، \LR{$\aconvone$}] \ollabel{defn:aconvone}
  \emph{د تړلي متغير د نوم بدلون} (\LR{$\aconvone$}) په استقرايي
  ډول داسې تعريفېږي:
  \begin{enumerate}
  \item که \LR{$N \aconvone N'$} وي، نو \LR{$\lambd[x][N] \aconvone
    \lambd[x][N']$} وي۔ \ollabel{defn:aconvone1}
  \item که \LR{$P \aconvone P'$} وي، نو \LR{$(PQ) \aconvone (P'Q)$} وي۔ \ollabel{defn:aconvone2}
  \item که \LR{$Q \aconvone Q'$} وي، نو \LR{$(PQ) \aconvone (PQ')$} وي۔ \ollabel{defn:aconvone3}
  \item که \LR{$x \neq y$}، \LR{$y \notin \FV{N}$} او \LR{$\Subst{N}{y}{x}$}
    ټاکلے وي، نو
    \LR{$\lambd[x][N] \redone[\alpha] \lambd[y][\Subst{N}{y}{x}]$} وي۔
    \ollabel{defn:aconvone4}
  \end{enumerate}
\end{defn}
\olpseganchor{OLP-0362-B010}{end}

\olpseganchor{OLP-0362-B011}{start}
دا تعريفونه معادل دي، خو ثبوت ئې تمرين ته پرېږدو۔ له دې وروسته
به استقرايي تعريف کاروو۔
\olpseganchor{OLP-0362-B011}{end}

\olpseganchor{OLP-0362-B012}{start}
\begin{defn}[\LR{$\alpha$}-بدلون، \LR{$\aconv$}]
  \emph{\LR{$\alpha$}-بدلون} (\LR{$\aconv$}) پر ترمونو تر ټولو کوچنۍ
  انعکاسي او انتقالي اړيکه ده چې \LR{$\aconvone$} پکې شامله وي۔
\end{defn}
\olpseganchor{OLP-0362-B012}{end}

\olpseganchor{OLP-0362-B013}{start}
لکه پورته، «تر ټولو کوچنۍ» معنا دا ده چې اړيکه يوازې هغه جوړې
لري چې د انعکاسيت، انتقاليت او \LR{$\aconvone$} له مخې لازمي دي۔
دا لاندې معادل تعريف ته رسوي:
\olpseganchor{OLP-0362-B013}{end}

\olpseganchor{OLP-0362-B014}{start}
\begin{defn}[\LR{$\alpha$}-بدلون، \LR{$\aconv$}] \ollabel{defn:aconv}
  \emph{\LR{$\alpha$}-بدلون} (\LR{$\aconv$}) په استقرايي ډول داسې
  تعريفېږي:
  \begin{enumerate}
  \item که \LR{$P \aconv Q$} او \LR{$Q \aconv R$} وي، نو \LR{$P \aconv R$} وي۔
    \ollabel{defn:aconv1}
  \item که \LR{$P \aconvone Q$} وي، نو \LR{$P \aconv Q$} وي۔ \ollabel{defn:aconv2}
  \item \LR{$P \aconv P$}. \ollabel{defn:aconv3}
  \end{enumerate}
\end{defn}
\olpseganchor{OLP-0362-B014}{end}

\olpseganchor{OLP-0362-B015}{start}
\begin{ex}
  \LR{$\lambd[x][f x]$} په \LR{$\lambd[y][f
  y]$} \LR{$\alpha$}-بدلېږي او برعکس هم۔ په غيررسمي وينا، دواړه داسې تابعګانې
  دي چې آرګومېنټ اخلي او د هغه آرګومېنټ د~\LR{$f$} پايله ورکوي؛
  \LR{$f$} د چاپېريال متغير ته اشاره کوي۔
\olpseganchor{OLP-0362-B015}{end}

\olpseganchor{OLP-0362-B016}{start}
  \LR{$\lambd[x][f x]$} په \LR{$\lambd[x][g
    x]$} د \LR{$\alpha$}-بدلون نۀ مومي۔ په غيررسمي وينا، لومړنۍ تابع د چاپېريال
  متغير~\LR{$f$} او دويمه د~\LR{$g$} حواله ورکوي؛ ځکه چې په کوم چاپېريال
  کښې \LR{$f$} او \LR{$g$} بېل وي، د دواړو تابعګانو چلند هم بېل وي۔
\end{ex}
\olpseganchor{OLP-0362-B016}{end}

\olpseganchor{OLP-0362-B017}{start}
\begin{prob}
  آيا د ترمونو لاندې جوړې \LR{$\alpha$}-بدلېدونکې دي؟
  \begin{enumerate}
  \item \LR{$\lambd[x][\lambd[y][x]]$} او \LR{$\lambd[y][\lambd[x][y]]$}
  \item \LR{$\lambd[x][\lambd[y][x]]$} او \LR{$\lambd[c][\lambd[b][a]]$}
  \item \LR{$\lambd[x][\lambd[y][x]]$} او \LR{$\lambd[c][\lambd[b][a]]$}
  \end{enumerate}
\end{prob}
\olpseganchor{OLP-0362-B017}{end}

\olpseganchor{OLP-0362-B018}{start}
\begin{lem}\ollabel{lem:fv-one}
  که \LR{$P \aconvone Q$} وي، نو \LR{$\FV{P} = \FV{Q}$} وي۔
\end{lem}
\olpseganchor{OLP-0362-B018}{end}

\olpseganchor{OLP-0362-B019}{start}
\begin{proof}
  د \LR{$P \aconvone Q$} د اشتقاق پر بنسټ استقرا کوو۔
  \begin{enumerate}
  \item که وروستۍ قاعده \olref{defn:aconvone4} وي، نو \LR{$P$} د
    \LR{$\lambd[x][N]$} او \LR{$Q$} د
    \LR{$\lambd[y][\Subst{N}{y}{x}]$} په بڼه دي؛ دلته \LR{$x \neq y$}،
    \LR{$y \notin \FV{N}$} او \LR{$\Subst{N}{y}{x}$} ټاکلے دے۔ اوس د
    \LR{$x \in \FV{N}$} او د دې د نفي حالتونه بېلوو:
    \begin{enumerate}
    \item که \LR{$x \in \FV{N}$} وي، نو:
      \begin{align*}
        \FV{\lambd[y][\Subst{N}{y}{x}]} &=
        \FV{\Subst{N}{y}{x}} \setminus \{y\} \\
        & = ((\FV{N} \setminus \{x\}) \cup \{y\}) \setminus \{y\}
         && \text{\RL{د \olref[sub]{thm:infv} له مخې}} \\
        & = \FV{N} \setminus \{x\} \\
        & = \FV{\lambd[x][N]}
      \end{align*}
    \item که \LR{$x \notin \FV{N}$} وي، نو:
      \begin{align*}
        \FV{\lambd[y][\Subst{N}{y}{x}]}
        & = \FV{\Subst{N}{y}{x}} \setminus \{y\} \\
        & = \FV{N} \setminus \{y\}
         && \text{\RL{د \olref[sub]{thm:notinfv} له مخې}} \\
        & = \FV{N} && y \notin \FV{N} \\
        & = \FV{N} \setminus \{x\} && x \notin \FV{N} \\
        & = \FV{\lambd[x][N]}.
      \end{align*}
    \end{enumerate}
  \item نور درې حالتونه د تمرين دپاره پرېښودل شوي دي۔
  \end{enumerate}
\end{proof}
\olpseganchor{OLP-0362-B019}{end}

\olpseganchor{OLP-0362-B020}{start}
\begin{prob}
  د \olref[lam][syn][alp]{lem:fv-one} ثبوت بشپړ کړئ۔
\end{prob}
\olpseganchor{OLP-0362-B020}{end}

\olpseganchor{OLP-0362-B021}{start}
\begin{lem}\ollabel{lem:inv}
  که \LR{$P \aconvone Q$} وي، نو \LR{$Q \aconvone P$} وي۔
\end{lem}
\olpseganchor{OLP-0362-B021}{end}

\olpseganchor{OLP-0362-B022}{start}
\begin{proof}
  د \LR{$P \aconvone Q$} د اشتقاق پر بنسټ استقرا کوو۔
  \begin{enumerate}
  \item که وروستۍ قاعده \olref{defn:aconvone4} وي، نو \LR{$P$} د
    \LR{$\lambd[x][N]$} او \LR{$Q$} د
    \LR{$\lambd[y][\Subst{N}{y}{x}]$} په بڼه دي؛ \LR{$x \neq y$}،
    \LR{$y \notin \FV{N}$} او \LR{$\Subst{N}{y}{x}$} ټاکلے دے۔ لومړے، د
    \olref[sub]{thm:clr} له مخې
    \LR{$x \notin \FV{\Subst{N}{y}{x}}$} لرو۔ د
    \olref[sub]{thm:inv} له مخې
    \LR{$\Subst{\Subst{N}{y}{x}}{x}{y}$} ټاکلے دے او له~\LR{$N$} سره
    برابر هم دے۔ بيا د \olref{defn:aconvone4} له مخې
    \LR{$\lambd[y][\Subst{N}{y}{x}]
    \aconvone \lambd[x][\Subst{\Subst{N}{y}{x}}{x}{y}] =
    \lambd[x][N]$}۔
  \end{enumerate}
\end{proof}
\olpseganchor{OLP-0362-B022}{end}

\olpseganchor{OLP-0362-B023}{start}
\begin{prob}
  د \olref[lam][syn][alp]{lem:inv} ثبوت بشپړ کړئ۔
\end{prob}
\olpseganchor{OLP-0362-B023}{end}

\olpseganchor{OLP-0362-B024}{start}
\begin{thm}
  \LR{$\alpha$}-بدلون پر ترمونو د معادلتوب اړيکه ده؛ يعنې انعکاسي،
  متناظره او انتقالي ده۔
\end{thm}
\olpseganchor{OLP-0362-B024}{end}

\olpseganchor{OLP-0362-B025}{start}
\begin{proof}
  \begin{enumerate}
  \item هر ترم~\LR{$M$} له تړلي متغيره د نوم له \emph{صفر} بدلونونو
    وروسته له~\LR{$M$} سره يو شان وي؛ يعنې هماغه~\LR{$M$} پاتې کېږي۔
  \item که \LR{$P$} د تړلو متغيرونو د نومونو د بدلونونو په لړۍ سره
    په~\LR{$Q$} \LR{$\alpha$}-بدلېږي، نو له~\LR{$Q$} څخه همدا بدلونونه د
    \olref{lem:inv} له مخې په معکوس ترتيب بېرته ترسره کولے شو
    او~\LR{$P$} تر لاسه کوو۔
  \item که \LR{$P$} د تړلو متغيرونو د نومونو په يوه لړۍ سره په~\LR{$Q$}
    او \LR{$Q$} په بله لړۍ سره په~\LR{$R$} \LR{$\alpha$}-بدلېږي، نو لومړۍ او بيا
    دويمه لړۍ په ترسره کولو \LR{$P$} په~\LR{$R$} بدلېږي۔
  \end{enumerate}
\end{proof}
\olpseganchor{OLP-0362-B025}{end}

\olpseganchor{OLP-0362-B026}{start}
له دې وروسته \LR{$M$} او \LR{$N$} ته \emph{\LR{$\alpha$}-معادل} وايو،
\LR{$M \aeq N$}، که \LR{$M$} په~\LR{$N$} \LR{$\alpha$}-بدلېږي؛ لکه ثابت مو چې کړل،
دا هغه وخت او يوازې هغه وخت کېږي چې \LR{$N$} هم په~\LR{$M$} \LR{$\alpha$}-بدلېږي۔
\olpseganchor{OLP-0362-B026}{end}

\olpseganchor{OLP-0362-B027}{start}
\begin{thm}\ollabel{thm:fv}
  که \LR{$M \aeq N$} وي، نو \LR{$\FV{M} = \FV{N}$} وي۔
\end{thm}
\olpseganchor{OLP-0362-B027}{end}

\olpseganchor{OLP-0362-B028}{start}
\begin{proof}
  دا سمدستي له \olref{lem:fv-one} څخه راځي۔
\end{proof}
\olpseganchor{OLP-0362-B028}{end}

\olpseganchor{OLP-0362-B029}{start}
\begin{lem}\ollabel{lem:sub:R}
  که \LR{$R \aeq R'$} او \LR{$\Subst{M}{R}{y}$} ټاکلے وي، نو
  \LR{$\Subst{M}{R'}{y}$} هم ټاکلے وي او له
  \LR{$\Subst{M}{R}{y}$} سره \LR{$\alpha$}-معادل وي۔
\end{lem}
\olpseganchor{OLP-0362-B029}{end}

\olpseganchor{OLP-0362-B030}{start}
\begin{proof}
  تمرين۔
\end{proof}
\olpseganchor{OLP-0362-B030}{end}

\olpseganchor{OLP-0362-B031}{start}
\begin{prob}
  \olref[lam][syn][alp]{lem:sub:R} ثابت کړئ۔
\end{prob}
\olpseganchor{OLP-0362-B031}{end}

\olpseganchor{OLP-0362-B032}{start}
په \olref[sub]{sec} کښې مو ولوستل چې تعويض په ځينو حالتونو
کښې ناټاکلے وي؛ خو د ترمونو د \LR{$\alpha$}-بدلون په وسيله تړلي
متغيرونه بيا نومولے شو او تعويض ټاکلے کولے شو۔ پايله ئې
\LR{$\alpha$}-معادلتوب ساتي، لکه لاندې قضيه چې وايي:
\olpseganchor{OLP-0362-B032}{end}

\olpseganchor{OLP-0362-B033}{start}
\begin{thm}\ollabel{thm:sub}
  د هر~\LR{$M$}، \LR{$R$} او~\LR{$y$} دپاره داسې~\LR{$M'$} شته چې \LR{$M \aeq M'$}
  او \LR{$\Subst{M'}{R}{y}$} ټاکلے وي۔ سربېره پر دې، که بله جوړه
  \LR{$M'' \aeq M$} او~\LR{$R''$} وي چې \LR{$\Subst{M''}{R''}{y}$} ټاکلے وي
  او \LR{$R'' \aeq R$} وي، نو
  \LR{$\Subst{M'}{R}{y} \aeq \Subst{M''}{R''}{y}$} وي۔
\end{thm}
\olpseganchor{OLP-0362-B033}{end}

\olpseganchor{OLP-0362-B034}{start}
\begin{proof}
  د~\LR{$M$} د جوړښت د غوټو پر شمېر قوي استقرا د وجود برخه
  داسې پر مخ وړي:
  \begin{enumerate}
  \item که \LR{$M$} متغير~\LR{$z$} وي، تعويض پر هماغه متغير ټاکلے دے۔
  \item فرض کړئ \LR{$M$} د \LR{$\lambd[x][N]$} په بڼه دے۔ له~\LR{$x$} او~\LR{$y$}
    څخه بېل داسې متغير~\LR{$z$} غوره کړئ چې \LR{$z \notin \FV{N}$} او
    \LR{$z \notin \FV{R}$} وي۔ د استقرا د فرضيې له مخې داسې~\LR{$N'$} شته
    چې \LR{$N' \aeq N$} او \LR{$\Subst{N'}{z}{x}$} ټاکلے وي۔ نو
    \LR{$\lambd[x][N] \aeq \lambd[x][N']$} هم د
    \olref{defn:aconvone} د \olref{defn:aconvone1} بند له مخې دے۔
    اوس \LR{$\lambd[x][N'] \aeq
    \lambd[z][\Subst{N'}{z}{x}]$} د
    \olref{defn:aconvone} د \olref{defn:aconvone4} بند له مخې دے؛ ځکه
    \LR{$z \ne x$}، \LR{$z \notin \FV{N'}$} او \LR{$\Subst{N'}{z}{x}$} ټاکلے دے۔
    خو له دې څخه لا دا نۀ راځي چې
    \LR{$\Subst{\lambd[z][\Subst{N'}{z}{x}]}{R}{y}$} ټاکلے دے؛
    دنننے تعويض هم بايد ټاکلے وي۔ د دې د تضمين دپاره د لومړي
    تعويض پر پايله د استقرا د وجود برخه بيا کاروو؛ دا پايله له
    اصلي تجريد څخه د جوړښت له پلوه کوچنۍ ده۔ الفا-بدلون
    او د يوه متغير پر ځاے د بل متغير تعويض د جوړښت د غوټو
    شمېر نۀ بدلوي؛ نو دغه بدنه د اصلي تجريد په پرتله يوه
    غوټه کمه لري، او د قوي استقرا کارونه صحيح ده۔ له هغې سره
    الفا-معادله بدنه غوره کوو چې د ټاکلي متغير پر ځاے د ټاکلي
    ترم تعويض پکې ټاکلے وي۔ د دې بدنې تجريد له اصلي ترم سره
    د سازګارتيا له مخې الفا-معادل پاتې کېږي؛ ځکه چې
    \LR{$z \neq y$} او \LR{$z \notin \FV{R}$} دي، د باندې تعويض هم
    ټاکلے وي۔
    اوس د استازي له ټاکنې خپلواکي ثابتوو۔ د هر ترم د جوړښت ونه په
    پام کښې ونيسئ۔ د متغير په هره پاڼه کښې يا د ازاد متغير نوم
    ثبت کړئ، يا د هغې غوټې ځاے چې دغه مورد تړي۔ الفا-بدلون دا
    ثبت نۀ بدلوي: د تړونکي نوم او د هغه تړلي موردونه يو ځاے
    بدلېږي، او د نوم بدلولو شرطونه د ازاد مورد د نيولو مخه نيسي۔
    برعکس، که د دوو ترمونو ونې له دغه ثبت سره يو شان وي، د
    دواړو ونو هرې جوړې هم‌ځايه تړونکو غوټو ته يو ګډ تازه نوم
    وټاکئ۔ دغه نومونه دې يو له بله او د دواړو ترمونو له ټولو
    نومونو بېل وي۔ نومونه له دننه نه بهر ته بدل کړئ۔ په هر ګام
    کښې تر دې غوټې دننني تړونکي لا په بېلو تازه نومونو نومول
    شوي دي؛ نو د پخواني نوم پر ځاے د نوي نوم تعويض په ټول
    بدن کښې ټاکلے دے۔ په دواړو ونو کښې بيا هماغه خام ترم
    جوړېږي، ځکه ازاد نومونه او د تړونکو ځايونه يو شان دي۔
    له دې امله دغه ثبت د الفا-معادلتوب د ټاکلو لپاره بس دے۔
\olpseganchor{OLP-0362-B034}{end}

\olpseganchor{OLP-0362-B035}{start}
    په عمومي حالت کښې د خام تعويض پر جوړښت استقرا وښيي چې
    هر ټاکلے تعويض د ونې د ثبت له مخې داسې عمل کوي۔ د ټاکلي
    متغير هر ازاد مورد د بديل ترم په يوې نسخې بدلېږي، او نور
    موردونه په خپل ځاے پاتې کېږي۔ په هرې ورګډې شوې نسخې کښې
    تړلي موردونه له خپلو دنننو تړونکو سره تړلي پاتې کېږي؛ د
    تعويض شرطونه د دې نسخې ازاد موردونه له نيول کېدو ساتي۔
    په تطبيق کښې دغه کار په دواړو برخو کښې کېږي، او په
    تجريد کښې د بدن د استقرا پايله کاروو۔ په دې توګه د اصلي
    ترم او د بديل د تړونکو ثبتونه د پايلې ثبت ټاکي۔ الفا-معادل
    استازي يو شان ثبتونه لري؛ نو د ټاکلو تعويضونو پايلې
    الفا-معادل دي۔
    د بېلګې په توګه، که بل~\LR{$N''$} او~\LR{$R''$} هماغه شرطونه پوره کړي
    او په دې حساب کښې ټول تعويضونه ټاکلي وي، نو لاندې لړۍ لرو۔
    \begin{multline*}
      \Subst{(\lambd[z][\Subst{N''}{z}{x}])}{R''}{y} \\
    \begin{aligned}
      &= \lambd[z][\Subst{\Subst{N''}{z}{x}}{R''}{y}] \\
      &\aeq \lambd[z][\Subst{\Subst{N''}{z}{x}}{R}{y}] && \text{\RL{د
        \olref{lem:sub:R} له مخې}}\\
      &\aeq\lambd[z][\Subst{\Subst{N'}{z}{x}}{R}{y}]
      && \text{\RL{د پورته ثابتې شوې يکتايۍ له مخې}}\\
      &=\Subst{(\lambd[z][\Subst{N'}{z}{x}])}{R}{y}
    \end{aligned}
    \end{multline*}
    \item که \LR{$M$} د \LR{$(PQ)$} په بڼه وي، د استقرا له مخې د دواړو برخو
    داسې استازي غوره کړئ چې تعويض پکې ټاکلے وي۔ د هغو تطبيق
    له اصلي تطبيق سره الفا-معادل دے او تعويض ئې ټاکلے دے۔
  \end{enumerate}
\end{proof}
\olpseganchor{OLP-0362-B035}{end}

\olpseganchor{OLP-0362-B036}{start}
\begin{prob}
  د \olref[lam][syn][alp]{thm:sub} د يکتايۍ برخه د خامو ترمونو
  پر جوړښت په نېغه استقرا ثابته کړئ۔
\end{prob}
\olpseganchor{OLP-0362-B036}{end}

\olpseganchor{OLP-0362-B037}{start}
\begin{cor}\ollabel{cor:sub}
  د هر~\LR{$M$}، \LR{$R$} او~\LR{$y$} دپاره داسې جوړه~\LR{$M'$} او~\LR{$R'$} شته
  چې \LR{$M' \aeq M$}، \LR{$R \aeq R'$} او \LR{$\Subst{M'}{R'}{y}$} ټاکلے
  وي۔ سربېره پر دې، که بله جوړه \LR{$M'' \aeq M$} او~\LR{$R''$}
  وي چې \LR{$R'' \aeq R$} او \LR{$\Subst{M''}{R''}{y}$} ټاکلے وي، نو
  \LR{$\Subst{M'}{R'}{y} \aeq
  \Subst{M''}{R''}{y}$} وي۔
\end{cor}
\olpseganchor{OLP-0362-B037}{end}

\olpseganchor{OLP-0362-B038}{start}
\begin{proof}
  دا سمدستي له \olref{thm:sub} څخه راځي۔
\end{proof}
\olpseganchor{OLP-0362-B038}{end}

\olpseganchor{OLP-0363-B004}{start}
\olfileid{lam}{syn}{deb}
\olpseganchor{OLP-0363-B004}{end}

\olpseganchor{OLP-0363-B005}{start}
\olsection{د دې بروين شاخص}
\olpseganchor{OLP-0363-B005}{end}

\olpseganchor{OLP-0363-B006}{start}
\LR{$\alpha$}-معادلتوب ډېر طبيعي دے، ځکه \LR{$\alpha$}-معادل ترمونه «يوه معنا»
لري۔ په حقيقت کښې د لامبډا ترمونو داسې نحو هم ورکول کېدے شي چې
هغه ترمونه سره نۀ بېلوي کوم چې د \LR{$\alpha$}-بدلون په وسيله يو بل ته
اوړي۔ تر ټولو پېژندل شوې لاره د متغيرونو پر ځاے د هغو د
\emph{دې بروين شاخص} کارول دي۔
\olpseganchor{OLP-0363-B006}{end}

\olpseganchor{OLP-0363-B007}{start}
کله چې \LR{$\lambd[x][M]$} ليکو، په څرګند ډول وايو چې \LR{$x$} د تابع
پاراميټر دے، څو په \LR{$M$} کښې د~\LR{$x$} په کارولو همدې پاراميټر ته
اشاره وکړو۔ خو د دې
بروين په شاخصي بڼه کښې پاراميټرونه نومونه نۀ لري؛ د تابع په بدنه
کښې د هغوی يادونه په داسې شمېر کېږي چې د کارېدونکي مورد او د هغه
د تړونکي تجريد تر منځ د تجريدونو شمېر ښيي۔ د بېلګې په توګه،
\LR{$\lambd[x][\lambd[y][y x]]$} ته وګورئ: دباندنے تجريد متغير~\LR{$x$}
تړي؛ دنننے تجريد متغير~\LR{$y$} تړي؛ فرعي ترم \LR{$y x$} د دننني تجريد
په ساحه کښې دے۔ د \LR{$y$} او د هغه د تړونکي \LR{$\lambd[y]$} تر منځ
هېڅ تجريد نشته، خو د \LR{$x$} او د هغه د تړونکي \LR{$\lambd[x]$} تر منځ
يو تجريد دے۔ نو د \LR{$y x$} دپاره \LR{$0\, 1$} ليکو، او د ټول ترم
دپاره \LR{$\lambd[][\lambd[][01]]$} ليکو۔ دلته وروستۍ دوه څنګ په
څنګ نښې تطبيق ښيي، يوه لسګونه شمېره نۀ ښيي۔
\olpseganchor{OLP-0363-B007}{end}

\olpseganchor{OLP-0363-B008}{start}
\begin{defn}
  د دې بروين ترمونه په استقرايي ډول داسې تعريفېږي:
  \begin{enumerate}
  \item \LR{$n$}، چې \LR{$n$} هر طبيعي عدد وي۔
  \item \LR{$PQ$}، چې \LR{$P$} او \LR{$Q$} دواړه د دې بروين ترمونه وي۔
  \item \LR{$\lambd[][N]$}، چې \LR{$N$} د دې بروين ترم وي۔
  \end{enumerate}
\end{defn}
\olpseganchor{OLP-0363-B008}{end}

\olpseganchor{OLP-0363-B009}{start}
د عادي لامبډا ترمونو د دې بروين په شاخصي ترمونو بدلون په رسمي ډول
داسې دے:
\begin{defn}
  \begin{align*}
    F_\Gamma(x) &= \Gamma(x) \\
    F_\Gamma(PQ) &= F_\Gamma(P)F_\Gamma(Q) \\
    F_\Gamma(\lambd[x][N]) &= \lambd[][F_{x,\Gamma}(N)]
  \end{align*}
  چې \LR{$\Gamma$} د متغيرونو له صفره شاخص لرونکے لړ دے، او
  \LR{$\Gamma(x)$} په \LR{$\Gamma$} کښې د متغير~\LR{$x$} ځای ښيي۔ د بېلګې په
  توګه، که \LR{$\Gamma$} د \LR{$x,y,z$} لړ وي، نو \LR{$\Gamma(x)$}، \LR{$0$}
  او \LR{$\Gamma(z)$}، \LR{$2$} دے۔
  که يو نوم په لړ کښې څو ځله راشي، له سره يې لومړے ځای اخلو۔
  همدغه ځای د تر ټولو دننني هم‌نومي تړونکي شاخص دے۔
\olpseganchor{OLP-0363-B009}{end}

\olpseganchor{OLP-0363-B010}{start}
  \LR{$x,\Gamma$} هغه لړ ښيي چې د \LR{$\Gamma$} سر ته \LR{$x$} ورزيات کړے
  شي؛ د بېلګې په توګه، د مخکنۍ بېلګې د \LR{$\Gamma$} په دوام کښې
  \LR{$w,\Gamma$} د \LR{$w,x,y,z$} لړ دے۔
\end{defn}
\olpseganchor{OLP-0363-B010}{end}

\olpseganchor{OLP-0363-B011}{start}
له دې بروين ترم څخه د معياري لامبډا ترم بېرته ترلاسه کول داسې
کېږي:
\olpseganchor{OLP-0363-B011}{end}

\olpseganchor{OLP-0363-B012}{start}
\begin{defn}
  \begin{align*}
    G_\Gamma(n) &= \Gamma[n] \\
    G_\Gamma(PQ) &= G_\Gamma(P) G_\Gamma(Q) \\
    G_\Gamma(\lambd[][N]) &= \lambd[x][G_{x,\Gamma}(N)]
  \end{align*}
  چې \LR{$\Gamma$} بيا د متغيرونو له صفره شاخص لرونکے لړ دے، او
  \LR{$\Gamma[n]$} د~\LR{$n$} په ځای کښې متغير ښيي۔ د بېلګې په توګه،
  که \LR{$\Gamma$} د \LR{$x,y,z$} لړ وي، نو \LR{$\Gamma[1]$}، \LR{$y$} دے۔
\olpseganchor{OLP-0363-B012}{end}

\olpseganchor{OLP-0363-B013}{start}
  په وروستۍ معادله کښې متغير~\LR{$x$} داسې هر متغير ټاکل کېږي چې
  په \LR{$\Gamma$} کښې نۀ وي۔
  که شاخص په اړوند لړ کښې ځای ونۀ لري، بېرته اخيستنه ناټاکلې ده۔
\end{defn}
\olpseganchor{OLP-0363-B013}{end}

\olpseganchor{OLP-0363-B014}{start}
دلته څو پايلې بې له ثبوته راوړو:
\olpseganchor{OLP-0363-B014}{end}

\olpseganchor{OLP-0363-B015}{start}
\begin{prop}
  که \LR{$M \aconvone M'$} وي، او \LR{$\Gamma$} داسې هر لړ وي چې
  \LR{$\FV{M}$} پکې شامل وي، نو \LR{$F_\Gamma(M) \eqs F_\Gamma(M')$}۔
\end{prop}
\olpseganchor{OLP-0363-B015}{end}

\olpseganchor{OLP-0364-B004}{start}
\olfileid{lam}{syn}{tr}
\olsection{ترمونه د \texorpdfstring{\LR{$\alpha$}}{α}-معادلتوب د ټولګيو په توګه}
\olpseganchor{OLP-0364-B004}{end}

\olpseganchor{OLP-0364-B005}{start}
له دې وروسته به ترمونه د \LR{$\alpha$}-معادلتوب تر کچې څېړو۔ يعنې کله
چې يو ترم ليکو، د هغه د \LR{$\alpha$}-معادلتوب ټولګے مو مراد وي۔
د بېلګې په توګه، له \LR{$\lambd[a][\lambd[b][a c]]$} څخه د ټولو هغو
ترمونو سټ مرادوو چې ورسره \LR{$\alpha$}-معادل وي؛ په دې کښې
\LR{$\lambd[a][\lambd[b][a
    c]]$} او \LR{$\lambd[b][\lambd[a][b c]]$} هم راځي، او داسې نور۔
\olpseganchor{OLP-0364-B005}{end}

\olpseganchor{OLP-0364-B006}{start}
همدارنګه، په پخوانيو برخو کښې د \LR{$N, Q$} په څېر توري د يوه ترم
دپاره کارېدل؛ له دې وروسته به هر يو د يوه ټولګي نښه وي۔
راتلونکې څېړنه به د ترمونو پر ځاے د دغو ټولګيو په باب وي۔
د \LR{$x, y$} په څېر توري بيا هم متغيرونه ښيي۔
\olpseganchor{OLP-0364-B006}{end}

\olpseganchor{OLP-0364-B007}{start}
د~\LR{$M$} د ټولګي د يوه خپل‌سري غړے د ښودلو لپاره
\LR{$\rep{M}$} ليکو۔ که له يوه زياتو ته اړتيا وي،
\LR{$\rep{M}[0], \rep{M}[1]$} او داسې نور ليکو۔
\olpseganchor{OLP-0364-B007}{end}

\olpseganchor{OLP-0364-B008}{start}
د وينا د اسانولو دپاره د ترمونو پخوانۍ نښې بيا کاروو۔ پر
ټولګيو لاندې تعريف لرو:
\begin{defn}
  \begin{enumerate}
  \item \LR{$\lambd[x][N]$} هغه ټولګے دے چې
    \LR{$\lambd[x][\rep{N}]$} پکې شامل وي۔
  \item \LR{$PQ$} هغه ټولګے دے چې \LR{$\rep{P}\rep{Q}$} پکې شامل وي۔
  \end{enumerate}
\end{defn}
\olpseganchor{OLP-0364-B008}{end}

\olpseganchor{OLP-0364-B009}{start}
دا تعريفونه د استازي له ټاکنې خپلواک دي؛ ځکه چې
\LR{$\alpha$}-بدلون له دغو جوړښتونو سره سازګار دے۔
\olpseganchor{OLP-0364-B009}{end}

\olpseganchor{OLP-0364-B010}{start}
\begin{defn} \ollabel{def:fv}
  د \LR{$\alpha$}-معادلتوب د ټولګي~\LR{$M$} \emph{ازاد متغيرونه}، يا
  \LR{$FV(M)$}، داسې تعريفېږي: \LR{$FV(\rep{M})$}۔
\end{defn}
\olpseganchor{OLP-0364-B010}{end}

\olpseganchor{OLP-0364-B011}{start}
دا تعريف د استازي له ټاکنې خپلواک دے، ځکه
\LR{$FV(\rep{M}[0]) =
FV(\rep{M}[1])$}؛ دا په \olref[alp]{thm:fv} کښې ښودل شوي دي۔
\olpseganchor{OLP-0364-B011}{end}

\olpseganchor{OLP-0364-B012}{start}
پر ټولګيو د تعويض دپاره هم پخوانۍ نښه کاروو:
\begin{defn} \ollabel{def:sub}
  په \LR{$M$} کښې د \LR{$y$} پر ځاے د \LR{$R$} \emph{تعويض}، يعنې
  \LR{$\Subst{M}{R}{y}$}، د داسې هر \LR{$\rep{M}$} او \LR{$\rep{R}$}
  دپاره چې تعويض ئې ټاکلے وي، په
  \LR{$\Subst{\rep{M}}{\rep{R}}{y}$} د الفا-معادلتوب ټولګے ټاکل کېږي۔
\end{defn}
\olpseganchor{OLP-0364-B012}{end}

\olpseganchor{OLP-0364-B013}{start}
له استازي څخه د تعريف خپلواکي په \olref[alp]{cor:sub} کښې ثابته شوې ده۔
\olpseganchor{OLP-0364-B013}{end}

\olpseganchor{OLP-0364-B014}{start}
وګورئ چې دا تعريف زمونږ استدلال څومره اسانوي۔ د بېلګې په توګه:
\begin{align}
  \Subst{\lambd[x][x]}{y}{x} & \ollabel{eq:1}\\
  &= \Subst{\lambd[z][z]}{y}{x} \ollabel{eq:2}\\
  &= \lambd[z][\Subst{z}{y}{x}] \\
  &= \lambd[z][z]
\end{align}
\olpseganchor{OLP-0364-B014}{end}

\olpseganchor{OLP-0364-B015}{start}
که دا لا هم پر خامو ترمونو تعويض وګڼو، نو \olref{eq:1}
ناټاکلے دے۔ خو لکه پورته چې وويل شول، اوس پر ټولګيو
تعويض څېړو۔ له همدې امله \olref{eq:2} روا دے:
\LR{$\lambd[x][x]$} په \LR{$\lambd[z][z]$} بدلولے شو، ځکه چې
دواړه په يوه ټولګي کښې دي۔
\olpseganchor{OLP-0364-B015}{end}

\olpseganchor{OLP-0364-B016}{start}
له همدې امله به له دې وروسته فرض کوو چې غوره شوي
استازي د تعويض اړين شرطونه پوره کوي۔ مثلاً کله چې
\LR{$\Subst{\lambd[x][N]}{R}{y}$} ووينو، فرض کوو چې
استازے \LR{$\lambd[x][N]$} داسې غوره شوے دے چې
\LR{$x \neq y$} او \LR{$x \notin FV(R)$} وي۔
\olpseganchor{OLP-0364-B016}{end}

\olpseganchor{OLP-0364-B017}{start}
څرنګه چې \LR{$\lambd[x][x]$} ته «ټولګے» ويل لږ نااشنا ښکاري،
له دې وروسته به دې ټولګيو ته \LR{$\Lambda$}-ترمونه وايو
(يا د دې برخې په پاتې متن کښې يوازې «ترمونه»)، څو
له هغو \LR{$\lambda$}-ترمونو سره ئې بېل وساتو چې له پخوا
ورسره اشنا يو۔
\olpseganchor{OLP-0364-B017}{end}

\olpseganchor{OLP-0364-B018}{start}
\begin{editorial}
  لا له خامو ترمونو سره مخه نۀ شو ښۀ کولے: د
  \LR{$\Lambda$}-ترمونو ټول تعريف پر \LR{$\lambda$}-ترمونو ولاړ دے،
  او تر اوسه مو پر \LR{$\Lambda$}-ترمونو د تابعګانو د تعريف
  طريقه نۀ ده ورکړې۔ له همدې امله دا ډول تابعګانې
  لومړے پر \LR{$\lambda$}-ترمونو تعريفوو او بيا ئې
  \LR{$\Lambda$}-ترمونو ته «لېږدوو»، لکه د تعويض په حالت کښې۔
  خو فرض کوو چې لوستونکے په بديهي ډول پوهېږي چې پر
  \LR{$\Lambda$}-ترمونو تابعګانې څنګه تعريفېدے شي۔
  پر ټولګيو دغه لېږد هغه وخت ښه تعريف لري چې د پايلې ټولګے
  د استازي له ټاکنې خپلواک وي۔
\end{editorial}
\olpseganchor{OLP-0364-B018}{end}

\olpseganchor{OLP-0365-B004}{start}
\olfileid{lam}{int}{bet}
\olsection{\texorpdfstring{\LR{$\beta$}}{β}-راکمول}
\olpseganchor{OLP-0365-B004}{end}

\olpseganchor{OLP-0365-B005}{start}
کله چې \LR{$(\lambd[m][(\lambd[y][y]) m])$} ووينو، طبيعي ده چې
ګومان وکړو له \LR{$\lambd[m][m]$} سره څه اړيکه لري؛ يعنې دويم ترم
بايد د لومړي د «اسانولو» پايله وي۔ د \emph{\LR{$\beta$}-راکمول}
مفهوم دا انګېرنه په رسمي ډول رانغاړي۔
\olpseganchor{OLP-0365-B005}{end}

\olpseganchor{OLP-0365-B006}{start}
\begin{defn}[\LR{$\beta$}-انقباض، \LR{$\bredone$}] \ollabel{defn:betacontr}
  \emph{\LR{$\beta$}-انقباض} (\LR{$\bredone$}) پر ترمونو تر ټولو کوچنۍ
  سازګاره اړيکه ده چې لاندې شرط پوره کوي:
  \[
    (\lambd[x][N])Q \bredone \Subst{N}{Q}{x} 
  \]
  که \LR{$P \bredone Q$} وي، وايو چې \LR{$P$}، \LR{$Q$} ته
  \emph{\LR{$\beta$}-منقبض} شوے دے۔ د \LR{$(\lambd[x][N])Q$}
  په بڼه ترم ته \emph{رېډکس} وايي۔
\end{defn}
\olpseganchor{OLP-0365-B006}{end}

\olpseganchor{OLP-0365-B007}{start}
\begin{prob} \ollabel{prob:def}
  د \LR{$\beta$}-انقباض معادل استقرايي تعريفونه په څرګنده وليکئ،
  لکه څنګه چې مو د تړلي متغير د نوم بدلون دپاره په
  \olref[lam][syn][alp]{defn:aconvone} کښې وکړل۔
\end{prob}
\olpseganchor{OLP-0365-B007}{end}
  
\olpseganchor{OLP-0365-B008}{start}
\begin{defn}[\LR{$\beta$}-راکمول، \LR{$\bred$}] \ollabel{defn:betared}
  \emph{\LR{$\beta$}-راکمول} (\LR{$\bred$}) پر ترمونو تر ټولو کوچنۍ
  انعکاسي او انتقالي اړيکه ده چې \LR{$\bredone$} پکې شامله وي۔
  که \LR{$P \bred Q$} وي، وايو چې \LR{$P$}، \LR{$Q$} ته
  \LR{$\beta$}-راکم شوے دے۔
\end{defn}
\olpseganchor{OLP-0365-B008}{end}

\olpseganchor{OLP-0365-B009}{start}
کله چې مطلب له متنه څرګند وي، د \LR{$\bredone$} پر ځاے
\LR{$\redone$} او د \LR{$\bred$} پر ځاے \LR{$\red$} ليکو۔
\olpseganchor{OLP-0365-B009}{end}

\olpseganchor{OLP-0365-B010}{start}
په غيررسمي وينا، \LR{$M \bred N$} هله او يوازې هله وي چې
\LR{$M$} د \LR{$\beta$}-انقباض په صفر يا څو ګامونو \LR{$N$} ته
بدلېدے شي۔
\olpseganchor{OLP-0365-B010}{end}

\olpseganchor{OLP-0365-B011}{start}
\begin{defn}[\LR{$\beta$}-عادي]
هغه ترم چې نور پرې \LR{$\beta$}-انقباض نۀ شي کېدے،
\emph{\LR{$\beta$}-عادي} بلل کېږي۔
\end{defn}
\olpseganchor{OLP-0365-B011}{end}

\olpseganchor{OLP-0365-B012}{start}
که \LR{$M \bred N$} او \LR{$N$}، \LR{$\beta$}-عادي وي، نو \LR{$N$}
د~\LR{$M$} \emph{عادي بڼه} بولو۔ هر ترم عادي بڼې ته نۀ رسېږي۔ که د يوه ترم عادي بڼه شته،
نو يکتا ده، لکه وروسته چې به ئې ووينو۔
\olpseganchor{OLP-0365-B012}{end}

\olpseganchor{OLP-0365-B013}{start}
راځئ څو بېلګې وګورو۔
\begin{enumerate}
\item لرو:
\begin{align*}
(\lambd[x][xxy]) \lambd[z][z] & \redone (\lambd[z][z])(\lambd[z][z]) y \\
& \redone (\lambd[z][z]) y \\
& \redone y
\end{align*}
\item د يوه ترم «اسانول» په حقيقت کښې هغه لا پېچلے کولے شي:
\begin{align*}
(\lambd[x][xxy])(\lambd[x][xxy]) & \redone (\lambd[x][xxy])(\lambd[x][xxy])y \\
& \redone (\lambd[x][xxy])(\lambd[x][xxy])yy \\
& \redone \dots
\end{align*}
\item دا کار ترم هماغسې بې بدله هم پرېښودے شي:
\[
(\lambd[x][xx])(\lambd[x][xx]) \redone (\lambd[x][xx])(\lambd[x][xx])
\]
\item ځينې ترمونه په يوه نه، بلکې په څو لارو هم راکمېږي؛
  د بېلګې په توګه،
\[
(\lambd[x][(\lambd[y][yx]) z]) v \redone (\lambd[y][yv]) z
\]
  د تر ټولو دباندني تطبيق د انقباض له لارې؛ او
\[
(\lambd[x][(\lambd[y][yx]) z]) v \redone (\lambd[x][zx]) v
\]
  د تر ټولو دننني تطبيق د انقباض له لارې۔ خو پام وکړئ
  چې په دې حالت کښې دواړه پايلې بيا هماغه يوه ترم،
  \LR{$zv$}، ته راکمېږي۔
\end{enumerate}
\olpseganchor{OLP-0365-B013}{end}

\olpseganchor{OLP-0365-B014}{start}
د وروستۍ بېلګې ګډه پايله تصادف نۀ دے؛ دا د لامبډا
حساب يو ژور او مهم خاصيت ښيي چې د چرچ--روسر خاصيت
بلل کېږي۔
\olpseganchor{OLP-0365-B014}{end}

\olpseganchor{OLP-0365-B015}{start}
\begin{digress}
  په عمومي ډول يو ترم د \LR{$\beta$}-راکمول څو لارې لرلے
  شي؛ له همدې امله د راکمول بېلابېلې تګلارې جوړې
  شوې دي۔ په هغو کښې تر ټولو عامه \emph{طبيعي تګلاره}
  ده۔ طبيعي تګلاره تل هغه \emph{تر ټولو کيڼ} رېډکس
  منقبضوي چې ځای ئې د ترم په ليکلي بڼه کښې د هغه
  د پيل له مخې ټاکل کېږي۔ د دې تګلارې ګټور خاصيت
  دا دے: که يو ترم په کومې تګلارې عادي بڼې ته
  راکمېدے شي، نو په طبيعي تګلارې هم هغې ته
  راکمېدے شي۔ له دې وروسته به، تر څو بل ډول
  څرګند نۀ شي، همدا طبيعي تګلاره کاروو۔
\end{digress}
\olpseganchor{OLP-0365-B015}{end}

\olpseganchor{OLP-0365-B016}{start}
\begin{defn}[\LR{$\beta$}-معادلتوب، \LR{$\equal$}]
  \emph{\LR{$\beta$}-معادلتوب} (\LR{$\equal$}) هغه اړيکه ده چې
  په استقرايي ډول داسې تعريفېږي:
  \begin{enumerate}
  \item \LR{$M \equal M$}.
  \item که \LR{$M \equal N$} وي، نو \LR{$N \equal M$}.
  \item که \LR{$M \equal N$} او \LR{$N \equal O$} وي، نو \LR{$M \equal O$}.
  \item که \LR{$M \equal N$} وي، نو \LR{$PM \equal PN$}.
  \item که \LR{$M \equal N$} وي، نو \LR{$MQ \equal NQ$}.
  \item که \LR{$M \equal N$} وي، نو \LR{$\lambd[x][M] \equal \lambd[x][N]$}.
  \item \LR{$(\lambd[x][N])Q \equal \Subst{N}{Q}{x}$}.
  \end{enumerate}
\end{defn}
\olpseganchor{OLP-0365-B016}{end}

\olpseganchor{OLP-0365-B017}{start}
لومړۍ درې قاعدې دا اړيکه د معادلتوب اړيکه ګرځوي؛
راتلونکې درې ئې سازګاره کوي؛ او وروستۍ قاعده
\LR{$\beta$}-انقباض پکې شاملوي۔
\olpseganchor{OLP-0365-B017}{end}

\olpseganchor{OLP-0365-B018}{start}
په غيررسمي وينا، دوه ترمونه هله او يوازې هله
\LR{$\beta$}-معادل دي چې يو ئې د \LR{$\beta$}-انقباض يا د
معکوس \LR{$\beta$}-انقباض په صفر يا څو ګامونو بل ته بدل شي۔
د \LR{$\beta$}-انقباض معکوس داسې تعريفېږي چې \LR{$M$}،
\LR{$N$} ته په معکوس \LR{$\beta$}-انقباض هله او يوازې هله
اوړي چې \LR{$N$}، \LR{$M$} ته \LR{$\beta$}-منقبضېږي۔
\olpseganchor{OLP-0365-B018}{end}

\olpseganchor{OLP-0365-B019}{start}
له پورته قواعدو سربېره به دې اړيکې ته نور قواعد
هم ورزياتوو او غځېدلې معادلتوبي اړيکه به
\LR{$\equal[X]$} ليکو؛ دلته \LR{$X$} هغه ورزياته شوې قاعده ده۔
\olpseganchor{OLP-0365-B019}{end}

\olpseganchor{OLP-0366-B004}{start}
\olfileid{lam}{syn}{eta}
\olsection{\texorpdfstring{\LR{$\eta$}}{η}-بدلون}
\olpseganchor{OLP-0366-B004}{end}

\olpseganchor{OLP-0366-B005}{start}
د \LR{$\lambda$} ترمونو ترمنځ يوه بله اړيکه هم شته۔ په \olref[fv]{sec} کښې مو
د \LR{$\lambd[x][(fx)]$} بېلګه کارولې وه: دا يو آرګومېنټ اخلي او \LR{$f$} پرې
تطبيقوي۔ په بله وينا، دا هماغه تابع ده چې~\LR{$f$} ده: \LR{$\lambd[x][(fx)]N$}
او \LR{$fN$} دواړه~\LR{$fN$} ته راکمېږي۔ د همدې مفکورې د څرګندولو دپاره
\LR{$\eta$}-راکمول او \LR{$\eta$}-پراختيا کاروو۔ د برابروالي په دې بېلګه کښې
د تجريد تړلے متغير د تابع په ازادو متغيرونو کښې نۀ وي؛ لاندې تعريف
دا شرط په څرګنده ټاکي۔
\olpseganchor{OLP-0366-B005}{end}

\olpseganchor{OLP-0366-B006}{start}
\begin{defn}[\LR{$\eta$}-انقباض، \LR{$\eredone$}]
  \ollabel{defn:beredone}
  \emph{\LR{$\eta$}-انقباض} (\LR{$\eredone$}) پر ترمونو تر ټولو کوچنۍ سازګاره
  اړيکه ده چې لاندې شرط پوره کوي:
  \[
    \lambd[x][M x] \eredone M \text{\RL{ که }} x \notin FV(M)
  \]
\end{defn}
\olpseganchor{OLP-0366-B006}{end}

\olpseganchor{OLP-0366-B007}{start}
\begin{defn}[\LR{$\beta\eta$}-راکمول، \LR{$\bered$}] \ollabel{defn:bered}
  \emph{\LR{$\beta\eta$}-راکمول} (\LR{$\bered$}) پر ترمونو تر ټولو کوچنۍ
  انعکاسي او انتقالي اړيکه ده چې \LR{$\bredone$} او \LR{$\eredone$} پکې
  شاملې دي؛ يعنې د انعکاس او انتقال د قواعدو ترڅنګ لاندې دوه
  قاعدې هم لري:
  \begin{enumerate}
  \item که \LR{$M \bredone N$} وي، نو \LR{$M \bered N$}۔ \ollabel{defn:bered3}
  \item که \LR{$M \eredone N$} وي، نو \LR{$M \bered N$}۔ \ollabel{defn:bered4}
  \end{enumerate}
\olpseganchor{OLP-0366-B007}{end}

\olpseganchor{OLP-0366-B008}{start}
\end{defn}
\olpseganchor{OLP-0366-B008}{end}

\olpseganchor{OLP-0366-B009}{start}
\begin{defn}
  د \LR{$\equal$} معادلتوبي اړيکه د \LR{$\eta$}-بدلون په دې قاعده غځوو:
  \[
  \lambd[x][f x] \equal f
  \]
  په دې قاعده کښې هم تجريد يوازې هغه وخت لرې کېږي
  چې تړلے متغير د ښي‌لاس ترم په ازادو متغيرونو کښې
  نۀ وي۔ په دې شرط له تجريد څخه د تابع د ازادو متغيرونو سټ نۀ بدلېږي۔
  غځېدلې اړيکه \LR{$\equal[\eta]$} ليکو۔
\end{defn}
\olpseganchor{OLP-0366-B009}{end}

\olpseganchor{OLP-0366-B010}{start}
\LR{$\eta$}-معادلتوب ځکه مهم دے چې د لامبډا ترمونو د قيمتونو له مخې
برابروالي سره تړاو لري:
\olpseganchor{OLP-0366-B010}{end}

\olpseganchor{OLP-0366-B011}{start}
\begin{defn}[د قيمتونو له مخې برابري]
  د \LR{$\equal$} معادلتوبي اړيکه د (\ext) قاعدې په زياتولو غځوو:
  \begin{center}
  که \LR{$Mx \equal Nx$} وي، نو \LR{$M \equal N$}، په دې شرط چې \LR{$x \notin FV(MN)$} وي۔
  \end{center}
  او غځېدلې اړيکه \LR{$\equal[\ext]$} ليکو۔
\end{defn}
\olpseganchor{OLP-0366-B011}{end}

\olpseganchor{OLP-0366-B012}{start}
په لنډه وينا، قاعده وايي چې که دوه ترمونه د تابعو په توګه د يوه او
هماغه آرګومېنټ دپاره يو شان چلند کوي، نو بايد برابر وګڼل شي۔
\olpseganchor{OLP-0366-B012}{end}

\olpseganchor{OLP-0366-B013}{start}
اوس ثابتوو چې د \LR{$\eta$} قاعده د قيمتونو له مخې برابروالي هماغه
اندازه برابروي او تر هغې زيات څۀ نۀ ورزياتوي۔
\olpseganchor{OLP-0366-B013}{end}

\olpseganchor{OLP-0366-B014}{start}
\begin{thm}
  \LR{$M \equal[\ext] N$} هله او يوازې هله چې \LR{$M \equal[\eta] N$}۔
\end{thm}
\olpseganchor{OLP-0366-B014}{end}

\olpseganchor{OLP-0366-B015}{start}
\begin{proof}
  لومړے ثابتوو چې \LR{$\equal[\eta]$} د قيمتونو له مخې برابروالي د
  قاعدې په وړاندې تړلې ده؛ يعنې د \LR{$ext$} قاعده په \LR{$\equal[\eta]$}
  کښې څۀ نۀ زياتوي۔ له دې امله \LR{$\equal[\eta]$}، \LR{$\equal[\ext]$}
  رانغاړي؛ نو که \LR{$M \equal[\ext] N$} وي، \LR{$M \equal[\eta] N$} هم وي۔
\olpseganchor{OLP-0366-B015}{end}

\olpseganchor{OLP-0366-B016}{start}
  د دې د ثابتولو دپاره چې \LR{$\equal[\eta]$} د \ext{} قاعدې په وړاندې
  تړلې ده، پام وکړئ: د هر هغه \LR{$M \equal N$} دپاره چې د \ext{} قاعدې
  له لارې ترلاسه شي، \LR{$Mx \equal[\eta] Nx$} د هغې مقدمه ده۔ بيا د
  \LR{$\equal$} د يوې قاعدې له مخې \LR{$\lambd[x][Mx] \equal[\eta]
  \lambd[x][Nx]$} لرو؛ دواړو خواوو ته د \LR{$\eta$} په کارولو سره
  \LR{$M \equal[\eta] N$} ترلاسه کوو۔
\olpseganchor{OLP-0366-B016}{end}

\olpseganchor{OLP-0366-B017}{start}
  په همدې ډول ثابتوو چې د \LR{$\eta$} قاعده په \LR{$\equal[ext]$} کښې
  شامله ده۔ د هر \LR{$\lambd[x][Mx]$} او \LR{$M$} دپاره چې \LR{$x \notin
  FV(M)$} وي، لرو چې \LR{$(\lambd[x][Mx])x \equal[\ext] Mx$}؛
  له دې څخه د \ext{} قاعدې له مخې \LR{$\lambd[x][Mx] \equal[\ext] M$}
  ترلاسه کېږي۔
\end{proof}
\olpseganchor{OLP-0366-B017}{end}

\olpseganchor{OLP-0367-B004}{start}
\olchapter{lam}{cr}{د چرچ--روسر خاصيت}
\olpseganchor{OLP-0367-B004}{end}

\olpseganchor{OLP-0368-B004}{start}
\olfileid{lam}{cr}{dap}
\olpseganchor{OLP-0368-B004}{end}

\olpseganchor{OLP-0368-B005}{start}
\olsection{تعريف او خاصيتونه}
\olpseganchor{OLP-0368-B005}{end}

\olpseganchor{OLP-0368-B006}{start}
په دې څپرکي کښې د چرچ--روسر خاصيت او د هغه ځينې عامې
ځانګړنې معرفي کوو۔ په لاندې تعريف کښې د اړيکې هر غشے د هماغې ټاکلې اړيکې يو ګام دے۔
\olpseganchor{OLP-0368-B006}{end}

\olpseganchor{OLP-0368-B007}{start}
\begin{defn}[د چرچ--روسر خاصيت، CR] پر ترمونو يوه اړيکه
  \LR{$\xredone$} هله او يوازې هله د \emph{چرچ--روسر خاصيت} لري
  چې کله هم \LR{$M \xredone P$} او \LR{$M \xredone Q$} وي، داسې يو~\LR{$N$}
  موجود وي چې \LR{$P \xredone N$} او \LR{$Q \xredone N$} وي۔
\end{defn}
\olpseganchor{OLP-0368-B007}{end}

\olpseganchor{OLP-0368-B008}{start}
لامبډا حساب د محاسبې د يوه ماډل په توګه ليدلے شو: عادي بڼې
لرونکي ترمونه يې «ارزښتونه» دي او پر ترمونو د راکمېدو اړيکه يې
«د محاسبې قاعدې» دي۔ د چرچ--روسر خاصيت وايي چې که محاسبه په
څو لارو پرمخ تللے شي، نو د هغې عادي پايله، که موجوده وي، يوازې
يوه وي۔
\olpseganchor{OLP-0368-B008}{end}

\olpseganchor{OLP-0368-B009}{start}
د ابتدايي الجبره د بېلګې په توګه، \LR{$4 \times (1+2) + 3$} په څو
لارو حسابېدے شي۔ يا ئې \LR{$4 \times 3+3$} ته راکمولاے شو، که
لومړے \LR{$1+2$}،~\LR{$3$} ته راکم کړو؛ يا ئې \LR{$4 \times 1+4
\times 2+3$} ته راکمولاے شو، که لومړے پر \LR{$4 \times (1+2)$}
د وېش قانون وکاروو۔ خو دواړه لارې بيا \LR{$12+3$} ته راکمېږي۔
\olpseganchor{OLP-0368-B009}{end}

\olpseganchor{OLP-0368-B010}{start}
که \LR{$\xredone$} د \LR{$\beta$}-راکمول اړيکه ونيسو، په اسانه وينو
چې د چرچ--روسر د خاصيت يوه پايله دا ده: که يو ترم عادي بڼه
ولري، هغه يکتا ده۔ فرض کړئ \LR{$M$}، \LR{$P$} او \LR{$Q$} ته راکمېږي
او دواړه عادي بڼې دي۔ د چرچ--روسر د خاصيت له مخې داسې يو
\LR{$N$} شته چې \LR{$P$} او \LR{$Q$} دواړه ورته راکمېږي۔ څرنګه چې د فرض
له مخې \LR{$P$} او \LR{$Q$} عادي بڼې دي، نو له \LR{$P$} او \LR{$Q$} څخه \LR{$N$}
ته راکمول يوازې بې‌بدلونه راکمول کېدے شي؛ يعنې \LR{$P$}، \LR{$Q$}
او \LR{$N$} يو شان دي۔ ځکه د يوه ترم د \emph{هماغې} عادي بڼې
په اړه خبرې کول سم دي۔
\olpseganchor{OLP-0368-B010}{end}

\olpseganchor{OLP-0368-B011}{start}
که لامبډا حساب د محاسبې ماډل وګڼو، د يوه ترم عادي بڼه د
هغه ترم له محاسبې څخه د راوتلې «وروستۍ پايلې» په توګه ليدلے
شو۔ د پورته پايلې معنا دا ده چې که وروستۍ پايله وي، نو يوازې
يوه وي؛ لکه د \LR{$4 \times (1+2)+3$} د حساب پايله چې~\LR{$15$} ده۔
\olpseganchor{OLP-0368-B011}{end}

\olpseganchor{OLP-0368-B012}{start}
\begin{thm} \ollabel{thm:str}
  که يوه اړيکه \LR{$\xredone$} د چرچ--روسر خاصيت ولري او \LR{$\xred$}
  تر ټولو کوچنۍ انتقالي اړيکه وي چې \LR{$\xredone$} پکې شامله ده،
  نو \LR{$\xred$} هم د چرچ--روسر خاصيت لري۔
\end{thm}
\olpseganchor{OLP-0368-B012}{end}

\olpseganchor{OLP-0368-B013}{start}
\begin{proof}
  فرض کړئ
  \begin{align*}
    M & \xredone P_1 \xredone \dots \xredone P_m \text{\RL{ او}}\\
    M & \xredone Q_1 \xredone \dots \xredone Q_n.
  \end{align*}
  د دواړو لړيو وروستي ترمونه د جال د دوو پولو پایونه دي۔ قضيه د ترمونو د يوه \LR{$N$} جال په جوړولو ثابتوو چې لوړوالے يې
  \LR{$m + 1$} او پلنوالے يې \LR{$n + 1$} دے۔ د \LR{$i$}-مې ليکې او
  \LR{$j$}-مې ستنې ترم په \LR{$N_{i,j}$} ښيو۔
\olpseganchor{OLP-0368-B013}{end}

\olpseganchor{OLP-0368-B014}{start}
  \LR{$N$} داسې جوړوو چې \LR{$N_{i,j} \xredone N_{i+1,j}$} او
  \LR{$N_{i,j} \xredone N_{i,j+1}$} وي۔ تعريف يې دا دے:
  \begin{align*}
    N_{0,0} &= M \\
    N_{i,0} &= P_i && \text{\RL{که }} 1 \le i \le m \\
    N_{0,j} &= Q_j && \text{\RL{که }} 1 \le j \le n \\
  \intertext{\RL{او په نورو حالاتو کښې:}}
    N_{i,j} &= R
  \end{align*}
  دلته \LR{$R$} داسې ترم دے چې \LR{$N_{i-1,j} \xredone R$} او
  \LR{$N_{i,j-1} \xredone R$} وي۔ د \LR{$\xredone$} د چرچ--روسر د
  خاصيت له مخې داسې ترم تل موجود وي، ځکه د دواړو مخکينيو ټکو ګډ پيل د جال د مخکنۍ ليکې او ستنې په تقاطع کښې دے۔
\olpseganchor{OLP-0368-B014}{end}

\olpseganchor{OLP-0368-B015}{start}
  اوس \LR{$N_{m,0} \xredone \dots \xredone N_{m,n}$} او \LR{$N_{0,n}
  \xredone \dots \xredone N_{m,n}$} لرو۔ پام وکړئ چې \LR{$N_{m,0}$}
  هماغه \LR{$P$} او \LR{$N_{0,n}$} هماغه \LR{$Q$} دے۔ د \LR{$\xred$} له تعريفه
  قضيه ثابتېږي۔
\end{proof}
\olpseganchor{OLP-0368-B015}{end}

\olpseganchor{OLP-0369-B004}{start}
\olfileid{lam}{cr}{pb}
\olpseganchor{OLP-0369-B004}{end}

\olpseganchor{OLP-0369-B005}{start}
\olsection{هممهاله \texorpdfstring{\LR{$\beta$}}{β}-راکمول}
\olpseganchor{OLP-0369-B005}{end}

\olpseganchor{OLP-0369-B006}{start}
د \emph{هممهاله \LR{$\beta$}-راکمول} مفهوم معرفي کوو او ثابتوو
چې دا د چرچ--روسر خاصيت لري۔
\olpseganchor{OLP-0369-B006}{end}

\olpseganchor{OLP-0369-B007}{start}
\begin{defn}[هممهاله \LR{$\beta$}-راکمول، \LR{$\bredpar$}] \ollabel{defn:bredpar}
  د ترمونو هممهاله راکمول (\LR{$\bredpar$}) په استقرايي ډول داسې تعريفېږي:
  \begin{enumerate}
    \item \ollabel{defn:bredpar1} \LR{$x \bredpar x$}.
    \item \ollabel{defn:bredpar2} که \LR{$N \bredpar N'$} وي، نو \LR{$\lambd[x][N] \bredpar
      \lambd[x][N']$}.
    \item \ollabel{defn:bredpar3} که \LR{$P \bredpar P'$} او \LR{$Q \bredpar Q'$} وي، نو \LR{$PQ \bredpar
      P'Q'$}.
    \item \ollabel{defn:bredpar4} که \LR{$N \bredpar N'$} او \LR{$Q \bredpar Q'$} وي، نو
      \LR{$(\lambd[x][N])Q \bredpar \Subst{N'}{Q'}{x}$}.
  \end{enumerate}
\end{defn}
\olpseganchor{OLP-0369-B007}{end}

\olpseganchor{OLP-0369-B008}{start}
هممهاله \LR{$\beta$}-راکمول موږ ته اجازه راکوي چې د يوه ترم هر
څو رېډکسونه په يوه ګام کښې راکم کړو۔ دا له \LR{$\beta$}-راکمول سره
په دې توپير لري چې يوازې په اصلي ترم کښې شته رېډکسونه منقبضولاے
شو، نۀ هغه رېډکسونه چې د هممهاله \LR{$\beta$}-راکمول په پايله کښې
پيدا کېږي۔ د بېلګې په توګه، \LR{$(\lambd[f][fx])(\lambd[y][y])$}
په يوه هممهاله \LR{$\beta$}-ګام کښې يا هماغه خپل ځان ته، يا
\LR{$(\lambd[y][y])x$} ته راکمېږي، خو له هغه ځايه لا~\LR{$x$} ته نۀ؛
حال دا چې د عادي \LR{$\beta$}-راکمول له لارې~\LR{$x$} ته رسېږي۔
سبب يې دا دے چې وروستے رېډکس د هممهاله \LR{$\beta$}-راکمول
له يوه ګام وروسته پيدا کېږي۔ دويم هممهاله \LR{$\beta$}-ګام
بيا~\LR{$x$} ورکوي۔
\olpseganchor{OLP-0369-B008}{end}

\olpseganchor{OLP-0369-B009}{start}
\begin{thm}\ollabel{thm:refl}
  \LR{$M \bredpar M$}.
\end{thm}
\olpseganchor{OLP-0369-B009}{end}

\olpseganchor{OLP-0369-B010}{start}
\begin{proof}
  تمرين۔
\end{proof}
\olpseganchor{OLP-0369-B010}{end}

\olpseganchor{OLP-0369-B011}{start}
\begin{prob}
  \olref[lam][cr][pb]{thm:refl} ثابت کړئ۔
\end{prob}
\olpseganchor{OLP-0369-B011}{end}

\olpseganchor{OLP-0369-B012}{start}
\begin{defn}[\LR{$\beta$}-بشپړ راکمول]\ollabel{defn:bcd}
  د \LR{$M$} \emph{\LR{$\beta$}-بشپړ راکمول}~\LR{$\bcd{M}$} په استقرايي ډول
  داسې تعريفېږي:
  \begin{align}
    \bcd{x} &= x \ollabel{defn:bcd1} \\
    \bcd{(\lambd[x][N])} &= \lambd[x][\bcd{N}] \ollabel{defn:bcd2}\\
    \bcd{(PQ)} &= \bcd{P}\bcd{Q} && \text{\RL{که \LR{$P$} د \LR{$\lambd$} تجريد نۀ وي}} 
    \ollabel{defn:bcd3} \\
    \bcd{((\lambd[x][N])Q)} &= \Subst{\bcd{N}}{\bcd{Q}}{x} \ollabel{defn:bcd4}
  \end{align}
\end{defn}
\olpseganchor{OLP-0369-B012}{end}

\olpseganchor{OLP-0369-B013}{start}
د ترم \LR{$\beta$}-بشپړ راکمول، لکه له نامه ئې چې ښکاري، يو
«بشپړ هممهاله راکمول» دے۔ په هممهاله \LR{$\beta$}-راکمول کښې
لا هم دا ټاکلے شو چې رېډکس منقبض نۀ کړو؛ خو په بشپړ
راکمول کښې بايد ټول موجود رېډکسونه منقبض کړو۔ نو د
\LR{$(\lambd[f][fx])(\lambd[y][y])$} بشپړ راکمول
\LR{$(\lambd[y][y])x$} دے، نۀ هماغه لومړنے ترم۔
\olpseganchor{OLP-0369-B013}{end}

\olpseganchor{OLP-0369-B014}{start}
\begin{editorial}
  په دې تعريف کښې دا ستونزه ده چې تر اوسه مو پر (\LR{$\lambd$}-)ترمونو
  د تابعو د بازګشتي تعريف لاره نۀ ده معرفي کړې۔ وروسته به ئې
  سم کړو۔
\end{editorial}
\olpseganchor{OLP-0369-B014}{end}

\olpseganchor{OLP-0369-B015}{start}
\begin{lem}\ollabel{lem:comp}
  که \LR{$M \bredpar M'$} او \LR{$R \bredpar R'$} وي، نو \LR{$\Subst{M}{R}{y}
  \bredpar \Subst{M'}{R'}{y}$}.
\end{lem}
\olpseganchor{OLP-0369-B015}{end}

\olpseganchor{OLP-0369-B016}{start}
\begin{proof}
  د \LR{$M \bredpar M'$} د اشتقاق پر جوړښت استقرا کوو۔ د الفا-بدلون په مرسته هر تړلے نوم د تعويض له متغير او د دواړو تعويضېدونکو ترمونو له ازادو نومونو تازه ټاکو۔
  \begin{enumerate}
    \item وروستے ګام \olref{defn:bredpar1} دے: تمرين۔
    \item وروستے ګام \olref{defn:bredpar2} دے: په دې حالت کښې \LR{$M$}،
      \LR{$\lambd[x][N]$} او \LR{$M'$}، \LR{$\lambd[x][N']$} دے، او
      \LR{$N \bredpar N'$} وي۔ غواړو ثابت کړو چې
      \LR{$\Subst{(\lambd[x][N])}{R}{y} \bredpar
      \Subst{(\lambd[x][N'])}{R'}{y}$}، يعنې
      \LR{$\lambd[x][\Subst{N}{R}{y}] \bredpar
      \lambd[x][\Subst{N'}{R'}{y}]$}۔ دا د
      \olref{defn:bredpar2} او د استقرا له فرضيې سملاسي راځي۔
    \item وروستے ګام \olref{defn:bredpar3} دے: تمرين۔
    \item وروستے ګام \olref{defn:bredpar4} دے: \LR{$M$}،
      \LR{$(\lambd[x][N])Q$} او \LR{$M'$}، \LR{$\Subst{N'}{Q'}{x}$} دے۔ غواړو
      ثابت کړو چې \LR{$\Subst{((\lambd[x][N])Q)}{R}{y} \bredpar
      \Subst{\Subst{N'}{Q'}{x}}{R'}{y}$}، يعنې
      \LR{$(\lambd[x][\Subst{N}{R}{y}])\Subst{Q}{R}{y} \bredpar
      \Subst{\Subst{N'}{R'}{y}}{\Subst{Q'}{R'}{y}}{x}$}۔ دا د
      \olref{defn:bredpar4} او د استقرا له فرضيې راځي۔
  \end{enumerate}
\end{proof}
\olpseganchor{OLP-0369-B016}{end}

\olpseganchor{OLP-0369-B017}{start}
\begin{prob}
  د \olref[lam][cr][pb]{lem:comp} ثبوت بشپړ کړئ۔
\end{prob}
\olpseganchor{OLP-0369-B017}{end}

\olpseganchor{OLP-0369-B018}{start}
\begin{lem}\ollabel{lem:cont}
  که \LR{$M \bredpar M'$} وي، نو \LR{$M' \bredpar \bcd{M}$}.
\end{lem}
\begin{proof}
  د \LR{$M \bredpar M'$} د اشتقاق پر جوړښت استقرا کوو۔
  \begin{enumerate}
    \item وروستۍ قاعده \olref{defn:bredpar1} ده: تمرين۔
    \item وروستۍ قاعده \olref{defn:bredpar2} ده: \LR{$M$}،
      \LR{$\lambd[x][N]$} او \LR{$M'$}، \LR{$\lambd[x][N']$} دے، او
      \LR{$N \bredpar N'$} وي۔ غواړو وښيو چې \LR{$\lambd[x][N'] \bredpar
      \bcd{(\lambd[x][N])}$}، يعنې د \olref{defn:bcd2} له مخې
      \LR{$\lambd[x][N'] \bredpar \lambd[x][\bcd{N}]$}۔ دا د
      \olref{defn:bredpar2} او د استقرا له فرضيې راځي۔
    \item وروستۍ قاعده \olref{defn:bredpar3} ده: \LR{$M$}، \LR{$PQ$} او \LR{$M'$}
      د ځينو \LR{$P$}، \LR{$Q$}، \LR{$P'$} او \LR{$Q'$} دپاره \LR{$P'Q'$} دے، چې \LR{$P \bredpar P'$}
      او \LR{$Q \bredpar Q'$} وي۔ د استقرا له فرضيې \LR{$P'
      \bredpar \bcd{P}$} او \LR{$Q' \bredpar \bcd{Q}$} لرو۔
      \begin{enumerate}
        \item که \LR{$P$} د ځينو \LR{$x$} او \LR{$N$} دپاره \LR{$\lambd[x][N]$} وي،
          نو \LR{$P'$} بايد د کوم \LR{$N'$} دپاره \LR{$\lambd[x][N']$} وي، او
          \LR{$N \bredpar N'$} وي۔ د استقرا له فرضيې \LR{$N' \bredpar \bcd{N}$} او
          \LR{$Q' \bredpar \bcd{Q}$} لرو۔ نو \LR{$(\lambd[x][N'])Q' \bredpar
          \Subst{\bcd{N}}{\bcd{Q}}{x}$} د \olref{defn:bredpar4} له مخې۔
        \item که \LR{$P$} د \LR{$\lambd$} تجريد نۀ وي، نو \LR{$P'Q' \bredpar
          \bcd{P}\bcd{Q}$} د \olref{defn:bredpar3} له مخې؛ ښۍ خوا
          د \olref{defn:bcd3} له مخې \LR{$\bcd{PQ}$} ده۔
      \end{enumerate}
    \item وروستۍ قاعده \olref{defn:bredpar4} ده: \LR{$M$}،
      \LR{$(\lambd[x][N])Q$} او \LR{$M'$}، \LR{$\Subst{N'}{Q'}{x}$} دے؛ دا د ځينو \LR{$x$}،
      \LR{$N$}، \LR{$Q$}، \LR{$N'$} او~\LR{$Q'$} دپاره دي، چې \LR{$N \bredpar N'$} او \LR{$Q \bredpar
      Q'$} وي۔ د استقرا له فرضيې \LR{$N' \bredpar \bcd{N}$} او
      \LR{$Q' \bredpar \bcd{Q}$} لرو۔ د \olref{lem:comp} له مخې
      \LR{$\Subst{N'}{Q'}{x} \bredpar \Subst{\bcd{N}}{\bcd{Q}}{x}$}؛
      ښۍ خوا کټ مټ \LR{$\bcd{((\lambd[x][N])Q)}$} ده۔
  \end{enumerate}
\end{proof}
\olpseganchor{OLP-0369-B018}{end}

\olpseganchor{OLP-0369-B019}{start}
\begin{prob}
  د \olref[lam][cr][pb]{lem:cont} ثبوت بشپړ کړئ۔
\end{prob}
\olpseganchor{OLP-0369-B019}{end}

\olpseganchor{OLP-0369-B020}{start}
\begin{thm}\ollabel{thm:cr}
  \LR{$\bredpar$} د چرچ--روسر خاصيت لري۔
\end{thm}
\olpseganchor{OLP-0369-B020}{end}

\olpseganchor{OLP-0369-B021}{start}
\begin{proof}
  دا د \olref{lem:cont} سملاسي پايله ده۔
\end{proof}
\olpseganchor{OLP-0369-B021}{end}

\olpseganchor{OLP-0370-B004}{start}
\olfileid{lam}{cr}{b}
\olpseganchor{OLP-0370-B004}{end}

\olpseganchor{OLP-0370-B005}{start}
\olsection{\texorpdfstring{\LR{$\beta$}}{β}-راکمول}
\olpseganchor{OLP-0370-B005}{end}

\olpseganchor{OLP-0370-B006}{start}
\begin{lem}\ollabel{lem:one-par}
  که \LR{$M \bredone M'$} وي، نو \LR{$M \bredpar M'$}.
\end{lem}
\begin{proof} که \LR{$M \bredone M'$} وي او انقباض د ترم په سر کښې
    وي، نو \LR{$M$} د ځينو \LR{$x$}، \LR{$N$} او~\LR{$Q$} دپاره
    \LR{$(\lambd[x][N])Q$} دے او \LR{$M'$}، \LR{$\Subst{N}{Q}{x}$} دے۔
    د \olref[pb]{thm:refl} له مخې \LR{$N \bredpar N$} او \LR{$Q \bredpar Q$}
    دي؛ نو له \olref[pb]{defn:bredpar} او
    \olref[pb]{defn:bredpar4} څخه سملاسي
    \LR{$(\lambd[x][N])Q \bredpar \Subst{N}{Q}{x}$} راځي۔
    که انقباض د تجريد په بدن کښې وي، د کوچني ترم د استقرايي فرضې راکمښت د تجريد په هممهاله قاعده کښې ږدو۔ که د تطبيق په کيڼ يا ښي غړي کښې وي، هماغه استقرايي فرض له بل غړي د انعکاس سره د تطبيق په قاعده کښې ږدو۔
\end{proof}
\olpseganchor{OLP-0370-B006}{end}

\olpseganchor{OLP-0370-B007}{start}
\begin{lem}\ollabel{lem:par-red}
  که \LR{$M \bredpar M'$} وي، نو \LR{$M \bred M'$}.
\end{lem}
\olpseganchor{OLP-0370-B007}{end}

\olpseganchor{OLP-0370-B008}{start}
\begin{proof} د \LR{$M \bredpar M'$} د اشتقاق پر جوړښت استقرا کوو۔
  \begin{enumerate}
  \item وروستۍ قاعده \olref[pb]{defn:bredpar1} ده: په دې حالت کښې
    \LR{$M$} او \LR{$M'$} دواړه هماغه \LR{$x$} دي، او \LR{$x \bred x$}.
  \item وروستۍ قاعده \olref[pb]{defn:bredpar2} ده: \LR{$M$} د ځينو
    \LR{$x$}، \LR{$N$} او \LR{$N'$} دپاره \LR{$\lambd[x][N]$} او \LR{$M'$}،
    \LR{$\lambd[x][N']$} دے، چې \LR{$N \bredpar N'$} وي۔ د استقرا له
    فرضيې \LR{$N \bred N'$} لرو۔ نو \LR{$\lambd[x][N] \bred
    \lambd[x][N']$}؛ دا د \LR{$N \bred N'$} د هماغو \LR{$\bredone$}
    انقباضونو لړۍ د تجريد دننه تطبيقول دي۔
  \item وروستۍ قاعده \olref[pb]{defn:bredpar3} ده: \LR{$M$}،
    \LR{$PQ$} او \LR{$M'$} د ځينو \LR{$P$}، \LR{$Q$}، \LR{$P'$} او \LR{$Q'$} دپاره \LR{$P'Q'$}
    دے، چې \LR{$P \bredpar P'$} او \LR{$Q \bredpar Q'$} وي۔ د استقرا
    له فرضيې \LR{$P \bred P'$} او \LR{$Q \bred Q'$} لرو۔ نو لومړے
    د \LR{$P \bred P'$} لړۍ او بيا د \LR{$Q \bred Q'$} لړۍ په کارولو
    \LR{$PQ \bred P'Q'$} ترلاسه کېږي۔
  \item وروستۍ قاعده \olref[pb]{defn:bredpar4} ده: \LR{$M$}،
    \LR{$(\lambd[x][N])Q$} او \LR{$M'$}، \LR{$\Subst{N'}{Q'}{x}$} دے؛ دا
    د ځينو \LR{$x$}، \LR{$N$}، \LR{$N'$}، \LR{$Q$} او \LR{$Q'$} دپاره دي، چې
    \LR{$N \bredpar N'$} او \LR{$Q \bredpar Q'$} وي۔ د استقرا له
    فرضيې \LR{$Q \bred Q'$} او \LR{$N \bred N'$} ترلاسه کوو۔ نو
    \LR{$(\lambd[x][N])Q \bred \Subst{N'}{Q'}{x}$}: لومړے
    \LR{$N \bred N'$}، بيا \LR{$Q \bred Q'$}، او په پاے کښې د
    \LR{$(\lambd[x][N'])Q'$} انقباض \LR{$\Subst{N'}{Q'}{x}$} ته۔
  \end{enumerate}
\end{proof}
\olpseganchor{OLP-0370-B008}{end}

\olpseganchor{OLP-0370-B009}{start}
\begin{lem}\ollabel{lem:str}
  \LR{$\bred$} تر ټولو کوچنۍ انتقالي اړيکه ده چې \LR{$\bredpar$} پکې
  شامله ده۔
\end{lem}
\olpseganchor{OLP-0370-B009}{end}

\olpseganchor{OLP-0370-B010}{start}
\begin{proof}
  \LR{$\xred$} دې تر ټولو کوچنۍ انتقالي اړيکه وي چې
  \LR{$\bredpar$} پکې شامله ده۔
\olpseganchor{OLP-0370-B010}{end}

\olpseganchor{OLP-0370-B011}{start}
  \LR{$\bred \subseteq \xred$}: فرض کړئ \LR{$M \bred M'$}، يعنې
  \LR{$M \ident M_1 \bredone \dots \bredone M_k \ident M'$}۔
  د \olref{lem:one-par} له مخې \LR{$M
  \ident M_1 \bredpar \dots \bredpar M_k \ident M'$}۔
  څرنګه چې \LR{$\xred$}، \LR{$\bredpar$} رانغاړي او انتقالي ده،
  \LR{$M \xred M'$} ترلاسه کېږي۔ که لړۍ صفر ګامونه ولري، د هممهاله راکمولو انعکاس د ګډې اړيکې دغه حالت ورکوي۔
\olpseganchor{OLP-0370-B011}{end}

\olpseganchor{OLP-0370-B012}{start}
  \LR{$\xred \subseteq \bred$}: فرض کړئ \LR{$M \xred M'$}، يعنې \LR{$M
  \ident M_1 \bredpar \dots \bredpar M_k \ident M'$}۔ د \olref{lem:par-red}
  له مخې \LR{$M \ident M_1 \bred \dots \bred M_k \ident M'$}۔
  څرنګه چې \LR{$\bred$} انتقالي ده، \LR{$M \bred M'$}.
\end{proof}
\olpseganchor{OLP-0370-B012}{end}

\olpseganchor{OLP-0370-B013}{start}
\begin{thm}\ollabel{thm:cr}
  \LR{$\bred$} د چرچ--روسر خاصيت لري۔
\end{thm}
\olpseganchor{OLP-0370-B013}{end}

\olpseganchor{OLP-0370-B014}{start}
\begin{proof}
  دا د \olref[dap]{thm:str}، \olref[pb]{thm:cr} او
  \olref{lem:str} سملاسي پايله ده۔
\end{proof}
\olpseganchor{OLP-0370-B014}{end}

\olpseganchor{OLP-0371-B004}{start}
\olfileid{lam}{cr}{pbe}
\olpseganchor{OLP-0371-B004}{end}

\olpseganchor{OLP-0371-B005}{start}
\olsection{هممهاله \texorpdfstring{\LR{$\beta\eta$}}{βη}-راکمول}
\olpseganchor{OLP-0371-B005}{end}

\olpseganchor{OLP-0371-B006}{start}
په دې برخه کښې د هممهاله \LR{$\beta\eta$}-راکمول د چرچ--روسر خاصيت
ثابتوو۔ دا د \LR{$\beta\eta$}-راکمول اړوند هممهاله راکمول دے۔
\olpseganchor{OLP-0371-B006}{end}

\olpseganchor{OLP-0371-B007}{start}
\begin{defn}[هممهاله \LR{$\beta\eta$}-راکمول، \LR{$\beredpar$}] \ollabel{defn:beredpar}
  پر ترمونو \emph{هممهاله \LR{$\beta\eta$}-راکمول} (\LR{$\beredpar$})
  په استقرايي ډول داسې تعريفېږي:
  \begin{enumerate}
    \item \ollabel{defn:beredpar1} \LR{$x \beredpar x$}.
    \item \ollabel{defn:beredpar2} که \LR{$N \beredpar N'$} وي، نو \LR{$\lambd[x][N] \beredpar
      \lambd[x][N']$}.
    \item \ollabel{defn:beredpar3} که \LR{$P \beredpar P'$} او \LR{$Q \beredpar Q'$} وي، نو \LR{$PQ \beredpar
      P'Q'$}.
    \item \ollabel{defn:beredpar4} که \LR{$N \beredpar N'$} او \LR{$Q \beredpar Q'$} وي، نو
      \LR{$(\lambd[x][N])Q \beredpar \Subst{N'}{Q'}{x}$}.
    \item \ollabel{defn:beredpar5} که \LR{$N \beredpar N'$} وي، نو \LR{$\lambd[x][Nx]
      \beredpar N'$}، په دې شرط چې \LR{$x \notin FV(N)$} وي۔
  \end{enumerate}
\end{defn}
\olpseganchor{OLP-0371-B007}{end}

\olpseganchor{OLP-0371-B008}{start}
\begin{thm}\ollabel{thm:refl}
  \LR{$M \beredpar M$}.
\end{thm}
\olpseganchor{OLP-0371-B008}{end}

\olpseganchor{OLP-0371-B009}{start}
\begin{proof}
  تمرين۔
\end{proof}
\olpseganchor{OLP-0371-B009}{end}

\olpseganchor{OLP-0371-B010}{start}
\begin{prob}
  \olref[lam][cr][pbe]{thm:refl} ثابت کړئ۔
\end{prob}
\olpseganchor{OLP-0371-B010}{end}

\olpseganchor{OLP-0371-B011}{start}
\begin{defn}[\LR{$\beta\eta$}-بشپړ راکمول]\ollabel{defn:becd}
  د~\LR{$M$} \emph{\LR{$\beta\eta$}-بشپړ راکمول}~\LR{$\becd{M}$} داسې
  تعريفېږي:
  \begin{align}
    \becd{x} &= x \ollabel{defn:becd1} \\
    \becd{(\lambd[x][N])} &= \lambd[x][\becd{N}] \ollabel{defn:becd2}\\
    \becd{(PQ)} &= \becd{P}\becd{Q} && \text{\RL{که \LR{$P$} د \LR{$\lambd$} تجريد نۀ وي}} 
    \ollabel{defn:becd3} \\
    \becd{((\lambd[x][N])Q)} &= \Subst{\becd{N}}{\becd{Q}}{x}
                             \ollabel{defn:becd4} \\
    \becd{(\lambd[x][Nx])} &= \becd{N} \ollabel{defn:becd5} & \text{\RL{که \LR{$x
                                                      \notin FV(N)$} وي}}
  \end{align}
  د تجريد دپاره که د اېټا شرط پوره وي، وروستۍ قاعده کاروو؛ په نورو تجريدونو کښې عمومي قاعده کاروو۔
\end{defn}
\olpseganchor{OLP-0371-B011}{end}

\olpseganchor{OLP-0371-B012}{start}
\begin{lem}\ollabel{lem:comp}
  که \LR{$M \beredpar M'$} او \LR{$R \beredpar R'$} وي، نو \LR{$\Subst{M}{R}{y}
  \beredpar \Subst{M'}{R'}{y}$}.
\end{lem}
\olpseganchor{OLP-0371-B012}{end}

\olpseganchor{OLP-0371-B013}{start}
\begin{proof}
  د \LR{$M \beredpar M'$} د اشتقاق پر جوړښت استقرا کوو۔
\olpseganchor{OLP-0371-B013}{end}

\olpseganchor{OLP-0371-B014}{start}
  په ټول ثبوت کښې تړلي نومونه د الفا-بدلون له لارې تازه ټاکو۔ هممهاله راکمول نوے ازاد نوم نۀ راولي: دا د قاعدو پر جوړښت له استقرا او د تعويض د ازادو نومونو له قاعدې راځي۔ لومړي څلور حالتونه د \olref[pb]{lem:comp} د حالتونو په شان دي۔
  که وروستۍ قاعده \olref{defn:beredpar5} وي، نو \LR{$M$}،
  \LR{$\lambd[x][Nx]$} او \LR{$M'$} د ځينو \LR{$x$} او~\LR{$N'$} دپاره \LR{$N'$}
  دے، چې \LR{$x \notin FV(N)$} او \LR{$N \beredpar N'$} وي۔ غواړو
  وښيو چې \LR{$\Subst{(\lambd[x][Nx])}{R}{y} \beredpar
  \Subst{N'}{R'}{y}$}، يعنې
  \LR{$\lambd[x][\Subst{N}{R}{y} x] \beredpar \Subst{N'}{R'}{y}$}۔
  دا د \olref{defn:beredpar}، \olref{defn:beredpar5} او
  د استقرا له فرضيې راځي۔
\end{proof}
\olpseganchor{OLP-0371-B014}{end}

\olpseganchor{OLP-0371-B015}{start}
\begin{lem}\ollabel{lem:cont}
  که \LR{$M \beredpar M'$} وي، نو \LR{$M' \beredpar \becd{M}$}.
\end{lem}
\olpseganchor{OLP-0371-B015}{end}

\olpseganchor{OLP-0371-B016}{start}
\begin{proof}
  د \LR{$M \beredpar M'$} د اشتقاق وروستۍ قاعده د استقرا په هر ګام کښې څېړو۔
\olpseganchor{OLP-0371-B016}{end}

\olpseganchor{OLP-0371-B017}{start}
  د دوو دعوو دپاره پر ترمونو يوځايي استقرا کوو، او د کوچني ترم دپاره دواړه دعوې فرضوو۔ لومړۍ دعوه همدا لمه ده۔ دويمه دعوه دا ده: که \LR{$P \ident \lambd[x][N]$}، \LR{$P \beredpar P'$} او \LR{$Q' \beredpar \becd{Q}$} وي، نو \LR{$P'Q' \beredpar \Subst{\becd{N}}{\becd{Q}}{x}$}۔ په دويمه دعوه کښې استقرا د \LR{$P$} پر جوړښت ده او د \LR{$Q$} او \LR{$Q'$} په ټاکلو هېڅ بنديز نشته۔
\par په ټول ثبوت کښې تړلي نومونه د الفا-بدلون له لارې تازه ټاکو۔ هممهاله راکمول نوے ازاد نوم نۀ راولي: دا د قاعدو پر جوړښت له استقرا او د تعويض د ازادو نومونو له قاعدې راځي۔ لومړے دويمه دعوه ثابته کړو۔ که \LR{$P'$} د تجريد له قاعدې \LR{$\lambd[x][N']$} وي، د لومړۍ دعوې استقرايي فرض له \LR{$N \beredpar N'$} څخه \LR{$N' \beredpar \becd{N}$} ورکوي۔ اوس د بېټا قاعده پر \LR{$P'Q'$} ولګوئ۔ که \LR{$P'$} د اېټا له قاعدې راغلے وي، نو \LR{$N \ident Sx$}، \LR{$x \notin FV(S)$} او \LR{$S \beredpar P'$}۔ که \LR{$S$} تجريد نۀ وي، \LR{$\becd{N} \ident \becd{S}x$}؛ د لومړۍ دعوې استقرايي فرض \LR{$P' \beredpar \becd{S}$} ورکوي او د تطبيق قاعده مطلوبه پايله ورکوي۔ که \LR{$S \ident \lambd[y][T]$} وي، \LR{$\becd{N} \ident \Subst{\becd{T}}{x}{y}$}۔ دويمه استقرايي دعوه پر کوچني \LR{$S$} ولګوئ: \LR{$P'Q' \beredpar \Subst{\becd{T}}{\becd{Q}}{y}$}۔ ځکه چې \LR{$x$} په \LR{$S$} کښې ازاد نۀ دے، دا هدف له \LR{$\Subst{\becd{N}}{\becd{Q}}{x}$} سره د تعويض له قاعدې برابر دے۔
\par د لومړۍ دعوې دپاره د متغير حالت انعکاسي دے۔ د تطبيق په حالت کښې که اصلي کيڼ ترم تجريد نۀ وي، د دواړو کوچنيو ترمونو د استقرايي فرضو او د تطبيق له قاعدې هدف راځي۔ که کيڼ ترم تجريد وي او وروستۍ قاعده تطبيق وي، د هغه کيڼ ترم دپاره دويمه استقرايي دعوه کاروو۔ که وروستۍ قاعده د بېټا انقباض وي، د تجريد د بدن او د ښي ترم استقرايي فرضونه د تعويض له لمې سره يوځاے مطلوبه پايله ورکوي۔
\par د تجريد په حالت کښې، که اصلي ترم د اېټا رېډکس نۀ وي، د بدن استقرايي فرض او د تجريد قاعده بس دي۔ که اصلي ترم \LR{$\lambd[x][Sx]$} وي او \LR{$x \notin FV(S)$}، هدف \LR{$\becd{S}$} دے۔ که د بدن هممهاله ګام د تطبيق له قاعدې وي، نو پايله \LR{$\lambd[x][S'x]$} ده او \LR{$S \beredpar S'$}۔ له استقرا او د ازادو نومونو له شرطه \LR{$S' \beredpar \becd{S}$} او \LR{$x \notin FV(S')$} لرو، نو د اېټا قاعده کار کوي۔ که د بدن هممهاله ګام په سر کښې بېټا منقبض کړي، \LR{$S \ident \lambd[y][T]$} او \LR{$T \beredpar T'$} وي۔ نو پايله \LR{$\lambd[x][\Subst{T'}{x}{y}]$} ده، چې د تازګۍ له امله له \LR{$\lambd[y][T']$} سره الفا-معادله ده۔ د تجريد قاعده دا له \LR{$S$} څخه هممهاله راکموي؛ نو د کوچني \LR{$S$} دپاره لومړۍ استقرايي فرضيه دا مستقيماً \LR{$\becd{S}$} ته راکموي۔
  که وروستۍ قاعده \olref{defn:beredpar5} وي، نو \LR{$M$}،
  \LR{$\lambd[x][Nx]$} او \LR{$M'$} د ځينو \LR{$x$}، \LR{$N$} او \LR{$N'$} دپاره
  \LR{$N'$} دے، چې \LR{$x \notin FV(N)$} او \LR{$N \beredpar N'$} وي۔
  غواړو وښيو چې \LR{$N' \beredpar
  \becd{(\lambd[x][Nx])}$}، يعنې \LR{$N' \beredpar \becd{N}$}؛
  دا د استقرا له فرضيې سملاسي راځي۔
\end{proof}
\olpseganchor{OLP-0371-B017}{end}

\olpseganchor{OLP-0371-B018}{start}
\begin{thm}\ollabel{thm:cr}
  \LR{$\beredpar$} د چرچ--روسر خاصيت لري۔
\end{thm}
\olpseganchor{OLP-0371-B018}{end}

\olpseganchor{OLP-0371-B019}{start}
\begin{proof}
  دا د \olref{lem:cont} سملاسي پايله ده۔
\end{proof}
\olpseganchor{OLP-0371-B019}{end}

\olpseganchor{OLP-0372-B004}{start}
\olfileid{lam}{cr}{be}
\olpseganchor{OLP-0372-B004}{end}

\olpseganchor{OLP-0372-B005}{start}
\olsection{\texorpdfstring{\LR{$\beta\eta$}}{βη}-راکمول}
\olpseganchor{OLP-0372-B005}{end}

\olpseganchor{OLP-0372-B006}{start}
د چرچ--روسر خاصيت د \LR{$\beta\eta$}-راکمول (\LR{$\bered$}) دپاره
هم صدق کوي۔ دلته ګډ يوګامی راکمول د بېټا او اېټا د يوګاميو انقباضونو يوځاے کول دي۔ انقباض په سر يا د يوه کوچني ترم دننه کېدے شي۔
\olpseganchor{OLP-0372-B006}{end}

\olpseganchor{OLP-0372-B007}{start}
\begin{lem}\ollabel{lem:one-par}
  که \LR{$M \beredone M'$} وي، نو \LR{$M \beredpar M'$}.
\end{lem}
\olpseganchor{OLP-0372-B007}{end}

\olpseganchor{OLP-0372-B008}{start}
\begin{proof} 
  د \LR{$M \beredone M'$} د اشتقاق پر جوړښت استقرا کوو۔
  که \LR{$M \eredone M'$} د \LR{$\eta$}-بدلون له لارې وي، يعنې
  \olref[syn][eta]{defn:beredone}، نو
  \olref[pbe]{thm:refl} کاروو۔ په سر کښې د بېټا حالت هم د بدن او آرګومېنټ له انعکاسه راځي۔ که انقباض دنننی وي، په راکمېدونکي غړي استقرا کوو او پاتې غړي انعکاسي نيسو؛ د تجريد او تطبيق قواعد دا ګام بېرته په ټول ترم کښې ږدي۔ د بېټا دغه حالتونه د
  \olref[b]{lem:one-par} په شان دي۔
\end{proof}
\olpseganchor{OLP-0372-B008}{end}


\olpseganchor{OLP-0372-B009}{start}
\begin{lem}\ollabel{lem:par-red}
  که \LR{$M \beredpar M'$} وي، نو \LR{$M \bered M'$}.
\end{lem}
\olpseganchor{OLP-0372-B009}{end}

\olpseganchor{OLP-0372-B010}{start}
\begin{proof} د \LR{$M \beredpar M'$} د اشتقاق پر جوړښت استقرا کوو۔
\olpseganchor{OLP-0372-B010}{end}

\olpseganchor{OLP-0372-B011}{start}
  که وروستۍ قاعده \olref[pbe]{defn:beredpar5} وي، نو \LR{$M$}،
  \LR{$\lambd[x][Nx]$} او \LR{$M'$} د ځينو \LR{$x$}، \LR{$N$} او \LR{$N'$} دپاره
  \LR{$N'$} دے، چې \LR{$x \notin FV(N)$} او \LR{$N \beredpar N'$} وي۔
  نو لومړے \LR{$\lambd[x][Nx]$} د \LR{$\eta$}-بدلون له لارې \LR{$N$}
  ته راکمولے شو؛ بيا د \LR{$\beredone$} د ګامونو هغه لړۍ
  تعقيبوو چې \LR{$N \bered N'$} ښيي۔ وروستۍ خبره د استقرا
  له فرضيې راځي۔ متغير د صفر ګامونو لړۍ اخلي۔ د تجريد او تطبيق حالتونه د استقرا د لړيو په هماغه ځايونو کښې په لګولو راځي۔ د بېټا د سر حالت کښې لومړے بدن او آرګومېنټ د استقرا له مخې راکموو، بيا د سر بېټا انقباض کوو۔
\end{proof}
\olpseganchor{OLP-0372-B011}{end}


\olpseganchor{OLP-0372-B012}{start}
\begin{lem}\ollabel{lem:str}
  \LR{$\bered$} تر ټولو کوچنۍ انتقالي اړيکه ده چې \LR{$\beredpar$}
  پکې شامله ده۔
\end{lem}
\olpseganchor{OLP-0372-B012}{end}

\olpseganchor{OLP-0372-B013}{start}
\begin{proof}
  لکه څنګه چې په \olref[b]{lem:str} کښې وشول۔
\end{proof}
\olpseganchor{OLP-0372-B013}{end}

\olpseganchor{OLP-0372-B014}{start}
\begin{thm}\ollabel{thm:cr}
  \LR{$\bered$} د چرچ--روسر خاصيت لري۔
\end{thm}
\olpseganchor{OLP-0372-B014}{end}

\olpseganchor{OLP-0372-B015}{start}
\begin{proof}
  دا د \olref[dap]{thm:str}، \olref[pbe]{thm:cr} او
  \olref{lem:str} له مخې راځي۔
\end{proof}
\olpseganchor{OLP-0372-B015}{end}

\olpseganchor{OLP-0373-B004}{start}
\olchapter{lam}{rep}{د لامبډا په وسيله تعريفېدنه}
\olpseganchor{OLP-0373-B004}{end}

\olpseganchor{OLP-0373-B005}{start}
\begin{editorial}
  دا څپرکے لا ازمايښتي دے۔ لا زياتې څرګندونې ته اړتيا لري؛
  مواد ئې هم بايد په تعريفونو، له ثبوتونو سره په قضيو، او
  نورو بېلګو کښې ښه منظم شي۔
\end{editorial}
\olpseganchor{OLP-0373-B005}{end}

\olpseganchor{OLP-0373-B006}{start}




%





\olpseganchor{OLP-0373-B006}{end}

\olpseganchor{OLP-0374-B004}{start}
\olfileid{lam}{ldf}{int}
\olsection{سريزه}
\olpseganchor{OLP-0374-B004}{end}

\olpseganchor{OLP-0374-B005}{start}
په لومړي نظر لامبډا حساب د هغو عبارتونو ډېر تجريدي حساب دے چې
تابعې او پر نورو شيانو د هغو تطبيق ښيي۔ د لامبډا حساب په نحو کښې
هېڅ داسې نښه نشته چې وښيي دا د کوم ځانګړي ډول شيانو تابعې دي؛ په
ځانګړې توګه د طبيعي عددونو يادونه پکې نشته۔ د دې حساب بنسټيز عملونه،
يعنې تطبيق او لامبډا تجريد، پر هرې تابع کار کوي، يوازې د طبيعي
عددونو پر تابعو نۀ۔
\olpseganchor{OLP-0374-B005}{end}

\olpseganchor{OLP-0374-B006}{start}
له دې سره سره، په لږ ځيرکتيا سره د طبيعي عددونو تابعې، يعنې حسابي
تابعې، په لامبډا حساب کښې تعريفولے شو۔ د دې دپاره د هر طبيعي
عدد~\LR{$n \in \Nat$} په مقابل کښې يو ځانګړے \LR{$\lambd$}-ترم~\LR{$\num n$}
تعريفوو؛ دا د~\LR{$n$} \emph{د چرچ عددنښه} ده۔ (د چرچ عددنښې د الونزو
چرچ په نوم نومول شوې دي۔)
\olpseganchor{OLP-0374-B006}{end}

\olpseganchor{OLP-0374-B007}{start}
\begin{defn}
  که \LR{$n \in \Nat$} وي، اړونده \emph{د چرچ عددنښه} \LR{$\num{n}$}
  د~\LR{$n$} استازيتوب کوي:
  \[
    \num{n} \ident \lambd[fx][f^n(x)]
  \]
  دلته \LR{$f^n(x)$} د~\LR{$x$} پر سر د~\LR{$f$} د~\LR{$n$} ځلو تطبيق پايله
  ښيي۔ د بېلګې په توګه، \LR{$\num{0}$}، \LR{$\lambd[fx][x]$} او
  \LR{$\num{3}$}، \LR{$\lambd[fx][f(f(f\,x))]$} دے۔
\end{defn}
\olpseganchor{OLP-0374-B007}{end}
  
\olpseganchor{OLP-0374-B008}{start}
د چرچ عددنښه \LR{$\num n$} د داسې لامبډا ترم په توګه کوډ شوې ده چې
دوه آرګومېنټونه \LR{$f$} او \LR{$x$} اخلي او \LR{$f^n(x)$} ورکوي۔ څرګنده ده چې
د چرچ عددنښې په عادي بڼه کښې دي۔
\olpseganchor{OLP-0374-B008}{end}

\olpseganchor{OLP-0374-B009}{start}
په لامبډا حساب کښې د طبيعي عددونو تمثيل هله ګټور دے چې پرې
محاسبه هم وکړے شو۔ د چرچ له عددنښو سره محاسبه دا معنا لري چې
يو \LR{$\lambd$}-ترم~\LR{$F$} پر داسې عددنښه تطبيق کړو او ګډ ترم~\LR{$F\,
\num n$} عادي بڼې ته راکم کړو۔ که هغه تل عادي بڼې ته راکمېږي
او دا بڼه تل د چرچ عددنښه~\LR{$\num m$} وي، د محاسبې وتون د~\LR{$m$}
عدد ګڼلے شو۔ بيا \LR{$F$} داسې تابع \LR{$f\colon \Nat \to \Nat$} تعريفوي
چې \LR{$f(n) = m$} هغه وخت او يوازې هغه وخت وي چې \LR{$F\, \num n
\red \num m$} وي۔ د چرچ--روسر د خاصيت له مخې، که عادي بڼه
موجوده وي، يکتا ده۔ نو که \LR{$F\, \num n \red \num m$} وي،
\LR{$F \, \num n$} بلې عادي بڼې، په ځانګړې توګه بلې د چرچ عددنښې ته،
نۀ شي راکمېدلے۔
\olpseganchor{OLP-0374-B009}{end}

\olpseganchor{OLP-0374-B010}{start}
پرعکس، که يوه تابع \LR{$f\colon \Nat \to \Nat$} راکړل شي، پوښتلے
شو چې آيا داسې ترم \LR{$F$} شته چې په دې ډول~\LR{$f$} تعريف کړي۔ په دې
حالت کښې وايو \LR{$F$}،~\LR{$f$} \emph{لامبډا-تعريفوي} او \LR{$f$}،
د لامبډا په وسيله د تعريف وړ ده۔ دا خبره څوځاييزو او جزوي تابعو ته هم
غځولے شو۔
\olpseganchor{OLP-0374-B010}{end}

\olpseganchor{OLP-0374-B011}{start}
\begin{defn}
فرض کړئ \LR{$f\colon \Nat^k \to \Nat$}۔ وايو يو لامبډا ترم~\LR{$F$}
هغه وخت \LR{$f$} \emph{لامبډا-تعريفوي} چې د ټولو
\LR{$n_0$}، \dots،~\LR{$n_{k-1}$} دپاره،
\[
F \, \num{n_0} \, \num{n_1} \dots \num{n_{k-1}} \red \num{f(n_0, n_1, \dots,
  n_{k-1})}
\]
وي، که \LR{$f(n_0, \dots, n_{k-1})$} تعريف شوې وي؛ او که نۀ وي،
\LR{$F \, \num{n_0} \,
\num{n_1} \dots \num{n_{k-1}}$} هېڅ عادي بڼه نۀ لري۔
\end{defn}
\olpseganchor{OLP-0374-B011}{end}

\olpseganchor{OLP-0374-B012}{start}
ثابتې تابعې ډېرې ساده بېلګې دي۔ ترم \LR{$C_k \ident
\lambd[x][\num{k}]$} هغه تابع \LR{$c_k\colon \Nat \to
\Nat$} لامبډا-تعريفوي چې \LR{$c(n) = k$} وي۔ ځکه د هر~\LR{$n$}
دپاره \LR{$C_k \, \num n \ident
(\lambd[x][\num{k}])\num n \redone \num{k}$}۔
د يووالي تابع \LR{$\lambd[x][x]$} لامبډا-تعريف کوي۔ د پېچلو تابعو
تعريفول، البته، ډېر ګران دي او ډېر ځله ځيرکتيا غواړي۔ نو دا
ښايي حيرانوونکې وي چې هره محاسبه کېدونکې تابع
د لامبډا په وسيله د تعريف وړ ده۔ برعکس خبره هم رښتيا ده: که يوه تابع
د لامبډا په وسيله د تعريف وړ وي، نو محاسبه کېدونکې ده۔
\olpseganchor{OLP-0374-B012}{end}

\olpseganchor{OLP-0375-B004}{start}
\olfileid{lam}{rep}{arf}
\olsection{د لامبډا په وسيله د تعريف وړ حسابي تابعې}
\olpseganchor{OLP-0375-B004}{end}

\olpseganchor{OLP-0375-B005}{start}
\begin{prop}
  \ollabel{prop:succ-ld}
  د ورپسې شمېرې تابع~\LR{$\Succ$} د لامبډا په وسيله د تعريف وړ ده۔
\end{prop}
\olpseganchor{OLP-0375-B005}{end}

\olpseganchor{OLP-0375-B006}{start}
\begin{proof}
هغه ترم چې د ورپسې شمېرې تابع لامبډا-تعريفوي، دا دے:
\begin{align*}
  \fn{Succ} & \ident \lambd[a][\lambd[fx][f ({a} f x)]].
  \intertext{\RL{زموږ د ليکنې د دود له مخې، دا د لاندې ترم لنډه بڼه ده:}}
  \fn{Succ} & \ident \lambd[a][\lambd[f][\lambd[x][(f (({a} f) x))]]].
  \intertext{\RL{\LR{$\fn{Succ}$} داسې تابع ده چې عدد~\LR{$a$} د آرګومېنټ په توګه
    اخلي او پايله ئې بله تابع، \LR{$\lambd[fx][f ({a}
      f x)]$}، ده۔ دا تابع په خپله د چرچ عددنښه نۀ ده۔ خو
    که آرګومېنټ \LR{$a$} د چرچ عددنښه وي، عددنښې ته راکمېږي۔ وګورئ:}}
   (\lambd[a][\lambd[fx][f ({a} f x)]])\,\num{n} & \redone
   \lambd[fx][f ({\num{n}} f x)].
   \intertext{\RL{دنننے ترم \LR{$\num{n}fx$} د بېټا راکمول وړ دے، ځکه
     \LR{$\num{n}$}، \LR{$\lambd[fx][f^nx]$} دے۔ نو \LR{$\num{n}fx \red f^nx$}؛
     په دې توګه د ټول ترم دپاره لرو:}}
   \fn{Succ}\, \num n & \red \lambd[fx][f(f^n(x))],
\end{align*}
د عددنښې دواړه آرګومېنټونه په پرله‌پسې بېټا ګامونو لګېږي۔
يعنې \LR{$\num{n+1}$}۔
\end{proof}
\olpseganchor{OLP-0375-B006}{end}

\olpseganchor{OLP-0375-B007}{start}
\begin{ex}
  وګورو چې د~\LR{$\num{0}$}، يعنې \LR{$\lambd[fx][x]$}، پر سر
  د~\LR{$\fn{Succ}$} په تطبيق څۀ پېښېږي۔ ترمونه په بشپړه بڼه ليکو:
  \begin{align*}
    \fn{Succ}\, \num{0} & \ident (\lambd[a][\lambd[f][\lambd[x][(f (({a} f) x))]]])(\lambd[f][\lambd[x][x]])\\
    & \redone \lambd[f][\lambd[x][(f (({(\lambd[f][\lambd[x][x]])} f) x))]] \\
    & \redone \lambd[f][\lambd[x][(f ({(\lambd[x][x])} x))]] \\
    & \redone \lambd[f][\lambd[x][(f x)]] \ident \num{1}\\
  \end{align*}
\end{ex}
\olpseganchor{OLP-0375-B007}{end}

\olpseganchor{OLP-0375-B008}{start}
\begin{prob}
  دا ترم
  \begin{align*}
    \fn{Succ}' & \ident \lambd[{n}][\lambd[fx][{n} f (f x)]]
  \end{align*}
  د ورپسې شمېرې تابع لامبډا-تعريفوي۔ څرګنده کړئ چې ولې۔
\end{prob}
\olpseganchor{OLP-0375-B008}{end}

\olpseganchor{OLP-0375-B009}{start}
\begin{prop}
  \ollabel{prop:add-ld}
  د جمع تابع~\LR{$\Add$} د لامبډا په وسيله د تعريف وړ ده۔
\end{prop}
\olpseganchor{OLP-0375-B009}{end}

\olpseganchor{OLP-0375-B010}{start}
\begin{proof}
جمع د دې ترمونو په وسيله لامبډا-تعريف شوې ده:
\begin{align*}
  \fn{Add} & \ident \lambd[{a}{b}][\lambd[fx][{a} f ({b} f x)]]
\intertext{\RL{يا په بل ډول:}}
  \fn{Add}' & \ident \lambd[{a}{b}][{a}\, \fn{Succ} \, {b}].
\intertext{\RL{لومړۍ جمع داسې کار کوي: \LR{$\fn{Add}$} لومړے دوه
  عددونه \LR{${a}$} او \LR{${b}$} اخلي۔ پايله ئې يوه تابع ده چې
  \LR{$f$} او \LR{$x$} اخلي او \LR{$af(bfx)$} ورکوي۔ که \LR{$a$} او \LR{$b$} د
  چرچ عددنښې \LR{$\num{n}$} او \LR{$\num{m}$} وي، دا \LR{$f^{n+m}(x)$}
  ته راکمېږي؛ دا له \LR{$f^{n}(f^{m}(x))$} سره يو شان دے۔
  په ورو ګامونو کښې:}}
  (\lambd[{a}{b}][\lambd[fx][{a} f ({b} f x)]])\num n\,\num m & \red
  \lambd[fx][\num{n}\, f (\num {m}\, f x)] \\
  & \red \lambd[fx][\num{n}\, f (f^m x)] \\
  & \red \lambd[fx][f^n (f^m x)] \ident \num {n+m}.
\intertext{\RL{د جمع دويم تمثيل \LR{$\fn{Add'}$} په بل ډول کار کوي:
  کله چې د چرچ پر دوو عددنښو \LR{$\num{n}$} او
  \LR{$\num{m}$} تطبيق شي،}}
\fn{Add}' \num n \,\num m
& \red \num{n}\, \fn{Succ}\, \num{m}.
\intertext{\RL{خو \LR{$\num{n} f x$} تل \LR{$f^n(x)$} ته راکمېږي؛ نو}}
  \num{n}\,  \fn{Succ}\, \num{m} & \red \fn{Succ}^n(\num{m}).
\end{align*}
د څو متغيرونو تجريد د يو متغيره تجريدونو لنډه ليکنه ده۔ له دې امله د عددنښو او جمع ترمونو د څو آرګومېنټونو تطبيق څو بېټا ګامونه اخلي۔
او ځکه چې \LR{$\fn{Succ}$} د ورپسې شمېرې تابع لامبډا-تعريفوي
او د ورپسې شمېرې تابع پر~\LR{$m$} باندې \LR{$n$} ځله تطبيق شي،
پايله \LR{$n+m$} وي؛ نو دا ترم بيا~\LR{$\num{n+m}$} ته راکمېږي۔
\end{proof}
\olpseganchor{OLP-0375-B010}{end}

\olpseganchor{OLP-0375-B011}{start}
\begin{prop}
  \ollabel{prop:mult-ld}
  ضرب د دې ترم په وسيله د لامبډا په وسيله د تعريف وړ دے:
  \[
  \fn{Mult} \ident \lambd[ab][\lambd[fx][a (b f) x]]
  \]
\end{prop}
\olpseganchor{OLP-0375-B011}{end}

\olpseganchor{OLP-0375-B012}{start}
\begin{proof}
  د دې د کار څرنګوالی داسې ګورو: فرض کړئ \LR{$\fn{Mult}$}
  د چرچ پر عددنښو \LR{$\num{n}$} او \LR{$\num{m}$} تطبيق کړو؛
  \LR{$\fn{Mult} \, \num{n} \, \num{m}$}،
  \LR{$\lambd[fx][\num{n}(\num{m}\, f)x]$} ته راکمېږي۔
  ترم \LR{$\num{m} f$} يوه تابع تعريفوي چې \LR{$f$} پر خپل آرګومېنټ
  \LR{$m$} ځله تطبيقوي۔ نو \LR{$\num{n} (\num{m} f) x$} همدا تابع،
  يعنې «\LR{$f$}، \LR{$m$} ځله تطبيق کړه»، پر~\LR{$x$} باندې \LR{$n$} ځله
  تطبيقوي۔ په بله وينا، \LR{$f$} پر~\LR{$x$} باندې \LR{$n\cdot
  m$} ځله تطبيقوو۔ خو ترلاسه شوې عادي بڼه هماغه د چرچ
  عددنښه \LR{$\num{nm}$} ده۔
\end{proof}
\olpseganchor{OLP-0375-B012}{end}

\olpseganchor{OLP-0375-B013}{start}
\begin{editorial}
موږ دا ترم د \LR{$\eta$}-راکمول په وسيله لا ساده کولے شو:
\[
  \fn{Mult} \ident \lambd[ab][\lambd[f][a (b f)]].
\]
\olpseganchor{OLP-0375-B013}{end}

\olpseganchor{OLP-0375-B014}{start}
خو بيا بايد لومړے \LR{$\eta$}-راکمول وڅرګندوو۔
\end{editorial}
\olpseganchor{OLP-0375-B014}{end}

\olpseganchor{OLP-0375-B015}{start}
\begin{prob}
ضرب د دې ترم په وسيله لامبډا-تعريف کېدے شي:
\[
  \fn{Mult}' \ident \lambd[ab][a (\fn{Add}\, b) \num{0}].
\]
لومړے آرګومېنټ د تکرار شمېر او دويم آرګومېنټ هر ځل جمع کېدونکے عدد دے۔
څرګنده کړئ چې دا ولې کار کوي۔
\end{prob}
\olpseganchor{OLP-0375-B015}{end}

\olpseganchor{OLP-0375-B016}{start}
د يو \LR{$\lambd$}-ترم په توګه د توان تعريف په حيرانوونکې توګه
ساده دے:
\[
  \fn{Exp} \ident \lambd[be][\lambd[fx][e b f x]].
\]
د صفر توان په حالت کښې، د دوو پاراميټرونو بدنه لومړنی پاراميټر پر دويم لګوي۔ دغه ترم د چرچ د يو عددنښه ده۔
لومړے آرګومېنټ \LR{$b$} بنسټ او دويم~\LR{$e$} توان دے۔ زموږ د
عددونو د کوډولو له مخې، په شهودي ډول \LR{$e f$}، \LR{$f^e$} دے۔
که دا پوهېدل ګران وي، توان د تکراري ضرب له لارې هم
تعريفولے شو:
\[
  \fn{Exp}' \ident \lambd[be][e (\fn{Mult}\, b) \num{1}].
\]
\olpseganchor{OLP-0375-B016}{end}

\olpseganchor{OLP-0375-B017}{start}
د چرچ پر عددنښو مخکنۍ شمېره او تفريق دومره ساده نۀ دي لکه
چې ښايي فکر وکړو: د جوړو کوډولو ته اړتيا لري۔
\olpseganchor{OLP-0375-B017}{end}

\olpseganchor{OLP-0376-B004}{start}
\olfileid{lam}{ldf}{pai}
\olsection{جوړې او مخکنۍ شمېره}
\olpseganchor{OLP-0376-B004}{end}

\olpseganchor{OLP-0376-B005}{start}
\begin{defn}
د \LR{$M$} او \LR{$N$} جوړه (چې \LR{$\tuple{M,N}$} ليکل کېږي) داسې تعريفېږي:
\[
\tuple{M,N} \ident \lambd[f][fMN].
\]
\end{defn}
\olpseganchor{OLP-0376-B005}{end}
  
\olpseganchor{OLP-0376-B006}{start}
په شهودي ډول دا يوه تابع ده چې بله تابع اخلي او هغه
د جوړې پر دوو غړو تطبيقوي۔ د همدې نظر له مخې لاندې
جوړوونکے لرو: دوه ترمونه اخلي او د هغو لرونکې جوړه ورکوي:
\[
  \fn{Pair} \ident \lambd[mn][\lambd[f][fmn]]
\]
که يوه جوړه راکړل شي، د غړو بېرته ترلاسه کول ئې هم غواړو۔
د دې دپاره دوو لاسرسي تابعو ته اړتيا ده چې جوړه د آرګومېنټ
په توګه اخلي او لومړے يا دويم غړے ئې ورکوي:
\begin{align*}
  \fn{Fst} & \ident \lambd[p][p(\lambd[mn][m])]\\
  \fn{Snd} & \ident \lambd[p][p(\lambd[mn][n])]
\end{align*}
\olpseganchor{OLP-0376-B006}{end}

\olpseganchor{OLP-0376-B007}{start}
\begin{prob}
  څرګنده کړئ چې د لاسرسي تابعې \LR{$\fn{Fst}$} او \LR{$\fn{Snd}$}
  ولې کار کوي۔
\end{prob}
\olpseganchor{OLP-0376-B007}{end}

\olpseganchor{OLP-0376-B008}{start}
اوس د جوړو په مرسته د مخکنۍ شمېرې تابع لامبډا-تعريف
کولے شو:
\[
  \fn{Pred} \ident \lambd[n][\fn{Fst}(n (\lambd[p][\tuple{\fn{Snd}\, {p}, \fn{Succ}(\fn{Snd}\, {p})}]) \tuple{\num 0, \num 0})]
\]
په ياد ولرئ چې \LR{$\num n\, f x$}، \LR{$f^{n}(x)$} ته راکمېږي؛
دلته \LR{$f$} داسې تابع ده چې جوړه \LR{$p$} اخلي او يوه نوې جوړه
ورکوي۔ د نوې جوړې لومړے غړے د~\LR{$p$} دويم غړے او دويم
غړے ئې د هماغه دويم غړي ورپسې شمېره ده؛ \LR{$x$}،
\LR{$\tuple{0,0}$} ده۔ نو که \LR{$n=0$} وي، پايله
\LR{$\tuple{0,0}$} ده، او که نۀ وي، \LR{$\tuple{\num{n-1}, \num n}$}
ده۔ بيا \LR{$\fn{Pred}$} د پايلې لومړے غړے ورکوي۔
\olpseganchor{OLP-0376-B008}{end}

\olpseganchor{OLP-0376-B009}{start}
تفريق داسې تعريفولے شو چې \LR{$\fn{Pred}$} پر \LR{$a$} باندې
\LR{$b$} ځله تطبيق کړو:
\[
\fn{Sub} \ident \lambd[ab][b \fn{Pred}\, a].
\]
\olpseganchor{OLP-0376-B009}{end}

\olpseganchor{OLP-0377-B004}{start}
\olfileid{lam}{ldf}{tvr}
\olsection{د رښتياوالي قيمتونه او اړيکې}
\olpseganchor{OLP-0377-B004}{end}

\olpseganchor{OLP-0377-B005}{start}
په خالص لامبډا حساب کښې د رښتياوالي قيمتونه داسې کوډولے شو:
\begin{align*}
  \fn{true} & \ident \lambd[x][\lambd[y][x]]\\
  \fn{false} & \ident \lambd[x][\lambd[y][y]]
\end{align*}
\olpseganchor{OLP-0377-B005}{end}

\olpseganchor{OLP-0377-B006}{start}
د رښتياوالي قيمتونه د \emph{ټاکونکو} په توګه تمثيلېږي؛ يعنې
داسې تابعې دي چې دوه آرګومېنټونه اخلي او يو ئې بېرته
ورکوي۔ د رښتياوالي قيمت \LR{$\fn{true}$} لومړے آرګومېنټ او
\LR{$\fn{false}$} دويم آرګومېنټ ټاکي۔ د بېلګې په توګه،
\LR{$\fn{true}\, M N$} تل \LR{$M$} ته او \LR{$\fn{false}\, M N$}
تل~\LR{$N$} ته راکمېږي۔
\olpseganchor{OLP-0377-B006}{end}

\olpseganchor{OLP-0377-B007}{start}
\begin{defn}
يوه اړيکه \LR{$R \subseteq \Nat^k$} هغه وخت د لامبډا په وسيله د تعريف وړ
بولو چې داسې ترم~\LR{$R$} وي چې
\begin{align*}
  R\, \num{n_1} \dots \num{n_k} & \bred \fn{true}
  \intertext{\RL{هر کله چې \LR{$R(n_1, \dots, n_k)$} وي؛ او}}
  R\, \num{n_1} \dots \num{n_k} & \bred \fn{false}
\end{align*}
په نورو حالتونو کښې۔
د اړيکې او د هغې د استازي ترم نوم په اصلي متن کښې يو
شان ليکل شوے؛ له سياقه ئې جلا معنا څرګندېږي۔
\end{defn}
\olpseganchor{OLP-0377-B007}{end}

\olpseganchor{OLP-0377-B008}{start}
د بېلګې په توګه، اړيکه \LR{$\fn{IsZero} = \{0\}$}، چې
پر \LR{$0$} او يوازې پر \LR{$0$} صدق کوي، د لاندې ترم په وسيله
د لامبډا په وسيله د تعريف وړ ده:
\[
  \fn{IsZero} \ident \lambd[n][n (\lambd[x][\fn{false}])\, \fn{true}].
\]
دا څنګه کار کوي؟ د چرچ عددنښې تکراروونکي دي، يعنې
خپل لومړے آرګومېنټ پر دويم \LR{$n$} ځله تطبيقوي۔ دلته
لومړنی ارزښت \LR{$\fn{true}$} ږدو او د هر تکرار په ګام کښې
د مخکني ګام له پايلې پرته \LR{$\fn{false}$} ورکوو۔ دا ګام
پر لومړني ارزښت \LR{$n$} ځله تطبيقېږي؛ پايله هغه وخت او
يوازې هغه وخت \LR{$\fn{true}$} وي چې ګام هېڅ تطبيق نۀ شي،
يعنې کله چې \LR{$n = 0$} وي۔
\olpseganchor{OLP-0377-B008}{end}

\olpseganchor{OLP-0377-B009}{start}
د رښتياوالي د قيمتونو د دې تمثيل پر بنسټ نورې د رښتياوالي تابعې
هم تعريفولے شو۔ دلته دوه دي: د نفي او عطف تمثيلونه:
\begin{align*}
  \fn{Not} & \ident \lambd[x][x\, \fn{false}\, \fn{true}]\\
  \fn{And} & \ident \lambd[x][\lambd[y][xy \,\fn{false}]]
\end{align*}
تابع «\LR{$\fn{Not}$}» يو آرګومېنټ اخلي؛ که آرګومېنټ
\LR{$\fn{false}$} وي، \LR{$\fn{true}$} ورکوي، او که آرګومېنټ
هم~\LR{$\fn{true}$} وي، \LR{$\fn{false}$} ورکوي۔ تابع
«\LR{$\fn{And}$}» دوه د رښتياوالي قيمتونه اخلي او بايد هله،
او يوازې هله، \LR{$\fn{true}$} ورکړي چې دواړه آرګومېنټونه
هم~\LR{$\fn{true}$} وي۔ د رښتياوالي قيمتونه لکه پورته د ټاکونکو
په توګه تمثيلېږي؛ نو که \LR{$x$} د رښتياوالي قيمت وي او پر دوو
آرګومېنټونو تطبيق شي، پايله لومړے آرګومېنټ وي که
\LR{$x$}، \LR{$\fn{true}$} وي، او کنه دويم آرګومېنټ۔ اوس
\LR{$\fn{And}$} خپل دوه آرګومېنټونه \LR{$x$} او \LR{$y$} اخلي او
بيا \LR{$y$} او \LR{$\fn{false}$} خپل لومړي آرګومېنټ~\LR{$x$} ته
ورکوي۔ که \LR{$x$} د رښتياوالي قيمت وي، پايله هغه وخت~\LR{$y$}
وي چې \LR{$x$}، \LR{$\fn{true}$} وي؛ او هغه وخت \LR{$\fn{false}$}
وي چې \LR{$x$}، \LR{$\fn{false}$} وي۔ همدا مو غوښتل۔
\olpseganchor{OLP-0377-B009}{end}

\olpseganchor{OLP-0377-B010}{start}
پام وکړئ دلته فرض کوو چې \LR{$\fn{And}$} ته يوازې رښتينې
ارزښتونه د آرګومېنټونو په توګه ورکول کېږي۔ که نور
ترمونه ورکړل شي، پايله ئې، يعنې عادي بڼه که موجوده
وي، ښايي د رښتياوالي قيمت نۀ وي۔
\olpseganchor{OLP-0377-B010}{end}

\olpseganchor{OLP-0377-B011}{start}
\begin{prob}
  د شاملې بېلېوالې او ځانګړې بېلېوالې د رښتياوالي د تابعو
  تمثيلوونکې تابعې \LR{$\fn{Or}$} او \LR{$\fn{Xor}$} تعريف کړئ؛
  د رښتياوالي قيمتونه د \LR{$\lambd$}-ترمونو په توګه کوډ کړئ۔
\end{prob}
\olpseganchor{OLP-0377-B011}{end}

\paragraph{د منجمدې سرچينې بشپړه اړونده برخه}
\olpseganchor{OLP-0650-B004}{start}
\olfileid{lam}{ldf}{lst-olp0650}
\olsection{لړونه}
\olpseganchor{OLP-0650-B004}{end}

\olpseganchor{OLP-0650-B005}{start}
لړ \LR{$[M_1, M_2, \ldots, M_n]$} په دې توګه تعريفولے شو:
\[
  [M_1, M_2, \ldots, M_n] \ident \lambd[sz][s M_1 (s M_2 (\ldots (s M_n z)))]
\]
په غير رسمي ژبه، د لړ تمثيلوونکے ترم د هغه لړ «راټولونکې» تابعه ده: داسې تابعه چې دوه ترمونه مني۔ لومړے يې داسې تابعه ده چې يو نوے عنصر او ورسره تر دې دمه راټوله شوې پايله اخلي، او نوې راټوله شوې پايله ورکوي۔ دويم يې د راټولېدنې لومړنے ارزښت دے۔ عناصر يو په يو راټولوي او په پاے کښې پايله ورکوي۔
\olpseganchor{OLP-0650-B005}{end}

\olpseganchor{OLP-0650-B006}{start}
\begin{digress}
که له تابعوي پروګرام جوړولو سره بلد ياست، نو دا تعريف به د \emph{ښي خوا راټولونې} سره ورته ومومئ: لړ د خپلو عناصرو د ښي خوا راټولونکې تابعې په توګه تعريفېږي۔
\end{digress}
\olpseganchor{OLP-0650-B006}{end}

\olpseganchor{OLP-0650-B007}{start}
د دې کوډونې په مرسته څو ګټورې تابعې په اسانۍ تعريفولے شو:
\begin{align*}
  \fn{Sum} & \ident
  \lambd[l][l \, (\lambd[xa][\fn{Add} \, x \, a])\,  \num{0}]\\
  \fn{Len} & \ident
  \lambd[l][l \, (\lambd[xa][\fn{Add} \, a \, \num{1}]) \num{0}]
\end{align*}
\olpseganchor{OLP-0650-B007}{end}

\olpseganchor{OLP-0650-B008}{start}
\LR{$\fn{Sum}$} د چرچ د عددنښو د يوه لړ مجموعه حسابوي۔ په لړ کښې راټولونه کوي: لومړنے ارزښت \LR{$\num{0}$} وي، او د هر عنصر دپاره \LR{$\fn{Add} \, x\, a$} د نوي ارزښت په توګه ورکوي؛ دلته \LR{$x$} اوسنے عنصر او \LR{$a$} تر دې دمه راټول شوے ارزښت دے۔ پايله د ټولو عناصرو مجموعه ده۔
\olpseganchor{OLP-0650-B008}{end}

\olpseganchor{OLP-0650-B009}{start}
\begin{prob}
  \LR{$\fn{Len}$} څۀ کوي؟ تشريح يې کړئ۔
\end{prob}
\olpseganchor{OLP-0650-B009}{end}

\olpseganchor{OLP-0650-B010}{start}
\begin{prob}
  داسې تابعه تعريف کړئ چې د يوه لړ ترتيب سرچپه کړي۔
\end{prob}
\olpseganchor{OLP-0650-B010}{end}

\olpseganchor{OLP-0378-B004}{start}
\olfileid{lam}{ldf}{prf}
\olsection{بنسټيزې بازګشتي تابعې د لامبډا په وسيله د تعريف وړ دي} 
\olpseganchor{OLP-0378-B004}{end}

\olpseganchor{OLP-0378-B005}{start}
په ياد ولرئ چې بنسټيزې بازګشتي تابعې هغه دي چې له
بنسټيزو تابعو \LR{$\Zero$}، \LR{$\Succ$} او \LR{$\Proj{n}{i}$} څخه د
ترکيب او بنسټيز بازګښت په وسيله تعريفېږي۔
\olpseganchor{OLP-0378-B005}{end}

\olpseganchor{OLP-0378-B006}{start}
\begin{lem}
  \ollabel{lem:basic}
  بنسټيزې بازګشتي تابعې \LR{$\Zero$}، \LR{$\Succ$} او
  پروجکشن تابعې~\LR{$\Proj{n}{i}$} د لامبډا په وسيله د تعريف وړ دي۔
\end{lem}
\olpseganchor{OLP-0378-B006}{end}

\olpseganchor{OLP-0378-B007}{start}
\begin{proof}
دا تابعې د لاندې ترمونو په وسيله لامبډا-تعريف کېږي:
\begin{align*}
  \fn{Zero} & \ident \lambd[a][\lambd[fx][x]]\\
  \fn{Succ} & \ident \lambd[a][\lambd[fx][f (a f x)]] \\
  \fn{Proj}^n_i & \ident \lambd[x_0\dots x_{n-1}][x_i]
\end{align*}
\end{proof}
\olpseganchor{OLP-0378-B007}{end}

\olpseganchor{OLP-0378-B008}{start}
\begin{lem}
  \ollabel{lem:comp} فرض کړئ چې \LR{$k$}-ځايه تابع \LR{$f$} او
  \LR{$n$}-ځايه تابعې \LR{$g_0, \dots, g_{k-1}$} په ترتيب سره د
  \LR{$F$}، \LR{$G_0$}، \dots، او \LR{$G_{k-1}$} ترمونو په وسيله
  د لامبډا په وسيله د تعريف وړ دي، او \LR{$h$} له هغو څخه د ترکيب
  په وسيله تعريف شوې ده۔ بيا \LR{$h$} هم د لامبډا په وسيله د تعريف وړ ده۔
\end{lem}
\olpseganchor{OLP-0378-B008}{end}

\olpseganchor{OLP-0378-B009}{start}
\begin{proof}
  تابع \LR{$h$} د لاندې ترم په وسيله لامبډا-تعريف کېږي:
  \[
  H \ident \lambd[x_0 \dots x_{n-1}][F\, (G_0 x_0 \dots
    x_{n-1}) \dots (G_{k-1} x_0 \dots x_{n-1})]
  \]
  د دې خبرې کره کول د تمرين په توګه پرېږدو۔
\end{proof}
\olpseganchor{OLP-0378-B009}{end}

\olpseganchor{OLP-0378-B010}{start}
\begin{prob}
  د \olref[lam][ldf][prf]{lem:comp} ثبوت پوره کړئ؛
  وښيئ چې \LR{$H\num{n_0}\dots\num{n_{n-1}} \red \num{h(n_0, \dots,
    n_{n-1})}$}۔
\end{prob}
\olpseganchor{OLP-0378-B010}{end}

\olpseganchor{OLP-0378-B011}{start}
پام وکړئ چې په \olref{lem:comp} کښې دا نۀ و شرط چې
\LR{$f$} او \LR{$g_0$}، \dots،~\LR{$g_{k-1}$} دې بنسټيزې بازګشتي
وي؛ يوازې دا شرط و چې بشپړې او د لامبډا په وسيله د تعريف وړ وي۔
\olpseganchor{OLP-0378-B011}{end}

\olpseganchor{OLP-0378-B012}{start}
\begin{lem}\ollabel{lem:prim}
  فرض کړئ چې \LR{$f$} يوه \LR{$n$}-ځايه تابع او~\LR{$g$} يوه
  \LR{$n+2$}-ځايه تابع ده، او په ترتيب سره د \LR{$F$} او~\LR{$G$}
  ترمونو په وسيله د لامبډا په وسيله د تعريف وړ دي۔ که تابع~\LR{$h$}
  له \LR{$f$} او \LR{$g$} څخه د بنسټيز بازګښت په وسيله تعريف شوې
  وي، نو \LR{$h$} هم د لامبډا په وسيله د تعريف وړ ده۔
\end{lem}
\olpseganchor{OLP-0378-B012}{end}

\olpseganchor{OLP-0378-B013}{start}
\begin{proof}
  په ياد ولرئ چې \LR{$h$} داسې تعريفېږي:
  \begin{align*}
    h(x_1, \dots, x_n, 0) &= f(x_1, \dots, x_n)\\
    h(x_1, \dots, x_n, y+1) & = g(x_1, \dots, x_n, y, h(x_1, \dots, x_n, y)).
  \end{align*}
په نارسمي ډول، د \LR{$h$}
  ارزښت د \LR{$y$} ګامونو په تېرېدو د ګامي تابع له مخکني
  ارزښت څخه ترلاسه کېږي، او پيلنی ارزښت
  \LR{$f(x_1, \dots, x_n)$} دے۔ د دې تکراري بڼې او د چرچ
  عددنښو د تکراروونکي تعريف تر منځ ورته‌والے شته۔
\olpseganchor{OLP-0378-B013}{end}

\olpseganchor{OLP-0378-B014}{start}
  د اسانۍ دپاره تعريف او ثبوت د يو اضافي آرګومېنټ~\LR{$x$}
  په حالت کښې ورکوو۔ تابع \LR{$h$} د لاندې ترم په وسيله
  لامبډا-تعريف کېږي:
  \begin{align*}
    H \ident & \lambd[x][\lambd[y][\fn{Snd} (y
        D \tuple{\num{0}, F x})]]
    \intertext{\RL{چې پکې }}
    D \ident & \lambd[p][\tuple{\fn{Succ} (\fn{Fst}\, p), (G x
          (\fn{Fst}\, p) (\fn{Snd}\, p))}]
  \end{align*}
  د تکرار ساتلے حالت يوه جوړه ده: لومړے غړے ئې
  د \LR{$y$}-ګامي تکرار اوسنے شمېر او دويم ئې د~\LR{$h$}
  اړوند ارزښت دے۔
  د تکرار په هر ګام کښې د \LR{$y$} او \LR{$h$} له نوو ارزښتونو
  يوه نوې جوړه جوړوو؛ له تکرار وروسته د جوړې دويم
  غړے اخلو، او هماغه د \LR{$h$} وروستنے ارزښت وي۔ اوس
  ثابتوو چې دا په رښتيا د بنسټيز بازګښت تمثيل دے۔
\olpseganchor{OLP-0378-B014}{end}

\olpseganchor{OLP-0378-B015}{start}
  غواړو ثابته کړو چې د هر \LR{$n$} او \LR{$m$} دپاره
  \LR{$H\,\num{n}\,\num{m} \red \num{h(n,m)}$}۔ د دې دپاره
  لومړے ښيو چې که \LR{$D_n \ident \Subst{D}{\num{n}}{x}$}
  وي، نو \LR{$D_n^m \tuple{\num{0}, F\, \num{n}} \red
  \tuple{\num{m}, \num{h(n, m)}}$}۔ پر~\LR{$m$} استقرا کوو۔
\olpseganchor{OLP-0378-B015}{end}

\olpseganchor{OLP-0378-B016}{start}
  که \LR{$m=0$} وي، بايد \LR{$D_n^0 \tuple{\num{0}, F\, \num{n}} \red
  \tuple{\num{0}, \num{h(n, 0)}}$} وښيو۔ خو
  \LR{$D_n^0 \tuple{\num{0}, F\, \num{n}}$} هماغه
  \LR{$\tuple{\num{0}, F\, \num{n}}$} ده۔ څرنګه چې \LR{$F$}
  تابع~\LR{$f$} لامبډا-تعريفوي، دا جوړه
  \LR{$\tuple{\num{0}, \num{f(n)}}$} ته راکمېږي؛ او ځکه
  چې \LR{$f(n) = h(n, 0)$}، دا هماغه
  \LR{$\tuple{\num{0}, \num{h(n,0)}}$} ده۔
\olpseganchor{OLP-0378-B016}{end}

\olpseganchor{OLP-0378-B017}{start}
  اوس فرض کړئ چې \LR{$D_n^m \tuple{\num{0}, F\, \num{n}} \red
  \tuple{\num{m}, \num{h(n, m)}}$}۔ غواړو وښيو چې
  \LR{$D_n^{m+1} \tuple{\num{0}, F\, \num{n}} \red
  \tuple{\num{m+1}, \num{h(n, m+1)}}$}۔
  \begin{align*}
    D_n^{m+1} \tuple{\num{0}, F\, \num{n}}
    & \ident D_n(D_n^{m} \tuple{\num{0}, F\, \num{n}})\\
    & \red D_n\,\tuple{\num{m}, \num{h(n, m)}} \text{\RL{\qquad (د استقرا د فرضيې له مخې)}}\\
    & \ident (\lambd[p][
      \tuple{
        \fn{Succ} (\fn{Fst}\, p), (G \, \num{n} (\fn{Fst}\, p) (\fn{Snd}\, p))
    }]) \tuple{\num{m}, \num{h(n,m)}}\\
    & \redone
    \tuple{\fn{Succ} (\fn{Fst}\, \tuple{\num{m}, \num{h(n,m)}}),\\
      & \qquad (G \, \num{n} (\fn{Fst}\, \tuple{\num{m}, \num{h(n,m)}}) (\fn{Snd}\, \tuple{\num{m}, \num{h(n,m)}}))}\\
    & \red 
    \tuple{\fn{Succ}\, \num{m}, (G \, \num{n} \,\num{m} \, \num{h(n,m)})}\\
    & \red \tuple{\num{m+1}, \num{g(n, m, h(n, m))}}
  \end{align*}
  ځکه چې \LR{$g(n, m, h(n, m)) = h(n, m+1)$}،
  غوښتل شوې پايله ترلاسه شوه۔
\olpseganchor{OLP-0378-B017}{end}

\olpseganchor{OLP-0378-B018}{start}
  په پاې کښې لاندې راکمېدنه وګورئ:
  \begin{align*}
    H\,\num n\,\num m &\ident \lambd[x][\lambd[y][\fn{Snd} (y
        (\lambd[p].\tuple{\fn{Succ} (\fn{Fst}\, p), (G\, x\,
          (\fn{Fst}\, p)\, (\fn{Snd}\, p))}) \tuple{\num{0}, F x})]]\\
    & \qquad\,\num n\, \num m\\
    & \red \fn{Snd} (\num m \,
    \underbrace{(\lambd[p].\tuple{\fn{Succ} (\fn{Fst}\, p), (G \,\num n\,
      (\fn{Fst}\, p) (\fn{Snd}\, p))})}_{D_n} \tuple{\num{0}, F \num n})\\
    & \ident \fn{Snd} (\num m \, D_n\, \tuple{\num{0}, F \num n})\\
    & \red \fn{Snd} \,(D_n^m  \tuple{\num{0}, F \num n})\\
    & \red \fn{Snd} \,\tuple{\num{m}, \num{h(n, m)}} \\
    & \red \num{h(n,m)}. 
  \end{align*}  
\end{proof}
\olpseganchor{OLP-0378-B018}{end}


\olpseganchor{OLP-0378-B019}{start}
\begin{prop}
  هره بنسټيزه بازګشتي تابع د لامبډا په وسيله د تعريف وړ ده۔
\end{prop}
\olpseganchor{OLP-0378-B019}{end}

\olpseganchor{OLP-0378-B020}{start}
\begin{proof}
  د \olref{lem:basic} له مخې ټولې بنسټيزې تابعې
  د لامبډا په وسيله د تعريف وړ دي، او د \olref{lem:comp}
  او \olref{lem:prim} له مخې د لامبډا په وسيله د تعريف وړ
  تابعې د ترکيب او بنسټيز بازګښت پر وړاندې تړلې دي۔
\end{proof}
\olpseganchor{OLP-0378-B020}{end}

\olpseganchor{OLP-0379-B004}{start}
\olfileid{lam}{ldf}{fp}
\olpseganchor{OLP-0379-B004}{end}

\olpseganchor{OLP-0379-B005}{start}
\olsection{ثابت ټکي}
\olpseganchor{OLP-0379-B005}{end}

\olpseganchor{OLP-0379-B006}{start}
فرض کړئ چې د فکتوريال تابع د داسې ترم \LR{$\fn{Fac}$} په توګه
په بازګشتي ډول تعريفول غواړو چې دا خاصيت ولري:
\[
\fn{Fac} \ident \lambd[n][\fn{IsZero}\, n\, \num 1 (\fn{Mult}\, n (\fn{Fac}(\fn{Pred} \, n)))]
\]
يعنې که \LR{$n = 0$} وي، د \LR{$n$} فکتوريال \LR{$1$} دے؛ او که نۀ وي،
د \LR{$n$} او د \LR{$n-1$} د فکتوريال حاصل‌ضرب دے۔ خو ترم
\LR{$\fn{Fac}$} په دې ډول نۀ شو تعريفولے، ځکه \LR{$\fn{Fac}$}
پخپله د تعريف ښي لوري ته راځي۔ د ځان‌مراجعې دغه ډول
بازګشتي تعريفونه د لامبډا حساب د ترمونو له عادي
تعريفونو څخه نۀ دي۔ د بېلګې په توګه، د يو ترم تعريف
داسې کېدے شي:
\[
\fn{Mult} \ident \lambd[ab][a (\fn{Add}\, b) 0]
\]
په ښي لوري کښې يوازې مخکې تعريف شوي ترمونه، لکه
\LR{$\fn{Add}$}، کاروي۔ \LR{$\fn{Add}$} تل د خپل تعريفوونکي
ترم په ايښودلو لرې کولے شو۔ په دې توګه، د عددنښې
لڼدې ليکنې له پرانيستلو سره، \LR{$\fn{Mult}$} د خالص
لامبډا ترم بڼه خپلوي۔ که \LR{$\fn{Add}$} هم نور تعريف شوي
ترمونه کاروي، لکه \LR{$\fn{Add}'$}، دغه پرله‌پسې ځاے پرځاے
کول غځوو او په پاې کښې خالص لامبډا ترم ته رسېږو۔
\olpseganchor{OLP-0379-B006}{end}

\olpseganchor{OLP-0379-B007}{start}
خو د پورته \LR{$\fn{Fac}$} په څېر بازګشتي تعريفونو کښې دا
چاره نۀ ترسره کېږي۔ که په ښي لوري کښې د \LR{$\fn{Fac}$}
پېښه پخپله د \LR{$\fn{Fac}$} په تعريف بدله کړو، لاندې
عبارت ترلاسه کېږي:
\begin{align*}
  \fn{Fac} & \ident \lambd[n].\bigl(\fn{IsZero}\, n\, \num{1} \\
    & \qquad (\fn{Mult}\, n
  ((\lambd[n][\fn{IsZero} \, n \, \num 1\, (\fn{Mult}\, n\, (\fn{Fac}
    (\fn{Pred}\, n)))]) (\fn{Pred}\, n)))\bigr)
\end{align*}
په دې ښودنيزه پرانيستنه کښې دواړه څانګې د بهرنۍ لامبډايي تجريد په بدن کښې دي۔
او \LR{$\fn{Fac}$} لا هم په ښي لوري کښې پاتې دے۔ که همدا
کار بيا بيا وکړو، عبارت اوږدېږي او هېڅکله خالص
لامبډا ترم نۀ راکوي۔ نو د فکتوريال، يا په عام ډول د
بازګشتي تابعو، تعريف په دې لاره ممکن نۀ دے۔
\olpseganchor{OLP-0379-B007}{end}

\olpseganchor{OLP-0379-B008}{start}
خو بازګشتي تعريف يوه ګټوره خبره راښيي: که \LR{$f$} د
فکتوريال تابع تمثيلوونکے ترم وي، نو لاندې ترم
\[
\fn{Fac}' \ident \lambd[g][\lambd[n][\fn{IsZero} \, n \, \num 1\, (\fn{Mult} \, n\, (g (\fn{Pred} n)))]]
\]
چې پر \LR{$f$} تطبيق شي، يعنې \LR{$\fn{Fac}'\,f$}، هم د
فکتوريال تابع تمثيل کوي۔ که \LR{$\fn{Fac}'$} د يوې تابعې
اخيستونکې او بله تابع ورکوونکې عمليې په توګه وګورو،
د فکتوريال پر يوه تمثيل ئې پايله بيا د هماغې تابعې
تمثيل وي۔ داسې ترم~\LR{$f$} چې \LR{$\fn{Fac}' \, f \equal[\beta] f$}
خاصيت ولري، د \LR{$\fn{Fac}'$} \emph{ثابت ټکے} بلل کېږي۔
په عددنښو کښې \LR{$\fn{Fac'}\, f$} هماغه ارزښتونه ورکوي چې \LR{$f$} ئې ورکوي، که \LR{$f$} د فکتوريال دپاره وي۔ د \LR{$\fn{Fac}'$} ثابت ټکے د بېټا-برابرۍ په هماغه شرط ټاکل کېږي۔
د يوې عددي تابعې دوه تمثيلونه په ټولو عددنښو کښې يو شان ارزښتونه ورکوي۔ د ثابت ټکي دپاره د ترمونو بېټا-برابري هم پکار ده۔
\olpseganchor{OLP-0379-B008}{end}

\olpseganchor{OLP-0379-B009}{start}
په لامبډا حساب کښې داسې ترمونه شته چې د ورکړل شوي
ترم ثابت ټکي پيدا کوي۔ په هغو سره د \LR{$\fn{Fac}'$} په
څېر ترم د فکتوريال په تعريف بدلولے شو۔
\olpseganchor{OLP-0379-B009}{end}

\olpseganchor{OLP-0379-B010}{start}
\begin{defn}
  \ollabel{defn:Turing-Y}
  \emph{Y-ترکيب کوونکے} دا ترم دے:
  \[
  Y \ident (\lambd[ux][x(uux)])(\lambd[ux][x(uux)]).
  \]
\end{defn}
\olpseganchor{OLP-0379-B010}{end}

\olpseganchor{OLP-0379-B011}{start}
\begin{thm}
  د هر ترم \LR{$g$} دپاره \LR{$Y$} دا خاصيت لري چې \LR{$Yg \red g(Yg)$}۔
  نو \LR{$Yg$} تل د~\LR{$g$} ثابت ټکے دے۔
\end{thm}
\olpseganchor{OLP-0379-B011}{end}

\olpseganchor{OLP-0379-B012}{start}
\begin{proof}
  \LR{$(\lambd[ux][x(uux)])$} په~\LR{$U$} لنډوو، نو \LR{$Y \ident UU$}۔ بيا
  \begin{align*}
    Y g &\ident (\lambd[ux][x(uux)])U\,g \\
    &\red (\lambd[x][x(UUx)])g\\
    & \red g(UUg) \ident g(Yg).
  \end{align*}
  څرنګه چې \LR{$g(Yg)$} او \LR{$Yg$} دواړه \LR{$g(Yg)$} ته
  راکمېږي، نو \LR{$g(Yg) \equal[\beta] Yg$}؛ له دې امله
  \LR{$Yg$} د~\LR{$g$} ثابت ټکے دے۔
\end{proof}
\olpseganchor{OLP-0379-B012}{end}

\olpseganchor{OLP-0379-B013}{start}
څرنګه چې \LR{$Yg$} راکمېدونکے عبارت دے، راکمول ئې
بې‌پاېه دوام موندلے شي:
\begin{align*}
  Y g &\red g (Y g) \\
    &\red g (g (Y g)) \\
    &\red g(g (g (Y g)))\\
    &\ldots
\end{align*}
نو \LR{$Yg$} داسې انګېرلے شو لکه \LR{$g$} چې بې‌شمېره ځله
پر خپل تېر حاصل تطبيقېږي۔ که \LR{$g$} پر \LR{$Yg$} يو ځل
نور هم تطبيق کړو، په دې شهودي انځور کښې ګواکې
څۀ نۀ زياتېږي: \LR{$g$} پر \LR{$g$} تطبيقېږي، خو بيا هم
د \LR{$Yg$} د بې‌شمېره تکرار هماغه لړۍ ده؛ د \LR{$g$} يوه
نوره کارونه پکې ورکېږي۔ دا يوازې شهودي بيان دے؛
پورته کره دعوه د راکمېدنې او بېټا-برابرۍ ده۔
\olpseganchor{OLP-0379-B013}{end}

\olpseganchor{OLP-0379-B014}{start}
پام وکړئ چې له~\LR{$Yg$} څخه پيلېدونکې د \LR{$\beta$}-راکمول
پورته لړۍ بې‌پاېه ده۔ که \LR{$Yg$} پر کوم ترم تطبيق کړو،
يعنې \LR{$(Yg)N$} وګورو، دا هم بې‌شمېره بېلابېلو ترمونو
ته راکمېدے شي: \LR{$(g(Yg))N$}، \LR{$(g(g(Yg)))N$}، \dots۔
سره له دې، کېدے شي د راکمولو کومه \emph{بله} لړۍ
په عادي بڼه پاے ته ورسېږي۔
\olpseganchor{OLP-0379-B014}{end}

\olpseganchor{OLP-0379-B015}{start}
د فکتوريال بېلګه واخلئ۔ \LR{$\fn{Fac}$} د \LR{$Y\, \fn{Fac}'$}
په توګه تعريف کړئ، يعنې د \LR{$\fn{Fac'}$} يو ثابت ټکے۔ بيا:
\begin{align*}
  \fn{Fac}\,\num 3
  &\red Y\, \fn{Fac}' \, \num 3 \\
  &\red \fn{Fac}' (Y \,\fn{Fac}') \, \num 3 \\
  &\ident (\lambd[x][\lambd[n][\fn{IsZero} \, n\, \num 1\,
      (\fn{Mult}\, n\, (x (\fn{Pred}\, n)))]]) \, \fn{Fac} \, \num 3\\
  & \red \fn{IsZero} \, \num 3\, \num 1\,
      (\fn{Mult}\, \num 3\, (\fn{Fac} (\fn{Pred}\, \num 3))) \\
  &\red  \fn{Mult}\, \num 3 \, (\fn{Fac} \, \num 2).
  \intertext{\RL{په همدې ډول،}}
  \fn{Fac}\,\num 2
  &\red  \fn{Mult}\, \num 2 \, (\fn{Fac} \, \num 1) \\
  \fn{Fac}\,\num 1
  &\red  \fn{Mult}\, \num 1 \, (\fn{Fac} \, \num 0)
  \intertext{\RL{خو}}
  \fn{Fac}\,\num 0
  &\red \fn{Fac}' (Y \,\fn{Fac}') \, \num 0 \\
  &\ident (\lambd[x][\lambd[n][\fn{IsZero} \, n\, \num 1\,
      (\fn{Mult}\, n\, (x (\fn{Pred}\, n)))]]) \, \fn{Fac} \, \num 0\\
  & \red \fn{IsZero} \, \num 0\, \num 1\,
      (\fn{Mult}\, \num 0\, (\fn{Fac} (\fn{Pred}\, \num 0))).\\
  &\red \num 1.
  \intertext{\RL{نو په ټوليز ډول}}
  \fn{Fac}\,\num 3
  &\red  \fn{Mult}\, \num 3 \,
  (\fn{Mult} \, \num 2\, (\fn{Mult} \, \num 1\, \num 1)).
\end{align*}
\olpseganchor{OLP-0379-B015}{end}

\olpseganchor{OLP-0379-B016}{start}
د \LR{$\fn{Fac'}$} په اړه دا چلند د هر بازګشتي تعريف
دپاره هم کار کوي۔ فرض کړئ لاندې بازګشتي معادله لرو:
\begin{align*}
g\,x_1\dots x_n & \equal[\beta] N
\intertext{\RL{چې \LR{$N$} ښايي~\LR{$g$} او \LR{$x_1$}، \dots،~\LR{$x_n$} ولري۔ بيا تل داسې
ترم~\LR{$G \ident (Y \lambd[g][\lambd[x_1\dots x_n][N]])$} شته چې}}
G\,x_1 \dots x_n & \equal[\beta] \Subst{N}{G}{g}.
\intertext{\RL{ځکه د ثابت ټکي د قضيې له مخې،}}
G \ident (Y \lambd[g][\lambd[x_1\dots x_n][N]]) & \red \lambd[g][\lambd[x_1\dots x_n][N]](Y \lambd[g][\lambd[x_1\dots x_n][N]])\\
& \ident (\lambd[g][\lambd[x_1\dots x_n][N]])\,G
\intertext{\RL{او په پايله کښې}}
G\,x_1 \dots x_n 
& \red (\lambd[g][\lambd[x_1\dots x_n][N]])\,G\,x_1\dots x_n\\
& \red (\lambd[x_1\dots x_n][\Subst{N}{G}{g}])\,x_1\dots x_n\\
  & \red \Subst{N}{G}{g}.
\end{align*}
\olpseganchor{OLP-0379-B016}{end}

\olpseganchor{OLP-0379-B017}{start}
د \olref{defn:Turing-Y} د \LR{$Y$} ترکيب کوونکے د الن
تورينګ دے۔ الونزو چرچ بله بڼه وړاندې کړې وه؛
هغې ته~\LR{$Y_C$} وايو:
\[
Y_C \ident \lambd[g][(\lambd[x][g(xx)])(\lambd[x][g(xx)])].
\]
د چرچ ترکيب کوونکے په يوه معنا د تورينګ له بڼې
څخه کمزورے دے: \LR{$Y_Cg \equal[\beta] g(Y_Cg)$}، خو
\LR{$Y_Cg \bred g(Y_Cg)$} نۀ دے۔ د دې فرق د کتلو لپاره \LR{$g$} يو متغير واخلئ۔ \LR{$V$} دې د
\LR{$\lambd[x][g(xx)]$} ترم وي، نو \LR{$Y_C \ident \lambd[g][VV]$}۔ بيا
\begin{align*}
  VV & \ident (\lambd[x][g(xx)])V \red g(VV)
  \text{\RL{ او له دې امله}}\\
  Y_C g & \ident (\lambd[g][VV])g \red VV \red g(VV), \text{\RL{ خو همدارنګه}}\\
  g(Y_C g) & \ident g((\lambd[g][VV])g) \red g(VV).
\end{align*}
په بله وينا، \LR{$Y_Cg$} او \LR{$g(Y_Cg)$} دواړه ګډ ترم
\LR{$g(VV)$} ته راکمېږي؛ نو \LR{$Y_Cg \equal[\beta] g(Y_Cg)$}۔
د ډېرو کارونو دپاره همدا بس وي۔
\olpseganchor{OLP-0379-B017}{end}

\olpseganchor{OLP-0380-B004}{start}
\olfileid{lam}{ldf}{min}
\olsection{کمينه‌موندنه}
\olpseganchor{OLP-0380-B004}{end}

\olpseganchor{OLP-0380-B005}{start}
عامې بازګشتي تابعې هغه دي چې له بنسټيزو تابعو
\LR{$\Zero$}، \LR{$\Succ$} او \LR{$\Proj{n}{i}$} څخه د ترکيب،
بنسټيز بازګښت او منظمې کمينه‌موندنې په وسيله
ترلاسه کېږي۔ د دې د ښودلو دپاره چې ټولې عامې
بازګشتي تابعې د لامبډا په وسيله د تعريف وړ دي، بايد ثابته
کړو چې له يوې د لامبډا په وسيله د تعريف وړ تابعې څخه د منظمې کمينه‌موندنې په
وسيله تعريف شوې تابع هم د لامبډا په وسيله د تعريف وړ ده۔
\olpseganchor{OLP-0380-B005}{end}

\olpseganchor{OLP-0380-B006}{start}
\begin{lem}
  \ollabel{lem:min} که \LR{$f(x_1, \dots, x_k, y)$} منظمه او
  د لامبډا په وسيله د تعريف وړ وي، نو هغه \LR{$h$} هم
  د لامبډا په وسيله د تعريف وړ ده چې داسې تعريفېږي:
  \[
  h(x_1, \dots, x_k) = \umin{y}{f(x_1,\dots,x_k, y) = 0}
  \]
\end{lem}
\olpseganchor{OLP-0380-B006}{end}

\olpseganchor{OLP-0380-B007}{start}
\begin{proof}
  فرض کړئ لامبډا ترم~\LR{$F$} منظمه تابع
  \LR{$f(\vec x, y)$} په \LR{$\lambda$} سره تعريفوي۔ د~\LR{$h$}
  د لامبډا-تعريف دپاره يوه لټونکې تابع او د ثابت
  ټکي ترکيب کوونکے کاروو:
  \begin{align*}
    \fn{Search} & \ident \lambd[g][\lambd[f\,\vec{x}\,y][
        \fn{IsZero} (f\, \vec{x}\, y)\, y\, (g\, f\, \vec{x} (\fn{Succ}\, y))]]\\
    H & \ident \lambd[\vec x][(Y \, \fn{Search}) F\, \vec{x}\, \num{0}],
  \end{align*}
  چې \LR{$Y$} پکې د مخکنۍ برخې د ټيورېنګ د ثابت ټکي ترکيب کوونکے دے۔ د هغه مخامخ راکمېدنه د لټون بل ګام هماغه بيا بلنې ته رسوي۔
  په نارسمي ډول، \LR{$\fn{Search}$} ځان‌مراجعه کوونکې
  تابع ده: له~\LR{$y$} څخه پيل کوي او ګوري چې آيا
  \LR{$f\, \vec x\, y$} صفر دے که نۀ؛ که وي،~\LR{$y$} ورکوي،
  او که نۀ وي، ځان له \LR{$\fn{Succ}\, y$} سره بيا
  رابولي۔ په بيا بلنه کښې ازمېيل کېدونکې تابع او ټول ټاکلي پارامترونه هماغه پاتې کېږي؛ يوازې د لټون شمېر يو زياتېږي۔ نو \LR{$(Y \, \fn{Search}) F
  \num{n_1}\dots\num{n_k}\,\num{0}$} هغه تر ټولو
  وړوکے~\LR{$m$} ورکوي چې \LR{$f(n_1,
  \dots, n_k, m) = 0$} وي۔
  د کمينه‌موندنې پايله هماغه تابع ده چې لمّه ئې بيانوي؛ د لټون دننه تړلے متغير د بيا بلنې دپاره کارېږي۔
\olpseganchor{OLP-0380-B007}{end}
  
\olpseganchor{OLP-0380-B008}{start}
  په ځانګړي ډول، پام وکړئ چې
  \begin{align*}
    (Y \, \fn{Search}) F \num{n_1}
    \dots\num{n_k}\, \num{m} & \red \num{m}
    \intertext{\RL{که \LR{$f(n_1, \dots,
      n_k, m) = 0$} وي؛ او کنه}}
    & \red (Y \, \fn{Search}) F\, \num{n_1} \dots\num{n_k}\, \num{m+1}
    \intertext{\RL{که نۀ وي۔ څرنګه چې \LR{$f$} منظمه ده،
      \LR{$f(n_1, \dots, n_k, y) = 0$} د کوم \LR{$y$} دپاره کېږي؛ نو}}
    (Y \, \fn{Search}) F \num{n_1} \dots\num{n_k}\,\num{0}
    & \red \num{h(n_1, \dots, n_k)}.
    \end{align*}
\end{proof}
\olpseganchor{OLP-0380-B008}{end}


\olpseganchor{OLP-0380-B009}{start}
\begin{prop}
  هره عامه بازګشتي تابع د لامبډا په وسيله د تعريف وړ ده۔
\end{prop}
\olpseganchor{OLP-0380-B009}{end}

\olpseganchor{OLP-0380-B010}{start}
\begin{proof}
 د \olref[prf]{lem:basic} له مخې ټولې بنسټيزې تابعې
 د لامبډا په وسيله د تعريف وړ دي، او د \olref[prf]{lem:comp}،
 \olref[prf]{lem:prim} او \olref{lem:min} له مخې
 د لامبډا په وسيله د تعريف وړ تابعې د ترکيب، بنسټيز بازګښت
 او منظمې کمينه‌موندنې پر وړاندې تړلې دي۔
\end{proof}
\olpseganchor{OLP-0380-B010}{end}

\olpseganchor{OLP-0381-B004}{start}
\olfileid{lam}{ldf}{par}
\olsection{جزوي بازګشتي تابعې د لامبډا په وسيله د تعريف وړ دي}
\olpseganchor{OLP-0381-B004}{end}

\olpseganchor{OLP-0381-B005}{start}
جزوي بازګشتي تابعې له بنسټيزو تابعو څخه د ترکيب،
بنسټيز بازګښت او نامحدودې کمينه‌موندنې په وسيله
ترلاسه کېږي۔ له عامو بازګشتي تابعو سره ئې توپير
دا دے چې په بې‌حده پلټنه کښې کارېدونکې تابعې
لازم نۀ دي چې منظمې وي۔ د منظم‌والي د شرط نشتوالے
په دې معنا دے چې د کمينه‌موندنې په وسيله تعريف شوې
تابعه کله ناکله ناتعريفه کېدے شي۔
\olpseganchor{OLP-0381-B005}{end}

\olpseganchor{OLP-0381-B006}{start}
په لومړي نظر ښايي داسې وبرېښي چې د ټولو هرځاے
تعريف شوو عامو بازګشتي تابعو د د لامبډا په وسيله د تعريف وړ
والي د ثبوت هماغه لارې د ټولو جزوي بازګشتي تابعو
د د لامبډا په وسيله د تعريف وړ والي د ثبوت دپاره هم کارېدے شي۔
د بېلګې په توګه، که \LR{$f$} او \LR{$g$} په ترتيب سره د \LR{$F$}
او~\LR{$G$} ترمونو په وسيله لامبډا-تعريف وي،
نو د \LR{$f$} او~\LR{$g$} ترکيب د
\LR{$\lambd[x][F (G x)]$} په وسيله لامبډا-تعريف
دے۔ خو کله چې تابعې جزوي وي، ستونزه راپيدا کېږي۔
که \LR{$g(x)$} ناتعريفه وي، د \LR{$G x$} عادي بڼه نۀ وي۔
په ډېرو حالاتو کښې له دې امله \LR{$F (G x)$} هم عادي
بڼه نۀ لري؛ همدا مو غوښتل۔ خو که \LR{$F$} د
\LR{$\lambd[x][\num{0}]$} ترم وي، بيا \LR{$F (G x)$}
عادي بڼه لري: \LR{$\num{0}$}۔
\olpseganchor{OLP-0381-B006}{end}

\olpseganchor{OLP-0381-B007}{start}
دا ستونزه د حل وړ ده۔ د ټولو جزوي بازګشتي تابعو
د لامبډا-تعريف داسې لارې شته چې ناتعريفه
ارزښتونه ئې په هغو ترمونو تمثيلېږي چې عادي بڼه نۀ لري۔ خو دا
لارې د هغو لارو په پرتله چې د عامو بازګشتي تابعو
دپاره مو کارولې، څه پېچلې او لږې شهودي دي۔ دا قضيه
دلته بې له ثبوته ثبتوو:
\olpseganchor{OLP-0381-B007}{end}

\olpseganchor{OLP-0381-B008}{start}
\begin{thm}
  ټولې جزوي بازګشتي تابعې د لامبډا په وسيله د تعريف وړ دي۔
\end{thm}
\olpseganchor{OLP-0381-B008}{end}

\olpseganchor{OLP-0382-B004}{start}
\olfileid{lam}{dfl}{ldr}
\olsection{د لامبډا په وسيله د تعريف وړ تابعې بازګشتي دي}
\olpseganchor{OLP-0382-B004}{end}

\olpseganchor{OLP-0382-B005}{start}
نۀ يوازې ټولې جزوي بازګشتي تابعې د لامبډا په وسيله د تعريف وړ
دي، بلکې برعکس خبره هم رښتيا ده: ټولې
د لامبډا په وسيله د تعريف وړ تابعې جزوي بازګشتي دي۔
\olpseganchor{OLP-0382-B005}{end}

\olpseganchor{OLP-0382-B006}{start}
\begin{thm}
  \ollabel{thm:lambda-computable} که يوه جزوي تابع~\LR{$f$}
  د لامبډا په وسيله د تعريف وړ وي، نو جزوي بازګشتي هم ده۔
\end{thm}
\olpseganchor{OLP-0382-B006}{end}

\olpseganchor{OLP-0382-B007}{start}
\begin{proof}
  د ثبوت يوازې لنډه طرحه وړاندې کوو۔ لومړے د
  \LR{$\lambd$}-ترمونو حسابي کول کوو: د لړيو د کوډولو
  دپاره د اوليو عددونو د توانونو په معمول طريقه هر
  \LR{$\lambd$}-ترم ته په منظم ډول د ګوډل شمېره ورکوو۔
  بيا يوه جزوي بازګشتي تابع \LR{$\fn{normalize}(t)$}
  تعريفوو چې د لامبډا ترم د ګوډل شمېره~\LR{$t$} د
  آرګومېنټ په توګه اخلي؛ که ترم عادي بڼه ولري، د
  هغې بڼې د ګوډل شمېره ورکوي، او که نۀ وي،
  ناتعريفه پاتې کېږي۔ ورپسې دوه جزوي بازګشتي
  تابعې \LR{$\fn{toChurch}$} او \LR{$\fn{fromChurch}$}
  تعريفوو چې طبيعي عددونه د اړوندو چرچ عددنښو
  د ګوډل شمېرو ته اړوي او له هغو شمېرو ئې بېرته
  طبيعي عددونو ته رااړوي۔
\olpseganchor{OLP-0382-B007}{end}

\olpseganchor{OLP-0382-B008}{start}
  له دغو بازګشتي تابعو څخه په کار اخيستو، تابع~\LR{$f$}
  د جزوي بازګشتي تابعې په توګه تعريفولے شو۔ داسې
  \LR{$\lambd$}-ترم~\LR{$F$} شته چې~\LR{$f$} لامبډا-تعريفوي۔
  د \LR{$f(n_1, \dots, n_k)$} د محاسبې دپاره لومړے د
  \LR{$\fn{toChurch}(n_i)$} په وسيله د اړوندو چرچ
  عددنښو د ګوډل شمېرې ترلاسه کوو۔ بيا دا شمېرې
  له \LR{$\Gn{F}$} سره يوځاے کوو، څو د
  \LR{$F \num{n_1}\dots\num{n_k}$} ترم د ګوډل شمېره
  ترلاسه شي۔ اوس پر دې ګوډل شمېره
  \LR{$\fn{normalize}$} کاروو۔ که \LR{$f(n_1, \dots,
  n_k)$} تعريف شوې وي، \LR{$F \num{n_1}\dots\num{n_k}$}
  عادي بڼه لري، او هغه بايد د چرچ عددنښه وي۔ که
  تابع ناتعريفه وي، ترم عادي بڼه
  نۀ لري؛ نو
  \[\fn{normalize}(\Gn{F\num{n_1}\dots\num{n_k}})\]
  هم ناتعريفه وي۔ په پاې کښې د عادي شوې بڼې د
  ګوډل پر شمېره \LR{$\fn{fromChurch}$} کاروو۔
\end{proof}
\olpseganchor{OLP-0382-B008}{end}

\olpseganchor{OLP-0383-B004}{start}
\olpart{mvl}{څو ارزښته منطق}
\olpseganchor{OLP-0383-B004}{end}

\olpseganchor{OLP-0383-B005}{start}
\begin{editorial}
  په دې برخه کښې د قضيوي څو ارزښته منطقونو د مسودې مطالب راغلي دي۔
\end{editorial}
\olpseganchor{OLP-0383-B005}{end}

\olpseganchor{OLP-0384-B004}{start}
\olchapter{mvl}{syn}{نحو او معناپوهنه}
\olpseganchor{OLP-0384-B004}{end}

\olpseganchor{OLP-0385-B004}{start}
\olfileid{mvl}{syn}{int}
\olpseganchor{OLP-0385-B004}{end}

\olpseganchor{OLP-0385-B005}{start}
\olsection{پېژندنه}
\olpseganchor{OLP-0385-B005}{end}

\olpseganchor{OLP-0385-B006}{start}
په کلاسيک منطق کښې له هغو فارمولونه سره کار
کوو چې له بياني متغيرونه څخه د
قضيوي رابطونو \LR{$\lnot$}، \LR{$\land$}، \LR{$\lor$}، \LR{$\lif$}
او \LR{$\liff$} په وسيله جوړېږي۔ د کلاسيک منطق د
معناپوهنې د تعريف دپاره د رښتياوالي دوه قيمتونه \LR{$\True$}
او \LR{$\False$} کاروو۔ بياني متغيرونه
په ارزښت ټاکنه~\LR{$\pAssign{v}$} کښې تفسيروو؛
دا له بياني متغيرونه څخه هر يوه ته له دغو دوو څخه يو قيمت،
\LR{$\True$} يا \LR{$\False$}، ورکوي۔ بيا هره
ارزښت ټاکنه د هر فارمول~\LR{$!A$} د رښتياوالي
قيمت \LR{$\pValue{v}(!A)$} ټاکي؛ او فارمول
په ارزښت ټاکنه~\LR{$\pAssign{v}$} کښې هله پوره
کېږي، \LR{$\pSat{v}{!A}$}، چې \LR{$\pValue{v}(!A) =
\True$} وي۔
\olpseganchor{OLP-0385-B006}{end}

\olpseganchor{OLP-0385-B007}{start}
څو ارزښته منطقونه د کلاسيک دوه ارزښته منطق
عموميتونه دي: د \LR{$\True$} او \LR{$\False$} تر څنګ
د رښتياوالي نور قيمتونه هم مني۔ په څو ارزښته منطق کښې
ارزښت ټاکنه~\LR{$\pAssign{v}$} داسې تابع ده چې
هر بياني متغير~\LR{$p$} ته د رښتياوالي
د منلو وړ قيمتونو له سټه يو قيمت ورکوي۔
د رښتياوالي د منل شوو قيمتونو سټ ته به عموماً~\LR{$V$}
وايو۔ کلاسيک منطق هم يو څو ارزښته منطق دے چې
\LR{$V = \{\True, \False\}$} لري؛ د فارمول د رښتياوالي
قيمت~\LR{$\pValue{v}(!A)$} د رابطونو د پېژندل شوو
ځانګړنو د رښتياوالي جدولونو له مخې محاسبه کېږي۔
\olpseganchor{OLP-0385-B007}{end}

\olpseganchor{OLP-0385-B008}{start}
کله چې د رښتياوالي نور قيمتونه ور زيات کړو، د هغو
رابطونو د \LR{$\pValue{v}(!A)$} د محاسبې دپاره چې
«او»، «يا»، «نۀ» او «که---نو» ئې لولو، له يوې
څخه زياتې طبيعي لارې شته۔ د څو ارزښته منطق د ټاکلو دپاره
هم د رښتياوالي د قيمتونو سټ او هم د هرې تړنې \emph{د صدق تابع}
ټاکو۔
خو چې د يوه څو ارزښته منطق~\LR{$\Log L$} دپاره دغه
تابعې وټاکل شي، د رښتياوالي قيمت
\LR{$\pValue{v}(!A)[\Log L]$} د ارزښت ټاکنې له مخې
په يوازيني ډول معلومېږي، لکه په کلاسيک منطق کښې۔
څو ارزښته منطقونه هم، د کلاسيک منطق په څېر،
\emph{صدق‌تابع} دي۔
\olpseganchor{OLP-0385-B008}{end}

\olpseganchor{OLP-0385-B009}{start}
له دې معنايي بنسټونو څخه د تاوتولوژۍ،
استلزام او د صدق وړتيا د معنايي مفهومونو سيالې
بڼې تعريفولے شو۔ په کلاسيک منطق کښې
فارمول هله تاوتولوژي وي چې د هرې
\LR{$\pAssign{v}$} دپاره ئې د رښتياوالي قيمت
\LR{$\pValue{v}(!A) = \True$} وي۔ په څو ارزښته
منطق کښې دا خبره لږ عمومي کوو۔ لومړے، لازمه نۀ
ده چې د رښتياوالي د قيمتونو سټ~\LR{$V$} پخپله \LR{$\True$}
ولري۔ د بېلګې په توګه، ځينې څو ارزښته منطقونه
عددونه کاروي: له \LR{$0$} څخه تر~\LR{$1$} پورې ټول ناطق
عددونه د رښتياوالي د قيمتونو سټ ګڼي۔ په دې حالت کښې
\LR{$1$} عموماً د~\LR{$\True$} رول لوبوي؛ په ځينو نورو
منطقونو کښې څو قيمتونه دا رول لري۔ نو له هر څو
ارزښته منطق سره د \emph{ټاکل شوو قيمتونو} يو سټ~\LR{$V^+$} اخلو۔ بيا ويلے شو چې
فارمول په ارزښت ټاکنه~\LR{$\pAssign{v}$}
کښې هله پوره کېږي، \LR{$\pSat{v}{!A}[\Log L]$}،
چې \LR{$\pValue{v}(!A)[\Log L] \in V^+$} وي۔
فارمول~\LR{$!A$} د دغه منطق تاوتولوژي،
\LR{$\Entails[\Log L] !A$}، هله ده چې د هرې
\LR{$\pAssign{v}$} دپاره \LR{$\pValue{v}(!A) \in V^+$}
وي۔ په پاې کښې وايو چې د فارمولونه سټ
د \LR{$!A$} استلزام کوي، \LR{$\Gamma \Entails[\Log L] !A$}،
که هره هغه ارزښت ټاکنه چې په~\LR{$\Gamma$} کښې
ټول فارمولونه پوره کوي، \LR{$!A$} هم پوره کوي۔
\olpseganchor{OLP-0385-B009}{end}

\olpseganchor{OLP-0386-B004}{start}
\olfileid{mvl}{syn}{con}
\olpseganchor{OLP-0386-B004}{end}

\olpseganchor{OLP-0386-B005}{start}
\olsection{ژبې او نښلوونکي}
\olpseganchor{OLP-0386-B005}{end}

\olpseganchor{OLP-0386-B006}{start}
کلاسيک بياني منطق، او ډېر نور منطقونه، د \emph{بياني ثابتونو} او
\emph{نښلوونکو} يو سټ کاروي۔ دلته لاندې نښې بنسټيزې اخلو:
\begin{enumerate}
\item د دروغوالی بياني ثابت~\LR{$\lfalse$}۔
\item د رښتياوالی بياني ثابت~\LR{$\ltrue$}۔
\item منطقي نښلوونکي:
  \startycommalist
  \ycomma \LR{$\lnot$} (نفي)%
  \ycomma \LR{$\land$} (عطف)%
  \ycomma \LR{$\lor$} (فصل)%
  \ycomma \LR{$\lif$} (شرطي)%
  \ycomma \LR{$\liff$} (دوه اړخيز شرطي)%
\end{enumerate}

\olpseganchor{OLP-0386-B006}{end}

\olpseganchor{OLP-0386-B007}{start}
په څو ارزښته منطقونو کښې هم همدغه نښلوونکي کارېږي۔ خو ډېر ځله
ګټوره وي چې د بېلګې په توګه د عطف بېلابېلې بڼې په يوه منطق کښې
شاملې کړو؛ بيا د دغو بڼو د بېلولو دپاره بېلو نښو ته اړتيا وي۔
ځينې څو ارزښته منطقونه داسې نښلوونکي هم لري چې په کلاسيک منطق
کښې ئې سيال نشته۔ له همدې امله دلته له معمول څخه لږ عمومي
تعريف کاروو۔
\olpseganchor{OLP-0386-B007}{end}

\olpseganchor{OLP-0386-B008}{start}
\begin{defn}
  يوه \emph{بياني ژبه} د \emph{نښلوونکو} له يوه سټ~\LR{$\Lang L$} څخه
  جوړه وي۔ هر نښلوونکے~\LR{$\star$} يو ټاکلے شمېر ځايونه لري؛ هغه
  نښلوونکے چې~\LR{$n$} ځايونه لري، \emph{\LR{$n$}-ځايه} دے۔
  د~\LR{$0$} ځايونو نښلوونکو ته \emph{ثابتونه} هم ويل کېږي؛
  د~\LR{$1$} ځاے نښلوونکي \emph{يو ځايه}، او د~\LR{$2$} ځايونو
  نښلوونکي \emph{دوه ځايه} بلل کېږي۔
\end{defn}
\olpseganchor{OLP-0386-B008}{end}

\olpseganchor{OLP-0386-B009}{start}
\begin{ex}
  د بياني منطق معياري ژبه~\LR{$\Lang L_0$} له لاندې نښلوونکو
  (او د هر يوه له ټاکلي شمېر ځايونو) څخه جوړه ده:
  \LR{$\lfalse$}~(\LR{$0$})،
  \LR{$\lnot$}~(\LR{$1$})،
  \LR{$\land$}~(\LR{$2$})،
  \LR{$\lor$}~(\LR{$2$})،
  \LR{$\lif$}~(\LR{$2$})۔ ډېری هغه منطقونه چې موږ ئې څېړو، همدا ژبه
  کاروي۔ ځينې منطقونه د دود له مخې د ځينو نښلوونکو دپاره
  بېلې نښې کاروي۔ د بېلګې په توګه، په حاصل‌ضرب منطق کښې
  د عطف نښه ډېر ځله د~\LR{$\land$} پر ځاے~\LR{$\odot$} وي۔ کله ناکله
  د نوي عملګر زياتول هم ګټور وي، لکه د ټاکلتيا عملګر~\LR{$\triangle$}
  (\LR{$1$}-ځايه)۔
\end{ex}
\olpseganchor{OLP-0386-B009}{end}

\olpseganchor{OLP-0387-B004}{start}
\olfileid{mvl}{syn}{fml}
\olpseganchor{OLP-0387-B004}{end}

\olpseganchor{OLP-0387-B005}{start}
\olsection{فارمولونه}
\olpseganchor{OLP-0387-B005}{end}

\olpseganchor{OLP-0387-B006}{start}
\begin{defn}[فارمول]
\ollabel{defn:formulas}
د بياني ژبې~\LR{$\Lang L$} د \emph{فارمولونه} سټ~\LR{$\Frm[L]$}
په استقرايي ډول داسې تعريفېږي:
\begin{enumerate}
\item هر بياني متغير~\LR{$\Obj p_i$} يو اټومي
  فارمول دے۔
\item د~\LR{$\Lang L$} هر \LR{$0$}-ځايه نښلوونکے (بياني ثابت)
  يو اټومي فارمول دے۔
\item که~\LR{$\star$} د~\LR{$\Lang L$} يو \LR{$n$}-ځايه نښلوونکے وي چې لږ تر لږه يو ځاے لري، او
  \LR{$!A_1$}، \dots، \LR{$!A_n$} فارمولونه وي، نو
  \LR{$\star(!A_1, \dots, !A_n)$} هم فارمول دے۔
\item له دې پرته هېڅ څيز فارمول نۀ دے۔
\end{enumerate}
که~\LR{$\star$} يو \LR{$1$}-ځايه نښلوونکے وي، نو
\LR{$\star(!A_1)$} به ډېر ځله په ساده ډول \LR{$\star !A_1$} ليکو۔
که~\LR{$\star$} يو \LR{$2$}-ځايه نښلوونکے وي، نو
\LR{$\star(!A_1,!A_2)$} به ډېر ځله د \LR{$(!A_1 \star !A_2)$}
په بڼه ليکو۔
\end{defn}
\olpseganchor{OLP-0387-B006}{end}

\olpseganchor{OLP-0387-B007}{start}
د معمول په څېر، ډېر ځله به تر ټولو بېروني قوسونه له ليکلو پرېږدو۔
\olpseganchor{OLP-0387-B007}{end}

\olpseganchor{OLP-0387-B008}{start}
\begin{ex}
  په معياري ژبه~\LR{$\Lang{L_0}$} کښې
  \LR{$\Obj p_1 \lif (\Obj p_1 \land \lnot \Obj p_2)$} يو فارمول دے۔
  د حاصل‌ضرب منطق په ژبه کښې به د دې پر ځاے
  \LR{$\Obj p_1 \lif (\Obj p_1 \odot \lnot \Obj p_2)$} وليکو۔
  که ژبې ته يو \LR{$1$}-ځايه~\LR{$\triangle$} هم ور زيات کړو، نو
  \LR{$\triangle (\Obj p_1 \land \Obj p_2) \lif (\triangle \Obj p_1 \land \triangle \Obj p_2)$}
  غوندې فارمولونه به هم ولرو۔
\end{ex}
\olpseganchor{OLP-0387-B008}{end}

\olpseganchor{OLP-0388-B004}{start}
\olfileid{mvl}{syn}{mat}
\olpseganchor{OLP-0388-B004}{end}

\olpseganchor{OLP-0388-B005}{start}
\olsection{ماتريسونه}
\olpseganchor{OLP-0388-B005}{end}

\olpseganchor{OLP-0388-B006}{start}
يو څو ارزښته منطق د خپلې ژبې، د رښتياوالي د قيمتونو د سټ~\LR{$V$}،
د ټاکل شوو رښتياوالي قيمتونو د يو فرعي سټ، او د خپلو نښلوونکو
د صدق تابعو له مخې تعريفېږي۔ دغو ټولو توکو ته په ګډه
\emph{ماتريس} ويل کېږي۔
\olpseganchor{OLP-0388-B006}{end}

\olpseganchor{OLP-0388-B007}{start}
\begin{defn}[ماتريس]
\ollabel{defn:matrix}
د منطق~\LR{$\Log L$} يو \emph{ماتريس} له لاندې توکو جوړ وي:
\begin{enumerate}
\item د نښلوونکو يو سټ چې ژبه~\LR{$\Lang L$} جوړوي؛
\item د رښتياوالي د قيمتونو يو ناتش سټ~\LR{$V \neq \emptyset$}؛
\item د ټاکل شوو رښتياوالي قيمتونو يو سټ~\LR{$V^+ \subseteq V$}؛
\item په~\LR{$\Lang L$} کښې د هر \LR{$n$}-ځايه نښلوونکي~\LR{$\star$}
دپاره د صدق تابعه~\LR{$\tf{\star} : V^n \to V$}۔ که~\LR{$n = 0$}
وي، نو~\LR{$\tf{\star}$} له مدخل پرته د~\LR{$V$} يو ټاکلے غړے ورکوي۔
\end{enumerate}
\end{defn}
\olpseganchor{OLP-0388-B007}{end}

\olpseganchor{OLP-0388-B008}{start}
\begin{ex}
د کلاسيک منطق~\LogCL{} ماتريس له لاندې توکو جوړ وي:
\begin{enumerate}
  \item معياري بياني ژبه~\LR{$\Lang L_0$}، چې
  \LR{$\lfalse$}، \LR{$\lnot$}، \LR{$\land$}، \LR{$\lor$}، \LR{$\lif$} لري۔
  \item د رښتياوالي د قيمتونو سټ~\LR{$V = \{\True, \False\}$}۔
  \item \LR{$\True$} يوازينے ټاکل شوے قيمت دے؛ يعنې~\LR{$V^+ = \{\True\}$}۔
  \item د~\LR{$\lfalse$} دپاره~\LR{$\tf{\lfalse} = \False$} لرو۔ د نورو
  صدق تابعې د معمول د رښتياوالي جدولونو له مخې ټاکل کېږي
  (\olref{fig:tf-CL} وګورئ)۔
\end{enumerate}
\begin{figure}
  \begin{center}
      \begin{tabular}{c|c}
        \LR{$\tf{\lnot}$} & \\
        \hline
        \LR{$\True$} & \LR{$\False$} \\
        \LR{$\False$} & \LR{$\True$}
      \end{tabular}
      \quad
      \begin{tabular}{c|cc}
        \LR{$\tf{\land}$} & \LR{$\True$} & \LR{$\False$} \\
        \hline
        \LR{$\True$} & \LR{$\True$} & \LR{$\False$} \\
        \LR{$\False$} & \LR{$\False$} & \LR{$\False$}
      \end{tabular}
      \quad
      \begin{tabular}{c|cc}
        \LR{$\tf{\lor}$} & \LR{$\True$} & \LR{$\False$} \\
        \hline
        \LR{$\True$} & \LR{$\True$} & \LR{$\True$} \\
        \LR{$\False$} & \LR{$\True$} & \LR{$\False$}
      \end{tabular}
      \quad
      \begin{tabular}{c|cc}
        \LR{$\tf{\lif}$} & \LR{$\True$} & \LR{$\False$} \\
        \hline
        \LR{$\True$} & \LR{$\True$} & \LR{$\False$} \\
        \LR{$\False$} & \LR{$\True$} & \LR{$\True$}
      \end{tabular}
    \end{center}
    \caption{د کلاسيک منطق~\LR{$\LogCL$} د صدق تابعې۔}
    \ollabel{fig:tf-CL}
  \end{figure}
  \end{ex}
\olpseganchor{OLP-0388-B008}{end}

\olpseganchor{OLP-0389-B004}{start}
\olfileid{mvl}{syn}{val}
\olpseganchor{OLP-0389-B004}{end}

\olpseganchor{OLP-0389-B005}{start}
\olsection{ارزښت ټاکنې او صدق}
\olpseganchor{OLP-0389-B005}{end}

\olpseganchor{OLP-0389-B006}{start}
\begin{defn}[ارزښت ټاکنې]
فرض کړئ~\LR{$V$} د رښتياوالي د قيمتونو يو سټ وي۔ د~\LR{$\Lang{L}$}
له ژبې څخه~\LR{$V$} ته يوه \emph{ارزښت ټاکنه} داسې
تابعه~\LR{$\pAssign{v}$} ده چې د ژبې له بياني متغيرونه
څخه هر يوه ته د~\LR{$V$} غړے ورکوي؛ يعنې
\LR{$\pAssign{v} \colon \PVar \to V$}۔
\end{defn}
\olpseganchor{OLP-0389-B006}{end}

\olpseganchor{OLP-0389-B007}{start}
\begin{defn}\ollabel{defn:pValue}
  د څو ارزښته منطق~\LR{$\Log L$} د رښتياوالي د قيمتونو
  سټ~\LR{$V$} ته د يوې ارزښت ټاکنه~\LR{$\pAssign{v}$} په لرلو،
  د ارزونې تابعه~\LR{$\pValue{v} \colon \Frm[L] \to V$}
  په استقرايي ډول داسې تعريف کړئ:
  \begin{enumerate}
    \item \LR{$\pValue{v}(\Obj p_n) = \pAssign{v}(\Obj p_n)$}؛
    \item که~\LR{$\star$} يو \LR{$0$}-ځايه نښلوونکے وي، نو
    \LR{$\pValue{v}(\star)  = \tf{\star}[\Log L]$}؛
    \item که~\LR{$\star$} يو \LR{$n$}-ځايه نښلوونکے وي او لږ تر لږه يو ځاے ولري، نو
    \[
      \pValue{v}(\star(!A_1, \dots, !A_n)) = \tf{\star}[\Log L]
      (\pValue{v}(!A_1), \dots, \pValue{v}(!A_n)).
    \]
  \end{enumerate}
\end{defn}
\olpseganchor{OLP-0389-B007}{end}

\olpseganchor{OLP-0389-B008}{start}
\begin{defn}[صدق]
\ollabel{defn:satisfaction} فارمول~\LR{$!A$} د
  ارزښت ټاکنه~\LR{$\pAssign{v}$} له مخې هله \emph{پوره کېږي}،
  \LR{$\pSat{v}{!A}[\Log L]$}، چې
  \LR{$\pValue{v}(!A)[\Log L] \in V^+$} وي؛ دلته~\LR{$V^+$}
  د~\LR{$\Log L$} د ټاکل شوو رښتياوالي قيمتونو سټ دے۔
\olpseganchor{OLP-0389-B008}{end}

\olpseganchor{OLP-0389-B009}{start}
  \LR{$\pSat/{v}{!A}[\Log L]$} د «\LR{$\pSat{v}{!A}[\Log
  L]$} نۀ دے» په معنا ليکو۔ که~\LR{$\Gamma$} د فارمولونه
  يو سټ وي، نو~\LR{$\pSat{v}{\Gamma}[\Log L]$} هله او يوازې
  هله وي چې د هر~\LR{$!A \in \Gamma$} دپاره
  \LR{$\pSat{v}{!A}[\Log L]$} وي۔
\end{defn}
\olpseganchor{OLP-0389-B009}{end}

\olpseganchor{OLP-0390-B004}{start}
\olfileid{mvl}{syn}{sem}
\olpseganchor{OLP-0390-B004}{end}

\olpseganchor{OLP-0390-B005}{start}
\olsection{معنايي مفهومونه}
\olpseganchor{OLP-0390-B005}{end}

\olpseganchor{OLP-0390-B006}{start}
فرض کړئ چې يو څو ارزښته منطق~\LR{$\Log L$} په يوه ماتريس ټاکل شوے
وي۔ بيا د~\LR{$\Log L$} دپاره معمول معنايي مفهومونه تعريفولے شو۔
\olpseganchor{OLP-0390-B006}{end}

\olpseganchor{OLP-0390-B007}{start}
\begin{defn}
\begin{enumerate}
\item فارمول~\LR{$!A$} هله \emph{د صدق وړ} دے چې د کوم
  \LR{$\pAssign{v}$} دپاره \LR{$\pSat{v}{!A}$} وي؛ او هله \emph{د صدق نۀ وړ}
  دے چې د هېڅ \LR{$\pAssign{v}$} دپاره \LR{$\pSat{v}{!A}$} نۀ وي؛
\item فارمول~\LR{$!A$} هله \emph{تاوتولوژي} ده چې د ټولو
  ارزښت ټاکنې~\LR{$v$} دپاره \LR{$\pSat{v}{!A}$} وي؛
\item که~\LR{$\Gamma$} د فارمولونه يو سټ وي، نو \LR{$\Gamma \Entails !A$}
  (``\LR{$\Gamma$}، \LR{$!A$} ته استلزام ورکوي'') هله او يوازې هله چې د هرې
  هغې ارزښت ټاکنه~\LR{$\pAssign{v}$} دپاره \LR{$\pSat{v}{!A}$} وي چې
  \LR{$\pSat{v}{\Gamma}$} وي۔
\item که~\LR{$\Gamma$} د فارمولونه يو سټ وي، نو~\LR{$\Gamma$} هله
  \emph{د صدق وړ} دے چې داسې ارزښت ټاکنه~\LR{$\pAssign{v}$} موجوده
  وي چې \LR{$\pSat{v}{\Gamma}$} وي؛ او که داسې نۀ وي، \LR{$\Gamma$}
  \emph{د صدق نۀ وړ} دے۔
\end{enumerate}
\end{defn}
\olpseganchor{OLP-0390-B007}{end}

\olpseganchor{OLP-0390-B008}{start}
د دغو مفهومونو ځينې خاصيتونه د کلاسيک منطق د حالت په څېر دي:
\olpseganchor{OLP-0390-B008}{end}

\olpseganchor{OLP-0390-B009}{start}
\begin{prop}
\ollabel{prop:semanticalfacts}
\begin{enumerate}
\item \LR{$!A$} هله او يوازې هله تاوتولوژي ده چې
  \LR{$\emptyset \Entails !A$} وي؛
\item که~\LR{$\Gamma$} د صدق وړ وي، نو د~\LR{$\Gamma$} هر متناهي فرعي سټ
  هم د صدق وړ دے؛
\item\ollabel{def:monotonicity}%
يکنواختي: که \LR{$\Gamma \subseteq \Delta$}
  او \LR{$\Gamma \Entails !A$} وي، نو \LR{$\Delta \Entails !A$} هم؛
\item\ollabel{def:Cut}%
انتقاليت: که \LR{$\Gamma \Entails !A$} او
  \LR{$\Delta \cup \{ !A\} \Entails !B$} وي، نو \LR{$\Gamma \cup \Delta \Entails
  !B$}؛
\end{enumerate}
\end{prop}
\olpseganchor{OLP-0390-B009}{end}

\olpseganchor{OLP-0390-B010}{start}
\begin{proof}
تمرين۔
\end{proof}
\olpseganchor{OLP-0390-B010}{end}

\olpseganchor{OLP-0390-B011}{start}
\begin{prob}
د \olref[mvl][syn][sem]{prop:semanticalfacts} ثبوت ورکړئ۔
\end{prob}
\olpseganchor{OLP-0390-B011}{end}

\olpseganchor{OLP-0390-B012}{start}
په کلاسيک منطق کښې استلزام له شرطي سره تړلے شو۔ د بېلګې په
توګه، د \emph{مودوس پوننس} اعتبار لرو: که \LR{$\Gamma
\Entails !A$} او \LR{$\Gamma \Entails !A \lif !B$} وي، نو \LR{$\Gamma \Entails
!B$} وي۔ په کلاسيک منطق کښې د \LR{$\Entails$} او \LR{$\lif$} تر منځ يوه بله
مهمه اړيکه د معنايي استنتاج قضيه ده: \LR{$\Gamma \Entails !A
\lif !B$} هله او يوازې هله چې \LR{$\Gamma \cup \{!A\} \Entails !B$}
وي۔ دا پايلې په څو ارزښته منطقونو کښې \emph{لازماً نۀ} صدق کوي۔
دا چې دغه پايلې صدق کوي که نۀ، د شرط د صدق پر تابعې~\LR{$\tf{\lif}$} او د رښتياوالي پر ټاکل شوو قيمتونو پورې اړه لري۔
\olpseganchor{OLP-0390-B012}{end}

\olpseganchor{OLP-0391-B004}{start}
\olfileid{mvl}{syn}{sub}
\olpseganchor{OLP-0391-B004}{end}

\olpseganchor{OLP-0391-B005}{start}
\olsection{د~\texorpdfstring{\LR{$\LogCL$}}{C} د فرعي منطقونو شرطونه}
\olpseganchor{OLP-0391-B005}{end}

\olpseganchor{OLP-0391-B006}{start}
دلته داسې څو ارزښته منطقونه څېړو چې په ماتريسونو کښې ئې د~\LR{$\{\True, \False\}$} د مدخلونو دپاره ئې د صدق
تابعې قيمت د کلاسيکو صدق تابعو له قيمت سره برابر وي۔ په کره توګه،
په دغو منطقونو کښې که~\LR{$x \in \{\True, \False\}$} وي، نو
\LR{$\tf{\lnot}[\Log L](x) = \tf{\lnot}[\LogCL](x)$}؛ او که~\LR{$\star$}
له~\LR{$\land$}، \LR{$\lor$} او~\LR{$\lif$} څخه هر يو وي او
\LR{$x, y \in \{\True, \False\}$} وي، نو
\LR{$\tf{\star}[\Log L](x,y) = \tf{\star}[\LogCL](x,y)$}۔ په بل عبارت،
د~\LR{$\lnot$}، \LR{$\land$}، \LR{$\lor$} او~\LR{$\lif$} صدق تابعې چې پر
\LR{$\{\True,\False\}$} محدودې شي، کټ مټ کلاسيکې صدق تابعې دي۔
\olpseganchor{OLP-0391-B006}{end}

\olpseganchor{OLP-0391-B007}{start}
\begin{prop}\ollabel{prop:mvl-cl}
  فرض کړئ چې د څو ارزښته منطق~\LR{$\Log L$} په ژبه کښې
  \LR{$\lnot$}، \LR{$\land$}، \LR{$\lor$} او~\LR{$\lif$} نښلوونکي شته،
  \LR{$\True, \False \in V$} دي او د دې منطق صدق تابعې
  لاندې شرطونه پوره کوي:
  \begin{enumerate}
    \item\ollabel{prop:not} \LR{$\tf{\lnot}[\Log L](x) = \tf{\lnot}[\LogCL](x)$} که~\LR{$x =
    \True$} يا~\LR{$x = \False$} وي؛
    \item\ollabel{prop:land} \LR{$\tf{\land}[\Log L](x,y) = \tf{\land}[\LogCL](x,y)$}،
    \item\ollabel{prop:lor} \LR{$\tf{\lor}[\Log L](x,y) = \tf{\lor}[\LogCL](x,y)$}،
    \item\ollabel{prop:lif} \LR{$\tf{\lif}[\Log L](x,y) = \tf{\lif}[\LogCL](x,y)$}،
    که~\LR{$x, y \in \{\True, \False\}$} وي۔
  \end{enumerate}
  نو د دواړو منطقونو د ګډې ژبې د هغو فارمولونو دپاره چې هر
  کارېدونکے ثابت او نښلوونکے ئې پر دغو دوو قيمتونو يو شان صدق
  تابعه لري، هره داسې ارزښت ټاکنه~\LR{$\pAssign v$} چې~\LR{$V$} ته ځي او د هر
  بياني متغير دپاره~\LR{$\pAssign v(p) \in \{\True,\False\}$} وي،
  \LR{$\pValue v[\Log L](!A) = \pValue v[\LogCL](!A)$} پوره کوي۔
\end{prop}
\olpseganchor{OLP-0391-B007}{end}

\olpseganchor{OLP-0391-B008}{start}
\begin{proof}
  پر~\LR{$!A$} د استقرا له مخې۔
  \begin{enumerate}
  \item که~\LR{$!A \ident p$} اتمي وي، نو
  \LR{$\pValue v[\Log L](!A) = \pAssign v(p) = \pValue
  v[\LogCL](!A)$} لرو۔
  \item که~\LR{$!A \ident \lnot !B$} وي، نو
  \begin{align*}
    \pValue v[\Log L](!A) & = \tf{\lnot}[\Log L](\pValue v[\Log L](!B)) & &\text{\RL{د \olref[val]{defn:pValue} له مخې}}\\
    & = \tf{\lnot}[\Log L](\pValue v[\LogCL](!B)) && \text{\RL{د استقرايي فرض له مخې}}\\
    & = \tf{\lnot}[\LogCL](\pValue v[\LogCL](!B))
&&\text{\RL{د \olref{prop:not} د فرض له مخې،}}\\
&&&\text{\RL{ځکه چې \LR{$\pValue v[\LogCL](!B) \in \{\True, \False\}$} دے،}}\\
& = \pValue v[\LogCL](!A) &&\text{\RL{د \olref[val]{defn:pValue} له مخې}}.
\end{align*}
\item که~\LR{$!A \ident (!B \land !C)$} وي، نو
\begin{align*}
  \pValue v[\Log L](!A) & = \tf{\land}[\Log L](\pValue v[\Log L](!B), \pValue v[\Log L](!C)) & &\text{\RL{د \olref[val]{defn:pValue} له مخې}}\\
  & = \tf{\land}[\Log L](\pValue v[\LogCL](!B),\pValue v[\LogCL](!C)) && \text{\RL{د استقرايي فرض له مخې}}\\
  & = \tf{\land}[\LogCL](\pValue v[\LogCL](!B),\pValue v[\LogCL](!C))
&&\text{\RL{د \olref{prop:land} د فرض له مخې،}}\\
&&&\text{\RL{ځکه چې \LR{$\pValue v[\LogCL](!B),\pValue v[\LogCL](!C) \in \{\True, \False\}$} دي،}}\\
& = \pValue v[\LogCL](!A) &&\text{\RL{د \olref[val]{defn:pValue} له مخې}}.
\end{align*}
\end{enumerate}
هغه حالتونه چې~\LR{$!A \ident (!B \lor !C)$} او
\LR{$!A \ident (!B \lif !C)$} وي، په همدې ډول دي۔
که په ګډه ژبه کښې بل نښلوونکے يا ثابت هم کارېدلے وي، د هغه
حالت هم د پورته اضافي شرط له مخې راځي: د صفر ځايه ثابت قيمت
په دواړو منطقونو کښې برابر دے؛ او د څو ځايه نښلوونکي د ټولو
فرعي فارمولونو قيمتونه د استقرايي فرض له مخې برابر او دوه ارزښته
دي، نو پر دغو مدخلونو د صدق تابعو برابري د ټول فارمول قيمت
هم برابر ساتي۔
\end{proof}
\olpseganchor{OLP-0391-B008}{end}

\olpseganchor{OLP-0391-B009}{start}
\begin{cor}
  که يو څو ارزښته منطق د~\olref{prop:mvl-cl} شرطونه پوره کوي،
  \LR{$\True \in V^+$} او~\LR{$\False \notin V^+$} وي، نو
  د پورته بيان شوې ګډې ژبې او برابرو صدق تابعو په اړه
  \LR{${\Entails[\Log L]} \subseteq {\Entails[\LogCL]}$}؛ يعنې که
  \LR{$\Gamma \Entails[\Log L] !B$} وي، نو
  \LR{$\Gamma \Entails[\LogCL] !B$} هم وي۔ په ځانګړې توګه،
  دلته ټولې مقدمې او پايله هم د همدغو شرطونو تابع دي۔
  په همدې ګډه ژبه کښې د~\LR{$\Log L$} هره تاوتولوژي يوه
  کلاسيکه تاوتولوژي هم ده۔
\end{cor}
\olpseganchor{OLP-0391-B009}{end}

\olpseganchor{OLP-0391-B010}{start}
\begin{proof}
  نقيض عکس ثابتوو۔ فرض کړئ~\LR{$\Gamma \Entails/[\LogCL] !B$}۔
  بيا يوه داسې ارزښت ټاکنه~\LR{$\pAssign v\colon \PVar \to
  \{\True, \False\}$} شته چې د هر~\LR{$!A \in \Gamma$} دپاره
  \LR{$\pValue v[\LogCL](!A) = \True$} او
  \LR{$\pValue v[\LogCL](!B) = \False$} وي۔ ځکه چې
  \LR{$\True, \False \in V$} دي، دا ارزښت ټاکنه~\LR{$\pAssign v$}
  د~\LR{$\Log L$} دپاره هم ارزښت ټاکنه ده۔ د~\olref{prop:mvl-cl}
  له مخې، د هر~\LR{$!A \in \Gamma$} دپاره
  \LR{$\pValue v[\Log L](!A) = \True$} او
  \LR{$\pValue v[\Log L](!B) = \False$} دي۔ ځکه چې
  \LR{$\True \in V^+$} او~\LR{$\False \notin V^+$} دي، له دې څخه
  \LR{$\pSat{v}{\Gamma}[\Log L]$} او
  \LR{$\pSat/{v}{!B}[\Log L]$} راوځي؛ يعنې
  \LR{$\Gamma \Entails/[\Log L] !B$}۔
\end{proof}
\olpseganchor{OLP-0391-B010}{end}

\olpseganchor{OLP-0392-B004}{start}
\olchapter{mvl}{thr}{درې ارزښته منطقونه}
\olpseganchor{OLP-0392-B004}{end}

\olpseganchor{OLP-0393-B004}{start}
\olfileid{mvl}{thr}{int}
\olpseganchor{OLP-0393-B004}{end}

\olpseganchor{OLP-0393-B005}{start}
\olsection{پېژندنه}
\olpseganchor{OLP-0393-B005}{end}

\olpseganchor{OLP-0393-B006}{start}
که موږ له~\LR{$\True$} او~\LR{$\False$} سره يوازې يو بل قيمت~\LR{$\Undef$}
ورزيات کړو، د درې ارزښته منطق د ماتريس د قيمتونو سټ ترلاسه کوو۔ که څه هم د رښتياوالي
يوازې يو قيمت ورزيات شوے، د~\LR{$\lnot$}، \LR{$\land$}، \LR{$\lor$} او
\LR{$\lif$} د صدق تابعو د تعريفولو امکانونه ډېر دي۔ بيا يو منطق د دغو
صدق تابعو هره ګډوله کارولے شي۔ په دې باب کښې به د ټاکلو قيمتونو دوه انتخابونه څېړو: يا يوازې
\LR{$\True$} ټاکل شوے وي، يا~\LR{$\True$} او~\LR{$\Undef$} دواړه ټاکل شوي وي۔
\olpseganchor{OLP-0393-B006}{end}

\olpseganchor{OLP-0393-B007}{start}
دلته موږ ځينې تر ټولو پېژندل شوي درې ارزښته منطقونه، د هغو
انګېزې او ځينې خاصيتونه وړاندې کوو۔
\olpseganchor{OLP-0393-B007}{end}

\olpseganchor{OLP-0394-B004}{start}
\olfileid{mvl}{thr}{luk}
\olpseganchor{OLP-0394-B004}{end}

\olpseganchor{OLP-0394-B005}{start}
\olsection{د لوکاشېوېچ منطق}
\olpseganchor{OLP-0394-B005}{end}

\olpseganchor{OLP-0394-B006}{start}
د څو ارزښته منطق له لومړنيو خپرو شوو او بشپړ تشريح شوو وړانديزونو
څخه يو د پولنډي فيلسوف يان لوکاشېوېچ و، چې په \href{https://kpbc.umk.pl/Content/247342/PDF/EOD_OPEN_006_12_nr_09.pdf}{۱۹۲۰} کال کښې ئې
وړاندې کړ۔ د ارسطو د سمندري جګړې مسئلې هغه دې ته وهڅاوه: داسې
ښکاري چې \emph{نن} د «سبا به سمندري جګړه وي» جمله نه رښتيا ده
او نه دروغ؛ د هغې د رښتياوالي قيمت لا نۀ دے ټاکل شوے۔ لوکاشېوېچ
وړانديز وکړ چې د دا ډول «راتلونکو اتفاقي» جملو دپاره درېيم
د رښتياوالي قيمت ورزيات شي۔
\begin{explain}
په وارسا کښې د راتلونکي کال د دسمبر پر يوويشتمه غرمه د
لوکاشېوېچ شتون په پام کښې ونيسئ۔ د بحث فرض دا دے چې په هغه
وخت کښې د ده شتون ممکن دے او لازمي نۀ دے۔ د هغه په استدلال،
که د شتون خبره اوس رښتيا وګڼو، پېښه به لازمه وي؛ که يې اوس
دروغ وګڼو، پېښه به ناشونې وي۔ دواړه حکمونه له فرض سره ټکر لري۔
نو د رښتيا (۱) او دروغو (۰) تر څنګ يو درېيم ارزښت ټاکي، او
په \LR{$\frac{1}{2}$} يې ليکي۔ دغه نښه د هغه په اصطلاح د امکان
ارزښت دے۔
\end{explain}
د دې استدلال لپاره د لوکاشېوېچ د ۱۹۳۰ مقالې شپږمه برخه وګورئ (\href{https://ci.nii.ac.jp/ncid/BA03950054}{د اېچ وېبر انګرېزي ژباړه، د لوډويک بورکوفسکي ۱۹۷۰ چاپ، مخونه ۱۶۵–۱۶۶})۔
موږ به د لوکاشېوېچ د درېيم رښتياوالي قيمت دپاره~\LR{$\Undef$}
کاروو۔\footnote{لوکاشېوېچ دلته «ممکن» په هغه معنا کاروي چې نن
عامه نۀ ده: ممکن، خو لازم نۀ۔}
\olpseganchor{OLP-0394-B006}{end}

\olpseganchor{OLP-0394-B007}{start}
په دې تفسير کښې د~\LR{$\lnot$}، \LR{$\land$} او~\LR{$\lor$} نښلوونکو
د صدق تابعې ټاکل اسانه دي: د راتلونکي اتفاقي بيان نفي هم
راتلونکے اتفاقي بيان دے، نو~\LR{$\tf{\lnot}(\Undef) = \Undef$}۔
که د عطف يو اړخ ناټاکلے او بل رښتيا وي، ټول عطف هم ناټاکلے
وي؛ ځکه چې د راتلونکي اتفاقي اړخ له پايلې سره سم عطف کېدے شي
رښتيا شي او کېدے شي دروغ شي۔ نو \[
    \tf{\land}(\True, \Undef) =
\tf{\land}(\Undef, \True) =
\Undef.
\]
خو که بل اړخ دروغ وي، عطف رښتيا کېدے نۀ شي؛ نو
\[\tf{\land}(\False, \Undef) =
\tf{\land}(\Undef, \False) = \False.\]
نور قيمتونه، چې مدخلونه ئې ټاکلي رښتياوالي قيمتونه~\LR{$\True$}
يا~\LR{$\False$} وي، د کلاسيک منطق په څېر دي۔
\olpseganchor{OLP-0394-B007}{end}

\olpseganchor{OLP-0394-B008}{start}
د شرطي په باب حالت لږ پېچلے دے۔ فرض کړئ~\LR{$q$} يو راتلونکے
اتفاقي بيان دے۔ که~\LR{$p$} دروغ وي، نو~\LR{$p \lif q$} به رښتيا وي،
د~\LR{$q$} له راتلونکې پايلې پرته؛ نو بايد
\LR{$\tf{\lif}(\False, \Undef) = \True$} وټاکو۔ او که~\LR{$p$}
رښتيا وي، نو~\LR{$q \lif p$} به د~\LR{$q$} له پايلې پرته رښتيا وي؛
نو~\LR{$\tf{\lif}(\Undef, \True) = \True$}۔ که~\LR{$p$} رښتيا
وي، نو~\LR{$p \lif q$} کېدے شي رښتيا يا دروغ شي؛ نو
\LR{$\tf{\lif}(\True, \Undef) = \Undef$}۔ په همدې ډول، که~\LR{$p$}
دروغ وي، نو~\LR{$q \lif p$} کېدے شي رښتيا يا دروغ شي؛ نو
\LR{$\tf{\lif}(\Undef, \False) = \Undef$}۔ هغه حالت پاتې کېږي
چې~\LR{$p$} او~\LR{$q$} دواړه راتلونکي اتفاقي بيانونه وي۔ د دې انګېزې
له مخې، په دې حالت کښې به مو~\LR{$\Undef$} ټاکلے و۔ خو په دې سره
\LR{$!A \lif !A$} به \emph{نۀ} تاوتولوژي وي۔ لوکاشېوېچ د
\LR{$!A \lor \lnot !A$} او~\LR{$\lnot(!A \land \lnot !A)$} له پرېښودو
سره ستونزه نۀ لرله، خو د~\LR{$!A \lif !A$} پرېښودل ئې نۀ منل۔
نو ئې~\LR{$\tf{\lif}(\Undef, \Undef) = \True$} وټاکه۔
\olpseganchor{OLP-0394-B008}{end}

\olpseganchor{OLP-0394-B009}{start}
\begin{defn}\ollabel{def:lukasiewicz}
درې ارزښته لوکاشېوېچ منطق د لاندې ماتريس په وسيله تعريفېږي:
\begin{enumerate}
  \item معياري بياني ژبه~\LR{$\Lang L_0$} چې
  \LR{$\lnot$}, \LR{$\land$}, \LR{$\lor$}, \LR{$\lif$}
  نښلوونکي لري۔
  \item د رښتياوالي د قيمتونو سټ~\LR{$V = \{\True, \Undef, \False\}$}۔
  \item \LR{$\True$} يوازينے ټاکل شوے قيمت دے؛ يعنې~\LR{$V^+ = \{\True\}$}۔
  \item د صدق تابعې په لاندې جدولونو کښې ورکړل شوې دي:
  \begin{center}
    \begin{tabular}{c|c} 
      \LR{$\tf{\lnot}$} & \\ 
      \hline  
      \LR{$\True$} & \LR{$\False$} \\ 
      \LR{$\Undef$} & \LR{$\Undef$} \\
      \LR{$\False$} & \LR{$\True$} 
    \end{tabular}
    \quad
    \begin{tabular}{c|ccc} 
      \LR{$\tf{\land}[\LogLuk[3]]$} & \LR{$\True$} & \LR{$\Undef$} & \LR{$\False$} \\ 
      \hline 
      \LR{$\True$} & \LR{$\True$} & \LR{$\Undef$} & \LR{$\False$} \\ 
      \LR{$\Undef$} & \LR{$\Undef$} & \LR{$\Undef$} & \LR{$\False$}\\ 
      \LR{$\False$} & \LR{$\False$} & \LR{$\False$} & \LR{$\False$} 
    \end{tabular}
    \\[2ex]
    \begin{tabular}{c|ccc} 
      \LR{$\tf{\lor}[\LogLuk[3]]$} & \LR{$\True$} & \LR{$\Undef$} & \LR{$\False$} \\ 
      \hline 
      \LR{$\True$} & \LR{$\True$} & \LR{$\True$} & \LR{$\True$} \\ 
      \LR{$\Undef$} & \LR{$\True$} & \LR{$\Undef$} & \LR{$\Undef$} \\
      \LR{$\False$} & \LR{$\True$} & \LR{$\Undef$} & \LR{$\False$} 
    \end{tabular}
    \quad
    \begin{tabular}{c|ccc} 
      \LR{$\tf{\lif}[\LogLuk[3]]$} & \LR{$\True$} & \LR{$\Undef$} & \LR{$\False$} \\ 
      \hline 
      \LR{$\True$} & \LR{$\True$} & \LR{$\Undef$} & \LR{$\False$} \\ 
      \LR{$\Undef$} & \LR{$\True$} & \LR{$\True$} & \LR{$\Undef$}  \\ 
      \LR{$\False$} & \LR{$\True$} & \LR{$\True$} & \LR{$\True$} 
    \end{tabular}
  \end{center} 
\end{enumerate}
\end{defn}
\olpseganchor{OLP-0394-B009}{end}

\olpseganchor{OLP-0394-B010}{start}
لکه په اسانه چې څرګندېږي، هر داسې فارمول~\LR{$!A$} چې يوازې
\LR{$\lnot$}، \LR{$\land$} او~\LR{$\lor$} لري، هغه وخت د رښتياوالي
قيمت~\LR{$\Undef$} اخلي چې ټولو بياني متغيرونه ته ئې
\LR{$\Undef$} ورکړل شوے وي۔ نو، د بېلګې په توګه، کلاسيکې
تاوتولوژۍ~\LR{$p \lor \lnot p$} او~\LR{$\lnot(p \land \lnot p)$}
په~\LR{$\LogLuk[3]$} کښې تاوتولوژۍ نۀ دي، ځکه چې کله
\LR{$\pAssign v(p) = \Undef$} وي، نو~\LR{$\pValue{v}(!A) = \Undef$}۔
\olpseganchor{OLP-0394-B010}{end}

\olpseganchor{OLP-0394-B011}{start}
په هغو ارزښت ټاکنې کښې چې د فارمول د هر بياني متغير دپاره
\LR{$\pAssign v(p) = \True$} يا~\LR{$\False$} وي، \LR{$\pValue v(!A)$}
له کلاسيک رښتياوالي قيمت سره برابر وي۔
\olpseganchor{OLP-0394-B011}{end}

\olpseganchor{OLP-0394-B012}{start}
\begin{prop}
  که د~\LR{$!A$} د هر~\LR{$p$} دپاره
  \LR{$\pAssign v(p) \in \{\True, \False\}$} وي، نو
  \LR{$\pValue v(!A)[\LogLuk[3]] = \pValue v(!A)[\LogCL]$}.
\end{prop}
\olpseganchor{OLP-0394-B012}{end}

\olpseganchor{OLP-0394-B013}{start}
\begin{prob}\label{mvl:thr:luk:prob:luk-iff} فرض کړئ
  \LR{$\pValue v(!A \liff !B) = \pValue v((!A \lif !B) \land (!B \lif
  !A))$} په~\LR{$\LogLuk[3]$} کښې داسې تعريفوو۔ د~\LR{$\liff$}
  د رښتياوالي جدول به څنګه وي؟
\end{prob}
\olpseganchor{OLP-0394-B013}{end}

\olpseganchor{OLP-0394-B014}{start}
ډېرې کلاسيکې تاوتولوژۍ په~\LogLuk[3] کښې هم \emph{تاوتولوژۍ دي}،
لکه~\LR{$\lnot p \lif (p \lif q)$}۔ د کلاسيک منطق په څېر، دا د
رښتياوالي په جدولونو هم تصديق کولے شو:
\begin{center}
\begin{tabular}{cc|cccccc}
  \LR{$p$} & \LR{$q$} & \LR{$\lnot$} & \LR{$p$} & \LR{$\lif$} & \LR{$\smash{(}p$} & \LR{$\lif$} & \LR{$q\smash{)}$} \\ \hline
  \True & \True & \False & \True & \True & \True & \True & \True \\
  \True & \Undef & \False & \True & \True & \True & \Undef & \Undef \\
  \True & \False & \False & \True & \True & \True & \False & \False \\
  \Undef & \True & \Undef & \Undef & \True & \Undef & \True & \True \\
  \Undef & \Undef & \Undef & \Undef & \True & \Undef & \True & \Undef \\
  \Undef & \False & \Undef & \Undef & \True & \Undef & \Undef & \False \\
  \False & \True & \True & \False & \True & \False & \True & \True \\
  \False & \Undef & \True & \False & \True & \False & \True & \Undef \\
  \False & \False & \True & \False & \True & \False & \True & \False \\  
\end{tabular}
\end{center}
\olpseganchor{OLP-0394-B014}{end}

\olpseganchor{OLP-0394-B015}{start}
\begin{prob}
  وښيئ چې لاندې فارمولونه په~\LogLuk[3] کښې تاوتولوژۍ دي:
  \begin{enumerate}
    \item \LR{$p \lif (q \lif p)$}
    \item\label{mvl:thr:luk:prob:luk-taut-2} \LR{$\lnot(p \land q) \liff (\lnot p \lor \lnot q)$}
    \item\label{mvl:thr:luk:prob:luk-taut-3} \LR{$\lnot(p \lor q) \liff (\lnot p \land \lnot q)$}
  \end{enumerate}
  (په~\olref[mvl][thr][luk]{prob:luk-taut-2} او
  \olref[mvl][thr][luk]{prob:luk-taut-3} کښې \LR{$!A \liff !B$}
  د~\LR{$(!A \lif !B) \land (!B \lif !A)$} د لنډيز په توګه واخلئ،
  يا د~\olref[mvl][thr][luk]{prob:luk-iff} خپل حل ته رجوع وکړئ۔)
\end{prob}
\olpseganchor{OLP-0394-B015}{end}

\olpseganchor{OLP-0394-B016}{start}
\begin{prob}
  وښيئ چې لاندې کلاسيکې تاوتولوژۍ په~\LogLuk[3] کښې
  تاوتولوژۍ نۀ دي:
  \begin{enumerate}
    \item \LR{$(\lnot p \land p) \lif q$}
    \item \LR{$((p \lif q) \lif p) \lif p$}
    \item \LR{$(p \lif (p \lif q)) \lif (p \lif q)$}
  \end{enumerate}
\end{prob}
\olpseganchor{OLP-0394-B016}{end}

\olpseganchor{OLP-0394-B017}{start}
له دې څخه ښايي څوک داسې وانګېري چې که څه هم ټولې کلاسيکې
تاوتولوژۍ په~\LR{$\LogLuk[3]$} کښې تاوتولوژۍ نۀ دي، لږ تر لږه بايد
په هره ارزښت ټاکنه کښې يا~\LR{$\True$} يا~\LR{$\Undef$} قيمت واخلي۔
خو داسې نۀ ده۔ لاندې فارمول ضد مثال دے:
\[
  \lnot(p \lif \lnot p) \lor \lnot(\lnot p \lif p)
\]
که~\LR{$p$}، \LR{$\Undef$} وي، د دې قيمت~\LR{$\False$} وي۔
\olpseganchor{OLP-0394-B017}{end}

\olpseganchor{OLP-0394-B018}{start}
\begin{prob}
  لاندې کومې اړيکې د لوکاشېوېچ په منطق کښې صدق کوي؟ د هرې
  يوې دپاره د رښتياوالي جدول ورکړئ۔
  \begin{enumerate}
    \item \LR{$p, p \lif q \Entails q$}
    \item \LR{$\lnot\lnot p \Entails p$}
    \item \LR{$p \land q \Entails p$}
    \item \LR{$p \Entails p \land p$}
    \item \LR{$p \Entails p \lor q$}
  \end{enumerate}
\end{prob}
\olpseganchor{OLP-0394-B018}{end}

\olpseganchor{OLP-0394-B019}{start}
لوکاشېوېچ هيله لرله چې د خپل درې ارزښته نظام پر بنسټ د امکان
منطق جوړ کړي۔ د دې دپاره ئې يوځايه نښلوونکے~\LR{$\Diamond !A$}
(د «\LR{$!A$} ممکن دے» دپاره) او ورسره اړوند~\LR{$\Box !A$}
(د «\LR{$!A$} لازم دے» دپاره) ورزيات کړل:
\begin{center}
  \begin{tabular}{c|c} 
    \LR{$\tf{\Diamond}$} & \\ 
    \hline  
    \LR{$\True$} & \LR{$\True$} \\ 
    \LR{$\Undef$} & \LR{$\True$} \\
    \LR{$\False$} & \LR{$\False$} 
  \end{tabular}
  \quad
  \begin{tabular}{c|c} 
    \LR{$\tf{\Box}$} & \\ 
    \hline  
    \LR{$\True$} & \LR{$\True$} \\ 
    \LR{$\Undef$} & \LR{$\False$} \\
    \LR{$\False$} & \LR{$\False$} 
  \end{tabular}
\end{center}
په بل عبارت، \LR{$p$} هله او يوازې هله ممکن دے چې لا ئې دروغوالے
نۀ وي ټاکل شوے؛ او~\LR{$p$} هله او يوازې هله لازم دے چې رښتياوالے
ئې لا ټاکل شوے وي۔
\olpseganchor{OLP-0394-B019}{end}

\olpseganchor{OLP-0394-B020}{start}
\begin{prob}
  وښيئ چې~\LR{$\Box p \liff \lnot\Diamond \lnot p$} او
  \LR{$\Diamond p \liff \lnot \Box \lnot p$} په هغه~\LogLuk[3]
  کښې تاوتولوژۍ دي چې د~\LR{$\Box$} او~\LR{$\Diamond$} د
  رښتياوالي جدولونه ورزيات شوي وي۔
\end{prob}
\olpseganchor{OLP-0394-B020}{end}

\olpseganchor{OLP-0394-B021}{start}
خو د دې وړانديز شوي مودالي منطق نيمګړتياوې ژر څرګندې شوې:
هر څه چې وشي، \LR{$p \land \lnot p$} هېڅکله رښتيا کېدے نۀ شي۔
نو که څه هم اوس ئې ارزښت نۀ وي ټاکل شوے (او ځکه ناټاکلے وي)،
بايد ناشونے وګڼل شي؛ يعنې~\LR{$\lnot \Diamond(p \land \lnot p)$}
بايد تاوتولوژي وي۔ خو که~\LR{$\pAssign v(p) = \Undef$} وي، نو
\LR{$\pValue v(\lnot \Diamond(p \land \lnot p)) = \False$}۔
د امکان عملګر پر ناټاکلي ارزښت رښتيا راولي، او د هغه نفي دروغ راوړي۔ ځکه دا فارمول تاوتولوژي نۀ ده۔ که څه هم
لوکاشېوېچ په دې کښې سم و چې د امکان او لزوم په شان مودالي
توپيرونو دپاره دوه د رښتياوالي قيمتونه بس نۀ دي، د درېيم
قيمت ورزياتول هم بس نۀ دي۔
\olpseganchor{OLP-0394-B021}{end}

\olpseganchor{OLP-0395-B004}{start}
\olfileid{mvl}{thr}{skl}
\olpseganchor{OLP-0395-B004}{end}

\olpseganchor{OLP-0395-B005}{start}
\olsection{د کليني منطقونه}
\olpseganchor{OLP-0395-B005}{end}

\olpseganchor{OLP-0395-B006}{start}
سټيفن کليني دوه درې ارزښته منطقونه وړاندې کړل۔ د هغو انګېزه داسې منطق
و چې د رښتياوالي قيمتونه د محاسبوي کړنلارو پايلې ګڼي: يوه کړنلاره
کېدے شي~\LR{$\True$} يا~\LR{$\False$} ورکړي، خو کېدے شي هېڅکله و نۀ درېږي۔
په دې حالت کښې اړوند د رښتياوالي قيمت ناټاکلے وي او د~\LR{$\Undef$}
په قيمت ښودل کېږي۔
\olpseganchor{OLP-0395-B006}{end}

\olpseganchor{OLP-0395-B007}{start}
د يوې قضيې~\LR{$!A$} د نفي د محاسبې دپاره به لومړے د~\LR{$!A$} قيمت
محاسبه کړئ او بيا به د پايلې مخالف قيمت ورکړئ۔ که د~\LR{$!A$} محاسبه
پای ته و نۀ رسېږي، نو ټوله کړنلاره هم پای ته نۀ رسېږي؛ نو د~\LR{$\Undef$}
نفي~\LR{$\Undef$} ده۔
\olpseganchor{OLP-0395-B007}{end}

\olpseganchor{OLP-0395-B008}{start}
د عطف~\LR{$!A \land !B$} د محاسبې دپاره دوه لارې شته: په يوه کښې
لومړے~\LR{$!A$} او بيا~\LR{$!B$} محاسبه کوو۔ که دواړو~\LR{$\True$} ورکړ، د عطف
پايله~\LR{$\True$} وي، او کنه~\LR{$\False$} وي۔ که له دغو دوو محاسبو يوه
هم و نۀ درېږي، ټوله کړنلاره پای ته نۀ رسېږي۔ نو په دې پرله‌پسې
تفسير کښې که د عطف يو اړخ ناټاکلے وي، ټول عطف هم ناټاکلے وي۔
د فصل په باب هم همداسې ده۔
\olpseganchor{OLP-0395-B008}{end}

\olpseganchor{OLP-0395-B009}{start}
خو که~\LR{$!A$} او~\LR{$!B$} په هممهاله ډول وارزوو، تر دې ښه پايله
تر لاسه کولے شو۔ که له دوو کړنلارو يوه ودرېږي او~\LR{$\False$}
ورکړي، درېدلے شو، ځکه چې د عطف ځواب به خامخا دروغ وي۔ نو په
دې حالت کښې د عطف يو اړخ چې دروغ وي، ټول عطف دروغ وي، که څه هم
بل اړخ ناټاکلے وي۔ په همدې ډول، د فصل د هممهالې محاسبې پر مهال
چې د يوه اړخ کړنلاره ودرېږي او رښتيا ورکړي، درېدلے شو: فصل به
خامخا رښتيا وي۔ نو دلته د مرکبې دعوې پايله معلومولے شو، که څه
هم يو اړخ ئې ناټاکلے وي۔ په دې تفسير کښې~\LR{$\Undef$} د «ناټاکلي»
پر ځای د «نامعلوم» په معنا لوستل کېدے شي۔
\olpseganchor{OLP-0395-B009}{end}

\olpseganchor{OLP-0395-B010}{start}
له دغو دوو تفسيرونو څخه د کليني قوي او کمزوري منطقونه جوړېږي۔
شرطي د~\LR{$\lnot !A \lor !B$} په معادل ډول تعريفېږي۔
\olpseganchor{OLP-0395-B010}{end}

\olpseganchor{OLP-0395-B011}{start}
\begin{defn}
\emph{د کليني قوي منطق}~\LR{$\LogKs$} د لاندې ماتريس په وسيله تعريفېږي:
\begin{enumerate}
  \item معياري بياني ژبه~\LR{$\Lang L_0$} چې
  \LR{$\lnot$}, \LR{$\land$}, \LR{$\lor$}, \LR{$\lif$} نښلوونکي لري۔
  \item د رښتياوالي د قيمتونو سټ~\LR{$V = \{\True, \Undef, \False\}$}۔
  \item \LR{$\True$} يوازينے ټاکل شوے قيمت دے؛ يعنې~\LR{$V^+ = \{\True\}$}۔
  \item د صدق تابعې په لاندې جدولونو کښې ورکړل شوې دي:
  \begin{center}
    \begin{tabular}{c|c} 
      \LR{$\tf{\lnot}$} & \\ 
      \hline  
      \LR{$\True$} & \LR{$\False$} \\ 
      \LR{$\Undef$} & \LR{$\Undef$} \\
      \LR{$\False$} & \LR{$\True$} 
    \end{tabular}
    \quad
    \begin{tabular}{c|ccc} 
      \LR{$\tf{\land}[\LogKs]$} & \LR{$\True$} & \LR{$\Undef$} & \LR{$\False$} \\ 
      \hline 
      \LR{$\True$} & \LR{$\True$} & \LR{$\Undef$} & \LR{$\False$} \\ 
      \LR{$\Undef$} & \LR{$\Undef$} & \LR{$\Undef$} & \LR{$\False$}\\ 
      \LR{$\False$} & \LR{$\False$} & \LR{$\False$} & \LR{$\False$} 
    \end{tabular}
    \\[2ex]
    \begin{tabular}{c|ccc} 
      \LR{$\tf{\lor}[\LogKs]$} & \LR{$\True$} & \LR{$\Undef$} & \LR{$\False$} \\ 
      \hline 
      \LR{$\True$} & \LR{$\True$} & \LR{$\True$} & \LR{$\True$} \\ 
      \LR{$\Undef$} & \LR{$\True$} & \LR{$\Undef$} & \LR{$\Undef$} \\
      \LR{$\False$} & \LR{$\True$} & \LR{$\Undef$} & \LR{$\False$} 
    \end{tabular}
    \quad
    \begin{tabular}{c|ccc} 
      \LR{$\tf{\lif}[\LogKs]$} & \LR{$\True$} & \LR{$\Undef$} & \LR{$\False$} \\ 
      \hline 
      \LR{$\True$} & \LR{$\True$} & \LR{$\Undef$} & \LR{$\False$} \\ 
      \LR{$\Undef$} & \LR{$\True$} & \LR{$\Undef$} & \LR{$\Undef$}  \\ 
      \LR{$\False$} & \LR{$\True$} & \LR{$\True$} & \LR{$\True$} 
    \end{tabular}
  \end{center} 
\end{enumerate}
\end{defn}
\olpseganchor{OLP-0395-B011}{end}

\olpseganchor{OLP-0395-B012}{start}
\begin{defn}
\emph{د کليني کمزورے منطق}~\LR{$\LogKw$} د لاندې ماتريس په وسيله تعريفېږي:
\begin{enumerate}
  \item معياري بياني ژبه~\LR{$\Lang L_0$} چې
  \LR{$\lnot$}, \LR{$\land$}, \LR{$\lor$}, \LR{$\lif$} نښلوونکي لري۔
  \item د رښتياوالي د قيمتونو سټ~\LR{$V = \{\True, \Undef, \False\}$}۔
  \item \LR{$\True$} يوازينے ټاکل شوے قيمت دے؛ يعنې~\LR{$V^+ = \{\True\}$}۔
  \item د صدق تابعې په لاندې جدولونو کښې ورکړل شوې دي:
  \begin{center}
    \begin{tabular}{c|c} 
      \LR{$\tf{\lnot}$} & \\ 
      \hline  
      \LR{$\True$} & \LR{$\False$} \\ 
      \LR{$\Undef$} & \LR{$\Undef$} \\
      \LR{$\False$} & \LR{$\True$} 
    \end{tabular}
    \quad
    \begin{tabular}{c|ccc} 
      \LR{$\tf{\land}[\LogKw]$} & \LR{$\True$} & \LR{$\Undef$} & \LR{$\False$} \\ 
      \hline 
      \LR{$\True$} & \LR{$\True$} & \LR{$\Undef$} & \LR{$\False$} \\ 
      \LR{$\Undef$} & \LR{$\Undef$} & \LR{$\Undef$} & \LR{$\Undef$}\\ 
      \LR{$\False$} & \LR{$\False$} & \LR{$\Undef$} & \LR{$\False$} 
    \end{tabular}
    \\[2ex]
    \begin{tabular}{c|ccc} 
      \LR{$\tf{\lor}[\LogKw]$} & \LR{$\True$} & \LR{$\Undef$} & \LR{$\False$} \\ 
      \hline 
      \LR{$\True$} & \LR{$\True$} & \LR{$\Undef$} & \LR{$\True$} \\ 
      \LR{$\Undef$} & \LR{$\Undef$} & \LR{$\Undef$} & \LR{$\Undef$} \\
      \LR{$\False$} & \LR{$\True$} & \LR{$\Undef$} & \LR{$\False$} 
    \end{tabular}
    \quad
    \begin{tabular}{c|ccc} 
      \LR{$\tf{\lif}[\LogKw]$} & \LR{$\True$} & \LR{$\Undef$} & \LR{$\False$} \\ 
      \hline 
      \LR{$\True$} & \LR{$\True$} & \LR{$\Undef$} & \LR{$\False$} \\ 
      \LR{$\Undef$} & \LR{$\Undef$} & \LR{$\Undef$} & \LR{$\Undef$}  \\ 
      \LR{$\False$} & \LR{$\True$} & \LR{$\Undef$} & \LR{$\True$} 
    \end{tabular}
  \end{center} 
\end{enumerate}
\end{defn}
\olpseganchor{OLP-0395-B012}{end}

\olpseganchor{OLP-0395-B013}{start}
\begin{prop}
  په~\LR{$\LogKs$} او~\LR{$\LogKw$} کښې هېڅ تاوتولوژي نشته۔
\end{prop}
\olpseganchor{OLP-0395-B013}{end}

\olpseganchor{OLP-0395-B014}{start}
\begin{proof}
  که ټولو بياني متغيرونه~\LR{$p$} ته
  \LR{$\pAssign v(p) = \Undef$} ورکړل شي، نو د هر فارمول~\LR{$!A$}
  د رښتياوالي قيمت~\LR{$\pValue v(!A) = \Undef$} وي؛ ځکه چې
  \[
    \tf{\lnot}(\Undef) = \tf{\lor}(\Undef, \Undef) = \tf{\land}(\Undef,
  \Undef) = \tf{\lif}(\Undef, \Undef) = \Undef
  \]
  په دواړو منطقونو کښې همداسې ده۔ ځکه چې په~\LR{$\LogKs$} او
  \LR{$\LogKw$} دواړو کښې~\LR{$\Undef \notin V^+$} دے، نو په دې
  ارزښت ټاکنه کښې~\LR{$!A$} ټاکل شوے قيمت نۀ اخلي۔
\end{proof}
\olpseganchor{OLP-0395-B014}{end}

\olpseganchor{OLP-0395-B015}{start}
که څه هم د کليني په قوي او کمزوري منطقونو کښې تاوتولوژۍ نشته،
خو د هغو د استلزام اړيکې بې‌محتوا نۀ دي۔
\olpseganchor{OLP-0395-B015}{end}

\olpseganchor{OLP-0395-B016}{start}
\begin{prob}
  لاندې کومې اړيکې (الف) د کليني په قوي او (ب) د کليني په
  کمزوري منطق کښې صدق کوي؟ د هرې يوې دپاره د رښتياوالي جدول ورکړئ۔
  \begin{enumerate}
    \item \LR{$p, p \lif q \Entails q$}
    \item \LR{$p \lor q, \lnot p \Entails q$}
    \item \LR{$p \land q \Entails p$}
    \item \LR{$p \Entails p \land p$}
    \item \LR{$p \Entails p \lor q$}
  \end{enumerate}
\end{prob}
\olpseganchor{OLP-0395-B016}{end}

\olpseganchor{OLP-0395-B017}{start}
دميتري بوچوار~\LR{$\Undef$} د «بې‌معنا» په توګه تعبير کړ او هڅه ئې وکړه
چې د دروغجن د پارادوکس په څېر پارادوکسونه پرې حل کړي: هغه وټاکله
چې پارادوکس لرونکې جملې به~\LR{$\Undef$} قيمت اخلي۔ ده داسې منطق
وړاندې کړ چې په اصل کښې د کليني کمزورے منطق دے، خو اضافي
نښلوونکي لري۔ له هغو څخه دوه «بهرنۍ نفي» او د «ناټاکلتيا»
عملګر دي:
\begin{center}
  \begin{tabular}{c|c} 
    \LR{$\tf{\sim}$} & \\ 
    \hline  
    \LR{$\True$} & \LR{$\False$} \\ 
    \LR{$\Undef$} & \LR{$\True$} \\
    \LR{$\False$} & \LR{$\True$} 
  \end{tabular}
\quad
  \begin{tabular}{c|c} 
    \LR{$\tf{+}$} & \\ 
    \hline  
    \LR{$\True$} & \LR{$\False$} \\ 
    \LR{$\Undef$} & \LR{$\True$} \\
    \LR{$\False$} & \LR{$\False$} 
  \end{tabular}
\end{center}
\olpseganchor{OLP-0395-B017}{end}

\olpseganchor{OLP-0395-B018}{start}
\begin{prob}
  آيا د بوچوار په منطق کښې~\LR{$\sim$} د~\LR{$\lnot$} او~\LR{$+$} په وسيله
  تعريفولے شئ؟ وڅېړئ چې داسې فارمول شته که نۀ چې يوازې
  بياني متغير~\LR{$p$} پکښې وي، \LR{$\sim$} پکښې نۀ وي،
  او تل د~\LR{$\mathord{\sim}p$} په شان د رښتياوالي قيمت واخلي۔
  خپل ځواب د رښتياوالي د جدولونو له مخې ثابت کړئ۔
\end{prob}
\olpseganchor{OLP-0395-B018}{end}

\olpseganchor{OLP-0396-B004}{start}
\olfileid{mvl}{thr}{god}
\olpseganchor{OLP-0396-B004}{end}

\olpseganchor{OLP-0396-B005}{start}
\olsection{د ګوډل منطقونه}
\olpseganchor{OLP-0396-B005}{end}

\olpseganchor{OLP-0396-B006}{start}
کرټ ګوډل د~\LR{$n$} ارزښته منطقونو يوه لړۍ وړاندې کړه۔ هر يو ئې
ټول هغه فارمولونه رانغاړي چې په شهودي منطق کښې معتبر دي، او
پخپله د کلاسيک منطق په دننه کښې راځي۔ دلته د دې لړۍ لومړنے
د پام وړ منطق دے:
\olpseganchor{OLP-0396-B006}{end}

\olpseganchor{OLP-0396-B007}{start}
\begin{defn}\ollabel{defn:goedel}
\emph{د ګوډل \LR{$3$} ارزښته منطق}~\LR{$\LogGod$} په دې ماتريس تعريفېږي:
\begin{enumerate}
  \item معياري بياني ژبه~\LR{$\Lang L_0$} چې
  \LR{$\lfalse$}، \LR{$\lnot$}، \LR{$\land$}، \LR{$\lor$} او~\LR{$\lif$} لري۔
  \item د رښتياوالي د قيمتونو سټ~\LR{$V = \{\True, \Undef, \False\}$}۔
  \item \LR{$\True$} يوازينے ټاکل شوے قيمت دے؛ يعنې~\LR{$V^+ = \{\True\}$}۔
  \item د~\LR{$\lfalse$} دپاره~\LR{$\tf{\lfalse} = \False$} دے۔ د پاتې
  نښلوونکو صدق تابعې په لاندې جدولونو کښې ورکړل شوې دي:
  \begin{center}
    \begin{tabular}{c|c} 
      \LR{$\tf{\lnot}[\LogGod]$} & \\ 
      \hline  
      \LR{$\True$} & \LR{$\False$} \\ 
      \LR{$\Undef$} & \LR{$\False$} \\
      \LR{$\False$} & \LR{$\True$} 
    \end{tabular}
    \quad
    \begin{tabular}{c|ccc} 
      \LR{$\tf{\land}[\LogGod]$} & \LR{$\True$} & \LR{$\Undef$} & \LR{$\False$} \\ 
      \hline 
      \LR{$\True$} & \LR{$\True$} & \LR{$\Undef$} & \LR{$\False$} \\ 
      \LR{$\Undef$} & \LR{$\Undef$} & \LR{$\Undef$} & \LR{$\False$}\\ 
      \LR{$\False$} & \LR{$\False$} & \LR{$\False$} & \LR{$\False$} 
    \end{tabular}
    \\[2ex]
    \begin{tabular}{c|ccc} 
      \LR{$\tf{\lor}[\LogGod]$} & \LR{$\True$} & \LR{$\Undef$} & \LR{$\False$} \\ 
      \hline 
      \LR{$\True$} & \LR{$\True$} & \LR{$\True$} & \LR{$\True$} \\ 
      \LR{$\Undef$} & \LR{$\True$} & \LR{$\Undef$} & \LR{$\Undef$} \\
      \LR{$\False$} & \LR{$\True$} & \LR{$\Undef$} & \LR{$\False$} 
    \end{tabular}
    \quad
    \begin{tabular}{c|ccc} 
      \LR{$\tf{\lif}[\LogGod]$} & \LR{$\True$} & \LR{$\Undef$} & \LR{$\False$} \\ 
      \hline 
      \LR{$\True$} & \LR{$\True$} & \LR{$\Undef$} & \LR{$\False$} \\ 
      \LR{$\Undef$} & \LR{$\True$} & \LR{$\True$} & \LR{$\False$}  \\ 
      \LR{$\False$} & \LR{$\True$} & \LR{$\True$} & \LR{$\True$} 
    \end{tabular}
  \end{center} 
\end{enumerate}
\end{defn}
\olpseganchor{OLP-0396-B007}{end}

\olpseganchor{OLP-0396-B008}{start}
تاسو به وګورئ چې د~\LR{$\land$} او~\LR{$\lor$} صدق جدولونه د
لوکاشېوېچ او د کليني د قوي منطق له جدولونو سره يو شان دي، خو
د~\LR{$\lnot$} جدول يوازې د ګوډل په منطق کښې له نورو بېل دے، او
د~\LR{$\lif$} جدولونه په درې واړو منطقونو کښې سره بېل دي۔
د ګوډل په منطق کښې~\LR{$\tf{\lnot}(\Undef) = \False$} دے۔ د لوکاشېوېچ
او کليني د منطقونو پر خلاف، دلته~\LR{$\tf{\lif}(\Undef, \False) = \False$}
دے؛ او د کليني د منطق پر خلاف (خو د لوکاشېوېچ د منطق په شان)،
\LR{$\tf{\lif}(\Undef,
\Undef) = \True$} دے۔
\olpseganchor{OLP-0396-B008}{end}

\olpseganchor{OLP-0396-B009}{start}
لکه چې د شهودي منطق سره پورته ياده شوې اړيکه ښيي،
\LR{$\LogGod[3]$} شهودي منطق ته نږدې دے۔ د شهودي منطق ټول معتبر
فارمولونه په~\LR{$\LogGod[3]$} کښې تاوتولوژۍ دي؛ او ډېرې کلاسيکې
تاوتولوژۍ چې په شهودي منطق کښې معتبرې نۀ دي، په~\LR{$\LogGod[3]$}
کښې هم تاوتولوژۍ نۀ دي۔ د بېلګې په توګه، لاندې فارمولونه
تاوتولوژۍ نۀ دي:
\begin{align*}
  & p \lor \lnot p && (p \lif q) \lif (\lnot p \lor q) \\
  & \lnot\lnot p \lif p && \lnot(\lnot p \land \lnot q) \lif (p \lor q) \\
  & ((p \lif q) \lif p) \lif p && \lnot(p \lif q) \lif (p \land \lnot q)
\end{align*}
خو د~\LR{$\LogGod[3]$} هره تاوتولوژي په شهودي منطق کښې معتبره
نۀ ده؛ د بېلګې په توګه~\LR{$\lnot\lnot p \lor \lnot p$} يا~\LR{$(p
\lif q) \lor (q \lif p)$}۔
\olpseganchor{OLP-0396-B009}{end}

\olpseganchor{OLP-0396-B010}{start}
\begin{prob}
  د صدق جدولونه ورکړئ او وښيئ چې لاندې فارمولونه د~\LR{$\LogGod[3]$}
  تاوتولوژۍ دي:
  \begin{align*}
    & \lnot\lnot p \lor \lnot p\\
    & (p \lif q) \lor (q \lif p) \\
    & \lnot(p \land q) \lif (\lnot p \lor \lnot q) \\
    & (p \lif q) \lor (q \lif r) \lor (r \lif s)
  \end{align*}
\end{prob}
\olpseganchor{OLP-0396-B010}{end}

\olpseganchor{OLP-0396-B011}{start}
\begin{prob}
  د صدق جدولونه ورکړئ او وښيئ چې لاندې فارمولونه د~\LR{$\LogGod[3]$}
  تاوتولوژۍ نۀ دي:
  \begin{align*}
    & (p \lif q) \lif (\lnot p \lor q) \\
    & \lnot(\lnot p \land \lnot q) \lif (p \lor q) \\
    & ((p \lif q) \lif p) \lif p \\
    & \lnot(p \lif q) \lif (p \land \lnot q)
  \end{align*}
\end{prob}
\olpseganchor{OLP-0396-B011}{end}

\olpseganchor{OLP-0396-B012}{start}
\begin{prob}
  له لاندې اړيکو کومې د ګوډل په منطق کښې صدق کوي؟ د هرې يوې
  دپاره د صدق جدول ورکړئ۔
  \begin{enumerate}
    \item \LR{$p, p \lif q \Entails q$}
    \item \LR{$p \lor q, \lnot p \Entails q$}
    \item \LR{$p \land q \Entails p$}
    \item \LR{$p \Entails p \land p$}
    \item \LR{$p \Entails p \lor q$}
  \end{enumerate}
\end{prob}
\olpseganchor{OLP-0396-B012}{end}

\olpseganchor{OLP-0397-B004}{start}
\olfileid{mvl}{thr}{mul}
\olpseganchor{OLP-0397-B004}{end}

\olpseganchor{OLP-0397-B005}{start}
\olsection{يوازې \texorpdfstring{\LR{$\True$}}{T} نۀ، نور قيمتونه هم ټاکل}
\olpseganchor{OLP-0397-B005}{end}

\olpseganchor{OLP-0397-B006}{start}
تر اوسه چې کوم منطقونه مو ليدلي، په ټولو کښې د ټاکل شوو
رښتياوالي قيمتونو سټ~\LR{$V^+ = \{\True\}$} و؛ يعنې يوه دعوه هله
رښتيا ګڼل کېده چې قيمت ئې~\LR{$\True$} و۔ خو کېدے شي هغه دعوه هم
رښتيا وګڼو چې يوازې~\LR{$\False$} نۀ وي۔ بيا د بېلګې په توګه په
ماتريس کښې~\LR{$V^+ = \{\True, \Undef\}$} ټاکل يو بل منطق جوړوي۔
\olpseganchor{OLP-0397-B006}{end}

\olpseganchor{OLP-0397-B007}{start}
\begin{defn}
\emph{د پارادوکس منطق}~\LR{$\LogLP$} په دې ماتريس تعريفېږي:
\begin{enumerate}
  \item معياري بياني ژبه~\LR{$\Lang L_0$} چې
  \LR{$\lnot$}, \LR{$\land$}, \LR{$\lor$}, \LR{$\lif$}
  نښلوونکي لري۔
  \item د رښتياوالي د قيمتونو سټ~\LR{$V = \{\True, \Undef, \False\}$}۔
  \item \LR{$\True$} او~\LR{$\Undef$} دواړه ټاکل شوي دي؛ يعنې~\LR{$V^+ = \{\True, \Undef\}$}۔
  \item صدق تابعې د کليني د قوي منطق له تابعو سره يو شان دي۔
\end{enumerate}
\end{defn}
\olpseganchor{OLP-0397-B007}{end}

\olpseganchor{OLP-0397-B008}{start}
\begin{defn}
د هالډېن \emph{د بې‌معنايۍ منطق}~\LR{$\LogHal$} په دې ماتريس تعريفېږي:
\begin{enumerate}
  \item معياري بياني ژبه~\LR{$\Lang L_0$} چې
  \LR{$\lnot$}, \LR{$\land$}, \LR{$\lor$}, \LR{$\lif$}
  او يو \LR{$1$} ځايه نښلوونکے~\LR{$+$} لري۔
  \item د رښتياوالي د قيمتونو سټ~\LR{$V = \{\True, \Undef, \False\}$}۔
  \item \LR{$\True$} او~\LR{$\Undef$} دواړه ټاکل شوي دي؛ يعنې~\LR{$V^+ = \{\True, \Undef\}$}۔
  \item صدق تابعې د کليني د کمزوري منطق له تابعو سره يو شان دي،
  خو د «بې‌معنا دے» عملګر هم ورسره شته:
  \begin{center}
    \begin{tabular}{c|c} 
    \LR{$\tf{+}$} & \\ 
    \hline  
    \LR{$\True$} & \LR{$\False$} \\ 
    \LR{$\Undef$} & \LR{$\True$} \\
    \LR{$\False$} & \LR{$\False$} 
    \end{tabular}
  \end{center}
\end{enumerate}
\end{defn}
\olpseganchor{OLP-0397-B008}{end}

\olpseganchor{OLP-0397-B009}{start}
دا منطقونه، سره له دې چې د کليني له منطقونو سره صدق جدولونه
شريکوي، د هغو پر خلاف \emph{تاوتولوژۍ لري}۔
\olpseganchor{OLP-0397-B009}{end}

\olpseganchor{OLP-0397-B010}{start}
\begin{prop}\ollabel{prop:LP-taut-CL} د~\LR{$\LogLP$} تاوتولوژۍ د کلاسيک
  بياني منطق له تاوتولوژيو سره يو شان دي۔
\end{prop}
\olpseganchor{OLP-0397-B010}{end}

\olpseganchor{OLP-0397-B011}{start}
\begin{proof}
  د~\olref[syn][sub]{prop:mvl-cl} له مخې، که~\LR{$\Entails[\LogLP] !A$}
  وي، نو~\LR{$\Entails[\LogCL] !A$} هم وي۔ د برعکس لوري د ښودلو دپاره
  ثابتوو چې که داسې ارزښت ټاکنه~\LR{$\pAssign v\colon \PVar \to \{\False, \True,
  \Undef\}$} وي چې~\LR{$\pValue v(!A)[\LogKs] = \False$}، نو داسې
  ارزښت ټاکنه~\LR{$\pAssign {v'}\colon \PVar \to \{\False, \True\}$}
  هم شته چې~\LR{$\pValue {v'}(!A)[\LogCL] = \False$}۔ دا د~\LR{$\LogLP$}
  دپاره مطلوبه پايله ورکوي، ځکه~\LR{$\LogKs$} او~\LR{$\LogLP$} يو شان
  ځانګړې صدق تابعې لري، او~\LR{$\False$} د~\LR{$\LogLP$} يوازينے
  نۀ ټاکل شوے قيمت دے (همدا د~\LR{$\LogLP$} او~\LR{$\LogKs$} يوازينے توپير دے)۔
  نو که~\LR{$\Entails/[\LogLP] !A$} وي، د يوه~ارزښت ټاکنه~\LR{$\pAssign v$}
  دپاره~\LR{$\pValue v(!A)[\LogLP] = \pValue
  v[\LogKs](!A) = \False$} وي۔ د ثابتېدونکې دعوې له مخې
  \LR{$\pValue{v'}[\LogCL](!A) = \False$}، يعنې~\LR{$\Entails/[\LogCL] !A$}۔
\olpseganchor{OLP-0397-B011}{end}
  
\olpseganchor{OLP-0397-B012}{start}
  د دې دعوې د ثبوت دپاره لومړے~\LR{$\pAssign {v'}$} داسې تعريفوو:
  \[
    \pAssign {v'}(p) = 
    \begin{cases}
    \True & \text{\RL{که }} \pAssign {v}(p) \in \{\True, \Undef\}\\
    \False & \text{\RL{په بل حالت کښې}}
  \end{cases}
  \]
  اوس پر~\LR{$!A$} په استقرا سره ښيو چې (الف) که~\LR{$\pValue v(!A)[\LogKs]
  = \False$} وي، نو~\LR{$\pValue {v'}(!A)[\LogCL] = \False$}؛ او (ب) که
  \LR{$\pValue v(!A)[\LogKs] = \True$} وي، نو~\LR{$\pValue {v'}(!A)[\LogCL] =
  \True$}۔
  \begin{enumerate}
    \item د استقرا بنسټ:~\LR{$!A \ident p$}۔ د~\olref[syn][val]{defn:pValue}
    له مخې~\LR{$\pValue v(!A)[\LogKs] = \pAssign v(p)$} دے۔ که دا قيمت
    دروغ وي، تعريف شوې نوې ټاکنه هم دروغ ورکوي؛ او که رښتيا وي،
    نوې ټاکنه هم رښتيا ورکوي۔ له همدې امله (الف) او (ب) دواړه پوره دي۔
    په ناټاکلي مدخل کښې نوې ټاکنه رښتيا ورکوي۔ (الف) او (ب) يوازې د
    رښتيا او دروغو مدخلونه رانغاړي۔
\olpseganchor{OLP-0397-B012}{end}
    
\olpseganchor{OLP-0397-B013}{start}
    د استقرايي ګام دپاره لاندې حالتونه وګورئ:
    \item \LR{$!A \ident \lnot !B$}. 
    \begin{enumerate}
      \item فرض کړئ~\LR{$\pValue v(\lnot!B)[\LogKs] = \False$}۔ د
      \LR{$\tf{\lnot}[\LogKs]$} د تعريف له مخې~\LR{$\pValue v(!B)[\LogKs]  =
      \True$}۔ د استقرايي فرض د (ب) حالت له مخې~\LR{$\pValue
      {v'}(!B)[\LogCL]  = \True$} کېږي؛ نو~\LR{$\pValue {v'}(\lnot !B)[\LogCL] = \False$}۔
      \item فرض کړئ~\LR{$\pValue v(\lnot!B)[\LogKs] = \True$}۔ د
      \LR{$\tf{\lnot}[\LogKs]$} د تعريف له مخې~\LR{$\pValue v(!B)[\LogKs]  =
      \False$}۔ د استقرايي فرض د (الف) حالت له مخې~\LR{$\pValue
      {v'}(!B)[\LogCL]  = \False$} کېږي؛ نو~\LR{$\pValue {v'}(\lnot !B)[\LogCL] = \True$}۔
    \end{enumerate}
\olpseganchor{OLP-0397-B013}{end}

\olpseganchor{OLP-0397-B014}{start}
    \item \LR{$!A \ident (!B \land !C)$}.
    \begin{enumerate}
      \item فرض کړئ~\LR{$\pValue v(!B \land !C)[\LogKs] = \False$}۔ د
      \LR{$\tf{\land}[\LogKs]$} د تعريف له مخې يا~\LR{$\pValue v(!B)[\LogKs]  =
      \False$} وي، يا~\LR{$\pValue v(!C)[\LogKs]  = \False$} وي۔ د
      استقرايي فرض د (الف) حالت له مخې يا~\LR{$\pValue {v'}(!B)[\LogCL]  = \False$}
      وي، يا~\LR{$\pValue {v'}(!C)[\LogCL]  = \False$}؛ نو~\LR{$\pValue {v'}(!B \land
      !C)[\LogCL] = \False$}۔
      \item فرض کړئ~\LR{$\pValue v(!B \land !C)[\LogKs] = \True$}۔ د
      \LR{$\tf{\land}[\LogKs]$} د تعريف له مخې~\LR{$\pValue v(!B)[\LogKs]  =
      \True$} او~\LR{$\pValue v(!C)[\LogKs]  = \True$}۔ د
      استقرايي فرض د (ب) حالت له مخې~\LR{$\pValue {v'}(!B)[\LogCL]  = \True$}
      او~\LR{$\pValue {v'}(!C)[\LogCL]  = \True$}؛ نو~\LR{$\pValue {v'}(!B \land
      !C)[\LogCL] = \True$}۔
    \end{enumerate}
  \end{enumerate}
\olpseganchor{OLP-0397-B014}{end}

\olpseganchor{OLP-0397-B015}{start}
  نور دوه حالتونه ورته دي او تمرين ته پرېښودل شوي دي۔ په بله لار،
  پورته ثبوت د هغو ټولو~فارمولونه دپاره پايله ثابتوي چې يوازې
  \LR{$\lnot$} او~\LR{$\land$} لري۔ بيا دا حقيقتونه کارولے شو چې په
  \LR{$\LogKs$} او~\LR{$\LogCL$} دواړو کښې د هر~\LR{$\pAssign v$} دپاره
  \LR{$\pValue v(!B \lor
  !C) = \pValue v(\lnot(\lnot!B \land \lnot !C))$} او~\LR{$\pValue v(!B \lif
  !C) = \pValue v(\lnot(!B \land \lnot !C))$} دي۔
\end{proof}
\olpseganchor{OLP-0397-B015}{end}

\olpseganchor{OLP-0397-B016}{start}
\begin{prob}
  د~\olref[mvl][thr][mul]{prop:LP-taut-CL} ثبوت بشپړ کړئ؛ يعنې
  په هغو حالتونو کښې (الف) او (ب) ثابت کړئ چې~\LR{$!A \ident (!B \lor !C)$}
  او~\LR{$!A \ident (!B \lif !C)$} وي۔
\end{prob}
\olpseganchor{OLP-0397-B016}{end}

\olpseganchor{OLP-0397-B017}{start}
\begin{prob}
ثابته کړئ چې هره کلاسيکه تاوتولوژي په~\LR{$\LogHal$} کښې هم تاوتولوژي ده۔
\end{prob}
\olpseganchor{OLP-0397-B017}{end}

\olpseganchor{OLP-0397-B018}{start}
په ګډه بياني ژبه کښې د دغو منطقونو تاوتولوژۍ له کلاسيک منطق سره يو شان دي،
خو د استلزام اړيکې ئې بېلې دي۔ د بېلګې په توګه،~\LR{$\LogLP$}
\emph{ناسازګاري زغمونکے} دے، ځکه چې~\LR{$\lnot p, p \Entails/ q$}؛
نو د انفجار اصل~\LR{$\lnot !A, !A \Entails !B$} په عمومي ډول نۀ صدق کوي۔
(د~\LR{$!A$} او~\LR{$!B$} په ځينو حالتونو کښې صدق کوي، لکه چې~\LR{$!B$}
تاوتولوژي وي۔)
\olpseganchor{OLP-0397-B018}{end}

\olpseganchor{OLP-0397-B019}{start}
\begin{prob}
  له لاندې اړيکو کومې (الف) په~\LR{$\LogLP$} او (ب) په~\LR{$\LogHal$}
  کښې صدق کوي؟ د هرې يوې دپاره د صدق جدول ورکړئ۔
  \begin{enumerate}
    \item \LR{$p, p \lif q \Entails q$}
    \item \LR{$\lnot q, p \lif q \Entails \lnot p$}
    \item \LR{$p \lor q, \lnot p \Entails q$}
    \item \LR{$\lnot p, p \Entails q$}
    \item \LR{$p \Entails p \lor q$}
    \item \LR{$p \lif q, q\lif r \Entails p \lif r$}
  \end{enumerate}
\end{prob}
\olpseganchor{OLP-0397-B019}{end}

\olpseganchor{OLP-0397-B020}{start}
که په~\LR{$\LogLuk[3]$} کښې~\LR{$\Undef$} هم ټاکلے وګرځوئ، نو څه کېږي؟
\olpseganchor{OLP-0397-B020}{end}

\olpseganchor{OLP-0397-B021}{start}
\begin{defn}
  \emph{دوه ټاکلي ارزښتونه لرونکے درې ارزښته لوکاشېوېچ منطق}~\LR{$\LogLuk[3]^{+}$} په دې ماتريس تعريفېږي:
  \begin{enumerate}
    \item معياري بياني ژبه~\LR{$\Lang L_0$} چې
    \LR{$\lfalse$}, \LR{$\lnot$}, \LR{$\land$}, \LR{$\lor$}, \LR{$\lif$} نښلوونکي لري۔
    \item د رښتياوالي د قيمتونو سټ~\LR{$V = \{\True, \Undef, \False\}$}۔
    \item \LR{$\True$} او~\LR{$\Undef$} دواړه ټاکل شوي دي؛ يعنې~\LR{$V^+ = \{\True, \Undef\}$}۔
    \item د دروغو ثابت~\LR{$\tf{\lfalse} = \False$} دے۔ د نورو نښلوونکو صدق تابعې د لوکاشېوېچ د منطق~\LR{$\LogLuk[3]$} له تابعو سره يو شان دي۔
  \end{enumerate}
\end{defn}
په \href{https://pure.unamur.be/ws/portalfiles/portal/53362118/2016_DegauquierV_article.pdf}{معياري درې ارزښته آر-مينګل} کښې د شرط تابعې \LR{$\tf{\lif}[\LogRM[3]](\Undef,\False)=\False$} وي۔
\olpseganchor{OLP-0397-B021}{end}

\olpseganchor{OLP-0397-B022}{start}
\begin{prob}
  له لاندې اړيکو کومې په~\LR{$\LogLuk[3]^{+}$} کښې صدق کوي؟
  \begin{enumerate}
    \item \LR{$p, p \lif q \Entails q$}
    \item \LR{$p \lor q, \lnot p \Entails q$}
    \item \LR{$\lnot p, p \Entails q$}
    \item \LR{$p \Entails p \lor q$}
  \end{enumerate}
\end{prob}
\olpseganchor{OLP-0397-B022}{end}

\olpseganchor{OLP-0397-B023}{start}
ځينې وخت بېلابېل صدق جدولونه يوازې د ټاکل شوو قيمتونو په بدلولو
سره هماغه يو منطق (هماغه د استلزام اړيکه) جوړولے شي۔ د بېلګې په
توګه، د ګوډل په منطق کښې که د~\LR{$\{\True\}$} پر ځای~\LR{$V^+ = \{\True,
\Undef\}$} واخلو، همداسې کېږي۔
\olpseganchor{OLP-0397-B023}{end}

\olpseganchor{OLP-0397-B024}{start}
\begin{prop}\ollabel{prop:gl-udes}
  هغه ماتريس چې~\LR{$V = \{\False, \Undef,\True\}$}، \LR{$V^+=\{\True,
  \Undef\}$} او د ګوډل د~\LR{$3$} ارزښته منطق صدق تابعې لري،
  کلاسيک منطق تعريفوي۔
\end{prop}
\olpseganchor{OLP-0397-B024}{end}

\olpseganchor{OLP-0397-B025}{start}
\begin{proof}
  تمرين۔
\end{proof}
\olpseganchor{OLP-0397-B025}{end}

\olpseganchor{OLP-0397-B026}{start}
\begin{prob}
  د~\olref[mvl][thr][mul]{prop:gl-udes} ثبوت ورکړئ۔ د دې دپاره
  هغه منطق~\LR{$\Log L$} وګورئ چې د ګوډل د منطق په شان تعريفېږي،
  خو~\LR{$V^+=\{\True,\Undef\}$} لري؛ وښيئ چې که
  \LR{$\Gamma \Entails/[\Log L] !B$} وي، نو~\LR{$\Gamma \Entails/[\LogCL] !B$}
  هم وي۔ د~\olref[mvl][thr][mul]{prop:LP-taut-CL} د ثبوت فکرونه
  وکاروئ، خو د (الف) او (ب) ځانګړنو د ثبوت پر ځای وښيئ چې
  \LR{$\pValue v(!A)[\LogGod] = \False$} هله او يوازې هله وي چې
  \LR{$\pValue {v'}(!A)[\LogCL] = \False$} وي (او له دې امله~\LR{$\pValue
  v(!A)[\LogGod] \in \{\True,\Undef\}$} هله او يوازې هله وي چې
  \LR{$\pValue {v'}(!A)[\LogCL] =
  \True$} وي)۔ روښانه کړئ چې دا ولې قضيه ثابتوي۔
\end{prob}
\olpseganchor{OLP-0397-B026}{end}

\olpseganchor{OLP-0398-B004}{start}
\olchapter{mvl}{inf}{نامتناهي ارزښته منطقونه}
\olpseganchor{OLP-0398-B004}{end}

\olpseganchor{OLP-0399-B004}{start}
\olfileid{mvl}{inf}{int}
\olpseganchor{OLP-0399-B004}{end}

\olpseganchor{OLP-0399-B005}{start}
\olsection{پېژندنه}
\olpseganchor{OLP-0399-B005}{end}

\olpseganchor{OLP-0399-B006}{start}
د يوه ماتريس د رښتياوالي د قيمتونو شمېر لازماً متناهي نۀ وي۔
د نامتناهي ډېرو رښتياوالي قيمتونو د سټ دپاره يوه څرګنده ټاکنه
د~\LR{$0$} او~\LR{$1$} تر منځ د ناطقو عددونو سټ دے، يعنې~\LR{$V_\infty = [0,1] \cap \Rat$}:
\begin{align*}
    V_\infty & = \Setabs{\frac{n}{m}}{n,m \in \Nat \text{\RL{ او }} 0<m \text{\RL{ او }} n\le m}.
\intertext{\RL{د رښتياوالي د دې نامتناهي سټ د کتنې په وخت کښې ډېر ځله
د لاندې فرعي سټونو کتل هم ګټور وي:}}
V_m & = \Setabs{\frac{n}{m-1}}{n \in \Nat \text{\RL{ او }} n<m}
\intertext{\RL{د بېلګې په توګه، \LR{$V_5$} د \LR{$5$} په برابر واټن پرتو رښتياوالي قيمتونو سټ دے:}}
V_5 & = \{0, \frac{1}{4}, \frac{1}{2}, \frac{3}{4}, 1\}.
\end{align*}
متناهي فرعي سټ په دويمه ښودنه کښې دوه يا تر دوو زيات قيمتونه لري؛
لومړے او وروستے ارزښت ئې په ترتيب صفر او يو دي۔
د دغو رښتياوالي قيمتونو پر سټونو ولاړ منطقونو کښې عموماً يوازې~\LR{$1$}
ټاکل شوے وي؛ يعنې~\LR{$V^+ = \{1\}$}۔ په بله وينا،~\LR{$1$} د (مطلق)
رښتيا او~\LR{$0$} د مطلق دروغ نقش لوبوي، خو فارمولونه په~\LR{$V$}
کښې هر منځنے قيمت هم اخيستلے شي۔
\olpseganchor{OLP-0399-B006}{end}

\olpseganchor{OLP-0399-B007}{start}
سره له دې، د~\LR{$0$} او~\LR{$1$} تر منځ د ټولو \emph{حقيقي} عددونو
سټ~\LR{$V_{[0,1]} = [0,1]$}، يا د~\LR{$[0,1]$} نور نامتناهي فرعي سټونه
هم کارولے شو۔ په دې رښتياوالي سټ ولاړ منطقونه ډېر ځله \emph{فازي}
بلل کېږي۔
\olpseganchor{OLP-0399-B007}{end}

\olpseganchor{OLP-0400-B004}{start}
\olfileid{mvl}{inf}{luk}
\olpseganchor{OLP-0400-B004}{end}

\olpseganchor{OLP-0400-B005}{start}
\olsection{د لوکاشېوېچ منطق}
\olpseganchor{OLP-0400-B005}{end}

\olpseganchor{OLP-0400-B006}{start}
\begin{editorial}
  دا برخه د لوکاشېوېچ د ناطقو ارزښتونو ماتريس، صدق تابعې او له متناهي ماتريسونو سره ئې اړيکه بيانوي۔
\end{editorial}
\olpseganchor{OLP-0400-B006}{end}

\olpseganchor{OLP-0400-B007}{start}
\begin{defn}\ollabel{def:lukasiewicz} نامتناهي ارزښته د لوکاشېوېچ
منطق~\LR{$\LogLuk[\infty]$} په دې ماتريس تعريفېږي:
\begin{enumerate}
  \item معياري بياني ژبه~\LR{$\Lang L_0$} چې
  \LR{$\lnot$}, \LR{$\land$}, \LR{$\lor$}, \LR{$\lif$} لري۔
  \item د رښتياوالي د قيمتونو سټ~\LR{$V_\infty$}۔
  \item يوازې~\LR{$1$} ټاکل شوے قيمت دے؛ يعنې~\LR{$V^+ = \{1\}$}۔
  \item صدق تابعې په لاندې تابعو ټاکل کېږي:
  \begin{align*}
    \tf{\lnot}[\LogLuk](x) & = 1 - x\\
    \tf{\land}[\LogLuk](x,y) & = \min(x,y)\\
    \tf{\lor}[\LogLuk](x,y) & = \max(x,y)\\
    \tf{\lif}[\LogLuk](x,y) & = \min(1,1-(x-y)) = \begin{cases}
      1 & \text{\RL{که }} x \le y\\
      1-(x-y) & \text{\RL{په بل حالت کښې.}}
    \end{cases}
    \end{align*}
\end{enumerate}
د لوکاشېوېچ \LR{$m$}-ارزښته منطق همداسې تعريفېږي، خو~\LR{$V = V_m$} لري۔
\end{defn}
\olpseganchor{OLP-0400-B007}{end}

\olpseganchor{OLP-0400-B008}{start}
\begin{prop}
  د~\olref[thr][luk]{def:lukasiewicz} له مخې تعريف شوے منطق~\LR{$\LogLuk[3]$}
  د~\olref{def:lukasiewicz} له مخې له تعريف شوي~\LR{$\LogLuk[3]$} سره يو شان دے۔
\end{prop}
\olpseganchor{OLP-0400-B008}{end}

\olpseganchor{OLP-0400-B009}{start}
\begin{proof}
  دا د~\olref[thr][luk]{def:lukasiewicz} د نښلوونکو صدق جدولونه
  د~\olref{def:lukasiewicz} له معادلو څخه له ترلاسه شوو جدولونو سره
  په پرتله کولو ليدل کېدے شي:
  \begin{center}
    \begin{tabular}{c|c} 
      \LR{$\tf{\lnot}$} & \\ 
      \hline  
      \LR{$1$} & \LR{$0$} \\ 
      \LR{$1/2$} & \LR{$1/2$} \\
      \LR{$0$} & \LR{$1$} 
    \end{tabular}
    \quad
    \begin{tabular}{c|ccc} 
      \LR{$\tf{\land}[\LogLuk[3]]$} & \LR{$1$} & \LR{$1/2$} & \LR{$0$} \\ 
      \hline 
      \LR{$1$} & \LR{$1$} & \LR{$1/2$} & \LR{$0$} \\ 
      \LR{$1/2$} & \LR{$1/2$} & \LR{$1/2$} & \LR{$0$}\\ 
      \LR{$0$} & \LR{$0$} & \LR{$0$} & \LR{$0$} 
    \end{tabular}
    \\[2ex]
    \begin{tabular}{c|ccc} 
      \LR{$\tf{\lor}[\LogLuk[3]]$} & \LR{$1$} & \LR{$1/2$} & \LR{$0$} \\ 
      \hline 
      \LR{$1$} & \LR{$1$} & \LR{$1$} & \LR{$1$} \\ 
      \LR{$1/2$} & \LR{$1$} & \LR{$1/2$} & \LR{$1/2$} \\
      \LR{$0$} & \LR{$1$} & \LR{$1/2$} & \LR{$0$} 
    \end{tabular}
    \quad
    \begin{tabular}{c|ccc} 
      \LR{$\tf{\lif}[\LogLuk[3]]$} & \LR{$1$} & \LR{$1/2$} & \LR{$0$} \\ 
      \hline 
      \LR{$1$} & \LR{$1$} & \LR{$1/2$} & \LR{$0$} \\ 
      \LR{$1/2$} & \LR{$1$} & \LR{$1$} & \LR{$1/2$}  \\ 
      \LR{$0$} & \LR{$1$} & \LR{$1$} & \LR{$1$} 
    \end{tabular}
  \end{center} 
\end{proof}
\olpseganchor{OLP-0400-B009}{end}

\olpseganchor{OLP-0400-B010}{start}
\begin{prop}\ollabel{prop:luk-infty-m}
  که~\LR{$\Gamma \Entails[\LogLuk[\infty]] !B$} وي، نو د هر~\LR{$m \ge 2$}
  دپاره~\LR{$\Gamma \Entails[\LogLuk[m]] !B$} هم وي۔
\end{prop}
\olpseganchor{OLP-0400-B010}{end}

\olpseganchor{OLP-0400-B011}{start}
\begin{proof}
  تمرين۔
\end{proof}
\olpseganchor{OLP-0400-B011}{end}

\olpseganchor{OLP-0400-B012}{start}
\begin{prob}
  \olref[mvl][inf][luk]{prop:luk-infty-m} ثابت کړئ۔
\end{prob}
\olpseganchor{OLP-0400-B012}{end}

\olpseganchor{OLP-0400-B013}{start}
که په~\LR{$\Gamma$} او~\LR{$!B$} کښې يوازې متناهي ډېر بياني متغيرونه
راشي، نو برعکس لورے هم صدق کوي۔ که يوه ناطقه ارزښت ټاکنه ټولې
مقدمې پوره کړي او پايله پوره نۀ کړي، د هغو متغيرونو د قيمتونو
يو ګډ مخرج واخلو۔ ټول دا قيمتونه په يوه متناهي سټ کښې راځي، او
هلته هم صدق تابعې هماغه دي؛ نو هماغه مخالف مثال په متناهي
ماتريس کښې هم شته۔
د نامتناهي مقدماتو لپاره عمومي برعکس لورے ناسم دے۔ د
\LR{$!A \oplus !B$} لنډيز د~\LR{$\lnot !A \lif !B$} لپاره وکاروئ او
د~\LR{$S_n(p)$} په ليکدود د~\LR{$p$} د~\LR{$n$} نقلونو دغه مجموعه وښيئ؛
\LR{$S_1(p)=p$}۔ که د~\LR{$p$} قيمت~\LR{$x$} وي، د دې مجموعې قيمت
\LR{$\min(1,nx)$} دے۔ د هر مثبت صحيح~\LR{$n$} لپاره دا مقدمه واخلئ:
\[
(p_n \lif \lnot S_n(p_n)) \land (\lnot S_n(p_n) \lif p_n).
\]
دا مقدمه هله رښتيا ده چې د دواړو خواوو قيمتونه برابر وي، يعنې
\LR{$\pAssign v(p_n)=1/(n+1)$}۔ په ناطق ماتريس کښې ټولې مقدمې
پوره کېدے شي، او جلا متغير~\LR{$q$} ته صفر ورکولے شو۔ خو په
\LR{$V_m$} کښې د~\LR{$n=m-1$} مقدمه د~\LR{$1/m$} قيمت غواړي، چې په دغه
سټ کښې نشته۔ نو هر متناهي ماتريس له دغو مقدماتو څخه~\LR{$q$}
لازموي، خو نامتناهي ناطق ماتريس ئې نۀ لازموي۔
\olpseganchor{OLP-0400-B013}{end}

\olpseganchor{OLP-0400-B014}{start}
نامتناهي ارزښته د لوکاشېوېچ منطق يو پېژندل شوے فازي منطق دے۔
د فازي منطق په ليکنو کښې شرطي نښلوونکے ډېر ځله په~\LR{$\lnot !A \lor !B$}
تعريفېږي۔ په دې سره به يو نامتناهي ارزښته قوي کليني منطق ترلاسه شي۔
\olpseganchor{OLP-0400-B014}{end}

\olpseganchor{OLP-0400-B015}{start}
\begin{prob}
  وښيئ چې~\LR{$(p \lif q) \lor (q \lif p)$} د~\LR{$\LogLuk[\infty]$}
  تاوتولوژي ده۔
\end{prob}
\olpseganchor{OLP-0400-B015}{end}

\olpseganchor{OLP-0401-B004}{start}
\olfileid{mvl}{inf}{god}
\olpseganchor{OLP-0401-B004}{end}

\olpseganchor{OLP-0401-B005}{start}
\olsection{د ګوډل منطقونه}
\olpseganchor{OLP-0401-B005}{end}

\olpseganchor{OLP-0401-B006}{start}
\begin{editorial}
  دا برخه د ګوډل ناطق ماتريس، صدق تابعې او له متناهي زنځيرونو او شهودي منطق سره ئې اړيکه بيانوي۔
\end{editorial}
\olpseganchor{OLP-0401-B006}{end}

\olpseganchor{OLP-0401-B007}{start}
\begin{defn}\ollabel{def:goedel} نامتناهي ارزښته د ګوډل
منطق~\LR{$\LogGod[\infty]$} په دې ماتريس تعريفېږي:
\begin{enumerate}
  \item معياري بياني ژبه~\LR{$\Lang L_0$} چې
  \LR{$\lfalse$}, \LR{$\lnot$}, \LR{$\land$}, \LR{$\lor$}, \LR{$\lif$} لري۔
  \item د رښتياوالي د قيمتونو سټ~\LR{$V_\infty$}۔
  \item يوازې~\LR{$1$} ټاکل شوے قيمت دے؛ يعنې~\LR{$V^+ = \{1\}$}۔
  \item صدق تابعې په لاندې تابعو ټاکل کېږي:
  \begin{align*}
    \tf{\lfalse} & = 0\\
    \tf{\lnot}[\LogGod](x) & = \begin{cases}
      1 & \text{\RL{که }} x =0\\
      0 & \text{\RL{په بل حالت کښې}}
    \end{cases}\\
    \tf{\land}[\LogGod](x,y) & = \min(x,y)\\
    \tf{\lor}[\LogGod](x,y) & = \max(x,y)\\
    \tf{\lif}[\LogGod](x,y) & = \begin{cases}
      1 & \text{\RL{که }} x \le y\\
      y & \text{\RL{په بل حالت کښې.}}
    \end{cases}
    \end{align*}
\end{enumerate}
د ګوډل \LR{$m$}-ارزښته منطق همداسې تعريفېږي، خو~\LR{$V = V_m$} لري۔
\end{defn}
\olpseganchor{OLP-0401-B007}{end}

\olpseganchor{OLP-0401-B008}{start}
\begin{prop}
  د~\olref[thr][god]{defn:goedel} له مخې تعريف شوے منطق~\LR{$\LogGod[3]$}
  د~\olref{def:goedel} له مخې له تعريف شوي~\LR{$\LogGod[3]$} سره يو شان دے۔
\end{prop}
\olpseganchor{OLP-0401-B008}{end}

\olpseganchor{OLP-0401-B009}{start}
\begin{proof}
  دا د~\olref[thr][god]{defn:goedel} د نښلوونکو صدق جدولونه
  د~\olref{def:goedel} له معادلو څخه له ترلاسه شوو جدولونو سره
  په پرتله کولو ليدل کېدے شي:
  \begin{center}
    \begin{tabular}{c|c} 
      \LR{$\tf{\lnot}[\LogGod[3]]$} & \\ 
      \hline  
      \LR{$1$} & \LR{$0$} \\ 
      \LR{$1/2$} & \LR{$0$} \\
      \LR{$0$} & \LR{$1$} 
    \end{tabular}
    \quad
    \begin{tabular}{c|ccc} 
      \LR{$\tf{\land}[\LogGod]$} & \LR{$1$} & \LR{$1/2$} & \LR{$0$} \\ 
      \hline 
      \LR{$1$} & \LR{$1$} & \LR{$1/2$} & \LR{$0$} \\ 
      \LR{$1/2$} & \LR{$1/2$} & \LR{$1/2$} & \LR{$0$}\\ 
      \LR{$0$} & \LR{$0$} & \LR{$0$} & \LR{$0$} 
    \end{tabular}
    \\[2ex]
    \begin{tabular}{c|ccc} 
      \LR{$\tf{\lor}[\LogGod]$} & \LR{$1$} & \LR{$1/2$} & \LR{$0$} \\ 
      \hline 
      \LR{$1$} & \LR{$1$} & \LR{$1$} & \LR{$1$} \\ 
      \LR{$1/2$} & \LR{$1$} & \LR{$1/2$} & \LR{$1/2$} \\
      \LR{$0$} & \LR{$1$} & \LR{$1/2$} & \LR{$0$} 
    \end{tabular}
    \quad
    \begin{tabular}{c|ccc} 
      \LR{$\tf{\lif}[\LogGod]$} & \LR{$1$} & \LR{$1/2$} & \LR{$0$} \\ 
      \hline 
      \LR{$1$} & \LR{$1$} & \LR{$1/2$} & \LR{$0$} \\ 
      \LR{$1/2$} & \LR{$1$} & \LR{$1$} & \LR{$0$}  \\ 
      \LR{$0$} & \LR{$1$} & \LR{$1$} & \LR{$1$} 
    \end{tabular}
  \end{center} 
\end{proof}
\olpseganchor{OLP-0401-B009}{end}

\olpseganchor{OLP-0401-B010}{start}
\begin{prop}\ollabel{prop:god-infty-m}
  که~\LR{$\Gamma \Entails[\LogGod[\infty]] !B$} وي، نو د هر~\LR{$m \ge 2$}
  دپاره~\LR{$\Gamma \Entails[\LogGod[m]] !B$} هم وي۔
\end{prop}
\olpseganchor{OLP-0401-B010}{end}

\olpseganchor{OLP-0401-B011}{start}
\begin{proof}
  تمرين۔
\end{proof}
\olpseganchor{OLP-0401-B011}{end}

\olpseganchor{OLP-0401-B012}{start}
\begin{prob}
  \olref[mvl][inf][god]{prop:god-infty-m} ثابت کړئ۔
\end{prob}
\olpseganchor{OLP-0401-B012}{end}

\olpseganchor{OLP-0401-B013}{start}
که په~\LR{$\Gamma$} او~\LR{$!B$} کښې يوازې متناهي ډېر بياني متغيرونه
راشي، نو برعکس لورے هم صدق کوي۔ د يو نامتناهي مخالف مثال د
متغيرونو قيمتونه له صفر او يو سره يوځاے کړئ او په يوه متناهي
زنځير کښې په هماغه ترتيب کښېږدئ۔ د صفر، يو او ترتيب په ساتلو
سره نفي، عطف، فصل او شرط هم ساتل کېږي؛ نو مخالف مثال متناهي
ماتريس ته لېږدېږي۔
د نامتناهي مقدماتو لپاره عمومي برعکس لورے ناسم دے۔ د هر
\LR{$n\ge 0$} لپاره دا مقدمه واخلئ:
\[
(p_n \lif p_{n+1}) \lif q.
\]
که~\LR{$\pAssign v(q)<1$} وي او ټولې مقدمې پوره وي، د هر~\LR{$n$}
لپاره~\LR{$\pAssign v(p_n)>\pAssign v(p_{n+1})$} پکار دي؛ که داسې
نۀ وي، دنننے شرط يو کېږي او بهرنۍ مقدمه د~\LR{$q$} هماغه قيمت
اخلي۔ په متناهي زنځير کښې دا نامتناهي نزولي لړۍ ناشونې ده۔
خو په ناطق ماتريس کښې~\LR{$\pAssign v(p_n)=1/(n+2)$} او
\LR{$\pAssign v(q)=1/2$} ټولې مقدمې پوره کوي، ځکه دننني شرطونه
د~\LR{$q$} تر قيمت زيات نۀ دي۔ دلته~\LR{$q$} پوره نۀ دے۔
\olpseganchor{OLP-0401-B013}{end}

\olpseganchor{OLP-0401-B014}{start}
د~\LR{$\LogGod[3]$} په شان، په~\LR{$\LogGod[\infty]$} کښې په شهودي منطق کښې
ټول معتبر فارمولونه تاوتولوژۍ دي، او هماغه د ناتاوتولوژيو بېلګې
په~\LR{$\LogGod[\infty]$} کښې هم ناتاوتولوژۍ دي:
\begin{align*}
  & p \lor \lnot p && (p \lif q) \lif (\lnot p \lor q) \\
  & \lnot\lnot p \lif p && \lnot(\lnot p \land \lnot q) \lif (p \lor q) \\
  & ((p \lif q) \lif p) \lif p && \lnot(p \lif q) \lif (p \land \lnot q)
\end{align*}
يو داسې فارمول چې په شهودي منطق کښې معتبر نۀ دے، خو
په~\LR{$\LogGod[3]$} کښې تاوتولوژي ده، يعنې~\LR{$(p \lif q) \lor (q \lif
p)$}، په~\LR{$\LogGod[\infty]$} کښې هم تاوتولوژي ده۔ په حقيقت کښې
\LR{$\LogGod[\infty]$} د شهودي منطق په توګه هم پېژندل کېدے شي چې
شېما~\LR{$(!A \lif !B) \lor (!B \lif !A)$} ورزياته شوې وي۔ دا مايکل
ډمېټ وښودله؛ نو~\LR{$\LogGod[\infty]$} ته ډېر ځله د ګوډل--ډمېټ
منطق~\LR{$\Log{LC}$} ويل کېږي۔
\olpseganchor{OLP-0401-B014}{end}

\olpseganchor{OLP-0401-B015}{start}
\begin{prob}
  وښيئ چې~\LR{$(p \lif q) \lor (q \lif p)$} د~\LR{$\LogGod[\infty]$}
  تاوتولوژي ده۔
\end{prob}
\olpseganchor{OLP-0401-B015}{end}

\olpseganchor{OLP-0401-B016}{start}
\begin{prob}
  وښيئ چې~\LR{$(p \lif q) \lor (q \lif r) \lor (r \lif s)$}، چې
  په~\LR{$\LogGod[3]$} کښې تاوتولوژي ده، په~\LR{$\LogGod[\infty]$}
  کښې تاوتولوژي نۀ ده۔
\end{prob}
\olpseganchor{OLP-0401-B016}{end}

\olpseganchor{OLP-0402-B004}{start}
\olchapter{mvl}{seq}{د سېکوېنټ حساب}
\olpseganchor{OLP-0402-B004}{end}

\olpseganchor{OLP-0403-B004}{start}
\olfileid{mvl}{seq}{int}
\olpseganchor{OLP-0403-B004}{end}

\olpseganchor{OLP-0403-B005}{start}
\olsection{پېژندنه}
\olpseganchor{OLP-0403-B005}{end}

\olpseganchor{OLP-0403-B006}{start}
د کلاسيک منطق د سېکوېنټ حساب د اشتقاق يو اغېزمن او ساده نظام
دے۔ که ډېرارزښته منطق په داسې ماتريس تعريف شوے وي چې د رښتياوالي
قيمتونه ئې متناهي وي، يعنې~\LR{$V$} متناهي وي، د هغه دپاره سېکوېنټ حساب
ورکول کېدے شي۔ د دې کار مفکوره د کلاسيک حالت د سېکوېنټونو له معناوو
او د استنتاج د قاعدو له بڼې څخه راځي۔
\olpseganchor{OLP-0403-B006}{end}

\olpseganchor{OLP-0403-B007}{start}
اوس راياد کړئ چې يو سېکوېنټ
\begin{align*}
    !A_1, \dots, !A_m & \Sequent !B_1, \dots, !B_n
\intertext{\RL{د لاندې فارمول په توګه تعبيرېدے شي:}}
    (!A_1 \land \cdots \land !A_m) & \lif (!B_1 \lor \cdots \lor
!B_n)
\end{align*}
په نورو ټکو کښې، ارزښت ټاکنه~\LR{$\pAssign{v}$}
يو سېکوېنټ~\LR{$\Gamma \Sequent \Delta$} هله او يوازې هله \emph{پوره کوي} چې
يا د کوم~\LR{$!A \in \Gamma$} دپاره~\LR{$\pValue{v}(!A) = \False$}
وي، يا د کوم~\LR{$!A \in \Delta$} دپاره~\LR{$\pValue{v}(!A) = \True$}
وي۔ د دې تعبير له مخې ابتدايي سېکوېنټونه~\LR{$!A \Sequent !A$} تل پوره
کېږي، ځکه چې يا~\LR{$\pValue v(!A) = \True$} وي يا~\LR{$\pValue v(!A) =
\False$} وي۔
\olpseganchor{OLP-0403-B007}{end}

\olpseganchor{OLP-0403-B008}{start}
دا په~\LR{$\Log{LK}$} کښې د شرطي نښلوونکي د استنتاج قاعدې دي، چې د څنګ
فارمولونه~\LR{$\Gamma$} او~\LR{$\Delta$} ترې اېستل شوي دي:
\olpseganchor{OLP-0403-B008}{end}

\olpseganchor{OLP-0403-B009}{start}
\begin{defish}
    \Axiom$ \fCenter !A$
    \Axiom$ !B \fCenter $
    \RightLabel{\LeftR{\lif}}
    \BinaryInf$ !A \lif !B \fCenter $
    \DisplayProof
    \hfill
    \Axiom$ !A \fCenter !B$
    \RightLabel{\RightR{\lif}}
    \UnaryInf$ \fCenter !A \lif !B $
    \DisplayProof
\end{defish}
\olpseganchor{OLP-0403-B009}{end}

\olpseganchor{OLP-0403-B010}{start}
که د سېکوېنټ پورته معنايي تعبير وکاروو، د~\LR{$\LeftR{\lif}$} قاعده
داسې لوستلے شو: که~\LR{$\pValue v(!A) = \True$} او~\LR{$\pValue v(!B) =
\False$} وي، نو~\LR{$\pValue v(!A \lif !B) = \False$} وي۔ همداسې،
د~\LR{$\RightR{\lif}$} قاعده وايي چې که يا~\LR{$\pValue v(!A) = \False$}
وي يا~\LR{$\pValue v(!B) = \True$} وي، نو~\LR{$\pValue v(!A \lif !B) =
\True$} وي۔ په حقيقت کښې دا شرطي بيانونه دوه‌لوريز هم دي۔ د
\LR{$\LeftR{\land}$} او~\LR{$\RightR{\lor}$} د قاعدو په معياري بڼه کښې اړوند
شرطي بيانونه دوه‌لوريز نۀ وي۔ خو د دغو قاعدو بديلې بڼې شته چې
دوه‌لوريز دي:
\olpseganchor{OLP-0403-B010}{end}

\olpseganchor{OLP-0403-B011}{start}
\begin{defish}
    \Axiom$!A, !B, \Gamma \fCenter \Delta$
    \RightLabel{\LeftR{\land}}
    \UnaryInf$!A \land !B, \Gamma \fCenter \Delta$
    \DisplayProof
    \hfill
    \Axiom$ \Gamma \fCenter \Delta, !A, !B$
    \RightLabel{\RightR{\lor}}
    \UnaryInf$ \Gamma \fCenter \Delta, !A \lor !B$
    \DisplayProof
\end{defish}
\olpseganchor{OLP-0403-B011}{end}

\olpseganchor{OLP-0403-B012}{start}
دا بنسټيزه مفکوره چې پر يو~\LR{$n$} ارزښته منطق پلې شي، بيا داسې
سېکوېنټ حساب ورکوي چې د دوو پر ځای~\LR{$n$} ځايونه لري، د رښتياوالي د
هر قيمت دپاره يو ځاے۔ د درې‌ارزښته منطق دپاره، چې~\LR{$V = \{\False,
\Undef, \True\}$} لري، سېکوېنټ د~\LR{$\Gamma \mid \Pi \mid \Delta$}
په بڼه عبارت دے۔ دا په~ارزښت ټاکنه~\LR{$\pAssign v$} کښې هله او يوازې هله پوره کېږي
چې يا د کوم~\LR{$!A \in \Gamma$} دپاره~\LR{$\pValue{v}(!A) =
\False$} وي، يا د کوم~\LR{$!A \in \Delta$} دپاره~\LR{$\pValue{v}(!A)
= \True$} وي، يا د کوم~\LR{$!A \in \Pi$} دپاره~\LR{$\pValue{v}(!A) =
\Undef$} وي۔ نو ابتدايي سېکوېنټونه~\LR{$!A \mid !A \mid !A$} تل پوره
کېږي۔
\olpseganchor{OLP-0403-B012}{end}

\olpseganchor{OLP-0404-B004}{start}
\olfileid{mvl}{seq}{rul}
\olpseganchor{OLP-0404-B004}{end}

\olpseganchor{OLP-0404-B005}{start}
\olsection{قاعدې او اشتقاقونه}
\olpseganchor{OLP-0404-B005}{end}

\olpseganchor{OLP-0404-B006}{start}
په لاندې بحث کښې دې~\LR{$\Gamma, \Delta, \Pi, \Lambda$} د
تړلي فارمولونه متناهي لړۍ تمثيل کړي۔
\olpseganchor{OLP-0404-B006}{end}

\olpseganchor{OLP-0404-B007}{start}
\begin{defn}[سېکوېنټ]
يو \emph{\LR{$n$} اړخيز سېکوېنټ} د لاندې بڼې عبارت دے:
\[
\Gamma_1 \nSequent \dots \nSequent \Gamma_n
\]
چې هره~\LR{$\Gamma_i$} د ژبې~\LR{$\Lang L$} د تړلي فارمولونه يوه متناهي
(کېدے شي تشه) لړۍ ده۔
\end{defn}
\olpseganchor{OLP-0404-B007}{end}

\olpseganchor{OLP-0404-B008}{start}
\begin{defn}[ابتدايي سېکوېنټ]
يو \emph{\LR{$n$} اړخيز ابتدايي سېکوېنټ} د هرې ژبنۍ تړلے فارمول~\LR{$!A$}
دپاره د~\LR{$!A \nSequent \dots \nSequent !A$} په بڼه~\LR{$n$} اړخيز
سېکوېنټ دے۔
\olpseganchor{OLP-0404-B008}{end}

\olpseganchor{OLP-0404-B009}{start}
که ژبه يو~\LR{$0$} ځايه نښلوونکے~\LR{$\star$}، يعنې بياني ثابت، ولري، نو
هغه سېکوېنټ~\LR{$\dots \nSequent \star \nSequent \dots$} هم ابتدايي
ګڼو چې~\LR{$\star$} پکې د~\LR{$\tf{\star} \in V$} له اړوند رښتياوالي
قيمت سره په برابر ځاے کښې راځي او نور ځايونه ئې تش وي۔
\end{defn}
\olpseganchor{OLP-0404-B009}{end}

\olpseganchor{OLP-0404-B010}{start}
د~\LR{$n$} ارزښته منطق~\LR{$\Log{L}$} د هر نښلوونکي دپاره، په~\LR{$\Log{L}$} کښې
د هغه د هر ممکن رښتياوالي قيمت يوه منطقي قاعده شته۔ د~\LR{$\Log{L}$} د~\LR{$n$} اړخيز
سېکوېنټ حساب اشتقاقونه د سېکوېنټونو ونې دي: تر ټولو بره
سېکوېنټونه ابتدايي وي، او که يو سېکوېنټ د يوه يا څو نورو لاندې
ولاړ وي، نو بايد د~\LR{$\Log{L}$} د سېکوېنټ حساب د يوې منطقي يا
جوړښتي قاعدې له مخې په سمه توګه ترې راوځي۔
\olpseganchor{OLP-0404-B010}{end}

\olpseganchor{OLP-0404-B011}{start}
\begin{defn}[قضيې]
تړلے فارمول~\LR{$!A$} د~\LR{$n$} ارزښته منطق~\LR{$\Log{L}$} \emph{قضيه}
ده که د هغه~\LR{$n$} اړخيز سېکوېنټ اشتقاق وي چې د~\LR{$\Log{L}$}
د هر ټاکل شوي رښتياوالي قيمت په اړوند ځاے کښې يوازې~\LR{$!A$} لري
او نور ځايونه ئې تش دي۔ که~\LR{$!A$}
قضيه وي،~\LR{$\Proves[\Log{L}] !A$} ليکو، او که نۀ وي،~\LR{$\Proves/[\Log{L}]
!A$} ليکو۔
\end{defn}
\olpseganchor{OLP-0404-B011}{end}

\olpseganchor{OLP-0404-B012}{start}
\begin{defn}[د اشتقاق وړتيا]
تړلے فارمول~\LR{$!A$} د~\LR{$n$} ارزښته منطق~\LR{$\Log{L}$} کښې د
تړلي فارمولونه له سټ~\LR{$\Gamma$} څخه \emph{د اشتقاق وړ ده}، يعنې~\LR{$\Gamma
\Proves[\Log{L}] !A$}، هله او يوازې هله چې يو متناهي
فرعي سټ~\LR{$\Gamma_0 \subseteq \Gamma$} او د~\LR{$\Gamma_0$} د
تړلي فارمولونه يوه لړۍ~\LR{$\Gamma_0'$} وي چې لاندې سېکوېنټ
اشتقاق ولري:
\[ \Lambda_1 \nSequent \dots \nSequent \Lambda_n \]
دلته~\LR{$\Lambda_i$} هغه مهال~\LR{$!A$} دے چې~\LR{$i$} ځاے له ټاکل شوي
رښتياوالي قيمت سره برابر وي، او په بل حالت کښې~\LR{$\Gamma_0'$} دے۔
که~\LR{$!A$} له~\LR{$\Gamma$} څخه د اشتقاق وړ نۀ وي،~\LR{$\Gamma \Proves/ !A$}
ليکو۔
\end{defn}
\olpseganchor{OLP-0404-B012}{end}

\olpseganchor{OLP-0404-B013}{start}
د بېلګې په توګه، د لوکاشېوېچ~\LR{$3$} ارزښته منطق يو~\LR{$3$} اړخيز
سېکوېنټ حساب لري۔ په~\LR{$3$} اړخيز سېکوېنټ~\LR{$\Gamma \nSequent \Pi
\nSequent \Delta$} کښې~\LR{$\Gamma$} له~\LR{$\False$}،~\LR{$\Delta$} له~\LR{$\True$}
او~\LR{$\Pi$} له~\LR{$\Undef$} سره برابر دي۔ بديهي سېکوېنټونه~\LR{$!A
\nSequent !A \nSequent !A$} دي۔ ځکه چې يوازې~\LR{$\True$} ټاکل شوے دے،
\LR{$\Gamma \Proves[\LogLuk[3]] !A$} هله او يوازې هله چې سېکوېنټ~\LR{$\Gamma
\nSequent \Gamma \nSequent !A$} يو~اشتقاق ولري۔ (که~\LR{$\Undef$}
هم ټاکل شوے وے، نو د~\LR{$\Gamma \nSequent !A \nSequent !A$}
اشتقاق به پکار و۔)
\olpseganchor{OLP-0404-B013}{end}

\olpseganchor{OLP-0405-B004}{start}
\olfileid{mvl}{seq}{str}
\olpseganchor{OLP-0405-B004}{end}

\olpseganchor{OLP-0405-B005}{start}
\olsection{جوړښتي قاعدې}
\olpseganchor{OLP-0405-B005}{end}

\olpseganchor{OLP-0405-B006}{start}
د~\LR{$n$} اړخيز سېکوېنټ حساب جوړښتي قاعدې د کلاسيک حالت په شان
کار کوي، خو د هر ځاے~\LR{$i$} دپاره بېلې پلې کېږي۔
\olpseganchor{OLP-0405-B006}{end}

\olpseganchor{OLP-0405-B007}{start}
\begin{defish}
\begin{center}
\AxiomC{\LR{$ \Gamma_1 \nSequent \dots \nSequent \phantom{!A,}\Gamma_i \nSequent \dots \nSequent \Gamma_n $}}
\RightLabel{\LR{$\iR{\Weakening}{i}$}}
\UnaryInfC{\LR{$\Gamma_1 \nSequent \dots \nSequent !A, \Gamma_i \nSequent \dots \nSequent \Gamma_n$}}
\DisplayProof
\\[2ex]
\AxiomC{\LR{$ \Gamma_1 \nSequent \dots \nSequent !A, !A, \Gamma_i \nSequent \dots \nSequent \Gamma_n $}}
\RightLabel{\LR{$\iR{\Contraction}{i}$}}
\UnaryInfC{\LR{$\Gamma_1 \nSequent \dots \nSequent \phantom{!A,}!A, \Gamma_i \nSequent
\dots \nSequent \Gamma_n$}}
\DisplayProof
\\[2ex]
\AxiomC{\LR{$ \Gamma_1 \nSequent \dots \nSequent \Gamma_i, !A, !B, \Gamma_i' \nSequent \dots \nSequent \Gamma_n $}}
\RightLabel{\LR{$\iR{\Exchange}{i}$}}
\UnaryInfC{\LR{$\Gamma_1 \nSequent \dots \nSequent \Gamma_i, !B, !A, \Gamma_i' \nSequent \dots
\nSequent \Gamma_n$}}
\DisplayProof
\end{center}
\end{defish}
\olpseganchor{OLP-0405-B007}{end}

\olpseganchor{OLP-0405-B008}{start}
د کمزوري کولو، انقباض او تبادلې د استنتاجونو يوه لړۍ ډېر ځله
په دوه‌ګونو استنتاجي کرښو ښودل کېږي۔
\olpseganchor{OLP-0405-B008}{end}

\olpseganchor{OLP-0405-B009}{start}
د~\Cut{} قاعده څو بڼې لري، د سېکوېنټ د بېلابېلو ځايونو د هرې
جوړې دپاره يوه، چې~\LR{$i \neq j$} وي:
\begin{defish}
\[
\AxiomC{\LR{$ \Gamma_1 \nSequent \dots \nSequent !A, \Gamma_i \nSequent \dots \nSequent \Gamma_n $}}
\AxiomC{\LR{$ \Delta_1 \nSequent \dots \nSequent !A, \Delta_j \nSequent \dots \nSequent \Delta_n $}}
\RightLabel{\LR{$\iR{\Cut}{i,j}$}}
\BinaryInfC{\LR{$\Gamma_1,\Delta_1 \nSequent \dots \nSequent \Gamma_n, \Delta_n$}}
\DisplayProof
\]
\end{defish}
\olpseganchor{OLP-0405-B009}{end}

\olpseganchor{OLP-0406-B004}{start}
\olfileid{mvl}{seq}{prl}
\olpseganchor{OLP-0406-B004}{end}

\olpseganchor{OLP-0406-B005}{start}
\olsection{د ټاکلو منطقونو بياني قاعدې}
\olpseganchor{OLP-0406-B005}{end}

\olpseganchor{OLP-0406-B006}{start}
په~\LR{$n$} اړخيز سېکوېنټ حساب کښې د يوه نښلوونکي د استنتاج قاعدې
يوازې د هغه نښلوونکي د رښتياوالي پر ځانګړې تابع پورې اړه لري۔
نو که يو نښلوونکے په بېلابېلو منطقونو کښې د رښتياوالي د هماغې تابع
له مخې تعريف شوے وي، د هغه د~\LR{$n$} اړخيز سېکوېنټ قاعدې هم په دغو
منطقونو کښې يو شان وي۔
\olpseganchor{OLP-0406-B006}{end}

\olpseganchor{OLP-0406-B007}{start}
\subsection{د~\texorpdfstring{\LR{$\lnot$}}{¬} قاعدې}
\olpseganchor{OLP-0406-B007}{end}

\olpseganchor{OLP-0406-B008}{start}
لاندې د~\LR{$\lnot$} قاعدې د لوکاشېوېچ او کليني په منطقونو او د
هغو په بېلابېلو بڼو کښې پلې کېږي۔
\olpseganchor{OLP-0406-B008}{end}

\olpseganchor{OLP-0406-B009}{start}
\begin{defish}
  \begin{center}
\AxiomC{\LR{$ \Gamma \nSequent \Pi \nSequent \Delta, !A $}}
\RightLabel{\iR{\lnot}{\False}}
\UnaryInfC{\LR{$ \lnot !A, \Gamma \nSequent \Pi \nSequent \Delta$}}
\DisplayProof
\\[2ex]
\AxiomC{\LR{$ \Gamma \nSequent !A, \Pi \nSequent \Delta $}}
\RightLabel{\iR{\lnot}{\Undef}}
\UnaryInfC{\LR{$\Gamma \nSequent \lnot !A, \Pi \nSequent \Delta$}}
\DisplayProof
\\[2ex]
\AxiomC{\LR{$!A, \Gamma \nSequent \Pi \nSequent \Delta$}}
\RightLabel{\iR{\lnot}{\True}}
\UnaryInfC{\LR{$\Gamma \nSequent \Pi \nSequent \Delta,  \lnot !A$}}
\DisplayProof
  \end{center}
\end{defish}
\olpseganchor{OLP-0406-B009}{end}

\olpseganchor{OLP-0406-B010}{start}
لاندې د~\LR{$\lnot$} قاعدې د ګوډل په منطق کښې پلې کېږي۔
\olpseganchor{OLP-0406-B010}{end}

\olpseganchor{OLP-0406-B011}{start}
\begin{defish}
\AxiomC{\LR{$ \Gamma \nSequent !A, \Pi \nSequent \Delta, !A $}}
\RightLabel{\iR{\lnot}{\False}[\LogGod]}
\UnaryInfC{\LR{$ \lnot !A, \Gamma \nSequent \Pi \nSequent \Delta$}}
\DisplayProof
\hfill
\AxiomC{\LR{$!A, \Gamma \nSequent \Pi \nSequent \Delta$}}
\RightLabel{\iR{\lnot}{\True}[\LogGod]}
\UnaryInfC{\LR{$\Gamma \nSequent \Pi \nSequent \Delta,  \lnot !A$}}
\DisplayProof
\hfill
\end{defish}
\olpseganchor{OLP-0406-B011}{end}

\olpseganchor{OLP-0406-B012}{start}
(د ګوډل په منطق کښې~\LR{$\lnot !A$} هېڅکله~\LR{$\Undef$} قيمت نشي اخيستلے؛
نو د منځني ځاے دپاره قاعده نشته۔)
\olpseganchor{OLP-0406-B012}{end}

\olpseganchor{OLP-0406-B013}{start}
\subsection{د~\texorpdfstring{\LR{$\land$}}{∧} قاعدې}
\olpseganchor{OLP-0406-B013}{end}

\olpseganchor{OLP-0406-B014}{start}
دا د~\LR{$\land$} قاعدې د لوکاشېوېچ، د کليني د قوي منطق او د ګوډل
د منطق دپاره دي۔
\olpseganchor{OLP-0406-B014}{end}

\olpseganchor{OLP-0406-B015}{start}
\begin{defish}
\begin{center}
  \AxiomC{\LR{$!A, !B, \Gamma \nSequent \Pi \nSequent \Delta$}}
  \RightLabel{\LR{$\iR{\land}{\False}$}}
  \UnaryInfC{\LR{$!A \land !B, \Gamma \nSequent \Pi \nSequent \Delta$}}
  \DisplayProof
\\[2ex]
  \AxiomC{\LR{$\Gamma \nSequent !A, \Pi \nSequent !A, \Delta$}}
  \AxiomC{\LR{$\Gamma \nSequent !B, \Pi \nSequent !B, \Delta$}}
  \AxiomC{\LR{$\Gamma \nSequent !A, !B, \Pi \nSequent \Delta$}}
  \RightLabel{\LR{$\iR\land\Undef$}}
  \TrinaryInfC{\LR{$\Gamma \nSequent !A \land !B, \Pi \nSequent \Delta$}}
  \DisplayProof
  \\[2ex]  
  \AxiomC{\LR{$\Gamma \nSequent \Pi \nSequent \Delta, !A$}}
  \AxiomC{\LR{$\Gamma \nSequent \Pi \nSequent \Delta, !B$}}
  \RightLabel{\LR{$\iR\land\True$}}
  \BinaryInfC{\LR{$\Gamma \nSequent \Pi \nSequent \Delta, !A \land !B$}}
  \DisplayProof
\end{center}
\end{defish}
\olpseganchor{OLP-0406-B015}{end}

\olpseganchor{OLP-0406-B016}{start}
\subsection{د~\texorpdfstring{\LR{$\lor$}}{∨} قاعدې}
\olpseganchor{OLP-0406-B016}{end}

\olpseganchor{OLP-0406-B017}{start}
دا د~\LR{$\lor$} قاعدې د لوکاشېوېچ، د کليني د قوي منطق او د ګوډل
د منطق دپاره دي۔
\olpseganchor{OLP-0406-B017}{end}

\olpseganchor{OLP-0406-B018}{start}
\begin{defish}
\begin{center}
 \AxiomC{\LR{$!A, \Gamma \nSequent \Pi \nSequent \Delta$}}
  \AxiomC{\LR{$!B, \Gamma \nSequent \Pi \nSequent \Delta$}}
  \RightLabel{\LR{$\iR\lor\False$}}
  \BinaryInfC{\LR{$!A \lor !B, \Gamma \nSequent \Pi \nSequent \Delta$}}
  \DisplayProof
\\[2ex]
  \AxiomC{\LR{$!A, \Gamma \nSequent !A, \Pi \nSequent \Delta$}}
  \AxiomC{\LR{$!B, \Gamma \nSequent !B, \Pi \nSequent \Delta$}}
  \AxiomC{\LR{$\Gamma \nSequent !A, !B, \Pi \nSequent \Delta$}}
  \RightLabel{\LR{$\iR\lor\Undef$}}
  \TrinaryInfC{\LR{$\Gamma \nSequent !A \lor !B, \Pi \nSequent \Delta$}}
  \DisplayProof
  \\[2ex]  
  \AxiomC{\LR{$\Gamma \nSequent \Pi \nSequent \Delta, !A, !B$}}
  \RightLabel{\LR{$\iR\lor\True$}}
  \UnaryInfC{\LR{$\Gamma \nSequent \Pi \nSequent \Delta, !A \lor !B$}}
  \DisplayProof
\end{center}
\end{defish}
\olpseganchor{OLP-0406-B018}{end}

\olpseganchor{OLP-0406-B019}{start}
\subsection{د~\texorpdfstring{\LR{$\lif$}}{→} قاعدې}
\olpseganchor{OLP-0406-B019}{end}

\olpseganchor{OLP-0406-B020}{start}
دا د~\LR{$\lif$} قاعدې د لوکاشېوېچ په منطق کښې دي۔
\olpseganchor{OLP-0406-B020}{end}

\olpseganchor{OLP-0406-B021}{start}
\begin{defish}
\begin{center}
  \AxiomC{\LR{$\Gamma \nSequent \Pi \nSequent \Delta, !A$}}
  \AxiomC{\LR{$!B, \Gamma \nSequent \Pi \nSequent \Delta$}}
  \RightLabel{\LR{$\iR\lif\False[\LogLuk[3]]$}}
  \BinaryInfC{\LR{$!A \lif !B, \Gamma \nSequent \Pi \nSequent \Delta$}}
  \DisplayProof
  \\[2ex]
  \AxiomC{\LR{$\Gamma \nSequent !A, !B, \Pi \nSequent \Delta$}}
  \AxiomC{\LR{$!B, \Gamma \nSequent \Pi \nSequent \Delta, !A$}}
  \RightLabel{\LR{$\iR\lif\Undef[\LogLuk[3]]$}}
  \BinaryInfC{\LR{$\Gamma \nSequent !A \lif !B, \Pi \nSequent \Delta$}}
  \DisplayProof
\\[2ex]
  \AxiomC{\LR{$!A, \Gamma \nSequent !B, \Pi \nSequent \Delta, !B$}}
  \AxiomC{\LR{$!A, \Gamma \nSequent !A, \Pi \nSequent \Delta, !B$}}
  \RightLabel{\LR{$\iR{\lif}{\True}[\LogLuk[3]]$}}
  \BinaryInfC{\LR{$\Gamma \nSequent \Pi \nSequent \Delta, !A \lif !B$}}
  \DisplayProof
\end{center}
\end{defish}
\olpseganchor{OLP-0406-B021}{end}

\olpseganchor{OLP-0406-B022}{start}
دا د~\LR{$\lif$} قاعدې د کليني په قوي منطق کښې دي۔
\olpseganchor{OLP-0406-B022}{end}

\olpseganchor{OLP-0406-B023}{start}
\begin{defish}
\begin{center}
  \AxiomC{\LR{$\Gamma \nSequent \Pi \nSequent \Delta, !A$}}
  \AxiomC{\LR{$!B, \Gamma \nSequent \Pi \nSequent \Delta$}}
  \RightLabel{\LR{$\iR\lif\False[\LogKs]$}}
  \BinaryInfC{\LR{$!A \lif !B, \Gamma \nSequent \Pi \nSequent \Delta$}}
  \DisplayProof
  \\[2ex]
  \AxiomC{\LR{$!B, \Gamma \nSequent !B, \Pi \nSequent \Delta$}}
  \AxiomC{\LR{$\Gamma \nSequent !A, !B, \Pi \nSequent \Delta$}}
  \AxiomC{\LR{$\Gamma \nSequent !A, \Pi \nSequent \Delta, !A$}}
  \RightLabel{\LR{$\iR\lif\Undef[\LogKs]$}}
  \TrinaryInfC{\LR{$\Gamma \nSequent !A \lif !B, \Pi \nSequent \Delta$}}
  \DisplayProof
\\[2ex]
  \AxiomC{\LR{$!A, \Gamma \nSequent \Pi \nSequent \Delta, !B$}}
  \RightLabel{\LR{$\iR{\lif}{\True}[\LogKs]$}}
  \UnaryInfC{\LR{$\Gamma \nSequent \Pi \nSequent \Delta, !A \lif !B$}}
  \DisplayProof
\end{center}
\end{defish}
\olpseganchor{OLP-0406-B023}{end}

\olpseganchor{OLP-0406-B024}{start}
دا د~\LR{$\lif$} قاعدې د ګوډل په منطق کښې دي۔
\olpseganchor{OLP-0406-B024}{end}

\olpseganchor{OLP-0406-B025}{start}
\begin{defish}
\begin{center}
  \AxiomC{\LR{$\Gamma \nSequent !A, \Pi \nSequent \Delta, !A$}}
  \AxiomC{\LR{$!B, \Gamma \nSequent \Pi \nSequent \Delta$}}
  \RightLabel{\LR{$\iR\lif\False[\LogGod[3]]$}}
  \BinaryInfC{\LR{$!A \lif !B, \Gamma \nSequent \Pi \nSequent \Delta$}}
  \DisplayProof
  \\[2ex]
  \AxiomC{\LR{$\Gamma \nSequent !B, \Pi \nSequent \Delta$}}
  \AxiomC{\LR{$\Gamma \nSequent \Pi \nSequent \Delta, !A$}}
  \RightLabel{\LR{$\iR\lif\Undef[\LogGod[3]]$}}
  \BinaryInfC{\LR{$\Gamma \nSequent !A \lif !B, \Pi \nSequent \Delta$}}
  \DisplayProof
\\[2ex]
  \AxiomC{\LR{$!A, \Gamma \nSequent !B, \Pi \nSequent \Delta, !B$}}
  \AxiomC{\LR{$!A, \Gamma \nSequent !A, \Pi \nSequent \Delta, !B$}}
  \RightLabel{\LR{$\iR{\lif}{\True}[\LogGod[3]]$}}
  \BinaryInfC{\LR{$\Gamma \nSequent \Pi \nSequent \Delta, !A \lif !B$}}
  \DisplayProof
\end{center}
\end{defish}
\olpseganchor{OLP-0406-B025}{end}


\olpseganchor{OLP-0406-B026}{start}
\begin{sidewaysfigure}
  \begin{center}  
  \AxiomC{\LR{$A \nSequent A \nSequent A$}}
  \RightLabel{\iR \Weakening \True}
  \UnaryInfC{\LR{$A \nSequent A \nSequent B, A$}}
  \RightLabel{\iR \Weakening \Undef}
  \UnaryInfC{\LR{$A \nSequent B, A \nSequent B, A$}}
  \RightLabel{\iR \Weakening \Undef}
  \UnaryInfC{\LR{$A \nSequent A, B, A \nSequent B, A$}}
  \AxiomC{\LR{$A \nSequent A \nSequent A$}}
  \RightLabel{\iR \Weakening \True}
  \UnaryInfC{\LR{$A \nSequent A \nSequent A, A$}}
  \RightLabel{\iR \Weakening \True}
  \UnaryInfC{\LR{$A \nSequent A \nSequent B, A, A$}}
  \RightLabel{\iR \Weakening \False}
  \UnaryInfC{\LR{$B, A \nSequent A \nSequent B, A, A$}}
  \RightLabel{\LR{$\iR\lif\Undef$}}
  \BinaryInfC{\LR{$A \nSequent A \lif B, A \nSequent B, A$}}
  \AxiomC{\LR{$B \nSequent B \nSequent B$}}
  \RightLabel{\iR \Weakening \Undef}
  \UnaryInfC{\LR{$B \nSequent A, B \nSequent B$}}
  \RightLabel{\iR \Exchange \Undef}
  \UnaryInfC{\LR{$B \nSequent B, A \nSequent B$}}
  \RightLabel{\iR \Weakening \Undef}
  \UnaryInfC{\LR{$B \nSequent A, B, A \nSequent B$}}
  \RightLabel{\iR \Weakening \False}
  \UnaryInfC{\LR{$A, B \nSequent A, B, A \nSequent B$}}
  \RightLabel{\iR \Exchange \False}
  \UnaryInfC{\LR{$B, A \nSequent A, B, A \nSequent B$}}
  \AxiomC{\LR{$A \nSequent A \nSequent A$}}
  \RightLabel{\iR \Weakening \True}
  \UnaryInfC{\LR{$A \nSequent A \nSequent B, A$}}
  \RightLabel{\iR \Weakening \False}
  \UnaryInfC{\LR{$B, A \nSequent A \nSequent B, A$}}
  \RightLabel{\iR \Weakening \False}
  \UnaryInfC{\LR{$B, B, A \nSequent A \nSequent B, A$}}
  \RightLabel{\LR{$\iR\lif\Undef$}}
  \BinaryInfC{\LR{$B, A \nSequent A \lif B, A \nSequent B$}}
  \RightLabel{\LR{$\iR\lif\False$}}
  \BinaryInfC{\LR{$A \lif B, A \nSequent A \lif B, A \nSequent B$}}
  \DisplayProof
  \end{center}
  \caption{په~\LR{$\LogLuk[3]$} کښې د اشتقاق بېلګه}
\end{sidewaysfigure}
\olpseganchor{OLP-0406-B026}{end}

\paragraph{د منجمدې سرچينې بشپړه اړونده برخه}
\olpseganchor{OLP-0652-B004}{start}
\olfileid{mvl}{seq}{prf-olp0652}
\olpseganchor{OLP-0652-B004}{end}

\olpseganchor{OLP-0652-B005}{start}
\olsection{د ثبوت-تيوريکي مفاهيم}
\olpseganchor{OLP-0652-B005}{end}

\olpseganchor{OLP-0652-B006}{start}
\begin{explain}
لکه څنګه چې مو يو شمېر مهم معنايي مفاهيم (اعتبار، معنايي استلزام
او د صدق وړتيا) تعريف کړي دي، اوس ورسره سم \emph{ثبوت-تيوريکي
مفاهيم} تعريفوو۔ دا مفاهيم په جوړښتونه کښې د تړلي فارمولونه
د صدق پر بنسټ نۀ، بلکې د ځينو سېکوېنټونو د اشتقاق وړتيا يا
د اشتقاق نۀ وړتيا پر بنسټ تعريفېږي۔ دا يوه مهمه موندنه وه چې دغه
مفاهيم سره سمون خوري۔ د دې سمون منځپانګه د \emph{سموالي} او
\emph{بشپړتيا قضيه} ده۔
\end{explain}
لاندې تعريفونه د پرتله کولو لپاره د کلاسيک سېکوېنټي حساب په بڼه دي۔ د څو ارزښته حساب قواعد د هغه د صدق ارزښتونو او څوځايي سېکوېنټونو له مخې ټاکل کېږي۔
\olpseganchor{OLP-0652-B006}{end}

\olpseganchor{OLP-0652-B007}{start}
\begin{defn}[قضيې]
تړلے فارمول~\LR{$!A$} هله \emph{قضيه} ده چې په~\LR{$\Log{LK}$} کښې د
سېکوېنټ \LR{$\quad \Sequent !A$} اشتقاق موجود وي۔ که \LR{$!A$}
قضيه وي، \LR{$\Proves !A$} ليکو؛ او که نۀ وي، \LR{$\Proves/ !A$} ليکو۔
\end{defn}
\olpseganchor{OLP-0652-B007}{end}

\olpseganchor{OLP-0652-B008}{start}
\begin{defn}[د اشتقاق وړتيا]
تړلے فارمول~\LR{$!A$} د تړلي فارمولونه له سټ~\LR{$\Gamma$} څخه
\emph{د اشتقاق وړ} ده، يعنې \LR{$\Gamma \Proves !A$}، هله او يوازې
هله چې يو متناهي فرعي سټ~\LR{$\Gamma_0 \subseteq \Gamma$} او په
\LR{$\Gamma_0$} کښې د تړلي فارمولونه يو لړ~\LR{$\Gamma_0'$} موجود وي، داسې
چې \LR{$\Log{LK}$} د \LR{$\Gamma_0' \Sequent !A$} سېکوېنټ اشتقاق کړي۔
که \LR{$!A$} له \LR{$\Gamma$} څخه د اشتقاق وړ نۀ وي، \LR{$\Gamma \Proves/ !A$}
ليکو۔
\end{defn}
\olpseganchor{OLP-0652-B008}{end}

\olpseganchor{OLP-0652-B009}{start}
د انقباض، کمزوري کولو او تبادلې د قاعدو له امله، په~\LR{$\Gamma_0'$}
کښې د تړلي فارمولونه ترتيب او شمېر اهميت نۀ لري: که سېکوېنټ
\LR{$\Gamma_0' \Sequent !A$} د اشتقاق وړ وي، نو د هر داسې
\LR{$\Gamma_0''$} دپاره \LR{$\Gamma_0'' \Sequent !A$} هم د اشتقاق وړ دے
چې د~\LR{$\Gamma_0'$} هماغه تړلي فارمولونه لري۔ مثلاً که
\LR{$\Gamma_0 = \{!B, !C\}$} وي، نو هم
\LR{$\Gamma_0' = \tuple{!B, !B, !C}$} او هم
\LR{$\Gamma_0'' = \tuple{!C, !C, !B}$} داسې لړونه دي چې يوازې د
\LR{$\Gamma_0$} تړلي فارمولونه لري۔ که د يوه لړ لرونکے سېکوېنټ
د اشتقاق وړ وي، د بل لړ لرونکے هم دے، مثلاً:
\begin{prooftree}
  \AxiomC{}
  \Deduce$!B, !B, !C \fCenter !A$
  \RightLabel{\LeftR{\Contraction}}
  \UnaryInf$!B, !C \fCenter !A$
  \RightLabel{\LeftR{\Exchange}}
  \UnaryInf$!C, !B \fCenter !A$
  \RightLabel{\LeftR{\Weakening}}
  \UnaryInf$!C, !C, !B \fCenter !A$
\end{prooftree}
له دې وروسته به وايو چې که \LR{$\Gamma_0$} د تړلي فارمولونه يو متناهي
سټ وي، نو \LR{$\Gamma_0 \Sequent !A$} هر هغه سېکوېنټ دے چې مقدم ئې
په~\LR{$\Gamma_0$} کښې د تړلي فارمولونه يو لړ وي؛ او د اړتيا په وخت کښې
انقباضونه، تبادلې او کمزوري کول په ناويلې توګه ورګډوو۔
\olpseganchor{OLP-0652-B009}{end}

\olpseganchor{OLP-0652-B010}{start}
\begin{defn}[سازګاري]
د جملو سټ~\LR{$\Gamma$} هله او يوازې هله \emph{ناسازګار} دے چې يو
متناهي فرعي سټ~\LR{$\Gamma_0 \subseteq \Gamma$} موجود وي، داسې چې
\LR{$\Log{LK}$} د \LR{$\Gamma_0 \Sequent \quad$} سېکوېنټ اشتقاق کړي۔
که \LR{$\Gamma$} ناسازګار نۀ وي، يعنې که د هر متناهي
\LR{$\Gamma_0 \subseteq \Gamma$} دپاره \LR{$\Log{LK}$} د
\LR{$\Gamma_0 \Sequent \quad$} سېکوېنټ اشتقاق نۀ کړي، نو هغه
\emph{سازګار} بولو۔
\end{defn}
\olpseganchor{OLP-0652-B010}{end}

\olpseganchor{OLP-0652-B011}{start}
\begin{prop}[انعکاسيت]
\ollabel{prop:reflexivity}
که \LR{$!A \in \Gamma$} وي، نو \LR{$\Gamma \Proves !A$}۔
\end{prop}
\olpseganchor{OLP-0652-B011}{end}

\olpseganchor{OLP-0652-B012}{start}
\begin{proof}
لومړنے سېکوېنټ \LR{$!A \Sequent !A$} د اشتقاق وړ دے، او
\LR{$\{!A\} \subseteq \Gamma$}۔
\end{proof}
\olpseganchor{OLP-0652-B012}{end}

\olpseganchor{OLP-0652-B013}{start}
\begin{prop}[يکنواختي]
\ollabel{prop:monotonicity}
که \LR{$\Gamma \subseteq \Delta$} او \LR{$\Gamma \Proves !A$} وي، نو
\LR{$\Delta \Proves !A$}۔
\end{prop}
\olpseganchor{OLP-0652-B013}{end}

\olpseganchor{OLP-0652-B014}{start}
\begin{proof}
فرض کړئ \LR{$\Gamma \Proves !A$}، يعنې يو متناهي
\LR{$\Gamma_0 \subseteq \Gamma$} شته چې
\LR{$\Gamma_0 \Sequent !A$} د اشتقاق وړ دے۔ ځکه چې
\LR{$\Gamma \subseteq \Delta$}، نو \LR{$\Gamma_0$} د~\LR{$\Delta$} هم يو متناهي
فرعي سټ دے۔ له دې امله د \LR{$\Gamma_0 \Sequent !A$} اشتقاق
هم \LR{$\Delta \Proves !A$} ښيي۔
\end{proof}
\olpseganchor{OLP-0652-B014}{end}

\olpseganchor{OLP-0652-B015}{start}
\begin{prop}[انتقاليت]
\ollabel{prop:transitivity}
که \LR{$\Gamma \Proves !A$} او \LR{$\{!A\} \cup \Delta \Proves !B$} وي، نو
\LR{$\Gamma \cup \Delta \Proves !B$}۔
\end{prop}
\olpseganchor{OLP-0652-B015}{end}

\olpseganchor{OLP-0652-B016}{start}
\begin{proof}
که \LR{$\Gamma \Proves !A$} وي، يو متناهي
\LR{$\Gamma_0 \subseteq \Gamma$} او د
\LR{$\Gamma_0 \Sequent !A$} اشتقاق~\LR{$\pi_0$} شته۔ که
\LR{$\{!A\} \cup \Delta \Proves !B$} وي، نو د کوم متناهي فرعي سټ
\LR{$\Delta_0 \subseteq \Delta$} دپاره د
\LR{$!A, \Delta_0 \Sequent !B$} اشتقاق~\LR{$\pi_1$} شته۔ لاندې
اشتقاق وګورئ:
\begin{prooftree}
  \AxiomC{}
  \RightLabel{\LR{$\pi_0$}}
  \Deduce$\Gamma_0 \fCenter !A$
  \AxiomC{}
  \RightLabel{\LR{$\pi_1$}}
  \Deduce$!A, \Delta_0 \fCenter !B$
  \RightLabel{\Cut}
  \BinaryInf$\Gamma_0, \Delta_0 \fCenter !B$
\end{prooftree}
ځکه چې \LR{$\Gamma_0 \cup \Delta_0 \subseteq \Gamma \cup \Delta$}، نو
دا \LR{$\Gamma \cup \Delta \Proves !B$} ښيي۔
\end{proof}
\olpseganchor{OLP-0652-B016}{end}

\olpseganchor{OLP-0652-B017}{start}
پام وکړئ، په ځانګړې توګه د دې معنٰی دا ده چې که
\LR{$\Gamma \Proves !A$} او \LR{$!A \Proves !B$} وي، نو
\LR{$\Gamma \Proves !B$}۔ همدارنګه ترې دا هم راځي چې که
\LR{$!A_1, \dots, !A_n \Proves !B$} وي او د هر~\LR{$i$} دپاره
\LR{$\Gamma \Proves !A_i$} وي، نو \LR{$\Gamma \Proves !B$}۔
\olpseganchor{OLP-0652-B017}{end}

\olpseganchor{OLP-0652-B018}{start}
\begin{prop}
\ollabel{prop:incons}
\LR{$\Gamma$} هله او يوازې هله ناسازګار دے چې د هرې جملې~\LR{$!A$} دپاره
\LR{$\Gamma \Proves {!A}$} وي۔
\end{prop}
\olpseganchor{OLP-0652-B018}{end}

\olpseganchor{OLP-0652-B019}{start}
\begin{proof}
د تمرين په توګه پرېښودل شوے دے۔
\end{proof}
\olpseganchor{OLP-0652-B019}{end}

\olpseganchor{OLP-0652-B020}{start}
\begin{prob}
\olref[fol][seq][ptn]{prop:incons} ثابت کړئ۔
\end{prob}
\olpseganchor{OLP-0652-B020}{end}

\olpseganchor{OLP-0652-B021}{start}
\begin{prop}[متناهي شاهد]
  \ollabel{prop:proves-compact}
  \begin{enumerate}
  \item که \LR{$\Gamma \Proves !A$} وي، نو يو متناهي فرعي سټ
    \LR{$\Gamma_0 \subseteq \Gamma$} شته چې \LR{$\Gamma_0 \Proves !A$}۔
  \item که د~\LR{$\Gamma$} هر متناهي فرعي سټ سازګار وي، نو \LR{$\Gamma$}
    سازګار دے۔
  \end{enumerate}
\end{prop}
\olpseganchor{OLP-0652-B021}{end}

\olpseganchor{OLP-0652-B022}{start}
\begin{proof}
  \begin{enumerate}
    \item که \LR{$\Gamma \Proves !A$} وي، نو يو متناهي فرعي سټ
      \LR{$\Gamma_0 \subseteq \Gamma$} شته، داسې چې د
      \LR{$\Gamma_0 \Sequent !A$} سېکوېنټ اشتقاق لري۔ له دې امله
      \LR{$\Gamma_0 \Proves !A$}۔
    \item که \LR{$\Gamma$} ناسازګار وي، نو يو متناهي فرعي سټ
      \LR{$\Gamma_0 \subseteq \Gamma$} شته، داسې چې \LR{$\Log{LK}$} د
      \LR{$\Gamma_0 \Sequent \quad$} سېکوېنټ اشتقاق کړي۔ خو بيا
      \LR{$\Gamma_0$} د~\LR{$\Gamma$} يو ناسازګار متناهي فرعي سټ دے۔
  \end{enumerate}
\end{proof}
\olpseganchor{OLP-0652-B022}{end}

\olpseganchor{OLP-0407-B004}{start}
\olpart{nml}{نورمال مودالي منطقونه}
\olpseganchor{OLP-0407-B004}{end}

\olpseganchor{OLP-0407-B005}{start}
\begin{editorial}
  دا برخه د نورمال مودالي منطقونو ميتاتيوري څېړي۔ د بنسټيز مودالي
  منطق د کلاسيکې مطابقت تيورۍ بحث د الډو انتونېلي پر يادښتونو ولاړ دے۔
\end{editorial}
\olpseganchor{OLP-0407-B005}{end}

\olpseganchor{OLP-0408-B004}{start}
\olchapter{nml}{syn}{نحو او معنٰی پوهنه}
\olpseganchor{OLP-0408-B004}{end}

\olpseganchor{OLP-0409-B004}{start}
\olfileid{nml}{syn}{int}
\olpseganchor{OLP-0409-B004}{end}

\olpseganchor{OLP-0409-B005}{start}
\olsection{پېژندنه}
\olpseganchor{OLP-0409-B005}{end}

\olpseganchor{OLP-0409-B006}{start}
مودالي منطق له \emph{مودالي قضيو} او د هغو تر منځ د استلزام له
اړيکو سره سروکار لري۔ د مودالي قضيو بېلګې دا دي:
\begin{enumerate}
\item دا لازمه ده چې \LR{$2+2=4$}.
\item دا لازمه ده چې سبا د باران کېدو امکان وي.
\item که دا لازمه وي چې~\LR{$!A$} ممکن وي، نو دا ممکنه ده چې~\LR{$!A$}.
\end{enumerate}
امکان او لزوم يوازيني موداليتونه نۀ دي: نور يوځايي نښلوونکي هم
موداليتونه ګڼل کېږي؛ لکه «پکار ده چې~\LR{$!A$} رښتيا وي»، «داسې به وي
چې~\LR{$!A$}»، «ډېنا پوهېږي چې~\LR{$!A$}»، يا «ډېنا باور لري چې~\LR{$!A$}».
\olpseganchor{OLP-0409-B006}{end}

\olpseganchor{OLP-0409-B007}{start}
د مودالي منطق له پخوانيو بحثونو څخه يو د ارسطو په
\emph{De Interpretatione} اثر کښې دے۔ په ديارلسم څپرکي کښې د امکان
يوه معنا داسې بيانوي چې لزوم ترې امکان راولي، او سرچپه لورے عمومي صدق نۀ کوي؛ امکان او لزوم يو د بل په وسيله تعريفېدلے شي؛ که
\LR{$!A \land !B$} په ممکنه توګه رښتيا وي، نو \LR{$!A$} او \LR{$!B$} هر يو په ممکنه
توګه رښتيا دي، خو برعکس نۀ؛ او که \LR{$!A \to !B$} لازم وي، نو که
\LR{$!A$} لازم وي، \LR{$!B$} هم لازم وي.
\olpseganchor{OLP-0409-B007}{end}

\olpseganchor{OLP-0409-B008}{start}
مودالي منطق ته لومړنۍ عصري لاره د سي.~آی. لېويس کار و، چې د لېويس
او لنګفورډ په \emph{سمبوليک منطق} (1932) کښې اوج ته ورسېد۔ لېويس او
لنګفورډ د مادي شرطيې په وسيله د استلزام له ښودنې راضي نۀ وو:
\LR{$!A \lif !B$} د «له \LR{$!A$} څخه \LR{$!B$} لازمېږي» کمزورے بديل دے۔ د هغې
پر ځاے ئې استلزام داسې مشخص کړ: «ضرورتاً، که \LR{$!A$} وي نو \LR{$!B$} وي»؛
نښه ئې \LR{$!A \strictif !B$} ده۔ د بېلابېلو خاصيتونو د روښانولو په هڅه
کښې لېويس پنځه بېل مودالي نظامونه، \Log{S1}, \ldots, \Log{S4},
\Log{S5}، وپېژندل، چې وروستي دوه لا هم کارېږي.
\olpseganchor{OLP-0409-B008}{end}

\olpseganchor{OLP-0409-B009}{start}
د لېويس او لنګفورډ لاره خالصه \emph{نحوي} وه: معقول بديهي اصول او
قاعدې ئې وټاکلې او وئې څېړل چې د هغو په مټ څه ثابتېدلے شي۔ معنايي
لار اوږده موده لاسرسۍ ته نۀ وه، تر هغه چې رودولف کارناپ په
\emph{معنٰی او لزوم} (1947) کښې \emph{د حالت توصيف} مفهوم وکاراوه،
يعنې د اټومي جملو يوه ټولګه (هغه جملې چې په هماغه توصيف کښې
«رښتيا» دي)۔ کارناپ چې د صدق تعريف هرې جملې \LR{$!A$} ته وغځاوه، نو
\LR{$!A$} هله او يوازې هله \emph{لازماً رښتيا} بولي چې د حالت په ټولو
توصيفونو کښې رښتيا وي۔ ځکه چې هر عامل د حالتونو هماغه ټوله مجموعه
ګوري، \emph{پرله‌پسې} لزوم او امکان د «ممکن، ضرورتاً، \ldots ممکن \LR{$!A$}»
په بڼه کښې د دننني مودالي عامل صدقي ارزښت نۀ بدلوي۔
\olpseganchor{OLP-0409-B009}{end}

\olpseganchor{OLP-0409-B010}{start}
ساول کرېپکي په ۱۹۵۹ کال کښې د «په مودالي منطق کښې د بشپړتيا
يوه قضيه» په مقاله کښې (د سمبوليک منطق ژورنال، ۲ مخ) د
کميه‌لرونکي اېس پنځه لپاره د ممکنه نړۍو يو مدل ورکړ۔ د لايبنېڅ له نظر سره سم، په دې مدل کښې يو بيان هله
او يوازې هله لازم رښتيا دے چې په ټولو ممکنه نړۍو کښې رښتيا وي۔
که د لزوم د معلومولو لپاره په يوه نړۍ \LR{$w$} (يا د کارناپ د حالت په
توصيف \LR{$s$}) کښې هماغه ټوله مجموعه ګورو، د لزوم د حکم صدق له \LR{$w$}
سره نۀ بدلېږي۔ په ۱۹۶۳ کال کښې د «د مودالي منطق معنايي شننه ۱»
په مقاله کښې (۶۹–۷۰ مخونه) کرېپکي د \emph{لاسرسي اړيکه} \LR{$R$}
وکاروله چې کومې نړۍ له کومې نړۍ څخه ممکنې دي؛ نو
«ضرورتاً \LR{$!A$}» په نړۍ \LR{$w$} کښې هله او يوازې هله رښتيا وي چې \LR{$!A$}
په ټولو هغو نړۍو \LR{$w'$} کښې رښتيا وي چې له \LR{$w$} څخه \emph{ورته لاسرسے
شته}۔ هغه معناپوهنې چې د دې لارې يوه بڼه کاروي، د کرېپکي معناپوهنې
بلل کېږي او د مودالي منطقونو د پراخې ودې لار ئې پرانيسته۔
\olpseganchor{OLP-0409-B010}{end}

\olpseganchor{OLP-0409-B011}{start}
کله چې مودالي منطق د کرېپکي د معناپوهنې له مخې تفسير شي، راته ښيي
چې \emph{اړيکيز جوړښتونه} له «دننه» څنګه ښکاري۔ اړيکيز جوړښت يوازې
يو سټ دے چې يوه دوه ځايه اړيکه ورباندې ټاکل شوې وي (د بېلګې په
توګه، په ټولګي کښې د زده کوونکو هغه سټ چې د ټولنيز امنيت د شمېرو له
مخې مرتب شوے وي، يو اړيکيز جوړښت دے)۔ خو اړيکيز جوړښتونه په
بېلابېلو ډګرونو کښې راځي: د نړۍ د حالتونو د نسبي امکان تر څنګ، د
يوه عامل معرفتي حالتونه د معرفتي امکان له مخې تړلے کېدلے شي، يا د
يوه خوځنده نظام حالتونه د هغو تر منځ د لېږدونو له مخې، او داسې نور۔
مودالي منطق د دې ټولو د مدل جوړولو دپاره کارېږي: لومړے عادي، د
حقيقت اړوند مودالي منطق راکوي؛ نور معرفتي منطق، خوځنده منطق، او
داسې نور.
\olpseganchor{OLP-0409-B011}{end}

\olpseganchor{OLP-0409-B012}{start}
موږ پر يوه ځانګړي اړخ تمرکز کوو چې مودالي منطقپوهان ئې «د مطابقت
تيوري» بولي۔ د کرېپکي له لومړنيو مهمو موندنو څخه يوه دا وه چې د
لاسرسي اړيکې~\LR{$R$} ډېر خاصيتونه (لکه تعدي، تقارن او داسې نور) پخپله
\emph{په مودالي ژبه کښې} د مناسبو «مودالي طرحو» په وسيله مشخصېدلے
شي۔ د بېلګې په توګه، مودالي منطقپوهان وايي چې د \LR{$R$} انعکاسي والی له
دې طرحې سره «مطابقت لري»: «که \LR{$!A$} لازماً رښتيا وي، نو \LR{$!A$} رښتيا
وي»۔ موږ په عمده توګه د مودالي منطق د هغو کلاسيکي نظامونو (لکه
\Log{S4} او \Log{S5}) د مطابقت تيوري څېړو چې د \Ax{D}, \Ax{T},
\Ax{B}, \Ax{4} او~\Ax{5} طرحو په ګډولو سره ترلاسه کېږي.
\olpseganchor{OLP-0409-B012}{end}

\olpseganchor{OLP-0410-B004}{start}
\olfileid{nml}{syn}{lan}
\olpseganchor{OLP-0410-B004}{end}

\olpseganchor{OLP-0410-B005}{start}
\olsection{د بنسټيز مودالي منطق ژبه}
\olpseganchor{OLP-0410-B005}{end}

\olpseganchor{OLP-0410-B006}{start}
\begin{defn}
د مودالي منطق بنسټيزه ژبه دا توکي لري:
\begin{enumerate}
  \item د دروغوالی قضيوي ثابت~\LR{$\lfalse$}۔
  \item د رښتياوالی قضيوي ثابت~\LR{$\ltrue$}۔
  \item د بياني متغيرونه يو شمېرېدونکے لامتناهي سټ: \LR{$\Obj
    p_0$}، \LR{$\Obj p_1$}، \LR{$\Obj p_2$}، \dots
  \item قضيوي نښلوونکي: \startycommalist
  \ycomma \LR{$\lnot$} (نفي)%
  \ycomma \LR{$\land$} (عطف)%
  \ycomma \LR{$\lor$} (فصل)%
  \ycomma \LR{$\lif$} (شرطي)%
  \ycomma \LR{$\liff$} (دوه اړخيز شرطي).
  \item مودالي عامل~\LR{$\Box$}۔
  \item مودالي عامل~\LR{$\Diamond$}۔
\end{enumerate}
\end{defn}
\olpseganchor{OLP-0410-B006}{end}

\olpseganchor{OLP-0410-B007}{start}
\begin{defn}
د بنسټيزې مودالي ژبې \emph{فارمولونه} په استقرايي ډول داسې
  تعريفېږي:
\begin{enumerate}
\item \LR{$\lfalse$} يو اټومي فارمول دے۔
\olpseganchor{OLP-0410-B007}{end}

\olpseganchor{OLP-0410-B008}{start}
\item \LR{$\ltrue$} يو اټومي فارمول دے۔
\olpseganchor{OLP-0410-B008}{end}

\olpseganchor{OLP-0410-B009}{start}
\item هر بياني متغير~\LR{$\Obj p_i$} يو اټومي فارمول دے۔
\olpseganchor{OLP-0410-B009}{end}

\olpseganchor{OLP-0410-B010}{start}
\item که \LR{$!A$} يو فارمول وي، نو \LR{$\lnot !A$} هم
  يو فارمول دے۔
\olpseganchor{OLP-0410-B010}{end}

\olpseganchor{OLP-0410-B011}{start}
\item که \LR{$!A$} او \LR{$!B$} فارمولونه وي، نو \LR{$(!A \land
  !B)$} يو فارمول دے۔
\olpseganchor{OLP-0410-B011}{end}

\olpseganchor{OLP-0410-B012}{start}
\item که \LR{$!A$} او \LR{$!B$} فارمولونه وي، نو \LR{$(!A \lor !B)$}
  يو فارمول دے۔
\olpseganchor{OLP-0410-B012}{end}

\olpseganchor{OLP-0410-B013}{start}
\item که \LR{$!A$} او \LR{$!B$} فارمولونه وي، نو \LR{$(!A \lif !B)$}
  يو فارمول دے۔
\olpseganchor{OLP-0410-B013}{end}

\olpseganchor{OLP-0410-B014}{start}
\item که \LR{$!A$} او \LR{$!B$} فارمولونه وي، نو \LR{$(!A \liff !B)$}
  يو فارمول دے۔
\olpseganchor{OLP-0410-B014}{end}

\olpseganchor{OLP-0410-B015}{start}
\item که \LR{$!A$} يو فارمول وي، نو \LR{$\Box !A$} هم
  يو فارمول دے۔
\olpseganchor{OLP-0410-B015}{end}

\olpseganchor{OLP-0410-B016}{start}
\item که \LR{$!A$} يو فارمول وي، نو \LR{$\Diamond !A$} هم
  يو فارمول دے۔
\olpseganchor{OLP-0410-B016}{end}

\olpseganchor{OLP-0410-B017}{start}
\item بل هېڅ شے فارمول نۀ دے۔
\end{enumerate}
\end{defn}
\olpseganchor{OLP-0410-B017}{end}

\olpseganchor{OLP-0410-B029}{start}
که فارمول~\LR{$!A$} کښې \LR{$\Box$} يا \LR{$\Diamond$} نۀ وي، هغه
\emph{له موداليت پرته} بولو۔
\olpseganchor{OLP-0410-B029}{end}

\olpseganchor{OLP-0411-B004}{start}
\olfileid{nml}{syn}{sub}
\olpseganchor{OLP-0411-B004}{end}

\olpseganchor{OLP-0411-B005}{start}
\olsection{هممهاله ځايناستي}
\olpseganchor{OLP-0411-B005}{end}

\olpseganchor{OLP-0411-B006}{start}
د \LR{$!A$} فارمول يوه \emph{بېلګه} هغه پايله ده چې په~\LR{$!A$} کښې د
بياني متغير هر راتګ د يو بل فارمول په راتګ بدل
شي۔ د فارمولونه د بېلګو يادونه به د اعتبار او د اشتقاق وړتيا د
بحث پر مهال بيا بيا کوو۔ نو د دې مفهوم دقيق تعريف ګټور دے۔
\olpseganchor{OLP-0411-B006}{end}

\olpseganchor{OLP-0411-B007}{start}
\begin{defn}\ollabel{def:subst-inst}
  که \LR{$!A$} يو مودالي فارمول وي او ټول بياني متغيرونه ئې له \LR{$p_1$}، \dots، \LR{$p_n$} څخه وي، او \LR{$!D_1$}، \dots،
  \LR{$!D_n$} هم مودالي فارمولونه وي، نو
  \LR{$\SSubst{!A}{\subst{!D_1}{p_1}, \dots, \subst{!D_n}{p_n}}$} داسې
  تعريفوو: په~\LR{$!A$} کښې د هر \LR{$p_i$} پر ځاے \LR{$!D_i$} په يو وخت ږدو۔ په
  رسمي ډول دا پر~\LR{$!A$} استقرايي تعريف دے:
  \begin{enumerate}
    \item \indcase{!A}{\lfalse}{%
        \LR{$\SSubst{\indfrm}{\subst{!D_1}{p_1}, \dots,
          \subst{!D_n}{p_n}}$} هماغه \LR{$\lfalse$} دے۔}
    \item \indcase{!A}{\ltrue}{%
        \LR{$\SSubst{\indfrm}{\subst{!D_1}{p_1}, \dots,
          \subst{!D_n}{p_n}}$} هماغه \LR{$\ltrue$} دے۔}
    \item \indcase{!A}{q}{\LR{$\SSubst{\indfrm}{\subst{!D_1}{p_1}, \dots,
        \subst{!D_n}{p_n}}$} هماغه \LR{$q$} دے، په دې شرط چې \LR{$q \not\ident p_i$}
      د \LR{$i = 1$}، \dots،~\LR{$n$} د هر ارزښت دپاره وي۔}
    \item \indcase{!A}{p_i}{\LR{$\SSubst{\indfrm}{\subst{!D_1}{p_1},
        \dots, \subst{!D_n}{p_n}}$} هماغه \LR{$!D_i$} دے۔}
    \item \indcase{!A}{\lnot !B}{%
        \LR{$\SSubst{\indfrm}{\subst{!D_1}{p_1}, \dots, \subst{!D_n}{p_n}}$} هماغه
        \LR{$\lnot \SSubst{!B}{\subst{!D_1}{p_1}, \dots, \subst{!D_n}{p_n}}$} دے۔}
    \item \indcase{!A}{(!B \land
        !C)}{\LR{$\SSubst{\indfrm}{\subst{!D_1}{p_1}, \dots,
          \subst{!D_n}{p_n}}$} دا دے:
        \[(\SSubst{!B}{\subst{!D_1}{p_1}, \dots, \subst{!D_n}{p_n}} \land
          \SSubst{!C}{\subst{!D_1}{p_1}, \dots, \subst{!D_n}{p_n}}).\]}
    \item \indcase{!A}{(!B \lor
          !C)}{\LR{$\SSubst{\indfrm}{\subst{!D_1}{p_1}, \dots,
          \subst{!D_n}{p_n}}$} دا دے:
        \[(\SSubst{!B}{\subst{!D_1}{p_1}, \dots, \subst{!D_n}{p_n}} \lor
          \SSubst{!C}{\subst{!D_1}{p_1}, \dots, \subst{!D_n}{p_n}}).\]}
    \item \indcase{!A}{(!B \lif
          !C)}{\LR{$\SSubst{\indfrm}{\subst{!D_1}{p_1}, \dots,
          \subst{!D_n}{p_n}}$} دا دے:
        \[(\SSubst{!B}{\subst{!D_1}{p_1}, \dots, \subst{!D_n}{p_n}} \lif
          \SSubst{!C}{\subst{!D_1}{p_1}, \dots, \subst{!D_n}{p_n}}).\]}
    \item \indcase{!A}{(!B \liff
          !C)}{\LR{$\SSubst{\indfrm}{\subst{!D_1}{p_1}, \dots,
          \subst{!D_n}{p_n}}$} دا دے:
        \[(\SSubst{!B}{\subst{!D_1}{p_1}, \dots, \subst{!D_n}{p_n}} \liff
          \SSubst{!C}{\subst{!D_1}{p_1}, \dots, \subst{!D_n}{p_n}}).\]}
    \item \indcase{!A}{\Box !B}{%
        \LR{$\SSubst{\indfrm}{\subst{!D_1}{p_1}, \dots, \subst{!D_n}{p_n}}$} داسې محاسبه کېږي:
        \LR{$\Box
        \SSubst{!B}{\subst{!D_1}{p_1}, \dots, \subst{!D_n}{p_n}}.$}}
    \item \indcase{!A}{\Diamond !B}{%
        \LR{$\SSubst{\indfrm}{\subst{!D_1}{p_1}, \dots, \subst{!D_n}{p_n}}$} داسې محاسبه کېږي:
        \LR{$\Diamond
        \SSubst{!B}{\subst{!D_1}{p_1}, \dots, \subst{!D_n}{p_n}}.$}}
  \end{enumerate}
  \LR{$\SSubst{!A}{\subst{!D_1}{p_1}, \dots,
    \subst{!D_n}{p_n}}$} فارمول د~\LR{$!A$} د \emph{ځايناستۍ بېلګه}
  بلل کېږي۔
\end{defn}
\olpseganchor{OLP-0411-B007}{end}

\olpseganchor{OLP-0411-B008}{start}
\begin{ex}
  فرض کړئ چې \LR{$!A$} د \LR{$p_1 \lif \Box(p_1 \land p_2)$}، \LR{$!D_1$} د
  \LR{$\Diamond(p_2 \lif p_3)$} او \LR{$!D_2$} د \LR{$\lnot\Box p_1$} سره برابر
  وي۔ نو \LR{$\SSubst{!A}{\subst{!D_1}{p_1}, \subst{!D_2}{p_2}}$} دا دے:
  \begin{align*}
    \Diamond(p_2 \lif p_3) &
    \lif \Box(\Diamond(p_2 \lif p_3) \land \lnot\Box p_1)
    \intertext{\RL{خو \LR{$\SSubst{!A}{\subst{!D_2}{p_1}, \subst{!D_1}{p_2}}$} دا دے:}}
    \lnot\Box p_1 &
    \lif \Box(\lnot\Box p_1 \land \Diamond(p_2 \lif p_3))
    \intertext{\RL{پام وکړئ چې هممهاله ځايناستي عموماً له پرله‌پسې
      ځايناستۍ سره يوه نۀ ده۔ د بېلګې په توګه، پورته
      \LR{$\SSubst{!A}{\subst{!D_1}{p_1}, \subst{!D_2}{p_2}}$} له
      \LR{$\Subst{(\Subst{!A}{!D_1}{p_1})}{!D_2}{p_2}$} سره پرتله کړئ، چې دا دے:}}
    \Subst{\Diamond(p_2 \lif p_3)}{\lnot\Box p_1}{p_2} &
    \lif \Subst{\Box(\Diamond(p_2 \lif p_3) \land p_2)}{\lnot\Box p_1}{p_2}, \text{\RL{ يعنې،}}\\
    \Diamond(\lnot\Box p_1 \lif p_3) &
    \lif \Box(\Diamond(\lnot\Box p_1
    \lif p_3) \land \lnot\Box p_1)\\
    \intertext{\RL{او له \LR{$\Subst{(\Subst{!A}{!D_2}{p_2})}{!D_1}{p_1}$} سره ئې هم پرتله کړئ:}}
    \Subst{p_1}{\Diamond(p_2 \lif p_3)}{p_1} &
    \lif \Subst{\Box(p_1 \land \lnot\Box p_1)}{\Diamond(p_2 \lif p_3)}{p_1}, \text{\RL{ يعنې،}}\\
    \Diamond(p_2 \lif p_3) &
    \lif \Box(\Diamond(p_2
    \lif p_3) \land \lnot\Box\Diamond(p_2 \lif p_3)).
  \end{align*}
\end{ex}
\olpseganchor{OLP-0411-B008}{end}

\olpseganchor{OLP-0412-B004}{start}
\olfileid{nml}{syn}{rel}
\olpseganchor{OLP-0412-B004}{end}

\olpseganchor{OLP-0412-B005}{start}
\olsection{اړيکيز مدلونه}
\olpseganchor{OLP-0412-B005}{end}

\olpseganchor{OLP-0412-B006}{start}
د نورمال مودالي منطقونو د معنٰی پوهنې بنسټيز مفهوم \emph{اړيکيز
مدل} دے۔ دا د نړۍو له يوه سټ څخه جوړ وي چې د دوه‌ځاييزې «لاسرسي
اړيکې» له لارې سره تړلې وي، او ورسره داسې يوه ټاکنه وي چې څرګندوي
کوم بياني متغيرونه په کومو نړۍو کښې «رښتيا» ګڼل کېږي۔
\olpseganchor{OLP-0412-B006}{end}

\olpseganchor{OLP-0412-B007}{start}
\begin{defn}
  د بنسټيزې مودالي ژبې يو \emph{مدل} ترتيب شوې درې‌ګونه \LR{$\mModel{M}
  = \tuple{W, R, V}$} ده، چې په دې کښې
  \begin{enumerate}
  \item \LR{$W$} د «نړۍو» يو نا تش سټ دے،
  \item \LR{$R$} پر~\LR{$W$} د لاسرسي يوه دوه‌ځاييزه اړيکه ده، او
  \item \LR{$V$} داسې تابع ده چې هر بياني متغير~\LR{$p$} ته د ممکنه نړۍو يو سټ \LR{$V(p)$} ټاکي۔
  \end{enumerate}
  که \LR{$Rww'$} وي، نو وايو چې \LR{$w'$} هغه نړۍ ده چې له~\LR{$w$} څخه \emph{ورته لاسرسے
    شته}۔ که \LR{$w \in V(p)$} وي، نو وايو چې \LR{$p$} په~\LR{$w$} کښې
  \emph{رښتيا دے}۔
\end{defn}
\olpseganchor{OLP-0412-B007}{end}

\olpseganchor{OLP-0412-B008}{start}
د اړيکيزې معنٰی پوهنې ستره ګټه دا ده چې مدلونه په ساده شکلونو
ښودل کېدے شي، لکه په \olref{fig:simple} کښې۔ نړۍوې د غوټو په بڼه
ښودل کېږي، او نړۍ~\LR{$w'$} ته هله او يوازې هله له~\LR{$w$} څخه لاسرسے شته
چې له~\LR{$w$} څخه~\LR{$w'$} ته غشے وي۔ سربېره پر دې، که \LR{$w \in V(p)$} وي،
غوټه (نړۍ) په~\LR{$\mTrue{p}$} نښه کوو؛ که نۀ وي، په~\mFalse{p} ئې
نښه کوو۔ \olref{fig:simple} هغه مدل ښيي چې
\LR{$W = \{w_1, w_2, w_3\}$}، \LR{$R = \{\tuple{w_1, w_2},\tuple{w_1,
  w_3}\}$}، \LR{$V(p) = \{w_1, w_2\}$} او \LR{$V(q) = \{w_2\}$} لري۔
\olpseganchor{OLP-0412-B008}{end}

\olpseganchor{OLP-0412-B009}{start}
\begin{figure}
  \begin{center}
    \begin{tikzpicture}[modal]
      \node[world] (w1) [label={[align=right]right:\mTrue{p}\\ \mFalse{q}}]{\LR{$w_1$}};
      \node[world] (w2) [label={[align=right]right:\mTrue{p}\\ \mTrue{q}},
        above right=of w1]{\LR{$w_2$}};
      \node[world] (w3) [label={[align=right]right:\mFalse{p}\\ \mFalse{q}},
        below right=of w1] {\LR{$w_3$}};
      \draw[->] (w1) to (w2);
      \draw[->] (w1) to (w3);
    \end{tikzpicture}
  \end{center}
  \caption{يو ساده مدل۔}
  \ollabel{fig:simple}
\end{figure}
\olpseganchor{OLP-0412-B009}{end}

\olpseganchor{OLP-0413-B004}{start}
\olfileid{nml}{syn}{trw}
\olpseganchor{OLP-0413-B004}{end}

\olpseganchor{OLP-0413-B005}{start}
\olsection{په يوه نړۍ کښې رښتياوالی}
\olpseganchor{OLP-0413-B005}{end}

\olpseganchor{OLP-0413-B006}{start}
هر مودالي مدل ټاکي چې کوم مودالي فارمولونه په کومو نړۍو کښې
رښتيا ګڼل کېږي۔ دا اړيکه چې «مدل~\LR{$\mModel{M}$}، فارمول~\LR{$!A$}
په نړۍ~\LR{$w$} کښې رښتيا کوي» د اړيکيزې معنٰی پوهنې بنسټيز مفهوم دے۔
دا اړيکه په استقرايي ډول تعريفېږي، او د نامودالي عاملونو دپاره د
رښتياوالي د جدولونو له دوديزو قاعدو سره سمون خوري۔
\olpseganchor{OLP-0413-B006}{end}

\olpseganchor{OLP-0413-B007}{start}
\begin{defn}\ollabel{defn:mmodels}
  په~\LR{$\mModel M$} کښې په~\LR{$w$} کښې د فارمول~\LR{$!A$} \emph{رښتياوالی}،
  په نښو کښې \LR{$\mSat{M}{!A}[w]$}، په استقرايي ډول داسې تعريفېږي:
  \begin{enumerate}
  \item \indcase{!A}{\lfalse}{هېڅکله \LR{$\mSat{M}{\lfalse}[w]$} نۀ وي}.
  \item \indcase{!A}{\ltrue}{تل \LR{$\mSat{M}{\ltrue}[w]$} وي}.
  \item \LR{$\mSat{M}{p}[w]$} هله او يوازې هله چې \LR{$w \in V(p)$} وي۔
  \item \indcase{!A}{\lnot !B}{\LR{$\mSat{M}{\indfrm}[w]$} هله او يوازې هله چې
    \LR{$\mSat/{M}{!B}[w]$} وي}.
  \item \indcase{!A}{(!B \land !C)}{\LR{$\mSat{M}{\indfrm}[w]$} هله او يوازې هله چې
    \LR{$\mSat{M}{!B}[w]$} او \LR{$\mSat{M}{!C}[w]$} دواړه وي}.
  \item \indcase{!A}{(!B \lor !C)}{\LR{$\mSat{M}{\indfrm}[w]$} هله او يوازې هله چې
    \LR{$\mSat{M}{!B}[w]$} يا \LR{$\mSat{M}{!C}[w]$} وي} (يا دواړه).
  \item \indcase{!A}{(!B \lif !C)}{\LR{$\mSat{M}{\indfrm}[w]$} هله او يوازې هله چې
    \LR{$\mSat/{M}{!B}[w]$} يا \LR{$\mSat{M}{!C}[w]$} وي}.
  \item \indcase{!A}{(!B \liff !C)}{\LR{$\mSat{M}{\indfrm}[w]$} هله او يوازې هله چې
    يا دواړه \LR{$\mSat{M}{!B}[w]$} او \LR{$\mSat{M}{!C}[w]$} وي، يا
    نۀ \LR{$\mSat{M}{!B}[w]$} وي او نۀ \LR{$\mSat{M}{!C}[w]$} وي}.
  \item\ollabel{defn:sub:mmodels-box}
    \indcase{!A}{\Box !B}{\LR{$\mSat{M}{\indfrm}[w]$} هله او يوازې هله چې
    \LR{$\mSat{M}{!B}[w']$} د هر هغه \LR{$w' \in W$} دپاره وي چې \LR{$Rww'$} وي۔}
  \item\ollabel{defn:sub:mmodels-diamond}
    \indcase{!A}{\Diamond !B}{\LR{$\mSat{M}{\indfrm}[w]$} هله او يوازې هله چې
    \LR{$\mSat{M}{!B}[w']$} لږ تر لږه د يوه هغه \LR{$w' \in W$} دپاره وي چې \LR{$Rww'$} وي۔}
  \end{enumerate}
\end{defn}
\olpseganchor{OLP-0413-B007}{end}

\olpseganchor{OLP-0413-B008}{start}
د \olref{defn:sub:mmodels-box} فقرې ته پام وکړئ: که داسې هېڅ~\LR{$w'$}
نۀ وي چې \LR{$Rww'$} وي، نو فارمول~\LR{$\Box !B$} په~\LR{$w$} کښې رښتيا
وي۔ په دې حالت کښې \LR{$\Box !B$} په~\LR{$w$} کښې \emph{په تش ډول} رښتيا وي۔ همدارنګه،
\LR{$\Box !B$} په~\LR{$w$} کښې رښتيا کېدے شي، حتٰی که \LR{$!B$} هلته رښتيا نۀ
وي۔ په~\LR{$w$} کښې د~\LR{$!B$} رښتياوالی دا نۀ تضمينوي چې \LR{$\Diamond !B$} به
هم په~\LR{$w$} کښې رښتيا وي۔ خو که \LR{$Rww$} وي، د بېلګې په توګه که \LR{$R$}
انعکاسي وي، نو دا نتيجه سمېږي۔ که داسې هېڅ \LR{$w'$} نۀ وي چې \LR{$Rww'$} وي،
نو د هر~\LR{$!A$} دپاره \LR{$\mSat/{M}{\Diamond !A}[w]$} وي۔
\olpseganchor{OLP-0413-B008}{end}

\olpseganchor{OLP-0413-B009}{start}
\begin{prob}
  د \olref[nml][syn][rel]{fig:simple} مدل وڅېړئ۔ له لاندې کومې
  جملې رښتيا دي؟
    \begin{enumerate}
    \item \LR{$\mSat{M}{q}[w_1]$};
    \item \LR{$\mSat{M}{\lnot q}[w_3]$};
    \item \LR{$\mSat{M}{p \lor q}[w_1]$};
    \item \LR{$\mSat{M}{\Box (p \lor q)}[w_1]$};
    \item \LR{$\mSat{M}{\Box q}[w_3]$};
    \item \LR{$\mSat{M}{\Box \bot}[w_3]$};
    \item \LR{$\mSat{M}{\Diamond q}[w_1]$};
    \item \LR{$\mSat{M}{\Box q}[w_1]$};
    \item \LR{$\mSat{M}{\lnot \Box\Box \lnot q}[w_1]$}.
    \end{enumerate}
\end{prob}
\olpseganchor{OLP-0413-B009}{end}

\olpseganchor{OLP-0413-B010}{start}
% Prove duality of \Box and \Diamond

\begin{prop}\ollabel{prop:dual}
  \begin{enumerate}
  \item \LR{$\mSat{M}{\Box !A}[w]$} هله او يوازې هله چې \LR{$\mSat{M}{\lnot\Diamond\lnot !A}[w]$} وي۔
  \item \LR{$\mSat{M}{\Diamond !A}[w]$} هله او يوازې هله چې \LR{$\mSat{M}{\lnot\Box\lnot !A}[w]$} وي۔
  \end{enumerate}
\end{prop}
\olpseganchor{OLP-0413-B010}{end}

\olpseganchor{OLP-0413-B011}{start}
\begin{proof}
  \begin{enumerate}
    \item \LR{$\mSat{M}{\lnot\Diamond\lnot !A}[w]$} هله او يوازې هله چې \LR{$\mSat/{M}{\Diamond\lnot
      !A}[w]$} وي، د \LR{$\mSat{M}{}[w]$} د تعريف له مخې۔ \LR{$\mSat{M}{\Diamond\lnot
      !A}[w]$} هله او يوازې هله چې د يوه داسې \LR{$w'$} دپاره چې \LR{$Rww'$} وي، \LR{$\mSat{M}{\lnot
      !A}[w']$} وي۔ نو \LR{$\mSat/{M}{\Diamond\lnot !A}[w]$} هله او يوازې هله چې د هر
    داسې \LR{$w'$} دپاره چې \LR{$Rww'$} وي، \LR{$\mSat/{M}{\lnot !A}[w']$} وي۔ همدارنګه
    \LR{$\mSat/{M}{\lnot !A}[w']$} هله او يوازې هله چې \LR{$\mSat{M}{!A}[w']$} وي۔
    له دې ټولو څخه \LR{$\mSat{M}{\lnot\Diamond\lnot !A}[w]$} هله او يوازې هله
    چې د هر داسې \LR{$w'$} دپاره چې \LR{$Rww'$} وي، \LR{$\mSat{M}{!A}[w']$} وي۔ بيا د
    \LR{$\mSat{M}{}[w]$} د تعريف له مخې دا هله او يوازې هله چې
      \LR{$\mSat{M}{\Box !A}[w]$} وي۔
    \item \LR{$\mSat{M}{\lnot\Box\lnot !A}[w]$} هله او يوازې هله چې
      \LR{$\mSat/{M}{\Box\lnot !A}[w]$} وي۔ \LR{$\mSat{M}{\Box\lnot !A}[w]$} هله او يوازې هله چې د هر
      هغه \LR{$w'$} دپاره چې \LR{$Rww'$} وي، \LR{$\mSat{M}{\lnot !A}[w']$} وي۔ نو
      \LR{$\mSat/{M}{\Box\lnot !A}[w]$} هله او يوازې هله چې د يوه داسې \LR{$w'$} دپاره چې \LR{$Rww'$} وي،
      \LR{$\mSat/{M}{\lnot !A}[w']$} وي۔ همدارنګه \LR{$\mSat/{M}{\lnot
        !A}[w']$} هله او يوازې هله چې \LR{$\mSat{M}{!A}[w']$} وي۔ له دې ټولو
      څخه \LR{$\mSat{M}{\lnot\Box\lnot !A}[w]$} هله او يوازې هله چې د يوه
      داسې~\LR{$w'$} دپاره چې \LR{$Rww'$} وي، \LR{$\mSat{M}{!A}[w']$} وي۔ بيا د
      \LR{$\mSat{M}{}[w]$} د تعريف له مخې دا هله او يوازې هله چې \LR{$\mSat{M}{\Diamond
        !A}[w]$} وي۔
  \end{enumerate}
\end{proof}
\olpseganchor{OLP-0413-B011}{end}



\olpseganchor{OLP-0413-B013}{start}
\begin{prob}
  که \LR{$\mModel{M} = \tuple{W, R, V}$} يو مدل وي، او فرض کړئ \LR{$w_1, w_2 \in
  W$} داسې وي چې:
  \begin{enumerate}
  \item \LR{$w_1 \in V(p)$} هله او يوازې هله چې \LR{$w_2 \in V(p)$} وي (د هر بياني متغير \LR{$p$} دپاره)؛ او
  \item د هر \LR{$w \in W$} دپاره: \LR{$Rw_1w$} هله او يوازې هله چې \LR{$Rw_2w$} وي۔
  \end{enumerate}
  پر فارمولونه په استقرا سره وښيئ چې د ټولو فارمولونه~\LR{$!A$} دپاره:
  \LR{$\mSat{M}{!A}[w_1]$} هله او يوازې هله چې \LR{$\mSat{M}{!A}[w_2]$} وي۔
\end{prob}
\olpseganchor{OLP-0413-B013}{end}

\olpseganchor{OLP-0413-B014}{start}
\begin{prob}
  که \LR{$\mModel{M} = \tuple{W, R, V}$} وي، وښيئ چې
  \LR{$\mSat{M}{\lnot\Diamond !A}[w]$} هله او يوازې هله چې
  \LR{$\mSat{M}{\Box\lnot!A}[w]$} وي۔
\end{prob}
\olpseganchor{OLP-0413-B014}{end}

\olpseganchor{OLP-0414-B004}{start}
\olfileid{nml}{syn}{tru}
\olpseganchor{OLP-0414-B004}{end}

\olpseganchor{OLP-0414-B005}{start}
\olsection{په يوه مدل کښې رښتياوالی}
\olpseganchor{OLP-0414-B005}{end}

\olpseganchor{OLP-0414-B006}{start}
کله کله مونږ دا کتل غواړو چې په يوه ټاکلي مدل کښې کوم فارمولونه
په هره نړۍ کښې رښتيا دي۔ د دې دپاره يوه نښه ټاکو۔
\olpseganchor{OLP-0414-B006}{end}

\olpseganchor{OLP-0414-B007}{start}
\begin{defn}
  فارمول~\LR{$!A$} په مدل \LR{$M = \tuple{W, R,
    V}$} کښې \emph{رښتيا} دے، او دا په \LR{$\mSat{M}{!A}$} ليکو، هله او
  يوازې هله چې د هر \LR{$w \in W$} دپاره \LR{$\mSat{M}{!A}[w]$} وي۔
\end{defn}
\olpseganchor{OLP-0414-B007}{end}

\olpseganchor{OLP-0414-B008}{start}
\begin{prop}\ollabel{prop:truthfacts}
  \begin{enumerate}
  \item که \LR{$\mSat{M}{!A}$} وي، نو \LR{$\mSat/{M}{\lnot !A}$} وي؛ خو سرچپه
    نتيجه يې \emph{نۀ} سمېږي۔
  \item که \LR{$\mSat{M}{!A \lif !B}$} وي، نو \LR{$\mSat{M}{!A}$} يوازې هله
    چې \LR{$\mSat{M}{!B}$} وي؛ خو سرچپه نتيجه يې \emph{نۀ} سمېږي۔
  \end{enumerate}
\end{prop}
\olpseganchor{OLP-0414-B008}{end}

\olpseganchor{OLP-0414-B009}{start}
\begin{proof}
  \begin{enumerate}
  \item که \LR{$\mSat{M}{!A}$} وي، نو \LR{$!A$} د~\LR{$W$} په ټولو نړۍو کښې رښتيا
    دے۔ څرنګه چې \LR{$W \neq \emptyset$} دے، \LR{$\mSat{M}{\lnot!A}$} نۀ شي
    کېدے؛ ګنې \LR{$!A$} به لږ تر لږه په يوه نړۍ کښې هم رښتيا او هم نارښتيا
    وي۔
\olpseganchor{OLP-0414-B009}{end}

\olpseganchor{OLP-0414-B010}{start}
    خو که \LR{$\mSat/{M}{\lnot !A}$} وي، نو \LR{$!A$} په کومه نړۍ \LR{$w \in W$}
    کښې رښتيا دے۔ له دې نۀ لازمېږي چې \LR{$\mSat{M}{!A}[w]$} د
    \emph{هر}~\LR{$w \in W$} دپاره وي۔ د بېلګې په توګه، د
    \olref[rel]{fig:simple} په مدل کښې \LR{$\mSat/{M}{\lnot p}$} هم وي او
    \LR{$\mSat/{M}{p}$} هم وي۔
  \item فرض کړئ \LR{$\mSat{M}{!A \lif !B}$} او \LR{$\mSat{M}{!A}$} دواړه وي۔
    د \LR{$\mSat{M}{!B}$} د ښودلو دپاره يوه نړۍ \LR{$w \in W$} په عمومي ډول په پام کښې ونيسئ۔ بيا
    \LR{$\mSat{M}{!A \lif !B}[w]$} او \LR{$\mSat{M}{!A}[w]$} دي، نو
    \LR{$\mSat{M}{!B}[w]$} دے۔ څرنګه چې \LR{$w$} هره نړۍ وه، \LR{$\mSat{M}{!B}$}
    دے۔
\olpseganchor{OLP-0414-B010}{end}

\olpseganchor{OLP-0414-B011}{start}
    د سرچپه ادعا د ناسم والي د ښودلو دپاره داسې مدل \LR{$\mModel{M}$}
    پکار دے چې \LR{$\mSat{M}{!A}$} يوازې هله چې \LR{$\mSat{M}{!B}$} وي، خو
    \LR{$\mSat/{M}{!A \lif !B}$} وي۔ بيا د \olref[rel]{fig:simple} مدل
    وڅېړئ: \LR{$\mSat/{M}{p}$} دے، نو (په تش ډول) \LR{$\mSat{M}{p}$} يوازې
    هله چې \LR{$\mSat{M}{q}$} وي۔ خو \LR{$\mSat/{M}{p \lif
      q}$} دے، ځکه په~\LR{$w_1$} کښې \LR{$p$} رښتيا او \LR{$q$} نارښتيا دے۔
  \end{enumerate}
\end{proof}
\olpseganchor{OLP-0414-B011}{end}

\olpseganchor{OLP-0414-B012}{start}
\begin{prob}
  لاندې مدل \LR{$\mModel{M}$} وڅېړئ؛ په ژبه کښې يوازې \LR{$p_1$}، \LR{$p_2$}،
  \LR{$p_3$} بياني متغيرونه دي:
  \begin{center}
    \begin{tikzpicture}[modal]
      \node[world] (w1) [label={[align=right]left:\mTrue{p_1}\\\mFalse{p_2}\\\mFalse{p_3}}]
            {\LR{$w_1$}} ; 
      \node[world] (w2) [label={[align=right]right:\mTrue{p_1}\\\mTrue{p_2}\\\mFalse{p_3}},
        below right=of w1]  {\LR{$w_2$}}; 
      \node[world] (w3) [label={[align=right]right:\mTrue{p_1}\\\mTrue{p_2}\\\mTrue{p_3}},
        above right=of w2] {\LR{$w_3$}};
      \draw[reflexive above] (w3) to (w3); 
      \draw[->] (w1) to (w2);
      \draw[->] (w2) to (w3);
      \draw[->] (w1) to (w3); 
    \end{tikzpicture}
  \end{center}
  له لاندې کوم فارمولونه او قالبونه په مدل \LR{$\mModel{M}$} کښې رښتيا
  دي، يعنې د \LR{$\mModel{M}$} په هره نړۍ کښې رښتيا دي؟ دليل يې ورکړئ۔
  \begin{enumerate}
  \item \LR{$p\lif \Diamond p$} (چې \LR{$p$} اټومي وي)؛
  \item \LR{$!A\lif \Diamond !A$} (چې \LR{$!A$} هره جمله وي)؛
  \item \LR{$\Box p \lif p$} (چې \LR{$p$} اټومي وي)؛
  \item \LR{$\lnot p \lif \Diamond \Box p$} (چې \LR{$p$} اټومي وي)؛
  \item \LR{$\Diamond \Box !A$} (چې \LR{$!A$} هره جمله وي)؛
  \item \LR{$\Box \Diamond p$} (چې \LR{$p$} اټومي وي)۔
  \end{enumerate}
\end{prob}
\olpseganchor{OLP-0414-B012}{end}

\olpseganchor{OLP-0415-B004}{start}
\olfileid{nml}{syn}{val}
\olpseganchor{OLP-0415-B004}{end}

\olpseganchor{OLP-0415-B005}{start}
\olsection{اعتبار}
\olpseganchor{OLP-0415-B005}{end}

\olpseganchor{OLP-0415-B006}{start}
\begin{explain}
  هغه فارمولونه چې په ټولو مدلونو کښې رښتيا وي، يعنې د هر مدل
  په هره نړۍ کښې رښتيا وي، ځانګړې دلچسپي لري۔ دا هغه مودالي قضيې
  څرګندوي چې د \LR{$\Box$} او \LR{$\Diamond$} د تفسير له څرنګوالي پرته رښتيا
  دي، تر هغې چې تفسير په دې مانا «نورمال» وي چې د ممکنه نړۍو ترمنځ
  له کومې لاسرسي اړيکې څخه جوړ شوے وي۔ داسې فارمولونه ته
  \emph{معتبر} وايو۔ د بېلګې په توګه، \LR{$\Box(p \land q) \lif \Box p$}
  معتبر دے۔ خو ځينې فارمولونه چې ښايي د \LR{$\Box$} د ضرورتي تفسير له
  مخې معتبر وګڼل شي، لکه \LR{$\Box p \lif p$}، معتبر نۀ دي۔ د اړيکيزو
  مدلونو د دلچسپۍ يوه وجه دا ده چې د \LR{$\Box$} او \LR{$\Diamond$} بېلابېل
  تفسيرونه د لاسرسي د بېلابېلو ډولونو اړيکو په وسيله ښودل کېدے شي۔
  دا راته وايي چې اعتبار يوازې د \emph{ټولو} مدلونو په نسبت نۀ، بلکې
  د \emph{يوه ټاکلي ډول} د ټولو مدلونو په نسبت هم تعريف کړو۔ د بېلګې
  په توګه، معلومه به شي چې \LR{$\Box p \lif p$} په ټولو هغو مدلونو کښې
  رښتيا دے چې هره نړۍ له خپل ځان څخه د لاسرسي وړ وي، يعنې \LR{$R$}
  انعکاسي وي۔ د مدلونو د ټولګيو په نسبت د اعتبار تعريف مونږ ته دا
  خبره په لنډ ډول د ويلو توان راکوي: \LR{$\Box p \lif p$} د انعکاسي
  مدلونو په ټولګي کښې معتبر دے۔
\end{explain}
\olpseganchor{OLP-0415-B006}{end}

\olpseganchor{OLP-0415-B007}{start}
\begin{defn}
  فارمول~\LR{$!A$} د مدلونو په ټولګي \LR{$\mClass{C}$} کښې \emph{معتبر}
  دے، هله او يوازې هله چې د~\LR{$\mClass{C}$} په هر مدل کښې رښتيا وي (يعنې د~\LR{$\mClass{C}$}
  د هر مدل په هره نړۍ کښې رښتيا وي)۔ که \LR{$!A$} په~\LR{$\mClass{C}$} کښې
  معتبر وي، نو \LR{$\mClass{C} \Entails !A$} ليکو؛ او که \LR{$!A$} د
  \emph{ټولو} مدلونو په ټولګي کښې معتبر وي، نو \LR{$\Entails !A$} ليکو۔
\end{defn}
\olpseganchor{OLP-0415-B007}{end}

\olpseganchor{OLP-0415-B008}{start}
\begin{prop}\ollabel{prop:subset-class}
  که \LR{$!A$} په \LR{$\mClass{C}$} کښې معتبر وي، نو په هر داسې ټولګي
  \LR{$\mClass{C}' \subseteq \mClass{C}$} کښې هم معتبر دے۔
\end{prop}
\olpseganchor{OLP-0415-B008}{end}

\olpseganchor{OLP-0415-B009}{start}
\begin{prop}\ollabel{prop:Nec-rule}
  که \LR{$!A$} معتبر وي، نو \LR{$\Box!A$} هم معتبر دے۔
\end{prop}
\olpseganchor{OLP-0415-B009}{end}

\olpseganchor{OLP-0415-B010}{start}
\begin{proof}
  فرض کړئ \LR{$\Entails !A$}۔ د \LR{$\Entails \Box!A$} د ښودلو دپاره يو
  مدل \LR{$\mModel{M} = \tuple{W, R, V}$} او يوه نړۍ \LR{$w \in W$} واخلئ۔
  که \LR{$Rww'$} وي، نو د \LR{$!A$} د اعتبار له امله \LR{$\mSat{M}{!A}[w']$} وي،
  او په دې توګه \LR{$\mSat{M}{\Box!A}[w]$} هم وي۔ څرنګه چې د
  \LR{$\mModel{M}$} او \LR{$w$} په ټاکنه کښې مو بل شرط نۀ و لګولے، \LR{$\Entails \Box!A$} دے۔
\end{proof}
\olpseganchor{OLP-0415-B010}{end}

\olpseganchor{OLP-0415-B011}{start}
\begin{prob}
  وښيئ چې لاندې جملې معتبرې دي:
  \begin{enumerate}
  \item \LR{$\Box p \lif \Box (q \lif p)$}؛
  \item \LR{$\Box \lnot \lfalse$}؛
  \item \LR{$\Box p \lif (\Box q \lif \Box p)$}۔
  \end{enumerate}
\end{prob}
\olpseganchor{OLP-0415-B011}{end}

\olpseganchor{OLP-0415-B012}{start}
\begin{prob}
  وښيئ چې \LR{$!A \lif \Box!A$} د هغو مدلونو \LR{$\mModel{M} = \tuple{W, R, V}$}
  په ټولګي \LR{$\mClass{C}$} کښې معتبر دے چې \LR{$W = \{w\}$} وي۔ همدارنګه،
  وښيئ چې \LR{$!B \lif \Box !A$} او \LR{$\Diamond !A \lif !B$} د هغو مدلونو
  \LR{$\mModel{M} = \tuple{W, R, V}$} په ټولګي کښې معتبر دي چې
  \LR{$R = \emptyset$} وي۔
\end{prob}
\olpseganchor{OLP-0415-B012}{end}

\olpseganchor{OLP-0416-B004}{start}
\olfileid{nml}{syn}{tau}
\olpseganchor{OLP-0416-B004}{end}

\olpseganchor{OLP-0416-B005}{start}
\olsection{د تاوتولوژۍ تعويضي بېلګې}
\olpseganchor{OLP-0416-B005}{end}

\olpseganchor{OLP-0416-B006}{start}
\begin{explain}
يو نامودالي فارمول هغه وخت تاوتولوژي وي چې د رښتياوالي د قيمت د هرې
ټاکنې لاندې رښتيا وي۔ ښکاره ده چې هره تاوتولوژي د هر مدل په هره نړۍ
کښې رښتيا ده۔ خو د هغو فارمولونه دپاره چې \LR{$\Box$} او~\LR{$\Diamond$}
لري، د تاوتولوژۍ مفهوم نۀ دے تعريف شوے۔ د بېلګې په توګه، آيا
\LR{$\Box p \lor \lnot \Box p$}، چې د درېم حالت د نفي د اصل يوه بېلګه
ده، معتبر دے؟ د \emph{تاوتولوژيکي تعويضي بېلګې} مفهوم مرسته کوي:
هغه فارمول چې د يوې (نامودالي) تاوتولوژۍ د تعويض بېلګه وي۔
دا چې هره داسې بېلګه معتبره ده، د حيرانتيا خبره نۀ ده؛ خو بيا هم
ثبوت غواړي۔
\end{explain}
\olpseganchor{OLP-0416-B006}{end}

\olpseganchor{OLP-0416-B007}{start}
\begin{defn}
  مودالي فارمول~\LR{$!B$} هله او يوازې هله يوه \emph{تاوتولوژيکي
    تعويضي بېلګه} ده چې داسې نامودالي تاوتولوژي~\LR{$!A$} د بياني متغيرونه \LR{$p_1$}، \dots،~\LR{$p_n$} سره، او داسې فارمولونه \LR{$!D_1$}،
  \dots،~\LR{$!D_n$} وي چې \LR{$!B \ident \SSubst{!A}{\subst{!D_1}{p_1},
    \dots, \subst{!D_n}{p_n}}$} وي۔
\end{defn}
\olpseganchor{OLP-0416-B007}{end}

\olpseganchor{OLP-0416-B008}{start}
\begin{lem}\ollabel{lem:valid-taut}
  فرض کړئ \LR{$!A$} داسې نامودالي فارمول وي چې بياني متغيرونه ئې \LR{$p_1$}، \dots، \LR{$p_n$} دي؛ او \LR{$!D_1$}، \dots،
  \LR{$!D_n$} مودالي فارمولونه وي۔ بيا د هرې ټاکنې \LR{$\pAssign{v}$}،
  هر مدل \LR{$\mModel{M} = \tuple{W, R, V}$}، او هر \LR{$w \in W$} دپاره،
  چې د هر شاخص دپاره \LR{$\pAssign{v}(p_i) = \True$} هله او يوازې هله چې
  \LR{$\mSat{M}{!D_i}[w]$} وي، \LR{$\pSat{v}{!A}$} هله او يوازې هله چې
  \LR{$\mSat{M}{\SSubst{!A}{\subst{!D_1}{p_1}, \dots,
      \subst{!D_n}{p_n}}}[w]$}.
\end{lem}
\olpseganchor{OLP-0416-B008}{end}

\olpseganchor{OLP-0416-B009}{start}
\begin{proof}
  پر \LR{$!A$} په استقرا سره ثبوت کوو۔
  \begin{enumerate}
  \item \indcase{!A}{\lfalse}{\LR{$\pSat/{v}{\lfalse}$}
      او \LR{$\mSat/{M}{\lfalse}[w]$} دواړه سم دي}.
  \item \indcase{!A}{\ltrue}{\LR{$\pSat{v}{\ltrue}$}
      او \LR{$\mSat{M}{\ltrue}[w]$} دواړه سم دي}.
  \item \indcase{!A}{p_i}{%
    \begin{align*}
      \pSat{v}{p_i} \Leftrightarrow {} & \pAssign{v}(p_i) = \True \\
      & \qquad \text{\RL{د \LR{$\pSat{v}{p_i}$} د تعريف له مخې}}\\
      \Leftrightarrow {} & \mSat{M}{!D_i}[w] \\
      &\qquad \text{\RL{د فرض له مخې}}\\
      \Leftrightarrow {} & \mSat{M}{\SSubst{p_i}{\subst{!D_1}{p_1}, \dots,
          \subst{!D_n}{p_n}}}[w]\\
      &\qquad \text{\RL{ځکه چې \LR{$\SSubst{p_i}{\subst{!D_1}{p_1}, \dots,
          \subst{!D_n}{p_n}} \ident !D_i$} دے}}.
      \end{align*}}
  \item \indcase{!A}{\lnot !B}{%
      \begin{align*}
        \pSat{v}{\lnot !B} \Leftrightarrow {} & \pSat/{v}{!B}\\
        &\qquad \text{\RL{د \LR{$\pSat{v}{}$} د تعريف له مخې}};\\
        \Leftrightarrow {} & \mSat/{M}{\SSubst{!B}{\subst{!D_1}{p_1}, \dots,
          \subst{!D_n}{p_n}}}[w]\\
        &\qquad \text{\RL{د استقرايي فرض له مخې}}\\
        \Leftrightarrow {} &
        \mSat{M}{\SSubst{\lnot !B}{\subst{!D_1}{p_1}, \dots,
          \subst{!D_n}{p_n}}}[w]\\
        &\qquad \text{\RL{د \LR{$\mSat{M}{}[w]$} د تعريف له مخې}}.
      \end{align*}}
  \item \indcase{!A}{(!B \land !C)}{%
      \begin{align*}
        \pSat{v}{!B \land !C} \Leftrightarrow {} &
        \pSat{v}{!B} \text{\RL{ او }} \pSat{v}{!C}\\
        &\qquad \text{\RL{د \LR{$\pSat{v}{}$} د تعريف له مخې}}\\
        \Leftrightarrow {} &
        \mSat{M}{\SSubst{!B}{\subst{!D_1}{p_1}, \dots,
            \subst{!D_n}{p_n}}}[w] \text{\RL{ او }} \\
        & \mSat{M}{\SSubst{!C}{\subst{!D_1}{p_1}, \dots,
          \subst{!D_n}{p_n}}}[w]\\
        &\qquad \text{\RL{د استقرايي فرض له مخې}}\\
        \Leftrightarrow{} &
        \mSat{M}{\SSubst{(!B \land !C)}{\subst{!D_1}{p_1}, \dots,
          \subst{!D_n}{p_n}}}[w]\\
        &\qquad \text{\RL{د \LR{$\mSat{M}{}[w]$} د تعريف له مخې}}.
      \end{align*}}
  \item \indcase{!A}{(!B \lor !C)}{%
      \begin{align*}
        \pSat{v}{!B \lor !C} \Leftrightarrow{} &
        \pSat{v}{!B} \text{\RL{ يا }} \pSat{v}{!C}\\
        &\qquad \text{\RL{د \LR{$\pSat{v}{}$} د تعريف له مخې}};\\
        \Leftrightarrow{} & \mSat{M}{\SSubst{!B}{\subst{!D_1}{p_1}, \dots,
            \subst{!D_n}{p_n}}}[w] \text{\RL{ يا }}\\
        & \mSat{M}{\SSubst{!C}{\subst{!D_1}{p_1}, \dots,
          \subst{!D_n}{p_n}}}[w]\\
        &\qquad \text{\RL{د استقرايي فرض له مخې}}\\
        \Leftrightarrow{} &
        \mSat{M}{\SSubst{(!B \lor !C)}{\subst{!D_1}{p_1}, \dots,
          \subst{!D_n}{p_n}}}[w]\\
        &\qquad \text{\RL{د \LR{$\mSat{M}{}[w]$} د تعريف له مخې}}.
      \end{align*}}
  \item \indcase{!A}{(!B \lif !C)}{%
      \begin{align*}
        \pSat{v}{!B \lif !C} \Leftrightarrow{} &
        \pSat/{v}{!B} \text{\RL{ يا }} \pSat{v}{!C}\\
        &\qquad \text{\RL{د \LR{$\pSat{v}{}$} د تعريف له مخې}}\\
        \Leftrightarrow{} &
        \mSat/{M}{\SSubst{!B}{\subst{!D_1}{p_1}, \dots,
            \subst{!D_n}{p_n}}}[w] \text{\RL{ يا }}\\
        & \mSat{M}{\SSubst{!C}{\subst{!D_1}{p_1}, \dots,
          \subst{!D_n}{p_n}}}[w]\\
        &\qquad \text{\RL{د استقرايي فرض له مخې}}\\
        \Leftrightarrow{} &
        \mSat{M}{\SSubst{(!B \lif !C)}{\subst{!D_1}{p_1}, \dots,
          \subst{!D_n}{p_n}}}[w]\\
        &\qquad \text{\RL{د \LR{$\mSat{M}{}[w]$} د تعريف له مخې}}.
      \end{align*}}
  \item \indcase{!A}{(!B \liff !C)}{%
      \begin{align*}
        \pSat{v}{!B \liff !C} \Leftrightarrow {} &
        \text{\RL{يا خو }} \pSat{v}{!B} \text{\RL{ او }} \pSat{v}{!C}\\
        & \text{\RL{ يا }} \pSat/{v}{!B} \text{\RL{ او }} \pSat/{v}{!C}\\
        &\qquad \text{\RL{د \LR{$\pSat{v}{}$} د تعريف له مخې}}\\
        \Leftrightarrow {} &
        \text{\RL{يا خو }} \mSat{M}{\SSubst{!B}{\subst{!D_1}{p_1}, \dots,
            \subst{!D_n}{p_n}}}[w] \text{\RL{ او }}\\
        & \qquad \mSat{M}{\SSubst{!C}{\subst{!D_1}{p_1}, \dots,
          \subst{!D_n}{p_n}}}[w]\\ 
        &  \text{\RL{ يا }} \mSat/{M}{\SSubst{!B}{\subst{!D_1}{p_1}, \dots,
            \subst{!D_n}{p_n}}}[w] \text{\RL{ او }}\\
        & \qquad \mSat/{M}{\SSubst{!C}{\subst{!D_1}{p_1}, \dots,
          \subst{!D_n}{p_n}}}[w]\\
        &\qquad\text{\RL{د استقرايي فرض له مخې}}\\
        \Leftrightarrow {} &
        \mSat{M}{\SSubst{(!B \liff !C)}{\subst{!D_1}{p_1}, \dots,
          \subst{!D_n}{p_n}}}[w]\\
        &\qquad \text{\RL{د \LR{$\mSat{M}{}[w]$} د تعريف له مخې}}.
      \end{align*}}
  \end{enumerate}
  \end{proof}
\olpseganchor{OLP-0416-B009}{end}

\olpseganchor{OLP-0416-B010}{start}
\begin{prop}\ollabel{prop:valid-taut}
  ټولې تاوتولوژيکي تعويضي بېلګې معتبرې دي۔
\end{prop}
\olpseganchor{OLP-0416-B010}{end}

\olpseganchor{OLP-0416-B011}{start}
\begin{proof}
  د مقابل نقيض له لارې: فرض کړئ \LR{$!A$} يو نامودالي فارمول دے او داسې تعويض لري چې
  \LR{$\mSat/{M}{\SSubst{!A}{\subst{!D_1}{p_1}, \dots,
  \subst{!D_n}{p_n}}}[w]$} د کوم مدل \LR{$\mModel{M}$} او نړۍ~\LR{$w$}
  دپاره وي۔ داسې ټاکنه \LR{$\pAssign{v}$} تعريف کړئ چې \LR{$\pAssign{v}(p_i) =
  \True$} هله او يوازې هله چې \LR{$\mSat{M}{!D_i}[w]$} وي (او \LR{$\pAssign{v}$}
  د \LR{$q \notin \{p_1, \dots, p_n \}$} دپاره هر ډول قيمتونه ټاکي)۔ بيا د
  \olref{lem:valid-taut} له مخې \LR{$\pSat/{v}{!A}$} وي، نو \LR{$!A$}
  تاوتولوژي نۀ ده۔
\end{proof}
\olpseganchor{OLP-0416-B011}{end}

\olpseganchor{OLP-0417-B004}{start}
\olfileid{nml}{syn}{sch}
\olpseganchor{OLP-0417-B004}{end}

\olpseganchor{OLP-0417-B005}{start}
\olsection{طرحې او اعتبار}
\olpseganchor{OLP-0417-B005}{end}

\olpseganchor{OLP-0417-B006}{start}
\begin{defn}
  يوه \emph{طرحه} د فارمولونه داسې مجموعه ده چې د يوه مودالي
  فارمول~\LR{$!C$} ټولې او يوازې هماغه تعويضي بېلګې لري؛ يعنې
  \[
  \Setabs{!B}{\lexists[!D_1], \dots, \lexists[!D_n] \left(!B =
  \SSubst{!C}{\subst{!D_1}{p_1}, \dots, \subst{!D_n}{p_n}} \right) }.
  \]
  فارمول~\LR{$!C$} د دې طرحې \emph{ځانګړونکے} فارمول بلل کېږي،
  او د بياني متغيرونه د نومونو تر بدلون پورې يکتا دے۔
  فارمول~\LR{$!A$} هله د يوې طرحې \emph{بېلګه} ده چې د هغې د
  مجموعې غړے وي۔
\end{defn}
\olpseganchor{OLP-0417-B006}{end}

\olpseganchor{OLP-0417-B007}{start}
اسانه ده چې يوه طرحه په هغه مابعدالژبني عبارت وښيو چې د~\LR{$!C$} د
بياني متغيرونو پر ځاے پکښې `\LR{$!A$}'، `\LR{$!B$}'،~\dots، اېښودل شوي وي۔ د هر متغير ټول راتګونه په يوه نښه بدلوو، او بياني
ثابتونه پر خپل ځاے ساتو۔
نو، د بېلګې په توګه، لاندې عبارتونه طرحې ښيي:
`\LR{$!A$}'، `\LR{$!A \lif \Box!A$}'، `\LR{$!A \lif (!B \lif !A)$}'۔ د دوى
ځانګړونکي فارمولونه په ترتيب سره \LR{$p$}، \LR{$p \lif \Box p$}،
او \LR{$p \lif (q \lif p)$} دي۔ د `\LR{$!A$}' طرحه د \emph{ټولو}
فارمولونه مجموعه ښيي۔
\olpseganchor{OLP-0417-B007}{end}

\olpseganchor{OLP-0417-B008}{start}
\begin{defn}
  يوه طرحه په يوه مدل کښې هله او يوازې هله \emph{رښتيا} ده چې ټولې
  بېلګې ئې پکې رښتيا وي؛ او هله او يوازې هله \emph{معتبره} ده چې په
  هر مدل کښې رښتيا وي۔
\end{defn}
\olpseganchor{OLP-0417-B008}{end}

\olpseganchor{OLP-0417-B009}{start}
\begin{prop}\ollabel{prop:Kvalid}
  لاندې \Ax{K} طرحه معتبره ده:
  \begin{equation*}
  \Box(!A \lif !B) \lif (\Box !A \lif \Box !B). \tag{\Ax{K}}
  \end{equation*}
\end{prop}
\olpseganchor{OLP-0417-B009}{end}

\olpseganchor{OLP-0417-B010}{start}
\begin{proof}
  بايد وښيو چې د طرحې ټولې بېلګې د هر مدل په هره نړۍ کښې رښتيا دي۔
  نو \LR{$\mModel{M} = \tuple{W,R,V}$} او \LR{$w \in W$} خپلسري ونيسئ۔ د دې
  دپاره چې وښيو يو شرطي عبارت په يوه نړۍ کښې رښتيا دے، فرض کوو چې
  مقدم ئې رښتيا دے، او بيا ښيو چې تالي ئې هم رښتيا دے۔ دلته فرض کړئ
  چې \LR{$\mSat{M}{\Box(!A \lif !B)}[w]$} او \LR{$\mSat{M}{\Box !A}[w]$}۔
  بايد \LR{$\mSat{M}{\Box !B}[w]$} وښيو۔ نو داسې خپلسرې نړۍ \LR{$w'$} ونيسئ
  چې \LR{$Rww'$} وي۔ د لومړي فرض له مخې \LR{$\mSat{M}{!A \lif !B}[w']$}، او
  د دويم فرض له مخې \LR{$\mSat{M}{!A}[w']$}۔ له دې څخه
  \LR{$\mSat{M}{!B}[w']$} راځي۔ ځکه چې \LR{$w'$} خپلسره وه، نو
  \LR{$\mSat{M}{\Box !B}[w]$}۔
\end{proof}
\olpseganchor{OLP-0417-B010}{end}

\olpseganchor{OLP-0417-B011}{start}
\begin{prop}\ollabel{prop:Dual-valid}
  لاندې \Dual{} طرحه معتبره ده:
  \begin{equation*}
  \Diamond !A \liff \lnot\Box\lnot !A. \tag{\Dual}
  \end{equation*}
\end{prop}
\olpseganchor{OLP-0417-B011}{end}

\olpseganchor{OLP-0417-B012}{start}
\begin{proof}
  ثبوت په لاندې تمرين کښې غوښتل شوے دے۔
\end{proof}
\olpseganchor{OLP-0417-B012}{end}

\olpseganchor{OLP-0417-B013}{start}
\begin{prob}
  \olref[nml][syn][sch]{prop:Dual-valid} ثابته کړئ۔
\end{prob}
\olpseganchor{OLP-0417-B013}{end}

\olpseganchor{OLP-0417-B014}{start}
\begin{prop}\ollabel{prop:soundMP}
  که \LR{$!A$} او \LR{$!A \lif !B$} په يوه مدل کښې په يوه نړۍ کښې رښتيا وي،
  نو \LR{$!B$} هم هلته رښتيا دے۔ نو د معتبره فارمولونه مجموعه د
  مودس پوننس تر قاعدې لاندې تړلې ده۔
\end{prop}
\olpseganchor{OLP-0417-B014}{end}

\olpseganchor{OLP-0417-B015}{start}
\begin{prop}\ollabel{prop:valid-instances}
  فارمول~\LR{$!A$} هله او يوازې هله معتبر دے چې ټولې تعويضي
  بېلګې ئې معتبرې وي۔ په بله وينا، يوه طرحه هله او يوازې هله معتبره
  ده چې ځانګړونکے فارمول ئې معتبر وي۔
\end{prop}
\olpseganchor{OLP-0417-B015}{end}

\olpseganchor{OLP-0417-B016}{start}
\begin{proof}
  د «که» لوري ثبوت څرګند دے، ځکه چې \LR{$!A$} د خپل ځان تعويضي بېلګه ده۔
\olpseganchor{OLP-0417-B016}{end}

\olpseganchor{OLP-0417-B017}{start}
  د «يوازې که» لوري د ثبوت دپاره لاندې خبره ښيو: فرض کړئ
  \LR{$\mModel{M} = \tuple{W, R, V}$} يو مودالي مدل دے، او
  \LR{$!B \ident \SSubst{!A}{\subst{!D_1}{p_1},\dots,\subst{!D_n}{p_n}}$}
  د~\LR{$!A$} يوه تعويضي بېلګه ده۔ \LR{$\mModel{M'} = \tuple{W, R,
    V'}$} داسې تعريف کړئ چې
  \LR{$V'(p_i) = \Setabs{w}{\mSat{M}{!D_i}[w]}$} وي۔ په دې سټ کښې
  يوازې د همدغه مدل نړۍوې اخلو۔ نورو بياني متغيرونو ته هر ډول ارزښتونه
  ټاکلے شو۔ بيا، د هر
  \LR{$w \in W$} دپاره، \LR{$\mSat{M}{!B}[w]$} هله او يوازې هله چې
  \LR{$\mSat{M'}{!A}[w]$}۔ (د دې خبرې ثبوت د تمرين دپاره پرېږدو۔)
  اوس فرض کړئ \LR{$!A$} معتبر وي، خو د~\LR{$!A$} کومه تعويضي بېلګه \LR{$!B$}
  معتبره نۀ وي۔ بيا د کوم \LR{$\mModel{M} = \tuple{W, R, V}$} او کوم
  \LR{$w \in W$} دپاره \LR{$\mSat/{M}{!B}[w]$} وي۔ خو د پورته خبرې له مخې
  بيا \LR{$\mSat/{M'}{!A}[w]$} وي، نو \LR{$!A$} معتبر نۀ دے؛ دا تناقض دے۔
\end{proof}
\olpseganchor{OLP-0417-B017}{end}

\olpseganchor{OLP-0417-B018}{start}
\begin{prob}
  د \olref[nml][syn][sch]{prop:valid-instances} د ثبوت د «يوازې که»
  په برخه کښې راغلې خبره ثابته کړئ۔ (لارښوونه: پر~\LR{$!A$} استقرا وکاروئ۔)
\end{prob}
\olpseganchor{OLP-0417-B018}{end}

\olpseganchor{OLP-0417-B019}{start}
خو پام وکړئ چې يوه طرحه په يوه مدل کښې هله او يوازې هله رښتيا ده چې
ځانګړونکے فارمول ئې هلته رښتيا وي، دا خبره سمه نۀ ده۔ البته د
«يوازې که» لور سم دے: که د~\LR{$!A$} هره بېلګه په~\LR{$\mModel{M}$} کښې
رښتيا وي، نو \LR{$!A$} پخپله هم په~\LR{$\mModel{M}$} کښې رښتيا دے۔ خو کېدے
شي چې \LR{$!A$} په~\LR{$\mModel{M}$} کښې رښتيا وي، حال دا چې د~\LR{$!A$}
کومه بېلګه په~\LR{$\mModel{M}$} کښې په کومه نړۍ کښې کاذبه وي۔ يوه ډېره ساده
ناقضه بېلګه په پام کښې ونيسئ: \LR{$p$} په داسې مدل کښې چې يوازې يوه
نړۍ~\LR{$w$} لري او \LR{$V(p) = \{w\}$} وي؛ نو \LR{$p$} په~\LR{$w$} کښې رښتيا دے۔
خو \LR{$\lfalse$} د~\LR{$p$} يوه بېلګه ده او په~\LR{$w$} کښې رښتيا نۀ ده۔
\olpseganchor{OLP-0417-B019}{end}

\olpseganchor{OLP-0417-B020}{start}
\begin{prob}
  وښيئ چې له لاندې فارمولونه څخه يو هم معتبر نۀ دے:
  \begin{enumerate}
    \item[\Ax{D}:] \quad \LR{$\Box p \lif \Diamond p$}؛
    \item[\Ax{T}:] \quad \LR{$\Box p \lif p$}؛
    \item[\Ax{B}:] \quad \LR{$p \lif \Box\Diamond p$}؛
    \item[\Ax{4}:] \quad \LR{$\Box p \lif \Box \Box p$}؛
    \item[\Ax{5}:] \quad \LR{$\Diamond p \lif \Box \Diamond
      p$}۔
  \end{enumerate}
\end{prob}
\olpseganchor{OLP-0417-B020}{end}

\olpseganchor{OLP-0417-B021}{start}
\begin{table}[t]
    \centering
    \begin{tabular}{| l || l |}
      \hline
      {\emph{معتبرې طرحې}} & {\emph{نامعتبرې طرحې}} \\
      \hline\hline
      \LR{$\Box(!A \lif !B) \lif (\Diamond !A \lif \Diamond !B)$}
      & \LR{$\Box (!A \lor !B) \lif (\Box !A \lor \Box !B)$} \\
      \LR{$\Diamond (!A \lif !B) \lif (\Box !A \lif \Diamond
      !B)$}
      & \LR{$(\Diamond !A \land \Diamond !B) \lif \Diamond (!A
      \land !B)$}\\
      \LR{$\Box (!A \land !B) \liff (\Box !A \land \Box !B)$}
      & \LR{$!A \lif \Box !A$} \\
      \LR{$\Box !A \lif \Box (!B \lif !A)$}
      & \LR{$\Box \Diamond !A \lif !B$} \\
      \LR{$\lnot \Diamond !A \lif \Box (!A \lif !B)$}
      & \LR{$\Box \Box !A \lif \Box !A$} \\
      \LR{$\Diamond (!A \lor !B) \liff (\Diamond !A \lor
      \Diamond !B)$}
      & \LR{$\Box \Diamond !A \lif \Diamond \Box !A$}. \\
      \hline
    \end{tabular}
    \caption{معتبرې او (يا؟) نامعتبرې طرحې۔}
    \ollabel{tab:valid-invalidSchemas}
\end{table}
\olpseganchor{OLP-0417-B021}{end}

\olpseganchor{OLP-0417-B022}{start}
\begin{prob}%\ollabel{ex:in/validSchemas}
  ثابته کړئ چې د \olref[nml][syn][sch]{tab:valid-invalidSchemas}
  په لومړۍ ستنه کښې طرحې معتبرې دي او په دويمه ستنه کښې طرحې
  معتبرې نۀ دي۔
\end{prob}
\olpseganchor{OLP-0417-B022}{end}

\olpseganchor{OLP-0417-B023}{start}
\begin{prob}
  معلومه کړئ چې لاندې طرحې معتبرې دي که نامعتبرې:
  \begin{enumerate}
  \item \LR{$(\Diamond !A \lif \Box !B) \lif (\Box !A \lif \Box
    !B)$}؛
  \item \LR{$\Diamond(!A \lif !B) \lor \Box(!B \lif !A)$}۔
  \end{enumerate}
\end{prob}
\olpseganchor{OLP-0417-B023}{end}

\olpseganchor{OLP-0417-B024}{start}
\begin{prob}
  د لاندې هرې طرحې دپاره داسې مدل \LR{$\mModel{M}$} ومومئ چې د
  فارمول هره بېلګه په~\LR{$\mModel{M}$} کښې رښتيا وي:
  \begin{enumerate}
  \item \LR{$p \lif \Diamond\Diamond p$}؛
  \item \LR{$\Diamond p \lif \Box p$}۔
  \end{enumerate}
\end{prob}
\olpseganchor{OLP-0417-B024}{end}

\paragraph{د منجمدې سرچينې بشپړه اړونده برخه}
\olpseganchor{OLP-0654-B004}{start}
\olfileid{nml}{syn}{nor-olp0654}
\olpseganchor{OLP-0654-B004}{end}

\olpseganchor{OLP-0654-B005}{start}
\olsection{عادي مودالي منطقونه}
\olpseganchor{OLP-0654-B005}{end}

\olpseganchor{OLP-0654-B006}{start}
مودالي منطق د فارمولونه هر داسې سټ ګڼلے شو چې مناسب خاصيتونه ولري۔ ځينې مودالي منطقونه ځانګړې دلچسپي لري، ځکه تعبيرونه او کارونې ئې په زړه پورې دي، خو په تېره بيا ځکه چې تيوري ئې په ښه توګه معلومه ده۔ دا هماغه منطقونه دي چې عادي مودالي منطقونه بللے کېږي۔
\olpseganchor{OLP-0654-B006}{end}

\olpseganchor{OLP-0654-B007}{start}
\begin{defn}\ollabel{defn:modal-logic}
\emph{مودالي منطق}~\LR{$\Log L$} د بنسټيزې مودالي ژبې د فارمولونه داسې يو سټ دے چې
\begin{enumerate}
\item ټولې قضيوي تاوتولوژۍ لري،
\item د يو‌شان ځاے ناستولو لاندې تړلے وي، او
\item د مودوس پوننس لاندې تړلے وي؛ يعنې هر کله چې \LR{$!A$} او \LR{$(!A \lif
  !B) \in \Log L$} وي، نو \LR{$!B \in \Log L$} هم وي۔
\end{enumerate}
\end{defn}
\olpseganchor{OLP-0654-B007}{end}

\olpseganchor{OLP-0654-B008}{start}
\begin{ex}
د فارمولونه ځينې داسې سټونه چې چندان په زړه پورې نۀ دي، خو دا تعريف پوره کوي او ځکه مودالي منطقونه ګڼلے کېږي، دا دي:
\begin{enumerate}
\item د ټولو تاوتولوژيو سټ،
\item د ټولو فارمولونه سټ (ناسازګار مودالي منطق)۔
\end{enumerate}
\end{ex}
\olpseganchor{OLP-0654-B008}{end}


\clearpage
\olpseganchor{OLP-0654-B009}{start}
\begin{defn}\label{defn:normal-modal-logic}
مودالي منطق~\LR{$\Log L$} هله او يوازې هله \emph{عادي} دے چې
\begin{enumerate}
\item دا دواړه ولري:
\begin{align}
\tag{K} \Box(p \lif q) & \lif (\Box p \lif \Box q) \\
\tag{Dual} \Diamond p & \liff \lnot \Box \lnot p
\end{align}
\item او د لزوم ورکولو لاندې تړلے وي؛ يعنې هر کله چې \LR{$!A \in
  \Log L$} وي، نو \LR{$\Box !A \in \Log L$} هم وي۔
\end{enumerate}
\end{defn}
\olpseganchor{OLP-0654-B009}{end}

\olpseganchor{OLP-0654-B010}{start}
د مودالي منطقونو د تعريف دوه عمده لارې به ووينو۔ يوه د ثبوت له تيورۍ څخه ده: د هغو فارمولونه د سټ په توګه چې په ټاکلو اشتقاق نظامونو کښې ترې اشتقاق کېدے شي۔ بله د مدل له تيورۍ څخه ده: د هغو فارمولونه د سټونو په توګه چې د مدلونو په ټاکلو ټولګيو کښې معتبر وي۔ دا د معمولي قضيوي منطق (او هم د لومړۍ درجې منطق) د حالت په شان دے۔ هلته هم د فارمولونه يو سټ په دوو بېلو لارو ځانګړے کولے شو؛ د بېلګې په توګه، د هغو تاوتولوژيو سټ چې د صدق ارزښتونو په ټولو ټاکنو کښې رښتينې دي، او د هغو ټولو د اشتقاق وړ فارمولونه سټ چې په يوه اشتقاق نظام کښې ترې اشتقاق کېږي۔ د عادي مودالي منطقونو دپاره دا دوه‌ګونې ټاکنه په اړيکيزو مدلونو او بديهي (د هېلبرټ په ډول) ثبوتي نظامونو کېږي۔
\olpseganchor{OLP-0654-B010}{end}
\clearpage


\olpseganchor{OLP-0418-B004}{start}
\olfileid{nml}{syn}{ent}
\olpseganchor{OLP-0418-B004}{end}

\olpseganchor{OLP-0418-B005}{start}
\olsection{استلزام}
\olpseganchor{OLP-0418-B005}{end}

\olpseganchor{OLP-0418-B006}{start}
\begin{explain}
  په يوه نړۍ کښې د رښتياوالي د تعريف په مرسته د فارمولونه ترمنځ
  د استلزام اړيکه تعريفولے شو۔ فارمول~\LR{$!B$} هله او يوازې هله
  \LR{$!A$} لازموي چې هر وخت \LR{$!B$} رښتيا وي، \LR{$!A$} هم رښتيا وي۔ دلته
  «هر وخت» دا دواړه معناوې لري: «هر کوم مدل چې وڅېړو» او «په هغه
  مدل کښې هره کومه نړۍ چې وڅېړو»۔
\end{explain}
\olpseganchor{OLP-0418-B006}{end}

\olpseganchor{OLP-0418-B007}{start}
\begin{defn}
  که \LR{$\Gamma$} د فارمولونه مجموعه وي او \LR{$!A$} فارمول وي، نو
  \LR{$\Gamma$}، \LR{$!A$} \emph{لازموي}، چې په نښو کښې \LR{$\Gamma \Entails !A$}
  ليکو، هله او يوازې هله چې د هر مدل \LR{$\mModel{M} = \tuple{W, R, V}$}
  او هرې نړۍ \LR{$w \in W$} دپاره، که د هر \LR{$!B \in \Gamma$} دپاره
  \LR{$\mSat{M}{!B}[w]$} وي، نو \LR{$\mSat{M}{!A}[w]$} وي۔ که په \LR{$\Gamma$}
  کښې يوازې يو فارمول~\LR{$!B$} وي، نو \LR{$!B \Entails !A$} ليکو۔
\end{defn}
\olpseganchor{OLP-0418-B007}{end}

\olpseganchor{OLP-0418-B008}{start}
\begin{ex}
  د دې د ښودلو دپاره چې يو فارمول بل لازموي، بايد د
  \LR{$\mSat{M}{}[w]$} د تعريف په کارولو سره د ټولو مدلونو په هکله
  استدلال وکړو۔ د بېلګې په توګه، د \LR{$p \lif \Diamond p \Entails
  \Box\lnot p \lif \lnot p$} د ښودلو دپاره داسې استدلال کېدے شي:
  يو مدل \LR{$\mModel{M} = \tuple{W, R, V}$} او \LR{$w \in W$} ونيسئ، او فرض
  کړئ چې \LR{$\mSat{M}{p \lif \Diamond p}[w]$}۔ بايد وښيو چې
  \LR{$\mSat{M}{\Box\lnot p \lif \lnot p}[w]$}۔ فرض کړئ داسې نۀ وي۔
  بيا \LR{$\mSat{M}{\Box\lnot p}[w]$} او \LR{$\mSat/{M}{\lnot p}[w]$} وي۔
  له \LR{$\mSat/{M}{\lnot p}[w]$} څخه \LR{$\mSat{M}{p}[w]$} راځي۔ د فرض
  له مخې \LR{$\mSat{M}{p \lif \Diamond p}[w]$}، نو
  \LR{$\mSat{M}{\Diamond p}[w]$}۔ د \LR{$\mSat{M}{\Diamond p}[w]$} د تعريف
  له مخې داسې \LR{$w'$} شته چې \LR{$Rww'$} او \LR{$\mSat{M}{p}[w']$} وي۔ خو
  \LR{$\mSat{M}{\Box \lnot p}[w]$} هم دے، نو
  \LR{$\mSat{M}{\lnot p}[w']$} وي؛ دا تناقض دے۔
\olpseganchor{OLP-0418-B008}{end}

\olpseganchor{OLP-0418-B009}{start}
  د دې د ښودلو دپاره چې يو فارمول~\LR{$!B$} بل~\LR{$!A$} نۀ لازموي،
  بايد ناقضه بېلګه ورکړو: داسې مدل \LR{$\mModel{M} = \tuple{W, R, V}$}
  چې په کومه نړۍ~\LR{$w \in W$} کښې \LR{$\mSat{M}{!B}[w]$} وي، خو
  \LR{$\mSat/{M}{!A}[w]$} وي۔ وښيو چې \LR{$p \lif \Diamond p \Entails/
  \Box p \lif p$}۔ د \olref{fig:counterex} مدل وڅېړئ۔ په همدغه نړۍ
  کښې \LR{$\mSat{M}{\Diamond p}[w_1]$} دے، نو
  \LR{$\mSat{M}{p \lif \Diamond p}[w_1]$} هم دے۔ خو
  \LR{$\mSat{M}{\Box p}[w_1]$} سره سره \LR{$\mSat/{M}{p}[w_1]$} دے؛ نو
  \LR{$\mSat/{M}{\Box p \lif p}[w_1]$}۔
  \begin{figure}
  \begin{center}
    \begin{tikzpicture}[modal]
      \node[world] (w1) [label=right:\mFalse{p}]{\LR{$w_1$}}; 
      \node[world] (w2) [label=right:\mTrue{p}, above left=of w1]{\LR{$w_2$}}; 
      \node[world] (w3) [label=right:\mTrue{p}, above right=of w1] {\LR{$w_3$}};
      \draw[->] (w1) to (w2);
      \draw[->] (w1) to (w3);
    \end{tikzpicture}
  \end{center}
  \caption{د \LR{$p \lif \Diamond p
  \Entails \Box p \lif p$} د استلزام ناقضه بېلګه.}
  \ollabel{fig:counterex}
  \end{figure}
\olpseganchor{OLP-0418-B009}{end}
  
\olpseganchor{OLP-0418-B010}{start}
  ډېر وخت ساده ناقضې بېلګې بس وي۔ مدل \LR{$\mModel{M'} =
  \tuple{W', R', V'}$} چې \LR{$W' = \{w\}$}، \LR{$R' = \emptyset$} او \LR{$V'(p) =
  \emptyset$} وي، هم ناقضه بېلګه ده: ځکه چې \LR{$\mSat/{M'}{p}[w]$}، نو
  \LR{$\mSat{M'}{p \lif \Diamond p}[w]$}۔ له~\LR{$w$} څخه هېڅ نړۍ ته لاسرسے
  نشته، نو \LR{$\mSat{M'}{\Box p}[w]$}؛ له دې امله
  \LR{$\mSat/{M'}{\Box p \lif p}[w]$}۔
\end{ex}
\olpseganchor{OLP-0418-B010}{end}

\olpseganchor{OLP-0418-B011}{start}
\begin{prob}
  وښيئ چې \LR{$\Box (!A \land !B) \Entails \Box !A$}۔
\end{prob}
\olpseganchor{OLP-0418-B011}{end}

\olpseganchor{OLP-0418-B012}{start}
\begin{prob}
  وښيئ چې \LR{$\Box (p \lif q) \Entails/ p \lif \Box q$} او \LR{$p \lif \Box
  q \Entails/\Box (p \lif q)$}۔
\end{prob}
\olpseganchor{OLP-0418-B012}{end}

\olpseganchor{OLP-0419-B004}{start}
\olchapter{nml}{frd}{د چوکاټ تعريف‌وړتيا}
\olpseganchor{OLP-0419-B004}{end}

\olpseganchor{OLP-0420-B004}{start}
\olfileid{nml}{frd}{int}
\olpseganchor{OLP-0420-B004}{end}

\olpseganchor{OLP-0420-B005}{start}
\olsection{پېژندنه}
\olpseganchor{OLP-0420-B005}{end}

\olpseganchor{OLP-0420-B006}{start}
مودالي منطقپوهان په دې پوښتنه کښې دلچسپي لري چې د لاسرسي اړيکه او
د هماغې اړيکې لرونکو مدلونو په دننه کښې د ځينو فارمولونه رښتياوالی
له يو بل سره څنګه تړاو لري۔ د بېلګې په توګه، فرض کړئ د لاسرسي اړيکه
انعکاسي وي؛ يعنې د هر \LR{$w \in W$} دپاره \LR{$Rww$} وي۔ په بله وينا، هره
نړۍ له خپل ځان څخه د لاسرسي وړ وي۔ دا معنا لري چې که \LR{$\Box !A$} په
کومه نړۍ~\LR{$w$} کښې رښتيا وي، نو \LR{$w$} پخپله هم د هغو د لاسرسي وړ نړۍو
په ډله کښې ده چې~\LR{$!A$} بايد پکې رښتيا وي۔ نو که د~\LR{$\mModel{M}$} د
لاسرسي اړيکه~\LR{$R$} انعکاسي وي، د هرې نړۍ \LR{$w$} او هر فارمول \LR{$!A$} دپاره
\LR{$\Box !A \lif !A$} هلته رښتيا دے (په بله وينا، د \LR{$\Box p \lif p$}
طرحه او د هغې ټولې تعويضي بېلګې په~\LR{$\mModel{M}$} کښې رښتيا دي)۔
\olpseganchor{OLP-0420-B006}{end}

\olpseganchor{OLP-0420-B007}{start}
خو سرچپه خبره ناسمه ده۔ د بېلګې په توګه، له دې چې \LR{$\Box p \lif p$}
په~\LR{$\mModel{M}$} کښې رښتيا وي، دا نۀ لازمېږي چې \LR{$R$} انعکاسي وي۔
ځکه داسې نامنعکسي مدل~\LR{$\mModel{M}$} په اسانۍ موندلے شو چې \LR{$\Box p
\lif p$} ئې په ټولو نړۍو کښې رښتيا وي: داسې مدل واخلئ چې يوازې يوه
نړۍ~\LR{$w$} ولري، هغې ته له خپل ځان څخه لاسرسے نۀ وي، خو \LR{$w \in V(p)$}
وي۔ د \LR{$p$} د رښتياوالي د قيمت په مناسبه ټاکنه سره کولے شو چې
\LR{$\Box p \lif p$} په داسې مدل کښې رښتيا کړو چې انعکاسي نۀ وي۔
\olpseganchor{OLP-0420-B007}{end}

\olpseganchor{OLP-0420-B008}{start}
د نړۍو سټ او د لاسرسي اړيکه پر خپل حال وساتئ او د متغيرونو ارزښت
ټاکنه~\LR{$V$} بدله کړئ۔ که \LR{$\Box p \lif p$} په~\LR{$\mModel{M}$} کښې په ټولو
نړۍو کښې د~\LR{$V(p)$} له هرې ممکنې ټاکنې سره رښتيا وي، نو اړينه ده چې
\LR{$R$} انعکاسي وي۔ ځکه په هر نامنعکسي مدل کښې به لږ تر لږه يوه نړۍ~\LR{$w$}
وي چې \LR{$Rww$} نۀ وي۔ که \LR{$V(p) = W \setminus \{w\}$} وټاکو، نو \LR{$p$}
به له~\LR{$w$} پرته په ټولو نړۍو کښې رښتيا وي؛ په ځانګړي ډول، له~\LR{$w$}
څخه د لاسرسي وړ په ټولو نړۍو کښې به رښتيا وي (ځکه چې \LR{$w$} له~\LR{$w$}
څخه د لاسرسي وړ نۀ دے او \LR{$w$} يوازينۍ نړۍ ده چې~\LR{$p$} پکښې نارښتيا دے)۔
خو \LR{$p$} په \LR{$w$} کښې نارښتيا دے؛ نو \LR{$\Box p \lif p$} په~\LR{$w$} کښې
نارښتيا دے۔
\olpseganchor{OLP-0420-B008}{end}

\olpseganchor{OLP-0420-B009}{start}
دا ښيي چې د ارزښت ټاکنې له پرته د مدل د جوړښتونو دپاره يوه نښه
وټاکو؛ دوى \emph{چوکاټونه} بولو۔ چوکاټ \LR{$\mModel{F}$} په ساده ډول د
نړۍو د مجموعې او د لاسرسي د اړيکې جوړه \LR{$\tuple{W, R}$} ده۔ بيا هر
مدل \LR{$\tuple{W, R, V}$}، لکه څنګه چې وايو،
پر چوکاټ \LR{$\tuple{W, R}$} \emph{ولاړ} دے۔ برعکس، هر چوکاټ د هغو
مدلونو ټولګے ټاکي چې پر هغه ولاړ وي؛ او د چوکاټونو يو ټولګے د هغو
مدلونو ټولګے ټاکي چې د دې ټولګي پر کوم چوکاټ ولاړ وي۔ همدارنګه،
\LR{$\mModel{F} \Entails !A$} تعريفولے شو، يعنې په يوه چوکاټ کښې د
فارمول د \emph{اعتبار} مفهوم داسې: د هر هغه \LR{$\mModel{M}$}
دپاره \LR{$\mSat{M}{!A}$} چې پر~\LR{$\mModel{F}$} ولاړ وي۔
\olpseganchor{OLP-0420-B009}{end}

\olpseganchor{OLP-0420-B010}{start}
په دې نښه سره د فارمولونه او د چوکاټونو د ټولګيو ترمنځ د مطابقت
اړيکې ټاکلے شو: د بېلګې په توګه، \LR{$\mModel{F} \Entails \Box p \lif p$}
هله او يوازې هله چې \LR{$\mModel{F}$} انعکاسي وي۔
\olpseganchor{OLP-0420-B010}{end}

\olpseganchor{OLP-0421-B004}{start}
\olfileid{nml}{frd}{acc}
\olpseganchor{OLP-0421-B004}{end}

\olpseganchor{OLP-0421-B005}{start}
\olsection{د لاسرسي د اړيکو خاصيتونه}
\olpseganchor{OLP-0421-B005}{end}

\olpseganchor{OLP-0421-B006}{start}
ډېر مودالي فارمولونه د لاسرسي د اړيکې د ساده او ان اشنا خاصيتونو
ځانګړنې څرګندوي۔ له يوې خوا دا معنا لري چې هر هغه مدل، چې يو ټاکلے
خاصيت ولري، اړوند فارمول (او د هغې ټولې تعويضي بېلګې) پکې رښتيا
کوي۔ لومړے د لاسرسي د اړيکو پنځه دوديز ډولونه او هغه فارمولونه ګورو
چې د دغو اړيکو له امله رښتيا وي۔
\olpseganchor{OLP-0421-B006}{end}

\olpseganchor{OLP-0421-B007}{start}
\begin{thm}\ollabel{thm:soundschemas}
  فرض کړئ \LR{$\mModel{M} = \tuple{W, R, V}$} يو مدل دے۔ که \LR{$R$} د
  \olref{tab:five} د چپې خوا خاصيت ولري، نو د ښي خوا د فارمول
  هره بېلګه په~\LR{$\mModel{M}$} کښې رښتيا ده۔
\end{thm}
\olpseganchor{OLP-0421-B007}{end}

\olpseganchor{OLP-0421-B008}{start}
\begin{table}[t]
    \begin{tabular}{| l || l |}
      \hline
      {\emph{که \LR{$R$} \dots وي}} & {\emph{نو \dots په~\LR{$\mModel{M}$} کښې رښتيا دے:}} \\
      \hline \hline
      \emph{مسلسل}: \LR{$\forall u \exists v Ruv$} & \hbox to.43\textwidth
           {\LR{$\Box p \lif \Diamond p$} \hfill (\Ax{D})} \\
      \hline
      \emph{انعکاسي}: \LR{$\forall w Rww$}  
      & \hbox to.43\textwidth
          {\LR{$\Box p \lif p$} \hfill (\Ax{T})} \\
      \hline
      \emph{متناظره}: &  \hbox to .43\textwidth
           {\LR{$p \lif \Box\Diamond p$} \hfill (\Ax{B})} \\
           \LR{$\forall u\forall v(Ruv \lif Rvu)$} & \\
      \hline
      \emph{انتقالي}: & \hbox to.43\textwidth
           {\LR{$\Box p \lif \Box \Box p$} \hfill (\Ax{4})} \\
      \LR{$\forall u \forall v \forall w ((Ruv \land Rvw) \lif Ruw)$} & \\
      \hline 
      \emph{اقليدسي}: & \hbox to .43\textwidth
        {\LR{$\Diamond p \lif \Box\Diamond p$} \hfill (\Ax{5})} \\
      \LR{$\forall w \forall u \forall v ((Rwu \land Rwv) \lif Ruv)$} &\\
      \hline
    \end{tabular}
    \caption{د مطابقت پنځه حقيقتونه۔}
    \ollabel{tab:five}
  \end{table} 
\olpseganchor{OLP-0421-B008}{end}


\olpseganchor{OLP-0421-B009}{start}
\begin{proof}
  د \Ax{B} مورد داسې دے: په يوه مدل کښې د طرحې د رښتياوالي دپاره
  بايد وښيو چې ټولې بېلګې ئې د مدل په ټولو نړۍو کښې رښتيا دي۔ نو د
  \Ax{B} يوه بېلګه \LR{$!A \lif \Box\Diamond !A$} او يوه خپلسرې نړۍ
  \LR{$w \in W$} ونيسئ۔ فرض کړئ مقدم~\LR{$!A$} په~\LR{$w$} کښې رښتيا دے؛ غواړو
  وښيو چې \LR{$\Box \Diamond !A$} هم په~\LR{$w$} کښې رښتيا دے۔ د دې دپاره بايد
  \LR{$\Diamond !A$} له~\LR{$w$} څخه د لاسرسي وړ په هرې نړۍ~\LR{$w'$} کښې رښتيا
  وي۔ اوس چې د کومې~\LR{$w'$} دپاره \LR{$Rww'$} وي، د تناظر له فرضه \LR{$Rw'w$}
  هم راځي (\olref{fig:Bsymm} وګورئ)۔ ځکه چې
  \LR{$\mSat{M}{!A}[w]$}، نو \LR{$\mSat{M}{\Diamond !A}[w']$}۔ او ځکه چې
  \LR{$w'$} د \LR{$Rww'$} لرونکو نړۍو له منځه خپلسره وه، نو
  \LR{$\mSat{M}{\Box\Diamond!A}[w]$}۔
\olpseganchor{OLP-0421-B009}{end}

\olpseganchor{OLP-0421-B010}{start}
  نور موردونه د تمرين دپاره پرېږدو۔
\end{proof}
\olpseganchor{OLP-0421-B010}{end}

\olpseganchor{OLP-0421-B011}{start}
\begin{prob}
  د \olref[nml][frd][acc]{thm:soundschemas} ثبوت بشپړ کړئ۔
\end{prob}
\olpseganchor{OLP-0421-B011}{end}

\olpseganchor{OLP-0421-B012}{start}
\begin{figure}
  \begin{center}
    \begin{tikzpicture}[modal]
      \node[world] (w1) [label={below:\LR{$\mSat{{}}{\formula{A}}$}\\
          \LR{$\mSat{{}}{\Box\Diamond \formula{A}}$}}] {\LR{$w$}}; 
      \node[world] (w2) [label=below:{\LR{$\mSat{{}}{\Diamond \formula{A}}$}},
        right=of w1] {\LR{$w'$}}; 
      \draw[->, bend left] (w1) to (w2); 
      \draw[->, bend left] (w2) to (w1); 
    \end{tikzpicture}
  \end{center}
\caption{له تناظر څخه استدلال۔}
\ollabel{fig:Bsymm}
\end{figure}
\olpseganchor{OLP-0421-B012}{end}

\olpseganchor{OLP-0421-B013}{start}
پام وکړئ چې د \olref{thm:soundschemas} سرچپه استلزامونه سم نۀ دي:
له دې چې يو مدل يوه طرحه رښتيا کوي، دا نۀ لازمېږي چې د هغه مدل د
لاسرسي اړيکه اړوند خاصيت ولري۔ د \Ax{T} او انعکاسي مدلونو په بحث کښې
داسې مدل په اسانۍ موندل کېږي چې پخپله د \Ax{T} يوه بېلګه پکښې
نارښتيا وي: \LR{$W = \{w\}$} او \LR{$V(p) = \emptyset$} وټاکئ، او د لاسرسي
اړيکه تشه ونيسئ۔ بيا \LR{$R$} انعکاسي نۀ ده، خو
\LR{$\mSat{M}{\Box p}[w]$} او \LR{$\mSat/{M}{p}[w]$}۔ په دې دليل مو د \Ax{T}
يوه نارښتيا بېلګه په~\LR{$\mModel{M}$} کښې وښوده؛ ځينې نورې بېلګې، لکه
\LR{$\Box \lnot p \lif \lnot p$}، رښتيا دي۔ د داسې مدل موندل ګران دي
چې د~\Ax{T} \emph{هره تعويضي بېلګه} پکې رښتيا وي او
\LR{$\mModel{M}$} انعکاسي نۀ وي۔ خو داسې مدل~\LR{$\mModel{M}$} هم شته:
\olpseganchor{OLP-0421-B013}{end}

\olpseganchor{OLP-0421-B014}{start}
\begin{prop}\ollabel{prop:reflexive}
  فرض کړئ \LR{$\mModel{M} = \tuple{W, R, V}$} داسې مدل دے چې
  \LR{$W = \{u, v \}$} وي، او نړۍوې \LR{$u$} او \LR{$v$} د~\LR{$R$} له مخې سره تړلې
  وي؛ يعنې \LR{$Ruv$} او \LR{$Rvu$} دواړه وي، او له دې پرته بله اړيکه نۀ وي۔
  فرض کړئ چې د هر \LR{$p$} دپاره \LR{$u \in V(p) \Leftrightarrow v
  \in V(p)$}۔ بيا:
  \begin{enumerate}
  \item د هر \LR{$!A$} دپاره: \LR{$\mSat{M}{!A}[u]$} هله او يوازې هله چې
    \LR{$\mSat{M}{!A}[v]$} (پر~\LR{$!A$} استقرا وکاروئ)۔
  \item د~\Ax{T} هره بېلګه په~\LR{$\mModel{M}$} کښې رښتيا ده۔
  \end{enumerate}
  ځکه چې \LR{$\mModel{M}$} انعکاسي نۀ دے (بلکې په حقيقت کښې
  \emph{نفي انعکاسي} دے)، د \olref{thm:soundschemas} سرچپه لور
  د \Ax{T} په مورد کښې ناسمه ده (د
  \olref{thm:soundschemas} د ځينو نورو طرحو دپاره هم ورته استدلال
  کېدلے شي، خو د ټولو دپاره نۀ)۔
\end{prop}
\olpseganchor{OLP-0421-B014}{end}

\olpseganchor{OLP-0421-B015}{start}
\begin{prob}
  د \olref[nml][frd][acc]{prop:reflexive} دعوې ثابتې کړئ۔
\end{prob}
\olpseganchor{OLP-0421-B015}{end}

\olpseganchor{OLP-0421-B016}{start}
که څه هم زموږ پام به پر پنځو دوديزو فارمولونه \Ax{D}،
\Ax{T}، \Ax{B}، \Ax{4} او \Ax{5} وي، په
\olref{tab:anotherfive} کښې د لاسرسي د اړيکو څو نور خاصيتونه هم
راخلو۔ د لاسرسي اړيکه~\LR{$R$} هله \emph{جزوي تابعي} ده چې له هرې نړۍ
څخه تر ډېره يوې نړۍ ته لاسرسے وي۔ که له هرې نړۍ څخه کټ مټ يوې نړۍ
ته لاسرسے وي، نو دا اړيکه \emph{تابعي} بولو۔ (نو تابعي اړيکې هماغه
دي چې مسلسل او جزوي تابعي دواړه وي۔) دا اړيکې ځکه «تابعي» بلل
کېږي چې د لاسرسي اړيکه د (جزوي) تابعې په شان کار کوي۔ يوه اړيکه
هله \emph{کمزورې ګڼه} ده چې هر کله \LR{$Ruv$} وي، د~\LR{$u$} او~\LR{$v$} ترمنځ
يوه نړۍ~\LR{$w$} وي۔ نو کمزورې ګڼې اړيکې په يوه معنا د انتقالي اړيکو
عکس دي: په انتقالي اړيکه کښې که له~\LR{$u$} څخه د~\LR{$w$} له لارې په چکر
سره~\LR{$v$} ته رسېدے شئ، نو نېغ هم ورته رسېدے شئ؛ په کمزورې ګڼې
اړيکه کښې که له~\LR{$u$} څخه نېغ~\LR{$v$} ته رسېدے شئ، نو له~\LR{$u$} څخه
د کومې~\LR{$w$} له لارې په چکر سره هم~\LR{$v$} ته رسېدے شئ۔ يوه اړيکه هله
\emph{کمزورې هم‌لوري} ده چې که له يوې نړۍ~\LR{$w$} څخه~\LR{$u$} او~\LR{$v$}
دواړو ته رسېدے شئ، نو له~\LR{$u$} او~\LR{$v$} دواړو څخه يوې ګډې نړۍ~\LR{$t$}
ته رسېدے شئ؛ دې ته کله کله «الماسي خاصيت» يا «يوځاے کېدنه» هم
وايي۔
\olpseganchor{OLP-0421-B016}{end}

\olpseganchor{OLP-0421-B017}{start}
\begin{table}[t]
    \begin{tabular}{| p{\dimexpr \linewidth-.43\textwidth-4\tabcolsep-4\arrayrulewidth-\doublerulesep\relax} || p{.43\textwidth} |}
      \hline
      {\emph{که \LR{$R$} \dots وي}} & {\emph{نو \dots په~\LR{$\mModel{M}$} کښې رښتيا دے:}} \\
      \hline \hline
      \emph{جزوي تابعي}: & 
      \multirow{2}{*}{\LR{$\Diamond p \lif \Box p$}} \\
      \LR{$ \forall w \forall u \forall v ((Rwu \land Rwv) \lif u =v )$} & \\
      \hline
      \multirow{2}{*}{\emph{تابعي}: \LR{$\forall w \exists v \forall u(Rwu \liff \eq[u][v]) $}} 
      & \multirow{2}{*}{\LR{$\Diamond p \liff \Box p$}} \\
      & \\
      \hline
      \emph{کمزورې ګڼه}: &  \multirow{2}{*}{\LR{$\Box \Box p \lif
        \Box p$}} \\
      \LR{$\forall u\forall v(Ruv \lif \exists w(Ruw \land Rwv))$} & \\
      \hline
      \emph{کمزورې تړلې}: &
      \multirow{3}{*}{\hbox to.43\textwidth{\LR{$
        \begin{array}{@{}l@{}}
          \Box ((p \land \Box p) \lif q) \lor {} \\
          \qquad \Box ((q \land \Box q) \lif p)
        \end{array}$} \hfill (\Ax{L})}} \\
      \LR{$\forall w \forall u \forall v ((Rwu \land Rwv) \lif {}$} &\\
      \qquad \LR{$(Ruv \lor u=v \lor Rvu))$} & \\
      \hline 
      \emph{کمزورې هم‌لوري}: & \multirow{3}{*}{\hbox to .43\textwidth
        {\LR{$\Diamond\Box p \lif \Box\Diamond p$}\hfill (\Ax{G})}} \\
      \LR{$\forall w \forall u \forall v ((Rwu \land Rwv) \lif {}$} & \\
      \qquad \LR{$\exists t (Rut \land Rvt))$} & \\
      \hline
    \end{tabular}
    \caption{د مطابقت پنځه نور حقيقتونه۔}
    \ollabel{tab:anotherfive}
  \end{table} 
\olpseganchor{OLP-0421-B017}{end}

\olpseganchor{OLP-0421-B018}{start}
\begin{prob}
  فرض کړئ \LR{$\mModel{M} = \tuple{W, R, V}$} يو مدل دے۔ وښيئ چې که
  \LR{$R$} د \olref[nml][frd][acc]{tab:anotherfive} د چپې خوا خاصيتونه
  ولري، نو د ښي خوا د هر اړوند فارمول هره بېلګه په~\LR{$\mModel{M}$}
  کښې رښتيا ده۔
\end{prob}
\olpseganchor{OLP-0421-B018}{end}

\olpseganchor{OLP-0422-B004}{start}
\olfileid{nml}{frd}{fra}
\olpseganchor{OLP-0422-B004}{end}

\olpseganchor{OLP-0422-B005}{start}
\olsection{چوکاټونه}
\olpseganchor{OLP-0422-B005}{end}

\olpseganchor{OLP-0422-B006}{start}
\begin{defn}
  \emph{چوکاټ} يوه جوړه \LR{$\mModel{F} = \tuple{W,R}$} ده، چې \LR{$W$} د
  نړۍو ناتشه مجموعه او \LR{$R$} پر~\LR{$W$} يوه دوه‌ځاييزه اړيکه ده۔ مدل
  \LR{$\mModel{M}$} هله او يوازې هله پر چوکاټ \LR{$\mModel{F} = \tuple{W,R}$}
  \emph{ولاړ} دے چې د کومې ارزښت ټاکنې \LR{$V$} دپاره
  \LR{$\mModel{M} = \tuple{W, R, V}$} وي۔
\end{defn}
\olpseganchor{OLP-0422-B006}{end}

\olpseganchor{OLP-0422-B007}{start}
\begin{defn}
  که \LR{$\mModel{F}$} يو چوکاټ وي، وايو چې \LR{$!A$}
  \emph{په~\LR{$\mModel{F}$} کښې معتبر} دے، \LR{$\mModel{F} \Entails !A$}،
  هله او يوازې هله چې \LR{$\mSat{M}{!A}$} په~\LR{$\mModel{F}$} ولاړ د هر مدل~\LR{$\mModel{M}$}
  دپاره وي۔
\olpseganchor{OLP-0422-B007}{end}
  
\olpseganchor{OLP-0422-B008}{start}
  که \LR{$\mClass{F}$} د چوکاټونو يو ټولګے وي، وايو چې \LR{$!A$}
  \emph{په~\LR{$\mClass{F}$} کښې معتبر} دے، \LR{$\mClass{F} \Entails !A$}،
  هله او يوازې هله چې د هر چوکاټ~\LR{$\mModel{F} \in \mClass{F}$}
  دپاره \LR{$\mModel{F} \Entails !A$} وي۔
\end{defn}
\olpseganchor{OLP-0422-B008}{end}

\olpseganchor{OLP-0422-B009}{start}
چوکاټونه ځکه مهم دي چې د طرحو او د لاسرسي د اړيکې~\LR{$R$} د خاصيتونو
ترمنځ مطابقت د چوکاټونو په کچه وي، \emph{د مدلونو په کچه نۀ}۔ د
بېلګې په توګه، که څه هم \Ax{T} په ټولو انعکاسي مدلونو کښې رښتيا ده،
خو هر هغه مدل انعکاسي نۀ دے چې \Ax{T} پکې رښتيا وي۔ البته دا
\emph{سمه} ده چې \Ax{T} په ټولو انعکاسي \emph{چوکاټونو} کښې
\emph{معتبره} ده، او هر هغه چوکاټ هم انعکاسي دے چې \Ax{T} پکې
معتبره وي۔
\olpseganchor{OLP-0422-B009}{end}

\olpseganchor{OLP-0422-B010}{start}
\begin{rem}
د چوکاټونو په يوه ټولګي کښې اعتبار د مدلونو په يوه ټولګي کښې د
اعتبار د مفهوم يو ځانګړے مورد دے: \LR{$\mClass{F} \Entails !A$} هله او
يوازې هله چې \LR{$\mClass{C} \Entails !A$}، چېرې چې \LR{$\mClass{C}$} د هغو
ټولو مدلونو ټولګے دے چې د~\LR{$\mClass{F}$} د کوم چوکاټ پر بنسټ ولاړ وي۔
\olpseganchor{OLP-0422-B010}{end}

\olpseganchor{OLP-0422-B011}{start}
څرګنده ده چې که فارمول يا يوه طرحه معتبره وي، يعنې د
\emph{ټولو} مدلونو د ټولګي په نسبت معتبره وي، نو د چوکاټونو د هر
ټولګي~\LR{$\mClass{F}$} په نسبت هم معتبره ده۔
\end{rem}
\olpseganchor{OLP-0422-B011}{end}

\olpseganchor{OLP-0423-B004}{start}
\olfileid{nml}{frd}{def}
\olpseganchor{OLP-0423-B004}{end}

\olpseganchor{OLP-0423-B005}{start}
\olsection{د چوکاټ تعريف‌وړتيا}
\olpseganchor{OLP-0423-B005}{end}

\olpseganchor{OLP-0423-B006}{start}
که څه هم د \olref[acc]{thm:soundschemas} سرچپه استلزامونه د مدلونو
دپاره سم نۀ دي، خو د «مدل» پر ځاے د «چوکاټ» په راوړلو سمېږي: د
\olref[acc]{thm:soundschemas} د خاصيتونو په اړه دا خبره \emph{سمه}
ده چې که فارمول په يوه \emph{چوکاټ} کښې معتبر وي، نو د هغه
چوکاټ د لاسرسي اړيکه اړوند خاصيت لري۔ نو څېړل شوي فارمولونه
د هغو چوکاټونو ټولګي \emph{تعريفوي} چې اړوند خاصيت لري۔
\olpseganchor{OLP-0423-B006}{end}

\olpseganchor{OLP-0423-B007}{start}
\begin{defn}
  که \LR{$\mClass{F}$} د چوکاټونو يو ټولګے وي، وايو چې \LR{$!A$}
  دا ټولګے~\LR{$\mClass{F}$} \emph{تعريفوي}، هله او يوازې هله چې
  \LR{$\mModel{F} \Entails !A$} د ټولو او يوازې هغو
  چوکاټونو~\LR{$\mModel{F} \in \mClass{F}$} دپاره وي۔
\end{defn}
\olpseganchor{OLP-0423-B007}{end}

\olpseganchor{OLP-0423-B008}{start}
اوس د چوکاټونو د تعريف‌وړتيا بشپړې پايلې ثابتوو۔
\olpseganchor{OLP-0423-B008}{end}

\olpseganchor{OLP-0423-B009}{start}
\begin{thm}\ollabel{thm:fullCorrespondence}
که د \olref[acc]{tab:five} د ښي خوا فارمول په يوه
چوکاټ~\LR{$\mModel{F}$} کښې معتبر وي، نو \LR{$\mModel{F}$} د چپې خوا
خاصيت لري۔
\end{thm}
\olpseganchor{OLP-0423-B009}{end}

\olpseganchor{OLP-0423-B010}{start}
\begin{proof}
  \begin{enumerate}
  \item فرض کړئ \Ax{D} په \LR{$\mModel{F} = \tuple{W, R}$} کښې معتبره
    وي؛ يعنې \LR{$\mModel{F} \Entails \Box p \lif \Diamond p$}۔
    \LR{$\mModel{M} = \tuple{W, R, V}$} پر~\LR{$\mModel{F}$} ولاړ يو مدل
    او \LR{$w \in W$} ونيسئ۔ بايد وښيو چې داسې \LR{$v$} شته چې~\LR{$Rwv$} وي۔
    فرض کړئ داسې نۀ وي: بيا د هر~\LR{$!A$} دپاره، او په ځانګړي ډول د~\LR{$p$}
    دپاره، \LR{$\mSat{M}{\Box !A}[w]$} او
    \LR{$\mSat/{M}{\Diamond !A}[w]$} دواړه وي۔ خو بيا
    \LR{$\mSat/{M}{\Box p \lif \Diamond p}[w]$}، چې له دې فرض سره
    په ټکر کښې دے چې \LR{$\mModel{F} \Entails \Box p \lif \Diamond p$}۔
  \item فرض کړئ \Ax{T} په \LR{$\mModel{F}$} کښې معتبره وي؛ يعنې
    \LR{$\mModel{F} \Entails \Box p \lif p$}۔ \LR{$w \in W$} يوه خپلسره
    نړۍ ونيسئ؛ بايد \LR{$Rww$} وښيو۔ \LR{$u \in V(p)$} هله او يوازې هله
    وټاکئ چې \LR{$Rwu$} وي (که \LR{$q$} له~\LR{$p$} څخه بېل وي، \LR{$V(q)$} خپلسرے
    دے؛ د بېلګې په توګه \LR{$V(q) = \emptyset)$}۔ فرض کړئ
    \LR{$\mModel{M} = \tuple{W, R, V}$}۔ د ټاکنې له مخې، هر هغه \LR{$u$} چې
    \LR{$Rwu$} وي \LR{$\mSat{M}{p}[u]$} پوره کوي؛ نو
    \LR{$\mSat{M}{\Box p}[w]$}۔ خو د فرض له مخې \LR{$\Box p \lif p$}
    په~\LR{$w$} کښې رښتيا دے؛ نو \LR{$\mSat{M}{p}[w]$}، او دا د~\LR{$V$} د
    تعريف له مخې يوازې هله شونې ده چې~\LR{$Rww$} وي۔
  \item د دعوې مقابل نقيض ثابتوو: فرض کړئ \LR{$\mModel{F}$} متناظره نۀ وي؛
    ښيو چې \Ax{B}، يعنې \LR{$p \lif \Box\Diamond p$}، په
    \LR{$\mModel{F}= \tuple{W, R}$} کښې معتبره نۀ ده۔ که
    \LR{$\mModel{F}$} متناظره نۀ وي، داسې \LR{$u$}، \LR{$v \in W$} شته چې
    \LR{$Ruv$} وي خو \LR{$Rvu$} نۀ وي۔ \LR{$V$} داسې تعريف کړئ چې
    \LR{$w \in V(p)$} هله او يوازې هله چې \LR{$Rvw$} نۀ وي (او \LR{$V$} د نورو بياني
    متغيرونو لپاره هر ډول ټاکل کېدے شي)۔ فرض کړئ \LR{$\mModel{M} =\tuple{W, R,
      V}$}۔ اوس د~\LR{$V$} له تعريفه د هر هغه \LR{$w$} دپاره چې \LR{$Rvw$} نۀ
    وي \LR{$\mSat{M}{p}[w]$}، په ځانګړي ډول \LR{$\mSat{M}{p}[u]$}، ځکه
    چې \LR{$Rvu$} نۀ وي۔ همدارنګه، ځکه چې \LR{$Rvw$} هله او يوازې هله چې
    \LR{$w \notin V(p)$}، داسې هېڅ~\LR{$w$} نشته چې \LR{$Rvw$} او
    \LR{$\mSat{M}{p}[w]$} دواړه وي؛ نو
    \LR{$\mSat/{M}{\Diamond p}[v]$}۔ ځکه چې~\LR{$Ruv$}، نو
    \LR{$\mSat/{M}{\Box\Diamond p}[u]$} هم۔ له دې څخه
    \LR{$\mSat/{M}{p \lif \Box\Diamond p}[u]$} راځي؛ نو \LR{$\Ax{B}$}
    په~\LR{$\mModel{F}$} کښې معتبره نۀ ده۔
  \item فرض کړئ \Ax{4} په \LR{$\mModel{F} = \tuple{W,R}$} کښې معتبره
    وي؛ يعنې \LR{$\mModel{F} \Entails \Box p \lif \Box\Box p$}۔
    خپلسرې نړۍوې \LR{$u$}، \LR{$v$}، \LR{$w \in W$} داسې ونيسئ چې \LR{$Ruv$} او
    \LR{$Rvw$} وي؛ بايد \LR{$Ruw$} وښيو۔ \LR{$V$} داسې تعريف کړئ چې
    \LR{$z \in V(p)$} هله او يوازې هله چې \LR{$Ruz$} وي (او \LR{$V$} د نورو بياني
    متغيرونو لپاره هر ډول ټاکل کېدے شي)۔ فرض کړئ \LR{$\mModel{M} =\tuple{W, R, V}$}۔
    د~\LR{$V$} د تعريف له مخې، د هر هغه \LR{$z$} دپاره چې \LR{$Ruz$} وي
    \LR{$\mSat{M}{p}[z]$}، نو \LR{$\mSat{M}{\Box p}[u]$}۔ خو د فرض له
    مخې \Ax{4}، \LR{$\Box p \lif \Box \Box p$}، په~\LR{$u$} کښې رښتيا
    ده؛ نو \LR{$\mSat{M}{\Box \Box p}[u]$}۔ ځکه چې \LR{$Ruv$} او \LR{$Rvw$}،
    نو \LR{$\mSat{M}{p}[w]$}؛ خو د~\LR{$V$} د تعريف له مخې دا يوازې
    هله شونې ده چې \LR{$Ruw$} وي، لکه څنګه چې مو غوښتل۔
  \item بيا د دعوې مقابل نقيض اخلو: فرض کړئ چوکاټ
    \LR{$\mModel{F} = \tuple{W, R}$} اقليدسي نۀ وي، او ښيو چې
    \Ax{5} پکې نارښتيا کېږي؛ يعنې
    \LR{$\mModel{F} \Entails/ \Diamond p \lif \Box\Diamond p$}۔
    فرض کړئ داسې نړۍوې \LR{$u$}، \LR{$v$}، \LR{$w \in W$} شته چې \LR{$Rwu$} او
    \LR{$Rwv$} وي، خو \LR{$Ruv$} نۀ وي۔ \LR{$V$} داسې تعريف کړئ چې د هرې نړۍ
    \LR{$z$} دپاره \LR{$z \in V(p)$} هله او يوازې هله چې \LR{$Ruz$}
    \emph{نۀ} وي۔ نورو بياني متغيرونو ته ارزښتونه خپلسري وټاکئ۔
    فرض کړئ \LR{$\mModel{M} =\tuple{W, R, V}$}۔ نو د
    فرض له مخې \LR{$\mSat{M}{p}[v]$}، او ځکه چې \LR{$Rwv$}، نو
    \LR{$\mSat{M}{\Diamond p}[w]$} هم۔ خو داسې نړۍ \LR{$y$} نشته چې
    \LR{$Ruy$} او \LR{$\mSat{M}{p}[y]$} دواړه پکې وي؛ نو
    \LR{$\mSat/{M}{\Diamond p}[u]$}۔ ځکه چې \LR{$Rwu$}، له دې څخه
    \LR{$\mSat/{M}{\Box\Diamond p}[w]$} راځي؛ نو \Ax{5}،
    \LR{$\Diamond p \lif \Box\Diamond p$}، په~\LR{$w$} کښې نارښتيا ده۔
  \end{enumerate}
\end{proof}
\olpseganchor{OLP-0423-B010}{end}

\olpseganchor{OLP-0423-B011}{start}
د \Ax{D} دليل د ارزښت ټاکنې~\LR{$V$} په هېڅ ځانګړي شرط ولاړ نۀ دے۔ په حقيقت کښې مو وښودله چې که
\LR{$\mSat{M}{\Ax{D}}$} وي، نو \LR{$\mModel{M}$} مسلسل دے۔ نو \Ax{D}
يوازې د چوکاټونو نۀ، بلکې د مسلسل \emph{مدلونو} ټولګے تعريفوي۔
\olpseganchor{OLP-0423-B011}{end}

\olpseganchor{OLP-0423-B012}{start}
\begin{cor}\ollabel{prop:D-serial}
  هر هغه مدل مسلسل دے چې \Ax{D} پکې رښتيا وي۔
\end{cor}
\olpseganchor{OLP-0423-B012}{end}

\olpseganchor{OLP-0423-B013}{start}
\begin{cor}
د \olref[acc]{tab:five} د ښي خوا هر فارمول د هغو چوکاټونو
ټولګے تعريفوي چې د چپې خوا خاصيت لري۔
\end{cor}
\olpseganchor{OLP-0423-B013}{end}

\olpseganchor{OLP-0423-B014}{start}
\begin{proof}
  په \olref[acc]{thm:soundschemas} کښې مو ثابته کړه چې که يو مدل
  د چپې خوا خاصيت ولري، د ښي خوا فارمول پکې رښتيا وي۔ نو که
  يو چوکاټ~\LR{$\mModel{F}$} د چپې خوا خاصيت ولري، د ښي خوا
  فارمول په~\LR{$\mModel{F}$} کښې معتبر دے۔ په
  \olref{thm:fullCorrespondence} کښې مو سرچپه استلزامونه ثابت کړل:
  که د ښي خوا فارمول په~\LR{$\mModel{F}$} کښې معتبر وي، نو
  \LR{$\mModel{F}$} د چپې خوا خاصيت لري۔
\end{proof}
\olpseganchor{OLP-0423-B014}{end}

\olpseganchor{OLP-0423-B015}{start}
\begin{prob}
وښيئ چې که د \olref[nml][frd][acc]{tab:anotherfive} د ښي خوا
فارمول په يوه چوکاټ~\LR{$\mModel{F}$} کښې معتبر وي، نو
\LR{$\mModel{F}$} د چپې خوا خاصيت لري۔ د دې دپاره داسې چوکاټ واخلئ
چې د چپې خوا خاصيت \emph{نۀ} پوره کوي، او مناسب~\LR{$V$} داسې
وټاکئ چې د ښي خوا فارمول په کومه نړۍ کښې نارښتيا وي۔
\end{prob}
\olpseganchor{OLP-0423-B015}{end}

\olpseganchor{OLP-0423-B016}{start}
\olref{thm:fullCorrespondence} دا هم ښيي چې خاصيتونه يوځاے کېدلے
شي: د بېلګې په توګه، که \Ax{B} او \Ax{4} دواړه په~\LR{$\mModel{F}$}
کښې معتبرې وي، نو چوکاټ هم متناظره او هم انتقالي دے، او داسې نور۔
ډېر مهم مودالي منطقونه د هغو ټولو چوکاټونو د معتبرو فارمولونه د
مجموعې په توګه پېژندل کېږي چې څو چوکاټي خاصيتونه يوځاے لري؛ نو دوى
د هغو ټولو چوکاټونو د معتبرو فارمولونه په توګه هم پېژندل کېدلے
شي چې اړوند تعريفوونکي فارمولونه پکې معتبر وي۔ د بېلګې په توګه،
دوديز نظام \Log{S4} د هغو ټولو فارمولونه مجموعه ده چې په ټولو
انعکاسي او انتقالي چوکاټونو کښې معتبر وي؛ يعنې په هغو ټولو کښې چې
\Ax{T} او~\Ax{4} دواړه معتبرې وي۔ \Log{S5} د هغو ټولو فارمولونو
مجموعه ده چې په ټولو انعکاسي، متناظرو او اقليدسي چوکاټونو کښې معتبر
وي؛ يعنې په هغو ټولو کښې چې \Ax{T}، \Ax{B} او \Ax{5} ټولې
معتبرې وي۔
\olpseganchor{OLP-0423-B016}{end}

\olpseganchor{OLP-0423-B017}{start}
د~\LR{$R$} د خاصيتونو ترمنځ منطقي اړيکې عموماً د اړوندو تعريفوونکو
فارمولونه ترمنځ اړيکو سره سمون لري۔ د بېلګې په توګه، هره انعکاسي
اړيکه مسلسل هم ده؛ نو چې کله \Ax{T} په يوه چوکاټ کښې معتبره وي،
\Ax{D} هم پکې معتبره وي۔ (پام وکړئ چې دا اړيکه د استلزام اړيکه
\emph{نۀ} ده۔ دا خبره سمه نۀ ده چې هر کله
\LR{$\mSat{M}{\Ax{T}}[w]$} وي، نو \LR{$\mSat{M}{\Ax{D}}[w]$} هم وي۔)
څو داسې اړيکې دلته ثبتوو۔
\olpseganchor{OLP-0423-B017}{end}

\olpseganchor{OLP-0423-B018}{start}
\begin{prop}\ollabel{prop:relation-facts}
  فرض کړئ \LR{$R$} پر يوه سټ \LR{$W$} دوه‌ځاييزه اړيکه ده؛ بيا:
  \begin{enumerate}
  \item که \LR{$R$} انعکاسي وي، نو مسلسل هم ده۔
  \item که \LR{$R$} متناظره وي، نو هله او يوازې هله انتقالي ده چې
    اقليدسي وي۔
  \item که \LR{$R$} متناظره يا اقليدسي وي، نو کمزورې هم‌لوري ده
    («الماسي خاصيت» لري)۔
  \item که \LR{$R$} اقليدسي وي، نو کمزورې تړلې ده۔
  \item که \LR{$R$} تابعي وي، نو مسلسل ده۔  
  \end{enumerate}
\end{prop}
\olpseganchor{OLP-0423-B018}{end}

\olpseganchor{OLP-0423-B019}{start}
\begin{prob}
  \olref[nml][frd][def]{prop:relation-facts} ثابته کړئ۔
\end{prob}
\olpseganchor{OLP-0423-B019}{end}

\olpseganchor{OLP-0424-B004}{start}
\olfileid{nml}{frd}{fol}
\olpseganchor{OLP-0424-B004}{end}

\olpseganchor{OLP-0424-B005}{start}
\olsection{د لومړۍ درجې تعريف‌وړتيا}
\olpseganchor{OLP-0424-B005}{end}

\olpseganchor{OLP-0424-B006}{start}
مونږ وليدل چې د چوکاټونو د لاسرسي اړيکو يو شمېر خاصيتونه د مودالي
فارمولونه په واسطه تعريفېدے شي۔ د بېلګې په توګه، د چوکاټونو
تناظر د~\Ax{B} په فارمول، \LR{$p \lif \Box\Diamond
p$}، تعريفېدے شي۔ تر اوسه چې کوم شرطونه مو ليدلي، ټول د لومړۍ درجې
په هغو فارمولونه بيانېدے شي چې په داسې ژبه کښې وي
چې يوازې يو دوه‌ځاييز پريديکات لري۔ د بېلګې په توګه، تناظر
د~\LR{$\lforall[x][\lforall[y][(\Atom{Q}{x,y} \lif \Atom{Q}{y, x})]]$}
په واسطه تعريفېږي: د لومړۍ درجې هغه جوړښت~\LR{$\Struct{M}$}
چې \LR{$\Domain{M} = W$} او \LR{$\Assign{Q}{M} = R$} وي، پورته فارمول
هله او يوازې هله پوره کوي چې \LR{$R$} متناظره وي۔ دا لاندې تعريف ته
لار هواروي:
\olpseganchor{OLP-0424-B006}{end}

\olpseganchor{OLP-0424-B007}{start}
\begin{defn}
  د چوکاټونو يو ټولګے~\LR{$\mClass{F}$} \emph{د لومړۍ درجې له مخې
  تعريف‌وړ} دے، که د لومړۍ درجې په داسې ژبه کښې چې يوازې يو
  دوه‌ځاييز پريديکات~\LR{$Q$} لري، يوه تړلے فارمول~\LR{$!A$} وي،
  داسې چې \LR{$\mModel{F} =
  \tuple{W, R} \in \mClass{F}$} هله او يوازې هله چې د لومړۍ درجې
  په هغه جوړښت~\LR{$\Struct{M}$} کښې \LR{$\Sat{M}{!A}$} وي چې
  \LR{$\Domain{M} = W$} او \LR{$\Assign{Q}{M}
  = R$} لري۔
\end{defn}
\olpseganchor{OLP-0424-B007}{end}

\olpseganchor{OLP-0424-B008}{start}
تر اوسه څېړل شوي خاصيتونه او هغه مودالي فارمولونه چې تعريفوي
ئې، استثنايي دي۔ داسې مودالي فارمول شته چې د چوکاټونو
تعريف کړے ټولګے ئې د لومړۍ درجې له مخې تعريف‌وړ نۀ دے۔ هم داسې
ټولګے شته چې د لومړۍ درجې له مخې تعريف‌وړ دے، خو په هېڅ مودالي
فارمول نۀ تعريفېږي۔
\olpseganchor{OLP-0424-B008}{end}

\olpseganchor{OLP-0424-B009}{start}
د لومړۍ دعوې مقابله بېلګه د لوب فارمول دے:
\begin{equation}
\Box(\Box p \lif p) \lif \Box p. \tag{\Ax{W}}
\end{equation}
\Ax{W} د هغو چوکاټونو ټولګے تعريفوي چې اړيکه ئې انتقالي او
سرچپه ښه‌بنسټه وي۔ يوه اړيکه هله ښه‌بنسټه ده چې داسې هېڅ
نامتناهي لړۍ \LR{$w_1$}، \LR{$w_2$}، \dots{} ونۀ لري چې \LR{$Rw_2w_1$}،
\LR{$Rw_3w_2$}، \dots وي۔ د بېلګې په توګه، پر~\LR{$\Nat$} د \LR{$<$} اړيکه
ښه‌بنسټه ده، خو پر~\LR{$\Int$} د \LR{$<$} اړيکه داسې نۀ ده۔ يوه اړيکه هله
او يوازې هله سرچپه ښه‌بنسټه ده چې د هغې سرچپه اړيکه
ښه‌بنسټه وي۔ نو سرچپه ښه‌بنسټې اړيکې هغه دي چې هېڅ داسې نامتناهي
لړۍ \LR{$w_1$}، \LR{$w_2$}، \dots{} ونۀ لري چې \LR{$Rw_1w_2$}،
\LR{$Rw_2w_3$}، \dots وي۔
\olpseganchor{OLP-0424-B009}{end}

\olpseganchor{OLP-0424-B010}{start}
خو د لومړۍ درجې هېڅ فارمول انتقالي او سرچپه ښه‌بنسټې
اړيکې نۀ تعريفوي۔ د دې د ليدلو دپاره فرض کړئ چې
\LR{$\Sat{M}{!F}$} هله او يوازې هله وي چې \LR{$R =
\Assign{Q}{M}$} انتقالي او سرچپه ښه‌بنسټه وي۔ \LR{$!A_n$} دا
فارمول ونيسئ
\[
(\Atom{Q}{a_1,a_2} \land \dots \land
    \Atom{Q}{a_{n-1},a_{n}})
\]
دلته د شمېر لرونکي ثابت سمبولونه د لومړۍ درجې د ژبې په موقتي
پراختيا کښې کاروو؛ د لومړي شاخص په مورد کښې عطف تش او ځکه رښتيا
ګڼل کېږي۔
اوس د فارمولونه دا سټ وڅېړئ
\[
\Gamma = \{!F, !A_1, !A_2, \dots\}.
\]
د~\LR{$\Gamma$} هر متناهي فرعي سټ د صدق وړ دے: \LR{$k$} هغه تر ټولو لوے
شاخص ونيسئ چې \LR{$!A_k$} په فرعي سټ کښې وي؛ که د لړۍ هېڅ غړے پکې
نۀ وي، لومړے شاخص ونيسئ۔ بيا \LR{$\Domain{M_k} = \{1, \dots, k\}$}،
\LR{$\Assign{a_i}{M_k} = i$} د همدغه حد پورې، او
\LR{$\Assign{Q}{M_k} = <$} وټاکئ۔ له دې حد څخه د پورته ثابتونو ارزښتونه
په دامنه کښې خپلسري وټاکئ۔ ځکه چې پر \LR{$\{1,
\dots, k\}$} د \LR{$<$} اړيکه انتقالي او سرچپه ښه‌بنسټه ده،
\LR{$\Sat{M_k}{!F}$}۔ د جوړښت له مخې، د هر \LR{$i \le k$} دپاره
\LR{$\Sat{M_k}{!A_i}$}۔
د لومړۍ درجې منطق د متناهي شاهد قضيې له مخې، \LR{$\Gamma$} په
کوم جوړښت~\LR{$\Struct{M}$} کښې د صدق وړ دے۔ د فرض له مخې،
ځکه چې \LR{$\Sat{M}{!F}$}، نو \LR{$\Assign{Q}{M}$} سرچپه
ښه‌بنسټه ده۔ خو څرګنده ده چې \LR{$\Assign{a_1}{M}$}،
\LR{$\Assign{a_2}{M}$}، \dots{} به هغه نامتناهي لړۍ جوړه کړي
چې سرچپه ښه‌بنسټوالی ئې ردوي۔
\olpseganchor{OLP-0424-B010}{end}

\olpseganchor{OLP-0424-B011}{start}
د دويمې دعوې مقابله بېلګه د کليت خاصيت دے: د هر \LR{$u$} او \LR{$v$}
دپاره \LR{$Ruv$}۔ کلي چوکاټونه د لومړۍ درجې د
\LR{$\lforall[x][\lforall[y][\Atom{Q}{x,y}]]$} فارمول په واسطه
تعريفېدے شي۔ خو داسې هېڅ مودالي فارمول نشته چې په ټولو او يوازې
کلي چوکاټونو کښې معتبر وي۔ دا د يوې په خپل ذات کښې په زړه پورې
پايلې نتيجه ده: په کلي چوکاټونو کښې معتبر فارمولونه عين هغه
فارمولونه دي چې په انعکاسي، متناظرو او انتقالي چوکاټونو کښې
معتبر دي۔ ځينې انعکاسي، متناظره او انتقالي چوکاټونه کلي
نۀ دي؛ نو هر فارمول چې په ټولو کلي چوکاټونو کښې معتبر وي،
په ځينو ناکلي چوکاټونو کښې هم معتبر وي۔
\olpseganchor{OLP-0424-B011}{end}

\olpseganchor{OLP-0425-B004}{start}
\olfileid{nml}{frd}{es5}
\olpseganchor{OLP-0425-B004}{end}

\olpseganchor{OLP-0425-B005}{start}
\olsection{د هم ارزښتۍ اړيکې او \texorpdfstring{\Log{S5}}{S5}}
\olpseganchor{OLP-0425-B005}{end}

\olpseganchor{OLP-0425-B006}{start}
مودالي منطق \Log{S5} د هغو فارمولونه د مجموعې په توګه پېژندل
کېږي چې په ټولو کلي چوکاټونو کښې معتبر وي؛ يعنې هرې نړۍ ته له هرې
نړۍ، د هغې له خپل ځان څخه هم، لاسرسے وي۔ په داسې حالت کښې
\LR{$\Box$} له لزوم او \LR{$\Diamond$} له امکان سره سمون لري: \LR{$\Box !A$}
هله رښتيا دے چې \LR{$!A$} په \emph{هرې} نړۍ کښې رښتيا وي، او
\LR{$\Diamond !A$} هله رښتيا دے چې \LR{$!A$} په \emph{کومې} نړۍ کښې
رښتيا وي۔ دا هم څرګندېږي چې \Log{S5} د هغو فارمولونه په
توګه هم پېژندل کېدلے شي چې په ټولو انعکاسي، متناظرو او انتقالي
چوکاټونو کښې معتبر وي؛ يعنې په هغو چوکاټونو کښې چې د لاسرسي
اړيکه ئې د \emph{هم ارزښتۍ اړيکه} وي۔
\olpseganchor{OLP-0425-B006}{end}

\olpseganchor{OLP-0425-B007}{start}
\begin{defn}
  پر \LR{$W$} دوه‌ځاييزه اړيکه \LR{$R$} هله او يوازې هله د
  \emph{هم ارزښتۍ اړيکه} ده چې انعکاسي، متناظره او انتقالي وي۔
  پر \LR{$W$} اړيکه \LR{$R$} هله او يوازې هله \emph{کلي} ده چې د ټولو
  \LR{$u,v \in W$} دپاره \LR{$Ruv$} وي۔
\end{defn}
\olpseganchor{OLP-0425-B007}{end}

\olpseganchor{OLP-0425-B008}{start}
ځکه چې \Ax{T}، \Ax{B} او \Ax{4} په ترتيب سره انعکاسي،
متناظره او انتقالي چوکاټونه پېژني، نو هغه چوکاټونه چې د لاسرسي
اړيکه ئې د هم ارزښتۍ اړيکه وي، کټ مټ هغه دي چې دا درې واړه
فارمولونه پکې معتبر وي۔ د هم ارزښتۍ اړيکې د فارمولونه
په نورو ترکيبونو هم پېژندل کېدلے شي، ځکه هغه شرطونه چې پرې مو
د هم ارزښتۍ اړيکې تعريف کړې دي، پر~\LR{$R$} د نورو پېژندل شوو
شرطونو له ترکيبونو سره هم‌ارز دي۔
\olpseganchor{OLP-0425-B008}{end}

\olpseganchor{OLP-0425-B009}{start}
\begin{prop}\ollabel{prop:equivalences}
  لاندې شرطونه له يو بل سره هم‌ارز دي:
  \begin{enumerate}
  \item \LR{$R$} د هم ارزښتۍ اړيکه ده؛
  \item \LR{$R$} انعکاسي او اقليدسي ده؛
  \item \LR{$R$} مسلسل، متناظره او اقليدسي ده؛
  \item \LR{$R$} مسلسل، متناظره او انتقالي ده۔
  \end{enumerate}
\end{prop}
\olpseganchor{OLP-0425-B009}{end}

\olpseganchor{OLP-0425-B010}{start}
\begin{proof}
  ثبوت په لاندې تمرين کښې غوښتل شوے دے۔
\end{proof}
\olpseganchor{OLP-0425-B010}{end}

\olpseganchor{OLP-0425-B011}{start}
\begin{prob}
  د لاندې خبرو په ښودلو \olref[nml][frd][es5]{prop:equivalences}
  ثابته کړئ:
  \begin{enumerate}
  \item که \LR{$R$} متناظره او انتقالي وي، نو اقليدسي ده۔
  \item که \LR{$R$} انعکاسي وي، نو مسلسل ده۔
  \item که \LR{$R$} انعکاسي او اقليدسي وي، نو متناظره ده۔
  \item که \LR{$R$} متناظره او اقليدسي وي، نو انتقالي ده۔
  \item که \LR{$R$} مسلسل، متناظره او انتقالي وي، نو انعکاسي ده۔
  \end{enumerate}
  څرګنده کړئ چې دا خبرې ولې د دغو شرطونو د هم‌ارزۍ د ثبوت
  دپاره بسنه کوي۔
\end{prob}
\olpseganchor{OLP-0425-B011}{end}

\olpseganchor{OLP-0425-B012}{start}
\olref{prop:equivalences} د \olref[prf][prs]{prop:S5} معنايي
سياله ده: د هغو چوکاټونو مودالي منطق په هم‌ارزه بڼه پېژني چې
په هغو کښې~\LR{$R$} د هم ارزښتۍ اړيکه وي؛ دې منطق ته په دوديزه توګه
\Log{S5} وايي۔
\olpseganchor{OLP-0425-B012}{end}

\olpseganchor{OLP-0425-B013}{start}
د کليو او د هم ارزښتۍ اړيکو ترمنځ څۀ اړيکه ده؟ که څۀ هم هره
کلي اړيکه د هم ارزښتۍ اړيکه ده، خو څرګنده ده چې د هم ارزښتۍ
هره اړيکه کلي نۀ ده۔ بيا هم هغه فارمولونه چې په ټولو
کلي اړيکو کښې معتبر وي، کټ مټ هماغه دي چې په ټولو د هم
ارزښتۍ اړيکو کښې معتبر وي۔
\olpseganchor{OLP-0425-B013}{end}

\olpseganchor{OLP-0425-B014}{start}
\begin{prop}
  فرض کړئ \LR{$R$} د هم ارزښتۍ اړيکه ده، او د هر \LR{$w \in W$} دپاره
  د~\LR{$w$} \emph{د هم ارزښتۍ ټولګے} د \LR{$[w] = \{w'\in W :
  Rww'\}$} په توګه تعريف کړئ۔ بيا:
  \begin{enumerate}
  \item \LR{$w \in [w]$}؛
  \item \LR{$R$} د هم ارزښتۍ پر هر ټولګي \LR{$[w]$} کلي ده؛
  \item د هم ارزښتۍ د ټولګيو مجموعه~\LR{$W$} په داسې فرعي سټونو
    وېشي چې يو له بل څخه بېل دي او په ګډه ټول سټ رانغاړي۔
  \end{enumerate}
\end{prop}
\olpseganchor{OLP-0425-B014}{end}

\olpseganchor{OLP-0425-B015}{start}
\begin{prop}\ollabel{prop:S5=univ}
  فارمول~\LR{$!A$} په ټولو هغو چوکاټونو
  \LR{$\mModel{F} =\tuple{W,R}$} کښې معتبر دے چې~\LR{$R$} پکې د هم
  ارزښتۍ اړيکه وي، هله او يوازې هله چې په ټولو هغو چوکاټونو
  \LR{$\mModel{F} =\tuple{W,R}$} کښې معتبر وي چې~\LR{$R$} پکې کلي
  وي۔ نو د کلي چوکاټونو منطق هماغه \Log{S5} دے۔
\end{prop}
\olpseganchor{OLP-0425-B015}{end}

\olpseganchor{OLP-0425-B016}{start}
\begin{proof}
  نېغ په نېغه ښکاري چې پر~\LR{$W$} هره کلي اړيکه~\LR{$R$} د هم
  ارزښتۍ اړيکه ده۔ نو که~\LR{$!A$} په ټولو هغو چوکاټونو کښې
  معتبر وي چې~\LR{$R$} پکې د هم ارزښتۍ اړيکه وي، په ټولو کلي
  چوکاټونو کښې هم معتبر دے۔ د بل لوري دپاره د دعوې مقابل
  نقيض ثابتوو: فرض کړئ \LR{$!B$} داسې فارمول وي چې په
  يوه نړۍ~\LR{$w$} کښې، په داسې مدل
  \LR{$\mModel{M} = \tuple{W, R, V}$} کښې نارښتيا وي چې پر
  چوکاټ~\LR{$\tuple{W,R}$} ولاړ دے او~\LR{$R$} پر~\LR{$W$} د هم ارزښتۍ
  اړيکه ده۔ نو \LR{$\mSat/{M}{!B}[w]$}۔ يو مدل
  \LR{$\mModel{M}' = \tuple{W',
    R', V'}$} په دې ډول تعريف کړئ:
  \begin{enumerate}
  \item \LR{$W' = [w]$}؛
  \item \LR{$R'$} پر \LR{$W'$} کلي ده؛
  \item \LR{$V'(p) = V(p) \cap W'$}۔
  \end{enumerate}
  (نو په~\LR{$\mModel{M}'$} کښې د نړۍو سټ~\LR{$W'$} د
  \olref{fig:partition} په خړ رنګ شوې برخه کښې ښودل شوے دے۔)
  په اسانه ښکاري چې \LR{$R$} او \LR{$R'$} پر~\LR{$W'$} يو شان دي۔ د انتقال
  له مخې، د دې ټولګي له هرې نړۍ څخه د لاسرسي وړ هره نړۍ هم
  په همدغه ټولګي کښې ده۔ بيا د فارمولونه پر جوړښت په استقرا ښودلے شو چې د هر
  \LR{$w' \in W'$} دپاره: د هر \LR{$!A$} دپاره
  \LR{$\mSat{M'}{!A}[w']$} هله او يوازې هله چې
  \LR{$\mSat{M}{!A}[w']$} وي (دا معنا لري، ځکه چې \LR{$W'
  \subseteq W$})۔ په ځانګړي ډول، \LR{$\mSat/{M'}{!B}[w]$}،
  او~\LR{$!B$} په يوه داسې مدل کښې نارښتيا دے چې پر کلي
  چوکاټ ولاړ دے۔
\end{proof}
\olpseganchor{OLP-0425-B016}{end}

\olpseganchor{OLP-0425-B017}{start}
\begin{figure}[t]
  \centering
  \begin{tikzpicture}[node distance=2cm, auto, thick]
    \clip [rounded corners] (0,0) -- (6,0) -- (6,4) -- (0,4) --  cycle;
    \begin{scope}
      \clip (0,0) -- (2,0)  .. controls (1.5,1.5) and (4.5,1.5)
      .. (4,4) -- (0,4) -- (0,0) -- cycle;
      \filldraw[fill=gray!40] (0,4) -- (0,2) .. controls (1.5,1.5)
      and (4.5,1.5) .. (4.5,0) -- (6,0) -- (6,4) -- (0,4) --  cycle;
    \end{scope}
    \draw [name path=line1] (2,0) .. controls (1.5,1.5) and (4.5,1.5) .. (4,4);
    \draw [name path=line2] (0,2)  .. controls (1.5,1.5) and (4.5,1.5) .. (4.5,0);
    \path [name intersections={of=line1 and line2}];
    \draw [rounded corners, use as bounding box, very thick] (0,0)
    -- (6,0) -- (6,4) -- (0,4) --  cycle; 
    \node at (2,3) {\LR{$[w]$}} ;
    \node at (1,1) {\LR{$[u]$}} ;
    \node at (3,0.75) {\LR{$[v]$}} ;
    \node at (4.75,2) {\LR{$[z]$}} ;
  \end{tikzpicture}
  \caption{د \LR{$W$} وېش د هم ارزښتۍ پر ټولګيو۔}
  \ollabel{fig:partition}
\end{figure}
\olpseganchor{OLP-0425-B017}{end}

\olpseganchor{OLP-0426-B004}{start}
\olfileid{nml}{frd}{st}
\olpseganchor{OLP-0426-B004}{end}

\olpseganchor{OLP-0426-B005}{start}
\olsection{د دويمې درجې تعريف‌وړتيا}
\olpseganchor{OLP-0426-B005}{end}

\olpseganchor{OLP-0426-B006}{start}
د مودالي فارمولونو په واسطه تعريفېدونکي د چوکاټونو ځينې
خاصيتونه د لومړۍ درجې له مخې نۀ تعريفېږي۔ خو که پر يوځاييزو پريديکاتونو
کميت ټاکنه ومنو (يعنې د دويمې درجې يوځاييزه کميت ټاکنه)، نو
ټول هغه چوکاټي خاصيتونه تعريفوو چې په مودالي ډول تعريفېدے شي۔
لار دا ده چې د يوه مودالي فارمول په کومه نړۍ کښې د صدق
شرطونه په منظم ډول د لومړۍ درجې له فارمولونه سره وتړو۔
دې ته له مودالي فارمولونو څخه د لومړۍ درجې فارمولونو ته
\emph{معياري ژباړه} وايي۔ د دې ژبې يوازې د لاسرسي اړيکې
دپاره دوه‌ځاييز پريديکات~\LR{$Q$} نۀ لري، بلکې
په~\LR{$!A$} کښې د شته بياني متغيرونه~\LR{$\Obj p_i$}
دپاره يوځاييز پريديکات~\LR{$P_i$} هم لري۔
\olpseganchor{OLP-0426-B006}{end}

\olpseganchor{OLP-0426-B007}{start}
\begin{defn}
  \emph{معياري ژباړه}~\LR{$\ST_x(!A)$} په استقرايي ډول
  په دې توګه تعريفېږي:
  \begin{enumerate}
  \item \indcase{!A}{\lfalse}{\LR{$\ST_x(\indfrm) = \lfalse$}.}
  \item \indcase{!A}{\ltrue}{\LR{$\ST_x(\indfrm) = \ltrue$}.}
  \item \indcase{!A}{\Obj p_i}{\LR{$\ST_x(\indfrm) = \Atom{P_i}{x}$}.}
  \item \indcase{!A}{\lnot !B}{\LR{$\ST_x(\indfrm) = \lnot \ST_x(!B)$}.}
  \item \indcase{!A}{(!B \land !C)}{\LR{$\ST_x(\indfrm) = (\ST_x(!B)
      \land \ST_x(!C))$}.}
  \item \indcase{!A}{(!B \lor !C)}{\LR{$\ST_x(\indfrm) = (\ST_x(!B)
      \lor \ST_x(!C))$}.}
  \item \indcase{!A}{(!B \lif !C)}{\LR{$\ST_x(\indfrm) = (\ST_x(!B)
      \lif \ST_x(!C))$}.}
  \item \indcase{!A}{(!B \liff !C)}{\LR{$\ST_x(\indfrm) = (\ST_x(!B)
      \liff \ST_x(!C))$}.}
  \item \indcase{!A}{\Box !B}{\LR{$\ST_x(\indfrm) =
      \lforall[y][(\Atom{Q}{x,y} \lif \ST_y(!B))]$}.}
  \item \indcase{!A}{\Diamond !B}{\LR{$\ST_x(\indfrm) =
      \lexists[y][(\Atom{Q}{x,y} \land \ST_y(!B))]$}.}
  \end{enumerate}
  په دوو وروستيو قواعدو کښې د کميت ټاکنې متغير تازه اخلو،
  يعنې د روانې نړۍ له متغيره بېل متغير ټاکو۔
\end{defn}
\olpseganchor{OLP-0426-B007}{end}

\olpseganchor{OLP-0426-B008}{start}
د بېلګې په توګه، \LR{$\ST_x(\Box p \lif p)$} دا دے:
\LR{$\lforall[y][(\Atom{Q}{x,y}
  \lif \Atom{P}{y})] \lif \Atom{P}{x}$}۔
د~\LR{$\ST_x(!A)$} د ژبې هر جوړښت يو ډومېن، د~\LR{$Q$}
دپاره ټاکل شوې دوه‌ځاييزه اړيکه، او د يوځاييزو
پريديکاتونه~\LR{$P_i$} دپاره د ډومېن فرعي سټونه غواړي۔
په بله وينا، د داسې جوړښت اجزا کټ مټ د~\LR{$!A$}
د يوه مدل اجزا دي: ډومېن د نړۍو سټ دے، د~\LR{$Q$} دپاره
ټاکل شوې دوه‌ځاييزه اړيکه د لاسرسي اړيکه ده، او
د~\LR{$P_i$} دپاره ټاکل شوي فرعي سټونه هماغه د
ارزښت ټاکنې~\LR{$V(\Obj p_i)$} دي۔ نو دا به د حېرانتيا
خبره نۀ وي چې په يوه مودالي مدل کښې د~\LR{$!A$} صدق او په
ورسره برابر جوړښت کښې د~\LR{$\ST_x(!A)$} صدق يو شان وي:
\olpseganchor{OLP-0426-B008}{end}

\olpseganchor{OLP-0426-B009}{start}
\begin{prop}\ollabel{prop:st}
  فرض کړئ \LR{$\mModel{M} = \tuple{W, R, V}$} وي،
  \LR{$\Struct{M'}$} داسې د لومړۍ درجې جوړښت وي
  چې \LR{$\Domain{M'} = W$}، \LR{$\Assign{Q}{M'} = R$} او
  \LR{$\Assign{P_i}{M'} = V(\Obj p_i)$} ولري، او
  \LR{$s(x) = w$} وي۔ بيا
  \[
  \mSat{M}{!A}[w] \text{\RL{ هله او يوازې هله چې }} \Sat{M'}{\ST_x(!A)}[s]
  \]
\end{prop}
\olpseganchor{OLP-0426-B009}{end}

\olpseganchor{OLP-0426-B010}{start}
\begin{proof}
  د~\LR{$!A$} پر جوړښت په استقرا۔
\end{proof}
\olpseganchor{OLP-0426-B010}{end}

\olpseganchor{OLP-0426-B011}{start}
\begin{prop}
  فرض کړئ \LR{$!A$} يو مودالي فارمول او
  \LR{$\mModel{F} = \tuple{W, R}$} يو چوکاټ وي۔
  \LR{$\Struct{F'}$} داسې د لومړۍ درجې جوړښت ونيسئ
  چې \LR{$\Domain{F'} = W$} او \LR{$\Assign{Q}{F'} = R$}
  ولري، او \LR{$!A'$} دا د دويمې درجې فارمول ونيسئ
  \[
  \lforall[X_1][\dots\lforall[X_n][\lforall[x][
        \SSubst{\ST_x(!A)}{\subst{X_1}{P_1}, \dots,
          \subst{X_n}{P_n}}]]],
  \]
  چې~\LR{$P_1$}، \dots، \LR{$P_n$} په~\LR{$\ST_x(!A)$} کښې ټول
  يوځاييز پريديکاتونه دي۔ بيا
  \[
  \mModel{F} \Entails !A \text{\RL{ هله او يوازې هله چې }} \Sat{F'}{!A'}
  \]
\end{prop}
\olpseganchor{OLP-0426-B011}{end}

\olpseganchor{OLP-0426-B012}{start}
\begin{proof}
  \LR{$\Sat{F'}{!A'}$} هله او يوازې هله چې د هر هغه
  جوړښت~\LR{$\mModel{M'}$} دپاره چې د چوکاټ هماغه ډومېن
  او د لاسرسي هماغه اړيکه لري، او
  \LR{$\Assign{P_i}{M'} \subseteq W$} د \LR{$i = 1$}،
  \dots،~\LR{$n$} دپاره ولري، او د هر داسې~\LR{$s$} دپاره
  چې \LR{$s(x) \in W$} وي،
  \LR{$\Sat{M'}{\ST_x(!A)}[s]$} وي۔ د
  \olref{prop:st} له مخې، دا هله او يوازې هله
  کېږي چې پر~\LR{$\mModel{F}$} ولاړ هر
  مدل~\LR{$\mModel{M}$} او هره نړۍ~\LR{$w \in W$}
  \LR{$\mSat{M}{!A}[w]$} پوره کړي؛ يعنې
  \LR{$\mModel{F} \Entails !A$}۔
\end{proof}
\olpseganchor{OLP-0426-B012}{end}

\olpseganchor{OLP-0426-B013}{start}
\begin{defn}
  د چوکاټونو يو ټولګے~\LR{$\mClass{F}$} \emph{د دويمې درجې
  له مخې تعريف‌وړ} دے، که د دويمې درجې په هغې ژبه کښې
  چې يوازې يو دوه‌ځاييز پريديکات~\LR{$P$} او يوازې
  پر يوځاييزو سټ‌متغيرونو کميت ټاکونکي لري، يوه
  تړلے فارمول~\LR{$!A$} وي، داسې چې
  \LR{$\mModel{F} = \tuple{W, R} \in \mClass{F}$}
  هله او يوازې هله چې په هغه
  جوړښت~\LR{$\Struct{M}$} کښې \LR{$\Sat{M}{!A}$}
  وي چې \LR{$\Domain{M} = W$} او
  \LR{$\Assign{P}{M} = R$} لري۔
\end{defn}
\olpseganchor{OLP-0426-B013}{end}

\olpseganchor{OLP-0426-B014}{start}
\begin{cor}
  که د چوکاټونو يو ټولګے په
  فارمول~\LR{$!A$} تعريفېدے شي،
  نو د لاسرسي د اړيکو اړوند ټولګے
  د يوې يوځاييزې د دويمې درجې جملې
  په واسطه تعريفېدے شي۔
\end{cor}
\olpseganchor{OLP-0426-B014}{end}

\olpseganchor{OLP-0426-B015}{start}
\begin{proof}
  د مخکيني ثبوت د دويمې درجې
  يوځاييزه جمله~\LR{$!A'$} اړوند
  خاصيت لري۔
\end{proof}
\olpseganchor{OLP-0426-B015}{end}

\olpseganchor{OLP-0426-B016}{start}
د بېلګې په توګه، بيا فارمول~\LR{$\Box p \lif p$}
وڅېړئ۔ دا انعکاسيت تعريفوي۔ انعکاسيت البته د
لومړۍ درجې په تړلے فارمول~\LR{$\lforall[x][\Atom{Q}{x,x}]$}
تعريفېدے شي۔ خو د دويمې درجې په دې يوځاييزه
تړلے فارمول هم تعريفېدے شي
\[
\lforall[X][\lforall[x][(\lforall[y][(\Atom{Q}{x,y} \lif \Atom{X}{y})]
    \lif \Atom{X}{x})]].
\]
دا په دې معنا ده چې دواړه جملې هم‌ارزې دي۔ دا خبره
په نېغه د ځان د قانع کولو دپاره داسې وڅېړئ: لومړے
فرض کړئ د دويمې درجې جمله په
جوړښت~\LR{$\Struct{M}$} کښې رښتيا ده۔ ځکه چې
\LR{$x$} او \LR{$X$} په عمومي ډول کميت ټاکل شوي دي، پاتې
برخه بايد د هر \LR{$x \in W$} او هر
سټ~\LR{$X \subseteq W$} دپاره سمه وي؛ د بېلګې په
توګه، د~\LR{$\Setabs{z}{Rxz}$} سټ دپاره، چې
\LR{$R = \Assign{Q}{M}$} وي۔ نو د هر
داسې~\LR{$s$} دپاره چې \LR{$s(x) \in W$}
او \LR{$s(X) = \Setabs{z}{Rxz}$} وي،
\LR{$\Sat{M}{\lforall[y][(\Atom{Q}{x,y} \lif
    \Atom{X}{y})] \lif \Atom{X}{x}}[s]$} لرو۔
خو د~\LR{$s(X)$} د ټاکنې له مخې دا معنا لري چې
\LR{$\Sat{M}{\lforall[y][(\Atom{Q}{x,y} \lif
    \Atom{Q}{x,y})] \lif \Atom{Q}{x,x}}[s]$}،
چې له~\LR{$\Atom{Q}{x,x}$} سره هم‌ارز دے، ځکه مقدم
تل رښتينی دے۔ ځکه چې \LR{$s(x)$} خپلسرے دے، نو
\LR{$\Sat{M}{\lforall[x][\Atom{Q}{x,x}]}$} لرو۔
\olpseganchor{OLP-0426-B016}{end}

\olpseganchor{OLP-0426-B017}{start}
اوس فرض کړئ \LR{$\Sat{M}{\lforall[x][\Atom{Q}{x,x}]}$}
او ښيو چې
\LR{$\Sat{M}{\lforall[X][\lforall[x][(\lforall[y][(\Atom{Q}{x,y} \lif
        \Atom{X}{y})] \lif \Atom{X}{x})]]}$}۔
هره ګومارنه~\LR{$s$} ونيسئ او فرض کړئ
\LR{$\Sat{M}{\lforall[y][(\Atom{Q}{x,y} \lif
    \Atom{X}{y})]}[s]$}۔
\LR{$s'$} د~\LR{$s$} هغه \LR{$y$}-بدله بڼه ونيسئ چې
\LR{$s'(y) = s(x)$} وي؛ بيا
\LR{$\Sat{M}{\Atom{Q}{x,y} \lif \Atom{X}{y}}[s']$}،
يعنې \LR{$\Sat{M}{\Atom{Q}{x,x} \lif \Atom{X}{x}}[s]$}
لرو۔ ځکه چې
\LR{$\Sat{M}{\lforall[x][\Atom{Q}{x,x}]}$}،
مقدم رښتينی دے او
\LR{$\Sat{M}{\Atom{X}{x}}[s]$} لرو؛ همدا مو
ښودل غوښتل۔
\olpseganchor{OLP-0426-B017}{end}

\olpseganchor{OLP-0426-B018}{start}
ځکه چې د چوکاټونو ځينې تعريفېدونکي ټولګي د
لومړۍ درجې له مخې نۀ تعريفېږي، نو د پورته
بڼې~\LR{$!A'$} لږ تر لږه يوه د دويمې درجې يوځاييزه
تړلے فارمول د لومړۍ درجې له هېڅ
تړلے فارمول سره هم‌ارزه نۀ ده۔
د دې پرېکړې دپاره، چې کومې داسې جملې هم‌ارزې
دي، اغېزمنه طريقه نشته۔
\olpseganchor{OLP-0426-B018}{end}

\olpseganchor{OLP-0427-B004}{start}
\olchapter{nml}{prf}{بديهي اشتقاقونه}
\olpseganchor{OLP-0427-B004}{end}

\olpseganchor{OLP-0428-B004}{start}
\olfileid{nml}{axs}{int}
\olpseganchor{OLP-0428-B004}{end}

\olpseganchor{OLP-0428-B005}{start}
\olsection{پېژندنه}
\olpseganchor{OLP-0428-B005}{end}

\olpseganchor{OLP-0428-B006}{start}
د بنسټيزې مودالي ژبې معناپوهنه مو د مودالي مدلونو په واسطه
لرو، او دا مفهوم هم لرو چې فارمول معتبر وي، يعنې په
ټولو مدلونو کښې په ټولو نړۍو کښې رښتيا وي؛ يا د مدلونو د
کوم ټولګي يا د چوکاټونو په نسبت معتبر وي، يعنې د ټولګي
په ټولو مدلونو کښې يا د هغه ټولګي پر چوکاټونو ولاړو
ټولو مدلونو کښې په ټولو نړۍو کښې رښتيا وي۔ منطق عموماً د اعتبار دا
معنايي ځانګړنې د ثبوت‌تيورۍ له د اشتقاق وړتيا سره تړي۔
موخه دا ده چې په يوه نظام کښې د اشتقاق وړتيا داسې
تعريف کړو چې فارمول هله او يوازې هله
د اشتقاق وړ وي چې معتبر وي۔
\olpseganchor{OLP-0428-B006}{end}

\olpseganchor{OLP-0428-B007}{start}
د اشتقاق تر ټولو ساده او له تاريخي پلوه تر ټولو
زاړۀ نظامونه د هېلبرټ په ډول يا بديهي اشتقاق
نظامونه دي۔ د ډېرو مودالي منطقونو دپاره د هېلبرټ په ډول
اشتقاق نظامونه جوړول نسبتاً اسانه دي: د ميتا-تيورۍ
د مطالعې د څيزونو په توګه ساده دي (د بېلګې په توګه، د
سموالي او بشپړتيا د ثبوت دپاره)۔ خو د طبيعي استنتاج د
نظامونو په پرتله په هغو کښې د فارمولونه ثبوت ډېر سخت دے۔
\olpseganchor{OLP-0428-B007}{end}

\olpseganchor{OLP-0428-B008}{start}
د هېلبرټ په ډول د اشتقاق په نظامونو کښې، د يوه
فارمول يو اشتقاق د فارمولونه هغه لړۍ ده
چې له ځينو بديهي اصلونو څخه، د څو استنتاجي قاعدو په
کارولو، تر اړوند فارمول رسېږي۔ ځکه چې غواړو د
اشتقاق نظام له معناپوهنې سره برابر وي، بايد دا
تضمين کړو چې ټول د اشتقاق وړ فارمولونه په ټولو مدلونو
کښې رښتيا وي (يا په هغو ټولو مدلونو کښې رښتيا وي چې
د هر بديهي اصل د ځايناستۍ هره بېلګه پکې رښتيا وي)۔ لومړے به د مودالي
منطقونو هغه خاصيتونه بېل کړو چې د دې کار دپاره اړين
دي: «عادي» مودالي منطقونه۔ د عادي مودالي منطقونو دپاره
يوازې دوه استنتاجي قاعدې فرضول غواړي: مودوس پوننس
او د لزوم ورکولو قاعده۔ د بديهي اصلونو په توګه د ټولو
تاوتولوژيو (د ځايناستۍ) ټول مثالونه اخلو، او د اړوند
مودالي منطق له مخې يو شمېر مودالي بديهي اصلونه هم۔
که يوازې د ټولو مدلونو ټولګي ته هم ګورو، بيا هم بايد
د~\LR{$\Ax{K}$} او~\LR{$\Ax{Dual}$}
د ځايناستۍ ټول مثالونه بديهي اصلونه وشمېرو۔ همدا
د لږ تر لږه عادي مودالي منطق~\LR{$\Log K$} بنسټ ږدي۔
\olpseganchor{OLP-0428-B008}{end}

\olpseganchor{OLP-0428-B009}{start}
\begin{defn}
د \emph{مودوس پوننس} قاعده دا استنتاجي طرحه ده:
\begin{prooftree}
\AxiomC{\LR{$!A$}}
\AxiomC{\LR{$!A \lif !B$}}
\RightLabel{\MP}
\BinaryInfC{\LR{$!B$}}
\end{prooftree}
وايو چې فارمول~\LR{$!B$} د مودوس پوننس له مخې
له فارمولونه \LR{$!A$}، \LR{$!C$} څخه \emph{راځي} هله
او يوازې هله چې \LR{$!C \ident !A \lif !B$} وي۔
\end{defn}
\olpseganchor{OLP-0428-B009}{end}

\olpseganchor{OLP-0428-B010}{start}
\begin{defn}
د \emph{لزوم ورکولو} قاعده دا استنتاجي طرحه ده:
\begin{prooftree}
\AxiomC{\LR{$!A$}}
\RightLabel{\Nec}
\UnaryInfC{\LR{$\Box !A$}}
\end{prooftree}
وايو چې فارمول~\LR{$!B$} د لزوم ورکولو له مخې
له فارمولونه \LR{$!A$} څخه راځي هله او يوازې هله
چې \LR{$!B \ident \Box !A$} وي۔
\end{defn}
\olpseganchor{OLP-0428-B010}{end}

\olpseganchor{OLP-0428-B011}{start}
\begin{defn}
د بديهي اصلونو له سټ~\LR{$\Sigma$} څخه يو \emph{اشتقاق}
د فارمولونه \LR{$!B_1$}، \LR{$!B_2$}، \dots، \LR{$!B_n$} داسې لړۍ
ده چې هر~\LR{$!B_i$} ئې يا
\begin{enumerate}
\item د يوې تاوتولوژۍ د ځايناستۍ مثال وي، يا
\item په~\LR{$\Sigma$} کښې د فارمول د ځايناستۍ مثال وي، يا
\item له دوو فارمولونه \LR{$!B_j$}، \LR{$!B_k$} څخه، چې \LR{$j$}، \LR{$k < i$} وي،
  د مودوس پوننس له مخې راځي، يا
\item له فارمول~\LR{$!B_j$} څخه، چې \LR{$j < i$} وي،
  د لزوم ورکولو له مخې راځي۔
\end{enumerate}
که داسې اشتقاق وي چې \LR{$!B_n \ident !A$}، نو وايو
چې~\LR{$!A$} \emph{له~\LR{$\Sigma$} څخه د اشتقاق وړ} دے؛ په نښو
\LR{$\Sigma \Proves !A$}۔
\end{defn}
\olpseganchor{OLP-0428-B011}{end}

\olpseganchor{OLP-0428-B012}{start}
له دې تعريف سره به څرګنده شي چې د ټولو د اشتقاق وړ
فارمولونو سټ يو عادي مودالي منطق جوړوي، او هر
د اشتقاق وړ فارمول په هر هغه مدل کښې رښتيا وي
چې د هر بديهي اصل د ځايناستۍ هره بېلګه پکې رښتيا وي۔
د اشتقاقونه دې
خاصيت ته \emph{سموالے} وايي۔ سرچپه لورې، يعنې
\emph{بشپړتيا}، ثابتول ګران دي۔
\olpseganchor{OLP-0428-B012}{end}

\olpseganchor{OLP-0429-B004}{start}
\olfileid{nml}{prf}{nor}
\olpseganchor{OLP-0429-B004}{end}

\olpseganchor{OLP-0429-B005}{start}
\olsection{عادي مودالي منطقونه}
\olpseganchor{OLP-0429-B005}{end}

\olpseganchor{OLP-0429-B006}{start}
د مودالي فارمولونه ځينې سټونه په اسانۍ سره د هغو
فارمولونه په توګه نۀ شي پېژندل کېدے چې د بديهي
اصلونو له يوه سټ څخه د اشتقاق وړ دي۔ غواړو مودالي
منطقونه ښې ځانګړنې ولري۔ تر هر څه وړاندې، هر هغه څه
چې په کلاسيک بياني منطق کښې ئې اشتقاق کولے شو،
بايد لا هم د اشتقاق وړ وي؛ البته دا په پام کښې نيسو
چې فارمولونه اوس \LR{$\Box$}
    او~\LR{$\Diamond$} هم لرلے شي۔
د دې دپاره غواړو چې مودالي منطق د تاوتولوژيو ټول
مثالونه ولري او د مودوس پوننس لاندې تړلے وي۔
\olpseganchor{OLP-0429-B006}{end}

\olpseganchor{OLP-0429-B007}{start}
\begin{defn}
  \emph{مودالي منطق} د مودالي فارمولونه داسې
  سټ~\LR{$\Sigma$} دے چې
  \begin{enumerate}
  \item ټولې تاوتولوژۍ لري، او
  \item د ځايناستۍ لاندې تړلے وي، يعنې که \LR{$!A \in \Sigma$}
    وي، او \LR{$!D_1$}، \dots، \LR{$!D_n$} فارمولونه وي، نو
    \[
    \SSubst{!A}{\subst{!D_1}{p_1}, \dots, \subst{!D_n}{p_n}} \in \Sigma,
    \]
    \item د \emph{مودوس پوننس} لاندې تړلے وي، يعنې که \LR{$!A$}
      او \LR{$!A \lif !B \in \Sigma$} وي، نو \LR{$!B \in \Sigma$} وي۔
  \end{enumerate}
\end{defn}
\olpseganchor{OLP-0429-B007}{end}

\olpseganchor{OLP-0429-B008}{start}
د مودالي منطقونو دپاره د اړيکو معناپوهنه کارولو په
موخه بايد دا هم وغواړو چې په ټولو مودالي مدلونو کښې
معتبر ټول فارمولونه پکې شامل وي۔ څرګندېږي چې دا
غوښتنه هغه وخت پوره کېږي چې د \Ax{K}
او~\Dual{} ټول مثالونه د اشتقاق وړ وي، او هر کله
چې فارمول~\LR{$!A$} د اشتقاق وړ وي، نو~\LR{$\Box !A$}
هم وي۔ هغه مودالي منطق چې دا شرطونه پوره کړي،
\emph{عادي} بلل کېږي۔ (البته ناعادي مودالي منطقونه هم
شته، خو د اړيکو معمول مدلونه د هغو دپاره مناسب نۀ دي۔)
\olpseganchor{OLP-0429-B008}{end}

\olpseganchor{OLP-0429-B009}{start}
\begin{defn}
  مودالي منطق~\LR{$\Sigma$} \emph{عادي} دے، که دا ولري:
  \begin{align*}
    \tag{\Ax{K}} & \Box(p \lif q) \lif (\Box p \lif \Box q),
    \\
      \tag{\Dual} & \Diamond p \liff \lnot\Box\lnot p
  \end{align*}
  او د \emph{لزوم ورکولو} لاندې تړلے وي، يعنې که \LR{$!A \in
  \Sigma$} وي، نو \LR{$\Box !A \in \Sigma$} وي۔
\end{defn}
\olpseganchor{OLP-0429-B009}{end}

\olpseganchor{OLP-0429-B010}{start}
پام وکړئ چې تاوتولوژيکي استلزام د \emph{يوې نړۍ په
کښې د صدق} ساتلو دپاره په کافي اندازه دقيق دے، خو
قاعده~\Nec{} يوازې \emph{په يوه مدل کښې صدق} ساتي
(او له دې سره په يوه چوکاټ يا د چوکاټونو په يوه ټولګي
کښې اعتبار هم ساتي)۔
\olpseganchor{OLP-0429-B010}{end}

\olpseganchor{OLP-0429-B011}{start}
\begin{prop}\ollabel{prop:rk}
  هر عادي مودالي منطق د~\RK قاعدې لاندې تړلے دے:
  \begin{prooftree}
    \AxiomC{\LR{$!A_1 \lif (!A_2 \lif \cdots (!A_{n-1} \lif !A_n)\cdots)$}}
    \RightLabel{\RK}
    \UnaryInfC{\LR{$\Box!A_1 \lif (\Box!A_2 \lif \cdots (\Box!A_{n-1}
      \lif \Box!A_n)\cdots).$}}
  \end{prooftree}
\end{prop}
\olpseganchor{OLP-0429-B011}{end}

\olpseganchor{OLP-0429-B012}{start}
\begin{proof}
  پر~\LR{$n$} استقرا کوو: که \LR{$n = 1$} وي، نو دا قاعده
  هماغه \Nec ده، او هر عادي مودالي منطق د~\Nec لاندې
  تړلے دے۔
\olpseganchor{OLP-0429-B012}{end}

\olpseganchor{OLP-0429-B013}{start}
  اوس فرض کړئ پايله د~\LR{$n-1$} دپاره سمه ده؛ ښيو چې
  د~\LR{$n$} دپاره هم سمه ده۔
\olpseganchor{OLP-0429-B013}{end}

\olpseganchor{OLP-0429-B014}{start}
  فرض کړئ
  \begin{align*}
  & !A_1 \lif (!A_2 \lif \cdots (!A_{n-1} \lif !A_n)\cdots) \in \Sigma
  \intertext{\RL{د استقرا د فرض له مخې لرو:}}
  & \Box!A_1 \lif (\Box!A_2 \lif \cdots \Box(!A_{n-1} \lif !A_n)\cdots)
  \in \Sigma
  \intertext{\RL{څرنګه چې \LR{$\Sigma$} عادي مودالي منطق دے، د~\LR{$\Ax{K}$}
    ټول مثالونه لري، په ځانګړې توګه:}}
  & \Box(!A_{n-1} \lif !A_n) \lif (\Box!A_{n-1} \lif \Box!A_n) \in \Sigma
  \intertext{\RL{د مودوس پوننس او د تاوتولوژۍ د مناسبو مثالونو له مخې ترلاسه کوو:}}
  & \Box!A_1 \lif (\Box!A_2 \lif \cdots (\Box!A_{n-1}
  \lif \Box!A_n)\cdots) \in \Sigma. 
  \end{align*}
\end{proof}
\olpseganchor{OLP-0429-B014}{end}

\olpseganchor{OLP-0429-B015}{start}
\begin{prop}\ollabel{prop:notDiamondBot}
  هر عادي مودالي منطق~\LR{$\Sigma$}، \LR{$\lnot\Diamond\lfalse$} لري۔
\end{prop}
\olpseganchor{OLP-0429-B015}{end}

\olpseganchor{OLP-0429-B016}{start}
\begin{prob}
  \olref[nml][prf][nor]{prop:notDiamondBot} ثابت کړئ۔
\end{prob}
\olpseganchor{OLP-0429-B016}{end}

\olpseganchor{OLP-0429-B017}{start}
\begin{prop}
  فرض کړئ \LR{$!A_1$}، \dots، \LR{$!A_n$} فارمولونه وي۔ بيا
  يو تر ټولو کوچنے عادي مودالي منطق~\LR{$\Sigma$} شته
  چې د~\LR{$!A_1$}، \dots، \LR{$!A_n$} ټول مثالونه لري۔
\end{prop}
\olpseganchor{OLP-0429-B017}{end}

\olpseganchor{OLP-0429-B018}{start}
\begin{proof}
  د ورکړل شويو~\LR{$!A_1$}، \dots، \LR{$!A_n$} دپاره،
  \LR{$\Sigma$} د ټولو هغو عادي مودالي منطقونو تقاطع
  وټاکئ چې د~\LR{$!A_1$}، \dots، \LR{$!A_n$} ټول مثالونه لري۔
  د تقاطع د سټونو کورنۍ تشه نۀ ده، ځکه چې
  \LR{$\Frm[L]$}، يعنې د ټولو فارمولونه سټ، يو داسې
  عادي مودالي منطق دے۔ تقاطع ټول اړين بديهي اصلونه
  او د ورکړل شويو فارمولونو ټولې بېلګې لري، ځکه چې په
  هر غړي کښې شته۔ که د يوې قاعدې مقدمې په تقاطع کښې وي،
  په هر غړي کښې هم وي؛ نو د قاعدې پايله هم په تقاطع کښې ده۔
  له دې امله تقاطع عادي ده او د کورنۍ د هر غړي فرعي سټ ده۔
\end{proof}
\olpseganchor{OLP-0429-B018}{end}

\olpseganchor{OLP-0429-B019}{start}
\begin{defn}
تر ټولو کوچني عادي مودالي منطق ته، چې \LR{$!A_1$}، \dots، \LR{$!A_n$}
لري، \emph{مودالي نظام} وايي، او په~\LR{$\Log{K} !A_1 \dots
!A_n$} ئې ښيو۔ تر ټولو کوچنے عادي مودالي منطق په~\Log{K}
ښودل کېږي۔
\end{defn}
\olpseganchor{OLP-0429-B019}{end}

\olpseganchor{OLP-0430-B004}{start}
\olfileid{nml}{prf}{prf}
\olpseganchor{OLP-0430-B004}{end}

\olpseganchor{OLP-0430-B005}{start}
\olsection{اشتقاقونه او مودالي نظامونه}
\olpseganchor{OLP-0430-B005}{end}

\olpseganchor{OLP-0430-B006}{start}
لومړے د عادي مودالي منطقونو دپاره اشتقاق
تعريفوو۔ په لنډ ډول، اشتقاق د فارمولونه
داسې لړۍ ده چې هر غړے ئې يا د څو
\emph{بديهي اصلونو} څخه د يوه (د ځايناستۍ مثال)
وي، يا له مخکنيو غړي څخه د څو استنتاجي
قاعدو له يوې راځي۔ د عادي مودالي منطقونو دپاره
د تاوتولوژيو، \Ax{K}، او~\Dual{}
ټول مثالونه بديهي اصلونه شمېرل کېږي۔ په دې سره
مودالي نظام~\LR{$\Log{K}$}، يعنې تر ټولو کوچنے عادي
مودالي منطق، ترلاسه کېږي۔ خو د نورو نظامونو د
ترلاسه کولو دپاره نور بديهي اصلونه هم زياتولے شو۔
قاعدې تل مودوس پوننس~\MP{} او د لزوم ورکولو قاعده~\Nec دي۔
\olpseganchor{OLP-0430-B006}{end}

\olpseganchor{OLP-0430-B007}{start}
\begin{defn}
  که مودالي نظام~\LR{$\Log{K} !A_1 \dots !A_n$} او
  فارمول~\LR{$!B$} راکړل شوي وي، وايو چې~\LR{$!B$}
  په~\LR{$\Log{K} !A_1 \dots !A_n$} کښې \emph{د اشتقاق وړ}
  دے، په نښو \LR{$\Log{K} !A_1 \dots !A_n \Proves !B$}،
  هله او يوازې هله چې داسې فارمولونه \LR{$!C_1$}،
  \dots، \LR{$!C_k$} وي چې \LR{$!C_k = !B$}، او هر~\LR{$!C_i$} يا
  د تاوتولوژۍ مثال وي، يا د \LR{$\Ax{K}$}،
  \LR{$\Dual$}، \LR{$!A_1$}، \dots، \LR{$!A_n$} څخه د يوه مثال
  وي، يا د \MP{} يا~\Nec د قاعدې له مخې له مخکنيو
  فارمولونه څخه راځي۔
\end{defn}
\olpseganchor{OLP-0430-B007}{end}

\olpseganchor{OLP-0430-B008}{start}
لاندې قضيه موږ ته اجازه راکوي چې د~\LR{$!B$} يو
\LR{$\Sigma$}-اشتقاق په ښودلو سره
\LR{$!B \in \Sigma$} ثابت کړو۔
\olpseganchor{OLP-0430-B008}{end}

\olpseganchor{OLP-0430-B009}{start}
\begin{prop}
  \LR{$\Log{K} !A_1 \dots !A_n = \Setabs{!B}{\Log{K} !A_1 \dots !A_n \Proves !B}$}۔
\end{prop}
\olpseganchor{OLP-0430-B009}{end}

\olpseganchor{OLP-0430-B010}{start}
\begin{proof}
  د اشتقاقونه پر اوږدوالي استقرا کوو، څو وښيو چې
  \LR{$\Setabs{!B}{\Log{K} !A_1 \dots !A_n \Proves !B} \subseteq
  \Log{K} !A_1 \dots !A_n$}۔
\olpseganchor{OLP-0430-B010}{end}

\olpseganchor{OLP-0430-B011}{start}
  که د~\LR{$!B$} د اشتقاق اوږدوالے~\LR{$1$} وي، يوازې
  يو فارمول لري۔ دا فارمول د \MP{} يا \Nec
  په واسطه له مخکنيو فارمولونو څخه نۀ شي راتلے، نو
  بايد د تاوتولوژۍ مثال، د \Ax{K} مثال،
  د \LR{$\Dual$} مثال، يا د \LR{$!A_1$}، \dots،~\LR{$!A_n$}
  څخه د يوه مثال وي۔ خو \LR{$\Log{K}!A_1\dots!A_n$} هم
  دا لري، نو \LR{$!B \in \Log{K}!A_1 \dots !A_n$}۔
\olpseganchor{OLP-0430-B011}{end}

\olpseganchor{OLP-0430-B012}{start}
  که د~\LR{$!B$} د اشتقاق اوږدوالے~\LR{$> 1$} وي، نو
  \LR{$!B$} سربېره پر دې د \MP{} يا \Nec{} په واسطه له
  هغو فارمولونه څخه هم ترلاسه کېدے شي چې د
  اشتقاق وروستۍ کرښه نۀ وي۔ که \LR{$!B$} له~\LR{$!C$}
  او~\LR{$!C \lif !B$} څخه (د \MP په واسطه) راځي، نو
  د استقرا د فرض له مخې \LR{$!C$} او \LR{$!C \lif !B \in
  \Log{K}!A_1 \dots !A_n$} دي۔ خو هر مودالي منطق
  د مودوس پوننس لاندې تړلے دے، نو
  \LR{$!B \in \Log{K}!A_1 \dots !A_n$}۔ که
  \LR{$!B \equiv \Box !C$} د~\LR{$!C$} څخه د \Nec په واسطه
  راځي، نو د استقرا د فرض له مخې \LR{$!C
  \in \Log{K}!A_1 \dots !A_n$}۔ خو هر عادي مودالي
  منطق د~\LR{$\Nec$} لاندې تړلے دے، نو
  \LR{$!B \in \Log{K}!A_1\dots!A_n$}۔
\olpseganchor{OLP-0430-B012}{end}

\olpseganchor{OLP-0430-B013}{start}
  سرچپه شموليت داسې ښيو چې
  \LR{$\Sigma = \Setabs{!B}{\Log{K} !A_1 \dots !A_n \Proves !B }$}
  يو عادي مودالي منطق دے چې د \LR{$!A_1$}، \dots، \LR{$!A_n$}
  ټول مثالونه لري؛ بيا دا کاروو چې
  \LR{$\Log{K} !A_1 \dots !A_n$} د تعريف له مخې تر ټولو
  کوچنے داسې منطق دے۔ د هر اضافي بديهي اصل هره تعويضي
  بېلګه يو‌کرښيز استنباط لري، نو په دې سټ کښې شامله ده۔
  \begin{enumerate}
    \item هره تاوتولوژي~\LR{$!B$} د تاوتولوژۍ يو مثال دے، نو
      \LR{$\Log{K}!A_1\dots!A_n \Proves !B$}، او \LR{$\Sigma$}
      ټولې تاوتولوژۍ لري۔
    \item که \LR{$\Log{K}!A_1\dots!A_n \Proves !C$} او
      \LR{$\Log{K}!A_1\dots!A_n \Proves !C \lif !B$}، نو
      \LR{$\Log{K}!A_1\dots!A_n \Proves !B$}: د~\LR{$!C$}
      اشتقاق د~\LR{$!C \lif !B$} له اشتقاق
      سره يوځاے کړئ او د~\LR{$!B$} کرښه ورزياته کړئ۔
      وروستۍ کرښه د \MP{} له مخې توجيه کېږي۔ نو
      \LR{$\Sigma$} د مودوس پوننس لاندې تړلے دے۔
    \item که \LR{$!B$} يو اشتقاق ولري، نو د~\LR{$!B$}
      هر د ځايناستۍ مثال هم اشتقاق لري:
      ځايناستي د اشتقاق پر هر فارمول
      تطبيق کړئ۔ (تمرین: د اشتقاقونه پر اوږدوالي
      په استقرا ثابت کړئ چې پايله لا هم سم
      اشتقاق دے۔) نو \LR{$\Sigma$} د يو شان
      ځايناستۍ لاندې تړلے دے۔ (اوس مو ثابته کړه
      چې \LR{$\Sigma$} د مودالي منطق ټول شرطونه پوره کوي۔)
    \item \LR{$\Log{K}!A_1\dots!A_n \Proves \Ax{K}$} لرو، نو
      \LR{$\Ax{K} \in \Sigma$}۔
    \item همدارنګه
      \LR{$\Log{K}!A_1\dots!A_n \Proves \Dual$} لرو، نو
      \LR{$\Dual \in \Sigma$}۔
    \item که \LR{$\Log{K}!A_1\dots!A_n \Proves !C$}، نو
      زياته شوې کرښه~\LR{$\Box !C$} د~\Nec له مخې
      توجيه کېږي۔ په پايله کښې \LR{$\Sigma$} د~\Nec
      لاندې تړلے دے۔ نو \LR{$\Sigma$} عادي دے۔
  \end{enumerate}
\end{proof}
\olpseganchor{OLP-0430-B013}{end}

\olpseganchor{OLP-0431-B004}{start}
\olfileid{nml}{prf}{prk}
\olpseganchor{OLP-0431-B004}{end}

\olpseganchor{OLP-0431-B005}{start}
\olsection{په \texorpdfstring{\Log{K}}{K} کښې ثبوتونه}
\olpseganchor{OLP-0431-B005}{end}

\olpseganchor{OLP-0431-B006}{start}
په تر ټولو کوچني مودالي نظام کښې د ثبوتونو د تمرين دپاره ښيو چې د
\olref[syn][sch]{tab:valid-invalidSchemas} په کيڼ اړخ کښې ټول معتبر
فارمولونه په \Log{K} کښې ثبوتونه لري۔
\olpseganchor{OLP-0431-B006}{end}

\olpseganchor{OLP-0431-B007}{start}
\begin{prop}
  \LR{$\Log{K} \Proves \Box!A \lif \Box (!B\lif !A)$}
\end{prop}
\olpseganchor{OLP-0431-B007}{end}

\olpseganchor{OLP-0431-B008}{start}
\begin{proof}
  \begin{derivation}
    1. & \LR{$!A \lif (!B \lif !A)$} & \Taut \\
    2. & \LR{$\Box(!A \lif (!B \lif !A))$} & \Nec, 1 \\
    3. & \LR{$\Box(!A \lif (!B \lif !A)) \lif
    (\Box!A \lif \Box(!B \lif !A))$} & \Ax{K}\\
    4. & \LR{$\Box!A \lif \Box (!B\lif !A)$} & \MP, 2, 3
  \end{derivation}
\end{proof}
\olpseganchor{OLP-0431-B008}{end}

\olpseganchor{OLP-0431-B009}{start}
\begin{prop}
  \LR{$\Log{K} \Proves \Box(!A \land !B) \lif (\Box !A \land \Box!B)$}
\end{prop}
\olpseganchor{OLP-0431-B009}{end}

\olpseganchor{OLP-0431-B010}{start}
\begin{proof}
\begin{derivation}
  1. & \LR{$(!A \land !B) \lif !A$} & \Taut \\
  2. & \LR{$\Box((!A \land !B) \lif !A)$} & \Nec, 1 \\
  3. & \LR{$\Box((!A \land !B) \lif !A) \lif (\Box(!A \land !B) \lif \Box!A)$}
  & \Ax{K} \\
  4. & \LR{$\Box(!A \land !B) \lif \Box!A$} & \MP, 2, 3 \\
  5. & \LR{$(!A \land !B) \lif !B$} & \Taut \\
  6. & \LR{$\Box((!A \land !B) \lif !B)$} & \Nec, 5 \\
  7. & \LR{$\Box((!A \land !B) \lif !B) \lif (\Box(!A \land !B) \lif \Box!B)$}
  & \Ax{K} \\
  8. & \LR{$\Box(!A \land !B) \lif \Box!B$} & \MP, 6, 7 \\
  9. & \LR{$(\Box(!A \land !B) \lif \Box!A) \lif{}$} \\
  & \qquad \LR{$((\Box(!A \land !B) \lif \Box!B) \lif{}$} \\
  & \qquad \LR{$(\Box(!A \land !B) \lif (\Box !A \land \Box!B)))$} & \Taut\\
  10. &  \LR{$(\Box(!A \land !B) \lif \Box!B) \lif{}$} \\
  & \qquad \LR{$(\Box(!A \land !B) \lif (\Box !A \land \Box!B))$} & \MP, 4, 9\\
  11. & \LR{$\Box(!A \land !B) \lif (\Box !A \land \Box!B)$} & \MP, 8, 10.
\end{derivation}
پام وکړئ چې د~\LR{$9$} کرښې فارمول د لاندې تاوتولوژۍ يو مثال دے:
\[
(p \lif q) \lif ((p \lif r) \lif (p \lif (q \land r))).
\]
\end{proof}
\olpseganchor{OLP-0431-B010}{end}

\olpseganchor{OLP-0431-B011}{start}
\begin{prop}
  \LR{$\Log{K}\Proves (\Box!A \land \Box!B) \lif \Box (!A \land !B)$}
\end{prop}
\olpseganchor{OLP-0431-B011}{end}

\olpseganchor{OLP-0431-B012}{start}
\begin{proof}
  \begin{derivation}
    1. & \LR{$!A \lif (!B \lif (!A \land !B))$} & \Taut \\
    2. & \LR{$\Box(!A \lif (!B \lif (!A \land !B)))$} & \Nec, 1 \\
    3. & \LR{$\Box(!A \lif (!B \lif (!A \land !B))) \lif (\Box!A \lif \Box(!B \lif (!A \land !B)))$} & \Ax{K}\\
    4. & \LR{$\Box!A \lif \Box(!B \lif (!A \land !B))$} & \MP, 2, 3\\
    5. & \LR{$\Box(!B \lif (!A \land !B)) \lif (\Box!B \lif \Box(!A \land !B))$} & \Ax{K} \\
    6. & \LR{$(\Box!A \lif \Box(!B \lif (!A \land !B))) \lif {}$} \\
    & \qquad \LR{$(\Box(!B \lif (!A \land !B)) \lif (\Box!B \lif \Box(!A \land !B))) \lif {}$}\\
    & \qquad \LR{$(\Box!A \lif (\Box !B \lif \Box(!A \land !B))))$} & \Taut\\
    7. & \LR{$(\Box(!B \lif (!A \land !B)) \lif (\Box!B \lif \Box(!A \land !B))) \lif {}$}\\
    & \qquad \LR{$(\Box!A \lif (\Box !B \lif \Box(!A \land !B)))$} & \MP, 4, 6\\
    8. & \LR{$\Box!A \lif (\Box !B \lif \Box(!A \land !B))$} & \MP, 5, 7\\
    9. & \LR{$(\Box!A \lif (\Box !B \lif \Box(!A \land !B)))) \lif {}$}\\
    & \qquad \LR{$((\Box!A \land \Box!B) \lif \Box(!A \land !B))$} & \Taut\\
    10. & \LR{$(\Box!A \land \Box!B) \lif \Box(!A \land !B)$} & \MP, 8, 9
  \end{derivation}
  د \LR{$6$} او \LR{$9$} کرښو فارمولونه د لاندې تاوتولوژيو مثالونه دي:
  \begin{align*}
    (p \lif q) & \lif ((q \lif r) \lif (p \lif r)) \\
    (p \lif (q \lif r)) & \lif ((p \land q) \lif r)
  \end{align*}
\end{proof}
\olpseganchor{OLP-0431-B012}{end}

\olpseganchor{OLP-0431-B013}{start}
\begin{prop}
  \LR{$\Log{K} \Proves \lnot\Box p \lif \Diamond \lnot
    p$}
\end{prop}
\olpseganchor{OLP-0431-B013}{end}

\olpseganchor{OLP-0431-B014}{start}
\begin{proof}
  % Box and Diamond are primitive
  \begin{derivation}
    1. & \LR{$\Diamond \lnot p \liff \lnot\Box\lnot\lnot p$} & \Dual\\
    2. & \LR{$(\Diamond \lnot p \liff \lnot\Box\lnot\lnot p) \lif {}$} \\
    & \qquad \LR{$(\lnot \Box \lnot\lnot p \lif \Diamond \lnot p)$}& \Taut\\
    3. & \LR{$\lnot \Box \lnot\lnot p \lif \Diamond \lnot p$} & \MP, 1, 2\\
    4. & \LR{$\lnot\lnot p \lif p$} & \Taut\\
    5. & \LR{$\Box(\lnot\lnot p \lif p)$} & \Nec, 4\\
    6. & \LR{$\Box(\lnot\lnot p \lif p) \lif (\Box \lnot\lnot p \lif \Box p)$} & \Ax{K}\\
    7. & \LR{$(\Box \lnot\lnot p \lif \Box p)$} & \MP, 5, 6\\
    8. & \LR{$(\Box \lnot\lnot p \lif \Box p) \lif (\lnot \Box p \lif \lnot\Box\lnot\lnot p)$} & \Taut\\
    9. & \LR{$\lnot \Box p \lif \lnot\Box\lnot\lnot p$} & \MP, 7, 8\\
    10. & \LR{$(\lnot \Box p \lif \lnot\Box\lnot\lnot p) \lif {}$} \\
    & \qquad \LR{$((\lnot \Box \lnot\lnot p \lif \Diamond \lnot p) \lif (\lnot \Box p \lif \Diamond\lnot p))$} & \Taut\\
    11. & \LR{$(\lnot \Box \lnot\lnot p \lif \Diamond \lnot p) \lif (\lnot \Box p \lif \Diamond\lnot p)$} & \MP, 9, 10\\
    12. & \LR{$\lnot\Box p \lif \Diamond \lnot p$} & \MP, 3, 11\\
  \end{derivation}
  د \LR{$8$} او \LR{$10$} کرښو فارمولونه د لاندې تاوتولوژيو مثالونه دي:
  \begin{align*}
    & (p \lif q) \lif (\lnot q \lif \lnot p) \\
    & (p \lif q) \lif ((q \lif r) \lif (p \lif r)).
  \end{align*}
  
\end{proof}
\olpseganchor{OLP-0431-B014}{end}

\olpseganchor{OLP-0431-B015}{start}
\begin{prob}
  په~\LR{$\Log{K}$} کښې د لاندې فارمولونه دپاره اشتقاقونه ومومئ:
  \begin{enumerate}
    \item \LR{$\Box \lnot p \lif \Box(p \lif q)$}
    \item \LR{$(\Box p \lor \Box q) \lif \Box(p \lor q)$}
    \item \LR{$\Diamond p \lif \Diamond(p \lor q)$}
  \end{enumerate}
\end{prob}
\olpseganchor{OLP-0431-B015}{end}

\olpseganchor{OLP-0432-B004}{start}
\olfileid{nml}{prf}{der}
\olpseganchor{OLP-0432-B004}{end}

\olpseganchor{OLP-0432-B005}{start}
\olsection{مشتقې قاعدې}
\olpseganchor{OLP-0432-B005}{end}

\olpseganchor{OLP-0432-B006}{start}
د اشتقاقونه موندل او ليکل څرګنده ده چې ګران، دروند او تکراري کار
دے۔ د بېلګې په توګه، ډېر ځله غواړو له \LR{$!A \lif !B$} څخه
\LR{$\Box !A \lif \Box !B$} ته لاړ شو، يعنې قاعده~\RK وکاروو۔ د دې دپاره
لومړے \Nec کارول، بيا د~\Ax{K} مناسب مثال ثبتول، او ورپسې~\MP کارول
پکار دي۔ له \LR{$!A \lif !B$} او \LR{$!B \lif !C$} څخه \LR{$!A \lif !C$} ته د تلو
دپاره دا (اوږد) تاوتولوژيک مثال ثبتول اړين دي:
\[
(!A \lif !B) \lif ((!B \lif !C) \lif (!A \lif !C))
\]
او بيا \MP{} دوه ځله کارول پکار دي۔ ډېر ځله غواړو يو فرعي
فارمول په داسې فارمول بدل کړو چې ورسره يې هم‌ارزښتيا راته معلومه
وي؛ د بېلګې په توګه، \LR{$\Diamond
  !A$} په \LR{$\lnot\Box\lnot !A$}، يا
\LR{$\lnot\lnot !A$} په \LR{$!A$}۔ نو د اصلي اشتقاق د ټولو کرښو د ليکلو
پر ځاے، دا ثبتول اسانه دي چې منځني ګامونه ولې د اشتقاق وړ دي۔ د همدې
دپاره د اشتقاق وړتيا په باب څو حقيقتونه راغونډوو۔
\olpseganchor{OLP-0432-B006}{end}

\olpseganchor{OLP-0432-B007}{start}
\begin{prop}
  که \LR{$\Log{K} \Proves !A_1$}، \dots، \LR{$\Log{K} \Proves !A_n$} وي، او
  \LR{$!B$} د قضيوي منطق له مخې له \LR{$!A_1$}، \dots، \LR{$!A_n$} څخه راځي، نو
  \LR{$\Log{K} \Proves !B$}۔
\end{prop}
\olpseganchor{OLP-0432-B007}{end}

\olpseganchor{OLP-0432-B008}{start}
\begin{proof}
  که \LR{$!B$} د قضيوي منطق له مخې له \LR{$!A_1$}، \dots،~\LR{$!A_n$} څخه راځي، نو
  \[
  !A_1 \lif (!A_2 \lif \cdots (!A_n \lif !B)\dots)
  \]
  د تاوتولوژۍ يو مثال دے۔ د \MP{} قاعده \LR{$n$} ځله کارول د~\LR{$!B$}
  اشتقاق راکوي۔
\end{proof}
\olpseganchor{OLP-0432-B008}{end}

\olpseganchor{OLP-0432-B009}{start}
د دې قضيې کارول به په~\PL نښه کوو۔
\olpseganchor{OLP-0432-B009}{end}

\olpseganchor{OLP-0432-B010}{start}
\begin{prop}
  که \LR{$\Log{K} \Proves !A_1 \lif (!A_2 \lif \cdots (!A_{n-1} \lif
  !A_n)\dots)$} وي، نو \LR{$\Log{K} \Proves \Box !A_1 \lif (\Box !A_2 \lif
  \cdots (\Box !A_{n-1} \lif \Box !A_n)\dots)$}۔
\end{prop}
\olpseganchor{OLP-0432-B010}{end}

\olpseganchor{OLP-0432-B011}{start}
\begin{proof}
  دا د~\LR{$n$} پر استقرا ثابتېږي، لکه د \olref[nor]{prop:rk} په ثبوت کښې۔
\end{proof}
\olpseganchor{OLP-0432-B011}{end}

\olpseganchor{OLP-0432-B012}{start}
د دې قضيې کارول به په~\RK نښه کوو۔ راځئ وښيو چې دا پايلې د
د اشتقاق وړتيا د ادعاوو ثابتول څنګه اسانه کوي۔
\olpseganchor{OLP-0432-B012}{end}

\olpseganchor{OLP-0432-B013}{start}
\begin{prop}
  \LR{$\Log{K} \Proves (\Box!A \land \Box!B) \lif \Box (!A \land !B)$}
\end{prop}
\olpseganchor{OLP-0432-B013}{end}

\olpseganchor{OLP-0432-B014}{start}
\begin{proof}
  \begin{derivation}
    1. & \LR{$\Log{K} \Proves !A \lif (!B \lif (!A \land !B))$} & \Taut \\
    2. & \LR{$\Log{K} \Proves \Box!A \lif (\Box !B \lif \Box(!A \land !B))$} & \RK, 1\\
    3. & \LR{$\Log{K} \Proves (\Box!A \land \Box!B) \lif \Box(!A \land !B)$} & \PL, 2
  \end{derivation}
\end{proof}
\olpseganchor{OLP-0432-B014}{end}

\olpseganchor{OLP-0432-B015}{start}
\begin{prop}\ollabel{prop:rewriting}
  که \LR{$\Log{K} \Proves !A \liff !B$} او \LR{$\Log{K} \Proves
  \Subst{!C}{!A}{q}$} وي، نو \LR{$\Log{K} \Proves \Subst{!C}{!B}{q}$}۔
\end{prop}
\olpseganchor{OLP-0432-B015}{end}

\olpseganchor{OLP-0432-B016}{start}
\begin{proof}
  ثبوت په لاندې تمرين کښې غوښتل شوے دے۔
\end{proof}
\olpseganchor{OLP-0432-B016}{end}

\olpseganchor{OLP-0432-B017}{start}
\begin{prob}
  \olref[nml][prf][der]{prop:rewriting} داسې ثابت کړئ: د~\LR{$!C$} پر
  پېچلتيا په استقرا وښيئ چې که \LR{$\Log{K} \Proves !A \liff !B$} وي، نو
  \LR{$\Log{K} \Proves \Subst{!C}{!A}{q} \liff \Subst{!C}{!B}{q}$}۔
\end{prob}
\olpseganchor{OLP-0432-B017}{end}

\olpseganchor{OLP-0432-B018}{start}
دا قضيه په ځانګړي ډول هغه وخت ګټوره ده چې غواړو \LR{$\Diamond$} په
\LR{$\Box$} (يا برعکس) واړوو، يا په فارمول کښې دوه‌ګونې نفي لرې
کړو۔ په راتلونکو برخو کښې، کله چې فارمول~\LR{$!C(!B)$} د~\LR{$!C(!A)$}
پر ځاے وليکو، د \olref{prop:rewriting} کارول به په «\LR{$!B$} د \LR{$!A$} پر
ځاے» نښه کوو۔ په بل عبارت، «\LR{$!B$} د \LR{$!A$} پر ځاے» د لاندې ګامونو
لنډيز دے:
  \begin{derivation}
    & \LR{$\Proves !C(!A)$} \\
    & \LR{$\Proves !A \liff !B$}\\
    & \LR{$\Proves !C(!B)$} & د \olref{prop:rewriting} له مخې
  \end{derivation}
د بېلګې په توګه:
\olpseganchor{OLP-0432-B018}{end}

\olpseganchor{OLP-0432-B019}{start}
\begin{prop}
  \LR{$\Log{K} \Proves \lnot\Box p \lif \Diamond \lnot p$}
\end{prop}
\olpseganchor{OLP-0432-B019}{end}

\olpseganchor{OLP-0432-B020}{start}
\begin{proof}
  % Diamond is primitive
  \begin{derivation}
    1. & \LR{$\Log{K} \Proves \Diamond \lnot p \liff \lnot\Box\lnot\lnot p$} &
    \Dual\\
    2. & \LR{$\Log{K} \Proves \lnot \Box \lnot\lnot p \lif \Diamond \lnot p$} & \PL, 1\\
    3. & \LR{$\Log{K} \Proves \lnot\Box p \lif \Diamond \lnot p$} & \LR{$p$} د \LR{$\lnot\lnot p$} پر ځاے\\
  \end{derivation}
  
\end{proof}
\olpseganchor{OLP-0432-B020}{end}

\olpseganchor{OLP-0432-B021}{start}
په پورته اشتقاق کښې وروستے ګام «\LR{$p$} د \LR{$\lnot\lnot p$} پر ځاے»
د لاندې ګامونو لنډيز دے:
  \begin{derivation}
    & \LR{$\Log{K} \Proves \lnot \Box \lnot\lnot p \lif \Diamond \lnot
    p$}\\
    & \LR{$\Log{K} \Proves \lnot\lnot p \liff p$} & \Taut\\
    & \LR{$\Log{K} \Proves \lnot\Box p \lif \Diamond \lnot p$} & د \olref{prop:rewriting} له مخې
  \end{derivation}
په دې بېلګه کښې د \olref{prop:rewriting} د \LR{$!C(q)$}، \LR{$!A$} او \LR{$!B$}
ځايونه په ترتيب سره \LR{$\lnot \Box q \lif \Diamond \lnot p$}،
\LR{$\lnot\lnot p$} او~\LR{$p$} نيسي۔
\olpseganchor{OLP-0432-B021}{end}

\olpseganchor{OLP-0432-B022}{start}
که په فارمول کښې فرعي فارمول \LR{$\lnot\Diamond !A$} وي،
د \olref{prop:rewriting} په کارولو يې په \LR{$\Box\lnot !A$} بدلولے شو، ځکه
چې \LR{$\Log{K} \Proves \lnot\Diamond !A \liff \Box\lnot !A$}۔ دا او دې ته
ورته بدلونونه به په لنډ ډول «\LR{$\Box\lnot$} د
\LR{$\lnot\Diamond$} پر ځاے» نښه کوو۔
\olpseganchor{OLP-0432-B022}{end}

\olpseganchor{OLP-0432-B023}{start}
لاندې قضيه دا توجيه کوي چې د اشتقاق وړتيا پايلې په طرحه‌يي
ډول هم ثابتولے شو۔ د بېلګې په توګه، مخکنۍ قضيه يوازې
\LR{$\Log{K} \Proves \lnot\Box p \lif \Diamond \lnot p$} نۀ ثابتوي، بلکې
د هر~\LR{$!A$} دپاره \LR{$\Log{K} \Proves \lnot\Box !A \lif \Diamond
\lnot !A$} هم ثابتوي۔
\olpseganchor{OLP-0432-B023}{end}

\olpseganchor{OLP-0432-B024}{start}
\begin{prop}
  که \LR{$!A$} د \LR{$!B$} د ځايناستۍ يو مثال وي او \LR{$\Log{K} \Proves !B$}،
  نو \LR{$\Log{K} \Proves !A$}۔
\end{prop}
\olpseganchor{OLP-0432-B024}{end}

\olpseganchor{OLP-0432-B025}{start}
\begin{proof}
  د~\LR{$!B$} د اشتقاق پر اوږدوالي په استقرا دا ثابتول ستړي
  کوونکي خو معمولي کار دے چې پر ټول اشتقاق ځايناستي تطبيقول
  بيا هم سم اشتقاق ورکوي۔ په ځانګړي ډول، د تاوتولوژيکو
  مثالونو د ځايناستۍ مثالونه بيا هم تاوتولوژيک مثالونه وي؛ د
  \Dual{} او~\Ax{K} د مثالونو د ځايناستۍ مثالونه
  بيا هم د~\Dual{}
    او~\Ax{K} مثالونه وي؛ او په مقدمې (يا مقدمو) او نتيجه دواړو کښې
  د بياني متغيرونه پر ځاے د فارمولونه په اېښودلو سره
  د \MP{} او~\Nec{} کارول سم پاتې کېږي۔
\end{proof}
\olpseganchor{OLP-0432-B025}{end}

\olpseganchor{OLP-0433-B004}{start}
\olfileid{nml}{prf}{mpr}
\olpseganchor{OLP-0433-B004}{end}

\olpseganchor{OLP-0433-B005}{start}
\olsection{په \texorpdfstring{\Log{K}}{K} کښې نور ثبوتونه}
\olpseganchor{OLP-0433-B005}{end}

\olpseganchor{OLP-0433-B006}{start}
راځئ په~\Log{K} کښې د اشتقاق وړتيا څو نورې بېلګې وګورو؛ دا ځل
د \olref[der]{sec} ساده شوې طريقه کاروو۔
\olpseganchor{OLP-0433-B006}{end}

\olpseganchor{OLP-0433-B007}{start}
\begin{prop}
  \LR{$\Log{K} \Proves \Box (!A \lif !B) \lif (\Diamond !A \lif \Diamond
  !B)$}
\end{prop}
\olpseganchor{OLP-0433-B007}{end}

\olpseganchor{OLP-0433-B008}{start}
\begin{proof}
\begin{derivation}
1. & \LR{$\Log{K} \Proves (!A \lif !B) \lif (\lnot !B \lif \lnot !A)$} & \PL \\
2. &  \LR{$\Log{K} \Proves \Box(!A \lif !B) \lif (\Box\lnot!B \lif \Box\lnot!A)$} & \RK, 1 \\
3. & \LR{$\Log{K} \Proves (\Box\lnot!B \lif \Box\lnot!A) \lif (\lnot
\Box\lnot!A \lif \lnot\Box\lnot!B)$} & \Taut\\
4. & \LR{$\Log{K} \Proves \Box(!A \lif !B) \lif (\lnot
\Box\lnot!A \lif \lnot\Box\lnot!B)$} & \PL, 2, 3 \\
5. &  \LR{$\Log{K} \Proves \Box(!A \lif !B) \lif (\Diamond!A \lif \Diamond!B)$} & \LR{$\Diamond$} د \LR{$\lnot\Box\lnot$} پر ځاے.
\end{derivation}
\end{proof}
\olpseganchor{OLP-0433-B008}{end}
  
\olpseganchor{OLP-0433-B009}{start}
\begin{prop}
\LR{$\Log{K}\Proves \Box!A \lif (\Diamond(!A \lif !B) \lif
  \Diamond !B)$}
\end{prop}
\olpseganchor{OLP-0433-B009}{end}

\olpseganchor{OLP-0433-B010}{start}
\begin{proof}
  \begin{derivation}
  1. & \LR{$\Log{K} \Proves !A \lif (\lnot!B \lif \lnot (!A \lif !B))$} & \Taut \\
  2. & \LR{$\Log{K} \Proves \Box!A \lif  (\Box\lnot!B \lif \Box\lnot (!A\lif !B))$} &
  \RK, 1 \\
  3. &  \LR{$\Log{K} \Proves \Box!A \lif (\lnot \Box\lnot (!A\lif !B) \lif \lnot
  \Box\lnot!B)$} & \PL, 2 \\
  4. & \LR{$\Log{K} \Proves \Box!A \lif (\Diamond(!A \lif !B) \lif
  \Diamond !B)$} & \LR{$\Diamond$} د \LR{$\lnot\Box\lnot$} پر ځاے.
  \end{derivation}
\end{proof}
\olpseganchor{OLP-0433-B010}{end}

\olpseganchor{OLP-0433-B011}{start}
\begin{prop}
  \LR{$\Log{K} \Proves (\Diamond!A \lor \Diamond!B) \lif \Diamond(!A \lor !B)$}
\end{prop}
\olpseganchor{OLP-0433-B011}{end}

\olpseganchor{OLP-0433-B012}{start}
\begin{proof}
  \begin{derivation}
    1. & \LR{$\Log{K} \Proves \lnot(!A \lor !B) \lif \lnot!A$} & \Taut \\
    2. & \LR{$\Log{K} \Proves \Box\lnot(!A \lor !B) \lif \Box\lnot!A$} & \RK, 1 \\
    3. & \LR{$\Log{K} \Proves \lnot\Box\lnot !A \lif \lnot\Box\lnot(!A \lor!B)$} &
    \PL, 2\\
    4. & \LR{$\Log{K} \Proves \Diamond!A \lif \Diamond(!A \lor !B)$} &
    \LR{$\Diamond$} د \LR{$\lnot\Box\lnot$} پر ځاے\\
    5. & \LR{$\Log{K} \Proves \Diamond!B \lif \Diamond(!A \lor !B)$} & په همدې ډول\\
    6. & \LR{$\Log{K} \Proves (\Diamond!A \lor\Diamond!B) \lif \Diamond(!A \lor !B)$}
    & \PL, 4, 5.
  \end{derivation}
\end{proof}
\olpseganchor{OLP-0433-B012}{end}

\olpseganchor{OLP-0433-B013}{start}
\begin{prop}
  \LR{$\Log{K} \Proves \Diamond(!A \lor!B) \lif (\Diamond!A \lor \Diamond!B)$}
\end{prop}
\olpseganchor{OLP-0433-B013}{end}

\olpseganchor{OLP-0433-B014}{start}
\begin{proof}
  \begin{derivation}
    1. & \LR{$\Log{K} \Proves \lnot !A \lif (\lnot !B \lif \lnot (!A \lor !B))$} & \Taut \\
    2. & \LR{$\Log{K} \Proves \Box\lnot !A \lif
    (\Box\lnot!B \lif \Box \lnot (!A \lor !B))$} & \RK, 1\\
    3. & \LR{$\Log{K} \Proves \Box\lnot !A \lif (\lnot \Box \lnot (!A \lor!B)
    \lif \lnot\Box\lnot !B)$} & \PL, 2\\
    4. & \LR{$\Log{K} \Proves \lnot \Box \lnot(!A \lor !B) \lif (\Box \lnot !A \lif 
    \lnot\Box\lnot !B)$} & \PL, 3\\
    5. & \LR{$\Log{K} \Proves \lnot \Box \lnot(!A\lor!B) \lif (\lnot
    \lnot\Box\lnot!B \lif \lnot\Box\lnot!A)$} & \PL, 4\\
    6. & \LR{$\Log{K} \Proves \Diamond(!A \lor !B) \lif (\lnot
    \Diamond!B \lif \Diamond!A)$} & \LR{$\Diamond$} د \LR{$\lnot\Box\lnot$} پر ځاے\\
    7. & \LR{$\Log{K} \Proves \Diamond(!A\lor!B) \lif (\Diamond!A \lor \Diamond!B)$} & \PL, 6. \\
  \end{derivation}
\end{proof}
\olpseganchor{OLP-0433-B014}{end}

\olpseganchor{OLP-0433-B015}{start}
\begin{prob}
  ثابته کړئ چې لاندې د اشتقاق وړتيا ادعاوې سمې دي:
  \begin{enumerate}
  \item \LR{$\Log{K} \Proves \Diamond \lnot \lfalse \lif (\Box !A \lif
    \Diamond !A)$};
  \item \LR{$\Log{K} \Proves \Box(!A \lor !B) \lif (\Diamond !A \lor \Box
    !B)$};
  \item \LR{$\Log{K} \Proves (\Diamond !A \lif \Box !B) \lif \Box(!A \lif
    !B)$}.
  \end{enumerate}
\end{prob}
\olpseganchor{OLP-0433-B015}{end}

\olpseganchor{OLP-0434-B004}{start}
\olfileid{nml}{prf}{dua}
\olpseganchor{OLP-0434-B004}{end}

\olpseganchor{OLP-0434-B005}{start}
\olsection{دوه‌ګوني فارمولونه}
\olpseganchor{OLP-0434-B005}{end}

\olpseganchor{OLP-0434-B006}{start}
\begin{defn}\ollabel{def:duals}
  د \Ax{T}، \Ax{B}، \Ax{4} او \Ax{5} له فارمولونه څخه هر يو
  \emph{دوه‌ګونے} لري؛ د دوه‌ګوني د ښودلو لپاره د امکان نښه
  د لاندېني شاخص په توګه ليکل کېږي:
  \begin{align}
    \tag{\Ax{T_\Diamond}} p & \lif \Diamond p\\
    \tag{\Ax{B_\Diamond}} \Diamond\Box p & \lif p\\
    \tag{\Ax{4_\Diamond}} \Diamond\Diamond p & \lif \Diamond p\\
    \tag{\Ax{5_\Diamond}} \Diamond\Box p & \lif \Box p
    \end{align}
\end{defn}
\olpseganchor{OLP-0434-B006}{end}

\olpseganchor{OLP-0434-B007}{start}
له پورتنيو دوه‌ګونو فارمولونه څخه هر يو له اړوند فارمول
څخه داسې لاس ته راځي: لومړے د \LR{$p$} پر ځاے \LR{$\lnot p$} کېږدو، بيا
مقابل نقيض واخلو، وروسته \LR{$\lnot\Box\lnot$} په \LR{$\Diamond$} او
\LR{$\lnot\Diamond\lnot$} په~\LR{$\Box$} بدل کړو۔ \Ax{D}، يعنې
\LR{$\Box!A \lif \Diamond!A$}، په همدې معنا خپله دوه‌ګونی دے۔
\olpseganchor{OLP-0434-B007}{end}

\olpseganchor{OLP-0434-B008}{start}
\begin{prop}\ollabel{prop:dualsys}
  په \olref[nml][prf][dua]{def:duals} کښې د هر فارمول~\LR{$!A$} دپاره:
  \LR{$\Log{K}!A = \Log{K}!A_{\Diamond}$}.
\end{prop}
\begin{proof}
  ثبوت په لاندې تمرين کښې غوښتل شوے دے۔
\end{proof}
\begin{prob}
  \olref[nml][prf][dua]{prop:dualsys} ثابته کړئ۔
\end{prob}
\olpseganchor{OLP-0434-B008}{end}

\olpseganchor{OLP-0435-B004}{start}
\olfileid{nml}{prf}{prs}
\olpseganchor{OLP-0435-B004}{end}

\olpseganchor{OLP-0435-B005}{start}
\olsection{په مودالي نظامونو کښې ثبوتونه}
\olpseganchor{OLP-0435-B005}{end}

\olpseganchor{OLP-0435-B006}{start}
اوس له~\Log{K} پرته په نورو مودالي نظامونو کښې ثبوتونه څېړو۔
\olpseganchor{OLP-0435-B006}{end}

\olpseganchor{OLP-0435-B007}{start}
\begin{prop}\ollabel{prop:S5facts}
لاندې ادعاوې په اړوندو نظامونو کښې ثابتېږي:
  \begin{enumerate}
  \item \LR{$\Log{KT5} \Proves \Ax{B}$};
  \item \LR{$\Log{KT5} \Proves \Ax{4}$};
  \item \LR{$\Log{KDB4} \Proves \Ax{T}$};
  \item \LR{$\Log{KB4} \Proves \Ax{5}$};
  \item \LR{$\Log{KB5} \Proves \Ax{4}$};
  \item \ollabel{prop:S5facts-KT-D}\LR{$\Log{KT} \Proves \Ax{D}$}.
  \end{enumerate}
\end{prop}
\olpseganchor{OLP-0435-B007}{end}

\olpseganchor{OLP-0435-B008}{start}
\begin{proof}
  د هرې ادعا ثبوت ښيو۔
  \begin{enumerate}
    \item \LR{$\Log{KT5} \Proves \Ax{B}$}:
      \begin{derivation}
        1. & \LR{$\Log{KT5} \Proves \Diamond!A \lif \Box\Diamond!A$} & \Ax{5}\\
        2. & \LR{$\Log{KT5} \Proves !A \lif \Diamond!A$} & \LR{$\Ax{T_\Diamond}$}\\
        3. & \LR{$\Log{KT5} \Proves !A \lif \Box\Diamond!A$} & \PL, 2, 1.
      \end{derivation}
    \item \LR{$\Log{KT5} \Proves \Ax{4}$}:
      \begin{derivation}
        1. &\LR{$\Log{KT5} \Proves \Diamond\Box!A \lif \Box\Diamond\Box!A$} & \Ax{5}
        چې \LR{$\Box!A$} د \LR{$p$} پر ځاے کارېدلے دے\\
        2. & \LR{$\Log{KT5} \Proves \Box!A \lif \Diamond\Box!A$} & \LR{$\Ax{T_\Diamond}$}
        چې \LR{$\Box!A$} د \LR{$p$} پر ځاے کارېدلے دے\\
        3. & \LR{$\Log{KT5} \Proves \Box!A \lif \Box\Diamond\Box!A$} & \PL, 2, 1\\
        4. & \LR{$\Log{KT5} \Proves \Diamond\Box!A \lif \Box!A$} & \LR{$\Ax{5_\Diamond}$}\\
        5. & \LR{$\Log{KT5} \Proves \Box\Diamond\Box!A \lif \Box\Box!A$} & \RK{}, 4 \\
        6. & \LR{$\Log{KT5} \Proves \Box!A \lif \Box\Box!A$} & \PL, 3, 5. \\
      \end{derivation}
    \item \LR{$\Log{KDB4} \Proves \Ax{T}$}:
      \begin{derivation}
        1. & \LR{$\Log{KDB4} \Proves \Diamond\Box!A \lif !A $} & \LR{$\Ax{B_\Diamond}$} \\
        2. & \LR{$\Log{KDB4} \Proves \Box\Box!A \lif \Diamond\Box!A$} & \LR{$\Ax{D}$}
        چې \LR{$\Box!A$} د \LR{$p$} پر ځاے کارېدلے دے\\
        3. & \LR{$\Log{KDB4} \Proves \Box\Box!A \lif !A$} & \PL, 1, 2\\
        4. & \LR{$\Log{KDB4} \Proves \Box!A \lif \Box\Box!A$} & \Ax{4} \\
        5. & \LR{$\Log{KDB4} \Proves \Box!A \lif !A$} & \PL, 4, 3. \\
      \end{derivation}
    \item \LR{$\Log{KB4} \Proves \Ax{5}$}:
      \begin{derivation}
        1. & \LR{$\Log{KB4} \Proves \Diamond!A \lif \Box \Diamond\Diamond!A$} & \Ax{B}
        چې \LR{$\Diamond!A$} د \LR{$p$} پر ځاے کارېدلے دے\\
        2. & \LR{$\Log{KB4} \Proves \Diamond\Diamond !A \lif \Diamond !A$} &
        \LR{$\Ax{4_\Diamond}$} \\
        3. & \LR{$\Log{KB4} \Proves \Box\Diamond\Diamond !A \lif \Box \Diamond!A$} &
        \RK{}, 2 \\
        4. & \LR{$\Log{KB4} \Proves \Diamond!A \lif \Box\Diamond!A$} & \PL, 1, 3.
      \end{derivation}
    \item \LR{$\Log{KB5} \Proves \Ax{4}$}:
      \begin{derivation}
        1. & \LR{$\Log{KB5} \Proves \Box!A \lif \Box\Diamond\Box !A$} & \Ax{B}
        چې \LR{$\Box!A$} د \LR{$p$} پر ځاے کارېدلے دے \\
        2. & \LR{$\Log{KB5} \Proves  \Diamond\Box!A \lif \Box !A$} &
        \LR{$\Ax{5_\Diamond}$}\\
        3. & \LR{$\Log{KB5} \Proves  \Box\Diamond\Box!A \lif \Box\Box !A$} & \RK{}, 2 \\
        4. & \LR{$\Log{KB5} \Proves  \Box!A \lif \Box\Box!A$} & \PL, 1, 3.
      \end{derivation}
    \item \LR{$\Log{KT} \Proves \Ax{D}$}:
      \begin{derivation}
        1. & \LR{$\Log{KT} \Proves \Box !A \lif !A$} & \Ax{T} \\
        2. & \LR{$\Log{KT} \Proves !A \lif \Diamond!A$} & \LR{$\Ax{T_\Diamond}$} \\
        3. & \LR{$\Log{KT} \Proves \Box !A \lif \Diamond !A$} &  \PL, 1, 2
      \end{derivation}
  \end{enumerate}
\end{proof}
\olpseganchor{OLP-0435-B008}{end}

\olpseganchor{OLP-0435-B009}{start}
\begin{defn}
  د دود له مخې، \Log{S4} د \Log{KT4} نظام او
  \Log{S5} د \Log{KTB4} نظام بولو۔
\end{defn}
\olpseganchor{OLP-0435-B009}{end}

\olpseganchor{OLP-0435-B010}{start}
لاندې قضيه ښيي چې کلاسيک نظام \Log{S5} څو يو له بل سره معادل
بديهي بنسټونه لري۔ د بديهي اصلونو دغه بېلابېلې ټولګې د هم ارزښتۍ
اړيکې ټاکي
(\olref[frd][es5]{prop:equivalences} وګورئ)۔
\olpseganchor{OLP-0435-B010}{end}

\olpseganchor{OLP-0435-B011}{start}
\begin{prop}\ollabel{prop:S5}
  \LR{$\Log{KTB4} = \Log{KT5} = \Log{KDB4} = \Log{KDB5}$}.
\end{prop}
\olpseganchor{OLP-0435-B011}{end}

\olpseganchor{OLP-0435-B012}{start}
\begin{proof}
  ثبوت په لاندې تمرين کښې غوښتل شوے دے۔
\end{proof}
\olpseganchor{OLP-0435-B012}{end}

\olpseganchor{OLP-0435-B013}{start}
\begin{prob}
  \olref[nml][prf][prs]{prop:S5} ثابته کړئ۔
\end{prob}
\olpseganchor{OLP-0435-B013}{end}

\olpseganchor{OLP-0436-B004}{start}
\olfileid{nml}{prf}{snd}
\olpseganchor{OLP-0436-B004}{end}

\olpseganchor{OLP-0436-B005}{start}
\olsection{سموالے}
\olpseganchor{OLP-0436-B005}{end}

\olpseganchor{OLP-0436-B006}{start}
اشتقاق نظام هله سم بلل کېږي چې هر هغه څه چې ترې
اشتقاق کېدلے شي، معتبر وي۔ د مودالي نظامونو په پام کښې نيولو سره،
يعنې داسې اشتقاقونه چې د \Ax{K} ترڅنګ د ځينو فارمولونه
\LR{$!A_1$}، \dots،~\LR{$!A_n$} بېلګې هم پکې کارولے شو، غواړو هر د اشتقاق وړ
فارمول په هر هغه مدل کښې رښتيا وي چې د \LR{$!A_1$}، \dots، \LR{$!A_n$}
د ځايناستۍ هره بېلګه پکې رښتيا وي۔
\olpseganchor{OLP-0436-B006}{end}

\olpseganchor{OLP-0436-B007}{start}
\begin{thm}[د سموالي قضيه]\ollabel{thm:soundness}
  که د \LR{$!A_1$}، \dots، \LR{$!A_n$} هره بېلګه په ترتيب سره د مدلونو په
  ټولګيو \LR{$\mClass{C}_1$}، \dots، \LR{$\mClass{C}_n$} کښې معتبره وي،
  نو له \LR{$\Log{K}!A_1\dots !A_n
  \Proves !B$} څخه لازمېږي چې \LR{$!B$} د مدلونو په ټولګي
  \LR{$\mClass{C}_1 \cap \dots \cap \mClass{C}_n$} کښې معتبر وي۔
\end{thm}
\olpseganchor{OLP-0436-B007}{end}

\olpseganchor{OLP-0436-B008}{start}
\begin{proof}
  د ثبوتونو پر اوږدوالي استقرا کوو۔ د لنډون دپاره \LR{$\mClass{C} =
  \mClass{C}_1 \cap \dots \cap \mClass{C}_n$} وټاکئ۔
  \begin{enumerate}
  \item د استقرا بنسټ: که د \LR{$!B$} د ثبوت اوږدوالے~\LR{$1$} وي، نو دا يا
    د همېش‌رښتيا فارمول يوه بېلګه وي، يا د~\Ax{K} يوه بېلګه،
     يا د~\Dual{} يوه بېلګه، يا د
    \LR{$!A_1$}، \dots،~\LR{$!A_n$} له ډلې د يوه فارمول بېلګه وي۔ په لومړي
    حالت کښې \LR{$!B$} په \LR{$\mClass{C}$} کښې معتبر وي، ځکه د همېش‌رښتيا
    فارمول بېلګې د مدلونو په \emph{هره} ټولګي کښې معتبرې دي؛
    \olref[syn][tau]{prop:valid-taut} وګورئ۔ په دويم حالت کښې هم
    د \olref[syn][sch]{prop:Kvalid} او
      \olref[syn][sch]{prop:Dual-valid} له مخې همدا پايله راځي۔
    په درېيم حالت کښې، څرنګه چې \LR{$!B$} په \LR{$\mClass{C}_i$} کښې معتبر دے
    او \LR{$\mClass{C} \subseteq
    \mClass{C}_i$}، نو د \olref[syn][val]{prop:subset-class} له مخې
    \LR{$!B$} په \LR{$\mClass{C}$} کښې هم معتبر دے۔
  \item د استقرا ګام: فرض کړئ \LR{$!B$} د \LR{$k>1$} اوږدوالي ثبوت لري۔ که
    \LR{$!B$} د همېش‌رښتيا فارمول يوه بېلګه يا د \LR{$!A_1$}،
    \dots، \LR{$!A_n$} له ډلې د يوه بېلګه وي، يا د بنسټيزو مودالي
    اصلونو يوه بېلګه وي، د تېر ګام په شان استدلال کوو۔ که وروستۍ
    کرښه د مودوس پوننس له مخې وي، نو \LR{$!B$} د پخوانيو فارمولونه \LR{$!C \lif !B$} او
    \LR{$!C$} څخه د \MP{} په واسطه لاس ته راځي۔ د \LR{$!C \lif !B$} او
    \LR{$!C$} د ثبوتونو اوږدوالے \LR{$<k$} دے، او د استقرا د فرضيې له مخې
    دواړه په~\LR{$\mClass{C}$} کښې معتبر دي۔ د
    \olref[syn][sch]{prop:soundMP} له مخې \LR{$!B$} هم په \LR{$\mClass{C}$}
    کښې معتبر دے۔ په پاے کښې فرض کړئ \LR{$!B$} له \LR{$!C$} څخه د \Nec{} په
    واسطه لاس ته راغلے وي (نو \LR{$!B = \Box!C$})۔ د استقرا د فرضيې له
    مخې \LR{$!C$} په \LR{$\mClass{C}$} کښې معتبر دے، او د
    \olref[syn][val]{prop:Nec-rule} له مخې~\LR{$!B$} هم معتبر دے۔
  \end{enumerate}
\end{proof}
\olpseganchor{OLP-0436-B008}{end}

\olpseganchor{OLP-0437-B004}{start}
\olfileid{nml}{prf}{dis}
\olpseganchor{OLP-0437-B004}{end}

\olpseganchor{OLP-0437-B005}{start}
\olsection{د نظامونو بېلوالے ښودل}
\olpseganchor{OLP-0437-B005}{end}

\olpseganchor{OLP-0437-B006}{start}
په \olref[prs]{sec} کښې مو وليدل چې د مودالي منطق دوه نظامونه
په حقيقت کښې څنګه يو شان ثابتولے شو۔ \olref[snd]{thm:soundness} مونږ ته
دا هم راښيي چې دوه مودالي نظامونه \LR{$\Sigma$} او \LR{$\Sigma'$} بېل دي: داسې
فارمول~\LR{$!A$} پيدا کوو چې \LR{$\Sigma' \Proves !A$} وي، خو د~\LR{$\Sigma$}
په يوه مدل کښې ناسم وي۔
\olpseganchor{OLP-0437-B006}{end}

\olpseganchor{OLP-0437-B007}{start}
\begin{prop}
  \LR{$\Log{KD} \subsetneq \Log{KT}$}
\end{prop}
\olpseganchor{OLP-0437-B007}{end}

\olpseganchor{OLP-0437-B008}{start}
\begin{proof} دا د هغه معنايي حقيقت نحوي سيال دے چې هره
  انعکاسي اړيکه مسلسله ده۔ د \LR{$\Log{KD} \subseteq
  \Log{KT}$} د ښودلو دپاره بايد ووينو چې له \LR{$\Log{KD} \Proves !B$}
  څخه \LR{$\Log{KT} \Proves !B$} لازمېږي۔ دا له \LR{$\Log{KT} \Proves
  \Ax{D}$} څخه راځي؛ د \olref[prs]{prop:S5facts} د
  \olref[prs]{prop:S5facts-KT-D} ادعا وګورئ۔ د دې د ښودلو دپاره
  چې شموليت سخت دے (مساوات نۀ دے)، د سموالي له مخې
  (\olref[snd]{thm:soundness}) بس ده چې د \Log{KD} داسې مدل وښيو
  چې \Ax{T}، يعنې \LR{$\Box p \lif p$}، پکې ناسم وي (دا اسانه کار
  د تمرين په توګه پرېښودل شوے)؛ ځکه بيا د سموالي له مخې \LR{$\Log{KD}
  \Proves/ \Box p \lif p$}.
\end{proof}
\olpseganchor{OLP-0437-B008}{end}

\olpseganchor{OLP-0437-B009}{start}
\begin{prop}
  \LR{$\Log{KB} \neq \Log{K4}$}. 
\end{prop}
\olpseganchor{OLP-0437-B009}{end}

\olpseganchor{OLP-0437-B010}{start}
\begin{proof}
  داسې متناظر مدل جوړوو چې د \Ax{4} يوه بېلګه پکې ناسمه وي؛ ځکه دا
  بېلګه په \Log{K4} کښې په څرګنده د اشتقاق وړ ده، خو په \Log{KB}
  کښې نۀ ده، نو \LR{$\Log{K4} \nsubseteq \Log{KB}$} به ترې لازم شي۔
  د \olref{fig:Bnot4} متناظر مدل~\LR{$\mModel{M}$} وګورئ۔ د
  \Ax{K} هره بېلګه په هر مدل کښې رښتيا ده، او د~\Ax{B} هره بېلګه
  د تناظر له کبله په~\LR{$\mModel{M}$} کښې رښتيا ده
  (په ترتيب سره د \olref[syn][sch]{prop:Kvalid} او
  \olref[frd][acc]{thm:soundschemas} له مخې)۔ له لومړۍ نړۍ يوازې
  دويمې نړۍ ته لاسرسے شته او هلته متغير رښتيا دے۔ خو له دويمې
  نړۍ بېرته لومړۍ ته لاسرسے شته، چې متغير پکې نارښتيا دے۔ نو
  \LR{$\mSat/{M}{\Box p \lif \Box\Box p}[w_1]$}۔
\end{proof}
\olpseganchor{OLP-0437-B010}{end}

\olpseganchor{OLP-0437-B011}{start}
\begin{figure}[htpb]
  \centering
  \begin{tikzpicture}[modal]
    \node[world] (w1) [label=above:\mFalse{p},
    label={below:\LR{$\mSat{{}}{\Box p}$}\\ \LR{$\mSat/{{}}{\Box\Box p}$}}] {\LR{$w_1$}}; 
    \node[world] (w2) [label=above:\mTrue{p},
      label=below:\LR{$\mSat/{{}}{\Box p}$}, right=of w1]{\LR{$w_2$}}; 
    \draw[->, bend left] (w1) to (w2); 
    \draw[->, bend left] (w2) to (w1); 
  \end{tikzpicture}
  \caption{يو متناظر مدل چې د \Ax{4} يوه بېلګه پکې ناسمه ده۔}
  \ollabel{fig:Bnot4}
\end{figure}
\olpseganchor{OLP-0437-B011}{end}

\olpseganchor{OLP-0437-B012}{start}
\begin{thm}\ollabel{thm:KTBnot45}
  \LR{$\Log{KTB} \Proves/ \Ax{4}$} او \LR{$\Log{KTB} \Proves/ \Ax{5}$}.
\end{thm}
\olpseganchor{OLP-0437-B012}{end}

\olpseganchor{OLP-0437-B013}{start}
\begin{proof}
  د \olref[frd][acc]{thm:soundschemas} له مخې پوهېږو چې د \Ax{T} او
  \Ax{B} ټولې بېلګې په هر انعکاسي او متناظر مدل کښې په ترتيب سره
  رښتيا دي۔ نو د سموالي له مخې بس ده چې يو انعکاسي متناظر مدل پيدا
  کړو چې په يوه نړۍ کښې د~\Ax{4} يوه بېلګه ناسمه وي، او د~\Ax{5}
  دپاره هم همداسې وکړو۔ د دواړو دعوو دپاره هماغه يو مدل کاروو۔ په
  \olref{fig:KTBnot45} کښې متناظر او انعکاسي مدل وګورئ۔ هلته
  \LR{$\mSat/{M}{\Box p \to \Box \Box
    p}[w_1]$}، نو \Ax{4} په \LR{$w_1$} کښې ناسم دے۔ همدارنګه
  \LR{$\mSat/{M}{\Diamond
    \lnot p \lif \Box \Diamond \lnot p}[w_2]$}، نو د~\Ax{5} هغه
  بېلګه چې \LR{$!A = \lnot p$} وي، په~\LR{$w_2$} کښې ناسمه ده۔
  که د پنځم اصل ځانګړونکے فارمول استنباط‌وړ وے، نو د يو شان
  ځايناستۍ له مخې به دا بېلګه هم استنباط‌وړ وه۔
\end{proof}
\olpseganchor{OLP-0437-B013}{end}

\olpseganchor{OLP-0437-B014}{start}
\begin{figure}[htpb]
  \centering
  \begin{tikzpicture}[modal]
    \node[world] (w1) [label={right:\mTrue{p}},
      label={below: \LR{$\mSat{{}}{\Box p}$}\\
        \LR{$\mSat/{{}}{\Box\Box p}$}\\
        \LR{$\mSat/{{}}{\Diamond\lnot p}$}}] {\LR{$w_1$}};
    \draw[reflexive above] (w1) to (w1);
    \node[world] (w2) [label={right:\mTrue{p}}, label={below:
        \LR{$\mSat{{}}{\Diamond\lnot p}$}\\
        \LR{$\mSat/{{}}{\Box\Diamond\lnot p}$}}, right=of w1] {\LR{$w_2$}} ;
    \draw[reflexive above] (w2) to (w2);
    \node[world] (w3) [label={right:\mFalse{p}},right=of w2] {\LR{$w_3$}};
    \draw[reflexive above] (w3) to (w3);
    \draw[->, bend left] (w1) to (w2);  
    \draw[->, bend left] (w2) to (w3);  
    \draw[->, bend left] (w3) to (w2);  
    \draw[->, bend left] (w2) to (w1);  
  \end{tikzpicture}
  \caption{د \olref{thm:KTBnot45} دپاره مدل۔}
  \ollabel{fig:KTBnot45}
\end{figure}
\olpseganchor{OLP-0437-B014}{end}

\olpseganchor{OLP-0437-B015}{start}
\begin{thm}\ollabel{thm:KD5not4}
  \LR{$\Log{KD5} \neq \Log{KT4} = \Log{S4}$}.
\end{thm}
\olpseganchor{OLP-0437-B015}{end}

\olpseganchor{OLP-0437-B016}{start}
\begin{proof}
  د \olref[frd][acc]{thm:soundschemas} له مخې پوهېږو چې د \Ax{D}
  او~\Ax{5} ټولې بېلګې په ټولو مسلسلو اقليدسي مدلونو کښې رښتيا دي۔
  نو بس ده چې يو مسلسل اقليدسي مدل پيدا کړو چې په يوه نړۍ کښې د~\Ax{4}
  يوه بېلګه ناسمه وي۔ د \olref{fig:KD5not4} مدل وګورئ او پام وکړئ
  چې \LR{$\mSat/{M}{\Box p
  \lif \Box\Box p}[w_1]$}.
\end{proof}
\olpseganchor{OLP-0437-B016}{end}

\olpseganchor{OLP-0437-B017}{start}
\begin{figure}[t]
  \centering
  \begin{tikzpicture}[modal]
    \node[world] (w2) [label=north west:\mTrue{p}]{\LR{$w_2$}} ;
    \draw[reflexive left] (w2) to (w2);
    \node[world] (w1) [label=right:\mFalse{p},
      label=below:{\LR{$\mSat{{}}{\Box p}, \mSat/{{}}{\Box\Box p}$}},
      below right=of w2]{\LR{$w_1$}};
    \node[world] (w3) [label=north east:\mTrue{p}, above right=of w1]{\LR{$w_3$}} ;
    \draw[reflexive right] (w3) to (w3);
    \node[world] (w4) [label=right:\mFalse{p},above right=of w2]{\LR{$w_4$}};
    \draw[reflexive above] (w4) to (w4);
    \draw[->] (w1) to (w2);  
    \draw[->] (w1) to (w3);  
    \draw[->, bend left=15] (w2) to (w3);
    \draw[->, bend left=15] (w3) to (w2);  
    \draw[->, bend left=15] (w2) to (w4);
    \draw[->, bend left=15] (w4) to (w2);  
    \draw[->, bend left=15] (w3) to (w4);
    \draw[->, bend left=15] (w4) to (w3);  
  \end{tikzpicture}
  \caption{د \olref{thm:KD5not4} دپاره مدل۔}
  \ollabel{fig:KD5not4}
\end{figure}
\olpseganchor{OLP-0437-B017}{end}

\olpseganchor{OLP-0437-B018}{start}
\begin{prob}
  د \olref[nml][prf][dis]{thm:KD5not4} يو بل ثبوت د \LR{$3$} نړۍو
  لرونکي مدل په مرسته ورکړئ۔
\end{prob}
\olpseganchor{OLP-0437-B018}{end}

\olpseganchor{OLP-0437-B019}{start}
\begin{prob}
  يو داسې انعکاسي او انتقالي مدل وړاندې کړئ چې دواړه
  \LR{$\Log{KT4} \Proves/ \Ax{B}$} او \LR{$\Log{KT4} \Proves/ \Ax{5}$}
  وښيي۔
\end{prob}
\olpseganchor{OLP-0437-B019}{end}

\olpseganchor{OLP-0438-B004}{start}
\olfileid{nml}{prf}{prg}
\olpseganchor{OLP-0438-B004}{end}

\olpseganchor{OLP-0438-B005}{start}
\olsection{د فارمولونه له يوه سټ څخه د اشتقاق وړتيا}
\olpseganchor{OLP-0438-B005}{end}

\olpseganchor{OLP-0438-B006}{start}
په \olref[prf][prf]{sec} کښې مو په يوه نظام~\LR{$\Sigma$} کښې د
فارمول د ثابتېدو مفهوم تعريف کړے و۔ اوس دا مفهوم په~\LR{$\Sigma$}
کښې د فارمولونه له سټ~\LR{$\Gamma$} څخه ثابتېدو ته غځوو۔
\olpseganchor{OLP-0438-B006}{end}

\olpseganchor{OLP-0438-B007}{start}
\begin{defn}\ollabel{defn:Gammaproves}
  فارمول~\LR{$!A$} په نظام~\LR{$\Sigma$} کښې د فارمولونه له سټ~\LR{$\Gamma$} څخه
  (چې په \LR{$\Gamma \Proves[\Sigma] !A$} ليکل کېږي) هله او يوازې هله
  د اشتقاق وړ وي، چې داسې \LR{$!B_1$}, \dots,
  \LR{$!B_n \in \Gamma$} وي چې
  \LR{$\Sigma \Proves !B_1 \lif (!B_2 \lif \cdots (!B_n \lif !A)
  \cdots)$}۔
  د فرضيو دا لړۍ تشه هم کېدے شي؛ په دې حالت کښې فارمول په
  خپله د نظام يو بې‌فرضيې استنباط‌وړ فارمول دے۔
\end{defn}
\olpseganchor{OLP-0438-B007}{end}

\olpseganchor{OLP-0439-B004}{start}
\olfileid{nml}{prf}{prp}
\olpseganchor{OLP-0439-B004}{end}

\olpseganchor{OLP-0439-B005}{start}
\olsection{د اشتقاق وړتيا خاصيتونه}
\olpseganchor{OLP-0439-B005}{end}

\olpseganchor{OLP-0439-B006}{start}
\begin{prop}\ollabel{prop:derivabilityfacts}
  \LR{$\Sigma$} دې يو مودالي نظام او \LR{$\Gamma$} دې د مودالي
  فارمولونه يو سټ وي۔ لاندې خاصيتونه لري:
  \begin{enumerate}
  \item\ollabel{prop:derivabilityfacts-monotonicity} \emph{يکنواختي}: که
    \LR{$\Gamma \Proves[\Sigma] !A$} او \LR{$\Gamma \subseteq \Delta$} وي، نو
    \LR{$\Delta \Proves[\Sigma] !A$};
  \item\ollabel{prop:derivabilityfacts-reflexivity}
    \emph{انعکاسيت}: که \LR{$!A \in \Gamma$} وي، نو \LR{$\Gamma
    \Proves[\Sigma] !A$};
  \item\ollabel{prop:derivabilityfacts-cut} \emph{قطع}: که \LR{$\Gamma
    \Proves[\Sigma] !A$} او \LR{$\Delta \cup \{!A\} \Proves[\Sigma] !B$}
    وي، نو \LR{$\Gamma \cup \Delta \Proves[\Sigma] !B$};
  \item\ollabel{prop:derivabilityfacts-deduction} \emph{د استنتاج قضيه}: \LR{$\Gamma \cup \{!B\}
    \Proves[\Sigma] !A$} هله او يوازې هله چې \LR{$\Gamma \Proves[\Sigma] !B \lif
      !A$} وي؛
  \item \ollabel{prop:derivabilityfacts-ruleT}% \emph{Rule T}: If
    که \LR{$\Gamma \Proves[\Sigma] !A_1$} او \dots او \LR{$\Gamma
    \Proves[\Sigma] !A_n$} وي، او \LR{$!A_1 \to (!A_2 \lif \cdots (!A_n \lif
    !B)\cdots)$} د تاوتولوژۍ يوه بېلګه وي، نو \LR{$\Gamma
    \Proves[\Sigma] !B$}.
  \end{enumerate}
\end{prop}
\olpseganchor{OLP-0439-B006}{end}

\olpseganchor{OLP-0439-B007}{start}
د دې ثبوت يو اسان تمرين دے۔ د
\olref{prop:derivabilityfacts} د
\olref{prop:derivabilityfacts-ruleT} برخه مونږ ته دا راکوي چې، د بېلګې
په توګه، که \LR{$\Gamma \Proves[\Sigma] !A \lor !B$}
او \LR{$\Gamma \Proves[\Sigma] \lnot !A$} وي، نو \LR{$\Gamma \Proves[\Sigma]
!B$} وي۔ همدارنګه په راتلونکي کښې به د
\LR{$\Gamma \cup \{ !A\} \Proves[\Sigma] !B$} پر ځای
\LR{$\Gamma, !A \Proves[\Sigma] !B$} ليکو۔
\olpseganchor{OLP-0439-B007}{end}

\olpseganchor{OLP-0439-B008}{start}
\begin{defn}
  يو سټ~\LR{$\Gamma$} د نظام~\LR{$\Sigma$} په نسبت \emph{له اشتقاقي پلوه تړلے} هله
  او يوازې هله وي چې له \LR{$\Gamma \Proves[\Sigma] !A$} څخه
  \LR{$!A \in \Gamma$} لازم شي۔
\end{defn}
\olpseganchor{OLP-0439-B008}{end}

\olpseganchor{OLP-0440-B004}{start}
\olfileid{nml}{prf}{con}
\olpseganchor{OLP-0440-B004}{end}

\olpseganchor{OLP-0440-B005}{start}
\olsection{سازګاري}
\olpseganchor{OLP-0440-B005}{end}

\olpseganchor{OLP-0440-B006}{start}
سازګاري د فارمولونه د سټونو يو مهم خاصيت دے۔ د فارمولونه يو
سټ هغه وخت ناسازګار وي چې يو تناقض، لکه~\LR{$\lfalse$}، ترې د اشتقاق وړ
وي؛ او که داسې نۀ وي، سازګار وي۔ که يو سټ د يوه نظام په نسبت
ناسازګار وي، نو ټول فارمولونه ئې د هغه نظام په يوه مدل کښې په
يوه نړۍ يوځاے رښتيا کېدے نۀ شي۔ په داسې مدل کښې د نظام د ټولو
بديهي اصلونو هره بېلګه په هره نړۍ کښې رښتيا وي۔ د بشپړتيا د
قضيې دپاره سرچپه لور ثابتوو: د نظام په نسبت هر سازګار سټ د
همدغه نظام د کانوني مدل په يوه نړۍ کښې رښتيا وي۔
\olpseganchor{OLP-0440-B006}{end}

\olpseganchor{OLP-0440-B007}{start}
\begin{defn}
  يو سټ~\LR{$\Gamma$} د نظام~\LR{$\Sigma$} په نسبت \emph{سازګار}، يا په لنډه
  «\LR{$\Sigma$}-سازګار»، هله او يوازې هله وي چې \LR{$\Gamma
  \Proves/[\Sigma] \lfalse$}.
\end{defn}
\olpseganchor{OLP-0440-B007}{end}

\olpseganchor{OLP-0440-B008}{start}
نو د بېلګې په توګه، سټ \LR{$\{ \Box(p \lif q), \Box p, \lnot\Box q \}$}
د بياني منطق په نسبت سازګار دے، خو له \Log{K} سره سازګار نۀ دے۔
همدارنګه، سټ \LR{$\{ \Diamond p, \Box\Diamond
p \lif q, \lnot q \}$} له \Log{K5} سره سازګار نۀ دے۔
\olpseganchor{OLP-0440-B008}{end}

\olpseganchor{OLP-0440-B009}{start}
\begin{prop}\ollabel{prop:consistencyfacts}
  \LR{$\Gamma$} دې د فارمولونه يو سټ وي۔ بيا:
  \begin{enumerate}
  \item \LR{$\Gamma$} هله او يوازې هله د~\LR{$\Sigma$} په نسبت سازګار وي چې
    داسې کوم فارمول~\LR{$!A$} وي چې \LR{$\Gamma \Proves/[\Sigma]
    !A$}.
  \item \ollabel{prop:consistencyfacts-b}%
    \LR{$\Gamma \Proves[\Sigma] !A$} هله او يوازې هله وي چې \LR{$\Gamma \cup \{
    \lnot!A \}$} د~\LR{$\Sigma$} په نسبت سازګار نۀ وي۔
  \item \ollabel{prop:consistencyfacts-c}%
    که \LR{$\Gamma$} د~\LR{$\Sigma$} په نسبت سازګار وي، نو د هر فارمول
    \LR{$!A$} دپاره يا \LR{$\Gamma \cup \{ !A \}$} د~\LR{$\Sigma$} په نسبت
    سازګار وي يا \LR{$\Gamma \cup \{ \lnot!A \}$} د~\LR{$\Sigma$} په نسبت
    سازګار وي۔
  \end{enumerate}
\end{prop}
\olpseganchor{OLP-0440-B009}{end}

\olpseganchor{OLP-0440-B010}{start}
\begin{proof}
  دا حقيقتونه له کلاسيک بياني منطق څخه په اسانۍ لازمېږي۔ د
  \olref{prop:consistencyfacts-c} دپاره استدلال وړاندې کوو۔ د
  دعوې مقابل نقيض اخلو او فرض کوو چې نۀ \LR{$\Gamma \cup \{ !A \}$} او نۀ
  \LR{$\Gamma \cup \{ \lnot!A \}$} د~\LR{$\Sigma$} په نسبت سازګار دي۔ نو د
  ناسازګارۍ له تعريفه دواړه \LR{$\Gamma, !A \Proves[\Sigma]
  \lfalse$} او \LR{$\Gamma, \lnot !A \Proves[\Sigma] \lfalse$} دي۔ د
  استنتاج د قضيې له مخې \LR{$\Gamma \Proves[\Sigma] !A \to \lfalse$} او
  \LR{$\Gamma \Proves[\Sigma] \lnot!A \lif \lfalse$} دي۔ خو \LR{$(!A \lif
  \lfalse) \lif ((\lnot!A \lif \lfalse) \lif \lfalse)$} د
  تاوتولوژۍ يوه بېلګه ده؛ نو د \olref[prp]{prop:derivabilityfacts}
  د \olref[prp]{prop:derivabilityfacts-ruleT} برخې له مخې \LR{$\Gamma \Proves[\Sigma] \lfalse$}.
\end{proof}
\olpseganchor{OLP-0440-B010}{end}

\olpseganchor{OLP-0441-B004}{start}
\olchapter{nml}{com}{بشپړتيا او کانوني مدلونه}
\olpseganchor{OLP-0441-B004}{end}

\olpseganchor{OLP-0442-B004}{start}
\olfileid{nml}{com}{int}
\olpseganchor{OLP-0442-B004}{end}

\olpseganchor{OLP-0442-B005}{start}
\olsection{پېژندنه}
\olpseganchor{OLP-0442-B005}{end}

\olpseganchor{OLP-0442-B006}{start}
که \LR{$\Sigma$} يو مودالي نظام وي، نو د سموالي قضيه ثابتوي چې که
\LR{$\Sigma \Proves !A$} وي، \LR{$!A$} د مدلونو په هر هغه ټولګي
\LR{$\mClass{C}$} کښې معتبر دے چې د~\LR{$\Sigma$} د ټولو فارمولونه
ټولې بېلګې پکې معتبرې وي۔ په ځانګړې توګه، که \LR{$\Log{K}
\Proves !A$} وي، نو \LR{$!A$} په ټولو مدلونو کښې رښتيا وي؛ که
\LR{$\Log{KT} \Proves !A$} وي، نو \LR{$!A$} په ټولو انعکاسي مدلونو کښې
رښتيا وي؛ او که \LR{$\Log{KD} \Proves !A$} وي، نو \LR{$!A$} په ټولو
مسلسلو مدلونو کښې رښتيا وي، او داسې نور۔
\olpseganchor{OLP-0442-B006}{end}

\olpseganchor{OLP-0442-B007}{start}
بشپړتيا د سموالي سرچپه لور دے: د \Log{K} بشپړتيا، د بېلګې په
توګه، دا معنا لري چې که فارمول~\LR{$!A$} معتبر وي، نو
\LR{$\Proves !A$}۔ د بشپړتيا ثابتول د سموالي له ثابتولو ډېر سخت دي۔
لومړے ئې مقابل نقيض ته وګورئ: \Log{K} هله او يوازې هله بشپړ
دے چې هر کله \LR{$\Proves/ !A$} وي، يو مقابل مدل، يعنې داسې
مدل~\LR{$\mModel{M}$}، وي چې \LR{$\mSat/{M}{!A}$}۔ په همارز ډول،
(د~\LR{$!A$} په نفي کولو) دا ثابتولے شو چې هر کله
\LR{$\Proves/ \lnot !A$} وي، د~\LR{$!A$} يو مدل شته۔ د داسې مدل په جوړولو
کښې په~\LR{$!A$} کښې له شته معلوماتو ګټه اخستلے شو۔ د ځانګړو
فارمولونه دپاره د مدلونو په موندلو کښې هم ډېر ځله همدا کوو:
د بېلګې په توګه، که د \LR{$p \lif \Box q$} دپاره مقابل مدل غواړو،
نو پوهېږو چې په هغه کښې بايد يوه نړۍ وي چې \LR{$p$} پکې رښتيا او
\LR{$\Box q$} پکې دروغ وي۔ او که \LR{$\Box q$} په يوه نړۍ کښې دروغ وي،
نو له هغې څخه د لاسرسي وړ داسې نړۍ بايد وي چې \LR{$q$} پکې دروغ
وي۔ همدومره پوهېدل بس دي: کومې نړۍوې کوم بياني متغيرونه
رښتيا کوي، او کومې نړۍوې له کومو نړۍو سره د لاسرسي اړيکه لري۔
\olpseganchor{OLP-0442-B007}{end}

\olpseganchor{OLP-0442-B008}{start}
خو د بشپړتيا د ثبوت په حالت کښې موږ کوم ځانګړے فارمول~\LR{$!A$}
نۀ لرو چې مدل ورته جوړ کړو۔ غواړو وښيو چې د هر هغه \LR{$!A$} دپاره
مدل شته چې \LR{$\Proves/[\Sigma] \lnot !A$} وي۔ دا لږ تر لږه اړينه
غوښتنه ده؛ ځکه \emph{که} \LR{$\Proves[\Sigma] \lnot !A$} وي، نو د
سموالي له مخې د~\LR{$!A$} دپاره هېڅ مدل نشته (چې \LR{$\Sigma$} پکې
رښتيا وي)۔ اوس پام وکړئ چې \LR{$\Proves/[\Sigma] \lnot !A$} هله او
يوازې هله وي چې \LR{$!A$} د~\LR{$\Sigma$} په نسبت سازګار وي۔ (ياد کړئ
چې \LR{$\Proves/[\Sigma] \lnot !A$} او
\LR{$!A \Proves/[\Sigma] \lfalse$} همارز دي۔) نو زموږ کار د~\LR{$\Sigma$}
په نسبت د هر سازګار فارمول دپاره مدل جوړول دي۔
\olpseganchor{OLP-0442-B008}{end}

\olpseganchor{OLP-0442-B009}{start}
هغه چل چې کاروو، دا دے: د فارمولونه داسې يو د~\LR{$\Sigma$} په نسبت
سازګار سټ پيدا کوو چې~\LR{$!A$} پکې وي، او نور فارمولونه هم ولري چې
موږ ته ووايي د~\LR{$!A$} رښتيا کوونکې نړۍ بايد څنګه وي۔ دا ډول سټونه
\emph{بشپړ} د~\LR{$\Sigma$} په نسبت سازګار سټونه دي۔ د~\LR{$!A$} د رښتيا کېدو لپاره د لاسرسي اړيکه هم
په پام کښې نيسو۔ ځينې فارمولونه د رښتياوالي لپاره تر يوې زياتې
نړۍ غواړي؛ زموږ په جوړونه کښې د لاسرسي اړيکه د نړۍو ترمنځ دا
اړيکې ټاکي۔ که هغه بشپړ د~\LR{$\Sigma$} په نسبت سازګار سټ چې~\LR{$!A$}
پکې دے، د \LR{$\Box !B$} يا~\LR{$\Diamond !C$} په بڼه نور فارمولونه
هم ولري، نو په لومړي حالت کښې \LR{$!B$} بايد په ټولو لاسرسي وړ نړۍو
کښې رښتيا وي؛ په دويم کښې \LR{$!C$} بايد لږ تر لږه په يوه کښې رښتيا وي۔ د دې دپاره به
\emph{ټول} ممکن بشپړ د~\LR{$\Sigma$} په نسبت سازګار سټونه د نړۍو د سټ
د بنسټ په توګه واخلو۔ ستونزمنه برخه به دا وي چې وټاکو زموږ په
مدل کښې يو بشپړ د~\LR{$\Sigma$} په نسبت سازګار سټ کله له بل څخه د
لاسرسي وړ ګڼل کېږي۔
\olpseganchor{OLP-0442-B009}{end}

\olpseganchor{OLP-0442-B010}{start}
ښيو به چې په همداسې تعريف شوي مدل کښې \LR{$!A$} په يوه نړۍ کښې---چې
هغه خپله هم يو بشپړ د~\LR{$\Sigma$} په نسبت سازګار سټ دے---هله او
يوازې هله رښتيا وي چې \LR{$!A$} د هغه سټ غړے وي۔ که
\LR{$!A$} د~\LR{$\Sigma$} په نسبت سازګار وي، د لږ تر لږه يوه بشپړ
د~\LR{$\Sigma$} په نسبت سازګار سټ غړے به وي (دا حقيقت به
ثابت کړو)، او په دې توګه به يوه نړۍ وي چې \LR{$!A$}~پکې رښتيا وي۔
نو يو داسې واحد مدل به ولرو چې د~\LR{$\Sigma$} په نسبت هر سازګار
فارمول~\LR{$!A$} پکې په کومه نړۍ کښې رښتيا وي۔ همدا واحد مدل
د~\LR{$\Sigma$} \emph{کانوني} مدل دے۔
\olpseganchor{OLP-0442-B010}{end}

\olpseganchor{OLP-0443-B004}{start}
\olfileid{nml}{com}{ccs}
\olpseganchor{OLP-0443-B004}{end}

\olpseganchor{OLP-0443-B005}{start}
\olsection{بشپړ \texorpdfstring{\LR{$\Sigma$}}{Σ}-سازګار سټونه}
\olpseganchor{OLP-0443-B005}{end}

\olpseganchor{OLP-0443-B006}{start}
فرض کړئ چې \LR{$\Sigma$} د مودالي فارمولونه يو سټ دے؛ دا د يوه
عادي مودالي منطق د بديهي اصلونو يا ټاکونکو اصولو سټ وګڼئ۔ يو
سټ~\LR{$\Gamma$} هله او يوازې هله \LR{$\Sigma$}-سازګار وي چې \LR{$\Gamma
\Proves/[\Sigma] \lfalse$}، يعنې له \LR{$\Sigma$} څخه د~\LR{$!A_1 \lif
(!A_2 \lif \cdots (!A_n \lif \lfalse)\dots)$} هېڅ اشتقاق
نۀ وي، چې هر \LR{$!A_i \in \Gamma$} وي۔ داسې يو «کانوني» مدل به جوړ
کړو چې هره نړۍ ئې د \LR{$\Sigma$}-سازګار سټ يو ځانګړے ډول وي: داسې
سټ چې يوازې \LR{$\Sigma$}-سازګار نۀ وي، بلکې د سازګارۍ له پلوه اعظمي وي،
په دې معنا چې د هر مودالي فارمول د صدق ارزښت ټاکي: د هر
\LR{$!A$} دپاره يا \LR{$!A \in \Gamma$} وي يا \LR{$\lnot !A \in \Gamma$}:
\olpseganchor{OLP-0443-B006}{end}

\olpseganchor{OLP-0443-B007}{start}
\begin{defn}
  يو سټ~\LR{$\Gamma$} هله او يوازې هله \emph{بشپړ \LR{$\Sigma$}-سازګار}
  وي چې \LR{$\Sigma$}-سازګار وي او د هر~\LR{$!A$} دپاره يا \LR{$!A \in \Gamma$}
  وي يا \LR{$\lnot !A \in \Gamma$}.
\end{defn}
\olpseganchor{OLP-0443-B007}{end}

\olpseganchor{OLP-0443-B008}{start}
بشپړ \LR{$\Sigma$}-سازګار سټونه~\LR{$\Gamma$} څو ګټور خاصيتونه لري۔ يو
دا چې له اشتقاقي پلوه تړلي وي، يعنې که \LR{$\Gamma \Proves[\Sigma]
!A$} وي، نو \LR{$!A \in \Gamma$} وي۔ له دې په ځانګړي ډول دا لازمېږي
چې د فارمول~\LR{$!A \in \Sigma$} هره بېلګه هم \LR{$\in \Gamma$} وي۔
سربېره پر دې، په \LR{$\Gamma$} کښې غړيتوب د بياني رابطو د صدق شرطونه
منعکسوي۔ دا به د «کانوني مدل» د تعريف په وخت مهم وي۔
\olpseganchor{OLP-0443-B008}{end}

\olpseganchor{OLP-0443-B009}{start}
\begin{prop}\ollabel{prop:ccs-properties}
  فرض کړئ چې \LR{$\Gamma$} بشپړ \LR{$\Sigma$}-سازګار دے۔ بيا:
  \begin{enumerate}
  \item \ollabel{prop:ccs-closed}%
    \LR{$\Gamma$} په~\LR{$\Sigma$} کښې له اشتقاقي پلوه تړلے دے۔
  \item \ollabel{prop:ccs-sigma}%
    \LR{$\Sigma \subseteq \Gamma$}.
  \item \ollabel{prop:ccs-lfalse}%
    \LR{$\lfalse \notin \Gamma$}
  \item \ollabel{prop:ccs-ltrue}%
    \LR{$\ltrue \in \Gamma$}
  \item \ollabel{prop:ccs-lnot}%
    \LR{$\lnot!A \in \Gamma$} هله او يوازې هله وي چې \LR{$!A \notin
    \Gamma$}.
  \item \ollabel{prop:ccs-land}%
    \LR{$!A \land !B \in \Gamma$} هله او يوازې هله وي چې \LR{$!A \in \Gamma$} او \LR{$!B \in \Gamma$} وي
  \item \ollabel{prop:ccs-lor}%
    \LR{$!A \lor !B \in \Gamma$} هله او يوازې هله وي چې \LR{$!A \in \Gamma$} يا \LR{$!B \in \Gamma$} وي
  \item \ollabel{prop:ccs-lif}%
    \LR{$!A \lif !B \in \Gamma$} هله او يوازې هله وي چې \LR{$!A \notin \Gamma$} يا \LR{$!B \in \Gamma$} وي
  \item \ollabel{prop:ccs-liff}%
    \LR{$!A \liff !B \in \Gamma$} هله او يوازې هله وي چې يا \LR{$!A \in \Gamma$} او \LR{$!B \in
    \Gamma$} دواړه وي، يا \LR{$!A \notin \Gamma$} او \LR{$!B \notin \Gamma$} دواړه وي
  \end{enumerate}
\end{prop}
\olpseganchor{OLP-0443-B009}{end}

\olpseganchor{OLP-0443-B010}{start}
\begin{proof}
  \begin{enumerate}
  \item فرض کړئ چې \LR{$\Gamma \Proves[\Sigma] !A$}، خو \LR{$!A \notin
    \Gamma$} دے۔ بيا چې \LR{$\Gamma$} بشپړ \LR{$\Sigma$}-سازګار دے، نو
    \LR{$\lnot!A \in \Gamma$} دے۔ د قطع د قاعدې له مخې به \LR{$\Gamma$}
    ناسازګار شي، ځکه چې \LR{$!A, \lnot !A \Proves[\Sigma] \lfalse$}۔
\olpseganchor{OLP-0443-B010}{end}
    
\olpseganchor{OLP-0443-B011}{start}
  \item که \LR{$!A \in \Sigma$} وي، نو \LR{$\Gamma \Proves[\Sigma] !A$} وي،
    او د اشتقاقي تړلتيا، يعنې د~\olref{prop:ccs-closed} حالت،
    له مخې \LR{$!A \in \Gamma$} وي۔
\olpseganchor{OLP-0443-B011}{end}
    
\olpseganchor{OLP-0443-B012}{start}
  \item که \LR{$\lfalse \in \Gamma$} وي، نو \LR{$\Gamma
    \Proves[\Sigma] \lfalse$} وي؛ نو \LR{$\Gamma$} به
    \LR{$\Sigma$}-ناسازګار وي۔
\olpseganchor{OLP-0443-B012}{end}
  
\olpseganchor{OLP-0443-B013}{start}
  \item \LR{$\Gamma \Proves[\Sigma] \ltrue$}؛ نو د اشتقاقي
    تړلتيا، يعنې د~\olref{prop:ccs-closed} حالت، له مخې \LR{$\ltrue \in
    \Gamma$} دے۔
\olpseganchor{OLP-0443-B013}{end}

\olpseganchor{OLP-0443-B014}{start}
\item که \LR{$\lnot!A \in \Gamma$} وي، نو د سازګارۍ له مخې \LR{$!A \notin
    \Gamma$} وي؛ او که \LR{$!A \notin \Gamma$} وي، نو \LR{$\lnot !A \in \Gamma$}
    وي، ځکه چې \LR{$\Gamma$} بشپړ \LR{$\Sigma$}-سازګار دے۔
\olpseganchor{OLP-0443-B014}{end}
    
\olpseganchor{OLP-0443-B015}{start}
  \item فرض کړئ چې \LR{$!A \land !B \in
      \Gamma$} دے۔ څرنګه چې \LR{$(!A \land !B) \lif !A$} د تاوتولوژۍ
      يوه بېلګه ده، نو د اشتقاقي تړلتيا، يعنې د~\olref{prop:ccs-closed}
      حالت، له مخې \LR{$!A \in \Gamma$} دے۔ همداسې \LR{$!B \in \Gamma$}
      هم دے۔ له بلې خوا، فرض کړئ چې \LR{$!A \in \Gamma$} او \LR{$!B \in
      \Gamma$} دواړه دي۔ بيا له اشتقاقي تړلتيا څخه \LR{$(!A \land !B) \in
      \Gamma$} لازمېږي، ځکه چې \LR{$!A \lif (!B \lif (!A \land !B))$}
      د تاوتولوژۍ يوه بېلګه ده۔
\olpseganchor{OLP-0443-B015}{end}

\olpseganchor{OLP-0443-B016}{start}
  \item فرض کړئ چې \LR{$!A \lor !B \in
      \Gamma$}، او \LR{$!A \notin \Gamma$} او \LR{$!B \notin \Gamma$} دي۔ چې
      \LR{$\Gamma$} بشپړ \LR{$\Sigma$}-سازګار دے، نو \LR{$\lnot!A \in \Gamma$}
      او \LR{$\lnot !B \in \Gamma$} دي۔ بيا \LR{$\lnot(!A \lor !B) \in \Gamma$}
      دے، ځکه چې \LR{$\lnot !A \lif (\lnot !B \lif \lnot (!A \lor !B))$}
      د تاوتولوژۍ يوه بېلګه ده۔ دا به په دې معنا وي چې \LR{$\Gamma$}
      \LR{$\Sigma$}-ناسازګار دے، چې تناقض دے۔ په سرچپه لوري کښې،
      که د فصل لږ تر لږه يوه خوا په سټ کښې وي، ټول فصل د هغې
      تاوتولوژيکي پايله ده؛ نو د اشتقاقي تړلتيا له مخې په سټ کښې دے۔
\olpseganchor{OLP-0443-B016}{end}
  
\olpseganchor{OLP-0443-B017}{start}
  \item فرض کړئ چې \LR{$!A \lif !B \in
      \Gamma$} او \LR{$!A \in \Gamma$} دي؛ بيا \LR{$\Gamma \Proves[\Sigma] !B$}
      دے، نو د اشتقاقي تړلتيا له مخې \LR{$!B \in \Gamma$} دے۔ په سرچپه
      لوري، که \LR{$!A \lif !B \notin \Gamma$} وي، نو چې \LR{$\Gamma$} بشپړ
      \LR{$\Sigma$}-سازګار دے، \LR{$\lnot (!A \lif !B) \in \Gamma$} دے۔ چې
      \LR{$\lnot(!A \lif !B) \lif !A$} د تاوتولوژۍ يوه بېلګه ده، نو د
      اشتقاقي تړلتيا له مخې \LR{$!A \in \Gamma$} دے۔ چې \LR{$\lnot(!A \lif !B) \lif
      \lnot !B$} د تاوتولوژۍ يوه بېلګه ده، نو \LR{$\lnot !B \in
      \Gamma$} دے۔ بيا \LR{$!B \notin \Gamma$} دے، ځکه چې \LR{$\Gamma$}
      \LR{$\Sigma$}-سازګار دے۔
\olpseganchor{OLP-0443-B017}{end}
  
\olpseganchor{OLP-0443-B018}{start}
  \item فرض کړئ چې \LR{$!A \liff !B \in
      \Gamma$} دے۔ که \LR{$!A \in \Gamma$} وي، نو \LR{$!B \in \Gamma$} وي،
      ځکه چې \LR{$(!A \liff !B) \lif (!A \lif !B)$} د تاوتولوژۍ يوه
      بېلګه ده۔ همدارنګه، که \LR{$!B \in \Gamma$} وي، نو \LR{$!A \in \Gamma$}
      وي۔ نو يا \LR{$!A \in \Gamma$} او \LR{$!B \in \Gamma$} دواړه وي، يا
      نۀ \LR{$!A \in \Gamma$} وي او نۀ \LR{$!B \in \Gamma$} وي۔
\olpseganchor{OLP-0443-B018}{end}
      
\olpseganchor{OLP-0443-B019}{start}
      په سرچپه لوري، فرض کړئ چې \LR{$!A \liff !B \notin \Gamma$} دے۔ چې
      \LR{$\Gamma$} بشپړ \LR{$\Sigma$}-سازګار دے، نو \LR{$\lnot (!A \liff !B) \in
      \Gamma$} دے۔ چې \LR{$\lnot(!A \liff !B) \lif (!A \lif \lnot !B)$} د
      تاوتولوژۍ يوه بېلګه ده، نو که \LR{$!A \in \Gamma$} وي، \LR{$\lnot !B \in
      \Gamma$} وي؛ او چې \LR{$\Gamma$} \LR{$\Sigma$}-سازګار دے، نو \LR{$!B \notin
      \Gamma$} وي۔ همدارنګه، که \LR{$!B \in \Gamma$} وي، نو \LR{$!A \notin \Gamma$}
      وي۔ د دوه‌اړخيز شرط له نفي څخه د دواړو خواوو فصل په
      تاوتولوژيکي ډول لازمېږي۔ نو د اشتقاقي تړلتيا او د فصل د
      غړيتوب له مخې لږ تر لږه يوه خوا په سټ کښې ده۔ نو نۀ
      \LR{$!A \in \Gamma$} او \LR{$!B \in \Gamma$} دواړه وي، او نۀ
      \LR{$!A \notin \Gamma$} او \LR{$!B \notin \Gamma$} دواړه وي۔
  \end{enumerate}
\end{proof}
\olpseganchor{OLP-0443-B019}{end}

\olpseganchor{OLP-0444-B004}{start}
\olfileid{nml}{com}{lin}
\olpseganchor{OLP-0444-B004}{end}

\olpseganchor{OLP-0444-B005}{start}
\olsection{د لينډنباوم لمه}
\olpseganchor{OLP-0444-B005}{end}

\olpseganchor{OLP-0444-B006}{start}
د لينډنباوم لمه ثابتوي چې د فارمولونه هر \LR{$\Sigma$}-سازګار سټ
لږ تر لږه په يوه \emph{بشپړ} \LR{$\Sigma$}-سازګار سټ کښې شامل دے۔
د کانوني مدل زموږ جوړونه به وښيي چې د هر بشپړ \LR{$\Sigma$}-سازګار
سټ~\LR{$\Delta$} دپاره په کانوني مدل کښې داسې يوه نړۍ شته چې ټول او
يوازې هغه فارمولونه پکې رښتيا وي چې په~\LR{$\Delta$} کښې دي۔ نو د
لينډنباوم لمه ډاډ راکوي چې هر \LR{$\Sigma$}-سازګار سټ د کانوني مدل
په کومه نړۍ کښې رښتيا دے۔
\olpseganchor{OLP-0444-B006}{end}

\olpseganchor{OLP-0444-B007}{start}
\begin{thm}[د لينډنباوم لمه]\ollabel{thm:lindenbaum}
  که \LR{$\Gamma$} \LR{$\Sigma$}-سازګار وي، نو داسې يو بشپړ \LR{$\Sigma$}-سازګار
  سټ~\LR{$\Delta$} شته چې~\LR{$\Gamma$} غځوي۔
\end{thm}
\olpseganchor{OLP-0444-B007}{end}

\olpseganchor{OLP-0444-B008}{start}
\begin{proof}
\LR{$!A_0$}، \LR{$!A_1$}، \dots{} دې د ژبې د ټولو فارمولونو يو بشپړ لړ
وي (تکرارونه روا دي)۔ د بېلګې په توګه، لومړے \LR{$\Obj p_0$} په لړ
کښې راوړئ، او په هر پړاو \LR{$n \ge 1$} کښې هغه متناهي شمېر فارمولونه
په لړ کښې راوړئ چې اوږدوالے ئې تر~\LR{$n$} زيات نۀ وي او يوازې له
\LR{$\Obj p_0$}، \dots،~\LR{$\Obj p_n$} څخه متغيرونه کاروي۔ د هر
فارمول د شاملولو لپاره داسې پړاو بس دے چې شمېر ئې د فارمول له
اوږدوالي او د متغيرونو له تر ټولو لوے شاخص څخه لږ نۀ وي۔
د فارمولونه سټونه~\LR{$\Delta_n$} پر~\LR{$n$} د استقرا له لارې تعريفوو،
او بيا \LR{$\Delta = \bigcup_n \Delta_n$} ږدو۔ لومړے \LR{$\Delta_0 = \Gamma$}
ږدو۔ که \LR{$\Delta_n$} تعريف شوے وي، نو \LR{$\Delta_{n+1}$} داسې تعريفوو:
\[
\Delta_{n+1} =
\begin{cases}
  \Delta_n \cup \{!A_n\}, & \text{\RL{که \LR{$\Delta_n \cup \{ !A_n\}$}
    \LR{$\Sigma$}-سازګار وي؛}} \\
  \Delta_n \cup \{ \lnot !A_n\}, & \text{\RL{که داسې نۀ وي۔}}
\end{cases}
\]
اوس \LR{$\Delta = \bigcup_{n=0}^\infty \Delta_n$} ږدو۔
\olpseganchor{OLP-0444-B008}{end}

\olpseganchor{OLP-0444-B009}{start}
بايد وښيو چې دا تعريف په رښتيا داسې يو سټ~\LR{$\Delta$} راکوي چې
غوښتل شوي خاصيتونه لري، يعنې \LR{$\Gamma \subseteq \Delta$} او
\LR{$\Delta$} بشپړ \LR{$\Sigma$}-سازګار دے۔
\olpseganchor{OLP-0444-B009}{end}

\olpseganchor{OLP-0444-B010}{start}
دا څرګنده ده چې \LR{$\Gamma \subseteq \Delta$}، ځکه د جوړونې له مخې
\LR{$\Delta_0 \subseteq \Delta$} او \LR{$\Delta_0 = \Gamma$} دي۔ په رښتيا،
د هر~\LR{$n$} دپاره \LR{$\Delta_n \subseteq \Delta$}، ځکه \LR{$\Delta$} د ټولو
سټونو~\LR{$\Delta_n$} اتحاد دے۔ (چې د جوړونې په هر ګام کښې له مخکې
جوړ شوي سټ سره فارمول زياتوو، نو \LR{$\Delta_n \subseteq
\Delta_{n+1}$}؛ او چې \LR{$\subseteq$} انتقالي دے، نو هر کله چې
\LR{$n \le m$} وي، \LR{$\Delta_n \subseteq \Delta_{m}$} وي۔) د جوړونې په
هر پړاو کښې يا \LR{$!A_n$} يا \LR{$\lnot !A_n$} زياتوو، او هر فارمول
د ټولو فارمولونو~\LR{$!A_n$} په لړ کښې (لږ تر لږه يو ځل) راځي۔ نو د
هر \LR{$!A$} دپاره يا \LR{$!A \in \Delta$} وي يا \LR{$\lnot !A \in \Delta$}؛
نو د تعريف له مخې \LR{$\Delta$} بشپړ دے۔
\olpseganchor{OLP-0444-B010}{end}

\olpseganchor{OLP-0444-B011}{start}
په پای کښې بايد وښيو چې \LR{$\Delta$} \LR{$\Sigma$}-سازګار دے۔ د دې دپاره
ښيو چې (a) که \LR{$\Delta$} \LR{$\Sigma$}-ناسازګار وے، نو کوم \LR{$\Delta_n$}
به \LR{$\Sigma$}-ناسازګار وے، او (b) ټول سټونه~\LR{$\Delta_n$}
\LR{$\Sigma$}-سازګار دي۔
\olpseganchor{OLP-0444-B011}{end}

\olpseganchor{OLP-0444-B012}{start}
نو فرض کړئ چې \LR{$\Delta$} \LR{$\Sigma$}-ناسازګار وے۔ بيا \LR{$\Delta
\Proves[\Sigma] \lfalse$}، يعنې داسې \LR{$!A_1$}، \dots،~\LR{$!A_k \in
\Delta$} شته چې \LR{$\Sigma \Proves !A_1 \lif (!A_2 \lif \cdots (!A_k
\lif \lfalse)\dots)$}۔ که د فرضيو لړۍ تشه وي، هماغه بې‌فرضيې
ثبوت لومړنے سټ هم ناسازګاروي۔ اوس د ناتشې لړۍ حالت وګورئ۔
چې \LR{$\Delta = \bigcup_{n=0}^\infty \Delta_n$}
دے، نو هر \LR{$!A_i \in \Delta_{n_i}$} د کوم~\LR{$n_i$} دپاره دے۔ \LR{$n$} دې
په دې شاخصونو کښې تر ټولو لوے وي۔ چې \LR{$n_i \le n$}، نو
\LR{$\Delta_{n_i} \subseteq \Delta_n$}۔ نو ټول \LR{$!A_i$} په کوم يوه
سټ~\LR{$\Delta_n$} کښې دي۔ دا به په دې معنا وي چې \LR{$\Delta_n
\Proves[\Sigma] \lfalse$}، يعنې \LR{$\Delta_n$} \LR{$\Sigma$}-ناسازګار دے۔
\olpseganchor{OLP-0444-B012}{end}

\olpseganchor{OLP-0444-B013}{start}
د دې د ښودلو دپاره چې هر \LR{$\Delta_n$} \LR{$\Sigma$}-سازګار دے، پر~\LR{$n$}
ساده استقرا کاروو۔ \LR{$\Delta_0 = \Gamma$}، او فرض مو کړے و چې
\LR{$\Gamma$} \LR{$\Sigma$}-سازګار دے۔ نو دعوه د \LR{$n = 0$} دپاره سمه ده۔
اوس فرض کړئ چې د \LR{$n$} دپاره سمه وي، يعنې \LR{$\Delta_n$}
\LR{$\Sigma$}-سازګار وي۔ \LR{$\Delta_{n+1}$} يا \LR{$\Delta_n \cup \{!A_n\}$}
دے که دا \LR{$\Sigma$}-سازګار وي، او که داسې نۀ وي، نو
\LR{$\Delta_n \cup \{\lnot!A_n\}$} دے۔ په لومړي حالت کښې
\LR{$\Delta_{n+1}$} څرګند \LR{$\Sigma$}-سازګار دے۔ خو د
\olref[prf][con]{prop:consistencyfacts} د
\olref[prf][con]{prop:consistencyfacts-c} برخې له مخې يا \LR{$\Delta_n \cup \{!A_n\}$} يا
\LR{$\Delta_n \cup \{\lnot!A_n\}$} سازګار دے؛ نو په دويم حالت کښې هم
\LR{$\Delta_{n+1}$} سازګار دے۔
\end{proof}
\olpseganchor{OLP-0444-B013}{end}

\olpseganchor{OLP-0444-B014}{start}
\begin{cor}\ollabel{cor:provability-characterization}
  \LR{$\Gamma \Proves[\Sigma] !A$} هله او يوازې هله وي چې \LR{$!A \in \Delta$}
  د هر هغه بشپړ \LR{$\Sigma$}-سازګار سټ~\LR{$\Delta$} دپاره وي چې
  \LR{$\Gamma$} غځوي (په دې کښې هغه حالت هم شامل دے چې
  \LR{$\Gamma = \emptyset$}، او په هغه حالت کښې د مودالي نظام~\LR{$\Sigma$}
  يوه بله ځانګړونه ترلاسه کوو)۔
\end{cor}
\olpseganchor{OLP-0444-B014}{end}

\olpseganchor{OLP-0444-B015}{start}
\begin{proof}
  فرض کړئ چې \LR{$\Gamma \Proves[\Sigma] !A$}، او \LR{$\Delta$} دې هر هغه
  بشپړ \LR{$\Sigma$}-سازګار سټ وي چې \LR{$\Gamma$} غځوي۔ که \LR{$!A
  \notin \Delta$} وي، نو د اعظميت له مخې \LR{$\lnot!A \in \Delta$}،
  او بيا \LR{$\Delta \Proves[\Sigma] !A$} (د يکنواختۍ له مخې) او
  \LR{$\Delta \Proves[\Sigma] \lnot!A$} (د انعکاس له مخې)؛ نو
  \LR{$\Delta$} ناسازګار دے۔ په سرچپه لوري، که
  \LR{$\Gamma \Proves/[\Sigma] !A$} وي، نو
  \LR{$\Gamma \cup \{ \lnot!A\}$} \LR{$\Sigma$}-سازګار دے، او د لينډنباوم
  د لمې له مخې داسې يو بشپړ سازګار سټ~\LR{$\Delta$} شته چې
  \LR{$\Gamma \cup \{ \lnot!A \}$} غځوي۔ د سازګارۍ له مخې \LR{$!A
  \notin \Delta$}.
\end{proof}
\olpseganchor{OLP-0444-B015}{end}

\olpseganchor{OLP-0445-B004}{start}
\olfileid{nml}{com}{mod}
\olpseganchor{OLP-0445-B004}{end}

\olpseganchor{OLP-0445-B005}{start}
\olsection{مودالي عملګر او بشپړ سازګار سټونه}
\olpseganchor{OLP-0445-B005}{end}

\olpseganchor{OLP-0445-B006}{start}
\begin{explain}
کله چې داسې يو مدل~\LR{$\mModel{M^\Sigma}$} جوړوو چې د نړۍو سټ ئې
په کوم عادي مودالي منطق~\LR{$\Sigma$} کښې له بشپړو
\LR{$\Sigma$}-سازګارو سټونو~\LR{$\Delta$} څخه جوړ وي، نو د دغو «نړۍو»
تر منځ د لاسرسي يوه اړيکه~\LR{$R^\Sigma$} هم تعريفول غواړي۔ غواړو
د لاسرسي اړيکه (او ټاکنه \LR{$V^\Sigma$}) داسې تعريف شي چې
\LR{$\mSat{M^\Sigma}{!A}[\Delta]$} هله او يوازې هله وي چې
\LR{$!A \in \Delta$} وي۔ دا به څنګه وکړو؟
\olpseganchor{OLP-0445-B006}{end}

\olpseganchor{OLP-0445-B007}{start}
د لاسرسي د اړيکې تر تعريف وروسته، په يوه نړۍ کښې د صدق تعريف
دا ډاډ راکوي چې \LR{$\mSat{M^\Sigma}{\Box!A}[\Delta]$}
  هله او يوازې هله وي چې \LR{$\mSat{M^\Sigma}{!A}[\Delta']$} د ټولو
  هغو \LR{$\Delta'$} دپاره وي چې \LR{$R^\Sigma\Delta\Delta'$} وي۔ د دې ثبوت چې
\LR{$\mSat{M^\Sigma}{!A}[\Delta]$} هله او يوازې هله وي چې
\LR{$!A \in \Delta$}، دا غواړي چې په ځانګړي ډول د هغو فارمولونه
دپاره هم دا سمه وي چې په يوه مودالي عملګر پيل کېږي، يعنې
\LR{$\mSat{M^\Sigma}{\Box!A}[\Delta]$} هله او يوازې هله
  وي چې \LR{$\Box !A \in \Delta$}۔ دا غوښتنه په يوه نړۍ کښې د
\LR{$\Box !A$} د صدق له تعريف سره يو
ځای کولو څخه دا ترلاسه کېږي:
\[
%
  \Box!A \in \Delta \text{\RL{ هله او يوازې هله چې }} !A \in \Delta' \text{\RL{
    د هر داسې \LR{$\Delta'$} دپاره چې }} R^\Sigma\Delta\Delta'
\]
له کيڼ نه ښي لوري ته پام وکړئ: دا وايي چې که
  \LR{$\Box !A \in \Delta$} وي، نو د هر \LR{$!A$} او هر داسې \LR{$\Delta'$}
  دپاره چې \LR{$R^\Sigma\Delta\Delta'$} وي، \LR{$!A \in \Delta'$} وي۔ که
  دا وټاکو چې \LR{$R^\Sigma\Delta\Delta'$} هله او يوازې هله وي چې
  \LR{$!A \in \Delta'$} د ټولو هغو \LR{$\Box!A \in \Delta$} دپاره وي، نو
  دا شرط پوره کېږي۔ د «هله او يوازې هله» ښي اړخ شرط داسې په لنډ
  ډول ليکلے شو: \LR{$\Setabs{!A}{\Box!A \in \Delta}
  \subseteq \Delta'$}۔
\olpseganchor{OLP-0445-B007}{end}

\olpseganchor{OLP-0445-B008}{start}
نو پوښتنه دا ده: آيا د \LR{$R^\Sigma$} دا تعريف په رښتيا ډاډ راکوي
چې \LR{$\Box !A \in \Delta$} هله او يوازې هله وي چې
  \LR{$\mSat{M^\Sigma}{\Box !A}[\Delta]$}؟ آيا دا
    هم ډاډ راکوي چې \LR{$\Diamond !A \in \Delta$}
  هله او يوازې هله وي چې \LR{$\mSat{M^\Sigma}{\Diamond !A}[\Delta]$}؟
راتلونکې څو پايلې به دا ثابتې کړي۔
\end{explain}
\olpseganchor{OLP-0445-B008}{end}

\olpseganchor{OLP-0445-B009}{start}
\begin{defn}
  که \LR{$\Gamma$} د فارمولونه يو سټ وي، نو دا ږدو:
  \begin{align*}
    \Box\Gamma & = \Setabs{\Box !B}{!B \in \Gamma}\\
    \Diamond\Gamma & = \Setabs{\Diamond !B}{!B \in \Gamma}\\
    \intertext{\RL{او}}
    \Box^{-1}\Gamma & = \Setabs{!B}{\Box !B \in \Gamma}\\
    \Diamond^{-1}\Gamma & = \Setabs{!B}{\Diamond !B \in \Gamma}\\
  \end{align*}
\end{defn}
\olpseganchor{OLP-0445-B009}{end}

\olpseganchor{OLP-0445-B010}{start}
په بل عبارت، \LR{$\Box\Gamma$} هغه \LR{$\Gamma$} دے چې په~\LR{$\Gamma$} کښې
د هر فارمول مخې ته ئې \LR{$\Box$} زيات کړے وي؛ \LR{$\Box^{-1}\Gamma$}
په~\LR{$\Gamma$} کښې له ټولو هغو فارمولونه څخه جوړ دے چې په
\LR{$\Box$} پيل کېږي، خو د پيل \LR{$\Box$} نښې ئې لرې شوې وي۔ دا تعريف په
خپل ځان کښې ډېر مهم نۀ دے، خو نښې به د پام وړ ساده کړي۔
\olpseganchor{OLP-0445-B010}{end}

\olpseganchor{OLP-0445-B011}{start}
پام وکړئ چې \LR{$\Box\Box^{-1}\Gamma \subseteq \Gamma$}:
\[
\Box\Box^{-1}\Gamma = \Setabs{\Box!B}{\Box!B \in \Gamma}
\]
يعنې دا يوازې په~\LR{$\Gamma$} کښې د ټولو هغو فارمولونه سټ دے
چې په~\LR{$\Box$} پيل کېږي۔
\olpseganchor{OLP-0445-B011}{end}

\olpseganchor{OLP-0445-B012}{start}
\begin{lem}\ollabel{lem:box1}
  که \LR{$\Gamma \Proves[\Sigma] !A$} وي، نو
  \LR{$\Box\Gamma \Proves[\Sigma] \Box !A$} وي۔
\end{lem}
\olpseganchor{OLP-0445-B012}{end}

\olpseganchor{OLP-0445-B013}{start}
\begin{proof}
  که \LR{$\Gamma \Proves[\Sigma] !A$} وي او د فرضيو لړۍ تشه وي،
  د لزوم قاعده د نظام پر خپله بې‌فرضيې نتيجې پلي کوو۔ په ناتشه
  لړۍ کښې داسې \LR{$!B_1$}، \dots،
  \LR{$!B_k \in \Gamma$} شته چې \LR{$\Sigma \Proves !B_1 \lif (!B_2 \lif \cdots
  (!B_k \lif !A)\cdots)$}۔ چې \LR{$\Sigma$} عادي دے، نو د \RK{}
  قاعدې له مخې \LR{$\Sigma \Proves \Box!B_1 \lif (\Box!B_2 \lif \cdots (\Box!B_k \lif
  \Box!A)\cdots)$}، چې په څرګند ډول \LR{$\Box!B_1$}، \dots،
  \LR{$\Box!B_k \in \Box\Gamma$} دي۔ نو د تعريف له مخې
  \LR{$\Box\Gamma \Proves[\Sigma] \Box !A$}۔
\end{proof}
\olpseganchor{OLP-0445-B013}{end}

\olpseganchor{OLP-0445-B014}{start}
\begin{lem}\ollabel{lem:box2}
  که \LR{$\Box^{-1} \Gamma \Proves[\Sigma] !A$} وي، نو
  \LR{$\Gamma \Proves[\Sigma] \Box!A$} وي۔
\end{lem}
\olpseganchor{OLP-0445-B014}{end}

\olpseganchor{OLP-0445-B015}{start}
\begin{proof}
  فرض کړئ چې \LR{$\Box^{-1}\Gamma \Proves[\Sigma] !A$}؛ بيا د
  \olref{lem:box1} له مخې
  \LR{$\Box\Box^{-1}\Gamma \Proves[\Sigma] \Box!A$}۔ څرنګه چې
  \LR{$\Box\Box^{-1}\Gamma \subseteq \Gamma$}، نو د يکنواختۍ له مخې
  \LR{$\Gamma \Proves[\Sigma] \Box!A$} هم دے۔
\end{proof}
\olpseganchor{OLP-0445-B015}{end}

\olpseganchor{OLP-0445-B016}{start}
%
% Box is primitive
\olpseganchor{OLP-0445-B016}{end}
  
\olpseganchor{OLP-0445-B017}{start}
\begin{prop}\ollabel{prop:box}
  که \LR{$\Gamma$} بشپړ \LR{$\Sigma$}-سازګار وي، نو \LR{$\Box !A \in \Gamma$}
  هله او يوازې هله وي چې د هر هغه بشپړ \LR{$\Sigma$}-سازګار
  سټ~\LR{$\Delta$} دپاره چې \LR{$\Box^{-1} \Gamma \subseteq \Delta$} وي،
  \LR{$!A \in \Delta$} وي۔
\end{prop}
\olpseganchor{OLP-0445-B017}{end}

\olpseganchor{OLP-0445-B018}{start}
\begin{proof}
  فرض کړئ چې \LR{$\Gamma$} بشپړ \LR{$\Sigma$}-سازګار دے۔ د «يوازې که» لور
  اسانه دے: فرض کړئ چې \LR{$\Box !A \in \Gamma$} او
  \LR{$\Box^{-1}\Gamma \subseteq \Delta$} دي۔ چې \LR{$\Box!A \in \Gamma$}،
  نو \LR{$!A \in \Box^{-1}\Gamma \subseteq \Delta$}؛ نو \LR{$!A \in \Delta$}.
\olpseganchor{OLP-0445-B018}{end}

د «که» لوري دپاره نقيض عکس ثابتوو: فرض کړئ چې
  \LR{$\Box!A \notin \Gamma$}۔ چې \LR{$\Gamma$} بشپړ \LR{$\Sigma$}-سازګار دے،
  نو له اشتقاقي پلوه تړلے دے، او له دې څخه
  \LR{$\Gamma \Proves/[\Sigma] \Box !A$} لازمېږي۔ د \olref{lem:box2}
  له مخې \LR{$\Box^{-1}\Gamma \Proves/[\Sigma] !A$}۔ د
  \LR{\olref[prf][con]{prop:consistencyfacts}(\olref[prf][con]{prop:consistencyfacts-b})}
  له مخې \LR{$\Box^{-1}\Gamma \cup \{ \lnot!A \}$} \LR{$\Sigma$}-سازګار
  دے۔ د لينډنباوم د لمې له مخې داسې يو بشپړ \LR{$\Sigma$}-سازګار
  سټ~\LR{$\Delta$} شته چې \LR{$\Box^{-1}\Gamma \cup \{ \lnot!A \} \subseteq
  \Delta$} وي۔ د سازګارۍ له مخې \LR{$!A \notin \Delta$}.
\end{proof}

\olpseganchor{OLP-0445-B023}{start}
\begin{lem}\ollabel{lem:box-iff-diamond}
  فرض کړئ چې \LR{$\Gamma$} او \LR{$\Delta$} بشپړ \LR{$\Sigma$}-سازګار دي۔ بيا
  \LR{$\Box^{-1}\Gamma \subseteq \Delta$} هله او يوازې هله وي چې
  \LR{$\Diamond\Delta \subseteq \Gamma$} وي۔
\end{lem}
\olpseganchor{OLP-0445-B023}{end}

\olpseganchor{OLP-0445-B024}{start}
\begin{proof}
  د «يوازې که» لور: فرض کړئ چې \LR{$\Box^{-1}\Gamma \subseteq \Delta$}،
  او فرض کړئ چې \LR{$\Diamond!A \in \Diamond\Delta$} (يعنې
  \LR{$!A \in \Delta$})۔ د دې د ښودلو دپاره چې \LR{$\Diamond!A \in \Gamma$}،
  د \LR{$\Box\lnot!A \notin \Gamma$} ښودل بس دي، ځکه بيا د اعظميت له
  مخې \LR{$\lnot\Box\lnot !A \in \Gamma$} وي۔ اوس که
  \LR{$\Box\lnot!A \in \Gamma$} وي، نو د فرضيې له مخې
  \LR{$\lnot!A \in \Delta$} وي، چې دا د \LR{$\Delta$} له سازګارۍ سره په
  ټکر کښې دے (ځکه \LR{$!A \in \Delta$})۔ نو
  \LR{$\Box\lnot!A \notin \Gamma$}، لکه چې غوښتل شوي وو۔
  بيا د دوه‌ګونتيا او اشتقاقي تړلتيا له مخې د امکان مطلوب
  فارمول هم د سټ غړے دے۔
\olpseganchor{OLP-0445-B024}{end}

\olpseganchor{OLP-0445-B025}{start}
  د «که» لور: فرض کړئ چې \LR{$\Diamond\Delta \subseteq \Gamma$}۔ د
  نقيض عکس له لارې استدلال کوو: د \LR{$\Box!A \notin \Gamma$} د ښودلو
  دپاره فرض کړئ چې \LR{$!A \notin \Delta$}۔ که \LR{$!A \notin \Delta$}
  وي، نو د اعظميت له مخې \LR{$\lnot!A \in \Delta$}، او د فرضيې له
  مخې \LR{$\Diamond\lnot!A \in \Gamma$}۔ خو په يوه عادي مودالي
  منطق کښې \LR{$\Diamond\lnot!A$} له \LR{$\lnot\Box !A$} سره همارز دے؛
  نو د اشتقاقي تړلتيا له مخې دا وروستے په \LR{$\Gamma$} کښې دے؛
  د سازګارۍ له مخې
  \LR{$\Box!A \notin\Gamma$}، لکه چې غوښتل شوي وو۔
\end{proof}
\olpseganchor{OLP-0445-B025}{end}

\olpseganchor{OLP-0445-B026}{start}
%
  %
% Box and Diamond are both primitive; prove
% prop:diamond using lem:box-iff-diamond
\olpseganchor{OLP-0445-B026}{end}
  
\olpseganchor{OLP-0445-B027}{start}
\begin{prop}\ollabel{prop:diamond}
  که \LR{$\Gamma$} بشپړ \LR{$\Sigma$}-سازګار وي، نو \LR{$\Diamond !A \in \Gamma$}
  هله او يوازې هله وي چې د کوم هغه بشپړ \LR{$\Sigma$}-سازګار
  سټ~\LR{$\Delta$} دپاره چې \LR{$\Diamond\Delta \subseteq \Gamma$} وي،
  \LR{$!A \in \Delta$} وي۔
\end{prop}
\olpseganchor{OLP-0445-B027}{end}

\olpseganchor{OLP-0445-B028}{start}
\begin{proof}
فرض کړئ چې \LR{$\Gamma$} بشپړ \LR{$\Sigma$}-سازګار دے۔ د \Dual{} او
تړلتيا له مخې \LR{$\Diamond!A \in \Gamma$} هله او يوازې هله وي چې
\LR{$\lnot\Box\lnot !A \in \Gamma$} وي۔ د
\LR{\olref[ccs]{prop:ccs-properties}(\olref[ccs]{prop:ccs-lnot})} له مخې
\LR{$\lnot\Box\lnot !A \in \Gamma$} هله او يوازې هله وي چې
\LR{$\Box\lnot !A \notin \Gamma$} وي، ځکه \LR{$\Gamma$} بشپړ
\LR{$\Sigma$}-سازګار دے۔ د \olref{prop:box} له مخې
\LR{$\Box\lnot !A \notin \Gamma$} هله او يوازې هله وي چې د کوم
بشپړ \LR{$\Sigma$}-سازګار سټ~\LR{$\Delta$} دپاره چې
\LR{$\Box^{-1}\Gamma \subseteq \Delta$} وي، \LR{$\lnot!A \notin \Delta$}
وي۔ اوس هر داسې سټ~\LR{$\Delta$} په نظر کښې ونيسئ۔ د
\olref{lem:box-iff-diamond} له مخې
\LR{$\Box^{-1}\Gamma \subseteq \Delta$} هله او يوازې هله وي چې
\LR{$\Diamond\Delta \subseteq \Gamma$} وي۔ همدارنګه، د
\LR{\olref[ccs]{prop:ccs-properties}(\olref[ccs]{prop:ccs-lnot})} له مخې
\LR{$\lnot !A \notin \Delta$} هله او يوازې هله وي چې \LR{$!A \in \Delta$}
وي۔ نو \LR{$\Diamond!A \in \Gamma$} هله او يوازې هله وي چې د کوم
بشپړ \LR{$\Sigma$}-سازګار سټ~\LR{$\Delta$} دپاره چې
\LR{$\Diamond\Delta \subseteq \Gamma$} وي، \LR{$!A \in \Delta$} وي۔
\end{proof}
\olpseganchor{OLP-0445-B028}{end}


  
\olpseganchor{OLP-0445-B030}{start}
\begin{prob}
  وښيئ چې که \LR{$\Gamma$} بشپړ \LR{$\Sigma$}-سازګار وي، نو
  \LR{$\Diamond !A \in \Gamma$} هله او يوازې هله وي چې
    داسې يو بشپړ \LR{$\Sigma$}-سازګار سټ~\LR{$\Delta$} شته چې
    \LR{$\Box^{-1}\Gamma \subseteq \Delta$} او \LR{$!A \in \Delta$} وي۔
  \emph{دا کار د \olref[nml][com][mod]{lem:box-iff-diamond} له
    کارولو پرته وکړئ}۔
\end{prob}
\olpseganchor{OLP-0445-B030}{end}

\olpseganchor{OLP-0446-B004}{start}
\olfileid{nml}{com}{cmd}
\olpseganchor{OLP-0446-B004}{end}

\olpseganchor{OLP-0446-B005}{start}
\olsection{کانوني مدلونه}
\olpseganchor{OLP-0446-B005}{end}

\olpseganchor{OLP-0446-B006}{start}
د يوه مودالي نظام~\LR{$\Sigma$} دپاره \emph{کانوني مدل} هغه ځانګړے
مدل~\LR{$\mModel{M^\Sigma}$} دے چې نړۍوې ئې ټول بشپړ
\LR{$\Sigma$}-سازګار سټونه دي۔ د لاسرسي اړيکه~\LR{$R^\Sigma$} او
ارزښت ټاکنه \LR{$V^\Sigma$} ئې داسې تعريف شوي دي چې په يوه
نړۍ~\LR{$\Delta$} کښې رښتيني فارمولونه په عين ډول هغه
فارمولونه وي چې~\LR{$\Delta$} ترې جوړ دے۔
\olpseganchor{OLP-0446-B006}{end}

\olpseganchor{OLP-0446-B007}{start}
\begin{defn}
  \LR{$\Sigma$} دې يو عادي مودالي منطق وي او تش سټ دې د همدې نظام
  په نسبت سازګار وي۔ د لينډنباوم د لمې له مخې د بشپړو سازګارو
  سټونو ټولګے ناتش دے۔ د \LR{$\Sigma$} دپاره
  \emph{کانوني مدل} \LR{$\mModel{M}^\Sigma = \tuple{W^\Sigma, R^\Sigma,
    V^\Sigma}$} دے، چې:
  \begin{enumerate}
  \item \LR{$W^\Sigma = \Setabs{ \Delta }{\Delta \text{\RL{
      بشپړ \LR{$\Sigma$}-سازګار دے}} }$}.
  \item \LR{$R^\Sigma \Delta\Delta'$} هله او يوازې هله وي چې
    \LR{$\Box^{-1}\Delta \subseteq
      \Delta'$}.
  \item \LR{$V^\Sigma(p) = \Setabs{\Delta}{p \in \Delta}$}.
  \end{enumerate}
\end{defn}
\olpseganchor{OLP-0446-B007}{end}

\olpseganchor{OLP-0447-B004}{start}
\olfileid{nml}{com}{tru}
\olpseganchor{OLP-0447-B004}{end}

\olpseganchor{OLP-0447-B005}{start}
\olsection{د صدق لمه}
\olpseganchor{OLP-0447-B005}{end}

\olpseganchor{OLP-0447-B006}{start}
کانوني مدل~\LR{$\mModel{M^\Sigma}$} داسې تعريف شوے دے چې
\LR{$\mSat{M^\Sigma}{!A}[\Delta]$} هله او يوازې هله چې \LR{$!A \in \Delta$}۔ د بياني متغيرونه لپاره د \LR{$V^\Sigma$} تعريف دا نتيجه نېغ په نېغه ورکوي۔ خو بايد
وښيو چې دا همارزي د ټولو فارمولونه لپاره سمه ده۔ دا کار په استقرا کوو۔
په استقرايي ګام کښې د هغو فارمولونه لپاره د همارزۍ ثبوت راځي چې بياني
رابطې لري (چې هلته \olref[ccs]{prop:ccs-properties} کارول پکار دي)، او د هغو لپاره
چې مودالي عملګر لري (چې هلته د \olref[mod]{sec} نتيجې کاروو)۔
\olpseganchor{OLP-0447-B006}{end}

\olpseganchor{OLP-0447-B007}{start}
\begin{prop}[د صدق لمه]
  \ollabel{prop:truthlemma}
  د هر فارمول~\LR{$!A$} لپاره، د کانوني مدل په هره نړۍ کښې
  \LR{$\mSat{M^\Sigma}{!A}[\Delta]$} هله او يوازې هله چې
  \LR{$!A \in \Delta$}۔
\end{prop}
\olpseganchor{OLP-0447-B007}{end}

\olpseganchor{OLP-0447-B008}{start}
\begin{proof}
  پر~\LR{$!A$} استقرا کوو۔
  \begin{enumerate}
    \item \indcase{!A}{\lfalse}{\LR{$\mSat/{M^\Sigma}{\lfalse}[\Delta]$}
        د \olref[syn][trw]{defn:mmodels} له مخې، او \LR{$\lfalse \notin \Delta$}
        د
        \LR{\olref[ccs]{prop:ccs-properties}(\olref[ccs]{prop:ccs-lfalse})} له مخې۔}
\olpseganchor{OLP-0447-B008}{end}

\olpseganchor{OLP-0447-B009}{start}
    \item \indcase{!A}{\ltrue}{\LR{$\mSat{M^\Sigma}{\ltrue}[\Delta]$}
        د \olref[syn][trw]{defn:mmodels} له مخې، او \LR{$\ltrue \in \Delta$} د
        \LR{\olref[ccs]{prop:ccs-properties}(\olref[ccs]{prop:ccs-ltrue})} له مخې۔}
\olpseganchor{OLP-0447-B009}{end}

\olpseganchor{OLP-0447-B010}{start}
    \item \indcase{!A}{p}{\LR{$\mSat{M^\Sigma}{p}[\Delta]$} هله او يوازې هله چې \LR{$\Delta \in
      V^\Sigma(p)$}، د \olref[syn][trw]{defn:mmodels} له مخې۔ همداراز، \LR{$\Delta
      \in V^\Sigma(p)$} هله او يوازې هله چې \LR{$p \in \Delta$}، د~\LR{$V^\Sigma$}
      د تعريف له مخې۔}
\olpseganchor{OLP-0447-B010}{end}

\olpseganchor{OLP-0447-B011}{start}
    \item \indcase{!A}{\lnot !B}{%
        \LR{$\mSat{M^\Sigma}{\lnot !B}[\Delta]$}
          هله او يوازې هله چې \LR{$\mSat/{M^\Sigma}{!B}[\Delta]$}
          (\olref[syn][trw]{defn:mmodels})، هله او يوازې هله چې \LR{$!B \notin \Delta$} (د
          استقرا فرضيې له مخې)، هله او يوازې هله چې \LR{$\lnot !B \in \Delta$} (د
          \LR{\olref[ccs]{prop:ccs-properties}(\olref[ccs]{prop:ccs-lnot})} له مخې)۔}
\olpseganchor{OLP-0447-B011}{end}

\olpseganchor{OLP-0447-B012}{start}
    \item \indcase{!A}{!B \land !C}{%
        \LR{$\mSat{M^\Sigma}{!B \land
            !C}[\Delta]$} هله او يوازې هله چې \LR{$\mSat{M^\Sigma}{!B}[\Delta]$} او
          \LR{$\mSat{M^\Sigma}{!C}[\Delta]$} (د
          \olref[syn][trw]{defn:mmodels} له مخې)، هله او يوازې هله چې \LR{$!B \in \Delta$} او \LR{$!C
          \in \Delta$} (د استقرا فرضيې له مخې)، هله او يوازې هله چې \LR{$!B \land !C \in
          \Delta$} (د
          \LR{\olref[ccs]{prop:ccs-properties}(\olref[ccs]{prop:ccs-land})} له مخې)۔}
\olpseganchor{OLP-0447-B012}{end}

\olpseganchor{OLP-0447-B013}{start}
    \item \indcase{!A}{!B \lor !C}{%
        \LR{$\mSat{M^\Sigma}{!B \lor
            !C}[\Delta]$} هله او يوازې هله چې \LR{$\mSat{M^\Sigma}{!B}[\Delta]$} يا
          \LR{$\mSat{M^\Sigma}{!C}[\Delta]$} (د
          \olref[syn][trw]{defn:mmodels} له مخې)، هله او يوازې هله چې \LR{$!B \in \Delta$} يا \LR{$!C
          \in \Delta$} (د استقرا فرضيې له مخې)، هله او يوازې هله چې \LR{$!B \lor !C \in
          \Delta$} (د
          \LR{\olref[ccs]{prop:ccs-properties}(\olref[ccs]{prop:ccs-lor})} له مخې)۔}
\olpseganchor{OLP-0447-B013}{end}

\olpseganchor{OLP-0447-B014}{start}
    \item \indcase{!A}{!B \lif !C}{%
        \LR{$\mSat{M^\Sigma}{!B \lif
            !C}[\Delta]$} هله او يوازې هله چې \LR{$\mSat/{M^\Sigma}{!B}[\Delta]$} يا
          \LR{$\mSat{M^\Sigma}{!C}[\Delta]$} (د
          \olref[syn][trw]{defn:mmodels} له مخې)، هله او يوازې هله چې \LR{$!B \notin \Delta$} يا
          \LR{$!C \in \Delta$} (د استقرا فرضيې له مخې)، هله او يوازې هله چې \LR{$!B \lif !C
          \in \Delta$} (د
          \LR{\olref[ccs]{prop:ccs-properties}(\olref[ccs]{prop:ccs-lif})} له مخې)۔}
\olpseganchor{OLP-0447-B014}{end}

\olpseganchor{OLP-0447-B015}{start}
    \item \indcase{!A}{!B \liff !C}{%
        \LR{$\mSat{M^\Sigma}{!B \liff
            !C}[\Delta]$} هله او يوازې هله چې يا \LR{$\mSat{M^\Sigma}{!B}[\Delta]$} او
          \LR{$\mSat{M^\Sigma}{!C}[\Delta]$}، يا
          \LR{$\mSat/{M^\Sigma}{!B}[\Delta]$} او
          \LR{$\mSat/{M^\Sigma}{!C}[\Delta]$} (د
          \olref[syn][trw]{defn:mmodels} له مخې)، هله او يوازې هله چې يا \LR{$!B \in \Delta$}
          او \LR{$!C \in \Delta$}، يا \LR{$!B \notin \Delta$} او \LR{$!C \notin
          \Delta$} (د استقرا فرضيې له مخې)، هله او يوازې هله چې \LR{$!B \liff !C \in
          \Delta$} (د
          \LR{\olref[ccs]{prop:ccs-properties}(\olref[ccs]{prop:ccs-liff})} له مخې)۔}
\olpseganchor{OLP-0447-B015}{end}

\olpseganchor{OLP-0447-B016}{start}
    \item \indcase{!A}{\Box !B}{%
        لومړے فرض کړو چې
          \LR{$\mSat{M^\Sigma}{\Box !B}[\Delta]$}۔ د
          \olref[syn][trw]{defn:mmodels} له مخې، د هر \LR{$\Delta'$} لپاره چې
          \LR{$R^\Sigma \Delta\Delta'$} وي،
          \LR{$\mSat{M^\Sigma}{!B}[\Delta']$}۔ د استقرا فرضيې له مخې، د
          هر \LR{$\Delta'$} لپاره چې \LR{$R^\Sigma \Delta\Delta'$} وي، \LR{$!B \in
          \Delta'$}۔ د \LR{$R^\Sigma$} د تعريف له مخې، د هر \LR{$\Delta'$}
          لپاره چې \LR{$\Box^{-1} \Delta \subseteq \Delta'$} وي، \LR{$!B \in
          \Delta'$}۔ د \olref[mod]{prop:box} له مخې، \LR{$\Box !B \in
          \Delta$}۔
\olpseganchor{OLP-0447-B016}{end}

\olpseganchor{OLP-0447-B017}{start}
          اوس فرض کړو چې \LR{$\Box !B \in \Delta$}۔ داسې \LR{$\Delta' \in W^\Sigma$}
          واخلو چې \LR{$R^\Sigma \Delta\Delta'$} وي، يعنې \LR{$\Box^{-1}\Delta
          \subseteq \Delta'$}۔ څرنګه چې \LR{$\Box !B \in \Delta$}، نو \LR{$!B \in
          \Box^{-1} \Delta$}۔ له دې امله، \LR{$!B \in \Delta'$}۔ د
          استقرا فرضيې له مخې، \LR{$\mSat{M^\Sigma}{!B}[\Delta']$}۔ څرنګه چې
          \LR{$\Delta'$} د \LR{$R^\Sigma \Delta\Delta'$} په شرط هره نړۍ وه، نو د
          ټولو \LR{$\Delta' \in W^\Sigma$} لپاره چې \LR{$R^\Sigma
          \Delta\Delta'$} وي، \LR{$\mSat{M^\Sigma}{!B}[\Delta']$}۔ د
          \olref[syn][trw]{defn:mmodels} له مخې، \LR{$\mSat{M^\Sigma}{\Box
            !B}[\Delta]$}۔}
\olpseganchor{OLP-0447-B017}{end}

\olpseganchor{OLP-0447-B018}{start}
    \item \indcase{!A}{\Diamond !B}{%
        لومړے فرض کړو چې
          \LR{$\mSat{M^\Sigma}{\Diamond !B}[\Delta]$}۔ د
          \olref[syn][trw]{defn:mmodels} له مخې، داسې يو \LR{$\Delta'$} شته چې
          \LR{$R^\Sigma \Delta\Delta'$} وي او \LR{$\mSat{M^\Sigma}{!B}[\Delta']$}۔
          د استقرا فرضيې له مخې، داسې يو \LR{$\Delta'$} شته چې
          \LR{$R^\Sigma \Delta\Delta'$} وي او \LR{$!B \in \Delta'$}۔ د
          \LR{$R^\Sigma$} د تعريف له مخې، داسې يو \LR{$\Delta'$} شته چې
            \LR{$\Box^{-1} \Delta \subseteq \Delta'$} وي او \LR{$!B \in \Delta'$}۔
            د \olref[mod]{lem:box-iff-diamond} له مخې، داسې يو \LR{$\Delta'$} شته
          چې \LR{$\Diamond \Delta' \subseteq \Delta$} وي او \LR{$!B \in
          \Delta'$}۔ څرنګه چې \LR{$!B \in \Delta'$}، نو \LR{$\Diamond !B \in \Diamond
          \Delta'$}، نو \LR{$\Diamond !B \in \Delta$}۔
\olpseganchor{OLP-0447-B018}{end}

\olpseganchor{OLP-0447-B019}{start}
          اوس فرض کړو چې \LR{$\Diamond !B \in \Delta$}۔ د
          \olref[mod]{prop:diamond} له مخې، داسې يو بشپړ
          \LR{$\Sigma$}-سازګار \LR{$\Delta' \in W^\Sigma$} شته چې
          \LR{$\Diamond \Delta' \subseteq \Delta$} وي او \LR{$!B \in
          \Delta'$}۔ د \olref[mod]{lem:box-iff-diamond} له مخې،
            داسې يو \LR{$\Delta' \in W^\Sigma$} شته چې \LR{$\Box^{-1}
            \Delta \subseteq \Delta'$} وي او \LR{$!B \in \Delta'$}۔ د
          \LR{$R^\Sigma$} د تعريف له مخې، \LR{$R^\Sigma \Delta\Delta'$}، نو داسې
          يو \LR{$\Delta' \in W^\Sigma$} شته چې \LR{$R^\Sigma
          \Delta\Delta'$} وي او \LR{$!B \in \Delta'$}۔ د استقرا د فرض
          له مخې په هماغه نړۍ کښې دنننے فارمول رښتيا دے۔ د
          \olref[syn][trw]{defn:mmodels} له مخې، \LR{$\mSat{M^\Sigma}{\Diamond
            !B}[\Delta]$}۔}
  \end{enumerate}
\end{proof}
\olpseganchor{OLP-0447-B019}{end}

\olpseganchor{OLP-0448-B004}{start}
\olfileid{nml}{com}{cmk}
\olpseganchor{OLP-0448-B004}{end}

\olpseganchor{OLP-0448-B005}{start}
\olsection{د \texorpdfstring{\Log{K}}{K} ټاکنه او بشپړتيا}
\olpseganchor{OLP-0448-B005}{end}

\olpseganchor{OLP-0448-B006}{start}
اوس چمتو يو چې کانوني مدل د بشپړتيا د ثابتولو لپاره وکاروو۔ بشپړتيا له دې واقعيت څخه
راځي چې د~\LR{$\Sigma$} په کانوني مدل کښې رښتيني فارمولونه عين هغه دي چې
د \LR{$\Sigma$} په نسبت د اشتقاق وړ دي۔ هغه مدلونه چې دا خاصيت لري، ويل کېږي چې
\LR{$\Sigma$} \emph{ټاکي}۔
\olpseganchor{OLP-0448-B006}{end}

\olpseganchor{OLP-0448-B007}{start}
\begin{defn}
  يو مدل \LR{$\mModel{M}$} يو عادي مودالي منطق~\LR{$\Sigma$} په عين هغه وخت
  \emph{ټاکي} چې، د ټولو فارمولونه~\LR{$!A$} لپاره، \LR{$\mSat{M}{!A}$} هله او يوازې هله
  چې \LR{$\Sigma \Proves !A$}۔
\end{defn}
\olpseganchor{OLP-0448-B007}{end}

\olpseganchor{OLP-0448-B008}{start}
\begin{thm}[ټاکنه]\ollabel{thm:determination}
   فرض کړئ چې تش سټ د مودالي نظام په نسبت سازګار دے۔ بيا
   \LR{$\mSat{M^\Sigma}{!A}$} هله او يوازې هله چې \LR{$\Sigma \Proves !A$}۔
\end{thm}
\olpseganchor{OLP-0448-B008}{end}

\olpseganchor{OLP-0448-B009}{start}
\begin{proof}
  که \LR{$\mSat{M^\Sigma}{!A}$} وي، نو د هر بشپړ
  \LR{$\Sigma$}-سازګار \LR{$\Delta$} لپاره
  \LR{$\mSat{M^\Sigma}{!A}[\Delta]$} لرو۔ نو د صدق لمې له مخې، د
  هر بشپړ \LR{$\Sigma$}-سازګار \LR{$\Delta$} لپاره \LR{$!A \in
  \Delta$}؛ بيا د \olref[lin]{cor:provability-characterization} له مخې (د \LR{$\Gamma =
  \emptyset$} په اخيستو)، \LR{$\Sigma \Proves !A$}۔
\olpseganchor{OLP-0448-B009}{end}

\olpseganchor{OLP-0448-B010}{start}
  په سرچپه لوري کښې، که \LR{$\Sigma \Proves !A$} وي، نو د
  \olref[ccs]{prop:ccs-properties} د \olref[ccs]{prop:ccs-closed} برخې له مخې هر
  بشپړ \LR{$\Sigma$}-سازګار \LR{$\Delta$}،~\LR{$!A$} لري؛ نو د
  صدق لمې له مخې، د هر~\LR{$\Delta \in W^\Sigma$} لپاره
  \LR{$\mSat{M^\Sigma}{!A}[\Delta]$}، يعنې \LR{$\mSat{M^\Sigma}{!A}$}۔
\end{proof}
\olpseganchor{OLP-0448-B010}{end}

\olpseganchor{OLP-0448-B011}{start}
څرنګه چې د \Log{K} کانوني مدل،~\Log{K} ټاکي، نو سمدستي
د \Log{K} بشپړتيا د يوې نتيجې په توګه لرو:
\olpseganchor{OLP-0448-B011}{end}

\olpseganchor{OLP-0448-B012}{start}
\begin{cor}\ollabel{cor:Kcomplete}
  بنسټيز مودالي منطق \Log{K} د ټولو مدلونو د ټولګي په نسبت
  بشپړ دے، يعنې که \LR{$\Entails !A$} وي، نو \LR{$\Log{K}
  \Proves !A$}۔
\end{cor}
\olpseganchor{OLP-0448-B012}{end}

\olpseganchor{OLP-0448-B013}{start}
\begin{proof}
  د مقابل نقيض له مخې، که \LR{$\Log{K} \Proves/ !A$} وي، نو د ټاکنې له مخې
  \LR{$\mSat/{M^{\Log{K}}}{!A}$}؛ نو \LR{$!A$} اعتبار نۀ لري۔
\end{proof}
\olpseganchor{OLP-0448-B013}{end}

\olpseganchor{OLP-0448-B014}{start}
د مدلونو د يوه ټولګي په نسبت د نظام \LR{$\Sigma$} د بشپړتيا په عمومي صورت کښې،
لکه د انعکاسي، متناظرو او انتقالي مدلونو د ټولګي په نسبت د \Log{KTB4} بشپړتيا،
يوازې ټاکنه بس نۀ ده۔ بايد دا هم وښيو چې د نظام
\LR{$\Sigma$} کانوني مدل د هماغه ټولګي غړے دے۔ دا خبره د کانوني مدل له جوړونې څخه
په څرګند ډول نۀ راځي---او تل رښتيا هم نۀ وي!
\olpseganchor{OLP-0448-B014}{end}

\olpseganchor{OLP-0449-B004}{start}
\olfileid{nml}{com}{fra}
\olpseganchor{OLP-0449-B004}{end}

\olpseganchor{OLP-0449-B005}{start}
\olsection{د چوکاټ بشپړتيا}
\olpseganchor{OLP-0449-B005}{end}

\olpseganchor{OLP-0449-B006}{start}
د \Log{K} د بشپړتيا قضيه نورو مودالي نظامونو ته غځېدے شي، کله چې وښيو
چې د يوه ټاکلي منطق کانوني مدل د چوکاټ اړوند خاصيت لري۔
\olpseganchor{OLP-0449-B006}{end}

\olpseganchor{OLP-0449-B007}{start}
\begin{thm}\ollabel{thm:completeframeprops}
  که يو سازګار عادي مودالي منطق \LR{$\Sigma$} د \olref{tab:correspondencetable} د کيڼ لوري
  له فارمولونه څخه يو ولري،
  نو د~\LR{$\Sigma$} کانوني مدل د ښي لوري اړوند خاصيت لري۔
\end{thm}
\olpseganchor{OLP-0449-B007}{end}

\olpseganchor{OLP-0449-B008}{start}
\begin{table}[htp]
  \centering
    \begin{tabular}{| l || l |}
      \hline
      \emph{که \LR{$\Sigma$} دا ولري \dots} & \emph{\dots نو د \LR{$\Sigma$} کانوني
        مدل دا خاصيت لري:} \\
      \hline \hline
      \Ax{D}:\;\; \LR{$\Box!A \lif \Diamond !A$} & \quad مسلسل؛ \\
      \hline
      \Ax{T}:\;\; \LR{$\Box!A \lif !A$} &\quad  انعکاسي؛\\
      \hline
      \Ax{B}: \;\;\LR{$!A \lif \Box\Diamond!A$} &\quad  متناظر؛ \\
      \hline
      \Ax{4}: \;\; \LR{$\Box!A \lif \Box\Box!A$} & \quad انتقالي؛ \\
      \hline 
      \Ax{5}: \;\; \LR{$\Diamond !A \lif \Box\Diamond!A$}& \quad اقليدسي۔\\
      \hline
    \end{tabular}
    \caption{د مطابقت بنسټيز واقعيتونه۔}
    \ollabel{tab:correspondencetable}
\end{table}
\olpseganchor{OLP-0449-B008}{end}

\olpseganchor{OLP-0449-B009}{start}
\begin{proof}
  دا هر يو په خپل وار اخلو۔
\olpseganchor{OLP-0449-B009}{end}

\olpseganchor{OLP-0449-B010}{start}
  فرض کړو چې \LR{$\Sigma$}، \Ax{D} لري، او \LR{$\Delta \in W^\Sigma$} واخلو؛ بايد
  وښيو چې داسې يو \LR{$\Delta'$} شته چې \LR{$R^\Sigma
  \Delta\Delta'$} وي۔ بس ده چې وښيو \LR{$\Box^{-1}\Delta$}
  \LR{$\Sigma$}-سازګار دے، ځکه بيا د لينډنباوم لمې له مخې يو
  بشپړ \LR{$\Sigma$}-سازګار سټ \LR{$\Delta' \supseteq
  \Box^{-1}\Delta$} شته؛ او د \LR{$R^\Sigma$} د تعريف
  له مخې
  \LR{$R^\Sigma \Delta\Delta'$} لرو۔ نو د تضاد لپاره فرض کړو چې
  \LR{$\Box^{-1}\Delta$} \LR{$\Sigma$}-سازګار \emph{نۀ} دے، يعنې
  \LR{$\Box^{-1}\Delta \Proves[\Sigma] \lfalse$}۔ د \olref[mod]{lem:box2} له مخې،
  \LR{$\Delta \Proves[\Sigma] \Box\lfalse$}؛ او چونکې \LR{$\Sigma$}،
  \Ax{D} لري، نو \LR{$\Delta \Proves[\Sigma] \Diamond\lfalse$} هم۔ خو \LR{$\Sigma$}
  عادي دے، نو \LR{$\Sigma \Proves \lnot\Diamond \lfalse$}
  (\olref[prf][nor]{prop:notDiamondBot})؛ له دې امله \LR{$\Delta
  \Proves[\Sigma] \lnot\Diamond \lfalse$} هم، چې د~\LR{$\Delta$}
  د سازګارۍ خلاف دے۔
\olpseganchor{OLP-0449-B010}{end}

\olpseganchor{OLP-0449-B011}{start}
  اوس فرض کړو چې \LR{$\Sigma$}، \Ax{T} لري، او \LR{$\Delta \in
  W^\Sigma$} واخلو۔ غواړو وښيو چې \LR{$R^\Sigma \Delta\Delta$}،
  يعنې
  \LR{$\Box^{-1}\Delta \subseteq \Delta$}۔ خو که \LR{$\Box!A \in \Delta$} وي، نو
  د \Ax{T} له مخې \LR{$!A \in \Delta$} هم، لکه چې غوښتل شوي وو۔
\olpseganchor{OLP-0449-B011}{end}

\olpseganchor{OLP-0449-B012}{start}
  اوس فرض کړو چې \LR{$\Sigma$}، \Ax{B} لري، او د \LR{$\Delta$}، \LR{$\Delta' \in W^\Sigma$}
  لپاره \LR{$R^\Sigma \Delta\Delta'$} فرض کړو۔ بايد وښيو
  چې \LR{$R^\Sigma \Delta'\Delta$}، يعنې
    \LR{$\Box^{-1}\Delta' \subseteq \Delta$}۔ د
    \olref[mod]{lem:box-iff-diamond} له مخې، دا له دې سره همارز دے چې
  \LR{$\Diamond\Delta \subseteq \Delta'$}۔ نو \LR{$!A \in \Delta$} فرض کړو۔ د
  \Ax{B} له مخې، \LR{$\Box\Diamond!A \in \Delta$} هم۔ له هغې فرضيې چې
  \LR{$R^\Sigma \Delta\Delta'$}، لرو چې \LR{$\Box^{-1}\Delta
  \subseteq \Delta'$}؛ نو \LR{$\Diamond!A \in \Delta'$}، لکه چې غوښتل شوي وو۔
\olpseganchor{OLP-0449-B012}{end}

\olpseganchor{OLP-0449-B013}{start}
  اوس فرض کړو چې \LR{$\Sigma$}، \Ax{4} لري، او \LR{$R^\Sigma
  \Delta_1\Delta_2$} او \LR{$R^\Sigma \Delta_2\Delta_3$} فرض کړو۔ بايد وښيو
  چې \LR{$R^\Sigma \Delta_1\Delta_3$}۔ له فرضيې دواړه \LR{$\Box^{-1}\Delta_1
  \subseteq \Delta_2$} او \LR{$\Box^{-1}\Delta_2 \subseteq \Delta_3$} لرو۔ د
  \LR{$R^\Sigma \Delta_1\Delta_3$} د ښودلو لپاره بس ده چې
  \LR{$\Box^{-1}\Delta_1 \subseteq \Delta_3$} وښيو۔ نو \LR{$!B \in
  \Box^{-1}\Delta_1$} واخلو، يعنې \LR{$\Box!B \in \Delta_1$}۔ د \Ax{4} له مخې،
  \LR{$\Box\Box!B \in \Delta_1$} هم؛ او له فرضيې لومړے دا ترلاسه کوو چې
  \LR{$\Box!B \in \Delta_2$}، او بيا دا چې \LR{$!B \in \Delta_3$}، لکه چې
  غوښتل شوي وو۔
\olpseganchor{OLP-0449-B013}{end}

\olpseganchor{OLP-0449-B014}{start}
  اوس فرض کړو چې \LR{$\Sigma$}، \Ax{5} لري، او \LR{$R^\Sigma
  \Delta_1\Delta_2$} او \LR{$R^\Sigma \Delta_1\Delta_3$} فرض کړو۔ بايد وښيو
  چې \LR{$R^\Sigma \Delta_2\Delta_3$}۔ لومړۍ فرضيه
  \LR{$\Box^{-1}\Delta_1 \subseteq \Delta_2$} راکوي، او دويمه فرضيه
  له دې سره همارزه ده چې \LR{$\Diamond\Delta_3 \subseteq \Delta_1$}،
    د \olref[mod]{lem:box-iff-diamond} له مخې۔ د \LR{$R^\Sigma
  \Delta_2\Delta_3$} د ښودلو لپاره، د
    \olref[mod]{lem:box-iff-diamond} له مخې، بس ده چې
  \LR{$\Diamond\Delta_3 \subseteq \Delta_2$} وښيو۔ نو \LR{$\Diamond!A \in
  \Diamond\Delta_3$} واخلو، يعنې \LR{$!A \in \Delta_3$}۔ د دويمې فرضيې له مخې
  \LR{$\Diamond!A \in \Delta_1$}؛ او د \Ax{5} له مخې، \LR{$\Box\Diamond!A \in
  \Delta_1$} هم۔ خو اوس لومړۍ فرضيه \LR{$\Diamond!A \in
  \Delta_2$} راکوي، لکه چې غوښتل شوي وو۔
  \end{proof}
\olpseganchor{OLP-0449-B014}{end}

\olpseganchor{OLP-0449-B015}{start}
د يوې نتيجې په توګه د څو نظامونو د بشپړتيا پايلې ترلاسه کوو۔
د بېلګې په توګه، پوهېږو چې \LR{$\Log{S5} = \Log{KT5} =
\Log{KTB4}$} د ټولو انعکاسي اقليدسي مدلونو د ټولګي په نسبت بشپړ دے؛
دا عين د ټولو انعکاسي، متناظرو او انتقالي مدلونو ټولګے دے۔
دغه پنځه چوکاټي خاصيتونه په يوې انعکاسي نړۍ لرونکي چوکاټ
کښې يوځاے پوره کېږي۔ نو د دغو اصلونو له هرې ټاکنې څخه جوړ
نظام د سموالي له مخې سازګار دے، او کانوني مدل ئې د اړوندو
ټولګيو په تقاطع کښې راځي۔
\olpseganchor{OLP-0449-B015}{end}

\olpseganchor{OLP-0449-B016}{start}
\begin{thm}\ollabel{thm:generaldet}
  پرېږدئ چې \LR{$\mClass{C}_\Ax{D}$}، \LR{$\mClass{C}_\Ax{T}$}،
  \LR{$\mClass{C}_\Ax{B}$}، \LR{$\mClass{C}_\Ax{4}$} او
  \LR{$\mClass{C}_\Ax{5}$} په ترتيب د ټولو مسلسل، انعکاسي،
  متناظرو، انتقالي او اقليدسي مدلونو ټولګي وي۔ بيا د
  \Ax{D}، \Ax{T}، \Ax{B}، \Ax{4} او \Ax{5} له منځه د هر ډول
  سکيماګانو \LR{$!A_1$}، \dots، \LR{$!A_n$} لپاره، نظام
  \LR{$\Log{K}!A_1 \dots !A_n$} د مدلونو د هغه ټولګي له خوا ټاکل کېږي چې
  \LR{$\mClass{C} = \mClass{C}_{!A_1} \cap \dots
  \cap \mClass{C}_{!A_n}$} دے۔
\end{thm}
\olpseganchor{OLP-0449-B016}{end}

\olpseganchor{OLP-0449-B017}{start}
\begin{prop}
  فرض کړئ چې \LR{$\Sigma$} يو سازګار عادي مودالي منطق دے۔ بيا:
  \begin{enumerate}
  \item\ollabel{prop:anotherfive-a}%
    که \LR{$\Sigma$} د \LR{$\Diamond!A \lif \Box
    !A$} سکيما ولري، نو د \LR{$\Sigma$} کانوني مدل جزوي تابعي دے۔
  \item که \LR{$\Sigma$} د \LR{$\Diamond!A \liff \Box
    !A$} سکيما ولري، نو د \LR{$\Sigma$} کانوني مدل تابعي دے۔
  \item که \LR{$\Sigma$} د \LR{$\Box\Box!A \lif \Box
    !A$} سکيما ولري، نو د \LR{$\Sigma$} کانوني مدل کمزورې ګڼه لري۔
  \end{enumerate}
(د چوکاټ د دغو خاصيتونو د تعريفونو لپاره \olref[frd][acc]{tab:anotherfive}
  وګورئ)۔
\end{prop}
\olpseganchor{OLP-0449-B017}{end}

\olpseganchor{OLP-0449-B018}{start}
\begin{proof}
  \begin{enumerate}
  \item فرض کړو چې \LR{$\Sigma$} د \LR{$\Diamond !A \lif
    \Box !A$} سکيما لري۔ د دې د ښودلو لپاره چې \LR{$R^\Sigma$} جزوي تابعي دے، بايد
    ثابت کړو چې د هر \LR{$\Delta_1$}، \LR{$\Delta_2$}، \LR{$\Delta_3 \in
    W^\Sigma$} لپاره، که \LR{$R^\Sigma \Delta_1\Delta_2$} او \LR{$R^\Sigma
    \Delta_1\Delta_3$} وي، نو \LR{$\Delta_2=\Delta_3$}۔ چونکې \LR{$R^\Sigma
    \Delta_1\Delta_2$}، نو \LR{$\Box^{-1}\Delta_1 \subseteq \Delta_2$} لرو؛
    او چونکې \LR{$R^\Sigma \Delta_1\Delta_3$}، نو \LR{$\Box^{-1}\Delta_1
    \subseteq \Delta_3$} هم۔ د \LR{$\Delta_2=\Delta_3$} عين‌والے
    هله ترلاسه کېږي چې دواړه شمولونه، \LR{$\Delta_2 \subseteq \Delta_3$} او \LR{$\Delta_3 \subseteq
    \Delta_2$}، ثابت کړو۔ د لومړي شمول لپاره، \LR{$!A \in \Delta_2$} واخلو؛ بيا
    \LR{$\Diamond!A \in \Delta_1$}، او د سکيما او د \LR{$\Delta_1$} د اشتقاقي تړلتيا
    له مخې \LR{$\Box!A \in \Delta_1$} هم؛ بيا له هغې فرضيې چې
    \LR{$R^\Sigma \Delta_1\Delta_3$}، \LR{$!A \in \Delta_3$}۔ دويم
    شمول هم په همدغه ډول دے۔
\olpseganchor{OLP-0449-B018}{end}

\olpseganchor{OLP-0449-B019}{start}
  \item د دوه‌اړخيزې طرحې يو لور د~\olref{prop:anotherfive-a}
    برخې له مخې جزوي تابعيت ورکوي۔ بل لور د مسلسل‌والي بديهي
    اصل دے؛ په \olref{thm:completeframeprops} کښې ئې ثبوت وګورئ۔
\olpseganchor{OLP-0449-B019}{end}

\olpseganchor{OLP-0449-B020}{start}
  \item فرض کړو چې \LR{$\Sigma$} د \LR{$\Box\Box!A \lif \Box!A$} سکيما لري۔
    د دې د ښودلو لپاره چې \LR{$R^\Sigma$} کمزورې ګڼه لري، \LR{$R^\Sigma
    \Delta_1\Delta_2$} واخلو۔ بايد وښيو چې داسې يو بشپړ
    \LR{$\Sigma$}-سازګار سټ \LR{$\Delta_3$} شته چې \LR{$R^\Sigma
    \Delta_1\Delta_3$} او \LR{$R^\Sigma \Delta_3\Delta_2$} وي۔ پرېږدئ چې:
    \[
    \Gamma = \Box^{-1}\Delta_1 \cup \Diamond\Delta_2.
    \]
    بس ده چې وښيو \LR{$\Gamma$} \LR{$\Sigma$}-سازګار دے، ځکه بيا
    د لينډنباوم لمې له مخې يو بشپړ
    \LR{$\Sigma$}-سازګار سټ~\LR{$\Delta_3$} ته غځېدے شي، داسې چې \LR{$\Box^{-1}\Delta_1
    \subseteq \Delta_3$} او \LR{$\Diamond\Delta_2 \subseteq \Delta_3$}،
    يعنې \LR{$R^\Sigma \Delta_1\Delta_3$} او \LR{$R^\Sigma
    \Delta_3\Delta_2$} (د
      \olref[mod]{lem:box-iff-diamond} له مخې)۔
\olpseganchor{OLP-0449-B020}{end}

\olpseganchor{OLP-0449-B021}{start}
    د تضاد لپاره فرض کړو چې \LR{$\Gamma$} سازګار نۀ دے۔ بيا
    داسې فارمولونه \LR{$\Box!A_1$}، \dots، \LR{$\Box!A_n \in \Delta_1$}
    او \LR{$!B_1$}، \dots،~\LR{$!B_m \in \Delta_2$} شته چې \[!A_1, \dots,
    !A_n, \Diamond!B_1, \dots, \Diamond!B_m \Proves[\Sigma]
    \lfalse.\] که د امکان لرونکې فرضيې هېڅ غړې پکې نه وي، د
    لومړۍ نړۍ د بکسونو دننني فارمولونه په دويمې نړۍ کښې دي،
    او د دروغو همدغه استنباط د دويمې نړۍ له سازګارۍ سره په ټکر
    کښې دے۔ د امکان د فرضيو په ناتشه حالت کښې،
    څرنګه چې \LR{$\Diamond (!B_1 \land \dots \land !B_m) \to
    (\Diamond!B_1 \land \dots \land \Diamond!B_m)$} په
    هر عادي مودالي منطق کښې د اشتقاق وړ دے، نو لاندې استدلال کوو، چې د~\LR{$\Delta_2$}
    له سازګارۍ سره تضاد لري:
    \begin{align*}
      !A_1, \dots, !A_n, & \Diamond!B_1, \dots,
      \Diamond!B_m \Proves[\Sigma] \lfalse \\
      !A_1, \dots,!A_n & \Proves[\Sigma] (\Diamond!B_1 \land
      \dots \land \Diamond!B_m) \lif \lfalse\\
      & \qquad\text{\RL{د استنتاج د قضيې له مخې}}\\
      & \qquad\text{\RL{\LR{\olref[prf][prp]{prop:derivabilityfacts}(\olref[prf][prp]{prop:derivabilityfacts-deduction})}، او \Taut}} \\
      !A_1, \dots,!A_n & \Proves[\Sigma]
      \Diamond(!B_1 \land
      \dots \land !B_m) \lif \lfalse\\
      & \qquad \text{\RL{چونکې \LR{$\Sigma$} عادي دے}} \\
      !A_1, \dots,!A_n & \Proves[\Sigma]
      \lnot\Diamond (!B_1 \land \dots \land !B_m)\\
      & \qquad\text{\RL{د \PL له مخې}} \\
      !A_1, \dots,!A_n & \Proves[\Sigma]
      \Box\lnot (!B_1 \land \dots \land !B_m)\\
      & \qquad\text{\RL{\LR{$\Box\lnot$} د \LR{$\lnot\Diamond$} په معنا}} \\
      \Box!A_1, \dots,\Box!A_n
      & \Proves[\Sigma] \Box\Box \lnot (!B_1 \land \dots \land !B_m)\\
      & \qquad\text{\RL{د \olref[mod]{lem:box1} له مخې}} \\
      \Box!A_1, \dots,\Box!A_n
      & \Proves[\Sigma]
      \Box\lnot (!B_1 \land \dots \land !B_m)\\
      &\qquad\text{\RL{د \LR{$\Box\Box!A \lif \Box!A$} سکيما له مخې}} \\
      \Delta_1 & \Proves[\Sigma]
      \Box\lnot (!B_1 \land \dots \land !B_m)\\
      &\qquad\text{\RL{د يو‌رنګه‌والي له مخې، \LR{\olref[prf][prp]{prop:derivabilityfacts}(\olref[prf][prp]{prop:derivabilityfacts-monotonicity})}}} \\
      & \Box\lnot (!B_1 \land \dots \land !B_m) \in \Delta_1\\
      &\qquad\text{\RL{د اشتقاقي تړلتيا له مخې}}; \\
      & \lnot (!B_1 \land \dots \land !B_m) \in \Delta_2\\
      &\qquad \text{\RL{چونکې }}
      R^\Sigma \Delta_1\Delta_2.
    \end{align*}
    د امکان د فرضيو ټول دننني فارمولونه د دويمې نړۍ غړي دي؛
    نو د اشتقاقي تړلتيا له مخې د هغو عطف هم غړے دے۔ وروستۍ
    کرښه د همدغه عطف د نفي غړيتوب ښيي، چې د دويمې نړۍ له
    سازګارۍ سره په ټکر کښې دے۔
  \end{enumerate}
\end{proof}
\olpseganchor{OLP-0449-B021}{end}

\olpseganchor{OLP-0449-B022}{start}
د دغو بېلګو پر بنسټ ښايي داسې وګڼل شي چې د مودالي منطق هر
نظام~\LR{$\Sigma$} \emph{بشپړ} دے، په دې معنا چې هر هغه فارمول ثابتوي چې په هر
هغه چوکاټ کښې اعتبار لري چې د \LR{$\Sigma$} هره قضيه پکښې اعتبار لري۔
له بده مرغه، ډېر نظامونه شته چې په دې معنا بشپړ نۀ دي۔
\olpseganchor{OLP-0449-B022}{end}

\olpseganchor{OLP-0450-B004}{start}
\olchapter{nml}{fil}{فلټريشنونه او د پرېکړې وړتيا}
\olpseganchor{OLP-0450-B004}{end}

\olpseganchor{OLP-0451-B004}{start}
\olfileid{nml}{fil}{int}
\olpseganchor{OLP-0451-B004}{end}

\olpseganchor{OLP-0451-B005}{start}
\olsection{سريزه}
\olpseganchor{OLP-0451-B005}{end}

\olpseganchor{OLP-0451-B006}{start}
د يوه منطق په اړه يوه مهمه پوښتنه دا وي چې آيا هغه د پرېکړې وړ دے، يعنې آيا
داسې اغېزمنه کړنلاره شته چې د «آيا دا فارمول اعتبار لري؟» پوښتنې ته ځواب
ورکړي۔ بياني منطق د پرېکړې وړ دے: د رښتياوالي د جدول په جوړولو په اغېزمن ډول
ازمېيلے شو چې آيا فارمول تاوتولوژي ده؛ او د ټاکلي فارمول لپاره
جدول متناهي دے۔ خو د مودالي فارمول د ټولو مدلونو کښې رښتياوالے د ټولو مدلونو په مستقيمې کتنې نۀ شو ازمېيلے،
ځکه چې مدلونه نامتناهي ډېر دي۔ د يوه ټاکلي فارمول لپاره
اړوند ټول متناهي مدلونه په لړ کښې راوستلے شو، ځکه يوازې د هغو
بياني متغيرونه لپاره د نړۍو د فرعي سټونو ټاکنه اړونده ده چې
رښتيا هم په فارمول کښې راځي۔ که د لاسرسي اړيکه ثابته وي، د \LR{$V(p)$}
بېلابېلې ممکنې ټاکنې د~\LR{$W$} ټول فرعي سټونه دي؛ او که \LR{$\card{W} = n$} وي،
نو د هغو شمېر \LR{$2^n$} دے۔ که زموږ فارمول~\LR{$!A$}، \LR{$m$} بياني متغيرونه ولري، نو پر ټاکلو~\LR{$n$} نړۍو او د لاسرسي پر همدې ثابتې اړيکې
\LR{$2^{nm}$} اړوندې ارزښت‌ټاکنې شته۔ د هر
مدل لپاره دا ازمېيلے شو چې آيا~\LR{$!A$} په ټولو نړۍو کښې رښتينے دے، يوازې په هره
نړۍ کښې د~\LR{$!A$} د رښتياوالي قيمت په محاسبه کولو سره۔ البته د لاسرسي ټولې ممکنې
اړيکې هم بايد وګورو، خو پر~\LR{$n$} نړۍو اړيکې هم متناهي ډېرې دي (په ځانګړي ډول
د \LR{$W \times W$} د فرعي سټونو شمېر، يعنې \LR{$2^{n^2}$})۔
\olpseganchor{OLP-0451-B006}{end}

\olpseganchor{OLP-0451-B007}{start}
که زموږ پام منطق~\Log{K} ته نۀ وي، بلکې داسې منطق ته وي چې د مدلونو د يوه
ټولګي په وسيله تعريف شوے وي (د بېلګې په توګه، انعکاسي متعدي مدلونه)، نو بايد
دا هم ازمېيلے شو چې د لاسرسي اړيکه د غوښتل شوي ډول ده که نه۔ کله چې زموږ
چوکاټونه په يوه ټاکلي متناهي لړ د مودالي فارمولونه تعريف شي، دا
ازموينه کولے شو
(د بېلګې په توګه، په چوکاټ کښې د \Ax{T} او \Ax{4} د اعتبار په ازمېيلو)۔ نو
لار دا ښکاري چې ټول متناهي چوکاټونه په لړ کښې واخلو، هر يو وازمېيو چې آيا
د غوښتل شوي ټولګي چوکاټ دے، بيا پر هغه چوکاټ ټول ممکن مدلونه راوړو او
وازمېيو چې آيا \LR{$!A$} په هر يوه کښې رښتينے دے۔ که داسې نۀ وي، ودرېږو: \LR{$!A$}
په غوښتل شوي مدل‌ټولګي کښې اعتبار نۀ لري۔
\olpseganchor{OLP-0451-B007}{end}

\olpseganchor{OLP-0451-B008}{start}
په دې لاره کښې ستونزه شته: نۀ پوهېږو چې لټون کله، او که هېڅکله، درولے شو۔
که فارمول متناهي ضدمدل ولري، کړنلاره به ئې ومومي۔ خو که متناهي ضدمدل
ونۀ لري، ځواب به رانۀ کړي۔ کېدے شي فارمول معتبر وي (يعنې هېڅ ضدمدل
ونۀ لري)، يا يوازې نامتناهي ضدمدل ولري چې موږ ئې هېڅکله نۀ ګورو۔ دا ستونزه
هله هوارېدے شي چې وښيو هر هغه فارمول چې ضدمدل لري، متناهي ضدمدل هم
لري۔ په دې صورت کښې وايو چې منطق د \emph{متناهي مدل خاصيت} لري۔
\olpseganchor{OLP-0451-B008}{end}

\olpseganchor{OLP-0451-B009}{start}
خو څنګه وښيو چې يو منطق د متناهي مدل خاصيت لري؟ يوه لار دا ده چې د~\LR{$!A$}
يو نامتناهي (ضد)مدل په متناهي مدل بدل کړو۔ که دا کار وشي، نو هر کله چې داسې
مدل وي چې \LR{$!A$} پکښې رښتينے نۀ وي، له هغه نه جوړ شوے متناهي مدل به هم \LR{$!A$}
نارښتينے کړي۔ دا متناهي مدل به زموږ د ټولو متناهي مدلونو په لړ کښې راشي؛ نو
د هر نامعتبر فارمول په اړه به بالاخره معلومه کړو چې نامعتبر دے۔ خو که
فارمول معتبر وي، دا لټونيزه کړنلاره به نۀ درېږي۔ که داسې محاسبه‌کېدونکے
حد هم ولرو چې زموږ له جوړونې راوتلے متناهي ضدمدل ترې لوے نۀ وي او دغه حد
يوازې په فارمول~\LR{$!A$} پورې تړلے وي، نو د ضدمدل د لټون لپاره به پوه شو
چې تر کومې کچې متناهي مدلونه وازمېيو۔ که تر هغه وخته ضدمدل ونۀ موندل شي،
هېڅ ضدمدل نشته۔ بيا به زموږ کړنلاره د هر فارمول~\LR{$!A$} لپاره د «آيا
\LR{$!A$} اعتبار لري؟» پوښتنه په رښتيا پرېکړه کړي۔
\olpseganchor{OLP-0451-B009}{end}

\olpseganchor{OLP-0451-B010}{start}
د نامتناهي جوړښتونو په متناهي جوړښتونو د بدلولو يوه ګټوره تګلاره دا ده چې
د جوړښت هغه غړي «يو شان» وګڼو چې له اړوندو پلوو يو شان چلند کوي۔
که په~\LR{$\mModel{M}$} کښې نامتناهي ډېرې نړۍ له اړوندو پلوو يو شان چلند کوي،
کېدے شي \emph{ټولې} نړۍ د دغو نړۍو په متناهي شمېر ټولګيو کښې راټولې کړو؛
هر ټولګے په خپله متناهي يا نامتناهي کېدے شي۔ د نړۍو سټ په داسې
متناهي ډېرو ټولګيو وېشو چې دوه دوه بېل وي او يو ځاے ټولې نړۍ راونغاړي۔ بيا يو نوے مدل~\LR{$\mModel{M^*}$} تعريفوو چې نړۍې
ئې همدا ټولګي وي۔ د پخواني مدل متناهي ډېر ټولګي په نوي مدل کښې متناهي ډېرې
نړۍې راکوي، يعنې متناهي مدل۔ د هغې وېش‌برخې نوم چې نړۍ~\LR{$w$} پکښې ده، \LR{$[w]$}
کېږدو۔ غواړو \LR{$\mSat{M}{!A}[w]$} هله او يوازې هله وي چې
\LR{$\mSat{M^*}{!A}[{[w]}]$} وي، ځکه که پخوانے مدل د~\LR{$!A$} ضدمدل وي، نو
نوے مدل هم بايد ضدمدل پاتې شي۔ د دې لپاره وېش او د~\LR{$\mModel{M^*}$}
د لاسرسي اړيکه په سمه توګه تعريفول پکار دي۔
\olpseganchor{OLP-0451-B010}{end}

\olpseganchor{OLP-0451-B011}{start}
د دې د څرنګوالي د ليدو لپاره، لومړے داسې ساده حالت وګڼو چې هره نړۍ هرې
نړۍ ته لاسرسی لري۔ په دې حالت کښې \LR{$\mSat{M}{\Box !B}[w]$} هله او يوازې
هله چې د هر \LR{$v \in W$} لپاره \LR{$\mSat{M}{!B}[v]$}؛ په~\LR{$\mModel{M^*}$}
کښې هم لاسرسے عمومي نيسو او نړۍې ئې د \LR{$[w]$} او \LR{$[v]$} په شان ټولګي
دي۔ دوه نړۍې \LR{$u$} او \LR{$v$} هله همارزې دي (په يوه وېش‌برخه کښې دي) چې
په~\LR{$\mModel{M}$} کښې د پام وړ فارمول پر ټولو راغلو بياني متغيرونه
يو شان رښتياوالے ولري، يعنې د دغو متغيرونو د هر يوه لپاره
\LR{$\mSat{M}{p}[u]$} هله او يوازې هله چې \LR{$\mSat{M}{p}[v]$}۔ پرېږدئ چې \LR{$V^*(p) = \Setabs{[w]}{\mSat{M}{p}[w]}$}۔
غواړو وښيو چې \LR{$\mSat{M}{!A}[w]$} هله او يوازې هله چې
\LR{$\mSat{M^*}{!A}[{[w]}]$}۔ دا په استقرا ثابتوو: بنسټيز صورت \LR{$!A
\equiv p$} دے۔ لومړے فرض کړو چې \LR{$\mSat{M}{p}[w]$}۔ نو د تعريف له مخې
\LR{$[w] \in V^*(p)$}، او ځکه \LR{$\mSat{M^*}{p}[{[w]}]$}۔ اوس فرض کړو چې
\LR{$\mSat{M^*}{p}[{[w]}]$}۔ دا معنا لري چې \LR{$[w] \in V^*(p)$}، يعنې د~\LR{$w$}
سره همارزې يوې نړۍ~\LR{$v$} لپاره \LR{$\mSat{M}{p}[v]$}۔ خو «\LR{$w$} له \LR{$v$}
سره همارز دے» دا معنا لري چې «\LR{$w$} او \LR{$v$} له اړوندو بياني متغيرونه
څخه هماغه متغيرونه رښتيني کوي»؛ نو \LR{$\mSat{M}{p}[w]$}۔ اوس استقرايي ګام، د بېلګې په توګه،
\LR{$!A \ident \lnot !B$}۔ بيا \LR{$\mSat{M}{\lnot !B}[w]$} هله او يوازې هله چې
\LR{$\mSat/{M}{!B}[w]$}، هله او يوازې هله چې \LR{$\mSat/{M^*}{!B}[{[w]}]$}
(د استقرا فرضيې له مخې)، هله او يوازې هله چې
\LR{$\mSat{M^*}{\lnot !B}[{[w]}]$}۔ نور نامودالي عملګر هم همداسې دي۔ د
\LR{$\Box$} لپاره هم کار کوي: فرض کړو \LR{$\mSat{M^*}{\Box
  !B}[{[w]}]$}۔ دا معنا لري چې د هر \LR{$[u]$} لپاره
\LR{$\mSat{M^*}{!B}[{[u]}]$}۔ د استقرا فرضيې له مخې، د هر \LR{$u$} لپاره
\LR{$\mSat{M}{!B}[u]$}۔ نو \LR{$\mSat{M}{\Box !B}[w]$}۔ په سرچپه
لوري کښې، د عمومي لاسرسي له مخې دنننے فارمول په هره اصلي نړۍ کښې
رښتينے دے۔ د استقرا فرضيه د هغه صدق هر ټولګي ته لېږدوي، نو بکس په
نوي مدل کښې هم رښتينے دے۔
\olpseganchor{OLP-0451-B011}{end}

\olpseganchor{OLP-0451-B012}{start}
په عمومي صورت کښې، چې د~\LR{$\mModel{M^*}$} د لاسرسي اړيکه هم بايد تعريف
کړو، چارې پېچلې دي۔ يو مدل~\LR{$\mModel{M^*}$} ته \emph{فلټريشن} وايو که
اړيکه~\LR{$R^*$} ئې هغه شرطونه پوره کړي چې پورته استقرايي ثبوت ته اړ دي۔
بيا هر فلټريشن~\LR{$\mModel{M^*}$}، \LR{$!A$} په \LR{$[w]$} کښې هله او يوازې هله
رښتينے کوي چې \LR{$\mModel{M}$}، \LR{$!A$} په~\LR{$w$} کښې رښتينے کوي۔ خو اوس بايد دا
هم وښيو چې فلټريشنونه \emph{شته}، يعنې~\LR{$R^*$} داسې تعريفولے شو چې اړين
شرطونه پوره کړي۔ د دې کار لپاره دا هم لازم دي چې نړۍې \LR{$u$} او \LR{$v$} يوازې
د اړوندو بياني متغيرونه د يو شان رښتياوالي له مخې
همارزې نۀ ګڼو؛ د~\LR{$!A$} د ټولو فرعي فارمولونه رښتياوالے هم بايد
پکښې يو شان وي۔ ځکه چې \LR{$!A$} يوازې متناهي ډېر فرعي فارمولونه
لري، دا بيا هم د فلټريشن متناهي‌والے تضمينوي۔ د فلټريشن د تعريف يوازې
يوه لار نشته؛ او د دې لپاره چې د فلټريشن د لاسرسي اړيکه غوښتل شوي
خاصيتونه (لکه انعکاس، تعدّي او نور) ولري، د~\LR{$R^*$} په تعريف کښې
له فکره کار اخيستل پکار دي۔
\olpseganchor{OLP-0451-B012}{end}

\olpseganchor{OLP-0452-B004}{start}
\olfileid{nml}{fil}{pre}
\olpseganchor{OLP-0452-B004}{end}

\olpseganchor{OLP-0452-B005}{start}
\olsection{لومړني تعريفونه}
\olpseganchor{OLP-0452-B005}{end}

\olpseganchor{OLP-0452-B006}{start}
فلټريشنونه راته دا لار برابروي چې د مودالي منطق د خپلو نظامونو د پرېکړې
وړتيا، د هغوی د \emph{متناهي مدل خاصيت} په ښودلو سره، ثابته کړو؛ يعنې هر
هغه فارمول چې د نظام د يوه مدل په يوه نړۍ کښې رښتينے (نارښتينے)
وي، د هماغه نظام د يوه \emph{متناهي} مدل په يوه نړۍ کښې هم رښتينے
(نارښتينے) وي۔ فلټريشنونه د هغو
فارمولونه د سټونو په نسبت تعريفېږي چې د فرعي فارمولونو لپاره تړلي وي۔
\olpseganchor{OLP-0452-B006}{end}

\olpseganchor{OLP-0452-B007}{start}
\begin{defn}\ollabel{defn:modallyclosed}
  د فارمولونه يو سټ~\LR{$\Gamma$} د \emph{فرعي فارمولونو لپاره تړلے} دے که
  په~\LR{$\Gamma$} کښې د هر فارمول هر فرعي فارمول هم ولري۔ سربېره پر
  دې، \LR{$\Gamma$} \emph{مودالي تړلے} دے که د فرعي فارمولونو لپاره تړلے وي
  او \LR{$!A \in \Gamma$} دا هم لازم کړي چې \LR{$\Box!A,
  \Diamond!A \in \Gamma$}۔
\end{defn}
\olpseganchor{OLP-0452-B007}{end}

\olpseganchor{OLP-0452-B008}{start}
د بېلګې په توګه، د يوه فارمول~\LR{$!A$} ټول فرعي فارمولونه
د فرعي فارمولونه لپاره تړلے سټ جوړوي۔ کله چې د~\LR{$!A$} د فرعي
فارمولونه د سټ له مخې د يوه مدل فلټريشن تعريفوو، نو دا به زموږ
غوښتلے خاصيت ولري: \LR{$!A$} په هر نوي ټولګي کښې هله او يوازې هله
رښتينے (نارښتينے) وي چې د هغه ټولګي په اصلي نړۍ کښې رښتينے (نارښتينے)
وي۔
\olpseganchor{OLP-0452-B008}{end}

\olpseganchor{OLP-0452-B009}{start}
د~\LR{$\Gamma$} له مخې د~\LR{$\mModel{M}$} د فلټريشن د نړۍو سټ د لاندې
هم‌ارزښتۍ د اړيکې د ټولو ټولګيو د سټ په توګه تعريفېږي۔
\olpseganchor{OLP-0452-B009}{end}

\olpseganchor{OLP-0452-B010}{start}
\begin{defn}
فرض کړئ \LR{$\mModel{M} =\tuple{W, R, V}$} او \LR{$\Gamma$} د فرعي
فارمولونه لپاره تړلے دے۔ پر~\LR{$W$} اړيکه~\LR{$\equiv$} داسې تعريفوو چې
دوه نړۍې هله پکښې په اړيکه وي چې د~\LR{$\Gamma$} هماغه فارمولونه
رښتيني کړي، يعنې:
\[
u \equiv v \quad \text{\RL{هله او يوازې هله چې }} \quad
\forall !A \in \Gamma : \mSat{M}{!A}[u] \Leftrightarrow \mSat{M}{!A}[v].
\]
د نړۍ~\LR{$w$} د هم‌ارزښتۍ ټولګے~\LR{$[w]_\equiv$}، يا په لنډ ډول \LR{$[w]$}،
د هغو ټولو نړۍو سټ دے چې له~\LR{$w$} سره د~\LR{$\equiv$} له مخې همارزې دي:
\[
[w] = \Setabs{v}{v \equiv w}.
\]
\end{defn}
\olpseganchor{OLP-0452-B010}{end}

\olpseganchor{OLP-0452-B011}{start}
\begin{prop}
  د~\LR{$\mModel{M}$} او~\LR{$\Gamma$} لپاره، پورته تعريف شوې اړيکه~\LR{$\equiv$}
  د هم‌ارزښتۍ اړيکه ده، يعنې انعکاسي، متناظره او متعديه ده۔
\end{prop}
\olpseganchor{OLP-0452-B011}{end}

\olpseganchor{OLP-0452-B012}{start}
\begin{proof}
  اړيکه~\LR{$\equiv$} انعکاسي ده، ځکه \LR{$w$} د~\LR{$\Gamma$} عين هماغه
  فارمولونه رښتيني کوي چې پخپله ئې رښتيني کوي۔ متناظره ده، ځکه
  که \LR{$u$} د~\LR{$\Gamma$} هماغه فارمولونه رښتيني کړي چې \LR{$v$} ئې
  رښتيني کوي، نو دا خبره د~\LR{$v$} او~\LR{$u$} لپاره هم سمه ده۔ متعديه هم ده:
  که \LR{$u$} د~\LR{$\Gamma$} هماغه فارمولونه رښتيني کړي چې \LR{$v$} ئې کوي،
  او \LR{$v$} هماغه فارمولونه رښتيني کړي چې \LR{$w$} ئې کوي، نو \LR{$u$} به
  د~\LR{$\Gamma$} هماغه فارمولونه رښتيني کړي چې \LR{$w$} ئې کوي۔
\end{proof}
\olpseganchor{OLP-0452-B012}{end}

\olpseganchor{OLP-0452-B013}{start}
اړيکه~\LR{$\equiv$}، لکه هره هم‌ارزښتۍ اړيکه، \LR{$W$} په \emph{وېش‌برخو}
وېشي، يعنې د~\LR{$W$} په هغو فرعي سټونو چې دوه دوه بېل وي او يوځای ټول~\LR{$W$}
وپوښي۔ هره نړۍ \LR{$w \in W$} د يوې وېش‌برخې غړے ده، يعنې
د~\LR{$[w]$}، ځکه \LR{$w \equiv w$}۔ نو وېش‌برخې \LR{$[w]$} ټول~\LR{$W$} پوښي۔ دا
وېش‌برخې دوه دوه بېلې دي: که \LR{$u \in [w]$} او \LR{$u \in [v]$} وي، نو
\LR{$u \equiv w$} او \LR{$u \equiv v$}؛ د تقارن او تعدّي له مخې
\LR{$w \equiv v$}، او ځکه \LR{$[w] = [v]$}۔
\olpseganchor{OLP-0452-B013}{end}

\olpseganchor{OLP-0453-B004}{start}
\olfileid{nml}{fil}{fil}
\olpseganchor{OLP-0453-B004}{end}

\olpseganchor{OLP-0453-B005}{start}
\olsection{فلټريشنونه}
\olpseganchor{OLP-0453-B005}{end}

\olpseganchor{OLP-0453-B006}{start}
د~\LR{$\Gamma$} له مخې د~\LR{$\mModel{M}$} «يوازينے» فلټريشن نۀ تعريفوو؛
بلکې دا تعريفوو چې يو مدل~\LR{$\mModel{M^*}$} کله د~\LR{$\mModel{M}$}
فلټريشن ګڼل کېږي۔ ټول فلټريشنونه د نړۍو عين سټ~\LR{$W^*$} او عين
ارزښت‌تابع~\LR{$V^*$} لري۔ خو بېلابېل فلټريشنونه بېلابېلې د لاسرسي
اړيکې~\LR{$R^*$} لرلے شي۔ سره له دې، د فلټريشن ګڼل کېدو لپاره \LR{$R^*$}
بايد څو شرطونه پوره کړي۔ دا عين هغه شرطونه دي چې د اصلي پايلې د
ثبوت لپاره ورته اړ يو: که \LR{$!A \in \Gamma$} وي، نو
\LR{$\mSat{M}{!A}[w]$} هله او يوازې هله وي چې
\LR{$\mSat{M^*}{!A}[{[w]}]$} وي۔
\olpseganchor{OLP-0453-B006}{end}

\olpseganchor{OLP-0453-B007}{start}
\begin{defn}\ollabel{defn:filtration}
  فرض کړئ \LR{$\Gamma$} د فرعي فارمولونو لپاره تړلے دے او
  \LR{$\mModel{M} = \tuple{W,
    R, V}$}۔ د~\LR{$\Gamma$} له مخې د~\LR{$\mModel{M}$} يو
  \emph{فلټريشن} هر هغه مدل
  \LR{$\mModel{M^*} = \tuple{W^*,R^*,V^*}$} دے چې:
  \begin{enumerate}
  \item \LR{$W^* = \Setabs{[w]}{w \in W}$}؛
  \item \ollabel{defn:filtration-R}%
    د هر \LR{$u,v \in W$} لپاره:
    \begin{enumerate}
    \item \ollabel{defn:filtration-R1}%
      که \LR{$Ruv$} وي، نو \LR{$R^*[u][v]$}؛
    \item \ollabel{defn:filtration-R2}%
      که \LR{$R^*[u][v]$} وي، نو د هر \LR{$\Box !A \in \Gamma$} لپاره، که
      \LR{$\mSat{M}{\Box!A}[u]$} وي، نو \LR{$\mSat{M}{!A}[v]$}
    \item \ollabel{defn:filtration-R3}%
      که \LR{$R^*[u][v]$} وي، نو د هر \LR{$\Diamond !A \in \Gamma$} لپاره، که
      \LR{$\mSat{M}{!A}[v]$} وي، نو \LR{$\mSat{M}{\Diamond!A}[u]$}۔
    \end{enumerate}
  \item \LR{$V^*(p) = \Setabs{[u]}{u \in V(p)}$}۔
  \end{enumerate}
\end{defn}
\olpseganchor{OLP-0453-B007}{end}

\olpseganchor{OLP-0453-B008}{start}
د \LR{$V^*(p)$} په معنا فکر کول ګټور دي: دا د هغو ټولو نړۍو~\LR{$w$} د
هم‌ارزښتۍ د ټولګيو~\LR{$[w]$} سټ دے چې \LR{$p$} په~\LR{$\mModel{M}$} کښې
پکښې رښتينے وي۔ که \LR{$w \in V(p)$} وي، نو د تعريف له مخې
\LR{$[w] \in V^*(p)$}۔ خو که \LR{$[w] \in V^*(p)$} وي، لازم نۀ ده چې
\LR{$w \in V(p)$} هم وي۔ له \LR{$[w] \in V^*(p)$} څخه يوازې دا ډاډېږي چې
د \emph{يوې} نړۍ~\LR{$u \in V(p)$} لپاره \LR{$[w] = [u]$} وي۔ البته
\LR{$[w] = [u]$} دا معنا لري چې \LR{$w \equiv u$}۔ نو کله چې
\LR{$[w] \in V^*(p)$} وي، يوازې دا نتيجه اخيستے شو چې د يوې
\LR{$u \in V(p)$} لپاره \LR{$w \equiv u$} وي۔
\olpseganchor{OLP-0453-B008}{end}

\olpseganchor{OLP-0453-B009}{start}
\begin{thm}\ollabel{thm:filtrations}
  که \LR{$\mModel{M^*}$} د~\LR{$\Gamma$} له مخې د~\LR{$\mModel{M}$} فلټريشن
  وي، نو د هر \LR{$!A \in \Gamma$} او \LR{$w \in W$} لپاره
  \LR{$\mSat{M}{!A}[w]$} هله او يوازې هله چې
  \LR{$\mSat{M^*}{!A}[{[w]}]$}۔
\end{thm}
\olpseganchor{OLP-0453-B009}{end}

\olpseganchor{OLP-0453-B010}{start}
\begin{proof}
  پر~\LR{$!A$} استقرا کوو او دا کاروو چې \LR{$\Gamma$} د فرعي فارمولونو
  لپاره تړلے دے۔ ځکه \LR{$!A \in \Gamma$} او \LR{$\Gamma$} د فرعي
  فارمولونه لپاره تړلے دے، د~\LR{$!A$} ټول فرعي فارمولونه هم
  \LR{$\in \Gamma$} دي۔ نو په هر استقرايي ګام کښې د~\LR{$!A$} پر فرعي
  فارمولونه د استقرا فرضيه پلي کېږي۔
  \begin{enumerate}
    \item \indcase{!A}{\lfalse}{نه
        \LR{$\mSat{M}{\indfrm}[w]$} او نه \LR{$\mSat{M^*}{\indfrm}[{[w]}]$}۔}
    \item \indcase{!A}{\ltrue}{دواړه
        \LR{$\mSat{M}{\indfrm}[w]$} او \LR{$\mSat{M^*}{\indfrm}[{[w]}]$}۔}
    \item \indcase{!A}{p}{له کيڼ نه ښي لوري ته خبره سمدستي ده:
      \LR{$\mSat{M}{\indfrm}[w]$} يوازې هله وي چې \LR{$w \in V(p)$} وي؛
      له دې نه \LR{$[w]
      \in V^*(p)$}، يعنې \LR{$\mSat{M^*}{\indfrm}[{[w]}]$} راځي۔
      برعکس، فرض کړئ \LR{$\mSat{M^*}{\indfrm}[{[w]}]$}، يعنې \LR{$[w] \in
      V^*(p)$}۔ نو د يوې \LR{$v \in V(p)$} لپاره \LR{$w \equiv v$}۔ البته
      \LR{$\mSat{M}{p}[v]$} هم وي۔ ځکه \LR{$w \equiv v$}، نړۍې \LR{$w$} او \LR{$v$}
      د~\LR{$\Gamma$} هماغه فارمولونه رښتيني کوي۔ د فرض له مخې \LR{$p
      \in \Gamma$} او \LR{$\mSat{M}{p}[v]$} دي؛ نو
      \LR{$\mSat{M}{\indfrm}[w]$}۔}
\olpseganchor{OLP-0453-B010}{end}

\olpseganchor{OLP-0453-B011}{start}
    \item \indcase{!A}{\lnot
          !B}{\LR{$\mSat{M}{\indfrm}[w]$} هله او يوازې هله چې
          \LR{$\mSat/{M}{!B}[w]$}۔ د استقرا فرضيې له مخې
          \LR{$\mSat/{M}{!B}[w]$} هله او يوازې هله چې
          \LR{$\mSat/{M^*}{!B}[{[w]}]$}۔ په پای کښې
          \LR{$\mSat/{M^*}{!B}[{[w]}]$} هله او يوازې هله چې
          \LR{$\mSat{M^*}{\indfrm}[{[w]}]$}۔}
\olpseganchor{OLP-0453-B011}{end}

\olpseganchor{OLP-0453-B012}{start}
    \item \indcase{!A}{(!B \land
          !C)}{\LR{$\mSat{M}{\indfrm}[w]$} هله او يوازې هله چې
          \LR{$\mSat{M}{!B}[w]$} او \LR{$\mSat{M}{!C}[w]$} دواړه وي۔ د استقرا
          فرضيې له مخې \LR{$\mSat{M}{!B}[w]$} هله او يوازې هله چې
          \LR{$\mSat{M^*}{!B}[{[w]}]$}، او \LR{$\mSat{M}{!C}[w]$} هله او يوازې
          هله چې \LR{$\mSat{M^*}{!C}[{[w]}]$}۔ او
          \LR{$\mSat{M^*}{\indfrm}[{[w]}]$} هله او يوازې هله چې
          \LR{$\mSat{M^*}{!B}[{[w]}]$} او
          \LR{$\mSat{M^*}{!C}[{[w]}]$} دواړه وي۔}
\olpseganchor{OLP-0453-B012}{end}

\olpseganchor{OLP-0453-B013}{start}
    \item \indcase{!A}{(!B \lor
          !C)}{\LR{$\mSat{M}{\indfrm}[w]$} هله او يوازې هله چې
          \LR{$\mSat{M}{!B}[w]$} يا \LR{$\mSat{M}{!C}[w]$} وي۔ د استقرا
          فرضيې له مخې \LR{$\mSat{M}{!B}[w]$} هله او يوازې هله چې
          \LR{$\mSat{M^*}{!B}[{[w]}]$}، او \LR{$\mSat{M}{!C}[w]$} هله او يوازې
          هله چې \LR{$\mSat{M^*}{!C}[{[w]}]$}۔ او
          \LR{$\mSat{M^*}{\indfrm}[{[w]}]$} هله او يوازې هله چې
          \LR{$\mSat{M^*}{!B}[{[w]}]$} يا
          \LR{$\mSat{M^*}{!C}[{[w]}]$} وي۔}
\olpseganchor{OLP-0453-B013}{end}

\olpseganchor{OLP-0453-B014}{start}
    \item \indcase{!A}{(!B \lif
          !C)}{\LR{$\mSat{M}{\indfrm}[w]$} هله او يوازې هله چې
          \LR{$\mSat/{M}{!B}[w]$} يا \LR{$\mSat{M}{!C}[w]$} وي۔ د استقرا
          فرضيې له مخې \LR{$\mSat{M}{!B}[w]$} هله او يوازې هله چې
          \LR{$\mSat{M^*}{!B}[{[w]}]$}، او \LR{$\mSat{M}{!C}[w]$} هله او يوازې
          هله چې \LR{$\mSat{M^*}{!C}[{[w]}]$}۔ او
          \LR{$\mSat{M^*}{\indfrm}[{[w]}]$} هله او يوازې هله چې
          \LR{$\mSat/{M^*}{!B}[{[w]}]$} يا
          \LR{$\mSat{M^*}{!C}[{[w]}]$} وي۔}
\olpseganchor{OLP-0453-B014}{end}

\olpseganchor{OLP-0453-B015}{start}
    \item \indcase{!A}{(!B \liff
          !C)}{\LR{$\mSat{M}{\indfrm}[w]$} هله او يوازې هله چې
          \LR{$\mSat{M}{!B}[w]$} او \LR{$\mSat{M}{!C}[w]$} دواړه وي، يا
          \LR{$\mSat/{M}{!B}[w]$} او \LR{$\mSat/{M}{!C}[w]$} دواړه وي۔ د
          استقرا فرضيې له مخې \LR{$\mSat{M}{!B}[w]$} هله او يوازې هله چې
          \LR{$\mSat{M^*}{!B}[{[w]}]$}، او \LR{$\mSat{M}{!C}[w]$} هله او يوازې
          هله چې \LR{$\mSat{M^*}{!C}[{[w]}]$}۔ او
          \LR{$\mSat{M^*}{\indfrm}[{[w]}]$} هله او يوازې هله چې
          \LR{$\mSat{M^*}{!B}[{[w]}]$} او
          \LR{$\mSat{M^*}{!C}[{[w]}]$} دواړه وي، يا
          \LR{$\mSat/{M^*}{!B}[{[w]}]$} او
          \LR{$\mSat/{M^*}{!C}[{[w]}]$} دواړه وي۔}
\olpseganchor{OLP-0453-B015}{end}

\olpseganchor{OLP-0453-B016}{start}
    \item \indcase{!A}{\Box
          !B}{فرض کړئ \LR{$\mSat{M}{\indfrm}[w]$}۔ د دې د ښودلو لپاره
          چې \LR{$\mSat{M^*}{\indfrm}[{[w]}]$}، داسې \LR{$v$} واخلئ چې
          \LR{$R^*[w][v]$} وي۔ د
          \LR{\olref{defn:filtration}(\olref{defn:filtration-R2})} له مخې
          \LR{$\mSat{M}{!B}[v]$} وي، او د استقرا فرضيې له مخې
          \LR{$\mSat{M^*}{!B}[{[v]}]$}۔ ځکه \LR{$v$} هره يوه نړۍ کېدے شوه،
          \LR{$\mSat{M^*}{\indfrm}[{[w]}]$} ثابت شو۔
\olpseganchor{OLP-0453-B016}{end}

\olpseganchor{OLP-0453-B017}{start}
          برعکس، فرض کړئ \LR{$\mSat{M^*}{\indfrm}[{[w]}]$} او داسې هره
          \LR{$v$} واخلئ چې \LR{$Rwv$} وي۔ د
          \LR{\olref{defn:filtration}(\olref{defn:filtration-R1})} له مخې
          \LR{$R^*[w][v]$}، نو \LR{$\mSat{M^*}{!B}[{[v]}]$} وي؛ د استقرا
          فرضيې له مخې \LR{$\mSat{M}{!B}[v]$}، او ځکه \LR{$v$} اختياري وه،
          \LR{$\mSat{M}{\indfrm}[w]$} ثابت شو۔}
\olpseganchor{OLP-0453-B017}{end}

\olpseganchor{OLP-0453-B018}{start}
    \item \indcase{!A}{\Diamond
          !B}{فرض کړئ \LR{$\mSat{M}{\indfrm}[w]$}۔ نو د يوې \LR{$v \in W$}
          لپاره \LR{$Rwv$} او \LR{$\mSat{M}{!B}[v]$} وي۔ د استقرا فرضيې له
          مخې \LR{$\mSat{M^*}{!B}[{[v]}]$}، او د
          \LR{\olref{defn:filtration}(\olref{defn:filtration-R1})} له مخې
          \LR{$R^*[w][v]$}۔ نو \LR{$\mSat{M^*}{\indfrm}[{[w]}]$}۔
\olpseganchor{OLP-0453-B018}{end}

\olpseganchor{OLP-0453-B019}{start}
          اوس فرض کړئ \LR{$\mSat{M^*}{\indfrm}[{[w]}]$}۔ نو د يوې \LR{$[v]
          \in W^*$} لپاره، چې \LR{$R^*[w][v]$} وي،
          \LR{$\mSat{M^*}{!B}[{[v]}]$} وي۔ د استقرا فرضيې له مخې
          \LR{$\mSat{M}{!B}[v]$}۔ د
          \LR{\olref{defn:filtration}(\olref{defn:filtration-R3})} له مخې
          \LR{$\mSat{M}{\indfrm}[w]$} وي۔}
  \end{enumerate}
\end{proof}
\olpseganchor{OLP-0453-B019}{end}

\olpseganchor{OLP-0453-B020}{start}
\begin{prob}
  د \olref[nml][fil][fil]{thm:filtrations} ثبوت بشپړ کړئ۔
\end{prob}
\olpseganchor{OLP-0453-B020}{end}

\olpseganchor{OLP-0453-B021}{start}
هغه څه چې په يوه مدل کښې د نړۍو د صدق په اړه سم دي، په يوه مدل کښې په ټولو نړۍو د صدق
او د مدلونو په يوه ټولګي کښې د اعتبار په اړه هم سم دي۔
\olpseganchor{OLP-0453-B021}{end}

\olpseganchor{OLP-0453-B022}{start}
\begin{cor}
  فرض کړئ \LR{$\Gamma$} د فرعي فارمولونو لپاره تړلے دے۔ بيا:
  \begin{enumerate}
  \item که \LR{$\mModel{M^*}$} د~\LR{$\Gamma$} له مخې د~\LR{$\mModel{M}$}
    فلټريشن وي، نو د هر \LR{$!A \in \Gamma$} لپاره: \LR{$\mSat{M}{!A}$} هله
    او يوازې هله چې \LR{$\mSat{M^*}{!A}$}۔
  \item که \LR{$\mClass{C}$} د مدلونو يو ټولگے وي، او
    \LR{$\Gamma(\mClass{C})$} په~\LR{$\mClass{C}$} کښې د مدلونو د
    \LR{$\Gamma$}-فلټريشنونو ټولگے وي، نو هر فارمول~\LR{$!A \in \Gamma$}
    په~\LR{$\mClass{C}$} کښې هله او يوازې هله معتبر دے چې
    په~\LR{$\Gamma(\mClass{C})$} کښې معتبر وي۔
  \end{enumerate}
\end{cor}
د لومړي حکم لپاره هره اصلي نړۍ د فلټريشن خپل ټولګے لري او هر
ټولګے د يوې اصلي نړۍ ټولګے دے۔ د دويم حکم په سرچپه لوري کښې
هر اصلي مدل لږ تر لږه يو فلټريشن لري، لکه چې په ورپسې برخه
کښې به ئې جوړ کړو۔ په دواړو لورو کښې لومړے حکم صدق لېږدوي۔
\olpseganchor{OLP-0453-B022}{end}

\olpseganchor{OLP-0454-B004}{start}
\olfileid{nml}{fil}{exf}
\olpseganchor{OLP-0454-B004}{end}

\olpseganchor{OLP-0454-B005}{start}
\olsection{د فلټريشنونو بېلګې}
\olpseganchor{OLP-0454-B005}{end}

\olpseganchor{OLP-0454-B006}{start}
تر اوسه مو دا نۀ دي ښودلي چې فلټريشنونه شته هم که نه۔ خو په حقيقت
کښې د هر مدل~\LR{$\mModel{M}$} لپاره، د~\LR{$\Gamma$} له مخې د~\LR{$\mModel{M}$}
لږ تر لږه يو فلټريشن شته۔ دلته دوه حدونه تعريفوو: تر ټولو نرے
او تر ټولو غټ فلټريشن، چې يو شان هم کېدے شي۔ د يوه مدل فلټريشنونه
يوازې د لاسرسي په اړيکه کښې توپير لرلے شي (ځکه
\olref[fil]{defn:filtration} نېغ په نېغه ټاکي چې
\LR{$W^*$} او~\LR{$V^*$} بايد څه وي)۔ د هر فلټريشن د لاسرسي اړيکه د تر ټولو
نري فلټريشن اړيکه رانغاړي او د تر ټولو غټ فلټريشن په اړيکه کښې راځي۔
\olpseganchor{OLP-0454-B006}{end}

\olpseganchor{OLP-0454-B007}{start}
\begin{defn}
  که \LR{$\Gamma$} د فرعي فارمولونو لپاره تړلے وي، د مدل~\LR{$\mModel{M}$}
  \emph{تر ټولو نرے} فلټريشن~\LR{$\mModel{M^*}$} داسې تعريفېږي:
  \[
  R^*[u][v] \quad \text{\RL{هله او يوازې هله چې }} \quad \exists u'\in [u] \;
  \exists v' \in [v] : Ru'v'.
  \]
\end{defn}
\olpseganchor{OLP-0454-B007}{end}

\olpseganchor{OLP-0454-B008}{start}
\begin{prop}\ollabel{prop:finest}
  تر ټولو نرے فلټريشن~\LR{$\mModel{M^*}$} په رښتيا هم فلټريشن دے۔
\end{prop}
\olpseganchor{OLP-0454-B008}{end}

\olpseganchor{OLP-0454-B009}{start}
\begin{proof}
  بايد وګورو چې په دې ډول تعريف شوې \LR{$R^*$} د
  \LR{\olref[fil]{defn:filtration}(\olref[fil]{defn:filtration-R})} شرط پوره
  کوي۔ درې واړه شرطونه په وار سره ګورو۔
\olpseganchor{OLP-0454-B009}{end}

\olpseganchor{OLP-0454-B010}{start}
  که \LR{$Ruv$} وي، نو ځکه چې \LR{$u \in [u]$} او \LR{$v \in [v]$}، \LR{$R^*[u][v]$}
  هم وي؛ نو \olref[fil]{defn:filtration-R1} پوره کېږي۔
\olpseganchor{OLP-0454-B010}{end}
  
\olpseganchor{OLP-0454-B011}{start}
  د
        \olref[fil]{defn:filtration-R2} لپاره، فرض کړئ
        \LR{$\Box!A \in \Gamma$}، \LR{$R^*[u][v]$} او
        \LR{$\mSat{M}{\Box!A}[u]$} وي۔ د~\LR{$R^*$} د تعريف له مخې داسې
        \LR{$u' \equiv u$} او \LR{$v' \equiv v$} شته چې \LR{$Ru'v'$} وي۔ ځکه \LR{$u$}
        او \LR{$u'$} د~\LR{$\Gamma$} پر فارمولونو سره يوه خوله دي،
        \LR{$\mSat{M}{\Box!A}[u']$} هم وي؛ نو
        \LR{$\mSat{M}{!A}[v']$} وي۔ د فرعي فارمولونه لپاره د~\LR{$\Gamma$}
        د تړلتيا له مخې، \LR{$v$} او \LR{$v'$} د~\LR{$!A$} په اړه سره يوه خوله
        دي؛ نو لکه څنګه چې غوښتل مو، \LR{$\mSat{M}{!A}[v]$}۔
\olpseganchor{OLP-0454-B011}{end}
  
\olpseganchor{OLP-0454-B012}{start}
  د
        \olref[fil]{defn:filtration-R3} د تصديق لپاره، فرض کړئ
        \LR{$\Diamond!A \in \Gamma$}، \LR{$R^*[u][v]$} او
        \LR{$\mSat{M}{!A}[v]$} وي۔ د~\LR{$R^*$} د تعريف له مخې داسې
        \LR{$u' \equiv u$} او \LR{$v' \equiv v$} شته چې \LR{$Ru'v'$} وي۔ ځکه \LR{$v$}
        او \LR{$v'$} د~\LR{$\Gamma$} پر فارمولونو سره يوه خوله دي، او
        \LR{$\Gamma$} د فرعي فارمولونه لپاره تړلے دے،
        \LR{$\mSat{M}{!A}[v']$} هم وي؛ نو
        \LR{$\mSat{M}{\Diamond !A}[u']$}۔ ځکه \LR{$u$} او \LR{$u'$} هم
        د~\LR{$\Gamma$} پر فارمولونو سره يوه خوله دي،
        \LR{$\mSat{M}{\Diamond !A}[u]$}۔
\end{proof}
\olpseganchor{OLP-0454-B012}{end}



\olpseganchor{OLP-0454-B014}{start}
\begin{defn}
  که \LR{$\Gamma$} د فرعي فارمولونو لپاره تړلے وي، د مدل~\LR{$\mModel{M}$}
  \emph{تر ټولو غټ} فلټريشن~\LR{$\mModel{M^*}$} داسې تعريفېږي چې
  \LR{$R^*[u][v]$} هله او يوازې هله وي چې
  له لاندې
    شرطونو څخه \emph{دواړه} پوره شي:
  \begin{enumerate}
  \item \ollabel{defn:coarsest-Box}د هر فارمول دپاره،
    که \LR{$\Box!A \in \Gamma$} او \LR{$\mSat{M}{\Box!A}[u]$} وي، نو
    \LR{$\mSat{M}{!A}[v]$}؛
  \item \ollabel{defn:coarsest-Diamond}د هر فارمول دپاره،
    که \LR{$\Diamond!A
    \in \Gamma$} او \LR{$\mSat{M}{!A}[v]$} وي، نو
    \LR{$\mSat{M}{\Diamond!A}[u]$}۔
  \end{enumerate}
\end{defn}
\olpseganchor{OLP-0454-B014}{end}

\olpseganchor{OLP-0454-B015}{start}
\begin{prop}
  تر ټولو غټ فلټريشن~\LR{$\mModel{M^*}$} په رښتيا هم فلټريشن دے۔
\end{prop}
\olpseganchor{OLP-0454-B015}{end}

\olpseganchor{OLP-0454-B016}{start}
\begin{proof}
  د~\LR{$R^*$} د تعريف له مخې، يوازې دا شرط پاتې دے چې له \LR{$Ruv$} څخه
  \LR{$R^*[u][v]$} راووځي۔ نو \LR{$Ruv$} فرض کړئ۔
  فرض کړئ \LR{$\Box!A \in \Gamma$} او
    \LR{$\mSat{M}{\Box!A}[u]$} وي؛ نو ښکاره ده چې
    \LR{$\mSat{M}{!A}[v]$}، او
      \olref{defn:coarsest-Box} پوره کېږي۔ فرض
    کړئ \LR{$\Diamond!A \in \Gamma$} او \LR{$\mSat{M}{!A}[v]$} وي۔ نو
    ځکه چې~\LR{$Ruv$}،
    \LR{$\mSat{M}{\Diamond!A}[u]$}، او
      \olref{defn:coarsest-Diamond} پوره کېږي۔
\end{proof}
\olpseganchor{OLP-0454-B016}{end}

\olpseganchor{OLP-0454-B017}{start}
\begin{ex}
  فرض کړئ \LR{$W = \PosInt$}، \LR{$Rnm$} هله او يوازې هله چې \LR{$m = n + 1$}،
  او \LR{$V(p) = \Setabs{2n}{n \in
    \PosInt}$}۔ مدل~\LR{$\mModel{M} = \tuple{W, R, V}$} په
  \olref{fig:ex-filtration} کښې انځور شوے دے۔ نړۍې \LR{$1$}، \LR{$2$} او
  داسې نورې دي؛ هره نړۍ يوازې يوې بلې نړۍ، يعنې خپل راتلونکي ته
  لاسرسے لري، او \LR{$p$} هله او يوازې هله رښتينے دے چې شمېره جفته وي۔
\olpseganchor{OLP-0454-B017}{end}

\olpseganchor{OLP-0454-B018}{start}
  \begin{figure}
    \centering
    \begin{tikzpicture}[modal]
      \node[world] (1) [label=below:\mFalse{p}] {\LR{$1$}};
      \node[world] (2) [label=below:\mTrue{p}, right=of 1] {\LR{$2$}};
      \node[world] (3) [label=below:\mFalse{p}, right=of 2] {\LR{$3$}};
      \node[world] (4) [label=below:\mTrue{p}, right=of 3] {\LR{$4$}};
      \node[phantom] (5) [right=of 4] {};
      \draw[->] (1) to (2);
      \draw[->] (2) to (3);
      \draw[->] (3) to (4);
      \draw[dotted] (4) to (5);
    \end{tikzpicture}
\olpseganchor{OLP-0454-B018}{end}

\olpseganchor{OLP-0454-B019}{start}
    \begin{tikzpicture}[modal]
      \node[world] (1) [label=below:\mFalse{p}] {\LR{$[1]$}};
      \node[world] (2) [label=below:\mTrue{p}, right=of 1] {\LR{$[2]$}};
      \draw[->,bend left] (1) to (2);
      \draw[->,bend left] (2) to (1);
\olpseganchor{OLP-0454-B019}{end}
      
\olpseganchor{OLP-0454-B020}{start}
      \node[world] (3) [label=below:\mFalse{p},right=of 2] {\LR{$[1]$}};
      \node[world] (4) [label=below:\mTrue{p}, right=of 3] {\LR{$[2]$}};
      \draw[->,bend left] (3) to (4);
      \draw[->,bend left] (4) to (3);
      \draw[->,reflexive above] (4) to (4);
    \end{tikzpicture}
    \caption{يو نامتناهي مدل او د هغه فلټريشنونه۔}
    \ollabel{fig:ex-filtration}
  \end{figure}
\olpseganchor{OLP-0454-B020}{end}
  
\olpseganchor{OLP-0454-B021}{start}
  اوس \LR{$\Gamma$} د~\LR{$\Box p \lif p$} د فرعي فارمولونه سټ، يعنې
  \LR{$\{p, \Box p, \Box p \lif p\}$}، وټاکئ۔ \LR{$p$} يوازې په جفتو
  شمېرو کښې رښتينے دے، او \LR{$\Box p$} يوازې په طاقو شمېرو کښې
  رښتينے دے؛ نو \LR{$\Box p \lif p$} هم يوازې په جفتو شمېرو کښې
  رښتينے دے۔ په بل عبارت، هره طاقه شمېره \LR{$\Box p$} رښتينے کوي،
  خو \LR{$p$} او \LR{$\Box p \lif p$} ناسم کوي؛ هره جفته شمېره \LR{$p$} او
  \LR{$\Box p \lif p$} رښتيني کوي، خو \LR{$\Box p$} ناسم کوي۔ نو
  \LR{$W^* = \{ [1], [2] \}$}، چېرته چې \LR{$[1] =
  \{1, 3, 5, \dots\}$} او \LR{$[2] = \{2, 4, 6, \dots\}$} دي۔ ځکه
  \LR{$2 \in V(p)$}، \LR{$[2] \in V^*(p)$}؛ او ځکه \LR{$1 \notin V(p)$}،
  \LR{$[1] \notin V^*(p)$}۔ نو \LR{$V^*(p) = \{[2]\}$}۔
\olpseganchor{OLP-0454-B021}{end}

\olpseganchor{OLP-0454-B022}{start}
  پر~\LR{$W^*$} ولاړ هر فلټريشن بايد داسې د لاسرسي اړيکه ولري چې
  \LR{$\tuple{[1], [2]}, \tuple{[2],[1]}$} پکښې وي: ځکه \LR{$R12$}،
  د \LR{\olref[fil]{defn:filtration}(\olref[fil]{defn:filtration-R1})} له
  مخې \LR{$R^*[1][2]$} اړين دے؛ او ځکه \LR{$R23$}، بايد
  \LR{$R^*[2][3]$} وي، او \LR{$[3]=[1]$}۔ خو
  \LR{$\tuple{[1],[1]}$} پکښې راتلے نۀ شي: که راغلے، نو
  \LR{$R^*[1][1]$} به وي، حال دا چې \LR{$\mSat{M}{\Box p}[1]$}
  خو \LR{$\mSat/{M}{p}[1]$}؛ دا له
  \olref[fil]{defn:filtration-R2} سره په ټکر کښې دے۔ هېڅ شرط
  \LR{$R^*[2][2]$} نۀ لازموي او نۀ ئې منع کوي۔ نو د~\LR{$\mModel{M}$}
  دوه ممکن فلټريشنونه دي چې له لاندې دوو د لاسرسي اړيکو سره سمون
  لري:
  \[
  \{\tuple{[1],[2]}, \tuple{[2],[1]}\} \text{\RL{ او }}
  \{\tuple{[1],[2]}, \tuple{[2],[1]}, \tuple{[2],[2]}\}.
  \]
  په دواړو حالاتو کښې، په~\LR{$[1]$} کښې \LR{$p$} او
  \LR{$\Box p \lif p$} ناسم، او \LR{$\Box p$} رښتينے دي؛ په~\LR{$[2]$}
  کښې \LR{$p$} او \LR{$\Box p \lif p$} رښتيني، او \LR{$\Box p$} ناسم دي۔
\end{ex}
\olpseganchor{OLP-0454-B022}{end}


\olpseganchor{OLP-0454-B023}{start}
\begin{prob}
  لاندې مدل~\LR{$\mModel{M} = \tuple{W, R, V}$} وګورئ، چې \LR{$W
  = \Setabs{0\sigma}{\sigma \in \Bin^*}$}، يعنې د~\LR{$0$} او~\LR{$1$} د
  هغو لړيو سټ دے چې په~\LR{$0$} پيلېږي؛ \LR{$R\sigma\sigma'$} هله او يوازې
  هله چې \LR{$\sigma' =
  \sigma 0$} يا \LR{$\sigma' = \sigma 1$} وي؛ او \LR{$V(p) = W \cap
  \Setabs{\sigma
    0}{\sigma \in \Bin^*}$} او \LR{$V(q) = W \cap \Setabs{\sigma 1}{\sigma \in
    \Bin^* \setminus \{1\}}$} دي۔ په دې حساب د مدل په نړۍو کښې
  لومړے متغير هله رښتينے دے چې لړۍ پر صفر ختمېږي، او دويم هله
  چې پر يو ختمېږي۔
  دا ئې انځور دے:
  \begin{center}
    \begin{tikzpicture}[modal,
        node distance=2cm
      ]
\olpseganchor{OLP-0454-B023}{end}

\olpseganchor{OLP-0454-B024}{start}
      \node[world] (0) [label={below:\mTrue{p}\\\mFalse{q}},font=\tiny] {\LR{$0$}};
      \node[world] (00) [label={below:\mTrue{p}\\\mFalse{q}},
        above right=of 0,font=\tiny] {\LR{$00$}};
      \node[world] (000) [label={below:\mTrue{p}\\\mFalse{q}},
        right=of 00, yshift=.8cm, font=\tiny] {\LR{$000$}};
      \node[phantom] (0000) [right=of 000, yshift=.5cm] {};
      \node[phantom] (0001) [right=of 000, yshift=-.5cm] {};
      \node[world] (001) [label={below:\mFalse{p}\\\mTrue{q}},
        right=of 00, yshift=-.8cm, font=\tiny] {\LR{$001$}};
      \node[phantom](0010) [right=of 001, yshift=.5cm] {};
      \node[phantom](0011) [right=of 001, yshift=-.5cm] {};
      \node[world] (01) [label={below:\mFalse{p}\\\mTrue{q}},
        below right=of 0,font=\tiny] {\LR{$01$}};
      \node[world] (010) [label={below:\mTrue{p}\\\mFalse{q}},
        right=of 01, yshift=.8cm,font=\tiny] {\LR{$010$}};
      \node[phantom] (0100) [right=of 010, yshift=.5cm] {};
      \node[phantom] (0101) [right=of 010, yshift=-.5cm] {};
      \node[world] (011) [label={below:\mFalse{p}\\\mTrue{q}},
        right=of 01, yshift=-.8cm,font=\tiny] {\LR{$011$}};
      \node[phantom] (0110) [right=of 011, yshift=.5cm] {};
      \node[phantom] (0111) [right=of 011, yshift=-.5cm] {};
\olpseganchor{OLP-0454-B024}{end}

\olpseganchor{OLP-0454-B025}{start}
      \draw[->] (0) to (00);
      \draw[->] (0) to (01);
      \draw[->] (00) to (000);
      \draw[dotted] (000) to (0000);
      \draw[dotted] (000) to (0001);
      \draw[->] (00) to (001);
      \draw[dotted] (001) to (0010);
      \draw[dotted] (001) to (0011);
      \draw[->] (01) to (010);
      \draw[dotted] (010) to (0100);
      \draw[dotted] (010) to (0101);
      \draw[->] (01) to (011);
      \draw[dotted] (011) to (0110);
      \draw[dotted] (011) to (0111);
    \end{tikzpicture}
  \end{center}
  د هرې نړۍ~\LR{$w$} لپاره
  \LR{$\mSat/{M}{\Box(p \lor q) \lif (\Box p \lor \Box q)}[w]$} وي۔
\olpseganchor{OLP-0454-B025}{end}

\olpseganchor{OLP-0454-B026}{start}
  \LR{$\Gamma$} د~\LR{$\Box(p \lor q) \lif
  (\Box p \lor \Box q)$} د فرعي فارمولونه سټ وټاکئ۔ \LR{$W^*$} او
  \LR{$V^*$} څه دي؟ د~\LR{$\mModel{M}$} د تر ټولو نري فلټريشن د لاسرسي
  اړيکه څه ده؟ د تر ټولو غټ فلټريشن اړيکه څه ده؟
\end{prob}
\olpseganchor{OLP-0454-B026}{end}

\olpseganchor{OLP-0455-B004}{start}
\olfileid{nml}{fil}{fin}
\olpseganchor{OLP-0455-B004}{end}

\olpseganchor{OLP-0455-B005}{start}
\olsection{فلټريشنونه متناهي دي}
\olpseganchor{OLP-0455-B005}{end}

\olpseganchor{OLP-0455-B006}{start}
موږ د هر هغه سټ~\LR{$\Gamma$} لپاره فلټريشنونه تعريف کړي چې د فرعي
فارمولونه لپاره تړلے وي۔ پخپله تعريف دا نۀ تضمينوي چې
فلټريشنونه متناهي دي۔ په حقيقت کښې، که \LR{$\Gamma$} نامتناهي وي (د
بېلګې په توګه، د ټولو فارمولونه سټ وي)، فلټريشن هم نامتناهي
کېدے شي۔ خو که \LR{$\Gamma$} متناهي وي (د بېلګې په توګه، د يوه ټاکلي
فارمول~\LR{$!A$} د فرعي فارمولونه سټ وي)، نو د~\LR{$\Gamma$} له
مخې هر فلټريشن هم متناهي وي۔
\olpseganchor{OLP-0455-B006}{end}

\olpseganchor{OLP-0455-B007}{start}
\begin{prop}\ollabel{prop:filt-are-finite}
  که \LR{$\Gamma$} متناهي وي، نو د مدل~\LR{$\mModel{M}$} هر
  فلټريشن~\LR{$\mModel{M^*}$} د~\LR{$\Gamma$} له مخې هم متناهي وي۔
\end{prop}
\olpseganchor{OLP-0455-B007}{end}

\olpseganchor{OLP-0455-B008}{start}
\begin{proof}
  د~\LR{$W^*$} کچه د هم‌ارزښتۍ اړيکې~\LR{$\equiv$} له مخې د بېلابېلو
  ټولګيو~\LR{$[w]$} شمېر دے۔ په داسې يوه ټولګي کښې هرې دوې نړۍې
  \LR{$u$} او \LR{$v$}، يعنې هرې هغه \LR{$u$} او \LR{$v$} چې~\LR{$u \equiv v$} وي،
  د~\LR{$\Gamma$} د ټولو فارمولونه~\LR{$!A$} په اړه سره يوه خوله دي:
  که \LR{$!A \in \Gamma$} وي، نو يا \LR{$!A$} په دواړو \LR{$u$} او \LR{$v$} کښې
  رښتينے وي، يا په هېڅ يوه کښې نۀ وي۔ نو هر ټولگے~\LR{$[w]$} د~\LR{$\Gamma$}
  له يوه فرعي سټ سره سمون لري، يعنې د ټولو هغو \LR{$!A \in
  \Gamma$} سټ سره چې \LR{$!A$} د~\LR{$[w]$} په نړۍو کښې رښتينے وي۔
  دوه بېلابېل ټولګي \LR{$[u]$} او \LR{$[v]$} د~\LR{$\Gamma$} له يوه او هماغه
  فرعي سټ سره سمون نۀ شي لرلے۔ ځکه که په~\LR{$u$} کښې د رښتينو
  فارمولونه سټ او په~\LR{$v$} کښې د رښتينو فارمولونه سټ يو
  وي، نو \LR{$u$} او \LR{$v$} د~\LR{$\Gamma$} په ټولو فارمولونو کښې سره يوه
  خوله دي، يعنې \LR{$u \equiv v$}۔ خو بيا \LR{$[u] = [v]$} وي۔ نو
  له~\LR{$W^*$} څخه \LR{$\Pow{\Gamma}$} ته يو پر يو تابع شته، او
  ځکه \LR{$\card{W^*} \le \card{\Pow{\Gamma}}$}۔ له دې امله که
  \LR{$\Gamma$} د \LR{$n$} فارمولونو لرونکے وي، د~\LR{$W^*$} شمېر له~\LR{$2^n$}
  څخه زيات نۀ دے۔
\end{proof}
\olpseganchor{OLP-0455-B008}{end}

\olpseganchor{OLP-0456-B004}{start}
\olfileid{nml}{fil}{fmp}
\olpseganchor{OLP-0456-B004}{end}

\olpseganchor{OLP-0456-B005}{start}
\olsection{\texorpdfstring{\Log{K}}{K} او \texorpdfstring{\Log{S5}}{S5} د متناهي مدل خاصيت لري}
\olpseganchor{OLP-0456-B005}{end}

\olpseganchor{OLP-0456-B006}{start}
\begin{defn}
  د مودالي منطق يو نظام~\LR{$\Sigma$} د \emph{متناهي
    مدل خاصيت} لري که هرکله چې کوم فارمول~\LR{$!A$} د~\LR{$\Sigma$}
  په يوه مدل کښې په يوه نړۍ رښتينے وي، نو \LR{$!A$} د~\LR{$\Sigma$} په
  يوه \emph{متناهي} مدل کښې هم په يوه نړۍ رښتينے وي۔
\end{defn}
\olpseganchor{OLP-0456-B006}{end}

\olpseganchor{OLP-0456-B007}{start}
\begin{prop}\ollabel{prop:K-fmp}
  \Log{K} د متناهي مدل خاصيت لري۔
\end{prop}
\olpseganchor{OLP-0456-B007}{end}

\olpseganchor{OLP-0456-B008}{start}
\begin{proof}
  \Log{K} د معتبرو فارمولونه سټ دے، يعنې هر مدل د~\Log{K}
  مدل دے۔ د \olref[fil]{thm:filtrations} له مخې، که
  \LR{$\mSat{M}{!A}[w]$} وي، نو د~\LR{$!A$} د فرعي فارمولونه د
  سټ~\LR{$\Gamma$} له مخې د~\LR{$\mModel{M}$} د هر فلټريشن لپاره
  \LR{$\mSat{M^*}{!A}[{[w]}]$} وي۔ تر ټولو نرے فلټريشن تل شته۔
  هر فارمول يوازې متناهي ډېر فرعي
  فارمولونه لري؛ نو \LR{$\Gamma$} متناهي دے۔ د
  \olref[fin]{prop:filt-are-finite} له مخې
  \LR{$\card{W^*} \le 2^n$}، چېرته چې \LR{$n$} په~\LR{$\Gamma$} کښې د
  فارمولونه شمېر دے۔ او ځکه \Log{K} پر مدلونو هېڅ قيد نۀ ږدي،
  \LR{$\mModel{M^*}$} د~\Log{K} مدل دے۔
\end{proof}
\olpseganchor{OLP-0456-B008}{end}

\olpseganchor{OLP-0456-B009}{start}
د فلټريشنونو له لارې د دې ښودلو لپاره چې يو منطق~\Log{L} د متناهي
مدل خاصيت لري، اړينه ده چې د~\Log{L} د يوه مدل فلټريشن پخپله هم
د~\Log{L} مدل وي۔ ډېر وخت دا ښۀ ډېر کار غواړي، او هر فلټريشن
د~\Log{L} مدل نۀ جوړوي۔ خو د عمومي مدلونو لپاره دا خاصيت پاتې
کېږي۔
\olpseganchor{OLP-0456-B009}{end}

\olpseganchor{OLP-0456-B010}{start}
\begin{prop}\ollabel{prop:univ-fin}
  فرض کړئ \LR{$\mClass{U}$} د عمومي مدلونو ټولگے وي (وګورئ
  \olref[frd][es5]{prop:S5=univ}) او \LR{$\mClass{U}_\mathrm{Fin}$} د
  ټولو متناهي عمومي مدلونو ټولگے وي۔ بيا هر فارمول~\LR{$!A$}
  په~\LR{$\mClass{U}$} کښې هله او يوازې هله معتبر دے چې
  په~\LR{$\mClass{U}_\mathrm{Fin}$} کښې معتبر وي۔
\end{prop}
\olpseganchor{OLP-0456-B010}{end}

\olpseganchor{OLP-0456-B011}{start}
\begin{proof}
  متناهي عمومي مدلونه عمومي مدلونه دي؛ نو له کيڼ نه ښي لوري ته
  خبره څرګنده ده۔ د سرچپه لوري لپاره فرض کړئ چې~\LR{$!A$} په يوه
  عمومي مدل~\LR{$\mModel{M}$} کښې په کومه نړۍ~\LR{$w$} ناسم وي۔
  \LR{$\Gamma$} د~\LR{$!A$} د ټولو فرعي فارمولونو سټ واخلئ، چې فارمول
  پخپله هم پکښې دے۔ دا سټ~\LR{$\Gamma$} د فرعي فارمولونو لپاره تړلے او متناهي دے۔ د~\LR{$\mModel{M}$} يو
  فلټريشن~\LR{$\mModel{M^*}$} واخلئ؛ بيا
  \LR{$\mModel{M^*}$} د \olref[fin]{prop:filt-are-finite} له مخې متناهي
  دے، او د \olref[fil]{thm:filtrations} له مخې \LR{$!A$}
  په~\LR{$\mModel{M^*}$} کښې په~\LR{$[w]$} ناسم دے۔ اوس يوازې دا پاتې دي
  چې \LR{$\mModel{M^*}$} هم عمومي دے: هر \LR{$u$} او \LR{$v$} چې واخلئ، د فرض
  له مخې \LR{$Ruv$} وي او د
  \LR{\olref[fil]{defn:filtration}(\olref[fil]{defn:filtration-R})} له مخې
  \LR{$R^*[u][v]$} هم وي۔
\end{proof}
\olpseganchor{OLP-0456-B011}{end}

\olpseganchor{OLP-0456-B012}{start}
\begin{cor}\ollabel{cor:S5fmp}
  \Log{S5} د متناهي مدل خاصيت لري۔
\end{cor}
\olpseganchor{OLP-0456-B012}{end}

\olpseganchor{OLP-0456-B013}{start}
\begin{proof}
  د \olref[frd][es5]{prop:S5=univ} له مخې، که \LR{$!A$} په کوم
  انعکاسي او اقليدسي مدل کښې په يوه نړۍ رښتينے وي، نو په يوه
  عمومي مدل کښې هم په يوه نړۍ رښتينے وي۔ د دغه فارمول نفي
  په عمومي مدلونو کښې اعتبار نۀ لري۔ نو د
  \olref{prop:univ-fin} له مخې دا نفي په متناهي عمومي مدلونو کښې
  هم اعتبار نۀ لري، يعنې اصلي فارمول په يوه متناهي عمومي مدل کښې هم
  په يوه نړۍ رښتينے وي (يعنې د~\LR{$!A$} د فرعي فارمولونه د سټ
  له مخې د هغه مدل په فلټريشن کښې)۔ هر عمومي مدل انعکاسي او
  اقليدسي هم دے؛ نو \LR{$!A$} په يوه متناهي انعکاسي اقليدسي مدل
  کښې په يوه نړۍ رښتينے دے۔
\end{proof}
\olpseganchor{OLP-0456-B013}{end}

\olpseganchor{OLP-0456-B014}{start}
\begin{prob}
  وښيئ چې د هر مسلسل يا انعکاسي مدل هر فلټريشن هم په ترتيب سره
  مسلسل يا انعکاسي وي۔
\end{prob}
\olpseganchor{OLP-0456-B014}{end}

\olpseganchor{OLP-0456-B015}{start}
\begin{prob}
  د يوه متناظر (متعدي، اقليدسي) مدل داسې يو فلټريشن ومومئ چې
  متناظر (متعدي، اقليدسي) نۀ وي۔
\end{prob}
\olpseganchor{OLP-0456-B015}{end}

\olpseganchor{OLP-0457-B004}{start}
\olfileid{nml}{fil}{dec}
\olpseganchor{OLP-0457-B004}{end}

\olpseganchor{OLP-0457-B005}{start}
\olsection{\texorpdfstring{\Log{S5}}{S5} د پرېکړې وړ دے}
د متناهي مدل خاصيت موږ ته د دې ښودلو اسانه لار راکوي چې د اغېزمنو سکيماوو
له مخې ورکړل شوي ځينې مودالي منطقي نظامونه \emph{د پرېکړې وړ} دي (يعنې
داسې محاسبه‌کېدونکې کړنلاره شته چې وټاکي آيا کوم فارمول په
نظام کښې د اشتقاق وړ دے که نه)۔
\olpseganchor{OLP-0457-B005}{end}

\olpseganchor{OLP-0457-B006}{start}
\begin{thm}
  \Log{S5} د پرېکړې وړ دے۔
\end{thm}
\olpseganchor{OLP-0457-B006}{end}

\olpseganchor{OLP-0457-B007}{start}
\begin{proof}
  \LR{$!A$} راکړل شوے وګڼئ، او فرض کړئ چې په~\LR{$!A$} کښې راغلي
  بياني متغيرونه له \LR{$p_1$}, \dots, \LR{$p_k$} څخه وي۔ ځکه
  د هر \LR{$n$} لپاره د~\LR{$n$} نړۍو او \LR{$p_1$}, \dots, \LR{$p_k$} ته د
  ارزښت ټاکلو لرونکي يوازې متناهي ډېر مدلونه شته، موږ په
  \emph{موازي ډول} دوه لړۍ شمېرلے شو: له يوې خوا په يوې منظمه
  طريقه د ثبوتونو په جوړولو د~\Log{S5} ټولې قضيې، او له بلې خوا
  هغه ټول عمومي مدلونه چې \LR{$1$}, \LR{$2$}, \dots نړۍې لري، او بيا
  وګورو چې آيا \LR{$!A$} په کوم داسې مدل کښې په يوه نړۍ ناسم دے۔
  عمومي لاسرسے د نړۍو ټولې جوړې رانغاړي؛ په دې دويمه پلټنه
  کښې د هرې کچې د نړۍو نومونه ثابت نيسو او د اړوندو متغيرونو
  ټولې ارزښت‌ټاکنې ازمېيو۔
  له دغو دوو موازي بهيرونو څخه يو به بالاخره ځواب ورکړي، ځکه د
  \olref[com][fra]{thm:generaldet} او
  \olref[fmp]{cor:S5fmp} له مخې، د عمومي او انعکاسي اقليدسي
  چوکاټونو د اعتبار له برابرۍ سره يو ځاے، يا \LR{$!A$} د اشتقاق وړ دے، يا
  په يوه متناهي عمومي مدل کښې ناسم دے۔
\end{proof}
\olpseganchor{OLP-0457-B007}{end}

\olpseganchor{OLP-0457-B008}{start}
پورتنے ثبوت د~\Log{S5} لپاره ځکه کار کوي چې د عمومي مدلونو
فلټريشنونه په خپله عمومي وي۔ د انعکاسيت او مسلسل‌والي لپاره هم
همداسې ده، خو د نورو خاصيتونو لپاره زيات کار پکار دے۔
\olpseganchor{OLP-0457-B008}{end}

\olpseganchor{OLP-0458-B004}{start}
\olfileid{nml}{fil}{acc}
\olpseganchor{OLP-0458-B004}{end}

\olpseganchor{OLP-0458-B005}{start}
\olsection{فلټريشنونه او د لاسرسي خاصيتونه}
\olpseganchor{OLP-0458-B005}{end}

\olpseganchor{OLP-0458-B006}{start}
لکه چې يادونه وشوه، د عمومي، مسلسل او انعکاسي مدلونو فلټريشنونه
هم تل په ترتيب سره عمومي، مسلسل يا انعکاسي وي۔ خو د متناظر يا
متعدي مدل هر فلټريشن په ترتيب سره متناظر يا متعدي نۀ وي۔ سره له
دې، په ځينو حالاتو کښې فلټريشنونه داسې تعريفېدے شي چې دا خاصيتونه
وساتل شي۔ د دې لپاره د تر ټولو غټ فلټريشن د تعريف په شان ځو، خو
د~\LR{$R^*$} په تعريف کښې نور شرطونه هم ورزياتوو۔ فرض کړئ \LR{$\Gamma$}
د فرعي فارمولونه لپاره تړلے دے۔ په مدل~\LR{$\mModel{M} =\tuple{W, R, V}$}
کښې د نړۍو \LR{$u$} او \LR{$v$} ترمنځ د
\olref{tab:Cn-filtrations} اړيکې~\LR{$C_i(u,v)$} وګورئ۔ د دغو شرطونو
له بېلابېلو ګډونو څخه \LR{$R^*[u][v]$} تعريفولے شو۔ د بېلګې په توګه،
که \LR{$R^*[u][v]$} هله او يوازې هله وټاکو چې \LR{$C_1(u,v)$} پوره وي،
نو عين تر ټولو غټ فلټريشن راځي۔ که \LR{$R^*[u][v]$} هله او يوازې هله
وټاکو چې \LR{$C_1(u,v)$} او \LR{$C_2(u, v)$} دواړه پوره وي، نو
لاندې قضيه ښيي چې په متناظر مدل کښې دا اړيکه فلټريشن جوړوي۔
هغه جوړې چې \LR{$C_1(u,v)$} او \LR{$C_2(u,v)$} دواړه پوره کوي د هغو
جوړو په سټ کښې راځي چې \LR{$C_1(u,v)$} پوره کوي؛ دواړه سټونه
برابر هم کېدے شي۔
\olpseganchor{OLP-0458-B006}{end}

\olpseganchor{OLP-0458-B007}{start}
\begin{table}[ht]
  \centering
  \begin{tabular}{|ll|}
    \hline
    \multirow{2}{*}{\LR{$C_1(u,v)$}:}
    &
      که \LR{$\Box!A \in \Gamma$} او \LR{$\mSat{M}{\Box!A}[u]$} وي،
      نو \LR{$\mSat{M}{!A}[v]$}؛ او\\
    &
      که \LR{$\Diamond!A \in \Gamma$} او \LR{$\mSat{M}{!A}[v]$} وي،
      نو \LR{$\mSat{M}{\Diamond!A}[u]$}؛ \\
    \hline
    \multirow{2}{*}{\LR{$C_2(u,v)$}:} 
    &
      که \LR{$\Box!A \in \Gamma$} او \LR{$\mSat{M}{\Box!A}[v]$} وي،
      نو \LR{$\mSat{M}{!A}[u]$}؛ او\\
    &
      که \LR{$\Diamond!A \in \Gamma$} او \LR{$\mSat{M}{!A}[u]$} وي،
      نو \LR{$\mSat{M}{\Diamond!A}[v]$}؛ \\
    \hline
    \multirow{2}{*}{\LR{$C_3(u,v)$}:}
    &
      که \LR{$\Box!A \in \Gamma$} او \LR{$\mSat{M}{\Box!A}[u]$} وي،
      نو \LR{$\mSat{M}{\Box!A}[v]$}؛ او\\
    &
      که \LR{$\Diamond!A \in \Gamma$} او \LR{$\mSat{M}{\Diamond!A}[v]$} وي،
      نو \LR{$\mSat{M}{\Diamond!A}[u]$}؛ \\
    \hline
    \multirow{2}{*}{\LR{$C_4(u,v)$}:}
    &
      که \LR{$\Box!A \in \Gamma$} او \LR{$\mSat{M}{\Box!A}[v]$} وي،
      نو \LR{$\mSat{M}{\Box!A}[u]$}؛ او\\
    &
      که \LR{$\Diamond!A \in \Gamma$} او \LR{$\mSat{M}{\Diamond!A}[u]$} وي،
      نو \LR{$\mSat{M}{\Diamond!A}[v]$}؛ \\
    \hline
    \end{tabular}
  \caption{د فلټريشنونو د تعريف لپاره پر ممکنو نړۍو شرطونه۔ د هرې ليکې شرط پر ټولو اړوندو فارمولونو پلي کېږي۔}
  \ollabel{tab:Cn-filtrations}
\end{table}
\olpseganchor{OLP-0458-B007}{end}

\olpseganchor{OLP-0458-B008}{start}
\begin{thm}\ollabel{thm:more-filtrations}
  فرض کړئ \LR{$\mModel{M} =\tuple{W,R,V}$} يو مدل وي او \LR{$\Gamma$} د
  فرعي فارمولونه لپاره تړلے وي۔ \LR{$W^*$} او \LR{$V^*$} د
  \olref[fil]{defn:filtration} په شان تعريف کړئ۔ بيا:
  \begin{enumerate}
  \item فرض کړئ \LR{$R^*[u][v]$} هله او يوازې هله چې \LR{$C_1(u, v) \land
    C_2(u,v)$}۔ بيا \LR{$R^*$} متناظره ده، او \LR{$\mModel{M^*} =
    \tuple{W^*,R^*,V^*}$} فلټريشن دے که \LR{$\mModel{M}$} متناظر وي۔
  \item فرض کړئ \LR{$R^*[u][v]$} هله او يوازې هله چې \LR{$C_1(u, v)
    \land C_3(u,v)$}۔ بيا \LR{$R^*$} متعديه ده، او
    \LR{$\mModel{M^*}=\tuple{W^*,R^*,V^*}$} فلټريشن دے که
    \LR{$\mModel{M}$} متعدي وي۔
  \item فرض کړئ \LR{$R^*[u][v]$} هله او يوازې هله چې \LR{$C_1(u, v) \land C_2(u,v)
    \land C_3(u,v) \land C_4(u,v)$}۔ بيا \LR{$R^*$} متناظره او متعديه
    ده، او \LR{$\mModel{M^*}=\tuple{W^*,R^*,V^*}$} فلټريشن دے که
    \LR{$\mModel{M}$} متناظر او متعدي وي۔
  \item فرض کړئ \LR{$R^*$} داسې تعريف شوې وي چې \LR{$R^*[u][v]$} هله او
    يوازې هله چې \LR{$C_1(u,
    v) \land C_3(u,v) \land C_4(u,v)$}۔ بيا \LR{$R^*$} متعديه او
    اقليدسي ده، او \LR{$\mModel{M^*}=\tuple{W^*,R^*,V^*}$} فلټريشن
    دے که \LR{$\mModel{M}$} متعدي او اقليدسي وي۔
  \end{enumerate}
\end{thm}
\olpseganchor{OLP-0458-B008}{end}

\olpseganchor{OLP-0458-B009}{start}
\begin{proof}
  \begin{enumerate}
    \item \LR{$R^*$} سمدستي متناظره ښکاري، ځکه \LR{$C_1(u,v)
      \Leftrightarrow C_2(v,u)$} او \LR{$C_2(u,v) \Leftrightarrow
      C_1(v,u)$}۔ نو دا ښودل پاتې دي چې که \LR{$\mModel{M}$} متناظر
      وي، \LR{$\mModel{M^*}$} د~\LR{$\Gamma$} له مخې فلټريشن دے۔
      شرط~\LR{$C_1(u,v)$} دا تضمينوي چې
      %
        %
          د \LR{\olref[fil]{defn:filtration}(\olref[fil]{defn:filtration-R2})} او
          \olref[fil]{defn:filtration-R3} پوره کېږي۔
      نو يوازې د
      \LR{\olref[fil]{defn:filtration}(\olref[fil]{defn:filtration-R1})},
      تصديق پکار دے، يعنې له \LR{$Ruv$} څخه \LR{$R^*[u][v]$} راووځي۔
\olpseganchor{OLP-0458-B009}{end}

\olpseganchor{OLP-0458-B010}{start}
      نو \LR{$Ruv$} فرض کړئ۔ د \LR{$R^*[u][v]$} د ښودلو لپاره بايد
      \LR{$C_1(u,v)$} او \LR{$C_2(u,v)$} ثابت کړو۔ د~\LR{$C_1$} لپاره:
      که \LR{$\Box!A
        \in\Gamma$} او \LR{$\mSat{M}{\Box!A}[u]$} وي، نو
        \LR{$\mSat{M}{!A}[v]$} هم وي (ځکه \LR{$Ruv$})۔
        همداسې، که
        \LR{$\Diamond!A \in \Gamma$} او \LR{$\mSat{M}{!A}[v]$} وي، نو
        د~\LR{$Ruv$} له امله \LR{$\mSat{M}{\Diamond!A}[u]$}۔
      د~\LR{$C_2$} لپاره: که \LR{$\Box!A \in \Gamma$} او
        \LR{$\mSat{M}{\Box!A}[v]$} وي، نو د متناظرۍ له مخې له \LR{$Ruv$}
        څخه \LR{$Rvu$} راځي؛ نو \LR{$\mSat{M}{!A}[u]$}۔
        همداسې، که
        \LR{$\Diamond!A \in\Gamma$} او \LR{$\mSat{M}{!A}[u]$} وي، نو
        \LR{$\mSat{M}{\Diamond!A}[v]$} (ځکه د متناظرۍ له مخې \LR{$Rvu$})۔
    \item تمرين۔
    \item تمرين۔
    \item تمرين۔
  \end{enumerate}
\end{proof}
\olpseganchor{OLP-0458-B010}{end}

\olpseganchor{OLP-0458-B011}{start}
\begin{prob}
  د \olref[nml][fil][acc]{thm:more-filtrations} ثبوت بشپړ کړئ۔
\end{prob}
\olpseganchor{OLP-0458-B011}{end}

\olpseganchor{OLP-0459-B004}{start}
\olfileid{nml}{fil}{euc}
\olpseganchor{OLP-0459-B004}{end}

\olpseganchor{OLP-0459-B005}{start}
\olsection{د اقليدسي مدلونو فلټريشنونه}
\olpseganchor{OLP-0459-B005}{end}

\olpseganchor{OLP-0459-B006}{start}
د \olref[acc]{sec} طريقه د هغو مدلونو لپاره کار نۀ ورکوي چې اقليدسي
وي، يا مسلسل او اقليدسي دواړه وي۔ د \olref{fig:ser-eucl} پاسنے مدل
وګورئ؛ هغه اقليدسي او مسلسل دواړه دے۔ فرض کړئ
\LR{$\Gamma = \{p, \Box p \}$}۔ د~\LR{$\Gamma$} له مخې فلټريشن کښې
\LR{$[w_1] = [w_3]$} کېږي، ځکه \LR{$w_1$} او \LR{$w_3$} يوازې هغه نړۍې دي
چې پر~\LR{$\Gamma$} سره يوه خوله دي۔ هر فلټريشن به له
\LR{$\mModel{M}$} څخه راغلي ټول غشي ولري، لکه په
\olref{fig:ser-eucl2} کښې چې ښودل شوي۔ دغه انځور تر ټولو نرے
فلټريشن ښيي، او اړيکه ئې اقليدسي نۀ ده۔
پر دې سربېره، د اقليدسي کولو لپاره ورسره نور غشي هم نۀ شو
زياتولے: د~\LR{$[w_2]$} او \LR{$[w_4]$} ترمنځ به په دواړو لوريو غشي
ورزياتوو، او بيا د~\LR{$w_2$} او~\LR{$w_5$} د ټولګيو ترمنځ هم۔ خو
\LR{$\Box p$} بايد په~\LR{$w_2$} کښې رښتينے وي، حال دا چې \LR{$p$} په~\LR{$w_5$}
کښې ناسم دے۔
\olpseganchor{OLP-0459-B006}{end}

\olpseganchor{OLP-0459-B007}{start}
\begin{figure}[htpb]
  \centering
  \begin{tikzpicture}[modal]
    \node[world] (w1)
          [label={left: \mFalse{p}},
            label={below: \LR{$\mSat{{}}{\Box p}$}}] {\LR{$w_1$}}; 
    \node[world] (w2)
          [label={right: \mTrue{p}},
            label = {below: \LR{$\mSat{{}}{\Box p}$}},
            right=of w1] {\LR{$w_2$}}; 
    \draw[->] (w1) to (w2);
    \draw[reflexive above] (w2) to (w2);
    \node[world] (w3)
          [label={left: \mFalse{p}},
            label={below: \LR{$\mSat{{}}{\Box p}$}},
            below=of w1] {\LR{$w_3$}}; 
    \node[world] (w4)
          [label={right:\mTrue{p}},
            label={below: \LR{$\mSat/{{}}{\Box p}$}},
            right=of w3] {\LR{$w_4$}};
    \draw[->] (w3) to (w4);
    \draw[reflexive above] (w4) to (w4);     
    \node[world] (w5)
          [label={right: \mFalse{p}},
            label=below:\LR{$\mSat/{{}}{\Box p}$},
            right=of w4] {\LR{$w_5$}}; 
    \draw[->, bend left] (w4) to (w5);
    \draw[->, bend left] (w5) to (w4);
    \draw[reflexive above] (w5) to (w5);     
  \end{tikzpicture}
  \caption{يو مسلسل او اقليدسي مدل۔}
  \ollabel{fig:ser-eucl}
\end{figure}
\olpseganchor{OLP-0459-B007}{end}

\olpseganchor{OLP-0459-B008}{start}
\begin{figure}[ht]
  \centering
  \begin{tikzpicture}[modal,every text node part/.style={align=center}]
    \node[world] (w1)
          [label={left: \mFalse{p}},
            label={right:\LR{$[w_1]=[w_3]$}},
            label={below: \LR{$\mSat{{}}{\Box p}$}}] {\LR{$[w_1]$}}; 
    \node[world] (w2)
         [label={right: \mTrue{p}},
           label = {below: \LR{$\mSat{{}}{\Box p}$}},
           right=of w1, yshift=1.5cm] {\LR{$[w_2]$}}; 
    \draw[->] (w1) to (w2);
    \draw[reflexive above] (w2) to (w2);
    \node[world] (w4)
         [label={right:\mTrue{p}},
           label={below: \LR{$\mSat/{{}}{\Box p}$}},
           right=of w1,yshift=-1.5cm] {\LR{$[w_4]$}};
    \draw[->] (w1) to (w4);
    \draw[reflexive above] (w4) to (w4);     
    \node[world] (w5)
         [label={right: \mFalse{p}},
             label=below:\LR{$\mSat/{{}}{\Box p}$},
             right=of w4] {\LR{$[w_5]$}}; 
    \draw[->, bend left] (w4) to (w5);
    \draw[->, bend left] (w5) to (w4);
    \draw[reflexive above] (w5) to (w5);     
  \end{tikzpicture}
  \caption{په \olref{fig:ser-eucl} کښې د مدل فلټريشن۔}
  \ollabel{fig:ser-eucl2}
\end{figure}
\olpseganchor{OLP-0459-B008}{end}

\olpseganchor{OLP-0459-B009}{start}
په تېره، د اقليدسي فلټريشن د ترلاسه کولو لپاره دا بس نۀ دي چې
د فرعي فارمولونه لپاره د هر تړلي سټ~\LR{$\Gamma$} له مخې
فلټريشنونه واخلو۔ پر ځای يې داسې سټونه~\LR{$\Gamma$} پکار دي چې
\emph{مودالي تړلي} وي (وګورئ
\olref[pre]{defn:modallyclosed})۔ دا ډول ناتش سټونه نامتناهي دي، ځکه پر يوه غړي په وار وار
بکس لګول هر وار نوے فارمول ورکوي۔ له همدې امله دا لار د اړوند
نظام د متناهي مدل خاصيت يا د پرېکړې وړتيا په خپله نۀ ثابتوي۔
\olpseganchor{OLP-0459-B009}{end}

\olpseganchor{OLP-0459-B010}{start}
\begin{thm}\ollabel{thm:modal-closed-filt}
  فرض کړئ \LR{$\Gamma$} مودالي تړلے وي،
  \LR{$\mModel{M}=\tuple{W,R,V}$} وي، او
  \LR{$\mModel{M^*} = \tuple{W^*,R^*,V^*}$} د
  \LR{$\mModel{M}$} تر ټولو غټ فلټريشن وي۔ د تناظر او اقليدسي‌والي
  په برخو کښې د تر ټولو غټې اړيکې د بکس او امکان دواړه شرطونه
  کاروو۔
  \begin{enumerate}
  \item که \LR{$\mModel{M}$} متناظر وي، \LR{$\mModel{M^*}$} هم متناظر دے۔
  \item که \LR{$\mModel{M}$} متعدي وي، \LR{$\mModel{M^*}$} هم متعدي دے۔
  \item که \LR{$\mModel{M}$} اقليدسي وي، \LR{$\mModel{M^*}$} هم اقليدسي دے۔
  \end{enumerate}
\end{thm}
\olpseganchor{OLP-0459-B010}{end}

\olpseganchor{OLP-0459-B011}{start}
\begin{proof}
\begin{enumerate}
  \item تمرين۔ له دې خبرې کار واخلئ چې \Ax{B} او
    \LR{$\Ax{B_\Diamond}$} دواړه په ټولو متناظرو مدلونو کښې معتبر دي۔
  \item که \LR{$\mModel{M^*}$} تر ټولو غټ فلټريشن وي، نو د تعريف له مخې
    \LR{$R^*[u][v]$} هله او يوازې هله وي چې \LR{$C_1(u,v)$} وي۔ د تعدي
    لپاره \LR{$C_1(u,v)$} او \LR{$C_1(v,w)$} فرض کړئ؛ بايد
    \LR{$C_1(u,w)$} وښيو۔ فرض کړئ
      \LR{$\mSat{M}{\Box !A}[u]$} وي؛ بيا
      \LR{$\mSat{M}{\Box\Box!A}[u]$} وي، ځکه \Ax{4} په ټولو
      متعديو مدلونو کښې معتبر دے۔ د تړلتيا له مخې
      \LR{$\Box\Box!A \in \Gamma$}؛ نو د \LR{$C_1(u,v)$} له مخې
      \LR{$\mSat{M}{\Box!A}[v]$}، او د \LR{$C_1(v,w)$} له مخې
      \LR{$\mSat{M}{!A}[w]$} هم وي۔ فرض کړئ
      \LR{$\mSat{M}{!A}[w]$} وي؛ بيا د \LR{$C_1(v,w)$} له مخې
      \LR{$\mSat{M}{\Diamond !A}[v]$} وي، ځکه د مودالي تړلتيا له مخې
      \LR{$\Diamond !A \in \Gamma$}۔ د \LR{$C_1(u,v)$} له مخې
      \LR{$\mSat{M}{\Diamond\Diamond !A}[u]$} راځي، ځکه د مودالي
      تړلتيا له مخې \LR{$\Diamond\Diamond!A \in \Gamma$}۔ ځکه چې
      \LR{$\Ax{4}_\Diamond$} په ټولو متعديو مدلونو کښې معتبر دے،
      \LR{$\mSat{M}{\Diamond!A}[u]$}۔
  \item تمرين۔ له دې خبرې کار واخلئ چې \Ax{5} او
    \LR{$\Ax{5_\Diamond}$} دواړه په ټولو اقليدسي مدلونو کښې معتبر دي۔
\end{enumerate}
\end{proof}
\olpseganchor{OLP-0459-B011}{end}

\olpseganchor{OLP-0459-B012}{start}
\begin{prob}
  د \olref[nml][fil][euc]{thm:modal-closed-filt} ثبوت بشپړ کړئ۔
\end{prob}
\olpseganchor{OLP-0459-B012}{end}

\olpseganchor{OLP-0460-B004}{start}
\olchapter{nml}{tab}{مودالي تابلوګانې}
\olpseganchor{OLP-0460-B004}{end}

\olpseganchor{OLP-0460-B005}{start}
\begin{editorial}
  دا باب د مودالي منطق مخکښ‌لرونکې تابلوګانې څېړي۔ د
  تابلوګانو قواعد، د ثبوت بېلګې، او له پرانيستو څانګو څخه
  د ضدمدلونو جوړول پکښې راځي۔
\end{editorial}
\olpseganchor{OLP-0460-B005}{end}

\olpseganchor{OLP-0461-B004}{start}
\olfileid{nml}{tab}{int}
\olpseganchor{OLP-0461-B004}{end}

\olpseganchor{OLP-0461-B005}{start}
\olsection{پېژندنه}
\olpseganchor{OLP-0461-B005}{end}

\olpseganchor{OLP-0461-B006}{start}
تابلوګانې د نښه لرونکي فارمولونه ځانګړې (ښکته څانګې جوړوونکې)
ونې دي؛ يعنې داسې جوړې چې د رښتياوالي نښه (\LR{$\True$} يا
\LR{$\False$}) او تړلے فارمول لري:
\[
\sFmla{\True}{!A} \text{\RL{ يا }} \sFmla{\False}{!A}.
\]
تابلو د يو شمېر \emph{فرضونو} څخه پيلېږي۔ هر بل
نښه لرونکے فارمول د يوې استنتاجي قاعدې په کارولو جوړېږي۔ ځينې
استنتاجي قاعدې د ونې يوې څوکې ته يو يا څو نښه لرونکي فارمولونه
زياتوي؛ نورې دوه نوې څوکې جوړوي، چې دوه څانګې ترې پيدا کېږي۔
د بياني رابطو په ماتولو داسې نښه لرونکي فارمولونه راځي چې د هر يوه
فارمول د هغې نښه لرونکے فارمول له فارمول څخه لږ پېچلے وي چې قاعده
پرې پلې شوې وه۔ کله چې په يوه څانګه کښې دواړه
\LR{$\sFmla{\True}{!A}$} او \LR{$\sFmla{\False}{!A}$} وي، څانګه
\emph{تړلې} بولو۔ که د تابلو هره څانګه تړلې وي، ټول
تابلو تړلے دے۔ تړلے تابلو داسې اشتقاق
دے چې ښيي د تابلو د پيل نښه لرونکي فارمولونه سټ نۀ شي پوره
کېدے۔ په دې سره د~\LR{$\Proves$} اړيکه تعريفولے شو:
\LR{$\Gamma \Proves !A$} هله او يوازې هله چې داسې متناهي
سټ~\LR{$\Gamma_0 = \{!B_1, \dots, !B_n\} \subseteq \Gamma$} وي
چې د لاندې فرضونو لپاره تړلے تابلو موجود وي:
\[
\{\sFmla{\False}{!A}, \sFmla{\True}{!B_1}, \dots, \sFmla{\True}{!B_n}\}.
\]
\olpseganchor{OLP-0461-B006}{end}

\olpseganchor{OLP-0461-B007}{start}
د مودالي منطقونو لپاره بايد د نښه لرونکے فارمول مفهوم وغځوو
او داسې قاعدې هم زياتې کړو چې
\LR{$\Box$} او
    \LR{$\Diamond$} راونغاړي۔ په مودالي تابلوګانې
کښې فارمولونه د نښې (\LR{$\True$} يا \LR{$\False$}) تر څنګ
\emph{مخکښونه}~\LR{$\sigma$} هم لري۔ مخکښونه د مثبتو صحيح عددونو
ناخالي لړۍ دي، يعنې \LR{$\sigma \in (\PosInt)^* \setminus
\{\emptyseq\}$}۔ موږ دا مخکښونه له شاوخوا
\LR{$\tuple{\ }$} پرته ليکو، او جلا غړي د~\LR{$,$} پر ځای
په~\LR{$.$} بېلوو۔ که \LR{$\sigma$} مخکښ وي، نو \LR{$\sigma.n$} له
\LR{$\sigma \concat \tuple{n}$} سره برابر دے؛ د بېلګې په توګه،
که \LR{$\sigma = 1.2.1$} وي، نو \LR{$\sigma.3$} برابر \LR{$1.2.1.3$}
دے۔ نو د بېلګې په توګه،
\[
\sFmla{\True}{\Box !A \lif !A}[1.2]
\]
يو \emph{مخکښ‌لرونکے نښه لرونکے فارمول} دے (يا په لنډ ډول
\emph{مخکښ‌لرونکے فارمول})۔ په مودالي تابلو کښې د
تړلو دپاره متضادې نښې بايد په عين فارمول او عين مخکښ پورې تړلې وي۔
\olpseganchor{OLP-0461-B007}{end}

\olpseganchor{OLP-0461-B008}{start}
د مفهوم له مخې، مخکښ په داسې مدل کښې د يوې نړۍ نوم دے چې ښايي
د تابلو د يوې څانګې فارمولونه پوره کړي؛ که \LR{$\sigma$}
د کومې نړۍ نوم وي، نو \LR{$\sigma.n$} د هغې نړۍ نوم دے چې له
(\LR{$\sigma$} نومې) نړۍ څخه د لاسرسي وړ وي۔ بېلابېل مخکښونه په يوه مدل کښې د يوې نړۍ نومونه
هم کېدے شي۔
\olpseganchor{OLP-0461-B008}{end}

\olpseganchor{OLP-0462-B004}{start}
\olfileid{nml}{tab}{rul}
\olpseganchor{OLP-0462-B004}{end}

\olpseganchor{OLP-0462-B005}{start}
\olsection{د \texorpdfstring{\Log{K}}{K} لپاره قاعدې}
\olpseganchor{OLP-0462-B005}{end}

\olpseganchor{OLP-0462-B006}{start}
د عامو قضيېيزو رابطو قاعدې د عامو نښه لرونکو قضيېيزو
تابلوګانې له قاعدو سره يو شان دي؛ يوازې مخکښونه ورسره
زياتېږي۔ په هر حالت کښې، پر نښه لرونکي فارمول
\LR{$\sFmla{S}{!A}[\sigma]$} پلې شوې قاعده نوي فارمولونه
جوړوي چې د هغو مخکښ هم~\LR{$\sigma$} وي۔ دا بايد له مفهومي پلوه
څرګنده وي: که \LR{$!A \land !B$} په~\LR{$\sigma$} نومې نړۍ کښې رښتينے
وي، نو \LR{$!A$} او \LR{$!B$} هم په~\LR{$\sigma$} کښې رښتيني دي۔ له دې
قاعدې د نورو نړۍو د رښتياوالي په اړه کوم حکم نۀ راووځي۔ قضيېيزې قاعدې په \olref{tab:prop-rules} کښې
راټولې شوې دي۔
\olpseganchor{OLP-0462-B006}{end}

\olpseganchor{OLP-0462-B007}{start}
\begin{table}
  \[\def\arraystretch{3}\begin{array}{|c|c|}
    \hline
    \AxiomC{\sFmla{\True}{\lnot !A}[\sigma]}
    \RightLabel{\TRule{\True}{\lnot}}
    \UnaryInfC{\sFmla{\False}{!A}[\sigma]}
    \DisplayProof
    &
    \AxiomC{\sFmla{\False}{\lnot !A}[\sigma]}
    \RightLabel{\TRule{\False}{\lnot}}
    \UnaryInfC{\sFmla{\True}{!A}[\sigma]}
    \DisplayProof
    \\[1ex]
    \hline
    \AxiomC{\sFmla{\True}{!A \land !B}[\sigma]}
    \RightLabel{\TRule{\True}{\land}}
    \UnaryInfC{\sFmla{\True}{!A}[\sigma]}
    \noLine
    \UnaryInfC{\sFmla{\True}{!B}[\sigma]}
    \DisplayProof
    &
    \AxiomC{\sFmla{\False}{!A \land !B}[\sigma]}
    \RightLabel{\TRule{\False}{\land}}
    \UnaryInfC{\LR{$\sFmla{\False}{!A}[\sigma] \quad \mid \quad
      \sFmla{\False}{!B}[\sigma]$}}
    \DisplayProof
    \\[2ex]
    \hline
    \AxiomC{\sFmla{\True}{!A \lor !B}[\sigma]}
    \RightLabel{\TRule{\True}{\lor}}
    \UnaryInfC{\LR{$\sFmla{\True}{!A}[\sigma] \quad \mid \quad
      \sFmla{\True}{!B}[\sigma]$}}
    \DisplayProof
    &
    \AxiomC{\sFmla{\False}{!A \lor !B}[\sigma]}
    \RightLabel{\TRule{\False}{\lor}}
    \UnaryInfC{\sFmla{\False}{!A}[\sigma]}
    \noLine
    \UnaryInfC{\sFmla{\False}{!B}[\sigma]}
    \DisplayProof
    \\[2ex]
    \hline
    \AxiomC{\sFmla{\True}{!A \lif !B}[\sigma]}
    \RightLabel{\TRule{\True}{\lif}}
    \UnaryInfC{\LR{$\sFmla{\False}{!A}[\sigma] \quad \mid
      \quad \sFmla{\True}{!B}[\sigma]$}}
    \DisplayProof
    &
    \AxiomC{\sFmla{\False}{!A \lif !B}[\sigma]}
    \RightLabel{\TRule{\False}{\lif}}
    \UnaryInfC{\sFmla{\True}{!A}[\sigma]}
    \noLine
    \UnaryInfC{\sFmla{\False}{!B}[\sigma]}
    \DisplayProof
    \\[2ex]
    \hline
  \end{array}\]
  \caption{د قضيېيزو رابطو لپاره د مخکښ‌لرونکي تابلو قاعدې}
  \ollabel{tab:prop-rules}
\end{table}
\olpseganchor{OLP-0462-B007}{end}

\olpseganchor{OLP-0462-B008}{start}
د تړلو شرط د عامو تابلوګانې د شرط په شان دے، خو دلته بايد
يوازې فارمولونه نه، بلکې مخکښونه هم سره يو شان وي۔ نو که
په يوه څانګه کښې دواړه
\[
\sFmla{\True}{!A}[\sigma] \quad\text{\RL{او}}\quad \sFmla{\False}{!A}[\sigma]
\]
د کوم مخکښ~\LR{$\sigma$} او فارمول~\LR{$!A$} لپاره وي، هغه څانګه
تړلې ده۔
\olpseganchor{OLP-0462-B008}{end}

\olpseganchor{OLP-0462-B009}{start}
د فرضونو د ايښودلو قاعدې هم د عامو تابلوګانې په شان دي، خو
د فرضونو لپاره تل مخکښ~\LR{$1$} کاروو۔ (د مخکښ ټاکنه مهمه نۀ ده،
تر هغو چې د ټولو فرضونو لپاره يو شان وي۔) نو، د بېلګې په توګه،
\[
!B_1, \dots, !B_n \Proves !A
\]
هله او يوازې هله وايو چې د لاندې فرضونو لپاره تړلے تابلو وي:
\[
\sFmla{\True}{!B_1}[1], \dots, \sFmla{\True}{!B_n}[1],
\sFmla{\False}{!A}[1].
\]
\olpseganchor{OLP-0462-B009}{end}

\olpseganchor{OLP-0462-B010}{start}
د مودالي نښو~\LR{$\Box$}
    او~\LR{$\Diamond$}
لپاره، پر هغه فارمول چې مخکښ يې~\LR{$\sigma$} دے د قاعدې
د پلي کېدو پايلې مخکښ \LR{$\sigma.n$} وي۔ خو کوم \LR{$n$} روا دے، دا
په دې پورې تړلې ده چې نښه~\LR{$\True$} ده که~\LR{$\False$}۔
\olpseganchor{OLP-0462-B010}{end}

\olpseganchor{OLP-0462-B011}{start}
د \LR{$\TRule{\True}{\Box}$} قاعده هغه څانګه چې
  \LR{$\sFmla{\True}{\Box !A}[\sigma]$} لري، په
  \LR{$\sFmla{\True}{!A}[\sigma.n]$} غځوي۔
    همداسې، د
  \LR{$\TRule{\False}{\Diamond}$} قاعده هغه څانګه چې
  \LR{$\sFmla{\False}{\Diamond !A}[\sigma]$} لري، په
  \LR{$\sFmla{\False}{!A}[\sigma.n]$} غځوي۔
دا قاعدې يوازې د
هغې~\LR{$\sigma.n$} لپاره پلې کېدے شي چې په اړونده څانګه کښې
\emph{لا مخکې} راغلې وي۔ داسې مخکښ په څانګه کښې «کارېدلے» بولو۔
\olpseganchor{OLP-0462-B011}{end}

\olpseganchor{OLP-0462-B012}{start}
د \LR{$\TRule{\False}{\Box}$} قاعده هغه څانګه چې
  \LR{$\sFmla{\False}{\Box !A}[\sigma]$} لري، په
  \LR{$\sFmla{\False}{!A}[\sigma.n]$} غځوي۔
    همداسې، د
  \LR{$\TRule{\True}{\Diamond}$} قاعده هغه څانګه چې
  \LR{$\sFmla{\True}{\Diamond !A}[\sigma]$} لري، په
  \LR{$\sFmla{\True}{!A}[\sigma.n]$} غځوي۔
خو دا قاعدې يوازې
د هغې~\LR{$\sigma.n$} لپاره پلې کېدے شي چې په اړونده څانګه کښې
\emph{لا مخکې نۀ} وي راغلې۔ داسې مخکښ په څانګه کښې «نوے» بولو۔
\olpseganchor{OLP-0462-B012}{end}

\olpseganchor{OLP-0462-B013}{start}
قاعدې په \olref{tab:rules-K} کښې ورکړل شوې دي۔
\olpseganchor{OLP-0462-B013}{end}

\olpseganchor{OLP-0462-B014}{start}
\begin{table}
  \begin{center}
    \def\arraystretch{3}\def\fCenter{}
    \begin{tabular}{|c|c|}
    \hline
    
      \AxiomC{\sFmla{\True}{\Box !A}[\sigma]}
      \RightLabel{\TRule{\True}{\Box}}
      \UnaryInfC{\sFmla{\True}{!A}[\sigma.n]}
      \DisplayProof
      &
      \AxiomC{\sFmla{\False}{\Box !A}[\sigma]}
      \RightLabel{\TRule{\False}{\Box}}
      \UnaryInfC{\sFmla{\False}{!A}[\sigma.n]}
      \DisplayProof\\
      \LR{$\sigma.n$} کارېدلے دے & \LR{$\sigma.n$} نوے دے
      \\[1ex]
      \hline
    
      \AxiomC{\sFmla{\True}{\Diamond !A}[\sigma]}
      \RightLabel{\TRule{\True}{\Diamond}}
      \UnaryInfC{\sFmla{\True}{!A}[\sigma.n]}
      \DisplayProof
      &
      \AxiomC{\sFmla{\False}{\Diamond !A}[\sigma]}
      \RightLabel{\TRule{\False}{\Diamond}}
      \UnaryInfC{\sFmla{\False}{!A}[\sigma.n]}
      \DisplayProof\\
      \LR{$\sigma.n$} نوے دے & \LR{$\sigma.n$} کارېدلے دے
      \\[1ex]
      \hline
    \end{tabular}
  \end{center}
  \caption{د \Ax{K} مودالي قاعدې۔}
  \ollabel{tab:rules-K}
\end{table}
\olpseganchor{OLP-0462-B014}{end}

\olpseganchor{OLP-0462-B015}{start}
د \TRule{\True}{\Box}
لپاره د مخکښ د مخکې کارېدلو قيد اړين دے؛ ګنې لاندې تابلو به
هم تړلے تابلو ګڼل کېده:

  \begin{oltableau}
    [\pFmla{\True}{\Box \formula{A}}{1}, just = \TAss
      [\pFmla{\False}{\Diamond \formula{A}}{1}, just = \TAss
        [\pFmla{\True}{\formula{A}}{1.1}, just = {\TRule{\True}{\Box}[1]}
          [\pFmla{\False}{\formula{A}}{1.1},
            just = {\TRule{\False}{\Diamond}[2]}, close]
        ]
      ]
    ]
\end{oltableau}
\olpseganchor{OLP-0462-B015}{end}

\olpseganchor{OLP-0462-B016}{start}
خو \LR{$\Box \formula{A} \Entails/ \Diamond \formula{A}$}؛ نو زموږ
ثبوتي نظام به ناسم وي۔ همداسې، \LR{$\Diamond \formula{A} \Entails/
\Box \formula{A}$}؛ خو که د
\TRule{\False}{\Box}
لپاره د نوي مخکښ قيد نۀ وي، لاندې تابلو به تړلے وي:

  \begin{oltableau}
    [\pFmla{\True}{\Diamond \formula{A}}{1}, just = \TAss,name=one
      [\pFmla{\False}{\Box \formula{A}}{1}, just = \TAss, name=two
        [\pFmla{\True}{\formula{A}}{1.1}, just = {\TRule{\True}{\Diamond}[1]}
          [\pFmla{\False}{\formula{A}}{1.1},
            just = {\TRule{\False}{\Box}[2]}, close]
        ]
      ]
    ]
  \end{oltableau}
\olpseganchor{OLP-0462-B016}{end}

\olpseganchor{OLP-0463-B004}{start}
\olfileid{nml}{tab}{prk}
\olpseganchor{OLP-0463-B004}{end}

\olpseganchor{OLP-0463-B005}{start}
\olsection{د \texorpdfstring{\Log{K}}{K} لپاره تابلوګانې}
\olpseganchor{OLP-0463-B005}{end}



\olpseganchor{OLP-0463-B007}{start}
\begin{ex}
  داسې تړلے تابلو ورکوو چې \LR{$\Proves (\Box!A \land \Box!B)
  \lif \Box (!A \land !B)$} ښيي۔ د نارښتيني بکس قاعده نوے
  مخکښ جوړوي۔ د دواړو رښتينو بکسونو قاعدې بيا همدا کارېدلے
  مخکښ کاروي او دواړه څانګې تړي۔
  \begin{oltableau}
    [\pFmla{\False}{(\Box\formula{A} \land \Box\formula{B}) \lif
        \Box (\formula{A} \land \formula{B})}{1},
      just =\TAss
      [\pFmla{\True}{\Box\formula{A} \land \Box\formula{B}}{1},
        just = {\TRule{\False}{\lif}[1]}
        [\pFmla{\False}{\Box(\formula{A} \land \formula{B})}{1},
          just = {\TRule{\False}{\lif}[1]}
          [\pFmla{\True}{\Box\formula{A}}{1},
            just = {\TRule{\True}{\land}[2]}
            [\pFmla{\True}{\Box\formula{B}}{1},
              just = {\TRule{\True}{\land}[2]}
              [\pFmla{\False}{\formula{A} \land \formula{B}}{1.1},
                just = {\TRule{\False}{\Box}[3]}
                [\pFmla{\False}{\formula{A}}{1.1},
                  just = {\TRule{\False}{\land}[6]}
                  [\pFmla{\True}{\formula{A}}{1.1},
                    just= {\TRule{\True}{\Box}[4]}, close]]
                [\pFmla{\False}{\formula{B}}{1.1},
                  just = {\TRule{\False}{\land}[6]}
                  [\pFmla{\True}{\formula{B}}{1.1},
                    just= {\TRule{\True}{\Box}[5]}, close]]
              ]
            ]
          ]
        ]
      ]
    ]
  \end{oltableau}
\end{ex}

\olpseganchor{OLP-0463-B007}{end}

\olpseganchor{OLP-0463-B008}{start}

\begin{ex}
  داسې تړلے تابلو ورکوو چې \LR{$\Proves \Diamond(!A \lor !B)
  \lif (\Diamond !A \lor \Diamond !B)$} ښيي۔ رښتينے امکان
  يو نوے مخکښ جوړوي، او د فصل دوه څانګې د نارښتينو امکانونو
  له قاعدو سره پر هماغه مخکښ تړل کېږي۔
  \begin{oltableau}
    [\pFmla{\False}{\Diamond(\formula{A} \lor \formula{B}) \lif
        (\Diamond \formula{A} \lor \Diamond \formula{B})}{1},
      just =\TAss
      [\pFmla{\True}{\Diamond(\formula{A} \lor \formula{B})}{1},
        just = {\TRule{\False}{\lif}[1]}
        [\pFmla{\False}{\Diamond\formula{A} \lor \Diamond\formula{B}}{1},
          just = {\TRule{\False}{\lif}[1]}
          [\pFmla{\False}{\Diamond\formula{A}}{1},
            just = {\TRule{\False}{\lor}[3]}
            [\pFmla{\False}{\Diamond\formula{B}}{1},
              just = {\TRule{\False}{\lor}[3]}
              [\pFmla{\True}{\formula{A} \lor \formula{B}}{1.1},
                just = {\TRule{\True}{\Diamond}[2]},
                [\pFmla{\True}{\formula{A}}{1.1},
                  just = {\TRule{\True}{\lor}[6]}
                  [\pFmla{\False}{\formula{A}}{1.1},
                    just= {\TRule{\False}{\Diamond}[4]}, close]]
                [\pFmla{\True}{\formula{B}}{1.1},
                  just = {\TRule{\True}{\lor}[6]}
                  [\pFmla{\False}{\formula{B}}{1.1},
                    just= {\TRule{\False}{\Diamond}[5]}, close]]
              ]
            ]
          ]
        ]
      ]
    ]
  \end{oltableau}
\end{ex}

\olpseganchor{OLP-0463-B008}{end}


\olpseganchor{OLP-0463-B009}{start}
\begin{prob}
  د لاندې فارمولونه لپاره په~\LR{$\Log{K}$} کښې تړلي
  تابلوګانې ومومئ:
  \begin{enumerate}
    \item \LR{$\Box \lnot p \lif \Box(p \lif q)$}
    \item \LR{$(\Box p \lor \Box q) \lif \Box(p \lor q)$}
    \item \LR{$\Diamond p \lif \Diamond(p \lor q)$}
    \item \LR{$\Box(p \land q) \lif \Box p$}
  \end{enumerate}
\end{prob}
\olpseganchor{OLP-0463-B009}{end}

\olpseganchor{OLP-0464-B004}{start}
\olfileid{nml}{tab}{sou}
\olpseganchor{OLP-0464-B004}{end}

\olpseganchor{OLP-0464-B005}{start}
\olsection{د \texorpdfstring{\Log{K}}{K} سموالی}
\olpseganchor{OLP-0464-B005}{end}

\olpseganchor{OLP-0464-B006}{start}
\begin{editorial}
  دا ثبوت د قضيېيزو رابطو او مودالي عملګرو قاعدې په يوه چوکاټ
  کښې څېړي۔ د مخکښونو تفسير پر څانګه د نړۍو د لاسرسي اړيکه ساتي۔
\end{editorial}
\olpseganchor{OLP-0464-B006}{end}

\olpseganchor{OLP-0464-B007}{start}
\begin{explain}
  د مخکښ‌لرونکو تابلوګانې د سموالي د ښودلو لپاره بايد ثابته
  کړو چې که
  \[
  \sFmla{\True}{!B_1}[1], \dots, \sFmla{\True}{!B_n}[1], \sFmla{\False}{!A}[1]
  \]
  تړلے تابلو ولري، نو \LR{$!B_1, \dots, !B_n \Entails !A$} وي۔
  د دې مقابل نقيض ثابتول اسانه دي: که په کوم~\LR{$\mModel{M}$} کښې د
  نړۍ~\LR{$w$} لپاره \LR{$\mSat{M}{!B_i}[w]$} د هر \LR{$i=1$}, \dots,~\LR{$n$}
  په حق کښې وي، خو \LR{$\mSat/{M}{!A}[w]$} وي، نو هېڅ
  تابلو تړلے نۀ شي جوړېدے۔
  داسې ضدمدل ښيي چې د تابلو پيلني فرضونه پوره کېدے شي۔
  د ثبوت لاره دا ده چې کله د تابلو پر يوه څانګه ټول
  مخکښ‌لرونکي فارمولونه پوره کېدے شي، د هرې قاعدې پلي کېدل
  لږ تر لږه يوه غځېدلې څانګه پرېږدي چې هغه هم پوره کېدے شي۔
  تړلې څانګې نۀ شي پوره کېدے؛ نو د مخکښ‌لرونکو فارمولونه
  د پوره کېدونکي سټ هر تابلو بايد لږ تر لږه يوه پرانيستې
  څانګه ولري۔
\olpseganchor{OLP-0464-B007}{end}

\olpseganchor{OLP-0464-B008}{start}
  د دې لارې د مودالي حالت لپاره کارولو ته بايد د «پوره کېدو»
  تعريف مودالي مدلونو او مخکښونو ته وغځوو۔
\end{explain}
\olpseganchor{OLP-0464-B008}{end}

\olpseganchor{OLP-0464-B009}{start}
\begin{defn}
  فرض کړئ \LR{$P$} د مخکښونو يو سټ وي، يعنې
  \LR{$P \subseteq (\PosInt)^* \setminus \{\emptyseq\}$}، او
  \LR{$\mModel{M}$} يو مدل وي۔ تابع~\LR{$f\colon P \to W$} په
  \LR{$\mModel{M}$} کښې د~\LR{$P$} \emph{تفسير} دے که هرکله چې
  \LR{$\sigma$} او \LR{$\sigma.n$} دواړه په~\LR{$P$} کښې وي، نو
  \LR{$Rf(\sigma)f(\sigma.n)$} وي۔
\olpseganchor{OLP-0464-B009}{end}

\olpseganchor{OLP-0464-B010}{start}
  د مخکښونو د~\LR{$P$} د تفسير په نسبت داسې تعريفولے شو:
  \begin{enumerate}
  \item \LR{$\mModel{M}$} هله او يوازې هله
    \LR{$\sFmla{\True}{!A}[\sigma]$} پوره کوي چې
    \LR{$\mSat{M}{!A}[f(\sigma)]$} وي۔
  \item \LR{$\mModel{M}$} هله او يوازې هله
    \LR{$\sFmla{\False}{!A}[\sigma]$} پوره کوي چې
    \LR{$\mSat/{M}{!A}[f(\sigma)]$} وي۔
  \end{enumerate}
\end{defn}
\olpseganchor{OLP-0464-B010}{end}

\olpseganchor{OLP-0464-B011}{start}
\begin{defn}
  فرض کړئ \LR{$\Gamma$} د مخکښ‌لرونکو فارمولونه سټ وي، او
  \LR{$P(\Gamma)$} د هغو مخکښونو سټ وي چې په دې سټ کښې راځي۔
  که \LR{$f$} په~\LR{$\mModel{M}$} کښې د~\LR{$P(\Gamma)$} تفسير وي، نو وايو
  چې \LR{$\mModel{M}$} د~\LR{$f$} په نسبت \LR{$\Gamma$} پوره کوي،
  \LR{$\mSat{M}{\Gamma}[f]$}، که \LR{$\mModel{M}$} په~\LR{$\Gamma$} کښې
  هر مخکښ‌لرونکے فارمول د~\LR{$f$} په نسبت پوره کوي۔
  \LR{$\Gamma$} هله او يوازې هله \emph{پوره کېدونکے} دے چې کوم
  مدل~\LR{$\mModel{M}$} او د~\LR{$P(\Gamma)$} تفسير~\LR{$f$} وي داسې چې
  \LR{$\mSat{M}{\Gamma}[f]$} وي۔
\end{defn}
\olpseganchor{OLP-0464-B011}{end}

\olpseganchor{OLP-0464-B012}{start}
\begin{prop}
  که \LR{$\Gamma$} د کوم فارمول~\LR{$!A$} او مخکښ~\LR{$\sigma$} لپاره
  دواړه \LR{$\sFmla{\True}{!A}[\sigma]$} او
  \LR{$\sFmla{\False}{!A}[\sigma]$} ولري، نو \LR{$\Gamma$} نۀ شي
  پوره کېدے۔
\end{prop}
\olpseganchor{OLP-0464-B012}{end}

\olpseganchor{OLP-0464-B013}{start}
\begin{proof}
  داسې مدل~\LR{$\mModel{M}$} او د~\LR{$P(\Gamma)$} تفسير~\LR{$f$} نۀ شي
  کېدے چې دواړه \LR{$\mSat{M}{!A}[f(\sigma)]$} او
  \LR{$\mSat/{M}{!A}[f(\sigma)]$} پوره وي۔
\end{proof}
\olpseganchor{OLP-0464-B013}{end}

\olpseganchor{OLP-0464-B014}{start}
\begin{thm}[سموالی]
  \ollabel{thm:tableau-soundness}
  فرض کړئ د~\LR{$\Gamma$} د هر مخکښ ټولې ناتشې پيلنۍ برخې هم په
  همدې سټ کښې د فارمولونو مخکښونه وي۔ که دا سټ تړلے تابلو
  ولري، نو \LR{$\Gamma$} نۀ شي پوره کېدے۔
\end{thm}
\olpseganchor{OLP-0464-B014}{end}

\olpseganchor{OLP-0464-B015}{start}
\begin{proof}
د تابلو يوه څانګه هله پوره کېدونکې بولو چې پر هغې د
نښه لرونکي فارمولونه سټ پوره کېدونکے وي؛ او تابلو هله
پوره کېدونکے بولو چې لږ تر لږه يوه پوره کېدونکې څانګه ولري۔
\olpseganchor{OLP-0464-B015}{end}

\olpseganchor{OLP-0464-B016}{start}
دا ښيو: د استنتاج په هره قاعده د پوره کېدونکي تابلو
غځول بيا هم پوره کېدونکے تابلو جوړوي۔ له دې نه قضيه
ثابتېږي: د تړلي تابلو پيلنيو فرضونو ته په کتو، د قاعدو
په پلي کولو له هماغه پيلني تابلو څخه د~\LR{$\Gamma$} تړلې
تابلو جوړېږي۔ قاعده يا هماغه مخکښ ساتي يا پر يوه موجود
مخکښ يوه نوې شمېره پسې نښلوي؛ نو د پيلنيو برخو شرط په هره
څانګه کښې پاتې کېږي۔ که \LR{$\Gamma$} پوره کېدونکے وای، د هغه
هر تابلو به هم پوره کېدونکے و۔ خو تړلے تابلو
ښکاره نۀ شي پوره کېدے، ځکه ټولې څانګې يې تړلې او تړلې
څانګې نۀ‌پوره کېدونکې دي۔
\olpseganchor{OLP-0464-B016}{end}

\olpseganchor{OLP-0464-B017}{start}
فرض کړئ يو پوره کېدونکے تابلو لرو، يعنې داسې
تابلو چې لږ تر لږه يوه پوره کېدونکې څانګه ولري۔
استنتاجي قاعده يا يوې څانګې ته نښه لرونکي فارمولونه زياتوي، يا
هغه په دوو څانګو وېشي۔ که په تابلو کښې يوه پوره کېدونکې
څانګه وي چې دا قاعده پرې نۀ پلې کېږي، هغه په غځېدلي
تابلو کښې هم پوره کېدونکې پاتې کېږي؛ نو غځېدلے تابلو
پوره کېدونکے دے۔ يوازې هغه حالت پکار دے چې قاعده پر پوره
کېدونکې څانګه پلې شي۔
\olpseganchor{OLP-0464-B017}{end}

\olpseganchor{OLP-0464-B018}{start}
فرض کړئ \LR{$\Gamma$} پر دې څانګه د نښه لرونکي فارمولونه سټ وي او
\LR{$\sFmla{S}{!A}[\sigma] \in \Gamma$} هغه نښه لرونکے فارمول
وي چې قاعده پرې پلې کېږي۔ که قاعده څانګه ونۀ وېشي، بايد وښيو
چې غځېدلې څانګه، يعنې \LR{$\Gamma$} او د قاعدې پايلې، لا
هم پوره کېدونکې ده۔ که قاعده څانګه ووېشي، بايد وښيو چې له
دوو نويو څانګو لږ تر لږه يوه پوره کېدونکې ده۔
لومړے هغه استنتاجونه ګورو چې يوازې يوه مقدمه لري۔
\begin{enumerate}
\item پر \LR{$\sFmla{\True}{\lnot !B}[\sigma] \in \Gamma$} د
  \LR{$\TRule{\True}{\lnot}$} په پلي کولو څانګه غځېږي۔ نو په
  غځېدلې څانګه کښې د نښه لرونکي فارمولونه سټ \LR{$\Gamma \cup
  \{\sFmla{\False}{!B}[\sigma]\}$} وي۔ فرض کړئ
  \LR{$\mSat{M}{\Gamma}[f]$}؛ په ځانګړې توګه
  \LR{$\mSat{M}{\lnot !B}[f(\sigma)]$} وي۔ نو
  \LR{$\mSat/{M}{!B}[f(\sigma)]$}؛ يعنې \LR{$\mModel{M}$} د~\LR{$f$} په
  نسبت \LR{$\sFmla{\False}{!B}[\sigma]$} پوره کوي۔
\item پر \LR{$\sFmla{\False}{\lnot !B}[\sigma] \in \Gamma$} د
  \LR{$\TRule{\False}{\lnot}$} په پلي کولو څانګه غځېږي: تمرين۔
\item پر \LR{$\sFmla{\True}{!B \land !C}[\sigma] \in \Gamma$} د
  \LR{$\TRule{\True}{\land}$} په پلي کولو څانګه غځېږي او دوه نوي
  نښه لرونکي فارمولونه ورکوي:
  \LR{$\sFmla{\True}{!B}[\sigma]$} او \LR{$\sFmla{\True}{!C}[\sigma]$}۔
  فرض کړئ \LR{$\mSat{M}{\Gamma}[f]$}، په ځانګړې توګه
  \LR{$\mSat{M}{!B \land !C}[f(\sigma)]$}۔ نو
  \LR{$\mSat{M}{!B}[f(\sigma)]$} او \LR{$\mSat{M}{!C}[f(\sigma)]$}۔
  يعنې \LR{$\mModel{M}$} د~\LR{$f$} په نسبت دواړه
  \LR{$\sFmla{\True}{!B}[\sigma]$} او
  \LR{$\sFmla{\True}{!C}[\sigma]$} پوره کوي۔
\item پر \LR{$\sFmla{\False}{!B \lor !C}[\sigma] \in \Gamma$} د
  \LR{$\TRule{\False}{\lor}$} په پلي کولو څانګه غځېږي: تمرين۔
\item پر \LR{$\sFmla{\False}{!B \lif !C}[\sigma] \in \Gamma$} د
  \LR{$\TRule{\False}{\lif}$} په پلي کولو څانګه غځېږي۔ په څانګه
  کښې دوه نوي نښه لرونکي فارمولونه راځي:
  \LR{$\sFmla{\True}{!B}[\sigma]$} او
  \LR{$\sFmla{\False}{!C}[\sigma]$}۔ فرض کړئ
  \LR{$\mSat{M}{\Gamma}[f]$}، په ځانګړې توګه
  \LR{$\mSat/{M}{!B \lif !C}[f(\sigma)]$}۔ نو
  \LR{$\mSat{M}{!B}[f(\sigma)]$} او
  \LR{$\mSat/{M}{!C}[f(\sigma)]$}۔ يعنې \LR{$\mModel{M}, f$} دواړه
  \LR{$\sFmla{\True}{!B}[\sigma]$} او
  \LR{$\sFmla{\False}{!C}[\sigma]$} پوره کوي۔
\olpseganchor{OLP-0464-B018}{end}

\olpseganchor{OLP-0464-B019}{start}
%
\item پر \LR{$\sFmla{\True}{\Box !B}[\sigma] \in \Gamma$} د
  \LR{$\TRule{\True}{\Box}$} په پلي کولو څانګه غځېږي:
  په څانګه کښې نوے نښه لرونکے فارمول~\LR{$\sFmla{\True}{!B}[\sigma.n]$} راځي، چېرته چې
    \LR{$\sigma.n \in P(\Gamma)$} (ځکه \LR{$\sigma.n$} بايد کارېدلے وي)۔
    فرض کړئ \LR{$\mSat{M}{\Gamma}[f]$}، په ځانګړې توګه
    \LR{$\mSat{M}{\Box !B}[f(\sigma)]$}۔ ځکه \LR{$f$} د مخکښونو تفسير دے او
    دواړه \LR{$\sigma$}, \LR{$\sigma.n \in P(\Gamma)$} دي، نو
    \LR{$Rf(\sigma)f(\sigma.n)$} وي۔ له دې نه
    \LR{$\mSat{M}{!B}[f(\sigma.n)]$}؛ يعنې \LR{$\mModel{M}, f$} د
    \LR{$\sFmla{\True}{!B}[\sigma.n]$} پوره کوي۔
\olpseganchor{OLP-0464-B019}{end}

\olpseganchor{OLP-0464-B020}{start}
\item پر \LR{$\sFmla{\False}{\Box !B}[\sigma] \in \Gamma$} د
  \LR{$\TRule{\False}{\Box}$} په پلي کولو څانګه غځېږي:
  نوے نښه لرونکے فارمول~\LR{$\sFmla{\False}{!B}[\sigma.n]$} راځي؛
    \LR{$\sigma.n$} په څانګه کښې نوے مخکښ دے، يعنې
    \LR{$\sigma.n \notin P(\Gamma)$}۔ ځکه \LR{$\Gamma$} پوره کېدونکے
    دے، داسې مودالي
    مدل~\LR{$\mModel{M}$} او د~\LR{$P(\Gamma)$} تفسير~\LR{$f$} شته چې
    \LR{$\mSat{M}{\Gamma}[f]$} وي؛ په ځانګړې توګه
    \LR{$\mSat/{M}{\Box !B}[f(\sigma)]$}۔ بايد وښيو چې
    \LR{$\Gamma \cup \{\sFmla{\False}{!B}[\sigma.n]\}$} پوره
    کېدونکے دے۔ د دې لپاره د \LR{$P(\Gamma) \cup \{\sigma.n\}$}
    تفسير داسې تعريفوو:
\olpseganchor{OLP-0464-B020}{end}

\olpseganchor{OLP-0464-B021}{start}
  ځکه \LR{$\mSat/{M}{\Box !B}[f(\sigma)]$}، داسې \LR{$w \in W$} شته چې
  \LR{$Rf(\sigma)w$} او \LR{$\mSat/{M}{!B}[w]$} وي۔ \LR{$f'$} د~\LR{$f$} په شان
  واخلئ، يوازې \LR{$f'(\sigma.n) = w$} دې وي۔ ځکه
  \LR{$f'(\sigma) = f(\sigma)$} او \LR{$Rf(\sigma)w$}، نو
  \LR{$Rf'(\sigma)f'(\sigma.n)$} وي۔ د نوې پيلنۍ برخې هېڅ اوږد غځون
له مخکې په څانګه کښې نشته؛ نو دا يوازينے نوے لاسرسے دے چې
تفسير ئې بايد وساتي۔ ځکه \LR{$f'$} د
  \LR{$P(\Gamma) \cup \{\sigma.n\}$} تفسير دے۔ ښکاره ده چې
  \LR{$\mSat/{M}{!B}[f'(\sigma.n)]$} وي۔ د هر
  \LR{$\sigma' \in P(\Gamma)$} لپاره \LR{$f(\sigma') = f'(\sigma')$}،
  نو \LR{$\mSat{M}{\Gamma}[f']$}۔ په دې توګه \LR{$\mModel{M}, f'$} د
  \LR{$\Gamma \cup \{\sFmla{\False}{!B}[\sigma.n]\}$} پوره کوي۔
  
\olpseganchor{OLP-0464-B021}{end}

\olpseganchor{OLP-0464-B022}{start}
%
\item پر \LR{$\sFmla{\True}{\Diamond !B}[\sigma] \in \Gamma$} د
  \LR{$\TRule{\True}{\Diamond}$} په پلي کولو څانګه غځېږي:
  نوے نښه لرونکے فارمول~\LR{$\sFmla{\True}{!B}[\sigma.n]$} راځي؛
    \LR{$\sigma.n$} په څانګه کښې نوے مخکښ دے، يعنې
    \LR{$\sigma.n \notin P(\Gamma)$}۔ ځکه \LR{$\Gamma$} پوره کېدونکے
    دے، داسې مودالي
    مدل~\LR{$\mModel{M}$} او د~\LR{$P(\Gamma)$} تفسير~\LR{$f$} شته چې
    \LR{$\mSat{M}{\Gamma}[f]$} وي؛ په ځانګړې توګه
    \LR{$\mSat{M}{\Diamond !B}[f(\sigma)]$}۔ بايد وښيو چې
    \LR{$\Gamma \cup \{\sFmla{\True}{!B}[\sigma.n]\}$} پوره
    کېدونکے دے۔ د دې لپاره د \LR{$P(\Gamma) \cup \{\sigma.n\}$}
    تفسير داسې تعريفوو:
\olpseganchor{OLP-0464-B022}{end}

\olpseganchor{OLP-0464-B023}{start}
  ځکه \LR{$\mSat{M}{\Diamond !B}[f(\sigma)]$}، داسې \LR{$w \in W$} شته چې
  \LR{$Rf(\sigma)w$} او \LR{$\mSat{M}{!B}[w]$} وي۔ \LR{$f'$} د~\LR{$f$} په شان
  واخلئ، يوازې \LR{$f'(\sigma.n) = w$} دې وي۔ ځکه
  \LR{$f'(\sigma) = f(\sigma)$} او \LR{$Rf(\sigma)w$}، نو
  \LR{$Rf'(\sigma)f'(\sigma.n)$} وي۔ د نوې پيلنۍ برخې هېڅ اوږد غځون
له مخکې په څانګه کښې نشته؛ نو دا يوازينے نوے لاسرسے دے چې
تفسير ئې بايد وساتي۔ ځکه \LR{$f'$} د
  \LR{$P(\Gamma) \cup \{\sigma.n\}$} تفسير دے۔ ښکاره ده چې
  \LR{$\mSat{M}{!B}[f'(\sigma.n)]$} وي۔ د هر
  \LR{$\sigma' \in P(\Gamma)$} لپاره \LR{$f(\sigma') = f'(\sigma')$}،
  نو \LR{$\mSat{M}{\Gamma}[f']$}۔ په دې توګه \LR{$\mModel{M}, f'$} د
  \LR{$\Gamma \cup \{\sFmla{\True}{!B}[\sigma.n]\}$} پوره کوي۔
\item پر \LR{$\sFmla{\False}{\Diamond !B}[\sigma] \in \Gamma$} د
  \LR{$\TRule{\False}{\Diamond}$} په پلي کولو څانګه غځېږي:
  په څانګه کښې نوے نښه لرونکے فارمول~\LR{$\sFmla{\False}{!B}[\sigma.n]$} راځي، چېرته چې
    \LR{$\sigma.n \in P(\Gamma)$} (ځکه \LR{$\sigma.n$} بايد کارېدلے وي)۔
    فرض کړئ \LR{$\mSat{M}{\Gamma}[f]$}، په ځانګړې توګه
    \LR{$\mSat/{M}{\Diamond !B}[f(\sigma)]$}۔ ځکه \LR{$f$} د مخکښونو
    تفسير دے او دواړه \LR{$\sigma$}, \LR{$\sigma.n \in P(\Gamma)$} دي،
    نو \LR{$Rf(\sigma)f(\sigma.n)$} وي۔ له دې نه
    \LR{$\mSat/{M}{!B}[f(\sigma.n)]$}؛ يعنې \LR{$\mModel{M}, f$} د
    \LR{$\sFmla{\False}{!B}[\sigma.n]$} پوره کوي۔
  
\end{enumerate}
اوس هغه استنتاجونه ګورو چې له يوې مقدمې دوه پايله‌څانګې جوړوي۔
\begin{enumerate}
\item پر \LR{$\sFmla{\False}{!B \land !C}[\sigma] \in \Gamma$} د
  \LR{$\TRule{\False}{\land}$} په پلي کولو څانګه غځېږي او دوه
  څانګې جوړوي: کيڼه له \LR{$\sFmla{\False}{!B}[\sigma]$} څخه او
  ښۍ له \LR{$\sFmla{\False}{!C}[\sigma]$} څخه دوام کوي۔ فرض کړئ
  \LR{$\mSat{M}{\Gamma}[f]$}، په ځانګړې توګه
  \LR{$\mSat/{M}{!B \land !C}[f(\sigma)]$}۔ نو يا
  \LR{$\mSat/{M}{!B}[f(\sigma)]$} وي يا
  \LR{$\mSat/{M}{!C}[f(\sigma)]$}۔ په لومړي حالت کښې
  \LR{$\mModel{M}, f$} د \LR{$\sFmla{\False}{!B}[\sigma]$} پوره کوي،
  يعنې کيڼه څانګه پوره کېدونکې ده۔ په دويم حالت کښې
  \LR{$\mModel{M}, f$} د \LR{$\sFmla{\False}{!C}[\sigma]$} پوره کوي،
  يعنې ښۍ څانګه پوره کېدونکې ده۔
\item پر \LR{$\sFmla{\True}{!B \lor !C}[\sigma] \in \Gamma$} د
  \LR{$\TRule{\True}{\lor}$} په پلي کولو څانګه غځېږي: تمرين۔
\item پر \LR{$\sFmla{\True}{!B \lif !C}[\sigma] \in \Gamma$} د
  \LR{$\TRule{\True}{\lif}$} په پلي کولو څانګه غځېږي: تمرين۔
\end{enumerate}
\end{proof}
\olpseganchor{OLP-0464-B023}{end}

\olpseganchor{OLP-0464-B024}{start}
\begin{prob}
د \olref[nml][tab][sou]{thm:tableau-soundness} ثبوت بشپړ کړئ۔
\end{prob}
\olpseganchor{OLP-0464-B024}{end}

\olpseganchor{OLP-0464-B025}{start}
\begin{cor}
\ollabel{cor:entailment-soundness}
که \LR{$\Gamma \Proves !A$} وي، نو \LR{$\Gamma \Entails !A$} هم وي۔
\end{cor}
\olpseganchor{OLP-0464-B025}{end}

\olpseganchor{OLP-0464-B026}{start}
\begin{proof}
  که \LR{$\Gamma \Proves !A$} وي، نو د ځينو \LR{$!B_1$}, \dots,
  \LR{$!B_n \in \Gamma$} لپاره
  \LR{$\Delta = \{\sFmla{\False}{!A}[1], \sFmla{\True}{!B_1}[1],
  \dots, \sFmla{\True}{!B_n}[1]\}$} تړلے تابلو لري۔
  غواړو وښيو چې \LR{$\Gamma \Entails !A$}۔ سرچپه فرض کړئ؛ نو د
  کوم \LR{$\mModel{M}$} او \LR{$w$} لپاره \LR{$\mSat{M}{!B_i}[w]$} د هر
  \LR{$i=1$}, \dots,~\LR{$n$} په حق کښې وي، خو
  \LR{$\mSat/{M}{!A}[w]$}۔ فرض کړئ \LR{$f(1) = w$}؛ بيا \LR{$f$} په
  \LR{$\mModel{M}$} کښې د~\LR{$P(\Delta)$} تفسير دے، او \LR{$\mModel{M}$}
  د~\LR{$f$} په نسبت \LR{$\Delta$} پوره کوي۔ خو د
  \olref{thm:tableau-soundness} له مخې، \LR{$\Delta$} نۀ شي پوره
  کېدے، ځکه تړلے تابلو لري؛ دا تناقض دے۔ نو
  \LR{$\Gamma \Entails !A$}۔
\end{proof}
\olpseganchor{OLP-0464-B026}{end}

\olpseganchor{OLP-0464-B027}{start}
\begin{cor}
\ollabel{cor:weak-soundness}
که \LR{$\Proves !A$} وي، نو \LR{$!A$} په ټولو مدلونو کښې رښتينے دے۔
\end{cor}
\olpseganchor{OLP-0464-B027}{end}

\olpseganchor{OLP-0465-B004}{start}
\olfileid{nml}{tab}{mru}
\olpseganchor{OLP-0465-B004}{end}

\olpseganchor{OLP-0465-B005}{start}
\olsection{د لاسرسي د نورو اړيکو لپاره قاعدې}
\olpseganchor{OLP-0465-B005}{end}

\olpseganchor{OLP-0465-B006}{start}
د هغو منطقونو د څېړلو لپاره چې ځانګړې د لاسرسي اړيکې يې ټاکي،
په \olref{tab:more-rules} کښې نورې قاعدې ګورو۔ د تناظر او
اقليدسي قاعدې د يو غځېدلي مخکښ له مقدمې نه د هغه نېغ مخکښ ته
پايله وړي؛ د پيلنيو فرضونو له ګډ مخکښ څخه په جوړو شوې څانګه
کښې دغه نېغ مخکښ لا شته۔
\olpseganchor{OLP-0465-B006}{end}

\olpseganchor{OLP-0465-B007}{start}
\begin{table}
  \begin{center}
    \def\arraystretch{3}\def\fCenter{}
    \begin{tabular}{|c|c|}
    \hline
    
      \AxiomC{\sFmla{\True}{\Box !A}[\sigma]}
      \RightLabel{T\LR{$\Box$}}
      \UnaryInfC{\sFmla{\True}{!A}[\sigma]}
      \DisplayProof
    &
    
      \AxiomC{\sFmla{\False}{\Diamond !A}[\sigma]}
      \RightLabel{T\LR{$\Diamond$}}
      \UnaryInfC{\sFmla{\False}{!A}[\sigma]}
      \DisplayProof
    \\[1ex]
    \hline
    
      \AxiomC{\sFmla{\True}{\Box !A}[\sigma]}
      \RightLabel{D\LR{$\Box$}}
      \UnaryInfC{\sFmla{\True}{\Diamond!A}[\sigma]}
      \DisplayProof
    &
    
      \AxiomC{\sFmla{\False}{\Diamond !A}[\sigma]}
      \RightLabel{D\LR{$\Diamond$}}
      \UnaryInfC{\sFmla{\False}{\Box!A}[\sigma]}
      \DisplayProof
    \\[1ex]
    \hline
    
      \AxiomC{\sFmla{\True}{\Box !A}[\sigma.n]}
      \RightLabel{B\LR{$\Box$}}
      \UnaryInfC{\sFmla{\True}{!A}[\sigma]}
      \DisplayProof
    &
    
      \AxiomC{\sFmla{\False}{\Diamond !A}[\sigma.n]}
      \RightLabel{B\LR{$\Diamond$}}
      \UnaryInfC{\sFmla{\False}{!A}[\sigma]}
      \DisplayProof
    \\[1ex]
    \hline
    
      \AxiomC{\sFmla{\True}{\Box !A}[\sigma]}
      \RightLabel{4\LR{$\Box$}}
      \UnaryInfC{\sFmla{\True}{\Box!A}[\sigma.n]}
      \DisplayProof
    &
    
      \AxiomC{\sFmla{\False}{\Diamond !A}[\sigma]}
      \RightLabel{4\LR{$\Diamond$}}
      \UnaryInfC{\sFmla{\False}{\Diamond!A}[\sigma.n]}
      \DisplayProof
    \\
    \LR{$\sigma.n$} په څانګه کښې کارېدلے دے & \LR{$\sigma.n$} په څانګه کښې کارېدلے دے
    \\[1ex]
    \hline
    
      \AxiomC{\sFmla{\True}{\Box !A}[\sigma.n]}
      \RightLabel{4r\LR{$\Box$}}
      \UnaryInfC{\sFmla{\True}{\Box!A}[\sigma]}
      \DisplayProof
    &
    
      \AxiomC{\sFmla{\False}{\Diamond !A}[\sigma.n]}
      \RightLabel{4r\LR{$\Diamond$}}
      \UnaryInfC{\sFmla{\False}{\Diamond!A}[\sigma]}
      \DisplayProof
    \\[1ex]
    \hline
    \end{tabular}
  \end{center}
  \caption{نورې مودالي قاعدې۔}
  \ollabel{tab:more-rules}
\end{table}
\olpseganchor{OLP-0465-B007}{end}

\olpseganchor{OLP-0465-B008}{start}
د دغو قاعدو په زياتولو داسې نظامونه ترلاسه کېږي چې په
\olref{tab:logics-rules} کښې ښودل شويو منطقونو لپاره سم او بشپړ دي۔
\olpseganchor{OLP-0465-B008}{end}

\olpseganchor{OLP-0465-B009}{start}
\begin{table}
  \begin{center}
    \begin{tabular}{lll}
      \hline
      منطق & \LR{$R$} \dots دے & قاعدې\\
      \hline
      \LR{$\Log{T} = \Log{KT}$} & انعکاسي &
      T\LR{$\Box$}%
      , %
      T\LR{$\Diamond$}
      \\ \hline
      \LR{$\Log{D} = \Log{KD}$} & مسلسل &
      D\LR{$\Box$}%
      , %
      D\LR{$\Diamond$}
      \\ \hline
      \LR{$\Log{K4}$} & متعدي &
      4\LR{$\Box$}%
      , %
      4\LR{$\Diamond$}
      \\ \hline
      \LR{$\Log{B} = \Log{KTB}$} & انعکاسي، &
      T\LR{$\Box$}%
      , %
      T\LR{$\Diamond$}\\
      & متناظره &
      B\LR{$\Box$}%
      , %
      B\LR{$\Diamond$}
      \\ \hline
      \LR{$\Log{S4} = \Log{KT4}$} & انعکاسي، &
      T\LR{$\Box$}%
      , %
      T\LR{$\Diamond$},\\
      & متعدي &
      4\LR{$\Box$}%
      , %
      4\LR{$\Diamond$}
      \\ \hline
      \LR{$\Log{S5} = \Log{KT4B}$} & انعکاسي، &
      T\LR{$\Box$}%
      , %
      T\LR{$\Diamond$},\\
      & متعدي، &
      4\LR{$\Box$}%
      , %
      4\LR{$\Diamond$},\\
      &  اقليدسي &
      4r\LR{$\Box$}%
      , %
      4r\LR{$\Diamond$}
      \\ \hline
    \end{tabular}
  \end{center}
  \caption{د بېلابېلو مودالي منطقونو لپاره تابلو قاعدې۔}
  \ollabel{tab:logics-rules}
\end{table}
\olpseganchor{OLP-0465-B009}{end}


\olpseganchor{OLP-0465-B010}{start}
\begin{ex}
  داسې تړلے تابلو ورکوو چې
  \LR{$\Log{S5} \Proves \Box!A \lif \Box\Diamond!A$} ښيي۔
  د \Ax{5} طرحه \LR{$\Diamond!A \lif \Box\Diamond!A$} ده۔
  په لاندې تابلو کښې د مخکښ غځون د نارښتيني بکس شاهد جوړوي،
  او اقليدسي قاعده د هماغه شاهد د امکان نفي نېغ مخکښ ته راوړي۔
  \begin{oltableau}
    [\pFmla{\False}{\Box\formula{A} \lif \Box\Diamond \formula{A}}{1},
      just = \TAss
      [\pFmla{\True}{\Box \formula{A}}{1}, just = {\TRule{\False}{\lif}[1]}
        [\pFmla{\False}{\Box\Diamond \formula{A}}{1},
          just = {\TRule{\False}{\lif}[1]}
          [\pFmla{\False}{\Diamond \formula{A}}{1.1},
            just = {\TRule{\False}{\Box}[3]}
            [\pFmla{\False}{\Diamond \formula{A}}{1},
              just = {4r\LR{$\Diamond$} 4}
              [\pFmla{\False}{\formula{A}}{1.1},
                just = {\TRule{\False}{\Diamond}[5]}
                [\pFmla{\True}{\formula{A}}{1.1},
                  just = {\TRule{\True}{\Box}[2]}, close]
              ]
            ]
          ]
        ]
      ]
    ]
  \end{oltableau}
\end{ex}
\olpseganchor{OLP-0465-B010}{end}

\olpseganchor{OLP-0465-B011}{start}
\begin{prob}
د لاندې استنتاجونو د ښودلو لپاره تړلي تابلوګانې ورکړئ:
  \begin{enumerate}
  \item \LR{$\Log{KT5} \Proves \Ax{B}$};
  \item \LR{$\Log{KT5} \Proves \Ax{4}$};
  \item \LR{$\Log{KDB4} \Proves \Ax{T}$};
  \item \LR{$\Log{KB4} \Proves \Ax{5}$};
  \item \LR{$\Log{KB5} \Proves \Ax{4}$};
  \item \LR{$\Log{KT} \Proves \Ax{D}$}.
  \end{enumerate}
\end{prob}
\olpseganchor{OLP-0465-B011}{end}

\olpseganchor{OLP-0466-B004}{start}
\olfileid{nml}{tab}{msn}
\olpseganchor{OLP-0466-B004}{end}
      
\olpseganchor{OLP-0466-B005}{start}
\olsection{د نورو قاعدو سموالی}
\olpseganchor{OLP-0466-B005}{end}

\olpseganchor{OLP-0466-B006}{start}
وايو چې يوه قاعده د مدلونو د يوې ټولګې لپاره سمه ده که د
تابلو هره څانګه چې د هغې ټولګې په کوم مدل کښې پوره کېږي،
د قاعدې له پلي کېدو وروسته ترې جوړه شوې څانګه هم د هماغې
ټولګې په کوم مدل کښې پوره شي۔ لاندې قاعدې په څانګه کښې
شته مخکښونه کاروي، او هماغه مدل او تفسير د پايلو صدق ښيي۔
\olpseganchor{OLP-0466-B006}{end}

\olpseganchor{OLP-0466-B007}{start}
\begin{prop}
  \ollabel{prop:soundness-T} 
  \Ax{T}\LR{$\Box$}%
   او %
  \Ax{T}\LR{$\Diamond$}
  د انعکاسي مدلونو لپاره سمې قاعدې دي۔
\end{prop}
\olpseganchor{OLP-0466-B007}{end}

\olpseganchor{OLP-0466-B008}{start}
\begin{proof}
\begin{enumerate}
\item د T\LR{$\Box$} قاعدې په پلي کولو څانګه غځېږي؛ مقدمه
  \LR{$\sFmla{\True}{\Box !B}[\sigma] \in \Gamma$} ده:
  پر څانګه نوی نښه لرونکے فارمول~\LR{$\sFmla{\True}{!B}[\sigma]$} راځي۔ فرض کړئ
    \LR{$\mSat{M}{\Gamma}[f]$}؛ په ځانګړې توګه، \LR{$\mSat{M}{\Box
      !B}[f(\sigma)]$}. څرنګه چې \LR{$R$} انعکاسي ده، نو
    \LR{$Rf(\sigma)f(\sigma)$}.  نو \LR{$\mSat{M}{!B}[f(\sigma)]$}, يعنې
    \LR{$\mModel{M}, f$} د \LR{$\sFmla{\True}{!B}[\sigma]$} نښه لرونکے فارمول پوره کوي۔
\item د T\LR{$\Diamond$} قاعدې په پلي کولو څانګه غځېږي؛ مقدمه
  \LR{$\sFmla{\False}{\Diamond !B}[\sigma] \in \Gamma$} ده:
  پر څانګه نوی نښه لرونکے فارمول~\LR{$\sFmla{\False}{!B}[\sigma]$} راځي۔ فرض کړئ
    \LR{$\mSat{M}{\Gamma}[f]$}؛ په ځانګړې توګه، \LR{$\mSat/{M}{\Diamond
      !B}[f(\sigma)]$}. څرنګه چې \LR{$R$} انعکاسي ده، نو
    \LR{$Rf(\sigma)f(\sigma)$}.  نو \LR{$\mSat/{M}{!B}[f(\sigma)]$}, يعنې
    \LR{$\mModel{M}, f$} د \LR{$\sFmla{\False}{!B}[\sigma]$} نښه لرونکے فارمول پوره کوي۔
\end{enumerate}
\end{proof}
\olpseganchor{OLP-0466-B008}{end}



\olpseganchor{OLP-0466-B010}{start}
\begin{prop}
  \ollabel{prop:soundness-D} 
  \Ax{D}\LR{$\Box$}%
   او %
  \Ax{D}\LR{$\Diamond$}
  د مسلسل مدلونو لپاره سمې قاعدې دي۔
\end{prop}
\olpseganchor{OLP-0466-B010}{end}

\olpseganchor{OLP-0466-B011}{start}
\begin{proof}
\begin{enumerate}
\item د D\LR{$\Box$} قاعدې په پلي کولو څانګه غځېږي؛ مقدمه
  \LR{$\sFmla{\True}{\Box !B}[\sigma] \in \Gamma$} ده:
  پر څانګه نوی نښه لرونکے فارمول~\LR{$\sFmla{\True}{\Diamond !B}[\sigma]$} راځي۔ فرض کړئ \LR{$\mSat{M}{\Gamma}[f]$}؛ په ځانګړې توګه،
    \LR{$\mSat{M}{\Box !B}[f(\sigma)]$}. څرنګه چې \LR{$R$} مسلسل ده، داسې \LR{$w
    \in W$} شته چې \LR{$Rf(\sigma)w$}. بيا \LR{$\mSat{M}{!B}[w]$}, او له دې امله
    \LR{$\mSat{M}{\Diamond !B}[f(\sigma)]$}.  نو \LR{$\mModel{M}, f$} د \LR{$\sFmla{\True}{\Diamond!B}[\sigma]$} نښه لرونکے فارمول پوره کوي۔
\item د D\LR{$\Diamond$} قاعدې په پلي کولو څانګه غځېږي؛ مقدمه
  \LR{$\sFmla{\False}{\Diamond !B}[\sigma] \in \Gamma$} ده:
  پر څانګه نوی نښه لرونکے فارمول~\LR{$\sFmla{\False}{\Box !B}[\sigma]$} راځي۔ فرض کړئ \LR{$\mSat{M}{\Gamma}[f]$}؛ په ځانګړې توګه،
    \LR{$\mSat/{M}{\Diamond !B}[f(\sigma)]$}. څرنګه چې \LR{$R$} مسلسل ده، داسې \LR{$w
    \in W$} شته چې \LR{$Rf(\sigma)w$}. بيا \LR{$\mSat/{M}{!B}[w]$}, او له دې امله
    \LR{$\mSat/{M}{\Box !B}[f(\sigma)]$}.  نو \LR{$\mModel{M}, f$} د \LR{$\sFmla{\False}{\Box!B}[\sigma]$} نښه لرونکے فارمول پوره کوي۔
\end{enumerate}
\end{proof}
\olpseganchor{OLP-0466-B011}{end}



\olpseganchor{OLP-0466-B013}{start}
\begin{prop}
  \ollabel{prop:soundness-B} 
  \Ax{B}\LR{$\Box$}%
   او %
  \Ax{B}\LR{$\Diamond$}
  د متناظرو مدلونو لپاره سمې قاعدې دي۔
\end{prop}
\olpseganchor{OLP-0466-B013}{end}

\olpseganchor{OLP-0466-B014}{start}
\begin{proof}
\begin{enumerate}
\item د B\LR{$\Box$} قاعدې په پلي کولو څانګه غځېږي؛ مقدمه
  \LR{$\sFmla{\True}{\Box !B}[\sigma.n] \in \Gamma$} ده:
  پر څانګه نوی نښه لرونکے فارمول~\LR{$\sFmla{\True}{!B}[\sigma]$} راځي۔ فرض کړئ
    \LR{$\mSat{M}{\Gamma}[f]$}؛ په ځانګړې توګه، \LR{$\mSat{M}{\Box
      !B}[f(\sigma.n)]$}. څرنګه چې \LR{$f$} په~\LR{$\mModel{M}$} کښې د څانګې د مخکښونو تفسير دے، نو
    \LR{$Rf(\sigma)f(\sigma.n)$}.  څرنګه چې \LR{$R$} متناظره ده، نو
    \LR{$Rf(\sigma.n)f(\sigma)$}.  او ځکه چې \LR{$\mSat{M}{\Box !B}[f(\sigma.n)]$} وي،
    \LR{$\mSat{M}{!B}[f(\sigma)]$}. نو \LR{$\mModel{M}, f$} د \LR{$\sFmla{\True}{!B}[\sigma]$} نښه لرونکے فارمول پوره کوي۔
\item د B\LR{$\Diamond$} قاعدې په پلي کولو څانګه غځېږي؛ مقدمه
  \LR{$\sFmla{\False}{\Diamond !B}[\sigma.n] \in \Gamma$} ده:
  پر څانګه نوی نښه لرونکے فارمول~\LR{$\sFmla{\False}{!B}[\sigma]$} راځي۔ فرض کړئ
    \LR{$\mSat{M}{\Gamma}[f]$}؛ په ځانګړې توګه، \LR{$\mSat/{M}{\Diamond
      !B}[f(\sigma.n)]$}. څرنګه چې \LR{$f$} په~\LR{$\mModel{M}$} کښې د څانګې د مخکښونو تفسير دے، نو
    \LR{$Rf(\sigma)f(\sigma.n)$}.  څرنګه چې \LR{$R$} متناظره ده، نو
    \LR{$Rf(\sigma.n)f(\sigma)$}.  او ځکه چې \LR{$\mSat/{M}{\Diamond !B}[f(\sigma.n)]$} وي،
    \LR{$\mSat/{M}{!B}[f(\sigma)]$}. نو \LR{$\mModel{M}, f$} د \LR{$\sFmla{\False}{!B}[\sigma]$} نښه لرونکے فارمول پوره کوي۔
\end{enumerate}
\end{proof}
\olpseganchor{OLP-0466-B014}{end}



\olpseganchor{OLP-0466-B016}{start}
\begin{prop}
  \ollabel{prop:soundness-4} 
  \Ax{4}\LR{$\Box$}%
   او %
  \Ax{4}\LR{$\Diamond$}
  د متعديو مدلونو لپاره سمې قاعدې دي۔
\end{prop}
\olpseganchor{OLP-0466-B016}{end}

\olpseganchor{OLP-0466-B017}{start}
\begin{proof}
\begin{enumerate}
\item د 4\LR{$\Box$} قاعدې په پلي کولو څانګه غځېږي؛ مقدمه
  \LR{$\sFmla{\True}{\Box !B}[\sigma] \in \Gamma$} ده:
  پر څانګه نوی نښه لرونکے فارمول~\LR{$\sFmla{\True}{\Box!B}[\sigma.n]$} راځي۔ فرض کړئ \LR{$\mSat{M}{\Gamma}[f]$}؛ په ځانګړې توګه،
    \LR{$\mSat{M}{\Box !B}[f(\sigma)]$}. څرنګه چې \LR{$f$} په~\LR{$\mModel{M}$} کښې د څانګې د مخکښونو تفسير دے او
    \LR{$\sigma.n$} بايد په څانګه کښې کارېدلے وي، نو \LR{$Rf(\sigma)f(\sigma.n)$}. اوس هر داسې \LR{$w$} ونيسئ چې \LR{$Rf(\sigma.n)w$}.  څرنګه چې \LR{$R$} متعديه ده، نو
    \LR{$Rf(\sigma)w$}.  ځکه چې \LR{$\mSat{M}{\Box !B}[f(\sigma)]$} وي، نو
    \LR{$\mSat{M}{!B}[w]$}. نو \LR{$\mSat{M}{\Box !B}[f(\sigma.n)]$}، او
    \LR{$\mModel{M}, f$} د \LR{$\sFmla{\True}{\Box !B}[\sigma.n]$} نښه لرونکے فارمول پوره کوي۔
\item د 4\LR{$\Diamond$} قاعدې په پلي کولو څانګه غځېږي؛ مقدمه
  \LR{$\sFmla{\False}{\Diamond !B}[\sigma] \in \Gamma$} ده:
  پر څانګه نوی نښه لرونکے فارمول~\LR{$\sFmla{\False}{\Diamond!B}[\sigma.n]$} راځي۔ فرض کړئ \LR{$\mSat{M}{\Gamma}[f]$}؛ په ځانګړې توګه،
    \LR{$\mSat/{M}{\Diamond !B}[f(\sigma)]$}. څرنګه چې \LR{$f$} په~\LR{$\mModel{M}$} کښې د څانګې د مخکښونو تفسير دے او
    \LR{$\sigma.n$} بايد په څانګه کښې کارېدلے وي، نو \LR{$Rf(\sigma)f(\sigma.n)$}. اوس هر داسې \LR{$w$} ونيسئ چې \LR{$Rf(\sigma.n)w$}.  څرنګه چې \LR{$R$} متعديه ده، نو
    \LR{$Rf(\sigma)w$}.  ځکه چې \LR{$\mSat/{M}{\Diamond !B}[f(\sigma)]$} وي، نو
    \LR{$\mSat/{M}{!B}[w]$}. نو \LR{$\mSat/{M}{\Diamond !B}[f(\sigma.n)]$}، او
    \LR{$\mModel{M}, f$} د \LR{$\sFmla{\False}{\Diamond !B}[\sigma.n]$} نښه لرونکے فارمول پوره کوي۔
\end{enumerate}
\end{proof}
\olpseganchor{OLP-0466-B017}{end}



\olpseganchor{OLP-0466-B019}{start}
\begin{prop}
  \ollabel{prop:soundness-4r} 
  \Ax{4r}\LR{$\Box$}%
   او %
  \Ax{4r}\LR{$\Diamond$}
  د اقليدسي مدلونو لپاره سمې قاعدې دي۔
\end{prop}
\olpseganchor{OLP-0466-B019}{end}

\olpseganchor{OLP-0466-B020}{start}
\begin{proof}
\begin{enumerate}
\item د 4r\LR{$\Box$} قاعدې په پلي کولو څانګه غځېږي؛ مقدمه
  \LR{$\sFmla{\True}{\Box !B}[\sigma.n] \in \Gamma$} ده:
  پر څانګه نوی نښه لرونکے فارمول~\LR{$\sFmla{\True}{\Box!B}[\sigma]$} راځي۔ فرض کړئ
    \LR{$\mSat{M}{\Gamma}[f]$}؛ په ځانګړې توګه، \LR{$\mSat{M}{\Box
      !B}[f(\sigma.n)]$}. څرنګه چې \LR{$f$} په~\LR{$\mModel{M}$} کښې د څانګې د مخکښونو تفسير دے، نو
    \LR{$Rf(\sigma)f(\sigma.n)$}. اوس هر داسې \LR{$w$} ونيسئ چې
    \LR{$Rf(\sigma)w$}.  څرنګه چې \LR{$R$} اقليدسي ده، نو \LR{$Rf(\sigma.n)w$}.  ځکه چې
    \LR{$\mSat{M}{\Box !B}[f(\sigma.n)]$} وي، نو \LR{$\mSat{M}{!B}[w]$}۔ نو
    \LR{$\mSat{M}{\Box !B}[f(\sigma)]$}، او \LR{$\mModel{M}, f$} د \LR{$\sFmla{\True}{\Box !B}[\sigma]$} نښه لرونکے فارمول پوره کوي۔
\item د 4r\LR{$\Diamond$} قاعدې په پلي کولو څانګه غځېږي؛ مقدمه
  \LR{$\sFmla{\False}{\Diamond !B}[\sigma.n] \in \Gamma$} ده:
  پر څانګه نوی نښه لرونکے فارمول~\LR{$\sFmla{\False}{\Diamond!B}[\sigma]$} راځي۔ فرض کړئ
    \LR{$\mSat{M}{\Gamma}[f]$}؛ په ځانګړې توګه، \LR{$\mSat/{M}{\Diamond
      !B}[f(\sigma.n)]$}. څرنګه چې \LR{$f$} په~\LR{$\mModel{M}$} کښې د څانګې د مخکښونو تفسير دے، نو
    \LR{$Rf(\sigma)f(\sigma.n)$}. اوس هر داسې \LR{$w$} ونيسئ چې
    \LR{$Rf(\sigma)w$}.  څرنګه چې \LR{$R$} اقليدسي ده، نو \LR{$Rf(\sigma.n)w$}.  ځکه چې
    \LR{$\mSat/{M}{\Diamond !B}[f(\sigma.n)]$} وي، نو \LR{$\mSat/{M}{!B}[w]$}۔ نو
    \LR{$\mSat/{M}{\Diamond !B}[f(\sigma)]$}، او \LR{$\mModel{M}, f$} د \LR{$\sFmla{\False}{\Diamond !B}[\sigma]$} نښه لرونکے فارمول پوره کوي۔
\end{enumerate}
\end{proof}
\olpseganchor{OLP-0466-B020}{end}



\olpseganchor{OLP-0466-B022}{start}
\begin{cor}
\ollabel{cor:soundness-logics} په \olref[mru]{tab:logics-rules} کښې ښودل شوي تابلو نظامونه
د مدلونو د خپلو اړوندو ټولګيو لپاره سم دي۔
\end{cor}
\olpseganchor{OLP-0466-B022}{end}

\olpseganchor{OLP-0467-B004}{start}
\olfileid{nml}{tab}{s5}
\olpseganchor{OLP-0467-B004}{end}

\olpseganchor{OLP-0467-B005}{start}
\olsection{د \texorpdfstring{\Log{S5}}{S5} لپاره ساده تابلوګانې}
\olpseganchor{OLP-0467-B005}{end}

\olpseganchor{OLP-0467-B006}{start}
\Log{S5} د عمومي مدلونو د ټولګي په نسبت سم او بشپړ دے؛
يعنې داسې مدلونه چې هره نړۍ ئې له هرې نړۍ څخه د لاسرسي وړ وي۔
په دغو عمومي مدلونو کښې د لاسرسي اړيکه توپير نۀ راولي:
«داسې نړۍ~\LR{$w$} شته چې \LR{$\mSat{M}{!A}[w]$} وي» هله او يوازې هله
رښتيا ده چې داسې يو~\LR{$w$} له~\LR{$u$} څخه د لاسرسي وړ وي۔
نو د \Log{S5} د دې معناپوهنې عمومي مدل د نړۍو په يوه ناتش
سټ او د ارزښت‌ټاکنې په تابع~\LR{$V$} ټاکل کېږي؛ د لاسرسي اړيکه
د نړۍو ټولې جوړې لري۔ دا ښيي چې د تابلو قاعدې هم
ساده کولے شو۔ په عمومي حالت کښې د مثبتو صحيح عددونو لړۍ
د مخکښونو په توګه اخلو، څو معلومه کړو چې کوم مخکښونه له
نورو نړۍو څخه د لاسرسي وړ نړۍې نوموي: \LR{$\sigma.n$} له~\LR{$\sigma$}
څخه د لاسرسي وړ نړۍ نوموي۔ خو د \Log{S5} د دغو عمومي
مدلونو هره نړۍ له هرې نړۍ څخه د لاسرسي وړ ده؛ نو د دې
څارنې اړتيا نۀ شته۔ د دې پر ځای مثبت صحيح عددونه د
مخکښونو په توګه کارولے شو۔ ساده شوې قاعدې په
\olref{tab:rules-S5} کښې راغلې دي۔
\olpseganchor{OLP-0467-B006}{end}

\olpseganchor{OLP-0467-B007}{start}
\begin{table}
  \begin{center}
    \def\arraystretch{3}\def\fCenter{}
    \begin{tabular}{|c|c|}
    \hline
    
      \AxiomC{\sFmla{\True}{\Box !A}[n]}
      \RightLabel{\TRule{\True}{\Box}}
      \UnaryInfC{\sFmla{\True}{!A}[m]}
      \DisplayProof
      &
      \AxiomC{\sFmla{\False}{\Box !A}[n]}
      \RightLabel{\TRule{\False}{\Box}}
      \UnaryInfC{\sFmla{\False}{!A}[m]}
      \DisplayProof\\
      \LR{$m$} په څانګه کښې کارېدلے دے & \LR{$m$} په څانګه کښې نوے دے
      \\[1ex]
      \hline
    
      \AxiomC{\sFmla{\True}{\Diamond !A}[n]}
      \RightLabel{\TRule{\True}{\Diamond}}
      \UnaryInfC{\sFmla{\True}{!A}[m]}
      \DisplayProof
      &
      \AxiomC{\sFmla{\False}{\Diamond !A}[n]}
      \RightLabel{\TRule{\False}{\Diamond}}
      \UnaryInfC{\sFmla{\False}{!A}[m]}
      \DisplayProof\\
      \LR{$m$} په څانګه کښې نوے دے & \LR{$m$} په څانګه کښې کارېدلے دے
      \\[1ex]
      \hline
    \end{tabular}
  \end{center}
  \caption{د \Log{S5} لپاره ساده شوې قاعدې۔}
  \ollabel{tab:rules-S5}
\end{table}
\olpseganchor{OLP-0467-B007}{end}

\olpseganchor{OLP-0467-B008}{start}
\begin{ex}
  داسې ساده تړلے تابلو ورکوو چې \LR{$\Log{S5} \Proves \Ax{5}$}،
  يعنې \LR{$\Diamond!A \lif \Box\Diamond!A$}، ثابتوي۔ د نارښتيني
  بکس او رښتيني امکان قاعدې جلا نوي مخکښونه جوړوي؛ په عمومي
  مدل کښې د نارښتيني امکان قاعده له خپل مخکښ څخه هر کارېدلے
  مخکښ اخيستلے شي، نو څانګه تړي۔
  \begin{oltableau}
    [\pFmla{\False}{\Diamond\formula{A} \lif \Box\Diamond \formula{A}}{1},
      just = \TAss
      [\pFmla{\True}{\Diamond \formula{A}}{1}, just = {\TRule{\False}{\lif}[1]}
        [\pFmla{\False}{\Box\Diamond \formula{A}}{1},
          just = {\TRule{\False}{\lif}[1]}
          [\pFmla{\False}{\Diamond \formula{A}}{2},
            just = {\TRule{\False}{\Box}[3]}
            [\pFmla{\True}{\formula{A}}{3},
              just = {\TRule{\True}{\Diamond}[2]}
                [\pFmla{\False}{\formula{A}}{3},
                  just = {\TRule{\False}{\Diamond}[4]}, close]
            ]
          ]
        ]
      ]
    ]
  \end{oltableau}
\end{ex}
\olpseganchor{OLP-0467-B008}{end}

\olpseganchor{OLP-0468-B004}{start}
\olfileid{nml}{tab}{cpl}
\olpseganchor{OLP-0468-B004}{end}

\olpseganchor{OLP-0468-B005}{start}
\olsection{د \texorpdfstring{\Log{K}}{K} بشپړتيا}
\olpseganchor{OLP-0468-B005}{end}

\olpseganchor{OLP-0468-B006}{start}
\begin{explain}
  د تابلوګانې د طريقې د بشپړتيا د ښودلو لپاره بايد ثابته کړو
  چې که د \LR{$\Gamma \Proves !A$} د ښودلو لپاره تړلے تابلو
  نۀ وي، نو \LR{$\Gamma \Entails/ !A$} وي، يعنې ضدمدل موجود وي۔
  خو «تړلے تابلو نۀ شته» معنا دا ده چې د هغه د جوړولو
  هره ممکنه هڅه په تړلو کښې پاتې راځي۔ مهمه خبره دا ده چې که
  هره هڅه داسې پاتې راځي، نو يوه ټاکلې، \emph{منظمه او هر اړخيزه}
  طريقه هم ټولې څانګې نۀ شي تړلے۔ او که تړلے تابلو موجود وي،
  همدغه منظمه او هر اړخيزه طريقه به ټولې څانګې وتړي۔ د دې يوې تابلو
  پرانيستې څانګې به د ضدمدل د تعريف لپاره ټول اړين معلومات ولري۔
  د داسې څانګې ضدمدل هغه ټول مخکښونه چې پر څانګه کارېدلي،
  د نړۍو په توګه اخلي؛ بياني متغير~\LR{$p$} په
  \LR{$\sigma$} کښې هله او يوازې هله رښتينے وي چې
  \LR{$\sFmla{\True}{p}[\sigma]$} په څانګه کښې راغلے وي۔
  په دې کړنلاره کښې دوه‌اړخيز شرط د دواړو لورو د شرطي فارمولونو
  په عطف بدلوو۔ د يوه ثابت قضيېيز متغير او د هغه د نفي فصل د
  رښتيا ثابت ښيي، او د دواړو عطف د دروغو ثابت ښيي۔ دا بدلونه
  په هره نړۍ کښې رښتياوالے ساتي، او پر فرعي فارمولونو په
  استقرا سره ئې ټول فارمول ته غځوو۔
\end{explain}
\olpseganchor{OLP-0468-B006}{end}

\olpseganchor{OLP-0468-B007}{start}
\begin{defn}
  د تابلو يوه څانګه هغه مهال بشپړه بولو چې هرکله پکښې
  داسې مخکښ‌لرونکے فارمول~\LR{$\sFmla{S}{!A}[\sigma]$} وي چې
  قاعده پرې پلې کېدے شي، نو څانګه دا هم ولري:
  \begin{enumerate}
    \item د قضيېيزو يو-پر-بل اېښودونکو قاعدو په حالت کښې د
      قاعدې ټول اړوند پايلې، يعنې مخکښ‌لرونکي فارمولونه؛
    \item د قضيېيزو څانګې وېشونکو قاعدو په حالت کښې له
      اړوندو پايلنو فارمولونه څخه لږ تر لږه يو؛
    \item د هغو مودالي قاعدو په حالت کښې چې نوے مخکښ غواړي،
      لږ تر لږه يوه ممکنه پايله؛
    \item د هغو مودالي قاعدو په حالت کښې چې کارېدلے مخکښ
      غواړي، د څانګې د هر پلي کېدونکي کارېدلي مخکښ لپاره
      اړونده پايله۔
  \end{enumerate}
\end{defn}
\olpseganchor{OLP-0468-B007}{end}

\olpseganchor{OLP-0468-B008}{start}
\begin{explain}
د بېلګې په توګه، که بشپړه څانګه \LR{$\sFmla{\True}{!B
  \land !C}[\sigma]$} ولري، نو \LR{$\sFmla{\True}{!B}[\sigma]$}
او \LR{$\sFmla{\True}{!C}[\sigma]$} هم لري۔ که
\LR{$\sFmla{\True}{!B \lor !C}[\sigma]$} ولري، نو لږ تر لږه
\LR{$\sFmla{\True}{!B}[\sigma]$} يا \LR{$\sFmla{\True}{!C}[\sigma]$}
لري۔
که څانګه \LR{$\sFmla{\False}{\Box !A}[\sigma]$}
ولري، نو د لږ تر لږه يوه~\LR{$n$} لپاره
\LR{$\sFmla{\False}{!A}[\sigma.n]$} هم لري۔ او که څانګه
\LR{$\sFmla{\True}{\Box !A}[\sigma]$} ولري، نو د هر هغه~\LR{$n$}
لپاره چې \LR{$\sigma.n$} په څانګه کښې کارېدلے وي،
\LR{$\sFmla{\True}{!A}[\sigma.n]$} هم لري۔
\end{explain}
\olpseganchor{OLP-0468-B008}{end}

\olpseganchor{OLP-0468-B009}{start}
\begin{prop}\ollabel{prop:complete-tableau}
  هر متناهي~\LR{$\Gamma$} داسې تابلو لري چې هره پرانيستې څانګه ئې بشپړه وي۔
\end{prop}
\olpseganchor{OLP-0468-B009}{end}

\olpseganchor{OLP-0468-B010}{start}
\begin{proof}
  د~\LR{$\Gamma$} د تابلو يوه پرانيستې څانګه واخلئ۔ پر دې
  څانګه يوازې متناهي ډېر داسې مخکښ‌لرونکي فارمولونه دي چې
  قاعده پرې پلې کېدے شي۔ په يوه ټاکلي ترتيب (مثلاً له پاسه
  تر ښکته) د دغو فارمولونه له هر يوه لپاره چې پورته (۱)–(۴) شرطونه
  يې لا نۀ وي پوره کړي، پلې کېدونکې قاعدې وکاروئ او څانګه
  وغځوئ۔ کله کله څانګه وېشل کېږي؛ د هرې نوې څانګې په سر کښې
  پر پاتې مخکښ‌لرونکو فارمولونه همدغه قاعدې پلې کړئ۔
  ځکه چې پر څانګه د مخکښ‌لرونکو فارمولونه او کارېدلو
  مخکښونو شمېر متناهي دے، دا کړنلاره په پايله کښې اصلي
  څانګه په (کېدے شي ګڼو) څانګو غځوي۔ په هره پرانيستې څانګه
  کښې همدا کړنلاره تکراروو۔ د هرې مقدمې لپاره د نوي شاهد قاعده
  يوازې يو ځل پکار ده؛ د کارېدلي مخکښ قاعده د هرې مقدمې او
  روا مخکښ د جوړې لپاره يوازې يو ځل پکار ده۔ ټول پيدا شوي
  فارمولونه د پيلنيو فارمولونو فرعي فارمولونه دي۔ د مخکښ د
  اوږدوالي او د فارمول د مودالي ژورتيا مجموعه د هرې قاعدې
  په پلي کېدو کښې نۀ زياتېږي۔ نو د مخکښ اوږدوالے محدود دے۔
  پر هر مخکښ د نويو شاهدانو شمېر هم محدود دے، ځکه د هرې
  نښه لرونکې مقدمې يو شاهد بس دے۔ په هره څانګه کښې د مخکښونو
  او فارمولونو محدود شمېر د قاعدو د پلي کولو شمېر محدودوي۔
  نو کړنلاره درېږي، او هره پاتې پرانيستې څانګه بشپړه وي۔
\end{proof}
\olpseganchor{OLP-0468-B010}{end}

\olpseganchor{OLP-0468-B011}{start}
\begin{thm}[بشپړتيا]
  \ollabel{thm:tableau-completeness}
  که متناهي سټ~\LR{$\Gamma$} تړلے تابلو ونۀ لري، نو~\LR{$\Gamma$} پوره کېدونکے دے۔
\end{thm}
\olpseganchor{OLP-0468-B011}{end}

\olpseganchor{OLP-0468-B012}{start}
\begin{proof}
که پيلنے سټ تش وي، يوې نړۍ والا مدل بس دے۔ د ناتش سټ لپاره
د پورته قضيې له مخې~\LR{$\Gamma$} داسې تابلو لري چې هره
پرانيستې څانګه ئې بشپړه ده۔ ځکه تړلے تابلو نۀ لري، نو په دغه
تابلو کښې لږ تر لږه يوه پرانيستې او بشپړه څانګه شته۔
\LR{$\Delta$} دې پر څانګه د مخکښ‌لرونکو فارمولونه سټ وي، او
\LR{$P(\Delta)$} دې پکښې د راغلو مخکښونو سټ وي۔
\olpseganchor{OLP-0468-B012}{end}

\olpseganchor{OLP-0468-B013}{start}
مدل~\LR{$\mModel{M(\Delta)} = \tuple{P(\Delta), R, V}$} داسې تعريفوو
چې نړۍې ئې په~\LR{$\Delta$} کښې راغلي مخکښونه دي، او د لاسرسي
اړيکه يې داسې ده:
\[
R\sigma\sigma' \quad \text{\RL{هله او يوازې هله}} \quad
\sigma'=\sigma.n \quad \text{\RL{د کوم~\LR{$n$} لپاره}}
\]
او
\[
V(p) = \Setabs{\sigma}{\sFmla{\True}{p}[\sigma] \in \Delta}.
\]
پر~\LR{$!A$} د استقرا له لارې ښيو چې که
\LR{$\sFmla{\True}{!A}[\sigma] \in \Delta$} وي، نو
\LR{$\mSat{M(\Delta)}{!A}[\sigma]$} وي؛ او که
\LR{$\sFmla{\False}{!A}[\sigma] \in \Delta$} وي، نو
\LR{$\mSat/{M(\Delta)}{!A}[\sigma]$} وي۔
\olpseganchor{OLP-0468-B013}{end}

\olpseganchor{OLP-0468-B014}{start}
\begin{enumerate}
  \item \indcase{!A}{p}{که
    \LR{$\sFmla{\True}{\indfrm}[\sigma] \in \Delta$} وي، نو د~\LR{$V$}
    د تعريف له مخې \LR{$\sigma \in V(p)$} وي؛ له دې امله
    \LR{$\mSat{M(\Delta)}{\indfrm}[\sigma]$} وي۔
\olpseganchor{OLP-0468-B014}{end}

\olpseganchor{OLP-0468-B015}{start}
    که \LR{$\sFmla{\False}{\indfrm}[\sigma] \in \Delta$} وي، نو
    \LR{$\sFmla{\True}{\indfrm}[\sigma] \notin \Delta$} وي؛ ځکه
    که داسې نۀ وي، څانګه به تړلې وي۔ نو \LR{$\sigma \notin V(p)$}
    او له دې امله \LR{$\mSat/{M(\Delta)}{\indfrm}[\sigma]$} وي۔}
  \item \indcase{!A}{\lnot !B}{که
      \LR{$\sFmla{\True}{\indfrm}[\sigma] \in \Delta$} وي، نو د
      څانګې د بشپړتيا له مخې
      \LR{$\sFmla{\False}{!B}[\sigma] \in \Delta$} وي۔ د استقرا د
      فرضيې له مخې \LR{$\mSat/{M(\Delta)}{!B}[\sigma]$} وي، نو
      \LR{$\mSat{M(\Delta)}{\indfrm}[\sigma]$} وي۔
\olpseganchor{OLP-0468-B015}{end}

\olpseganchor{OLP-0468-B016}{start}
      که \LR{$\sFmla{\False}{\indfrm}[\sigma] \in \Delta$} وي، نو د
      څانګې د بشپړتيا له مخې
      \LR{$\sFmla{\True}{!B}[\sigma] \in \Delta$} وي۔ د استقرا د
      فرضيې له مخې \LR{$\mSat{M(\Delta)}{!B}[\sigma]$} وي، نو
      \LR{$\mSat/{M(\Delta)}{\indfrm}[\sigma]$} وي۔}
  \item \indcase{!A}{!B \land !C}{که
      \LR{$\sFmla{\True}{\indfrm}[\sigma] \in \Delta$} وي، نو د
      څانګې د بشپړتيا له مخې دواړه
      \LR{$\sFmla{\True}{!B}[\sigma] \in \Delta$} او
      \LR{$\sFmla{\True}{!C}[\sigma] \in \Delta$} وي۔ د استقرا د
      فرضيې له مخې \LR{$\mSat{M(\Delta)}{!B}[\sigma]$} او
      \LR{$\mSat{M(\Delta)}{!C}[\sigma]$} وي؛ نو
      \LR{$\mSat{M(\Delta)}{\indfrm}[\sigma]$} وي۔
\olpseganchor{OLP-0468-B016}{end}

\olpseganchor{OLP-0468-B017}{start}
      که \LR{$\sFmla{\False}{\indfrm}[\sigma] \in \Delta$} وي، نو د
      څانګې د بشپړتيا له مخې يا
      \LR{$\sFmla{\False}{!B}[\sigma] \in \Delta$} يا
      \LR{$\sFmla{\False}{!C}[\sigma] \in \Delta$} وي۔ د استقرا د
      فرضيې له مخې يا \LR{$\mSat/{M(\Delta)}{!B}[\sigma]$} يا
      \LR{$\mSat/{M(\Delta)}{!C}[\sigma]$} وي؛ نو
      \LR{$\mSat/{M(\Delta)}{\indfrm}[\sigma]$} وي۔}
  \item \indcase{!A}{!B \lor !C}{که
      \LR{$\sFmla{\True}{\indfrm}[\sigma] \in \Delta$} وي، نو د
      څانګې د بشپړتيا له مخې يا
      \LR{$\sFmla{\True}{!B}[\sigma] \in \Delta$} يا
      \LR{$\sFmla{\True}{!C}[\sigma] \in \Delta$} وي۔ د استقرا د
      فرضيې له مخې يا \LR{$\mSat{M(\Delta)}{!B}[\sigma]$} يا
      \LR{$\mSat{M(\Delta)}{!C}[\sigma]$} وي؛ نو
      \LR{$\mSat{M(\Delta)}{\indfrm}[\sigma]$} وي۔
\olpseganchor{OLP-0468-B017}{end}

\olpseganchor{OLP-0468-B018}{start}
      که \LR{$\sFmla{\False}{\indfrm}[\sigma] \in \Delta$} وي، نو د
      څانګې د بشپړتيا له مخې دواړه
      \LR{$\sFmla{\False}{!B}[\sigma] \in \Delta$} او
      \LR{$\sFmla{\False}{!C}[\sigma] \in \Delta$} وي۔ د استقرا د
      فرضيې له مخې \LR{$\mSat/{M(\Delta)}{!B}[\sigma]$} او
      \LR{$\mSat/{M(\Delta)}{!C}[\sigma]$} وي؛ نو
      \LR{$\mSat/{M(\Delta)}{\indfrm}[\sigma]$} وي۔}
  \item \indcase{!A}{!B \lif !C}{که
      \LR{$\sFmla{\True}{\indfrm}[\sigma] \in \Delta$} وي، نو د
      څانګې د بشپړتيا له مخې يا
      \LR{$\sFmla{\False}{!B}[\sigma] \in \Delta$} يا
      \LR{$\sFmla{\True}{!C}[\sigma] \in \Delta$} وي۔ د استقرا د
      فرضيې له مخې يا \LR{$\mSat/{M(\Delta)}{!B}[\sigma]$} يا
      \LR{$\mSat{M(\Delta)}{!C}[\sigma]$} وي؛ نو
      \LR{$\mSat{M(\Delta)}{\indfrm}[\sigma]$} وي۔
\olpseganchor{OLP-0468-B018}{end}

\olpseganchor{OLP-0468-B019}{start}
      که \LR{$\sFmla{\False}{\indfrm}[\sigma] \in \Delta$} وي، نو د
      څانګې د بشپړتيا له مخې دواړه
      \LR{$\sFmla{\True}{!B}[\sigma] \in \Delta$} او
      \LR{$\sFmla{\False}{!C}[\sigma] \in \Delta$} وي۔ د استقرا د
      فرضيې له مخې \LR{$\mSat{M(\Delta)}{!B}[\sigma]$} او
      \LR{$\mSat/{M(\Delta)}{!C}[\sigma]$} وي؛ نو
      \LR{$\mSat/{M(\Delta)}{\indfrm}[\sigma]$} وي۔}
  \item \indcase{!A}{\Box !B}{که
      \LR{$\sFmla{\True}{\indfrm}[\sigma] \in \Delta$} وي، نو د
      څانګې د بشپړتيا له مخې د هر کارېدلي~\LR{$\sigma.n$} لپاره
      \LR{$\sFmla{\True}{!B}[\sigma.n] \in \Delta$} وي؛ يعنې د هر
      هغه~\LR{$\sigma' \in P(\Delta)$} لپاره چې \LR{$R\sigma\sigma'$} وي۔
      د استقرا د فرضيې له مخې د هر هغه~\LR{$\sigma'$} لپاره چې
      \LR{$R\sigma\sigma'$} وي،
      \LR{$\mSat{M(\Delta)}{!B}[\sigma']$} وي؛ نو
      \LR{$\mSat{M(\Delta)}{\indfrm}[\sigma]$} وي۔
\olpseganchor{OLP-0468-B019}{end}

\olpseganchor{OLP-0468-B020}{start}
      که \LR{$\sFmla{\False}{\indfrm}[\sigma] \in \Delta$} وي، نو د
      څانګې د بشپړتيا له مخې د کوم~\LR{$\sigma.n$} لپاره
      \LR{$\sFmla{\False}{!B}[\sigma.n] \in \Delta$} وي۔ د استقرا د
      فرضيې له مخې \LR{$\mSat/{M(\Delta)}{!B}[\sigma.n]$} وي۔
      ځکه چې \LR{$R\sigma(\sigma.n)$}، داسې~\LR{$\sigma'$} شته چې
      \LR{$\mSat/{M(\Delta)}{!B}[\sigma']$} وي؛ نو
      \LR{$\mSat/{M(\Delta)}{\indfrm}[\sigma]$} وي۔}
  \item \indcase{!A}{\Diamond !B}{که
      \LR{$\sFmla{\True}{\indfrm}[\sigma] \in \Delta$} وي، نو د
      څانګې د بشپړتيا له مخې د کوم~\LR{$\sigma.n$} لپاره
      \LR{$\sFmla{\True}{!B}[\sigma.n] \in \Delta$} وي۔ د استقرا د
      فرضيې له مخې \LR{$\mSat{M(\Delta)}{!B}[\sigma.n]$} وي۔
      ځکه چې \LR{$R\sigma(\sigma.n)$}، داسې~\LR{$\sigma'$} شته چې
      \LR{$\mSat{M(\Delta)}{!B}[\sigma']$} وي؛ نو
      \LR{$\mSat{M(\Delta)}{\indfrm}[\sigma]$} وي۔
\olpseganchor{OLP-0468-B020}{end}

\olpseganchor{OLP-0468-B021}{start}
      که \LR{$\sFmla{\False}{\indfrm}[\sigma] \in \Delta$} وي، نو د
      څانګې د بشپړتيا له مخې د هر کارېدلي~\LR{$\sigma.n$} لپاره
      \LR{$\sFmla{\False}{!B}[\sigma.n] \in \Delta$} وي؛ يعنې د هر
      هغه~\LR{$\sigma' \in P(\Delta)$} لپاره چې \LR{$R\sigma\sigma'$} وي۔
      د استقرا د فرضيې له مخې د هر هغه~\LR{$\sigma'$} لپاره چې
      \LR{$R\sigma\sigma'$} وي،
      \LR{$\mSat/{M(\Delta)}{!B}[\sigma']$} وي؛ نو
      \LR{$\mSat/{M(\Delta)}{\indfrm}[\sigma]$} وي۔}
\end{enumerate}
ځکه چې \LR{$\Gamma \subseteq \Delta$}، نو
\LR{$\mSat{M(\Delta)}{\Gamma}[\mathrm{id}]$} وي۔ دلته د مخکښونو
تفسير هر مخکښ هماغې نړۍ ته لېږي چې په همدې مخکښ نومول شوې ده۔
\end{proof}
\olpseganchor{OLP-0468-B021}{end}

\olpseganchor{OLP-0468-B022}{start}
\begin{prob}
د \olref[nml][tab][cpl]{thm:tableau-completeness} ثبوت بشپړ کړئ۔
\end{prob}
\olpseganchor{OLP-0468-B022}{end}

\olpseganchor{OLP-0468-B023}{start}
\begin{cor}
\ollabel{cor:entailment-completeness}
که \LR{$\Gamma \Entails !A$} وي، نو \LR{$\Gamma \Proves !A$} وي۔
\end{cor}
له \olref[nml][com][lin]{cor:provability-characterization} او د
رښتياوالي له لمې نه د هيلبرټي اشتقاق بشپړتيا د فرضيو د هر سټ لپاره
راځي: که پايله اشتقاقېدونکې نۀ وي، يو بشپړ سازګار غځون شته چې ټول
فرضونه لري او پايله نۀ لري؛ د کانوني مدل اړونده نړۍ ضدمدل دے۔
د هيلبرټي اشتقاق تعريف متناهي ډېر فرضونه اخلي، او د اشتقاق سموالی
ښيي چې همدغه متناهي فرضونه پايله لازموي۔ دغه فرضونه د پيل په يو
مخکښ رښتيني ونيسئ او پايله پر هماغه مخکښ نارښتينې ونيسئ۔ که تړلې
تابلو نۀ وای، پورتنۍ قضيه به د دې نښو لپاره مدل ورکاوه۔ نو تړلې
تابلو شته، او د ټول سټ لپاره مطلوب اشتقاق ترې راځي۔
\olpseganchor{OLP-0468-B023}{end}

\olpseganchor{OLP-0468-B024}{start}
\begin{cor}
\ollabel{cor:weak-completeness}
که~\LR{$!A$} په ټولو مدلونو کښې رښتينے وي، نو \LR{$\Proves !A$} وي۔
\end{cor}
\olpseganchor{OLP-0468-B024}{end}

\olpseganchor{OLP-0469-B004}{start}
\olfileid{nml}{tab}{cou}
\olpseganchor{OLP-0469-B004}{end}
      
\olpseganchor{OLP-0469-B005}{start}
\olsection{له تابلوګانې څخه ضدمدلونه}
\olpseganchor{OLP-0469-B005}{end}

\olpseganchor{OLP-0469-B006}{start}
\begin{explain}
  د بشپړتيا د قضيې ثبوت يوازې دا نۀ ښيي چې که
  \LR{$\Entails !A$} وي، نو \LR{$\Proves !A$} وي؛ که
  \LR{$\Entails/ !A$} وي، نو د~\LR{$!A$} د ضدمدل د جوړولو لاره هم راکوي۔
  د ضدمدل د لټون پيلنۍ نښه د~\LR{$!A$} لپاره نارښتيا ده، او
  همدغه ازموينه د هر فارمول \LR{$A$} لپاره کاروو، چې دلته ئې~\LR{$!A$} ليکو۔
  د~\LR{$\Log{K}$} په حالت کښې دا لاره د \emph{پرېکړې کړنلاره}
  جوړوي۔ فرض کړئ \LR{$\Entails/ !A$} وي۔ د
  \olref[cpl]{prop:complete-tableau} ثبوت د بشپړ
  تابلو د جوړولو طريقه راکوي، او دا طريقه تل درېږي۔
  د~\LR{$\Log{K}$} قضيېيزې قاعدې يوازې د لږې پېچلتيا
  مخکښ‌لرونکي فارمولونه ورزياتوي؛ يعنې هره قضيېيزه قاعده
  پر څانګه د هر نښه لرونکي فارمول
  \LR{$\sFmla{S}{!A}[\sigma]$} لپاره يوازې يو ځل پلې کېدے شي۔
  نوي مخکښونه يوازې د
  \LR{$\TRule{\False}{\Box}$}
    او \LR{$\TRule{\True}{\Diamond}$}
  قاعدې جوړوي؛ دا هم
  يوازې يو ځل پلې کېږي او يو نوے مخکښ جوړوي۔
  \LR{$\TRule{\True}{\Box}$}
    او \LR{$\TRule{\False}{\Diamond}$}
  قاعدې کېدے شي څو
  ځله پلې شي، خو د هرې مقدمې او کارېدلي مخکښ د جوړې لپاره
  يو ځل بس دے۔ د بشپړې تابلو د جوړولو په ثبوت کښې مو وليدل
  چې د مودالي ژورتيا حد د نويو مخکښونو شمېر محدودوي۔ نو
  کړنلاره يا ټولې څانګې تړي يا يوه پرانيستې بشپړه څانګه پرېږدي۔
\olpseganchor{OLP-0469-B006}{end}

\olpseganchor{OLP-0469-B007}{start}
  کله چې داسې تابلو جوړه شي چې يوه پرانيستې بشپړه څانګه ولري،
  د \olref[cpl]{thm:tableau-completeness} ثبوت څرګند مدل
  راکوي چې د اصلي مخکښ‌لرونکو فارمولونه سټ پوره کوي۔ نو
  يوازې دا نۀ ده چې که \LR{$\Gamma \Entails !A$} وي، تړلے
  تابلو او \LR{$\Gamma \Proves !A$} هم وي؛ که د تړلې
  تابلو لټون په سمه منظمه طريقه وکړو او «بشپړې»
  تابلو ته
  ورسېږو چې پرانيستې څانګه ولري، نو \LR{$\Gamma \Entails/ !A$} به هم وپېژنو او ضدمدل
  به په رښتيا جوړولے شو۔
\end{explain}
\olpseganchor{OLP-0469-B007}{end}

\olpseganchor{OLP-0469-B008}{start}

\begin{ex}
  پوهېږو چې \LR{$\Proves/ \Box(p \lor q) \lif (\Box p \lor \Box
  q)$}۔ د تابلو جوړول داسې پيلېږي:
  \begin{oltableau}
    [\pFmla{\False}{\Box(p \lor q) \lif (\Box p \lor \Box q)}{1},
      just = \TAss, checked
      [\pFmla{\True}{\Box(p \lor q)}{1},
        just = {\TRule{\False}{\lif}[1]}, 
        [\pFmla{\False}{\Box p \lor \Box q}{1},
          just = {\TRule{\False}{\lif}[1]}, checked
          [\pFmla{\False}{\Box p}{1},
            just = {\TRule{\False}{\lor}[3]}, checked
            [\pFmla{\False}{\Box q}{1},
              just = {\TRule{\False}{\lor}[3]}, checked
              [\pFmla{\False}{p}{1.1}, 
                just = {\TRule{\False}{\Box}[4]}, checked
                [\pFmla{\False}{q}{1.2}, 
                  just = {\TRule{\False}{\Box}[5]}, checked
                ]
              ]
            ]
          ]
        ]
      ]
    ]
  \end{oltableau}
  دا تابلو لا نۀ ده بشپړه۔ په ورپسې پړاو کښې يوازې
  هغه کرښه ګورو چې د کتنې نښه نۀ لري: په~\LR{$2$}مه کرښه کښې
  مخکښ‌لرونکے فارمول
  \LR{$\sFmla{\True}{\Box(p \lor q)}[1]$}۔ د بشپړې تابلو د
  جوړولو له مخې د \LR{$\TRule{\True}{\Box}$} قاعده بايد د څانګې
  د هر کارېدلي مخکښ لپاره، يعنې د \LR{$1.1$} او \LR{$1.2$} دواړو لپاره
  پلې شي، څو پرانيستې څانګه بشپړه شي۔
  \begin{oltableau}
    [\pFmla{\False}{\Box(p \lor q) \lif (\Box p \lor \Box q)}{1},
      just = \TAss, checked
      [\pFmla{\True}{\Box(p \lor q)}{1},
        just = {\TRule{\False}{\lif}[1]}, 
        [\pFmla{\False}{\Box p \lor \Box q}{1},
          just = {\TRule{\False}{\lif}[1]}, checked
          [\pFmla{\False}{\Box p}{1},
            just = {\TRule{\False}{\lor}[3]}, checked
            [\pFmla{\False}{\Box q}{1},
              just = {\TRule{\False}{\lor}[3]}, checked
              [\pFmla{\False}{p}{1.1}, 
                just = {\TRule{\False}{\Box}[4]}, checked
                [\pFmla{\False}{q}{1.2}, 
                  just = {\TRule{\False}{\Box}[5]}, checked
                  [\pFmla{\True}{p \lor q}{1.1}, 
                    just = {\TRule{\True}{\Box}[2]}
                    [\pFmla{\True}{p \lor q}{1.2}, 
                      just = {\TRule{\True}{\Box}[2]}
                    ]
                  ]
                ]
              ]
            ]
          ]
        ]
      ]
    ]
  \end{oltableau}
  اوس ۲مه، ۸مه او ۹مه کرښې د کتنې نښه نۀ لري۔ خو نوے
  مخکښ نۀ دے ورزيات شوے؛ نو \LR{$\TRule{\True}{\lor}$}
  په ۸مه او ۹مه کرښو او ټولو راپيدا شويو څانګو پلې کوو،
  تر هغو چې څانګه تړلې نۀ شي:
  \begin{oltableau}
    [\pFmla{\False}{\Box(p \lor q) \lif (\Box p \lor \Box q)}{1},
      just = \TAss, checked
      [\pFmla{\True}{\Box(p \lor q)}{1},
        just = {\TRule{\False}{\lif}[1]}, checked
        [\pFmla{\False}{\Box p \lor \Box q}{1},
          just = {\TRule{\False}{\lif}[1]}, checked
          [\pFmla{\False}{\Box p}{1},
            just = {\TRule{\False}{\lor}[3]}, checked
            [\pFmla{\False}{\Box q}{1},
              just = {\TRule{\False}{\lor}[3]}, checked
              [\pFmla{\False}{p}{1.1}, 
                just = {\TRule{\False}{\Box}[4]}, checked
                [\pFmla{\False}{q}{1.2}, 
                  just = {\TRule{\False}{\Box}[5]}, checked
                  [\pFmla{\True}{p \lor q}{1.1}, 
                    just = {\TRule{\True}{\Box}[2]}, checked
                    [\pFmla{\True}{p \lor q}{1.2}, 
                      just = {\TRule{\True}{\Box}[2]}, checked
                      [\pFmla{\True}{p}{1.1},
                        just = {\TRule{\True}{\lor}[8]}, checked, close
                      ]
                      [\pFmla{\True}{q}{1.1},
                        just = {\TRule{\True}{\lor}[8]}, checked
                        [\pFmla{\True}{p}{1.2},
                          just = {\TRule{\True}{\lor}[9]}, checked]
                        [\pFmla{\True}{q}{1.2},
                          just = {\TRule{\True}{\lor}[9]}, checked, close]
                      ]
                    ]
                  ]
                ]
              ]
            ]
          ]
        ]
      ]
    ]
  \end{oltableau}
  يوه پرانيستې څانګه پاتې ده او بشپړه ده۔ له هغې داسې مدل
  تعريفوو چې نړۍې يې \LR{$W = \{1, 1.1, 1.2\}$} دي، يعنې
  په پرانيستې څانګه کښې راغلي يوازيني مخکښونه؛ د لاسرسي
  اړيکه يې \LR{$R = \{\tuple{1, 1.1}, \tuple{1, 1.2}\}$} ده؛
  او ارزښت ټاکنه يې \LR{$V(p) = \{1.2\}$} ده، ځکه په ۱۱مه کرښه
  کښې \LR{$\sFmla{\True}{p}[1.2]$} شته، او \LR{$V(q) = \{1.1\}$}
  ده، ځکه په ۱۰مه کرښه کښې
  \LR{$\sFmla{\True}{q}[1.1]$} شته۔ مدل په
  \olref{fig:counter-Box} کښې انځور شوے؛ کتلے شئ چې دا د
  \LR{$\Box(p \lor q) \lif (\Box p \lor \Box q)$} ضدمدل دے۔
  \begin{figure}
  \begin{center}
    \begin{tikzpicture}[modal]
      \node[world] (w1) [label={[align=right]right:\mFalse{p}\\ \mFalse{q}}]{\LR{$1$}}; 
      \node[world] (w2) [label={[align=right]right:\mFalse{p}\\ \mTrue{q}},
        above left=of w1]{\LR{$1.1$}}; 
      \node[world] (w3) [label={[align=right]right:\mTrue{p}\\ \mFalse{q}},
        above right=of w1] {\LR{$1.2$}};
      \draw[->] (w1) to (w2);
      \draw[->] (w1) to (w3);
    \end{tikzpicture}
  \end{center}
  \caption{د \LR{$\Box(p \lor q) \lif (\Box p \lor \Box
    q)$} ضدمدل۔}
  \ollabel{fig:counter-Box}
\end{figure}
\end{ex}          


\begin{ex}
  پوهېږو چې \LR{$\Proves/ (\Diamond p \land \Diamond q) \lif \Diamond(p
  \land q)$}۔ د تابلو جوړول داسې پيلېږي:
  \begin{oltableau}
    [\pFmla{\False}{(\Diamond p \land \Diamond q) \lif \Diamond(p \land q)}{1},
      just = \TAss, checked
      [\pFmla{\True}{\Diamond p \land \Diamond q}{1},
        just = {\TRule{\False}{\lif}[1]}, checked
        [\pFmla{\False}{\Diamond(p \land q)}{1},
          just = {\TRule{\False}{\lif}[1]}
          [\pFmla{\True}{\Diamond p}{1},
            just = {\TRule{\True}{\land}[2]}, checked
            [\pFmla{\True}{\Diamond q}{1},
              just = {\TRule{\True}{\land}[2]}, checked
              [\pFmla{\True}{p}{1.1}, 
                just = {\TRule{\True}{\Diamond}[4]}, checked
                [\pFmla{\True}{q}{1.2}, 
                  just = {\TRule{\True}{\Diamond}[5]}, checked
                ]
              ]
            ]
          ]
        ]
      ]
    ]
  \end{oltableau}
  دا تابلو لا نۀ ده بشپړه۔ په ورپسې پړاو کښې يوازې
  هغه کرښه ګورو چې د کتنې نښه نۀ لري: په~\LR{$3$}مه کرښه کښې
  مخکښ‌لرونکے فارمول
  \LR{$\sFmla{\False}{\Diamond(p \land q)}[1]$}۔ د بشپړې تابلو
  د جوړولو له مخې د \LR{$\TRule{\False}{\Diamond}$} قاعده بايد
  د څانګې د هر کارېدلي مخکښ لپاره، يعنې د \LR{$1.1$} او \LR{$1.2$}
  دواړو لپاره پلې شي، څو پرانيستې څانګه بشپړه شي۔
  \begin{oltableau}
    [\pFmla{\False}{(\Diamond p \land \Diamond q) \lif \Diamond(p \land q)}{1},
      just = \TAss, checked
      [\pFmla{\True}{\Diamond p \land \Diamond q}{1},
        just = {\TRule{\False}{\lif}[1]}, checked
        [\pFmla{\False}{\Diamond (p \land q)}{1},
          just = {\TRule{\False}{\lif}[1]}
          [\pFmla{\True}{\Diamond p}{1},
            just = {\TRule{\True}{\land}[2]}, checked
            [\pFmla{\True}{\Diamond q}{1},
              just = {\TRule{\True}{\land}[2]}, checked
              [\pFmla{\True}{p}{1.1}, 
                just = {\TRule{\True}{\Diamond}[4]}, checked
                [\pFmla{\True}{q}{1.2}, 
                  just = {\TRule{\True}{\Diamond}[5]}, checked
                  [\pFmla{\False}{p \land q}{1.1}, 
                    just = {\TRule{\False}{\Diamond}[3]}
                    [\pFmla{\False}{p \land q}{1.2}, 
                      just = {\TRule{\False}{\Diamond}[3]}
                    ]
                  ]
                ]
              ]
            ]
          ]
        ]
      ]
    ]
  \end{oltableau}
  اوس ۳مه، ۸مه او ۹مه کرښې د کتنې نښه نۀ لري۔ خو نوے
  مخکښ نۀ دے ورزيات شوے؛ نو \LR{$\TRule{\False}{\land}$}
  په ۸مه او ۹مه کرښو او ټولو راپيدا شويو څانګو پلې کوو،
  تر هغو چې څانګه تړلې نۀ شي:
  \begin{oltableau}
    [\pFmla{\False}{(\Diamond p \land \Diamond q) \lif \Diamond(p \land q)}{1},
      just = \TAss, checked
      [\pFmla{\True}{\Diamond p \land \Diamond q}{1},
        just = {\TRule{\False}{\lif}[1]}, checked
        [\pFmla{\False}{\Diamond(p \land q)}{1},
          just = {\TRule{\False}{\lif}[1]}
          [\pFmla{\True}{\Diamond p}{1},
            just = {\TRule{\True}{\land}[2]}, checked
            [\pFmla{\True}{\Diamond q}{1},
              just = {\TRule{\True}{\land}[2]}, checked
              [\pFmla{\True}{p}{1.1}, 
                just = {\TRule{\True}{\Diamond}[4]}, checked
                [\pFmla{\True}{q}{1.2}, 
                  just = {\TRule{\True}{\Diamond}[5]}, checked
                  [\pFmla{\False}{p \land q}{1.1}, 
                    just = {\TRule{\False}{\Diamond}[3]}, checked
                    [\pFmla{\False}{p \land q}{1.2}, 
                      just = {\TRule{\False}{\Diamond}[3]}, checked
                      [\pFmla{\False}{p}{1.1},
                        just = {\TRule{\False}{\land}[8]}, checked, close
                      ]
                      [\pFmla{\False}{q}{1.1},
                        just = {\TRule{\False}{\land}[8]}, checked
                        [\pFmla{\False}{p}{1.2},
                          just = {\TRule{\False}{\land}[9]}, checked]
                        [\pFmla{\False}{q}{1.2},
                          just = {\TRule{\False}{\land}[9]}, checked, close]
                      ]
                    ]
                  ]
                ]
              ]
            ]
          ]
        ]
      ]
    ]
  \end{oltableau}
  يوه پرانيستې څانګه پاتې ده او بشپړه ده۔ له هغې داسې مدل
  تعريفوو چې نړۍې يې \LR{$W = \{1, 1.1, 1.2\}$} دي، يعنې
  په پرانيستې څانګه کښې راغلي يوازيني مخکښونه؛ د لاسرسي
  اړيکه يې \LR{$R = \{\tuple{1, 1.1}, \tuple{1, 1.2}\}$} ده؛
  او ارزښت ټاکنه يې \LR{$V(p) = \{1.1\}$} ده، ځکه په ۶مه کرښه
  کښې \LR{$\sFmla{\True}{p}[1.1]$} شته، او \LR{$V(q) = \{1.2\}$}
  ده، ځکه په ۷مه کرښه کښې
  \LR{$\sFmla{\True}{q}[1.2]$} شته۔ مدل په \olref{fig:counter-Diamond} کښې انځور
  شوے؛ کتلے شئ چې دا د
  \LR{$(\Diamond p \land \Diamond q) \lif \Diamond (p \land q)$}
  ضدمدل دے۔
  \begin{figure}
  \begin{center}
    \begin{tikzpicture}[modal]
      \node[world] (w1) [label={[align=right]right:\mFalse{p}\\ \mFalse{q}}]{\LR{$1$}}; 
      \node[world] (w2) [label={[align=right]right:\mTrue{p}\\ \mFalse{q}},
        above left=of w1]{\LR{$1.1$}}; 
      \node[world] (w3) [label={[align=right]right:\mFalse{p}\\ \mTrue{q}},
        above right=of w1] {\LR{$1.2$}};
      \draw[->] (w1) to (w2);
      \draw[->] (w1) to (w3);
    \end{tikzpicture}
  \end{center}
  \caption{د \LR{$(\Diamond p \land \Diamond q) \lif
    \Diamond (p \land q)$} ضدمدل۔}
  \ollabel{fig:counter-Diamond}
\end{figure}
\end{ex}          

\olpseganchor{OLP-0469-B008}{end}

\olpseganchor{OLP-0470-B004}{start}
\olchapter{nml}{seq}{د مودالي سېکوېنټ حساب}
%\tagfalse{prvDiamond}
\begin{editorial}
  په دې باب کښې د مودالي سېکوېنټ حساب قاعدې او د ثبوتونو
  بېلګې راځي۔ ځينې قاعدې د لاسرسي پر اړيکه ځانګړي شرطونه
  لري، او د پرېکون کارول په يوه بېلګه کښې څېړل کېږي۔
\end{editorial}
\olpseganchor{OLP-0470-B004}{end}

\olpseganchor{OLP-0470-B005}{start}



%

%
%
\olpseganchor{OLP-0470-B005}{end}

\olpseganchor{OLP-0471-B004}{start}
\olfileid{nml}{seq}{int}
\olpseganchor{OLP-0471-B004}{end}

\olpseganchor{OLP-0471-B005}{start}
\olsection{سريزه}
\olpseganchor{OLP-0471-B005}{end}

\olpseganchor{OLP-0471-B006}{start}
د قضيېيز منطق د سېکوېنټ حساب په هغو اضافي قاعدو غځېدلے شي چې
له~\LR{$\Box$} او~%
\LR{$\Diamond$} سره کار کوي۔ مثلاً د~\LR{$\Log{K}$}
لپاره له \Log{LK} سره دا قاعدې هم شته:
\[
  % <> and [] primitive
    \Axiom$\Gamma \fCenter \Delta, !A$
    \RightLabel{\LR{$\Box$}}
    \UnaryInf$\Box\Gamma \fCenter \Diamond\Delta, \Box!A$
    \DisplayProof  
    \qquad
    \Axiom$!A, \Gamma \fCenter \Delta$
    \RightLabel{\LR{$\Diamond$}}
    \UnaryInf$\Diamond!A, \Box\Gamma \fCenter \Diamond\Delta$
    \DisplayProof
\]
د~\LR{$\Log{K}$} د غځونو لپاره نورې قاعدې هم زياتول پکار دي۔
\olpseganchor{OLP-0471-B006}{end}

\olpseganchor{OLP-0471-B007}{start}
د مودالي منطقونو د سېکوېنټ قواعد د هغوی له معناپوهنې سره بدلېږي۔
د~\LR{$\Log{S5}$} معناپوهنه په عمومي مدلونو بيانېدے شي، چې لاسرسے
پکښې د نړۍو ټولې جوړې رانغاړي۔ په دې باب کښې له~\LR{$\Log{LK}$} نه
راوتلے حساب د~\Cut قاعده کاروي۔ د پرېکون له قاعدې پرته د عادي
سېکوېنټونو يو حساب د مودالي معياري بڼو د اړولو قواعد کاروي؛
د پرېکون له قاعدې پرته بشپړ \emph{هايپر سېکوېنټ} حساب هم شته۔\footnote{د عادي سېکوېنټونو لپاره
\href{https://doi.org/10.18778/0138-0680.45.2.05}{Indrzejczak (2016)}،
او د هايپر سېکوېنټونو لپاره \href{https://arxiv.org/abs/1905.09039v1}{Aghaei او Mohammadi (2019)} وګورئ۔}
\olpseganchor{OLP-0471-B007}{end}

\olpseganchor{OLP-0472-B004}{start}
\olfileid{nml}{seq}{rul}
\olpseganchor{OLP-0472-B004}{end}

\olpseganchor{OLP-0472-B005}{start}
\olsection{د \texorpdfstring{\Log{K}}{K} لپاره قاعدې}
\olpseganchor{OLP-0472-B005}{end}

\olpseganchor{OLP-0472-B006}{start}
د عادي قضيېيزو اړيکو قاعدې د عادي سېکوېنټ حساب~\LR{$\Log{LK}$}
له قاعدو سره يو شان دي۔ بديهي اصول هم هماغه دي: د
\LR{$!A \Sequent !A$} په بڼه هر سېکوېنټ بديهي اصل شمېرل کېږي۔
\olpseganchor{OLP-0472-B006}{end}

\olpseganchor{OLP-0472-B007}{start}
د مودالي عملګرونو~\LR{$\Box$}
او~\LR{$\Diamond$} لپاره دا
اضافي قاعدې لرو:
\[
  % <> and [] primitive
    \Axiom$\Gamma \fCenter \Delta, !A$
    \RightLabel{\LR{$\Box$}}
    \UnaryInf$\Box\Gamma \fCenter \Diamond\Delta, \Box!A$
    \DisplayProof  
    \qquad
    \Axiom$!A, \Gamma \fCenter \Delta$
    \RightLabel{\LR{$\Diamond$}}
    \UnaryInf$\Diamond!A, \Box\Gamma \fCenter \Diamond\Delta$
    \DisplayProof
\]
دلته \LR{$\Box\Gamma$} د هغو فارمولونه لړۍ ده
چې د~\LR{$\Gamma$} د هر فارمول مخې ته \LR{$\Box$} په اېښودلو
له~\LR{$\Gamma$} څخه تر لاسه کېږي، او
\LR{$\Diamond\Delta$} د هغو فارمولونه لړۍ ده
چې د~\LR{$\Delta$} د هر فارمول مخې ته \LR{$\Diamond$} په اېښودلو
له~\LR{$\Delta$} څخه تر لاسه کېږي۔ %
%
\LR{$\Gamma$} او\LR{$\Delta$}
تشه کېدے شي؛ په دې حالت کښې د پايلې سېکوېنټ اړونده برخه
\LR{$\Box\Gamma$} او\LR{$\Diamond\Delta$}
هم تشه وي۔ که د بکس د قاعدې پايله په کومه نړۍ کښې ناسمه
وي، نارښتينے بکس په يوه لاسرسي وړ نړۍ کښې د دننني فارمول
نارښتيا ښيي۔ په هماغې نړۍ کښې ټول کيڼ څنګ‌فارمولونه رښتيني
او ټول ښي څنګ‌فارمولونه نارښتيني وي، ځکه په پايله کښې د
هغوی بکسونه رښتيني او امکانونه نارښتيني دي۔ نو مقدمه هم ناسمه
ده۔ د امکان په قاعده کښې رښتينے امکان په لاسرسي وړ نړۍ کښې
رښتينے دنننے فارمول ورکوي، او هماغه څنګ‌فارمولونه بيا د مقدمې
د ناسموالي شاهدان وي۔
\olpseganchor{OLP-0472-B007}{end}

\olpseganchor{OLP-0472-B008}{start}
يوازې يوه فارمول~\LR{$!A$} ته
په ښۍ خوا کښې د \LR{$\Box$} زياتولو او
په کيڼه خوا کښې د \LR{$\Diamond$} زياتولو
قيد اړين دے۔ که مو
د ښۍ خوا هر څومره فارمولونه ته د \LR{$\Box$}
زياتولو يا د کيڼې خوا
هر څومره فارمولونه ته د \LR{$\Diamond$} زياتولو اجازه ورکړې
واے، نو دا مو اشتقاق کولے شول:
  \[
      \Axiom$!A \fCenter !A$
      \RightLabel{\RightR{\lnot}}
      \UnaryInf$\fCenter !A, \lnot !A$
      \RightLabel{\LR{$\Box*$}}
      \UnaryInf$\fCenter \Box!A, \Box\lnot !A$
      \doubleLine
      \RightLabel{\RightR{\lor}}
      \UnaryInf$\fCenter \Box!A \lor \Box \lnot !A$
      \DisplayProof
    \qquad
      \Axiom$!A \fCenter !A$
      \RightLabel{\LeftR{\lnot}}
      \UnaryInf$\lnot !A, !A \fCenter$
      \RightLabel{\LR{$\Diamond*$}}
      \UnaryInf$\Diamond\lnot !A,\Diamond!A \fCenter $
      \RightLabel{\RightR{\lnot}}
      \UnaryInf$\Diamond!A \fCenter \lnot\Diamond\lnot !A$
      \RightLabel{\RightR{\lif}}
      \UnaryInf$ \fCenter \Diamond!A \lif \lnot\Diamond\lnot !A$
      \DisplayProof
  \]
خو په~\LR{$\Log{K}$} کښې \LR{$\Box!A \lor \Box \lnot !A$} او
\LR{$\Diamond!A \lif \lnot\Diamond\lnot !A$} دواړه معتبر نۀ دي۔ د دواړو بېلګو لپاره
داسې نړۍ واخلئ چې دوو نړۍو ته لاسرسے لري۔ دنننے فارمول
په يوه کښې رښتينے او په بله کښې نارښتينے ونيسئ۔
\olpseganchor{OLP-0472-B008}{end}

\olpseganchor{OLP-0472-B009}{start}
که مو په مقدمه کښې له~\LR{$!A$} سره نور څنګ‌فارمولونه هم پرېښي، او
د \LR{$\Box$} قاعدې ته مو اجازه ورکړې واے چې په ښۍ خوا
کښې يوازې د~\LR{$!A$} مخې ته \LR{$\Box$} کېږدي، يا 
د \LR{$\Diamond$} قاعدې ته مو اجازه ورکړې واے چې
په کيڼه خوا کښې يوازې د~\LR{$!A$} مخې ته \LR{$\Diamond$} کېږدي
(خو څنګ‌فارمولونه بدل نۀ کړي)، نو دا مو اشتقاق کولے شول:
\[
    \Axiom$!A \fCenter !A$
    \RightLabel{\RightR{\lnot}}
    \UnaryInf$\fCenter !A, \lnot !A$
    \RightLabel{\RightR{\Exchange}}
    \UnaryInf$\fCenter \lnot !A, !A$
    \RightLabel{\LR{$\Box*$}}
    \UnaryInf$\fCenter \lnot!A, \Box !A$
    \doubleLine
    \RightLabel{\RightR{\lor}}
    \UnaryInf$\fCenter \lnot!A \lor \Box !A$
    \DisplayProof
  \qquad
    \Axiom$!A \fCenter !A$
    \RightLabel{\LeftR{\lnot}}
    \UnaryInf$\lnot !A, !A \fCenter$
    \RightLabel{\LR{$\Diamond*$}}
    \UnaryInf$\Diamond\lnot !A, !A \fCenter $
    \RightLabel{\RightR{\lnot}}
    \UnaryInf$!A \fCenter \lnot\Diamond\lnot !A$
    \RightLabel{\RightR{\lif}}
    \UnaryInf$ \fCenter !A \lif \lnot\Diamond\lnot !A$
    \DisplayProof
\]
خو په~\LR{$\Log{K}$} کښې \LR{$\lnot!A \lor \Box !A$} (چې له \LR{$!A
\lif \Box!A$} سره همارز دے) او \LR{$!A \lif \lnot\Diamond\lnot !A$}
دواړه معتبر نۀ دي۔
د دواړو بېلګو لپاره دنننے فارمول په لومړۍ نړۍ کښې رښتينے
او په هغې ته د لاسرسي په يوه نړۍ کښې نارښتينے ونيسئ۔
\olpseganchor{OLP-0472-B009}{end}

\olpseganchor{OLP-0473-B004}{start}
\olfileid{nml}{seq}{prk}
\olpseganchor{OLP-0473-B004}{end}

\olpseganchor{OLP-0473-B005}{start}
\olsection{د \texorpdfstring{\Log{K}}{K} لپاره سېکوېنټ اشتقاقونه}
\olpseganchor{OLP-0473-B005}{end}



\olpseganchor{OLP-0473-B007}{start}
\begin{ex}
  دلته د سېکوېنټ حساب داسې اشتقاق وړاندې کوو چې
  \LR{$\Proves (\Box!A \land \Box!B)
  \lif \Box (!A \land !B)$} ښيي۔
  \begin{prooftree}
    \Axiom$!A \fCenter !A$
    \doubleLine
    \UnaryInf$!B, !A \fCenter !A$   
    \Axiom$!B \fCenter !B$
    \doubleLine
    \UnaryInf$!B, !A \fCenter !B$   
    \RightLabel{\RightR{\land}}
    \BinaryInf$!B, !A \fCenter !A \land !B$    
    \RightLabel{\LR{$\Box$}}
    \UnaryInf$\Box!B, \Box!A \fCenter \Box
    (!A \land !B)$
    \RightLabel{\LeftR{\land}}
    \UnaryInf$\Box!A \land \Box!B, \Box!A \fCenter \Box
    (!A \land !B)$
    \RightLabel{\LeftR{\Exchange}}
    \UnaryInf$\Box!A, \Box!A \land \Box!B \fCenter \Box
    (!A \land !B)$
    \RightLabel{\LeftR{\land}}
    \UnaryInf$\Box!A \land \Box!B, \Box!A \land \Box!B \fCenter \Box (!A \land !B)$
    \RightLabel{\LeftR{\Contraction}}
    \UnaryInf$\Box!A \land \Box!B \fCenter \Box (!A \land !B)$
    \RightLabel{\RightR{\lif}}
    \UnaryInf$\fCenter (\Box!A \land \Box!B)
    \lif \Box (!A \land !B)$
  \end{prooftree}
\end{ex}

\olpseganchor{OLP-0473-B007}{end}

\olpseganchor{OLP-0473-B008}{start}

  \begin{ex}
    دلته د سېکوېنټ حساب داسې اشتقاق وړاندې کوو چې
    \LR{$\Proves \Diamond(!A
    \lor !B) \lif (\Diamond !A \lor \Diamond!B)$} ښيي۔
    \begin{prooftree}
      \Axiom$!A \fCenter !A$
      \doubleLine
      \UnaryInf$!A \fCenter !A, !B$   
      \Axiom$!B \fCenter !B$
      \doubleLine
      \UnaryInf$!B \fCenter !A, !B$   
      \RightLabel{\LeftR{\lor}}
      \BinaryInf$!A \lor !B \fCenter !A, !B$    
      \RightLabel{\LR{$\Diamond$}}
      \UnaryInf$\Diamond(!A \lor !B) \fCenter 
      \Diamond !A, \Diamond!B$
      \RightLabel{\RightR{\lor}}
      \UnaryInf$\Diamond(!A \lor !B) \fCenter 
      \Diamond !A, \Diamond !A \lor \Diamond!B$
      \RightLabel{\RightR{\Exchange}}
      \UnaryInf$\Diamond(!A \lor !B) \fCenter \Diamond !A \lor \Diamond!B, \Diamond !A$
      \RightLabel{\RightR{\lor}}
      \UnaryInf$\Diamond(!A \lor !B) \fCenter 
      \Diamond !A \lor \Diamond!B, \Diamond !A \lor \Diamond!B$
      \RightLabel{\RightR{\Contraction}}
      \UnaryInf$\Diamond(!A
      \lor !B) \fCenter \Diamond !A \lor \Diamond!B$
      \RightLabel{\RightR{\lif}}
      \UnaryInf$\fCenter \Diamond(!A
      \lor !B) \lif (\Diamond !A \lor \Diamond!B)$
    \end{prooftree}
  \end{ex}

\olpseganchor{OLP-0473-B008}{end}

\olpseganchor{OLP-0473-B009}{start}

  دلته د~\Dual يو اشتقاق دے۔ په وروستي ګام کښې
  دوه‌اړخيز شرط د دوو شرطي فارمولونو د عطف لنډون دے۔
  \begin{prooftree}
    \Axiom$!A \fCenter !A$
    \RightLabel{\LeftR{\lnot}}
    \UnaryInf$\lnot !A, !A \fCenter $
    \RightLabel{\LR{$\Diamond$}}
    \UnaryInf$\Diamond\lnot !A, \Box !A \fCenter $
    \RightLabel{\RightR{\lnot}}
    \UnaryInf$\Box !A \fCenter \lnot\Diamond\lnot !A$
    \RightLabel{\RightR{\lif}}
    \UnaryInf$\fCenter \Box !A \lif \lnot\Diamond\lnot !A$
    \Axiom$!A \fCenter !A$
    \RightLabel{\RightR{\lnot}}
    \UnaryInf$\fCenter !A, \lnot !A $
    \RightLabel{\RightR{\Exchange}}
    \UnaryInf$\fCenter \lnot !A, !A $
    \RightLabel{\LR{$\Box$}}
    \UnaryInf$\fCenter \Diamond\lnot !A, \Box !A$
    \RightLabel{\RightR{\Exchange}}
    \UnaryInf$\fCenter \Box !A, \Diamond\lnot !A$
    \RightLabel{\LeftR{\lnot}}
    \UnaryInf$\lnot\Diamond\lnot !A \fCenter \Box !A$
    \RightLabel{\RightR{\lif}}
    \UnaryInf$\fCenter \lnot \Diamond \lnot !A \lif \Box !A$
    \RightLabel{\RightR{\land}}
    \BinaryInf$\fCenter \Box !A \liff \lnot\Diamond\lnot !A$ 
  \end{prooftree} 
\olpseganchor{OLP-0473-B009}{end}

\olpseganchor{OLP-0473-B010}{start}
% this does not work if problems are deferred
% \begin{prob}
%     Give اشتقاق of $\Diamond!A \liff \lnot\Box\lnot !A$
%     in~$\Log{K}$.
% \end{prob}

\olpseganchor{OLP-0473-B010}{end}

\olpseganchor{OLP-0473-B011}{start}
\begin{prob}
  په~\LR{$\Log{K}$} کښې د لاندې فارمولونه لپاره د سېکوېنټ حساب
  ثبوتونه ومومئ:
  \begin{enumerate}
    \item \LR{$\Box \lnot p \lif \Box(p \lif q)$}
    \item \LR{$(\Box p \lor \Box q) \lif \Box(p \lor q)$}
    \item \LR{$\Diamond p \lif \Diamond(p \lor q)$}
    \item \LR{$\Box(p \land q) \lif \Box p$}
  \end{enumerate}
\end{prob}
\olpseganchor{OLP-0473-B011}{end}

\olpseganchor{OLP-0474-B004}{start}
\olfileid{nml}{seq}{mru}
\olpseganchor{OLP-0474-B004}{end}

\olpseganchor{OLP-0474-B005}{start}
\olsection{د لاسرسي د نورو اړيکو لپاره قاعدې}
\olpseganchor{OLP-0474-B005}{end}

\olpseganchor{OLP-0474-B006}{start}
د هغو منطقونو لپاره چې د لاسرسي ځانګړې اړيکې يې ټاکي،
په \olref{tab:more-rules} کښې اضافي قاعدې ګورو۔ د پنځم
اصل لپاره لاندې قواعد په متعديو اقليدسي چوکاټونو کښې کارېږي:
له يوې نړۍ نه د هغې هرې لاسرسي وړ نړۍ ته د تلو په وخت
د دواړو د لاسرسي وړ نړۍو سټ برابر وي۔ تعدي د يوه لوري
شموليت او اقليدسي‌والے د بل لوري شموليت ورکوي۔
\olpseganchor{OLP-0474-B006}{end}

\olpseganchor{OLP-0474-B007}{start}
% [] and <> primitive
\begin{table}
  \begin{center}
    \def\arraystretch{3}
    \begin{tabular}{cc}
    \hline
      \Axiom$!A, \Gamma \fCenter \Delta$
      \RightLabel{T\LR{$\Box$}}
      \UnaryInf$\Box!A, \Gamma \fCenter \Delta$
      \DisplayProof
    &
      \Axiom$\Gamma \fCenter \Delta, !A$
      \RightLabel{T\LR{$\Diamond$}}
      \UnaryInf$\Gamma \fCenter \Delta, \Diamond!A$
      \DisplayProof
    \\[1ex]
%    \hline
    \multicolumn{2}{c}{
      \Axiom$\Gamma \fCenter \Delta$
      \RightLabel{D}
      \UnaryInf$\Box\Gamma \fCenter \Diamond\Delta$
      \DisplayProof}
    \\[1ex]
%    \hline
      \Axiom$\Gamma, \Diamond\Pi \fCenter \Box\Delta, \Lambda, !A$
      \RightLabel{B\LR{$\Box$}}
      \UnaryInf$\Box\Gamma, \Pi \fCenter \Delta, \Diamond\Lambda, \Box!A$
      \DisplayProof
      &
      \Axiom$!A, \Diamond\Gamma, \Pi \fCenter \Box\Lambda, \Delta$
      \RightLabel{B\LR{$\Diamond$}}
      \UnaryInf$\Diamond!A, \Gamma, \Box\Pi \fCenter \Lambda, \Diamond\Delta$
      \DisplayProof
    \\[1ex]
%    \hline
      \Axiom$\Box\Gamma \fCenter \Diamond\Delta, !A$
      \RightLabel{4\LR{$\Box$}}
      \UnaryInf$\Box\Gamma \fCenter \Diamond\Delta, \Box!A$
      \DisplayProof
      &
      \Axiom$!A, \Box\Gamma \fCenter \Diamond\Delta$
      \RightLabel{4\LR{$\Diamond$}}
      \UnaryInf$\Diamond!A, \Box\Gamma \fCenter \Diamond\Delta$
      \DisplayProof
    \\[1ex]
%    \hline
      \Axiom$\Box\Gamma, \Diamond\Pi \fCenter \Box\Delta,
      \Diamond\Lambda, !A$
      \RightLabel{5\LR{$\Box$}}
      \UnaryInf$\Box\Gamma, \Diamond\Pi \fCenter \Box\Delta,
      \Diamond\Lambda, \Box!A$
      \DisplayProof
&
      \Axiom$!A, \Diamond\Gamma, \Box\Pi \fCenter \Diamond\Delta, \Box\Lambda$
      \RightLabel{5\LR{$\Diamond$}}
      \UnaryInf$\Diamond!A,\Diamond\Gamma,\Box\Pi \fCenter \Diamond\Delta,\Box\Lambda$
      \DisplayProof
    \\[1ex]
    \hline
    \end{tabular}
  \end{center}
  \caption{نورې مودالي قاعدې۔}
  \ollabel{tab:more-rules}
\end{table}

\olpseganchor{OLP-0474-B007}{end}

\olpseganchor{OLP-0474-B008}{start}
د دې قاعدو په ورزياتولو داسې نظامونه تر لاسه کېږي چې د
\olref{tab:logics-rules} د منطقونو لپاره سم او بشپړ دي۔
په جدول کښې راغلې قاعدې د K د سېکوېنټ حساب له قاعدو سره
يو ځاے کارول کېږي۔
\olpseganchor{OLP-0474-B008}{end}

\olpseganchor{OLP-0474-B009}{start}
\begin{table}
  \begin{center}
    \begin{tabular}{lll}
      \hline
      منطق & د~\LR{$R$} خاصيت \dots & قاعدې\\
      \hline
      \LR{$\Log{T} = \Log{KT}$} & انعکاسي &
      T\LR{$\Box$}%
      , %
      T\LR{$\Diamond$}
      \\ \hline
      \LR{$\Log{D} = \Log{KD}$} & مسلسل &
      %
        D
      \\ \hline
      \LR{$\Log{K4}$} & متعدي &
      4\LR{$\Box$}% 
      , %
      4\LR{$\Diamond$}
      \\ \hline
      \LR{$\Log{B} = \Log{KTB}$} & انعکاسي، &
      T\LR{$\Box$}%
      , %
      T\LR{$\Diamond$}\\
      & متناظره &
      B\LR{$\Box$}%
      , %
      B\LR{$\Diamond$}
      \\ \hline
      \LR{$\Log{S4} = \Log{KT4}$} & انعکاسي، &
      T\LR{$\Box$}%
      , %
      T\LR{$\Diamond$}\\
      & متعدي &
      4\LR{$\Box$}%
      , %
      4\LR{$\Diamond$}
      \\ \hline
      \LR{$\Log{S5} = \Log{KT5}$} & انعکاسي، &
      T\LR{$\Box$}%
      , %
      T\LR{$\Diamond$}\\
      & متعدي، &
      5\LR{$\Box$}%
      , %
      5\LR{$\Diamond$}\\
      &  اقليدسي &
      \\ \hline
    \end{tabular}
  \end{center}
  \caption{د بېلابېلو مودالي منطقونو سېکوېنټ قاعدې۔}
  \ollabel{tab:logics-rules}
\end{table}
\olpseganchor{OLP-0474-B009}{end}

\olpseganchor{OLP-0474-B010}{start}

\begin{ex}
  دلته داسې سېکوېنټ اشتقاق وړاندې کوو چې \LR{$\Log{K4} \Proves \Ax{4}$}، يعنې
  \LR{$\Box!A \lif \Box\Box!A$}.
  \begin{prooftree}
    \Axiom$\Box!A \fCenter \Box!A$
    \RightLabel{4\LR{$\Box$}}
    \UnaryInf$\Box!A \fCenter \Box\Box!A$
    \RightLabel{\RightR{\lif}}
    \UnaryInf$\fCenter \Box!A \lif \Box\Box!A$
  \end{prooftree}
\end{ex}


\begin{ex}
  دلته داسې سېکوېنټ اشتقاق وړاندې کوو چې \LR{$\Log{K4} \Proves \Ax{4}$}، يعنې
  \LR{$\Diamond\Diamond!A \lif \Diamond!A$}.
  \begin{prooftree}
    \Axiom$\Diamond!A \fCenter \Diamond!A$
    \RightLabel{4\LR{$\Diamond$}}
    \UnaryInf$\Diamond\Diamond!A \fCenter \Diamond!A$
    \RightLabel{\RightR{\lif}}
    \UnaryInf$\fCenter \Diamond\Diamond!A \lif \Diamond!A$
  \end{prooftree}
\end{ex}

\olpseganchor{OLP-0474-B010}{end}

\olpseganchor{OLP-0474-B011}{start}
% <> and [] primitive
\begin{ex}
  دلته داسې سېکوېنټ اشتقاق وړاندې کوو چې \LR{$\Log{S5} \Proves \Ax{5}$}، يعنې
  \LR{$\Diamond!A \lif \Box\Diamond!A$}.
  \begin{prooftree}
    \Axiom$\Diamond!A \fCenter \Diamond!A$
    \RightLabel{5\LR{$\Box$}}
    \UnaryInf$\Diamond!A \fCenter \Box\Diamond!A$
    \RightLabel{\RightR{\lif}}
    \UnaryInf$\fCenter \Diamond!A \lif \Box\Diamond!A$
  \end{prooftree}
\end{ex}

\begin{ex}
د~\Log{S5} همدا سېکوېنټ حساب د~\Cut{} له قاعدې پرته بشپړ نۀ دے؛
مثلاً \LR{$\Diamond\Box !A \lif !A$}، چې په~\LR{$\Log{S5}$} کښې معتبر دے، بې له~\Cut
څخه ثبوت نۀ لري که دنننے فارمول يو قضيېيز متغير وي۔ د پرېکون
پرته د وروستي شرط له قاعدې نه شاته چې لاړ شو، په اړونده
سېکوېنټ کښې ساده متغير د مودالي قاعدې د څنګ‌فارمول شرط ماتوي۔
د جوړښتي قاعدو له لارې د دواړو اړخونو هره ممکنه کمونه چې
يو اړخ تش کړي، يو نامعتبر سېکوېنټ ورکوي؛ تکراري نسخې هم
د مودالي قاعدې شرط نۀ شي پوره کولے۔ دلته د~\Cut په کارولو
يو اشتقاق دے:
\begin{prooftree}
% both [] and <> primitive
\Axiom$\Box !A \fCenter \Box!A$
\RightLabel{5\LR{$\Diamond$}}
\UnaryInf$\Diamond\Box!A \fCenter \Box!A$
\olpseganchor{OLP-0474-B011}{end}

\Axiom$!A \fCenter !A$
\RightLabel{T\LR{$\Box$}}
\UnaryInf$\Box!A \fCenter !A$
\RightLabel{\Cut}
\BinaryInf$\Diamond\Box!A \fCenter !A$
\RightLabel{\RightR{\lif}}
\UnaryInf$\fCenter \Diamond\Box!A \lif !A$

\end{prooftree}
\end{ex}
\olpseganchor{OLP-0474-B016}{start}
\begin{prob}
د لاندې ادعاوو د ښودلو لپاره سېکوېنټ اشتقاقونه وړاندې کړئ:
  \begin{enumerate}
  \item \LR{$\Log{KT5} \Proves \Ax{B}$};
  \item \LR{$\Log{KT5} \Proves \Ax{4}$};
  \item \LR{$\Log{KDB4} \Proves \Ax{T}$};
  \item \LR{$\Log{KB4} \Proves \Ax{5}$};
  \item \LR{$\Log{KB5} \Proves \Ax{4}$};
  \item \LR{$\Log{KT} \Proves \Ax{D}$}.
  \end{enumerate}
\end{prob}
\olpseganchor{OLP-0474-B016}{end}

\olpseganchor{OLP-0475-B004}{start}
\olpart{aml}{تطبيقي مودالي منطق}
\olpseganchor{OLP-0475-B004}{end}

\olpseganchor{OLP-0475-B005}{start}
\begin{editorial}
په دې برخه کښې د مودالي منطق زماني او معرفتي کارونې راځي۔
\end{editorial}
\olpseganchor{OLP-0475-B005}{end}

\olpseganchor{OLP-0476-B004}{start}
\olchapter{aml}{tl}{زماني منطقونه}
\olpseganchor{OLP-0476-B004}{end}

\olpseganchor{OLP-0476-B005}{start}
\begin{editorial}
  دا څپرکے زماني منطقونه څېړي۔
\end{editorial}
\olpseganchor{OLP-0476-B005}{end}

\olpseganchor{OLP-0477-B004}{start}
\olfileid{aml}{tl}{int}
\olpseganchor{OLP-0477-B004}{end}

\olpseganchor{OLP-0477-B005}{start}
\olsection{سريزه}
\olpseganchor{OLP-0477-B005}{end}

\olpseganchor{OLP-0477-B006}{start}
زماني منطقونه د هغو دعوو څېړنه کوي چې يو څۀ به په راتلونکي کښې واقع شي يا
په تېر وخت کښې واقع شوي وي۔ د زماني منطق بنسټ ايښودونکے آرتر پرايور ګڼل
کېږي؛ هغه ورته \emph{د زمانو منطق} وايۀ۔ دلته به د زماني منطق په څېړنه کښې
تر ډېره د پرايور هغه لومړنۍ مودالي طريقه تعقيبوو چې د قضيېيز منطق بنسټيز
چوکاټ ته زماني عاملونه ورزياتوي۔ د قضيېيز منطق په بنسټيزه ژبه کښې د
دعوې د تېر يا راتلونکي وخت د ښودلو ځانګړے عملګر نشته۔
\olpseganchor{OLP-0477-B006}{end}

\olpseganchor{OLP-0477-B007}{start}
د بېلګې په توګه، زه ښايي د بيزي (Beezie) په نوم د يوه سپي په اړه خبرې
وکړم چې کله کښېني او کله نۀ کښېني، لکه د سپيو عادت چې وي۔ په کلاسيک
منطق کښې دا ويل متناقض دي چې بيزي ناست دے او بيزي ناست نۀ دے۔ خو ښکاره
ده چې دواړه خبرې رښتيا کېدے شي، البته په يوه وخت کښې نۀ۔ ژبې ته د زماني
عاملونو په ورزياتولو سره دا خبره په نسبتاً اسانه توګه څرګندولے شو۔ د راتلونکي د يوې شېبې د شتون په معنا، زماني عاملونه دا هم راښيي
چې ولې له «بيزي به انعام يا توپ ترلاسه کړي» څخه «بيزي به انعام
ترلاسه کړي يا بيزي به توپ ترلاسه کړي» ته استدلال معتبر دے۔
\olpseganchor{OLP-0477-B007}{end}

\olpseganchor{OLP-0477-B008}{start}
خو زماني منطق ډېرې فلسفي پوښتنې راولاړوي چې ښايي موږ د زماني منطق يوه
چوکاټ پر بل غوره کولو ته وهڅوي۔ د بېلګې په توګه، راتلونکے اتفاقي
بيان د راتلونکي په اړه هغه بيان دے چې نۀ لازم دے او نۀ ناممکن۔ که ووايو
«ريچارډ (Richard) به سبا د خوراکي توکو دوکان ته لاړ شي»، نو د يوه داسې
کار په اړه دعوه کوو چې لا نۀ دے شوے او د رښتيا ارزښت ئې د بحث وړ دے۔
ان دا هم د بحث وړ ده چې آيا په لومړي سر کښې دې دعوې ته د رښتيا ارزښت
\emph{ټاکل} کېدے شي که نۀ۔ که موږ سخت جبرپال يو، نو ښايي دا ومنو چې
دا جمله له پېښې وړاندې هم رښتينې يا نارښتينې ده، که څه هم موږ يې
ارزښت لا نۀ پېژنو۔ که راتلونکے پرانيستے وګڼو، د راتلونکي د اتفاقي
بيانونو د رښتيا ارزښتونه لا ټاکل شوي نۀ وي۔
\olpseganchor{OLP-0477-B008}{end}

\olpseganchor{OLP-0477-B009}{start}
په پايله کښې، د وخت د جوړښت او ماهيت په اړه ډېرې داسې ژمنې د زماني
منطقونو د مدلونو او چوکاټونو په ټاکلو کښې نغښتې وي۔ د بېلګې په توګه،
له ځانه پوښتلے شو چې آيا داسې مدلونه جوړ کړو چې وخت پکښې خطي، څوڅانګيز
يا ان دوراني وي۔ ښايي دا هم وټاکو چې زموږ زماني مدلونه د پيل او پاې
ټکي لري که نۀ، او وخت د بېلو شېبو په توګه وښيو که د يوه پيوسته امتداد
په توګه۔
\olpseganchor{OLP-0477-B009}{end}

\olpseganchor{OLP-0478-B004}{start}
\olfileid{aml}{tl}{sem}
\olpseganchor{OLP-0478-B004}{end}

\olpseganchor{OLP-0478-B005}{start}
\olsection{د زماني منطق معناپوهنه}
\olpseganchor{OLP-0478-B005}{end}

\olpseganchor{OLP-0478-B006}{start}
\begin{defn}
د زماني منطق بنسټيزه ژبه دا توکي لري:
\begin{enumerate}
  \item د دروغوالی قضيوي ثابت~\LR{$\lfalse$}۔
  \item د رښتياوالی قضيوي ثابت~\LR{$\ltrue$}۔
  \item د بياني متغيرونه يو شمېرېدونکے لامتناهي سټ: \LR{$\Obj
    p_0$}، \LR{$\Obj p_1$}، \LR{$\Obj p_2$}، \dots
  \item قضيوي نښلوونکي: \startycommalist
  \ycomma \LR{$\lnot$} (نفي)%
  \ycomma \LR{$\land$} (عطف)%
  \ycomma \LR{$\lor$} (فصل)%
  \ycomma \LR{$\lif$} (شرطي)%
  \ycomma \LR{$\liff$} (دوه اړخيز شرطي).
  \item د تېر وخت عاملونه \LR{$\Ptemp$} او \LR{$\Htemp$}۔
  \item د راتلونکي وخت عاملونه \LR{$\Ftemp$} او \LR{$\Gtemp$}۔
\end{enumerate}
\end{defn}
\olpseganchor{OLP-0478-B006}{end}

\olpseganchor{OLP-0478-B007}{start}
وروسته به د مودالي عاملونو د نورو ډولونو د ورزياتولو امکان هم وڅېړو۔
\olpseganchor{OLP-0478-B007}{end}

\olpseganchor{OLP-0478-B008}{start}
\begin{defn}
د زماني ژبې \emph{فارمولونه} په استقرايي ډول داسې
  تعريفېږي:
\begin{enumerate}
\item \LR{$\lfalse$} يو اټومي فارمول دے۔
\olpseganchor{OLP-0478-B008}{end}

\olpseganchor{OLP-0478-B009}{start}
\item \LR{$\ltrue$} يو اټومي فارمول دے۔
\olpseganchor{OLP-0478-B009}{end}

\olpseganchor{OLP-0478-B010}{start}
\item هر قضيوي متغير~\LR{$\Obj p_i$} يو اټومي فارمول دے۔
\olpseganchor{OLP-0478-B010}{end}

\olpseganchor{OLP-0478-B011}{start}
\item که \LR{$!A$} يو فارمول وي، نو \LR{$\lnot !A$} هم
  يو فارمول دے۔
\olpseganchor{OLP-0478-B011}{end}

\olpseganchor{OLP-0478-B012}{start}
\item که \LR{$!A$} او \LR{$!B$} فارمولونه وي، نو \LR{$(!A \land
  !B)$} يو فارمول دے۔
\olpseganchor{OLP-0478-B012}{end}

\olpseganchor{OLP-0478-B013}{start}
\item که \LR{$!A$} او \LR{$!B$} فارمولونه وي، نو \LR{$(!A \lor !B)$}
  يو فارمول دے۔
\olpseganchor{OLP-0478-B013}{end}

\olpseganchor{OLP-0478-B014}{start}
\item که \LR{$!A$} او \LR{$!B$} فارمولونه وي، نو \LR{$(!A \lif !B)$}
  يو فارمول دے۔
\olpseganchor{OLP-0478-B014}{end}

\olpseganchor{OLP-0478-B015}{start}
\item که \LR{$!A$} او \LR{$!B$} فارمولونه وي، نو \LR{$(!A \liff !B)$}
  يو فارمول دے۔
\olpseganchor{OLP-0478-B015}{end}

\olpseganchor{OLP-0478-B016}{start}
\item که \LR{$!A$} يو فارمول وي، نو \LR{$\Ptemp  !A$}، \LR{$\Htemp  !A$}، \LR{$\Ftemp !A$}، \LR{$\Gtemp  !A$} ټول
  فارمولونه دي۔
\olpseganchor{OLP-0478-B016}{end}

\olpseganchor{OLP-0478-B017}{start}
\item بل هېڅ شے فارمول نۀ دے۔
\end{enumerate}
\end{defn}
\olpseganchor{OLP-0478-B017}{end}

\olpseganchor{OLP-0478-B018}{start}
د هغو نورو منطقونو په څېر چې د يوې دعوې د ارزولو لپاره نور
اړوند حالتونه هم ګوري، د زماني منطقونو معناپوهنه په اړيکيزو
مدلونو ولاړه ده۔
\olpseganchor{OLP-0478-B018}{end}

\olpseganchor{OLP-0478-B019}{start}
\begin{defn}
  د زماني ژبې يو \emph{مدل} درې‌ګونې
  \LR{$\mModel{M} = \tuple{T, \prec, V}$} ده، چې په دې کښې
  \begin{enumerate}
  \item \LR{$T$} يو ناتش سټ دے، چې غړي ئې د وخت ټکي دي۔
  \item \LR{$\prec$} پر~\LR{$T$} يوه دوه‌ځاييزه اړيکه ده۔
  \item \LR{$V$} داسې تابع ده چې هر بياني متغير~\LR{$p$} ته د وخت د ټکو يو سټ \LR{$V(p)$} ټاکي۔
  \end{enumerate}
  که \LR{$t \prec t'$} وي، نو وايو چې \LR{$t$} له~\LR{$t'$} څخه \emph{مخکښې راځي}۔
  که \LR{$t \in V(p)$} وي، نو وايو چې \LR{$p$} په~\LR{$t$} کښې \emph{رښتيا دے}۔
\end{defn}
\olpseganchor{OLP-0478-B019}{end}

\olpseganchor{OLP-0478-B020}{start}
تر دې ځايه موږ د مخکښې‌والي پر اړيکه~\LR{$\prec$} هېڅ شرط نۀ دے ايښے۔
يعنې اوس د زماني مدلونو پر جوړښت هېڅ محدوديت نشته: وخت پکښې خطي،
څوڅانګيز، دوراني يا د بل هر ډول جوړښت لرلے شي۔
\olpseganchor{OLP-0478-B020}{end}

\olpseganchor{OLP-0478-B021}{start}
د عادي مودالي منطق په څېر، هر زماني مدل ټاکي چې کوم فارمولونه
ئې په کومو ټکو کښې رښتيا دي۔ د اړيکيزې معناپوهنې د بنسټيز مفهوم
دپاره هماغه نښه کاروو: «مدل~\LR{$\mModel{M}$} د ټکي~\LR{$t$} په حالت کښې
فارمول~\LR{$!A$} رښتيا کوي»۔ د رښتياوالي دا اړيکه په استقرايي ډول
تعريفېږي او د ټولو نا‌مودالي عاملونو دپاره له عادي مودالي حالت سره
يو شان ده۔
\olpseganchor{OLP-0478-B021}{end}

\olpseganchor{OLP-0478-B022}{start}
\begin{defn}\ollabel{defn:tmodels}
  په~\LR{$\mModel M$} کښې د~\LR{$t$} په حالت کښې \emph{د فارمول~\LR{$!A$}
  رښتياوالی}، چې په نښو کښې ئې \LR{$\mSat{M}{!A}[t]$} ليکو، په استقرايي
  ډول داسې تعريفېږي:
  \begin{enumerate}
  \item \indcase{!A}{\lfalse}{\LR{$\mSat{M}{\lfalse}[t]$} هېڅکله نۀ وي}۔
  \item \indcase{!A}{\ltrue}{\LR{$\mSat{M}{\ltrue}[t]$} تل وي}۔
  \item \LR{$\mSat{M}{p}[t]$} هله او يوازې هله چې \LR{$t \in V(p)$} وي۔
  \item \indcase{!A}{\lnot !B}{\LR{$\mSat{M}{\indfrm}[t]$} هله او يوازې هله چې
    \LR{$\mSat/{M}{!B}[t]$} وي}۔
  \item \indcase{!A}{(!B \land !C)}{\LR{$\mSat{M}{\indfrm}[t]$} هله او يوازې هله چې
    \LR{$\mSat{M}{!B}[t]$} او \LR{$\mSat{M}{!C}[t]$} دواړه وي}۔
  \item \indcase{!A}{(!B \lor !C)}{\LR{$\mSat{M}{\indfrm}[t]$} هله او يوازې هله چې
    \LR{$\mSat{M}{!B}[t]$} يا \LR{$\mSat{M}{!C}[t]$} وي} (يا دواړه)۔
  \item \indcase{!A}{(!B \lif !C)}{\LR{$\mSat{M}{\indfrm}[t]$} هله او يوازې هله چې
    \LR{$\mSat/{M}{!B}[t]$} يا \LR{$\mSat{M}{!C}[t]$} وي}۔
  \item \indcase{!A}{(!B \liff !C)}{\LR{$\mSat{M}{\indfrm}[t]$} هله او يوازې هله چې
    يا \LR{$\mSat{M}{!B}[t]$} او \LR{$\mSat{M}{!C}[t]$} دواړه وي، يا
    نۀ \LR{$\mSat{M}{!B}[t]$} وي او نۀ \LR{$\mSat{M}{!C}[t]$}}۔
  \item\ollabel{defn:sub:mmodels-p}
    \indcase{!A}{\Ptemp  !B}{\LR{$\mSat{M}{\indfrm}[t]$} هله او يوازې هله چې
    د کوم \LR{$t' \in T$} دپاره چې \LR{$t' \prec t$} وي، \LR{$\mSat{M}{!B}[t']$} وي}
\item\ollabel{defn:sub:mmodels-h}
    \indcase{!A}{\Htemp  !B}{\LR{$\mSat{M}{\indfrm}[t]$} هله او يوازې هله چې
    د هر \LR{$t' \in T$} دپاره چې \LR{$t' \prec t$} وي، \LR{$\mSat{M}{!B}[t']$} وي}
   \item\ollabel{defn:sub:mmodels-f}
    \indcase{!A}{\Ftemp  !B}{\LR{$\mSat{M}{\indfrm}[t]$} هله او يوازې هله چې
    د کوم \LR{$t' \in T$} دپاره چې \LR{$t \prec t'$} وي، \LR{$\mSat{M}{!B}[t']$} وي}
\item\ollabel{defn:sub:mmodels-g}
    \indcase{!A}{\Gtemp  !B}{\LR{$\mSat{M}{\indfrm}[t]$} هله او يوازې هله چې
    د هر \LR{$t' \in T$} دپاره چې \LR{$t \prec t'$} وي، \LR{$\mSat{M}{!B}[t']$} وي}
  \end{enumerate}
\end{defn}
\olpseganchor{OLP-0478-B022}{end}

\olpseganchor{OLP-0478-B023}{start}
له دې معناپوهنې څخه څرګندېږي چې \LR{$\Ptemp$} او~\LR{$\Htemp$} دوه‌ګوني
عاملونه دي؛ همداسې \LR{$\Ftemp$} او \LR{$\Gtemp$} هم دوه‌ګوني دي۔ نو
\LR{$\Htemp  !A$} د \LR{$\lnot \Ptemp \lnot !A$} په توګه تعريفولے شو،
او همدا خبره د \LR{$\Gtemp$} او~\LR{$\Ftemp$} په اړه هم کېږي۔
\olpseganchor{OLP-0478-B023}{end}

\olpseganchor{OLP-0479-B004}{start}
\olfileid{aml}{tl}{acc}
\olpseganchor{OLP-0479-B004}{end}

\olpseganchor{OLP-0479-B005}{start}
\olsection{د زماني چوکاټونو خاصيتونه}
\olpseganchor{OLP-0479-B005}{end}

\olpseganchor{OLP-0479-B006}{start}
زموږ زماني مدلونه پر اړيکه~\LR{$\prec$} هېڅ شرط نۀ ږدي۔ د يوه
لوري لپاره د عادي مودالي اکسيومونو له ډلې~\LR{$K$} په ټولو مدلونو
کښې سم دے؛ د راتلونکي او تېر لوري بڼې ئې \LR{$K_{\Gtemp}$} او
~\LR{$K_{\Htemp}$} دي۔ د دواړو لورو د يوې اړيکې په توګه کارول نورې
ګډې معتبرې طرحې هم ورکوي: د هرې راتلونکې شېبې په تېر کښې
اوسنۍ شېبه شته، او د هرې تېرې شېبې په راتلونکي کښې اوس شته۔
\begin{align*}
\tag{\LR{$K_{\Gtemp}$}}  \Gtemp (p \to q) & \to (\Gtemp p \to \Gtemp q)\\
\tag{\LR{$K_{\Htemp}$}}  \Htemp (p \to q) & \to (\Htemp p \to \Htemp q)
\end{align*}
\olpseganchor{OLP-0479-B006}{end}

\olpseganchor{OLP-0479-B007}{start}
خو که غواړو زموږ مدلونه د وخت د ښودلو په طريقه سخت شرطونه ولري،
لکه دا چې~\LR{$\prec$} يو خطي ترتيب وي، نو په مدلونو کښې به نورې
معتبرې طرحې هم ولرو۔
\olpseganchor{OLP-0479-B007}{end}

\olpseganchor{OLP-0479-B008}{start}
\begin{table}[t]
  \begin{tabular}{| p{\dimexpr .5\linewidth-2\tabcolsep-2\arrayrulewidth-\doublerulesep/2\relax} || p{\dimexpr .5\linewidth-2\tabcolsep-2\arrayrulewidth-\doublerulesep/2\relax} |}
    \hline
    {\emph{که \LR{$\prec$} \dots وي}} & {\emph{نو \dots په~\LR{$\mModel{M}$} کښې رښتيا دے:}} \\
    \hline \hline
    \emph{انتقالي}: \newline
    \LR{$\forall u \forall v \forall w ((u \prec v \land v \prec w) \lif u \prec w)$} &
    \LR{$\Ftemp \Ftemp p \lif \Ftemp p$}  \\
    \hline
    \emph{د هرې جوړې پرتله کېدنه}: \newline
    \LR{$\forall w \forall v (w \prec v \lor w = v \lor v \prec w)$} &
    \LR{$(\Ftemp \Ptemp  p \lor \Ptemp \Ftemp p) \lif (\Ptemp p \lor p \lor \Ftemp p)$}\\
    \hline
    \emph{ګڼه}: \newline
    \LR{$\forall w \forall v (w \prec v \to \exists u(w \prec u \land u \prec v))$} &
    \LR{$\Ftemp p \lif \Ftemp \Ftemp p$} \\
    \hline
    \emph{په تېر لوري بې سرحده}: \newline
    \LR{$\forall w \exists v( v \prec w)$} &
    \LR{$\Htemp p \to \Ptemp p$} \\
    \hline
    \emph{په راتلونکي لوري بې سرحده}: \newline
    \LR{$\forall w \exists v( w \prec v)$} &
    \LR{$\Gtemp p \to \Ftemp p$} \\
    \hline
  \end{tabular}
  \caption{د زماني چوکاټونو د مطابقت ځينې خاصيتونه۔}
  \ollabel{tab:correspondence}
\end{table}
\olpseganchor{OLP-0479-B008}{end}

\olpseganchor{OLP-0479-B009}{start}
د \olref{tab:correspondence} ځينې خاصيتونه ښايي د وخت د ښودونکي
مدل لپاره مطلوب ښکاره شي۔ د جدول طرحه دا غواړي چې هره جوړه چې
يوه ګډه مخکنۍ يا وروستۍ شېبه لري، سره د پرتله کېدو وړ وي۔
مثلاً په دوو جلا شېبو او تشې اړيکې کښې طرحه معتبره ده،
او دواړه شېبې سره د پرتله کېدو اړيکه نۀ لري۔ خو که څۀ هم پر~\LR{$\prec$} اړيکه هر ډول
شرط ايښودلے شو، ټول شرطونه د زماني منطق په ژبه کښې د څرګندېدونکو
فارمولونه سره نۀ برابرېږي۔ د بېلګې په توګه، نفي انعکاسيت،
يعنې دا چې \LR{$\forall w \lnot (w \prec w)$}، په همدې بنسټيزه
يو‌ځاييزه ژبه کښې د چوکاټ د اعتبار له لارې نۀ تعريفېږي۔
په يوه ځان ته رسېدونکې شېبه کښې د متغيرونو هره ارزښت‌ټاکنه
د صحيح عددونو پر هر ټکي يو شان وټاکئ او هلته کلک ترتيب واخلئ۔
د تېر او راتلونکي لوري په دواړو چوکاټونو کښې ناتش دي۔ نو په
استقرا سره د هر فارمول رښتياوالے په دواړو کښې يو شان وي۔
هره طرحه چې پر کلک ترتيب معتبره وي، پر ځان ته رسېدونکې
يوې شېبې هم معتبره ده۔
\olpseganchor{OLP-0479-B009}{end}

\olpseganchor{OLP-0480-B004}{start}
\olfileid{aml}{tl}{ext}
\olpseganchor{OLP-0480-B004}{end}

\olpseganchor{OLP-0480-B005}{start}
\olsection{د زماني منطق نور عاملونه}
\olpseganchor{OLP-0480-B005}{end}

\olpseganchor{OLP-0480-B006}{start}
د تېر او راتلونکي وخت د يو-ځاييزو عاملونو ترڅنګ، زماني منطقونه کله کله
دوه‌ځاييز عاملونه \LR{$\Since $} او~\LR{$\Until$} هم لري چې په ترتيب سره
«له هغه وخته راهيسې» او «تر هغه وخته پورې» ښيي۔ دا معنا لري چې
\LR{$\Since $} او~\LR{$\Until$} د زماني منطق ژبې ته ورزيات کړو، او د زماني
فارمول په تعريف کښې لاندې ماده ورزياته کړو:
\olpseganchor{OLP-0480-B006}{end}

\olpseganchor{OLP-0480-B007}{start}
\begin{itemize}
\item[] که \LR{$!A$} او \LR{$!B$} فارمولونه وي، نو \LR{$(\Since  !A !B)$} او
  \LR{$(\Until !A !B)$} دواړه فارمولونه دي۔
\end{itemize}
\olpseganchor{OLP-0480-B007}{end}

\olpseganchor{OLP-0480-B008}{start}
د دغو عاملونو معناپوهنه بيا داسې ورکول کېږي:
\olpseganchor{OLP-0480-B008}{end}

\olpseganchor{OLP-0480-B009}{start}
\begin{defn}\ollabel{defn:since-until}
  په~\LR{$\mModel M$} کښې د~\LR{$t$} په حالت کښې \emph{د فارمول~\LR{$!A$} رښتياوالی}:
  \begin{enumerate}
  \item\ollabel{defn:sub:mmodels-since} \indcase{!A}{\Since  !B
    !C}{\LR{$\mSat{M}{\indfrm}[t]$} هله او يوازې هله چې د کوم \LR{$t' \in
    T$} دپاره \LR{$t' \prec t$} او \LR{$\mSat{M}{!B}[t']$} وي، او د هر هغه
    \LR{$s$} دپاره چې \LR{$t' \prec s \prec t$} وي، \LR{$\mSat{M}{!C}[s]$} وي}
   \item\ollabel{defn:sub:mmodels-until} \indcase{!A}{\Until !B
    !C}{\LR{$\mSat{M}{\indfrm}[t]$} هله او يوازې هله چې د کوم \LR{$t' \in
    T$} دپاره \LR{$t \prec t'$} او \LR{$\mSat{M}{!B}[t']$} وي، او د هر هغه
    \LR{$s$} دپاره چې \LR{$t \prec s \prec t'$} وي، \LR{$\mSat{M}{!C}[s]$} وي}
  \end{enumerate}
\end{defn}
\olpseganchor{OLP-0480-B009}{end}

\olpseganchor{OLP-0480-B010}{start}
د \LR{$\Since  !B !C$} د پوهېدو عادي لاره دا ده: «له هغه وخته چې
\LR{$!B$} واقع شوے، \LR{$!C$} روان پاتې شوے دے»۔ د \LR{$\Until !B !C$}
عادي لوست بيا دا دے: «تر هغه وخته چې \LR{$!B$} به واقع شي، \LR{$!C$}
به روان وي»۔ په دواړو تعريفونو کښې يوه داسې ټاکلې شېبه شته چې
لومړے فارمول پکښې رښتينے وي۔ د دويم فارمول شرط د هماغې
ټاکلې شېبې او اوسنۍ شېبې تر منځ د مخکښوالي له دواړو شرطونو
سره سم پلي کېږي۔ که اړيکه نفي انعکاسي وي، دواړه سرونه له
دې شرطه بهر دي۔ که له سرونو څخه يو له خپله ځانه مخکښ وي،
نو دويم فارمول بايد په هغه سر کښې هم رښتينے وي۔
\olpseganchor{OLP-0480-B010}{end}

\olpseganchor{OLP-0481-B004}{start}
\olfileid{aml}{tl}{poss}
\olpseganchor{OLP-0481-B004}{end}

\olpseganchor{OLP-0481-B005}{start}
\olsection{ممکنې تاريخچې}
\olpseganchor{OLP-0481-B005}{end}

\olpseganchor{OLP-0481-B006}{start}
د زماني منطق هغه اړيکيز مدلونه چې تر اوسه مو کارولي، ډېر انعطاف لري؛
ځکه موږ پر لاسرسي اړيکه هېڅ قيد نۀ ږدو۔ نو زماني مدلونه په تېر او
راتلونکي دواړو لوريو کښې څانګې لرلے شي۔ خو ښايي وغواړو د څانګې يو
لا «مودالي» تصور وڅېړو چې پکښې د پېښو لړۍ د ممکنو تاريخچو په توګه
ګڼو۔ د دې دپاره د ژبې بدلول لازمي نۀ دي، که څۀ هم خپل «عادي» مودالي
عاملونه \LR{$\Box$} او~\LR{$\Diamond$} هم ورزياتولے شو؛ همدارنګه د پوهېدنې
د لاسرسي اړيکې ورزياتولے شو، څو د وخت په تېرېدو سره د عاملانو د
پوهې بدلون وښيو۔
\olpseganchor{OLP-0481-B006}{end}

\olpseganchor{OLP-0481-B007}{start}
\begin{defn}
  د زماني ژبې يو \emph{د ممکنو تاريخچو مدل} درې‌ګونې
  \LR{$\mModel{M} = \tuple{T, C, V}$} ده، چې په دې کښې
  \begin{enumerate}
    \item \LR{$T$} يو ناتش سټ دے، چې غړي ئې د وخت حالتونه دي۔
    \item \LR{$C$} د يوه نظام د محاسبوي لارو، يا \emph{ممکنو
      تاريخچو}، سټ دے۔ په بله وينا، \LR{$C$} د حالتونو \LR{$s_1$}، \LR{$s_2$}،
      ~\LR{$s_3$}، \dots د لړيو~\LR{$\sigma$} سټ دے، چې هر \LR{$s_i \in T$} وي۔
      تاريخچې متناهي يا نامتناهي ناتشې لړۍ دي۔ په هره تاريخچه
      کښې د وخت يو ټاکلے حالت يوازې په يو ځاے راځي۔ د يوه
      محاسبوي حالت د تکرار بېلابېلې شېبې جلا وختي حالتونه دي۔
    \item \LR{$V$} داسې تابع ده چې هر بياني متغير~\LR{$p$} ته د وخت د ټکو يو سټ~\LR{$V(p)$} ټاکي۔
  \end{enumerate}
  د اسانتيا دپاره عموماً دا هم فرضوو چې که يوه تاريخچه په~\LR{$C$} کښې
  وي، نو د هغې ټولې وروستۍ برخې هم پکښې وي۔ د بېلګې په توګه،
  که \LR{$s_1$}، \LR{$s_2$}،~\LR{$s_3$} په~\LR{$C$} کښې يوه لړۍ وي، نو د \LR{$s_2$} او \LR{$s_3$}
  دوه‌غړې لړۍ او د~\LR{$s_3$} يوغړې لړۍ هم پکښې دي۔ همدارنګه، که دوه حالتونه \LR{$s_i$} او~\LR{$s_j$}
  په لړۍ~\LR{$\sigma$} کښې راشي، نو د \LR{$i < j$} په حالت کښې وايو چې
  \LR{$s_i \prec_\sigma s_j$}۔ که \LR{$t \in V(p)$} وي، نو وايو چې \LR{$p$}
  په~\LR{$t$} کښې \emph{رښتيا دے}۔
\end{defn}
\olpseganchor{OLP-0481-B007}{end}

\olpseganchor{OLP-0481-B008}{start}
اړوند بدلون دا دے چې په مدل~\LR{$\mModel M$} کښې د وخت په ټکي~\LR{$t$}
د فارمول د رښتياوالي ارزونه د تاريخچې~\LR{$\sigma$} په نسبت کوو، چې
\LR{$t$} پکښې يو حالت وي۔ د بياني متغيرونه او د
رښتياوالي تابع نښلوونکو لپاره د \olref[sem]{defn:tmodels}
قاعدې کاروو۔ د متغير خپل ارزښت له~\LR{$\sigma$} سره نۀ بدلېږي۔
په رښتياوالي تابع نښلوونکو کښې د فرعي فارمولونو صدق په
هماغې ټاکلې تاريخچه کښې اخلو۔ اوس د راتلونکي عامل~\LR{$\Ftemp$}
معنا بيا تعريفوو او د دې تاريخچو په نسبت مودالي عامل~\LR{$\Diamond$}
ورزياتوو۔
\olpseganchor{OLP-0481-B008}{end}

\olpseganchor{OLP-0481-B009}{start}
\begin{defn}\ollabel{defn:phmodels}
  په~\LR{$\mModel M$} کښې په~\LR{$t, \sigma$} باندې \emph{د فارمول~\LR{$!A$}
  رښتياوالے}، چې په نښو کښې ئې \LR{$\mSat{M}{!A}[t, \sigma]$} ليکو:
  \begin{enumerate}
  \item\ollabel{defn:sub:phmodels-f} \indcase{!A}{\Ftemp
    !B}{\LR{$\mSat{M}{\indfrm}[t, \sigma]$} هله او يوازې هله چې د کوم
    \LR{$t' \in T$} دپاره \LR{$t \prec_\sigma t'$} او
    \LR{$\mSat{M}{!B}[t', \sigma]$} وي۔}
  \item\ollabel{defn:sub:phmodels-diamond} \indcase{!A}{\Diamond
    !B}{\LR{$\mSat{M}{\indfrm}[t, \sigma]$} هله او يوازې هله چې د
    کومې \LR{$\sigma' \in C$} دپاره چې \LR{$t$} پکښې راځي،
    \LR{$\mSat{M}{!B}[t, \sigma']$} وي۔}
  \end{enumerate}
\end{defn}
\olpseganchor{OLP-0481-B009}{end}

\olpseganchor{OLP-0481-B010}{start}
نور زماني او مودالي عاملونه هم په ورته ډول تعريفولے شو۔ اوس بيا هغه
دعوې هم څرګندولے شو چې زمان او موداليت سره يوځاے کوي۔ د بېلګې په
توګه، «\LR{$p$} به واقع نۀ شي، خو په بله ممکنه تاريخچه کښې ښايي واقع شي»
په \LR{$\lnot \Ftemp p \land \Diamond \Ftemp p$} ښودلے شو۔ دا په هغه
ټکي او تاريخچه کښې رښتيا وي چې \LR{$p$} پکښې په راتلونکي حالت کښې رښتيا
نۀ شي، او له همدې ټکي تېرېدونکې بله ممکنه تاريخچه شته چې \LR{$p$}
به پکښې په راتلونکي حالت کښې رښتيا شي۔
\olpseganchor{OLP-0481-B010}{end}

\olpseganchor{OLP-0482-B004}{start}
\olchapter{aml}{el}{معرفتي منطقونه}
\olpseganchor{OLP-0482-B004}{end}

\olpseganchor{OLP-0482-B005}{start}
\begin{editorial}
  دا باب د معرفتي منطقونو فرا تيوري څېړي۔ د باب جوړښت د کلاسيک
  بنسټيز مودالي منطق په اړه د الدو انتونيلي \mbox{\LR{\latinfont (Aldo Antonelli)}} د يادښتونو
  جوړښت ته ورته دے۔ آډري ياپ \mbox{\LR{\latinfont (Audrey Yap)}} باب له سره ليکلے او
  د باي‌سيميولېشن او خوځنده معرفتي منطقونو بحثونه ئې ورزيات کړي دي۔
\end{editorial}
\olpseganchor{OLP-0482-B005}{end}

\olpseganchor{OLP-0483-B004}{start}
\olfileid{aml}{el}{int}
\olpseganchor{OLP-0483-B004}{end}

\olpseganchor{OLP-0483-B005}{start}
\olsection{سريزه}
\olpseganchor{OLP-0483-B005}{end}

\olpseganchor{OLP-0483-B006}{start}
لکه څنګه چې مودالي منطق له \emph{مودالي قضيو} او د هغو تر منځ د لازم‌والي
له اړيکو سره کار لري، معرفتي منطق له \emph{معرفتي قضيو} او د هغو تر منځ
د لازم‌والي له اړيکو سره کار لري۔ په عادي مودالي لوست کښې
يو‌ځاييز عملګر د امکان او ضرورت معنا ورکوي۔ په معرفتي لوست
کښې دغه عملګر د عامل پوهه يا باور ښيي۔ د بېلګې په توګه، ښايي لاندې دعوې څرګندول وغواړو:
\olpseganchor{OLP-0483-B006}{end}

\olpseganchor{OLP-0483-B007}{start}
\begin{enumerate}
\item ريچارډ پوهېږي چې کالګري په البرټا کښې ده۔
\item آډري فکر کوي چې شونې ده يو سپے پر صوفه وي۔
\item ريچارډ پوهېږي چې آډري پوهېږي د هغې ټولګے د نهې په ورځو وي۔
\item هر څوک پوهېږي چې يو کال ۱۲ مياشتې لري۔
\end{enumerate}
د اوسني معرفتي منطق پيل ډېر ځله د ياکو هينتيکا \mbox{\LR{\latinfont (Jaakko Hintikka)}}
د ۱۹۶۲ کال کتاب \emph{پوهه او باور} \mbox{\LR{\latinfont (Knowledge and Belief)}} ته ورګرځول
کېږي۔ په هغه وخت کښې په منطق کښې د ممکنو نړۍو معناپوهنه ورځ تر بلې
ډېره کارېده۔ په حقيقت کښې معرفتي منطقونه تر ډېره د نورو مودالي
منطقونو هماغه معنايي وسيلې کاروي، خو بل ډول ئې تفسيروي۔ اصلي بدلون
په دې کښې دے چې \emph{د لاسرسي اړيکه} د څۀ نماينده ګڼو۔ په معرفتي
منطقونو کښې دا اړيکه د \emph{معرفتي امکان} يوه بڼه ښيي۔ وبه ګورو چې
مدل کېدونکے معرفتي تصور د لاسرسي پر اړيکه زموږ د مطلوبو شرطونو په
ټاکلو څۀ اغېز کوي۔ همدارنګه به وګورو چې د مطابقت تيوري ته د معرفتي
تفسير په ورکولو سره څۀ پېښېږي۔ په پورته بېلګو کښې دوه عاملان، ريچارډ
او آډري، ياد شوي، او د دوى د پوهې د دعوو تړاو هم څرګندېږي۔ هغه معرفتي منطقونه چې دلته ئې څېړو څو-عامله منطقونه به وي؛
په هغو کښې دا ډول خبرې څرګندېدے شي۔ برعکس، يو-عامله معرفتي منطق
يوازې د يوه کس د پوهې يا باور په اړه خبرې کوي۔
\olpseganchor{OLP-0483-B007}{end}

\olpseganchor{OLP-0484-B004}{start}
\olfileid{aml}{el}{lan}
\olpseganchor{OLP-0484-B004}{end}

\olpseganchor{OLP-0484-B005}{start}
\olsection{د معرفتي منطق ژبه}
\olpseganchor{OLP-0484-B005}{end}

\olpseganchor{OLP-0484-B006}{start}
\begin{defn}
فرض کړئ \LR{$G$} د عاملانو د نښو يو سټ وي۔ د څو-عامله معرفتي منطق
بنسټيزه ژبه دا توکي لري:
\begin{enumerate}
  \item د دروغوالی قضيوي ثابت~\LR{$\lfalse$}۔
  \item د رښتياوالی قضيوي ثابت~\LR{$\ltrue$}۔
  \item د بياني متغيرونه يو شمېرېدونکے لامتناهي سټ: \LR{$\Obj
    p_0$}، \LR{$\Obj p_1$}، \LR{$\Obj p_2$}، \dots
  \item قضيوي نښلوونکي: \startycommalist
  \ycomma \LR{$\lnot$} (نفي)%
  \ycomma \LR{$\land$} (عطف)%
  \ycomma \LR{$\lor$} (فصل)%
  \ycomma \LR{$\lif$} (شرطي)%
  \ycomma \LR{$\liff$} (دوه اړخيز شرطي).
  \item د پوهې عامل \LR{$\Knows_a$}، چې \LR{$a \in G$} وي۔
\end{enumerate}
\end{defn}
\olpseganchor{OLP-0484-B006}{end}

\olpseganchor{OLP-0484-B007}{start}
که په خپل نظام کښې يوازې د يوه عامل پوهه څېړو، د سټ~\LR{$G$} او د
جلا جلا عاملانو يادونه پرېښودلے شو۔ په دې حالت کښې يوازې بنسټيز
عامل~\LR{$\Knows$} لرو۔
\olpseganchor{OLP-0484-B007}{end}

\olpseganchor{OLP-0484-B008}{start}
\begin{defn}
د معرفتي ژبې \emph{فارمولونه} په استقرايي ډول داسې
  تعريفېږي:
\begin{enumerate}
\item \LR{$\lfalse$} يو اټومي فارمول دے۔
\olpseganchor{OLP-0484-B008}{end}

\olpseganchor{OLP-0484-B009}{start}
\item \LR{$\ltrue$} يو اټومي فارمول دے۔
\olpseganchor{OLP-0484-B009}{end}

\olpseganchor{OLP-0484-B010}{start}
\item هر قضيوي متحول~\LR{$\Obj p_i$} يو اټومي فارمول دے۔
\olpseganchor{OLP-0484-B010}{end}

\olpseganchor{OLP-0484-B011}{start}
\item که \LR{$!A$} يو فارمول وي، نو \LR{$\lnot !A$} هم
  يو فارمول دے۔
\olpseganchor{OLP-0484-B011}{end}

\olpseganchor{OLP-0484-B012}{start}
\item که \LR{$!A$} او \LR{$!B$} فارمولونه وي، نو \LR{$(!A \land
  !B)$} يو فارمول دے۔
\olpseganchor{OLP-0484-B012}{end}

\olpseganchor{OLP-0484-B013}{start}
\item که \LR{$!A$} او \LR{$!B$} فارمولونه وي، نو \LR{$(!A \lor !B)$}
  يو فارمول دے۔
\olpseganchor{OLP-0484-B013}{end}

\olpseganchor{OLP-0484-B014}{start}
\item که \LR{$!A$} او \LR{$!B$} فارمولونه وي، نو \LR{$(!A \lif !B)$}
  يو فارمول دے۔
\olpseganchor{OLP-0484-B014}{end}

\olpseganchor{OLP-0484-B015}{start}
\item که \LR{$!A$} او \LR{$!B$} فارمولونه وي، نو \LR{$(!A \liff !B)$}
  يو فارمول دے۔
\olpseganchor{OLP-0484-B015}{end}

\olpseganchor{OLP-0484-B016}{start}
\item که \LR{$!A$} يو فارمول وي او \LR{$a \in G$} وي، نو \LR{$\Knows_a !A$}
  هم يو فارمول دے۔
\olpseganchor{OLP-0484-B016}{end}

\olpseganchor{OLP-0484-B017}{start}
\item بل هېڅ شے فارمول نۀ دے۔
\end{enumerate}
\end{defn}
\olpseganchor{OLP-0484-B017}{end}

\olpseganchor{OLP-0484-B018}{start}
که فارمول~\LR{$!A$} کښې د هېڅ عامل د پوهې عملګر \LR{$\Knows_a$} رانشي،
هغه \emph{له موداليت پرته} بولو۔
\olpseganchor{OLP-0484-B018}{end}

\olpseganchor{OLP-0484-B019}{start}
\begin{defn}
  د \LR{$\Knows$} عامل د يوه کس د پوهې د ښودلو لپاره دے؛ \LR{$\EKnows$}،
  چې ډېر ځله «هر څوک پوهېږي» لوستل کېږي، د ډلې د هر غړي پوهه ښيي۔
  که \LR{$G' \subseteq G$} يوه متناهي ډله وي، نو \LR{$\EKnows_{G'} !A$} د
  \[\bigwedge_{b \in G'} \Knows_b !A.\] لنډه نښه تعريفوو۔
  تش عطف تل رښتينے دے۔ که ډله نامتناهي وي، د ډلې د هر غړي لپاره
  د فردي پوهې شرط اخلو؛ دا په بنسټيزه ژبه کښې د متناهي عطف لنډون
  نۀ دے۔
\end{defn}
\olpseganchor{OLP-0484-B019}{end}

\olpseganchor{OLP-0484-B020}{start}
د عاملانو د ډلې~\LR{$G$} په منځ کښې د پوهې يوه لا پياوړې معنا هم
تعريفولے شو، يعنې \emph{مشترکه پوهه}۔ که د معلوماتو يوه برخه د
ډلې مشترکه پوهه وي، نو د عاملانو د هرې ناتشې متناهي لړۍ لپاره
د لومړي عامل پوهه اخلو چې دويم پوهېږي، هغه بيا پوهېږي چې درېيم
پوهېږي، \dots تر وروستي عامل پورې چې پر هماغو معلوماتو پوهېږي۔
په دې لړيو کښې عامل تکرارېدے هم شي۔ مشترکه پوهه د ډلې د هر غړي
پوهه غواړي۔ داسې اړيکيز مدلونه شته چې پکښې يو فارمول د ډلې
د هر غړي پوهه وي، خو مشترکه پوهه نۀ وي۔ \LR{$\CKnows_G !A$} به د دې
ښودلو لپاره کاروو چې «په~\LR{$G$} کښې دا مشترکه پوهه ده چې~\LR{$!A$}»۔
\olpseganchor{OLP-0484-B020}{end}

\olpseganchor{OLP-0485-B004}{start}
\olfileid{aml}{el}{rel}
\olpseganchor{OLP-0485-B004}{end}

\olpseganchor{OLP-0485-B005}{start}
\olsection{اړيکيز مدلونه}
\olpseganchor{OLP-0485-B005}{end}

\olpseganchor{OLP-0485-B006}{start}
د معرفتي منطقونو بنسټيز معنايي مفهوم د عادي مودالي منطقونو له مفهوم سره
يو شان دے۔ اړيکيز مدلونه لا هم د نړۍو له يوه سټ او داسې ټاکنې څخه جوړ
دي چې معلوموي کوم بياني متغيرونه په کومو نړۍو کښې «رښتيا» ګڼل کېږي۔ که يوازې له يوه عامل
سره کار لرو، د معمول په څېر د لاسرسي يوه اړيکه لرو۔ خو د څو-عامله
معرفتي منطق لپاره د لاسرسي اړيکو داسې کورنۍ کاروو چې
د عامل‌نښو په سټ~\LR{$G$} کښې د هر~\LR{$a$} لپاره يوه اړيکه لري۔
\olpseganchor{OLP-0485-B006}{end}

\olpseganchor{OLP-0485-B007}{start}
يو \emph{اړيکيز مدل} د نړۍو له يوه سټ، د هغو تر منځ له دوه‌ځاييزو
لاسرسي اړيکو—د هر عامل لپاره له يوې—او داسې ټاکنې څخه جوړ دے چې
معلوموي کوم بياني متغيرونه په کومو نړۍو کښې رښتيا دي۔
\olpseganchor{OLP-0485-B007}{end}

\olpseganchor{OLP-0485-B008}{start}
\begin{defn}
  د څو-عامله معرفتي ژبې يو \emph{مدل} درې‌ګونې
  \LR{$\mModel{M} = \tuple{W, R, V}$} ده، چې په دې کښې
  \begin{enumerate}
  \item \LR{$W$} د «نړۍو» يو ناتش سټ دے،
  \item د هر \LR{$a \in G$} لپاره \LR{${R}_a$} پر~\LR{$W$} د لاسرسي دوه‌ځاييزه
  اړيکه ده؛ د دوو عاملانو اړيکې سره برابرېدې هم شي، او
  \item \LR{$V$} داسې تابع ده چې هر بياني متغير~\LR{$p$} ته د ممکنو نړۍو يو سټ \LR{$V(p)$} ټاکي۔
  \end{enumerate}
  که \LR{$R_a ww'$} وي، نو وايو چې \LR{$w'$} د همدغه عامل لپاره
  له~\LR{$w$} څخه \emph{د لاسرسي وړ} ده۔ که \LR{$w \in V(p)$} وي، نو وايو چې
  \LR{$p$} په~\LR{$w$} کښې \emph{رښتيا دے}۔
\end{defn}
\olpseganchor{OLP-0485-B008}{end}

\olpseganchor{OLP-0485-B009}{start}
چاره د نورمال مودالي منطق د چارې په شان ده، يوازې د لاسرسي اړيکې
پکښې ډېرې دي۔ د يوه ټاکلي عامل د لاسرسي اړيکه عموماً د هغه د
معلوماتي حالتونو د ښودلو په توګه تفسيروو۔ د بېلګې په توګه،
\LR{$R_a ww'$} ډېر ځله داسې لولو چې \LR{$w'$} په~\LR{$w$} کښې د~\LR{$a$} له
معلوماتو سره سازګاره ده۔ يا په بله وينا، په~\LR{$w$} کښې عامل د
نړۍ~\LR{$w$} او نړۍ~\LR{$w'$} تر منځ توپير نۀ شي پېژندلے۔
\olpseganchor{OLP-0485-B009}{end}

\olpseganchor{OLP-0486-B004}{start}
\olfileid{aml}{el}{trw}
\olpseganchor{OLP-0486-B004}{end}

\olpseganchor{OLP-0486-B005}{start}
\olsection{په يوه نړۍ کښې رښتياوالی}
\olpseganchor{OLP-0486-B005}{end}

\olpseganchor{OLP-0486-B006}{start}
لکه د نورمال مودالي منطق په شان، هر معرفتي مدل ټاکي چې کوم فارمولونه
ئې په کومو نړۍو کښې رښتيا ګڼل کېږي۔ د اړيکيزې معناپوهنې د بنسټيز
مفهوم لپاره هماغه نښه کاروو: «مدل~\LR{$\mModel{M}$} په نړۍ~\LR{$w$} کښې
فارمول~\LR{$!A$} رښتيا کوي»۔ دا اړيکه په استقرايي ډول تعريفېږي، او د
ټولو غيرمودالي عاملونو لپاره ئې تعريف د نورمال مودالي حالت په شان عين دے۔
\olpseganchor{OLP-0486-B006}{end}

\olpseganchor{OLP-0486-B007}{start}
\begin{defn}\ollabel{defn:mmodels}
  په~\LR{$\mModel M$} کښې \emph{په~\LR{$w$} د فارمول~\LR{$!A$} رښتياوالی}،
  چې په نښو کښې ئې \LR{$\mSat{M}{!A}[w]$} ليکو، په استقرايي ډول داسې تعريفېږي:
  \begin{enumerate}
  \item \indcase{!A}{\lfalse}{هېڅکله \LR{$\mSat{M}{\lfalse}[w]$} نۀ وي}.
  \item \indcase{!A}{\ltrue}{تل \LR{$\mSat{M}{\ltrue}[w]$} وي}.
  \item \LR{$\mSat{M}{p}[w]$} هله او يوازې هله چې \LR{$w \in V(p)$} وي۔
  \item \indcase{!A}{\lnot !B}{\LR{$\mSat{M}{\indfrm}[w]$} هله او يوازې هله چې
    \LR{$\mSat/{M}{!B}[w]$} وي}.
  \item \indcase{!A}{(!B \land !C)}{\LR{$\mSat{M}{\indfrm}[w]$} هله او يوازې هله چې
    \LR{$\mSat{M}{!B}[w]$} او \LR{$\mSat{M}{!C}[w]$} دواړه وي}.
  \item \indcase{!A}{(!B \lor !C)}{\LR{$\mSat{M}{\indfrm}[w]$} هله او يوازې هله چې
    \LR{$\mSat{M}{!B}[w]$} يا \LR{$\mSat{M}{!C}[w]$} وي} (يا دواړه)۔
  \item \indcase{!A}{(!B \lif !C)}{\LR{$\mSat{M}{\indfrm}[w]$} هله او يوازې هله چې
    \LR{$\mSat/{M}{!B}[w]$} يا \LR{$\mSat{M}{!C}[w]$} وي}.
  \item \indcase{!A}{(!B \liff !C)}{\LR{$\mSat{M}{\indfrm}[w]$} هله او يوازې هله چې
    يا دواړه \LR{$\mSat{M}{!B}[w]$} او \LR{$\mSat{M}{!C}[w]$} وي، يا
    نۀ \LR{$\mSat{M}{!B}[w]$} وي او نۀ \LR{$\mSat{M}{!C}[w]$} وي}.
  \item\ollabel{defn:sub:mmodels-box}
    \indcase{!A}{\Knows_a !B}{\LR{$\mSat{M}{\indfrm}[w]$} هله او يوازې هله چې
    د هر هغه \LR{$w' \in W$} لپاره \LR{$\mSat{M}{!B}[w']$} وي چې \LR{$R_a ww'$} وي}
  \end{enumerate}
\end{defn}
\olpseganchor{OLP-0486-B007}{end}

\olpseganchor{OLP-0486-B008}{start}
خو دلته د لاسرسي پر اړيکو د قيدونو په اړه فکر کول پکار دي۔ د
\olref{defn:sub:mmodels-box} د مادې له مخې، که هېڅ داسې~\LR{$w'$}
نۀ وي چې \LR{$R_a ww'$} وي، نو فارمول~\LR{$\Knows_a !B$} په~\LR{$w$}
کښې رښتيا دے۔ په نورمال مودالي منطق کښې هم که له يوې نړۍ هېڅ
نړۍ د لاسرسي وړ نۀ وي، نو د ټولو \LR{$\Box$}-فارمولونو عمومي شرط
پکښې په تش ډول پوره کېږي۔ که \LR{$\Knows$} د \emph{پوهې} د ښودلو عامل وګڼو، دا ډېره ناعادي
پايله ښکاري۔ ځکه موږ عموماً فکر کوو چې داسې \emph{هېڅ} حالت نشته
چې يو عامل په يوه وخت کښې پر \LR{$!A$} هم پوه شي او پر \LR{$\lnot !A$} هم۔
\olpseganchor{OLP-0486-B008}{end}

\olpseganchor{OLP-0486-B009}{start}
يوه لاره دا ده چې په معرفتي منطق کښې د لاسرسي اړيکه تل
\emph{انعکاسي} وټاکو۔ دا په نژدې معنا له دې نظر سره سمون لري چې
واقعي نړۍ د عامل له معلوماتو سره سازګاره ده۔ په حقيقت کښې معرفتي
منطقونه عموماً S5 کاروي، خو نور يې ښايي کمزوري نظامونه وکاروي،
چې دا د \LR{$\Knows_a$} د عملګر له مطلوب تفسير سره تړاو لري۔ د
انعکاسيت شرط په هره نړۍ کښې لږ تر لږه خپله هماغه نړۍ د لاسرسي
وړ ګرځوي او د پوهې رښتينولي هم تضمينوي۔
\olpseganchor{OLP-0486-B009}{end}

\olpseganchor{OLP-0486-B010}{start}
\begin{figure}
  \begin{center}
    \begin{tikzpicture}[modal]
      \node[world] (w1) [label={[align=right]right:\mTrue{p}\\ \mFalse{q}}]{\LR{$w_1$}}; 
      \node[world] (w2) [label={[align=right]right:\mTrue{p}\\ \mTrue{q}},
        above right=of w1]{\LR{$w_2$}}; 
      \node[world] (w3) [label={[align=right]right:\mFalse{p}\\ \mFalse{q}},
        below right=of w1] {\LR{$w_3$}};
      \draw[<->] (w1) to node {\LR{$a$}} (w2);
      \draw[->,loop left] (w1) to node {\LR{$a,b$}} (w1);
      \draw[<->] (w1) to [swap] node {\LR{$b$}} (w3);
      \draw[->,loop above] (w2) to node {\LR{$a,b$}} (w2) ;
      \draw[->,loop below] (w3) to node {\LR{$a,b$}} (w3) ;
    \end{tikzpicture}
  \end{center}
  \caption{يو ساده معرفتي مدل۔}
  \ollabel{fig:simple}
\end{figure}
\olpseganchor{OLP-0486-B010}{end}

\olpseganchor{OLP-0486-B011}{start}
\begin{prob}
  په \olref[aml][el][trw]{fig:simple} کښې وګورئ چې له لاندې کومې
  دعوې سمې دي:
  \begin{enumerate}
    \item \LR{$\mSat{M}{\lnot q}[w_1]$};
    \item \LR{$\mSat{M}{\Knows_a \lnot q}[w_1]$};
    \item \LR{$\mSat{M}{\Knows_b \lnot q}[w_1]$};
    \item \LR{$\mSat{M}{\Knows_b q \lor \Knows_b \lnot q}[w_2]$};
    \item \LR{$\mSat{M}{\Knows_a( \Knows_b q \lor \Knows_b \lnot q)}[w_2]$};
    \item \LR{$\mSat{M}{\EKnows_{\{a,b\}} \lnot q}[w_3]$};
    \end{enumerate}
\end{prob}
\olpseganchor{OLP-0486-B011}{end}

\olpseganchor{OLP-0486-B012}{start}
اوس چې په يوه نړۍ کښې مو د رښتياوالي بنسټيز تعريف ورکړ، د نورمال
مودالي منطق نور معنايي مفاهيم، لکه مودالي اعتبار او لازم راتلل، هم
دغه نوي تفسير ته د مودالي عاملونو په پلي کولو سره راځي۔
\olpseganchor{OLP-0486-B012}{end}

\olpseganchor{OLP-0486-B013}{start}
اوس د مشترکې پوهې د عامل~\LR{$\CKnows_G$} لپاره هم د صدق شرطونه ټاکلے
شو۔ د \olref[sfr][rel][ops]{sec} له مخې ياد کړئ چې د اړيکې~\LR{$R$}
\emph{انتقالي تړون}~\LR{$R^+$} داسې تعريفېږي:
\begin{align*}
  R^+ &= \bigcup_{0 < n \in \mathbb{N}} R^n, 
\intertext{\RL{چې}}
  R^1 & = R \text{\RL{ او}}\\ 
  R^{n+1} & = \Setabs{\tuple{x, z}}{\exists y (R^n xy \land Ryz)}.
\end{align*}
بيا، که \LR{$G$} د عاملانو يوه ډله وي، نو \LR{$R_G = ( \bigcup_{b
\in G} R_b )^+$} د ټولو عاملانو د لاسرسي اړيکو د اتحاد انتقالي
تړون تعريفوو۔ دا اړيکه د ډلې د اړيکو له يو يا زياتو پړاوونو
جوړې شوې لارې نيسي۔ په داسې لار کښې د هر پړاو عامل جلا ټاکل
کېدے شي۔
\olpseganchor{OLP-0486-B013}{end}

\olpseganchor{OLP-0486-B014}{start}
\begin{defn}
که \LR{$G' \subseteq G$} وي، نو \LR{$\mSat{M}{\CKnows_{G'} !A}[ w]$}
هله او يوازې هله وي چې د هر هغه \LR{$w'$} لپاره چې \LR{$R_{G'} w w'$}
وي، \LR{$\mSat{M}{!A}[w']$} وي۔ په دې تعريف کښې صفر پړاوه لاره
نۀ شمېرل کېږي؛ خو که د ډلې د کوم عامل اړيکه انعکاسي وي،
له اوسنۍ نړۍ څخه بېرته هماغې نړۍ ته يوه يوپړاوه لاره شته۔
\end{defn}
\olpseganchor{OLP-0486-B014}{end}

\olpseganchor{OLP-0487-B004}{start}
\olfileid{aml}{el}{acc}
\olpseganchor{OLP-0487-B004}{end}

\olpseganchor{OLP-0487-B005}{start}
\olsection{د لاسرسي اړيکې او معرفتي اصول}
\olpseganchor{OLP-0487-B005}{end}

\olpseganchor{OLP-0487-B006}{start}
د نورمال مودالي منطقونو د چوکاټ د مطابقت په اړه له خپلو پخوانيو
معلوماتو سره، ښايي وغواړو وګورو چې ځانګړي فارمولونه د معرفتي
تفسير له مخې څۀ بڼه لري۔ مخکې مو ويلي چې معرفتي منطقونه عموماً
په S5 کښې تفسيرېږي۔ نو وګورو چې بېلابېل معرفتي اصول څنګه ښودل
کېږي او له بېلابېلو چوکاټي شرطونو سره يې مطابقت څنګه دے۔
\olpseganchor{OLP-0487-B006}{end}

\olpseganchor{OLP-0487-B007}{start}
له نورمال مودالي منطق څخه ياد کړئ چې بېلابېلو مودالي فارمولونه
د لاسرسي د اړيکو بېلابېل خاصيتونه مشخصول۔ دا جدول له هغو څخه څو
خاصيتونه رااخلي چې له ټاکلو معرفتي اصولو سره تړاو لري۔
\olpseganchor{OLP-0487-B007}{end}

\olpseganchor{OLP-0487-B008}{start}
\begin{table}[t]
    \begin{tabular}{| p{\dimexpr .5\linewidth-2\tabcolsep-2\arrayrulewidth-\doublerulesep/2\relax} || p{\dimexpr .5\linewidth-2\tabcolsep-2\arrayrulewidth-\doublerulesep/2\relax} |}
      \hline
      {\emph{که \LR{$R$} \dots وي}} & {\emph{نو \dots په~\LR{$\mModel{M}$} کښې رښتيا دے:}} \\
      \hline \hline
\olpseganchor{OLP-0487-B008}{end}

\olpseganchor{OLP-0487-B009}{start}
      \emph{هره دوه‌ځاييزه اړيکه}
      & \LR{$\Knows (p \lif q) \lif (\Knows p \lif \Knows q)$}
      \hfill \newline{(تړلتيا)} \\
      \hline
      \emph{انعکاسي}: \LR{$\forall w Rww$}
      & \LR{$\Knows p \lif p$} \hfill \newline{(رښتينولي)} \\
      \hline
      \emph{متعدي}: \newline
      \LR{$\forall u \forall v \forall w ((Ruv \land Rvw) \lif Ruw)$} &
      \LR{$\Knows p \lif \Knows \Knows p$} \hfill
           \newline{(مثبت ځان‌پوهاوی)} \\
      \hline
      \emph{اقليدسي}: \newline
      \LR{$\forall w \forall u \forall v ((Rwu \land Rwv) \lif Ruv)$} &
      \LR{$\lnot \Knows p \lif \Knows \neg \Knows p$} \hfill
        \newline{(منفي ځان‌پوهاوی)} \\
      \hline
    \end{tabular}
    \caption{څلور معرفتي اصول۔}
    \ollabel{tab:four}
  \end{table}
\olpseganchor{OLP-0487-B009}{end}

\olpseganchor{OLP-0487-B010}{start}
رښتينولي، چې د \LR{$T$} اکسيوم سره مطابقت لري، ډېر ځله د دغو اصولو
تر ټولو لږ د اختلاف وړ ګڼل کېږي؛ ځکه دا څرګندوي چې که يو
فارمول د عامل په پوهه کښې وي، نو بايد رښتينے هم وي۔ تړلتيا او مثبت
او منفي ځان‌پوهاوی تر دې ډېر د بحث وړ دي۔
\olpseganchor{OLP-0487-B010}{end}

\olpseganchor{OLP-0487-B011}{start}
تړلتيا، چې د \LR{$K$} اکسيوم سره مطابقت لري، دا نظر ښيي چې د يوه
عامل پوهه د استلزام له مخې تړلې ده۔ په ځينو حالاتو کښې ښايي دا
موږ ته معقوله ښکاره شي۔ د بېلګې په توګه، زه ښايي پوهېږم چې که
په ويکټوريا کښې يم، نو د وېنکوور پر ټاپو يم۔ له عجيبو شکاکانه
حالتونو پرته، زه پوهېږم چې په ويکټوريا کښې يم۔ د دې اصل له مخې
بايد په دې هم پوه شم چې د وېنکوور پر ټاپو يم۔ په دې حالت کښې زما د پوهې
منطقي تړلتيا ښايي تر يوه بريده طبيعي ښکاره شي۔ خو موږ تل د
خپلې پوهې په ټولو پايلو فکر نۀ کوو، او له دې امله په نورو
حالتونو کښې لږ طبيعي نتيجې هم راوتلے شي۔
\olpseganchor{OLP-0487-B011}{end}

\olpseganchor{OLP-0487-B012}{start}
مثبت ځان‌پوهاوی، چې کله د KK-اصل هم بلل کېږي، داسې بيانېږي: که
زه پر يوه څۀ پوهېږم، نو پوهېږم چې پرې پوهېږم۔ دا د 4 اکسيوم
معرفتي سيال دے۔ په ورته ډول، منفي ځان‌پوهاوی داسې بيانېږي: که
زه پر يوه څۀ \emph{نۀ} پوهېږم، نو پوهېږم چې پرې نۀ پوهېږم؛ دا د
5 اکسيوم سيال دے۔ دواړو ته نسبتاً عادي ضدبېلګې پيدا کېدے شي:
مثبت اصل هغه وخت تر پوښتنې لاندې راځي چې زه پر يوه خبره
پوهېږم، خو په دې نۀ پوهېږم چې پرې پوهېږم۔ د منفي اصل اړونده
بېلګه دا ده چې پر يوه خبره نۀ پوهېږم او له خپلې ناپوهۍ هم خبر
نۀ يم۔
\olpseganchor{OLP-0487-B012}{end}

\olpseganchor{OLP-0488-B004}{start}
\olfileid{aml}{el}{bsd}
\olpseganchor{OLP-0488-B004}{end}

\olpseganchor{OLP-0488-B005}{start}
\olsection{باي‌سيميولېشنونه}
\olpseganchor{OLP-0488-B005}{end}

\olpseganchor{OLP-0488-B006}{start}
د معرفتي ژبې د څرګندونې د ځواک په اړه يوه پاتې پوښتنه د مدلونو او
په هغو کښې د رښتينو فارمولونو تر منځ له اړيکې سره تړلې ده۔ د چوکاټ
د مطابقت له پايلو مو ليدلي چې کله په يوه چوکاټ کښې ځينې فارمولونه
معتبر وي، نو هغه چوکاټ ځينې خاصيتونه هم پوره کوي۔ خو آيا زموږ
مودالي ژبه، د بېلګې په توګه، هغه نړۍ چې انعکاسي غشے لري له
هغې د نړۍو له نامتناهي لړۍ څخه بېلولے شي چې هره نړۍ ئې بلې ته
لار لري؟ يعنې آيا کوم فارمول \LR{$A$} شته چې يوازې له دغو دوو نړۍو
څخه په يوه کښې رښتيا وي؟ که قضيېيز متغيرونه په دواړو خواوو کښې
يو شان ارزښت ولري او همدغه غشي يوازينۍ د لاسرسي اړيکه وټاکي، نو
ځان ته رسېدونکې نړۍ د لړۍ له هرې نړۍ سره د باي‌سيميولېشن له لارې
تړل کېدے شي۔
\olpseganchor{OLP-0488-B006}{end}

\olpseganchor{OLP-0488-B007}{start}
باي‌سيميولېشن د اړيکيزو مدلونو تر منځ داسې اړيکه ده چې په وسيله
ئې ټاکلې نړۍ سره تړو۔ سره تړلې نړۍ د معرفتي ژبې عين فارمولونه
رښتيني کوي۔
\olpseganchor{OLP-0488-B007}{end}

\olpseganchor{OLP-0488-B008}{start}
\begin{defn}[باي‌سيميولېشن]
فرض کړئ \LR{$M_1 = \tuple{W_1, R_1, V_1}$} او \LR{$M_2 = \tuple{W_2, R_2, V_2}$}
دوه اړيکيز مدلونه وي۔ همدارنګه، \LR{$\mathcal{R} \subseteq W_1 \times W_2$}
دې يوه دوه‌ځاييزه اړيکه وي۔ \LR{$\mathcal{R}$} هله
\emph{باي‌سيميولېشن} بولو چې د هر \LR{$\tuple{w_1, w_2} \in \mathcal{R}$}
لپاره لاندې شرطونه وي:
\olpseganchor{OLP-0488-B008}{end}

\olpseganchor{OLP-0488-B009}{start}
\begin{enumerate}
  \item د ټولو بياني متغيرونه~\LR{$p$} لپاره
  \LR{$w_1 \in V_1(p)$} هله او يوازې هله چې \LR{$w_2 \in V_2(p)$} وي۔
  \item د ټولو عاملانو \LR{$a \in G$} او نړۍو \LR{$v_1 \in W_1$} لپاره، که
  \LR{$R_{1_a} w_1 v_1$} وي، نو داسې \LR{$v_2 \in W_2$} شته چې
  \LR{$R_{2_a} w_2 v_2$} او \LR{$\tuple{v_1, v_2} \in
  \mathcal{R}$} وي۔
  \item د ټولو عاملانو \LR{$a \in G$} او نړۍو \LR{$v_2 \in W_2$} لپاره، که
   \LR{$R_{2_a} w_2 v_2$} وي، نو داسې \LR{$v_1 \in W_1$} شته چې
   \LR{$R_{1_a} w_1 v_1$} او \LR{$\tuple{v_1, v_2} \in
   \mathcal{R}$} وي۔
\end{enumerate}
د دواړو مدلونو د عاملانو نښې له يوه ګډ سټه اخلو۔
\olpseganchor{OLP-0488-B009}{end}

\olpseganchor{OLP-0488-B010}{start}
که د \LR{$M_1$} او \LR{$M_2$} تر منځ باي‌سيميولېشن نړۍوې \LR{$w_1$} او \LR{$w_2$}
سره وتړي، نو \LR{$\tuple{M_1, w_1} \leftrightarroweq
\tuple{M_2, w_2}$} هم ليکلے شو، او \LR{$\tuple{M_1, w_1}$} او
\LR{$\tuple{M_2, w_2}$} \emph{باي‌سيميولر} بولو۔
\end{defn}
\olpseganchor{OLP-0488-B010}{end}

\olpseganchor{OLP-0488-B011}{start}
د باي‌سيميولېشن د اړيکې بېلابېلې مادې بېلابېل شيان تضمينوي۔
لومړۍ ماده ډاډ راکوي چې باي‌سيميولر نړۍوې به هماغه له موداليت
پرته فارمولونه پوره کوي، ځکه د ټولو بياني متغيرونه
په اړه يې يو شان حالت ټاکي۔ نورې دوه مادې، چې په ترتيب سره کله
«مخ ته» او «شاته» بلل کېږي۔ د هرې تړلې جوړې او هر عامل لپاره
په يوه مدل کښې هره د لاسرسي وړ نړۍ په بل مدل کښې يوه تړلې
ملګرې نړۍ لري۔ يوه نړۍ د څو نړۍو ملګرې کېدے شي۔
\olpseganchor{OLP-0488-B011}{end}

\olpseganchor{OLP-0488-B012}{start}
\begin{thm}
که \LR{$\tuple{M_1, w_1} \leftrightarroweq \tuple{M_2, w_2}$} وي، نو د
هر فارمول~\LR{$!A$} لپاره \LR{$\mSat{M_1}{!A}[w_1]$} هله او يوازې
هله چې \LR{$\mSat{M_2}{!A}[w_2]$} وي۔
\end{thm}
\olpseganchor{OLP-0488-B012}{end}

\olpseganchor{OLP-0488-B013}{start}
\begin{figure}
  \begin{center}
    \begin{tikzpicture}[modal]
      \node[world] (w1) {\LR{$w_1$}}; 
      \node[world] (w2) [above left=of w1]{\LR{$w_2$}}; 
      \node[world] (w3) [above right=of w1] {\LR{$w_3$}};
      \draw[<->] (w1) to node {\LR{$a$}} (w2);
      \draw[->,loop below] (w1) to node {\LR{$a$}} (w1);
      \draw[<->] (w1) to [swap] node {\LR{$a$}} (w3);
      \draw[->,loop above] (w2) to node {\LR{$a$}} (w2) ;
      \draw[->,loop above] (w3) to node {\LR{$a$}} (w3) ;
\olpseganchor{OLP-0488-B013}{end}
      
\olpseganchor{OLP-0488-B014}{start}
      \node[world](v1)[right of=w1, xshift=1.5in]{\LR{$v_1$}};
      \node[world](v2)[above of=v1]{\LR{$v_2$}};
      \draw[<->] (v1) to node {\LR{$a$}} (v2);
      \draw[->,loop below] (v1) to node {\LR{$a$}} (v1);
      \draw[->,loop above] (v2) to node {\LR{$a$}} (v2) ;
\olpseganchor{OLP-0488-B014}{end}

\olpseganchor{OLP-0488-B015}{start}
      \draw[-,dotted] (w1) to (v1) ;
      \draw[-,dotted,bend left=45] (w2) to (v2) ;
      \draw[-,dotted,bend left=30] (w3) to (v2) ;
    \end{tikzpicture}
  \end{center}
  \caption{دوه باي‌سيميولر مدلونه۔}
  \ollabel{fig:bisimilar}
\end{figure}
\olpseganchor{OLP-0488-B015}{end}

\olpseganchor{OLP-0488-B016}{start}
که څۀ هم په \olref{fig:bisimilar} کښې ښودل شوي دوه مدلونه عين
يو شان نۀ دي، خو داسې باي‌سيميولېشن شته چې نړۍوې \LR{$w_1$} او~\LR{$v_1$}
سره تړي۔ دا اړيکه \LR{$w_2$} او \LR{$w_3$} دواړه هم له~\LR{$v_2$} سره تړي؛
مطلب دا چې زموږ په مودالي ژبه کښې داسې څۀ نۀ شي څرګندېدے چې
په رښتيا يې له يو بل څخه بېل کړي۔ خو که \LR{$w_2$} او~\LR{$w_3$} بېلابېل
بياني متغيرونه پوره کول، نو دواړه له عينې درېيمې
نړۍ سره تړل ناشوني وو؛ ځکه هره جوړه بايد په اټومي ارزښتونو
کښې سره سمه وي۔
\olpseganchor{OLP-0488-B016}{end}

\olpseganchor{OLP-0489-B004}{start}
\olfileid{aml}{el}{pal}
\olpseganchor{OLP-0489-B004}{end}

\olpseganchor{OLP-0489-B005}{start}
\olsection{د عام اعلان منطق}
\olpseganchor{OLP-0489-B005}{end}

\olpseganchor{OLP-0489-B006}{start}
خوځنده معرفتي منطقونه موږ ته دا ښودلے شي چې د عاملانو پوهه د وخت
په تېرېدو يا د نويو معلوماتو په ترلاسه کولو سره څنګه بدلېږي۔ ډېر
داسې منطقونه د پوهې بدلون د معلوماتي \emph{پېښو} يا
\emph{تازه‌کونو} په وسيله ښيي۔ د تازه‌کونې تر ټولو بنسټيز ډول
عام اعلان دے۔ اعلان شوې دعوه د اعلان په شېبه کښې رښتينې وي او
ټول عاملان په ګډه د اعلان شاهدان وي۔ د دې دپاره ژبه داسې پراخوو:
\olpseganchor{OLP-0489-B006}{end}

\olpseganchor{OLP-0489-B007}{start}
\begin{defn}
فرض کړئ \LR{$G$} د عاملانو د نښو يو سټ وي۔ د عامو اعلانونو په ګډون
د څو-عامله معرفتي منطق بنسټيزه ژبه دا توکي لري:
\begin{enumerate}
  \item د دروغوالی قضيوي ثابت~\LR{$\lfalse$}۔
  \item د رښتياوالی قضيوي ثابت~\LR{$\ltrue$}۔
  \item د بياني متغيرونه يو شمېرېدونکے لامتناهي سټ: \LR{$\Obj
    p_0$}، \LR{$\Obj p_1$}، \LR{$\Obj p_2$}، \dots
  \item قضيوي نښلوونکي: \startycommalist
  \ycomma \LR{$\lnot$} (نفي)%
  \ycomma \LR{$\land$} (عطف)%
  \ycomma \LR{$\lor$} (فصل)%
  \ycomma \LR{$\lif$} (شرطي)%
  \item د پوهې عامل \LR{$\Knows_a$}، چې \LR{$a \in G$} وي۔
  \item د عام اعلان عامل \LR{$[!B]$}، چې \LR{$!B$} يو فارمول وي۔
\end{enumerate}
\end{defn}
\olpseganchor{OLP-0489-B007}{end}

\olpseganchor{OLP-0489-B008}{start}
د اعلان له شرط سره تړلے عامل د بکس په شان کار کوي۔ اعلان خپله
يو فارمول دے او په اعلان پسې بل فارمول راځي۔ د دوه‌اړخيز شرط ماده
هم لاندې د هغې ژبې د بڼې لپاره راغلې ده چې دا نښلوونکے لري۔
\olpseganchor{OLP-0489-B008}{end}

\olpseganchor{OLP-0489-B009}{start}
\begin{defn}
  د معرفتي ژبې \emph{فارمولونه} په استقرايي ډول داسې
  تعريفېږي:
  \begin{enumerate}
  \item \LR{$\lfalse$} يو اټومي فارمول دے۔
\olpseganchor{OLP-0489-B009}{end}

\olpseganchor{OLP-0489-B010}{start}
  \item \LR{$\ltrue$} يو اټومي فارمول دے۔
\olpseganchor{OLP-0489-B010}{end}

\olpseganchor{OLP-0489-B011}{start}
  \item هر قضيوي متحول \LR{$\Obj p_i$} يو اټومي
  فارمول دے۔
\olpseganchor{OLP-0489-B011}{end}

\olpseganchor{OLP-0489-B012}{start}
  \item که \LR{$!A$} يو فارمول وي، نو \LR{$\lnot !A$} هم
  يو فارمول دے۔
\olpseganchor{OLP-0489-B012}{end}

\olpseganchor{OLP-0489-B013}{start}
  \item که \LR{$!A$} او \LR{$!B$} فارمولونه وي، نو \LR{$(!A \land
  !B)$} يو فارمول دے۔
\olpseganchor{OLP-0489-B013}{end}

\olpseganchor{OLP-0489-B014}{start}
  \item که \LR{$!A$} او \LR{$!B$} فارمولونه وي، نو \LR{$(!A \lor !B)$}
  يو فارمول دے۔
\olpseganchor{OLP-0489-B014}{end}

\olpseganchor{OLP-0489-B015}{start}
  \item که \LR{$!A$} او \LR{$!B$} فارمولونه وي، نو \LR{$(!A \lif !B)$}
  يو فارمول دے۔
\olpseganchor{OLP-0489-B015}{end}

\olpseganchor{OLP-0489-B016}{start}
  \item که \LR{$!A$} او \LR{$!B$} فارمولونه وي، نو \LR{$(!A \liff !B)$}
  يو فارمول دے۔
\olpseganchor{OLP-0489-B016}{end}

\olpseganchor{OLP-0489-B017}{start}
  \item که \LR{$!A$} يو فارمول وي او \LR{$a \in G$} وي، نو \LR{$\Knows_a !A$}
  يو فارمول دے۔
\olpseganchor{OLP-0489-B017}{end}
  
\olpseganchor{OLP-0489-B018}{start}
  \item که \LR{$!A$} او \LR{$!B$} فارمولونه وي، نو \LR{$[!A] !B$} يو فارمول دے۔
\olpseganchor{OLP-0489-B018}{end}

\olpseganchor{OLP-0489-B019}{start}
  \item بل هېڅ شے فارمول نۀ دے۔
\end{enumerate}
\end{defn}
\olpseganchor{OLP-0489-B019}{end}

\olpseganchor{OLP-0489-B020}{start}
د فارمول~\LR{$[!A] !B$} موخه شوې معنا دا ده: «وروسته له دې چې
\LR{$!A$} په رښتيا اعلان شي، \LR{$!B$} رښتيا وي»۔ کله کله د عامو اعلانونو
په بحث کښې د مشترکې پوهې يادونه هم ګټوره وي؛ نو ژبه ښايي د مشترکې
پوهې عامل~\LR{$\CKnows_G
!A$} هم ولري۔
\olpseganchor{OLP-0489-B020}{end}

\olpseganchor{OLP-0490-B004}{start}
\olfileid{aml}{el}{psm}
\olpseganchor{OLP-0490-B004}{end}

\olpseganchor{OLP-0490-B005}{start}
\olsection{د عام اعلان د منطق معناپوهنه}
\olpseganchor{OLP-0490-B005}{end}

\olpseganchor{OLP-0490-B006}{start}
د عام اعلان د منطق اړيکيز مدلونه د معرفتي منطق د مدلونو په شان
دي۔ د اعلان عامل د همدغه مدل پر محدودې شوې برخې ارزول کېږي۔
\olpseganchor{OLP-0490-B006}{end}

\olpseganchor{OLP-0490-B007}{start}
\begin{defn}\ollabel{defn:mmodels} په~\LR{$\mModel M = \tuple{W, R, V}$}
  کښې \emph{په~\LR{$w$} د فارمول~\LR{$!A$} رښتياوالی}، چې په نښو
  کښې ئې \LR{$\mSat{M}{!A}[w]$} ليکو، په استقرايي ډول داسې تعريفېږي:
  \begin{enumerate}
  \item \indcase{!A}{\lfalse}{هېڅکله \LR{$\mSat{M}{\lfalse}[w]$} نۀ وي}.
  \item \indcase{!A}{\ltrue}{تل \LR{$\mSat{M}{\ltrue}[w]$} وي}.
  \item \LR{$\mSat{M}{p}[w]$} هله او يوازې هله چې \LR{$w \in V(p)$} وي۔
  \item \indcase{!A}{\lnot !B}{\LR{$\mSat{M}{\indfrm}[w]$} هله او يوازې هله چې
    \LR{$\mSat/{M}{!B}[w]$} وي}.
  \item \indcase{!A}{(!B \land !C)}{\LR{$\mSat{M}{\indfrm}[w]$} هله او يوازې هله چې
    \LR{$\mSat{M}{!B}[w]$} او \LR{$\mSat{M}{!C}[w]$} دواړه وي}.
  \item \indcase{!A}{(!B \lor !C)}{\LR{$\mSat{M}{\indfrm}[w]$} هله او يوازې هله چې
    \LR{$\mSat{M}{!B}[w]$} يا \LR{$\mSat{M}{!C}[w]$} وي} (يا دواړه)۔
  \item \indcase{!A}{(!B \lif !C)}{\LR{$\mSat{M}{\indfrm}[w]$} هله او يوازې هله چې
    \LR{$\mSat/{M}{!B}[w]$} يا \LR{$\mSat{M}{!C}[w]$} وي}.
  \item \indcase{!A}{(!B \liff !C)}{\LR{$\mSat{M}{\indfrm}[w]$} هله او يوازې هله چې
    يا دواړه \LR{$\mSat{M}{!B}[w]$} او \LR{$\mSat{M}{!C}[w]$} وي، يا
    نۀ \LR{$\mSat{M}{!B}[w]$} وي او نۀ \LR{$\mSat{M}{!C}[w]$} وي}.
  \item\ollabel{defn:sub:mmodels-box}
    \indcase{!A}{\Knows_a !B}{\LR{$\mSat{M}{\indfrm}[w]$} هله او يوازې هله چې
    د هر هغه \LR{$w' \in W$} لپاره \LR{$\mSat{M}{!B}[w']$} وي چې \LR{$R_a ww'$} وي}
  \item\ollabel{defn:sub:mmodels-pal}
    \indcase{!A}{[!B] !C}{\LR{$\mSat{M}{\indfrm}[w]$} هله او يوازې هله چې
    \LR{$\mSat{M}{!B}[w]$} له مخې \LR{$\mSat{M \mid !B}{!C}[w]$} لازم راځي}
\olpseganchor{OLP-0490-B007}{end}

  
\olpseganchor{OLP-0490-B008}{start}
  دلته \LR{$\mModel M \mid !B = \tuple{W', R', V'}$} داسې تعريفېږي:
  \begin{enumerate}
    \item \LR{$W' = \Setabs{ u \in W}{\mSat{M}{!B}[u]}$}۔ نو د
      \LR{$\mModel M \mid !B$} نړۍوې د~\LR{$\mModel M$} هغه نړۍوې دي چې
      \LR{$!B$} پکښې رښتيا وي۔
    \item \LR{$R'_a = R_a \cap (W' \times W')$}۔ د هر عامل د لاسرسي
      اړيکه يوازې هغو نړۍو ته محدودېږي چې په~\LR{$W'$} کښې پاتې وي۔
    \item \LR{$V'(p) = \Setabs{ u \in W'}{u \in V(p) }$}۔ په ورته ډول،
      په هره پاتې نړۍ کښې قضيوي ارزښت ټاکنه هماغسې پاتې کېږي۔ دا هغه
      نظر ښيي چې معلوماتي پېښې د قضيوي متحولونو د رښتياوالي
      قيمت نۀ بدلوي۔
  \end{enumerate}
\olpseganchor{OLP-0490-B008}{end}
  
\olpseganchor{OLP-0490-B009}{start}
  \end{enumerate}
\end{defn}
\olpseganchor{OLP-0490-B009}{end}

\olpseganchor{OLP-0490-B010}{start}
نو د عام اعلان د منطق ځانګړتيا دا ده چې په~\LR{$\mModel M$} کښې د
فارمول رښتياوالے کله کله يوازې د~\LR{$\mModel M$} له محدود
مدل څخه معلومولے شو۔
\olpseganchor{OLP-0490-B010}{end}

\olpseganchor{OLP-0490-B011}{start}
دلته د اعلان عامل د \LR{$\Box$} په شان چلند کوي۔ که فارمول~\LR{$!A$}
په اوسنۍ نړۍ کښې رښتينے نۀ وي، نو \LR{$[!A]!B$} په هماغې نړۍ کښې
په تش ډول رښتينے دے۔ په دې حالت کښې محدود مدل ته تګ او هلته
ارزونه نۀ کېږي۔ همداسې، که له يوې نړۍ هېڅ نړۍ ته لار نۀ وي،
د \LR{$\Box$} د فارمولونه عمومي شرط په تش ډول پوره کېږي۔
\olpseganchor{OLP-0490-B011}{end}

\olpseganchor{OLP-0490-B012}{start}
\begin{figure}
  \begin{center}
    \begin{tikzpicture}[modal]
      \node[world] (w1) [label={below left:\LR{$p, \lnot q$}}]{\LR{$w_1$}}; 
      \node[world] (w2) [left=of w1, label={left:\LR{$\lnot p, \lnot q$}}]{\LR{$w_2$}}; 
      \node[world] (w3) [above=of w1, label={right:\LR{$p, q$}}] {\LR{$w_3$}};
      \draw[<->] (w1) to [swap] node {\LR{$b$}} (w2);
      \draw[->,loop below] (w1) to node {\LR{$a, b$}} (w1);
      \draw[<->] (w1) to [swap] node {\LR{$a$}} (w3);
      \draw[->,loop above] (w2) to node {\LR{$a, b$}} (w2) ;
      \draw[->,loop above] (w3) to node {\LR{$a, b$}} (w3) ;
\olpseganchor{OLP-0490-B012}{end}
      
\olpseganchor{OLP-0490-B013}{start}
      \node[world](v1)[right of=w1, xshift=1.25in, label={right:\LR{$p, \lnot q$}}]{\LR{$w'_1$}};
      \node[world](v3)[above of=v1,label={right:\LR{$p,q$}}]{\LR{$w'_3$}};
      \draw[<->] (v1) to node {\LR{$a$}} (v3);
      \draw[->,loop below] (v1) to node {\LR{$a,b$}} (v1);
      \draw[->,loop above] (v3) to node {\LR{$a,b$}} (v3) ;
\olpseganchor{OLP-0490-B013}{end}
      
\olpseganchor{OLP-0490-B014}{start}
      \node(m)[below=of w1]{\LR{$\mModel M$}};
      \node(m')[below=of v1]{\LR{$\mModel M \mid p$}}; 
\olpseganchor{OLP-0490-B014}{end}

\olpseganchor{OLP-0490-B015}{start}
      \draw[-,dotted] (w1) to node {\footnotesize \RL{د~\LR{$p$} اعلان}}(v1) ;
\olpseganchor{OLP-0490-B015}{end}
     
\olpseganchor{OLP-0490-B016}{start}
    \end{tikzpicture}
  \end{center}
  \caption{د~\LR{$p$} له عام اعلان څخه مخکې او وروسته۔}
  \ollabel{fig:announcement-example}
\end{figure}
\olpseganchor{OLP-0490-B016}{end}

\olpseganchor{OLP-0490-B017}{start}
د فارمول عام اعلان داسې ليدلے شو چې مدل کوچنے کوي، يا يې
يوازې هغو نړۍو ته محدودوي چې فارمول پکښې رښتيا و۔
\olref{fig:announcement-example} د~\LR{$p$} د عام اعلان د اغېز بېلګه
ښيي۔ د دغه مدل يوه د پام وړ خبره دا ده چې عامل~\LR{$b$} د اعلان په
پايله کښې د~\LR{$p$} په رښتياوالي پوهېږي۔ عامل~\LR{$a$} پر دې له مخکې
پوهېدۀ۔ د \LR{$a$} لپاره له اوسنۍ نړۍ د لاسرسي وړ هرې نړۍ کښې~\LR{$p$}
رښتيا و۔
\olpseganchor{OLP-0490-B017}{end}

\olpseganchor{OLP-0490-B018}{start}
په رسمي ډول، \LR{$\mSat{M}{\lnot \Knows_b p}[w_1]$} لرو، خو \LR{$\mSat{M
\mid p}{\Knows_b p}[w'_1]$} هم لرو۔ له دې څخه \LR{$\mSat{M}{[p] \Knows_b
p}[w_1]$} لازمېږي۔ خو د اعلان د پايلې په اړه تر دې پياوړې دعوې
هم کولے شو۔ په حقيقت کښې \LR{$\mSat{M}{[p]\CKnows_{\{a,b\}} p}[w_1]$}
هم رښتينے دے۔ د شکل پر ښي لوري د ټکو پريمونه د پخوانيو ټکو
نوې نښې دي۔ له اوسنۍ نړۍ څخه د دواړو عاملانو په اړيکو تر لاسرسي
لاندې ټولې پاتې نړۍ همدا دعوه رښتينې کوي؛ ځکه د~\LR{$p$} له اعلان
وروسته دا خبره \emph{مشترکه پوهه} کېږي۔
\olpseganchor{OLP-0490-B018}{end}

\olpseganchor{OLP-0490-B019}{start}
خو ښايي وپوښتو چې آيا دا په عمومي حالت کښې هم رښتيا دي، او آيا
د~\LR{$!A$} رښتينے اعلان به \emph{تل} \LR{$!A$} په مشترکه پوهه بدلوي؟
ځواب نفي دے، که څۀ هم ښايي حيرانوونکے وي۔ ان داسې فارمولونه
په رښتيا اعلانېدل ممکن دي چې تر اعلان وروسته نور رښتيا نۀ وي۔
د بېلګې په توګه، په \olref{fig:announcement-example} کښې په~\LR{$w_1$}
د \LR{$p \land \lnot \Knows_b p$} د اعلان اغېز وګورئ۔ په حقيقت کښې
\LR{$\mModel M \mid p$} او \LR{$\mModel M \mid (p \land \lnot \Knows_b p)$}
بېلابېلې نړۍ لري۔ په لومړي محدود شوي مدل کښې د ښي لوري
دواړه نړۍ پاتې دي۔ په دويم کښې يوازې اوسنۍ نړۍ پاتې ده، ځکه
په برنۍ نړۍ کښې عامل ب له مخکې پر اټومي دعوه پوهېږي۔ په دواړو
محدوديتونو کښې له اوسنۍ نړۍ د عامل ب د لاسرسي يوازينۍ نړۍ همدا
خپله اوسنۍ نړۍ ده، او اټومي دعوه پکښې رښتينې ده۔ نو په لومړي
محدود شوي مدل کښې \LR{$\mSat{M \mid p}{\Knows_b
p}[w'_1]$} دے۔ نو \LR{$\mSat{M \mid (p \land \lnot \Knows_b p)}{\lnot
(p \land \lnot \Knows_b p)}[w'_1]$} دے۔ په دې توګه دا داسې
فارمول دے چې له اعلان وروسته نارښتيا کېږي۔
\olpseganchor{OLP-0490-B019}{end}

\olpseganchor{OLP-0491-B004}{start}
\olpart{int}{شهودي منطق}
\olpseganchor{OLP-0491-B004}{end}

\olpseganchor{OLP-0492-B004}{start}
\olchapter{int}{int}{سريزه}
\olpseganchor{OLP-0492-B004}{end}

\olpseganchor{OLP-0493-B004}{start}
\olfileid{int}{int}{cr}
\olpseganchor{OLP-0493-B004}{end}

\olpseganchor{OLP-0493-B005}{start}
\section{تعميري استدلال}
\olpseganchor{OLP-0493-B005}{end}

\olpseganchor{OLP-0493-B006}{start}
په مودالي عاملونو يا د دويمې درجې کميت‌ټاکونکو سره د کلاسيک منطق
د پراخولو پر خلاف، شهودي منطق په دې معنا «غيرکلاسيک» دے چې کلاسيک
منطق محدودوي۔ کلاسيک منطق له بېلابېلو اړخونو \emph{ناتعميري} دے۔
شهودي منطق غواړي هغه ډېر «تعميري» استدلال ونيسي چې د تعميري رياضياتو
له يوه ډول سره تړاو لري۔ لاندې بېلګې د دې نظر ځينې بنسټيزې انګېزې
روښانولے شي۔
\olpseganchor{OLP-0493-B006}{end}

\olpseganchor{OLP-0493-B007}{start}
فرض کړئ يو څوک دعوه وکړي چې داسې طبيعي عدد~\LR{$n$} ئې ټاکلے دے چې که
\LR{$n$} جوړ وي، د ريمان فرضيه رښتيا ده، او که \LR{$n$} طاق وي، د ريمان فرضيه
نارښتيا ده۔ دا خو زېرے شو!{} د ريمان د فرضيې رښتياوالی د رياضياتو له
سترو پرانيستو پوښتنو څخه دے، او ښکاري چې دغه کس ستونزه يوازې يوې
محاسبې ته راکمه کړې ده: دا معلومول چې يو ځانګړے عدد جوړ دے که نه۔
\olpseganchor{OLP-0493-B007}{end}

\olpseganchor{OLP-0493-B008}{start}
د~\LR{$n$} دا جادويي قيمت څو دے؟ هغه ئې داسې بيانوي: \LR{$n$} هغه طبيعي عدد
دے چې که د ريمان فرضيه رښتيا وي، له \LR{$2$} سره برابر دے، او که نه، له
\LR{$3$} سره برابر دے۔
\olpseganchor{OLP-0493-B008}{end}

\olpseganchor{OLP-0493-B009}{start}
تاسو په غصه خپلې پيسې بېرته غواړئ۔ له کلاسيک نظره پورته بيان په رښتيا
د \LR{$n$} يو يوازينے قيمت ټاکي؛ خو هغه څۀ چې تاسو ئې غواړئ، دا دي چې د
\LR{$n$} قيمت \emph{په څرګنده توګه} درکړے شي۔
\olpseganchor{OLP-0493-B009}{end}

\olpseganchor{OLP-0493-B010}{start}
يوه بله، ښايي لږه تصنعي بېلګه واخلئ۔ لاندې پوښتنه وګورئ۔ موږ پوهېږو
چې د يو غير ناطق عدد ناطق توان ناطقه پايله ورکولے شي۔ د
بېلګې په توګه، \LR{$\sqrt{2}^2 = 2$}۔ خو دا دومره څرګنده نۀ ده چې آيا
د يو غير ناطق عدد \emph{غير ناطق} توان هم ناطقه پايله ورکولے
شي که نه۔ لاندې قضيه مثبت ځواب ورکوي:
\olpseganchor{OLP-0493-B010}{end}

\olpseganchor{OLP-0493-B011}{start}
\begin{thm}
داسې غير ناطق عددونه \LR{$a$} او \LR{$b$} شته چې \LR{$a^b$} ناطق وي۔
\end{thm}
\olpseganchor{OLP-0493-B011}{end}

\olpseganchor{OLP-0493-B012}{start}
\begin{proof}
\LR{$\sqrt{2}^{\sqrt{2}}$} وګورئ۔ که دا ناطق وي، نو خبره ثابته شوه:
\LR{$a = b = \sqrt{2}$} اخلو۔ که نه، نو دا غير ناطق دے۔ بيا لرو:
\[
(\sqrt{2}^{\sqrt{2}})^{\sqrt{2}} = \sqrt{2}^{\sqrt{2} \cdot
  \sqrt{2}} = \sqrt{2}^2 = 2,
\]
چې ناطق دے۔ نو په دې حالت کښې \LR{$a$} د
\LR{$\sqrt{2}^{\sqrt{2}}$} او \LR{$b$} د~\LR{$\sqrt 2$} په توګه اخلو۔
\end{proof}
\olpseganchor{OLP-0493-B012}{end}

\olpseganchor{OLP-0493-B013}{start}
آيا دا معتبر ثبوت دے؟ ډېری رياضيپوهان ئې معتبر ګڼي۔ خو بيا هم يوه
ناخوښه نکته شته: موږ د حقيقي عددونو د داسې يوې جوړې شتون ثابت کړے دے
چې ټاکلے خاصيت لري، خو دا نۀ شو وئيلے چې \emph{کومه} جوړه ده۔ همدا
پايله په داسې ډول هم ثابتولے شو چې د \LR{$a$} او \LR{$b$} جوړه \emph{په ثبوت کښې
ورکړل شوې} وي: \LR{$a = \sqrt{3}$} او \LR{$b = \log_3 4$} واخلئ۔ بيا:
\[
a^b = \sqrt{3}^{\log_3 4} = 3^{1/2 \cdot \log_3 4} = (3^{\log_3
  4})^{1/2} = 4^{1/2}= 2,
\]
ځکه چې \LR{$3^{\log_3 x} = x$}۔ د درې مربع ريښه ناطقه نۀ ده؛
د ناطق کسر د مربع په صورت کښې د درې لومړني عامل توان جفت وي،
خو خپله درې دا عامل په طاق توان لري۔ که د درې پر بنسټ د څلورو
لوګارتم ناطق واے، د کسر د مخرج په توان اخيستلو سره به د درې
مثبت صحيح توان د څلورو له مثبت صحيح توان سره برابر شوے واے۔
په لومړنيو عاملونو د تجزيې يووالی دا برابري ناشونې ګرځوي۔
\olpseganchor{OLP-0493-B013}{end}

\olpseganchor{OLP-0493-B014}{start}
په لومړي ثبوت کښې د دوو حالتونو وېش د منځني حالت د ايستلو
پر عام اصل ولاړ دے۔ شهودي منطق داسې وېش هغه وخت مني چې د
اړوندو حالتونو فصل په لاس کښې ولرو۔ د داسې \LR{$x$} شتون ثابتول چې~\LR{$!A(x)$} پوره
کوي، دا معنا لري چې بايد يو ځانګړے~\LR{$x$} او دا ثبوت ورکړئ چې \LR{$!A$} پوره
کوي، لکه په دويم ثبوت کښې۔ د \LR{$!A$} يا \LR{$!B$} ثابتول غواړي چې له دواړو
څخه يو ثابت کړے شئ۔
\olpseganchor{OLP-0493-B014}{end}

\olpseganchor{OLP-0493-B015}{start}
په رسمي ژبه، شهودي منطق د کلاسيک منطق د اشتقاق نظام له يوې
ټاکلې محدودونې څخه لاس ته راځي۔ له رياضيکي نظره دا يوازې رسمي استنتاجي
نظامونه دي؛ خو لکه چې وويل شول، موخه ئې د رياضيکي استدلال يو ډول نيول
دي۔ يو څوک دا هغه استدلال ګڼلے شي چې د رياضياتو د يو ځانګړي فلسفي
نظر، لکه د براور د شهوديت، له مخې موجه وي؛ بل څوک ئې د بېشپ د تعميريت
په شان لا «محسوس» او قانع کوونکے رياضيکي استدلال ګڼلے شي؛ او دا بحث
هم کېدے شي چې آيا رسمي تعريف غيررسمي انګېزه په سمه توګه نيسي که نه۔
خو هر فلسفي دريځ چې ولرو، شهودي منطق د يوه په رسمي توګه وړاندې شوي
منطق په توګه څېړلے شو؛ او ډېر رياضيکي منطقپوهان، د هر دليل له مخې،
په دې څېړنه کښې دلچسپي لري۔
\olpseganchor{OLP-0493-B015}{end}

\olpseganchor{OLP-0494-B004}{start}
\olfileid{int}{int}{syn}
\olpseganchor{OLP-0494-B004}{end}

\olpseganchor{OLP-0494-B005}{start}
\olsection{د شهودي منطق نحو}
\olpseganchor{OLP-0494-B005}{end}

\olpseganchor{OLP-0494-B006}{start}
د شهودي منطق نحو د بياني منطق د نحو په شان ده۔ په کلاسيک بياني
منطق کښې ځينې نښلوونکي د نورو په وسيله تعريفېدے شي؛ د بېلګې په توګه،
\LR{$!A \lif !B$} د \LR{$\lnot !A \lor !B$} په وسيله، يا \LR{$!A \lor !B$} د
\LR{$\lnot(\lnot !A \land \lnot !B)$} په وسيله تعريفولے شو۔ نو د کلاسيک
منطق په ځينو وړاندې‌کونو کښې ځينې نښلوونکي د دغو تعريفونو لنډې نښې
وي۔ په شهودي منطق کښې داسې نۀ ده، مګر دوه استثناوې شته: \LR{$\lnot !A$}
د \LR{$!A \lif \lfalse$} لنډه نښه کېدے شي، او ډېر ځله همداسې تعريفېږي۔
په دې ژبه کښې \LR{$\lfalse$} لومړنۍ نښه ده، څو د نفي تعريف پر
خپله نفي تکيه ونۀ کړي!{} همدارنګه،
\LR{$!A \liff !B$} د کلاسيک منطق په شان د \LR{$(!A \lif !B) \land (!B \lif
!A)$} په توګه تعريفېدے شي۔
\olpseganchor{OLP-0494-B006}{end}

\olpseganchor{OLP-0494-B007}{start}
د بياني شهودي منطق فارمولونه له \emph{بياني متغيرونه}،
د قضيوي ثابت~\LR{$\lfalse$}، او \emph{منطقي نښلوونکو} څخه جوړېږي۔ موږ لرو:
\olpseganchor{OLP-0494-B007}{end}

\olpseganchor{OLP-0494-B008}{start}
\begin{enumerate}
\item د بياني متغيرونه يو شمېرېدونکے لامتناهي سټ~\LR{$\PVar$}:
  \LR{$\Obj p_0$}، \LR{$\Obj p_1$}، \dots
  \item د دروغوالی قضيوي ثابت~\LR{$\lfalse$}۔
\item منطقي نښلوونکي: \LR{$\land$} (عطف)، \LR{$\lor$} (فصل)، \LR{$\lif$} (شرطي)۔
\item د ليک نښې: پرانيستونکې قوس «(»، تړونکې قوس «)»، او کامه۔
\end{enumerate}
\olpseganchor{OLP-0494-B008}{end}

\olpseganchor{OLP-0494-B009}{start}
\begin{defn}[فارمول]
\ollabel{defn:formulas} د بياني شهودي منطق د \emph{فارمولونه}
سټ~\LR{$\Frm[L_0]$} په استقرايي ډول داسې تعريفېږي:
\begin{enumerate}
\item \LR{$\lfalse$} يو اټومي فارمول دے۔
\olpseganchor{OLP-0494-B009}{end}

\olpseganchor{OLP-0494-B010}{start}
\item هر بياني متغير~\LR{$\Obj p_i$} يو اټومي
  فارمول دے۔
\olpseganchor{OLP-0494-B010}{end}

\olpseganchor{OLP-0494-B011}{start}
\item که \LR{$!A$} او \LR{$!B$} فارمولونه وي، نو \LR{$(!A \land
  !B)$} يو فارمول دے۔
\olpseganchor{OLP-0494-B011}{end}

\olpseganchor{OLP-0494-B012}{start}
\item که \LR{$!A$} او \LR{$!B$} فارمولونه وي، نو \LR{$(!A \lor !B)$}
  يو فارمول دے۔
\olpseganchor{OLP-0494-B012}{end}

\olpseganchor{OLP-0494-B013}{start}
\item که \LR{$!A$} او \LR{$!B$} فارمولونه وي، نو \LR{$(!A \lif !B)$}
  يو فارمول دے۔
\olpseganchor{OLP-0494-B013}{end}

\olpseganchor{OLP-0494-B014}{start}
\item له دې پرته بل هېڅ څۀ فارمول نۀ دے۔
\end{enumerate}
\end{defn}
\olpseganchor{OLP-0494-B014}{end}

\olpseganchor{OLP-0494-B015}{start}
د پورته لومړنيو نښلوونکو تر څنګ، لاندې \emph{تعريف شوې} نښې هم
کاروو: \LR{$\lnot$} (نفي) او \LR{$\liff$} (دوه اړخيز شرطي)۔ په دې تعريف
شويو عاملونو جوړ شوي فارمولونه داسې لوستل کېږي:
\begin{enumerate}
\item \LR{$\lnot !A$} د \LR{$!A \lif \lfalse$} لنډه نښه ده۔
\olpseganchor{OLP-0494-B015}{end}

\olpseganchor{OLP-0494-B016}{start}
\item \LR{$!A \liff !B$} د \LR{$(!A \lif !B) \land (!B
  \lif !A)$} لنډه نښه ده۔
\end{enumerate}
\olpseganchor{OLP-0494-B016}{end}

\olpseganchor{OLP-0494-B017}{start}
که څۀ هم \LR{$\lnot$} په رسمي ډول لنډه نښه ګڼل کېږي، کله کله به په
تعريفونو کښې د~\LR{$\lnot$} لپاره ښکاره قاعدې او مادې داسې ورکوو لکه
چې لومړنۍ نښه وي۔ دا زياتره د تمريني مسئلو د بيان لپاره کوو۔
\olpseganchor{OLP-0494-B017}{end}

\olpseganchor{OLP-0495-B004}{start}
\olfileid{int}{int}{bhk}
\olpseganchor{OLP-0495-B004}{end}

\olpseganchor{OLP-0495-B005}{start}
\olsection{د براور--هايتنګ--کولموګوروف تفسير}
\olpseganchor{OLP-0495-B005}{end}

\olpseganchor{OLP-0495-B006}{start}
\begin{editorial}
  د بي‌اېچ‌کې په تفسير کښې د هر نښلوونکي لپاره ګورو چې د هغه د
  جوړونې معلومات څۀ دي۔ د شرطي فارمول جوړونه يوې جوړونې ته بله
  جوړونه ورکوي، او د فصل جوړونه له خپلې نښې سره اړوند اړخ ټاکي۔
\end{editorial}
\olpseganchor{OLP-0495-B006}{end}

\olpseganchor{OLP-0495-B007}{start}
د شهودي نښلوونکو يو غيررسمي تعميري تفسير شته چې عموماً د
براور--هايتنګ--کولموګوروف تفسير بلل کېږي۔ دا د «جوړونې» مفهوم کاروي۔
جوړونه دلته د يوې دعوې تعميري ثبوت يا د ثبوت د جوړولو طريقه ده۔ (په BHK تفسير کښې د «ثبوت»
پر ځاے «جوړونه» وايو، څو په رسمي اشتقاق نظام کښې له
اشتقاق سره ګډ نۀ شي۔) د همدې شهودي مفهوم په بنسټ BHK تفسير
د شهودي نښلوونکو معناوې بيانوي۔
\olpseganchor{OLP-0495-B007}{end}

\olpseganchor{OLP-0495-B008}{start}
\begin{enumerate}
\item فرض کوو چې د يوې اټومي جملې جوړونه څۀ شی دے، دا موږ ته معلومه ده۔
\item د \LR{$!A_1 \land !A_2$} جوړونه يوه جوړه \LR{$\tuple{M_1, M_2}$} ده
  چې \LR{$M_1$} د~\LR{$!A_1$} جوړونه او \LR{$M_2$} د~\LR{$!A_2$} جوړونه وي۔
\item د \LR{$!A_1 \lor !A_2$} جوړونه يوه جوړه \LR{$\tuple{s, M}$} ده
  چې يا \LR{$s$} له~\LR{$1$} سره برابر وي او \LR{$M$} د~\LR{$!A_1$} جوړونه وي، يا \LR{$s$}
  له~\LR{$2$} سره برابر وي او \LR{$M$} د~\LR{$!A_2$} جوړونه وي۔
\item د \LR{$!A \lif !B$} جوړونه يو تابع دے چې د~\LR{$!A$} جوړونه
  د~\LR{$!B$} په جوړونه اړوي۔
\item د~\LR{$\lfalse$} (محال) هېڅ جوړونه نۀ شته۔
\item \LR{$\lnot !A$} د \LR{$!A \lif \lfalse$} مترادفه نښه تعريفېږي۔ يعنې
  د \LR{$\lnot !A$} جوړونه يو تابع دے چې د~\LR{$!A$} جوړونه د~\LR{$\lfalse$}
  په جوړونه اړوي۔
\end{enumerate}
\olpseganchor{OLP-0495-B008}{end}

\olpseganchor{OLP-0495-B009}{start}
\begin{ex}
د بېلګې په توګه \LR{$\lnot \lfalse$} واخلئ۔ د دې جوړونه داسې تابع دے
چې د \LR{$\lfalse$} هره ورکړل شوې جوړونه د~\LR{$\lfalse$} په جوړونه اړوي۔
ښکاره ده چې عيني تابع~\LR{$\Id{}$} داسې جوړونه ده: که د~\LR{$\lfalse$}
جوړونه~\LR{$M$} ورکړل شي، \LR{$\Id{}(M) = M$} بيا د~\LR{$\lfalse$} جوړونه ورکوي۔
\end{ex}
\olpseganchor{OLP-0495-B009}{end}

\olpseganchor{OLP-0495-B010}{start}
په عمومي ډول، \LR{$\lnot !A$} دا معنا لري چې «د \LR{$!A$} جوړونه ناشونې ده»۔
\olpseganchor{OLP-0495-B010}{end}

\olpseganchor{OLP-0495-B011}{start}
\begin{ex}
د هرې قضيې~\LR{$!A$} لپاره \LR{$!A \lif \lnot \lnot !A$} ثابتوو، چې
\LR{$!A \lif ((!A \lif \lfalse) \lif \lfalse)$} دے۔ جوړونه بايد داسې
تابع~\LR{$f$} وي چې د \LR{$!A$} يوه جوړونه~\LR{$M$} واخلي او د
\LR{$(!A \lif \lfalse) \lif \lfalse$} يوه جوړونه~\LR{$f(M)$} ورکړي۔
تابع~\LR{$f$} د \LR{$(!A \lif \lfalse) \lif \lfalse$} جوړونه داسې جوړوي:
بايد داسې تابع \LR{$g$} تعريف کړو چې د \LR{$!A \lif \lfalse$} جوړونه~\LR{$h$}
واخلي او د~\LR{$\lfalse$} جوړونه ورکړي۔ \LR{$g$} داسې تعريفولے شو: ورکړل شوے
\LR{$h$} د \LR{$!A$} پر هغې جوړونې~\LR{$M$} تطبيق کړئ چې مخکې مو اخيستې وه۔ ځکه د
\LR{$h$} پايله~\LR{$h(M)$} د \LR{$\lfalse$} جوړونه ده، نو که \LR{$M$} د~\LR{$!A$} جوړونه
وي، \LR{$f(M)(h) = h(M)$} د~\LR{$\lfalse$} جوړونه ده۔
\end{ex}
\olpseganchor{OLP-0495-B011}{end}

\olpseganchor{OLP-0495-B012}{start}
\begin{ex}
د \LR{$\lnot(!A \land \lnot !A)$} لپاره يوه جوړونه ورکوو، يعنې د \LR{$(!A
\land (!A \lif \lfalse)) \lif \lfalse$} لپاره۔ دا داسې تابع~\LR{$f$}
دے چې د \LR{$!A \land (!A \lif \lfalse)$} جوړونه~\LR{$M$} واخلي او د
~\LR{$\lfalse$} جوړونه ورکړي۔ د عطف \LR{$!B_1 \land !B_2$} جوړونه يوه
جوړه \LR{$\tuple{N_1, N_2}$} ده چې \LR{$N_1$} د~\LR{$!B_1$} او \LR{$N_2$} د~\LR{$!B_2$}
جوړونه وي۔ موږ \LR{$p_1$} او \LR{$p_2$} داسې تابعې تعريفولے شو چې د
\LR{$!B_1 \land !B_2$} له جوړونې څخه په ترتيب د \LR{$!B_1$} او \LR{$!B_2$}
جوړونې راوباسي:
\begin{align*}
  p_1(\tuple{N_1, N_2}) &= N_1\\
  p_2(\tuple{N_1, N_2}) &= N_2
\end{align*}
تابع \LR{$f$} داسې کار کوي: لومړے \LR{$p_1$} په خپل ورودي~\LR{$M$} تطبيقوي۔
په دې سره د~\LR{$!A$} جوړونه لاس ته راځي۔ بيا \LR{$p_2$} په~\LR{$M$}
تطبيقوي او د~\LR{$!A \lif \lfalse$} جوړونه ترې لاس ته راځي۔ دا
جوړونه پخپله يو تابع~\LR{$p_2(M)$} دے چې د~\LR{$!A$} جوړونه واخلي او د
~\LR{$\lfalse$} جوړونه ورکړي۔ په بله وينا، که \LR{$p_2(M)$} په \LR{$p_1(M)$}
تطبيق کړو، د~\LR{$\lfalse$} جوړونه تر لاسه کوو۔ نو
\LR{$f(M) = p_2(M)(p_1(M))$} تعريفولے شو۔
\end{ex}
\olpseganchor{OLP-0495-B012}{end}

\olpseganchor{OLP-0495-B013}{start}
\begin{ex}
د \LR{$((!A \land !B) \lif !C) \lif (!A \lif
(!B \lif !C))$} جوړونه ورکوو۔ يعنې داسې تابع \LR{$f$} ورکوو چې د
\LR{$(!A \land !B) \lif !C$} جوړونه~\LR{$g$} د \LR{$(!A \lif (!B \lif
!C))$} په جوړونه واړوي۔ جوړونه \LR{$g$} پخپله يو تابع دے: د
\LR{$!A \land !B$} جوړونې د~\LR{$!C$} په جوړونو اړوي۔ د \LR{$f(g)$} پايله
يو تابع~\LR{$h_g$} دے چې د \LR{$!A$} جوړونې د هغو تابعو په جوړونو اړوي
چې د~\LR{$!B$} جوړونې د~\LR{$!C$} په جوړونو اړوي۔
ورپسې دا تابع د دوو پرله‌پسې ورودي جوړونو په کارولو ټاکو۔
\olpseganchor{OLP-0495-B013}{end}

\olpseganchor{OLP-0495-B014}{start}
اوس \LR{$h_g$} د ورودي~\LR{$g$} لپاره د \LR{$f$} د پايلې په توګه جوړوو۔ د~\LR{$h_g$} ورودي د~\LR{$!A$} جوړونه~\LR{$M$} ده۔ د
\LR{$h_g(M)$} پايله بايد داسې تابع~\LR{$k_M$} وي چې د~\LR{$!B$} جوړونې~\LR{$N$}
د~\LR{$!C$} په جوړونو اړوي۔ څرنګه چې لومړنۍ تابع ټاکلې ده، لنډ
شاخص يې د هغې ټاکنه په ضمني ډول ښيي۔ په بشپړ شاخص سره
\LR{$k_{g,M}(N) = g(\tuple{M, N})$} وټاکئ۔
په ياد ولرئ چې \LR{$\tuple{M, N}$} د~\LR{$!A \land !B$} جوړونه ده۔ نو
\LR{$k_{g,M}$} د \LR{$!B \lif !C$} جوړونه ده: د~\LR{$!B$} جوړونې~\LR{$N$}
د~\LR{$!C$} په جوړونو اړوي۔ اوس \LR{$h_g(M) = k_{g,M}$} وټاکئ۔ دا داسې
تابع دے چې د~\LR{$!A$} جوړونې~\LR{$M$} د \LR{$!B
\lif !C$} په جوړونو~\LR{$k_{g,M}$} اړوي۔ اوس \LR{$f(g) = h_g$} وټاکئ۔
دا داسې تابع ده چې د \LR{$(!A \land !B) \lif !C$} جوړونې~\LR{$g$}
د \LR{$!A \lif (!B \lif !C)$} په جوړونو اړوي۔ همدا مطلوبه جوړونه ده۔
\end{ex}
\olpseganchor{OLP-0495-B014}{end}

\olpseganchor{OLP-0495-B015}{start}
جمله \LR{$!A \lor \lnot !A$} د منځني حالت د ايستلو اصل بلل کېږي۔ د ځينو
ځانګړو \LR{$!A$} لپاره ئې ثابتولے شو، لکه \LR{$\lfalse \lor
\lnot \lfalse$}، خو عمومي طرحه په شهودي منطق کښې د ثبوت وړ نۀ ده۔
د ګولډباخ حدسيه ښيي چې د فصل د هېڅ يوې خوا جوړونه ښايي
زموږ په لاس کښې نۀ وي۔ د عمومي طرحې د ناسم اعتبار شاهد په راتلونکي
باب کښې د کرېپکي د معناپوهنې له مخې جوړېږي۔ دوه پرله‌پسې نړۍ
واخلئ چې يو قضيېيز متغير يوازې په دويمه کښې رښتينے وي۔ په
لومړۍ نړۍ کښې متغير نۀ تحميلېږي، او د دويمې نړۍ له امله يې
نفي هم نۀ تحميلېږي۔ سره له دې، د طرحې دوه‌ګونې نفي ثابتېږي،
يعنې \LR{$\lnot \lnot (!A
\lor \lnot !A)$} رښتيا دے۔
\olpseganchor{OLP-0495-B015}{end}

\olpseganchor{OLP-0495-B016}{start}
\begin{ex}
د \LR{$\lnot \lnot (!A \lor \lnot !A)$} د ثبوت لپاره داسې تابع~\LR{$f$}
غواړو چې د \LR{$\lnot(!A \lor \lnot !A)$} جوړونه، يعنې د \LR{$(!A
\lor (!A \lif \lfalse)) \lif \lfalse$} جوړونه، د~\LR{$\lfalse$}
په جوړونه واړوي۔ په بله وينا، داسې تابع \LR{$f$} غواړو چې که \LR{$g$} د
\LR{$\lnot(!A \lor \lnot !A)$} جوړونه وي، نو \LR{$f(g)$} د~\LR{$\lfalse$}
جوړونه وي۔
\olpseganchor{OLP-0495-B016}{end}

\olpseganchor{OLP-0495-B017}{start}
فرض کړئ \LR{$g$} د \LR{$\lnot(!A \lor \lnot !A)$} جوړونه ده؛ يعنې داسې
تابع چې د \LR{$!A \lor \lnot !A$} جوړونه د~\LR{$\lfalse$} په جوړونه
اړوي۔ د \LR{$!A \lor \lnot !A$} جوړونه يوه جوړه \LR{$\tuple{s, M}$} ده
چې يا \LR{$s = 1$} وي او \LR{$M$} د~\LR{$!A$} جوړونه وي، يا \LR{$s = 2$} وي او
\LR{$M$} د \LR{$\lnot !A$} جوړونه وي۔ \LR{$h_1$} هغه تابع وټاکئ چې د \LR{$!A$}
جوړونه~\LR{$M_1$} د \LR{$!A \lor \lnot !A$} په جوړونه اړوي: \LR{$M_1$}
د \LR{$\tuple{1, M_1}$} جوړې ته لېږي۔ \LR{$h_2$} هغه تابع وټاکئ چې د
\LR{$\lnot !A$} جوړونه~\LR{$M_2$} د \LR{$!A \lor \lnot !A$} په جوړونه
اړوي: \LR{$M_2$} د \LR{$\tuple{2, M_2}$} جوړې ته لېږي۔
هره تابع خپل ورودي ثبوت له هماغه اړوند نښان سره يوځاے کوي۔
\olpseganchor{OLP-0495-B017}{end}

\olpseganchor{OLP-0495-B018}{start}
\LR{$k$} د \LR{$\comp{h_1}{g}$} په توګه وټاکئ: دا داسې تابع دے چې د~\LR{$!A$}
جوړونه واخلي او د~\LR{$\lfalse$} جوړونه ورکړي؛ يعنې د
~\LR{$!A \lif \lfalse$} يا \LR{$\lnot !A$} جوړونه ده۔ اوس \LR{$l$} د
\LR{$\comp{h_2}{g}$} په توګه وټاکئ۔ دا داسې تابع ده چې د \LR{$\lnot !A$}
جوړونه واخلي او د~\LR{$\lfalse$} جوړونه ورکړي۔ څرنګه چې \LR{$k$} د
\LR{$\lnot !A$} جوړونه ده، نو \LR{$l(k)$} د~\LR{$\lfalse$} جوړونه ده۔
\olpseganchor{OLP-0495-B018}{end}

\olpseganchor{OLP-0495-B019}{start}
په ګډه مو وښودله چې د \LR{$\lnot(!A \lor \lnot !A)$} جوړونه~\LR{$g$}
څنګه د~\LR{$\lfalse$} په جوړونه اړولے شو۔ يعنې هغه تابع \LR{$f$} چې د
\LR{$\lnot(!A \lor \lnot !A)$} جوړونه~\LR{$g$} د~\LR{$\lfalse$} په
جوړونه \LR{$l(k)$} اړوي، د \LR{$\lnot\lnot(!A \lor \lnot !A)$} جوړونه ده۔
\end{ex}
\olpseganchor{OLP-0495-B019}{end}

\olpseganchor{OLP-0495-B020}{start}
د بي‌اېچ‌کې تفسير د فارمولونه جوړونې ښيي۔ په راتلونکو
برخو کښې به د شهودي منطق رسمي اشتقاق نظامونه او د
اعتبار د څېړلو لپاره دقيق معنايي تفسيرونه وکاروو۔
\olpseganchor{OLP-0495-B020}{end}

\olpseganchor{OLP-0496-B004}{start}
\olfileid{int}{int}{ntd}
\olpseganchor{OLP-0496-B004}{end}

\olpseganchor{OLP-0496-B005}{start}
\olsection{طبيعي استنتاج}
\olpseganchor{OLP-0496-B005}{end}

\olpseganchor{OLP-0496-B006}{start}
له \LR{$\FalseCl$} قاعدو پرته طبيعي استنتاج د شهودي منطق يو معياري
اشتقاق نظام دے۔ دلته قاعدې بيا راوړو او د BHK تفسير له مخې
ئې انګېزه څرګندوو۔ په هر حالت کښې قاعده داسې ګڼلے شو چې دا پايله
راوباسي: که مقدمې جوړونې ولري، نو نتيجه ئې هم جوړونه لري۔
\olpseganchor{OLP-0496-B006}{end}

\olpseganchor{OLP-0496-B007}{start}
څرنګه چې د طبيعي استنتاج په اشتقاقونه کښې نۀ ختم شوي
فرضونه وي، نو د بېلګې په توګه له نۀ ايستل شوې فرضونو~\LR{$\Gamma$}
څخه د \LR{$!A$} داسې اشتقاق يوه تابع ګڼو۔ ورودي ئې د
هر \LR{$!B \in \Gamma$} لپاره يوه جوړونه لري، او پايله ئې د~\LR{$!A$}
جوړونه ده۔ که د~\LR{$!A$}
داسې اشتقاق وي چې هېڅ نۀ ايستل شوې فرض ونۀ لري،
نو د BHK تفسير په معنا د~\LR{$!A$} جوړونه شته۔ خو د بحث د اسانۍ لپاره
به چې اړتيا نۀ وي، \LR{$\Gamma$} پرېږدو۔
\olpseganchor{OLP-0496-B007}{end}

\olpseganchor{OLP-0496-B008}{start}
يو فرض~\LR{$!A$} پخپله له نۀ ايستل شوې فرض~\LR{$!A$} څخه د~\LR{$!A$}
اشتقاق دے۔ دا له BHK تفسير سره سمون لري: د جوړونو عيني
تابع د~\LR{$!A$} هره جوړونه د~\LR{$!A$} په جوړونه اړوي۔
\olpseganchor{OLP-0496-B008}{end}

\olpseganchor{OLP-0496-B009}{start}
\subsection{عطف}
\olpseganchor{OLP-0496-B009}{end}

\olpseganchor{OLP-0496-B010}{start}
\begin{defish}
\AxiomC{\LR{$!A$}}
\AxiomC{\LR{$!B$}}
\RightLabel{\Intro{\land}}
\BinaryInfC{\LR{$!A \land !B$}}
\DisplayProof
\hfill
\begin{tabular}{l}
\AxiomC{\LR{$!A \land !B$}}
\RightLabel{\Elim{\land}}
\UnaryInfC{\LR{$!A$}}
\DisplayProof
\\[3ex]
\AxiomC{\LR{$!A \land !B$}}
\RightLabel{\Elim{\land}}
\UnaryInfC{\LR{$!B$}}
\DisplayProof
\end{tabular}
\end{defish}
\olpseganchor{OLP-0496-B010}{end}

\olpseganchor{OLP-0496-B011}{start}
فرض کړئ چې په ترتيب د \LR{$!A_1$} او \LR{$!A_2$} جوړونې \LR{$N_1$} او \LR{$N_2$}
لرو۔ نو د \LR{$!A_1 \land !A_2$} جوړونه هم لرو، يعنې جوړه
\LR{$\tuple{N_1, N_2}$}۔
\olpseganchor{OLP-0496-B011}{end}

\olpseganchor{OLP-0496-B012}{start}
د BHK تفسير له مخې د \LR{$!A_1 \land !A_2$} جوړونه يوه جوړه
\LR{$\tuple{N_1, N_2}$} ده۔ نو فرض کړئ چې داسې جوړه لرو۔ بيا د عطف
د دواړو غړو جوړونې هم لرو: \LR{$N_1$} د~\LR{$!A_1$} او \LR{$N_2$} د~\LR{$!A_2$}
جوړونه ده۔ هره ايستونکې قاعده د جوړې هماغه اړوند غړے اخلي۔
\olpseganchor{OLP-0496-B012}{end}

\olpseganchor{OLP-0496-B013}{start}
\subsection{شرط}
\olpseganchor{OLP-0496-B013}{end}

\olpseganchor{OLP-0496-B014}{start}
\begin{defish}
  \AxiomC{\LR{$\Discharge{!A}{u}$}}
  \DeduceC{\LR{$!B$}}
  \DischargeRule{\LR{$\Intro{\lif}$}}{u}
  \UnaryInfC{\LR{$!A \lif !B$}}
  \DisplayProof
\hfill
  \AxiomC{\LR{$!A \lif !B$}}
  \AxiomC{\LR{$!A$}}
  \RightLabel{\LR{$\Elim{\lif}$}}
  \BinaryInfC{\LR{$!B$}}
  \DisplayProof
\end{defish}
\olpseganchor{OLP-0496-B014}{end}

\olpseganchor{OLP-0496-B015}{start}
که له نۀ ايستل شوې فرض~\LR{$!A$} څخه د~\LR{$!B$} اشتقاق
ولرو، نو داسې تابع~\LR{$f$} شته چې د~\LR{$!A$} جوړونې د~\LR{$!B$} په جوړونو
اړوي۔ همدا تابع د \LR{$!A \lif !B$} جوړونه ده۔ نو که د \LR{$\Intro\lif$}
مقدمې ته د~\LR{$!A$} پر جوړونه مشروطه جوړونه وي، نتيجه \LR{$!A
\lif !B$} هم جوړونه لري۔
\olpseganchor{OLP-0496-B015}{end}

\olpseganchor{OLP-0496-B016}{start}
بل خوا فرض کړئ چې \LR{$N$} د \LR{$!A$} جوړونه او \LR{$f$} د \LR{$!A \lif !B$}
جوړونه ده۔ د \LR{$!A \lif !B$} جوړونه داسې تابع ده چې د \LR{$!A$}
جوړونې د~\LR{$!B$} په جوړونو اړوي۔ نو \LR{$f(N)$} د~\LR{$!B$} جوړونه ده؛
يعنې د \LR{$\Elim\lif$} نتيجه جوړونه لري۔
\olpseganchor{OLP-0496-B016}{end}

\olpseganchor{OLP-0496-B017}{start}
\subsection{فصل}
\olpseganchor{OLP-0496-B017}{end}

\olpseganchor{OLP-0496-B018}{start}
\begin{defish}
\begin{tabular}{l}
\AxiomC{\LR{$!A$}}
\RightLabel{\Intro{\lor}}
\UnaryInfC{\LR{$!A \lor !B$}}
\DisplayProof
\\[3ex]
\AxiomC{\LR{$!B$}}
\RightLabel{\Intro{\lor}}
\UnaryInfC{\LR{$!A \lor !B$}}
\DisplayProof
\end{tabular}
\hfill
\AxiomC{\LR{$!A \lor !B$}}
\AxiomC{\LR{$\Discharge{!A}{n}$}}
\DeduceC{\LR{$!C$}}
\AxiomC{\LR{$\Discharge{!B}{n}$}}
\DeduceC{\LR{$!C$}}
\DischargeRule{\Elim{\lor}}{n}
\TrinaryInfC{\LR{$!C$}}
\DisplayProof
\end{defish}
\olpseganchor{OLP-0496-B018}{end}

\olpseganchor{OLP-0496-B019}{start}
که د~\LR{$!A_i$} جوړونه~\LR{$N_i$} ولرو، هغې ته د~\LR{$!A_1 \lor !A_2$}
جوړونې~\LR{$\tuple{i, N_i}$} بڼه ورکولے شو۔ بل خوا فرض کړئ چې د
~\LR{$!A_1 \lor !A_2$} جوړونه لرو، يعنې يوه جوړه~\LR{$\tuple{i, N_i}$}
چې \LR{$N_i$} د~\LR{$!A_i$} جوړونه ده۔ همدارنګه تابعې \LR{$f_1$} او \LR{$f_2$}
لرو چې په ترتيب د \LR{$!A_1$} او \LR{$!A_2$} جوړونې د~\LR{$!C$} په جوړونو
اړوي۔ نو \LR{$f_i(N_i)$} د~\LR{$!C$} جوړونه، يعنې د~\LR{$\Elim\lor$} نتيجه ده۔
\olpseganchor{OLP-0496-B019}{end}

\olpseganchor{OLP-0496-B020}{start}
\subsection{محال}
\olpseganchor{OLP-0496-B020}{end}

\olpseganchor{OLP-0496-B021}{start}
\begin{defish}
  \AxiomC{\LR{$\lfalse$}}
  \RightLabel{\FalseInt}
  \UnaryInfC{\LR{$!A$}}
  \DisplayProof
\end{defish}
\olpseganchor{OLP-0496-B021}{end}

\olpseganchor{OLP-0496-B022}{start}
که له نۀ ايستل شوې فرضونو~\LR{$!B_1$}، \dots، \LR{$!B_n$} څخه د
~\LR{$\lfalse$} اشتقاق ولرو، نو تابع~\LR{$f(M_1,
\dots, M_n)$} د \LR{$!B_1$}، \dots، \LR{$!B_n$} جوړونې د~\LR{$\lfalse$}
په جوړونه اړوي۔ څرنګه چې \LR{$\lfalse$} هېڅ جوړونه نۀ لري، نو د
\LR{$!B_1$}، \dots، \LR{$!B_n$} ټولو جوړونې هم يوځاے نۀ شي موجودېدے۔
له دې امله د همدې ورودي جوړونو له تشې ساحې څخه يوې بلې
تابعې \LR{$g$} ته مخه کوو۔ \emph{که} \LR{$M_1$}، \dots، \LR{$M_n$} په ترتيب
د \LR{$!B_1$}، \dots، \LR{$!B_n$} جوړونې وي، \emph{نو}
\LR{$g(M_1, \dots, M_n)$} د~\LR{$!A$} جوړونه ده۔ دا تابع د محال له
جوړونې څخه د مطلوب فارمول جوړونې ته له تشې تابعې سره د
لومړۍ تابعې په يوځاے کولو ترلاسه کېږي۔
\olpseganchor{OLP-0496-B022}{end}

\olpseganchor{OLP-0496-B023}{start}
\subsection{د \texorpdfstring{\LR{$\lnot$}}{¬} قاعدې}
\olpseganchor{OLP-0496-B023}{end}

\olpseganchor{OLP-0496-B024}{start}
ځکه چې \LR{$\lnot !A$} د \LR{$!A \lif \lfalse$} په توګه تعريفېږي، په
دقيقه معنا د~\LR{$\lnot$} قاعدو ته اړتيا نۀ لرو۔ خو که ئې وغواړو،
بڼه به ئې داسې وي:
\olpseganchor{OLP-0496-B024}{end}

\olpseganchor{OLP-0496-B025}{start}
\begin{defish}
\AxiomC{\LR{$\Discharge{!A}{n}$}}
\noLine
\DeduceC{\LR{$\lfalse$}}
\DischargeRule{\Intro{\lnot}}{n}
\UnaryInfC{\LR{$\lnot !A$}}
\DisplayProof
\hfill
\AxiomC{\LR{$\lnot !A$}}
\AxiomC{\LR{$!A$}}
\RightLabel{\Elim{\lnot}}
\BinaryInfC{\LR{$\lfalse$}}
\DisplayProof
\end{defish}
\olpseganchor{OLP-0496-B025}{end}


\olpseganchor{OLP-0496-B026}{start}
\subsection{د اشتقاقونه بېلګې}
\olpseganchor{OLP-0496-B026}{end}

\olpseganchor{OLP-0496-B027}{start}
\begin{enumerate}
\item \LR{$\Proves !A \lif (\lnot !A \lif \lfalse)$}، يعنې \LR{$\Proves !A
  \lif ((!A \lif \lfalse) \lif \lfalse)$}
  \begin{prooftree}
    \AxiomC{\LR{$\Discharge{!A}{2}$}}
    \AxiomC{\LR{$\Discharge{!A \lif \lfalse}{1}$}}
    \RightLabel{\Elim\lif}
    \BinaryInfC{\LR{$\lfalse$}}
    \DischargeRule{\Intro\lif}{1}
    \UnaryInfC{\LR{$(!A \lif \lfalse)\lif \lfalse$}}
    \DischargeRule{\Intro\lif}{2}
    \UnaryInfC{\LR{$!A \lif (!A \lif \lfalse) \lif \lfalse$}}
  \end{prooftree}
\olpseganchor{OLP-0496-B027}{end}

\olpseganchor{OLP-0496-B028}{start}
\item \LR{$\Proves ((!A \land !B) \lif !C) \lif (!A \lif (!B \lif !C))$}
  \begin{prooftree}
    \AxiomC{\LR{$\Discharge{(!A \land !B) \lif !C}{3}$}}
    \AxiomC{\LR{$\Discharge{!A}{2}$}}
    \AxiomC{\LR{$\Discharge{!B}{1}$}}
    \RightLabel{\Intro\land}
    \BinaryInfC{\LR{$!A \land !B$}}
    \RightLabel{\Elim\lif}
    \BinaryInfC{\LR{$!C$}}
    \DischargeRule{\Intro\lif}{1}
    \UnaryInfC{\LR{$!B \lif !C$}}
    \DischargeRule{\Intro\lif}{2}
    \UnaryInfC{\LR{$!A \lif (!B \lif !C)$}}
    \DischargeRule{\Intro\lif}{3}
    \UnaryInfC{\LR{$((!A \land !B) \lif !C) \lif (!A \lif (!B \lif !C))$}}
  \end{prooftree}
\olpseganchor{OLP-0496-B028}{end}

\olpseganchor{OLP-0496-B029}{start}
\item \LR{$\Proves \lnot(!A \land \lnot !A)$}، يعنې \LR{$\Proves (!A \land (!A
  \lif \lfalse))\lif \lfalse$}
  \begin{prooftree}
    \AxiomC{\LR{$\Discharge{!A \land (!A \lif \lfalse)}{1}$}}
    \RightLabel{\Elim\land}
    \UnaryInfC{\LR{$!A \lif \lfalse$}}
    \AxiomC{\LR{$\Discharge{!A \land (!A \lif \lfalse)}{1}$}}
    \RightLabel{\Elim\land}
    \UnaryInfC{\LR{$!A$}}
    \RightLabel{\Elim\lif}
    \BinaryInfC{\LR{$\lfalse$}}
    \DischargeRule{\Intro\lif}{1}
    \UnaryInfC{\LR{$(!A \land (!A \lif \lfalse))\lif \lfalse$}}
  \end{prooftree}
\olpseganchor{OLP-0496-B029}{end}

\olpseganchor{OLP-0496-B030}{start}
\item \LR{$\Proves \lnot\lnot(!A \lor \lnot !A)$}، يعنې \LR{$\Proves ((!A \lor
  (!A \lif \lfalse))\lif \lfalse)\lif \lfalse$}
  \begin{prooftree}\hskip-3em
    \AxiomC{\LR{$\Discharge{(!A \lor (!A \lif \lfalse))\lif \lfalse}{2}$}}
    \AxiomC{\LR{$\Discharge{(!A \lor (!A \lif \lfalse))\lif \lfalse}{2}$}}
    \AxiomC{\LR{$\Discharge{!A}{1}$}}
    \RightLabel{\Intro\lor}
    \UnaryInfC{\LR{$!A \lor (!A \lif \lfalse)$}}
    \RightLabel{\Elim\lif}
    \BinaryInfC{\LR{$\lfalse$}}
    \DischargeRule{\Intro\lif}{1}
    \UnaryInfC{\LR{$!A \lif \lfalse$}}
    \RightLabel{\Intro\lor}
    \UnaryInfC{\LR{$!A \lor (!A \lif \lfalse)$}}
    \RightLabel{\Elim\lif}
    \insertBetweenHyps{\hskip -4em}
    \BinaryInfC{\LR{$\lfalse$}}
    \DischargeRule{\Intro\lif}{2}
    \UnaryInfC{\LR{$((!A \lor (!A \lif \lfalse))\lif \lfalse)\lif \lfalse$}}
  \end{prooftree}
\end{enumerate}
\olpseganchor{OLP-0496-B030}{end}

\olpseganchor{OLP-0496-B031}{start}
\begin{prop}
  که په شهودي منطق کښې د طبيعي استنتاج له مخې \LR{$\Gamma \Proves !A$}
  وي، نو په کلاسيک منطق کښې هم \LR{$\Gamma \Proves !A$} وي۔ په تېره بيا،
  که \LR{$!A$} شهودي قضيه وي، نو کلاسيکه قضيه هم ده۔
\end{prop}
\olpseganchor{OLP-0496-B031}{end}

\olpseganchor{OLP-0496-B032}{start}
\begin{proof}
  د طبيعي استنتاج هره قاعده په کلاسيک طبيعي استنتاج کښې هم قاعده ده؛
  نو د شهودي منطق هر اشتقاق په کلاسيک منطق کښې هم
  اشتقاق دے۔
\end{proof}
\olpseganchor{OLP-0496-B032}{end}

\olpseganchor{OLP-0496-B033}{start}
\begin{prob}
  په شهودي منطق کښې د لاندې فارمولونه اشتقاقونه ورکړئ:
  \begin{enumerate}
    \item \LR{$(\lnot !A \lor !B) \lif (!A \lif !B)$}
    \item \LR{$\lnot\lnot\lnot !A \lif \lnot !A$}
    \item \LR{$\lnot\lnot (!A \land !B) \liff (\lnot\lnot !A\land \lnot \lnot !B)$}
    \item \LR{$\lnot (!A \lor !B) \liff (\lnot !A \land \lnot !B)$}
    \item \LR{$(\lnot !A \lor \lnot !B) \lif \lnot (!A \land !B)$}
    \item \LR{$\lnot\lnot ( !A \land !B) \lif  (\lnot\lnot !A \lor
    \lnot\lnot !B)$}
  \end{enumerate}
\end{prob}
\olpseganchor{OLP-0496-B033}{end}

\olpseganchor{OLP-0497-B004}{start}
\olfileid{int}{int}{axd}
\olpseganchor{OLP-0497-B004}{end}

\olpseganchor{OLP-0497-B005}{start}
\olsection{بديهي اشتقاقونه}
\olpseganchor{OLP-0497-B005}{end}

\olpseganchor{OLP-0497-B006}{start}
د شهودي قضييز منطق بديهي اشتقاقونه د تصور له پلوه تر ټولو ساده
او په تاريخي لحاظ لومړني اشتقاق نظامونه دي۔ د ثبوت
متناهي لړۍ او د مودوس پوننس استعمال د کلاسيک قضييز منطق په
څېر دي۔ د بديهي شېماګانو ټاکنه د شهودي منطق له مخې کېږي۔
\olpseganchor{OLP-0497-B006}{end}

\olpseganchor{OLP-0497-B007}{start}
\begin{defn}[د اشتقاق وړتيا]
که \LR{$\Gamma$} د \LR{$\Lang L$} د فارمولونه يو سټ وي، نو له
\LR{$\Gamma$} څخه \emph{اشتقاق} د فارمولونه يوه متناهي
لړۍ \LR{$!A_1$}، \dots،~\LR{$!A_n$} ده چې د هر \LR{$i \le n$} لپاره لاندې
يو شرط پوره کېږي:
\begin{enumerate}
\item \LR{$!A_i \in \Gamma$} وي؛ يا
\item \LR{$!A_i$} يو بديهي اصل وي؛ يا
\item \LR{$!A_i$} د ځينو \LR{$!A_j$} او \LR{$!A_k$} څخه، چې \LR{$j < i$} او \LR{$k <
  i$} وي، د مودوس پوننس له مخې راشي؛ يعنې \LR{$!A_k \ident !A_j \lif !A_i$}۔
\end{enumerate}
\end{defn}
\olpseganchor{OLP-0497-B007}{end}

\olpseganchor{OLP-0497-B008}{start}
\begin{defn}[بديهي اصول]
د شهودي قضييز منطق د \emph{بديهي اصولو} سټ~\LR{$\PAx$} د لاندې
بڼو ټول فارمولونه دي:
\begin{align}
  & (!A \land !B) \lif !A \ollabel{ax:land1}\\
  & (!A \land !B) \lif !B \ollabel{ax:land2}\\
  & !A \lif (!B \lif (!A \land !B)) \ollabel{ax:land3}\\
  & !A \lif (!A \lor !B) \ollabel{ax:lor1}\\
  & !A \lif (!B \lor !A) \ollabel{ax:lor2}\\
  & (!A \lif !C) \lif ((!B \lif !C) \lif ((!A \lor !B) \lif !C)) \ollabel{ax:lor3}\\
  & !A \lif (!B \lif !A) \ollabel{ax:lif1}\\
  & (!A \lif (!B \lif !C)) \lif ((!A \lif !B) \lif (!A \lif !C)) \ollabel{ax:lif2}\\
  & \lfalse \lif !A \ollabel{ax:lfalse1}
\end{align}
\end{defn}
\olpseganchor{OLP-0497-B008}{end}

\olpseganchor{OLP-0497-B009}{start}
\begin{defn}[د اشتقاق وړتيا]
فارمول~\LR{$!A$} له \LR{$\Gamma$} څخه هله \emph{د اشتقاق وړ}
ګڼل کېږي، او دا په \LR{$\Gamma \Proves !A$} ليکو، چې له
\LR{$\Gamma$} څخه داسې اشتقاق وي چې په \LR{$!A$} پاے ته رسېږي۔
\end{defn}
\olpseganchor{OLP-0497-B009}{end}

\olpseganchor{OLP-0497-B010}{start}
\begin{defn}[قضيې]
فارمول~\LR{$!A$} هله \emph{قضيه} ده چې له تش سټ څخه د~\LR{$!A$}
اشتقاق وي۔ که \LR{$!A$} قضيه وي، \LR{$\Proves !A$} ليکو، او که
نه وي، \LR{$\Proves/ !A$} ليکو۔
\end{defn}
\olpseganchor{OLP-0497-B010}{end}

\olpseganchor{OLP-0497-B011}{start}
\begin{prop}
  که د شهودي منطق په بديهي اشتقاق نظام کښې
  \LR{$\Gamma \Proves !A$} وي، نو په کلاسيک منطق کښې هم
  \LR{$\Gamma \Proves !A$} وي۔ په تېره بيا، که \LR{$!A$} شهودي قضيه وي،
  نو کلاسيکه قضيه هم ده۔
\end{prop}
\olpseganchor{OLP-0497-B011}{end}

\olpseganchor{OLP-0497-B012}{start}
\begin{proof}
  هر شهودي بديهي اصل په کلاسيک منطق کښې معتبر دے، او په هر
  بشپړ کلاسيک حساب کښې ثابتېږي۔ د شهودي اشتقاق په هره کرښه
  کښې يا هماغه مقدمه اخلو، يا د بديهي اصل کلاسيک ثبوت کاروو،
  او يا مودوس پوننس پلي کوو۔ په دې توګه کلاسيک اشتقاق
  ترلاسه کوو۔
\end{proof}
\olpseganchor{OLP-0497-B012}{end}

\olpseganchor{OLP-0498-B004}{start}
\olchapter{int}{sem}{معنايي پوهنه}
\olpseganchor{OLP-0498-B004}{end}

\olpseganchor{OLP-0498-B005}{start}
\begin{editorial}
  دا باب د شهودي منطق دوه معنايي پوهنې بيانوي۔ د کرېپکي په
  مدلونو کښې رښتياوالی د پوهې له پراخېدو سره ساتل کېږي۔ په
  توپولوژيکي مدلونو کښې فارمولونه د پرانيستو سټونو په توګه ارزول کېږي۔
\end{editorial}
\olpseganchor{OLP-0498-B005}{end}

\olpseganchor{OLP-0499-B004}{start}
\olfileid{int}{sem}{int}
\olpseganchor{OLP-0499-B004}{end}

\olpseganchor{OLP-0499-B005}{start}
\olsection{سريزه}
\olpseganchor{OLP-0499-B005}{end}

\olpseganchor{OLP-0499-B006}{start}
له معنايي پوهنې پرته د هېڅ منطق بشپړ قناعت وړ تشريح نۀ کېږي، او
شهودي منطق هم استثنا نۀ دے۔ د کلاسيک منطق لپاره په ارزښت ټاکنې
ولاړه معنايي پوهنه معياري ده، خو د شهودي منطق لپاره څو بېلابېلې
معنايي پوهنې سيالي کوي۔ هېڅ يوه ئې په دې معنا بشپړه قناعت وړ نۀ ده
چې د نښلوونکو د معناوو د شهودي منلو وړ تشريح ورکړي۔ د دې بحث لپاره د جوړونې له پلوه تفسير او
د مدل پر بنسټ د رښتياوالي تعريف بېلې پوښتنې دي۔
\olpseganchor{OLP-0499-B006}{end}

\olpseganchor{OLP-0499-B007}{start}
د اړيکيز مدلونه پر بنسټ معنايي پوهنه، چې د وجهي منطقونو له
معنايي پوهنې سره ورته ده، ښايي تر ټولو ډېره دود وي۔ په دې معنايي
پوهنه کښې بياني متغيرونه نړۍو ته ټاکل کېږي او دغه
نړۍې د لاسرسي په اړيکه تړلې وي۔ په دې وړاندې‌کونه کښې دا اړيکه جزئي ترتيب وي، يعنې
انعکاسي، ضد متقارن او متعدي وي۔
\olpseganchor{OLP-0499-B007}{end}

\olpseganchor{OLP-0499-B008}{start}
په شهودي ډول دا نړۍې د پوهې حالتونه يا «د شواهدو حالتونه» ګڼلے شئ۔
حالت~\LR{$w'$} له~\LR{$w$} څخه هله او يوازې هله لاسرسي وړ دے چې زموږ د
اوسنۍ پوهې له مخې \LR{$w'$} د پوهې يو ممکن غځېدلے حالت وي؛ يعنې
هغه څۀ پکښې پاتې وي چې په~\LR{$w$} کښې معلوم دي۔ کله چې يوه قضيه
معلومه شي، بيا نامعلومه کېدے نۀ شي: که \LR{$!A$} په~\LR{$w$} کښې معلوم وي
او~\LR{$Rww'$} وي، نو \LR{$!A$} په~\LR{$w'$} کښې هم معلوم وي۔ نو «پوهه» د
لاسرسي د اړيکې له مخې يورخيزه زياتېدونکې ده۔
\olpseganchor{OLP-0499-B008}{end}

\olpseganchor{OLP-0499-B009}{start}
که د معرفتي منطق په څېر «\LR{$!A$} معلوم دے» داسې تعريف کړو چې «په ټولو
معرفتي بديلو کښې رښتيا دے»، نو \LR{$!A \land !B$} په~\LR{$w$} کښې هله
معلوم وي چې په ټولو معرفتي بديلو کښې \LR{$!A$} او~\LR{$!B$} دواړه معلوم وي۔
د انعکاسيت له مخې اوسنۍ نړۍ پخپله هم د لاسرسي وړ ده، او
د يورخيزې زياتېدنې له مخې هر معلوم غړے په ټولو دغو نړۍو کښې
معلوم پاتې کېږي۔ د \LR{$R$} په اړه له همدې دوو خاصيتونو راځي
چې \LR{$!A \land !B$} په \LR{$w$} کښې هله او يوازې هله معلوم وي چې \LR{$!A$}
او \LR{$!B$} په~\LR{$w$} کښې دواړه معلوم وي۔ د همدې دليل له مخې
\LR{$!A \lor !B$} په \LR{$w$} کښې هله او يوازې هله معلوم وي چې لږ تر لږه
يو غړی ئې معلوم وي۔ نو د \LR{$\land$} او \LR{$\lor$} د رښتياوالي شرطونه
د کلاسيک منطق له شرطونو سره يو شان دي۔
\olpseganchor{OLP-0499-B009}{end}

\olpseganchor{OLP-0499-B010}{start}
خو د شرطي د رښتياوالي شرطونه له کلاسيک منطق سره توپير لري۔
\LR{$!A \lif !B$} په~\LR{$w$} کښې هله او يوازې هله معلوم وي چې داسې هېڅ
\LR{$w'$}، چې \LR{$Rww'$} ولري، نۀ وي چې هلته \LR{$!A$} معلوم وي خو \LR{$!B$}
معلوم نۀ وي۔ يوازې د اوسنۍ نړۍ دا شرط کافي نۀ دے چې \LR{$!A$} نامعلوم وي
يا \LR{$!B$} په~\LR{$w$} کښې معلوم وي۔ ځکه که په~\LR{$w$} کښې نۀ \LR{$!A$}
او نۀ \LR{$!B$} معلوم وي، ښايي داسې راتلونکے معرفتي حالت~\LR{$w'$} وي
چې \LR{$Rww'$} ولري او په \LR{$w'$} کښې \LR{$!A$} معلوم شي خو~\LR{$!B$} معلوم نۀ شي۔
\olpseganchor{OLP-0499-B010}{end}

\olpseganchor{OLP-0499-B011}{start}
موږ \LR{$\lnot !A$} يوازې هغه وخت معلومولے شو چې داسې هېڅ ممکن
راتلونکے معرفتي حالت نۀ وي چې پکې~\LR{$!A$} معلوم وي۔ دلته تصور دا دے
چې که \LR{$!A$} د معلومېدو وړ وي، نو په کوم ممکن راتلونکي معرفتي حالت
کښې به~\LR{$!A$} معلوم شي۔ څرنګه چې \LR{$\lfalse$} معلومولے نۀ شو، نو
په هغه راتلونکي معرفتي حالت کښې به \LR{$!A$} معلوم وي خو~\LR{$\lfalse$}
به معلوم نۀ وي۔
\olpseganchor{OLP-0499-B011}{end}

\olpseganchor{OLP-0499-B012}{start}
په دې تفسير کښې د منځني حالت د ايستلو اصل نۀ چلېږي۔ ځکه ځينې
\LR{$!A$} شته چې لا مو نۀ دي معلوم، خو وروسته ښايي معلوم شي۔ د داسې
فارمول~\LR{$!A$} لپاره نۀ \LR{$!A$} معلوم دے او نۀ~\LR{$\lnot !A$}؛
نو \LR{$!A \lor \lnot !A$} معلوم نۀ دے۔ خو د بېلګې په توګه موږ
پوهيږو چې \LR{$\lnot(!A \land \lnot !A)$} دے۔ ځکه داسې راتلونکے
حالت ممکن نۀ دے چې پکې \LR{$!A$} او \LR{$\lnot !A$} دواړه معلوم وي؛
دا خبره له دې پرته هم معلومه ده چې~\LR{$!A$} يا \LR{$\lnot !A$} معلوموو
که نه۔
\olpseganchor{OLP-0499-B012}{end}

\olpseganchor{OLP-0499-B013}{start}
اړيکيز مدلونه د شهودي منطق يوازينۍ شونې معنايي پوهنه نۀ ده۔
توپولوژيکي معنايي پوهنه بله لاره ده: دلته قضیې په توپولوژيکي فضا
کښې د پرانيستو سټونو په توګه او نښلوونکي پر دغو سټونو د عملونو په
توګه تفسيرېږي (مثلاً \LR{$\land$} له تقاطع سره سمون لري)۔
\olpseganchor{OLP-0499-B013}{end}

\olpseganchor{OLP-0500-B004}{start}
\olfileid{int}{sem}{rel}
\olpseganchor{OLP-0500-B004}{end}

\olpseganchor{OLP-0500-B005}{start}
\olsection{اړيکيز مدلونه}
\olpseganchor{OLP-0500-B005}{end}

\olpseganchor{OLP-0500-B006}{start}
د شهودي قضييز منطق د دقيقې معنايي پوهنې لپاره بايد تعريف کړو چې
کوم څيز داسې مدل ګڼل کېږي چې فارمولونه د هغه له مخې وارزوو۔
د دغه تعريف پر بنسټ بيا د اعتبار او استلزام په شان معنايي مفاهيم
هم تعريفولے شو۔ يوه داسې معنايي پوهنه د اړيکيز مدلونه له
لارې ورکول کېږي۔
\olpseganchor{OLP-0500-B006}{end}

\olpseganchor{OLP-0500-B007}{start}
\begin{defn}
  د شهودي قضييز منطق اړيکيز مدل يوه درېګونې
  \LR{$\mModel{M} = \tuple{W, R, V}$} ده، چې پکې
  \begin{enumerate}
  \item \LR{$W$} يو ناتش سټ دے،
  \item \LR{$R$} پر~\LR{$W$} جزئي ترتيب دے (يعنې انعکاسي، ضد متقارن او
    متعديه دوه‌ځاييزه اړيکه)، او
  \item \LR{$V$} داسې تابع ده چې هر بياني متغير~\LR{$p$} ته د~\LR{$W$} يو فرعي سټ ټاکي،
  \item \LR{$V$} د \LR{$R$} له مخې يورخيزه زياتېدونکې وي، يعنې که \LR{$w
    \in V(p)$} او \LR{$Rww'$} وي، نو \LR{$w' \in V(p)$} وي۔
  \end{enumerate}
\end{defn}
\olpseganchor{OLP-0500-B007}{end}

\olpseganchor{OLP-0500-B008}{start}
\begin{defn}\ollabel{defn:true-at-w}
  دا مفهوم چې \LR{$!A$} \emph{په \LR{$\mModel{M}$} کښې په \LR{$w$} رښتيا دے}،
  چې په \LR{$\mSat{M}{!A}[w]$} ئې ليکو، په لاندې ډول په استقرايي توګه
  تعريفوو:
  \begin{enumerate}
  \item \indcase{!A}{p}{\LR{$\mSat{M}{\indfrm}[w]$} هله او يوازې هله چې \LR{$w \in V(p)$} وي۔}
  \item \indcase{!A}{\lfalse}{\LR{$\mSat{M}{\indfrm}[w]$} نۀ وي}۔
  \item \indcase{!A}{\lnot !B}{\LR{$\mSat{M}{\indfrm}[w]$} هله او يوازې هله چې
    داسې هېڅ \LR{$w'$} نۀ وي چې \LR{$Rww'$} او \LR{$\mSat{M}{!B}[w']$} وي}۔
  \item \indcase{!A}{!B \land !C}{\LR{$\mSat{M}{\indfrm}[w]$} هله او يوازې هله چې
    \LR{$\mSat{M}{!B}[w]$} او \LR{$\mSat{M}{!C}[w]$} وي}۔
  \item \indcase{!A}{!B \lor !C}{\LR{$\mSat{M}{\indfrm}[w]$} هله او يوازې هله چې
    \LR{$\mSat{M}{!B}[w]$} يا \LR{$\mSat{M}{!C}[w]$} وي (يا دواړه)}۔
  \item \indcase{!A}{!B \lif !C}{\LR{$\mSat{M}{\indfrm}[w]$} هله او يوازې هله چې
    د هر داسې \LR{$w'$} لپاره چې \LR{$Rww'$} وي، يا \LR{$\mSat{M}{!B}[w']$}
    نۀ وي يا \LR{$\mSat{M}{!C}[w']$} وي} (يا دواړه)۔
  \end{enumerate}
  که \LR{$\mSat{M}{!A}[w]$} نۀ وي، \LR{$\mSat/{M}{!A}[w]$} ليکو۔ که
  \LR{$\Gamma$} د فارمولونه يو سټ وي، \LR{$\mSat{M}{\Gamma}[w]$} دا معنا
  لري چې د هر \LR{$!B \in \Gamma$} لپاره \LR{$\mSat{M}{!B}[w]$} وي۔
\end{defn}
\olpseganchor{OLP-0500-B008}{end}

\olpseganchor{OLP-0500-B009}{start}
\begin{prob}
  وښيئ چې د \olref[int][sem][rel]{defn:true-at-w} له مخې
  \LR{$\mSat{M}{\lnot !A}[w]$} هله او يوازې هله چې
  \LR{$\mSat{M}{!A \lif \lfalse}[w]$} وي۔
\end{prob}
\olpseganchor{OLP-0500-B009}{end}

\olpseganchor{OLP-0500-B010}{start}
\begin{prop}\ollabel{prop:true-monotonic}
  په نړۍو کښې رښتياوالے د~\LR{$R$} له مخې يورخيزه زياتېږي، يعنې که
  \LR{$\mSat{M}{!A}[w]$} او \LR{$Rww'$} وي، نو \LR{$\mSat{M}{!A}[w']$} وي۔
\end{prop}
\olpseganchor{OLP-0500-B010}{end}

\olpseganchor{OLP-0500-B011}{start}
\begin{proof}
  د فارمول پر جوړښت استقرا کوو۔ د متغير لپاره دعوه د ارزښت
  ټاکنې له يورخيزې زياتېدنې راځي، او محال په هېڅ نړۍ کښې
  رښتيا نۀ دے۔ په عطف او فصل کښې د غړو د استقرا فرض کاروو۔
  د نفي او شرط لپاره، له نوې نړۍ څخه هره د لاسرسي وړ نړۍ له
  لومړۍ نړۍ هم د لاسرسي وړ ده؛ دا د اړيکې له تعدي راځي۔ نو
  د نفي د رښتياوالي شرط او د شرط عمومي حکم دواړه په نوې نړۍ کښې
  هم پوره کېږي۔
\end{proof}
\olpseganchor{OLP-0500-B011}{end}

\olpseganchor{OLP-0500-B012}{start}
\begin{prob}
  \olref[int][sem][rel]{prop:true-monotonic} ثابت کړئ۔
\end{prob}
\olpseganchor{OLP-0500-B012}{end}

\paragraph{د منجمدې سرچينې بشپړه اړونده برخه}
\olpseganchor{OLP-0649-B004}{start}
\olfileid{int}{sem}{pro-olp0649}
\olpseganchor{OLP-0649-B004}{end}

\olpseganchor{OLP-0649-B005}{start}
\olsection{بيانونه}
\olpseganchor{OLP-0649-B005}{end}

\olpseganchor{OLP-0649-B006}{start}
ځينې وخت په اړيکيز مدل~\LR{$\mModel{M}$} کښې د هغو نړۍو سټ ته نوم ورکول ګټور وي چې فارمول~\LR{$!A$} پکې رښتيا وي۔ دغه سټ د \LR{$!A$} له خوا تعريف شوے \emph{بيان} بولو، او \LR{$\Prop{M}{!A}$} يې ليکو۔
\olpseganchor{OLP-0649-B006}{end}

\olpseganchor{OLP-0649-B007}{start}
\begin{defn}
تابعه \LR{$V$} په استقرايي ډول داسې تابعې \LR{$\Prop{M}{!A} \colon \Frm
  \to \Pow{W}$} ته غځوو چې هر فارمول ته د هغه د رښتينوالي د نړۍو سټ ګوماري:
  \begin{enumerate}
  \item \LR{$\Prop{M}{\lfalse} = \emptyset$}
  \item \LR{$\Prop{M}{p} = V(p)$}
  \item \LR{$\Prop{M}{\lnot !B} = {}$}
    \[
      \Setabs{w \in W}{\text{\RL{د هر داسې \LR{$w'$} دپاره چې \LR{$Rww'$} وي، \LR{$w'
          \notin \Prop{M}{!B}$}}}}.
    \]
  \item \LR{$\Prop{M}{!B \land !C} = \Prop{M}{!B} \cap \Prop{M}{!C}$}.
  \item \LR{$\Prop{M}{!B \lor !C} = \Prop{M}{!B} \cup \Prop{M}{!C}$}.
  \item \LR{$\Prop{M}{!B \lif !C} = {}$}
    \[
      \Setabs{w \in W}{\text{\RL{د هر داسې \LR{$w'$} دپاره چې \LR{$Rww'$} وي، که \LR{$w'
          \in \Prop{M}{!B}$} وي، نو \LR{$w' \in \Prop{M}{!C}$}}}}.
    \]
  \end{enumerate}
\end{defn}
\olpseganchor{OLP-0649-B007}{end}

\olpseganchor{OLP-0649-B008}{start}
\begin{prop}
  \ollabel{prop:propositions-true}
  \LR{$\mSat{M}{!A}[w]$} هغه او يوازې هغه وخت چې \LR{$w \in \Prop{M}{!A}$}۔
\end{prop}
\olpseganchor{OLP-0649-B008}{end}

\olpseganchor{OLP-0649-B009}{start}
\begin{proof}
  د تمرين په توګه پرېښودل شوے دے۔
\end{proof}
\olpseganchor{OLP-0649-B009}{end}

\olpseganchor{OLP-0649-B010}{start}
\begin{prob}
  \olref[int][sem][pro-olp0649]{prop:propositions-true} ثابت کړئ۔
\end{prob}
\olpseganchor{OLP-0649-B010}{end}

\olpseganchor{OLP-0501-B004}{start}
\olfileid{int}{sem}{sem}
\olpseganchor{OLP-0501-B004}{end}

\olpseganchor{OLP-0501-B005}{start}
\olsection{معنايي مفاهيم}
\olpseganchor{OLP-0501-B005}{end}

\olpseganchor{OLP-0501-B006}{start}
\begin{defn}
  وايو \LR{$!A$} \emph{په مدل \LR{$\mModel{M} = \tuple{W,R,V}$} کښې رښتيا دے،} \LR{$
  \mSat{M}{!A}$}، هله او يوازې هله چې د هر \LR{$w \in W$} لپاره
  \LR{$\mSat{M}{!A}[w]$} وي۔ \LR{$!A$} هله \emph{معتبر} دے، \LR{$\Entails !A$}،
  چې په ټولو مدلونو کښې رښتيا وي۔ وايو د فارمولونه سټ~\LR{$\Gamma$}
  \LR{$!A$} په \emph{معنايي ډول لازموي}، \LR{$\Gamma
  \Entails !A$}، هله او يوازې هله چې د هر مدل~\LR{$\mModel{M}$} او هر
  داسې~\LR{$w$} لپاره چې \LR{$\mSat{M}{\Gamma}[w]$} وي،
  \LR{$\mSat{M}{!A}[w]$} وي۔
\end{defn}
\olpseganchor{OLP-0501-B006}{end}

\olpseganchor{OLP-0501-B007}{start}
\begin{prop}\ollabel{prop:sat-entails}
  \begin{enumerate}
  \item\ollabel{prop:sat-entails1} که \LR{$\mSat{M}{\Gamma}[w]$} او
    \LR{$\Gamma \Entails !A$} وي، نو \LR{$\mSat{M}{!A}[w]$} وي۔
  \item\ollabel{prop:sat-entails2} که \LR{$\mSat{M}{\Gamma}$} او \LR{$\Gamma
    \Entails !A$} وي، نو \LR{$\mSat{M}{!A}$} وي۔
  \end{enumerate}
\end{prop}
\olpseganchor{OLP-0501-B007}{end}

\olpseganchor{OLP-0501-B008}{start}
\begin{proof}
  \begin{enumerate}
  \item فرض کړئ \LR{$\mSat{M}{\Gamma}[w]$}۔ ځکه \LR{$\Gamma \Entails !A$}
    دے، د استلزام له تعريفه \LR{$\mSat{M}{!A}[w]$} راځي۔
  \item فرض کړئ \LR{$\mSat{M}{\Gamma}$}۔ نو د هر \LR{$u \in W$} لپاره
    \LR{$\mSat{M}{\Gamma}[u]$} وي۔ په تېره بيا د هرې نړۍ لپاره
    \LR{$\mSat{M}{\Gamma}[w]$} وي؛ د \olref{prop:sat-entails1} له مخې
    \LR{$\mSat{M}{!A}[w]$} هم وي۔ نو قضيه په ټول مدل کښې رښتيا ده۔
  \end{enumerate}
\end{proof}
\olpseganchor{OLP-0501-B008}{end}


\olpseganchor{OLP-0501-B009}{start}
\begin{defn}\ollabel{defn:restrict}
  فرض کړئ \LR{$\mModel{M}$} اړيکيز مدل او \LR{$w \in W$} وي۔ تر~\LR{$w$}
  پورې د~\LR{$\mModel{M}$} \emph{محدودونه}
  \LR{$\mModel{M}_w=\tuple{W_w, R_w, V_w}$} داسې ټاکل کېږي:
  \begin{align*}
    W_w & = \Setabs{u \in W}{Rwu},\\
    R_w  & = R \cap (W_w)^2, \text{\RL{ او}}\\
    V_w(p) & = V(p) \cap W_w.
  \end{align*}
  انعکاسيت د پيل نړۍ په محدود مدل کښې ساتي۔ تعدي دا تضمينوي چې
  له هرې پاتې نړۍ څخه هره د لاسرسي وړ نړۍ هم په محدود مدل کښې
  پاتې کېږي۔
\end{defn}
\olpseganchor{OLP-0501-B009}{end}

\olpseganchor{OLP-0501-B010}{start}
\begin{prop}\ollabel{prop:restrict}
  \LR{$\mSat{M}{!A}[w]$} هله او يوازې هله چې \LR{$\mSat{M_w}{!A}$} وي۔
\end{prop}
\olpseganchor{OLP-0501-B010}{end}

\olpseganchor{OLP-0501-B011}{start}
\begin{prob}
  \olref[int][sem][sem]{prop:restrict} ثابت کړئ۔ د فارمول پر
  جوړښت په استقرا وښيئ چې په هره پاتې نړۍ کښې له محدودونې مخکې
  او وروسته رښتياوالے يو شان دے۔ بيا د رښتياوالي يورخيزه
  زياتېدنه او په محدود مدل کښې د پيل نړۍ غړيتوب وکاروئ۔
\end{prob}
\olpseganchor{OLP-0501-B011}{end}

\olpseganchor{OLP-0501-B012}{start}
\begin{prop}
  فرض کړئ چې د هر داسې مدل~\LR{$\mModel{M}$} لپاره چې
  \LR{$\mSat{M}{\Gamma}$} وي، \LR{$\mSat{M}{!A}$} هم وي۔ نو
  \LR{$\Gamma \Entails !A$} دے۔
\end{prop}
\olpseganchor{OLP-0501-B012}{end}

\olpseganchor{OLP-0501-B013}{start}
\begin{proof}
  فرض کړئ \LR{$\mSat{M}{\Gamma}[w]$} وي۔ د
  \olref{prop:restrict} په هر~\LR{$!B \in \Gamma$} تطبيقولو سره
  \LR{$\mSat{M_w}{\Gamma}$} لرو۔ د فرض له مخې
  \LR{$\mSat{M_w}{!A}$} لرو۔ بيا د \olref{prop:restrict} له مخې
  \LR{$\mSat{M}{!A}[w]$} راځي۔
\end{proof}
\olpseganchor{OLP-0501-B013}{end}

\olpseganchor{OLP-0502-B004}{start}
\olfileid{int}{sem}{top}
\olpseganchor{OLP-0502-B004}{end}

\olpseganchor{OLP-0502-B005}{start}
\olsection{توپولوژيکي معنايي پوهنه}
\olpseganchor{OLP-0502-B005}{end}

\olpseganchor{OLP-0502-B006}{start}
د شهودي منطق د معنايي پوهنې بله لار د توپولوژۍ د رياضيکي مفهوم
کارول دي۔
\olpseganchor{OLP-0502-B006}{end}

\olpseganchor{OLP-0502-B007}{start}
\begin{defn}
  فرض کړئ \LR{$X$} يو سټ دے۔ پر~\LR{$X$} \emph{توپولوژي} داسې سټ
  \LR{$\Top{O} \subseteq \Pow{X}$} دے چې لاندې خاصيتونه پوره کوي۔
  د~\LR{$\Top{O}$} غړي د توپولوژۍ \emph{پرانيستي سټونه}
  بلل کېږي۔ سټ \LR{$X$} له \LR{$\Top{O}$} سره په ګډه \emph{توپولوژيکي فضا}
  بلل کېږي۔
  \begin{enumerate}
  \item تش سټ او ټوله فضا پرانيستي دي: \LR{$\emptyset$}، \LR{$X \in
    \Top{O}$}۔
  \item پرانيستي سټونه د متناهي تقاطع له مخې تړلي دي: که \LR{$U$}، \LR{$V \in
    \Top{O}$} وي، نو \LR{$U \cap V \in \Top{O}$}۔
  \item پرانيستي سټونه د هرې اجتماع له مخې تړلي دي: که د هر \LR{$i \in I$}
    لپاره \LR{$U_i \in \Top{O}$} وي، نو \LR{$\bigcup \Setabs{U_i}{i \in
      I} \in \Top{O}$}۔
  \end{enumerate}
\end{defn}
\olpseganchor{OLP-0502-B007}{end}

\olpseganchor{OLP-0502-B008}{start}
که د پرانيستو سټونو ټولګه له سياق څخه معلومېدے شي، توپولوژيکي
فضا په لنډ ډول په \LR{$X$} هم نومولے شو۔ خو پام وکړئ چې \LR{$X$} يوازې
هغه وخت توپولوژيکي فضا بلل کېږي چې د پرانيستو سټونو ټولګه ورته
ټاکل شوې وي۔ پر يوه سټ بېلابېلې توپولوژۍ ټاکل کېدې شي۔
\olpseganchor{OLP-0502-B008}{end}

\olpseganchor{OLP-0502-B009}{start}
\begin{defn}
  د شهودي قضييز منطق \emph{توپولوژيکي مدل} يوه درېګونې
  \LR{$\mModel{X} = \tuple{X, \Top{O}, V}$} ده، چې \LR{$\Top{O}$}
  پر~\LR{$X$} توپولوژي ده او \LR{$V$} داسې تابع ده چې هر بياني متغير ته
  په \LR{$\Top{O}$} کښې يو پرانيستے سټ ټاکي۔
\olpseganchor{OLP-0502-B009}{end}

\olpseganchor{OLP-0502-B010}{start}
  د توپولوژيکي مدل~\LR{$\mModel{X}$} په لرلو سره \LR{$\Prop{X}{!A}$} په
  استقرايي ډول داسې تعريفولے شو:
  \begin{enumerate}
  \item \LR{$\Prop{X}{\lfalse} = \emptyset$}
  \item \LR{$\Prop{X}{p} = V(p)$}
  \item \LR{$\Prop{X}{!A \land !B} = \Prop{X}{!A} \cap \Prop{X}{!B}$}
  \item \LR{$\Prop{X}{!A \lor !B} = \Prop{X}{!A} \cup \Prop{X}{!B}$}
  \item \LR{$\Prop{X}{!A \lif !B} = \Interior{(X \setminus \Prop{X}{!A}) \cup
    \Prop{X}{!B}}$}
  \end{enumerate}
  دلته \LR{$\Interior{V}$} هغه تابع ده چې سټ \LR{$V \subseteq X$} ئې خپل
  \emph{دنننۍ برخه} ته وړي، يعنې د هغو ټولو پرانيستو سټونو اجتماع
  ته چې په دغه سټ کښې شامل دي۔ په بله وينا،
  \[
  \Interior{V} = \bigcup \Setabs{U}{U \subseteq V \text{\RL{ او }} U \in \Top{O}}.
  \]
\end{defn}
\olpseganchor{OLP-0502-B010}{end}

\olpseganchor{OLP-0502-B011}{start}
د هر سټ دنننۍ برخه تل پرانيستې وي، ځکه د پرانيستو سټونو اجتماع ده۔
نو \LR{$\Prop{X}{!A}$} تل پرانيستے سټ دے۔
\olpseganchor{OLP-0502-B011}{end}

\olpseganchor{OLP-0502-B012}{start}
که څۀ هم توپولوژيکي معنايي پوهنه ډېره انتزاعي ده، د هغې د انګېزې
لپاره تصوراتي لارې شته۔ فرض کړئ د \LR{$X$} غړي، يا «ټکي»،
هغه ځايونه دي چې ويناوې پکې ارزول کېدے شي۔ د هغو ټولو ټکو سټ
چې \LR{$!A$} پکې رښتيا وي، د~\LR{$!A$} څرګنده کړې قضيه ده۔ د ټکو هر سټ
ممکنه قضيه نۀ ده؛ يوازې د \LR{$\Top{O}$} غړي داسې دي۔
\LR{$!A \Entails !B$} هله او يوازې هله چې \LR{$!B$} په هر هغه ټکي کښې
رښتيا وي چې~\LR{$!A$} پکې رښتيا وي؛ يعنې د~\LR{$X$} پر هر توپولوژيکي مدل، له هرې پرانيستې ارزښت ټاکنې سره،
\LR{$\Prop{X}{!A} \subseteq \Prop{X}{!B}$} وي۔ محاله وينا~\LR{$\lfalse$}
هېڅکله رښتيا نۀ ده، نو \LR{$\Prop{X}{\lfalse} = \emptyset$}۔
\olpseganchor{OLP-0502-B012}{end}

\olpseganchor{OLP-0502-B013}{start}
د \LR{$!B \land !C$}، \LR{$!B \lor !C$} او \LR{$!B \lif !C$} څرګندې کړې
قضيې بايد د \LR{$!B$} او~\LR{$!C$} له قضيو سره څنګه تړلې وي، څو شهودي
معتبر اصول پوره شي؟ يعنې څو \LR{$!A \Proves !B$} هله او يوازې هله
وي چې \LR{$\Prop{X}{!A} \subseteq \Prop{X}{!B}$} په هر توپولوژيکي
مدل کښې وي؟ برابري هم روا ده، ځکه هر فارمول خپل ځان لازموي۔ د هر \LR{$!A$}
لپاره \LR{$\lfalse \Proves !A$} غواړو، او دا پوره کېږي ځکه د هر
\LR{$U$} لپاره \LR{$\emptyset \subseteq U$} دے۔ څرنګه چې
\LR{$!B \land !C \Proves !B$}، نو
\LR{$\Prop{X}{!B \land !C} \subseteq \Prop{X}{!B}$} غواړو؛ په
همدې ډول \LR{$\Prop{X}{!B \land !C} \subseteq \Prop{X}{!C}$} هم۔
تر ټولو لوے سټ چې \LR{$W \subseteq U$} او \LR{$W \subseteq V$} پوره کوي،
\LR{$U \cap V$} دے۔ بل خوا \LR{$!B \Proves !B \lor !C$} او
\LR{$!C \Proves !B \lor !C$} دي، نو
\LR{$\Prop{X}{!B} \subseteq \Prop{X}{!B \lor !C}$} او
\LR{$\Prop{X}{!C} \subseteq \Prop{X}{!B \lor !C}$} غواړو۔ تر ټولو
کوچنے سټ~\LR{$W$} چې \LR{$U \subseteq W$} او \LR{$V \subseteq W$} پوره کوي،
\LR{$U \cup V$} دے۔
\olpseganchor{OLP-0502-B013}{end}

\olpseganchor{OLP-0502-B014}{start}
د \LR{$\lif$} تعريف لږ پېچلے دے: \LR{$!A \lif !B$} تر ټولو کمزورې هغه
قضيه څرګندوي چې له \LR{$!A$} سره يوځاے \LR{$!B$} ته استلزام ورکوي۔
له \LR{$(!A \lif !B) \land !A \Proves !B $} څخه څرګنده ده چې \LR{$!A
\lif !B$} له \LR{$!A$} سره يوځاے~\LR{$!B$} ته استلزام ورکوي۔ نو
\LR{$\Prop{X}{!A \lif !B}$} بايد تر ټولو لوے پرانيستے سټ وي چې
\LR{$\Prop{X}{!A \lif !B} \cap \Prop{X}{!A} \subseteq \Prop{X}{!B}$} پوره کړي۔ دا شرط هله پوره کېږي چې
د نوموړي پرانيستي سټ هر ټکے يا د مقدمې له سټه بهر وي او يا د
پايلې په سټ کښې وي۔ د دغو ټکو د سټ دنننۍ برخه تر ټولو لوے داسې
پرانيستے سټ دے؛ له همدې څخه زموږ تعريف راځي۔
\olpseganchor{OLP-0502-B014}{end}

\olpseganchor{OLP-0503-B004}{start}
\olchapter{int}{sc}{سموالے او بشپړتيا}
\olpseganchor{OLP-0503-B004}{end}

\olpseganchor{OLP-0503-B005}{start}
\begin{editorial}
  په دې باب کښې سموالے ښيي چې استنتاج معنايي استلزام ساتي۔
  د بشپړتيا ثبوت ښيي چې که له يوې ډلې فرضونو استنتاج ممکن نۀ وي،
  يو اړيکيز مخالف مدل جوړولے شو۔
\end{editorial}
\olpseganchor{OLP-0503-B005}{end}

\olpseganchor{OLP-0504-B004}{start}
\olfileid{int}{sc}{sax}
\olpseganchor{OLP-0504-B004}{end}

\olpseganchor{OLP-0504-B005}{start}
\olsection{د بديهي اشتقاقونه سموالے}
\olpseganchor{OLP-0504-B005}{end}

\olpseganchor{OLP-0504-B006}{start}
\begin{editorial}
  په بديهي حساب کښې هره کرښه يا مقدمه ده، يا د بديهي بڼې بېلګه،
  يا له دوو پخوانيو کرښو د مودوس پوننس له مخې راځي۔
\end{editorial}
\olpseganchor{OLP-0504-B006}{end}

\olpseganchor{OLP-0504-B007}{start}
\begin{thm}[سموالے]
  \ollabel{thm:soundness}
  که \LR{$\Gamma \Proves !A$} وي، نو \LR{$ \Gamma \Entails !A$} وي۔
\end{thm}
\olpseganchor{OLP-0504-B007}{end}

\olpseganchor{OLP-0504-B008}{start}
\begin{proof}
  ثابتوو چې که \LR{$\Gamma \Proves !A$} وي، نو \LR{$\Gamma \Entails !A$}
  وي۔ ثبوت له~\LR{$\Gamma$} څخه د~\LR{$!A$} په اشتقاق کښې د
  فارمولونه پر شمېر~\LR{$n$} استقرا ده۔ ښيو چې که \LR{$!A_1$}،
  \dots، \LR{$!A_n = !A$} له~\LR{$\Gamma$} څخه اشتقاق وي، نو
  \LR{$\Gamma \Entails !A_n$} وي۔ پام وکړئ چې که \LR{$!A_1$}، \dots،
  \LR{$!A_n$} اشتقاق وي، نو د هر \LR{$k<n$} لپاره \LR{$!A_1$}،
  \dots، \LR{$!A_k$} هم استنتاج دے۔
\olpseganchor{OLP-0504-B008}{end}

\olpseganchor{OLP-0504-B009}{start}
  د اوږدوالي~\LR{$0$} هېڅ اشتقاقونه نشته، نو د \LR{$n=0$} لپاره ادعا
  تشه رښتيا ده۔ اوس فرض کوو چې ادعا د اوږدوالي~\LR{$<n$} د ټولو
  اشتقاقونه لپاره سمه ده۔ د~\LR{$!A_n$} د توجيه له مخې حالتونه
  بېلوو۔
\olpseganchor{OLP-0504-B009}{end}
  
\olpseganchor{OLP-0504-B010}{start}
  \begin{enumerate}
  \item \LR{$!A_n$} بديهي اصل دے۔ ټول بديهي اصول معتبر دي، نو
    د هر~\LR{$\Gamma$} لپاره \LR{$\Gamma \Entails !A_n$} وي۔
    د عطف لومړۍ دوې بڼې د جوړې غړي راباسي۔ د عطف درېيمه بڼه
    په دوو پرله‌پسې د لاسرسي وړ نړۍو کښې دواړه فرضونه سره يوځاے
    کوي، ځکه لومړے فرض تر دويمې نړۍ پورې ساتل کېږي۔ د فصل لومړۍ
    دوې بڼې د اړوند غړي له مخې فصل رښتينے کوي۔ د فصل درېيمه
    بڼه په وروستۍ نړۍ کښې پر دواړو شرطونو عمل کوي، او د فصل
    له رښتيني غړي څخه پايله اخلي۔ د شرط په لومړۍ بڼه کښې
    لومړے فرض په ټولو ورپسې نړۍو کښې پاتې کېږي۔ په دويمه
    بڼه کښې په هره وروستۍ نړۍ کښې د لومړي فرض په منلو سره
    د دويم فرض او ورپسې پايلې تر منځ شرط لرو۔ د بلې مقدمې
    له مخې دويم فرض هم لرو، او انعکاسيت دواړه شرطونه په
    هماغې نړۍ کښې تطبيقوي۔ د محال لپاره هېڅ نړۍ مقدمه نۀ
    رښتيا کوي، نو د شرطي بڼې عمومي حکم پوره کېږي۔
  \item \LR{$!A_n \in \Gamma$} وي۔ نو د هر \LR{$\mModel{M}$} او \LR{$w$}
    لپاره، که \LR{$\mSat{M}{\Gamma}[w]$} وي، څرګنده ده چې
    \LR{$\mSat{M}{!A_n}[w]$} وي؛ يعنې \LR{$\Gamma \Entails !A$} وي۔
  \item \LR{$!A_n$} له \LR{$!A_i$} او \LR{$!A_j \ident !A_i \lif !A_n$} څخه د
    \MP{} له مخې راځي۔ \LR{$!A_1$}، \dots، \LR{$!A_i$} او \LR{$!A_1$}،
    \dots، \LR{$!A_j$} له~\LR{$\Gamma$} څخه اشتقاقونه دي، نو د
    استقرايي فرض له مخې \LR{$\Gamma \Entails !A_i$} او
    \LR{$\Gamma \Entails !A_i \lif !A_n$} وي۔
\olpseganchor{OLP-0504-B010}{end}
    
\olpseganchor{OLP-0504-B011}{start}
    فرض کړئ \LR{$\mSat{M}{\Gamma}[w]$} وي۔ ځکه
    \LR{$\mSat{M}{\Gamma}[w]$} او \LR{$\Gamma \Entails !A_i \lif !A_n$}
    دي، نو \LR{$\mSat{M}{!A_i \lif !A_n}[w]$} وي۔ د تعريف له مخې
    د هر داسې \LR{$w'$} لپاره چې \LR{$Rww'$} وي، که
    \LR{$\mSat{M}{!A_i}[w']$} وي، نو \LR{$\mSat{M}{!A_n}[w']$} وي۔
    څرنګه چې \LR{$R$} انعکاسي ده، \LR{$w$} هم په هغو \LR{$w'$} کښې دے چې
    \LR{$Rww'$} لري؛ يعنې که \LR{$\mSat{M}{!A_i}[w]$} وي، نو
    \LR{$\mSat{M}{!A_n}[w]$} وي۔ ځکه چې \LR{$\Gamma \Entails !A_i$}،
    \LR{$\mSat{M}{!A_i}[w]$} وي۔ نو \LR{$\mSat{M}{!A_n}[w]$} شو،
    لکه چې غوښتل مو۔
  \end{enumerate}
\end{proof}
\olpseganchor{OLP-0504-B011}{end}

\olpseganchor{OLP-0505-B004}{start}
\olfileid{int}{sc}{snd}
\olpseganchor{OLP-0505-B004}{end}

\olpseganchor{OLP-0505-B005}{start}
\olsection{د طبيعي استنتاج سموالے}
\olpseganchor{OLP-0505-B005}{end}

\olpseganchor{OLP-0505-B006}{start}
اوس به د اړيکيزې معنايي پوهنې له مخې د طبيعي استنتاج سموالے ثابت
کړو: يعنې وښيو چې که فارمول د فرضونو له يوه سټ څخه
د اشتقاق وړ وي، نو د فرضونو سټ هغه فارمول لازموي۔
\olpseganchor{OLP-0505-B006}{end}

\olpseganchor{OLP-0505-B007}{start}
\begin{thm}[سموالے]
  \ollabel{thm:soundness}
  که \LR{$\Gamma \Proves !A$} وي، نو \LR{$ \Gamma \Entails !A$} وي۔
\end{thm}
\olpseganchor{OLP-0505-B007}{end}

\olpseganchor{OLP-0505-B008}{start}
\begin{proof}
  ثابتوو چې که \LR{$\Gamma \Proves !A$} وي، نو \LR{$\Gamma \Entails !A$}
  وي۔ ثبوت له~\LR{$\Gamma$} څخه د~\LR{$!A$} پر اشتقاق استقرا ده۔
\olpseganchor{OLP-0505-B008}{end}
  
\olpseganchor{OLP-0505-B009}{start}
  \begin{enumerate}
  \item که اشتقاق يوازې له فرض~\LR{$!A$} څخه جوړ وي،
    \LR{$!A \Proves !A$} لرو او بايد \LR{$!A \Entails !A$} وښيو۔ فرض کړئ
    \LR{$\mSat{M}{!A}[w]$} وي۔ نو په څرګنده \LR{$\mSat{M}{!A}[w]$} وي۔
\olpseganchor{OLP-0505-B009}{end}

\olpseganchor{OLP-0505-B010}{start}
  \item اشتقاق په \LR{$\Intro{\land}$} پاے ته رسېږي:
    له نۀ ايستل شوې فرضونو~\LR{$\Gamma$}
      څخه د مقدمې~\LR{$!B$} او له نۀ ايستل شوې فرضونو~\LR{$\Delta$}
      څخه د مقدمې~\LR{$!C$} اشتقاقونه ښيي چې \LR{$\Gamma \Proves
      !B$} او \LR{$\Delta \Proves !C$}۔ د استقرا د فرض له مخې
      \LR{$\Gamma \Entails !B$} او \LR{$\Delta \Entails !C$} لرو۔ بايد
      \LR{$\Gamma \cup \Delta \Entails !B \land !C$} وښيو، ځکه د ټول
      اشتقاق نۀ ايستل شوي فرضونه د \LR{$\Gamma$} او
      \LR{$\Delta$} ګډ سټ دے۔ نو \LR{$\mSat{M}{\Gamma
        \cup \Delta}[w]$} فرض کړئ۔ بيا \LR{$\mSat{M}{\Gamma}[w]$}
      هم وي۔ ځکه چې \LR{$\Gamma \Entails !B$}، نو \LR{$\mSat{M}{!B}[w]$}
      وي۔ په ورته ډول \LR{$\mSat{M}{!C}[w]$} وي۔ نو
      \LR{$\mSat{M}{!B \land !C}[w]$} وي۔
\olpseganchor{OLP-0505-B010}{end}

\olpseganchor{OLP-0505-B011}{start}
  \item اشتقاق په \LR{$\Elim{\land}$} پاے ته رسېږي:
    له نۀ ايستل شوې فرضونو \LR{$\Gamma$}
        څخه د مقدمې \LR{$!B \land !C$} اشتقاق ښيي چې
        \LR{$\Gamma \Proves !B \land !C$}۔ د استقرا د فرض له مخې
        \LR{$\Gamma \Entails !B \land !C$}۔ بايد \LR{$\Gamma
        \Entails !B$} وښيو۔ نو \LR{$\mSat{M}{\Gamma}[w]$} فرض کړئ۔
        ځکه \LR{$\Gamma \Entails !B \land !C$}، نو
        \LR{$\mSat{M}{!B \land !C}[w]$} وي۔ بيا
        \LR{$\mSat{M}{!B}[w]$} هم وي۔ په ورته ډول که \LR{$\Elim\land$}
        په~\LR{$!C$} پاے ته ورسېږي، نو \LR{$\Gamma \Entails !C$} وي۔
\olpseganchor{OLP-0505-B011}{end}

\olpseganchor{OLP-0505-B012}{start}
  \item اشتقاق په \LR{$\Intro{\lor}$} پاے ته رسېږي:
    فرض کړئ مقدمه \LR{$!B$} ده او په~\LR{$!B$}
      ختمېدونکي اشتقاق نۀ ايستل شوي فرضونه \LR{$\Gamma$}
      دي۔ نو \LR{$\Gamma \Proves !B$} لرو او د استقرا د فرض له مخې
      \LR{$\Gamma \Entails !B$}۔ بايد \LR{$\Gamma \Entails !B \lor !C$}
      وښيو۔ \LR{$\mSat{M}{\Gamma}[w]$} فرض کړئ۔ ځکه
      \LR{$\Gamma \Entails !B$}، نو \LR{$\mSat{M}{!B}[w]$} وي۔ بيا
      \LR{$\mSat{M}{!B \lor !C}[w]$} هم وي۔ په ورته ډول، که مقدمه
      ~\LR{$!C$} وي، نو \LR{$\Gamma \Entails !C$} لرو۔
\olpseganchor{OLP-0505-B012}{end}
    
\olpseganchor{OLP-0505-B013}{start}
  \item اشتقاق په~\LR{$\Elim{\lor}$} پاے ته رسېږي:
    په مقدمو ختمېدونکي اشتقاقونه دا دي:
      \LR{$!B \lor !C$} له نۀ ايستل شوې فرضونو~\LR{$\Gamma$} څخه،
      \LR{$!D$} له نۀ ايستل شوې فرضونو
      \LR{$\Delta_1 \cup \{!B\}$} څخه، او \LR{$!D$} له نۀ ايستل شوې
      فرضونو~\LR{$\Delta_2 \cup \{!C\}$} څخه۔ نو \LR{$\Gamma \Proves
      !B \lor !C$}، \LR{$\Delta_1 \cup \{!B\} \Proves !D$} او
      \LR{$\Delta_2 \cup \{!C\} \Proves !D$} لرو۔ د استقرا د فرض له
      مخې \LR{$\Gamma \Entails !B \lor !C$}،
      \LR{$\Delta_1 \cup \{!B\} \Entails !D$} او
      \LR{$\Delta_2 \cup \{!C\} \Entails !D$} دي۔ بايد
      \LR{$\Gamma \cup \Delta_1 \cup \Delta_2 \Entails !D$} ثابت کړو۔
\olpseganchor{OLP-0505-B013}{end}

\olpseganchor{OLP-0505-B014}{start}
    \LR{$\mSat{M}{\Gamma \cup \Delta_1 \cup \Delta_2}[w]$} فرض کړئ۔
    نو \LR{$\mSat{M}{\Gamma}[w]$} وي، او ځکه چې
    \LR{$\Gamma \Entails !B \lor !C$}، نو
    \LR{$\mSat{M}{!B \lor !C}[w]$} وي۔ د \LR{$\mSat{M}{}$} د تعريف له
    مخې يا \LR{$\mSat{M}{!B}[w]$} وي يا \LR{$\mSat{M}{!C}[w]$}۔ نو دوه
    حالتونه بېلوو: (الف)~\LR{$\mSat{M}{!B}[w]$} وي۔ بيا
    \LR{$\mSat{M}{\Delta_1 \cup \{!B\}}[w]$} وي۔ ځکه
    \LR{$\Delta_1 \cup \{!B\} \Entails !D$}، نو
    \LR{$\mSat{M}{!D}[w]$} وي۔ (ب)~\LR{$\mSat{M}{!C}[w]$} وي۔ بيا
    \LR{$\mSat{M}{\Delta_2 \cup \{!C\}}[w]$} وي۔ ځکه
    \LR{$\Delta_2 \cup \{!C\} \Entails !D$}، نو
    \LR{$\mSat{M}{!D}[w]$} وي۔ په دواړو حالاتو کښې
    \LR{$\mSat{M}{!D}[w]$} لرو، لکه چې غوښتل مو۔
    
\olpseganchor{OLP-0505-B014}{end}

\olpseganchor{OLP-0505-B015}{start}
  \item اشتقاق په \LR{$\Intro{\lif}$} پاے ته رسېږي او
    پايله ئې~\LR{$!B\lif !C$} ده۔ نو مقدمه~\LR{$!C$} ده، او په دې مقدمه
    ختمېدونکے اشتقاق نۀ ايستل شوي فرضونه~\LR{$\Gamma \cup
    \{!B\}$} لري۔ په دې توګه \LR{$\Gamma \cup \{!B\} \Proves !C$}،
    او د استقرا د فرض له مخې \LR{$\Gamma \cup \{!B\} \Entails !C$}
    لرو۔ بايد \LR{$\Gamma \Entails !B \lif !C$} وښيو۔
\olpseganchor{OLP-0505-B015}{end}

\olpseganchor{OLP-0505-B016}{start}
    \LR{$\mSat{M}{\Gamma}[w]$} فرض کړئ۔ بايد وښيو چې د هر داسې
    \LR{$w'$} لپاره چې \LR{$Rww'$} وي، که \LR{$\mSat{M}{!B}[w']$} وي، نو
    \LR{$\mSat{M}{!C}[w']$} وي۔ نو \LR{$Rww'$} او
    \LR{$\mSat{M}{!B}[w']$} فرض کړئ۔ د
    \olref[sem][rel]{prop:true-monotonic} له مخې
    \LR{$\mSat{M}{\Gamma}[w']$} وي۔ ځکه
    \LR{$\Gamma \cup \{!B\} \Entails !C$}، نو
    \LR{$\mSat{M}{!C}[w']$} وي؛ همدا مو غوښتل۔
\olpseganchor{OLP-0505-B016}{end}

\olpseganchor{OLP-0505-B017}{start}
  \item اشتقاق په \LR{$\Elim{\lif}$} پاے ته رسېږي او
    پايله ئې~\LR{$!C$} ده۔ مقدمې \LR{$!B \lif !C$} او \LR{$!B$} دي چې په
    ترتيب له نۀ ايستل شوې فرضونو \LR{$\Gamma$} او \LR{$\Delta$}
    څخه اشتقاقونه لري۔ نو \LR{$\Gamma \Proves !B \lif !C$} او
    \LR{$\Delta \Proves !B$} لرو۔ د استقرا د فرض له مخې
    \LR{$\Gamma \Entails !B \lif !C$} او \LR{$\Delta \Entails !B$}
    دي۔ بايد \LR{$\Gamma \cup \Delta \Entails !C$} وښيو۔
\olpseganchor{OLP-0505-B017}{end}

\olpseganchor{OLP-0505-B018}{start}
    \LR{$\mSat{M}{\Gamma \cup \Delta}[w]$} فرض کړئ۔ ځکه
    \LR{$\mSat{M}{\Gamma}[w]$} او \LR{$\Gamma \Entails !B \lif !C$}
    دي، نو \LR{$\mSat{M}{!B \lif !C}[w]$} وي۔ د تعريف له مخې د هر
    داسې \LR{$w'$} لپاره چې \LR{$Rww'$} وي، که \LR{$\mSat{M}{!B}[w']$} وي،
    نو \LR{$\mSat{M}{!C}[w']$} وي۔ څرنګه چې \LR{$R$} انعکاسي ده، \LR{$w$}
    هم په هغو \LR{$w'$} کښې دے چې \LR{$Rww'$} لري؛ يعنې که
    \LR{$\mSat{M}{!B}[w]$} وي، نو \LR{$\mSat{M}{!C}[w]$} وي۔ ځکه
    \LR{$\mSat{M}{\Delta}[w]$} او \LR{$\Delta \Entails !B$} دي، نو
    \LR{$\mSat{M}{!B}[w]$} وي۔ نو \LR{$\mSat{M}{!C}[w]$} شو؛ همدا مو
    غوښتل۔
\olpseganchor{OLP-0505-B018}{end}

\olpseganchor{OLP-0505-B019}{start}
  \item اشتقاق په \LR{$\FalseInt$} پاے ته رسېږي او پايله ئې
    ~\LR{$!A$} ده۔ مقدمه \LR{$\lfalse$} ده او د مقدمې د اشتقاق
    نۀ ايستل شوي فرضونه~\LR{$\Gamma$} دي۔ نو \LR{$\Gamma \Proves
    \lfalse$} وي۔ د استقرا د فرض له مخې
    \LR{$\Gamma \Entails \lfalse$} وي۔ بايد \LR{$\Gamma \Entails !A$}
    وښيو۔
\olpseganchor{OLP-0505-B019}{end}

\olpseganchor{OLP-0505-B020}{start}
    په نامستقيم ډول ځو۔ که \LR{$\Gamma \Entails/ !A$} وي، داسې
    مدل~\LR{$\mModel{M}$} او نړۍ~\LR{$w$} شته چې
    \LR{$\mSat{M}{\Gamma}[w]$} او \LR{$\mSat/{M}{!A}[w]$} وي۔ ځکه
    \LR{$\Gamma \Entails \lfalse$}، نو \LR{$\mSat{M}{\lfalse}[w]$}
    وي۔ خو دا ناشونې ده، ځکه د تعريف له مخې
    \LR{$\mSat/{M}{\lfalse}[w]$} وي۔ نو \LR{$\Gamma \Entails !A$} دے۔
  \item اشتقاق په \LR{$\Intro\lnot$} پاے ته رسېږي۔ له فرضونو
    څخه يوه نړۍ واخلئ چې ټول پاتې فرضونه پکښې رښتيا وي۔
    دغه فرضونه په هرې د لاسرسي وړ نړۍ کښې پاتې کېږي۔ که په داسې
    نړۍ کښې تړل شوے فرض هم رښتينے وي، د استقرا د فرض له مخې
    محال رښتينے کېږي۔ نو داسې نړۍ نشته، او نفي رښتينې ده۔
  \item اشتقاق په \LR{$\Elim\lnot$} پاے ته رسېږي۔ که په يوه
    نړۍ کښې د دواړو مقدمو فرضونه رښتيا وي، د استقرا له مخې
    فارمول او نفي ئې دواړه رښتيني وي۔ انعکاسيت او د نفي
    تعريف د هماغه فارمول صدق په اوسنۍ نړۍ کښې هم ردوي۔ نو
    داسې نړۍ نشته او محال د دغو فرضونو معنايي پايله ده۔
  \end{enumerate}
\end{proof}
\olpseganchor{OLP-0505-B020}{end}

\olpseganchor{OLP-0505-B021}{start}
\begin{prob}
  د \olref[int][sc][snd]{thm:soundness} ثبوت بشپړ کړئ۔ د
  \LR{$\Intro\lnot$} او \LR{$\Elim\lnot$} د حالتونو لپاره د
  \olref[int][sem][rel]{defn:true-at-w} له مخې د
  \LR{$\mSat{M}{\lnot !A}[w]$} تعريف وکاروئ؛ يعنې \LR{$\lnot !A$} د
  \LR{$!A \lif \lfalse$} په توګه تعريف شوے مۀ ګڼئ۔ د دواړو حالتونو
  دليل په خپلو ټکو، د پرانيستو او تړل شويو فرضونو په څرګندولو
  سره وليکئ۔
\end{prob}
\olpseganchor{OLP-0505-B021}{end}

\olpseganchor{OLP-0505-B022}{start}
\begin{prob}
  وښيئ چې لاندې فارمولونه په شهودي منطق کښې د اشتقاق وړ نۀ دي:
  \begin{enumerate}
    \item \LR{$(!A \lif !B) \lor (!B \lif !A)$}
    \item \LR{$(\lnot\lnot !A \lif !A) \lif (!A \lor \lnot !A)$}
    \item \LR{$(!A \lif !B \lor !C) \lif \bigl((!A \lif !B)\lor(!A \lif !C)\bigr)$}
  \end{enumerate}
\end{prob}
\olpseganchor{OLP-0505-B022}{end}

\olpseganchor{OLP-0506-B004}{start}
\olfileid{int}{sc}{lin}
\olpseganchor{OLP-0506-B004}{end}

\olpseganchor{OLP-0506-B005}{start}
\olsection{د لينډنباوم لمه}
\olpseganchor{OLP-0506-B005}{end}

\olpseganchor{OLP-0506-B006}{start}
د شهودي منطق د بشپړتيا قضيه د دې په فرضولو ثابتېږي چې
\LR{$\Gamma \Proves/ !A$}؛ بيا داسې مدل جوړوو چې
\LR{$\mSat{M}{\Gamma}$} او~\LR{$\mSat/{M}{!A}$} ولري۔
\olpseganchor{OLP-0506-B006}{end}

\olpseganchor{OLP-0506-B007}{start}
په کلاسيک منطق کښې د اشتقاق وړتيا اړيکه سازګارۍ ته راکمولے
شو، ځکه فارمول~\LR{$!A$} د فارمولونه له يوه سټ څخه هله او
يوازې هله د اشتقاق وړ وي چې دغه سټ د~\LR{$!A$} له نفي سره يوځاے
ناسازګار وي۔ په شهودي منطق کښې دا کار نۀ شو کولے۔ په شهودي منطق
کښې که \LR{$\lnot!A$} ناسازګار وي، يوازې \LR{$\Proves \lnot\lnot !A$}
ترلاسه کوو۔ څرنګه چې \LR{$\lnot\lnot!A \lif !A$} په عمومي ډول په
شهودي منطق کښې نۀ منل کېږي، نو \LR{$\Proves !A$} نۀ شو اخيستلے۔
\olpseganchor{OLP-0506-B007}{end}

\olpseganchor{OLP-0506-B008}{start}
نو د مدل~\LR{$\mModel{M}$} د جوړولو پر مهال بايد د
فارمول~\LR{$!A$} د اشتقاق وړتيا نشتون په پام کښې وساتو؛ له دې امله د مدل
\LR{$\mModel{M}$} د جوړولو دپاره بشپړ سټ \LR{$\Gamma^* \supseteq \Gamma$}
نۀ شو کارولے، ځکه په هر بشپړ سټ~\LR{$\Gamma^*$} کښې
\LR{$\Gamma^* \Proves !A \lor \lnot !A$} وي۔
\olpseganchor{OLP-0506-B008}{end}

\olpseganchor{OLP-0506-B009}{start}
د بشپړ سټ~\LR{$\Gamma^*$} پر ځاے د فارمولونو د اولني سټ مفهوم کاروو:
\olpseganchor{OLP-0506-B009}{end}

\olpseganchor{OLP-0506-B010}{start}
\begin{defn}\ollabel{defn:prime}
  د فارمولونه سټ~\LR{$\Gamma$} هله او يوازې هله \emph{اولنی} دے چې
  \begin{enumerate}
  \item\ollabel{defn:prime1} \LR{$\Gamma$} سازګار وي، يعنې \LR{$\Gamma \Proves/ \lfalse$}؛
  \item\ollabel{defn:prime2} که \LR{$\Gamma \Proves !A$} وي، نو \LR{$!A \in \Gamma$} وي؛ او
  \item\ollabel{defn:prime3} که \LR{$!A \lor !B \in \Gamma$} وي، نو \LR{$!A \in
    \Gamma$} يا \LR{$!B \in \Gamma$} وي۔
  \end{enumerate}
\end{defn}
\olpseganchor{OLP-0506-B010}{end}

\olpseganchor{OLP-0506-B011}{start}
\begin{lem}[د لينډنباوم لمه]
  \ollabel{lem:lindenbaum} که \LR{$\Gamma \Proves/ !A$} وي، نو داسې
  \LR{$\Gamma^* \supseteq \Gamma$} شته چې \LR{$\Gamma^*$} اولنی وي او
  \LR{$\Gamma^* \Proves/ !A$} وي۔
\end{lem}
\olpseganchor{OLP-0506-B011}{end}

\olpseganchor{OLP-0506-B012}{start}
\begin{proof}
  \LR{$!B_1 \lor !C_1$}، \LR{$!B_2 \lor !C_2$}، \dots، دې د
  \LR{$!B \lor !C$} په بڼه د ټولو فارمولونه يوه شمېرنه وي۔ د
  فارمولونه د سټونو يوه زياتېدونکې لړۍ~\LR{$\Gamma_n$} تعريفوو،
  چې که لاندې شرطونه پوره کوونکے فصل وي، په \LR{$\Gamma_{n+1}$} کښې
  \LR{$\Gamma_n$} او يو نوے فارمول راځي۔ \LR{$\Gamma^*$} د ټولو~\LR{$\Gamma_n$}
  اتحاد دے۔ نوي فارمولونه داسې ټاکو چې \LR{$\Gamma^*$} اولنی شي او لا
  هم \LR{$\Gamma^* \Proves/ !A$} وي۔ په هر ګام کښې، که شته وي، لومړنے
  داسې فصل~\LR{$!B_i \lor !C_i$} اخلو چې:
  \begin{enumerate}
  \item\ollabel{gamma-1} \LR{$\Gamma_n \Proves !B_i \lor !C_i$}
  \item\ollabel{gamma-2} \LR{$!B_i \notin \Gamma_n$} او \LR{$!C_i \notin \Gamma_n$}
  \end{enumerate}
  يا \LR{$!B_i$} په \LR{$\Gamma_n$} کښې زياتوو، که \LR{$\Gamma_n \cup \{!B_i\}
  \Proves/ !A$} وي؛ او که داسې نۀ وي، \LR{$!C_i$} زياتوو۔ بايد وښيو
  چې دا کار سم دے۔ اوس \LR{$i(n)$} هغه تر ټولو کوچنی \LR{$i$} ټاکو چې
  \olref{gamma-1} او~\olref{gamma-2} پوره کوي۔
\olpseganchor{OLP-0506-B012}{end}
  
\olpseganchor{OLP-0506-B013}{start}
  \LR{$\Gamma_0 = \Gamma$} او
  \[
  \Gamma_{n+1} = \begin{cases}
    \Gamma_n \cup \{!B_{i(n)}\} &
    \text{\RL{که \LR{$\Gamma_n \cup \{!B_{i(n)}\} \Proves/ !A$} وي}} \\
    \Gamma_n \cup \{!C_{i(n)}\} & \text{\RL{که داسې نۀ وي}}
    \end{cases}
  \]
  تعريف کړئ۔ که \LR{$i(n)$} نۀ وي ټاکلے، يعنې هر کله چې
  \LR{$\Gamma_n \Proves !B \lor !C$} وي، نو يا \LR{$!B \in \Gamma_n$}
  يا \LR{$!C \in \Gamma_n$} وي، \LR{$\Gamma_{n+1}
  = \Gamma_n$} ږدو۔ اوس \LR{$\Gamma^* = \bigcup_{n=0}^\infty \Gamma_n$}
  ږدو۔
\olpseganchor{OLP-0506-B013}{end}

\olpseganchor{OLP-0506-B014}{start}
  لومړے ښيو چې د هر \LR{$n$} دپاره \LR{$\Gamma_n \Proves/ !A$} وي۔ پر~\LR{$n$}
  استقرا کاروو۔ د \LR{$n = 0$} لپاره دعوه د لمې له فرضه، يعنې
  \LR{$\Gamma \Proves/ !A$}، راځي۔ د استقرا په ګام کښې د هر \LR{$n\geq 0$} دپاره ښيو چې له
  \LR{$\Gamma_n \Proves/ !A$} څخه \LR{$\Gamma_{n+1} \Proves/ !A$}
  راځي۔ که \LR{$i(n)$} نۀ وي ټاکلے، \LR{$\Gamma_{n+1} = \Gamma_n$} او
  د ثابتولو څۀ نشته۔ نو فرض کړئ \LR{$i(n)$} ټاکلے وي۔ د اسانتيا دپاره
  \LR{$i = i(n)$} ږدو۔
\olpseganchor{OLP-0506-B014}{end}
  
\olpseganchor{OLP-0506-B015}{start}
  له پايلې د نفي څخه د فرض نفي اخلو۔ فرض کړئ \LR{$\Gamma_{n+1}
  \Proves !A$}۔ د جوړونې له مخې \LR{$\Gamma_{n+1} = \Gamma_n \cup
  \{!B_i\}$} وي که \LR{$\Gamma_n \cup \{!B_i\} \Proves/ !A$} وي؛
  او که داسې نۀ وي، \LR{$\Gamma_{n+1} = \Gamma_n \cup \{!C_i\}$}
  وي۔ لومړے حالت څرګند ناشونی دے، ځکه په هغه کښې به
  \LR{$\Gamma_{n+1} \Proves/ !A$} وي۔ نو \LR{$\Gamma_n \cup
  \{!B_i\} \Proves !A$} او \LR{$\Gamma_{n+1} = \Gamma_n \cup
  \{!C_i\}$} دي۔ د \LR{$i(n)$} له تعريفه \LR{$\Gamma _n \Proves !B_i \lor
  !C_i$} لرو۔ \LR{$\Gamma_n \cup \{!B_i\} \Proves !A$} هم لرو، او
  \LR{$\Gamma_n \cup \{!C_i\} \Proves !A$} هم، ځکه دا هماغه د بل
  پړاو سټ دے۔ د فصل د ايستلو له مخې \LR{$\Gamma_n \Proves !A$} راځي۔
\olpseganchor{OLP-0506-B015}{end}

\olpseganchor{OLP-0506-B016}{start}
  که \LR{$\Gamma^* \Proves !A$} وي، نو داسې متناهي فرعي سټ
  \LR{$\Gamma' \subseteq \Gamma^*$} شته چې \LR{$\Gamma' \Proves !A$}
  وي۔ که فرعي سټ تش وي، لومړے پړاو بس دے۔ که ناتش وي، د هر
  \LR{$!D \in \Gamma'$} دپاره داسې يو~\LR{$i$} اخلو چې دغه غړے په \LR{$\Gamma_i$}
  کښې وي۔ \LR{$n$} دې د اخيستل شويو شاخصونو تر ټولو لوے وي۔ څرنګه چې د
  \LR{$i \le n$} لپاره \LR{$\Gamma_i \subseteq \Gamma_n$} وي، نو
  \LR{$\Gamma' \subseteq \Gamma_n$}۔ خو بيا \LR{$\Gamma_n \Proves !A$}
  کېږي؛ دا زموږ له پورته ثبوت شوي \LR{$\Gamma_n \Proves/ !A$} سره
  په ټکر کښې دے۔
\olpseganchor{OLP-0506-B016}{end}

\olpseganchor{OLP-0506-B017}{start}
  په پاے کښې ښيو چې \LR{$\Gamma^*$} اولنی دے، يعنې د
  \olref{defn:prime} درې شرطونه \olref{defn:prime1}،
  \olref{defn:prime2} او~\olref{defn:prime3} پوره کوي۔
\olpseganchor{OLP-0506-B017}{end}

\olpseganchor{OLP-0506-B018}{start}
  لومړے، \LR{$\Gamma^* \Proves/ !A$} دے، نو \LR{$\Gamma^*$} سازګار
  دے۔ که محال ترې ثابتېدے، له محاله به هماغه فارمول هم
  ثابتېدے۔ نو \olref{defn:prime1} پوره کېږي۔
\olpseganchor{OLP-0506-B018}{end}

\olpseganchor{OLP-0506-B019}{start}
  اوس ښيو چې که \LR{$\Gamma^* \Proves !B \lor !C$} وي، نو يا
  \LR{$!B \in \Gamma^*$} يا \LR{$!C \in \Gamma^*$} وي۔ په دې سره
  \olref{defn:prime3} ثابتېږي، ځکه که \LR{$!B \lor !C \in \Gamma^*$}
  وي، نو \LR{$\Gamma^* \Proves !B \lor !C$} هم وي۔ نو فرض کړئ
  \LR{$\Gamma^* \Proves !B \lor !C$} وي، خو \LR{$!B \notin \Gamma^*$}
  او \LR{$!C \notin \Gamma^*$} وي۔ له \LR{$\Gamma^* \Proves !B \lor !C$}
  څخه د کوم~\LR{$n$} دپاره \LR{$\Gamma_n \Proves !B \lor !C$} راځي۔
  \LR{$!B \lor !C$} د ټولو فصلونو په شمېرنه کښې، ووايئ د
  \LR{$!B_j \lor !C_j$} په توګه، راځي۔ دغه \LR{$!B_j \lor !C_j$} د
  \LR{$i(n)$} د تعريف شرطونه پوره کوي: \LR{$\Gamma_n \Proves !B_j
  \lor !C_j$} او \LR{$!B_j \notin \Gamma_n$} او \LR{$!C_j \notin \Gamma_n$}
  دي۔ له دې پړاوه د دغه فصل ثبوت په ټولو ورپسې سټونو کښې
  پاتې کېږي، ځکه سټونه زياتېږي۔ د فرض له مخې يې هېڅ غړے اتحاد
  ته نۀ ننوځي، نو دغه فصل په هر ورپسې پړاو کښې دواړه شرطونه
  پوره کوي۔ کېدے شي نوي فصلونه هم ثابت شي، خو له دغه شاخصه
  ټيټ شاخصونه متناهي دي۔ هر ټيټ فصل \LR{$!B_i \lor !C_i$} چې
  غوره شي، د يوه غړي \LR{$!B_i$} يا \LR{$!C_i$} په زياتېدو سره
  تل د ټاکنې له شرطه وځي۔ ځکه په هر پړاو کښې تر ټولو ټيټ
  وړ شاخص اخلو، په يو پړاو~\LR{$m$} کښې به \LR{$j = i(m)$} شي۔
  نو يا \LR{$!B \in \Gamma_{m+1}$} يا \LR{$!C \in \Gamma_{m+1}$}
  کېږي، چې د \LR{$!B \notin \Gamma^*$} او \LR{$!C \notin \Gamma^*$}
  له فرض سره ټکر لري۔
\olpseganchor{OLP-0506-B019}{end}

\olpseganchor{OLP-0506-B020}{start}
  اوس فرض کړئ \LR{$\Gamma^* \Proves !B$}۔ نو \LR{$\Gamma^* \Proves !B \lor
  !B$} وي۔ خو تازه مو ثابت کړل چې که \LR{$\Gamma^* \Proves !B \lor !B$}
  وي، نو \LR{$!B \in \Gamma^*$} وي۔ نو \LR{$\Gamma^*$} د
  \olref{defn:prime} شرط~\olref{defn:prime2} پوره کوي۔
\end{proof}
\olpseganchor{OLP-0506-B020}{end}

\olpseganchor{OLP-0506-B021}{start}
\begin{prob}
وښيئ چې که \LR{$\Gamma \Proves/ \bot$} وي، نو \LR{$\Gamma$} په کلاسيک
منطق کښې سازګار دے؛ يعنې داسې ټاکنه شته چې د~\LR{$\Gamma$} ټول
فارمولونه رښتيا کوي۔
\end{prob}
\olpseganchor{OLP-0506-B021}{end}

\olpseganchor{OLP-0507-B004}{start}
\olfileid{int}{sc}{mod}
\olpseganchor{OLP-0507-B004}{end}

\olpseganchor{OLP-0507-B005}{start}
\olsection{کانوني مدل}
\olpseganchor{OLP-0507-B005}{end}

\olpseganchor{OLP-0507-B006}{start}
زموږ په مدل کښې نړۍ به د طبيعي عددونو متناهي لړۍ~\LR{$\sigma$}
وي، يعنې \LR{$\sigma \in \Nat^*$}۔ پام وکړئ چې \LR{$\Nat^*$} په استقرايي
ډول داسې تعريفېږي:
\begin{enumerate}
\item \LR{$\emptyseq \in \Nat^*$}۔
\item که \LR{$\sigma \in \Nat^*$} او \LR{$n \in \Nat$} وي، نو \LR{$\sigma.n \in
  \Nat^*$} وي (دلته \LR{$\sigma.n$} د \LR{$\sigma \concat \tuple{n}$} معنا
  لري او \LR{$\sigma \concat \sigma'$} د \LR{$\sigma$} او \LR{$\sigma'$} يوځاے
  کول دي)۔
\item له دې پرته بل څۀ په \LR{$\Nat^*$} کښې نشته۔
\end{enumerate}
نو د استقرايي تعريفونو دپاره \LR{$\Nat^*$} کارولے شو۔
\olpseganchor{OLP-0507-B006}{end}

\olpseganchor{OLP-0507-B007}{start}
\LR{$\tuple{!B_0, !C_0}$}، \LR{$\tuple{!B_1, !C_1}$}، \dots، دې د
فارمولونه د ټولو جوړو يوه شمېرنه وي۔ د فارمولونه د يوه
سټ~\LR{$\Delta$} په لرلو سره، \LR{$\Delta(\sigma)$} په استقرايي ډول
داسې تعريفوو:
\begin{enumerate}
\item \LR{$\Delta(\emptyseq) = \Delta$}
\item \LR{$\Delta(\sigma.n) = {}$}
  \[
  \begin{cases}
    (\Delta(\sigma) \cup \{!B_n\})^* &
    \text{\RL{که \LR{$\Delta(\sigma) \cup \{!B_n\} \Proves/ !C_n$} وي}} \\
    \Delta(\sigma) & \text{\RL{که داسې نۀ وي}}
  \end{cases}
  \]
\end{enumerate}
دلته له \LR{$(\Delta(\sigma) \cup \{!B_n\})^*$} څخه د
فارمولونه هغه اولنی سټ مراد دے چې د
\olref[lin]{lem:lindenbaum} له مخې، د \LR{$\Delta(\sigma) \cup
\{!B_n\}$} پر سټ او فارمول~\LR{$!C_n$} د لمې په تطبيقولو شته۔
پام وکړئ چې د دې تعريف له مخې، که \LR{$\Delta(\sigma) \cup
\{!B_n\} \Proves/ !C_n$} وي، نو \LR{$\Delta(\sigma.n)
\Proves !B_n$} او \LR{$\Delta(\sigma.n) \Proves/ !C_n$} وي۔ دا هم
پام وکړئ چې د هر~\LR{$n$} دپاره \LR{$\Delta(\sigma) \subseteq
\Delta(\sigma.n)$} وي۔ که \LR{$\Delta$} اولنی وي، نو د هر~\LR{$\sigma$}
دپاره \LR{$\Delta(\sigma)$} هم اولنی وي۔
\olpseganchor{OLP-0507-B007}{end}

\olpseganchor{OLP-0507-B008}{start}
\begin{defn}\ollabel{defn:canonical-model}
  فرض کړئ \LR{$\Delta$} اولنی دے۔ نو د~\LR{$\Delta$} دپاره
  \emph{کانوني مدل} \LR{$\mModel{M(\Delta)}$} داسې تعريفېږي:
  \begin{enumerate}
  \item \LR{$W = \Nat^*$}، د طبيعي عددونو د متناهي لړيو سټ۔
  \item \LR{$R$} داسې جزئي ترتيب دے چې \LR{$R\sigma\sigma'$} هله او
    يوازې هله وي چې \LR{$\sigma$} د~\LR{$\sigma'$} لومړنۍ برخه وي
    (يعنې د کومې لړۍ~\LR{$\sigma''$} دپاره \LR{$\sigma' =
    \sigma \concat \sigma''$} وي)۔
  \item \LR{$V(p) = \Setabs{\sigma}{p \in \Delta(\sigma)}$}۔
  \end{enumerate}
\end{defn}
\olpseganchor{OLP-0507-B008}{end}

\olpseganchor{OLP-0507-B009}{start}
د پيلنۍ برخې اړيکه~\LR{$R$} جزئي ترتيب دے۔ هره لړۍ خپله پيلنۍ
برخه ده؛ د پيلنۍ برخې پيلنۍ برخه هم پيلنۍ برخه ده؛ او که
دوې لړۍ د يو بل پيلنۍ برخې وي، يو شان اوږدوالے او يو شان غړي
لري۔ د~\LR{$V$} د يورخيزې زياتېدنې شرط هم پوره دے۔ ځکه چې \LR{$\Delta(\sigma)
\subseteq \Delta(\sigma.n)$} وي، نو له~\LR{$\sigma$} وروسته د ورزياتې
شوې برخې پر اوږدوالي په استقرا ښيو چې هر کله \LR{$R\sigma\sigma'$} وي،
\LR{$\Delta(\sigma)
\subseteq \Delta(\sigma')$} راځي۔ که ورزياته
شوې برخه تشه وي، سټونه برابر دي۔ هر بل زياتېدلے غړے د سټ
شمول ساتي۔ نو هر اټومي متغير چې په لومړي سټ کښې وي، په
دويم کښې هم وي۔
\olpseganchor{OLP-0507-B009}{end}

\olpseganchor{OLP-0508-B004}{start}
\olfileid{int}{sc}{tru}
\olpseganchor{OLP-0508-B004}{end}

\olpseganchor{OLP-0508-B005}{start}
\olsection{د صدق لمه}
\olpseganchor{OLP-0508-B005}{end}

\olpseganchor{OLP-0508-B006}{start}
لاندې لمه په کانوني مدل کښې د هرې نړۍ صدق له دې سره تړي
چې کوم فارمولونه د هماغې نړۍ د اولني سټ~\LR{$\Delta(\sigma)$}
غړي دي۔
\olpseganchor{OLP-0508-B006}{end}

\olpseganchor{OLP-0508-B007}{start}
\begin{lem}\ollabel{lem:truth}
  که \LR{$\Delta$} اولنی وي، نو \LR{$\mSat{M(\Delta)}{!A}[\sigma]$} هله
  او يوازې هله چې \LR{$\Delta(\sigma) \Proves !A$} وي۔
\end{lem}
\olpseganchor{OLP-0508-B007}{end}

\olpseganchor{OLP-0508-B008}{start}
\begin{proof}
  پر~\LR{$!A$} استقرا کوو۔
  \begin{enumerate}
    \item \indcase{!A}{\lfalse}{څرنګه چې \LR{$\Delta(\sigma)$} اولنی
      دے، سازګار دے؛ نو \LR{$\Delta(\sigma) \Proves/ \indfrm$} وي۔
      د تعريف له مخې \LR{$\mSat/{M(\Delta)}{\indfrm}[\sigma]$} وي۔}
      \item \indcase{!A}{p}{د \LR{$\mSat{{}}{}$} د تعريف له مخې
        \LR{$\mSat{M(\Delta)}{\indfrm}[\sigma]$} هله او يوازې هله چې
        \LR{$\sigma \in V(p)$} وي۔ د ټاکنې د تعريف له مخې دا
        اټومي متغير د هماغې نړۍ د سټ غړے وي۔ هر غړے له سټه
        د فرض په توګه ثابتېدے شي، او اولنے سټ هر ترې ثابتېدونکے
        فارمول رانغاړي۔ نو \LR{$\Delta(\sigma) \Proves
        \indfrm$} هم له دې سره برابر دے۔}
      \item \indcase!{!A}{\lnot !B}{}
      \item \indcase{!A}{!B \land
        !C}{\LR{$\mSat{M(\Delta)}{\indfrm}[\sigma]$} هله او يوازې
        هله چې \LR{$\mSat{M(\Delta)}{!B}[\sigma]$} او
        \LR{$\mSat{M(\Delta)}{!C}[\sigma]$} وي۔ د استقرا د فرض له
        مخې، \LR{$\mSat{M(\Delta)}{!B}[\sigma]$} هله او يوازې هله
        چې \LR{$\Delta(\sigma) \Proves !B$} وي، او د~\LR{$!C$} لپاره هم
        همداسې ده۔ خو \LR{$\Delta(\sigma) \Proves !B$} او
        \LR{$\Delta(\sigma) \Proves !C$} هله او يوازې هله چې
        \LR{$\Delta(\sigma) \Proves \indfrm$} وي۔}
      \item \indcase{!A}{!B \lor
        !C}{\LR{$\mSat{M(\Delta)}{\indfrm}[\sigma]$} هله او يوازې
        هله چې \LR{$\mSat{M(\Delta)}{!B}[\sigma]$} يا
        \LR{$\mSat{M(\Delta)}{!C}[\sigma]$} وي۔ د استقرا د فرض له
        مخې دا هله او يوازې هله وي چې \LR{$\Delta(\sigma) \Proves
        !B$} يا \LR{$\Delta(\sigma) \Proves !C$} وي۔ بايد وښيو چې
        دا په خپل وار هله او يوازې هله وي چې
        \LR{$\Delta(\sigma) \Proves \indfrm$} وي۔ له کيڼ څخه ښي ته
        لوری څرګند دے؛ له ښي څخه کيڼ ته لوری د
        \LR{$\Delta(\sigma)$} له اولنيوالي راځي۔}
      \item \indcase{!A}{!B \lif !C}{د لومړي لوري دپاره
        له پايلې د نفي څخه د فرض نفي اخلو۔ فرض کړئ \LR{$\Delta(\sigma)
        \Proves/ !B \lif !C$} وي۔ نو \LR{$\Delta(\sigma) \cup
        \{!B\} \Proves/ !C$} هم وي۔ که پايله له زيات شوي
        فرض سره ثابتېدے، د شرط د ورګډولو قاعدې به دغه فرض
        تړلے او له پخواني سټه به يې شرط ثابت کړے و۔ څرنګه چې
        \LR{$\tuple{!B, !C}$} د کوم~\LR{$n$} دپاره \LR{$\tuple{!B_n, !C_n}$}
        دے، نو \LR{$\Delta(\sigma.n) = (\Delta(\sigma) \cup
        \{!B\})^*$} وي، او \LR{$\Delta(\sigma.n) \Proves !B$} خو
        \LR{$\Delta(\sigma.n) \Proves/ !C$} وي۔ د استقرا د فرض له
        مخې \LR{$\mSat{M(\Delta)}{!B}[\sigma.n]$} او
        \LR{$\mSat/{M(\Delta)}{!C}[\sigma.n]$} دي۔ ځکه
        \LR{$R\sigma(\sigma.n)$} دے، له دې څخه
        \LR{$\mSat/{M(\Delta)}{\indfrm}[\sigma]$} راځي۔
\olpseganchor{OLP-0508-B008}{end}

\olpseganchor{OLP-0508-B009}{start}
        اوس فرض کړئ \LR{$\Delta(\sigma) \Proves !B \lif !C$} وي،
        او \LR{$R\sigma\sigma'$} دې وي۔ څرنګه چې \LR{$\Delta(\sigma)
        \subseteq \Delta(\sigma')$} وي، شرط له ورپسې سټه هم
        ثابتېږي۔ نو که \LR{$\Delta(\sigma') \Proves !B$} وي، د
        شرط د ايستلو له مخې \LR{$\Delta(\sigma') \Proves !C$} هم وي۔
        په بله وينا، د هر
        داسې \LR{$\sigma'$} دپاره چې \LR{$R\sigma\sigma'$} وي، يا
        \LR{$\Delta(\sigma') \Proves/ !B$} يا
        \LR{$\Delta(\sigma') \Proves !C$} وي۔ د استقرا د فرض له
        مخې، دا په دې معنا ده چې هر کله \LR{$R\sigma\sigma'$} وي،
        يا \LR{$\mSat/{M(\Delta)}{!B}[\sigma']$} يا
        \LR{$\mSat{M(\Delta)}{!C}[\sigma']$} وي؛ يعنې
        \LR{$\mSat{M(\Delta)}{\indfrm}[\sigma]$} وي۔}
  \end{enumerate}
\end{proof}
\olpseganchor{OLP-0508-B009}{end}

\olpseganchor{OLP-0509-B004}{start}
\olfileid{int}{sc}{cpl}
\olpseganchor{OLP-0509-B004}{end}

\olpseganchor{OLP-0509-B005}{start}
\olsection{د بشپړتيا قضيه}
\olpseganchor{OLP-0509-B005}{end}

\olpseganchor{OLP-0509-B006}{start}
\begin{thm}\ollabel{thm:completeness}
  که \LR{$\Gamma \Entails !A$} وي، نو \LR{$\Gamma \Proves !A$} وي۔
\end{thm}
\olpseganchor{OLP-0509-B006}{end}

\olpseganchor{OLP-0509-B007}{start}
\begin{proof}
  له پايلې د نفي څخه د فرض نفي اخلو۔ فرض کړئ \LR{$\Gamma \Proves/ !A$} وي۔ نو د
  \olref[lin]{lem:lindenbaum} له مخې داسې اولنی سټ~\LR{$\Gamma^*
  \supseteq \Gamma$} شته چې \LR{$\Gamma^* \Proves/ !A$} وي۔ د
  \LR{$\Gamma^*$} دپاره هغه کانوني مدل~\LR{$\mModel{M(\Gamma^*)}$}
  وګورئ چې په \olref[mod]{defn:canonical-model} کښې تعريف شوے
  دے۔ د هر \LR{$!B \in \Gamma$} دپاره \LR{$\Gamma^* \Proves !B$}
  وي، ځکه غځېدلے سټ دا فارمول د فرض په توګه لري۔
  \LR{$\Gamma^*(\emptyseq) = \Gamma^*$} دے۔ د صدق
  د لمې (\olref[tru]{lem:truth}) له مخې د هر \LR{$!B \in \Gamma$}
  دپاره \LR{$\mSat{M(\Gamma^*)}{!B}[\emptyseq]$} او
  \LR{$\mSat/{M(\Gamma^*)}{!A}[\emptyseq]$} لرو۔ دا ښيي چې
  \LR{$\Gamma \Entails/ !A$} دے۔
\end{proof}
\olpseganchor{OLP-0509-B007}{end}

\olpseganchor{OLP-0509-B008}{start}
\begin{prob}
  وښيئ چې که \LR{$!A$} يوازې بياني متغيرونه، \LR{$\lor$}
  او~\LR{$\land$} ولري، نو \LR{$\Entails/ !A$} وي۔ له دې څخه نتيجه
  واخلئ چې \LR{$\lif$} په شهودي منطق کښې له \LR{$\lor$} او~\LR{$\land$}
  څخه نۀ تعريفېږي۔
\end{prob}
\olpseganchor{OLP-0509-B008}{end}
  
\olpseganchor{OLP-0509-B009}{start}
\begin{prob}
  د بشپړتيا د قضيې په کارولو ثابت کړئ چې که \LR{$\Proves !A \lor !B$}
  وي، نو \LR{$\Proves !A$} يا \LR{$\Proves !B$} وي۔ (لارښوونه: فرض
  کړئ \LR{$\mSat/{M_1}{!A}$} او \LR{$\mSat/{M_2}{!B}$} وي او داسې
  نوے مدل~\LR{$\mModel{M}$} جوړ کړئ چې \LR{$\mSat/{M}{!A \lor !B}$} وي۔)
\end{prob}
\olpseganchor{OLP-0509-B009}{end}
  
\olpseganchor{OLP-0509-B010}{start}
\begin{prob} 
  وښيئ چې که \LR{$\mModel{M}$} داسې اړيکيز مدل وي چې خطي ترتيب
  کاروي، نو \LR{$\mSat{M}{(!A \lif !B)\lor(!B \lif !A)}$} وي۔
\end{prob}
\olpseganchor{OLP-0509-B010}{end}

\olpseganchor{OLP-0510-B004}{start}
\olfileid{int}{sc}{dec}
\olpseganchor{OLP-0510-B004}{end}

\olpseganchor{OLP-0510-B005}{start}
\olsection{پرېکړې وړتيا}
پام وکړئ چې د بشپړتيا د قضيې ثبوت د هر
\LR{$\Gamma \Proves/ !A$} دپاره د نامتناهي شمېر نړۍو لرونکے داسې مدل
راکوي چې \LR{$\Gamma \Entails/ !A$} پرې ښيي۔ لاندې قضيه وايي چې
د \LR{$\Entails !A$} د ثابتولو دپاره د ټولو متناهي مدلونو (يعنې د
متناهي نړۍو سټ لرونکو مدلونو) لپاره د \LR{$\mSat{M}{!A}$} ثابتول
کافي دي۔
\olpseganchor{OLP-0510-B005}{end}

\olpseganchor{OLP-0510-B006}{start}
\begin{thm}\ollabel{thm:decidability}
  که \LR{$\Entails/ !A$} وي، نو داسې متناهي مدل شته چې
  \LR{$\mSat/{M'}{!A}$} وي۔
\end{thm}
\olpseganchor{OLP-0510-B006}{end}

\olpseganchor{OLP-0510-B007}{start}
\begin{proof}
  فرض کړئ \LR{$\mModel{M}=\tuple{W, R, V}$} داسې وي چې
  \LR{$\mSat/{M}{!A}$}۔ تعريف شوي نښلوونکي پر خپل تعريف پرانيزو او
  د~\LR{$!A$} د ټولو فرعي فارمولونو سټ~\LR{$S$} اخلو؛ په دې سټ کښې
  فارمول پخپله او ټول راغلي بياني متغيرونه هم دي۔
  د هرې نړۍ د رښتينو فرعي فارمولونو سټ داسې ليکو:
  \[
    t(w) = \Setabs{!B \in S}{\mSat{M}{!B}[w]}.
  \]
  \LR{$\mModel{M'}=\tuple{W', R', V'}$} داسې تعريفوو چې
  \LR{$W' = \Setabs{t(w)}{w \in W}$} وي، \LR{$R'tu$} هله او يوازې
  هله وي چې \LR{$t \subseteq u$} وي، او
  \LR{$V'(p)=\Setabs{t \in W'}{p \in t}$} وي۔ \LR{$W'$} ناتش دے او
  تر~\LR{$2^{|S|}$} زيات غړي نۀ لري۔ د سټونو شمول جزئي ترتيب دے
  او په دې ترتيب د اټومي متغير غړيتوب پاتې کېږي۔ نو
  \LR{$\mModel{M'}$} اړيکيز مدل دے۔
\olpseganchor{OLP-0510-B007}{end}

\olpseganchor{OLP-0510-B008}{start}
  پر فرعي فارمولونه~\LR{$!B \in S$} استقرا کوو او په هرې
  نړۍ کښې دا برابري ثابتوو:
  \[
    \mSat{M}{!B}[w] \text{\RL{ هله او يوازې هله چې }} \mSat{M'}{!B}[t(w)].
  \]
  د بياني متغيرونه حالت د ټاکنې له تعريفه راځي،
  او محال په دواړو مدلونو کښې په هېڅ نړۍ کښې صدق نۀ لري۔
  د عطف او فصل دواړه غړي هم په~\LR{$S$} کښې دي، نو د هغو د صدق
  په اړه د استقرا فرض د عطف او فصل صدق هم ساتي۔ اوس
  \LR{$!B = !C \lif !D$} وګورئ۔ که \LR{$\mSat{M}{!B}[w]$} وي او
  \LR{$t(w)\subseteq t(v)$} وي، نو شرط په~\LR{$t(v)$} کښې هم غړے
  دے او په اصلي نړۍ~\LR{$v$} کښې رښتينے دے۔ که مخکښ فارمول
  په نوي مدل کښې په~\LR{$t(v)$} کښې رښتينے وي، د استقرا له مخې
  په اصلي نړۍ~\LR{$v$} کښې هم رښتينے دے۔ د انعکاسيت له مخې
  شرط په هماغې نړۍ کښې پلي کوو او د استقرا په کارولو
  يې پايله بېرته~\LR{$t(v)$} ته وړو۔ نو شرط په نوې نړۍ~\LR{$t(w)$}
  کښې رښتينے دے۔ که \LR{$\mSat/{M}{!B}[w]$} وي، داسې \LR{$v$} شته
  چې \LR{$Rwv$}، \LR{$\mSat{M}{!C}[v]$} او \LR{$\mSat/{M}{!D}[v]$} وي۔
  د صدق د پاتې کېدو له مخې \LR{$t(w)\subseteq t(v)$} دے او
  د استقرا له مخې په~\LR{$t(v)$} کښې مخکښ رښتينے او پايله
  نارښتينې ده۔ نو شرط په~\LR{$t(w)$} کښې هم نارښتينے دے۔
  اوس \LR{$!A \in S$} دے۔ په اصلي مدل کښې هره نړۍ چې \LR{$!A$}
  پکښې نارښتينے وي، په متناهي مدل کښې يې انځور هم
  \LR{$!A$} نارښتينے کوي۔
\end{proof}
\olpseganchor{OLP-0510-B008}{end}

\olpseganchor{OLP-0510-B009}{start}
\begin{prob}
د \olref[int][sc][dec]{thm:decidability} په استقرا کښې
د محال، اټومي متغير، عطف او فصل حالتونه په تفصيل وليکئ۔ د
دغو فارمولونه~\LR{$!B \in S$} دپاره وښيئ چې
\LR{$\mSat{M}{!B}[w]$} هله او يوازې هله وي چې
\LR{$\mSat{M'}{!B}[t(w)]$} وي۔ بيا داسې دوې نړۍ ومومئ چې
يو شان اټومي ارزښتونه لري، خو يو شرط په يوه نړۍ کښې
رښتينے او په بله کښې نارښتينے وي۔
\end{prob}
\olpseganchor{OLP-0510-B009}{end}

\olpseganchor{OLP-0510-B010}{start}
له \olref{thm:decidability} څخه دا راځي چې د~\LR{$\Entails !A$}
د پرېکړې دپاره الګوريتم شته۔ تر پورته متناهي حده ټول
جزئي ترتيبونه او د اړوندو متغيرونو ټولې يورخيزې ټاکنې
شمېرئ، او په هر مدل کښې د فارمول صدق وګورئ۔
\olpseganchor{OLP-0510-B010}{end}

\olpseganchor{OLP-0511-B004}{start}
\olchapter{int}{tab}{شهودي تابلوګانې}
\olpseganchor{OLP-0511-B004}{end}

\olpseganchor{OLP-0511-B005}{start}
\begin{editorial}
  مخکښ‌لرونکې تابلوګانې د فارمول صدق او ناصدق د مدل له
  نړۍو سره تړي۔ په دې باب کښې يې قاعدې او بېلګې راځي، او
  سموالے يې ښيي چې تړلې تابلو د پيل د فرضونو ناشونتيا ثابتوي۔
\end{editorial}
\olpseganchor{OLP-0511-B005}{end}

\olpseganchor{OLP-0511-B006}{start}




%
%
\olpseganchor{OLP-0511-B006}{end}

\olpseganchor{OLP-0512-B004}{start}
\olfileid{int}{tab}{int}
\olpseganchor{OLP-0512-B004}{end}

\olpseganchor{OLP-0512-B005}{start}
\olsection{سريزه}
\olpseganchor{OLP-0512-B005}{end}

\olpseganchor{OLP-0512-B006}{start}
عادي تابلوګانې د نښه لرونکي فارمولونه ځانګړې ونې دي چې لاندې خوا ته څانګې کوي۔ هره
نښه‌لرونکې جوړه د صدق نښه (\LR{$\True$} يا \LR{$\False$}) او تړلے فارمول لري:
\[
\sFmla{\True}{!A} \text{\RL{ يا }} \sFmla{\False}{!A}.
\]
تابلو له څو \emph{فرضونو} پيلېږي۔ هر بل
نښه لرونکے فارمول د استنتاج د يوې قاعدې په تطبيقولو جوړېږي۔
ځينې قاعدې د ونې يوه څوکه ته يو يا څو نښه لرونکي فارمولونه
ورزياتوي؛ نورې دوه نوې څوکې ورزياتوي او دوه څانګې جوړوي۔ د
قاعدې پايله داسې نښه لرونکي فارمولونه وي چې فارمول يې تر
هغه نښه لرونکے فارمول ساده وي چې قاعده پرې تطبيق شوې وي۔
کله چې په يوه څانګه کښې \LR{$\sFmla{\True}{!A}$} او
\LR{$\sFmla{\False}{!A}$} دواړه وي، هغه څانګه \emph{تړلې} بولو۔
که د تابلو هره څانګه تړلې وي، ټول تابلو تړلے
دے۔ تړلے تابلو داسې اشتقاق جوړوي چې ښيي د
تابلو د پيل نښه لرونکي فارمولونه سټ نۀ شي پوره کېدے۔ له دې څخه د
\LR{$\Proves$} اړيکه تعريفولے شو: \LR{$\Gamma \Proves !A$} هله او
يوازې هله چې داسې متناهي سټ~\LR{$\Gamma_0 = \{!B_1,
\dots, !B_n\} \subseteq \Gamma$} شته چې د لاندې فرضونو دپاره
تړلے تابلو وي:
\[
\{\sFmla{\False}{!A}, \sFmla{\True}{!B_1}, \dots, \sFmla{\True}{!B_n}\}.
\]
\olpseganchor{OLP-0512-B006}{end}

\olpseganchor{OLP-0512-B007}{start}
د شهودي منطق دپاره بايد د نښه لرونکے فارمول مفهوم هم پراخ
کړو او د نښلوونکو قاعدې هم بدلې کړو۔ د نښې (\LR{$\True$} يا
\LR{$\False$}) تر څنګ، په شهودي تابلوګانې کښې
فارمولونه \emph{مخکښونه}~\LR{$\sigma$} هم لري۔ مخکښونه د
مثبتو صحيح عددونو ناتشې لړۍ دي، يعنې \LR{$\sigma \in
(\PosInt)^* \setminus \{\emptyseq\}$}۔ مخکښونه له چاپېرې
\LR{$\tuple{\ }$} پرته ليکو او بېل غړي يې د
\LR{$.$} په وسيله جلا کوو، نه د \LR{$,$} په وسيله۔ که \LR{$\sigma$} يو
مخکښ وي، نو \LR{$\sigma.n$} د \LR{$\sigma \concat \tuple{n}$}
معنا لري؛ مثلاً که \LR{$\sigma = 1.2.1$} وي، نو \LR{$\sigma.3$}
به \LR{$1.2.1.3$} وي۔ د بېلګې په توګه،
\[
\sFmla{\True}{!A \lif (!B \lif !C)}[1.2]
\]
يو \emph{مخکښ‌لرونکے نښه لرونکے فارمول} دے (يا په لنډ ډول
يو \emph{مخکښ‌لرونکے فارمول})۔
\olpseganchor{OLP-0512-B007}{end}

\olpseganchor{OLP-0512-B008}{start}
په تصور کښې مخکښ د مدل يوه نړۍ نوموي چې ښايي د
تابلو د يوې څانګې فارمولونه پوره کړي؛ که
\LR{$\sigma$} کومه نړۍ ونوموي، نو \LR{$\sigma.n$} له~\LR{$\sigma$} نومولې نړۍ څخه
لاس‌رسي وړ نړۍ نوموي۔
\olpseganchor{OLP-0512-B008}{end}

\olpseganchor{OLP-0512-B009}{start}
په شهودي مدلونو کښې د لاسرسي اړيکه انعکاسي او متعديه ده۔
د مخکښونو په ترتيب کښې له~\LR{$\sigma$} څخه \LR{$\sigma$} پخپله او
هر مخکښ چې~\LR{$\sigma$} غځوي لاسرسي وړ دے، يعنې هر مخکښ چې د
\LR{$\sigma.n_1.\cdots.n_k$} په بڼه وي۔ د~\LR{$\sigma$} خپل ځان
او د هغه د هر غځول شوي مخکښ دپاره \LR{$\sigma.*$} ليکو۔
په بله وينا، د \LR{$\sigma.*$} مخکښونه ټول هغه او يوازې هغه
دي چې د مخکښونو په ترتيب کښې له~\LR{$\sigma$} څخه لاسرسي وړ دي۔
په يوه مدل کښې نومول شوې نړۍ نورې لاسرسي وړ نړۍ هم لرلے شي۔
\olpseganchor{OLP-0512-B009}{end}

\olpseganchor{OLP-0513-B004}{start}
\olfileid{int}{tab}{rul}
\olpseganchor{OLP-0513-B004}{end}

\olpseganchor{OLP-0513-B005}{start}
\olsection{د شهودي منطق قاعدې}
\olpseganchor{OLP-0513-B005}{end}

\olpseganchor{OLP-0513-B006}{start}
د \LR{$\land$} او \LR{$\lor$} نښلوونکو قاعدې د عادي قضييزو نښه‌لرونکو
تابلوګانې له قاعدو سره يو شان دي؛ يوازې مخکښونه ورزيات شوي
دي۔ په هر حالت کښې پر نښه‌لرونکي فارمول
\LR{$\sFmla{S}{!A}[\sigma]$} پلې شوې قاعده داسې نوي
فارمولونه جوړوي چې مخکښ يې هم~\LR{$\sigma$} وي۔ دا له مفهومي
پلوه څرګنده ده: مثلاً که \LR{$!A \land !B$} په~\LR{$\sigma$} نومولې
نړۍ کښې رښتيا وي، نو \LR{$!A$} او \LR{$!B$} هم په~\LR{$\sigma$} کښې رښتيا
دي۔ په ټولو لاسرسي وړ نړۍو کښې يې صدق هم پاتې کېږي، خو
د عطف د ماتولو دا قاعده دواړه غړي په هماغه مخکښ ليکي۔
د \LR{$\land$} او \LR{$\lor$} قاعدې
په \olref{tab:prop-rules} کښې راټولوو۔
\olpseganchor{OLP-0513-B006}{end}

\olpseganchor{OLP-0513-B007}{start}
\begin{table}
  \[\def\arraystretch{3}\begin{array}{|c|c|}
    \hline
    \AxiomC{\sFmla{\True}{!A \land !B}[\sigma]}
    \RightLabel{\TRule{\True}{\land}}
    \UnaryInfC{\sFmla{\True}{!A}[\sigma]}
    \noLine
    \UnaryInfC{\sFmla{\True}{!B}[\sigma]}
    \DisplayProof
    &
    \AxiomC{\sFmla{\False}{!A \land !B}[\sigma]}
    \RightLabel{\TRule{\False}{\land}}
    \UnaryInfC{\LR{$\sFmla{\False}{!A}[\sigma] \quad \mid \quad
      \sFmla{\False}{!B}[\sigma]$}}
    \DisplayProof
    \\[2ex]
    \hline
    \AxiomC{\sFmla{\True}{!A \lor !B}[\sigma]}
    \RightLabel{\TRule{\True}{\lor}}
    \UnaryInfC{\LR{$\sFmla{\True}{!A}[\sigma] \quad \mid \quad
      \sFmla{\True}{!B}[\sigma]$}}
    \DisplayProof
    &
    \AxiomC{\sFmla{\False}{!A \lor !B}[\sigma]}
    \RightLabel{\TRule{\False}{\lor}}
    \UnaryInfC{\sFmla{\False}{!A}[\sigma]}
    \noLine
    \UnaryInfC{\sFmla{\False}{!B}[\sigma]}
    \DisplayProof
    \\[2ex]
    \hline
  \end{array}\]
  \caption{د \LR{$\land$} او \LR{$\lor$} لپاره د مخکښ‌لرونکي تابلو قاعدې}
  \ollabel{tab:prop-rules}
\end{table}
\olpseganchor{OLP-0513-B007}{end}

\olpseganchor{OLP-0513-B008}{start}
د تړل کېدو شرط د عادي تابلوګانې له شرط سره ورته دے، خو
يوازې د فارمولونه برابري کافي نۀ ده؛ مخکښونه هم بايد سره
برابر يا د لاسرسي له مخې اړوند وي۔ په حقيقت کښې لږ پراخ شرط
اخيستلے شو: په شهودي مدلونو کښې فارمولونه چې يو ځل رښتيا
شي، بيا رښتيا پاتې کېږي؛ نو ناشونې ده چې \LR{$!A$} په~\LR{$\sigma$}
کښې رښتيا خو په کوم لاس‌رسي وړ مخکښ~\LR{$\sigma.{*}$} کښې کاذب
وي۔ نو څانګه هله تړلې وي چې دواړه لاندې نښې ولري:
\[
\sFmla{\True}{!A}[\sigma] \quad\text{\RL{او}}\quad \sFmla{\False}{!A}[\sigma.{*}]
\]
د کوم مخکښ~\LR{$\sigma$} او فارمول~\LR{$!A$} دپاره۔ پام وکړئ چې
که نښې سرچپه وي، يعنې څانګه دا ولري:
\[
\sFmla{\False}{!A}[\sigma] \quad\text{\RL{او}}\quad \sFmla{\True}{!A}[\sigma.{*}]
\]
نو دغه دوې کرښې پخپله څانګه يوازې هغه وخت تړي چې \LR{$*$} تشه
لړۍ وي۔ په دوو پرله‌پسې نړۍو کښې يو اټومي فارمول په لومړۍ
نړۍ کښې نارښتينے او په دويمه کښې رښتينے کېدے شي۔
\olpseganchor{OLP-0513-B008}{end}

\olpseganchor{OLP-0513-B009}{start}
پر دې سربېره، که څانګه~\LR{$\sFmla{\True}{\bot}[\sigma]$} ولري،
هغه هم تړلې ده۔
\olpseganchor{OLP-0513-B009}{end}

\olpseganchor{OLP-0513-B010}{start}
د فرضونو د ايښودلو قاعدې هم د عادي تابلوګانې په شان دي،
خو د فرضونو دپاره تل مخکښ~\LR{$1$} کاروو۔ (کوم مخکښ چې کاروو مهم
نۀ دے، بس د ټولو فرضونو دپاره يو شان وي۔) نو مثلاً وايو چې
\[
!B_1, \dots, !B_n \Proves !A
\]
هله او يوازې هله چې د لاندې فرضونو دپاره تړلے تابلو وي:
\[
\sFmla{\True}{!B_1}[1], \dots, \sFmla{\True}{!B_n}[1],
\sFmla{\False}{!A}[1].
\]
\olpseganchor{OLP-0513-B010}{end}

\olpseganchor{OLP-0513-B011}{start}
د شرط~\LR{$\lif$} قاعدې له کلاسيکو او مودالي قاعدو بېلې دي۔ د
\LR{$\TRule{\True}{\lif}$} قاعده هغه څانګه چې
\LR{$\sFmla{\True}{!A \lif !B}[\sigma]$} لري، په دوو بېلو
څانګو کښې په ترتيب \LR{$\sFmla{\False}{!A}[\sigma.{*}]$} يا
\LR{$\sFmla{\True}{!B}[\sigma.{*}]$} ورزياتولو غځوي۔ د رښتيني
شرط په هر لاسرسي وړ مخکښ کښې مخکښ فارمول يا نارښتينے وي،
يا وروستے فارمول رښتينے وي؛ هره څانګه له دغو حالاتو يو اخلي۔
قاعده يوازې پر داسې مخکښ~\LR{$\sigma.{*}$} پلې کېدے شي چې په اړونده څانګه کښې
\emph{لا مخکې} راغلے وي۔ داسې مخکښ «کارېدلے» بولو۔ (څرنګه چې
\LR{$\sigma.{*}$} خپله \LR{$\sigma$} هم رانغاړي، قاعده تل د
\LR{$\sFmla{\False}{!A}[\sigma]$} او
\LR{$\sFmla{\True}{!B}[\sigma]$} په بېلو څانګو کښې په زياتولو
پلې کېدے شي۔)
\olpseganchor{OLP-0513-B011}{end}

\olpseganchor{OLP-0513-B012}{start}
د \LR{$\TRule{\False}{\lif}$} قاعده هغه څانګه چې
\LR{$\sFmla{\False}{!A \lif !B}[\sigma]$} لري، په هماغه څانګه
کښې د \LR{$\sFmla{\True}{!A}[\sigma.n]$} او
\LR{$\sFmla{\False}{!B}[\sigma.n]$} دواړو په زياتولو غځوي؛
\LR{$\sigma.n$} بايد په څانګه کښې نوے مخکښ وي۔
\olpseganchor{OLP-0513-B012}{end}

\olpseganchor{OLP-0513-B013}{start}
د \LR{$\lnot$} قاعدې په ورته ډول تعريفېږي (د \LR{$\lnot !A$} د
\LR{$!A \lif \lfalse$} په توګه تعريف کاروو)۔
\olpseganchor{OLP-0513-B013}{end}

\olpseganchor{OLP-0513-B014}{start}
قاعدې په \olref{tab:rules-lif-lnot} کښې ورکړل شوې دي۔
\olpseganchor{OLP-0513-B014}{end}

\olpseganchor{OLP-0513-B015}{start}
\begin{table}
  \[\def\arraystretch{3}\begin{array}{|c|c|}
    \hline
    \AxiomC{\sFmla{\True}{\lnot !A}[\sigma]}
    \RightLabel{\TRule{\True}{\lnot}}
    \UnaryInfC{\LR{$\sFmla{\False}{!A}[\sigma.{*}]$}}
    \DisplayProof
    &
    \AxiomC{\sFmla{\False}{\lnot !A}[\sigma]}
    \RightLabel{\TRule{\False}{\lnot}}
    \UnaryInfC{\sFmla{\True}{!A}[\sigma.n]}
    \DisplayProof
    \\[1ex]
    \text{\RL{\LR{$\sigma.{*}$} کارېدلے دے}} & \text{\RL{\LR{$\sigma.n$} نوے دے}}\\
    \hline
    \AxiomC{\sFmla{\True}{!A \lif !B}[\sigma]}
    \RightLabel{\TRule{\True}{\lif}}
    \UnaryInfC{\LR{$\sFmla{\False}{!A}[\sigma.{*}] \quad \mid \quad
      \sFmla{\True}{!B}[\sigma.{*}]$}}
    \DisplayProof
    &
    \AxiomC{\sFmla{\False}{!A \lif !B}[\sigma]}
    \RightLabel{\TRule{\False}{\lif}}
    \UnaryInfC{\sFmla{\True}{!A}[\sigma.n]}
    \noLine
    \UnaryInfC{\sFmla{\False}{!B}[\sigma.n]}
    \DisplayProof
    \\[1ex]
    \text{\RL{\LR{$\sigma.{*}$} کارېدلے دے}} & \text{\RL{\LR{$\sigma.n$} نوے دے}}\\
    \hline
  \end{array}\]
  \caption{د \LR{$\lnot$} او \LR{$\lif$} لپاره د مخکښ‌لرونکي تابلو قاعدې}
  \ollabel{tab:rules-lif-lnot}
\end{table}
\olpseganchor{OLP-0513-B015}{end}

\olpseganchor{OLP-0514-B004}{start}
\olfileid{int}{tab}{prf}
\olpseganchor{OLP-0514-B004}{end}

\olpseganchor{OLP-0514-B005}{start}
\olsection{د شهودي منطق لپاره تابلوګانې}
\olpseganchor{OLP-0514-B005}{end}

\olpseganchor{OLP-0514-B006}{start}
\begin{ex}
  داسې تړلے تابلو وړاندې کوو چې \LR{$(!A \land !B) \lif !C \Proves
  !A \lif (!B \lif !C)$} ښيي۔ په اوومه کرښه کښې نارښتينے
  عطف دوه څانګې جوړوي۔ په لومړۍ کښې له پخواني مخکښ څخه
  د عطف د لومړي غړي صدق پاتې کېږي، او په دويمه کښې په هماغه
  مخکښ دويم غړے رښتينے دے۔ په دريمه څانګه کښې پايله په يوه
  مخکښ هم رښتينې او هم نارښتينې ده۔
  \begin{oltableau}
    [\pFmla{\True}{(\formula{A} \land \formula{B}) \lif \formula{C}}{1},
      just =\TAss
      [\pFmla{\False}{\formula{A} \lif (\formula{B} \lif \formula{C})}{1},
        just =\TAss
        [\pFmla{\True}{\formula{A}}{1.1},
          just = {\TRule{\False}{\lif}[2]}
          [\pFmla{\False}{\formula{B} \lif \formula{C}}{1.1},
            just = {\TRule{\False}{\lif}[2]}
            [\pFmla{\True}{\formula{B}}{1.1.1},
              just = {\TRule{\False}{\lif}[4]}
              [\pFmla{\False}{\formula{C}}{1.1.1},
                just = {\TRule{\False}{\lif}[4]}
                [\pFmla{\False}{\formula{A} \land \formula{B}}{1.1.1},
                  just = {\TRule{\True}{\lif}[1]}
                  [\pFmla{\False}{\formula{A}}{1.1.1},
                    just= {\TRule{\False}{\land}[7]}, close]
                  [\pFmla{\False}{\formula{B}}{1.1.1},
                    just= {\TRule{\False}{\land}[7]}, close]]
                [\pFmla{\True}{\formula{C}}{1.1.1},
                  just = {\TRule{\True}{\lif}[1]}, close]]
            ]
          ]
        ]
      ]
    ]
  \end{oltableau}
\end{ex}
\olpseganchor{OLP-0514-B006}{end}



\olpseganchor{OLP-0514-B007}{start}
\begin{prob}
  د لاندې بيانونو د ښودلو لپاره تړلي شهودي تابلوګانې ومومئ:
  \begin{enumerate}
    \item \LR{$\Proves !A \lif (!B \lif !A)$}
    \item \LR{$\Proves \lnot(!A \land \lnot !A)$}
    \item \LR{$!A \lif (!B \lif !C) \Proves (!A \land !B) \lif !C$}
    \item \LR{$\lnot !A \lor \lnot !B \Proves \lnot(!A \land !B)$}
  \end{enumerate}
\end{prob}
\olpseganchor{OLP-0514-B007}{end}

\olpseganchor{OLP-0515-B004}{start}
\olfileid{int}{tab}{sou}
\olpseganchor{OLP-0515-B004}{end}

\olpseganchor{OLP-0515-B005}{start}
\olsection{د شهودي تابلوګانې سموالے}
\olpseganchor{OLP-0515-B005}{end}

\olpseganchor{OLP-0515-B006}{start}
\begin{explain}
  د دې ښودلو لپاره چې شهودي تابلوګانې سم دي، بايد وښيو چې که
  \[
  \sFmla{\True}{!B_1}[1], \dots, \sFmla{\True}{!B_n}[1], \sFmla{\False}{!A}[1]
  \]
  تړلے تابلو ولري، نو \LR{$!B_1, \dots, !B_n \Entails !A$}۔
  له پايلې د نفي څخه د فرض نفي اخلو۔ که د کوم \LR{$\mModel{M}$} او نړۍ~\LR{$w$}
  لپاره د ټولو \LR{$i=1$}، \dots,~\LR{$n$} دپاره \LR{$\mSat{M}{!B_i}[w]$}
  وي خو \LR{$\mSat/{M}{!A}[w]$} وي، نو هېڅ تابلو نۀ شي تړلے۔
  داسې مخالف مدل ښيي چې د تابلو لومړني
  فرضونه د صدق وړ دي۔ د ثبوت تګلاره دا ده چې وښيو: کله هم
  چې د تابلو پر څانګه ټول مخکښ‌لرونکي فارمولونه په
  ګډه د صدق وړ وي، د هرې قاعدې پلې کول لږ تر لږه يوه غځېدلې
  څانګه ورکوي چې هغه هم د صدق وړ وي۔ تړلې څانګې د صدق وړ نۀ
  دي؛ نو د مخکښ‌لرونکو فارمولونه د صدق وړ سټ هر
  تابلو بايد لږ تر لږه يوه پرانيستې څانګه ولري۔
\olpseganchor{OLP-0515-B006}{end}

\olpseganchor{OLP-0515-B007}{start}
  په شهودي حالت کښې د دې تګلارې د پلې کولو لپاره بايد د
  «د صدق وړ» تعريف اړيکيزو مدلونو او مخکښونو ته وغځوو۔
  د مخکښونو تفسير ټاکي چې هره نښه په کومه نړۍ کښې وکتل شي۔
\end{explain}
\olpseganchor{OLP-0515-B007}{end}

\olpseganchor{OLP-0515-B008}{start}
\begin{defn}
  \LR{$P$} دې د مخکښونو يو سټ وي، يعنې \LR{$P \subseteq (\PosInt)^*
  \setminus \{\emptyseq\}$}، او \LR{$\mModel{M}$} دې يو مدل وي۔
  تابع~\LR{$f\colon P \to W$} په~\LR{$\mModel{M}$} کښې د~\LR{$P$}
  \emph{تفسير} دے که، هر کله چې \LR{$\sigma$} او د هغه يو
  غځېدلے مخکښ~\LR{$\sigma.{*}$} دواړه په~\LR{$P$} کښې وي، نو
  \LR{$Rf(\sigma)f(\sigma.{*})$} هم وي۔ د تفسير شرط په دغو دوو
  مخکښونو پورې تړلے دے، که منځني مخکښونه موجود وي او که نه۔
\olpseganchor{OLP-0515-B008}{end}

\olpseganchor{OLP-0515-B009}{start}
  د مخکښونو~\LR{$P$} د يوه تفسير په نسبت دا تعريفولے شو:
  \begin{enumerate}
  \item \LR{$\mModel{M}$} په \LR{$\sFmla{\True}{!A}[\sigma]$} کښې صدق کوي هله او يوازې هله چې
    \LR{$\mSat{M}{!A}[f(\sigma)]$}.
  \item \LR{$\mModel{M}$} په \LR{$\sFmla{\False}{!A}[\sigma]$} کښې صدق کوي هله او يوازې هله چې
    \LR{$\mSat/{M}{!A}[f(\sigma)]$}.
  \end{enumerate}
\end{defn}
\olpseganchor{OLP-0515-B009}{end}

\olpseganchor{OLP-0515-B010}{start}
\LR{$R$} انعکاسي او متعدي دے، او \LR{$\sigma.{*}$} د
\LR{$\sigma$}، \LR{$\sigma.n_1$}، \LR{$\sigma.n_1.n_2$}، \dots نښه ده۔
نو که دواړه مخکښونه د تفسير په ساحه کښې وي، د تعريف له
مخې \LR{$Rf(\sigma)f(\sigma.{*})$} لرو۔
\olpseganchor{OLP-0515-B010}{end}

\olpseganchor{OLP-0515-B011}{start}
\begin{defn}
  \LR{$\Gamma$} دې د مخکښ‌لرونکو فارمولونه سټ وي، او \LR{$P(\Gamma)$}
  دې په هغه کښې د راغلو مخکښونو سټ وي۔ که \LR{$f$} په
  \LR{$\mModel{M}$} کښې د~\LR{$P(\Gamma)$} تفسير وي، وايو چې
  \LR{$\mModel{M}$} د~\LR{$f$} په نسبت په \LR{$\Gamma$} کښې صدق کوي،
  \LR{$\mSat{M}{\Gamma}[f]$}، که \LR{$\mModel{M}$} د~\LR{$f$} په نسبت د
  \LR{$\Gamma$} په هر مخکښ‌لرونکي فارمول کښې صدق کوي۔
  \LR{$\Gamma$} \emph{د صدق وړ} دے هله او يوازې هله چې يو
  مدل~\LR{$\mModel{M}$} او د \LR{$P(\Gamma)$} تفسير~\LR{$f$} داسې وي چې
  \LR{$\mSat{M}{\Gamma}[f]$}۔
\end{defn}
\olpseganchor{OLP-0515-B011}{end}

\olpseganchor{OLP-0515-B012}{start}
\begin{prop}
  که \LR{$\Gamma$} د کوم فارمول~\LR{$!A$} او مخکښ~\LR{$\sigma$} دپاره
  \LR{$\sFmla{\True}{!A}[\sigma]$} او
  \LR{$\sFmla{\False}{!A}[\sigma.{*}]$} دواړه ولري، يا
  \LR{$\sFmla{\True}{\lfalse}[\sigma]$} ولري، نو \LR{$\Gamma$} د صدق
  وړ نۀ دے۔
\end{prop}
\olpseganchor{OLP-0515-B012}{end}

\olpseganchor{OLP-0515-B013}{start}
\begin{proof}
  تل \LR{$\mSat/{M}{\lfalse}[f(\sigma)]$} وي؛ نو هغه \LR{$\Gamma$} چې
  \LR{$\sFmla{\True}{\lfalse}[\sigma]$} لري، د صدق وړ نۀ دے۔
\olpseganchor{OLP-0515-B013}{end}

\olpseganchor{OLP-0515-B014}{start}
  د \LR{$P(\Gamma)$} کوم تفسير~\LR{$f$} په مدل~\LR{$\mModel{M}$} کښې
  فرض کړئ۔ که
  \LR{$\mSat{M}{!A}[f(\sigma)]$} وي، نو د
  \olref[sem][rel]{prop:true-monotonic} له مخې، څرنګه چې
  \LR{$Rf(\sigma)f(\sigma.{*})$} دے،
  \LR{$\mSat{M}{!A}[f(\sigma.{*})]$} هم دے۔ نو
  \LR{$\mSat{M}{!A}[f(\sigma)]$} او
  \LR{$\mSat/{M}{!A}[f(\sigma.{*})]$} دواړه نۀ شو لرلے۔
\end{proof}
\olpseganchor{OLP-0515-B014}{end}

\olpseganchor{OLP-0515-B015}{start}
\begin{thm}[سموالے]
  \ollabel{thm:tableau-soundness}
  که د \LR{$\Gamma$} ټول لومړني فرضونه يو ګډ مخکښ ولري او
  تړلے تابلو ترې جوړ شي، نو \LR{$\Gamma$} د صدق وړ نۀ دے۔
\end{thm}
\olpseganchor{OLP-0515-B015}{end}

\olpseganchor{OLP-0515-B016}{start}
\begin{proof}
د تابلو څانګه د صدق وړ بولو هله او يوازې هله چې پر هغې
د نښه لرونکي فارمولونه سټ د صدق وړ وي؛ او تابلو د صدق
وړ بولو که لږ تر لږه يوه د صدق وړ څانګه ولري۔
\olpseganchor{OLP-0515-B016}{end}

\olpseganchor{OLP-0515-B017}{start}
دا ښيو: د استنتاج د يوې قاعدې په وسيله د صدق وړ تابلو
غځول تل د صدق وړ تابلو ورکوي۔ په دې سره قضيه ثابتېږي:
هر تړلے تابلو پر هغه تابلو د استنتاج د قاعدو په
پلې کولو جوړېږي چې يوازې د~\LR{$\Gamma$} فرضونه لري۔ نو که
\LR{$\Gamma$} د صدق وړ وے، د هغه هر تابلو به د صدق وړ وے۔
خو تړلے تابلو په ښکاره د صدق وړ نۀ دے، ځکه ټولې څانګې
يې تړلې دي او تړلې څانګې د صدق وړ نۀ دي۔ ټول فرضونه يو
ګډ مخکښ لري۔ په هره څانګه کښې له هماغه لومړني مخکښ څخه تر
هر بل راغلي مخکښ پورې ټول منځني مخکښونه هم په څانګه کښې
شته۔ دا خاصيت په پيل کښې پوره دے۔ کومه قاعده چې نوے مخکښ
راوړي، هغه د يوه موجود مخکښ په يوه غړي غځوي، نو خاصيت پاتې
کېږي۔
\olpseganchor{OLP-0515-B017}{end}

\olpseganchor{OLP-0515-B018}{start}
فرض کړئ د صدق وړ تابلو لرو، يعنې داسې تابلو چې
لږ تر لږه يوه د صدق وړ څانګه لري۔ د استنتاج قاعده يا څانګې ته
نښه لرونکي فارمولونه ورزياتوي يا څانګه په دوو وېشي۔ که د
تابلو د صدق وړ يوه څانګه د اړوندې قاعدې په پلې کولو نۀ
غځېږي، هغه په غځېدلي تابلو کښې هم د صدق وړ پاتې کېږي؛
نو غځېدلے تابلو د صدق وړ دے۔ ځکه يوازې هغه حالت کتل غواړي
چې قاعده پر د صدق وړ څانګې پلې کېږي۔
\olpseganchor{OLP-0515-B018}{end}

\olpseganchor{OLP-0515-B019}{start}
\LR{$\Gamma$} دې پر هماغې څانګې د نښه لرونکي فارمولونه سټ وي، او
\LR{$\sFmla{S}{!A}[\sigma] \in \Gamma$} دې هغه نښه لرونکے فارمول
وي چې قاعده پرې پلې کېږي۔ که قاعده څانګه نۀ وېشي، بايد وښيو
چې غځېدلې څانګه، يعنې \LR{$\Gamma$} له قاعده‌يي پايلو سره، لا د
صدق وړ ده۔ که قاعده څانګه وېشي، بايد وښيو چې د جوړو شويو دوو
څانګو لږ تر لږه يوه د صدق وړ ده۔
\begin{enumerate}
\item پر \LR{$\sFmla{\True}{\lnot !B}[\sigma] \in \Gamma$} د
  \LR{$\TRule{\True}{\lnot}$} په پلې کولو څانګه غځېږي۔ نو غځېدلې
  څانګه نښه لرونکي فارمولونه \LR{$\Gamma \cup
  \{\sFmla{\False}{!B}[\sigma.{*}]\}$} لري۔ فرض کړئ
  \LR{$\mSat{M}{\Gamma}[f]$} دے۔ په ځانګړي ډول
  \LR{$\mSat{M}{\lnot !B}[f(\sigma)]$} دے۔ نو د هر هغه \LR{$w$} لپاره
  چې \LR{$Rf(\sigma)w$} وي، \LR{$\mSat/{M}{!B}[w]$} دے؛
  \LR{$f(\sigma.{*})$} هم پکښې شامل دے۔ ځکه \LR{$\mModel{M}$} د~\LR{$f$}
  په نسبت په \LR{$\sFmla{\False}{!B}[\sigma.{*}]$} کښې صدق کوي۔
\item پر \LR{$\sFmla{\False}{\lnot !B}[\sigma] \in \Gamma$} د
  \LR{$\TRule{\False}{\lnot}$} په پلې کولو څانګه غځېږي۔ د نفي
  د ناصدق له مخې يوه لاسرسي وړ نړۍ شته چې نفي شوے غړے پکښې
  رښتينے وي۔ د نوي مخکښ تفسير هماغه نړۍ ته غځوو، لکه لاندې
  د نارښتيني شرط په حالت کښې۔
\item پر \LR{$\sFmla{\True}{!B \land !C}[\sigma] \in \Gamma$} د
  \LR{$\TRule{\True}{\land}$} په پلې کولو څانګه غځېږي؛ په څانګه
  کښې دوه نوي نښه لرونکي فارمولونه جوړېږي:
  \LR{$\sFmla{\True}{!B}[\sigma]$} او \LR{$\sFmla{\True}{!C}[\sigma]$}۔
  فرض کړئ \LR{$\mSat{M}{\Gamma}[f]$} دے؛ په ځانګړي ډول
  \LR{$\mSat{M}{!B \land !C}[f(\sigma)]$} دے۔ نو
  \LR{$\mSat{M}{!B}[f(\sigma)]$} او \LR{$\mSat{M}{!C}[f(\sigma)]$}
  دي۔ دا معنا لري چې \LR{$\mModel{M}$} د~\LR{$f$} په نسبت په
  \LR{$\sFmla{\True}{!B}[\sigma]$} او
  \LR{$\sFmla{\True}{!C}[\sigma]$} دواړو کښې صدق کوي۔
\item پر \LR{$\sFmla{\False}{!B \lor !C}[\sigma] \in \Gamma$} د
  \LR{$\TRule{\False}{\lor}$} په پلې کولو څانګه غځېږي۔ په
  هماغې نړۍ کښې د فصل دواړه غړي نارښتيني دي، نو هماغه تفسير
  دواړه زياتې کرښې پوره کوي۔
\item پر \LR{$\sFmla{\False}{!B \lif !C}[\sigma] \in \Gamma$} د
  \LR{$\TRule{\False}{\lif}$} په پلې کولو څانګه غځېږي؛ په څانګه
  کښې دوه نوي نښه لرونکي فارمولونه جوړېږي:
  \LR{$\sFmla{\True}{!B}[\sigma.n]$} او
  \LR{$\sFmla{\False}{!C}[\sigma.n]$}؛ \LR{$\sigma.n$} په څانګه کښې
  نوے مخکښ دے، يعنې \LR{$\sigma.n \notin P(\Gamma)$}۔
\olpseganchor{OLP-0515-B019}{end}
  
\olpseganchor{OLP-0515-B020}{start}
  \LR{$\Gamma$} د صدق وړ دے؛ نو يو \LR{$\mModel{M}$} او د \LR{$P(\Gamma)$}
  تفسير~\LR{$f$} داسې دي چې \LR{$\mSat{M}{\Gamma}[f]$}، او په ځانګړي
  ډول \LR{$\mSat/{M}{!B \lif !C}[f(\sigma)]$} وي۔ بايد وښيو چې
  \LR{$\Gamma \cup \{\sFmla{\True}{!B}[\sigma.n],\sFmla{\False}{!C}[\sigma.n]\}$}
  د صدق وړ دے۔ د دې لپاره د \LR{$P(\Gamma)
  \cup \{\sigma.n\}$} تفسير په دې ډول تعريفوو:
\olpseganchor{OLP-0515-B020}{end}
  
\olpseganchor{OLP-0515-B021}{start}
  \LR{$\mSat/{M}{!B \lif !C}[f(\sigma)]$} دے؛ نو يو \LR{$w \in W$}
  داسې شته چې \LR{$Rf(\sigma)w$}، \LR{$\mSat{M}{!B}[w]$} او
  \LR{$\mSat/{M}{!C}[w]$} وي۔ \LR{$f'$} دې د \LR{$f$} په شان وي، يوازې
  \LR{$f'(\sigma.n) = w$} دے۔ څرنګه چې \LR{$f'(\sigma) = f(\sigma)$}
  او \LR{$Rf(\sigma)w$} دي، \LR{$Rf'(\sigma)f'(\sigma.n)$} لرو؛ نو
  \LR{$f'$} د \LR{$P(\Gamma) \cup \{\sigma.n\}$} تفسير دے۔ د دې دپاره
  پام وکړئ چې د نوي مخکښ هېڅ غځونه لا په څانګه کښې نشته،
  ځکه بيا به نوے مخکښ هم د هغې غځونې منځنۍ برخه وه او
  په څانګه کښې به موجود و۔ د نوي مخکښ نور زاړۀ پلرني
  مخکښونه د هغه د سمدستي پلار هم پلرني دي۔ د زاړه تفسير
  له مخې د هغو نړۍو نه د پلار نړۍ ته لاسرسے شته؛ د تعدي
  له مخې نوې نړۍ ته هم لاسرسے شته۔ په ښکاره
  \LR{$\mSat{M}{!B}[f'(\sigma.n)]$} او
  \LR{$\mSat/{M}{!C}[f'(\sigma.n)]$} دي۔ د ټولو مخکښونو
  \LR{$\sigma' \in P(\Gamma)$} لپاره \LR{$f(\sigma') = f'(\sigma')$}
  دے؛ نو \LR{$\mSat{M}{\Gamma}[f']$}۔ ځکه \LR{$\mModel{M}, f'$} په
  \LR{$\Gamma \cup \{\sFmla{\True}{!B}[\sigma.n],\sFmla{\False}{!C}[\sigma.n]\}$}
  کښې صدق کوي۔
\end{enumerate}
اوس هغه ممکن استنتاجونه ګورو چې دوه څانګې جوړوي۔
\begin{enumerate}
\item پر \LR{$\sFmla{\False}{!B \land !C}[\sigma] \in \Gamma$} د
  \LR{$\TRule{\False}{\land}$} په پلې کولو څانګه غځېږي؛ دوه
  څانګې جوړېږي: کيڼه يې پر \LR{$\sFmla{\False}{!B}[\sigma]$} او
  ښي يې پر \LR{$\sFmla{\False}{!C}[\sigma]$} روانېږي۔ فرض کړئ
  \LR{$\mSat{M}{\Gamma}[f]$} دے؛ په ځانګړي ډول
  \LR{$\mSat/{M}{!B \land !C}[f(\sigma)]$} دے۔ نو
  \LR{$\mSat/{M}{!B}[f(\sigma)]$} يا \LR{$\mSat/{M}{!C}[f(\sigma)]$}
  دے۔ په لومړي حالت کښې \LR{$\mModel{M}, f$} په
  \LR{$\sFmla{\False}{!B}[\sigma]$} کښې صدق کوي، يعنې کيڼه څانګه
  د صدق وړ ده۔ په دويم کښې \LR{$\mModel{M}, f$} په
  \LR{$\sFmla{\False}{!C}[\sigma]$} کښې صدق کوي، يعنې ښي څانګه
  د صدق وړ ده۔
\item پر \LR{$\sFmla{\True}{!B \lor !C}[\sigma] \in \Gamma$} د
  \LR{$\TRule{\True}{\lor}$} په پلې کولو څانګه غځېږي۔ په هماغې
  نړۍ کښې د فصل لږ تر لږه يو غړے رښتينے دے، نو د هغه غړي
  څانګه د هماغه تفسير له مخې د صدق وړ ده۔
\item پر \LR{$\sFmla{\True}{!B \lif !C}[\sigma] \in \Gamma$} د
  \LR{$\TRule{\True}{\lif}$} په پلې کولو څانګه غځېږي۔ کارېدلے
  غځېدلے مخکښ د تفسير له مخې لاسرسي وړ نړۍ نوموي۔ هلته يا
  لومړے غړے نارښتينے وي يا پايله رښتينې وي، نو د دوو څانګو
  لږ تر لږه يوه د هماغه تفسير له مخې د صدق وړ ده۔
\end{enumerate}
\end{proof}
\olpseganchor{OLP-0515-B021}{end}

\olpseganchor{OLP-0515-B022}{start}
\begin{prob}
د \olref[int][tab][sou]{thm:tableau-soundness} د ثبوت دپاره
د نارښتينې نفي، نارښتيني فصل، رښتيني فصل او رښتيني شرط د
حالتونو ټولې نښه‌لرونکې کرښې او د تفسير شرطونه په تفصيل وليکئ۔
\end{prob}
\olpseganchor{OLP-0515-B022}{end}

\olpseganchor{OLP-0515-B023}{start}
\begin{cor}
\ollabel{cor:entailment-soundness}
که \LR{$\Gamma \Proves !A$} وي، نو \LR{$\Gamma \Entails !A$}۔
\end{cor}
\olpseganchor{OLP-0515-B023}{end}

\olpseganchor{OLP-0515-B024}{start}
\begin{proof}
  که \LR{$\Gamma \Proves !A$} وي، نو د ځينو \LR{$!B_1$}، \dots,
  \LR{$!B_n \in \Gamma$} دپاره
  \LR{$\Delta = \{\sFmla{\False}{!A}[1], \sFmla{\True}{!B_1}[1],
  \dots, \sFmla{\True}{!B_n}[1]\}$} تړلے تابلو لري۔
  غواړو وښيو چې \LR{$\Gamma \Entails !A$}۔ نفي يې فرض کړئ؛ نو
  د کوم \LR{$\mModel{M}$} او \LR{$w$} دپاره د \LR{$i=1$}، \dots,~\LR{$n$}
  لپاره \LR{$\mSat{M}{!B_i}[w]$} خو \LR{$\mSat/{M}{!A}[w]$} دے۔
  \LR{$f(1) = w$} دې وي؛ نو \LR{$f$} په~\LR{$\mModel{M}$} کښې د~\LR{$P(\Delta)$}
  تفسير دے، او \LR{$\mModel{M}$} د~\LR{$f$} په نسبت په~\LR{$\Delta$} کښې
  صدق کوي۔ خو د \olref{thm:tableau-soundness} له مخې
  \LR{$\Delta$} د صدق وړ نۀ دے، ځکه تړلے تابلو لري؛ تناقض
  دے۔ نو په پاے کښې بايد \LR{$\Gamma \Entails !A$} ولرو۔
\end{proof}
\olpseganchor{OLP-0515-B024}{end}

\olpseganchor{OLP-0515-B025}{start}
\begin{cor}
\ollabel{cor:weak-soundness}
که \LR{$\Proves !A$} وي، نو \LR{$!A$} په ټولو مدلونو کښې رښتيا دے۔
\end{cor}
\olpseganchor{OLP-0515-B025}{end}

\olpseganchor{OLP-0516-B004}{start}
\olpart{cnt}{د واقعيت خلاف شرطونه}
\olpseganchor{OLP-0516-B004}{end}

\olpseganchor{OLP-0517-B004}{start}
\olchapter{cnt}{int}{پېژندنه}
\olpseganchor{OLP-0517-B004}{end}

\olpseganchor{OLP-0518-B004}{start}
\olfileid{cnt}{int}{mat}
\olpseganchor{OLP-0518-B004}{end}

\olpseganchor{OLP-0518-B005}{start}
\olsection{مادي شرطي}
\olpseganchor{OLP-0518-B005}{end}

\olpseganchor{OLP-0518-B006}{start}
په انګريزۍ کښې د شرطي تر ټولو ساده بڼه د «که \dots نو
\dots» په جوړښت جمله ده؛ د \dots{} پر ځاے خپله جملې راځي،
لکه «که د کور خدمتګار دا کار کړے وي، نو باغوان بې‌ګناه دے»۔
د منطق په لومړنيو کورسونو کښې د \LR{$\lif$} نښلوونکي په کارولو د
شرطي جملې په نښو ليکل زده کوو۔ د \dots پر ځاے راغلې
برخې مثلاً په فارمولونه \LR{$!A$} او~\LR{$!B$} ښيو، او ټول شرطي
په \LR{$!A \lif !B$} ليکو۔
\olpseganchor{OLP-0518-B006}{end}

\olpseganchor{OLP-0518-B007}{start}
نښلوونکے \LR{$\lif$} \emph{د صدق تابع} دے، يعنې د \LR{$!A \lif !B$}
د رښتياوالي قيمت---\LR{$\True$} يا \LR{$\False$}---د \LR{$!A$} او~\LR{$!B$} د
رښتياوالي په قيمتونو ټاکل کېږي: \LR{$!A \lif !B$} رښتيا دے هله او
يوازې هله چې \LR{$!A$}~کاذب يا \LR{$!B$}~رښتيا وي؛ که داسې نۀ وي،
کاذب دے۔ د رښتياوالي د قيمتونو د ګومارنې~\LR{$\pAssign v$} په
نسبت، \LR{$\pSat{v}{!A \lif !B}$} تعريفوو هله او يوازې هله چې
\LR{$\pSat/{v}{!A}$} يا \LR{$\pSat{v}{!B}$} وي۔ د دې معنايي پوهنې
سره نښلوونکے \LR{$\lif$} \emph{مادي شرطي} بلل کېږي۔
\olpseganchor{OLP-0518-B007}{end}

\olpseganchor{OLP-0518-B008}{start}
له دې تعريف څخه څو لومړني منطقي حقيقتونه راوځي۔ تر هر څه
وړاندې، د مادي شرطي لپاره د استنتاج قضيه سمه ده۔ په هره
ارزونه کښې چې فرضونه رښتيا وي، که مقدم نارښتينے وي نو شرطي
رښتينے دے۔ که مقدم رښتينے وي، له استلزامه يې پايله هم رښتينې ده۔
\begin{align}
  \text{\RL{که }} \Gamma, !A \Entails !B \text{\RL{ وي، نو }} \Gamma \Entails !A \lif !B
\end{align}
هغه د صدق تابع دے: \LR{$!A \lif !B$} او \LR{$\lnot !A \lor !B$} هم‌ارزښته دي:
\begin{align}
  !A \lif !B & \Entails \lnot !A \lor !B\\
  \lnot !A \lor !B & \Entails !A \lif !B
  \intertext{\RL{د مادي شرطي پايله او د هغه د مقدم نفي دواړه دغه
    شرطي په معنايي ډول لازموي:}}
  !B & \Entails !A \lif !B\\
  \lnot !A & \Entails !A \lif !B
  \intertext{\RL{د مادي شرطي نفي د هغه د مقدم او د پايلې د نفي له
    عطف سره هم‌ارزښته ده: که \LR{$!A \lif !B$} کاذب وي،
    \LR{$!A \land \lnot !B$} رښتيا دے، او سرچپه هم همداسې ده:}}
  \lnot(!A \lif !B) & \Entails !A \land \lnot !B\\
  !A \land \lnot !B & \Entails \lnot(!A \lif !B)
  \intertext{\RL{مادي شرطي مودوس پوننس مني:}}
  !A, !A \lif !B & \Entails !B
  \intertext{\RL{د مادي شرطيو پايلې په عطف کښې يوځاے کېدے شي:}}
  !A \lif !B,  !A \lif !C & \Entails !A \lif (!B \land !C)
  \intertext{\RL{مقدم تل پياوړے کولے شو، يعنې شرطي \emph{يکنواخت} دے:}}
  !A \lif !B & \Entails (!A \land !C) \lif !B
  \intertext{\RL{مادي شرطي متعدي دے، يعنې د ځنځير قاعده معتبره ده:}}
  !A \lif !B, !B \lif !C & \Entails !A \lif !C
  \intertext{\RL{مادي شرطي له خپل نقيض عکس سره هم‌ارزښته دے:}}
  !A \lif !B & \Entails \lnot !B \lif \lnot !A\\
  \lnot !B \lif \lnot !A & \Entails !A \lif !B
\end{align}
\olpseganchor{OLP-0518-B008}{end}

\olpseganchor{OLP-0518-B009}{start}
په رياضيکي استدلال کښې دا ټول ګټور او بې‌ستونزې استنتاجونه
دي۔ خو په فلسفي او ژبپوهنې اثارو کښې په غيررياضيکي سياقونو
کښې د ورته استنتاجونو ډېرې ادعا شوې ضدبېلګې شته۔ دا وړانديز
کوي چې د انګريزۍ د «که \dots نو \dots» ويناوو په
نښه‌ليکنه کښې مادي شرطي~\LR{$\lif$} مناسب نښلوونکے نۀ دے---يا لږ
تر لږه تل مناسب نۀ دے۔
\olpseganchor{OLP-0518-B009}{end}

\olpseganchor{OLP-0519-B004}{start}
\olfileid{cnt}{int}{par}
\olpseganchor{OLP-0519-B004}{end}

\olpseganchor{OLP-0519-B005}{start}
\olsection{د مادي شرطي پاراډوکسونه}
\olpseganchor{OLP-0519-B005}{end}

\olpseganchor{OLP-0519-B006}{start}
د انګريزۍ د «که \dots نو \dots» ويناوو د نښه‌ليکنې لپاره د
\LR{$!A \lif !B$} پر کارولو له لومړنيو نيوکه کوونکو څخه سي. آي.
لويس و۔ لويس د وايټ هېډ او رسل په \emph{پرنسيپيا متماتيکا}
کښې د مادي شرطي پر کارولو نيوکه کوله؛ هغوے \LR{$\lif$} د
«لازموي» په توګه لوسته۔ لويس په سمه توګه اعتراض وکړ چې که
\LR{$\lif$} د «لازموي» معنا ولري، نو هره کاذبه قضيه~\LR{$p$} دا
لازموي چې \LR{$p$}،~\LR{$q$} لازموي، ځکه که~\LR{$p$} کاذب وي،
\LR{$p \lif (p \lif q)$} رښتيا دے؛ او هره رښتينې قضيه~\LR{$q$} دا
لازموي چې \LR{$p$}،~\LR{$q$} لازموي، ځکه که \LR{$q$}~رښتيا وي،
\LR{$q \lif (p \lif q)$} رښتيا دے۔ دويم ټول شرطي په هرې ارزونې
کښې رښتينے دے: که مقدم يې نارښتينے وي، شرطي رښتينے دے؛ او
که مقدم يې رښتينے وي، دنننۍ پايله يې رښتينې ده۔
\olpseganchor{OLP-0519-B006}{end}

\olpseganchor{OLP-0519-B007}{start}
منطقپوهان البته پوهېږي چې \emph{لازمول}، يعنې منطقي
استلزام، نښلوونکے نۀ دے، بلکې د فارمولونه يا ويناوو
ترمنځ اړيکه ده۔ نو د ګډوډۍ د مخنيوي لپاره بايد \LR{$\lif$} د
«لازموي» په توګه ونۀ لولو۔\footnote{رياضيپوهان او د
  کمپيوټرپوهې پوهان لا هم ډېر ځله «\LR{$\lif$}» د «لازموي» په
  توګه لولي؛ خو فيلسوفان يې «يوازې که» بولي، څو د لويس له
  ښودلې ګډوډۍ څخه ډډه وکړي۔} تر څو چې داسې يې نۀ لولو، د
لويس ځانګړې اندېښنه نۀ راپيدا کېږي: \LR{$p$}، \LR{$q$} نۀ «لازموي»،
که څه هم \LR{$p$} د يوې کاذبې انګريزي جملې نښه وګڼو۔ د دې معلومولو
لپاره چې آيا \LR{$p \Entails q$}، بايد \emph{ټولې}
ارزښت ټاکنې وګورو۔ که دا دوه بېل قضييوي متغيرونه وي،
داسې ارزونه شته چې لومړے رښتينے او دويم نارښتينے کوي۔ نو
\LR{$p \Entails/ q$} دے، ان که \LR{$p$} د داسې جملې نښه وي چې په
اتفاق سره نارښتينې وي۔
\olpseganchor{OLP-0519-B007}{end}

\olpseganchor{OLP-0519-B008}{start}
خو بيا هم د لويس په شان د «که \dots نو\dots» ويناوو کښې يو
څه عجيب شته:
\begin{quote}
که سپوږمۍ له شنې پنيرې جوړه وي، نو \LR{$2+2=4$}۔
\end{quote}
او په دغو استنتاجونو کښې هم:
\begin{quote}
  سپوږمۍ له شنې پنيرې جوړه نۀ ده۔ له دې امله که سپوږمۍ له شنې
  پنيرې جوړه وي، نو \LR{$2+2=4$}۔
\olpseganchor{OLP-0519-B008}{end}

\olpseganchor{OLP-0519-B009}{start}
  \LR{$2+2 = 4$}۔ له دې امله که سپوږمۍ له شنې پنيرې جوړه وي، نو
  \LR{$2+2=4$}۔
\end{quote}
خو که «که \dots نو \dots» يوازې \LR{$\lif$} وے، جمله به بې له
ستونزې رښتيا، او استنتاجونه به بې له ستونزې معتبر وو۔
\olpseganchor{OLP-0519-B009}{end}

\olpseganchor{OLP-0519-B010}{start}
بله بېلګه له تاوتولوژۍ \LR{$(!A \lif !B) \lor (!B \lif !A)$} سره
تړلې ده۔ دا وړانديز کوي چې که له ورځپاڼې څخه په ناټاکلي ډول
دوه خبري جملې~\LR{$S$} او \LR{$T$} واخلئ، نو جمله «که \LR{$S$}، نو \LR{$T$}؛
يا که \LR{$T$}، نو~\LR{$S$}» بايد رښتيا وي۔
\olpseganchor{OLP-0519-B010}{end}

\olpseganchor{OLP-0520-B004}{start}
\olfileid{cnt}{int}{str}
\olpseganchor{OLP-0520-B004}{end}

\olpseganchor{OLP-0520-B005}{start}
\olsection{کلک شرطي}
\olpseganchor{OLP-0520-B005}{end}

\olpseganchor{OLP-0520-B006}{start}
لويس \emph{کلک شرطي}~\LR{$\strictif$} معرفي کړ او استدلال يې
وکړ چې له لازمولو سره د همدغه شرطي مطابقت شته، د مادي شرطي
نۀ۔ د صدق د امکان او لزوم په مودالي منطق کښې
\LR{$!A \strictif !B$} د \LR{$\Box(!A \lif !B)$} په توګه تعريفېدے
شي۔ نو کلک شرطي (په يوه نړۍ کښې) رښتيا دے هله او يوازې هله
چې اړوند مادي شرطي لازم وي۔
\olpseganchor{OLP-0520-B006}{end}

\olpseganchor{OLP-0520-B007}{start}
کلک شرطي د مادي شرطي د پاراډوکسونو په وړاندې څه حالت لري؟
د کاذب مقدم لرونکے کلک شرطي او د رښتينې پايلې لرونکے کلک
شرطي کېدے شي رښتيا وي او کېدے شي کاذب وي۔ پر دې سربېره
\LR{$(!A \strictif !B) \lor (!B \strictif !A)$} معتبر نۀ دے۔
کلک شرطي \LR{$!A \strictif !B$} له \LR{$\lnot !A \lor !B$} سره هم
هم‌ارزښته نۀ دے۔ په اوسنۍ نړۍ کښې د دواړو غړو صدق د کلک
شرطي صدق نۀ ټاکي، ځکه د نورو لاسرسي وړ نړۍو ارزښتونه هم
پرې اغېز لري۔ نو د صدق تابع نۀ دے۔
\olpseganchor{OLP-0520-B007}{end}

\olpseganchor{OLP-0520-B008}{start}
دا استلزامونه او ناستلزامونه د \Log{S5} په چوکاټونو کښې
ګورو، چې لاسرسے پکې د هم‌ارزۍ اړيکه ده۔ د انعکاسيت له
مخې لازم مادي شرطي په اوسنۍ نړۍ کښې هم رښتينے وي۔ د کلک
شرطي نفي بيا په کومه لاسرسي وړ نړۍ کښې د مادي شرطي
نارښتيا غواړي۔
\begin{align}
  !A \strictif !B & \Entails \lnot !A \lor !B
  \text{\RL{ خو:}}\\
  \lnot !A \lor !B & \Entails/ !A \strictif !B\\
  !B & \Entails/ !A \strictif !B\\
  \lnot !A & \Entails/ !A \strictif !B\\
  \lnot(!A \strictif !B) & \Entails/ !A \land \lnot !B
  \text{\RL{ خو:}}\\
  !A \land \lnot !B & \Entails \lnot(!A \strictif !B)
  \intertext{\RL{خو کلک شرطي لا هم مودوس پوننس مني:}}
  !A, !A \strictif !B & \Entails !B
  \intertext{\RL{د کلکو شرطيو پايلې په عطف کښې يوځاے کېدے شي:}}
  !A \strictif !B,  !A \strictif !C & \Entails !A \strictif (!B \land !C)
  \intertext{\RL{د کلک شرطي لپاره د مقدم پياوړي کول سم دي:}}
  !A \strictif !B & \Entails (!A \land !C) \strictif !B
  \intertext{\RL{کلک شرطي هم متعدي دے:}}
  !A \strictif !B, !B \strictif !C & \Entails !A \strictif !C
  \intertext{\RL{په پاے کښې، کلک شرطي له خپل نقيض عکس سره
    هم‌ارزښته دے:}}
  !A \strictif !B & \Entails \lnot !B \strictif \lnot !A\\
  \lnot !B \strictif \lnot !A & \Entails !A \strictif !B
\end{align}
\olpseganchor{OLP-0520-B008}{end}

\olpseganchor{OLP-0520-B009}{start}
\begin{prob}
  د هغو استلزامي اړيکو لپاره د \Log{S5} ضدبېلګې ورکړئ چې د
  کلک شرطي لپاره نۀ سمېږي، يعنې د دې لپاره:
  \begin{enumerate}
  \item \LR{$\lnot p \Entails/ \Box(p \lif q)$}
  \item \LR{$q \Entails/ \Box(p \lif q)$}
  \item \LR{$\lnot\Box(p \lif q) \Entails/ p \land \lnot q$}
  \item \LR{$\Entails/ \Box(p \lif q) \lor \Box(q \lif p)$}
  \end{enumerate}
\end{prob}
\olpseganchor{OLP-0520-B009}{end}

\olpseganchor{OLP-0520-B010}{start}
\begin{prob}
  د لاندې بيانونو د \Log{S5} ثبوتونو په ورکولو وښيئ چې
  معتبرې استلزامي اړيکې د کلک شرطي لپاره سمې دي:
  \begin{enumerate}
  \item  \LR{$\Box(!A \lif !B) \Entails \lnot !A \lor !B$}
  \item \LR{$!A \land \lnot !B \Entails \lnot\Box(!A \lif !B)$}
  \item \LR{$!A, \Box(!A \lif !B) \Entails !B$}
  \item \LR{$\Box(!A \lif !B), \Box(!A \lif !C) \Entails \Box(!A \lif (!B
    \land !C))$}
  \item \LR{$\Box(!A \lif !B) \Entails \Box((!A \land !C) \lif !B)$}
  \item \LR{$\Box(!A \lif !B), \Box(!B \lif !C) \Entails \Box(!A
    \lif !C)$}
  \item \LR{$\Box(!A \lif !B) \Entails \Box(\lnot !B \lif \lnot !A)$}
  \item \LR{$\Box(\lnot !B \lif \lnot !A) \Entails \Box(!A \lif !B)$}
  \end{enumerate}
  \end{prob}
\olpseganchor{OLP-0520-B010}{end}

\olpseganchor{OLP-0520-B011}{start}
خو کلک شرطي لا هم خپل «پاراډوکسونه» لري۔ لکه څرنګه چې د
کاذب مقدم يا رښتينې پايلې لرونکے مادي شرطي رښتيا وي، د
\emph{په لازمي ډول} کاذب مقدم يا په لازمي ډول رښتينې پايلې
لرونکے کلک شرطي هم رښتيا وي۔ پر دې سربېره هر رښتينے کلک
شرطي په لازمي ډول رښتيا او هر کاذب کلک شرطي په لازمي ډول
کاذب دے۔ په همدغو چوکاټونو کښې، له يوې نړۍ څخه لاسرسي وړ
نړۍ ټولې د هم‌ارزۍ په يوه ټولګي کښې دي۔ د ټولګي هره نړۍ
هماغو ټولو نړۍو ته لاسرسے لري، نو د کلک شرطي صدق يا
ناصدق په ټول ټولګي کښې يو شان وي۔
\begin{align}
  \Box \lnot !A & \Entails !A \strictif !B\\
  \Box !B & \Entails !A \strictif !B\\
  !A \strictif !B& \Entails \Box(!A \strictif !B)\\
  \lnot(!A \strictif !B) & \Entails \Box\lnot(!A \strictif !B)
\end{align}
که \LR{$\strictif$} د «لازموي» په توګه وګڼئ، دا ستونزې نۀ دي۔
منطقي استلزامي اړيکې په هر حال رياضيکي حقيقتونه دي؛ نو
اتفاقي نۀ شي کېدے۔ خو که غواړئ \LR{$\strictif$} داسې منطقي
نښلوونکے وکاروئ چې «که \dots نو \dots» ويناوې رانغاړي، دا
دعوې، په تېره وروستۍ دوه، ستونزې راولاړوي۔ ځکه خامخا داسې
«که \dots نو \dots» ويناوې شته چې په اتفاق رښتيا يا په
اتفاق کاذبې دي---په حقيقت کښې عموماً نۀ لازمې دي او نۀ
ناممکنې۔
\olpseganchor{OLP-0520-B011}{end}

\olpseganchor{OLP-0520-B012}{start}
\begin{prob}
  په \Log{S5} کښې د دې ثبوتونه ورکړئ:
  \begin{enumerate}
    \item  \LR{$\Box \lnot !A \Entails !A \strictif !B$}
    \item \LR{$!A \strictif !B \Entails \Box(!A \strictif !B)$}
    \item \LR{$\lnot(!A \strictif !B)  \Entails \Box\lnot(!A \strictif !B)$}
  \end{enumerate}
  د دې لپاره د \LR{$\strictif$} تعريف وکاروئ۔
\end{prob}
\olpseganchor{OLP-0520-B012}{end}

\olpseganchor{OLP-0521-B004}{start}
\olfileid{cnt}{int}{cnt}
\olpseganchor{OLP-0521-B004}{end}

\olpseganchor{OLP-0521-B005}{start}
\olsection{د واقعيت خلاف شرطونه}
\olpseganchor{OLP-0521-B005}{end}

\olpseganchor{OLP-0521-B006}{start}
په انګرېزۍ کښې د «که \dots نو \dots» جوړښتونو يوه ډېره عامه او
مهمه بڼه د \emph{to be} («وي/وے») فعل د تېر التزامي صورت په کارولو جوړېږي:
«که داسې وے چې \dots نو داسې به وے چې \dots»۔ څرنګه چې د داسې
شرطي مقدم عموماً کاذب، يعنې له واقعيت خلاف وي، دې ته
\emph{د واقعيت خلاف شرطونه} وايي (او څرنګه چې د \emph{to be}
فعل التزامي صورت کاروي، \emph{التزامي شرطونه} هم ورته وايي)۔
دا له \emph{خبري} شرطونو بېل دي، چې بڼه يې داسې
ده: «که داسې ده چې \dots نو داسې ده چې \dots»۔ د واقعيت خلاف او
خبري شرطونه په صدق-شرطونو کښې توپير لري۔ د اېډمز د بېلګې دوه شرطونه وګورئ۔
\begin{explain}
  په خبري بڼه کښې اېډمز د کېنېډي وژنه يوه شوې پېښه ګڼي۔
  که اوسوالډ قاتل نۀ وي، د هماغې پېښې عامل بل څوک دے۔
\olpseganchor{OLP-0521-B006}{end}

\olpseganchor{OLP-0521-B007}{start}
  په التزامي بڼه کښې د اېډمز پوښتنه دا ده چې د پېښو په بديل
  بهير کښې به څه شوي وو۔ د اوسوالډ له خوا د کېنېډي وژل له
  هغه بهير څخه وباسئ؛ د شرطي دعوه دا ده چې کېنېډي به بيا هم
  د بل کس له لاسه وژل شوے وے۔
\end{explain}
لومړۍ خبري ده، دويمه د واقعيت خلاف شرطي ده۔ لومړۍ په څرګنده
رښتينې ده: مونږ پوهېږو چې ولسمشر جان اېف کېنېډي \emph{چا} وژلے
دے، او که هغه کس (د وارېن د راپور خلاف) لي هاروي اوسوالډ نۀ و، نو
بل چا کېنېډي وژلے دے۔ دويمه بيا بېل څه وايي۔ دا دعوه کوي چې که
اوسوالډ کېنېډي نۀ وے وژلے، يعنې که په ډالاس کښې د ډزو د پېښې
مخنيوے شوے وے يا ناکامه شوې وے، نو تاريخ به تر هغې وروسته داسې
مخ ته تلے وے چې بله د وژنې هڅه به بريالۍ شوې وے۔ د داسې
پايلې د تصور يوه لاره هغه فرضيه ده چې ځواکمنو قوتونو يې د
مرګ د يقيني کولو سازش کړے و، لکه د سازش د نظريو ډېر پلويان
چې باور لري۔ د شرطي د صدق دپاره دا ځانګړې فرضيه اړينه نۀ
ده؛ خبره دا ده چې په ورته بديل حالت کښې به څه شوي وو۔
\olpseganchor{OLP-0521-B007}{end}

\olpseganchor{OLP-0521-B008}{start}
لا بحث روان دے چې آيا \emph{خبري} شرطي د مادي شرطي په وسيله سم
بيانېږي، په تېره بيا دا چې آيا د مادي شرطي پاراډوکسونه په داسې
طريقه «تشريح» کېدے شي چې له دې سره برابره وي چې مادي شرطي د
انګرېزۍ د خبري شرطونو صدق-شرطونه ورکوي۔ خو په دې کښې اختلاف نشت
چې د واقعيت خلاف شرطونه د مادي شرطونو په وسيله سم نۀ شي
نښه‌ليکل کېدے۔ دا ځکه څرګنده ده چې که څه هم د واقعيت خلاف شرطونو
مقدمونه عموماً نارښتيني وي، ځينې داسې شرطونه له نارښتيني
مقدم سره هم نارښتيني وي۔ د بېلګې په توګه، که د وارېن پر راپور
باور لرئ او په نژدې بديلو حالاتو کښې د لومړۍ پېښې له
مخنيوي سره بله داسې پېښه هم نۀ رامنځته کېده، نو د اېډمز د
واقعيت خلاف شرطي يې يوه بېلګه ده۔
\olpseganchor{OLP-0521-B008}{end}

\olpseganchor{OLP-0521-B009}{start}
د واقعيت خلاف شرطونه په علّي استدلال کښې مهم رول لري: د هغو د
کارولو يوه څرګنده بېلګه د علّي اړيکو بيانول دي۔ د بېلګې په توګه،
د اورلګيت د لرګي وهل د هغه د بلېدو سبب کېږي، او دا په دې وينا
بيانولے شئ: «که دا د اورلګيت لرګے وهل شوے وے، نو بل شوے به وے»۔
مادي شرطونه، او په عموم کښې خبري شرطونه، د دې د بيانولو لپاره
نۀ شي کارېدے: «د اورلګيت لرګے وهل کېږي \LR{$\lif$} د اورلګيت لرګے
بلېږي» هله رښتينې ده چې لرګے هېڅکله ونۀ وهل شي، پرته له دې چې
که وهل شوے وے نو څه به شوي وو۔ تر دې هم بده دا ده چې «د
اورلګيت لرګے وهل کېږي \LR{$\lif$} د اورلګيت لرګے د ګلونو په ګېډۍ
بدلېږي» هم هله رښتينې ده چې لرګے هېڅکله ونۀ وهل شي، خو که لرګے
وهل شوے وے نو په يقين سره به د ګلونو په ګېډۍ نۀ وے بدل شوے۔
\olpseganchor{OLP-0521-B009}{end}

\olpseganchor{OLP-0521-B010}{start}
لا بحث روان دے چې د واقعيت خلاف شرطونو سم منطق په دقيقه توګه
څه دے۔ سټالنکر او لويس د واقعيت خلاف شرطونو يوه اغېزمنه شننه
وړاندې کړه۔ په دواړو شننو کښې د واقعيت خلاف شرطي «که داسې
وے چې~\LR{$S$} نو داسې به وے چې~\LR{$T$}» د ارزونې لپاره ګورو چې \LR{$T$} په
اړوندو بديلو حالاتو («ممکنو نړۍو») کښې رښتينے دے که نۀ۔
سټالنکر د مقدم د صدق لپاره يوه تر ټولو نژدې نړۍ ټاکي، يعنې
داسې نړۍ چې~\LR{$S$} پکښې رښتينے وي۔ د لويس په شننه کښې، که تر
ټولو نژدې نړۍ وي، پايله بايد په ټولو يو شان نژدې نړۍو کښې
رښتينې وي چې مقدم پکښې رښتينے دے۔ د هغه کره تعريف
هغه حالت هم رانغاړي چې د مقدم تر ټولو نژدې نړۍ پکې نشته۔
دې ته «هستي‌اړونده» شننه وايي، ځکه چې د ممکنو نړۍو هستي‌پوهنې
ته رجوع کوي۔ نورې شننې شرطي احتمالونه يا د باور د بياکتنې
نظريې کاروي۔ د واقعيت خلاف شرطونو ډېر بېلابېل منطقونه وړانديز
شوي دي۔ آن د واقعيت خلاف شرطونو يو يوازينے د لويس--سټالنکر
منطق هم نشت: که څه هم سټالنکر او لويس د تر ټولو نژدې ممکنو
نړۍو په حواله په ورته طريقو شننې وړاندې کړې، د هغو فرضونه د
معتبرو استنتاجونو بېلابېلې پايلې ورکوي۔
\olpseganchor{OLP-0521-B010}{end}

\olpseganchor{OLP-0522-B004}{start}
\olchapter{cnt}{min}{د لږ تر لږه بدلون معنٰی پوهنه}
\olpseganchor{OLP-0522-B004}{end}

\olpseganchor{OLP-0523-B004}{start}
\olfileid{cnt}{min}{int}
\olpseganchor{OLP-0523-B004}{end}

\olpseganchor{OLP-0523-B005}{start}
\olsection{پېژندنه}
\olpseganchor{OLP-0523-B005}{end}

\olpseganchor{OLP-0523-B006}{start}
سټالنکر او لويس د واقعيت خلاف شرطونو، لکه «که د اورلګيت لرګے
وهل شوے وے، نو بل شوے به وے»، شننې وړاندې کړې۔ د هغو شننې د
داسې جملو د صدق-شرطونو د سم پوهېدو وړانديزونه وو۔ د دواړو
وړانديزونو تر شا نظر دا دے: د دې د ارزولو لپاره چې د واقعيت
خلاف شرطي رښتينې ده که نۀ، بايد هغه ممکنې نړۍ په پام کښې
ونيسو چې د مقدم د رښتيني کولو لپاره د واقعي نړۍ له حالت څخه
تر ټولو لږې بېلې دي۔ کله چې دغسې تر ټولو نژدې نړۍ موجودې
وي، د کروي سيمو په شننه کښې شرطي هله رښتينې وي چې په دغو ټولو
نړۍو کښې پايله رښتينې وي۔ د بېلګې په توګه، فرض کړئ
چې زما په لاس کښې د اورلګيت يو لرګے او د اورلګيتونو يوه کتابچه
ده۔ په واقعي نړۍ کښې زه يوازې ورته ګورم او سوچ کوم چې که
لرګے ووهم نو څه به وشي۔ له واقعي نړۍ څخه هغه لږ تر لږه بدلون
چې پکښې لرګے وهم، دا دے چې د عمل پرېکړه وکړم او لرګے ووهم۔
بدلون په دې معنا لږ تر لږه دے چې نور هېڅ نۀ بدلېږي: زه ورسره
په هوا کښې ټوپ نۀ وهم، د لرګي وهل زما وېښتۀ هم نۀ بلوي، زه
ناڅاپه په ګوتو کښې ټول ځواک نۀ بايلم، په هماغه وخت د اوبو د
لوبې د سوپر سوکر ټوپک په کمين کښې پر ما اوبۀ نۀ اچول کېږي، او داسې نور۔
په دغه بديل امکان کښې د اورلګيت لرګے بلېږي۔ نو دا رښتينې ده
چې که زه لرګے ووهم، نو بل به شي۔
\olpseganchor{OLP-0523-B006}{end}

\olpseganchor{OLP-0523-B007}{start}
د دې لومړني تصور شننه د واقعيت خلاف شرطونو د منطقونو له رسمي
معنٰی پوهنې سره يو ځاے کېدے شي۔ لويس د واقعيت خلاف شرطي لپاره
«\LR{$\boxright$}» نښه معرفي کړه، او سټالنکر «\LR{$>$}» نښه کاروله۔
مونږ به~\LR{$\cif$} کاروو او د ويناوو په منطق کښې به يې د دوه‌ځايي
نښلوونکي په توګه ورزياتوو۔ نو د \LR{$!A \lif
!B$} د بڼې پر فارمولونه سربېره، د~\LR{$!A \cif !B$} د بڼې
فارمولونه هم لرو۔ رسمي معنٰی پوهنه، د مودالي منطق د اړيکو د
معنٰی پوهنې په څېر، پر داسې مدلونو ولاړه ده چې فارمولونه پکښې
په نړۍو کښې ارزول کېږي؛ او د \LR{$\mSat{M}{!A \cif !B}[w]$} د
تعريفوونکي تحقق شرط د ځينو (نورو) نړۍو~\LR{$w'$} لپاره د
\LR{$\mSat{M}{!A}[w']$} او \LR{$\mSat{M}{!B}[w']$} له مخې ورکول کېږي۔
کومې \LR{$w'$}؟ د لومړني تصور له مخې، هغه يو يا څو چې~\LR{$w$} ته تر
ټولو نژدې وي او~\LR{$\mSat{M}{!A}[w']$} پکښې برقرار وي۔ دا غواړي
چې په مدل کښې د «نژدې‌والي» اړيکه هم شامله وي۔ راتلونکے
تعريف هغه حالت هم رانغاړي چې تر ټولو نژدې نړۍ موجودې نۀ وي۔
\olpseganchor{OLP-0523-B007}{end}

\olpseganchor{OLP-0523-B008}{start}
لويس د واقعيت خلاف حالتونو د انځوريز ښودلو يوه ښوونکې طريقه
معرفي کړه۔ هره ممکنه نړۍ د يو داسې سټ په مرکز کښې ده چې يو
په بل کښې دننه کروي سيمې لري، او نورې نړۍ رانغاړي---مونږ دغه
کروي سيمې د هم‌مرکزو دايرو په توګه رسموو۔ د دوو کروي سيمو
ترمنځ نړۍ د مرکز نړۍ ته يو شان نژدې دي؛ په دننۍ کروي سيمه
کښې نړۍ نژدې، او په چاپېره کروي سيمه کښې نړۍ ليرې دي۔
\begin{center}
\begin{tikzpicture}[scale=.7]\small
  \spheresystem{5}
  \draw (0,0) node {\LR{$w$}};
  \propositionintersect{0}{4}{40}{3.5}
  {\tikzset{proposition/.append={smooth,tension=1.4}}
  \proposition[shift={(1.6,-1)}]{90}{1}{30}{4}}
  \path (1.5,3) node[below] {\LR{$\formula{B}$}};
  \path (3,0) node {\LR{$\formula{A}$}};
\end{tikzpicture}
\end{center}
په دې انځور کښې تر ټولو نژدې \LR{$!A$}-نړۍ هغه نړۍ \LR{$w'$} دي
چې~\LR{$!A$} پکښې تحقق لري او د مرکز نړۍ~\LR{$w$} چاپېره په تر ټولو
کوچنۍ مقدم‌منونکې کروي سيمه کښې پرتې دي (خړه سيمه)۔
په دې حالت کښې \LR{$!A \cif !B$}
په~\LR{$w$} کښې هله تحقق لري چې \LR{$!B$} په ټولو تر ټولو نژدې
\LR{$!A$}-نړۍو کښې رښتينے وي۔
\olpseganchor{OLP-0523-B008}{end}

\olpseganchor{OLP-0524-B004}{start}
\olfileid{con}{min}{sph}
\olpseganchor{OLP-0524-B004}{end}

\olpseganchor{OLP-0524-B005}{start}
\olsection{د کروي سيمو مدلونه}
\olpseganchor{OLP-0524-B005}{end}

\olpseganchor{OLP-0524-B006}{start}
د واقعيت خلاف شرطونو د رسمي معنٰی پوهنې د ورکولو يوه طريقه دا
ده چې د لويس غيررسمي شننه په رياضيکي جوړښت بدله کړو۔ بيا د
يوې نړۍ~\LR{$w$} چاپېره کروي سيمې د نړۍو سټونه دي۔ څرنګه چې
کروي سيمې يو په بل کښې دننه دي، د~\LR{$w$} چاپېره د نړۍو سټونه
بايد د فرعي سټ د اړيکې له مخې خطي ترتيب ولري۔
\olpseganchor{OLP-0524-B006}{end}

\olpseganchor{OLP-0524-B007}{start}
\begin{defn}
  \emph{د کروي سيمو مدل} يو درې‌توب~\LR{$\mModel{M} = \tuple{W, O, V}$}
  دے، چې \LR{$W$} د نړۍو يو ناتش سټ، \LR{$V\colon \PVar \to \Pow{W}$}
  يوه ارزښت ټاکنه، او \LR{$O\colon W \to \Pow{\Pow{W}}$} داسې تابع ده
  چې هرې نړۍ~\LR{$w$} ته \emph{د کروي سيمو يو نظام}~\LR{$O_w$} ټاکي۔
  د هر \LR{$w$} لپاره، \LR{$O_w$} د نړۍو د سټونو يو سټ دے او بايد دا شرطونه پوره کړي:
  \begin{enumerate}
  \item \LR{$O_w$} پر~\LR{$w$} \emph{مرکوز} دے: \LR{$\{w\} \in O_w$}۔
  \item \LR{$O_w$} \emph{يو په بل کښې دننه} دے: هر کله چې \LR{$S_1$}، \LR{$S_2 \in O_w$}، نو \LR{$S_1
    \subseteq S_2$} يا \LR{$S_2 \subseteq S_1$}؛ يعنې \LR{$O_w$} د~\LR{$\subseteq$}
    له مخې خطي ترتيب لري۔
  \item \LR{$O_w$} د ناتشو کورنيو د اتحاد پر وړاندې تړلے دے۔
  \item \LR{$O_w$} د ناتشو کورنيو د تقاطع پر وړاندې تړلے دے۔
  \end{enumerate}
\end{defn}
\olpseganchor{OLP-0524-B007}{end}

\olpseganchor{OLP-0524-B008}{start}
د \LR{$O_w$} تر شا تصور دا دے چې د \LR{$w$} «چاپېره» نړۍ له~\LR{$w$} څخه
د واټن له مخې په پوړونو کښې وېشل شوې دي۔ تر ټولو دننۍ ناتشه
کروي سيمه يوازې \LR{$w$}، يعنې سټ~\LR{$\{w\}$}، ده: \LR{$w$} له~\LR{$w$} سره
تر هغو نورو نړۍو نژدې دے چې د سيمو په نظام کښې دي۔
د تشې کروي سيمې شتون د نړۍو په دغه ترتيب اغېز نۀ لري۔
که \LR{$S \subsetneq S'$}، نو په \LR{$S' \setminus S$}
کښې نړۍ له \LR{$w$} څخه تر هغو نړۍو ليرې دي چې په~\LR{$S$} کښې دي:
\LR{$S' \setminus S$} د \LR{$S$} او د~\LR{$S'$} د بهرنۍ سطحې ترمنځ «پوړ»
دے۔ دغه پوړ له دننۍ سيمې بهر او له بهرنۍ پولې دننه دے۔
په تېره بيا، بايد کروي سيمې داسې وګڼو چې د خپلې
بهرنۍ سطحې دننه ټولې نړۍ رانغاړي؛ يوازې جلا پوړونه نۀ دي۔
\olpseganchor{OLP-0524-B008}{end}

\olpseganchor{OLP-0524-B009}{start}
\begin{figure}
\begin{center}
\begin{tikzpicture}[layerwidth=1.5,scale=.6]\tiny
  \spheresystem{3}
  \propositionintersect{0}{3}{60}{6}
  \path (0,0) node[world] {\LR{$w$}};
  \spherepos{-30}{2}{node[world] {\LR{$w_2$}}}
  \spherepos{220}{2}{node[world] {\LR{$w_3$}}}
  \spherepos{90}{2}{node[world] {\LR{$w_1$}}}
  \spherepos[shift={(.1,0)}]{10}{3}{node[world] {\LR{$w_5$}}}
  \spherepos[shift={(.1,0)}]{-10}{3}{node[world] {\LR{$w_6$}}}
  \spherepos{225}{3}{node[world] {\LR{$w_4$}}}
  \spherepos{20}{4}{node[world] {\LR{$w_7$}}}
  \path (5,0) node[left] {\LR{$p$}};
\end{tikzpicture}
\caption{د کروي سيمو د يو مدل انځور}
\ollabel{fig:sphere-model}
\end{center}
\end{figure}
\olpseganchor{OLP-0524-B009}{end}

\olpseganchor{OLP-0524-B010}{start}
په \olref{fig:sphere-model} کښې انځور له هغه د کروي سيمو له
مدل سره برابر دے چې \LR{$W = \{w, w_1, \dots, w_7\}$} او
\LR{$V(p) = \{w_5, w_6,
w_7\}$} لري۔ تر ټولو دننۍ کروي سيمه \LR{$S_1 = \{w\}$} ده۔
له \LR{$w$} سره تر ټولو نژدې نړۍ \LR{$w_1, w_2, w_3$} دي، نو ورپسې
لويه کروي سيمه \LR{$S_2 = \{w, w_1,
w_2, w_3\}$} ده۔ نورې ليرې نړۍ \LR{$w_4$}، \LR{$w_5$} او \LR{$w_6$} دي،
نو تر ټولو بهرنۍ کروي سيمه \LR{$S_3 = \{w, w_1, \dots, w_6\}$}
ده۔ د \LR{$w$} چاپېره د کروي سيمو نظام \LR{$O_w = \{S_1, S_2, S_3\}$}
دے۔ نړۍ~\LR{$w_7$} د~\LR{$w$} چاپېره په هېڅ کروي سيمه کښې نۀ ده۔
هغه تر ټولو نژدې نړۍ چې \LR{$p$}~پکښې رښتينے دے \LR{$w_5$} او~\LR{$w_6$}
دي، نو تر ټولو کوچنۍ \LR{$p$}-منونکې کروي سيمه~\LR{$S_3$} ده۔
\olpseganchor{OLP-0524-B010}{end}

\olpseganchor{OLP-0524-B011}{start}
د کروي سيمو په مدل~\LR{$\mModel M$} کښې په نړۍ~\LR{$w$} کښې د
فورمول~\LR{$!A$} تحقق، \LR{$\mSat{M}{!A}[w]$}، د تعريفولو لپاره د
مودالي فارمولونه تعريف داسې غځوو چې د \LR{$!B \cif !C$} لپاره
يو شرط ورشامل شي:
\olpseganchor{OLP-0524-B011}{end}

\olpseganchor{OLP-0524-B012}{start}
\begin{defn}
  \LR{$\mSat{M}{!B \cif !C}[w]$} هله او يوازې هله چې له دې دواړو يو وي:
  \begin{enumerate}
  \item\ollabel{sphere-vac} د ټولو \LR{$u \in \bigcup O_w$} لپاره، \LR{$\mSat/{M}{!B}[u]$}، يا
  \item\ollabel{sphere-nonvac} د يوې \LR{$S \in O_w$} لپاره،
    \begin{enumerate}
      \item د يوه \LR{$u \in
        S$} لپاره \LR{$\mSat{M}{!B}[u]$}، او
      \item د ټولو \LR{$v \in S$} لپاره، يا
        \LR{$\mSat/{M}{!B}[v]$} يا \LR{$\mSat{M}{!C}[v]$}۔
    \end{enumerate}
  \end{enumerate}
\end{defn}
\olpseganchor{OLP-0524-B012}{end}

\olpseganchor{OLP-0524-B013}{start}
د دې تعريف له مخې، \LR{$\mSat{M}{!B \cif !C}[w]$} هله او يوازې
هله چې يا مقدم~\LR{$!B$} د~\LR{$w$} چاپېره په کروي سيمو کښې هر چېرته
کاذب وي، يا يوه کروي سيمه~\LR{$S$} وي چې \LR{$!B$} پکښې رښتينے وي او
مادي شرطي \LR{$!B \lif !C$} په دغه «\LR{$!B$}-منونکې» کروي سيمه کښې
په ټولو نړۍو کښې رښتينې وي۔ پام وکړئ چې په تعريف کښې مو دا
شرط نۀ دے ايښے چې \LR{$S$} دې \emph{تر ټولو دننۍ} \LR{$!B$}-منونکې
کروي سيمه وي، که څه هم د لومړني تصور له تشريح داسې تمه کېدے
شي۔ خو که په~\olref{sphere-nonvac} کښې شرط د يوې کروي
سيمې~\LR{$S$} لپاره پوره وي، نو د هغې هرې کروي سيمې لپاره هم
پوره دے چې د~\LR{$S$} دننه وي او مقدم پکښې په کومه نړۍ کښې
رښتينے وي۔ په داسې دننۍ سيمه کښې ټولې نړۍ لا هم د مادي
شرطي شرط پوره کوي۔ نو دا خبره پر تر ټولو کوچنۍ
مقدم‌منونکې سيمې هم صدق کوي، که موجوده وي۔
\olpseganchor{OLP-0524-B013}{end}

\olpseganchor{OLP-0524-B014}{start}
دا هم په پام کښې ونيسئ چې د کروي سيمو د مدلونو تعريف دا
شرط نۀ ږدي چې يوه تر ټولو دننۍ \LR{$!B$}-منونکې کروي سيمه
\emph{موجوده وي}: کېدے شي د \LR{$!B$}-منونکو کروي سيمو يو
نامتناهي لړ \LR{$S_1 \supsetneq S_2 \supsetneq \dots \supsetneq
\{w\}$} ولرو۔ که هره مقدم‌منونکې کروي سيمه د لړۍ لږ تر
لږه يوه سيمه رانغاړي، نو تر ټولو دننۍ \LR{$!B$}-منونکې کروي
سيمه نشته: که يوه تر ټولو کوچنۍ وای، د لړۍ يوه سيمه به
پکښې پرته وه او د هغې لا ورپسې سيمه به ترې کوچنۍ وه۔
په همدغه شرط، \LR{$\mSat{M}{!B \cif !C}[w]$} هله او يوازې
هله چې \LR{$!B \lif !C$} د يوه~\LR{$i$} لپاره په ټولو کروي سيمو
\LR{$S_i$}، \LR{$S_{i+1}$}، \dots کښې برقرار وي۔ که د لړۍ له بهر
يوه مقدم‌منونکې سيمه د صدق شاهد وي، د لړۍ يوه سيمه د
شاهد دننه راځي او هماغه شرط پوره کوي۔ په عام حالت کښې
له پورته تعريف سره سم هره مقدم‌منونکې کروي سيمه د صدق
شاهد کېدے شي چې مادي شرطي يې په ټولو نړۍو کښې رښتينے وي۔
\olpseganchor{OLP-0524-B014}{end}

\olpseganchor{OLP-0525-B004}{start}
\olfileid{con}{min}{tf}
\olpseganchor{OLP-0525-B004}{end}

\olpseganchor{OLP-0525-B005}{start}
\olsection{د واقعيت خلاف شرطونو صدق او کذب}
\olpseganchor{OLP-0525-B005}{end}

\olpseganchor{OLP-0525-B006}{start}
که تر ټولو کوچنۍ مقدم‌منونکې کروي سيمه موجوده وي، د
واقعيت خلاف شرطي \LR{$!A \cif !B$} (په ناتشه توګه) هله رښتينې
ده چې تر ټولو نژدې \LR{$!A$}-نړۍ ټولې \LR{$!B$}-نړۍ وي، لکه څنګه چې په
\olref{fig:true} کښې ښودل شوې ده۔
\begin{figure}
\begin{center}
\begin{tikzpicture}[scale=.7]\small
  \spheresystem{5}
  \draw (0,0) node {\LR{$w$}};
  \propositionintersect{0}{4}{40}{3.5}
  {\tikzset{proposition/.append={smooth,tension=1.4}}
  \proposition[shift={(1.6,-1)}]{90}{1}{30}{4}}
  \path (1.5,3) node[below] {\LR{$\formula{B}$}};
  \path (3,0) node {\LR{$\formula{A}$}};
\end{tikzpicture}
\caption{په ناتشه توګه رښتينې د واقعيت خلاف شرطي}
\ollabel{fig:true}
\end{center}
\end{figure}
د واقعيت خلاف شرطي په \LR{$w$} کښې هله هم رښتينې ده چې
د~\LR{$w$} چاپېره د کروي سيمو نظام هېڅ \LR{$!A$}-منونکې کروي سيمه
ونۀ لري۔ په دې حالت کښې هغه \emph{په تشه توګه} رښتينې ده
(\olref{fig:vacuous} وګورئ)۔
\begin{figure}
\begin{center}
\begin{tikzpicture}[scale=.7]\small
  \spheresystem{5}
  \draw (0,0) node {\LR{$w$}};
  \proposition{0}{6.5}{40}{3.5}
  {\tikzset{proposition/.append={smooth,tension=1.4}}
  \proposition[shift={(1.2,-1)}]{90}{1}{30}{4}}
  \path (1.5,3) node[below] {\LR{$\formula{B}$}};
  \spherepos{0}{7}{node {\LR{$\formula{A}$}}}
\end{tikzpicture}
\caption{په تشه توګه رښتينې د واقعيت خلاف شرطي}
\ollabel{fig:vacuous}
\end{center}
\end{figure}
\olpseganchor{OLP-0525-B006}{end}

\olpseganchor{OLP-0525-B007}{start}
کله چې تر ټولو کوچنۍ مقدم‌منونکې کروي سيمه موجوده وي،
شرطي په دوو طريقو نارښتينې کېدے شي۔ لومړۍ دا ده چې په تر
ټولو نژدې \LR{$!A$}-نړۍو کښې ځينې \LR{$!B$}-نړۍ وي او ځينې نۀ وي۔ په دې
حالت کښې، \LR{$!A \cif
\lnot !B$} هم کاذبه ده (\olref{fig:false} وګورئ)۔
\begin{figure}
\begin{center}
\begin{tikzpicture}[scale=.7]\small
  \spheresystem{5}
  \draw (0,0) node {\LR{$w$}};
  \propositionintersect{0}{4}{40}{3.5}
  {\tikzset{proposition/.append={smooth,tension=1.4}}
  \proposition[shift={(1.6,0)}]{90}{1}{30}{3}}
  \path (1.5,3) node[below] {\LR{$\formula{B}$}};
  \path (3,0) node {\LR{$\formula{A}$}};
\end{tikzpicture}
\caption{کاذبه د واقعيت خلاف شرطي، کاذب مقابل}
\ollabel{fig:false}
\end{center}
\end{figure}
که تر ټولو نژدې \LR{$!A$}-نړۍ له \LR{$!B$}-نړۍو سره هېڅ ګډه نړۍ
ونۀ لري، نو \LR{$!A \cif !B$} کاذبه ده۔ خو په دې حالت کښې ټولې
تر ټولو نژدې \LR{$!A$}-نړۍ \LR{$\lnot !B$}-نړۍ دي، نو
\LR{$!A \cif \lnot !B$} رښتينې ده (\olref{fig:false-opposite} وګورئ)۔
\begin{figure}
\begin{center}
\begin{tikzpicture}[scale=.7]\small
  \spheresystem{5}
  \draw (0,0) node {\LR{$w$}};
  \propositionintersect{0}{4}{40}{3.5}
  {\tikzset{proposition/.append={smooth,tension=1.4}}
  \proposition[shift={(-1,0)}]{90}{1}{30}{3}}
  \path (-1,3) node[below] {\LR{$\formula{B}$}};
  \path (3,0) node {\LR{$\formula{A}$}};
\end{tikzpicture}
\caption{کاذبه د واقعيت خلاف شرطي، رښتينے مقابل}
\ollabel{fig:false-opposite}
\end{center}
\end{figure}
\olpseganchor{OLP-0525-B007}{end}

\olpseganchor{OLP-0525-B008}{start}
د هم‌ارزۍ په چوکاټونو کښې د کلک شرطي ارزښت د لاسرسي
د يوه ټولګي په دننه کښې يو شان وي۔ د واقعيت خلاف شرطونو
ارزښت د ارزونې له نړۍ سره بدلېدے شي۔
په \olref{fig:contingent} کښې د کروي سيمو مدل وګورئ۔
له~\LR{$u$} سره تر ټولو نژدې \LR{$!A$}-نړۍ ټولې \LR{$!B$}-نړۍ دي، نو
\LR{$\mSat{M}{!A \cif
  !B}[u]$}۔ خو له~\LR{$v$} سره ځينې تر ټولو نژدې \LR{$!A$}-نړۍ
\LR{$!B$}-نړۍ نۀ دي، نو \LR{$\mSat/{M}{!A \cif !B}[v]$}۔
\begin{figure}
\begin{center}
\begin{tikzpicture}[scale=.6]\small
  \spheresystem{7}
  \spheresystem[shift={(0:3.1)},dashed]{4}
  \propositionintersect{45}{4}{40}{4.5}
  \begin{scope}
    \clip \propositionplot{45}{4}{40}{4.5} ;
    \spherelayer[shift={(0:3.1)}]{3}
  \end{scope}
  \proposition[shift={(1.4,-.5)}]{90}{1}{35}{5}
  \path (0,0) node {\LR{$u$}};
  \path (0:3.1) node {\LR{$v$}};
  \spherepos{35}{9}{node {\LR{$\formula{A}$}}}
  \spherepos{80}{9}{node {\LR{$\formula{B}$}}}
\end{tikzpicture}
\end{center}
\caption{اتفاقي د واقعيت خلاف شرطي}
\ollabel{fig:contingent}
\end{figure}
\olpseganchor{OLP-0525-B008}{end}

\olpseganchor{OLP-0526-B004}{start}
\olfileid{cnt}{min}{agg}
\olpseganchor{OLP-0526-B004}{end}

\olpseganchor{OLP-0526-B005}{start}
\olsection{د مقدم پياوړي کول}
\olpseganchor{OLP-0526-B005}{end}

\olpseganchor{OLP-0526-B006}{start}
«د مقدم پياوړي کول» دې استنتاج ته وايي:
\LR{$!A \lif !C \Entails (!A \land !B) \lif !C$}۔ دا د مادي
شرطي لپاره معتبر دے، خو د واقعيت خلاف شرطونو لپاره نامعتبر
دے۔ فرض کړئ چې دا رښتينې ده چې که زه د اورلګيت دا لرګے ووهم،
نو بل به شي۔ (يعنې په لرګي يا د اورلګيت د کتابچې په سطحه کښې
څه خرابوالے نشت، زه به لرګے مات نۀ کړم، او داسې نور۔) خو دا
رښتينې نۀ ده چې که زه دا لرګے په بهرنۍ خلا کښې ووهم، نو بل
به شي۔ د دې بېلګې په دويم حالت کښې د بلېدو عادي شرطونه
پوره نۀ دي۔ په لومړۍ جمله کښې يوازې د لرګي وهل فرض شوي دي؛
په دويمه کښې ځاے هم ټاکل شوے دے۔ نو دا
استنتاج نامعتبر دے:
\begin{quote}
  که د اورلګيت لرګے وهل شوے وے، نو بل شوے به وے۔
\olpseganchor{OLP-0526-B006}{end}

\olpseganchor{OLP-0526-B007}{start}
  نو که د اورلګيت لرګے په بهرنۍ خلا کښې وهل شوے وے، نو بل شوے به وے۔
\end{quote}
\olpseganchor{OLP-0526-B007}{end}

\olpseganchor{OLP-0526-B008}{start}
د شرطونو د لويس--سټالنکر شننه دا تشريح کوي: هغه تر ټولو
نژدې نړۍ چې زه پکښې لرګے وهم او دا په بهرنۍ خلا کښې کوم،
له واقعي نړۍ څخه د هغې تر ټولو نژدې نړۍ په پرتله ډېره ليرې
ده چې زه پکښې لرګے وهم۔ نو که څه هم په وروستۍ نړۍ کښې د
لرګي بلېدل رښتيني دي، په لومړۍ کښې نۀ دي۔ او بايد همداسې وي۔
\olpseganchor{OLP-0526-B008}{end}

\olpseganchor{OLP-0526-B009}{start}
\begin{ex}
  د کروي سيمو معنٰی پوهنه استنتاج نامعتبر ثابتوي، يعنې \LR{$p
  \cif r \Entails/ (p \land q) \cif r$} لرو۔ مدل \LR{$\mModel{M}
  = \tuple{W, O, V}$} وګورئ چې \LR{$W = \{w, w_1, w_2\}$}، \LR{$O_w = \{\{w\}, \{w,
  w_1\}, \{w, w_1, w_2\}\}$}، \LR{$V(p) = \{w_1, w_2\}$}، \LR{$V(q) = \{w_2\}$}
  او \LR{$V(r) = \{w_1\}$} لري۔ د دوو نورو نړۍو چاپېره يوازې د
  هماغې نړۍ يوغړې سيمه ټاکو۔ يوه \LR{$p$}-منونکې کروي سيمه \LR{$S = \{w,
  w_1\}$} شته او \LR{$p \lif r$} پکښې په ټولو نړۍو کښې رښتينې ده،
  نو \LR{$\mSat{M}{p
    \cif r}[w]$}۔ يوه \LR{$(p \land q)$}-منونکې کروي سيمه \LR{$S' =
  \{w, w_1, w_2\}$} ده۔ د عطف د رښتياوالي يوازينۍ نړۍ په
  همدې بهرنۍ سيمه کښې ده، نو همدا يوه سيمه عطف مني۔ هلته
  \LR{$\mSat/{M}{(p \land q) \lif r}[w_2]$}، نو
  \LR{$\mSat/{M}{(p \land q) \cif r}[w]$}
  (\olref{fig:antecedent-strengthening} وګورئ)۔
\olpseganchor{OLP-0526-B009}{end}

\olpseganchor{OLP-0526-B010}{start}
\begin{figure}
\begin{center}
\begin{tikzpicture}[layerwidth=1.5,scale=.6]\tiny
  \spheresystem{3}
  \propositionintersect{-10}{3}{30}{5}
  \begin{scope}
    \clip plot[smooth,tension=1] coordinates
          {(0,4.5) (30:.95) (4.5,-3.5)} -- (4.5,4.5) ;
    \spherefill{2}
  \end{scope}
  \draw[proposition] plot[dashed,smooth,tension=1] coordinates
          {(0,4.5) (30:.95) (4.5,-3.5)} ;
  \proposition{30}{2}{30}{5}
  \path (0,0) node[world] {\LR{$w$}};
  \spherepos{30}{2}{node[world] {\LR{$w_1$}}}
  \spherepos{-10}{3}{node[world] {\LR{$w_2$}}}
  \spherepos{-10}{4}{node {\LR{$q$}}}
  \spherepos{30}{4}{node {\LR{$r$}}}
  \draw (1,4.5) node[below] {\LR{$p$}};
\end{tikzpicture}
\caption{د مقدم د پياوړي کولو ضدبېلګه}
\ollabel{fig:antecedent-strengthening}
\end{center}
\end{figure}
\end{ex}
\olpseganchor{OLP-0526-B010}{end}

\olpseganchor{OLP-0527-B004}{start}
\olfileid{cnt}{min}{tra}
\olpseganchor{OLP-0527-B004}{end}

\olpseganchor{OLP-0527-B005}{start}
\olsection{انتقاليت}
\olpseganchor{OLP-0527-B005}{end}

\olpseganchor{OLP-0527-B006}{start}
د مادي شرطي لپاره د زنځير قاعده برقرار ده: \LR{$!A \lif !B, !B
\lif !C \Entails !A \lif !C$}۔ په بله وينا، مادي شرطي انتقالي
ده۔ د واقعيت خلاف شرطونو لپاره د سټالنکر په بېلګه کښې دوه مقدمې او يوه پايله لرو۔
\begin{explain}
  لومړۍ مقدمه د هوور د زېږېدو ځاے بدلوي۔ په دې فرضي حالت
  کښې د شوروي اتحاد په چاپېريال کښې د هغه کمونست کېدل منو۔
\olpseganchor{OLP-0527-B006}{end}

\olpseganchor{OLP-0527-B007}{start}
  دويمه مقدمه له هوور څخه په خپل امريکايي چاپېريال کښې
  کمونست جوړوي۔ په دې جلا فرضي حالت کښې له خپل هېواد سره د
  هغه خيانت منو۔
\olpseganchor{OLP-0527-B007}{end}

\olpseganchor{OLP-0527-B008}{start}
  د دواړو مقدماتو له يوځاے کولو به دا نتيجه راووځي چې د
  هوور په روس کښې زېږېدل هم له خپل هېواد سره د هغه خيانت
  لازموي۔
\end{explain}
که هوور (په هماغه وخت چې په واقعيت کښې زېږېدلے و) د
متحده ايالاتو پر ځاے په روس کښې زېږېدلے وے، نو په شوروي
اتحاد کښې به لوے شوے وے او کمونست به شوے وے (دا فرض کړو)۔
د دې فرض له مخې لومړے مقدم رښتينے دے۔ په بېلګه کښې
دويم مقدم هم، چې په جلا توګه وکتل شي، رښتينے ګڼو۔ خو
د همدې فرضونو له مخې پايله نارښتينې ده: په دغه بديل حالت کښې
هوور که په شوروي اتحاد کښې زېږېدلے وے، نو کلک کمونست به وے
او (له خپل هېواد سره) غدار به نۀ وے۔ د سټالنکر--لويس شننه
د صدق د دغو ارزښتونو همدا وېش تاييدوي۔ زمونږ نړۍ ته هغه تر ټولو
نژدې ممکنه نړۍ چې يوازې د هوور د زېږون ځاے پکښې بدلېږي، هغه
ده چې هوور پکښې د شوروي اتحاد ښه وګړے لوېږي۔ دا هغه تر
ټولو نژدې ممکنه نړۍ ده چې د لومړي مقدم او پايلې مقدم پکښې
رښتينے دے؛ په هغې نړۍ کښې هوور د کمونست ګوند وفادار غړے
دے، نو غدار نۀ دے۔ خو د دويم مقدم د ارزولو لپاره بايد بله
نړۍ وګورو: تر ټولو نژدې نړۍ چې هوور پکښې کمونست دے، هغه
ده چې په متحده ايالاتو کښې زېږېدلے، بيا يې لورے بدل کړے او
په دې توګه غدار شوے دے۔\footnote{البته، د بېلګې د قوت د
  پوهېدو لپاره بايد ځينې مابعدالطبيعي او سياسي فرضونه ومنو؛
  د بېلګې په توګه، دا چې ممکنه وه هوور له روسي مور و پلار
  زېږېدلے وے، او دا سياسي فرض چې د نولس سوه پنځوسمو کلونو
  په متحده ايالاتو کښې کمونست کېدل خپل هېواد ته غداري وه۔}
\olpseganchor{OLP-0527-B008}{end}

\olpseganchor{OLP-0527-B009}{start}
\begin{prob}
  د واقعيت خلاف شرطونو د انتقاليت د ناکامۍ لپاره يوه
  قانع کوونکې، له لومړني تصور سره برابره بېلګه ومومئ۔
\end{prob}
\olpseganchor{OLP-0527-B009}{end}

\olpseganchor{OLP-0527-B010}{start}
\begin{ex}\ollabel{ex:trans-counterex}
  د کروي سيمو معنٰی پوهنه استنتاج نامعتبر ثابتوي، يعنې \LR{$p
  \cif q, q \cif r \Entails/ p \cif r$} لرو۔ مدل \LR{$\mModel{M}
  = \tuple{W, O, V}$} وګورئ چې \LR{$W = \{w, w_1, w_2\}$}، \LR{$O_w = \{\{w\}, \{w,
  w_1\}, \{w, w_1, w_2\}\}$}، \LR{$V(p) = \{w_2\}$}، \LR{$V(q) = \{w_1, w_2\}$}
  او \LR{$V(r) = \{w_1\}$} لري۔ د پاتې دوو نړۍو د سيمو نظام
  يوازې د هماغې نړۍ له يوغړي سټه جوړوو۔ يوه \LR{$p$}-منونکې کروي سيمه \LR{$S = \{w, w_1,
  w_2\}$} شته او \LR{$p \lif q$} پکښې په ټولو نړۍو کښې رښتينې ده،
  نو \LR{$\mSat{M}{p
    \cif q}[w]$}۔ يوه \LR{$q$}-منونکې کروي سيمه \LR{$S' = \{w,
  w_1\}$} هم شته، او \LR{$q \lif r$} يې په ټولو نړۍو کښې رښتينے
  دے، نو \LR{$\mSat{M}{q \cif r}[w]$}۔ خو يوازې
  بهرنۍ کروي سيمه د لومړي مقدم د رښتياوالي نړۍ لري۔ دغه
  \LR{$p$}-منونکې کروي سيمه \LR{$\{w,
  w_1, w_2\}$} يوه نړۍ، يعنې~\LR{$w_2$}، لري چې
  \LR{$\mSat/{M}{p \lif
    r}[w_2]$} پکښې برقرار دے۔ نو له لومړي مقدم څخه تر
  وروستۍ پايلې د واقعيت خلاف شرطي نارښتينې ده۔
\end{ex}
\olpseganchor{OLP-0527-B010}{end}

\olpseganchor{OLP-0527-B011}{start}
\begin{prob}
په \olref[cnt][min][tra]{ex:trans-counterex} کښې د ضدبېلګې
سره برابر د کروي سيمو انځور رسم کړئ۔
\end{prob}
\olpseganchor{OLP-0527-B011}{end}

\olpseganchor{OLP-0527-B012}{start}
\begin{prob}
  په \olref[cnt][min][tra]{ex:trans-counterex} کښې نړۍ \LR{$w_2$}
  هغه ده چې هوور پکښې په روس کښې زېږېدلے، کمونست دے او
  غدار نۀ دے؛ او \LR{$w_1$} هغه نړۍ ده چې هوور پکښې په متحده
  ايالاتو کښې زېږېدلے، کمونست دے او غدار دے۔ په دې مدل کښې
  \LR{$w_1$} له~\LR{$w$} سره تر \LR{$w_2$} نژدې ده۔ آيا دا لازمي ده؟ داسې
  ضدبېلګه ورکولے شئ چې دا فرض نۀ کړي چې په روس کښې د هوور
  زېږېدل د هغه تر کمونست کېدو ډېر ليرې امکان دے؟ لارښوونه:
  د دواړو رښتينو شرطونو د شاهدې سيمې سره پرتله کړئ۔ که د
  دويم شرط په شاهده سيمه کښې د لومړي شرط مقدم رښتينے وي،
  د دغو دوو سيمو په تقاطع کښې به وروستۍ پايله څۀ ارزښت ولري؟
\end{prob}
\olpseganchor{OLP-0527-B012}{end}

\olpseganchor{OLP-0528-B004}{start}
\olfileid{cnt}{min}{cpo}
\olpseganchor{OLP-0528-B004}{end}

\olpseganchor{OLP-0528-B005}{start}
\olsection{نقيض عکس}
\olpseganchor{OLP-0528-B005}{end}

\olpseganchor{OLP-0528-B006}{start}
مادي او کلک شرطونه له خپلو نقيض عکسونو سره هم‌ارزښته دي۔
د واقعيت خلاف شرطونه داسې نۀ دي۔ د کراټزر په بېلګه کښې دوه فرضي حالتونه دي۔
\begin{explain}
  په لومړي شرطي کښې د 1832 د مرګ پېښه بدلېږي۔ د ګويټه د
  ژوند عادي موده دومره ده چې د اوسني وخت تر رارسېدو وړاندې
  به يې ژوند پاے ته رسېدلے وے۔
\olpseganchor{OLP-0528-B006}{end}

\olpseganchor{OLP-0528-B007}{start}
  نقيض عکس د ګويټه اوسنے ژوند فرضوي او بيا د هغه مرګ په
  1832 کښې غواړي۔ دغه دوه حالتونه يو ځاے نۀ شي کېدے۔
\end{explain}
لومړۍ جمله رښتينې ده: انسانان سلګونه کلونه ژوند نۀ کوي۔
دويمه په څرګنده کاذبه ده: که ګويټه اوس مړ نۀ وے، نو لا به
ژوندے وے، نو په 1832 کښې مړ شوے نۀ شو کېدے۔
\olpseganchor{OLP-0528-B007}{end}

\olpseganchor{OLP-0528-B008}{start}
\begin{ex}\ollabel{ex:contraposition-counterex}
  د کروي سيمو معنٰی پوهنه نقيض عکس نامعتبر ثابتوي، يعنې \LR{$p
  \cif q \Entails/ \lnot q \cif \lnot p$} لرو۔ \LR{$p$} د دې په معنا
  وګڼئ چې «ګويټه په 1832 کښې مړ نۀ شو» او \LR{$q$} د دې په معنا
  چې «ګويټه اوس مړ دے»۔ دا په مدل \LR{$\mModel{M_1} = \tuple{W, O, V}$}
  کښې بيانولے شو چې \LR{$W =
  \{w, w_1, w_2\}$}، \LR{$O_w = \{\{w\}, \{w, w_1\}, \{w, w_1, w_2\}\}$}،
  \LR{$V(p) = \{w_1, w_2\}$} او \LR{$V(q) = \{w, w_1\}$} لري۔ په نورو
  دوو نړۍو کښې د سيمو نظام يوازې د هماغې نړۍ له يوغړې سيمې
  جوړېږي۔ نو \LR{$w$} واقعي نړۍ ده چې ګويټه پکښې په 1832 کښې مړ شوے
  او لا مړ دے؛ \LR{$w_1$} هغه (نژدې) نړۍ ده چې ګويټه پکښې، د
  بېلګې په توګه، په 1833 کښې مړ شوے او لا مړ دے؛ او \LR{$w_2$}
  هغه (ليرې) نړۍ ده چې ګويټه پکښې لا ژوندے دے۔
\begin{figure}
\begin{center}
\begin{tikzpicture}[scale=.6,layerwidth=1.5]\tiny
  \spheresystem{3}
  \begin{scope}
    \clip \propositionplot{0}{.8}{50}{4.5} -- (4.5,-4.5) -- (-4.5,-4.5) -- (-4.5,4.5) -- (4.5,4.5);
    \spherelayer{3}
  \end{scope}
  \propositionintersect{0}{2}{110}{5.5}
  \proposition{0}{.8}{50}{4.5}
  \path (0,0) node[world] {\LR{$w$}};
  \spherepos{0}{2}{node[world] {\LR{$w_1$}}}
  \spherepos{-50}{3}{node[world] {\LR{$w_2$}}}
  \spherepos{28}{3}{node {\LR{$q$}}}
  \spherepos{42}{3}{node {\LR{$\lnot q$}}}
  \spherepos{-55}{4}{node {\LR{$p$}}}
  \spherepos{-67}{4}{node {\LR{$\lnot p$}}}
\end{tikzpicture}
\caption{د نقيض عکس ضدبېلګه}
\ollabel{fig:contraposition}
\end{center}
\end{figure}
يوه \LR{$p$}-منونکې کروي سيمه \LR{$S =
\{w, w_1\}$} شته او \LR{$p \lif q$} پکښې په ټولو نړۍو کښې رښتينې
ده، نو \LR{$\mSat{M_1}{p
  \cif q}[w]$}۔ خو يوازې بهرنۍ سيمه د پايلې نفي مني۔
دغه \LR{$\lnot q$}-منونکې
کروي سيمه \LR{$\{w, w_1,
w_2\}$} يوه نړۍ، يعنې~\LR{$w_2$}، لري چې \LR{$q$} پکښې کاذب او \LR{$p$}
رښتينے دے، نو \LR{$\mSat/{M_1}{\lnot q \lif \lnot p}[w_2]$}۔
په دې يوازينۍ منونکې سيمه کښې د مادي شرطي د ناکامۍ له امله
د واقعيت خلاف نقيض عکس هم نارښتينے دے۔
\end{ex}
\olpseganchor{OLP-0528-B008}{end}

\olpseganchor{OLP-0529-B004}{start}
\olpart{sth}{سټ تيوري}
\olpseganchor{OLP-0529-B004}{end}

\olpseganchor{OLP-0530-B004}{start}
\olchapter{sth}{story}{تکراري تصور}
\olpseganchor{OLP-0530-B004}{end}

\olpseganchor{OLP-0531-B003}{start}
\olfileid{sth}{story}{extensionality}
\olsection{د غړو له مخې برابري}
\olpseganchor{OLP-0531-B003}{end}

\olpseganchor{OLP-0531-B004}{start}
تر ټولو لومړنۍ خبره دا ده چې سټونه د خپلو غړي له مخې
له يو بل څخه بېلېږي۔ په لا کره توګه:
\olpseganchor{OLP-0531-B004}{end}

\olpseganchor{OLP-0531-B005}{start}
\begin{axiom}[د غړو له مخې برابري]
که سټونه \LR{$A$} او \LR{$B$} يو شان غړي ولري، نو \LR{$A$} او \LR{$B$}
هماغه يو سټ دے۔
\[
  \lforall[A][\lforall[B][(\lforall[x][(x \in A \liff x \in B)] \lif
  \eq[A][B])]]
\]
په دې بڼه کښې متغيرونه پر سټونو ګرځي۔ که د بحث په ساحه
کښې داسې شيان هم وي چې سټونه نۀ دي، اصل يوازې پر هغو دوو
شيانو پلې کېږي چې دواړه سټونه وي۔
\end{axiom}
\olpseganchor{OLP-0531-B005}{end}

\olpseganchor{OLP-0531-B006}{start}
دا مو د \olref[sfr][][]{part} په ټوله برخه کښې فرض کړې وه۔
خو بيا ويل يې ګټه لري۔ د غړو له مخې د برابرۍ اصل دا بنسټيز
تصور بيانوي چې يو سټ د خپلو غړي له مخې ټاکل کېږي۔
(نو سټونه له \emph{مفهومونو} سره پرتله کولے شو، چې عين هماغه
شيان ښايي د ډېرو بېلو مفهومونو لاندې راشي۔)
\olpseganchor{OLP-0531-B006}{end}

\olpseganchor{OLP-0531-B007}{start}
ولې دا اصل ومنو؟ د غړو له مخې برابري د ټولګې د هويت
معيار دے۔ په دې معنا \emph{سټ تيوري} د هغو ټولګو تيوري
ده چې هويت يې په خپلو غړو پورې تړلے دے۔ د دغه اصل له
پرېښودو سره د ټولګې د هويت دپاره بل معيار پکار وي۔
\olpseganchor{OLP-0531-B007}{end}

\olpseganchor{OLP-0531-B008}{start}
خو په \olref[sth][][]{part} کښې اصلي ننګونه د هغو اصلونو
ټاکل دي چې راته ووايي \emph{کوم سټونه وجود لري}۔ او څرګندېږي
چې د دې پوښتنې يوازې هغه ځواب چې په رښتيا «څرګند» برېښي،
په ثبوت سره ناسم دے۔
\olpseganchor{OLP-0531-B008}{end}

\olpseganchor{OLP-0532-B003}{start}
\olfileid{sth}{story}{rus}
\olsection{د رسل تناقض (بيا)}
\olpseganchor{OLP-0532-B003}{end}

\olpseganchor{OLP-0532-B004}{start}
په \olref[sfr][][]{part} کښې مو له ساده سټ تيورۍ سره کار
کړے و۔ خو د يو \emph{ډېر} ساده تصور له مخې، سټونه يوازې
د پريديکاتونو د پوره کوونکو شيانو ټولګې دي۔ له دې ډېر
ساده تصور څخه لاندې اصل راځي:
\olpseganchor{OLP-0532-B004}{end}

\olpseganchor{OLP-0532-B005}{start}
\begin{defish}
  \emph{ساده ادراک۔} د هر فارمول \LR{$\phi$} لپاره \LR{$\Setabs{x}{\phi(x)}$} وجود لري۔
\end{defish}
\olpseganchor{OLP-0532-B005}{end}

\olpseganchor{OLP-0532-B006}{start}
که څه هم دا اصل زړه راکښونکے دے، په ثبوت سره ناسازګاره دے۔
دا مو په \olref[sfr][set][rus]{sec} کښې ليدلي وو، خو نتيجه
دومره مهمه او دومره نېغه ده چې دلته يې بشپړ دليل بيا راوړو۔
\olpseganchor{OLP-0532-B006}{end}

\olpseganchor{OLP-0532-B007}{start}
\begin{thm}[د رسل تناقض]
سټ \LR{$R = \Setabs{x}{x \notin x}$} نشت۔
\end{thm}
\olpseganchor{OLP-0532-B007}{end}

\olpseganchor{OLP-0532-B008}{start}
\begin{proof}
که \LR{$R = \Setabs{x}{x \notin x}$} وجود ولري، نو
\LR{$R \in R$} هله او يوازې هله چې \LR{$R \notin R$} وي۔
که سټ خپل غړے وي، د تعريف له مخې خپل غړے نۀ دے؛ نو خپل
غړے نۀ دے۔ د شرط بل لور بيا وايي چې خپل غړے دے۔ دا تناقض
د خارجِ منځني پر اصل تکيه نۀ کوي۔
\end{proof}
\olpseganchor{OLP-0532-B008}{end}

\olpseganchor{OLP-0532-B009}{start}
رسل دا نتيجه د 1901 په جون کښې وموندله۔ (خو تناقض يې په
کټ مټ هغې بڼې کښې نۀ و وړاندې کړے چې اوس مو وړاندې کړ؛ ځکه
هغه د فرېګه سټ تيوري څېړله، چې په \emph{Grundgesetze} کښې
بيان شوې ده۔ دې ته به په \olref[blv]{sec} کښې بيا راوګرځو۔)
رسل د 1902 د جون په 16مه فرېګه ته وليکل او د فرېګه د نظام
ناسازګاري يې ورته تشريح کړه۔ د خطونو او يو څه شاليد لپاره
\mbox{\LR{\latinfont \citet[pp.~124--8]{Heijenoort1967}}} وګورئ۔
\olpseganchor{OLP-0532-B009}{end}

\olpseganchor{OLP-0532-B010}{start}
ټينګار پکار دے چې دا د دوو کرښو ثبوت د \emph{خالص منطق}
يوه نتيجه ده۔ دا منو چې په نښه \LR{$\Setabs{x}{x \notin x}$} کښې
مو په پټه يو (غير منطقي؟)\ اصل، د غړو له مخې برابري، کارولے
دے؛ ځکه \LR{$\Setabs{x}{\phi(x)}$} د \LR{$\phi$} شرط پوره کوونکو شيانو د \emph{هماغه يوازيني}
سټ لپاره دے (د غړو له مخې د برابرۍ په وجه)، که وجود ولري۔ خو
د غړو له مخې د برابرۍ ان لږه اشاره هم لرې کولے شو، که نتيجه
داسې بيان کړو: \emph{داسې سټ نشت چې غړي يې کټ مټ هغه سټونه
وي چې خپل ځان د غړي په توګه نۀ لري}۔ او دا له سټونو سره ډېره
ځانګړې اړيکه نۀ لري۔ لکه رسل چې پخپله وليدل، کټ مټ همداسې
استدلال به دې نتيجې ته مو ورسوي: \emph{هيڅ سړے کټ مټ د هغو
سړو ږيرې نۀ خريي چې خپلې ږيرې پخپله نۀ خريي}۔ يا:
\emph{د پګ نسل هيڅ سپے کټ مټ هغه پګ سپي نۀ بوي کوي چې خپل
ځانونه نۀ بوي کوي}۔ او داسې نور۔ د شېما په بڼه د نتيجې جوړښت
يوازې دا دے:
\[
\lnot \exists x \forall z(Rzx \liff \lnot R zz).
\]
دا د لومړۍ درجې منطق يوه قضيه ده؛ هره دوه‌ځايې اړيکه
پکې اخيستے شو۔ که د لومړي کمي کوونکي فرضي شاهد موجود وي، د
عمومي مادې تطبيق پر هماغه شاهد هغه رابطه له خپلې نفي سره
هم‌ارزښته کوي۔ نو
د خپلې سټ تيورۍ په لږ بدلولو د رسل له تناقضه نۀ شو بچ کېدے؛
دا تر هغې مخکې راولاړېږي چې سټ تيورۍ ته \emph{ورسېږو}۔ که
(کلاسيکي) د لومړۍ درجې منطق کاروو، نو دا خبره بايد
\emph{ومنو} چې سټ \LR{$R = \Setabs{x}{x\notin x}$} نشت۔
\olpseganchor{OLP-0532-B010}{end}

\olpseganchor{OLP-0532-B011}{start}
پايله دا ده: که ساده ادراک منل غواړئ او له ناسازګارۍ څخه
\emph{بچ کېدل} هم غواړئ، نو يوازې په \emph{سټ تيوري} کښې
لږ بدلون بس نۀ دے۔ بلکې خپل \emph{منطق} به له بنسټه بيا جوړوئ۔
\olpseganchor{OLP-0532-B011}{end}

\olpseganchor{OLP-0532-B012}{start}
البته، له غير کلاسيکي منطقونو سره سټ تيورۍ وړاندې شوې دي۔
خو، په لږو ټکو، معياري نۀ دي۔ د رسل د تناقض معياري چلند دا
دے چې د نۀ شتون د نېغ ثبوت په توګه يې ومنو او بيا زده کړو
چې څنګه ورسره کار پر مخ يوسو۔ مونږ به همدا لاره ونيسو۔
\olpseganchor{OLP-0532-B012}{end}

\olpseganchor{OLP-0533-B003}{start}
\olfileid{sth}{story}{predicative}	
\olpseganchor{OLP-0533-B003}{end}

\olpseganchor{OLP-0533-B004}{start}
\olsection{پريديکاتيف او امپريديکاتيف}
\olpseganchor{OLP-0533-B004}{end}

\olpseganchor{OLP-0533-B005}{start}
د رسل سټ \LR{$R$} د \LR{$\Setabs{x}{x \notin
x}$} له لارې تعريف شوے و۔ په بشپړه تشريح، \LR{$R$} به هغه سټ وے
چې ټول او يوازې هغه سټونه لري چې خپل ځان د غړي په توګه نۀ
لري۔ نو د \LR{$R$} په تعريف کښې پر داسې ساحه کمي کول کوو چې \LR{$R$} به
پکښې شامل وے (که وجود يې لرلے)۔
\olpseganchor{OLP-0533-B005}{end}

\olpseganchor{OLP-0533-B006}{start}
دا يو \emph{امپريديکاتيف} تعريف دے۔ په عمومي توګه، ويلے شو
چې يو تعريف هله او يوازې هله امپريديکاتيف دے چې د
کمي کولو په ساحه کښې تعريفېدونکے شے هم راځي۔
\olpseganchor{OLP-0533-B006}{end}

\olpseganchor{OLP-0533-B007}{start}
د تناقضونو له راولاړېدو وروسته، وايټ هېډ، رسل، پوانکارې او
وايل دغسې امپريديکاتيف تعريفونه د «بدې دايرې» په توګه رد کړل۔
\begin{explain}
  د وايټ هېډ او رسل په شننه کښې ستونزه هغه وخت پيدا کېږي
  چې د يو غړي د ټاکلو لپاره د ټولې ټولګې يادونه ضروري وي
  او ټاکل شوے څيز بيا د هماغې ټولګې يو غړے وګڼل شي۔
\olpseganchor{OLP-0533-B007}{end}

\olpseganchor{OLP-0533-B008}{start}
  د هغوي د اصل له مخې، کوم څيز چې تعريف يې د ټولګې
  \emph{ټول} غړي په بر کښې نيسي، پخپله د هغې ټولګې غړے نۀ
  دے۔ په بله بڼه، که د يوې ټولګې د کليت منل د داسې غړو
  منلو ته اړ باسي چې هماغه کليت د هغوي د تعريف شرط وي،
  نو د ټولګې د کليت فرض پرېږدو۔ دواړه بيانونه د بدې دايرې
  د مخنيوي يوه لار ښيي۔
  \mbox{\LR{\latinfont \citep[pp.~39--40]{WhiteheadRussell1910}}}
\end{explain}
که د هغوي په څېر \emph{د بدې دايرې اصل} ومنو، نو ښايي د
ساده ادراک زيانمنه شېما (د \olref[sth][story][rus]{sec})
په لاندې اصل بدله کړو۔ په دې اصل کښې پخواني سټونه د
تعريف ساحه جوړوي او نوې ټولګه بېل ډول لري:
\olpseganchor{OLP-0533-B008}{end}

\olpseganchor{OLP-0533-B009}{start}
\begin{defish}
\emph{پريديکاتيف ادراک۔} د هر فارمول \LR{$\phi$} لپاره چې يوازې پر سټونو کمي کول کوي: سټ\LR{$^\prime$} \LR{$\Setabs{x}{\phi(x)}$} وجود لري۔
\end{defish}
\olpseganchor{OLP-0533-B009}{end}

\olpseganchor{OLP-0533-B010}{start}
تر څو چې سټونه\LR{$^{\prime}$} له پخوانيو سټونو بېل ډول وي،
د رسل دغه دليل په دې قاعده نۀ چلېږي۔ د نوې ټولګې په تعريف
کښې متغير يوازې پر پخوانيو سټونو ګرځي، نو نوې ټولګه خپله
د متغير يو روا ارزښت نۀ دے۔
\olpseganchor{OLP-0533-B010}{end}

\olpseganchor{OLP-0533-B011}{start}
له بده مرغه، پريديکاتيف ادراک ډېر \emph{جامع} نۀ دے۔ ځکه
دا له نويو شيانو، سټونو\LR{$^\prime$}، سره مو اشنا کوي۔ نو بايد
هغه فارمولونه وڅېړو چې پر سټونو\LR{$^\prime$} کمي کول کوي۔ که دا
تل يو سټ\LR{$^\prime$} راکوي، نو د رسل تناقض به بيا راولاړ شي؛
يوازې د هغو ټولو سټونو\LR{$^\prime$} د سټ\LR{$^\prime$} په څېړلو چې
خپل ځان د غړي په توګه نۀ لري۔ نو د همدغه فکر په تعقيب بايد
ووايو چې پر سټونو\LR{$^\prime$} کمي کول کوونکے فارمول يو اړوند
سټ\LR{$^{\prime\prime}$} راکوي۔ بيا به سټونه\LR{$^{\prime\prime\prime}$}،
سټونه\LR{$^{\prime\prime\prime\prime}$}، او داسې نور پکار وي۔ د
پرايمونو د ډېرېدو د مخنيوي لپاره به دا اسانه وي چې د
سټونو\LR{$_0$}، سټونو\LR{$_1$}، سټونو\LR{$_2$}، سټونو\LR{$_3$}، سټونو\LR{$_4$}،\ldots
په توګه يې وګڼو۔ دا به د ډولونو (ساده) تيورۍ ته د ننوتلو
يوه لاره راکړي۔
\olpseganchor{OLP-0533-B011}{end}

\olpseganchor{OLP-0533-B012}{start}
پر داسې تيورۍ څو څرګند اعتراضونه شته (که څه هم دا څرګنده
نۀ ده چې \emph{پرېکنده} اعتراضونه دي)۔ په لږو ټکو: له پيدا
شوې تيورۍ سره کار کول ستونزمن دي؛ د شيانو د بېلو ډولونو په
فرض کولو کښې ډېرښت کوي؛ او په پاے کښې دا روښانه نۀ ده چې
امپريديکاتيف تعريفونه ان \emph{دومره بد} دي۔
\olpseganchor{OLP-0533-B012}{end}

\olpseganchor{OLP-0533-B013}{start}
د وروستي ټکي د روښانولو لپاره د \mbox{\LR{\latinfont \citeauthor{Ramsey1925}}}
دا خبره وګورئ:
\begin{quote}
  يو سړے په دې پېژندلے شو چې په خپله ډله کښې تر ټولو
  دنګ دے۔ په دې ډول يې د داسې کليت په وسيله ټاکو چې هغه
  پخپله يې يو غړے دے، او په دې کښې هېڅ بده دايره نشته۔
  \mbox{\LR{\latinfont \citep[p.~368]{Ramsey1925}}}
\end{quote}
که په ډله کښې يوازې يو کس تر نورو ټولو دنګ وي، د رمزي
وينا يو ټاکلے کس ښيي۔ «په ډله کښې تر ټولو دنګ سړے» يو
امپريديکاتيف تعريف \emph{دے}، او په همدې مثال کښې بده
دايره نۀ لري۔
\olpseganchor{OLP-0533-B013}{end}

\olpseganchor{OLP-0533-B014}{start}
ځواب دا کېدلے شي چې په دې حالت کښې تر ټولو دنګ کس په
\emph{پريديکاتيف} وسيله هم بېلولے شو۔ د بېلګې په توګه، ښايي
هغه سړي ته يوازې ګوته ونيسو۔ نو پر امپريديکاتيف تعريفونو
اعتراض به بايد په څرګنده هغو شيانو پورې محدود وي چې
\emph{يوازې} په امپريديکاتيف ډول بېلېدلے شي۔ خو بيا هم لا
زياته تشريح پکار ده چې ولې دغسې «لازمي امپريديکاتيف‌والے»
دومره بد دے۔\footnote{د نورو معلوماتو لپاره \mbox{\LR{\latinfont \citet{Linnebo2010}}} وګورئ۔}
\olpseganchor{OLP-0533-B014}{end}

\olpseganchor{OLP-0533-B015}{start}
دا منو چې امپريديکاتيف تعريفونه ډېرې بدې پايلې لري، که غواړو
زمونږ تعريفونه د يو شي د \emph{جوړولو} د نسخې په څېر څه
راکړي۔ ځکه له يوه امپريديکاتيف تعريف سره به په رښتيا په بده
دايره کښې ايسار يو: د امپريديکاتيف ډول ټاکل شوي شي د جوړولو
لپاره به \emph{لومړے} ټول شيان جوړوو (هماغه امپريديکاتيف ډول
ټاکل شوے شے هم)، ځکه دغه شے د ټولو شيانو له مخې ټاکل شوے
دے؛ نو د امپريديکاتيف ډول ټاکل شوي شي تر جوړولو مخکې به
هماغه شے جوړوو۔ خو بيا هم دا يوازې هغه وخت د «لازمي
امپريديکاتيف» ډول ټاکل شويو سټونو پر وړاندې يو جدي اعتراض
دے چې سټونه داسې شيان وګڼو چې مونږ يې \emph{جوړوو}۔ او
مونږ يې (ښايي) داسې نۀ ګڼو۔
\olpseganchor{OLP-0533-B015}{end}

\olpseganchor{OLP-0533-B016}{start}
نو، د ښه يا بد لپاره، هغه لاره چې عامه شوه، د
(ام)پريديکاتيف‌والي په باب سخت دريځ نۀ نيسي۔ بلکې هغه څه
رانغاړي چې اوس د تجمعي-تکراري لارې په توګه ګڼل کېږي۔ په
پاے کښې به دا اجازه راکړي چې خپل سټونه په «پړاوونو» کښې
طبقه‌بندي کړو، \emph{يو څه} لکه پريديکاتيف لاره چې شيان په
سټونو\LR{$_0$}، سټونو\LR{$_1$}، سټونو\LR{$_2$}، \ldots کښې طبقه‌بندي کوي؛
خو د هغو تر منځ به د نوعيت هيڅ توپير فرض نۀ کړو۔
\olpseganchor{OLP-0533-B016}{end}

\olpseganchor{OLP-0534-B003}{start}
\olfileid{sth}{story}{approach}
\olsection{تجمعي-تکراري لاره}
\olpseganchor{OLP-0534-B003}{end}

\olpseganchor{OLP-0534-B004}{start}
دا يې يو څه بشپړ بيان دے چې څنګه به سټونه په پړاوونو کښې
طبقه‌بندي کړو:
\begin{explain}
  د شوينفيلډ په تصور کښې د هر پړاو \LR{$S$} دپاره ټاکو چې کوم پړاوونه
  تر \LR{$S$} \emph{وړاندې} دي۔ د پړاو \LR{$S$} د غړو د انتخاب ساحه
  د همدغو پخوانيو پړاوونو جوړ شوي سټونه دي۔ د دې ساحې هره فرعي
  ټولګه په پړاو \LR{$S$} کښې يو سټ جوړوي۔ هر سټ لږ تر لږه په يوه
  پړاو کښې جوړېږي۔ \mbox{\LR{\latinfont \citep[p.~323]{Shoenfield:AST}}}
\end{explain}
دا \emph{د سټ د تجمعي-تکراري تصور} يوه خاکه ده۔ دا به د هغې
رسمي سټ تيورۍ بنسټ وي چې په \olref[sth][][]{part} کښې يې
وړاندې کوو۔
\olpseganchor{OLP-0534-B004}{end}

\olpseganchor{OLP-0534-B005}{start}
دا يو څه نور وڅېړو۔ لکه شوينفيلډ چې دا بهير بيانوي، په هر
پړاو کښې له هغو {سټونو} نوي سټونه جوړوو چې له مخکينيو
پړاوونو راپاتې وو او زمونږ لپاره موجود وو۔ نو د شوينفيلډ
د انځور له مخې، په لومړني پړاو، پړاو \LR{$0$}، کښې هيڅ
\emph{مخکيني} پړاوونه نشت؛ نو \emph{په لا قوي دليل} له
مخکينيو پړاوونو زمونږ لپاره هيڅ سټونه نشت۔\footnote{ولې
  بايد فرض کړو چې لومړنے پړاو \emph{شته}؟ په
  \olref[z][story]{sec} کښې د \stagesord{} پښه‌ليک وګورئ۔}
نو يوازې يو سټ جوړوو: بې له غړي سټ \LR{$\emptyset$}۔
په پړاو \LR{$1$} کښې له مخکينيو پړاوونو کټ مټ يو سټ راته موجود
دے، نو يوازې يو نوے سټ \LR{$\{\emptyset\}$} دے۔ په پړاو \LR{$2$} کښې
له مخکينيو پړاوونو دوه سټونه راته موجود دي او دوه نوي
سټونه \LR{$\{\{\emptyset\}\}$} او \LR{$\{\emptyset, \{\emptyset\}\}$} جوړوو۔
په پړاو \LR{$3$} کښې له مخکينيو پړاوونو څلور سټونه راته موجود
دي، نو دولس نوي سټونه جوړوو\ldots۔ نو د سټونو تجمعي-تکراري
انځور به يو څه داسې ښکاري (عددونه پړاوونه ښيي):
\begin{center}
	\begin{tikzpicture}[scale=0.6]
	\tikzset{cut_here/.style={densely dotted}}
	\draw (-4,6) -- (-1, 0)-- (2,6); 
	\draw[cut_here] (-5,8)--(-4,6);
	\draw[cut_here] (2,6)--(3,8);
	\draw(-1.5, 1)--(-0.5, 1);
	\draw(-2, 2)--(0, 2);
	\draw(-2.5, 3)--(0.5,3);
	\draw(-3, 4)--(1, 4);
	\draw(-3.5, 5)--(1.5, 5);
	\draw(-4, 6)--(2, 6);
	\node[label] at (0, 0) {\small 0};
	\node[label] at (0.5, 1) {\small 1};
	\node[label] at (1, 2) {\small 2};
	\node[label] at (1.5, 3) {\small 3};
	\node[label] at (2, 4) {\small 4};
	\node[label] at (2.5, 5) {\small 5};
	\node[label] at (3, 6) {\small 6};
	\end{tikzpicture}
\end{center}
نو ولې دا بيان ومنو؟
\olpseganchor{OLP-0534-B005}{end}

\olpseganchor{OLP-0534-B006}{start}
يوه وجه دا ده چې دا يو ښه او د کار وړ بيان دے۔ د تر ټولو
څرګند بيان، يعنې ساده ادراک، تر ناکامېدو وروسته يو ښه
بيان راته پکار دے۔
\olpseganchor{OLP-0534-B006}{end}

\olpseganchor{OLP-0534-B007}{start}
خو دا بيان \emph{يوازې} ښه نۀ دے۔ د دې باور لپاره ښه وجه
لرو چې د سټونو دا تصور د رسل له تناقض سره مخ نۀ شي۔
لاندې دليل دا ځانګړې خبره ثابتوي۔ د ټولې تيورۍ
\emph{سازګاري} جلا ثبوت غواړي۔
\olpseganchor{OLP-0534-B007}{end}

\olpseganchor{OLP-0534-B008}{start}
د سټ د تجمعي-تکراري تصور له مخې، سټونه په پړاوونو کښې
جوړوو؛ او د هغو غړي بايد داسې شيان وي چې
\emph{له مخکې} موجود وو۔ نو د هر پړاو~\LR{$S$} لپاره دا سټ جوړولے شو:
\[
  R_S = \Setabs{x}{x \notin x \text{\RL{ او \LR{$x$} تر \LR{$S$} وړاندې موجود و}}}
\]
د رسل د تناقض په ثبوت کښې راغلے استدلال به اوس دا ثابت کړي
چې \LR{$R_S$} پخپله تر پړاو \LR{$S$} وړاندې موجود نۀ دے۔ که له مخکې
موجود وے، د تعريف له مخې به يې غړيتوب د خپل غړيتوب له نفي
سره هم‌ارزښته وے۔ په ټاکلي پړاو کښې د نوي سټ راڅرګندېدل له
دې پايلې سره سم دي۔ سربېره پر دې، که د سټ تجمعي-تکراري تصور ومنو، نو
بايد ان دا \emph{تمه} مو هم نۀ واے لرلې چې پخپله د رسل سټ
جوړ کړو۔ ځکه هغه به د هغو ټولو خپل ځان د غړي په توګه نۀ
لرونکو سټونو سټ وي چې «هر کله به موجود شي»۔ په لږو ټکو:
دا حقيقت چې پخپله د رسل سټ (په ثبوت سره) نۀ شو جوړولے، د
تجمعي-تکراري بيان له مخې \emph{حيرانوونکے} نۀ دے؛ دا هماغه
څه دي چې مونږ به يې \emph{وړاندوينه} کوله۔
\olpseganchor{OLP-0534-B008}{end}

\olpseganchor{OLP-0534-B009}{start}
%In one sense, then, the cumulative-iterative conception of set yields a response to Russell's Paradox which is rather like the \emph{predicativist}'s response. After all, the predicativist said that the collection of all non-self-membered sets$_n$ is a set$_{n+1}$, and this set$_n$ will not be self-membered. (Indeed, most predicativists will treat it as \emph{ungrammatical} to try to ask whether a set is self-membered.)
%
%But there is an important difference between the cumulative-iterative approach and the predicativist approach: the cumulative-iterative approach treats all of our entities as being \emph{of the same kind}. The predicativist will presumably have an empty set$_0$ (i.e., a set$_0$ with no غړي), and an empty set$_1$ (i.e., a set$_1$ with no غړي), and these will be \emph{different entities}. But on the cumulative-iterative approach, there is just one empty set, $\emptyset$, and it will be available at every stage of the hierarchy. 
\olpseganchor{OLP-0534-B009}{end}

\olpseganchor{OLP-0534-B010}{start}
%Indeed, the Axiom of Extensionality, stated at the start of this chapter, will hold true of the elements in this hierarchy (so that there can be \emph{only one} ``empty'' set). But I do not intend to use the cumulative-iterative conception to justify Extensionality (see \mbox{\LR{\latinfont \cite{Boolos1971}}}). Again: I suggest that we should accept Extensionality, just because we are interested in extensional collections. 
\olpseganchor{OLP-0534-B010}{end}

\olpseganchor{OLP-0535-B003}{start}
\olfileid{sth}{story}{urelements}
\olsection{بنسټيز غير سټ غړي ومنو که نۀ؟}
\olpseganchor{OLP-0535-B003}{end}

\olpseganchor{OLP-0535-B004}{start}
په راتلونکو څو څپرکو کښې به هڅه وکړو چې د سټ له
تجمعي-تکراري تصور څخه اصلونه راوباسو۔ خو تر نور وړاندې
تګ مخکې بايد د \emph{بنسټيزو غير سټ غړو} په باب څه نور ووايو۔
\olpseganchor{OLP-0535-B004}{end}

\olpseganchor{OLP-0535-B005}{start}
د \olref[approach]{sec} انځور يوازې دا اجازه راکوله چې له
زړو \emph{سټونو} نوي سټونه جوړ کړو۔ خو ښايي دا اجازه هم
وغواړو چې ځينې \emph{غير سټونه}، غواګانې، خوګان، د شګو
دانې يا بل هر څه، د سټونو غړي شي۔ په دې حالت کښې
به له ځينو بنسټيزو غړو، \emph{بنسټيزو غير سټ غړو}، پيل
وکړو او بيا به ووايو چې په هر پړاو \LR{$S$} کښې هغه «ټول ممکن»
سټونه جوړوو چې غړي يې له بنسټيزو غير سټ غړو، له هغو سټونو
چې تر \LR{$S$} په مخکنيو پړاوونو کښې جوړ شوي وي، يا له دواړو څخه
اخيستل شوي وي۔ پيدا شوے انځور به دې ته ډېر ورته وي:
\begin{center}
	\begin{tikzpicture}[scale=0.6]
	\tikzset{cut_here/.style={densely dotted}}
	\draw (-4,6) -- (-1, 0) -- (1,0) -- (4,6); 
	\draw[cut_here] (-5,8)--(-4,6);
	\draw[cut_here] (5,8)--(4,6);
	\draw(-1.5, 1)--(1.5, 1);
	\draw(-2, 2)--(2, 2);
	\draw(-2.5, 3)--(2.5,3);
	\draw(-3, 4)--(3, 4);
	\draw(-3.5, 5)--(3.5, 5);
	\draw(-4, 6)--(4, 6);
	\node[label] at (2, 0) {\small 0};
	\node[label] at (2.5, 1) {\small 1};
	\node[label] at (3, 2) {\small 2};
	\node[label] at (3.5, 3) {\small 3};
	\node[label] at (4, 4) {\small 4};
	\node[label] at (4.5, 5) {\small 5};
	\node[label] at (5, 6) {\small 6};
	\end{tikzpicture}
\end{center}
نو اوس يوه پرېکړه راته پکار ده: \emph{آيا بنسټيز غير سټ غړي ومنو؟}
\olpseganchor{OLP-0535-B005}{end}

\olpseganchor{OLP-0535-B006}{start}
له فلسفي پلوه په خپله تيوري کښې د بنسټيزو غير سټ غړو
راګډول معنا لري۔ اصلي وجه يې د سټ تيورۍ د \emph{کارولو وړ}
ګرځول دي۔ د دې ټکي د څرګندولو لپاره د \olref[sfr][siz][]{chap}
دا خبره راياده کړئ چې وايو دوه سټونه \LR{$A$} او~\LR{$B$} يو شان
اندازه لري، يعنې \LR{$\cardeq{A}{B}$}، که او يوازې که تر منځ يې
يوه يو-پر-يو او پر-ټوله تابع وي۔ اوس که د پټي غواګانې او
د خوګانو په ځاے کښې خوګان دواړه سټونه جوړوي، نو د دې
ادعا لپاره يو سټ-تيوريکي چلند ورکولے شو چې «د غواګانو او
خوګانو شمېر يو شان دے»۔ خو که بنسټيز غير سټ غړي منع کړو،
نو غواګانې او خوګان پخپله د دغو سټونو غړي \emph{نۀ} شي
کېدے، او همدا نېغه سټ-تيوريکي طريقه به راته شونې نۀ وي۔
په خالصه سټ تيورۍ کښې به د داسې شيانو د ښودلو لپاره
لومړے ورته سټي استازي ټاکو۔ (د بنسټيزو غير سټ
غړو د راګډولو د نورو وجهو لپاره \mbox{\LR{\latinfont \citealt[pp.~vi, 24,
50--1]{Potter2004}}} وګورئ۔)
\olpseganchor{OLP-0535-B006}{end}

\olpseganchor{OLP-0535-B007}{start}
خو له رياضيکي پلوه د بنسټيزو غير سټ غړو منل ډېر کم دي۔
يوه وجه يې دا ده چې بې له بنسټيزو غير سټ غړو د سټ تيورۍ
بيانول \emph{ډېر لږ} اسانه دي۔ خو کله کله له سټ تيورۍ د
بنسټيز غير سټ غړي د ايستلو لپاره لا په زړه پورې دليلونه
موندل کېږي:
\begin{explain}
  د کونن د بنسټيزې موخې له مخې د رياضياتو څيزونه په سټونو
  تمثيلوو، نو د تيورۍ د متغيرونو د ګرځېدو دپاره د سټونو
  ساحه بس ده۔ د باندې نړۍ څيزونه، لکه څاروي، د دې ساحې
  غړي نۀ دي۔ \mbox{\LR{\latinfont \citep[p.~8]{Kunen1980}}}
\end{explain}
نو د کارېدنې پام د بنسټيزو غير سټ غړو \emph{راګډول} وړانديز
کوي؛ د بنسټ لپاره تحويلي موخه (خالصې سټ تيورۍ ته د رياضۍ
راکمولو موخه) د هغو \emph{ايستل} وړانديز کولے شي۔ لږه
سستي هم د بنسټيزو غير سټ غړو د ايستلو لوري ته اشاره کوي۔
\olpseganchor{OLP-0535-B007}{end}

\olpseganchor{OLP-0535-B008}{start}
مونږ به تر ټولو اسانه لاره ونيسو۔ خو يوه وجه يې تعليمي هم
ده۔ موخه مو دا ده چې د سټ تيورۍ له هغو بنسټونو سره مو اشنا
کړو چې د سټ تيورۍ د فلسفې د پيل لپاره يې اړتيا لرئ۔ او د
دغې فلسفې زياتره ليکنې هغه سټ تيورۍ څېړي چې \emph{بې له}
بنسټيزو غير سټ غړو بيان شوې دي۔ نو دا کتاب به ښايي ډېر
ګټور وي که له دغو ليکنو سره نژدې پاتې شي۔
\olpseganchor{OLP-0535-B008}{end}

\olpseganchor{OLP-0536-B003}{start}
\olfileid{sth}{story}{blv}
\olpseganchor{OLP-0536-B003}{end}

\olpseganchor{OLP-0536-B004}{start}
\olsection{ضميمه: د فرېګه پنځم بنسټيز قانون}
\olpseganchor{OLP-0536-B004}{end}

\olpseganchor{OLP-0536-B005}{start}
په \olref[rus]{sec} کښې مو بيان کړل چې رسل خپل تناقض د هغه
نظام د ستونزې په توګه وړاندې کړے و چې فرېګه په خپل
\emph{Grundgesetze} کښې بيان کړے و۔ د فرېګه په نظام کښې د
ساده ادراک نېغ بيان نۀ و شامل۔ نو په دې ضميمه کښې به ډېر
لنډ بيان کړو چې د فرېګه نظام څه \emph{لرل}، له ساده ادراک
سره څنګه تړاو لري او د رسل له تناقض سره يې څه اړيکه ده۔
\olpseganchor{OLP-0536-B005}{end}

\olpseganchor{OLP-0536-B006}{start}
د فرېګه نظام \emph{د دويمې درجې} دے او د \emph{مفهوم د
غځونې} د تصور د بيانولو لپاره طرح شوے و۔\footnote{په کره
  توګه، فرېګه د لا عمومي تصور د رسمي کولو هڅه کوي: د تابع
  «د ارزښتونو غځونه»۔ دا د تابع د حاصلېدونکو ارزښتونو ساده
  ټولګه نۀ ده۔ د مفهومونو غځونې د دې لا عمومي تصور
  ځانګړې بېلګه ده۔ د تفصيل لپاره \mbox{\LR{\latinfont \citet[pp.\ 8--9]{Heck2012}}}
  وګورئ۔} د فرېګه له نښو الهام اخيستے، \LR{$\fregeext{x}{F(x)}$}
به د \emph{مفهوم~\LR{$F$} د غځونې} لپاره وليکو۔ دا يوه وسيله
ده چې يو \emph{پريديکات}، «\LR{$F$}»، اخلي او په يوه (د لومړۍ
درجې) \emph{ترم}، «\LR{$\fregeext{x}{F(x)}$}»، يې بدلوي۔ فرېګه
په دې وسيله د غړيتوب دا \emph{تعريف} وړاندې کړ:
\[
  a \in b =_\text{\RL{df}} \exists G(b = \fregeext{x}{G(x)} \land Ga)
\]
په لنډه توګه: \LR{$a \in b$} که او يوازې که \LR{$a$} د داسې مفهوم
لاندې راځي چې غځونه يې \LR{$b$} ده۔ (ياد ولرئ چې کمي کوونکے
«\LR{$\exists G$}» د دويمې درجې دے۔) فرېګه دا اصل هم مانه چې
د \emph{پنځم بنسټيز قانون} په نوم پېژندل کېږي:
$$\fregeext{x}{F(x)} = \fregeext{x}{G(x)} \liff \forall x (Fx \liff Gx)$$
په لنډه توګه: مفهومونه عين غځونې لري که او يوازې که پر عين
شيانو صدق کوي۔ (بيا هم، دواړه «\LR{$F$}» او «\LR{$G$}» د پريديکات
په ځاے کښې دي۔) اوس يو ساده اصل غړيتوب د خاصيت له
پوره کولو سره تړي:
\olpseganchor{OLP-0536-B006}{end}

\olpseganchor{OLP-0536-B007}{start}
\begin{lem}[په \emph{Grundgesetze} کښې]\ollabel{lem:Fregeextensions}
\LR{$\forall F \forall a(a \in \fregeext{x}{F(x)} \liff Fa)$}
\end{lem}
\olpseganchor{OLP-0536-B007}{end}

\olpseganchor{OLP-0536-B008}{start}
\begin{proof} 
\LR{$F$} او \LR{$a$} ثابت ونيسئ۔ اوس \LR{$a \in \fregeext{x}{F(x)}$} که او يوازې که
\LR{$\exists G(\fregeext{x}{F(x)}
= \fregeext{x}{G(x)} \land Ga)$} (د غړيتوب د تعريف له مخې)، که او يوازې که
\LR{$\exists G(\forall x(Fx \liff Gx) \land Ga)$} (د پنځم بنسټيز قانون له مخې)، که او يوازې که \LR{$Fa$}
(د دويمې درجې د بنسټيز منطق له مخې)۔ په وروستي ګام کښې
که د وجودي جملې شاهد ومنو، د دواړو پريديکاتونو له يو شان
صدق څخه مطلوبه نتيجه راځي۔ په بل لوري کښې هماغه لومړے
پريديکات د وجودي کمي کوونکي شاهد اخلو۔
\end{proof}
\olpseganchor{OLP-0536-B008}{end}

\olpseganchor{OLP-0536-B009}{start}
او دا نژدې سملاسي ساده ادراک راکوي:
\olpseganchor{OLP-0536-B009}{end}

\olpseganchor{OLP-0536-B010}{start}
\begin{lem}[په \emph{Grundgesetze} کښې۔]
\LR{$\forall F \exists s \forall a (a \in s \liff Fa)$}
\end{lem}
\olpseganchor{OLP-0536-B010}{end}

\olpseganchor{OLP-0536-B011}{start}
\begin{proof}
\LR{$F$} ثابت ونيسئ؛ اوس \olref{lem:Fregeextensions} دا راکوي چې \LR{$\forall a (a \in
\fregeext{x}{F(x)} \liff Fa)$}؛ نو د وجودي عمومي کولو له مخې \LR{$\exists s\forall a(a \in s \liff Fa)$}۔
نتيجه ترې راځي، ځکه \LR{$F$} هر يو ټاکل کېدلے شو۔
\end{proof}
\olpseganchor{OLP-0536-B011}{end}

\olpseganchor{OLP-0536-B012}{start}
د دويمې درجې د ادراک له مخې داسې پريديکات ټاکلے شو:
که \LR{$F$} د \LR{$\forall x(Fx \liff x \notin x)$} له مخې ونيسو، د رسل تناقض ترې راځي۔
\olpseganchor{OLP-0536-B012}{end}

\olpseganchor{OLP-0537-B004}{start}
\olchapter{sth}{z}{د \texorpdfstring{\LR{$\Z$}}{Z} په لور ګامونه}
\olpseganchor{OLP-0537-B004}{end}

\olpseganchor{OLP-0538-B003}{start}
\olfileid{sth}{z}{story}
\olsection{د کيسې نور تفصيل}
\olpseganchor{OLP-0538-B003}{end}
	
\olpseganchor{OLP-0538-B004}{start}
په \olref[story][approach]{sec} کښې مو د سټونو د جوړېدو د بهير په اړه د
شوينفيلډ بيان رانقل کړے و. اوس غواړو يو څو نور اصول وليکو، څو دا
کيسه لږه نوره دقيقه شي. اصول دا دي:
\begin{enumerate}
\item[] \stageshier. هر سټ په يوه پړاو کښې جوړېږي.
\item[] \stagesord. پړاوونه ترتيب لري: ځينې د نورو \emph{نه مخکې}
راځي.\footnote{په حقيقت کښې به په پټه دا فرض کوو چې پړاوونه
\emph{ښه ترتيب شوي} دي. د دې مانا په \olref[ordinals][]{chap}
کښې تشريح کېږي. دا يو مهم فرض دے. په حقيقت کښې، د \mbox{\LR{\latinfont \citet{Scott1974}}}
د يوې ډېرې ځيرکې طريقې په کارولو سره، د ښه ترتيب له
\emph{فرضولو} ډډه کېدلے شي او ښه ترتيب بيا د نورو شرطونو
له مخې د \emph{نتيجې} په توګه راايستل کېدلے شي. (دا به دا هم
تشريح کړي چې ولې بايد يو لومړنے پړاو ومنو.) دلته دې بحث ته نۀ شو
ننوتلے؛ د نورو معلوماتو دپاره \mbox{\LR{\latinfont \citet{ButtonLT1}}} وګورئ.}
\item[] \stagesacc. د هر پړاو \LR{$S$} او د هغو سټونو د هرې ټولګې
دپاره چې له پړاو~\LR{$S$} \emph{نه مخکې} جوړ شوي وي: په پړاو~\LR{$S$}
کښې يو سټ جوړېږي چې غړي يې کټ مټ هماغه سټونه وي۔
په پړاو~\LR{$S$} کښې جوړېدونکي ټول سټونه همداسې جوړېږي۔
\end{enumerate}
دا غير رسمي اصول دي، خو د زرمېلو د سټ تيورۍ د څو بديهي اصلونو
د توجيه دپاره به ئې وکارولے شو.
\olpseganchor{OLP-0538-B004}{end}

\olpseganchor{OLP-0538-B005}{start}
(دلته يو احتياط بيانول پکار دي. که څه هم هغه بديهي اصلونه چې
وړاندې کوو ئې، او د هغو نومونه، په بشپړه توګه معياري دي، هغه په
مايل ليک ليکل شوي اصول چې اوس مو وړاندې کړل په اړوندو ليکنو کښې
ځانګړي نومونه نۀ لري. موږ ورته يوازې داسې نومونه غوره کړي دي
چې هيله لرو ګټور به وي.)
\olpseganchor{OLP-0538-B005}{end}

\olpseganchor{OLP-0539-B003}{start}
\olfileid{sth}{z}{sep}
\olsection{بېلونه}
\olpseganchor{OLP-0539-B003}{end}

\olpseganchor{OLP-0539-B004}{start}
له داسې يو اصل نه پېل کوو چې د ساده ادراک ځاے ونيسي:
\olpseganchor{OLP-0539-B004}{end}

\olpseganchor{OLP-0539-B005}{start}
\begin{axiom}[د بېلونې شېما] د هر فارمول \LR{$\phi(x)$} دپاره دا يو
بديهي اصل دے: د هر \LR{$A$} دپاره سټ \LR{$\Setabs{x \in A}{\phi(x)}$} شتون لري.
\end{axiom}
\olpseganchor{OLP-0539-B005}{end}

\olpseganchor{OLP-0539-B006}{start}
پام وکړئ چې دا يو يوازينے بديهي اصل نۀ دے. دا د بديهي اصلونو
يوه \emph{شېما} ده. د بېلونې \emph{بې‌نهايته ډېر} بديهي اصلونه
شته؛ د هر فارمول \LR{$\phi(x)$} دپاره يو اصل. شېما په همدې مانا
په لاندې ډول هم ليکل کېدلے شي (او معمولاً همداسې ليکل کېږي):
\olpseganchor{OLP-0539-B006}{end}

\olpseganchor{OLP-0539-B007}{start}
\begin{defish}
د هر هغه فارمول \LR{$\phi(x)$} دپاره چې «\LR{$S$}» پکښې نۀ وي، دا يو
بديهي اصل دے:
\[
	\forall A \exists S \forall x(x \in S \liff (\phi(x) \land x \in A)).
\]
\end{defish}
\olpseganchor{OLP-0539-B007}{end}

\olpseganchor{OLP-0539-B008}{start}
د \olref[sth][][]{part} په پېل کښې د ياد شوي دود له مخې،
د بېلونې د اصلونو فارمولونه~\LR{$\phi$} پارامترونه لرلے شي.\footnote{د دې
خبرې د مانا د تشريح دپاره هغه بحث وګورئ چې له
\olref[sfr][infinite][induction]{natinductionschema} نه سمدستي وروسته راځي.}
\olpseganchor{OLP-0539-B008}{end}

\olpseganchor{OLP-0539-B009}{start}
د سټونو هغه تجمعي-تکراري تصور چې بيان کړے مو دے، بېلونه په
نېغه توجيه کوي. د دې د ليدلو دپاره \LR{$A$} يو سټ وګرځوئ. نو \LR{$A$}
په يوه پړاو~\LR{$S$} کښې جوړېږي (د \stageshier له مخې). ځکه چې \LR{$A$}
په پړاو~\LR{$S$} کښې جوړ شوے دے، د \LR{$A$} ټول غړي له پړاو \LR{$S$} نه
مخکې جوړ شوي دي (د \stagesacc له مخې). اوس په ځانګړې توګه
هغه ټول سټونه په پام کښې ونيسئ چې د \LR{$A$} غړي دي او \LR{$\phi$} هم
پوره کوي؛ څرګنده ده چې دا ټول سټونه هم له پړاو~\LR{$S$} نه مخکې
جوړ شوي دي. نو دا هم په سټ \LR{$\Setabs{x \in A}{\phi(x)}$} کښې په پړاو~\LR{$S$}
کښې راټولېږي (د \stagesacc له مخې).
\olpseganchor{OLP-0539-B009}{end}

\olpseganchor{OLP-0539-B010}{start}
د ساده ادراک په خلاف، دا د رسل له تناقضه ځان ساتي. ځکه چې د سټ
\LR{$\Setabs{x}{x \notin
x}$} د شتون ادعا په ساده توګه نۀ شو کولے. بلکې، د يو سټ~\LR{$A$}
د \emph{راکړل کېدو} په صورت کښې، د سټ
\LR{$R_A = \Setabs{x \in A}{x \notin x}$} د شتون ادعا کولے شو.
که دغه سټ خپل غړے وے، د تعريف له مخې به خپل غړے نۀ و۔
نو \LR{$R_A \notin R_A$}۔ که د راکړل شوي سټ غړے وے، د همدې
تعريف له مخې به خپل غړے هم و۔ نو \LR{$R_A \notin A$}۔
\olpseganchor{OLP-0539-B010}{end}

\olpseganchor{OLP-0539-B011}{start}
خو بېلونه يوه سمدستي او د پام وړ پايله لري:
\olpseganchor{OLP-0539-B011}{end}

\olpseganchor{OLP-0539-B012}{start}
\begin{thm}\ollabel{thm:NoUniversalSet}
هېڅ \emph{عمومي} سټ نۀ شته، يعنې \LR{$\Setabs{x}{x = x}$} شتون نۀ لري.
\end{thm}
\olpseganchor{OLP-0539-B012}{end}

\olpseganchor{OLP-0539-B013}{start}
\begin{proof}
د تناقض له لارې د ثبوت دپاره فرض کړئ چې \LR{$V$} يو عمومي سټ دے.
نو د بېلونې له مخې \LR{$R =
\Setabs{x \in V}{x \notin x} = \Setabs{x}{x \notin x}$} شتون لري،
چې د رسل له تناقض سره ټکر کوي.
\end{proof}
\olpseganchor{OLP-0539-B013}{end}

\olpseganchor{OLP-0539-B014}{start}
د عمومي سټ نۀ شتون---بلکې د سټونو د مراتبو د لړۍ پرانيستے پاتې
کېدل---د تجمعي-تکراري تصور له تر ټولو بنسټيزو مفکورو څخه دي.
نو دا ليدل ګټور دي چې په شهودي توګه دې پايلې ته له بلې لارې هم
رسېدلے شو. يو عمومي سټ بايد د خپل ځان غړے وي.
خو زموږ د تجمعي-تکراري تصور له مخې، هر سټ د مراتبو په لړۍ کښې
لومړے په داسې پړاو کښې راڅرګندېږي چې ټول غړي يې
ترې وړاندې جوړ شوي وي۔ نو \emph{هېڅ} سټ د خپل ځان غړے نۀ
دے۔ که يو سټ خپل غړے وے، بايد د خپل لومړي پړاو نه مخکې
لا موجود وے، چې د لومړي پړاو له تعريف سره ټکر لري۔
(په \olref[sth][spine][rank]{defnsetrank} کښې به ووينو چې د سټ
د «رتبې» د مفهوم په کارولو سره دا تشريح څنګه لا دقيقه کولے شو.
خو د دې دپاره به لومړے يو څو نور بديهي اصلونه هم ولرو.)
\olpseganchor{OLP-0539-B014}{end}

\olpseganchor{OLP-0539-B015}{start}
دلته د بېلونې او د غړو له مخې د برابرۍ يو څو نورې پايلې دي.
\olpseganchor{OLP-0539-B015}{end}

\olpseganchor{OLP-0539-B016}{start}
\begin{prop}\ollabel{prop:emptyexists}
که کوم سټ شتون ولري، نو \LR{$\emptyset$} شتون لري.
\end{prop}
\olpseganchor{OLP-0539-B016}{end}

\olpseganchor{OLP-0539-B017}{start}
\begin{proof}
که \LR{$A$} يو سټ وي، نو \LR{$\emptyset = \Setabs{x \in A}{x \neq x}$} د بېلونې له مخې شتون لري.
\end{proof}
\olpseganchor{OLP-0539-B017}{end}

\olpseganchor{OLP-0539-B018}{start}
\begin{prop}
\LR{$A \setminus B$} د هر دوو سټونو \LR{$A$} او \LR{$B$} دپاره شتون لري.
\end{prop}
\olpseganchor{OLP-0539-B018}{end}

\olpseganchor{OLP-0539-B019}{start}
\begin{proof}
\LR{$A \setminus B = \Setabs{x \in A}{x \notin B}$} د بېلونې له مخې شتون لري.
\end{proof}
\olpseganchor{OLP-0539-B019}{end}

\olpseganchor{OLP-0539-B020}{start}
دا هم څرګندېږي چې (نږدې) خپل‌سري تقاطعګانې شتون لري:
\olpseganchor{OLP-0539-B020}{end}

\olpseganchor{OLP-0539-B021}{start}
\begin{prop}\ollabel{prop:intersectionsexist}
که \LR{$A \neq \emptyset$} وي، نو \LR{$\bigcap A = \Setabs{x}{(\forall y \in A)x \in y}$} شتون لري.
\end{prop}
\olpseganchor{OLP-0539-B021}{end}

\olpseganchor{OLP-0539-B022}{start}
\begin{proof}
\LR{$A \neq \emptyset$} وګرځوئ؛ نو يو \LR{$c \in A$} شته. بيا
\LR{$\bigcap A =
\Setabs{x}{(\forall y \in A)x \in y} = \Setabs{x \in c}{(\forall y \in
A)x \in y}$} د بېلونې له مخې شتون لري.
\end{proof}
\olpseganchor{OLP-0539-B022}{end}

\olpseganchor{OLP-0539-B023}{start}
خو شرط \LR{$A \neq \emptyset$} ته پام وکړئ؛ ځکه چې \LR{$\bigcap
\emptyset$} به د تشې رښتينولۍ له مخې عمومي سټ وي، چې له
\olref{thm:NoUniversalSet} سره ټکر کوي.
\olpseganchor{OLP-0539-B023}{end}

\olpseganchor{OLP-0540-B003}{start}
\olfileid{sth}{z}{union}
\olsection{اتحاد}
\olpseganchor{OLP-0540-B003}{end}

\olpseganchor{OLP-0540-B004}{start}
\olref[sth][z][sep]{prop:intersectionsexist} تقاطعګانې راکړې.
خو که غواړو چې خپل‌سري اتحادونه شتون ولري، نو يو بل بديهي
اصل هم ټاکل پکار دي:
\olpseganchor{OLP-0540-B004}{end}

\olpseganchor{OLP-0540-B005}{start}
\begin{axiom}[اتحاد] د هر سټ \LR{$A$} دپاره سټ \LR{$\bigcup A =
\Setabs{x}{(\exists b \in A) x \in b}$} شتون لري.
\[
	\forall A \exists U \forall x(x \in U \liff (\exists b \in A)x \in b)
\]
\end{axiom}
\olpseganchor{OLP-0540-B005}{end}

\olpseganchor{OLP-0540-B006}{start}
د تجمعي-تکراري تصور له مخې دا بديهي اصل هم توجيه کېږي.
\LR{$A$} يو سټ وګرځوئ؛ نو \LR{$A$} په يوه پړاو \LR{$S$} کښې جوړېږي
(د \stageshier له مخې). د \LR{$A$} هر غړے له \LR{$S$} \emph{نه مخکې}
جوړ شوے دے (د \stagesacc له مخې)۔ د هر داسې غړي خپل
غړي د هغه له پړاوه هم وړاندې جوړ شوي وي۔ د پړاوونو د
ترتيب د انتقاليت له مخې، د \LR{$A$} د هر غړي هر غړے له \LR{$S$}
نه مخکې جوړ شوے دے۔
په دې توګه دا \emph{ټول} سټونه له \LR{$S$} نه مخکې موجود دي،
څو په \LR{$S$} کښې په يوه سټ کښې راټول شي. او همدغه سټ \LR{$\bigcup A$} دے.
\olpseganchor{OLP-0540-B006}{end}

\olpseganchor{OLP-0541-B003}{start}
\olfileid{sth}{z}{pairs}
\olsection{جوړې}
\olpseganchor{OLP-0541-B003}{end}

\olpseganchor{OLP-0541-B004}{start}
راتلونکے بديهي اصل چې په پام کښې ئې نيسو، دا دے:
\olpseganchor{OLP-0541-B004}{end}

\olpseganchor{OLP-0541-B005}{start}
\begin{axiom}[جوړې]
د هر دوو سټونو \LR{$a, b$} دپاره سټ \LR{$\{a, b\}$} شتون لري.
\[
	\forall a \forall b \exists P \forall x (x \in P \liff (x = a \lor x = b))
\]
\end{axiom}
\olpseganchor{OLP-0541-B005}{end}

\olpseganchor{OLP-0541-B006}{start}
د تکراري تصور په کارولو سره دا اصل په لاندې ډول توجيه کولے شو.
فرض کړئ چې \LR{$a$} په پړاو \LR{$S$} کښې موجود دے او \LR{$b$} په پړاو \LR{$T$}
کښې موجود دے. \LR{$M$} د \LR{$S$} او \LR{$T$} له پړاوونو څخه هغه يو وګرځوئ
چې وروسته راځي؛ که دواړه يو پړاو وي، هماغه واخلئ۔ نو
\LR{$a$} او \LR{$b$} دواړه تر پړاو \LR{$M$} پورې موجود دي۔ نو سټ
\LR{$\{a,b\}$} له \LR{$M$} نه په هر وروستي پړاو کښې جوړېدلے شي،
ځکه چې هلته يې دواړه غړي له مخکې موجود دي۔
\olpseganchor{OLP-0541-B006}{end}

\olpseganchor{OLP-0541-B007}{start}
خو لږ تم شئ!{} ولې فرض کړو چې له \LR{$M$} نه وروسته کوم پړاوونه
\emph{شته}؟ که نۀ وي، نو زموږ توجيه ناکامېږي. نو د جوړو د اصل
د توجيه دپاره به د \olref[sth][z][story]{sec} په کيسه کښې يو بل
اصل هم زياتوو، يعنې:
\begin{enumerate}
	\item[] \stagessucc. وروستے پړاو نۀ شته.
\end{enumerate}
ايا دا اصل توجيه لري؟ د شوينفيلډ په کيسه کښې دا خبره
\emph{په ښکاره} نۀ وه بيان شوې چې وروستے پړاو نۀ شته.
خو بيا هم، که څه هم دا (په دقيقه مانا) زموږ په کيسه کښې يو
اضافي اصل دے، له هغې بنسټيزې مفکورې سره ښه سمون خوري چې
سټونه په پړاوونو کښې جوړېږي. په راتلونکي بحث کښې به ئې
يوازې ومنو. او په دې توګه د جوړو بديهي اصل هم منو.
\olpseganchor{OLP-0541-B007}{end}

\olpseganchor{OLP-0541-B008}{start}
په دې نوي اصل سره د ډېرو نورو سټونو شتون ثابتولے شو. د بېلګې په توګه:
\olpseganchor{OLP-0541-B008}{end}

\olpseganchor{OLP-0541-B009}{start}
\begin{prop}\ollabel{prop:pairsconsequences}
د هر دوو سټونو \LR{$a$} او \LR{$b$} دپاره لاندې سټونه شتون لري:
	\begin{enumerate}
		\item\ollabel{singleton} \LR{$\{a\}$}
		\item\ollabel{binunion} \LR{$a \cup b$}
		\item\ollabel{tuples} \LR{$\tuple{a, b}$}
	\end{enumerate}
\end{prop}
\olpseganchor{OLP-0541-B009}{end}

\olpseganchor{OLP-0541-B010}{start}
\begin{proof}
\olref{singleton}. د جوړو له مخې \LR{$\{a, a\}$} شتون لري، چې د
غړو له مخې د برابرۍ په اساس هماغه \LR{$\{a\}$} دے.
\olpseganchor{OLP-0541-B010}{end}

\olpseganchor{OLP-0541-B011}{start}
\olref{binunion}. د جوړو له مخې \LR{$\{a, b\}$} شتون لري. اوس
\LR{$a \cup b = \bigcup
\{a, b\}$} د اتحاد له مخې شتون لري.
\olpseganchor{OLP-0541-B011}{end}

\olpseganchor{OLP-0541-B012}{start}
\olref{tuples}. د \olref{singleton} له مخې \LR{$\{a\}$} شتون لري.
د جوړو له مخې \LR{$\{a,
b\}$} شتون لري. اوس \LR{$\{\{a\}, \{a, b\}\} = \tuple{a, b}$}
د جوړو د اصل د بيا کارولو له مخې شتون لري.
\end{proof}
\olpseganchor{OLP-0541-B012}{end}

\olpseganchor{OLP-0541-B013}{start}
\begin{prob}
وښيئ چې د هر درېو سټونو \LR{$a, b, c$} دپاره سټ \LR{$\{a, b, c\}$} شتون لري.
\end{prob}
\olpseganchor{OLP-0541-B013}{end}

\olpseganchor{OLP-0541-B014}{start}
\begin{prob}
وښيئ چې د هرې لړۍ سټونو \LR{$a_1, \ldots, a_n$} دپاره سټ \LR{$\{a_1, \ldots,
a_n\}$} شتون لري.
\end{prob}
\olpseganchor{OLP-0541-B014}{end}

\olpseganchor{OLP-0541-B015}{start}
%\begin{proof} By two applications of Pairs, $\{\{a_1, a_2\}, \{a_1,
%   a_3\}\}$ exists. By Union and Extensionality, $\bigcup \{\{a_1,
%   a_2\}, \{a_1, a_3\}\} = \{a_1, a_2, a_3\}$ exists. Repeat this
%   trick as often as necessary. \end{proof}
\olpseganchor{OLP-0541-B015}{end}

\olpseganchor{OLP-0542-B003}{start}
\olfileid{sth}{z}{power}
\olsection{د فرعي سټونو سټونه}
\olpseganchor{OLP-0542-B003}{end}

\olpseganchor{OLP-0542-B004}{start}
د يو بل بديهي اصل په وړاندې کولو سره مخته ځو:
\olpseganchor{OLP-0542-B004}{end}

\olpseganchor{OLP-0542-B005}{start}
\begin{axiom}[د فرعي سټونو سټونه]
د هر سټ \LR{$A$} دپاره سټ \LR{$\Pow{A} = \Setabs{x}{x \subseteq A}$} شتون لري.
\[
	\forall A \exists P \forall x(x \in P \liff (\forall z \in x)z \in A)
\]
\end{axiom}
\olpseganchor{OLP-0542-B005}{end}

\olpseganchor{OLP-0542-B006}{start}
د دې اصل توجيه په نسبتاً نېغه توګه کېږي. فرض کړئ چې \LR{$A$} په
پړاو \LR{$S$} کښې جوړ شوے دے. نو د \LR{$A$} ټول غړي له \LR{$S$} نه مخکې
موجود وو (د \stagesacc له مخې). په دې توګه، د بېلونې په توجيه
کښې زموږ د استدلال په شان، د \LR{$A$} هر فرعي سټ تر پړاو \LR{$S$} پورې
جوړېږي. نو دا ټول له \LR{$S$} نه په هر وروسته پړاو کښې موجود دي،
څو په يوه سټ کښې راټول شي. او پوهېږو چې داسې يو پړاو شته،
ځکه چې \LR{$S$} وروستے پړاو نۀ دے (د \stagessucc له مخې).
نو \LR{$\Pow{A}$} شتون لري.
\olpseganchor{OLP-0542-B006}{end}

\olpseganchor{OLP-0542-B007}{start}
دلته د فرعي سټونو د سټونو يوه ښکلې پايله ده:
\olpseganchor{OLP-0542-B007}{end}

\olpseganchor{OLP-0542-B008}{start}
\begin{prop}\label{thm:Products}
د هر دوو راکړل شويو سټونو \LR{$A, B$} دپاره د هغو کارتيزي حاصل ضرب \LR{$A \times B$} شتون لري.
\end{prop}
\olpseganchor{OLP-0542-B008}{end}

\olpseganchor{OLP-0542-B009}{start}
\begin{proof}
سټ \LR{$\Pow{\Pow{A \cup B}}$} د فرعي سټونو د سټونو او
\olref[sth][z][pairs]{prop:pairsconsequences} له مخې شتون لري.
نو د بېلونې له مخې دا سټ شتون لري:
\[
	C = \Setabs{z \in \Pow{\Pow{A \cup B}}}{(\exists x \in A)(\exists y \in B) z = \tuple{x, y}}.
\]
اوس د هر \LR{$x \in A$} او \LR{$y \in B$} دپاره سټ \LR{$\tuple{x, y}$} د
\olref[sth][z][pairs]{prop:pairsconsequences} له مخې شتون لري.
پر دې سربېره، ځکه چې \LR{$x, y
\in A \cup B$} دي، نو \LR{$\{x\}, \{x, y\} \in \Pow{A \cup B}$} او
\LR{$\tuple{x,y} \in \Pow{\Pow{A \cup B}}$} لرو. نو \LR{$A \times B = C$}.
\end{proof}
\olpseganchor{OLP-0542-B009}{end}

\olpseganchor{OLP-0542-B010}{start}
په دې ثبوت کښې د فرعي سټونو اصل له بېلونې سره يوځاے کار کوي.
او دا د حيرانتيا خبره نۀ ده. بې له بېلونې به د فرعي سټونو اصل
دومره \emph{پياوړے} اصل نۀ و. ځکه چې بېلونه تضمينوي چې د يوه سټ هغه فرعي سټونه شتون
لري چې د بېلونې د شېما په فارمولونو ټاکل کېدلے شي۔ له همدې
لارې د فرعي سټونو د سټ دننه ټاکلې فرعي ټولګې جوړوو۔
\olpseganchor{OLP-0542-B010}{end}

\olpseganchor{OLP-0542-B011}{start}
\begin{prob}
وښيئ چې د هر دوو سټونو \LR{$A, B$} دپاره: (i) د هغو ټولو اړيکو سټ
شتون لري چې د تعريف ساحه ئې \LR{$A$} او د قيمتونو سټ ئې \LR{$B$} وي؛ او
(ii) له \LR{$A$} څخه \LR{$B$} ته د ټولو تابعو سټ شتون لري.
\end{prob}
\olpseganchor{OLP-0542-B011}{end}

\olpseganchor{OLP-0542-B012}{start}
\begin{prob}
\LR{$A$} يو سټ او \LR{$\sim$} په \LR{$A$} باندې د هم‌ارزۍ اړيکه وګرځوئ.
ثابت کړئ چې د \LR{$\sim$} له مخې په \LR{$A$} باندې د هم‌ارزۍ د ټولګيو
سټ، يعنې \LR{$\equivclass{A}{\sim}$}، شتون لري.
\end{prob}
\olpseganchor{OLP-0542-B012}{end}

\olpseganchor{OLP-0543-B003}{start}
\olfileid{sth}{z}{infinity-again}	
\olsection{نامتناهيت}
\olpseganchor{OLP-0543-B003}{end}

\olpseganchor{OLP-0543-B004}{start}
تر اوسه دومره بديهي اصلونه لرو چې د بې‌نهايته ډېرو سټونو شتون
تضمين کړي (که کوم سټونه وي). ځکه فرض کړئ چې کوم سټ شتون لري؛
نو \LR{$\emptyset$} شتون لري (د \olref[sth][z][sep]{prop:emptyexists} له مخې).
اوس د هر سټ~\LR{$x$} دپاره سټ \LR{$x \cup \{x\}$} د
\olref[sth][z][pairs]{prop:pairsconsequences} له مخې شتون لري.
نو د دې اصل په يو څو ځله کارولو سره به لاندې سټونه ترلاسه کړو:
\begin{enumerate}
	\item[0.] \LR{$\emptyset$} %& $0$ & stage $0$\\
	\item[1.] \LR{$ \{\emptyset\}$} %& $1$ & stage $1$\\
	\item[2.] \LR{$\{\emptyset, \{\emptyset\}\}$} %& $2$ & stage $2$\\
	\item[3.] \LR{$\{\emptyset, \{\emptyset\}, \{\emptyset, \{\emptyset\}\}\}$} %&$3$ & stage $3$\\
	\item[4.] \LR{$\{\emptyset, \{\emptyset\}, \{\emptyset, \{\emptyset\}\}, \{\emptyset, \{\emptyset\}, \{\emptyset, \{\emptyset\}\}\}\}$}% & $4$ & stage $4$ 
\end{enumerate}
او کتلے شو چې دا سټونه ټول له يو بل څخه بېل دي.
\olpseganchor{OLP-0543-B004}{end}

\olpseganchor{OLP-0543-B005}{start}
د څو وجوهاتو له مخې مو شمېرنه له \LR{$0$} څخه پېل کړې ده. خو يو
سبب ئې دا دے. دا کتل ډېر ګران نۀ دي چې هغه سټ چې «\LR{$n$}» نوم مو
ورکړے دے، په کره توګه \LR{$n$} غړي لري او (په شهودي توګه) په
\LR{$n$}م پړاو کښې جوړېږي.
\olpseganchor{OLP-0543-B005}{end}

\olpseganchor{OLP-0543-B006}{start}
خو دا \emph{بې‌نهايته ډېر} سټونه راکوي، مګر د يوه \emph{نامتناهي سټ}
شتون نۀ تضمينوي؛ يعنې د داسې سټ شتون چې بې‌نهايته ډېر غړي ولري.
او دا خبره ډېره مهمه ده: تر څو د ډېډېکېنډ له مخې يو نامتناهي سټ
ونۀ مومو، د ډېډېکېنډ الجبر نۀ شو جوړولے. خو د ډېډېکېنډ الجبر
غواړو، څو ئې د طبيعي عددونو د سټ په توګه واخلو. (له
\olref[sfr][infinite][dedekindsproof]{sec} سره پرتله کړئ.)
\olpseganchor{OLP-0543-B006}{end}

\olpseganchor{OLP-0543-B007}{start}
هغه بديهي اصلونه چې تر اوسه مو ټاکلي دي، د هېڅ نامتناهي سټ
شتون \emph{نۀ} تضمينوي. نو يو نوے بديهي اصل ټاکل پکار دي:
\olpseganchor{OLP-0543-B007}{end}

\olpseganchor{OLP-0543-B008}{start}
\begin{axiom}[نامتناهيت]
يو سټ \LR{$I$} شته چې \LR{$\emptyset \in I$} وي او \LR{$x \cup \{x\} \in I$} وي، هر کله چې \LR{$x \in I$} وي.
\begin{align*}
	\exists I( & (\exists o \in I)\forall x\ x \notin o \land {}\\
	& (\forall x \in I)(\exists s \in I)\forall z(z \in s \liff (z \in x \lor z = x)))
\end{align*}
\end{axiom}
\olpseganchor{OLP-0543-B008}{end}

\olpseganchor{OLP-0543-B009}{start}
د نامتناهيت بديهي اصل داسې يو سټ \LR{$I$} راکوي چې
\LR{$\emptyset$} لري او په \LR{$I$} کښې د \LR{$s(x) = x\cup
\{x\}$} د قاعدې له مخې د هر غړي تالي هم شته۔ د طبيعي عددونو
دپاره يې تر ټولو کوچنۍ داسې برخه اخلو چې تش سټ لري او تر
تالي لاندې بنده ده۔ دا برخه د هغو ټولو فرعي سټونو تقاطع
ده چې دا دواړه شرطونه پوره کوي۔ دغه کورنۍ ناتشه ده، ځکه
خپله راکړل شوے سټ هم پکې راځي۔ د تقاطع شتون بېلونه تضمينوي۔
تش سټ په هر غړي کښې شته، او د تقاطع د هر شي تالي هم په
هر غړي کښې شته۔ نو تقاطع هم بنده ده۔
په \olref[sfr][infinite][dedekind]{thm:DedekindInfiniteAlgebra} کښې
د ډېډېکېنډ الجبر عمومي جوړښت بيان شوے دے۔ دلته يې ټاکلے
سټ په لاندې ډول دے:
\begin{defn}\ollabel{defnomega}
\LR{$I$} هغه هر سټ وګرځوئ چې د نامتناهيت بديهي اصل ئې راکړي.
\LR{$s$} تابع \LR{$s(x) = x \cup \{x\}$} وګرځوئ.
\LR{$\omega =
\closureofunder{s}{\emptyset}$} وګرځوئ. د \LR{$\omega$} غړو ته
\emph{طبيعي عددونه} وايو، او وايو چې \LR{$n$} د \LR{$n$} ځله د \LR{$s$} د
کارولو پايله ده، چې پر \LR{$\emptyset$} باندې ترسره شوي دي.
\end{defn}
په دې کوچنۍ برخه کښې تالي يو پر يو دے۔ د دې د ثبوت لپاره
پر تالي استقرا کوو۔ هره برخه چې تش سټ لري او تر تالي
لاندې بنده وي، ټوله دا کوچنۍ برخه رانغاړي۔ تش سټ متعدي
دے او خپل غړے نۀ دے۔
که يو سټ متعدي وي او خپل غړے نۀ وي، تالي يې هم متعدي دے۔
که تالي يې خپل غړے وے، نو يا به له اصلي سټ سره برابر وے،
يا به يې غړے وے۔ په لومړي حالت کښې به اصلي سټ خپل غړے وے،
ځکه چې هر سټ د خپل تالي غړے دے۔ په دويم حالت کښې هم
د اصلي سټ د انتقاليت له مخې هماغه تناقض راځي۔ نو په دې
برخه کښې هر سټ متعدي دے او خپل غړے نۀ دے۔ که دوه بېل غړي
عين تالي ولري، هر يو به د بل غړے وي؛ انتقاليت به بيا يو
له هغو څخه د خپل ځان غړے کړي۔ له دې څرګندېږي چې تالي يو
پر يو دے۔ تش سټ د تالي په حاصل کښې نۀ راځي، ځکه چې هر تالي
لږ تر لږه خپل مخکښ د غړي په توګه لري۔ نو دا لږ تر لږه
بنده برخه د ډېډېکېنډ له مخې نامتناهي ده۔ د تالي د تابعې
ګراف د همدې برخې له خپل ځان سره په کارتيزي حاصل ضرب کښې
د بېلونې له مخې شته۔ په ټول راکړل شوي سټ کښې هم يو پر يو
تابع لرو: په دې کوچنۍ برخه کښې تالي اخلو او پاتې غړي په
خپل ځاے پرېږدو۔ دا تابع تش سټ د ارزښت په توګه نۀ اخلي۔
\olpseganchor{OLP-0543-B009}{end}

\olpseganchor{OLP-0543-B010}{start}
اوس بېرته کتلے شئ چې هغه سټ چې يو څو پراګرافه مخکې «\LR{$n$}»
نوم ورکړل شوے و، د عدد \LR{$n$} \emph{په توګه} به اخيستل کېږي.
\olpseganchor{OLP-0543-B010}{end}

\olpseganchor{OLP-0543-B011}{start}
د دې ټاکنې د اهميت په اړه به په \olref[sth][z][nat]{sec} کښې
بحث وکړو. د اوس دپاره دا موږ ته د يوې شهودي پايلې ثبوت ممکنوي:
\olpseganchor{OLP-0543-B011}{end}

\olpseganchor{OLP-0543-B012}{start}
\begin{prop}\ollabel{naturalnumbersarentinfinite}
هېڅ طبيعي عدد د ډېډېکېنډ له مخې نامتناهي نۀ دے.
\end{prop}
\olpseganchor{OLP-0543-B012}{end}

\olpseganchor{OLP-0543-B013}{start}
\begin{proof}
ثبوت د استقرا په وسيله دے، يعنې د
\olref[sfr][infinite][induction]{thm:dedinfiniteinduction} له مخې.
څرګنده ده چې \LR{$0
= \emptyset$} د ډېډېکېنډ له مخې نامتناهي نۀ دے. د استقرا د ګام
دپاره به نقيض عکس ثابت کړو: که (په نامعقوله توګه) \LR{$s(n)$} د
ډېډېکېنډ له مخې نامتناهي وي، نو \LR{$n$} هم د ډېډېکېنډ له مخې نامتناهي دے.
\olpseganchor{OLP-0543-B013}{end}

\olpseganchor{OLP-0543-B014}{start}
نو فرض کړئ چې \LR{$s(n)$} د ډېډېکېنډ له مخې نامتناهي دے، يعنې
يوه يو پر يو تابع \LR{$f$} شته چې \LR{$\ran{f}\subsetneq \dom{f} = s(n) = n \cup
\{n\}$} وي. دوه حالتونه په پام کښې نيول پکار دي.
\olpseganchor{OLP-0543-B014}{end}

\olpseganchor{OLP-0543-B015}{start}
\emph{لومړے حالت: \LR{$n \notin \ran{f}$}.} نو \LR{$\ran{f} \subseteq n$} او
\LR{$f(n) \in n$} دي. \LR{$g = \funrestrictionto{f}{n}$} وګرځوئ؛ اوس \LR{$\ran{g} =
\ran{f} \setminus \{f(n)\} \subsetneq n = \dom{g}$} لرو. نو \LR{$n$}
د ډېډېکېنډ له مخې نامتناهي دے.
\olpseganchor{OLP-0543-B015}{end}

\olpseganchor{OLP-0543-B016}{start}
\emph{دويم حالت: \LR{$n \in \ran{f}$}.} يو \LR{$m \in \dom{f} \setminus \ran{f}$}
وټاکئ، او تابع \LR{$h$} د تعريف ساحې \LR{$s(n) = n \cup \{n\}$} سره تعريف کړئ:
\[
h(x) =
	\begin{cases}
		f(x) & \text{\RL{که }}f(x) \neq n\\
		m & \text{\RL{که }}f(x)=n
	\end{cases}
\]
نو \LR{$h$} او \LR{$f$} په هر ځاے کښې يو شان دي، پرته له دې چې \LR{$h(f^{-1}(n)) = m
\neq n = f(f^{-1}(n))$} دي. ځکه چې \LR{$f$} يو پر يو تابع ده او
نوے ارزښت له مخکيني حاصل څخه بهر ټاکل شوے، نو بدله شوې
تابع هم يو پر يو ده۔ همداراز \LR{$n \notin
\ran{h}$} او \LR{$\ran{h} \subsetneq \dom{h} = s(n)$} دي. اوس \LR{$n$}
د لومړي حالت د استدلال په کارولو سره د ډېډېکېنډ له مخې نامتناهي دے.
\end{proof}
\olpseganchor{OLP-0543-B016}{end}

\olpseganchor{OLP-0543-B017}{start}
خو دا پوښتنه لا پاتې ده چې د نامتناهيت بديهي اصل څنګه
\emph{توجيه} کولے شو. لنډ ځواب دا دے چې هغې کيسې ته چې بيان کړې
مو ده، يو بل اصل زياتول پکار دي. هغه اصل دا دے:
\begin{enumerate}
	\item[] \stagesinf. يو نامتناهي پړاو شته. يعنې داسې يو پړاو شته
چې (a) لومړنے پړاو نۀ دے، او (b) ترې مخکې ځينې پړاوونه شته،
خو (c) سمدستي مخکښ پړاو نۀ لري.
\end{enumerate}
له دې اصل څخه د نامتناهيت بديهي اصل په نېغه نتيجه کېږي.
پوهېږو چې طبيعي عدد \LR{$n$} په پړاو \LR{$n$} کښې جوړېږي. نو سټ
\LR{$\omega$} په لومړي نامتناهي پړاو کښې جوړېږي. او همدا \LR{$\omega$}
د نامتناهيت بديهي اصل شاهد دے.
\olpseganchor{OLP-0543-B017}{end}

\olpseganchor{OLP-0543-B018}{start}
خو دا موږ بېرته هغې پوښتنې ته ستنوي چې \stagesinf څنګه توجيه کولے
شو. د \stagessucc په شان، دا هم د تجمعي-تکراري مراتبو په اړه زموږ
د کيسې ښکاره برخه نۀ وه. خو خبره تر دې هم وړاندې ده: د تکراري
مراتبو په لړۍ کښې، چې سټونه پړاو په پړاو جوړېږي، هېڅ شے دا
نۀ لازموي چې په دې بهير کښې يو \emph{نامتناهي} پړاو هم ومنو.
دا تصور پوره سازګار ښکاري چې پړاوونه د طبيعي عددونو په شان مرتب وي.
\olpseganchor{OLP-0543-B018}{end}

\olpseganchor{OLP-0543-B019}{start}
خو له دې نه يوه څرګنده ستونزه پيدا کېږي. په
\olref[sfr][infinite][dedekindsproof]{sec} کښې مو د ډېډېکېنډ هغه
«ثبوت» وڅېړه چې د (فکرونو) يو د ډېډېکېنډ له مخې نامتناهي سټ شته.
کېدلے شي هغه ثبوت مو چندان قانع کوونکے نۀ وي ګڼلے. خو که
\stagesinf{} د سټ د تکراري تصور (يا «د فکر د قوانينو») له مخې
«راپورې نۀ تپل کېږي»، نو لا هم د ډېډېکېنډ له مخې د يوه نامتناهي
سټ د شتون د ادعا ذاتي توجيه نۀ لرو.
\olpseganchor{OLP-0543-B019}{end}

\olpseganchor{OLP-0543-B020}{start}
بېشکه دلته ډېر نور څه هم ويل کېدلے شي. خو هيله ده اوس دې ځاے ته
رسېدلي يئ چې دا فکر پېل کړئ چې د يوه بديهي اصل (يا اصل) د
توجيه دپاره \emph{څه پکار} دي. په راتلونکي بحث کښې به
\stagesinf{} يوازې منلے فرض کړو.
\olpseganchor{OLP-0543-B020}{end}

\olpseganchor{OLP-0544-B003}{start}
\olfileid{sth}{z}{milestone}
\olsection{\texorpdfstring{\LR{$\Z^-$}}{Z⁻}: يو مهم پړاو}
\olpseganchor{OLP-0544-B003}{end}

\olpseganchor{OLP-0544-B004}{start}
په راتلونکې برخه کښې به \stagesinf{} ته بيا ستانۀ شو.
خو د نامتناهيت بديهي اصل په لرلو سره يو مهم پړاو ته رسېدلي يو.
اوس د تيورۍ \LR{$\Zminus$} دپاره ټول اړين بديهي اصلونه لرو.
په تفصيل سره:
\olpseganchor{OLP-0544-B004}{end}

\olpseganchor{OLP-0544-B005}{start}
\begin{defn}
تيوري \LR{$\Zminus$} دا بديهي اصلونه لري: د غړو له مخې برابري،
اتحاد، جوړې، د فرعي سټونو سټونه، نامتناهيت، او د بېلونې د
شېما ټولې بېلګې.
\end{defn}
\olpseganchor{OLP-0544-B005}{end}

\olpseganchor{OLP-0544-B006}{start}
دا نوم د \emph{زرمېلو} د سټ تيورۍ ښکارندويي کوي (يو شے ترې
\emph{کم} دے چې وروسته به ورته راشو). دا وياړ د زرمېلو حق دے،
ځکه چې همدا تيوري ئې په اساسي توګه په خپل
\mbox{\LR{\latinfont \citeyear{Zermelo1908Untersuchungen}}} کښې جوړه کړې وه.\footnote{د تاريخ
او تخنيکي جزئياتو د په زړه پورې بحث دپاره \mbox{\LR{\latinfont \citet[Appendix
A]{Potter2004}}} وګورئ.}
\olpseganchor{OLP-0544-B006}{end}

\olpseganchor{OLP-0544-B007}{start}
دا تيوري دومره پياوړې ده چې ډېر زيات رياضيکي کار پکښې کولے شو.
په ځانګړې توګه، \emph{پکار ده} چې \olref[sfr][][]{part} ته بېرته
وګورئ او د هغو دليلونو رسمي بڼه په~\LR{$\Zminus$} کښې وازمويئ۔
ځينې نتيجې اضافي اصل غواړي۔ د بېلګې په توګه، د هرې پر-ټوله
تابع د ښي معکوس عمومي قضيه د ټاکنې له اصل سره هم‌ارزښته ده۔ (يو څه د دې کار تر
کولو وروسته ښايي وغواړئ چې مخته لاړ شئ او
\olref[sth][z][arbintersections]{sec} ولولئ.) نو له دې وروسته،
بې له نورو څرګندونو، ځانونه په \LR{$\Zminus$} کښې (لږ تر لږه) کار کوونکي ګڼو.
\olpseganchor{OLP-0544-B007}{end}

\olpseganchor{OLP-0545-B003}{start}
\olfileid{sth}{z}{nat}
\olsection{د خپلو طبيعي عددونو ټاکنه}
\olpseganchor{OLP-0545-B003}{end}

\olpseganchor{OLP-0545-B004}{start}
%Let us take it for granted that the informal principle \stagesinf{}
%can be justified. It is also, though, worth commenting on the
%particular axiom we extracted from that informal principle. 
\olpseganchor{OLP-0545-B004}{end}

\olpseganchor{OLP-0545-B005}{start}
په \olref[infinity-again]{defnomega} کښې مو د «طبيعي عددونو»
عبارت په ښکاره تعريف کړے دے. دا ټاکنه څنګه درک کول پکار دي؟
دا يوه مابعدالطبيعي ادعا نۀ ده، بلکې يوازې دا پرېکړه ده چې ځينې
سټونه د طبيعي عددونو په توګه \emph{واخلو}. د دې نظر پر وجوهاتو مو
په \olref[sfr][rel][ref]{sec}، \olref[sfr][arith][ref]{sec} او
\olref[sfr][infinite][dedekindsproof]{sec} کښې لږ بحث کړے و.
خو دا وجوهات لا هم دقيقولے شو.
\olpseganchor{OLP-0545-B005}{end}

\olpseganchor{OLP-0545-B006}{start}
زموږ د نامتناهيت بديهي اصل د \mbox{\LR{\latinfont \citet{VonNeumann1925}}} پيروي کوي.
خو دلته يو بل اصل دے چې پر ځاے ئې منلے شو:
\olpseganchor{OLP-0545-B006}{end}

\olpseganchor{OLP-0545-B007}{start}
\begin{defish}
\emph{د زرمېلو د \mbox{\LR{\latinfont \citeyear{Zermelo1908Untersuchungen}}} د نامتناهيت
بديهي اصل.} يو سټ \LR{$A$} شته چې \LR{$\emptyset \in A$} او
\LR{$(\forall x \in A)\{x\} \in A$} وي.
\end{defish}
\olpseganchor{OLP-0545-B007}{end}

\olpseganchor{OLP-0545-B008}{start}
که د خپل (له فون نويمان څخه الهام اخيستي) نامتناهيت اصل پر ځاے
مو د زرمېلو اصل کارولے و، نو په همدې ډول به د ډېډېکېنډ له مخې
يو نامتناهي سټ او ورسره د ډېډېکېنډ الجبر راکړل شوے و. د زرمېلو
په لاره کښې به زموږ د الجبر ځانګړے غړے بيا هم \LR{$\emptyset$} و
(زموږ د \LR{$0$} استازے)، او تالي به د تابعې
\LR{$x \mapsto \{x\}$} له مخې ټاکل کېدو۔ دا تابع يو پر يو ده،
ځکه چې د دوو يوغړيزو سټونو برابري د هغو د غړو برابري
لازموي۔ تش سټ د هغې په حاصل کښې نۀ شته۔ د فون نويمان
تالي د \LR{$x \mapsto x \cup \{x\}$} له مخې ټاکل کېږي۔
د دې تر ټولو ساده پايله دا ده
چې زرمېلو \LR{$2$} د \LR{$\{\{\emptyset\}\}$} په توګه اخلي، خو موږ
(له فون نويمان سره) \LR{$2$} د \LR{$\{\emptyset, \{\emptyset\}\}$} په توګه اخلو.
\olpseganchor{OLP-0545-B008}{end}

\olpseganchor{OLP-0545-B009}{start}
ولې د نامتناهيت يو اصل د بل پر ځاے غوره کړو؟ مهم عملي سبب دا
دے چې د فون نويمان لاره د نامتناهيت هاخوا عددونو د سمبالولو
دپاره ښه غځېدلے شي. دا به له \olref[ordinals][]{chap} څخه
وروسته وڅېړو. د \emph{حساب کولو} لپاره د دواړو الجبرو د
صفر استازي يو له بل سره مقابله کوو او د هر استازي تالي د
بل لوري له تالي سره يو ځاے ساتو۔ دغه مطابقت د تالي د تکرار
له مخې د دواړو د لږ تر لږه بندو برخو ټول غړي سره تړي۔
نو که څوک درته وايي چې طبيعي عددونه سټونه
\emph{دي}، څرګنده پوښتنه دا ده: \emph{کوم سټونه دي؟}
\olpseganchor{OLP-0545-B009}{end}

\olpseganchor{OLP-0545-B010}{start}
همدا دقيقه پوښتنه \mbox{\LR{\latinfont \citet{Benacerraf1965}}} مشهوره کړه.
خو دا ويل پکار دي چې دا د هغې ښکارندې يوازې تر ټولو مشهوره
بېلګه ده چې مخکې هم ډېر ځله ورسره مخامخ شوي يو. بنسټيزه خبره
دا ده: د سټ تيوري د ډېرو «شهودي» ډوله شيانو د \emph{تمثيل}
يوه لاره راکوي: حقيقي، ناطق، صحيح او طبيعي عددونه، هو؛ خو
مرتبه جوړې، تابعې او اړيکې هم. مګر د سټ تيوري د تمثيل يو
\emph{يوازينے} انتخاب هېڅکله نۀ راکوي. \emph{تل} داسې بديلونه
شته چې په نېغه به همدومره ښه کار راکړي.
\olpseganchor{OLP-0545-B010}{end}

\olpseganchor{OLP-0546-B003}{start}
\olfileid{sth}{z}{arbintersections}
\olsection{ضميمه: بندښت، ادراک او تقاطع}
\olpseganchor{OLP-0546-B003}{end}

\olpseganchor{OLP-0546-B004}{start}
په \olref[sth][z][milestone]{sec} کښې مو وړانديز وکړ چې
د \olref[sfr][][]{part} ساده کار ته بېرته وګورئ او ويې څېړئ
چې په~\LR{$\Zminus$} کښې ترسره کېدلے شي. که مو دا مشوره تعقيب کړې وي،
ښايي \emph{يوه} خبره مو ستونزمنه ليدلې وي: د ډېډېکېنډ د
\emph{بندښتونو} په بحث کښې د \emph{تقاطع} کارول.
\olpseganchor{OLP-0546-B004}{end}

\olpseganchor{OLP-0546-B005}{start}
د \olref[sfr][infinite][dedekind]{Closure} له مخې ياد ولرئ چې
\begin{align*}
  \closureofunder{f}{o} & = \bigcap\Setabs{X}{o \in X \text{\RL{ او \LR{$X$} د
\LR{$f$} له مخې بند دے}}}.
\intertext{\RL{د دې عمومي بڼه د لاندې ډول يو تعريف دے:}}
  C & = \bigcap\Setabs{X}{\phi(X)}.
\end{align*}
خو دا د خبرداري زنګ دے: ځکه چې ساده ادراک ناکام دے، نو هېڅ
تضمين نۀ شته چې \LR{$\Setabs{X}{\phi(X)}$} شتون لري۔
د داسې تعريف \emph{رواوالے} په دې پورې اړه لري چې د
تقاطع د غړو لپاره ټاکلې ساحه پيدا کړو۔
\olpseganchor{OLP-0546-B005}{end}

\olpseganchor{OLP-0546-B006}{start}
د شرط پوره کوونکے هر يو سټ د تقاطع د غړو لپاره
\emph{ساحه} جوړولے شي۔ همدا لاره په
\olref[sth][z][sep]{prop:intersectionsexist} کښې کارول شوې وه،
چې ولې \LR{$\bigcap A$} د هر \LR{$A \neq \emptyset$} دپاره شتون لري۔
اوس يې د يوې ټاکونکې فارمولې لپاره بيانوو۔
\olpseganchor{OLP-0546-B006}{end}

\olpseganchor{OLP-0546-B007}{start}
د غړو له مخې د برابرۍ په لرلو سره، که وغواړو \LR{$C$} د
\LR{$\bigcap\Setabs{X}{\phi(X)}$} په توګه تعريف کړو، په حقيقت کښې
يوازې يو داسې شے \LR{$C$} غواړو چې لاندې شرط پوره کړي:
\begin{equation}\ollabel{bicondelimarbintersection}
	\forall x(x \in C \liff \forall X(\phi(X) \lif x \in X))\tag{*}
\end{equation}
اوس فرض کړئ چې \emph{کوم} سټ \LR{$S$} شته چې \LR{$\phi(S)$} وي.
نو د \olref{bicondelimarbintersection} د ترلاسه کولو دپاره \LR{$C$}
د \emph{بېلونې} په کارولو سره په ساده توګه داسې تعريفولے شو:
\[
	C = \Setabs{x \in S}{\forall X(\phi(X) \lif x \in X)}.
\]
د تمرين په توګه وښيئ چې دا تعريف هماغه غوښتل شوے
\olref{bicondelimarbintersection} ورکوي۔ په يوه لوري کښې
د بېلونې له تعريفه عمومي شرط راځي۔ په بل لوري کښې هماغه
عمومي شرط پر ټاکلي شاهد هم پلې کړئ؛ په دې توګه د بېلونې
لپاره د هغه شاهد غړيتوب هم ترلاسه کېږي۔
%  fix $x$. First suppose that $x \in C$, as (re)defined; then both $x
%  \in S$ and $\forall X(\phi(X) \lif x \in X)$; but since
%  $\phi(S)$, this reduces to the condition that $\forall X(\phi(X)
%  \lif x \in X)$. Conversely, suppose $\forall X(\phi(X) \lif
%  x \in X)$; then, since $\phi(S)$, we also have that $x \in S$; so
%  $x \in c$, as (re)defined. 
او همدا عمومي لاره په تقاطعګانو تعريفولو کښې د ساده ادراک له هر
ظاهري استعماله ځان ساتل ممکنوي. په هغې ځانګړې بېلګه کښې چې
دا بحث ترې پېل شو، يعنې \LR{$\closureofunder{f}{o}$}، دا کار داسې
کېږي. د \olref[sfr][infinite][dedekind]{closureproperties} ثبوت مو
په دې يادونه پېل کړے و چې \LR{$o \in \ran{f}\cup\{o\}$} او
\LR{$\ran{f} \cup \{o\}$} د \LR{$f$} له مخې بند دے۔ دا د هغې تابعې
په سياق کښې دے چې له خپلې تعريف ساحې څخه هماغې ساحې
ته ځي او ټاکلے پيل غړے هم هلته دے۔ نو غوښتل شوے
تعريف داسې ورکولے شو:
\[
	\closureofunder{f}{o} = \Setabs{x \in \ran{f} \cup \{o\}}{(\forall X \ni o)(X \text{\RL{ د \LR{$f$} له مخې بند دے}} \lif x \in X)}.
\]
\olpseganchor{OLP-0546-B007}{end}

\olpseganchor{OLP-0547-B004}{start}
\olchapter{sth}{ordinals}{ترتيبي عددونه}
\olpseganchor{OLP-0547-B004}{end}

\olpseganchor{OLP-0548-B004}{start}
\olfileid{sth}{ordinals}{intro}
\olpseganchor{OLP-0548-B004}{end}

\olpseganchor{OLP-0548-B005}{start}
\olsection{سريزه}
\olpseganchor{OLP-0548-B005}{end}

\olpseganchor{OLP-0548-B006}{start}
په \olref[z][]{chap} کښې مو د مراتبو د لړۍ د يوه نامتناهي پړاو
شتون د \stagesinf{} په بڼه فرض کړے و (زموږ د نامتناهيت بديهي
اصل هم وګورئ). خو د \stagessucc{} په لرلو سره په نامتناهي پړاو
نۀ شو درېدلے؛ مخته تلل پکار دي. نو د لومړي نامتناهي پړاو نه په
ورپسې پړاو کښې د هغو سټونو ټولې ممکنې ټولګې جوړوو چې تر لومړي
نامتناهي پړاو پورې موجود وو۔ په هر ورپسې پړاو کښې همدا قاعده د
ټولو هغو سټونو لپاره کاروو چې تر هغه وړاندې جوړ شوي وو۔ \dots
\olpseganchor{OLP-0548-B006}{end}

\olpseganchor{OLP-0548-B007}{start}
دلته په پټه مو د عدد د يوه «شهودي» تصور کارول پېل کړي دي، چې
له مخې ئې د ټولو طبيعي عددونو \emph{وروسته} هم عددونه کېدلے شي.
په ځانګړې توګه، دلته مفهوم د \emph{نامتناهيت هاخوا ترتيبي عدد}
دے. د دې څپرکي مقصد دا دے چې دا تصور لا کره کړو. د ترتيبي عدد
عمومي مفهوم به وڅېړو، او بيا به په ښکاره ځينې سټونه د خپلو
ترتيبي عددونو په توګه تعريف کړو.
\olpseganchor{OLP-0548-B007}{end}

\olpseganchor{OLP-0549-B003}{start}
\olfileid{sth}{ordinals}{idea}	
\olsection{د ترتيبي عدد عمومي تصور}
\olpseganchor{OLP-0549-B003}{end}

\olpseganchor{OLP-0549-B004}{start}
طبيعي عددونه په خپل معمولي ترتيب کښې په پام کښې ونيسئ:
\begin{center}
	\begin{tikzpicture}
	\foreach \x/\xtext in {0, 1, 2, 3, 4, 5}
	{
		\node (\x a) at (\x, 1) {\small{\LR{$\x$}}};
		\node (\x a) at (\x.5, 1) {\small{\LR{$<$}}};
	}
	\node (ldots) at (6, 1) {\small{\LR{$\ldots$}}};
	\end{tikzpicture}
\end{center}
دې ته په اصطلاحي ژبه کښې \LR{$\omega$}-لړۍ وايو. او په رښتيا هم
دا عمومي ترتيب زموږ د سټونو د مراتبو د پړاوونو په لومړني جوړښت
کښې څرګندېږي. خو اوس فرض کړئ چې \LR{$0$} د دې لړۍ پاے ته يوسو،
داسې چې له نورو ټولو عددونو وروسته راشي:
\begin{center}
	\begin{tikzpicture}
	\foreach \x/\xtext in {1, 2, 3, 4, 5}
	{
		\node (\x a) at (\x, 1) {\small{\LR{$\x$}}};
		\node (\x a) at (\x.5, 1) {\small{\LR{$<$}}};
	}
	\node (ldots) at (6, 1) {\small{\LR{$\ldots$}}};
	\node (bea) at (6.5, 1) {\small{\LR{$<$}}};
	\node (noa) at (7, 1) {\small{\LR{$0$}}};
	\end{tikzpicture}
\end{center}
دلته هماغه شيان لرو، خو په بنسټيز ډول بل شان مرتب دي: زموږ
لومړني ترتيب وروستے غړے نۀ لرلو؛ نوے ترتيب ئې لري. په رښتيا
نوے ترتيب د شيانو له يوې \LR{$\omega$}-لړۍ (\LR{$1, 2, 3, 4, 5, \ldots$})
څخه جوړ دے چې ورپسې يو بل شے راځي. دا به \LR{$\omega+1$}-لړۍ وي.
\olpseganchor{OLP-0549-B004}{end}

\olpseganchor{OLP-0549-B005}{start}
يوازې د همدې شيانو په کارولو سره د ترتيب نور ډولونه هم جوړولے
شو. مثلاً ټول جفت عددونه (په خپل طبيعي ترتيب کښې) په پام کښې
ونيسئ، او ورپسې ټول طاق عددونه (په خپل طبيعي ترتيب کښې):
\begin{center}
	\begin{tikzpicture}
	\node(a1) at (1,1) {\small{\LR{$0$}}};
	\node(a2) at (1.5,1) {\small{\LR{$<$}}};
	\node(a3) at (2,1) {\small{\LR{$2$}}};
	\node(a4) at (2.5,1) {\small{\LR{$<$}}};
	\node(a5) at (3,1) {\small{\LR{$4$}}};
	\node(a6) at (3.5,1) {\small{\LR{$<$}}};
	\node(dots) at (4,1) {\small{\LR{$\ldots$}}};
	\node(aaoe) at (4.5,1) {\small{\LR{$<$}}};
	\node(b1) at (5,1) {\small{\LR{$1$}}};
	\node(b2) at (5.5,1) {\small{\LR{$<$}}};
	\node(b3) at (6,1) {\small{\LR{$3$}}};
	\node(b4) at (6.5,1) {\small{\LR{$<$}}};
	\node(b5) at (7,1) {\small{\LR{$\ldots$}}};
	\end{tikzpicture}
\end{center}
دا يوه \LR{$\omega$}-لړۍ ده چې ورپسې بله \LR{$\omega$}-لړۍ راځي؛ يعنې
يوه \LR{$\omega+\omega$}-لړۍ. د طاقو عددونو لومړے غړے دلته له
ټولو جفتو عددونو وروسته راځي۔ په درې واړو بېلګو کښې غړي
هماغه طبيعي عددونه دي؛ د هغو ترتيبونه بېل دي۔ 
\olpseganchor{OLP-0549-B005}{end}

\olpseganchor{OLP-0549-B006}{start}
همدغه شان نور هم مخته تللے شو. خو داسې يوه عمومي لاره غواړو
چې د \emph{ترتيبونو} په اړه دا خبرې پرې درک کړو. 
\olpseganchor{OLP-0549-B006}{end}

\olpseganchor{OLP-0550-B003}{start}
\olfileid{sth}{ordinals}{wo}
\olsection{ښه ترتيبونه}
\olpseganchor{OLP-0550-B003}{end}

\olpseganchor{OLP-0550-B004}{start}
بنسټيز مفهوم داسې دے:
\olpseganchor{OLP-0550-B004}{end}

\olpseganchor{OLP-0550-B005}{start}
\begin{defn}
اړيکه \LR{$<$} د \LR{$A$} \emph{ښه ترتيب} ورکوي، که او يوازې که دا دواړه شرطونه پوره کړي:
\begin{enumerate}
	\item \LR{$<$} متصله وي، يعنې د ټولو \LR{$a, b \in A$} دپاره يا \LR{$a < b$}
	يا \LR{$a = b$} يا \LR{$b < a$} وي؛
	\item د \LR{$A$} هر ناتشه فرعي سټ د \LR{$<$} له مخې يو مينيمال غړے ولري،
	يعنې که \LR{$\emptyset \neq X \subseteq A$} وي، نو \LR{$(\exists m \in
	X)(\forall z \in X)z \nless m$} وي.
\end{enumerate}
\end{defn}
\olpseganchor{OLP-0550-B005}{end}

\olpseganchor{OLP-0550-B006}{start}
په اسانۍ ليدلے شو چې هغه درې بېلګې چې اوس مو وڅېړلې، په رښتيا
د ښه ترتيب اړيکې وې.
\olpseganchor{OLP-0550-B006}{end}

\olpseganchor{OLP-0550-B007}{start}
\begin{prob}
\Olref[sth][ordinals][idea]{sec} په طبيعي عددونو باندې د ترتيب درې
بېلګې وړاندې کړې. وڅېړئ چې هر يو ښه ترتيب دے.
\end{prob}
\olpseganchor{OLP-0550-B007}{end}

\olpseganchor{OLP-0550-B008}{start}
دلته د ښه ترتيب په اړه ځينې ابتدايي، خو ډېرې مهمې مشاهدې دي.
\olpseganchor{OLP-0550-B008}{end}

\olpseganchor{OLP-0550-B009}{start}
\begin{prop}\ollabel{wo:strictorder}
که \LR{$<$} د \LR{$A$} ښه ترتيب ورکوي، نو د \LR{$A$} هر ناتشه فرعي سټ د \LR{$<$} له
مخې يو يوازينے تر ټولو کوچنے غړے لري؛ او \LR{$<$} غير انعکاسي، نامتقارنه او متعدي ده.
\end{prop}
\olpseganchor{OLP-0550-B009}{end}

\olpseganchor{OLP-0550-B010}{start}
\begin{proof}
که \LR{$X$} د \LR{$A$} ناتشه فرعي سټ وي، نو د \LR{$<$} له مخې يو مينيمال
غړے \LR{$m$} لري، يعنې \LR{$(\forall z \in X)z \nless m$} وي.
ځکه چې \LR{$<$} متصله ده، نو \LR{$(\forall z \in X)m \leq z$} دي.
نو \LR{$m$} د \LR{$<$} له مخې د \LR{$X$} تر ټولو کوچنے غړے دے.
\olpseganchor{OLP-0550-B010}{end}

\olpseganchor{OLP-0550-B011}{start}
د غير انعکاسيت دپاره \LR{$a \in A$} وټاکئ؛ د \LR{$<$} له مخې د \LR{$\{a\}$}
مينيمال غړے \LR{$a$} دے، نو \LR{$a \nless a$} دي۔
د تعديت دپاره، که \LR{$a < b < c$} وي، نو \LR{$\{a, b, c\}$} د
\LR{$<$} له مخې تر ټولو کوچنے غړے لري۔ منځنے او وروستے
غړي هر يو تر خپل ځان مخکې يو غړے لري، نو يو هم تر ټولو
کوچنے غړے نۀ شي کېدے۔ که لومړے او وروستے يو وي، بيا
به هېڅ تر ټولو کوچنے غړے نۀ وي۔ نو لومړے غړے له وروستي
غړي سره بېل او تر هغه مخکې دے، يعنې \LR{$a < c$} دي۔
نامتقارن توب د غير انعکاسيت او تعديت څخه نتيجه کېږي۔
که دوه بېل تر ټولو کوچني غړي وے، هر يو به تر بل مخکې
وے۔ دا د نامتقارن توب له مخې ناشونې ده۔
\end{proof}
\olpseganchor{OLP-0550-B011}{end}

\olpseganchor{OLP-0550-B012}{start}
\begin{prop}\ollabel{propwoinduction}
که \LR{$<$} د \LR{$A$} ښه ترتيب ورکوي، نو د هر فارمول \LR{$\phi(x)$} دپاره:
\[
	\text{\RL{که }}(\forall a \in A)((\forall b < a)\phi(b) \lif
		\phi(a))\text{\RL{ وي، نو }}(\forall a \in A)\phi(a).
\]
\end{prop}
\olpseganchor{OLP-0550-B012}{end}

\olpseganchor{OLP-0550-B013}{start}
\begin{proof}
نقيض عکس به ثابت کړو. فرض کړئ چې \LR{$\lnot(\forall a \in
A)\phi(a)$}، يعنې \LR{$X = \Setabs{x \in A}{\lnot\phi(x)} \neq
\emptyset$} دے. نو \LR{$X$} د \LR{$<$} له مخې يو مينيمال غړے \LR{$a$}
لري. په دې توګه \LR{$(\forall b < a)\phi(b)$}، خو \LR{$\lnot \phi(a)$} دي.
دلته تر هغه مخکيني غړي د راکړل شوي سټ غړي دي۔ نو دغه
تر ټولو لومړے مخالف مثال د استقرا له فرض سره ټکر لري۔
\end{proof}\noindent دا وروستے خاصيت بايد په طبيعي عددونو باندې
د قوي استقرا اصل در په ياد کړي، يعنې: که
\LR{$(\forall n \in \omega)((\forall m < n)\phi(m)
\lif \phi(n))$} وي، نو \LR{$(\forall n \in \omega)\phi(n)$} دي.
او همدا خاصيت ښه ترتيب يو ډېر \emph{ټينګ} مفهوم ګرځوي.\footnote{يوه
يادونه: ټول فارمولونه پارامترونه لرلے شي (مګر که په ښکاره بل ډول ويل شوي وي).}
\olpseganchor{OLP-0550-B013}{end}

\olpseganchor{OLP-0551-B003}{start}
\olfileid{sth}{ordinals}{iso}
\olsection{د ترتيب آيزومورفيزمونه}
\olpseganchor{OLP-0551-B003}{end}

\olpseganchor{OLP-0551-B004}{start}
د دې د بيانولو دپاره چې ښه ترتيب \emph{څومره} ټينګ مفهوم دے،
لومړے به د ښه ترتيبونو د پرتله کولو يوه لاره معرفي کړو.
\olpseganchor{OLP-0551-B004}{end}

\olpseganchor{OLP-0551-B005}{start}
\begin{defn}
يو \emph{ښه ترتيب} جوړه \LR{$\tuple{A, <}$} ده، داسې چې \LR{$<$} د \LR{$A$}
ښه ترتيب ورکوي. ښه ترتيبونه \LR{$\tuple{A, <}$} او \LR{$\tuple{B,
\lessdot}$} \emph{د ترتيب له مخې آيزومورف} دي، {که او يوازې که}
يوه دوه اړخيزه يو پر يو تابع \LR{$f \colon A \to B$} شته، داسې چې:
\LR{$x < y$}، که او يوازې که \LR{$f(x) \lessdot  f(y)$} وي.
په دې حالت کښې \LR{$\ordeq{\tuple{A, <}}{\tuple{B, \lessdot}}$} ليکو،
او وايو چې \LR{$f$} يو \emph{د ترتيب آيزومورفيزم} دے.
\end{defn}
\noindent
په راتلونکي بحث کښې به د لنډون دپاره د «د ترتيب آيزومورفيزمونو»
پر ځاے «آيزومورفيزمونه» وايو. په شهودي توګه آيزومورفيزمونه
هغه دوه اړخيزې يو پر يو تابعې دي چې جوړښت ساتي. دلته د آيزومورفيزمونو ځينې ساده حقيقتونه دي.
\olpseganchor{OLP-0551-B005}{end}

\olpseganchor{OLP-0551-B006}{start}
\begin{lem}\ollabel{isoscompose}
د آيزومورفيزمونو ترکيبونه آيزومورفيزمونه دي، يعنې: که
\LR{$f \colon A \to B$} او \LR{$g \colon B \to C$} آيزومورفيزمونه وي، نو
\LR{$(g \circ f)\colon A \to C$} يو آيزومورفيزم دے.
\end{lem}
\begin{prob}
	\olref[sth][ordinals][iso]{isoscompose} ثابت کړئ.
\end{prob}
\begin{proof}
د تمرين په توګه پرېښودل شوے.
\end{proof}
\begin{cor}\ollabel{ordisoisequiv}
	\LR{$\ordeq{X}{Y}$} د هم‌ارزۍ اړيکه ده.
\end{cor}
\begin{prop}\ollabel{ordisounique}
که \LR{$\tuple{A, <}$} او \LR{$\tuple{B, \lessdot}$} آيزومورف ښه ترتيبونه
وي، نو د هغو ترمنځ آيزومورفيزم يوازينے دے.
\end{prop}
\olpseganchor{OLP-0551-B006}{end}

\olpseganchor{OLP-0551-B007}{start}
\begin{proof}
\LR{$f$} او \LR{$g$} آيزومورفيزمونه \LR{$A \to B$} وګرځوئ. پايله به د
استقرا په وسيله ثابت کړو، يعنې د \olref[wo]{propwoinduction} په کارولو سره.
\LR{$a\in A$} وټاکئ او (د استقرا دپاره) فرض کړئ چې
\LR{$(\forall b < a)f(b) = g(b)$} دي. \LR{$x \in B$} وټاکئ.
\olpseganchor{OLP-0551-B007}{end}

\olpseganchor{OLP-0551-B008}{start}
که \LR{$x \lessdot f(a)$} وي، نو \LR{$f^{-1}(x) < a$} دي؛ په دې توګه
\LR{$g(f^{-1}(x)) \lessdot g(a)$} دي، ځکه چې \LR{$f$} او \LR{$g$}
آيزومورفيزمونه دي. خو ځکه چې \LR{$f^{-1}(x) < a$} دي، زموږ د فرض
له مخې \LR{$x =f(f^{-1}(x)) = g(f^{-1}(x))$} دي.
نو \LR{$x \lessdot g(a)$} دي. همدغه شان، که \LR{$x \lessdot g(a)$} وي،
نو \LR{$x \lessdot f(a)$} دي.
\olpseganchor{OLP-0551-B008}{end}

\olpseganchor{OLP-0551-B009}{start}
په عمومي توګه \LR{$(\forall x \in B)(x \lessdot f(a) \liff x \lessdot
g(a))$} لرو. د \olref[sfr][rel][ord]{prop:extensionality-strictlinearorders}
له مخې ترې \LR{$f(a) = g(a)$} نتيجه کېږي. نو د
\olref[wo]{propwoinduction} له مخې \LR{$(\forall a \in A)f(a) = g(a)$} دي.
\end{proof}\noindent
دا تر يوه حده ښيي چې ښه ترتيبونه ټينګ دي. خو د دې د نور بيان
دپاره به د نورو نښو معرفي کول ګټور وي.
\olpseganchor{OLP-0551-B009}{end}

\olpseganchor{OLP-0551-B010}{start}
\begin{defn}
کله چې \LR{$\tuple{A, <}$} يو ښه ترتيب او \LR{$a \in A$} وي،
\LR{$A_a = \Setabs{x \in A}{x < a}$} وګرځوئ. وايو چې \LR{$A_a$} د \LR{$A$}
يوه خاصه \emph{پيلنۍ برخه} ده (او اجازه ورکوو چې خپله \LR{$A$} د
\LR{$A$} يوه غير خاصه پيلنۍ برخه وي). \LR{$<_a$} د \LR{$<$} هماغې پيلنۍ
برخې ته محدودول وګرځوئ، يعنې
\LR{$\funrestrictionto{\mathord{<}}{A_a^2}$}.
\end{defn}
\noindent
د دې نښو په کارولو سره دا بيان او ثابتولے شو چې هېڅ ښه ترتيب
د خپلې هېڅ خاصې پيلنۍ برخې سره آيزومورف نۀ دے.
\olpseganchor{OLP-0551-B010}{end}

\olpseganchor{OLP-0551-B011}{start}
\begin{lem}\ollabel{wellordnotinitial}
که \LR{$\tuple{A, <}$} يو ښه ترتيب او \LR{$a \in A$} وي، نو
\LR{$\ordneq{\tuple{A, <}}{\tuple{A_a, <_a}}$} دي.
\end{lem}
\olpseganchor{OLP-0551-B011}{end}

\olpseganchor{OLP-0551-B012}{start}
\begin{proof}
د نقيض د ردولو دپاره فرض کړئ چې \LR{$f \colon A \to A_a$} يو
آيزومورفيزم دے. ځکه چې \LR{$f$} يوه بيجکشن او \LR{$A_a \subsetneq A$}
دے، نو د \olref[wo]{wo:strictorder} په کارولو سره \LR{$b \in A$}
د \LR{$<$} له مخې د \LR{$A$} تر ټولو کوچنے غړے وګرځوئ چې
\LR{$b \neq f(b)$} وي. دا به وښيو چې
\LR{$(\forall x \in A)(x<b \liff x < f(b))$} دي؛ له دې به د
\olref[sfr][rel][ord]{prop:extensionality-strictlinearorders}
له مخې \LR{$b = f(b)$} نتيجه شي، چې د نقيض ردول بشپړوي.
\olpseganchor{OLP-0551-B012}{end}

\olpseganchor{OLP-0551-B013}{start}
فرض کړئ چې \LR{$x < b$} دي. نو \LR{$x = f(x)$} دي، د \LR{$b$} د ټاکنې له مخې.
او \LR{$f(x) < f(b)$} دي، ځکه چې \LR{$f$} آيزومورفيزم دے.
نو \LR{$x < f(b)$} دي.
\olpseganchor{OLP-0551-B013}{end}

\olpseganchor{OLP-0551-B014}{start}
فرض کړئ چې \LR{$x < f(b)$} دي. نو \LR{$f^{-1}(x) < b$} دي، ځکه چې
\LR{$f$} آيزومورفيزم دے۔ د انځور اصلي مخکښ شته، ځکه چې له
انځور شوي غړي مخکې هر غړے په هماغې پيلنۍ برخې کښې دے۔
په دې توګه \LR{$f^{-1}(x) = x$} دي، د \LR{$b$} د ټاکنې له مخې. نو \LR{$x < b$} دي.
\end{proof}
\olpseganchor{OLP-0551-B014}{end}

\olpseganchor{OLP-0551-B015}{start}
زموږ راتلونکې پايله په تقريبي توګه دا ښيي چې د يوه آيزومورفيزم
«پيلنۍ برخه» آيزومورفيزم ده:
\olpseganchor{OLP-0551-B015}{end}

\olpseganchor{OLP-0551-B016}{start}
\begin{lem}\ollabel{wellordinitialsegment}
\LR{$\tuple{A, <}$} او \LR{$\tuple{B, \lessdot}$} ښه ترتيبونه وګرځوئ.
که \LR{$f\colon A \to B$} آيزومورفيزم او \LR{$a \in A$} وي، نو
\LR{$\funrestrictionto{f}{A_{a}} : A_a \to B_{f(a)}$} آيزومورفيزم دے.
\end{lem}
\olpseganchor{OLP-0551-B016}{end}

\olpseganchor{OLP-0551-B017}{start}
\begin{proof}
ځکه چې \LR{$f$} آيزومورفيزم دے:
	%Since $f$ is an isomorphism, $b < a$ iff $f(b) \lessdot f(a)$, so that 
\begin{align*}
	\funimage{f}{A_a} &= \funimage{f}{\Setabs{x \in A}{x < a}}\\
	&= \funimage{f}{\Setabs{f^{-1}(y) \in A}{f^{-1}(y) < a}} \\
	&= \Setabs{y \in B}{y \lessdot f(a)} \\
	&=B_{f(a)}
\end{align*}
او \LR{$\funrestrictionto{f}{A_a}$} ترتيب ساتي، ځکه چې \LR{$f$} ئې ساتي.
\end{proof}
\olpseganchor{OLP-0551-B017}{end}

\olpseganchor{OLP-0551-B018}{start}
زموږ راتلونکې دوه پايلې ثابتوي چې ښه ترتيبونه تل د پرتله وړ دي:
\olpseganchor{OLP-0551-B018}{end}

\olpseganchor{OLP-0551-B019}{start}
\begin{lem}\ollabel{lemordsegments}
\LR{$\tuple{A, <}$} او \LR{$\tuple{B, \lessdot}$} ښه ترتيبونه وګرځوئ.
که \LR{$\ordeq{\tuple{A_{a_1}, <_{a_1}}}{\tuple{B_{b_1}, \lessdot_{b_1}}}$}
او \LR{$\ordeq{\tuple{A_{{a_2}}, <_{a_2}}}{\tuple{B_{{b_2}},
\lessdot_{b_2}}}$} وي، نو \LR{${a_1}  < {a_2} \text{\RL{ که او يوازې که }}{b_1} \lessdot
{b_2}$} دي.
\end{lem}
\olpseganchor{OLP-0551-B019}{end}

\olpseganchor{OLP-0551-B020}{start}
\begin{proof}
\emph{له کيڼې خوا نه ښۍ خوا ته} به ثابت کړو؛ بل لورے هم ورته دے.
فرض کړئ چې دواړه
\LR{$\ordeq{\tuple{A_{a_1}, <_{a_1}}}{\tuple{B_{b_1},
\lessdot_{b_1}}}$} او \LR{$\ordeq{\tuple{A_{{a_2}},
<_{a_2}}}{\tuple{B_{{b_2}}, \lessdot_{b_2}}}$} دي، او زموږ
آيزومورفيزم \LR{$f \colon A_{{a_2}} \to B_{{b_2}}$} دے.
\LR{${a_1} < {a_2}$} وګرځوئ؛ نو د \olref{wellordinitialsegment} له مخې
\LR{$\ordeq{\tuple{A_{a_1}, <_{a_1}}}{\tuple{B_{f({a_1})},
\lessdot_{f({a_1})}}}$} دي. په دې توګه
\LR{$\ordeq{\tuple{B_{b_1}, \lessdot_{b_1}}}{\tuple{B_{f({a_1})},
\lessdot_{f({a_1})}}}$} دي؛ او د \olref{wellordnotinitial} له مخې
\LR{${b_1} = f({a_1})$} دي. اوس \LR{${b_1} \lessdot {b_2}$} دي، ځکه چې
د \LR{$f$} د قيمتونو سټ \LR{$B_{{b_2}}$} دے۔
\end{proof}
\olpseganchor{OLP-0551-B020}{end}

\olpseganchor{OLP-0551-B021}{start}
\begin{thm}\ollabel{thm:woalwayscomparable}
د هر دوو راکړل شويو ښه ترتيبونو دپاره يو د بل له يوې پيلنۍ برخې
(چې خاصه کېدل ئې ضروري نۀ دي) سره آيزومورف دے.
\end{thm}
\olpseganchor{OLP-0551-B021}{end}

\olpseganchor{OLP-0551-B022}{start}
\begin{proof}
\LR{$\tuple{A, <}$} او \LR{$\tuple{B, \lessdot}$} ښه ترتيبونه وګرځوئ.
د بېلونې په کارولو سره
\[
	f = \Setabs{\tuple{a, b} \in A \times B}{
		\ordeq{\tuple{A_a, <_a}}{\tuple{B_b, \lessdot_b}}}.
\]
وګرځوئ. د \olref{lemordsegments} له مخې
\LR{$a_1 < a_2$}، که او يوازې که \LR{$b_1 \lessdot b_2$} دي، د ټولو
\LR{$\tuple{a_1, b_1}, \tuple{a_2, b_2} \in f$} دپاره.
د دواړو ترتيبونو د اتصال له مخې، که دوه جوړې يوه خوا
يوه شان ولري، بله خوا يې هم يوه شان ده۔ نو دا يوه يو پر
يو تابع ده، او د خپلې ساحې څخه خپل حاصل ته پر ټوله ده۔
په دې توګه \LR{$f \colon \dom{f} \to \ran{f}$} آيزومورفيزم دے۔
\olpseganchor{OLP-0551-B022}{end}

\olpseganchor{OLP-0551-B023}{start}
که \LR{$a_2 \in \dom{f}$} او \LR{$a_1 < a_2$} وي، نو د
\olref{wellordinitialsegment} له مخې د اړوندو پيلنيو برخو
آيزومورفيزم تر مخکيني غړي پورې محدودوو او له همدې ځايه
\LR{$a_1 \in \dom{f}$} دي۔ نو \LR{$\dom{f}$} د \LR{$A$} پيلنۍ برخه ده۔ همدغه شان \LR{$\ran{f}$} د \LR{$B$}
پيلنۍ برخه ده. د نقيض د ردولو دپاره فرض کړئ چې دواړه
\emph{خاصې} پيلنۍ برخې دي. نو \LR{$a$} د \LR{$<$} له مخې د \LR{$A \setminus \dom{f}$}
تر ټولو کوچنے غړے وګرځوئ، داسې چې
\LR{$\dom{f} = A_a$} وي؛ او \LR{$b$} د \LR{$\lessdot$} له مخې د \LR{$B \setminus \ran{f}$}
تر ټولو کوچنے غړے وګرځوئ، داسې چې \LR{$\ran{f} = B_b$}
وي. نو \LR{$f \colon A_a \to B_b$} آيزومورفيزم دے، او له دې امله
\LR{$\tuple{a, b} \in f$} دي، چې تضاد دے.
\end{proof}
\olpseganchor{OLP-0551-B023}{end}

\olpseganchor{OLP-0552-B003}{start}
\olfileid{sth}{ordinals}{vn}	
\olsection[د فون نويمن جوړونه]{د فون نويمن له خوا د ترتيبي عددونو جوړونه}
\olpseganchor{OLP-0552-B003}{end}

\olpseganchor{OLP-0552-B004}{start}
\olref[sth][ordinals][iso]{thm:woalwayscomparable} يو فکر راولاړوي.
موږ ځينې شيان معرفي کولے شو چې \emph{د ترتيب ډولونه} ئې وبولو او د
ښه ترتيبونو استازيتوب پرې وکړو. که \LR{$\ordtype{A, <}$} د ښه ترتيب \LR{$\tuple{A, <}$} د ترتيب ډول
په توګه وليکو، نو غواړو دا دوه اصلونه ترلاسه کړو:
\begin{align*}
	\ordtype{A, <} = \ordtype{B, \lessdot} &
	\text{\RL{ هله او يوازې هله چې }} \ordeq{\tuple{A, <}}{\tuple{B, \lessdot}}\\
	\ordtype{A, <} < \ordtype{B, \lessdot}&
	\text{\RL{ هله او يوازې هله چې }}\ordeq{\tuple{A, <}}{\tuple{B_b, \lessdot_b}}\text{\RL{ د يو مناسب }}b \in B
\end{align*}
همدارنګه، هيله لرلے شو چې د ترتيب ډولونه \emph{د ځينو ټاکلو سټونو په توګه}
معرفي کړو، لکه څنګه چې طبيعي عددونه د ځينو ټاکلو سټونو په توګه معرفي کولے شو.
\olpseganchor{OLP-0552-B004}{end}

\olpseganchor{OLP-0552-B005}{start}
د دې تر ټولو عامه طريقه، او هغه لاره چې موږ به ئې تعقيب کړو، دا ده چې
د ترتيب دا ډولونه د ځينو \emph{کانوني} ښه مرتب شويو سټونو له لارې تعريف کړو.
دا کانوني سټونه په لومړي ځل فون نويمن معرفي کړي وو:
\olpseganchor{OLP-0552-B005}{end}

\olpseganchor{OLP-0552-B006}{start}
\begin{defn}
سټ \LR{$A$} \emph{متعدي} دے، {که او يوازې که} \LR{$(\forall x \in A)x \subseteq
A$} وي. بيا \LR{$A$} يو \emph{ترتيبي عدد} دے، {که او يوازې که} \LR{$A$} متعدي وي او
\LR{$\in$} ئې ښه ترتيب ورکړي۔ دلته غړيتوب د همدغه سټ د غړو
تر منځ اړيکه ده؛ د متعدي سټ د هر غړي ټول غړي بيا د اصلي
سټ غړي دي۔
\end{defn}
\noindent
په لاندې بحث کښې به د ترتيبي عددونو دپاره يوناني توري کاروو. له تعريف څخه
سمدستي نتيجه اخلو چې، که \LR{$\alpha$} ترتيبي عدد وي، نو
\LR{$\tuple{\alpha, \in_\alpha}$} يو ښه ترتيب دے، چې پکښې \LR{$\in_\alpha =
\Setabs{\tuple{x, y} \in \alpha^2}{x \in y}$} دے۔ د لنډون لپاره
وايو چې \LR{$\alpha$} \emph{پخپله} ښه ترتيب دے۔
\olpseganchor{OLP-0552-B006}{end}

\olpseganchor{OLP-0552-B007}{start}
زموږ لومړني څو ترتيبي عددونه دا دي:
\[
	\emptyset, \{\emptyset\},
	\{\emptyset, \{\emptyset\}\},
	\{\emptyset, \{\emptyset\}, \{\emptyset, \{\emptyset\}\}\}, \ldots
\]
و به وينئ چې دا هماغه لومړني څو ترتيبي عددونه دي چې زموږ د نامتناهيت په
بديهي اصل کښې مو ليدلي وو، يعنې د فون نويمن د \LR{$\omega$} په تعريف کښې
(\olref[sth][z][infinity-again]{sec} وګورئ). دا تصادف نۀ دے.
د فون نويمن د ترتيبي عددونو تعريف طبيعي عددونه ترتيبي عددونه ګڼي، خو
د نامتناهيت هاخوا ترتيبي عددونو ته هم ځای ورکوي.
\olpseganchor{OLP-0552-B007}{end}

\olpseganchor{OLP-0552-B008}{start}
د تل په شان اوس پوښتلے شو: ايا دا \emph{هماغه} ترتيبي عددونه دي؟
يا فون نويمن يوازې ځينې داسې سټونه راکړي دي چې د ترتيبي عددونو په توګه
ئې \emph{کارولے} شو؟ د دې پوښتنې بحثونه هغو بحثونو ته ورته دي چې په
\olref[sfr][rel][ref]{sec}،
\olref[sfr][arith][ref]{sec}،
\olref[sfr][infinite][dedekindsproof]{sec} او
\olref[sth][z][nat]{sec} کښې مو کړي وو؛ نو خبره به نۀ اوږدوو.
پر ځای ئې، په لاندې بحث کښې به «ترتيبي عددونه» د «فون نويمن ترتيبي
عددونو» دپاره کاروو.
\olpseganchor{OLP-0552-B008}{end}

\olpseganchor{OLP-0553-B003}{start}
\olfileid{sth}{ordinals}{basic}
\olsection{د ترتيبي عددونو بنسټيز خاصيتونه}
\olpseganchor{OLP-0553-B003}{end}

\olpseganchor{OLP-0553-B004}{start}
ومو ليدل چې لومړني څو ترتيبي عددونه طبيعي عددونه دي. د ترتيبي عددونو
د نظريې د جوړولو اصلي وجه دا ده چې د استقرا هغه اصل پراخ کړو چې په
طبيعي عددونو باندې عملي کېږي. دې مقصد ته به د ځينو ابتدايي نتيجو په
لړۍ ورسېږو.
\olpseganchor{OLP-0553-B004}{end}

\olpseganchor{OLP-0553-B005}{start}
\begin{lem}\ollabel{ordmemberord}
د ترتيبي عدد هر غړے ترتيبي عدد دے.
\end{lem}
\olpseganchor{OLP-0553-B005}{end}

\olpseganchor{OLP-0553-B006}{start}
\begin{proof}
فرض کړئ \LR{$\alpha$} ترتيبي عدد دے او \LR{$b \in \alpha$} دے. ځکه چې \LR{$\alpha$}
متعدي دے، \LR{$b \subseteq \alpha$} دے. نو \LR{$\in$} د \LR{$b$} ښه ترتيب ورکوي، ځکه
چې \LR{$\in$} د \LR{$\alpha$} ښه ترتيب ورکوي.
\olpseganchor{OLP-0553-B006}{end}

\olpseganchor{OLP-0553-B007}{start}
د دې د ليدلو دپاره چې \LR{$b$} متعدي دے، فرض کړئ \LR{$x \in c \in b$} دي.
نو \LR{$c \in \alpha$} دے، ځکه چې \LR{$b
\subseteq \alpha$} دے. بيا، ځکه چې \LR{$\alpha$} متعدي دے، \LR{$c \subseteq
\alpha$} دے، نو \LR{$x \in \alpha$} دے. په دې توګه \LR{$x, c, b \in \alpha$} دي.
خو \LR{$\in$} د \LR{$\alpha$} ښه ترتيب ورکوي، نو د
\olref[sth][ordinals][wo]{wo:strictorder} له مخې \LR{$\in$} په
\LR{$\alpha$} باندې متعدي اړيکه ده. نو له \LR{$x
\in c \in b$} څخه \LR{$x \in b$} ترلاسه کوو. په تعميمولو سره \LR{$c \subseteq b$}
\end{proof}
\olpseganchor{OLP-0553-B007}{end}

\olpseganchor{OLP-0553-B008}{start}
\begin{cor}\ollabel{ordissetofsmallerord}
\LR{$\alpha = \Setabs{\beta \in \alpha}{\beta \text{\RL{ ترتيبي عدد دے}}}$}، د
هر ترتيبي عدد~\LR{$\alpha$} دپاره.
\end{cor}
\olpseganchor{OLP-0553-B008}{end}

\olpseganchor{OLP-0553-B009}{start}
\begin{proof}
له \olref{ordmemberord} څخه سمدستي نتيجه اخلو.
\end{proof}
\olpseganchor{OLP-0553-B009}{end}

\olpseganchor{OLP-0553-B010}{start}
د لاندې دوو مهمو نتيجو،
\olref{ordinductionschema} او \olref{ordtrichotomy}، عمومي مطلب دا دے
چې ترتيبي عددونه پخپله د غړيتوب له مخې ښه مرتب دي:
\olpseganchor{OLP-0553-B010}{end}

\olpseganchor{OLP-0553-B011}{start}
\begin{thm}[د نامتناهيت هاخوا استقرا]\ollabel{ordinductionschema}
د هر فارمول \LR{$\phi(x)$} دپاره:
\[
	\text{\RL{که }}\exists \alpha \phi(\alpha)\text{\RL{ وي، نو }}\exists \alpha(\phi(\alpha)
	\land  (\forall \beta \in \alpha) \lnot \phi(\beta))
\]
دلته ښودل شوي کمّي کوونکي په ضمني توګه ترتيبي عددونو ته محدود دي.
\end{thm}
\olpseganchor{OLP-0553-B011}{end}

\olpseganchor{OLP-0553-B012}{start}
\begin{proof}
%\olref{ordinductionschema1}. 
فرض کړئ \LR{$\phi(\alpha)$} د کوم ترتيبي عدد \LR{$\alpha$} دپاره صدق کوي.
که \LR{$ (\forall \beta
\in \alpha) \lnot \phi(\beta)$} وي، نو مقصد ترلاسه شو. که داسې نۀ وي، نو ځکه
چې \LR{$\alpha$} ترتيبي عدد دے، د \LR{$\in$} له مخې يو تر ټولو کوچنے داسې غړے
لري چې \LR{$\phi$} پرې صدق کوي۔ د بېلونې له مخې د دغو غړو ناتشه
فرعي برخه جوړوو او د هغې تر ټولو کوچنے غړے اخلو۔ د
\olref{ordmemberord} له مخې دا ترتيبي عدد دے او هېڅ غړے يې
غوښتلے شرط نۀ پوره کوي۔
%	\olref{ordinductionschema2}. Suppose there is some ordinal
%	$\gamma$ such that $\lnot\phi(\gamma)$. Then by
%	\olref{ordinductionschema1}, there is an $\in$-minimal ordinal
%	$\alpha$ for which $\lnot\phi(\alpha)$. So $(\forall \beta <
%	\alpha) \phi(\beta)$, rendering the antecedent of the conditional
%	false.
\end{proof}
\noindent
وګورئ چې \olref{ordinductionschema} په هماغه معنا د دې شېما په بڼه هم
بيانولے شو:
\[
\text{\RL{که }}\forall \alpha((\forall \beta \in \alpha)\phi(\beta) \lif
\phi(\alpha))\text{\RL{ وي، نو }}\forall \alpha\phi(\alpha)
\]
يوازې په \olref{ordinductionschema} کښې \LR{$\lnot\phi(\alpha)$} اخلو او بيا
ابتدايي منطقي بدلونونه ترسره کوو.
\olpseganchor{OLP-0553-B012}{end}

\olpseganchor{OLP-0553-B013}{start}
\begin{thm}[درې ګونتوب]\ollabel{ordtrichotomy}
\LR{$\alpha \in \beta \lor \alpha = \beta \lor \beta \in \alpha$}، د هر دوو
ترتيبي عددونو \LR{$\alpha$} او \LR{$\beta$} دپاره.
\end{thm}
\olpseganchor{OLP-0553-B013}{end}

\olpseganchor{OLP-0553-B014}{start}
\begin{proof}
ثبوت په دوه ګونې استقرا ولاړ دے، يعنې
\olref{ordinductionschema} دوه ځله کاروو. وايو چې \LR{$x$} له \LR{$y$} سره
\emph{د پرتله کېدو وړ} دے، که او يوازې که \LR{$x \in y \lor x = y \lor y \in x$} وي.
\olpseganchor{OLP-0553-B014}{end}

\olpseganchor{OLP-0553-B015}{start}
د استقرا دپاره فرض کړئ چې په~\LR{$\alpha$} کښې هر ترتيبي عدد له
\emph{هر} ترتيبي عدد سره د پرتله کېدو وړ دے. د دويمې استقرا دپاره فرض کړئ
چې \LR{$\alpha$} په~\LR{$\beta$} کښې له هر ترتيبي عدد سره د پرتله کېدو وړ دے.
و به ښيو چې \LR{$\alpha$} له~\LR{$\beta$} سره د پرتله کېدو وړ دے. په~\LR{$\beta$} باندې
له استقرا څخه به نتيجه واخلو چې \LR{$\alpha$} له هر ترتيبي عدد سره د پرتله کېدو
وړ دے؛ او بيا په~\LR{$\alpha$} باندې له استقرا څخه چې \emph{هر} ترتيبي عدد له
\emph{هر} ترتيبي عدد سره د پرتله کېدو وړ دے، لکه چې پکار دي.
بس دا فرض کول کافي دي چې \LR{$\alpha \notin \beta$} او \LR{$\beta \notin
\alpha$} دي، او وښيو چې \LR{$\alpha = \beta$} دے.
\olpseganchor{OLP-0553-B015}{end}

\olpseganchor{OLP-0553-B016}{start}
د دې د ښودلو دپاره چې \LR{$\alpha \subseteq \beta$} دے، \LR{$\gamma \in \alpha$}
ثابت وټاکئ؛ د \olref{ordmemberord} له مخې دا ترتيبي عدد دے. نو د لومړۍ
استقرا د فرض له مخې \LR{$\gamma$} له \LR{$\beta$} سره د پرتله کېدو وړ دے.
خو که يا \LR{$\gamma
= \beta$} يا \LR{$\beta \in \gamma$} وي، نو \LR{$\beta \in \alpha$} دے
(که اړتيا وي، د \LR{$\alpha$} له تعديت څخه کار اخلو)، چې زموږ د فرض خلاف دے؛
نو \LR{$\gamma \in \beta$} دے. په تعميمولو سره \LR{$\alpha \subseteq
\beta$} دے.
\olpseganchor{OLP-0553-B016}{end}

\olpseganchor{OLP-0553-B017}{start}
په عين ورته استدلال، د دويمې استقرا د فرض په کارولو، ښيو چې
\LR{$\beta \subseteq \alpha$} دے. نو \LR{$\alpha = \beta$} دے.
\end{proof}\noindent له دې امله، کله ناکله به د \LR{$\alpha <\beta$}
نښه د \LR{$\alpha \in \beta$} پر ځای وليکو، ځکه چې \LR{$\in$} د ترتيب د اړيکې
په شان کار کوي. له اشناوالي پرته کومه ژوره وجه نۀ لري؛ او هم دا چې
\LR{$\alpha \leq \beta$} ليکل تر \LR{$\alpha \in
\beta \lor \alpha = \beta$} اسانه دي.\footnote{موږ \LR{$\alpha
\mathrel{\underline{\in}} \beta$} ليکلے شو؛ خو دا به بالکل غير معياري وي.}
\olpseganchor{OLP-0553-B017}{end}

\olpseganchor{OLP-0553-B018}{start}
زموږ د وروستيو نتيجو دوه لنډې پايلې دا دي؛ لومړۍ پايله نوې نښې کاروي:
\olpseganchor{OLP-0553-B018}{end}

\olpseganchor{OLP-0553-B019}{start}
\begin{cor}\ollabel{ordordered}
که \LR{$\exists \alpha\phi(\alpha)$} وي، نو \LR{$\exists \alpha(\phi(\alpha)
\land \forall \beta(\phi(\beta) \lif \alpha \leq \beta))$} دے.
همدارنګه، د هر درې ترتيبي عددونو \LR{$\alpha, \beta, \gamma$} دپاره دواړه
\LR{$\alpha
\notin \alpha$} او \LR{$\alpha \in \beta \in \gamma \lif \alpha \in
\gamma$} دي.
\end{cor}
\olpseganchor{OLP-0553-B019}{end}

\olpseganchor{OLP-0553-B020}{start}
\begin{proof}
د \olref[wo]{wo:strictorder} په شان، د درې ګونتوب له مخې
مينيمال شاهد له هر بل شاهد وړاندې راځي۔ هر ترتيبي عدد خپل
غړے کېدلے نۀ شي، ځکه چې په خپل غړيتوب ښه مرتب دے او
په خپل ځان کښې خپل غړيتوب به د هماغه سخت ترتيب خلاف وي۔
د غړيتوب تعديت د وروستي ترتيبي عدد له متعدي‌والي راځي۔
\end{proof}
\olpseganchor{OLP-0553-B020}{end}

\olpseganchor{OLP-0553-B021}{start}
\begin{prob}
د \olref[sth][ordinals][basic]{ordtrichotomy} په ثبوت کښې «په عين ورته
استدلال» بشپړ کړئ.
\end{prob}
\olpseganchor{OLP-0553-B021}{end}

\olpseganchor{OLP-0553-B022}{start}
\begin{cor}\ollabel{corordtransitiveord}
\LR{$A$} ترتيبي عدد دے، که او يوازې که \LR{$A$} د ترتيبي عددونو متعدي سټ وي.
\end{cor}
\olpseganchor{OLP-0553-B022}{end}

\olpseganchor{OLP-0553-B023}{start}
\begin{proof}
\emph{له چپ نه ښي خوا ته.} د \olref{ordmemberord} له مخې.
\emph{له ښي نه چپ خوا ته.} که \LR{$A$} د ترتيبي عددونو متعدي سټ وي، نو
\LR{$\in$} د \LR{$A$} ښه ترتيب د \olref{ordinductionschema} او
\olref{ordtrichotomy} له مخې ورکوي.
\end{proof}
\olpseganchor{OLP-0553-B023}{end}

\olpseganchor{OLP-0553-B024}{start}
اوس، موږ د
\olref{ordinductionschema} او \olref{ordtrichotomy} معنا داسې بيان کړه چې
\LR{$\in$} د ترتيبي عددونو ښه ترتيب ورکوي. خو د لاندې نتيجې له امله بايد
د دې ډول ادعا په اړه \emph{ډېر احتياط} وکړو:
\olpseganchor{OLP-0553-B024}{end}

\olpseganchor{OLP-0553-B025}{start}
\begin{thm}[د بورالي--فورتي تناقض]\ollabel{buraliforti}
د ټولو ترتيبي عددونو سټ نشته.
\end{thm}
\olpseganchor{OLP-0553-B025}{end}

\olpseganchor{OLP-0553-B026}{start}
\begin{proof}
د تناقض د راوستلو دپاره فرض کړئ \LR{$O$} د ټولو ترتيبي عددونو سټ دے.
که \LR{$\alpha \in
\beta \in O$} وي، نو د \olref{ordmemberord} له مخې \LR{$\alpha$} ترتيبي عدد دے،
نو \LR{$\alpha \in O$} دے. په دې توګه \LR{$O$} متعدي دے، او له
\olref{corordtransitiveord} څخه \LR{$O$} ترتيبي عدد دے. نو \LR{$O \in O$} دے، چې له
\olref{ordordered} سره په تناقض کښې دے.
\end{proof}
\noindent
دا نتيجه د \mbox{\LR{\latinfont \citeauthor{Burali-Forti1897}}} په نوم نومول شوې ده. خو کانتور په
۱۸۹۹ کښې، ډېډېکنډ ته په يوه ليک کښې، په لومړي ځل په روښانه توګه وليدل
چې د ټولو ترتيبي عددونو د يوه سټ فرض کول \emph{تناقض} راولاړوي.
\begin{explain}
  د فان هاينورټ د تاريخي شننې له مخې، بورالي--فورتي له دې تناقض
  څخه د ترتيبي عددونو د مقايسې په باب نتيجه اخيسته۔ ده استدلال کاوه
  چې که طبيعي ترتيب ټولوالے ولري، تناقض راځي، نو دغه ترتيب
  \emph{جزوي} پاتې کېږي۔ \mbox{\LR{\latinfont \citep[p.~105]{Heijenoort1967}}}
\end{explain}
د بورالي--فورتي تېروتنه يوې خوا ته پرېږدو؛ پورته بحث داسې لنډولے شو:
هر ترتيبي عدد د غړيتوب له مخې ښه مرتب دے۔ ټول ترتيبي
عددونه يوه حقيقي طبقه جوړوي۔ د درې ګونتوب له مخې يې هر دوه
غړي د غړيتوب پر بنسټ پرتله کېدلے شي، او د استقرا شېما د
هرې ناتشې تعريف وړ فرعي طبقې تر ټولو کوچنے غړے ورکوي۔
\olpseganchor{OLP-0553-B026}{end}

\olpseganchor{OLP-0553-B027}{start}
د دې بحث په پای کښې، د ترتيبي عددونو ځينې نور بنسټيز خاصيتونه وړاندې
کوو چې له \olref{ordinductionschema} او
\olref{ordtrichotomy} څخه نتيجه اخلو.
\olpseganchor{OLP-0553-B027}{end}

\olpseganchor{OLP-0553-B028}{start}
\begin{prop}
د ترتيبي عددونو هره په کلکه نزولي لړۍ متناهي ده.
\end{prop}
\olpseganchor{OLP-0553-B028}{end}

\olpseganchor{OLP-0553-B029}{start}
\begin{proof}
د ترتيبي عددونو هره نامتناهي په کلکه نزولي لړۍ \LR{$\alpha_0 > \alpha_1 > \alpha_2 > \ldots$}
د \LR{$<$} له مخې مينيمال غړے نۀ لري، چې له \olref{ordinductionschema} سره
په تناقض کښې ده.
\end{proof}
\olpseganchor{OLP-0553-B029}{end}

\olpseganchor{OLP-0553-B030}{start}
\begin{prop}\ollabel{ordinalsaresubsets}
\LR{$\alpha \subseteq \beta \lor \beta \subseteq \alpha$}، د هر دوو ترتيبي
عددونو \LR{$\alpha, \beta$} دپاره.
\end{prop}
\olpseganchor{OLP-0553-B030}{end}

\olpseganchor{OLP-0553-B031}{start}
\begin{proof}
که \LR{$\alpha \in \beta$} وي، نو \LR{$\alpha \subseteq \beta$} دے، ځکه چې \LR{$\beta$}
متعدي دے. همدارنګه، که \LR{$\beta \in \alpha$} وي، نو \LR{$\beta \subseteq
\alpha$} دے. او که \LR{$\alpha = \beta$} وي، نو \LR{$\alpha \subseteq \beta$} او
\LR{$\beta \subseteq \alpha$} دي. نو د \olref{ordtrichotomy} له مخې مقصد ترلاسه شو.
\end{proof}
\olpseganchor{OLP-0553-B031}{end}

\olpseganchor{OLP-0553-B032}{start}
\begin{prop}\ollabel{ordisoidentity}
\LR{$\alpha = \beta$}، که او يوازې که \LR{$\ordeq{\alpha}{\beta}$} وي، د هر دوو
ترتيبي عددونو \LR{$\alpha, \beta$} دپاره.
\end{prop}
\olpseganchor{OLP-0553-B032}{end}

\olpseganchor{OLP-0553-B033}{start}
\begin{proof}
ترتيبي عددونه ښه ترتيبونه دي؛ نو دا له درې ګونتوب
(\olref[basic]{ordtrichotomy}) او
\olref[iso]{wellordnotinitial} څخه سمدستي نتيجه اخلو.
\end{proof}
\olpseganchor{OLP-0553-B033}{end}

\olpseganchor{OLP-0553-B034}{start}
\begin{prob}\ollabel{probunionordinalsordinal}
ثابت کړئ چې، که د \LR{$X$} هر غړے ترتيبي عدد وي، نو \LR{$\bigcup X$} ترتيبي عدد دے.
\end{prob}
\olpseganchor{OLP-0553-B034}{end}

\olpseganchor{OLP-0554-B003}{start}
\olfileid{sth}{ordinals}{replacement}
\olsection{تعويض}
\olpseganchor{OLP-0554-B003}{end}

\olpseganchor{OLP-0554-B004}{start}
په \olref[sth][ordinals][vn]{sec} کښې مو د ترتيبي عددونو د معرفي کولو
وجه دا وړاندې کړه چې د ترتيب د ډولونو په توګه، يعنې د ښه ترتيبونو د
کانوني استازو په توګه، ئې کارولے شو. د دې د عملي کېدو دپاره به ثابتول
پکار وي چې \emph{هر ښه ترتيب له کوم ترتيبي عدد سره آيزومورف دے}.
په دې سره \LR{$\ordtype{A, <}$} د هغه ترتيبي عدد \LR{$\alpha$} په توګه تعريفولے
شو چې \LR{$\tuple{A, <} \isomorphic \alpha$} وي.
\olpseganchor{OLP-0554-B004}{end}

\olpseganchor{OLP-0554-B005}{start}
له بده مرغه، يوازې د هغو بديهي اصلونو په کارولو چې تر اوسه مو معرفي
کړي دي، مطلوبه نتيجه \emph{نۀ شو} ثابتولے. (وجه به ئې په
\olref[replacement][strength]{sec} کښې ووينو؛ خو اوس مطلب دا دے چې
نۀ ئې شو ثابتولے.) نوے فکر پکار دے، او هغه دا دے:
\olpseganchor{OLP-0554-B005}{end}

\olpseganchor{OLP-0554-B006}{start}
\begin{axiom}[د تعويض شېما]
د هر فارمول \LR{$\phi(x, y)$} دپاره، لاندې خبره بديهي اصل دے:
\begin{quote}
	د هر \LR{$A$} دپاره، که \LR{$(\forall x \in A)\lexists![y][\phi(x,y)]$} وي، نو
	\LR{$\Setabs{y}{(\exists x \in A)\phi(x,y)}$} شتون لري.
\end{quote}
\end{axiom}
\noindent
د بېلولو په شان، دا يوه شېما ده: د بې شمېره بېلو \LR{$\phi$} ګانو د هر يوه
دپاره يو بديهي اصل ورکوي، نو نامتناهي ډېر بديهي اصلونه ترې جوړېږي.
په هماغه معنا داسې هم ليکل کېدلے شي، او عموماً همداسې ليکل کېږي:
\olpseganchor{OLP-0554-B006}{end}

\olpseganchor{OLP-0554-B007}{start}
\begin{defish}
د هر داسې فارمول \LR{$\phi(x,y)$} دپاره چې «\LR{$B$}» پکښې نۀ راځي، لاندې خبره بديهي اصل دے:
\[
\forall A[(\forall x \in A)\lexists![y][\phi(x,y)] \lif \exists B\forall y (y \in B \liff (\exists x \in A)\phi(x,y))]
\]
\end{defish}
\olpseganchor{OLP-0554-B007}{end}

\olpseganchor{OLP-0554-B008}{start}
خو په لومړي ځل دا فارمول څه پېچلے ښکاري. د تعويض لاندې لنډه پايله ښايي
زموږ مطلوب شهودي تصور په \emph{لا روښانه} توګه بيان کړي:\footnote{ترم هغه
عبارت دے چې د خپلو ازادو متغيرونو په هر بشپړ ټاکلو کښې عين يو شے ټاکي.
د بېلګې په توګه، «\LR{$\emptyset$}» هغه ترم دے چې تش سټ ټاکي؛ «\LR{$\{x\}$}» هغه
ترم دے چې د \LR{$x$} يوغړيز سټ ټاکي، هر څه چې \LR{$x$} وي؛ او «\LR{$x \cup y$}» هغه
ترم دے چې د \LR{$x$} او \LR{$y$} اتحاد ټاکي، هر څه چې دا دواړه وي.}
\olpseganchor{OLP-0554-B008}{end}

\olpseganchor{OLP-0554-B009}{start}
\begin{cor}
د هر ترم \LR{$\tau(x)$} او هر سټ \LR{$A$} دپاره، دا سټ شتون لري:
\[
	\Setabs{\tau(x)}{x \in A} = \Setabs{y}{(\exists x \in A)y = \tau(x)}.
\]
\end{cor}
\olpseganchor{OLP-0554-B009}{end}

\olpseganchor{OLP-0554-B010}{start}
\begin{proof}
ځکه چې \LR{$\tau$} يو \emph{ترم} دے، \LR{$\forall x \lexists![y][\tau(x) = y]$} دے.
په تېره بيا، \LR{$(\forall x \in A)\lexists![y][\tau(x) = y]$} دے. نو
\LR{$\Setabs{y}{(\exists x \in A)\tau(x) = y}$} د تعويض له مخې شتون لري.
\end{proof}
\noindent
له دې ښکاري چې «تعويض» د دې بديهي اصل دپاره ښه نوم دے: که سټ \LR{$A$} درکړے
شي، نو د \LR{$\Setabs{\tau(x)}{x \in A}$} په توګه نوے سټ جوړولے شئ، په دې
چې د \LR{$A$} هر غړے د~\LR{$\tau$} له مخې د هغه په تصوير بدل کړئ. په رښتيا، د
تابعې له مخې د سټ د تصوير د نښو په تعقيب،
\LR{$\funimage{\tau}{A}$} د \LR{$\Setabs{\tau(x)}{x \in A}$} دپاره ليکلے شو.
\olpseganchor{OLP-0554-B010}{end}

\olpseganchor{OLP-0554-B011}{start}
خو دا ډېره مهمه ده چې \LR{$\tau$} يو \emph{ترم} دے. ضرور نۀ ده چې د
\olref[sfr][fun][rel]{sec} په معنا د کومې \emph{تابعې} نوم وي، يعنې د
مرتبو جوړو د کوم ټاکلي سټ نوم. ځکه که \LR{$f$} په دې معنا تابعه وي او ټول ورودي غړي يې د تعريف په ساحه
کښې وي، نو سټ
\LR{$\funimage{f}{A} = \Setabs{f(x)}{x \in
A}$} د \LR{$\ran{f}$} يو ټاکلے فرعي سټ دے، او د~\LR{$\Zminus$} بديهي
اصلونه لا له وړاندې ئې شتون تضمينوي.\footnote{يوازې
\LR{$\Setabs{y \in \bigcup \bigcup f}{(\exists x \in A)y =
f(x)}$} په پام کښې ونيسئ.} خو تعويض زموږ بديهي اصلونو ته يوه
\emph{قوي} اضافه ده، لکه چې په \olref[sth][replacement][]{chap} کښې به ووينو۔
دلته د فارمول له مخې ټاکل کېدونکي حاصلونه راټولوو؛ د قاعدې
د کارولو لپاره د هغې ټوليز ګراف ته اړتيا نشته۔
\olpseganchor{OLP-0554-B011}{end}

\olpseganchor{OLP-0555-B002}{start}
\olfileid{sth}{ordinals}{zfm}	
\olsection{\texorpdfstring{\LR{$\ZFminus$}}{ZF⁻}: يو مهم پړاو}
\olpseganchor{OLP-0555-B002}{end}

\olpseganchor{OLP-0555-B003}{start}
دا پوښتنه چې تعويض څنګه توجيه کړو، که ئې توجيه کېدلے شي، اسانه نۀ ده.
نو دا به د \olref[replacement][]{chap} دپاره پرېږدو. خو د تعويض په اضافه
کولو سره يو بل مهم پړاو ته رسېدلي يو. اوس زموږ سره د \LR{$\ZFminus$} د
نظريې ټول اړين بديهي اصلونه شته. په تفصيل سره:
\olpseganchor{OLP-0555-B003}{end}

\olpseganchor{OLP-0555-B004}{start}
\begin{defn}
نظريه \LR{$\ZFminus$} دا بديهي اصلونه لري: د غړو له مخې برابري، اتحاد،
جوړې، د فرعي سټونو سټونه، نامتناهيت، او د بېلونې او تعويض د شېماګانو
ټولې بېلګې. په بله وينا، \LR{$\ZFminus$} په \LR{$\Zminus$} کښې تعويض اضافه کوي.
\end{defn}
\noindent
دا د \emph{زرمېلو--فرېنکل} د سټونو د نظريې نوم دے، چې يو څه ترې
\emph{کم} دي؛ هغې خبرې ته به وروسته راشو. فرېنکل ته دا وياړ ځکه ورکول
کېږي چې په \mbox{\LR{\latinfont \citeyear{Fraenkel1922}}} کښې د تعويض فورمولبندي ده ته منسوبه
ده، که څه هم لومړۍ دقيقه فورمولبندي د \mbox{\LR{\latinfont \citet{Skolem1922}}} وه.
\olpseganchor{OLP-0555-B004}{end}

\olpseganchor{OLP-0556-B003}{start}
\olfileid{sth}{ordinals}{ordtype}
\olsection{ترتيبي عددونه د ترتيب د ډولونو په توګه}
\olpseganchor{OLP-0556-B003}{end}

\olpseganchor{OLP-0556-B004}{start}
اوس چې تعويض لرو او په \LR{$\ZFminus$} کښې کار کوو، بالاخره هغه نتيجه
ثابتولے شو چې مقصد مو وه:
\olpseganchor{OLP-0556-B004}{end}

\olpseganchor{OLP-0556-B005}{start}
\begin{thm}\ollabel{thmOrdinalRepresentation}
هر ښه ترتيب له يوه يوازيني ترتيبي عدد سره آيزومورف دے.
\end{thm}
\olpseganchor{OLP-0556-B005}{end}

\olpseganchor{OLP-0556-B006}{start}
\begin{proof}
فرض کړئ \LR{$\tuple{A, <}$} ښه ترتيب دے. د
\olref[basic]{ordisoidentity} له مخې، له زيات نه زيات يوه ترتيبي عدد
سره آيزومورف دے. نو د تناقض د راوستلو دپاره فرض کړئ \LR{$\tuple{A, <}$} له
\emph{هېڅ} ترتيبي عدد سره آيزومورف نۀ دے. لومړے به «\LR{$\tuple{A, <}$} تر
ممکن حده کوچنے کړو». په تفصيل سره: که کومه خاصه پيلنۍ برخه \LR{$\tuple{A_a,
<_a}$} له هېڅ ترتيبي عدد سره آيزومورف نۀ وي، نو يو تر ټولو کوچنے \LR{$a \in A$}
دا خاصيت لري؛ بيا \LR{$B = A_a$} او \LR{$\mathord{\lessdot} =
\mathord{<_a}$} واخلو. که داسې نۀ وي، نو \LR{$B = A$} او \LR{$\mathord{\lessdot} =
\mathord{<}$} واخلو.
\olpseganchor{OLP-0556-B006}{end}

\olpseganchor{OLP-0556-B007}{start}
له تعريفه، د \LR{$B$} هره خاصه پيلنۍ برخه له کوم ترتيبي عدد سره آيزومورف
ده، او د پورته بحث له مخې هغه يوازينے دے۔ د تعويض لپاره
هر غړي ته د هغه د پيلنۍ برخې يوازينے ترتيبي عدد ټاکو،
او بيا دغه دواړه د لاندې ترتيب له مخې په يوه مرتبه جوړه
کښې ليکو۔ په دې توګه د ټولو جوړو سټ شتون لري:
\[
	f = \Setabs{\tuple{\beta, b}}{b \in B\text{\RL{ او }}\ordeq{\beta}{\tuple{B_b, \lessdot_b}}}
\]
که د دوو غړو پيلنۍ برخې عين ترتيبي عدد ولري، نو يوه
له بلې سره آيزومورف ده۔ دوه بېل غړي د ښه ترتيب له مخې يو
تر بل وړاندې وي، او د مخکيني غړي پيلنۍ برخه د بل يوه
خاصه پيلنۍ برخه ده۔ له دغسې برخې سره آيزومورفيزم ناشونی
دے۔ نو دا اړيکه تابع ده۔ د تناقض د بشپړولو دپاره به وښيو
چې \LR{$f$} يو آيزومورفيزم \LR{$\alpha \to B$} دے، د کوم ترتيبي عدد
\LR{$\alpha$} دپاره.
\olpseganchor{OLP-0556-B007}{end}

\olpseganchor{OLP-0556-B008}{start}
ښکاره ده چې \LR{$\ran{f}
= B$} دے۔ د \olref[iso]{lemordsegments} د استدلال په شان، \LR{$f$} ترتيب ساتي۔
د دوو بېلو ترتيبي عددونو يو د بل غړے دے۔ کوچنے ترتيبي
عدد د لوے يوه پيلنۍ برخه ده۔ هغه آيزومورفيزم چې لوے
ترتيبي عدد د خپل اړوند غړي له پيلنۍ برخې سره تړي، کوچنے
عدد يوې وړې پيلنۍ برخې ته وړي۔ د پيلنۍ برخې د ډول
يوازينيت ښيي چې هماغه برخه د کوچني عدد له اړوند غړي
سره تړلې ده۔ په دې توګه \LR{$\gamma \in \beta$} هله او يوازې هله
چې \LR{$f(\gamma) \lessdot f(\beta)$} وي۔
د دې د ښودلو دپاره چې \LR{$\dom{f}$} ترتيبي عدد دے، د
\olref[basic]{corordtransitiveord} له مخې دا ښودل کافي دي چې
\LR{$\dom{f}$} متعدي دے. نو \LR{$\beta \in \dom{f}$} ثابت وټاکئ، يعنې
\LR{$\ordeq{\beta}{\tuple{B_b, \lessdot_b}}$} د کوم \LR{$b$} دپاره.
که \LR{$\gamma \in
\beta$} وي، نو \LR{$\gamma \in \dom{f}$} د
\olref[iso]{wellordinitialsegment} له مخې دے؛ په تعميمولو سره، \LR{$\beta \subseteq
\dom{f}$} دے۔ نو د تابعې ساحه يو ترتيبي عدد دے، او د
ترتيب دوه اړخيز ساتل او پر ټوله توب ترې آيزومورفيزم
جوړوي۔ دا د ټاکلي ښه ترتيب د نۀ استازيتوب له فرض سره ټکر لري۔
\end{proof}
\olpseganchor{OLP-0556-B008}{end}

\olpseganchor{OLP-0556-B009}{start}
دا نتيجه د لاندې تعريف جواز راکوي، چې له \olref[vn]{sec} راهيسې مو
وړاندې کول غوښتل:
\olpseganchor{OLP-0556-B009}{end}

\olpseganchor{OLP-0556-B010}{start}
\begin{defn}
که \LR{$\tuple{A, <}$} ښه ترتيب وي، نو د هغه د ترتيب ډول،
\LR{$\ordtype{A, <}$}، هغه يوازينے ترتيبي عدد \LR{$\alpha$} دے چې
\LR{$\ordeq{\tuple{A, < }}{\alpha}$} وي.
\end{defn}
\olpseganchor{OLP-0556-B010}{end}

\olpseganchor{OLP-0556-B011}{start}
همدارنګه، دا تعريف د دوو ښو اصلونو جواز راکوي:
\olpseganchor{OLP-0556-B011}{end}

\olpseganchor{OLP-0556-B012}{start}
\begin{cor}\ollabel{ordtypesworklikeyouwant}
که \LR{$\tuple{A, <}$} او \LR{$\tuple{B, \lessdot}$} ښه ترتيبونه وي:
\begin{align*}
	\ordtype{A, <} = \ordtype{B, \lessdot}&\text{\RL{ هله او يوازې هله چې }}\ordeq{\tuple{A, <}}{\tuple{B, \lessdot}}\\
	\ordtype{A, <} \in \ordtype{B, \lessdot}&\text{\RL{ هله او يوازې هله چې }}\ordeq{\tuple{A, <}}{\tuple{B_b, \lessdot_b}}\text{\RL{ د يو مناسب }}b \in B
\end{align*}
\end{cor}
\olpseganchor{OLP-0556-B012}{end}

\olpseganchor{OLP-0556-B013}{start}
\begin{proof}
برابري د \olref[basic]{ordisoidentity} له مخې ده. د دويمې ادعا د
ثابتولو دپاره، \LR{$\ordtype{A, <} = \alpha$} او \LR{$\ordtype{B,
\lessdot} = \beta$} واخلو، او \LR{$f \colon \beta \to B$}
د ورکړل شويو ترتيبونو آيزومورفيزم وټاکو۔ که \LR{$\alpha \in \beta$}
وي، نو د تابعې محدوديت دا آيزومورفيزم ورکوي:
\begin{align*}
	&\funrestrictionto{f}{\alpha}\colon \alpha \to B_{f(\alpha)}.\\
	&\ordeq{\tuple{A, <}}{\tuple{B_{f(\alpha)}, \lessdot_{f(\alpha)}}}.
\end{align*}
دويمه کرښه د ترتيب د ډول له تعريفه راځي۔ په بل لوري کښې،
فرض کړئ چې \LR{$\ordeq{\tuple{A, <}}{\tuple{B_b, \lessdot_b}}$}
د کوم \LR{$b \in B$} دپاره۔ د هغه مخکنی انځور په \LR{$\beta$} کښې دے،
او د تابعې محدوديت يې له هماغې پيلنۍ برخې سره آيزومورف
کوي۔ دا مخکنی انځور ترتيبي عدد دے، او د ترتيبي استازي
د يوازينيت له مخې هماغه \LR{$\alpha$} دے۔ د دې دوو لورو لپاره
\olref[iso]{ordisounique}،
\olref[iso]{wellordinitialsegment} او
\olref[basic]{ordissetofsmallerord} کاروو۔
\end{proof}
\olpseganchor{OLP-0556-B013}{end}

\olpseganchor{OLP-0557-B003}{start}
\olfileid{sth}{ordinals}{opps}
\olsection{تالي او حدي ترتيبي عددونه}
\olpseganchor{OLP-0557-B003}{end}

\olpseganchor{OLP-0557-B004}{start}
په ورپسې څو څپرکو کښې به ترتيبي عددونه ډېر کاروو. نو ځينې ساده
مفهومونه معرفي کول به ګټور وي.
\olpseganchor{OLP-0557-B004}{end}

\olpseganchor{OLP-0557-B005}{start}
\begin{defn}
د هر ترتيبي عدد \LR{$\alpha$} \emph{تالي} \LR{$\ordsucc{\alpha}
=\alpha \cup \{\alpha\}$} دے. وايو چې \LR{$\alpha$} \emph{تالي} ترتيبي عدد دے،
که \LR{$\ordsucc{\beta} = \alpha$} د کوم ترتيبي عدد~\LR{$\beta$} دپاره وي.
وايو چې \LR{$\alpha$} \emph{حدي} ترتيبي عدد دے، که او يوازې که \LR{$\alpha$}
نۀ تش وي او نۀ تالي ترتيبي عدد وي.
\end{defn}
\noindent
لاندې نتيجه ښيي چې دا د \emph{تالي} مناسب مفهوم دے:
\olpseganchor{OLP-0557-B005}{end}

\olpseganchor{OLP-0557-B006}{start}
\begin{prop}
د هر ترتيبي عدد \LR{$\alpha$} دپاره:
\begin{enumerate}
	\item \LR{$\alpha \in \ordsucc{\alpha}$}؛
	\item \LR{$\ordsucc{\alpha}$} ترتيبي عدد دے؛
	\item داسې ترتيبي عدد \LR{$\beta$} نشته چې \LR{$\alpha \in \beta \in
	\ordsucc{\alpha}$} وي.
	\end{enumerate}
\end{prop}
\olpseganchor{OLP-0557-B006}{end}

\olpseganchor{OLP-0557-B007}{start}
\begin{proof}
ښکاره ده چې \LR{$\alpha \in \alpha \cup \{\alpha\} = \ordsucc{\alpha}$} دے.
همدارنګه، \LR{$\ordsucc{\alpha}$} د ترتيبي عددونو متعدي سټ دے۔
د نوي غړي ټول غړي پخپله په پخواني ترتيبي عدد کښې دي، او د
پخوانيو غړو غړي هلته د تعديت له مخې شته۔ نو د
\olref[basic]{corordtransitiveord} له مخې ترتيبي عدد دے۔
او \LR{$\alpha \in \beta \in \ordsucc{\alpha}$} ناشوني دي، ځکه چې بيا به يا
\LR{$\beta \in \alpha$} يا \LR{$\beta = \alpha$} وي، چې له
\olref[basic]{ordordered} سره په تناقض کښې دي.
\end{proof}
\noindent
دا د نامتناهيت هاخوا استقرا د ثبوت يوې بلې بڼې ته هم جواز ورکوي:
\olpseganchor{OLP-0557-B007}{end}

\olpseganchor{OLP-0557-B008}{start}
\begin{thm}[د نامتناهيت هاخوا ساده استقرا]\ollabel{simpletransrecursion}
فرض کړئ \LR{$\phi(x)$} داسې فارمول دے چې:
\begin{enumerate}
	\item \LR{$\phi(\emptyset)$}؛ او
	\item د هر ترتيبي عدد \LR{$\alpha$} دپاره، که \LR{$\phi(\alpha)$} وي، نو
	\LR{$\phi(\ordsucc{\alpha})$} وي؛ او
	\item که \LR{$\alpha$} حدي ترتيبي عدد وي او \LR{$(\forall \beta \in
	\alpha)\phi(\beta)$} وي، نو \LR{$\phi(\alpha)$} وي.
\end{enumerate}
بيا \LR{$\forall \alpha \phi(\alpha)$} دے.
\end{thm}
\olpseganchor{OLP-0557-B008}{end}

\olpseganchor{OLP-0557-B009}{start}
\begin{proof}
نقيض عکس ثابتوو. نو فرض کړئ کوم ترتيبي عدد \LR{$\lnot\phi$} لري؛
\LR{$\gamma$} تر ټولو کوچنے داسې ترتيبي عدد وټاکئ. بيا يا
\LR{$\gamma = \emptyset$} دے، يا \LR{$\gamma = \ordsucc{\alpha}$} د کوم \LR{$\alpha$}
دپاره دے چې \LR{$\phi(\alpha)$} لري؛ يا \LR{$\gamma$} حدي ترتيبي عدد دے او
\LR{$(\forall
\beta \in \gamma)\phi(\beta)$} دي۔
په درې واړو حالاتو کښې اړوند فرض ښيي چې شرط پر همدې
ترتيبي عدد صدق کوي۔ دا د مخالف مثال له ټاکنې سره ټکر لري۔
\end{proof}
\noindent
يوه وروستۍ نښه به وروسته ګټوره ثابته شي:
\olpseganchor{OLP-0557-B009}{end}

\olpseganchor{OLP-0557-B010}{start}
\begin{defn}\ollabel{defsupstrict}
که \LR{$X$} د ترتيبي عددونو سټ وي، نو \LR{$\supstrict(X) = \bigcup_{\alpha
\in X} \ordsucc{\alpha}$} دے.
\end{defn}
\noindent
دلته «lsub» د «تر ټولو لږ په کلکه پاسني حد» نښه ده.\footnote{ځينې کتابونه
دې ته «\LR{$\text{\RL{sup}}(X)$}» کاروي. خو نور کتابونه «\LR{$\text{\RL{sup}}(X)$}» د تر ټولو
لږ \emph{غير اکيد} پاسني حد دپاره کاروي، يعنې يوازې د \LR{$\bigcup X$} دپاره.
که \LR{$X$} تر ټولو لوے غړے \LR{$\alpha$} ولري، نو دا دوه مفهومونه بېلېږي:
تر ټولو لږ \emph{اکيد} پاسنے حد \LR{$\ordsucc{\alpha}$} دے، خو تر ټولو لږ
\emph{غير اکيد} پاسنے حد يوازې \LR{$\alpha$} دے.} لاندې نتيجه دا څرګندوي:
\olpseganchor{OLP-0557-B010}{end}

\olpseganchor{OLP-0557-B011}{start}
\begin{prop}
که \LR{$X$} د ترتيبي عددونو سټ وي، نو \LR{$\supstrict(X)$} هغه تر ټولو کوچنے
ترتيبي عدد دے چې د \LR{$X$} له هر ترتيبي عدد څخه لوے وي.
\end{prop}
\olpseganchor{OLP-0557-B011}{end}

\olpseganchor{OLP-0557-B012}{start}
\begin{proof}
د تعويض له مخې \LR{$Y = \Setabs{\ordsucc{\alpha}}{\alpha \in X}$} واخلو، نو
\LR{$\supstrict(X) = \bigcup Y$} دے. ځکه چې ترتيبي عددونه متعدي دي او د هر
ترتيبي عدد هر غړے ترتيبي عدد دے، \LR{$\supstrict(X)$} د ترتيبي عددونو متعدي
سټ دے، نو د \olref[basic]{corordtransitiveord} له مخې ترتيبي عدد دے.
\olpseganchor{OLP-0557-B012}{end}

\olpseganchor{OLP-0557-B013}{start}
که \LR{$\alpha \in X$} وي، نو \LR{$\ordsucc{\alpha} \in Y$} دے، نو \LR{$\ordsucc{\alpha}
\subseteq \bigcup Y = \supstrict(X)$} دے، او له دې \LR{$\alpha \in
\supstrict(X)$} دے. نو \LR{$\supstrict(X)$} د \LR{$X$} له هر ترتيبي عدد څخه په کلکه
لوے دے.
\olpseganchor{OLP-0557-B013}{end}

\olpseganchor{OLP-0557-B014}{start}
په معکوسه توګه، که \LR{$\alpha \in \supstrict(X)$} وي، نو \LR{$\alpha \in
\ordsucc{\beta} \in Y$} د کوم \LR{$\beta \in X$} دپاره دي، نو \LR{$\alpha \leq
\beta \in X$} دي۔ د راکړل شوي سټ هر اکيد پاسنے حد د
همدې شاهد نه لوے دے، نو د اتحاد هر غړے هم تر هغه حده
کوچنے دے۔ ځکه چې ترتيبي عددونه د فرعي سټ له مخې هم
پرتله کېدلے شي، \LR{$\supstrict(X)$} د \LR{$X$} تر ټولو \emph{لږ}
اکيد پاسنے حد دے۔
\end{proof}
\olpseganchor{OLP-0557-B014}{end}

\olpseganchor{OLP-0558-B004}{start}
\olchapter{sth}{spine}{پړاوونه او رتبې}
\olpseganchor{OLP-0558-B004}{end}

\olpseganchor{OLP-0559-B003}{start}
\olfileid{sth}{spine}{valpha}
\olpseganchor{OLP-0559-B003}{end}

\olpseganchor{OLP-0559-B004}{start}
\olsection{پړاوونه د \texorpdfstring{\LR{$V_\alpha$}}{V\_α} ګانو په توګه تعريفول}
\olpseganchor{OLP-0559-B004}{end}

\olpseganchor{OLP-0559-B005}{start}
په \olref[sth][ordinals][]{chap} کښې مو ښه ترتيبونه او د فون نويمن
ترتيبي عددونه تعريف کړل. په دې څپرکي کښې به دا د سټونو د مراتبو د لړۍ
\emph{پخپله} د مشخصولو دپاره وکاروو. د دې دپاره په ياد کړئ چې په
\olref[sth][ordinals][opps]{sec} کښې مو تالي او حدي ترتيبي عددونه تعريف
کړل. همدا مفهومونه په لاندې تعريف کښې کاروو:
\olpseganchor{OLP-0559-B005}{end}

\olpseganchor{OLP-0559-B006}{start}
\begin{defn}\ollabel{defValphas}
	\begin{align*}
	V_\emptyset &\defis \emptyset\\
	V_{\ordsucc{\alpha}} &\defis \Pow{V_\alpha} & &
	\text{\RL{د هر ترتيبي عدد دپاره }}\alpha\\
	V_{\alpha} &\defis \bigcup_{\gamma < \alpha} V_\gamma & &
	\text{\RL{کله چې }}\alpha\text{\RL{ حدي ترتيبي عدد وي}}
\end{align*}
\end{defn}
\noindent
دا به په ترتيبي عددونو باندې د \emph{نامتناهيت هاخوا بازګښت} له لارې
تعريف وي. په دې اړه ئې بايد په طبيعي عددونو باندې د تابعو له بازګشتي
تعريفونو سره پرتله کړو.\footnote{په
\olref[sfr][infinite][dedekind]{sec} کښې د جمع، ضرب او د قوې تعريفونه
وګورئ.} د طبيعي عددونو په شان يو بنسټيز حالت او تالي حالتونه تعريفوو؛
خو د ترتيبي عددونو په کارولو کښې د \emph{حدي} حالتونو چلند هم بيانول
پکار دي.
\olpseganchor{OLP-0559-B006}{end}

\olpseganchor{OLP-0559-B007}{start}
د \LR{$V_\alpha$} ګانو دا تعريف به يو مهم پړاو وي. موږ د سټونو د مراتبو لړۍ
په غير رسمي توګه داسې توجيه کړې ده چې سټونه په \emph{پړاوونو} جوړېږي.
هر \LR{$V_\alpha$} تر ټاکلي پړاو وړاندې جوړ شوي سټونه رانغاړي۔
په ورپسې پړاو کښې د دې ټولو موجودو سټونو د ټولګې
هر فرعي سټ راټولېږي۔ دا د پړاوونو \emph{دنننے} توصيف دے۔
په دې توګه د نظريې د ممکن \emph{ماډل} د بيان پر ځاے،
پړاوونه د سټونو د خپلې نظريې \emph{دننه} تعريفوو۔
\olpseganchor{OLP-0559-B007}{end}

\olpseganchor{OLP-0560-B003}{start}
\olfileid{sth}{spine}{recursion}
\olsection{د نامتناهيت هاخوا بازګښت قضيې}
\olpseganchor{OLP-0560-B003}{end}

\olpseganchor{OLP-0560-B004}{start}
خو لومړے کار دا دے چې تاييد کړو
\olref[valpha]{defValphas} يو بريالے تعريف دے. په عمومي توګه، ثابتول
پکار دي چې د نامتناهيت هاخوا بازګښت د داسې قاعدې لپاره ثابت کړو
چې د هرې ټاکلې مخکنۍ برخې لپاره يو يوازينے ارزښت ټاکي۔
د دې برخې مقصد همدا دے۔
\olpseganchor{OLP-0560-B004}{end}

\olpseganchor{OLP-0560-B005}{start}
\emph{خبردار: دا مواد پېچلي دي}. خو عمومي مطلب ئې ساده دے:
د نامتناهيت هاخوا استقرا او تعويض په ګډه د نامتناهيت هاخوا بازګښت د
څو بڼو جواز تضمينوي.\footnote{يوه يادونه: ټول فارمولونه او ترمونه
پارامترونه لرلے شي، مګر که په ښکاره بل ډول ويل شوي وي.}
\olpseganchor{OLP-0560-B005}{end}

\olpseganchor{OLP-0560-B006}{start}
\begin{defn}
	فرض کړئ \LR{$\tau(x)$} ترم دے، \LR{$f$} تابعه ده، او \LR{$\alpha$} ترتيبي عدد دے.
وايو چې \LR{$f$} يو \emph{\LR{$\alpha$}-تقريب} د \LR{$\tau$} دپاره دے، که او يوازې که
دواړه \LR{$\dom{f} = \alpha$} او \LR{$(\forall \beta \in \alpha)f(\beta) = \tau(\funrestrictionto{f}{\beta})$} وي.
\end{defn}
\olpseganchor{OLP-0560-B006}{end}

\olpseganchor{OLP-0560-B007}{start}
\begin{lem}[محدود بازګښت]\ollabel{transrecursionfun}
د هر ترم \LR{$\tau(x)$} او هر ترتيبي عدد \LR{$\alpha$} دپاره، يو يوازينے
\LR{$\alpha$}-تقريب د \LR{$\tau$} دپاره شته.
\end{lem}
\begin{proof}
و به ښيو چې د هر \LR{$\gamma \leq \alpha$} دپاره يو يوازينے
\LR{$\gamma$}-تقريب شته.
\olpseganchor{OLP-0560-B007}{end}

\olpseganchor{OLP-0560-B008}{start}
لومړے يکتايي ثابتوو. فرض کړئ \LR{$g$} او \LR{$h$} په ترتيب سره \LR{$\gamma$}- او
\LR{$\delta$}-تقريبونه دي. د دواړو تابعو په ورننوتونکو ترتيبي عددونو باندې د نامتناهيت هاخوا استقرا ښيي چې
\LR{$g(\beta) = h(\beta)$} د هر \LR{$\beta \in \dom{g} \cap \dom{h} =
\gamma \cap \delta = \min(\gamma, \delta)$} دپاره دي۔ که دواړه
تقريبونه په ټولو مخکنيو ځايونو کښې برابر وي، نو هماغو
ځايونو ته يې محدوديتونه عين تابعې دي۔ قاعده عين محدوديت ته
عين ارزښت ورکوي۔ نو تقريبونه پر ګډه ساحه سره برابر دي؛ په
ځانګړې توګه، د يوې ټاکلې ساحې لپاره تقريب يوازينے دے۔
\olpseganchor{OLP-0560-B008}{end}

\olpseganchor{OLP-0560-B009}{start}
د شتون د ثابتولو دپاره اوس د نامتناهيت هاخوا ساده استقرا
(\olref[ordinals][opps]{simpletransrecursion}) په هغو ترتيبي عددونو
کاروو چې \LR{$\delta \leq \alpha$} وي.
\olpseganchor{OLP-0560-B009}{end}

\olpseganchor{OLP-0560-B010}{start}
تشه تابعه په ښکاره \LR{$\emptyset$}-تقريب ده.
\olpseganchor{OLP-0560-B010}{end}

\olpseganchor{OLP-0560-B011}{start}
که \LR{$g$} يو \LR{$\gamma$}-تقريب وي، نو \LR{$g \cup
\{\tuple{\gamma, \tau(g)}\}$} يو \LR{$\ordsucc{\gamma}$}-تقريب دے.
\olpseganchor{OLP-0560-B011}{end}

\olpseganchor{OLP-0560-B012}{start}
که \LR{$\gamma$} حدي ترتيبي عدد وي او \LR{$g_\delta$} يو
\LR{$\delta$}-تقريب د هر \LR{$\delta < \gamma$} دپاره وي، نو تعويض د
دې يوازينيو تقريبونو ټولګه راکوي۔ د هغې اتحاد
\LR{$g = \bigcup_{\delta \in \gamma} g_\delta$} واخلو۔
دا تابعه ده، ځکه چې بېل \LR{$g_\delta$} ګانې په ټولو ګډو قيمتونو سره سمون
لري۔ د حدي شاخص هر مخکنی شاخص تر خپل تالي کم دے، او
هغه تالي هم تر حدي شاخص کم دے۔ نو دغه اتحاد پر ټوله
غوښتل شوې ساحه تعريف دے۔ او که \LR{$\delta \in \gamma$} وي، نو \LR{$g(\delta) = g_{\ordsucc{\delta}}(\delta) =
\tau(\funrestrictionto{g_{\ordsucc{\delta}}}{\delta}) =
\tau(\funrestrictionto{g}{\delta})$} دے.
\olpseganchor{OLP-0560-B012}{end}

\olpseganchor{OLP-0560-B013}{start}
په دې سره د نامتناهيت هاخوا استقرا ثبوت بشپړېږي.
\end{proof}
\olpseganchor{OLP-0560-B013}{end}

\olpseganchor{OLP-0560-B014}{start}
که ځان ته د تابعې پر ځای د يوه \emph{ترم} د تعريفولو اجازه ورکړو، نو
د پخوانۍ نتيجې نه د \LR{$\alpha$} حد لرې کولے شو. د لاندې نتيجې په بيان او
ثبوت کښې، کله چې \LR{$\sigma$} ترم وي، نو
\LR{$\funrestrictionto{\sigma}{\alpha} = \Setabs{\tuple{\beta,
\sigma(\beta)}}{\beta \in \alpha}$} اخلو.
\olpseganchor{OLP-0560-B014}{end}

\olpseganchor{OLP-0560-B015}{start}
\begin{thm}[عمومي بازګښت]\ollabel{transrecursionschema}
د هر ترم \LR{$\tau(x)$} دپاره يو ترم
\LR{$\sigma(x)$} په ښکاره تعريفولے شو، داسې چې
\LR{$\sigma(\alpha) = \tau(\funrestrictionto{\sigma}{\alpha})$} د هر
ترتيبي عدد~\LR{$\alpha$} دپاره وي.
\end{thm}
\olpseganchor{OLP-0560-B015}{end}

\olpseganchor{OLP-0560-B016}{start}
\begin{proof}
د هر \LR{$\alpha$} دپاره د \olref{transrecursionfun} له مخې يو يوازينے
\LR{$\alpha$}-تقريب \LR{$f_\alpha$} د~\LR{$\tau$} دپاره شته. \LR{$\sigma(\alpha)$} د
\LR{$f_{\ordsucc{\alpha}}(\alpha)$} په توګه تعريف کړو۔ پر هغو ورودي
شيانو چې ترتيبي عددونه نۀ دي، ارزښت تش سټ ټاکلے شو۔ اوس
د ترتيبي عددونو لپاره دا برابري لرو:
	\begin{align*}
		\sigma(\alpha) &=
		f_{\ordsucc{\alpha}}(\alpha) \\&=
		\tau(\funrestrictionto{f_{\ordsucc{\alpha}}}{\alpha}) \\&=
		\tau(\Setabs{\tuple{\beta, f_{\ordsucc{\alpha}}(\beta)}}{\beta \in \alpha}) \\&=
		\tau(\Setabs{\tuple{\beta, f_{\ordsucc{\beta}}(\beta)}}{\beta \in \alpha})\\&=
		\tau(\funrestrictionto{\sigma}{\alpha})
	\end{align*}
په ياد ولرئ چې \LR{$f_{\ordsucc{\beta}}(\beta) = f_{\ordsucc{\alpha}}(\beta)$}
د ټولو \LR{$\beta < \alpha$} دپاره دي، لکه په \olref{transrecursionfun} کښې.
\end{proof}
\noindent
وګورئ چې \olref{transrecursionschema} يوه \emph{شېما} ده. دا مهمه ده چې
د \LR{$\sigma$} څخه د تابعې د تعريفولو تمه نۀ شو لرلے، يعنې د يوه ټاکلي
ډول \emph{سټ} د تعريفولو، ځکه چې بيا \LR{$\dom{\sigma}$} به د ټولو ترتيبي
عددونو سټ وي، چې د بورالي--فورتي له تناقض سره په ټکر کښې دے
(\olref[ordinals][basic]{buraliforti}).
\olpseganchor{OLP-0560-B016}{end}

\olpseganchor{OLP-0560-B017}{start}
خو لا ښودل پکار دي چې \olref{transrecursionschema}
زموږ د \LR{$V_\alpha$} ګانو تعريف توجيه کوي. ښايي دا سمدستي څرګند نۀ وي؛
خو د نامتناهيت هاخوا بازګښت په يوه وروستۍ ساده بڼه به څرګند شي.
\olpseganchor{OLP-0560-B017}{end}

\olpseganchor{OLP-0560-B018}{start}
\begin{thm}[ساده بازګښت]\ollabel{simplerecursionschema}
د هر دوو ترمونو \LR{$\tau(x)$} او \LR{$\theta(x)$} او هر سټ \LR{$A$} دپاره، يو ترم
\LR{$\sigma(x)$} په ښکاره تعريفولے شو چې:
\begin{align*}
	\sigma(\emptyset) &= A\\
	\sigma(\ordsucc{\alpha}) &= \tau(\sigma(\alpha)) &&
		\text{\RL{د هر ترتيبي عدد دپاره }}\alpha\\
	\sigma(\alpha) &= \theta(\ran{\funrestrictionto{\sigma}{\alpha}})&&
	\text{\RL{کله چې }}\alpha\text{\RL{ حدي ترتيبي عدد وي}}
\end{align*}
\end{thm}
\olpseganchor{OLP-0560-B018}{end}

\olpseganchor{OLP-0560-B019}{start}
\begin{proof}
لومړے ترم \LR{$\xi(x)$} داسې تعريفوو:
%t $\xi(x) = A$ if $x$ is not a function whose domain is an ordinal;
%otherwise let:
\[
	\xi(x) =
	\begin{cases}
			A &\text{\RL{که \LR{$x$} تابعه نۀ وي، يا ئې د تعريف}}\\
			&\hspace{1em}\text{\RL{ساحه ترتيبي عدد نۀ وي يا تشه وي؛ په بل صورت:}}\\
			\tau(x(\alpha)) & \text{\RL{که \LR{$\dom{x} = \ordsucc{\alpha}$} وي}}\\
			\theta(\ran{x}) & \text{\RL{که \LR{$\dom{x}$} حدي ترتيبي عدد وي}}
	\end{cases}
\]
د \olref{transrecursionschema} له مخې يو ترم \LR{$\sigma(x)$} شته چې
\LR{$\sigma(\alpha) = \xi(\funrestrictionto{\sigma}{\alpha})$} د هر
ترتيبي عدد \LR{$\alpha$} دپاره وي؛ همدارنګه، \LR{$\funrestrictionto{\sigma}{\alpha}$}
د \LR{$\alpha$} د تعريف ساحې لرونکې تابعه ده. ښيو چې \LR{$\sigma$} مطلوب خاصيتونه
لري، د نامتناهيت هاخوا ساده استقرا په وسيله
(\olref[ordinals][opps]{simpletransrecursion})۔ په لومړي شرط کښې
تشه تابع بنسټيز ارزښت ورکوي، او پاتې دوه شرطونه د هرې
ناتشې ترتيبي ساحې لپاره ارزښت ټاکي۔
\olpseganchor{OLP-0560-B019}{end}

\olpseganchor{OLP-0560-B020}{start}
لومړے، \LR{$\sigma(\emptyset) = \xi(\emptyset) = A$} دے.
\olpseganchor{OLP-0560-B020}{end}

\olpseganchor{OLP-0560-B021}{start}
بيا، \LR{$\sigma(\ordsucc{\alpha}) = \xi(\funrestrictionto{\sigma}{\ordsucc{\alpha}}) = \tau(\funrestrictionto{\sigma}{\ordsucc{\alpha}}(\alpha)) = \tau(\sigma(\alpha))$} دے.
\olpseganchor{OLP-0560-B021}{end}

\olpseganchor{OLP-0560-B022}{start}
په پای کښې، \LR{$\sigma(\alpha) = \xi(\funrestrictionto{\sigma}{\alpha}) = \theta(\ran{\funrestrictionto{\sigma}{\alpha}})$}، کله چې \LR{$\alpha$} حدي عدد وي.
\end{proof}
\noindent
اوس، د \olref[valpha]{defValphas} د توجيه کولو دپاره يوازې \LR{$A
= \emptyset$} او \LR{$\tau(x) = \Pow{x}$} او \LR{$\theta(x) = \bigcup x$} واخلو.
بالاخره په دې سره د \LR{$V_\alpha$} ګانو تعريف توجيه کېږي!{}
\olpseganchor{OLP-0560-B022}{end}

\olpseganchor{OLP-0561-B003}{start}
\olfileid{sth}{spine}{Valphabasic}
\olsection{د پړاوونو بنسټيز خاصيتونه}
\olpseganchor{OLP-0561-B003}{end}

\olpseganchor{OLP-0561-B004}{start}
د \LR{$V_\alpha$} ګانو د تعريف د بنسټيز اهميت د څرګندولو دپاره به ځينې
بنسټيزې نتيجې وړاندې کړو. په يوه تعريف پېل کوو:\footnote{د دې خاصيت
دپاره «توانمن» د \mbox{\LR{\latinfont \citet{ButtonLT1}}} د «potent» اصطلاح
ژباړه ده۔ د دې خاصيت لپاره يو منلے معياري نوم نشته۔}
\begin{defn}
	سټ \LR{$A$} \emph{توانمن} دے، که او يوازې که \LR{$\forall x((\exists y \in A)x \subseteq y \lif x \in A)$} وي.
\end{defn}
\begin{lem}\ollabel{Valphabasicprops}
د هر ترتيبي عدد \LR{$\alpha$} دپاره:
\begin{enumerate}
	\item\ollabel{Valphatrans} هر \LR{$V_\alpha$} متعدي دے.
	\item\ollabel{Valphapotent} هر \LR{$V_\alpha$} توانمن دے.
	\item\ollabel{Valphacum} که \LR{$\gamma \in \alpha$} وي، نو \LR{$V_\gamma
	\in V_\alpha$} دے (او له دې \LR{$V_\gamma \subseteq V_\alpha$} هم د
	\olref{Valphatrans} له مخې دے).
\end{enumerate}
\end{lem}
\olpseganchor{OLP-0561-B004}{end}

\olpseganchor{OLP-0561-B005}{start}
\begin{proof}
دا د نامتناهيت هاخوا په هممهاله استقرا ثابتوو. د استقرا دپاره فرض کړئ
چې \olref{Valphatrans}--\olref{Valphacum} د هر ترتيبي عدد \LR{$\beta < \alpha$}
دپاره صدق کوي.
\olpseganchor{OLP-0561-B005}{end}

\olpseganchor{OLP-0561-B006}{start}
د \LR{$\alpha = \emptyset$} حالت څرګند دے.
\olpseganchor{OLP-0561-B006}{end}

\olpseganchor{OLP-0561-B007}{start}
فرض کړئ \LR{$\alpha = \ordsucc{\beta}$} دے. د \olref{Valphacum} د ښودلو
دپاره، که \LR{$\gamma \in \alpha$} وي، نو \LR{$V_\gamma \subseteq V_\beta$} دے۔
که شاخصونه برابر وي، دا د سټ له ځان سره د احتوا له مخې
دے؛ او که لومړے شاخص کوچنے وي، د استقرا له فرضه راځي۔
نو \LR{$V_\gamma \in \Pow{V_\beta} = V_\alpha$} دے۔
د \olref{Valphapotent} د ښودلو دپاره فرض کړئ \LR{$A \subseteq B \in V_\alpha$}
دي، يعنې \LR{$A
\subseteq B \subseteq V_\beta$}؛ نو \LR{$A \subseteq V_\beta$} دے، نو \LR{$A \in
V_\alpha$} دے. د \olref{Valphatrans} د ښودلو دپاره وګورئ چې، که \LR{$x \in A \in
V_\alpha$} وي، نو \LR{$A \subseteq V_\beta$} لرو، نو \LR{$x \in V_\beta$} دے، نو
\LR{$x
\subseteq V_\beta$} دے، ځکه چې \LR{$V_\beta$} د فرض له مخې متعدي دے، نو \LR{$x
\in V_\alpha$} دے.
\olpseganchor{OLP-0561-B007}{end}

\olpseganchor{OLP-0561-B008}{start}
فرض کړئ \LR{$\alpha$} حدي ترتيبي عدد دے. د \olref{Valphacum} د ښودلو
دپاره، که \LR{$\gamma \in \alpha$} وي، نو \LR{$\gamma \in \ordsucc{\gamma} \in \alpha$}
دي، نو \LR{$V_\gamma \in V_{\ordsucc{\gamma}}$} د فرض له مخې دے، او له دې
\LR{$V_\gamma \in \bigcup_{\beta \in \alpha} V_\beta = V_\alpha$} دے.
د \olref{Valphatrans} او \olref{Valphapotent} د ښودلو دپاره يوازې وګورئ
چې په اتحاد کښې د يوه غړي لپاره يو شاهد سټ شته چې هغه
غړے رانغاړي۔ که شاهد متعدي وي، د غړي نور غړي هم په شاهد
او اتحاد کښې دي۔ که شاهد توانمن وي، د غړي فرعي سټونه هم
په شاهد او اتحاد کښې دي۔
\end{proof}
\olpseganchor{OLP-0561-B008}{end}

\olpseganchor{OLP-0561-B009}{start}
\begin{lem}\ollabel{Valphanotref}
د هر ترتيبي عدد \LR{$\alpha$} دپاره، \LR{$V_\alpha \notin V_\alpha$} دے.
\end{lem}
\olpseganchor{OLP-0561-B009}{end}

\olpseganchor{OLP-0561-B010}{start}
\begin{proof}
د نامتناهيت هاخوا استقرا کاروو. ښکاره ده چې \LR{$V_\emptyset \notin V_\emptyset$} دے.
\olpseganchor{OLP-0561-B010}{end}

\olpseganchor{OLP-0561-B011}{start}
که \LR{$V_{\ordsucc{\alpha}} \in V_{\ordsucc{\alpha}} = \Pow{V_\alpha}$} وي،
نو \LR{$V_{\ordsucc{\alpha}} \subseteq V_\alpha$} دے؛ او ځکه چې \LR{$V_\alpha
\in V_{\ordsucc{\alpha}}$} د \olref{Valphabasicprops} له مخې دے،
\LR{$V_\alpha \in V_\alpha$} لرو. په معکوسه توګه: که \LR{$V_\alpha \notin V_\alpha$}
وي، نو \LR{$V_{\ordsucc{\alpha}} \notin V_{\ordsucc{\alpha}}$} دے.
\olpseganchor{OLP-0561-B011}{end}

\olpseganchor{OLP-0561-B012}{start}
که \LR{$\alpha$} حدي عدد وي او \LR{$V_\alpha \in V_\alpha = \bigcup_{\beta \in
\alpha}V_\beta$} وي، نو \LR{$V_\alpha \in V_\beta$} د کوم \LR{$\beta \in
\alpha$} دپاره دے؛ خو بيا \LR{$V_\beta \in V_\alpha$} هم دے، نو \LR{$V_\beta \in
V_\beta$} د \olref{Valphabasicprops} په دوه ځله کارولو دے.
په معکوسه توګه، که \LR{$V_\beta
\notin V_\beta$} د ټولو \LR{$\beta \in \alpha$} دپاره وي، نو \LR{$V_\alpha \notin
V_\alpha$} دے.
\end{proof}
\olpseganchor{OLP-0561-B012}{end}

\olpseganchor{OLP-0561-B013}{start}
\begin{cor}
د هر دوو ترتيبي عددونو \LR{$\alpha, \beta$} دپاره: \LR{$\alpha \in \beta$}، که او يوازې که \LR{$V_\alpha \in V_\beta$} وي.
\end{cor}
\olpseganchor{OLP-0561-B013}{end}

\olpseganchor{OLP-0561-B014}{start}
\begin{proof}
\Olref{Valphabasicprops} يو لورے ورکوي. په معکوسه توګه فرض کړئ
\LR{$V_\alpha \in V_\beta$} دے. نو \LR{$\alpha \neq \beta$} د \olref{Valphanotref}
له مخې دي؛ او \LR{$\beta \notin \alpha$}، ځکه چې که داسې نۀ وي، نو
\LR{$V_\beta \in V_\alpha$} به لرو او له دې \LR{$V_\beta \in V_\beta$} د
\olref{Valphabasicprops} په دوه ځله کارولو لرو، چې له
\olref{Valphanotref} سره په تناقض کښې دي. نو \LR{$\alpha \in \beta$} د درې ګونتوب له مخې دے.
\end{proof}
\noindent
دا ټول موږ ته اجازه راکوي چې هر \LR{$V_\alpha$} د مراتبو د لړۍ د \LR{$\alpha$}
پړاو په توګه وګڼو. وجه ئې دا ده:
\olpseganchor{OLP-0561-B014}{end}

\olpseganchor{OLP-0561-B015}{start}
بې شکه \LR{$V_\alpha$} ګانې د يوه \emph{تکراري} بهير په وسيله جوړې شوې
ګڼلے شو، ځکه چې د ترتيبي عددونو کارول مو د تکرار تصور تعقيبوي.
همدارنګه، که يو پړاو تر بل مخکې جوړېږي، يعنې
\LR{$V_\beta \in V_\alpha$}، يعنې \LR{$\beta \in \alpha$}، نو د جوړولو بهير مو
\emph{تجمعي} دے، ځکه چې \LR{$V_\beta \subseteq V_\alpha$} دے.
په پای کښې، په رښتيا د هغو سټونو \emph{ټولې} ممکنې ټولګې جوړوو چې
په کوم مخکيني پړاو کښې موجود وو، ځکه چې هر تالي پړاو
\LR{$V_{\ordsucc{\alpha}}$} د خپل مخکيني پړاو \LR{$V_\alpha$} د فرعي سټونو سټ دے.
\olpseganchor{OLP-0561-B015}{end}

\olpseganchor{OLP-0561-B016}{start}
په دې توګه، په \LR{$\ZFminus$} کښې د سټونو د تجمعي--تکراري مراتبو د لړۍ
د تصور په کره بيانولو کښې \emph{تقريباً} کار بشپړ شوے دے.
خو، بې شکه، لا د تعويض توجيه کول پکار دي.
\olpseganchor{OLP-0561-B016}{end}

\olpseganchor{OLP-0562-B003}{start}
\olfileid{sth}{spine}{foundation}
\olpseganchor{OLP-0562-B003}{end}

\olpseganchor{OLP-0562-B004}{start}
\olsection{د بنسټ اصل}
\olpseganchor{OLP-0562-B004}{end}

\olpseganchor{OLP-0562-B005}{start}
زموږ کار يوازې \emph{نږدې} پاې ته رسېدلے دے، لا \emph{بشپړ} نۀ دے؛ ځکه په
\LR{$\ZFminus$} کښې هېڅ داسې څه نشته چې دا تضمين کړي چې \emph{هر} سټ په کوم
\LR{$V_\alpha$} کښې دے؛ يعنې هر سټ په کوم پړاو کښې جوړېږي۔
\olpseganchor{OLP-0562-B005}{end}

\olpseganchor{OLP-0562-B006}{start}
اوس په يوه نسبتاً څرګنده رياضياتي معنٰی کښې موږ ته دا \emph{پروا نۀ لري}
چې له دې سلسلې نه بهر سټونه شته که نۀ۔ (که هلته سټونه وي، په اسانۍ سره ئې
له پامه غورځولے شو۔) خو د سټ خپل \emph{مفهوم} مو په دې فکر ولاړ کړے دے چې
هر سټ په کوم پړاو کښې جوړېږي (په \olref[z][story]{sec} کښې \stageshier{}
وګورئ)۔ نو غواړو چې له دې سلسلې نه بهر د سټونو امکان رد کړو۔ په دې خاطر
بايد يو نوے اصل ورزيات کړو، څو ډاډمن شو چې هر سټ د سلسلې په کوم ځاے کښې
راځي۔
\olpseganchor{OLP-0562-B006}{end}

\olpseganchor{OLP-0562-B007}{start}
څرنګه چې \LR{$V_\alpha$} زموږ پړاوونه دي، په ساده توګه دا لاندې وينا د يوه
اصل په توګه ورزياتولے شو:
\olpseganchor{OLP-0562-B007}{end}

\olpseganchor{OLP-0562-B008}{start}
\begin{defish}
\emph{منظموالی۔} \LR{$\forall A \exists \alpha\, A \subseteq V_\alpha$}
\end{defish}
\olpseganchor{OLP-0562-B008}{end}

\olpseganchor{OLP-0562-B009}{start}
دا به يو بالکل معقول چلند وي۔ خو د هغو دليلونو له مخې چې په راتلونکې
برخه کښې به بيان شي، د دې پر ځاے يو بل اصل غوره کوو:
\olpseganchor{OLP-0562-B009}{end}

\olpseganchor{OLP-0562-B010}{start}
\begin{axiom}[د بنسټ اصل]
\LR{$(\forall A \neq \emptyset)(\exists B \in A)A \cap B = \emptyset$}.
\end{axiom}
\olpseganchor{OLP-0562-B010}{end}

\olpseganchor{OLP-0562-B011}{start}
په يو څه هڅه سره په \LR{$\ZFminus$} کښې ښودلے شو چې د بنسټ اصل منظموالي ته
استلزام ورکوي:
\begin{defn}
د هر سټ \LR{$A$} دپاره داسې واخلئ:
\begin{align*}
	\text{\RL{cl}}_0(A) &= A,\\
	\text{\RL{cl}}_{n+1}(A) &= \bigcup \text{\RL{cl}}_n(A),\\
	\text{\RL{trcl}}(A) &= \bigcup_{n < \omega} \text{\RL{cl}}_{n}(A).
\end{align*}
\LR{$\text{\RL{trcl}}(A)$} ته د \LR{$A$} \emph{متعدي بندښت} وايو۔
\end{defn}
\noindent
د ``متعدي بندښت'' نوم مناسب دے:
\begin{prop}\ollabel{subsetoftrcl}
\LR{$A \subseteq \trcl{A}$} او \LR{$\trcl{A}$} يو متعدي سټ دے۔
\end{prop}
\olpseganchor{OLP-0562-B011}{end}

\olpseganchor{OLP-0562-B012}{start}
\begin{proof}
په ښکاره \LR{$A = \text{\RL{cl}}_0(A) \subseteq \text{\RL{trcl}}(A)$} دے۔ او که \LR{$x
\in b \in \trcl{A}$} وي، نو \LR{$b \in \text{\RL{cl}}_n(A)$} د کوم \LR{$n$} دپاره دے؛
له دې امله \LR{$x \in \text{\RL{cl}}_{n+1}(A) \subseteq \text{\RL{trcl}}(A)$} دے۔
\end{proof}
\olpseganchor{OLP-0562-B012}{end}

\olpseganchor{OLP-0562-B013}{start}
\begin{lem}\ollabel{lem:TransitiveWellFounded}
که \LR{$A$} يو متعدي سټ وي، نو داسې \LR{$\alpha$} شته چې \LR{$A
\subseteq V_\alpha$} دے۔
\end{lem}
\olpseganchor{OLP-0562-B013}{end}

\olpseganchor{OLP-0562-B014}{start}
\begin{proof}
له \olref[ordinals][opps]{defsupstrict} نه د ``\LR{$\supstrict(X)$}'' تعريف راياد
کړئ او دا دوه سټونه تعريف کړئ:
\begin{align*}
	D  &= \Setabs{x \in A}{\forall \delta\ x \nsubseteq V_\delta}\\
	\alpha &= \supstrict\Setabs{\delta}{(\exists x \in A)
	(x \subseteq V_\delta \land (\forall \gamma \in \delta)x \nsubseteq V_\gamma)}
\end{align*}
د بېلونې له مخې له راکړل شوي سټه د هغو غړو فرعي سټ
اخلو چې لږ تر لږه يو داسې شاخص لري۔ د هر داسې غړي
لپاره لومړنے شاخص يوازينے دے؛ تعويض د ټولو دغو شاخصونو
سټ راکوي۔
فرض کړئ \LR{$D = \emptyset$} دے۔ نو که \LR{$x \in A$} وي، داسې \LR{$\delta$} شته چې
\LR{$x \subseteq V_\delta$} دے او د ترتيبي عددونو د ښه ترتيب له مخې
\LR{$(\forall \gamma \in \delta)x \nsubseteq V_\gamma$} دے؛ نو \LR{$\delta \in \alpha$}
دے او د \olref[spine][Valphabasic]{Valphabasicprops} له مخې اړوند
پخوانی پړاو د وروستي پړاو غړے دے۔ د وروستي پړاو د
توانمن‌والي له مخې \LR{$x \in V_\alpha$} دے۔ له دې امله
\LR{$A \subseteq V_{\alpha}$} دے۔
\olpseganchor{OLP-0562-B014}{end}

\olpseganchor{OLP-0562-B015}{start}
نو يوازې دا ښودل بس دي چې \LR{$D = \emptyset$} دے۔ د تناقض له لارې د ثبوت
دپاره ئې عکس فرض کړئ۔ د بنسټ د اصل له مخې داسې \LR{$B \in D \subseteq A$}
شته چې \LR{$D \cap B = \emptyset$} دے۔ که \LR{$x \in B$} وي، نو \LR{$x \in A$} دے، ځکه
\LR{$A$} متعدي دے؛ او څرنګه چې \LR{$x \notin D$} دے، نو \LR{$\exists \delta\ x \subseteq
V_\delta$} دے۔ د لږ تر لږه شاخصونو سټ بيا په تعويض جوړوو۔
اوس داسې واخلئ:
\[
	\beta = \supstrict\Setabs{\delta}{(\exists x \in B)(x \subseteq V_\delta \land (\forall \gamma < \delta)x \nsubseteq V_\gamma)}.
\]
لکه د مخکښې په شان، \LR{$B \subseteq V_\beta$} دے؛ دا له دې وينا سره تناقض
لري چې \LR{$B \in D$} دے۔
\end{proof}
\olpseganchor{OLP-0562-B015}{end}

\olpseganchor{OLP-0562-B016}{start}
\begin{thm}\ollabel{zfentailsregularity}
منظموالی صدق کوي۔
\end{thm}
\olpseganchor{OLP-0562-B016}{end}

\olpseganchor{OLP-0562-B017}{start}
\begin{proof}
\LR{$A$} وټاکئ؛ د \olref{subsetoftrcl} له مخې \LR{$A \subseteq \trcl{A}$} دے، او
دا متعدي سټ دے۔ نو د \olref{lem:TransitiveWellFounded} له مخې داسې
\LR{$\alpha$} شته چې \LR{$A \subseteq \trcl{A} \subseteq V_\alpha$} دے۔
\end{proof}
\olpseganchor{OLP-0562-B017}{end}

\olpseganchor{OLP-0562-B018}{start}
دا پايلې ښيي چې \LR{$\ZFminus$} شرطي وينا
\LR{$\emph{د بنسټ اصل}\Rightarrow\emph{منظموالی}$} ثابتوي۔ په
\olref[spine][rank]{zfminusregularityfoundation} کښې به وښيو چې
\LR{$\ZFminus$} د \LR{$\emph{منظموالی}\Rightarrow\emph{د بنسټ اصل}$} وينا ثابتوي۔
نو د بنسټ اصل او منظموالی د \LR{$\ZFminus$} په شتون کښې \emph{همارزښته} دي۔
دا معنا لري چې د \LR{$\ZFminus$} په فرضولو سره د بنسټ اصل په دې توجيه کولے
شو چې له منظموالي سره همارزښته دے۔ او منظموالی د \stageshier{} پر بنسټ
نېغ په نېغه توجيه کولے شو۔
\olpseganchor{OLP-0562-B018}{end}

\olpseganchor{OLP-0563-B003}{start}
\olfileid{sth}{spine}{zf}
\olsection{\texorpdfstring{\LR{$\Z$}}{Z} او \texorpdfstring{\LR{$\ZF$}}{ZF}: يو مهم پړاو}
\olpseganchor{OLP-0563-B003}{end}

\olpseganchor{OLP-0563-B004}{start}
د بنسټ له اصل سره يو بل مهم پړاو ته رسېږو۔ \LR{$\Zminus$} او \LR{$\ZFminus$}
تيورۍ مو څېړلې وې؛ ويلي مو وو چې دا د کوم ټاکلي شي له ايستلو وروسته
پاتې تيورۍ دي۔ هماغه ټاکلے شوے شے د بنسټ اصل دے۔ نو:
\olpseganchor{OLP-0563-B004}{end}

\olpseganchor{OLP-0563-B005}{start}
\begin{defn}
\LR{$\Z$} تيوري \LR{$\Zminus$} ته د بنسټ اصل ورزياتوي۔ نو اصول ئې دا دي: د غړو
له مخې برابري، اتحاد، جوړې، د فرعي سټونو سټونه، نامتناهيت، د بنسټ اصل،
او د بېلونې د شېما ټولې بېلګې۔
\olpseganchor{OLP-0563-B005}{end}

\olpseganchor{OLP-0563-B006}{start}
\LR{$\ZF$} تيوري \LR{$\ZFminus$} ته د بنسټ اصل ورزياتوي۔ په بله وينا، \LR{$\ZF$}
د تعويض ټولې بېلګې \LR{$\Z$} ته ورزياتوي۔
\end{defn}
\olpseganchor{OLP-0563-B006}{end}

\olpseganchor{OLP-0563-B007}{start}
خو ښايي بيا هم يو سوال درته پيدا شوے وي۔ که منظموالی په \LR{$\ZFminus$}
کښې د بنسټ له اصل سره همارزښته دے او د منظموالي توجيه څرګنده ده، نو
ولې دا اوږده لاره ونيسو او د منظموالي پر ځاے د بنسټ اصل خپل بنسټيز
اصل وګرځوو؟
\olpseganchor{OLP-0563-B007}{end}

\olpseganchor{OLP-0563-B008}{start}
تاريخي دليلونه، يعنې دا چې چا کومه وينا کله جوړه کړې وه، يوې خوا ته
پرېږدئ۔ اصلي دليل دا دے چې د بنسټ اصل د \LR{$V_\alpha$} د تعريف له کارولو
پرته وړاندې کېدلے شي۔ هغه تعريف د \olref[recursion]{sec} پر ټول کار
ولاړ و: د نامتناهيت هاخوا بازګښت ثبوت ته اړ وو څو د تعريف توجيه وښيو۔ خو
د نامتناهيت هاخوا بازګښت په ثبوت کښې \emph{تعويض} کارول کېږي۔
د بنسټ اصل او منظموالی د \LR{$\ZFminus$} په شتون کښې همارزښته دي۔
په \LR{$\Zminus$} کښې د دغې هم‌ارزۍ د همدې دليل لپاره تعويض پکار وي۔
\olpseganchor{OLP-0563-B008}{end}

\olpseganchor{OLP-0563-B009}{start}
په حقيقت کښې موضوع تر دې ساده څرګندونې هم جدي ده۔ که څه هم دا بحث
د دې کتاب له حدودو ډېر بهر دے، معلومېږي چې \LR{$\Zminus$} او \LR{$\Z$} دواړه
د ټولو \LR{$V_\alpha$} ګانو د بازګشتي جوړولو د تضمين لپاره بس نۀ دي۔
نو په \LR{$\Z$} کښې د منظموالي پورته بڼه د پړاوونو پر داسې
جوړښت تکيه کوي چې دلته يې \emph{شتون} نۀ دے تضمين شوے۔
د بنسټ اصل مستقيماً د سټونو او غړيتوب په ژبه بيانېږي۔
\olpseganchor{OLP-0563-B009}{end}

\olpseganchor{OLP-0563-B010}{start}
له دې وروسته به، که بل ډول مو نۀ وي ويلي، له کومې اضافي يادونې پرته
په \LR{$\ZF$} کښې کار کوو۔
\olpseganchor{OLP-0563-B010}{end}

\olpseganchor{OLP-0564-B003}{start}
\olfileid{sth}{spine}{rank}
\olsection{رتبه}
\olpseganchor{OLP-0564-B003}{end}

\olpseganchor{OLP-0564-B004}{start}
اوس چې پړاوونه مو د \LR{$V_\alpha$} په توګه تعريف کړي او پوهېږو چې هر سټ
د کوم پړاو فرعي سټ دے، د سټ \emph{رتبه} تعريفولے شو۔ په بداهتي توګه،
د \LR{$A$} رتبه هغه لومړۍ شېبه ده چې \LR{$A$} پکښې جوړېږي۔ په کره توګه:
\olpseganchor{OLP-0564-B004}{end}

\olpseganchor{OLP-0564-B005}{start}
\begin{defn}\ollabel{defnsetrank}
د هر سټ \LR{$A$} دپاره \LR{$\setrank{A}$} تر ټولو وړوکے هغه ترتيبي عدد \LR{$\alpha$}
دے چې \LR{$A \subseteq V_\alpha$} دے۔
\end{defn}
\begin{prop}\ollabel{ranksexist}
	\LR{$\setrank{A}$} د هر \LR{$A$} دپاره شته۔
\end{prop}
\begin{proof}
	د تمرين لپاره اشاره دا ده چې منظموالی لږ تر لږه يو مناسب
	پړاو راکوي۔ د مناسب شاخصونو په منځ کښې د تر ټولو کوچني
	شاخص د شتون لپاره د ترتيبي عددونو ښه ترتيب وکاروئ۔
\end{proof}
\begin{prob}
	د \olref[sth][spine][rank]{ranksexist} ثبوت ورکړئ۔
\end{prob}
\noindent
د رتبو ښه ترتيب موږ ته د ځينو مهمو پايلو د ثابتولو امکان راکوي:
\begin{prop}\ollabel{valphalowerrank}
د هر ترتيبي عدد \LR{$\alpha$} دپاره \LR{$V_\alpha = \Setabs{x}{\setrank{x} \in \alpha}$} دے۔
\end{prop}
\olpseganchor{OLP-0564-B005}{end}

\olpseganchor{OLP-0564-B006}{start}
\begin{proof}
که \LR{$\setrank{x} \in \alpha$} وي، نو \LR{$x \subseteq V_{\setrank{x}} \in
V_\alpha$} دے؛ نو \LR{$x \in V_\alpha$} دے، ځکه چې \LR{$V_\alpha$} توانمن دے
(دلته \olref[Valphabasic]{Valphabasicprops} څو ځلې کارول کېږي)۔ بالعکس،
که \LR{$x \in V_\alpha$} وي، نو \LR{$x \subseteq V_\alpha$} دے؛ له دې امله
\LR{$\setrank{x} \leq \alpha$} دے۔ خو \LR{$x \notin V_{\setrank{x}}$}
دے۔ که رتبه صفر وي، د هغې پړاو تش دے۔ که رتبه تالي وي،
د هغې په پړاو کښې غړيتوب د مخکني پړاو فرعي سټ کېدل
غواړي۔ که رتبه حدي وي، د هغې په پړاو کښې هر غړے د يوه
مخکني پړاو غړے هم دے، او د تعديت له مخې د هماغه مخکني
پړاو فرعي سټ دے۔ په دواړو وروستيو حالتونو کښې به يو کوچنے
شاخص د رتبې شرط پوره کړي، چې دا د رتبې له لږوالي سره ټکر
لري۔ نو د سټ رتبه له ورکړل شوي ترتيبي عدد څخه کمه ده۔
\end{proof}
\begin{prob}
	په \olref[sth][spine][rank]{valphalowerrank} کښې د خپل پړاو
	د غړيتوب د نشتون دليل په صفر، تالي او حدي حالتونو کښې په
	تفصيل وليکئ۔
\end{prob}
\olpseganchor{OLP-0564-B006}{end}

\olpseganchor{OLP-0564-B007}{start}
\begin{prop}\ollabel{rankmemberslower}
که \LR{$B \in A$} وي، نو \LR{$\setrank{B} \in \setrank{A}$} دے۔
\end{prop}
\olpseganchor{OLP-0564-B007}{end}

\olpseganchor{OLP-0564-B008}{start}
\begin{proof}
د \olref{valphalowerrank} له مخې
\LR{$A \subseteq V_{\setrank{A}} = \Setabs{x}{\setrank{x} \in
\setrank{A}}$} دے۔
\end{proof}
\noindent
په دې حقيقت داسې پايله ثابتولے شو چې د استقرا د يوې بڼې له لارې د
\emph{ټولو سټونو} په اړه د ويناوو د ثابتولو امکان راکوي:
\olpseganchor{OLP-0564-B008}{end}

\olpseganchor{OLP-0564-B009}{start}
\begin{thm}[د \LR{$\in$}-استقرا شېما]
د هر فارمول \LR{$\phi$} دپاره:
\[
	\forall A((\forall x \in A)\phi(x) \lif \phi(A)) \lif \forall A \phi(A).
\]
\end{thm}
\olpseganchor{OLP-0564-B009}{end}

\olpseganchor{OLP-0564-B010}{start}
\begin{proof}
نقيض عکس به ثابتوو۔ نو فرض کړئ \LR{$\lnot \forall A
\phi(A)$} دے۔ د نامتناهيت هاخوا استقرا
(\olref[ordinals][basic]{ordinductionschema}) له مخې يو داسې سټ شته
چې \LR{$\phi$} نۀ دے او تر ټولو وړوکې ممکنه رتبه لري؛ يعنې داسې \LR{$A$} شته
چې \LR{$\lnot \phi(A)$} دے او \LR{$\forall x(\setrank{x} \in \setrank{A}
\lif \phi(x))$} دے۔ اوس که \LR{$x \in A$} وي، نو د
\olref{rankmemberslower} له مخې \LR{$\setrank{x} \in
\setrank{A}$} دے؛ له دې امله \LR{$\phi(x)$} دے؛ يعنې \LR{$(\forall x \in
A)\phi(x) \land \lnot \phi(A)$} دے۔
\end{proof}\noindent دا ځواکمنه پايله په غير رسمي توګه داسې بيانولے
شو: \LR{$\phi$} هله او يوازې هله \emph{وراثتي} دے چې، هر کله د يوه سټ
هر غړے د \LR{$\phi$} خاصيت ولري، هماغه سټ پخپله هم \LR{$\phi$} وي۔ نو
\LR{$\in$}-استقرا دا درته وايي: که \LR{$\phi$} وراثتي وي، هر سټ \LR{$\phi$} دے۔
\olpseganchor{OLP-0564-B010}{end}

\olpseganchor{OLP-0564-B011}{start}
د رتبو بحث د اوس دپاره په دې توګه پاې ته رسوو چې هغه څو ادعاوې
ثابتې کړو چې څو ځلې مو ئې مخکښې اشاره کړې وه۔
\olpseganchor{OLP-0564-B011}{end}

\olpseganchor{OLP-0564-B012}{start}
\begin{prop}\ollabel{ranksupstrict}
\LR{$\setrank{A} = \supstrict_{x \in A}\setrank{x}$}.
\end{prop}
\olpseganchor{OLP-0564-B012}{end}

\olpseganchor{OLP-0564-B013}{start}
\begin{proof}
تعويض د غړو د رتبو سټ راکوي۔ د دغو رتبو سخت پاسنے حد
\LR{$\alpha = \supstrict_{x \in A}\setrank{x}$} واخلئ۔ د
\olref{rankmemberslower} له مخې \LR{$\alpha \leq \setrank{A}$} دے۔ خو که
\LR{$x \in A$} وي، نو \LR{$\setrank{x} \in \alpha$} دے؛ نو د
\olref{valphalowerrank} له مخې \LR{$x \in V_\alpha$} دے، او له دې امله
\LR{$A \subseteq V_\alpha$} دے؛ يعنې \LR{$\setrank{A} \leq \alpha$} دے۔ نو
\LR{$\setrank{A} = \alpha$} دے۔
\end{proof}
\olpseganchor{OLP-0564-B013}{end}

\olpseganchor{OLP-0564-B014}{start}
\begin{cor}\ollabel{ordsetrankalpha}
د هر ترتيبي عدد \LR{$\alpha$} دپاره \LR{$\setrank{\alpha} = \alpha$} دے۔
\end{cor}
\olpseganchor{OLP-0564-B014}{end}

\olpseganchor{OLP-0564-B015}{start}
\begin{proof}
د نامتناهيت هاخوا استقرا دپاره فرض کړئ چې \LR{$\setrank{\beta} = \beta$}
د هر \LR{$\beta \in \alpha$} دپاره دے۔ اوس د \olref{ranksupstrict} له مخې
\LR{$\setrank{\alpha} = \supstrict_{\beta \in
\alpha}\setrank{\beta} = \supstrict_{\beta \in \alpha}\beta = \alpha$}
دے۔
%	First note that $\setrank{\alpha} \neq \beta$ for any $\beta \in
%	\alpha$. For otherwise we would have $\beta = \setrank{\beta} \in
%	\setrank{\alpha} = \beta$ by \olref{rankmemberslower}, i.e.,
%	$\beta \in \beta$, a contradiction. 
%
%	Now note that $\alpha \subseteq V_\alpha$. For $\alpha =
%	\Setabs{\beta}{\beta \in \alpha}$ by
%	\olref[sfr][ordinals][basic]{ordissetofsmallerord}. And if
%	$\beta \in \alpha$ then $\beta \subseteq V_{\setrank{\beta}} =
%	V_\beta \in V_\alpha$, hence $\beta \in V_\alpha$, by
%	\olref[sfr][spine][Valphabasic]{Valphabasicprops} twice. 
\end{proof}
\olpseganchor{OLP-0564-B015}{end}

\olpseganchor{OLP-0564-B016}{start}
په پاې کښې، د هغې پايلې يو لنډ ثبوت وړاندې کوو چې د
\olref[foundation]{sec} په پاې کښې مو ئې ژمنه کړې وه: \LR{$\ZFminus$} د
\LR{$\emph{منظموالی} \Rightarrow \emph{د بنسټ اصل}$} شرطي وينا ثابتوي۔
(ياد ولرئ چې په دې ثبوت کښې د ``رتبې'' مفهوم او
\olref{rankmemberslower} کارولے شو؛ ځکه لکه څنګه چې د دې برخې په
پيل کښې مو ياد کړل، دا د \LR{$\ZFminus + \text{\RL{منظموالی}}$} له مخې وړاندې
کېدلے شي۔)
\olpseganchor{OLP-0564-B016}{end}

\olpseganchor{OLP-0564-B017}{start}
\begin{prop}[په \LR{$\ZFminus + \text{\RL{منظموالی}}$} کښې کار کوو]\ollabel{zfminusregularityfoundation} د بنسټ اصل صدق کوي۔
\end{prop}
\olpseganchor{OLP-0564-B017}{end}

\olpseganchor{OLP-0564-B018}{start}
\begin{proof}
\LR{$A \neq \emptyset$} وټاکئ او داسې \LR{$B \in A$} واخلئ چې تر ټولو وړوکې
ممکنه رتبه لري۔ که \LR{$c \in B$} وي، نو د \olref{rankmemberslower} له مخې
\LR{$\setrank{c} \in \setrank{B}$} دے؛ نو \LR{$c \notin A$} دے، د \LR{$B$} د
انتخاب له امله۔ نو د دواړو سټونو تقاطع تشه ده۔
\end{proof}
\olpseganchor{OLP-0564-B018}{end}

\olpseganchor{OLP-0565-B004}{start}
\olchapter{sth}{replacement}{تعويض}
\olpseganchor{OLP-0565-B004}{end}

\olpseganchor{OLP-0566-B003}{start}
\olfileid{sth}{replacement}{intro}
\olsection{سريزه}
\olpseganchor{OLP-0566-B003}{end}

\olpseganchor{OLP-0566-B004}{start}
تعويض هغه د اصلونو شېما ده چې د \LR{$\ZF$} او~\LR{$\Z$} تر منځ توپير
راولي۔ له \crefrange{sth:ordinals::chap}{sth:spine::chap} پورې مو ترې
کار واخيست۔ په دې څپرکي کښې به بالاخره دا پوښتنه وڅېړو: آيا تعويض
توجيه کېدلے شي؟
\olpseganchor{OLP-0566-B004}{end}

\olpseganchor{OLP-0566-B005}{start}
د پوښتنې د کره کولو دپاره بايد پام وکړو چې تعويض په رښتيا هم ډېر
\emph{قوي} دے۔ په دې څپرکي کښې، او بيا په
\olref[sth][card-arithmetic][fix]{sec} کښې، به د دې قوت اندازه
تر يوه حده څرګنده شي۔ خو همدا قوت ښيي چې توجيې ته ئې رښتينې اړتيا شته۔
\olpseganchor{OLP-0566-B005}{end}

\olpseganchor{OLP-0566-B006}{start}
د توجيې دوه ډولونه به وڅېړو۔ په لنډه توګه، \emph{بهرنۍ} توجيه
هڅه کوي چې يو اصل د هغه د پايلو په مټ توجيه کړي؛ \emph{دننۍ} توجيه
هڅه کوي چې وښيي د بحث رياضيکي مفهومونه پخپله د هغه اصل ملاتړ کوي۔ د
دې توپير معنا به په دې څپرکي کښې لا ښه روښانه شي، خو دا د يوه پراخ
بحث يوازې سر دے۔ د زياتو معلوماتو دپاره په ځانګړي ډول
\mbox{\LR{\latinfont \citeauthor{Maddy1988a}}} (\mbox{\LR{\latinfont \citeyear{Maddy1988a}}} او
\mbox{\LR{\latinfont \citeyear{Maddy1988b}}}) وګورئ۔
\olpseganchor{OLP-0566-B006}{end}

\olpseganchor{OLP-0567-B003}{start}
\olfileid{sth}{replacement}{strength}
\olsection{د تعويض قوت}
\olpseganchor{OLP-0567-B003}{end}

\olpseganchor{OLP-0567-B004}{start}
که زرمېلو–فرېنکل تيوري سازګاره وي، نو په \LR{$\Z$} کښې د
\LR{$\omega + \omega$} سره برابر يا تر هغه لوے هېڅ د فون نويمن
ترتيبي عدد شتون نۀ شي ثابتېدلے۔
\olpseganchor{OLP-0567-B004}{end}

\olpseganchor{OLP-0567-B005}{start}
د دې دليل لنډه طرحه داسې ده۔ په~\LR{$\ZF$} کښې کار وکړئ او
\LR{$V_{\omega+\omega}$} سټ په پام کښې ونيسئ۔ دا سټ د~\LR{$\Z$} د يوه
\emph{مدل} د ساحې کار کوي۔ د دې د ليدلو دپاره د فارمول د
\emph{نسبي کولو} نښه تعريفوو:
\begin{defn}\ollabel{formularelativization}
	د هر سټ \LR{$M$} او هر فارمول \LR{$\phi$} دپاره، \LR{$\phi^M$} هغه فارمول دے
چې د \LR{$\phi$} ټول کميت ټاکونکي په \LR{$M$} پورې محدود کړو۔ يعنې
``\LR{$\lexists[x]$}'' په ``\LR{$(\lexists[x \in M])$}'' او
``\LR{$\lforall[x]$}'' په ``\LR{$(\lforall[x \in M])$}'' بدل کړئ۔
\end{defn}\noindent
د هر اصل \LR{$\phi$} د~\LR{$\Z$} دپاره
\LR{$\ZF \vdash \phi^{V_{\omega+\omega}}$} دے۔ دغه ساحه متعديه
ده، نو د غړو له مخې برابري پکښې صدق کوي او د بنسټ اصل له ټول
جهانه ورته رالېږدي۔ ساحه تش سټ او د طبيعي عددونو سټ لري۔
د سټونو جوړه، اتحاد او د فرعي سټونو سټ هم پکښې پاتې کېږي، ځکه
د رتبې د شاخص هر متناهي زياتوالے بيا هم تر حدي شاخص لاندې
دے۔ بېلونه د ساحې په نسبت کاروو؛ د يوه موجود سټ هر فرعي
سټ بيا په ساحه کښې دے۔ د \olref[spine][rank]{ordsetrankalpha}
له مخې \LR{$\omega+\omega$} په \LR{$V_{\omega+\omega}$} کښې
\emph{غړے نۀ دے}۔ نو د زرمېلو–فرېنکل د سازګارۍ په فرض،
\LR{$\Z$} د \LR{$\omega+\omega$} له نشتون سره سازګار دے۔
\olpseganchor{OLP-0567-B005}{end}

\olpseganchor{OLP-0567-B006}{start}
له همدې امله مو په \olref[ordinals][replacement]{sec} کښې ويلي وو
چې \olref[ordinals][ordtype]{thmOrdinalRepresentation} په زرمېلو
تيوري کښې په يوازې ځان نۀ شي ثابتېدلے۔ ځکه په~\LR{$\Z$} کښې يو څرګند ښه ترتيب په اسانۍ تعريفولے
شو چې په بداهتي توګه ئې د ترتيب ډول بايد \LR{$\omega+\omega$} وي۔ د دې
يوه غير رسمي مثال مو په \olref[ordinals][idea]{sec} کښې د طبيعي
عددونو په دې ترتيب سره ورکړے و:
\begin{align*}
	n \lessdot m \text{\RL{ هله او يوازې هله چې }}&\text{\RL{يا }}n < m\text{\RL{ او }}m-n\text{\RL{ جفت وي،}}\\
	& \text{\RL{يا \LR{$n$} جفت او \LR{$m$} طاق وي.}}
\end{align*}
په دې ښه ترتيب کښې ټول جفت عددونه تر طاقو مخکې دي او
هره ډله خپل طبيعي ترتيب لري۔ دا اړيکه د طبيعي عددونو د سټ د ديکارتي ضرب فرعي سټ دے او په همدې ساحه کښې شته۔ له بهره يې ترتيبي ډول
\LR{$\omega+\omega$} دے۔ هر ترتيبي عدد چې په متعديه ساحه کښې
شته، تر دې ډول کوچنے دے۔ که په ساحه کښې له داسې عدد سره
د ترتيب آيزومورفيزم موجود وے، له بهره به هم آيزومورفيزم وے،
ځکه ساحه د تابعې د دواړو خواوو ټول غړي لري۔ دا د ترتيبي
ډول له يوازينيت سره ټکر لري۔ د پورته سازګارۍ په فرض،
\LR{$\Z$} د \olref[ordinals][ordtype]{thmOrdinalRepresentation}
ثبوت \emph{نۀ} ورکوي۔
\olpseganchor{OLP-0567-B006}{end}

\olpseganchor{OLP-0567-B007}{start}
اوس موضوع بل اړخ ته واړوو: تعويض موږ ته د \LR{$\omega+\omega$} شتون
ثابتوي؛ نو د \LR{$V_{\omega+\omega}$} شتون هم ورسره ثابتولے شي۔ دا يوازې
همدومره نۀ دي۔ د هر هغه ښه ترتيب دپاره چې تعريفولے ئې شو،
\olref[ordinals][ordtype]{thmOrdinalRepresentation} وايي چې له هغه
ښه ترتيب سره آيزومورف يو \LR{$\alpha$} شته، او له دې امله \LR{$V_\alpha$} هم
شته۔ نو تعويض په څرګنده توګه دا تضمينوي چې د سټونو سلسله بايد
\emph{ډېره لوړه} وي۔
\olpseganchor{OLP-0567-B007}{end}

\olpseganchor{OLP-0567-B008}{start}
په راتلونکو څو برخو کښې، او بيا په
\olref[card-arithmetic][fix]{sec} کښې، به لا ښه پوه شو چې تعويض د
سټونو سلسله \emph{څومره} لوړېدو ته اړ باسي۔ اوس دپاره ساده خبره دا ده
چې تعويض په رښتيا \emph{توجيه} غواړي!
\olpseganchor{OLP-0567-B008}{end}

\olpseganchor{OLP-0568-B003}{start}
\olfileid{sth}{replacement}{extrinsic}
\olsection[بهرني دليلونه]{د تعويض په اړه بهرني دليلونه}
\olpseganchor{OLP-0568-B003}{end}

\olpseganchor{OLP-0568-B004}{start}
لومړے د تعويض د توجيې يوه \emph{بهرنۍ} هڅه څېړو۔
\begin{explain}
  بولوس د تعويض د منلو دليل د هغو قضيو په کار کښې ويني چې په
  دغه اصل ثابتېږي۔ تعويض د سټونو د تکراري تصور په اړه ځينې
  قضيې ثابتوي، د نامتناهي عددونو د يوې د منلو وړ تيورۍ د جوړولو
  وس راکوي او پر ښه‌بنسټيزو اړيکو استقرايي تعريفونه توجيه کوي۔
  د بولوس په ارزونه کښې دغه ګټې څرګندې دي او کومه ناغوښتې
  پايله ورته نۀ ښکاري۔ \mbox{\LR{\latinfont \citep[229]{Boolos1971}}}
\end{explain}
د بولوس د نظر لنډيز دا دے چې تعويض بايد د هغه د پايلو په مټ توجيه
کړو۔ هغه مشخصې پايلې همدا دي چې په تېرو څو څپرکو کښې مو وڅېړلې۔
تعويض موږ ته دا ثابته کړه چې د فون نويمن ترتيبي عددونه د ښه ترتيب د
ډول ډېر ښه استازي دي (دا زموږ «د نامتناهي عددونو د قناعت وړ، که څه
هم نۀ بشپړه، تيوري» ده)۔ تعويض مو د \LR{$V_\alpha$} ګانو تعريف، د رتبې
مفهوم جوړول، او د \LR{$\in$}-استقرا ثبوت هم ممکن کړ (دا «د تکراري تصور په
اړه قضيې» دي)۔ بالاخره، تعويض د نامتناهيت هاخوا بازګښت د قضیې ثبوت
ممکن کړ (دا «په ښه‌بنسټيزو اړيکو باندې استقرايي تعريفونه» دي)۔
\olpseganchor{OLP-0568-B004}{end}

\olpseganchor{OLP-0568-B005}{start}
دا بېشکه غوښتل شوې پايلې دي۔ خو آيا همدغې پايلې د تعويض د
\emph{توجيې} دپاره بس دي؟ \emph{نۀ}۔ يا لږ تر لږه، نېغ په نېغه نۀ۔
\olpseganchor{OLP-0568-B005}{end}

\olpseganchor{OLP-0568-B006}{start}
يوه ساده ستونزه دا ده: که څه هم مو د تعويض ځينې مطلوبې پايلې
وښودلې، ډېرې ئې له نورو لارو هم ترلاسه کولے شو۔ دا خبره دومره مشهوره
نۀ ده څومره چې بايد وي؛ نو د حالت د روښانولو دپاره لږ تم کېږو۔
\olpseganchor{OLP-0568-B006}{end}

\olpseganchor{OLP-0568-B007}{start}
د سټونو يوه ساده تيوري شته چې د پړاوونو تيوري، يا په لنډه \LR{$\LT$}
بلل کېږي۔\footnote{د \LR{$\LT$} لومړنۍ بڼې \mbox{\LR{\latinfont \citet{Montague1965}}} او
\mbox{\LR{\latinfont \citet{Scott1974}}} وړاندې کړې دي؛ \mbox{\LR{\latinfont \citet{Potter2004}}} ساده کړه او
د کتاب په کچه ئې وڅېړله؛ او \mbox{\LR{\latinfont \citet{ButtonLT1}}} په وروستيو کښې \LR{$\LT$}
لا ساده کړې ده۔} د \LR{$\LT$} اصول يوازې د غړو له مخې برابري، بېلونه او دا
وينا دي چې هر سټ د کوم \emph{پړاو} فرعي سټ دے۔ «پړاو» په داسې کره ډول
تعريف شوے چې پړاوونه زموږ د \LR{$V_\alpha$} ګانو په شان چلند کوي۔ نو \LR{$\ZF$}
د \LR{$\LT$} ثبوت ورکوي، خو \LR{$\LT$} له \LR{$\ZF$} څخه \emph{ډېره} کمزورې ده۔
په حقيقت کښې \LR{$\LT$} جوړې، د فرعي سټونو سټونه، نامتناهيت يا تعويض نۀ
راکوي۔ \LR{$\Zr$} دې د \LR{$\LT$} سره د نامتناهيت او د فرعي سټونو د سټونو د
اصلو د ورزياتولو پايله وي؛ له دې جوړې هم راځي، نو \LR{$\Zr$} لږ تر لږه د
\LR{$\Z$} هومره قوي ده۔ خو \LR{$\Zr$} په حقيقت کښې له \LR{$\Z$} څخه په کلکه قوي
ده، ځکه دا وينا ورزياتوي چې هر سټ يوه رتبه لري (له همدې امله ئې
\LR{$\Zr$} بولم)۔ په رښتيا \LR{$\Zr$} دا پايلې راکوي: د ترتيبي عددونو بشپړه د
قناعت وړ تيوري؛ هغه پايلې چې سلسله په ښه مرتب پړاوونو وېشي؛ د
\LR{$\in$}-استقرا ثبوت؛ او د نامتناهيت هاخوا بازګښت يوه \emph{بڼه}۔
\olpseganchor{OLP-0568-B007}{end}

\olpseganchor{OLP-0568-B008}{start}
نو په لنډه توګه: که څه هم بولوس په دې نۀ پوهېده، د هغو مطلوبو پايلو مناسبې بڼې چې يادوي ئې له تعويض
\emph{پرته} هم ترلاسه کېدې؛ يوازې \LR{$\Zr$} ئې د \LR{$\Z$} پر ځاے
کارول پکار وو۔
\olpseganchor{OLP-0568-B008}{end}

\olpseganchor{OLP-0568-B009}{start}
(دې ټولو خبرو سره سره، ولې مو تاسو ته د \LR{$\ZF$} د ښوولو دوديزه لاره غوره کړه،
د \LR{$\LT$} او \LR{$\Zr$} د ښوولو پر ځاے؟ دوه دليلونه شته۔ لومړے، د
تاريخي لاملونو له امله له \LR{$\LT$} څخه پيلول ډېر نااشنا دي؛ غوښتل مو چې
تاسو د سټ تيورۍ نور دوديز بحثونه هم ولوستلای شئ۔ دويم، کله چې د
\LR{$\LT$} او \LR{$\Zr$} د ژورې مطالعې لپاره چمتو شئ، د \mbox{\LR{\latinfont \citealt{Potter2004}}} او
\mbox{\LR{\latinfont \citealt{ButtonLT1}}} ليکنې پخپله لوستلے شئ۔)
\olpseganchor{OLP-0568-B009}{end}

\olpseganchor{OLP-0568-B010}{start}
البته، څرنګه چې \LR{$\Zr$} له \LR{$\ZF$} څخه په کلکه کمزورې ده، ځينې
داسې پايلې شته چې \LR{$\ZF$} ئې ثابتوي خو \LR{$\Zr$} ئې بې پرېکړې پرېږدي۔
نو څوک د دغو پايلو په يادولو د تعويض بهرنۍ توجيه هڅولے شي۔ خو کله
چې د \LR{$\Zr$} کارول زده کړئ، بيا د هغو خبرو ډېر مثالونه موندل ګران وي
چې (الف) تعويض ئې فيصله کوي او له هغه پرته نۀ فيصله کېږي، او (ب) په
بداهتي توګه رښتيا وي۔ (په دې اړه \mbox{\LR{\latinfont \citealt[\S13.2]{Potter2004}}}
وګورئ۔)
\olpseganchor{OLP-0568-B010}{end}

\olpseganchor{OLP-0568-B011}{start}
پايله دا ده: د تعويض د قناعت وړ بهرنۍ توجيې دپاره بايد داسې
پايله وموندل شي چې له تعويض \emph{پرته} نۀ شي ترلاسه کېدے۔ دا اسانه
کار نۀ دے۔
\olpseganchor{OLP-0568-B011}{end}

\olpseganchor{OLP-0568-B012}{start}
اوس يوه بله ستونزه وڅېړو چې د تعويض د هرې خالصې بهرنۍ توجيې
پر وړاندې درېږي۔ (ښايي دا له لومړۍ ستونزې هم ژوره وي۔) بولوس يوازې
دا نۀ وايي چې تعويض ډېرې مطلوبې پايلې لري؛ هغه دا هم وايي چې
«(ظاهراً) هېڅ ناغوښتلې» پايله نۀ لري۔ خو دلته د «ظاهراً» قيد بېخي
اساسي دے۔
\olpseganchor{OLP-0568-B012}{end}

\olpseganchor{OLP-0568-B013}{start}
راياد کړئ چې دلته څنګه راورسېدو: د ناپام جامعيت اصل له تناقض سره
مخ شو، او موږ په ځواب کښې د سټ تجمعي-تکراري تصور غوره کړ۔ له دې تصور
سره يوه کيسه تړلې ده چې هيله لرو د هغه سازګاري راته ډاډمنه کړي۔ خو
که تعويض له همدې کيسې نه توجيه نۀ کړو، تر اوسه د \LR{$\ZF$} د سازګارۍ
د باور دپاره دليل نۀ لرو۔ يا په کره توګه: د \LR{$\ZF$} د سازګارۍ د باور دپاره يوازې دا
(ښايي تصادفي) واقعيت راسره دے چې چا تر \emph{اوسه} تناقض نۀ دے موندلے۔ د بولوس خبره
په همدغه معنٰی دې ته راټيټېږي: «(ظاهراً) \LR{$\ZF$} سازګاره ده»۔ تر دې
د سازګارۍ ډېر ډاډ بايد وغواړو۔
\olpseganchor{OLP-0568-B013}{end}

\olpseganchor{OLP-0568-B014}{start}
دا ستونزه به د تعويض د هرې \emph{خالصې} بهرنۍ توجيې سره وي؛
يعنې له هغې توجيې سره چې يوازې د \LR{$\ZF$} په (معلومو) پايلو ولاړه وي۔
نو د تعويض \emph{دننۍ} توجيه غواړو: داسې توجيه چې وښيي د سټونو په اړه
زموږ کيسه موږ په يوه ډول تعويض ته «له مخکښې» ژمن کړي وو۔
\olpseganchor{OLP-0568-B014}{end}

\olpseganchor{OLP-0569-B003}{start}
\olfileid{sth}{replacement}{limofsize}
\olsection{د اندازې محدوديت}
\olpseganchor{OLP-0569-B003}{end}

\olpseganchor{OLP-0569-B004}{start}
ښايي د تعويض د «دننۍ» توجيې تر ټولو عامه هڅه د لاندې تصور له لارې وي:
\begin{enumerate}
	\item[] \limofsize. هر څيزونه يو سټ جوړوي، په دې شرط چې ډېر زيات نۀ وي.
\end{enumerate}
دا اصل د تعويض د توجيې لپاره داسې کارول کېږي۔ فرض کړئ
د تعويض په کارولو سره
\LR{$\funimage{\tau}{A} = \Setabs{\tau(x)}{x \in A}$} جوړوئ۔
ټاکلې تابعه له خپلې ساحې څخه پر خپل انځور پوره ده۔ د ټاکنې
له اصل سره، يا که ساحه ښه ترتيب ولري، د انځور د هر غړي
لپاره يو مخکنے غړے ټاکلے شو؛ په ښه ترتيب کښې تر ټولو کوچنے
مخکنے غړے اخلو۔ دا ټاکنه يو پر يو تابعه ورکوي، نو
\LR{$\cardle{\funimage{\tau}{A}}{A}$}۔ که د اندازې شرط په خپله د
پوره تابعې د انځور له مخې بيان کړو، د يو پر يو تابعې دغه
ټاکنه پکار نۀ ده۔ په دې توګه، که د \LR{$A$} غړي
د سټ جوړولو دپاره تر اندازې ډېر نۀ وو، نو د هغو انځورونه هم د
\LR{$\funimage{\tau}{A}$} د جوړولو دپاره تر اندازې ډېر نۀ دي۔
\olpseganchor{OLP-0569-B004}{end}

\olpseganchor{OLP-0569-B005}{start}
د تعويض د توجيې دپاره د \limofsize{} په کارولو کښې ښکاره ستونزه دا ده
چې تر اوسه مو د \limofsize په څېر هېڅ اصل \emph{نۀ} دے اېښے۔ سربېره
پر دې، کله چې مو د سټ د تجمعي-تکراري تصور کيسه په
\crefrange{sth:story::chap}{sth:z::chap} کښې وکړه، په هغې کښې
د \limofsize لوري ته هېڅ \emph{اشاره} هم نۀ وه۔ بولوس دا احتمال وړاندې کړ چې د سټ تيورۍ بېلابېل اصول له لږ تر لږه دوو جلا مفکورو څخه توجيه اخلي \mbox{\LR{\latinfont \citeyearpar[p.~19]{Boolos1989}}}۔ له يوې خوا، د سټ
د تجمعي-تکراري تصور شاوخوا مفکورې د~\LR{$\Z$} توجيه کول غواړي۔ له بلې خوا،
\limofsize{} د تعويض توجيه کول غواړي۔
\olpseganchor{OLP-0569-B005}{end}

\olpseganchor{OLP-0569-B006}{start}
خو ستونزه يوازې دا نۀ ده چې تر اوسه د \limofsize په اړه \emph{چپ}
وو۔ اصلي ستونزه دا ده چې \limofsize{}، په همدغه بڼه چې بيان شو، د سټ
له تجمعي-تکراري تصور سره چندان سم نۀ ښکاري۔ ځکه په دې اصل کښې په
\emph{پړاوونو} کښې د سټونو د جوړېدو تصور هېڅ نۀ يادېږي۔
\olpseganchor{OLP-0569-B006}{end}

\olpseganchor{OLP-0569-B007}{start}
د دغو «دوو فکرونو» تاريخ ته په کتو دا حيرانوونکې نۀ ده؛ يعنې د سټ
تجمعي-تکراري تصور او \limofsize دوه بېلې لارې دي چې د سټ تيورۍ
تناقضونه پرې بندېږي۔ تجمعي-تکراري تصور د رسل تناقض په لنډه توګه داسې
بندوي: \emph{موږ بايد د رسل د سټ د شتون تمه نۀ واے کړې، ځکه په هېڅ
پړاو کښې به «جوړ» شوے نۀ واے}۔ په مقابل کښې، \limofsize{} د رسل سټ
په دې وينا ردوي: \emph{موږ بايد د رسل د سټ د شتون تمه نۀ واے کړې،
ځکه هغه به تر اندازې ډېر لوے واے}۔
\olpseganchor{OLP-0569-B007}{end}

\olpseganchor{OLP-0569-B008}{start}
نو په دې بڼه ئې په ښکاره وايو: د تناقضونو د ځواب په توګه،
\limofsize{} \emph{ډېرې} زياتې توجيې ته اړ دے۔ د بېلګې په توګه، د رسل
د تناقض دا بڼه وګورئ: \emph{هېڅ داسې پګ سپے نشته چې عيناً هغه او يوازې
هغه پګ سپي بوي کړي چې خپل ځان نۀ بوي کوي} (\olref[sth][story][rus]{sec}
وګورئ)۔ که پوښتنه وکړئ «ولې داسې پګ سپے نشته؟»، دا ښه ځواب نۀ دے چې
هغه به ډېر زيات پګ سپي بوي کول غواړي۔ نو د رسل د سټ د نۀ شتون دا به
ولې ښه بداهتي توضيح وي چې د شتون دپاره به «تر اندازې ډېر لوے» واے؟
\olpseganchor{OLP-0569-B008}{end}

\olpseganchor{OLP-0569-B009}{start}
په لنډه توګه، که د \limofsize د «بداهتي» انګېزې په اړه لږ حيران ياست،
دا د پوهېدو وړ ده۔
\olpseganchor{OLP-0569-B009}{end}

\olpseganchor{OLP-0570-B003}{start}
\olfileid{sth}{replacement}{absinf}
\olsection{تعويض او «مطلق نامتناهيت»}
\olpseganchor{OLP-0570-B003}{end}

\olpseganchor{OLP-0570-B004}{start}
اوس \limofsize{} شاته پرېږدو او د تعويض د توجيې د نورو
(دننيو) هڅو يوه ډله څېړو۔ دغه هڅې دا خبره جدي نيسي چې سټونه په
پړاوونو کښې جوړېږي۔
\olpseganchor{OLP-0570-B004}{end}

\olpseganchor{OLP-0570-B005}{start}
کله چې مو لومړے تکراري بهير بيان کړ، داسې اصول مو وړاندې کړل چې په
هر پړاو کښې پېښېدونکے کار ئې روښانولو۔ دا \stageshier، \stagesord او
\stagesacc وو۔ وروسته مو داسې اصول ورزيات کړل چې د پړاوونو د شمېر
په اړه څه وايي: \stagessucc{} وايي چې د سټ جوړولو بهير هېڅکله نۀ
ختمېږي، او \stagesinf{} وايي چې بهير تر نامتناهيتم پړاو پورې رسېږي۔
\olpseganchor{OLP-0570-B005}{end}

\olpseganchor{OLP-0570-B006}{start}
دا مناسبه ده چې ووايو دغه وروستي دوه اصول ښايي له يوه لا پراخ اصل
څخه راووځي، لکه:
\begin{enumerate}
	\item[] \stagesinex. په مطلق ډول نامتناهي ډېر پړاوونه شته؛
	سلسله تر هغه حده لوړه ده چې امکان لري.
\end{enumerate}
ښکاره ده چې دا يو غير رسمي اصل دے۔ خو که دا د سټ له تجمعي-تکراري
تصور څخه سمدستي \emph{نۀ هم لازمېږي}، له هغه سره خو \emph{سمون}
لري۔ لږ تر لږه، او د \limofsize پر خلاف، دا تصور ساتي چې سټونه
يو پړاو په بل پسې جوړېږي۔
\olpseganchor{OLP-0570-B006}{end}

\olpseganchor{OLP-0570-B007}{start}
اوس هيله دا ده چې له \stagesinex{} څخه د تعويض توجيه ترلاسه کړو۔
راځئ وګورو چې دا څنګه کېدے شي۔
\olpseganchor{OLP-0570-B007}{end}

\olpseganchor{OLP-0570-B008}{start}
په \olref[ordinals][idea]{sec} کښې مو وليدل چې داسې ښه ترتيب جوړول
اسانه دي چې د مفهوم له مخې بايد له \LR{$\omega+\omega$} سره آيزومورف
وي۔ په بل عبارت، په اسانه داسې پړاو په پړاو تکراري بهير تصورولے شو
چې د ترتيب ډول ئې د مفهوم له مخې \LR{$\omega+\omega$} وي۔ نو که مو
\stagesinex منلے وي، خامخا بايد ومنو چې د سلسلې لږ تر لږه
\LR{$\omega+\omega$}-م پړاو، يعنې \LR{$V_{\omega+\omega}$}، شته؛ ځکه سلسله
تر دې ځايه \emph{غځېدې} شي۔
\olpseganchor{OLP-0570-B008}{end}

\olpseganchor{OLP-0570-B009}{start}
دا فکر داسې عامېږي: د هر ښه ترتيب دپاره بايد د تکراري سلسلې د
جوړولو بهير لږ تر لږه د هماغه ښه ترتيب تر اوږدوالي ورسېږي۔ دا خبره
په دې هم تضمينولے شو چې \olref[ordinals][ordtype]{thmOrdinalRepresentation}
د يو \emph{اصل} په توګه واخلو۔ نو هر ښه ترتيب به د فون نويمن له
کوم ترتيبي عدد سره آيزومورف وي۔ څرنګه چې هر فون نويمن ترتيبي عدد د
خپلې رتبې سره برابر وي، \olref[ordinals][ordtype]{thmOrdinalRepresentation}
به راته ووايي چې کله هم په خپله سټ تيورۍ کښې يو ښه ترتيب بيانولے
شو، د سټ جوړولو تکراري بهير بايد تر هغه ښه ترتيب هاخوا ولاړ شي۔
\olpseganchor{OLP-0570-B009}{end}

\olpseganchor{OLP-0570-B010}{start}
دا فکر بېشکه د \stagesinex يوه پايله ښکاري۔ له بده مرغه، که مو موخه
له دې فکره د تعويض راايستل وي، يوه ساده تخنيکي خنډ مخې ته راځي:
تعويض له \olref[ordinals][ordtype]{thmOrdinalRepresentation} څخه په
کلکه قوي دے۔ (دا خبره \mbox{\LR{\latinfont \citet[\S13.2]{Potter2004}}} کړې ده؛ موږ به ئې
په \olref[sth][replacement][finiteaxiomatizability]{sec} کښې ثابته کړو۔)
\olpseganchor{OLP-0570-B010}{end}

\olpseganchor{OLP-0570-B011}{start}
نو که \stagesinex{} داسې وپېژنو چې تعويض ترې راووځي، دا اصل
\emph{يوازې} دومره نۀ شي ويلے چې سلسله تر هر ښه ترتيب هاخوا غځېږي۔
بايد تر دې لا قوي خبره وکړي۔ د همدې دپاره \mbox{\LR{\latinfont \cite{Shoenfield:AST}}} د
فکر يوه ډېره طبيعي پياوړې بڼه وړاندې کړه: سلسله له هېڅ سټ سره
\emph{هم‌پايه} نۀ ده۔\footnote{ښايي G\"odel ورته فکر وړاندې کړے وي؛
\mbox{\LR{\latinfont \citet[p.~223]{Potter2004}}} وګورئ۔ د G\"odel او
\mbox{\LR{\latinfont \citeauthor{Shoenfield:AST}}} د بحث دپاره \mbox{\LR{\latinfont \citet[90--5]{Incurvati2020}}}
وګورئ۔} په لږ تفصيل: که \LR{$\tau$} داسې نگاشت وي چې سټونه د سلسلې
پړاوونو ته لېږي، د هر سټ \LR{$A$} لپاره د \LR{$\tau$} ټول ارزښتونه
تر يوه پړاو لاندې راځي۔ په بل عبارت، په شېمايي ډول:
\begin{enumerate}
	\item[] \stagescofin. که \LR{$A$} يو سټ وي او \LR{$\tau(x)$} د هر
	\LR{$x \in A$} دپاره يو پړاو وي، نو داسې پړاو شته چې د هر
	\LR{$\tau(x)$} نه وروسته راځي، د هر \LR{$x \in A$} دپاره.
\end{enumerate}
په \LR{$\ZF$} کښې تعويض د دغو پړاوونو د شاخصونو سټ راکوي۔
د شاخصونو تر سخت پاسني حد پورې پړاو د \stagescofin
يوه مناسبه رسمي بڼه ثابتوي۔
برعکس، په غير رسمي استدلال سره ښودلای شو چې \stagescofin{} تعويض
توجيه کوي۔\footnote{د \stagescofin د کومې رسمي بڼې نه تعويض ثابتول
به ګران وي، ځکه \LR{$\Z$} په خپله د پړاوونو د تعريف دپاره کافي قوي نۀ
ده؛ نو څرګنده نۀ ده چې \stagescofin به څنګه رسمي کړو۔ يوه لاره دا
ده چې د \LR{$\LT$} په کوم پراخ نظام کښې کار وکړو، لکه په
\olref[sth][replacement][extrinsic]{sec} کښې چې بحث شوے دے۔}
فرض کړئ \LR{$(\forall x \in A)\lexists![y][\phi(x,y)]$}۔ نو د هر
\LR{$x \in A$} دپاره \LR{$\sigma(x)$} هغه \LR{$y$} وټاکئ چې \LR{$\phi(x,y)$} پوره
کوي، او \LR{$\tau(x)$} هغه پړاو وټاکئ چې \LR{$\sigma(x)$} په لومړي ځل پکښې
جوړېږي۔ د \stagescofin له مخې داسې پړاو \LR{$V$} شته چې
\LR{$(\forall x \in A)\tau(x)\in V$}۔ اوس، څرنګه چې هر
\LR{$\tau(x) \in V$} دے او \LR{$\sigma(x) \subseteq \tau(x)$}، د
وروستي پړاو توانمن‌والے وايي چې جوړ شوے سټ د هماغه پړاو
غړے دے۔ نو د بېلونې په اصل سره
\LR{$\Setabs{y
\in V}{(\exists x \in A)\sigma(x) = y} = \Setabs{y}{(\exists x \in
A)\phi(x,y)}$} ترلاسه کولے شو۔
\olpseganchor{OLP-0570-B011}{end}

\olpseganchor{OLP-0570-B012}{start}
\begin{prob}
	\stagescofin{} د \LR{$\ZF$} دننه په رسمي ډول بيان کړئ.
\end{prob}
\olpseganchor{OLP-0570-B012}{end}

\olpseganchor{OLP-0570-B013}{start}
نو \stagescofin{} تعويض توجيه کوي۔ او لږ تر لږه دا معقوله ښکاري
چې \stagesinex{}، \stagescofin توجيه کوي۔ فرض کړئ \stagescofin{}
نۀ پوره کېږي۔ نو سلسله له کوم سټ~\LR{$A$} سره هم‌پايه ده: داسې نگاشت
\LR{$\tau$} لرو چې د هر پړاو~\LR{$S$} دپاره کوم \LR{$x \in A$} شته چې
\LR{$S \in \tau(x)$}۔ په دې حالت کښې د سلسلې د ادعا شوي «مطلق نامتناهيت»
د اندازه کولو يوه لاره لرو: د \LR{$\tau$} په \LR{$A$} باندې
ترلاسه شوي پړاوونه په سلسله کښې \emph{هر څومره هاخوا}
رسېږي؛ هر ورکړل شوے پړاو له دغو پړاوونو څخه د يوه غړے دے۔ دا د سلسلې د «مطلق نامتناهيت» فکر کمزورے
کوي۔ په معکوس استدلال: \stagesinex{}، \stagescofin لازموي، او هغه
بيا تعويض توجيه کوي۔
\olpseganchor{OLP-0570-B013}{end}

\olpseganchor{OLP-0570-B014}{start}
دا د تعويض د دننۍ توجيې يوه رښتيا هيله‌بښونکې هڅه ده۔ خو په پای کښې
دا کار کوي که نۀ، پرېکړه ئې تاسو ته پرېږدو۔
\olpseganchor{OLP-0570-B014}{end}

\olpseganchor{OLP-0571-B003}{start}
\olfileid{sth}{replacement}{ref}
\olsection{تعويض او انعکاس}
\olpseganchor{OLP-0571-B003}{end}

\olpseganchor{OLP-0571-B004}{start}
د \stagesinex له لارې د تعويض د توجيې زموږ وروستۍ هڅه له يوې ژورې
او ښکلې پايلې پيلېږي:\footnote{يادونه: ټول فورمولونه پاراميټرونه
لرلے شي، مګر چې په څرګند ډول بل څه ويل شوي وي.}
%In this, I will use overlining, such as $x_1, \ldots, x_n$, to
%abbreviate ``$x_1, \ldots, x_n$'': 
\begin{thm}[د انعکاس شېما]\ollabel{reflectionschema}
د هر فورمول \LR{$\phi$} دپاره:
\[
\forall \alpha \exists \beta > \alpha (\forall x_1 \ldots, x_n \in
V_\beta)(\phi(x_1, \ldots, x_n) \liff \phi^{V_\beta}(x_1, \ldots, x_n))
\]
\end{thm}
\noindent 
لکه په \olref[sth][replacement][strength]{formularelativization} کښې،
\LR{$\phi^{V_\beta}$} د \LR{$\phi$} د ټولو کميت ټاکونکو د \LR{$V_\beta$} سټ ته
د محدودولو پايله ده۔ انعکاس وايي چې له هر ټاکلي پړاو څخه وروسته
داسې \emph{پيلنۍ برخه} شته چې د پاراميټرونو د هرې هغې ټاکنې
لپاره چې غړي يې په هماغه برخه کښې وي، د \LR{$\phi$} د رښتيا يا دروغو ارزښت پکښې د ټولې سلسلې د
هماغه \LR{$\phi$} له ارزښت سره يو شان وي۔
\olpseganchor{OLP-0571-B004}{end}

\olpseganchor{OLP-0571-B005}{start}
\mbox{\LR{\latinfont \citet{Montague1961}}} او \mbox{\LR{\latinfont \citet{Levy1960}}} وښودله چې د تعويض او
انعکاس مناسبې رسمي بڼې د~\LR{$\Z$} په نسبت همارزښته دي؛ نو له دغو دوو
هر يو چې ورزيات کړو، \LR{$\ZF$} ترلاسه کېږي۔ (دا پايلې په
\olref[sth][replacement][refproofs]{sec} کښې ثابتوو۔) له دې
همارزښت څخه هيله کېدے شي چې د \stagesinex{} په مټ انعکاس او تعويض
داسې توجيه کړو: د \stagesinex په منلو سره سلسله بايد ډېره، ډېره
لوړه وي؛ دومره لوړه چې د هغې په اړه هېڅ وينا ئې د لوړوالي د حد
ټاکلو دپاره بس نۀ وي۔ د پاراميټرونو پرته د جملې په حالت کښې،
که جمله~\LR{$\phi$} په ټوله سلسله کښې رښتيا وي، نو له هر ټاکلي
پړاو وروسته په يوه پيلنۍ برخه کښې هم رښتيا وي۔ دا د انعکاس
د عمومي شېما د جملو ځانګړے حالت دے۔
\olpseganchor{OLP-0571-B005}{end}

\olpseganchor{OLP-0571-B006}{start}
بيا هم، دا د تعويض د دننۍ توجيې يوه رښتيا هيله‌بښونکې هڅه ښکاري۔
خو دلته د دې په اړه تر دې ډېر څه ويل پکار دي۔ اوس پخپله پرېکړه
وکړئ چې دا هڅه بريالۍ ده که نۀ۔\footnote{کېدے شي د نورو معلوماتو
دپاره \mbox{\LR{\latinfont \citet[95--100]{Incurvati2020}}} هم ولولئ.}
\olpseganchor{OLP-0571-B006}{end}

\olpseganchor{OLP-0572-B003}{start}
\olfileid{sth}{replacement}{refproofs}
\olsection{ضميمه: د تعويض اړوندې پايلې}
\olpseganchor{OLP-0572-B003}{end}

\olpseganchor{OLP-0572-B004}{start}
په دې برخه کښې به د \LR{$\ZF$} دننه انعکاس ثابت کړو۔ دا به هم وښيو چې
انعکاس په يوه معنا له تعويض سره همارزښته دے۔ بيا به د دې ټولو يوه
په زړه پورې پايله د انعکاس/تعويض د قوت په اړه ثابته کړو۔
\emph{خبردارے: دا په دې درسي کتاب کښې تر ټولو پرمختللې رياضيکي برخه ده۔}
\olpseganchor{OLP-0572-B004}{end}

\olpseganchor{OLP-0572-B005}{start}
له يوې لمې پيل کوو چې د لنډيز دپاره د متغيرونو يا څيزونو د لړيو
د ښودلو \emph{پر سر کرښې} نښه کاروي۔ يعنې «\LR{$\overline{a}_k$}» د
«\LR{$a_{k_1}$}، \dots، \LR{$a_{k_n}$}» لنډه نښه ده، او \LR{$n$} له سياقه
ټاکل کېږي۔ په دې بحث کښې يو پړاو ته د داسې لړۍ غړيتوب
دا معنا لري چې د لړۍ هره برخه د هماغه پړاو غړې ده۔
\olpseganchor{OLP-0572-B005}{end}

\olpseganchor{OLP-0572-B006}{start}
\begin{lem}\ollabel{lemreflection}
د هر \LR{$1 \leq i \leq k$} دپاره دې \LR{$\phi_i(\overline{v}_i, x)$}
يو فورمول وي۔ نو د هر \LR{$\alpha$} دپاره داسې \LR{$\beta > \alpha$} شته چې
د هر \LR{$\overline{a}_1,
\ldots, \overline{a}_k \in V_\beta$} او هر \LR{$1 \leq i \leq k$}
دپاره:
\[
	\exists x\phi_i(\overline{a}_i, x) \rightarrow (\exists x \in V_\beta) \phi_i(\overline{a}_i, x)
\]
\end{lem}
\olpseganchor{OLP-0572-B006}{end}

\olpseganchor{OLP-0572-B007}{start}
\begin{proof}
يو ترم \LR{$\mu$} داسې تعريفوو: \LR{$\mu(\overline{a}_1, \ldots,
\overline{a}_k)$} تر ټولو لږ پړاو \LR{$V$} دے چې د هر \LR{$1 \leq i \leq k$}
دپاره دا ټول شرطي بيانونه پوره کوي:
\[
\exists x\phi_i(\overline{a}_i, x) \rightarrow (\exists x \in V) \phi_i(\overline{a}_i, x)
\]
د هر رښتيني وجودي شرط لپاره يو شاهد واخلئ۔ شرطونه
متناهي دي او هر شاهد په يوه پړاو کښې غړے دے۔ داسې پړاو
واخلئ چې دا ټول شاهدان پکښې وي۔ نو د پړاوونو د شاخصونو
له ښه ترتيبه
\LR{$\mu(\overline{a}_1, \ldots, \overline{a}_k)$} د ټولو
\LR{$\overline{a}_1, \ldots, \overline{a}_k$} دپاره شته۔ د يوه ټاکلي
پړاو د متناهيو مرتب پاراميټري لړيو ټولګه سټ دے۔ تعويض د هغو
اړوند لږ پړاوونه په يو سټ کښې راټولوي۔ اوس د تعويض
او زموږ د بازګښت د قضيې په مټ تعريف کړئ:
\begin{align*}
	S_0 & = V_{\alpha+1}\\
	S_{n+1} & = S_n \cup \bigcup
	\Setabs{\mu(\overline{a}_1, \ldots, \overline{a}_k)}
	{\overline{a}_1, \ldots, \overline{a}_k \in S_n} \\
	S &= \bigcup_{m < \omega} S_m.
\end{align*}
د هرې ناتشې پړاويزې کورنۍ اتحاد د هغې د شاخصونو د
تر ټولو لږ پاسني حد پړاو دے۔ د تالي په معادله کښې مخکنے
پړاو هم په اتحاد کښې شامل دے، نو له هماغه پړاوه شاته
نۀ ځو۔ ځکه هر \LR{$S_n$}، او له دې امله پخپله \LR{$S$} هم،
له \LR{$V_\alpha$} وروسته يو پړاو دے۔ اوس
\LR{$\overline{a}_1$}، \dots،~\LR{$\overline{a}_k \in S$} وټاکئ؛ نو
داسې \LR{$n < \omega$} شته چې \LR{$\overline{a}_1$}، \dots،
\LR{$\overline{a}_k \in S_n$}۔ کوم \LR{$1 \leq i \leq k$} وټاکئ او فرض
کړئ چې \LR{$\exists x \phi_i(\overline{a}_i,x)$}۔ د جوړونې له مخې
\LR{$(\exists x \in \mu(\overline{a}_1,\ldots, \overline{a}_k))
\phi_i(\overline{a}_i, x)$}؛ نو
\LR{$(\exists x \in S_{n+1})\phi_i(\overline{a}_i, x)$}، او بيا
\LR{$(\exists x \in S)\phi_i(\overline{a}_i, x)$}۔ نو \LR{$S$} زموږ
\LR{$V_\beta$} دے۔
\end{proof}
\noindent
اوس \olref[sth][replacement][ref]{reflectionschema} په نېغه توګه
ثابتولے شو:
\olpseganchor{OLP-0572-B007}{end}

\olpseganchor{OLP-0572-B008}{start}
\begin{proof}[ثبوت]
\LR{$\alpha$} وټاکئ۔ بې له دې چې عموميت کم شي، فرضولے شو چې د \LR{$\phi$}
يوازيني وصلي نښانې \LR{$\exists$}، \LR{$\lnot$} او \LR{$\land$} دي (ځکه د بيان
دپاره بس دي)۔ \LR{$\psi_1, \ldots, \psi_k$} دې د \LR{$\phi$} ټول فرعي
فورمولونه د پېچلتيا په ترتيب وښيي، څو \LR{$\psi_k = \phi$} شي۔
د \olref{lemreflection} له مخې داسې \LR{$\beta > \alpha$} شته چې د هر
\LR{$\overline{a}_i \in V_\beta$} او هر \LR{$1 \leq i \leq k$} دپاره:
\begin{align}\label{reflectionnicelybehaved}
	\exists x\psi_i(\overline{a}_i, x) \rightarrow
	(\exists x \in V_\beta) \psi_i(\overline{a}_i, x)\tag{*}
\end{align}
د \LR{$\psi_i$} د پېچلتيا په استقرا به وښيو چې
\LR{$\psi_i(\overline{a}_i) \leftrightarrow
\psi_i^{V_\beta}(\overline{a}_i)$} د هر
\LR{$\overline{a}_i \in V_\beta$} دپاره۔ که \LR{$\psi_i$} اټومي وي،
خبره ښکاره ده۔ دغه دوه‌اړخيزه همدارنګه ښيي چې که \LR{$\psi_i$} د هغو
فرعي فورمولونو نفي يا عطف وي چې دا خاصيت لري، نو \LR{$\psi_i$} هم دا
خاصيت لري۔ يوازې د کميت ټاکلو حالت پاتې کېږي۔
\LR{$\overline{a}_i \in V_\beta$} وټاکئ؛ بيا:
\begin{align*}
	(\exists x \psi_i(\overline{a}_i, x))^{V_\beta}
	&\text{\RL{ هله او يوازې هله چې }}
	(\exists x \in V_\beta)\psi_i^{V_\beta}(\overline{a}_i, x)
	&&\text{\RL{د تعريف له مخې}}\\
	&\text{\RL{ هله او يوازې هله چې }}
	(\exists x \in V_\beta)\psi_i(\overline{a}_i,  x)
	&&\text{\RL{د استقرا د فرض له مخې}}\\
	&\text{\RL{ هله او يوازې هله چې }}
	\exists x \psi_i(\overline{a}_i, x)
	&&\text{\RL{د \eqref{reflectionnicelybehaved} له مخې}}
\end{align*}
په دې وجودي ګام کښې فرعي فورمول د کميت ټاکونکي دننه
فورمول دے؛ د استقرا فرض د هغه په ټولو پاراميټرونو، د
تړل کېدونکي متغير د ارزښت په ګډون، پلي کېږي۔ د وروستي
دوه اړخيز شرط بل لور د محدود شاهد له نامحدود شاهد
کېدو راځي۔ استقرا بشپړه شوه؛ ځکه \LR{$\psi_k = \phi$}،
غوښتل شوې پايله راځي۔
\end{proof}
\olpseganchor{OLP-0572-B008}{end}

\olpseganchor{OLP-0572-B009}{start}
موږ په \LR{$\ZF$} کښې انعکاس ثابت کړ۔ زموږ ثبوت په اصل کښې د
\mbox{\LR{\latinfont \citet{Montague1961}}} لاره تعقيبوي۔ اوس غواړو په \LR{$\Z$} کښې ثابت کړو
چې له انعکاسه تعويض راځي۔ ثبوت د \mbox{\LR{\latinfont \citet{Levy1960}}} لاره په ساده
بڼه تعقيبوي۔
\olpseganchor{OLP-0572-B009}{end}

\olpseganchor{OLP-0572-B010}{start}
څرنګه چې په \LR{$\Z$} کښې کار کوو، انعکاس په عين مخکنۍ بڼه نۀ شو
وړاندې کولے۔ هغه مو د «\LR{$V_\alpha$}» په نښه بيان کړے و، او دغه
پړاوونه په \LR{$\Z$} کښې نۀ شو تعريفولے
(\olref[sth][spine][zf]{sec} وګورئ)۔ نو د انعکاس يوه ظاهراً
کمزورې بڼه داسې وړاندې کوو:
\olpseganchor{OLP-0572-B010}{end}

\olpseganchor{OLP-0572-B011}{start}
\begin{defish}
\emph{کمزورے انعکاس.} د هر فورمول \LR{$\phi$} دپاره داسې متعدي سټ \LR{$S$}
شته چې \LR{$0$}، \LR{$1$} او د \LR{$\phi$} هر پاراميټر د \LR{$S$} غړي وي،
او \LR{$(\forall \overline{x} \in S)(\phi \liff \phi^S)$}۔
\end{defish}
\olpseganchor{OLP-0572-B011}{end}

\olpseganchor{OLP-0572-B012}{start}
له دې څخه د تعويض د ثبوت دپاره لومړے د
\mbox{\LR{\latinfont \citet[first
part of Theorem 2]{Levy1960}}} لاره تعقيبوو او ښيو
چې دوه فورمولونه په يو وخت «منعکس» کولے شو:
\olpseganchor{OLP-0572-B012}{end}

\olpseganchor{OLP-0572-B013}{start}
\begin{lem}[په \LR{$\Z + \text{\RL{کمزوري انعکاس}}$} کښې.]\ollabel{lem:reflect}
د هر دوو فورمولونو \LR{$\psi, \chi$} دپاره داسې متعدي سټ \LR{$S$} شته چې
\LR{$0$} او \LR{$1$} (او د فورمولونو هر پاراميټر) د \LR{$S$} غړي وي،
او \LR{$(\forall \overline{x} \in S)((\psi \liff \psi^S) \land (\chi
\liff \chi^S))$}۔
\end{lem}
\olpseganchor{OLP-0572-B013}{end}

\olpseganchor{OLP-0572-B014}{start}
\begin{proof}
\LR{$\phi$} دې \LR{$(z = 0 \land \psi) \lor (z = 1 \land \chi)$} فورمول وي۔
دوه‌حالته متغير له دواړو فورمولونو څخه تازه اخلو، څو د
هغو خپل ازاد متغيرونه بدل نۀ شي۔
\olpseganchor{OLP-0572-B014}{end}

\olpseganchor{OLP-0572-B015}{start}
دلته مو لنډه نښه کارولې ده؛ «\LR{$z = 0$}» بايد په بشپړه بڼه
«\LR{$\forall t\, t \notin z$}» او «\LR{$z =1$}» بايد
«\LR{$\forall s(s \in z \liff \forall t\, t \notin s)$}»
وليکو۔ خو څرنګه چې \LR{$0, 1 \in S$} او \LR{$S$} متعدي دے، دا فورمولونه
د \LR{$S$} دپاره \emph{مطلق} دي: که کميت ټاکونکي ئې \LR{$S$} ته محدود هم
کړو، پر هماغه څيز به پلي کېږي۔\footnote{په لا رسمي ډول، که
\LR{$\xi$} له دغو دوو فورمولونو يو وي، نو
\LR{$\xi(z) \liff \xi^S(z)$}۔}
\olpseganchor{OLP-0572-B015}{end}

\olpseganchor{OLP-0572-B016}{start}
د کمزوري انعکاس له مخې داسې مناسب \LR{$S$} شته چې:
\begin{align*}
	(\forall z, \overline{x} \in S)(&\phi \liff \phi^S)\\
	\text{\RL{يعنې }}(\forall z, \overline{x} \in S)(&((z = 0 \land \psi) \lor (z = 1 \land \chi)) \liff {}\\
	&\phantom{(}((z = 0 \land \psi) \lor (z = 1 \land \chi))^S)\\
	\text{\RL{يعنې }}(\forall z, \overline{x} \in S)(&((z = 0 \land \psi) \lor (z = 1 \land \chi))\liff {}\\
	&\phantom{(}((z = 0 \land \psi^S) \lor (z = 1 \land \chi^S)))\\
	\text{\RL{يعنې }}(\forall \overline{x} \in S)(&(\psi \liff \psi^S) \land (\chi \liff \chi^S))
\end{align*}
له دويمې دعوې درېيمه راځي، ځکه «\LR{$z = 0$}» او «\LR{$z=1$}»
د \LR{$S$} دپاره مطلق دي؛ څلورمه دعوه بيا له \LR{$0 \neq 1$} څخه راځي۔ د
دوه‌حالته متغير لپاره په وار سره صفر او يو وټاکئ۔ په
هر وار يوازې اړوند فورمول پاتې کېږي۔
\end{proof}\noindent اوس د \mbox{\LR{\latinfont \citet[Theorem 6]{Levy1960}}} د لارې
په تعقيب او ساده کولو سره تعويض ترلاسه کولے شو:
\olpseganchor{OLP-0572-B016}{end}

\olpseganchor{OLP-0572-B017}{start}
\begin{thm}[په \LR{$\Z$} + کمزوري انعکاس کښې]\label{thm:replacement}
د هر فورمول \LR{$\phi(v,w)$} او هر \LR{$A$} دپاره، که
\LR{$(\forall x \in A)\lexists![y][\phi(x,y)]$} وي، نو
\LR{$\Setabs{y}{(\exists x \in A)\phi(x,y)}$} شته۔
\end{thm}
\olpseganchor{OLP-0572-B017}{end}

\olpseganchor{OLP-0572-B018}{start}
\begin{proof}
داسې \LR{$A$} وټاکئ چې \LR{$(\forall x \in A)\lexists![y][\phi(x,y)]$}
وي، او دا فورمولونه تعريف کړئ:
\begin{align*}
	\psi &\text{\RL{ دا دے }} (\phi(x, z) \land A = A)\\
	\chi &\text{\RL{ دا دے }} \lexists[y][\phi(x, y)]
\end{align*}
د \olref{lem:reflect} په کارولو، څرنګه چې \LR{$A$} د \LR{$\psi$}
پاراميټر دے، داسې متعدي~\LR{$S$} شته چې \LR{$0, 1, A \in S$} وي
(د نورو پاراميټرونو ترڅنګ)، او:
\[
	(\forall x,z \in S)((\psi \liff \psi^S) \land (\chi \liff \chi^S))
\]
نو په ځانګړې توګه:
\begin{align*}
	(\forall  x, z \in S)(&\phi(x, z) \liff \phi^S(x, z))\\
	(\forall x \in S)(&\exists y\phi(x, y) \liff (\exists y \in S)\phi^S(x, y))
\end{align*}
دا سره يوځاے کړئ، او دا هم وګورئ چې \LR{$A \subseteq S$} دے، ځکه
\LR{$A \in S$} او \LR{$S$} متعدي دے:
\begin{align*}
	(\forall x \in A)(&\exists y\phi(x, y) \liff (\exists y \in S)\phi(x, y))
\end{align*}
اوس \LR{$(\forall x \in A)(\lexists![y \in S])\phi(x, y)$}؛ ځکه
\LR{$(\forall x \in A)\lexists![y][\phi(x, y)]$}۔ بيا بېلونه دا راکوي:
\LR{$\Setabs{y \in S}{(\exists x \in A) \phi(x, y)} = \Setabs{y}{(\exists
x \in A) \phi(x, y)}$}۔
\end{proof}
\olpseganchor{OLP-0572-B018}{end}

\olpseganchor{OLP-0573-B003}{start}
\olfileid{sth}{replacement}{finiteaxiomatizability}
\olsection{ضميمه: متناهي اصل‌بندي کېدنه}
دا څپرکے له تعويض څخه د ځينو پايلو په راايستلو پاے ته رسوو۔ لومړۍ
پايله د \mbox{\LR{\latinfont \citet{Montague1961}}} ده؛ پام وکړئ چې دا د \LR{$\ZF$} \emph{په دننه}
کښې ثبوت نه، بلکې د \LR{$\ZF$} \emph{په اړه} ثبوت دے:
\begin{thm}\ollabel{zfnotfinitely}
	د سازګارۍ په فرض، \LR{$\ZF$} په متناهيو اصلونو نۀ شي اصل‌بندي کېدلے۔ په عمومي ډول: که \LR{$\Th{T}$}
	متناهي وي او \LR{$\Th{T} \Proves \ZF$} وي، نو \LR{$\Th{T}$} ناسازګاره دے۔
\olpseganchor{OLP-0573-B003}{end}
	
\olpseganchor{OLP-0573-B004}{start}
	(دلته مو په پټه خپل بحث هغو لومړۍ مرتبې جملو ته محدود کړے چې يوازېنے
	غيرمنطقي لومړنی نښان يې \LR{$\in$} دے؛ او \LR{$\Th{T} \Proves \ZF$} په دې
	معنا ليکو چې \LR{$\Th{T} \Proves \phi$} د هر \LR{$\phi \in \ZF$}
	دپاره دے۔)
\end{thm}
\begin{proof}
	داسې متناهي \LR{$\Th{T}$} وټاکئ چې \LR{$\Th{T} \Proves \ZF$} وي۔ نو
	\LR{$\Th{T}$} انعکاس ثابتوي، يعنې \olref[sth][replacement][ref]{reflectionschema}۔
	څرنګه چې \LR{$\Th{T}$} متناهي دے، د يوې عطفې \LR{$\theta$} په بڼه يې بيا ليکلے شو۔
	د همدې فورمول په انعکاس سره \LR{$\Th{T} \Proves \exists \beta(\theta \liff \theta^{V_\beta})$}۔
	او ځکه چې په څرګنده \LR{$\Th{T} \Proves \theta$} دے، نو
	\LR{$\Th{T} \Proves \exists \beta\ \theta^{V_\beta}$} ترلاسه کوو۔
\olpseganchor{OLP-0573-B004}{end}
	
\olpseganchor{OLP-0573-B005}{start}
	اوس دې \LR{$\psi(X)$} د لاندې فورمول لنډه نښه وي:
	\[
		\theta^X \land X\text{\RL{ متعدي دے}} \land (\forall Y \in X)(Y\text{\RL{ متعدي دے}}\lif \lnot \theta^{Y})
	\]
	په لنډ ډول، دا وايي چې \LR{$X$} د \LR{$\theta$} متعدي مدل دے، او په دې اړه
	د \LR{$\in$} له مخې لږتر لږه دے۔ اوس دا په ياد ولرئ چې
	\LR{$\Th{T} \Proves \exists \beta\ \theta^{V_\beta}$}؛ د \LR{$\ZF$}
	دننه له هغو متعديو مدلونو يو د تر ټولو وړې رتبې سره ټاکلے شو۔
	د ټاکلو لپاره د لومړي موندل شوي مدل تر رتبې پورې مناسب شاخصونه
	په بېلونه راټول کړئ۔ د هغه هر غړے ټيټه رتبه لري، نو د
	\LR{$\Th{T}$} دننه لرو:
	\begin{equation}
		\Th{T} \Proves \exists M \psi(M). \tag{*}\label{Mpsi}
	\end{equation}
	د \LR{$\psi(X)$} د لومړي عطف له مخې، هرکله چې \LR{$\Th{T} \Proves \sigma$}
	وي، \LR{$\Th{T} \Proves \forall X(\psi(X) \lif \sigma^X)$} هم لرو۔ نو
	د \eqref{Mpsi} له مخې:
	\begin{align*}
		\Th{T} &\Proves \forall X(\psi(X) \lif (\exists N \psi(N))^X)\\
	\intertext{\RL{له دې او بيا له \eqref{Mpsi} څخه په ګټې اخيستو:}}
		\Th{T} &\Proves \exists M(\psi(M) \land (\exists N \psi(N))^M)
	\intertext{\RL{نو په ځانګړې توګه:}}
		\Th{T} &\Proves \exists M(\psi(M) \land (\exists N \in M)((N\text{\RL{ متعدي دے}})^M \land (\theta^N)^M))
	\intertext{\RL{نو د متعدي والي په اړه له بنسټيز استدلاله:}}
		\Th{T} &\Proves \exists M(\psi(M) \land (\exists N \in M)(N\text{\RL{ متعدي دے}} \land \theta^N))
	\end{align*} 
	دا د ټاکل شوي مدل له لږوالي سره ټکر لري، نو \LR{$\Th{T}$} ناسازګاره
	دے۔\footnote{د متعدي لوی مدل دننه د يوه غړي د متعدي والي
	کتنه د باندې هماغه نتيجه لري، ځکه د غړيتوب په دوو پرله‌پسې
	اړيکو کښې ښکېل ټول څيزونه په لوی مدل کښې دي۔ د کوچني مدل
	پر غړو محدود شوي کميت ټاکونکي بيا دننه او بهر عين ساحه لري۔
	نو د جملې پر جوړښت استقرا د هغې د نسبي رښتينولۍ مطلق‌والے ورکوي۔}
\end{proof}
\olpseganchor{OLP-0573-B005}{end}

\olpseganchor{OLP-0573-B006}{start}
ورته پايله \mbox{\LR{\latinfont \citet[223]{Potter2004}}} ياده کړې ده:
\olpseganchor{OLP-0573-B006}{end}

\olpseganchor{OLP-0573-B007}{start}
\begin{prop}\ollabel{finiteextensionofZ}
	فرض کړئ \LR{$\Th{T}$} د \LR{$\Z$} په نسبت د متناهيو نويو اصلونو پراختيا ده۔
	که \LR{$\Th{T} \Proves \ZF$} وي، نو \LR{$\Th{T}$} ناسازګاره دے۔
	(دلته د \olref{zfnotfinitely} هماغه ضمني محدوديتونه کاروو۔)
\end{prop}
\begin{proof}
	\LR{$\theta$} دې د \LR{$\Th{T}$} د ټولو اصلونو عطف وي، خو د بېلونې
	(نامتناهي ډېرو) بېلګو پرته۔ \LR{$\psi$} له \LR{$\theta$} څخه د
	\olref{zfnotfinitely} په شان تعريف کړئ، له دې زيات شرط سره
	چې مدل توانمن وي او د هغه په غړو کښې له هېڅ متعدي توانمن
	مدل سره د لږوالي شرط مات نۀ شي۔ توانمن سټ د هر غړي ټول فرعي
	سټونه هم د غړو په توګه لري۔ انعکاس يو پړاو راکوي چې د عطف
	مدل دے۔ هر پړاو متعدي او توانمن دے؛ نو د رتبې په لږوالي
	\LR{$\Th{T} \Proves \exists M \psi(M)$} ترلاسه کوو۔
\olpseganchor{OLP-0573-B007}{end}
	
\olpseganchor{OLP-0573-B008}{start}
	د \olref{zfnotfinitely} په شان، دا شېما ټينګولے شو: هرکله چې
	\LR{$\Th{T} \Proves \sigma$} وي، نو
	\LR{$\Th{T} \Proves \forall X(\psi(X) \lif \sigma^X)$}۔ بيا ثبوت عين
	د \olref{zfnotfinitely} په شان پاے ته رسوو۔
\olpseganchor{OLP-0573-B008}{end}
	
\olpseganchor{OLP-0573-B009}{start}
	خو د دې شېما ټينګول د \olref{zfnotfinitely} تر حالت لږ ډېر کار غواړي۔
	د بېلونې بېلګې په \LR{$\Th{T}$} کښې شته، خو د \LR{$\theta$} عطفونه نۀ دي۔
	د راکړل شوي غړي پر سر د بېلونې اصل له لوی جهان څخه
	پلي کړئ، او د بېلونې په شرط کښې ټول کميت ټاکونکي تر مدل
	محدود کړئ۔ پايله د توانمن‌والي له مخې د مدل غړې ده۔ د
	متعدي والي له مخې د راکړل شوي غړي ټول غړي هم هلته دي۔ نو
	\LR{$\Th{T} \Proves \forall X(X\text{\RL{ متعدي او توانمن دے}} \lif \sigma^X)$}
	د بېلونې د هرې بېلګې \LR{$\sigma$} دپاره۔ که کوچنے مدل په لوی
	توانمن مدل کښې غړے وي، د کوچني مدل د هر غړي ټول فرعي سټونه
	په لوی مدل کښې موندل کېږي۔ له دې امله د کوچني مدل دنننی
	توانمن‌والے د لوی مدل له باندې هم صدق کوي۔ دا د لږوالي
	د وروستي ټکر لپاره د متعدي والي او نسبي رښتينولۍ ترڅنګ
	اړين درېيم مطلق خاصيت دے۔
\end{proof}
\begin{prob}
	وښيئ چې د بېلونې د هرې بېلګې \LR{$\sigma$} دپاره دا رښتيا ده:
	\LR{$\Z \Proves \forall X(X\text{\RL{ متعدي او توانمن دے}} \lif \sigma^X)$}۔
	(دا شېما مو په \olref[sth][replacement][finiteaxiomatizability]{finiteextensionofZ}
	کښې کارولې وه۔)
\end{prob}
\begin{prob}
	وښيئ چې د هر \LR{$\phi \in \Z$} دپاره
	\LR{$\ZF \Proves \phi^{V_{\omega+\omega}}$} دے۔
\end{prob}
\begin{prob}
	په \olref[sth][replacement][finiteaxiomatizability]{zfnotfinitely}
	او \olref[sth][replacement][finiteaxiomatizability]{finiteextensionofZ}
	کښې په ثبوتونو کښې کارول شوې پاتې شېمايي پايلې تائيد کړئ۔
\end{prob}
\olpseganchor{OLP-0573-B009}{end}

\olpseganchor{OLP-0573-B010}{start}
لکه په \olref[sth][replacement][absinf]{sec} کښې چې وويل شول، دا ښيي
چې تعويض له \olref[ordinals][ordtype]{thmOrdinalRepresentation} څخه
په کلکه قوي دے۔ يا په لږ زيات دقت: که \LR{$\Z$} + «هر ښه ترتيب د يوه
يوازيني ترتيبي عدد سره آيزومورف دے» سازګار وي، نو لږ تر لږه د
تعويض يوه بېلګه نۀ شي ثابتولے۔
\olpseganchor{OLP-0573-B010}{end}


\olpseganchor{OLP-0573-B011}{start}
	%	By assumption, $\psi^{V_\alpha}$ has the form: 
	%	\[
	%		(\forall A \in V_\alpha)(\exists S \in V_\alpha)(\forall x \in V_\alpha)(x \in S \liff (\phi^{V_\alpha}(x) \land x \in A))
	%	\]
	%	To establish this holds, fix $A \in V_\alpha$. Using Separation, obtain: 
	%	\[
	%		S = \Setabs{x \in A}{\phi^{V_\alpha}(x)}
	%	\]
	%	Now $S \in V_\alpha$, since $S \subseteq A \in V_\alpha$, and clearly $(\forall x \in V_\alpha)(x \in S \liff (\phi^{V_\alpha}(x) \land x \in A)$.
\olpseganchor{OLP-0573-B011}{end}

\olpseganchor{OLP-0574-B004}{start}
\olchapter{sth}{ord-arithmetic}{د ترتيبي عددونو حساب}
\olpseganchor{OLP-0574-B004}{end}

\olpseganchor{OLP-0575-B002}{start}
\olfileid{sth}{ord-arithmetic}{intro}
\olpseganchor{OLP-0575-B002}{end}

\olpseganchor{OLP-0575-B003}{start}
\olsection{سريزه}
\olpseganchor{OLP-0575-B003}{end}

\olpseganchor{OLP-0575-B004}{start}
په \olref[ordinals][]{chap} کښې مو د ترتيبي عددونو تيوري جوړه کړه۔
په \olref[spine][]{chap} کښې مو وليدل چې ترتيبي عددونه د هغې
ملاتير په توګه ګڼلے شو چې د سلسلې پاتې برخه پرې رغېږي۔ خو د
ترتيبي عددونو يوازينۍ ونډه دا نه ده۔ د ترتيبي عددونو حساب هم بايد وکړو۔
\olpseganchor{OLP-0575-B004}{end}

\olpseganchor{OLP-0575-B005}{start}
مخکښې په \olref[ordinals][idea]{sec} کښې مو د \LR{$\omega$}،
\LR{$\omega+1$} او \LR{$\omega+\omega$} په يادولو دې خبرې ته اشاره کړې وه۔
هلته مو غير رسمي خبرې کړې وې؛ اوس ئې د کره بيانولو وخت دے۔ خو دا
هم ووايو چې په دې څپرکي کښې ډېر فلسفي بحث نشته؛ تر ډېره تخنيکي
پرمختګونه دي، او ورسره دا يو (لږ څه) په زړه پورې مشاهده چې يو کار
په دوو بېلو لارو کولے شو۔
\olpseganchor{OLP-0575-B005}{end}

\olpseganchor{OLP-0576-B003}{start}
\olfileid{sth}{ord-arithmetic}{add}
\olsection{د ترتيبي عددونو جمع}
\olpseganchor{OLP-0576-B003}{end}

\olpseganchor{OLP-0576-B004}{start}
فرض کړئ چې \LR{$\alpha$} او \LR{$\beta$} جمع کول غواړو۔ د \LR{$\beta$} يوه
کاپي د \LR{$\alpha$} د يوې کاپۍ سملاسي وروسته کېښودلے شو۔ (کاپيو ته اړ يو،
ځکه له \olref[ordinals][basic]{ordinalsaresubsets} څخه پوهېږو چې يا
\LR{$\alpha \subseteq \beta$} وي يا \LR{$\beta \subseteq \alpha$}۔) په بداهتي
ډول دا داسې ده لکه لومړے د \LR{$\alpha$} په اوږدوالي ګامونه واخلو،
\emph{او بيا} د \LR{$\beta$} په اوږدوالي ګامونه واخلو۔ پايلې ته رسېدلې
لړۍ به ښه ترتيب شوې وي؛ نو د
\olref[ordinals][ordtype]{thmOrdinalRepresentation} له مخې له يوه
(يوازيني) ترتيبي عدد سره هم‌شکله ده۔ هماغه ترتيبي عدد د \LR{$\alpha$}
او \LR{$\beta$} \emph{مجموعه} ګڼلے شو۔
\olpseganchor{OLP-0576-B004}{end}

\olpseganchor{OLP-0576-B005}{start}
دا د ترتيبي جمع تر شا بداهتي تصور دے۔ د کره تعريف دپاره لومړے
د سټونو د \emph{کاپيو} جوړولو تصور اخلو۔ دلته د \LR{$0$} او \LR{$1$} خپل‌سري
ټاګونه کاروو، څو څرګنده شي چې کوم څيز له کومه راغلے دے:
\olpseganchor{OLP-0576-B005}{end}

\olpseganchor{OLP-0576-B006}{start}
\begin{defn}\ollabel{defdissum}
د \LR{$A$} او \LR{$B$} \emph{ناګډه مجموعه} دا ده:
\LR{$A \disjointsum B = (A\times \{0\}) \cup (B \times \{1\})$}۔
\end{defn}
\olpseganchor{OLP-0576-B006}{end}

\olpseganchor{OLP-0576-B007}{start}
اوس د ترتيبي عددونو پر جوړو يو ترتيب تعريفوو:
\olpseganchor{OLP-0576-B007}{end}

\olpseganchor{OLP-0576-B008}{start}
\begin{defn}
د هر ډول ترتيبي عددونو \LR{$\alpha_1, \alpha_2, \beta_1, \beta_2$}
دپاره وايو چې:
\begin{align*}
	\tuple{\alpha_1, \alpha_2} \rlexless \tuple{\beta_1, \beta_2}\text{\RL{ هله او يوازې هله چې }}&
	\text{\RL{يا \LR{$\alpha_2 \in \beta_2$} وي}}\\
	& \text{\RL{يا هم \LR{$\alpha_2 = \beta_2$} او \LR{$\alpha_1 \in \beta_1$} دواړه وي}}
\end{align*}
\end{defn}
\olpseganchor{OLP-0576-B008}{end}

\olpseganchor{OLP-0576-B009}{start}
دا \emph{برعکس قاموسي} ترتيب دے، ځکه لومړے د جوړې دويم غړے او
بيا لومړے غړے پرتله کېږي۔ اوس په ياد ولرئ چې \LR{$\alpha \ordplus \beta$}
مو د \LR{$\alpha$} د کاپۍ ورپسې د \LR{$\beta$} د کاپۍ د ترتيب ډول غوښت۔ د دې
دپاره وايو:
\olpseganchor{OLP-0576-B009}{end}

\olpseganchor{OLP-0576-B010}{start}
\begin{defn}\ollabel{defordplus}
د هر ډول ترتيبي عددونو \LR{$\alpha$}، \LR{$\beta$} دپاره د دوی مجموعه دا ده:
\LR{$\alpha \ordplus \beta = \ordtype{\alpha \disjointsum \beta, \rlexless}$}۔
\end{defn}
\noindent
پام وکړئ چې دلته مو د نښو په کارونه کښې لږ تساهل کړے دے؛ په کره
ډول بايد «\LR{$\Setabs{\tuple{x,y}}{x,y\in\alpha\disjointsum\beta\land x \rlexless y}$}»
د «\LR{$\rlexless$}» پر ځاے وليکو۔ خو د لنډيز دپاره به لاندې همدا
تساهل روان ساتو۔
\olpseganchor{OLP-0576-B010}{end}

\olpseganchor{OLP-0576-B011}{start}
لاندې پايله، له \olref[ordinals][ordtype]{thmOrdinalRepresentation}
سره يوځاے، ښيي چې زموږ تعريف په سمه توګه جوړ شوے دے:
\olpseganchor{OLP-0576-B011}{end}

\olpseganchor{OLP-0576-B012}{start}
\begin{lem}\ollabel{ordsumlessiswo}
\LR{$\tuple{\alpha \disjointsum \beta, \rlexless}$} ښه ترتيب دے،
د هر ډول ترتيبي عددونو \LR{$\alpha$} او \LR{$\beta$} دپاره۔
\end{lem}
\olpseganchor{OLP-0576-B012}{end}

\olpseganchor{OLP-0576-B013}{start}
\begin{proof}
\LR{$\rlexless$} پر \LR{$\alpha \disjointsum \beta$} يو کلک کلي
ترتيب دے۔ د جوړې دويمې برخې لومړے پرتله کېږي؛ که برابرې
وي، لومړۍ برخې پرېکړه کوي۔ د دواړو مختصونو کلک ترتيبي
خاصيتونه د ځان پر وړاندې نۀ صدق، انتقالي والى او د بېلو
جوړو پرتله تضمينوي۔ د دې د ښه‌بنسټه والي د ښودلو دپاره ناتش
\LR{$X \subseteq \alpha \disjointsum \beta$} وټاکئ۔ \LR{$Y$} دې د \LR{$X$}
هغه فرعي سټ وي چې د غړو دويم مختص ئې تر ټولو کوچنے وي، يعنې
\LR{$Y = \Setabs{\tuple{\gamma, i} \in X}{(\forall \tuple{\delta, j} \in X)i \leq j}$}۔
اوس د \LR{$Y$} هغه غړے وټاکئ چې لومړے مختص ئې تر ټولو کوچنے وي۔
هغه په ټول فرعي سټ کښې هم تر ټولو کوچنے دے، ځکه هر بل غړے يا
لوے دويم مختص لري، يا په همدې مختص کښې لوے يا برابر لومړے مختص لري۔
\end{proof}
\noindent
په دې توګه د ترتيبي جمع يو ښه ښکاره تعريف لرو۔ يوه د تعجب وړ
نه ده پايله دا ده (\LR{$1 = \{0\}$} په ياد ولرئ، د
\olref[z][infinity-again]{defnomega} له مخې):
\begin{prop}
\LR{$\alpha \ordplus  1 = \ordsucc{\alpha}$}، د هر ترتيبي عدد
\LR{$\alpha$} دپاره۔
\end{prop}
\olpseganchor{OLP-0576-B013}{end}

\olpseganchor{OLP-0576-B014}{start}
\begin{proof}
د \LR{$f$} هم‌شکلتيا له \LR{$\ordsucc{\alpha} = \alpha \cup \{\alpha\}$} څخه تر
\LR{$\alpha\disjointsum1 = (\alpha \times \{0\})
\cup (\{0\} \times \{1\})$} پورې وګورئ؛ دا
\LR{$f(\gamma) = \tuple{\gamma, 0}$} د \LR{$\gamma \in \alpha$} دپاره،
او \LR{$f(\alpha) = \tuple{0, 1}$} ټاکي۔
په لومړۍ برخه کښې ترتيب د لومړي مختص له مخې ساتل کېږي۔
وروستے غړے دويمې برخې ته ځي، چې يوازې همدا يو غړے لري۔
نو دا تابعه يو پر يو او پوره ده او ترتيب په دواړو لورو ساتي۔
\end{proof}
\noindent
سربېره پر دې، په اسانه ښودل کېږي چې جمع ځينې بازګشتي شرطونه
پوره کوي:
\olpseganchor{OLP-0576-B014}{end}

\olpseganchor{OLP-0576-B015}{start}
\begin{lem}\ollabel{ordadditionrecursion}
د هر ډول ترتيبي عددونو \LR{$\alpha, \beta$} دپاره:
\begin{align*}
	\alpha\ordplus 0 &= \alpha\\
	\alpha \ordplus  (\beta\ordplus 1) &= (\alpha \ordplus  \beta) \ordplus  1\\
	\alpha  \ordplus  \beta &= \supstrict_{\delta < \beta}(\alpha \ordplus  \delta) && \text{\RL{که \LR{$\beta $} حدي ترتيبي عدد وي}}
\end{align*}
\end{lem}
\olpseganchor{OLP-0576-B015}{end}

\olpseganchor{OLP-0576-B016}{start}
\begin{proof}
هر حالت جلا ګورو؛ لومړے:
\begin{align*}
	\alpha \ordplus  0
	& = \ordtype{(\alpha \times \{0\}) \cup (0 \times \{1\}), \rlexless} \\
	&= \ordtype{(\alpha \times \{0\}) \cup \emptyset, \rlexless}\\
	&= \alpha\\
	\alpha \ordplus (\beta \ordplus  1)
	%&= \ordtype{(\alpha\times \{0\}) \cup (\ordtype{\beta \ordplus  1}\times \{1\}), \rlexless} \\
	&= \ordtype{(\alpha\times \{0\}) \cup (\ordsucc{\beta}\times \{1\}), \rlexless} \\
	&= \ordtype{(\alpha\times \{0\}) \cup (\beta \times \{1\}), \rlexless} \ordplus 1\\
	&= (\alpha \ordplus  \beta) \ordplus  1
\end{align*}
له تشې برخې سره اتحاد لومړۍ برخه پر خپل حال پرېږدي۔
د تالي په حالت کښې دويمه برخه يوازې يو وروستے غړے زياتوي۔
اوس دې \LR{$\beta \neq \emptyset$} حدي وي۔ که \LR{$\delta < \beta$} وي،
نو \LR{$\delta\ordplus 1 < \beta$} هم دے؛ له دې امله
\LR{$\alpha \ordplus  \delta$} د \LR{$\alpha \ordplus  \beta$} يوه
مناسبه پيلنۍ برخه ده۔ نو \LR{$\alpha \ordplus  \beta$} د
\LR{$X = \Setabs{\alpha \ordplus  \delta}{\delta < \beta}$} کلکه
پاسنۍ کرانه ده۔ سربېره پر دې، که
\LR{$\alpha \leq \gamma < \alpha \ordplus  \beta$} وي، نو په ښکاره
\LR{$\gamma = \alpha \ordplus \delta$} د کوم \LR{$\delta < \beta$}
دپاره دے؛ ځکه له لومړۍ برخې وروسته هر مناسب پيلنے پرېکون
د دويمې برخې يو پيلنے پرېکون اخلي۔ تر لومړۍ برخې کوچني
شاخصونه د کورنۍ د لومړي غړي لاندې دي؛ نور کوچني شاخصونه
پخپله د کورنۍ غړي دي۔ له دغو دواړو ډلو هېڅ يو شاخص کلک
پاسنے حد نۀ شي کېدے۔ نو
\LR{$\alpha \ordplus  \beta = \supstrict_{\delta< \beta}(\alpha\ordplus \delta)$}۔
\end{proof}
\olpseganchor{OLP-0576-B016}{end}

\olpseganchor{OLP-0576-B017}{start}
خو دلته يوه په زړه پورې خبره ده۔ د ترتيبي جمع د تعريف دپاره مو
\emph{پر ځاے} يوازې د نامتناهيت هاخوا بازګښت قضيه کارولے شوه،
او د \olref{ordadditionrecursion} هماغه بازګشتي معادلې مو ايښودلے
شوې (البته «\LR{$\ordsucc{\beta}$}» په کارولو، د
«\LR{$\beta \ordplus 1$}» پر ځاے)۔
\olpseganchor{OLP-0576-B017}{end}

\olpseganchor{OLP-0576-B018}{start}
نو د ترتيبي عددونو عمليې په دوو بېلو لارو تعريفولے شو۔ په
\emph{ترکيبي} لاره يو ښه ترتيب شوے سټ په ښکاره جوړوو او د هغه
د ترتيب ډول اخلو۔ يا ئې \emph{بازګشتي} لاره تعريفولے شو، يوازې
د بازګښت معادلې په ايښودلو۔ دواړه طريقې يوه پايله ورکوي۔
\olref[ordinals][ordtype]{thmOrdinalRepresentation} په ترکيبي
طريقه کښې د ترتيب ډول يو ترتيبي عدد ټاکي۔ د بازګښت د قضيې
\emph{يوازينيت} بيا وايي چې د دغو معادلو هره پوره کوونکې
تابعه هماغه ترتيبي جمع ورکوي۔
\olpseganchor{OLP-0576-B018}{end}

\olpseganchor{OLP-0576-B019}{start}
په ډېرو اړخونو کښې د ترتيبي عددونو حساب د طبيعي عددونو د جمع
په شان دے۔ د بېلګې په توګه، لاندې ثابتولے شو:
\olpseganchor{OLP-0576-B019}{end}

\olpseganchor{OLP-0576-B020}{start}
\begin{lem}\ollabel{ordinaladditionisnice}
که \LR{$\alpha, \beta, \gamma$} ترتيبي عددونه وي، نو:
\begin{enumerate}
	\item\ollabel{ordaddition1} که \LR{$\beta < \gamma$} وي، نو \LR{$\alpha
	\ordplus  \beta < \alpha \ordplus  \gamma$}
	\item\ollabel{ordaddition2} که \LR{$\alpha \ordplus  \beta =
	\alpha\ordplus \gamma$} وي، نو \LR{$\beta = \gamma$}
	\item\ollabel{ordaddition3} \LR{$\alpha \ordplus  (\beta \ordplus
	\gamma) = (\alpha \ordplus  \beta) \ordplus  \gamma$}، يعنې
	جمع اتحادي ده
	\item\ollabel{ordaddition4} که \LR{$\alpha \leq \beta$} وي، نو \LR{$\alpha
	\ordplus  \gamma \leq \beta \ordplus \gamma$}
\end{enumerate}
\end{lem}
\olpseganchor{OLP-0576-B020}{end}

\olpseganchor{OLP-0576-B021}{start}
\begin{proof}
\olref{ordaddition3} ثابتوو او پاتې برخې د تمرين دپاره پرېږدو۔
ثبوت د \LR{$\gamma$} په اړه د ساده نامتناهيت هاخوا استقرا له لارې،
د \olref{ordadditionrecursion} په کارولو کېږي۔ که \LR{$\gamma = 0$} وي:
\[
(\alpha \ordplus  \beta) \ordplus  0 = \alpha \ordplus  \beta  = \alpha \ordplus  (\beta \ordplus  0)
\]
که \LR{$\gamma = \delta\ordplus 1$} وي، د استقرا له مخې فرض کړئ چې
\LR{$(\alpha \ordplus  \beta) \ordplus  \delta = \alpha \ordplus  (\beta \ordplus
\delta)$}؛ اوس \olref{ordadditionrecursion} درې وارې کاروو:
\begin{align*}
	(\alpha \ordplus  \beta) \ordplus  (\delta \ordplus  1) & = ((\alpha \ordplus  \beta) \ordplus  \delta)\ordplus 1\\
	& = (\alpha \ordplus  (\beta \ordplus  \delta)) \ordplus  1\\
	& = \alpha \ordplus  ((\beta \ordplus  \delta)\ordplus 1)\\
	& = \alpha \ordplus  (\beta \ordplus  (\delta\ordplus 1))
\end{align*}
که \LR{$\gamma$} حدي ترتيبي عدد وي، د استقرا له مخې فرض کړئ چې که
\LR{$\delta \in \gamma$} وي، نو
\LR{$(\alpha \ordplus  \beta) \ordplus  \delta =
\alpha \ordplus  (\beta \ordplus  \delta)$}۔ د لاندې محاسبې
درېيمه برابري د جمع له حدي معادلې راځي۔ د ښي اړخ د جمع دويم
حد پخپله يو حدي عدد دے، ځکه د مخکنيو پړاوونو جمعې په کلکه
زياتېږي او تر هغه حده هم‌پايه دي۔ نو د ټولو کوچنيو شاخصونو
پر ځاے همدغه هم‌پايه شاخصونه د کلک پاسني حد لپاره بس دي:
\begin{align*}
	(\alpha \ordplus  \beta) \ordplus  \gamma & = \supstrict_{\delta < \gamma}((\alpha \ordplus  \beta) \ordplus  \delta) \\
	&= \supstrict_{\delta < \gamma}(\alpha \ordplus  (\beta \ordplus  \delta))\\
	&= \alpha \ordplus  \supstrict_{\delta < \gamma}(\beta \ordplus  \delta)\\
	& = \alpha \ordplus  (\beta \ordplus  \gamma)
\end{align*}
\end{proof}
\olpseganchor{OLP-0576-B021}{end}

\olpseganchor{OLP-0576-B022}{start}
\begin{prob}
د \olref[sth][ord-arithmetic][add]{ordinaladditionisnice} پاتې
برخې ثابتې کړئ۔
\end{prob}
\olpseganchor{OLP-0576-B022}{end}

\olpseganchor{OLP-0576-B023}{start}
تر دې ځايه ترتيبي جمع بايد ډېره اشنا ښکاره شي۔ خو يوه مهمه
لاره شته چې ترتيبي جمع پکې د طبيعي عددونو له جمع سره \emph{نه}
ورته کېږي۔
\olpseganchor{OLP-0576-B023}{end}

\olpseganchor{OLP-0576-B024}{start}
\begin{prop}\ollabel{ordsumnotcommute}
ترتيبي جمع تبادلي \emph{نه} ده؛ \LR{$1 \ordplus  \omega = \omega <
\omega \ordplus  1$}۔
\end{prop}
\olpseganchor{OLP-0576-B024}{end}

\olpseganchor{OLP-0576-B025}{start}
\begin{proof}
پام وکړئ چې \LR{$1 \ordplus  \omega = \supstrict_{n < \omega} (1 \ordplus n)
= \omega \in \omega \cup \{\omega\} = \ordsucc{\omega} = \omega
\ordplus  1$}۔
\end{proof}
\noindent
که څه هم دا په لومړي نظر حيرانولے شي، \emph{بايد مو حيران نه کړي}۔
له يوې خوا، په \LR{$1 \ordplus  \omega$} کښې له ټولو طبيعي عددونو
\emph{مخکښې} يو اضافي غړے ږدئ۔ لکه د هلبرټ د هوټل په استدلال کښې
چې مو په \olref[sfr][infinite][hilbert]{sec} کښې وکړل، په بداهتي
ډول دا اضافي لومړنے غړے بايد د ټول ترتيب په ډول کښې توپير رانه ولي۔
له بلې خوا، په \LR{$\omega \ordplus  1$} کښې له ټولو طبيعي عددونو
\emph{وروسته} يو اضافي غړے ږدئ۔ دا بيا بيخي بل ډول څيز دے!
\olpseganchor{OLP-0576-B025}{end}

\olpseganchor{OLP-0577-B003}{start}
\olfileid{sth}{ord-arithmetic}{using-addition}
\olsection{د ترتيبي جمع کارول}
\olpseganchor{OLP-0577-B003}{end}

\olpseganchor{OLP-0577-B004}{start}
د ترتيبي عددونو په جمع سره د بېلابېلو سټونو رتبې په ښکاره حسابولے شو،
د \olref[spine][rank]{defnsetrank} په معنا:
\olpseganchor{OLP-0577-B004}{end}

\olpseganchor{OLP-0577-B005}{start}
\begin{lem}\ollabel{rankcomputation}
که \LR{$\setrank{A} = \alpha$} او \LR{$\setrank{B} = \beta$} وي، نو:
\begin{enumerate}
	\item\ollabel{exrankpow} \LR{$\setrank{\Pow{A}} = \alpha\ordplus 1$}
	\item\ollabel{exrankpair} \LR{$\setrank{\{A, B\}} = \max(\alpha,
	\beta) \ordplus 1$}
	\item\ollabel{exrankcup} \LR{$\setrank{A \cup B} = \max(\alpha,
	\beta)$}
	\item\ollabel{exranktuple} \LR{$\setrank{\tuple{A,B}} = \max(\alpha,
	\beta) \ordplus  2$}
	\item\ollabel{exranktimes} \LR{$\setrank{A \times B} \leq \max(\alpha,
	\beta) \ordplus  2$}
	\item\ollabel{exrankunion} \LR{$\setrank{\bigcup A} = \alpha$} چې
	\LR{$\alpha$} تش يا حدي وي؛ \LR{$\setrank{\bigcup A} = \gamma$} چې
	\LR{$\alpha = \gamma\ordplus 1$} وي
\end{enumerate}
\end{lem}
\olpseganchor{OLP-0577-B005}{end}

\olpseganchor{OLP-0577-B006}{start}
\begin{proof}
په ټول ثبوت کښې \olref[spine][rank]{ranksupstrict} بيا بيا کاروو۔
\olpseganchor{OLP-0577-B006}{end}

\olpseganchor{OLP-0577-B007}{start}
\emph{\olref{exrankpow}.} که \LR{$x \subseteq A$} وي، نو \LR{$\setrank{x} \leq
\setrank{A}$}۔ نو \LR{$\setrank{\Pow{A}} \leq \alpha \ordplus  1$}۔
څرنګه چې په ځانګړې توګه \LR{$A \in \Pow{A}$}، نو
\LR{$\setrank{\Pow{A}} = \alpha \ordplus  1$}۔
\olpseganchor{OLP-0577-B007}{end}

\olpseganchor{OLP-0577-B008}{start}
\emph{\olref{exrankpair}.} د \olref[spine][rank]{ranksupstrict} له مخې۔
\olpseganchor{OLP-0577-B008}{end}

\olpseganchor{OLP-0577-B009}{start}
\emph{\olref{exrankcup}.} د \olref[spine][rank]{ranksupstrict} له مخې۔
\olpseganchor{OLP-0577-B009}{end}

\olpseganchor{OLP-0577-B010}{start}
\emph{\olref{exranktuple}.} \olref{exrankpair} دوه ځله وکاروئ۔
\olpseganchor{OLP-0577-B010}{end}

\olpseganchor{OLP-0577-B011}{start}
\emph{\olref{exranktimes}.} پام وکړئ چې
\LR{$A \times B \subseteq \Pow{\Pow{A \cup B}}$}، او \olref{exranktuple}
وکاروئ۔ د هرې مرتبې جوړې د غړو رتبې د دواړو سټونو د رتبو
له لوی حده ټيټې دي۔ د جوړې د رتبې له پاسه يو تالي واخلو؛
د دغو تاليانو پاسنے حد د حاصل ضرب رتبه ده۔ که لوی حد تالي
وي او دواړه سټونه ناتش وي، يو عامل تر هغه سمدستي مخکنۍ
رتبه لرونکے غړے لري؛ بل عامل د جوړې لپاره يو غړے ورکوي۔
که لوی حد حدي وي، له هرې مخکنۍ رتبې څخه هاخوا د يوه عامل
غړے موندل کېږي، او متناهي تاليان هم تر حدي رتبې لاندې پاتې
کېږي۔ که يو عامل تش وي، حاصل ضرب تش دے۔
\olpseganchor{OLP-0577-B011}{end}

\olpseganchor{OLP-0577-B012}{start}
\emph{\olref{exrankunion}.} که \LR{$\alpha = \gamma\ordplus 1$} وي، نو
داسې \LR{$c \in A$} شته چې \LR{$\setrank{c} = \gamma$} وي، او د \LR{$A$}
هېڅ غړے تر دې لوړه رتبه نه لري؛ نو
\LR{$\setrank{\bigcup A} = \gamma$}۔ که \LR{$\alpha$} حدي ترتيبي عدد وي،
نو \LR{$A$} داسې غړي لري چې رتبې يې \LR{$\alpha$} ته هر څومره نږدې،
خو په کلکه ترې ټيټې، رسېږي؛ نو \LR{$\bigcup A$} هم داسې غړي
لري چې رتبې يې \LR{$\alpha$} ته هر څومره نږدې، خو په کلکه ترې ټيټې،
رسېږي؛ له دې امله \LR{$\setrank{\bigcup A} = \alpha$}۔ د وروستي
استلزام لپاره له هر ټاکلي شاخص څخه دوه تاليان پورته شئ۔
ځکه چې حد حدي دے، دا تاليان لا تر حده ټيټ دي؛ د يوه
غړي دومره لوړه رتبه د هغه دننه هم تر ټاکلي شاخص د لوړې
رتبې غړے راکوي۔
\end{proof}
\noindent
دا د تمرين دپاره پرېږدو چې ولې \olref{exranktimes} يو
\emph{نا}مساوات دے۔
\olpseganchor{OLP-0577-B012}{end}

\olpseganchor{OLP-0577-B013}{start}
\begin{prob}
داسې سټونه \LR{$A$} او \LR{$B$} جوړ کړئ چې \LR{$\setrank{A \times B}=
\max(\setrank{A}, \setrank{B})$} وي۔ داسې سټونه \LR{$A$} او \LR{$B$}
جوړ کړئ چې
\LR{$\setrank{A \times B}=\max(\setrank{A}, \setrank{B}) \ordplus  2$}
وي۔ ايا نورې رتبې هم ممکنې دي؟
\end{prob}
\olpseganchor{OLP-0577-B013}{end}

\olpseganchor{OLP-0577-B014}{start}
اوس دا هم ښودلے شو چې د يو ترتيبي عدد د «متناهي» يا «نامتناهي»
بللو څو معقولې معناوې سره برابرې دي:
\olpseganchor{OLP-0577-B014}{end}

\olpseganchor{OLP-0577-B015}{start}
\begin{lem}\ollabel{ordinfinitycharacter}
د هر ترتيبي عدد \LR{$\alpha$} دپاره لاندې شرطونه همارزښته دي:
\begin{enumerate}
	\item\ollabel{ord:notinomega} \LR{$\alpha\notin \omega$}، يعنې
	\LR{$\alpha$} طبيعي عدد نه دے
	\item\ollabel{ord:omegaplus} \LR{$\omega \leq \alpha$}
	\item\ollabel{ord:oneplus} \LR{$1 \ordplus  \alpha = \alpha$}
	%\alpha \approx \alpha \ordplus 1$
	\item\ollabel{ord:plusone} \LR{$\alpha \approx \alpha\ordplus 1$}،
	يعنې \LR{$\alpha$} او \LR{$\alpha\ordplus 1$} هم‌شمېره دي
	\item\ollabel{ord:infinite} \LR{$\alpha$} د ډېډېکېنډ په معنا نامتناهي دے
	\end{enumerate}
\end{lem}
\noindent
نو د ترتيبي عدد د (نا)متناهي والي د پوهېدو پنځه په ثبوت برابرې
لارې لرو۔
\olpseganchor{OLP-0577-B015}{end}

\olpseganchor{OLP-0577-B016}{start}
\begin{proof}
\emph{\olref{ord:notinomega} \LR{$\Rightarrow$} \olref{ord:omegaplus}.}
د درې‌ګوني وېش له مخې۔
\olpseganchor{OLP-0577-B016}{end}

\olpseganchor{OLP-0577-B017}{start}
\emph{\olref{ord:omegaplus} \LR{$\Rightarrow$} \olref{ord:oneplus}.}
\LR{$\alpha \geq \omega$} وټاکئ۔ د نامتناهيت هاخوا استقرا له مخې داسې
تر ټولو کوچنے ترتيبي عدد \LR{$\gamma$} (ښايي \LR{$0$}) شته چې يو حدي
ترتيبي عدد \LR{$\beta$} د \LR{$\alpha = \beta \ordplus \gamma$} شرط پوره
کوي۔ د دې وېش لپاره له ټاکلي عدد څخه د هغه سمدستي مخکنيو
عددونو ته ستنېږو، تر څو حدي عدد ته ورسېږو۔ دا تګ بې‌پايه
نۀ شي کېدے، ځکه ترتيبي عددونه ښه مرتب دي؛ او صفر ته هم
نۀ رسېږي، ځکه ټاکلے عدد نامتناهي دے۔ د حدي عدد په لومړنۍ
نامتناهي برخه کښې ورزيات شوے لومړے ټکی صفر ته لېږو او هر
طبيعي عدد خپل تالي ته لېږو۔ له طبيعي عددونو وروسته هر ټکی
پر خپل ځاے پاتې کېږي۔ نو د حدي عدد د ترتيب ډول د لومړي
ټکي په زياتولو سره نۀ بدلېږي۔ اوس:
\[
	1 \ordplus \alpha =
	1 \ordplus (\beta \ordplus \gamma) =
	(1 \ordplus \beta) \ordplus \gamma =
	\supstrict_{\delta < \beta} (1 \ordplus  \delta) \ordplus  \gamma =
	\beta \ordplus  \gamma =
	\alpha.
\]
\emph{\olref{ord:oneplus} \LR{$\Rightarrow$} \olref{ord:plusone}.}
په څرګنده يوه دوه اړخيزه يو پر يو تابع
\LR{$f \colon (\alpha \disjointsum 1) \to (1 \disjointsum \alpha)$}
شته۔ که \LR{$1 \ordplus \alpha = \alpha$} وي، نو يوه هم‌شکلتيا
\LR{$g \colon (1 \disjointsum \alpha) \to \alpha$} شته۔ اوس
\LR{$\comp{f}{g}$} وګورئ۔
\olpseganchor{OLP-0577-B017}{end}

\olpseganchor{OLP-0577-B018}{start}
\emph{\olref{ord:plusone} \LR{$\Rightarrow$} \olref{ord:infinite}.}
که \LR{$\alpha \approx \alpha \ordplus 1$} وي، يوه دوه اړخيزه يو پر يو تابع
\LR{$f \colon (\alpha \disjointsum 1) \to \alpha$} شته۔
\LR{$g(\gamma) = f(\gamma, 0)$} د هر \LR{$\gamma < \alpha$} دپاره تعريف کړئ؛
دا يو پر يو تابع ښيي چې \LR{$\alpha$} د ډېډېکېنډ په معنا نامتناهي دے،
ځکه \LR{$f(0,1) \in \alpha \setminus \ran{g}$}۔
\olpseganchor{OLP-0577-B018}{end}

\olpseganchor{OLP-0577-B019}{start}
\emph{\olref{ord:infinite} \LR{$\Rightarrow$} \olref{ord:notinomega}.}
دا \olref[z][infinity-again]{naturalnumbersarentinfinite} دے۔
\end{proof}
\olpseganchor{OLP-0577-B019}{end}

\olpseganchor{OLP-0578-B003}{start}
\olfileid{sth}{ord-arithmetic}{mult}
\olpseganchor{OLP-0578-B003}{end}

\olpseganchor{OLP-0578-B004}{start}
\olsection{د ترتيبي عددونو ضرب}
\olpseganchor{OLP-0578-B004}{end}

\olpseganchor{OLP-0578-B005}{start}
اوس د ترتيبي عددونو ضرب ته راځو، او د جمع په شان يې څېړو۔ فرض
کړئ چې \LR{$\alpha$} له~\LR{$\beta$} سره ضربول غواړو۔ د دې د تصور دپاره
يو مستطيل جال په پام کښې ونيسئ چې سور يې \LR{$\alpha$} او لوړوالے يې
\LR{$\beta$} وي؛ د \LR{$\alpha$} او \LR{$\beta$} حاصل ضرب دا دے چې د هر قطار
په اوږدو کښې لاړ شئ، بيا بل قطار ته ورشئ\ldots تر څو ټول جال
تېر کړئ۔ په بله وينا، د \LR{$\alpha$} او \LR{$\beta$} حاصل ضرب داسې جوړېږي
چې د \LR{$\beta$} \emph{هر} غړے د \LR{$\alpha$} په يوه کاپۍ بدل کړئ۔
\olpseganchor{OLP-0578-B005}{end}

\olpseganchor{OLP-0578-B006}{start}
د رسمي کولو دپاره د \LR{$\alpha$} او \LR{$\beta$} پر کارتيزي حاصل ضرب
هماغه برعکس قاموسي ترتيب کاروو:
\olpseganchor{OLP-0578-B006}{end}

\olpseganchor{OLP-0578-B007}{start}
\begin{defn}
د هر ډول ترتيبي عددونو \LR{$\alpha, \beta$} دپاره د دوی حاصل ضرب
\LR{$\alpha \ordtimes \beta = \ordtype{\alpha \times \beta, \rlexless}$} دے۔
\end{defn}
\noindent
بيا بايد پخلی وکړو چې دا تعريف په سمه توګه جوړ شوے دے:
\olpseganchor{OLP-0578-B007}{end}

\olpseganchor{OLP-0578-B008}{start}
\begin{lem}\ollabel{ordtimeslessiswo}
\LR{$\tuple{\alpha \times \beta, \rlexless}$} د هر ډول ترتيبي عددونو
\LR{$\alpha$} او \LR{$\beta$} دپاره ښه ترتيب دے۔
\end{lem}
\olpseganchor{OLP-0578-B008}{end}

\olpseganchor{OLP-0578-B009}{start}
\begin{proof}
عين د \olref[add]{ordsumlessiswo} په شان۔
\end{proof}
\noindent
دا هم ثابتول ګران نه دي چې ضرب لاندې چلند کوي:
\olpseganchor{OLP-0578-B009}{end}

\olpseganchor{OLP-0578-B010}{start}
\begin{lem}\ollabel{ordtimesrecursion}
د هر ډول ترتيبي عددونو \LR{$\alpha, \beta$} دپاره:
\begin{align*}
	\alpha \ordtimes 0 &= 0\\
	\alpha \ordtimes (\beta \ordplus 1) &=
		(\alpha \ordtimes \beta) \ordplus \alpha\\
	\alpha  \ordtimes \beta &=
		\bigcup_{\delta < \beta}(\alpha \ordtimes \delta) &&
		\text{\RL{چې \LR{$\beta$} حدي ترتيبي عدد وي}}.
\end{align*}
\end{lem}
\olpseganchor{OLP-0578-B010}{end}

\olpseganchor{OLP-0578-B011}{start}
\begin{proof}
د تمرين لپاره اشاره: د صفر ښي عامل جال تشوي، او تالي
يو بشپړ وروستے قطار زياتوي۔ په حدي حالت کښې د ټولو
مخکنيو پړاوونو جالونه پيلنۍ برخې دي او په اتحاد کښې
ټول جال جوړوي۔ د ترتيب ډولونه يې هم په هماغه ترتيب
يو د بل پيلنۍ برخې دي؛ نو د ډولونو اتحاد د ټول جال
ډول ورکوي۔ دا استدلال د صفر کيڼ عامل لپاره هم کار کوي۔
\end{proof}
\olpseganchor{OLP-0578-B011}{end}

\olpseganchor{OLP-0578-B012}{start}
په رښتيا، د جمع په شان مو ترتيبي ضرب هم د نېغ تعريف پر ځاے د
دغو بازګشتي معادلو له لارې تعريفولے شو۔ د جمع په څېر دلته هم
ځينې چلندونه اشنا دي:
\olpseganchor{OLP-0578-B012}{end}

\olpseganchor{OLP-0578-B013}{start}
\begin{lem}\ollabel{ordinalmultiplicationisnice}
که \LR{$\alpha, \beta, \gamma$} ترتيبي عددونه وي، نو:
\begin{enumerate}
	\item\ollabel{ordtimes1} که \LR{$\alpha \neq 0$} او \LR{$\beta < \gamma$}
	وي، نو \LR{$\alpha \ordtimes \beta < \alpha \ordtimes \gamma$}؛
	\item\ollabel{ordtimes2} که \LR{$\alpha \neq 0$} او \LR{$\alpha \ordtimes
	\beta = \alpha\ordtimes\gamma$} وي، نو \LR{$\beta = \gamma$}؛
	\item\ollabel{ordtimes3} \LR{$\alpha \ordtimes (\beta \ordtimes
	\gamma) = (\alpha \ordtimes \beta) \ordtimes \gamma$}؛
	\item\ollabel{ordtimes4} که \LR{$\alpha \leq \beta$} وي، نو \LR{$\alpha
	\ordtimes \gamma \leq \beta \ordtimes \gamma$}؛
	\item\ollabel{ordtimes5} \LR{$\alpha \ordtimes (\beta \ordplus
	\gamma) = (\alpha \ordtimes \beta )\ordplus (\alpha\ordtimes
	\gamma)$}۔
\end{enumerate}
\end{lem}
\olpseganchor{OLP-0578-B013}{end}

\olpseganchor{OLP-0578-B014}{start}
\begin{proof}
دا هم د تمرين دپاره پرېږدو۔
\end{proof}
\olpseganchor{OLP-0578-B014}{end}

\olpseganchor{OLP-0578-B015}{start}
نورې پايلې هم د خپلې خوښې تر بريده ثابتولے يا لوستلے شئ۔ خو د
\olref[ord-arithmetic][add]{ordsumnotcommute} له مخې لاندې خبره
بايد حيرانوونکې نه وي:
\olpseganchor{OLP-0578-B015}{end}

\olpseganchor{OLP-0578-B016}{start}
\begin{prop}
ترتيبي ضرب تبادلي نه دے: \LR{$2 \ordtimes \omega  =
\omega < \omega \ordtimes 2$}۔
\end{prop}
\olpseganchor{OLP-0578-B016}{end}

\olpseganchor{OLP-0578-B017}{start}
\begin{proof}
\LR{$2 \ordtimes \omega = \supstrict_{n < \omega}(2\ordtimes  n) = \omega \in \supstrict_{n < \omega}(\omega \ordplus n) = \omega \ordplus \omega = \omega \ordtimes 2$}۔
\end{proof}
\noindent
دلته هم بداهتي دليل ډېر روښانه دے۔ د \LR{$2
\ordtimes \omega$} د حساب دپاره هر طبيعي عدد په دوو څيزونو بدل کړئ۔
هماغه د ترتيب ډول ترلاسه کوئ که په طبيعي عددونو کښې د ټولو
«نيمو» عددونو ځايونه ورزيات کړئ؛ په هره جوړه کښې لومړے
بشپړ عدد، بيا تر هغه نيم زيات عدد راځي۔ يعنې پر
\LR{$\Setabs{\nicefrac{n}{2}}{n \in \omega}$} طبيعي ترتيب وګورئ۔
په دې بڼه ئې د ترتيب ډول په څرګنده پخپله د \LR{$\omega$} په شان دے۔
خو د \LR{$\omega \ordtimes 2$} د حساب دپاره د \LR{$\omega$} دوې کاپۍ
يوه د بلې ورپسې ږدئ۔
\olpseganchor{OLP-0578-B017}{end}

\olpseganchor{OLP-0578-B018}{start}
\begin{prob}
\olref[sth][ord-arithmetic][mult]{ordtimeslessiswo}،
\olref[sth][ord-arithmetic][mult]{ordtimesrecursion}،
او \olref[sth][ord-arithmetic][mult]{ordinalmultiplicationisnice}
ثابت کړئ۔
\end{prob}
\olpseganchor{OLP-0578-B018}{end}

\olpseganchor{OLP-0579-B003}{start}
\olfileid{sth}{ord-arithmetic}{expo}
\olsection{د ترتيبي عددونو توان}
\olpseganchor{OLP-0579-B003}{end}

\olpseganchor{OLP-0579-B004}{start}
اوس د ترتيبي عددونو توان ته راځو۔ له بده مرغه، د ترتيبي عددونو د
توان دپاره کوم \emph{ښکلے} ترکيبي تعريف نشته۔
\olpseganchor{OLP-0579-B004}{end}

\olpseganchor{OLP-0579-B005}{start}
البته څرګند ترکيبي تعريفونه \emph{شته}۔ يو يې دا دے۔
پرېږدئ چې \LR{$\text{\RL{finfun}}(\alpha,\beta)$} د ټولو هغو تابعو \LR{$f
\colon \alpha \to \beta$} سټ وي چې \LR{$\Setabs{\gamma \in
\alpha}{f(\gamma) \neq 0}$} يې له کوم طبيعي عدد سره هم‌شمېره وي۔
پر \LR{$\text{\RL{finfun}}(\alpha,\beta)$} دا ښه ترتيب وټاکئ: \LR{$f
\sqsubset g$} هغه وخت او يوازې هغه وخت چې \LR{$f \neq g$} او
\LR{$f(\gamma_0) < g(\gamma_0)$} وي، چېرته چې
\LR{$\gamma_0 = \text{\RL{max}}\Setabs{\gamma \in \alpha}{f(\gamma) \neq
g(\gamma)}$} دے۔ که \LR{$\alpha\neq0$} وي، بيا \LR{$\ordexpo{\alpha}{\beta}$} د
\LR{$\ordtype{\text{\RL{finfun}}(\beta, \alpha), \sqsubset}$} په توګه تعريفولے
شو۔ پاټر هم دغه څرګند تعريف کاروي۔
\begin{explain}
  د دې ترتيب موخه دا ده چې د ترتيبي توان د بازګښت قاعدې
  پوره کړي۔ په تالي پړاو کښې وروستے مدخل د مخکنيو تابعو
  مرتبې ډلې جوړوي؛ په حدي پړاو کښې د متناهي ملاتړ هره
  تابعه له يوه مخکيني پړاو څخه راځي۔ د دې جوړښت انځور
  د ترتيبي جمع او ضرب تر انځور زيات کار غواړي۔ \mbox{\LR{\latinfont \citep[p.~199]{Potter2004}}}
\end{explain}
\noindent د تابعې ساحه د توان شاخص دے، او د ارزښتونو
سټ يې بنسټ دے۔ د صفر بنسټ لپاره لاندې بازګښتي تعريف کاروو۔
د بېلګې په توګه، که \LR{$\alpha=0$} او \LR{$\beta=\omega$} وي،
\LR{$\text{\RL{finfun}}(\omega,0)$} تشه ده او ترتيب ډول يې \LR{$0$} دے،
خو لاندې بازګښت \LR{$\ordexpo{0}{\omega}=1$} ورکوي، ځکه په
حدي اتحاد کښې د صفر شاخص د توان ارزښت هم شامل دے۔
جمع مو داسې تشريح کړه چې «د \LR{$\alpha$} يوه کاپي او
ورپسې د \LR{$\beta$} يوه کاپي» ده، او ضرب مو داسې تشريح کړ چې «يوه
\LR{$\beta$}-لړۍ د \LR{$\alpha$} له کاپيو جوړه ده»۔ خو د
\LR{$\text{\RL{finfun}}(\gamma, \alpha)$} په اړه دومره لنډه او څرګنده خبره
نه لرو۔ نو د ترتيبي عددونو توان به د نامتناهيت هاخوا بازګښت \emph{له لارې}
تعريف کړو، يعنې:
\olpseganchor{OLP-0579-B005}{end}

\olpseganchor{OLP-0579-B006}{start}
\begin{defn}\ollabel{ordexporecursion}
\begin{align*}
	\ordexpo{\alpha}{0} &= 1\\
	\ordexpo{\alpha}{\beta\ordplus 1} &=\ordexpo{\alpha}{\beta} \ordtimes \alpha\\
	\ordexpo{\alpha}{\beta} &= \bigcup_{\delta < \beta}\ordexpo{\alpha}{\delta}& & \text{\RL{چې \LR{$\beta$} حدي ترتيبي عدد وي}}
\end{align*}
\end{defn}
\olpseganchor{OLP-0579-B006}{end}

\olpseganchor{OLP-0579-B007}{start}
که مو د سټ د نظريې د څېړونکو په توګه کار کاوه، ښايي د ترتيبي
عددونو د توان ځينې ځانګړنې مو نورې هم څېړلې واے۔ خو دلته له دې
څخه زيات څه نه ورزياتوو؛ يوازې دې اشنا حقيقت ته پام کوو چې د
ترتيبي عددونو توان تبادلي نه دے۔ ځکه
\LR{$\ordexpo{2}{\omega} =
\bigcup_{\delta < \omega}\ordexpo{2}{\delta} = \omega$}، خو
\LR{$\ordexpo{\omega}{2} = \omega \ordtimes \omega$}۔ په اصل کښې بايد
د توان د تبادلي والي \emph{تمه} ونه کړو، ځکه چې د طبيعي عددونو
دپاره هم تبادلي نه دے: \LR{$\ordexpo{2}{3} = 8 < 9 = \ordexpo{3}{2}$}۔
\olpseganchor{OLP-0579-B007}{end}

\olpseganchor{OLP-0579-B008}{start}
\begin{prob}
د نامتناهيت هاخوا استقرا په کارولو ثابته کړئ چې که \LR{$\alpha\neq0$} وي او
\LR{$\ordexpo{\alpha}{\beta} = \ordtype{\text{\RL{finfun}}(\beta, \alpha),
\sqsubset}$} تعريف کړو، نو د
\olref[sth][ord-arithmetic][expo]{ordexporecursion} بازګشتي معادلې
ترلاسه کوو۔ اشاره: پر تشه ساحه يوازې تشه تابعه شته۔ د
تالي په حالت کښې تابعې د وروستي مختص د ارزښت له مخې په
پرله‌پسې ډلو ووېشئ۔ په هره ډله کښې د مخکني پړاو ترتيب
راګرځي۔ په حدي حالت کښې د ناصفر ارزښتونو ځايونه متناهي دي؛
نو هره تابعه د کوم کوچني پړاو له تابعې څخه په صفرونو
غځېدلې ده۔ د همدغو غځونو پيلنۍ برخې د ټول ترتيب اتحاد ورکوي۔
\end{prob}
\olpseganchor{OLP-0579-B008}{end}

\olpseganchor{OLP-0580-B004}{start}
\olchapter{sth}{cardinals}{اصلي عددونه}
\olpseganchor{OLP-0580-B004}{end}

\olpseganchor{OLP-0581-B003}{start}
\olfileid{sth}{cardinals}{cp}
\olsection{د کانتور اصل}
\olpseganchor{OLP-0581-B003}{end}

\olpseganchor{OLP-0581-B004}{start}
خپل ذهن بېرته \olref[ordinals][vn]{sec} ته يوسئ۔ هلته مو د ښه
ترتيب شويو سټونو په اړه خبرې کولې او وړانديز مو وکړ چې ښه به وي
داسې څيزونه ولرو چې د ښه ترتيبونو استازيتوب وکړي۔ په همدې فکر مو
ترتيبي عددونه وپېژندل، او بيا مو په
\olref[ordinals][ordtype]{ordtypesworklikeyouwant} کښې وښودل چې دوی
هماغسې چلند کوي لکه چې غواړو، يعنې:
\[
	\ordtype{A, <} = \ordtype{B, \lessdot}
	\text{\RL{ هغه وخت او يوازې هغه وخت چې }} \tuple{A, <} \isomorphic \tuple{B, \lessdot}.
\]
ذهن لا پخوا، \olref[sfr][siz][equ]{sec} ته يوسئ۔ هلته مو په ساده
بڼه د سټ د «اندازې» مفهوم وپېژانده۔ په ځانګړې توګه مو وويل چې دوه
سټونه، \LR{$\cardeq{A}{B}$}، هغه وخت او يوازې هغه وخت هم‌شمېره دي چې
دوه اړخيزه يو پر يو تابع \LR{$f \colon A \to B$} موجود وي۔ د هم‌شمېرۍ
شرط تر ترتيبي هم‌شکلتيا {ساده} دے: د اندازې لپاره دوه اړخيزې يو پر يو تابعې
بس دي۔ د ترتيبي هم‌شکلتيا دوه اړخيزې يو پر يو تابعې بيا د ترتيب اړيکه هم
ساتي۔
\olpseganchor{OLP-0581-B004}{end}

\olpseganchor{OLP-0581-B005}{start}
له دې ټولو يو څرګند فکر راولاړېږي۔ لکه څنګه مو چې د ښه ترتيبونو د
سنجولو دپاره ځانګړي څيزونه، \emph{ترتيبي عددونه}، وپېژندل، همداسې
د اندازې د سنجولو دپاره هم ځانګړي څيزونه، \emph{اصلي عددونه}،
پېژندلے شو۔ د دې څپرکي موخه همدا ده۔
\olpseganchor{OLP-0581-B005}{end}

\olpseganchor{OLP-0581-B006}{start}
مخکښې له دې چې ووايو اصلي عددونه به څه وي، بايد هغه اصل وټاکو
چې دوی يې پوره کړي۔ که \LR{$\card{X}$} د سټ \LR{$X$} د اصلي شمېر دپاره
وليکو، غواړو چې لاندې شرط پوره کړي:
\[
	\card{A} = \card{B} \text{\RL{ هغه وخت او يوازې هغه وخت چې }} \cardeq{A}{B}.
\]
دا به د \emph{کانتور} اصل وبولو، ځکه چې ښايي کانتور لومړنے کس و
چې دا يې په څرګنده توګه په ذهن کښې درلود۔ (له \emph{هيوم} اصل
سره د دې د اړيکې په اړه به په \olref[hp]{sec} کښې نور ووايو۔) نو
موخه مو دا ده چې \LR{$\card{X}$} د هر \LR{$X$} دپاره داسې تعريف کړو چې د
کانتور اصل ترې ترلاسه شي۔
\olpseganchor{OLP-0581-B006}{end}

\olpseganchor{OLP-0582-B003}{start}
\olfileid{sth}{cardinals}{cardsasords}
\olsection{اصلي عددونه د ترتيبي عددونو په توګه}
\olpseganchor{OLP-0582-B003}{end}

\olpseganchor{OLP-0582-B004}{start}
په حقيقت کښې به زموږ د اصلي عددونو نظريه په څرګنده توګه د ترتيبي
عددونو له نظريې کار اخلي۔ يعنې: اصلي عددونه به د ځانګړو ترتيبي
عددونو په توګه تعريف کړو۔ په تېره بيا، لاندې تعريف وړاندې کوو:
\olpseganchor{OLP-0582-B004}{end}

\olpseganchor{OLP-0582-B005}{start}
\begin{defn}\ollabel{defcardinalasordinal}
که \LR{$A$} ښه ترتيبېدے شي، نو \LR{$\card{A}$} تر ټولو کوچنے ترتيبي عدد
\LR{$\gamma$} دے چې \LR{$\cardeq{A}{\gamma}$} وي۔ د هر ترتيبي عدد \LR{$\gamma$}
دپاره وايو چې \LR{$\gamma$} يو \emph{اصلي عدد} دے که او يوازې که
\LR{$\gamma = \card{\gamma}$} وي۔
\end{defn}
\olpseganchor{OLP-0582-B005}{end}

\olpseganchor{OLP-0582-B006}{start}
موږ اوس «\LR{$A$} ښه ترتيبېدے شي» وويل۔ لکه په رياضياتو کښې چې نږدې
تل وي، دلته د امکان عبارت يوازې د يوې وجودي ادعا لنډ بيان دے:
«\LR{$A$} ښه ترتيبېدے شي» دا وايي چې «داسې يوه اړيکه شته چې پر
\LR{$A$} ښه ترتيب جوړوي»۔
\olpseganchor{OLP-0582-B006}{end}

\olpseganchor{OLP-0582-B007}{start}
خو \olref{defcardinalasordinal} يوه ستونزه لري۔ غواړو چې \emph{هر}
سټ اندازه ولري، يعنې \LR{$\card{A}$} د هر~\LR{$A$} دپاره موجود وي۔ خو زموږ
تعريف په يوه شرط پيلېږي: «که \LR{$A$} ښه ترتيبېدے شي\ldots»۔ که
کوم سټ \LR{$A$} ښه ترتيب نه شي، نو دا تعريف د هغه دپاره \LR{$\card{A}$}
څيز نه ټاکي۔
\olpseganchor{OLP-0582-B007}{end}

\olpseganchor{OLP-0582-B008}{start}
نو د \olref{defcardinalasordinal} د کارولو دپاره پکار ده چې هر سټ
ښه ترتيبېدے شي۔ خو له بده مرغه، په~\LR{$\ZF$} کښې دا تضمين نه لرو۔
له همدې امله، که \olref{defcardinalasordinal} کاروو، بايد يو نوے
اصل هم ور زيات کړو، لکه:
\begin{axiom}[ښه ترتيب]
هر سټ ښه ترتيبېدے شي۔
\end{axiom}
په \olref[choice][]{chap} کښې به وڅېړو چې د ښه ترتيب اصل د منلو
وړ دے که نه۔ خو له دې ځايه به يې کاروو۔ په دې اصل سره دا ثابتول
اسانه دي چې اصلي عددونه (د \olref{defcardinalasordinal} د تعريف
له مخې) موجود دي او ښه چلند کوي:
\olpseganchor{OLP-0582-B008}{end}

\olpseganchor{OLP-0582-B009}{start}
\begin{lem}\ollabel{lem:CardinalsExist}
د هر سټ \LR{$A$} دپاره:
\begin{enumerate}
	\item\ollabel{cardaexists} \LR{$\card{A}$} موجود او يوازينے دے؛
	\item\ollabel{cardaapprox}  \LR{$\cardeq{\card{A}}{A}$}؛
	\item\ollabel{cardaidem}  \LR{$\card{A}$} اصلي عدد دے، يعنې
	\LR{$\card{A} = \card{\card{A}}$}؛
\end{enumerate}
\end{lem}
\olpseganchor{OLP-0582-B009}{end}

\olpseganchor{OLP-0582-B010}{start}
\begin{proof}
\LR{$A$} وټاکئ۔ د ښه ترتيب له مخې \LR{$\tuple{A, R}$} يو ښه ترتيب
شته۔ د \olref[ordinals][ordtype]{thmOrdinalRepresentation} له مخې
\LR{$\tuple{A, R}$} له يوه يوازيني ترتيبي عدد \LR{$\beta$} سره هم‌شکل دے۔
نو \LR{$\cardeq{A}{\beta}$}۔ تر دې شاهد پورې ټول هم‌شمېره
ترتيبي عددونه په بېلونه راټول کړئ۔ دا ناتش سټ د ښه ترتيب
له مخې يو تر ټولو کوچنے غړے لري۔ يعنې داسې يو
تر ټولو کوچنے ترتيبي عدد \LR{$\gamma$} شته چې \LR{$\cardeq{A}{\gamma}$}
وي۔ نو \LR{$\card{A} = \gamma$}، او \olref{cardaexists} او
\olref{cardaapprox} ثابت شول۔ د \olref{cardaidem} د ثبوت دپاره
وګورئ چې که \LR{$\delta \in \gamma$} وي، نو \LR{$\cardless{\delta}{A}$} د
\LR{$\gamma$} د ټاکنې له مخې: کوچنے ترتيبي عدد په لوی کښې
فرعي سټ دے، نو د ورګډولو تابعه له هغه څخه راځي۔ د لږوالي
له مخې دوه اړخيز يو پر يو سمون نشته۔ دا دواړه شرطونه د
هم‌شمېرۍ په دواړو لورو د تابعو له يوځاے کولو سره ساتل کېږي
(\olref[sfr][siz][equ]{equinumerosityisequi})، نو
\LR{$\cardless{\delta}{\gamma}$} هم دے۔ نو \LR{$\gamma =
\card{\gamma}$}۔
%So, for reductio, suppose that there is some ordinal $\delta \in \gamma$ such that $\cardeq{\gamma}{\delta}$. Then, $\cardeq{\cardeq{A}{\gamma}}{\delta}$ so that $\cardeq{A}{\delta}$ by \olref[sfr][set][equ]{equinumerosityisequi}, which contradicts the choice of $\gamma$ as the least ordinal such that $\cardeq{A}{\gamma}$. 
\end{proof}
\olpseganchor{OLP-0582-B010}{end}

\olpseganchor{OLP-0582-B011}{start}
راتلونکې پايله د کانتور اصل، او تر هغه زيات، تضمينوي۔ (په پام
کښې ونيسئ چې اصلي عددونه خپل ترتيب له ترتيبي عددونو اخلي، يعنې
\LR{$\cardfont{a} < \cardfont{b}$} هغه وخت او يوازې هغه وخت چې
\LR{$\cardfont{a} \in \cardfont{b}$} وي۔ د دې د بيان دپاره به د هغو
څيزونو دپاره چې اصلي عددونه يې ګڼو، فرکټور توري کاروو۔ دا عام
دود دے۔ يو بل عام دود د يوناني تورو کارول دي، ځکه چې اصلي عددونه
ترتيبي عددونه هم دي؛ خو د الفبا د منځ توري ټاکل کېږي، لکه
\LR{$\kappa, \lambda$}۔):
\begin{lem}\ollabel{lem:CardinalsBehaveRight}
د هر ډول سټونو \LR{$A$} او \LR{$B$} دپاره:
\begin{align*}
	\cardeq{A}{B} &\text{\RL{ هغه وخت او يوازې هغه وخت چې }} \card{A} = \card{B}\\
	\cardle{A}{B} &\text{\RL{ هغه وخت او يوازې هغه وخت چې }} \card{A} \leq \card{B}\\
	\cardless{A}{B}&\text{\RL{ هغه وخت او يوازې هغه وخت چې }} \card{A} < \card{B}
\end{align*}
\end{lem}
\olpseganchor{OLP-0582-B011}{end}

\olpseganchor{OLP-0582-B012}{start}
\begin{proof}
د دويمې ادعا له چپې خوا ښۍ خوا ته لوري نه پيل کوو۔
د پاتې لورو دليلونه ورپسې راځي۔ لومړے لاندې شکل وګورئ۔
\begin{center}
	\begin{tikzpicture}
	\node (nodea) {\LR{$A$}};
	\node[right = 6em of nodea] (nodeb) {\LR{$B$}};
	\node[below = 2em of nodea] (nodecarda) {\LR{$\card{A}$}};
	\node[below = 2em of nodeb] (nodecardb) {\LR{$\card{B}$}};
	\draw[->] (nodea)--(nodeb);
	\draw[<->] (nodea)--(nodecarda);
	\draw[<->] (nodeb)--(nodecardb);
	\draw[->, dashed] (nodecarda)--(nodecardb);
\end{tikzpicture}
\end{center}
دواړه‌سري غشي دوه اړخيزې يو پر يو تابعې ښيي چې د شتون تضمين يې
\olref{lem:CardinalsExist} کوي۔ له \LR{$\cardle{A}{B}$} څخه
يو پر يو تابع \LR{$A\to B$} شته۔ اوس د غشو په تعقيب له \LR{$\card{A}$}
څخه \LR{$A$} ته، بيا \LR{$B$} ته، او بيا \LR{$\card{B}$} ته ځو؛ په دې توګه
يو پر يو تابع \LR{$\card{A} \to \card{B}$} ترلاسه کوو (مات غشے)۔
د دې تابعې حاصل د ښي ترتيبي عدد فرعي سټ دے؛ د القا شوي
ترتيب ډول يې تر ښي ترتيبي عدد لوی نۀ دے۔ ځکه له فرعي سټ
څخه د ترلاسه شوي ترتيب هر غړے په اصلي ترتيب کښې لږ تر لږه
د خپل نوي شاخص هومره شاخص لري۔ دا په نامتناهيت هاخوا استقرا
ثابتېږي۔ حاصل له کيڼ سټ سره هم‌شمېره دے، نو د کيڼ اصلي
عدد لږوالے غوښتل شوې نابرابري ورکوي۔ په مقابل لوري کښې
د کوچني ترتيبي عدد د ورګډولو تابع واخلئ او له دواړو
دوه اړخيزو تابعو سره يې يوځاے کړئ۔ د برابرۍ ادعا د
هم‌شمېرۍ د انتقالي والي او د لږ ترتيبي استازي له ټاکلو راځي۔
کلکه نابرابري بيا د نابرابرۍ او د برابرۍ د نشتون ګډ شرط دے۔
\end{proof}\noindent له \olref{lem:CardinalsBehaveRight} څخه د
شرودر--برنشتاين قضيه هم بيا ثابتولے شو۔ دا ادعا ده چې که
\LR{$\cardle{A}{B}$} او \LR{$\cardle{B}{A}$} وي، نو \LR{$\cardeq{A}{B}$}۔
دا مو د \olref[sfr][siz][sb]{thm:schroder-bernstein} په توګه
بيان کړې وه، خو لومړے مو په \olref[sfr][infinite][card-sb]{sec}
کښې په يو څه هڅه ثابته کړه۔ اوس وګورئ:
\olpseganchor{OLP-0582-B012}{end}

\olpseganchor{OLP-0582-B013}{start}
\begin{proof}[د شرودر--برنشتاين بيا ثبوت]
که \LR{$\cardle{A}{B}$} او \LR{$\cardle{B}{A}$} وي، نو د
\olref{lem:CardinalsBehaveRight} له مخې \LR{$\card{A} \leq \card{B}$}
او \LR{$\card{B} \leq \card{A}$}۔ نو د درې‌ګوني وېش او
\olref{lem:CardinalsBehaveRight} له مخې \LR{$\card{A} = \card{B}$}
او \LR{$\cardeq{A}{B}$}۔
\end{proof}
\noindent
که څه هم دا ثبوت ډېر ساده دے، خو په ناڅرګنده توګه پر تعويض
(د \olref[ordinals][ordtype]{thmOrdinalRepresentation} د تضمين
دپاره) او پر ښه ترتيب (د \olref{lem:CardinalsBehaveRight} د
تضمين دپاره) دواړو تکيه کوي۔ پر خلاف يې، د
\olref[sfr][infinite][card-sb]{sec} ثبوت ډېر خپلواک و (په رښتيا
هم په~\LR{$\Zminus$} کښې ترسره کېدے شي)۔
\olpseganchor{OLP-0582-B013}{end}

\olpseganchor{OLP-0583-B003}{start}
\olfileid{sth}{cardinals}{zfc}
\olsection{\texorpdfstring{\LR{$\ZFC$}}{ZFC}: يو مهم پړاو}
\olpseganchor{OLP-0583-B003}{end}

\olpseganchor{OLP-0583-B004}{start}
د ښه ترتيب د اصل په ور زياتولو سره وروستي نظري پړاو ته رسېدلي
يو۔ اوس د~\LR{$\ZFC$} دپاره ټول اړين اصلونه لرو۔ په تفصيل سره:
\olpseganchor{OLP-0583-B004}{end}

\olpseganchor{OLP-0583-B005}{start}
\begin{defn}
د \LR{$\ZFC$} نظريه دا اصلونه لري: د غړو له مخې برابري، اتحاد، جوړې،
د فرعي سټونو سټونه، نامتناهيت، د بنسټ اصل، د ښه ترتيب اصل،
او د بېلونې او تعويض د شېمو ټولې بېلګې۔ په بله وينا، \LR{$\ZFC$} پر~\LR{$\ZF$} د ښه
ترتيب اصل ور زياتوي۔
\end{defn}
\olpseganchor{OLP-0583-B005}{end}

\olpseganchor{OLP-0583-B006}{start}
\LR{$\ZFC$} له \emph{ټاکنې} سره د \emph{زرمېلو--فرېنکل} د سټونو
تيوري ښيي۔ زموږ په لست کښې ورزيات شوے اصل د ښه ترتيب اصل دے۔
کله چې وروسته د {ټاکنې اصل} په رسمي ډول بيان کړو، څرګنده
به شي چې ښه ترتيب (د~\LR{$\ZF$} په چوکاټ کښې) له ټاکنې سره
همارزښته دے (وګورئ
\olref[choice][woproblem]{thmwochoice})۔ نو کوم يو چې خپل «بنسټيز»
اصل وګڼو، توپير نه کوي۔ او د «\LR{$\ZFC$}» نوم په ادبياتو کښې پوره
عام دے۔
\olpseganchor{OLP-0583-B006}{end}

\olpseganchor{OLP-0584-B003}{start}
\olfileid{sth}{cardinals}{classing}
\olsection{متناهي، د شمېر وړ او د شمېر نۀ وړ}
\olpseganchor{OLP-0584-B003}{end}

\olpseganchor{OLP-0584-B004}{start}
اوس چې اصلي عددونه مو وپېژندل، ښه ده د اصلي عددونو د بېلابېلو
ډولونو په اړه لږ وغږېږو؛ په ځانګړې توګه د متناهي،
د شمېر وړ او د شمېر نۀ وړ اصلي عددونو په اړه۔
\olpseganchor{OLP-0584-B004}{end}

\olpseganchor{OLP-0584-B005}{start}
زموږ له لومړنيو دوو پايلو څخه څرګندېږي چې متناهي اصلي عددونه به
عين متناهي ترتيبي عددونه وي؛ هماغه چې پخوا مو په
\olref[z][infinity-again]{defnomega} کښې د خپلو \emph{طبيعي
عددونو} په توګه تعريف کړل:
\olpseganchor{OLP-0584-B005}{end}

\olpseganchor{OLP-0584-B006}{start}
\begin{prop}\ollabel{finitecardisoequal}
فرض کړئ \LR{$n, m \in \omega$}۔ نو \LR{$n = m$} هغه وخت او يوازې هغه وخت
چې \LR{$\cardeq{n}{m}$} وي۔
\end{prop}
\olpseganchor{OLP-0584-B006}{end}

\olpseganchor{OLP-0584-B007}{start}
\begin{proof}
له \emph{چپې خوا ښۍ خوا ته} لورے څرګند دے۔ د \emph{ښۍ خوا
چپې خوا ته} لوري د ثبوت دپاره فرض کړئ چې \LR{$\cardeq{n}{m}$}، خو
\LR{$n \neq m$}۔ د درې‌ګوني وېش له مخې يا \LR{$n \in m$} يا \LR{$m \in n$}؛
له عموميت له کمېدو پرته \LR{$n \in m$} فرض کړئ۔ بيا \LR{$n \subsetneq m$}
او يوه دوه اړخيزه يو پر يو تابع \LR{$f \colon m \to n$} شته، نو \LR{$m$} د
ډېډېکېنډ په معنا نامتناهي دے؛ دا له
\olref[z][infinity-again]{naturalnumbersarentinfinite} سره ټکر لري۔
\end{proof}
\olpseganchor{OLP-0584-B007}{end}

\olpseganchor{OLP-0584-B008}{start}
\begin{cor}\ollabel{naturalsarecardinals}
که \LR{$n \in \omega$} وي، نو \LR{$n$} اصلي عدد دے۔
\end{cor}
\olpseganchor{OLP-0584-B008}{end}

\olpseganchor{OLP-0584-B009}{start}
\begin{proof}
سمدستي پايله ده۔
\end{proof}
\noindent
له دې دا هم څرګندېږي چې د يوه اصلي عدد د «متناهي» يا
«نامتناهي» بللو څو معقولې معناوې سره برابرې دي:
\begin{thm}\ollabel{generalinfinitycharacter}د هر سټ \LR{$A$} دپاره لاندې
شرطونه همارزښته دي:
\begin{enumerate}
	\item\ollabel{card:notinomega} \LR{$\card{A} \notin \omega$}، يعنې
	د \LR{$A$} اصلي شمېر طبيعي عدد نۀ دے؛
	\item\ollabel{card:omegaplus} \LR{$\omega \leq \card{A}$}؛
	\item\ollabel{card:infinite} \LR{$A$} د ډېډېکېنډ په معنا نامتناهي دے۔
\end{enumerate}
\end{thm}
\olpseganchor{OLP-0584-B009}{end}

\olpseganchor{OLP-0584-B010}{start}
\begin{proof}
دا د \olref[ord-arithmetic][using-addition]{ordinfinitycharacter}،
\olref[cardinals][cardsasords]{lem:CardinalsBehaveRight} او
\olref{naturalsarecardinals} له مخې دي۔ د سټ او د هغه د
اصلي عدد تر منځ دوه اړخيز يو پر يو سمون د ډېډېکېنډ په معنا
نامتناهي والى دواړو لورو ته لېږدوي: يو پر يو ځان‌تابعه د
دغه سمون او د هغه د معکوس په وسيله بلې خوا ته وړو؛ له
حاصل څخه بهر غړے هم ورسره لېږدېږي۔
\end{proof}
\olpseganchor{OLP-0584-B010}{end}

\olpseganchor{OLP-0584-B011}{start}
له دې سره هغه ځينې مفهومونه په \emph{رسمي ډول تعريفولے} شو چې
په \olref[sfr][][]{part} کښې مو لږ په غيررسمي توګه کارول:
\olpseganchor{OLP-0584-B011}{end}

\olpseganchor{OLP-0584-B012}{start}
\begin{defn}\ollabel{defnfinite}
وايو چې \LR{$A$} \emph{متناهي} دے که او يوازې که \LR{$\card{A}$} يو
طبيعي عدد وي، يعنې \LR{$\card{A} \in \omega$}۔ کنه وايو چې \LR{$A$}
\emph{نامتناهي} دے۔
\end{defn}
\noindent
خو پام وکړئ چې دا تعريف د \LR{$\ZFC$} په چوکاټ کښې وړاندې شوے دے۔
بالاخره، د هر سټ د اصلي شمېر د شتون د تضمين دپاره مو ښه ترتيب
ته اړتيا درلوده۔ په رښتيا، له ښه ترتيب پرته داسې سټ شته کېدے شي
چې نه متناهي وي او نه د ډېډېکېنډ په معنا نامتناهي۔ دې ډول مسئلې
ته به په \olref[choice][]{chap} کښې بېرته راوګرځو۔ اوس مهال پر
ښه ترتيب تکيه کوو۔
\olpseganchor{OLP-0584-B012}{end}

\olpseganchor{OLP-0584-B013}{start}
اوس له متناهي اصلي عددونو څخه نامتناهي اصلي عددونو ته راځو۔
دا دوه لومړنۍ پايلې دي:
\olpseganchor{OLP-0584-B013}{end}

\olpseganchor{OLP-0584-B014}{start}
\begin{cor}\ollabel{omegaisacardinal}
\LR{$\omega$} تر ټولو کوچنے نامتناهي اصلي عدد دے۔
\end{cor}
\olpseganchor{OLP-0584-B014}{end}

\olpseganchor{OLP-0584-B015}{start}
\begin{proof}
\LR{$\omega$} اصلي عدد دے، ځکه چې \LR{$\omega$} د ډېډېکېنډ په معنا
نامتناهي دے او که \LR{$\cardeq{\omega}{n}$} د کوم \LR{$n \in \omega$}
دپاره وي، نو \LR{$n$} به هم د ډېډېکېنډ په معنا
نامتناهي وي؛ دا له
\olref[z][infinity-again]{naturalnumbersarentinfinite} سره ټکر لري۔
اوس د تعريف له مخې \LR{$\omega$} تر ټولو کوچنے نامتناهي اصلي عدد دے۔
\end{proof}
\olpseganchor{OLP-0584-B015}{end}

\olpseganchor{OLP-0584-B016}{start}
\begin{cor}
هر نامتناهي اصلي عدد يو حدي ترتيبي عدد دے۔
\end{cor}
\olpseganchor{OLP-0584-B016}{end}

\olpseganchor{OLP-0584-B017}{start}
\begin{proof}
فرض کړئ \LR{$\alpha$} يو نامتناهي تالي ترتيبي عدد وي، نو
\LR{$\alpha = \beta \ordplus 1$} د کوم \LR{$\beta$} دپاره دے۔ د
\olref{finitecardisoequal} په رڼا کښې د متناهي ترتيبي
عدد تالی هم متناهي دے۔ نو \LR{$\beta$} هم نامتناهي دے؛ له دې امله د
\olref[ord-arithmetic][using-addition]{ordinfinitycharacter} له مخې
\LR{$\cardeq{\beta}{\beta \ordplus 1}$}۔ اوس د
\olref[cardinals][cardsasords]{lem:CardinalsBehaveRight} له مخې
\LR{$\card{\beta} = \card{\beta\ordplus 1} = \card{\alpha}$}؛ خو
د مخکني ترتيبي عدد اصلي عدد تر هماغه مخکني عدد لوی نۀ دے۔
ځکه مخکنی ترتيبي عدد پخپله يو هم‌شمېره شاهد دے۔ نو
\LR{$\alpha \neq \card{\alpha}$}۔
\end{proof}
\olpseganchor{OLP-0584-B017}{end}

\olpseganchor{OLP-0584-B018}{start}
موږ له وړاندې، په \olref[sfr][siz][enm-alt]{defn:enumerable} کښې،
دې ته پام اړولے و چې د نامتناهي سټونو تر منځ د شمېر وړ او
د شمېر نۀ وړ توپيرولے شو۔ هغه تعريف په طبيعي ډول لاندې
پايله راکوي:
\olpseganchor{OLP-0584-B018}{end}

\olpseganchor{OLP-0584-B019}{start}
\begin{prop}
\LR{$A$} هغه وخت او يوازې هغه وخت د شمېر وړ دے چې
\LR{$\card{A} \leq \omega$} وي؛ او \LR{$A$} هغه وخت او يوازې هغه وخت
د شمېر نۀ وړ دے چې \LR{$\omega < \card{A}$} وي۔
\end{prop}
\olpseganchor{OLP-0584-B019}{end}

\olpseganchor{OLP-0584-B020}{start}
\begin{proof}
د درې‌ګوني وېش له مخې دواړه ادعاوې همارزښې دي، نو بس ده دا
ثابته کړو چې \LR{$A$} هغه وخت او يوازې هغه وخت د شمېر وړ دے
چې \LR{$\card{A} \leq \omega$} وي۔ له \emph{ښۍ خوا چپې خوا ته}:
که \LR{$\card{A} \leq \omega$} وي، نو د
\olref[cardinals][cardsasords]{lem:CardinalsBehaveRight} او
\olref{omegaisacardinal} له مخې \LR{$\cardle{A}{\omega}$}۔ له
دې څخه شمېر وړتيا راځي، ځکه طبيعي عددونو ته يو پر يو تابعه
د خپل حاصل پر سر يو دوه اړخيز سمون ورکوي۔ تش حاصل، متناهي
حاصل او نامتناهي حاصل د اړوندې شمېرنې درې حالتونه دي۔ له
\emph{چپې خوا ښۍ خوا ته}: که \LR{$A$} د شمېر وړ وي، نو د
\olref[sfr][siz][enm-alt]{defn:enumerable} له مخې درې حالتونه دي:
\begin{enumerate}
	\item که \LR{$A = \emptyset$} وي، نو د \olref{naturalsarecardinals}
	او \olref[cardinals][cardsasords]{lem:CardinalsBehaveRight} له مخې
	\LR{$\card{A} = 0 \in \omega$}؛
	\item که \LR{$\cardeq{n}{A}$} وي، نو د \olref{naturalsarecardinals}
	او \olref[cardinals][cardsasords]{lem:CardinalsBehaveRight} له مخې
	\LR{$\card{A} = n \in \omega$}؛
	\item که \LR{$\cardeq{\omega}{A}$} وي، نو د \olref{omegaisacardinal}
	له مخې \LR{$\card{A} = \omega$}۔
\end{enumerate}
نو په ټولو حالتونو کښې \LR{$\card{A} \leq \omega$}۔
\end{proof}
\noindent
په رښتيا، \LR{$\omega$} ځانګړے ځاے لري۔ که څه هم شمېر وړ ترتيبي
عددونه ډېر دي:
\olpseganchor{OLP-0584-B020}{end}

\olpseganchor{OLP-0584-B021}{start}
\begin{cor}
\LR{$\omega$} يوازينے د شمېر وړ نامتناهي اصلي عدد دے۔
\end{cor}
\olpseganchor{OLP-0584-B021}{end}

\olpseganchor{OLP-0584-B022}{start}
\begin{proof}
فرض کړئ \LR{$\cardfont{a}$} يو د شمېر وړ نامتناهي اصلي عدد
وي۔ ځکه چې \LR{$\cardfont{a}$} نامتناهي دے، \LR{$\omega \leq
\cardfont{a}$}۔ ځکه چې \LR{$\cardfont{a}$} يو د شمېر وړ اصلي
عدد دے، \LR{$\cardfont{a} = \card{\cardfont{a}} \leq \omega$}۔ نو
د درې‌ګوني وېش له مخې \LR{$\cardfont{a} = \omega$}۔
\end{proof}
\olpseganchor{OLP-0584-B022}{end}

\olpseganchor{OLP-0584-B023}{start}
د اصلي عددونو لړۍ نامتناهي ده۔ نو پوښتنه کولے شو: \emph{څو
اصلي عددونه شته؟} لاندې پايلې ښيي چې ښايي دا پوښتنه له سره
وڅېړو۔
\olpseganchor{OLP-0584-B023}{end}

\olpseganchor{OLP-0584-B024}{start}
\begin{prop}\ollabel{unioncardinalscardinal}
که د \LR{$X$} هر غړے اصلي عدد وي، نو \LR{$\bigcup X$} هم اصلي عدد دے۔
\end{prop}
\olpseganchor{OLP-0584-B024}{end}

\olpseganchor{OLP-0584-B025}{start}
\begin{proof}
دا په اسانه کتل کېدے شي چې \LR{$\bigcup X$} ترتيبي عدد دے۔
فرض کړئ \LR{$\alpha \in \bigcup X$} يو ترتيبي عدد وي؛ نو
\LR{$\alpha \in \cardfont{b} \in X$} د کوم اصلي عدد \LR{$\cardfont{b}$}
دپاره وي۔ ځکه چې \LR{$\cardfont{b}$} اصلي عدد دے،
\LR{$\cardless{\alpha}{\cardfont{b}}$}۔ ځکه چې \LR{$\cardfont{b}
\subseteq \bigcup X$}، نو \LR{$\cardle{\cardfont{b}}{\bigcup X}$}،
او له دې \LR{$\cardneq{\alpha}{\bigcup X}$}۔ که دواړه
هم‌شمېره واے، له لوی اتحاد څخه د سمون معکوس به د منځني
اصلي عدد لپاره په کوچني عدد کښې يو پر يو تابعه ورکړې واے۔
له دواړو يو پر يو تابعو به دوه اړخيز سمون راوتلے واے، چې
له کلکې نابرابرۍ سره ټکر لري۔ په عمومي کولو سره
\LR{$\bigcup X$} اصلي عدد دے؛ د تشې کورنۍ اتحاد هم صفر دے۔
\end{proof}
\olpseganchor{OLP-0584-B025}{end}

\olpseganchor{OLP-0584-B026}{start}
\begin{thm}\ollabel{lem:NoLargestCardinal}
تر ټولو لوے اصلي عدد نشته۔
\end{thm}
\olpseganchor{OLP-0584-B026}{end}

\olpseganchor{OLP-0584-B027}{start}
\begin{proof}
د هر اصلي عدد \LR{$\cardfont{a}$} دپاره د کانتور قضيه
(\olref[sfr][siz][car]{thm:cantor}) او
\olref[cardinals][cardsasords]{lem:CardinalsExist} دا راکوي چې
\LR{$\cardfont{a} < \card{\Pow{\cardfont{a}}}$}۔
\end{proof}
\olpseganchor{OLP-0584-B027}{end}

\olpseganchor{OLP-0584-B028}{start}
\begin{thm}
د ټولو اصلي عددونو سټ نشته۔
\end{thm}
\olpseganchor{OLP-0584-B028}{end}

\olpseganchor{OLP-0584-B029}{start}
\begin{proof}
د تناقض دپاره فرض کړئ \LR{$C = \Setabs{\cardfont{a}}{\cardfont{a} \text{\RL{
يو اصلي عدد دے}}}$}۔ اوس د \olref{unioncardinalscardinal} له مخې
\LR{$\bigcup C$} اصلي عدد دے، نو د \olref{lem:NoLargestCardinal} له مخې
يو اصلي عدد \LR{$\cardfont{b} > \bigcup C$} شته۔ د تعريف له مخې
\LR{$\cardfont{b} \in C$}، نو \LR{$\cardfont{b} \subseteq \bigcup{C}$}،
او له دې \LR{$\cardfont{b} \leq \bigcup C$}؛ دا تناقض دے۔
\end{proof}
\olpseganchor{OLP-0584-B029}{end}

\olpseganchor{OLP-0584-B030}{start}
دا د رسل له تناقض او بورالي--فورتي دواړو سره پرتله کړئ۔
\olpseganchor{OLP-0584-B030}{end}

\olpseganchor{OLP-0585-B003}{start}
\olfileid{sth}{cardinals}{hp}
\olsection{ضميمه: د هيوم اصل}
\olpseganchor{OLP-0585-B003}{end}

\olpseganchor{OLP-0585-B004}{start}
په \olref[cp]{sec} کښې مو د کانتور اصل بيان کړ۔ هغه دا و:
\begin{align*}
	\card{A} = \card{B} & \text{\RL{ هغه وخت او يوازې هغه وخت چې }} A \approx B.
\intertext{\RL{دا له هغه څه سره ډېر ورته دے چې اوس ورته
\emph{د هيوم اصل} ويل کېږي؛ هغه وايي:}}
	\fregenum{x} {F(x)} = \fregenum{x}{G(x)} & \text{\RL{ هغه وخت او يوازې هغه وخت چې }} F \sim G
\end{align*}
دلته «\LR{$F \sim G$}» په دې معنا لنډه نښه ده چې د \LR{$F$}ګانو شمېر
عين د \LR{$G$}ګانو له شمېر سره برابر دے؛ يعنې \LR{$F$}ګانې له \LR{$G$}ګانو
سره په يوه دوه اړخيزه يو پر يو سمون کښې اېښودل کېدے شي، يعنې:
\begin{align*}
	\exists R(&\forall v\forall y(Rvy \lif (Fv \land Gy)) \land {}\\
		&\forall v(Fv \lif \lexists![y][Rvy]) \land {}\\
		&\forall y(Gy \lif \lexists![v][Rvy]))
\end{align*}
خو د هيوم اصل او د کانتور اصل تر منځ د ډول توپير شته۔ د کانتور
د اصل په بيان کښې «\LR{$A$}» او «\LR{$B$}» د لومړۍ درجې ترمونه دي چې د
\emph{سټونو} استازيتوب کوي۔ د هيوم د اصل په بيان کښې «\LR{$F$}»،
«\LR{$G$}» او «\LR{$R$}» د لومړۍ درجې ترمونه \emph{نه} دي؛ بلکې د
\emph{پريديکات په ځاے} کښې دي۔ (ښايي د \emph{خاصيتونو}
استازيتوب کوي۔) نو د هيوم اصل په نثري ډول داسې بيانولے شو: د
\LR{$F$}ګانو شمېر د \LR{$G$}ګانو له شمېر سره هغه وخت او يوازې هغه وخت
برابر دے چې \LR{$F$}ګانې له~\LR{$G$}ګانو سره دوه اړخيزه يو پر يو سمون
ولري۔ دې ته \emph{د هيوم اصل} ويل کېږي، ځکه چې هيوم يو وخت
داسې ليکلي وو:
\begin{quote}
  کله چې دوه عددونه داسې سره جوړ کړو چې د يوه د هر واحد په مقابل
  کښې د بل يوه واحد تل وي، دواړه سره برابر بولو۔
  \mbox{\LR{\latinfont \citep[Bk.~I, Pt.~III, \S1, \P5]{Hume1740}}}
\end{quote}
او \mbox{\LR{\latinfont \citet[\S63]{Frege1884}}} د هيوم له دې عبارت څخه نقل وکړ،
بيا يې (په اصل کښې) هغه څه تشريح کړل چې موږ د هيوم اصل نومولے
دے؛ له دې سره د هيوم اصل په اوسني رياضيکي منطق کښې مشهور شو۔
\olpseganchor{OLP-0585-B004}{end}

\olpseganchor{OLP-0585-B005}{start}
د هيوم د اصل او د فرېګه د پنځم بنسټيز قانون تر منځ جوړښتي ورته
والي ته پام وکړئ۔ موږ هغه په \olref[story][blv]{sec} کښې داسې
بيان کړے و:
\[
	\fregeext{x}{F(x)} = \fregeext{x}{G(x)} \text{\RL{هغه وخت او يوازې هغه وخت چې }} \lforall[x][(F(x) \liff G(x))].
\]
دلته يو څه تشريح او پرتله ګټوره ده۔
\olpseganchor{OLP-0585-B005}{end}

\olpseganchor{OLP-0585-B006}{start}
د هيوم د اصل يا د پنځم بنسټيز قانون په شان اصل په دوو ډولونو
اخيستلے شو: \emph{پريديکاتيف} يا \emph{امپريديکاتيف} (وګورئ
\olref[story][predicative]{sec})۔ د پنځم بنسټيز قانون په
امپريديکاتيف لوست کښې، د هر~\LR{$F$} دپاره څيز
\LR{$\fregeext{x}{F(x)}$} د کمي کولو په هماغه دامن کښې راځي چې د
پنځم بنسټيز قانون د بيان دپاره مو کارولې وه۔ په ورته ډول، د
هيوم د اصل په امپريديکاتيف لوست کښې، د هر~\LR{$F$} دپاره څيز
\LR{$\fregenum{x}{F(x)}$} د کمي کولو په هماغه دامن کښې راځي چې د
هيوم د اصل د بيان دپاره مو کارولې وه۔ برعکس، په
\emph{پريديکاتيف} لوست کښې به څيزونه
\LR{$\fregeext{x}{F(x)}$} او~\LR{$\fregenum{x}{F(x)}$} د کوم
\emph{بل} دامن موجودات وي۔
\olpseganchor{OLP-0585-B006}{end}

\olpseganchor{OLP-0585-B007}{start}
که پنځم بنسټيز قانون امپريديکاتيف ولولو، د ساده ادراک له لارې
تناقض ته رسېږي (د تفصيل دپاره \olref[story][blv]{sec} وګورئ)۔
د ساده ادراک په څېر، د هغه پريديکاتيف لوست يې له تناقضه
ساتلے شي۔ خو ښايي بيا هغه ټول کارونه ونه کړي چې موږ ترې غوښتل۔
\olpseganchor{OLP-0585-B007}{end}

\olpseganchor{OLP-0585-B008}{start}
خو د هيوم اصل \emph{کېدے شي} په امپريديکاتيف ډول بې له تناقضه
ولوستل شي۔ د بېلګې په توګه، د څيزونو ساحه طبيعي عددونه
واخلئ او د هر يوځايي پريديکات لپاره يې د پوره کوونکو سټ
وګورئ۔ که دغه سټ متناهي وي، هغه ته د غړو له شمېر څخه يو
زيات طبيعي عدد ورکړئ۔ که نامتناهي وي، هغه ته صفر ورکړئ۔
د طبيعي عددونو دوه نامتناهي فرعي سټونه هم‌شمېره دي، او
متناهي فرعي سټونه هله هم‌شمېره دي چې شمېر يې برابر وي۔
نو د شمېر ټاکونکې تابعه د اصل دواړه لوري پوره کوي او ټول
ارزښتونه يې په هماغه ساحه کښې دي۔
\olpseganchor{OLP-0585-B008}{end}

\olpseganchor{OLP-0585-B009}{start}
د روښانولو دپاره پريديکات «\LR{$x \neq x$}» وګورئ؛ څرګنده ده چې هېڅ
څيز يې نه پوره کوي۔ د هيوم اصل اوس يو څيز \LR{$\# x( x\neq x)$}
راکوي۔ دا ښايي د~\LR{$0$} عدد وګڼو۔ اوس د \emph{امپريديکاتيف} لوست
له مخې---خو \emph{يوازې} له همدې لوست سره---دا څيز \LR{$0$} زموږ
د کمي کولو په اصلي دامن کښې شامل دے۔ نو د هيوم اصل پر پريديکات
«\LR{$x = 0$}» په معنا لرونکي ډول کارولے شو او څيز
\LR{$\#x (x = 0)$} ترلاسه کولے شو۔ دا ښايي د~\LR{$1$} عدد وګڼو۔ سربېره
پر دې، د هيوم اصل وايي چې \LR{$0 \neq 1$}، ځکه د هغو څيزونو چې له
ځان سره برابر نه دي او هغو چې له \LR{$0$} سره برابر دي تر منځ دوه
اړخيز يو پر يو سمون نشته (له لومړيو هېڅ نشته، خو له دويمو يو
شته)۔ اوس بيا په امپريديکاتيف لوست کښې \LR{$1$} زموږ د کمي کولو په
اصلي دامن کښې شامل دے۔ نو د هيوم اصل پر پريديکات
«\LR{$(x = 0 \lor x = 1)$}» کارولے شو او څيز
\LR{$\#x(x = 0 \lor x = 1)$} ترلاسه کولے شو۔ دا ښايي د~\LR{$2$} عدد
وګڼو، او ښودلے شو چې \LR{$0\neq 2$} او \LR{$1 \neq 2$}، او همداسې نور۔
\olpseganchor{OLP-0585-B009}{end}

\olpseganchor{OLP-0585-B010}{start}
په لنډه، د هيوم اصل په امپريديکاتيف لوست کښې د \emph{نامتناهي
ډېرو څيزونو} شتون راکوي۔ دې خبرې \emph{نويو فرېګه‌پالو
منطق‌پوهانو} ته هڅونه ورکړې چې د هيوم اصل د حساب بنسټ ونيسي۔
\olpseganchor{OLP-0585-B010}{end}

\olpseganchor{OLP-0585-B011}{start}
خو فرېګه \emph{پخپله} د هيوم اصل د حساب بنسټ نۀ ګرځاوه۔ د دې
پر ځاے يې د هيوم اصل له يوه څرګند تعريفه ثابت کړ:
\LR{$\fregenum{x}{F(x)}$} د مفهوم \LR{$F \sim \Phi$} د غځونې په توګه
تعريفېږي۔ په اوسنۍ ژبه کښې ښايي دا په
\LR{$\fregenum{x}{F(x)} = \Setabs{G}{F \sim G}$} بيان کړو؛ خو دا مو
بيا د ساده ادراک ستونزو ته راګرځوي۔
\olpseganchor{OLP-0585-B011}{end}

\olpseganchor{OLP-0586-B004}{start}
\olchapter{sth}{card-arithmetic}{د اصلي عددونو حساب}
\olpseganchor{OLP-0586-B004}{end}

\olpseganchor{OLP-0587-B003}{start}
\olfileid{sth}{card-arithmetic}{opps}
\olsection{د بنسټيزو عملياتو تعريف}
\olpseganchor{OLP-0587-B003}{end}

\olpseganchor{OLP-0587-B004}{start}
څرنګه چې دلته د ترتيب ساتلو ته اړتيا نه لرو، د اصلي عددونو حساب
د ترتيبي عددونو تر حسابه يو څه په اسانه تعريفېږي۔ جمع، ضرب او
توان به يو ځاے تعريف کړو۔
\olpseganchor{OLP-0587-B004}{end}

\olpseganchor{OLP-0587-B005}{start}
\begin{defn}
کله چې \LR{$\cardfont{a}$} او \LR{$\cardfont{b}$} اصلي عددونه وي:
\begin{align*}
	\cardfont{a} \cardplus \cardfont{b} &\defis
	\card{\cardfont{a} \disjointsum \cardfont{b}}\\
	\cardfont{a} \cardtimes \cardfont{b} &\defis
	\card{\cardfont{a} \times \cardfont{b}}\\
	\cardexpo{\cardfont{a}}{\cardfont{b}} &\defis
	\card{\funfromto{\cardfont{b}}{\cardfont{a}}}
\end{align*}
چې دلته \LR{$\funfromto{X}{Y} = \Setabs{f}{f \text{\RL{تابع ده }} X \to Y}$}۔
(دا ښودل اسانه دي چې \LR{$\funfromto{X}{Y}$} د هر دوو سټونو \LR{$X$}
او \LR{$Y$} دپاره شته؛ د تابعو ګرافونه د کارتيزي حاصل ضرب
فرعي سټونه دي، او له هغو څخه ټولې بشپړې تابعې په بېلونه
اخلو۔ د دې جوړښت تفصيل د تمرين په توګه پرېږدو۔)
\end{defn}
\olpseganchor{OLP-0587-B005}{end}

\olpseganchor{OLP-0587-B006}{start}
\begin{prob}
په \LR{$\Zminus$} کښې ثابته کړئ چې \LR{$\funfromto{X}{Y}$} د هر دوو
سټونو \LR{$X$} او~\LR{$Y$} دپاره شته۔ په \LR{$\ZF$} کښې کار وکړئ او
\LR{$\setrank{\funfromto{X}{Y}}$} د \LR{$\setrank{X}$} او \LR{$\setrank{Y}$}
له مخې محاسبه کړئ، لکه څنګه چې په
\olref[sth][ord-arithmetic][using-addition]{rankcomputation} کښې شوي دي۔
\end{prob}
\olpseganchor{OLP-0587-B006}{end}

\olpseganchor{OLP-0587-B007}{start}
د دې تعريف تشريح به مرسته وکړي۔ د جمعې په برخه کښې: دلته د
بېلې مجموعې مفهوم، \LR{$\disjointsum$}، کارول کېږي، چې په
\olref[ord-arithmetic][add]{defdissum} کښې تعريف شوے دے؛ او
په اسانه ښکاري چې دا تعريف د متناهي حالتونو دپاره سمه پايله
ورکوي۔ د ضرب په برخه کښې: \olref[sfr][set][pai]{cardnmprod}
موږ ته وايي چې که \LR{$A$} د \LR{$n$} غړي او \LR{$B$} د \LR{$m$} غړي ولري، نو
\LR{$A \times B$} د \LR{$n \cdot m$} غړي لري؛ نو زموږ تعريف همدا تصور
له متناهي څخه نامتناهي ضرب ته غځوي۔ توان همداسې دے: د متناهي
حالت تصور نامتناهي حالت ته غځوو۔ په رښتيا، له ځينو اړخونو د
نامتناهي اصلي عددونو حساب د نامتناهي ترتيبي عددونو تر حسابه
«عادي» حساب ته ډېر ورته دے:
\olpseganchor{OLP-0587-B007}{end}

\olpseganchor{OLP-0587-B008}{start}
\begin{prop}\ollabel{cardplustimescommute}
\LR{$\cardplus$} او \LR{$\cardtimes$} دواړه تبادلي او اتحادي دي۔
\end{prop}
\olpseganchor{OLP-0587-B008}{end}

\olpseganchor{OLP-0587-B009}{start}
\begin{proof}
د تبادلي والي دپاره، د
\olref[cardinals][cardsasords]{lem:CardinalsBehaveRight} له مخې
بس ده پام وکړو چې \LR{$\cardeq{(\cardfont{a} \disjointsum
\cardfont{b})}{(\cardfont{b} \disjointsum \cardfont{a})}$} او
\LR{$\cardeq{(\cardfont{a} \times \cardfont{b})}{(\cardfont{b} \times
\cardfont{a})}$}۔ اتحادي والى د تمرين په توګه پرېږدو۔
\end{proof}
\olpseganchor{OLP-0587-B009}{end}

\olpseganchor{OLP-0587-B010}{start}
\begin{prob}
ثابته کړئ چې \LR{$\cardplus$} او \LR{$\cardtimes$} اتحادي دي۔
\end{prob}
\olpseganchor{OLP-0587-B010}{end}

\olpseganchor{OLP-0587-B011}{start}
\begin{prop}
\LR{$A$} هغه وخت او يوازې هغه وخت نامتناهي دے چې
\LR{$\card{A} \cardplus 1 = 1 \cardplus \card{A} = \card{A}$} وي۔
\end{prop}
\olpseganchor{OLP-0587-B011}{end}

\olpseganchor{OLP-0587-B012}{start}
\begin{proof}
لکه په
\olref[cardinals][classing]{generalinfinitycharacter} کښې،
دا له \olref[ord-arithmetic][using-addition]{ordinfinitycharacter} او
\olref[cardinals][cardsasords]{lem:CardinalsBehaveRight} څخه
راځي۔
\end{proof}
\olpseganchor{OLP-0587-B012}{end}

\olpseganchor{OLP-0587-B013}{start}
دا څرګندوي چې ولې د ترتيبي او اصلي عددونو د جمعې او ضرب دپاره
بېلابېلو نښو ته اړتيا لرو: دا په رښتيا \emph{بېلابېل} عمليات
دي۔ د پايلو راتلونکې جوړه ښيي چې ترتيبي او اصلي توان هم
بېلابېل عمليات دي۔ (په ياد ولرئ چې له
\olref[z][infinity-again]{defnomega} څخه \LR{$2 = \{0,
1\}$} راځي):
\olpseganchor{OLP-0587-B013}{end}

\olpseganchor{OLP-0587-B014}{start}
\begin{lem}\ollabel{lem:SizePowerset2Exp}
\LR{$\card{\Pow{A}} = \cardexpo{2}{\card{A}}$}، د هر \LR{$A$} دپاره۔
\end{lem}
\olpseganchor{OLP-0587-B014}{end}

\olpseganchor{OLP-0587-B015}{start}
\begin{proof}
د هر فرعي سټ \LR{$B \subseteq A$} دپاره، \LR{$\chi_B \in
\funfromto{A}{2}$} داسې وټاکئ:
\begin{align*}
	\chi_{B}(x) &\defis
	\begin{cases}
		1 & \text{\RL{که }}x\in B\\
		0 & \text{\RL{کنه.}}
	\end{cases}
\end{align*}
اوس \LR{$f(B) = \chi_B$} وټاکئ؛ دا يوه دوه اړخيزه يو پر يو تابع
\LR{$f \colon \Pow{A} \to \funfromto{A}{2}$} تعريفوي۔ نو
\LR{$\cardeq{\Pow{A}}{\funfromto{A}{2}}$}۔ له دې
\LR{$\cardeq{\Pow{A}}{\funfromto{\card{A}}{2}}$} هم راځي، او بيا
\LR{$\card{\Pow{A}} = \card{\funfromto{\card{A}}{2}} =
2^{\card{A}}$}۔
\end{proof}
\olpseganchor{OLP-0587-B015}{end}

\olpseganchor{OLP-0587-B016}{start}
دا لنډ ثبوت په اصل کښې د \olref[sfr][siz][red-alt]{sec} بحث
هم رانغاړي۔ هلته مو وښودله چې څنګه د \LR{$\Pow{\omega}$}
نۀ‌شمېر وړتيا د نامتناهي دوه‌ګونو لړيو د سټ، \LR{$\Bin^\omega$}،
نۀ‌شمېر وړتيا ته «راوګرځوو»۔ په عمل کښې، \LR{$\Bin^{\omega}$}
هماغه \LR{$\funfromto{\omega}{2}$} دے؛ او مخکيني ثبوت وښودله چې
په \olref[sfr][siz][red-alt]{sec} کښې زموږ استدلال د
هر سټ~\LR{$A$} دپاره، د~\LR{$\omega$} پر ځاے، هم کار کوي۔ دا پايله د
اصلي توان په اړه يو ژر تر لاسه کېدونکے حقيقت هم راکوي:
\olpseganchor{OLP-0587-B016}{end}

\olpseganchor{OLP-0587-B017}{start}
\begin{cor}\ollabel{cantorcor}
\LR{$\cardfont{a} < \cardexpo{2}{\cardfont{a}}$} د هر اصلي عدد~\LR{$\cardfont{a}$} دپاره۔
\end{cor}
\olpseganchor{OLP-0587-B017}{end}

\olpseganchor{OLP-0587-B018}{start}
\begin{proof}
دا د کانتور له قضيې (\olref[sfr][siz][car]{thm:cantor}) او
\olref{lem:SizePowerset2Exp} څخه راځي۔
\end{proof}
\noindent
نو \LR{$\omega < \cardexpo{2}{\omega}$}۔ خو پام وکړئ: دا پايله
د \emph{اصلي عددونو} د توان په اړه ده۔ له \emph{ترتيبي
عددونو} د توان سره بايد پرتله شي، ځکه په وروستي حالت کښې
\LR{$\omega = \ordexpo{2}{\omega}$} (وګورئ
\olref[ord-arithmetic][expo]{sec})۔
\olpseganchor{OLP-0587-B018}{end}

\olpseganchor{OLP-0587-B019}{start}
څرنګه چې اوس د اصلي توان په اړه غږېږو، د هغې «لارې» په اړه هم
لږ ډېر کره کېدے شو چې \LR{$\Real$} پرې د شمېر نۀ وړ دے۔
\olpseganchor{OLP-0587-B019}{end}

\olpseganchor{OLP-0587-B020}{start}
\begin{thm}\ollabel{continuumis2aleph0}
\LR{$\card{\Real} = \cardexpo{2}{\omega}$}
\end{thm}
\olpseganchor{OLP-0587-B020}{end}

\olpseganchor{OLP-0587-B021}{start}
\begin{proof}[د ثبوت چوکاټ]
د دې د ثبوت ډېرې لارې شته۔ تر ټولو نېغه لاره دا ده چې
\LR{$\cardle{\Pow{\omega}}{\Real}$} او
\LR{$\cardle{\Real}{\Pow{\omega}}$} وښيو، بيا د
شرودر--برنشتاين له قضيې څخه \LR{$\cardeq{\Real}{\Pow{\omega}}$}
واخلو، او له \olref[card-arithmetic][opps]{lem:SizePowerset2Exp}
څخه \LR{$\card{\Real} = \cardexpo{2}{\omega}$} ترلاسه کړو۔ د
\LR{$f \colon \Pow{\omega} \to \Real$} او
\LR{$g \colon \Real \to \Pow{\omega}$} غوندې يو پر يو تابعو
تعريف د تمرين په توګه پرېږدو۔ اشاره: حقيقي عدد د هغو
ناطقو عددونو په واسطه وپېژنئ چې تر هغه کم دي، او ناطق
عددونه په طبيعي عددونو ونوموئ۔ د بل لوري لپاره يو فرعي
سټ د صفرونو او دوو د درې‌بنسټه اعشاري لړۍ په توګه وليکئ:
د هر طبيعي شاخص رقم دوه دے که هغه شاخص غړے وي، او کنه
صفر دے۔ د لومړي بېل رقم ونډه د ټولو وروسته رقمونو له
ګډ توپير څخه زياته ده۔
\end{proof}
\olpseganchor{OLP-0587-B021}{end}

\olpseganchor{OLP-0587-B022}{start}
\begin{prob}
د \olref[sth][card-arithmetic][opps]{continuumis2aleph0}
ثبوت بشپړ کړئ؛ يعنې وښيئ چې \LR{$\cardle{\Pow{\omega}}{\Real}$}
او \LR{$\cardle{\Real}{\Pow{\omega}}$}۔
\end{prob}
\olpseganchor{OLP-0587-B022}{end}

\olpseganchor{OLP-0588-B002}{start}
\olfileid{sth}{card-arithmetic}{simp}
\olpseganchor{OLP-0588-B002}{end}

\olpseganchor{OLP-0588-B003}{start}
\olsection{د جمعې او ضرب ساده کول}
\olpseganchor{OLP-0588-B003}{end}

\olpseganchor{OLP-0588-B004}{start}
د اصلي عددونو نامتناهي جمع او ضرب په پايله کښې \emph{ډېر} اسانه
ثابتېږي۔ دا له دې راځي چې اصلي عددونه (ځينې) ترتيبي عددونه دي،
نو ښه ترتيب شوي دي، او له همدې امله په يوه ټاکلې توګه کار پرې
کېدے شي۔ خو د دې ښودل \emph{هغومره} اسانه نه دي۔ د پيل دپاره
يو څه ځيرک تعريف ته اړتيا لرو:
\olpseganchor{OLP-0588-B004}{end}

\olpseganchor{OLP-0588-B005}{start}
\begin{defn}
پر ترتيبي عددونو جوړو يو \emph{معياري ترتيب}، \LR{$\canonord$}،
داسې تعريفوو چې \LR{$\tuple{\alpha_1, \alpha_2} \canonord
\tuple{\beta_1, \beta_2}$} هغه وخت او يوازې هغه وخت وي چې له
لاندې حالتونو يو پېښ شي:
\begin{enumerate}
	\item \LR{$\max(\alpha_1, \alpha_2) < \max(\beta_1, \beta_2)$}؛ يا
	\item \LR{$\max(\alpha_1, \alpha_2) = \max(\beta_1, \beta_2)$} او
	\LR{$\alpha_1 < \beta_1$}؛ يا
	\item \LR{$\max(\alpha_1, \alpha_2) = \max(\beta_1, \beta_2)$} او
	\LR{$\alpha_1 = \beta_1$} او \LR{$\alpha_2 < \beta_2$}
\end{enumerate}
\end{defn}
\olpseganchor{OLP-0588-B005}{end}

\olpseganchor{OLP-0588-B006}{start}
\begin{lem}
\LR{$\tuple{\alpha \times \alpha, \canonord}$} د هر ترتيبي عدد
\LR{$\alpha$} دپاره ښه ترتيب دے۔
\end{lem}
\olpseganchor{OLP-0588-B006}{end}

\olpseganchor{OLP-0588-B007}{start}
\begin{proof}
\LR{$\canonord$} پر \LR{$\alpha \times \alpha$} د درې‌برخيزې کلي
له مخې يو کلک ترتيب دے: لومړے تر ټولو لوے مختص،
بيا لومړے مختص او بيا دويم مختص پرتله کېږي۔ د دغو
درې برخو کلکه قاموسي پرتله انتقالي او بې‌انعکاسه ده او
هرې دوو بېلابېلو جوړو ته ترتيب ورکوي۔ فرض کړئ نه \LR{$\tuple{\alpha_1,
\alpha_2}$} له \LR{$\tuple{\beta_1,
\beta_2}$} څخه د \LR{$\canonord$} له مخې کوچنۍ وي، نه دويمه له
لومړۍ څخه۔ بيا \LR{$\max(\alpha_1,
\alpha_2) = \max(\beta_1, \beta_2)$}، \LR{$\alpha_1 = \beta_1$}
او \LR{$\alpha_2 = \beta_2$}؛ نو \LR{$\tuple{\alpha_1, \alpha_2} =
\tuple{\beta_1,  \beta_2}$}۔
\olpseganchor{OLP-0588-B007}{end}

\olpseganchor{OLP-0588-B008}{start}
د ښه ترتيب د ښودلو دپاره دې \LR{$X \subseteq \alpha\times\alpha$}
ناتش وي۔ څرنګه چې \LR{$\alpha$} ترتيبي عدد دے، يو \LR{$\delta$} د
\LR{$\Setabs{\max(\gamma_1, \gamma_2)}{\tuple{\gamma_1,
\gamma_2} \in X}$} تر ټولو کوچنے غړے دے۔ اوس له
\LR{$\Setabs{\tuple{\gamma_1,\gamma_2} \in X}{\max(\gamma_1, \gamma_2) =
\delta}$} څخه ټولې جوړې لرې کړئ، پرته له هغو چې لومړے مختص يې
تر ټولو کوچنے وي؛ له پاتې جوړو څخه هغه چې دويم مختص يې تر
ټولو کوچنے وي، د \LR{$\canonord$} له مخې د \LR{$X$} تر ټولو کوچنے غړے دے۔
\end{proof}
\noindent
اوس يوې وړوکې، ساده کتنې ته راځو:
\olpseganchor{OLP-0588-B008}{end}

\olpseganchor{OLP-0588-B009}{start}
\begin{prop}\ollabel{simplecardproduct}
که \LR{$\cardeq{\alpha}{\beta}$} وي، نو \LR{$\cardeq{\alpha \times \alpha}{\beta
\times \beta}$}۔
\end{prop}
\olpseganchor{OLP-0588-B009}{end}

\olpseganchor{OLP-0588-B010}{start}
\begin{proof}
د هم‌شمېرۍ يوه دوه اړخيزه تابعه \LR{$f \colon \alpha \to \beta$}
واخلئ۔ هغه دا تابع ټاکي:
\LR{$\tuple{\gamma_1,
\gamma_2} \mapsto \tuple{f(\gamma_1), f(\gamma_2)}$}۔
په دواړو مختصونو د هغې معکوس وکاروئ، څو د دې نوې تابعې
معکوس ترلاسه شي۔
\end{proof}
\noindent
اوس دا ټول د يوه مهمې لمې په ثبوت کښې کاروو:
\begin{lem}\ollabel{alphatimesalpha}
\LR{$\cardeq{\alpha}{\alpha \times \alpha}$}، د هر نامتناهي ترتيبي
عدد \LR{$\alpha$} دپاره۔
\end{lem}
\olpseganchor{OLP-0588-B010}{end}

\olpseganchor{OLP-0588-B011}{start}
\begin{proof}
د تناقض دپاره دې \LR{$\alpha$} تر ټولو کوچنے نامتناهي ترتيبي عدد
وي چې دا ادعا يې په اړه ناسمه وي۔
\olref[sfr][siz][zigzag]{natsquaredenumerable} ښيي چې
\LR{$\cardeq{\omega}{\omega\times\omega}$}، نو \LR{$\omega \in \alpha$}۔
سربېره پر دې، \LR{$\alpha$} اصلي عدد دے: د تناقض دپاره فرض کړئ
چې نه دے؛ بيا \LR{$\card{\alpha} \in \alpha$}، او د
راکړل شوي نامتناهي عدد سره د هم‌شمېرۍ له مخې دا کوچنے
اصلي عدد هم نامتناهي دے۔ نو د فرضيې له مخې
\LR{$\cardeq{\card{\alpha}}{\card{\alpha} \times \card{\alpha}}$}؛
او د تعريف له مخې \LR{$\cardeq{\card{\alpha}}{\alpha}$}؛ نو له
\olref{simplecardproduct} څخه
\LR{$\cardeq{\alpha}{\alpha\times\alpha}$} راځي۔
\olpseganchor{OLP-0588-B011}{end}

\olpseganchor{OLP-0588-B012}{start}
اوس د هر \LR{$\tuple{\gamma_1, \gamma_2} \in \alpha \times \alpha$}
دپاره لاندې برخه وګورئ:
\begin{align*}
	\text{\RL{Seg}}(\gamma_1, \gamma_2) &= \Setabs{\tuple{\delta_1, \delta_2} \in \alpha \times \alpha}{\tuple{\delta_1, \delta_2} \canonord \tuple{\gamma_1, \gamma_2}}
\end{align*}
که \LR{$\gamma = \max(\gamma_1, \gamma_2)$} ونيسو، نو
\LR{$\tuple{\gamma_1, \gamma_2} \canonord \tuple{\gamma+1, \gamma + 1}$}۔
له ټاکلې جوړې مخکښې د هرې جوړې دواړه مختصونه تر لوی
مختص يا له هغه سره برابر دي۔ که دا مختص متناهي وي،
نو مخکنۍ برخه متناهي ده او اندازه يې تر اصلي نامتناهي
عدد کمه ده۔ نو کله چې \LR{$\gamma$} نامتناهي وي، لاندې ته پام وکړئ:
\begin{align*}
	\text{\RL{Seg}}(\gamma_1, \gamma_2) &
	\precsim ((\gamma \ordplus 1)\ordtimes (\gamma \ordplus 1))\\
	&\approx (\gamma \ordtimes \gamma)
	\text{\RL{، د \olref[ord-arithmetic][using-addition]{ordinfinitycharacter} او
	\olref{simplecardproduct} له مخې}}\\
	&\approx \gamma \text{\RL{، د استقرا د فرضيې له مخې}}\\
	& \prec \alpha\text{\RL{، ځکه \LR{$\alpha$} اصلي عدد دے}}
\end{align*}
د ترتيب د ډول هره خاصه پيلنۍ برخه د يوې ټاکلې جوړې
مخکنۍ برخه ده۔ که ډول تر اصلي عدد لوے وے، نو خپله
اصلي عدد به يوه داسې پيلنۍ برخه وه چې اندازه يې تر
ځان کمه وي۔ نو \LR{$\ordtype{\alpha\times \alpha, \canonord} \leq \alpha$}،
او له دې \LR{$\cardle{\alpha \times \alpha}{\alpha}$}۔ څرنګه چې
طبعاً \LR{$\cardle{\alpha}{\alpha \times \alpha}$}، پايله د
شرودر--برنشتاين له قضيې راځي۔
\end{proof}
\olpseganchor{OLP-0588-B012}{end}

\olpseganchor{OLP-0588-B013}{start}
بالاخره د ساده کولو پايلې ته رسېږو:
\olpseganchor{OLP-0588-B013}{end}

\olpseganchor{OLP-0588-B014}{start}
\begin{thm}\ollabel{cardplustimesmax}
که \LR{$\cardfont{a}, \cardfont{b}$} نامتناهي اصلي عددونه وي، نو:
\[
	\cardfont{a}
\cardtimes \cardfont{b} = \cardfont{a} \cardplus \cardfont{b} =
\text{\RL{max}}(\cardfont{a}, \cardfont{b}).
\]
\end{thm}
\olpseganchor{OLP-0588-B014}{end}

\olpseganchor{OLP-0588-B015}{start}
\begin{proof}
له عموميت له کمېدو پرته فرض کړئ \LR{$\cardfont{a} = \max(\cardfont{a},
\cardfont{b})$}۔ بيا د \olref{alphatimesalpha} په کارولو سره
\LR{$\cardfont{a}\cardtimes\cardfont{a} = \cardfont{a} \leq \cardfont{a}
\cardplus \cardfont{b} \leq \cardfont{a} \cardplus \cardfont{a} \leq
\cardfont{a} \cardtimes \cardfont{a}$}۔ د ضرب د
نابرابرۍ لپاره له دويم نامتناهي عامل څخه يو ثابت غړے
وټاکئ؛ په دغه مختص لومړے عامل په حاصل ضرب کښې ورګډېږي۔
بل لور د کوچني عامل د ورګډولو تابع ورکوي، ځکه د دواړو
سترو عاملونو حاصل ضرب د لومړي عامل هومره غړي لري۔
\end{proof}\noindent په ورته ډول،
که \LR{$\cardfont{a}$} نامتناهي وي، د \LR{$\cardfont{a}$} په اندازه د
داسې سټونو اتحاد چې د هر يوه اندازه \LR{$\leq\cardfont{a}$} وي،
اندازه يې هم \LR{$\leq\cardfont{a}$} وي:
\olpseganchor{OLP-0588-B015}{end}

\olpseganchor{OLP-0588-B016}{start}
\begin{prop}\ollabel{kappaunionkappasize}
دې \LR{$\cardfont{a}$} يو نامتناهي اصلي عدد وي۔ د هر ترتيبي عدد
\LR{$\beta \in \cardfont{a}$} دپاره دې \LR{$X_\beta$} داسې سټ وي چې
\LR{$\card{X_\beta} \leq \cardfont{a}$}۔ بيا
\LR{$\card{\bigcup_{\beta \in \cardfont{a}} X_\beta}
\leq \cardfont{a}$}۔
\end{prop}
\olpseganchor{OLP-0588-B016}{end}

\olpseganchor{OLP-0588-B017}{start}
\begin{proof}
د هر \LR{$\beta \in \cardfont{a}$} دپاره يوه يو پر يو تابع
\LR{$f_\beta \colon X_\beta \to \cardfont{a}$} وټاکئ۔\footnote{دا تابعې
څنګه يوځاے «ټاکل» کېږي؟ \olref[sth][choice][countablechoice]{sec}
وګورئ۔} يوه يو پر يو تابع \LR{$g \colon
\bigcup_{\beta \in \cardfont{a}} X_\beta \to \cardfont{a} \times
\cardfont{a}$} په \LR{$g(v) = \tuple{\beta, f_\beta(v)}$} سره تعريف
کړئ، چې دلته \LR{$v \in X_\beta$}، او \LR{$v \notin X_\gamma$} د هر
\LR{$\gamma \in \beta$} دپاره۔ اوس
\LR{$\bigcup_{\beta \in \cardfont{a}} X_\beta \preceq \cardfont{a} \times
\cardfont{a} \approx \cardfont{a}$} د \olref{cardplustimesmax}
له مخې۔
\end{proof}
\olpseganchor{OLP-0588-B017}{end}

\olpseganchor{OLP-0589-B002}{start}
\olfileid{sth}{card-arithmetic}{expotough}
\olpseganchor{OLP-0589-B002}{end}

\olpseganchor{OLP-0589-B003}{start}
\olsection[ځينې ساده کول]{د اصلي توان ځينې ساده کول}
\olpseganchor{OLP-0589-B003}{end}

\olpseganchor{OLP-0589-B004}{start}
که څه هم د \LR{$\canonord$} تعريف لږ پېچلے و، پايله يې د اصلي جمعې
او ضرب په اړه ګټوره ده، \olref[simp]{cardplustimesmax}۔ خو
نامتناهي توان په دومره نېغه توګه نه ساده کېږي۔ د دې د علت د
روښانولو دپاره له يوې پايلې پيل کوو چې د متناهي حالت يوه اشنا
بڼه غځوي (که څه هم ثبوت يې په لوړه انتزاعي کچه دے):
\olpseganchor{OLP-0589-B004}{end}

\olpseganchor{OLP-0589-B005}{start}
\begin{prop}\ollabel{simplecardexpo}
\LR{$\cardexpo{\cardfont{a}}{\cardfont{b} \cardplus \cardfont{c}} =
\cardexpo{\cardfont{a}}{\cardfont{b}} \cardtimes
\cardexpo{\cardfont{a}}{\cardfont{c}}$} او
\LR{$\cardexpo{(\cardexpo{\cardfont{a}}{\cardfont{b}})}{\cardfont{c}} =
\cardexpo{\cardfont{a}}{\cardfont{b} \cardtimes \cardfont{c}}$}، د
هر درېيو اصلي عددونو \LR{$\cardfont{a}, \cardfont{b}, \cardfont{c}$} دپاره۔
\end{prop}
\olpseganchor{OLP-0589-B005}{end}

\olpseganchor{OLP-0589-B006}{start}
\begin{proof}
د لومړۍ ادعا دپاره يوه تابع \LR{$f \colon
(\cardfont{b}\disjointsum\cardfont{c}) \to \cardfont{a}$} وګورئ۔ اوس
دا «ووېشئ»: \LR{$f_\cardfont{b}(\beta) = f(\beta, 0)$} د هر
\LR{$\beta \in \cardfont{b}$} دپاره، او \LR{$f_\cardfont{c}(\gamma) = f(\gamma, 1)$}
د هر \LR{$\gamma \in \cardfont{c}$} دپاره تعريف کړئ۔ نقشه \LR{$f \mapsto
\tuple{f_{\cardfont{b}}, f_\cardfont{c}}$} يوه دوه اړخيزه يو پر يو تابع
\LR{$\funfromto{\cardfont{b} \disjointsum \cardfont{c}}{\cardfont{a}} \to
(\funfromto{\cardfont{b}}{\cardfont{a}} \times
\funfromto{\cardfont{c}}{\cardfont{a}})$} ده۔
\olpseganchor{OLP-0589-B006}{end}

\olpseganchor{OLP-0589-B007}{start}
د دويمې ادعا دپاره يوه تابع \LR{$f \colon \cardfont{c} \to
(\funfromto{\cardfont{b}}{\cardfont{a}})$} وګورئ؛ نو د هر
\LR{$\gamma \in
\cardfont{c}$} دپاره يوه تابع \LR{$f(\gamma) \colon \cardfont{b} \to
\cardfont{a}$} لرو۔ اوس
\LR{$f^*(\beta, \gamma) = (f(\gamma))(\beta)$} د هر
\LR{$\tuple{\beta, \gamma} \in \cardfont{b} \times \cardfont{c}$}
دپاره تعريف کړئ۔ نقشه
\LR{$f \mapsto f^*$} يوه دوه اړخيزه يو پر يو تابع
\LR{$\funfromto{\cardfont{c}}{(\funfromto{\cardfont{b}}{\cardfont{a}})}
\to \funfromto{\cardfont{b} \times \cardfont{c}}{\cardfont{a}}$} ده۔
دواړه لوري د مختصونو د يوځاے کولو او بېلولو له مخې يو د بل
معکوس دي۔ د حاصل ضرب اصلي عدد د تعريف له مخې د همدغه
کارتيزي حاصل ضرب اندازه ده۔
\end{proof}
\olpseganchor{OLP-0589-B007}{end}

\olpseganchor{OLP-0589-B008}{start}
اوس زموږ \emph{غوښتنه} دا ده چې د نامتناهي اصلي عددونو په حالت
کښې \LR{$\cardexpo{\cardfont{a}}{\cardfont{b}}$} په اسانه محاسبه کړو۔
لاندې پايله په دې لوري يو ښه ګام دے:
\olpseganchor{OLP-0589-B008}{end}

\olpseganchor{OLP-0589-B009}{start}
\begin{prop}\ollabel{cardexpo2reduct}
که \LR{$2 \leq \cardfont{a} \leq \cardfont{b}$} وي او \LR{$\cardfont{b}$}
نامتناهي وي، نو \LR{$\cardexpo{\cardfont{a}}{\cardfont{b}} =
\cardexpo{2}{\cardfont{b}}$}۔
\end{prop}
\olpseganchor{OLP-0589-B009}{end}

\olpseganchor{OLP-0589-B010}{start}
\begin{proof}
\begin{align*}
	\cardexpo{2}{\cardfont{b}} &\leq
	\cardexpo{\cardfont{a}}{\cardfont{b}}\text{\RL{، ځکه \LR{$2 \leq \cardfont{a}$}}}\\
	&\leq \cardexpo{(2^\cardfont{a})}{\cardfont{b}}
	\text{\RL{، د \olref[opps]{lem:SizePowerset2Exp} له مخې}}\\
	&= \cardexpo{2}{\cardfont{a} \cardtimes \cardfont{b}}
	\text{\RL{، د \olref{simplecardexpo} له مخې}} \\
	&= \cardexpo{2}{\cardfont{b}}
	\text{\RL{، د \olref[simp]{cardplustimesmax} له مخې}}
\end{align*}
که بنسټ متناهي وي، هماغه وروستۍ برابري بيا هم سمه ده:
دوه يا زيات، خو متناهي، کاپۍ د نامتناهي ښي عامل په اندازه دي۔
له يوې خوا يوه کاپي ورګډېږي؛ له بلې خوا ټولې کاپۍ د نامتناهي
عامل د مربع فرعي سټ جوړوي۔
\end{proof}
\olpseganchor{OLP-0589-B010}{end}

\olpseganchor{OLP-0589-B011}{start}
له دې څخه نور ساده کول په رښتيا تمه نه شو کولے، ځکه
\LR{$\cardfont{b} < \cardexpo{2}{\cardfont{b}}$} د کانتور د
قضيې له مخې، د \olref[card-arithmetic][opps]{lem:SizePowerset2Exp}
د تواني سټ د اندازې له برابرۍ سره يوځاے۔
خو دا موږ ته د \LR{$\cardexpo{\cardfont{a}}{\cardfont{b}}$} په اړه
نه وايي چې کله \LR{$\cardfont{b} < \cardfont{a}$} وي، څه ووايو۔
البته که \LR{$\cardfont{b}$} \emph{متناهي} وي،
لار يې پېژنو۔
\olpseganchor{OLP-0589-B011}{end}

\olpseganchor{OLP-0589-B012}{start}
\begin{prop}
که \LR{$\cardfont{a}$} نامتناهي وي او \LR{$n \in \omega$} وي، او نوموړے
طبيعي عدد صفر نه وي، نو \LR{$\cardexpo{\cardfont{a}}{n} = \cardfont{a}$}۔
د صفر توان ارزښت يو دے، ځکه پر تشه ساحه يوازې تشه
تابعه شته۔
\end{prop}
\olpseganchor{OLP-0589-B012}{end}

\olpseganchor{OLP-0589-B013}{start}
\begin{proof}
د ناصفر متناهي توان دپاره
\LR{$\cardexpo{\cardfont{a}}{n} = \cardfont{a} \cardtimes \cardfont{a}
\cardtimes \ldots \cardtimes \cardfont{a} = \cardfont{a}$}، د
\olref[simp]{cardplustimesmax} له مخې۔
\end{proof}
\noindent
پر دې سربېره، په ځينو نورو حالتونو کښې هم د
\LR{$\cardexpo{\cardfont{a}}{\cardfont{b}}$} اندازه ټاکلے شو:
\olpseganchor{OLP-0589-B013}{end}

\olpseganchor{OLP-0589-B014}{start}
\begin{prop}
که \LR{$2 \leq \cardfont{b} < \cardfont{a} \leq
\cardexpo{2}{\cardfont{b}}$} وي او \LR{$\cardfont{b}$} نامتناهي وي، نو
\LR{$\cardexpo{\cardfont{a}}{\cardfont{b}} = \cardexpo{2}{\cardfont{b}}$}۔
\end{prop}
\olpseganchor{OLP-0589-B014}{end}

\olpseganchor{OLP-0589-B015}{start}
\begin{proof}
\LR{$\cardexpo{2}{\cardfont{b}}\leq \cardexpo{\cardfont{a}}{\cardfont{b}}
\leq \cardexpo{(\cardexpo{2}{\cardfont{b}})}{\cardfont{b}} =
\cardexpo{2}{\cardfont{b}\cardtimes\cardfont{b}} =
\cardexpo{2}{\cardfont{b}}$}، لکه په \olref{cardexpo2reduct} کښې
استدلال وشو۔
\end{proof}
\noindent
خو له دې ځايه وروسته مسئلې ډېرې نازکې کېږي۔
\olpseganchor{OLP-0589-B015}{end}

\olpseganchor{OLP-0590-B003}{start}
\olfileid{sth}{card-arithmetic}{ch}
\olpseganchor{OLP-0590-B003}{end}

\olpseganchor{OLP-0590-B004}{start}
\olsection{د پيوستار فرضيه}
\olpseganchor{OLP-0590-B004}{end}

\olpseganchor{OLP-0590-B005}{start}
تېرې پايلې (په سمه توګه) دې ته اشاره وکړه چې د اصلي عددونو توان
به ډېر \emph{اسان} و، که نامتناهي اصلي عددونه د \LR{$2$} له توانونو سره
«په نېغه توګه» چلند وکړي؛ يعنې، د \olref[opps]{lem:SizePowerset2Exp}
له مخې، د توان سټونو له اخيستلو سره۔ خو دا نه شو فرضولے چې
نامتناهي اصلي عددونه په رښتيا له توان سټونو سره داسې نېغ چلند کوي۔
\olpseganchor{OLP-0590-B005}{end}

\olpseganchor{OLP-0590-B006}{start}
د دې خبرې د سپړلو د پيل دپاره څو ګټورې نښې راوړو۔
\olpseganchor{OLP-0590-B006}{end}

\olpseganchor{OLP-0590-B007}{start}
\begin{defn}
چې \LR{$ \cardsucc{\cardfont{a}} $} له \LR{$\cardfont{a}$} څخه په کلکه لوے
تر ټولو کوچنے اصلي عدد دے، لاندې دوه نامتناهي لړۍ تعريفوو:
\begin{align*}
	\aleph_{0} &\defis \omega & 		
	\beth_{0} &\defis \omega\\
	\aleph_{\alpha \ordplus 1} &\defis \cardsucc{(\aleph_{\alpha})} &
	\beth_{\alpha+1} &\defis \cardexpo{2}{\beth_{\alpha}}\\
	\aleph_{\alpha} &\defis \bigcup_{\beta< \alpha} \aleph_{\beta} &
	\beth_{\alpha} &\defis \bigcup_{\beta < \alpha}\beth_{\beta} & \text{\RL{چې \LR{$\alpha$} حدي ترتيبي عدد وي}}.
	\end{align*}
\end{defn}
\olpseganchor{OLP-0590-B007}{end}

\olpseganchor{OLP-0590-B008}{start}
د \LR{$\cardsucc{\cardfont{a}}$} تعريف پر ځاے دے، ځکه
\olref[cardinals][classing]{lem:NoLargestCardinal} وايي چې د هر اصلي
عدد \LR{$\cardfont{a}$} دپاره له \LR{$\cardfont{a}$} څخه لوے يو اصلي عدد
شته، او د نامتناهيت هاخوا
استقرا دا تضمينوي چې له \LR{$\cardfont{a}$} څخه لوے يو \emph{تر ټولو کوچنے}
اصلي عدد شته۔ د دغو لړيو پاتې تعريفونه د نامتناهيت هاخوا بازګښت
له لارې ورکول کېږي۔
\olpseganchor{OLP-0590-B008}{end}

\olpseganchor{OLP-0590-B009}{start}
کانتور د «\LR{$\aleph$}» نښه راوړه؛ دې ته \emph{اَلېف} وايي۔ دا د
عبراني الفبا لومړے تورى او د «نامتناهي» د عبراني کلمې لومړے
تورے دے۔ پيرس د «\LR{$\beth$}» نښه راوړه؛ دې ته \emph{بېت} وايي،
چې د عبراني الفبا دويم تورى دے۔\footnote{پيرس دا نښه د ۱۹۰۰م کال
د دسمبر په يوه ليک کښې کانتور ته کارولې وه۔ خو په هماغه ليک
کښې ئې د دې دپاره ناسم استدلال هم وکړ چې \LR{$\beth_\alpha$} د
\LR{$\alpha \geq \omega$} دپاره شتون نه لري۔}
اوس دا نښې موږ ته نامتناهي اصلي عددونه راکوي۔
\olpseganchor{OLP-0590-B009}{end}

\olpseganchor{OLP-0590-B010}{start}
\begin{prop}
\LR{$\aleph_\alpha$} او \LR{$\beth_\alpha$} د هر ترتيبي عدد \LR{$\alpha$}
دپاره اصلي عددونه دي۔
\end{prop}
\olpseganchor{OLP-0590-B010}{end}

\olpseganchor{OLP-0590-B011}{start}
\begin{proof}
دواړه پايلې د نامتناهيت هاخوا په ساده استقرا ثابتېږي۔
\LR{$\aleph_0 = \beth_0 = \omega$} د
\olref[cardinals][classing]{omegaisacardinal} له مخې اصلي عدد دے۔
که فرض کړو چې \LR{$\aleph_\alpha$} او \LR{$\beth_\alpha$} دواړه اصلي
عددونه دي، نو \LR{$\aleph_{\alpha+1}$} او
\LR{$\beth_{\alpha+1}$} په څرګند ډول د اصلي عددونو په توګه تعريف
شوي دي۔ او د اصلي عددونو د يوه سټ اتحاد د
\olref[cardinals][classing]{unioncardinalscardinal} له مخې اصلي
عدد دے۔
\end{proof}
\noindent
سربېره پر دې، هر نامتناهي اصلي عدد يو \LR{$\aleph$} دے:
\olpseganchor{OLP-0590-B011}{end}

\olpseganchor{OLP-0590-B012}{start}
\begin{prop}
که \LR{$\cardfont{a}$} نامتناهي اصلي عدد وي، نو
\LR{$\cardfont{a} = \aleph_\gamma$}، د يوه يوازيني \LR{$\gamma$} دپاره۔
\end{prop}
\olpseganchor{OLP-0590-B012}{end}

\olpseganchor{OLP-0590-B013}{start}
\begin{proof}
پر نامتناهي اصلي عددونو د نامتناهيت هاخوا استقرا کوو۔
لومړنے حالت \LR{$\omega=\aleph_0$} دے۔ د استقرا دپاره فرض
کړئ چې که \LR{$\omega\leq\cardfont{b} < \cardfont{a}$} وي، نو
\LR{$\cardfont{b} = \aleph_{\gamma_\cardfont{b}}$}۔ که
\LR{$\cardfont{a} = \cardsucc{\cardfont{b}}$} د کوم نامتناهي
\LR{$\cardfont{b}$} دپاره وي، نو \LR{$\cardfont{a} =
\cardsucc{(\aleph_{\gamma_\cardfont{b}})}=
\aleph_{\gamma_\cardfont{b}+1}$}۔ په پاتې حالت کښې د کوچنيو
نامتناهي اصلي عددونو کورنۍ ناتشه ده او تر ټولو لوے غړے نۀ لري۔
اتحاد يې پخپله اصلي عدد دے۔ که دا اتحاد تر راکړل شوي اصلي عدد
کوچنے وے، نو د کورنۍ تر ټولو لوے غړے به وے۔ ځکه
\LR{$\cardfont{a} = \bigcup_{\omega\leq\cardfont{b}<\cardfont{a}} \cardfont{b} =
\bigcup_{\omega\leq\cardfont{b}<\cardfont{a}}{\aleph_{\gamma_\cardfont{b}}}$}۔
د اَلېف لړۍ په کلکه زياتېږي: هر تالی تر خپل مخکني لوے دے
او حدي اتحاد له هر مخکني تالي څخه لږ نۀ دے۔ نو د کوچنيو
اصلي عددونو اړوند شاخصونه هم تر ټولو لوے غړے نۀ لري۔ تعويض
د هغوی سټ راکوي، او اتحاد يې يو حدي شاخص دے۔ دغه شاخصونه
تر خپل اتحاد پورې هم‌پايه دي؛ د اَلېف د حدي معادلې په حساب
کښې د همدې کورنۍ اتحاد بس دے۔ نو
\LR{$\cardfont{a} = \aleph_\gamma$}، چېرته چې \LR{$\gamma =
\bigcup_{\omega\leq\cardfont{b}<\cardfont{a}}\gamma_\cardfont{b}$}۔
د لړۍ کلک زياتوالے د شاخص يوازينيت هم ثابتوي۔
\end{proof}
\olpseganchor{OLP-0590-B013}{end}

\olpseganchor{OLP-0590-B014}{start}
څرنګه چې هر نامتناهي اصلي عدد يو \LR{$\aleph$} دے، دا پوښتنه راولاړېږي:
ايا هر نامتناهي اصلي عدد يو \LR{$\beth$} هم دے؟ که داسې \emph{وے}،
نو نامتناهي اصلي عددونو به د توان سټونو د اخيستلو له عمل سره
«په نېغه توګه» چلند کاوه۔ په رښتيا، بيا به مو لاندې ادعا درلوده:
\olpseganchor{OLP-0590-B014}{end}

\olpseganchor{OLP-0590-B015}{start}
\begin{defish}
\emph{د پيوستار عمومي فرضيه} (GCH)۔
\LR{$\aleph_\alpha  = \beth_\alpha$}، د هر \LR{$\alpha$} دپاره۔
\end{defish}
\olpseganchor{OLP-0590-B015}{end}

\olpseganchor{OLP-0590-B016}{start}
سربېره پر دې، که GCH سمه وے، نو د اصلي عددونو توان مو د پام وړ
ساده کولے شو۔ په ځانګړې توګه، ثابتولے مو شول چې کله
\LR{$0<\cardfont{b} < \cardfont{a}$} وي، د
\LR{$\cardexpo{\cardfont{a}}{\cardfont{b}}$} اندازه د
\LR{$\cardfont{a}\leq \cardexpo{\cardfont{a}}{\cardfont{b}} \leq
\cardsucc{\cardfont{a}}$} تر منځ راځي۔ بيا مو دقيق شرطونه هم
ورکولے شول چې له دغو دوو امکانونو کوم يو پېښېږي، يعنې
\LR{$\cardfont{a} = \cardexpo{\cardfont{a}}{\cardfont{b}}$}
او که \LR{$\cardexpo{\cardfont{a}}{\cardfont{b}} =
\cardsucc{\cardfont{a}}$}۔\footnote{دا شرط د
\emph{هم‌پايۍ} له مخې ټاکل کېږي۔}
\olpseganchor{OLP-0590-B016}{end}

\olpseganchor{OLP-0590-B017}{start}
خو GCH يوه \emph{فرضيه} ده، \emph{قضيه} نه ده۔ په رښتيا،
\mbox{\LR{\latinfont \citet{Godel1938}}} ثابته کړه چې که \LR{$\ZFC$} سازګاره وي، نو
\LR{$\ZFC + \text{\RL{GCH}}$} هم سازګاره ده۔ خو وروسته څرګنده شوه
چې \LR{$\lnot$}GCH هم له \LR{$\ZFC$} سره زياتولے شو۔ د GCH تر ټولو
ساده غير بديهي \emph{ځانګړے حالت} دا دے:
\olpseganchor{OLP-0590-B017}{end}

\olpseganchor{OLP-0590-B018}{start}
\begin{defish}
\emph{د پيوستار فرضيه} (CH)۔ \LR{$\aleph_1 = \beth_1$}۔
\end{defish}
\olpseganchor{OLP-0590-B018}{end}

\olpseganchor{OLP-0590-B019}{start}
\mbox{\LR{\latinfont \citet{Cohen1963}}} ثابته کړه چې که \LR{$\ZFC$} سازګاره وي، نو
\LR{$\ZFC + \lnot\text{\RL{CH}}$} هم سازګاره ده۔ نو د پيوستار فرضيه
له \LR{$\ZFC$} څخه خپلواکه ده۔
\olpseganchor{OLP-0590-B019}{end}

\olpseganchor{OLP-0590-B020}{start}
دې ته د پيوستار فرضيه ځکه ويل کېږي چې «پيوستار» د حقيقي عددونو
د کرښې، \LR{$\Real$}، بل نوم دے۔
\olref[opps]{continuumis2aleph0} وايي چې \LR{$\card{\Real} =
\beth_1$}۔ نو د پيوستار فرضيه وايي چې د طبيعي عددونو د سټ د
اصلي شمېر، \LR{$\aleph_0 = \beth_0$}، او د پيوستار د اصلي شمېر،
\LR{$\beth_1$}، تر منځ بل اصلي عدد نشته۔
\olpseganchor{OLP-0590-B020}{end}

\olpseganchor{OLP-0590-B021}{start}
له \LR{$\ZFC$} څخه د (G)CH د \emph{خپلواکۍ} له مخې، د دې د
\emph{رښتياوالي} په باب څه ووايو؟ په دې باب \emph{ډېرې}
خبرې شته۔ د سټونو په نظريه کښې د وروستيو څېړنو ډېره ګټوره
برخه هم د دغو پوښتنو څېړلو ته متوجه ده۔ خو دوه لنډ ټکي
خامخا د يادولو وړ دي۔
\olpseganchor{OLP-0590-B021}{end}

\olpseganchor{OLP-0590-B022}{start}
لومړے: له دغو رسمي خپلواکۍ پايلو څخه دا \emph{سمدستي}
نه راځي چې د GCH يا CH د رښتيا ارزښت \emph{ناتاکلے}
دے۔ کېدای شي يوازې نورو داسې اصلونو ته اړتيا ولرو چې موږ
ته طبيعي ښکاري او پوښتنه په يوه يا بله خوا هواره کړي۔
ګېډل پخپله همدا غبرګون مناسب ګاڼه۔
\olpseganchor{OLP-0590-B022}{end}

\olpseganchor{OLP-0590-B023}{start}
دويم: له \LR{$\ZFC$} څخه د CH خپلواکي بېشکه \emph{حيرانوونکې}
ده، خو په لفظي معنا \emph{ناباورېدونکې} نه ده۔ خبره دا ده چې
\LR{$\ZFC$}، تر کومه ځايه چې موږ ته وايي، د يوه اصلي عدد څخه د
هغه تالي ته تګ ښايي د توان سټونو له اخيستلو څخه ډېر نازک
عمل وغواړي۔
 %The operation of taking powersets moves us from one stage of the
 %hierarchy to its successor stage (from $V_{\alpha}$ to
 %$V_{\alpha+1}$). 
\olpseganchor{OLP-0590-B023}{end}

\olpseganchor{OLP-0590-B024}{start}
له دې دوو خبرو وروسته، که نور معلومات غواړئ، اوس به د دې
بحث د بېلابېلو دريځونو لرونکو فيلسوفانو او رياضي‌پوهانو
ليکنې ولولئ۔\footnote{کېدای شي لوستل د \mbox{\LR{\latinfont \citet[\S15.6]{Potter2004}}}
څخه پيل کړئ۔}
\olpseganchor{OLP-0590-B024}{end}

\olpseganchor{OLP-0591-B003}{start}
\olfileid{sth}{card-arithmetic}{fix}
\olsection{د \texorpdfstring{\LR{$\aleph$}}{ℵ} ثابت ټکي}
\olpseganchor{OLP-0591-B003}{end}

\olpseganchor{OLP-0591-B004}{start}
په \olref[spine][]{chap} کښې مو ويلي وو چې د تعويض اصل توجيه
غواړي، ځکه دا د سټونو مراتبي جوړښت ډېر لوړوي۔ اوس چې مو د اصلي
عددونو يو څه حساب وکړ، د دې جوړښت د لوړوالي يوه بېلګه وړاندې
کولے شو۔
\olpseganchor{OLP-0591-B004}{end}

\olpseganchor{OLP-0591-B005}{start}
ښکاره ده چې \LR{$0 < \aleph_0$}، او \LR{$1 < \aleph_1$}، او
\LR{$2 < \aleph_2$}\ldots؛ په رښتيا، د اندازې توپير په هر پړاو کښې
لا \emph{لوېږي}۔ نو دا ګومان راولاړېږي چې
\LR{$\kappa< \aleph_\kappa$} به د هر ترتيبي عدد \LR{$\kappa$} دپاره وي۔
\olpseganchor{OLP-0591-B005}{end}

\olpseganchor{OLP-0591-B006}{start}
خو د \LR{$\ZFC$} په منلو دا ګومان \emph{ناسم} دے۔ په رښتيا،
\emph{د \LR{$\aleph$} ثابت ټکي} هم شته؛ يعنې داسې اصلي عددونه
\LR{$\kappa$} چې \LR{$\kappa=\aleph_\kappa$} وي۔
\olpseganchor{OLP-0591-B006}{end}

\olpseganchor{OLP-0591-B007}{start}
\begin{prop}\ollabel{alephfixed}
د \LR{$\aleph$} يو ثابت ټکی شته۔
\end{prop}
\olpseganchor{OLP-0591-B007}{end}

\olpseganchor{OLP-0591-B008}{start}
\begin{proof}
د بازګښت په کارولو لاندې تعريف کړئ:
\begin{align*}
	\kappa_0 &= 0\\
	\kappa_{n+1} &= \aleph_{\kappa_n}\\
	\kappa&= \bigcup_{n < \omega}\kappa_n
\end{align*}
اوس \LR{$\kappa$} د
\olref[cardinals][classing]{unioncardinalscardinal} له مخې
اصلي عدد دے۔ لړۍ په کلکه زياتېږي: لومړے ګام له صفر
څخه اومېګا ته ځي، او هر بل ګام د اَلېف د کلک زياتوالي له
مخې راځي۔ اتحاد يې يو حدي ترتيبي عدد دے، او د لړۍ غړي
پکښې هم‌پايه دي۔ د اَلېف د حدي معادلې له مخې:
\[
	\kappa= \bigcup_{n < \omega} \kappa_{n+1} = 
	\bigcup_{n < \omega}\aleph_{\kappa_n} = 
	\bigcup_{\alpha < \kappa}\aleph_\alpha = \aleph_\kappa
\]
%	By construction, $\kappa$ is the least cardinal greater than each
%	$\kappa_n$. So $\aleph_\kappa$ is the least cardinal greater than
%	each $\aleph_{\kappa_n}$, and hence greater than each
%	$\kappa_{n+1}$. But equally $\kappa$ is the least cardinal greater
%	than each $\kappa_{n+1} = \aleph_{\kappa_n}$. So $\kappa=
%	\aleph_\kappa$.
\end{proof}
\olpseganchor{OLP-0591-B008}{end}

\olpseganchor{OLP-0591-B009}{start}
بولوس يو وخت د همدغې، زموږ له خوا جوړې شوې، د \LR{$\aleph$}
د ثابت ټکي په اړه يې وليکل۔ د بحث په پيل کښې د \LR{$\kappa$}
له شتون څخه راولاړېدونکې دعوه څېړي۔
\begin{explain}
  د \LR{$\kappa$} او د سټ تيورۍ د دوديزو لويو اصلي عددونو د پرتلې په
  وخت کښې دا ثابت ټکی وړوکے ښکاري۔ بولوس بيا پوښتي چې له
  دغې تيورۍ بهر لوستونکي ته د لږ تر لږه \LR{$\kappa$} څيزونو
  د شتون د منلو دپاره کوم دليل شته۔ \mbox{\LR{\latinfont \citep[p.~257]{Boolos2000}}}
\end{explain}
د ليکنې په پاے کښې دا پوښتنه بيا رااخلي۔
\begin{explain}
  بولوس پوښتي چې ايا د سټونو له لومړني تصور څخه تر دې
  لرې پايلې پورې د شتون د دعوې توجيه ورسره پاتې کېږي، که
  يوازې د ټاکل شويو اصولو صوري پايلې راباسو۔ \mbox{\LR{\latinfont \citep[p.~268]{Boolos2000}}}
\end{explain}
که د شتون دغه دعوه کافي توجيه ونۀ لري، نو د تعويض د اصل
رول بايد وڅېړو۔ ځکه همدا اصل زموږ ثبوت ته
لار برابروي۔ که داسې وي، نو ګومان کېږي چې بولوس به اړ شي
خپله هغه څو لسيزې پخوانۍ ادعا بيا وګوري چې تعويض «هېڅ
ناغوښتې» پايله نه لري (وګورئ \olref[replacement][extrinsic]{sec})۔
\olpseganchor{OLP-0591-B009}{end}

\olpseganchor{OLP-0591-B010}{start}
خو ايا د \LR{$\kappa$} شتون دومره بد دے؟ دلته ښايي د رسل
\emph{د ټرېسټرم شېنډي تناقض} په پام کښې نيول مرسته وکړي۔
ټرېسټرم شېنډي خپل ژوند په ورځپاڼه کښې ثبتوي، خو د يوې
ورځې د ثبتولو دپاره يو کال وخت اخلي۔ د هر کال په تېرېدو
لا شاته پاتې کېږي: له يو کال وروسته يې يوازې يوه ورځ
ثبت کړې او ۳۶۴ تېرې ورځې يې نه دي ثبت کړې؛ له دوو کلونو
وروسته يې يوازې دوې ورځې ثبت کړې او ۷۲۸ ورځې نه دي ثبت
کړې؛ له درېيو کلونو وروسته يې يوازې درې ورځې ثبت کړې
او ۱۰۹۲ ورځې نه دي ثبت کړې \dots\footnote{کبيسې کلونه
له پامه غورځوو۔} بيا هم، که ټرېسټرم \emph{تلپاتې} وي،
هره ورځ به ثبت کړے شي، ځکه د خپل ژوند \LR{$n$}مه ورځ
به په \LR{$n$}م کال کښې ثبت کړي۔ په دې ډول به «د وخت په پاے
کښې» د هغه ورځپاڼه بشپړه وي۔
\olpseganchor{OLP-0591-B010}{end}

\olpseganchor{OLP-0591-B011}{start}
اوس ولې دا له دې فکره دومره جلا ده چې \LR{$\alpha$} له
\LR{$\aleph_\alpha$} څخه کوچنے، بلکې په زياتېدونکي او ډېر
ښکاره ډول کوچنے، وي، تر هغې چې \LR{$\kappa$} ته ورسېږو،
او هلته دواړه سره برابر شي او \LR{$\kappa
= \aleph_\kappa$} وګرځي؟
\olpseganchor{OLP-0591-B011}{end}

\olpseganchor{OLP-0591-B012}{start}
دا خبره يوې خوا ته پرېږدو، او که \LR{$\ZFC$} ومنو، نو د
ثابت ټکي د جوړښتونو په اړه له يوې بلې په زړه پورې
خبرې سره پای ته رسېږو۔ راتلونکې درې پايلې په وجداني
ډول دا ثابتوي چې يو (ناابتدايي) پړاو شته چې د مراتبي
جوړښت پلنوالے يې له لوړوالي سره برابر دے:
\olpseganchor{OLP-0591-B012}{end}

\olpseganchor{OLP-0591-B013}{start}
\begin{prop}\ollabel{bethfixed}
د \LR{$\beth$} يو ثابت ټکی شته، يعنې داسې \LR{$\kappa$} چې
\LR{$\kappa=\beth_\kappa$} وي۔
\end{prop}
\olpseganchor{OLP-0591-B013}{end}

\olpseganchor{OLP-0591-B014}{start}
\begin{proof}
لکه په \olref{alephfixed} کښې، خو «\LR{$\beth$}» د
«\LR{$\aleph$}» پر ځاے وکاروئ۔
\end{proof}
\olpseganchor{OLP-0591-B014}{end}

\olpseganchor{OLP-0591-B015}{start}
\begin{prop}\ollabel{stagesize}
\LR{$\card{V_{\omega+\alpha}} = \beth_{\alpha}$}۔ که
\LR{$\omega \ordtimes \omega \leq \alpha$} وي، نو
\LR{$\card{V_\alpha} = \beth_\alpha$}۔
\end{prop}
\olpseganchor{OLP-0591-B015}{end}

\olpseganchor{OLP-0591-B016}{start}
\begin{proof}
لومړۍ ادعا د نامتناهيت هاخوا په استقرا ثابتېږي۔
په بنسټيز پړاو کښې د متناهيو پړاوونو اتحاد شمارېدونکے
نامتناهي دے: هر متناهي پړاو متناهي دے، د دغو پړاوونو
شمارېدونکے اتحاد زياته تر زياته شمېر وړ دے، او ټول طبيعي
عددونه پکښې دي۔ په تالي کښې تواني سټ د بېت د تالي معادله
ورکوي۔ په ناتش حدي حالت کښې د مخکنيو پړاوونو اندازې د
اړوند بېت تر حده هم‌پايه دي، نو د اتحاد اندازه تر دغه حد
کمه نۀ ده۔ د بېت لړۍ په کلکه زياتېږي او د هر شاخص ارزښت
لږ تر لږه د هغه شاخص هومره وي؛ نو د شاخصونو د سټ اندازه
تر همدغه بېت حد زياته نۀ ده۔ د نامتناهي سټونو د محدود اتحاد
قاعده بيا هماغه حد پاسنۍ کرانه ګرځوي۔
دويمه ادعا له دې راځي چې که \LR{$\omega \ordtimes \omega \leq
\alpha$} وي، نو \LR{$\omega + \alpha = \alpha$}۔ د دې د
ثابتولو دپاره د ترتيبي حساب هغه حقيقتونه کاروو چې په
\olref[ord-arithmetic][]{chap} کښې راغلي دي۔ لومړے پام
وکړئ چې \LR{$\omega \ordtimes \omega = \omega \ordtimes (1
\ordplus \omega) = ( \omega  \ordtimes 1) \ordplus
(\omega\ordtimes\omega) = \omega
\ordplus (\omega \ordtimes \omega)$}۔ اوس که
\LR{$\omega \ordtimes \omega \leq \alpha$} وي، يعنې
\LR{$\alpha = (\omega\ordtimes\omega) \ordplus \beta$}
د کوم \LR{$\beta$} دپاره وي، نو \LR{$\omega \ordplus \alpha = \omega
\ordplus ((\omega \ordtimes \omega) \ordplus \beta) =
(\omega \ordplus (\omega \ordtimes \omega)) \ordplus
\beta = (\omega \ordtimes \omega) \ordplus \beta =
\alpha$}۔
\end{proof}
\olpseganchor{OLP-0591-B016}{end}

\olpseganchor{OLP-0591-B017}{start}
\begin{cor}
داسې \LR{$\kappa$} شته چې \LR{$\card{V_\kappa} = \kappa$}۔
\end{cor}
\olpseganchor{OLP-0591-B017}{end}

\olpseganchor{OLP-0591-B018}{start}
\begin{proof}
هغه \LR{$\kappa$} واخلئ چې د \olref{bethfixed} له مخې
د \LR{$\beth$} ثابت ټکی دے۔ د بېت هر ارزښت لږ تر لږه
اومېګا هومره دے، نو ثابت ټکی هم لږ تر لږه اومېګا
هومره دے۔ د کلک زياتوالي له مخې ثابت ټکی لږ تر
لږه د اومېګا په شاخص کښې د بېت هومره دے؛ دا ارزښت
ناشمارېدونکے دے، ځکه د يو په شاخص کښې بېت د کانتور
له قضيې ناشمارېدونکے دے۔ خو د اومېګا ترتيبي مربع
شمارېدونکې ده، نو \LR{$\omega \ordtimes
\omega < \kappa$}۔ نو د \olref{stagesize} له مخې
\LR{$\card{V_\kappa} = \beth_\kappa= \kappa$}۔
\end{proof}
\olpseganchor{OLP-0591-B018}{end}

\olpseganchor{OLP-0591-B019}{start}
د \LR{$V_\kappa$} لاندې دومره پړاوونه شته، څومره چې په
\LR{$V_\kappa$} کښې غړي شته۔ نو په وجداني توګه
\LR{$V_\kappa$} هماغومره پلن دے څومره چې لوړ دے۔
دا د ټرېسټرم شېنډي د کيسې په شان ده: له يوه پړاو څخه
بل ته د \emph{توان سټونو} په اخيستلو ځو، او په هر
ګام کښې مو مراتبي جوړښت \emph{ډېر} لوېږي۔ خو «په
پاے کښې»، يعنې د \LR{$\kappa$} په پړاو کښې، د جوړښت
پلنوالے له لوړوالي سره برابر شي۔
\olpseganchor{OLP-0591-B019}{end}

\olpseganchor{OLP-0591-B020}{start}
ښايي څوک وپوښتي: \emph{د مراتبي جوړښت پلنوالے
څو ځله له لوړوالي سره برابرېږي؟} ځواب دا دے:
\emph{د ترتيبي عددونو هومره ځله۔} خو دا لږې
تشريح ته اړتيا لري۔
\olpseganchor{OLP-0591-B020}{end}

\olpseganchor{OLP-0591-B021}{start}
يو ترم \LR{$\tau$} په لاندې ډول تعريفوو۔ د هر \LR{$A$} دپاره:
\begin{align*}
	\tau_0(A) & \defis \cardsucc{\card{A}}\\
	\tau_{n+1}(A) & \defis \beth_{\tau_n(A)}\\
	\tau(A) & \defis \bigcup_{n < \omega}\tau_n(A)
\intertext{\RL{لکه په \olref{bethfixed} کښې، \LR{$\tau(A)$} د
\LR{$\beth$} ثابت ټکی دے، د هر \LR{$A$} دپاره۔ د بېت لړۍ هر شاخص
لږ تر لږه د هماغه شاخص په اندازه ارزښت لري، نو تکرار
نۀ کمېږي۔ که کوم ګام برابر شي، هماغه ارزښت ثابت ټکی
دے؛ که تل زياتېږي، د لړۍ په حدي اتحاد کښې د بېت د حدي
معادلې له مخې ثابت ټکی ترلاسه کېږي۔ ځکه له اصلي تالي
څخه پيل کوو، \LR{$\card{A} < \tau(A)$}۔ نو اوس دا بازګښتي
تعريف وګورئ:}}
	W_0 &\defis \tau(0)\\
	W_{\alpha + 1} & \defis \tau(W_\alpha)\\
	W_\alpha & \defis \bigcup_{\beta < \alpha} W_\beta
	\text{\RL{، چې \LR{$\alpha$} حدي وي}}
\end{align*}
دا جوړښت د ټولو ترتيبي عددونو دپاره تعريف شوے دے۔
لومړنے غړے ثابت ټکی دے، او هر تالی تر مخکني ثابت ټکي
لوے ثابت ټکی دے، ځکه د اصلي عدد اصلي شمېر هماغه عدد
دے۔ په ناصفر حدي حالت کښې د کوچنيو ثابتو ټکو اتحاد يو نوی
ثابت ټکی دے: هغوی په خپل اتحاد کښې هم‌پايه دي او د
بېت حدي معادله د هغوی پر ارزښتونو هماغه اتحاد ورکوي۔
نو \LR{$W$} د ترتيبي عددونو څخه د \LR{$\beth$} ثابتو ټکو ته
يو پر يو تابع ده۔ او لکه مخکې، \LR{$V_{W_\alpha}$} د هر
\LR{$\alpha$} دپاره له خپل لوړوالي سره هم‌پلن دے۔
\olpseganchor{OLP-0591-B021}{end}

\olpseganchor{OLP-0592-B004}{start}
\olchapter{sth}{choice}{ټاکنه}
\olpseganchor{OLP-0592-B004}{end}

\olpseganchor{OLP-0593-B003}{start}
\olfileid{sth}{choice}{intro}
\olpseganchor{OLP-0593-B003}{end}

\olpseganchor{OLP-0593-B004}{start}
\olsection{سريزه}
\olpseganchor{OLP-0593-B004}{end}

\olpseganchor{OLP-0593-B005}{start}
په \crefrange{sth:cardinals::chap}{sth:card-arithmetic::chap} کښې مو
اصلي عددونه د ترتيبي عددونو په توګه په ګڼلو سره د اصلي عددونو
نظريه جوړه کړه۔ دا لاره د ښه ترتيب پر اصل ولاړه ده۔ څرګندېږي
چې ښه ترتيب له يوه بل اصل، يعنې د ټاکنې له اصل سره، همارز
دے؛ د دې اصل د منلو په اړه جدي فلسفي بحث هم شوے دے۔
په دې څپرکي کښې مو پوښتنې دا دي: دا اصل څنګه کارول کېږي،
او ايا توجيه کېدلے شي؟
\olpseganchor{OLP-0593-B005}{end}

\olpseganchor{OLP-0594-B003}{start}
\olfileid{sth}{choice}{tarskiscott}
\olsection{د تارسکي--سکاټ چاره}
\olpseganchor{OLP-0594-B003}{end}

\olpseganchor{OLP-0594-B004}{start}
په \olref[cardinals][cardsasords]{defcardinalasordinal} کښې مو
اصلي عددونه د ترتيبي عددونو په توګه تعريف کړي وو۔ د دې دپاره
مو د ښه ترتيب اصل فرض کړے و۔ دا کار مو يوازې ځکه وکړ چې دا
«عام منل شوے دود» دے۔
\olpseganchor{OLP-0594-B004}{end}

\olpseganchor{OLP-0594-B005}{start}
نو مخکې له دې چې د ښه ترتيب په اړه فلسفي مسئلې وڅېړو،
بايد روښانه کړو چې موږ له دغه عام دود څخه \emph{وتلے}
شو او د ښه ترتيب له فرضولو \emph{پرته} هم د اصلي عددونو
نظريه جوړولے شو۔ د \LR{$\cardeq{A}{B}$}، \LR{$\cardle{A}{B}$}
او \LR{$\cardless{A}{B}$} تعريفونه لا کارولے شو، لکه څنګه
چې په \olref[sfr][siz][]{chap} کښې راغلي دي۔ يوازې
د \emph{اصلي عدد} نوي تصور ته اړتيا لرو۔
\olpseganchor{OLP-0594-B005}{end}

\olpseganchor{OLP-0594-B006}{start}
يو ساده فکر دا دے چې د \LR{$A$} اصلي شمېر داسې تعريف کړو:
\[
	\Setabs{x}{\cardeq{A}{x}}.
\]
دا د فرېګه له هغې تعريف \LR{$\# x Fx$} سره پرتله کولے
شئ چې د \olref[cardinals][hp]{sec} په پاے کښې مو
په لنډه بڼه راوړے و۔ خو د هغو دلايلو له مخې چې هلته
مو ورته اشاره وکړه، دا تعريف نه سمېږي۔ هر يو غړيز
سټ له \LR{$\{\emptyset\}$} سره هم‌شمېره دے۔ خو د مراتبي
جوړښت په هر تالي پړاو کښې نوي يو غړيز سټونه جوړېږي
(يوازې د تېر پړاو يو غړيز سټ په پام کښې ونيسئ)۔
د يو غړي سټ دپاره نو \LR{$\Setabs{x}{\cardeq{A}{x}}$}
يو سټ نۀ دے، ځکه غړي يې تر هرې رتبې هاخوا موندل کېږي۔
همدغه ستونزه د هر ناتش سټ دپاره شته: له هر يو غړيز سټ
سره يې کارتيزي حاصل له خپله سټ سره هم‌شمېره دے، او
د يو غړي په لوړه رتبه اخيستلو سره د حاصل رتبه هم
هر څومره لوړېدے شي۔ له تش سټ سره يوازې تش سټ هم‌شمېره
دے۔
\olpseganchor{OLP-0594-B006}{end}

\olpseganchor{OLP-0594-B007}{start}
د دې مسئلې د هوارولو دپاره د تارسکي او سکاټ
چاره کاروو:\footnote{يادونه: په ټولو فورمولونو کښې
پارامترونه کېدے شي، مګر که په ښکاره بل څه وويل شي۔}
\olpseganchor{OLP-0594-B007}{end}

\olpseganchor{OLP-0594-B008}{start}
\begin{defn}[تارسکي--سکاټ]
د هر فورمول \LR{$\phi(x)$} دپاره،
\LR{$[ x : \phi(x)] $} د ټولو هغو \LR{$x$} سټ وبولئ چې
تر ټولو ټيټه ممکنه رتبه لري او \LR{$\phi(x)$} پوره کوي
(يا \LR{$\emptyset$}، که د فورمول \LR{$\phi$} هېڅ شاهد نۀ وي)۔
\end{defn}
\olpseganchor{OLP-0594-B008}{end}

\olpseganchor{OLP-0594-B009}{start}
بايد وګورو چې دا تعريف روا دے۔ په \LR{$\ZF$} کښې
\olref[spine][foundation]{zfentailsregularity} تضمينوي
چې \LR{$\setrank{x}$} د هر \LR{$x$} دپاره شتون لري۔ اوس که
يو داسې شاهد ولرو چې \LR{$\phi$} پوره کوي، تر هغې رتبې
پورې د بېلونې له لارې د هغو رتبو سټ اخلو چې يو شاهد
لري، او لږ غړے يې \LR{$\alpha$} بولو۔ په دې رتبه کښې
\LR{$(\exists x\subseteq V_\alpha)\phi(x)$}، او په هره
ترې ټيټه رتبه کښې
\LR{$(\forall \beta
\in \alpha)(\forall x \subseteq V_\beta)\lnot \phi(x)$}۔
بيا \LR{$[x : \phi(x)]$} د بېلونې په کارولو سره د
\LR{$\Setabs{x \in
V_{\alpha+1}}{\phi(x)}$} په توګه تعريفولے شو۔ د دې
پړاو غړي تر ټاکلې رتبې لوړه رتبه نۀ لري، او ټيټه
رتبه لرونکے غړے فورمول نۀ پوره کوي۔
\olpseganchor{OLP-0594-B009}{end}

\olpseganchor{OLP-0594-B010}{start}
د تارسکي--سکاټ چاره چې مو توجيه کړه، اوس يې د
اصلي شمېر د يوه تصور د تعريف دپاره کارولے شو:
\olpseganchor{OLP-0594-B010}{end}

\olpseganchor{OLP-0594-B011}{start}
\begin{defn}
د \LR{$A$} د \textsc{ts}-اصلي شمېر تعريف
\LR{$\text{\RL{tsc}}(A) = [x :
\cardeq{A}{x}]$} دے۔
\end{defn}
\olpseganchor{OLP-0594-B011}{end}

\olpseganchor{OLP-0594-B012}{start}
د \textsc{ts}-اصلي عدد تعريف ښه ترتيب نه کاروي۔
خو بې له دې اصله هم ښودلے شو چې
\emph{\textsc{ts}-اصلي عددونه} د هغو
\emph{اصلي عددونو} په شان چلند کوي چې په
\olref[cardinals][cardsasords]{defcardinalasordinal}
کښې تعريف شوي دي۔ په
\olref[cardinals][cardsasords]{lem:CardinalsBehaveRight}
کښې د برابرو او نابرابرو اندازو دعوې داسې بيا بيانولے
شو: دوه تارسکي--سکاټ اصلي عددونه هله برابر دي چې
استازي يې هم‌شمېره وي۔ په دې حالت کښې د هم‌شمېرو
سټونو طبقه يوه ده، نو د لږې رتبې استازي يې هم يو شان
دي۔ د لومړي اصلي عدد اندازه تر دويمې زياته نۀ ده که
له يوه استازي نه بل ته يو پر يو تابع وي۔ هر استازے
د خپلې طبقې له نورو استازو سره هم‌شمېره دے، نو دا
اړيکه د استازي په ټاکلو پورې نۀ تړل کېږي۔ د
شرودر--برنشتاين قضيه يې جزوي ترتيب ثابتوي؛ کلک
ترتيب يې برابري وباسي۔ د ټولو جوړو پرتله د ښه
ترتيب له اصل سره همارزه ده، لکه په راتلونکې برخه
کښې چې به يې ووينو۔ د
\olref[card-arithmetic][opps]{lem:SizePowerset2Exp}
د تواني سټ دوه اړخيزه تابع په \LR{$\ZF$} کښې هم
هماغه د اندازې برابري ورکوي۔
\olpseganchor{OLP-0594-B012}{end}

\olpseganchor{OLP-0594-B013}{start}
په همدې ترڅ کښې، د تارسکي--سکاټ د چارې په وسيله
د ترتيبي عددونو نظريه هم جوړولے شو۔ که
\LR{$\tuple{A, <}$} ښه ترتيب وي، نو
\LR{$\text{\RL{tso}}(A, <) = [\tuple{X,
R} : \ordeq{\tuple{A, <}}{\tuple{X, R}}]$}
وبولئ۔ د اصلي او ترتيبي عددونو د دې لارې په اړه
نور په \mbox{\LR{\latinfont \citet[chs.~9--12]{Potter2004}}} کښې وګورئ۔
\olpseganchor{OLP-0594-B013}{end}

\olpseganchor{OLP-0595-B003}{start}
\olfileid{sth}{choice}{hartogs}
\olsection{د پرتلې وړتيا او د هارتوګس لمه}
\olpseganchor{OLP-0595-B003}{end}

\olpseganchor{OLP-0595-B004}{start}
بې له ټاکنې د هرې جوړې د اندازو پرتله نۀ شي ثابتېدې۔
د \mbox{\LR{\latinfont \cite{Hartogs1915}}} لاندې پايله د دې خبرې په ثبوت
کښې کارېږي۔
\olpseganchor{OLP-0595-B004}{end}

\olpseganchor{OLP-0595-B005}{start}
\begin{lem}[\emph{په \LR{$\ZF$} کښې}]\ollabel{HartogsLemma}
د هر سټ \LR{$A$} دپاره داسې ترتيبي عدد \LR{$\alpha$} شته
چې \LR{$\cardnless{\alpha}{A}$} وي۔
\end{lem}
\olpseganchor{OLP-0595-B005}{end}

\olpseganchor{OLP-0595-B006}{start}
\begin{proof}
که \LR{$B \subseteq A$} او \LR{$R \subseteq B^2$} وي، نو
د \olref[ord-arithmetic][using-addition]{rankcomputation}
له مخې \LR{$\tuple{B, R} \subseteq
V_{\setrank{A}+4}$} وي۔ نو د بېلونې په
کارولو دا سټ په پام کښې ونيسئ:
\[
	C = \Setabs{\tuple{B, R} \in V_{\setrank{A}+5}}{B\subseteq A 
	\text{\RL{ او \LR{$\tuple{B, R}$} ښه ترتيب دے}}}
\]
د تعويض او
\olref[ordinals][ordtype]{thmOrdinalRepresentation}
په کارولو دا سټ جوړ کړئ:
\[
	\alpha = \Setabs{\ordtype{B, R}}{\tuple{B, R} \in C}.
\]
د \olref[ordinals][basic]{corordtransitiveord} له مخې
\LR{$\alpha$} ترتيبي عدد دے، ځکه دا د ترتيبي عددونو انتقالي
سټ دے۔ ځکه که \LR{$\gamma \in \beta \in \alpha$} وي، نو
\LR{$\beta = \ordtype{B, R}$} د داسې کوم
\LR{$B \subseteq A$} دپاره وي، او بيا
\LR{$\gamma = \ordtype{B_b, R_b}$} د کوم \LR{$b
\in B$} دپاره وي، د
\olref[ordinals][iso]{wellordinitialsegment} له مخې۔
دا پيلنۍ برخه هم د اصلي سټ پر فرعي سټ ښه ترتيب دے،
نو \LR{$\gamma \in \alpha$}۔
\olpseganchor{OLP-0595-B006}{end}

\olpseganchor{OLP-0595-B007}{start}
د تناقض دپاره فرض کړئ چې
يو پر يو تابع \LR{$f \colon \alpha \to A$} شته۔
بيا چې:
\begin{align*}
	B &= \ran{f}\\
	R &= \Setabs{\tuple{f(\xi), f(\eta)} \in A \times A}{\xi,\eta\in\alpha \text{\RL{ او }} \xi\in\eta}.
\end{align*}
دا اړيکه د غړيتوب ترتيب د تابعې حاصل ته لېږدوي۔
ځکه تابع يو پر يو ده او پر خپل حاصل باندې ده، د دواړو
ترتيبونو هم‌شکلي ده۔ نو \LR{$\alpha = \ordtype{B, R}$} او
\LR{$\tuple{B, R} \in C$}۔ نو \LR{$\alpha
\in \alpha$}، چې تناقض دے۔
\end{proof}
\olpseganchor{OLP-0595-B007}{end}

\olpseganchor{OLP-0595-B008}{start}
له دې څخه يوه ژوره پايله راځي:
\olpseganchor{OLP-0595-B008}{end}

\olpseganchor{OLP-0595-B009}{start}
\begin{thm}[\emph{په \LR{$\ZF$} کښې}]
لاندې دعوې همارزې دي:
\begin{enumerate}
	\item\ollabel{equivwo} د ښه ترتيب اصل
	\item\ollabel{equivcompare} يا \LR{$\cardle{A}{B}$} وي يا
	\LR{$\cardle{B}{A}$} وي، د هر دوو سټونو \LR{$A$} او \LR{$B$}
	دپاره
\end{enumerate}
\end{thm}
\olpseganchor{OLP-0595-B009}{end}

\olpseganchor{OLP-0595-B010}{start}
\begin{proof}
\emph{\olref{equivwo} \LR{$\Rightarrow$} \olref{equivcompare}.}
\LR{$A$} او \LR{$B$} وټاکئ۔ د \olref{equivwo} په کارولو
ښه ترتيبونه \LR{$\tuple{A, R}$} او \LR{$\tuple{B, S}$}
شته۔ د
\olref[ordinals][ordtype]{thmOrdinalRepresentation}
له مخې \LR{$f \colon
\alpha \to \tuple{A, R}$} او \LR{$g \colon \beta \to \tuple{B, S}$}
هم‌شکلي وګڼئ۔ د
\olref[sth][ordinals][basic]{ordinalsaresubsets} له مخې
يا \LR{$\alpha \subseteq \beta$} وي يا
\LR{$\beta \subseteq \alpha$} وي۔ که
\LR{$\alpha \subseteq \beta$} وي، نو
\LR{$\comp{f^{-1}}{g} \colon A \to B$} يوه
يو پر يو تابع ده، نو \LR{$\cardle{A}{B}$}؛ او که
\LR{$\beta \subseteq \alpha$} وي، په همدې ډول
\LR{$\cardle{B}{A}$}۔
\olpseganchor{OLP-0595-B010}{end}

\olpseganchor{OLP-0595-B011}{start}
\emph{\olref{equivcompare} \LR{$\Rightarrow$} \olref{equivwo}.}
\LR{$A$} وټاکئ؛ د \olref{HartogsLemma} له مخې داسې
ترتيبي عدد \LR{$\beta$} شته چې
\LR{$\cardnless{\beta}{A}$} وي۔ د
\olref{equivcompare} په کارولو
\LR{$\cardle{A}{\beta}$} لرو۔ نو يوه
يو پر يو تابع \LR{$f \colon A \to
\beta$} شته، او د \LR{$A$} غړي د دې يو پر يو تابع په وسيله
د \LR{$\Setabs{\tuple{a, b} \in A \times A}{f(a)
\in f(b)}$} اړيکې په تعريفولو سره ښه ترتيبولے شو۔
\end{proof}
\noindent
سمدستي پايله دا ده: که ښه ترتيب ناکام شي، ځينې سټونه
د اندازې له مخې \emph{په رښتيا نه پرتله کېږي}۔ نو که
ښه ترتيب ناکام شي، د اصلي عددونو نامتناهيت هاخوا حساب
به ګډ وډ شي۔ د بېلګې په توګه، دا فکر به پرېږدو چې
که \LR{$A$} او \LR{$B$} نامتناهي وي، نو
\LR{$\cardeq{\cardeq{A \disjointsum B}{A \times
B}}{M}$}، چې \LR{$M$} د \LR{$A$} او \LR{$B$} له منځه لوے
سټ دے (وګورئ
\olref[card-arithmetic][simp]{cardplustimesmax})۔
مسئله ساده ده: که د \LR{$A$} او \LR{$B$} اندازه
\emph{پرتله} نه شو کړے، نو دا پوښتنه بې‌معنا ده
چې کوم يو لوے دے۔
\olpseganchor{OLP-0595-B011}{end}

\olpseganchor{OLP-0596-B003}{start}
\olfileid{sth}{choice}{woproblem}
\olsection{د ښه ترتيب مسئله}
\olpseganchor{OLP-0596-B003}{end}

\olpseganchor{OLP-0596-B004}{start}
ښکاره ده چې د ښه ترتيب منل يا نه منل ډېرې پايلې لري۔
خو د دې اصل بحث تر ډېره د هغه پر يوه همارز اصل، يعنې
د ټاکنې پر اصل، راڅرخي۔ نو اوس همدغه اصل څېړو
(او همارزي به يې ثابته کړو)۔
\olpseganchor{OLP-0596-B004}{end}

\olpseganchor{OLP-0596-B005}{start}
په \mbox{\LR{\latinfont \citeyear{Cantor1883}}} کښې کانتور د ښه ترتيب
له اصله ملاتړ وکړ۔ هغه ښه ترتيب د استدلال له بنسټونو څخه ګاڼه، ځکه د اصل کارونه عامه وه او ګڼې پايلې ترې راوتلې۔ دا دريځ \mbox{\LR{\latinfont \citeauthor{Potter2004}}} په خپل کتاب کښې رااخلي
\mbox{\LR{\latinfont \citeyear[p.~243]{Potter2004}}}۔
خو کانتور بالاخره دې باور ته ورسېد چې دا «اصل»
ثبوت غواړي۔ د رياضي پوهانو ټولنه هم ورته نظر ته ورسېده۔
\olpseganchor{OLP-0596-B005}{end}

\olpseganchor{OLP-0596-B006}{start}
دا مسئله په \mbox{\LR{\latinfont \citeyear{Zermelo1904}}} کښې زېرملو
«هواره» کړه۔ د هغه د حل د بيانولو دپاره يو څو
تعريفونو ته اړتيا لرو۔
\olpseganchor{OLP-0596-B006}{end}

\olpseganchor{OLP-0596-B007}{start}
\begin{defn}
تابعه \LR{$f$} د \emph{انتخاب تابعه} هله او يوازې هله ده
چې \LR{$f(x) \in x$} د هر \LR{$x \in \dom{f}$} دپاره وي۔
وايو چې \LR{$f$} \emph{د \LR{$A$} دپاره د ټاکنې تابعه} ده
هله او يوازې هله چې \LR{$f$} د ټاکنې تابعه وي او
\LR{$\dom{f} = A \setminus \{\emptyset\}$} وي۔
\end{defn}
\olpseganchor{OLP-0596-B007}{end}

\olpseganchor{OLP-0596-B008}{start}
په وجداني توګه، د هر (ناخالي) سټ
\LR{$x \in A$} دپاره د \LR{$A$} د ټاکنې تابعه
يو ځانګړے غړے، \LR{$f(x)$}، له \LR{$x$} څخه
\emph{ټاکي}۔ نو د ټاکنې اصل دا دے:
\olpseganchor{OLP-0596-B008}{end}

\olpseganchor{OLP-0596-B009}{start}
\begin{axiom}[ټاکنه]
هر سټ د ټاکنې تابعه لري۔
\end{axiom}
\olpseganchor{OLP-0596-B009}{end}

\olpseganchor{OLP-0596-B010}{start}
زېرملو وښودله چې ټاکنه ښه ترتيب لازموي، او
ښه ترتيب هم ټاکنه لازموي:
\olpseganchor{OLP-0596-B010}{end}

\olpseganchor{OLP-0596-B011}{start}
\begin{thm}[په \LR{$\ZF$} کښې]\ollabel{thmwochoice}
ښه ترتيب او ټاکنه همارز دي۔
\end{thm}
\olpseganchor{OLP-0596-B011}{end}

\olpseganchor{OLP-0596-B012}{start}
\begin{proof}
\emph{له چپ څخه ښي ته۔} \LR{$A$} د سټونو يو سټ
ونيسئ۔ نو \LR{$\bigcup A$} د اتحاد د اصل له مخې
شتون لري، او د ښه ترتيب له مخې داسې
\LR{$<$} شته چې \LR{$\bigcup A$} ښه ترتيبوي۔ اوس
\LR{$f(x) = \text{\RL{د \LR{$<$} له مخې تر ټولو کوچنے غړے د }}x$}
د هر ناتش غړي دپاره وټاکئ۔ دا تابع د بېلونې له لارې
د اړوند کارتيزي حاصل دننه جوړېږي، او د \LR{$A$} دپاره
د ټاکنې تابعه ده۔
\olpseganchor{OLP-0596-B012}{end}

\olpseganchor{OLP-0596-B013}{start}
\emph{له ښي څخه چپ ته۔} \LR{$A$} وټاکئ۔ د ټاکنې له مخې
د ټاکنې تابعه \LR{$f$} د \LR{$\Pow{A} \setminus \{\emptyset\}$}
دپاره شته۔ د \olref[hartogs]{HartogsLemma} له مخې د \LR{$A$}
دپاره يو داسې ترتيبي عدد \LR{$h$} واخلو چې \LR{$A$} ته ترې يو پر
يو تابع نۀ وي۔ د نامتناهيت هاخوا بازګښت پر \LR{$h$} باندې
تابع په دې ډول تعريفوي:
\begin{align*}
	g(\alpha) &=
		\begin{cases}
			A &\text{\RL{که }} A \subseteq \funimage{g}{\alpha}\\
			f(A \setminus \funimage{g}{\alpha}) & \text{\RL{په بل صورت کښې.}}
		\end{cases}
\end{align*}
د رتبې له مخې \LR{$A\notin A$}، نو خپله \LR{$A$} د تمېدو نښه ده۔
بازګښت هر نوی ارزښت د هماغه پړاو د مخکنيو ارزښتونو له
حاصله ټاکي۔ \olref[sth][spine][recursion]{transrecursionfun}
په دې ټاکلې ساحه کښې د تابعې شتون او يوازينيوالے راکوي۔
\olpseganchor{OLP-0596-B013}{end}

\olpseganchor{OLP-0596-B014}{start}
تر تمېدو مخکې، د هر \LR{$\alpha$} دپاره
\LR{$g(\alpha) =
f(A \setminus \funimage{g}{\alpha}) \in A \setminus
\funimage{g}{\alpha}$}؛ نو
\LR{$g(\alpha) \notin \funimage{g}{\alpha}$}۔
که \LR{$g(\alpha) = g(\beta)$} وي، نو
\LR{$g(\beta) \notin \funimage{g}{\alpha}$}، يعنې
\LR{$\beta \notin \alpha$}؛ او په همدې ډول
\LR{$\alpha \notin \beta$}۔ د درې‌ګونيتوب له مخې
\LR{$\alpha = \beta$}۔ نو د تمېدو نه مخکې محدوده شوې تابعه
يو پر يو ده۔
\olpseganchor{OLP-0596-B014}{end}

\olpseganchor{OLP-0596-B015}{start}
که تر \LR{$h$} مخکې هېڅ پړاو کښې تم نۀ شو، نو \LR{$g$} به
له \LR{$h$} نه \LR{$A$} ته يو پر يو وي، چې له
\olref[hartogs]{HartogsLemma} سره تناقض لري۔ نو داسې
لږ ترتيبي عدد \LR{$\alpha$} شته چې \LR{$A = g[\alpha]$}؛
تر دغه پړاو مخکې ټول ارزښتونه د \LR{$A$} غړي دي او په
همدغه پړاو کښې يې حاصل ټول \LR{$A$} دے۔ که \LR{$A$} تش وي،
\LR{$\alpha=0$} دے۔
\olpseganchor{OLP-0596-B015}{end}

\olpseganchor{OLP-0596-B016}{start}
د \LR{$g$} لومړۍ \LR{$\alpha$} برخې ته محدودول له \LR{$\alpha$}
څخه \LR{$A$} ته دوه اړخيزه يو پر يو تابع ورکوي۔ د ترتيبي عدد ښه
ترتيب د همدې محدودې تابعې په وسيله \LR{$A$} ته لېږدوو۔
\end{proof}
\olpseganchor{OLP-0596-B016}{end}

\olpseganchor{OLP-0596-B017}{start}
نو ښه ترتيب او ټاکنه دواړه يو ځاے ټينګېږي
يا يو ځاے لوېږي۔ خو پوښتنه لا پاتې ده: ايا
دواړه ټينګېږي که لوېږي؟
\olpseganchor{OLP-0596-B017}{end}

\olpseganchor{OLP-0597-B003}{start}
\olfileid{sth}{choice}{countablechoice}
\olsection{د شمېر وړه ټاکنه}
\olpseganchor{OLP-0597-B003}{end}

\olpseganchor{OLP-0597-B004}{start}
بې له ټاکنې يا ښه ترتيب څخه دا په اسانۍ ثابتېږي چې:
\olpseganchor{OLP-0597-B004}{end}

\olpseganchor{OLP-0597-B005}{start}
\begin{lem}[په \LR{$\Zminus$} کښې]
هر متناهي سټ د ټاکنې تابعه لري۔
\end{lem}
\olpseganchor{OLP-0597-B005}{end}

\olpseganchor{OLP-0597-B006}{start}
\begin{proof}
\LR{$a = \{b_1, \ldots, b_n\}$} ونیسئ۔ له نوملړ نه تش سټ وباسئ او هر پاتې غړے يوازې يو ځل
وليکئ۔ اوس هر \LR{$b_i
\neq \emptyset$} دے۔ که هېڅ غړے پاتې نۀ وي، تشه تابعه
بس ده۔ نو داسې څيزونه
\LR{$c_1, \ldots, c_n$} شته چې
\LR{$c_1 \in b_1, \ldots, c_n \in b_n$} وي۔
اوس د
\olref[z][pairs]{prop:pairsconsequences} له مخې
سټ \LR{$\{\langle b_1,
c_1\rangle , \ldots, \langle b_n, c_n\rangle\}$}
شتون لري؛ او دا د~\LR{$a$} دپاره د ټاکنې تابعه ده۔
\end{proof}
\olpseganchor{OLP-0597-B006}{end}

\olpseganchor{OLP-0597-B007}{start}
خو د نامتناهي سټونو په پام کښې نيولو سره چارې
سمدستي لا مبهمېږي۔ د بېلګې په توګه، د پورته
پايلې دا «تر ټولو کوچنے» پراخول وګورئ:
\olpseganchor{OLP-0597-B007}{end}

\olpseganchor{OLP-0597-B008}{start}
\begin{defish}
\emph{د شمېر وړه ټاکنه۔} هر \emph{د شمېر وړ}
سټ د ټاکنې تابعه لري۔
\end{defish}
\olpseganchor{OLP-0597-B008}{end}

\olpseganchor{OLP-0597-B009}{start}
دا د ټاکنې د اصل يو ځانګړے حالت دے۔ او څرګندېږي
چې دا اصل ډېر ځله بې له دې کارول شوے چې کاروونکي
يې په ښکاره پوه شوي وي۔ دوه ښې بېلګې يې دلته دي۔
\footnote{د \mbox{\LR{\latinfont \citet[\S9.4]{Potter2004}}} او لوکا
انکورواتي له مخې۔}
\olpseganchor{OLP-0597-B009}{end}

\olpseganchor{OLP-0597-B010}{start}
\begin{ex}
دا يو طبيعي فکر دے: د هر سټ \LR{$A$} دپاره يا
\LR{$\cardle{\omega}{A}$} وي، يا
\LR{$\cardeq{A}{n}$} د کوم \LR{$n \in \omega$}
دپاره وي۔ دا د دې وجداني مفکورې د بيان يوه
لاره ده چې هر سټ يا متناهي دے يا نامتناهي۔
کانتور او ډېرو نورو رياضيپوهانو دا دعوه بې له
ثبوته کړې وه۔ موږ په احتياط سره دا په
\olref[cardinals][classing]{generalinfinitycharacter}
کښې ثابته کړه۔ خو په هغه ثبوت کښې مو په
\LR{$\ZFC$} کښې کار کاوه، ځکه دا مو فرض کړې
وه چې هر سټ \LR{$A$} ښه ترتيبېږي، او ځکه
\LR{$\card{A}$} خامخا شتون لري۔ يعنې موږ
ټاکنه په ښکاره فرض کړے و۔
\olpseganchor{OLP-0597-B010}{end}

\olpseganchor{OLP-0597-B011}{start}
په حقيقت کښې \mbox{\LR{\latinfont \citet{Dedekind1888}}} د
دې دعوې خپل ثبوت په لاندې ډول وړاندې کړ:
\olpseganchor{OLP-0597-B011}{end}

\olpseganchor{OLP-0597-B012}{start}
\begin{thm}[په \LR{$\Zminus + \text{\RL{د شمېر وړه ټاکنه}}$} کښې]
د هر \LR{$A$} دپاره يا \LR{$\cardle{\omega}{A}$}
وي يا \LR{$\cardeq{A}{n}$} د کوم
\LR{$n \in \omega$} دپاره وي۔
\end{thm}
\olpseganchor{OLP-0597-B012}{end}

\olpseganchor{OLP-0597-B013}{start}
\begin{proof}
فرض کړئ چې \LR{$\cardneq{A}{n}$} د هر
\LR{$n \in \omega$} دپاره وي۔ نو په ځانګړي
ډول د هر \LR{$n < \omega$} دپاره داسې
فرعي سټ \LR{$A_n \subseteq A$} شته چې
عين \LR{$\cardexpo{2}{n}$} غړي لري۔ د دې
لړۍ \LR{$A_0, A_1, A_2,
\ldots$} په کارولو د هر \LR{$n$} دپاره دا
تعريفوو:
\[
	B_n = A_n \setminus \bigcup_{i < n} A_i.
\]
اوس لاندې ته پام وکړئ:
\begin{align*}
	\card{\bigcup_{i < n}A_i} 
	&\leq \card{A_0} + \card{A_1} + \ldots + \card{A_{n-1}}\\
	&=1 + 2 + \ldots + 2^{n-1}\\
	& = 2^n - 1\\
	& < 2^n = \card{A_n}
\end{align*}
د صفر شاخص دپاره پورتنۍ مجموعه تشه ده او جمع صفر ده۔
نو هر \LR{$B_n$} لږ تر لږه يو غړے،
\LR{$c_n$}، لري۔ سربېره پردې، \LR{$B_n$}-ګانې
دوه په دوه بېلې دي؛ نو که \LR{$c_n = c_m$}
وي، بيا \LR{$n = m$} وي۔ خو هر
\LR{$c_n
\in A$} دے۔ نو تابعه \LR{$f(n) = c_n$}
له \LR{$\omega \to A$} څخه يوه يو پر يو
تابعه ده۔
\end{proof}
\noindent
دېدېکيند په ښکاره ونه ويل چې شمېر وړې ټاکنې نه يې کار اخيستے دے۔ خو \emph{تاسو} يې
کارونه وليدله؟ بيا يې وګورئ۔ (په رښتيا:
\emph{بيا يې وګورئ}۔)
\olpseganchor{OLP-0597-B013}{end}

\olpseganchor{OLP-0597-B014}{start}
په ثبوت کښې د شمېر وړه ټاکنه دوه ځله
کارول شوے۔ يو ځل مو د سټونو لړۍ
\LR{$A_0$}, \LR{$A_1$}, \LR{$A_2$}, \dots\@
ترلاسه کولو ته وکاراوه۔ بيا مو د غړو
\LR{$c_n$} د هر~\LR{$B_n$} څخه ټاکلو ته وکاراوه۔ په لومړي
کار کښې د ټاکنې تابع د ټاکلي متناهي شمېر لرونکو فرعي
سټونو له هرې کورنۍ يو سټ اخلي؛ په دويم کښې د ناتشو
ټوټو له هرې يو غړے اخلي۔ د دواړو لړيو ګرافونه په ټاکلو
کارتيزي حاصلونو کښې د بېلونې له لارې جوړېږي۔
سربېره پردې، دلته د ټاکنې کارونه نه شي
لرې کېدے۔ \mbox{\LR{\latinfont \citet[p.~138]{Cohen1966}}}
ثابته کړې چې دا پايله بې له هر ډول
ټاکنې ناکامېږي۔ يعنې له \LR{$\ZF$} سره
دا سازګاره ده چې ځينې سټونه له~\LR{$\omega$}
سره \emph{نه پرتله کېږي}۔
\end{ex}
\olpseganchor{OLP-0597-B014}{end}

\olpseganchor{OLP-0597-B015}{start}
\begin{ex}
په \mbox{\LR{\latinfont \citeyear{Cantor1878}}} کښې کانتور وويل
چې د شمېر وړ سټونو د شمېر وړ اتحاد هم
د شمېر وړ دے۔ هغه ثبوت نه و وړاندې کړے؛
ښايي ثبوت يې ښکاره ګڼلے و۔ موږ بيا په
احتياط سره د دې پايلې يوه عامه بڼه په
\olref[card-arithmetic][simp]{kappaunionkappasize}
کښې ثابته کړې۔ خو زموږ ثبوت ټاکنه
په ښکاره فرض کړے و۔ ان د لږې عامې
پايلې ثبوت هم د شمېر وړه ټاکنه غواړي۔
\olpseganchor{OLP-0597-B015}{end}

\olpseganchor{OLP-0597-B016}{start}
\begin{thm}[په \LR{$\Zminus + \text{\RL{د شمېر وړه ټاکنه}}$} کښې]
که \LR{$A_n$} د هر \LR{$n \in \omega$}
دپاره د شمېر وړ وي، نو \LR{$\bigcup_{n <
\omega} A_n$} هم د شمېر وړ دے۔
\end{thm}
\olpseganchor{OLP-0597-B016}{end}

\olpseganchor{OLP-0597-B017}{start}
\begin{proof}
که ټول اتحاد تش وي، پايله څرګنده ده۔ که يو غړے
ولري، يو داسې غړے واخلو او هر تش پړاو د همدغه غړي
په يو غړيز سټ بدل کړو۔ اتحاد هماغسې پاتې کېږي او اوس
هر \LR{$A_n \neq \emptyset$} دے۔ نو د هر \LR{$n \in \omega$} دپاره
پر هدف سټ بشپړه تابع \LR{$f_n \colon \omega
\to A_n$} شته۔ تابعه \LR{$f \colon \omega
\times \omega \to \bigcup_{n <
\omega} A_n$} د \LR{$f(m, n) = f_n(m)$}
په وسيله تعريف کړئ۔ پايله له دې راځي چې
\LR{$\omega
\times \omega$} د شمېر وړ دے
(\olref[sfr][siz][zigzag]{natsquaredenumerable})
او \LR{$f$} يوه پر هدف سټ بشپړه تابع ده۔
\end{proof}
\noindent
ايا تاسو د شمېر وړې ټاکنې کارونه وليدله؟
دا د تابعو لړۍ \LR{$f_0$}, \LR{$f_1$}, \LR{$f_2$},
\dots\footnote{د ټاکنې ورته کارونه په
\olref[card-arithmetic][simp]{kappaunionkappasize}
کښې هم شوې وه، چې هلته مو وويل «د هر
\LR{$\beta \in \cardfont{a}$} دپاره يوه
يو پر يو تابع \LR{$f_\beta$} وټاکئ»۔}
د ټاکلو دپاره کارېږي۔ بيا هم، دا پايله
بې له هر ډول ټاکنيز اصل څخه ناکامېږي۔
په ځانګړي ډول،
\mbox{\LR{\latinfont \citet{FefermanLevy1963}}} ثابته کړې
چې له \LR{$\ZF$} سره دا سازګاره ده چې
د شمېر وړ سټونو د شمېر وړ اتحاد اصلي
شمېر~\LR{$\beth_1$} ولري۔ خو د همدې خبرې
توپير د رسل د بوټانو او جرابو په بېلګه کښې ليدل کېږي۔
\begin{explain}
  د رسل په بېلګه کښې يو شتمن کس د بوټانو \LR{$\aleph_0$} جوړې
  او د جرابو \LR{$\aleph_0$} جوړې لري۔ د بوټانو د ټولو جوړو
  لپاره يوه قاعده بس ده: د هرې جوړې ښے بوټ واخلو، يا د
  هرې جوړې کيڼ بوټ۔ د جرابو په هره جوړه کښې داسې ګډه
  نښه نشته چې يوه جرابه ترې په همدې ډول وټاکو۔ د ضرب
  اصل د هرې جوړې لپاره د يوه غړي د ټاکنې د ټولګې شتون
  تضمينوي۔ \mbox{\LR{\latinfont \citep[p.~126]{Russell1919}}}
\end{explain}
په لنډه، د لاندې خبرې د ثبوت دپاره
د ټاکنې يو ډول ته اړتيا ده: که تاسو
د شمېر وړ ډېرې د جرابو جوړې لرئ،
نو يوازې د شمېر وړ ډېرې جرابې لرئ۔
په حقيقت کښې، بې له شمېر وړې ټاکنې
(يا يې له کوم همارز اصل څخه) د شمېر وړ
سټونو د شمېر وړ اتحاد کېدے شي د شمېر
وړ نه وي۔
\end{ex}
\olpseganchor{OLP-0597-B017}{end}

\olpseganchor{OLP-0597-B018}{start}
پايله دا ده چې د شمېر وړه ټاکنه ډېری
ځله بې له دې کارول کېده چې کاروونکي
يې پوره خبر وي۔ فلسفي پوښتنه دا ده:
د شمېر وړه ټاکنه څنګه \emph{توجيه}
کولے شو؟
\olpseganchor{OLP-0597-B018}{end}

\olpseganchor{OLP-0597-B019}{start}
د وجداني توجيې يوه هڅه ښايي
له فوق‌العمل څخه مرسته وغواړي۔ فرض
کړئ لومړۍ ټاکنه په \LR{$\nicefrac{1}{2}$}
دقيقه کښې، دويمه په
\LR{$\nicefrac{1}{4}$} دقيقه کښې، \dots،
\LR{$n$}مه ټاکنه په
\LR{$\nicefrac{1}{2^n}$} دقيقه کښې،
\dots\@ کوو۔ نو په \LR{$1$}~دقيقه کښې
به د ټاکنو يوه \LR{$\omega$}-لړۍ
ترسره کړو او د ټاکنې تابعه به
تعريف کړو۔
\olpseganchor{OLP-0597-B019}{end}

\olpseganchor{OLP-0597-B020}{start}
خو دا ډول فکري تجربه په رښتيا څه
راښودلے شي؟ لومړے دا چې د «ټاکلو»
فکر تر ډېره په لفظي مانا اخلي۔
بل دا چې رياضي له مابعدالطبيعي
امکان سره تړي۔
\olpseganchor{OLP-0597-B020}{end}

\olpseganchor{OLP-0597-B021}{start}
تر دې مهمه خبره دا ده چې دا د
\emph{مطلقې} ټاکنې توجيه نه شي
کولے؛ يوازې \emph{شمېر وړې} ټاکنې
ته رسېږي۔ ځکه که \emph{هر} سټ
د ټاکنې تابعه غواړي، نو د «هرې
ترتيبي اوږدوالي فوق‌العمل» ترسره
کولو توان به هم پکار وي۔ په ډاګه
ووايو، دا فکر د خندا وړ دے۔
\olpseganchor{OLP-0597-B021}{end}

\olpseganchor{OLP-0598-B003}{start}
\olfileid{sth}{choice}{justifications}
\olsection{د ټاکنې په اړه دننۍ کتنې}
\olpseganchor{OLP-0598-B003}{end}

\olpseganchor{OLP-0598-B004}{start}
نو پراخه پوښتنه دا ده چې ايا ښه ترتيب،
يا ټاکنه، يا د ټولو سټونو د اندازې
پرتله کېدنه — هر يو يې واخلئ، توپير
نه کوي — توجيه کېدلے شي؟
\olpseganchor{OLP-0598-B004}{end}

\olpseganchor{OLP-0598-B005}{start}
دا يې د \emph{دننۍ} توجيې يوه هڅه
ده۔ مخکې په \olref[z][story]{sec}
کښې مو د مراتبي جوړښت څو اصلونه
وړاندې کړي وو۔ يو يې بيا يادول
ارزي:
\begin{enumerate}
	\item[] \stagesacc. د هر پړاو \LR{$S$} دپاره، او د هغو
	هر ډول سټونو دپاره چې د \LR{$S$} له پړاو
	\emph{مخکې} جوړ شوي وو: په \LR{$S$} پړاو
	کښې داسې سټ جوړېږي چې غړي يې
	عين هماغه سټونه دي۔ په
	پړاو~\LR{$S$} کښې نور هېڅ نه جوړېږي۔
\end{enumerate}
په حقيقت کښې ګڼو ليکوالانو ويلي چې
د ټاکنې اصل د (همدې ته ورته) اصل
په وسيله توجيه کېدلے شي۔ د دې لارې
لنډه تشريح وړاندې کوو۔
\olpseganchor{OLP-0598-B005}{end}

\olpseganchor{OLP-0598-B006}{start}
له يوې کوچنۍ ساده پايلې پيل کوو چې
د ټاکنې يو \emph{بل} همارز اصل
وړاندې کوي:
\olpseganchor{OLP-0598-B006}{end}

\olpseganchor{OLP-0598-B007}{start}
\begin{thm}[په \LR{$\ZF$} کښې]\ollabel{choiceset}
ټاکنه له لاندې اصل سره همارز دے۔
که د \LR{$A$} غړي دوه په دوه بېل
او ناخالي وي، نو داسې \LR{$C$} شته چې
\LR{$C
\cap x$} د هر \LR{$x \in A$} دپاره
يو غړيز سټ وي۔ (داسې \LR{$C$} ته د
\LR{$A$} دپاره {د ټاکنې سټ} وايو۔)
\end{thm}
\olpseganchor{OLP-0598-B007}{end}

\olpseganchor{OLP-0598-B008}{start}
د دې پايلې ثبوت نېغ دے، او لوستونکي
ته يې د تمرين په توګه پرېږدو۔
\olpseganchor{OLP-0598-B008}{end}

\olpseganchor{OLP-0598-B009}{start}
\begin{prob}
\olref[sth][choice][justifications]{choiceset}
ثابت کړئ۔ که ستونزه لرئ، ثبوت په
\mbox{\LR{\latinfont \cite[pp.~242--3]{Potter2004}}} کښې
موندلے شئ۔
\end{prob}
\olpseganchor{OLP-0598-B009}{end}

\olpseganchor{OLP-0598-B010}{start}
بنسټيزه خبره دا ده چې د \LR{$A$} دپاره
د ټاکنې سټ د \LR{$A$} دپاره د ټاکنې
تابعې د ارزښتونو سټ دے، که د ټاکنې سټ د اصلي کورنۍ په اتحاد کښې محدود کړو۔ د ټاکنې له تابعې هر غړے په خپلې ټولګې کښې راځي، او د بېلوالي له مخې د بلې ټولګې له ټاکلي غړي سره نۀ ګډېږي۔ که ټاکنيز سټ راکړل شي، له هر غړي سره د هغه يوازينی ګډ غړے تړو۔ نو د ټاکنې
د توجيې دپاره يوازې د هغه همارزه
بڼه، د ټاکنې سټونو د شتون له مخې،
توجيه کولے شو۔ اوس همدا هڅه کوو۔
\olpseganchor{OLP-0598-B010}{end}

\olpseganchor{OLP-0598-B011}{start}
فرض کړئ د \LR{$A$} غړي دوه په
دوه بېل او ناخالي دي۔ د \stageshier{}
له مخې (وګورئ \olref[z][story]{sec})
\LR{$A$} په کوم پړاو~\LR{$S$} کښې جوړېږي۔
پام وکړئ چې د \LR{$\bigcup A$} ټول
غړي د \LR{$S$} پړاو څخه
مخکې شته۔ اوس د \stagesacc{}
له مخې، د پړاو~\LR{$S$} څخه مخکې
جوړ شويو \emph{هر ډول} سټونو
دپاره داسې سټ جوړېږي چې غړي
يې عين هماغه سټونه دي۔ په بله
وينا، د مخکينيو شته سټونو هره
\emph{ممکنه} ټولګه به په~\LR{$S$}
کښې شتون ولري۔ خو په يقين سره
داسې څيزونه ټاکل \emph{ممکن}
دي چې د~\LR{$A$} دپاره د ټاکنې سټ
ترې جوړ شي؛ دا د \LR{$\bigcup A$}
يو ډېر ځانګړے فرعي سټ دے۔ نو
د دغې ټاکنې د امکان له فرضه داسې سټ شته۔ د پړاو اصل د هرې ټاکل شوې ټولګې د سټ شتون ورکوي۔ له هرې بېلې ټولګې څخه د يوه غړي يوځای ټاکلو امکان دلته يو جلا فرض دے۔
\olpseganchor{OLP-0598-B011}{end}

\olpseganchor{OLP-0598-B012}{start}
دا د دننيو دلايلو پر بنسټ د
ټاکنې د توجيې \emph{ډېره}
لنډه هڅه وه۔ خو د دې فکر د لا
څېړلو دپاره د پوټر
(\mbox{\LR{\latinfont \citeyear[\S14.8]{Potter2004}}})
ښه پراختيا ولولئ۔
\olpseganchor{OLP-0598-B012}{end}

\olpseganchor{OLP-0599-B003}{start}
\olfileid{sth}{choice}{banach}
\olsection{د باناخ--تارسکي پاراډوکس}
\olpseganchor{OLP-0599-B003}{end}

\olpseganchor{OLP-0599-B004}{start}
د ټاکنې د توجيې دپاره، لکه څنګه چې بولوس د تعويض د
توجيې هڅه کړې وه، کېدے شي \emph{بهرنيو} دلايلو ته هم مخه کړو
(وګورئ \olref[replacement][extrinsic]{sec})۔ بالاخره د ټاکنې منل
ډېرې خوښې پايلې لري: د هر دوو اصلي عددونو د پرتله کولو وړتيا؛
د هر سټ د ښه ترتيبولو وړتيا؛ د اصلي عددونو د ترتيبي عددونو د
يوه ځانګړي ډول په توګه ګڼلو وړتيا؛ او داسې نور۔
\olpseganchor{OLP-0599-B004}{end}

\olpseganchor{OLP-0599-B005}{start}
خو کله کله ويل کېږي چې ټاکنه \emph{ناخوښې} پايلې هم لري۔
دا خبره تر ډېره د \mbox{\LR{\latinfont \cite{BanachTarski1924}}} د يوې پايلې له امله کېږي۔
\olpseganchor{OLP-0599-B005}{end}

\olpseganchor{OLP-0599-B006}{start}
\begin{thm}[د باناخ--تارسکي پاراډوکس (په \LR{$\ZFC$} کښې)]
په درې بُعدي فضا کښې د مثبتې شعاع هر ډک توپ په متناهي شمېر ټوټو وېشل کېدے شي؛ بيا دا ټوټې
(په ګرځولو او لېږدولو سره) داسې بيا يوځاے کېدے شي چې
د هماغه توپ دوه نقلونه ترې جوړ شي۔
\end{thm}
\noindent
په لومړي نظر دا لږه حيرانوونکې ده۔ ښکاره ده چې دوه توپونه
د اصلي توپ د حجم \emph{دوه چنده} حجم لري۔ خو صُلب حرکتونه ---
ګرځول او لېږدول --- حجم نه بدلوي۔ نو داسې برېښي لکه
باناخ--تارسکي چې موږ ته له هېڅه نوې ماده راپيدا کوي۔
\olpseganchor{OLP-0599-B006}{end}

\olpseganchor{OLP-0599-B007}{start}
تر دې هم عجيبه خبره شته۔\footnote{وګورئ \mbox{\LR{\latinfont \citet[Theorem 3.12]{Wagon2016}}}۔}
ورته استدلال ښيي چې يوه وړه د نخود دانه په متناهي شمېر ټوټو
پرې کېدے شي، او بيا همدغه ټوټې (په ګرځولو او لېږدولو سره)
د بېګ بېن د شکل او اندازې يو څيز جوړولے شي۔
\olpseganchor{OLP-0599-B007}{end}

\olpseganchor{OLP-0599-B008}{start}
خو له دې څخه هېڅ يوه خبره يوازې په \LR{$\ZF$} کښې نه ثابتېږي۔\footnote{که څه هم
باناخ--تارسکي د ټاکنې له اصله په ټينګه کمزورو اصلونو سره
ثابتېږي؛ وګورئ \mbox{\LR{\latinfont \citet[303]{Wagon2016}}}۔} نو له يوې پرېکړې سره
مخامخ يو: يا ټاکنه رد کړو، يا له دې «پاراډوکس» سره ژوند ومنو۔
\olpseganchor{OLP-0599-B008}{end}

\olpseganchor{OLP-0599-B009}{start}
موږ وايو چې له دې «پاراډوکس» سره ژوند منل په کار دي۔ په
حقيقت کښې موږ ته دا ډېر پاراډوکس نه ښکاري۔ په ځانګړې توګه
نه وينو چې دا د لاندې پايلو په پرتله ولې زيات يا کم پاراډوکس
وګڼل شي:\footnote{\mbox{\LR{\latinfont \citet[276--7]{Potter2004}}}،
\mbox{\LR{\latinfont \citet[16]{Weston2003}}}، او \mbox{\LR{\latinfont \citet[31, 308--9]{Wagon2016}}}
له نورو بېلګو سره ورته ټکي بيانوي۔}
\begin{enumerate}
	\item په \LR{$(0,1)$} وقفه کښې د ټکو شمېر د \LR{$\Real$} د ټکو هومره دے۔
	\\\emph{ثبوت}: \LR{$\tan(\pi(r-\nicefrac{1}{2}))$} په پام کښې ونيسئ۔ دا تابع په ټوله وقفه کښې په کلکه زياتېږي، او د دواړو سرونو حدونه يې منفي او مثبت بې‌نهايت دي۔ نو له وقفې څخه ټولې حقيقي کرښې ته دوه اړخيزه تابعه ده۔
	\item په يوه کرښه کښې د ټکو شمېر د يوه مربع د ټکو هومره دے۔
	\\وګورئ \olref[his][set][pathology]{sec} او \olref[his][set][cantorplane]{sec}۔
	\item فضا ډکوونکې منحنيانې شته۔
	\\وګورئ \olref[his][set][pathology]{sec} او \olref[his][set][hilbertcurve]{sec}۔
\end{enumerate}
له دې درېو پايلو څخه يوه هم ټاکنې ته اړتيا نه لري۔ اوس
يې موږ د رياضياتو حيرانوونکې، ښکلې برخې ګڼو۔ ښايي د
باناخ--تارسکي پاراډوکس په اړه هم همداسې نظر ولرو۔
\olpseganchor{OLP-0599-B009}{end}


\olpseganchor{OLP-0599-B010}{start}
البته دلته يوه تخنيکي خبره يادول پکار دي، خو په ارامه
ذهن پوهېدل پرې بس دي۔ صُلب حرکتونه حجم ساتي۔ نو هغه
پنځه\footnote{موږ پاراډوکس د «متناهي شمېر ټوټو» په ژبه بيان کړ۔
په حقيقت کښې \mbox{\LR{\latinfont \citet{Robinson1947}}} ثابته کړه چې وېش په
\emph{پنځو} ټوټو کېدے شي (خو تر پنځو په کمو نه)۔ د ثبوت
دپاره \mbox{\LR{\latinfont \citet[pp.~66--7]{Wagon2016}}} وګورئ۔} ټوټې چې توپ
ورباندې وېشل کېږي، ټولې \emph{د اندازې وړ} نه شي کېدے۔
نو په لنډو ټکو، دې هرې ټوټې ته په جلا توګه حجم ټاکل
معنا نه لري۔ دا ټوټې د ټکو داسې ناليدل کېدونکې «نامتناهي
خورې ورې ټولګې» وګڼئ۔ کېدے شي د دغسې «نامتناهي
خورو ورو» سټونو تصور عجيب وبرېښي۔ خو د هغو شتون له
هغه غوښتنې څخه راځي چې په \stagesacc{} کښې راغلې: د
مخکښې موجودو سټونو \emph{ټولې ممکنې} ټولګې جوړې کړئ۔
\olpseganchor{OLP-0599-B010}{end}

\olpseganchor{OLP-0599-B011}{start}
که بيا هم دا خبرې تاسو نه قانع کوي، د باناخ--تارسکي
پاراډوکس د منلو په ګټه دا وروستے (بهرنے) دليل واورئ۔ له دې
سمدستي د رياضياتو تر ټولو ښه ټوکه جوړېږي:
\begin{enumerate}
	\item[] \emph{پوښتنه}. د «باناخ--تارسکي» د تورو بل ترتيب څه دے؟
	\item[] \emph{ځواب}. «باناخ--تارسکي باناخ--تارسکي»۔
\end{enumerate}
\olpseganchor{OLP-0599-B011}{end}

\olpseganchor{OLP-0600-B003}{start}
\olfileid{sth}{choice}{vitali}
\olsection{ضميمه: د ويتالي پاراډوکس}
\olpseganchor{OLP-0600-B003}{end}

\olpseganchor{OLP-0600-B004}{start}
د باناخ--تارسکي د جوړښت د منلو يا نه منلو د سمې ارزونې
دپاره بايد د هغه \emph{ثبوت} وګورو۔ له بده مرغه، دا تر هغه
ډېر جبر غواړي چې دلته يې بيانولے شو۔ خو د لاندې پايلې
په رڼا کښې څو لنډې خبرې کولے شو چې ښايي د باناخ--تارسکي
ثبوت روښانه کړي۔\footnote{د ډېر بشپړ بحث دپاره
\mbox{\LR{\latinfont \cite{Weston2003}}} يا \mbox{\LR{\latinfont \cite{Wagon2016}}} وګورئ۔}
\olpseganchor{OLP-0600-B004}{end}

\olpseganchor{OLP-0600-B005}{start}
\begin{thm}[د ويتالي پاراډوکس (په \LR{$\ZFC$} کښې)]\ollabel{vitaliparadox}
هره دايره د شمېر وړ شمېر ټوټو ته وېشل کېدے شي؛ بيا دا
ټوټې (په ګرځولو او لېږدولو سره) داسې يوځاے کېدے شي
چې د هماغې دايرې دوه نقلونه ترې جوړ شي۔
\end{thm}
\olpseganchor{OLP-0600-B005}{end}

\olpseganchor{OLP-0600-B006}{start}
د ويتالي پاراډوکس ثبوت د باناخ--تارسکي د پاراډوکس تر
ثبوت ډېر اسانه دے۔ موږ دا د «ويتالي پاراډوکس» بولو، ځکه
چې د ويتالي د \mbox{\LR{\latinfont \citeyear{Vitali1905}}} د نه اندازېدونکي سټ
له جوړښت څخه راځي۔ خو د دواړو ثبوتونو سټي اړخونه ډېر
ورته دي۔ بنسټيز توپير يوازې دا دے چې باناخ--تارسکي
\emph{متناهي} وېش کاروي، او د ويتالي پاراډوکس
\emph{د شمېر وړ نامتناهي} وېش۔ د \mbox{\LR{\latinfont \citet{Weston2003}}}
د بحث له مخې، د ويتالي مثال ښيي چې د داسې هندسي وېش
دپاره دايره بس ده، حال دا چې باناخ--تارسکي د درې‌بعدي
توپ پر صُلبو حرکتونو ولاړ دے۔
\olpseganchor{OLP-0600-B006}{end}

\olpseganchor{OLP-0600-B007}{start}
د ويتالي پاراډوکس له دوه-بُعدي شکل، يعنې دايرې، سره
اړيکه لري۔ نو په مستوي \LR{$\Real^2$} کښې کار کوو۔
\LR{$\rotationsgroup$} دې د مبدأ شاوخوا د ټکو د هغو
(د ساعت د ستنې په لوري) ګرځونو سټ وي چې د راډيان
زاويې يې په \LR{$[0,2\pi)$} کښې د بشپړ چورلښت
\emph{ناطقي برخې} دي۔
د \LR{$\rotationsgroup$} په اړه څو جبري حقيقتونه وړاندې کوو
(که د پايلې وينا اوس نه پوهېږئ، ثبوت به يې معنا روښانه کړي):
\olpseganchor{OLP-0600-B007}{end}

\olpseganchor{OLP-0600-B008}{start}
\begin{lem}\ollabel{rotationsgroupabelian}
\LR{$\rotationsgroup$} د تابعو د ترکيب له مخې يو تبديلېدونکے {ګروپ} جوړوي۔
\end{lem}
\olpseganchor{OLP-0600-B008}{end}

\olpseganchor{OLP-0600-B009}{start}
\begin{proof}
د \LR{$0_{\rotationsgroup}$} نښه د \LR{$0$} راډيانو د ګرځون دپاره
کاروو؛ دا د
\LR{$\rotationsgroup$} دپاره د هويت غړے دے، ځکه چې
\LR{$\comp{0_{\rotationsgroup}}{\rho} = \comp{\rho}{0_{\rotationsgroup}} =
\rho$} د هر \LR{$\rho \in \rotationsgroup$} دپاره۔
\olpseganchor{OLP-0600-B009}{end}

\olpseganchor{OLP-0600-B010}{start}
هر غړے معکوس لري۔ که \LR{$\rho \in \rotationsgroup$} په
\LR{$r$} راډيانو ګرځوي، نو \LR{$\rho^{-1} \in \rotationsgroup$}
په \LR{$2\pi - r$} راډيانو ګرځوي، چې بشپړ چورلښت د صفر
زاويې په توګه اخلو۔ دواړه زاويې د بشپړ چورلښت ناطقي
برخې دي۔ ځکه \LR{$\rho \circ \rho^{-1} = 0_\rotationsgroup$}۔
\olpseganchor{OLP-0600-B010}{end}

\olpseganchor{OLP-0600-B011}{start}
د دوو ګرځونو ترکيب د هغوی زاويې جمع کوي؛ که مجموعه
يو بشپړ چورلښت ته ورسېږي، له هغې يو بشپړ چورلښت کموو۔
په دواړو صورتونو کښې پاتې زاويه بيا د بشپړ چورلښت ناطقي
برخه ده۔ ترکيب تجميعي دے: \LR{$\comp{\rho}{(\comp{\sigma}{\tau})} =
\comp{(\comp{\rho}{\sigma})}{\tau}$} د هر \LR{$\rho, \sigma, \tau \in
\rotationsgroup$} دپاره۔
\olpseganchor{OLP-0600-B011}{end}

\olpseganchor{OLP-0600-B012}{start}
ترکيب تبديلېدونکے دے: \LR{$\comp{\rho}{\sigma} =
\comp{\sigma}{\rho}$} د هر \LR{$\rho, \sigma \in \rotationsgroup$} دپاره۔
\end{proof}
\olpseganchor{OLP-0600-B012}{end}

\olpseganchor{OLP-0600-B013}{start}
په حقيقت کښې موږ خپل ګروپ \LR{$\rotationsgroup$} په دوو برخو
وېشلے شو، او بيا له هرې برخې ټول ګروپ بېرته ترلاسه کولے شو:
\olpseganchor{OLP-0600-B013}{end}

\olpseganchor{OLP-0600-B014}{start}
\begin{lem}\ollabel{disjointgroup}
\LR{$\rotationsgroup$} په دوو بېلو سټونو،
\LR{$\rotationsgroup_{1}$} او \LR{$\rotationsgroup_{2}$}، وېشل کېدے شي
چې د هرې برخې د ګرځونو دوه ځله ترکيب ټول
\LR{$\rotationsgroup$} په دوه اړخيز ډول ورکوي۔
\end{lem}
\olpseganchor{OLP-0600-B014}{end}

\olpseganchor{OLP-0600-B015}{start}
\begin{proof}
په \LR{$\rotationsgroup_{1}$} کښې دې د \LR{$[0, \pi)$} هغه
ګرځونونه وي چې زاويې يې د بشپړ چورلښت ناطقي برخې دي؛ او
\LR{$\rotationsgroup_{2} = \rotationsgroup
\setminus \rotationsgroup_1$} ونيسئ۔ د لومړني جبر له مخې،
\LR{$\Setabs{\comp{\rho}{\rho}}{\rho \in \rotationsgroup_1} =
\rotationsgroup$}۔ د ورکړې زاويې نيمه په لومړۍ برخه
کښې يوازينۍ مخکينۍ زاويه ده؛ د دويمې برخې لپاره هماغې
نيمې ته نيم چورلښت ورزيات کړئ۔ ځکه د هرې برخې دوه
چنده کول د ټول ګروپ سره دوه اړخيز سمون جوړوي۔
\end{proof}
\olpseganchor{OLP-0600-B015}{end}

\olpseganchor{OLP-0600-B016}{start}
موږ به د ګروپونو په اړه همدا خبره د \olref{vitaliparadox}
د ثابتولو دپاره وکاروو۔ \LR{$\onesphere$} دې واحده دايره وي،
يعنې د مستوي له مبدأ څخه عيناً~\LR{$1$} واحد لرې د ټکو سټ،
يعنې \LR{$\Setabs{\tuple{r,s} \in \Real^2}{\sqrt{r^2+s^2}=1}$}۔
د \LR{$\onesphere$} د وېشلو دپاره په \LR{$\onesphere$} دا اړيکه
په پام کښې نيسو:
\[
	r \sim s \emph{ هله او يوازې هله چې }(\exists \rho \in \rotationsgroup)\rho(r) = s.
\]
يعنې د \LR{$\onesphere$} دوه ټکي د دې اړيکې له مخې سره
تړلي دي که او يوازې که يو له بله د مبدأ شاوخوا په
د بشپړ چورلښت د يوې ناطقي برخې په اندازه ګرځون ترلاسه کېدے شي۔ نو تمه کېږي چې:
\olpseganchor{OLP-0600-B016}{end}

\olpseganchor{OLP-0600-B017}{start}
\begin{lem}
\LR{$\sim$} د همارزۍ اړيکه ده۔
\end{lem}
\olpseganchor{OLP-0600-B017}{end}

\olpseganchor{OLP-0600-B018}{start}
\begin{proof}
د \olref{rotationsgroupabelian} له مخې د هويت ګرځون
انعکاسي والى، معکوس تقارني والى، او ترکيب انتقالي
والى ورکوي۔
\end{proof}
\olpseganchor{OLP-0600-B018}{end}

\olpseganchor{OLP-0600-B019}{start}
اوس د ټاکنې اصل کاروو، څو داسې سټ \LR{$C$} واخلو چې د
\LR{$\onesphere$} د \LR{$\sim$} اړيکې د هرې همارزۍ طبقې څخه
عين يو غړے ولري۔ يعنې د همارزۍ د طبقو پر سټ د ټاکنې
تابعه \LR{$f$} په پام کښې نيسو،\footnote{ځکه چې
\LR{$\rotationsgroup$} د شمېر وړ دے، د \LR{$E$} هر
غړے هم د شمېر وړ دے۔ او ځکه چې
\LR{$\onesphere$} د شمېر نۀ وړ دے، د
\olref{vitalicover} او
\olref[card-arithmetic][simp]{kappaunionkappasize}
له مخې \LR{$E$} هم د شمېر نۀ وړ دے۔ نو دا د
\emph{نه}شمېر وړې ټاکنې کارونه ده۔}
\[
	E = \Setabs{\equivrep{r}{\sim}}{r \in \onesphere},
\]
او \LR{$C = \ran{f}$} ونيسئ۔ د هر ګرځون
\LR{$\rho \in \rotationsgroup$} دپاره، په \LR{$\funimage{\rho}{C}$}
کښې هغه ټکي دي چې د \LR{$\rho$} ګرځون د \LR{$C$} پر هر ټکي
له کارولو ترلاسه کېږي۔ لاندې دوه پايلې ښيي چې دا سټونه
دايره بشپړه او بې له ګډو غړو پوښي:
\olpseganchor{OLP-0600-B019}{end}

\olpseganchor{OLP-0600-B020}{start}
\begin{lem}\ollabel{vitalicover}
\LR{$\onesphere = \bigcup_{\rho \in \rotationsgroup} \funimage{\rho}{C}$}.
\end{lem}
\olpseganchor{OLP-0600-B020}{end}

\olpseganchor{OLP-0600-B021}{start}
\begin{proof}
\LR{$s \in \onesphere$} ثابت ونيسئ؛ داسې \LR{$r \in C$} شته چې
\LR{$r \in \equivrep{s}{\sim}$}، يعنې \LR{$r \sim s$}، يعنې
\LR{$\rho(r) = s$} د کوم \LR{$\rho \in \rotationsgroup$} دپاره۔
\end{proof}
\olpseganchor{OLP-0600-B021}{end}

\olpseganchor{OLP-0600-B022}{start}
\begin{lem}\ollabel{vitalinooverlap}
که \LR{$\rho_1 \neq \rho_2$}، نو
\LR{$\funimage{\rho_1}{C} \cap \funimage{\rho_2}{C} = \emptyset$}۔
\end{lem}
\olpseganchor{OLP-0600-B022}{end}

\olpseganchor{OLP-0600-B023}{start}
\begin{proof}
فرض کړئ \LR{$s \in \funimage{\rho_{1}}{C} \cap \funimage{\rho_{2}}{C}$}۔
نو \LR{$s = \rho_{1}(r_{1}) = \rho_{2}(r_{2})$} د ځينو
\LR{$r_{1}, r_{2} \in C$} دپاره۔ له دې
\LR{$\rho^{-1}_2(\rho_1(r_1)) = r_2$}، او
\LR{$\comp{\rho_1}{\rho^{-1}_2} \in \rotationsgroup$}؛ نو
\LR{$r_{1} \sim r_{2}$}۔ ځکه \LR{$r_1 = r_2$}، چې \LR{$C$} د
\LR{$\sim$} د هرې همارزۍ طبقې څخه عين يو غړے ټاکي۔
نو \LR{$s = \rho_1(r_1) = \rho_2(r_1)$}، او له دې
\LR{$\rho_1 = \rho_2$}۔
\end{proof}
\olpseganchor{OLP-0600-B023}{end}

\olpseganchor{OLP-0600-B024}{start}
اوس خپل مخکيني جبري حقيقتونه په دايره کاروو:
\olpseganchor{OLP-0600-B024}{end}

\olpseganchor{OLP-0600-B025}{start}
\begin{lem}\ollabel{pseudobanachtarski}
\LR{$\onesphere$} په دوو بېلو سټونو، \LR{$D_{1}$} او \LR{$D_{2}$}،
وېشل کېدے شي؛ \LR{$D_{1}$} بيا د شمېر وړ شمېر سټونو ته
وېشل کېدے شي چې په ګرځولو يې د \LR{$\onesphere$} يو نقل
جوړېږي (او د \LR{$D_{2}$} دپاره هم همداسې)۔
\end{lem}
\olpseganchor{OLP-0600-B025}{end}

\olpseganchor{OLP-0600-B026}{start}
\begin{proof}
د \olref{disjointgroup} څخه \LR{$\rotationsgroup_{1}$} او
\LR{$\rotationsgroup_{2}$} وکاروئ، او ونيسئ:
\begin{align*}
	D_{1} &= \bigcup_{\rho \in \rotationsgroup_1} \funimage{\rho}{C} & 
	D_{2} &= \bigcup_{\rho \in \rotationsgroup_2} \funimage{\rho}{C}
\end{align*}
دا د \olref{vitalicover} له مخې د \LR{$\onesphere$} وېش دے،
او \LR{$D_1$} او \LR{$D_2$} د \olref{vitalinooverlap} له مخې بېل دي۔
د جوړښت له مخې \LR{$D_1$} د شمېر وړ شمېر سټونو ته وېشل
کېدے شي: \LR{$\funimage{\rho}{C}$} د هر
\LR{$\rho \in R_1$} دپاره۔ دا سټونه په ګرځولو سره د
\LR{$\onesphere$} يو نقل جوړوي، ځکه چې د
\olref{disjointgroup} او \olref{vitalicover} له مخې
\LR{$\onesphere = \bigcup_{\rho \in
\rotationsgroup}\funimage{\rho}{C} = \bigcup_{\rho \in
\rotationsgroup_1}\funimage{(\comp{\rho}{\rho})}{C}$}۔
د \LR{$D_2$} دپاره هم همدا استدلال چلېږي۔ \end{proof}\noindent
له دې سمدستي د ويتالي پاراډوکس راځي۔ ځکه د
\LR{$\onesphere$} څخه د \LR{$\onesphere$} \emph{دوه} نقلونه
جوړوو: په د شمېر وړ شمېر ټوټو (بېلابېلو
\LR{$\funimage{\rho}{C}$}) يې وېشو، بيا يې په صُلب
حرکتونو خوځوو: د \LR{$D_1$} هره ټوټه په خپل ټاکلي ګرځون اړوو۔ د \LR{$D_2$}
هره ټوټه په خپل ټاکلي ګرځون اړوو او بيا يې په يوه ګډ
لېږد له لومړۍ دايرې بېلوو۔
\olpseganchor{OLP-0600-B026}{end}

\olpseganchor{OLP-0600-B027}{start}
د ثبوت تګلاره بيا وګورو۔ لومړے مو په مستوي کښې د
ګرځونو د ګروپ ځينې جبري حقيقتونه وکارول۔ په دې ګروپ
مو \LR{$\onesphere$} د همارزۍ په طبقو ووېشه۔ بيا مو د
ټاکنې په کارولو له هرې طبقې يو غړے وټاکه او
«پاراډوکس» ته ورسېدو۔
\olpseganchor{OLP-0600-B027}{end}

\olpseganchor{OLP-0600-B028}{start}
د باناخ--تارسکي د ثبوت دپاره هم عين همدا تګلاره
کاروو۔ لوی توپير دا دے چې د باناخ--تارسکي د ثبوت
جبري حقيقتونه د ويتالي د پاراډوکس تر جبري حقيقتونو
ډېر پېچلي دي۔ خو دا جبري حقيقتونه له ټاکنې سره
اړيکه نه لري۔ په لنډه توګه يې بيانوو۔
\olpseganchor{OLP-0600-B028}{end}

\olpseganchor{OLP-0600-B029}{start}
د باناخ--تارسکي د ثبوت دپاره د \olref{disjointgroup}
په څېر لومړے د دوو ازادو مولدونو ګروپ په پام کښې نيسو؛ دا په
څلورو ټوټو وېشل کېدے شي، چې په تصوري ډول يې
«خوځولو» سره د ټول ګروپ دوه نقلونه ترلاسه کوو۔\footnote{د
\emph{څلورو} ټوټو بسيا کېدل د \mbox{\LR{\latinfont \cite{Robinson1947}}}
پايله ده۔ د نوي ثبوت دپاره \mbox{\LR{\latinfont \citet[Theorem 5.2]{Wagon2016}}}
وګورئ۔ د ګروپ د ټوټو د «خوځولو» په بيان کښې د
\mbox{\LR{\latinfont \citet[p.~3]{Weston2003}}} پيروي کوو۔} بيا ښيو چې
د \LR{$\Real^3$} د مبدأ شاوخوا دوه ځانګړي ګرځونونه د
ګرځونو يو ازاد ګروپ، \LR{$F$}، جوړولے شي۔\footnote{وګورئ
\mbox{\LR{\latinfont \citet[Theorem 2.1]{Wagon2016}}}۔} (تر اوسه ټاکنه نشته۔)
اوس د کرې د سطحې دوه ټکي «ورته» بولو که او يوازې
که يو يې په~\LR{$F$} کښې په کوم ګرځون له بله ترلاسه کېږي۔ لومړے هغه
ټکي وباسو چې د هويت پرته کوم ګرځون يې پر خپل ځاے
ساتي۔ هر داسې ګرځون د خپلې محوري کرښې پر سر دوه ثابت
ټکي لري، او د دغو ګرځونو د ثابتو ټکو شمېر وړ اتحاد هم شمېر وړ دے۔
پاتې سطحه د ګروپ له عمل لاندې تړلې ده او هلته هر ټکی
يوازې د هويت په وسيله پر خپل ځاے پاتې کېږي۔
بيا له «ورته» ټکو د همارزۍ د هرې طبقې څخه د عين يو
ټکي د ټاکلو دپاره \emph{ټاکنه کاروو}۔ د
\olref{pseudobanachtarski} په څېر د \LR{$F$} وېش د کرې
پر پاتې سطحه کاروو او هغه په څلورو ټوټو وېشو۔
د هرې طبقې ټکي له ګروپ سره دوه اړخيز سمون لري؛ نو د
ګروپ د ټوټو ګرځول د پاتې سطحې دوه نقلونه ورکوي۔
د هاوسدورف اصلي استدلال هم د ثابتو ټکو شمېر وړ سټ
بېلوي \mbox{\LR{\latinfont \citep[p.~432]{Hausdorff1914}}}۔
\olpseganchor{OLP-0600-B029}{end}

\olpseganchor{OLP-0600-B030}{start}
\begin{thm}[د هاوسدورف پاراډوکس (په \LR{$\ZFC$} کښې)]
د مثبتې شعاع د هرې کرې له سطحې څخه يو شمېر وړ سټ
اېستل کېدے شي، داسې چې پاتې سطحه په متناهي شمېر ټوټو
ووېشل شي او د صُلبو حرکتونو له لارې ترې د هماغې پاتې
سطحې دوه بېل نقلونه جوړ شي۔
\end{thm}
\olpseganchor{OLP-0600-B030}{end}

\olpseganchor{OLP-0600-B031}{start}
د باناخ--تارسکي بشپړې قضيې دپاره (چې يوازې د کرې
سطحه نه، بلکې دننه برخه يې هم رانغاړي) څو نورې جبري
لارې پکار دي، په تېره د ثابتو ټکو شمېر وړ سټ بېرته
ورګډول او له سطحې نه د ډک توپ دننه ته تلل۔ د ازاد ګروپ
هماغه تجزيه د دې پراختيا بنسټ برابروي۔
\begin{explain}
  د ازاد ګروپ تجزيه د باناخ--تارسکي د ثبوت جبري برخه
  برابروي۔ دغه تجزيه هغه مهال هندسي نتيجه ورکوي چې
  ګروپ د \LR{$\Real^3$} پر ټکو د ګرځونو او لېږدونو په
  وسيله عمل وکړي۔ په دې توګه د ټوټو د بيا راغونډولو
  امکان په \LR{$\Real^3$} کښې د صُلبو حرکتونو د ګروپ له جوړښت څخه راځي،
  لکه وېسټن چې تشريح کوي۔ \mbox{\LR{\latinfont \cite[p.~16]{Weston2003}}}
\end{explain}
يعنې متناهي وېش (لکه باناخ--تارسکي) او د شمېر وړ
نامتناهي وېش (لکه د ويتالي پاراډوکس) د دوه-بُعدي
او درې-بُعدي فضا د ګروپونو په ځينو حقيقتونو پورې تړلي دي۔
\olpseganchor{OLP-0600-B031}{end}

\olpseganchor{OLP-0600-B032}{start}
منو چې دا وروستۍ خبره د \olref[banach]{sec} د پای ټوکه
لږه خرابوي۔ ځکه چې ليکل په دوه بُعدونو کښې وي، نو که
«باناخ--تارسکي» په «باناخ--تارسکي باناخ--تارسکي»
بدلوو، بايد په \emph{د شمېر وړ نامتناهي} شمېر ټوټو
ووېشل شي۔ د ټوکې د سمولو دپاره بايد په درې بُعدونو
کښې وليکل شي۔ دا د لوستونکي د تمرين دپاره پرېږدو۔
\olpseganchor{OLP-0600-B032}{end}

\olpseganchor{OLP-0600-B033}{start}
يوه وروستۍ خبره۔ په \olref[banach]{sec} کښې مو وويل
چې د کرې ترلاسه شوې «ټوټې» ټولې \emph{د اندازې وړ}
نه شي کېدے؛ دا ناليدل کېدونکې «نامتناهي خورې ورې
ټولګې» دي۔ د \olref{pseudobanachtarski} د دايرې په وېش کښې هم
ځينې ټوټې د ګرځون‌ساتونکې اندازې په ساحه کښې نۀ راځي۔ دا خبره تشريح
ارزي۔
\olpseganchor{OLP-0600-B033}{end}

\olpseganchor{OLP-0600-B034}{start}
بيا لږ مخينه بيانوو (خو دا \emph{يوازې} يوه لنډه نقشه
ده؛ د \emph{اندازې} په اړه د درسي کتاب برخه هم کتلے شئ)۔
د سټ \LR{$X$} دپاره د اندازې تعريف دا دے: ارزښت
\LR{$\mu(E) \in [0,\infty]$} د هر \LR{$E$} دپاره په کوم
«\LR{$\sigma$}-الجبر» کښې، چې پر \LR{$X$} وي، ټاکل کېږي۔
دلته تفصيل اړين نه دے،
خو تابعه \LR{$\mu$} بايد د شمېر وړ جمعيت اصل ومني: د
بېلو سټونو د شمېر وړ اتحاد اندازه د هغو د بېلو اندازو
مجموعه ده، يعنې
\LR{$\mu(\bigcup_{n < \omega} X_n) = \sum_{n < \omega}\mu(X_n)$}
کله چې \LR{$X_n$} سټونه دوه په دوه بېل وي۔ يو سټ «نه
اندازېدونکے» بلل دا دي چې سټ د ټاکلې اندازې د
تعريف په ساحه کښې نۀ وي۔ اوس خپل مخکينے \LR{$\rotationsgroup$} کاروو:
\olpseganchor{OLP-0600-B034}{end}

\olpseganchor{OLP-0600-B035}{start}
\begin{cor}[ويتالي]
\LR{$\mu$} دې داسې اندازه وي چې \LR{$\mu(\onesphere) = 1$}،
او د تعريف ساحه يې د دايرې د ګرځونو لاندې تړلې وي۔
که په دې ساحه کښې \LR{$X$} او \LR{$Y$} هم‌نهشت وي، نو
\LR{$\mu(X) = \mu(Y)$} دې وي۔
بيا \LR{$\funimage{\rho}{C}$} د
هر \LR{$\rho \in \rotationsgroup$} دپاره نه اندازېدونکے دے۔
\end{cor}
\olpseganchor{OLP-0600-B035}{end}

\olpseganchor{OLP-0600-B036}{start}
\begin{proof}
د تناقض له لارې فرض کړئ چې داسې نه ده۔ نو
\LR{$\mu(\funimage{\sigma}{C}) = r$} د کوم
\LR{$\sigma \in \rotationsgroup$} او کوم \LR{$r \in \Real$}
دپاره ونيسئ۔ د هر \LR{$\rho \in \rotationsgroup$} دپاره
\LR{$\funimage{\rho}{C}$} او \LR{$\funimage{\sigma}{C}$}
هم‌نهشت دي؛ نو \LR{$\mu(\funimage{\rho}{C}) = r$}
د هر \LR{$\rho \in \rotationsgroup$} دپاره۔ د \olref{vitalicover} او
\olref{vitalinooverlap} له مخې
\LR{$\onesphere = \bigcup_{\rho \in \rotationsgroup}\funimage{\rho}{C}$}
د دوه په دوه بېلو سټونو د شمېر وړ اتحاد دے۔ نو د شمېر
وړ جمعيت له مخې \LR{$\mu(\onesphere) = 1$} د هر
\LR{$\funimage{\rho}{C}$} د اندازو مجموعه ده، يعنې
\[
	1 = \mu(\onesphere) = \sum_{\rho \in \rotationsgroup}\mu(\funimage{\rho}{C}) = \sum_{\rho \in \rotationsgroup}r
\]
د اندازې ارزښت منفي نۀ دے او د ټولې دايرې تر اندازې
لوی نۀ دے۔ که \LR{$r = 0$} وي، نو \LR{$\sum_{\rho \in \rotationsgroup}r = 0$}؛
او که \LR{$r > 0$} وي، نو
\LR{$\sum_{\rho \in \rotationsgroup}r = \infty$}۔
\end{proof}
\olpseganchor{OLP-0600-B036}{end}

\olpseganchor{OLP-0601-B004}{start}
\olpart{mth}{لارې چارې}
\olpseganchor{OLP-0601-B004}{end}

\olpseganchor{OLP-0601-B005}{start}
\begin{editorial}
  په دې برخه کښې عمومي او د لارو چارو اړوند مواد راځي،
  په ځانګړې توګه د ثبوت د بېلابېلو طريقو تشريح، چې هغه
  زده‌کوونکے به ورسره نااشنا وي چې رياضيات يې اصلي
  څانګه نه ده۔ دلته د ثبوتونو د ليکلو طريقې او استقرا څېړو۔
\end{editorial}
\olpseganchor{OLP-0601-B005}{end}

\olpseganchor{OLP-0602-B004}{start}
\olchapter{mth}{prf}{ثبوتونه}
\olpseganchor{OLP-0602-B004}{end}

\olpseganchor{OLP-0603-B004}{start}
\olfileid{mth}{prf}{int}
\olpseganchor{OLP-0603-B004}{end}

\olpseganchor{OLP-0603-B005}{start}
\olsection{سريزه}
\olpseganchor{OLP-0603-B005}{end}

\olpseganchor{OLP-0603-B006}{start}
د منطق د لومړنيو درسونو د تجربې له مخې ښايي تاسو د
اشتقاق له نظام سره بلد ياست --- غالباً د طبيعي
استنتاج يا د فېچ په بڼه اشتقاق نظام، يا د ثبوت
د ونې نظام۔ ښايي په ياد مو وي چې په دغو نظامونو کښې
مو ثبوتونه جوړول: يا مو فارمول ثابتوله، يا مو
ښودله چې يو ټاکلے استدلال معتبر دے۔ د دې دپاره به مو
د نظام قاعدې تر هغې کارولې چې غوښتل شوې پايلې ته
ورسېدئ۔ د منطق \emph{په اړه} په استدلال کښې هم څيزونه
ثابتوو، خو په ډېرو حالاتو کښې د اشتقاق نظام
نه کاروو۔ دلته ډېر ثبوتونه په عادي نثر ليکو او کله
کله ورسره نښې کاروو۔ په دې پښتو متن کښې استدلال د
پښتو په جملو او د رياضياتو په نښو وړاندې کېږي۔
د داسې ثبوتونو په جوړولو کښې کېدے شي په لومړي سر
کښې حيران شئ: بې له اشتقاق نظامه يو څيز
څنګه ثابت کړم؟ له کومه ځايه پيل وکړم؟ څنګه پوه شم
چې زما ثبوت سم دے؟
\olpseganchor{OLP-0603-B006}{end}

\olpseganchor{OLP-0603-B007}{start}
د ثبوت له هڅې مخکې، دا پېژندل مهم دي چې ثبوت څه
دے او څنګه جوړېږي۔ لکه له نومه يې چې ښکاري،
\emph{ثبوت} د دې د ښودلو دپاره وي چې يو څيز رښتيا
دے۔ دا د خبرو اترو په بڼه هم تصورولے شئ: څوک
له تاسو پوښتنه کوي چې ايا کومه خبره رښتيا ده، د
بېلګې په توګه ايا له دوه پرته هر اولي عدد طاق دے۔
يوازې «هو» ويل بس نه دي؛ ښايي هغه وغواړي پوه شي
چې \emph{ولې}۔ په دې حالت کښې به ثبوت ورکوئ۔
\olpseganchor{OLP-0603-B007}{end}

\olpseganchor{OLP-0603-B008}{start}
په ورځنيو خبرو کښې ښايي ځواب ته لنډه اشاره يا
نيمګړے ځواب بس وي۔ خو په منطق او رياضياتو کښې
کره ثبوت غواړو: غواړو وښيو چې خبره له \emph{هر}
شک څخه ورهاخوا رښتيا ده۔ معنا يې دا ده چې د
ثبوت هر ګام بايد توجيه ولري، او توجيه يې بايد
سمه او قانع کوونکې وي (يعنې کارېدونکې فرضيه
بايد په هغه قضيه کښې په رښتيا فرض شوې وي چې
تاسو يې ثابتوئ؛ تعريفونه بايد سم وکارول شي؛
هغه استنتاجونه چې دليل پرې ولاړ وي بايد سم وي؛
او داسې نور)۔
\olpseganchor{OLP-0603-B008}{end}

\olpseganchor{OLP-0603-B009}{start}
عموماً موږ يوه وينا ثابتوو۔ دغو ويناوو ته بېلابېل
نومونه ورکوو: بيانونه، قضيې، لمونه، يا فرعي پايلې۔
بيان داسې بنسټيزه وينا ده چې ثبوت يې ارزښت لري:
د ثبتولو هومره مهمه وي، خو ښايي ډېره ژوره نه وي
يا ډېر ځله ونه کارول شي۔ قضيه يو مهم او د پام وړ
بيان دے۔ ثبوت يې ډېری وخت په څو پړاوونو وېشل کېږي،
او کله کله د هغه چا په نوم يادېږي چې لومړے يې
ثابته کړې وي (لکه د کانتور قضيه او د
لوېنهايم--سکولم قضيه)، يا د هغه حقيقت په نوم چې
اړيکه ورسره لري (لکه د بشپړتيا قضيه)۔ لمه هغه
بيان يا قضيه ده چې د لا مهمې پايلې په ثبوت کښې
کارېږي۔ د پام وړ خبره دا ده چې ځينې لمونه
پخپله هم مهمې پايلې وي او د وړاندې کوونکي په نوم
هم يادېږي (لکه د زورن لمه)۔ فرعي پايله هغه نتيجه
ده چې له بلې پايلې په اسانۍ راځي۔
\olpseganchor{OLP-0603-B009}{end}

\olpseganchor{OLP-0603-B010}{start}
هغه وينا چې ثبوت يې غواړو، ډېر ځله فرضونه لري
چې څرګندوي د کومو څيزونو په اړه څه ثابتوو۔ ښايي
په دې پيل شي: «\LR{$!A$} دې د \LR{$!B \lif !C$} په بڼه
فارمول وي»، يا «فرض کړئ \LR{$\Gamma \Proves !A$}»،
يا ورته څه۔ دا د بيان، قضيې، يا لمې \emph{فرضونه}
دي، او په خپل ثبوت کښې يې رښتيا ګڼلے شئ۔ دا
هغه څه محدودوي چې ثابتوو، او د بحث څيزونو ته
ځينې نومونه هم ورپېژني۔ د بېلګې په توګه، که ستاسو
بيان داسې پيل شي: «\LR{$!A$} دې د \LR{$!B \lif !C$} په بڼه
فارمول وي»، نو تاسو يوازې د يوه ځانګړي ډول
فارمولونو (يعنې شرطي فارمولونو) په اړه څه ثابتوئ،
او څرګنده ده چې \LR{$!B \lif !C$} يو هر ډول شرطي
فارمول دے چې ستاسو ثبوت پرې خبرې کوي۔
\olpseganchor{OLP-0603-B010}{end}

\olpseganchor{OLP-0604-B004}{start}
\olfileid{mth}{prf}{str}
\olpseganchor{OLP-0604-B004}{end}

\olpseganchor{OLP-0604-B005}{start}
\olsection{د ثبوت پيل}
\olpseganchor{OLP-0604-B005}{end}

\olpseganchor{OLP-0604-B006}{start}
خو له کومه ځايه پيل کوئ؟
\olpseganchor{OLP-0604-B006}{end}

\olpseganchor{OLP-0604-B007}{start}
تاسو ته د ثابتولو دپاره يوه خبره درکړل شوې؛ نو همدا
بايد د ثبوت په پای کښې وروستۍ خبره وي (البته په سر
کښې \emph{اعلانولے} شئ چې دا به ثابتوئ، خو د فرض
په توګه يې نه شئ کارولے)۔ هغه څه چې ثابتول يې
غواړئ د کاغذ د نوې پاڼې په ښکته برخه کښې وليکئ،
څو مو هدف له سترګو پناه نه شي۔
\olpseganchor{OLP-0604-B007}{end}

\olpseganchor{OLP-0604-B008}{start}
بيا ښايي ځينې داسې فرضونه ولرئ چې کارولے يې شئ
(دا خبره به په راتلونکې برخه کښې، د خپل ثبوت د
\emph{ډول} د بحث پر مهال، نوره هم روښانه شي)۔ دا
فرضونه د پاڼې په سر کښې وليکئ او څرګند کړئ چې
فرضونه دي (يعنې که \LR{$p$} فرضوئ، وليکئ «فرض کړئ
چې \LR{$p$}»، يا «ونيسئ چې \LR{$p$}»)۔ په پای کښې ښايي
په پوښتنه کښې داسې تعريفونه وي چې پېژندل يې پکار
دي۔ کېدے شي د يوه ځانګړي تعريف د کارولو غوښتنه
درڅخه شوې وي، يا په فرضونو يا هغه پايله کښې چې
ورته رسېدل غواړئ بېلابېل تعريفونه راغلي وي۔
\emph{دا وليکئ او ډاډ ترلاسه کړئ چې په معنا يې پوهېږئ۔}
\olpseganchor{OLP-0604-B008}{end}

\olpseganchor{OLP-0604-B009}{start}
د ثبوت د برابرولو بڼه د پوښتنې په شکل هم تړلې ده۔
راتلونکې برخه ښيي چې د جملې د ډول له مخې خپل
ثبوت څنګه برابر کړئ۔
\olpseganchor{OLP-0604-B009}{end}

\olpseganchor{OLP-0605-B004}{start}
\olfileid{mod}{prf}{def}
\olpseganchor{OLP-0605-B004}{end}

\olpseganchor{OLP-0605-B005}{start}
\olsection{د تعريفونو کارول}
\olpseganchor{OLP-0605-B005}{end}

\olpseganchor{OLP-0605-B006}{start}
مخکې مو وويل چې له هغو ټولو تعريفونو سره بلد اوسئ چې
په ثبوت کښې کارېدلے شي، او سم يې وکارولے شئ۔ دا ډېر
مهم ټکے دے؛ ښه ده چې لږ نور يې هم وڅېړو۔ تعريفونه
ځانګړنې او اړيکې لنډوي، څو پرې په لنډه توګه خبرې
وکړو۔ هغه نوے راوړل شوے لنډ نوم \emph{تعريفېدونکے}
(\texttt{definiendum}) بلل کېږي، او هغه څه چې نوم يې
لنډوي \emph{تعريفوونکے بيان} (\texttt{definiens}) دے۔
په ثبوتونو کښې ډېری وخت
بايد بېرته وګورو چې تعريفېدونکے څنګه راوړل شوے و،
ځکه د ثبوت د پرمخ بيولو دپاره د تعريفوونکي بيان
منطقي جوړښت کاروو: هماغه اوږده بڼه چې ټاکلے نوم
يې لنډيز دے۔ د تعريفونو په پرانيستلو سره د استدلال
اصل منطقي جوړښت څرګندوئ۔
\olpseganchor{OLP-0605-B006}{end}

\olpseganchor{OLP-0605-B007}{start}
راځئ په يوه بېلګه پيل وکړو۔ فرض کړئ دا خبره ثابتول غواړئ:
\olpseganchor{OLP-0605-B007}{end}

\olpseganchor{OLP-0605-B008}{start}
\begin{prop}
د هر دوو سټونو \LR{$A$} او \LR{$B$} دپاره، \LR{$A \cup B = B \cup A$}۔
\end{prop}
\olpseganchor{OLP-0605-B008}{end}

\olpseganchor{OLP-0605-B009}{start}
د ثبوت له پيلولو مخکې بايد پوه شو چې د دوو سټونو
يو شان کېدل څه معنا لري؛ يعنې په دې مساوات کښې د
سټونو دپاره «\LR{$=$}» څه معنا لري۔ سټونه هغه وخت يو
شان تعريفېږي چې يو شان غړي ولري۔ نو دا
تعريف بايد پرانيزو:
\olpseganchor{OLP-0605-B009}{end}

\olpseganchor{OLP-0605-B010}{start}
\begin{defn}
سټونه \LR{$A$} او \LR{$B$} \emph{يو شان} دي، \LR{$A = B$}، که او يوازې
که د~\LR{$A$} هر غړے د~\LR{$B$} غړے وي، او برعکس۔
\end{defn}
\olpseganchor{OLP-0605-B010}{end}

\olpseganchor{OLP-0605-B011}{start}
دا تعريف \LR{$A$} او~\LR{$B$} د هر ډول سټونو د ځايناستو
نښو په توګه کاروي۔ هغه څه چې تعريفوي --- \emph{تعريفېدونکے}
--- «\LR{$A = B$}» عبارت دے؛ تعريف هغه شرط ښيي چې
له مخې يې \LR{$A = B$} رښتيا وي۔ همدا شرط --- «د~\LR{$A$} هر
غړے د~\LR{$B$} غړے وي، او برعکس» ---
\emph{تعريفوونکے بيان} دے۔\footnote{په همدې ځانګړې
  بېلګه کښې --- که څه هم دا خبره مغشوشوونکې ښکاري! ---
  کله چې \LR{$A = B$} وي، سټونه \LR{$A$} او \LR{$B$} په رښتيا يو
  او هماغه سټ وي، سره له دې چې د مساوات په کيڼ
  او ښي اړخ کښې يې په بېلو تورو ښيو۔ خو د هماغه
  سټ د پېژندلو لارې بېلې کېدے شي؛ له همدې امله
  تعريف بې ګټې نه دے۔} تعريف ټاکي چې \LR{$A = B$}
هغه وخت، او يوازې هغه وخت، رښتيا دے چې ياد شرط
پوره شي؛ «که او يوازې که» ته په لنډه \texttt{iff} هم وايي۔
\olpseganchor{OLP-0605-B011}{end}

\olpseganchor{OLP-0605-B012}{start}
د تعريف د کارولو پر مهال بايد د تعريف \LR{$A$} او \LR{$B$}
له خپلې ځانګړې بېلګې سره برابر کړئ۔ زموږ په بېلګه
کښې، د \LR{$A \cup B = B \cup A$} د رښتيا کېدو دپاره د
هر \LR{$z \in A \cup B$} بايد په \LR{$B \cup A$} کښې هم
وي، او برعکس۔ په بيان کښې \LR{$A \cup B$} د تعريف
د~\LR{$A$} رول لري، او \LR{$B \cup A$} د~\LR{$B$} رول۔ ځکه چې
\LR{$A$} او \LR{$B$} د تعريف او د ثابتېدونکي بيان په دواړو
ځايونو کښې راغلي، خو بېل رولونه لري، پام وکړئ
چې سره يې ګډ نه کړئ۔ د بېلګې په توګه، دا به
تېروتنه وي چې بيان د دې په ښودلو ثابت کړئ چې
د~\LR{$A$} هر غړے د~\LR{$B$} غړے دے، او
برعکس: دا به \LR{$A = B$} وښيي، نه دا چې
\LR{$A \cup B = B \cup A$}۔ (څرنګه چې \LR{$A$} او \LR{$B$}
هر دوه سټونه کېدے شي، له \LR{$A$} او \LR{$B$} په اړه
بې له نورو فرضونو هغوی په رښتيا بېل هم
کېدے شي؛ نو دغه لار
مو وړاندې نه بيايي۔)
\olpseganchor{OLP-0605-B012}{end}

\olpseganchor{OLP-0605-B013}{start}
په ثبوت کښې د سټونو د نظريې مفاهيم، لکه اتحاد،
کاروو؛ نو د ثبوت د لارې د پوهېدو دپاره د \LR{$\cup$}
نښې په معنا هم پوهېدل پکار دي۔ کله کله د يوه
تعريف له پرانيستلو وروسته نور تعريفونه هم
پرانيستل غواړي۔ د بېلګې په توګه، \LR{$A \cup B$}
د \LR{$\Setabs{z}{z \in A
  \text{\RL{ يا }} z \in B}$} په توګه تعريف شوے دے۔ نو که
غواړئ \LR{$x \in A \cup B$} ثابت کړئ، د \LR{$\cup$} تعريف
پرانيستل درته وايي چې
\LR{$x \in \Setabs{z}{z \in A \text{\RL{ يا }} z \in B}$}
ثابت کړئ۔ اوس دا هم په ياد ولرئ چې
\LR{$x \in \Setabs{z}{\dots z\dots}$} هغه وخت او يوازې
هغه وخت وي چې \LR{$\dots
x\dots$}۔ له همدې امله، د
\LR{$\Setabs{z}{\dots z \dots}$} نښې تعريف نور هم
پرانيزو: بايد وښيو چې \LR{$x \in A$} يا \LR{$x \in B$}۔
نو «د \LR{$A \cup B$} هر
غړے د \LR{$B \cup A$} غړے هم دے»
په رښتيا دا معنا لري: «د هر \LR{$x$} دپاره، که \LR{$x \in
A$} يا \LR{$x \in B$} وي، نو \LR{$x \in B$} يا \LR{$x \in A$}
وي»۔ که د بيان ټول تعريفونه پرانيزو، دا
راڅرګندېږي چې لاندې خبره بايد ثابته کړو:
\olpseganchor{OLP-0605-B013}{end}

\olpseganchor{OLP-0605-B014}{start}
\begin{prop}
د هر دوو سټونو \LR{$A$} او \LR{$B$} دپاره: (الف) د هر \LR{$x$}
دپاره، که \LR{$x \in A$} يا \LR{$x \in B$} وي، نو \LR{$x \in B$}
يا \LR{$x \in A$} وي؛ او (ب) د هر \LR{$x$} دپاره، که
\LR{$x \in B$} يا \LR{$x \in A$} وي، نو \LR{$x \in A$} يا
\LR{$x \in B$} وي۔
\end{prop}
\olpseganchor{OLP-0605-B014}{end}

\olpseganchor{OLP-0605-B015}{start}
مهمه خبره دا ده چې د تعريفونو پرانيستل د ثبوت
د جوړولو اړينه برخه ده۔ سم کارول يې کله کله
ګران وي: د تعريف متغيرونه او د ثابتېدونکې دعوې
عبارتونه بايد سره بېل کړئ او په دقت يې سره
برابر کړئ۔ د برياليتوب دپاره بايد پوه شئ چې
پوښتنه څه غواړي او په هغې کښې هر کارېدلے
عبارت څه معنا لري؛ ډېر ځله به تر يوه زيات
تعريف پرانيستل غواړئ۔ په لاندې ساده ثبوتونو
کښې حل نږدې په نېغه له تعريفونو راځي۔ البته
تل به خبره دومره ساده نه وي۔
\olpseganchor{OLP-0605-B015}{end}

\olpseganchor{OLP-0605-B016}{start}
\begin{prob}
فرض کړئ له تاسو څخه غوښتل کېږي چې
\LR{$A \cap B \neq \emptyset$} ثابت کړئ۔ دلته راغلي
ټول تعريفونه پرانيزئ؛ يعنې خبره داسې بيا
ووايئ چې «\LR{$\cap$}»، «=»، يا «\LR{$\emptyset$}»
پکښې ياد نه شي۔
\end{prob}
\olpseganchor{OLP-0605-B016}{end}

\olpseganchor{OLP-0606-B004}{start}
\olfileid{mth}{prf}{pat}
\olpseganchor{OLP-0606-B004}{end}

\olpseganchor{OLP-0606-B005}{start}
\olsection{د استنتاج بڼې}
\olpseganchor{OLP-0606-B005}{end}

\olpseganchor{OLP-0606-B006}{start}
ثبوتونه له بېلو بېلو استنتاجونو جوړېږي۔ کله چې استنتاج
کوو، عموماً يې په «نو»، «له دې امله»، يا «ځکه نو» غوندې
ټکي څرګندوو۔ استنتاج ډېر ځله په هغو يو يا دوو خبرو
ولاړ وي چې په ثبوت کښې مو لا په لاس کښې دي: ښايي
فرض مو کړې وي، يا مو مخکې له بل استنتاج څخه
اخيستې وي۔ د روڼوالي دپاره دغو خبرو ته نښې
ورکولے شو، او د استنتاج پر مهال وايو چې کومې
نورې خبرې کاروو۔ استنتاج ډېر ځله دا هم تشريح کوي
چې نوې پايله مو له مخکېنيو خبرو څخه \emph{ولې}
راوځي۔ په ثبوتونو کښې د استنتاج ځينې عامې بڼې
بيا بيا کارېږي؛ لاندې يې څو بڼې څېړو۔ ځينې
بڼې، لکه د استقرا له لارې ثبوت، زياتې پېچلې دي
او وروسته به ورباندې خبرې وکړو۔
\olpseganchor{OLP-0606-B006}{end}

\olpseganchor{OLP-0606-B007}{start}
د استنتاج يوه بڼه مو لا څېړلې ده: د تعريف پرانيستل
يا کارول۔ د تعريف په پرانيستلو کښې هغه خبره چې
تعريفېدونکے پکې راغلے وي، د تعريفوونکي بيان په
کارولو بيا وايو۔ د بېلګې په توګه، فرض کړئ چې
د ثبوت په ترڅ کښې مو \LR{$D = E$} لا ثابته کړې ده
(الف)۔ بيا د سټونو دپاره د \LR{$=$} تعريف کاروو
او استنتاج کوو: «نو، له (الف) څخه د تعريف له مخې،
د~\LR{$D$} هر غړے د~\LR{$E$} غړے دے، او برعکس»۔
\olpseganchor{OLP-0606-B007}{end}

\olpseganchor{OLP-0606-B008}{start}
يو څه مغشوشوونکې دا ده چې موږ د استنتاج توجيه
ډېر ځله د هماغه استنتاج د کولو پر مهال نه،
بلکې تر هغه مخکې ليکو۔ فرض کړئ \LR{$D = E$} مو لا
نه دي ثابت کړي، خو ثابتول يې غواړو۔ که \LR{$D = E$}
زموږ غوښتل شوې پايله وي، هدف د تعريف په کارولو
بيا ويلے شو: د \LR{$D = E$} د ثابتولو دپاره بايد
ثابته کړو چې د~\LR{$D$} هر غړے د~\LR{$E$}
غړے دے، او برعکس۔ نو د ثبوت بڼه به
داسې وي: (الف) د~\LR{$D$} هر غړے د~\LR{$E$}
غړے ثابت کړئ؛ (ب) د~\LR{$E$} هر
غړے د~\LR{$D$} غړے ثابت کړئ؛
(ج) ځکه نو، له (الف) او (ب) څخه د \LR{$=$} د تعريف
له مخې، \LR{$D = E$}۔ خو عموماً يې داسې نه ليکو۔
پر ځاے يې ښايي داسې وليکو:
\begin{quote}
غواړو وښيو چې \LR{$D = E$}۔ د~\LR{$=$} له تعريفه، د دې
دپاره بايد وښيو چې د~\LR{$D$} هر غړے د~\LR{$E$}
غړے دے، او برعکس۔
\olpseganchor{OLP-0606-B008}{end}

\olpseganchor{OLP-0606-B009}{start}
(الف) \dots (دا ثبوت چې د~\LR{$D$} هر غړے
د~\LR{$E$} غړے دے) \dots
\olpseganchor{OLP-0606-B009}{end}

\olpseganchor{OLP-0606-B010}{start}
(ب) \dots (دا ثبوت چې د~\LR{$E$} هر غړے
د~\LR{$D$} غړے دے) \dots
\end{quote}
\olpseganchor{OLP-0606-B010}{end}

\olpseganchor{OLP-0606-B011}{start}
\subsection{د عطف کارول}
\olpseganchor{OLP-0606-B011}{end}

\olpseganchor{OLP-0606-B012}{start}
ښايي د استنتاج تر ټولو ساده بڼه دا وي چې د عطف
له دوو غړو څخه يو د پايلې په توګه واخلو۔ يعنې که
مو فرض کړي يا مخکې ثابت کړي وي چې \LR{$p$} او~\LR{$q$}،
نو \LR{$p$} (او همدارنګه~\LR{$q$}) استنتاجولے شو۔ دا
دومره بنسټيز استنتاج دے چې ډېر ځله يې جلا يادونه
نه کېږي۔ د بېلګې په توګه، کله چې د \LR{$D = E$} تعريف
پرانيزو، دا مو ثابته کړې وي چې د~\LR{$D$} هر
غړے د~\LR{$E$} غړے دے، او برعکس۔
له دې څخه دا اخيستلے شو چې د~\LR{$E$} هر
غړے د~\LR{$D$} غړے دے؛ همدا د
«برعکس» برخه ده۔
\olpseganchor{OLP-0606-B012}{end}

\olpseganchor{OLP-0606-B013}{start}
\subsection{د عطف ثابتول}
\olpseganchor{OLP-0606-B013}{end}

\olpseganchor{OLP-0606-B014}{start}
کله کله هغه خبره چې ثابتول يې غواړئ د عطف
بڼه لري: له تاسو څخه غوښتل کېږي چې «\LR{$p$} او
\LR{$q$} ثابت کړئ»۔ دلته يوازې دوه کارونه کوئ:
لومړے \LR{$p$} ثابت کړئ، بيا \LR{$q$} ثابت کړئ۔ ثبوت
په دوو برخو وېشلے شئ او د روڼوالي دپاره يې
ونوموئ۔ د لومړنيو يادښتونو پر مهال ښايي د پاڼې
په سر کښې «(۱) \LR{$p$} ثابت کړئ» او د پاڼې په منځ
کښې «(۲) \LR{$q$} ثابت کړئ» وليکئ۔ (کېدے شي له
تاسو څخه د عطف د ثابتولو غوښتنه په څرګنده
نه وي شوې، خو د ثبوت دپاره يې ثابتول اړين
ومومئ۔ د بېلګې په توګه، که د \LR{$D = E$} د ثابتولو
غوښتنه وي، د~\LR{$=$} د تعريف له پرانيستلو وروسته
بايد دا دواړه ثابت کړئ: د~\LR{$D$} هر غړے
د~\LR{$E$} غړے دے \emph{او} د~\LR{$E$} هر
غړے د~\LR{$D$} غړے دے۔)
\olpseganchor{OLP-0606-B014}{end}

\olpseganchor{OLP-0606-B015}{start}
\subsection{د فصل ثابتول}
\olpseganchor{OLP-0606-B015}{end}

\olpseganchor{OLP-0606-B016}{start}
که ثابتېدونکې خبره د فصل بڼه ولري (يعنې
«\LR{$p$} يا \LR{$q$}» وي)، د دوو بديلو غړو له ډلې
د يوه رښتياوالے ښودل بس دے۔ خو په عمل کښې
تقريباً هېڅکله داسې نه وي چې يو بديل غړے
په نېغه د قضيې له فرضونو راووځي۔ ډېر ځله
د قضيې فرضونه پخپله په فصل کښې وي، يا تاسو
ښيئ چې د يوه ډول هر څيز له دوو ځانګړنو څخه
يوه لري، خو ځينې څيزونه لومړۍ او نور دويمه
ځانګړنه لري۔ همدلته د حالتونو له مخې ثبوت
ګټور وي (لاندې يې وګورئ)۔
\olpseganchor{OLP-0606-B016}{end}

\olpseganchor{OLP-0606-B017}{start}
\subsection{شرطي ثبوت}
\olpseganchor{OLP-0606-B017}{end}

\olpseganchor{OLP-0606-B018}{start}
ډېرې قضيې چې ورسره مخامخ کېږئ شرطي بڼه
لري (يعنې وښيئ چې که \LR{$p$} وي، نو \LR{$q$} هم
رښتيا دے)۔ د دغه ډول ثبوت برابرول اسان دي:
د شرطي مقدم فرض کړئ (دلته \LR{$p$})، او له هغه
څخه پايله~\LR{$q$} ثابته کړئ۔ نو که قضيه وايي
«که \LR{$p$}، نو \LR{$q$}»، ثبوت په «\LR{$p$} فرض کړئ»
پيلوئ او په پای کښې بايد~\LR{$q$} مو ثابت کړے وي۔
\olpseganchor{OLP-0606-B018}{end}

\olpseganchor{OLP-0606-B019}{start}
شرطي خبرې په بېلابېلو لارو ويل کېږي۔ نو د
«که \LR{$p$}، نو \LR{$q$}» پر ځاے قضيه ښايي ووايي:
«\LR{$p$} يوازې که \LR{$q$}»، «\LR{$q$} که \LR{$p$}»، يا
«\LR{$q$}، په دې شرط چې \LR{$p$}»۔ دا ټول يوه معنا
لري: \LR{$p$} فرض کړئ او له هغه څخه~\LR{$q$} ثابت
کړئ۔ په ياد ولرئ چې دوه اړخيز شرطي («\LR{$p$}
که او يوازې که \LR{$q$}») په حقيقت کښې دوه
شرطي خبرې دي: که \LR{$p$}، نو \LR{$q$}؛ او که \LR{$q$}،
نو~\LR{$p$}۔ نو دوه شرطي ثبوتونه پکار دي: يو
د لومړي لوري او بل د دويم لوري دپاره۔
کله کله بيا د «که او يوازې که» خبره د
همداسې دوه اړخيزو خبرو په لړۍ ثابتولے شئ،
چې له «\LR{$p$}» څخه پيل او «\LR{$q$}» ته رسېږي؛
خو هر ګام بايد په رښتيا د دواړو لورو
شرط وي۔
\olpseganchor{OLP-0606-B019}{end}

\olpseganchor{OLP-0606-B020}{start}
\subsection{عمومي دعوې}
\olpseganchor{OLP-0606-B020}{end}

\olpseganchor{OLP-0606-B021}{start}
د عمومي دعوې کارول اسان دي: که خبره د
هر څيز دپاره رښتيا وي، د هر ځانګړي څيز
دپاره هم رښتيا ده۔ نو، د بېلګې په توګه،
که ستاسو د ثبوت فرض \LR{$A \subseteq B$} وي،
د~\LR{$\subseteq$} د تعريف له پرانيستلو سره
دا معنا لري چې د هر \LR{$x \in A$} دپاره
\LR{$x \in B$}۔ ځکه نو که لا پوهېږئ چې
\LR{$z \in A$}، دا پايله اخيستلے شئ چې
\LR{$z \in B$}۔
\olpseganchor{OLP-0606-B021}{end}

\olpseganchor{OLP-0606-B022}{start}
د عمومي دعوې ثابتول ښايي لږ ستونزمن
وبرېښي۔ دا ډول خبرې عموماً داسې وي:
«که \LR{$x$} د~\LR{$P$} ځانګړنه ولري، نو د~\LR{$Q$}
ځانګړنه هم لري»، يا «ټول \LR{$P$}s د \LR{$Q$}s
دي»۔ البته تل به يې بڼه کټ مټ داسې نه
وي؛ د پوښتنې د کره غوښتنې پېژندل لږ
تمرین غواړي۔ خو ډېر ځله بايد ثابته کړو
چې د يوې ځانګړنې لرونکي ټول څيزونه
يوه بله ځانګړنه هم لري۔
\olpseganchor{OLP-0606-B022}{end}

\olpseganchor{OLP-0606-B023}{start}
د عمومي دعوې د ثابتولو لار دا ده چې
د لومړۍ ځانګړنې لرونکو څيزونو دپاره
نومونه يا متغيرونه راوړئ، او بيا وښيئ
چې دويمه ځانګړنه هم لري۔ ويلے شو چې
د \emph{ټولو}~\LR{$P$}s دپاره د خبرې د
ثابتولو دپاره بايد يې د يوه
\emph{هر ډول}~\LR{$P$} دپاره ثابت کړئ۔
راوړل شوے نوم د همداسې هر ډول~\LR{$P$}
نوم دے۔ عموماً د دغو هر ډول څيزونو
د نومونو دپاره يو توری کاروو، او
توري له دود سره سم ټاکو: لکه \LR{$n$}
د طبيعي عددونو، \LR{$!A$} د فارمولونه،
\LR{$A$} د سټونو، او \LR{$f$} د تابعو دپاره۔
\olpseganchor{OLP-0606-B023}{end}

\olpseganchor{OLP-0606-B024}{start}
اصل چل دا دے چې په ټول ثبوت کښې
عموميت وساتئ۔ په پيل کښې فرض کړئ چې
يو هر ډول څيز («\LR{$x$}») د~\LR{$P$}
ځانګړنه لري؛ بيا يوازې د تعريفونو
او اجازه شوو فرضونو له مخې وښيئ
چې \LR{$x$} د~\LR{$Q$} ځانګړنه هم لري۔
ځکه مو د \LR{$x$} په اړه، له \LR{$P$} لرلو
پرته، بل ځانګړے شرط نه دے ايښے،
ويلے شئ چې هر هغه څيز چې \LR{$P$}
لري، \LR{$Q$} هم لري۔ لنډه دا چې \LR{$x$}
د~\LR{$P$} ځانګړنې د \emph{ټولو}
څيزونو ځايناستے دے۔
\olpseganchor{OLP-0606-B024}{end}

\olpseganchor{OLP-0606-B025}{start}
\begin{prop}
  د هر دوو سټونو \LR{$A$} او \LR{$B$} دپاره، \LR{$A \subseteq A \cup B$}۔
\end{prop}
\olpseganchor{OLP-0606-B025}{end}

\olpseganchor{OLP-0606-B026}{start}
\begin{proof}
  \LR{$A$} او \LR{$B$} دې هر ډول سټونه وي۔ غواړو وښيو چې \LR{$A
  \subseteq A \cup B$}۔ د \LR{$\subseteq$} له تعريفه دا
  معنا لري: د هر \LR{$x$} دپاره، که \LR{$x \in A$} وي، نو
  \LR{$x \in A \cup B$}۔ نو \LR{$x \in A$} دې د~\LR{$A$} يو
  هر ډول غړے وي۔ بايد وښيو چې \LR{$x \in A
  \cup B$}۔ څرنګه چې \LR{$x \in A$}، نو \LR{$x \in A$}
  يا \LR{$x \in B$}۔ له دې څخه \LR{$x \in
  \Setabs{x}{x \in A \lor x \in B}$}۔ خو دا،
  د \LR{$\cup
  $} له تعريفه، \LR{$x \in A \cup B$} معنا لري۔
\end{proof}
\olpseganchor{OLP-0606-B026}{end}

\olpseganchor{OLP-0606-B027}{start}
\subsection{د حالتونو له مخې ثبوت}
\olpseganchor{OLP-0606-B027}{end}

\olpseganchor{OLP-0606-B028}{start}
فرض کړئ فصل د يوه فرض يا مخکې ثابتې شوې پايلې
په توګه لرئ: فرض مو کړي يا ثابت کړي دي چې \LR{$p$}
يا \LR{$q$} رښتيا دے۔ غواړئ \LR{$r$} ثابت کړئ۔ دا کار
په دوو پړاوونو کښې کوئ: لومړے \LR{$p$} رښتيا
فرض کړئ او~\LR{$r$} ثابت کړئ؛ بيا \LR{$q$} رښتيا
فرض کړئ او~\LR{$r$} يو ځل بيا ثابت کړئ۔ دا
لار ځکه کار کوي چې له دوو بديلو څخه
يو رښتيا فرض شوے يا پېژندل شوے دے۔
دواړه پړاوونه ښيي چې هر يو په جلا
توګه د~\LR{$r$} د رښتياوالي دپاره بس دے۔
(که دواړه رښتيا وي، د \LR{$r$} د رښتياوالي
دپاره يو نه، بلکې دوه دليلونه لرو۔
اړينه نه ده چې \LR{$r$} د \LR{$p$} او~\LR{$q$}
دواړو په يو وخت فرضولو سره بيا جلا
ثابت کړو۔) د خپلې لارې د څرګندولو
دپاره وايو چې «حالتونه بېلوو»۔
د بېلګې په توګه، فرض کړئ پوهېږو
چې \LR{$x \in B \cup C$}۔ \LR{$B \cup C$}
د \LR{$\Setabs{x}{x \in B \text{\RL{ or }} x \in C}$}
په توګه تعريف شوے دے۔ يعنې د تعريف
له مخې، \LR{$x \in B$} يا \LR{$x \in C$}۔
له دې څخه \LR{$x \in A$} داسې ثابتوو:
لومړے \LR{$x \in B$} فرض کړئ او له هغه
څخه \LR{$x \in A$} ثابت کړئ؛ بيا
\LR{$x \in C$} فرض کړئ او له هغه
څخه يو ځل بيا \LR{$x \in A$} ثابت
کړئ۔ د فرض تر لاندې «حالتونه
بېلوو» وليکئ، بيا «حالت (۱):
\LR{$x \in B$}»؛ او د پاڼې په منځ
کښې «حالت (۲): \LR{$x \in C$}»۔
وروسته د پاڼې پاسنۍ او لاندنۍ
نيمايي د ثبوت په ګامونو ډکې کړئ۔
\olpseganchor{OLP-0606-B028}{end}

\olpseganchor{OLP-0606-B029}{start}
د حالتونو له مخې ثبوت په ځانګړې توګه هغه
وخت ګټور دے چې غوښتل شوې پايله پخپله په
فصل کښې وي۔ دا يې يوه ساده بېلګه ده:
\olpseganchor{OLP-0606-B029}{end}

\olpseganchor{OLP-0606-B030}{start}
\begin{prop}
فرض کړئ \LR{$B \subseteq D$} او \LR{$C \subseteq E$}۔
نو \LR{$B \cup C \subseteq
D \cup E$}۔
\end{prop}
\olpseganchor{OLP-0606-B030}{end}

\olpseganchor{OLP-0606-B031}{start}
\begin{proof}
  فرض کړئ (الف) \LR{$B \subseteq D$} او (ب)
  \LR{$C \subseteq E$}۔ د تعريف له مخې،
  هر \LR{$x \in B$}، \LR{$ \in D$} هم دے
  (ج)، او هر \LR{$x \in C$}، \LR{$ \in E$}
  هم دے (د)۔
  د \LR{$B \cup C \subseteq D \cup E$}
  د ښودلو دپاره بايد وښيو چې که
  \LR{$x \in B \cup C$} وي، نو
  \LR{$x \in D \cup E$} وي (د
  \LR{$\subseteq$} له تعريفه)۔
  \LR{$x \in B \cup C$} هغه وخت او
  يوازې هغه وخت وي چې \LR{$x \in B$}
  يا \LR{$x \in C$} وي (د~\LR{$\cup$}
  له تعريفه)۔ همدارنګه،
  \LR{$x \in D \cup E$} هغه وخت او
  يوازې هغه وخت وي چې \LR{$x \in D$}
  يا \LR{$x \in E$} وي۔ نو بايد وښيو:
  د هر \LR{$x$} دپاره، که \LR{$x \in B$}
  يا \LR{$x \in C$} وي، نو \LR{$x \in D$}
  يا \LR{$x \in E$} وي۔
\olpseganchor{OLP-0606-B031}{end}

\olpseganchor{OLP-0606-B032}{start}
  \begin{quote}
  تر دې ځايه مو يوازې تعريفونه پرانيستي
  دي!{} خپل بيان مو د \LR{$\subseteq$} او
  \LR{$\cup$} له يادولو پرته بيا وايه، او اوس
  بايد يوه عمومي شرطي خبره ثابته کړو۔
  د پورته بحث له مخې، د دې دپاره
  \LR{$x$} داسې هر ډول څيز نيسو چې د
  «که» برخه يې رښتيا فرضوو؛ بيا
  ښيو چې د «نو» برخه هم رښتيا ده۔
  په بل عبارت، فرضوو چې \LR{$x
  \in B$} يا \LR{$x \in C$}، او ښيو
  چې \LR{$x \in D$} يا \LR{$x \in
  E$}۔\footnote{دا پراګراف يوازې
    زموږ کار تشريح کوي؛ د ثبوت
    برخه نه ده۔ د خپلو ثبوتونو
    د ليکلو پر مهال دومره تفصيل
    اړين نه دے۔}
  \end{quote}
\olpseganchor{OLP-0606-B032}{end}

\olpseganchor{OLP-0606-B033}{start}
  فرض کړئ \LR{$x \in B$} يا \LR{$x \in C$}۔
  بايد وښيو چې \LR{$x \in D$} يا
  \LR{$x \in E$}۔ حالتونه بېلوو۔
\olpseganchor{OLP-0606-B033}{end}

\olpseganchor{OLP-0606-B034}{start}
  حالت ۱: \LR{$x \in B$}۔ د (ج)
  له مخې، \LR{$x \in D$}۔ ځکه نو،
  \LR{$x \in D$} يا \LR{$x \in
  E$}۔ (دلته مو په مخکېنۍ
  فرعي برخه کښې څېړل شوے
  استنتاج وکړ!)
\olpseganchor{OLP-0606-B034}{end}

\olpseganchor{OLP-0606-B035}{start}
  حالت ۲: \LR{$x \in C$}۔ د (د)
  له مخې، \LR{$x \in E$}۔ ځکه نو،
  \LR{$x \in D$} يا \LR{$x \in E$}۔
 \end{proof}
\olpseganchor{OLP-0606-B035}{end}

\olpseganchor{OLP-0606-B036}{start}
\subsection{د وجود د دعوې ثابتول}
\olpseganchor{OLP-0606-B036}{end}

\olpseganchor{OLP-0606-B037}{start}
کله چې د وجود د دعوې د ثابتولو غوښتنه
درڅخه کېږي، پوښتنه عموماً داسې وي:
«ثابت کړئ چې داسې يو~\LR{$x$} شته چې
\LR{$\dots x \dots$}»؛ يعنې داسې يو
څيز شته چې د «\LR{$\dots x
\dots$}» په عبارت کښې راغلې
ځانګړنه لري۔ دلته بايد مناسب
څيز وټاکئ او وښيئ چې غوښتل
شوې ځانګړنه لري۔ دا کار
ساده ښکاري، خو دا ډول ثبوت
ستونزمن کېدے شي۔ عموماً بايد
يو څيز \emph{جوړ} يا
\emph{تعريف} کړئ او بيا
ثابته کړئ چې هغه ټاکلې
ځانګړنه لري۔ د سم څيز
موندل، د ځانګړنې ثابتول،
او کله کله ان دا ښودل
چې څيز مو په رښتيا تعريف
کړے دے، ټول ګران کېدے شي!{}
\olpseganchor{OLP-0606-B037}{end}

\olpseganchor{OLP-0606-B038}{start}
عموماً دا کار د څيز په ټاکلو ليکئ،
لکه «\LR{$x$} دې \dots وي» (دلته
\dots{} څرګندوي چې کوم څيز
مو په پام کښې دے)۔ کېدے شي
لومړے ثابته کړئ چې \LR{$\dots$}
په رښتيا يو موجود څيز ښيي؛
بيا وښيئ چې \LR{$x$} د~\LR{$Q$}
ځانګړنه لري۔ دا يې يوه ساده
بېلګه ده۔
\olpseganchor{OLP-0606-B038}{end}

\olpseganchor{OLP-0606-B039}{start}
\begin{prop}
  فرض کړئ \LR{$x \in B$}۔ نو داسې يو~\LR{$A$}
  شته چې \LR{$A \subseteq
  B$} او \LR{$A \neq \emptyset$}۔
\end{prop}
\olpseganchor{OLP-0606-B039}{end}

\olpseganchor{OLP-0606-B040}{start}
\begin{proof}
  فرض کړئ \LR{$x \in B$}۔ \LR{$A = \{x\}$} ونيسئ۔
  \begin{quote}
    دلته مو د~\LR{$A$} سټ د
    غړي په شمېرلو تعريف کړے دے۔
    څرنګه چې \LR{$x$} يو څيز
    فرضوو، او تل د غړي
    په شمېرلو سټ جوړولے
    شو، دلته دا جلا ثابتول
    نه دي پکار چې د~\LR{$A$}
    سټ مو په برياليتوب
    تعريف کړے دے۔ خو دا
    لا بايد وښيو چې \LR{$A$}
    د بيان غوښتل شوې
    ځانګړنې لري؛ بې له
    دې ثبوت بشپړ نه دے!{}
  \end{quote}
  څرنګه چې \LR{$x \in A$}،
  نو \LR{$A \neq \emptyset$}۔
  \begin{quote}
    دا د \LR{$A$} له تعريفه
    د \LR{$\{x\}$} په توګه
    او له ښکاره واقعيتونو
    څخه راځي چې
    \LR{$x \in \{x\}$}
    او \LR{$x \notin \emptyset$}۔
  \end{quote}
  ځکه چې \LR{$x$} د~\LR{$\{x\}$}
  يوازينے غړے دے،
  او \LR{$x \in B$}، د~\LR{$A$}
  هر غړے د~\LR{$B$}
  غړے هم دے۔
  د~\LR{$\subseteq$} له تعريفه،
  \LR{$A \subseteq B$}۔
\end{proof}
\olpseganchor{OLP-0606-B040}{end}

\olpseganchor{OLP-0606-B041}{start}
\subsection{د وجود د دعوو کارول}
\olpseganchor{OLP-0606-B041}{end}

\olpseganchor{OLP-0606-B042}{start}
فرض کړئ پوهېږئ چې يوه وجودي دعوه
رښتيا ده (يا مو ثابته کړې، يا يې
د کارولو فرض لرئ)، لکه «د ځينو~\LR{$x$}
دپاره \LR{$x \in A$}» يا «يو
\LR{$x \in A$} شته»۔ که يې په ثبوت
کښې کاروئ، داسې چلند وکړئ لکه
د هغو څيزونو له ډلې د يوه نوم
درسره وي چې فرض يې شتون ښيي۔
څرنګه چې په \LR{$A$} کښې لږ تر لږه
يو څيز شته، نوم پرې ايښودل
ممکن دي۔ ښايي تاسو ونه شئ
کړے هغه په ځانګړې توګه
وټاکئ يا يې، له دې پرته چې
\LR{$\in A$} دے، نور تشريح کړئ۔
خو د ثبوت دپاره يې د ټاکلو
په توګه ونيسئ او يو نوم
ورکړئ۔ مهمه ده چې داسې نوم
وټاکئ چې مخکې مو نه وي
کارولے او په فرضونو کښې
هم نه وي راغلے، ګنې تېروتنه
پېښېدے شي۔ په ثبوت کښې دا
ګام د «د ځينو \LR{$x$} دپاره
\LR{$x \in A$}» څخه «\LR{$a
\in A$} دې وي» ته په تګ
ښيئ۔ اوس د~\LR{$a$} په اړه
استدلال او نور فرضونه
کارولے شئ، تر څو يوې
پايلې \LR{$p$} ته ورسېږئ۔
که \LR{$p$} نور~\LR{$a$} نه يادوي،
نو \LR{$p$} د \LR{$a \in A$}
له فرضه خپلواکه ده؛ ښودلې
مو ده چې يوازې د «د ځينو
\LR{$x$} دپاره \LR{$x \in A$}»
له فرضه راځي۔
\olpseganchor{OLP-0606-B042}{end}

\olpseganchor{OLP-0606-B043}{start}
\begin{prop}
که \LR{$A \neq \emptyset$} وي، نو
\LR{$A \cup B \neq \emptyset$}۔
\end{prop}
\olpseganchor{OLP-0606-B043}{end}

\olpseganchor{OLP-0606-B044}{start}
\begin{proof}
  فرض کړئ \LR{$A \neq \emptyset$}۔
  نو د ځينو \LR{$x$} دپاره \LR{$x \in A$}۔
  \begin{quote}
    دلته مو لومړے د بيان
    فرض بيا ووايه۔ دا فرض،
    يعنې \LR{$A \neq \emptyset$}،
    يوه وجودي دعوه پټوي
    چې د څو تعريفونو له
    پرانيستلو وروسته
    راښکاره کېږي۔ د \LR{$=$}
    تعريف وايي چې \LR{$A = \emptyset$}
    هغه وخت او يوازې هغه
    وخت وي چې هر \LR{$x \in A$}
    په \LR{$\in \emptyset$} کښې
    هم وي او هر \LR{$x \in
    \emptyset$} په \LR{$\in A$}
    کښې هم وي۔ د دواړو
    خواوو په نفي سره،
    \LR{$A \neq
    \emptyset$} هغه وخت او
    يوازې هغه وخت وي چې
    يا ځينې \LR{$x \in A$} په
    \LR{$\notin \emptyset$} کښې وي،
    يا ځينې \LR{$x \in \emptyset$}
    په \LR{$\notin A$} کښې وي۔
    څرنګه چې هېڅ شے په
    \LR{$\in \emptyset$} کښې نه
    دے، دويم بديل هېڅکله
    رښتيا نه شي کېدے؛ او
    «\LR{$x \in A$} او \LR{$x
    \notin \emptyset$}» يوازې
    \LR{$x \in A$} ته راټيټېږي۔
    نو \LR{$A \neq
    \emptyset$} هغه وخت او
    يوازې هغه وخت وي چې
    د ځينو \LR{$x$} دپاره \LR{$x \in A$}۔
    دا وجودي دعوه ده۔
    اوس د~\LR{$A$} له
    غړي څخه
    يوه ته نوم ورکوو:
  \end{quote}
  \LR{$a \in A$} دې وي۔
  \begin{quote}
    اوس مو په \LR{$\in A$}
    کښې يوه څيز ته نوم
    ورکړ۔ د~\LR{$a$} په
    اړه استدلال ته دوام
    ورکوو، خو پام کوو چې
    له \LR{$a \in A$} پرته
    نور څه فرض نه کړو:
  \end{quote}
  ځکه چې \LR{$a \in A$}،
  \LR{$a \in A \cup B$}،
  د~\LR{$\cup$} له تعريفه۔
  نو د ځينو \LR{$x$} دپاره
  \LR{$x \in A \cup B$}،
  يعنې \LR{$A \cup B \neq \emptyset$}۔
  \begin{quote}
    په دغه وروستي ګام
    کښې له «\LR{$a \in A \cup B$}»
    څخه «د ځينو \LR{$x$}
    دپاره \LR{$x \in A \cup B$}»
    ته لاړو۔ وروستۍ
    خبره نور \LR{$a$} نه يادوي؛
    نو پوهېږو چې «د ځينو
    \LR{$x$} دپاره \LR{$x \in A \cup B$}»
    يوازې له «د ځينو
    \LR{$x$} دپاره \LR{$x \in A$}»
    څخه راځي۔ همدا معنا
    لري چې \LR{$A \cup B \neq
    \emptyset$}۔
  \end{quote}
\end{proof}
\olpseganchor{OLP-0606-B044}{end}

\olpseganchor{OLP-0606-B045}{start}
ښايي ښه تمرين وي چې تړلي متغيرونه،
لکه «\LR{$x$}»، له فرضي نومونو،
لکه \LR{$a$}، بېل وساتو؛ موږ هم
همداسې وکړل۔ خو په عمل کښې
ډېر ځله دا بېلوالے نه ساتو
او \LR{$x$} کاروو، لکه دلته:
\begin{quote}
فرض کړئ \LR{$A \neq \emptyset$}،
يعنې يو \LR{$x \in A$} شته۔
د \LR{$\cup$} له تعريفه،
\LR{$x \in A \cup B$}۔
نو \LR{$A \cup B \neq \emptyset$}۔
\end{quote}
خو په دې توګه استدلال
کښې بايد په ډېر پام
د بېلابېلو وجودي دعوو
دپاره بېل \LR{$x$} او \LR{$y$}
وکاروئ۔ د بېلګې په توګه،
لاندې د «که \LR{$A \neq
\emptyset$} او \LR{$B \neq \emptyset$}
وي، نو \LR{$A \cap B \neq \emptyset$}»
ثبوت \emph{سم نه دے} (ځکه
دا دعوه رښتيا نه ده):
\begin{quote}
فرض کړئ \LR{$A \neq \emptyset$}
او \LR{$B \neq \emptyset$}۔
نو د ځينو \LR{$x$} دپاره
\LR{$x \in A$}، او هم د ځينو
\LR{$x$} دپاره \LR{$x \in B$}۔
څرنګه چې \LR{$x \in A$}
او \LR{$x \in
B$}، \LR{$x \in A \cap B$}،
د~\LR{$\cap$} له تعريفه۔
نو \LR{$A \cap B \neq
\emptyset$}۔
\end{quote}
ايا هغه ناسم ګام پېژندلے
شئ او ويلے شئ چې دا
پايله ولې نه راځي؟
\olpseganchor{OLP-0606-B045}{end}

\olpseganchor{OLP-0607-B004}{start}
\olfileid{mth}{prf}{ex1}
\olpseganchor{OLP-0607-B004}{end}

\olpseganchor{OLP-0607-B005}{start}
\olsection{يوه بېلګه}
\olpseganchor{OLP-0607-B005}{end}

\olpseganchor{OLP-0607-B006}{start}
زموږ لومړۍ بېلګه د سټونو د اتحاد او تقاطع
په اړه لاندې ساده خبره ده۔ دا د تعريفونو
پرانيستل، د عطف او عمومي دعوو ثابتول،
او د حالتونو له مخې ثبوت روښانوي۔
\olpseganchor{OLP-0607-B006}{end}

\olpseganchor{OLP-0607-B007}{start}
\begin{prop}
د هر ډول سټونو \LR{$A$}، \LR{$B$}، او \LR{$C$} دپاره،
\LR{$A \cup (B \cap C) = (A \cup B)
\cap (A \cup C)$}۔
\end{prop}
\olpseganchor{OLP-0607-B007}{end}

\olpseganchor{OLP-0607-B008}{start}
راځئ ثابته يې کړو!{}
\olpseganchor{OLP-0607-B008}{end}

\olpseganchor{OLP-0607-B009}{start}
\begin{proof}
غواړو وښيو چې د هر ډول سټونو \LR{$A$}، \LR{$B$}، او
\LR{$C$} دپاره \LR{$A \cup (B \cap
C) = (A \cup B) \cap (A \cup C)$}۔
\begin{quote}
لومړے د بيان د «\LR{$=$}» تعريف پرانيزو۔
په ياد ولرئ چې د سټونو يو شان کېدل
ثابتول دا معنا لري چې دواړه يو شان
غړي لري۔ يعنې د \LR{$A \cup
(B \cap C)$} ټول غړي
د \LR{$(A \cup B) \cap (A \cup C)$}
غړي هم دي، او برعکس۔
«برعکس» معنا دا ده چې د \LR{$(A
\cup B) \cap (A \cup C)$} هر
غړے بايد د \LR{$A \cup (B \cap
C)$} غړے هم وي۔
نو د تعريف له پرانيستلو ښکاري
چې يو عطف ثابتول پکار دي۔
دا وليکو:
\end{quote}
د تعريف له مخې،
\LR{$A \cup (B \cap C) = (A \cup B) \cap (A \cup C)$}
هغه وخت او يوازې هغه وخت وي چې د
\LR{$A \cup (B \cap C)$} هر غړے
د \LR{$(A
\cup B) \cap (A \cup C)$}
غړے هم وي، او د
\LR{$(A \cup B) \cap (A
\cup C)$} هر غړے د
\LR{$A \cup (B \cap C)$}
غړے وي۔
\begin{quote}
څرنګه چې دا عطف دے، د دواړو غړو
ثبوت جلا پکار دے۔ له لومړي څخه
پيلوو: ثابتوو چې د
\LR{$A \cup (B \cap C)$} هر
غړے د \LR{$(A
\cup B) \cap (A \cup C)$}
غړے هم دے۔
\olpseganchor{OLP-0607-B009}{end}

\olpseganchor{OLP-0607-B010}{start}
دا عمومي دعوه ده؛ نو د
\LR{$A \cup (B \cap C)$} يو هر ډول
غړے نيسو او ښيو چې
د \LR{$(A \cup B) \cap (A \cup C)$}
غړے هم دے۔ د دغه
هر ډول غړے دپاره يو
متغير، لکه~\LR{$z$}، ټاکو۔ ثبوت
مو داسې دوام کوي:
\end{quote}
لومړے ثابتوو چې د
\LR{$A \cup (B \cap C)$} هر
غړے د
\LR{$(A \cup B) \cap (A \cup C)$}
غړے هم دے۔ \LR{$z \in A \cup (B
\cap C)$} دې وي۔ بايد وښيو چې
\LR{$z \in (A \cup B) \cap (A \cup C)$}۔
\begin{quote}
اوس د \LR{$\cup$} او~\LR{$\cap$} تعريفونه
پرانيزو۔ د بېلګې په توګه، د
\LR{$\cup$} تعريف دا دے:
\LR{$A \cup B = \Setabs{z}{z \in A
  \text{\RL{ يا }} z \in B}$}۔ کله چې
تعريف په «\LR{$A \cup (B
\cap C)$}» کاروو، د تعريف د
«\LR{$B$}» رول اوس «\LR{$B \cap C$}»
لوبوي؛ نو
\LR{$A \cup (B \cap C) = \Setabs{z}{z \in A \text{\RL{ يا }}
  z \in B \cap C}$}۔
ځکه زموږ د \LR{$z \in A \cup (B \cap C)$}
فرض دا معنا لري:
\LR{$z \in \Setabs{z}{z \in A \text{\RL{ يا }} z \in B \cap
  C}$}۔ او \LR{$z \in \Setabs{z}{\dots z\dots}$}
هغه وخت او يوازې هغه وخت وي چې
\dots \LR{$z$} \dots؛ يعنې دلته
\LR{$z \in A$} يا \LR{$z \in B \cap C$}۔
\end{quote}
د \LR{$\cup$} له تعريفه،
\LR{$z \in A$} يا \LR{$z \in B \cap C$}۔
\begin{quote}
دا فصل دے؛ نو د حالتونو له مخې
ثبوت ګټور دے۔ دواړه حالتونه اخلو
او په هر يوه کښې غوښتل شوې پايله،
يعنې «\LR{$z \in (A \cup B) \cap (A \cup
C)$}»، ښيو۔
\end{quote}
حالت ۱: فرض کړئ \LR{$z \in A$}۔
\begin{quote}
له خپلو فرضونو څخه نور ډېر څه
په لاس کښې نه لرو۔ نو هغه څه
وګورو چې په پايله کښې ترې کار
اخيستلے شو۔ غواړو وښيو چې
\LR{$z \in (A \cup B) \cap (A \cup C)$}۔
د \LR{$\cap$} له تعريفه، د
\LR{$z \in (A \cup B) \cap (A \cup C)$}
د ښودلو دپاره بايد وښيو چې
دا هم په \LR{$(A \cup B)$} او هم
په \LR{$(A \cup C)$} کښې دے۔
خو \LR{$z
\in A \cup B$} هغه وخت او يوازې
هغه وخت وي چې \LR{$z \in A$} يا
\LR{$z \in B$} وي، او د حالت~۱
له فرضه \LR{$z \in A$} راسره شته۔
په همداسې استدلال، \LR{$C$} د \LR{$B$}
پر ځاے راوړلو سره، \LR{$z \in A \cup C$}
راځي۔ دا استدلال مو په معکوس
لوري وکړ؛ اوس يې د ثبوت د
اړين لوري په ترتيب ليکو۔
\end{quote}
څرنګه چې \LR{$z \in A$}، نو
\LR{$z \in A$} يا \LR{$z \in B$}؛
له دې امله د~\LR{$\cup$} له تعريفه
\LR{$z \in A \cup B$}۔ همدارنګه
\LR{$z \in A \cup C$}۔ دا د
\LR{$z \in (A \cup B) \cap (A \cup C)$}
معنا لري، د~\LR{$\cap$} له تعريفه۔
\begin{quote}
دا د حالتونو له مخې د ثبوت
لومړے حالت بشپړوي۔ اوس په
دويم حالت کښې، چې
\LR{$z \in B \cap C$} دے، هم
غوښتل شوې پايله اخلو۔
\end{quote}
حالت ۲: فرض کړئ \LR{$z \in B \cap C$}۔
\begin{quote}
بيا د دوو سټونو له تقاطع
سره کار لرو۔ د~\LR{$\cap$}
تعريف کاروو:
\end{quote}
څرنګه چې \LR{$z \in B \cap C$}،
\LR{$z$} بايد هم د \LR{$B$} او هم د \LR{$C$}
غړے وي، د~\LR{$\cap$}
له تعريفه۔
\begin{quote}
بيا خپلې پايلې ته ګورو۔
بايد وښيو چې \LR{$z$} هم
په \LR{$(A \cup B)$} او هم
په \LR{$(A \cup C)$} کښې
دے۔ حل بيا سمدستي راځي۔
\end{quote}
څرنګه چې \LR{$z \in B$}،
\LR{$z \in (A \cup B)$}۔
څرنګه چې \LR{$z \in C$}،
\LR{$z \in (A
\cup C)$} هم۔ نو
\LR{$z \in (A \cup B) \cap (A \cup C)$}۔
\begin{quote}
دلته مو د \LR{$\cup$} او \LR{$\cap$}
تعريفونه بيا وکارول؛ خو ځکه
چې مخکې مو ياد کړي او ښودلې
مو ده چې د دوو سټونو په يوه کښې د \LR{$z$}
غړيتوب په اتحاد کښې يې غړيتوب ورکوي،
اړتيا نشته
چې دا ګامونه بيا په هماغه
تفصيل ووايو۔
\olpseganchor{OLP-0607-B010}{end}

\olpseganchor{OLP-0607-B011}{start}
د حالتونو له مخې د ثبوت
دويم حالت مو هم بشپړ کړ؛
اوس لومړۍ پايله ويلے شو۔
\end{quote}
نو که \LR{$z \in A \cup (B \cap C)$}،
بيا \LR{$z \in (A \cup B) \cap (A \cup C)$}۔
\begin{quote}
اوس يوازې بل لوری ښيو:
د \LR{$(A \cup B) \cap (A \cup C)$}
هر غړے د
\LR{$A \cup (B \cap
C)$} غړے هم دے۔
لکه مخکې، دا عمومي دعوه
د لومړي سټ د يوه هر ډول
غړے په فرضولو
او په دويم سټ کښې يې
د غړيتوب په ښودلو ثابتوو۔
راځئ دا ګام وليکو۔
\end{quote}
اوس فرض کړئ
\LR{$z \in (A \cup B) \cap (A \cup C)$}۔
غواړو وښيو چې
\LR{$z \in A \cup (B \cap C)$}۔
\begin{quote}
اوس زموږ فرض
\LR{$z \in (A \cup B) \cap (A
\cup C)$} دے۔ هيله ده
دا خبره به مغشوشوونکې
نه وي چې دلته هم د ثبوت
د لومړۍ برخې هماغه~\LR{$z$}
کاروو۔ کله چې هغه برخه
بشپړه شوه، هلته شوي
فرضونه نور نه پلي کېږي؛
نو اوس د \LR{$z$} په اړه
نوي فرضونه کولے شو۔
که دا درته مغشوش وي،
په لاندې برخه کښې \LR{$z$}
په بل متغير بدل کړئ۔
\olpseganchor{OLP-0607-B011}{end}

\olpseganchor{OLP-0607-B012}{start}
پوهېږو چې \LR{$z$} هم په
\LR{$A \cup B$} او هم په
\LR{$A \cup C$} کښې دے،
د~\LR{$\cap$} له تعريفه۔
د \LR{$\cup$} له تعريفه
دا نور هم پرانيستلے شو:
يا \LR{$z \in A$} يا
\LR{$z \in B$}؛ او هم
يا \LR{$z \in A$} يا \LR{$z \in
C$}۔ دا بيا د حالتونو له مخې
ثبوت ښکاري، خو «او» يې
يو څه مغشوشوي۔ ښايي فکر
وکړئ چې درې امکانونه دي:
\LR{$z$} يا په \LR{$A$}، يا په
\LR{$B$}، يا په \LR{$C$} کښې دے۔
خو دا تېروتنه ده۔ بايد
پام وکړو او هر فصل
جلا وڅېړو۔
\end{quote}
د \LR{$\cap$} له تعريفه،
\LR{$z \in A \cup B$} او
\LR{$z \in A \cup C$}۔
د \LR{$\cup$} له تعريفه،
\LR{$z \in A$} يا \LR{$z \in B$}۔
حالتونه بېلوو۔
\begin{quote}
ځکه چې اوس په لومړي فصل
تمرکز کوو، دويم فصل
(چې د \LR{$A \cup C$} له
پرانيستلو راځي) مو لا
نه دے راوړے؛ په حقيقت
کښې لا ورته اړتيا هم
نشته۔ لومړے حالت
\LR{$z \in A$} دے، او د يوه
سټ غړے د هماغه
سټ او هر بل سټ د اتحاد
غړے هم وي۔
نو لومړے حالت اسان دے:
\end{quote}
حالت ۱: فرض کړئ
\LR{$z \in A$}۔ له دې څخه
\LR{$z \in A \cup (B \cap
C)$} راځي۔
\begin{quote}
اوس دويم حالت،
\LR{$z \in B$}، اخلو۔
دلته دويم \LR{$\cup$}
پرانيزو او يو بل
د حالتونو له مخې
ثبوت کوو:
\end{quote}
حالت ۲: فرض کړئ
\LR{$z \in B$}۔ څرنګه چې
\LR{$z \in A \cup C$}،
يا \LR{$z \in A$} يا
\LR{$z \in C$}۔ حالتونه
نور هم بېلوو:
\olpseganchor{OLP-0607-B012}{end}

\olpseganchor{OLP-0607-B013}{start}
حالت ۲الف: \LR{$z \in A$}۔
نو بيا هم \LR{$z \in A \cup (B \cap C)$}۔
\begin{quote}
دا يو څه عجيبه و: په دې
حالت کښې مو د~\LR{$z \in B$}
فرض په رښتيا نه و پکار؛
خو ستونزه نشته۔
\end{quote}
حالت ۲ب: \LR{$z \in C$}۔
نو \LR{$z \in B$} او
\LR{$z \in C$}، ځکه
\LR{$z \in B \cap C$}؛
او له دې څخه
\LR{$z \in A \cup (B \cap C)$}۔
\begin{quote}
دا د حالتونو دواړه
ثبوتونه بشپړوي؛ نو
دويمه نيمايي هم
ثابته شوه۔
\end{quote}
نو که
\LR{$z \in (A \cup B) \cap (A \cup C)$}،
بيا \LR{$z \in A \cup (B \cap C)$}۔
\end{proof}
\olpseganchor{OLP-0607-B013}{end}

\olpseganchor{OLP-0608-B004}{start}
\olfileid{mth}{prf}{ex2}
\olpseganchor{OLP-0608-B004}{end}

\olpseganchor{OLP-0608-B005}{start}
\olsection{بله بېلګه}
\olpseganchor{OLP-0608-B005}{end}

\olpseganchor{OLP-0608-B006}{start}
\begin{prop}
که \LR{$A \subseteq C$} وي، نو \LR{$A \cup (C \setminus A) = C$}۔
\end{prop}  
\olpseganchor{OLP-0608-B006}{end}

\olpseganchor{OLP-0608-B007}{start}
\begin{proof}
  فرض کړئ چې \LR{$A \subseteq C$}۔ غواړو وښيو چې \LR{$A \cup (C
  \setminus A) = C$}۔
  \begin{quote}
    لومړے ګورو چې دا يوه شرطي دعوه ده۔ په ضمني ډول پرې عمومي
    کمیت ټاکل شوے دے: دا قضيه د ټولو سټونو \LR{$A$} او \LR{$C$} دپاره
    ټينګېږي۔ نو \LR{$A$} او \LR{$C$} د هر ډول سټونو متغيرونه دي۔ د داسې
    دعوې د ثابتولو دپاره مقدم فرضوو او نتيجه ثابتوو۔
\olpseganchor{OLP-0608-B007}{end}

\olpseganchor{OLP-0608-B008}{start}
    اوس د \LR{$A \subseteq C$} فرض کاروو۔ د~\LR{$\subseteq$} تعريف
    پرانيزو: فرض دا معنا لري چې د~\LR{$A$} ټول غړي د~\LR{$C$}
    غړي هم دي۔ دا وليکو؛ په ټول ثبوت کښې به ترې کار اخلو۔
  \end{quote}
  د~\LR{$\subseteq$} له تعريفه، څرنګه چې \LR{$A \subseteq C$}، د هر
  \LR{$z$} دپاره، که \LR{$z \in A$} وي، نو \LR{$z \in C$} وي۔
  \begin{quote}
    د فرض ټول ورکړل شوي تعريفونه مو پرانيستل۔ اوس پايلې ته ځو۔
    غواړو وښيو چې \LR{$A \cup (C \setminus A) = C$}؛ نو د تېرې
    بېلګې په شان ثبوت ترتيب کوو: ښيو چې د
    \LR{$A \cup (C \setminus A)$} هر غړے د~\LR{$C$} غړے
    هم دے، او برعکس د \LR{$C$} هر غړے د \LR{$A \cup (C
    \setminus A)$} غړے هم دے۔ په لنډه يې داسې ليکو:
    \LR{$A \cup (C \setminus A) \subseteq C$} او
    \LR{$C \subseteq A \cup (C \setminus A)$}۔ (دلته تعريف د
    پرانيستلو پر ځاے په يوه لنډه بڼه راغونډوو، څو ثبوت اسان
    ولوستل شي۔) دا عطف دے؛ نو دواړه برخې ثابتول پکار دي۔
    د لومړۍ برخې، يعنې دا چې د \LR{$A \cup (C \setminus A)$} هر
    غړے د~\LR{$C$} غړے هم دے، د ثابتولو دپاره
    \LR{$z \in A \cup (C \setminus A)$} د يوه هر ډول~\LR{$z$} دپاره
    فرضوو او ښيو چې \LR{$z \in C$}۔
    د \LR{$\cup$} له تعريفه اخلو چې \LR{$z \in A$} يا
    \LR{$z \in C \setminus A$}؛ دا له
    \LR{$z \in A \cup (C \setminus A)$} څخه راځي۔ اوس به مو د
    دې لارې چل زده شوے وي۔
    \end{quote}
  \LR{$A \cup (C \setminus A) = C$} هغه وخت او يوازې هغه وخت وي
  چې \LR{$A \cup (C \setminus A) \subseteq
  C$} او \LR{$C \subseteq (A \cup (C \setminus A))$} وي۔
  لومړے ثابتوو چې \LR{$A \cup (C \setminus A) \subseteq C$}۔
  \LR{$z \in A \cup (C
  \setminus A)$} دې وي۔ نو يا \LR{$z \in A$} يا \LR{$z \in (C \setminus A)$}۔
  \begin{quote}
    يوه فصل ته ورسېدو؛ له دې څخه غواړو \LR{$z \in C$} ثابت کړو۔
    دا د حالتونو له مخې ثبوت کوو۔
  \end{quote}
  حالت ۱: \LR{$z \in A$}۔ ځکه چې د هر \LR{$z$} دپاره، که \LR{$z \in A$} وي،
  نو \LR{$z \in C$} وي، دلته هم \LR{$z \in C$} لرو۔
  \begin{quote}
    دلته مو هغه مخکې ليکلے حقيقت وکاراوه چې د قضیې له
    \LR{$A \subseteq C$} فرضه راغلے وو۔ لومړے حالت بشپړ شو؛ اوس
    دويم حالت، \LR{$z \in (C
    \setminus A)$}، ګورو۔ په ياد ولرئ چې \LR{$C \setminus A$} د
    دوو سټونو \emph{تفاضل} ښيي؛ يعنې د~\LR{$C$} د هغو ټولو
    غړي سټ چې د~\LR{$A$} غړي نه دي۔ خو د \LR{$C$}
    هر غړے چې په~\LR{$A$} کښې نه وي، په ځانګړې توګه
    د~\LR{$C$} غړے خو وي۔
  \end{quote}
  حالت ۲: \LR{$z \in (C \setminus A)$}۔ دا معنا لري چې \LR{$z \in C$}
  او \LR{$z \notin A$}۔ نو په ځانګړې توګه \LR{$z \in C$}۔
  \begin{quote}
    ښه، لومړے لوری مو ثابت کړ۔ اوس دويم لوری: دلته
    \LR{$C \subseteq A \cup (C \setminus
    A)$} ثابتوو۔ نو فرضوو چې \LR{$z \in C$} او ښيو چې
    \LR{$z \in A \cup (C \setminus A)$}۔
  \end{quote}
  اوس \LR{$z \in C$} دې وي۔ غواړو وښيو چې \LR{$z \in A$} يا
  \LR{$z \in C \setminus A$}۔
  \begin{quote}
    د \LR{$A$} ټول غړي د \LR{$C$} غړي هم دي؛ او
    \LR{$C \setminus A$} د هغو ټولو شيانو سټ دے چې د \LR{$C$}
    غړي دي خو د \LR{$A$} نه دي۔ له دې ښکاري چې \LR{$z$}
    يا په \LR{$A$} کښې دے يا په \LR{$C
    \setminus A$} کښې۔ که لا مو د دې پايلې وجه نه وي پېژندلې،
    دا ګام ښايي روښانه نه وي۔ غوره ده چې ګام په ګام يې ثابت
    کړو۔ يو ساده حقيقت کار راکوي چې بې له ثبوته يې ويلے شو:
    \LR{$z \in A$} يا \LR{$z \notin A$}۔ دې ته «د منځني حالت د ايستلو
    اصل» وايي: د هر بيان~\LR{$p$} دپاره، يا \LR{$p$} رښتيا وي يا يې نفي
    رښتيا وي۔ (دلته \LR{$p$} دا بيان دے چې \LR{$z \in A$}۔) دا فصل دے؛
    نو بيا د حالتونو له مخې ثبوت کارولے شو۔
  \end{quote}
  يا \LR{$z \in A$} يا \LR{$z \notin A$}۔ په لومړي حالت کښې
  \LR{$z \in A \cup
  (C \setminus A)$}۔ په دويم حالت کښې \LR{$z \in C$} او
  \LR{$z \notin A$}؛ نو \LR{$z \in C \setminus A$}۔ بيا
  \LR{$z \in A \cup (C \setminus A)$} هم دے۔
  \begin{quote}
    ثبوت بشپړ شو: ومو ښودل چې \LR{$A \cup (C \setminus A) = C$}۔
  \end{quote}
\end{proof}
\olpseganchor{OLP-0608-B008}{end}

\olpseganchor{OLP-0609-B004}{start}
\olfileid{mth}{prf}{con}
\olpseganchor{OLP-0609-B004}{end}

\olpseganchor{OLP-0609-B005}{start}
\olsection{د تناقض له لارې ثبوت}
\olpseganchor{OLP-0609-B005}{end}

\olpseganchor{OLP-0609-B006}{start}
په لومړي سر کښې د تناقض له لارې ثبوت د استنتاج يوه بڼه
ده چې منفي دعوې پرې ثابتېږي۔ فرض کړئ غواړئ وښيئ چې کومه
دعوه~\LR{$p$} \emph{ناسمه} ده؛ يعنې غواړئ~\LR{$\lnot p$} ثابت کړئ۔
تر ټولو غوره لاره دا ده: (الف) فرض کړئ چې \LR{$p$} رښتيا ده،
او (ب) وښيئ چې له دې فرضه داسې پايله راځي چې ناسموالی
يې درته څرګند دے۔ «د څرګند ناسموالي پايله» ښايي له
خپلې~\LR{$p$} سره يا د کتل کېدونکې ټولې دعوې له کوم بل
فرض سره په ټکر کښې وي، يعنې تناقض جوړ کړي۔ د بېلګې
په توګه، د «که \LR{$q$} وي نو \LR{$\lnot p$}» ثبوت کښې \LR{$q$}
رښتيا فرضېږي او له هغې~\LR{$\lnot
p$} ثابتېږي۔ که \LR{$\lnot p$} د تناقض له لارې ثابتوئ،
دا د \LR{$p$} فرضول دي، د~\LR{$q$} ترڅنګ۔ که \LR{$\lnot q$}
له \LR{$p$} څخه ثابت کړے شئ، ښيئ چې د~\LR{$p$} له فرضه
ستاسې له بل فرض~\LR{$q$} سره په ټکر کښې پايله راځي، ځکه
\LR{$q$} او \LR{$\lnot q$} دواړه په يو وخت رښتيا کېدلے نه شي۔
البته، د تناقض د ثبوت په بهير کښې نورې استنتاجي بڼې او
د تعريفونو پرانيستل هم کارول پکار دي۔ يوه بېلګه ګورو۔
\olpseganchor{OLP-0609-B006}{end}

\olpseganchor{OLP-0609-B007}{start}
\begin{prop}
  که \LR{$A \subseteq B$} او \LR{$B = \emptyset$} وي، نو \LR{$A$} هېڅ
  غړي نه لري۔
\end{prop}
\olpseganchor{OLP-0609-B007}{end}

\olpseganchor{OLP-0609-B008}{start}
\begin{proof}
  فرض کړئ \LR{$A \subseteq B$} او \LR{$B = \emptyset$}۔ غواړو وښيو چې
  \LR{$A$} هېڅ غړي نه لري۔
  \begin{quote}
    دا شرطي دعوه ده؛ نو مقدم فرضوو او نتيجه ثابتوو۔ نتيجه
    دا ده چې \LR{$A$} هېڅ غړي نه لري۔ په څرګنده بڼه:
    داسې~\LR{$x \in A$} نشته۔
  \end{quote}
  \LR{$A$} هېڅ غړي نه لري هغه وخت او يوازې هغه وخت
  چې داسې~\LR{$x$} نشته چې \LR{$x \in A$} وي۔
  \begin{quote}
    نو اصلاً يوه منفي دعوه~\LR{$\lnot p$} ثابتول غواړو: داسې
    \LR{$x \in A$} نشته۔ د تناقض له لارې ثبوت دپاره يې مقابل
    مثبت حکم~\LR{$p$} فرضوو، يعنې داسې \LR{$x \in A$} شته، او
    له دې څخه تناقض ثابتوو۔ دا لاره په «د تناقض دپاره
    فرض کړئ» يا لنډ «برعکس فرض کړئ» پېلوو او بيا
    فرض~\LR{$p$} ليکو۔
  \end{quote}
  برعکس فرض کړئ: داسې \LR{$x \in A$} شته۔
  \begin{quote}
    دا اوس هغه نوے فرض دے چې تناقض ترې اخلو۔ دوه نور
    فرضونه هم لرو: \LR{$A \subseteq B$} او \LR{$B = \emptyset$}۔
    له لومړي څخه \LR{$x \in B$} راځي:
  \end{quote}
  څرنګه چې \LR{$A \subseteq B$}، نو \LR{$x \in B$}۔
  \begin{quote}
    خو څرنګه چې \LR{$B = \emptyset$}، د \LR{$B$} هر غړے،
    لکه \LR{$x$}، بايد د~\LR{$\emptyset$} غړے هم وي۔
  \end{quote}
  څرنګه چې \LR{$B = \emptyset$}، نو \LR{$x \in \emptyset$}۔ دا تناقض
  دے، ځکه د تعريف له مخې \LR{$\emptyset$} هېڅ غړي
  نه لري۔
  \begin{quote}
    ثبوت همدلته بشپړېږي: له ټاکلو فرضونو څخه تناقض ته
    ورسېدو، نو ټول فرضونه يو ځای رښتيا نه شي کېدلے۔
    لومړي دوه فرضونه، \LR{$A \subseteq B$} او \LR{$B = \emptyset$}،
    د قضیې ورکړل شوي شرطونه دي؛ نو وروستی زيات شوے فرض،
    دا چې داسې \LR{$x \in A$} شته، بايد ناسم وي۔ که بشپړ
    تفصيل غواړو، دا پايله ښکاره ليکلے شو۔
  \end{quote}
  نو زموږ دا فرض چې داسې \LR{$x \in A$} شته ناسم دے؛ له دې امله
  \LR{$A$} د تناقض له لارې د ثبوت له مخې هېڅ غړي نه لري۔
\end{proof}
\olpseganchor{OLP-0609-B008}{end}

\olpseganchor{OLP-0609-B009}{start}
هره مثبته دعوه په څرګنده توګه له يوې منفي
دعوې سره برابره ده: \LR{$p$} هغه وخت او يوازې هغه وخت چې
\LR{$\lnot\lnot p$}۔ نو د تناقض له لارې ثبوت په نامستقيم ډول د
مثبتو دعوو د ثابتولو دپاره هم کارېدلے شي۔ د~\LR{$p$} د ثبوت
دپاره يې د منفي دعوې \LR{$\lnot\lnot p$} په توګه واخلئ۔
که له \LR{$\lnot p$} څخه تناقض ثابت کړو، د تناقض له لارې
مو \LR{$\lnot\lnot p$} ثابت کړ، او له هغې~\LR{$p$}۔
\olpseganchor{OLP-0609-B009}{end}

\olpseganchor{OLP-0609-B010}{start}
په تېرې بېلګې کښې مو يوه منفي دعوه ثابتول غوښتل:
دا چې \LR{$A$} هېڅ غړي نه لري۔ نو د تناقض د ثبوت
دپاره مو مثبت فرض ومانۀ، يعنې داسې \LR{$x \in A$} شته۔
دې موږ ته د کار دپاره فرضي \LR{$x \in A$} راکړ؛ ورپسې مو
تر \LR{$x \in \emptyset$} پورې د هغه په اړه استدلال وکړ۔
\olpseganchor{OLP-0609-B010}{end}

\olpseganchor{OLP-0609-B011}{start}
کله چې يوه مثبته دعوه په نامستقيم ډول ثابتوئ، د
تناقض دپاره منفي فرض نيسئ۔ خو ډېر ځله مثبته دعوه په
منفي بڼه او منفي دعوه په مثبته بڼه په اسانۍ اړولے شئ۔
زموږ تېر ثبوت به تقريباً هماغسې وو که مو
«\LR{$A = \emptyset$}» ثابتول غوښتل، د «\LR{$A$} هېڅ
غړي نه لري» د منفي نتيجې پر ځاے۔
(د \LR{$=$} له تعريفه «\LR{$A = \emptyset$}» يوه
عمومي دعوه ده: «د~\LR{$A$} هر غړے د~\LR{$\emptyset$}
غړے دے او برعکس».) خو په اسانۍ ښکاري چې
دا له «نه: داسې \LR{$x \in A$} شته» سره برابره ده۔
\olpseganchor{OLP-0609-B011}{end}

\olpseganchor{OLP-0609-B012}{start}
نو کله ناکله له \LR{$\lnot p$} سره د فرض په توګه کار
کول د~\LR{$p$} له نېغ ثبوته اسان وي۔ که نېغ
ثبوت همدومره اسان يا تر دې اسان هم وي، لکه راتلونکو
بېلګو کښې، ځينې خلک نامستقيمه لاره خوښوي۔ که دوه
ځله نفي مو مغشوشوي، د يوې دعوې د تناقض له لارې ثبوت
د هغې د \emph{مقابلې} دعوې له فرضه د تناقض اخيستل وګڼئ۔
نو د \LR{$\lnot
p$} د تناقض له لارې ثبوت د~\LR{$p$} له فرضه تناقض
ثابتوي؛ او د~\LR{$p$} د تناقض له لارې ثبوت له~\LR{$\lnot p$}
څخه تناقض ثابتوي۔
\olpseganchor{OLP-0609-B012}{end}

\olpseganchor{OLP-0609-B013}{start}
\begin{prop}
\LR{$A \subseteq A \cup B$}۔
\end{prop}
\olpseganchor{OLP-0609-B013}{end}

\olpseganchor{OLP-0609-B014}{start}
\begin{proof}
  غواړو وښيو چې \LR{$A \subseteq A \cup B$}۔
  \begin{quote}
    په ظاهر کښې دا مثبته دعوه ده: هر \LR{$x \in A$} د
    \LR{$A \cup B$} غړے هم دے۔ نفي يې دا ده: ځينې
    \LR{$x \in A$}، \LR{$\notin A \cup B$} دي۔ نو دا منفي دعوه فرضولے شو
    او ښودلے شو چې تناقض ته رسېږي؛ په دې توګه دعوه
    په نامستقيم ډول ثابتېږي۔
    \end{quote}
  برعکس فرض کړئ، يعنې \LR{$A \nsubseteq A \cup B$}۔
  \begin{quote}
    د \LR{$A \subseteq A \cup B$} تعريف لرو: هر \LR{$x \in A$}،
    \LR{$\in A \cup B$} هم دے۔ د
    \LR{$A \nsubseteq A \cup B$} د معنا د پوهېدو دپاره
    د پرانيستي تعريف ساده منطقي اړونه کاروو: دا خبره
    ناسمه ده چې هر \LR{$x \in A$}، \LR{$\in
    A \cup B$} هم دے، هغه وخت او يوازې هغه
    وخت چې \emph{يو}~\LR{$x \in A$}، \LR{$\notin A \cup B$} وي۔
    (په دې ځاے کښې ډېر پام پکار دے:
    د زده کوونکو د تناقض د ثبوت ډېرې هڅې د «ټول
    \LR{$A$}s د \LR{$B$}s دي» په څېر دعوې د نفي له ناسمه
    شننه لوېږي۔) په بل عبارت،
    \LR{$A \nsubseteq A \cup B$} هغه وخت او يوازې هغه وخت
    چې داسې \LR{$x$} وي چې \LR{$x \in A$} او
    \LR{$x \notin A \cup B$}۔ له دې وروسته کار اسان دے۔
  \end{quote}
  نو داسې \LR{$x \in A$} شته چې \LR{$x \notin A \cup B$}۔
  د \LR{$\cup$} له تعريفه، \LR{$x \in A \cup B$} هغه وخت او
  يوازې هغه وخت چې \LR{$x \in A$} يا \LR{$x \in
  B$} وي۔ څرنګه چې \LR{$x \in A$}، نو \LR{$x \in A \cup B$}۔
  دا له هغه فرض سره په ټکر کښې دے چې
  \LR{$x \notin A \cup B$}۔
\end{proof}
\olpseganchor{OLP-0609-B014}{end}

\olpseganchor{OLP-0609-B015}{start}
\begin{prob}
په \emph{نامستقيم ډول} ثابت کړئ چې \LR{$A \cap B \subseteq A$}۔
\end{prob}
\olpseganchor{OLP-0609-B015}{end}

\olpseganchor{OLP-0609-B016}{start}
\begin{prop}
که \LR{$A \subseteq B$} او \LR{$B \subseteq C$} وي، نو \LR{$A \subseteq C$}۔
\end{prop}
\olpseganchor{OLP-0609-B016}{end}

\olpseganchor{OLP-0609-B017}{start}
\begin{proof}
  فرض کړئ \LR{$A \subseteq B$} او \LR{$B \subseteq C$}۔ غواړو وښيو چې
  \LR{$A
  \subseteq C$}۔
  \begin{quote}
    نامستقيمه لاره نيسو: د غوښتل شوې پايلې نفي فرضوو۔
  \end{quote}
  برعکس فرض کړئ، يعنې \LR{$A \nsubseteq C$}۔
  \begin{quote}
    د تېر په شان، \LR{$A \nsubseteq C$} معنا دا ده چې دا خبره
    رښتيا نۀ ده چې هر \LR{$x \in A$}، \LR{$\in C$} هم وي؛
    يعنې ځينې \LR{$x \in A$}، \LR{$\notin C$} دي۔ اندېښنه مه کوئ:
    په تمرين سره به د داسې نفي د پرانيستلو دپاره
    ډېر فکر ته اړ نه شئ۔
  \end{quote}
  په بل عبارت، داسې~\LR{$x$} شته چې \LR{$x \in A$} او
  \LR{$x \notin C$}۔
  \begin{quote}
    اوس ترې تناقض اخيستلے شو۔ البته، د دې دپاره
    نور دوه فرضونه هم کاروو۔
  \end{quote}
  څرنګه چې \LR{$A \subseteq B$}، نو \LR{$x \in B$}۔ څرنګه چې
  \LR{$B \subseteq C$}، نو \LR{$x \in
  C$}۔ خو دا له \LR{$x \notin C$} سره تناقض دے۔
\end{proof}
\olpseganchor{OLP-0609-B017}{end}

\olpseganchor{OLP-0609-B018}{start}
\begin{prop}
که \LR{$A \cup B = A \cap B$} وي، نو \LR{$A = B$}۔
\end{prop}
\olpseganchor{OLP-0609-B018}{end}

\olpseganchor{OLP-0609-B019}{start}
\begin{proof}
  فرض کړئ \LR{$A \cup B = A \cap B$}۔ غواړو وښيو چې \LR{$A = B$}۔
  \begin{quote}
    پيل اوس عادي دے:
  \end{quote}
  د تناقض دپاره فرض کړئ چې \LR{$A \neq B$}۔
  \begin{quote}
    د تناقض د ثبوت فرض دا دے چې \LR{$A \neq
    B$}۔ څرنګه چې \LR{$A = B$} هغه وخت او يوازې هغه وخت
    چې \LR{$A \subseteq B$} او \LR{$B \subseteq A$}، نو
    \LR{$A \neq B$} هغه وخت او يوازې هغه وخت چې
    \LR{$A \nsubseteq B$} \emph{يا} \LR{$B \nsubseteq
    A$}۔ (وګورئ چې د نفيو په اړولو کښې پام څومره اړين
    دے!{}) له دې فصله د تناقض د اخيستلو دپاره د حالتونو
    له مخې ثبوت کاروو او ښيو چې په هر حالت کښې تناقض
    راځي۔
  \end{quote}
  \LR{$A \neq B$} هغه وخت او يوازې هغه وخت چې
  \LR{$A \nsubseteq B$} يا \LR{$B \nsubseteq A$}۔ حالتونه بېلوو۔
  \begin{quote}
    په لومړي حالت کښې \LR{$A \nsubseteq B$} فرضوو؛ يعنې
    د ځينو \LR{$x$} دپاره \LR{$x \in A$} خو \LR{$\notin B$}۔ \LR{$A \cap B$} د هغو
    غړي سټ دے چې \LR{$A$} او \LR{$B$} دواړه يې لري؛
    نو که يو څيز په يوه کښې نه وي، په تقاطع کښې هم
    نه وي۔ \LR{$A \cup B$} د \LR{$A$} او \LR{$B$} يوځای کول دي؛
    نو د دواړو له هر يوه غړے په اتحاد کښې هم وي۔
    له دې اخلو چې \LR{$x \in A \cup B$} خو
    \LR{$x \notin A \cap
    B$}، او په پايله کښې \LR{$A \cap B \neq A \cup B$}۔
  \end{quote}
\olpseganchor{OLP-0609-B019}{end}

\olpseganchor{OLP-0609-B020}{start}
  حالت ۱: \LR{$A \nsubseteq B$}۔ نو د ځينو \LR{$x$} دپاره
  \LR{$x \in A$} خو \LR{$x \notin
  B$}۔ څرنګه چې \LR{$x \notin B$}، نو
  \LR{$x \notin A \cap B$}۔ څرنګه چې \LR{$x \in A$}،
  \LR{$x \in A \cup B$}۔ نو \LR{$A \cap B \neq A \cup B$}؛
  دا له هغه فرض سره په ټکر کښې دے چې
  \LR{$A \cap B = A \cup B$}۔
\olpseganchor{OLP-0609-B020}{end}

\olpseganchor{OLP-0609-B021}{start}
  حالت ۲: \LR{$B \nsubseteq A$}۔ نو د ځينو \LR{$y$} دپاره
  \LR{$y \in B$} خو \LR{$y \notin
  A$}۔ لکه مخکې، \LR{$y \in A \cup B$} لرو خو
  \LR{$y \notin A \cap B$}؛ نو \LR{$A \cap B \neq A \cup B$}،
  او دا بيا له \LR{$A \cap B = A \cup
  B$} سره په ټکر کښې دے۔
\end{proof}
\olpseganchor{OLP-0609-B021}{end}

\olpseganchor{OLP-0610-B004}{start}
\olfileid{mth}{prf}{rea}
\olsection{د ثبوتونو لوستل}
\olpseganchor{OLP-0610-B004}{end}

\olpseganchor{OLP-0610-B005}{start}
په درسي کتابونو او مقالو کښې موندل کېدونکي ثبوتونه ډېر کم هغه
ټول تفصيلات ليکي چې موږ تر اوسه په خپلو بېلګو کښې شامل کړي دي۔
ليکوال ډېر ځله نۀ څرګندوي چې کله حالتونه بېلوي، کله نامستقيم
ثبوت ورکوي، يا دا هم نۀ يادوي چې يو تعريف کاروي۔ نو په درسي
کتاب کښې د ثبوت د لوستلو پر وخت به ډېر ځله د هغه د پوهېدو
دپاره دا تفصيلات پخپله بشپړوئ۔ دا کار د هغو بېلابېلو ګامونو
د زده کولو ښه تمرين هم دے چې په ثبوت کښې ئې اخيستل پکار دي۔
يوه بېلګه ګورو۔
\olpseganchor{OLP-0610-B005}{end}

\olpseganchor{OLP-0610-B006}{start}
\begin{prop}[جذب]
د ټولو سټونو \LR{$A$}، \LR{$B$} دپاره:
\[
A \cap (A \cup B) = A
\]
\end{prop}
\olpseganchor{OLP-0610-B006}{end}

\olpseganchor{OLP-0610-B007}{start}
\begin{proof}
که \LR{$z \in A \cap (A \cup B)$} وي، نو \LR{$z \in A$}؛ له دې امله
\LR{$A \cap (A \cup B)
\subseteq A$}۔ اوس \LR{$z \in A$} فرض کړئ۔ بيا \LR{$z \in A \cup B$} هم
دے، او ځکه \LR{$z \in A \cap (A \cup B)$} هم دے۔
\end{proof}
\olpseganchor{OLP-0610-B007}{end}

\olpseganchor{OLP-0610-B008}{start}
د جذب د قانون پورتنے ثبوت ډېر لنډ شوے دے۔ نۀ د کارول شوو
تعريفونو يادونه لري، نۀ د يوې خبرې تر ثابتولو وړاندې «دا
ثابتول پکار دي» وايي، او داسې نور۔ راځئ پرانيزو ئې۔ ثابته
شوې قضيه د هر ډول سټونو \LR{$A$} او \LR{$B$} په اړه عمومي دعوه ده؛
نو کله چې ثبوت \LR{$A$} يا \LR{$B$} يادوي، دا د اختياري سټونو متغيرونه
دي۔ ثبوت هغه عمومي دعوې ثابتوي چې د سټونو د برابرۍ دپاره
اړينې دي: يعنې د برابرۍ د چپ اړخ هر غړے د ښي اړخ
غړے دے، او برعکس۔
\olpseganchor{OLP-0610-B008}{end}

\olpseganchor{OLP-0610-B009}{start}
\begin{quote}
«که \LR{$z \in A \cap (A \cup B)$} وي، نو \LR{$z \in A$}؛ له دې امله
  \LR{$A \cap (A \cup B)
  \subseteq A$}۔»
\end{quote}
\olpseganchor{OLP-0610-B009}{end}

\olpseganchor{OLP-0610-B010}{start}
دا د برابرۍ د ثبوت لومړے نيم دے: ثابتوي چې که اختياري~\LR{$z$}
د چپ اړخ غړے وي، نو د ښي اړخ غړے هم دے؛
يعنې \LR{$A \cap (A \cup B) \subseteq A$}۔ فرض کړئ
\LR{$z \in A \cap (A \cup B)$}۔ څرنګه چې \LR{$z$} د دوو سټونو د تقاطع
غړے هغه وخت او يوازې هغه وخت دے چې د دواړو سټونو
غړے وي، نو \LR{$z \in A$} او \LR{$z \in A \cup B$} دواړه اخلو۔
په ځانګړي ډول \LR{$z \in A$}، چې همدا مو ثابتول غوښتل۔ څرنګه چې
د لومړي نيم دپاره همدومره کار پکار دے، پوهېږو چې پاتې ثبوت
بايد د دويم نيم ثبوت وي؛ يعنې دا ثابت کړي چې
\LR{$A \subseteq A \cap
(A \cup B)$}۔
\olpseganchor{OLP-0610-B010}{end}

\olpseganchor{OLP-0610-B011}{start}
\begin{quote}
«اوس \LR{$z \in A$} فرض کړئ۔ بيا \LR{$z \in A \cup B$} هم دے،
او ځکه \LR{$z \in A \cap (A \cup B)$} هم دے۔»
\end{quote}
\olpseganchor{OLP-0610-B011}{end}

\olpseganchor{OLP-0610-B012}{start}
په \LR{$z \in A$} فرضولو پيل کوو، ځکه ښيو چې د هر~\LR{$z$} دپاره
که \LR{$z \in A$} وي، نو \LR{$z \in A \cap (A \cup B)$}۔ د دې
د ښودلو دپاره چې \LR{$z
\in A \cap (A \cup B)$}، د «\LR{$\cap$}» د تعريف له مخې بايد
دواړه وښيو: (i) \LR{$z \in A$} او (ii) \LR{$z \in A \cup B$}۔
دلته (i) هماغه زموږ فرض دے؛ نو نور ثبوت ته اړتيا نۀ لري،
او همدا وجه ده چې ثبوت ئې بيا نۀ يادوي۔ د (ii) دپاره
راياد کړئ چې \LR{$z$} د سټونو د اتحاد غړے هغه وخت او
يوازې هغه وخت دے چې لږ تر لږه د يوه سټ غړے وي۔
څرنګه چې \LR{$z \in A$} او \LR{$A \cup B$} د \LR{$A$} او \LR{$B$} اتحاد دے،
دا شرط دلته پوره دے۔ نو \LR{$z \in A \cup B$}۔ موږ دواړه
ښودلي دي: (i) \LR{$z \in A$} او (ii) \LR{$z \in A \cup B$}؛ ځکه
د «\LR{$\cap$}» له تعريفه \LR{$z \in A \cap (A \cup B)$}۔
ثبوت دغه تعريفونه نۀ يادوي؛ فرض شوې ده چې لوستونکے له
مخکې پرې ښه پوهېږي۔ که تاسو پرې نۀ پوهېږئ، بايد بېرته
ورته لاړ شئ او راياد ئې کړئ۔ بيا به دا هم پېژنئ چې ولې
له \LR{$z
\in A$} څخه \LR{$z \in A \cup B$} راځي، او له \LR{$z \in A$} او
\LR{$z \in A \cup B$} څخه \LR{$z \in A \cap (A \cup B)$} راځي۔
\olpseganchor{OLP-0610-B012}{end}

\olpseganchor{OLP-0610-B013}{start}
دا د پورتني ثبوت بله بڼه ده، چې هر څه پکښې ښکاره شوي دي:
\begin{proof}{}
[د سټونو د \LR{$=$} له تعريفه، د \LR{$A \cap (A \cup B) = A$} دپاره
بايد (a) \LR{$A \cap (A \cup B) \subseteq A$} او
(b) \LR{$A \subseteq A \cap (A \cup B)$} وښيو۔ (a): د \LR{$\subseteq$} له تعريفه
بايد وښيو چې که \LR{$z \in A \cap (A \cup B)$} وي، نو
\LR{$z \in A$}۔] که \LR{$z \in A
\cap (A \cup B)$} وي، نو \LR{$z \in A$} [ځکه د \LR{$\cap$} له تعريفه
\LR{$z
  \in A \cap (A \cup B)$} هغه وخت او يوازې هغه وخت چې
\LR{$z \in A$} او \LR{$z \in A \cup B$}]؛ نو
\LR{$A
\cap (A \cup B) \subseteq A$}۔ [(b): د \LR{$\subseteq$} له تعريفه
بايد وښيو چې که \LR{$z \in A$} وي، نو
\LR{$z \in A \cap (A \cup B)$}۔] اوس [(1)] \LR{$z \in A$} فرض کړئ۔
بيا [(2)] \LR{$z \in A \cup B$} هم دے [ځکه له (1) څخه
\LR{$z \in A$} يا \LR{$z \in B$}، چې د \LR{$\cup$} له تعريفه
د \LR{$z
  \in A \cup B$} معنا لري]؛ او ځکه \LR{$z \in A \cap (A \cup B)$}
هم دے [ځکه د \LR{$\cap$} تعريف دواړه غواړي: \LR{$z \in A$}، يعنې
(1)، او \LR{$z
  \in A \cup B$}، يعنې (2)۔]
\end{proof}
\olpseganchor{OLP-0610-B013}{end}

\olpseganchor{OLP-0610-B014}{start}
\begin{prob}
د \LR{$A \cup (A \cap B) = A$} لاندې ثبوت پراخ کړئ: ټولې
کارول شوې استنتاجي بڼې يادې کړئ، وښيئ چې هر ګام ولې
له فرضونو يا له مخکې ثابتو شوو دعوو راځي، او په کوم
ځاے کښې کوم تعريف ته اړتيا لرو۔
\begin{proof}
  که \LR{$z \in A \cup (A \cap B)$} وي، نو \LR{$z \in A$} يا
  \LR{$z \in A \cap B$}۔ که \LR{$z \in A \cap B$} وي، \LR{$z \in A$}۔
  هر \LR{$z \in A$}، \LR{$\in A \cup (A
  \cap B)$} هم دے۔
\end{proof}
\end{prob}
\olpseganchor{OLP-0610-B014}{end}

\olpseganchor{OLP-0611-B004}{start}
\olfileid{mth}{prf}{cnt}
\olpseganchor{OLP-0611-B004}{end}

\olpseganchor{OLP-0611-B005}{start}
\olsection{زه ئې نۀ شم کولے!{}}
\olpseganchor{OLP-0611-B005}{end}

\olpseganchor{OLP-0611-B006}{start}
موږ ټول داسې پړاو ته رسېږو چې د پرېښودلو زړه مو کېږي۔
خو تاسو دا کار \emph{کولے شئ}۔ استاد، د تدريس مرستيال
او ستاسو ملګري زده کوونکي مرسته کولے شي۔ مرسته ترې
وغواړئ!{} دلته څو مشورې دي چې له سخت حالت څخه مو وساتي،
او دا هم ښيي چې د پرېښودلو د احساس پر وخت څه وکړئ۔
\olpseganchor{OLP-0611-B006}{end}

\olpseganchor{OLP-0611-B007}{start}
د دې دپاره چې مسئلې په برياليتوب حل کړے شئ، دا کارونه وکړئ:
\begin{enumerate}
\item \emph{څومره چې کېدلے شي، هومره وختي پيل وکړئ۔}
  په ټول سمسټر کښې بوختېږو او زموږ ډېرو ته د کار ځنډول
  ستونزه وي؛ يو له غوره کارونو دا دے چې کورنۍ دندې
  وختي پيل کړئ۔ په دې توګه، که ايسار شئ، د حل د لټولو
  وخت به لرئ (چې ژړا نۀ وي)۔
\item \emph{له خپلو ټولګيوالو سره خبرې وکړئ}۔ تاسو يوازې
  نۀ يئ۔ ښايي په ټولګي کښې نورو ته هم ستونزې وي، خو
  کېدے شي د هغوى ستونزې په نورو څيزونو کښې وي۔ له
  ملګرو سره پرې خبرې کول د مسئلې بله زاويه درښودلے شي
  چې د حل لاره پرانيزي۔ البته، يوازې د هغوى حل مه نقلوئ:
  اشاره ترې وغواړئ، يا ووايئ چې چېرته ايسار يئ او د بل
  ګام پوښتنه وکړئ۔ او کله چې تاسو پرې پوه شئ، بيا د
  هغوى مرسته وکړئ۔ د بل چا مرسته او د مطلب تشريح کول
  ستاسو خپله پوهه هم ښه کوي۔
\item \emph{مرسته وغواړئ۔} ډېرې وسيلې درته شته؛ استاد
  او د تدريس مرستيال ستاسو دپاره حاضر دي او ستاسو
  برياليتوب \emph{غواړي}۔ هغوى بايد د مسئلې په حلولو او
  د بهير د هغه ځاے په پېژندلو کښې مرسته وکړے شي چې
  تاسو پکښې ستونزه لرئ۔
\item \emph{دمه وکړئ۔} که ايسار يئ، \emph{کېدے شي}
  وجه ئې دا وي چې له ډېرې مودې راهيسې هماغې مسئلې ته
  ګورئ۔ لنډه دمه وکړئ، يوه پياله چاے وڅښئ، يا يو څه
  وخت په بله مسئله کار وکړئ؛ بيا په تازه ذهن هغې ته
  راوګرځئ۔ يوه شپه ئې پرېږدئ او تر خوب وروسته بيا پرې
  فکر وکړئ۔
\end{enumerate}
\olpseganchor{OLP-0611-B007}{end}

\olpseganchor{OLP-0611-B008}{start}
ګورئ چې دا لارې له تاسو غواړي ثبوت له مخکې نه وختي
پيل کړے وي؟ که د سپارلو له نېټې يوه ورځ وړاندې د شپې
په دوو بجو مو ثبوت پيل کړے وي، ښايي دا لارې دومره
مرستندويې نۀ وي۔
\olpseganchor{OLP-0611-B008}{end}

\olpseganchor{OLP-0611-B009}{start}
ښايي دا خبرې ډېرې ناامېدوونکې ښکاره شي، خو د ثبوت
حلول داسې ننګونه ده چې په پاې کښې ګټه لري۔ ځينې
خلک همدا کار د مسلک په توګه کوي؛ نو ضرور به پکښې
د خوند يو څه وي۔ د نږدې هر کار په شان د مسئلو
حلول او د ثبوتونو کول تمرين غواړي۔ ښايي داسې
ټولګيوال ووينئ چې دا کار ورته اسان وي؛ غالباً
هغوى له مخکې ډېر تمرين کړے دے۔ هڅه وکړئ چې ډېر
په آسانۍ لاس ترې وانۀ خلئ۔
\olpseganchor{OLP-0611-B009}{end}

\olpseganchor{OLP-0611-B010}{start}
که په کومه ځانګړې مسئله وخت يا زغم درنه ختم شي،
پروا نۀ کوي۔ دا معنا نۀ لري چې تاسو ناپوه يئ يا
هېڅکله به پرې پوه نۀ شئ۔ له استاد يا بل زده
کوونکي معلومه کړئ چې څنګه حلېږي، او دا وپېژنئ
چې چېرته مو تېروتنه کړې يا ايسار شوي يئ، څو
راتلونکے ځل له ورته ستونزې سره د مخ کېدو پر
وخت ترې ځان وساتئ۔ بيا هڅه وکړئ چې حل ته له
کتلو پرته ئې پخپله وکړئ۔ او راتلونکے ځل وختي
پيل وکړئ او وختي مرسته وغواړئ۔
\olpseganchor{OLP-0611-B010}{end}

\olpseganchor{OLP-0612-B004}{start}
\olfileid{mth}{prf}{res}
\olpseganchor{OLP-0612-B004}{end}

\olpseganchor{OLP-0612-B005}{start}
\section{نورې سرچينې}
\olpseganchor{OLP-0612-B005}{end}

\olpseganchor{OLP-0612-B006}{start}
په رياضياتو کښې د ثبوتونو د کولو په اړه ډېر کتابونه
شته چې ګټور کېدے شي۔ په ځانګړي ډول \emph{ثبوتونه
څنګه ولولو او وکړو: د رياضيکي فکر د بهيرونو پېژندنه}
\mbox{\LR{\latinfont \citep{Solow2013}}} او \emph{څنګه ئې ثابت کړو: يوه منظمه
لار} \mbox{\LR{\latinfont \citep{Velleman2019}}} وګورئ۔
\href{http://www.people.vcu.edu/~rhammack/BookOfProof/BookOfProof.pdf}{\emph{د ثبوت کتاب}}
\mbox{\LR{\latinfont \citep{Hammack2013}}} او
\href{https://scholarworks.gvsu.edu/books/7/}{\emph{رياضيکي استدلال}}
\mbox{\LR{\latinfont \citep{Sandstrum2019}}} د ثبوت په اړه داسې کتابونه دي
چې په انټرنېټ کښې وړيا شته۔ د فلسفې لوستونکي ښايي
\emph{لا په دقت: هغه رياضيات چې د فلسفې د کولو
دپاره ئې اړتيا لرئ} \mbox{\LR{\latinfont \citep{Steinhart2018}}} د رياضيکي
استدلال ښه لومړنے لارښود وګڼي۔
\olpseganchor{OLP-0612-B006}{end}

\olpseganchor{OLP-0612-B007}{start}
په انټرنېټ کښې د ثبوتونو بېلابېل لنډ لارښودونه هم
شته؛ د بېلګې په توګه
\href{https://math.berkeley.edu/~hutching/teach/proofs.pdf}{«د رياضيکي استدلالونو پېژندنه»}
\mbox{\LR{\latinfont \citep{Hutchings2003}}} او
\href{https://eugeniacheng.com/wp-content/uploads/2017/02/cheng-proofguide.pdf}{«ثبوتونه څنګه وليکو»}
\mbox{\LR{\latinfont \citep{Cheng2004}}}۔
\olpseganchor{OLP-0612-B007}{end}

\olpseganchor{OLP-0612-B008}{start}
\subsection{هڅوونکې ويډيوګانې}
\olpseganchor{OLP-0612-B008}{end}

\olpseganchor{OLP-0612-B009}{start}
د کورنۍ دندې د کولو زړه مو نۀ کېږي؟ خپه يئ؟
ښايي دا ويډيوګانې مرسته وکړي!
\olpseganchor{OLP-0612-B009}{end}

\olpseganchor{OLP-0612-B010}{start}
\begin{itemize}
\item \url{https://www.youtube.com/watch?v=ZXsQAXx_ao0}
\item \url{https://www.youtube.com/watch?v=BQ4yd2W50No}
\item \url{https://www.youtube.com/watch?v=StTqXEQ2l-Y}
\end{itemize}
\olpseganchor{OLP-0612-B010}{end}

\olpseganchor{OLP-0613-B004}{start}
\olchapter{mth}{ind}{استقرا}
\olpseganchor{OLP-0613-B004}{end}

\olpseganchor{OLP-0614-B004}{start}
\olfileid{mth}{ind}{int}
\olpseganchor{OLP-0614-B004}{end}

\olpseganchor{OLP-0614-B005}{start}
\olsection{پېژندګلو}
\olpseganchor{OLP-0614-B005}{end}

\olpseganchor{OLP-0614-B006}{start}
استقرا د ثبوت يو مهم تخنيک دے چې په بېلابېلو بڼو کښې د منطق،
نظري کمپيوټرپوهنې او رياضياتو په تقريباً ټولو څانګو کښې کارېږي۔
د منطق د ډېرو نتيجو د ثابتولو دپاره ورته اړتيا وي۔
\olpseganchor{OLP-0614-B006}{end}

\olpseganchor{OLP-0614-B007}{start}
استقرا ډېر ځله د استنتاج په مقابله کښې يادېږي، او له ځانګړي څخه
عام ته د استدلال په توګه بيانېږي۔ د بېلګې په توګه، که ډېر شين
زمرد ووينو، او داسې څه و نۀ وينو چې زمرد ئې بولو خو شين نۀ وي،
نو ښايي دې نتيجې ته ورسېږو چې ټول زمرد شين دي۔ دا استقرايي
استدلال دے، ځکه چې له ډېرو ځانګړو حالتونو څخه، لکه دا زمرد شين
دے، هغه زمرد شين دے، او داسې نور، يوې عامې دعوې ته ځي چې ټول
زمرد شين دي۔ \emph{رياضي} استقرا هم داسې استدلال دے چې پايله ئې
يوه عامه دعوه وي، خو ډول ئې له دې «ساده استقرا» څخه ډېر بېل دے۔
\olpseganchor{OLP-0614-B007}{end}

\olpseganchor{OLP-0614-B008}{start}
په ډېر عمومي ډول، په رياضياتو کښې استقرايي ثبوت دې نتيجې ته
رسېږي چې د يو ټاکلي ډول ټول رياضيکي څيزونه يو ټاکلے خاصيت لري۔
په تر ټولو ساده حالت کښې هغه رياضيکي څيزونه چې استقرايي ثبوت
ورسره کار لري، طبيعي عددونه دي۔ په دې حالت کښې استقرايي ثبوت
دا ثابتوي چې ټول طبيعي عددونه يو خاصيت لري، او د دې دپاره ښيي چې
\begin{enumerate}
    \item \LR{$0$} هغه خاصيت لري، او
    \item هر کله چې يو عدد~\LR{$k$} هغه خاصيت ولري، نو~\LR{$k+1$} ئې هم لري۔
\end{enumerate}
پر طبيعي عددونو استقرا بيا ډېر ځله د داسې رياضيکي څيزونو په
هکله د عامو دعوو د ثابتولو دپاره هم کارېږي چې عددونه ورته
ټاکل کېداے شي۔ د بېلګې په توګه، هر متناهي سټ د غړو يو متناهي
شمېر~\LR{$n$} لري۔ که په استقرا وښيو چې هر عدد~\LR{$n$} دا خاصيت لري چې
«د~\LR{$n$} اندازې ټول متناهي سټونه \dots دي»، نو د ټولو متناهي
سټونو په هکله به مو يوه خبره ثابته کړې وي۔
\olpseganchor{OLP-0614-B008}{end}

\olpseganchor{OLP-0614-B009}{start}
استقرا هغو رياضيکي څيزونو ته هم تعميمېداے شي چې
\emph{په استقرايي ډول تعريف شوي} وي۔ د بېلګې په توګه، د يوې
رسمي ژبې عبارتونه، لکه د لومړي ترتيب د منطق عبارتونه، په
استقرايي ډول تعريفېږي۔ \emph{جوړښتي استقرا} د ټولو داسې
عبارتونو په هکله د نتيجو د ثابتولو يوه طريقه ده۔ په تېره بيا
جوړښتي استقرا په منطق کښې ډېره ګټوره ده او ډېر کارېږي۔
\olpseganchor{OLP-0614-B009}{end}

\olpseganchor{OLP-0718-B004}{start}
\olfileid{sfr}{ind}{int}
\olpseganchor{OLP-0718-B004}{end}

\olpseganchor{OLP-0718-B005}{start}
\olsection{د طبيعي عددونو او فورمولونو استقرا}
\olpseganchor{OLP-0718-B005}{end}

\olpseganchor{OLP-0718-B006}{start}
\begin{explain}
استقرا د رياضياتو د ثبوت يوه عامه طريقه ده چې په منطق کښې هم کارېږي۔ ښايي له رياضياتو سره ئې اشنا يئ، چې پر طبيعي عددونو استقرا کوئ۔ د رياضيکي او منطقي څيزونو ډېر نور ډولونه هم داسې دي چې استقرا پرې کېدے شي۔
\olpseganchor{OLP-0718-B006}{end}

\olpseganchor{OLP-0718-B007}{start}
په عمومي ډول په استقرا يوه کلي دعوه ثابتوو؛ يعنې ښيو چې د يو ټاکلي ډول هر څيز يو خاصيت لري۔ استقرا په تېره بيا هغه وخت ګټوره وي چې د کلي حکم د مستقيم داخلولو طريقه ډېره پېچلې وي۔ د جوړښتي استقرا دپاره اړوند څيزونه په ځانګړي ډول جوړېږي: که د سټ هر غړے يا بنسټيز وي يا په متناهي ګامونو له بنسټيزو غړو جوړ شوے وي، پر دغه جوړښت استقرا کولے شو۔ لاندې تعريفونه دا خبره څرګندوي۔
\end{explain}
\olpseganchor{OLP-0718-B007}{end}

\olpseganchor{OLP-0718-B008}{start}
\begin{defn}
سټ \LR{$S$} د تابعې \LR{$f$} لاندې هله \emph{تړلے} دے چې تابع ئې پر ټولو غړو تعريف وي او، هر کله چې \LR{$x \in
S$}، نو \LR{$f(x) \in S$}۔
\end{defn}
\olpseganchor{OLP-0718-B008}{end}

\olpseganchor{OLP-0718-B009}{start}
\begin{explain}
د بېلګې په توګه، طبيعي عددونه (\LR{$\mathbb{N}$}) د ځاے‌ناستې تابعې \LR{$f(x) = x+1$} لاندې تړلي دي۔ که \LR{$x$} طبيعي عدد وي، \LR{$x+1$} هم طبيعي عدد دے۔ خو د مخکيني عدد د تابعې \LR{$f(x) = x-1$} لاندې تړلي نۀ دي، ځکه \LR{$0$} طبيعي عدد دے خو \LR{$-1$} طبيعي عدد نۀ دے۔
\end{explain}
\olpseganchor{OLP-0718-B009}{end}

\olpseganchor{OLP-0718-B010}{start}
\begin{defn}
وايو چې خاصيت \LR{$P$} د تابعې \LR{$f(x)$} لاندې \emph{ساتل کېږي} که د \LR{$f$} د تعريف په ساحه کښې د هر څيز \LR{$a$} دپاره له \LR{$P(a)$} څخه \LR{$P(f(a))$} راځي؛ يعنې که \LR{$a$}، \LR{$P$} خاصيت ولري، \LR{$f(a)$} ئې هم لري۔
\end{defn}
\olpseganchor{OLP-0718-B010}{end}

\olpseganchor{OLP-0718-B011}{start}
\begin{explain}
د «جفت عدد کېدو» خاصيت د تابعې \LR{$f(x) = x+2$} لاندې ساتل کېږي، ځکه که \LR{$x$} جفت وي، \LR{$x+2$} هم جفت دے۔ خو د ځاے‌ناستې تابعې لاندې نۀ ساتل کېږي؛ که \LR{$x$} جفت وي، \LR{$x+1$} طاق دے۔
\olpseganchor{OLP-0718-B011}{end}

\olpseganchor{OLP-0718-B012}{start}
پر طبيعي عددونو استقرا ځکه کېږي چې له صفر نه د ځاے‌ناستې تابعې په متناهي ځلو کارولو هر طبيعي عدد ته رسېږو۔ په همدې ډول د جوړښتي استقرا وړ سټونو غړي هم له خپل بنسټ نه په متناهي ګامونو جوړېږي۔
\end{explain}
\olpseganchor{OLP-0718-B012}{end}

\olpseganchor{OLP-0718-B013}{start}
\begin{defn}[پر \LR{$\Nat$} استقرا]
که صفر \LR{$P$} خاصيت ولري او \LR{$P$} د ځاے‌ناستې تابعې لاندې وساتل شي، نو هر طبيعي عدد \LR{$P$} خاصيت لري۔
\end{defn}
\olpseganchor{OLP-0718-B013}{end}

\olpseganchor{OLP-0718-B014}{start}
\begin{ex} د لومړنيو \LR{$n$} مثبتو طبيعي عددونو جمع \LR{$\frac{n(n+1)}{2}$} ده۔ \\
\olpseganchor{OLP-0718-B014}{end}

\olpseganchor{OLP-0718-B015}{start}
\textbf{بنسټيز حالت۔} د صفر عددونو جمع \LR{$0=
\frac{0\cdot 1}{2}$} ده۔\\
\olpseganchor{OLP-0718-B015}{end}

\olpseganchor{OLP-0718-B016}{start}
\textbf{استقرايي ګام۔} فرض کړئ چې خاصيت د \LR{$k$} دپاره سم دے۔ ښيو چې د \LR{$k+1$} دپاره هم سم دے۔ يعنې فرض کړئ \LR{$P(k)$} صدق کوي:
\[ 0 + 1 + \ldots + k = \frac{k(k+1)}{2} \]
ښيو چې \LR{$P(k+1)$} صدق کوي، يعنې:
\[ 0 + 1 + \ldots + k + (k+1) = \frac{(k+1)(k+2)}{2} \]
د دې دپاره هغه فرض کاروو چې \LR{$P(k)$}۔
\olpseganchor{OLP-0718-B016}{end}

\olpseganchor{OLP-0718-B017}{start}
\begin{align*}
0 + 1 + \ldots + k &= \frac{k(k+1)}{2} \tag{\RL{\normalfont فرض}} \\
0 + 1 + \ldots + k + (k+1) &= \frac{k(k+1)}{2} + (k+1)
\tag{\RL{\normalfont دواړو خواوو ته \LR{$k+1$} زيات کړئ}} \\
&= \frac{k(k+1) + 2(k+1)}{2} \\
&= \frac{k^2 + k + 2k +2}{2} \\
&=\frac{(k+1)(k+2)}{2}
\end{align*}
نو استقرايي ګام پوره شو او دعوه ثابته شوه۔
\end{ex}
\olpseganchor{OLP-0718-B017}{end}

\olpseganchor{OLP-0718-B018}{start}
\begin{explain}
د فارمولونه بنسټيز غړي اټومي فارمولونه دي۔ پېچلي فارمولونه د لږ پېچلو فارمولونه نه د منطقي عملګرو په کارولو جوړېږي۔ هر عملګر له داسې تابعې سره سمون لري چې له اړينو فارمولونه څخه يو نوے فارمول جوړوي او هماغه عملګر پرې کاروي۔ د بېلګې په توګه، هغه تابع چې فارمولونه \LR{$!A$} او \LR{$!B$} د او په تړنه تړي، \LR{$\mathcal E _\land (!A,
!B) = !A \land !B$} ليکلے شو۔ دې تابعو ته \emph{فورمول جوړوونکې تابعې} وايو۔ د فارمولونه سټ د دغو تابعو لاندې تړلے دے؛ هر کله چې \LR{$!A$} او \LR{$!B$} فارمولونه وي، \LR{$\lnot !A$}، \LR{$!A \land
!B$} او داسې نور هم فورمولونه دي۔
\end{explain}
\olpseganchor{OLP-0718-B018}{end}

\olpseganchor{OLP-0718-B019}{start}
\begin{defn}[پر فارمولونه استقرا]
که هر اټومي فارمول، \LR{$P$} خاصيت ولري او \LR{$P$} د فارمول جوړوونکو تابعو لاندې وساتل شي، نو هر فارمول، \LR{$P$} خاصيت لري۔
\end{defn}
\olpseganchor{OLP-0718-B019}{end}

\olpseganchor{OLP-0718-B020}{start}
\begin{explain}
دا تعريف پر فارمولونه د استقرايي ثبوت يوه طريقه راښيي۔ که ښودل غواړو چې هر فارمول، \LR{$P$} خاصيت لري، لاندې ګامونه اخلو۔\\
\olpseganchor{OLP-0718-B020}{end}

\olpseganchor{OLP-0718-B021}{start}
\textbf{بنسټيز حالت۔} فرض کړئ \LR{$!A$} يو اټومي فارمول دے۔ [\ldots]
نو \LR{$!A$}، \LR{$P$} خاصيت لري۔
\olpseganchor{OLP-0718-B021}{end}

\olpseganchor{OLP-0718-B022}{start}
\textbf{استقرايي ګام۔} فرض کړئ \LR{$!A$} او \LR{$!B$} داسې فارمولونه دي چې دواړه \LR{$P$} خاصيت لري۔
\olpseganchor{OLP-0718-B022}{end}

\olpseganchor{OLP-0718-B023}{start}
د \LR{$\lnot$} حالت۔ [\ldots] نو \LR{$\lnot !A$}، \LR{$P$} خاصيت لري۔
\olpseganchor{OLP-0718-B023}{end}

\olpseganchor{OLP-0718-B024}{start}
د \LR{$\land$} حالت۔ [\ldots] نو \LR{$!A \land !B$}، \LR{$P$} خاصيت لري۔
\olpseganchor{OLP-0718-B024}{end}

\olpseganchor{OLP-0718-B025}{start}
د \LR{$\lor$} حالت۔ [\ldots] نو \LR{$!A \lor !B$}، \LR{$P$} خاصيت لري۔
\olpseganchor{OLP-0718-B025}{end}

\olpseganchor{OLP-0718-B026}{start}
د \LR{$\lif$} حالت۔ [\ldots] نو \LR{$!A \lif !B$}، \LR{$P$} خاصيت لري۔
\olpseganchor{OLP-0718-B026}{end}

\olpseganchor{OLP-0718-B027}{start}
د \LR{$\liff$} حالت۔ [\ldots] نو \LR{$!A \liff !B$}، \LR{$P$} خاصيت لري۔
\olpseganchor{OLP-0718-B027}{end}

\olpseganchor{OLP-0718-B028}{start}
د \LR{$\lforall[]$} حالت۔ [\ldots] نو \LR{$\lforall[x] !A$}، \LR{$P$} خاصيت لري۔
\olpseganchor{OLP-0718-B028}{end}

\olpseganchor{OLP-0718-B029}{start}
د \LR{$\lexists[]$} حالت۔ [\ldots] نو \LR{$\lexists[x] !A$}، \LR{$P$} خاصيت لري۔ \\
\olpseganchor{OLP-0718-B029}{end}

\olpseganchor{OLP-0718-B030}{start}
نو هر فارمول، \LR{$P$} خاصيت لري۔
\end{explain}
\olpseganchor{OLP-0718-B030}{end}

\olpseganchor{OLP-0615-B004}{start}
\olfileid{mth}{ind}{inN}
\olpseganchor{OLP-0615-B004}{end}

\olpseganchor{OLP-0615-B005}{start}
\olsection{پر~\texorpdfstring{\LR{$\Nat$}}{ℕ} استقرا}
\olpseganchor{OLP-0615-B005}{end}

\olpseganchor{OLP-0615-B006}{start}
په تر ټولو ساده بڼه کښې استقرا يو داسې تخنيک دے چې د ټولو
طبيعي عددونو دپاره د نتيجو د ثابتولو دپاره کارېږي۔ دا پر دې
حقيقت تکيه کوي چې له \LR{$0$} څخه په پيل کولو او بيا بيا د~\LR{$1$} په
زياتولو بالاخره هر طبيعي عدد ته رسېږو۔ نو د دې د ثابتولو
دپاره چې يوه خبره د هر عدد دپاره رښتيا ده، کولاے شو (1)~ثابته
کړو چې د \LR{$0$} دپاره رښتيا ده، او (2)~وښيو چې هر کله د يو
عدد~\LR{$n$} دپاره رښتيا وي، نو د راتلونکي عدد~\LR{$n+1$} دپاره هم
رښتيا ده۔ که «عدد~\LR{$n$} خاصيت \LR{$P$} لري» په \LR{$P(n)$} لنډ کړو، او
«عدد~\LR{$k$} خاصيت \LR{$P$} لري» په \LR{$P(k)$}، او داسې نور، نو د \LR{$P(n)$}
استقرايي ثبوت، د ټولو \LR{$n \in \Nat$} دپاره، دا برخې لري:
\begin{enumerate}
\item د \LR{$P(0)$} ثبوت، او
\item دا ثبوت چې د هر~\LR{$k$} دپاره، که \LR{$P(k)$} وي نو \LR{$P(k+1)$} هم وي۔
\end{enumerate}
د دې د پوره روښانولو دپاره فرض کړئ چې (1) او~(2) دواړه لرو۔
بيا (1) راته وايي چې \LR{$P(0)$} رښتيا دے۔ که~(2) هم ولرو، نو
په ځانګړي ډول پوهېږو چې که \LR{$P(0)$} وي نو \LR{$P(0+1)$}، يعنې
\LR{$P(1)$} هم وي۔ دا له دې عام بيان څخه چې «د هر~\LR{$k$} دپاره، که
\LR{$P(k)$} وي نو \LR{$P(k+1)$} هم وي»، د~\LR{$0$} په اېښودلو، د~\LR{$k$} پر ځاے،
راوځي۔ نو د مودوس پوننس په قاعده~\LR{$P(1)$} لرو۔ له (2) څخه
بيا، اوس \LR{$1$} په اخيستلو، د~\LR{$n$} پر ځاے،، لرو چې که \LR{$P(1)$} وي
نو~\LR{$P(2)$} هم وي۔ ځکه چې همدا اوس مو~\LR{$P(1)$} ثابت کړ، د مقدم
د اثبات په قاعده~\LR{$P(2)$} لرو۔ او همداسې مخ په وړاندې ځو۔
د هر عدد~\LR{$n$} دپاره د دې کار له \LR{$n$}~ځله کولو وروسته بالاخره
تر~\LR{$P(n)$} رسېږو۔ نو (1) او~(2) په ګډه~\LR{$P(n)$} ثابتوي، د هر \LR{$n
\in \Nat$} دپاره۔
\olpseganchor{OLP-0615-B006}{end}

\olpseganchor{OLP-0615-B007}{start}
يوه بېلګه ګورو۔ فرض کړئ غواړو معلومه کړو چې د نرد په \LR{$n$}
دانو څو بېلابېلې مجموعې غورځولاے شو۔ که څه هم ښايي بې‌معنا
ښکاره شي، له \LR{$0$} دانو پيل کوو۔ که هېڅ دانه نۀ وي، نو يوازې
يوه ممکنه مجموعه «غورځولاے» شئ: هېڅ ټکے نشته، چې مجموعه ئې~\LR{$0$}
ده۔ نو د غورځولو د بېلابېلو ممکنو مجموعو شمېر~\LR{$1$} دے۔ که
يوازې يوه دانه ولرئ، يعنې \LR{$n=1$}، نو له \LR{$1$} څخه تر~\LR{$6$} پورې
شپږ ممکنه ارزښتونه دي۔ په دوو دانو له \LR{$2$} څخه تر \LR{$12$} پورې
هره مجموعه غورځولاے شو، چې \LR{$11$} امکانات دي۔ په دريو دانو
له \LR{$3$} څخه تر \LR{$18$} پورې هر عدد غورځولاے شو، يعنې \LR{$16$}
بېلابېل امکانات۔ \LR{$1$}، \LR{$6$}، \LR{$11$}، \LR{$16$}: داسې ښکاري چې يو
نظم لري؛ ښايي ځواب \LR{$5n+1$} وي؟ البته \LR{$5n+1$} تر ټولو زيات
ممکن شمېر دے، ځکه چې يوازې \LR{$5n+1$} عددونه شته، د \LR{$n$}، يعنې په \LR{$n$} دانو د غورځولو
تر ټولو ټيټ ارزښت چې ټولې \LR{$1$} وي، او د \LR{$6n$}، يعنې تر ټولو
لوړ ارزښت چې ټولې \LR{$6$} وي، تر منځ۔
\olpseganchor{OLP-0615-B007}{end}

\olpseganchor{OLP-0615-B008}{start}
\begin{thm}
په \LR{$n$} نرد دانو ټول \LR{$5n+1$} ممکنه ارزښتونه غورځولاے شو،
د \LR{$n$} او \LR{$6n$} تر منځ۔
\end{thm}
\olpseganchor{OLP-0615-B008}{end}

\olpseganchor{OLP-0615-B009}{start}
\begin{proof}
\LR{$P(n)$} دې دا دعوه وي چې «د \LR{$n$} او~\LR{$6n$} تر منځ هر عدد
غورځول ممکن دي، په \LR{$n$}~نرد دانو»۔ د استقرا د کارولو دپاره
دا ثابتوو:
\begin{enumerate}
\item \emph{د استقرا بنسټ} \LR{$P(1)$}، يعنې په يوازې يوه دانه
  د \LR{$1$} او \LR{$6$} تر منځ هر عدد غورځولاے شئ۔
\item \emph{د استقرا ګام}، چې د ټولو \LR{$k$} دپاره، که \LR{$P(k)$}
  وي نو~\LR{$P(k+1)$} هم وي۔
\end{enumerate}
\olpseganchor{OLP-0615-B009}{end}

\olpseganchor{OLP-0615-B010}{start}
(1) د نرد د يوې \LR{$6$}-مخيزې دانې په کتلو ثابتېږي۔ دا ټول
6 مخونه لري، او د \LR{$1$} او~\LR{$6$} تر منځ هر عدد ئې پر يو مخ
ښکاري۔ نو په يوه دانه د \LR{$1$} او~\LR{$6$} تر منځ هر عدد غورځول
ممکن دي۔
\olpseganchor{OLP-0615-B010}{end}

\olpseganchor{OLP-0615-B011}{start}
د (2) د ثابتولو دپاره د شرط مقدم، يعنې \LR{$P(k)$}، فرض کوو۔
دې فرض ته \emph{استقرايي فرض} وايي۔ له دې څخه د~\LR{$P(k+1)$}
د ثابتولو دپاره کار اخلو۔ سخته برخه دا ده چې د \LR{$k+1$} دانو
د غورځولو د ممکنو ارزښتونو په هکله د فکر کولو داسې لاره
ومومو چې د \LR{$k$}~دانو د غورځولو ممکن ارزښتونه او د اضافي
\LR{$k+1$}-مې دانې د غورځولو ارزښتونه سره يوځاے کړي۔ خو که
استقرايي فرض کارول غواړو، نو همدا کار مو پکار دے۔
\olpseganchor{OLP-0615-B011}{end}

\olpseganchor{OLP-0615-B012}{start}
استقرايي فرض وايي چې د \LR{$k$} او~\LR{$6k$} تر منځ هر عدد ترلاسه
کولاے شو، په \LR{$k$}~دانو۔ که~\LR{$1$} وغورځوو، په خپلې \LR{$(k+1)$}-مې
دانې، نو مجموعې ته \LR{$1$} ورزياتوي۔ نو د \LR{$k+1$} او \LR{$6k+1$}
تر منځ هر ارزښت غورځولاے شو، د \LR{$k$}~دانو په غورځولو او
بيا د~\LR{$1$} په غورځولو، په \LR{$(k+1)$}-مه دانه۔ څه پاتې شول؟ له \LR{$6k+2$}
څخه تر \LR{$6k+6$} پورې ارزښتونه۔ دا په \LR{$k$} دانو د \LR{$6$} په
غورځولو او بيا د \LR{$2$} او \LR{$6$} تر منځ د يو عدد په غورځولو،
په خپلې \LR{$(k+1)$}-مې دانې، ترلاسه کولاے شو۔ په ګډه دا
معنا لري چې په \LR{$k+1$} دانو د \LR{$k+1$} او \LR{$6(k+1)$} تر منځ
هر عدد غورځولاے شو۔ يعنې~\LR{$P(k+1)$} مو ثابت کړ، د~\LR{$P(k)$}
فرض په کارولو، چې استقرايي فرض دے۔
\end{proof}
\olpseganchor{OLP-0615-B012}{end}

\olpseganchor{OLP-0615-B013}{start}
ډېر ځله استقرا هغه وخت کاروو چې د داسې څيزونو، لکه
عددونو، سټونو او نورو، د يوې لړۍ په هکله يوه خبره ثابتول
غواړو چې پخپله «په استقرايي ډول» تعريف شوې وي؛ يعنې
\LR{$(n+1)$}-م څيز ئې د \LR{$n$}-م له مخې تعريف شوے وي۔ د بېلګې
په توګه، د طبيعي عددونو مجموعه~\LR{$s_n$}، تر~\LR{$n$} پورې، داسې
تعريفولاے شو:
\begin{align*}
  s_0 & = 0\\
  s_{n+1} & = s_n + (n+1)
\end{align*}
له دې تعريفه ترلاسه کېږي:
\begin{align*}
  s_0 & = 0,\\
  s_1 & = s_0 + 1 && = 1,\\
  s_2 & = s_1 + 2 && = 1 + 2 = 3\\
  s_3 & = s_2 + 3 && = 1 + 2 + 3 = 6, \text{\RL{ او داسې نور۔}}
\end{align*}
اوس په استقرا ثابتولاے شو چې \LR{$s_n = n(n+1)/2$}۔
\olpseganchor{OLP-0615-B013}{end}

\olpseganchor{OLP-0615-B014}{start}
\begin{prop}
  \LR{$s_n = n(n+1)/2$}.
\end{prop}
\olpseganchor{OLP-0615-B014}{end}

\olpseganchor{OLP-0615-B015}{start}
\begin{proof}
بايد (1) دا ثابت کړو چې \LR{$s_0 = 0\cdot(0 + 1)/2$}، او (2)
که \LR{$s_k =
  k(k+1)/2$} وي نو \LR{$s_{k+1} = (k+1)(k+2)/2$} هم وي۔ (1)
ښکاره دے۔ د (2) د ثابتولو دپاره استقرايي فرض اخلو چې
\LR{$s_k = k(k+1)/2$}۔ د دې په کارولو بايد وښيو چې
\LR{$s_{k+1} = (k+1)(k+2)/2$}۔
\olpseganchor{OLP-0615-B015}{end}

\olpseganchor{OLP-0615-B016}{start}
\LR{$s_{k+1}$} څه دے؟ د تعريف له مخې، \LR{$s_{k+1} = s_k + (k+1)$}۔
د استقرايي فرض له مخې، \LR{$s_k = k(k+1)/2$}۔ دا په تېره
معادله کښې ځايناستي کولاے شو، او بيا يوازې د کسرونو لږ
حساب ته اړتيا لرو:
  \begin{align*}
    s_{k+1} & = \frac{k(k+1)}{2} + (k+1) = {}\\
    & = \frac{k(k+1)}{2} + \frac{2(k+1)}{2} = {}\\
    & = \frac{k(k+1) + 2(k+1)}{2} = {}\\
    & = \frac{(k+2)(k+1)}{2}.
  \end{align*}
\end{proof}
\olpseganchor{OLP-0615-B016}{end}

\olpseganchor{OLP-0615-B017}{start}
دلته مهمه زده کړه دا ده چې که د يوې په استقرايي ډول
تعريف شوې لړۍ \LR{$a_n$} په هکله يوه خبره ثابتوئ، نو استقرا
ښکاره مناسبه لاره ده۔ او که لړۍ داسې تعريف شوې هم نۀ
وي، لکه د نرد د غورځولو د امکاناتو په بېلګه کښې، بيا
هم استقرا کارولاے شئ، که په څه ډول د~\LR{$k+1$} حالت د~\LR{$k$}
له حالت سره تړلاے شئ۔
\olpseganchor{OLP-0615-B017}{end}

\olpseganchor{OLP-0616-B004}{start}
\olfileid{mth}{ind}{str}
\olpseganchor{OLP-0616-B004}{end}

\olpseganchor{OLP-0616-B005}{start}
\olsection{قوي استقرا}
\olpseganchor{OLP-0616-B005}{end}

\olpseganchor{OLP-0616-B006}{start}
د استقرا په هغه اصل کښې چې پورته مو بيان کړ، \LR{$P(0)$}
ثابتوو، او دا هم چې که \LR{$P(k)$} وي نو \LR{$P(k+1)$} هم وي۔ په
دويمه برخه کښې فرض کوو چې \LR{$P(k)$} رښتيا دے، او د دې
فرض په کارولو \LR{$P(k+1)$} ثابتوو۔ البته، په معادل ډول
\LR{$P(k-1)$} هم فرضولاے شو او د \LR{$P(k)$} د ثابتولو دپاره
ترې کار اخيستلاے شو۔ مهمه خبره دا ده چې له هر عدد
څخه ئې راتلونکي ته استدلال کولاے شو؛ يعنې د هر عدد
دپاره اړونده دعوه په دې فرض ثابتولاے شو چې د هغه د
مخکيني عدد دپاره سمه ده۔
\olpseganchor{OLP-0616-B006}{end}

\olpseganchor{OLP-0616-B007}{start}
د استقرا د اصل يوه بله بڼه شته چې په هغې کښې يوازې
دا نۀ فرضوو چې دعوه د مخکيني \LR{$k-1$} دپاره سمه
ده، چې د \LR{$k$} مخکينے عدد دے، بلکې د~\LR{$k$} څخه د ټولو کوچنيو عددونو دپاره ئې
فرضوو، او د دې فرض په کارولو دعوه د~\LR{$k$} دپاره
ثابتوو۔ دا هم دعوه \LR{$P(n)$} راکوي، د ټولو~\LR{$n \in
\Nat$} دپاره۔ ځکه چې يو ځل \LR{$P(0)$}
ثابت کړو، نو په دې سره مو ثابته کړې وي چې \LR{$P$} له~\LR{$1$}
څخه د ټولو کوچنيو عددونو دپاره سم دے۔ او که پوهېږو
چې له \LR{$P(l)$} څخه، د ټولو \LR{$l<k$} دپاره، \LR{$P(k)$} راوځي، نو
په ځانګړي ډول د \LR{$k=1$} دپاره هم په دې پوهېږو۔ نو
\LR{$P(1)$} نتيجه اخيستلاے شو۔ په دې سره مو \LR{$P(0)$} او
\LR{$P(1)$} ثابت کړي، يعنې \LR{$P(l)$}، د ټولو \LR{$l<2$} دپاره۔ او
ځکه چې دا شرط هم لرو چې که \LR{$P(l)$} وي، د ټولو \LR{$l<2$} دپاره،
نو \LR{$P(2)$} هم وي، نو~\LR{$P(2)$} نتيجه اخيستلاے شو، او
همداسې مخ په وړاندې ځو۔
\olpseganchor{OLP-0616-B007}{end}

\olpseganchor{OLP-0616-B008}{start}
په حقيقت کښې، که دا عام شرط ثابتولاے شو چې «د ټولو
\LR{$k$} دپاره، که \LR{$P(l)$} وي، د ټولو \LR{$l<k$} دپاره، نو \LR{$P(k)$}
هم وي»، بيا د \LR{$P(0)$} جلا ثابتولو ته اړتيا نشته، ځکه
چې له همدغه شرطه راوځي۔ په ياد ولرئ چې داسې عامه
دعوه لکه «د ټولو \LR{$l<k$} دپاره \LR{$P(l)$}» هغه وخت رښتيا
ده چې هېڅ \LR{$l<k$} نۀ وي۔ دا د تشې کميت‌ښودنې يو حالت
دے: «ټول \LR{$A$}، \LR{$B$} دي» هغه وخت رښتيا ده چې هېڅ \LR{$A$}
نۀ وي؛ \LR{$\lforall[x][(!A(x) \lif !B(x))]$} هغه وخت
رښتيا دے چې هېڅ \LR{$x$}،~\LR{$!A(x)$} نۀ پوره کوي۔ په دې
حالت کښې ئې رسمي بڼه «\LR{$\lforall[l][(l < k \lif P(l))]$}»
ده، او دا هغه وخت رښتيا ده چې هېڅ \LR{$l < k$} نۀ وي۔
او که \LR{$k=0$} وي، همدا حالت دے: هېڅ \LR{$l<0$} نشته؛ نو
«د ټولو \LR{$l<0$} دپاره \LR{$P(0)$}» رښتيا دے، که \LR{$P$} هر
څه وي۔ له دې امله د دې شرط ثبوت چې «که \LR{$P(l)$} وي، د ټولو
\LR{$l<k$} دپاره، نو \LR{$P(k)$} هم وي»، په خپله~\LR{$P(0)$}
ثابتوي۔
\olpseganchor{OLP-0616-B008}{end}

\olpseganchor{OLP-0616-B009}{start}
دا بڼه هغه وخت ګټوره ده چې د~\LR{$k$} دپاره د دعوې
ثابتول يوازې د \LR{$k-1$} په دعوه تکيه نۀ شي کولاے،
بلکې ښايي دې فرض ته اړتيا ولري چې د يوه يا څو
\LR{$l<k$} دپاره رښتيا ده۔
\olpseganchor{OLP-0616-B009}{end}

\olpseganchor{OLP-0617-B004}{start}
\olfileid{mth}{ind}{idf}
\olpseganchor{OLP-0617-B004}{end}

\olpseganchor{OLP-0617-B005}{start}
\olsection{استقرايي تعريفونه}
\olpseganchor{OLP-0617-B005}{end}

\olpseganchor{OLP-0617-B006}{start}
په منطق کښې ډېر ځله د څيزونو ډولونه \emph{په استقرايي ډول}
تعريفوو؛ يعنې د دې دپاره چې څه د تعريفېدونکي ډول څيز ګڼل
کېږي، داسې قواعد ټاکو چې بيانوي څنګه د همدغه ډول له پخوانيو
څيزونو څخه نوي څيزونه ترلاسه کړو۔ د بېلګې په توګه، ډېر
ځله د نښو د لړيو ځانګړي ډولونه، لکه د يوې ژبې ترمونه او
فارمولونه، په استقرا تعريفوو۔ د يوې ساده بېلګې دپاره
داسې تارونه په نظر کښې ونيسئ چې له \LR{$\mathrm{a}$}، \LR{$\mathrm{b}$}،
\LR{$\mathrm{c}$}، \LR{$\mathrm{d}$} تورو، له~\LR{$\circ$} نښې او له \LR{$[$} او
\LR{$]$} قوسونو جوړ دي، لکه «\LR{$[[\mathrm{c} \circ \mathrm{d}][$}»،
«\LR{$[\mathrm{a}[]\circ]$}»، «\LR{$\mathrm{a}$}» يا
«\LR{$[[\mathrm{a} \circ \mathrm{b}]\circ
\mathrm{d}]$}»۔ ښايي احساس وکړئ چې په لومړيو دوو تارونو
کښې څه «ناسم» دي: په لومړي کښې قوسونه هېڅ «برابر» نۀ
دي، او ښايي احساس وکړئ چې «\LR{$\circ$}» بايد داسې عبارتونه
«سره وتړي» چې پخپله معنا ولري۔ درېيم او څلورم تار ښه
ښکاري: د هر «\LR{$[$}» دپاره يو تړونکے «\LR{$]$}» شته، که قوسونه
پکښې وي، او د هر \LR{$\circ$} دواړو خواوو ته «ښه» عبارتونه
موندلاے شو چې د يوې جوړې قوسونو دننه دي۔
\olpseganchor{OLP-0617-B006}{end}

\olpseganchor{OLP-0617-B007}{start}
غواړو په کره توګه وټاکو چې څه «ښه‌جوړ ترم» ګڼل کېږي۔
تر ټولو وړاندې، هر تورے په يوازې ځان ښه‌جوړ دے۔ هر څه
چې يوازې يو تورے نۀ وي، بايد د «\LR{$[t \circ s]$}» په بڼه
وي، چې \LR{$s$} او \LR{$t$} پخپله ښه‌جوړ وي۔ برعکس، که \LR{$t$} او
\LR{$s$} ښه‌جوړ وي، نو د دواړو تر منځ د يو \LR{$\circ$} په
اېښودلو او د يوې جوړې قوسونو دننه ئې په راوستلو يو
نوی ښه‌جوړ ترم جوړولاے شو۔ له دې عملونو څخه د ښه‌جوړو
ترمونو د سټ د \emph{تعريفولو} دپاره کار اخيستلاے شو۔
دا يو \emph{استقرايي تعريف} دے۔
\olpseganchor{OLP-0617-B007}{end}

\olpseganchor{OLP-0617-B008}{start}
\begin{defn}[ښه‌جوړ ترمونه]
د \emph{ښه‌جوړو ترمونو} سټ په استقرايي ډول داسې تعريفېږي:
  \begin{enumerate}
  \item هر \LR{$\mathrm{a}$}، \LR{$\mathrm{b}$}، \LR{$\mathrm{c}$}،
    \LR{$\mathrm{d}$} تورے يو ښه‌جوړ ترم دے۔
  \item که \LR{$s_1$} او \LR{$s_2$} ښه‌جوړ ترمونه وي، نو
    \LR{$[s_1 \circ s_2]$} هم ښه‌جوړ ترم دے۔
  \item نور هېڅ څه ښه‌جوړ ترم نۀ دي۔
  \end{enumerate}
\end{defn}
\olpseganchor{OLP-0617-B008}{end}

\olpseganchor{OLP-0617-B009}{start}
دا تعريف راته وايي چې يو څيز هله او يوازې هله ښه‌جوړ
ترم ګڼل کېږي چې د (1) او~(2) دوو شرطونو له مخې په يو
متناهي شمېر ګامونو کښې جوړېداے شي۔ په لومړي ګام کښې
ټول هغه ښه‌جوړ ترمونه جوړوو چې يوازې تورې دي، يعنې:
\[
\mathrm{a}, \mathrm{b}, \mathrm{c}, \mathrm{d}
\]
په دويم ګام کښې (2) پر جوړو شوو ترمونو کاروو۔ ترلاسه کوو:
\[
[\mathrm{a} \circ \mathrm{a}], [\mathrm{a} \circ \mathrm{b}],
[\mathrm{b} \circ \mathrm{a}], \dots, [\mathrm{d} \circ \mathrm{d}]
\]
دا د دوو تورو د ټولو يوځاے کېدو دپاره دي۔ په درېيم ګام
کښې بيا (2) کاروو، پر هر دوو ښه‌جوړو ترمونو چې تر اوسه
مو جوړ کړي دي۔ نوي ښه‌جوړ ترمونه ترلاسه کوو، لکه
\LR{$[\mathrm{a} \circ [\mathrm{a} \circ
    \mathrm{a}]]$}، چې \LR{$t$} پکښې \LR{$\mathrm{a}$} د لومړي ګام
څخه او \LR{$s$} پکښې \LR{$[\mathrm{a} \circ \mathrm{a}]$} د دويم
ګام څخه دے؛ او \LR{$[[\mathrm{b} \circ
    \mathrm{c}] \circ [\mathrm{d} \circ \mathrm{b}]]$}، چې
له دوو ترمونو \LR{$[\mathrm{b} \circ \mathrm{c}]$} او \LR{$[\mathrm{d}
  \circ \mathrm{b}]$} څخه د دويم ګام په بنسټ جوړ شوے دے۔
او همداسې نور۔ (3) شرط دا امکان ردوي چې کوم څيز چې
په دې ډول نۀ وي جوړ شوے، د ښه‌جوړو ترمونو سټ ته ننوځي۔
\olpseganchor{OLP-0617-B009}{end}

\olpseganchor{OLP-0617-B010}{start}
په ياد ولرئ چې تر اوسه مو دا نۀ دي ثابت کړي چې د نښو
هره لړۍ چې ښه‌جوړه «احساسېږي»، د دې تعريف له مخې هم
ښه‌جوړه ده۔ خو دا بايد ښکاره وي چې هر څه چې جوړولاے
شو، په رښتيا ښه‌جوړ «احساسېږي»: قوسونه برابر دي، او
\LR{$\circ$} داسې برخې سره تړي چې پخپله ښه‌جوړې دي۔
\olpseganchor{OLP-0617-B010}{end}

\olpseganchor{OLP-0617-B011}{start}
د استقرايي تعريفونو مهمه ځانګړنه دا ده چې که د ټولو
ښه‌جوړو ترمونو په هکله څه ثابتول غواړئ، تعريف راته
وايي چې کوم حالتونه بايد په نظر کښې ونيسو۔ د بېلګې
په توګه، که درته وويل شي چې \LR{$t$} يو ښه‌جوړ ترم دے، نو
استقرايي تعريف راته وايي چې \LR{$t$} څنګه کېداے شي: \LR{$t$}
يو تورے کېداے شي، يا \LR{$[s_1 \circ s_2]$} کېداے شي، د
ښه‌جوړو ترمونو د يوې جوړې \LR{$s_1$} او~\LR{$s_2$} دپاره۔ د
(3) شرط له امله همدا يوازيني امکانات دي۔
\olpseganchor{OLP-0617-B011}{end}

\olpseganchor{OLP-0617-B012}{start}
د يو په استقرايي ډول تعريف شوي سټ د ټولو غړو په هکله
د دعوو د ثابتولو پر وخت د استقرا قوي بڼه په ځانګړي
ډول مهمه کېږي۔ د بېلګې په توګه، فرض کړئ غواړو ثابت
کړو چې د اوږدوالي~\LR{$n$} په هر ښه‌جوړ ترم کښې د \LR{$[$}
شمېر~\LR{$< n/2$} دے۔ دا د ټولو~\LR{$n$} په هکله يوه دعوه
ګڼلاے شو: د هر \LR{$n$} دپاره، د \LR{$[$} شمېر په هر هغه
ښه‌جوړ ترم کښې چې اوږدوالے ئې~\LR{$n$} دے، \LR{$< n/2$} دے۔
\olpseganchor{OLP-0617-B012}{end}

\olpseganchor{OLP-0617-B013}{start}
\begin{prop}
د هر \LR{$n$} دپاره، د \LR{$[$} شمېر په هغه ښه‌جوړ ترم کښې چې
اوږدوالے ئې~\LR{$n$} دے، \LR{$< n/2$} دے۔
\end{prop}
\olpseganchor{OLP-0617-B013}{end}

\olpseganchor{OLP-0617-B014}{start}
\begin{proof}
د قوي استقرا په ذريعه د دې نتيجې د ثابتولو دپاره
بايد وښيو چې دا شرطي دعوه رښتيا ده:
\begin{quote}
که د هر \LR{$l < k$} دپاره، د اوږدوالي~\LR{$l$} په هر ښه‌جوړ ترم
کښې د نښو شمېر \LR{$< l/2$} وي، د \LR{$[$} نښو، نو د
اوږدوالي~\LR{$k$} په هر ښه‌جوړ ترم کښې ئې شمېر \LR{$< k/2$}
دے، د \LR{$[$} نښو۔
\end{quote}
د دې شرط د ښودلو دپاره فرض کړئ چې مقدم ئې رښتيا
دے؛ يعنې د هر \LR{$l<k$} دپاره، د اوږدوالي~\LR{$l$} په ښه‌جوړو
ترمونو کښې د نښو شمېر \LR{$< l/2$} دے، د \LR{$[$} نښو۔ دې فرض ته
استقرايي فرض وايو۔ غواړو وښيو چې همدا خبره د
اوږدوالي~\LR{$k$} ښه‌جوړو ترمونو دپاره هم رښتيا ده۔
\olpseganchor{OLP-0617-B014}{end}

\olpseganchor{OLP-0617-B015}{start}
نو فرض کړئ چې \LR{$t$} يو ښه‌جوړ ترم دے چې اوږدوالے
ئې~\LR{$k$} دے۔ ځکه چې ښه‌جوړ ترمونه په استقرايي ډول
تعريف شوي دي، دوه حالتونه لرو: (1)~\LR{$t$} يوازې يو
تورے دے، يا (2)~\LR{$t$}، \LR{$[s_1 \circ s_2]$} دے، د ځينو
ښه‌جوړو ترمونو \LR{$s_1$} او~\LR{$s_2$} دپاره۔
\begin{enumerate}
\item \LR{$t$} يو تورے دے۔ بيا \LR{$k = 1$}، او د \LR{$[$} شمېر په
\LR{$t$} کښې~\LR{$0$} دے۔ ځکه چې \LR{$0 < 1/2$}، دعوه سمه ده۔
\item \LR{$t$}، \LR{$[s_1 \circ s_2]$} دے، د ځينو ښه‌جوړو ترمونو
  \LR{$s_1$} او~\LR{$s_2$} دپاره۔ \LR{$l_1$} دې د~\LR{$s_1$} اوږدوالے
  او \LR{$l_2$} دې د~\LR{$s_2$} اوږدوالے وي۔ بيا اوږدوالے~\LR{$k$}،
  د \LR{$t$} دپاره، \LR{$l_1+l_2+3$} دے؛ د \LR{$s_1$} او \LR{$s_2$}
  اوږدوالے او درې نښې \LR{$[$}، \LR{$\circ$}، \LR{$]$}۔ ځکه چې
  \LR{$l_1+l_2+3$} تل له \LR{$l_1$} څخه زيات دے، \LR{$l_1 < k$}۔
  همداسې \LR{$l_2
  < k$}۔ دا معنا لري چې استقرايي فرض پر ترمونو
  \LR{$s_1$} او~\LR{$s_2$} کارېږي: شمېر~\LR{$m_1$}، د \LR{$[$} دپاره
  په \LR{$s_1$} کښې، \LR{$< l_1/2$} دے، او شمېر~\LR{$m_2$}، د
  \LR{$[$} دپاره په~\LR{$s_2$} کښې، \LR{$< l_2/2$} دے۔
\olpseganchor{OLP-0617-B015}{end}

\olpseganchor{OLP-0617-B016}{start}
د \LR{$[$} شمېر په \LR{$t$} کښې د \LR{$[$} د شمېر په~\LR{$s_1$}
کښې، د \LR{$[$} د شمېر په~\LR{$s_2$} کښې، او د~\LR{$1$} مجموعه
ده؛ يعنې \LR{$m_1 + m_2 + 1$} دے۔ ځکه چې \LR{$m_1
  < l_1/2$} او \LR{$m_2 < l_2/2$}، نو لرو:
  \[
  m_1 + m_2 + 1 < \frac{l_1}{2} + \frac{l_2}{2} + 1 = \frac{l_1+l_2+2}{2} < \frac{l_1+l_2+3}{2} = k/2.
  \]
\end{enumerate}
په هر حالت کښې مو ښودلې چې د \LR{$[$} شمېر په \LR{$t$} کښې
\LR{$< k/2$} دے، د استقرايي فرض په بنسټ۔ په قوي استقرا
مطلوبه دعوه ثابتېږي۔
\end{proof}
\olpseganchor{OLP-0617-B016}{end}

\olpseganchor{OLP-0617-B017}{start}
\begin{prob}
د ډېرو ښه‌جوړو ترمونو سټ داسې تعريف کړئ:
  \begin{enumerate}
  \item هر \LR{$\mathrm{a}$}، \LR{$\mathrm{b}$}، \LR{$\mathrm{c}$}،
    \LR{$\mathrm{d}$} تورے يو ډېر ښه‌جوړ ترم دے۔
  \item که \LR{$s$} ډېر ښه‌جوړ ترم وي، نو \LR{$[s]$} هم دے۔
  \item که \LR{$s_1$} او \LR{$s_2$} ډېر ښه‌جوړ ترمونه وي، نو
    \LR{$[s_1 \circ s_2]$} هم دے۔
  \item نور هېڅ څه ډېر ښه‌جوړ ترم نۀ دي۔
  \end{enumerate}
وښيئ چې د \LR{$[$} شمېر په ډېر ښه‌جوړ ترم~\LR{$t$} کښې، چې
اوږدوالے ئې~\LR{$n$} دے، \LR{$\le n/2 +1$} دے۔
\end{prob}
\olpseganchor{OLP-0617-B017}{end}

\olpseganchor{OLP-0618-B004}{start}
\olfileid{mth}{ind}{sti}
\olpseganchor{OLP-0618-B004}{end}

\olpseganchor{OLP-0618-B005}{start}
\olsection{جوړښتي استقرا}
\olpseganchor{OLP-0618-B005}{end}

\olpseganchor{OLP-0618-B006}{start}
تر اوسه مو استقرا د ټولو طبيعي عددونو په هکله د
نتيجو د ثابتولو دپاره کارولې ده۔ خو يو ورته اصل
مستقيم د يو په استقرايي ډول تعريف شوي سټ د ټولو
غړي په هکله د نتيجو د ثابتولو دپاره
کارېداے شي۔ دې ته ډېر ځله \emph{جوړښتي} استقرا
وايي، ځکه چې د په استقرايي ډول تعريف شوو څيزونو
پر جوړښت تکيه کوي۔
\olpseganchor{OLP-0618-B006}{end}

\olpseganchor{OLP-0618-B007}{start}
په عمومي ډول، يو استقرايي تعريف په (الف)~د سټ د
«لومړنيو» غړي په لړليک او (ب)~د داسې عملونو
په لړليک ورکول کېږي چې د سټ له پخوانيو غړو
څخه ئې نوي غړي جوړوي۔ د بېلګې په توګه، د
ښه‌جوړو ترمونو په حالت کښې لومړني څيزونه تورې دي۔
يوازې يو عمل لرو؛ عملونه دا دي:
\begin{align*}
  o(s_1, s_2) = & [s_1 \circ s_2]
\end{align*}
آن طبيعي عددونه~\LR{$\Nat$} پخپله هم د يو استقرايي
تعريف په ذريعه ورکړل شوي ګڼلاے شئ: لومړنے څيز~\LR{$0$}
دے، او عمل د راتلونکي تابع~\LR{$x + 1$} ده۔
\olpseganchor{OLP-0618-B007}{end}

\olpseganchor{OLP-0618-B008}{start}
د يو په استقرايي ډول تعريف شوي سټ د ټولو غړو په
هکله د څه ثابتولو دپاره، يعنې دا چې د سټ هر
غړے خاصيت~\LR{$P$} لري، بايد:
\begin{enumerate}
\item ثابت کړو چې لومړني څيزونه~\LR{$P$} لري۔
\item ثابت کړو چې د هر عمل~\LR{$o$} دپاره، که ارګومېنټونه
  ئې~\LR{$P$} ولري، نو نتيجه ئې هم لري۔
\end{enumerate}
د بېلګې په توګه، د ټولو ښه‌جوړو ترمونو په هکله
د څه ثابتولو دپاره به وښيو چې د ټولو تورو دپاره
رښتيا دي، او د \LR{$[s_1 \circ s_2]$} دپاره هم رښتيا
دي، په دې شرط چې د \LR{$s_1$} او \LR{$s_2$} هر يوه دپاره
په جلا توګه رښتيا وي۔
\olpseganchor{OLP-0618-B008}{end}

\olpseganchor{OLP-0618-B009}{start}
\begin{prop}
د \LR{$[$} شمېر د \LR{$]$} له شمېر سره برابر دے، په هر
ښه‌جوړ ترم~\LR{$t$} کښې۔
\end{prop}
\olpseganchor{OLP-0618-B009}{end}

\olpseganchor{OLP-0618-B010}{start}
\begin{proof}
جوړښتي استقرا کاروو۔ ښه‌جوړ ترمونه په استقرايي
ډول تعريف شوي، چې تورې ئې لومړني څيزونه دي او
عمل \LR{$o$} ئې له پخوانيو څخه نوي ښه‌جوړ ترمونه جوړوي۔
\begin{enumerate}
\item دعوه د هر توري دپاره رښتيا ده، ځکه چې د \LR{$[$}
  شمېر په يوازې يوه توري کښې~\LR{$0$} دے او د \LR{$]$} شمېر
  پکښې هم~\LR{$0$} دے۔
\item فرض کړئ چې د \LR{$[$} شمېر په \LR{$s_1$} کښې د \LR{$]$}
  له شمېر سره برابر دے، او همدا خبره د \LR{$s_2$}
  دپاره هم رښتيا ده۔ د \LR{$[$} شمېر په \LR{$o(s_1, s_2)$}
  کښې، يعنې په \LR{$[s_1 \circ s_2]$} کښې، د \LR{$[$} د
  شمېرونو مجموعه ده په \LR{$s_1$} او \LR{$s_2$} کښې، له
  يو سره۔ د \LR{$]$} شمېر په \LR{$o(s_1, s_2)$} کښې د \LR{$]$}
  د شمېرونو مجموعه ده په \LR{$s_1$} او \LR{$s_2$} کښې، له
  يو سره۔ نو د \LR{$[$} شمېر په \LR{$o(s_1, s_2)$} کښې د
  \LR{$]$} له شمېر سره برابر دے په \LR{$o(s_1,s_2)$} کښې۔
\end{enumerate}
\end{proof}
\olpseganchor{OLP-0618-B010}{end}

\olpseganchor{OLP-0618-B011}{start}
\begin{prob}
په جوړښتي استقرا ثابت کړئ چې هېڅ ښه‌جوړ ترم
په~\LR{$]$} نۀ پيلېږي۔
\end{prob}
\olpseganchor{OLP-0618-B011}{end}

\olpseganchor{OLP-0618-B012}{start}
په جوړښتي استقرا يو بل ثبوت ورکوو: د نښو د
يو تار~\LR{$t$} خاصه پيلنۍ برخه هر هغه تار~\LR{$s$} دے
چې له کيڼې خوا په لوستلو د نښې په نښه له \LR{$t$}
سره يو شان وي، خو \LR{$t$} اوږد وي۔ نو، د بېلګې په
توګه، \LR{$[a \circ {}$} خاصه پيلنۍ برخه ده د
\LR{$[a \circ b]$}؛ خو نۀ \LR{$[b \circ {}$}، ځکه چې په
دويمه نښه کښې توپير لري، او نۀ \LR{$[a
\circ b]$}، ځکه چې دواړه يو شان اوږد دي۔
\olpseganchor{OLP-0618-B012}{end}

\olpseganchor{OLP-0618-B013}{start}
\begin{prop}\ollabel{prop:initial}
د يوه ښه‌جوړ ترم~\LR{$t$} هره ناتشه خاصه پيلنۍ برخه
د \LR{$[$} نښې د \LR{$]$} له نښو زياتې لري۔
\end{prop}
\olpseganchor{OLP-0618-B013}{end}

\olpseganchor{OLP-0618-B014}{start}
\begin{proof}
پر~\LR{$t$} په استقرا:
  \begin{enumerate}
  \item \LR{$t$} يوازې يو تورے دے: بيا \LR{$t$} هېڅ خاصه پيلنۍ برخه نۀ لري۔
  \item \LR{$t = [s_1 \circ s_2]$}، د ځينو ښه‌جوړو ترمونو
  \LR{$s_1$} او~\LR{$s_2$} دپاره۔ که \LR{$r$} د \LR{$t$} يوه خاصه
  پيلنۍ برخه وي، نو څو امکانات شته:
    \begin{enumerate}
    \item \LR{$r$} يوازې \LR{$[$} دے: بيا \LR{$r$} يو \LR{$[$} زيات لري
      له خپلو~\LR{$]$} څخه۔
    \item \LR{$r$}، \LR{$[r_1$} دے، چې \LR{$r_1$} خاصه پيلنۍ برخه
      ده د~\LR{$s_1$}۔ ځکه چې \LR{$s_1$} ښه‌جوړ ترم دے، د
      استقرايي فرض له مخې \LR{$r_1$} د \LR{$[$} نښې له
      \LR{$]$} څخه زياتې لري، او همدا خبره د \LR{$[r_1$}
      دپاره هم رښتيا ده۔
    \item \LR{$r$}، \LR{$[s_1$} يا \LR{$[s_1 \circ {}$} دے: د
      تېرې نتيجې له مخې د \LR{$[$} او \LR{$]$} شمېرونه
      په~\LR{$s_1$} کښې برابر دي؛ نو د \LR{$[$} شمېر
      په \LR{$[s_1$} يا \LR{$[s_1 \circ {}$} کښې يو زيات
      دے د~\LR{$]$} له شمېر څخه۔
    \item \LR{$r$}، \LR{$[s_1 \circ r_2$} دے، چې \LR{$r_2$} خاصه
      پيلنۍ برخه ده د~\LR{$s_2$}۔ د استقرايي فرض
      له مخې \LR{$r_2$} د \LR{$[$} نښې له \LR{$]$} څخه زياتې
      لري۔ د تېرې نتيجې له مخې د \LR{$[$} او
      \LR{$]$} شمېرونه په~\LR{$s_1$} کښې برابر دي۔ نو
      د \LR{$[$} شمېر په \LR{$[s_1 \circ r_2$} کښې زيات
      دے د~\LR{$]$} له شمېر څخه۔
    \item \LR{$r$}، \LR{$[s_1 \circ s_2$} دے: د تېرې نتيجې
      له مخې د \LR{$[$} او \LR{$]$} شمېرونه په \LR{$s_1$}
      کښې برابر دي، او همدا د~\LR{$s_2$} دپاره
      هم ده۔ نو يو \LR{$[$} زيات دے په \LR{$[s_1 \circ s_2$}
      کښې، د~\LR{$]$} له شمېر څخه۔
    \end{enumerate}
  \end{enumerate}
\end{proof}
\olpseganchor{OLP-0618-B014}{end}

\olpseganchor{OLP-0619-B004}{start}
\olfileid{mth}{ind}{rel}
\olpseganchor{OLP-0619-B004}{end}

\olpseganchor{OLP-0619-B005}{start}
\olsection{اړيکې او تابعې}
\olpseganchor{OLP-0619-B005}{end}

\olpseganchor{OLP-0619-B006}{start}
کله چې د څيزونو يو سټ، لکه طبيعي عددونه يا ښه‌جوړ
ترمونه، په استقرايي ډول تعريف کړو، نو \emph{پر دغو
څيزونو اړيکې} هم په استقرا تعريفولاے شو۔ د بېلګې
په توګه، دا نظر په پام کښې ونيسئ: يو ښه‌جوړ
ترم~\LR{$t_1$} د ښه‌جوړ ترم~\LR{$t_2$} فرعي ترم دے، که د
هغه د يوې برخې په توګه پکښې راځي۔ يوه نښه ورته
کاروو: \LR{$t_1 \sqsubseteq t_2$}۔ البته هر ښه‌جوړ
ترم د خپل ځان فرعي ترم دے: \LR{$t \sqsubseteq t$}۔
د دې اړيکې استقرايي تعريف داسې ورکولاے شو:
\olpseganchor{OLP-0619-B006}{end}

\olpseganchor{OLP-0619-B007}{start}
\begin{defn}
د ښه‌جوړ ترم~\LR{$t_1$} فرعي ترم کېدل د~\LR{$t_2$}، يعنې
\LR{$t_1
  \sqsubseteq t_2$} اړيکه، پر~\LR{$t_2$} په استقرا
داسې تعريفېږي:
  \begin{enumerate}
  \item که \LR{$t_2$} يو تورے وي، نو \LR{$t_1 \sqsubseteq t_2$}
    هله او يوازې هله چې \LR{$t_1 = t_2$}۔
  \item که \LR{$t_2$}، \LR{$[s_1 \circ s_2]$} وي، نو
    \LR{$t_1 \sqsubseteq t_2$} هله او يوازې هله چې
    \LR{$t_1 =
    t_2$}، \LR{$t_1 \sqsubseteq s_1$}، يا \LR{$t_1 \sqsubseteq s_2$}۔
  \end{enumerate}
\end{defn}
\olpseganchor{OLP-0619-B007}{end}

\olpseganchor{OLP-0619-B008}{start}
د بېلګې په توګه، دا تعريف به راته ووايي چې
\LR{$\mathrm{a}
\sqsubseteq [\mathrm{b} \circ \mathrm{a}]$}۔ ځکه چې
(2) وايي \LR{$\mathrm{a} \sqsubseteq [\mathrm{b} \circ \mathrm{a}]$}
هله او يوازې هله چې \LR{$\mathrm{a}
= [\mathrm{b} \circ \mathrm{a}]$}، يا \LR{$\mathrm{a} \sqsubseteq b$}،
يا \LR{$\mathrm{a} \sqsubseteq \mathrm{a}$}۔ لومړني دوه
ناسم دي: \LR{$\mathrm{a}$} ښکاره له \LR{$[\mathrm{b} \circ
  \mathrm{a}]$} سره يو شان نۀ دے، او د~(1) له مخې
\LR{$\mathrm{a} \sqsubseteq \mathrm{b}$} هله او يوازې
هله چې \LR{$\mathrm{a} = \mathrm{b}$}، چې دا هم ناسم
دے۔ خو بيا هم د~(1) له مخې \LR{$\mathrm{a} \sqsubseteq \mathrm{a}$}
هله او يوازې هله چې \LR{$\mathrm{a} = \mathrm{a}$}،
چې دا رښتيا دے۔
\olpseganchor{OLP-0619-B008}{end}

\olpseganchor{OLP-0619-B009}{start}
دا مهمه ده چې پوه شو د دې تعريف برياليتوب پر
يو داسې حقيقت تکيه کوي چې تر اوسه مو نۀ دے
ثابت کړے: هر ښه‌جوړ ترم~\LR{$t$} يا يوازې يو تورے
دے، يا \emph{په يکتا ډول ټاکل شوي} ښه‌جوړ
ترمونه \LR{$s_1$} او \LR{$s_2$} شته چې \LR{$t = [s_1 \circ s_2]$}۔
دلته «په يکتا ډول ټاکل شوي» دا معنا لري چې
که \LR{$t = [s_1 \circ s_2]$} وي، نو دا \emph{هم}
\LR{$= [r_1 \circ r_2]$} نۀ وي، له \LR{$s_1 \neq r_1$}
يا \LR{$s_2 \neq r_2$} سره۔ که داسې کېداے، نو (2)
شرط ښايي له خپل ځان سره ټکر وکړي: \LR{$t_2$} چې
د \LR{$[s_1 \circ s_2]$} په توګه ولولو، ښايي
\LR{$t_1
\sqsubseteq t_2$} ترلاسه کړو، خو که \LR{$t_2$} د
\LR{$[r_1 \circ r_2]$} په توګه ولولو، ښايي
\LR{$t_1 \sqsubseteq t_2$} نۀ وي، دا نتيجه ترلاسه کړو۔ مخکې
له دې چې ثابت کړو دا نۀ شي کېداے، يوه داسې
بېلګه ګورو چې دا پکښې \emph{کېداے} شي۔
\olpseganchor{OLP-0619-B009}{end}

\olpseganchor{OLP-0619-B010}{start}
\begin{defn}
\emph{بې‌قوسه ترمونه} په استقرا داسې تعريف کړئ:
  \begin{enumerate}
  \item هر تورے يو بې‌قوسه ترم دے۔
    \item که \LR{$s_1$} او \LR{$s_2$} بې‌قوسه ترمونه وي، نو
      \LR{$s_1 \circ s_2$} بې‌قوسه ترم دے۔
    \item نور هېڅ څه بې‌قوسه ترم نۀ دي۔
  \end{enumerate}
\end{defn}
\olpseganchor{OLP-0619-B010}{end}

\olpseganchor{OLP-0619-B011}{start}
بې‌قوسه ترمونه، د بېلګې په توګه، \LR{$\mathrm{a}$}،
\LR{$\mathrm{b} \circ
\mathrm{d}$}، \LR{$\mathrm{b} \circ \mathrm{a} \circ \mathrm{b}$}
دي۔ اوس که د بې‌قوسه ترمونو دپاره «فرعي ترم»
په هماغه ډول تعريف کړو لکه پورته مو چې وکړ،
دويم شرط به داسې وي:
\begin{center}
که \LR{$t_2 = s_1 \circ s_2$} وي، نو \LR{$t_1 \sqsubseteq t_2$}
هله او يوازې هله چې \LR{$t_1 = t_2$}، \LR{$t_1
  \sqsubseteq s_1$}، يا \LR{$t_1 \sqsubseteq s_2$}۔
\end{center}
\olpseganchor{OLP-0619-B011}{end}

\olpseganchor{OLP-0619-B012}{start}
اوس \LR{$\mathrm{b} \circ \mathrm{a} \circ \mathrm{b}$}
د \LR{$s_1
\circ s_2$} په بڼه دے، په دې ډول چې:
\begin{align*}
  s_1 & = \mathrm{b} \text{\RL{ او}} & s_2 & = \mathrm{a} \circ \mathrm{b}.
\intertext{\RL{دا د \LR{$r_1 \circ r_2$} په بڼه هم دے، په دې ډول چې:}}
  r_1 & = \mathrm{b} \circ \mathrm{a} \text{\RL{ او}} & r_2 &= \mathrm{b}.
\end{align*}
اوس آيا \LR{$\mathrm{a} \circ \mathrm{b}$} فرعي ترم دے
د \LR{$\mathrm{b} \circ
\mathrm{a} \circ \mathrm{b}$}؟ که لومړے لوست
تعقيب کړو، ځواب هو دے، او که دويم تعقيب
کړو، ځواب نۀ دے۔
\olpseganchor{OLP-0619-B012}{end}

\olpseganchor{OLP-0619-B013}{start}
دا خاصيت چې له نورو ښه‌جوړو ترمونو څخه د يوه
ښه‌جوړ ترم د جوړېدو طريقه يکتا ده،
\emph{يکتا لوستونتيا} نومېږي۔ ځکه چې د داسې
په استقرايي ډول تعريف شوو څيزونو دپاره د
اړيکو استقرايي تعريفونه مهم دي، نو بايد
ثابت کړو چې دا خاصيت سم دے۔
\olpseganchor{OLP-0619-B013}{end}

\olpseganchor{OLP-0619-B014}{start}
\begin{prop}
فرض کړئ چې \LR{$t$} يو ښه‌جوړ ترم دے۔ بيا يا \LR{$t$}
يوازې يو تورے دے، يا په يکتا ډول ټاکل شوي
ښه‌جوړ ترمونه \LR{$s_1$}، \LR{$s_2$} شته چې~\LR{$t =
  [s_1 \circ s_2]$}۔
\end{prop}
\olpseganchor{OLP-0619-B014}{end}

\olpseganchor{OLP-0619-B015}{start}
\begin{proof}
که \LR{$t$} يوازې يو تورے وي، شرط پوره دے۔ نو
فرض کړئ چې \LR{$t$} يوازې يو تورے نۀ دے۔ له
استقرايي تعريفه معلومولاے شو چې بيا \LR{$t$}
خامخا د \LR{$[s_1 \circ s_2]$} په بڼه وي، د
ځينو ښه‌جوړو ترمونو \LR{$s_1$} او~\LR{$s_2$} دپاره۔
دا ثابتول پاتې دي چې په يکتا ډول ټاکل
شوي دي؛ يعنې که \LR{$t = [r_1 \circ r_2]$}، نو
\LR{$s_1 = r_1$} او \LR{$s_2 = r_2$}۔
\olpseganchor{OLP-0619-B015}{end}

\olpseganchor{OLP-0619-B016}{start}
نو فرض کړئ چې \LR{$t = [s_1 \circ s_2]$} او
هم \LR{$t = [r_1 \circ r_2]$}، د ښه‌جوړو
ترمونو \LR{$s_1$}، \LR{$s_2$}، \LR{$r_1$}، \LR{$r_2$}
دپاره۔ بايد وښيو چې \LR{$s_1 = r_1$} او
\LR{$s_2 = r_2$}۔ لومړے، \LR{$s_1$} او \LR{$r_1$}
خامخا يو شان دي، ځکه که نۀ وي، نو
يو ئې د بل خاصه پيلنۍ برخه ده۔ خو
د \olref[sti]{prop:initial} له مخې دا
ناممکنه ده، که \LR{$s_1$} او~\LR{$r_1$} دواړه
ښه‌جوړ ترمونه وي۔ خو که \LR{$s_1 = r_1$}
وي، نو ښکاره \LR{$s_2 = r_2$} هم دے۔
\end{proof}
\olpseganchor{OLP-0619-B016}{end}

\olpseganchor{OLP-0619-B017}{start}
تابعې هم په استقرايي ډول تعريفولاے شو۔
د بېلګې په توګه، تابع~\LR{$f$} چې هر ښه‌جوړ
ترم پکښې د يو بل دننه د راغلو \LR{$[\dots]$}
قوسونو تر ټولو زيات ژوروالي ته لېږي،
داسې تعريفولاے شو:
\olpseganchor{OLP-0619-B017}{end}

\olpseganchor{OLP-0619-B018}{start}
\begin{defn}
\ollabel{defn:depth} د يوه ښه‌جوړ ترم
\emph{ژوروالے}، \LR{$f(t)$}، په استقرايي
ډول داسې تعريفېږي:
  \[
  f(t) = \begin{cases}
    0 & \text{\RL{ که \LR{$t$} يو تورے وي}}\\
    \max(f(s_1), f(s_2)) + 1 & \text{\RL{ که \LR{$t = [s_1 \circ s_2]$} وي۔}}
  \end{cases}
  \]
\end{defn}
\olpseganchor{OLP-0619-B018}{end}

\olpseganchor{OLP-0619-B019}{start}
د بېلګې په توګه:
\begin{align*}
  f([\mathrm{a} \circ \mathrm{b}]) & =
    \max(f(\mathrm{a}),f(\mathrm{b})) + 1 = \\
  &= \max(0, 0) + 1 = 1, \text{\RL{ او}}\\
  f([[\mathrm{a} \circ \mathrm{b}] \circ \mathrm{c}]) & =
    \max(f([\mathrm{a} \circ \mathrm{b}]), f(\mathrm{c})) + 1 = \\
  & = \max(1,0) + 1 = 2.
\end{align*}
\olpseganchor{OLP-0619-B019}{end}

\olpseganchor{OLP-0619-B020}{start}
البته دلته فرض کوو چې \LR{$s_1$} او \LR{$s_2$}
ښه‌جوړ ترمونه دي، او له دې حقيقته
کار اخلو چې هر ښه‌جوړ ترم يا يو
تورے دے يا د \LR{$[s_1 \circ s_2]$} په
بڼه دے۔ بيا دا مهمه ده چې په دغه
بڼه يوازې په يوه طريقه کېداے شي۔
د دې د علت د ليدلو دپاره هغه
بې‌قوسه ترمونه بيا په نظر کښې
ونيسئ چې مخکې مو تعريف کړل۔
ورته «تعريف» به داسې وي:
\[
  g(t) =
  \begin{cases}
    0 & \text{\RL{ که \LR{$t$} يو تورے وي}}\\
   \max(g(s_1), g(s_2)) + 1 & \text{\RL{ که \LR{$t = s_1 \circ s_2$} وي۔}}
  \end{cases}
\]
اوس بې‌قوسه ترم \LR{$\mathrm{a} \circ \mathrm{b} \circ
\mathrm{c} \circ \mathrm{d}$} په نظر
کښې ونيسئ۔ دا له يوې څخه په زياتو
طريقو لوستلاے شو؛ د بېلګې په توګه
د \LR{$s_1 \circ s_2$} په توګه، په دې
ډول چې:
\begin{align*}
  s_1 & = \mathrm{a} \text{\RL{ او}} &
  s_2 & = \mathrm{b} \circ \mathrm{c} \circ \mathrm{d},
\intertext{\RL{يا د \LR{$r_1 \circ r_2$} په توګه، په دې ډول چې:}}
  r_1 & = \mathrm{a} \circ b \text{\RL{ او}} &
  r_2 &= \mathrm{c} \circ \mathrm{d}.
\end{align*}
د لومړي لوست له مخې د \LR{$g$} حساب
به راکړي:
\begin{align*}
  g(s_1 \circ s_2) & =
  \max(g(\mathrm{a}), g(\mathrm{b} \circ \mathrm{c} \circ \mathrm{d})) + 1 =\\ & =
  \max(0,2) + 1 = 3
  \intertext{\RL{خو د بل لوست له مخې ترلاسه کوو:}}
  g(r_1 \circ r_2) & =
  \max(g(\mathrm{a} \circ \mathrm{b}), g(\mathrm{c} \circ \mathrm{d})) + 1=\\ &
  = \max(1,1) + 1 = 2
\end{align*}
خو يوه تابع بايد تل يو يکتا ارزښت
ورکړي؛ نو د~\LR{$g$} دغه «تعريف» هېڅ
تابع نۀ تعريفوي۔
\olpseganchor{OLP-0619-B020}{end}

\olpseganchor{OLP-0619-B021}{start}
\begin{prob}
د تابع \LR{$l$} استقرايي تعريف ورکړئ،
چې \LR{$l(t)$} د نښو شمېر دے په
ښه‌جوړ ترم~\LR{$t$} کښې۔
\end{prob}
\olpseganchor{OLP-0619-B021}{end}

\olpseganchor{OLP-0619-B022}{start}
\begin{prob}
پر ښه‌جوړو ترمونو \LR{$t$} په جوړښتي
استقرا ثابت کړئ چې \LR{$f(t) < l(t)$}،
چې \LR{$l(t)$} د نښو شمېر دے په~\LR{$t$}
کښې او \LR{$f(t)$} د \LR{$t$} ژوروالے دے،
لکه چې په \olref[mth][ind][rel]{defn:depth}
کښې تعريف شوے دے۔
\end{prob}
\olpseganchor{OLP-0619-B022}{end}

\olpseganchor{OLP-0685-B004}{start}
\olpart{pt}{د ثبوتونو تيوري}
\olpseganchor{OLP-0685-B004}{end}

\olpseganchor{OLP-0685-B005}{start}
\begin{editorial}
  د ثبوتونو د تيورۍ مواد لا نيمګړي او تجربوي دي، او د اصلي
  PDF په جوړونه کښې لا نۀ دي ورګډ شوي۔
\end{editorial}
\olpseganchor{OLP-0685-B005}{end}

\olpseganchor{OLP-0712-B004}{start}
\olchapter{pt}{seq}{د سېکوېنټ حساب}
\olpseganchor{OLP-0712-B004}{end}

\olpseganchor{OLP-0700-B004}{start}
\olfileid{pt}{seq}{int}
\olpseganchor{OLP-0700-B004}{end}

\olpseganchor{OLP-0700-B005}{start}
\olsection{سريزه}
\olpseganchor{OLP-0700-B005}{end}

\olpseganchor{OLP-0700-B006}{start}
د سېکوېنټ حسابونه د ثبوت نظامونه دي چې په 1930مو کلونو
کښې ګېرهرډ ګېنټسن (Gerhard Gentzen) معرفي کړل۔ موخه ئې تر
ډېره د شيانو د ثابتولو کولو عملي طريقه نۀ وه۔ بلکې ګېنټسن
دا د دې دپاره کارول چې د ثبوتونه د ممکنه جوړښت او خاصيتونو
په اړه ځينې پايلې ثابتې کړے شي۔ دا ډول پوښتنې کټ مټ هغه
پوښتنې دي چې د ثبوت تيوري ورسره سروکار لري۔
\olpseganchor{OLP-0700-B006}{end}

\olpseganchor{OLP-0700-B007}{start}
د سېکوېنټ حسابونه، لکه نوم ئې چې ښيي، پر \emph{سېکوېنټونو}
کار کوي، پر فارمولونه نۀ۔ سېکوېنټ د فارمولونه د دوو
لړيو يا څوګونو مجموعو يوه جوړه ده۔ د ګېنټسن اصلي نظامونو
(\Log{LK} او \Log{LJ}) لړۍ کارولې؛ اوس د ثبوت د تيورۍ
څېړونکي څوګونې مجموعې غوره ګڼي۔ څوګونې مجموعه د مجموعې
په شان ده، خو يو غړے له يو ځل څخه زيات لرلے شي۔
خو لکه په مجموعو کښې، د غړي ترتيب اهميت نۀ لري۔
نو عبارتونه \LR{$\{!A, !A, !B\}$} او \LR{$\{!A, !B, !A\}$} هماغه يوه
څوګونې مجموعه بيانوي، خو دا څوګونې مجموعه له \LR{$\{!A,
!B\}$} او له \LR{$\{!A, !A, !B, !B\}$} څخه توپير لري، ځکه چې
په وروستنيو کښې د \LR{$!A$} او \LR{$!B$} شمېر د لومړۍ په نسبت بېل دے۔
\olpseganchor{OLP-0700-B007}{end}

\olpseganchor{OLP-0700-B008}{start}
په دوديزه توګه، د ثبوت په تيورۍ کښې د فارمولونه څوګونې
مجموعې په لويو يوناني تورو ليکل کېږي۔ نو کله چې \LR{$\Gamma$}، \LR{$\Delta$}،
\LR{$\Pi$}، \LR{$\Lambda$} يا \LR{$\Theta$} ليکو، موخه مو د هغه منطق د
فارمولونه څوګونې مجموعې وي چې کار پرې کوو۔ همداراز، په
دوديزه توګه د ثبوت د تيورۍ څېړونکي د \LR{$\tuple{\Gamma, \Delta}$}
د ليکلو پر ځاے دواړه برخې د سېکوېنټ په يوه نښه بېلوي؛
موږ~\LR{$\Sequent$} کاروو۔
\olpseganchor{OLP-0700-B008}{end}

\olpseganchor{OLP-0700-B009}{start}
\begin{defn}[سېکوېنټ]
\emph{سېکوېنټ} د دې بڼې يو عبارت دے:
\[
\Gamma \Sequent \Delta
\]
چې \LR{$\Gamma$} او \LR{$\Delta$} د فارمولونه متناهي، او ښايي
تشې، څوګونې مجموعې دي۔ \LR{$\Gamma$} ته \emph{مقدم} او
\LR{$\Delta$} ته \emph{تالي} ويل کېږي۔
\end{defn}
\olpseganchor{OLP-0700-B009}{end}

\olpseganchor{OLP-0700-B010}{start}
فرض کړئ \LR{$!G$}~په~\LR{$\Gamma$} کښې د فارمولونه عطفي ترکيب دے
(تش عطفي ترکيب~\LR{$\ltrue$} ګڼو)، او \LR{$!D$}~په~\LR{$\Delta$} کښې د
فارمولونه فصلي ترکيب دے (تش فصلي ترکيب~\LR{$\lfalse$} ګڼو)۔
نو سېکوېنټ \LR{$\Gamma \Sequent \Delta$} په يو جوړښت~\LR{$\Struct{M}$}
کښې رښتينے دے، که او يوازې که \LR{$!G \lif !D$} رښتينے وي۔
په کلاسيکي منطق کښې دا له دې خبرې سره يو شان ده چې يا
په~\LR{$\Gamma$} کښې د فارمولونه څخه يوه دروغ وي، يا
په~\LR{$\Delta$} کښې د فارمولونه څخه يوه رښتينې وي۔
\olpseganchor{OLP-0700-B010}{end}

\olpseganchor{OLP-0700-B011}{start}
که \LR{$\Gamma$} د فارمولونه يوه څوګونې مجموعه وي، نو
\LR{$\Gamma, !A$} (يا \LR{$!A, \Gamma$}) د هغې پايلې دپاره ليکو
چې \LR{$!A$} پر~\LR{$\Gamma$} ورزيات کړو، يعنې يوه څرګندېدنه ورزياتوو۔
که \LR{$\Delta$} هم د فارمولونه يوه څوګونې مجموعه
وي، نو \LR{$\Gamma, \Delta$} د دواړو لړيو د څوګونو مجموعو اتحاد
دے؛ که \LR{$\Gamma$} د \LR{$!A$} د \emph{تکرار شمېر}~n سره ولري
(يعنې \LR{$n$} د~\LR{$!A$} نقلونه ولري)، او \LR{$\Delta$} ئې د تکرار
شمېر~\LR{$m$} سره ولري، نو \LR{$\Gamma, \Delta$} د \LR{$!A$} د تکرار
شمېر~\LR{$n+m$} سره لري۔ که د سېکوېنټ په يوه خوا کښې څوګونې
مجموعه تشه وي، عموماً هلته تش ځاے پرېږدو۔ د بېلګې په توګه،
\LR{$\Sequent !A$} هغه سېکوېنټ دے چې مقدم ئې تش دے او تالي
ئې \LR{$!A$} د تکرار شمېر~\LR{$1$} سره لري۔
\olpseganchor{OLP-0700-B011}{end}

\olpseganchor{OLP-0700-B012}{start}
ثبوت د سېکوېنټ په حساب کښې د سېکوېنټونو يوه ونه ده:
تر ټولو پورته، يا ابتدايي، سېکوېنټونه په اړوند نظام کښې
بديهي اصول وي؛ او که يو سېکوېنټ له يوه يا دوو نورو سېکوېنټونو
لاندې وي، بايد په نظام کښې د استنتاج د يوې قاعدې له مخې
ترې په سم ډول راوتلے وي۔ لومړے د کلاسيکي منطق دپاره د
سېکوېنټ دوه بېلابېل نظامونه، \LR{$\Log{G1c}$} او~\LR{$\Log{G3c}$}،
څېړو۔ بديهي اصول او قاعدې ئې په
\cref{pt:seq:int:tab:G1c,pt:seq:int:tab:G3c} کښې ورکړل شوي دي۔
 (د ګېنټسن اصلي نظام \Log{LK}
په \olref[fol][seq][]{chap} کښې بيان شوے دے۔)
\olpseganchor{OLP-0700-B012}{end}

\olpseganchor{OLP-0700-B013}{start}


\olpseganchor{OLP-0700-B013}{end}

\olpseganchor{OLP-0700-B014}{start}
نظامونه نازک توپيرونه لري، او له دې توپيرونو څخه راځي چې
د ثبوت د تيورۍ له مخې بېلابېل خاصيتونه لري۔ \Log{G1c} د
\LR{$!A
\Sequent !A$} بديهي اصول لري چې د نورو فارمولونه اجازه
پکښې نۀ شته۔ \Log{G3c} د \LR{$!C, \Gamma \Sequent \Delta, !C$}
په بڼه بديهي اصول لري چې په \LR{$\Gamma$} او~\LR{$\Delta$} کښې ئې
«څنګ» فارمولونه هر ډول کېدے شي؛ خو په دواړو خواوو کښې
راغلے فارمول~\LR{$!C$} بايد \emph{اتومي} وي۔ \Log{G1c} د
\LR{$\LeftR{\land}$} او \LR{$\RightR{\lor}$} له قواعدو څخه د هرې يوې
دوه بڼې لري، خو \Log{G3c} د هرې يوې يوازې يوه بڼه لري۔
او \Log{G1c} اضافي «جوړښتي قاعدې» لري چې «انقباض»
(\LeftR{\Contraction}، \RightR{\Contraction}) او «کمزوري کول»
(\LeftR{\Weakening}، \RightR{\Weakening}) بلل کېږي۔
\olpseganchor{OLP-0700-B014}{end}

\olpseganchor{OLP-0700-B015}{start}
نظامونه \Log{G1c} او \Log{G3c} د کلاسيکي منطق دپاره سم
او بشپړ دي۔ په بله وينا، هر سېکوېنټ چې په دې نظامونو کښې
د ثابتولو وړ وي، په هر جوړښت کښې رښتينے وي؛ او هر
سېکوېنټ چې په هر جوړښت کښې رښتينے وي، ثبوت
لري۔ نو دا پوښتنه چې \emph{آيا} يو سېکوېنټ ثبوت لري،
د منطق په نورو طريقو، لکه د موډل د تيورۍ په طريقو، هم
ځوابېدې شي۔ او ځکه چې دواړه نظامونه کټ مټ هغه سېکوېنټونه
ثابت کوي چې په هر جوړښت کښې رښتيني وي، نو په
ځانګړې توګه هماغه سېکوېنټونه ثابتوي۔
\olpseganchor{OLP-0700-B015}{end}

\olpseganchor{OLP-0700-B016}{start}
خو د ثبوت تيوري يوازې له دې سره سروکار نۀ لري چې \emph{آيا}
يو څۀ ثابتېدے شي، بلکې دا هم څېړي چې \emph{څنګه}۔ د ثبوت د
تيورۍ څېړونکي دا ډول پوښتنې څېړي: «آيا د يوې ځانګړې قاعدې
له کارولو څخه تل ډډه کېدے شي؟»، «آيا دا قاعده نظام ته داسې
ورزياتېدې شي چې نوي سېکوېنټونه د ثبوت وړ نۀ کړي؟»، يا
«آيا په يوه نظام کښې ثبوتونه په بل نظام کښې په ثبوتونه
بدلېدې شي؟» که د داسې يوې پوښتنې ځواب «هو» وي، دا حقيقت
اکثر د ثبوتونه پر پېچلتيا د استقرا په وسيله، يا په معادل
ډول د ثبوتونه پر جوړښت د استقرا په وسيله، ثابتېږي۔
له دې څخه يوازې د حقيقت ثبوت نۀ ترلاسه کېږي، لکه دا چې که
\Log{G1c} يو سېکوېنټ ثابتوي نو \Log{G3c} هم هماغه سېکوېنټ
ثابتوي، بلکې يوه کړنلاره هم ترلاسه کېږي، لکه په~\Log{G1c}
کښې ثبوتونه په~\Log{G3c} کښې په ثبوتونه بدلول۔
\olpseganchor{OLP-0700-B016}{end}

\olpseganchor{OLP-0700-B017}{start}
د ثبوت په تيورۍ کښې يوه مرکزي پايله \emph{د پرېکون لرې کولو
قضيه} ده۔ دا ښيي چې د پرېکون قاعده
\[
\Axiom$\Gamma \fCenter \Delta, !A$
\Axiom$!A, \Pi \fCenter \Lambda$
\RightLabel{\Cut}
\BinaryInf$\Gamma, \Pi \fCenter \Delta, \Lambda$
\DisplayProof
\]
د سېکوېنټ نظامونو ته ورزياتولے شو، پرته له دې چې دا بدل شي
چې کوم سېکوېنټونه د ثابتولو وړ دي۔ سربېره پر دې، ثبوت يوه
طريقه راکوي چې هر هغه ثبوت چې د \Log{G1c} (يا \Log{G3c})
قواعد له \Cut{} سره يوځاے کاروي، په داسې ثبوت بدلوي چې
يوازې د \Log{G1c} (يا \Log{G3c}) قواعد کاروي۔ په بله وينا،
دا کړنلاره له هغو ثبوتونه څخه چې \Cut{} کاروي، دغه قاعده
«لرې کوي»؛ نوم ئې هم له همدې امله دے۔
\olpseganchor{OLP-0700-B017}{end}

\olfileid{pt}{seq}{int}
\olpseganchor{OLP-0704-B004}{start}
\begin{table}
\begin{tabular}{@{}c@{}c@{}}
\multicolumn{2}{@{}c@{}}{بديهي اصول}\\[2ex]
\LR{$!A \Sequent !A$} & \LR{$\lfalse \Sequent \quad$}
\\[2ex]
\multicolumn{2}{@{}c@{}}{جوړښتي قواعد}\\[2ex]
\Axiom$ \Gamma \fCenter \Delta $
\RightLabel{\LeftR{\Weakening}}
\UnaryInf$ !A, \Gamma \fCenter \Delta$
\DisplayProof
&
\Axiom$ \Gamma \fCenter \Delta$
\RightLabel{\RightR{\Weakening}}
\UnaryInf$ \Gamma \fCenter \Delta, !A$
\DisplayProof
\\[4ex]
\Axiom$ !A, !A, \Gamma \fCenter \Delta $
\RightLabel{\LeftR{\Contraction}}
\UnaryInf$ !A, \Gamma \fCenter \Delta$
\DisplayProof
&
\Axiom$ \Gamma \fCenter \Delta, !A, !A$
\RightLabel{\RightR{\Contraction}}
\UnaryInf$ \Gamma \fCenter \Delta, !A$
\DisplayProof
\\[4ex]
\multicolumn{2}{@{}c@{}}{منطقي قواعد}\\[2ex]
\Axiom$ \Gamma \fCenter \Delta, !A $
\RightLabel{\LeftR{\lnot}}
\UnaryInf$ \lnot !A, \Gamma \fCenter \Delta$
\DisplayProof
&
\Axiom$!A, \Gamma \fCenter \Delta$
\RightLabel{\RightR{\lnot}}
\UnaryInf$ \Gamma \fCenter \Delta, \lnot !A $
\DisplayProof
\\[4ex]\begin{tabular}[t]{@{}l@{}}
\Axiom$ !A, \Gamma \fCenter \Delta$
\RightLabel{\LeftR{\land}}
\UnaryInf$ !A \land !B, \Gamma \fCenter \Delta$
\DisplayProof
\\[3ex]
\Axiom$!B, \Gamma \fCenter \Delta$
\RightLabel{\LeftR{\land}}
\UnaryInf$!A \land !B, \Gamma \fCenter \Delta$
\DisplayProof\\[3ex]
\end{tabular}
&
\Axiom$\Gamma \fCenter \Delta, !A$
\Axiom$ \Gamma \fCenter \Delta, !B$
\RightLabel{\RightR{\land}}
\BinaryInf$ \Gamma \fCenter \Delta, !A \land !B $
\DisplayProof
\\[4ex]\Axiom$!A, \Gamma \fCenter \Delta$
\Axiom$!B, \Gamma \fCenter \Delta$
\RightLabel{\LeftR{\lor}}
\BinaryInf$!A \lor !B, \Gamma \fCenter \Delta$
\DisplayProof
&
\begin{tabular}[t]{@{}r@{}}
\Axiom$\Gamma \fCenter \Delta, !A$
\RightLabel{\RightR{\lor}}
\UnaryInf$ \Gamma \fCenter \Delta, !A \lor !B$
\DisplayProof
\\[3ex]
\Axiom$ \Gamma \fCenter \Delta, !B$
\RightLabel{\RightR{\lor}}
\UnaryInf$ \Gamma \fCenter \Delta, !A \lor !B$
\DisplayProof\\[3ex]
\end{tabular}
\\[4ex]
\Axiom$ \Gamma \fCenter \Delta, !A$
\Axiom$ !B, \Gamma \fCenter \Delta$
\RightLabel{\LeftR{\lif}}
\BinaryInf$ !A \lif !B, \Gamma \fCenter \Delta$
\DisplayProof
&
\Axiom$ !A, \Gamma \fCenter \Delta, !B$
\RightLabel{\RightR{\lif}}
\UnaryInf$ \Gamma \fCenter \Delta, !A \lif !B $
\DisplayProof
\\[4ex]
\Axiom$ !A(t), \Gamma \fCenter \Delta$
\RightLabel{\LeftR{\lforall}}
\UnaryInf$ \lforall[x][!A(x)],\Gamma \fCenter \Delta$
\DisplayProof
&
\Axiom$ \Gamma \fCenter \Delta, !A(a) $
\RightLabel{\RightR{\lforall}}
\UnaryInf$ \Gamma \fCenter \Delta, \lforall[x][!A(x)]$
\DisplayProof
\\[4ex]
\Axiom$ !A(a), \Gamma \fCenter \Delta $
\RightLabel{\LeftR{\lexists}}
\UnaryInf$ \lexists[x][!A(x)], \Gamma \fCenter \Delta$
\DisplayProof
&
\Axiom$ \Gamma \fCenter \Delta, !A(t) $
\RightLabel{\RightR{\lexists}}
\UnaryInf$ \Gamma \fCenter \Delta, \lexists[x][!A(x)]$
\DisplayProof
\end{tabular}
\caption{د \Log{G1c} قواعد۔ په \LR{$\RightR{\forall}$} او \LR{$\LeftR{\exists}$} کښې ثابت سمبول~\LR{$a$} بايد په وروستي سېکوېنټ کښې رانۀ شي۔ \LR{$\Gamma$} او \LR{$\Delta$} د فارمولونه څوګونې مجموعې دي۔}
\ollabel{tab:G1c}
\end{table}
\olpseganchor{OLP-0704-B004}{end}

\olfileid{pt}{seq}{int}
\olpseganchor{OLP-0707-B004}{start}
\begin{table}
\begin{tabular}{@{}c@{}c@{}}
\multicolumn{2}{@{}c@{}}{بديهي اصول}\\[2ex]
\LR{$!C, \Gamma \Sequent \Delta, !C$} & \LR{$\lfalse, \Gamma \Sequent \Delta$}
\\[2ex]
\multicolumn{2}{@{}c@{}}{منطقي قواعد}\\[2ex]
\Axiom$ \Gamma \fCenter \Delta, !A $
\RightLabel{\LeftR{\lnot}}
\UnaryInf$ \lnot !A, \Gamma \fCenter \Delta$
\DisplayProof
&
\Axiom$!A, \Gamma \fCenter \Delta$
\RightLabel{\RightR{\lnot}}
\UnaryInf$ \Gamma \fCenter \Delta, \lnot !A $
\DisplayProof
\\[4ex]
\Axiom$ !A, !B, \Gamma \fCenter \Delta$
\RightLabel{\LeftR{\land}}
\UnaryInf$ !A \land !B, \Gamma \fCenter \Delta$
\DisplayProof
&
\Axiom$\Gamma \fCenter \Delta, !A$
\Axiom$ \Gamma \fCenter \Delta, !B$
\RightLabel{\RightR{\land}}
\BinaryInf$ \Gamma \fCenter \Delta, !A \land !B $
\DisplayProof
\\[4ex]\Axiom$!A, \Gamma \fCenter \Delta$
\Axiom$!B, \Gamma \fCenter \Delta$
\RightLabel{\LeftR{\lor}}
\BinaryInf$!A \lor !B, \Gamma \fCenter \Delta$
\DisplayProof
&
\Axiom$\Gamma \fCenter \Delta, !A, !B$
\RightLabel{\RightR{\lor}}
\UnaryInf$ \Gamma \fCenter \Delta, !A \lor !B$
\DisplayProof
\\[4ex]
\Axiom$ \Gamma \fCenter \Delta, !A$
\Axiom$ !B, \Gamma \fCenter \Delta$
\RightLabel{\LeftR{\lif}}
\BinaryInf$ !A \lif !B, \Gamma \fCenter \Delta$
\DisplayProof
&
\Axiom$ !A, \Gamma \fCenter \Delta, !B$
\RightLabel{\RightR{\lif}}
\UnaryInf$ \Gamma \fCenter \Delta, !A \lif !B $
\DisplayProof
\\[4ex]
\Axiom$ !A(t), \lforall[x][!A(x)], \Gamma \fCenter \Delta$
\RightLabel{\LeftR{\lforall}}
\UnaryInf$ \lforall[x][!A(x)],\Gamma \fCenter \Delta$
\DisplayProof
&
\Axiom$ \Gamma \fCenter \Delta, !A(a) $
\RightLabel{\RightR{\lforall}}
\UnaryInf$ \Gamma \fCenter \Delta, \lforall[x][!A(x)]$
\DisplayProof
\\[4ex]
\Axiom$ !A(a), \Gamma \fCenter \Delta $
\RightLabel{\LeftR{\lexists}}
\UnaryInf$ \lexists[x][!A(x)], \Gamma \fCenter \Delta$
\DisplayProof
&
\Axiom$ \Gamma \fCenter \Delta, \lexists[x][!A(x)], !A(t) $
\RightLabel{\RightR{\lexists}}
\UnaryInf$ \Gamma \fCenter \Delta, \lexists[x][!A(x)]$
\DisplayProof
\end{tabular}
\caption{د \Log{G3c} قواعد۔ په \LR{$\RightR{\forall}$} او \LR{$\LeftR{\exists}$} کښې ثابت سمبول~\LR{$a$} بايد په وروستي سېکوېنټ کښې رانۀ شي۔ \LR{$\Gamma$} او \LR{$\Delta$} د فارمولونه څوګونې مجموعې دي۔ د پېژندنې په بديهي اصولو کښې فارمول~\LR{$!C$} بايد اتومي وي۔}
\ollabel{tab:G3c}
\end{table}
\olpseganchor{OLP-0707-B004}{end}

\olfileid{pt}{seq}{variant-g1i}
\olpseganchor{OLP-0705-B004}{start}
\begin{table}
\begin{tabular}{@{}c@{}c@{}}
\multicolumn{2}{@{}c@{}}{بديهي اصول}\\[2ex]
\LR{$!A \Sequent !A$} & \LR{$\lfalse \Sequent \quad$}
\\[2ex]
\multicolumn{2}{@{}c@{}}{جوړښتي قواعد}\\[2ex]
\Axiom$ \Gamma \fCenter \Delta $
\RightLabel{\LeftR{\Weakening}}
\UnaryInf$ !A, \Gamma \fCenter \Delta$
\DisplayProof
&
\Axiom$ \Gamma \fCenter $
\RightLabel{\RightR{\Weakening}}
\UnaryInf$ \Gamma \fCenter !A$
\DisplayProof
\\[4ex]
\Axiom$ !A, !A, \Gamma \fCenter \Delta $
\RightLabel{\LeftR{\Contraction}}
\UnaryInf$ !A, \Gamma \fCenter \Delta$
\DisplayProof
&
\\[4ex]
\multicolumn{2}{@{}c@{}}{منطقي قواعد}\\[2ex]
\Axiom$ \Gamma \fCenter !A $
\RightLabel{\LeftR{\lnot}}
\UnaryInf$ \lnot !A, \Gamma \fCenter $
\DisplayProof
&
\Axiom$!A, \Gamma \fCenter $
\RightLabel{\RightR{\lnot}}
\UnaryInf$ \Gamma \fCenter \lnot !A $
\DisplayProof
\\[4ex]\begin{tabular}[t]{@{}l@{}}
\Axiom$ !A, \Gamma \fCenter \Delta$
\RightLabel{\LeftR{\land}}
\UnaryInf$ !A \land !B, \Gamma \fCenter \Delta$
\DisplayProof
\\[3ex]
\Axiom$!B, \Gamma \fCenter \Delta$
\RightLabel{\LeftR{\land}}
\UnaryInf$!A \land !B, \Gamma \fCenter \Delta$
\DisplayProof\\[3ex]
\end{tabular}
&
\Axiom$\Gamma \fCenter !A$
\Axiom$ \Gamma \fCenter !B$
\RightLabel{\RightR{\land}}
\BinaryInf$ \Gamma \fCenter !A \land !B $
\DisplayProof
\\[4ex]\Axiom$!A, \Gamma \fCenter \Delta$
\Axiom$!B, \Gamma \fCenter \Delta$
\RightLabel{\LeftR{\lor}}
\BinaryInf$!A \lor !B, \Gamma \fCenter \Delta$
\DisplayProof
&
\begin{tabular}[t]{@{}r@{}}
\Axiom$\Gamma \fCenter !A$
\RightLabel{\RightR{\lor}}
\UnaryInf$ \Gamma \fCenter !A \lor !B$
\DisplayProof
\\[3ex]
\Axiom$ \Gamma \fCenter !B$
\RightLabel{\RightR{\lor}}
\UnaryInf$ \Gamma \fCenter !A \lor !B$
\DisplayProof\\[3ex]
\end{tabular}
\\[4ex]
\Axiom$ \Gamma \fCenter !A$
\Axiom$ !B, \Gamma \fCenter \Delta$
\RightLabel{\LeftR{\lif}}
\BinaryInf$ !A \lif !B, \Gamma \fCenter \Delta$
\DisplayProof
&
\Axiom$ !A, \Gamma \fCenter !B$
\RightLabel{\RightR{\lif}}
\UnaryInf$ \Gamma \fCenter !A \lif !B $
\DisplayProof
\\[4ex]
\Axiom$ !A(t), \Gamma \fCenter \Delta$
\RightLabel{\LeftR{\lforall}}
\UnaryInf$ \lforall[x][!A(x)],\Gamma \fCenter \Delta$
\DisplayProof
&
\Axiom$ \Gamma \fCenter !A(a) $
\RightLabel{\RightR{\lforall}}
\UnaryInf$ \Gamma \fCenter \lforall[x][!A(x)]$
\DisplayProof
\\[4ex]
\Axiom$ !A(a), \Gamma \fCenter \Delta $
\RightLabel{\LeftR{\lexists}}
\UnaryInf$ \lexists[x][!A(x)], \Gamma \fCenter \Delta$
\DisplayProof
&
\Axiom$ \Gamma \fCenter !A(t) $
\RightLabel{\RightR{\lexists}}
\UnaryInf$ \Gamma \fCenter \lexists[x][!A(x)]$
\DisplayProof
\end{tabular}
\caption{د شهودي منطق د \Log{G1i} قواعد۔ په \LR{$\RightR{\forall}$} او \LR{$\LeftR{\exists}$} کښې ثابت سمبول~\LR{$a$} بايد په وروستي سېکوېنټ کښې رانۀ شي۔ \LR{$\Gamma$} د فارمولونه څوګونې مجموعه ده؛ \LR{$\Delta$} تشه کېدے شي يا يو فارمول لرلے شي۔ د \RightR{\Contraction} قاعده نشته۔ \Log{G1m} هماغه \Log{G1i} دے، خو \RightR{\Weakening} او د \LR{$\lfalse,
\Gamma \Sequent \Delta$} په بڼه بديهي اصول نۀ لري۔}
\ollabel{tab:G1c}
\end{table}
\olpseganchor{OLP-0705-B004}{end}

\olfileid{pt}{seq}{variant-g3i}
\olpseganchor{OLP-0708-B004}{start}
\begin{table}
\begin{tabular}{@{}c@{}c@{}}
\multicolumn{2}{@{}c@{}}{بديهي اصول}\\[2ex]
\LR{$!C, \Gamma \Sequent !C$} & \LR{$\lfalse, \Gamma \Sequent \Delta$}
\\[2ex]
\multicolumn{2}{@{}c@{}}{منطقي قواعد}\\[2ex]
\Axiom$\lnot !A, \Gamma \fCenter !A $
\RightLabel{\LeftR{\lnot}}
\UnaryInf$\lnot !A, \Gamma \fCenter \Delta$
\DisplayProof
&
\Axiom$!A, \Gamma \fCenter$
\RightLabel{\RightR{\lnot}}
\UnaryInf$ \Gamma \fCenter \lnot !A $
\DisplayProof
\\[4ex]
\Axiom$ !A, !B, \Gamma \fCenter \Delta$
\RightLabel{\LeftR{\land}}
\UnaryInf$ !A \land !B, \Gamma \fCenter \Delta$
\DisplayProof
&
\Axiom$\Gamma \fCenter !A$
\Axiom$ \Gamma \fCenter !B$
\RightLabel{\RightR{\land}}
\BinaryInf$ \Gamma \fCenter !A \land !B $
\DisplayProof
\\[4ex]\Axiom$!A, \Gamma \fCenter \Delta$
\Axiom$!B, \Gamma \fCenter \Delta$
\RightLabel{\LeftR{\lor}}
\BinaryInf$!A \lor !B, \Gamma \fCenter \Delta$
\DisplayProof
&
\Axiom$\Gamma \fCenter !A$
\RightLabel{\RightR{\lor}}
\UnaryInf$ \Gamma \fCenter !A \lor !B$
\DisplayProof
\\[3ex]
&
\Axiom$\Gamma \fCenter !B$
\RightLabel{\RightR{\lor}}
\UnaryInf$ \Gamma \fCenter !A \lor !B$
\DisplayProof
\\[4ex]
\Axiom$!A \lif !B, \Gamma \fCenter !A$
\Axiom$!B, \Gamma \fCenter \Delta$
\RightLabel{\LeftR{\lif}}
\BinaryInf$!A \lif !B, \Gamma \fCenter \Delta$
\DisplayProof
&
\Axiom$!A, \Gamma \fCenter !B$
\RightLabel{\RightR{\lif}}
\UnaryInf$ \Gamma \fCenter !A \lif !B $
\DisplayProof
\\[4ex]
\Axiom$ !A(t), \lforall[x][!A(x)], \Gamma \fCenter \Delta$
\RightLabel{\LeftR{\lforall}}
\UnaryInf$ \lforall[x][!A(x)],\Gamma \fCenter \Delta$
\DisplayProof
&
\Axiom$ \Gamma \fCenter !A(a) $
\RightLabel{\RightR{\lforall}}
\UnaryInf$ \Gamma \fCenter \lforall[x][!A(x)]$
\DisplayProof
\\[4ex]
\Axiom$ !A(a), \Gamma \fCenter \Delta $
\RightLabel{\LeftR{\lexists}}
\UnaryInf$ \lexists[x][!A(x)], \Gamma \fCenter \Delta$
\DisplayProof
&
\Axiom$ \Gamma \fCenter !A(t) $
\RightLabel{\RightR{\lexists}}
\UnaryInf$ \Gamma \fCenter \lexists[x][!A(x)]$
\DisplayProof
\end{tabular}
\caption{د شهودي منطق د \Log{G3i} قواعد۔ په \LR{$\RightR{\forall}$} او \LR{$\LeftR{\exists}$} کښې ثابت سمبول~\LR{$a$} بايد په وروستي سېکوېنټ کښې رانۀ شي۔ \LR{$\Gamma$} د فارمولونه څوګونې مجموعه ده؛ \LR{$\Delta$} تشه کېدے شي يا يو فارمول لرلے شي۔ د پېژندنې په بديهي اصولو کښې فارمول~\LR{$!C$} بايد اتومي وي۔ \Log{G1m} د \Log{G1i} هغه بڼه ده چې \RightR{\Weakening} او د \LR{$\lfalse \Sequent \quad$} بديهي اصل نۀ لري۔}
\ollabel{tab:G3c}
\end{table}
\olpseganchor{OLP-0708-B004}{end}

\olfileid{pt}{seq}{variant-lk}
\olpseganchor{OLP-0709-B004}{start}
\begin{table}
\begin{tabular}{@{}c@{}c@{}}
\multicolumn{2}{@{}c@{}}{بديهي اصول}\\[2ex]
\LR{$!A \Sequent !A$} & \LR{$\lfalse \Sequent \quad$}
\\[2ex]
\multicolumn{2}{@{}c@{}}{جوړښتي قواعد}\\[2ex]
\Axiom$ \Gamma \fCenter \Delta $
\RightLabel{\LeftR{\Weakening}}
\UnaryInf$ !A, \Gamma \fCenter \Delta$
\DisplayProof
&
\Axiom$ \Gamma \fCenter \Delta$
\RightLabel{\RightR{\Weakening}}
\UnaryInf$ \Gamma \fCenter \Delta, !A$
\DisplayProof
\\[4ex]
\Axiom$ !A, !A, \Gamma \fCenter \Delta $
\RightLabel{\LeftR{\Contraction}}
\UnaryInf$ !A, \Gamma \fCenter \Delta$
\DisplayProof
&
\Axiom$ \Gamma \fCenter \Delta, !A, !A$
\RightLabel{\RightR{\Contraction}}
\UnaryInf$ \Gamma \fCenter \Delta, !A$
\DisplayProof
\\[4ex]
\Axiom$ \Gamma, !A, !B, \Pi \fCenter \Delta $
\RightLabel{\LeftR{\Exchange}}
\UnaryInf$\Gamma, !B, !A, \Pi \fCenter \Delta$
\DisplayProof
&
\Axiom$ \Gamma \fCenter \Delta, !A, !B, \Lambda$
\RightLabel{\RightR{\Exchange}}
\UnaryInf$ \Gamma \fCenter \Delta, !B, !A, \Lambda$
\DisplayProof
\\[4ex]
\multicolumn{2}{@{}c@{}}{منطقي قواعد}\\[2ex]
\Axiom$ \Gamma \fCenter \Delta, !A $
\RightLabel{\LeftR{\lnot}}
\UnaryInf$ \lnot !A, \Gamma \fCenter \Delta$
\DisplayProof
&
\Axiom$!A, \Gamma \fCenter \Delta$
\RightLabel{\RightR{\lnot}}
\UnaryInf$ \Gamma \fCenter \Delta, \lnot !A $
\DisplayProof
\\[4ex]\begin{tabular}[t]{@{}l@{}}
\Axiom$ !A, \Gamma \fCenter \Delta$
\RightLabel{\LeftR{\land}}
\UnaryInf$ !A \land !B, \Gamma \fCenter \Delta$
\DisplayProof
\\[3ex]
\Axiom$!B, \Gamma \fCenter \Delta$
\RightLabel{\LeftR{\land}}
\UnaryInf$!A \land !B, \Gamma \fCenter \Delta$
\DisplayProof\\[3ex]
\end{tabular}
&
\Axiom$\Gamma \fCenter \Delta, !A$
\Axiom$ \Gamma \fCenter \Delta, !B$
\RightLabel{\RightR{\land}}
\BinaryInf$ \Gamma \fCenter \Delta, !A \land !B $
\DisplayProof
\\[4ex]\Axiom$!A, \Gamma \fCenter \Delta$
\Axiom$!B, \Gamma \fCenter \Delta$
\RightLabel{\LeftR{\lor}}
\BinaryInf$!A \lor !B, \Gamma \fCenter \Delta$
\DisplayProof
&
\begin{tabular}[t]{@{}r@{}}
\Axiom$\Gamma \fCenter \Delta, !A$
\RightLabel{\RightR{\lor}}
\UnaryInf$ \Gamma \fCenter \Delta, !A \lor !B$
\DisplayProof
\\[3ex]
\Axiom$ \Gamma \fCenter \Delta, !B$
\RightLabel{\RightR{\lor}}
\UnaryInf$ \Gamma \fCenter \Delta, !A \lor !B$
\DisplayProof\\[3ex]
\end{tabular}
\\[4ex]
\Axiom$ \Gamma \fCenter \Delta, !A$
\Axiom$ !B, \Gamma \fCenter \Delta$
\RightLabel{\LeftR{\lif}}
\BinaryInf$ !A \lif !B, \Gamma \fCenter \Delta$
\DisplayProof
&
\Axiom$ !A, \Gamma \fCenter \Delta, !B$
\RightLabel{\RightR{\lif}}
\UnaryInf$ \Gamma \fCenter \Delta, !A \lif !B $
\DisplayProof
\\[4ex]
\Axiom$ !A(t), \Gamma \fCenter \Delta$
\RightLabel{\LeftR{\lforall}}
\UnaryInf$ \lforall[x][!A(x)],\Gamma \fCenter \Delta$
\DisplayProof
&
\Axiom$ \Gamma \fCenter \Delta, !A(a) $
\RightLabel{\RightR{\lforall}}
\UnaryInf$ \Gamma \fCenter \Delta, \lforall[x][!A(x)]$
\DisplayProof
\\[4ex]
\Axiom$ !A(a), \Gamma \fCenter \Delta $
\RightLabel{\LeftR{\lexists}}
\UnaryInf$ \lexists[x][!A(x)], \Gamma \fCenter \Delta$
\DisplayProof
&
\Axiom$ \Gamma \fCenter \Delta, !A(t) $
\RightLabel{\RightR{\lexists}}
\UnaryInf$ \Gamma \fCenter \Delta, \lexists[x][!A(x)]$
\DisplayProof
\end{tabular}
\caption{د \Log{LK} قواعد۔ په \LR{$\RightR{\forall}$} او \LR{$\LeftR{\exists}$} کښې ثابت سمبول~\LR{$a$} بايد په وروستي سېکوېنټ کښې رانۀ شي۔ \LR{$\Gamma$}، \LR{$\Delta$} او~\LR{$\Pi$} د فارمولونه مرتبې لړۍ دي۔}
\ollabel{tab:G1c}
\end{table}
\olpseganchor{OLP-0709-B004}{end}

\olfileid{pt}{seq}{variant-mg3i}
\olpseganchor{OLP-0710-B004}{start}
\begin{table}
\begin{tabular}{@{}c@{}c@{}}
\multicolumn{2}{@{}c@{}}{بديهي اصول}\\[2ex]
\LR{$!C, \Gamma \Sequent \Delta, !C$} & \LR{$\lfalse, \Gamma \Sequent \Delta$}
\\[2ex]
\multicolumn{2}{@{}c@{}}{منطقي قواعد}\\[2ex]
\Axiom$\lnot !A, \Gamma \fCenter \Delta, !A $
\RightLabel{\LeftR{\lnot}}
\UnaryInf$\lnot !A, \Gamma \fCenter \Delta$
\DisplayProof
&
\Axiom$!A, \Gamma \fCenter$
\RightLabel{\RightR{\lnot}}
\UnaryInf$ \Gamma \fCenter \Delta, \lnot !A $
\DisplayProof
\\[4ex]
\Axiom$ !A, !B, \Gamma \fCenter \Delta$
\RightLabel{\LeftR{\land}}
\UnaryInf$ !A \land !B, \Gamma \fCenter \Delta$
\DisplayProof
&
\Axiom$\Gamma \fCenter \Delta, !A$
\Axiom$ \Gamma \fCenter \Delta, !B$
\RightLabel{\RightR{\land}}
\BinaryInf$ \Gamma \fCenter \Delta, !A \land !B $
\DisplayProof
\\[4ex]\Axiom$!A, \Gamma \fCenter \Delta$
\Axiom$!B, \Gamma \fCenter \Delta$
\RightLabel{\LeftR{\lor}}
\BinaryInf$!A \lor !B, \Gamma \fCenter \Delta$
\DisplayProof
&
\Axiom$\Gamma \fCenter \Delta, !A, !B$
\RightLabel{\RightR{\lor}}
\UnaryInf$ \Gamma \fCenter \Delta, !A \lor !B$
\DisplayProof
\\[4ex]
\Axiom$!A \lif !B, \Gamma \fCenter \Delta, !A$
\Axiom$ !B, \Gamma \fCenter \Delta$
\RightLabel{\LeftR{\lif}}
\BinaryInf$!A \lif !B, \Gamma \fCenter \Delta$
\DisplayProof
&
\Axiom$ !A, \Gamma \fCenter !B$
\RightLabel{\RightR{\lif}}
\UnaryInf$ \Gamma \fCenter \Delta, !A \lif !B $
\DisplayProof
\\[4ex]
\Axiom$ !A(t), \lforall[x][!A(x)], \Gamma \fCenter \Delta$
\RightLabel{\LeftR{\lforall}}
\UnaryInf$ \lforall[x][!A(x)],\Gamma \fCenter \Delta$
\DisplayProof
&
\Axiom$ \Gamma \fCenter !A(a) $
\RightLabel{\RightR{\lforall}}
\UnaryInf$ \Gamma \fCenter \lforall[x][!A(x)]$
\DisplayProof
\\[4ex]
\Axiom$ !A(a), \Gamma \fCenter \Delta $
\RightLabel{\LeftR{\lexists}}
\UnaryInf$ \lexists[x][!A(x)], \Gamma \fCenter \Delta$
\DisplayProof
&
\Axiom$ \Gamma \fCenter \Delta, \lexists[x][!A(x)], !A(t) $
\RightLabel{\RightR{\lexists}}
\UnaryInf$ \Gamma \fCenter \Delta, \lexists[x][!A(x)]$
\DisplayProof
\end{tabular}
\caption{د شهودي منطق د څو تاليز \Log{mG3i} قواعد۔ په \LR{$\RightR{\forall}$} او \LR{$\LeftR{\exists}$} کښې ثابت سمبول~\LR{$a$} بايد په وروستي سېکوېنټ کښې رانۀ شي۔ \LR{$\Gamma$} او \LR{$\Delta$} د فارمولونه څوګونې مجموعې دي۔ د پېژندنې په بديهي اصولو کښې فارمول~\LR{$!C$} بايد اتومي وي۔ د اصلي عبارت په پای کښې \Log{mG1m} د \Log{mG1i} په نسبت د \LR{$\lfalse,
\Gamma \Sequent \Delta$} بديهي اصولو له حذف سره يادېږي۔ دا نومونه د دې جدول له نوم سره يو شان نۀ دي۔ د نفي کيڼه قاعده له يوه اتومي فورمول او د هغه له نفي څخه بل فورمول لا هم ورکوي؛ نو يوازې د دې بديهي اصل حذف د لږ تر لږه منطق کره تعريف نۀ دے۔}
\ollabel{tab:G3c}
\end{table}
\olpseganchor{OLP-0710-B004}{end}

\olpseganchor{OLP-0711-B004}{start}
\olfileid{pt}{seq}{rp}
\olpseganchor{OLP-0711-B004}{end}

\olpseganchor{OLP-0711-B005}{start}
\olsection{قواعد او ثبوتونه}
\olpseganchor{OLP-0711-B005}{end}

\olpseganchor{OLP-0711-B006}{start}
لومړے ځينې تعريفونه ورکوو او د اړوندو اصطلاحاتو معنا ټاکو۔
\olpseganchor{OLP-0711-B006}{end}

\olpseganchor{OLP-0711-B007}{start}
د بېلګې په توګه، د \Log{G3c} هغه دوه قواعد وګورئ چې په \LR{$!A \land !B$} لرونکو سېکوېنټونو کښې له \LR{$!A$} او~\LR{$!B$} لرونکو سېکوېنټونو څخه د ثابتولو اجازه ورکوي:
\[
\Axiom$ !A, !B, \Gamma \fCenter \Delta$
\RightLabel{\LeftR{\land}}
\UnaryInf$ !A \land !B, \Gamma \fCenter \Delta$
\DisplayProof
\qquad
\Axiom$\Gamma \fCenter \Delta, !A$
\Axiom$ \Gamma \fCenter \Delta, !B$
\RightLabel{\RightR{\land}}
\BinaryInf$ \Gamma \fCenter \Delta, !A \land !B $
\DisplayProof
\]
د کرښې د پاسه سېکوېنټونو ته \emph{مقدمې} او تر کرښې لاندې سېکوېنټ ته \emph{پايله} وايو۔ په هره قاعده کښې د \LR{$\Gamma$} او \LR{$\Delta$} فارمولونه \emph{سياق} يا \emph{اړخيز فارمولونه} بلل کېږي۔ په پايله کښې هغه فارمول چې د سياق برخه نۀ ده، \emph{اصلي} فارمول بلل کېږي؛ او په مقدماتو کښې له سياق څخه بهر فارمولونه ته \emph{فعال} يا \emph{مرستيال} فارمولونه وايو۔
\olpseganchor{OLP-0711-B007}{end}

\olpseganchor{OLP-0711-B008}{start}
د کميت‌ښودونکو قواعد لږ پېچلي دي۔ د بېلګې په توګه، په \Log{G1c} کښې د \LR{$\lforall$} قواعد داسې دي:
\[\Axiom$ !A(t), \Gamma \fCenter \Delta$
\RightLabel{\LeftR{\lforall}}
\UnaryInf$ \lforall[x][!A(x)],\Gamma \fCenter \Delta$
\DisplayProof
\qquad
\Axiom$ \Gamma \fCenter \Delta, !A(a) $
\RightLabel{\RightR{\lforall}}
\UnaryInf$ \Gamma \fCenter \Delta, \lforall[x][!A(x)]$
\DisplayProof
\]
په \LeftR{\lforall} کښې فعال فارمول~\LR{$!A(t)$} دے، چې \LR{$t$}~تړلے ترم دے، يعنې متغير نۀ لري۔ دا استنتاج هغه وخت صحيح، يعنې د \LeftR{\lforall} له شېما سره برابر دے، چې داسې فارمول~\LR{$!A$}، متغير~\LR{$x$} او تړلے ترم~\LR{$t$} شته چې د پايلې اصلي فارمول \LR{$\lforall[x][!A]$} او د مقدمې فعال فارمول~\LR{$\Subst{!A}{t}{x}$} وي۔ د \RightR{\lforall} قاعده هم همداسې ده، خو د هر ترم~\LR{$t$} پر ځاے په مقدمه کښې فعال فارمول بايد~\LR{$\Subst{!A}{a}{x}$} وي، چې \LR{$a$} يو ثابت سمبول دے۔ دغه ثابت د \RightR{\lforall} د استنتاج \emph{ځانګړے متغير} بلل کېږي۔ استنتاج يوازې هغه وخت صحيح دے چې په لاندې سېکوېنټ کښې~\LR{$a$} نۀ وي؛ يعنې \LR{$a$} په \LR{$\Gamma$}، \LR{$\Delta$} يا~\LR{$!A$} کښې رانۀ شي۔ دې ته \emph{د ځانګړي متغير شرط} وايو۔
\olpseganchor{OLP-0711-B008}{end}

\olpseganchor{OLP-0711-B009}{start}
د \Log{G1c} او \Log{G3c} په څېر سېکوېنټي نظامونه د شېما په بڼه له ټاکلو قواعدو او ځينو سېکوېنټونو څخه جوړ دي چې بديهي اصول ګڼل کېږي او هغه هم د شېما په بڼه ټاکل شوي دي۔ په دغه نظام کښې ثبوتونه د سېکوېنټونو متناهي ونې دي چې له بديهي اصولو څخه د استنتاج په قواعدو، په استقرايي ډول جوړېږي۔ هره پاڼه بديهي اصل دے؛ هر بل سېکوېنټ د نظام په يوه قاعده د پاسه نېغو سېکوېنټونو څخه اخيستل کېږي۔
\olpseganchor{OLP-0711-B009}{end}

\olpseganchor{OLP-0711-B010}{start}
\begin{defn}\ollabel{defn:proof}
په سېکوېنټي نظام کښې \emph{ثبوتونه} د سېکوېنټونو متناهي ونې دي چې داسې استقرايي تعريفېږي:
\begin{enumerate}
  \item\ollabel{defn:prf:axiom} هر بديهي اصل ثبوت دے۔
  \item\ollabel{defn:prf:one-prem} که \LR{$\pi_1$} داسې ثبوت وي چې په سېکوېنټ~\LR{$S_1$} ختمېږي، او \LR{$S$} هغه سېکوېنټ وي چې له~\LR{$S_1$} څخه په يوه‌مقدمه‌ييزه قاعده~\LR{$R$} اخيستل کېږي، نو
  \[
  \AxiomC{}
  \RightLabel{\LR{$\pi_1$}}
  \DeduceC{\LR{$S_1$}}
  \RightLabel{\LR{$R$}}
  \UnaryInfC{\LR{$S$}}
  \DisplayProof
  \]
  ثبوت دے۔
  \item\ollabel{defn:prf:two-prem} که \LR{$\pi_1$} او \LR{$\pi_2$} داسې ثبوتونه وي چې په ترتيب سره په سېکوېنټونو~\LR{$S_1$} او~\LR{$S_2$} ختمېږي، او \LR{$S$} هغه سېکوېنټ وي چې له~\LR{$S_1$} او \LR{$S_2$} څخه په دوه‌مقدمه‌ييزه قاعده~\LR{$R$} اخيستل کېږي، نو
  \[
  \AxiomC{}
  \RightLabel{\LR{$\pi_1$}}
  \DeduceC{\LR{$S_1$}}
  \AxiomC{}
  \RightLabel{\LR{$\pi_2$}}
  \DeduceC{\LR{$S_2$}}
  \RightLabel{\LR{$R$}}
  \BinaryInfC{\LR{$S$}}
  \DisplayProof
  \]
  ثبوت دے۔
\end{enumerate}
\end{defn}
\olpseganchor{OLP-0711-B010}{end}

\olpseganchor{OLP-0711-B011}{start}
\Olref{defn:proof} ثبوتونه د سېکوېنټونو د ونو په توګه تعريفوي۔ ونې داسې انګېرو چې «پورته» لوري ته غځېږي۔ تر ټولو پورته غوټې، يعنې پاڼې، بديهي اصول دي او \emph{لومړني سېکوېنټونه} بلل کېږي۔ تر ټولو لاندې سېکوېنټ ته \emph{وروستے سېکوېنټ} وايو۔ که \LR{$\pi$} د وروستي سېکوېنټ~\LR{$S$} ثبوت وي، نو \LR{$\pi \Proves S$} هم ليکو۔ که د سېکوېنټ~\LR{$S$} ثبوت شته، \LR{$\Proves S$} ليکو؛ کله کله د \LR{$\Proves$} نښې کيڼ لوري ته هغه نظام هم ليکو چې ثبوت پکې شته، لکه \LR{$\Log{G1c} \Proves !A, !B \Sequent !A \land
!B$}۔ له لومړنيو سېکوېنټونو پرته په ثبوت کښې هر سېکوېنټ د کومې قاعدې~\LR{$R$} د استنتاج \emph{پايله} ده۔ يو يا دوه سېکوېنټونه چې په ثبوت کښې نېغ د هغه د پاسه دي، يعنې د غوټې مورني غړي، د استنتاج \emph{مقدمې} دي۔
\olpseganchor{OLP-0711-B011}{end}

\olpseganchor{OLP-0711-B012}{start}
هغه ثبوتونه چې د ثبوت په استقرايي جوړولو کښې کارېږي، د هغه \emph{فرعي ثبوتونه} بلل کېږي۔ هغه فرعي ثبوتونه چې د استنتاج په مقدماتو ختمېږي، د هغه استنتاج \emph{نېغ فرعي ثبوتونه} دي۔
\olpseganchor{OLP-0711-B012}{end}

\olpseganchor{OLP-0711-B013}{start}
ډېر ځله د ثبوتونه پېچلتيا په يوه طبيعي شمېره اندازه کوو۔ په ثبوت کښې د سېکوېنټونو شمېر اخيستلے شو، خو ډېر ځله د ثبوت لوړوالے غوره اندازه وي: د وروستي سېکوېنټ او يوه لومړني سېکوېنټ تر منځ په تر ټولو اوږده څانګه کښې د استنتاجونو شمېر دے۔ استقرايي تعريف ئې هم ورکولے شو۔
\olpseganchor{OLP-0711-B013}{end}

\olpseganchor{OLP-0711-B014}{start}
\begin{defn}
د ثبوت~\LR{$\pi$} \emph{لوړوالے}~\LR{$\pheight{\pi}$} داسې استقرايي تعريفېږي:
\begin{enumerate}
  \item که \LR{$\pi$} يوازې يو بديهي اصل وي، نو \LR{$\pheight{\pi} = 0$}۔
  \item که \LR{$\pi$} په يوه‌مقدمه‌ييز استنتاج ختمېږي چې نېغ فرعي ثبوت~\LR{$\pi_1$} لري، نو \LR{$\pheight{\pi} = \pheight{\pi_1} + 1$}۔
  \item که \LR{$\pi$} په دوه‌مقدمه‌ييز استنتاج ختمېږي چې نېغ فرعي ثبوتونه~\LR{$\pi_1$} او~\LR{$\pi_2$} لري، نو \LR{$\pheight{\pi} =
  \max(\pheight{\pi_1}, \pheight{\pi_2}) + 1$}۔
\end{enumerate}
که سېکوېنټ~\LR{$S$} داسې ثبوت ولري چې لوړوالے ئې \LR{$\le n$} وي، نو \LR{$\Proves[n] S$} ليکو۔
\end{defn}
\olpseganchor{OLP-0711-B014}{end}

\olpseganchor{OLP-0711-B015}{start}
\begin{ex}
په \Log{G3c} کښې د \LR{$!C, \Gamma \Sequent \Delta, !C$} په بڼه هر سېکوېنټ بديهي اصل دے، چې \LR{$!C$} بايد اتومي وي۔ فرض کړئ \LR{$!D$} او \LR{$!E$} اتومي جملې دي۔ نو \LR{$!D, !E \Sequent !D$} او \LR{$!D, !E \Sequent !E$} د \LR{$!C, \Gamma \Sequent \Delta, !C$} بديهي اصل بېلګې دي۔ د بېلګې په توګه، که \LR{$!C \ident !E$}، \LR{$\Gamma = \{!D\}$} او \LR{$\Delta = \emptyset$} په \LR{$!C, \Gamma \Sequent \Delta, !C$} کښې واخلو، نو عين \LR{$!D, !E
\Sequent !E$} سېکوېنټ ترلاسه کوو۔ (په ياد ولرئ چې دلته څوګونې \emph{مجموعې} دي، نو \LR{$!D, !E \Sequent !E$} او \LR{$!E, !D \Sequent !E$} يو سېکوېنټ دي۔)
\olpseganchor{OLP-0711-B015}{end}

\olpseganchor{OLP-0711-B016}{start}
څرنګه چې \LR{$!D, !E \Sequent !D$} د~\Log{G3c} د يوه بديهي اصل بېلګه ده، نو پخپله په~\Log{G3c} کښې ثبوت شمېرل کېږي؛ د \LR{\olref{defn:proof}(\olref{defn:prf:axiom})} له مخې \LR{$!D, !E
\Sequent !E$} هم همداسې دے۔ دې دواړو ته \LR{$\pi_1$} او~\LR{$\pi_2$} وايو۔ د \olref{defn:prf:two-prem} له مخې له دغو دوو ثبوتونه څخه د دوه‌مقدمه‌ييزې قاعدې~\LR{$\RightR{\land}$} په کارولو يو نوے ثبوت (\LR{$\pi_3$}) جوړولے شو:
\[
\Axiom$!D, !E \fCenter !D$
\Axiom$!D, !E \fCenter !E$
\RightLabel{\RightR{\land}}
\BinaryInf$!D, !E  \fCenter !D \land !E $
\DisplayProof
\]
بيا له \LR{$\pi_3$} څخه د يوه‌مقدمه‌ييزې قاعدې~\LR{$\LeftR{\land}$} په کارولو (د \LR{\olref{defn:proof}(\olref{defn:prf:one-prem})} له مخې) دا ثبوت~\LR{$\pi_4$} اخلو:
\[
\Axiom$!D, !E \fCenter !D$
\Axiom$!D, !E \fCenter !E$
\RightLabel{\RightR{\land}}
\insertBetweenHyps{\hspace{5ex}}
\BinaryInf$!D, !E  \fCenter !D \land !E $
\LeftSubproofLabel{\LR{$\pi_3$}}
\RightLabel{\LeftR{\land}}
\UnaryInf$!D \land !E  \fCenter !D \land !E$
\RightSubproofLabel{\LR{$\pi_4$}}
\DisplayProof
\]
په~\LR{$\pi_4$} کښې د وروستي استنتاج نېغ فرعي ثبوت~\LR{$\pi_3$} دے۔ د \LR{$\RightR{\land}$} د استنتاج نېغ فرعي ثبوتونه \LR{$\pi_1$} او~\LR{$\pi_2$} دي۔ د \LR{$\pi_4$} فرعي ثبوتونه \LR{$\pi_1$}، \LR{$\pi_2$}، \LR{$\pi_3$} او پخپله \LR{$\pi_4$} دي۔ لرو: \LR{$\pheight{\pi_1} =
\pheight{\pi_2} = 0$}، \LR{$\pheight{\pi_3} = 1$} او \LR{$\pheight{\pi_4} = 2$}۔ ومو ښودل چې د هر اتومي \LR{$!D$} او~\LR{$!E$} دپاره \LR{$\Log{G3c} \Proves[4] !D \land !E  \fCenter !D
\land !E$}۔
\end{ex}
\olpseganchor{OLP-0711-B016}{end}

\olpseganchor{OLP-0702-B004}{start}
\olfileid{pt}{seq}{ex}
\olpseganchor{OLP-0702-B004}{end}

\olpseganchor{OLP-0702-B005}{start}
\olsection{د ثبوتونو بېلګې}
\olpseganchor{OLP-0702-B005}{end}

\olpseganchor{OLP-0702-B006}{start}
د سېکوېنټونو د ثبوتونه د ورکولو دپاره عموماً دا اسانه وي چې له وروستي سېکوېنټ څخه پورته خوا ته لاړ شو: په هغه کښې فارمول د فعالې فارمول په توګه وټاکو، د اړوندې قاعدې مقدمې د سېکوېنټ د پاسه وليکو، او بيا دغه کړنه تر هغې تکرار کړو چې بديهي اصولو ته ورسېږو۔ د بېلګې په توګه، فرض کړئ غواړو دا سېکوېنټ ثابت کړو:
\[(!C \land !D) \lif !E \Sequent !C \lif (!D \lif !E)\]
\LR{$!C \lif (!D \lif !E)$} ته وګورئ: دا د \LR{$!A \lif !B$} په بڼه دے او د سېکوېنټ په ښۍ خوا کښې راځي۔ که ټول سېکوېنټ د \Log{G1c} له \LR{$\RightR{\lif}$} قاعدې سره برابر کړو،
\[\Axiom$ !A, \Gamma \fCenter \Delta, !B$
\RightLabel{\RightR{\lif}}
\UnaryInf$ \Gamma \fCenter \Delta, !A \lif !B $
\DisplayProof\]
نو راښيي چې بايد \LR{$\Gamma = \{(!C \land !D) \lif !E\}$}، \LR{$\Delta =
\emptyset$}، \LR{$!A \ident !C$} او \LR{$!B \ident !D \lif !E$} واخلو۔ په دې توګه د ثبوت وروستے استنتاج داسې کېدے شي:
\[\Axiom$ !C, (!C \land !D) \lif !E \fCenter !D \lif !E$
\RightLabel{\RightR{\lif}}
\UnaryInf$(!C \land !D) \lif !E \fCenter !C \lif (!D \lif !E)$
\DisplayProof\]
که همدغه لاره بيا وکاروو او \LR{$!D \lif !E$} د فعال فورمول په توګه وټاکو، نو دا ترلاسه کوو:
\[\Axiom$ !D, !C, (!C \land !D) \lif !E \fCenter !E$
\RightLabel{\RightR{\lif}}
\UnaryInf$!C, (!C \land !D) \lif !E \fCenter !D \lif !E$
\RightLabel{\RightR{\lif}}
\UnaryInf$(!C \land !D) \lif !E \fCenter !C \lif (!D \lif !E)$
\DisplayProof\]
اوس بايد \LR{$(!C \land !D) \lif !E$} ته پام وکړو او د \Log{G1c} \LR{$\LeftR{\lif}$} قاعده وکاروو:
\[\Axiom$ \Gamma \fCenter \Delta, !A$
\Axiom$ !B, \Gamma \fCenter \Delta$
\RightLabel{\LeftR{\lif}}
\BinaryInf$ !A \lif !B, \Gamma \fCenter \Delta$
\DisplayProof\]
له دې څخه دې بڼې ته رسېږو:
\[\Axiom$!D, !C \fCenter !E, !C \land !D$
\Axiom$!D, !C, !E \fCenter !E$
\RightLabel{\LeftR{\lif}}
\BinaryInf$ !D, !C, (!C \land !D) \lif !E \fCenter !E$
\RightLabel{\RightR{\lif}}
\UnaryInf$!C, (!C \land !D) \lif !E \fCenter !D \lif !E$
\RightLabel{\RightR{\lif}}
\UnaryInf$(!C \land !D) \lif !E \fCenter !C \lif (!D \lif !E)$
\DisplayProof\]
په کيڼ پورته سېکوېنټ کښې فارمول \LR{$!C \land !D$} د اصلي فورمول په توګه او \LR{$\RightR{\land}$} قاعده کارول دا راکوي:
\[
\Axiom$!D, !C \fCenter !E, !C$
\Axiom$!D, !C \fCenter !E, !D$
\RightLabel{\RightR{\land}}
\BinaryInf$!D, !C \fCenter !E, !C \land !D$
\Axiom$!D, !C, !E \fCenter !E$
\RightLabel{\LeftR{\lif}}
\BinaryInf$ !D, !C, (!C \land !D) \lif !E \fCenter !E$
\RightLabel{\RightR{\lif}}
\UnaryInf$!C, (!C \land !D) \lif !E \fCenter !D \lif !E$
\RightLabel{\RightR{\lif}}
\UnaryInf$(!C \land !D) \lif !E \fCenter !C \lif (!D \lif !E)$
\DisplayProof
\]
تر دې ځايه هر څۀ سم دي۔ پام وکړئ چې ټولې کارولې لارې په \Log{G3c} کښې هم کار کوي، ځکه تر دې ځايه کارولې قاعدې په \Log{G1c} او~\Log{G3c} کښې يو شان دي۔ د ثبوت داسې ټوټې ته رسېدلي يو چې په ټولو پورته سېکوېنټونو کښې هماغه فارمول په کيڼه او ښۍ خوا دواړو کښې شته۔ خو دا لا پوره بديهي اصول نۀ دي۔
\olpseganchor{OLP-0702-B006}{end}

\olpseganchor{OLP-0702-B007}{start}
په \Log{G1c} کښې بايد د ټولو پورته سېکوېنټونو ثبوتونه د \LR{$!A \Sequent !A$} په بڼه له بديهي اصولو څخه ورکړو۔ دا د کمزوري کولو د قواعدو په وسيله په اسانۍ کېږي؛ د بېلګې په توګه، تر ټولو کيڼ پورته سېکوېنټ \LR{$!D, !C \Sequent
!E, !C$} په لاندې ډول ثابت کېدے شي:
\[
\Axiom$!C \fCenter !C$
\RightLabel{\LeftR{\Weakening}}
\UnaryInf$!D, !C \fCenter !C$
\RightLabel{\RightR{\Weakening}}
\UnaryInf$!D, !C \fCenter !E, !C$
\DisplayProof
\]
په ټوليزه توګه، په \Log{G1c} کښې ثبوت لرو چې داسې ښکاري:
\[
\Axiom$!C \fCenter !C$
\RightLabel{\LeftR{\Weakening}}
\UnaryInf$!D, !C \fCenter !C$
\RightLabel{\RightR{\Weakening}}
\UnaryInf$!D, !C \fCenter !E, !C$
\Axiom$!D \fCenter !D$
\RightLabel{\LeftR{\Weakening}}
\UnaryInf$!D, !C \fCenter !D$
\RightLabel{\RightR{\Weakening}}
\UnaryInf$!D, !C \fCenter !E, !D$
\RightLabel{\RightR{\land}}
\BinaryInf$!D, !C \fCenter !E, !C \land !D$
\Axiom$!E \fCenter !E$
\RightLabel{\LeftR{\Weakening}}
\UnaryInf$!C, !E \fCenter !E$
\RightLabel{\LeftR{\Weakening}}
\UnaryInf$!D, !C, !E \fCenter !E$
\RightLabel{\LeftR{\lif}}
\BinaryInf$ !D, !C, (!C \land !D) \lif !E \fCenter !E$
\RightLabel{\RightR{\lif}}
\UnaryInf$!C, (!C \land !D) \lif !E \fCenter !D \lif !E$
\RightLabel{\RightR{\lif}}
\UnaryInf$(!C \land !D) \lif !E \fCenter !C \lif (!D \lif !E)$
\DisplayProof
\]
د دې ثبوت جوړښت، په تېره بيا لوړوالے ئې، په دې پورې نۀ تړلے چې فارمولونه \LR{$!C$}، \LR{$!D$} او~\LR{$!E$} څۀ دي۔ په پورته ثبوت کښې دغه نښې په هر ډول فارمولونه بدلولے شو، او پايله به بيا هم په~\Log{G1c} کښې ثبوت وي۔ په \Log{G1c} کښې دا ثبوت \emph{سکيماتي} دے۔
\olpseganchor{OLP-0702-B007}{end}

\olpseganchor{OLP-0702-B008}{start}
په \Log{G3c} کښې د پېژندنې بديهي اصول د \LR{$!A, \Gamma \Sequent
\Delta, !A$} په بڼه سېکوېنټونه دي چې \LR{$!A$} بايد اتومي وي۔ نو که څۀ هم \LR{$!D, !C \Sequent
!E, !C$} د \Log{G3c} بديهي اصل ښکاري (\LR{$\Gamma = \{!D\}$}؛ \LR{$\Delta =
\{!E\}$})، خو د پېژندنې دا بڼه هغه وخت \emph{په رښتيا} بديهي اصل ورکوي چې \LR{$!C$} په رښتيا اتومي وي۔ زموږ د ثبوت ټوټه په \Log{G3c} کښې بشپړ ثبوت دے چې \LR{$!C$}، \LR{$!D$} او \LR{$!E$} اتومي فارمولونه وي۔
د دې درو فورمولونو اتومي کېدل کافي دي، لازم شرط نۀ دے؛ د پاتې سياق يو ګډ اتومي فورمول هم هر پورته سېکوېنټ بديهي کولے شي۔ د کيڼې تشې نښې بديهي اصول هم په نظام کښې شته۔
\olpseganchor{OLP-0702-B008}{end}

\olpseganchor{OLP-0702-B009}{start}
که څۀ هم په کره توګه مو لا نۀ دي ښودلي چې
\[\Log{G3c} \Proves (!C \land !D) \lif !E \fCenter !C \lif (!D \lif
!E),\]
خو په عمل کښې همدومره کول راته بس دي او همدغه کولے شو۔ دا ځکه چې د \LR{$!A, \Gamma \Sequent \Delta, !A$} په بڼه هر سېکوېنټ په \Log{G3c} کښې د ثبوت وړ دے، که پخپله \LR{$!A$} اتومي هم نۀ وي۔ دا د~\LR{$!A$} پر ژوروالي~\LR{$\depth{!A}$} د استقرا په وسيله ښودلے شو۔
\olpseganchor{OLP-0702-B009}{end}

\olpseganchor{OLP-0702-B010}{start}
\begin{defn}\ollabel{defn:depth}
د فارمول \emph{ژوروالے}~\LR{$\depth{!A}$} په لاندې ډول تعريفېږي:
\begin{enumerate}
  \item که \LR{$!A$} اتومي وي، نو \LR{$\depth{!A} = 0$}۔
  \item که \LR{$!A \ident \lnot !B$}، \LR{$!A \ident \lforall[x][!B]$} يا \LR{$!A \ident \lexists[x][!B]$} وي، نو \LR{$\depth{!A} = \depth{!B} + 1$}۔
  \item که \LR{$!A \ident (!B \land !C)$}، \LR{$(!B \lor !C)$} يا \LR{$(!B \lif
  !C)$} وي، نو \LR{$\depth{!A} = \max(\depth{!B},\depth{!C}) + 1$}۔
\end{enumerate}
\end{defn}
\olpseganchor{OLP-0702-B010}{end}

\olpseganchor{OLP-0702-B011}{start}
دا ثابتول ګران نۀ دي چې د هر ترم~\LR{$t$} دپاره \LR{$\depth{A(x)} = \depth{A(t)}$} دے۔ زموږ دپاره مهمه خبره دا ده چې په هره قاعده کښې د اصلي فارمول ژوروالے تل د فعالو فارمولونه له ژوروالي څخه زيات وي۔ د بېلګې په توګه، لرو:
\begin{align*}
\depth{P(x)} = \depth{P(a)} & = 0,\\
\depth{\lforall[x][P(x)]} & = 1,\\
\depth{(\lforall[x][P(x)] \lif P(a))} & = \max(1,0) + 1 = 2.
\end{align*}
له دې څخه پر~\LR{$\depth{!A}$} د استقرا په وسيله د ځينو پايلو د ثابتولو دپاره کار اخيستلے شو، لکه لاندې پايله۔
\olpseganchor{OLP-0702-B011}{end}

\olpseganchor{OLP-0702-B012}{start}
\begin{prop}\ollabel{prop:G1c-axioms}
د هر~\LR{$!A$} دپاره، په \LR{$\Log{G1c}$} کښې د \LR{$!A \Sequent !A$} ثبوت شته چې ټول بديهي اصول ئې اتومي دي۔
\end{prop}
\olpseganchor{OLP-0702-B012}{end}

\olpseganchor{OLP-0702-B013}{start}
\begin{proof}
  پر~\LR{$\depth{!A}$} د استقرا په وسيله۔ که \LR{$\depth{!A} = 0$} وي، نو \LR{$!A$} اتومي دے او \LR{$!A \Sequent !A$} له وړاندې په~\Log{G1c} کښې د غوښتل شوې بڼې ثبوت دے۔ که \LR{$\depth{!A} > 0$} وي، نو \LR{$!A$} د لږ ژوروالي له فارمولونه څخه جوړ شوے دے، لکه \LR{$!A \ident \lnot !B$}، \LR{$!A \ident
  (!B \land !C)$} يا \LR{$!A \ident \lforall[x][!B(x)]$}۔ د استقرا فرضيه دا ده چې د ټولو هغو \LR{$!D$} دپاره چې \LR{$\depth{!D} < \depth{!A}$} وي، \LR{$\Log{G1c} \Proves !D \Sequent !D$} له اتومي بديهي اصولو څخه حاصلېږي۔
\olpseganchor{OLP-0702-B013}{end}

\olpseganchor{OLP-0702-B014}{start}
  که \LR{$!A \ident \lnot !B$} وي، نو \LR{$\depth{!B} < \depth{!A}$} دے او د استقرا له فرضيې څخه په \Log{G1c} کښې د \LR{$!B \Sequent !B$} ثبوت~\LR{$\pi_1$} له اتومي بديهي اصولو څخه شته۔ دا د \LR{$\lnot !B \Sequent \lnot !B$} ثبوت ته په لاندې ډول غځولے شو:
  \[
  \AxiomC{}
  \RightLabel{\LR{$\pi_1$}}
  \Deduce$!B \fCenter !B$
  \RightLabel{\LeftR{\lnot}}
  \UnaryInf$\lnot !B, !B \fCenter $
  \RightLabel{\RightR{\lnot}}
  \UnaryInf$\lnot !B \fCenter \lnot !B$
  \DisplayProof
\]
\olpseganchor{OLP-0702-B014}{end}

\olpseganchor{OLP-0702-B015}{start}
  که \LR{$!A \ident (!B \land !C)$} وي، نو \LR{$\depth{!B} < \depth{!A}$} او \LR{$\depth{!C} < \depth{!A}$} دي۔ د استقرا له فرضيې څخه په \Log{G1c} کښې د \LR{$!B
  \Sequent !B$} او \LR{$!C \Sequent !C$} ثبوتونه~\LR{$\pi_1$} او~\LR{$\pi_2$} له اتومي بديهي اصولو څخه شته۔ دا د \LR{$!B \land !C \Sequent !B \land !C$} ثبوت ته په لاندې ډول غځولے شو:
  \[
  \AxiomC{}
  \RightLabel{\LR{$\pi_1$}}
  \Deduce$!B, \fCenter !B$
  \RightLabel{\LeftR{\Weakening}}
  \UnaryInf$!B, !C, \fCenter !B$
  \RightLabel{\LeftR{\land}}
  \UnaryInf$!B \land !C, !C\fCenter !B$
  \RightLabel{\LeftR{\land}}
  \UnaryInf$!B \land !C, !B \land !C \fCenter !B$
  \AxiomC{}
  \RightLabel{\LR{$\pi_2$}}
  \Deduce$!C, \fCenter !C$
  \RightLabel{\LeftR{\Weakening}}
  \UnaryInf$!B, !C, \fCenter !C$
  \RightLabel{\LeftR{\land}}
  \UnaryInf$!B \land !C, !C \fCenter !C$
  \RightLabel{\LeftR{\land}}
  \UnaryInf$!B \land !C, !B \land !C \fCenter !C$
  \RightLabel{\RightR{\land}}
  \BinaryInf$!B \land !C, !B \land !C \fCenter !B \land !C$
  \RightLabel{\LeftR{\Contraction}}
  \UnaryInf$!B \land !C \fCenter !B \land !C$
  \DisplayProof
  \]
\olpseganchor{OLP-0702-B015}{end}

\olpseganchor{OLP-0702-B016}{start}
  اوس فرض کړئ \LR{$!A \ident \lforall[x][B(x)]$}۔ داسې ثابت سمبول~\LR{$c$} واخلئ چې په \LR{$\lforall[x][B(x)]$} کښې نۀ راځي۔ بيا \LR{$\depth{B(c)} =
  \depth{!B(x)} < \depth{!A}$} دے او د استقرا له فرضيې څخه په \Log{G1c} کښې د \LR{$!B(c)
  \Sequent !B(c)$} ثبوت~\LR{$\pi_1$} له اتومي بديهي اصولو څخه شته۔ دا د \LR{$\lforall[x][B(x)] \Sequent \lforall[x][B(x)]$} ثبوت ته په لاندې ډول غځولے شو:
  \[
  \AxiomC{}
  \RightLabel{\LR{$\pi_1$}}
  \Deduce$!B(c) \fCenter !B(c)$
  \RightLabel{\LeftR{\lforall}}
  \UnaryInf$\lforall[x][!B(x)] \fCenter !B(c)$
  \RightLabel{\RightR{\lforall}}
  \UnaryInf$\lforall[x][!B(x)] \fCenter \lforall[x][!B(x)]$
  \DisplayProof
\]
\olpseganchor{OLP-0702-B016}{end}

\olpseganchor{OLP-0702-B017}{start}
  پاتې صورتونه د تمرينونو په توګه پرېښودل کېږي۔
\end{proof}
\olpseganchor{OLP-0702-B017}{end}

\olpseganchor{OLP-0702-B018}{start}
\begin{prob}
  د \olref{prop:G1c-axioms} ثبوت بشپړ کړئ: هغه صورتونه سر ته ورسوئ چې \LR{$!A \ident (!B \lor !C)$}، \LR{$!A \ident (!B \lif !C)$} او \LR{$!A \ident \lexists[x][B(x)]$} وي۔
\end{prob}
\olpseganchor{OLP-0702-B018}{end}

\olpseganchor{OLP-0702-B019}{start}
زموږ د موخې دپاره، په هغه صورت کښې چې \LR{$!A \ident \lnot !B$} وي، دا ثبوت هم کارولے شو:
\[
\AxiomC{}
\RightLabel{\LR{$\pi_1$}}
\Deduce$!B \fCenter !B$
\RightLabel{\RightR{\lnot}}
\UnaryInf$\fCenter !B, \lnot !B$
\RightLabel{\LeftR{\lnot}}
\UnaryInf$\lnot !B \fCenter \lnot !B$
\DisplayProof
\]
خو داسې سېکوېنټي نظامونه هم شته چې دغه لاره پکې کار نۀ کوي۔ شهودي سېکوېنټي حساب \Log{G1i} د \Log{G1c} په شان دے، او ټول سېکوېنټونه داسې محدودوي چې په ښۍ خوا کښې ئې تر يوه زياته فارمول نۀ وي۔ که \LR{$\Delta = \emptyset$} وي، نو لومړنے ثبوت به په~\Log{G1i} کښې هم ثبوت وي، خو دويم به نۀ وي، ځکه په يوه سېکوېنټ کښې \LR{$!B$} او~\LR{$\lnot !B$} دواړه په ښۍ خوا کښې دي۔
له دې سره د شرط د کيڼې قاعدې لومړۍ مقدمه هم بدلېږي: د هغې په ښي لوري کښې يوازې د شرط مقدم وي؛ د پاتې ښي سياق کاپي پکې نۀ راځي۔ يوازې د سېکوېنټونو د شمېر شرط د شهودي نظام د تعريف دپاره کافي نۀ دے۔
\olpseganchor{OLP-0702-B019}{end}

\olpseganchor{OLP-0702-B020}{start}
پام وکړئ چې ان په \Log{G1c} کښې هم، په هغه صورت کښې چې \LR{$!A \ident \lforall[x][!B(x)]$} وي، په عمومي ډول د قواعدو بل ترتيب نۀ شو کارولے۔ که ټاکلے ثابت په کيڼ فورمول کښې په رښتيا راځي، لاندې بڼه ثبوت \emph{نۀ ده}:
  \[
  \AxiomC{}
  \RightLabel{\LR{$\pi_1$}}
  \Deduce$!B(c) \fCenter !B(c)$
  \RightLabel{\RightR{\lforall}}
  \UnaryInf$!B(c) \fCenter \lforall[x][!B(x)]$
  \RightLabel{\LeftR{\lforall}}
  \UnaryInf$\lforall[x][!B(x)] \fCenter \lforall[x][!B(x)]$
  \DisplayProof
\]
ځکه په~\LR{$\RightR{\lforall}$} کښې د ځانګړي متغير شرط ماتېږي۔
که کمي متغير په خپل بدن کښې هېڅ پېښه ونۀ لري، نو ټاکلے ثابت په کيڼ فورمول کښې هم نۀ راځي او د همدې ځانګړي متغير شرط نۀ ماتېږي؛ دا د پورته عمومي ممانعت استثنا ده۔
\olpseganchor{OLP-0702-B020}{end}

\olpseganchor{OLP-0702-B021}{start}
\begin{prop}\ollabel{prop:G3c-axioms}
د \LR{$!A, \Gamma
\Sequent \Delta, !A$} په بڼه هر سېکوېنټ (چې \LR{$!A$} ئې \emph{ضروري نۀ ده} اتومي وي) په~\Log{G3c} کښې د ثبوت وړ دے۔
\end{prop}
\olpseganchor{OLP-0702-B021}{end}

\olpseganchor{OLP-0702-B022}{start}
\begin{proof}
دا ثبوت د \olref{prop:G1c-axioms} ثبوت ته ډېر ورته دے، خو د استقرا په فرضيې کښې بايد احتياط وکړو۔ هغه صورت چې \LR{$n =
\depth{!A} = 0$} وي، بيا هم اسانه دے۔ که \LR{$n=0$} وي، نو \LR{$!A$} اتومي دے او \LR{$!A,
\Gamma \Sequent \Delta, !A$} د~\Log{G3c} بديهي اصل دے۔
\olpseganchor{OLP-0702-B022}{end}

\olpseganchor{OLP-0702-B023}{start}
اوس فرض کړئ \LR{$n>0$}۔ بيا بېل بېل صورتونه څېړو۔ د استقرا فرضيه دا ده چې د ټولو هغو \LR{$!D$} دپاره چې \LR{$\depth{!D} < n$} وي، او د ټولو \LR{$\Pi$} او \LR{$\Lambda$} دپاره، \LR{$\Log{G3c} \Proves !D, \Pi \Sequent \Lambda, !D$} دے۔ دلته د \LR{$!A \ident (!B \land !C)$} صورت ښيو۔ د استقرا فرضيه دوه ځله کاروو: يو ځل د \LR{$!D = !B$}، \LR{$\Pi = !C, \Gamma$} او \LR{$\Lambda
= \Delta$} دپاره، څو د \LR{$!B, !C, \Gamma \Sequent
\Delta, !B$} ثبوت~\LR{$\pi_1$} ترلاسه کړو؛ او بل ځل د \LR{$!D = !C$}، \LR{$\Pi = !B, \Gamma$} او \LR{$\Lambda
= \Delta$} دپاره، څو د \LR{$!B, !C, \Gamma \Sequent
\Delta, !C$} ثبوت~\LR{$\pi_2$} ترلاسه کړو۔ دلته د~\LR{$!B \land !C, \Gamma \Sequent
\Delta, !B \land !C$} ثبوت دے:
\[
\AxiomC{}
\RightLabel{\LR{$\pi_1$}}
\Deduce$!B, !C, \Gamma \fCenter \Delta, !B$
\RightLabel{\LeftR{\land}}
\UnaryInf$!B \land !C, \Gamma \fCenter \Delta, !B$
\AxiomC{}
\RightLabel{\LR{$\pi_2$}}
\Deduce$!B, !C, \Gamma \fCenter \Delta, !C$
\RightLabel{\LeftR{\land}}
\UnaryInf$!B \land !C, \Gamma \fCenter \Delta, !C$
\RightLabel{\RightR{\land}}
\BinaryInf$!B \land !C, \Gamma \fCenter \Delta, !B \land !C$
\DisplayProof
\]
\olpseganchor{OLP-0702-B023}{end}

\olpseganchor{OLP-0702-B024}{start}
د \LR{$!A \ident \lforall[x][!B(x)]$} دپاره، داسې ثابت سمبول واخلئ چې په~\LR{$\lforall[x][!B(x)]$}، \LR{$\Gamma$} يا~\LR{$\Delta$} کښې نۀ وي۔ د استقرا فرضيه د \LR{$!D = !B(c)$}، \LR{$\Pi = \lforall[x][!B(x)], \Gamma$} او \LR{$\Lambda = \Delta$} دپاره تطبيق کړئ۔ دا د \LR{$!B(c),
\lforall[x][!B(x)], \Gamma \Sequent \Delta, !B(c)$} ثبوت~\LR{$\pi_1$} راکوي چې له هغې څخه دا ترلاسه کوو:
\[
\AxiomC{}
\RightLabel{\LR{$\pi_1$}}
\Deduce$!B(c), \lforall[x][!B(x)], \Gamma \fCenter \Delta, !B(c)$
\RightLabel{\LeftR{\lforall}}
\UnaryInf$\lforall[x][!B(x)], \Gamma \fCenter \Delta, !B(c)$
\RightLabel{\RightR{\lforall}}
\UnaryInf$\lforall[x][!B(x)], \Gamma \fCenter \Delta, \lforall[x][!B(x)]$
\DisplayProof
\]
په \LR{$\RightR{\lforall}$} کښې د ځانګړي متغير شرط د~\LR{$c$} د ټاکلو د لارې له مخې پوره کېږي۔
\end{proof}
\olpseganchor{OLP-0702-B024}{end}

\olpseganchor{OLP-0702-B025}{start}
\begin{prob}
  د \olref{prop:G3c-axioms} ثبوت ته دوام ورکړئ: هغه صورتونه سر ته ورسوئ چې \LR{$!A \ident (!B \lif !C)$} او \LR{$!A \ident
  \lexists[x][B(x)]$} وي۔
\end{prob}
\olpseganchor{OLP-0702-B025}{end}

\olpseganchor{OLP-0699-B004}{start}
\olfileid{pt}{seq}{irl}
\olpseganchor{OLP-0699-B004}{end}

\olpseganchor{OLP-0699-B005}{start}
\olsection{د قواعدو تعبير}
\olpseganchor{OLP-0699-B005}{end}

\olpseganchor{OLP-0699-B006}{start}
د سېکوېنټ تعبير راياد کړئ: سېکوېنټ \LR{$\Gamma \Sequent
\Delta$} په جوړښت~\LR{$\Struct{M}$} کښې رښتينے دے، که يا لږ
تر لږه يو \LR{$!A \in \Gamma$} دروغ وي، يا لږ تر لږه يو~\LR{$!A \in \Delta$}
رښتينے وي۔ د~\Log{G3c} منطقي قاعدې داسې لوستلے شو چې د خپلو
اصلي فارمولونه د رښتينولۍ شرطونه په قاعدو کښې راټولوي۔
\RightR{\lif} ته وګورئ۔ اصلي فارمول ئې \LR{$!A \lif !B$} دے۔
دا رښتينے دے، که يا \LR{$!A$} دروغ وي، يا~\LR{$!B$} رښتينے وي۔ نو
سېکوېنټ \LR{$\ \Sequent !A \lif !B$} رښتينے دے، که او يوازې که
\LR{$!A \Sequent !B$} رښتينے وي۔ که د سياق فارمولونه
\LR{$\Gamma$} او \LR{$\Delta$} د دغو سېکوېنټونو کيڼو او ښيو خواوو ته
ورزيات کړو، نو کټ مټ د \RightR{\lif} قاعدې پايله او مقدمه
ترلاسه کوو۔ د~\LeftR{\lif} صورت هم دې ته ورته دے۔ \LR{$!A
\lif !B$} دروغ دے، که \LR{$!A$}~رښتينے او \LR{$!B$}~دروغ وي۔ نو
\LR{$!A \lif !B
\Sequent \ $} رښتينے دے، که او يوازې که دواړه \LR{$\ \Sequent !A$}
او \LR{$!B \Sequent \ $} رښتيني وي۔ د \LeftR{\lif} پايله د خپلو
مقدماتو له عطفي ترکيب سره معادله ده۔ دا بڼه ډاډمنوي چې
د~\Log{G3c} قاعدې سمې دي: که ټول مقدمات رښتيني وي، پايله
رښتينې ده۔ دغه محلي معادلت د معنايي بشپړتيا په دليل کښې کارېږي۔
\olpseganchor{OLP-0699-B006}{end}

\olpseganchor{OLP-0699-B007}{start}
په حقيقت کښې، د هر هغه تړونکي د رښتينولۍ شرطونه، چې ارزښت
ئې يوازې د خپلو غړو د رښتينولۍ په ارزښتونو پورې تړلے وي،
په همدې ډول د سېکوېنټونو د يوې مجموعې په توګه بيانېدې شي۔
دا له دې حقيقت څخه راځي چې په کلاسيکي منطق کښې د رښتينولۍ
هره تابع په عطفي عادي بڼه د فارمول په وسيله بيانېدې شي:
د اتومي او نفي شويو اتومي فارمولونه د فصلي ترکيبونو يو
عطفي ترکيب۔ د بېلګې په توګه، \LR{$!A
\oplus !B$} هغه وخت رښتينے وګڼئ چې د \LR{$!A$} او \LR{$!B$} څخه کټ مټ
يو رښتينے وي، يعنې «انحصاري يا»۔ نو \LR{$!A \oplus !B$} رښتينے دے،
که او يوازې که \LR{$(!A \lor !B) \land
(\lnot !A \lor \lnot !B)$} رښتينے وي؛ او دروغ دے، که
\LR{$(\lnot !A \lor !B)
\land (!A \lor \lnot !B)$} رښتينے وي۔ نو د \LR{$\oplus$} دپاره د
سېکوېنټ د حساب قاعدې داسې معرفي کولے شو:
\begin{align*}
&
\Axiom$\Gamma \fCenter \Delta, !A, !B$
\Axiom$!A, !B, \Gamma \fCenter \Delta$
\RightLabel{\RightR{\oplus}}
\BinaryInf$\Gamma \fCenter \Delta, !A \oplus !B$
\DisplayProof
\\
& \Axiom$!A, \Gamma \fCenter \Delta, !B$
\Axiom$!B, \Gamma \fCenter \Delta, !A$
\RightLabel{\LeftR{\oplus}}
\BinaryInf$!A \oplus !B, \Gamma \fCenter \Delta$
\DisplayProof
\end{align*}
او دا قاعدې به هم سمې او بشپړې وي۔ پايلې ئې رښتينې دي،
که او يوازې که ټول مقدمات رښتيني وي۔
په لومړۍ قاعده کښې اصلي فورمول په ښي اړخ او په دويمه کښې په کيڼ اړخ کښې دے۔
\olpseganchor{OLP-0699-B007}{end}

\olpseganchor{OLP-0699-B008}{start}
\begin{prob}
فرض کړئ \LR{$!A \downarrow !B$} رښتينے دے، که او يوازې که نۀ
\LR{$!A$} او نۀ~\LR{$!B$} رښتينے وي («نۀ دا او نۀ هغه»، يعنې NOR)۔
د دې د \Log{G3c} قاعدې به څۀ وي؟ (داسې دوه سېکوېنټونه ومومئ
چې دواړه په يو وخت رښتيني وي، که او يوازې که \LR{$!A \downarrow !B$}
رښتينے وي۔ له دې څخه د \RightR{\downarrow} قاعده ترلاسه کېږي۔
بيا داسې يو سېکوېنټ ومومئ چې رښتينے وي، که او يوازې که
\LR{$!A \downarrow !B$} دروغ وي۔ له دې څخه د \LeftR{\downarrow}
قاعده ترلاسه کېږي۔)
\end{prob}
\olpseganchor{OLP-0699-B008}{end}

\olpseganchor{OLP-0699-B009}{start}
په \Log{G3c} کښې هر بنسټيز ناثابت تړونکے او هر کميت‌ښودونکے کټ مټ دوه قاعدې
لري: يوه کيڼه قاعده او يوه ښۍ قاعده۔ خو په \Log{G1c} کښې
دوه \LeftR{\land} او دوه \RightR{\lor} قاعدې دي۔ لامل دا دے
چې د ګېنټسن اصلي نظام~\Log{LK}، چې \Log{G1c} د هغې په بڼه
جوړ شوے دے، د يو مهم غير کلاسيکي منطق، يعنې شهودي منطق،
دپاره بله بڼه~\Log{LJ} لرله۔ د شهودي منطق دپاره د يو سم او
بشپړ نظام د ترلاسه کولو په موخه، په~\Log{LJ} کښې د هر
سېکوېنټ ښۍ خوا په دې محدوده ده چې له يوې فارمول څخه
زياتې ونۀ لري۔ په دې سره د~\Log{G3c} د \RightR{\lor} قاعدې
کارول ناشوني کېږي، ځکه مقدمه ئې په ښۍ خوا کښې دواړه \LR{$!A$}
او~\LR{$!B$} لري۔ نو ګېنټسن د \Log{LJ} دپاره د \RightR{\lor}
قواعدو يوه جوړه وکاروله، او هماغه ئې د~\Log{LK} دپاره هم
وکاروله۔ شهودي بڼه~\Log{G1i} د~\Log{G1c} داسې بڼه ده
چې ټولې ښۍ خواوې ئې تر يوې فارمول پورې محدودې دي،
خو د شرط د کيڼې قاعدې لومړۍ مقدمه هم بدله شوې ده۔
(\Log{G3c} هم يوه شهودي بڼه لري، خو د هغې ځينې قاعدې
په~\Log{G3c} کښې له اړوندو قواعدو څخه توپير لري۔ د بېلګې
په توګه، هغه هم د \RightR{\lor} قواعدو يوه جوړه کاروي۔)
په~\Log{G1i} کښې د شرط د کيڼې قاعدې د لومړۍ مقدمې
تالي يوازې د شرط مقدمه ده۔ د پايلې تالي سياق په دويمه مقدمه کښې وي۔
په دې ډول قاعده له غيرتش تالي سره هم کارېدے شي۔
\olpseganchor{OLP-0699-B009}{end}

\olpseganchor{OLP-0699-B010}{start}
دا ليدل اسانه دي چې د~\Log{G1c} جوړښتي قاعدې سمې دي۔
د بېلګې په توګه، که \LR{$\Gamma \Sequent \Delta$} په کيڼه خوا کښې
يوه دروغ فارمول يا په ښۍ خوا کښې يوه رښتينې فارمول
ولري، نو \LR{$\Gamma \Sequent
\Delta, !A$} هم هماغه يوه لري؛ دا چې \LR{$!A$} رښتينے دے که دروغ،
څۀ اغېز نۀ لري۔ نو که د~\RightR{\Weakening} مقدمه رښتينې وي،
پايله ئې هم رښتينې ده۔ په ياد ولرئ چې دلته معکوسه خبره سمه
نۀ ده۔ ښايي پايله ځکه رښتينې وي چې \LR{$!A$} رښتينے دے، خو بيا
دا تضمين نۀ شته چې مقدمه به هم رښتينې وي۔
\olpseganchor{OLP-0699-B010}{end}


\olpseganchor{OLP-0699-B011}{start}
جوړښتي قاعدې ولې ولرو؟ انقباض (\RightR{\Contraction}،
\LeftR{\Contraction}) د بېلګې په توګه د \LR{$\ \Sequent !A
\lor \lnot !A$} د ثابت کولو دپاره اړين دے۔ که له
\LR{$!A \Sequent !A$} څخه پيل وکړئ، لومړے د~\RightR{\lnot} په وسيله
\LR{$\Sequent !A, \lnot !A$} ترلاسه کوئ۔ بيا د \RightR{\lor} دواړه
بڼې پرله‌پسې کارولے شئ: يو ځل د \LR{$!A \lor \lnot
!A$} د ترلاسه کولو دپاره له \LR{$\lnot !A$} څخه، او بل ځل د
\LR{$!A \lor \lnot !A$} د ترلاسه کولو دپاره له~\LR{$!A$} څخه۔ په دې توګه
\LR{$\ \Sequent !A \lor \lnot !A, !A \lor \lnot !A$} ترلاسه کېږي۔
نو په~\Log{G1c} کښې د \LR{$\Sequent !A \lor \lnot !A$} د ترلاسه
کولو دپاره بايد د يوې فارمول دوه څرګندېدنې په ثبوت
کښې «منقبض» کولے شئ:
\[
\Axiom$!A \fCenter !A$
\RightLabel{\RightR{\lnot}}
\UnaryInf$\fCenter !A, \lnot !A$
\RightLabel{\RightR{\lor}}
\UnaryInf$\fCenter !A, !A \lor \lnot !A$
\RightLabel{\RightR{\lor}}
\UnaryInf$\fCenter !A \lor \lnot !A, !A \lor \lnot !A$
\RightLabel{\RightR{\Contraction}}
\UnaryInf$\fCenter !A \lor \lnot !A$
\DisplayProof
\]
که د اټومي فورمول په توګه يو بياني متغير واخلو، له ښي انقباض پرته د پورتني
سېکوېنټ د ثبوت وروستۍ قاعده يا ښے کمزوري کول وي يا د يا
ښۍ قاعده۔ په لومړي حالت کښې مقدمه له دواړو اړخونو تشه ده۔
په دويم کښې تشه کيڼه خوا او په ښۍ خوا کښې يوازې اتوم يا
نفي ئې وي۔ دغه مقدمې نامعتبرې دي۔
په~\Log{G1c} کښې داسې معتبر سېکوېنټونه هم شته چې له
جوړښتي قواعدو پرته ئې ثابت کولے نۀ شو۔ د بېلګې په توګه،
که دوه بېل بياني متغيرونه واخلو، د سېکوېنټ
\LR{$!A \Sequent !B \lif !A$} جوړښتي‌قاعدې‌خالي ثبوت بايد په
\RightR{\lif} پاے ته ورسېږي، ځکه يوازې \LR{$!B \lif !A$} اصلي
کېدے شي۔ مقدمه ئې \LR{$!B, !A \Sequent !A$} ده۔ دا نۀ بديهي
اصل دے او نۀ منطقي عملګر لري۔ په \LeftR{\Weakening} ئې
له بديهي اصل \LR{$!A \Sequent !A$} څخه اخلو:
\[
\Axiom$!A \fCenter !A$
\RightLabel{\LeftR{\Weakening}}
\UnaryInf$!B, !A \fCenter !A$
\RightLabel{\RightR{\lif}}
\UnaryInf$!A \fCenter !B \lif !A$
\DisplayProof
\]
\olpseganchor{OLP-0699-B011}{end}

\olpseganchor{OLP-0699-B012}{start}
په \Log{G3c} کښې د انقباض او کمزوري کولو قواعدو ته هېڅ
اړتيا نۀ شته۔ کمزوري کول په بديهي اصولو کښې له مخکې شامل
دي۔ د انقباض له کارولو څخه په زياتره صورتونو کښې داسې
ډډه کولے شو چې په دوو مقدمو لرونکو قواعدو کښې د سياق \LR{$\Gamma$} او \LR{$\Delta$}
دوه نقلونه په يوه کښې يوځاے کړو۔ دواړه \Log{G1c} او \Log{G3c}
د «سياق شريکولو» داسې قاعدې لري۔ په يوه مقدمه لرونکو
قواعدو کښې انقباض يوازې د \RightR{\lor}، \LeftR{\land}، \RightR{\lexists}
او~\LeftR{\lforall} دپاره اړين دے۔ په~\Log{G3c} کښې د دې
قواعدو بڼې انقباض په بشپړه توګه بې‌اړتيا کوي۔ يو نظام
\Log{G2c} شته (\olref{tab:G2c} وګورئ) چې دوې مقدمې لرونکې
قاعدې ئې بېل بېل سياقونه لري۔ په \Log{G2c} کښې انقباض
له~\Log{G1c} څخه لا زيات مهم دے۔
\olpseganchor{OLP-0699-B012}{end}

\olpseganchor{OLP-0699-B013}{start}

\olpseganchor{OLP-0699-B013}{end}

\olfileid{pt}{seq}{irl}
\olpseganchor{OLP-0706-B004}{start}
\begin{table}
\begin{tabular}{@{}c@{\hskip-5ex}c@{}}
\multicolumn{2}{@{}c@{}}{بديهي اصول}\\[2ex]
\LR{$!A \Sequent !A$} & \LR{$\lfalse \Sequent \quad$}
\\[2ex]
\multicolumn{2}{@{}c@{}}{جوړښتي قواعد}\\[2ex]
\Axiom$ \Gamma \fCenter \Delta $
\RightLabel{\LeftR{\Weakening}}
\UnaryInf$ !A, \Gamma \fCenter \Delta$
\DisplayProof
&
\Axiom$ \Gamma \fCenter \Delta$
\RightLabel{\RightR{\Weakening}}
\UnaryInf$ \Gamma \fCenter \Delta, !A$
\DisplayProof
\\[4ex]
\Axiom$ !A, !A, \Gamma \fCenter \Delta $
\RightLabel{\LeftR{\Contraction}}
\UnaryInf$ !A, \Gamma \fCenter \Delta$
\DisplayProof
&
\Axiom$ \Gamma \fCenter \Delta, !A, !A$
\RightLabel{\RightR{\Contraction}}
\UnaryInf$ \Gamma \fCenter \Delta, !A$
\DisplayProof
\\[4ex]
\multicolumn{2}{@{}c@{}}{منطقي قواعد}\\[2ex]
\Axiom$ \Gamma \fCenter \Delta, !A $
\RightLabel{\LeftR{\lnot}}
\UnaryInf$ \lnot !A, \Gamma \fCenter \Delta$
\DisplayProof
&
\Axiom$!A, \Gamma \fCenter \Delta$
\RightLabel{\RightR{\lnot}}
\UnaryInf$ \Gamma \fCenter \Delta, \lnot !A $
\DisplayProof
\\[4ex]\begin{tabular}[t]{@{}l@{}}
\Axiom$ !A, \Gamma \fCenter \Delta$
\RightLabel{\LeftR{\land}}
\UnaryInf$ !A \land !B, \Gamma \fCenter \Delta$
\DisplayProof
\\[3ex]
\Axiom$!B, \Gamma \fCenter \Delta$
\RightLabel{\LeftR{\land}}
\UnaryInf$!A \land !B, \Gamma \fCenter \Delta$
\DisplayProof\\[3ex]
\end{tabular}
&
\Axiom$\Gamma_1 \fCenter \Delta_1, !A$
\Axiom$ \Gamma_2 \fCenter \Delta_2, !B$
\RightLabel{\RightR{\land}}
\BinaryInf$ \Gamma_1, \Gamma_2 \fCenter \Delta_1, \Delta_2, !A \land !B $
\DisplayProof
\\[4ex]\Axiom$!A, \Gamma_1 \fCenter \Delta_1$
\Axiom$!B, \Gamma_2 \fCenter \Delta_2$
\RightLabel{\LeftR{\lor}}
\BinaryInf$!A \lor !B, \Gamma_1, \Gamma_2 \fCenter \Delta_1, \Delta_2$
\DisplayProof
&
\begin{tabular}[t]{@{}r@{}}
\Axiom$\Gamma \fCenter \Delta, !A$
\RightLabel{\RightR{\lor}}
\UnaryInf$ \Gamma \fCenter \Delta, !A \lor !B$
\DisplayProof
\\[3ex]
\Axiom$ \Gamma \fCenter \Delta, !B$
\RightLabel{\RightR{\lor}}
\UnaryInf$ \Gamma \fCenter \Delta, !A \lor !B$
\DisplayProof\\[3ex]
\end{tabular}
\\[4ex]
\Axiom$ \Gamma_1 \fCenter \Delta_1, !A$
\Axiom$ !B, \Gamma_2 \fCenter \Delta_2$
\RightLabel{\LeftR{\lif}}
\BinaryInf$ !A \lif !B, \Gamma_1, \Gamma_2 \fCenter \Delta_1, \Delta_2$
\DisplayProof
&
\Axiom$ !A, \Gamma \fCenter \Delta, !B$
\RightLabel{\RightR{\lif}}
\UnaryInf$ \Gamma \fCenter \Delta, !A \lif !B $
\DisplayProof
\\[4ex]
\Axiom$ !A(t), \Gamma \fCenter \Delta$
\RightLabel{\LeftR{\lforall}}
\UnaryInf$ \lforall[x][!A(x)],\Gamma \fCenter \Delta$
\DisplayProof
&
\Axiom$ \Gamma \fCenter \Delta, !A(a) $
\RightLabel{\RightR{\lforall}}
\UnaryInf$ \Gamma \fCenter \Delta, \lforall[x][!A(x)]$
\DisplayProof
\\[4ex]
\Axiom$ !A(a), \Gamma \fCenter \Delta $
\RightLabel{\LeftR{\lexists}}
\UnaryInf$ \lexists[x][!A(x)], \Gamma \fCenter \Delta$
\DisplayProof
&
\Axiom$ \Gamma \fCenter \Delta, !A(t) $
\RightLabel{\RightR{\lexists}}
\UnaryInf$ \Gamma \fCenter \Delta, \lexists[x][!A(x)]$
\DisplayProof
\end{tabular}
\caption{د \Log{G2c} قواعد۔ په \LR{$\RightR{\forall}$} او \LR{$\LeftR{\exists}$} کښې ثابت سمبول~\LR{$a$} بايد په وروستي سېکوېنټ کښې رانۀ شي۔ \LR{$\Gamma$} او \LR{$\Delta$} د فارمولونه څوګونې مجموعې دي۔}
\ollabel{tab:G2c}
\end{table}
\olpseganchor{OLP-0706-B004}{end}

\olpseganchor{OLP-0703-B004}{start}
\olfileid{pt}{seq}{qua}
\olpseganchor{OLP-0703-B004}{end}

\olpseganchor{OLP-0703-B005}{start}
\olsection{منظم ثبوتونه او تعويض}
\olpseganchor{OLP-0703-B005}{end}

\olpseganchor{OLP-0703-B006}{start}
د کميت‌ښودونکو قواعد د «ځانګړي متغير د شرطونو» له کبله ځينې ستونزې پيدا کوي۔ دغه شرطونه پر \LeftR{\lexists} او~\RightR{\lforall} تطبيقېږي۔ ډېرې ستونزې داسې لرې کولے شو چې په دغو استنتاجونو کښې ثابت سمبول~\LR{$a$} بيا ونۀ کاروو۔ لاندې تعريف دا خبره کره کوي۔
\olpseganchor{OLP-0703-B006}{end}

\olpseganchor{OLP-0703-B007}{start}
\begin{defn}
ثبوت~\LR{$\pi$} په \Log{G1c} يا \Log{G3c} کښې \emph{منظم} دے که:
\begin{enumerate}
  \item هېڅ ځانګړے متغير~\LR{$a$} د \LeftR{\lexists} يا~\RightR{\lforall} تر يوه زياتو استنتاجونو کښې د ځانګړي متغير په توګه نۀ کارېږي؛ او
  \item که \LR{$a$} د \LeftR{\lexists} يا~\RightR{\lforall} د کوم استنتاج ځانګړے متغير وي، نو په~\LR{$\pi$} کښې يوازې د هماغه استنتاج د پاسه راځي۔
\end{enumerate}
\end{defn}
\olpseganchor{OLP-0703-B007}{end}

\olpseganchor{OLP-0703-B008}{start}
\begin{lem}\ollabel{lem:subst-regular}
که \LR{$\pi$} منظم وي او په \LR{$\Gamma \Sequent \Delta$} پاے ته رسېږي، \LR{$c$} داسې ثابت سمبول وي چې په~\LR{$\pi$} کښې د ځانګړي متغير په توګه نۀ کارېږي، او \LR{$t$} داسې تړلے ترم وي چې د~\LR{$\pi$} هېڅ ځانګړے متغير ونۀ لري، نو \LR{$\Subst{\pi}{t}{c}$} هم منظم ثبوت دے۔ دا له~\LR{$\pi$} څخه داسې ترلاسه کېږي چې د~\LR{$c$} هره پېښه په~\LR{$t$} بدله کړو۔ تړلے بدل ترم ازاد متغير نۀ لري؛ نو د کميت‌ښودونکي تر تړون لاندې د ازاد متغير د نيولو ستونزه هم نشته۔ د عمومي کميت کيڼه او د وجودي کميت ښۍ قاعده د نمونې دپاره تړلے ترم غواړي؛ دا شرط د ځانګړي متغير له شرط څخه بېل دے۔
د \LR{$\Subst{\pi}{t}{c}$} وروستے سېکوېنټ \LR{$\Subst{\Gamma}{t}{c}
\Sequent \Subst{\Delta}{t}{c}$} دے، چې \LR{$\Subst{\Gamma}{t}{c} =
\Setabs{!A(t)}{!A(c) \in \Gamma}$}۔ په سياق کښې د هر فورمول هر نقل جلا بدلېږي؛ که بدل شوي فورمولونه يو شان شي، د هغو د نقلونو شمېر ساتل کېږي۔
\end{lem}
\olpseganchor{OLP-0703-B008}{end}

\olpseganchor{OLP-0703-B009}{start}
\begin{proof}
  پر \LR{$\pheight{\pi}$} استقرا کوو۔ که \LR{$\pheight{\pi} = 0$} وي، نو ثبوت يوازې يو بديهي اصل دے۔ بيا \LR{$\Subst{\pi}{t}{c}$} هم بديهي اصل دے، ځکه په دې کښې هر فارمول~\LR{$!A(c)$} په \LR{$!A(t)$} بدلېږي۔
\olpseganchor{OLP-0703-B009}{end}

\olpseganchor{OLP-0703-B010}{start}
  که \LR{$\pheight{\pi} > 0$} وي، د~\LR{$\pi$} د وروستي استنتاج له مخې صورتونه بېلوو۔ د جوړښتي قواعدو صورتونه (په \Log{G1c} کښې) او د قضيتي تړونکو د منطقي قواعدو صورتونه اسانه دي او تمرين ته پرېښودل کېږي۔ هغه صورتونه څېړو چې \LR{$\pi$} په يوه کمي استنتاج ختمېږي۔
  \begin{enumerate}
    \item که وروستے استنتاج \LR{$\RightR{\lforall}$} وي، نو \LR{$\pi$} داسې بڼه لري:
    \[
    \AxiomC{}
    \RightLabel{\LR{$\pi'$}}
    \Deduce$\Gamma(c) \fCenter \Delta'(c), !A(a,c)$
    \RightLabel{\RightR{\lforall}}
    \UnaryInf$\Gamma(c) \fCenter \Delta'(c), \lforall[x][!A(x,c)]$
    \DisplayProof
    \]
    د استقرا د فرضيې له مخې \LR{$\Subst{\pi'}{t}{c}$} د \LR{$\Gamma(t) \fCenter \Delta'(t), !A(a,t)$} منظم ثبوت دے۔ څرنګه چې \LR{$a$} د~\LR{$\pi$} ځانګړے متغير دے، نو \LR{$a$} په~\LR{$t$} کښې نۀ راځي۔ له دې امله په \LR{$\Subst{\pi}{t}{c}$} کښې وروستے استنتاج
    \[
    \AxiomC{}
    \RightLabel{\LR{$\Subst{\pi'}{t}{c}$}}
    \Deduce$\Gamma(t) \fCenter \Delta'(t), !A(a,t)$
    \RightLabel{\RightR{\lforall}}
    \UnaryInf$\Gamma(t) \fCenter \Delta'(t), \lforall[x][!A(x,t)]$
    \DisplayProof
    \]
    صحيح دے: په پايله کښې~\LR{$a$} نشته۔
    \item که وروستے استنتاج په \Log{G3c} کښې \LR{$\LeftR{\lforall}$} وي، نو \LR{$\pi$} داسې بڼه لري:
    \[
    \AxiomC{}
    \RightLabel{\LR{$\pi'$}}
    \Deduce$!A(s(c),c), \lforall[x][!A(x,c)], \Gamma(c) \fCenter \Delta'(c)$
    \RightLabel{\LeftR{\lforall}}
    \UnaryInf$\lforall[x][!A(x,c)], \Gamma(c) \fCenter \Delta'(c)$
    \DisplayProof
    \]
    د استقرا د فرضيې له مخې \LR{$\Subst{\pi'}{t}{c}$} د \LR{$!A(s(t),t), \lforall[x][!A(x,t)], \Gamma(t) \fCenter
    \Delta'(t)$} منظم ثبوت دے۔ په \LR{$\Subst{\pi}{t}{c}$} کښې وروستے استنتاج
    \[
    \AxiomC{}
    \RightLabel{\LR{$\Subst{\pi'}{t}{c}$}}
    \Deduce$!A(s(t),t), \lforall[x][!A(x,t)], \Gamma(t) \fCenter \Delta'(t)$
    \RightLabel{\LeftR{\lforall}}
    \UnaryInf$\lforall[x][!A(x,t)], \Gamma(t) \fCenter \Delta'(t)$
    \DisplayProof
    \]
    صحيح دے۔ په~\Log{G1c} کښې د \LeftR{\lforall} حالت هم ورته، خو لږ اسانه دے۔
    \item که وروستے استنتاج \RightR{\lexists} يا \LeftR{\lexists} وي، د تمرين په توګه ئې وګورئ۔
  \end{enumerate}
  په هر صورت کښې مو د يوه منظم ثبوت په پاې کښې يو سېکوېنټ ورزيات کړے (يا د دوو منظمو ثبوتونه په پاې کښې، که قاعده دوه مقدمې لري)۔ نوے سېکوېنټ د \LR{$\pi$} له وروستي سېکوېنټ څخه يوازې په دې کښې بېل دے چې \LR{$c$} په ترم~\LR{$t$} بدل شوے۔ د \LR{$\pi$} او \LR{$\Subst{\pi}{t}{c}$} ځانګړي متغيرونه هماغه دي۔ څرنګه چې \LR{$t$} د~\LR{$\pi$} هېڅ ځانګړے متغير نۀ لري، نو په نوي وروستي سېکوېنټ کښې د \LR{$\Subst{\pi}{t}{c}$} هېڅ ځانګړے متغير نۀ راځي۔ تعويض د څانګو ځانګړي متغيرونه هم نۀ يو شان کوي۔ له دې امله \LR{$\Subst{\pi}{t}{c}$} منظم دے۔
\end{proof}
\olpseganchor{OLP-0703-B010}{end}

\olpseganchor{OLP-0703-B011}{start}
\begin{prop}
که \LR{$\pi$} په \Log{G1c} يا \Log{G3c} کښې ثبوت وي، نو يو منظم ثبوت~\LR{$\pi'$} شته چې هماغه جوړښت، په تېره بيا هماغه لوړوالے، لري او له \LR{$\pi$} څخه يوازې د ځانګړو متغيرونو د نومونو په بدلولو بېلېږي۔
\end{prop}
\olpseganchor{OLP-0703-B011}{end}

\olpseganchor{OLP-0703-B012}{start}
\begin{proof}
يوازې د دې دليل دپاره د~\LR{$\pi$} د ځانګړي متغير يو استنتاج هغه وخت \emph{پاک} بولو چې ځانګړے متغير ئې د بل استنتاج ځانګړے متغير نۀ وي، او په~\LR{$\pi$} کښې د همدغه استنتاج له نېغ فرعي ثبوت پرته بل ځاے رانۀ شي۔ که داسې نۀ وي، \emph{ناپاک} ئې بولو۔ د \LR{$n$} په وسيله په~\LR{$\pi$} کښې د ناپاکو ځانګړي‌متغيري استنتاجونو شمېر ښيو۔ پر~\LR{$n$} په قوي استقرا ښيو چې \LR{$\pi$} په غوښتلي منظم ثبوت~\LR{$\pi'$} بدلېدے شي۔ که \LR{$n=0$} وي، په~\LR{$\pi$} کښې ټول د ځانګړي متغير استنتاجونه پاک دي؛ نو \LR{$\pi$} منظم دے او بل کار نشته۔
که نۀ وي، داسې ناپاک استنتاج وټاکئ چې د پاسه ئې بل ناپاک استنتاج نۀ وي۔ د بېلګې په توګه د \RightR{\lforall} حالت ليکو؛ د کيڼ وجودي کميت صورت هم په همدې ډول دے۔ د~\LR{$\pi$} هغه فرعي ثبوت~\LR{$\pi_1$} چې په دې استنتاج ختمېږي، داسې بڼه لري:
    \[
    \AxiomC{}
    \RightLabel{\LR{$\pi_2$}}
    \Deduce$\Gamma \fCenter \Delta, !A(c)$
    \RightLabel{\RightR{\lforall}}
    \UnaryInf$\Gamma \fCenter \Delta, \lforall[x][!A(x)]$
    \DisplayProof
    \]
څرنګه چې \LR{$c$} ځانګړے متغير دے، په~\LR{$\Gamma$}، \LR{$\Delta$} يا \LR{$\lforall[x][!A(x)]$} کښې نۀ راځي۔ فرعي ثبوت \LR{$\pi_2$} منظم دے، ځکه د ټاکلي استنتاج د پاسه ټول ځانګړي‌متغيري استنتاجونه پاک دي۔ د ټاکلي استنتاج ځانګړے متغير په دغه فرعي ثبوت کښې د بل استنتاج ځانګړے متغير نۀ دے؛ که داسې وے، هغه بل استنتاج به هم ناپاک و۔ داسې ثابت سمبول~\LR{$a$} واخلو چې په~\LR{$\pi$} کښې نۀ راځي۔ د \olref{lem:subst-regular} له مخې \LR{$\Subst{\pi_2}{a}{c}$} د \LR{$\Gamma \fCenter \Delta,
!A(a)$} منظم ثبوت دے چې پرې~\LR{$\RightR{\lforall}$} تطبيقولے شو۔ که \LR{$\pi_1$} په~\LR{$\pi$} کښې په ثبوت \LR{$\pi_1'$} بدل کړو،
    \[
    \AxiomC{}
    \RightLabel{\LR{$\Subst{\pi_2}{a}{c}$}}
    \Deduce$\Gamma \fCenter \Delta, !A(a)$
    \RightLabel{\RightR{\lforall}}
    \UnaryInf$\Gamma \fCenter \Delta, \lforall[x][!A(x)]$
    \DisplayProof
    \]
نو يو ناپاک استنتاج په پاک استنتاج بدلېږي۔ يوازې د ځانګړي متغير \LR{$c$} نوم په~\LR{$a$} بدلوو؛ وروستے سېکوېنټ، جوړښت او لوړوالے هماغه دي۔ هېڅ پخوا پاک استنتاج ناپاکېږي نه، او ښايي د زاړه نوم ځينې نور استنتاجونه هم پاک شي۔ نو په نوي ثبوت کښې د ناپاکو ځانګړي‌متغيري استنتاجونو شمېر~\LR{$< n$} دے۔ پايله د استقرا له فرضيې څخه ترلاسه کېږي۔
\end{proof}
\olpseganchor{OLP-0703-B012}{end}

\olpseganchor{OLP-0703-B013}{start}
د پورته قضيې څرګنده حواله به بيا نۀ ورکوو، خو فرضوو به چې ټول تر بحث لاندې ثبوتونه منظم دي۔ که منظم نۀ وي، د ځانګړو متغيرونو د نومونو په بدلولو ئې تل منظمولے شو۔
\olpseganchor{OLP-0703-B013}{end}

\olpseganchor{OLP-0698-B004}{start}
\olfileid{pt}{seq}{adm}
\olpseganchor{OLP-0698-B004}{end}

\olpseganchor{OLP-0698-B005}{start}
\olsection{مجازې او د اشتقاق وړ قاعدې}
\olpseganchor{OLP-0698-B005}{end}

\olpseganchor{OLP-0698-B006}{start}
د ثبوت د تيورۍ ډېرې پايلې د دې پوښتنې په اړه دي، يا د هغې په اړه
پايلې کاروي، چې آيا هر څۀ چې د يوې ځانګړې قاعدې په کارولو سره
ثابت کېدے شي، له هغې قاعدې پرته هم ثابت کېدے شي۔
د پرېکون د له منځه وړلو قضيه ئې تر ټولو پېژندلے مثال دے۔ دا ښيي
چې هر سېکوېنټ چې د يو سېکوېنټي نظام د قواعدو او~\Cut{} په
کارولو سره ثابت کېدے شي، له~\Cut پرته هم ثابت کېدے شي۔
ځکه چې دا ډول خاصيت ډېر مهم دے، خپل نوم لري۔
\olpseganchor{OLP-0698-B006}{end}

\olpseganchor{OLP-0698-B007}{start}
\begin{defn}
قاعده~\LR{$R$} په يو سېکوېنټي نظام کښې \emph{مجاز} ده، که هر هغه
سېکوېنټ چې په دغه نظام کښې له قاعدې~\LR{$R$} سره ثابتېدے شي،
له~\LR{$R$} پرته هم ثابتېدے شي۔
\end{defn}
\olpseganchor{OLP-0698-B007}{end}

\olpseganchor{OLP-0698-B008}{start}
\begin{ex}\ollabel{ex:land-deriv}
قاعده \LR{$\LeftR{\land}$} په \Log{G1c} او
\Log{G3c} کښې بېلابېلې بڼې لري۔ په \Log{G1c} کښې د دې قاعدې
دوه بڼې دي: په يوه کښې د او تړنې کيڼ غړے \LR{$!A$} د اصلي
فارمول~\LR{$!A \land !B$} په مقدمه کښې دے، او په بله کښې ئې ښے
د او تړنې غړے~\LR{$!B$} دے:
\[
\Axiom$ !A, \Gamma \fCenter \Delta$
\RightLabel{\LR{$\LeftR{\land}_1$}}
\UnaryInf$ !A \land !B, \Gamma \fCenter \Delta$
\DisplayProof
\qquad
\Axiom$!B, \Gamma \fCenter \Delta$
\RightLabel{\LR{$\LeftR{\land}_2$}}
\UnaryInf$!A \land !B, \Gamma \fCenter \Delta$
\DisplayProof
\]
په مقابل کښې، \Log{G3c} يوازې يوه قاعده لري:
\[
\Axiom$ !A, !B, \Gamma \fCenter \Delta$
\RightLabel{\LR{$\LeftR{\land}_3$}}
\UnaryInf$ !A \land !B, \Gamma \fCenter \Delta$
\DisplayProof
\]
اوس دا پوښتنه کولے شو: «آيا د \Log{G3c} قاعده
\LR{$\LeftR{\land}_3$} په \Log{G1c} کښې مجاز ده؟» هو، ده:
هر هغه استنتاج چې \LR{$\LeftR{\land}_3$} کاروي، يوازې د~\Log{G1c}
د قواعدو په کارولو سره نېغ په نېغه تمثيلولے شو۔ د
\LR{$\LeftR{\land}_1$} او \LR{$\LeftR{\land}_2$} ترڅنګ به جوړښتي
قاعدې~\LR{$\LeftR{\Contraction}$} ته هم اړتيا ولرو:
\[
\Axiom$!A, !B, \Gamma \fCenter \Delta$
\RightLabel{\LR{$\LeftR{\land}_1$}}
\UnaryInf$!A \land !B, !B, \Gamma \fCenter \Delta$
\RightLabel{\LR{$\LeftR{\land}_2$}}
\UnaryInf$!A \land !B, !A \land !B, \Gamma \fCenter \Delta$
\RightLabel{\LeftR{\Contraction}}
\UnaryInf$!A \land !B, \Gamma \fCenter \Delta$
\DisplayProof
\]
\end{ex}
\olpseganchor{OLP-0698-B008}{end}

\olpseganchor{OLP-0698-B009}{start}
د يوې قاعدې په کارولو سره شوي استنتاجونه د نورو قواعدو په وسيله
تمثيلول د دې ښودلو يوه لاره ده چې قاعده مجاز ده۔ خو دا يوازينۍ
لاره نۀ ده؛ ځينې قاعدې مجازې دي، خو داسې نۀ شي تمثيلېدې۔
هغه قاعدې چې داسې تمثيلېدې شي، \emph{د اشتقاق وړ} بلل کېږي۔
\olpseganchor{OLP-0698-B009}{end}

\olpseganchor{OLP-0698-B010}{start}
\begin{defn}
قاعده~\LR{$R$}،
\[
\AxiomC{\LR{$S_1$}}
\AxiomC{\LR{$\dots$}}
\AxiomC{\LR{$S_n$}}
\RightLabel{\LR{$R$}}
\TrinaryInfC{\LR{$S$}}
\DisplayProof
\]
په يو نظام کښې \emph{د اشتقاق وړ} بلل کېږي، که د سېکوېنټ~\LR{$S$}
يو سکيماتي ثبوت موجود وي چې د نظام قواعد او بديهيات
او د~\LR{$S_1$}، \dots،~\LR{$S_n$} د بڼو اضافي ابتدايي سېکوېنټونه کاروي۔
\end{defn}
\olpseganchor{OLP-0698-B010}{end}

\olpseganchor{OLP-0698-B011}{start}
په \olref{ex:land-deriv} کښې ورکړے شوے ثبوت ښيي چې
\LR{$\LeftR{\land}_3$} په~\Log{G1c} کښې د اشتقاق وړ قاعده ده؛ ځکه
دا ثبوت د \LR{$\LeftR{\land}_3$} د پايلې يو سکيماتي ثبوت دے
چې د \LR{$\LeftR{\land}_3$} له مقدماتو څخه، د ابتدايي سېکوېنټونو په
توګه، يوازې د \Log{G1c} د قواعدو په کارولو جوړ شوے دے۔
(د \Log{G1c} هېڅ بديهيه ئې نۀ ده کارولې۔)
\olpseganchor{OLP-0698-B011}{end}

\olpseganchor{OLP-0698-B012}{start}
\begin{prop}\ollabel{prop:lor-G3-adm}
د~\Log{G3c} قاعدې \LR{$\LeftR{\land}$} او \LR{$\RightR{\lor}$}
په~\Log{G1c} کښې د اشتقاق وړ دي۔
\end{prop}
\olpseganchor{OLP-0698-B012}{end}

\olpseganchor{OLP-0698-B013}{start}
\begin{proof}
  د \LR{$\LeftR{\land}$} ثبوت په \olref{ex:land-deriv} کښې ورکړے شوے دے۔
  د \LR{$\RightR{\lor}$} ثبوت د تمرين په توګه پرېښودل شوے دے۔
\end{proof}
\olpseganchor{OLP-0698-B013}{end}

\olpseganchor{OLP-0698-B014}{start}
\begin{prob}
وښيئ چې د \Log{G3c} قاعده \LR{$\RightR{\lor}$}
په~\Log{G1c} کښې د اشتقاق وړ ده۔
\end{prob}
\olpseganchor{OLP-0698-B014}{end}

\olpseganchor{OLP-0698-B015}{start}
\begin{prob}
فرض کړئ \LR{$\liff$} يو تعريف شوے منطقي عملګر دے: \LR{$!A \liff !B$}
د \LR{$(!A \lif !B) \land (!B \lif !A)$} لنډه ليکنه ده۔ وښيئ چې
لاندې قاعدې
\begin{enumerate}
\item \Axiom$\Gamma, !A, !B \fCenter \Delta$
\Axiom$\Gamma \fCenter \Delta, !A, !B$
\RightLabel{\LeftR{\liff}}
\BinaryInf$!A \liff !B, \Gamma \fCenter \Delta$
\DisplayProof
\item
\Axiom$!A, \Gamma \fCenter \Delta, !B$
\Axiom$!B, \Gamma \fCenter \Delta, !A$
\RightLabel{\RightR{\liff}}
\BinaryInf$\Gamma \fCenter \Delta, !A \liff !B$
\DisplayProof
\end{enumerate}
په (a)~\Log{G1c} او (b)~\Log{G3c} کښې د اشتقاق وړ دي۔
\end{prob}
\olpseganchor{OLP-0698-B015}{end}

\olpseganchor{OLP-0698-B016}{start}
\begin{ex}\ollabel{ex:weak-adm}
د \Log{G1c} قاعده \LR{$\RightR{\Weakening}$}،
\[
\Axiom$ \Gamma \fCenter \Delta$
\RightLabel{\RightR{\Weakening}}
\UnaryInf$ \Gamma \fCenter \Delta, !A$
\DisplayProof
\]
په~\Log{G3c} کښې مجاز ده۔ د ثبوت نظر ساده دے: فرض کړئ يو
ثبوت~\LR{$\pi$} د مقدمې \LR{$\Gamma \Sequent \Delta$} دپاره په~\Log{G3c}
کښې موجود دے۔ لومړے د ټول
ثبوت ځانګړي متغيرونه داسې تازه نومونو ته واړوئ چې
په ورزياتېدونکي فورمول کښې نۀ راځي۔ بيا سياق پراخ
کړئ؛ ګنې د کمي قواعدو د ځانګړي متغير شرط ماتېدلے شي۔
دا نوم بدلول د ثبوت لوړوالے نۀ بدلوي۔
فارمول~\LR{$!A$} په~\LR{$\pi$} کښې د هر سېکوېنټ ښۍ
خوا ته ورزيات کړئ۔
بديهيات بديهيات پاتې کېږي، او سم استنتاجونه لا هم سم استنتاجونه
دي۔ د نوي ثبوت وروستنے سېکوېنټ، چې اصل ئې وروستۍ فارمول بولي، دا دے: \LR{$\Gamma \Sequent
\Delta, !A$}۔
\olpseganchor{OLP-0698-B016}{end}

\olpseganchor{OLP-0698-B017}{start}
خو قاعده \LR{$\RightR{\Weakening}$} په~\Log{G3c} کښې د اشتقاق وړ
نۀ ده۔ ځکه فرض کړئ چې د اشتقاق وړ وه؛ يعنې فرض کړئ د
\LR{$\Gamma \Sequent
\Delta, !A$} يو ثبوت موندل کېدے شو چې يوازې د~\Log{G3c}
بديهيات او قواعد او، ښايي، اضافي ابتدايي سېکوېنټ~\LR{$\Gamma \Sequent \Delta$}
کاروي۔ که دا ثبوت په بشپړه توګه سکيماتي وي، لکه د
\LR{$\RightR{\Weakening}$} د اشتقاق وړ کېدو دپاره چې اړين دي، نو دا په
يو ثبوت بدلولے شو چې د~\LR{$\Sequent !A$} دپاره وي، که
\LR{$\Gamma$} او~\LR{$\Delta$} لرې کړو۔ په دې سره اضافي ابتدايي
سېکوېنټونه~\LR{$\Gamma \Sequent \Delta$} په تش
سېکوېنټ~\LR{$\quad \Sequent \quad$} بدلېږي۔ ځکه چې تش سېکوېنټ هېڅ
فارمولونه نۀ لري، خو د~\Log{G3c} ټول قواعد په خپلو مقدماتو کښې
لږ تر لږه يوه فعاله فارمول غواړي، نو تش سېکوېنټ په دې سکيماتي
ثبوت کښې د هېڅ استنتاج مقدمه نۀ شي کېدے۔ په نتيجه کښې،
ثبوت يوازې هغه وخت سکيماتي کېدے شي چې په رښتيا اضافي
سېکوېنټ \LR{$\Gamma \Sequent \Delta$} د ابتدايي سېکوېنټ په توګه نۀ کاروي۔
په بله وينا، دا به د \LR{$\Gamma \Sequent
\Delta, !A$} يو سکيماتي اشتقاق وي چې يوازې د~\Log{G3c} بديهيات
او قواعد کاروي؛ او موږ به ښودلې وي چې د دې بڼې هر سېکوېنټ
په~\Log{G3c} کښې يو ثبوت لري۔ خو دا ناشونې ده، ځکه \Log{G3c}
سم دے او نو يوازې هغه سېکوېنټونه ثابت کولے شي چې په هر
جوړښت کښې رښتيني وي۔ که \LR{$\Gamma = \Delta
= \emptyset$} او \LR{$!A \ident \lfalse$} واخلو، \Log{G3c} به د بېلګې
په توګه سېکوېنټ \LR{$\quad \Sequent \lfalse$} ثابت کولے شي، حال دا
چې داسې نۀ شي کولے۔
\end{ex}
\olpseganchor{OLP-0698-B017}{end}

\olpseganchor{OLP-0698-B018}{start}
راځئ د دې پايلې يو ښکلے استقرايي ثبوت ورکړو چې کمزوري کول
په~\Log{G3c} کښې يوه مجازه قاعده ده۔ په حقيقت کښې، لکه شهودي دليل
چې ښيي، يو څۀ قوي پايله هم سمه ده: که \LR{$!A$} د~\LR{$\pi$} د هر
سېکوېنټ ښۍ خوا ته ورزيات کړو، نو د ثبوت جوړښت نۀ بدلوو؛
په ځانګړې توګه، د ثبوت لوړوالے نۀ زياتوو۔ ځکه چې دا مهمه ده،
خپل تعريف غواړي۔
\olpseganchor{OLP-0698-B018}{end}

\olpseganchor{OLP-0698-B019}{start}
\begin{defn}
قاعده~\LR{$R$}،
\[
\AxiomC{\LR{$S_1$}}
\AxiomC{\LR{$\dots$}}
\AxiomC{\LR{$S_n$}}
\RightLabel{\LR{$R$}}
\TrinaryInfC{\LR{$S$}}
\DisplayProof
\]
په يو نظام کښې \emph{لوړوالے ساتونکې مجازه} (يا
\emph{د لوړوالي ساتونکې-مجازه}) ده، که هر کله \LR{$\Proves[h_i] S_i$} د
\LR{$i = 1$}، \dots،~\LR{$n$} دپاره وي، نو \LR{$\Proves_m S$} هم وي، چې
\LR{$m = \max(h_1, \dots,
h_n)$} دے۔
\end{defn}
\olpseganchor{OLP-0698-B019}{end}

\olpseganchor{OLP-0698-B020}{start}
په ياد ولرئ چې د يوې قاعدې د لوړوالے ساتونکې مجازې کېدو دپاره
دومره بس دي چې، هر کله مقدمې~\LR{$S_i$} ثبوتونه~\LR{$\pi_i$} ولري
چې لوړوالے~\LR{$h_i = \pheight{\pi_i}$} لري،
د پايلې يو ثبوت~\LR{$\pi$} موجود وي چې لوړوالے~\LR{$\pheight{\pi}$}
ئې د~\LR{$h_i$} له اعظمي ارزښت څخه زيات نۀ وي۔ په ځانګړې توګه،
که يوازې يوه مقدمه~\LR{$S_1$} وي، نو دومره بس دي چې
\LR{$\pheight{\pi} \le \pheight{\pi_1}$} وي۔ د \LR{$\RightR{\Weakening}$}
او \LR{$\LeftR{\Weakening}$} په صورت کښې په رښتيا
\LR{$\pheight{\pi} = \pheight{\pi_1}$} لرو، خو دا اړينه نۀ ده۔
\olpseganchor{OLP-0698-B020}{end}

\olpseganchor{OLP-0698-B021}{start}
\begin{prop}\ollabel{prop:weak-G3c-adm}
د \Log{G1c} قاعدې \LR{$\RightR{\Weakening}$} او \LR{$\LeftR{\Weakening}$}
په~\Log{G3c} کښې لوړوالے ساتونکې مجازې دي۔
\end{prop}
\olpseganchor{OLP-0698-B021}{end}

\olpseganchor{OLP-0698-B022}{start}
\begin{proof}
  بايد وښيو چې که \LR{$\Log{G3c} \Proves \Gamma \Sequent \Delta$}
  له يو ثبوت~\LR{$\pi$} سره وي چې لوړوالے~\LR{$\pheight{\pi} = n$} لري،
  نو يو \Log{G3c}-ثبوت \LR{$\pi'$} د \LR{$\Gamma \Sequent \Delta, !A$} دپاره
  موجود دے چې \LR{$\pheight{\pi'} = n$} وي۔ لومړے د ثبوت ځانګړي متغيرونه داسې تازه کوو چې په ورزياتېدونکي فورمول کښې نۀ راځي۔ دا پر~\LR{$n$} د استقرا په
  وسيله ثابتوو۔ که \LR{$n=0$} وي، نو \LR{$\pi$} يوازې بديهيه
  \LR{$\Gamma \Sequent \Delta$} ده (يعنې يو اتومي \LR{$!C$} په دواړو
  \LR{$\Gamma$} او~\LR{$\Delta$} کښې راځي، يا \LR{$\lfalse$} په \LR{$\Gamma$} کښې راځي)، او \LR{$\Gamma
  \Sequent \Delta, !A$} هم بديهيه ده۔ که \LR{$n = \pheight{\pi} >0$}
  وي، نو ثبوت~\LR{$\pi$} يو وروستے استنتاج لري۔ د کارول شوې
  قاعدې له مخې بېلابېل صورتونه جلا کوو۔ يوازې يوه بېلګه ورکوو:
  فرض کړئ وروستے استنتاج \LR{$\RightR{\land}$} و۔ نو
  \LR{$\Delta = \Delta', !B \land !C$} دے، او د وروستي استنتاج
  نېغې فرعي ثبوتونه \LR{$\pi_1$} او~\LR{$\pi_2$} په ترتيب سره وروستي
  سېکوېنټونه \LR{$\Gamma \Sequent \Delta', !B$} او \LR{$\Gamma \Sequent
  \Delta', !C$} لري۔ په بله وينا، \LR{$\pi$} داسې ښکاري:
  \[
  \AxiomC{}
  \RightLabel{\LR{$\pi_1$}}
  \Deduce$\Gamma \fCenter \Delta', !B$
  \AxiomC{}
  \RightLabel{\LR{$\pi_2$}}
  \Deduce$\Gamma \fCenter \Delta', !C$
  \RightLabel{\RightR{\land}}
  \BinaryInf$\Gamma \fCenter \Delta', !B \land !C$
  \DisplayProof
  \]
  همداراز \LR{$n =
  \max(\pheight{\pi_1}, \pheight{\pi_2}) + 1$} لرو۔ نو \LR{$\pheight{\pi_1} <
  n$} او \LR{$\pheight{\pi_2} < n$} دي، او د استقرا فرضيه تطبيقېږي:
  ثبوتونه \LR{$\pi_1'$} او~\LR{$\pi_2'$} د \LR{$\Gamma \Sequent
  \Delta', !B, !A$} او \LR{$\Gamma \Sequent \Delta', !C, !A$} دپاره موجود دي۔
  د \LR{$\RightR{\land}$} قاعدې په کارولو سره يو ثبوت~\LR{$\pi'$} د
  \LR{$\Gamma
  \Sequent \Delta', !B \land !C, !A$} دپاره ترلاسه کوو:
  \[
  \AxiomC{}
  \RightLabel{\LR{$\pi_1'$}}
  \Deduce$\Gamma \fCenter \Delta', !B, !A$
  \AxiomC{}
  \RightLabel{\LR{$\pi_2'$}}
  \Deduce$\Gamma \fCenter \Delta', !C, !A$
  \RightLabel{\RightR{\land}}
  \BinaryInf$\Gamma \fCenter \Delta', !B \land !C, !A$
  \DisplayProof
  \]
  موږ \LR{$\pheight{\pi'} =
  \max(\pheight{\pi_1'}, \pheight{\pi_2'}) + 1 = \max(\pheight{\pi_1},
  \pheight{\pi_2}) + 1 = \pheight{\pi}$} لرو۔
\end{proof}
\olpseganchor{OLP-0698-B022}{end}

\olpseganchor{OLP-0698-B023}{start}
\olref{prop:weak-G3c-adm} ډېر ګټور دے۔ د دې د تکراري تطبيق دپاره
يوه ليکنښه معرفي کوو: که \LR{$\pi$} د~\LR{$\Gamma \Sequent \Delta$}
يو ثبوت وي، نو \LR{$\pi[\Pi \Sequent \Lambda]$} هغه پايله ده
چې د~\LR{$\Pi$} له ټولو فارمولونه سره په کيڼ اړخ کښې د کمزوري کولو
مجازتيا او د~\LR{$\Lambda$} له ټولو فارمولونه سره په ښي اړخ کښې
د کمزوري کولو مجازتيا تطبيق کړو۔ دا د \LR{$\Gamma, \Pi
\Sequent \Delta, \Lambda$} يو ثبوت دے۔
\olpseganchor{OLP-0698-B023}{end}

\olpseganchor{OLP-0701-B004}{start}
\olfileid{pt}{seq}{inv}
\olpseganchor{OLP-0701-B004}{end}

\olpseganchor{OLP-0701-B005}{start}
\olsection{د قواعدو د معکوسولو وړتيا}
\olpseganchor{OLP-0701-B005}{end}

\olpseganchor{OLP-0701-B006}{start}
\begin{defn}
قاعده~\LR{$R$}،
\[
\AxiomC{\LR{$S_1$}}
\AxiomC{\LR{$\dots$}}
\AxiomC{\LR{$S_n$}}
\RightLabel{\LR{$R$}}
\TrinaryInfC{\LR{$S$}}
\DisplayProof
\]
په يو نظام کښې \emph{د معکوسولو وړ} ده، که هر کله \LR{$\Proves S$}
وي، نو \LR{$\Proves[h_i] S_i$} د \LR{$i = 1$}، \dots،~\LR{$n$} دپاره وي۔
دا په يو نظام کښې \emph{لوړوالے ساتونکې او د معکوسولو وړ}
(يا \emph{د لوړوالي ساتونکې-د معکوسولو وړ}) ده، که هر کله \LR{$\Proves_n S$}
وي، نو \LR{$\Proves[n] S_i$} د \LR{$i =
1$}، \dots،~\LR{$n$} دپاره وي۔
\end{defn}
\olpseganchor{OLP-0701-B006}{end}

\olpseganchor{OLP-0701-B007}{start}
\begin{lem}\ollabel{lem:G3c-invert}
د~\Log{G3c} قاعدې د \LR{$\lnot$}، \LR{$\land$}، \LR{$\lor$} او \LR{$\lif$} دپاره
لوړوالے ساتونکې او د معکوسولو وړ دي، يعنې:
\begin{enumerate}
  \item \RightR{\lnot}: که \LR{$\Log{G3c} \Proves[n] \Gamma \Sequent
  \Delta, \lnot !A$} وي، نو \LR{$\Log{G3c} \Proves[n] !A, \Gamma \Sequent
  \Delta$} دے۔
  \item \LeftR{\lnot}: که \LR{$\Log{G3c} \Proves[n] \Gamma, \lnot !A
  \Sequent \Delta$} وي، نو \LR{$\Log{G3c} \Proves[n] \Gamma \Sequent \Delta,
  !A$} دے۔
  \item \RightR{\land}: که \LR{$\Log{G3c} \Proves[n] \Gamma \Sequent
  \Delta, !A \land !B$} وي، نو \LR{$\Log{G3c} \Proves[n] \Gamma \Sequent
  \Delta, !A$} او \LR{$\Log{G3c} \Proves[n] \Gamma \Sequent
  \Delta, !B$} دي۔
  \item \LeftR{\land}: که \LR{$\Log{G3c} \Proves[n] \Gamma, !A \land !B
  \Sequent \Delta$} وي، نو \LR{$\Log{G3c} \Proves[n] !A, !B, \Gamma \Sequent
  \Delta$} دے۔
  \item \RightR{\lor}: که \LR{$\Log{G3c} \Proves[n] \Gamma \Sequent
  \Delta, !A \lor !B$} وي، نو \LR{$\Log{G3c} \Proves[n] \Gamma \Sequent
  \Delta, !A, !B$} دے۔
  \item \LeftR{\lor}: که \LR{$\Log{G3c} \Proves[n] !A \lor !B, \Gamma \Sequent
  \Delta$} وي، نو \LR{$\Log{G3c} \Proves[n] !A, \Gamma \Sequent
  \Delta$} او \LR{$\Log{G3c} \Proves[n] !B, \Gamma \Sequent
  \Delta$} دي۔
  \item \RightR{\lif}: که \LR{$\Log{G3c} \Proves[n] \Gamma \Sequent
  \Delta, !A \lif !B$} وي، نو \LR{$\Log{G3c} \Proves[n] !A, \Gamma \Sequent
  \Delta, !B$} دے۔
  \item \LeftR{\lif}: که \LR{$\Log{G3c} \Proves[n] !A \lif !B, \Gamma
  \Sequent \Delta$} وي، نو \LR{$\Log{G3c} \Proves[n] \Gamma \Sequent \Delta,
  !A$} او \LR{$\Log{G3c} \Proves[n] !B, \Gamma \Sequent \Delta$} دي۔
\end{enumerate}
\end{lem}
\olpseganchor{OLP-0701-B007}{end}

\olpseganchor{OLP-0701-B008}{start}
\begin{proof}
بايد وښيو چې د هرې قاعدې~\LR{$R$} د~\Log{G3c} دپاره، که
ثبوت~\LR{$\pi$} د پايلې~\LR{$S$} دپاره وي چې د~\LR{$R$} ده، نو
ثبوتونه~\LR{$\pi_i$} د مقدماتو \LR{$S_i$} دپاره هم شته چې د~\LR{$R$} دي،
او \LR{$\pheight{\pi_i}
\le \pheight{\pi}$} وي۔ قاعده \LeftR{\land} وګورئ: فرض کړئ
ثبوت~\LR{$\pi$} شته چې وروستے سېکوېنټ ئې د \LR{$!A \land !B, \Gamma
\Sequent \Delta$} په بڼه دے۔ بايد وښيو چې ثبوت \LR{$\pi'$}
د \LR{$!A, !B, \Gamma \Sequent \Delta$} دپاره شته چې \LR{$\pheight{\pi'} \le
\pheight{\pi}$} وي۔ دا پر \LR{$n = \pheight{\pi}$} د استقرا په وسيله کوو۔
\olpseganchor{OLP-0701-B008}{end}

\olpseganchor{OLP-0701-B009}{start}
که \LR{$n = 0$} وي، نو \LR{$\pi$} يوازې يو بديهي اصل لري: \LR{$!A \land !B, \Gamma
\Sequent \Delta$}۔ په \Log{G3c} کښې د پېژندنې بديهي اصول د \LR{$!C, \Gamma
\Sequent \Delta, !C$} په بڼه دي چې \LR{$!C$} اتومي وي۔ په بله وينا،
\LR{$!C$}، \LR{$!A \land !B$} نۀ شي کېدے۔ که \LR{$!A \land !B, \Gamma \Sequent \Delta$}
د پېژندنې بديهي اصل وي، نو يو اتومي \LR{$!C$} بايد د دواړو \LR{$\Gamma$}
او~\LR{$\Delta$} غړے وي۔ خو بيا \LR{$!A, !B, \Gamma \Sequent \Delta$}
هم بديهي اصل دے، او خپل ثبوت~\LR{$\pi'$} مو وموند
(چې لوړوالے ئې هم~\LR{$0$} دے)۔
\textbf{د بديهي حالت دويمه بڼه:} که \LR{$\lfalse$} په \LR{$\Gamma$} کښې وي، نو د کيڼې تشې نښې بديهي اصل هم شته؛ د تړنې د دواړو غړو تر بېلولو وروسته هم دغه بديهي اصل پاتې کېږي او لوړوالے ئې صفر دے۔
\olpseganchor{OLP-0701-B009}{end}

\olpseganchor{OLP-0701-B010}{start}
که \LR{$n>0$} وي، په يو استنتاج پاے ته رسېږي۔ بايد د~\LR{$n$} دپاره
دا ادعا ثابته کړو، خو فرض کولے شو چې د هر هغه ثبوت
دپاره سمه ده چې د ثبوت لوړوالے ئې~\LR{$<n$} وي، که څۀ هم
چاپېرچل ئې بېل وي۔ په بله وينا، د هر \LR{$k<n$} او هر
\LR{$\Pi$} او \LR{$\Lambda$} دپاره، که \LR{$!A \land !B, \Pi \Sequent \Lambda$}
ثبوت ولري چې لوړوالے ئې~\LR{$k$} وي، نو ثبوت شته
چې لوړوالے ئې~\LR{$\le k$} وي او د \LR{$!A, !B, \Pi \Sequent \Lambda$} دپاره وي۔
\olpseganchor{OLP-0701-B010}{end}

\olpseganchor{OLP-0701-B011}{start}
دوه صورتونه لرو: يا په کيڼه خوا کښې \LR{$!A \land
!B$} اصلي دے، يا نۀ دے۔ که اصلي وي، نو وروستے استنتاج
\LeftR{\land} دے او نېغ فرعي ثبوت~\LR{$\pi_1$} په
\LR{$!A, !B, \Gamma \Sequent \Delta$} پاے ته رسېږي۔ په دې صورت
کښې پايله سمه ده، ځکه \LR{$\pheight{\pi_1} < \pheight{\pi}$} دے۔
\olpseganchor{OLP-0701-B011}{end}

\olpseganchor{OLP-0701-B012}{start}
که \LR{$!A \land !B$} اصلي نۀ وي، نو بله فارمول اصلي ده، او \LR{$!A
\land !B$} د وروستي استنتاج په چاپېرچل کښې دے۔ بيا د استقرا
فرضيه د وروستۍ قاعدې~\LR{$R$} پر نېغو فرعي ثبوتونه تطبيقېږي؛
له دې څخه ثبوت~\LR{$\pi_1'$}، يا دوه ثبوتونه~\LR{$\pi_1'$} او~\LR{$\pi_2'$}
ترلاسه کېږي چې لوړوالے ئې د~\LR{$\pi$} د نېغو فرعي ثبوتونه
د ثبوت له لوړوالي څخه زيات نۀ وي۔ د \LR{$\pi_1'$} او \LR{$\pi_2'$}
وروستي سېکوېنټونه په کيڼ چاپېرچل کښې \LR{$!A, !B$} د
\LR{$!A \land !B$} پر ځاے لري۔ قاعده~\LR{$R$} تطبيق کړئ، څو ثبوت~\LR{$\pi'$} ترلاسه کړئ۔
د بېلګې په توګه، فرض کړئ \LR{$\pi$} په~\LeftR{\lif} پاے ته رسېږي:
\[
\AxiomC{}
\RightLabel{\LR{$\pi_1$}}
\Deduce$!A \land !B, \Gamma' \fCenter \Delta, !C$
\AxiomC{}
\RightLabel{\LR{$\pi_2$}}
\Deduce$!A \land !B, \Gamma', !D \fCenter \Delta$
\RightLabel{\LeftR{\lif}}
\BinaryInf$!A \land !B, \Gamma', !C \lif !D \fCenter \Delta$
\DisplayProof
\]
دلته په کيڼه خوا کښې \LR{$!C \lif !D$} اصلي دے، او کيڼ چاپېرچل
\LR{$!A \land !B, \Gamma'$} دے (يعنې \LR{$\Gamma = \Gamma', !C \lif !D$})۔
د استقرا فرضيه ثبوتونه \LR{$\pi_1'$} او \LR{$\pi_2'$} په ترتيب سره د
\LR{$!A, !B, \Gamma' \Sequent \Delta, !C$} او \LR{$!A, !B, \Gamma',
!D \Sequent \Delta$} دپاره راکوي۔
دا دوه سېکوېنټونه بيا د \LeftR{\lif} قاعدې د مقدماتو په
توګه کارېدے شي، څو~\LR{$\pi'$} ترلاسه کړو:
\[
\AxiomC{}
\RightLabel{\LR{$\pi_1'$}}
\Deduce$!A, !B, \Gamma' \fCenter \Delta, !C$
\AxiomC{}
\RightLabel{\LR{$\pi_2'$}}
\Deduce$!A, !B, \Gamma', !D \fCenter \Delta$
\RightLabel{\LeftR{\lif}}
\BinaryInf$!A, !B, \Gamma', !C \lif !D \fCenter \Delta$
\DisplayProof
\]
لرو:
\begin{align*}
\pheight{\pi} & = \max(\pheight{\pi_1}, \pheight{\pi_2}) + 1\\
\pheight{\pi'} & = \max(\pheight{\pi_1'}, \pheight{\pi_2'}) + 1
\intertext{\RL{د تعريف له مخې، او}}
\pheight{\pi_1'} & \le \pheight{\pi_1}\\
\pheight{\pi_2'} & \le \pheight{\pi_2}
\intertext{\RL{د استقرا د فرضيې له مخې۔ نو}}
\pheight{\pi'} & \le \pheight{\pi}.
\end{align*}
\end{proof}
\olpseganchor{OLP-0701-B012}{end}

\olpseganchor{OLP-0701-B013}{start}
\begin{prob}
د \olref{lem:G3c-invert} ثبوت د \RightR{\lif} دپاره ورکړئ۔
\end{prob}
\olpseganchor{OLP-0701-B013}{end}

\olpseganchor{OLP-0701-B014}{start}
\begin{lem}\ollabel{lem:invert-quant}
قاعدې \RightR{\lforall} او \LeftR{\lexists}
په لاندې معنا لوړوالے ساتونکې او د معکوسولو وړ دي:
\begin{enumerate}
  \item که \LR{$\Log{G3c} \Proves[n] \Gamma \Sequent \Delta,
  \lforall[x][!A(x)]$} وي، نو \LR{$\Log{G3c} \Proves[n] \Gamma \Sequent
  \Delta, !A(c)$} دے، او
  \item که \LR{$\Log{G3c} \Proves[n] \lexists[x][!A(x)], \Gamma \Sequent
  \Delta$} وي، نو \LR{$\Log{G3c} \Proves[n] !A(c), \Gamma \Sequent \Delta$} دے،
\end{enumerate}
د هر هغه \LR{$c$} دپاره چې په \LR{$\Gamma$}، \LR{$\Delta$} او په ترتيب سره
په \LR{$\lforall[x][!A(x)]$} يا \LR{$\lexists[x][!A(x)]$} کښې نۀ راځي۔
\end{lem}
\olpseganchor{OLP-0701-B014}{end}

\olpseganchor{OLP-0701-B015}{start}
\begin{proof}
  (1) ښيو او (2) د تمرين په توګه پرېږدو۔ دليل د
  \olref{lem:G3c-invert} ثبوت ته ورته دے۔ د استقرا له پيل څخه ټول ثبوت منظم فرضوو او د ټاکلي نوي ثابت څخه د ټولو داخلي ځانګړو متغيرونو نومونه تازه کوو۔ دا نوم بدلول وروستے سېکوېنټ او لوړوالے نۀ بدلوي؛ په نورو کمي قواعدو کښې هم د ځانګړي متغير شرط ساتي۔ د استقرا په ګام کښې
  بيا دوه صورتونه دي۔ هغه صورت چې \LR{$\lforall[x][!A(x)]$}
  اصلي وي، بيا ژر حلېږي: د وروستي \LR{$\RightR{\lforall}$} استنتاج
  مقدمه په غوښتل شوې بڼه \LR{$\Gamma
  \Sequent \Delta, !A(a)$} ده، او په هغې پاے ته رسېدونکے نېغ
  فرعي ثبوت~\LR{$\pi_1$} د \LR{$\pheight{\pi_1} < \pheight{\pi}$} شرط
  پوره کوي۔ سربېره پر دې، د ځانګړي متغير شرط راکوي چې \LR{$a$}
  په \LR{$\Gamma$}، \LR{$\Delta$} يا \LR{$\lforall[x][!A(x)]$} کښې نۀ راځي۔
  ځکه فرض کولے شو چې \LR{$\pi_1$} منظم دے، نو
  \LR{$\Subst{\pi_1}{c}{a}$} د \LR{$\Gamma \Sequent \Delta, !A(c)$} يو
  ثبوت دے چې لوړوالے ئې هماغه دے، د
  \olref[qua]{lem:subst-regular} له مخې۔
\olpseganchor{OLP-0701-B015}{end}

\olpseganchor{OLP-0701-B016}{start}
  په دويم صورت کښې، \LR{$\lforall[x][!A(x)]$} په وروستي استنتاج
  کښې اصلي نۀ دے؛ بل فورمول اصلي دے۔ بيا يو صورت د بېلګې
  په توګه څېړو۔ فرض کړئ وروستے استنتاج~\RightR{\land} دے۔
  په \LR{$\Delta$} کښې يوه فارمول د~\LR{$!B \land !C$} په بڼه ده
  او اصلي ده۔ ثبوت~\LR{$\pi$} داسې پاے ته رسېږي:
  \[
\AxiomC{}
\RightLabel{\LR{$\pi_1$}}
\Deduce$\Gamma \fCenter \Delta', \lforall[x][!A(x)], !B$
\AxiomC{}
\RightLabel{\LR{$\pi_2$}}
\Deduce$\Gamma \fCenter \Delta', \lforall[x][!A(x)], !C$
\RightLabel{\RightR{\land}}
\BinaryInf$\Gamma \fCenter \Delta', \lforall[x][!A(x)], !B \land !C$
\DisplayProof
\]
چې \LR{$\Delta = \Delta', !B \land !C$} دے۔ فرض کړئ \LR{$c$} داسې
ثابت سمبول دے چې په \LR{$\Delta$} کښې نۀ ښکاري؛ نو څرګنده ده
چې په \LR{$\Delta', !B$} يا په \LR{$\Delta', !C$} کښې هم نۀ راځي۔
نو د استقرا فرضيه پر \LR{$\pi_1$} او~\LR{$\pi_2$} تطبيقېږي؛
ثبوتونه~\LR{$\pi_1'$} او \LR{$\pi_2'$} لرو چې
\LR{$\pheight{\pi_1'} \le \pheight{\pi_1}$} او \LR{$\pheight{\pi_2'} \le
\pheight{\pi_2}$} وي۔ غوښتل شوے ثبوت~\LR{$\pi'$} د \LR{$\Gamma \Sequent
\Delta, !A(c)$} دپاره دا دے:
\[
\AxiomC{}
\RightLabel{\LR{$\pi_1'$}}
\Deduce$\Gamma \fCenter \Delta', !A(c), !B$
\AxiomC{}
\RightLabel{\LR{$\pi_2'$}}
\Deduce$\Gamma \fCenter \Delta', !A(c), !C$
\RightLabel{\RightR{\land}}
\BinaryInf$\Gamma \fCenter \Delta', !A(c), !B \land !C$
\DisplayProof
\]
چې \LR{$\pheight{\pi'} = \max(\pheight{\pi_1'}, \pheight{\pi_2'})+1 \le \pheight{\pi}$} دے۔
\end{proof}
\olpseganchor{OLP-0701-B016}{end}

\olpseganchor{OLP-0701-B017}{start}
\Olref{lem:G3c-invert} د پرېکون لرې کولو د قضيې په ثبوت
کښې مهم رول لري۔ خو د لاندې خبرې د ښودلو دپاره هم کارېدے شي:
\olpseganchor{OLP-0701-B017}{end}

\olpseganchor{OLP-0701-B018}{start}
\begin{prop}\ollabel{prop:G3c-cont-adm}
د~\Log{G1c} د انقباض قاعدې \LeftR{\Contraction} او \RightR{\Contraction}
په~\Log{G3c} کښې لوړوالے ساتونکې مجازې دي۔
\end{prop}
\olpseganchor{OLP-0701-B018}{end}

\olpseganchor{OLP-0701-B019}{start}
\begin{proof}
د \RightR{\Contraction} او \LeftR{\Contraction} دپاره دليل
يوځاے ورکوو۔ فرض کړئ \LR{$\pi$} د \LR{$\Gamma \Sequent \Delta, !A, !A$}
يا د \LR{$!A, !A, \Gamma \Sequent
\Delta$} ثبوت دے، په~\Log{G3c} کښې۔ بايد وښيو چې يو
\Log{G3c}-ثبوت~\LR{$\pi'$} هم شته چې په ترتيب سره د
\LR{$\Gamma \Sequent \Delta, !A$} يا \LR{$!A,
\Gamma \Sequent \Delta$} دپاره وي، او \LR{$\pheight{\pi'} \le
\pheight{\pi}$} وي۔ دا پر \LR{$n = \pheight{\pi}$} د استقرا په وسيله کوو۔
\olpseganchor{OLP-0701-B019}{end}

\olpseganchor{OLP-0701-B020}{start}
که \LR{$n=0$} وي، نو \LR{$\pi$} يوازې يو بديهي اصل دے۔ خو که
\LR{$\Gamma \Sequent \Delta, !A,
!A$} د~\Log{G3c} بديهي اصل وي، نو \LR{$\Gamma \Sequent \Delta, !A$}
هم بديهي اصل دے۔ د پېژندنې د بديهي اصل په حالت کښې يا يو اتومي~\LR{$!C$} د دواړو \LR{$\Gamma$} او~\LR{$\Delta$}
غړے دے، يا \LR{$!A$} اتومي دے او د~\LR{$\Gamma$} غړے
دے۔ هماغه دليل هغه وخت هم تطبيقېږي چې وروستے سېکوېنټ
\LR{$!A, !A, \Gamma \Sequent
\Delta$} وي۔
\textbf{د تشې نښې حالت:} که \LR{$\lfalse$} په \LR{$\Gamma$} کښې وي، ښے او کيڼ انقباض دواړه دا بديهي اصل ساتي۔ په کيڼ انقباض کښې که \LR{$!A \ident \lfalse$} هم وي، د دوو نقلونو له يوځای کولو وروسته يو نقل پاتې وي؛ پايله بيا هم بديهي اصل دے۔
\olpseganchor{OLP-0701-B020}{end}

\olpseganchor{OLP-0701-B021}{start}
که \LR{$n>0$} وي، نو \LR{$\pi$} په يوه قاعده~\LR{$R$} پاے ته رسېږي۔ که \LR{$!A$}
په دې وروستي استنتاج کښې اصلي نۀ وي، نو ثبوت~\LR{$\pi'$}
د \olref{lem:G3c-invert} د ثبوت په شان ترلاسه کولے شو: د~\LR{$\pi$}
پر نېغو فرعي ثبوتونه د استقرا فرضيه تطبيقوو، او پر ترلاسه
شويو ثبوت(s) هماغه قاعده~\LR{$R$} بيا تطبيقوو۔ د استقرا فرضيه
دا ده چې که ثبوت شته چې لوړوالے ئې~\LR{$k<n$} وي او يا
د \LR{$\Pi \Sequent \Lambda, !D,
!D$} يا د \LR{$!D, !D, \Pi \Sequent \Lambda$} دپاره وي، نو
ثبوت شته چې لوړوالے ئې~\LR{$\le k$} وي او په ترتيب سره
د \LR{$\Pi \Sequent \Lambda, !D$} يا \LR{$!D, \Pi \Sequent
\Lambda$} دپاره وي۔ (په ياد ولرئ چې د استقرا فرضيه تطبيقېږي،
که چاپېرچلونه يا منقبضه فارمول له \LR{$\Gamma$}، \LR{$\Delta$}
يا~\LR{$!A$} څخه بېل هم وي!)
\olpseganchor{OLP-0701-B021}{end}

\olpseganchor{OLP-0701-B022}{start}
د بېلګې په توګه، فرض کړئ \LR{$R$}، \LR{$\RightR{\land}$} دے،
او \LR{$\pi$}~داسې پاے ته رسېږي:
\[
\AxiomC{}
\RightLabel{\LR{$\pi_1$}}
\Deduce$\Gamma \fCenter \Delta', !B, !A, !A$
\AxiomC{}
\RightLabel{\LR{$\pi_2$}}
\Deduce$\Gamma \fCenter \Delta', !C, !A, !A$
\RightLabel{\RightR{\land}}
\BinaryInf$\Gamma \fCenter \Delta', !B \land !C, !A, !A$
\DisplayProof
\]
چې \LR{$\Delta = \Delta', !B \land !C$} دے۔ د استقرا فرضيه
ثبوتونه \LR{$\pi_1'$} او \LR{$\pi_2'$} په ترتيب سره د \LR{$\Gamma \fCenter \Delta', !B, !A$}
او \LR{$\Gamma \fCenter \Delta', !C, !A$} دپاره راکوي، چې
\LR{$\pheight{\pi_1'} \le  \pheight{\pi_1}$} او \LR{$\pheight{\pi_2'} \le
\pheight{\pi_2}$} وي۔ ثبوت~\LR{$\pi'$} دا دے:
\[
\AxiomC{}
\RightLabel{\LR{$\pi_1'$}}
\Deduce$\Gamma \fCenter \Delta', !B, !A$
\AxiomC{}
\RightLabel{\LR{$\pi_2'$}}
\Deduce$\Gamma \fCenter \Delta', !C, !A$
\RightLabel{\RightR{\land}}
\BinaryInf$\Gamma \fCenter \Delta', !B \land !C, !A$
\DisplayProof
\]
\olpseganchor{OLP-0701-B022}{end}

\olpseganchor{OLP-0701-B023}{start}
که \LR{$!A$} په وروستي استنتاج کښې اصلي وي، نو د معکوسولو
لمّه (\olref{lem:G3c-invert}) د~\LR{$\pi$} پر نېغو فرعي ثبوتونه
تطبيقوو، بيا د استقرا فرضيه، او بيا قاعده~\LR{$R$} تطبيقوو۔
د بېلګې په توګه، فرض کړئ \LR{$\pi$}، \LR{$\Gamma \Sequent \Delta, !A, !A$}
ثابتوي، \LR{$!A
\ident (!B \land !C)$} دے، او په وروستي استنتاج کښې اصلي دے۔
نو \LR{$\pi$} داسې ښکاري:
\[
\AxiomC{}
\RightLabel{\LR{$\pi_1$}}
\Deduce$\Gamma \fCenter \Delta, !B \land !C, !B$
\AxiomC{}
\RightLabel{\LR{$\pi_2$}}
\Deduce$\Gamma \fCenter \Delta, !B \land !C, !C$
\RightLabel{\RightR{\land}}
\BinaryInf$\Gamma \fCenter \Delta, !B \land !C, !B \land !C$
\DisplayProof
\]
\olref{lem:G3c-invert} پر \LR{$\pi_1$} تطبيق کړئ، څو ثبوت~\LR{$\pi_1'$}
د \LR{$\Gamma \Sequent \Delta, !B, !B$} دپاره ترلاسه کړئ۔ (دا
ثبوت د \LR{$\Gamma \Sequent \Delta, !C, !B$} دپاره هم راکوي،
خو هغې ته اړتيا نۀ لرو۔) په همدې ډول، د
\olref{lem:G3c-invert} پر~\LR{$\pi_2$} د تطبيق په وسيله
ثبوت~\LR{$\pi_2'$} د \LR{$\Gamma \Sequent \Delta, !C, !C$} دپاره
ترلاسه کوو۔ ځکه د \RightR{\land} معکوسول لوړوالے ساتي، نو
\LR{$\pheight{\pi_1'} \le  \pheight{\pi_1} < \pheight{\pi}$} او
\LR{$\pheight{\pi_2'} \le  \pheight{\pi_2} < \pheight{\pi}$} دي۔
نو د استقرا فرضيه پر \LR{$\pi_1'$} او \LR{$\pi_2'$} تطبيقېږي:
ثبوتونه~\LR{$\pi_1''$} او \LR{$\pi_2''$} په ترتيب سره د \LR{$\Gamma \Sequent \Delta, !B$}
او \LR{$\Gamma \Sequent \Delta, !C$} دپاره شته، چې \LR{$\pheight{\pi_1''}
\le \pheight{\pi_1'} \le \pheight{\pi_1}$} او \LR{$\pheight{\pi_2''} \le
\pheight{\pi_2'} \le \pheight{\pi_2}$} وي۔ بيا ثبوت~\LR{$\pi'$} دا دے:
\[
\AxiomC{}
\RightLabel{\LR{$\pi_1''$}}
\Deduce$\Gamma \fCenter \Delta, !B$
\AxiomC{}
\RightLabel{\LR{$\pi_2''$}}
\Deduce$\Gamma \fCenter \Delta, !C$
\RightLabel{\RightR{\land}}
\BinaryInf$\Gamma \fCenter \Delta, !B \land !C$
\DisplayProof
\]
چې \LR{$\pheight{\pi'} \le \pheight{\pi}$} دے۔
\olpseganchor{OLP-0701-B023}{end}

\olpseganchor{OLP-0701-B024}{start}
د کميت‌ښودونکو د قواعدو په هغو صورتونو کښې ځانګړے پام
ضروري دے چې فارمول اصلي وي۔
\olpseganchor{OLP-0701-B024}{end}

\olpseganchor{OLP-0701-B025}{start}
فرض کړئ \LR{$!A \ident \lforall[x][!B(x)]$} په ښۍ خوا کښې اصلي
دے۔ نو \LR{$\pi$} داسې پاے ته رسېږي:
\[
\AxiomC{}
\RightLabel{\LR{$\pi_1$}}
\Deduce$\Gamma \fCenter \Delta, \lforall[x][!B(x)], !B(a)$
\RightLabel{\RightR{\lforall}}
\UnaryInf$\Gamma \fCenter \Delta, \lforall[x][!B(x)], \lforall[x][!B(x)]$
\DisplayProof
\]
د \olref{lem:invert-quant} له مخې، ثبوت~\LR{$\pi_1'$} شته
چې لوړوالے ئې \LR{$\le \pheight{\pi_1}$} وي او د \LR{$\Gamma \fCenter \Delta, !B(c),
!B(a)$} دپاره وي، د هر هغه~\LR{$c$} دپاره چې په \LR{$\Gamma \fCenter \Delta,
\lforall[x][!B(x)], !B(a)$} کښې نۀ راځي۔ فرض کولے شو چې
\LR{$\pi_1'$} منظم دے او د بدلولو دواړه ثابتونه پکې د ځانګړو متغيرونو نومونه نۀ دي؛ دا د نومونو په تازه کولو سره بې له دې چې وروستے سېکوېنټ يا لوړوالے بدل شي، حاصلېږي۔ نو د \olref[qua]{lem:subst-regular} له مخې،
ثبوت \LR{$\Subst{\pi_1'}{a}{c}$} د \LR{$\Gamma \fCenter \Delta, !B(a),
!B(a)$} ثبوت دے چې لوړوالے ئې هم \LR{$\le \pheight{\pi_1}$} دے۔
اوس د استقرا فرضيه تطبيقېږي او ثبوت~\LR{$\pi_1''$} د \LR{$\Gamma
\fCenter \Delta, !B(a)$} دپاره راکوي۔ د \RightR{\lforall} قاعده
تطبيق کړئ، څو~\LR{$\pi'$} ترلاسه کړئ:
\[
\AxiomC{}
\RightLabel{\LR{$\pi_1''$}}
\Deduce$\Gamma \fCenter \Delta, !B(a)$
\RightLabel{\RightR{\lforall}}
\UnaryInf$\Gamma \fCenter \Delta, \lforall[x][!B(x)]$
\DisplayProof
\]
چې \LR{$\pheight{\pi'} \le \pheight{\pi}$} دے۔
\olpseganchor{OLP-0701-B025}{end}

\olpseganchor{OLP-0701-B026}{start}
اوس فرض کړئ \LR{$!A \ident \lexists[x][!B(x)]$} دے
او په ښۍ خوا کښې اصلي دے۔ نو \LR{$\pi$} داسې ښکاري:
\[
\AxiomC{}
\RightLabel{\LR{$\pi_1$}}
\Deduce$\Gamma \fCenter \Delta, \lexists[x][!B(x)], \lexists[x][!B(x)], !B(t)$
\RightLabel{\RightR{\lexists}}
\UnaryInf$\Gamma \fCenter \Delta, \lexists[x][!B(x)], \lexists[x][!B(x)]$
\DisplayProof
\]
د استقرا فرضيه پر~\LR{$\pi_1$} تطبيقېږي، او
ثبوتونه~\LR{$\pi_1'$} د \LR{$\Gamma \Sequent \Delta, \lexists[x][!B(x)],
!B(t)$} دپاره شته۔ (په ياد ولرئ چې دلته د معکوسولو لمّې ته
هېڅ اړتيا نۀ شته!) بيا ثبوت~\LR{$\pi'$} دا دے:
\[
\AxiomC{}
\RightLabel{\LR{$\pi_1'$}}
\Deduce$\Gamma \fCenter \Delta,  \lexists[x][!B(x)], !B(t)$
\RightLabel{\RightR{\lexists}}
\UnaryInf$\Gamma \fCenter \Delta, \lexists[x][!B(x)]$
\DisplayProof
\]
چې \LR{$\pheight{\pi'} \le \pheight{\pi}$} دے۔
\end{proof}
\olpseganchor{OLP-0701-B026}{end}

\olpseganchor{OLP-0701-B027}{start}
\begin{prob}
د \olref{prop:G3c-cont-adm} د ثبوت خاکه د \LeftR{\Contraction}
دپاره ورکړئ۔ هغه صورتونه بيان کړئ چې \LR{$!A \ident (!B \land !C)$}
او \LR{$!A \ident (!B
\lif !C)$} وي۔
\end{prob}
\olpseganchor{OLP-0701-B027}{end}

\olpseganchor{OLP-0701-B028}{start}
\begin{prob}
د \olref{prop:G3c-cont-adm} د ثبوت خاکه د \LeftR{\Contraction}
دپاره د کميت‌ښودونکو په پاتې صورتونو کښې ورکړئ، يعنې هلته
چې \LR{$!A \ident
\lforall[x][!B(x)]$} او \LR{$!A \ident \lexists[x][!B(x)]$} وي۔
\end{prob}
\olpseganchor{OLP-0701-B028}{end}

\olpseganchor{OLP-0713-B004}{start}
\olfileid{pt}{seq}{tr}
\olpseganchor{OLP-0713-B004}{end}

\olpseganchor{OLP-0713-B005}{start}
\olsection{د \texorpdfstring{\Log{G1c}}{G1c} او \texorpdfstring{\Log{G3c}}{G3c} تر منځ د ثبوتونو اړول}
\olpseganchor{OLP-0713-B005}{end}

\olpseganchor{OLP-0713-B006}{start}
يادونه مو کړې وه چې \Log{G1c} او \Log{G3c} دواړه د کلاسيکي منطق دپاره سم او بشپړ دي، يعنې هر يو ټول هغه سېکوېنټونه او يوازې هماغه ثابتوي چې په هر جوړښت کښې رښتيا وي۔ په تېره، دواړه هماغه سېکوېنټونه ثابتوي۔ اوس دا خبره د سموالي او بشپړتيا له قضيو پرته، يوازې د ثبوت تيورۍ له لارې ثابتولے شو۔ دا ډول دليل د \Log{G1c} ثبوتونه د \Log{G3c} په ثبوتونه او برعکس د اړولو \emph{کړنلارې} ورکوي۔
\olpseganchor{OLP-0713-B006}{end}

\olpseganchor{OLP-0713-B007}{start}
\begin{prop}
که \LR{$\Log{G1c} \Proves S$} وي، نو \LR{$\Log{G3c} \Proves S$}۔
\end{prop}
\olpseganchor{OLP-0713-B007}{end}

\olpseganchor{OLP-0713-B008}{start}
\begin{proof}
  ښيو چې که \LR{$\pi$} په \Log{G1c} کښې ثبوت وي، نو په \Log{G3c} کښې د~\LR{$S$} ثبوت~\LR{$\pi'$} شته۔ پر \LR{$n
  = \pheight{\pi}$} استقرا کوو۔
\olpseganchor{OLP-0713-B008}{end}

\olpseganchor{OLP-0713-B009}{start}
  که \LR{$n = 0$} وي، نو \LR{$\pi$} بديهي اصل دے۔ \LR{$\lfalse \Sequent
  \quad$} په~\Log{G3c} کښې هم بديهي اصل دے؛ سياق \LR{$\Gamma,
  \Delta$} تش واخلئ۔ په \Log{G1c} کښې \LR{$!A \Sequent !A$} د هر ډول، نۀ يوازې د اتومي~\LR{$!A$}، دپاره بديهي اصول دي۔ خو په \olref[ex]{prop:G3c-axioms} کښې مو ثابته کړې چې دا ټول سېکوېنټونه په~\Log{G3c} کښې د ثابتولو وړ دي۔
\olpseganchor{OLP-0713-B009}{end}

\olpseganchor{OLP-0713-B010}{start}
  که \LR{$n>0$} وي، نو \LR{$\pi$} په قاعده~\LR{$R$} ختمېږي۔ د استقرا فرضيه راکوي چې په~\Log{G3c} کښې د دې قاعدې د مقدماتو ثبوتونه شته؛ يعنې نېغ فرعي ثبوتونه د \Log{G3c} په ثبوتونو اړول کېدے شي۔ که قاعده~\LR{$R$} په \Log{G3c} کښې هم وي، نو د مقدماتو پر اړول شوو ثبوتونه ئې تطبيقوو۔ د هغو قواعدو دپاره چې په~\Log{G3c} کښې عينې نۀ دي، مجازتيا کاروو۔ که قاعده~\LR{$R$} \LR{$\LeftR{\land}$} يا \LR{$\RightR{\lor}$} وي، نو \olref[adm]{prop:lor-G3-adm} تطبيق کړئ۔ که د کمزوري کولو قاعده وي، \olref[adm]{prop:weak-G3c-adm} او که د انقباض قاعده وي، \olref[inv]{prop:G3c-cont-adm} تطبيق کړئ۔ د عمومي کميت د کيڼې او وجودي کميت د ښۍ قاعدې دپاره، لومړے اصلي کمي فورمول د مجاز کمزوري کولو په وسيله د اړولې مقدمې په هماغه خوا کښې ورزيات کړئ؛ بيا هغه قاعده کار کړئ چې دغه نقل په مقدمه کښې ساتي۔
\end{proof}
\olpseganchor{OLP-0713-B010}{end}

\olpseganchor{OLP-0713-B011}{start}
\begin{prop}
که \LR{$\Log{G3c} \Proves S$} وي، نو \LR{$\Log{G1c} \Proves S$}۔
\end{prop}
\olpseganchor{OLP-0713-B011}{end}

\olpseganchor{OLP-0713-B012}{start}
\begin{proof}
  بيا د \LR{$S$} د ثبوت~\LR{$\pi$} پر لوړوالے استقرا کوو۔ د \Log{G3c} هر بديهي اصل په~\Log{G1c} کښې له يوه بديهي اصل څخه د \LeftR{\Weakening} او \RightR{\Weakening} په څو تطبيقونو ثابت کېدے شي:
  \[
  \Axiom$!C \fCenter !C$
  \RightLabel{\LeftR{\Weakening}}
  \Deduce$!C, \Gamma \fCenter !C$
  \RightLabel{\RightR{\Weakening}}
  \Deduce$!C, \Gamma \fCenter \Delta, !C$
  \DisplayProof
\qquad
  \Axiom$\lfalse \fCenter$
  \RightLabel{\LeftR{\Weakening}}
  \Deduce$\lfalse, \Gamma \fCenter$
  \RightLabel{\RightR{\Weakening}}
  \Deduce$\lfalse, \Gamma \fCenter \Delta$
  \DisplayProof
  \]
  که \LR{$\pi$} په قاعده~\LR{$R$} ختمېږي، د استقرا فرضيه په \Log{G1c} کښې د مقدماتو ثبوتونه راکوي؛ چېرته چې قاعده ګډه وي، هماغه قاعده~\LR{$R$} پر دغو ثبوتونه تطبيقوو۔ د عمومي کميت د کيڼې او وجودي کميت د ښۍ قاعدې په مقدمه کښې د اصلي کمي فورمول يو نقل شته۔ کله چې په \Log{G1c} کښې اړونده قاعده تطبيق کړو، په پايله کښې ئې دوه نقلونه راځي؛ په هماغه خوا کښې انقباض يو نقل پرېږدي۔ د \LeftR{\land} او \RightR{\lor} قواعد هم د \Log{G1c} او~\Log{G3c} تر منځ بېل دي۔ د \Log{G3c} بڼې په \Log{G1c} کښې داسې اشتقاقېږي:
  \[
  \Axiom$!A, !B, \Gamma \fCenter \Delta$
  \RightLabel{\LeftR{\land}}
  \UnaryInf$!A \land !B, !B, \Gamma \fCenter \Delta$
  \RightLabel{\LeftR{\land}}
  \UnaryInf$!A \land !B, !A \land !B, \Gamma \fCenter \Delta$
  \RightLabel{\LeftR{\Contraction}}
  \UnaryInf$!A \land !B, \Gamma \fCenter \Delta$
    \DisplayProof
\qquad
  \Axiom$\Gamma \fCenter \Delta, !A, !B$
  \RightLabel{\RightR{\lor}}
  \UnaryInf$\Gamma \fCenter \Delta, !A \lor !B, !B$
  \RightLabel{\RightR{\lor}}
  \UnaryInf$\Gamma \fCenter \Delta, !A \lor !B, !A \lor !B$
  \RightLabel{\RightR{\Contraction}}
  \UnaryInf$\Gamma \fCenter \Delta, !A \lor !B$
    \DisplayProof
  \]
\end{proof}
\olpseganchor{OLP-0713-B012}{end}

\olpseganchor{OLP-0657-B004}{start}
\olchapter{pt}{cut}{د پرېکون لرې کول}
\olpseganchor{OLP-0657-B004}{end}

\olpseganchor{OLP-0659-B004}{start}
\olfileid{pt}{cut}{int}
\olpseganchor{OLP-0659-B004}{end}

\olpseganchor{OLP-0659-B005}{start}
\olsection{سريزه}
\olpseganchor{OLP-0659-B005}{end}

\olpseganchor{OLP-0659-B006}{start}
د \Cut{} قاعده دا ده:
\[
\Axiom$\Gamma \fCenter \Delta, !A$
\Axiom$!A, \Pi \fCenter \Lambda$
\RightLabel{\Cut}
\BinaryInf$\Gamma, \Pi \fCenter \Delta, \Lambda$
\DisplayProof
\]
(فارمول~\LR{$!A$} ته \emph{د پرېکون فارمول} وايي۔)
\olpseganchor{OLP-0659-B006}{end}

\olpseganchor{OLP-0659-B007}{start}
\Cut{} د جېنڅن په اصلي سېکوېنټ حساب~\Log{LK} کښې شامله وه۔
د جېنڅن \emph{هاوپټ‌ساتڅ} («اصلي قضيه») وايي چې په~\Log{LK} کښې
هر د ثابتولو وړ سېکوېنټ د \Cut{} له قاعدې کار اخستلو پرته هم
د ثابتولو وړ دے؛ يعنې \Cut{} د \Log{LK} له ثبوتونو څخه
«لرې کېدے شي»۔ زمونږ نظامونه \Log{G1c} او \Log{G3c}، \Cut نۀ لري،
خو د هغو دپاره هم پوښتنه کولے شو: آيا \Cut{} له هغو ثبوتونه
څخه لرې کېدے شي چې \Cut{} او د \Log{G1c} نور قواعد، يا د
\Log{G3c} قواعد، کاروي؟ په زمونږ ټاکل شوې اصطلاح کښې هم‌ارزښته
پوښتنه دا ده: آيا \Cut{} د \Log{G1c}، يا \Log{G3c}، يوه د منلو
وړ قاعده ده؟
\olpseganchor{OLP-0659-B007}{end}

\olpseganchor{OLP-0659-B008}{start}
د پرېکون لرې کولو قضيه ښايي د ثبوت پوهنې تر ټولو مهمه او
مشهوره پايله وي۔ تر هر څه مخکې ښيي چې \Log{G1c} او \Log{G3c}
هغې قاعدې ته اړ نۀ دي چې په لومړي نظر درته لازمي ښکاري۔ کله
چې واقعي ثبوتونه رسمي کوئ، د مرستندويه قضيو له استعمال سره د
کار طريقه غواړئ۔ رياضيپوهان پايلې ثابتوي او بيا د نورو پايلو
په ثبوتونو کښې ورته رجوع کوي۔ نو لومړے د پايلې \LR{$L$}، يعنې
مرستندويې قضيې، ثبوت د سېکوېنټ \LR{$\Sequent L$} د ثبوت په توګه
رسمي کولے شئ۔ بيا د دويمې پايلې~\LR{$R$} ثبوت، چې~\LR{$L$} کاروي، د
سېکوېنټ \LR{$L \Sequent
R$} د ثبوت په توګه رسمي کوئ۔ څنګه پوهېږئ چې~\LR{$R$} ثبوت لري؛
يعنې يوازې د \LR{$\Sequent R$} يو ثبوت شته؟ د دواړو ثبوتونه
د يوځاے کولو طريقه غواړئ، او د \Cut{} قاعده دا اجازه درکوي۔
\olpseganchor{OLP-0659-B008}{end}

\olpseganchor{OLP-0659-B009}{start}
يادونه مو کړې وه چې \Log{G1c} او \Log{G3c} بشپړ دي، نو ټول
هغه سېکوېنټونه ثابتوي چې د ثابتېدو تمه ئې کوئ؛ يعنې ټول معتبر
سېکوېنټونه۔ دا نېغ په نېغه ثابتېدے شي۔ خو د جېنڅن د کار په
وخت کښې يوازې بديهياتي اشتقاقي نظامونه بشپړ معلوم وو۔ د
\Log{LK} د بشپړتيا د ښودلو دپاره جېنڅن د بديهياتي نظام
اشتقاقونه د \Log{LK} په ثبوتونه بدل کړل۔ په دې بدلولو
کښې \Cut{} اساسي نقش درلود۔ د پرېکون لرې کول يوازې د کلاسيکي
منطق مهمه پايله نۀ ده: ډېر نور منطقونه هم سېکوېنټ حسابونه لري۔
د ځينو منطقونو د سېکوېنټ حسابونو د بشپړتيا نېغ ثبوت سخت يا
ناممکن وي، نو د نورو نظامونو څخه، د \Cut په کارولو، بدلول ئې
د بشپړتيا يوازېنۍ طريقه وي۔ بيا په ثبوت پوهنه د پرېکون لرې
کول ثابتول اړين دي، څو وښيي چې \Cut{} اضافي ده او له \Cut{}
پرته نظام بشپړ دے۔
\olpseganchor{OLP-0659-B009}{end}

\olpseganchor{OLP-0659-B010}{start}
د پرېکون لرې کول ځکه هم مهم دي چې د نورو، په خپل ځاے په زړه
پورې، پايلو د تر لاسه کولو دپاره کارېدے شي۔ دوه بېلګې چې
ګورو ئې د کرېګ د منځګړيتوب مرستندويه قضيه او د هېربراند قضيه دي۔
\olpseganchor{OLP-0659-B010}{end}

\olpseganchor{OLP-0659-B011}{start}
د پرېکون لرې کولو د ثبوت طريقې عموماً ښيي چې له \Cut{} څخه
کار اخستونکے ثبوت څنګه په بې‌پرېکونه ثبوت بدلېدے شي۔ دا
بدلونونه اکثر «ځايي» وي؛ يعنې د يو ثبوت يوه برخه، عموماً په \Cut{} استنتاج ختمېدونکے فرعي
ثبوت، د هماغه سېکوېنټ په بل فرعي ثبوت بدلوو چې پکښې
د \Cut{} استنتاج لرې يا په نورو بدل شوے وي۔ داسې استدلالونه
له \Cut{} سره ثبوتونه په بې‌پرېکونه ثبوتونو د بدلولو ګام
په ګام الګوريتمونه ورکوي۔ دا د پرېکون لرې کولو د طريقو په
استعمال کښې نور معلومات راکوي۔ مثلاً د مدل پوهنې په طريقه
د منځګړيتوب د مرستندويې قضيې داسې ثبوتونه شته: که \LR{$!A \lif !B$}
معتبر وي، نو منځګړی شته (اوس دا پرېږدئ چې منځګړی څه ته وايي)۔
د ثبوت پوهنې استدلال د پرېکون لرې کولو په وسيله داسې الګوريتم
ورکوي چې د \LR{$!A
\lif !B$} يو ثبوت اخلي او منځګړی راګرځوي۔ د مدل پوهنې
له استدلاله بېل، د ثبوت پوهنې استدلال \emph{رغنده} دے۔
\olpseganchor{OLP-0659-B011}{end}

\olpseganchor{OLP-0659-B012}{start}
د پرېکون لرې کولو د قضيې د ثابتولو دوه مشهورې طريقې دي۔
يوه، چې جېنڅن ته ورګرځي، ښيي چې د ثبوت تر ټولو پاسني پرېکونونه
لرې کولے شئ، او بيا ئې تر هغه لرې کوي چې يو هم پاتې نۀ شي۔
بله، چې لومړے شوتې او ټېټ وړاندې کړې ده، په يوه ثبوت
کښې تر ټولو پېچلو پرېکونونو ته ګوري او هغه پرېکونونه په وار
لرې کوي چې د دې معنا له مخې اعظمي دي: هېڅ بل پرېکون تر هغو
ژور د پرېکون فارمول نۀ لري۔ هره طريقه خپلې ګټې او ستونزې
لري۔ د ځينو نظامونو دپاره يوه طريقه کار کوي او بله سختېږي يا
ناممکنه وي۔ نو دواړه طريقې بيانوو۔ د~\Log{G3c} دپاره د
پرېکون لرې کول ثابتوو۔ په حقيقت کښې ښيو چې يوازې \Cut{} نۀ،
بلکې \emph{د ګډ سياق پرېکون}~\CutCS لرې کېدے شي:
\[
\Axiom$\Gamma \fCenter \Delta, !A$
\Axiom$!A, \Gamma \fCenter \Delta$
\RightLabel{\CutCS}
\BinaryInf$\Gamma \fCenter \Delta$
\DisplayProof
\]
د ګډ سياق په پرېکون کښې دواړه مقدمې يو سياق لري او پايله ئې يو نقل ساتي۔
که سېکوېنټ حساب کمزوري کول او انقباض مني، که د \Log{G1c} په
شان اصلي قواعد وي يا د \Log{G3c} په شان د منلو وړ قواعد، نو
\Cut{}، \CutCS تمثيلولے شي او معکوس ئې هم شته۔ د \Cut پر
ځاے \CutCS{} غوره کوو، ځکه د پرېکون لرې کولو لومړۍ طريقه
د~\CutCS دپاره اسانه بيانېږي او دويمه طريقه يوازې د~\CutCS
دپاره کار کوي۔
\olpseganchor{OLP-0659-B012}{end}

\olpseganchor{OLP-0656-B004}{start}
\olfileid{pt}{cut}{top}
\olpseganchor{OLP-0656-B004}{end}

\olpseganchor{OLP-0656-B005}{start}
\olsection{د تر ټولو پاسنيو پرېکونونو لرې کول}
\olpseganchor{OLP-0656-B005}{end}

\olpseganchor{OLP-0656-B006}{start}
\begin{defn}
داسې ثبوت~\LR{$\pi$} وګورئ چې د \CutCS{} په قاعده ختمېږي:
\[
\AxiomC{}
\RightLabel{\LR{$\pi_1$}}
\Deduce$\Gamma \fCenter \Delta, !A$
\AxiomC{}
\RightLabel{\LR{$\pi_2$}}
\Deduce$!A, \Gamma \fCenter \Delta$
\RightLabel{\CutCS}
\BinaryInf$\Gamma \fCenter \Delta$
\DisplayProof
\]
او د \CutCS{} بل هېڅ استنتاج نۀ لري۔ د~\LR{$\pi$} \emph{د پرېکون لوړوالے} دا دے: \LR{$\cheight{\pi} = \pheight{\pi_1} + \pheight{\pi_2}$}۔ د~\LR{$\pi$} \emph{د پرېکون رتبه} دا ده: \LR{$\cutrank{\pi} = \depth{!A}$}۔
\end{defn}
\olpseganchor{OLP-0656-B006}{end}

\olpseganchor{OLP-0656-B007}{start}
\begin{lem}\ollabel{lem:cut-adm-G3c}
د \CutCS{} قاعده په~\Log{G3c} کښې د منلو وړ ده۔ په نورو ټکو، که \LR{$\pi$} يو ثبوت وي چې په يوه~\CutCS{} ختمېږي او په پاتې برخه کښې يوازې د~\LR{$\Log{G3c}$} قاعدې کاروي، نو د هماغه وروستي سېکوېنټ يو ثبوت~\LR{$\pi'$} په~\Log{G3c} کښې شته۔
\end{lem}
\olpseganchor{OLP-0656-B007}{end}

\olpseganchor{OLP-0656-B008}{start}
\begin{proof}
ثبوت د پرېکون پر لوړوالے او رتبه په دوه‌ګونې استقرا کوو؛ رتبه لومړے د کمېدو معيار دے، او په عين رتبه کښې لوړوالے کمېږي۔ ثبوت~\LR{$\pi$} داسې ختمېږي:
\[
\AxiomC{}
\RightLabel{\LR{$\pi_1$}}
\Deduce$\Gamma \fCenter \Delta, !A$
\AxiomC{}
\RightLabel{\LR{$\pi_2$}}
\Deduce$!A, \Gamma \fCenter \Delta$
\RightLabel{\CutCS}
\BinaryInf$\Gamma \fCenter \Delta$
\DisplayProof
\]
\olpseganchor{OLP-0656-B008}{end}

\olpseganchor{OLP-0656-B009}{start}
د استقرا بنسټ: \LR{$\cheight{\pi} = 0$} او \LR{$\cutrank{\pi} = 0$}۔ څرنګه چې \LR{$\cheight{\pi} = \pheight{\pi_1} + \pheight{\pi_2} = 0$}، نو \LR{$\Gamma \Sequent \Delta, !A$} او \LR{$!A, \Gamma \Sequent \Delta$} دواړه بديهي اصول دي۔ دوه حالتونه شته: (الف)~يا \LR{$!A$} په يوه کښې اصلي فارمول نۀ دے، يا (ب)~\LR{$!A$} په دواړو کښې اصلي فارمول دے۔ په دواړو حالاتو کښې لاندې ښيو چې \LR{$\Gamma \Sequent \Delta$} بايد بديهي اصل وي، او د دې دپاره د استقرا فرض نۀ پکارېږي۔
\olpseganchor{OLP-0656-B009}{end}

\olpseganchor{OLP-0656-B010}{start}
اوس د دې له مخې بېلابېل ممکن حالتونه جلا کوو چې \LR{$\pi_1$} او \LR{$\pi_2$} بديهي اصول دي که په استنتاج ختمېږي، او دا چې \LR{$!A$}~په هغه استنتاج کښې اصلي فارمول دے که نۀ۔ لومړے او دويم حالت، \LR{A} او~\LR{B}، د استقرا بنسټ رانغاړي۔ د استقرا په ګام کښې فرض کوو چې \LR{$\cheight{\pi} > 0$} يا \LR{$\cutrank{\pi} > 0$}، يا دواړه۔ د استقرا فرض کارولے شو: د هر هغه ثبوت وروستۍ~\CutCS{} لرې کېدے شي چې په يوه~\CutCS{} ختمېږي او يا ئې د پرېکون رتبه \LR{$= \cutrank{\pi}$} او د پرېکون لوړوالے \LR{$< \cheight{\pi}$} وي، يا ئې د پرېکون رتبه~\LR{$<
\cutrank{\pi}$} وي۔ د \LR{$\cheight{\pi} = 0$} حالت، چې دواړه مقدمې بديهي اصول وي، لومړے او دويم حالت رانغاړي۔ د \LR{$\cheight{\pi} >
0$} حالتونه درېم، څلورم او پنځم حالت، \LR{C} تر~\LR{E}، رانغاړي۔
\olpseganchor{OLP-0656-B010}{end}

\olpseganchor{OLP-0656-B011}{start}
\paragraph{لومړے حالت \LR{A}:}
کوم \LR{$!B$} هم په \LR{$\Gamma$} او هم په~\LR{$\Delta$} کښې راځي۔ نو \LR{$\Gamma \Sequent \Delta$} بديهي اصل دے، که \LR{$!B$}~اټومي وي؛ کنه، د \olref[seq][ex]{prop:G3c-axioms} له مخې په \Log{G3c} کښې بې له \CutCS{} يو ثبوت~\LR{$\pi'$} لري۔ دا د استقرا د بنسټ (الف) حالت رانغاړي: که \LR{$!A$} په \LR{$\Gamma
\Sequent \Delta, !A$} کښې يا په \LR{$!A, \Gamma \Sequent \Delta$} کښې اصلي فارمول نۀ وي او دا بديهي اصول وي، نو کوم بل اټومي~\LR{$!B$} بايد هم په \LR{$\Gamma$} او هم په~\LR{$\Delta$} کښې راشي۔ د استقرا هغه ګام هم رانغاړي چې يوه مقدمه ئې بديهي اصل وي او \LR{$!A$} پکښې اصلي فارمول نۀ وي، که څه هم بله مقدمه د يوې قاعدې پايله وي۔
\olpseganchor{OLP-0656-B011}{end}

\olpseganchor{OLP-0656-B012}{start}
\paragraph{دويم حالت \LR{B}:} دواړه مقدمې بديهي اصول دي او \LR{$!A$} اصلي فارمول دے؛ نو اټومي دے۔ کيڼه مقدمه \LR{$\Gamma \Sequent
\Delta, !A$} وګورئ۔ څرنګه چې \LR{$!A$} اصلي فارمول دے، بايد په \LR{$\Gamma$} کښې راشي، کنه \LR{$\Gamma \Sequent \Delta, !A$} به بديهي اصل نۀ وي۔ همداسې \LR{$!A$} بايد په \LR{$\Delta$} کښې راشي، کنه \LR{$!A, \Gamma \Sequent \Delta$} به بديهي اصل نۀ وي۔ نو \LR{$!A$} اټومي دے او هم په \LR{$\Gamma$} او هم په \LR{$\Delta$} کښې راځي؛ يعنې \LR{$\Gamma \Sequent \Delta$} بديهي اصل دے۔ دا د استقرا د بنسټ (ب) حالت رانغاړي۔
\olpseganchor{OLP-0656-B012}{end}

\olpseganchor{OLP-0656-B013}{start}
\begin{center}
د لومړي او دويم حالت بديلې بڼې
\end{center}
\olpseganchor{OLP-0656-B013}{end}

\olpseganchor{OLP-0656-B014}{start}
\paragraph{لومړے حالت \LR{A}:} د پرېکون يوه مقدمه د \LR{$!C, \Gamma' \Sequent \Delta', !C$} په بڼه بديهي اصل دے، چې \LR{$!C$}~اټومي دے۔ که \LR{$!C$} د پرېکون فارمول~\LR{$!A$} نۀ وي، نو \LR{$!C$} بايد هم په \LR{$\Gamma$} او هم په~\LR{$\Delta$} کښې راشي۔ په دې حالت کښې \LR{$\Gamma \Sequent \Delta$} بديهي اصل دے، او همدا د~\LR{$\pi'$} په توګه اخيستلے شو۔ که \LR{$!C \ident !A$} د پرېکون فارمول وي، نو يا کيڼه مقدمه بديهي اصل ده، \LR{$!A \in \Gamma$} او \LR{$\Gamma = \Gamma', !A$}؛ يا ښۍ مقدمه بديهي اصل ده، \LR{$!A \in
\Delta$} او \LR{$\Delta = \Delta', !A$}۔ په لومړي صورت کښې \LR{$\pi_2$} د \LR{$!A, !A, \Gamma' \Sequent \Delta$} يو ثبوت دے، او د \LeftR{\Contraction} د منلو وړتيا، \olref[seq][inv]{prop:G3c-cont-adm}، د \LR{$\Gamma \Sequent \Delta$} يو ثبوت راکوي۔ په دويم صورت کښې \LR{$\pi_1$} د \LR{$\Gamma \Sequent \Delta', !A, !A$} يو ثبوت دے، او د \RightR{\Contraction} د منلو وړتيا په وسيله د \LR{$\Gamma \Sequent
\Delta$} يو ثبوت ترلاسه کوو۔
\olpseganchor{OLP-0656-B014}{end}

\olpseganchor{OLP-0656-B015}{start}
\paragraph{دويم حالت \LR{B}:} يوه مقدمه د~\LR{$\lfalse, \Gamma' \Sequent \Delta'$} په بڼه بديهي اصل ده۔ که کيڼه مقدمه داسې بديهي اصل وي، نو \LR{$\lfalse \in \Gamma$}، او \LR{$\Gamma \Sequent \Delta$} هم بديهي اصل دے۔ که ښۍ مقدمه داسې بديهي اصل وي، نو يا \LR{$\lfalse \in \Gamma$}، چې بيا \LR{$\Gamma \Sequent \Delta$} بديهي اصل دے، يا \LR{$!A \ident \lfalse$}۔ په وروستي صورت کښې \LR{$\pi_1$} د \LR{$\Gamma \Sequent \Delta, \lfalse$} يو ثبوت دے۔ په~\LR{$\pi_1$} کښې د هر سېکوېنټ له تالي څخه د~\LR{$\lfalse$} د يوه نقل په لرې کولو ئې د~\LR{$\Gamma \Sequent \Delta$} په يوه ثبوت بدلولے شو۔ (تمرين: دا پر~\LR{$\pheight{\pi_1}$} د استقرا په وسيله وڅېړئ۔)
\olpseganchor{OLP-0656-B015}{end}

\olpseganchor{OLP-0656-B016}{start}
اوس فرض کړئ چې نۀ لومړے حالت او نۀ دويم حالت لګېږي۔ په پاتې حالاتو کښې \LR{$\cheight{\pi} > 0$}؛ يعنې لږ تر لږه يوه مقدمه د استنتاج پايله ده۔
\olpseganchor{OLP-0656-B016}{end}

\olpseganchor{OLP-0656-B017}{start}
\paragraph{درېم او څلورم حالت، \LR{C} او~\LR{D}: د پرېکونونو ځاے بدلول۔}
په دې دواړو حالاتو کښې څېړل شوي امکانونه يوه ګډه بڼه لري۔ له داسې ثبوت څخه پيل کوو چې په~\CutCS{} ختمېږي او يوه مقدمه ئې د قاعدې~\LR{$R$} په وسيله ترلاسه شوې وي، خو د پرېکون فارمول~\LR{$!A$} په هغې کښې اصلي فارمول نۀ وي؛ کوم بل فارمول~\LR{$!D$} اصلي وي۔ په دې ثبوت کښې \CutCS{} له قاعدې~\LR{$R$} وروسته راځي، او لنډه بڼه ئې داسې ده:
\[
\AxiomC{}
\RightLabel{\LR{$\pi_1$}}
\Deduce$\Gamma \fCenter \Delta', !D, !A$
\AxiomC{}
\RightLabel{\LR{$\pi_3$}}
\Deduce$!A, \Gamma, \Pi \fCenter \Delta', \Lambda$
\RightLabel{\LR{$R$}}
\UnaryInf$!A, \Gamma \fCenter \Delta', !D$
\RightSubproofLabel{\LR{$\pi_2$}}
\RightLabel{\CutCS}
\BinaryInf$\Gamma \fCenter \Delta', !D$
\DisplayProof
\]
دا يوازې هغه حالت رانغاړي چې \LR{$R$}~يوه يو‌مقدمې ښۍ قاعده وي، \LR{$A$}~د~\LR{$R$} اصلي فارمول نۀ وي، او اصلي فارمول~\LR{$!D$} د~\CutCS د ښۍ مقدمې په تالي کښې راځي۔ که \LR{$!D$}~په مقدم کښې راشي، يعنې \LR{$R$} کيڼه قاعده وي، يا د~\CutCS په کيڼه مقدمه کښې راشي، نو بڼه به ئې بېله وي۔ په \LR{$\Pi$} او~\LR{$\Lambda$} کښې فارمولونه د قاعدې~\LR{$R$} فعال فارمولونه ښيي۔
\olpseganchor{OLP-0656-B017}{end}

\olpseganchor{OLP-0656-B018}{start}
دا ترتيب سرچپه کوو او داسې ثبوت ګورو چې \LR{$R$} په کښې له~\CutCS{} وروسته راځي:
\[
\AxiomC{}
\RightLabel{\LR{$\pi_1'$}}
\Deduce$\Gamma, \Pi \fCenter \Delta', !D, \Lambda, !A$
\AxiomC{}
\RightLabel{\LR{$\pi_3'$}}
\Deduce$!A, \Gamma, \Pi \fCenter \Delta', !D, \Lambda$
\RightLabel{\CutCS}
\BinaryInf$\Gamma, \Pi \fCenter \Delta', !D, \Lambda$
\RightLabel{\LR{$R$}}
\UnaryInf$\Gamma, \fCenter \Delta', !D, !D$
\DisplayProof
\]
د لوړوالے ساتونکي کمزوري کولو په وسيله له \LR{$\pi_1$} او \LR{$\pi_3$} څخه په ترتيب سره \LR{$\pi_1'$} او \LR{$\pi_3'$} ترلاسه کوو۔
\olpseganchor{OLP-0656-B018}{end}

\olpseganchor{OLP-0656-B019}{start}
څرنګه چې د \CutCS{} او~\LR{$R$} ترتيب بدلوو، دې ته «د~\LR{$R$} څخه پورته د پرېکون ځاے بدلول» وايي۔ دا د هغه ثبوت لوړوالے کموي چې د~\CutCS په ښۍ مقدمه ختمېږي؛ نو د پرېکون لوړوالے هم کموي۔ يعنې په نوي \CutCS{} ختمېدونکي فرعي ثبوتونه داسې اخيستلے شو چې لوړوالے ئې له \LR{$\pi_1$} او \LR{$\pi_3$} څخه زيات نۀ وي۔ د پرېکون لوړوالے کم شو، ځکه په ښي اړخ کښې مو يو استنتاج لرې کړ۔ اوس د استقرا فرض لګېږي او وايي چې د \CutCS{} استنتاج لرې کېدے شي۔ په پاے کښې د~\LR{$\RightR{\Contraction}$} د منلو وړتيا په وسيله د~\LR{$!D$} زيات نقل لرې کوو۔
\olpseganchor{OLP-0656-B019}{end}

\olpseganchor{OLP-0656-B020}{start}
\paragraph{درېم حالت \LR{C}:}
لږ تر لږه \LR{$\pi_1$} يا \LR{$\pi_2$} په يوه يو‌مقدمې قاعده~\LR{$R$} ختمېږي چې \LR{$!A$}~پکښې اصلي فارمول نۀ دے۔ ټول اتلس حالتونه شته: د دې له مخې چې \LR{$!A$}~د~\CutCS په کيڼه مقدمه کښې اصلي نۀ دے که په ښۍ، او دا چې \LR{$R$} له نهو يو‌مقدمې قاعدو، \LeftR{\lnot}، \RightR{\lnot}، \RightR{\lor}، \LeftR{\land}، \RightR{\lif} او د کمي کوونکو له څلورو قاعدو څخه کومه يوه ده۔ يوازې دوه بېلګې ګورو: \LR{$!A$}~د~\LR{$\pi_2$} په وروستي استنتاج کښې اصلي فارمول نۀ دے، او هغه په~\LR{$\RightR{\lif}$} يا په~\LeftR{\lexists} ختمېږي۔ نور حالتونه ورته دي او د تمرينونو په توګه پاتې دي۔
\olpseganchor{OLP-0656-B020}{end}

\olpseganchor{OLP-0656-B021}{start}
فرض کړئ چې \LR{$\pi$} داسې ختمېږي:
\[
\AxiomC{}
\RightLabel{\LR{$\pi_1$}}
\Deduce$\Gamma \fCenter \Delta', !B \lif !C, !A$
\AxiomC{}
\RightLabel{\LR{$\pi_3$}}
\Deduce$!A, !B, \Gamma \fCenter \Delta', !C$
\RightLabel{\LR{$\RightR{\lif}$}}
\UnaryInf$!A, \Gamma \fCenter \Delta', !B \lif !C$
\RightSubproofLabel{\LR{$\pi_2$}}
\RightLabel{\CutCS}
\BinaryInf$\Gamma \fCenter \Delta', !B \lif !C$
\DisplayProof
\]
چې \LR{$\Delta = \Delta', !B \lif !C$}۔ پر \LR{$\pi_1$} د لوړوالے ساتونکي کمزوري کولو، \olref[seq][adm]{prop:weak-G3c-adm}، په لګولو د \LR{$!B, \Gamma \fCenter \Delta', !B \lif !C, !C,
!A$} يو ثبوت \LR{$\pi_1'$} ترلاسه کوو؛ په کيڼ اړخ کښې \LR{$!B$} او په ښي اړخ کښې \LR{$!C$} ورزياتوو۔ همداسې پر~\LR{$\pi_3$} د \olref[seq][adm]{prop:weak-G3c-adm} په لګولو د \LR{$!A, !B, \Gamma \fCenter
\Delta', !B \lif !C, !C$} يو ثبوت~\LR{$\pi_3'$} ترلاسه کوو؛ په ښي اړخ کښې \LR{$!B \lif !C$} ورزياتوو۔ د استقرا فرض پر دې لګېږي:
\[
\AxiomC{}
\RightLabel{\LR{$\pi_1'$}}
\Deduce$!B, \Gamma \fCenter \Delta', !B \lif !C, !C, !A$
\AxiomC{}
\RightLabel{\LR{$\pi_3'$}}
\Deduce$!A, !B, \Gamma \fCenter \Delta', !B \lif !C, !C$
\RightLabel{\CutCS}
\BinaryInf$!B, \Gamma \fCenter \Delta', !B \lif !C, !C$
\DisplayProof
\]
د پرېکون فارمول~\LR{$!A$} دے۔ دا ثبوت له~\LR{$\pi$} سره د پرېکون عين رتبه لري، خو د پرېکون لوړوالے ئې کم دے۔ په \Log{G3c} کښې بې له \CutCS{} د~\LR{$!B, \Gamma \fCenter \Delta', !B \lif
!C, !C$} يو ثبوت~\LR{$\pi_4$} راکوي۔ د \LR{$\RightR{\lif}$} په لګولو د \LR{$\Gamma \fCenter
\Delta', !B \lif !C, !B \lif !C$} يو ثبوت ترلاسه کوو:
\[
\AxiomC{}
\RightLabel{\LR{$\pi_4$}}
\Deduce$!B, \Gamma \fCenter \Delta', !B \lif !C, !C$
\RightLabel{\LR{$\RightR{\lif}$}}
\UnaryInf$\Gamma \fCenter \Delta', !B \lif !C, !B \lif !C$
\DisplayProof
\]
د \LR{$\Gamma \Sequent \Delta$}، يعنې د \LR{$\Gamma
\fCenter \Delta', !B \lif !C$}، د يوه ثبوت~\LR{$\pi'$} د ترلاسه کولو دپاره د انقباض د منلو وړتيا، \olref[seq][inv]{prop:G3c-cont-adm}، ولګوئ۔
\olpseganchor{OLP-0656-B021}{end}

\olpseganchor{OLP-0656-B022}{start}
اوس فرض کړئ چې \LR{$\pi$} داسې ختمېږي:
\[
\AxiomC{}
\RightLabel{\LR{$\pi_1$}}
\Deduce$\lexists[x][!B(x)], \Gamma' \fCenter \Delta, !A$
\AxiomC{}
\RightLabel{\LR{$\pi_3$}}
\Deduce$!A, !B(c), \Gamma' \fCenter \Delta$
\RightLabel{\LR{$\LeftR{\lexists}$}}
\UnaryInf$!A, \lexists[x][!B(x)], \Gamma' \fCenter \Delta$
\RightSubproofLabel{\LR{$\pi_2$}}
\RightLabel{\CutCS}
\BinaryInf$\lexists[x][!B(x)], \Gamma' \fCenter \Delta$
\DisplayProof
\]
د استقرا فرض بيا پر دې لګوو:
\[
\AxiomC{}
\RightLabel{\LR{$\pi_1[!B(c) \Sequent {}]$}}
\Deduce$\lexists[x][!B(x)], !B(c), \Gamma' \fCenter \Delta, !A$
\AxiomC{}
\LeftLabel{\LR{$\pi_3[\lexists[x][!B(x)] \Sequent {}]$}}
\Deduce$!A, \lexists[x][!B(x)], !B(c), \Gamma' \fCenter \Delta$
\RightLabel{\CutCS}
\BinaryInf$\lexists[x][!B(x)], !B(c), \Gamma' \fCenter \Delta$
\RightLabel{\LeftR{\lexists}}
\UnaryInf$\lexists[x][!B(x)], \lexists[x][!B(x)], \Gamma' \fCenter \Delta$
\DisplayProof
\]
د ځانګړي متغير شرط پوره دے، ځکه په~\LR{$\pi_2$} کښې د~\LeftR{\lexists} هماغه شرط تضمينوي چې \LR{$c$} په \LR{$\lexists[x][!B(x)]$}، \LR{$\Gamma'$} يا~\LR{$\Delta$} کښې نۀ راځي۔ په پاے کښې \olref[seq][inv]{prop:G3c-cont-adm} ولګوئ۔
\olpseganchor{OLP-0656-B022}{end}

\olpseganchor{OLP-0656-B023}{start}
\paragraph{څلورم حالت \LR{D}:}
لږ تر لږه \LR{$\pi_1$} يا \LR{$\pi_2$} په يوه دوه‌مقدمې قاعده~\LR{$R$} ختمېږي چې \LR{$!A$} پکښې اصلي فارمول نۀ دے۔ شپږ حالتونه شته۔ هغه حالت ګورو چې \LR{$!A$} د~\LR{$\pi_2$} په وروستي استنتاج کښې اصلي نۀ دے او هغه په~\LR{$\RightR{\land}$} ختمېږي؛ نور حالتونه ورته دي۔
\olpseganchor{OLP-0656-B023}{end}

\olpseganchor{OLP-0656-B024}{start}
فرض کړئ چې \LR{$\pi$} داسې ختمېږي:
\[
\AxiomC{}
\RightLabel{\LR{$\pi_1$}}
\Deduce$\Gamma \fCenter \Delta', !B \land !C, !A$
\AxiomC{}
\RightLabel{\LR{$\pi_3$}}
\Deduce$!A, \Gamma \fCenter \Delta', !B$
\AxiomC{}
\RightLabel{\LR{$\pi_4$}}
\Deduce$!A, \Gamma \fCenter \Delta', !C$
\RightLabel{\LR{$\RightR{\land}$}}
\BinaryInf$!A, \Gamma \fCenter \Delta', !B \land !C$
\RightSubproofLabel{\LR{$\pi_2$}}
\RightLabel{\CutCS}
\BinaryInf$\Gamma \fCenter \Delta$
\DisplayProof
\]
د استقرا فرض پر دې ثبوتونه لګېږي:
\[
\AxiomC{}
\RightLabel{\LR{$\pi_1'$}}
\Deduce$\Gamma \fCenter \Delta', !B \land !C, !B, !A$
\AxiomC{}
\RightLabel{\LR{$\pi_3'$}}
\Deduce$!A, \Gamma \fCenter \Delta', !B \land !C, !B$
\RightLabel{\CutCS}
\BinaryInf$\Gamma \fCenter \Delta', !B \land !C, !B$
\DisplayProof
\]
او
\[
\AxiomC{}
\RightLabel{\LR{$\pi_1''$}}
\Deduce$\Gamma \fCenter \Delta', !B \land !C, !C, !A$
\AxiomC{}
\RightLabel{\LR{$\pi_4'$}}
\Deduce$!A, \Gamma \fCenter \Delta', !B \land !C, !C$
\RightLabel{\CutCS}
\BinaryInf$\Gamma \fCenter \Delta', !B \land !C, !C$
\DisplayProof
\]
دلته د \Log{G3c} ثبوتونه \LR{$\pi_1'$} او \LR{$\pi_1''$} له~\LR{$\pi_1$} څخه د لوړوالے ساتونکي کمزوري کولو، \olref[seq][adm]{prop:weak-G3c-adm}، په وسيله ترلاسه کېږي؛ يو ځل په ښي اړخ کښې \LR{$!B$} او بل ځل \LR{$!C$} ورزياتوو۔ ثبوتونه \LR{$\pi_3'$} او~\LR{$\pi_4'$} هم له \LR{$\pi_3$} او~\LR{$\pi_4$} څخه په ترتيب سره د لوړوالے ساتونکي کمزوري کولو په وسيله ترلاسه کېږي؛ په \LR{$\pi_3$} او~\LR{$\pi_4$} کښې په ترتيب په ښي اړخ کښې \LR{$!B \land !C$} ورزياتوو۔
\olpseganchor{OLP-0656-B024}{end}

\olpseganchor{OLP-0656-B025}{start}
د استقرا فرض په \Log{G3c} کښې د \LR{$ \Gamma \fCenter \Delta', !B \land !C, !B$} او \LR{$\Gamma
\fCenter \Delta', !B \land !C, !C$} ثبوتونه~\LR{$\pi_5$} او~\LR{$\pi_6$} راکوي۔ د قاعدې~\LR{$\RightR{\land}$} په وسيله ئې د \LR{$\Gamma \Sequent \Delta', !B \land !C, !B \land !C$} په يوه ثبوت کښې يوځاے کوو:
\[
\AxiomC{}
\RightLabel{\LR{$\pi_5$}}
\Deduce$\Gamma \fCenter \Delta', !B \land !C, !B$
\AxiomC{}
\RightLabel{\LR{$\pi_6$}}
\Deduce$\Gamma \fCenter \Delta', !B \land !C, !C$
\RightLabel{\RightR{\land}}
\BinaryInf$\Gamma \fCenter \Delta', !B \land !C, !B \land !C$
\DisplayProof
\]
د \LR{$\Gamma \Sequent \Delta$}، يعنې \LR{$\Gamma \fCenter
\Delta', !B \land !C$}، د يوه ثبوت د ترلاسه کولو دپاره د انقباض د منلو وړتيا، \olref[seq][inv]{prop:G3c-cont-adm}، ولګوئ۔
\olpseganchor{OLP-0656-B025}{end}

\olpseganchor{OLP-0656-B026}{start}
\paragraph{پنځم حالت \LR{E}: د پرېکونونو راکمول۔}
وروستے حالت هغه دے چې \LR{$\pi_1$} او \LR{$\pi_2$} دواړه په داسې استنتاجونو ختمېږي چې \LR{$!A$} پکښې اصلي فارمول دے۔ شپږ حالتونه شته، د~\LR{$!A$} د هر ممکن اصلي عملګر دپاره يو۔
\olpseganchor{OLP-0656-B026}{end}

\olpseganchor{OLP-0656-B027}{start}
په هر حالت کښې د پرېکون فارمول~\LR{$!A$} لرونکے \CutCS{} د~\LR{$!A$} پر فرعي فارمولونه په يوه يا ډېرو پرېکونونو بدلوو؛ يعنې د هغو استنتاجونو پر فعالو فارمولونه چې~\LR{$!A$} ئې اصلي فارمول دے۔ له همدې کبله دې حالتونو ته ډېر ځله د پرېکون د استنتاج «راکمونه» يا «بدلونه» وايي۔
\olpseganchor{OLP-0656-B027}{end}

\olpseganchor{OLP-0656-B028}{start}
\begin{enumerate}
\item که \LR{$!A \ident \lnot !B$} وي، نو \LR{$\pi$} داسې ختمېږي:
\[
\AxiomC{}
\RightLabel{\LR{$\pi_1'$}}
\Deduce$!B, \Gamma \fCenter \Delta$
\RightLabel{\RightR{\lnot}}
\UnaryInf$\Gamma \fCenter \Delta, \lnot !B$
\LeftSubproofLabel{\LR{$\pi_1$}}
\AxiomC{}
\RightLabel{\LR{$\pi_2'$}}
\Deduce$\Gamma \fCenter \Delta, !B$
\RightLabel{\LeftR{\lnot}}
\UnaryInf$\lnot !B, \Gamma \fCenter \Delta$
\RightSubproofLabel{\LR{$\pi_2$}}
\RightLabel{\CutCS}
\BinaryInf$\Gamma \fCenter \Delta$
\DisplayProof
\]
د استقرا فرض له مخې له دې څخه \CutCS{} لرې کېدے شي:
\[
\AxiomC{}
\RightLabel{\LR{$\pi_2'$}}
\Deduce$\Gamma \fCenter \Delta, !B$
\AxiomC{}
\RightLabel{\LR{$\pi_1'$}}
\Deduce$!B, \Gamma \fCenter \Delta$
\RightLabel{\CutCS}
\BinaryInf$\Gamma \fCenter \Delta$
\DisplayProof
\]
\olpseganchor{OLP-0656-B028}{end}

\olpseganchor{OLP-0656-B029}{start}
\item
که \LR{$!A \ident (!B \lor !C)$} وي، نو \LR{$\pi$} داسې ختمېږي:
\[
\AxiomC{}
\RightLabel{\LR{$\pi_1'$}}
\Deduce$\Gamma \fCenter \Delta, !B, !C$
\RightLabel{\RightR{\lor}}
\UnaryInf$\Gamma \fCenter \Delta, !B \lor !C$
\LeftSubproofLabel{\LR{$\pi_1$}}
\AxiomC{}
\RightLabel{\LR{$\pi_2'$}}
\Deduce$!B, \Gamma \fCenter \Delta$
\AxiomC{}
\RightLabel{\LR{$\pi_2''$}}
\Deduce$!C, \Gamma \fCenter \Delta$
\RightLabel{\LeftR{\lor}}
\BinaryInf$!B \lor !C, \Gamma \fCenter \Delta$
\RightSubproofLabel{\LR{$\pi_2$}}
\RightLabel{\CutCS}
\BinaryInf$\Gamma \fCenter \Delta$
\DisplayProof
\]
دا ثبوت وګورئ چې په \CutCS ختمېږي:
\[
\AxiomC{}
\RightLabel{\LR{$\pi_1'$}}
\Deduce$\Gamma \fCenter \Delta, !B, !C$
\AxiomC{}
\RightLabel{\LR{$\pi_2'''$}}
\Deduce$!C, \Gamma \fCenter \Delta, !B$
\RightLabel{\CutCS}
\BinaryInf$\Gamma \fCenter \Delta, !B$
\DisplayProof
\]
چې \LR{$\pi_2'''$} له \LR{$\pi_2''$} څخه د لوړوالے ساتونکي کمزوري کولو په وسيله ترلاسه شوے دے۔ د پرېکون رتبه ئې \LR{$< \cutr{\pi}$} ده، ځکه \LR{$\depth{!C} < \depth{!B \lor !C}$}؛ نو د استقرا فرض په \Log{G3c} کښې د \LR{$\Gamma \fCenter \Delta, !B$} يو ثبوت \LR{$\pi_1''$} راکوي۔ دا ثبوت چې په \CutCS ختمېږي:
\[
\AxiomC{}
\RightLabel{\LR{$\pi_1''$}}
\Deduce$\Gamma \fCenter \Delta, !B$
\AxiomC{}
\RightLabel{\LR{$\pi_2'$}}
\Deduce$!B, \Gamma \fCenter \Delta$
\RightLabel{\CutCS}
\BinaryInf$\Gamma \fCenter \Delta$
\DisplayProof
\]
د پرېکون رتبه \LR{$= \depth{!B} < \depth{!B \lor !C}$} لري؛ نو د استقرا فرض بيا لګېږي او د \LR{$\Gamma
\fCenter \Delta$} يو ثبوت \LR{$\pi'$} راکوي۔
\olpseganchor{OLP-0656-B029}{end}

\olpseganchor{OLP-0656-B030}{start}
\item \LR{$!A \ident (!B \land !C)$}۔ تمرين۔
\olpseganchor{OLP-0656-B030}{end}

\olpseganchor{OLP-0656-B031}{start}
\item که \LR{$!A \ident (!B \lif !C)$} وي، نو \LR{$\pi$} داسې ختمېږي:
\[
\AxiomC{}
\RightLabel{\LR{$\pi_1'$}}
\Deduce$!B, \Gamma \fCenter \Delta, !C$
\RightLabel{\RightR{\lif}}
\UnaryInf$\Gamma \fCenter \Delta, !B \lif !C$
\LeftSubproofLabel{\LR{$\pi_1$}}
\AxiomC{}
\RightLabel{\LR{$\pi_2'$}}
\Deduce$\Gamma \fCenter \Delta, !B$
\AxiomC{}
\RightLabel{\LR{$\pi_2''$}}
\Deduce$!C, \Gamma \fCenter \Delta$
\RightLabel{\LeftR{\lif}}
\BinaryInf$!B \lif !C, \Gamma \fCenter \Delta$
\RightSubproofLabel{\LR{$\pi_2$}}
\RightLabel{\CutCS}
\BinaryInf$\Gamma \fCenter \Delta$
\DisplayProof
\]
دا ثبوت چې په \CutCS ختمېږي:
\[
\AxiomC{}
\RightLabel{\LR{$\pi_1'$}}
\Deduce$!B, \Gamma \fCenter \Delta, !C$
\AxiomC{}
\RightLabel{\LR{$\pi_2''$}}
\Deduce$!C, \Gamma \fCenter \Delta$
\RightLabel{\CutCS}
\BinaryInf$!B, \Gamma \fCenter \Delta$
\DisplayProof
\]
له~\LR{$\pi$} څخه د پرېکون ټيټه رتبه لري۔ د استقرا فرض د \LR{$!B, \Gamma \fCenter \Delta$} يو ثبوت~\LR{$\pi_1''$} راکوي۔ اوس دا ثبوت وګورئ چې په \CutCS ختمېږي:
\[
\AxiomC{}
\RightLabel{\LR{$\pi_2'$}}
\Deduce$\Gamma \fCenter \Delta, !B$
\AxiomC{}
\RightLabel{\LR{$\pi_1''$}}
\Deduce$!B, \Gamma \fCenter \Delta$
\RightLabel{\CutCS}
\BinaryInf$\Gamma \fCenter \Delta$
\DisplayProof
\]
د پرېکون رتبه ئې بيا هم د~\LR{$\pi$} د پرېکون له رتبې څخه ټيټه ده؛ نو د استقرا فرض د \LR{$\Gamma \fCenter
\Delta$} يو \Log{G3c}-ثبوت~\LR{$\pi'$} راکوي۔
\item که \LR{$!A \ident \lforall[x][!B(x)]$} وي، نو \LR{$\pi$} داسې ختمېږي:
\[
\AxiomC{}
\RightLabel{\LR{$\pi_1'$}}
\Deduce$\Gamma \fCenter \Delta, !B(c)$
\RightLabel{\RightR{\lforall}}
\UnaryInf$\Gamma \fCenter \Delta, \lforall[x][!B(x)]$
\LeftSubproofLabel{\LR{$\pi_1$}}
\AxiomC{}
\RightLabel{\LR{$\pi_2'$}}
\Deduce$!B(t), \lforall[x][!B(x)], \Gamma \fCenter \Delta$
\RightLabel{\LeftR{\lforall}}
\UnaryInf$\lforall[x][!B(x)], \Gamma \fCenter \Delta$
\RightSubproofLabel{\LR{$\pi_2$}}
\RightLabel{\CutCS}
\BinaryInf$\Gamma \fCenter \Delta$
\DisplayProof
\]
د استقرا فرض له مخې له دې څخه \CutCS{} لرې کېدے شي:
\[
\AxiomC{}
\RightLabel{\LR{$\pi_1''$}}
\Deduce$!B(t), \Gamma \fCenter \Delta, \lforall[x][!B(x)]$
\AxiomC{}
\RightLabel{\LR{$\pi_2'$}}
\Deduce$!B(t), \lforall[x][!B(x)], \Gamma \fCenter \Delta$
\RightLabel{\CutCS}
\BinaryInf$!B(t), \Gamma \fCenter \Delta$
\DisplayProof
\]
چې \LR{$\pi_1''$} له \LR{$\pi_1$} څخه، نۀ له~\LR{$\pi_1'$} څخه، د لوړوالے ساتونکي کمزوري کولو په وسيله ترلاسه شوے دے۔ دا ثبوت د پرېکون عين رتبه، خو د پرېکون ټيټ لوړوالے لري۔ د \LR{$!B(t), \Gamma \fCenter \Delta$} يو ثبوت~\LR{$\pi_2''$} راکوي۔
\olpseganchor{OLP-0656-B031}{end}

\olpseganchor{OLP-0656-B032}{start}
ثبوت~\LR{$\pi_1'$} وروستے سېکوېنټ \LR{$\Gamma \Sequent \Delta,
!B(c)$} لري، چې \LR{$c$}~د~\LR{$\RightR{\lforall}$} د يوه استنتاج ځانګړے متغير دے۔ نو \LR{$c$}~په~\LR{$!B(x)$}، \LR{$\Gamma$} يا~\LR{$\Delta$} کښې نۀ راځي۔ فرض کوو چې ثبوت~\LR{$\pi$} منظم دے: \LR{$c$}~يوازې د ورپسې~\LR{$\RightR{\lforall}$} د استنتاج ځانګړے متغير دے او يوازې له هغه پورته راځي؛ يعنې يوازې په~\LR{$\pi_1'$} کښې راځي او په~\LR{$\pi_1'$} کښې د کوم استنتاج ځانګړے متغير نۀ دے۔ څرنګه چې \LR{$c$} په~\LR{$\pi_2$} کښې نۀ راځي، په~\LR{$t$} کښې هم نۀ شي راتلے۔ د \olref[seq][qua]{lem:subst-regular} له مخې ثبوت~\LR{$\Subst{\pi_1'}{t}{c}$} د \LR{$\Gamma \Sequent \Delta,
!B(t)$} يو ثبوت دے۔ اوس د استقرا فرض پر دې ثبوت لګوو:
\[
\AxiomC{}
\RightLabel{\LR{$\Subst{\pi_1'}{t}{c}$}}
\Deduce$\Gamma \fCenter \Delta, !B(t)$
\AxiomC{}
\RightLabel{\LR{$\pi_2''$}}
\Deduce$!B(t), \Gamma \fCenter \Delta$
\RightLabel{\CutCS}
\BinaryInf$\Gamma \fCenter \Delta$
\DisplayProof
\]
د استقرا فرض لګېږي، ځکه د پرېکون فارمول~\LR{$!B(t)$} دے؛ نو دا ثبوت له~\LR{$\pi$} څخه د پرېکون ټيټه رتبه لري۔
\item د \LR{$!A \ident \lexists[x][!B(x)]$} حالت د تمرين په توګه پرېږدو۔
\end{enumerate}
\end{proof}
\olpseganchor{OLP-0656-B032}{end}

\olpseganchor{OLP-0656-B033}{start}
\begin{prob}
د \olref{lem:cut-adm-G3c} په ثبوت کښې د څلورم حالت \LR{D} هغه فرعي حالت وڅېړئ چې \LR{$!A$} د~\LR{$\pi_1$} په وروستي استنتاج کښې اصلي فارمول نۀ وي او هغه په~\LR{$\LeftR{\lif}$} ختمېږي۔ پام وکړئ چې د دې تمرين يوه برخه د هغه څه څرګند بيان دے چې ثابتول ئې غواړئ؛ يعنې په دې حالت کښې د \LR{$\pi$} بڼه څه ده؟
\end{prob}
\olpseganchor{OLP-0656-B033}{end}

\olpseganchor{OLP-0656-B034}{start}
\begin{prob}
د \olref{lem:cut-adm-G3c} په ثبوت کښې د درېم حالت \LR{C} هغه فرعي حالت وڅېړئ چې \LR{$!A$} د~\LR{$\pi_1$} په وروستي استنتاج کښې اصلي فارمول نۀ وي او هغه په~\LR{$\LeftR{\lforall}$} ختمېږي۔
\end{prob}
\olpseganchor{OLP-0656-B034}{end}

\olpseganchor{OLP-0656-B035}{start}
\begin{prob}
د \olref{lem:cut-adm-G3c} په ثبوت کښې د پنځم حالت \LR{E} هغه فرعي حالت وڅېړئ چې \LR{$!A$} د \LR{$\pi_1$} او \LR{$\pi_2$} په دواړو وروستيو استنتاجونو کښې اصلي فارمول وي او \LR{$!A \ident (!B \land !C)$}۔
\end{prob}
\olpseganchor{OLP-0656-B035}{end}

\olpseganchor{OLP-0656-B036}{start}
\begin{prob}
د \olref{lem:cut-adm-G3c} په ثبوت کښې د پنځم حالت \LR{E} هغه فرعي حالت وڅېړئ چې \LR{$!A$} د \LR{$\pi_1$} او \LR{$\pi_2$} په دواړو وروستيو استنتاجونو کښې اصلي فارمول وي او \LR{$!A \ident \lexists[x][!B(x)]$}۔
\end{prob}
\olpseganchor{OLP-0656-B036}{end}

\olpseganchor{OLP-0656-B037}{start}
پام وکړئ چې په پنځم حالت \LR{E} کښې د هغه ثبوت د پرېکون لوړوالے چې په \CutCS{} ختمېږي او د استقرا فرض پرې لګوو، د اصلي ثبوت د پرېکون له لوړوالے څخه لوړېدے شي۔ مثلاً، د \LR{$!A
\ident (!B \lif !C)$} په حالت کښې مو \LR{$\pi$} په داسې يوه ثبوت بدل کړ چې دوه \CutCS{} استنتاجونه لري:
\[
\AxiomC{}
\RightLabel{\LR{$\pi_2'$}}
\Deduce$\Gamma \fCenter \Delta, !B$
\AxiomC{}
\RightLabel{\LR{$\pi_1'$}}
\Deduce$!B, \Gamma \fCenter \Delta, !C$
\AxiomC{}
\RightLabel{\LR{$\pi_2''$}}
\Deduce$!C, \Gamma \fCenter \Delta$
\RightLabel{\CutCS}
\BinaryInf$!B, \Gamma \fCenter \Delta$
\RightSubproofLabel{\LR{$\pi_1''$}}
\RightLabel{\CutCS}
\BinaryInf$\Gamma \fCenter \Delta$
\DisplayProof
\]
د \LR{$\pi_1'$} او \LR{$\pi_2''$} تر منځ پاسنے \CutCS{} له~\LR{$\pi$} څخه هم د پرېکون ټيټه رتبه او هم ټيټ لوړوالے لري۔ خو ثبوت~\LR{$\pi_1''$}، چې د استقرا د فرض د لګولو پايله ده، نور له~\LR{$\pi$} څخه د پرېکون د ټيټ لوړوالے شرط نۀ تضمينوي۔ بيا هم څرنګه چې د لاندې \CutCS{} د پرېکون فارمول \LR{$!B$} دے، د پرېکون \emph{رتبه} ئې د~\LR{$\pi$} د پرېکون له رتبې څخه ټيټه ده۔
\olpseganchor{OLP-0656-B037}{end}

\olpseganchor{OLP-0656-B038}{start}
\olref{lem:cut-adm-G3c} ثابتوي چې له يوه ثبوت څخه د~\CutCS{} يو تر ټولو پاسنے استنتاج لرې کولے شو۔ د همدې لېمې په پرله‌پسې لګولو پر هغو تر ټولو پاسنيو فرعي ثبوتونه چې په~\CutCS ختمېږي، له يوه ثبوت څخه هر شمېر \CutCS{} استنتاجونه لرې کولے شو۔
\olpseganchor{OLP-0656-B038}{end}

\olpseganchor{OLP-0656-B039}{start}
\begin{thm}\ollabel{thm:cut-elim-G3c-topmost} په \LR{$ \Log{G3c} + \CutCS$} کښې هر ثبوت \LR{$\pi$} په~\Log{G3c} کښې په يوه ثبوت بدلولے شو۔
\end{thm}
\olpseganchor{OLP-0656-B039}{end}

\olpseganchor{OLP-0656-B040}{start}
\begin{proof}
په~\LR{$\pi$} کښې د \CutCS{} د استنتاجونو پر شمېر~\LR{$n$} استقرا کوو۔ که \LR{$n=0$} وي، څه کول نۀ دي پکار: \LR{$\pi$} لا په~\Log{G3c} کښې يو ثبوت دے۔ کنه، هغه فرعي ثبوت~\LR{$\pi'$} وګورئ چې د \CutCS{} په يو تر ټولو پاسني استنتاج ختمېږي۔ په وروستي استنتاج کښې يوه د \CutCS{} قاعده لري او بله هېڅ نۀ لري۔ د \olref{lem:cut-adm-G3c} په لګولو ئې په بې‌\CutCS{} ثبوت~\LR{$\pi''$} بدل کړئ، او فرعي ثبوت~\LR{$\pi'$} په~\LR{$\pi''$} بدل کړئ۔ پايله يو ثبوت دے چې \LR{$n-1$} د \CutCS{} استنتاجونه لري، او د استقرا فرض پرې لګېږي۔
\end{proof}
\olpseganchor{OLP-0656-B040}{end}

\olpseganchor{OLP-0655-B004}{start}
\olfileid{pt}{cut}{inv}
\olpseganchor{OLP-0655-B004}{end}

\olpseganchor{OLP-0655-B005}{start}
\olsection{د تر ټولو لوړې رتبې پرېکونونو لرې کول}
\olpseganchor{OLP-0655-B005}{end}

\olpseganchor{OLP-0655-B006}{start}
د \olref[seq][inv]{lem:G3c-invert} ثبوت ښودلې ده چې که \Log{G3c} د يوې منطقي قاعدې (له \RightR{\lexists} او \LeftR{\lforall} پرته) پايله ثابت کړي، نو د ثبوت د لوړوالے له زياتولو پرته د قاعدې هره مقدمه هم ثابت کوي۔ لا قوي پايله ښودلے شو: دا خاصيت د \LR{$\Log{G3c} +
\CutCS$} دپاره هم صدق کوي۔
\olpseganchor{OLP-0655-B006}{end}

\olpseganchor{OLP-0655-B007}{start}
لکه مخکې، د \CutCS{} د استنتاج \emph{د پرېکون رتبه} د هغه د پرېکون د فارمول ژوروالے ګڼو۔
\olpseganchor{OLP-0655-B007}{end}

\olpseganchor{OLP-0655-B008}{start}
\begin{lem}\ollabel{lem:inv-G3c-cut}
که \LR{$\Log{G3c} + \CutCS$} د يوې منطقي قاعدې~\LR{$R$} (له \RightR{\lexists} او \LeftR{\lforall} پرته، او د ځانګړي متغير له اصلي شرط سره) پايله د يوه ثبوت~\LR{$\pi$} په وسيله ثابت کړي، نو هره مقدمه هم د ثبوت~\LR{$\pi'$} په وسيله ثابت کوي (يا د ثبوتونه~\LR{$\pi'$}، \LR{$\pi''$} په وسيله، د دې له مخې چې قاعده يوه مقدمه لري که دوه)۔ سربېره پر دې،
(الف)~\LR{$\pheight{\pi'}\le\pheight{\pi}$} (او \LR{$\pheight{\pi''}\le\pheight{\pi}$})، او (ب)~په~\LR{$\pi'$} (او \LR{$\pi''$}) کښې د \CutCS{} د استنتاجونو شمېر په~\LR{$\pi$} کښې له شمېر څخه زيات نۀ دے؛ هر پاتې پرېکون د پخواني پرېکون په هماغه رتبه راځي۔
\end{lem}
\olpseganchor{OLP-0655-B008}{end}

\olpseganchor{OLP-0655-B009}{start}
\begin{proof}
د \cref{pt:seq:inv:lem:G3c-invert,pt:seq:inv:lem:invert-quant} ثبوتونه وګورئ۔ هغه ثبوت پر~\LR{$\pheight{\pi}$} استقرا وه۔ په بنسټيز حالت کښې مو يو بديهي اصل په بل بديهي اصل بدل کړ؛ نو (الف) او (ب) شرطونه څرګند دي۔ که د~\LR{$\pi$} وروستۍ قاعده~\LR{$R$} وي، نو د هغې بې‌واسطې فرعي ثبوتونه هماغه مطلوب ثبوتونه~\LR{$\pi_i$} دي، او (الف) او (ب) دواړه پوره کېږي۔ که وروستۍ قاعده~\LR{$R$} وي او د هماغې ټاکلې پېښې اصلي فارمول ولري، يو مطلوب ثبوت د هماغه استنتاج يوه فرعي څانګه ده۔ پرېکونونه د هماغې څانګې دي، نو شمېر ئې نۀ زياتېږي او هر پاتې پرېکون خپله رتبه ساتي۔ په نورو ټولو حالاتو کښې مو د~\LR{$\pi$} پر بې‌واسطې فرعي ثبوتونه، يعنې د وروستي استنتاج مقدمې يا مقدماتو ته رسېدونکو ثبوتونو، د استقرا فرض ولګاوه او بيا مو هماغه وروستۍ قاعده پر پايلو ولګوله۔ (الف) او (ب) د نوي ثبوت د بې‌واسطې فرعي ثبوتونه دپاره صدق کوي، نو د ټول ثبوت دپاره هم۔ هغه حالت لا پکار دے چې وروستے استنتاج~\CutCS وي۔ دلته بيا يوازې يو مثال ګورو: فرض کړئ چې \LR{$\pi$} د \LeftR{\land} د پايلې \LR{$!B \land !C, \Gamma \Sequent
\Delta$} يو ثبوت دے او په~\CutCS ختمېږي:
\[
\AxiomC{}
\RightLabel{\LR{$\pi_1$}}
\Deduce$!B \land !C, \Gamma \fCenter \Delta, !A$
\AxiomC{}
\RightLabel{\LR{$\pi_2$}}
\Deduce$!A, !B \land !C, \Gamma \fCenter \Delta$
\RightLabel{\CutCS}
\BinaryInf$!B \land !C, \Gamma \fCenter \Delta$
\DisplayProof
\]
د استقرا فرض پر \LR{$\pi_1$} او~\LR{$\pi_2$} لګېږي (دا هم د \LR{$\LeftR{\land}$} د پايلې د بېلګو ثبوتونه دي)، او په ترتيب سره د \LR{$!B, !C, \Gamma \fCenter
\Delta, !A$} او \LR{$!A, !B, !C, \Gamma \fCenter \Delta$} ثبوتونه \LR{$\pi_1'$} او~\LR{$\pi_2'$} راکوي۔ دا د \CutCS{} په وسيله يوځاے کوو او~\LR{$\pi'$} ترلاسه کوو:
\[
\AxiomC{}
\RightLabel{\LR{$\pi_1'$}}
\Deduce$!B, !C, \Gamma \fCenter \Delta, !A$
\AxiomC{}
\RightLabel{\LR{$\pi_2'$}}
\Deduce$!A, !B, !C, \Gamma \fCenter \Delta$
\RightLabel{\CutCS}
\BinaryInf$!B, !C, \Gamma \fCenter \Delta$
\DisplayProof
\]
څرګنده ده چې (الف) او (ب) پوره کېږي۔ په دې بيان کښې د معکوسولو د قاعدې او د ثبوت د وروستي استنتاج نومونه جلا وساتئ۔ په دې رسم کښې د معکوسېدو وړ عطف د پرېکون په ګډ سياق کښې دے۔ همدا کار د هرې معکوسېدونکې قاعدې لپاره کېږي۔ د دوو مقدمو د قاعدې د معکوسولو لپاره هره مطلوبه مقدمه جلا جوړوو۔ د کميتي قاعدې تازه ثابت د پرېکون له فارمول او له دواړو فرعي ثبوتونو څخه هم تازه ټاکو۔ که د مطلوبې مقدمې ثابت له مخکې ټاکل شوے وي، نو لومړے يو بل تازه ثابت کاروو۔ د نوې ونې د ځانګړو متغيرونو تر تازه نومولو وروسته هماغه ټاکل شوے ثابت د نوي ثابت په ځاے اخلو۔ دا تعويض د ثبوت لوړوالے او د پرېکون رتبې ساتي۔
\end{proof}
\olpseganchor{OLP-0655-B009}{end}

\olpseganchor{OLP-0655-B010}{start}
پام وکړئ چې که د \CutCS پر ځاے \Cut{} وکاروو، دا استدلال نۀ چلېږي؛ ځکه د وروستي ثبوت په پايله کښې به په مقدم کښې د \LR{$!B, !C, \Gamma$} دوه نقلونه او په تالي کښې د \LR{$\Delta$} دوه نقلونه وي۔ بيا به د انقباض د منلو وړتيا وکاروو، خو د \olref[seq][inv]{prop:G3c-cont-adm} ثبوت د قاعدو د معکوسولو لېمه کاروي؛ نو دا به دوراني استدلال شي۔
\olpseganchor{OLP-0655-B010}{end}

\olpseganchor{OLP-0655-B011}{start}
د \cref{pt:cut:inv:lem:inv-G3c-cut} په وسيله د پرېکون د لرې کولو يو بديل ثبوت ورکولے شو۔ په دې ثبوت کښې د \CutCS{} تر ټولو پاسنۍ پېښې يو په بل پسې نۀ لرې کوو؛ د \CutCS{} د تر ټولو لوړې رتبې استنتاجونه لرې کوو۔ يعنې په \LR{$\Log{G3c} + \CutCS$} کښې له يوه ثبوت~\LR{$\pi$} څخه پيل کوو او پکښې هر چېرې يو پرېکون لرې کوو (ښايي د ټيټو رتبو په پرېکونونو ئې بدل کړو)، بيا همدا کار تر هغې کوو چې هېڅ پرېکون پاتې نۀ شي۔ يوازې دا شرط دے چې له هغه پرېکون څخه پورته چې لرې کوو ئې، د هماغې يا لوړې رتبې پرېکون نۀ وي۔
\olpseganchor{OLP-0655-B011}{end}

\olpseganchor{OLP-0655-B012}{start}
\begin{lem}\ollabel{lem:max-cut-red-G3c}
که \LR{$\pi$} په \LR{$\Log{G3c}+\CutCS$} کښې يو ثبوت وي چې په~\LR{$n$} رتبې د~\CutCS{} په استنتاج ختمېږي، او په پاتې برخه کښې يوازې د~\LR{$\Log{G3c}$} قاعدې او د \LR{$<n$} رتبې \CutCS{} استنتاجونه کاروي، نو د هماغه وروستي سېکوېنټ يو ثبوت~\LR{$\pi'$} په~\LR{$\Log{G3c} + \CutCS$} کښې شته چې هر \CutCS{} استنتاج ئې د~\LR{$<n$} رتبې وي۔
\end{lem}
\olpseganchor{OLP-0655-B012}{end}

\olpseganchor{OLP-0655-B013}{start}
\begin{proof}
ثبوت~\LR{$\pi$} داسې ختمېږي:
\[
\AxiomC{}
\RightLabel{\LR{$\pi_1$}}
\Deduce$\Gamma \fCenter \Delta, !A$
\AxiomC{}
\RightLabel{\LR{$\pi_2$}}
\Deduce$!A, \Gamma \fCenter \Delta$
\RightLabel{\CutCS}
\BinaryInf$\Gamma \fCenter \Delta$
\DisplayProof
\]
کمزوري کول، د ځانګړي متغير بدلول او د ترم تعويض هم پر پرېکون لرونکو ثبوتونو لګېږي۔ د پرېکون په دواړو مقدمو کښې ګډ سياق يو شان بدلېږي، او تعويض د پرېکون د فارمول ژوروالے ساتي۔ د انقباض د منلو وړتيا هم د پرېکون د رتبې حد ساتي۔ د انقباض په ثبوت کښې د معکوسولو د پورته لېمې همدا بڼه کاروو۔ که وروستے استنتاج پرېکون وي، د ګډ سياق دواړه ټاکلي نقلونه په دواړو مقدمو کښې راکموو او پر هماغه پرېکون فارمول بيا پرېکون کوو۔
د~\LR{$!A$} د بڼې له مخې حالتونه بېلوو۔
\olpseganchor{OLP-0655-B013}{end}

\olpseganchor{OLP-0655-B014}{start}
فرض کړئ چې \LR{$!A$} اټومي دے۔ د دې شرط له مخې چې په~\LR{$\pi_1$} او \LR{$\pi_2$} کښې ټول \CutCS{} استنتاجونه د \LR{$<
\depth{!A} = 0$} رتبې دي، په~\LR{$\pi_1$} يا~\LR{$\pi_2$} کښې هېڅ \CutCS{} استنتاج نۀ شي کېدے۔ پر \LR{$\pheight{\pi_2}$} د استقرا په وسيله ښيو چې \LR{$\Gamma \Sequent \Delta$} بې‌پرېکونه ثبوت لري۔ که \LR{$\pheight{\pi_2} = 0$} وي، نو \LR{$!A, \Gamma \Sequent \Delta$} بديهي اصل دے۔ که بديهي اصل د دروغو له کبله وي، يا \LR{$\lfalse \in \Gamma$} او مطلوب سېکوېنټ پخپله بديهي اصل دے، يا \LR{$!A \ident \lfalse$} دے۔ په دويم صورت کښې د~\LR{$\pi_1$} له تالي څخه د~\LR{$\lfalse$} ټاکلې پېښه او په پاسنيو سياقونو کښې ئې اړوندې پېښې لرې کړئ۔ دروغ په ښي اړخ کښې د کومې قاعدې اصلي فارمول نۀ دے۔ که د هويت په کوم بديهي اصل کښې همدغه پېښه وکارېږي، په مقدم کښې دروغ پاتې دي او د دروغو بديهي اصل لګېږي۔ نو بې‌پرېکونه ثبوت پاتې کېږي۔ د هويت په بديهي اصل کښې که \LR{$!A$} اصلي فارمول نۀ وي، نو کوم اټومي \LR{$!C$} هم د~\LR{$\Gamma$} او هم د \LR{$\Delta$} يو غړے دے، او \LR{$\Gamma \Sequent \Delta$} پخپله بديهي اصل دے۔ که \LR{$!A$} اصلي فارمول وي، نو \LR{$\Delta = \Delta', !A$}۔ په دې حالت کښې \LR{$\pi_1$} په \LR{$\Gamma \Sequent \Delta', !A, !A$} ختمېږي۔ د \olref[seq][inv]{prop:G3c-cont-adm} (د انقباض د منلو وړتيا) په لګولو د \LR{$\Gamma \Sequent \Delta', !A$} بې‌پرېکونه ثبوت ترلاسه کوو۔ که \LR{$\pheight{\pi_2} > 0$} وي، په يوه قاعده~\LR{$R$} ختمېږي۔ ځانګړي متغيرونه تازه نوموو۔ د دې قاعدې يو فرعي ثبوت~\LR{$\rho$} واخلو۔ د مقدمې بڼه \LR{$!A, \Gamma', \Pi \Sequent \Delta', \Lambda$} ده، چې \LR{$\Pi$} او \LR{$\Lambda$} د~\LR{$R$} فعال فارمولونه دي۔ پر~\LR{$\pi_1$} او \LR{$\rho$} د \olref[seq][adm]{prop:weak-G3c-adm} په لګولو د \LR{$\Gamma, \Gamma', \Pi \Sequent \Delta, \Delta', \Lambda, !A$} او \LR{$!A, \Gamma, \Gamma', \Pi \Sequent \Delta, \Delta', \Lambda$} ثبوتونه ترلاسه کېږي، چې د استقرا فرض پرې لګېږي۔ د~\LR{$R$} د هرې مقدمې دپاره د \LR{$\Gamma, \Gamma', \Pi \Sequent \Delta, \Delta', \Lambda$} يو ثبوت راکوي۔ د \LR{$R$} په لګولو \LR{$\Gamma, \Gamma \Sequent \Delta, \Delta$} ترلاسه کېږي، او \olref[seq][inv]{prop:G3c-cont-adm} د \LR{$\Gamma \Sequent \Delta$} بې‌پرېکونه ثبوت راکوي۔
\olpseganchor{OLP-0655-B014}{end}

\olpseganchor{OLP-0655-B015}{start}
که \LR{$!A$} اټومي نۀ وي، نو د~\LR{$!A$} د اصلي عملګر له مخې حالتونه بېلوو۔ په هر حالت کښې \olref{lem:inv-G3c-cut} لګوو او د هغو کيڼو او ښيو قاعدو د مقدماتو ثبوتونه ترلاسه کوو چې \LR{$!A$} ئې اصلي فارمول دے او قاعده ئې معکوسېدونکې ده۔ بيا د \olref[top]{lem:cut-adm-G3c} د ثبوت د پنځم حالت په شان مخ ته ځو۔
\begin{enumerate}
  \item \LR{$!A \ident \lnot !B$}۔ ثبوتونه \LR{$\pi_1$} او \LR{$\pi_2$} په \LR{$\Gamma \Sequent \Delta, \lnot !B$} او \LR{$\lnot !B, \Gamma \Sequent \Delta$} ختمېږي۔ د \olref{lem:inv-G3c-cut} له مخې د \LR{$!B, \Gamma \Sequent \Delta$} او \LR{$\Gamma \Sequent \Delta, !B$} ثبوتونه \LR{$\pi_1'$} او \LR{$\pi_2'$} شته۔ ثبوت~\LR{$\pi'$} دا دے:
\[
  \AxiomC{}
  \RightLabel{\LR{$\pi_2'$}}
  \Deduce$\Gamma \fCenter \Delta, !B$
  \AxiomC{}
  \RightLabel{\LR{$\pi_1'$}}
  \Deduce$!B, \Gamma \fCenter \Delta$
  \RightLabel{\CutCS}
  \BinaryInf$\Gamma \fCenter \Delta$
  \DisplayProof
  \]
  \item \LR{$!A \ident (!B \lor !C)$}۔ ثبوتونه \LR{$\pi_1$} او \LR{$\pi_2$} په \LR{$\Gamma \Sequent \Delta, !B \lor !C$} او \LR{$!B \lor !C, \Gamma \Sequent \Delta$} ختمېږي۔ د \olref{lem:inv-G3c-cut} له مخې (پر~\LR{$\pi_1$} د \RightR{\lor} معکوسول) د \LR{$\Gamma \Sequent \Delta, !B, !C$} يو ثبوت \LR{$\pi_1'$} شته۔ د \olref{lem:inv-G3c-cut} له مخې (پر~\LR{$\pi_2$} د \LeftR{\lor} معکوسول) د \LR{$!B, \Gamma \Sequent \Delta$} يو ثبوت \LR{$\pi_2'$} او د \LR{$!C, \Gamma \Sequent \Delta$} يو ثبوت \LR{$\pi_2''$} شته۔ پر \LR{$\pi_2''$} د \olref[seq][adm]{prop:weak-G3c-adm} په لګولو د \LR{$!C, \Gamma \Sequent \Delta, !B$} يو ثبوت~\LR{$\pi_2'''$} هم ترلاسه کوو۔ ثبوت~\LR{$\pi'$} دا دے:
\[
  \AxiomC{}
  \RightLabel{\LR{$\pi_1'$}}
  \Deduce$\Gamma \fCenter \Delta, !B, !C$
  \AxiomC{}
  \RightLabel{\LR{$\pi_2'''$}}
  \Deduce$!C, \Gamma \fCenter \Delta, !B$
  \RightLabel{\CutCS}
  \BinaryInf$\Gamma \fCenter \Delta, !B$
  \AxiomC{}
  \RightLabel{\LR{$\pi_2'$}}
  \Deduce$!B, \Gamma \fCenter \Delta$
  \RightLabel{\CutCS}
  \BinaryInf$\Gamma \fCenter \Delta$
  \DisplayProof
  \]
  \item \LR{$!A \ident (!B \land !C)$}۔ د~\LR{$\pi_1$} د ښي عطف په معکوسولو د \LR{$\Gamma \Sequent \Delta, !B$} او \LR{$\Gamma \Sequent \Delta, !C$} ثبوتونه ترلاسه کوو۔ د~\LR{$\pi_2$} د کيڼ عطف په معکوسولو د \LR{$!B, !C, \Gamma \Sequent \Delta$} ثبوت ترلاسه کوو۔ لومړے د ښي عطف د لومړۍ مقدمې ثبوت په کيڼ اړخ په~\LR{$!C$} کمزورے کړئ او پر~\LR{$!B$} پرېکون وکړئ۔ پايله \LR{$!C, \Gamma \Sequent \Delta$} ده۔ بيا ئې د ښي عطف له دويمې مقدمې سره پر~\LR{$!C$} پرې کړئ۔ دواړه پرېکونونه له~\LR{$!B \land !C$} څخه ټيټه رتبه لري۔
  \item \LR{$!A \ident (!B \lif !C)$}۔ ثبوتونه \LR{$\pi_1$} او \LR{$\pi_2$} په \LR{$\Gamma \Sequent \Delta, !B \lif !C$} او \LR{$!B \lif !C, \Gamma \Sequent \Delta$} ختمېږي۔ د \olref{lem:inv-G3c-cut} له مخې (پر~\LR{$\pi_1$} د \RightR{\lif} معکوسول) د \LR{$!B, \Gamma \Sequent \Delta, !C$} يو ثبوت \LR{$\pi_1'$} شته۔ د \olref{lem:inv-G3c-cut} له مخې (پر~\LR{$\pi_2$} د \LeftR{\lif} معکوسول) د \LR{$\Gamma \Sequent \Delta, !B$} يو ثبوت \LR{$\pi_2'$} او د \LR{$!C, \Gamma \Sequent \Delta$} يو ثبوت \LR{$\pi_2''$} شته۔ پر \LR{$\pi_2''$} د \olref[seq][adm]{prop:weak-G3c-adm} په لګولو د \LR{$!C, !B, \Gamma \Sequent \Delta$} يو ثبوت~\LR{$\pi_2'''$} هم ترلاسه کوو۔ ثبوت~\LR{$\pi'$} دا دے:
\[
  \AxiomC{}
  \RightLabel{\LR{$\pi_2'$}}
  \Deduce$\Gamma \fCenter \Delta, !B$
  \AxiomC{}
  \RightLabel{\LR{$\pi_1'$}}
  \Deduce$!B, \Gamma \fCenter \Delta, !C$
  \AxiomC{}
  \RightLabel{\LR{$\pi_2'''$}}
  \Deduce$!C, !B, \Gamma \fCenter \Delta$
  \RightLabel{\CutCS}
  \BinaryInf$!B, \Gamma \fCenter \Delta$
  \RightLabel{\CutCS}
  \BinaryInf$\Gamma \fCenter \Delta$
  \DisplayProof
  \]
  \item \LR{$!A \ident \lforall[x][!B(x)]$}۔ ثبوتونه \LR{$\pi_1$} او \LR{$\pi_2$} په \LR{$\Gamma \Sequent \Delta, \lforall[x][!B(x)]$} او \LR{$\lforall[x][!B(x)], \Gamma \Sequent \Delta$} ختمېږي۔ د \olref{lem:inv-G3c-cut} له مخې د \LR{$\Gamma \Sequent \Delta, !B(c)$} يو ثبوت~\LR{$\pi_1'$} شته۔ \LR{$c$} د دواړو اصلي ثبوتونو له ثابتونو څخه تازه ټاکو۔
\olpseganchor{OLP-0655-B015}{end}

\olpseganchor{OLP-0655-B016}{start}
اوس ثبوت~\LR{$\pi_2$} وګورئ۔ په \LR{$\lforall[x][!B(x)], \Gamma \Sequent \Delta$} ختمېږي۔ د~\LR{$\pi_2$} له وروستي سېکوېنټ څخه پورته په تلو، په کيڼ اړخ کښې د \LR{$\lforall[x][!B(x)]$} هره هغه پېښه په نښه کړئ چې د~\LR{$\pi_2$} په وروستي سېکوېنټ کښې د \LR{$\lforall[x][!B(x)]$} پېښې ته «رسېږي»؛ يعنې:
(الف)~د~\LR{$\pi_2$} په وروستي سېکوېنټ کښې \LR{$\lforall[x][!B(x)]$} په نښه کړئ۔ (ب)~که \LR{$\lforall[x][!B(x)]$} د يوې قاعدې د پايلې په سياق کښې وي او په نښه شوے وي، نو د قاعدې په مقدمه يا مقدماتو کښې ئې اړونده پېښه هم په نښه کړئ۔ (ج)~که \LR{$\lforall[x][!B(x)]$} د \LeftR{\lforall} قاعدې په پايله کښې اصلي فارمول وي او په نښه شوے وي، نو په مقدمه کښې د \LR{$\lforall[x][!B(x)]$} او \LR{$!B(t)$} اړوندې فعالې پېښې په نښه کړئ۔
\olpseganchor{OLP-0655-B016}{end}

\olpseganchor{OLP-0655-B017}{start}
اوس د \LR{$\lforall[x][!B(x)]$} ټولې په نښه شوې پېښې وګورئ۔ هېڅ يوه له \LeftR{\lforall} پرته په بله قاعده کښې فعاله نۀ ده (يوازې د \LeftR{\lforall} فعال فارمولونه په نښه کېږي)۔ په \LR{$\pi_2$} کښې د \LR{$\lforall[x][!B(x)]$} ټولې په نښه شوې پېښې لرې کړئ۔ که دا سم ثبوت وي، نو د \LR{$\lforall[x][!B(x)]$} په نښه شوې پېښې هېڅکله فعالې نۀ وې؛ ټولې نېغ له بديهي اصولو راغلې وې۔ ثبوت په~\LR{$\Gamma \Sequent \Delta$} ختمېږي، او \LR{$\pi$} په همدې بدلولے شو۔ د تر ټولو لوړې رتبې يو پرېکون مو لرې کړ۔
\olpseganchor{OLP-0655-B017}{end}

\olpseganchor{OLP-0655-B018}{start}
کنه، ترلاسه شوے شي سم ثبوت نۀ دے، ځکه د \LR{$\lforall[x][!B(x)]$} هغه فارمولونه مو لرې کړي چې د \LeftR{\lforall} په ځينو استنتاجونو کښې فعال وو۔ دغه استنتاجونه مو په دې بڼه «استنتاجونو» بدل کړي دي:
\[
  \AxiomC{}
  \RightLabel{\LR{$\pi_t$}}
  \Deduce$!B(t), \Pi \fCenter \Lambda$
  \RightLabel{\LeftR{\lforall}}
  \UnaryInf$\Pi \fCenter \Lambda$
  \DisplayProof
  \]
له پاسه تر ښکته د ثبوت ونې ته راځو۔ په هره څانګه کښې د زياتېدونکو سياقونو له نښو څخه ځانګړي متغيرونه تازه کوو۔ تر تعويض مخکې د~\LR{$t$} له نښو څخه د~\LR{$\pi_1'$} ځانګړي متغيرونه هم تازه کوو۔ په هرې قاعدې کښې يو شان اضافي سياقونه ساتو او هر داسې «استنتاج» په دې بدلوو:
\[
  \AxiomC{}
  \RightLabel{\LR{$\Subst{\pi_1'}{t}{c}[\Pi \Sequent \Lambda]$}}
  \Deduce$\Gamma, \Pi \fCenter \Delta, \Lambda, !B(t)$
  \AxiomC{}
  \RightLabel{\LR{$\pi_t[\Gamma \Sequent \Delta]$}}
  \Deduce$!B(t), \Gamma, \Pi \fCenter \Delta, \Lambda$
  \RightLabel{\CutCS}
  \BinaryInf$\Gamma, \Pi \fCenter \Delta, \Lambda$
  \DisplayProof
  \]
د ونې په هر سېکوېنټ کښې د~\LR{$\Gamma$} او~\LR{$\Delta$} يو يو نقل ورزياتېږي۔ په رسم کښې ښے فرعي ثبوت د پخوانۍ مقدمې له پاسه تر ښکته رغول شوې بڼه ده، چې اضافي سياق هم لري۔ دا بدلون د \LeftR{\lforall} په هر هغه ځاے کښې کوو چې اصلي پېښه ئې په نښه وي۔ نورې قاعدې خپل ګډ سياق ساتي۔ پاييز سېکوېنټ \LR{$\Gamma, \Gamma \Sequent \Delta, \Delta$} دے۔ د انقباض د منلو وړتيا ولګوئ، چې د~\LR{$\Gamma \Sequent \Delta$} يو ثبوت شي، او \LR{$\pi$} پرې بدل کړئ۔ پر \LR{$\lforall[x][!B(x)]$} يو پرېکون مو پر \LR{$!B(t)$} په (ښايي ډېرو) پرېکونونو بدل کړ، خو دا ټول د ټيټې رتبې دي۔ په \LR{$\pi_1$} يا \LR{$\pi_2$} کښې پخوا موجود پرېکونونه هم پاتې وي، خو هغه هم ټول د ټيټې رتبې دي۔
  \item \LR{$!A \ident \lexists[x][!B(x)]$}۔ د~\LR{$\pi_2$} په پای کښې د کيڼ وجودي کميت په معکوسولو د \LR{$!B(c), \Gamma \Sequent \Delta$} يو ثبوت~\LR{$\pi_2'$} واخلو۔ \LR{$c$} په دواړو اصلي ثبوتونو کښې تازه ثابت اخلو۔ د~\LR{$\pi_1$} په تالي کښې هغه پېښې په نښه کوو چې د پای د ټاکلي وجودي فارمول اړوندې دي۔ نښه د هرې قاعدې ګډ سياق ته، او د \RightR{\lexists} د اصلي فارمول له پېښې څخه د هماغه وجودي فارمول د مقدمې پېښې ته رسېږي۔ په ټول ثبوت کښې په نښه شوې وجودي پېښې لرې کوو او د هر سېکوېنټ مقدم ته~\LR{$\Gamma$} او تالي ته~\LR{$\Delta$} يو ځل ورزياتوو۔ د هويت او دروغو بديهي اصول او نورې قاعدې پر خپل ځاے لګېږي۔ د ځانګړو متغيرونو نومونه له مخکې بدلوو، څو د ورزياتو شوو سياقونو له نښو سره ګډ نۀ شي۔
  د~\RightR{\lexists} په هر هغه استنتاج کښې چې اصلي پېښه ئې په نښه شوې وي، بدله مقدمه \LR{$\Gamma, \Pi \Sequent \Delta, \Lambda, !B(t)$} ده۔ په~\LR{$\pi_2'$} کښې د~\LR{$c$} پر ځاے~\LR{$t$} واخلو، وروسته له هغې چې ځانګړي متغيرونه مو د~\LR{$t$} له نښو بېل کړي وي۔ د~\LR{$\Pi$} او~\LR{$\Lambda$} په ورزياتولو د \LR{$!B(t), \Gamma, \Pi \Sequent \Delta, \Lambda$} ثبوت تر لاسه کېږي۔ د ګډ سياق پرېکون پر~\LR{$!B(t)$} د \LR{$\Gamma, \Pi \Sequent \Delta, \Lambda$} ثبوت راکوي، چې د پخواني استنتاج بدله پايله ده۔ د دې ثبوت پاييز سېکوېنټ دوه ځله اصلي سياقونه لري۔ انقباض پرې لګوو او \LR{$\Gamma \Sequent \Delta$} اخلو۔ تعويض او کمزوري کول د پرېکونونو رتبه نۀ زياتوي، او انقباض هم د رتبې همدا حد ساتي۔ پر~\LR{$!B(t)$} هر نوے پرېکون له اصلي وجودي فارمول څخه ټيټه رتبه لري۔
\end{enumerate}
\end{proof}
\olpseganchor{OLP-0655-B018}{end}

\olpseganchor{OLP-0655-B019}{start}
\begin{prob}
د \olref{lem:max-cut-red-G3c} د ثبوت هغه حالت وڅېړئ چې \LR{$!A \ident (!B \land !C)$} وي۔ يعنې وښيئ چې که يو ثبوت~\LR{$\pi$} داسې ختمېږي:
\[
\AxiomC{}
\RightLabel{\LR{$\pi_1$}}
\Deduce$\Gamma \fCenter \Delta, !B \land !C$
\AxiomC{}
\RightLabel{\LR{$\pi_2$}}
\Deduce$!B \land !C, \Gamma \fCenter \Delta$
\RightLabel{\CutCS}
\BinaryInf$\Gamma \fCenter \Delta$
\DisplayProof
\]
او \LR{$\pi_1$} او \LR{$\pi_2$} يوازې د \LR{$<\depth{!B \land !C}$} رتبې پرېکونونه لري، نو د \LR{$\Gamma \Sequent \Delta$} يو ثبوت~\LR{$\pi'$} شته چې يوازې د \LR{$<\depth{!B \land !C}$} رتبې پرېکونونه لري۔
\end{prob}
\olpseganchor{OLP-0655-B019}{end}

\olpseganchor{OLP-0655-B020}{start}
% \begin{prob}
% Verify the case where $!A \ident \lexists[x][!B(x)]$ in the proof of
% \olref{lem:inv-G3c-cut}.
% \end{prob}
\olpseganchor{OLP-0655-B020}{end}

\olpseganchor{OLP-0655-B021}{start}
\olref{lem:max-cut-red-G3c} ثابتوي چې په يوه ثبوت کښې د تر ټولو لوړې رتبې د \CutCS{} استنتاج د ټيټو رتبو د \CutCS{} په استنتاجونو بدلولے شو۔ بيا د همدې لېمې په وسيله ښودلے شو چې له يوه ثبوت څخه هر شمېر \CutCS{} استنتاجونه لرې کولے شو: لېمه پر هغو تر ټولو پاسنيو فرعي ثبوتونه يو په بل پسې لګوو چې د تر ټولو لوړې رتبې په~\CutCS ختمېږي۔
\olpseganchor{OLP-0655-B021}{end}

\olpseganchor{OLP-0655-B022}{start}
\begin{thm}\ollabel{thm:cut-elim-G3c-largest}
په \LR{$\Log{G3c} + \CutCS$} کښې هر ثبوت \LR{$\pi$} په~\Log{G3c} کښې په يوه ثبوت بدلولے شو۔
\end{thm}
\olpseganchor{OLP-0655-B022}{end}

\olpseganchor{OLP-0655-B023}{start}
\begin{proof}
لومړے ثابتوو چې په~\LR{$\pi$} کښې د تر ټولو لوړې رتبې ټول پرېکونونه لرې کولے شو۔ \LR{$d$} دې په~\LR{$\pi$} کښې د پرېکونونو تر ټولو لوړه رتبه وي، او \LR{$n$} دې د~\LR{$d$} رتبې د پرېکونونو شمېر وي۔ ثبوت پر~\LR{$n$} استقرا ده۔ که \LR{$n=0$} وي، د ثابتولو څه نشته۔ که \LR{$n>0$} وي، د~\LR{$d$} رتبې يو تر ټولو پاسنے پرېکون واخلو، او \LR{$\pi_1$} دې هغه فرعي ثبوت وي چې په همدې ختمېږي۔ د \olref{lem:max-cut-red-G3c} له مخې د هماغه وروستي سېکوېنټ يو ثبوت~\LR{$\pi_1'$} شته چې ټول پرېکونونه ئې د \LR{$<d$} رتبې دي۔ په~\LR{$\pi$} کښې \LR{$\pi_1$} په \LR{$\pi_1'$} بدل کړو۔ يو ثبوت ترلاسه کېږي چې د~\LR{$d$} رتبې \LR{$n-1$} پرېکونونه لري، او د استقرا فرض پرې لګېږي۔
\olpseganchor{OLP-0655-B023}{end}

\olpseganchor{OLP-0655-B024}{start}
اوس پايله پر~\LR{$d$} د استقرا له مخې راځي۔ که \LR{$d=0$} وي، ټول پرېکونونه اټومي دي، او د مخکينۍ پايلې له مخې ټول لرې کېدے شي۔ که \LR{$d>0$} وي، مخکينۍ پايله داسې ثبوت راکوي چې د \LR{$d$} رتبې ټول پرېکونونه ئې د~\LR{$<d$} رتبې په پرېکونونو بدل شوي دي۔ څرنګه چې \LR{$d$} په~\LR{$\pi$} کښې د پرېکونونو تر ټولو لوړه رتبه وه، اوس داسې ثبوت لرو چې د پرېکونونو تر ټولو لوړه رتبه ئې \LR{$<d$} ده، او د استقرا فرض پرې لګېږي۔
\end{proof}
\olpseganchor{OLP-0655-B024}{end}

\olfileid{pt}{cut}{implication-alternate}
\paragraph{د شرط د پرېکون بله بڼه}
\olpseganchor{OLP-0660-B001}{start}
که \LR{$!A \ident (!B \lif !C)$} وي، نو \LR{$\pi$} په دې ختمېږي:
\[
\AxiomC{}
\RightLabel{\LR{$\pi_1'$}}
\Deduce$!B, \Gamma \fCenter \Delta, !C$
\RightLabel{\RightR{\lif}}
\UnaryInf$\Gamma \fCenter \Delta, !B \lif !C$
\LeftSubproofLabel{\LR{$\pi_1$}}
\AxiomC{}
\RightLabel{\LR{$\pi_2'$}}
\Deduce$!B \lif !C, \Pi \fCenter \Lambda, !B$
\AxiomC{}
\RightLabel{\LR{$\pi_2''$}}
\Deduce$!C, \Pi \fCenter \Lambda$
\RightLabel{\LeftR{\lif}}
\BinaryInf$!B \lif !C, \Pi \fCenter \Lambda$
\RightSubproofLabel{\LR{$\pi_2$}}
\RightLabel{\Cut}
\BinaryInf$\Gamma, \Pi \fCenter \Delta, \Lambda$
\DisplayProof
\]
په \Cut{} ختمېدونکے دا ثبوت
\[
\AxiomC{}
\RightLabel{\LR{$\pi_1'$}}
\Deduce$!B, \Gamma \fCenter \Delta, !C$
\RightLabel{\RightR{\lif}}
\UnaryInf$\Gamma \fCenter \Delta, !B \lif !C$
\LeftSubproofLabel{\LR{$\pi_1$}}
\AxiomC{}
\RightLabel{\LR{$\pi_2'$}}
\Deduce$!B \lif !C, \Pi \fCenter \Lambda, !B$
\RightLabel{\Cut}
\BinaryInf$\Gamma, \Pi \fCenter \Delta, \Lambda, !B$
\DisplayProof
\]
د~\LR{$\pi$} په پرتله د پرېکون لږ لوړوالے لري۔ استقرايي فرضيه
د \LR{$\Gamma, \Pi \fCenter \Delta, \Lambda, !B$} يو ثبوت~\LR{$\pi_1''$} راکوي۔
په \Cut{} ختمېدونکے دا ثبوت
\[
\AxiomC{}
\RightLabel{\LR{$\pi_1'$}}
\Deduce$!B, \Gamma \fCenter \Delta, !C$
\AxiomC{}
\RightLabel{\LR{$\pi_2''$}}
\Deduce$!C, \Pi \fCenter \Lambda$
\RightLabel{\Cut}
\BinaryInf$!B, \Gamma, \Pi \fCenter \Delta, \Lambda$
\DisplayProof
\]
د دې پرېکون رتبه او د پرېکون لوړوالے دواړه د~\LR{$\pi$} د
پرېکون له رتبې او لوړوالے څخه لږ دي، نو استقرايي فرضيه د
\LR{$!B, \Gamma, \Pi, \Pi \fCenter \Delta, \Lambda,
\Lambda$} يو \Log{G3c}-ثبوت~\LR{$\pi_2'''$} راکوي۔ اوس استقرايي
فرضيه په \Cut{} ختمېدونکي دې ثبوت باندې وکاروئ:
\[
\AxiomC{}
\RightLabel{\LR{$\pi_1''$}}
\Deduce$\Gamma, \Pi \fCenter \Delta, \Lambda, !B$
\AxiomC{}
\RightLabel{\LR{$\pi_2'''$}}
\Deduce$!B, \Gamma, \Pi, \Pi \fCenter \Delta, \Lambda,
\Lambda$
\RightLabel{\Cut}
\BinaryInf$\Gamma, \Gamma, \Pi, \Pi, \Pi \fCenter \Delta, \Delta, \Lambda, \Lambda, \Lambda$
\DisplayProof
\]
په دې عام پرېکون کښې د پرېکون فارمول لرې کېږي، او د دواړو مقدمو
پاتې سياقونه له خپلو نقلونو سره يوځاے کېږي۔
(د دې د پرېکون رتبه \LR{$< \cutr{\pi}$} ده۔) په دې سره د \LR{$\Gamma,
\Gamma, \Pi, \Pi, \Pi \fCenter \Delta, \Delta, \Lambda, \Lambda, \Lambda$}
يو ثبوت تر لاسه کېږي۔ د \olref[seq][inv]{prop:G3c-cont-adm}
په کارولو د \LR{$\Gamma, \Pi \fCenter \Delta, \Lambda$} ثبوت~\LR{$\pi'$} تر لاسه کوو۔
\olpseganchor{OLP-0660-B001}{end}

\olpseganchor{OLP-0661-B004}{start}
\olfileid{pt}{cut}{mid}
\olpseganchor{OLP-0661-B004}{end}

\olpseganchor{OLP-0661-B005}{start}
\olsection{د منځني سېکوېنټ قضيه او د هېربراند قضيه}
\olpseganchor{OLP-0661-B005}{end}

\olpseganchor{OLP-0661-B006}{start}
د پرېکون لرې کولو د قضيې يوه مهمه پايله د منځني سېکوېنټ
قضيه ده۔ دا وايي: که يو ثبوت~\LR{$\pi$} د~\LR{$\Gamma \Sequent \Delta$}
وي چې پکښې هر فارمول په مقدم کميتي بڼه وي، نو يو ثبوت~\LR{$\pi'$}
شته چې يو سېکوېنټ \LR{$S \ident \Pi \Sequent
\Lambda$} (\emph{منځنے سېکوېنټ}) لري، داسې چې په~\LR{$\pi'$} کښې د~\LR{$S$}
پاس ټول استنتاجونه گزاره‌يي او د~\LR{$S$} لاندې ټول استنتاجونه د کميت
اچوونکو وي۔ (فارمول~\LR{$!A$} ته \emph{په مقدم کميتي بڼه}
وايي که دا بڼه ولري:
\[\mathsf{Q}_1x_1\dots\mathsf{Q}_nx_n\,!B(x_1,\dots,x_n)\]
چې هر \LR{$\mathsf{Q}_i$} پکښې يا \LR{$\lforall$} يا~\LR{$\lexists$} وي او وروستنے
متن له کميت اچوونکو خالي وي۔) د منځني سېکوېنټ د قضيې يوه مهمه پايله د هېربراند
قضيه ده۔ عام حالت ئې بيانول لږ سخت دي، خو په پراخه توګه دا بڼه
لري: د لومړۍ درجې منطق د هرې د ثابتولو وړ تړلے فارمول~\LR{$!A$} دپاره
يوه گزاره‌يي توتولوژي~\LR{$!A'$} شته چې ترې \LR{$!A$} ثابت کېدے شي۔
دا ئې د هغو تړلي فارمولونه دپاره بڼه ده چې د~\LR{$\lexists[x][!B(x)]$}
په شکل وي، چې \LR{$!B$} له کميت اچوونکو خالي وي:
\olpseganchor{OLP-0661-B006}{end}


\olpseganchor{OLP-0661-B007}{start}
\begin{thm}[د هېربراند قضيه]
فرض کړئ چې \LR{$!B(x)$} له کميت اچوونکو خالي دے او \LR{$\Proves 
\lexists[x][!B(x)]$}۔ نو ترمونه~\LR{$t_1$}، \dots، \LR{$t_n$} شته داسې چې
\LR{$\Proves !B(t_1) \lor \dots \lor !B(t_n)$}۔
\end{thm}
\olpseganchor{OLP-0661-B007}{end}

\olpseganchor{OLP-0661-B008}{start}
فارمول~\LR{$!B(t_1) \lor \dots \lor !B(t_n)$} ته د~\LR{$\lexists[x][!B(x)]$}
\emph{د هېربراند فصل} وايي۔ څرنګه چې \LR{$!B$} له کميت اچوونکو خالي
دے، دا فصل هم له هغو خالي دے۔ نو که دا د ثابتولو وړ وي، په
بې‌پرېکونه ثبوت د ثابتولو وړ دے، او په نتيجه کښې په داسې
ثبوت چې يوازې گزاره‌يي استنتاجونه کاروي۔ ځکه نو توتولوژي ده۔
\olpseganchor{OLP-0661-B008}{end}

\olpseganchor{OLP-0661-B009}{start}
د هېربراند د قضيې سخته برخه دا ښودل دي چې د هر د ثابتولو وړ
فارمول~\LR{$\lexists[x][!B(x)]$} دپاره د هېربراند يو فصل شته۔
معکوس ئې، يعنې که \LR{$\lexists[x][!B(x)]$} د هېربراند فصل ولري نو
ثبوت لري، اسان دے۔ په حقيقت کښې د هېربراند فصل، \LR{$\lexists[x][!B(x)]$}
لازموي۔ مثلاً په~\LR{$\Log{G3c}$} کښې د \LR{$n=2$} حالت دا ثبوت دے:
\[
\Axiom$!B(t_1) \fCenter \lexists[x][!B(x)], !B(t_1), !B(t_2)$
\Axiom$!B(t_2) \fCenter \lexists[x][!B(x)], !B(t_1), !B(t_2)$
\RightLabel{\LeftR{\lor}}
\BinaryInf$!B(t_1) \lor !B(t_2) \fCenter \lexists[x][!B(x)], !B(t_1), !B(t_2)$
\RightLabel{\RightR{\lexists}}
\UnaryInf$!B(t_1) \lor !B(t_2) \fCenter \lexists[x][!B(x)], !B(t_2)$
\RightLabel{\RightR{\lexists}}
\UnaryInf$!B(t_1) \lor !B(t_2) \fCenter \lexists[x][!B(x)]$
\DisplayProof
\]
\olpseganchor{OLP-0661-B009}{end}

\olpseganchor{OLP-0661-B010}{start}
د \LR{$\Sequent
\lexists[x][!B(x)]$} له ثبوت~\LR{$\pi$} څخه د هېربراند د فصل
راايستلو دپاره هغه حالت وګورئ چې په يوه بې‌پرېکونه ثبوت~\LR{$\pi$}
کښې ټول \LR{$\RightR{\lexists}$} استنتاجونه په پای کښې وي:
\[
\AxiomC{}
\RightLabel{\LR{$\pi_1$}}
\Deduce$ \fCenter \lexists[x][!B(x)], !B(t_1), \dots, !B(t_n)$
\RightLabel{\RightR{\lexists}}
\UnaryInf$ \fCenter \lexists[x][!B(x)], !B(t_2), \dots, !B(t_n)$
\RightLabel{\LR{$(n-1) \times \RightR{\lexists}$}}
\Deduce$ \fCenter \lexists[x][!B(x)]$
\DisplayProof
\]
که په~\LR{$\pi$} کښې دا ټول هغه \LR{$\RightR{\lexists}$} استنتاجونه وي چې
\LR{$\lexists[x][!B(x)]$} پکښې فعال وي، نو په~\LR{$\pi_1$} کښې نور کميتي
استنتاجونه نشته، ځکه چې \LR{$!B(x)$}
نور کميت اچوونکي نۀ لري او پای سېکوېنټ په کيڼ اړخ کښې هېڅ
فارمولونه نۀ لري۔ په نورو ټکو، \LR{$ \fCenter \lexists[x][!B(x)], !B(t_1),
\dots, !B(t_n)$} د~\LR{$\pi$} منځنے سېکوېنټ دے۔ همدارنګه، د منځني
سېکوېنټ پاس کميتي استنتاجونه نشته، نو پاس ټول سېکوېنټونه
\LR{$\lexists[x][!B(x)]$} لري او دا نۀ فعال دي او نۀ اصلي۔ ځکه دا له
ثبوت~\LR{$\pi_1$} څخه لرې کولے شو او د~\LR{$\fCenter !B(t_1), \dots, !B(t_n)$}
يو ثبوت تر لاسه کوو؛ بيا د \LR{$\RightR{\lor}$} د \LR{$n-1$}
استعمالونو په وسيله د هېربراند د فصل~\LR{$\Sequent !B(t_1) \lor \dots \lor !B(t_n)$}
يو ثبوت تر لاسه کوو۔
\olpseganchor{OLP-0661-B010}{end}

\olpseganchor{OLP-0661-B011}{start}
ستونزه دا ده چې فرض کولے نۀ شو چې د \LR{$\Sequent \lexists[x][!B(x)]$}
يو بې‌پرېکونه ثبوت هم خپل ټول \RightR{\lexists} استنتاجونه
په پای کښې په يوه ډله کښې لري۔ د دې \RightR{\lexists} استنتاجونو
تر منځ داسې استنتاجونه کېدے شي چې ځينې \LR{$!B(t_i)$} پکښې اصلي وي۔
ځکه نو د منځني سېکوېنټ قضيې ته اړتيا ده۔
\olpseganchor{OLP-0661-B011}{end}

\olpseganchor{OLP-0661-B012}{start}
\begin{thm}[د منځني سېکوېنټ قضيه]\ollabel{thm:midsequent}
که د \LR{$\Gamma \Sequent \Delta$} يو \Log{G3c}-ثبوت~\LR{$\pi$} وي، چې
په \LR{$\Gamma$} او \LR{$\Delta$} کښې ټول فارمولونه په مقدم کميتي بڼه
وي، نو يو ثبوت~\LR{$\pi'$} شته چې سېکوېنټ \LR{$\Pi \Sequent \Lambda$}
لري، داسې چې لاندې ترې ټول استنتاجونه کميتي او پاس ترې ټول
استنتاجونه گزاره‌يي وي۔
\end{thm}
\olpseganchor{OLP-0661-B012}{end}

\olpseganchor{OLP-0661-B013}{start}
\begin{proof}
د پرېکون لرې کولو په وسيله لومړے بې‌پرېکونه ثبوت اخلو۔ په~\LR{$\pi$}
کښې د يوه کميتي استنتاج \emph{ترتيبي شمېر} د لاندې گزاره‌يي
استنتاجونو شمېر تعريفوو، او د \LR{$\pi$} ترتيبي شمېر~\LR{$o(\pi)$} د ټولو
کميتي استنتاجونو د ترتيبي شمېرونو مجموعه ده۔ که \LR{$o(\pi) = 0$}
وي، نو د هېڅ کميتي استنتاج لاندې گزاره‌يي استنتاج نشته او کار
ختم دے: که کميتي استنتاج نۀ وي، پای سېکوېنټ منځنے سېکوېنټ
اخلو؛ که وي، د ټولو کميتي استنتاجونو يوه مقدمه ده، نو يو واحد
تر ټولو پاسنے کميتي استنتاج شته۔ د هغه مقدمه منځنے سېکوېنټ ده۔
\olpseganchor{OLP-0661-B013}{end}

\olpseganchor{OLP-0661-B014}{start}
که \LR{$o(\pi) > 0$} وي، د \LR{$>0$} ترتيبي شمېر لرونکو کميتي
استنتاجونو تر ټولو لاندې يو غوره کړئ۔ څرنګه چې ترتيبي شمېر
ئې \LR{$>0$} دے، لاندې ترې يو گزاره‌يي استنتاج ضرور شته۔ دا چې تر
ټولو لاندې يو دے، پايله ئې د بل کميتي استنتاج مقدمه کېدے نۀ
شي: ګنې د هغه دويم لاندې کميتي استنتاج لاندې به هم يو گزاره‌يي
استنتاج وے۔ د کميتي استنتاج او د لاندې گزاره‌يي استنتاج له
ډول سره سم (ډېر!) حالتونه بېلوو۔ څرنګه چې پای سېکوېنټ يوازې
د مقدم کميتي بڼې فارمولونه لري، د کميتي استنتاج اصلي فارمول
په ورپسې گزاره‌يي استنتاج کښې فعال کېدے نۀ شي۔ ګنې د هغه
استنتاج د اصلي فارمول په يوه گزاره‌يي رابط کښې به کميت اچوونکے
راشي۔ د فرعي فارمول د خاصيت له مخې به دا د پای سېکوېنټ
د يوې فارمول فرعي فارمول وے۔ نو د کميتي استنتاج
اصلي فارمول د لاندې گزاره‌يي استنتاج د سياق يو غړے دے۔
\olpseganchor{OLP-0661-B014}{end}

\olpseganchor{OLP-0661-B015}{start}
دا دوه بېلګې دي:
\olpseganchor{OLP-0661-B015}{end}

\olpseganchor{OLP-0661-B016}{start}
لومړے فرض کړئ چې يو \LR{$\RightR{\lforall}$} د \LR{$\LeftR{\lif}$} په
ښۍ مقدمه ختمېږي۔ نو د~\LR{$\pi$} اړوند فرعي ثبوت په فرعي
ثبوت~\LR{$\pi_1$} ختمېږي:
\[
\AxiomC{}
\RightLabel{\LR{$\pi_2$}}
\Deduce$\Gamma' \fCenter \Delta', \lforall[x][!A(x)], !B$
\AxiomC{}
\RightLabel{\LR{$\pi_3$}}
\Deduce$!C, \Gamma' \fCenter \Delta', !A(c)$
\RightLabel{\RightR{\lforall}}
\UnaryInf$!C, \Gamma' \fCenter \Delta', \lforall[x][!A(x)]$
\RightLabel{\LeftR{\lif}}
\BinaryInf$!B \lif !C, \Gamma' \fCenter \Delta', \lforall[x][!A(x)]$
\DisplayProof
\]
فرض کولے شو چې \LR{$\pi$} منظم دے، نو ځانګړے متغير~\LR{$c$} له~\LR{$\pi_3$}
بهر نۀ راځي؛ په ځانګړې توګه په \LR{$!B
\lif !C$} او په~\LR{$\pi_2$} کښې نۀ راځي۔ اوس ثبوت~\LR{$\pi_1'$} وګورئ:
\[
\AxiomC{}
\RightLabel{\LR{$\pi_2'$}}
\Deduce$\Gamma' \fCenter \Delta', \lforall[x][!A(x)], !A(c), !B$
\AxiomC{}
\RightLabel{\LR{$\pi_3$}}
\Deduce$!C, \Gamma' \fCenter \Delta', \lforall[x][!A(x)], !A(c)$
\RightLabel{\LeftR{\lif}}
\BinaryInf$!B \lif !C, \Gamma' \fCenter \Delta', \lforall[x][!A(x)], !A(c)$
\RightLabel{\RightR{\lforall}}
\UnaryInf$!B \lif !C, \Gamma' \fCenter \Delta', \lforall[x][!A(x)], \lforall[x][!A(x)]$
\DisplayProof
\]
چې \LR{$\pi_2'$} له~\LR{$\pi_2$} څخه په ښي اړخ کښې د~\LR{$!A(c)$} په
کمزوري کولو تر لاسه کېږي۔ (دا هېڅ استنتاج نۀ زياتوي، په
ځانګړې توګه داسې کميتي استنتاج نۀ زياتوي چې ترتيبي شمېر
بدلولے شي۔ په~\LR{$\pi_2$} کښې د ځانګړے متغير شرطونه هم نۀ ماتوي،
ځکه په~\LR{$!A(c)$} کښې هېڅ ثابت په~\LR{$\pi_2$} کښې د ځانګړے متغير په
توګه نۀ کارېږي۔) پر \RightR{\lforall} د ځانګړے متغير شرط لا هم
پوره دے، ځکه \LR{$c$} په~\LR{$!B \lif !C$} کښې نۀ راځي۔
\olpseganchor{OLP-0661-B016}{end}

\olpseganchor{OLP-0661-B017}{start}
د انقباض د منلو وړتيا کاروو چې د \LR{$!B \lif !C, \Gamma' \fCenter \Delta',
\lforall[x][!A(x)]$} يو ثبوت~\LR{$\pi_1''$} تر لاسه کړو۔ د
\LR{$\lforall[x][!A(x)]$} دپاره د \LR{$\RightR{\Contraction}$} د منلو
وړتيا ثبوت او په هغه کښې د \LR{$\RightR{\lforall}$} د معکوسولو وړتيا
ثبوت د هېڅ استنتاج ترتيبي شمېر نۀ بدلولو: تر ډېره ئې استنتاجونه
لرې کړي وو۔ نو په~\LR{$\pi_1''$} کښې کميتي استنتاجونه تر~\LR{$\pi_1$}
زيات نۀ دي، او په~\LR{$\pi_1''$} کښې د هر کميتي استنتاج لاندې د
گزاره‌يي استنتاجونو شمېر په~\LR{$\pi_1$} کښې د ورته استنتاج د
لاندې استنتاجونو هومره دے؛ پرته له وروستي \RightR{\lforall}:
په~\LR{$\pi$} کښې يو \LeftR{\lif} لاندې ترې و، چې په~\LR{$\pi_1''$} کښې
مو پاس ترې يووړ۔ په \LR{$\pi$} کښې \LR{$\pi_1$} په \LR{$\pi_1''$} بدل کړئ۔
د نوي ثبوت ترتيبي شمېر لږ دے، ځکه لږ تر لږه د
\LR{$\RightR{\lforall}$} استنتاج ترتيبي شمېر (لږ تر لږه)~\LR{$1$} کم
شوے دے۔ اوس استقرايي فرضيه وکاروئ۔
\olpseganchor{OLP-0661-B017}{end}

\olpseganchor{OLP-0661-B018}{start}
هغه حالتونه چې کميتي استنتاج \RightR{\lexists} وي نور هم
اسان دي، ځکه انقباض نۀ غواړي او د ځانګړے متغير شرط ته اندېښنه
نۀ وي۔ مثلاً فرض کړئ چې له \LR{$\RightR{\lexists}$} وروسته
\LR{$\RightR{\land}$} راځي (دا ځل د کيڼې مقدمې په توګه)۔ اړوند فرعي ثبوت
\LR{$\pi_1$} دے:
\[
\AxiomC{}
\RightLabel{\LR{$\pi_2$}}
\Deduce$\Gamma' \fCenter \Delta', \lexists[x][!A(x)], !A(t), !B$
\RightLabel{\RightR{\lexists}}
\UnaryInf$\Gamma' \fCenter \Delta', \lexists[x][!A(x)], !B$
\AxiomC{}
\RightLabel{\LR{$\pi_3$}}
\Deduce$\Gamma' \fCenter \Delta', \lexists[x][!A(x)], !C$
\RightLabel{\RightR{\land}}
\BinaryInf$\Gamma' \fCenter \Delta', \lexists[x][!A(x)], !B \land !C$
\DisplayProof
\]
په~\LR{$\pi$} کښې \LR{$\pi_1$} په~\LR{$\pi_1'$} بدل کړئ:
\[
\AxiomC{}
\RightLabel{\LR{$\pi_2$}}
\Deduce$\Gamma' \fCenter \Delta', \lexists[x][!A(x)], !A(t), !B$
\AxiomC{}
\RightLabel{\LR{$\pi_3'$}}
\Deduce$\Gamma' \fCenter \Delta', \lexists[x][!A(x)], !A(t), !C$
\RightLabel{\RightR{\land}}
\BinaryInf$\Gamma' \fCenter \Delta', \lexists[x][!A(x)], !A(t), !B \land !C$
\RightLabel{\RightR{\lexists}}
\UnaryInf$\Gamma' \fCenter \Delta', \lexists[x][!A(x)], !B \land !C$
\DisplayProof
\]
چې \LR{$\pi_3'$} له \LR{$\pi_3$} څخه په ښي اړخ کښې د~\LR{$!A(t)$} په
کمزوري کولو تر لاسه کېږي۔ دا کمزوري کول يوازې د~\LR{$\pi_3$} د
هر سېکوېنټ ښي اړخ ته \LR{$!A(t)$} زياتوي، نو د هېڅ کميتي استنتاج
ترتيبي شمېر نۀ بدلوي۔ په~\LR{$\pi'$} کښې د کميتي استنتاجونو ترتيبي
شمېرونه د~\LR{$\pi$} په شان دي، پرته له هغه \RightR{\lexists}
استنتاج څخه چې له~\RightR{\land} سره بدل شوے او ترتيبي شمېر
ئې~\LR{$1$} کم شوے دے۔ استقرايي فرضيه کارولے شو۔
\end{proof}
\olpseganchor{OLP-0661-B018}{end}

\olpseganchor{OLP-0661-B019}{start}
\begin{prob}
د \olref{thm:midsequent} په ثبوت کښې هغه حالتونه بيان کړئ چې
\LeftR{\lexists} د~\RightR{\lif} په مقدمه او \LeftR{\lforall}
د~\LeftR{\lor} په کيڼه مقدمه ختمېږي۔
\end{prob}
\olpseganchor{OLP-0661-B019}{end}

\olpseganchor{OLP-0661-B020}{start}
\begin{proof}[د هېربراند د قضيې ثبوت]
فرض کړئ چې \LR{$\pi$} په \LR{$\Sequent \lexists[x][!A(x)]$} ختمېږي۔ د
\olref{thm:midsequent} له مخې \LR{$\pi$} په داسې ثبوت~\LR{$\pi'$}
بدلولے شو چې منځنے سېکوېنټ~\LR{$S$} ولري، او هغه ضرور دا بڼه لري:
\[\Sequent
\lexists[x][!A(x)], !A(t_1), \dots, !A(t_n).\]
څرنګه چې د \LR{$S$} پاس ټول استنتاجونه گزاره‌يي دي، \LR{$\lexists[x][!A(x)]$}
د~\LR{$S$} پاس د کوم استنتاج اصلي فارمول کېدے نۀ شي۔ دا چې اټومي
نۀ دے، په بديهيه کښې هم اصلي کېدے نۀ شي۔ څرنګه چې
\LR{$\lexists[x][!A(x)]$} د \LR{$!A(t_1)$} اصلي فرعي فارمول کېدے نۀ شي
(ياد مو وي چې \LR{$!A(x)$} له کميت اچوونکو خالي دے)، د~\LR{$S$} پاس
فعال هم کېدے نۀ شي۔ ځکه نو په~\LR{$S$} ختمېدونکي فرعي ثبوت
کښې د \LR{$\lexists[x][!A(x)]$} لرې کول د \LR{$\Sequent !A(t_1), \dots, !A(t_n)$}
يو کره ثبوت ورکوي، چې ترې د \LR{$\RightR{\lor}$} د \LR{$n-1$}
استنتاجونو په وسيله د \LR{$\Sequent !A(t_1) \lor \dots \lor !A(t_n)$}
يو ثبوت تر لاسه کېدے شي۔
\end{proof}
\olpseganchor{OLP-0661-B020}{end}

\olpseganchor{OLP-0661-B021}{start}
% \begin{prob}
% Find a Herbrand disjunction of $\lexists[x][\lforall[y][(!A(x) \lif
% !A(y))]]$ by finding a cut-free ثبوت of $\Sequent
% \lexists[x][\lforall[y][(!A(x) \lif !A(y))]]$ and applying the
% transformation in \olref{thm:midsequent} to find a midsequent (if
% necessary).
% \end{prob}
\olpseganchor{OLP-0661-B021}{end}

\olpseganchor{OLP-0658-B004}{start}
\olfileid{pt}{cut}{itp}
\olpseganchor{OLP-0658-B004}{end}

\olpseganchor{OLP-0658-B005}{start}
\olsection{د مايېهارا مرستندويه قضيه او د کرېګ د منځګړيتوب قضيه}
\olpseganchor{OLP-0658-B005}{end}

\olpseganchor{OLP-0658-B006}{start}
د منځګړيتوب قضيه، چې ويليم کرېګ وړاندې کړې ده، وايي: که \LR{$!A
\Entails !B$} وي، نو داسې تړلے فارمول~\LR{$!C$} (\emph{منځګړی}) شته
چې \LR{$!A \Entails !C$} او \LR{$!C \Entails !B$} وي۔ منځګړی~\LR{$!C$} داسې
غوره کېدے شي چې په~\LR{$!C$} کښې راغلي ټول ثابت سمبولونه او پريديکاتونه
په~\LR{$!A$} او~\LR{$!B$} دواړو کښې راغلي وي۔ (فرض کوو چې هېڅ تابعې
نۀ کارېږي۔) له دې قيد پرته، \LR{$!A$}~او \LR{$!B$} په اسانۍ منځګړي کېدے شي۔
\olpseganchor{OLP-0658-B006}{end}

\olpseganchor{OLP-0658-B007}{start}
د منځګړيتوب قضيه د کلاسيکي لومړۍ درجې منطق دپاره صدق کوي او
دا د لومړۍ درجې منطق «وروستے مهم کشف شوے خاصيت» بلل شوے دے۔
په مدل پوهنه او اتومات استدلال کښې مهم استعمالونه لري۔ د دې
لومړنے محرک او استعمال د علم په فلسفه کښې ؤ، ځکه د دې په وسيله
کتلے شو چې يوه تيوري په کومو لومړنيو مفاهيمو بديهياتي کېدے شي
يا نۀ شي کېدے۔ دا د بېت د تعريف‌وړتيا قضيه هم لازموي: که يوه
تيوري يو مفهوم په ضمني توګه تعريفولے شي، نو په څرګند ډول ئې هم
تعريفولے شي۔
\olpseganchor{OLP-0658-B007}{end}

\olpseganchor{OLP-0658-B008}{start}
لکه په رياضيکي منطق کښې ډېرې پايلې، د منځګړيتوب قضيه هم د
مدل پوهنې او هم د ثبوت پوهنې په طريقو ثابتېدے شي۔ د ثبوت پوهنې
د طريقو ګټه دا ده چې د منځګړي شتون ښيي او داسې الګوريتم هم راکوي
چې منځګړی د دې خبرې له ثبوت څخه محاسبه کوي چې \LR{$!A \Entails !B$}
دے؛ مثلاً په~\Log{G3c} کښې د \LR{$!A \Sequent !B$} له ثبوت څخه۔
\olpseganchor{OLP-0658-B008}{end}

\olpseganchor{OLP-0658-B009}{start}
دې \LR{$\Lang L(!A)$} په~\LR{$!A$} کښې د ثابت سمبولونه او پريديکاتونه
سټ وي او \LR{$\Lang L(\Gamma) = \bigcup \Setabs{\Lang L(!A)}{!A \in
\Gamma}$} دې وي۔ په دې ټوله برخه کښې فرض کوو چې له داسې
فارمولونه سره کار کوو چې تابعې نۀ لري۔
\olpseganchor{OLP-0658-B009}{end}

\olpseganchor{OLP-0658-B010}{start}
\begin{thm}[د کرېګ د منځګړيتوب قضيه]\ollabel{thm:interpolation}
دې د~\LR{$!A$} ژبه \LR{$\Lang L(!A)$} او د \LR{$!B$} ژبه
\LR{$\Lang L(!B)$} وي۔ که \LR{$\Log{G3c} \Proves !A \Sequent !B$} وي،
نو په ژبه~\LR{$\Lang L(!C) \subseteq \Lang
L(!A) \cap \Lang L(!B)$} کښې داسې فارمول~\LR{$!C$} شته چې
\LR{$\Log{G3c} \Proves !A \Sequent !C$} او
\LR{$\Log{G3c} \Proves !C \Sequent !B$} وي۔
\end{thm}
\olpseganchor{OLP-0658-B010}{end}

\olpseganchor{OLP-0658-B011}{start}
د منځګړيتوب د قضيې يوه عمومي بڼه ثابتوو، څو په~\Log{G3c} کښې
د ثبوتونه له جوړښت څخه کار واخلو۔ د دې دپاره د سېکوېنټونو
مقدم او تالي هر يو په دوو برخو وېشل شوي وګڼئ۔
\LR{$\maeh{\Gamma_1}{\Gamma_2} \Sequent
\maeh{\Delta_1}{\Delta_2}$} داسې ليکو چې \LR{$\Gamma_1$} او \LR{$\Delta_1$}
لومړۍ برخې ته او \LR{$\Gamma_2$} او \LR{$\Delta_2$} دويمې برخې ته اړه ولري۔
بيا د منځګړيتوب قضيه له لاندې مرستندويې قضيې څخه نېغ په نېغه راځي۔
\olpseganchor{OLP-0658-B011}{end}

\olpseganchor{OLP-0658-B012}{start}
\begin{lem}[د مايېهارا مرستندويه قضيه]\ollabel{lem:maehara}
که \LR{$\Log{G3c}
\Proves \maeh{\Gamma_1}{\Gamma_2} \Sequent \maeh{\Delta_1}{\Delta_2}$}
وي، نو په ژبه~\LR{$\Lang L(!C)
\subseteq \Lang L(\Gamma_1,
\Delta_1) \cap \Lang L(\Gamma_2, \Delta_2)$} کښې داسې فارمول~\LR{$!C$}
شته چې \LR{$\Log{G3c}
\Proves \Gamma_1 \Sequent \Delta_1, !C$} او \LR{$\Log{G3c} \Proves !C,
\Gamma_2 \Sequent \Delta_2$} وي۔
\end{lem}
\olpseganchor{OLP-0658-B012}{end}

\olpseganchor{OLP-0658-B013}{start}
\begin{proof}
  فرض کړئ چې \LR{$\pi \Proves \maeh{\Gamma_1}{\Gamma_2} \Sequent
  \maeh{\Delta_1}{\Delta_2}$} دے۔ دعوه په
  \LR{$n=\pheight{\pi}$} باندې په استقرا ثابتوو۔
\olpseganchor{OLP-0658-B013}{end}

\olpseganchor{OLP-0658-B014}{start}
  که \LR{$n=0$} وي، نو \LR{$\maeh{\Gamma_1}{\Gamma_2} \Sequent
  \maeh{\Delta_1}{\Delta_2}$} يو بديهي اصل دے او کوم اټومي
  فارمول~\LR{$!A$} په \LR{$\Gamma_1, \Gamma_2$} او \LR{$\Delta_1,
  \Delta_2$} دواړو کښې راغلے دے۔ څلور حالتونه دي: په کيڼ اړخ
  کښې \LR{$!A$}~لومړۍ که دويمې برخې ته اړه لري، او په ښي اړخ کښې
  لومړۍ که دويمې برخې ته۔
  \begin{enumerate}
    \item \LR{$!A \in \Gamma_1$} او \LR{$!A \in \Delta_1$}۔ داسې~\LR{$!C$}
    غواړو چې \LR{$\Gamma_1 \Sequent \Delta_1, !C$} او \LR{$!C, \Gamma_2
    \Sequent \Delta_2$} دواړه د ثابتولو وړ وي۔ لومړنے له وړاندې
    بديهي اصل دے، که~\LR{$!C$} هر ډول غوره کړو، ځکه \LR{$!A$} په کيڼ او ښي
    دواړو اړخونو کښې شته۔ د دې دپاره چې \LR{$!C, \Gamma_2
    \Sequent \Delta_2$} خامخا د ثابتولو وړ وي، دلته عمومي خوندي
    انتخاب دا دے چې \LR{$!C \ident \lfalse$} واخلو۔
    \item \LR{$!A \in \Gamma_1$} او \LR{$!A \in \Delta_2$}۔ نو \LR{$\Gamma_1
    \Sequent \Delta_1, !A$} او \LR{$!A, \Gamma_2 \Sequent \Delta_2$}
    دواړه بديهي اصلونه دي؛ ځکه \LR{$!C \ident !A$} اخستلے شو۔
    \item \LR{$!A \in \Gamma_2$} او \LR{$!A \in \Delta_1$}۔ نو \LR{$!A, \Gamma_1
    \Sequent \Delta_1$} او \LR{$\Gamma_2 \Sequent \Delta_2, !A$} به
    بديهي اصلونه وي، خو~\LR{$!A$} ئې په ناسم اړخونو کښې دے۔ نو \LR{$!A$}
    منځګړی نۀ دے۔ خو د \LR{$\RightR{\lnot}$} او \LR{$\LeftR{\lnot}$} په
    وسيله \LR{$\Gamma_1
    \Sequent \Delta_1, \lnot!A$} او \LR{$\lnot !A, \Gamma_2 \Sequent
    \Delta_2$} ثابت کولے شو۔
    \item \LR{$!A \in \Gamma_2$} او \LR{$!A \in \Delta_2$}۔ تمرين۔
  \end{enumerate}
  هغه حالتونه هم تمرين ته پرېږدو چې \LR{$\maeh{\Gamma_1}{\Gamma_2} \Sequent
  \maeh{\Delta_1}{\Delta_2}$} ځکه بديهي اصل وي چې \LR{$\lfalse \in
  \Gamma_1$} يا \LR{$\lfalse \in \Gamma_2$} دے۔
\olpseganchor{OLP-0658-B014}{end}

\olpseganchor{OLP-0658-B015}{start}
  که \LR{$n>0$} وي، \LR{$\pi$} په يوه منطقي استنتاج ختمېږي۔ د کارول شوې
  قاعدې او د اصلي فارمول د برخې له مخې بېل حالتونه ګورو۔ په هر
  حالت کښې دا استقرايي فرض کارولے شو چې مرستندويه قضيه د ټولو
  هغو ثبوتونه~\LR{$\pi_1$} دپاره صدق کوي چې
  \LR{$\pheight{\pi_1} < n$} وي، که د~\LR{$\pi_1$} وروستے سېکوېنټ هر ډول
  په دوو برخو وېشل شوے وي۔ څو ځانګړي حالتونه وګورو۔
  \begin{enumerate}
    \item \LR{$\pi$} په \LR{$\RightR{\lnot}$} ختمېږي او اصلي فارمول ئې~\LR{$\lnot
    !A$} دے۔ دوه حالتونه دي:~\LR{$\lnot !A$} لومړۍ که دويمې برخې ته
    اړه لري؛ د دې له مخې \LR{$\pi$} په دې ختمېږي:
\[
    \AxiomC{}
    \RightLabel{\LR{$\pi_1$}}
    \Deduce$\maeh{!A, \Gamma_1}{\Gamma_2} \fCenter
      \maeh{\Delta_1}{\Delta_2}$
    \RightLabel{\RightR{\lnot}}
    \UnaryInf$\maeh{\Gamma_1}{\Gamma_2} \fCenter
      \maeh{\Delta_1, \lnot !A}{\Delta_2}$
    \DisplayProof
    \]
    يا
\[
    \AxiomC{}
    \RightLabel{\LR{$\pi_1$}}
    \Deduce$\maeh{\Gamma_1}{!A, \Gamma_2} \fCenter
      \maeh{\Delta_1}{\Delta_2}$
    \RightLabel{\RightR{\lnot}}
    \UnaryInf$\maeh{\Gamma_1}{\Gamma_2} \fCenter
      \maeh{\Delta_1}{\Delta_2, \lnot !A}$
    \DisplayProof
    \]
    پام وکړئ: که اصلي فارمول~\LR{$\lnot !A$} لومړۍ برخې ته اړه
    لري، نو په مقدمه کښې مو فعال فارمول~\LR{$!A$} هم لومړۍ برخې
    ته ورکړے دے۔ دا انتخاب زمونږ په لاس کښې دے۔ استقرايي فرض
    به هله هم کارېدۀ چې فعال فارمول مو بلې برخې ته ورکړے وے، خو
    بيا به مو منځګړی نۀ شو موندلے (تمرين: وې ازمويئ)۔
\olpseganchor{OLP-0658-B015}{end}

\olpseganchor{OLP-0658-B016}{start}
    لومړے هغه حالت وګورئ چې~\LR{$\lnot !A$} لومړۍ برخې ته اړه لري۔
    د استقرايي فرض له مخې، د~\LR{$\pi_1$} د وروستي سېکوېنټ منځګړی
    شته؛ يعنې داسې فارمول~\LR{$!C_1$} چې دواړه لاندې سېکوېنټونه
\begin{align*}
    !A, \Gamma_1 & \Sequent \Delta_1, !C_1 &&\text{\RL{او}} &
    !C_1, \Gamma_2 & \Sequent \Delta_2
    \intertext{\RL{د ثابتولو وړ دي۔ مونږ داسې فارمول~\LR{$!C$} غواړو چې
    لاندې دواړه سېکوېنټونه د ثابتولو وړ وي:}}
    \Gamma_1 & \Sequent \Delta_1, \lnot !A, !C &&\text{\RL{او}} &
    !C, \Gamma_2 & \Sequent \Delta_2
    \end{align*}
    څرګنده ده چې \LR{$!C \ident !C_1$} اخستلے شو: استقرايي فرض له
    وړاندې وايي چې د ښي اړخ سېکوېنټ د ثابتولو وړ دے؛ د کيڼ اړخ
    سېکوېنټ له هغه سېکوېنټ څخه په~\LR{$\RightR{\lnot}$} راځي چې له
    استقرايي فرضه مو تر لاسه کړے ؤ۔ استقرايي فرض دا هم وايي چې
    \LR{$\Lang L(!C_1) \subseteq \Lang L(!A, \Gamma_1,
    \Delta_1) \cap \Lang L(\Gamma_2, \Delta_2)$}۔ ځکه چې \LR{$\Lang L(\lnot
    !A) = \Lang L(!A)$} او \LR{$\Lang L(!C) = \Lang L(!C_1)$} دي، نو
    \LR{$\Lang L(!C_1) \subseteq \Lang L(\Gamma_1, \Delta_1, \lnot !A) \cap \Lang
    L(\Gamma_2, \Delta_2)$} لرو۔
\olpseganchor{OLP-0658-B016}{end}

\olpseganchor{OLP-0658-B017}{start}
    په دويم حالت کښې وضع لږه بېله ده، ځکه~\LR{$\lnot !A$} دويمې
    برخې ته اړه لري۔ په مقدمه کښې فعال فارمول هم دويمې
    برخې ته ورکوو او استقرايي فرض کاروو، څو لاندې سېکوېنټونه
\begin{align*}
    \Gamma_1 & \Sequent \Delta_1, !C_1 &&\text{\RL{او}} &
    !C_1, !A, \Gamma_2 & \Sequent \Delta_2
    \intertext{\RL{د ثابتولو وړ وي۔ که بيا د دويم سېکوېنټ دپاره د \RightR{\lnot} قاعده وکاروو، د لاندې سېکوېنټونو ثبوتونه تر لاسه کوو}}
    \Gamma_1 & \Sequent \Delta_1, !C_1 &&\text{\RL{او}} &
    !C_1, \Gamma_2 & \Sequent \Delta_2, \lnot !A
    \end{align*}
    دا بيا عين هغه سېکوېنټونه دي چې د \LR{$!C_1$} د منځګړي کېدو
    دپاره ئې د ثابتولو وړ کېدل غواړو؛ بيا \LR{$!C \ident !C_1$} اخلو۔
    لکه مخکې، \LR{$\Lang
    L(!C_1) \subseteq \Lang L(\Gamma_1, \Delta_1) \cap \Lang L(\Gamma_2,
    \Delta_2, \lnot !A)$} دے۔
    \item \LR{$\pi$} په \LR{$\RightR{\lif}$} ختمېږي او اصلي فارمول
    ئې~\LR{$!A \lif !B$} دے۔ بيا دوه حالتونه دي: \LR{$!A \lif
    !B$} لومړۍ که دويمې برخې ته اړه لري۔ ثبوت~\LR{$\pi$} په دې
    ختمېږي:
\[
    \AxiomC{}
    \RightLabel{\LR{$\pi_1$}}
    \Deduce$\maeh{!A, \Gamma_1}{\Gamma_2} \fCenter
      \maeh{\Delta_1, !B}{\Delta_2}$
    \RightLabel{\RightR{\lif}}
    \UnaryInf$\maeh{\Gamma_1}{\Gamma_2} \fCenter
      \maeh{\Delta_1, !A \lif !B}{\Delta_2}$
    \DisplayProof
    \]
    يا
\[
    \AxiomC{}
    \RightLabel{\LR{$\pi_1$}}
    \Deduce$\maeh{\Gamma_1}{!A, \Gamma_2} \fCenter
      \maeh{\Delta_1}{\Delta_2, !B}$
    \RightLabel{\RightR{\lif}}
    \UnaryInf$\maeh{\Gamma_1}{\Gamma_2} \fCenter
      \maeh{\Delta_1}{\Delta_2, !A \lif !B}$
    \DisplayProof
    \]
    بيا مو په مقدمه کښې فعال فارمولونه هماغې برخې ته ورکړي
    دي چې په پايله کښې اصلي فارمول ورته اړه لري۔ په لومړي
    حالت کښې استقرايي فرض کاروو او لاندې ثبوتونه تر لاسه کوو:
\begin{align*}
    !A, \Gamma_1 & \Sequent \Delta_1, !B, !C_1 &&\text{\RL{او}} & !C_1, \Gamma_2 & \Sequent \Delta_2
    \intertext{\RL{د کيڼ اړخ سېکوېنټ د \RightR{\lif} دپاره مناسبه مقدمه ده۔ نو لاندې سېکوېنټونه د ثابتولو وړ دي:}}
    \Gamma_1 & \Sequent \Delta_1, !A \lif !B, !C_1 &&\text{\RL{او}} & !C_1, \Gamma_2 & \Sequent \Delta_2
    \end{align*}
    نو \LR{$!C_1$} د~\LR{$\pi$} د وروستي سېکوېنټ دپاره هم منځګړی دے او
    يو ځل بيا \LR{$!C \ident !C_1$} اخستلے شو۔
\olpseganchor{OLP-0658-B017}{end}

\olpseganchor{OLP-0658-B018}{start}
    بل حالت هم په همدې ډول څېړل کېږي۔
    \item \LR{$\pi$} په \LR{$\LeftR{\lif}$} ختمېږي او اصلي فارمول
    ئې~\LR{$!A \lif !B$} دے۔ که \LR{$!A \lif !B$} لومړۍ برخې ته اړه
    ولري، نو ثبوت~\LR{$\pi$} په دې ختمېږي:
\[
    \AxiomC{}
    \RightLabel{\LR{$\pi_1$}}
    \Deduce$\maeh{\Gamma_1}{\Gamma_2} \fCenter
      \maeh{\Delta_1, !A}{\Delta_2}$
    \AxiomC{}
    \RightLabel{\LR{$\pi_2$}}
    \Deduce$\maeh{!B, \Gamma_1}{\Gamma_2} \fCenter
      \maeh{\Delta_1}{\Delta_2}$
    \BinaryInf$\maeh{!A \lif !B, \Gamma_1}{\Gamma_2} \fCenter
      \maeh{\Delta_1}{\Delta_2}$
    \DisplayProof
    \]
    استقرايي فرض اوس څلور د ثابتولو وړ سېکوېنټونه راکوي؛ دوه
    د کيڼې مقدمې د منځګړي \LR{$!C_1$} دپاره او دوه د دويمې مقدمې
    د منځګړي \LR{$!C_2$} دپاره:
\begin{align*}
    \Gamma_1 & \Sequent \Delta_1, !A, !C_1 &&\text{\RL{(1)}} &
    !C_1, \Gamma_2 & \Sequent \Delta_2 && \text{\RL{(2)}}\\
    !B, \Gamma_1 & \Sequent \Delta_1, !C_2 &&\text{\RL{(3)}} &
    !C_2, \Gamma_2 & \Sequent \Delta_2 && \text{\RL{(4)}}
    \end{align*}
    سېکوېنټونه (1) او (3) د~\LeftR{\lif} مقدمې کېدے شي، که
    کمزوري کول وکاروو او د~(1) په ښي اړخ \LR{$!C_2$} او د~(3) په
    ښي اړخ \LR{$!C_1$} ورزيات کړو:
\[
    \AxiomC{}
    \Deduce$\Gamma_1 \fCenter \Delta_1, !A, !C_1, !C_2$
    \AxiomC{}
    \Deduce$!B, \Gamma_1 \fCenter \Delta_1, !C_1, !C_2$
    \RightLabel{\LeftR{\lif}}
    \BinaryInf$!A \lif !B, \Gamma_1 \fCenter \Delta_1, !C_1, !C_2$
    \DisplayProof
    \]
    د \LR{$\RightR{\lor}$} په کارولو دا سېکوېنټ تر لاسه کوو:
\[
    !A \lif !B, \Gamma_1 \fCenter \Delta_1, !C_1 \lor !C_2
    \]
    دا ښيي چې \LR{$!C_1 \lor !C_2$} په دې حالت کښې زمونږ منځګړی
    کېدے شي۔ او په رښتيا، بل اړين سېکوېنټ له پاتې سېکوېنټونو
    (2) او~(4) څخه ثابت کولے شو:
\[
    \AxiomC{}
    \Deduce$!C_1, \Gamma_2 \fCenter \Delta_2$
    \AxiomC{}
    \Deduce$!C_2, \Gamma_2 \fCenter \Delta_2$
    \RightLabel{\LeftR{\lor}}
    \BinaryInf$!C_1 \lor !C_2, \Gamma_2 \fCenter \Delta_2$
    \DisplayProof
    \]
    دا کتنه پاتې ده چې \LR{$!C_1 \lor !C_2$} په سمه ژبه کښې دے۔
    استقرايي فرض راکوي چې
\begin{align*}
    \Lang L(!C_1) & \subseteq
    \Lang L(\Gamma_1, \Delta_1, !A) \cap \Lang L(\Gamma_2, \Delta_2) &\text{\RL{او}}\\
    \Lang L(!C_2) & \subseteq
    \Lang L(!B, \Gamma_1, \Delta_1) \cap \Lang L(\Gamma_2, \Delta_2)
    \intertext{\RL{ځکه چې}}
    \Lang L(!A) & \subseteq \Lang L(!A \lif !B) &\text{\RL{او}}\\
    \Lang L(!B) & \subseteq \Lang L(!A \lif !B)
    \intertext{\RL{دواړه لاندې اړيکې لرو}}
    \Lang L(!C_1) & \subseteq
    \Lang L(\Gamma_1, \Delta_1, !A \lif !B) \cap \Lang L(\Gamma_2, \Delta_2) &\text{\RL{او}}\\
    \Lang L(!C_2) & \subseteq
    \Lang L(\Gamma_1, \Delta_1, !A \lif !B) \cap \Lang L(\Gamma_2, \Delta_2)
    \intertext{\RL{او په پايله کښې، \LR{$\Lang L (!C_1 \lor !C_2) = {}$}}}
    \Lang L(!C_1) \cup \Lang L(!C_2) & \subseteq
    \Lang L(\Gamma_1, \Delta_1, !A \lif !B) \cap \Lang L(\Gamma_2, \Delta_2),
    \end{align*}
    لکه چې غوښتل شوي وو۔
\olpseganchor{OLP-0658-B018}{end}

\olpseganchor{OLP-0658-B019}{start}
    که \LR{$!A \lif !B$} دويمې برخې ته اړه لري، وضع بېله ده، خو
    په ورته ډول څېړل کېږي۔ استقرايي فرض بيا څلور د ثابتولو وړ
    سېکوېنټونه راکوي:
\begin{align*}
    \Gamma_1 & \Sequent \Delta_1, !C_1 &&\text{\RL{(1)}} &
    !C_1, \Gamma_2 & \Sequent \Delta_2, !A && \text{\RL{(2)}}\\
    \Gamma_1 & \Sequent \Delta_1, !C_2 &&\text{\RL{(3)}} &
    !C_2, !B, \Gamma_2 & \Sequent \Delta_2 && \text{\RL{(4)}}
    \end{align*}
    اوس سېکوېنټونه (2) او (4)، له ځينو د کمزوري کولو په وسيله
    زياتو شوو فارمولونه سره، د~\LeftR{\lif} مقدمې جوړوي:
\[
    \AxiomC{}
    \Deduce$!C_1, !C_2, \Gamma_2 \fCenter \Delta_2, !A$
    \AxiomC{}
    \Deduce$!B, !C_1, !C_2, \Gamma_2 \fCenter \Delta_2$
    \RightLabel{\LeftR{\lif}}
    \BinaryInf$!A \lif !B, !C_1, !C_2, \Gamma_2 \fCenter \Delta_2$
    \DisplayProof
    \]
    بيا غواړو دا داسې سېکوېنټ کړو چې له \LR{$!C_1$} او \LR{$!C_2$} څخه
    جوړ يو واحد فارمول ولري؛ خو دا ځل ئې په مقدم کښې
    غواړو، ځکه اوس له دويمې برخې سره کار کوو۔ د \LR{$\LeftR{\land}$}
    په کارولو دا ترلاسه کېږي؛ منځګړی بايد \LR{$!C \ident
    (!C_1 \land !C_2)$} واخلو۔ پاتې سېکوېنټونه (1) او (3) هم
    د بلې برخې اړوند سېکوېنټ د ثابت کولو دپاره کارولے شو:
\[
    \AxiomC{}
    \Deduce$\Gamma_1 \fCenter \Delta_1, !C_1$
    \AxiomC{}
    \Deduce$\Gamma_1 \fCenter \Delta_1, !C_2$
    \RightLabel{\RightR{\land}}
    \BinaryInf$\Gamma_1 \fCenter \Delta_1, !C_1 \land !C_2$
    \DisplayProof
    \]
    لکه مخکې، \LR{$\Lang L(!C_1 \land !C_2) \subseteq \Lang
    L(\Gamma_1, \Delta_1) \cap \Lang L(\Gamma_2, \Delta_2, !A \lif
    !B)$} لرو۔
\olpseganchor{OLP-0658-B019}{end}

\olpseganchor{OLP-0658-B020}{start}
    \item \LR{$\pi$} په \LR{$\RightR{\lforall}$} ختمېږي او اصلي فارمول
    ئې~\LR{$\lforall[x][!A(x)]$} دے:
\[
    \AxiomC{}
    \RightLabel{\LR{$\pi_1$}}
    \Deduce$\maeh{\Gamma_1}{\Gamma_2} \fCenter
    \maeh{\Delta_1, !A(c)}{\Delta_2}$
    \RightLabel{\RightR{\lforall}}
    \UnaryInf$\maeh{\Gamma_1}{\Gamma_2} \fCenter
      \maeh{\Delta_1, \lforall[x][!A(x)]}{\Delta_2}$
    \DisplayProof
    \]
    يا
\[
    \AxiomC{}
    \RightLabel{\LR{$\pi_1$}}
    \Deduce$\maeh{\Gamma_1}{\Gamma_2} \fCenter
      \maeh{\Delta_1}{\Delta_2, !A(c)}$
    \RightLabel{\RightR{\lforall}}
    \UnaryInf$\maeh{\Gamma_1}{\Gamma_2} \fCenter
      \maeh{\Delta_1}{\Delta_2, \lforall[x][!A(x)]}$
    \DisplayProof
    \]
    په لومړي حالت کښې استقرايي فرض د \LR{$\Gamma_1 \Sequent \Delta_1, !A(c), !C_1$}
    او \LR{$!C_1, \Gamma_2
    \Sequent \Delta_2$} ثبوتونه راکوي۔ په لومړي سېکوېنټ
    \LR{$\RightR{\lforall}$} کارولے شو، څو \LR{$\Gamma_1 \Sequent \Delta_1, \lforall[x][!A(x)], !C_1$}
    تر لاسه کړو۔ (د ځانګړي متغير شرط پوره دے، ځکه د~\LR{$\pi$} په
    پای کښې د \RightR{\lforall} په قاعده کښې له وړاندې پوره ؤ۔)
    نو \LR{$!C_1$} به منځګړی وي، که دا شرط پوره وي چې \LR{$\Lang
    L(!C_1) \subseteq \Lang L(\Gamma_1, \Delta_1, \lforall[x][!A(x)])
    \cap \Lang L(\Gamma_2, \Delta_2)$}۔ د استقرايي فرض له مخې،
    \LR{$\Lang L(!C_1) \subseteq \Lang L(\Gamma_1,
    \Delta_1, !A(c)) \cap \Lang L(\Gamma_2, \Delta_2)$} لرو۔ نو
    د~\LR{$c$} په هکله بايد پام وکړو: آيا \LR{$c$} په~\LR{$!C_1$} کښې راتلے
    شي؟ نۀ شي راتلے۔ که \LR{$c$} په \LR{$!C_1$} کښې راشي، نو هم
    \LR{$c \in \Lang L(\Gamma_1, \Delta_1,
    !A(c))$} (چې رښتيا ده) او هم \LR{$c \in \Lang L(\Gamma_2,
    \Delta_2)$} به وي۔ خو بيا به \LR{$c$} په~\LR{$\pi$} کښې د
    \LR{$\RightR{\lforall}$} د استنتاج په پايله کښې راشي او د ځانګړي
    متغير شرط به مات کړي۔
\olpseganchor{OLP-0658-B020}{end}

\olpseganchor{OLP-0658-B021}{start}
    بل حالت هم ورته دے۔
\olpseganchor{OLP-0658-B021}{end}

\olpseganchor{OLP-0658-B022}{start}
    \item \LR{$\pi$} په \LR{$\RightR{\lexists}$} ختمېږي او اصلي
    فارمول ئې~\LR{$\lexists[x][!A(x)]$} دے۔ بيا دوه حالتونه دي۔
    هغه حالت ګورو چې \LR{$\lexists[x][!A(x)]$} لومړۍ برخې ته اړه لري؛
    بل حالت تمرين ته پرېږدو۔
\olpseganchor{OLP-0658-B022}{end}

\olpseganchor{OLP-0658-B023}{start}
    په لومړي حالت کښې \LR{$\pi$} په دې ختمېږي:
\[
    \AxiomC{}
    \RightLabel{\LR{$\pi_1$}}
    \Deduce$\maeh{\Gamma_1}{\Gamma_2} \fCenter
      \maeh{\Delta_1, \lexists[x][!A(x)], !A(t)}{\Delta_2}$
    \RightLabel{\RightR{\lexists}}
    \UnaryInf$\maeh{\Gamma_1}{\Gamma_2} \fCenter
      \maeh{\Delta_1, \lexists[x][!A(x)]}{\Delta_2}$
    \DisplayProof
    \]
    استقرايي فرض د \LR{$\Gamma_1 \Sequent
    \Delta_1, \lexists[x][!A(x)], !A(t), !C_1$} او \LR{$!C_1, \Gamma_2
    \Sequent \Delta_2$} ثبوتونه راکوي۔ په لومړي سېکوېنټ
    \LR{$\RightR{\lexists}$} کارولے شو، څو \LR{$\Gamma_1 \Sequent \Delta_1, \lexists[x][!A(x)], !C_1$}
    تر لاسه کړو۔ نو \LR{$!C_1$} به بيا منځګړی وي، که دا شرط پوره
    وي چې \LR{$\Lang L(!C_1) \subseteq \Lang L(\Gamma_1, \Delta_1,
    \lexists[x][!A(x)]) \cap \Lang L(\Gamma_2, \Delta_2)$}۔ دلته
    استقرايي فرض \LR{$\Lang L(!C_1)
    \subseteq \Lang L(\Gamma_1, \Delta_1, \lexists[x][!A(x)], !A(t))
    \cap \Lang L(\Gamma_2, \Delta_2)$} راکوي۔ اوس ياد کړئ چې
    فرض مو کړے ؤ چې هېڅ تابعې نۀ شته۔ نو د دې قواعدو
    تړلے ترم~\LR{$t$} يوازې ثابت سمبول کېدے شي؛ يعنې \LR{$t \ident b$}۔
    که \LR{$b \notin \Lang L(!C_1)$} يا \LR{$b \in \Lang L(\Gamma_1, \Delta_1, \lexists[x][!A(x)]) \cap
    \Lang L(\Gamma_2, \Delta_2)$} وي، اندېښنه نۀ شته:
    \LR{$!C_1$} منځګړی اخستلے شو۔ خو که \LR{$b \in \Lang
    L(!C_1)$} او \LR{$b \notin \Lang L(\Gamma_1, \Delta_1,
    \lexists[x][!A(x)]) \cap \Lang L(\Gamma_2, \Delta_2)$} وي، بيا؟
    لومړے، په~\LR{$!C_1$} کښې د~\LR{$b$} هر راتګ په يوه تازه
    متغير~\LR{$y$} بدلولے شو او \LR{$!C_1 \ident \Subst{!C_1'}{b}{y}$}
    ليکلے شو، چې~\LR{$!C_1'$} پکښې~\LR{$b$} نۀ لري۔ سربېره پردې، يا
    \LR{$b \notin
    \Lang L(\Gamma_1, \Delta_1, \lexists[x][!A(x)])$} يا \LR{$b \notin
    \Lang L(\Gamma_2, \Delta_2)$} دے۔ په ترتيب سره \LR{$!C \ident
    \lforall[y][!C_1'(y)]$} يا \LR{$!C \ident \lexists[y][!C_1'(y)]$}
    اخستلے شو۔ په لومړي حالت کښې لرو:
\[
    \AxiomC{}
    \Deduce$\Gamma_1 \fCenter
      \Delta_1, \lexists[x][!A(x)], !A(b), !C_1'(b)$
    \RightLabel{\RightR{\lexists}}
    \UnaryInf$\Gamma_1 \fCenter \Delta_1, \lexists[x][!A(x)], !C_1'(b)$
    \RightLabel{\RightR{\lforall}}
    \UnaryInf$\Gamma_1 \fCenter
      \Delta_1, \lexists[x][!A(x)], \lforall[y][!C_1'(y)]$
    \DisplayProof
    \]
    او
\[
    \AxiomC{}
    \Deduce$!C_1'(b), \lforall[y][!C_1'(y)], \Gamma_2 \fCenter \Delta_2$
    \RightLabel{\LeftR{\lforall}}
    \UnaryInf$\lforall[y][!C_1'(y)], \Gamma_2 \fCenter \Delta_2$
    \DisplayProof
    \]
    د پاسنيو سېکوېنټونو ثبوتونه له استقرايي فرضه راځي؛ په
    دويم کښې په مقدم کښې د \LR{$\lforall[y][!C_1'(y)]$} د زياتولو
    دپاره د کمزوري کولو منل کېدل کاروو۔ په لومړي ثبوت
    کښې د \RightR{\lforall} د ځانګړي متغير شرط پوره دے، ځکه
    \LR{$b \notin \Lang L(\Gamma_1, \Delta_1,
    \lexists[x][!A(x)])$} او~\LR{$!C_1'$} داسې تعريف شوے چې~\LR{$b$}
    نۀ لري۔
\olpseganchor{OLP-0658-B023}{end}

\olpseganchor{OLP-0658-B024}{start}
    په دويم حالت کښې، له استقرايي فرض او د~\LR{$\lexists[y][!C_1'(y)]$}
    په زياتولو د کمزوري کولو له منل کېدو څخه، لرو:
\[
    \AxiomC{}
    \Deduce$\Gamma_1 \fCenter
      \Delta_1, \lexists[x][!A(x)], !A(b), \lexists[y][!C_1'(y)], !C_1'(b)$
    \RightLabel{\RightR{\lexists}}
    \UnaryInf$\Gamma_1 \fCenter
      \Delta_1, \lexists[x][!A(x)], \lexists[y][!C_1'(y)], !C_1'(b)$
    \RightLabel{\RightR{\lexists}}
    \UnaryInf$\Gamma_1 \fCenter
      \Delta_1, \lexists[x][!A(x)], \lexists[y][!C_1'(y)]$
    \DisplayProof
    \]
    او
\[
    \AxiomC{}
    \Deduce$!C_1'(b), \Gamma_2 \fCenter \Delta_2$
    \RightLabel{\LeftR{\lexists}}
    \UnaryInf$\lexists[y][!C_1'(y)], \Gamma_2 \fCenter \Delta_2$
    \DisplayProof
    \]
    په دويم ثبوت کښې د \LeftR{\lexists} د ځانګړي متغير
    شرط پوره دے، ځکه \LR{$b \notin \Lang L(\Gamma_2,
    \Delta_2)$} او \LR{$!C_1'$} داسې تعريف شوے چې~\LR{$b$} نۀ لري۔
  \end{enumerate}
  نور حالتونه هم ورته دي؛ تمرينونو ته ئې پرېږدو۔
\end{proof}
\olpseganchor{OLP-0658-B024}{end}

\olpseganchor{OLP-0658-B025}{start}
\begin{prob}
فرض کړئ چې د \LR{$\lnand$} په نوم يو تړونکے د دې قواعدو سره لرو:
\[
\Axiom$\Gamma \fCenter \Delta, !A$
\Axiom$\Gamma \fCenter \Delta, !B$
\RightLabel{\LeftR{\lnand}}
\BinaryInf$!A \lnand !B, \Gamma \fCenter \Delta$
\DisplayProof
\quad
\Axiom$!A, !B, \Gamma \fCenter \Delta$
\RightLabel{\RightR{\lnand}}
\UnaryInf$\Gamma \fCenter \Delta, !A \lnand !B$
\DisplayProof
\]
د \olref{lem:maehara} د ثبوت هغه حالتونه ترسره کړئ چې \LR{$\pi$}
په دغو قواعدو کښې په يوه ختمېږي، او وښيئ چې د مايېهارا
مرستندويه قضيه لا هم صدق کوي۔
\end{prob}
\olpseganchor{OLP-0658-B025}{end}

\olpseganchor{OLP-0658-B026}{start}
اوس چې د مايېهارا مرستندويه قضيه مو ثابته کړه، د کرېګ د
منځګړيتوب د قضيې د ثابتولو دپاره ئې کارولے شو۔
\olpseganchor{OLP-0658-B026}{end}

\olpseganchor{OLP-0658-B027}{start}
\begin{proof}[د \olref{thm:interpolation} ثبوت]
  فرض کړئ چې \LR{$!A \Sequent !B$} د ثابتولو وړ دے۔ په
  \LR{$\maeh{!A}{\ } \Sequent \maeh{\ }{!B}$} باندې \olref{lem:maehara}
  وکاروئ، څو منځګړی~\LR{$!C$} تر لاسه کړئ۔ \LR{$\Proves !A \Sequent !C$}
  او \LR{$\Proves !C
  \Sequent !B$} لرو۔
\end{proof}
\olpseganchor{OLP-0658-B027}{end}

\olpseganchor{OLP-0658-B028}{start}
فرض کړئ چې \LR{$\Gamma$} په ژبه~\LR{$\Lang L$} کښې يوه تيوري
(د تړلي فارمولونه سټ) ده او ژبه د \LR{$n$} ځايونو {مسند}~\LR{$R$} لري۔
وايو چې \LR{$\Gamma$}،~\LR{$R$} \emph{په ضمني توګه تعريفوي}، هله او يوازې هله چې
\begin{align*}
\Gamma, \Gamma'
& \Entails \lforall[x_1][\dots\lforall[x_n][(R(x_1,\dots,x_n) \liff
R'(x_1, \dots, x_n))]],
\intertext{\RL{دلته \LR{$\Gamma'$} له \LR{$\Gamma$} څخه د~\LR{$R$} هر راتګ په~\LR{$R' \notin \Lang L$}
بدلولو تر لاسه کېږي۔ \LR{$\Gamma$}،~\LR{$R$} \emph{په څرګند ډول تعريفوي}،
هله او يوازې هله چې داسې~\LR{$!A(x_1, \dots, x_n)$} شته چې~\LR{$R$} نۀ لري او}}
\Gamma &\Entails
\lforall[x_1][\dots\lforall[x_n][(R(x_1,\dots,x_n) \liff !A(x_1,
\dots, x_n))]].
\end{align*}
څرګنده ده: که \LR{$\Gamma$}، \LR{$R$} په څرګند ډول تعريفوي، نو~\LR{$R$} په
ضمني توګه هم تعريفوي۔ معکوسه خبره څرګنده نۀ ده، خو بيا هم صدق کوي:
\olpseganchor{OLP-0658-B028}{end}

\olpseganchor{OLP-0658-B029}{start}
\begin{lem}[د بېت د تعريف‌وړتيا قضيه]
  که \LR{$\Gamma$}،~\LR{$R$} په ضمني توګه تعريفوي، نو په څرګند ډول ئې هم
  تعريفوي۔
\end{lem}
\olpseganchor{OLP-0658-B029}{end}

\olpseganchor{OLP-0658-B030}{start}
\begin{proof}
  فرض کړئ چې \LR{$\Gamma$}،~\LR{$R$} په ضمني توګه تعريفوي۔ دې \LR{$R'$} او
  \LR{$\Gamma'$} د پورته په شان وي او \LR{$\Lang L' = \Lang L(\Gamma') = \Lang L \setminus \{R\}
  \cup \{R'\}$} دې وي۔ د فرض له مخې لرو:
\begin{align*}
    \Gamma, \Gamma' & 
    \Sequent \lforall[x_1][\dots\lforall[x_n][(R(x_1,\dots,x_n) 
    \liff R'(x_1, \dots, x_n))]]
  \intertext{\RL{دې \LR{$c_1$}، \dots،~\LR{$c_n$} داسې ثابت سمبولونه وي چې په~\LR{$\Gamma$} کښې نۀ راځي۔ د معکوسولو له مرستندويې قضيې څخه لرو}}
    \Gamma, \Gamma', R(c_1,\dots,c_n) & \Sequent  R'(c_1, \dots, c_n)
  \intertext{\RL{د مايېهارا مرستندويه قضيه په دې وکاروئ}}
    \maeh{\Gamma, R(c_1,\dots,c_n)}{\Gamma'} & \Sequent
      \maeh{\ }{R'(c_1, \dots, c_n)}
  \intertext{\RL{څو داسې~\LR{$!A \ident !A(x_1, \dots, x_n)$} ومومو چې}}
    \Gamma, R(c_1,\dots,c_n) & \Sequent !A(c_1, \dots, c_n) \\
    !A(c_1, \dots, c_n), \Gamma' & \Sequent R'(c_1, \dots, c_n)
  \intertext{\RL{ثبوتونه لري۔ د مايېهارا مرستندويه قضيه منځګړے د دواړو خواوو په ګډو نښو محدودوي۔ تازه ثابتونه په دواړو خواوو کښې راتلے شي۔ چې د هغو په ځاے متغيرونه واخلو، \LR{$\Lang L(!A) \subseteq \Lang L(\Gamma) \setminus \{R\}$} دے، نو~\LR{$!A$}،~\LR{$R$} نۀ لري۔ په دويم سېکوېنټ کښې د ثبوت په هر ځاے \LR{$R'$} په~\LR{$R$} بدل کړئ، څو د لاندې سېکوېنټ ثبوت تر لاسه شي}}
    !A(c_1, \dots, c_n), \Gamma & \Sequent R(c_1, \dots, c_n)
  \end{align*}
  \RightR{\lif} او \RightR{\lforall} وکاروئ، څو د لاندې
  سېکوېنټ ثبوت تر لاسه کړئ:
\[
    \Gamma \Sequent \lforall[x_1][\dots\lforall[x_n][(R(x_1,\dots,x_n) \liff !A(x_1, \dots, x_n))]].
  \]
  دا ښيي چې \LR{$!A(x_1, \dots, x_n)$} په~\LR{$\Gamma$} کښې د \LR{$!R$}
  څرګند تعريف دے۔
\end{proof}
\olpseganchor{OLP-0658-B030}{end}

\olpseganchor{OLP-0658-B031}{start}
درېيمه پايله، چې منځګړيتوب ته هم‌ارزښته ده او د هغې په شان
د مايېهارا په مرستندويه قضيه ثابتېدے شي، دا قضيه ده: که دوه
سازګارې تيورۍ~\LR{$\Gamma_1$} او \LR{$\Gamma_2$} د خپلې ټولې ګډې ژبې
يوه بشپړه تيوري \LR{$\Gamma$} ولري، يعنې
\LR{$\Lang L(\Gamma) = \Lang L(\Gamma_1)
\cap \Lang L(\Gamma_2)$} وي، نو \LR{$\Gamma_1 \cup \Gamma_2$} سازګاره ده۔
په لاندې استدلال کښې د ټولې ګډې ژبې د هرې جملې د فيصلې لپاره
د ګډې تيورۍ بشپړتيا کاروو۔
\olpseganchor{OLP-0658-B031}{end}

\olpseganchor{OLP-0658-B032}{start}
په ياد ولرئ چې بشپړه تيوري هغه ده چې په خپله ژبه کښې د هرې
تړلے فارمول~\LR{$!A$} دپاره يا \LR{$\Gamma \Proves !A$} يا
\LR{$\Gamma \Proves \lnot !A$} وي۔ ځکه چې \LR{$\Gamma_1$} او~\LR{$\Gamma_2$}
د~\LR{$\Gamma$} سازګارې پراختياوې دي، نو \LR{$\Gamma$} بايد سازګاره وي۔
له دې راځي: که \LR{$!A$} د \LR{$\Lang L(\Gamma_1) \cap \Lang
L(\Gamma_2)$} په ژبه کښې تړلے فارمول وي، نو \LR{$\Gamma_1 \Entails !A$}
او \LR{$\Gamma_2 \Entails \lnot !A$} دواړه نۀ شي لرلے۔ د دې د کتلو
دپاره فرض کړئ داسې~\LR{$!A$} شته۔ که \LR{$\Gamma_1 \Entails !A$} وي، نو
\LR{$\Gamma \Entails !A$} دے۔ که داسې نۀ وي، د~\LR{$\Gamma$} د بشپړتيا
له مخې \LR{$\Gamma \Entails \lnot !A$} دے؛ او ځکه چې
\LR{$\Gamma \subseteq \Gamma_1$} ده، \LR{$\Gamma_1 \Entails \lnot !A$}
بۀ وي او \LR{$\Gamma_1$} به ناسازګاره شي۔ نو \LR{$\Gamma \Entails !A$}
دے؛ او ځکه چې \LR{$\Gamma
\subseteq \Gamma_2$} ده، \LR{$\Gamma_2 \Entails !A$} دے۔ نو
\LR{$\Gamma \Entails/
\lnot !A$}، ګنې \LR{$\Gamma_2$} به ناسازګاره وي۔
\olpseganchor{OLP-0658-B032}{end}

\olpseganchor{OLP-0658-B033}{start}
د ثبوت دپاره همدا کافي ده چې په ګډه ژبه~\LR{$\Lang L(\Gamma_1) \cap \Lang L(\Gamma_2)$}
کښې داسې تړلے فارمول~\LR{$!A$} نۀ وي چې \LR{$\Gamma_1
\Entails !A$} او \LR{$\Gamma_2 \Entails \lnot !A$} وي۔ دلته دا
ثبوت د ثبوت پوهنې په طريقه، د مايېهارا په مرستندويه قضيه، ورکوو۔
\olpseganchor{OLP-0658-B033}{end}

\olpseganchor{OLP-0658-B034}{start}
\begin{thm}[د رابنسن د ګډې سازګارۍ قضيه]
فرض کړئ چې \LR{$\Gamma_1$} او \LR{$\Gamma_2$} سازګارې تيورۍ دي او په
ژبه \LR{$\Lang L(\Gamma_1) \cap \Lang L(\Gamma_2)$} کښې
د هېڅ~\LR{$!A$} دپاره \LR{$\Gamma_1 \Entails !A$} او
\LR{$\Gamma_2 \Entails \lnot !A$} دواړه نۀ لرو۔ نو
\LR{$\Gamma_1 \cup \Gamma_2$} سازګاره ده۔
\end{thm}
\olpseganchor{OLP-0658-B034}{end}

\olpseganchor{OLP-0658-B035}{start}
\begin{proof}
  د خلاف فرض وکړئ چې \LR{$\Gamma_1 \cup \Gamma_2$} ناسازګاره ده۔
  نو سېکوېنټ \LR{$\Gamma_1, \Gamma_2 \Sequent \ $} يو
  ثبوت لري۔ په \LR{$\maeh{\Gamma_1}{\Gamma_2}
  \Sequent \maeh{\ }{\ }$} د مايېهارا مرستندويه قضيه وکاروئ،
  څو په ژبه~\LR{$\Lang L(\Gamma_1) \cap \Lang L(\Gamma_2)$}
  کښې داسې تړلے فارمول~\LR{$!A$} تر لاسه کړئ چې
  \LR{$\Proves \Gamma_1 \Sequent !A$} او \LR{$\Proves !A, \Gamma_2 \Sequent
  \quad$} وي۔ له وروستي څخه \LR{$\Proves \Gamma_2 \Sequent \lnot !A$}
  تر لاسه کوو۔ خو فرض مو کړے ؤ چې داسې تړلے فارمول نۀ شته۔
\end{proof}
\olpseganchor{OLP-0658-B035}{end}

\olpseganchor{OLP-0658-B036}{start}
پورته مو په پټه فرض کړے ؤ چې تيورۍ \LR{$\Gamma$}، \LR{$\Gamma_1$} او
\LR{$\Gamma_2$} متناهي دي۔ د متناهي شاهد خاصيت له مخې، پايلې د
نامتناهي تيوريو دپاره هم صدق کوي۔
\olpseganchor{OLP-0658-B036}{end}

\olpseganchor{OLP-0681-B004}{start}
\olchapter{pt}{ps}{د ثبوت لټون}
\olpseganchor{OLP-0681-B004}{end}

\olpseganchor{OLP-0680-B004}{start}
\olfileid{pt}{ps}{int}
\olpseganchor{OLP-0680-B004}{end}

\olpseganchor{OLP-0680-B005}{start}
\olsection{پېژندنه}
\olpseganchor{OLP-0680-B005}{end}

\olpseganchor{OLP-0680-B006}{start}
د \Log{G3c} يا \Log{N1i} په شان د ثبوت نظامونو
بديهيتونه او استنتاجي قواعد ټاکي چې د سيکوېنټونو يا
فارمولونه کومې ونې صحيح ثبوتونه ګڼل کېږي۔ که د
سيکوېنټونو يا فارمولونه داسې يوه ونه راکړل شي، کېدے شي
ورسره نور معلومات هم وي: په هر ګام کښې کوم قواعد کارېدلي،
يا د فرضونو او هغو استنتاجي قواعدو نښې چې دغه فرضونه باسي۔
د بديهيتونو او قواعدو تعريفونه بيا دا کتنه ممکنوي چې
راکړې ونه صحيح ثبوت ده که نۀ۔ خو د يو راکړي سيکوېنټ،
يا له راکړو فرضونو څخه د يو راکړي فارمول دپاره
ثبوت \emph{موندل} بېله او عموماً ګرانه مسئله ده۔
همدغه د \emph{ثبوت د لټون} مسئله ده۔
\olpseganchor{OLP-0680-B006}{end}

\olpseganchor{OLP-0680-B007}{start}
د ثبوت د لټون يوه ډېره ساده طريقه شته۔ فارمولونه د
متناهي الفبا د نښو لړۍ ګڼلے شو۔ سيکوېنټونه او د
سيکوېنټونو يا فارمولونه ونې هم د فارمولونه لړۍ
ګڼلے شو؛ کېدے شي د ونې د څانګېدو دپاره ځانګړي قوسونه
وکاروو۔ بديهيتونه او قواعد اجازه راکوي چې په ميخانيکي
ډول وګورو آيا د نښو يوه لړۍ صحيح ثبوت ښيي۔ نو په
ميخانيکي ډول \emph{ټول ممکن} صحيح ثبوتونه جوړولے شو:
د ثبوتونه د ښودلو د الفبا ټولې نښيزې لړۍ جوړوو او د هرې
ممکنې لړۍ دپاره ګورو چې صحيح ثبوت ښيي که نۀ۔ د
ثبوت د لټون ساده طريقه دا ده: ټول صحيح ثبوتونه تر
هغه وخته جوړ کړئ چې له هغو يو د راکړي سيکوېنټ، يا د
راکړي فارمول، ثبوت وي او پرانيستي فرضونه ئې په
راکړو فرضونو کښې شامل وي۔
\olpseganchor{OLP-0680-B007}{end}

\olpseganchor{OLP-0680-B008}{start}
غوره طريقه هغه ساده لاره ده چې په لاس د ثبوت د موندلو
په وخت ئې کاروئ: له پاييز سيکوېنټ څخه پيل کوئ، يو
فارمول ټاکئ او استنتاجي قاعده ``شاته'' کاروئ، څو
يوه يا څو مقدمې جوړې شي۔ که د دغو مقدمو ثبوتونه موجود
وي، له وروستي استنتاج سره يوځاے د راکړي سيکوېنټ
ثبوت جوړوي۔ په طبيعي استنتاجي نظام کښې هم د
ثبوت د موندلو دپاره ورته، خو يو څه پېچلې، طريقه
کاروئ۔
\olpseganchor{OLP-0680-B008}{end}

\olpseganchor{OLP-0680-B009}{start}
خو دغه طريقه بې‌خطا نۀ ده: که ناسمه قاعده وټاکئ، يا
فارمولونه په ناسم ترتيب وټاکئ، بې‌لارې پاے ته رسېدې
شئ۔ دا طريقه پوره ټاکلې هم نۀ ده: په هر ځاے کښې کېدے
شي د فارمول د ټاکلو يا د قاعدې د شاته کارولو څو
امکانونه وي۔ د \emph{ثبوت د لټون الګوريتم} يوه پوره
ټاکلې کړنلاره ده چې که ثبوت موجود وي، هغه مومي؛
يعنې بې‌خطا ده۔
\olpseganchor{OLP-0680-B009}{end}

\olpseganchor{OLP-0680-B010}{start}
د ځينو ثبوت نظامونو دپاره د ثبوت د لټون ساده
الګوريتمونه جوړېږي۔ په سيکوېنټي نظامونو کښې د
فارمول اصلي عملګر او دا چې په کيڼ يا ښي اړخ
کښې دے، ټاکي چې کومه قاعده شاته وکارول شي۔ د بېلګې په
توګه، که د \LR{$\Gamma \Sequent \Delta, !A \land !B$}
ثبوت غواړئ، چې \LR{$!A \land !B$} ئې په ښي اړخ کښې
دے، نو \RightR{\land} شاته کاروئ او دوه اړوندې مقدمې
\LR{$\Gamma \Sequent \Delta, !A$} او \LR{$\Gamma \Sequent \Delta, !B$}
جوړوئ۔ خو که نظام مو \Log{G1c} وي، د راغونډولو موجودګي
کار ګرانوي، ځکه نور امکانونه درکوي: \RightR{\Contraction}
هم شاته کارولے او مقدمه
\LR{$\Gamma \Sequent \Delta, !A \land !B, !A \land !B$}
جوړولے شئ؛ بيا
\LR{$\Gamma \Sequent \Delta, !A \land !B, !A \land !B, !A \land !B$}
او همداسې نور۔ په \Log{G2c} کښې کار لا ګران دے۔ د
\RightR{\land} شاته کارونه دا هم غواړي چې \LR{$\Gamma_1$}،
\LR{$\Gamma_2$} داسې اټکل کړئ چې \LR{$\Gamma = \Gamma_1 \cup \Gamma_2$}
وي، او \LR{$\Delta_1$}، \LR{$\Delta_2$} داسې چې
\LR{$\Delta = \Delta_1 \cup \Delta_2$} وي؛ بيا د مقدمو
\LR{$\Gamma_1 \Sequent \Delta_1, !A$} او
\LR{$\Gamma_2 \Sequent \Delta_2, !B$} هڅه وکړئ۔ ناسم اټکل
وروسته بې‌لارې پاے ته رسولے شي؛ بيا بايد پيل ته بېرته
لاړ شئ او نور \LR{$\Gamma_1$}، \LR{$\Gamma_2$}، \LR{$\Delta_1$}،
\LR{$\Delta_2$} وټاکئ۔ که فارمول د \LR{$!A \lor !B$} په بڼه
وي، بايد د \RightR{\lor} له دوو امکانونو څخه يو وټاکئ؛
دا يا مقدمه \LR{$\Gamma \Sequent \Delta, !A$} او يا مقدمه
\LR{$\Gamma \Sequent \Delta, !B$} جوړوي۔ ناسم انتخاب بيا هم
شاته تګ ته اړکولے شي۔
\olpseganchor{OLP-0680-B010}{end}

\olpseganchor{OLP-0680-B011}{start}
نظام \Log{G3c} دغه ستونزې نۀ لري۔ جوړښتي قواعد نۀ لري،
او د استنتاجي قاعدې انتخاب د فارمول په
اصلي عملګر او په هغه اړخ ټاکل کېږي چې فارمول
پکښې موجود دے۔ نو \Log{G3c} د ثبوت د لټون دپاره له
\Log{G1c} يا \Log{G2c} څخه زيات مناسب دے۔ بريالے لټون
په \Log{G3c} کښې ثبوت ورکوي؛ دغه ثبوت بيا په
\Log{G1c} يا \Log{G2c}، او حتٰی په \Log{N1c}، کښې
ثبوتونه ته اړولے شو۔ نو که د \Log{G3c} د ثبوت د
لټون الګوريتم او د \Log{G3c} ثبوتونه نورو نظامونو ته
د اړولو الګوريتمونه ولرو، د دغو نورو نظامونو د ثبوت د
لټون طريقې هم ترلاسه کېږي۔
\olpseganchor{OLP-0680-B011}{end}

\olpseganchor{OLP-0680-B012}{start}
څنګه پوهېږو چې د ثبوت د لټون الګوريتم بې‌خطا دے؟ بايد
وښيو چې که الګوريتم ثبوت نۀ مومي، هېڅ ثبوت
نۀ شي موجودېدې۔ که ناسم الګوريتم وټاکو، کېدے شي داسې
حالتونه راشي چې د موجود ثبوت له موندلو ناکام شي۔ ساده
الګوريتم دغه ستونزه نۀ لري، ځکه د ټولو ممکنو
ثبوتونه لټون کوي۔ خو د ثبوت د لټون الګوريتم بايد
اغېزمن وي او لږ تر لږه هڅه وغواړي۔
\olpseganchor{OLP-0680-B012}{end}

\olpseganchor{OLP-0680-B013}{start}
د ثبوت د لټون د بې‌خطا والي د ښودلو بنسټيزه مفکوره
دا ده: که الګوريتم ناکام شي، د ناکامۍ له طريقې ښيو چې
سيکوېنټ هېڅ ثبوت نۀ شي لرلے۔ دا په دې ښيو چې
سيکوېنټ معتبر نۀ دے۔ په کلاسيکي لومړۍ درجې منطق کښې
معتبر سيکوېنټونه \LR{$\Gamma \Sequent \Delta$} هغه دي چې
په هر جوړښت کښې يا د \LR{$\Gamma$} لږ تر لږه يو
فارمول دروغ وي، يا د~\LR{$\Delta$} لږ تر لږه يو
فارمول رښتيا وي۔ \Log{G3c} \emph{سم} دے؛ يعنې
يوازې معتبر سيکوېنټونه ثابت کوي۔ دا د ثبوتونه پر
سر په استقرا ثابتېږي: بديهيتونه ښکاره معتبر دي، او که
د يوې قاعدې مقدمې معتبرې وي، پايله ئې هم معتبره ده۔
د الګوريتم د بې‌خطا والي د ښودلو دپاره ښيو چې که د
\LR{$\Gamma \Sequent \Delta$} د ثبوت لټون ناکام شي،
جوړښت تعريفولے شو چې پکښې د \LR{$\Gamma$} ټول
فارمولونه رښتيا او د~\LR{$\Delta$} ټول فارمولونه دروغ
وي۔ دا د \Log{G3c} \emph{بشپړتيا} هم ثابتوي؛ يعنې که
\LR{$\Gamma \Sequent \Delta$} معتبر وي، د هغه يو
\Log{G3c}-ثبوت موجود وي۔ ځکه چې د ثبوت هر ناکام
لټون د \LR{$\Gamma \Sequent \Delta$} ناسم والے ښيي، د معتبر
سيکوېنټ هر لټون بايد بريالے شي۔
\olpseganchor{OLP-0680-B013}{end}

\olpseganchor{OLP-0683-B004}{start}
\olfileid{pt}{ps}{sal}
\olpseganchor{OLP-0683-B004}{end}

\olpseganchor{OLP-0683-B005}{start}
\olsection{د \texorpdfstring{\Log{G3c}}{G3c} د لټون الګوريتم}
\olpseganchor{OLP-0683-B005}{end}

\olpseganchor{OLP-0683-B006}{start}
د \Log{G3c} دپاره د لټون يو الګوريتم بيانوو۔ دا يوازېنې
ممکن الګوريتم نۀ دے، او تر ټولو اغېزمن هم نۀ دے۔ خو بشپړ دے:
که سيکوېنټ \LR{$\Gamma \Sequent \Delta$} ثبوت ولري، نو
بالاخره به ودرېږي او ثبوت به ومومي۔ له دې امله،
که بې له ثبوت موندلو ودرېږي، يا هېڅکله ونۀ درېږي، نو
\LR{$\Gamma \Sequent \Delta$} له پيله معتبر نۀ دے۔
\olpseganchor{OLP-0683-B006}{end}

\olpseganchor{OLP-0683-B007}{start}
د الګوريتم يو ضروري شرط دا دے چې په
\LR{$\Gamma \Sequent \Delta$} کښې، يا د ثبوت د لټون په جوړ شوي
سيکوېنټ کښې، هر ناعټومي فارمول «راکم» کړي؛ يعنې د
استنتاجي قاعدې په سرچپه کارونه کښې ئې د اصلي فارمول په توګه
وکاروي، او دا کار هر څومره چې اړين وي هغومره تکرار کړي۔ دې
شرط ته \emph{انصاف} وايو۔
\olpseganchor{OLP-0683-B007}{end}

\olpseganchor{OLP-0683-B008}{start}
الګوريتم په پړاوونو کښې کار کوي۔ د هر پړاو پايله د
سيکوېنټونو يوه ونه ده؛ د بل پړاو پايله په هر هغه تر ټولو
پاسني سيکوېنټ کښې د فارمول په راکمولو جوړېږي چې لا
بديهيت نۀ دے۔
\olpseganchor{OLP-0683-B008}{end}

\olpseganchor{OLP-0683-B009}{start}
په پړاو \LR{$0$} کښې له \LR{$\Gamma \Sequent \Delta$} څخه
پيل کوو۔ د انصاف د تضمين دپاره په \LR{$\Gamma$} او \LR{$\Delta$} کښې
فارمولونه ته شمېرې د شاخصونو په توګه داسې ورکوو چې بېلابېل
فارمولونه بېلابېلې شمېرې واخلي۔ مثلاً، له کيڼې خوا تر ښي
خوا د فارمولونه شمېرل په پام کښې ونيسئ۔ شمېرې د پاسنيو
نښو په توګه ليکو۔ که د \LR{$!A, !B \Sequent !C$} دپاره ثبوت
لټوو، د پړاو~\LR{$0$} پايله \LR{$!A^1, !B^2 \Sequent !C^3$} ده۔ په
داسې شاخص لرونکي سيکوېنټ کښې د پاسنۍ نښې په توګه راغلے تر
ټولو لوے~\LR{$n$} د هغه \emph{اعظمي شاخص} دے (په بېلګه کښې~\LR{$3$})۔
\olpseganchor{OLP-0683-B009}{end}

\olpseganchor{OLP-0683-B010}{start}
که سيکوېنټ کميت ټاکونکي لري، بايد دا هم معلومه وي چې
په ښي اړخ د \LR{$\lexists[x][!A(x)]$} يا په کيڼ اړخ د
\LR{$\lforall[x][!A(x)]$} د راکمولو دپاره کوم ترمونه وکاروو۔ يوازې
له ثابت سمبولونه \LR{$\Obj c_0$}، \LR{$\Obj c_1$}، \LR{$\Obj c_2$}، \dots
څخه جوړ شوي ترمونه پکار دي، چې په~\LR{$\Gamma$} او \LR{$\Delta$} کښې
راغلي هر تابع سمبول پکښې کارېدے شي۔ فرض کوو چې د دغو ټولو
ترمونو سېټ په يوه ټاکلې شمېرنه کښې ورکړے شوے، څو مثلاً «هغه
لومړے ثابت سمبول چې په \LR{$\Gamma \Sequent \Delta$} کښې نۀ
راځي» وښيو۔ دا شمېرنه په اغېزمنه توګه کوو، مثلاً د ترمونو د نښو د کوډ د اوږدوالي او بيا د الفبا د ترتيب له مخې۔ د الګوريتم د ونې له هر شاخص لرونکي سيکوېنټ سره
د ثابت سمبولونه يو سېټ تړلے وي۔ په پړاو~\LR{$0$} کښې سيکوېنټ ته
ورکړے شوے د ثابت سمبولونه سېټ د هغه ټول موجود
ثابت سمبولونه او هم \LR{$\{\Obj c_0\}$} رانغاړي؛ نو که سيکوېنټ
هېڅ ثابت سمبولونه ونۀ لري، همدغه وروستے سېټ پاتې کېږي۔
دغه پيليز ثابت په هرې ورپسې څانګه کښې پاتې کېږي۔
\olpseganchor{OLP-0683-B010}{end}

\olpseganchor{OLP-0683-B011}{start}
فرض کړئ چې الګوريتم په پړاو~\LR{$n$} کښې د شاخص لرونکو
سيکوېنټونو يوه ونه د ثابت سمبولونه له تړلو سېټونو سره جوړه
کړې ده۔ په پړاو~\LR{$n+1$} کښې د هر تر ټولو پاسني سيکوېنټ دپاره
لاندې کار کوو۔
\olpseganchor{OLP-0683-B011}{end}

\olpseganchor{OLP-0683-B012}{start}
که يو اټومي فارمول~\LR{$!A$} د سيکوېنټ په کيڼ او
ښي دواړو اړخونو کښې راځي، يا کيڼ اړخ~\LR{$\lfalse$} لري، هېڅ
کار نۀ کوو: داسې سيکوېنټونه په~\Log{G3c} کښې ثبوتونه لري،
او د دغه ځانګړي پاسني سيکوېنټ دپاره نور څۀ پاتې نۀ دي۔
\olpseganchor{OLP-0683-B012}{end}

\olpseganchor{OLP-0683-B013}{start}
که سيکوېنټ هېڅ ناعټومي فارمول ونۀ لري، الګوريتم
بند کړئ۔ د \LR{$\Gamma \Sequent \Delta$} هېڅ ثبوت نشته۔
\olpseganchor{OLP-0683-B013}{end}

\olpseganchor{OLP-0683-B014}{start}
کنه، تر ټولو کوچني شاخص~\LR{$i$} لرونکے ناعټومي
فارمول~\LR{$!A$} وټاکئ۔ اړوند پاسنے سيکوېنټ يا
\LR{$!A^i, \Pi \Sequent \Lambda$} دے، يا
\LR{$\Pi \Sequent \Lambda, !A^i$}۔ \LR{$k$} دې د سيکوېنټ اعظمي شاخص
او \LR{$C$} دې د ثابت سمبولونه هغه سېټ وي چې ورسره تړلے دے۔
\olpseganchor{OLP-0683-B014}{end}

\olpseganchor{OLP-0683-B015}{start}
\LR{$!A^i$} «راکوو»؛ يعنې د~\LR{$!A$} د اصلي عملګر او
په کيڼ يا ښي اړخ د \LR{$!A^i$} د ځاے له مخې د استنتاجي قاعدې
مطابق نوي پاسني سيکوېنټونه جوړوو۔
\olpseganchor{OLP-0683-B015}{end}

\olpseganchor{OLP-0683-B016}{start}
\begin{enumerate}
  \item \LR{$!A \ident !B \land !C$} او په کيڼ اړخ راځي: د
  \LR{$(!B \land !C)^i, \Pi \Sequent \Lambda$} پاس يو نوے
  پاسنے سيکوېنټ ورزيات کړئ:
  \[
  \Axiom$!B^{k+1}, !C^{k+2}, \Pi \fCenter \Lambda$
  \RightLabel{\LeftR{\land}}
  \UnaryInf$(!B \land !C)^i, \Pi \fCenter \Lambda$
  \DisplayProof
  \]
  نوي سيکوېنټ ته ورکړے شوے د ثابت سمبولونه سېټ~\LR{$C$} دے۔
  \item \LR{$!A \ident !B \land !C$} او په ښي اړخ راځي: د
  \LR{$\Pi \Sequent \Lambda, (!B \land !C)^i$} پاس دوه نوي
  پاسني سيکوېنټونه ورزيات کړئ:
  \[
  \Axiom$\Pi \fCenter \Lambda, !B^{k+1}$
  \Axiom$\Pi \fCenter \Lambda, !C^{k+1}$
  \RightLabel{\RightR{\land}}
  \BinaryInf$\Pi \fCenter \Lambda, (!B \land !C)^i$
  \DisplayProof
  \]
  نوي سيکوېنټ ته ورکړے شوے د ثابت سمبولونه سېټ~\LR{$C$} دے۔
  په هرې مقدمې کښې د نوو پېښو شاخصونه د پخواني اعظمي شاخص څخه لوے وي۔
\olpseganchor{OLP-0683-B016}{end}

\olpseganchor{OLP-0683-B017}{start}
  \item د \LR{$!A \ident \lnot !B$}، \LR{$!A \ident !B \lor !C$}
  او \LR{$!A \ident !B \lif !C$} حالتونه ورته دي او د مشقونو
  دپاره پرېښودل شوي دي۔
\olpseganchor{OLP-0683-B017}{end}

\olpseganchor{OLP-0683-B018}{start}
  \item \LR{$!A \ident \lexists[x][!B(x)]$} او په کيڼ اړخ
  راځي: د \LR{$\lexists[x][!B(x)]^i, \Pi \Sequent \Lambda$}
  پاس يو نوے پاسنے سيکوېنټ ورزيات کړئ:
  \[
  \Axiom$!B(c)^{k+1}, \Pi \fCenter \Lambda$
  \RightLabel{\LeftR{\lexists}}
  \UnaryInf$\lexists[x][!B(x)]^i, \Pi \fCenter \Lambda$
  \DisplayProof
  \]
  دلته \LR{$c$} هغه لومړے ثابت سمبول دے چې په~\LR{$C$} کښې نۀ وي۔
  نوي سيکوېنټ ته ورکړے شوے د ثابت سمبولونه سېټ
  \LR{$C \cup \{c\}$} دے۔
  \item \LR{$!A \ident \lexists[x][!B(x)]$} او په ښي اړخ راځي:
  د \LR{$\Pi \Sequent \Lambda, \lexists[x][!B(x)]$} پاس يو
  نوے پاسنے سيکوېنټ ورزيات کړئ:
  \[
  \Axiom$\Pi \fCenter \Lambda, \lexists[x][!B(x)]^{k+1}, !B(t)^{k+2}$
  \RightLabel{\RightR{\lexists}}
  \UnaryInf$\Pi \fCenter \Lambda, \lexists[x][!B(x)]^i$
  \DisplayProof
  \]
  دلته \LR{$t$} هغه لومړے ترم دے چې يوازې د~\LR{$C$} ثابت سمبولونه
  لري او تر دې سيکوېنټ لاندې په ښي اړخ د
  \LR{$\lexists[x][!B(x)]$} په راکمولو کښې لا نۀ دے کارول شوے
  (يا \LR{$\Obj c_0$}، که ټول داسې ترمونه کارول شوي وي)۔ نوي
  سيکوېنټ ته ورکړے شوے د ثابت سمبولونه سېټ \LR{$C$} دے۔
  \item د \LR{$!A \ident \lforall[x][!B(x)]$} حالتونه ورته دي
  او د مشقونو دپاره پرېښودل شوي دي۔
\end{enumerate}
\olpseganchor{OLP-0683-B018}{end}

\olpseganchor{OLP-0683-B019}{start}
که په پړاو \LR{$n+1$} کښې هېڅ پاسني سيکوېنټ ته نوے
سيکوېنټ ورزيات نۀ کړو، الګوريتم په برياليتوب درېږي۔ په دې
حالت کښې مو د~\LR{$\Gamma \Sequent \Delta$} دپاره ثبوت
موندلے، چې د پړاو \LR{$n+1$} له ونې څخه د ټولو شاخصونو په لرې
کولو ترلاسه کېږي۔ ټول پاسني سيکوېنټونه يا
\LR{$!A, \Pi \Sequent \Lambda, !A$} بڼه لري، چې ګډ فارمول ئې
اټومي دے، يا \LR{$\lfalse, \Pi \Sequent \Lambda$} بڼه لري۔ ټول
استنتاجونه د~\Log{G3c} سم استنتاجونه دي۔ دا د بياني
استنتاجونو دپاره څرګنده ده۔ وګورئ چې په هر ګام کښې له
سيکوېنټ سره تړلے د ثابت سمبولونه سېټ~\LR{$C$} د هغه ټول موجود
ثابت سمبولونه رانغاړي۔ نو د \LeftR{\lexists} او
\RightR{\lforall} په کارونه کښې د ځانګړي متغير شرط پوره
دے: ځانګړے متغير~\LR{$c$} داسې ثابت سمبول و چې د نتيجې په
تړلي سېټ~\LR{$C$} کښې نۀ و، نو \LR{$c$} په نتيجه کښې هم نۀ راتلو۔ لرو:
\olpseganchor{OLP-0683-B019}{end}

\olpseganchor{OLP-0683-B020}{start}
\begin{prop}
د \Log{G3c} د ثبوت د لټون الګوريتم سم دے: که په برياليتوب
ودرېږي، د پيل سيکوېنټ~\LR{$\Gamma \Sequent \Delta$}
ثبوت شته۔
\end{prop}
\olpseganchor{OLP-0683-B020}{end}

\olpseganchor{OLP-0679-B004}{start}
\olfileid{pt}{ps}{cpl}
\olpseganchor{OLP-0679-B004}{end}

\olpseganchor{OLP-0679-B005}{start}
\olsection{بشپړتيا}
\olpseganchor{OLP-0679-B005}{end}

\olpseganchor{OLP-0679-B006}{start}
اوس ښيو چې د ثبوت د لټون الګوريتم بې‌خطا دے: يعنې که
درېږي او ناکامېږي يا هېڅکله نۀ درېږي، د پيل
سيکوېنټ~\LR{$\Gamma \Sequent \Delta$} هېڅ ثبوت نشته۔ د دې
دپاره جوړښت~\LR{$\Struct M$} داسې تعريفوو چې د هر
\LR{$!A \in \Gamma$} دپاره \LR{$\Sat{M}{!A}$}، او د هر~\LR{$!A \in \Delta$}
دپاره \LR{$\Sat/{M}{!A}$} وي۔ په بله وينا، په~\LR{$\Struct M$} کښې
د~\LR{$\Gamma$} ټول فارمولونه رښتيا او د~\LR{$\Delta$} ټول
فارمولونه دروغ دي؛ يعنې \LR{$\Gamma \Sequent \Delta$} معتبر
نۀ دے۔ ځکه چې \Log{G3c} سم دے، \LR{$\Gamma \Sequent \Delta$}
هېڅ ثبوت نۀ لري۔
که په پيل سيکوېنټ کښې ازاد متغيرونه وي، هر يو په يوه بېل تازه ثابت بدلوو۔ دا يوازې د ازادو پېښو تعويض دے۔ الګوريتم په دغو ثابتونو لرونکي سيکوېنټ چلوو۔ که يو ثبوت ومومي، تړلي متغيرونه لومړے تازه نوموو، څو له پخوانيو ازادو متغيرونو سره ګډ نۀ شي۔ بيا هر پيلنی تازه ثابت په خپل پخواني ازاد متغير بدلوو۔ د کميتي قواعدو ځانګړي ثابتونه په خپل ځاے پاتې کېږي، ځکه د لټون په پيل کښې اخيستل شوي ثابتونه له دې ځانګړو ثابتونو بېل دي۔ نو بدله ونه د اصلي سيکوېنټ ثبوت دے۔ په لاندې د ناکامۍ استدلال کښې د سيکوېنټ نښې د همدې تړلې کاري بڼې لپاره کاروو۔
\olpseganchor{OLP-0679-B006}{end}

\olpseganchor{OLP-0679-B007}{start}
\LR{$\Struct{M}$} د فارمولونه له يوې جوړې \LR{$\Theta, \Xi$}
او د ثابت سمبولونه له سېټ~\LR{$C$} څخه تعريفوو۔ \LR{$\Theta, \Xi$} او~\LR{$C$}
د ناکام شوي، يعنې درېدلي يا نۀ درېدېدونکي، ثبوت لټون له
\emph{ناکامې څانګې} څخه اخلو۔ ناکامه څانګه د سيکوېنټونو
\LR{$\Pi_n \Sequent \Lambda_n$} او د ثابت سمبولونه د سېټونو~\LR{$C_n$}
متناهي يا نامتناهي لړۍ ده، چې \LR{$\Pi_0 \Sequent \Lambda_0$}
ئې پاييز سيکوېنټ \LR{$\Gamma \Sequent \Delta$} دے۔ هر چې راتلونکے
مورد موجود وي، \LR{$\Pi_{n+1} \Sequent \Lambda_{n+1}$} د
\LR{$\Pi_n \Sequent \Lambda_n$} هغه مقدمه ده چې په پړاو~\LR{$n+1$}
کښې جوړه شوې او الګوريتم ئې ثبوت نۀ مومي۔ د هر~\LR{$n$} په داسې
ناتمام پړاو کښې بايد لږ تر لږه يوه داسې مقدمه وي؛ کنه
الګوريتم به ثبوت موندلے و۔ \LR{$C_n$} د ثابت سمبولونه هغه
سېټ دے چې له~\LR{$\Pi_n \Sequent \Lambda_n$} سره تړلے دے۔
\olpseganchor{OLP-0679-B007}{end}

\olpseganchor{OLP-0679-B008}{start}
که الګوريتم پر هغه تر ټولو پاسني سيکوېنټ ناکام ودرېږي
چې يوازې اټومي فارمولونه لري، نو په همدغه سيکوېنټ ختمه
څانګه ناکامه څانګه ده۔ که الګوريتم هېڅکله نۀ درېږي، پرله‌پسې
لويې ونې جوړوي۔ دا د يوې نامتناهي ونې وده بللے شئ: هغه ونه
چې د الګوريتم له نامتناهي ډېرو ګامونو وروسته به موجوده وه۔
دا ونه نامتناهي، خو متناهي څانګېدونکې وي؛ هر غوټه يوازې يو
يا دوه اولادونه لري، يعنې هغه سيکوېنټونه چې سمدستي ئې د
پاسه دي۔ د کونيګ د لمې له مخې هره نامتناهي، متناهي
څانګېدونکې ونه يوه نامتناهي څانګه لري۔ که د لټون الګوريتم
تل روان وي، همدغه نامتناهي څانګه ناکامه څانګه ده۔
\olpseganchor{OLP-0679-B008}{end}

\olpseganchor{OLP-0679-B009}{start}
که د \LR{$\Pi_n \Sequent \Lambda_n$} او \LR{$C_n$} لړۍ ناکامه څانګه
وي، \LR{$\Theta = \bigcup_n \Pi_n$} او \LR{$\Xi = \bigcup_n \Lambda_n$}
واخلئ، چې شاخصونه او د فارمولونه تکراري موردونه ترې لرې
کوو؛ همدارنګه \LR{$C = \bigcup_n C_n$} واخلئ۔
\olpseganchor{OLP-0679-B009}{end}

\olpseganchor{OLP-0679-B010}{start}
اوس~\LR{$\Struct{M}$} تعريفوو۔
\begin{enumerate}
\item د تعريف ساحه \LR{$\Domain{M}$} د هغو ترمونو سېټ دے چې له~\LR{$C$}
او په~\LR{$\Gamma \Sequent \Delta$} کښې له موجودو تابعې
څخه جوړېږي، خو هېڅ متغيرونه نۀ لري۔
\item د هر ثابت \LR{$c \in C$} دپاره، \LR{$\Assign{c}{M} = c$}۔
\item که \LR{$f$} په \LR{$\Gamma \Sequent \Delta$} کښې يو \LR{$m$}-ځايه
تابع وي، او \LR{$t_1$}، \dots، \LR{$t_m \in \Domain{M}$} وي، نو
\LR{$\Assign{f}{M}(t_1, \dots, t_m) = f(t_1, \dots, t_m)$}۔
\item که \LR{$R$} په \LR{$\Gamma \Sequent \Delta$} کښې يو \LR{$m$}-ځايه
پريديکات وي، نو \LR{$\Assign{R}{M} = \Setabs{\tuple{t_1, \dots, t_m} \in
\Domain{M}^m}{R(t_1, \dots, t_m) \in \Theta}$}۔
\end{enumerate}
\LR{$\Struct{M}$} يو \emph{ترمي مدل} دے: د تعريف ساحه ئې د ترمونو
سېټ دے او ثابت سمبولونه او تابعې داسې تفسير شوي چې
\LR{$\Value{t}{M} = t$} تضمين کړي۔ په~\LR{$\Theta$} کښې اټومي
فارمولونه د \LR{$\Assign{R}{M}$} د تعريف دپاره داسې کارېږي چې
\LR{$\Sat{M}{R(t_1, \dots, t_m)}$} هله او يوازې هله چې
\LR{$R(t_1, \dots, t_m) \in \Theta$} وي۔
\olpseganchor{OLP-0679-B010}{end}

\olpseganchor{OLP-0679-B011}{start}
\begin{lem}\ollabel{lem:truth}
د هر \LR{$!A \in \Theta \cup \Xi$} دپاره:
\begin{enumerate}
  \item که \LR{$!A \in \Theta$} وي، نو \LR{$\Sat{M}{!A}$}۔
  \item که \LR{$!A \in \Xi$} وي، نو \LR{$\Sat/{M}{!A}$}۔
\end{enumerate}
\end{lem}
\olpseganchor{OLP-0679-B011}{end}

\olpseganchor{OLP-0679-B012}{start}
\begin{proof}
  د \LR{$\depth{!A}$} پر سر په استقرا۔
\olpseganchor{OLP-0679-B012}{end}

\olpseganchor{OLP-0679-B013}{start}
لومړے فرض کړئ چې \LR{$!A$} اټومي دے؛ يعنې يا
\LR{$!A \ident \lfalse$} او يا \LR{$!A \ident R(t_1, \dots, t_m)$}۔
\olpseganchor{OLP-0679-B013}{end}

\olpseganchor{OLP-0679-B014}{start}
ځکه چې \LR{$\Pi_n \Sequent \Lambda_n$} ناکامه څانګه ده، هېڅ
\LR{$\Pi_n$} د \LR{$\lfalse$} فارمول نۀ شي لرلے؛ کنه
\LR{$\Pi_n \Sequent \Lambda_n$} به بديهيت و۔ د~\LR{$\Struct{M}$} د
تعريف له مخې \LR{$\Sat{M}{R(t_1, \dots, t_m)}$} هله او يوازې هله
چې \LR{$R(t_1, \dots, t_m) \in \Theta$} وي۔ نو دواړه خبرې يوځاے
دا راکوي: که \LR{$!A \in \Theta$} وي، نو \LR{$\Sat{M}{!A}$}۔
\olpseganchor{OLP-0679-B014}{end}

\olpseganchor{OLP-0679-B015}{start}
که \LR{$!A$} اټومي او \LR{$!A \in \Pi_n$} وي، نو د راتلونکي
مورد د موجودېدو په صورت کښې \LR{$!A \in \Pi_{n+1}$}؛ او که
\LR{$!A \in \Lambda_n$} وي، نو \LR{$!A \in \Lambda_{n+1}$}۔ دا د لټون
الګوريتم د راتلونکو سيکوېنټونو د جوړولو له طريقې راځي۔ نو
که \LR{$!A \in \Theta \cap \Xi$} وي، د يو~\LR{$n$} دپاره
\LR{$!A \in \Pi_n \cap \Lambda_n$} وي۔ دا مانا لري چې
\LR{$\Pi_n \Sequent \Lambda_n$} بديهيت دے، خو ناکامه څانګه هېڅ
بديهيت نۀ لري۔ نو د اټومي~\LR{$!A$} دپاره، د جوړونې له مخې
\LR{$!A \notin \Theta \cap \Xi$}؛ ځکه که \LR{$!A \in \Xi$} وي، نو
\LR{$!A \notin \Theta$}۔ د~\LR{$\Struct{M}$} د تعريف له مخې، له دې
راځي چې که \LR{$!A \in \Xi$} وي، نو \LR{$\Sat/{M}{!A}$}۔
\olpseganchor{OLP-0679-B015}{end}

\olpseganchor{OLP-0679-B016}{start}
د استقرا د ګام دپاره فرض کړئ چې (۱) او (۲) د هغو ټولو
فارمولونه دپاره صحيح دي چې ژوروالے ئې له~\LR{$!A$} څخه لږه
ده۔ د حالتونو په بېلولو د \LR{$!A$} دپاره (۱) او (۲) ثابتوو۔
\olpseganchor{OLP-0679-B016}{end}

\olpseganchor{OLP-0679-B017}{start}
لومړے د ناعټومي فارمول د ټاکلو انصاف په پام کښې ونيسئ۔
کله چې \LR{$!A^i$} په \LR{$\Pi_n \Sequent \Lambda_n$} کښې موجود وي،
تر څو ټاکل شوے نۀ وي، \LR{$!A^i$} په
\LR{$\Pi_{n+1} \Sequent \Lambda_{n+1}$} کښې هم پاتې کېږي۔ ټاکنه
په \LR{$\Pi_n \Sequent \Lambda_n$} کښې د ټولو شاخصونو پر ځاے د ناعټومي فارمولونو تر ټولو کوچنے
شاخص~\LR{$i$} کاروي۔ د منصفانه لټون له مخې بالاخره داسې
\LR{$\Pi_n \Sequent \Lambda_n$} ته رسېږو چې \LR{$!A^i$} پکښې موجود
وي او \LR{$i$} د ناعټومي فارمولونو تر ټولو کوچنے شاخص وي۔ په
پړاو \LR{$n+1$} کښې الګوريتم~\LR{$!A^i$} راکموي۔ په بله وينا، که
\LR{$!A \in \Theta$} وي، د يو \LR{$n$} دپاره \LR{$\Pi_n = \Pi_n', !A^i$}
وي او \LR{$i$} د ناعټومي فارمولونو تر ټولو کوچنے شاخص وي۔ او
که \LR{$!A \in \Xi$} وي، د يو \LR{$n$} دپاره
\LR{$\Lambda_n = \Lambda_n', !A^i$} وي او \LR{$i$}~ئې همدغه تر ټولو
کوچنے ناعټومي شاخص وي۔ په هر ګام کښې نوي شاخصونه د اوسني اعظمي شاخص څخه لوے دي۔ تر يوې ورکړې ناعټومي پېښې کوچني ناعټومي شاخصونه محدود شمېر لري۔ تر څو چې دغه پېښه نۀ وي ټاکل شوې، په هر ګام کښې له هغو څخه يو کمېږي۔ نو دغه پېښه بالاخره راکمېږي۔
\olpseganchor{OLP-0679-B017}{end}

\olpseganchor{OLP-0679-B018}{start}
\begin{enumerate}
  \item \LR{$!A \in \Theta$} او \LR{$!A \ident !B \land !C$}۔ نو د يو
  \LR{$n$} دپاره \LR{$\Pi_n = \Pi_n', (!B \land !C)^i$}۔
  \LR{$\Pi_{n+1} \Sequent \Lambda_{n+1}$} د~\LeftR{\land} اړونده
  مقدمه ده؛ يعنې \LR{$\Pi_{n+1} = \Pi_n', !B^{k+1}, !C^{k+2}$}۔ نو
  \LR{$!B, !C \in \Theta$}۔ د استقرا د فرضيې له مخې
  \LR{$\Sat{M}{!B}$} او \LR{$\Sat{M}{!C}$}۔ د \LR{$\Sat{M}{}$} د تعريف له
  مخې، \LR{$\Sat{M}{!B \land !C}$} لرو۔
\olpseganchor{OLP-0679-B018}{end}

\olpseganchor{OLP-0679-B019}{start}
  \item \LR{$!A \in \Xi$} او \LR{$!A \ident !B \land !C$}۔ نو د يو
  \LR{$n$} دپاره \LR{$\Lambda_n = \Lambda_n', (!B \land !C)^i$}۔
  \LR{$\Pi_{n+1} \Sequent \Lambda_{n+1}$} د~\RightR{\land} له
  دوو مقدمو څخه يوه ده، چې پايله ئې
  \LR{$\Pi_n \Sequent \Lambda'_n, (!B \land !C)^i$} ده؛ يعنې
  \LR{$\Lambda_{n+1} = \Lambda_n', !B^{k+1}$} يا
  \LR{$\Lambda_{n+1} = \Lambda_n', !C^{k+1}$}۔ نو يا \LR{$!B \in \Xi$} يا
  \LR{$!C \in \Xi$}۔ د استقرا د فرضيې له مخې، يا \LR{$\Sat/{M}{!B}$}
  يا \LR{$\Sat/{M}{!C}$}۔ د \LR{$\Sat{M}{}$} د تعريف له مخې
  \LR{$\Sat/{M}{!B \land !C}$} لرو۔
\olpseganchor{OLP-0679-B019}{end}

\olpseganchor{OLP-0679-B020}{start}
  \item هغه حالتونه چې \LR{$!A \ident \lnot !B$}،
  \LR{$!A \ident !B \lor !C$} او \LR{$!A \ident !B \lif !C$} وي، ورته
  دي او د مشقونو په توګه پاتې دي۔
\olpseganchor{OLP-0679-B020}{end}

\olpseganchor{OLP-0679-B021}{start}
  \item \LR{$!A \in \Theta$} او \LR{$!A \ident \lexists[x][!B(x)]$}۔
  نو د يو \LR{$n$} دپاره \LR{$\Pi_n = \Pi_n', \lexists[x][!B(x)]^i$}۔
  \LR{$\Pi_{n+1} \Sequent \Lambda_{n+1}$} د~\LeftR{\lexists}
  اړونده مقدمه ده؛ يعنې د يو~\LR{$c$} دپاره
  \LR{$\Pi_{n+1} = \Pi_n', !B(c)^{k+1}$}۔ نو \LR{$!B(c) \in \Theta$} او
  \LR{$c \in C$}۔ د استقرا د فرضيې له مخې \LR{$\Sat{M}{!B(c)}$}۔ د
  \LR{$\Sat{M}{}$} د تعريف له مخې، \LR{$\Sat{M}{\lexists[x][!B(x)]}$}
  لرو۔
\olpseganchor{OLP-0679-B021}{end}

\olpseganchor{OLP-0679-B022}{start}
  \item \LR{$!A \in \Xi$} او \LR{$!A \ident \lexists[x][!B(x)]$}۔ نو
  د يو \LR{$n$} دپاره \LR{$\Lambda_n = \Lambda_n', \lexists[x][!B(x)]^i$}۔
  هر ځل چې \LR{$\lexists[x][!B(x)]$} راکمېږي،
  \LR{$\lexists[x][!B(x)]$} د مقدمې سيکوېنټ په ښي اړخ کښې له نوي
  شاخص سره پاتې کېږي؛ نو که \LR{$\lexists[x][!B(x)]$} په ناکامه څانګه کښې موجود وي،
  په ښي اړخ کښې نامتناهي ځله راکمېږي۔ د ناکامې څانګې پر
  سر هر ترم \LR{$t \in \Domain{M}$} بالاخره هغه لومړنے ترم شي
  چې يوازې د~\LR{$C_n$} ثابتونه لري او تر دې وړاندې په ښي اړخ
  کښې د \LR{$\lexists[x][!B(x)]$} په راکمولو کښې نۀ دے کارېدلے۔ که ترم له مخکې کارېدلے وي، مطلوب مورد له مخکې په څانګه کښې شته۔ که نۀ وي، د شمېرنې تر هغه مخکيني ترمونه محدود شمېر لري او ترم بيا بيا د ټاکلو وار ته رسېږي۔
  نو د هر \LR{$t \in \Domain{M}$} دپاره، د يو~\LR{$n$} لپاره
  \LR{$!B(t) \in \Lambda_n$} وي او ځکه \LR{$!B(t) \in \Xi$}۔ د استقرا
  د فرضيې له مخې، د هر \LR{$t \in \Domain{M}$} دپاره
  \LR{$\Sat/{M}{!B(t)}$}۔ د \LR{$\Sat{M}{}$} د تعريف له مخې،
  \LR{$\Sat/{M}{\lexists[x][!B(x)]}$} لرو۔
\olpseganchor{OLP-0679-B022}{end}

\olpseganchor{OLP-0679-B023}{start}
  \item هغه حالتونه چې \LR{$!A \ident \lforall[x][!B(x)]$} وي،
  ورته دي او د مشقونو په توګه پاتې دي۔
\end{enumerate}
\end{proof}
\olpseganchor{OLP-0679-B023}{end}

\olpseganchor{OLP-0679-B024}{start}
\begin{cor}\ollabel{cor:complete} د \Log{G3c} د ثبوت د
لټون الګوريتم بشپړ دے: که ناکام ودرېږي يا هېڅکله نۀ
درېږي، د پيل سيکوېنټ~\LR{$\Gamma \Sequent \Delta$} معتبر نۀ دے
او ځکه هېڅ ثبوت نۀ لري۔
\end{cor}
\olpseganchor{OLP-0679-B024}{end}

\olpseganchor{OLP-0679-B025}{start}
\begin{proof}
  که الګوريتم ناکام ودرېږي يا هېڅکله نۀ درېږي، يوه ناکامه
  څانګه شته چې~\LR{$\Struct{M}$} ترې تعريفېږي۔ د
  \olref{lem:truth} له مخې، د هر \LR{$!A \in \Gamma$} دپاره
  \LR{$\Sat{M}{!A}$} او د هر \LR{$!A \in \Delta$} دپاره
  \LR{$\Sat/{M}{!A}$}۔ په ياد ولرئ چې \LR{$\Pi_0 = \Gamma$} او
  \LR{$\Lambda_0 = \Delta$}، نو \LR{$\Gamma \subseteq \Theta$} او
  \LR{$\Delta \subseteq \Xi$}۔ که تازه ثابتونه مو د ازادو متغيرونو لپاره ورزيات کړي وي، جوړښت د اصلي ژبې پر نښو محدودوو۔ هر پخواني ازاد متغير ته د هغه د تازه ثابت ارزښت ورکوو، او نورو متغيرونو ته د ډومېن هر ټاکلے غړے ورکوو۔ د تعويض د لمې له مخې د اصلي فارمول صدق په دې ګومارنه کښې د هغه د تړلې بدلې بڼې له صدق سره يو شان دے۔ نو د اصلي سيکوېنټ کيڼ فارمولونه رښتيا او ښي فارمولونه دروغ دي۔
  ځکه \LR{$\Gamma \Sequent \Delta$} معتبر نۀ دے۔ د~\Log{G3c} د سموالي له مخې،
  \LR{$\Gamma \Sequent \Delta$} هېڅ ثبوت نۀ لري۔
\end{proof}
\olpseganchor{OLP-0679-B025}{end}

\olpseganchor{OLP-0679-B026}{start}
\begin{cor}
\Log{G3c} بشپړ دے: که \LR{$\Gamma \Sequent \Delta$} معتبر وي، نو
\LR{$\Log{G3c} \Proves \Gamma \Sequent \Delta$}۔
\end{cor}
\olpseganchor{OLP-0679-B026}{end}

\olpseganchor{OLP-0679-B027}{start}
\begin{proof}
  عکس نقيض ثابتوو۔ فرض کړئ چې \LR{$\Log{G3c} \Proves/ \Gamma
  \Sequent \Delta$}۔ نو د ثبوت د لټون الګوريتم بايد ناکام
  ودرېږي يا هېڅکله نۀ درېږي، ځکه هېڅ ثبوت نۀ شته چې
  وموندل شي۔ د \olref{cor:complete} له مخې،
  \LR{$\Gamma \Sequent \Delta$} معتبر نۀ دے۔
\end{proof}
\olpseganchor{OLP-0679-B027}{end}

\olpseganchor{OLP-0684-B004}{start}
\olfileid{pt}{ps}{tab}
\olpseganchor{OLP-0684-B004}{end}

\olpseganchor{OLP-0684-B005}{start}
\olsection{تابلوګانې}
\olpseganchor{OLP-0684-B005}{end}

\olpseganchor{OLP-0684-B006}{start}
د تابلو سيستمونه د ثبوت هغه سيستمونه دي چې د
سيکوېنټ حسابونو سره نژدې تړاو لري او د ثبوت د لټون دپاره،
هم په لاس او هم په سافټوير کښې، ځانګړې آسانتيا لري۔ په
تابلو کښې د سيکوېنټونو د ونې پر ځاے ثبوت د
فارمولونه ونه ده۔ خو د طبيعي استنتاج پر خلاف ئې قواعد د
سيکوېنټ حساب د قواعدو په شان دي۔ د تابلو ثبوت په اصل
کښې د يوه مدل څرګند لټون دے؛ ځکه کله کله ورته
\emph{معنايي} تابلو يا \emph{د صدق ونه} هم وايي۔
\olpseganchor{OLP-0684-B006}{end}

\olpseganchor{OLP-0684-B007}{start}
\emph{نښه لرونکے} تابلو د نښه لرونکو فارمولونه ونه
ده، چې \LR{$\sFmla{\True}{!A}$} يا \LR{$\sFmla{\False}{!A}$} بڼه
لرلے شي۔ ونې ښکته لور ته غځېږي۔ په هغو فارمولونه پيلېږي
چې زموږ د ثابت کولو مقصد ټاکي۔ مثلاً، د سيکوېنټ
\LR{$!A, !B \Sequent !C, !D$} د ثابت کولو دپاره به په
\LR{$\sFmla{\True}{!A}, \sFmla{\True}{!B},
\sFmla{\False}{!C}, \sFmla{\False}{!D}$} پيل کېږي۔ داسې
جوړښت چې \LR{$!A$} او \LR{$!B$} رښتيا او \LR{$!C$} او \LR{$!D$} دروغ
کړي، هغه جوړښت دے چې ښيي
\LR{$!A, !B \Sequent !C, !D$} معتبر نۀ دے۔
\olpseganchor{OLP-0684-B007}{end}

\olpseganchor{OLP-0684-B008}{start}
د تابلو پاڼه، يعنې تر ټولو لاندې غوټه، د څانګې د
غځولو په قواعدو غځېدے شي: په هماغه څانګه نورې غوټې زياتوي
يا لاندې دوې هم‌مورې غوټې ورزياتوي او څانګه په دوو نويو
څانګو وېشي۔ قواعد په هماغه څانګه کښې پر هر پاسني نښه
لرونکي فارمول کارېدلے شي۔ د کلاسيکي تابلو سيستم
\Log{Tc} د څانګې د غځولو قواعد په \olref{tab:Tc} کښې دي۔
\olpseganchor{OLP-0684-B008}{end}

\olpseganchor{OLP-0684-B009}{start}

\olpseganchor{OLP-0684-B009}{end}

\olpseganchor{OLP-0684-B010}{start}
لکه چې ښکاري، قواعد هماغه د \Log{G3c} قواعد دي خو سرچپه
شوي؛ \LR{$\True$} او \LR{$\False$} نښې په ترتيب ښيي چې
فارمول د سيکوېنټ په کيڼ يا ښي اړخ اصلي فارمول دے۔
\olpseganchor{OLP-0684-B010}{end}

\olpseganchor{OLP-0684-B011}{start}
د تابلو يوه څانګه هله \emph{تړلې} ده چې دواړه
\LR{$\sFmla{\True}{!A}$} او \LR{$\sFmla{\False}{!A}$} ولري؛ يا
\LR{$\sFmla{\True}{\lfalse}$} ولري۔ که هره څانګه تړلې وي، ټول
تابلو تړلے بولو۔
مثلاً، د \LR{$!A \lor (!B \land !C) \Sequent (!A \lor !B) \land
(!A \lor !C)$} دپاره دلته يو تړلے تابلو دے؛ يعنې په
\LR{$\sFmla{\True}{!A \lor (!B \land
!C)}$} او \LR{$\sFmla{\False}{(!A \lor !B) \land (!A \lor !C)}$}
پيلېږي:
\begin{oltableau}
  [\sFmla{\True}{\formula{A} \lor (\formula{B} \land \formula{C})},
    just=\TAss, checked
    [\sFmla{\False}{(\formula{A} \lor \formula{B}) \land
        (\formula{A} \lor \formula{C})}, just=\TAss, checked
      [\sFmla{\True}{\formula{A}}, just={\TRule{\True}{\lor}[1]}
        [\sFmla{\False}{\formula{A} \lor \formula{B}},
          just={\TRule{\False}{\land}[2]}, checked
          [\sFmla{\False}{\formula{A}},
            just={\TRule{\False}{\lor}[4]}
            [\sFmla{\False}{\formula{B}},
              just={\TRule{\False}{\lor}[4]}, close]
          ]
        ]
        [\sFmla{\False}{\formula{A} \lor \formula{C}},
          just={\TRule{\False}{\land}[2]}, checked
          [\sFmla{\False}{\formula{A}},
            just={\TRule{\False}{\lor}[4]}
            [\sFmla{\False}{\formula{C}},
              just={\TRule{\False}{\lor}[4]}, close]
          ]
        ]
      ]
      [\sFmla{\True}{\formula{B} \land \formula{C}},
        just={\TRule{\True}{\lor}[1]}, checked
        [\sFmla{\False}{\formula{A} \lor \formula{B}},
          just={\TRule{\False}{\land}[2]}, checked
          [\sFmla{\True}{\formula{B}},
            just={\TRule{\True}{\land}[3]}, move by=2
            [\sFmla{\True}{\formula{C}},
              just={\TRule{\True}{\land}[3]}
              [\sFmla{\False}{\formula{A}},
                just={\TRule{\False}{\lor}[4]}
                [\sFmla{\False}{\formula{B}},
                  just={\TRule{\False}{\lor}[4]},close
                ]
              ]
            ]
          ]
        ]
        [\sFmla{\False}{\formula{A} \lor \formula{C}},
          just={\TRule{\False}{\land}[2]}, checked
          [\sFmla{\True}{\formula{B}},
            just={\TRule{\True}{\land}[3]}, move by=2
              [\sFmla{\True}{\formula{C}},
                just={\TRule{\True}{\land}[3]}
                [\sFmla{\False}{\formula{A}},
                  just={\TRule{\False}{\lor}[4]}
                  [\sFmla{\False}{\formula{C}},
                  just={\TRule{\False}{\lor}[4]},close
                ]
              ]
            ]
          ]
        ]
      ]
    ]
  ]
\end{oltableau}
\olpseganchor{OLP-0684-B011}{end}

\olpseganchor{OLP-0684-B012}{start}
د کتنې نښې ښيي چې اړونده قاعده په فارمول باندې
په ټولو لاندې څانګو کښې کارول شوې ده۔ د \LR{$\otimes$} نښه د
څانګې تړل ښيي۔ د قواعدو ورپسې د کرښو شمېرې ښيي چې قاعده
په کومو نښه لرونکو فارمولونه کارول شوې ده؛ يعنې په
هماغه څانګه کښې پر هغه فارمول چې ښودل شوې کرښه ئې لري۔
\olpseganchor{OLP-0684-B012}{end}

\olpseganchor{OLP-0684-B013}{start}
په تابلو کښې د ثبوت لټون آسان دے۔ تابلو په پړاوونو
جوړوو۔ په پړاو \LR{$0$} کښې پيليز فارمولونه پاس وليکئ۔ په
پړاو \LR{$n+1$} کښې د هرې نۀ تړلې پاڼې غوټې دپاره داسې وکړئ:
لومړے، يعنې تر ټولو پاسنے، ناعټومي نښه لرونکے فارمول ومومئ
چې پر دې څانګه لا نۀ وي کارول شوے او د لاندې تکراري کميت ټاکنې له
نوعې نۀ وي۔ په دغه فارمول د څانګې د غځولو قاعده وکاروئ۔
کله چې قاعده په هره څانګه کښې وکارول شي چې همدغه نښه
لرونکے فارمول لري، د کتنې نښه پرې کېږدئ۔ تر هغې مخکې د هرې څانګې د قاعدو کارونه بېله ثبت کړئ۔ (پورته بېلګه په
همدې الګوريتم جوړه شوې ده۔)
\olpseganchor{OLP-0684-B013}{end}

\olpseganchor{OLP-0684-B014}{start}
د \LR{$\sFmla{\True}{\lforall}$} او
\LR{$\sFmla{\False}{\lexists}$} بڼې نښه لرونکي فارمولونه
هېڅکله د کتنې نښه نۀ اخلي، نو ځانګړې کارونه غواړي۔ که
موجود وي، بايد يقيني کړو چې \LR{$\TRule{\True}{\forall}$} او
\LR{$\TRule{\False}{\lexists}$} قواعد بالاخره پر ټولو د کارونې
وړ ترمونو وکارول شي۔ مثلاً، دا په دې تضمينېدے شي چې هر
پړاو کښې په هره څانګه پر دغو ټولو فارمولونو قاعده وکاروو،
وروسته له دې چې د دغې بڼې پرته تر ټولو پاسنے ناعټومي نښه
لرونکے فارمول وکارول شي۔ د کارونې ترم \LR{$t$} هغه لومړے
ترم~\LR{$t$} دے چې په څانګه کښې له مخکې موجودو ثابت سمبولونه
(يا \LR{$\Obj
c_0$}، که هېڅ ثابت سمبول موجود نۀ وي) او په پيليزو
فارمولونه کښې له راغلو تابعې څخه جوړېدے شي، او
اړوند نښه لرونکے فارمول \LR{$\sFmla{\True}{!B(t)}$} يا
~\LR{$\sFmla{\False}{!B(t)}$} لا په څانګه کښې موجود نۀ وي۔
که په دغه پړاو کښې داسې ترم نۀ وي، نوې غوټه نۀ زياتوو۔ په راتلونکي پړاو کښې که يو نوے ثابت راشي، د هماغې قاعدې کارونه بيا ګورو۔
\olpseganchor{OLP-0684-B014}{end}

\olpseganchor{OLP-0684-B015}{start}
که د ثبوت دا لټون تړلے تابلو جوړ نۀ کړي، نو يا يوه
متناهي پرانيستې مشبوع څانګه لري چې ټول عادي ناعټومي
فارمولونه ئې د کتنې نښې لري او تکراري کميت ټاکنه ئې پر
ټولو اوسني د کارونې وړ ترمونو کارول شوې ده؛ يا لټون يوه
نامتناهي، متناهي څانګېدونکې ونه جوړوي چې بايد يوه نامتناهي
څانګه ولري۔ لکه په \olref[cpl]{sec} کښې، داسې
جوړښت~\LR{$\Struct{M}$} تعريفولے شو چې که
\LR{$\sFmla{\True}{!A}$} په څانګه کښې راځي، \LR{$\Sat{M}{!A}$}، او
که \LR{$\sFmla{\False}{!A}$} راځي، \LR{$\Sat/{M}{!A}$} وي۔
(\LR{$\Theta =
\Setabs{!A}{\sFmla{\True}{!A} \text{\RL{ په څانګه کښې}}}$} او
\LR{$\Xi =
\Setabs{!A}{\sFmla{\False}{!A} \text{\RL{ په څانګه کښې}}}$}
واخلئ۔)
که په پيليزو فارمولونو کښې ازاد متغيرونه وي، د ثبوت د لټون د بشپړتيا په شان هر يو په بېل تازه ثابت بدلوو۔ په پرانيستې څانګه کښې د همدغو ثابتونو ارزښتونه د اصلي متغيرونو ګومارنه راکوي۔
\olpseganchor{OLP-0684-B015}{end}

\olpseganchor{OLP-0684-B016}{start}
د ثبوت د لټون په دې طريقه جوړ شوي تابلوګانې هرې غوټې
ته د سيکوېنټ په ورکولو په آسانۍ د \Log{G3c} په ثبوتونو
بدلېدے شي۔ که پيليز فارمولونه د سيکوېنټ
\LR{$\Gamma \Sequent
\Delta$} سره برابر وي، هر پيليز فارمول په
\LR{$\Gamma \Sequent
\Delta$} ونوموئ۔ که څانګه پر نښه لرونکي
فارمول \LR{$\sFmla{\True}{!A}$} د غځولو د قاعدې په کارونه
وغځېږي، د پاڼې نوم \LR{$!A, \Pi \Sequent \Lambda$} دے؛ او که
پر نښه لرونکي فارمول \LR{$\sFmla{\False}{!A}$} کارول شوې
وي، د پاڼې نوم \LR{$\Pi \Sequent
\Lambda, !A$} دے۔ هر چې تابلو د \TRule{\True}{\star} قاعده
کاروي، پر سيکوېنټ د \LeftR{\star} قاعده وکاروئ، چې \LR{$!A$}
ئې اصلي فارمول وي؛ او هر چې د \TRule{\False}{\star} قاعده
کاروي، د \RightR{\star} قاعده وکاروئ۔ د قاعدې مقدمې هغه
سيکوېنټونه دي چې په تابلو کښې ورزياتې شوې اړوندې غوټې
نوموي۔ کله چې تابلو په \LR{$\sFmla{\True}{!A}$} او
\LR{$\sFmla{\False}{!A}$} وتړل شي، د وروستۍ پاڼې نوم به
\LR{$!A, \Pi \Sequent \Lambda,
!A$} بڼه ولري۔ د تابلو ځينې پرله‌پسې غوټې به يو شان
سيکوېنټ واخلي؛ هغه پرله‌پسې ګامونه يو ځاے کړئ چې يوازې يوه قاعده تمثيلوي۔
د ګډ اټومي فارمول تړل د هويت بديهي اصل ورکوي۔ د رښتيا په نښه دروغ د کيڼ دروغ بديهي اصل ورکوي۔ که ګډ فارمول مرکب وي، د \olref[seq][ex]{prop:G3c-axioms} د هويت ثبوت د پاڼې پر ځاے ږدو۔ که يوه کارول شوې پېښه په تابلو کښې پاتې وي، د هغې اضافي نقل په سېکوېنټ کښې د کمزوري کولو او انقباض په منلو وړ قاعدو سمبالوو۔
مثلاً، پورته تابلو په دغه ثبوت بدلېږي:
\[\small
\Axiom$!A \fCenter !A, !B$
\RightLabel{\RightR{\lor}}
\UnaryInf$!A \fCenter !A \lor !B$
\Axiom$!A \fCenter !A, !C$
\RightLabel{\RightR{\lor}}
\UnaryInf$!A \fCenter !A \lor !C$
\RightLabel{\RightR{\land}}
\BinaryInf$!A \fCenter (!A \lor !B) \land (!A \lor !C)$
\Axiom$!B, !C \fCenter !A, !B$
\RightLabel{\RightR{\lor}}
\UnaryInf$!B, !C \fCenter !A \lor !B$
\RightLabel{\LeftR{\land}}
\UnaryInf$!B \land !C \fCenter !A \lor !B$
\Axiom$!B, !C \fCenter !A, !C$
\RightLabel{\RightR{\lor}}
\UnaryInf$!B, !C \fCenter !A \lor !C$
\RightLabel{\LeftR{\land}}
\UnaryInf$!B \land !C \fCenter !A \lor !C$
\RightLabel{\RightR{\land}}
\BinaryInf$!B \land !C \fCenter (!A \lor !B) \land
(!A \lor !C)$
\RightLabel{\LeftR{\lor}}
\BinaryInf$!A \lor (!B \land !C) \fCenter (!A \lor !B) \land
(!A \lor !C)$
\DisplayProof
\]
\olpseganchor{OLP-0684-B016}{end}

\olfileid{pt}{ps}{tab}
\olpseganchor{OLP-0682-B004}{start}
\begin{table}
\begin{center}
\begin{tabular}{@{}cc@{}}
\AxiomC{\sFmla{\True}{\lnot !A}}
\RightLabel{\TRule{\True}{\lnot}}
\UnaryInfC{\sFmla{\False}{!A}}
\DisplayProof
&
\AxiomC{\sFmla{\False}{\lnot !A}}
\RightLabel{\TRule{\False}{\lnot}}
\UnaryInfC{\sFmla{\True}{!A}}
\DisplayProof
\\[6ex]
\AxiomC{\sFmla{\True}{!A \land !B}}
\RightLabel{\TRule{\True}{\land}}
\UnaryInfC{\sFmla{\True}{!A}}
\noLine
\UnaryInfC{\sFmla{\True}{!B}}
\DisplayProof
&
\AxiomC{\sFmla{\False}{!A \land !B}}
\RightLabel{\TRule{\False}{\land}}
\UnaryInfC{\LR{$\sFmla{\False}{!A} \quad \mid \quad \sFmla{\False}{!B}$}}
\DisplayProof
\\[6ex]
\AxiomC{\sFmla{\True}{!A \lor !B}}
\RightLabel{\TRule{\True}{\lor}}
\UnaryInfC{\LR{$\sFmla{\True}{!A} \quad \mid \quad \sFmla{\True}{!B}$}}
\DisplayProof
&
\AxiomC{\sFmla{\False}{!A \lor !B}}
\RightLabel{\TRule{\False}{\lor}}
\UnaryInfC{\sFmla{\False}{!A}}
\noLine
\UnaryInfC{\sFmla{\False}{!B}}
\DisplayProof
\\[6ex]
\AxiomC{\sFmla{\True}{!A \lif !B}}
\RightLabel{\TRule{\True}{\lif}}
\UnaryInfC{\LR{$\sFmla{\False}{!A} \quad \mid \quad \sFmla{\True}{!B}$}}
\DisplayProof
&
\AxiomC{\sFmla{\False}{!A \lif !B}}
\RightLabel{\TRule{\False}{\lif}}
\UnaryInfC{\sFmla{\True}{!A}}
\noLine
\UnaryInfC{\sFmla{\False}{!B}}
\DisplayProof
\\[6ex]
\AxiomC{\sFmla{\True}{\lforall[x][!A(x)]}}
\RightLabel{\TRule{\True}{\forall}}
\UnaryInfC{\sFmla{\True}{!A(t)}}
\DisplayProof
&
\AxiomC{\sFmla{\False}{\lforall[x][!A(x)]}}
\RightLabel{\TRule{\False}{\lforall}}
\UnaryInfC{\sFmla{\False}{!A(a)}}
\DisplayProof
\\[6ex]
\AxiomC{\sFmla{\True}{\lexists[x][!A(x)]}}
\RightLabel{\TRule{\True}{\lexists}}
\UnaryInfC{\sFmla{\True}{!A(a)}}
\DisplayProof
&
\AxiomC{\sFmla{\False}{\lexists[x][!A(x)]}}
\RightLabel{\TRule{\False}{\lexists}}
\UnaryInfC{\sFmla{\False}{!A(t)}}
\DisplayProof
\end{tabular}
\end{center}
\caption{د \Log{Tc} قواعد۔ په \LR{$\TRule{\False}{\lforall}$} او
\LR{$\TRule{\True}{\lexists}$} کښې، ثابت سمبول~\LR{$a$} بايد په پاسنۍ
څانګه کښې موجود نۀ وي؛ يعنې د دغې څانګې دپاره نوے وي۔}
\ollabel{tab:Tc}
\end{table}
\olpseganchor{OLP-0682-B004}{end}

\olpseganchor{OLP-0664-B004}{start}
\olchapter{pt}{nat}{فطري استنتاج}
\olpseganchor{OLP-0664-B004}{end}

\olpseganchor{OLP-0663-B004}{start}
\olfileid{pt}{nat}{int}
\olpseganchor{OLP-0663-B004}{end}

\olpseganchor{OLP-0663-B005}{start}
\olsection{پېژندنه}
\olpseganchor{OLP-0663-B005}{end}

\olpseganchor{OLP-0663-B006}{start}
فطري استنتاج د ثبوت د سيستمونو يوه کورنۍ ده چې په هغې کښې هر
منطقي تړونکے يا کميت ټاکونکے د استنتاجونو يوه جوړه لري: يو استنتاج
تړونکے «داخلوي» (د داخلولو قاعده) او بل ئې «وباسي»
(د ايستلو قاعده)۔ د بېلګې په توګه، د \Intro{\lif}
پايله د~\LR{$!A \lif !B$} په بڼه فارمولونه لري او د \Elim{\lif}
مقدمې د~\LR{$!A \lif !B$} په بڼه فارمولونه لري۔
\olpseganchor{OLP-0663-B006}{end}

\olpseganchor{OLP-0663-B007}{start}
د فطري استنتاج لومړۍ دوه بڼې په ۱۹۳۰مو کلونو کښې ګېرهارډ
ګېنټسن او ستانيسواف ياشکووسکي وړاندې کړې۔ د ثبوت د تيورۍ پوهان
تر ډېره د ګېنټسن سيستم څېړي؛ د ياشکووسکي سيستم بيا هغو سيستمونو
ته لار هواره کړه چې په تدريس کښې تر ټولو ډېر کارېږي۔ په دواړو
سيستمونو کښې له \emph{فرضيو} څخه پيل کولے شو: داسې فارمولونه
چې ثابت کېدو ته اړتيا نۀ لري، خو له هغو څخه د نورو
فارمولونو په ثابتولو کښې کار اخيستلے شو۔ ټول ثبوت په
داسې فرضيو پورې تړلے کېدے شي۔ دا چې فرضيې ثابتې نۀ دي، نو
ثبوت خپله پايله (خپل \emph{پاييز فارمول}) نۀ ثابتوي،
بلکې يوازې دا ثابتوي چې پايله له فرضيو څخه راځي۔
\olpseganchor{OLP-0663-B007}{end}

\olpseganchor{OLP-0663-B008}{start}
د استنتاج ځينې قاعدې داسې دي چې پايلې ئې بيا په ځينو فرضيو
پورې نۀ تړل کېږي: دغه فرضيې «وباسي»۔ د بېلګې په توګه،
\Intro{\lif} د فرضيې~\LR{$!A$} په مرسته د پايلې~\LR{$!B$} يو
ثبوت اخلي او پايله ئې \LR{$!A \lif !B$} وي۔ په~\LR{$!A \lif !B$} ختمېدونکے
ثبوت بيا په~\LR{$!A$} پورې نۀ تړل کېږي؛ فرضيه \LR{$!A$} ايستل
شوې ده۔ د ګېنټسن په سيستمونو کښې ثبوتونه د
فارمولونو ونې دي، فرضيې تر ټولو بره فارمولونه
دي، او د~\Intro{\lif} غوندې قاعدې داسې نښې لري چې ايستل شوې
فرضيې ښيي۔ د ياشکووسکي په سيستمونو کښې د ثبوت هغه برخې
چې په يوه فرضيه پورې تړلې وي، په يو ډول بېلې ښودل کېږي؛ مثلاً
له يوې ولاړې کرښې وروسته لږ دننه ليکل کېږي يا په يوه بکس کښې
راځي۔ ايستل کېدونکې فرضيه~\LR{$!A$} او د شرط تړنې تالي \LR{$!B$}
په بکس کښې لومړۍ او وروستۍ کرښه وي، او د \Intro{\lif}
پايله~\LR{$!A \lif !B$} له بکس څخه بهر وي۔
\olpseganchor{OLP-0663-B008}{end}

\olpseganchor{OLP-0663-B009}{start}
دغه سيستمونه ځکه «فطري» دي چې د استنتاج قاعدې ئې زموږ
(يا لږ تر لږه د رياضي پوهانو) د فطري استدلال ډول ښيي۔ د يوه
فرضي حاصل (يوې شرط تړنې) ثابتولو دپاره مقدم فرضوو او په هغې
تالي ثابتوو۔ همدا \Intro{\lif} ده۔ چې شرط تړنه \LR{$!A \lif !B$} او
د هغې مقدم~\LR{$!A$} راته معلوم وي، پايله~\LR{$!B$} اخلو۔ همدا
\Elim{\lif} ده۔ د حالتونو له مخې ثبوت په دې کاروو چې وښيو
د فصل تړنې~\LR{$!A \lor !B$} په فرضولو يوه پايله راځي۔ همدا \Elim{\lor}
ده۔ د عمومي دعوې~\LR{$\lforall[x][!A(x)]$} ثابتولو دپاره ثابتوو چې د يوه
خپل‌سري~\LR{$c$} دپاره \LR{$!A(c)$} رښتيا ده۔ همدا \Intro{\lforall}
ده، او داسې نور۔
\olpseganchor{OLP-0663-B009}{end}

\olpseganchor{OLP-0663-B010}{start}
د بېلګې په توګه، د ګېنټسن په طرز فطري استنتاج کښې
د~\LR{$!A \lif (!A \lor !B)$} يو ساده ثبوت داسې دے:
\[
\AxiomC{\LR{$\Discharge{!A}{x}$}}
\RightLabel{\Intro{\lor}}
\UnaryInfC{\LR{$!A \lor !B$}}
\DischargeRule{\Intro{\lif}}{x}
\UnaryInfC{\LR{$!A \lif (!A \lor !B)$}}
\DisplayProof
\]
دا له فرضيې~\LR{$!A$} څخه پيلېږي، په \Intro{\lor} ئې
پايله~\LR{$!A
\lor !B$} اخلي، او بيا په \Intro{\lif} ئې
پايله~\LR{$!A \lif (!A \lor !B)$} اخلي۔ د \Intro{\lif} استنتاج
فرضيه~\LR{$!A$} وباسي؛ نښه~\LR{$x$} همدا ښيي۔
\olpseganchor{OLP-0663-B010}{end}

\olpseganchor{OLP-0663-B011}{start}
د ثبوت د تيورۍ پوهان د ګېنټسن اصلي سيستمونه غوره کوي چې د
فارمولونو پر ونو کار کوي، خو يو بدل ډول ئې هم
د ثبوت په تيورۍ کښې مهم دے۔ ثبوت چې له فرضيو~\LR{$\Gamma$}
څخه~\LR{$!A$} ثابتوي، ښيي چې \LR{$!A$} له~\LR{$\Gamma$} څخه راځي؛ يعنې
\LR{$\Gamma \Entails !A$}۔ د فطري استنتاج قاعدې په فطري ډول د دې ډول منطقي
لازميتونو په هکله لوستل کېدے شي۔ مثلاً \Intro{\lif} وايي چې
که \LR{$\Gamma + !A
\Entails !B$} وي، نو \LR{$\Gamma \Entails !A \lif !B$} دے۔ دا فطري ده چې قاعدې نېغ په دې
ډول جوړې کړو چې د استنتاج د قاعدې په مقدمو او پايله کښې پر
فرضيو او له هغو پر راوتونکې پايله دواړو کار وکړي؛ يعنې پر
سېکوېنټونو \LR{$\Gamma \Sequent !A$} کار وکړي۔ پورته ساده ثبوت په داسې سيستم
کښې داسې ښکاري:
\[
\Axiom$x:!A \fCenter !A$
\RightLabel{\Intro{\lor}}
\UnaryInf$x:!A \fCenter !A \lor !B$
\DischargeRule{\Intro{\lif}}{x}
\UnaryInf$\fCenter !A \lif (!A \lor !B)$
\DisplayProof
\]
لکه څنګه چې وينئ، قاعدې هماغه دي: د~\LR{$\Sequent$} ښي لاس پر فارمول
کار کوي؛ چپ لاس فارمولونه يوازې هغه فرضيې ښيي چې
هر ګام په هغو پورې تړلے دے۔
\olpseganchor{OLP-0663-B011}{end}

\olpseganchor{OLP-0663-B012}{start}
د داخلولو/ايستلو د قاعدو ټولې جوړې په
خپل سر د کلاسيکي منطق دپاره بشپړې نۀ دي۔ دوه نورې قاعدې
ورزياتول غواړي چې په~\LR{$\lfalse$} پورې اړه لري۔ لومړۍ
«د شهودي محال قاعده»~\FalseInt، يا «چاودنه» ده، چې وايي
له \LR{$\lfalse$} څخه هر څه استنتاجولے شو۔ دويمه د غيرمستقيم ثبوت
کلاسيکي اصل~\FalseCl{} (يا «د کلاسيکي محال قاعده») ده،
چې وايي که له~\LR{$\lnot !A$} څخه~\LR{$\lfalse$} ثابت کولے شو، نو
\LR{$!A$} مو ثابت کړے دے۔ که \FalseCl پرېږدو، د شهودي
منطق مناسب سيستم ترې جوړېږي۔ که دواړه پرېږدو، هغه سيستم
ترې جوړېږي چې «لږ تر لږه منطق» نومېږي۔
\olpseganchor{OLP-0663-B012}{end}

\olpseganchor{OLP-0663-B013}{start}
موږ به د کلاسيکي او شهودي منطق دپاره د ګېنټسن په طرز
فطري استنتاج وڅېړو۔ دا د \Log{N1c} او \Log{N1i} سيستمونه
دي۔ د \Log{N2c} او \Log{N2i} سيستمونه ئې ورته بڼې دي چې
د يوازېنيو فارمولونو~\LR{$!A$} پر ځاے پر
سېکوېنټونو \LR{$\Gamma \Sequent !A$} کار کوي۔
\olpseganchor{OLP-0663-B013}{end}

\olpseganchor{OLP-0668-B004}{start}
\olfileid{pt}{nat}{rp}
\olpseganchor{OLP-0668-B004}{end}

\olpseganchor{OLP-0668-B005}{start}
\olsection{قاعدې او ثبوتونه}
\olpseganchor{OLP-0668-B005}{end}

\olpseganchor{OLP-0668-B006}{start}
لومړے ځينې تعريفونه ورکوو او ځينې اصطلاحات ټاکو۔
\olpseganchor{OLP-0668-B006}{end}

\olpseganchor{OLP-0668-B007}{start}
د \Log{N1c} او~\Log{N1i} هغه دوه قاعدې وګورئ چې د \LR{$!A
\lif !B$} يادونه کوي۔
\[
\bottomAlignProof
\AxiomC{\LR{$\Discharge{!A}{x}$}}
\shortDeduce
\DeduceC{\LR{$!B$}}
\DischargeRule{\Intro{\lif}}{x}
\UnaryInfC{\LR{$!A \lif !B$}}
\DisplayProof
\qquad
\bottomAlignProof
\AxiomC{\LR{$!A \lif !B$}}
\AxiomC{\LR{$!A$}}
\RightLabel{\Elim{\lif}}
\BinaryInfC{\LR{$!B$}}
\DisplayProof
\]
د \Elim{\lif} قاعده آسانه ده۔ له کرښې بره فارمولونه
\emph{مقدمې} دي۔ کيڼ اړخ ته فارمول هماغه \emph{اصلي}
فارمول~\LR{$!A \lif !B$} دے۔ ټولې د ايستلو قاعدې
يوه داسې مقدمه لري چې تل په تر ټولو کيڼ اړخ کښې وي۔ دې ته د
استنتاج \emph{لويه} مقدمه وايي۔ که نورې مقدمې هم وي، لکه په دې
حالت کښې، هغو ته \emph{کوچنۍ} مقدمې وايي۔ د داخلولو
قاعدو مقدمه يا مقدمې هم لويې مقدمې بلل کېږي۔ له کرښې لاندې
فارمول \emph{پايله} ده۔
\olpseganchor{OLP-0668-B007}{end}

\olpseganchor{OLP-0668-B008}{start}
د \Intro{\lif} قاعده يوازې يوه مقدمه لري، او دلته پايله
اصلي فارمول~\LR{$!A \lif !B$} دے۔ د داخلولو ټولې قاعدې
همداسې دي: پايله اصلي فارمول دے؛ د هغه لويه عملګره هماغه ده چې
``داخلېږي''۔
\olpseganchor{OLP-0668-B008}{end}

\olpseganchor{OLP-0668-B009}{start}
د \LR{$\Discharge{!A}{x}$} نښه دا معنا لري۔ په ثبوت کښې
د \Intro{\lif} استنتاج فارمول~\LR{$!B$} د مقدمې په توګه اخلي۔
په~\LR{$!B$} ختمېدونکے فرعي ثبوت يو شمېر پيلني فارمولونه
لري۔ دې ته \emph{فرضونه} وايي۔ ثبوت چې له يوه يا ډېرو
د~\LR{$!A$} په بڼه فرضونو څخه~\LR{$!B$} ثابتوي، ثبوت دے چې \LR{$!A$}
څخه~\LR{$!B$} نتيجه کېږي۔ چې د \Intro{\lif} قاعده وکاروو،~\LR{$!A
\lif !B$} نتيجه اخلو او پايله نور په فرض~\LR{$!A$} پورې نۀ تړل کېږي۔
د دې ښودلو دپاره، په هغه ثبوت کښې چې د~\Intro{\lif}
استنتاج هم لري، د~\LR{$!A$} په بڼه فرضونه \emph{ايستل کېږي} يا
\emph{بندېږي}۔ قاعده ټاکي چې کوم فرضونه ايستل کېدے شي۔ د
\Intro{\lif} په حالت کښې چې پايله ئې \LR{$!A \lif !B$} وي، دا
د~\LR{$!A$} په بڼه فرضونه دي، يعنې د شرطي فارمول له مخکينۍ برخې سره
يو شان دي۔ کله کله بايد وښيو چې په کومه قاعده کوم فرضونه ايستل
کېږي؛ د فرض ايستلو نښه~\LR{$x$} همدا ښيي۔ د فرضونو ايستونکې قاعدې
\emph{اړ نۀ دي} چې د اړينې بڼې \emph{ټول} فرضونه يا ان يو فرض
هم وباسي۔ مثلاً، دا سم دے، که څه هم لومړے \Intro{\lif}
هېڅ فرض نۀ باسي:
\[
\AxiomC{\LR{$\Discharge{!A}{y}$}}
\DischargeRule{\Intro{\lif}}{x}
\UnaryInfC{\LR{$!B \lif !A$}}
\DischargeRule{\Intro{\lif}}{y}
\UnaryInfC{\LR{$!A \lif (!B \lif !A)$}}
\DisplayProof
\]
چې قاعده هېڅ فرض نۀ باسي، د فرض ايستلو نښه اختياري ده۔ په دې
رسم کښې لومړۍ قاعده له خپلې نښې سره هم هېڅ فرض نۀ باسي۔
\olpseganchor{OLP-0668-B009}{end}

\olpseganchor{OLP-0668-B010}{start}
په دې ثبوت کښې
\[
\AxiomC{\LR{$\Discharge{!A}{x}$}}
\AxiomC{\LR{$\Discharge{!A}{y}$}}
\RightLabel{\Intro{\land}}
\BinaryInfC{\LR{$!A \land !A$}}
\RightLabel{\Elim{\land}}
\UnaryInfC{\LR{$!A$}}
\DischargeRule{\Intro{\lif}}{x}
\UnaryInfC{\LR{$!A \lif !A$}}
\DischargeRule{\Intro{\lif}}{y}
\UnaryInfC{\LR{$!A \lif (!A \lif !A)$}}
\DisplayProof
\]
لومړے \Intro{\lif} استنتاج کيڼ فرض \LR{$!A^x$} باسي او دويم
ښي فرض~\LR{$!A^y$} باسي۔ يعنې ټول په~\LR{$x$} نښه شوي فرضونه هغه
\Intro{\lif} باسي چې نښه ئې~\LR{$x$} وي، او ټول په~\LR{$y$} نښه شوي
فرضونه هغه باسي چې نښه ئې~\LR{$y$} وي۔ د دې د سم کار دپاره بايد
ټول په يوه نښه شوي فرضونه هماغه يو فارمول وي۔
\olpseganchor{OLP-0668-B010}{end}

\olpseganchor{OLP-0668-B011}{start}
د کميت قاعدې لږې پېچلې دي۔ مثلاً، په \Log{N1i} کښې د
\LR{$\lexists$} قاعدې دا دي:
\[
\bottomAlignProof
\AxiomC{\LR{$!A(t)$}}
\RightLabel{\Intro{\lexists}}
\UnaryInfC{\LR{$\lexists[x][!A(x)]$}}
\DisplayProof
\qquad
\bottomAlignProof
\AxiomC{\LR{$\lexists[x][!A(x)]$}} 
\AxiomC{\LR{$\Discharge{!A(c)}{x}$}}
\shortDeduce
\DeduceC{\LR{$!B$}}
\DischargeRule{\Elim{\lexists}}{x}
\BinaryInfC{\LR{$!B$}} \DisplayProof 
\bottomAlignProof
\] په
\Intro{\lexists} قاعده کښې مقدمه~\LR{$!A(t)$} ده، چې \LR{$t$}~يو
تړلے ترم دے، يعنې داسې ترم چې متغيرونه نۀ لري۔ دا استنتاج سم
دے، يعنې د \Intro{\lexists} له شېمايي قاعدې سره سمون لري، که
يو فارمول~\LR{$!A$}، يو متغير~\LR{$x$} او يو تړلے ترم~\LR{$t$} داسې وي
چې پايله \LR{$\lexists[x][!A]$} او مقدمه~\LR{$\Subst{!A}{t}{x}$} وي۔
د \Elim{\lexists} قاعده د~\LR{$!A(c)$} په بڼه فرضونه ايستلو ته
اجازت راکوي؛ يعنې د هر استنتاج دپاره يو ثابت سمبول~\LR{$c$} وي او
د \LR{$\Subst{!A}{c}{x}$} په بڼه فرضونه ايستل کېدے شي۔ د يو
استنتاج ټول ايستل شوي فرضونه بايد هماغه~\LR{$c$} وکاروي۔ دې ثابت
ته د \Elim{\lexists} استنتاج \emph{ځانګړے متغير} وايي۔
استنتاج يوازې هغه وخت سم دے چې \LR{$c$}~له ايستل شوي~\LR{$!A(c)$} پرته
په بل پرانيستي فرض کښې نۀ راځي، په~\LR{$!B$} کښې نۀ راځي او
په~\LR{$!A$} کښې نۀ راځي۔ دې ته \emph{د ځانګړي متغير شرط} وايي۔
\olpseganchor{OLP-0668-B011}{end}

\olpseganchor{OLP-0668-B012}{start}
په \Log{N1c} او~\Log{N1i} کښې ثبوتونه د فارمولونو
متناهي ونې دي چې له نښه لرونکو فرضونو څخه په استقرايي ډول د
استنتاج د قواعدو په وسيله جوړېږي۔ هره پاڼه يو فرض دے، چې کېدے
شي د فرض ايستلو نښه ولري، او هر بل فارمول د نظام د يوې
استنتاجي قاعدې په وسيله له هغو فارمولونه څخه نتيجه کېږي چې
نېغ ئې بره وي۔ که له استنتاج بره په ونه کښې د فرض په توګه د
فارمول يو راتګ له هغه بره استنتاج نۀ وي ايستلے، هغه
په دغه استنتاج کښې \emph{پرانيستے} بلل کېږي۔
\olpseganchor{OLP-0668-B012}{end}

\olpseganchor{OLP-0668-B013}{start}
\begin{defn}\ollabel{defn:proof}
په \Log{N1c} او~\Log{N1i} کښې \emph{ثبوتونه} د
فارمولونه متناهي ونې دي چې په استقرايي ډول داسې جوړېږي:
\begin{enumerate}
  \item\ollabel{defn:prf:assumption} هر فارمول چې د فرض
  ايستلو نښه لکه \LR{$x$} يا \LR{$y$} ولري، په ځانله ډول ثبوت دے۔
  په هغه ثبوت کښې چې يوازې يو فارمول لري، دغه
  فارمول پرانيستے فرض ګڼل کېږي۔
  \item\ollabel{defn:prf:one-prem} که \LR{$R$} يوه قاعده وي چې يوه،
  دوه يا درې مقدمې \LR{$!A_1$}، \LR{$!A_2$}، \LR{$!A_3$} او پايله~\LR{$!A$}
  ولري، او \LR{$\delta_1$}، \LR{$\delta_2$}، \LR{$\delta_3$} هغه ثبوتونه
  وي چې په ترتيب په \LR{$!A_1$}، \LR{$!A_2$}، \LR{$!A_3$} ختمېږي، نو
  په ترتيب دا
  \[
  \AxiomC{}
  \RightLabel{\LR{$\delta_1$}}
  \DeduceC{\LR{$!A_1$}}
  \RightLabel{\LR{$R$}}
  \UnaryInfC{\LR{$!A$}}
  \DisplayProof
\qquad
  \AxiomC{}
  \RightLabel{\LR{$\delta_1$}}
  \DeduceC{\LR{$!A_1$}}
  \AxiomC{}
  \RightLabel{\LR{$\delta_2$}}
  \DeduceC{\LR{$!A_2$}}
  \RightLabel{\LR{$R$}}
  \BinaryInfC{\LR{$!A$}}
  \DisplayProof
\qquad
  \AxiomC{}
  \RightLabel{\LR{$\delta_1$}}
  \DeduceC{\LR{$!A_1$}}
  \AxiomC{}
  \RightLabel{\LR{$\delta_2$}}
  \DeduceC{\LR{$!A_2$}}
  \AxiomC{}
  \RightLabel{\LR{$\delta_3$}}
  \DeduceC{\LR{$!A_3$}}
  \RightLabel{\LR{$R$}}
  \TrinaryInfC{\LR{$!A$}}
  \DisplayProof
  \]
  ثبوت دے، په دې شرط چې د ثبوتونه د فرضونو نښې
  سازګارې وي، يعنې دوه بېل فارمولونه د فرض ايستلو يوه نښه
  و نۀ لري، او استنتاجونه د~\LR{$R$} سمې کارونې وي۔
  \item په نوي ثبوت کښې د تر ټولو بره فارمول راتګونه
  د ثبوت \emph{پرانيستي فرضونه} دي، که د \LR{$\delta_1$}،
  \LR{$\delta_2$} يا~\LR{$\delta_3$} پرانيستي فرضونه وي، پرته له هغو
  راتګونو چې قاعده ئې باسي۔ که د قاعدې نښه~\LR{$x$} يا \LR{$x\, y$}
  وي، نو د \Intro{\lif} او~\Intro{\lnot} دپاره دا په
  \LR{$\delta_1$} کښې ټول په~\LR{$x$} نښه شوي پرانيستي فرضونه دي؛ د
  \Elim{\lexists} دپاره په~\LR{$\delta_2$} کښې ټول په~\LR{$x$} نښه
  شوي پرانيستي فرضونه دي؛ او د \Elim{\lor} دپاره په
  \LR{$\delta_2$} کښې ټول په~\LR{$x$} نښه شوي او په~\LR{$\delta_3$} کښې
  ټول په~\LR{$y$} نښه شوي پرانيستي فرضونه دي۔
\end{enumerate}
\end{defn}
\olpseganchor{OLP-0668-B013}{end}

\olpseganchor{OLP-0668-B014}{start}
هغه ثبوتونه چې د ثبوت په استقرايي جوړښت کښې
کارول کېږي، د هغه \emph{فرعي ثبوتونه} بلل کېږي۔ هغه فرعي
ثبوتونه چې د يو استنتاج په مقدمو ختمېږي، د هغه استنتاج
\emph{سمدستي فرعي ثبوتونه} بلل کېږي۔ د \Log{N1c} او
\Log{N1i} قاعدې په \olref{tab:N1} کښې ورکړل شوې دي۔
\olpseganchor{OLP-0668-B014}{end}

\olpseganchor{OLP-0668-B015}{start}

\olpseganchor{OLP-0668-B015}{end}

\olpseganchor{OLP-0668-B016}{start}
\begin{defn}
د ثبوت~\LR{$\delta$} \emph{لوړوالے}~\LR{$\pheight{\delta}$}
په استقرايي ډول داسې تعريفېږي۔
\begin{enumerate}
  \item که \LR{$\delta$} يوازې يو فرض ولري، \LR{$\pheight{\delta} = 0$}۔
  \item که \LR{$\delta$} په يوه مقدمه لرونکي استنتاج ختمېږي چې
  سمدستي فرعي ثبوت ئې~\LR{$\delta_1$} وي، نو
  \LR{$\pheight{\delta} = \pheight{\delta_1} + 1$}۔
  \item که \LR{$\delta$} په دوو مقدمو لرونکي استنتاج ختمېږي چې
  سمدستي فرعي ثبوتونه ئې~\LR{$\delta_1$} او~\LR{$\delta_2$} وي، نو
  \LR{$\pheight{\delta} =
  \max(\pheight{\delta_1}, \pheight{\delta_2}) + 1$}۔
  \item که \LR{$\delta$} په \Elim{\lor} ختمېږي چې سمدستي فرعي
  ثبوتونه ئې~\LR{$\delta_1$}، \LR{$\delta_2$} او~\LR{$\delta_3$} وي، نو
  \LR{$\pheight{\delta} =
  \max(\pheight{\delta_1}, \pheight{\delta_2}, \pheight{\delta_3}) + 1$}۔
\end{enumerate}
که سېکوېنټ~\LR{$S$} د \LR{$\le n$} لوړوالي ثبوت ولري،
\LR{$\Proves[n] S$} ليکو۔
\end{defn}
\olpseganchor{OLP-0668-B016}{end}

\olpseganchor{OLP-0668-B017}{start}
\begin{ex}
په \Log{N1i} کښې هر فارمول له خپل ځانه د پرانيستي فرض
په توګه ثبوت ګڼل کېږي۔ نو
\[
!A \land !B^x
\]
ثبوت~\LR{$\delta_1$} دے۔ لوړوالے ئې~\LR{$0$} دے۔ له
\LR{$\delta_1$} څخه د~\Elim{\land} د دوو ډولونو په کارولو دوه
ثبوتونه~\LR{$\delta_2$} او~\LR{$\delta_3$} ترلاسه کولے شو:
\[
\AxiomC{\LR{$!A \land !B^x$}}
\RightLabel{\Elim{\land}[2]}
\UnaryInfC{\LR{$!B$}}
\DisplayProof
\qquad
\AxiomC{\LR{$!A \land !B^x$}}
\RightLabel{\Elim{\land}[1]}
\UnaryInfC{\LR{$!A$}}
\DisplayProof
\]
په~\LR{$\delta_2$} او~\LR{$\delta_3$} کښې د وروستي استنتاج سمدستي
فرعي ثبوتونه هماغه يو دي: \LR{$!A \land !B$}۔ دواړه
ثبوتونه د \LR{$!A\land !B$} راتګ د پرانيستي فرض په توګه لري۔
\LR{$\pheight{\delta_2} = \pheight{\delta_3} = 1$} لرو۔
\olpseganchor{OLP-0668-B017}{end}

\olpseganchor{OLP-0668-B018}{start}
د \Intro{\land} قاعده د~\LR{$!B \land !A$} په بڼه پايلې
دپاره د \LR{$!B$} او \LR{$!A$} په بڼه دوه مقدمې غواړي۔ نو دا
ثبوت~\LR{$\delta_4$} دے:
\[
\AxiomC{\LR{$!A \land !B^x$}}
\RightLabel{\Elim{\land}[2]}
\UnaryInfC{\LR{$!B$}}
\LeftSubproofLabel{\LR{$\delta_2$}}
\AxiomC{\LR{$!A \land !B^x$}}
\RightLabel{\Elim{\land}[1]}
\UnaryInfC{\LR{$!A$}}
\RightSubproofLabel{\LR{$\delta_3$}}
\RightLabel{\Intro{\land}}
\BinaryInfC{\LR{$!B \land !A$}}
\DisplayProof
\]
لوړوالے ئې \LR{$\pheight{\delta_4} = \max(\pheight{\delta_2},
\pheight{\delta_3})+1 = \max(1,1) + 1 = 2$} دے۔ د فرضونو
په توګه د \LR{$!A \land !B$} دوه راتګونه لري او دواړه پرانيستي دي۔
\olpseganchor{OLP-0668-B018}{end}

\olpseganchor{OLP-0668-B019}{start}
په پاے کښې د \Intro{\lif} په کارولو \LR{$\delta_5$} ترلاسه
کولے شو:
\[
\AxiomC{\LR{$\Discharge{!A \land !B}{x}$}}
\RightLabel{\Elim{\land}[2]}
\UnaryInfC{\LR{$!B$}}
\LeftSubproofLabel{\LR{$\delta_2$}}
\AxiomC{\LR{$\Discharge{!A \land !B}{x}$}}
\RightLabel{\Elim{\land}[1]}
\UnaryInfC{\LR{$!A$}}
\RightSubproofLabel{\LR{$\delta_3$}}
\RightLabel{\Intro{\land}}
\BinaryInfC{\LR{$!B \land !A$}}
\DischargeRule{\Intro{\lif}}{x}
\UnaryInfC{\LR{$(!A \land !B) \lif (!B \land !A)$}}
\DisplayProof
\]
لوړوالے ئې \LR{$\pheight{\delta_4} + 1 = 3$} دے۔ د
\Intro{\lif} استنتاج د \LR{$!A \land !B$} په بڼه ټول په~\LR{$x$}
نښه شوي فرضونه باسي، يعنې د سمدستي فرعي ثبوت دواړه پخوا
پرانيستي فرضونه۔ دا چې د \LR{$\delta_5$} يوازې همدغه فرضونه دي،
نو پرانيستے فرض نۀ لري، او د خپل پاييز فارمول
\LR{$(!A \land !B) \lif (!B \land !A)$} \emph{ثبوت} دے، يعنې
ثبوت دے۔
\end{ex}
\olpseganchor{OLP-0668-B019}{end}

\olfileid{pt}{nat}{rp}
\olpseganchor{OLP-0666-B004}{start}
\begin{table}
\begin{center}
\begin{tabular}[t]{@{}cc@{}}
\bottomAlignProof
\AxiomC{\LR{$\Discharge{!A}{x}$}}
\shortDeduce
\DeduceC{\LR{$!B$}}
\DischargeRule{\Intro{\lif}}{x}
\UnaryInfC{\LR{$!A \lif !B$}}
\DisplayProof
&
\bottomAlignProof
\AxiomC{\LR{$!A \lif !B$}}
\AxiomC{\LR{$!A$}}
\RightLabel{\Elim{\lif}}
\BinaryInfC{\LR{$!B$}}
\DisplayProof
\\[4ex]
\bottomAlignProof
\AxiomC{\LR{$!A$}}
\AxiomC{\LR{$!B$}}
\RightLabel{\Intro{\land}}
\BinaryInfC{\LR{$!A \land !B$}}
\DisplayProof
&
\begin{tabular}[b]{@{}l@{}}
\AxiomC{\LR{$!A \land !B$}}
\RightLabel{\Elim{\land}[1]}
\UnaryInfC{\LR{$!A$}}
\DisplayProof
\\[3ex]
\bottomAlignProof
\AxiomC{\LR{$!A \land !B$}}
\RightLabel{\Elim{\land}[2]}
\UnaryInfC{\LR{$!B$}}
\DisplayProof
\end{tabular}
\\[4ex]
\begin{tabular}{@{}l@{}}
\AxiomC{\LR{$!A$}}
\RightLabel{\Intro{\lor}[1]}
\UnaryInfC{\LR{$!A \lor !B$}}
\DisplayProof
\\[3ex]
\AxiomC{\LR{$!B$}}
\RightLabel{\Intro{\lor}[2]}
\UnaryInfC{\LR{$!A \lor !B$}}
\DisplayProof
\end{tabular}
&
\AxiomC{\LR{$!A \lor !B$}}
\AxiomC{\LR{$\Discharge{!A}{x}$}}
\shortDeduce
\DeduceC{\LR{$!C$}}
\AxiomC{\LR{$\Discharge{!B}{y}$}}
\shortDeduce
\DeduceC{\LR{$!C$}}
\DischargeRule{\Elim{\lor}}{x\,y}
\TrinaryInfC{\LR{$!C$}}
\DisplayProof
\bottomAlignProof
\\[6ex]
\bottomAlignProof
\AxiomC{\LR{$\Discharge{!A}{x}$}}
\shortDeduce
\DeduceC{\LR{$\lfalse$}}
\DischargeRule{\Intro{\lnot}}{x}
\UnaryInfC{\LR{$\lnot !A$}}
\DisplayProof
&
\bottomAlignProof
\AxiomC{\LR{$\lnot !A$}}
\AxiomC{\LR{$!A$}}
\RightLabel{\Elim{\lnot}}
\BinaryInfC{\LR{$\lfalse$}}
\DisplayProof
\\[4ex]
\bottomAlignProof
\AxiomC{\LR{$!A(a)$}}
\RightLabel{\Intro{\lforall}}
\UnaryInfC{\LR{$\lforall[x][!A(x)]$}}
\DisplayProof
&
\bottomAlignProof
\AxiomC{\LR{$\lforall[x][!A(x)]$}}
\RightLabel{\Elim{\lforall}}
\UnaryInfC{\LR{$!A(t)$}}
\DisplayProof
\\[4ex]
\bottomAlignProof
\AxiomC{\LR{$!A(t)$}}
\RightLabel{\Intro{\lexists}}
\UnaryInfC{\LR{$\lexists[x][!A(x)]$}}
\DisplayProof
&
\bottomAlignProof
\AxiomC{\LR{$\lexists[x][!A(x)]$}}
\AxiomC{\LR{$\Discharge{!A(a)}{x}$}}
\shortDeduce
\DeduceC{\LR{$!C$}}
\DischargeRule{\Elim{\lexists}}{x}
\BinaryInfC{\LR{$!C$}}
\DisplayProof
\\[4ex]
\bottomAlignProof
\AxiomC{\LR{$\lfalse$}}
\RightLabel{\FalseInt}
\UnaryInfC{\LR{$!A$}}
\DisplayProof
&
\bottomAlignProof
\AxiomC{\LR{$\Discharge{\lnot!A}{x}$}}
\shortDeduce
\DeduceC{\LR{$\lfalse$}}
\DischargeRule{\FalseCl}{x}
\UnaryInfC{\LR{$!A$}}
\DisplayProof
\end{tabular}
\end{center}
\caption{د \Log{N1i} او \Log{N1c} قاعدې۔ په \LR{$\Intro{\lforall}$} کښې
ثابت سمبول~\LR{$a$} بايد نۀ په پايله کښې راشي او نۀ په کوم پرانيستي
فرض کښې۔ په \LR{$\Elim{\lexists}$} کښې \LR{$a$}~بايد نۀ په مقدمو کښې
راشي او نۀ په کوم پاتې پرانيستي فرض کښې۔ \Log{N1i} د \FalseCl
لرونکے نۀ دے۔}
\ollabel{tab:N1}
\end{table}
\olpseganchor{OLP-0666-B004}{end}

\olpseganchor{OLP-0669-B004}{start}
\olfileid{pt}{nat}{seq}
\olpseganchor{OLP-0669-B004}{end}

\olpseganchor{OLP-0669-B005}{start}
\olsection{له سېکوېنټونو سره فطري استنتاج: \texorpdfstring{\Log{N2c}}{N2c} او~\texorpdfstring{\Log{N2i}}{N2i}}
\olpseganchor{OLP-0669-B005}{end}

\olpseganchor{OLP-0669-B006}{start}
په \Log{N1c} يا \Log{N1i} کښې ثبوت~\LR{$\delta$}
ښيي چې پاييز فارمول ئې~\LR{$!A$} له پرانيستو فرضونو څخه
نتيجه کېږي۔ که \LR{$\Gamma$} د هغو فارمولونه~\LR{$!B$} سټ وي چې
\LR{$!B^x$} ئې په~\LR{$\delta$} کښې د پرانيستي فرض په توګه راځي، نو
دا په \LR{$\delta \Proves \Gamma
\Sequent !A$} ليکلے شو۔ دا ښيي چې د سېکوېنټونو دپاره هم يو
فطري استنتاجي نظام جوړولے شو۔ په داسې نظام کښې هر سېکوېنټ دا
معلومات لري چې ښي اړخ ته فارمول په کومو پرانيستو فرضونو
پورې تړلے دے، يعنې کوم فرضونه په هغه فرعي ثبوت کښې
پرانيستي دي چې په دغه سېکوېنټ ختمېږي۔ دا له سېکوېنټونو سره
فطري استنتاج، يا په ``سېکوېنټي طرز'' فطري استنتاج دے؛ پايله
شوي نظامونه \Log{N2c} او~\Log{N2i} بولو۔
\olpseganchor{OLP-0669-B006}{end}

\olpseganchor{OLP-0669-B007}{start}
د \Log{N1} او \Log{N2} د لا نېغ سمون دپاره، کيڼ اړخ ته
سټونه~\LR{$\Gamma$} د نښه لرونکو فارمولونه~\LR{$x:!B$} څخه جوړ
اخلو۔ په \Log{N1} کښې قاعدې يوازې د مقدمو پر فارمولونه
کار کوي، يعنې د سمدستي فرعي ثبوتونه پر پاييزو
فارمولونه؛ نو د \Log{N2} قاعدې يوازې د سېکوېنټ پر ښي
فارمول کار کوي۔ په دې توګه د سېکوېنټونو کيڼ اړخ ته نښه لرونکي
فارمولونه يوازې سياق بللے شو۔
\olpseganchor{OLP-0669-B007}{end}

\olpseganchor{OLP-0669-B008}{start}
د فرضونو پر ځاے، د \Log{N2} ثبوتونو تر ټولو بره سېکوېنټونه
د دې بڼې بديهيات دي:
\[x:!A \Sequent !A.\]
قاعدې د \Log{N1} له قواعدو سره کره سمون لري او په
\olref{tab:N2} کښې ورکړل شوې دي۔
\olpseganchor{OLP-0669-B008}{end}

\olpseganchor{OLP-0669-B009}{start}

\olpseganchor{OLP-0669-B009}{end}

\olpseganchor{OLP-0669-B010}{start}
لکه په \Log{N1c} او \Log{N1i} کښې، بايد دا ومنو چې
\Intro{\lif}، \Intro{\lnot}، \Elim{\lor} او \Elim{\lexists}
قاعدې فرضونه ايستلے شي، خو اړې نۀ دي چې فرضونه وباسي۔ نو په
\Log{N2c} او~\Log{N2i} کښې د اړوندو قواعدو سمې کارونې هغه
وخت هم منو چې اړوند سياقي فارمولونه \LR{$!A$}، \LR{$!B$} او
\LR{$!A(a)$} د اړوندې مقدمې په سياق کښې نۀ وي۔
\olpseganchor{OLP-0669-B010}{end}

\olpseganchor{OLP-0669-B011}{start}
\begin{prop}
که \LR{$\delta$} په~\Log{N1c} يا \Log{N1i} کښې د~\LR{$!A$}
ثبوت وي، نو په~\Log{N2c} (\Log{N2i}) کښې د
\LR{$\Gamma \Sequent !A$} ثبوت~\LR{$\delta'$} شته، چې
\LR{$\Gamma$} د ټولو هغو~\LR{$x:!B$} سټ دے چې \LR{$!B^x$} ئې د~\LR{$\delta$}
پرانيستے فرض وي۔
\end{prop}
\olpseganchor{OLP-0669-B011}{end}

\olpseganchor{OLP-0669-B012}{start}
\begin{proof}
  په \LR{$n=\pheight{\delta}$} استقرا کوو۔ که \LR{$n=0$} وي، \LR{$\delta$}
  يوازې يو پرانيستے فرض~\LR{$!A^x$} دے۔ نو \LR{$\Gamma = \{x:!A\}$}
  او \LR{$\Gamma
  \Sequent !A$} د \Log{N2c} يا~\Log{N2i} يو بديهيت دے۔
\olpseganchor{OLP-0669-B012}{end}

\olpseganchor{OLP-0669-B013}{start}
  که \LR{$n>0$} وي، د~\LR{$\delta$} د وروستي استنتاج له مخې حالتونه
  بېلوو۔ يوازې يو حالت بيان کوو؛ پاتې حالتونه تمرينونه دي۔
\olpseganchor{OLP-0669-B013}{end}

\olpseganchor{OLP-0669-B014}{start}
  فرض کړئ وروستے استنتاج \LR{$\Intro{\lif}$} وي۔ نو \LR{$\delta$}
  په دې ختمېږي:
  \[
  \AxiomC{\LR{$\Discharge{!B}{x}$}}
  \RightLabel{\LR{$\delta_1$}}
  \DeduceC{\LR{$!C$}}
  \DischargeRule{\Intro{\lif}}{x}
  \UnaryInfC{\LR{$!B \lif !C$}}
  \DisplayProof
  \]
  د استقرا له فرضيې، په~\Log{N2c} (\Log{N2i}) کښې د
  \LR{$\Gamma' \Sequent !C$} ثبوت~\LR{$\delta_1'$} شته، چې
  \LR{$\Gamma'$} هغه نښه لرونکي فرضونه دي چې په~\LR{$\delta_1$}
  کښې پرانيستي وي۔ که \LR{$!B^x$} د \LR{$\delta_1$} پرانيستے فرض وي،
  نو د \Intro{\lif} په استنتاج ايستل کېږي او د~\LR{$\delta$}
  پرانيستي فرضونه \LR{$\Gamma =
  \Gamma' \setminus \{x:!B\}$} دي۔ يعنې \LR{$\delta_1'$} د
  \LR{$x:!B, \Gamma \Sequent !C$} ثبوت دے؛ نو \LR{$\delta'$}
  داسې اخستلے شو:
  \[
  \AxiomC{}
  \RightLabel{\LR{$\delta_1'$}}
  \Deduce$x:!B, \Gamma \fCenter !C$
  \RightLabel{\Intro{\lif}}
  \UnaryInf$\Gamma \fCenter !B \lif !C$
  \DisplayProof
  \]
  که \LR{$!B^x$} د~\LR{$\delta_1$} پرانيستے فرض نۀ وي، نو
  \Intro{\lif} هېڅ فرض نۀ باسي او د \LR{$\delta$} او \LR{$\delta_1$}
  پرانيستي فرضونه يو شان دي۔ په~\Log{N2c} او~\Log{N2i}
  کښې د \Intro{\lif} د مقدمې په سياق کښې د \LR{$x:!B$} شته‌والے
  اړين نۀ دے؛ نو \LR{$\delta'$} داسې اخستلے شو:
  \[
  \AxiomC{}
  \RightLabel{\LR{$\delta_1'$}}
  \Deduce$\Gamma \fCenter !C$
  \RightLabel{\Intro{\lif}}
  \UnaryInf$\Gamma \fCenter !B \lif !C$
  \DisplayProof
  \]
\olpseganchor{OLP-0669-B014}{end}

\olpseganchor{OLP-0669-B015}{start}
  هغه حالتونه چې \LR{$\delta$} په بله قاعده ختمېږي، همداسې دي۔
\end{proof}
\olpseganchor{OLP-0669-B015}{end}

\olpseganchor{OLP-0669-B016}{start}
\begin{prop}
که \LR{$\delta$} په \Log{N2c} يا \Log{N2i} کښې د \LR{$\Gamma
\Sequent !A$} ثبوت وي، نو په \Log{N1c} (\Log{N1i})
کښې ثبوت~\LR{$\delta'$} شته چې پاييز فارمول ئې~\LR{$!A$}
وي او د هر \LR{$x:!B \in \Gamma$} دپاره پرانيستي فرضونه ئې
\LR{$!B^x$} وي۔
\end{prop}
\olpseganchor{OLP-0669-B016}{end}

\olpseganchor{OLP-0669-B017}{start}
\begin{proof}
  بيا د~\LR{$\delta$} پر لوړوالے~\LR{$n$} استقرا کوو۔ که \LR{$n=0$}
  وي، \LR{$\delta$} يو بديهيت دے، \LR{$x:!A \Sequent !A$}۔
  ثبوت~\LR{$\delta'$} هماغه فرض~\LR{$!A^x$} دے۔
\olpseganchor{OLP-0669-B017}{end}

\olpseganchor{OLP-0669-B018}{start}
  که \LR{$n>0$} وي، د~\LR{$\delta$} د وروستي استنتاج له مخې حالتونه
  بېلوو۔ مثلاً، که دغه استنتاج~\LR{$\Intro{\lif}$} وي، نو
  \LR{$\delta$}~په دې ختمېږي:
  \[
  \bottomAlignProof
  \AxiomC{}
  \RightLabel{\LR{$\delta_1$}}
  \Deduce$x:!B, \Gamma \fCenter !C$
  \RightLabel{\Intro{\lif}}
  \UnaryInf$\Gamma \fCenter !B \lif !C$
  \DisplayProof
  \quad\text{\RL{ يا }}\quad
  \bottomAlignProof
  \AxiomC{}
  \RightLabel{\LR{$\delta_1$}}
  \Deduce$\Gamma \fCenter !C$
  \RightLabel{\Intro{\lif}}
  \UnaryInf$\Gamma \fCenter !B \lif !C$
  \DisplayProof
  \]
  د استقرا له فرضيې، په~\Log{N1c} (\Log{N1i}) کښې د~\LR{$!C$}
  ثبوت~\LR{$\delta_1'$} شته۔ په لومړي حالت کښې \LR{$!B^x$}
  د~\LR{$\delta_1'$} پرانيستے فرض دے؛ په دويم کښې نۀ دے۔
  ثبوت~\LR{$\delta'$} په ترتيب داسې اخستلے شو:
  \[
  \bottomAlignProof
  \AxiomC{\LR{$\Discharge{!B}{x}$}}
  \RightLabel{\LR{$\delta_1'$}}
  \DeduceC{\LR{$!C$}}
  \DischargeRule{\Intro{\lif}}{x}
  \UnaryInfC{\LR{$!B \lif !C$}}
  \DisplayProof
  \quad\text{\RL{ يا }}\quad
  \bottomAlignProof
  \AxiomC{}
  \RightLabel{\LR{$\delta_1'$}}
  \DeduceC{\LR{$!C$}}
  \RightLabel{\Intro{\lif}}
  \UnaryInfC{\LR{$!B \lif !C$}}
  \DisplayProof
  \]
  په لومړي رسم کښې وروستۍ قاعده نښه شوے فرض باسي؛ په دويم کښې
  پاييزه پايله له هماغو پخوانيو فرضونو څخه ثابته ده۔
\end{proof}
\olpseganchor{OLP-0669-B018}{end}

\olfileid{pt}{nat}{seq}
\olpseganchor{OLP-0667-B004}{start}
\begin{table}
\begin{tabular}{@{}c@{}c@{}}
\Axiom$x: !A, \Gamma \fCenter !B$
\RightLabel{\Intro{\lif}}
\UnaryInf$\Gamma \fCenter !A \lif !B $
\DisplayProof
&
\Axiom$\Gamma_1 \fCenter !A \lif !B$
\Axiom$\Gamma_2 \fCenter !A$
\RightLabel{\Elim{\lif}}
\BinaryInf$\Gamma_1, \Gamma_2 \fCenter !B$
\DisplayProof
\\[4ex]
\Axiom$\Gamma_1 \fCenter !A$
\Axiom$\Gamma_2 \fCenter !B$
\RightLabel{\Intro{\land}}
\BinaryInf$\Gamma_1, \Gamma_2 \fCenter !A \land !B $
\DisplayProof
&
\begin{tabular}[b]{@{}l@{}}
\Axiom$\Gamma \fCenter !A \land !B$
\RightLabel{\Elim{\land}[1]}
\UnaryInf$\Gamma \fCenter !A$
\DisplayProof
\\[3ex]
\Axiom$\Gamma \fCenter !A \land !B$
\RightLabel{\Elim{\land}[2]}
\UnaryInf$\Gamma \fCenter !B$
\DisplayProof\\[3ex]
\end{tabular}
\\[4ex]
\Axiom$\Gamma \fCenter !A$
\RightLabel{\Intro{\lor}[1]}
\UnaryInf$\Gamma \fCenter !A \lor !B$
\DisplayProof
&
\Axiom$ \Gamma \fCenter !B$
\RightLabel{\Intro{\lor}[2]}
\UnaryInf$\Gamma \fCenter !A \lor !B$
\DisplayProof\\[3ex]
\multicolumn{2}{c}{
\begin{tabular}[b]{@{}l@{}}
\Axiom$\Gamma_1 \fCenter !A \lor !B$
\Axiom$x:!A, \Gamma_2 \fCenter !C$
\Axiom$y:!B, \Gamma_3 \fCenter !C$
\RightLabel{\Elim{\lor}}
\TrinaryInf$\Gamma_1, \Gamma_2, \Gamma_3 \fCenter !C$
\DisplayProof
\end{tabular}
}
\\[4ex]
\Axiom$x:!A, \Gamma \fCenter \lfalse$
\RightLabel{\Intro{\lnot}}
\UnaryInf$ \Gamma \fCenter \lnot !A $
\DisplayProof
&
\Axiom$\Gamma_1 \fCenter \lnot !A $
\Axiom$\Gamma_2 \fCenter !A $
\RightLabel{\Elim{\lnot}}
\BinaryInf$\Gamma_1, \Gamma_2 \fCenter \lfalse$
\DisplayProof
\\[4ex]
\Axiom$\Gamma \fCenter !A(a) $
\RightLabel{\Intro{\lforall}}
\UnaryInf$\Gamma \fCenter \lforall[x][!A(x)]$
\DisplayProof
&
\Axiom$\Gamma \fCenter \lforall[x][!A(x)]$
\RightLabel{\Elim{\lforall}}
\UnaryInf$\Gamma \fCenter !A(t)$
\DisplayProof
\\[4ex]
\Axiom$\Gamma \fCenter !A(t) $
\RightLabel{\Intro{\lexists}}
\UnaryInf$\Gamma \fCenter \lexists[x][!A(x)]$
\DisplayProof
&
\Axiom$\Gamma_1 \fCenter  \lexists[x][!A(x)]$
\Axiom$x: !A(a), \Gamma_2 \fCenter !C$
\RightLabel{\Elim{\lexists}}
\BinaryInf$\Gamma_1, \Gamma_2 \fCenter !C$
\DisplayProof
\\[4ex]
\Axiom$\Gamma \fCenter \lfalse$
\RightLabel{\FalseInt}
\UnaryInf$\Gamma \fCenter !A$
\DisplayProof
&
\Axiom$x: \lnot !A, \Gamma \fCenter \lfalse$
\RightLabel{\FalseCl}
\UnaryInf$\Gamma \fCenter !A$
\DisplayProof
\end{tabular}
\caption{د \Log{N2c} او \Log{N2i} قاعدې۔ په \LR{$\Intro{\forall}$} او
\LR{$\Elim{\exists}$} کښې ثابت سمبول~\LR{$a$} بايد په پاييز سېکوېنټ کښې
نۀ راځي؛ د موجود کميت د ايستلو په لويه مقدمه کښې هم دغه ثابت
نۀ شي راتلے۔ \LR{$\Gamma_i$}~د نښه لرونکو فارمولونو~\LR{$x :
!A$} سټونه دي۔ د \Intro{\lforall} او \Elim{\lexists} قاعدې بايد
د ځانګړي متغير شرط پوره کړي: \LR{$a$} بايد په پاييز سېکوېنټ کښې
نۀ راځي او د موجود کميت د ايستلو په لويه مقدمه کښې هم نۀ راځي۔
\Log{N2i} د \FalseCl لرونکے نۀ دے۔}
\ollabel{tab:N2}
\end{table}
\olpseganchor{OLP-0667-B004}{end}

\olpseganchor{OLP-0665-B004}{start}
\olfileid{pt}{nat}{qua}
\olpseganchor{OLP-0665-B004}{end}

\olpseganchor{OLP-0665-B005}{start}
\olsection{منظم ثبوتونه او تعويض}
\olpseganchor{OLP-0665-B005}{end}

\olpseganchor{OLP-0665-B006}{start}
د کميت ټاکونکو قاعدې ځينې ستونزې راولي، ځکه د «ځانګړي متغير
شرطونه» د \Elim{\lexists} او~\Intro{\lforall} پر قاعدو لګېږي۔
د دغو ډېرو ستونزو حل دا دے چې يقيني کړو په دې استنتاجونو کښې
ثابت سمبول~\LR{$c$} بيا نۀ کارول کېږي۔ لاندې تعريف وړاندې کوو۔
\olpseganchor{OLP-0665-B006}{end}

\olpseganchor{OLP-0665-B007}{start}
\begin{defn}
د فطري استنتاج ثبوت~\LR{$\delta$} هله \emph{منظم} دے چې
\begin{enumerate}
  \item هېڅ ځانګړے متغير~\LR{$a$} د \Elim{\lexists} يا
  \Intro{\lforall} له يوه څخه د ډېرو استنتاجونو ځانګړے متغير نۀ وي؛ او
  \item که \LR{$a$} د \Elim{\lexists} يا~\Intro{\lforall} د کوم
  استنتاج ځانګړے متغير وي، نو يوازې په هغه نېغ فرعي ثبوت کښې
  راځي چې په ترتيب ئې کوچنۍ مقدمه يا يوازېنۍ مقدمه ورکوي۔
\end{enumerate}
\end{defn}
\olpseganchor{OLP-0665-B007}{end}

\olpseganchor{OLP-0665-B008}{start}
\begin{lem}\ollabel{lem:subst-regular}
که \LR{$\delta$} منظم وي او په~\LR{$!C$} ختمېږي، \LR{$c$} داسې ثابت سمبول
وي چې په~\LR{$\delta$} کښې د ځانګړي متغير په توګه نۀ کارېږي، او
\LR{$t$} داسې ترم وي چې د~\LR{$\delta$} هېڅ ځانګړے متغير پکښې نۀ وي،
نو \LR{$\Subst{\delta}{t}{c}$} چې له~\LR{$\delta$} څخه د~\LR{$c$} د هر راتګ په~\LR{$t$}
بدلولو جوړېږي، يو منظم ثبوت دے۔ د \LR{$\Subst{\delta}{t}{c}$} پاييز
فارمول د \LR{$\Subst{!C}{t}{c}$} سره برابر دے۔ که \LR{$!A^x$} د~\LR{$\delta$}
يو پرانيستے فرض وي، نو \LR{$\Subst{!A}{t}{c}^x$} د~\LR{$\Subst{\delta}{t}{c}$} يو پرانيستے فرض دے۔
\end{lem}
\olpseganchor{OLP-0665-B008}{end}

\olpseganchor{OLP-0665-B009}{start}
\begin{proof}
د \LR{$\pheight{\delta}$} پر لوړوالي په استقرا۔ که \LR{$\pheight{\delta} = 0$} وي، نو ثبوت يوازې
له فرض~\LR{$!A^x$} څخه جوړ دے۔ بيا \LR{$\Subst{\delta}{t}{c}$} د~\LR{$\Subst{!A}{t}{c}^x$} سره برابر دے۔
\olpseganchor{OLP-0665-B009}{end}

\olpseganchor{OLP-0665-B010}{start}
که \LR{$\pheight{\delta} > 0$} وي، د~\LR{$\delta$} د وروستي استنتاج له مخې حالتونه بېلوو۔
د قضييزو تړونکو د قاعدو حالتونه اسانه دي او د تمرينونو په توګه
پرېښودل شوي دي۔ هغه حالتونه څېړو چې \LR{$\delta$} په يوه کميتي استنتاج ختمېږي۔
\begin{enumerate}
  \item که وروستے استنتاج \LR{$\Intro{\lforall}$} وي، نو \LR{$\delta$} دا بڼه لري:
  \[
    \AxiomC{}
    \RightLabel{\LR{$\delta'$}}
    \DeduceC{\LR{$!A(a,c)$}}
    \RightLabel{\Intro{\lforall}}
    \UnaryInfC{\LR{$\lforall[x][!A(x,c)]$}}
    \DisplayProof
    \]
  د استقرا د فرض له مخې، \LR{$\Subst{\delta'}{t}{c}$} د \LR{$!A(a,t)$} يو منظم ثبوت
  دے۔ دا چې \LR{$a$} د~\LR{$\delta$} ځانګړے متغير دے، نو \LR{$a$}
  په~\LR{$t$} کښې نۀ راځي۔ ځکه په \LR{$\Subst{\delta}{t}{c}$} کښې وروستے استنتاج
  \[
    \AxiomC{}
    \RightLabel{\LR{$\Subst{\delta'}{t}{c}$}}
    \DeduceC{\LR{$!A(a,t)$}}
    \RightLabel{\Intro{\lforall}}
    \UnaryInfC{\LR{$\lforall[x][!A(x,t)]$}}
    \DisplayProof
    \]
  سم دے: \LR{$t$} د~\LR{$a$} لرونکے نۀ دے، نو نۀ په پايله
  کښې~\LR{$a$} راځي او نۀ په کوم پرانيستي فرض کښې۔
  \item که وروستے استنتاج \LR{$\Elim{\lforall}$} وي، نو \LR{$\delta$} دا بڼه لري:
  \[
    \AxiomC{}
    \RightLabel{\LR{$\delta'$}}
    \DeduceC{\LR{$\lforall[x][!A(x,c)]$}}
    \RightLabel{\Elim{\lforall}}
    \UnaryInfC{\LR{$!A(s(c),c)$}}
    \DisplayProof
    \]
  د استقرا د فرض له مخې، \LR{$\Subst{\delta'}{t}{c}$} د \LR{$\lforall[x][!A(x,t)]$} يو منظم ثبوت
  دے۔ (په پايله کښې ترم~\LR{$s(c)$} کېدے شي د~\LR{$c$} لرونکے وي،
  خو اړين نۀ دے۔) په \LR{$\Subst{\delta}{t}{c}$} کښې وروستے استنتاج
  \[
    \AxiomC{}
    \RightLabel{\LR{$\Subst{\delta'}{t}{c}$}}
    \DeduceC{\LR{$\lforall[x][!A(x,t)]$}}
    \RightLabel{\Elim{\lforall}}
    \UnaryInfC{\LR{$!A(s(t),t)$}}
    \DisplayProof
    \]
  سم دے، او \LR{$!A(s(t),t) \ident \Subst{!A(s(c),c)}{t}{c}$} دے۔ د \Intro{\lexists} حالت هم ورته دے۔
  \item که وروستے استنتاج \LR{$\Elim{\lexists}$} وي، نو \LR{$\delta$} دا بڼه لري:
  \[
    \AxiomC{}
    \RightLabel{\LR{$\delta_1'$}}
    \DeduceC{\LR{$\lexists[x][!A(x,c)]$}}
    \AxiomC{\LR{$\Discharge{!A(a,c)}{x}$}}
    \RightLabel{\LR{$\delta_2'$}}
    \DeduceC{\LR{$!C$}}
    \DischargeRule{\Elim{\lexists}}{x}
    \BinaryInfC{\LR{$!C$}}
    \DisplayProof
    \]
  د استقرا د فرض له مخې، \LR{$\Subst{\delta_1'}{t}{c}$} او \LR{$\Subst{\delta_2'}{t}{c}$} د \LR{$\lexists[x][!A(x,t)]$} او
  د \LR{$\Subst{!C}{t}{c}$} منظم ثبوتونه دي؛ دويم له پرانيستي
  فرض~\LR{$!A(a,t)$} څخه راځي۔ دا چې \LR{$a$} د~\LR{$\delta$} ځانګړے متغير
  دے، نو \LR{$a$} په~\LR{$t$} کښې نۀ راځي۔ ځکه په \LR{$\Subst{\delta}{t}{c}$} کښې
  وروستے استنتاج
  \[
    \AxiomC{}
    \RightLabel{\LR{$\Subst{\delta_1'}{t}{c}$}}
    \DeduceC{\LR{$\lexists[x][!A(x,t)]$}}
    \AxiomC{\LR{$\Discharge{!A(a,t)}{x}$}}
    \RightLabel{\LR{$\Subst{\delta_2'}{t}{c}$}}
    \DeduceC{\LR{$\Subst{!C}{t}{c}$}}
    \DischargeRule{\Elim{\lexists}}{x}
    \BinaryInfC{\LR{$\Subst{!C}{t}{c}$}}
    \DisplayProof
    \]
  سم دے: \LR{$\Subst{!A(x,t)}{a}{x} \ident \Subst{!A(a,c)}{t}{c}$} دے، او نۀ په پايله~\LR{$\Subst{!C}{t}{c}$} کښې~\LR{$a$}
  راځي او نۀ په کوم پاتې پرانيستي فرض کښې۔
  همدا د موجود کميت د ايستلو دوه‌مقدمي استنتاج دے۔
  تعويض د دواړو مقدمو او پايلې په ټولو اړوندو ځايونو کښې کېږي۔
\end{enumerate}
\end{proof}
\olpseganchor{OLP-0665-B010}{end}

\olpseganchor{OLP-0665-B011}{start}
\begin{prop}\ollabel{prop:regular}
که \LR{$\delta$} په \Log{N1c} يا \Log{N1i} کښې ثبوت وي، نو
يو منظم ثبوت~\LR{$\delta'$} شته چې جوړښت (په ځانګړي ډول هماغه
لوړوالے) ئې عين دے او له \LR{$\delta$} څخه يوازې د ځانګړو متغيرونو
په بدلولو توپير لري۔
\end{prop}
\olpseganchor{OLP-0665-B011}{end}

\olpseganchor{OLP-0665-B012}{start}
\begin{proof}
يوازې د دې ثبوت دپاره د~\LR{$\delta$} د ځانګړي متغير يو استنتاج هله
\emph{پاک} ګڼو چې ځانګړے متغير ئې د کوم بل استنتاج ځانګړے
متغير نۀ وي، او په~\LR{$\delta$} کښې د همدغه استنتاج له نېغ فرعي
ثبوت څخه پرته بل چېرته نۀ راځي؛ که داسې نۀ وي، نو
\emph{ناپاک} ئې ګڼو۔ \LR{$n$} دې په~\LR{$\delta$} کښې د ناپاکو ځانګړي
متغيرونو د استنتاجونو شمېر وي۔ د~\LR{$n$} پر شمېر په قوي استقرا
ښيو چې \LR{$\delta$} په اړين منظم ثبوت~\LR{$\delta'$} بدلېدے شي۔ که
\LR{$n=0$} وي، په~\LR{$\delta$} کښې د ځانګړي متغير ټول استنتاجونه پاک
دي؛ نو \LR{$\delta$} منظم دے او د ثابتولو څه پاتې نۀ دي۔ که داسې
نۀ وي، تر ټولو بره يو ناپاک د ځانګړي متغير استنتاج غوره کوو؛
مثلاً \Intro{\lforall}۔ د~\LR{$\delta$} هغه فرعي ثبوت~\LR{$\delta_1$}
چې په دې استنتاج ختمېږي، دا بڼه لري:
\[
    \AxiomC{}
    \RightLabel{\LR{$\delta_2$}}
    \DeduceC{\LR{$!A(c)$}}
    \RightLabel{\Intro{\lforall}}
    \UnaryInfC{\LR{$\lforall[x][!A(x)]$}}
    \DisplayProof
    \]
دا چې \LR{$c$} ځانګړے متغير دے، نو په~\LR{$\lforall[x][!A(x)]$} يا په کوم
پرانيستي فرض کښې نۀ راځي۔ ثبوت \LR{$\delta_2$} منظم دے، ځکه د
ځانګړي متغير ټول استنتاجونه ئې پاک دي؛ د غوره شوي استنتاج
ځانګړے متغير په دې فرعي ثبوت کښې د کوم استنتاج ځانګړے
متغير نۀ دے۔ \LR{$a$} دې داسې ثابت سمبول وي چې په~\LR{$\delta$}
کښې نۀ راځي۔ د \olref{lem:subst-regular} له مخې، \LR{$\Subst{\delta_2}{a}{c}$}
د \LR{$!A(a)$} يو منظم ثبوت دے۔ د ځانګړي متغير د شرط له مخې،
د~\LR{$\delta_2$} هېڅ پرانيستے فرض د~\LR{$c$} لرونکے نۀ دے؛ نو
د~\LR{$\Subst{\delta_2}{a}{c}$} هېڅ پرانيستے فرض د~\LR{$a$} لرونکے نۀ دے۔ ځکه
\LR{$\Intro{\lforall}$} لګولے شو۔ که په~\LR{$\delta$} کښې \LR{$\delta_1$} په ثبوت \LR{$\delta_1'$}
بدل کړو،
\[
    \AxiomC{}
    \RightLabel{\LR{$\Subst{\delta_2}{a}{c}$}}
    \DeduceC{\LR{$!A(a)$}}
    \RightLabel{\Intro{\lforall}}
    \UnaryInfC{\LR{$\lforall[x][!A(x)]$}}
    \DisplayProof
    \]
نو د ځانګړي متغير يو ناپاک استنتاج مو په پاک استنتاج بدل کړے
دے؛ يوازې دا توپير راغلے دے چې ځانګړے متغير \LR{$c$} په~\LR{$a$}
بدل شوے دے۔ په راوتلي ثبوت کښې له~\LR{$n$} څخه کم د ناپاکو
ځانګړي متغيرونو استنتاجونه دي؛ نو د استقرا له فرض څخه پايله راځي۔
د زاړه ځانګړي متغير د ځينو موردونو په لرې کېدو بل
ناپاک استنتاج هم پاکېدے شي۔ د استقرا لپاره د شمېر کمېدل بس دي۔
\end{proof}
\olpseganchor{OLP-0665-B012}{end}

\olpseganchor{OLP-0665-B013}{start}
\begin{prob}
د \olref{prop:regular} په ثبوت کښې هغه حالت بيان کړئ چې
تر ټولو بره د ځانګړي متغير استنتاج~\Elim{\lexists} وي۔
\end{prob}
\olpseganchor{OLP-0665-B013}{end}

\olpseganchor{OLP-0665-B014}{start}
ثابته شوې قضيه به په څرګند ډول د دليل په توګه نۀ کاروو،
خو فرض به کړو چې ټول څېړل کېدونکي ثبوتونه منظم دي؛
که منظم نۀ وي، ځانګړي متغيرونه ئې تل بدلولے شو چې منظم شي۔
\olpseganchor{OLP-0665-B014}{end}

\olpseganchor{OLP-0665-B015}{start}
\Olref{lem:subst-regular} او \olref{prop:regular} د
\Log{N2c} او~\Log{N2i} دپاره هم کره دي۔ ثبوتونه ئې
د تمرينونو په توګه پرېږدو۔
\olpseganchor{OLP-0665-B015}{end}

\olpseganchor{OLP-0662-B004}{start}
\olfileid{pt}{nat}{gra}
\olpseganchor{OLP-0662-B004}{end}

\olpseganchor{OLP-0662-B005}{start}
\olsection{ثبوتونه پېوندول}
\olpseganchor{OLP-0662-B005}{end}

\olpseganchor{OLP-0662-B006}{start}
په \Log{N1c} او~\Log{N1i} کښې ثبوتونه له فرضونو~\LR{$!B$}
څخه د فارمولونه~\LR{$!A$} ثبوتونه دي۔ فرض کړئ چې پخپله \LR{$!B$}
يو ثبوت لري۔ د~\LR{$!B$} ثبوت د \LR{$!A$} له هغه ثبوت
سره څنګه يوځاے کړو چې له~\LR{$!B$} څخه دے، څو د~\LR{$!A$} داسې
يو ثبوت تر لاسه کړو چې فرض~\LR{$!B$} نۀ کاروي؟
\olpseganchor{OLP-0662-B006}{end}

\olpseganchor{OLP-0662-B007}{start}
يوه طريقه دا ده چې د~\Intro{\lif} په کارولو د~\LR{$!A$} له
ثبوت څخه فرض~\LR{$!B$} وايستل شي۔ بيا د~\LR{$!B \lif !A$} نوے
ثبوت د~\LR{$!B$} له ثبوت سره يوځاے کولے شو چې دا تر
لاسه کړو:
\[
\AxiomC{\LR{$\Discharge{!B}{x}$}}
\DeduceC{\LR{$!A$}}
\DischargeRule{\Intro{\lif}}{x}
\UnaryInfC{\LR{$!B \lif !A$}}
\AxiomC{}
\DeduceC{\LR{$!B$}}
\RightLabel{\Elim{\lif}}
\BinaryInfC{\LR{$!A$}}
\DisplayProof
\]\LR{$\lif$} داخل کړو چې سم
لاسي ئې بيا وايستل شي؟ د \LR{$!B \lif !A$} له لارې دا ډول
«تاو راتاو» بايد پرېښودلے شي۔ (دا چې تل پرېښودلے شي، په
حقيقت کښې د عادي کولو د قضيې مضمون دے۔)
\olpseganchor{OLP-0662-B007}{end}

\olpseganchor{OLP-0662-B008}{start}
په \Log{N1c} او~\Log{N1i} کښې په دې حالت کښې دا په اسانۍ
پرېښودلے شو: فرضونه~\LR{$!B^x$} يوازې د~\LR{$!B$} په ټول ثبوت
بدلولے شو۔ د نورو د فرضونو پر ځاے د ثبوتونه دا «پېوندول»
ډېر ګټور دي، نو د تعريف په توګه ئې ثبتوو۔
\olpseganchor{OLP-0662-B008}{end}

\olpseganchor{OLP-0662-B009}{start}
\begin{defn}
که \LR{$\delta_1$} د پرانيستي فرض~\LR{$!B^x$} څخه د~\LR{$!A$} يو
\Log{N1c}-ثبوت (يا \Log{N1i}-ثبوت) وي، او \LR{$\delta_2$}
د~\LR{$!B$} يو \Log{N1c}-ثبوت (\Log{N1i}-ثبوت) وي، نو
\LR{$\Subst{\delta_1}{\delta_2}{!B^x}$} په~\LR{$\delta_1$} کښې د~\LR{$!B^x$}
د هر نۀ ايستل شوي راتګ په~\LR{$\delta_2$} د بدلولو پايله ده۔
\end{defn}
\olpseganchor{OLP-0662-B009}{end}

\olpseganchor{OLP-0662-B010}{start}
په دې شرط چې \LR{$\delta_1$} او \LR{$\delta_2$} په عين نښه د بېلابېلو
فرض فارمولونه نۀ لري او د~\LR{$\delta_2$} هېڅ پرانيستے فرض
د~\LR{$\delta_1$} ځانګړے متغير نۀ لري، نو
\LR{$\Subst{\delta_1}{\delta_2}{!B^x}$} يو کره ثبوت دے۔ هر
پرانيستے فرض ئې يا د \LR{$\delta_1$} يو نۀ بدل شوے پرانيستے فرض
دے، يا د~\LR{$\delta_2$} داسې پرانيستے فرض دے چې له پېوندولو وروسته
نۀ دے ايستل شوے۔ بدل شوے فرض له خپل پخواني ځاے څخه لرې کېږي،
او د ورګډ شوي ثبوت نښه لرونکے فرض د نوي ځاے په ورپسې قاعده
کښې ايستل کېدے شي۔ د ثبوتونه د پېوندولو په وخت دا شرطونه
عموماً پوره وي۔ لومړنے شرط هله پوره وي چې \LR{$\delta_1$} او
\LR{$\delta_2$} د يوه لوے ثبوت فرعي ثبوتونه وي۔ که دويم
شرط پوره نۀ وي، د~\LR{$\delta_1$} ځانګړي متغيرونه بيا نومولے شو
چې پوره شي۔
\olpseganchor{OLP-0662-B010}{end}

\olpseganchor{OLP-0662-B011}{start}
ورته تعريف په \Log{N2c} او~\Log{N2i} هم لګېږي۔
\olpseganchor{OLP-0662-B011}{end}

\olpseganchor{OLP-0662-B012}{start}
\begin{defn}
که \LR{$\delta_1$} د~\LR{$x:!B, \Gamma_1 \Sequent !A$} يو
\Log{N2c}-ثبوت (يا \Log{N2i}-ثبوت) وي، او \LR{$\delta_2$}
د~\LR{$\Gamma_2 \Sequent !B$} يو \Log{N2c}-ثبوت
(\Log{N2i}-ثبوت) وي، نو \LR{$\Subst{\delta_1}{\delta_2}{x:!B}$}
په~\LR{$\delta_1$} کښې د هرې بديهيې \LR{$x:!B \Sequent !B$} په~\LR{$\delta_2$}
د بدلولو او د داسې بديهياتو د لاندې هر سېکوېنټ سياق ته د
\LR{$\Gamma_2$} د زياتولو پايله ده۔
\end{defn}
\olpseganchor{OLP-0662-B012}{end}

\olpseganchor{OLP-0662-B013}{start}
\begin{prop}
که \LR{$\delta_1 \Proves x:!B, \Gamma_1 \Sequent
!A$} او \LR{$\delta_2 \Proves \Gamma_2 \Sequent !B$} وي، نو
\LR{$\Subst{\delta_1}{\delta_2}{x:!B} \Proves \Gamma_1, \Gamma_2 \Sequent
!A$}، په دې شرط چې د~\LR{$\delta_1$} هېڅ ځانګړے متغير په~\LR{$\Gamma_2$}
کښې رانۀ شي۔
\end{prop}
\olpseganchor{OLP-0662-B013}{end}

\olpseganchor{OLP-0662-B014}{start}
\begin{proof}
د \LR{$\pheight{\delta_1}$} پر لوړوالي په استقرا۔
\end{proof}
\olpseganchor{OLP-0662-B014}{end}

\olpseganchor{OLP-0670-B004}{start}
\olfileid{pt}{nat}{tgi}
\olpseganchor{OLP-0670-B004}{end}

\olpseganchor{OLP-0670-B005}{start}
\olsection{له \texorpdfstring{\Log{G2i}}{G2i} څخه \texorpdfstring{\Log{N2i}}{N2i} ته ژباړه}
\olpseganchor{OLP-0670-B005}{end}

\olpseganchor{OLP-0670-B006}{start}
په \Log{G2i} کښې د سېکوېنټونو کيڼ سياقونه د فارمولونه
د تکرار شمېر لرونکي سټونه~\LR{$\Gamma$} دي، خو په \Log{N2i} کښې
د نښه لرونکو فارمولونو سټونه~\LR{$\Gamma'$} دي، چې هره نښه ئې
يوازې يو ځل راځي۔ وايو چې د نښه لرونکو فارمولونه سټ~\LR{$\Gamma'$}
د تکرار شمېر لرونکي سټ~\LR{$\Gamma$} \emph{سره سمون لري}، که د هر
فارمول~\LR{$!A$} دپاره په~\LR{$\Gamma'$} کښې د \LR{$x : !A$} په بڼه
د غړي شمېر په~\LR{$\Gamma$} کښې د~\LR{$!A$} د تکرار له شمېر
څخه کم يا ورسره برابر وي، يعنې د~\LR{$!A$} د نقلونو له شمېر څخه۔
\olpseganchor{OLP-0670-B006}{end}

\olpseganchor{OLP-0670-B007}{start}
\begin{prop}\ollabel{prop:G2i-to-N2i}
که \LR{$\Log{G2i} + \Cut \Proves \Gamma \Sequent \Delta$} وي، نو
\LR{$\Log{N2i} \Proves \Gamma' \Sequent !C$}، چې \LR{$!C \ident \lfalse$}
که \LR{$\Delta$} تش وي، او که تش نۀ وي \LR{$\Delta = \{!C\}$} وي، او
\LR{$\Gamma'$} له~\LR{$\Gamma$} سره سمون لري۔
\end{prop}
\olpseganchor{OLP-0670-B007}{end}

\olpseganchor{OLP-0670-B008}{start}
\begin{proof}
  ښيو چې که \LR{$\pi$} په \Log{G2i} کښې، د پرېکون قاعدې سره، د
  \LR{$\Gamma \Sequent \Delta$} ثبوت وي، نو له~\LR{$\Gamma$} سره
  سمون لرونکے يو~\LR{$\Gamma'$} او په~\Log{N2i} کښې د
  \LR{$\Gamma' \Sequent !C$} ثبوت~\LR{$\delta$} شته۔ دا په
  \LR{$n = \pheight{\pi}$} استقرا ثابتوو۔
\olpseganchor{OLP-0670-B008}{end}

\olpseganchor{OLP-0670-B009}{start}
  که \LR{$n = 0$} وي، \LR{$\pi$} يو بديهيت دے۔ که د
  \LR{$\lfalse \Sequent \quad$} بديهيت وي، نو \LR{$\delta$} د
  \LR{$x:\lfalse \Sequent \lfalse$} بديهيت دے، يعنې
  \LR{$\Gamma' = \{x:\lfalse\}$} او \LR{$!C \ident \lfalse$}۔ که \LR{$\pi$}
  د \LR{$!A \Sequent !A$} په بڼه يو بديهيت وي، نو \LR{$\delta$} د
  \LR{$x: !A \Sequent !A$} بديهيت دے۔
\olpseganchor{OLP-0670-B009}{end}

\olpseganchor{OLP-0670-B010}{start}
  که \LR{$n>0$} وي، \LR{$\pi$} په يوه قاعده~\LR{$R$} ختمېږي۔ د استقرا
  فرضيه په~\Log{N2i} کښې د دې قاعدې له مقدمو سره سمون لرونکو
  سېکوېنټونو ثبوتونه راکوي۔ فرضولے شو چې په دغو ثبوتونو کښې
  د فرض ايستلو ټولې نښې بېلې دي، او که يوه نښه \LR{$x$} کله ايستل
  کېږي، نو يوازې د هغې ايستونکي استنتاج پورته راځي؛ که اړتيا
  وي، د نښو نومونه بدلوو۔ بيا ښيو چې دغه ثبوتونه څنګه يو
  ځاے کېدے شي چې د يوه مناسب \LR{$\Gamma' \Sequent !C$} ثبوتونه
  جوړ کړي۔ د~\LR{$R$} له مخې حالتونه بېلوو۔
\olpseganchor{OLP-0670-B010}{end}

\olpseganchor{OLP-0670-B011}{start}
  \begin{enumerate}
    \item که \LR{$R$} د \RightR{\Weakening} قاعده وي، \LR{$\pi$} داسې ختمېږي:
    \[
    \AxiomC{}
    \RightLabel{\LR{$\pi_1$}}
    \Deduce$\Gamma \fCenter$
    \RightLabel{\RightR{\Weakening}}
    \UnaryInf$\Gamma \fCenter !A$
    \DisplayProof
    \]
    د استقرا فرضيه پر \LR{$\pi_1$} د تطبيق له لارې د
    \LR{$\Gamma' \Sequent \lfalse$} ثبوت راکوي۔ د \LR{$\FalseInt$}
    په کارولو د \LR{$\Gamma' \Sequent !A$} ثبوت ترې جوړوو۔
    \item که \LR{$R$} د \LeftR{\Weakening} قاعده وي، \LR{$\pi$} داسې ختمېږي:
    \[
    \AxiomC{}
    \RightLabel{\LR{$\pi_1$}}
    \Deduce$\Gamma \fCenter \Delta$
    \RightLabel{\LeftR{\Weakening}}
    \UnaryInf$!A, \Gamma \fCenter \Delta$
    \DisplayProof
    \]
    د استقرا فرضيه پر \LR{$\pi_1$} د تطبيق له لارې د
    \LR{$\Gamma' \Sequent !C$} ثبوت راکوي، چې د \LR{$\delta$} په
    توګه ئې اخلو۔ که \LR{$\Gamma'$} له \LR{$\Gamma$} سره سمون لري، نو
    له \LR{$!A, \Gamma$} سره هم سمون لري، ځکه په \LR{$\Gamma'$} کښې د
    \LR{$x : !A$} شمېر په \LR{$\Gamma$} کښې د~\LR{$!A$} د نقلونو له شمېر
    څخه زيات نۀ دے، نو په \LR{$\Gamma \cup \{!A\}$} کښې د \LR{$!A$}
    له نقلونو هم زيات نۀ دے۔ دلته د تکرار شمېر لرونکو سياقونو
    يو ځاے کول د دواړو د نقلونو شمېر سره جمع کوي۔
    \item که \LR{$R$} د \LeftR{\Contraction} قاعده وي، \LR{$\pi$} داسې ختمېږي:
    \[
    \AxiomC{}
    \RightLabel{\LR{$\pi_1$}}
    \Deduce$!A, !A, \Gamma \fCenter \Delta$
    \RightLabel{\LeftR{\Contraction}}
    \UnaryInf$!A, \Gamma \fCenter \Delta$
    \DisplayProof
    \]
    د استقرا فرضيه پر \LR{$\pi_1$} د تطبيق له لارې د
    \LR{$\Pi, \Gamma' \Sequent !C$} ثبوت~\LR{$\delta_1$} راکوي،
    چې \LR{$\Pi \subseteq \{x : !A, y:!A\}$} دے۔ که
    \LR{$\Pi = \{x : !A, y:!A\}$} وي، په~\LR{$\delta_1$} کښې ټول
    نښه لرونکي فارمولونه \LR{$y:!A$} په \LR{$x:!A$} بدلوو۔ ځکه
    \LR{$x$}~د~\LR{$\delta_1$} د پاييز سېکوېنټ په سياق کښې راځي،
    فرضولے شو چې په~\LR{$\delta_1$} کښې هېڅ استنتاج د~\LR{$x:!B$}
    په بڼه فارمول نۀ باسي، يعنې په~\LR{$\delta_1$} کښې د
    سېکوېنټ پر کيڼ اړخ هېڅ \LR{$x:!B$} د ايستلو له قاعدې سره تړلے نۀ دے۔ نو پايله لا
    په~\Log{N2i} کښې ثبوت ده۔
    د يوځاے کولو پر وخت د پخوانيو بېلو فرضونو معنا او د ايستلو ځايونه ساتو۔
    \item که \LR{$R$} د \LR{$\RightR{\lif}$} قاعده وي، \LR{$\pi$} داسې ختمېږي:
    \[
    \AxiomC{}
    \RightLabel{\LR{$\pi_1$}}
    \Deduce$!A, \Gamma \fCenter !B$
    \RightLabel{\RightR{\lif}}
    \UnaryInf$\Gamma \fCenter !A \lif !B$
    \DisplayProof
    \]
    د استقرا له فرضيې د \LR{$x: !A, \Gamma' \Sequent !B$} يا
    \LR{$\Gamma' \Sequent !B$} ثبوت~\LR{$\delta_1$} شته، چې
    \LR{$\Gamma'$} له~\LR{$\Gamma$} سره سمون لري۔ \LR{$\delta$} د
    \LR{$\delta_1$} په پاے کښې د \Intro{\lif} د يوې کارونې په
    زياتولو اخستلے شو۔
    \item که \LR{$R$} د \LeftR{\lif} قاعده وي، \LR{$\pi$} داسې ختمېږي:
    \[
    \AxiomC{}
    \RightLabel{\LR{$\pi_1$}}
    \Deduce$\Gamma_1 \fCenter !A$
    \AxiomC{}
    \RightLabel{\LR{$\pi_2$}}
    \Deduce$!B, \Gamma_2 \fCenter \Delta$
    \RightLabel{\LeftR{\lif}}
    \BinaryInf$!A \lif !B, \Gamma_1, \Gamma_2 \fCenter \Delta$
    \DisplayProof
    \]
    د استقرا فرضيه د \LR{$\Gamma_1' \Sequent !A$} ثبوتونه
    \LR{$\delta_1$} او د \LR{$x: !B, \Gamma_2' \Sequent !C$} يا د
    \LR{$\Gamma_2' \Sequent !C$} ثبوت \LR{$\delta_2$} راکوي۔ که
    \LR{$\delta_2$} د وروستۍ بڼې وي، \LR{$\delta_2$} د~\LR{$\delta$} کار
    کولے شي۔ که داسې نۀ وي، په~\LR{$\delta_1$} کښې د ټولو
    فارمولونو او د فرض ايستلو نښو نومونه بدلولے شو، څو هېڅ
    نښه په دواړو \LR{$\delta_1$} او~\LR{$\delta_2$} کښې رانۀ شي۔ نو
    \LR{$\Gamma_1' \cup \Gamma_2'$} له \LR{$\Gamma_1 \cup \Gamma_2$}
    سره سمون لري۔ ځکه \LR{$x:!B$} د~\LR{$\delta_2$} په پاييز سېکوېنټ
    کښې راځي، \LR{$x:!B$} په~\LR{$\delta_2$} کښې د~\LR{$x$} د فرض ايستلو
    نښې لرونکي هېڅ استنتاج کښې فعال نۀ شي کېدے۔ په~\LR{$\delta_2$}
    کښې هر بديهيت \LR{$x:!B \Sequent !B$} په دې ثبوت~\LR{$\delta_3$}
    بدلوو:
    \[
    \Axiom$x: !A \lif !B \fCenter !A \lif !B$
    \AxiomC{}
    \RightLabel{\LR{$\delta_1$}}
    \Deduce$\Gamma_1' \fCenter !A$
    \RightLabel{\Elim{\lif}}
    \BinaryInf$x: !A \lif !B, \Gamma_1' \fCenter !B$
    \DisplayProof
    \]
    او په~\LR{$\delta_2$} کښې د \LR{$x:!B$} هر مورد په \LR{$x:!A \lif !B$}
    بدلوو۔ د ګرافټ کولو په ترڅ کښې د نوي فرعي ثبوت پاتې
    پرانيستي فرضونه هم پر اړوندو پورته سياقونو ورزياتېږي۔
    پايله \LR{$\Subst{\delta_2}{\delta_3}{x:!B}$} د
    \LR{$x:!A \lif !B, \Gamma_1', \Gamma_2' \Sequent !C$} ثبوت ده۔
    په هغه څانګه کښې چې د نوي فرعي ثبوت ځاے نيسي، د هغه پاتې سياق هم حسابېږي۔
    \item که \LR{$R$} د \RightR{\land} قاعده وي، \LR{$\pi$} داسې ختمېږي:
    \[
    \AxiomC{}
    \RightLabel{\LR{$\pi_1$}}
    \Deduce$\Gamma_1 \fCenter !A$
    \AxiomC{}
    \RightLabel{\LR{$\pi_2$}}
    \Deduce$\Gamma_2 \fCenter !B$
    \RightLabel{\RightR{\land}}
    \BinaryInf$\Gamma_1, \Gamma_2 \fCenter !A \land !B$
    \DisplayProof
    \]
    د استقرا فرضيه د \LR{$\Gamma_1' \Sequent !A$} ثبوت \LR{$\delta_1$}
    او د \LR{$\Gamma_2' \Sequent !B$} ثبوت \LR{$\delta_2$} راکوي۔
    په~\LR{$\delta_2$} کښې د ټولو فارمولونو او د فرض ايستلو نښو
    نومونه په بدلولو دا باوري کولے شو چې هېڅ نښه په دواړو
    \LR{$\delta_1$} او \LR{$\delta_2$} کښې رانۀ شي۔ نو
    \LR{$\Gamma_1' \cup \Gamma_2'$} له \LR{$\Gamma_1 \cup \Gamma_2$}
    سره سمون لري۔ د \LR{$\delta_1$} او~\LR{$\delta_2$} پر پاييزو
    سېکوېنټونو د \LR{$\Intro{\land}$} په کارولو
    ثبوت~\LR{$\delta$} تر لاسه کوو۔
    \item که \LR{$R$} د \LR{$\LeftR{\land}$} قاعده وي، \LR{$\pi$} مثلاً داسې ختمېږي:
    \[
    \AxiomC{}
    \RightLabel{\LR{$\pi_1$}}
    \Deduce$!A, \Gamma \fCenter \Delta$
    \RightLabel{\LeftR{\land}}
    \UnaryInf$!A \land !B, \Gamma \fCenter \Delta$
    \DisplayProof
    \]
    د استقرا له فرضيې د \LR{$x: !A, \Gamma' \Sequent !C$} يا
    \LR{$\Gamma' \Sequent !C$} ثبوت~\LR{$\delta_1$} شته، چې
    \LR{$\Gamma'$} له~\LR{$\Gamma$} سره سمون لري۔ په وروستي حالت کښې
    \LR{$\delta_1$} د \LR{$\delta$} په توګه اخستلے شو۔ په لومړي کښې،
    ځکه \LR{$x:!A$} په پاييز سېکوېنټ کښې راځي، \LR{$x:!A$} په~\LR{$\delta_1$}
    کښې د~\LR{$x$} د فرض ايستلو نښې لرونکي هېڅ استنتاج کښې فعال
    نۀ دے۔ هر بديهيت \LR{$x:!A \Sequent !A$} په دې ثبوت~\LR{$\delta_2$}
    بدلوو:
    \[
    \Axiom$x: !A \land !B \fCenter !A \land !B$
    \RightLabel{\Elim{\land}}
    \UnaryInf$x: !A \land !B \fCenter !A$
    \DisplayProof
    \]
    او د \LR{$x:!A$} هر مورد په \LR{$x:!A \land !B$} بدلوو، يعنې
    \LR{$\Subst{\delta_1}{\delta_2}{x:!A}$} اخلو۔ دا د
    \LR{$x:!A \land !B, \Gamma' \Sequent !C$} ثبوت دے۔
    \item که \LR{$R$} د \Cut قاعده وي، \LR{$\pi$} داسې ختمېږي:
    \[
    \AxiomC{}
    \RightLabel{\LR{$\pi_1$}}
    \Deduce$\Gamma_1 \fCenter !A$
    \AxiomC{}
    \RightLabel{\LR{$\pi_2$}}
    \Deduce$!A, \Gamma_2 \fCenter \Delta$
    \RightLabel{\Cut}
    \BinaryInf$\Gamma_1, \Gamma_2 \fCenter \Delta$
    \DisplayProof
    \]
    د استقرا فرضيه د \LR{$\Gamma_1' \Sequent !A$} ثبوت \LR{$\delta_1$}
    او د \LR{$x:!A, \Gamma_2' \Sequent !C$} يا د
    \LR{$\Gamma_2' \Sequent !C$} ثبوت \LR{$\delta_2$} راکوي۔ په وروستي
    حالت کښې \LR{$\delta$} د~\LR{$\delta_2$} په توګه اخلو۔ که داسې نۀ وي،
    په \LR{$\delta_2$} کښې هر بديهيت \LR{$x:!A \Sequent !A$} په
    \LR{$\delta_1$} بدلوو، څو~\LR{$\delta$} د \LR{$\Gamma_1', \Gamma_2' \Sequent !C$}
    ثبوت~\LR{$\Subst{\delta_2}{\delta_1}{x:!A}$} په توګه تر لاسه کړو۔
  \end{enumerate}
  د پاتې قاعدو لپاره هم له استقرايي ثبوتونو څخه کار اخلو۔
  له استقرا څخه اخيستي ثبوتونه داسې يو ځاے کوو چې د بېلو
  فرضونو نښې سره ګډې نۀ شي۔ د ورګډولو پر وخت محلي ځانګړي
  ثابتونه داسې نوموو چې د نويو پرانيستو فرضونو له ثابتونو
  سره ټکر ونۀ لري۔ په ښي يا کښې پر اړول شوې مقدمې د هماغې څانګې د يا داخلول لګوو۔
  په کيڼ يا کښې اصلي يا فارمول د يوه تازه نښه شوي فرض په توګه
  اخلو۔ دوه فرعي ثبوتونه د يا د ايستلو کوچنۍ مقدمې ګرځوو او
  د هرې څانګې نښه شوے فرض باسو، که هغه فرض پکې کارېدلے وي۔
  د دوو څانګو د پخوانيو نښو له يو بل سره د ګډېدو مخه نيسو۔
  دواړه فرعي ثبوتونه بايد د پاييز ښي اړخ هماغه پايله ولري۔
  که ښے اړخ د يوې مقدمې تش وي او پاييز ښے اړخ فارمول ولري،
  له دروغو څخه د هماغه فارمول د ثابتولو قاعده کاروو۔
  په ښۍ نفي کښې اړول شوې مقدمه د دروغ ثبوت دے؛ نفي داخلوو
  او نښه شوے فرض باسو که استعمال شوے وي۔ په کيڼه نفي کښې
  د نفي شوے فارمول د تازه فرض په توګه اخلو او له اړول شوې
  مقدمې سره د نفي په ايستلو دروغ تر لاسه کوو۔
  په ښي عمومي کميت کښې د عمومي کميت داخلول کاروو۔ په اړول
  شوي کيڼ سياق کښې يوازې له اصلي سياق څخه فارمولونه پاتې دي،
  نو ځانګړے ثابت پکې نۀ راځي۔ په کيڼ عمومي کميت کښې له يوه
  تازه نښه شوي عمومي فرض څخه د عمومي کميت په ايستلو اړوند
  مورد ثابتوو او د هماغه مورد د فرض پر ځاے يې ږدو۔ که د مورد
  فرض استعمال شوے نۀ وي، اړول شوے ثبوت بس دے۔
  په ښي موجود کميت کښې د هماغه مورد له ثبوت څخه موجود کميت
  داخلوو۔ په کيڼ موجود کميت کښې اصلي موجود فارمول د تازه فرض
  په توګه اخلو او له اړول شوې مقدمې سره موجود کميت باسو؛
  د مورد نښه شوے فرض، که موجود وي، په همدې قاعده کښې ايستل کېږي۔
  ځانګړے ثابت په پاييز سېکوېنټ او په لويه مقدمه کښې نۀ راځي؛
  که په فرعي ثبوتونو کښې د محلي ثابتونو د نوم ټکر وي، د هغو
  نومونه له لګولو وړاندې بدلوو۔ ښي انقباض دلته نشي لګېدے،
  ځکه په مقدمه کښې به دوه ښي فارمولونه وغواړي۔
\end{proof}
\olpseganchor{OLP-0670-B011}{end}

\olpseganchor{OLP-0671-B004}{start}
\olfileid{pt}{nat}{tni}
\olpseganchor{OLP-0671-B004}{end}

\olpseganchor{OLP-0671-B005}{start}
\olsection{له \texorpdfstring{\Log{N2}}{N2} څخه \texorpdfstring{\Log{G2}}{G2} ته ژباړه}
\olpseganchor{OLP-0671-B005}{end}

\olpseganchor{OLP-0671-B006}{start}
\begin{prop}\ollabel{prop:N2-to-G2}
که \LR{$\Log{N2c}\ (\Log{N2i}) \Proves \Gamma \Sequent !A$} وي، نو
\LR{$\Log{G2c}\ (\Log{G2i}) + \Cut \Proves \Gamma' \Sequent \Delta$}،
چې \LR{$\Gamma'$} د فارمولونه د تکرار شمېر لرونکے هغه سټ دے
چې له \LR{$\Gamma$} څخه د نښو په لرې کولو جوړېږي؛ که \LR{$!A$} د
\LR{$\lfalse$} فارمول نۀ وي \LR{$\Delta = \{!A\}$}، او که وي
\LR{$\Delta = \emptyset$}۔
\end{prop}
\olpseganchor{OLP-0671-B006}{end}

\olpseganchor{OLP-0671-B007}{start}
\begin{proof}
ثبوت بيا د \LR{$\Gamma \Sequent !A$} د ثبوت~\LR{$\delta$} پر
لوړوالي استقرا ده۔ که \LR{$\delta$} د \LR{$x:!A \Sequent !A$}
بديهيت وي، \LR{$\pi$} اړوند بديهيت \LR{$!A \Sequent !A$} دے، يا
\LR{$\lfalse \Sequent \ $} که \LR{$!A \ident \lfalse$} وي۔
\olpseganchor{OLP-0671-B007}{end}

\olpseganchor{OLP-0671-B008}{start}
که \LR{$\delta$} په استنتاجي قاعده ختمېږي، حالتونه بېلوو:
\begin{enumerate}
  \item که \LR{$R$} د \Intro{\lif} قاعده وي، اصلي فارمول ئې
  \LR{$!B \lif !C$} وي، او \LR{$x:!B$} د مقدمې په سياق کښې موجود خو
  د پايلې په سياق کښې نۀ وي، نو \LR{$\delta$} داسې ختمېږي:
  \[
  \AxiomC{}
  \RightLabel{\LR{$\delta_1$}}
  \Deduce$x:!B, \Gamma \fCenter !C$
  \RightLabel{\Intro{\lif}}
  \UnaryInf$\Gamma \fCenter !B \lif !C$
  \DisplayProof
  \]
  د استقرا فرضيه د \LR{$!B, \Gamma' \Sequent !C$}
  ثبوت~\LR{$\pi_1$} راکوي، يا د \LR{$!B, \Gamma' \Sequent \ $}
  که \LR{$!C \ident \lfalse$} وي۔ له دې څخه \LR{$\pi$} د
  \RightR{\lif} په کارولو اخلو؛ که \LR{$!C \ident \lfalse$} وي،
  لومړے \RightR{\Weakening} کاروو۔
  \item که \LR{$R$} د \Intro{\lif} قاعده وي، خو \LR{$x:!B$} د
  مقدمې په سياق کښې نۀ وي، نو \LR{$\delta$} داسې ختمېږي:
  \[
  \AxiomC{}
  \RightLabel{\LR{$\delta_1$}}
  \Deduce$\Gamma \fCenter !C$
  \RightLabel{\Intro{\lif}}
  \UnaryInf$\Gamma \fCenter !B \lif !C$
  \DisplayProof
  \]
  د استقرا فرضيه د \LR{$\Gamma' \Sequent !C$}
  ثبوت~\LR{$\pi_1$} هغه وخت راکوي چې پاييز فارمول دروغ
  نۀ وي۔ که دروغ وي، د استقرا ثبوت تش ښے اړخ لري؛ لومړے
  ښے اړخ کمزوري کوو چې هماغه دروغ فارمول ښي اړخ ته راشي۔
  بيا په دواړو حالتونو کښې \LR{$\pi$} د~\LR{$!B$} سره د
  \LeftR{\Weakening} او ورپسې د~\RightR{\lif} په کارولو اخلو۔
  \item که \LR{$R$} د \Elim{\lif} قاعده وي، \LR{$\delta$} داسې ختمېږي:
  \[
  \AxiomC{}
  \RightLabel{\LR{$\delta_1$}}
  \Deduce$\Gamma_1 \fCenter !B \lif !A$
  \AxiomC{}
  \RightLabel{\LR{$\delta_2$}}
  \Deduce$\Gamma_2 \fCenter !B$
  \RightLabel{\Elim{\lif}}
  \BinaryInf$\Gamma_1, \Gamma_2 \fCenter !A$
  \DisplayProof
  \]
  چې \LR{$\Gamma = \Gamma_1 \cup \Gamma_2$} دے۔ د استقرا فرضيه
  د \LR{$\Gamma_1' \Sequent !B \lif !A$} ثبوتونه \LR{$\pi_1$} او
  د \LR{$\Gamma_2' \Sequent !B$} ثبوت \LR{$\pi_2$} راکوي، په دې
  وروستي بيان کښې هغه وخت چې کوچنۍ مقدمه دروغ نۀ وي۔ هغه
  حالت چې کوچنۍ مقدمه دروغ وي، لاندې بېل بيانېږي۔ کله چې
  پاييز فارمول هم دروغ نۀ وي، ثبوت~\LR{$\pi$} داسې اخلو:
  \[
  \AxiomC{}
  \RightLabel{\LR{$\pi_1$}}
  \Deduce$\Gamma_1' \fCenter !B \lif !A$
  \AxiomC{}
  \RightLabel{\LR{$\pi_2$}}
  \Deduce$\Gamma_2' \fCenter !B$
  \Axiom$!A \fCenter !A$
  \RightLabel{\LeftR{\lif}}
  \BinaryInf$!B \lif !A, \Gamma_2' \fCenter !A$
  \RightLabel{\Cut}
  \BinaryInf$\Gamma_1', \Gamma_2' \fCenter !A$
  \DisplayProof
  \]
  که پاييز فارمول دروغ وي خو کوچنۍ مقدمه دروغ نۀ وي، په
  همدې رسم کښې د هويت بديهيت پر ځاے د دروغ-تش سېکوېنټ
  بديهيت اخلو؛ د شرطي کيڼې قاعدې او ورپسې د پرېکون پايلې
  هم تش ښے اړخ لري۔
  د تکرار شمېر لرونکے سټ \LR{$\Gamma_1' \cup \Gamma_2'$} ښايي
  له \LR{$\Gamma' = (\Gamma_1 \cup \Gamma_2)'$} سره برابر نۀ وي۔
  مثلاً، که دواړه \LR{$\Gamma_1$} او \LR{$\Gamma_2$} د \LR{$x:!C$} فارمول
  ولري، نو \LR{$\Gamma_1' \cup \Gamma_2'$} د \LR{$!C$} دوه نقلونه لري،
  خو \LR{$(\Gamma_1 \cup \Gamma_2)'$} يوازې يو نقل لري۔ د
  \LeftR{\Contraction} د مناسبو استنتاجونو په زياتولو دا سمولے شو۔
\olpseganchor{OLP-0671-B008}{end}

\olpseganchor{OLP-0671-B009}{start}
  که \LR{$!B \ident \lfalse$} وي، پر \LR{$\delta_2$} د استقرا
  فرضيه د \LR{$\Gamma_2' \Sequent \ $} ثبوت راکوي۔ له دې
  څخه ثبوت~\LR{$\pi$} د \LeftR{\Weakening} د مناسبو
  استنتاجونو په زياتولو جوړوو چې د لوې مقدمې اړين سياق
  هم ولري؛ بيا د~\LR{$!A$} سره \RightR{\Weakening} يوازې هغه
  وخت کاروو چې پاييز فارمول دروغ نۀ وي۔ که پاييز فارمول
  دروغ وي، ښے اړخ هماغه تش پرېږدو۔
\item که \LR{$R$} د \FalseCl قاعده وي، \LR{$\delta$} د ښودل شوي
  نښه لرونکي فرض د موجوديت په حالت کښې داسې ختمېږي:
  \[
  \AxiomC{}
  \RightLabel{\LR{$\delta_1$}}
  \Deduce$x: \lnot !A, \Gamma \fCenter \lfalse$
  \RightLabel{\FalseCl}
  \UnaryInf$\Gamma \fCenter !A$
  \DisplayProof
  \]
  د استقرا له فرضيې د \LR{$\lnot !A, \Gamma' \Sequent \ $}
  ثبوت \LR{$\pi_1$} شته۔ که د ايستلو نښه لرونکے فرض په
  اصل کښې موجود نۀ وي، لومړے کيڼ اړخ کمزوري کوو چې نفي
  شوې مقدمه ورزياته شي۔
  ثبوت~\LR{$\pi$} داسې اخلو:
  \[
  \Axiom$!A \fCenter !A$
  \RightLabel{\RightR{\lnot}}
  \UnaryInf$\fCenter !A, \lnot !A$
  \AxiomC{}
  \RightLabel{\LR{$\pi_1$}}
  \Deduce$\lnot !A, \Gamma' \fCenter $
  \RightLabel{\Cut}
  \BinaryInf$\Gamma' \fCenter !A$
  \DisplayProof
  \]
  که پاييز فارمول دروغ وي، د دې رسم په پاے کښې له
  دروغ-تش بديهيت سره يو بل پرېکون کوو، څو ښے اړخ د
  قضیې له شرط سره سم تش شي۔
\end{enumerate}
د \FalseCl{} په کلاسيکي حالت کښې ثبوت~\LR{$\pi$} منځنے
سېکوېنټ لري چې ښي اړخ ته له يوه فارمول څخه زيات لري۔
د~\Log{N2i} د ثبوت په ژباړه کښې ټول سېکوېنټونه تر يوه
زيات ښي فارمول نۀ لري؛ نو دا يو \Log{G2i}-ثبوت دے۔
د نورو قاعدو د اړولو لپاره په ياد ولرئ چې د دروغ پايله
تش ښے اړخ لري۔ که د يوې ښۍ قاعدې د لګولو لپاره هماغه دروغ
فارمول پکار وي، په ښي کمزوري کولو يې ورزياتوو۔ کله چې د
څو مقدمو د سياقونو په يو ځاے کولو د هماغې پخوانۍ نښې يو
فارمول څو ځله راشي، په کيڼ انقباض يې شمېر راکموو۔ د بېلو
نښو يو شان فارمولونه په خپل شمېر پاتې کېږي۔
د «او» په داخلولو کښې د دواړو اړول شويو مقدمو له پايلو څخه
ښۍ قاعده کاروو۔ د «او» په ايستلو کښې له اړوندې هويتي بديهيې
څخه په کيڼې قاعدې مطلوب جز اخلو او په لويه مقدمه پرېکون کوو۔
که دغه جز دروغ وي، د هويت پر ځاے دروغ-تش بديهيه کاروو۔
د «يا» په داخلولو کښې د اړوندې څانګې ښۍ قاعده کاروو۔
د «يا» په ايستلو کښې په هره کوچنۍ مقدمه کښې د هغې څانګې فرض
له کيڼ کمزوري کولو ورزياتوو که هلته نۀ وي۔ بيا پر دغو
دوو مقدمو د يا کيڼه قاعده کاروو او له لويې مقدمې سره
پرېکون کوو۔ په حدسي نظام کښې د يا دواړه مقدمې ګډ ښے اړخ
لري، چې په پايله کښې هماغسې پاتې کېږي۔ په کلاسيکي نظام
کښې که ګډ فارمول دوه ځله راغلے وي، په ښي انقباض يې يو
ځل ساتو۔ د دروغ ګډې پايلې ښے اړخ په دواړو نظامونو کښې
تش پاتې کېږي۔
د نفي په داخلولو کښې که نښه شوے فرض نۀ وي، په کيڼ کمزوري
کولو يې ورزياتوو او له تش ښي اړخه د نفي ښۍ قاعده لګوو۔
د نفي په ايستلو کښې پر هغې مقدمې چې د نفي دنننے فارمول
ثابتوي، د نفي کيڼه قاعده لګوو۔ که د دغې مقدمې ښے اړخ تش وي،
لومړے په ښي کمزوري کولو اړين فارمول ورزياتوو۔ د نفي له بلې مقدمې سره
پرېکون کوو۔ پايله تش ښے اړخ لري۔
د کميتونو داخلول د اړوندې ښۍ قاعدې له لارې کېږي۔ په عمومي
ايستلو کښې پر هويتي بديهيه يا دروغ-تش بديهيه د عمومي کميت
کيڼه قاعده لګوو او له اړول شوې لويې مقدمې سره پرېکون کوو۔
په موجود ايستلو کښې پر اړول شوې کوچنۍ مقدمه د موجود کميت
کيڼه قاعده لګوو او له اړول شوې لويې مقدمې سره پرېکون کوو۔
که د مورد فرض په کوچنۍ مقدمه کښې نۀ وي، په کيڼ کمزوري کولو
يې ورزياتوو۔ د عمومي کميت په داخلولو او د موجود کميت په
ايستلو کښې د ځانګړي ثابت شرط د فطري قاعدې له شرط څخه راځي۔
په حدسي دروغ ايستلو کښې که نوې پايله دروغ نۀ وي، په ښي
کمزوري کولو يې ورزياتوو؛ که دروغ وي، تش ښے اړخ ساتو۔
\end{proof}
\olpseganchor{OLP-0671-B009}{end}

\olpseganchor{OLP-0674-B004}{start}
\olchapter{pt}{nor}{عادي کول}
\olpseganchor{OLP-0674-B004}{end}

\olpseganchor{OLP-0672-B004}{start}
\olfileid{pt}{nor}{int}
\olpseganchor{OLP-0672-B004}{end}

\olpseganchor{OLP-0672-B005}{start}
\olsection{پېژندنه}
\olpseganchor{OLP-0672-B005}{end}

\olpseganchor{OLP-0672-B006}{start}
که مو يو شرطي فارمول~\LR{$!A \lif !B$} ثابت کړے وي،
نو ښايي د~\Elim{\lif} په کارولو ئې د~\LR{$!B$} په
ثبوت کښې وکاروئ۔ د دې دپاره البته د~\LR{$!A$}
ثبوت~\LR{$\delta_1$} هم پکار دے۔ د~\LR{$!A \lif !B$} ستاسو
ثبوت~\LR{$\delta_3$} ډېر امکان لري چې په~\Intro{\lif}
ختمېږي۔ دا قاعده له پرانيستي فرض~\LR{$!A$} څخه د~\LR{$!B$}
ثبوت~\LR{$\delta_2$} د~\LR{$!A \lif !B$} په ثبوت بدلوي۔ که
همداسې وي، ستاسو وروستے ثبوت دا بڼه لري:
\[
\AxiomC{\LR{$\Discharge{!A}{x}$}}
\RightLabel{\LR{$\delta_2$}}
\DeduceC{\LR{$!B$}}
\DischargeRule{\Intro{\lif}}{x}
\UnaryInfC{\LR{$!A \lif !B$}}
\AxiomC{}
\RightLabel{\LR{$\delta_1$}}
\DeduceC{\LR{$!A$}}
\RightLabel{\Elim{\lif}}
\BinaryInfC{\LR{$!B$}}
\DisplayProof
\]
ستاسو تګلاره په دې اغېزمنه ده چې له مخکې موجود يو څيز،
د~\LR{$!A \lif !B$} ثبوت، مو بيا وکاروۀ، خو خپله
ثبوت مو اغېزمن نۀ دے۔ د \LR{$!A \lif !B$} داخلول او
سمدستي ئې ايستل يو تاو راتاو ښکاري۔ د~\LR{$!B$} د ثابت
کولو نېغه لار بايد موجوده وي۔
\olpseganchor{OLP-0672-B006}{end}

\olpseganchor{OLP-0672-B007}{start}
په حقيقت کښې، نېغه لار تل شته۔ پورته په څېر
``تاو راتاو''، يعنې داخلولو قاعده چې ورپسې
سمدستي ايستلو قاعده وي، د فطري استنتاج له
ثبوتونه څخه تل لرې کېدے شي۔ دې ته \emph{عادي کول}
وايو، او پايله يو \emph{عادي} ثبوت، يا ثبوت
\emph{په عادي بڼه کښې} وي۔
\olpseganchor{OLP-0672-B007}{end}

\olpseganchor{OLP-0672-B008}{start}
د دې پايلې ثبوت، لکه په سېکوېنټي حساب کښې د پرېکون
لرې کولو قضیه، يو څه پېچلے دے او ډېر حالتونه څېړل غواړي۔
په کلاسيکي \Log{N1c} کښې ئې ثابتول د حدسي \Log{N1i}
په پرتله سخت دي۔ لکه چې د پرېکون لرې کول ښيي چې په
سېکوېنټي حساب کښې د \Cut{} قاعدې په کارولو نور څه
ثابت کولے نۀ شو، عادي کول ښيي چې د تاو راتاو په
کارولو د هغه له مخنيوي څخه زيات څه نۀ شو ثابتولے۔ نور
ورته‌والي هم شته: يو عادي ثبوت د هماغې پايلې له
تر ټولو لنډ تاو راتاو لرونکي ثبوت څخه ډېر اوږد
کېدے شي؛ لکه چې سېکوېنټي ثبوتونه د پرېکونونو سره د
تر ټولو لنډ بې‌پرېکونه ثبوت څخه ډېر لنډ کېدے شي۔
د فرعي فارمول خاصيت، چې يوه پايله تل په داسې
ثبوت کښې ثابت کولے شو چې يوازې د پايلې او
پاتې پرانيستو فرضونو فرعي فارمولونه ولري، په حدسي
د فطري استنتاج په نظام کښې له عادي کولو راځي، لکه په
سېکوېنټي حساب کښې چې له پرېکون لرې کولو راځي۔ په لومړي
ترتيب کښې د کميت‌لرونکو فرعي فارمولونو اړوند تعويضي
موردونه هم د همدې خاصيت په معنا کښې شامل دي۔
د بېلګې په توګه، له يوه شرطي فرض او د هغه له مقدمې څخه پايله
په نېغه ايستلو ثابتېږي؛ شرطي فرض په عمومي ډول د پايلې فرعي
فارمول نۀ وي۔ په دې باب کښې د عادي کولو قضيه د حدسي نظام
دپاره ثابتوو۔
\olpseganchor{OLP-0672-B008}{end}

\olpseganchor{OLP-0677-B004}{start}
\olfileid{pt}{nor}{seg}
\olpseganchor{OLP-0677-B004}{end}

\olpseganchor{OLP-0677-B005}{start}
\olsection{برخې او پرېکونونه}
\olpseganchor{OLP-0677-B005}{end}

\olpseganchor{OLP-0677-B006}{start}
د داخلولو قاعدې ورپسې ايستلو قاعده په
ثبوت کښې يو تاو راتاو دے، او د ثبوت عادي کول د
داسې تاو راتاوونو لرې کول دي۔ خو \Elim{\lor} او
\Elim{\lexists} د ``هغه تاو راتاو چې لرې کول ئې غواړو''
تعريف هم لږ پېچلے کوي۔ ستونزه دا ده چې د دغو قواعدو له امله
په تاو راتاو کښې ايستلو قاعده ضرور نۀ ده چې له
داخلولو قاعدې \emph{سمدستي} وروسته راشي۔ په منځ
کښې يو \Elim{\lor} يا \Elim{\lexists} راتلے شي، مثلاً:
\[
\AxiomC{}
\DeduceC{\LR{$\lexists[x][!C(x)]$}}
\AxiomC{\LR{$\Discharge{!A}{x}$}}
\DeduceC{\LR{$!B$}}
\DischargeRule{\Intro{\lif}}{x}
\UnaryInfC{\LR{$!A \lif !B$}}
\DischargeRule{\Elim{\lexists}}{y}
\BinaryInfC{\LR{$!A \lif !B$}}
\AxiomC{}
\DeduceC{\LR{$!A$}}
\RightLabel{\Elim{\lif}}
\BinaryInfC{\LR{$!B$}}
\DisplayProof
\]
په حقيقت کښې هر څومره \Elim{\lor} يا \Elim{\lexists}
کولے شي \Intro{\lif} له \Elim{\lif} څخه بېل کړي، مثلاً:
\[
\AxiomC{}
\DeduceC{\LR{$\lexists[x][!C(x)]$}}
\AxiomC{}
\DeduceC{\LR{$!D \lor !E$}}
\AxiomC{\LR{$\Discharge{!A}{x}$}}
\DeduceC{\LR{$!B$}}
\DischargeRule{\Intro{\lif}}{x}
\UnaryInfC{\LR{$!A \lif !B$}}
\AxiomC{\LR{$\Discharge{!A}{x}$}}
\DeduceC{\LR{$!B$}}
\DischargeRule{\Intro{\lif}}{x}
\UnaryInfC{\LR{$!A \lif !B$}}
\RightLabel{\Elim{\lor}}
\TrinaryInfC{\LR{$!A \lif !B$}}
\DischargeRule{\Elim{\lexists}}{y}
\BinaryInfC{\LR{$!A \lif !B$}}
\AxiomC{}
\DeduceC{\LR{$!A$}}
\RightLabel{\Elim{\lif}}
\BinaryInfC{\LR{$!B$}}
\DisplayProof
\]
دا بېلګه ښيي چې يو~\Elim{\lif} په څو تاو راتاوونو کښې
راتلے شي، ځکه له يوه څخه د زياتو~\Intro{\lif} وروسته دے۔
د دغو ټولو امکانونو د رانغاړلو دپاره تاو راتاوونه د
فارمول د موردونو په هغو لړيو تعريفوو چې په
\Log{N1i}-ثبوت کښې د \emph{برخو} په نوم يادېږي۔ زموږ په
بېلګه کښې د~\LR{$!A \lif !B$} د موردونو يوه داسې لړۍ ده چې د
\LR{$\Elim{\lor}$} يا \LR{$\Elim{\lexists}$} له کوچنۍ مقدمې وروسته
د هماغه استنتاج پايله راځي۔ په دې بېلګه کښې دوه داسې برخې
شته چې هره يوه د~\LR{$!A \lif !B$} درې موردونه لري۔ اوس ئې کره
تعريف وړاندې کوو۔
\olpseganchor{OLP-0677-B006}{end}

\olpseganchor{OLP-0677-B007}{start}
\begin{defn}
په \Log{N1i} ثبوت~\LR{$\delta$} کښې \emph{برخه} د هماغه
فارمول~\LR{$!A$} د موردونو يوه لړۍ \LR{$!A_1$}، \dots، \LR{$!A_n$}
ده چې په \LR{$\delta$} کښې لاندې شرطونه پوره کوي:
\begin{enumerate}
  \item \LR{$!A_1$} د \Elim{\lor} يا \Elim{\lexists} پايله نۀ ده؛
  \item \LR{$!A_n$} د \Elim{\lor} يا \Elim{\lexists} کوچنۍ مقدمه نۀ ده؛
  \item د \LR{$i = 1$}، \dots،~\LR{$n-1$} دپاره، \LR{$!A_i$} د
  \Elim{\lor} يا \Elim{\lexists} کوچنۍ مقدمه ده او \LR{$A_{i+1}$}
  د هماغه استنتاج پايله ده۔
\end{enumerate}
د دغې برخې \emph{اوږدوالے}~\LR{$n$} دے۔
\end{defn}
\olpseganchor{OLP-0677-B007}{end}

\olpseganchor{OLP-0677-B008}{start}
هره برخه د لرې کولو هدف، يعنې تاو راتاو، نۀ ده۔ هغه
برخې چې داسې هدفونه وي، د \emph{پرېکونونو} په نوم يادېږي۔
\olpseganchor{OLP-0677-B008}{end}

\olpseganchor{OLP-0677-B009}{start}
\begin{defn}
\emph{پرېکون} هغه برخه ده چې \LR{$!A_n$} د ايستلو
استنتاج لويه مقدمه وي، او يا \LR{$n=1$} او \LR{$!A_1 = !A_n$} د
داخلولو استنتاج يا~\FalseInt پايله وي، يا~\LR{$n>1$}
وي۔
\end{defn}
\olpseganchor{OLP-0677-B009}{end}

\olpseganchor{OLP-0677-B010}{start}
پام وکړئ چې د پرېکون برخه ضرور نۀ ده چې د
داخلولو قاعدې په پايله پيل شي۔ لازمه دا ده چې
د ايستلو قاعدې په لويه مقدمه پاے ته ورسېږي۔ خو
د~\LR{$1$} اوږدوالي برخه هغه وخت پرېکون ګڼل کېږي چې د
داخلولو قاعدې يا د دروغو ايستلو پايله وي۔
\olpseganchor{OLP-0677-B010}{end}

\olpseganchor{OLP-0677-B011}{start}
\begin{defn}
د پرېکون \emph{رتبه}~\LR{$\depth{!A}$} ده۔ د
ثبوت~\LR{$\delta$} د \emph{پرېکون رتبه}~\LR{$\cutrank{\delta}$}
په~\LR{$\delta$} کښې د پرېکونونو تر ټولو لوړه رتبه ده۔ پرېکون په
~\LR{$\delta$} کښې د \emph{تر ټولو لوړې رتبې} بلل کېږي که رتبه
ئې~\LR{$\cutrank{\delta}$} وي، يعنې تر ټولو لوړه رتبه ولري۔
د~\LR{$\delta$} \emph{پرېکون اوږدوالے} د هغې د تر ټولو لوړې رتبې
د ټولو پرېکونونو د اوږدوالو مجموعه ده۔
\end{defn}
\olpseganchor{OLP-0677-B011}{end}

\olpseganchor{OLP-0675-B004}{start}
\olfileid{pt}{nor}{per}
\olpseganchor{OLP-0675-B004}{end}

\olpseganchor{OLP-0675-B005}{start}
\olsection{د ځاے بدلولو بدلونونه}
\olpseganchor{OLP-0675-B005}{end}

\olpseganchor{OLP-0675-B006}{start}
په \Log{N1i} کښې پرېکونونه هغه برخې دي چې د
ايستلو قاعدې په لويه مقدمه پاے ته رسېږي۔ د عادي کولو
په ثبوت کښې يو حالت دا دے چې د تر ټولو لوړې رتبې د پرېکونونو
او ورسره د ثبوت د پرېکون اوږدوالے کم کړو۔ د \LR{$>1$} اوږدوالي
پرېکون په بېلو طريقو پاے ته رسېدے شي: دا په دې اړه لري چې په
پرېکون کښې د~\LR{$!A$} وروستے فارمول مورد د کومې قاعدې
(\Elim{\lor} يا \Elim{\lexists}) پايله ده، او د کوم
ايستلو استنتاج لويه مقدمه ده۔
\olpseganchor{OLP-0675-B006}{end}

\olpseganchor{OLP-0675-B007}{start}
مثلاً، که اړوند استنتاجونه \Elim{\lor} او \Elim{\land}
وي، \LR{$!A \ident !B \land !C$} وي او داسې فرعي
ثبوت~\LR{$\delta_1$} وي چې په لاندې رسم پاے ته رسېږي:
\[
\AxiomC{}
\RightLabel{\LR{$\delta_2$}}
\DeduceC{\LR{$!D \lor !E$}}
\AxiomC{\LR{$\Discharge{!D}{x}$}}
\RightLabel{\LR{$\delta_3$}}
\DeduceC{\LR{$!B \land !C$}}
\AxiomC{\LR{$\Discharge{!E}{y}$}}
\RightLabel{\LR{$\delta_4$}}
\DeduceC{\LR{$!B \land !C$}}
\DischargeRule{\Elim{\lor}}{x\,y}
\TrinaryInfC{\LR{$!B \land !C$}}
\RightLabel{\Elim{\land}[1]}
\UnaryInfC{\LR{$!B$}}
\DisplayProof
\]فارمول مورد~\LR{$!A_n$} د~\Elim{\lor}
پايله، يعنې لاندې \LR{$!B \land !C$} دے۔ له هغې يو مخکې
فارمول مورد~\LR{$!A_{n-1}$} د~\Elim{\lor} له کوچنيو مقدمو
څخه يوه ده۔
\olpseganchor{OLP-0675-B007}{end}

\olpseganchor{OLP-0675-B008}{start}
د دې پرېکون د اوږدوالي کمولو دپاره د \Elim{\lor} او
\Elim{\land} استنتاجونو ترتيب بدلولے، يعنې ``ځاے بدلولے'' شو۔
په~\LR{$\delta$} کښې فرعي ثبوت~\LR{$\delta_1$} په لاندې رسم بدلوو:
\[
\AxiomC{}
\RightLabel{\LR{$\delta_2$}}
\DeduceC{\LR{$!D \lor !E$}}
\AxiomC{\LR{$\Discharge{!D}{x}$}}
\RightLabel{\LR{$\delta_3$}}
\DeduceC{\LR{$!B \land !C$}}
\RightLabel{\Elim{\land}[1]}
\UnaryInfC{\LR{$!B$}}
\AxiomC{\LR{$\Discharge{!E}{y}$}}
\RightLabel{\LR{$\delta_4$}}
\DeduceC{\LR{$!B \land !C$}}
\RightLabel{\Elim{\land}[1]}
\UnaryInfC{\LR{$!B$}}
\DischargeRule{\Elim{\lor}}{x\,y}
\TrinaryInfC{\LR{$!B$}}
\DisplayProof
\]
هغه پرېکونونه چې د \Elim{\lor} لاندې د \Elim{\land} په
مقدمه کښې ختمېدل، اوس د \Elim{\lor} د کوچنيو مقدمو د پاسه
د يوه \Elim{\land} استنتاج په مقدمه کښې ختمېږي۔ يعنې د هر
يو اوږدوالے په~\LR{$1$} کم شوے دے۔
\olpseganchor{OLP-0675-B008}{end}

\olpseganchor{OLP-0675-B009}{start}
سربېره پر دې، هر هغه پرېکون چې بشپړ په \LR{$\delta_2$}،
\LR{$\delta_3$} يا \LR{$\delta_4$} کښې وي، يا په~\LR{$\delta$} کښې بشپړ له
\LR{$\delta_1$} څخه بهر وي، نۀ بدلېږي: هماغه فارمولونه رانغاړي
او له هماغو استنتاجونو تېرېږي۔ نو په ځانګړي ډول د هغې رتبه
او اوږدوالے په~\LR{$\delta$} کښې د اړوند پرېکون سره عين دي۔ له دې
امله د نوي ثبوت د پرېکون رتبه نۀ ده لوړه شوې او د نوي
ثبوت د پرېکون اوږدوالے د زاړه ثبوت څخه کم دے۔
\olpseganchor{OLP-0675-B009}{end}

\olpseganchor{OLP-0675-B010}{start}
\begin{prob}
د ځاے بدلولو له بدلون وروسته لږ تر لږه د يوه پرېکون اوږدوالے
په~\LR{$1$} کمېږي۔ که \Elim{\land} له يوه~\Elim{\lor} څخه تېر
کړو، په حقيقت کښې د دوو پرېکون برخو اوږدوالے کمېږي؛ نو په دې
حالت کښې د ثبوت د پرېکون اوږدوالے لږ تر لږه په~\LR{$2$}
کمېږي۔ آيا يو داسې بدلون د ثبوت د پرېکون اوږدوالے له~\LR{$2$}
څخه زيات هم کمولے شي؟ آيا د پرېکون \emph{رتبه} هم ټيټېدې شي؟
\end{prob}
\olpseganchor{OLP-0675-B010}{end}

\olpseganchor{OLP-0675-B011}{start}
د \Elim{\lor} يا \Elim{\lexists} قاعدې د پاسه د
ايستلو قاعدې د وړلو عمل د \emph{ځاے بدلولو بدلون}
بلل کېږي۔ د ځينو قاعده‌جوړو بڼې په
\olref{tab:permutation-conversions} کښې ورکړې دي۔ پاتې بڼې
د مشقونو په توګه پرېښودل شوې دي۔
\olpseganchor{OLP-0675-B011}{end}

\olpseganchor{OLP-0675-B012}{start}
% \begin{table}
% \begin{multline*}
% \shoveleft{\AxiomC{$!D \lor !E$}
% \AxiomC{$\Discharge{!D}{x}$}
% \shortDeduce
% \DeduceC{$!B \land !C$}
% \AxiomC{$\Discharge{!E}{y}$}
% \shortDeduce
% \DeduceC{$!B \land !C$}
% \DischargeRule{\Elim{\lor}}{x\,y}
% \TrinaryInfC{$!B \land !C$}
% \RightLabel{\Elim{\land}[1]}
% \UnaryInfC{$!B$}
% \DisplayProof}\\
% \shoveright{\longrightarrow\ 
% \AxiomC{$!D \lor !E$}
% \AxiomC{$\Discharge{!D}{x}$}
% \shortDeduce
% \DeduceC{$!B \land !C$}
% \RightLabel{\Elim{\land}[1]}
% \UnaryInfC{$!B$}
% \AxiomC{$\Discharge{!E}{y}$}
% \shortDeduce
% \DeduceC{$!B \land !C$}
% \RightLabel{\Elim{\land}[1]}
% \UnaryInfC{$!B$}
% \DischargeRule{\Elim{\lor}}{x\,y}
% \TrinaryInfC{$!B$}
% \DisplayProof}
% \\
% \shoveleft{\AxiomC{$!D \lor !E$}
% \AxiomC{$\Discharge{!D}{x}$}
% \shortDeduce
% \DeduceC{$!B \lif !C$}
% \AxiomC{$\Discharge{!E}{y}$}
% \shortDeduce
% \DeduceC{$!B \lif !C$}
% \DischargeRule{\Elim{\lor}}{x\,y}
% \TrinaryInfC{$!B \lif !C$}
% \AxiomC{$!B$}
% \RightLabel{\Elim{\lif}}
% \BinaryInfC{$!C$}
% \DisplayProof}\\
% \shoveright{\longrightarrow\ 
% \AxiomC{$!D \lor !E$}
% \AxiomC{$\Discharge{!D}{x}$}
% \shortDeduce
% \DeduceC{$!B \lif !C$}
% \AxiomC{$!B$}
% \RightLabel{\Elim{\lif}}
% \BinaryInfC{$!C$}
% \AxiomC{$\Discharge{!E}{y}$}
% \shortDeduce
% \DeduceC{$!B \lif !C$}
% \AxiomC{$!B$}
% \RightLabel{\Elim{\lif}}
% \BinaryInfC{$!C$}
% \DischargeRule{\Elim{\lor}}{x\,y}
% \TrinaryInfC{$!C$}
% \DisplayProof}
% \\
% \shoveleft{\AxiomC{$!D \lor !E$}
% \AxiomC{$\Discharge{!D}{x}$}
% \shortDeduce
% \DeduceC{$!B \lif !C$}
% \AxiomC{$\Discharge{!E}{y}$}
% \shortDeduce
% \DeduceC{$!B \lif !C$}
% \DischargeRule{\Elim{\lor}}{x\,y}
% \TrinaryInfC{$!B \lif !C$}
% \AxiomC{$!B$}
% \RightLabel{\Elim{\lif}}
% \BinaryInfC{$!C$}
% \DisplayProof}\\
% \shoveright{\longrightarrow\ 
% \AxiomC{$!D \lor !E$}
% \AxiomC{$\Discharge{!D}{x}$}
% \shortDeduce
% \DeduceC{$!B \lif !C$}
% \AxiomC{$!B$}
% \RightLabel{\Elim{\lif}}
% \BinaryInfC{$!C$}
% \AxiomC{$\Discharge{!E}{y}$}
% \shortDeduce
% \DeduceC{$!B \lif !C$}
% \AxiomC{$!B$}
% \RightLabel{\Elim{\lif}}
% \BinaryInfC{$!C$}
% \DischargeRule{\Elim{\lor}}{x\,y}
% \TrinaryInfC{$!C$}
% \DisplayProof}
% \\
% \shoveleft{\AxiomC{$!D \lor !E$}
% \AxiomC{$\Discharge{!D}{x}$}
% \shortDeduce
% \DeduceC{$\lforall[x][!B(x)]$}
% \AxiomC{$\Discharge{!E}{y}$}
% \shortDeduce
% \DeduceC{$\lforall[x][!B(x)]$}
% \DischargeRule{\Elim{\lor}}{x\,y}
% \TrinaryInfC{$\lforall[x][!B(x)]$}
% \RightLabel{\Elim{\lforall}}
% \UnaryInfC{$!B(t)$}
% \DisplayProof}\\
% \shoveright{\longrightarrow\ 
% \AxiomC{$!D \lor !E$}
% \AxiomC{$\Discharge{!D}{x}$}
% \shortDeduce
% \DeduceC{$\lforall[x][!B(x)]$}
% \RightLabel{\Elim{\lforall}}
% \UnaryInfC{$!B(t)$}
% \AxiomC{$\Discharge{!E}{y}$}
% \shortDeduce
% \DeduceC{$\lforall[x][!B(x)]$}
% \RightLabel{\Elim{\lforall}}
% \UnaryInfC{$!B(t)$}
% \DischargeRule{\Elim{\lor}}{x\,y}
% \TrinaryInfC{$!B(t)$}
% \DisplayProof}
% \end{multline*}
% \caption{Some permutation conversions for \Log{N1i}.}
% \ollabel{tab:permutation-conversions}
% \end{table}
\olpseganchor{OLP-0675-B012}{end}

\olpseganchor{OLP-0675-B013}{start}
{\small
\begin{longtable}{@{}l@{}}
\AxiomC{\LR{$!D \lor !E$}}
\AxiomC{\LR{$\Discharge{!D}{x}$}}
\shortDeduce
\DeduceC{\LR{$!B \land !C$}}
\AxiomC{\LR{$\Discharge{!E}{y}$}}
\shortDeduce
\DeduceC{\LR{$!B \land !C$}}
\DischargeRule{\Elim{\lor}}{x\,y}
\TrinaryInfC{\LR{$!B \land !C$}}
\RightLabel{\Elim{\land}[1]}
\UnaryInfC{\LR{$!B$}}
\DisplayProof
\LR{$\longrightarrow$}\\[6ex]
\quad\AxiomC{\LR{$!D \lor !E$}}
\AxiomC{\LR{$\Discharge{!D}{x}$}}
\shortDeduce
\DeduceC{\LR{$!B \land !C$}}
\RightLabel{\Elim{\land}[1]}
\UnaryInfC{\LR{$!B$}}
\AxiomC{\LR{$\Discharge{!E}{y}$}}
\shortDeduce
\DeduceC{\LR{$!B \land !C$}}
\RightLabel{\Elim{\land}[1]}
\UnaryInfC{\LR{$!B$}}
\DischargeRule{\Elim{\lor}}{x\,y}
\TrinaryInfC{\LR{$!B$}}
\DisplayProof
\\[10ex]
\AxiomC{\LR{$!D \lor !E$}}
\AxiomC{\LR{$\Discharge{!D}{x}$}}
\shortDeduce
\DeduceC{\LR{$!B \lif !C$}}
\AxiomC{\LR{$\Discharge{!E}{y}$}}
\shortDeduce
\DeduceC{\LR{$!B \lif !C$}}
\DischargeRule{\Elim{\lor}}{x\,y}
\TrinaryInfC{\LR{$!B \lif !C$}}
\AxiomC{\LR{$!B$}}
\RightLabel{\Elim{\lif}}
\BinaryInfC{\LR{$!C$}}
\DisplayProof
\LR{$\longrightarrow$}\\[6ex]
\AxiomC{\LR{$!D \lor !E$}}
\AxiomC{\LR{$\Discharge{!D}{x}$}}
\shortDeduce
\DeduceC{\LR{$!B \lif !C$}}
\AxiomC{\LR{$!B$}}
\RightLabel{\Elim{\lif}}
\BinaryInfC{\LR{$!C$}}
\AxiomC{\LR{$\Discharge{!E}{y}$}}
\shortDeduce
\DeduceC{\LR{$!B \lif !C$}}
\AxiomC{\LR{$!B$}}
\RightLabel{\Elim{\lif}}
\BinaryInfC{\LR{$!C$}}
\DischargeRule{\Elim{\lor}}{x\,y}
\TrinaryInfC{\LR{$!C$}}
\DisplayProof
\\[10ex]
\AxiomC{\LR{$!D \lor !E$}}
\AxiomC{\LR{$\Discharge{!D}{x}$}}
\shortDeduce
\DeduceC{\LR{$!B \lif !C$}}
\AxiomC{\LR{$\Discharge{!E}{y}$}}
\shortDeduce
\DeduceC{\LR{$!B \lif !C$}}
\DischargeRule{\Elim{\lor}}{x\,y}
\TrinaryInfC{\LR{$!B \lif !C$}}
\AxiomC{\LR{$!B$}}
\RightLabel{\Elim{\lif}}
\BinaryInfC{\LR{$!C$}}
\DisplayProof
\LR{$\longrightarrow$}\\[6ex]
\AxiomC{\LR{$!D \lor !E$}}
\AxiomC{\LR{$\Discharge{!D}{x}$}}
\shortDeduce
\DeduceC{\LR{$!B \lif !C$}}
\AxiomC{\LR{$!B$}}
\RightLabel{\Elim{\lif}}
\BinaryInfC{\LR{$!C$}}
\AxiomC{\LR{$\Discharge{!E}{y}$}}
\shortDeduce
\DeduceC{\LR{$!B \lif !C$}}
\AxiomC{\LR{$!B$}}
\RightLabel{\Elim{\lif}}
\BinaryInfC{\LR{$!C$}}
\DischargeRule{\Elim{\lor}}{x\,y}
\TrinaryInfC{\LR{$!C$}}
\DisplayProof
\\[10ex]
\AxiomC{\LR{$!D \lor !E$}}
\AxiomC{\LR{$\Discharge{!D}{x}$}}
\shortDeduce
\DeduceC{\LR{$\lforall[x][!B(x)]$}}
\AxiomC{\LR{$\Discharge{!E}{y}$}}
\shortDeduce
\DeduceC{\LR{$\lforall[x][!B(x)]$}}
\DischargeRule{\Elim{\lor}}{x\,y}
\TrinaryInfC{\LR{$\lforall[x][!B(x)]$}}
\RightLabel{\Elim{\lforall}}
\UnaryInfC{\LR{$!B(t)$}}
\DisplayProof
\quad\LR{$\longrightarrow$}\\[8ex]
\quad\AxiomC{\LR{$!D \lor !E$}}
\AxiomC{\LR{$\Discharge{!D}{x}$}}
\shortDeduce
\DeduceC{\LR{$\lforall[x][!B(x)]$}}
\RightLabel{\Elim{\lforall}}
\UnaryInfC{\LR{$!B(t)$}}
\AxiomC{\LR{$\Discharge{!E}{y}$}}
\shortDeduce
\DeduceC{\LR{$\lforall[x][!B(x)]$}}
\RightLabel{\Elim{\lforall}}
\UnaryInfC{\LR{$!B(t)$}}
\DischargeRule{\Elim{\lor}}{x\,y}
\TrinaryInfC{\LR{$!B(t)$}}
\DisplayProof
\\
\caption{د \Log{N1i} د ځاے بدلولو ځينې بدلونونه۔}
\ollabel{tab:permutation-conversions}
\end{longtable}
}
\olpseganchor{OLP-0675-B013}{end}

\olpseganchor{OLP-0675-B014}{start}
\begin{prob}
د \Elim{\lor} او ورپسې \Elim{\lor} يا \Elim{\lnot} د ځاے
بدلولو بدلونونه ورکړئ۔
\end{prob}
\olpseganchor{OLP-0675-B014}{end}

\olpseganchor{OLP-0675-B015}{start}
\begin{prob}
د \Elim{\lexists} او ورپسې \Elim{\exists} د ځاے بدلولو
بدلونونه ورکړئ۔ ووايئ چې په وروستي حالت کښې ولې د ځانګړي
متغير هېڅ شرط نۀ ماتېږي۔
\end{prob}
\olpseganchor{OLP-0675-B015}{end}

\olpseganchor{OLP-0675-B016}{start}
بايد ځان ډاډه کړو چې د ځاے بدلولو د بدلون په تطبيق سره
د ځانګړي متغير کوم شرط نۀ ماتوو۔ دا رسم وګورئ:
\[
\AxiomC{}
\RightLabel{\LR{$\delta_2$}}
\DeduceC{\LR{$!D \lor !E$}}
\AxiomC{\LR{$\Discharge{!D}{x}$}}
\RightLabel{\LR{$\delta_3$}}
\DeduceC{\LR{$\lexists[x][!B(x)]$}}
\AxiomC{\LR{$\Discharge{!E}{y}$}}
\RightLabel{\LR{$\delta_4$}}
\DeduceC{\LR{$\lexists[x][!B(x)]$}}
\DischargeRule{\Elim{\lor}}{x\,y}
\TrinaryInfC{\LR{$\lexists[x][!B(x)]$}}
\AxiomC{\LR{$\Discharge{!B(c)}{z}$}}
\RightLabel{\LR{$\delta_5$}}
\DeduceC{\LR{$!C$}}
\DischargeRule{\Elim{\lexists}}{z}
\BinaryInfC{\LR{$!C$}}
\DisplayProof
\]
که دا په لاندې رسم بدل کړو:
\[
\AxiomC{}
\RightLabel{\LR{$\delta_2$}}
\DeduceC{\LR{$!D \lor !E$}}
\AxiomC{\LR{$\Discharge{!D}{x}$}}
\RightLabel{\LR{$\delta_3$}}
\DeduceC{\LR{$\lexists[x][!B(x)]$}}
\AxiomC{\LR{$\Discharge{!B(c)}{z}$}}
\RightLabel{\LR{$\delta_5$}}
\DeduceC{\LR{$!C$}}
\DischargeRule{\Elim{\lexists}}{z}
\BinaryInfC{\LR{$!C$}}
\AxiomC{\LR{$\Discharge{!E}{y}$}}
\RightLabel{\LR{$\delta_4$}}
\DeduceC{\LR{$\lexists[x][!B(x)]$}}
\AxiomC{\LR{$\Discharge{!B(c)}{z}$}}
\RightLabel{\LR{$\delta_5$}}
\DeduceC{\LR{$!C$}}
\DischargeRule{\Elim{\lexists}}{z}
\BinaryInfC{\LR{$!C$}}
\DischargeRule{\Elim{\lor}}{x\,y}
\TrinaryInfC{\LR{$!C$}}
\DisplayProof
\]
اندېښنه کېدے شي چې ځانګړے متغير~\LR{$c$} مثلاً په~\LR{$!D$} کښې
راشي۔ ځکه فرض \LR{$!D$} د اصلي~\LR{$\Elim{\lexists}$} د پاسه ايستل
کېږي؛ نو په اصلي ثبوت کښې دا د~\Elim{\exists} د لويې
مقدمې د پاسه په کوم پرانيستي فرض کښې نۀ راځي او د ځانګړي
متغير شرط هلته پوره دے۔ خو په نوي ثبوت کښې \LR{$!D$}~تر
هغې نۀ ايستل کېږي چې د~\Elim{\lexists} استنتاجونه تېر شي؛
نو د ځانګړي متغير شرط به مات شي۔ خو په ياد ولرئ چې زموږ
ثبوتونه منظم فرض شوي دي: هر ځانګړے متغير يوازې د يوه
استنتاج ځانګړے متغير دے او په ثبوت کښې يوازې د
\Elim{\lexists} د کوچنۍ مقدمې د پاسه راځي۔ په ځانګړي ډول،
\LR{$c$}~په \LR{$\delta_3$} يا~\LR{$\delta_4$} کښې نۀ شي راتلے۔
\olpseganchor{OLP-0675-B016}{end}

\olpseganchor{OLP-0675-B017}{start}
دا مثال بله خبره هم ښيي: په نوي ثبوت کښې
د~\LR{$\delta_5$} دوه نقلونه شته۔ نو که \LR{$\delta_5$} د تر ټولو لوړې
رتبې پرېکون ولري، د پرېکون اوږدوالے مو \emph{نۀ} دے کم
کړے! ځکه بايد يوازې هغه وخت د ځاے بدلولو بدلون تطبيق کړو
چې په \LR{$\delta_5$} کښې د تر ټولو لوړې رتبې پرېکون نۀ وي۔ دا
د تل لپاره د تر ټولو \emph{پاسني} لوړې رتبې پرېکون په
غوره کولو يقيني کوو۔ يعنې داسې د تر ټولو لوړې رتبې پرېکون
غوره کوو چې فرعي ثبوت، دلته~\LR{$\delta_1$}، داسې بل د تر ټولو
لوړې رتبې پرېکون نۀ لري چې د همدغه فرعي ثبوت د وروستي
استنتاج د پاسه ختمېږي۔ نو په~\LR{$\delta_5$}
کښې، او هم په \LR{$\delta_2$}، \LR{$\delta_3$} يا~\LR{$\delta_4$} کښې،
د نورو د تر ټولو لوړې رتبې پرېکونونو امکان له منځه ځي۔
\olpseganchor{OLP-0675-B017}{end}

\olpseganchor{OLP-0676-B004}{start}
\olfileid{pt}{nor}{red}
\olpseganchor{OLP-0676-B004}{end}

\olpseganchor{OLP-0676-B005}{start}
\olsection{د کمولو بدلونونه}
\olpseganchor{OLP-0676-B005}{end}

\olpseganchor{OLP-0676-B006}{start}
لومړے د معرفي کولو د قاعدې حالت اخلو۔ که په \Log{N1i} کښې د
ثبوت د پرېکون اوږدوالے~\LR{$1$} وي، دا يوه داسې برخه ده چې د
فارمول يو مورد لري، او همدغه فارمول هم د
داخلولو استنتاج پايله او هم د ايستلو استنتاج
لويه مقدمه ده۔ د عادي کولو په ثبوت کښې دا ډول پرېکونونه داسې
``کمول'' کېږي چې په دغې داخلولو/ايستلو جوړې
پاے ته رسېدونکے فرعي ثبوت په نوي ثبوت بدل شي، چې
د استنتاجونو دا جوړه ترې لرې شوې وي۔ داسې بدلول نوي پرېکونونه
پيدا کولے شي۔ خو کتلے شو چې د نويو پرېکونونو رتبه د کمېدونکي
پرېکون له رتبې ټيټه ده۔ نو د نوي ثبوت يا د پرېکون اوږدوالے
کم دے، ځکه د تر ټولو لوړې رتبې يو د~\LR{$1$} اوږدوالي پرېکون لرې شو،
يا د ثبوت د پرېکون رتبه ټيټه شوې ده، که همدا د تر ټولو
لوړې رتبې يوازينے پرېکون و۔
\olpseganchor{OLP-0676-B006}{end}

\olpseganchor{OLP-0676-B007}{start}
مثلاً، که اړوند استنتاجونه \Intro{\lif} او \Elim{\lif}
وي، \LR{$!A \ident !B \lif !C$} وي او داسې فرعي
ثبوت~\LR{$\delta_1$} وي چې په لاندې رسم پاے ته رسېږي:
\[
\AxiomC{\LR{$\Discharge{!B}{x}$}}
\RightLabel{\LR{$\delta_2$}}
\DeduceC{\LR{$!C$}}
\DischargeRule{\Intro{\lif}}{x}
\UnaryInfC{\LR{$!B \lif !C$}}
\AxiomC{}
\RightLabel{\LR{$\delta_3$}}
\DeduceC{\LR{$!B$}}
\RightLabel{\Elim{\lif}}
\BinaryInfC{\LR{$!C$}}
\DisplayProof
\]
د دې پرېکون د لرې کولو دپاره په~\LR{$\delta$} کښې فرعي
ثبوت~\LR{$\delta_1$} په لاندې رسم بدلوو:
\[
\AxiomC{}
\RightLabel{\LR{$\delta_3$}}
\DeduceC{\LR{$!B$}}
\RightLabel{\LR{$\Subst{\delta_2}{\delta_3}{!B^x}$}}
\DeduceC{\LR{$!C$}}
\DisplayProof
\]
په \LR{$\Subst{\delta_2}{\delta_3}{!B^x}$} هغه ثبوت مراد دے چې
له~\LR{$\delta_2$} څخه د هر~\LR{$!B$} فرض، چې نښه ئې~\LR{$x$} ده، د بشپړ
ثبوت~\LR{$\delta_3$} په ځاے کولو حاصلېږي۔ په دې ثبوت کښې
بدل شوي اصلي پرانيستي \LR{$!B$} فرضونه، چې نښه ئې~\LR{$x$} ده، نۀ دي
پاتې؛ خو دننه شوي ورته نښه لرونکي فرضونه پاتې کېدے شي۔ د~\LR{$x$}
نښه لرونکے بل بېل فارمول نۀ شته، ځکه د معمول په شان فرض کوو
چې دوه بېل فارمولونه د يوې نښې لرونکي فرضونه نۀ وي۔ د نوي
ثبوت هر پرانيستے فرض يا له~\LR{$\delta_1$} څخه پاتې نۀ بدل شوے
پرانيستے فرض دے، يا له داخل شوي ثبوت څخه داسې فرض دے چې له
داخلېدو وروسته نۀ دے ايستل شوے۔ په~\LR{$\delta_2$} کښې د فرضونو
ايستونکے هر استنتاج په \LR{$\Subst{\delta_2}{\delta_3}{!B^x}$} کښې
د هماغې نښې فرضونه باسي؛ په هغو کښې نوي دننه شوي فرضونه هم
راتلے شي۔ په~\LR{$\delta_3$} کښې د فرضونو هر ايستونکے استنتاج ښايي
څو نقلونه ولري، چې په \LR{$\Subst{\delta_2}{\delta_3}{!B^x}$} کښې
اړوند فرضونه باسي۔
که د ايستل کېدونکي فرض هېڅ مورد نۀ وي، د هغه د ثبوت د
ورګډولو ځاے هم نشته۔ په ورګډ شوي ثبوت کښې ځينې پرانيستي
فرضونه د نوي ځاے د ايستلو قاعدو تر واک لاندې راتلے شي۔
\olpseganchor{OLP-0676-B007}{end}

\olpseganchor{OLP-0676-B008}{start}
هر پرېکون چې بشپړ په \LR{$\delta_2$} يا \LR{$\delta_3$} کښې وي، يا
په~\LR{$\delta$} کښې بشپړ له~\LR{$\delta_1$} څخه بهر وي، نۀ بدلېږي:
هماغه فارمولونه رانغاړي او له هماغو استنتاجونو تېرېږي۔ نو
په ځانګړي ډول ئې رتبه او اوږدوالے په~\LR{$\delta$} کښې د اړوند
پرېکون سره عين دي۔ د تر ټولو لوړې رتبې يو د~\LR{$1$} اوږدوالي
پرېکون مو لرې کړے دے، نو يا د پرېکون اوږدوالے کم شوے يا ئې
رتبه ټيټه شوې ده، که د~\LR{$\delta$} د پرېکون اوږدوالے~\LR{$1$} و او
همدا د تر ټولو لوړې رتبې يوازينے پرېکون و۔
\olpseganchor{OLP-0676-B008}{end}

\olpseganchor{OLP-0676-B009}{start}
خو دوه خبرې کېدے شي:
\begin{enumerate}
  \item د ټولو \LR{$!B^x$} فرضونو په \LR{$\delta_3$} بدلولو سره مو
  په~\LR{$\delta_3$} کښې د ټولو پرېکونونو نقلونه زيات کړي دي۔ که په
  هغو کښې د تر ټولو لوړې رتبې پرېکونونه وي، د نوي ثبوت د
  پرېکون اوږدوالے د زاړه ثبوت له پرېکون اوږدوالي زياتېدے
  شي۔ بايد ډاډه شو چې داسې نۀ کېږي۔ دا په دې يقيني کوو چې د
  کمولو بدلونونه يوازې په تر ټولو \emph{ښي اړخ} د تر ټولو لوړې
  رتبې پرېکونونو تطبيق کړو: داسې پرېکونونه چې په~\LR{$\delta$} کښې
  ئې ښي اړخ ته په کومه څانګه د تر ټولو لوړې رتبې بل پرېکون نۀ
  وي۔ که په~\LR{$\delta_3$} کښې داسې پرېکون وي، د دې کمېدونکي
  پرېکون ښي اړخ ته به وي، نو دلته به مو تر ټولو ښي پرېکون نۀ
  و غوره کړے۔ تل د تر ټولو ښي پرېکون په غوره کولو سره يقيني
  کوو چې يا د ثبوت د پرېکون اوږدوالے کم کړو يا ئې د پرېکون
  رتبه هم ټيټه کړو۔
  \item که په \LR{$\delta_2$} کښې فرض \LR{$!B^x$} د ايستلو
  استنتاج لويه مقدمه وي او د~\LR{$\delta_3$} وروستے استنتاج
  \Elim{\lor}، \Elim{\lexists} يا داخلولو وي، نو
  په~\LR{$\delta_2$} کښې پر داسې فرض د~\LR{$\delta_3$} په پېوندولو مو
  نوی پرېکون جوړ کړے دے۔ هغه څۀ چې په~\LR{$\delta_2$} کښې يوازې
  فرض و، اوس د~\LR{$1$} اوږدوالي پرېکون دے، يعنې د
  داخلولو استنتاج پايله او د ايستلو استنتاج
  لويه مقدمه، يا د پرېکون د برخې وروستے فارمول مورد دے،
  يعنې د ايستلو استنتاج لويه مقدمه او د \Elim{\lor}
  يا \Elim{\lexists} پايله۔ که \LR{$B^x$} د \Elim{\lor} يا
  \Elim{\lexists} کوچنۍ مقدمه وه، له مخکې د يوې برخې، او ښايي
  د پرېکون د برخې، لومړے فارمول مورد و۔ په نوي ثبوت
  کښې دا برخه اوس اوږدېدې شي۔ خو پام وکړئ چې د دې نوي پرېکون
  يا پخوانۍ پرېکون برخې فارمول~\LR{$!B$} دے او
  \LR{$\depth{!B} < \depth{!B \lif !C}$}۔ نو که مو پرېکونونه زيات
  کړي يا ئې اوږدوالے زيات کړے هم وي، په هغو کښې يو هم د تر
  ټولو لوړې رتبې پرېکون نۀ دے۔ په دې توګه مو يا د پرېکون رتبه
  ټيټه کړې، يا که د هماغې رتبې نور پرېکونونه پاتې وي، د پرېکون
  اوږدوالے مو کم کړے دے۔
\end{enumerate}
\olpseganchor{OLP-0676-B009}{end}

\olpseganchor{OLP-0676-B010}{start}
د~\LR{$1$} اوږدوالي پرېکون د کمولو عمل د \emph{کمولو بدلون}
يا د \emph{تاو راتاو بدلون} بلل کېږي۔ د ځينو قاعده‌جوړو بڼې
په \olref{tab:reduction-conversions} کښې ورکړې دي؛ پاتې بڼې د
مشقونو په توګه پرېښودل شوې دي۔
\olpseganchor{OLP-0676-B010}{end}

\olpseganchor{OLP-0676-B011}{start}
\begin{table}
\begin{multline*}
\bottomAlignProof
\AxiomC{\LR{$\Discharge{!B}{x}$}}
\RightLabel{\LR{$\delta_2$}}
\shortDeduce
\DeduceC{\LR{$!C$}}
\DischargeRule{\Intro{\lif}}{x}
\UnaryInfC{\LR{$!B \lif !C$}}
\AxiomC{}
\RightLabel{\LR{$\delta_3$}}
\shortDeduce
\DeduceC{\LR{$!B$}}
\RightLabel{\Elim{\lif}}
\BinaryInfC{\LR{$!C$}}
\DisplayProof
\quad \longrightarrow\quad
\shoveright{\bottomAlignProof
\AxiomC{}
\RightLabel{\LR{$\delta_3$}}
\shortDeduce
\DeduceC{\LR{$!B$}}
\RightLabel{\LR{$\Subst{\delta_2}{\delta_3}{!B^x}$}}
\shortDeduce
\DeduceC{\LR{$!C$}}
\DisplayProof}
\\[4ex]
\shoveleft{\bottomAlignProof
\AxiomC{}
\RightLabel{\LR{$\delta_2$}}
\shortDeduce
\DeduceC{\LR{$!B$}}
\AxiomC{}
\RightLabel{\LR{$\delta_3$}}
\shortDeduce
\DeduceC{\LR{$!C$}}
\RightLabel{\Intro{\land}}
\BinaryInfC{\LR{$!B \land !C$}}
\RightLabel{\Elim{\land}[1]}
\UnaryInfC{\LR{$!B$}}
\DisplayProof
\quad \longrightarrow\quad
\bottomAlignProof
\AxiomC{}
\RightLabel{\LR{$\delta_2$}}
\shortDeduce
\DeduceC{\LR{$!B$}}
\DisplayProof}\qquad
\shoveright{\bottomAlignProof
\AxiomC{}
\RightLabel{\LR{$\delta_2$}}
\shortDeduce
\DeduceC{\LR{$!B$}}
\AxiomC{}
\RightLabel{\LR{$\delta_3$}}
\shortDeduce
\DeduceC{\LR{$!C$}}
\RightLabel{\Intro{\land}}
\BinaryInfC{\LR{$!B \land !C$}}
\RightLabel{\Elim{\land}[2]}
\UnaryInfC{\LR{$!C$}}
\DisplayProof
\quad \longrightarrow\quad
\bottomAlignProof
\AxiomC{}
\RightLabel{\LR{$\delta_3$}}
\shortDeduce
\DeduceC{\LR{$!C$}}
\DisplayProof}
\\[4ex]
\bottomAlignProof
\AxiomC{}
\RightLabel{\LR{$\delta_2$}}
\shortDeduce
\DeduceC{\LR{$!B$}}
\RightLabel{\Intro{\lor}[1]}
\UnaryInfC{\LR{$!B \lor !C$}}
\AxiomC{\LR{$\Discharge{!B}{x}$}}
\RightLabel{\LR{$\delta_3$}}
\shortDeduce
\DeduceC{\LR{$!D$}}
\AxiomC{\LR{$\Discharge{!C}{y}$}}
\RightLabel{\LR{$\delta_4$}}
\shortDeduce
\DeduceC{\LR{$!D$}}
\DischargeRule{\Elim{\lor}}{x\,y}
\TrinaryInfC{\LR{$!D$}}
\DisplayProof
\quad \longrightarrow\quad
\shoveright{\bottomAlignProof
\AxiomC{}
\RightLabel{\LR{$\delta_2$}}
\shortDeduce
\DeduceC{\LR{$!B$}}
\RightLabel{\LR{$\Subst{\delta_3}{\delta_2}{!B^x}$}}
\shortDeduce
\DeduceC{\LR{$!D$}}
\DisplayProof}
\\[2ex]
\shoveleft{\bottomAlignProof
\AxiomC{}
\RightLabel{\LR{$\delta_2$}}
\shortDeduce
\DeduceC{\LR{$!B(c)$}}
\RightLabel{\Intro{\lforall}}
\UnaryInfC{\LR{$\lforall[x][!B(x)]$}}
\RightLabel{\Elim{\lforall}}
\UnaryInfC{\LR{$!B(t)$}}
\DisplayProof
\quad \longrightarrow\quad
\bottomAlignProof
\AxiomC{}
\RightLabel{\LR{$\Subst{\delta_2}{t}{c}$}}
\shortDeduce
\DeduceC{\LR{$!B(t)$}}
\DisplayProof}
\\[4ex]
\shoveright{\bottomAlignProof
\AxiomC{}
\RightLabel{\LR{$\delta_2$}}
\shortDeduce
\DeduceC{\LR{$!B(t)$}}
\RightLabel{\Intro{\lexists}}
\UnaryInfC{\LR{$\lexists[x][!B(x)]$}}
\AxiomC{\LR{$\Discharge{!B(c)}{x}$}}
\RightLabel{\LR{$\delta_3$}}
\shortDeduce
\DeduceC{\LR{$!D$}}
\DischargeRule{\Elim{\lexists}}{x}
\BinaryInfC{\LR{$!D$}}
\DisplayProof
\quad \longrightarrow\quad
\bottomAlignProof
\AxiomC{}
\RightLabel{\LR{$\delta_2$}}
\shortDeduce
\DeduceC{\LR{$!B(t)$}}
\RightLabel{\LR{$\Subst{\Subst{\delta_3}{t}{c}}{\delta_2}{!B(t)^x}$}}
\shortDeduce
\DeduceC{\LR{$!D$}}
\DisplayProof}
\end{multline*}
\caption{د \Log{N1i} د کمولو ځينې بدلونونه۔}
\ollabel{tab:reduction-conversions}
\end{table}
\olpseganchor{OLP-0676-B011}{end}

\olpseganchor{OLP-0676-B012}{start}
\begin{prob}
د \Intro{\lnot} او ورپسې \Elim{\lnot} د کمولو بدلون ورکړئ۔
\end{prob}
\olpseganchor{OLP-0676-B012}{end}

\olpseganchor{OLP-0676-B013}{start}
يو حالت لا پاتې دے۔ د پرېکون تعريف هغه حالت هم رانغاړي
چې د پرېکون اوږدوالے~\LR{$1$} وي او \LR{$!A$} هم د~\FalseInt{} پايله
او هم د ايستلو استنتاج لويه مقدمه وي۔ په دې حالت
کښې په هغه شان چلند کوو لکه چې \LR{$!A$} د داخلولو
قاعدې پايله وي۔ مثلاً، دا رسم
\[
\bottomAlignProof
\AxiomC{}
\RightLabel{\LR{$\delta_2$}}
\DeduceC{\LR{$\lfalse$}}
\RightLabel{\FalseInt}
\UnaryInfC{\LR{$!B \lif !C$}}
\AxiomC{}
\RightLabel{\LR{$\delta_3$}}
\DeduceC{\LR{$!B$}}
\RightLabel{\Elim{\lif}}
\BinaryInfC{\LR{$!C$}}
\DisplayProof
\quad\text{\RL{په دې}}\quad
\bottomAlignProof
\AxiomC{}
\RightLabel{\LR{$\delta_2$}}
\DeduceC{\LR{$\lfalse$}}
\RightLabel{\FalseInt}
\UnaryInfC{\LR{$!C$}}
\DisplayProof
\]
په ښي اړخ رسم بدلوو۔ په هغو حالتونو کښې لږ احتياط پکار دے چې
\LR{$!A \ident (!B \lor !C)$} يا \LR{$!A \ident \lexists[x][!B(x)]$}
وي۔ مثلاً، که لاندې چپ رسم په ښي رسم بدل کړو،
\[
\bottomAlignProof
\AxiomC{}
\RightLabel{\LR{$\delta_2$}}
\DeduceC{\LR{$\lfalse$}}
\RightLabel{\FalseInt}
\UnaryInfC{\LR{$!B \lor !C$}}
\AxiomC{\LR{$\Discharge{!B}{x}$}}
\RightLabel{\LR{$\delta_3$}}
\DeduceC{\LR{$!D$}}
\AxiomC{\LR{$\Discharge{!C}{y}$}}
\RightLabel{\LR{$\delta_4$}}
\DeduceC{\LR{$!D$}}
\RightLabel{\Elim{\lor}}
\TrinaryInfC{\LR{$!D$}}
\DisplayProof
\quad\text{\RL{په دې}}\quad
\bottomAlignProof
\AxiomC{}
\RightLabel{\LR{$\delta_2$}}
\DeduceC{\LR{$\lfalse$}}
\RightLabel{\FalseInt}
\UnaryInfC{\LR{$!D$}}
\DisplayProof
\]
نو د~\LR{$!D$} په پرېکون فارمول نوی پرېکون جوړېدے شي، او هېڅ
ضمانت نشته چې \LR{$\depth{!D} < \depth{!B \lor !C}$}!{} نو دلته
هم د \Intro{\lor}/\Elim{\lor} د کمولو په شان عمل کوو او
په لاندې يوه بڼه ئې بدلوو:
\[
\bottomAlignProof
\AxiomC{}
\RightLabel{\LR{$\delta_2$}}
\DeduceC{\LR{$\lfalse$}}
\RightLabel{\FalseInt}
\UnaryInfC{\LR{$!B$}}
\RightLabel{\LR{$\delta_3$}}
\DeduceC{\LR{$!D$}}
\DisplayProof
\quad\text{\RL{يا}}\quad
\bottomAlignProof
\AxiomC{}
\RightLabel{\LR{$\delta_2$}}
\DeduceC{\LR{$\lfalse$}}
\RightLabel{\FalseInt}
\UnaryInfC{\LR{$!C$}}
\RightLabel{\LR{$\delta_4$}}
\DeduceC{\LR{$!D$}}
\DisplayProof
\]
\olpseganchor{OLP-0676-B013}{end}

\olpseganchor{OLP-0673-B004}{start}
\olfileid{pt}{nor}{thm}
\olpseganchor{OLP-0673-B004}{end}

\olpseganchor{OLP-0673-B005}{start}
\olsection{د عادي کولو قضیه}
\olpseganchor{OLP-0673-B005}{end}

\olpseganchor{OLP-0673-B006}{start}
ثبوت په \emph{عادي} بڼه کښې دے که د پرېکون
برخې نۀ لري۔ په څرګند ډول دا هغه وخت، او يوازې هغه وخت،
داسې وي چې نۀ پرې د ځاے بدلولو بدلون او نۀ د راکمولو
بدلون تطبيقولے شو۔ د ثبوت \emph{عادي کولو} معنا
دا ده چې د ځاے بدلولو او راکمولو بدلونونه پرې تطبيق
کړو، تر څو هېڅ پرېکون پاتې نۀ شي، او په عادي ثبوت ئې
واړوو۔ دا چې دا کار تل شونے دے، د \emph{عادي کولو قضیه} ده۔
\olpseganchor{OLP-0673-B006}{end}

\olpseganchor{OLP-0673-B007}{start}
\begin{thm}\ollabel{thm:normalization}
هر \Log{N1i}-ثبوت \LR{$\delta$} عادي کېږي، يعنې داسې
ثبوت~\LR{$\delta'$} ته اوړېدے شي چې پرېکونونه نۀ لري۔
\end{thm}
\olpseganchor{OLP-0673-B007}{end}

\olpseganchor{OLP-0673-B008}{start}
\begin{proof}
  د~\LR{$\delta$} د پرېکون پر رتبه او د پرېکون پر اوږدوالي
  استقرا کوو۔ د استقرا فرضيه دا ده چې هر ثبوت چې
  يا د پرېکون ټيټه رتبه لري، يا په هماغه رتبه کښې د
  پرېکون کم اوږدوالے لري، عادي کېږي۔
  دويمه اندازه د پرېکون اوږدوالے دے، يعنې د تر ټولو لوړې
  رتبې د پرېکون برخو د اوږدوالو جمع۔
\olpseganchor{OLP-0673-B008}{end}

\olpseganchor{OLP-0673-B009}{start}
  که د پرېکون اوږدوالے~\LR{$0$} وي، پرېکونونه نۀ شته او
  \LR{$\delta$} له مخکې عادي دے۔ نو فرض کړئ چې د~\LR{$\delta$} د
  پرېکون رتبه او اوږدوالے~\LR{$> 0$} دي۔ د تر ټولو پاسنۍ او
  ښي خوا ته د تر ټولو لوړې رتبې د پرېکون برخه غوره
  کړئ۔ که د هغې اوږدوالے \LR{$>1$} وي، د ځاے بدلولو بدلون
  تطبيق کړئ۔ که اوږدوالے ئې~\LR{$1$} وي، اړوند د راکمولو
  بدلون تطبيق کړئ۔ په دې يو نوے ثبوت جوړېږي چې يا
  د پرېکون هماغه رتبه او کم اوږدوالے لري، يا د پرېکون
  ټيټه رتبه لري؛ نو د استقرا فرضيه پرې تطبيقېږي۔
\end{proof}
\olpseganchor{OLP-0673-B009}{end}

\olpseganchor{OLP-0673-B010}{start}
د پورته ثبوت تګلاره، چې تل د تر ټولو پاسني او ښي
خوا ته د تر ټولو لوړې رتبې پرېکون سره کار کوي، داسې
ځانګړے ترتيب راښيي چې بدلونونه په هغې کښې پر
ثبوت تطبيق شي او عادي ثبوت ورکړي۔ خو د
حيرانتيا خبره دا ده چې په حقيقت کښې ترتيب اهميت نۀ
لري۔ د ځاے بدلولو او راکمولو بدلونونه په هر ترتيب
تطبيقول په پاے کښې يو عادي ثبوت ورکوي۔
\olpseganchor{OLP-0673-B010}{end}

\olpseganchor{OLP-0673-B011}{start}
د برخې او پرېکون تعريف د \Log{N2i} دپاره هم په
آسانۍ جوړېږي، او د \Log{N1i} د ځاے بدلولو او راکمولو
بدلونونه نېغ \Log{N2i} ته ژباړل کېږي۔ نو د عادي کولو
قضیه د \Log{N2i} دپاره هم صحيح ده۔
\olpseganchor{OLP-0673-B011}{end}

\olpseganchor{OLP-0673-B012}{start}
\begin{thm}\ollabel{thm:normalization-N2i}
هر \Log{N2i}-ثبوت \LR{$\delta$} عادي کېږي، يعنې داسې
ثبوت~\LR{$\delta'$} ته اوړېدے شي چې پرېکونونه نۀ لري۔
\end{thm}
\olpseganchor{OLP-0673-B012}{end}

\olpseganchor{OLP-0678-B004}{start}
\olfileid{pt}{nor}{tr}
\olpseganchor{OLP-0678-B004}{end}

\olpseganchor{OLP-0678-B005}{start}
\olsection{د عادي \texorpdfstring{\Log{N2i}}{N2i}-ثبوتونه او \texorpdfstring{\Log{G2i}}{G2i} تر منځ اړول}
\olpseganchor{OLP-0678-B005}{end}

\olpseganchor{OLP-0678-B006}{start}
په \Log{N2i} کښې ثبوتونه په \Log{G2i} کښې
ثبوتونه ته اړولے شو او برعکس هم۔ لومړے، د بې‌پرېکونه
\Log{G2i} ثبوتونه اړول عادي \Log{N2i}-ثبوتونه ورکوي۔
\olpseganchor{OLP-0678-B006}{end}

\olpseganchor{OLP-0678-B007}{start}
د دې نتيجې د بيانولو دپاره لږ اصطلاحي تيارے پکار دے، څو
په~\Log{N2i} او~\Log{G2i} کښې د سيکوېنټونو کيڼ سياقونه په
اسانۍ سره ونښلوو۔ د نښه لرونکو فارمولونه سېټ~\LR{$\Gamma'$}
له بې‌نښو فارمولونه د څوګوني سېټ~\LR{$\Gamma$} سره
\emph{مطابقت لري} که د هر فارمول~\LR{$!A$} دپاره په~\LR{$\Gamma'$}
کښې د \LR{$x : !A$} په بڼه د غړي شمېر په~\LR{$\Gamma$} کښې
د~\LR{$!A$} له تکراري شمېر څخه کم يا ورسره برابر وي۔
\olpseganchor{OLP-0678-B007}{end}

\olpseganchor{OLP-0678-B008}{start}
\begin{cor}
که بې‌پرېکونه \LR{$\Log{G2i} \Proves \Gamma \Sequent \Delta$}
وي، نو د \LR{$\Gamma' \Sequent !C$} يو عادي
\Log{N2i}-ثبوت~\LR{$\delta$} شته، چې دلته د پايلې فارمول
د مخکني حکم د ښي اړخ يوازنے فارمول، او په تش ښي اړخ کښې
د دروغو فارمول دے۔ د نښه لرونکو فارمولونه سېټ~\LR{$\Gamma'$}
له څوګوني سېټ~\LR{$\Gamma$} سره مطابقت لري: يعنې د هر
فارمول~\LR{$!A$} دپاره په~\LR{$\Gamma'$} کښې د \LR{$x : !A$} په بڼه
د غړي شمېر په~\LR{$\Gamma$} کښې د~\LR{$!A$} له تکراري شمېر،
يعنې د~\LR{$!A$} د نقلونو له شمېر، څخه کم يا ورسره برابر وي۔
\end{cor}
\olpseganchor{OLP-0678-B008}{end}

\olpseganchor{OLP-0678-B009}{start}
\begin{proof}
د \olref[nat][tgi]{prop:G2i-to-N2i} د ثبوت په کتنه: د
ثبوتونه په پاے کښې داخلولو قواعد ورزياتوو، او
ايستلو قواعد يوازې د ثبوتونه په پاسنۍ برخه کښې
ورزياتوو۔ نو هېڅکله داخلولو قاعده د
ايستلو قاعدې د پاسه نۀ راولو۔
\end{proof}
\olpseganchor{OLP-0678-B009}{end}

\olpseganchor{OLP-0678-B010}{start}
خو په \olref[nat][tgi]{prop:G2i-to-N2i} کښې د \Cut{} قاعدې
ورکړل شوے اړول په \Log{N2i}-ثبوت~\LR{$\delta$} کښې
پرېکونونه پيدا کولے شي۔ که \Log{G2i} ثبوت په~\Cut{}
پاے ته رسېږي او نږدې فرعي ثبوتونه ئې \LR{$\pi_1$} او \LR{$\pi_2$}
په ترتيب سره \LR{$\Gamma_1 \Sequent !A$} او
\LR{$!A, \Gamma_2 \Sequent \Delta_2$} ثابتوي، نو د~\LR{$\pi_1$}
اړول شوے \LR{$\delta_1$} د~\LR{$\pi_2$} په اړول شوي~\LR{$\delta_2$} کښې
د \LR{$x:!A \Sequent !A$} بديهيتونو پر ځاے پېوندول کېږي۔ که
\LR{$\delta_1$} د داخلولو په قاعده پاے ته رسېږي چې
اصلي فارمول ئې~\LR{$!A$} دے، او په~\LR{$\delta_2$} کښې بديهيت
\LR{$x:!A \Sequent !A$} د ايستلو قاعدې لويه مقدمه
وي، په~\LR{$\delta$} کښې پرېکون پيدا کېږي۔
\olpseganchor{OLP-0678-B010}{end}

\olpseganchor{OLP-0678-B011}{start}
په \Log{N2i} کښې ثبوتونه په~\Log{G2i} کښې
ثبوتونه ته هم اړولے شو۔ د \olref[nat][tni]{prop:N2-to-G2}
په ثبوت کښې د \Elim{\land}، \Elim{\lif}، \Elim{\lnot}
او \Elim{\lforall} اړول په حاصل شوي \Log{G2i} ثبوت
کښې~\Cut{} استنتاجونه راولي۔ خو که هغه \Log{N2i}-ثبوت
چې ترې پيل کوو عادي وي، د دې مخنيوے کېدے شي۔
\olpseganchor{OLP-0678-B011}{end}

\olpseganchor{OLP-0678-B012}{start}
په \Log{N2i}-ثبوت~\LR{$\delta$} کښې د سيکوېنټ د
موردونو لړۍ \LR{$S_1$}، \dots،~\LR{$S_n$} \emph{څانګه} بولو که
\LR{$S_1$} بديهيت وي، او هر وخت چې \LR{$S_i$} د استنتاج لويه مقدمه
وي، \LR{$S_{i+1}$} ئې پايله وي۔ په بله وينا، څانګه په~\LR{$\delta$}
کښې له بديهيتونو سيکوېنټونه لاندې تعقيبوي، خو د
ايستلو قواعدو پر کوچنيو مقدمو درېږي۔ د~\LR{$\delta$}
\emph{اصلي څانګه} هغه څانګه ده چې \LR{$S_n$} ئې پاييز سيکوېنټ
وي۔ هر ثبوت~\LR{$\delta$} لږ تر لږه يوه اصلي څانګه لري،
ځکه هره قاعده لږ تر لږه يوه لويه مقدمه لري؛ يوازې
\Intro{\land} دوه لويې مقدمې لري۔
\olpseganchor{OLP-0678-B012}{end}

\olpseganchor{OLP-0678-B013}{start}
\begin{prop}\ollabel{prop:normal-N2i-to-G2i}
که \LR{$\delta$} د~\LR{$\Gamma \Sequent !A$} عادي \LR{$\Log{N2c}$} يا
\LR{$\Log{N2i}$}-ثبوت وي، نو د \LR{$\Gamma' \Sequent \Delta$}
بې‌پرېکونه \LR{$\Log{G2c}$}-ثبوت، يا په ترتيب سره
\LR{$\Log{G2i}$}-ثبوت~\LR{$\pi$}، شته۔ دلته \LR{$\Gamma'$} هغه
د فارمولونه څوګونے سېټ دے چې له \LR{$\Gamma$} څخه د نښو په
لرې کولو حاصلېږي؛ که \LR{$!A$}~د~\LR{$\lfalse$} فارمول نۀ وي،
\LR{$\Delta = \{!A\}$}، او که وي، \LR{$\Delta = \emptyset$} دے۔
لاندې د اصلي څانګې استدلال حدسي حالت ثابتوي۔ په کلاسيکي
نظام کښې د دروغو قاعده فرضونه هم باسي، نو د فرضونو د
ساتلو استدلال په هغه حالت کښې جلا ثبوت غواړي۔
\end{prop}
\olpseganchor{OLP-0678-B013}{end}

\olpseganchor{OLP-0678-B014}{start}
\begin{proof}
دا ځل استقرا د \LR{$\Gamma \Sequent !A$} په ثبوت~\LR{$\delta$}
کښې د استنتاجونو پر شمېر کوو۔ که په~\LR{$\delta$} کښې استنتاج
نۀ وي، دا بديهيت \LR{$x:!A \Sequent !A$} دے، او \LR{$\pi$} اړوند
بديهيت \LR{$!A \Sequent !A$} دے، يا که \LR{$!A \ident \lfalse$}
وي، \LR{$\lfalse \Sequent \ $} دے۔
\olpseganchor{OLP-0678-B014}{end}

\olpseganchor{OLP-0678-B015}{start}
که \LR{$\delta$} په استنتاجي قاعده پاے ته رسېږي، حالتونه بېلوو:
\begin{enumerate}
  \item که \LR{$R$} د داخلولو قاعده يا~\FalseInt وي،
  اړول کټ مټ د \olref[nat][tni]{prop:N2-to-G2} د ثبوت په
  شان کېږي او~\Cut ته اړتيا نۀ لري۔
\olpseganchor{OLP-0678-B015}{end}

\olpseganchor{OLP-0678-B016}{start}
  که د~\LR{$\delta$} وروستۍ قاعده~\LR{$R$} د ايستلو
  قاعده وي، لاندې خبرې وګورئ:
  \begin{enumerate}
    \item د~\LR{$\delta$} له اصلي څانګې هېڅ داخلولو
    قاعده نۀ تېرېږي؛ کنه د اصلي څانګې په اوږدو کښې به د
    داخلولو قاعدې يا~\FalseInt{} ورپسې
    ايستلو قاعده راتله، او \LR{$\delta$} به عادي نۀ و۔
    \item ځکه د~\LR{$\delta$} په اصلي څانګو کښې
    داخلولو قاعدې نشته، يوازې يوه اصلي څانګه ده۔
    يوازې \Intro{\land} دوه لويې مقدمې لري، نو څو اصلي
    څانګې به د \Intro{\land} قاعدې له مقدمو تېرېدې۔
    \item که په اصلي څانګه کښې د \Elim{\lor} يا
    \Elim{\lexists} لويه مقدمه وي، داسې قاعده يوازې يوه ده
    او همدا وروستۍ، يعنې~\LR{$R$} ده۔ کنه په اصلي څانګه کښې به
    داسې \Elim{\lor} يا \Elim{\lexists} قاعده وه چې پايله
    ئې د ايستلو قاعدې لويه مقدمه وي؛ نو ثبوت به
    عادي نۀ و، ځکه د ځاے بدلولو بدلون به تطبيقېده۔
    \item له \Elim{\lor} او \Elim{\lexists} پرته نور
    ايستلو قواعد فرضونه نۀ باسي، يعنې د سيکوېنټ
    کيڼ سياق نۀ کموي۔ \Elim{\lor} او \Elim{\lexists}
    يوازې د کوچنۍ مقدمې د پاسه فرضونه باسي، يعنې يوازې د
    خپلو کوچنيو مقدمو کيڼ سياقونه کموي۔ نو که د~\LR{$\delta$}
    په اصلي څانګه کښې سيکوېنټ \LR{$S_i$} سياق~\LR{$\Gamma_1$} ولري،
    د هغه استنتاج پايله~\LR{$S_{i+1}$} چې \LR{$S_i$} ئې لويه مقدمه
    ده، سياق \LR{$\Gamma_1' \supseteq \Gamma_1$} لري۔
    \item له دې امله، که \LR{$x: !C \Sequent !C$} د اصلي
    څانګې تر ټولو پاسنے سيکوېنټ~\LR{$S_1$} وي، نو
    \LR{$x:!C \in \Gamma$} دے۔
  \end{enumerate}
  اوس د~\LR{$!C$} د بڼې له مخې حالتونه بېلوو۔ ځکه په اصلي
  څانګه کښې هر \LR{$S_i$} د ايستلو قاعدې لويه مقدمه
  ده، دا خبره پر \LR{$x:!C \Sequent !C$} هم تطبيقېږي۔
\olpseganchor{OLP-0678-B016}{end}

\olpseganchor{OLP-0678-B017}{start}
  \item \LR{$!C \ident !D \land !E$}۔ بيا \LR{$\delta$} دا بڼه لري:
\[
  \Axiom$x:!D \land !E \fCenter !D \land !E$
  \RightLabel{\Elim{\land}[1]}
  \UnaryInf$x:!D \land !E \fCenter !D$
  \Deduce$x:!D \land !E, \Gamma_1 \fCenter !A$
  \DisplayProof
  \]
  په اصلي څانګه کښې \LR{$x:!D \land !E \Sequent !D$} راځي۔
  هلته د \Intro{\lif} يا \Intro{\lnot} لويه مقدمه، يا د
  \Elim{\lor} يا \Elim{\lexists} کوچنۍ مقدمه نشته؛ نو
  د \LR{$x:!D \land !E$} د سياق فارمولونه د اصلي څانګې په
  اوږدو کښې پاتې کېږي، که څۀ هم له کوچنيو څانګو نور فرضونه
  ورزياتېدې شي۔ له همدې امله پاييز سياق~\LR{$x:!D \land !E$}
  لري۔ سرچينه وړانديز کوي چې ثبوت \LR{$\delta$} په
  ثبوت~\LR{$\delta_1$} داسې بدل کړو: په \Elim{\land}
  پاے ته رسېدونکے فرعي ثبوت په \LR{$x:!D \Sequent !D$}
  بدل کړو او د اصلي څانګې د هر سيکوېنټ په سياق کښې
  \LR{$x:!D \land !E$} په \LR{$x:!D$} بدل کړو، يعنې دا بدلون:
\[
  \bottomAlignProof
  \Axiom$x:!D \land !E \fCenter !D \land !E$
  \RightLabel{\Elim{\land}[1]}
  \UnaryInf$x:!D \land !E \fCenter !D$
  \Deduce$x:!D \land !E, \Gamma_1 \fCenter !A$
  \DisplayProof
  \quad\text{\RL{په دې}}\quad
  \bottomAlignProof
  \Axiom$x:!D \fCenter !D$
  \Deduce$x:!D, \Gamma_1 \fCenter !A$
  \DisplayProof
  \]
  که وړانديز شوے بدلون صحيح فرعي ثبوت راکړي، د استقرا
  فرضيه پر \LR{$\delta_1$} تطبيقېږي، ځکه په \LR{$\delta$}
  کښې د استنتاجونو شمېر په \LR{$\delta_1$} کښې له شمېر څخه
  په~\LR{$1$} زيات دے۔
  د دې بدلولو پر وخت د کوچنيو څانګو پرانيستي فرضونه هم په
  پام کښې نيسو۔ که پخوانے نښه لرونکے فرض هلته هم کارېدلے
  وي، د هغه حذف د کوچنۍ څانګې يوه لازمه مقدمه له منځه وړلے
  شي۔ د نوې مقدمې د نښې او د پخواني فرض د نښې ټکر هم بايد
  حل شي۔ په داسې حالت کښې پورته رسم د بدلون د بشپړ ثبوت
  دپاره بس نۀ دے۔
\olpseganchor{OLP-0678-B017}{end}

\olpseganchor{OLP-0678-B018}{start}
  د استقرا د فرضيې له مخې، د \LR{$!D, \Gamma_1' \Sequent !A$}
  يو \Log{G2i}-ثبوت \LR{$\pi_1$} شته۔
  د \LeftR{\land} په تطبيق ثبوت~\LR{$\pi$} اخلو:
\[
  \AxiomC{}
  \RightLabel{\LR{$\pi_1$}}
  \Deduce$!D, \Gamma_1' \fCenter !A$
  \RightLabel{\LeftR{\land}}
  \UnaryInf$!D \land !E, \Gamma_1' \fCenter !A$
  \DisplayProof
  \]
  که \LR{$!A \ident \lfalse$} وي، \RightR{\Weakening}
  ورزياتوو، مثلاً:
\[
  \AxiomC{}
  \RightLabel{\LR{$\pi_1$}}
  \Deduce$!D, \Gamma_1' \fCenter $
  \RightLabel{\LeftR{\land}}
  \UnaryInf$!D \land !E, \Gamma_1' \fCenter $
  \RightLabel{\RightR{\Weakening}}
  \UnaryInf$!D \land !E, \Gamma_1' \fCenter !A$
  \DisplayProof
  \]
  \item که \LR{$!C \ident !D \lor !E$} وي، ضرور د
  \Elim{\lor} لويه مقدمه ده او دا استنتاج په اصلي څانګه
  کښې وروستے استنتاج~\LR{$R$} دے۔ يعنې \LR{$\delta$} دا بڼه لري:
\[
  \Axiom$x:!D \lor !E \fCenter !D \lor !E$
  \AxiomC{}
  \RightLabel{\LR{$\delta_1$}}
  \Deduce$y:!D, \Gamma_1 \fCenter!A$
  \AxiomC{}
  \RightLabel{\LR{$\delta_2$}}
  \Deduce$z:!E, \Gamma_2 \fCenter !A$
  \RightLabel{\Elim{\lor}}
  \TrinaryInf$x:!D \lor !E, \Gamma_1, \Gamma_2 \fCenter !A$
  \DisplayProof
  \]
  د استقرا فرضيه پر ثبوتونه \LR{$\delta_1$} او \LR{$\delta_2$}
  تطبيقېږي؛ د \LR{$!D, \Gamma_1' \Sequent !A$} او
  \LR{$!E, \Gamma_2' \Sequent !A$} دپاره \Log{G2i}-ثبوتونه
  \LR{$\pi_1$} او \LR{$\pi_2$} اخلو۔ د~\LR{$\pi$} د ترلاسه کولو دپاره
  \LR{$\LeftR{\lor}$} تطبيقوو:
\[
  \AxiomC{}
  \RightLabel{\LR{$\pi_1$}}
  \Deduce$!D, \Gamma_1' \fCenter !A$
  \AxiomC{}
  \RightLabel{\LR{$\pi_2$}}
  \Deduce$!E, \Gamma_2' \fCenter !A$
  \RightLabel{\LeftR{\lor}}
  \BinaryInf$!D \lor !E, \Gamma_1', \Gamma_2' \fCenter !A$
  \DisplayProof
  \]
  که \LR{$\Elim{\lor}$} فرض \LR{$y: !D$} يا \LR{$z: !E$} نۀ باسي، يعنې
  د کوچنيو مقدمو په سياق کښې دا فرضونه نۀ وي، يا که~\LR{$!A$}
  د~\LR{$\lfalse$} فارمول وي، بايد ځينې \LeftR{\Weakening}
  او \RightR{\Weakening} ورزياتې کړو، مثلاً:
\[
  \AxiomC{}
  \RightLabel{\LR{$\pi_1$}}
  \Deduce$\Gamma_1' \fCenter $
  \RightLabel{\LeftR{\Weakening}}
  \UnaryInf$!D, \Gamma_1' \fCenter $
  \AxiomC{}
  \RightLabel{\LR{$\pi_2$}}
  \Deduce$\Gamma_2' \fCenter $
  \RightLabel{\LeftR{\Weakening}}
  \UnaryInf$!E, \Gamma_2' \fCenter $
  \RightLabel{\LeftR{\lor}}
  \BinaryInf$!D \lor !E, \Gamma_1', \Gamma_2' \fCenter $
  \RightLabel{\RightR{\Weakening}}
  \UnaryInf$!D \lor !E, \Gamma_1', \Gamma_2' \fCenter !A$
  \DisplayProof
  \]
\olpseganchor{OLP-0678-B018}{end}

\olpseganchor{OLP-0678-B019}{start}
  \item \LR{$!C \ident !D \lif !E$}۔ بيا \LR{$\delta$} دا بڼه لري:
\[
  \Axiom$x:!D \lif !E \fCenter !D \lif !E$
  \AxiomC{}
  \RightLabel{\LR{$\delta_1$}}
  \Deduce$\Gamma_1 \fCenter!D$
  \RightLabel{\Elim{\lif}}
  \BinaryInf$x:!D \lif !E, \Gamma_1 \fCenter !E$
  \RightLabel{\LR{$\delta_2$}}
  \Deduce$x:!D \lif !E, \Gamma_1, \Gamma_2 \fCenter !A$
  \DisplayProof
  \]
  پام وکړئ چې په اصلي څانګه کښې \LR{$x:!D \lif !E,
  \Gamma_1 \Sequent !E$} راځي؛ د \Intro{\lif} يا
  \Intro{\lnot} لويه مقدمه، يا د \Elim{\lor} يا
  \Elim{\lexists} کوچنۍ مقدمه نۀ لري۔ نو د
  \LR{$x:!D \lif !E, \Gamma_1$} د سياق فارمولونه د اصلي
  څانګې په اوږدو کښې پاتې کېږي؛ نور فرضونه ورزياتېدې شي۔
  له دې امله پاييز سياق~\LR{$x:!D \lif !E, \Gamma_1$} لري۔
  سرچينه وړانديز کوي چې د ثبوت ټوټه \LR{$\delta_2$} په
  صحيح ثبوت~\LR{$\delta_2'$} داسې بدله شي: په
  \Elim{\lif} پاے ته رسېدونکے فرعي ثبوت په
  \LR{$x:!E \Sequent !E$} بدل شي او د اصلي څانګې د هر سيکوېنټ
  په سياق کښې \LR{$x:!D \lif !E, \Gamma_1$} په \LR{$x:!E$}
  بدل شي، يعنې:
\[
  \bottomAlignProof
  \Axiom$x:!D \lif !E, \Gamma_1 \fCenter !E$
  \RightLabel{\LR{$\delta_2$}}
  \Deduce$x:!D \lif !E, \Gamma_1, \Gamma_2 \fCenter !A$
  \DisplayProof
  \quad\text{\RL{په دې}}\quad
  \bottomAlignProof
  \Axiom$x:!E \fCenter !E$
  \RightLabel{\LR{$\delta_2'$}}
  \Deduce$x:!E, \Gamma_2 \fCenter !A$
  \DisplayProof
  \]
  که وړانديز شوے بدلون صحيح فرعي ثبوت راکړي، د استقرا
  فرضيه پر \LR{$\delta_1$} او \LR{$\delta_2'$} تطبيقېږي،
  ځکه په \LR{$\delta$} کښې د استنتاجونو شمېر په \LR{$\delta_1$} او
  ~\LR{$\delta_2$} کښې د استنتاجونو د شمېرونو مجموعه ده، او
  ~\LR{$1$} د اصلي څانګې په پيل کښې د \Elim{\lif} استنتاج
  دپاره ورسره زياتېږي۔ که همدغه استنتاج په \LR{$\delta$} کښې
  وروستے استنتاج~\LR{$R$} هم وي، نو \LR{$\delta_2$} د \LR{$0$}
  استنتاجونو ټوټه ده۔
  د کوچنۍ مقدمې د استنتاجونو شمېر له دې څخه نۀ ټاکل کېږي۔
  د کوچنيو څانګو د پرانيستو فرضونو په اړه د مخکيني حالت
  هماغه شرط دلته هم پکار دے۔
\olpseganchor{OLP-0678-B019}{end}

\olpseganchor{OLP-0678-B020}{start}
  د استقرا د فرضيې له مخې، د \LR{$\Gamma_1' \Sequent !D$}
  او \LR{$!E, \Gamma_2' \Sequent !A$} دپاره په ترتيب سره
  \Log{G2i}-ثبوتونه \LR{$\pi_1$} او \LR{$\pi_2$} شته۔
  د \LeftR{\lif} په تطبيق ثبوت~\LR{$\pi$} اخلو:
\[
  \AxiomC{}
  \RightLabel{\LR{$\pi_1$}}
  \Deduce$\Gamma_1' \fCenter !D$
  \AxiomC{}
  \RightLabel{\LR{$\pi_2$}}
  \Deduce$!E, \Gamma_2' \fCenter !A$
  \RightLabel{\LeftR{\lif}}
  \BinaryInf$!D \lif !E, \Gamma_1', \Gamma_2' \fCenter !A$
  \DisplayProof
  \]
  که \LR{$!D$} د دروغو فارمول وي، او/يا \LR{$!A \ident \lfalse$}
  وي، يو يا دوه \RightR{\Weakening} ورزياتوو، مثلاً:
\[
  \AxiomC{}
  \RightLabel{\LR{$\pi_1$}}
  \Deduce$\Gamma_1' \fCenter $
  \RightLabel{\RightR{\Weakening}}
  \UnaryInf$\Gamma_1' \fCenter !D$
  \AxiomC{}
  \RightLabel{\LR{$\pi_2$}}
  \Deduce$!E, \Gamma_2' \fCenter $
  \RightLabel{\RightR{\Weakening}}
  \UnaryInf$!E, \Gamma_2' \fCenter !A$
  \RightLabel{\LeftR{\lif}}
  \BinaryInf$!D \lif !E, \Gamma_1', \Gamma_2' \fCenter !A$
  \DisplayProof
  \]
\olpseganchor{OLP-0678-B020}{end}

\olpseganchor{OLP-0678-B021}{start}
\end{enumerate}
پاتې حالتونه په ورته طريقه څېړل کېږي او د مشقونو په توګه
پرېښودل شوي دي۔
\end{proof}
\olpseganchor{OLP-0678-B021}{end}

\olpseganchor{OLP-0678-B022}{start}
\begin{cor}
په عادي \Log{N1i} يا \Log{N2i}-ثبوت کښې هر راغلے
فارمول د پاييز فارمول يا د پاتې پرانيستو فرضونو
فرعي‌فارمول، يا د پاييز سيکوېنټ اړوند فرعي فارمول دے؛
په لومړي ترتيب کښې اړوند کمي تعويضي موردونه هم مراد دي۔
\end{cor}
\olpseganchor{OLP-0678-B022}{end}

\olpseganchor{OLP-0678-B023}{start}
\begin{proof}
د \olref{prop:normal-N2i-to-G2i} له مخې، هر عادي
\Log{N2i}-ثبوت په \Log{G2i}-ثبوت اړولے شو۔ د
ثبوت کتنه ښيي چې په \Log{N2i}-ثبوت کښې هر راغلے
فارمول، د دروغو د پايلې له استثنا سره، په
\Log{G2i}-ثبوت کښې هم راځي؛ د دروغو پايله دلته په
تش ښي اړخ ښودل کېږي۔ \Log{G2i} د \Cut{} قاعده نۀ لري،
نو د اړوند فرعي‌فارمول خاصيت لري۔ د اصلي جوړونې
پورته څرګند شوي پاتې حالتونه د بشپړې کتنې ځاے نۀ نيسي۔
\end{proof}
\olpseganchor{OLP-0678-B023}{end}

\olpseganchor{OLP-0690-B004}{start}
\olchapter{int}{pty}{بيانونه د ټايپونو په توګه}
\olpseganchor{OLP-0690-B004}{end}

\olpseganchor{OLP-0690-B005}{start}
\begin{editorial}
  دا د کري--هاورډ د مطابقت په باب د يوه فصل
  \emph{ډېره تجربوي} مسوده ده۔ نور وضاحت او د
  موضوع د اړتيا بيان غواړي، او ښايي تېروتنې او
  نيمګړتياوې ولري۔ د عادي کولو ثبوت بياکتنه او
  پراخول غواړي۔ د ضرب ټايپ بېلګه د عطف د برخو د بدلولو په ثبوت کښې راځي۔
  د پروجکشن او د يا د ايستلو ترتيب د بدلېدو يوه بېلګه هم شته،
  خو د دغو بدلونونو بشپړه کتنه او د ساده کولو
  بدلونونه نۀ دي رانغښتل شوي۔ کله چې د ټايپ‌لرونکي
  لامبډا حساب مواد هم ورسره وي، ډېر به ئې مطلب
  څرګند شي، ځکه دلته اساساً د هغه پوهه له مخه
  فرض شوې ده۔ په ډېر احتياط ئې وکاروئ۔
\end{editorial}
\olpseganchor{OLP-0690-B005}{end}

\olpseganchor{OLP-0690-B006}{start}

%







\olpseganchor{OLP-0690-B006}{end}

\olpseganchor{OLP-0686-B004}{start}
\olfileid{int}{pty}{int}
\olpseganchor{OLP-0686-B004}{end}

\olpseganchor{OLP-0686-B005}{start}
\olsection{پېژندنه}
\olpseganchor{OLP-0686-B005}{end}

\olpseganchor{OLP-0686-B006}{start}
لامبډا حساب د محاسبې يو ساده ماډل دے۔ دا له متغيرونو جوړو
ترمونو سره کار کوي۔ که \LR{$N$} او \LR{$M$} ترمونه وي، نو \LR{$(NM)$} هم
ترم دے («په~\LR{$M$} د \LR{$N$} تطبيق»)۔ که \LR{$N$} ترم وي، نو
\LR{$\lambd[x][N]$} هم ترم دے۔ که \LR{$N$} متغير~\LR{$x$} ولري، نو په
\LR{$\lambd[x][N]$} کښې د~\LR{$x$} اړوند آزاد موردونه د \LR{$\lambd[x]$}
عملګر تړي، او بيا آزاد نۀ پاتې کېږي۔ د \LR{$(\lambd[x][N]M)$}
په بڼه ترم \LR{$\Subst{N}{M}{x}$} ته \LR{$\beta$}-انقباض کوي
(په~\LR{$N$} کښې د \LR{$x$} د هر آزاد مورد پر ځاے د \LR{$M$} د اېښودلو
پايله؛ د آزاد متغير د نيولو د مخنيوي دپاره تړلي نومونه د اړتيا
پر وخت مخکې بيا نومول کېږي)۔ ترم \LR{$N$} په يوه ګام کښې~\LR{$N'$}
ته \LR{$\beta$}-بدلون کوي که \LR{$N'$} په \LR{$N$} کښې د يوه فرعي ترم
په \LR{$\beta$}-انقباض تر لاسه شي؛ \LR{$N \redone N'$} ليکو۔ په
\LR{$\lambd$}-حساب کښې محاسبه د \LR{$N \redone N_1
\redone N_2 \redone \dots$} په بڼه لړۍ ده۔ که دغه لړۍ په
داسې ترم \LR{$N_k$} پاے ته ورسي چې هېڅ فرعي ترم ئې
\LR{$\beta$}-انقباض نۀ شي کولے، نو دغه ترم \emph{عادي} يا
د~\LR{$N$} عادي بڼه بلل کېږي۔ په \LR{$\lambd$}-حساب کښې محاسبه د
همداسې لړيو په توګه وګڼئ۔ محاسبه هغه وخت درېږي چې عادي بڼې
ته ورسي، او همدغه عادي بڼه د محاسبې پايله ده۔ د محاسبې دغه
ساده ماډل د ټيورينګ له پلوه بشپړ دے: هره جزوي بازګشتي تابع په
\LR{$\lambd$}-حساب کښې د يوه ترم په وسيله محاسبه کېدے شي۔
\olpseganchor{OLP-0686-B006}{end}

\olpseganchor{OLP-0686-B007}{start}
هاسکل کَري او ويليم هاورډ په خپلواک ډول د فطري استنتاج او
\LR{$\lambd$}-حساب تر منځ نږدې ورته‌والے وموند۔ په ابتدايي معنا،
ترم \LR{$\lambd[x][N]$} يوه تابع ده: \LR{$x$} ئې آرګومېنټ دے؛ \LR{$N$}
ښيي چې د ورکړل شوي~\LR{$x$} دپاره ارزښت څنګه محاسبه کړو۔ نو
ترم \LR{$(\lambd[x][N]M)$} په آرګومېنټ~\LR{$M$} د تابعې
\LR{$\lambd[x][N]$} تطبيق دے؛ \LR{$\beta$}-انقباض ئې د ارزښت د
موندلو طريقه ښيي: په \LR{$N$} کښې د \LR{$x$} د آزادو موردونو پر ځاے
\LR{$M$} کېږدئ (او د پايلې ترم ارزونه جاري وساتئ)۔ اوس دا د
ټاکلي داخلې او وتلو سېټ لرونکې تابعې وګڼئ۔ که د
\LR{$\lambd[x][N]$} د داخلې سېټ \LR{$A$} او د وتلو سېټ~\LR{$B$} وي،
نو متغير~\LR{$x$} يوازې له~\LR{$A$} څخه آرګومېنټونه اخلي او \LR{$N$}
له~\LR{$B$} څخه ارزښتونه ورکوي۔ له دې کبله، \LR{$(\lambd[x][N]M)$}
يوازې هغه وخت معنا لري چې \LR{$\lambd[x][N]$} د \LR{$A \to B$}
تابع او \LR{$M \in A$} وي۔ په \LR{$\lambd$}-حساب کښې د دغو
معلوماتو په ورزياتولو \emph{ټايپ‌لرونکے} \LR{$\lambd$}-حساب
تر لاسه کېږي۔ په ټايپ‌لرونکي \LR{$\lambd$}-حساب کښې هره څرګندونه
\LR{$(\lambd[x][N]M)$} روا نۀ ده؛ يوازې هغه روا ده چې د
\LR{$\lambd[x][N]$} او \LR{$M$} «ټايپونه» سره سمون خوري، يعنې
\LR{$\lambd[x][N]$} ټايپ \LR{$A \lif B$} او \LR{$M$} ټايپ~\LR{$A$} ولري۔
د \LR{$\lambd[x][N]$} ټايپ يوازې هغه وخت \LR{$A \to B$} وي چې
د \LR{$x$} ټايپ \LR{$A$} او د \LR{$N$} ټايپ~\LR{$B$} وي۔ په عمومي ډول، هره
څرګندونه \LR{$(M_1M_2)$} روا نۀ ده۔ يوازې هغه ترمونه \LR{$(M_1M_2)$}
روا دي چې \LR{$M_1$} ټايپ \LR{$A \to B$} او \LR{$M_2$} ټايپ~\LR{$A$}
ولري؛ په دغه صورت کښې د \LR{$(M_1M_2)$} ټايپ~\LR{$B$} دے۔
\olpseganchor{OLP-0686-B007}{end}

\olpseganchor{OLP-0686-B008}{start}
د ټايپ ټاکلو دغه قواعد اوس په څرګند ډول د فطري استنتاج له
اشتقاقونه سره سمون لري: ټايپونه د فارمولونه په وسيله
ښودل کېږي، او خپله اشتقاقونه له \LR{$\lambd$}-ترمونو سره
سمون لري۔ د بېلګې په توګه، د \LR{$!A \lif !B$} اشتقاق
له داسې ترم \LR{$\lambd[\typeof{x}{!A}][N]$} سره سمون لري چې
د تابعې ټايپ ئې \LR{$!A \lif !B$} دے۔ د فطري استنتاج استنتاجي
قواعد د ټايپ‌لرونکي لامبډا حساب د ټايپ ټاکلو له قواعدو سره
سمون لري؛ د بېلګې په توګه،
\begin{prooftree}
  \AxiomC{\LR{$\Discharge{!A}{x}$}}
  \DeduceC{\LR{$!B$}}
  \DischargeRule{\Intro\lif}{x}
  \UnaryInfC{\LR{$!A \lif !B$}}
  \DisplayProof\bottomAlignProof
  \quad\text{\RL{سره سمون لري}}\quad
  \Axiom$x: !A \fCenter N: !B$
  \UnaryInf$ \fCenter \lambd[\typeof{x}{!A}][N] : !A \lif !B$
\end{prooftree}
د ښي اړخ قاعده وايي چې که \LR{$x$} ټايپ~\LR{$!A$} او \LR{$N$}
ټايپ~\LR{$!B$} ولري، نو \LR{$\lambd[\typeof{x}{!A}][N]$}
ټايپ~\LR{$!A \lif !B$} لري۔
\olpseganchor{OLP-0686-B008}{end}

\olpseganchor{OLP-0686-B009}{start}
د \LR{$\Elim\lif$} قاعده د تطبيق د ترمونو د ټايپ ټاکلو له
قاعدې سره سمون لري، يعنې
\[
  \AxiomC{\LR{$!A \lif !B$}}
  \AxiomC{\LR{$!A$}}
  \RightLabel{\Elim\lif}
  \BinaryInfC{\LR{$!B$}}
  \DisplayProof
  \quad\text{\RL{سره سمون لري}}\quad
  \Axiom$\fCenter M_1: !A \lif !B$
  \Axiom$\fCenter M_2: !A$
  \BinaryInf$ \fCenter (M_1M_2) : !B$
  \DisplayProof
\]
که د \Intro\lif{} قاعدې سمدستي ورپسې د \Elim\lif{}
قاعده راشي، نو اشتقاق ساده کېدے شي:
\begin{gather*}
  \AxiomC{\LR{$\Discharge{!A}{x}$}}
  \DeduceC{\LR{$!B$}}
  \DischargeRule{\Intro\lif}{x}
  \UnaryInfC{\LR{$!A \lif !B$}}
  \AxiomC{}
  \DeduceC{\LR{$!A$}}
  \RightLabel{\Elim\lif}
  \BinaryInfC{\LR{$!B$}}
  \DisplayProof
  \quad\redone\quad
  \AxiomC{}
  \DeduceC{\LR{$!A$}}
  \DeduceC{\LR{$!B$}}
  \DisplayProof
\end{gather*}
دا د لامبډا ترمونو له \LR{$\beta$}-انقباض سره سمون لري:
\[
(\lambd[\typeof{x}{!A}][N]M) \quad\redone\quad
\Subst{N}{M}{x}.
\]
\olpseganchor{OLP-0686-B009}{end}

\olpseganchor{OLP-0686-B010}{start}
د \LR{$\land$} د قواعدو او د «ضرب» د ټايپونو تر منځ، او د
\LR{$\lor$} د قواعدو او د «جمع» د ټايپونو تر منځ هم ورته سمونونه شته۔
\olpseganchor{OLP-0686-B010}{end}

\olpseganchor{OLP-0686-B011}{start}
په ساده ټايپ‌لرونکي لامبډا حساب کښې د ترمونو او د فطري
استنتاج د اشتقاقونه دغه سمون «کَري--هاورډ»، «ثبوتونه د
پروګرامونو په توګه» يا «بيانونه د ټايپونو په توګه» سمون بلل
کېږي۔ د فارمولونه (بيانونو) له ټايپونو او د ثبوتونو له
ترمونو سره د سمون تر څنګ، سمونونه داسې لنډ ښودلے شو:
\begin{center}
  \begin{tabular}{c c c}
    منطق & پروګرام \\
    \hline
    بيان/فارمول & ټايپ \\
    ثبوت & ترم \\
    فرض & متغير \\
    ايستل شوے فرض & تړلے متغير \\
    پاتې پرانيستے فرض & آزاد متغير \\
    \LR{$!A \lif !B$} & د تابعې ټايپ \\
    \LR{$!A \land !B$} & د ضرب ټايپ \\
    \LR{$!A \lor !B$} & د جمع ټايپ \\
    \LR{$\lfalse$} & تش ټايپ \\
    \hline
  \end{tabular}
\end{center}
\olpseganchor{OLP-0686-B011}{end}

\olpseganchor{OLP-0686-B012}{start}
د کَري--هاورډ سمون د اتومات ثبوت مرستيالانو او د پروګرامونو د
ټايپ کتونکو يو بنسټ دے، ځکه د يوه بيان شاهد ثبوت کتل (لکه
پورته مو چې وکړل) د دې کتلو سره برابر دي چې پروګرام (ترم)
خپل اعلان شوے ټايپ لري که نه۔
\olpseganchor{OLP-0686-B012}{end}

\olpseganchor{OLP-0694-B004}{start}
\olfileid{int}{pty}{snd}
\olpseganchor{OLP-0694-B004}{end}

\olpseganchor{OLP-0694-B005}{start}
\olsection{سېکوېنټي فطري استنتاج}
\olpseganchor{OLP-0694-B005}{end}

\olpseganchor{OLP-0694-B006}{start}
که داسې فطري استنتاجي اشتقاق وي چې
نۀ ايستل شوي فرضونه ئې \LR{$\Gamma$} او پايله ئې
\LR{$!A$} وي، نو \LR{$\Gamma \Sequent !A$} ليکو؛ که
\LR{$\Gamma$} تش وي، \LR{$\Sequent !A$} ليکو۔
\olpseganchor{OLP-0694-B006}{end}

\olpseganchor{OLP-0694-B007}{start}
د \LR{$\Gamma \cup \{!A_1, \dots, !A_n\}$} دپاره
\LR{$\Gamma, !A_1, \dots, !A_n$}، او د
\LR{$\Gamma \cup \Delta$} دپاره \LR{$\Gamma, \Delta$} ليکو۔
\olpseganchor{OLP-0694-B007}{end}

\olpseganchor{OLP-0694-B008}{start}
پام وکړئ: که \LR{$\Gamma \Sequent !A \land !A$}
ولرو، معنا ئې دا ده چې اشتقاق شته چې
نۀ ايستل شوي فرضونه ئې \LR{$\Gamma$} او وروستے
فارمول ئې \LR{$!A \land !A$} دے۔ په پاې کښې د
\LR{$\Elim{\land}$} په کارولو سره اشتقاق
اخيستلے شو چې نۀ ايستل شوي فرضونه ئې عين
وي او پايله ئې \LR{$!A$} وي۔ په بله وينا، که
\LR{$\Gamma \Sequent !A \land !B$} وي، نو
\LR{$\Gamma \Sequent !A$} هم وي۔
\begin{prooftree}
  \Axiom$\Gamma \fCenter !A \land !B$
  \RightLabel{\LR{$\Elim{\land}$}}
  \UnaryInf$\Gamma \fCenter !A$
  \DisplayProof\qquad\bottomAlignProof
  \Axiom$\Gamma \fCenter !A \land !B$
  \RightLabel{\LR{$\Elim{\land}$}}
  \UnaryInf$\Gamma \fCenter !B$
\end{prooftree}
د \LR{$\Elim{\land}$} نښه په فطري استنتاج کښې د
همدغه نوم له قاعدې سره تړاو ښيي۔
\olpseganchor{OLP-0694-B008}{end}

\olpseganchor{OLP-0694-B009}{start}
په همدې ډول، فرض کړئ \LR{$\Gamma, !A \Sequent !B$}
لرو؛ يعنې اشتقاق شته چې
نۀ ايستل شوي فرضونه ئې \LR{$\Gamma, !A$} او وروستے
فارمول ئې~\LR{$!B$} دے۔ که د \LR{$\Intro{\lif}$} قاعده
وکاروو، اشتقاق ترلاسه کېږي چې
نۀ ايستل شوي فرضونه ئې \LR{$\Gamma$} او وروستے
فارمول ئې \LR{$!A \lif !B$} وي؛ يعنې
\LR{$\Gamma \Sequent !A \lif !B$}۔ پام وکړئ چې د
فرضونو ايستل ئې څومره څرګند کړل۔
\begin{prooftree}
  \Axiom $\Gamma, !A \fCenter !B$
  \RightLabel{\LR{$\Intro{\lif}$}}
  \UnaryInf $\Gamma \fCenter !A \lif !B$
\end{prooftree}
\olpseganchor{OLP-0694-B009}{end}

\olpseganchor{OLP-0694-B010}{start}
په همدې طريقه له نورو قواعدو هم پايلې راځي؛
تفصيل ئې دا دے:
\begin{gather*}
  \Axiom $\Gamma \fCenter !A$
  \Axiom $\Delta \fCenter !B$
  \RightLabel{\LR{$\Intro{\land}$}}
  \BinaryInf $\Gamma,\Delta \fCenter !A \land !B$
  \DisplayProof\\
  \Axiom $\Gamma \fCenter !A \land !B$
  \RightLabel{\LR{$\Elim{\land}_1$}}
  \UnaryInf $\Gamma \fCenter !A$
  \DisplayProof
  \qquad
  \Axiom $\Gamma \fCenter !A \land !B$
  \RightLabel{\LR{$\Elim{\land}_2$}}
  \UnaryInf $\Gamma \fCenter !B$
  \DisplayProof
  \\
  \Axiom $\Gamma \fCenter !A$
  \RightLabel{\LR{$\Intro{\lor}_1$}}
  \UnaryInf $\Gamma \fCenter !A \lor !B$
  \DisplayProof\qquad
  \Axiom $\Gamma \fCenter !B$
  \RightLabel{\LR{$\Intro{\lor}_2$}}
  \UnaryInf $\Gamma \fCenter !A \lor !B$
  \DisplayProof\\
  \Axiom $\Gamma \fCenter !A \lor !B$
  \Axiom $\Delta, !A \fCenter !C$
  \Axiom $\Delta', !B \fCenter !C$
  \RightLabel{\LR{$\Elim{\lor}$}}
  \TrinaryInf $\Gamma, \Delta, \Delta' \fCenter !C$
  \DisplayProof
  \\
  \Axiom $\Gamma, !A \fCenter !B$
  \RightLabel{\LR{$\Intro{\lif}$}}
  \UnaryInf $\Gamma \fCenter !A \lif !B$
  \DisplayProof
  \qquad
  \Axiom $\Delta \fCenter !A \lif !B$
  \Axiom $\Gamma \fCenter !A$
  \RightLabel{\LR{$\Elim{\lif}$}}
  \BinaryInf $\Gamma, \Delta \fCenter !B$
  \DisplayProof
  \\
  \Axiom $\Gamma \fCenter \lfalse$
  \RightLabel{\LR{$\FalseInt$}}
  \UnaryInf $\Gamma \fCenter !A$
  \DisplayProof
\end{gather*}
\olpseganchor{OLP-0694-B010}{end}

\olpseganchor{OLP-0694-B011}{start}
هر فرض په خپله له~\LR{$!A$} څخه د \LR{$!A$}
اشتقاق دے؛ يعنې تل \LR{$!A \Sequent !A$} لرو۔
\olpseganchor{OLP-0694-B011}{end}

\olpseganchor{OLP-0694-B012}{start}
\begin{prooftree}
  \AxiomC{\LR{$ $}}
  \UnaryInfC{\LR{$!A \fCenter !A$}}
\end{prooftree}
\olpseganchor{OLP-0694-B012}{end}

\olpseganchor{OLP-0694-B013}{start}
دغه قواعد يوځاے د هغو فطري استنتاجي
اشتقاقونه د شتون حساب ګڼلے شو چې شته۔ د
فطري استنتاج د نښو يوه بله بڼه ئې هم ګڼلے شو،
چې هر ګام ئې يوازې اشتقاق شوے فارمول
نۀ، بلکې هغه نۀ ايستل شوي فرضونه هم ثبتوي
چې ترې اشتقاق شوے دے۔
\olpseganchor{OLP-0694-B013}{end}

\olpseganchor{OLP-0694-B014}{start}
\begin{prooftree}
  \Axiom$!A \fCenter !A$
  \UnaryInf $!A \fCenter !A \lor (!A \lif \lfalse)$
  \Axiom $!B \fCenter !B$
  \BinaryInf$!A, !B \fCenter \lfalse$
  \UnaryInf$!B \fCenter !A \lif \lfalse$
  \UnaryInf$!B \fCenter !A \lor (!A \lif \lfalse)$
  \Axiom$!B \fCenter !B$
  \BinaryInf$!B \fCenter \lfalse$
  \UnaryInf$\fCenter !B \lif \lfalse$
\end{prooftree}
دلته \LR{$!B$} د \LR{$(!A \lor (!A \lif \lfalse))\lif \lfalse$}
لنډ نوم دے۔
\olpseganchor{OLP-0694-B014}{end}

\olpseganchor{OLP-0688-B004}{start}
\olfileid{int}{pty}{ter}
\olpseganchor{OLP-0688-B004}{end}

\olpseganchor{OLP-0688-B005}{start}
\olsection{ثبوت ترمونه}
\olpseganchor{OLP-0688-B005}{end}

\olpseganchor{OLP-0688-B006}{start}
په \Log{N1i} کښې فرضونه په متغيرونو نښه کوو (مثلاً
\LR{$!A^x$})، څو د فرض ايستلو نښې وي۔ په \Log{N2i} کښې هر
سېکوېنټ \LR{$\Gamma \Sequent !A$} په کيڼ سياق~\LR{$\Gamma$}
کښې د متغير--فارمول جوړو لړ لري۔ نو په \Log{N2i}
کښې د ثبوت هر سېکوېنټ يوازې دا نۀ ښيي چې ثبوت
تر اوسه څه ثابت کړي (يعنې \LR{$!A$})، بلکې دا هم ښيي
چې په هر ځای کښې کوم فرضونه پرانيستي دي (يعنې
\LR{$\Gamma$})۔ که په هر سېکوېنټ کښې دا معلومات هم شامل
کړو چې \LR{$!A$} \emph{څنګه} ثابت شوے، نو څه به وشي؟
د دې دپاره يو \emph{په ترمونو نښه شوے} نظام \Log{tN2i}
تعريفولے شو، چې سېکوېنټونه ئې د \LR{$\Gamma \Sequent N : !A$}
په بڼه وي۔ \LR{$\Gamma$}~او \LR{$!A$} هماغه پخواني دي: د
\LR{$x: !B$} په بڼه جوړو سېټ، او هغه فارمول چې په
دغه سېکوېنټ پاې ته رسېدونکي فرعي-ثبوت ئې له
فرضونو~\LR{$\Gamma$} څخه پيروي ښودلې ده۔ \LR{$N$}~به داسې
ترم وي چې په دغه سېکوېنټ پاې ته رسېدونکے ثبوت
کوډ کوي۔
\olpseganchor{OLP-0688-B006}{end}

\olpseganchor{OLP-0688-B007}{start}
هر سېکوېنټ ته ټاکل شوي ترمونه له متغيرونو \LR{$x$}، \LR{$y$}،
\dots، او داسې تابع نښو جوړېږي چې تر اوسه ترسره شوي
استنتاجونه کوډ کوي۔ د هرې استنتاجي قاعدې دپاره بېله
تابع نښه پکار ده۔ د داخلولو له قواعدو سره
تړلې نښې «جوړوونکې» او د ايستلو له قواعدو
سره تړلې نښې «ماتوونکې» بلل کېږي۔ هره تابع نښه د
قاعدې د مقدمو په شمېر آرګومېنټونه اخلي۔ مثلاً د
\Intro{\land} او \Elim{\land}[1] نښې \LR{$c^\land$}
(دوه‌ځايه) او \LR{$d^\land_1$} (يوځايه) کېدے شي۔ په
\Log{tN2i} کښې دغه قواعد داسې کېږي:
\[
\Axiom$\Gamma_1 \fCenter N: !A$
\Axiom$\Gamma_2 \fCenter M: !B$
\RightLabel{\Intro{\land}}
\BinaryInf$\Gamma_1, \Gamma_2 \fCenter c^\land(N, M) : !A \land !B $
\DisplayProof
\quad
\Axiom$\Gamma \fCenter N: !A \land !B$
\RightLabel{\Elim{\land}[1]}
\UnaryInf$\Gamma \fCenter d^\land_1(N) : !A$
\DisplayProof
\]
که استنتاجي قاعده يو فرض باسي، يعنې په يوه مقدمه کښې
نښه لرونکے فارمول~\LR{$x:!A$} وي چې د پايلې په سياق
کښې نشته، نو په ثبوت ترم کښې هم دغه خبره ښودل پکار
دي۔ له دغو قواعدو سره تړلې تابع نښې اړوند متغيرونه
\emph{تړي}۔ مثلاً د \Intro{\lif} قاعده له مقدمې
\LR{$x:!A, \Gamma \Sequent !B$} څخه پايلې
\LR{$\Gamma \Sequent !A \lif !B$} ته ځي۔ جوړوونکے
\LR{$c^{\lif}$} يو آرګومېنټ اخلي، خو دا هم بايد وښيي چې
متغير \LR{$x$} تړل کېږي۔ ثبوت ترم په دې توګه ښيي چې د
نښې~\LR{$x$} فرض ايستل شوے دے۔ مثلاً که د مقدمې ثبوت
ترم~\LR{$N$} وي، د پايلې ثبوت ترم \LR{$c^{\lif}([x]N)$}
او قاعده داسې ليکلے شو:
\[
\Axiom$x:!A, \Gamma \fCenter N: !B$
\RightLabel{\Intro{\lif}}
\UnaryInf$\Gamma \fCenter c^{\lif}([x]N) : !A \lif !B$
\DisplayProof
\]
\olpseganchor{OLP-0688-B007}{end}

\olpseganchor{OLP-0688-B008}{start}
د ثبوت د بشپړې بيا جوړونې دپاره نور معلومات
هم پکار دي۔ که د قاعدې د مقدمو فارمولونه د پايلې
فارمول په بشپړ ډول نۀ ټاکي، دغه معلومات بايد
ترم ته ورزيات کړو۔ په \Intro{\lif} کښې همداسې دي،
نو \LR{$c^{\lif}([x:!A]N)$} ليکو۔ په \Intro{\lor} او
~\FalseInt کښې همداسې دي، نو داسې تابع نښې کاروو
چې دغه معلومات ولري؛ مثلاً:
\[
\Axiom$\Gamma \fCenter N:!A$
\RightLabel{\Intro{\lor}[1]}
\UnaryInf$\Gamma \fCenter c^{\lor{!B}}_1(N):!A \lor !B$
\DisplayProof
\quad
\Axiom$\Gamma \fCenter N:\lfalse$
\RightLabel{\FalseInt}
\UnaryInf$\Gamma \fCenter d^\lfalse_{!A}(N):!A$
\DisplayProof
\]
\olpseganchor{OLP-0688-B008}{end}

\olpseganchor{OLP-0688-B009}{start}
په دغو طريقو داسې سېکوېنټي فطري استنتاجي نظام
جوړولے شو چې هر سېکوېنټ ئې د \LR{$\Gamma \Sequent N : !A$}
په بڼه وي او \LR{$N$}~يو \emph{ثبوت ترم} وي۔
\olpseganchor{OLP-0688-B009}{end}

\olpseganchor{OLP-0688-B010}{start}
دغه ثبوت ترمونه له هغه حساب د ترمونو سره نږدې
اړيکه لري چې \emph{ټايپ‌لرونکے لامبډا حساب} بلل
کېږي۔ د دې اړيکې د ساتلو دپاره به د لامبډا حساب
په ادبياتو کښې کارېدونکې نښې وکاروو، د پورته
عمومي جوړوونکو او ماتوونکو نښو پر ځاے۔ مثلاً د
\LR{$c^{\lif}([x:!A]N)$} پر ځاے \LR{$\lambd[\typeof{x}{!A}][N]$}،
د \LR{$d^{\lif}(N,M)$} پر ځاے \LR{$(NM)$}، د \LR{$c^\land(N,M)$}
پر ځاے \LR{$\pair{N}{M}$}، او د~\LR{$d^\land_i(N)$} پر ځاے
\LR{$\proj{i}{N}$} ليکو۔
\olpseganchor{OLP-0688-B010}{end}

\olpseganchor{OLP-0688-B011}{start}
\begin{defn}
  ثبوت ترمونه په استقرايي ډول د لاندې قواعدو په وسيله جوړېږي:
  \begin{enumerate}
  \item هر متغير \LR{$x$} ثبوت ترم دے (بديهيتونه)۔
  \item که \LR{$x$} متغير، \LR{$!A$}، فارمول او \LR{$N$} ثبوت ترم
    وي، نو \LR{$\lambd[\typeof{x}{!A}][N]$} هم ثبوت ترم دے~(\Intro\lif)۔
  \item که \LR{$N$} او \LR{$M$} ثبوت ترمونه وي، نو \LR{$(NM)$} هم ثبوت
    ترم دے~(\Elim\lif)۔
  \item که \LR{$N$} او \LR{$M$} ثبوت ترمونه وي، نو \LR{$\pair{N}{M}$}
    ثبوت ترم دے~(\Intro\land)۔
  \item که \LR{$N$} ثبوت ترم وي، نو \LR{$\proj{i}{N}$} هم ثبوت
    ترم دے؛ دلته \LR{$i$}، \LR{$1$} يا~\LR{$2$} دے~(\Elim\land[i])۔
  \item که \LR{$N$} ثبوت ترم او \LR{$!A$} يو فورمول وي، نو
    \LR{$\inj[!A]{i}{N}$} ثبوت ترم دے؛ دلته \LR{$i$}، \LR{$1$} يا~\LR{$2$}
    دے~(\Intro\lor[i])۔
  \item که \LR{$M$}، \LR{$N_1$} او \LR{$N_2$} ثبوت ترمونه، او \LR{$x_1$} او
    \LR{$x_2$} متغيرونه وي، نو \LR{$\dcase{M}{x_1}{N_1}{x_2}{N_2}$}
    ثبوت ترم دے~(\Elim\lor)۔
  \item که \LR{$N$} ثبوت ترم او \LR{$!A$}، فارمول وي، نو
    \LR{$\abort{!A}{N}$} ثبوت ترم دے~(\FalseInt)۔
  \end{enumerate}
\end{defn}
\olpseganchor{OLP-0688-B011}{end}

\olpseganchor{OLP-0688-B012}{start}
هر ثبوت ترم د ثبوت ترم نۀ دے۔ مثلاً ترم
\LR{$\pair{N}{M}$} داسې ثبوت کوډ کوي چې په يوه~\Intro{\land}
استنتاج پاے ته رسېږي، نو پاې-فارمول ئې د او
تړنې فورمول~\LR{$!A \land !B$} دے۔ دا د يوې~\Elim{\lif}
قاعدې لويه مقدمه نۀ شي کېدے، نو \LR{$(\pair{N}{M}P)$} د
صحيح ثبوت ترم نۀ شي کېدے۔ يوازې هغه ثبوت ترمونه
\emph{صحيح} دي چې د \Log{tN2ip} قواعد تعقيبوي۔ دغه
قواعد په \olref{tab:tN2ip} کښې ورکړل شوي دي۔
\olpseganchor{OLP-0688-B012}{end}

\olpseganchor{OLP-0688-B013}{start}

\olpseganchor{OLP-0688-B013}{end}

\olfileid{int}{pty}{ter}
\olpseganchor{OLP-0692-B004}{start}
\begin{table}
\begin{tabular}{@{}c@{}c@{}}
\multicolumn{2}{l}{\textbf{اکسيومونه:}}\\
\multicolumn{2}{c}{\LR{$x:!A \Sequent x:!A$}}\\[3ex]
\multicolumn{2}{l}{\textbf{د استنتاج قواعد:}}\\[3ex]
\multicolumn{2}{c}{\Axiom$x: !A, \Gamma \fCenter N:!B$
\RightLabel{\Intro{\lif}}
\UnaryInf$\Gamma \fCenter \lambd[\typeof{x}{!A}][N] : !A \lif !B $
\DisplayProof}
\\[4ex]
\multicolumn{2}{c}{\Axiom$\Gamma_1 \fCenter N:!A \lif !B$
\Axiom$\Gamma_2 \fCenter M:!A$
\RightLabel{\Elim{\lif}}
\BinaryInf$\Gamma_1, \Gamma_2 \fCenter (NM) : !B$
\DisplayProof}
\\[4ex]
\multicolumn{2}{c}{\Axiom$\Gamma_1 \fCenter N_1 : !A$
\Axiom$\Gamma_2 \fCenter N_2 : !B$
\RightLabel{\Intro{\land}}
\BinaryInf$\Gamma_1, \Gamma_2 \fCenter \pair{N_1}{N_2} : !A \land !B $
\DisplayProof}
\\[4ex]
\Axiom$\Gamma \fCenter N: !A \land !B$
\RightLabel{\Elim{\land}[1]}
\UnaryInf$\Gamma \fCenter \proj{1}{N} : !A$
\DisplayProof
&
\Axiom$\Gamma \fCenter N:!A \land !B$
\RightLabel{\Elim{\land}[2]}
\UnaryInf$\Gamma \fCenter \proj{2}{N} : !B$
\DisplayProof\\[3ex]
\\[4ex]
\Axiom$\Gamma \fCenter N: !A$
\RightLabel{\Intro{\lor}[1]}
\UnaryInf$\Gamma \fCenter \inj[!B]{1}{N} : !A \lor !B$
\DisplayProof
&
\Axiom$ \Gamma \fCenter N : !B$
\RightLabel{\Intro{\lor}[2]}
\UnaryInf$\Gamma \fCenter \inj[!A]{2}{N} : !A \lor !B$
\DisplayProof\\[3ex]
\multicolumn{2}{c}{
\begin{tabular}[b]{@{}l@{}}
\Axiom$\Gamma_1 \fCenter M : !A \lor !B$
\Axiom$x:!A, \Gamma_2 \fCenter N_1 : !C$
\Axiom$y:!B, \Gamma_3 \fCenter N_2: !C$
\RightLabel{\Elim{\lor}}
\TrinaryInf$\Gamma_1, \Gamma_2, \Gamma_3 \fCenter
\dcase{M}{x}{N_1}{y}{N_2} : !C$
\DisplayProof
\end{tabular}
}
\\[4ex]
\multicolumn{2}{c}{\Axiom$\Gamma \fCenter N: \lfalse$
\RightLabel{\FalseInt}
\UnaryInf$\Gamma \fCenter \abort{!A}{N} : !A$
\DisplayProof}
\end{tabular}
\caption{د بيانونو د حدسي فطري استنتاج د په ترمونو نښه شوي
نظام \Log{tN2ip} قواعد۔}
\ollabel{tab:tN2ip}
\end{table}
\olpseganchor{OLP-0692-B004}{end}

\olpseganchor{OLP-0689-B004}{start}
\olfileid{int}{pty}{pt}
\olpseganchor{OLP-0689-B004}{end}

\olpseganchor{OLP-0689-B005}{start}
\olsection{اشتقاقونه په ثبوت ترمونو بدلول}
\olpseganchor{OLP-0689-B005}{end}

\olpseganchor{OLP-0689-B006}{start}
د فطري استنتاج ثبوتونه په ثبوت ترمونو بدلول، يعنې
د \Log{N2i}~ثبوتونه د \Log{tN2i}~په ثبوتونه
اړول، ډېر نېغ بهير دے۔ يوازې دومره کوو چې په
اشتقاق کښې د هر سېکوېنټ په ښي‌لاس اړخ کښې
د هر فارمول کيڼې خوا ته يو ثبوت ترم وليکو؛ په
پايله کښې د \LR{$\Gamma \Sequent M : !A$} په بڼه
سېکوېنټونه ترلاسه کېږي۔ بيا وايو چې \LR{$M$} په
سياق~\LR{$\Gamma$} کښې د~\LR{$!A$} \emph{شاهد} دے۔ کوم ثبوت
ترم ليکل پکار دي، دا د \Log{tN2i} قواعد په بشپړ
ډول ټاکي۔
\olpseganchor{OLP-0689-B006}{end}

\olpseganchor{OLP-0689-B007}{start}
\begin{prop}
په \Log{N2i} کښې د \LR{$\Gamma \Sequent !A$} هر
ثبوت~\LR{$\delta$} ته په \Log{tN2i} کښې د
\LR{$\Gamma \Sequent N : !A$} ثبوت~\LR{$\delta'$}
ورسره برابر وي۔ ثبوت ترم \LR{$N$} په يکتا ډول د~\LR{$\delta$}
له مخې ټاکل کېږي۔
\end{prop}
\olpseganchor{OLP-0689-B007}{end}

\olpseganchor{OLP-0689-B008}{start}
\begin{proof}
د~\LR{$\delta$} په لوړوالي استقرا کوو۔ که \LR{$\delta$} د
\LR{$x:!A \Sequent !A$} اکسيوم وي، نو \LR{$\delta'$} د
\LR{$x:!A \Sequent x:!A$} اکسيوم دے۔ که \LR{$\delta$} په
يوه استنتاجي قاعده پاې ته رسېږي، استقرايي فرضيه
هرې مقدمې ته ثبوت ترم ورکوي او د دغو مقدمو سره
برابر د سېکوېنټونو \Log{tN2i}-ثبوتونه ورکوي چې
دغه ثبوت ترمونه پکې شامل وي۔ د \Log{tN2i} اړونده
قاعده کاروو؛ اړوند ثبوت ترم د \Log{tN2i} استنتاجي
قاعده ټاکي۔
\end{proof}
\olpseganchor{OLP-0689-B008}{end}

\olpseganchor{OLP-0689-B009}{start}
\begin{ex}
  په \Log{N1i} کښې د \LR{$(!A \land !B) \lif (!B \land !A)$}
  دا ډېر ساده ثبوت وګورئ:
\[
  \AxiomC{\LR{$\Discharge{!A\land !B}{x}$}}
  \RightLabel{\Elim\land[2]}
  \UnaryInfC{\LR{$!B$}}
  \AxiomC{\LR{$\Discharge{!A\land !B}{x}$}}
  \RightLabel{\Elim\land[1]}
  \UnaryInfC{\LR{$!A$}}
  \RightLabel{\Intro\land}
  \BinaryInfC{\LR{$!B \land !A$}}
  \DischargeRule{\Intro\lif}{x}
  \UnaryInfC{\LR{$(!A \land !B) \lif (!B \land !A)$}}
  \DisplayProof
  \]
  په \Log{N2i} کښې ئې بڼه داسې ده:
\[
  \Axiom$x : !A\land !B \fCenter !A\land !B$
  \RightLabel{\Elim\land[2]}
  \UnaryInf$x : !A\land !B \fCenter !B$
  \Axiom$x : !A\land !B \fCenter !A\land !B$
  \RightLabel{\Elim\land[1]}
  \UnaryInf$x : !A\land !B \fCenter !A$
  \RightLabel{\Intro\land}
  \BinaryInf$x : !A\land !B \fCenter !B \land !A$
  \RightLabel{\Intro\lif}
  \UnaryInf$\fCenter (!A \land !B) \lif (!B \land !A)$
  \DisplayProof
  \]
  ثبوت ترمونه له پاسه ښکته ټاکو۔ د اکسيومونو ثبوت
  ترمونه هماغه د سياق متغير~\LR{$x$} دے۔ بيا له
  \LR{$\Elim\land[i]$} سره تړلي ثبوت ترمونه
  \LR{$\proj{2}{x}$} او~\LR{$\proj{1}{x}$} ټاکل کېږي۔ کله چې
  \Intro{\land} کاروو، دغه دوه ثبوت ترمونه يوځاے
  کوو: \LR{$\pair{\proj{2}{x}}{\proj{1}{x}}$}۔ په پاې
  کښې د \Intro{\lif} قاعده يو \LR{$\lambd$}-تجريد
  کاروي،~\LR{$x$} تړي او ثبتوي چې \LR{$x$} د~\LR{$!A \land !B$}
  فرض نښه وه: \LR{$\lambd[\typeof{x}{!A \land
  !B}][\pair{\proj{2}{x}}{\proj{1}{x}}]$}۔ نو په
  \Log{tN2i} کښې ثبوت داسې کېږي:
\[
  \Axiom$x : !A\land !B \fCenter x: !A\land !B$
  \RightLabel{\Elim\land[2]}
  \UnaryInf$x : !A\land !B \fCenter \proj{2}{x} : !B$
  \Axiom$x : !A\land !B \fCenter x: !A\land !B$
  \RightLabel{\Elim\land[1]}
  \UnaryInf$x : !A\land !B \fCenter \proj{1}{x} : !A$
  \RightLabel{\Intro\land}
  \BinaryInf$x : !A\land !B \fCenter \pair{\proj{2}{x}}{\proj{1}{x}} : 
    !B \land !A$
  \RightLabel{\Intro\lif}
  \UnaryInf$\fCenter \lambd[\typeof{x}{!A \land
  !B}][\pair{\proj{2}{x}}{\proj{1}{x}}] : (!A \land !B) \lif (!B \land !A)$
  \DisplayProof
  \]
\end{ex}
\olpseganchor{OLP-0689-B009}{end}

\olpseganchor{OLP-0695-B004}{start}
\olfileid{int}{pty}{tp}
\olpseganchor{OLP-0695-B004}{end}

\olpseganchor{OLP-0695-B005}{start}
\olsection{له ثبوت ترمونو څخه اشتقاقونه بېرته اخيستل}
\olpseganchor{OLP-0695-B005}{end}

\olpseganchor{OLP-0695-B006}{start}
اوس بل لور ته ګورو: يو ثبوت ترم~\LR{$N$} بېرته
په \Log{N2i} کښې ثبوت ته اړول۔ په عمومي
ډول، دا يوازې د داسې سياق په نسبت شوني دي
چې د ورکړي ثبوت ترم د هر آزاد متغير~\LR{$x$} دپاره
د \LR{$x : !A$} جوړه ولري۔ خو لومړے داسې بېلګه
اخلو چې آزاد متغيرونه نۀ لري؛ يعنې دا ترم:
\[
\lambd[\typeof{x}{!A \land
  !B}][\pair{\proj{2}{x}}{\proj{1}{x}}]
\]
دا د \LR{$\lambd[\typeof{x}{!A \land !B}][N_1]$} په
بڼه دے۔ په \Log{tN2i} کښې يوازينۍ قاعده چې
داسې ترم جوړوي،~\Intro{\lif} ده۔ نو پوهېږو چې
د~\LR{$N$} کوډ شوے ثبوت بايد په~\Intro{\lif}
پاې ته ورسېږي؛ يعنې په دې پاې ته رسېږي:
\[
  \Axiom$x : !A\land !B \fCenter \pair{\proj{2}{x}}{\proj{1}{x}} : 
    !C$
  \RightLabel{\Intro\lif}
  \UnaryInf$\fCenter \lambd[\typeof{x}{!A \land
  !B}][\pair{\proj{2}{x}}{\proj{1}{x}}] : (!A \land !B) \lif !C$
  \DisplayProof
  \]
تر اوسه نۀ پوهېږو چې \LR{$!C$} څه دے، خو دا پوهېږو
چې مخکنۍ برخه \LR{$!A \land !B$} ده، او د مقدمې په
سياق کښې \LR{$x:!A \land !B$} شاملېدل پکار دي۔
\olpseganchor{OLP-0695-B006}{end}

\olpseganchor{OLP-0695-B007}{start}
اوس د مقدمې اړوند ثبوت ترم له مخه ټاکل شوے
دے: \LR{$\pair{\proj{2}{x}}{\proj{1}{x}}$}۔ هغه
ثبوت چې \emph{همدغه} ترم ئې کوډ کوي، بايد
په \Intro{\land} پاې ته ورسېږي۔ لا نۀ پوهېږو
چې \LR{$!C$} څه دے، خو څرنګه چې د~\Intro{\land}
پايله ده، بايد د او تړنه وي؛ يعنې
\LR{$!C \ident (!C_1 \land !C_2)$}۔ بياجوړونه
داسې پرمخ وړو:
\[
  \Axiom$x : !A\land !B \fCenter \proj{2}{x} : !C_1$
  \Axiom$x : !A\land !B \fCenter \proj{1}{x} : !C_2$
  \RightLabel{\Intro\land}
  \BinaryInf$x : !A\land !B \fCenter \pair{\proj{2}{x}}{\proj{1}{x}} : 
    !C_1 \land !C_2$
  \RightLabel{\Intro\lif}
  \UnaryInf$\fCenter \lambd[\typeof{x}{!A \land
  !B}][\pair{\proj{2}{x}}{\proj{1}{x}}] : (!A \land !B) \lif (!C_1 \land !C_2)$
  \DisplayProof
  \]
ثبوت ترم \LR{$\pi_2(x)$} ښيي چې پاسنے کيڼ سېکوېنټ
د \Elim{\land}[2] په وسيله ترلاسه شوے؛ يعنې
مقدمه د \LR{$!D \land !C_1$} په بڼه ده۔ په ښي‌لاس
اړخ کښې دا ټاکلے شو چې \LR{$!C_2$} ته رسوونکے
استنتاج \Elim{\land}[1] دے، او مقدمه ئې د
\LR{$!C_2 \land !E$} په بڼه ده۔
\[
  \Axiom$x : !A\land !B \fCenter x: !D\land !C_1$
  \RightLabel{\Elim\land[2]}
  \UnaryInf$x : !A\land !B \fCenter \proj{2}{x} : !C_1$
  \Axiom$x : !A\land !B \fCenter x: !C_2\land !E$
  \RightLabel{\Elim\land[1]}
  \UnaryInf$x : !A\land !B \fCenter \proj{1}{x} : !C_2$
  \RightLabel{\Intro\land}
  \BinaryInf$x : !A\land !B \fCenter \pair{\proj{2}{x}}{\proj{1}{x}} : 
    !C_1 \land !C_2$
  \RightLabel{\Intro\lif}
  \UnaryInf$\fCenter \lambd[\typeof{x}{!A \land
  !B}][\pair{\proj{2}{x}}{\proj{1}{x}}] : (!A \land !B) \lif (!C_1 \land !C_2)$
  \DisplayProof
  \]
څرنګه چې د تر ټولو پاسنيو سېکوېنټونو ثبوت
ترمونه اوس متغيرونه دي، پوهېږو چې اکسيومونه
دي۔ دا \LR{$!C_1$}، \LR{$!C_2$}، \LR{$!D$} او \LR{$!E$} ټاکي:
\LR{$!D \ident !A$} او \LR{$!C_1 \ident !B$}، کنه کيڼ
سېکوېنټ اکسيوم نۀ وي؛ او \LR{$!C_2 \ident !A$} او
\LR{$!E \ident !B$}، کنه ښے سېکوېنټ اکسيوم نۀ وي۔
ثبوت مو په بشپړ ډول بېرته ترلاسه کړ:
\[
  \Axiom$x : !A\land !B \fCenter x: !A\land !B$
  \RightLabel{\Elim\land[2]}
  \UnaryInf$x : !A\land !B \fCenter \proj{2}{x} : !B$
  \Axiom$x : !A\land !B \fCenter x: !A\land !B$
  \RightLabel{\Elim\land[1]}
  \UnaryInf$x : !A\land !B \fCenter \proj{1}{x} : !A$
  \RightLabel{\Intro\land}
  \BinaryInf$x : !A\land !B \fCenter \pair{\proj{2}{x}}{\proj{1}{x}} : 
    !B \land !A$
  \RightLabel{\Intro\lif}
  \UnaryInf$\fCenter \lambd[\typeof{x}{!A \land
  !B}][\pair{\proj{2}{x}}{\proj{1}{x}}] : (!A \land !B) \lif (!B \land !A)$
  \DisplayProof
  \]
\olpseganchor{OLP-0695-B007}{end}

\olpseganchor{OLP-0697-B004}{start}
\olfileid{int}{pty}{typ}
\olpseganchor{OLP-0697-B004}{end}

\olpseganchor{OLP-0697-B005}{start}
\olsection{ټايپونه}
\olpseganchor{OLP-0697-B005}{end}

\olpseganchor{OLP-0697-B006}{start}
ثبوت ترم، مثلاً \LR{$(MN)$}، پخپله په يکتا ډول هغه ثبوت
نۀ ټاکي چې دا ئې ثبوت ترم دے۔ ځکه آزاد متغيرونه، يعنې
د اکسيومونو ثبوت ترمونه، د اړوندو اکسيومونو فارمولونه
نۀ ښيي۔ خو که دا په سياق~\LR{$\Gamma$} کښې وټاکو، نو
که ترم ټايپ ولري، د ټايپ ټاکنې استدلال ئې \emph{په يکتا ډول} جوړېږي۔
\olpseganchor{OLP-0697-B006}{end}

\olpseganchor{OLP-0697-B007}{start}
\begin{defn}
د ثبوت ترم~\LR{$N$} دپاره \emph{سياق} داسې مجموعه \LR{$\Gamma$}
ده چې د~\LR{$N$} د هر آزاد متغير دپاره په \LR{$\Gamma$} کښې
د \LR{$x : !A$} کټ مټ يوه جوړه وي۔
\end{defn}
\olpseganchor{OLP-0697-B007}{end}

\olpseganchor{OLP-0697-B008}{start}
په ياد ولرئ چې تعريف دا شرط نۀ ږدي چې \LR{$\Gamma$}
دې \LR{$x:!A$} \emph{يوازې} د~\LR{$N$} د آزادو متغيرونو دپاره
ولري۔ مثلاً \LR{$\Gamma = \{x:!A, y:!B, z:!C\}$}
د \LR{$\pair{x}{y}$} دپاره يو سياق دے۔
\olpseganchor{OLP-0697-B008}{end}

\olpseganchor{OLP-0697-B009}{start}
\begin{defn}\ollabel{defn:type}
  فرض کړئ \LR{$\Gamma$} د~\LR{$N$} دپاره يو سياق دے۔ د~\LR{$N$}
  \emph{ټايپ} د~\LR{$\Gamma$} په نسبت په استقرايي ډول داسې
  تعريفېږي:
  \begin{enumerate}
  \item د \LR{$x$} ټايپ \LR{$!A$} هغه وخت او يوازې هغه وخت دے
  چې \LR{$x:!A \in \Gamma$} وي۔
  \item که \LR{$N$} د \LR{$x:!A, \Gamma$} په نسبت ټايپ~\LR{$!B$}
  ولري، نو د \LR{$\lambd[\typeof{x}{!A}][N]$} ټايپ \LR{$!A \lif !B$} دے۔
  \item که \LR{$N$} ټايپ \LR{$!A \lif !B$} او \LR{$M$} ټايپ~\LR{$!A$}
  ولري، نو \LR{$(NM)$} ټايپ~\LR{$!B$} لري۔
  \item که \LR{$N$} ټايپ \LR{$!A$} او \LR{$M$} ټايپ~\LR{$!B$} ولري،
  نو \LR{$\pair{N}{M}$} ټايپ \LR{$!A \land !B$} لري۔
  \item که \LR{$N$} ټايپ~\LR{$!A_1 \land !A_2$} ولري، نو
  \LR{$\proj{i}{N}$} ټايپ~\LR{$!A_i$} لري۔
  \item که \LR{$N$} ټايپ \LR{$!A$} ولري، نو
  \LR{$\inj[!B]{i}{N}$} ټايپ \LR{$!A \lor !B$} لري که \LR{$i=1$}
  وي، او ټايپ \LR{$!B \lor !A$} لري که \LR{$i=2$} وي۔
  \item که \LR{$M$} ټايپ \LR{$!A \lor !B$} ولري، او \LR{$N_1$}
  د \LR{$x:!A, \Gamma$} په نسبت ټايپ \LR{$!C$} او \LR{$N_2$}
  د \LR{$y:!B, \Gamma$} په نسبت ټايپ~\LR{$!C$} ولري، نو
  \LR{$\dcase{M}{x}{N_1}{y}{N_2}$} ټايپ~\LR{$!C$} لري۔
  \item که \LR{$N$} ټايپ~\LR{$\lfalse$} ولري، نو
  \LR{$\abort{!A}{N}$} ټايپ~\LR{$!A$} لري۔
  \end{enumerate}
\end{defn}
\olpseganchor{OLP-0697-B009}{end}

\olpseganchor{OLP-0697-B010}{start}
\begin{prop}
ثبوت ترم~\LR{$N$} د~\LR{$N$} لپاره په ټاکلي سياق~\LR{$\Gamma$} کښې
تر ډېره يو ټايپ لري۔
\end{prop}
\olpseganchor{OLP-0697-B010}{end}

\olpseganchor{OLP-0697-B011}{start}
\begin{proof}
  پر~\LR{$N$} استقرا کوو۔ د متغير ټايپ د سياق له يوازينۍ جوړې
  ټاکل کېږي۔ په تجريد کښې د فرض ټايپ پر خپله نښه ليکل شوے
  دے او د بدن ټايپ د استقرا له مخې ټاکلے دے۔ په تطبيق کښې
  د تابع ټايپ تالي پايله ټاکي او مقدم ئې د آرګومېنټ له
  ټايپ سره برابرېدل غواړي۔ د جوړې ټايپ د دواړو غړو له
  ټايپونو جوړېږي، او پروجکشن هماغه ټاکلے جز اخلي۔ د يا
  جوړوونکے د خپل ترم ټايپ له ليکل شوي ټايپ سره يو ځاے
  کوي۔ د يا په ايستلو کښې دواړه څانګې، له ټاکلو فرضونو
  سره، يو ګډ ټايپ غواړي، او له دروغو د ايستلو ټايپ
  پخپله په ترم کښې ليکل شوے دے۔ په هر حالت کښې دغه
  شرطونه که پوره شي، يوازې يو ټايپ ورکوي۔
\end{proof}
\olpseganchor{OLP-0697-B011}{end}

\olpseganchor{OLP-0697-B012}{start}
که د \LR{$N$} ټايپ د \LR{$\Gamma$} په نسبت~\LR{$!A$} وي،
نو \LR{$\Gamma \Sequent N : !A$} ليکو۔ د \olref{defn:type}
مواد بايد اشنا ښکاره شي: دا کټ مټ د \Log{tN2i} قواعد
دي، يوازې په يوه لږ توپير سره؛ سياق~\LR{$\Gamma$} په ټوله
لړۍ کښې هماغه دے۔ دا تعريف د \Log{tN3i} په حساب کښې
د \olref{tab:tN2ip} د قواعدو په وسيله ليکلے شو۔
\olpseganchor{OLP-0697-B012}{end}

\olpseganchor{OLP-0697-B013}{start}

\olpseganchor{OLP-0697-B013}{end}

\olpseganchor{OLP-0697-B014}{start}
ثبوت ترمونه د \LR{$\lambd$}-حساب د يوه ډول ترمونه دي؛
يعنې \emph{له حاصل‌ضربونو، مجموعو او تش ټايپ سره
ټايپ‌لرونکے \LR{$\lambd$}-حساب}۔ دوديز \LR{$\lambd$}-حساب يوازې
له متغيرونو، \emph{\LR{$\lambd$}-تجريدونو}~\LR{$\lambd[x][N]$}
او تطبيق \LR{$(MN)$} څخه جوړ ترمونه لري، او د ټايپ‌لرونکي
\LR{$\lambd$}-تجريد \LR{$\lambd[\typeof{x}{!A}][N]$} د~\LR{$!A$}
نښه پکې نشته۔ \LR{$\lambd$}-حساب د محاسبې داسې موډل دے
چې ټيورېنګ بشپړ دے؛ يعنې هره جزوي محاسبه کېدونکې
تابع پکې د يو ترم په وسيله تمثيلېدلے شي۔ د \LR{$\lambd$}-حساب
هغه ډول چې ثبوت ترمونه او د ټايپ ټاکلو قواعد ئې ورکوي،
د محاسبې محدود موډل دے، خو بنسټيزې مفکورې هماغه دي۔
د ټايپ ټاکنه ئې محدودوي: يوازې \emph{صحيح ټايپ‌لرونکي}
ترمونه، يعنې هغه ثبوت ترمونه چې د سياق~\LR{$\Gamma$} په
نسبت ټايپونه لري، روا ترمونه دي۔ مثلاً \LR{$\lambd[x][(xx)]$}
په دوديز بې‌ټايپه \LR{$\lambd$}-حساب کښې روا ترم دے، خو
\LR{$\lambd[\typeof{x}{!A}][(xx)]$} د هېڅ \LR{$!A$} دپاره
صحيح ټايپ نۀ لري۔ ځکه د ټايپ قواعد دا شرط ږدي چې د
\LR{$(xx)$} د ټايپ لرلو دپاره لومړنے~\LR{$x$} بايد ټايپ \LR{$!A \lif !B$}
او دويم \LR{$x$} ټايپ~\LR{$!A$} ولري؛ او دا ناشونې ده، ځکه
\LR{$!A$} د~\LR{$!A \lif !B$} يوه حقيقي فرعي فارمول ده۔
\olpseganchor{OLP-0697-B014}{end}

\olpseganchor{OLP-0697-B015}{start}
په وجداني ډول ټايپونه د هغو شيانو ډول ګڼلے شو چې
د هماغه ټايپ ترم ئې نوم اخلي۔ نو د ټايپ \LR{$!A \land !B$}
ترم داسې جوړه ټاکي چې لومړۍ برخه ئې ټايپ~\LR{$!A$} او
دويمه برخه ئې ټايپ~\LR{$!B$} ولري؛ يعنې د حاصل‌ضرب~\LR{$!A \times !B$}
غړے۔ د ټايپ \LR{$!A \lif !B$} ترم له هغو شيانو
چې ټايپ \LR{$!A$} لري، هغو شيانو ته تابع ده چې ټايپ~\LR{$!B$}
لري۔ د ټايپ \LR{$!A \lor !B$} ترم يا د ټايپ \LR{$!A$} يو شے
دے يا د ټايپ~\LR{$!B$} يو شے، او دا معلومات هم ورسره وي
چې کوم يو دے۔ دا د دوو مجموعو \LR{$!A$} او \LR{$!B$} د
بې‌اشتراکه اتحاد په شان دے، چې داسې تعريفېږي:
\LR{$\{\tuple{i,x} : i = 1 \text{\RL{ يا }} 2, x \in !A \text{\RL{ که }} i = 1, x \in !B \text{\RL{ که }} i=2\}$}۔
ټايپ \LR{$\lfalse$} تشه مجموعه ده؛ يعنې که يو ترم دا ټايپ
ولري، د هېڅ شي نوم نۀ شي اخيستلے۔ څرنګه چې طبيعي
اشتقاق، د \Log{tN3i} په ګډون، سازګار دے، د~\LR{$\lfalse$}
ثبوت هله شته چې خلاص فرضونه موجود وي؛ نو داسې
تړلے ثبوت ترم، يعنې بې له آزادو متغيرونو ترم، نشته
چې ټايپ~\LR{$\lfalse$} ولري۔
\olpseganchor{OLP-0697-B015}{end}

\olpseganchor{OLP-0697-B016}{start}
د ټايپ‌لرونکي \LR{$\lambd$}-حساب عملګرونه داسې ترمونه
جوړوي چې په وجداني ډول د سم ډول شيانو نوم اخلي، په
دې شرط چې هغه ترمونه چې عملګر پرې پلي کېږي، د ټاکلو
شيانو نوم واخلي۔ مثلاً دا خبره چې ``که \LR{$N_1 : !A$} او
\LR{$N_2 : !B$} وي، نو \LR{$\pair{N_1}{N_2} : !A \land !B$}''
دا وايي: که \LR{$N_1$} د ټايپ~\LR{$!A$} د يو شي او \LR{$N_2$}
د ټايپ~\LR{$!B$} د يو شي نوم اخلي، نو \LR{$\pair{N_1}{N_2}$}
د ټايپ \LR{$!A \land !B$} د يو شي نوم اخلي۔
\olpseganchor{OLP-0697-B016}{end}

\olfileid{int}{pty}{typ}
\olpseganchor{OLP-0693-B004}{start}
\begin{table}
\begin{tabular}{@{}c@{}c@{}}
\multicolumn{2}{l}{\textbf{اکسيومونه:}}\\
\multicolumn{2}{c}{\LR{$\Gamma \Sequent x:!A$} که \LR{$x:!A \in \Gamma$} وي}\\[3ex]
\multicolumn{2}{l}{\textbf{د استنتاج قواعد:}}\\[3ex]
\multicolumn{2}{c}{\Axiom$x: !A, \Gamma \fCenter N:!B$
\RightLabel{\Intro{\lif}}
\UnaryInf$\Gamma \fCenter \lambd[\typeof{x}{!A}][N] : !A \lif !B $
\DisplayProof}
\\[4ex]
\multicolumn{2}{c}{\Axiom$\Gamma \fCenter N:!A \lif !B$
\Axiom$\Gamma \fCenter M:!A$
\RightLabel{\Elim{\lif}}
\BinaryInf$\Gamma \fCenter (NM) : !B$
\DisplayProof}
\\[4ex]
\multicolumn{2}{c}{\Axiom$\Gamma \fCenter N_1 : !A$
\Axiom$\Gamma \fCenter N_2 : !B$
\RightLabel{\Intro{\land}}
\BinaryInf$\Gamma \fCenter \pair{N_1}{N_2} : !A \land !B $
\DisplayProof}
\\[4ex]
\Axiom$\Gamma \fCenter N: !A \land !B$
\RightLabel{\Elim{\land}[1]}
\UnaryInf$\Gamma \fCenter \proj{1}{N} : !A$
\DisplayProof
&
\Axiom$\Gamma \fCenter N:!A \land !B$
\RightLabel{\Elim{\land}[2]}
\UnaryInf$\Gamma \fCenter \proj{2}{N} : !B$
\DisplayProof\\[3ex]
\\[4ex]
\Axiom$\Gamma \fCenter N: !A$
\RightLabel{\Intro{\lor}[1]}
\UnaryInf$\Gamma \fCenter \inj[!B]{1}{N} : !A \lor !B$
\DisplayProof
&
\Axiom$ \Gamma \fCenter N : !A$
\RightLabel{\Intro{\lor}[2]}
\UnaryInf$\Gamma \fCenter \inj[!B]{2}{N} : !B \lor !A$
\DisplayProof\\[3ex]
\multicolumn{2}{c}{
\begin{tabular}[b]{@{}l@{}}
\Axiom$\Gamma \fCenter M : !A \lor !B$
\Axiom$x:!A, \Gamma \fCenter N_1 : !C$
\Axiom$y:!B, \Gamma \fCenter N_2: !C$
\RightLabel{\Elim{\lor}}
\TrinaryInf$\Gamma \fCenter
\dcase{M}{x}{N_1}{y}{N_2} : !C$
\DisplayProof
\end{tabular}
}
\\[4ex]
\multicolumn{2}{c}{\Axiom$\Gamma \fCenter N: \lfalse$
\RightLabel{\FalseInt}
\UnaryInf$\Gamma \fCenter \abort{!A}{N} : !A$
\DisplayProof}
\end{tabular}
\caption{د بيانونو د حدسي فطري استنتاج د په ترمونو نښه شوي
نظام \Log{tN3ip} قواعد، چې د ټايپ ټاکلو د نظام په توګه
کار کوي۔}
\ollabel{tab:tN2ip}
\end{table}
\olpseganchor{OLP-0693-B004}{end}

\olpseganchor{OLP-0691-B004}{start}
\olfileid{int}{pty}{red}
\olpseganchor{OLP-0691-B004}{end}

\olpseganchor{OLP-0691-B005}{start}
\olsection{راکمول}
\olpseganchor{OLP-0691-B005}{end}

\olpseganchor{OLP-0691-B006}{start}
د فطري استنتاج په اشتقاقونه کښې د
داخلولو قاعدې پسې ايستلو قاعده
داسې تاو راتاو دے چې «راکمېدل» ئې شوني دي۔ مثلاً
د \LR{$\Intro{\land}$} پسې \LR{$\Elim{\land}$} يو بل «لرې
کوي»، ځکه د \LR{$\Elim{\land}$} پايله تل د هغه او د
دواړو برخو يوه وي چې \Intro{\land} ئې جوړوي، او
د \LR{$\Intro{\land}$} دواړه مقدمې هماغه دواړه برخې
دي۔ نو که د يوې برخې اشتقاق پکار وي،
هماغه اشتقاق نېغ کارولے شو، د او د
داخلولو او بيا ايستلو اړتيا نشته۔ داسې لرو:
\begin{gather*}
  \bottomAlignProof
  \AxiomC{}
  \RightLabel{\LR{$\delta_1$}}
  \DeduceC{\LR{$!A_1$}}
  \AxiomC{}
  \RightLabel{\LR{$\delta_2$}}
  \DeduceC{\LR{$!A_2$}}
  \RightLabel{\LR{$\Intro{\land}$}}
  \BinaryInfC{\LR{$!A_1 \land !A_2$}}
  \RightLabel{\LR{$\Elim{\land}_i$}}
  \UnaryInfC{\LR{$!A_i$}}
  \DisplayProof\bottomAlignProof
  \quad
  \redone
  \quad
  \AxiomC{}
  \RightLabel{\LR{$\delta_i$}}
  \DeduceC{\LR{$!A_i$}}
  \DisplayProof
\end{gather*}
\olpseganchor{OLP-0691-B006}{end}

\olpseganchor{OLP-0691-B007}{start}
د دې ساده کولو له مخې په اړوندو ثبوت ترمونو
کښې هم د راکمول ورته اړيکه اخيستلے شو۔ که \LR{$N_1$}
د~\LR{$\delta_1$} ثبوت ترم وي او \LR{$N_2$} د~\LR{$\delta_2$}
وي، نو ويلے شو چې د کيڼ ثبوت ثبوت ترم په يو ګام
کښې د ښي ثبوت ثبوت ترم ته راکمېږي:
\[
  \proj{i}{\pair{N_1}{N_2}} \redone N_i
\]
په ټايپ‌لرونکي لامبډا حساب کښې دا د ضرب ټايپ
\emph{د بېټا راکمول} قاعده ده۔
\olpseganchor{OLP-0691-B007}{end}

\olpseganchor{OLP-0691-B008}{start}
په فطري استنتاج کښې د استنتاجونو داسې جوړه،
لکه د \Intro{\land} پسې \Elim{\land}، چې راکمېدل
پرې عملي کېږي، \emph{پرېکون} بلل کېږي۔ په
ټايپ‌لرونکي لامبډا حساب کښې د \LR{$\redone$} کيڼ ترم
\emph{رېډکس} او ښي ترم ئې \emph{راکم شوے ترم}
بلل کېږي۔ په بې‌ټايپه لامبډا حساب کښې يوازې
\LR{$(\lambd[x][N])Q$} رېډکس ګڼل کېږي؛ په ټايپ‌لرونکي
حساب کښې نحو داسې ترمونه هم رانغاړي چې
\LR{$\pair{N}{M}$}، \LR{$\proj{i}{N}$}، \LR{$\inj[!A]{i}{N}$}،
\LR{$\dcase{N}{x_1}{M_1}{x_2}{M_2}$} او \LR{$\abort{!A}{N}$}
پکې وي، او خپل اړوند رېډکسونه لري۔
د يا جوړوونکے د خپلې څانګې شاخص او ثبوت ترم اخلي؛ په اختياري
ټايپي نښه کښې د بلې څانګې ټايپ ليکو۔
\olpseganchor{OLP-0691-B008}{end}

\olpseganchor{OLP-0691-B009}{start}
په همدې ډول د يا دپاره راکمول لرو:
\begin{gather*}
  \bottomAlignProof
  \AxiomC{}
  \RightLabel{\LR{$\delta$}}
  \DeduceC{\LR{$!A_i$}}
  \RightLabel{\LR{$\Intro{\lor}$}}
  \UnaryInfC{\LR{$!A_1 \lor !A_2$}}
  \AxiomC{\LR{$\Discharge{!A_1}{x_1}$}}
  \RightLabel{\LR{$\delta_1$}}
  \DeduceC{\LR{$!C$}}
  \AxiomC{\LR{$\Discharge{!A_2}{x_2}$}}
  \RightLabel{\LR{$\delta_2$}}
  \DeduceC{\LR{$!C$}}
  \DischargeRule{\LR{$\Elim{\lor}$}}{x_1,x_2}
  \TrinaryInfC{\LR{$!C$}}
  \DisplayProof\bottomAlignProof
  \quad
  \redone
  \quad
  \AxiomC{}
  \RightLabel{\LR{$\delta$}}
  \DeduceC{\LR{$!A_i$}}
  \RightLabel{\LR{$\delta_i$}}
  \DeduceC{\LR{$!C$}}
  \DisplayProof
\end{gather*}
دې ته د ثبوت ترمونو لاندې راکمول برابر دي:
\begin{gather*}
  \dcase{\inj[!A_{3-i}]{i}{M}}{x_1}{N_1}{x_2}{N_2} \redone \Subst{N_i}{M}{x_i}
\end{gather*}
دلته \LR{$M$} د~\LR{$\delta$} ثبوت ترم او \LR{$N_i$} د~\LR{$\delta_i$}
ثبوت ترم دے۔ دا د مجموعې ټايپونو د بېټا راکمول
قاعده ده۔ په عمومي نښه کښې \LR{$\Subst{M}{N_i}{x}$}
دا معنا لري چې په \LR{$M$} کښې د متغير~\LR{$x$} ټول آزاد
موردونه په~\LR{$N_i$} بدل کړو۔
ورګډېدونکے ترم د ټاکلې څانګې ثبوت دے، او ليکل شوې
ټايپي نښه د بلې څانګې ده۔ د ايستلو پر وخت يوازې د ټاکلې
څانګې د فرض پر ځاے همدا ترم ورګډېږي۔
\olpseganchor{OLP-0691-B009}{end}

\olpseganchor{OLP-0691-B010}{start}
د تابعې ټايپ راکمول د \LR{$\Intro\lif$} پسې د
\LR{$\Elim\lif$} د تاو راتاو لرې کولو سره برابر دي۔
\begin{gather*}
  \bottomAlignProof
  \AxiomC{\LR{$\Discharge{!A}{x}$}}
  \RightLabel{\LR{$\delta$}}
  \DeduceC{\LR{$!B$}}
  \DischargeRule{\LR{$\Intro{\lif}$}}{x}
  \UnaryInfC{\LR{$!A \lif !B$}}
  \AxiomC{}
  \RightLabel{\LR{$\delta'$}}
  \DeduceC{\LR{$!A$}}
  \RightLabel{\LR{$\Elim{\lif}$}}
  \BinaryInfC{\LR{$!B$}}
  \DisplayProof\bottomAlignProof
  \quad
  \redone
  \quad
  \AxiomC{}
  \RightLabel{\LR{$\delta'$}}
  \DeduceC{\LR{$!A$}}
  \RightLabel{\LR{$\delta$}}
  \DeduceC{\LR{$!B$}}
  \DisplayProof
\end{gather*}
د ثبوت ترمونو دپاره دا هماغه عادي بېټا راکمول دي:
\begin{gather*}
  (\lambd[\typeof{x}{!A}][N]M)
  \redone \Subst{N}{M}{x}
\end{gather*}
\olpseganchor{OLP-0691-B010}{end}

\olpseganchor{OLP-0691-B011}{start}
په ثبوتونه کښې د راکمول د بدلونونو ترڅنګ
د ځاے بدلولو بدلونونه هم اخيستلے شو۔ دا پر هغو
(فرعي-)ثبوتونه عملي کېږي چې په \Elim{\lor}
پاې ته رسېږي او ورپسې د ايستلو بله
قاعده وي؛ مثلاً:
\begin{multline*}
  \bottomAlignProof
  \AxiomC{}
  \RightLabel{\LR{$\delta$}}
  \DeduceC{\LR{$!A_i$}}
  \RightLabel{\LR{$\Intro{\lor}$}}
  \UnaryInfC{\LR{$!A_1 \lor !A_2$}}
  \AxiomC{\LR{$\Discharge{!A_1}{x_1}$}}
  \RightLabel{\LR{$\delta_1$}}
  \DeduceC{\LR{$!C_1 \land !C_2$}}
  \AxiomC{\LR{$\Discharge{!A_2}{x_2}$}}
  \RightLabel{\LR{$\delta_2$}}
  \DeduceC{\LR{$!C_1 \land !C_2$}}
  \DischargeRule{\LR{$\Elim{\lor}$}}{x_1,x_2}
  \TrinaryInfC{\LR{$!C_1 \land !C_2$}}
  \RightLabel{\Elim{\land}[i]}
  \UnaryInfC{\LR{$!C_i$}}
  \DisplayProof\bottomAlignProof
  \quad
  \redone
  \\
  \AxiomC{}
  \RightLabel{\LR{$\delta$}}
  \DeduceC{\LR{$!A_i$}}
  \RightLabel{\LR{$\Intro{\lor}$}}
  \UnaryInfC{\LR{$!A_1 \lor !A_2$}}
  \AxiomC{\LR{$\Discharge{!A_1}{x_1}$}}
  \RightLabel{\LR{$\delta_1$}}
  \DeduceC{\LR{$!C_1 \land !C_2$}}
  \RightLabel{\Elim{\land}[i]}
  \UnaryInfC{\LR{$!C_i$}}
  \AxiomC{\LR{$\Discharge{!A_2}{x_2}$}}
  \RightLabel{\LR{$\delta_2$}}
  \DeduceC{\LR{$!C_1 \land !C_2$}}
  \RightLabel{\Elim{\land}[i]}
  \UnaryInfC{\LR{$!C_i$}}
  \DischargeRule{\LR{$\Elim{\lor}$}}{x_1,x_2}
  \TrinaryInfC{\LR{$!C_i$}}
  \DisplayProof
\end{multline*}
که \LR{$M$} د هغه بشپړ يا-ثبوت ثبوت ترم وي چې
د~\LR{$\delta$} پسې د يا د داخلولو ګام هم پکې شامل
وي، او \LR{$N_1$} او \LR{$N_2$} په ترتيب د \LR{$\delta_1$} او
\LR{$\delta_2$} ثبوت ترمونه وي، نو د ثبوت ترمونو،
يعنې \LR{$\lambd$} ترمونو، اړوند راکمول قاعده دا ده:
\[
  \proj{i}{\dcase{M}{x_1}{N_1}{x_2}{N_2}} \redone 
\dcase{M}{x_1}{\proj{i}{N_1}}{x_2}{\proj{i}{N_2}}
\]
\olpseganchor{OLP-0691-B011}{end}

\olpseganchor{OLP-0691-B012}{start}
تاو راتاو يوازې د ثبوتونه له پاې څخه نۀ،
بلکې د ثبوتونه له منځه هم لرې کولے شو: په
داسې فرعي-ثبوت د راکمول بدلون عملي کوو چې
په تاو راتاو پاې ته رسېږي۔ همدغه شان
\LR{$\beta$}-راکمول د \LR{$\lambd$}-ترمونو په فرعي ترمونو
هم عملي کېږي۔ که \LR{$N$} د \LR{$M$} فرعي ترم وي او
\LR{$N \redone N'$} وي، نو په~\LR{$M$} کښې فرعي ترم \LR{$N$}
په \LR{$N'$} بدلولے شو، څو \LR{$M'$} ترلاسه کړو۔ په دې
حالت کښې \LR{$M \redone M'$} هم ليکو۔ که داسې لړۍ
\LR{$M = M_1 \redone M_2 \redone \dots \redone M_k = M'$}
وي (د \LR{$k=1$}، يعنې \LR{$M = M'$} حالت هم پکې شامل دے)،
نو \LR{$M \red M'$} ليکو۔
\olpseganchor{OLP-0691-B012}{end}

\olpseganchor{OLP-0691-B013}{start}
هغه ثبوت چې په هېڅ فرعي-ثبوت کښې
ئې کوم بدلون نۀ شي عملي کېدے، \emph{عادي} بلل
کېږي۔ په همدې ډول هغه \LR{$\lambd$}-ترم چې په هېڅ
فرعي ترم کښې ئې کوم بدلون نۀ شي عملي کېدے،
\emph{عادي} دے۔ عادي \LR{$\lambd$}-ترم ته
\emph{عادي بڼه} هم وايي۔
\olpseganchor{OLP-0691-B013}{end}

\olpseganchor{OLP-0696-B004}{start}
\olfileid{int}{pty}{tpr}
\olpseganchor{OLP-0696-B004}{end}

\olpseganchor{OLP-0696-B005}{start}
\olsection{د ټايپ ساتنه}
\olpseganchor{OLP-0696-B005}{end}

\olpseganchor{OLP-0696-B006}{start}
هغه راکموونکي او ترتيبي بدلونونه چې ومو کتل، د
\Log{tN2i} او \Log{tN3i} نظامونو دپاره هم نېغ په نېغه
تعريفېدلے شي۔ د \Log{tN3i} نظام د سياق~\LR{$\Gamma$} په
نسبت د \LR{$\lambd$}-ترم~\LR{$N$} ټايپ هغه فارمول~\LR{$!A$}
تعريفوي چې \LR{$\Gamma \Sequent N : !A$} په \Log{tN3i}
کښې ثبوت ولري۔ دا چې د \LR{$\lambd$}-ترمونو د
\LR{$\beta$}-راکمول قواعد د ثبوتونه راکموونکي او ترتيبي
بدلونونه کټ مټ منعکسوي، يوه مهمه پايله لري: د سياق~\LR{$\Gamma$}
په نسبت د \LR{$\lambd$}-ترم ټايپ له \LR{$\beta$}-بدلون وروسته
هم هماغه پاتې کېږي۔ دې خاصيت ته \emph{د ټايپ ساتنه} وايي۔
\olpseganchor{OLP-0696-B006}{end}

\olpseganchor{OLP-0696-B007}{start}
\begin{prop}
که \LR{$N$} د \LR{$\Gamma$} په نسبت ټايپ~\LR{$!A$} ولري او
\LR{$N \redone N'$} وي، نو \LR{$N'$} هم د~\LR{$\Gamma$} په نسبت
ټايپ~\LR{$!A$} لري۔
\end{prop}
\olpseganchor{OLP-0696-B007}{end}

\olpseganchor{OLP-0696-B008}{start}
\begin{proof}
  پر~\LR{$N$} او د \LR{$\Gamma \Sequent N : !A$} د ثبوتونه~\LR{$\delta$}
  پر لوړوالي په استقرا دا خبره په اسانه ثابتولے شو: که \LR{$M$}
  د~\LR{$N$} يو فرعي ترم وي، نو \LR{$\delta$} يو فرعي ثبوت~\LR{$\delta_1$}
  لري چې په \LR{$\Gamma' \Sequent M : !B$} پاې ته رسېږي۔
  که \LR{$M$} رېډکس وي، نو پر~\LR{$\delta_1$} يو بدلون پلي کېږي۔
  پر \LR{$\delta_1$} د راکمول په پلي کولو د \LR{$\Gamma' \Sequent M' : !B$}
  ثبوت~\LR{$\delta_1'$} ترلاسه کوو۔ د \Log{tN3i} د بدلونونو
  له کتنې څرګندېږي چې \LR{$M'$} د~\LR{$M$} راکم شوے ترم دے۔
  په~\LR{$\delta$} کښې د \LR{$\delta_1$} پر ځای د \LR{$\delta_1'$} په ايښودلو
  د \LR{$\Gamma \Sequent N' : !A$} ثبوت~\LR{$\delta'$} جوړېږي؛
  دلته \LR{$N'$} په~\LR{$N$} کښې د فرعي ترم~\LR{$M$} پر ځای د~\LR{$M'$}
  په ايښودلو ترلاسه کېږي۔
\olpseganchor{OLP-0696-B008}{end}

\olpseganchor{OLP-0696-B009}{start}
  مثلاً د~\Intro{\land} هغه راکموونکے بدلون وګورئ
  چې ورپسې \Elim{\land} راځي:
\[
  \bottomAlignProof
  \AxiomC{}
  \Deduce$\Gamma \fCenter N_1 : !B_1$
  \AxiomC{}
  \Deduce$\Gamma \fCenter N_2 : !B_2$
  \RightLabel{\Intro{\land}}
  \BinaryInf$\Gamma \fCenter \pair{N_1}{N_2} : !B_1 \land !B_2$
  \RightLabel{\Elim{\land}[i]}
  \UnaryInf$\Gamma \fCenter \proj{i}{\pair{N_1}{N_2}} : !B_i$
  \DisplayProof
  \quad\redone\quad
  \bottomAlignProof
  \AxiomC{}
  \Deduce$\Gamma \fCenter N_i : !B_i$
  \DisplayProof
  \]
  وينو چې د ښي‌لاس ثبوت ثبوت ترم، يعنې
  \LR{$\delta_1'$}، \LR{$N_i$} دے۔ دا کټ مټ هغه پايله ده چې د
  کيڼ‌لاس ثبوت~(\LR{$\delta_1$}) پر ثبوت ترم
  \LR{$\proj{i}{\pair{N_1}{N_2}}$} د \LR{$\beta$}-بدلون له
  پلي کولو ترلاسه کېږي۔
\olpseganchor{OLP-0696-B009}{end}

\olpseganchor{OLP-0696-B010}{start}
  يا د \Intro{\lif} هغه استنتاج وګورئ چې ورپسې
  د \Elim{\lif} استنتاج راځي، او دا پرې پلي کېدونکے
  راکموونکے بدلون:
\begin{multline*}
  \bottomAlignProof
  \AxiomC{\LR{$\delta_2$}}
  \Deduce$x: !B_1, \Gamma \fCenter N : !B_2$
  \RightLabel{\Intro{\lif}}
  \UnaryInf$\Gamma \fCenter \lambd[\typeof{x}{!B_1}][N] : !B_1 \lif !B_2$
  \AxiomC{}
  \RightLabel{\LR{$\delta_3$}}
  \Deduce$\Gamma \fCenter M : !B_1$
  \RightLabel{\Elim{\lif}}
  \BinaryInf$\Gamma \fCenter (\lambd[\typeof{x}{!B_1}][N] M) : !B_2$
  \DisplayProof
  \quad\redone\\
  \bottomAlignProof
  \AxiomC{}
  \RightLabel{\LR{$\delta_3$}}
  \Deduce$\Gamma \fCenter M : !B_1$
  \RightLabel{\LR{$\Subst{\delta_2}{\delta_3}{x:!B_1}$}}
  \Deduce$\Gamma \fCenter \Subst{N}{M}{x} : !B_2$
  \DisplayProof
  \end{multline*}
  يو ځل بيا د ښي‌لاس ثبوت ثبوت ترم کټ مټ د
  کيڼ‌لاس ثبوت د ثبوت ترم د \LR{$\beta$}-راکمول پايله ده۔
  البته، دا خبره به په احتياط سره، د~\LR{$\delta_2$} پر لوړوالي
  په استقرا، کتل پکار وي: که په~\LR{$\delta_2$} کښې د
  \LR{$\Gamma' \Sequent x:!B_1$} په بڼه د ټولو اکسيومونو
  پر ځای د \LR{$\Gamma \Sequent M:!B_1$} ثبوت~\LR{$\delta_3$}
  واچوو، نو په رښتيا د \LR{$\Gamma \Sequent \Subst{N}{M}{x} : !B_2$}
  ثبوت \LR{$\Subst{\delta_2}{\delta_3}{x:!B_1}$} جوړېږي۔
\end{proof}
\olpseganchor{OLP-0696-B010}{end}

\olpseganchor{OLP-0696-B011}{start}
د ثبوتونه او \LR{$\lambd$}-ترمونو ترمنځ ټينګه مطابقت
بله پايله هم لري۔ که \LR{$N$} د سياق~\LR{$\Gamma$} په نسبت
د ټايپ \LR{$!A$} يو \LR{$\lambd$}-ترم وي او رېډکس~\LR{$M$} ولري
ـ يعنې عادي بڼه نۀ وي ـ نو د \LR{$\Gamma \Sequent N : !A$}
اړوند \Log{tN3i} ثبوت~\LR{$\delta$} يو فرعي ثبوت
لري چې په \LR{$\Gamma' \Sequent M: !B$} پاې ته رسېږي
او بدلون پرې پلي کېدلے شي۔ که \LR{$N \redone N'$} د
\LR{$M$} پر ځای د هغه د راکم شوي ترم په ايښودلو وي، نو
پر~\LR{$\delta$} د راکمول د پلي کولو اړونده پايله د
\LR{$\Gamma \Sequent N' : !A$} ثبوت ده۔ نو هر
\LR{$\lambd$}-ترم چې په عادي بڼه کښې نۀ وي، په \LR{$\beta$}
يو بل ترم ته راکمېږي۔ دې خاصيت ته \emph{پرمختګ} وايي۔
\olpseganchor{OLP-0696-B011}{end}

\olpseganchor{OLP-0696-B012}{start}
پرمختګ او د ټايپ ساتنه په ګډه دا خبره باوري کوي چې
د \LR{$\lambd$}-ترمونو \LR{$\beta$}-راکمول هېڅکله ``نۀ بندېږي'':
که يو ترم صحيح ټايپ ولري او په عادي بڼه کښې نۀ وي، نو
بل صحيح ټايپ لرونکي ترم ته راکمېږي، او د هغه ترم ټايپ
هم هماغه وي۔
\olpseganchor{OLP-0696-B012}{end}

\olpseganchor{OLP-0687-B004}{start}
\olfileid{int}{pty}{nor}
\olpseganchor{OLP-0687-B004}{end}

\olpseganchor{OLP-0687-B005}{start}
\olsection{عادي کول}
\olpseganchor{OLP-0687-B005}{end}

\olpseganchor{OLP-0687-B006}{start}
په دې برخه کښې ښيو چې د راکمولو د يوه ترتيب په وسيله هر
صحيح ټايپ‌لرونکے ثبوت عادي ثبوت ته راکمېږي؛ دې ته
\emph{د عادي کولو خاصيت} وايي۔ د بيانونو-په-ټايپونو سمون
کاروو: ښيو چې هر صحيح ټايپ‌لرونکے ثبوت ترم عادي بڼې ته
راکمېږي؛ بيا د فطري استنتاج د اشتقاقونه عادي کول ترې
راځي۔ دلته د درې بېټا قاعدو او د پروجکشن او يا د ترتيب د
بدلون پر هماغه ورکړې قاعدې بحث کوو۔
\olpseganchor{OLP-0687-B006}{end}

\olpseganchor{OLP-0687-B007}{start}
لومړے د ترمونو د پېچلتيا د اندازې دپاره ځينې تابعې تعريفوو۔
د فارمولونه \emph{اوږدوالے}~\LR{$\len{!A}$} داسې تعريفېږي:
\begin{align*}
  \len{p} &= 0\\
  \len{\lfalse} &= 0\\
  \len{!A \land !B} &= \len{!A} + \len{!B} + 1\\
  \len{!A \lor !B} &= \len{!A} + \len{!B} + 1 \\
  \len{!A \lif !B} &= \len{!A} + \len{!B} + 1.
\end{align*}
د رېډکس~\LR{$M$} پېچلتيا د هغه په \emph{پرېکون رتبه}
~\LR{$\cutrank{M}$} اندازه کوو:
\begin{align*}
  \cutrank{(\lambd[\typeof{x}{!A}][\typeof{N}{!B}])Q} &=
  \len{!A} + \len{!B} + 1 \\
  \cutrank{\ande{i}{\andi{\typeof{M}{!A}}{\typeof{N}{!B}}}} &=
  \len{!A} + \len{!B} + 1 \\
  \cutrank{\ore{\ori{i}{!A_{i}}{\typeof{M}{!A_i}}}
    {\typeof{x_1}{!A_1}}{\typeof{N_1}{!C}}
    {\typeof{x_2}{!A_2}}{\typeof{N_2}{!C}}} &=
  \len{!A_1} + \len{!A_2} + 1
\end{align*}
د ثبوت ترم پېچلتيا په هغه کښې د تر ټولو پېچلي رېډکس رتبه ده؛
که عادي وي، \LR{$0$} ده:
\[
  \maxrank{M} = \max (\{0\} \cup \{\cutrank{N} | N \text{\RL{ د }} M \text{\RL{
    فرعي ترم او رېډکس دے}}\})
\]
په اعظمي اندازه کښې صفر د بې‌رېډکسه ترم ارزښت ټاکي۔ د
پرېکون رتبه دلته د ټايپ په اوږدوالي اندازه کېږي۔
\olpseganchor{OLP-0687-B007}{end}

\olpseganchor{OLP-0687-B008}{start}
\begin{lem}\ollabel{lem:subst}
  که \LR{$\Subst{M}{\typeof{N}{!A}}{\typeof{x}{!A}}$} رېډکس وي او
  \LR{$M \not\ident x$} وي، نو له لاندې حالاتو څخه يو پېښېږي:
  \begin{enumerate}
  \item خپله \LR{$M$} رېډکس دے، يا
  \item \LR{$M$} د \LR{$\ande{i}{x}$} په بڼه او \LR{$N$} د
    \LR{$\andi{P_1}{P_2}$} په بڼه دے۔
  \item \LR{$M$} د \LR{$\ore{x}{x_1}{P_1}{x_2}{P_2}$} په بڼه او \LR{$N$}
    د \LR{$\ori{i}{}{Q}$} په بڼه دے۔
  \item \LR{$M$} د \LR{$x Q$} په بڼه او \LR{$N$} د
    \LR{$\lambd[x][P]$} په بڼه دے۔
  \end{enumerate}
  په لومړي حالت کښې \LR{$\cutrank{\Subst{M}{N}{x}} = \cutrank{M}$}؛
  په نورو حالاتو کښې \LR{$\cutrank{\Subst{M}{N}{x}} = \len{!A}$}۔
\end{lem}
\begin{proof}
  پر~\LR{$M$} د استقرا ثبوت دے۔
  \begin{enumerate}
  \item که \LR{$M$} يوازې متغير \LR{$y$} وي او \LR{$y \not\ident x$} وي، نو
    \LR{$\Subst{y}{N}{x}$} هم \LR{$y$} دے، نو رېډکس نۀ دے۔
  \item که \LR{$M$} د \LR{$\andi{N_1}{N_2}$}، \LR{$\lambd[x][N]$} يا
    \LR{$\ori{i}{!A}{N}$} په بڼه وي، نو
    \LR{$\Subst{M}{\typeof{N}{!A}}{\typeof{x}{!A}}$} هم په هماغه
    بڼه دے، نو رېډکس نۀ دے۔
  \item که \LR{$M$} د \LR{$\ande{i}{P}$} په بڼه وي، دوه حالتونه ګورو۔
    \begin{enumerate}
      \item که \LR{$P$} د \LR{$\andi{P_1}{P_2}$} په بڼه وي، نو
        \LR{$M \ident \ande{i}{\andi{P_1}{P_2}}$} رېډکس دے؛ په
        څرګند ډول
        \[
        \Subst{M}{N}{x} \ident
        \ande{i}{\andi{\Subst{P_1}{N}{x}}{\Subst{P_2}{N}{x}}}
        \]
        هم رېډکس دے۔ د پرېکون رتبې برابرې دي۔
      \item که \LR{$P$} يوازې متغير وي، بايد \LR{$x$} وي، څو تعويض
        رېډکس شي؛ او \LR{$N$} بايد د \LR{$\andi{P_1}{P_2}$} په بڼه وي۔
        اوس
        \[
        \Subst{M}{N}{x} \ident \Subst{\ande{i}{x}}{\andi{P_1}{P_2}}{x},
        \]
        وګورئ؛ دا \LR{$\ande{i}{\andi{P_1}{P_2}}$} دے۔ د پرېکون
        رتبه ئې د تعويض کېدونکي جوړې د ټايپ اوږدوالے، يعنې
        \LR{$\len{!A}$}، دے۔
    \end{enumerate}
  \end{enumerate}
  د \LR{$\ore{N}{x_1}{N_1}{x_2}{N_2}$} او \LR{$P Q$} حالتونه ورته دي۔
\end{proof}
\olpseganchor{OLP-0687-B008}{end}

\olpseganchor{OLP-0687-B009}{start}
\begin{lem}
  که \LR{$M$} ترم \LR{$M'$} ته انقباض وکړي، او د \LR{$M$} د هر مناسب
  رېډکس فرعي ترم \LR{$N$} دپاره \LR{$\cutrank{M} > \cutrank{N}$} وي،
  نو \LR{$\cutrank{M} > \maxrank{M'}$}۔
\end{lem}
\olpseganchor{OLP-0687-B009}{end}

\olpseganchor{OLP-0687-B010}{start}
\begin{proof}
  د حالاتو له مخې ثبوت دے۔
  \begin{enumerate}
  \item که \LR{$M$} د \LR{$\ande{i}{\andi{M_1}{M_2}}$} په بڼه وي، نو
    \LR{$M'$}، \LR{$M_i$} دے؛ د \LR{$M_i$} هر فرعي ترم د \LR{$M$} مناسب
    فرعي ترم هم دے، نو دعوه صحيح ده۔
  \item که \LR{$M$} د
    \LR{$(\lambd[\typeof{x}{!A}][N])\typeof{Q}{!A}$} په بڼه وي، نو
    \LR{$M'$}، \LR{$\Subst{N}{\typeof{Q}{!A}}{\typeof{x}{!A}}$} دے۔ په
    \LR{$M'$} کښې يو رېډکس وګورئ۔ يا په \LR{$N$} کښې ورته رېډکس
    د برابرې پرېکون رتبې سره شته، چې د فرض له مخې له
    \LR{$\cutrank{M}$} څخه کمه ده؛ يا رتبه \LR{$\len{!A}$} ده، چې
    د تعريف له مخې له \LR{$\cutrank{(\lambd[x^{!A}][N])Q}$} څخه
    کمه ده۔ د آرګومېنټ په نقل شوو برخو کښې د پخوا شته
    رېډکسونو حالت هم شته: هغوی د اصلي ټول رېډکس مناسب
    فرعي رېډکسونه وو، نو د فرض له مخې ئې رتبې ټيټې دي۔
  \item که \LR{$M$} د
    \[
    \ore{\ori{i}{}{\typeof{N}{!A_i}}}
        {\typeof{x_1}{!A_1}}{\typeof{N_1}{!C}}
        {\typeof{x_2}{!A_2}}{\typeof{N_2}{!C}},
    \]
    په بڼه وي، نو \LR{$M' \ident \Subst{N_i}{N}{\typeof{x_i}{!A_i}}$}۔
    په~\LR{$M'$} کښې يو رېډکس وګورئ۔ يا په \LR{$N_i$} کښې ورته
    رېډکس د برابرې پرېکون رتبې سره شته، چې د فرض له مخې
    له \LR{$\cutrank{M}$} څخه کمه ده؛ يا رتبه \LR{$\len{!A_i}$} ده،
    چې د تعريف له مخې له
    \LR{$\cutrank{\ore{\ori{i}{}{\typeof{N}{!A_i}}}
      {\typeof{x_1}{!A_1}}{\typeof{N_1}{!C}}
      {\typeof{x_2}{!A_2}}{\typeof{N_2}{!C}}}$} څخه کمه ده۔
    د ورننوتونکي آرګومېنټ په نقلونو کښې پخواني رېډکسونه هم
    د اصلي رېډکس مناسب فرعي ترمونه دي، نو د فرض له مخې
    ټيټې رتبې لري۔
  \end{enumerate}
\end{proof}
\olpseganchor{OLP-0687-B010}{end}

\olpseganchor{OLP-0687-B011}{start}
\begin{thm}
  ټول صحيح ټايپ‌لرونکي ثبوت ترمونه عادي بڼې ته راکمېږي؛
  د دې باب په اړوند محاسبوي معنا ټول اشتقاقونه عادي
  اشتقاقونه ته راکمېږي۔
\end{thm}
\olpseganchor{OLP-0687-B011}{end}

\olpseganchor{OLP-0687-B012}{start}
\begin{proof}
  دويمه پايله له لومړۍ راځي۔ د لومړۍ دپاره د
  \LR{$m=\maxrank{M}$} پر اندازه بشپړه استقرا کوو، چې \LR{$M$}
  صحيح ټايپ‌لرونکے ثبوت ترم دے۔
  \begin{enumerate}
\olpseganchor{OLP-0687-B012}{end}

\olpseganchor{OLP-0687-B013}{start}
  \item که \LR{$m=0$} وي، \LR{$M$} له مخکې عادي دے۔
  \item
    که نه وي، پر~\LR{$n$} استقرا کوو؛ دا په \LR{$M$} کښې د هغو
    رېډکسونو شمېر دے چې د پرېکون رتبه ئې~\LR{$m$} ده۔
    \begin{enumerate}
    \item که \LR{$n=1$} وي، داسې هر رېډکس \LR{$N$} غوره کړئ چې
      \LR{$m = \cutrank{N} > \cutrank{P}$} د هر مناسب فرعي ترم
      \LR{$P$} دپاره چې رېډکس هم وي، صحيح وي۔ داسې رېډکس شته،
      ځکه هر ترم يوازې متناهي شمېر فرعي ترمونه لري۔
\olpseganchor{OLP-0687-B013}{end}

\olpseganchor{OLP-0687-B014}{start}
      د \LR{$N$} راکمې شوې پايلې ته \LR{$N'$} وايو۔ د لمې له مخې
      \LR{$\maxrank{N'} < \maxrank{N}$}۔ د دې تګلارې مطلوب ګام دا دے چې
      \LR{$n$}، يعنې د \LR{$\cutrank=m$} لرونکو رېډکسونو شمېر، کمېږي۔
      له دې سره \LR{$m$} هم (په \LR{$1$} يا زيات) کمېږي، او د~\LR{$m$}
      استقرايي فرضيه کارولے شو۔
    \item د استقرا په ګام کښې \LR{$n>1$} فرض کړئ۔ بهير ورته
      دے، خو \LR{$n$} يوازې يوه مثبته شمېره ته کمېږي، نو \LR{$m$}
      نۀ بدلېږي۔ د~\LR{$n$} استقرايي فرضيه کاروو۔
    \end{enumerate}
\end{enumerate}
دغه استقرا هغه وخت بشپړېږي چې د ټول ترم اندازه کمه شي۔
د رېډکس په دننه کښې د رتبې کمېدل دغه شرط نۀ تضمينوي۔
د يا د ايستلو پايله کېدے شي له خپله يا ټايپ څخه اوږد
تابعي ټايپ ولري۔ که دواړه څانګې تجريدونه وي او ټول ترم
د تابع په توګه تطبيق شي، د يا د رېډکس له کمولو وروسته
په باندې تطبيق کښې نوی رېډکس پيدا کېږي۔ د هغه د پرېکون
رتبه د تابع اوږد ټايپ ټاکي۔ نو د ټول ترم اعظمي رتبه
په دغه ګام کښې لوړېدے شي۔ لاندې قوي عادي کېدل د عادي
بڼې شتون هم ثابتوي۔
\end{proof}
\olpseganchor{OLP-0687-B014}{end}

\olpseganchor{OLP-0687-B015}{start}
د ټايپ‌لرونکو \LR{$\lambd$}-ترمونو عادي کېدل وايي چې هر داسې
ترم \emph{عادي بڼه لري}۔ که د \LR{$\lambd$}-ترمونو عادي بڼې
ته \LR{$\beta$}-راکمول د محاسبې اجرا او عادي بڼه ئې د محاسبې
ارزښت وګڼو، نو په ټايپ‌لرونکي \LR{$\lambd$}-حساب کښې د ټاکلي مناسب ترتيب په وسيله محاسبه درېږي
او ارزښت ورکوي۔ د ټايپ ساتنه هم د ارزښت صحيح ټايپ
تضمينوي۔ مثلاً که \LR{$N$} ټايپ \LR{$!A \lif !B$} ولري (يعنې د
ټايپ~\LR{$!A$} آرګومېنټونو دپاره تابع ده او د ټايپ~\LR{$!B$}
ارزښتونه ورکوي)، او په د ټايپ~\LR{$!A$} ترم~\LR{$M$} (داخلې)
ئې تطبيق کړو، نو \LR{$(NM)$} ترم \LR{$N'$} (وت) ته راکمېږي او
دغه وت په يقيني ډول ټايپ~\LR{$!B$} لري۔
\olpseganchor{OLP-0687-B015}{end}

\olpseganchor{OLP-0687-B016}{start}
سرچينه يو بل قوي خاصيت هم بيانوي: په ټايپ‌لرونکي
\LR{$\lambd$}-حساب کښې ترمونه يوازې پر فرعي ترمونو د
\LR{$\beta$}-راکمول د يوه ترتيب له مخې عادي بڼې ته نۀ رسېږي؛
بلکې د \LR{$\beta$}-راکمول \emph{هر} ترتيب بالاخره په عادي
بڼه کښې درېږي۔ دغه خاصيت \emph{قوي عادي کول} بلل کېږي۔
دا وايي چې يوازې د محاسبې د درېدنې او ارزښت د ورکولو
يوه تضمين شوې لاره نشته، بلکې د محاسبې \emph{هره}
اجرا بالاخره ارزښت ورکوي۔ د دې پايلې لپاره د
\href{https://members.loria.fr/PdeGroote/papers/tlca01.pdf}{د دي ګروټ د ۲۰۰۱ کال د مقالې دويمه برخه او اتمه قضيه} وګورئ۔
د دې کتاب ترمونه د هغه حساب ته داسې اړوو: د متغيرونو،
تجريدونو، تطبيق، جوړو، پروجکشن او د يا د دواړو جوړوونکو
او ايستلو جوړښتونه ساتو او ټايپي ليبلونه د ټايپ ټاکنې
په ثبوت کښې کاروو۔ له دروغو د وتلو ترم په داسې تش‌تړونکي
ميو تجريد اړوو چې نوی تړونکے نوم ئې په بدن کښې نۀ راځي۔
په دې ډول ټاکلي ټايپ ته د دروغو له بدن څخه رسېږو او د
پايلې د ټايپ د ثبوت په سياق کښې نوے پرانيستے نوم نۀ پاتې
کېږي۔ د جوړو، تابعو او د يا هر بېټا ګام په اړول شوي ترم
کښې يو د تاو راتاو د لرې کولو ګام دے۔ د يا له ايستلو سره
د پروجکشن د ترتيب بدلول هم په هغه حساب کښې يو ګام دے۔
دا اړول په فرعي ترمونو کښې همدا ګامونه ساتي او تازه نومونه
د متغير له نيولو مخه نيسي۔ که په دې کتاب کښې له يوه
صحيح ټايپ‌لرونکي ترم څخه بې‌پايه راکمونه شونې وای، په
هغه حساب کښې به هم بې‌پايه راکمونه شونې وه، چې د يادې
قضيې خلاف ده۔ له همدې امله هر بشپړ راکمېدونکے لړ په
عادي بڼه کښې پای ته رسېږي۔
\olpseganchor{OLP-0687-B016}{end}

\olpseganchor{OLP-0687-B017}{start}
څنګه پوهېږو چې د محاسبې د اجرا بېلابېل ترتيبونه
\emph{يو شان} ارزښت ورکوي؟ تر اوسه د ټايپ ساتنه يوازې
وايي چې ټول ورکړي ارزښتونه يو ټايپ لري۔ ټايپ‌لرونکے
\LR{$\lambd$}-حساب يو بل مهم خاصيت لري: دا \emph{هم‌پايلے}
يا \emph{چرچ--روسر} دے۔ که \LR{$N \red N_1$} او \LR{$N \red N_2$}
وي، نو داسې ترم~\LR{$N'$} شته چې \LR{$N_1 \red N'$} او \LR{$N_2 \red N'$}
وي۔ په ياد ولرئ چې \LR{$N \red N'$} معنا لري چې \LR{$N'$} له پيلني ترم
څخه د صفر يا زياتو \LR{$\redone$}-راکمول ګامونو په وسيله راځي۔
که \LR{$N_1$} عادي بڼه وي، نو يوازې هغه ترم \LR{$N'$} چې
\LR{$N_1 \red N'$} ورته صحيح وي، خپله \LR{$N_1$} دے (د صفر
\LR{$\redone$}-راکمول ګامونو په وسيله)۔ نو که \LR{$N_1$} او \LR{$N_2$}
دواړه عادي بڼې وي، \LR{$N_1 = N' = N_2$}۔ په بل عبارت،
د \LR{$\red$} د قوي عادي کېدو له خاصيته هر ترتيب بالاخره
عادي بڼې ته رسېږي؛ او د \LR{$\red$} د چرچ--روسر له خاصيته
هرې دوې عادي بڼې برابرې دي۔ نو په ټايپ‌لرونکي
\LR{$\lambd$}-حساب کښې محاسبه تل درېږي او ارزښت (عادي بڼه)
ئې د اجرا له ترتيب څخه خپلواک دے۔ دلته دغه پايله د
يادو دواړو خاصيتونو پر رښتينوالي ولاړه ده؛ د مخکني
نيمګړي دليل په وسيله د دواړو بشپړ ثبوت نۀ ادعا کېږي۔
\olpseganchor{OLP-0687-B017}{end}

\olpseganchor{OLP-0687-B018}{start}
د يوې اړيکې د هم‌پايلتوب ثبوت عموماً ګران وي۔ خو
لاندې کمزورے شرط په اسانه ښودلے شو:
\olpseganchor{OLP-0687-B018}{end}

\olpseganchor{OLP-0687-B019}{start}
\begin{prop}[د چرچ--روسر کمزورے خاصيت]
که \LR{$N \redone N_1$} او \LR{$N \redone N_2$} وي، نو داسې ترم
\LR{$N'$} شته چې \LR{$N_1 \red N'$} او \LR{$N_2 \red N'$} وي۔
\end{prop}
\olpseganchor{OLP-0687-B019}{end}

\olpseganchor{OLP-0687-B020}{start}
\begin{proof}
  دوه ترمونه \LR{$N_1$} او \LR{$N_2$} په ترتيب سره د \LR{$N$} د رېډکس
  \LR{$M_1$} يا \LR{$M_2$} په راکمې شوې پايله بدلولو څخه راځي۔
  رېډکسونه \LR{$M_1$} او \LR{$M_2$} د \LR{$N$} فرعي ترمونه دي، نو د
  ترم په ونه کښې جزوي متقاطع نۀ دي: يا يو د بل فرعي
  ترم دے، يا د \LR{$N$} په بېلو نۀ پوښل کېدونکو ځایونو کښې
  واقع دي۔ وروستے حالت وګورئ: \LR{$N$} د \LR{$N(M_1,M_2)$} په
  بڼه دے۔ د \LR{$M_1$} پر ځاے د هغه راکمې پايلې \LR{$M_1'$}
  اېښودل \LR{$N_1$} ورکوي، چې \LR{$N(M_1',M_2)$} دے؛ او د
  \LR{$M_2$} پر ځاے د هغه راکمې پايلې \LR{$M_2'$} اېښودل \LR{$N_2$}
  ورکوي، چې \LR{$N(M_1,M_2')$} دے۔ دواړه \LR{$N(M_1',M_2')$}
  ته راکمېږي۔ د بل حالت دپاره فرض کړئ چې \LR{$M_2$} د~\LR{$M_1$}
  فرعي ترم دے (بل حالت متناظر دے)، او د \LR{$M_1$} د رېډکس
  د ډول له مخې حالتونه بېل کړئ۔ مثلاً که
  \LR{$M_1 \ident \proj{1}{\pair{O_1}{O_2}}$} وي، نو \LR{$O_1$}
  ته راکمېږي۔ که \LR{$M_2$} د \LR{$O_1$} فرعي ترم وي، د \LR{$M_2$}
  بدلول \LR{$O_1'$} ورکوي۔ نو لومړے د \LR{$M_1$} او ورپسې د
  \LR{$M_2$} بدلول \LR{$O_1'$} ورکوي۔ خو لومړے د \LR{$M_2$} بدلول
  \LR{$\proj{1}{\pair{O_1'}{O_2}}$} ورکوي، او دا هم~\LR{$O_1'$}
  ته راکمېږي۔ دا د يوې جوړې دننه د رېډکسونو د لګېدو حالت ښيي۔
  د ځايي هم‌پايلتوب لپاره د نورو جوړښتونو د نښتو حالتونه هم
  بايد وکتل شي۔
\end{proof}
\olpseganchor{OLP-0687-B020}{end}

\olpseganchor{OLP-0687-B021}{start}
\begin{lem}[د نيومن لمه]
که \LR{$\red$} قوي عادي کېدونکے وي او د چرچ--روسر کمزورے
خاصيت ولري، نو هم‌پايلے دے۔
\end{lem}
\olpseganchor{OLP-0687-B021}{end}

\olpseganchor{OLP-0687-B022}{start}
\begin{proof}
  څرنګه چې \LR{$\red$} قوي عادي کېدونکے دے، هره تر پايه
  وړل شوې د راکمولو لړۍ عادي بڼې ته رسېږي۔ د دې
  فرض لاندې عادي بڼې يکتا دي هله او يوازې هله چې
  \LR{$\red$} هم‌پايلے وي۔ فرض کړئ چې يو ترم \LR{$N$} دوې
  بېلې عادي بڼې \LR{$N'$} او \LR{$N''$} ولري، چې د لاندې
  راکمولو لړيو پايلې دي:
  \begin{align*}
  N & = N_1' \redone N_2' \redone \dots \redone N_k' = N' \text{\RL{ او}}\\
  N & = N_1'' \redone N_2'' \redone \dots \redone N_l'' = N''
  \end{align*}
  (د راکمولو لړۍ بايد اوږدوالے \LR{$>0$} ولري، ځکه که نه
  وي، \LR{$N$} به له مخکې عادي بڼه وه۔)
\olpseganchor{OLP-0687-B022}{end}

\olpseganchor{OLP-0687-B023}{start}
  که \LR{$N_2' = N_2''$} وي، همدغه سمدستي پايله دوې
  بېلې عادي بڼې (\LR{$N'$} او \LR{$N''$}) لري۔ که
  \LR{$N_2' \neq N_2''$} وي، د چرچ--روسر د کمزوري خاصيت
  له مخې داسې \LR{$N'''$} شته چې \LR{$N_2' \red N'''$} او
  \LR{$N_2'' \red N'''$} وي۔ دغه ګډه پايله د درېدنې په
  وسيله يوې عادي بڼې \LR{$K$} ته ورسوئ؛ نو
  \LR{$N_2' \red K$} او \LR{$N_2'' \red K$}۔ څرنګه چې
  \LR{$N' \neq N''$} دي، يا \LR{$K \neq N'$} يا \LR{$K \neq N''$}۔
  نو يا \LR{$N_2'$} يا \LR{$N_2''$} دوې بېلې عادي بڼې لري۔
  ومو ښودله چې که \LR{$N$} دوې بېلې عادي بڼې ولري، داسې
  \LR{$N_1$} شته چې \LR{$N \redone N_1$} وي او \LR{$N_1$} دوې
  بېلې عادي بڼې ولري۔ دغه ګام تکرارول
  \LR{$N \redone N_1 \redone N_2 \redone \dots$} ورکوي،
  خو دا ناممکن دے، ځکه \LR{$\redone$} قوي عادي کېدونکے دے۔
\end{proof}
\olpseganchor{OLP-0687-B023}{end}

\olpseganchor{OLP-0620-B004}{start}
\olpart{his}{تاريخ}
\olpseganchor{OLP-0620-B004}{end}

\olpseganchor{OLP-0621-B004}{start}
\olchapter{his}{bio}{ژوندليکونه}
\olpseganchor{OLP-0621-B004}{end}

\olpseganchor{OLP-0622-B004}{start}
\olfileid{his}{bio}{can}
\olpseganchor{OLP-0622-B004}{end}

\olpseganchor{OLP-0622-B005}{start}
\olsection{ګېورګ کانتور}
\olpseganchor{OLP-0622-B005}{end}

\olpseganchor{OLP-0622-B006}{start}
\olphoto{cantor-georg}{ګېورګ کانتور}
\olpseganchor{OLP-0622-B006}{end}

\olpseganchor{OLP-0622-B007}{start}
د ګېورګ کانتور، چې نوم ئې (\textsc{ګې}-اورګ
\textsc{کان}-تور) ويل کېږي، يوه پخواني ژوندليک
ادعا کړې وه چې هغه په يوه بېړۍ کښې زېږېدلے او
موندل شوے و، چې د روسيې سېنټ پېټرزبرګ ته روانه
وه، او مور او پلار ئې نۀ وو معلوم۔ خو دا رښتيا
نۀ دي؛ که څه هم هغه په 1845 کښې په سېنټ پېټرزبرګ
کښې زېږېدلے و۔
\olpseganchor{OLP-0622-B007}{end}

\olpseganchor{OLP-0622-B008}{start}
کانتور په 1867 کښې د برلين له پوهنتون څخه په
رياضياتو کښې دوکتورا ترلاسه کړه۔ هغه د سټونو
په تيورۍ کښې د خپل کار له امله پېژندل کېږي،
او د سټونو د تيورۍ د يوې بېلې څېړنيزې څانګې
په توګه بنسټ اېښودل ورته منسوبېږي۔ هغه لومړنے
کس و چې ثابت ئې کړل د بېلابېلو اندازو
نامتناهي سټونه شته۔ د هغه تيوريو، او په تېره
بيا د نامتناهياتو تيورۍ ئې، د هغه وخت د
رياضيپوهانو تر منځ ډېر بحث پيدا کړ، او کار
ئې د اختلاف وړ و۔
\olpseganchor{OLP-0622-B008}{end}

\olpseganchor{OLP-0622-B009}{start}
د کانتور مذهبي عقيدې او رياضيکي کار يو له بل
سره نۀ بېلېدونکي تړلي وو؛ آن هغه دا ادعا هم
کوله چې د ترانسفينيټ عددونو تيوري خداے نېغ
په نېغه ورته رسولې وه۔ د ژوند په وروستيو
کلونو کښې کانتور په ذهني ناروغۍ اخته و۔
له 1894 څخه پيل او په وروستيو کلونو کښې
ئې په زيات تکرار، کانتور روغتون ته داخلېدو۔
پر کار ئې سختو نيوکو، چې له رياضيپوه
لېوپولډ کرونېکر سره اختلاف هم پکښې و،
هغه ژور خپګان او له رياضياتو سره د شوق
نشتوالي ته ورساوۀ۔ د ژور خپګان په دورو
کښې کانتور فلسفې او ادب ته مخه کوله، او
آن يوه تيوري ئې هم خپره کړه چې فرانسس
بېکن د شېکسپير د ډرامو ليکوال و۔
\olpseganchor{OLP-0622-B009}{end}

\olpseganchor{OLP-0622-B010}{start}
کانتور د 1918 د جنورۍ پر 6مه په هالې
کښې په يوه درملنځاے کښې وفات شو۔
\olpseganchor{OLP-0622-B010}{end}

\olpseganchor{OLP-0622-B011}{start}
\begin{reading}
د کانتور د بشپړو ژوندليکونو دپاره
\mbox{\LR{\latinfont \citet{Dauben1990}}} او \mbox{\LR{\latinfont \citet{Grattan-Guinness1971}}}
وګورئ۔ د کانتور انقلابي نظرونه د بي بي سي
راډيو 4 په پروګرام \emph{د رياضياتو لنډ تاريخ}
کښې هم بيان شوي دي \mbox{\LR{\latinfont \citep{Sautoy2014}}}۔ که
غواړئ د کانتور د تيوريو په هکله د رېپ
په بڼه واورئ، \mbox{\LR{\latinfont \citet{Rose2012}}} وګورئ۔
\end{reading}
\olpseganchor{OLP-0622-B011}{end}

\olpseganchor{OLP-0623-B004}{start}
\olfileid{his}{bio}{chu}
\olpseganchor{OLP-0623-B004}{end}

\olpseganchor{OLP-0623-B005}{start}
\olsection{الونزو چرچ}
\olpseganchor{OLP-0623-B005}{end}

\olpseganchor{OLP-0623-B006}{start}
\olphoto{church-alonzo}{الونزو چرچ}
\olpseganchor{OLP-0623-B006}{end}

\olpseganchor{OLP-0623-B007}{start}
الونزو چرچ د 1903 د جون پر 14مه په واشنګټن، ډي سي
کښې زېږېدلے و۔ په لومړني ماشومتوب کښې د يو هوايي
ټوپک پېښې د چرچ يوه سترګه ړنده کړه۔ هغه په 1920
کښې په کنېکټيکټ کښې د پوهنتون د چمتووالي ښوونځے
بشپړ کړ او په هماغه کال ئې په پرنسټن کښې د
پوهنتون زده کړه پيل کړه۔ دوکتورا ئې په 1927
کښې بشپړه کړه۔ له دوو کلونو بهر اوسېدو وروسته
چرچ پرنسټن ته راوګرځېد۔ چرچ د ډېر ادب او
احتياط له امله پېژندل کېدو۔ پر تخته ليکل ئې
بې‌عيبه وو، او مهمې مقالې به ئې په ډېر احتياط
په ډوکو سېمنټ، چې يو روڼ سريښ دے، پوښلې او
ساتلې۔ له علمي کارونو بهر ئې د ساينسي خيالي ادب
مجلې په شوق لوستلې، او که په ليکنو کښې ئې
ناسمې خبرې وليدې، مديرانو ته ئې په ليکلو
کښې هېڅ وېره نۀ کوله۔
\olpseganchor{OLP-0623-B007}{end}

\olpseganchor{OLP-0623-B008}{start}
د چرچ علمي لاسته راوړنې لويې وې۔ له خپلو
زده کوونکو سټيفن کليني او بارکلي روسر سره
ئې په ګډه د اغېزمنې محاسبې يوه تيوري،
لامبډا حساب، جوړه کړه، د الَن ټيورينګ له
خوا د ټيورينګ ماشين له جوړولو څخه په
خپلواک ډول۔ د محاسبه کېدنې دا دوه
تعريفونه معادل دي، او هغه څه ترې
راولاړېږي چې اوس \emph{د چرچ--ټيورينګ اصل}
نومېږي: د طبيعي عددونو يوه تابع په
اغېزمنه توګه محاسبه کېدونکې ده هله او
يوازې هله چې د ټيورينګ ماشين، يا
لامبډا حساب، په ذريعه محاسبه کېدونکې
وي۔ هغه دا هم ثابت کړل چې اوس
\emph{د چرچ قضيه} نومېږي: د لومړي
ترتيب د فارمولونو د اعتبار د پرېکړې
مسئله نۀ شي حل کېداے۔
\olpseganchor{OLP-0623-B008}{end}

\olpseganchor{OLP-0623-B009}{start}
چرچ تر زړښت پورې خپل کار جاري وساتۀ۔
په 1967 کښې ئې پرنسټن د UCLA دپاره
پرېښود، او هلته تر خپل تقاعد، په
1990 کښې، پروفېسر و۔ چرچ د 1995
د اګست پر 11مه د 92 کلونو په عمر
وفات شو۔
\olpseganchor{OLP-0623-B009}{end}

\olpseganchor{OLP-0623-B010}{start}
\begin{reading}
د چرچ د لنډ ژوندليک دپاره
\mbox{\LR{\latinfont \citet{EndertonND}}} وګورئ۔ د لامبډا
حساب او د پرېکړې مسئلې، يعنې
Entscheidungsproblem، په هکله د
چرچ اصلي ليکنې
\mbox{\LR{\latinfont \citet{Church1936,Church1936a}}} دي۔
\mbox{\LR{\latinfont \citet{Aspray1984}}} له چرچ سره يوه
مرکه ثبتوي چې په 1930s کښې د
پرنسټن د رياضياتو د ټولنې په
هکله ده۔ چرچ له 1936 څخه تر
1979 پورې د \emph{د سمبوليک
منطق مجلې} د کتابونو د کتنو
يوه لړۍ وليکله۔ دا ټولې د
جان مکفارلېن پر وېبپاڼه
خوندي شوې دي \mbox{\LR{\latinfont \citep{MacFarlane2015}}}۔
\end{reading}
\olpseganchor{OLP-0623-B010}{end}

\olpseganchor{OLP-0624-B004}{start}
\olfileid{his}{bio}{gen}
\olpseganchor{OLP-0624-B004}{end}

\olpseganchor{OLP-0624-B005}{start}
\olsection{ګېرهارډ ګېنټسن}
\olpseganchor{OLP-0624-B005}{end}

\olpseganchor{OLP-0624-B006}{start}
\olphoto{gentzen-gerhard}{ګېرهارډ ګېنټسن}
\olpseganchor{OLP-0624-B006}{end}

\olpseganchor{OLP-0624-B007}{start}
ګېرهارډ ګېنټسن په عمده توګه د ثبوتونو د جوړښتي تيورۍ
د بنسټ ايښودونکي په توګه پېژندل کېږي، او په ځانګړي
ډول د طبيعي استنتاج او د سېکوېنټ حساب د
اشتقاق سيسټمونو د جوړولو له امله۔ هغه د
1909 د نومبر پر 24مه د جرمني په ګرايفسوالډ کښې
زېږېدلے و۔ ګېرهارډ درې کاله په کور کښې زده کړه
وکړه او بيا د پوهنتون د چمتووالي ښوونځي ته لاړ،
چې هلته د زده کړې له مخې له خپلو ډېرو ټولګيوالو
وروسته و۔ له دې سره سره هغه تکړه زده کوونکے و
او په رياضياتو کښې ئې پياوړې وړتيا څرګنده کړه۔
شوقونه ئې بېلابېل وو؛ د بېلګې په توګه، خپلې مور
ته ئې شعرونه او د ښوونځي د تياتر دپاره ئې
ډرامې هم ليکلې۔
\olpseganchor{OLP-0624-B007}{end}

\olpseganchor{OLP-0624-B008}{start}
ګېنټسن خپلې پوهنتوني زده کړې د ګرايفسوالډ په
پوهنتون کښې پيل کړې، خو بيا ګوټينګن، ميونيخ او
برلين ته لاړ۔ په 1933 کښې ئې د ګوټينګن له
پوهنتون څخه د هېرمان وايل تر لارښوونې لاندې
دوکتورا واخيسته۔ (پاول برنايس ئې د کار د زياتې
برخې لارښوونه کړې وه، خو نازيانو هغه له پوهنتونه
ګوښه کړ۔) په 1934 کښې ګېنټسن د ډېوېډ هېلبرټ د
مرستيال په توګه کار پيل کړ۔ په هماغه کال ئې په
خپلو ليکنو \emph{Untersuchungen \"{u}ber das logische Schlie\ss en I--II
  [د منطقي استنتاج څېړنې I--II]} کښې د سېکوېنټ
حساب او طبيعي استنتاج اشتقاق سيسټمونه
جوړ کړل۔ په 1936 کښې ئې د پيانو د بديهياتو
سازګاري ثابته کړه۔
\olpseganchor{OLP-0624-B008}{end}

\olpseganchor{OLP-0624-B009}{start}
له نازيانو سره د ګېنټسن اړيکه پېچلې ده۔ په
هماغه وخت کښې چې د هغه لارښوونکے برنايس د
جرمني پرېښودلو ته اړ کړے شو، ګېنټسن د نازي
نيمه‌پوځي سازمان SA په پوهنتوني څانګه کښې شامل
شو۔ د ډېرو جرمنيانو په شان هغه د نازي ګوند
غړے و۔ د جګړې پر وخت ئې د هوايي استخباراتو
د يوې څانګې د مخابراتي افسر په توګه کار کاوۀ۔
خو په 1942 کښې د عصبي ماتې له امله له دندې
خلاص کړے شو۔ دا څرګنده نۀ ده چې د ګېنټسن
وفاداري له نازي ګوند سره وه، که په ګوند کښې
د دې دپاره شامل شوے و چې علمي برياليتوب
ځان ته يقيني کړي۔
\olpseganchor{OLP-0624-B009}{end}

\olpseganchor{OLP-0624-B010}{start}
په 1943 کښې ګېنټسن ته د پراګ د جرمني پوهنتون
په رياضيکي انسټيټيوټ کښې علمي دنده وړاندې شوه،
چې هغه ومنله۔ خو په 1945 کښې د پراګ اوسېدونکو
د جرمني د اشغال پر ضد پاڅون وکړ۔ شوروي ځواکونه
ښار ته راورسېدل او د پوهنتون ټول پروفېسران ئې
ونيول۔ په نازي سازمانونو کښې د غړيتوب له امله
ګېنټسن د جبري کار يوه کېمپ ته بوتلل شو۔ هغه
په خپله حجره کښې د خوارځواکۍ له امله د 1945
د اګست پر 4مه د 35 کلونو په عمر وفات شو۔
\olpseganchor{OLP-0624-B010}{end}

\olpseganchor{OLP-0624-B011}{start}
\begin{reading}
د ګېنټسن د بشپړ ژوندليک دپاره
\mbox{\LR{\latinfont \citet{Menzler-Trott2007}}} وګورئ۔ د نازي واکمنۍ
لاندې د رياضيپوهانو په هکله يوه په زړه پورې
ليکنه، چې د ګېنټسن د ژوند لنډ يادښت هم ورکوي،
\mbox{\LR{\latinfont \citet{Segal2014}}} وړاندې کوي۔ د منطقي استنتاج
په هکله د ګېنټسن ليکنې په اصلي جرمني ژبه
شته \mbox{\LR{\latinfont \citep{Gentzen1935a,Gentzen1935b}}}۔ د ګېنټسن
د ليکنو انګرېزي ژباړې \mbox{\LR{\latinfont \citet{Gentzen1969}}} په
يو ټوک کښې راټولې کړې دي، چې د هغه د ژوند
يوه لنډه خاکه هم لري۔
\end{reading}
\olpseganchor{OLP-0624-B011}{end}

\olpseganchor{OLP-0625-B004}{start}
\olfileid{his}{bio}{god}
\olpseganchor{OLP-0625-B004}{end}

\olpseganchor{OLP-0625-B005}{start}
\olsection{کرټ ګوډل}
\olpseganchor{OLP-0625-B005}{end}

\olpseganchor{OLP-0625-B006}{start}
\olphoto{goedel-kurt}{کرټ ګوډل}
\olpseganchor{OLP-0625-B006}{end}

\olpseganchor{OLP-0625-B007}{start}
کرټ ګوډل (\textsc{ګېړ}-ډل) د 1906 د اپرېل پر
28مه د اتريش--هنګري په امپراتورۍ کښې په برون
کښې زېږېدلے و (چې اوس د چېک جمهوريت برنو دے)۔
د ځيرک او پوښتونکي طبيعت له امله به ئې کورنۍ
ځوان کرټلې ته ډېر ځله ``Der kleine Herr Warum''
(کوچنے ښاغلے ولې) ويل۔ هغه له لومړني ښوونځي
څخه په زده کړو کښې غوره و، او يوازې په
رياضياتو کښې ئې کله تر ټولو لوړې درجې کمه
درجه اخيسته۔ ګوډل د خرابې روغتيا له امله
ډېر ځله له ښوونځي غيرحاضر و او له بدني روزنې
څخه مستثنا و۔ په ماشومتوب کښې ئې د روماتيکي
تبې تشخيص شوے و۔ د ټول ژوند په اوږدو کښې ئې
دا باور درلود چې دې تبې ئې زړه د تل دپاره
اغېزمن کړے دے، که څه هم طبي ارزونې د دې
خلاف ويل۔
\olpseganchor{OLP-0625-B007}{end}

\olpseganchor{OLP-0625-B008}{start}
ګوډل په 1924 کښې د ويانا په پوهنتون کښې زده
کړه پيل کړه او په~1929 کښې ئې د دوکتورا زده
کړې بشپړې کړې۔ لومړے ئې د فزيک د زده کړې
اراده وه، خو شوق ئې ژر رياضياتو او په ځانګړي
ډول منطق ته واوښت، چې د فيلسوف روډولف کارناپ
اغېز هم پکې برخه لرله۔ د هانس هان تر
لارښوونې لاندې ليکلې رسالې ئې د عينيت لرونکي
لومړي ترتيب د پريديکات منطق د بشپړتيا قضيه
ثابته کړه \mbox{\LR{\latinfont \citep{Godel1929}}}۔ يوازې يو کال
وروسته ئې خپلې تر ټولو مشهورې نتيجې ترلاسه
کړې: د ناتکميلۍ لومړۍ او دويمه قضيه
(په \mbox{\LR{\latinfont \citealt{Godel1931}}} کښې خپرې شوې)۔ په
ويانا کښې د اوسېدو پر وخت ګوډل د ويانا له
کړۍ سره ډېر بوخت و، چې د علمي فکر لرونکو
فيلسوفانو يوه ډله وه او کارناپ هم پکې شامل
و۔ د کارناپ کار په ځانګړي ډول د ګوډل د
نتيجو تر اغېز لاندې و۔
\olpseganchor{OLP-0625-B008}{end}

\olpseganchor{OLP-0625-B009}{start}
په 1938 کښې ګوډل له اډېلې نيمبورسکي سره واده
وکړ۔ مور او پلار ئې خوشحاله نۀ وو: نۀ يوازې
دا چې هغه تر ګوډل شپږ کاله مشره وه او له
وړاندې ئې طلاق اخيستے و، بلکې په يوه شپني
کلب کښې ئې د نڅاګرې په توګه کار کاوۀ۔ خو
ټولنيزو فشارونو ګوډل اغېزمن نۀ کړ، او تر
وفاته ئې خوشحاله واده‌وال پاتې شول۔
\olpseganchor{OLP-0625-B009}{end}

\olpseganchor{OLP-0625-B010}{start}
وروسته له دې چې نازي جرمني په 1938 کښې
اتريش له ځان سره يوځاے کړ، ګوډل او اډېلې
متحده ايالاتو ته کډوال شول۔ هلته ئې د
نيو جرسي په پرنسټن کښې د لوړو څېړنو په
انسټيټيوټ کښې دنده واخيسته۔ له خپل درونګرا
او غيرمعمول طبيعت سره سره په پرنسټن کښې
د ګوډل وخت له ګډ کاره ډک او ګټور و۔ هغه
د سټونو په تيورۍ، فلسفه او فزيک کښې مقالې
خپرې کړې۔ په ځانګړي ډول ئې په IAS کښې
له خپل همکار البرټ اينشټاين سره ډېره
پياوړې ملګرتيا جوړه کړه۔
\olpseganchor{OLP-0625-B010}{end}

\olpseganchor{OLP-0625-B011}{start}
په وروستيو کلونو کښې د ګوډل ذهني روغتيا
خرابه شوه۔ په 1977 کښې د مېرمنې روغتون ته
تلل دا معنا لرله چې نور ئې هغه دپاره
خواړۀ نۀ شول پخولے۔ هغه چې د ټول ژوند
په اوږدو کښې له ذهني روغتيايي ستونزو سره
مخ و، په وهمي بدګمانۍ اخته شو۔ د زهر ورکولو
له سختې وېرې ګوډل له خوراکه انکار وکړ۔
هغه د لوږې له امله د 1978 د جنورۍ پر
14مه په پرنسټن کښې وفات شو۔
\olpseganchor{OLP-0625-B011}{end}

\olpseganchor{OLP-0625-B012}{start}
\begin{reading}
د ګوډل د ژوند د بشپړ ژوندليک دپاره
\mbox{\LR{\latinfont \citet{Dawson1997}}} وګورئ۔ د نورو ژوندليکي
ليکنو، او منطق او فلسفې ته د ګوډل د
مرستو په هکله د مقالو دپاره
\mbox{\LR{\latinfont \citet{Wang1990}}}، \mbox{\LR{\latinfont \citet{Baaz2011}}}،
\mbox{\LR{\latinfont \citet{Takeuti2003}}} او \mbox{\LR{\latinfont \citet{Sigmund2007}}}
وګورئ۔
\olpseganchor{OLP-0625-B012}{end}

\olpseganchor{OLP-0625-B013}{start}
د ګوډل د دوکتورا رساله په اصلي جرمني ژبه
شته \mbox{\LR{\latinfont \citep{Godel1929}}}۔ د ناتکميلۍ د قضيو
اصلي متن \mbox{\LR{\latinfont \citep{Godel1931}}} دے۔ د ګوډل
ټولې خپرې شوې او ناخپرې ليکنې، او د
ليکونو يوه ټاکلې برخه، په انګرېزي ژبه
د هغه په \emph{راټولو شويو ليکنو}
کښې شته \mbox{\LR{\latinfont \citet{Godel1986,Godel1990}}}۔
\olpseganchor{OLP-0625-B013}{end}

\olpseganchor{OLP-0625-B014}{start}
د ګوډل د ناتکميلۍ د قضيو د مفصل بيان
دپاره \mbox{\LR{\latinfont \citet{Smith2013}}} وګورئ۔ د ګوډل
د قضيو د غيرصوري، فلسفي بحث دپاره
د مارک لينسېنماير پوډکاسټ وګورئ
\mbox{\LR{\latinfont \citep{Linsenmayer2014}}}۔
\end{reading}
\olpseganchor{OLP-0625-B014}{end}

\olpseganchor{OLP-0626-B004}{start}
\olfileid{his}{bio}{noe}
\olpseganchor{OLP-0626-B004}{end}

\olpseganchor{OLP-0626-B005}{start}
\olsection{اېمي نوېتر}
\olpseganchor{OLP-0626-B005}{end}

\olpseganchor{OLP-0626-B006}{start}
\olphoto{noether-emmy}{اېمي نوېتر}
\olpseganchor{OLP-0626-B006}{end}

\olpseganchor{OLP-0626-B007}{start}
اېمي نوېتر (\textsc{نېړ}-ټر) د 1882 د مارچ پر
23مه د جرمني په اېرلانګن کښې د لوړې منځنۍ
طبقې په يوه علمي کورنۍ کښې زېږېدلې وه۔
نوېتر، چې «د نوي الجبر مور» بلل کېږي، د
ښځو د زده کړو له لويو خنډونو سره سره په
رياضياتو او فزيک دواړو کښې بنسټيزې نوې
مرستې وکړې۔ په هغه وخت کښې په جرمني کښې
ځوانو نجونو ته د هنرونو زده کړه ټاکل کېده
او د پوهنتون د چمتووالي ښوونځيو ته د
تللو اجازه ئې نۀ لرله۔ خو د ګوټينګن او
اېرلانګن په پوهنتونونو کښې د ټولګيو له
اورېدو وروسته، چې پلار ئې په اېرلانګن
کښې د رياضياتو پروفېسر و، نوېتر ته بالاخره
په 1904 کښې هلته د زده کوونکې په توګه
د نوم‌ليکنې شونتيا برابره شوه، کله چې د پوهنتون
تګلاره د ښځينه زده کوونکو د منلو دپاره
بدله شوه۔ هغې په 1907 کښې په رياضياتو
کښې دوکتورا واخيسته۔
\olpseganchor{OLP-0626-B007}{end}

\olpseganchor{OLP-0626-B008}{start}
له خپل اهليت سره سره نوېتر په خپل کاري
ژوند کښې له ډېر مخالفت سره مخ شوه۔ له
1908--1915 پورې ئې په اېرلانګن کښې بې
معاشه درس ورکاوۀ۔ په دې وخت کښې ئې د
ډېوېډ هېلبرټ پام ځان ته راواړوۀ، چې د
هغه وخت د نړۍ له مخکښو رياضيپوهانو څخه
و او هغه ئې ګوټينګن ته راوبلله۔ خو ښځو
ته د پروفېسر د مقام ترلاسه کول منع وو،
او هغې يوازې د هېلبرټ تر نامه لاندې
درس ورکولے شو، بيا هم بې معاشه۔ په دې
وخت کښې ئې هغه څه ثابت کړل چې اوس د
نوېتر قضيه نومېږي او نن هم په نظري
فزيک کښې کارېږي۔ نوېتر ته بالاخره په
1919 کښې د درس ورکولو حق ورکړے شو۔
د يوې مشهورې کيسې له مخې، د همکارانو د مخالفت په ځواب کښې هېلبرټ د پوهنځي مجلس له حمام سره پرتله کړ، څو څرګنده کړي چې د تدريس د حق په پرېکړه کښې د نوماندې جنسيت دليل نۀ شي کېدے۔
\olpseganchor{OLP-0626-B008}{end}

\olpseganchor{OLP-0626-B009}{start}
د 1920s په وروستيو کښې هغې په تجريدي الجبر
کار ته پام واړاوۀ، او مرستو ئې دا ډګر
واړوۀ۔ په خپلو ثبوتونو کښې ئې ډېر ځله
د لوړېدونکې لړۍ په نوم شرط کاراوۀ، چې
وايي د ځينو سټونو هېڅ نامتناهي سخته
زياتېدونکې لړۍ نشته۔ د بېلګې په توګه،
ځينې جبري جوړښتونه چې اوس د نوېتر
حلقې نومېږي، دا خاصيت لري چې د ايډيالونو
هېڅ نامتناهي لړۍ \LR{$I_1
\subsetneq I_2 \subsetneq \dots$} نشته۔
دا شرط هر جزوي ترتيب ته عامېدلے شي
(په الجبر کښې دا د فرعي سټ د اړيکې
له مخې د مرتب شوو ايډيالونو ځانګړے
حالت دے)، او موږ ئې دوه‌ګون، د
ښکته کېدونکې لړۍ شرط هم څېړلے شو،
چې په جزوي ترتيب کښې هره سخته
\emph{کمېدونکې} لړۍ بالاخره ختمېږي۔
که يو جزوي ترتيب د ښکته کېدونکې لړۍ
شرط پوره کوي، نو پر همدغه ترتيب
استقرا په هغه شان کارولے شو چې
د~\LR{$<$} ترتيب پر~\LR{$\Nat$} استقرا کاروو۔
داسې ترتيبونه \emph{ښه‌بنسټه} يا
\emph{د نوېتر ترتيبونه} نومېږي، او
اړوند ثبوتي اصل \emph{د نوېتر استقرا}
نومېږي۔
\olpseganchor{OLP-0626-B009}{end}

\olpseganchor{OLP-0626-B010}{start}
نوېتر يهودۍ وه، او کله چې نازيان په 1933
کښې واک ته ورسېدل، هغه له خپلې دندې
ګوښه شوه۔ له نېکه مرغه نوېتر په
پنسلوانيا کښې په برين ماور کښې د
يوې موقتي دندې دپاره متحده ايالاتو
ته کډواله کېدلے شوه۔ هلته د اوسېدو
پر وخت ئې په پرنسټن کښې هم درسونه
ورکول، که څه هم هغې وليدل چې پوهنتون
د ښځو هرکلے نۀ کاوۀ
\mbox{\LR{\latinfont \citep[81]{Dick1981}}}۔ په 1935 کښې
د نوېتر د رحم د تومور د لرې کولو
عمليات وشول۔ هغه د همدغو عملياتو
په پايله کښې د عفونت له امله وفات
شوه او په برين ماور کښې خاورو ته
وسپارل شوه۔
\olpseganchor{OLP-0626-B010}{end}

\olpseganchor{OLP-0626-B011}{start}
\begin{reading}
د نوېتر د ژوندليک دپاره \mbox{\LR{\latinfont \citet{Dick1981}}}
وګورئ۔ د نظري فزيک پېريميټر انسټيټيوټ
د نوېتر د ژوند او اغېز درسونه پر
انټرنېټ وړاندې کوي \mbox{\LR{\latinfont \citep{Perimeter2015}}}۔
که له لوستلو ستړي شوي يئ، \emph{هغه څه
چې د تاريخ په ټولګي کښې مو هېر کړي}
د نوېتر د ژوند او اغېز يو پوډکاسټ لري
\mbox{\LR{\latinfont \citep{Frey2015}}}۔ د نوېتر راټولې شوې
ليکنې په اصلي جرمني ژبه شته
\mbox{\LR{\latinfont \citep{Noether1983}}}۔
\end{reading}
\olpseganchor{OLP-0626-B011}{end}

\olpseganchor{OLP-0627-B004}{start}
\olfileid{his}{bio}{pet} 
\olpseganchor{OLP-0627-B004}{end}

\olpseganchor{OLP-0627-B005}{start}
\olsection{روژا پېتر}
\olpseganchor{OLP-0627-B005}{end}

\olpseganchor{OLP-0627-B006}{start}
\olphoto{peter-rozsa}{روژا پېتر}
\olpseganchor{OLP-0627-B006}{end}

\olpseganchor{OLP-0627-B007}{start}
روژا پېتر د روژا پوليتسر په نوم د 1905 د
فبرورۍ پر 17مه د هنګري په بوداپېست کښې
زېږېدلې وه۔ هغه تر ټولو زيات په بازګشتي
تابعو د خپل کار له امله پېژندل کېږي،
چې د بازګښت د تيورۍ د ډګر د جوړېدو
دپاره بنسټيز و۔
\olpseganchor{OLP-0627-B007}{end}

\olpseganchor{OLP-0627-B008}{start}
پېتر په سختو سياسي وختونو کښې لويه شوه؛
په تنکۍ ځوانۍ کښې ئې لومړۍ نړيواله جګړه
روانه وه۔ خو په بوداپېست کښې د ماريا
تېرېزيا د نجونو شتمن ښوونځي ته تللے
شوه او په 1922 کښې ترې فارغه شوه۔ بيا
ئې په بوداپېست کښې د پازمان پېتر په
پوهنتون کښې زده کړه وکړه، چې وروسته
ئې نوم لورانډ اېوتووش پوهنتون شو۔ د
پلار په ټينګار ئې د کيميا زده کړه
پيل کړه، خو وروسته رياضياتو ته
واوښته او په 1927 کښې فارغه شوه۔
که څه هم د ثانوي ښوونځي د رياضياتو
د تدريس اهليت ئې درلود، د هغه وخت
اقتصادي حالت ډېر خراب و، ځکه ستر
اقتصادي رکود د نړۍ اقتصاد اغېزمن
کړے و۔ په دې وخت کښې پېتر د
رياضياتو د مرستندويې ښوونکې او
شخصي ښوونکې په توګه بېلابېل واړۀ
کارونه کول۔ بالاخره پوهنتون ته
ستنه شوه، څو په رياضياتو کښې لوړې
زده کړې وکړي۔ اصل کښې ئې د عددونو
په تيورۍ د کار اراده وه، خو کله
چې پوهه شوه چې نتيجې ئې له مخکې
ثابتې شوې دي، نږدې وه چې رياضيات
بيخي پرېږدي۔ هغه وهڅول شوه چې د
ګوډل د ناتکميلۍ پر قضيو کار وکړي،
او بې له دې چې پوهه وي، د هغه
څو نتيجې ئې په بېلو لارو ثابتې
کړې۔ دې ئې باور بېرته راژوندے
کړ، او پېتر د ډېوېډ هېلبرټ له
بنسټيز پروګرامه په الهام د بازګښت
په تيورۍ خپلې لومړۍ ليکنې وليکلې۔
په 1935 کښې ئې دوکتورا واخيسته،
او په 1937 کښې د \emph{د سمبوليک
منطق مجلې} مديره شوه۔
\olpseganchor{OLP-0627-B008}{end}

\olpseganchor{OLP-0627-B009}{start}
د پېتر لومړنۍ ليکنې په پراخه توګه د
بازګشتي تابعو د تيورۍ بنسټ ايښودونکې
مرستې ګڼل کېږي۔ په \mbox{\LR{\latinfont \cite{Peter1935a}}}
کښې ئې د بازګښت د بېلابېلو ډولونو
اړيکه وڅېړله۔ په \mbox{\LR{\latinfont \cite{Peter1935b}}}
کښې ئې وښودله چې يوه ځانګړې په
بازګشتي ډول تعريف شوې تابع بنسټيزه
بازګشتي تابع نۀ ده۔ دې د وېلهېلم
اکېرمان يوه پخوانۍ نتيجه آسانه
کړه۔ د پېتر ساده شوې تابع هغه
ده چې اوس ډېر ځله د اکېرمان
تابع نومېږي، او کله، په لا
کره نوم، د اکېرمان--پېتر تابع۔
هغې د بازګشتي تابعو د تيورۍ
لومړے کتاب وليکۀ \mbox{\LR{\latinfont \citep{Peter1951}}}۔
\olpseganchor{OLP-0627-B009}{end}

\olpseganchor{OLP-0627-B010}{start}
د خپل کار له اهميت او اغېز سره سره پېتر
تر 1945 پورې د بشپړ وخت د تدريس دنده
ترلاسه نۀ کړه۔ د دويمې نړيوالې جګړې
پر وخت د هنګري د نازي اشغال لاندې،
پېتر ته د يهودانو ضد قوانينو له
امله د تدريس اجازه نۀ وه۔ په 1944
کښې حکومت په بوداپېست کښې د
يهودانو يوه محصوره سيمه جوړه کړه؛
دا له پاتې ښاره بېله وه او وسلوالو
ساتونکو ساتله۔ پېتر تر 1945 پورې
هلته ژوند کولو ته اړ وه، کله چې
دا سيمه آزاده شوه۔ بيا ئې د
بوداپېست د ښوونکو د روزنې په
کالج کښې تدريس وکړ، او له 1955
څخه وروسته د اېوتووش لورانډ
په پوهنتون کښې۔ هغه لومړۍ
هنګرۍ رياضيپوهه وه چې د
رياضياتو اکاډميکي دوکتوره شوه،
او لومړۍ ښځه وه چې د هنګري
د علومو اکاډمۍ ته وټاکل شوه۔
\olpseganchor{OLP-0627-B010}{end}

\olpseganchor{OLP-0627-B011}{start}
پېتر د رياضياتو د شوقمنې ښوونکې په توګه
پېژندل کېده۔ هغې له خپلو زده کوونکو
سره د رياضيکي مسئلو د طبيعت او
ښکلا څېړلو ته تر تش درس ورکولو
ترجيح ورکوله۔ له همدې امله به
زده کوونکو په مينه ورته «روزا
ترور» ويل۔ پېتر په 1977 کښې
د~71 کلونو په عمر وفات شوه۔
\olpseganchor{OLP-0627-B011}{end}

\olpseganchor{OLP-0627-B012}{start}
\begin{reading}
د نورو ژوندليکي ليکنو دپاره
\mbox{\LR{\latinfont \citep{Oconnor2014}}} او \mbox{\LR{\latinfont \citep{Andrasfai1986}}}
وګورئ۔ \mbox{\LR{\latinfont \citet{Tamassy1994}}} له پېتر سره
يوه لنډه مرکه کړې ده۔ د رياضياتو
د يوه په زړه پورې لوست دپاره د
پېتر کتاب \emph{له نامتناهيت سره
لوبې} وګورئ \mbox{\LR{\latinfont \citep{Peter2010}}}۔
\end{reading}
\olpseganchor{OLP-0627-B012}{end}

\olpseganchor{OLP-0628-B003}{start}
\olfileid{his}{bio}{rob}
\olpseganchor{OLP-0628-B003}{end}

\olpseganchor{OLP-0628-B004}{start}
\olsection{جوليا رابنسن}
\olpseganchor{OLP-0628-B004}{end}

\olpseganchor{OLP-0628-B005}{start}
\olphoto{robinson-julia}{جوليا رابنسن}
\olpseganchor{OLP-0628-B005}{end}

\olpseganchor{OLP-0628-B006}{start}
جوليا بومن رابنسن يوه امريکايۍ رياضيپوهه
وه۔ هغه په عمده توګه د پرېکړې پر مسئلو
د خپل کار، او تر ټولو زيات د هېلبرټ
د لسمې مسئلې په حل کښې د خپلو مرستو
له امله پېژندل کېږي۔ رابنسن د 1919
د دسمبر پر 8مه د ميزوري په سېنټ
لويس کښې زېږېدلې وه۔ رابنسن يادوي
چې له ماشومتوبه د عددونو سره شوق
لرلو \mbox{\LR{\latinfont \citep[4]{Reid1986}}}۔ په نهه
کلنۍ کښې په سرې تبې اخته شوه او
د روماتيکي تبې څو ځله تکراري
حملې ئې تېرې کړې۔ دې هغه اړ کړه
چې ډېر وخت په کټ کښې تېر کړي،
او په زده کړه کښې وروسته شوه۔
که څه هم د شخصي ښوونکو په مرسته
بېرته له نورو سره برابره شوه،
د ناروغۍ بدني اغېزو ئې پر ژوند
تلپاتې اثر وکړ۔
\olpseganchor{OLP-0628-B006}{end}

\olpseganchor{OLP-0628-B007}{start}
د ماشومتوب له ستونزو سره سره رابنسن
له ثانوي ښوونځي څخه په رياضياتو او
ساينس کښې له څو انعامونو سره فارغه
شوه۔ پوهنتوني زده کړې ئې د سېن
ډياګو په دولتي کالج کښې پيل کړې،
او په وروستي کال کښې د کاليفورنيا
برکلي پوهنتون ته لاړه۔ هلته د
رياضيپوه رافايل رابنسن تر اغېز
لاندې راغله۔ ښه ملګري شول او په
1941 کښې ئې واده وکړ۔ د يو استاد
د مېرمنې په توګه رابنسن په برکلي
کښې د رياضياتو په څانګه کښې له
تدريسه منع وه۔ که څه هم د
رياضياتو ټولګي ئې لا هم اورېدل،
هيله ئې لرله چې پوهنتون پرېږدي
او کورنۍ جوړه کړي۔ خو له واده
څه موده وروسته رابنسن په سينه‌بغل
اخته شوه۔ ورته وويل شول چې د
ماشومتوب د روماتيکي تبې له امله
ئې پر زړه ډېر داغي نسج راټول
شوے دے۔ د دې نسج د سختوالي له
امله ډاکټر وړاندوينه وکړه چې
له څلوېښتو کلونو زيات ژوند به
نۀ کوي، او ورته مشوره ورکړے
شوه چې اولاد ونۀ لري
\mbox{\LR{\latinfont \citep[13]{Reid1986}}}۔
\olpseganchor{OLP-0628-B007}{end}

\olpseganchor{OLP-0628-B008}{start}
رابنسن د ډېر وخت دپاره خپه وه، خو
بالاخره ئې پرېکړه وکړه چې د
رياضياتو زده کړه جاري وساتي۔ برکلي
ته ستنه شوه او په 1948 کښې ئې
د الفرېډ تارسکي تر لارښوونې
لاندې دوکتورا بشپړه کړه۔ تارسکي
ښودلې وه چې د حقيقي عددونو د
لومړي ترتيب تيوري د پرېکړې وړ
ده، او د ګوډل له کاره دا پايله
راتله چې د طبيعي عددونو د
لومړي ترتيب تيوري د پرېکړې
وړ نۀ ده۔ دا يوه لويه خلاصه
مسئله وه چې د نسبي عددونو
د لومړي ترتيب تيوري د پرېکړې
وړ ده که نۀ۔ رابنسن په خپله
رساله \mbox{\LR{\latinfont \citeyearpar{Robinson1949}}}
کښې ثابته کړه چې نۀ ده۔
\olpseganchor{OLP-0628-B008}{end}

\olpseganchor{OLP-0628-B009}{start}
رابنسن چې د پرېکړې له مسئلو سره ئې
شوق لرلو، بيا ئې د هېلبرټ د لسمې
مسئلې د حل هڅه وکړه۔ دا مسئله
د هغو 23 مشهورو رياضيکي مسئلو
له لړ څخه وه چې ډېوېډ هېلبرټ په
1900 کښې وړاندې کړې وې۔ لسمه
مسئله پوښتي چې آيا داسې الګوريتم
شته چې په متناهي وخت کښې ځواب
ورکړي چې د صحيح عددي ضريبونو
يوه کثيرالحدي معادله، لکه
\LR{$3x^2 - 2y +3 = 0$}، په صحيح
عددونو کښې حل لري که نۀ۔ داسې
پوښتنې \emph{ديوفانټي مسئلې}
نومېږي۔ له ځينو لومړنيو برياوو
وروسته رابنسن له مارټن ډېوس او
هېلري پټنم سره ګډ کار پيل کړ،
چې هغوی هم پر دې مسئله کار
کاوۀ۔ دوی وښودله چې نمايي
ديوفانټي مسئلې، چې نامعلومونه
پکې د نماوو په توګه هم راتلے
شي، د پرېکړې وړ نۀ دي۔ او
وښودله چې يوه ځانګړې اټکلي
دعوه، چې وروسته «J.R.» ونومول
شوه، دا پايله لري چې د
هېلبرټ لسمه مسئله د پرېکړې
وړ نۀ ده \mbox{\LR{\latinfont \citep{DavisPutnamRobinson1961}}}۔
رابنسن په ټولو 1960s کلونو
کښې پر دې مسئله کار جاري
وساتۀ۔ په 1970 کښې ځوان
روسي رياضيپوه يوري ماتياسېوېچ
بالاخره د J.R. فرضيه ثابته
کړه۔ ګډه نتيجه اوس د
ماتياسېوېچ--رابنسن--ډېوس--پټنم
قضيه، يا په لنډه MRDP قضيه،
نومېږي۔ ماتياسېوېچ او رابنسن
ملګري شول او پر څو مقالو ئې
ګډ کار وکړ۔ رابنسن په يوه ليک کښې ماتياسېوېچ ته وويل چې د دومره لرې واټن سره سره د دوی ګډ کار د دواړو له بېل کار څخه ډېر پرمختګ کړے و \mbox{\LR{\latinfont \citep[45]{Matijasevich1992}}}۔
\olpseganchor{OLP-0628-B009}{end}

\olpseganchor{OLP-0628-B010}{start}
رابنسن د امريکا د رياضيکي ټولنې
لومړۍ ښځينه صدره وه، او د علومو ملي اکاډمۍ د رياضياتو څانګې ته ټاکل شوې لومړۍ ښځه وه \mbox{\LR{\latinfont \citep{Feferman1994}}}۔ هغه د وينې د
سرطان له تشخيص وروسته د 1985
د جولای پر 30مه د 65 کلونو
په عمر وفات شوه۔
\olpseganchor{OLP-0628-B010}{end}

\olpseganchor{OLP-0628-B011}{start}
\begin{reading}
د رابنسن رياضيکي ليکنې د هغې په
\textit{راټولو شويو ليکنو} کښې
شته \mbox{\LR{\latinfont \citep{Robinson1996}}}۔ په دې
کښې د علومو ملي اکاډمۍ ژوندليکي
يادښت \mbox{\LR{\latinfont \citep{Feferman1994}}} هم
بيا چاپ شوے دے۔ د رابنسن مشرې
خور کانسټنس ريډ د مرکو پر بنسټ
«د جوليا ځانليک» خپور کړے دے
\mbox{\LR{\latinfont \citep{Reid1986}}}، او يو بشپړ
ژوندليکي يادښت هم
\mbox{\LR{\latinfont \citep{Reid1996}}}۔ د رابنسن او
د هېلبرټ د لسمې مسئلې په هکله
د يو لنډ مستند فلم لارښوونه
جورج چيچري کړې ده
\mbox{\LR{\latinfont \citep{Csicsery2016}}}۔ له رابنسن
سره د يوري ماتياسېوېچ د ګډ
کار، او د هغه پر کار د هغې
د اغېز د لنډ يادښت دپاره
\mbox{\LR{\latinfont \citep{Matijasevich1992}}} وګورئ۔
\end{reading}
\olpseganchor{OLP-0628-B011}{end}

\olpseganchor{OLP-0629-B004}{start}
\olfileid{his}{bio}{rus} 
\olpseganchor{OLP-0629-B004}{end}

\olpseganchor{OLP-0629-B005}{start}
\olsection{برټرېنډ رسل}
\olpseganchor{OLP-0629-B005}{end}

\olpseganchor{OLP-0629-B006}{start}
\olphoto{russell-bertrand}{برټرېنډ رسل}
\olpseganchor{OLP-0629-B006}{end}

\olpseganchor{OLP-0629-B007}{start}
برټرېنډ رسل د نوې تجزيه‌يي فلسفې د بنسټ
ايښودونکو له ډلې ګڼل کېږي۔ رسل د 1872
د مې پر 18مه زېږېدلے و۔ هغه يوازې په
فلسفه او منطق کښې د خپل کار له امله
نۀ پېژندل کېده، بلکې په بېلابېلو
موضوعاتو ئې ډېر عام‌فهم کتابونه هم
ليکلي وو۔ هغه د خپل ټول ژوند په
اوږدو کښې يو تود سياسي فعال هم و۔
\olpseganchor{OLP-0629-B007}{end}

\olpseganchor{OLP-0629-B008}{start}
رسل د وېلز د مونماوتشاير په ټرېلېک کښې
زېږېدلے و۔ مور او پلار ئې د برتانوي
اشرافو غړي وو۔ هغوي آزادخياله وو،
او د هغه وخت د بوسټن له تندلارو سره
ئې هم ملګرتيا کړې وه۔ له بده مرغه،
د رسل مور او پلار د هغه په کم عمرۍ
کښې وفات شول، او رسل د خپل نيکۀ
او نيا سره اوسېدو ته ولېږل شو۔ هلته
ئې مذهبي روزنه وشوه؛ دا هغه څه وو
چې مور او پلار ئې په هر قيمت ترې
ځان ساتل غوښتل۔ نيا ئې د اخلاقو
په ټولو چارو کښې ډېره سخته وه۔ په
تنکۍ ځوانۍ کښې ئې زياتره په کور
کښې له شخصي ښوونکو زده کړه کوله۔
\olpseganchor{OLP-0629-B008}{end}

\olpseganchor{OLP-0629-B009}{start}
په تجزيه‌يي فلسفه، او په تېره په منطق،
کښې د رسل اغېز ډېر ستر دے۔ هغه د
کېمبرېج په ټرينيټي کالج کښې رياضيات
او فلسفه ولوستل، او هلته د رياضيپوه
او فيلسوف الفرېډ نورټ وېټ هېډ تر
اغېز لاندې راغے۔ په 1910 کښې رسل
او وېټ هېډ د \emph{د رياضياتو اصول}
لومړے ټوک خپور کړ، او پکښې ئې د
دې نظر ملاتړ وکړ چې رياضيات منطق
ته راکمول کېدلے شي۔ وروسته ئې په
سلګونو کتابونه، مقالې او سياسي
رسالې خپرې کړې۔ په 1950 کښې ئې
د ادبياتو نوبېل جايزه وګټله۔
\olpseganchor{OLP-0629-B009}{end}

\olpseganchor{OLP-0629-B010}{start}
رسل په سياست او ټولنيز فعاليت کښې ژور
ښکېل و۔ د لومړۍ نړيوالې جګړې پر
وخت، د سوله‌پالو فعاليتونو او احتجاج
له امله ونيول شو او د شپږو مياشتو
دپاره زندان ته ولېږل شو۔ په زندان
کښې ئې ليکل او لوستل کولے شول،
او وايي چې دا تجربه ئې «خورا د
خوښې وړ» وموندله۔ هغه د سولې لپاره اوږد فعاليت وکړ۔ د دويمې نړيوالې جګړې پر مهال ئې د هټلر پر ضد د متحدينو ملاتړ وکړ۔ په 1961 کښې د اټومي وسلو د
لرې کولو په يوه غونډه کښې د ګډون
له امله بيا بندي شو۔ هغه په 1948
کښې د الوتکې له يوې پېښې هم ژوندے
پاتې شو؛ په دې پېښه کښې يوازې د
سګرټ څښونکو په برخه کښې ناست کسان
ژوندي پاتې شوي وو۔ له دې امله رسل
ويل چې ژوند ئې د سګرټ څښلو مرهون
دے۔ رسل څلور ځله وادۀ کړے و، خو
د وادۀ څخه بهر د عاشقانه اړيکو
شهرت ئې هم درلود۔ هغه د 1970
د فبرورۍ پر 2مه د 97 کلونو په
عمر د وېلز په پېنريندېودراېت کښې
وفات شو۔
\olpseganchor{OLP-0629-B010}{end}

\olpseganchor{OLP-0629-B011}{start}
\begin{reading}
رسل په درېو برخو کښې خپل ژوندليک
وليکۀ، چې له 1872--1967 پورې ئې
ژوند رانغاړي
\mbox{\LR{\latinfont \citep{Russell1967,Russell1968,Russell1969}}}۔
په مکماسټر پوهنتون کښې د برټرېنډ
رسل د څېړنې مرکز د هغه د آرشيف
کور دے۔ د راټولو آثارو د ټوکونو
په هکله د معلوماتو، د لټون وړ
فهرستونو او آرشيفي پروژو دپاره
ئې په \mbox{\LR{\latinfont \citet{Duncan2015}}} کښې وېبپاڼه
وګورئ۔ د رسل مقاله \emph{د دلالت
په هکله} \mbox{\LR{\latinfont \citep{Russell1905}}} د
شلمې پېړۍ د تجزيه‌يي فلسفې له
کلاسيکو آثارو ده۔
\olpseganchor{OLP-0629-B011}{end}

\olpseganchor{OLP-0629-B012}{start}
د فلسفې د سټېنفورډ پوهنغونډ په
\mbox{\LR{\latinfont \citep{Irvine2015}}} کښې د رسل مدخل
د خواهش او سياسي تيورۍ په هکله
د هغه د خبرو غږيزې ټوټې لري۔ د
رسل ډېرې ويډيويي مرکې په انټرنېټ
کښې شته۔ د بېلګې په توګه، د
سګرټ څښلو او د الوتکې په پېښه
کښې د ښکېلتيا په هکله ئې خبرې
په \mbox{\LR{\latinfont \citet{RussellND}}} کښې وګورئ۔
د هغه ځينې آثار، چې \emph{د
رياضيکي فلسفې پېژندنه} هم پکښې
ده، په \mbox{\LR{\latinfont \citet{LibriVoxND}}} کښې د
وړيا غږيزو کتابونو په توګه شته۔
\end{reading}
\olpseganchor{OLP-0629-B012}{end}

\olpseganchor{OLP-0630-B004}{start}
\olfileid{his}{bio}{tar}
\olpseganchor{OLP-0630-B004}{end}

\olpseganchor{OLP-0630-B005}{start}
\olsection{الفرېډ تارسکي}
\olpseganchor{OLP-0630-B005}{end}

\olpseganchor{OLP-0630-B006}{start}
\olphoto{tarski-alfred}{الفرېډ تارسکي}
\olpseganchor{OLP-0630-B006}{end}

\olpseganchor{OLP-0630-B007}{start}
الفرېډ تارسکي د 1901 د جنورۍ پر 14مه
د پولېنډ په وارسا کښې زېږېدلے و؛
هغه وخت دا د روسيې د سترواکۍ برخه
وه۔ تارسکي چې «ناپوليوني» بلل شوے
دے، پُرجوش، ډېر غږېدونکے او توند
و۔ توان ئې ډېر ځله په درسونو کښې
هم ښکاره کېده؛ يو ځل ئې د درس
پر وخت د سګرټ په غورځولو د کثافاتو
ټوکرۍ اور واخيست، او بيا په هغه
ودانۍ کښې له درس ورکولو منع شو۔
\olpseganchor{OLP-0630-B007}{end}

\olpseganchor{OLP-0630-B008}{start}
تارسکي له کم عمره د پوهې تنده لرله۔
که څه هم وروسته به ئې زده کوونکو
ته ويل چې منطق ئې ځکه ولوست چې
دا يوازينے مضمون و چې B درجه ئې
پکښې اخيستې وه، د ثانوي ښوونځي
اسناد ئې ښيي چې په ټولو مضامينو،
آن په منطق کښې، ئې A درجې اخيستې
وې۔ له 1918 څخه تر 1924 پورې
ئې د وارسا په پوهنتون کښې زده
کړه وکړه۔ لومړۍ اراده ئې د
بيولوجۍ لوستل وو، خو په رياضياتو،
فلسفه او منطق کښې ئې شوق پيدا
شو، ځکه پوهنتون د منطق او فلسفې
د وارسا مکتب مرکز و۔ تارسکي په
1924 کښې د ستانيسواف لېشنيېفسکي
تر لارښوونې لاندې دوکتورا واخيسته۔
\olpseganchor{OLP-0630-B008}{end}

\olpseganchor{OLP-0630-B009}{start}
په 1939 کښې امريکا ته له کډې وړلو
مخکښې، تارسکي د خپل ځينې تر ټولو
مهم کار بشپړول په وارسا کښې د
ثانوي ښوونځي د ښوونکي په توګه
وکړل۔ د منطقي پايلې او منطقي
رښتياوالي آثار ئې په همدې وخت
کښې وليکل۔ په 1939 کښې تارسکي
د درسونو د يو سفر دپاره امريکا
ته تلے و۔ د هغه د سفر پر وخت
جرمني پر پولېنډ يرغل وکړ، او د
يهودي اصل له امله تارسکي بېرته
نۀ شو ستنېدلے۔ مېرمن او ماشومان
ئې د جګړې تر پايه په پولېنډ کښې
پاتې شول، خو بيا هغوي هم امريکا
ته کډه يووړلے شوه۔ تارسکي په
هارورډ، د نيويارک ښار په کالج،
د پرينسټن د لوړو څېړنو په موسسه
او بالاخره د برکلي د کليفورنيا
په پوهنتون کښې تدريس وکړ۔ هلته
ئې د منطق او د علم د طريقه‌پوهې
ګڼ‌څانګيز پروګرام جوړ کړ۔ تارسکي
د 1983 د اکتوبر پر 26مه د 82
کلونو په عمر وفات شو۔
\olpseganchor{OLP-0630-B009}{end}

\olpseganchor{OLP-0630-B010}{start}
\begin{reading}
د تارسکي د ژوند په هکله د نورو
معلوماتو دپاره ژوندليک \emph{الفرېډ
تارسکي: ژوند او منطق}
\mbox{\LR{\latinfont \citep{Feferman2004}}} وګورئ۔ د منطقي
پايلې او رښتياوالي په هکله ئې
بنسټيز آثار په انګليسي کښې په
\mbox{\LR{\latinfont \citep{Tarski1983}}} کښې شته۔ د
تارسکي ټول اصلي آثار په څلور
ټوکي لړۍ \mbox{\LR{\latinfont \citep{Tarski1981}}} کښې
راټول شوي دي۔
\end{reading}
\olpseganchor{OLP-0630-B010}{end}

\olpseganchor{OLP-0631-B004}{start}
\olfileid{his}{bio}{tur} 
\olpseganchor{OLP-0631-B004}{end}

\olpseganchor{OLP-0631-B005}{start}
\olsection{اېلن ټيورينګ}
\olpseganchor{OLP-0631-B005}{end}

\olpseganchor{OLP-0631-B006}{start}
اېلن ټيورينګ د 1912 د جون پر 23مه
د لندن په مېډا وېل کښې زېږېدلے و۔
هغه د نظري کمپيوټرپوهنې پلار ګڼل
کېږي۔ په طبيعي علومو او رياضياتو
کښې ئې شوق له کم عمره پيل شو۔
خو په ماشومتوب کښې ئې ښوونځيو
دې شوقونو ته مناسب پام نۀ کولو،
ځکه هلته په ادبياتو او کلاسيکو
زده کړو ټينګار کېده۔ په نتيجه
کښې په ښوونځي کښې کمزورے و او
ډېرو ښوونکو به ترټلو۔
\olpseganchor{OLP-0631-B006}{end}

\olpseganchor{OLP-0631-B007}{start}
\olphoto{turing-alan}{اېلن ټيورينګ}
\olpseganchor{OLP-0631-B007}{end}

\olpseganchor{OLP-0631-B008}{start}
ټيورينګ د کېمبرېج په کينګز کالج کښې
د لومړۍ پوهنتوني درجې دپاره
رياضيات ولوستل۔ په 1936 کښې ئې
هغه څه جوړ کړل چې اوس د ټيورينګ
ماشين بلل کېږي، څو د محاسبه
کېدونکې تابع مفهوم په کره ډول
تعريف کړي او د پرېکړې د مسئلې
ناپرېکړه کېدنه ثابته کړي۔ الونزو
چرچ تر هغه مخکښې دې نتيجې ته
رسېدلے و، او په خپل لامبډا حساب
ئې ثابته کړې وه۔ د ټيورينګ مقاله
بيا هم د چرچ نتيجې ته په حوالې
سره خپره شوه۔ چرچ ټيورينګ
پرينسټن ته راوبالۀ؛ هغه له
1936--1938 پورې هلته و او د
چرچ تر لارښوونې لاندې ئې
دوکتورا واخيسته۔
\olpseganchor{OLP-0631-B008}{end}

\olpseganchor{OLP-0631-B009}{start}
په منطق کښې له شوق سره سره، په طبيعي
علومو کښې د ټيورينګ پخوانے شوق هم
غښتلے پاتې شو۔ د دويمې نړيوالې
جګړې پر وخت د بلېچلي پارک په
برتانوي د رمزشننې څانګه کښې د
کار پر مهال ئې عملي مهارتونه
وکارول شول۔ ټيورينګ د جرمني د
سمندري پوځي اړيکو د رمز، يعنې
اېنيګما، په ماتولو کښې مرکزي
شخص و۔ د هغه د رمزشننې کار او د ګورډن وېلچمن طرحو د برقي او ميخانيکي بامب په جوړېدو کښې برخه لرله۔ دغه ماشين د اېنيګما د رمز په ماتولو کښې کارېدۀ۔ په بلېچلي پارک کښې د لورېنز د رمز لپاره بل ماشين، کولوسس، وکارول شو۔ ټامي فلاورز او د هغه ډلې دغه الېکټروني ماشين جوړ کړ، چې د ځانګړو دندو لپاره پروګرام کېدے شو۔
\olpseganchor{OLP-0631-B009}{end}

\olpseganchor{OLP-0631-B010}{start}
ټيورينګ په 1941 کښې د بلېچلي پارک له خپلې همکارې جون کلارک سره کوژده وکړه۔ هغې ته ئې څرګنده کړه چې همجنس‌پال دے، او د هماغه کال په وروستيو کښې يې کوژده ماته شوه۔ په ژوند کښې ئې
څو معشوقان لرل، که څه هم هغه
وخت همجنسي اعمال په برتانيه کښې
جرمي سرغړونې وې۔ په 1952 کښې
ئې د هغه وخت د معشوق يو ملګري
کور غلا کړ، او د پوليسو د راپور
د ورکولو پر وخت ټيورينګ ومنله
چې همجنسي اړيکه لري؛ هغه ګڼله
چې حکومت دا اعمال قانوني کوي۔
دا سمه نۀ وه، او د سختې
بې‌حيايۍ تور پرې ولګېد۔ زندان
ته د تلو پر ځای ټيورينګ هغه
هورموني درملنه غوره کړه چې
جنسي ميل ئې کمولو۔ ټيورينګ د
1954 د جون پر 8مه مړ وموندل
شو؛ علت ئې د ساينايډ زيات مقدار
و، او غالباً ځان‌وژنه وه۔ په
2013 کښې ملکې الېزابت~II ورته
شاهي بخښنه ورکړه۔
\olpseganchor{OLP-0631-B010}{end}

\olpseganchor{OLP-0631-B011}{start}
\begin{reading}
د اېلن ټيورينګ د يوه هراړخيز
ژوندليک دپاره \mbox{\LR{\latinfont \citet{Hodges2014}}}
وګورئ۔ د هغه ژوند او کار د
\emph{د رمز ماتول} ډرامې الهام
شو؛ دا په 1996 کښې د ټېلېويژن
دپاره جوړه شوه، او ډېريک جکوبي
پکښې د ټيورينګ رول ولوباوه۔
\emph{د تقليد لوبه} هم، چې
بېنېډېکټ کمبربېچ او کېرا نايتلي
پکښې لوبېږي او د اکاډمۍ جايزې
ته نومول شوے فلم دے، د اېلن
ټيورينګ پر ژوند او په بلېچلي
پارک کښې پر وخت په آزاد ډول
ولاړ دے \mbox{\LR{\latinfont \citep{Imitation2014}}}۔
\olpseganchor{OLP-0631-B011}{end}

\olpseganchor{OLP-0631-B012}{start}
\mbox{\LR{\latinfont \citet{Radiolab2012}}} د ټيورينګ د
ژوند او کار په هکله څو غږيزې
خپرونې لري۔ د بي بي سي د
هورايزن مستند فلم \emph{د ډاکټر
ټيورينګ عجيب ژوند او مرګ} په
انټرنېټ کښې کتل کېدلے شي
\mbox{\LR{\latinfont \citep{Sykes1992}}}۔ \mbox{\LR{\latinfont \citep{Theelen2012}}}
د ليګو د يوه کارکوونکي ټيورينګ
ماشين لنډه ويډيو ده؛ دا په
2012 کښې د هغه د زېږون د
سلمې کليزې د درناوي دپاره
جوړ شوے و۔
\olpseganchor{OLP-0631-B012}{end}

\olpseganchor{OLP-0631-B013}{start}
د ټيورينګ د ماشينونو او د پرېکړې
د مسئلې په هکله د هغه اصلي
مقاله \mbox{\LR{\latinfont \citet{Turing1937}}} ده۔
\end{reading}
\olpseganchor{OLP-0631-B013}{end}

\olpseganchor{OLP-0632-B004}{start}
\olfileid{his}{bio}{zer}
\olpseganchor{OLP-0632-B004}{end}

\olpseganchor{OLP-0632-B005}{start}
\olsection{اېرنست زېر مېلو}
\olpseganchor{OLP-0632-B005}{end}

\olpseganchor{OLP-0632-B006}{start}
\olphoto{zermelo-ernst}{اېرنست زېر مېلو}
\olpseganchor{OLP-0632-B006}{end}

\olpseganchor{OLP-0632-B007}{start}
اېرنست زېر مېلو د 1871 د جولای پر 27مه
د جرمني په برلين کښې زېږېدلے و۔
هغه پنځۀ خويندې لرلې، خو کورنۍ
ئې د خرابې روغتيا له امله کړېده
او يوازې درې ئې تر بالغ عمره
ژوندۍ پاتې شوې۔ مور او پلار ئې
هم په کم عمرۍ کښې وفات شول؛
کله چې هغه اووۀ‌لس کلن و، هغه
او وروڼه خويندې ئې يتيمان شول۔
زېر مېلو په هنرونو، او په تېره
په شاعرۍ، کښې ژور شوق لرلو۔
هغه په ځيرکتيا، حاضرځوابۍ او
نيوکه‌کوونکي خوي پېژندل کېده۔
د هغه په تر ټولو مشهورو رياضيکي
لاسته راوړنو کښې په 1904 کښې
د ټاکنې د بديهي اصل معرفي کول
او په 1908 کښې د سټونو د
تيورۍ اصل‌بندي کول شامل دي۔
\olpseganchor{OLP-0632-B007}{end}

\olpseganchor{OLP-0632-B008}{start}
په پوهنتون کښې د زېر مېلو شوقونه
بېلابېل وو۔ هغه د فزيکس،
رياضياتو او فلسفې درسونه واخيستل۔
د هېرمان شوارڅ تر لارښوونې لاندې
ئې په 1894 کښې د برلين په
پوهنتون کښې خپله رساله
\emph{د تغيراتو په حساب کښې
څېړنې} بشپړه کړه۔ په 1897
کښې ئې د ګوټينګن په پوهنتون
کښې د نورو زده کړو پرېکړه وکړه،
او هلته د ډېوېډ هېلبرټ بنسټيز
کار ډېر اغېزمن کړ۔ په 1899
کښې د پروفېسرۍ اهل شو، خو
دنده ئې تر يوولسو کلونو وروسته
ترلاسه نۀ کړه؛ ښايي علت ئې
عجيب چلند او «عصبي تلوار» و۔
\olpseganchor{OLP-0632-B008}{end}

\olpseganchor{OLP-0632-B009}{start}
زېر مېلو بالاخره په 1910 کښې د زيوريخ
په پوهنتون کښې معاش لرونکې
پروفېسري ترلاسه کړه، خو په
1916 کښې د نري رنځ له امله
تقاعد ته اړ شو۔ له رغېدو وروسته
ئې په 1921 کښې د فرايبورګ په
پوهنتون کښې افتخاري پروفېسري
ترلاسه کړه۔ په دې وخت کښې ئې
د رياضياتو پر بنسټونو کار کولو۔
د تورالف سکولېم او کُرټ ګوډل
آثار ئې په غصه کړل، او په خپلو
مقالو کښې ئې په ښکاره د هغوي
پر لارو نيوکې وکړې۔ په 1935
کښې په جرمني کښې د هټلر د
واک ته رسېدو د مخالفت او د
خپل نامقبوليت له امله د
فرايبورګ له دندې لرې کړے شو۔
\olpseganchor{OLP-0632-B009}{end}

\olpseganchor{OLP-0632-B010}{start}
د زېر مېلو د ژوند وروستي کلونه په
ګوښه‌والي کښې تېر شول۔ په 1935
کښې له دندې له لرې کېدو وروسته
ئې رياضيات پرېښودل۔ کليوالې
سيمې ته ولاړ او په ساده ډول
ئې ژوند کولو۔ په 1944 کښې ئې
وادۀ وکړ، او د سترګو د نظر د
کمېدو له امله بيخي په مېرمن
متکي شو۔ تر 1951 پورې ئې نظر
په بشپړ ډول لاړو۔ هغه د 1953
د مې پر 21مه د جرمني په
ګينترشتال کښې وفات شو۔
\olpseganchor{OLP-0632-B010}{end}

\olpseganchor{OLP-0632-B011}{start}
\begin{reading}
د زېر مېلو د بشپړ ژوندليک دپاره
\mbox{\LR{\latinfont \citet{Ebbinghaus2015}}} وګورئ۔ د
هغه د 1904 او 1908 بنسټيزې
مقالې په اصلي جرمني کښې لوستل
کېدلے شي
\mbox{\LR{\latinfont \citep{Zermelo1904,Zermelo1908}}}۔ د
زېر مېلو راټول آثار، چې د فزيکس
ليکنې ئې هم پکښې دي،
په انګليسي ژباړه کښې په
\mbox{\LR{\latinfont \citep{Ebbinghaus2010,Ebbinghaus2013}}}
کښې شته۔
\end{reading}
\olpseganchor{OLP-0632-B011}{end}

\olpseganchor{OLP-0633-B004}{start}
\olchapter{his}{set}{د سټ نظريې تاريخ او اسطوري بيانونه}
\olpseganchor{OLP-0633-B004}{end}

\olpseganchor{OLP-0633-B005}{start}
\begin{editorial}
په دې باب کښې د ټېم بټن د «خلاصه سټ نظريه» د متن تاريخي سريزه شامله ده۔
\end{editorial}
\olpseganchor{OLP-0633-B005}{end}

\olpseganchor{OLP-0634-B003}{start}
\olfileid{his}{set}{infinitesimals}
\olsection{انفينټېسيمالونه او مشتق اخيستل}
\olpseganchor{OLP-0634-B003}{end}

\olpseganchor{OLP-0634-B004}{start}
نيوټن او لايبنيڅ د 17مې پېړۍ په پای کښې، له يو بل څخه په خپلواکه توګه، حسابان وموندل۔ د حسابان يوه په ځانګړې توګه مهمه کارونه \emph{مشتق اخيستل} وو۔ په لنډه توګه، د مشتق اخيستلو موخه دا ده چې په يوه نقطه کښې د يوې تابع د «بدلون د کچې»، يا ميلان، يو مفهوم راکړي۔
\olpseganchor{OLP-0634-B004}{end}

\olpseganchor{OLP-0634-B005}{start}
د دې نظر د څرګندولو يوه روښانه لاره دا ده۔ تابع
\LR{$f(x) = \nicefrac{x^2}{4} + \nicefrac{1}{2}$} په پام کښې ونيسئ، چې لاندې په تور رنګ ښودل شوې ده:
\begin{center}
	\begin{tikzpicture}[scale=1]
	\draw[->, gray] (-1,0) -- (4.25,0) node[right] {\LR{$x$}};
	\draw[->, gray] (0,-0.25) -- (0,5.25) node[above] {\LR{$f(x)$}};
\olpseganchor{OLP-0634-B005}{end}

\olpseganchor{OLP-0634-B006}{start}
	\foreach \x/\xtext in {1/1, 2/2, 3/3, 4/4}
	\draw[shift={(\x,0)}, gray] (0pt,2pt) -- (0pt,-2pt) node[below] {\LR{$\xtext$}};
\olpseganchor{OLP-0634-B006}{end}

\olpseganchor{OLP-0634-B007}{start}
	\foreach \y/\ytext in {1/1, 2/2, 3/3, 4/4, 5/5}
	\draw[shift={(0,\y)}, gray] (2pt,0pt) -- (-2pt,0pt) node[left] {\LR{$\ytext$}};
\olpseganchor{OLP-0634-B007}{end}

\olpseganchor{OLP-0634-B008}{start}
	\draw[oldiagcolorC] (0.5, 9/16) -- (3.5, 57/16) -- (3.5, 9/16)--cycle;
	\draw[oldiagcolorD] (.5, 9/16) -- (2.5, 33/16) -- (2.5, 9/16) -- cycle;
	\draw[oldiagcolorE] (.5, 9/16) -- (1.5, 17/16) -- (1.5, 9/16) -- cycle;
	\draw[thick] (-1,.75) parabola bend (0,0.5) (4,4.5);
	\end{tikzpicture}
\end{center}
فرض کړئ چې غواړو په \LR{$c =
\nicefrac{1}{2}$} کښې د تابع ميلان پيدا کړو۔ پيل په داسې مثلث کوو چې وتر يې په هغه نقطه کښې د ميلان تقريب راکوي؛ کېدای شي دا هماغه پورتنے سور مثلث وي۔ که \LR{$\beta$} د مثلث د قاعدې اوږدوالے وي، نو لوړوالے يې
\LR{$f(\nicefrac{1}{2}+\beta) - f(\nicefrac{1}{2})$} دے؛ ځکه نو د وتر ميلان دا دے:
\[
\frac{f(\nicefrac{1}{2}+\beta) - f(\nicefrac{1}{2})}{\beta}.
\]
نو زموږ د سور مثلث ميلان، چې د قاعدې اوږدوالے يې~\LR{$3$} دے، عين~\LR{$1$} دے۔ د يو کوچني مثلث وتر، يعنې د آبي مثلث چې د قاعدې اوږدوالے يې~\LR{$2$} دے، غوره تقريب راکوي؛ ميلان يې \LR{$\nicefrac{3}{4}$} دے۔ تر دې هم کوچنے زرغون مثلث، چې د قاعدې اوږدوالے يې~\LR{$1$} دے، لا غوره تقريب راکوي؛ ميلان يې \LR{$\nicefrac{1}{2}$} دے۔
\olpseganchor{OLP-0634-B008}{end}

\olpseganchor{OLP-0634-B009}{start}
څومره چې مثلثونه کوچني کېږي، هومره غوره تقريبونه راکوي۔ نو داسې څه ويلے شو: د هغه مثلث وتر چې د قاعدې اوږدوالے يې \emph{انفينټېسيمال}، يعنې بې‌نهايت کوچنے، وي، په خپله د \LR{$c =
\nicefrac{1}{2}$} ميلان راکوي۔ په دې توګه به د تابع \LR{$f$} د (لومړي) مشتق لپاره، په نقطه \LR{$c$} کښې، دا فارمول ترلاسه کړو:
\[
{f'}(c) = \frac{f(c+\beta) - f(c)}{\beta} \text{\RL{ چې \LR{$\beta$} انفينټېسيمال دے۔}}
\]
او په تقريباً همدغه معنا نيوټن او لايبنيڅ دا خبره کړې وه۔
\olpseganchor{OLP-0634-B009}{end}

\olpseganchor{OLP-0634-B010}{start}
خو چې هغوي دا خبره کړې ده، نو بايد ترې وپوښتو: \emph{انفينټېسيمال} څه دے؟ يوه جدي دوه‌لارې مسئله پيدا کېږي۔ که \LR{$\beta = 0$} وي، نو د \LR{$f'$} لپاره ليکل شوے خارج قسمت تعريف نۀ لري، ځکه په \LR{$0$} وېش پکې راځي۔ مشتق بيا د همدې خارج قسمت د حد له لارې موندل کېږي۔ خو که \LR{$\beta >
0$} وي، نو يوازې د ميلان يو \emph{تقريب} ترلاسه کوو، خپله ميلان نۀ ترلاسه کوو۔
\olpseganchor{OLP-0634-B010}{end}

\olpseganchor{OLP-0634-B011}{start}
دا د پخواني وخت په اړه د وروستي وخت يوه بې‌ځايه اندېښنه نۀ ده۔ برکلي د نيوټن پر پيروانو داسې نيوکه کوي:
\begin{quote}
زه دا منم چې په نښو هر څيز يا هېڅ څيز ښودلے شو۔ نو په ليکنه \LR{$c + \beta$} کښې \LR{$\beta$} د زياتوالي يا د هېڅ شي نښه کېدے شوه۔ له دې دوو معناوو چې هره يوه ورته ټاکئ، استدلال بايد هماغې معنا ته وفادار پاتې شي۔ په دوو معناوو استدلال څرګنده مغالطه ده۔ (\mbox{\LR{\latinfont \citealt[\S{}XV]{Berkeley1734}}}، د پورتنۍ بېلګې په نښو۔)
\end{quote}
له برکلي څخه د انفينټېسيمال حسابان د دفاع لپاره څوک ويلے شي چې د «انفينټېسيمالونو» خبرې يوازې مجازي دي۔ ويلے شي چې تر څو يو \emph{په رښتيا کوچنے} مثلث واخلو، د مماس \emph{کافي ښه} تقريب به ترلاسه کړو۔ برکلي دې ته هم ځواب درلود: دا که د انجينيرۍ لپاره کافي ښه وي، خو د رياضياتو \emph{حيثيت} کمزورے کوي، ځکه
\begin{quote}
موږ ته ويل کېږي چې \emph{\mbox{\LR{\latinfont in rebus mathematicis errores qu\`{a}m minimi non sunt contemnendi}}}. [د رياضياتو په چارو کښې ان تر ټولو وړې تېروتنې هم بايد له پامه ونۀ غورځوو۔] \mbox{\LR{\latinfont \citep[\S{}IX]{Berkeley1734}}}
\end{quote}
په مايله بڼه ليکل شوے عبارت د نيوټن د خپل کتاب \emph{د منحني‌ګانو مساحت موندل} (1704) تقريباً لفظ په لفظ نقل دے۔
\olpseganchor{OLP-0634-B011}{end}

\olpseganchor{OLP-0634-B012}{start}
د برکلي فلسفي اعتراضونه ډېر ژور او تېز دي۔ سره له دې، حسابان بېخي بريالے کار و، او رياضيپوهانو له تېروتنې سره د مخ کېدو پرته د هغه کارول جاري وساتل۔
\olpseganchor{OLP-0634-B012}{end}

\olpseganchor{OLP-0635-B004}{start}
\olfileid{his}{set}{limits}	
\olpseganchor{OLP-0635-B004}{end}

\olpseganchor{OLP-0635-B005}{start}
\olsection{د حدونو کره تعريف}
\olpseganchor{OLP-0635-B005}{end}

\olpseganchor{OLP-0635-B006}{start}
نن سبا د پورتنۍ مسئلې معياري حل د انفينټېسيمالونو پرېښودل دي۔ دا يې طريقه ده۔
\olpseganchor{OLP-0635-B006}{end}

\olpseganchor{OLP-0635-B007}{start}
ومو ليدل چې څومره \LR{$|\beta|$} کوچنے کېږي، د ميلان هومره ښه تقريبونه ترلاسه کوو۔ په رښتيا، چې \LR{$\beta$} له \LR{$0$} سره هر څومره نژدې کېږي، د \LR{$\frac{f(c+\beta)-f(c)}{\beta}$} ارزښت هغه ميلان ته «بې له حده ورنژدې کېږي» چې موږ يې غواړو۔ نو د دې پر ځای چې د \LR{$\beta = 0$} \emph{په حالت کښې} پېښې په پام کښې ونيسو، يوازې د \LR{$\frac{f(c+\beta)-f(c)}{\beta}$} \emph{بهير} په پام کښې نيولو ته اړتيا لرو، کله چې \LR{$\beta \ne 0$} وي او له \LR{$0$} سره نژدې کېږي۔
\olpseganchor{OLP-0635-B007}{end}

\olpseganchor{OLP-0635-B008}{start}
په دې ډول چې يې ووايو، عمومي ستونزه د داسې دعوو په معنا پوهېدل دي:
\begin{center}
چې \LR{$x$} له \LR{$c$} سره نژدې کېږي، \LR{$g(x)$} له \LR{$\ell$} سره بې له حده نژدې کېږي۔
\end{center}
دا په لا لنډه بڼه داسې ليکلے شو:
\[
\lim_{x \rightarrow c}g(x) = \ell.
\]
په 19مه پېړۍ کښې وايرشټراس، د کوشي د پخواني کار پر بنسټ، د دې عبارت يو پوره کره تعريف وړاندې کړ۔ نظر دا دے چې \LR{$g(x)$} له~\LR{$\ell$} سره هر څومره چې وغواړو نژدې کولے شو، که \LR{$x$} له~\LR{$c$} سره په مناسبه توګه نژدې کړو۔ په لا کره توګه، دا ټاکو چې
\LR{$\lim_{x \rightarrow c} g(x) = \ell$} به دا معنا لري:
\[
(\forall\epsilon > 0)(\exists \delta > 0)\forall x \left(0 < |x - c| < \delta \lif |g(x) - \ell| < \epsilon \right).
\]
دلته عمودي کرښې مطلقه اندازه ښيي۔ يعنې \LR{$|x| = x$} که \LR{$x \geq 0$} وي، او \LR{$|x| =-x$} که \LR{$x < 0$} وي؛ \emph{دا} تابع داسې ښودلے شئ:
\begin{center}
	\begin{tikzpicture}[scale=1]
	\draw[->, gray] (-2.5,0) -- (2.5,0) node[right] {\LR{$x$}};
	\draw[->, gray] (0,-.2) -- (0,2.5) node[above] {\LR{$|x|$}};
\olpseganchor{OLP-0635-B008}{end}

\olpseganchor{OLP-0635-B009}{start}
	\foreach \x/\xtext in {-2/-2, -1,1, 1/1, 2/2}
	\draw[shift={(\x,0)}] (0pt,2pt) -- (0pt,-2pt) node[below] {{\LR{$\xtext$}}};
\olpseganchor{OLP-0635-B009}{end}

\olpseganchor{OLP-0635-B010}{start}
	\foreach \y/\ytext in {1/1, 2/2}
	\draw[shift={(0,\y)}] (2pt,0pt) -- (-2pt,0pt) node[left] {{\LR{$\ytext$}}};
\olpseganchor{OLP-0635-B010}{end}

\olpseganchor{OLP-0635-B011}{start}
	\draw[thick] (-2.5, 2.5)--(0,0)--(2.5, 2.5);
	\end{tikzpicture}
\end{center}
نو تعريف تقريباً دا وايي: خپل «خطا» له \LR{$\epsilon$} څخه کمه کولے شئ (يعنې \LR{$|g(x) - \ell| < \epsilon$})، که داسې دليلونه وټاکئ چې له~\LR{$c$} څخه ئې واټن مثبت او تر \LR{$\delta$} کم وي (يعنې \LR{$0 < |x - c| < \delta$})۔
\olpseganchor{OLP-0635-B011}{end}

\olpseganchor{OLP-0635-B012}{start}
چې د حد مفهوم مو تعريف کړ، نو د انفينټېسيمالونو د پوره پرېښودلو لپاره يې کارولے شو، او دا ټاکلے شو چې د \LR{$f$} ميلان په \LR{$c$} کښې داسې راکول کېږي:
\[
{f}'(c) = \lim_{x \rightarrow 0}\left(\frac{f(c +x) - f(c)}{x}\right) \text{\RL{ چې حد موجود وي}}.
\]
خو دا پوهېدل مهم دي چې ولې زموږ تعريف د «چې حد موجود وي» قيد ته اړتيا لري۔ د يوې ساده بېلګې لپاره \LR{$f(x) = |x|$} په پام کښې ونيسئ، چې ګراف مو يې اوس وليد۔ ښکاره ده چې \LR{$f'(0)$} سم نۀ دے تعريف شوے: که له «ښي لوري» \LR{$0$} ته نژدې شو، ميلان تل \LR{$1$} وي؛ که له «کيڼ لوري» \LR{$0$} ته نژدې شو، ميلان تل \LR{$-1$} وي؛ نو حد نۀ دے تعريف شوے۔ له دې امله زياتولے شو چې يوه تابع~\LR{$f$} په~\LR{$x$} کښې \emph{د مشتق اخيستلو وړ} ده که او يوازې که داسې حد موجود وي۔
\olpseganchor{OLP-0635-B012}{end}

\olpseganchor{OLP-0635-B013}{start}
ومو ليدل چې د \emph{حد} په مفهوم څنګه مشتق اخيستل سمبال کړو۔ همدغه مفهوم د \emph{پيوستې} تابع د نظر د تعريف لپاره کارولے شو۔ (بولڅانو تر 1817 پورې په اصل کښې دا درک کړے و۔) د پيوستون په اړه د کوشي--وايرشټراس طريقه دا ده۔ په لنډه توګه: يوه تابع~\LR{$f$} (په يوه نقطه کښې) پيوسته ده، که د تابع د راوتلي ارزښت لپاره د يو ټاکلي دقت غوښتنه په دې تضمين کړے شئ چې د تابع د داخلېدونکي دليل لپاره پر يو ټاکلي دقت ټينګار وکړئ۔ په لا کره توګه: \LR{$f$} په \LR{$c$} کښې پيوسته ده، که، چې \LR{$x$} صفر ته ورنژدې کېږي، د \LR{$f(c + x)$} او \LR{$f(c)$} توپير خپله \LR{$0$} ته ورنژدې شي۔ په نورو ټکو: \LR{$f$} په \LR{$c$} کښې \emph{پيوسته} ده که او يوازې که \LR{$f(c) = \lim_{x \rightarrow c} f(x)$} وي۔
\olpseganchor{OLP-0635-B013}{end}

\olpseganchor{OLP-0635-B014}{start}
نور وړاندې تګ به مو حقيقي تحليل ته بوځي، حال دا چې زموږ موضوع سټ نظريه ده۔ نو اوس بايد تم شو او حاصل يې بيان کړو۔ په 19مه پېړۍ کښې رياضيپوهانو زده کړل چې د \emph{حد} د کره تعريف شوي مفهوم په کارولو څنګه له انفينټېسيمالونو پرته کار وکړي۔
\olpseganchor{OLP-0635-B014}{end}

\olpseganchor{OLP-0636-B003}{start}
\olfileid{his}{set}{pathology}
\olsection{غيرمعمولي جوړښتونه}
\olpseganchor{OLP-0636-B003}{end}

\olpseganchor{OLP-0636-B004}{start}
خو د \emph{حد} تعريف د ځينو داسې جوړښتونو اجازه ورکړه چې تر ډېره «غيرمعمولي» بلل کېدل۔
\olpseganchor{OLP-0636-B004}{end}

\olpseganchor{OLP-0636-B005}{start}
د 1830مو کلونو په شاوخوا کښې بولڅانو داسې يوه تابع وموندله چې \emph{په هره نقطه کښې پيوسته} وه، خو \emph{په هېڅ نقطه کښې يې مشتق نۀ درلود}۔ (له بده مرغه، بولڅانو دا کار هېڅکله خپور نۀ کړ؛ رياضيپوهان په لومړي ځل په 1872 کښې له دې نظر سره اشنا شول، کله چې وايرشټراس په خپلواکه توګه همدغه نظر وموند۔)\footnote{دا تاريخ د ويکيپېډيا د \href{http://en.wikipedia.org/wiki/Weierstrass_function}{وايرشټراس د تابع} په مقاله کښې په ډېر مفصلو لمنليکونو کښې بيان شوے دے۔} لږ تر لږه دا خبره ډېره حيرانوونکې وه۔ داسې تابعې موندل اسانه دي، لکه \LR{$|x|$}، چې په هره نقطه کښې پيوستې وي خو په يوه ځانګړې نقطه کښې مشتق نۀ لري۔ خو داسې تابع چې په هره نقطه کښې پيوسته وي او په \emph{هېڅ} نقطه کښې مشتق نۀ لري، ډېر بل ډول څيز دے۔ يوه شېبه فکر وکړئ چې د داسې تابع د رسمولو هڅه به څنګه کوئ۔ د پيوستون لپاره بايد دا د پاڼې څخه د قلم پورته کولو پرته رسم کړے شئ؛ خو د دې لپاره چې په هېڅ نقطه کښې مشتق ونۀ لري، بايد د خپل قلم لورے په پرله‌پسې توګه ناڅاپه بدل کړئ۔
\olpseganchor{OLP-0636-B005}{end}

\olpseganchor{OLP-0636-B006}{start}
نور «غيرمعمولي جوړښتونه» هم راپسې راغلل۔ د 1874 د جنورۍ په 5مه نېټه کانتور ډېډېکنډ ته يو ليک وليکه او دا مسئله يې وړاندې کړه:
\begin{quote}
ايا يوه سطحه، لکه له خپلې څنډې سره يوه مربع، له يوه
خط سره، لکه له دواړو پای‌ټکو سره يوه مستقيمه ټوټه، په داسې
يو پر يو تړاو کښې راتلے شي چې د سطحې له هرې نقطې سره د
خط يوه نقطه برابره وي او د خط له هرې نقطې سره د سطحې يوه
نقطه؟
\olpseganchor{OLP-0636-B006}{end}

\olpseganchor{OLP-0636-B007}{start}
اوس هم داسې راته ښکاري چې دې پوښتنې ته ځواب ويل ډېر ګران
دي۔ که څه هم دلته هم سړے دومره \emph{نه} ويلو ته مايل کېږي
چې د ثبوت اړتيا ورته نژدې له منځه تللې ښکاري۔ د دې ليک د
تاريخي بحث لپاره \mbox{\LR{\latinfont \citealt{Gouvea2011}}} وګورئ۔
\end{quote}
خو په 1877 کښې کانتور ثابته کړه چې د هغه پخوانے ګومان ناسم و۔ په رښتيا، يوه مستقيمه ټوټه او يوه مربع عين يو هومره نقطې لري۔ د 1877 د جون په 29مه نېټه يې ډېډېکنډ ته وليکل: «\emph{\mbox{\LR{\latinfont je le vois, mais je ne le crois pas}}}»؛ يعنې «وينم يې، خو باور پرې نۀ کوم»۔ د رياضياتو په «رواجي تاريخ» کښې دا خبره ډېر ځله د دې نښې په توګه اخيستل کېږي چې دا نوې نتيجې د هغه وخت د رياضيپوهانو لپاره څومره \emph{په لفظي معنا د باور وړ نۀ وې}۔ (دا ليکونه په \mbox{\LR{\latinfont \citet{Gouvea2011}}} کښې وړاندې شوي، او په \olref[his][set][mythology]{sec} کښې ورته بېرته راځو۔ د کانتور د ثبوت خاکه په \olref[his][set][cantorplane]{sec} کښې شته۔)
\olpseganchor{OLP-0636-B007}{end}

\olpseganchor{OLP-0636-B008}{start}
د کانتور له نتيجې څخه په الهام، پېانو وپوښتل چې ايا له يوې
تړلې عددي وقفې څخه داسې \emph{پيوسته} تصوير شته چې د مربع
هرې نقطې ته ورسېږي۔ داسې تصوير ته \emph{ځای ډکوونکے منحني خط}
وايو۔ په \mbox{\LR{\latinfont \citeyear{Peano1890}}} کښې پېانو همداسې منحني جوړه کړه۔
پيوستون د دې جوړښت شرط دے۔ که د منحني مشتق هم پيوسته وے، نو
انځور به يې صفر مساحت لرلے او مربع به يې نۀ شوے ډکولے۔
اقليدس خط «بې پلنوالي اوږدوالے» تعريف کړے و (لومړے کتاب،
تعريف 2)۔ د پېانو منحني بيا د مربع پر هرې نقطې تېرېږي۔ په \mbox{\LR{\latinfont \citeyear{Hilbert1891}}} کښې هېلبرټ يو څه زيات له وجدان سره سم ځای ډکوونکے منحني خط بيان کړ او د هغه د څرګندولو لپاره يې ځينې انځورونه هم وړاندې کړل۔ دا منحني خط په پرله‌پسې پړاوونو کښې جوړېږي، او دلته د جوړښت لومړني شپږ پړاوونه دي:
\begin{center}
\begin{tikzpicture}
\begin{scope}[xshift=20pt, yshift=20pt]
	\draw [oldiagcolorC, l-system={Hilbert curve, step=40pt, angle=90, axiom=L, order=1}] lindenmayer system;
\end{scope}
\begin{scope}[xshift=110pt, yshift=10pt]
	\draw [oldiagcolorC, l-system={Hilbert curve, step=20pt, angle=90, axiom=L, order=2}] lindenmayer system;
\end{scope}
\begin{scope}[xshift=205pt, yshift=5pt]
	\draw [oldiagcolorC, l-system={Hilbert curve, step=10pt, angle=90, axiom=L, order=3}] lindenmayer system;
\end{scope}
\begin{scope}[xshift=2.5pt, yshift=-97.5pt]
	\draw [oldiagcolorC, l-system={Hilbert curve, step=5pt, angle=90, axiom=L, order=4}] lindenmayer system;
\end{scope}
\begin{scope}[xshift=101.25pt, yshift=-98.75pt]
	\draw [oldiagcolorC, l-system={Hilbert curve, step=2.5pt, angle=90, axiom=L, order=5}] lindenmayer system;
\end{scope}
\begin{scope}[xshift=200.625pt, yshift=-99.375pt]
	\draw [oldiagcolorC, l-system={Hilbert curve, step=1.25pt, angle=90, axiom=L, order=6}] lindenmayer system;
\end{scope}
\draw[gray] (0pt,0pt) rectangle (80pt, 80pt);
\draw[gray] (100pt,0pt) rectangle (180pt, 80pt);
\draw[gray] (200pt,0pt) rectangle (280pt, 80pt);
\draw[gray] (0pt,-100pt) rectangle (80pt, -20pt);
\draw[gray] (100pt,-100pt) rectangle (180pt, -20pt);
\draw[gray] (200pt,-100pt) rectangle (280pt, -20pt);
\end{tikzpicture}  
\end{center}
په حد کښې — هغه مفهوم چې تر دې وخته يې کره تعريف ترلاسه کړے و — ټوله مربع په بشپړ ډول په سور رنګ ډکېږي۔ د يادونې وړ دا هم ده چې د هېلبرټ منحني خط په هره نقطه کښې پيوسته دے، خو په هېڅ نقطه کښې مشتق نۀ لري؛ د وجدان له مخې يې وجه دا ده چې په نامتناهي حد کښې تابع په هره شېبه کښې ناڅاپه خپل لورے بدلوي۔ (د هېلبرټ د جوړښت لا مفصله خاکه به په \olref[his][set][hilbertcurve]{sec} کښې وړاندې کړو۔)
\olpseganchor{OLP-0636-B008}{end}

\olpseganchor{OLP-0636-B009}{start}
د ښه يا بد لپاره، دا «غيرمعمولي» هندسي جوړښتونه پر هندسي وجدان د تکيه کولو د شک لپاره د دليل په توګه واخيستل شول۔ دا د رياضياتو د يوې نوې طريقې لپاره د \emph{تبليغاتو} تر کچې ورنژدې شول؛ داسې طريقه چې په پای کښې به سټ نظريې ته ورسېږي۔ وروسته، د دې مضمون په اسطوري بيانونو کښې، دا خبره بيا بيا وشوه چې دا نتيجې هم پوره کره وې او هم پوره ټکان ورکوونکې۔ ځکه نو دوه موخې يې ترسره کولې: پر هندسي وجدان د تکيه کولو پر ضد خبردارے، او د فکر د نويو طريقو د حاصل ورکولو څرګندونه۔
\olpseganchor{OLP-0636-B009}{end}

\olpseganchor{OLP-0637-B003}{start}
\olfileid{his}{set}{mythology}
\olpseganchor{OLP-0637-B003}{end}

\olpseganchor{OLP-0637-B004}{start}
\olsection{تر تاريخ زيات اسطوري بيان؟}
\olpseganchor{OLP-0637-B004}{end}

\olpseganchor{OLP-0637-B005}{start}
چې له يوې پېړۍ څخه د زيات وخت له واټنه دې پېښو ته بېرته ګورو، بايد پام وکړو چې دا پرېکړې بې له پوښتنې ونۀ منو۔ نتيجې بې‌شکه نوې، په زړه پورې او حيرانوونکې وې۔ خو په رښتيا څومره ټکان ورکوونکې وې؟ او آيا په رښتيا يې ښودله چې پر هندسي وجدان بايد تکيه ونۀ کړو؟
\olpseganchor{OLP-0637-B005}{end}

\olpseganchor{OLP-0637-B006}{start}
د ټکان په مسئله کښې، \mbox{\LR{\latinfont \citet{Gouvea2011}}} د کانتور د دې مشهور عبارت سياق څېړي۔
«\emph{\mbox{\LR{\latinfont je le vois, mais je ne le crois pas}}}» هماغه عبارت دے چې کانتور ډېډېکنډ ته وليکه۔
دا عبارت ډېر ځله له خپل سياق څخه بهر اخيستل کېږي۔ د 1877 د جون د 29مې نېټې په ليک کښې کانتور داسې ليکي (د دې سياق د شننې لپاره \mbox{\LR{\latinfont \citeauthor{Gouvea2011}}} وګورئ):
\begin{quote}
که زه وار په وار ستاسو له مهربانۍ او زيار څخه کار اخلم، نو
د دې مسئلې لپاره زما شوق راوبخښئ۔ هغه مطالب چې په دې نژدې
وخت کښې مې دراستولي دي، زما خپل ځان ته هم دومره ناڅاپي او
دومره نوي دي چې زړه به مې تر هغې يو څه ډاډه نۀ شي، څو له
تاسو، قدرمنه ملګريه، د هغو د سموالي په باب فيصله ترلاسه
نۀ کړم۔ تر څو چې تاسو راسره موافق نۀ يئ، همدا ويلے شم۔
\emph{\mbox{\LR{\latinfont je le vois, mais je ne le crois pas.}}} [وينم يې، خو باور پرې نۀ کوم۔]
\end{quote}
کانتور پوهېده چې نتيجه يې «دومره ناڅاپي، دومره نوې» وه۔ خو دا شکمنه ده چې هغه کله هم خپله نتيجه \emph{د باور وړ نۀ} ګڼلې وي۔ لکه چې \mbox{\LR{\latinfont \citeauthor{Gouvea2011}}} څرګندوي، هغه يوازې له ډېډېکنډ څخه غوښتل چې وړاندې کړے ثبوت يې وګوري۔
\olpseganchor{OLP-0637-B006}{end}

\olpseganchor{OLP-0637-B007}{start}
پېانو خپل ځای ډکوونکے منحني خط بې له انځورونو خپور کړ۔
هېلبرټ په 1891 کښې وليکل چې د دغسې تصوير دپاره اړينې تابعې
په هندسي طريقه په لا څرګند ډول جوړېدے شي۔ هغه خپله هندسي
طريقه په يو لړ \emph{انځورونو} کښې وښوده، لکه د
\olref[his][set][pathology]{sec} انځورونه۔
\olpseganchor{OLP-0637-B007}{end}

\olpseganchor{OLP-0637-B008}{start}
په عمومي توګه: که څه هم په خپرو شويو ثبوتونو کښې د انځورونو کارول يو څه له رواجه لوېدلي، خو له دې حقيقت څخه انکار نۀ شي کېدای چې رياضيپوهان د ثبوت پر وخت \emph{ډېر ځله} انځورونه کاروي۔ (په تقريبي ټکو: ښه رياضيپوهان پوهېږي چې کله پر هندسي وجدان تکيه کولے شي۔)
\olpseganchor{OLP-0637-B008}{end}

\olpseganchor{OLP-0637-B009}{start}
په لنډه توګه: پر تشې شورماشور باور مۀ کوئ؛ يا لږ تر لږه يې بې له پوښتنې مۀ منئ۔ د دې په اړه د نور لوست لپاره \mbox{\LR{\latinfont \citet{Giaquinto2007}}} کتلے شئ۔
\olpseganchor{OLP-0637-B009}{end}

\olpseganchor{OLP-0638-B003}{start}
\olfileid{his}{set}{cantorplane}
\olsection{د خط او مستوي په اړه د کانتور کار}
\olpseganchor{OLP-0638-B003}{end}

\olpseganchor{OLP-0638-B004}{start}
د شروډر--برنشتاين د ثبوت ځينې اړوند حالات له هغه تاريخ سره تړاو لري چې په \olref[his][set][pathology]{sec} کښې مو بيان کړ۔ په ياد ولرئ چې په 1877 کښې کانتور ثابته کړه چې په يوه مربع کښې عين يو هومره نقطې دي څومره چې يې په يوې څنډې کښې دي۔ دلته به د هغه ثبوت (لومړۍ هڅه) وړاندې کړو۔
\olpseganchor{OLP-0638-B004}{end}

\olpseganchor{OLP-0638-B005}{start}
فرض کړئ چې \LR{$\unitline$} واحد خط وي، يعنې د نقطو سټ \LR{$[0,1]$}۔ فرض کړئ چې \LR{$\unitsquare$} واحده مربع وي، يعنې د نقطو سټ \LR{$\unitline
\times \unitline$}۔ په دې ټکو، کانتور ثابته کړه چې \LR{$\cardeq{\unitline}{\unitsquare}$}۔ کانتور ډېډېکنډ ته په خپل ليک کښې د اعشاري رقمونو د اوبدلو دليل وړاندې کړ۔ دلته يې مفکوره په دوه‌ييزه نښه‌ليکنه کښې بيانوو۔
\olpseganchor{OLP-0638-B005}{end}

\olpseganchor{OLP-0638-B006}{start}
\begin{thm}\ollabel{thm:cantorplane}
\LR{$\cardeq{\unitline}{\unitsquare}$}
\end{thm}
\olpseganchor{OLP-0638-B006}{end}

\olpseganchor{OLP-0638-B007}{start}
\begin{proof}[ثبوت: لومړۍ برخه۔]
\LR{$a, b \in \unitline$} ثابت کړئ۔ په دوه‌ييزه نښه‌ليکنه يې وليکئ، څو د \LR{$0$}ونو او \LR{$1$}ونو بې‌پايه لړۍ، \LR{$a_1$}، \LR{$a_2$}، \dots، او \LR{$b_1$}، \LR{$b_2$}، \dots، ولرو۔ د \LR{$1$} نه د کوچني عدد لپاره هغه يوازينۍ پراختيا اخلو چې له کوم ځايه وروسته ئې ټول رقمونه \LR{$1$} نۀ وي۔ د بېلګې په توګه، د نيمې لپاره \LR{$0.1000\ldots$} اخلو۔ د \LR{$1$} لپاره \LR{$0.1111\ldots$} اخلو۔ په دې توګه:
\begin{align*}
a &= 0.a_1a_2a_3a_4\dots\\
b &= 0.b_1b_2b_3b_4\dots
\intertext{\RL{اوس تابع \LR{$f \colon \unitsquare \to \unitline$} په پام کښې ونيسئ، چې داسې راکول کېږي:}} 
f(a, b) & = 0.a_1b_1a_2b_2a_3b_3a_4b_4\dots
\end{align*}
که د \LR{$f(a,b)$} له کوم رقم وروسته ټول رقمونه \LR{$1$} وي، نو د \LR{$a$} او \LR{$b$} دواړو په پراختياوو کښې هم له کوم ځايه وروسته ټول رقمونه \LR{$1$} دي؛ له دې امله \LR{$a=b=1$} او \LR{$f(a,b)=1$}۔ په بل هر حالت کښې راوتلې پراختيا هماغه ټاکل شوې پراختيا ده۔ اوس \LR{$f$} يو پر يو تابع دے، ځکه که \LR{$f(a, b) = f(c,d)$} وي، نو \LR{$a_n =
c_n$} او \LR{$b_n = d_n$} د هر \LR{$n \in \Nat$} لپاره دي؛ ځکه نو \LR{$a = c$} او \LR{$b =
d$} دي۔
\end{proof}
\olpseganchor{OLP-0638-B007}{end}

\olpseganchor{OLP-0638-B008}{start}
ډېډېکنډ کانتور ته څرګنده کړه چې د رقمونو د اوبدلو پر مهال د يو عدد دوه پراختياوې ستونزه جوړوي۔ د پورته دوه‌ييز انتخاب لپاره \LR{$0.0\overline{01}=0.00101010\ldots$} په پام کښې ونيسئ۔ دغه عدد يوازې يوه دوه‌ييزه پراختيا لري، ځکه رقمونه ئې په پرله‌پسې صفرونو يا پرله‌پسې يوونو نۀ ختمېږي۔ که د \LR{$f(a,b)$} ارزښت همدا وای، نو جفت او طاق ځايونه به په ترتيب د \LR{$b$} او \LR{$a$} رقمونه ټاکل:
\begin{align*}
a&=0.0111\ldots=\frac{1}{2}\\
b&=0.
\end{align*}
خو د نيمې لپاره مو \LR{$0.1000\ldots$} ټاکلې ده؛ نو دغه راوتلے عدد د کومې جوړې تصوير نۀ دے۔
\olpseganchor{OLP-0638-B008}{end}

\olpseganchor{OLP-0638-B009}{start}
په دې دوه‌ييزه بڼه کښې د رقمونو د اوبدلو تابع \LR{$f$} يو پر يو تابع دے، خو پر هدف سټ بشپړه تابع \emph{نۀ} دے۔ نو کانتور يوازې \LR{$\cardle{\unitsquare}{\unitline}$} ښودلې ده، \emph{نۀ} دا چې \LR{$\cardeq{\unitsquare}{\unitline}$}۔
\olpseganchor{OLP-0638-B009}{end}

\olpseganchor{OLP-0638-B010}{start}
کانتور نژدې سمدستي ډېډېکنډ ته ځواب وليکه او په اصل کښې يې وړانديز وکړ چې ثبوت په لاندې ډول بشپړېدای شي:
\olpseganchor{OLP-0638-B010}{end}

\olpseganchor{OLP-0638-B011}{start}
\begin{proof}[ثبوت: بشپړه برخه۔]
نو موږ ښودلې ده چې \LR{$\cardle{\unitsquare}{\unitline}$}۔ خو ښکاره ده چې له \LR{$\unitline$} څخه \LR{$\unitsquare$} ته يو پر يو تابع شته: يوازې خط د مربع د يوې څنډې پر اوږدوالي کېږدئ۔ نو \LR{$\cardle{\unitline}{\unitsquare}$} او \LR{$\cardle{\unitsquare}{\unitline}$} دي۔ د شروډر--برنشتاين له مخې (\olref[sfr][siz][sb]{thm:schroder-bernstein})، \LR{$\cardeq{\unitline}{\unitsquare}$}۔
\end{proof}
\olpseganchor{OLP-0638-B011}{end}

\olpseganchor{OLP-0638-B012}{start}
خو، څرګنده ده چې کانتور وروستۍ کرښه په دې ټکو نۀ شوه بشپړولے، ځکه د شروډر--برنشتاين قضيه لا نۀ وه ثابته شوې۔ په رښتيا، که څه هم کانتور وروسته دا د يوې عمومي فرضيې په توګه وړاندې کړه، خو تر 1897 پورې يې د قناعت وړ ثبوت نۀ و ورکړل شوے۔ (ځکه نو، وروسته په 1877 کښې، کانتور د \olref{thm:cantorplane} يو بل ثبوت وړاندې کړ چې د شروډر--برنشتاين له لارې نۀ و۔)
\olpseganchor{OLP-0638-B012}{end}

\olpseganchor{OLP-0639-B003}{start}
\olfileid{his}{set}{hilbertcurve}
\olpseganchor{OLP-0639-B003}{end}

\olpseganchor{OLP-0639-B004}{start}
\olsection{ضميمه: د هلبرټ ځای ډکوونکي منحني خطونه}
\olpseganchor{OLP-0639-B004}{end}

\olpseganchor{OLP-0639-B005}{start}
په \olref[pathology]{sec} څپرکي کښې مو يادونه وکړه چې د کانتور ثبوت، چې خط او مربع عين يو هومره نقطې لري (\olref[his][set][cantorplane]{thm:cantorplane})، پيانو دې پوښتنې ته وهڅاوه چې آيا ځای ډکوونکے \emph{منحني خط} به موجود وي۔ هغه په 1890 کښې مثبت ځواب ترلاسه کړ۔ په دې برخه کښې به (په غيررسمي او لنډ وضاحتي ډول) وښيو چې د هلبرټ ځای ډکوونکے منحني خط څنګه جوړېږي، له يو ډېر وړوکي بدلون سره۔\footnote{د لا کره تشريح لپاره \mbox{\LR{\latinfont \cite{Rose2010}}} وګورئ۔ دا بدلون د لاندې منحني خطونو د سرو برخو زياتول دي۔ له دې سره د منحني خط پيوسته والے کتل لږ آسانه کېږي۔}
\olpseganchor{OLP-0639-B005}{end}

\olpseganchor{OLP-0639-B006}{start}
بايد تابع \LR{$h$} د تابعو د يوې لړۍ \LR{$h_1$}، \LR{$h_2$}، \LR{$h_3$}، \dots\@ د حد په توګه تعريف کړو۔ لومړے به يې جوړول بيان کړو۔ بيا به وښيو چې حد شته، ځای ډکوي او پيوسته دے۔
\olpseganchor{OLP-0639-B006}{end}

\olpseganchor{OLP-0639-B007}{start}
د \LR{$h$} د قيمتونو سټ به واحده مربع، \LR{$\unitsquare$}، واخلو۔ د \LR{$h$} لومړنے تقريب، يعنې \LR{$h_1$}، دا دے:
\begin{center}
	\begin{tikzpicture}
	\draw[gray] (0pt, 0pt) rectangle (80pt, 80pt);
	\draw[step = 40pt, gray, very thin] (0pt, 0pt) grid (80pt, 80pt);
	\draw[oldiagcolorC, thick] (20pt, 0)--(20pt, 20pt);
	\draw[oldiagcolorC, thick] (60pt, 0)--(60pt, 20pt);		
	\begin{scope}[xshift=20pt, yshift=20pt]
	\draw [black, thick, l-system={Hilbert curve, step=40pt, angle=90, axiom=L, order=1}]  lindenmayer system;
	\end{scope}
	\end{tikzpicture}
\end{center}
د جوړښت د څارلو لپاره مو په مربع د \LR{$2 \times 2$} خانه‌جال ايښے دے۔ داسې يې تصور کولے شو چې منحني خط په لاندې کيڼې ربع کښې پيلېږي، پورته کيڼې ته ځي، بيا پورته ښۍ ته، او په پای کښې لاندې ښۍ ته۔ د جوړښت دويم پړاو، يعنې \LR{$h_2$}، دا دے:
\begin{center}
	\begin{tikzpicture}
	\draw[gray] (0pt, 0pt) rectangle (80pt, 80pt);
	\draw[step = 20pt, gray, very thin] (0pt, 0pt) grid (80pt, 80pt);
	\draw[oldiagcolorC, thick] (0pt, 10pt)--(10pt, 10pt);
	\draw[oldiagcolorC, thick] (70pt, 10pt)--(80pt, 10pt);
	\draw[thick, oldiagcolorE] (10pt, 30pt)--(10pt, 50pt); 
	\draw[thick, oldiagcolorE] (30pt, 50pt)--(50pt, 50pt); 
	\draw[thick, oldiagcolorE] (70pt, 50pt)--(70pt, 30pt); \begin{scope}[xshift=10pt, yshift=30pt]
	\draw [thick, rotate=270,black, l-system={Hilbert curve, step=20pt, angle=90, axiom=L, order=1}]  lindenmayer system;
	\end{scope}
	\begin{scope}[xshift=10pt, yshift=50pt]
	\draw [thick, black, l-system={Hilbert curve, step=20pt, angle=90, axiom=L, order=1}]  lindenmayer system;
	\end{scope}
	\begin{scope}[xshift=50pt, yshift=50pt]
	\draw [thick, black, l-system={Hilbert curve, step=20pt, angle=90, axiom=L, order=1}]  lindenmayer system;
	\end{scope}
	\begin{scope}[xshift=70pt, yshift=10pt]
	\draw [thick, rotate=90,black, l-system={Hilbert curve, step=20pt, angle=90, axiom=L, order=1}]  lindenmayer system;
	\end{scope}
	\end{tikzpicture}
\end{center}
بېل بېل رنګونه به د \LR{$h_2$} د جوړولو په تشريح کښې مرسته وکړي۔ لومړے د \LR{$h_1$} د غير سور برخې په کوچنۍ کچه نقلونه د مربع په لاندې کيڼې، پورته کيڼې، پورته ښۍ او لاندې ښۍ برخو کښې ږدو؛ په شکل کښې په تور رنګ رسم شوي دي۔ بيا دا څلور شکلونه په زرغون خطونو نښلوو۔ په پای کښې خپل شکل د مربع له پولې سره په سور خطونو نښلوو۔
\olpseganchor{OLP-0639-B007}{end}

\olpseganchor{OLP-0639-B008}{start}
اوس \LR{$h_3$} ته راځو۔ لکه چې \LR{$h_2$} د \LR{$h_1$} د غيرسرې برخې له څلورو نښلول شويو، په کوچنۍ کچه نقلونو جوړ شوے و، همداسې \LR{$h_3$} د \LR{$h_2$} د غيرسرې برخې له څلورو په کوچنۍ کچه نقلونو جوړېږي؛ په تور رنګ رسم شوي دي۔ بيا په زرغون خطونو يو بل سره نښلول کېږي او په پای کښې د سور خطونو په مرسته د مربع له پولې سره نښلول کېږي۔
\begin{center}
	\begin{tikzpicture}
	\draw[gray] (0pt, 0pt) rectangle (80pt, 80pt);
	\draw[step = 10pt, gray, very thin] (0pt, 0pt) grid (80pt, 80pt);
	\draw[thick, oldiagcolorC](5pt, 0pt)--(5pt, 5pt); 
	\draw[thick, oldiagcolorC] (75pt, 0pt)--(75pt, 5pt); 
	\draw[thick, oldiagcolorE] (75pt, 45pt)--(75pt, 35pt);
	\draw[thick, oldiagcolorE] (5pt, 35pt)--(5pt, 45pt); 
	\draw[thick, oldiagcolorE] (35pt, 45pt)--(45pt, 45pt); 
	\draw[thick, oldiagcolorE] (75pt, 45pt)--(75pt, 35pt); \begin{scope}[xshift=5pt, yshift=35pt]
	\draw [rotate=270, thick, black, l-system={Hilbert curve, step=10pt, angle=90, axiom=L, order=2}]  lindenmayer system;
	\end{scope}
	\begin{scope}[xshift=5pt, yshift=45pt]
	\draw [thick, black, l-system={Hilbert curve, step=10pt, angle=90, axiom=L, order=2}]  lindenmayer system;
	\end{scope}
	\begin{scope}[xshift=45pt, yshift=45pt]
	\draw [thick, black, l-system={Hilbert curve, step=10pt, angle=90, axiom=L, order=2}]  lindenmayer system;
	\end{scope}
	\begin{scope}[xshift=75pt, yshift=5pt]
	\draw [thick, rotate=90,black, l-system={Hilbert curve, step=10pt, angle=90, axiom=L, order=2}]  lindenmayer system;
	\end{scope}
	\end{tikzpicture}
\end{center}
د \LR{$h_n$} د پارامتري کولو لپاره واحده وقفه په \LR{$4^n$} برابرو تړلو برخو وېشو۔ هره برخه د منحني خط د تېرېدو په ترتيب له يوې خانې سره اړخ لګوو۔ د هغې برخې په نيمه کښې له ورودي ټکي څخه د خانې مرکز ته په مستقيم خط ځو، او په بله نيمه کښې له مرکز څخه وتونکي ټکي ته ځو۔ د دوو پرله‌پسې خانو تر منځ ګډ ټکے د ګډې څنډې منځ دے۔ د لومړۍ او وروستۍ خانې دباندې ټکي د رسم د سرو برخو په شان د مربع پر پوله اخلو۔ په ګډو پای ټکو کښې دواړه پارامتري ټوټې سره يو ځای کېږي؛ نو هر \LR{$h_n$} پيوسته دے۔
د هر \LR{$x\in\unitline$} لپاره به وښيو چې لاندې حد شته:
\begin{align*}
 h(x)&=\lim_{n\rightarrow\infty} h_n(x).
\end{align*}
په بل پړاو کښې د هرې برخې څلور وړې برخې په ترتيب د هماغې خانې څلورو وړو خانو ته ورځي۔ که \LR{$x$} د دوو برخو ګډ پای وي، د اړوندو تړلو خانو ګډه پوله ئې تصوير رانغاړي۔ نو که \LR{$m\ge n$} وي، \LR{$h_n(x)$} او \LR{$h_m(x)$} دواړه د \LR{$n$}م پړاو په يوه تړلې خانه کښې دي۔ د هغې خانې د قطر اوږدوالے دے:
\[
\sqrt{\left(\nicefrac{1}{2^n}\right)^2+\left(\nicefrac{1}{2^n}\right)^2}=2^{\frac{1}{2}-n}.
\]
نو د دواړو تصويرونو واټن تر دغه حده زيات نۀ دے۔ دغه اټکل د \LR{$x$} له ټاکنې سره تړلے نۀ دے، او
\[
\lim_{n\rightarrow\infty}2^{\frac{1}{2}-n}=0.
\]
د \LR{$\Real^2$} د بشپړوالي له امله د هر \LR{$x\in\unitline$} لپاره حد شته او تقارب يوشان دے۔
د ځای ډکولو لپاره په \LR{$\unitsquare$} کښې يو ټکے \LR{$p$} واخلئ۔ په لومړي پړاو کښې يوه تړلې خانه وټاکئ چې \LR{$p$} پکښې وي، او په هر بل پړاو کښې د ټاکل شوې خانې يوه وړه خانه وټاکئ چې \LR{$p$} هم پکښې وي۔ د دغو خانو اړوندې تړلې پارامتري برخې يوې بلې کښې پرتې دي او اوږدوالے ئې صفر ته نژدې کېږي۔ په ټولو برخو کښې يو ګډ \LR{$x$} شته۔ د هر \LR{$n$} لپاره \LR{$p$} او \LR{$h_n(x)$} په هماغې ټاکل شوې خانې کښې دي؛ د دواړو واټن تر \LR{$2^{\frac{1}{2}-n}$} زيات نۀ دے۔ نو \LR{$h(x)=p$}۔
\olpseganchor{OLP-0639-B008}{end}

\olpseganchor{OLP-0639-B009}{start}
دا ښودل پاتې دي چې \LR{$h$} په رښتيا يو \emph{منحني خط} دے۔ د دې لپاره بايد دغه مفهوم تعريف کړو۔ نوے تعريف پر هغه تعريف ولاړ دے چې ژورډان په 1887 کښې ورکړے و، يعنې د لومړي ځای ډکوونکي منحني خط له وړاندې کولو څخه يوازې څو کاله مخکې:
\olpseganchor{OLP-0639-B009}{end}

\olpseganchor{OLP-0639-B010}{start}
\begin{defn}
منحني خط له \LR{$\unitline$} څخه \LR{$\Real^2$} ته يو پيوسته تصوير دے۔
\end{defn}
\olpseganchor{OLP-0639-B010}{end}

\olpseganchor{OLP-0639-B011}{start}
زموږ تابع \LR{$h$} له \LR{$\unitline$} څخه \LR{$\Real^2$} ته تصوير دے۔ د پيوستون د کتلو لپاره، لکه په \olref[limits]{sec} کښې، يو ټاکلے \LR{$x_0\in\unitline$} او يو \LR{$\epsilon>0$} واخلئ۔ په يوشان تقارب سره داسې \LR{$n$} ټاکو چې د \LR{$h_n$} او \LR{$h$} د تصويرونو واټن په ټوله وقفه کښې تر \LR{$\epsilon/3$} کم وي۔ د \LR{$h_n$} د پيوستون له امله داسې \LR{$\delta>0$} شته چې د \LR{$|x-x_0|<\delta$} په صورت کښې د \LR{$h_n(x)$} او \LR{$h_n(x_0)$} واټن تر \LR{$\epsilon/3$} کم وي۔
\olpseganchor{OLP-0639-B011}{end}

\olpseganchor{OLP-0639-B012}{start}
د مثلث نابرابري د \LR{$h(x)$} او \LR{$h(x_0)$} واټن په درې واټنونو محدودوي: له \LR{$h(x)$} تر \LR{$h_n(x)$}، له \LR{$h_n(x)$} تر \LR{$h_n(x_0)$}، او له \LR{$h_n(x_0)$} تر \LR{$h(x_0)$}۔ هر يو له \LR{$\epsilon/3$} څخه کم دے، نو ټول واټن له \LR{$\epsilon$} څخه کم دے۔
\olpseganchor{OLP-0639-B012}{end}

\olpseganchor{OLP-0639-B013}{start}
دغه دليل د واحدې وقفې پر هر ټکي چلېږي۔ نو \LR{$h$} يو پيوسته ځای ډکوونکے منحني خط دے۔
\olpseganchor{OLP-0639-B013}{end}

\olpseganchor{OLP-0640-B004}{start}
\olpart{ref}{مرستندويه معلومات}
\olpseganchor{OLP-0640-B004}{end}

\olpseganchor{OLP-0640-B005}{start}
\begin{editorial}
  دا برخه د الفباوو، نښه‌ليکنو، تعريفونو، قاعدو او نورو داسې څيزونو بېل بېل لستونه راغونډوي، څو د درسي کتاب د ضميمو په توګه پکې شامل شي۔
\end{editorial}
\olpseganchor{OLP-0640-B005}{end}

\olpseganchor{OLP-0641-B004}{start}
\olchapter{ref}{grk}{يوناني الفبا}
\olpseganchor{OLP-0641-B004}{end}

\olpseganchor{OLP-0641-B005}{start}
\begin{center}
    \begin{tabular}{lll@{\qquad}lll}
    الفا &	\LR{$\alpha$} & \LR{$A$} &
    نيو & \LR{$\nu$} & \LR{$N$} \\
    بېټا & \LR{$\beta$} & \LR{$B$} &
    کسي & \LR{$\xi$} & \LR{$\Xi$} \\
    ګاما &	\LR{$\gamma$} & \LR{$\Gamma$} &
    اوميکرون & \LR{$o$} & \LR{$O$} \\
    ډېلټا & \LR{$\delta$} & \LR{$\Delta$} &
    پاي & \LR{$\pi$} & \LR{$\Pi$} \\
    اېپسيلون & \LR{$\epsilon$} & \LR{$E$} &
    رو & \LR{$\rho$} & \LR{$P$} \\
    زېټا & \LR{$\zeta$} & \LR{$Z$} &
    سېګما &	\LR{$\sigma$} & \LR{$\Sigma$} \\
    اېټا & \LR{$\eta$} & \LR{$H$} &
    تاو & \LR{$\tau$} & \LR{$T$} \\
    تېټا & \LR{$\theta$} & \LR{$\Theta$} &
    اوپسيلون & \LR{$\upsilon$} & \LR{$\Upsilon$} \\
    ايوټا & \LR{$\iota$} & \LR{$I$} &
    فاي & \LR{$\varphi$} & \LR{$\Phi$} \\
    کاپا & \LR{$\kappa$} & \LR{$K$} &
    کاي & \LR{$\chi$} & \LR{$X$} \\
    لامبډا & \LR{$\lambda$} & \LR{$\Lambda$} &
    پسي & \LR{$\psi$} & \LR{$\Psi$} \\
    ميو & \LR{$\mu$} & \LR{$M$} &
    اوميګا & \LR{$\omega$} & \LR{$\Omega$}
    \end{tabular}
\end{center}
\olpseganchor{OLP-0641-B005}{end}

\olpseganchor{OLP-0642-B004}{start}
\olchapter{ref}{frk}{د فراکتور الفبا}
\olpseganchor{OLP-0642-B004}{end}

\olpseganchor{OLP-0642-B005}{start}
\begin{center}
    \begin{tabular}{llll@{\qquad}llll}
    A &	\LR{$\mathfrak{A}$} & a & \LR{$\mathfrak{a}$} &
    N &	\LR{$\mathfrak{N}$} & n & \LR{$\mathfrak{n}$} \\
    B &	\LR{$\mathfrak{B}$} & b & \LR{$\mathfrak{b}$} &
    O &	\LR{$\mathfrak{O}$} & o & \LR{$\mathfrak{o}$} \\
    C &	\LR{$\mathfrak{C}$} & c & \LR{$\mathfrak{c}$} &
    P &	\LR{$\mathfrak{P}$} & p & \LR{$\mathfrak{p}$} \\
    D &	\LR{$\mathfrak{D}$} & d & \LR{$\mathfrak{d}$} &
    Q &	\LR{$\mathfrak{Q}$} & q & \LR{$\mathfrak{q}$} \\
    E &	\LR{$\mathfrak{E}$} & e & \LR{$\mathfrak{e}$} &
    R &	\LR{$\mathfrak{R}$} & r & \LR{$\mathfrak{r}$} \\
    F &	\LR{$\mathfrak{F}$} & f & \LR{$\mathfrak{f}$} &
    S &	\LR{$\mathfrak{S}$} & s & \LR{$\mathfrak{s}$} \\
    G &	\LR{$\mathfrak{G}$} & g & \LR{$\mathfrak{g}$} &
    T &	\LR{$\mathfrak{T}$} & t & \LR{$\mathfrak{t}$} \\
    H &	\LR{$\mathfrak{H}$} & h & \LR{$\mathfrak{h}$} &
    U &	\LR{$\mathfrak{U}$} & u & \LR{$\mathfrak{u}$} \\
    I &	\LR{$\mathfrak{I}$} & i & \LR{$\mathfrak{i}$} &
    V &	\LR{$\mathfrak{V}$} & v & \LR{$\mathfrak{v}$} \\
    J &	\LR{$\mathfrak{J}$} & j & \LR{$\mathfrak{j}$} &
    W &	\LR{$\mathfrak{W}$} & w & \LR{$\mathfrak{w}$} \\
    K &	\LR{$\mathfrak{K}$} & k & \LR{$\mathfrak{k}$} &
    X &	\LR{$\mathfrak{X}$} & x & \LR{$\mathfrak{x}$} \\
    L &	\LR{$\mathfrak{L}$} & l & \LR{$\mathfrak{l}$} &
    Y &	\LR{$\mathfrak{Y}$} & y & \LR{$\mathfrak{y}$} \\
    M &	\LR{$\mathfrak{M}$} & m & \LR{$\mathfrak{m}$} &
    Z &	\LR{$\mathfrak{Z}$} & z & \LR{$\mathfrak{z}$} 
    \end{tabular}
    \end{center}
\olpseganchor{OLP-0642-B005}{end}

\clearpage
\begin{LTR}\latinfont
\nocite{Frege1953,Peter1967}
\bibliographystyle{plainnat}
\bibliography{open-logic}
\end{LTR}

\clearpage\chapter*{د انځورونو سرچينې او د کارولو شرطونه}
\addcontentsline{toc}{chapter}{د انځورونو سرچينې او د کارولو شرطونه}
په دې بڼه کښې د څېرو انځوري فايلونه نشته۔ د څېرو سرليکونه، نسبتونه او د اصلي انځورونو د کارولو شرطونه پکې شته۔

دا انځورونه د اوپن لاجک پروژې له جلا انځوري مجموعې څخه دي۔ د هر انځور بشپړ اصلي انتساب او ځانګړي شرطونه لاندې ساتل شوي، او پښتو بيان ورسره دے۔ د انځورونو اجازه له يوه انځور څخه بل ته توپير لري؛ د کتاب د متن عمومي \LR{\latinfont CC BY 4.0} جواز د ټولو انځورونو لپاره يو شان جواز نۀ دے۔ د OLP-بنسټ غيرتجارتي ښوونيزو موادو اجازه د اړوندو انځورونو په اصلي شرطونو کښې محدوده ده۔ د بل استعمال لپاره د هماغه انځور ځانګړي شرطونه وګورئ۔

د انځورونو پښتو انتسابونه \LR{\latinfont OpenAI Codex — GPT-6.1 Sol, Ultra} په مصنوعي ځيرکتيا ژباړلي دي۔ د قانوني متن اصلي بڼه لاندې ورسره شته۔

\par\medskip\noindent \textbf{ګېورګ کانتور۔} د ګېورګ کانتور انځور اوټو زېت اخيستے دے؛ د هاله--وېټنبرګ د مارټين لوتر پوهنتون د آرشيف په مننه کارول کېږي۔ د آرشيف کره نښه: \LR{\latinfont UAHW Rep.~40-VI, Nr.~3 Bild~102}۔

\textbf{د اصلي سرچينې بشپړ انتساب او شرطونه:}
\begin{LTR}\latinfont\small Portrait of Georg Cantor by Otto Zeth courtesy of the
\href{http://www.archiv.uni-halle.de/}{Universit\"atsarchiv,
Martin-Luther Universit\"at Halle--Wittenberg}. UAHW Rep.~40-VI, Nr.~3
Bild~102.
\end{LTR}

\par\medskip\noindent \textbf{الونزو چرچ۔} د الونزو چرچ د دې انځور نېټه او انځور اخيستونکے نۀ دي معلوم۔ انځور د چرچ د 1924--1995 کاغذونو له مجموعې څخه دے: \LR{\latinfont C0948, Box~60, Folder~3}۔ ساتونکے يې د پرنسټن پوهنتون د کتابتون د نادره کتابونو او ځانګړو مجموعو د ادارې د قلمي نسخو څانګه ده؛ د انځور حقونه د پرنسټن پوهنتون دي۔ اوپن لاجک پروژې په هغو غيرتجارتي موادو کښې د دې انځور د شاملولو اجازه اخيستې چې له OLP څخه جوړېږي۔ د بل هر استعمال لپاره د پرنسټن پوهنتون اجازه ضروري ده۔

\textbf{د اصلي سرچينې بشپړ انتساب او شرطونه:}
\begin{LTR}\latinfont\small Portrait of Alonzo Church, undated, photographer unknown. Alonzo
Church Papers; 1924--1995, (C0948) Box~60,
Folder~3. \href{http://rbsc.princeton.edu/divisions/manuscripts-division}{Manuscripts
Division, Department of Rare Books and Special Collections, Princeton
University Library}. \copyright{} Princeton University. The Open Logic
Project has obtained permission to use this image for inclusion in
non-commercial OLP-derived materials. Permission from Princeton
University is required for any other use.
\end{LTR}

\par\medskip\noindent \textbf{ګېرهارډ ګېنټسن۔} دا د ګېرهارډ ګېنټسن انځور دے چې مېزي ټېنس کوي؛ د اېکهارت مېنټسلر--ټروټ په مننه کارول کېږي۔

\textbf{د اصلي سرچينې بشپړ انتساب او شرطونه:}
\begin{LTR}\latinfont\small Portrait of Gerhard Gentzen playing ping-pong courtesy of Ekhart Mentzler-Trott.
\end{LTR}

\par\medskip\noindent \textbf{کورت ګوډل۔} د کورت ګوډل انځور تقريباً په 1925 کښې اخيستل شوے؛ انځور اخيستونکے نۀ دے معلوم۔ انځور د امريکې د نيوجرسي په پرنسټن کښې د پرمختللو څېړنو د انسټيټيوټ د شېلبي وايټ او ليون لېوي د آرشيف له مرکز څخه دے او د پرنسټن پوهنتون د کتابتون د نادره کتابونو او ځانګړو مجموعو د ادارې د قلمي نسخو په څانګه کښې امانت ساتل کېږي۔ د کورت ګوډل د کاغذونو کره نښه: \LR{\latinfont C0282, Box~14b, \#110000}۔ اوپن لاجک پروژې د انسټيټيوټ د آرشيف له مرکز څخه په OLP-بنسټ غيرتجارتي موادو کښې د انځور د شاملولو اجازه اخيستې۔ د بل هر استعمال لپاره د هماغه مرکز اجازه ضروري ده۔

\textbf{د اصلي سرچينې بشپړ انتساب او شرطونه:}
\begin{LTR}\latinfont\small Portrait of Kurt G\"odel, ca.~1925, photographer unknown. From
the \href{https://library.ias.edu/archives}{Shelby White and Leon Levy
Archives Center, Institute for Advanced Study}, Princeton, NJ, USA, on
deposit at Princeton University
Library, \href{http://rbsc.princeton.edu/divisions/manuscripts-division}{Manuscript
Division, Department of Rare Books and Special Collections}, Kurt
G\"odel Papers, (C0282), Box 14b, \#110000.  The Open Logic Project
has obtained permission from the Institute's Archives Center to use
this image for inclusion in non-commercial OLP-derived materials.
Permission from the Archives Center is required for any other use.
\end{LTR}

\par\medskip\noindent \textbf{اېمي نوېتر۔} د اېمي نوېتر انځور تقريباً د 1922 دے؛ د ګوټېنګن د نيدرزاكسن د دولتي او پوهنتوني کتابتون د قلمي نسخو او نادره چاپونو د څانګې په مننه کارول کېږي۔ د مجموعې کره نښه: \LR{\latinfont Cod.\ Ms.\ D.~Hilbert~754, Bl.~14 Nr.~73}۔ جوېل فولر د اصلي سکېن انځور بيا رغولے دے۔

\textbf{د اصلي سرچينې بشپړ انتساب او شرطونه:}
\begin{LTR}\latinfont\small Portrait of Emmy Noether, ca.~1922, courtesy of
the \href{http://www.sub.uni-goettingen.de/en/contact/departments-a-z/departmental-and-unit-details/abteilunggruppe/handschriften-und-seltene-drucke/}{Abteilung
f\"ur Handschriften und Seltene Drucke, Nieders\"achsische Staats- und
Universit\"atsbibliothek G\"ottingen}, Cod.\ Ms.\ D.~Hilbert~754,
Bl.~14 Nr.~73. Restored from an original scan by Joel Fuller.
\end{LTR}

\par\medskip\noindent \textbf{روژا پېتر۔} د روژا پېتر د دې انځور نېټه او انځور اخيستونکے نۀ دي معلوم؛ د بېلا اندراشفاي په مننه کارول کېږي۔

\textbf{د اصلي سرچينې بشپړ انتساب او شرطونه:}
\begin{LTR}\latinfont\small Portrait of R\'ozsa P\'eter, undated, photographer unknown. Courtesy of
B\'ela Andr\'asfai.
\end{LTR}

\par\medskip\noindent \textbf{جوليا رابنسن۔} د جوليا رابنسن انځور د نيل ډي ريډ په مننه کارول کېږي؛ انځور اخيستونکے نۀ دے معلوم۔ اوپن لاجک پروژې په OLP-بنسټ غيرتجارتي موادو کښې د دې انځور د شاملولو اجازه اخيستې۔ د بل هر استعمال لپاره اجازه ضروري ده۔

\textbf{د اصلي سرچينې بشپړ انتساب او شرطونه:}
\begin{LTR}\latinfont\small Portrait of Julia Robinson, unknown photographer, courtesy of Neil
D.~Reid. The Open Logic Project has obtained permission to use this
image for inclusion in non-commercial OLP-derived materials.
Permission is required for any other use.
\end{LTR}

\par\medskip\noindent \textbf{برټرېنډ رسل۔} د برټرېنډ رسل انځور تقريباً د 1907 دے؛ د مک‌ماسټر پوهنتون د کتابتون د ويليم رېډي د آرشيفونو او څېړنيزو مجموعو د څانګې په مننه کارول کېږي۔ د برټرېنډ رسل د آرشيف کره نښه: \LR{\latinfont Box~2, f.~4}۔

\textbf{د اصلي سرچينې بشپړ انتساب او شرطونه:}
\begin{LTR}\latinfont\small Portrait of Bertrand Russell, ca.~1907, courtesy of the William Ready
Division of Archives and Research Collections, McMaster University
Library. \href{http://www.mcmaster.ca/russdocs/russell.htm}{Bertrand Russell Archives}, Box 2, f.~4.
\end{LTR}

\par\medskip\noindent \textbf{الفرېډ تارسکي۔} دا د الفرېډ تارسکي د 1939 د پاسپورټ انځور دے۔ جوېل فولر د پاسپورټ له سکېن څخه د انځور برخه بېله کړې او بيا رغولې ده۔ اصلي انځور د برکلي د کليفورنيا پوهنتون د بېنکرافټ کتابتون په مننه کارول کېږي؛ د تارسکي د کاغذونو کره نښه: \LR{\latinfont Banc MSS~84/49}۔ اوپن لاجک پروژې په OLP-بنسټ غيرتجارتي موادو کښې د دې انځور د شاملولو اجازه اخيستې۔ د بل هر استعمال لپاره د بېنکرافټ کتابتون اجازه ضروري ده۔ دا بيا رغونه د ساتل شوې سرچينې ده؛ په دې پښتو نسخه کښې انځور بيا نۀ دے بدل شوے۔

\textbf{د اصلي سرچينې بشپړ انتساب او شرطونه:}
\begin{LTR}\latinfont\small Passport photo of Alfred Tarski, 1939. Cropped and restored from a
scan of Tarski's passport by Joel Fuller. Original courtesy of
\href{http://www.lib.berkeley.edu/libraries/bancroft-library}{Bancroft
  Library, University of California, Berkeley}. Alfred Tarski Papers,
Banc MSS~84/49. The Open Logic Project has obtained permission to use
this image for inclusion in non-commercial OLP-derived
materials. Permission from Bancroft Library is required for any other
use.
\end{LTR}

\par\medskip\noindent \textbf{اېلن ټيورينګ۔} د اېلن مېتېسن ټيورينګ انځور اېليټ او فراي د 1951 د مارچ په 29 اخيستے دے۔ کره نښه: \LR{\latinfont NPG x82217}۔ د انځور حقونه د لندن د انځورونو ملي ګېلرۍ دي۔ انځور د \href{http://creativecommons.org/licenses/by-nc-nd/3.0/}{\LR{\latinfont Creative Commons BY-NC-ND 3.0}} د جواز له مخې کارول کېږي: انتساب ضروري دے، تجارتي استعمال او د بدلې شوې بڼې وېش اجازه نۀ لري۔ دلته د سرچينې ساتل شوې بڼه هماغسې ده؛ يوازې په مخ کښې يې د ښودلو اندازه ټاکل شوې۔

\textbf{د اصلي سرچينې بشپړ انتساب او شرطونه:}
\begin{LTR}\latinfont\small Portrait of
\href{http://www.npg.org.uk/collections/search/portrait/mw63680/Alan-Mathison-Turing?}{Alan
Mathison Turing} by Elliott \& Fry, 29 March 1951, NPG
x82217, \copyright{} National Portrait Gallery, London. Used under
a \href{http://creativecommons.org/licenses/by-nc-nd/3.0/}{Creative
Commons BY-NC-ND 3.0 license.}
\end{LTR}

\par\medskip\noindent \textbf{اېرنست زېر مېلو۔} د اېرنست زېر مېلو انځور تقريباً د 1922 دے؛ د ګوټېنګن د نيدرزاكسن د دولتي او پوهنتوني کتابتون د قلمي نسخو او نادره چاپونو د څانګې په مننه کارول کېږي۔ د مجموعې کره نښه: \LR{\latinfont Cod.\ Ms.\ D.~Hilbert~754, Bl.~6 Nr.~25}۔

\textbf{د اصلي سرچينې بشپړ انتساب او شرطونه:}
\begin{LTR}\latinfont\small Portrait of Ernst Zermelo, ca.~1922, courtesy of
the \href{http://www.sub.uni-goettingen.de/en/contact/departments-a-z/departmental-and-unit-details/abteilunggruppe/handschriften-und-seltene-drucke/}{Abteilung
f\"ur Handschriften und Seltene Drucke, Nieders\"achsische Staats- und
Universit\"atsbibliothek G\"ottingen}, Cod.\ Ms.\ D.~Hilbert~754,
Bl.~6 Nr.~25.
\end{LTR}


\end{document}
